% Studio parametrico sul parametro \lambda: le cinque sotto-figure della tesi
% (SDF1, SDF2, SDF3, SVsD, SE) composte in un unico PDF di cinque pagine
% (tesi pp. 63-67). Le legende «Class k=...» sono riscritte in «Class =...»:
% la legenda si compone fuori dalla moda matematica, dove un \xi darebbe
% «Missing $ inserted».
%
% Compilare con LuaLaTeX da questa cartella (la CWD deve essere la cartella di
% questo file):
%   lualatex -shell-escape -interaction=nonstopmode -output-directory=..\..\TikZpdf -jobname=figSSERParametricStudyLambda StudioParametricol.tex
% Così il PDF finisce in Figure/.fnl/Articolo/TikZpdf/ come tutti gli altri.
%
% Perché LuaLaTeX e non pdfLaTeX: ogni pagina è una sola tikzpicture con cinque
% groupplot (quindici grafici) e il soft path di TikZ sopravvive fino a
% \end{tikzpicture}; standalone, a sua volta, trattiene tutte le pagine fino
% alla fine del documento. La sola pagina SDF2 contiene cinquanta layer
% «only marks, mark=*» (circa 23000 cerchi, quattro curve di Bézier ciascuno) e
% da sé supera le 5000000 parole di main_memory, che in pdfTeX è fissata nel
% formato e non è modificabile da riga di comando. LuaTeX alloca invece la
% memoria dinamicamente.
%
% Non è una deroga alla pipeline dell'articolo: main.tex compila già con
% \tikzexternalizedfigure[lualatex] le quattro Infopagina di configurazione di
% riferimento, che hanno la stessa struttura a cinque groupplot di SDF2. Questo
% master resta comunque un documento supplementare isolato, non entra in
% main.tex, e la pipeline pdfLaTeX delle altre figure resta invariata.

\documentclass[
  class=cas-sc,
  tikz
]{standalone}

\def\isodd{even}


%\begingroup For image generation 
  \usepackage[strict]{changepage}

  \usepackage{shellesc}

  \usepackage{etoolbox}
  \newbool{forcetikzrecompile}

  \usepackage{tikz}
  \usetikzlibrary{
    positioning,
    shapes.geometric,
    plotmarks,
    arrows.meta,
    calc,
    decorations.pathreplacing,
    backgrounds,
    patterns.meta,
    fpu, % «floating point unit»
    % external,
  } 

  \usepackage{pgfplots}
  \pgfplotsset{compat=newest}
  \usepgfplotslibrary{
    patchplots,
    groupplots,
    colorbrewer
  }
  \pgfdeclarelayer{background}
  \pgfdeclarelayer{foreground}
  \pgfsetlayers{background,main,foreground}
%\endgroup

%\begingroup For the paper
  \usepackage{microtype}

  \usepackage{amsmath}
  \usepackage{amsthm}
  \usepackage{accents}

  \usepackage{zref-savepos}
  \newcounter{figposcount}

  \usepackage{xstring}

  %\usepackage[numbers]{natbib}
  %\usepackage[authoryear]{natbib}
  % \usepackage[authoryear,longnamesfirst]{natbib}
  \usepackage[numbers,square]{natbib}

  \usepackage{
    graphicx,
    caption,
    subcaption,
    stfloats,
    float
  }

  \usepackage[%
    ruled,%
    vlined,%
    linesnumbered%
  ]{algorithm2e}

  \usepackage{ragged2e}

  \usepackage{enumitem}
  \setlist[enumerate]{%
    label=\arabic*.,
    topsep=0.5em,%
    parsep=0em,%
    itemsep=0.2em,%
    % leftmargin=2em,%
    rightmargin=1.5em,%
  }
  \setlist[itemize]{%
    label=\(\diamond\),
    topsep=0.5em,%
    parsep=0em,%
    itemsep=0.2em,%
    % leftmargin=2em,%
    rightmargin=1.5em,%
  }

  \usepackage{hyperref}
  \usepackage{xurl}

  \usepackage[english]{cleveref}

  \usepackage{siunitx}
  \sisetup{%
    exponent-mode=threshold,%
    exponent-thresholds=-3:3,%
    round-mode=places,% figures
    round-precision=3,%
    round-pad=false,%
    separate-uncertainty=true,%
    print-unity-mantissa=false,%
    tight-spacing=true,%
  }
%\endgroup


\newif\ifcasscclass
\makeatletter
\@ifclassloaded{cas-sc}{\casscclasstrue}{\casscclassfalse}
\makeatother

\newif\ifpageisodd % È usato (nelle macro delle infopagini) dalla meccanica

\ifcasscclass
%\begingroup For image generation
  %\begingroup Tipografia
    \newcommand\textsize[1]{%
      \pgfmathparse{#1*1.2}%
      \fontsize{#1pt}{\pgfmathresult pt}\selectfont%
    }

    \DeclareMathOperator\FRC{\mathcal R}

    % (il rilevamento della classe è in cima al file)

    \providecommand\Ifthispageodd[2]{\checkoddpage\ifoddpage#1\else#2\fi}
  %\endgroup
\fi

  %\begingroup TikZ
    %\begingroup Didascalie
      \newcommand\newcaption[2]{%
        \captionsetup{width=#2} % aboveskip=-.2cm,
        \caption{#1}%
      }
      
      \newcommand\newsubcaption[2][ ]{%
        \refstepcounter{subfigure}%
        (\thesubfigure)#1#2%
      }
    %\endgroup
    
\ifcasscclass
    %\begingroup RedefineTikZbounds
      % La differenza è che colla prima bisogna catturare anche le etichette perché si crea un pdf isolato mentre colla non è importante perché viene tutto compilato sul momento

      %\begingroup Versione valida con esternalizzazione [ma anche senza]
        \newcommand\redefineTikZbounds[2]{
          \coordinate (RealSW) at (current bounding box.south west);
          \coordinate (RealNE) at (current bounding box.north east);
          \coordinate (Center) at ($(#1.west)!0.5!(#2.east)$);
    
          \pgfresetboundingbox 

          \useasboundingbox 
            let \p1=(RealSW),\p2=(RealNE),\p3=(Center),
            \n1={max(\x3-\x1,\x2-\x3)} in
            (\x3-\n1,\y1) rectangle (\x3+\n1,\y2+.25em);

          % \path[use as bounding box] 
          %   let \p1=(RealSW),\p2=(RealNE),\p3=(Center),
          %   \n1={max(\x3-\x1,\x2-\x3)} in
          %   (\x3-\n1,\y1) rectangle (\x3+\n1,\y2+.25em);
          % Lo «.25em» discende dalla traslazione imposta all'etichetta della scala verticale mediante «yScaleLabelShift»
        } % Scatola delimitante simmetrica
      %\endgroup

      %\begingroup Versione valida senza esternalizzazione
        % \newcommand\redefineTikZbounds[2]{
        %   % Salva le coordinate della cima e del fondo della figura corrente (cosí da mantenere i titoli e le didascalie)
        %   \coordinate (Cima) at (current bounding box.north);
        %   \coordinate (Fondo) at (current bounding box.south);
    
        %   \pgfresetboundingbox % Resetta la scatola dei limiti
    
        %   \useasboundingbox % Definisce la nuova scatola come un rettangolo
        %     (#1.west |- Cima)  % con ascissa sinistra data da «Col1 c1r1.west», 
        %     rectangle      % ascissa destra data da «Col2 c1r1.east»
        %     (#2.east |- Fondo); % ordinate date dalla cima e fondo originali
        %     % L'operatore «|-» in «(A |- B)» prende come ascissa quella del punto A mentre come coordinata quella del punto B
        % }
      %\endgroup
    %\endgroup

    %\begingroup TikZBox
      % Parzialmente ispirato da https://gist.github.com/Caramdir/899107

      \makeatletter
      \newsavebox\tikzBufferBox
      \xdef\offsetX{0pt}
      \xdef\offsetY{0pt}
      
      % \tikzbox[#1]{#2}{#3}[#4]{#5}
      % #1 -> Ancora rispetto cui posizionare l'intera scatola (di base è «south west», ossia l'angolo in basso a sinistra)
      % #2 -> Nodo dove posizionare l'origine (per la natura di come LaTeX riceve gli argomenti è buon pratica passare questo circondato da parentesi graffe)
      % #3 -> Nome della scatola limitante definita dall'ambiente «scope»
      % #4 -> Argomenti opzionali da passare all'ambiente «scope»; di base sono vuoti
      % #5 -> Contenuto della scatola circondato dall'ambiente «scope»
      % Es. \tikzbox{{(0,0)}}{figuraTikZ}{blocco TikZ}

      % Macro di apertura che controlla se v'è l'argomento facoltativo dell'ancora
      \def\tikzbox{\pgfutil@ifnextchar[{\tikzbox@anchor}{\tikzbox@anchor[south west]}}
      % La funzione «\pgfutil@ifnextchar» serve per controllare se il prossimo carattere è l'apertura di una paretensi graffa «[»
      
      % Macro intermedia per cogliere gli argomenti obbligatori e controllare quelli successivi facoltativi
      \def\tikzbox@anchor[#1]#2#3{\pgfutil@ifnextchar[{\tikzbox@opt[#1]#2#3}{\tikzbox@opt[#1]#2#3[]}}

      % Macro degli argomenti obbligatori
      \long\def\tikzbox@opt[#1]#2#3[#4]#5{%
        \node[ % Misurazione della scatola in un nodo invisibile
          overlay,%
          opacity=0,%
          inner sep=0pt,%
          outer sep=0pt%
        ] {%
          \sbox{\tikzBufferBox}{%
            \begin{tikzpicture}%
              \begin{scope}[local bounding box=measuringTikZBox]%
                #5%
              \end{scope}%
              \pgfpointanchor{measuringTikZBox}{#1}%
              \pgfgetlastxy{\tempX}{\tempY}% 
              \xdef\offsetX{\tempX}%
              \xdef\offsetY{\tempY}%
            \end{tikzpicture}%
          }%
        };%
        %
        % Disegno reale traslato
        \begin{scope}[%
          local bounding box=#3,%
          shift={#2},%
          xshift=-\offsetX,%
          yshift=-\offsetY,%
          #4%
        ]%
          #5%
        \end{scope}%
      }%
      \makeatother
    %\endgroup

    %\begingroup Chiavi-Valori
      %\begingroup Testo
        % La dimensione del testo di base è 11pt
        \newcommand\titleSize{11}
        \newcommand\subCaptionSize{10}
        \newcommand\ParameterSize{9}
        \newcommand\labelSize{10}
        \newcommand\tickSize{9}
        \newcommand\legendSize{9}
      %\endgroup
    
      %\begingroup Infopagine
        %\begingroup Chiavi condivise
        \newcounter{lpc} % «loopcounter»
        \newcommand\infopageSharedKeys{
          %\begingroup Testo
            \renewcommand\titleSize{10}
            \renewcommand\subCaptionSize8
            \renewcommand\ParameterSize7
            \renewcommand\labelSize9
            \renewcommand\tickSize6
            \renewcommand\legendSize7
          %\endgroup
    
          %\begingroup Dimensioni
            \pgfmathsetmacro\plotWidth7
            \pgfmathsetmacro\plotHeight{.618*\plotWidth}
            \pgfmathsetmacro\halfPlotHeight{\plotHeight/2}
            %
            \pgfmathparse{0.825*\plotWidth}
            \pgfmathsetmacro\plotWidth\pgfmathresult
            \pgfmathsetmacro\halfWidth{\plotWidth/2}
            \pgfmathsetmacro\plotLastWidth{1.25*\plotWidth}
    
            \pgfmathsetmacro\xPlotSep{1}
            \pgfmathsetmacro\yPlotSep{1.25}
            \pgfmathsetmacro\yGroupSep{2.75}
          %\endgroup
    
          %\begingroup Legende
            \newcommand\placeLegend[5]{%
              /pgfplots/plotLegend,
              at={([shift={(##4,##5)}]axis description cs:##1,##2)},
              anchor=##3,
            }
            %
            \pgfplotsset{%
              plotLegend/.style={%
                font=\textsize{\legendSize},
                /pgfplots/legend cell align=left,
                % align=left,
                % draw=white!15!black,
                draw=none, % Bordo invisibile
                fill=white,
                fill opacity=0.5,
                text opacity=1,
                % text width=2cm,
              },
              plotLegendNW/.style={legend style={\placeLegend{0}{1}{north west}{0em}{0em}}},
              plotLegendNE/.style={legend style={\placeLegend{1}{1}{north east}{0em}{0em}}},
              plotLegendSW/.style={legend style={\placeLegend{0}{0}{south west}{0em}{0em}}},
              plotLegendSE/.style={legend style={\placeLegend{1}{0}{south east}{0em}{0em}}},
              %
              % plotLegendW/.style={legend style={plotLegend,at={(0,.5)},anchor=west}},
              %
              % plotLegendNW/.style={legend style={plotLegend,legend pos=north west}},
              % plotLegendNE/.style={legend style={plotLegend,legend pos=north east}},
              % plotLegendSW/.style={legend style={plotLegend,legend pos=south west}},
              % plotLegendSE/.style={legend style={plotLegend,legend pos=south east}},
              %
              % plotLegendNWcloser/.style={legend style={plotLegend,at={(0,1)},anchor=north west}},
              % plotLegendNEcloser/.style={legend style={plotLegend,at={(1,1)},anchor=north east}},
              % plotLegendSWcloser/.style={legend style={plotLegend,at={(0,0)},anchor=south west}},
              % plotLegendSEcloser/.style={legend style={plotLegend,at={(0,1)},anchor=south east}},
            }
            %
            \pgfmathsetmacro\lineLegendLength{1}
            \pgfmathsetmacro\areaLegendHeigth{.25}
            \pgfplotsset{
              legend image code/.code={
                \draw [mark repeat=2,mark phase=2,##1]
                  plot coordinates {
                    (0em,0em)
                    (.5*\lineLegendLength em,0em)
                    (\lineLegendLength em,0em)
                  };
              },
              /pgfplots/area legend/.style={
                legend image code/.code={
                  \draw [##1,fill,draw=none] 
                    (0em,-\areaLegendHeigth em) 
                    rectangle 
                    (\lineLegendLength em,\areaLegendHeigth em);
                },
              },
            }
            %
            \newcommand\ExactHistogramLegend{Exact histogram}
            \newcommand\ApproxHistogramLegend{Approx. histogram}
            \newcommand\RealHistogramLegend{Real histogram}
            %
            \newcommand\ExactBLNFitLegend{Exact BLN fit}
            \newcommand\ApproxBLNFitLegend{Approx BLN fit}
            \newcommand\RealBLNFitLegend{Real BLN fit}
            %
            \newcommand\ExactEmpiricCCDFLegend{Exact empiric CCDF}
            \newcommand\ApproxEmpiricCCDFLegend{Approx. empiric CCDF}
            \newcommand\RealEmpiricCCDFLegend{Real empiric CCDF}
            %
            \newcommand\ExactParetoFitLegend{Exact Pareto fit}
            \newcommand\ApproxParetoFitLegend{Approx. Pareto fit}
            \newcommand\RealParetoFitLegend{Real Pareto fit}
            %
            \newcommand\ExactScatterLegend{Exact scatter}
            \newcommand\ApproxScatterLegend{Approx. scatter}
            \newcommand\RealScatterLegend{Real scatter}
            %
            \newcommand\ClasskILegend{Class k=283}
            \newcommand\ClasskIILegend{Class k=83}
            \newcommand\ClasskIIILegend{Class k=43}
            \newcommand\ClasskIVLegend{Class k=10}
          %\endgroup
    
          %\begingroup Etichette vuote
            \pgfplotsset{
              yLabelEmpty/.style={ylabel={}},
              yTickLabelsEmpty/.style={yticklabels={}},
              yTickScaleLabelEmpty/.style={ytick scale label code/.code={}},
              xTickScaleLabelEmpty/.style={xtick scale label code/.code={}},
              emptyLeged/.style={
                legend style={draw=none,fill=none,nodes={draw=none}},
                legend image post style={opacity=0},
              },
            }
          %\endgroup
    
          %\begingroup Estetica
            \newcommand\smallminus{\scalebox{0.5}[1]{$-$}\hspace{-.11em}}
            \pgfplotsset{
              /pgfplots/log number format basis/.code 2 args={%
                \pgfmathsetmacro\exponent{##2}%
                \pgfmathparse{sign(\exponent)}%
                \ifnum\pgfmathresult<0%
                  \pgfmathparse{abs(\exponent)}%
                  $##1^{\smash{\smallminus\pgfmathprintnumber{\pgfmathresult}}}$%
                \else%
                  $##1^{\smash{\pgfmathprintnumber{\exponent}}}$%
                \fi%
              },
              ytick scale label code/.code={%
                \pgfmathsetmacro\exponent{##1}%
                \pgfmathparse{sign(\exponent)}%
                \ifnum\pgfmathresult<0
                  \pgfmathparse{abs(\exponent)}%
                  $\cdot 10^{\smash{\smallminus\pgfmathprintnumber{\pgfmathresult}}}$%
                \else
                  $\cdot 10^{\smash{\pgfmathprintnumber{\exponent}}}$%
                \fi
              },
            }
          %\endgroup
    
          %\begingroup Colori
            \definecolor{blueish}{rgb}{0.00000,0.44700,0.74100}%
            \definecolor{reddish}{rgb}{0.83529,0.36863,0.00000}%
          %\endgroup
    
          %\begingroup Commenti
            %\begingroup Impostazioni
            \tikzset{
              title/.style={
                align=center,
                anchor=south,
                font=\bfseries\textsize{\titleSize},
              },
              fittedParameters/.style={
                align=center,
                font=\textsize{\ParameterSize},
              },
              varyingParameter/.style={
                align=center,
                % rotate=90,
                font=\textsize{\subCaptionSize},
              },
              subCaption/.style={
                % text width=\plotWidth cm,
                align=center,
                font=\textsize{\subCaptionSize},
              },
            }
            \newcommand\codetocolor[1]{\ifcase##1blueish\or reddish\else black\fi}
            %\endgroup
    
            %\begingroup Parametri Sardegna
            \newcommand\betaReal{\num{1.273}}
            \newcommand\zetaReal{\num{0.929}}
            \newcommand\mIReal{\num{3.227}}
            \newcommand\sIReal{\num{0.408}}
            \newcommand\mIIReal{\num{4.099}}
            \newcommand\sIIReal{\num{0.556}}
            %\endgroup
    
            %\begingroup Parametri Pareto
            \newcommand\codetoParetovalue[1]{
              \ifcase##1\betaExact\or\betaAprx\else\betaReal\fi
            }
            \newcommand\writeParetoParameters[7][center]{
              \node[fittedParameters,anchor=##1] at ([shift={(##5,##6)}]##4) {%
                \(\begin{aligned}
                  \color{\codetocolor{##2}}\beta
                  &\color{\codetocolor{##2}}=\codetoParetovalue{##2}\\[##7]
                  \color{\codetocolor{##3}}\beta
                  &\color{\codetocolor{##3}}=\codetoParetovalue{##3}
                \end{aligned}\)
              };
            }
            %\endgroup

            %\begingroup % Vecchia versione in cui si scrivevano esplicitamente tutte le 9 combinazioni possibili
            % \newcommand\writeParameters[6]{%
            %   \node[parameter] at ([shift={(##4,##5)}]##1) {%
            %     \ifcase##2% Caso 0
            %       \ifcase##3% Caso 0
            %         \(\begin{aligned}
            %           \color{blueish}\beta&\color{blueish}=\betaExact\\[##6]
            %           \color{blueish}\beta&\color{blueish}=\betaExact
            %         \end{aligned}\)
            %       \or% Caso 1
            %         \(\begin{aligned}
            %           \color{blueish}\beta&\color{blueish}=\betaExact\\[##6]
            %           \color{reddish}\beta&\color{reddish}=\betaAprx% Caso 1
            %         \end{aligned}\)
            %       \else% Caso 2
            %         \(\begin{aligned}
            %           \color{blueish}\beta&\color{blueish}=\betaExact\\[##6]
            %           \color{black}\beta&\color{black}=\betaReal% Caso 2
            %         \end{aligned}\)
            %       \fi
            %     \or% Caso 1
            %       \ifcase##3% Caso 0
            %         \(\begin{aligned}
            %           \color{reddish}\beta&\color{reddish}=\betaAprx\\[##6]
            %           \color{blueish}\beta&\color{blueish}=\betaExact
            %         \end{aligned}\)
            %       \or% Caso 1
            %         \(\begin{aligned}
            %           \color{reddish}\beta&\color{reddish}=\betaAprx\\[##6]
            %           \color{reddish}\beta&\color{reddish}=\betaAprx
            %         \end{aligned}\)
            %       \else% Caso 2
            %         \(\begin{aligned}
            %           \color{reddish}\beta&\color{reddish}=\betaAprx\\[##6]
            %           \color{black}\beta&\color{black}=\betaReal
            %         \end{aligned}\)
            %       \fi
            %     \else% Caso 2
            %       \ifcase##3% Caso 0
            %         \(\begin{aligned}
            %           \color{black}\beta&\color{black}=\betaReal\\[##6]
            %           \color{blueish}\beta&\color{blueish}=\betaExact
            %         \end{aligned}\)
            %       \or% Caso 1
            %         \(\begin{aligned}
            %           \color{black}\beta&\color{black}=\betaReal\\[##6]
            %           \color{reddish}\beta&\color{reddish}=\betaAprx
            %         \end{aligned}\)
            %       \else% Caso 2
            %         \(\begin{aligned}
            %           \color{black}\beta&\color{black}=\betaReal\\[##6]
            %           \color{black}\beta&\color{black}=\betaReal
            %         \end{aligned}\)
            %       \fi
            %     \fi
            %   };
            % }
            %\endgroup
    
            %\begingroup Parametri bilognormale
            \newcommand\codetoparameter[1]{
              \ifcase##1\zeta  % Caso 0
              \or\mu_1     % Caso 1
              \or\sigma_1   % Caso 2
              \or\mu_2     % Caso 3
              \else\sigma_2\fi % Caso 4
            }
            \newcommand\codetoBiLognormalvalue[2]{
              \ifcase##1 % Caso 0
                \ifcase##2\zetaExact % Caso 0
                \or\mIExact    % Caso 1
                \or\sIExact    % Caso 2
                \or\mIIExact    % Caso 3
                \else\sIIExact\fi % Caso 4
              \or % Caso 1
                \ifcase##2\zetaAprx % Caso 0
                \or\mIAprx     % Caso 1
                \or\sIAprx     % Caso 2
                \or\mIIAprx    % Caso 3
                \else\sIIAprx\fi  % Caso 4
              \else % Caso 2
                \ifcase##2\zetaReal % Caso 0
                \or\mIReal     % Caso 1
                \or\sIReal     % Caso 2
                \or\mIIReal    % Caso 3
                \else\sIIReal\fi  % Caso 4
              \fi
            }
            \newcommand\matrixline[2]{\codetoparameter{##2}&=\codetoBiLognormalvalue{##1}{##2}}
            \def\lineloop##1##2{%
              \matrixline{##1}{\number\value{lpc}}\stepcounter{lpc}%
              \ifnum\number\value{lpc}<5\\[##2]\lineloop{##1}{##2}\fi%
            } % https://tex.stackexchange.com/questions/491358/tikz-foreach-to-work-with-the-aligned-environment
            \newcommand\writeBiLognormalParameters[6][center]{
              \node[fittedParameters,anchor=##1] at ([shift={(##4,##5)}]##3) {%
                \setcounter{lpc}{0}
                \(\color{\codetocolor{##2}}
                \begin{aligned}
                  \lineloop{##2}{##6}
                \end{aligned}\)
              };
            }
            %\endgroup
    
            %\begingroup Altre etichette
            \newcommand\writeSizeEvolutionDescriptions[1][5]{%
              \foreach \col/\clr/\txt in {%
                1/blueish/Esatto,
                2/reddish/Approssimato
              }{
                \node[
                  fittedParameters,
                  anchor=south,
                  yshift=.1em,
                ] at (Row##1 c\col r1.south) {\color{\clr}\txt};
              }
            }
    
            \newcommand\writeVaryingParameter[5]{%
              \foreach \r in {1,...,5}{
                \node[
                  varyingParameter,
                  anchor=##3,
                ] at ([shift={(##4,##5)}]Row\r\space c##1r1.##2)
                {\(\varyingParameterSymbol=\varyingParameterList\r\)};
              }
            }
            %\endgroup
          %\endgroup
    
          %\begingroup Sottodidascalie
            \newcommand\writeSubCaption[5]{%
              \node[subCaption,anchor=##3]
              at ([shift={(##4,##5)}]##1.##2)
              {\newsubcaption[]{}};
            }
            \newcommand\writeSubCaptions[5][3]{%
              \foreach \r in {1,...,5}{
                \foreach \c in {1,...,##1}{
                  \writeSubCaption{Row\r\space c\c r1}{##2}{##3}{##4}{##5}
                }
              }
            }
          %\endgroup
        }
        %\endgroup
    
        %\begingroup Chiavi del primo tipo d'infopagina
        \newcommand\infopageKeysI{
          \infopageSharedKeys
    
          %\begingroup Gruppi
            \ifpageisodd
              \pgfplotsset{/pgfplots/group/yDescPos/.style={y descriptions at=edge right}}
            \else
              \pgfplotsset{/pgfplots/group/yDescPos/.style={y descriptions at=edge left}}
            \fi
            \pgfplotsset{
              /pgfplots/group/plotGroup/.style={
                horizontal sep=\xPlotSep em,
                vertical sep=\yPlotSep em,
                yDescPos,
                group size=3 by 1,
              },
              /pgfplots/group/plotLastGroup/.style={
                plotGroup,
                x descriptions at=edge bottom,
              },
            }
          %\endgroup
    
          %\begingroup Stili
            %\begingroup Selezione «recto»/«verso»
            \ifpageisodd
              \pgfplotsset{%
                yLineAxis/.style={axis y line*=right},
                /tikz/yLabelShift/.style={%
                  yshift=-1.25em,
                  at={(axis description cs:1,0.5)},
                  anchor=north,
                  % rotate=180,
                },
                /tikz/yScaleLabelShift/.style={%
                  yshift=.25em,
                  at={(1,1)},
                  anchor=south east,
                },
                /tikz/yTickLabelShift/.style={xshift=-0.15em},
              }
            \else
              \pgfplotsset{%
                yLineAxis/.style={axis y line*=left},
                /tikz/yLabelShift/.style={%
                  yshift=1.25em,
                  at={(axis description cs:0,0.5)},
                  anchor=south,
                  % rotate=180,
                },
                /tikz/yScaleLabelShift/.style={%
                  yshift=.25em,
                  at={(0,1)},
                  anchor=south west,
                },
                /tikz/yTickLabelShift/.style={xshift=+0.15em},
              }
            \fi
            %\endgroup
            %
            %\begingroup Grafici
            \pgfplotsset{
              plotRow/.style={
                scale only axis,
                width=\plotWidth cm,height=\plotHeight cm,
                xlabel style={font=\color{white!15!black}\textsize{\labelSize}},
                ylabel style={font=\color{white!15!black}\textsize{\labelSize},yLabelShift},
                xticklabel style={font=\textsize{\tickSize}},
                yticklabel style={font=\textsize{\tickSize},yTickLabelShift},
                y tick scale label style={yScaleLabelShift},
                % axis x line*=bottom,yLineAxis,
                enlarge x limits=0.1,enlarge y limits=0.1,
                xmajorgrids,ymajorgrids,
                xminorticks=true,yminorticks=true,
                scale ticks above exponent=2,
                axis background/.style={fill=white},
                grid style={dashed}
              },
              % xlabel123/.style={xlabel={\(h\)}},
              xlabel3/.style={xlabel={\(h\)}},
              ylabel23/.style={ylabel={\(\FRC(h)\)}},
              row1Keys/.style={
                % ylabel={\(\underhat{\bar l}_N(h)\)},
                ylabel={\(\bar l_N(h)\)},
                % enlarge x limits={lower,value=.05},
                enlarge x limits=.05,
                enlarge y limits={upper,value=.05},
                xmode=linear,
                log origin=infty,
              },
              row2KeysWithLabels/.style={
                ylabel23,
                xmode=linear,ymode=log
              },
              row2Keys/.style={
                ylabel23,
                xmode=linear,ymode=log,
                xticklabels={},
              },
              row3Keys/.style={
                xlabel3,ylabel23,
                xmode=linear,ymode=log
              },
              row4Keys/.style={
                xlabel={\(\xi\)},ylabel={\(s(\xi)\)},
                xmode=log,ymode=log
              },
              row5KeysLin/.style={
                width=\plotLastWidth cm,
                xlabel={\(t\)},ylabel={\(s(t)\)},
                enlarge x limits=0,
                enlarge y limits=0.05,
                x tick scale label style={
                  yshift=-1em,
                  at={(.95,0)},
                  anchor=north,
                },
              },
              row5Keys/.style={
                row5KeysLin,
                ymode=log,
              },
            }
            %\endgroup
          %\endgroup
        }
        %\endgroup
    
        %\begingroup Chiavi del secondo tipo d'infopagina
        \newcommand\infopageKeysII{
          \infopageSharedKeys
    
          %\begingroup Gruppi
            \ifpageisodd
              \pgfplotsset{/pgfplots/group/yDescPos/.style={y descriptions at=edge right}}
            \else
              \pgfplotsset{/pgfplots/group/yDescPos/.style={y descriptions at=edge left}}
            \fi
            \pgfplotsset{
              /pgfplots/group/plotGroup/.style={
                horizontal sep=\xPlotSep em,
                vertical sep=\yPlotSep em,
                y descriptions at=all,
                group size=3 by 1,
              },
              /pgfplots/group/plotLastGroup/.style={
                plotGroup,
                x descriptions at=edge bottom,
              },
            }
          %\endgroup
    
          %\begingroup Stili
            %\begingroup Selezione «recto»/«verso»
            \ifpageisodd
              \pgfplotsset{
                yLabel/.style={
                  ylabel style={
                    font=\color{white!15!black}\textsize{\labelSize},
                    yshift=-1.25em,
                    at={(axis description cs:1,0.5)},
                    anchor=north,
                    % rotate=180,
                  },
                },
                yLabelInv/.style={yLabelEmpty},
                yScaleLabelShift/.style={
                  y tick scale label style={
                    yshift=.25em,
                    % xshift=1em,
                    at={(1,1)},
                    anchor=south east,
                  },
                },
                yScaleLabelShiftInv/.style={
                  y tick scale label style={
                    yshift=.25em,
                    % xshift=-1em,
                    at={(0,1)},
                    anchor=south west,
                  },
                },
                % yLineAxis/.style={axis y line*=right},
              }
            \else
              \pgfplotsset{
                yLabel/.style={yLabelEmpty},
                yLabelInv/.style={
                  ylabel style={
                    font=\color{white!15!black}\textsize{\labelSize},
                    yshift=1.25em,
                    at={(axis description cs:0,0.5)},
                    anchor=south,
                    % rotate=180,
                  },
                },
                yScaleLabelShift/.style={
                  y tick scale label style={
                    yshift=.25em,
                    % xshift=1em,
                    at={(1,1)},
                    anchor=south east,
                  },
                },
                yScaleLabelShiftInv/.style={
                  y tick scale label style={
                    yshift=.25em,
                    % xshift=-1em,
                    at={(0,1)},
                    anchor=south west,
                  },
                },
                % yLineAxis/.style={axis y line*=left},
              }
            \fi
            %\endgroup
            %
            %\begingroup Grafici
            \pgfplotsset{
              plotRow/.style={
                scale only axis,
                width=\plotWidth cm,height=\plotHeight cm,
                xlabel style={font=\color{white!15!black}\textsize{\labelSize}},
                % ylabel style={font=\color{white!15!black}\textsize{\labelSize},yLabelShift},
                xticklabel style={font=\textsize{\tickSize}},
                % axis x line*=bottom,yLineAxis,
                enlarge x limits=0.1,enlarge y limits=0.1,
                xmajorgrids,ymajorgrids,
                xminorticks=true,yminorticks=true,
                scale ticks above exponent=2,
                axis background/.style={fill=white},
                grid style={dashed},
              },
              plotLeftCol/.style={
                yLabelInv,
                yticklabel style={font=\textsize{\tickSize},xshift=0.15em},
                yticklabel pos=left,
                yScaleLabelShiftInv,
              },
              plotMiddleCol/.style={
                yLabelEmpty,
                yTickLabelsEmpty,
                yTickScaleLabelEmpty,
              },
              plotRightCol/.style={
                yLabel,
                yticklabel style={font=\textsize{\tickSize},xshift=-0.15em},
                yticklabel pos=right,
                yScaleLabelShift,
              },
              xlabel3/.style={xlabel={\(h\)}},
              ylabel23/.style={ylabel={\(\FRC(h)\)}},
              row1Keys/.style={
                ylabel={\(\bar l_N(h)\)},
                enlarge x limits=.05,
                enlarge y limits={upper,value=.05},
                xmode=linear,
                log origin=infty,
              },
              row2Keys/.style={
                ylabel23,
                xmode=linear,ymode=log,
                % xticklabels={},
              },
              row3Keys/.style={
                xlabel3,ylabel23,
                xmode=linear,ymode=log,
              },
              row4KeysLin/.style={
                xlabel={\(\xi\)},ylabel={\(s(\xi)\)},
                xmode=log,
              },
              row4Keys/.style={row4KeysLin,ymode=log},
              row5KeysLin/.style={
                width=\plotLastWidth cm,
                xlabel={\(t\)},ylabel={\(s(t)\)},
                enlarge x limits=0,
                enlarge y limits=0.05,
                x tick scale label style={
                  yshift=-1em,
                  at={(.95,0)},
                  anchor=north,
                },
              },
              row5Keys/.style={
                row5KeysLin,
                ymode=log,
              },
            }
            %\endgroup
          %\endgroup
        }
        %\endgroup
      %\endgroup
    %\endgroup

    %\begingroup Diagrammi dei nodi
    \newcommand\nodeDiagramSettings{
      %\begingroup Dimensioni
        \pgfmathsetmacro\xPlotSep{10}
        \pgfmathsetmacro\yPlotSep{4}
        \pgfmathsetmacro\plotScale{.8}

        \pgfmathsetmacro\titleSize{11}
        \pgfmathsetmacro\commentSize{10}

        \pgfmathsetmacro\nodeSize{20pt} % Diametro del nodo
        \pgfmathsetmacro\starSize{2.25*\nodeSize}
        \pgfmathsetmacro\numberSize{.53*\nodeSize}
      %\endgroup

      %\begingroup Colori
        % \definecolor{cRed}{HTML}{D68490}
        % \definecolor{cPurple}{HTML}{7E6EAC}
        % \definecolor{cTeal}{HTML}{87BDB1}
        % \definecolor{cOrange}{HTML}{F2D0A4}
        % \definecolor{cEdgeLight}{HTML}{EEE8D5} 
        % \definecolor{cEdgeDark}{HTML}{998888}
        % \definecolor{cStar}{HTML}{FDF6E3}
        %
        \definecolor{color1}{HTML}{800000} % Deep Red
        \definecolor{color2}{HTML}{4682B4} % Steel Blue
        \definecolor{color3}{HTML}{6B8E23} % Olive Green
        \definecolor{color4}{HTML}{708090} % Slate Gray
        \definecolor{color5}{HTML}{BC8F8F} % Rosy Brown
        \definecolor{color6}{HTML}{008B8B} % Dark Cyan
        \definecolor{color7}{HTML}{B8860B} % Dark Goldenrod
      %\endgroup

      %\begingroup Testo
        \tikzset{
          titleStyle/.style={
            font=\textsize{\titleSize},
            align=center,
            anchor=south,
          },
          commentStyle/.style={font=\textsize{\commentSize}},
          arrowCommentStyle/.style={
            midway,
            commentStyle,
            left,
            sloped,
            below
          },
          numberStyle/.style={font=\textsize{\numberSize}}
        }
      %\endgroup

      %\begingroup Frecce
        \tikzset{
          transformingArrow/.style={
            % =Latex,
            ->,
            >=Stealth,
            dashed,
            line width=.8pt,
            black!90,
            shorten >=5pt,
            shorten <=10pt
          },
          indirectArrow/.style={
            draw=black!95,
            line width=.5pt,
          },
          nodeArrowExact/.style={
            indirectArrow,
            % ->,
            % -{Stealth[scale=1.2]},
            -{Triangle[scale=.75]},
            % draw=cEdgeLight,
            % draw=red!50,
            % shorten >=2pt,
            % shorten <=2pt
          },
          nodeArrowApprox/.style={
            indirectArrow,
            -,
            dashed,
          },
          transparentNodeArrow/.style={
            opacity=.4,
          },
          arrowAxis/.style={
            ->,
            line width=.3pt,
            >=Latex,
          },
        }
      %\endgroup

      %\begingroup Nodi
        \tikzset{
          nodeStyle/.style={
            circle,
            minimum size=\nodeSize,
            % draw=black!80,
            draw=black,
            line width=.5pt,
            inner sep=0pt,
          },
          starStyle/.style={
            star,
            star points=5,
            star point ratio=2,
            % fill=cStar,
            fill=black!50,
            fill opacity=.25,
            minimum size=\starSize,
            draw=none
          },
        }
      %\endgroup
    }
    %\endgroup

    %\begingroup GNU-Plot
      % \IfFileExists{\sourceauxpath.aux}{
      %   \externaldocument{\sourceauxpath}
      % }{\relax}
      
      \pgfkeys{/pgf/plot/gnuplot call={cd /d "\sourceoutdir" && gnuplot}}
    %\endgroup
  %\endgroup
\fi
%\endgroup

%\begingroup For the paper
  \def\tsc#1{\csdef{#1}{\textsc{\lowercase{#1}}\xspace}}
  \tsc{WGM}
  \tsc{QE}
  
  \newcommand\etal{et al.}
  \newcommand\KTMAS{\hyperlink{acrKineticTheoryMultiAgentSystems}{KTMAS}}
  
  % \newcommand\MathSpacing[3]{
  %   \thinmuskip = #1mu\relax
  %   \medmuskip = #2mu\relax
  %   \thickmuskip = #3mu\relax
  % }
  % \MathSpacing123
  
  \newcommand\RPlus{\mathbb R_+}
  \newcommand\NPlus{\mathbb N_+}
  \DeclareMathAlphabet\pazocal{OMS}{zplm}{m}{n}
  
  \DeclareMathOperator\Lognormal{Lognormal}
  \DeclareMathOperator\BiLognormal{BiLognormal}
  \DeclareMathOperator\Bernoulli{Bernoulli}
  \DeclareMathOperator\Exponential{Exp}
  \DeclareMathOperator\Student{Student}
  \DeclareMathOperator\Variance{var}
  \DeclareMathOperator\IC{IC}
  \DeclareMathOperator\Histogram{\mathcal B}
  \DeclareMathOperator\CCDF{\mathcal R}
    
  \newcommand\ThreeQuarters{{3\hspace{-0.15em}/\hspace{-0.15em}4}}
  \newcommand\OneQuarter{{1\hspace{-0.15em}/\hspace{-0.15em}4}}
  
  \newcommand\de[1][]{\mathrm{d}#1}
  \newcommand\derS[3][d]{\frac{\mathrm{#1}#2}{\mathrm{#1}#3}}
  \newcommand\derP[2]{\frac{\partial#1}{\partial#2}}
  \newcommand\underhat[1]{\underaccent{\hat}{#1}}
  \renewcommand\undertilde[1]{\underaccent{\tilde}{#1}}
  
  \makeatletter
  \newcommand\eqtag{\@ifstar{\eqtag@star}{\eqtag@nostar}}
  
  \newcounter{eqtag}
  \newcommand\eqtag@nostar[2]{%
    \begin{equation*}%
      \refstepcounter{eqtag}%
      % \protected@edef\@currentlabeltype{eqtag}%
      \protected@edef\@currentlabel{#1}%
      #2%
      \tag*{#1}%
    \end{equation*}%
  }
  \newcommand\eqtag@star[2]{%
    \begin{equation*}%
      \refstepcounter{eqtag}%
      % \protected@edef\@currentlabeltype{eqtag}%
      \protected@edef\@currentlabel{#1}%
      \let\label\ltx@label%
      #2%
      % \notag%
    \end{equation*}%
  }
  \makeatother
  
  \newcommand\crefsubfigInt[4][-]{Figures~\ref{#2}#3#1\ref{#2}#4}
  \newcommand\crefsubfigSep[4]{Figures~\ref{#1}#2\foreach \l in#3{,~\ref{#1}\l}~and~\ref{#1}#4}
  
  \crefname{section}{§}{§§}
  \crefrangeformat{section}{§§~#3#1#4--#5#2#6}
  \crefmultiformat{section}{§§~#2#1#3}{ and~#2#1#3}{, #2#1#3}{ and~#2#1#3}
  \crefrangemultiformat{section}{§§~#3#1#4--#5#2#6}{ and~#3#1#4--#5#2#6}{, #3#1#4--#5#2#6}{ and~#3#1#4--#5#2#6}
  
  \crefname{equation}{}{}
  \crefrangeformat{equation}{(#3#1#4--#5#2#6)}
  \crefmultiformat{equation}{(#2#1#3}{,~#2#1#3)}{,~#2#1#3}{,~#2#1#3)}
  \crefrangemultiformat{equation}{(#3#1#4--#5#2#6}{,~#3#1#4--#5#2#6)}{,~#3#1#4--#5#2#6}{,~#3#1#4--#5#2#6)}
  
  \crefname{figure}{Figure}{Figures}
  \crefrangeformat{figure}{Figures~#3#1#4--#5#2#6}
  \crefmultiformat{figure}{Figures~#2#1#3}{ and~#2#1#3}{, #2#1#3}{ and~#2#1#3}
  \crefrangemultiformat{figure}{Figures~#3#1#4--#5#2#6}{ and~#3#1#4--#5#2#6}{, #3#1#4--#5#2#6}{ and~#3#1#4--#5#2#6}
  
  \crefname{table}{Table}{Tables}
  \crefrangeformat{table}{Tables~#3#1#4--#5#2#6}
  \crefmultiformat{table}{Tables~#2#1#3}{ and~#2#1#3}{, #2#1#3}{ and~#2#1#3}
  \crefrangemultiformat{table}{Tables~#3#1#4--#5#2#6}{ and~#3#1#4--#5#2#6}{, #3#1#4--#5#2#6}{ and~#3#1#4--#5#2#6}
  
  \crefname{algorithm}{Algorithm}{Algorithms}
  \crefrangeformat{algorithm}{Algorithms~#3#1#4--#5#2#6}
  \crefmultiformat{algorithm}{Algorithms~#2#1#3}{ and~#2#1#3}{, #2#1#3}{ and~#2#1#3}
  \crefrangemultiformat{algorithm}{Algorithms~#3#1#4--#5#2#6}{ and~#3#1#4--#5#2#6}{, #3#1#4--#5#2#6}{ and~#3#1#4--#5#2#6}
  
  %\newtheorem{theorem}{Theorem}
  %\newtheorem{lemma}[theorem]{Lemma}
  \ifcasscclass
    % «cas-sc» fornisce \newdefinition
    \newdefinition{dfn}{Definition}
    \newdefinition{rmk}{Remark}
    \newdefinition{hyp}{Hypotesis}
  \else
    % Con «article» (Articolo) si usano gli ambienti del kernel, con
    % contatore condiviso e «Assumption» come titolo di «hyp»
    \newtheorem{dfn}{Definition}
    \newtheorem{rmk}[dfn]{Remark}
    \newtheorem{hyp}[dfn]{Assumption}
  \fi
  %\newproof{pf}{Proof}
  %\newproof{pot}{Proof of Theorem \ref{thm}}
  
  \crefname{dfn}{Definition}{Definitions}
  \crefrangeformat{dfn}{Definitions~#3#1#4--#5#2#6}
  \crefmultiformat{dfn}{Definitions~#2#1#3}{ and~#2#1#3}{, #2#1#3}{ and~#2#1#3}
  \crefrangemultiformat{dfn}{Definitions~#3#1#4--#5#2#6}{ and~#3#1#4--#5#2#6}{, #3#1#4--#5#2#6}{ and~#3#1#4--#5#2#6}
  \crefname{hyp}{Hypothesis}{Hypotheses}
  
  \crefrangeformat{hyp}{Hypotheses~#3#1#4--#5#2#6}
  \crefmultiformat{hyp}{Hypotheses~#2#1#3}{ and~#2#1#3}{, #2#1#3}{ and~#2#1#3}
  \crefrangemultiformat{hyp}{Hypotheses~#3#1#4--#5#2#6}{ and~#3#1#4--#5#2#6}{, #3#1#4--#5#2#6}{ and~#3#1#4--#5#2#6}
  
  \crefname{rmk}{Remark}{Remarks}
  \crefrangeformat{rmk}{Remarks~#3#1#4--#5#2#6}
  \crefmultiformat{rmk}{Remarks~#2#1#3}{ and~#2#1#3}{, #2#1#3}{ and~#2#1#3}
  \crefrangemultiformat{rmk}{Remarks~#3#1#4--#5#2#6}{ and~#3#1#4--#5#2#6}{, #3#1#4--#5#2#6}{ and~#3#1#4--#5#2#6}
  
  % \newcommand\tikzfigure[5][0em]{%
  %   \begin{figure}[%
  %     pos=#2,%
  %     width=#4,%
  %     align=\centering%
  %   ]%
  %     \centering%
  %     \vspace*{#1}%
  %     \IfSubStr{#2}{p}{%
  %       \leavevmode\stepcounter{figposcount}%
  %       \zsaveposx{figpos-\thefigposcount}%
  %       \makebox[0pt][l]{%
  %         \hspace*{-\dimexpr\zposx{figpos-\thefigposcount}sp\relax}%
  %         \hspace*{\dimexpr0.5\paperwidth\relax}%
  %         \makebox[0pt][c]{%
  %           \scalebox{0.86}{\includegraphics{figs/#5.pdf}}%
  %         }%
  %       }
  %     }{%
  %       \includegraphics{figs/#5.pdf}%
  %     }%
  %     \caption{%
  %       \hyphenpenalty=0%
  %       \exhyphenpenalty=0%
  %       #3%
  %     }%
  %     \protect\label{#5}%
  %   \end{figure}%
  % }
%\endgroup

%\begingroup Versione per la classe «cas-sc»
  % 
  % %\begingroup For image generation
  %   %\begingroup Tipografia
  %     \newcommand\textsize[1]{%
  %       \pgfmathparse{#1*1.2}%
  %       \fontsize{#1pt}{\pgfmathresult pt}\selectfont%
  %     }
  % 
  %     \DeclareMathOperator\FRC{\mathcal R}
  % 
  %     \newif\ifcasscclass
  %     \makeatletter
  %     \@ifclassloaded{cas-sc}{\casscclasstrue}{\casscclassfalse}
  %     \makeatother
  % 
  %     \providecommand\Ifthispageodd[2]{\checkoddpage\ifoddpage#1\else#2\fi}
  %   %\endgroup
  % 
  %   %\begingroup TikZ
  %     %\begingroup Didascalie
  %       \newcommand\newcaption[2]{%
  %         \captionsetup{width=#2} % aboveskip=-.2cm,
  %         \caption{#1}%
  %       }
  %       
  %       \newcommand\newsubcaption[2][ ]{%
  %         \refstepcounter{subfigure}%
  %         (\thesubfigure)#1#2%
  %       }
  %     %\endgroup
  %     
  %     %\begingroup RedefineTikZbounds
  %       % La differenza è che colla prima bisogna catturare anche le etichette perché si crea un pdf isolato mentre colla non è importante perché viene tutto compilato sul momento
  % 
  %       %\begingroup Versione valida con esternalizzazione [ma anche senza]
  %         \newcommand\redefineTikZbounds[2]{
  %           \coordinate (RealSW) at (current bounding box.south west);
  %           \coordinate (RealNE) at (current bounding box.north east);
  %           \coordinate (Center) at ($(#1.west)!0.5!(#2.east)$);
  %     
  %           \pgfresetboundingbox 
  % 
  %           \useasboundingbox 
  %             let \p1=(RealSW),\p2=(RealNE),\p3=(Center),
  %             \n1={max(\x3-\x1,\x2-\x3)} in
  %             (\x3-\n1,\y1) rectangle (\x3+\n1,\y2+.25em);
  % 
  %           % \path[use as bounding box] 
  %           %   let \p1=(RealSW),\p2=(RealNE),\p3=(Center),
  %           %   \n1={max(\x3-\x1,\x2-\x3)} in
  %           %   (\x3-\n1,\y1) rectangle (\x3+\n1,\y2+.25em);
  %           % Lo «.25em» discende dalla traslazione imposta all'etichetta della scala verticale mediante «yScaleLabelShift»
  %         } % Scatola delimitante simmetrica
  %       %\endgroup
  % 
  %       %\begingroup Versione valida senza esternalizzazione
  %         % \newcommand\redefineTikZbounds[2]{
  %         %   % Salva le coordinate della cima e del fondo della figura corrente (cosí da mantenere i titoli e le didascalie)
  %         %   \coordinate (Cima) at (current bounding box.north);
  %         %   \coordinate (Fondo) at (current bounding box.south);
  %     
  %         %   \pgfresetboundingbox % Resetta la scatola dei limiti
  %     
  %         %   \useasboundingbox % Definisce la nuova scatola come un rettangolo
  %         %     (#1.west |- Cima)  % con ascissa sinistra data da «Col1 c1r1.west», 
  %         %     rectangle      % ascissa destra data da «Col2 c1r1.east»
  %         %     (#2.east |- Fondo); % ordinate date dalla cima e fondo originali
  %         %     % L'operatore «|-» in «(A |- B)» prende come ascissa quella del punto A mentre come coordinata quella del punto B
  %         % }
  %       %\endgroup
  %     %\endgroup
  % 
  %     %\begingroup TikZBox
  %       % Parzialmente ispirato da https://gist.github.com/Caramdir/899107
  % 
  %       \makeatletter
  %       \newsavebox\tikzBufferBox
  %       \xdef\offsetX{0pt}
  %       \xdef\offsetY{0pt}
  %       
  %       % \tikzbox[#1]{#2}{#3}[#4]{#5}
  %       % #1 -> Ancora rispetto cui posizionare l'intera scatola (di base è «south west», ossia l'angolo in basso a sinistra)
  %       % #2 -> Nodo dove posizionare l'origine (per la natura di come LaTeX riceve gli argomenti è buon pratica passare questo circondato da parentesi graffe)
  %       % #3 -> Nome della scatola limitante definita dall'ambiente «scope»
  %       % #4 -> Argomenti opzionali da passare all'ambiente «scope»; di base sono vuoti
  %       % #5 -> Contenuto della scatola circondato dall'ambiente «scope»
  %       % Es. \tikzbox{{(0,0)}}{figuraTikZ}{blocco TikZ}
  % 
  %       % Macro di apertura che controlla se v'è l'argomento facoltativo dell'ancora
  %       \def\tikzbox{\pgfutil@ifnextchar[{\tikzbox@anchor}{\tikzbox@anchor[south west]}}
  %       % La funzione «\pgfutil@ifnextchar» serve per controllare se il prossimo carattere è l'apertura di una paretensi graffa «[»
  %       
  %       % Macro intermedia per cogliere gli argomenti obbligatori e controllare quelli successivi facoltativi
  %       \def\tikzbox@anchor[#1]#2#3{\pgfutil@ifnextchar[{\tikzbox@opt[#1]#2#3}{\tikzbox@opt[#1]#2#3[]}}
  % 
  %       % Macro degli argomenti obbligatori
  %       \long\def\tikzbox@opt[#1]#2#3[#4]#5{%
  %         \node[ % Misurazione della scatola in un nodo invisibile
  %           overlay,%
  %           opacity=0,%
  %           inner sep=0pt,%
  %           outer sep=0pt%
  %         ] {%
  %           \sbox{\tikzBufferBox}{%
  %             \begin{tikzpicture}%
  %               \begin{scope}[local bounding box=measuringTikZBox]%
  %                 #5%
  %               \end{scope}%
  %               \pgfpointanchor{measuringTikZBox}{#1}%
  %               \pgfgetlastxy{\tempX}{\tempY}% 
  %               \xdef\offsetX{\tempX}%
  %               \xdef\offsetY{\tempY}%
  %             \end{tikzpicture}%
  %           }%
  %         };%
  %         %
  %         % Disegno reale traslato
  %         \begin{scope}[%
  %           local bounding box=#3,%
  %           shift={#2},%
  %           xshift=-\offsetX,%
  %           yshift=-\offsetY,%
  %           #4%
  %         ]%
  %           #5%
  %         \end{scope}%
  %       }%
  %       \makeatother
  %     %\endgroup
  % 
  %     %\begingroup Chiavi-Valori
  %       %\begingroup Testo
  %         % La dimensione del testo di base è 11pt
  %         \newcommand\titleSize{11}
  %         \newcommand\subCaptionSize{10}
  %         \newcommand\ParameterSize{9}
  %         \newcommand\labelSize{10}
  %         \newcommand\tickSize{9}
  %         \newcommand\legendSize{9}
  %       %\endgroup
  %     
  %       %\begingroup Infopagine
  %         %\begingroup Chiavi condivise
  %         \newcounter{lpc} % «loopcounter»
  %         \newcommand\infopageSharedKeys{
  %           %\begingroup Testo
  %             \renewcommand\titleSize{10}
  %             \renewcommand\subCaptionSize8
  %             \renewcommand\ParameterSize7
  %             \renewcommand\labelSize9
  %             \renewcommand\tickSize6
  %             \renewcommand\legendSize7
  %           %\endgroup
  %     
  %           %\begingroup Dimensioni
  %             \pgfmathsetmacro\plotWidth7
  %             \pgfmathsetmacro\plotHeight{.618*\plotWidth}
  %             \pgfmathsetmacro\halfPlotHeight{\plotHeight/2}
  %             %
  %             \pgfmathparse{0.825*\plotWidth}
  %             \pgfmathsetmacro\plotWidth\pgfmathresult
  %             \pgfmathsetmacro\halfWidth{\plotWidth/2}
  %             \pgfmathsetmacro\plotLastWidth{1.25*\plotWidth}
  %     
  %             \pgfmathsetmacro\xPlotSep{1}
  %             \pgfmathsetmacro\yPlotSep{1.25}
  %             \pgfmathsetmacro\yGroupSep{2.75}
  %           %\endgroup
  %     
  %           %\begingroup Legende
  %             \newcommand\placeLegend[5]{%
  %               /pgfplots/plotLegend,
  %               at={([shift={(##4,##5)}]axis description cs:##1,##2)},
  %               anchor=##3,
  %             }
  %             %
  %             \pgfplotsset{%
  %               plotLegend/.style={%
  %                 font=\textsize{\legendSize},
  %                 /pgfplots/legend cell align=left,
  %                 % align=left,
  %                 % draw=white!15!black,
  %                 draw=none, % Bordo invisibile
  %                 fill=white,
  %                 fill opacity=0.5,
  %                 text opacity=1,
  %                 % text width=2cm,
  %               },
  %               plotLegendNW/.style={legend style={\placeLegend{0}{1}{north west}{0em}{0em}}},
  %               plotLegendNE/.style={legend style={\placeLegend{1}{1}{north east}{0em}{0em}}},
  %               plotLegendSW/.style={legend style={\placeLegend{0}{0}{south west}{0em}{0em}}},
  %               plotLegendSE/.style={legend style={\placeLegend{1}{0}{south east}{0em}{0em}}},
  %               %
  %               % plotLegendW/.style={legend style={plotLegend,at={(0,.5)},anchor=west}},
  %               %
  %               % plotLegendNW/.style={legend style={plotLegend,legend pos=north west}},
  %               % plotLegendNE/.style={legend style={plotLegend,legend pos=north east}},
  %               % plotLegendSW/.style={legend style={plotLegend,legend pos=south west}},
  %               % plotLegendSE/.style={legend style={plotLegend,legend pos=south east}},
  %               %
  %               % plotLegendNWcloser/.style={legend style={plotLegend,at={(0,1)},anchor=north west}},
  %               % plotLegendNEcloser/.style={legend style={plotLegend,at={(1,1)},anchor=north east}},
  %               % plotLegendSWcloser/.style={legend style={plotLegend,at={(0,0)},anchor=south west}},
  %               % plotLegendSEcloser/.style={legend style={plotLegend,at={(0,1)},anchor=south east}},
  %             }
  %             %
  %             \pgfmathsetmacro\lineLegendLength{1}
  %             \pgfmathsetmacro\areaLegendHeigth{.25}
  %             \pgfplotsset{
  %               legend image code/.code={
  %                 \draw [mark repeat=2,mark phase=2,##1]
  %                   plot coordinates {
  %                     (0em,0em)
  %                     (.5*\lineLegendLength em,0em)
  %                     (\lineLegendLength em,0em)
  %                   };
  %               },
  %               /pgfplots/area legend/.style={
  %                 legend image code/.code={
  %                   \draw [##1,fill,draw=none] 
  %                     (0em,-\areaLegendHeigth em) 
  %                     rectangle 
  %                     (\lineLegendLength em,\areaLegendHeigth em);
  %                 },
  %               },
  %             }
  %             %
  %             \newcommand\ExactHistogramLegend{Exact histogram}
  %             \newcommand\ApproxHistogramLegend{Approx. histogram}
  %             \newcommand\RealHistogramLegend{Real histogram}
  %             %
  %             \newcommand\ExactBLNFitLegend{Exact BLN fit}
  %             \newcommand\ApproxBLNFitLegend{Approx BLN fit}
  %             \newcommand\RealBLNFitLegend{Real BLN fit}
  %             %
  %             \newcommand\ExactEmpiricCCDFLegend{Exact empiric CCDF}
  %             \newcommand\ApproxEmpiricCCDFLegend{Approx. empiric CCDF}
  %             \newcommand\RealEmpiricCCDFLegend{Real empiric CCDF}
  %             %
  %             \newcommand\ExactParetoFitLegend{Exact Pareto fit}
  %             \newcommand\ApproxParetoFitLegend{Approx. Pareto fit}
  %             \newcommand\RealParetoFitLegend{Real Pareto fit}
  %             %
  %             \newcommand\ExactScatterLegend{Exact scatter}
  %             \newcommand\ApproxScatterLegend{Approx. scatter}
  %             \newcommand\RealScatterLegend{Real scatter}
  %             %
  %             \newcommand\ClasskILegend{Class k=283}
  %             \newcommand\ClasskIILegend{Class k=83}
  %             \newcommand\ClasskIIILegend{Class k=43}
  %             \newcommand\ClasskIVLegend{Class k=10}
  %           %\endgroup
  %     
  %           %\begingroup Etichette vuote
  %             \pgfplotsset{
  %               yLabelEmpty/.style={ylabel={}},
  %               yTickLabelsEmpty/.style={yticklabels={}},
  %               yTickScaleLabelEmpty/.style={ytick scale label code/.code={}},
  %               xTickScaleLabelEmpty/.style={xtick scale label code/.code={}},
  %               emptyLeged/.style={
  %                 legend style={draw=none,fill=none,nodes={draw=none}},
  %                 legend image post style={opacity=0},
  %               },
  %             }
  %           %\endgroup
  %     
  %           %\begingroup Estetica
  %             \newcommand\smallminus{\scalebox{0.5}[1]{$-$}\hspace{-.11em}}
  %             \pgfplotsset{
  %               /pgfplots/log number format basis/.code 2 args={%
  %                 \pgfmathsetmacro\exponent{##2}%
  %                 \pgfmathparse{sign(\exponent)}%
  %                 \ifnum\pgfmathresult<0%
  %                   \pgfmathparse{abs(\exponent)}%
  %                   $##1^{\smash{\smallminus\pgfmathprintnumber{\pgfmathresult}}}$%
  %                 \else%
  %                   $##1^{\smash{\pgfmathprintnumber{\exponent}}}$%
  %                 \fi%
  %               },
  %               ytick scale label code/.code={%
  %                 \pgfmathsetmacro\exponent{##1}%
  %                 \pgfmathparse{sign(\exponent)}%
  %                 \ifnum\pgfmathresult<0
  %                   \pgfmathparse{abs(\exponent)}%
  %                   $\cdot 10^{\smash{\smallminus\pgfmathprintnumber{\pgfmathresult}}}$%
  %                 \else
  %                   $\cdot 10^{\smash{\pgfmathprintnumber{\exponent}}}$%
  %                 \fi
  %               },
  %             }
  %           %\endgroup
  %     
  %           %\begingroup Colori
  %             \definecolor{blueish}{rgb}{0.00000,0.44700,0.74100}%
  %             \definecolor{reddish}{rgb}{0.83529,0.36863,0.00000}%
  %           %\endgroup
  %     
  %           %\begingroup Commenti
  %             %\begingroup Impostazioni
  %             \tikzset{
  %               title/.style={
  %                 align=center,
  %                 anchor=south,
  %                 font=\bfseries\textsize{\titleSize},
  %               },
  %               fittedParameters/.style={
  %                 align=center,
  %                 font=\textsize{\ParameterSize},
  %               },
  %               varyingParameter/.style={
  %                 align=center,
  %                 % rotate=90,
  %                 font=\textsize{\subCaptionSize},
  %               },
  %               subCaption/.style={
  %                 % text width=\plotWidth cm,
  %                 align=center,
  %                 font=\textsize{\subCaptionSize},
  %               },
  %             }
  %             \newcommand\codetocolor[1]{\ifcase##1blueish\or reddish\else black\fi}
  %             %\endgroup
  %     
  %             %\begingroup Parametri Sardegna
  %             \newcommand\betaReal{\num{1.273}}
  %             \newcommand\xiReal{\num{0.929}}
  %             \newcommand\mIReal{\num{3.227}}
  %             \newcommand\sIReal{\num{0.408}}
  %             \newcommand\mIIReal{\num{4.099}}
  %             \newcommand\sIIReal{\num{0.556}}
  %             %\endgroup
  %     
  %             %\begingroup Parametri Pareto
  %             \newcommand\codetoParetovalue[1]{
  %               \ifcase##1\betaExact\or\betaAprx\else\betaReal\fi
  %             }
  %             \newcommand\writeParetoParameters[7][center]{
  %               \node[fittedParameters,anchor=##1] at ([shift={(##5,##6)}]##4) {%
  %                 \(\begin{aligned}
  %                   \color{\codetocolor{##2}}\beta
  %                   &\color{\codetocolor{##2}}=\codetoParetovalue{##2}\\[##7]
  %                   \color{\codetocolor{##3}}\beta
  %                   &\color{\codetocolor{##3}}=\codetoParetovalue{##3}
  %                 \end{aligned}\)
  %               };
  %             }
  %             %\endgroup
  % 
  %             %\begingroup % Vecchia versione in cui si scrivevano esplicitamente tutte le 9 combinazioni possibili
  %             % \newcommand\writeParameters[6]{%
  %             %   \node[parameter] at ([shift={(##4,##5)}]##1) {%
  %             %     \ifcase##2% Caso 0
  %             %       \ifcase##3% Caso 0
  %             %         \(\begin{aligned}
  %             %           \color{blueish}\beta&\color{blueish}=\betaExact\\[##6]
  %             %           \color{blueish}\beta&\color{blueish}=\betaExact
  %             %         \end{aligned}\)
  %             %       \or% Caso 1
  %             %         \(\begin{aligned}
  %             %           \color{blueish}\beta&\color{blueish}=\betaExact\\[##6]
  %             %           \color{reddish}\beta&\color{reddish}=\betaAprx% Caso 1
  %             %         \end{aligned}\)
  %             %       \else% Caso 2
  %             %         \(\begin{aligned}
  %             %           \color{blueish}\beta&\color{blueish}=\betaExact\\[##6]
  %             %           \color{black}\beta&\color{black}=\betaReal% Caso 2
  %             %         \end{aligned}\)
  %             %       \fi
  %             %     \or% Caso 1
  %             %       \ifcase##3% Caso 0
  %             %         \(\begin{aligned}
  %             %           \color{reddish}\beta&\color{reddish}=\betaAprx\\[##6]
  %             %           \color{blueish}\beta&\color{blueish}=\betaExact
  %             %         \end{aligned}\)
  %             %       \or% Caso 1
  %             %         \(\begin{aligned}
  %             %           \color{reddish}\beta&\color{reddish}=\betaAprx\\[##6]
  %             %           \color{reddish}\beta&\color{reddish}=\betaAprx
  %             %         \end{aligned}\)
  %             %       \else% Caso 2
  %             %         \(\begin{aligned}
  %             %           \color{reddish}\beta&\color{reddish}=\betaAprx\\[##6]
  %             %           \color{black}\beta&\color{black}=\betaReal
  %             %         \end{aligned}\)
  %             %       \fi
  %             %     \else% Caso 2
  %             %       \ifcase##3% Caso 0
  %             %         \(\begin{aligned}
  %             %           \color{black}\beta&\color{black}=\betaReal\\[##6]
  %             %           \color{blueish}\beta&\color{blueish}=\betaExact
  %             %         \end{aligned}\)
  %             %       \or% Caso 1
  %             %         \(\begin{aligned}
  %             %           \color{black}\beta&\color{black}=\betaReal\\[##6]
  %             %           \color{reddish}\beta&\color{reddish}=\betaAprx
  %             %         \end{aligned}\)
  %             %       \else% Caso 2
  %             %         \(\begin{aligned}
  %             %           \color{black}\beta&\color{black}=\betaReal\\[##6]
  %             %           \color{black}\beta&\color{black}=\betaReal
  %             %         \end{aligned}\)
  %             %       \fi
  %             %     \fi
  %             %   };
  %             % }
  %             %\endgroup
  %     
  %             %\begingroup Parametri bilognormale
  %             \newcommand\codetoparameter[1]{
  %               \ifcase##1\xi  % Caso 0
  %               \or\mu_1     % Caso 1
  %               \or\sigma_1   % Caso 2
  %               \or\mu_2     % Caso 3
  %               \else\sigma_2\fi % Caso 4
  %             }
  %             \newcommand\codetoBiLognormalvalue[2]{
  %               \ifcase##1 % Caso 0
  %                 \ifcase##2\xiExact % Caso 0
  %                 \or\mIExact    % Caso 1
  %                 \or\sIExact    % Caso 2
  %                 \or\mIIExact    % Caso 3
  %                 \else\sIIExact\fi % Caso 4
  %               \or % Caso 1
  %                 \ifcase##2\xiAprx % Caso 0
  %                 \or\mIAprx     % Caso 1
  %                 \or\sIAprx     % Caso 2
  %                 \or\mIIAprx    % Caso 3
  %                 \else\sIIAprx\fi  % Caso 4
  %               \else % Caso 2
  %                 \ifcase##2\xiReal % Caso 0
  %                 \or\mIReal     % Caso 1
  %                 \or\sIReal     % Caso 2
  %                 \or\mIIReal    % Caso 3
  %                 \else\sIIReal\fi  % Caso 4
  %               \fi
  %             }
  %             \newcommand\matrixline[2]{\codetoparameter{##2}&=\codetoBiLognormalvalue{##1}{##2}}
  %             \def\lineloop##1##2{%
  %               \matrixline{##1}{\number\value{lpc}}\stepcounter{lpc}%
  %               \ifnum\number\value{lpc}<5\\[##2]\lineloop{##1}{##2}\fi%
  %             } % https://tex.stackexchange.com/questions/491358/tikz-foreach-to-work-with-the-aligned-environment
  %             \newcommand\writeBiLognormalParameters[6][center]{
  %               \node[fittedParameters,anchor=##1] at ([shift={(##4,##5)}]##3) {%
  %                 \setcounter{lpc}{0}
  %                 \(\color{\codetocolor{##2}}
  %                 \begin{aligned}
  %                   \lineloop{##2}{##6}
  %                 \end{aligned}\)
  %               };
  %             }
  %             %\endgroup
  %     
  %             %\begingroup Altre etichette
  %             \newcommand\writeSizeEvolutionDescriptions[1][5]{%
  %               \foreach \col/\clr/\txt in {%
  %                 1/blueish/Esatto,
  %                 2/reddish/Approssimato
  %               }{
  %                 \node[
  %                   fittedParameters,
  %                   anchor=south,
  %                   yshift=.1em,
  %                 ] at (Row##1 c\col r1.south) {\color{\clr}\txt};
  %               }
  %             }
  %     
  %             \newcommand\writeVaryingParameter[5]{%
  %               \foreach \r in {1,...,5}{
  %                 \node[
  %                   varyingParameter,
  %                   anchor=##3,
  %                 ] at ([shift={(##4,##5)}]Row\r\space c##1r1.##2)
  %                 {\(\varyingParameterSymbol=\varyingParameterList\r\)};
  %               }
  %             }
  %             %\endgroup
  %           %\endgroup
  %     
  %           %\begingroup Sottodidascalie
  %             \newcommand\writeSubCaption[5]{%
  %               \node[subCaption,anchor=##3]
  %               at ([shift={(##4,##5)}]##1.##2)
  %               {\newsubcaption[]{}};
  %             }
  %             \newcommand\writeSubCaptions[5][3]{%
  %               \foreach \r in {1,...,5}{
  %                 \foreach \c in {1,...,##1}{
  %                   \writeSubCaption{Row\r\space c\c r1}{##2}{##3}{##4}{##5}
  %                 }
  %               }
  %             }
  %           %\endgroup
  %         }
  %         %\endgroup
  %     
  %         %\begingroup Chiavi del primo tipo d'infopagina
  %         \newcommand\infopageKeysI{
  %           \infopageSharedKeys
  %     
  %           %\begingroup Gruppi
  %             \ifpageisodd
  %               \pgfplotsset{/pgfplots/group/yDescPos/.style={y descriptions at=edge right}}
  %             \else
  %               \pgfplotsset{/pgfplots/group/yDescPos/.style={y descriptions at=edge left}}
  %             \fi
  %             \pgfplotsset{
  %               /pgfplots/group/plotGroup/.style={
  %                 horizontal sep=\xPlotSep em,
  %                 vertical sep=\yPlotSep em,
  %                 yDescPos,
  %                 group size=3 by 1,
  %               },
  %               /pgfplots/group/plotLastGroup/.style={
  %                 plotGroup,
  %                 x descriptions at=edge bottom,
  %               },
  %             }
  %           %\endgroup
  %     
  %           %\begingroup Stili
  %             %\begingroup Selezione «recto»/«verso»
  %             \ifpageisodd
  %               \pgfplotsset{%
  %                 yLineAxis/.style={axis y line*=right},
  %                 /tikz/yLabelShift/.style={%
  %                   yshift=-1.25em,
  %                   at={(axis description cs:1,0.5)},
  %                   anchor=north,
  %                   % rotate=180,
  %                 },
  %                 /tikz/yScaleLabelShift/.style={%
  %                   yshift=.25em,
  %                   at={(1,1)},
  %                   anchor=south east,
  %                 },
  %                 /tikz/yTickLabelShift/.style={xshift=-0.15em},
  %               }
  %             \else
  %               \pgfplotsset{%
  %                 yLineAxis/.style={axis y line*=left},
  %                 /tikz/yLabelShift/.style={%
  %                   yshift=1.25em,
  %                   at={(axis description cs:0,0.5)},
  %                   anchor=south,
  %                   % rotate=180,
  %                 },
  %                 /tikz/yScaleLabelShift/.style={%
  %                   yshift=.25em,
  %                   at={(0,1)},
  %                   anchor=south west,
  %                 },
  %                 /tikz/yTickLabelShift/.style={xshift=+0.15em},
  %               }
  %             \fi
  %             %\endgroup
  %             %
  %             %\begingroup Grafici
  %             \pgfplotsset{
  %               plotRow/.style={
  %                 scale only axis,
  %                 width=\plotWidth cm,height=\plotHeight cm,
  %                 xlabel style={font=\color{white!15!black}\textsize{\labelSize}},
  %                 ylabel style={font=\color{white!15!black}\textsize{\labelSize},yLabelShift},
  %                 xticklabel style={font=\textsize{\tickSize}},
  %                 yticklabel style={font=\textsize{\tickSize},yTickLabelShift},
  %                 y tick scale label style={yScaleLabelShift},
  %                 % axis x line*=bottom,yLineAxis,
  %                 enlarge x limits=0.1,enlarge y limits=0.1,
  %                 xmajorgrids,ymajorgrids,
  %                 xminorticks=true,yminorticks=true,
  %                 scale ticks above exponent=2,
  %                 axis background/.style={fill=white},
  %                 grid style={dashed}
  %               },
  %               % xlabel123/.style={xlabel={\(h\)}},
  %               xlabel3/.style={xlabel={\(h\)}},
  %               ylabel23/.style={ylabel={\(\FRC(h)\)}},
  %               row1Keys/.style={
  %                 % ylabel={\(\underhat{\bar l}_N(h)\)},
  %                 ylabel={\(\bar l_N(h)\)},
  %                 % enlarge x limits={lower,value=.05},
  %                 enlarge x limits=.05,
  %                 enlarge y limits={upper,value=.05},
  %                 xmode=linear,
  %                 log origin=infty,
  %               },
  %               row2KeysWithLabels/.style={
  %                 ylabel23,
  %                 xmode=linear,ymode=log
  %               },
  %               row2Keys/.style={
  %                 ylabel23,
  %                 xmode=linear,ymode=log,
  %                 xticklabels={},
  %               },
  %               row3Keys/.style={
  %                 xlabel3,ylabel23,
  %                 xmode=linear,ymode=log
  %               },
  %               row4Keys/.style={
  %                 xlabel={\(\xi\)},ylabel={\(s(\xi)\)},
  %                 xmode=log,ymode=log
  %               },
  %               row5KeysLin/.style={
  %                 width=\plotLastWidth cm,
  %                 xlabel={\(t\)},ylabel={\(s(t)\)},
  %                 enlarge x limits=0,
  %                 enlarge y limits=0.05,
  %                 x tick scale label style={
  %                   yshift=-1em,
  %                   at={(.95,0)},
  %                   anchor=north,
  %                 },
  %               },
  %               row5Keys/.style={
  %                 row5KeysLin,
  %                 ymode=log,
  %               },
  %             }
  %             %\endgroup
  %           %\endgroup
  %         }
  %         %\endgroup
  %     
  %         %\begingroup Chiavi del secondo tipo d'infopagina
  %         \newcommand\infopageKeysII{
  %           \infopageSharedKeys
  %     
  %           %\begingroup Gruppi
  %             \ifpageisodd
  %               \pgfplotsset{/pgfplots/group/yDescPos/.style={y descriptions at=edge right}}
  %             \else
  %               \pgfplotsset{/pgfplots/group/yDescPos/.style={y descriptions at=edge left}}
  %             \fi
  %             \pgfplotsset{
  %               /pgfplots/group/plotGroup/.style={
  %                 horizontal sep=\xPlotSep em,
  %                 vertical sep=\yPlotSep em,
  %                 y descriptions at=all,
  %                 group size=3 by 1,
  %               },
  %               /pgfplots/group/plotLastGroup/.style={
  %                 plotGroup,
  %                 x descriptions at=edge bottom,
  %               },
  %             }
  %           %\endgroup
  %     
  %           %\begingroup Stili
  %             %\begingroup Selezione «recto»/«verso»
  %             \ifpageisodd
  %               \pgfplotsset{
  %                 yLabel/.style={
  %                   ylabel style={
  %                     font=\color{white!15!black}\textsize{\labelSize},
  %                     yshift=-1.25em,
  %                     at={(axis description cs:1,0.5)},
  %                     anchor=north,
  %                     % rotate=180,
  %                   },
  %                 },
  %                 yLabelInv/.style={yLabelEmpty},
  %                 yScaleLabelShift/.style={
  %                   y tick scale label style={
  %                     yshift=.25em,
  %                     % xshift=1em,
  %                     at={(1,1)},
  %                     anchor=south east,
  %                   },
  %                 },
  %                 yScaleLabelShiftInv/.style={
  %                   y tick scale label style={
  %                     yshift=.25em,
  %                     % xshift=-1em,
  %                     at={(0,1)},
  %                     anchor=south west,
  %                   },
  %                 },
  %                 % yLineAxis/.style={axis y line*=right},
  %               }
  %             \else
  %               \pgfplotsset{
  %                 yLabel/.style={yLabelEmpty},
  %                 yLabelInv/.style={
  %                   ylabel style={
  %                     font=\color{white!15!black}\textsize{\labelSize},
  %                     yshift=1.25em,
  %                     at={(axis description cs:0,0.5)},
  %                     anchor=south,
  %                     % rotate=180,
  %                   },
  %                 },
  %                 yScaleLabelShift/.style={
  %                   y tick scale label style={
  %                     yshift=.25em,
  %                     % xshift=1em,
  %                     at={(1,1)},
  %                     anchor=south east,
  %                   },
  %                 },
  %                 yScaleLabelShiftInv/.style={
  %                   y tick scale label style={
  %                     yshift=.25em,
  %                     % xshift=-1em,
  %                     at={(0,1)},
  %                     anchor=south west,
  %                   },
  %                 },
  %                 % yLineAxis/.style={axis y line*=left},
  %               }
  %             \fi
  %             %\endgroup
  %             %
  %             %\begingroup Grafici
  %             \pgfplotsset{
  %               plotRow/.style={
  %                 scale only axis,
  %                 width=\plotWidth cm,height=\plotHeight cm,
  %                 xlabel style={font=\color{white!15!black}\textsize{\labelSize}},
  %                 % ylabel style={font=\color{white!15!black}\textsize{\labelSize},yLabelShift},
  %                 xticklabel style={font=\textsize{\tickSize}},
  %                 % axis x line*=bottom,yLineAxis,
  %                 enlarge x limits=0.1,enlarge y limits=0.1,
  %                 xmajorgrids,ymajorgrids,
  %                 xminorticks=true,yminorticks=true,
  %                 scale ticks above exponent=2,
  %                 axis background/.style={fill=white},
  %                 grid style={dashed},
  %               },
  %               plotLeftCol/.style={
  %                 yLabelInv,
  %                 yticklabel style={font=\textsize{\tickSize},xshift=0.15em},
  %                 yticklabel pos=left,
  %                 yScaleLabelShiftInv,
  %               },
  %               plotMiddleCol/.style={
  %                 yLabelEmpty,
  %                 yTickLabelsEmpty,
  %                 yTickScaleLabelEmpty,
  %               },
  %               plotRightCol/.style={
  %                 yLabel,
  %                 yticklabel style={font=\textsize{\tickSize},xshift=-0.15em},
  %                 yticklabel pos=right,
  %                 yScaleLabelShift,
  %               },
  %               xlabel3/.style={xlabel={\(h\)}},
  %               ylabel23/.style={ylabel={\(\FRC(h)\)}},
  %               row1Keys/.style={
  %                 ylabel={\(\bar l_N(h)\)},
  %                 enlarge x limits=.05,
  %                 enlarge y limits={upper,value=.05},
  %                 xmode=linear,
  %                 log origin=infty,
  %               },
  %               row2Keys/.style={
  %                 ylabel23,
  %                 xmode=linear,ymode=log,
  %                 % xticklabels={},
  %               },
  %               row3Keys/.style={
  %                 xlabel3,ylabel23,
  %                 xmode=linear,ymode=log,
  %               },
  %               row4KeysLin/.style={
  %                 xlabel={\(\xi\)},ylabel={\(s(\xi)\)},
  %                 xmode=log,
  %               },
  %               row4Keys/.style={row4KeysLin,ymode=log},
  %               row5KeysLin/.style={
  %                 width=\plotLastWidth cm,
  %                 xlabel={\(t\)},ylabel={\(s(t)\)},
  %                 enlarge x limits=0,
  %                 enlarge y limits=0.05,
  %                 x tick scale label style={
  %                   yshift=-1em,
  %                   at={(.95,0)},
  %                   anchor=north,
  %                 },
  %               },
  %               row5Keys/.style={
  %                 row5KeysLin,
  %                 ymode=log,
  %               },
  %             }
  %             %\endgroup
  %           %\endgroup
  %         }
  %         %\endgroup
  %       %\endgroup
  %     %\endgroup
  % 
  %     %\begingroup Diagrammi dei nodi
  %     \newcommand\nodeDiagramSettings{
  %       %\begingroup Dimensioni
  %         \pgfmathsetmacro\xPlotSep{10}
  %         \pgfmathsetmacro\yPlotSep{4}
  %         \pgfmathsetmacro\plotScale{.8}
  % 
  %         \pgfmathsetmacro\titleSize{11}
  %         \pgfmathsetmacro\commentSize{10}
  % 
  %         \pgfmathsetmacro\nodeSize{20pt} % Diametro del nodo
  %         \pgfmathsetmacro\starSize{2.25*\nodeSize}
  %         \pgfmathsetmacro\numberSize{.53*\nodeSize}
  %       %\endgroup
  % 
  %       %\begingroup Colori
  %         % \definecolor{cRed}{HTML}{D68490}
  %         % \definecolor{cPurple}{HTML}{7E6EAC}
  %         % \definecolor{cTeal}{HTML}{87BDB1}
  %         % \definecolor{cOrange}{HTML}{F2D0A4}
  %         % \definecolor{cEdgeLight}{HTML}{EEE8D5} 
  %         % \definecolor{cEdgeDark}{HTML}{998888}
  %         % \definecolor{cStar}{HTML}{FDF6E3}
  %         %
  %         \definecolor{color1}{HTML}{800000} % Deep Red
  %         \definecolor{color2}{HTML}{4682B4} % Steel Blue
  %         \definecolor{color3}{HTML}{6B8E23} % Olive Green
  %         \definecolor{color4}{HTML}{708090} % Slate Gray
  %         \definecolor{color5}{HTML}{BC8F8F} % Rosy Brown
  %         \definecolor{color6}{HTML}{008B8B} % Dark Cyan
  %         \definecolor{color7}{HTML}{B8860B} % Dark Goldenrod
  %       %\endgroup
  % 
  %       %\begingroup Testo
  %         \tikzset{
  %           titleStyle/.style={
  %             font=\textsize{\titleSize},
  %             align=center,
  %             anchor=south,
  %           },
  %           commentStyle/.style={font=\textsize{\commentSize}},
  %           arrowCommentStyle/.style={
  %             midway,
  %             commentStyle,
  %             left,
  %             sloped,
  %             below
  %           },
  %           numberStyle/.style={font=\textsize{\numberSize}}
  %         }
  %       %\endgroup
  % 
  %       %\begingroup Frecce
  %         \tikzset{
  %           transformingArrow/.style={
  %             % =Latex,
  %             ->,
  %             >=Stealth,
  %             dashed,
  %             line width=.8pt,
  %             black!90,
  %             shorten >=5pt,
  %             shorten <=10pt
  %           },
  %           indirectArrow/.style={
  %             draw=black!95,
  %             line width=.5pt,
  %           },
  %           nodeArrowExact/.style={
  %             indirectArrow,
  %             % ->,
  %             % -{Stealth[scale=1.2]},
  %             -{Triangle[scale=.75]},
  %             % draw=cEdgeLight,
  %             % draw=red!50,
  %             % shorten >=2pt,
  %             % shorten <=2pt
  %           },
  %           nodeArrowApprox/.style={
  %             indirectArrow,
  %             -,
  %             dashed,
  %           },
  %           transparentNodeArrow/.style={
  %             opacity=.4,
  %           },
  %           arrowAxis/.style={
  %             ->,
  %             line width=.3pt,
  %             >=Latex,
  %           },
  %         }
  %       %\endgroup
  % 
  %       %\begingroup Nodi
  %         \tikzset{
  %           nodeStyle/.style={
  %             circle,
  %             minimum size=\nodeSize,
  %             % draw=black!80,
  %             draw=black,
  %             line width=.5pt,
  %             inner sep=0pt,
  %           },
  %           starStyle/.style={
  %             star,
  %             star points=5,
  %             star point ratio=2,
  %             % fill=cStar,
  %             fill=black!50,
  %             fill opacity=.25,
  %             minimum size=\starSize,
  %             draw=none
  %           },
  %         }
  %       %\endgroup
  %     }
  %     %\endgroup
  % 
  %     %\begingroup GNU-Plot
  %       % \IfFileExists{\sourceauxpath.aux}{
  %       %   \externaldocument{\sourceauxpath}
  %       % }{\relax}
  %       
  %       \pgfkeys{/pgf/plot/gnuplot call={cd /d "\sourceoutdir" && gnuplot}}
  %     %\endgroup
  %   %\endgroup
  % %\endgroup
  % 
  % %\begingroup For the paper
  %   \def\tsc#1{\csdef{#1}{\textsc{\lowercase{#1}}\xspace}}
  %   \tsc{WGM}
  %   \tsc{QE}
  %   
  %   \newcommand\etal{et al.}
  %   \newcommand\KTMAS{\hyperlink{acrKineticTheoryMultiAgentSystems}{KTMAS}}
  %   
  %   % \newcommand\MathSpacing[3]{
  %   %   \thinmuskip = #1mu\relax
  %   %   \medmuskip = #2mu\relax
  %   %   \thickmuskip = #3mu\relax
  %   % }
  %   % \MathSpacing123
  %   
  %   \newcommand\RPlus{\mathbb R_+}
  %   \newcommand\NPlus{\mathbb N_+}
  %   \DeclareMathAlphabet\pazocal{OMS}{zplm}{m}{n}
  %   
  %   \DeclareMathOperator\Lognormal{Lognormal}
  %   \DeclareMathOperator\BiLognormal{BiLognormal}
  %   \DeclareMathOperator\Bernoulli{Bernoulli}
  %   \DeclareMathOperator\Exponential{Exp}
  %   \DeclareMathOperator\Student{Student}
  %   \DeclareMathOperator\Variance{var}
  %   \DeclareMathOperator\IC{IC}
  %   \DeclareMathOperator\Histogram{\mathcal B}
  %   \DeclareMathOperator\CCDF{\mathcal R}
  %     
  %   \newcommand\ThreeQuarters{{3\hspace{-0.15em}/\hspace{-0.15em}4}}
  %   \newcommand\OneQuarter{{1\hspace{-0.15em}/\hspace{-0.15em}4}}
  %   
  %   \newcommand\de[1][]{\mathrm{d}#1}
  %   \newcommand\derS[3][d]{\frac{\mathrm{#1}#2}{\mathrm{#1}#3}}
  %   \newcommand\derP[2]{\frac{\partial#1}{\partial#2}}
  %   \newcommand\underhat[1]{\underaccent{\hat}{#1}}
  %   \renewcommand\undertilde[1]{\underaccent{\tilde}{#1}}
  %   
  %   \makeatletter
  %   \newcommand\eqtag{\@ifstar{\eqtag@star}{\eqtag@nostar}}
  %   
  %   \newcounter{eqtag}
  %   \newcommand\eqtag@nostar[2]{%
  %     \begin{equation*}%
  %       \refstepcounter{eqtag}%
  %       % \protected@edef\@currentlabeltype{eqtag}%
  %       \protected@edef\@currentlabel{#1}%
  %       #2%
  %       \tag*{#1}%
  %     \end{equation*}%
  %   }
  %   \newcommand\eqtag@star[2]{%
  %     \begin{equation*}%
  %       \refstepcounter{eqtag}%
  %       % \protected@edef\@currentlabeltype{eqtag}%
  %       \protected@edef\@currentlabel{#1}%
  %       \let\label\ltx@label%
  %       #2%
  %       % \notag%
  %     \end{equation*}%
  %   }
  %   \makeatother
  %   
  %   \newcommand\crefsubfigInt[4][-]{Figures~\ref{#2}#3#1\ref{#2}#4}
  %   \newcommand\crefsubfigSep[4]{Figures~\ref{#1}#2\foreach \l in#3{,~\ref{#1}\l}~and~\ref{#1}#4}
  %   
  %   \crefname{section}{§}{§§}
  %   \crefrangeformat{section}{§§~#3#1#4--#5#2#6}
  %   \crefmultiformat{section}{§§~#2#1#3}{ and~#2#1#3}{, #2#1#3}{ and~#2#1#3}
  %   \crefrangemultiformat{section}{§§~#3#1#4--#5#2#6}{ and~#3#1#4--#5#2#6}{, #3#1#4--#5#2#6}{ and~#3#1#4--#5#2#6}
  %   
  %   \crefname{equation}{}{}
  %   \crefrangeformat{equation}{(#3#1#4--#5#2#6)}
  %   \crefmultiformat{equation}{(#2#1#3}{,~#2#1#3)}{,~#2#1#3}{,~#2#1#3)}
  %   \crefrangemultiformat{equation}{(#3#1#4--#5#2#6}{,~#3#1#4--#5#2#6)}{,~#3#1#4--#5#2#6}{,~#3#1#4--#5#2#6)}
  %   
  %   \crefname{figure}{Figure}{Figures}
  %   \crefrangeformat{figure}{Figures~#3#1#4--#5#2#6}
  %   \crefmultiformat{figure}{Figures~#2#1#3}{ and~#2#1#3}{, #2#1#3}{ and~#2#1#3}
  %   \crefrangemultiformat{figure}{Figures~#3#1#4--#5#2#6}{ and~#3#1#4--#5#2#6}{, #3#1#4--#5#2#6}{ and~#3#1#4--#5#2#6}
  %   
  %   \crefname{table}{Table}{Tables}
  %   \crefrangeformat{table}{Tables~#3#1#4--#5#2#6}
  %   \crefmultiformat{table}{Tables~#2#1#3}{ and~#2#1#3}{, #2#1#3}{ and~#2#1#3}
  %   \crefrangemultiformat{table}{Tables~#3#1#4--#5#2#6}{ and~#3#1#4--#5#2#6}{, #3#1#4--#5#2#6}{ and~#3#1#4--#5#2#6}
  %   
  %   \crefname{algorithm}{Algorithm}{Algorithms}
  %   \crefrangeformat{algorithm}{Algorithms~#3#1#4--#5#2#6}
  %   \crefmultiformat{algorithm}{Algorithms~#2#1#3}{ and~#2#1#3}{, #2#1#3}{ and~#2#1#3}
  %   \crefrangemultiformat{algorithm}{Algorithms~#3#1#4--#5#2#6}{ and~#3#1#4--#5#2#6}{, #3#1#4--#5#2#6}{ and~#3#1#4--#5#2#6}
  %   
  %   %\newtheorem{theorem}{Theorem}
  %   %\newtheorem{lemma}[theorem]{Lemma}
  %   \newdefinition{dfn}{Definition}
  %   \newdefinition{rmk}{Remark}
  %   \newdefinition{hyp}{Hypotesis}
  %   %\newproof{pf}{Proof}
  %   %\newproof{pot}{Proof of Theorem \ref{thm}}
  %   
  %   \crefname{dfn}{Definition}{Definitions}
  %   \crefrangeformat{dfn}{Definitions~#3#1#4--#5#2#6}
  %   \crefmultiformat{dfn}{Definitions~#2#1#3}{ and~#2#1#3}{, #2#1#3}{ and~#2#1#3}
  %   \crefrangemultiformat{dfn}{Definitions~#3#1#4--#5#2#6}{ and~#3#1#4--#5#2#6}{, #3#1#4--#5#2#6}{ and~#3#1#4--#5#2#6}
  %   \crefname{hyp}{Hypothesis}{Hypotheses}
  %   
  %   \crefrangeformat{hyp}{Hypotheses~#3#1#4--#5#2#6}
  %   \crefmultiformat{hyp}{Hypotheses~#2#1#3}{ and~#2#1#3}{, #2#1#3}{ and~#2#1#3}
  %   \crefrangemultiformat{hyp}{Hypotheses~#3#1#4--#5#2#6}{ and~#3#1#4--#5#2#6}{, #3#1#4--#5#2#6}{ and~#3#1#4--#5#2#6}
  %   
  %   \crefname{rmk}{Remark}{Remarks}
  %   \crefrangeformat{rmk}{Remarks~#3#1#4--#5#2#6}
  %   \crefmultiformat{rmk}{Remarks~#2#1#3}{ and~#2#1#3}{, #2#1#3}{ and~#2#1#3}
  %   \crefrangemultiformat{rmk}{Remarks~#3#1#4--#5#2#6}{ and~#3#1#4--#5#2#6}{, #3#1#4--#5#2#6}{ and~#3#1#4--#5#2#6}
  %   
  %   % \newcommand\tikzfigure[5][0em]{%
  %   %   \begin{figure}[%
  %   %     pos=#2,%
  %   %     width=#4,%
  %   %     align=\centering%
  %   %   ]%
  %   %     \centering%
  %   %     \vspace*{#1}%
  %   %     \IfSubStr{#2}{p}{%
  %   %       \leavevmode\stepcounter{figposcount}%
  %   %       \zsaveposx{figpos-\thefigposcount}%
  %   %       \makebox[0pt][l]{%
  %   %         \hspace*{-\dimexpr\zposx{figpos-\thefigposcount}sp\relax}%
  %   %         \hspace*{\dimexpr0.5\paperwidth\relax}%
  %   %         \makebox[0pt][c]{%
  %   %           \scalebox{0.86}{\includegraphics{figs/#5.pdf}}%
  %   %         }%
  %   %       }
  %   %     }{%
  %   %       \includegraphics{figs/#5.pdf}%
  %   %     }%
  %   %     \caption{%
  %   %       \hyphenpenalty=0%
  %   %       \exhyphenpenalty=0%
  %   %       #3%
  %   %     }%
  %   %     \protect\label{#5}%
  %   %   \end{figure}%
  %   % }
  % %\endgroup
%\endgroup


% Tutte e cinque le pagine alla stessa larghezza, con la figura centrata.
% L'opzione «tikz» di standalone implica «multi=tikzpicture»: ogni pagina è la
% scatola delimitante della propria tikzpicture, perciò standalone ritaglia
% ciascuna pagina sul proprio contenuto e le larghezze naturali vanno da
% 502pt (SE) a 569pt (SDF1); sfogliando il PDF la figura sembra spostarsi.
% Non si può rimediare incapsulando l'\input in un \hbox, che standalone
% scarterebbe. Si ridefinisce invece \redefineTikZbounds, la macro che
% ciascuna Infopagina chiama per ultima, con due modifiche: una semilarghezza
% minima, e la centratura sul centro della scatola delimitante reale anziché
% sul punto medio della griglia dei grafici. Quest'ultima serve perché su SE
% l'inchiostro sborda a destra della griglia: a parità di larghezza di pagina
% la figura restava spostata (69.6bp di margine a sinistra contro 19.2bp).
% In cambio le colonne dei grafici non cadono più esattamente alla stessa
% ascissa in tutte le pagine. Gli argomenti #1 e #2 non servono più, ma la
% firma resta a due argomenti per compatibilità colle chiamate esistenti.
% La ridefinizione è locale a questo master; definizioni.tex non viene
% toccato, così le figure dell'articolo conservano le scatole attuali.
\def\semilarghezzaMinima{286} % pt: metà della larghezza comune di pagina (572pt)
% «current bounding box» registra a sinistra circa 36bp che non contengono
% inchiostro (l'etichetta della scala verticale, traslata da yScaleLabelShift:
% cf. definizioni.tex:59). Centrando sulla scatola, la figura resterebbe
% spostata a destra di 36bp esatti — lo scarto misurato è identico su tutte e
% dieci le pagine dei due studi — perciò si sposta il centro di metà di esso.
% Per riverificarlo: rendere le pagine con «pdftoppm -r 30 -gray» e
% confrontare i margini d'inchiostro a sinistra e a destra.
\def\correzioneCentratura{18.07pt} % = 18bp, metà dello scarto misurato
\renewcommand\redefineTikZbounds[2]{
  \coordinate (RealSW) at (current bounding box.south west);
  \coordinate (RealNE) at (current bounding box.north east);
  \coordinate (Center) at ($(RealSW)!0.5!(RealNE)+(\correzioneCentratura,0)$);

  \pgfresetboundingbox

  \useasboundingbox
    let \p1=(RealSW),\p2=(RealNE),\p3=(Center),
    % Si misura solo verso destra: \x2 è il bordo reale dell'inchiostro,
    % mentre \x1 comprende i 36bp fantasma e, col centro corretto, gonfierebbe
    % la semilarghezza rompendo l'uniformità delle pagine.
    \n1={max(\x2-\x3,\semilarghezzaMinima)} in
    (\x3-\n1,\y1) rectangle (\x3+\n1,\y2+.25em);
}

\begin{document}
  \def\even{even}
  \newif\ifpageisodd

  \ifx\isodd\even
    \pageisoddfalse
  \else
    \pageisoddtrue
  \fi

  \let\originfopageKeysI\infopageKeysI
  \renewcommand\infopageKeysI{%
    \originfopageKeysI
    \renewcommand\ClasskILegend{Class =283}%
    \renewcommand\ClasskIILegend{Class =83}%
    \renewcommand\ClasskIIILegend{Class =43}%
    \renewcommand\ClasskIVLegend{Class =10}%
  }

  \begingroup
    \setcounter{subfigure}{0}
    
% !TEX root = ../../../../../Esperimenti/.tex/MEF.tex
% LTeX: language=it

\begin{tikzpicture}%
  %\begingroup Impostazioni
    \infopageKeysI
    \pgfplotsset{
      row1Keys/.append style={enlarge x limits=.05}
    }
  %\endgroup

  %\begingroup Righe
    \begin{groupplot}[
      group style={
        group name=Row1,
        plotGroup,
      },
      plotRow,
      row1Keys,
      plotLegendNW,
      %
      % xmin=1,xmax=2.3788e+05,
      ymin=0,ymax=1.2,
    ]
      \nextgroupplot[%
        xmin=0.57363,xmax=5.4115,
        % xmin=2.248,xmax=5.3456,
        % ymin=0,ymax=1.2,
      ] % This file was created by matlab2tikz.
%
\definecolor{mycolor1}{rgb}{0.00000,0.44706,0.69804}%
\definecolor{mycolor2}{rgb}{0.83529,0.36863,0.00000}%
%
\addplot[ybar interval, fill=mycolor1, fill opacity=0.35, area legend, draw=none] table[row sep=crcr, x=Lower, y=Count] {%
Lower	Upper	Count\\
2.3888	2.6054	0.041733\\
2.6054	2.822	0.2014\\
2.822	3.0386	0.54955\\
3.0386	3.2553	1.0394\\
3.2553	3.4719	1.1082\\
3.4719	3.6885	0.79207\\
3.6885	3.9051	0.39148\\
3.9051	4.1217	0.22615\\
4.1217	4.3383	0.14391\\
4.3383	4.5549	0.073002\\
4.5549	4.7715	0.027453\\
4.7715	4.9882	0.0098486\\
4.9882	5.2048	0.012311\\
5.2048	5.2048	0.012311\\
};
\addlegendentry{\ExactHistogramLegend}

\addplot [color=mycolor1, only marks, every error bar/.append style={opacity=0.45}, mark=*, mark size=0pt, draw=none, forget plot]
 plot [error bars/.cd, y dir=both, y explicit, error bar style={line width=1pt, color=mycolor1}, error mark options={mark=none,mark size=0pt}]
 table[row sep=crcr, y error plus index=2, y error minus index=3]{%
2.4971	0.041733	0.002687	0.002687\\
2.7137	0.2014	0.0050135	0.0050135\\
2.9303	0.54955	0.0094032	0.0094032\\
3.1469	1.0394	0.011933	0.011933\\
3.3636	1.1082	0.011346	0.011346\\
3.5802	0.79207	0.0094842	0.0094842\\
3.7968	0.39148	0.0071828	0.0071828\\
4.0134	0.22615	0.0058268	0.0058268\\
4.23	0.14391	0.0035355	0.0035355\\
4.4466	0.073002	0.0023097	0.0023097\\
4.6632	0.027453	0.0019901	0.0019901\\
4.8799	0.0098486	0.0011515	0.0011515\\
5.0965	0.012311	2.0756e-18	1.7347e-18\\
};
\addplot [color=mycolor1]
  table[row sep=crcr]{%
2.3888	0.011106\\
2.417	0.014703\\
2.4424	0.018815\\
2.4649	0.023309\\
2.4875	0.028736\\
2.5072	0.03437\\
2.5269	0.040947\\
2.5466	0.048587\\
2.5636	0.056076\\
2.5805	0.064523\\
2.5974	0.074012\\
2.6143	0.084631\\
2.6312	0.096468\\
2.6481	0.10961\\
2.6622	0.12162\\
2.6763	0.13464\\
2.6904	0.14872\\
2.7045	0.1639\\
2.7186	0.18021\\
2.7355	0.20134\\
2.7524	0.22419\\
2.7693	0.24881\\
2.7862	0.27521\\
2.8032	0.30338\\
2.8201	0.33332\\
2.837	0.36497\\
2.8567	0.40399\\
2.8764	0.44512\\
2.899	0.49447\\
2.9244	0.55255\\
2.9526	0.61952\\
2.9977	0.72946\\
3.0399	0.83169\\
3.0653	0.89066\\
3.085	0.93446\\
3.1048	0.97592\\
3.1217	1.0092\\
3.1358	1.0352\\
3.1499	1.0594\\
3.164	1.0816\\
3.1752	1.0979\\
3.1865	1.1128\\
3.1978	1.1262\\
3.2091	1.1381\\
3.2203	1.1485\\
3.2288	1.1553\\
3.2373	1.1611\\
3.2457	1.1661\\
3.2542	1.1701\\
3.2626	1.1732\\
3.2711	1.1753\\
3.2795	1.1765\\
3.288	1.1768\\
3.2964	1.1762\\
3.3049	1.1746\\
3.3134	1.1721\\
3.3218	1.1688\\
3.3303	1.1645\\
3.3387	1.1594\\
3.3472	1.1534\\
3.3585	1.1441\\
3.3697	1.1334\\
3.381	1.1214\\
3.3923	1.108\\
3.4036	1.0934\\
3.4177	1.0735\\
3.4318	1.052\\
3.4487	1.0242\\
3.4656	0.99448\\
3.4853	0.95787\\
3.5079	0.91398\\
3.5389	0.85131\\
3.6206	0.68456\\
3.646	0.63527\\
3.6685	0.59329\\
3.6911	0.55331\\
3.7108	0.52016\\
3.7305	0.48882\\
3.7503	0.45934\\
3.7672	0.43557\\
3.7841	0.41319\\
3.801	0.39217\\
3.8179	0.37248\\
3.8348	0.35408\\
3.8518	0.33691\\
3.8687	0.32091\\
3.8856	0.30601\\
3.9053	0.28992\\
3.925	0.27512\\
3.9448	0.26149\\
3.9645	0.24891\\
3.9871	0.23566\\
4.0096	0.22348\\
4.035	0.21086\\
4.0632	0.19796\\
4.0942	0.18488\\
4.128	0.17165\\
4.1675	0.15727\\
4.2126	0.14189\\
4.2605	0.12652\\
4.3112	0.11117\\
4.362	0.096759\\
4.4099	0.08405\\
4.4578	0.072284\\
4.5029	0.062126\\
4.548	0.052895\\
4.5931	0.044608\\
4.6382	0.037258\\
4.6833	0.030819\\
4.7312	0.024925\\
4.7791	0.019938\\
4.8299	0.015553\\
4.8863	0.011632\\
4.9454	0.0084353\\
5.0131	0.0057235\\
5.0892	0.0036041\\
5.1822	0.0019723\\
5.2048	0.0016935\\
};
\addlegendentry{\ExactBLNFitLegend}


\addplot[area legend, draw=none, fill=mycolor1, fill opacity=0.15, forget plot]
table[row sep=crcr] {%
x	y\\
2.3888	0.011357\\
2.3916	0.011682\\
2.3944	0.012015\\
2.3973	0.012358\\
2.4001	0.012709\\
2.4029	0.013069\\
2.4057	0.013438\\
2.4085	0.013817\\
2.4113	0.014206\\
2.4142	0.014604\\
2.417	0.015013\\
2.4198	0.015432\\
2.4226	0.015862\\
2.4254	0.016303\\
2.4283	0.016754\\
2.4311	0.017217\\
2.4339	0.017692\\
2.4367	0.018178\\
2.4395	0.018676\\
2.4424	0.019187\\
2.4452	0.01971\\
2.448	0.020246\\
2.4508	0.020795\\
2.4536	0.021357\\
2.4564	0.021933\\
2.4593	0.022522\\
2.4621	0.023126\\
2.4649	0.023744\\
2.4677	0.024377\\
2.4705	0.025025\\
2.4734	0.025688\\
2.4762	0.026366\\
2.479	0.027061\\
2.4818	0.027771\\
2.4846	0.028499\\
2.4875	0.029242\\
2.4903	0.030003\\
2.4931	0.030781\\
2.4959	0.031577\\
2.4987	0.032391\\
2.5015	0.033224\\
2.5044	0.034075\\
2.5072	0.034945\\
2.51	0.035834\\
2.5128	0.036743\\
2.5156	0.037672\\
2.5185	0.038621\\
2.5213	0.039592\\
2.5241	0.040583\\
2.5269	0.041595\\
2.5297	0.04263\\
2.5326	0.043686\\
2.5354	0.044765\\
2.5382	0.045867\\
2.541	0.046992\\
2.5438	0.048141\\
2.5466	0.049313\\
2.5495	0.05051\\
2.5523	0.051732\\
2.5551	0.052979\\
2.5579	0.054251\\
2.5607	0.05555\\
2.5636	0.056874\\
2.5664	0.058225\\
2.5692	0.059603\\
2.572	0.061009\\
2.5748	0.062442\\
2.5777	0.063904\\
2.5805	0.065394\\
2.5833	0.066913\\
2.5861	0.068462\\
2.5889	0.070041\\
2.5917	0.071649\\
2.5946	0.073289\\
2.5974	0.074959\\
2.6002	0.076661\\
2.603	0.078394\\
2.6058	0.08016\\
2.6087	0.081959\\
2.6115	0.08379\\
2.6143	0.085655\\
2.6171	0.087554\\
2.6199	0.089487\\
2.6228	0.091455\\
2.6256	0.093457\\
2.6284	0.095496\\
2.6312	0.09757\\
2.634	0.09968\\
2.6368	0.10183\\
2.6397	0.10401\\
2.6425	0.10623\\
2.6453	0.10849\\
2.6481	0.11079\\
2.6509	0.11312\\
2.6538	0.1155\\
2.6566	0.11791\\
2.6594	0.12037\\
2.6622	0.12286\\
2.665	0.1254\\
2.6679	0.12797\\
2.6707	0.13059\\
2.6735	0.13324\\
2.6763	0.13594\\
2.6791	0.13869\\
2.682	0.14147\\
2.6848	0.1443\\
2.6876	0.14717\\
2.6904	0.15008\\
2.6932	0.15304\\
2.696	0.15604\\
2.6989	0.15909\\
2.7017	0.16218\\
2.7045	0.16532\\
2.7073	0.1685\\
2.7101	0.17173\\
2.713	0.175\\
2.7158	0.17832\\
2.7186	0.18168\\
2.7214	0.1851\\
2.7242	0.18856\\
2.7271	0.19206\\
2.7299	0.19562\\
2.7327	0.19922\\
2.7355	0.20287\\
2.7383	0.20656\\
2.7411	0.21031\\
2.744	0.2141\\
2.7468	0.21794\\
2.7496	0.22183\\
2.7524	0.22577\\
2.7552	0.22976\\
2.7581	0.2338\\
2.7609	0.23788\\
2.7637	0.24202\\
2.7665	0.2462\\
2.7693	0.25043\\
2.7722	0.25471\\
2.775	0.25905\\
2.7778	0.26343\\
2.7806	0.26786\\
2.7834	0.27234\\
2.7862	0.27686\\
2.7891	0.28144\\
2.7919	0.28607\\
2.7947	0.29075\\
2.7975	0.29547\\
2.8003	0.30025\\
2.8032	0.30507\\
2.806	0.30994\\
2.8088	0.31486\\
2.8116	0.31983\\
2.8144	0.32485\\
2.8173	0.32991\\
2.8201	0.33503\\
2.8229	0.34019\\
2.8257	0.34539\\
2.8285	0.35065\\
2.8313	0.35595\\
2.8342	0.3613\\
2.837	0.36669\\
2.8398	0.37213\\
2.8426	0.37762\\
2.8454	0.38315\\
2.8483	0.38873\\
2.8511	0.39435\\
2.8539	0.40001\\
2.8567	0.40572\\
2.8595	0.41147\\
2.8624	0.41727\\
2.8652	0.4231\\
2.868	0.42898\\
2.8708	0.4349\\
2.8736	0.44086\\
2.8764	0.44686\\
2.8793	0.45289\\
2.8821	0.45897\\
2.8849	0.46509\\
2.8877	0.47124\\
2.8905	0.47743\\
2.8934	0.48366\\
2.8962	0.48992\\
2.899	0.49621\\
2.9018	0.50254\\
2.9046	0.50891\\
2.9075	0.5153\\
2.9103	0.52173\\
2.9131	0.52819\\
2.9159	0.53467\\
2.9187	0.54119\\
2.9215	0.54774\\
2.9244	0.55431\\
2.9272	0.56091\\
2.93	0.56753\\
2.9328	0.57418\\
2.9356	0.58085\\
2.9385	0.58755\\
2.9413	0.59427\\
2.9441	0.601\\
2.9469	0.60776\\
2.9497	0.61454\\
2.9526	0.62133\\
2.9554	0.62814\\
2.9582	0.63496\\
2.961	0.6418\\
2.9638	0.64865\\
2.9666	0.65551\\
2.9695	0.66239\\
2.9723	0.66927\\
2.9751	0.67616\\
2.9779	0.68306\\
2.9807	0.68996\\
2.9836	0.69687\\
2.9864	0.70378\\
2.9892	0.71069\\
2.992	0.71761\\
2.9948	0.72452\\
2.9977	0.73143\\
3.0005	0.73834\\
3.0033	0.74524\\
3.0061	0.75214\\
3.0089	0.75903\\
3.0117	0.76591\\
3.0146	0.77278\\
3.0174	0.77964\\
3.0202	0.78648\\
3.023	0.79331\\
3.0258	0.80013\\
3.0287	0.80692\\
3.0315	0.8137\\
3.0343	0.82046\\
3.0371	0.8272\\
3.0399	0.83391\\
3.0428	0.8406\\
3.0456	0.84726\\
3.0484	0.8539\\
3.0512	0.86051\\
3.054	0.86708\\
3.0569	0.87363\\
3.0597	0.88014\\
3.0625	0.88662\\
3.0653	0.89306\\
3.0681	0.89946\\
3.0709	0.90582\\
3.0738	0.91214\\
3.0766	0.91842\\
3.0794	0.92466\\
3.0822	0.93085\\
3.085	0.937\\
3.0879	0.94309\\
3.0907	0.94914\\
3.0935	0.95514\\
3.0963	0.96108\\
3.0991	0.96697\\
3.102	0.97281\\
3.1048	0.97859\\
3.1076	0.98431\\
3.1104	0.98997\\
3.1132	0.99557\\
3.116	1.0011\\
3.1189	1.0066\\
3.1217	1.012\\
3.1245	1.0173\\
3.1273	1.0226\\
3.1301	1.0278\\
3.133	1.033\\
3.1358	1.038\\
3.1386	1.043\\
3.1414	1.048\\
3.1442	1.0528\\
3.1471	1.0576\\
3.1499	1.0623\\
3.1527	1.0669\\
3.1555	1.0714\\
3.1583	1.0759\\
3.1611	1.0803\\
3.164	1.0845\\
3.1668	1.0887\\
3.1696	1.0929\\
3.1724	1.0969\\
3.1752	1.1009\\
3.1781	1.1047\\
3.1809	1.1085\\
3.1837	1.1122\\
3.1865	1.1158\\
3.1893	1.1193\\
3.1922	1.1227\\
3.195	1.126\\
3.1978	1.1292\\
3.2006	1.1323\\
3.2034	1.1354\\
3.2062	1.1383\\
3.2091	1.1412\\
3.2119	1.1439\\
3.2147	1.1465\\
3.2175	1.1491\\
3.2203	1.1515\\
3.2232	1.1539\\
3.226	1.1561\\
3.2288	1.1583\\
3.2316	1.1603\\
3.2344	1.1623\\
3.2373	1.1641\\
3.2401	1.1658\\
3.2429	1.1675\\
3.2457	1.169\\
3.2485	1.1704\\
3.2513	1.1718\\
3.2542	1.173\\
3.257	1.1741\\
3.2598	1.1751\\
3.2626	1.176\\
3.2654	1.1768\\
3.2683	1.1775\\
3.2711	1.1781\\
3.2739	1.1786\\
3.2767	1.179\\
3.2795	1.1793\\
3.2824	1.1795\\
3.2852	1.1796\\
3.288	1.1795\\
3.2908	1.1794\\
3.2936	1.1792\\
3.2964	1.1788\\
3.2993	1.1784\\
3.3021	1.1779\\
3.3049	1.1772\\
3.3077	1.1765\\
3.3105	1.1756\\
3.3134	1.1747\\
3.3162	1.1737\\
3.319	1.1725\\
3.3218	1.1713\\
3.3246	1.1699\\
3.3275	1.1685\\
3.3303	1.167\\
3.3331	1.1653\\
3.3359	1.1636\\
3.3387	1.1618\\
3.3415	1.1599\\
3.3444	1.1578\\
3.3472	1.1557\\
3.35	1.1535\\
3.3528	1.1513\\
3.3556	1.1489\\
3.3585	1.1464\\
3.3613	1.1439\\
3.3641	1.1412\\
3.3669	1.1385\\
3.3697	1.1357\\
3.3726	1.1328\\
3.3754	1.1298\\
3.3782	1.1267\\
3.381	1.1236\\
3.3838	1.1203\\
3.3866	1.117\\
3.3895	1.1136\\
3.3923	1.1102\\
3.3951	1.1066\\
3.3979	1.103\\
3.4007	1.0993\\
3.4036	1.0955\\
3.4064	1.0917\\
3.4092	1.0878\\
3.412	1.0838\\
3.4148	1.0798\\
3.4177	1.0757\\
3.4205	1.0715\\
3.4233	1.0673\\
3.4261	1.063\\
3.4289	1.0586\\
3.4318	1.0542\\
3.4346	1.0497\\
3.4374	1.0451\\
3.4402	1.0405\\
3.443	1.0359\\
3.4458	1.0312\\
3.4487	1.0264\\
3.4515	1.0216\\
3.4543	1.0167\\
3.4571	1.0118\\
3.4599	1.0068\\
3.4628	1.0018\\
3.4656	0.99679\\
3.4684	0.99169\\
3.4712	0.98655\\
3.474	0.98137\\
3.4769	0.97616\\
3.4797	0.9709\\
3.4825	0.9656\\
3.4853	0.96027\\
3.4881	0.95491\\
3.4909	0.94951\\
3.4938	0.94408\\
3.4966	0.93862\\
3.4994	0.93314\\
3.5022	0.92762\\
3.505	0.92208\\
3.5079	0.91651\\
3.5107	0.91092\\
3.5135	0.90531\\
3.5163	0.89967\\
3.5191	0.89402\\
3.522	0.88834\\
3.5248	0.88265\\
3.5276	0.87695\\
3.5304	0.87123\\
3.5332	0.86549\\
3.536	0.85975\\
3.5389	0.85399\\
3.5417	0.84822\\
3.5445	0.84244\\
3.5473	0.83666\\
3.5501	0.83087\\
3.553	0.82508\\
3.5558	0.81928\\
3.5586	0.81348\\
3.5614	0.80767\\
3.5642	0.80187\\
3.5671	0.79607\\
3.5699	0.79027\\
3.5727	0.78447\\
3.5755	0.77868\\
3.5783	0.77289\\
3.5811	0.76711\\
3.584	0.76133\\
3.5868	0.75557\\
3.5896	0.74981\\
3.5924	0.74406\\
3.5952	0.73832\\
3.5981	0.7326\\
3.6009	0.72689\\
3.6037	0.72119\\
3.6065	0.71551\\
3.6093	0.70985\\
3.6122	0.7042\\
3.615	0.69856\\
3.6178	0.69295\\
3.6206	0.68736\\
3.6234	0.68178\\
3.6262	0.67623\\
3.6291	0.6707\\
3.6319	0.66519\\
3.6347	0.6597\\
3.6375	0.65424\\
3.6403	0.64881\\
3.6432	0.64339\\
3.646	0.63801\\
3.6488	0.63265\\
3.6516	0.62731\\
3.6544	0.62201\\
3.6573	0.61673\\
3.6601	0.61149\\
3.6629	0.60627\\
3.6657	0.60108\\
3.6685	0.59592\\
3.6713	0.5908\\
3.6742	0.5857\\
3.677	0.58064\\
3.6798	0.57561\\
3.6826	0.57061\\
3.6854	0.56565\\
3.6883	0.56072\\
3.6911	0.55583\\
3.6939	0.55096\\
3.6967	0.54614\\
3.6995	0.54135\\
3.7024	0.53659\\
3.7052	0.53187\\
3.708	0.52719\\
3.7108	0.52254\\
3.7136	0.51793\\
3.7164	0.51336\\
3.7193	0.50882\\
3.7221	0.50432\\
3.7249	0.49986\\
3.7277	0.49543\\
3.7305	0.49105\\
3.7334	0.4867\\
3.7362	0.48239\\
3.739	0.47812\\
3.7418	0.47388\\
3.7446	0.46969\\
3.7475	0.46553\\
3.7503	0.46141\\
3.7531	0.45733\\
3.7559	0.45329\\
3.7587	0.44928\\
3.7616	0.44532\\
3.7644	0.44139\\
3.7672	0.4375\\
3.77	0.43365\\
3.7728	0.42984\\
3.7756	0.42607\\
3.7785	0.42233\\
3.7813	0.41864\\
3.7841	0.41498\\
3.7869	0.41136\\
3.7897	0.40778\\
3.7926	0.40423\\
3.7954	0.40073\\
3.7982	0.39726\\
3.801	0.39382\\
3.8038	0.39043\\
3.8067	0.38707\\
3.8095	0.38375\\
3.8123	0.38047\\
3.8151	0.37722\\
3.8179	0.37401\\
3.8207	0.37083\\
3.8236	0.36769\\
3.8264	0.36459\\
3.8292	0.36152\\
3.832	0.35848\\
3.8348	0.35548\\
3.8377	0.35252\\
3.8405	0.34959\\
3.8433	0.34669\\
3.8461	0.34383\\
3.8489	0.341\\
3.8518	0.33821\\
3.8546	0.33544\\
3.8574	0.33271\\
3.8602	0.33001\\
3.863	0.32735\\
3.8658	0.32472\\
3.8687	0.32211\\
3.8715	0.31954\\
3.8743	0.317\\
3.8771	0.31449\\
3.8799	0.31201\\
3.8828	0.30956\\
3.8856	0.30714\\
3.8884	0.30475\\
3.8912	0.30239\\
3.894	0.30005\\
3.8969	0.29775\\
3.8997	0.29547\\
3.9025	0.29322\\
3.9053	0.29099\\
3.9081	0.2888\\
3.9109	0.28663\\
3.9138	0.28448\\
3.9166	0.28236\\
3.9194	0.28027\\
3.9222	0.2782\\
3.925	0.27616\\
3.9279	0.27414\\
3.9307	0.27214\\
3.9335	0.27017\\
3.9363	0.26822\\
3.9391	0.26629\\
3.942	0.26439\\
3.9448	0.2625\\
3.9476	0.26064\\
3.9504	0.2588\\
3.9532	0.25699\\
3.956	0.25519\\
3.9589	0.25341\\
3.9617	0.25165\\
3.9645	0.24992\\
3.9673	0.2482\\
3.9701	0.2465\\
3.973	0.24482\\
3.9758	0.24315\\
3.9786	0.24151\\
3.9814	0.23988\\
3.9842	0.23827\\
3.9871	0.23668\\
3.9899	0.2351\\
3.9927	0.23354\\
3.9955	0.23199\\
3.9983	0.23047\\
4.0011	0.22895\\
4.004	0.22745\\
4.0068	0.22597\\
4.0096	0.2245\\
4.0124	0.22305\\
4.0152	0.22161\\
4.0181	0.22018\\
4.0209	0.21877\\
4.0237	0.21736\\
4.0265	0.21598\\
4.0293	0.2146\\
4.0322	0.21324\\
4.035	0.21189\\
4.0378	0.21055\\
4.0406	0.20922\\
4.0434	0.2079\\
4.0462	0.2066\\
4.0491	0.2053\\
4.0519	0.20402\\
4.0547	0.20275\\
4.0575	0.20148\\
4.0603	0.20023\\
4.0632	0.19898\\
4.066	0.19775\\
4.0688	0.19652\\
4.0716	0.19531\\
4.0744	0.1941\\
4.0773	0.1929\\
4.0801	0.19171\\
4.0829	0.19053\\
4.0857	0.18936\\
4.0885	0.18819\\
4.0913	0.18703\\
4.0942	0.18588\\
4.097	0.18474\\
4.0998	0.1836\\
4.1026	0.18248\\
4.1054	0.18135\\
4.1083	0.18024\\
4.1111	0.17913\\
4.1139	0.17803\\
4.1167	0.17693\\
4.1195	0.17585\\
4.1224	0.17476\\
4.1252	0.17369\\
4.128	0.17261\\
4.1308	0.17155\\
4.1336	0.17049\\
4.1365	0.16943\\
4.1393	0.16838\\
4.1421	0.16734\\
4.1449	0.1663\\
4.1477	0.16527\\
4.1505	0.16424\\
4.1534	0.16321\\
4.1562	0.1622\\
4.159	0.16118\\
4.1618	0.16017\\
4.1646	0.15916\\
4.1675	0.15816\\
4.1703	0.15717\\
4.1731	0.15617\\
4.1759	0.15518\\
4.1787	0.1542\\
4.1816	0.15322\\
4.1844	0.15224\\
4.1872	0.15127\\
4.19	0.1503\\
4.1928	0.14934\\
4.1956	0.14837\\
4.1985	0.14742\\
4.2013	0.14646\\
4.2041	0.14551\\
4.2069	0.14456\\
4.2097	0.14362\\
4.2126	0.14268\\
4.2154	0.14174\\
4.2182	0.14081\\
4.221	0.13988\\
4.2238	0.13895\\
4.2267	0.13803\\
4.2295	0.13711\\
4.2323	0.13619\\
4.2351	0.13527\\
4.2379	0.13436\\
4.2407	0.13345\\
4.2436	0.13255\\
4.2464	0.13165\\
4.2492	0.13075\\
4.252	0.12985\\
4.2548	0.12896\\
4.2577	0.12807\\
4.2605	0.12718\\
4.2633	0.1263\\
4.2661	0.12541\\
4.2689	0.12454\\
4.2718	0.12366\\
4.2746	0.12279\\
4.2774	0.12192\\
4.2802	0.12105\\
4.283	0.12019\\
4.2858	0.11932\\
4.2887	0.11847\\
4.2915	0.11761\\
4.2943	0.11676\\
4.2971	0.11591\\
4.2999	0.11506\\
4.3028	0.11422\\
4.3056	0.11338\\
4.3084	0.11254\\
4.3112	0.1117\\
4.314	0.11087\\
4.3169	0.11004\\
4.3197	0.10921\\
4.3225	0.10839\\
4.3253	0.10757\\
4.3281	0.10675\\
4.3309	0.10594\\
4.3338	0.10512\\
4.3366	0.10432\\
4.3394	0.10351\\
4.3422	0.10271\\
4.345	0.10191\\
4.3479	0.10111\\
4.3507	0.10032\\
4.3535	0.099524\\
4.3563	0.098736\\
4.3591	0.097951\\
4.362	0.097169\\
4.3648	0.096391\\
4.3676	0.095615\\
4.3704	0.094843\\
4.3732	0.094073\\
4.376	0.093307\\
4.3789	0.092544\\
4.3817	0.091784\\
4.3845	0.091027\\
4.3873	0.090274\\
4.3901	0.089524\\
4.393	0.088777\\
4.3958	0.088033\\
4.3986	0.087292\\
4.4014	0.086555\\
4.4042	0.085821\\
4.4071	0.085091\\
4.4099	0.084363\\
4.4127	0.08364\\
4.4155	0.082919\\
4.4183	0.082202\\
4.4211	0.081488\\
4.424	0.080778\\
4.4268	0.080071\\
4.4296	0.079367\\
4.4324	0.078667\\
4.4352	0.07797\\
4.4381	0.077277\\
4.4409	0.076587\\
4.4437	0.075901\\
4.4465	0.075219\\
4.4493	0.074539\\
4.4522	0.073864\\
4.455	0.073192\\
4.4578	0.072523\\
4.4606	0.071858\\
4.4634	0.071197\\
4.4663	0.070539\\
4.4691	0.069885\\
4.4719	0.069235\\
4.4747	0.068588\\
4.4775	0.067944\\
4.4803	0.067305\\
4.4832	0.066669\\
4.486	0.066037\\
4.4888	0.065408\\
4.4916	0.064783\\
4.4944	0.064162\\
4.4973	0.063544\\
4.5001	0.06293\\
4.5029	0.06232\\
4.5057	0.061714\\
4.5085	0.061111\\
4.5114	0.060512\\
4.5142	0.059917\\
4.517	0.059326\\
4.5198	0.058738\\
4.5226	0.058154\\
4.5254	0.057573\\
4.5283	0.056997\\
4.5311	0.056424\\
4.5339	0.055855\\
4.5367	0.05529\\
4.5395	0.054728\\
4.5424	0.054171\\
4.5452	0.053617\\
4.548	0.053066\\
4.5508	0.05252\\
4.5536	0.051977\\
4.5565	0.051438\\
4.5593	0.050903\\
4.5621	0.050371\\
4.5649	0.049843\\
4.5677	0.049319\\
4.5705	0.048799\\
4.5734	0.048282\\
4.5762	0.04777\\
4.579	0.04726\\
4.5818	0.046755\\
4.5846	0.046253\\
4.5875	0.045755\\
4.5903	0.045261\\
4.5931	0.04477\\
4.5959	0.044284\\
4.5987	0.0438\\
4.6016	0.043321\\
4.6044	0.042845\\
4.6072	0.042373\\
4.61	0.041904\\
4.6128	0.041439\\
4.6156	0.040978\\
4.6185	0.040521\\
4.6213	0.040067\\
4.6241	0.039616\\
4.6269	0.03917\\
4.6297	0.038727\\
4.6326	0.038287\\
4.6354	0.037851\\
4.6382	0.037419\\
4.641	0.03699\\
4.6438	0.036565\\
4.6467	0.036143\\
4.6495	0.035725\\
4.6523	0.03531\\
4.6551	0.034899\\
4.6579	0.034491\\
4.6607	0.034087\\
4.6636	0.033687\\
4.6664	0.033289\\
4.6692	0.032896\\
4.672	0.032505\\
4.6748	0.032118\\
4.6777	0.031735\\
4.6805	0.031355\\
4.6833	0.030978\\
4.6861	0.030605\\
4.6889	0.030235\\
4.6918	0.029868\\
4.6946	0.029505\\
4.6974	0.029145\\
4.7002	0.028789\\
4.703	0.028435\\
4.7058	0.028085\\
4.7087	0.027739\\
4.7115	0.027395\\
4.7143	0.027055\\
4.7171	0.026718\\
4.7199	0.026384\\
4.7228	0.026053\\
4.7256	0.025726\\
4.7284	0.025401\\
4.7312	0.02508\\
4.734	0.024762\\
4.7369	0.024447\\
4.7397	0.024135\\
4.7425	0.023827\\
4.7453	0.023521\\
4.7481	0.023218\\
4.7509	0.022919\\
4.7538	0.022622\\
4.7566	0.022328\\
4.7594	0.022038\\
4.7622	0.02175\\
4.765	0.021465\\
4.7679	0.021183\\
4.7707	0.020905\\
4.7735	0.020629\\
4.7763	0.020355\\
4.7791	0.020085\\
4.782	0.019818\\
4.7848	0.019553\\
4.7876	0.019291\\
4.7904	0.019032\\
4.7932	0.018776\\
4.796	0.018523\\
4.7989	0.018272\\
4.8017	0.018024\\
4.8045	0.017778\\
4.8073	0.017536\\
4.8101	0.017296\\
4.813	0.017058\\
4.8158	0.016824\\
4.8186	0.016591\\
4.8214	0.016362\\
4.8242	0.016135\\
4.8271	0.01591\\
4.8299	0.015688\\
4.8327	0.015469\\
4.8355	0.015252\\
4.8383	0.015038\\
4.8412	0.014826\\
4.844	0.014616\\
4.8468	0.014409\\
4.8496	0.014204\\
4.8524	0.014002\\
4.8552	0.013802\\
4.8581	0.013604\\
4.8609	0.013409\\
4.8637	0.013216\\
4.8665	0.013025\\
4.8693	0.012836\\
4.8722	0.01265\\
4.875	0.012466\\
4.8778	0.012284\\
4.8806	0.012105\\
4.8834	0.011927\\
4.8863	0.011752\\
4.8891	0.011579\\
4.8919	0.011408\\
4.8947	0.011239\\
4.8975	0.011072\\
4.9003	0.010907\\
4.9032	0.010745\\
4.906	0.010584\\
4.9088	0.010425\\
4.9116	0.010269\\
4.9144	0.010114\\
4.9173	0.0099611\\
4.9201	0.0098103\\
4.9229	0.0096613\\
4.9257	0.0095143\\
4.9285	0.0093692\\
4.9314	0.0092259\\
4.9342	0.0090845\\
4.937	0.0089449\\
4.9398	0.0088071\\
4.9426	0.0086711\\
4.9454	0.0085369\\
4.9483	0.0084044\\
4.9511	0.0082737\\
4.9539	0.0081447\\
4.9567	0.0080175\\
4.9595	0.0078919\\
4.9624	0.007768\\
4.9652	0.0076457\\
4.968	0.0075251\\
4.9708	0.0074061\\
4.9736	0.0072887\\
4.9765	0.0071729\\
4.9793	0.0070587\\
4.9821	0.006946\\
4.9849	0.0068349\\
4.9877	0.0067252\\
4.9905	0.0066171\\
4.9934	0.0065105\\
4.9962	0.0064054\\
4.999	0.0063017\\
5.0018	0.0061995\\
5.0046	0.0060987\\
5.0075	0.0059993\\
5.0103	0.0059013\\
5.0131	0.0058047\\
5.0159	0.0057094\\
5.0187	0.0056155\\
5.0216	0.005523\\
5.0244	0.0054317\\
5.0272	0.0053418\\
5.03	0.0052531\\
5.0328	0.0051658\\
5.0356	0.0050797\\
5.0385	0.0049948\\
5.0413	0.0049112\\
5.0441	0.0048287\\
5.0469	0.0047475\\
5.0497	0.0046675\\
5.0526	0.0045887\\
5.0554	0.004511\\
5.0582	0.0044344\\
5.061	0.004359\\
5.0638	0.0042848\\
5.0667	0.0042116\\
5.0695	0.0041395\\
5.0723	0.0040685\\
5.0751	0.0039986\\
5.0779	0.0039297\\
5.0807	0.0038618\\
5.0836	0.003795\\
5.0864	0.0037292\\
5.0892	0.0036644\\
5.092	0.0036006\\
5.0948	0.0035378\\
5.0977	0.0034759\\
5.1005	0.003415\\
5.1033	0.003355\\
5.1061	0.003296\\
5.1089	0.0032379\\
5.1118	0.0031806\\
5.1146	0.0031243\\
5.1174	0.0030689\\
5.1202	0.0030143\\
5.123	0.0029606\\
5.1258	0.0029077\\
5.1287	0.0028557\\
5.1315	0.0028045\\
5.1343	0.0027541\\
5.1371	0.0027045\\
5.1399	0.0026557\\
5.1428	0.0026077\\
5.1456	0.0025604\\
5.1484	0.0025139\\
5.1512	0.0024682\\
5.154	0.0024232\\
5.1569	0.0023789\\
5.1597	0.0023354\\
5.1625	0.0022926\\
5.1653	0.0022504\\
5.1681	0.002209\\
5.171	0.0021682\\
5.1738	0.0021281\\
5.1766	0.0020887\\
5.1794	0.00205\\
5.1822	0.0020118\\
5.185	0.0019743\\
5.1879	0.0019375\\
5.1907	0.0019012\\
5.1935	0.0018656\\
5.1963	0.0018305\\
5.1991	0.0017961\\
5.202	0.0017622\\
5.2048	0.0017289\\
5.2048	0.0016581\\
5.202	0.0016904\\
5.1991	0.0017233\\
5.1963	0.0017567\\
5.1935	0.0017907\\
5.1907	0.0018253\\
5.1879	0.0018605\\
5.185	0.0018963\\
5.1822	0.0019327\\
5.1794	0.0019698\\
5.1766	0.0020074\\
5.1738	0.0020458\\
5.171	0.0020847\\
5.1681	0.0021244\\
5.1653	0.0021647\\
5.1625	0.0022056\\
5.1597	0.0022473\\
5.1569	0.0022897\\
5.154	0.0023327\\
5.1512	0.0023765\\
5.1484	0.0024211\\
5.1456	0.0024663\\
5.1428	0.0025124\\
5.1399	0.0025591\\
5.1371	0.0026067\\
5.1343	0.002655\\
5.1315	0.0027042\\
5.1287	0.0027541\\
5.1258	0.0028048\\
5.123	0.0028564\\
5.1202	0.0029088\\
5.1174	0.0029621\\
5.1146	0.0030162\\
5.1118	0.0030712\\
5.1089	0.003127\\
5.1061	0.0031838\\
5.1033	0.0032415\\
5.1005	0.0033001\\
5.0977	0.0033596\\
5.0948	0.00342\\
5.092	0.0034815\\
5.0892	0.0035438\\
5.0864	0.0036072\\
5.0836	0.0036716\\
5.0807	0.0037369\\
5.0779	0.0038033\\
5.0751	0.0038707\\
5.0723	0.0039392\\
5.0695	0.0040087\\
5.0667	0.0040793\\
5.0638	0.004151\\
5.061	0.0042237\\
5.0582	0.0042976\\
5.0554	0.0043726\\
5.0526	0.0044488\\
5.0497	0.0045261\\
5.0469	0.0046045\\
5.0441	0.0046842\\
5.0413	0.004765\\
5.0385	0.004847\\
5.0356	0.0049303\\
5.0328	0.0050148\\
5.03	0.0051006\\
5.0272	0.0051876\\
5.0244	0.0052759\\
5.0216	0.0053656\\
5.0187	0.0054565\\
5.0159	0.0055488\\
5.0131	0.0056424\\
5.0103	0.0057373\\
5.0075	0.0058336\\
5.0046	0.0059314\\
5.0018	0.0060305\\
4.999	0.0061311\\
4.9962	0.0062331\\
4.9934	0.0063365\\
4.9905	0.0064414\\
4.9877	0.0065478\\
4.9849	0.0066557\\
4.9821	0.0067652\\
4.9793	0.0068761\\
4.9765	0.0069886\\
4.9736	0.0071027\\
4.9708	0.0072184\\
4.968	0.0073357\\
4.9652	0.0074546\\
4.9624	0.0075751\\
4.9595	0.0076973\\
4.9567	0.0078211\\
4.9539	0.0079467\\
4.9511	0.0080739\\
4.9483	0.0082029\\
4.9454	0.0083336\\
4.9426	0.0084661\\
4.9398	0.0086003\\
4.937	0.0087364\\
4.9342	0.0088742\\
4.9314	0.0090139\\
4.9285	0.0091555\\
4.9257	0.0092989\\
4.9229	0.0094441\\
4.9201	0.0095913\\
4.9173	0.0097404\\
4.9144	0.0098915\\
4.9116	0.010044\\
4.9088	0.010199\\
4.906	0.010356\\
4.9032	0.010515\\
4.9003	0.010676\\
4.8975	0.01084\\
4.8947	0.011005\\
4.8919	0.011172\\
4.8891	0.011341\\
4.8863	0.011513\\
4.8834	0.011686\\
4.8806	0.011862\\
4.8778	0.01204\\
4.875	0.01222\\
4.8722	0.012402\\
4.8693	0.012587\\
4.8665	0.012774\\
4.8637	0.012963\\
4.8609	0.013155\\
4.8581	0.013348\\
4.8552	0.013544\\
4.8524	0.013743\\
4.8496	0.013944\\
4.8468	0.014147\\
4.844	0.014352\\
4.8412	0.01456\\
4.8383	0.014771\\
4.8355	0.014984\\
4.8327	0.015199\\
4.8299	0.015417\\
4.8271	0.015638\\
4.8242	0.015861\\
4.8214	0.016086\\
4.8186	0.016315\\
4.8158	0.016545\\
4.813	0.016779\\
4.8101	0.017015\\
4.8073	0.017253\\
4.8045	0.017495\\
4.8017	0.017739\\
4.7989	0.017986\\
4.796	0.018235\\
4.7932	0.018487\\
4.7904	0.018743\\
4.7876	0.019\\
4.7848	0.019261\\
4.782	0.019524\\
4.7791	0.019791\\
4.7763	0.02006\\
4.7735	0.020332\\
4.7707	0.020607\\
4.7679	0.020885\\
4.765	0.021165\\
4.7622	0.021449\\
4.7594	0.021736\\
4.7566	0.022026\\
4.7538	0.022318\\
4.7509	0.022614\\
4.7481	0.022913\\
4.7453	0.023215\\
4.7425	0.023519\\
4.7397	0.023827\\
4.7369	0.024139\\
4.734	0.024453\\
4.7312	0.02477\\
4.7284	0.025091\\
4.7256	0.025414\\
4.7228	0.025741\\
4.7199	0.026071\\
4.7171	0.026404\\
4.7143	0.026741\\
4.7115	0.027081\\
4.7087	0.027424\\
4.7058	0.02777\\
4.703	0.028119\\
4.7002	0.028472\\
4.6974	0.028828\\
4.6946	0.029188\\
4.6918	0.029551\\
4.6889	0.029917\\
4.6861	0.030287\\
4.6833	0.03066\\
4.6805	0.031036\\
4.6777	0.031416\\
4.6748	0.031799\\
4.672	0.032186\\
4.6692	0.032576\\
4.6664	0.032969\\
4.6636	0.033366\\
4.6607	0.033767\\
4.6579	0.034171\\
4.6551	0.034578\\
4.6523	0.034989\\
4.6495	0.035404\\
4.6467	0.035822\\
4.6438	0.036244\\
4.641	0.036669\\
4.6382	0.037097\\
4.6354	0.03753\\
4.6326	0.037965\\
4.6297	0.038405\\
4.6269	0.038848\\
4.6241	0.039294\\
4.6213	0.039745\\
4.6185	0.040198\\
4.6156	0.040656\\
4.6128	0.041117\\
4.61	0.041581\\
4.6072	0.04205\\
4.6044	0.042522\\
4.6016	0.042997\\
4.5987	0.043476\\
4.5959	0.043959\\
4.5931	0.044445\\
4.5903	0.044936\\
4.5875	0.045429\\
4.5846	0.045927\\
4.5818	0.046428\\
4.579	0.046932\\
4.5762	0.047441\\
4.5734	0.047953\\
4.5705	0.048468\\
4.5677	0.048988\\
4.5649	0.049511\\
4.5621	0.050037\\
4.5593	0.050567\\
4.5565	0.051101\\
4.5536	0.051639\\
4.5508	0.05218\\
4.548	0.052725\\
4.5452	0.053273\\
4.5424	0.053825\\
4.5395	0.054381\\
4.5367	0.05494\\
4.5339	0.055503\\
4.5311	0.05607\\
4.5283	0.05664\\
4.5254	0.057213\\
4.5226	0.057791\\
4.5198	0.058372\\
4.517	0.058956\\
4.5142	0.059544\\
4.5114	0.060136\\
4.5085	0.060731\\
4.5057	0.06133\\
4.5029	0.061932\\
4.5001	0.062538\\
4.4973	0.063147\\
4.4944	0.06376\\
4.4916	0.064377\\
4.4888	0.064996\\
4.486	0.06562\\
4.4832	0.066247\\
4.4803	0.066877\\
4.4775	0.067511\\
4.4747	0.068148\\
4.4719	0.068789\\
4.4691	0.069433\\
4.4663	0.070081\\
4.4634	0.070732\\
4.4606	0.071386\\
4.4578	0.072044\\
4.455	0.072706\\
4.4522	0.07337\\
4.4493	0.074038\\
4.4465	0.074709\\
4.4437	0.075384\\
4.4409	0.076062\\
4.4381	0.076744\\
4.4352	0.077428\\
4.4324	0.078116\\
4.4296	0.078807\\
4.4268	0.079502\\
4.424	0.0802\\
4.4211	0.080901\\
4.4183	0.081605\\
4.4155	0.082312\\
4.4127	0.083023\\
4.4099	0.083737\\
4.4071	0.084454\\
4.4042	0.085174\\
4.4014	0.085897\\
4.3986	0.086624\\
4.3958	0.087354\\
4.393	0.088086\\
4.3901	0.088822\\
4.3873	0.089561\\
4.3845	0.090303\\
4.3817	0.091048\\
4.3789	0.091796\\
4.376	0.092548\\
4.3732	0.093302\\
4.3704	0.094059\\
4.3676	0.094819\\
4.3648	0.095582\\
4.362	0.096349\\
4.3591	0.097118\\
4.3563	0.09789\\
4.3535	0.098665\\
4.3507	0.099443\\
4.3479	0.10022\\
4.345	0.10101\\
4.3422	0.10179\\
4.3394	0.10258\\
4.3366	0.10338\\
4.3338	0.10417\\
4.3309	0.10497\\
4.3281	0.10577\\
4.3253	0.10658\\
4.3225	0.10738\\
4.3197	0.10819\\
4.3169	0.109\\
4.314	0.10982\\
4.3112	0.11064\\
4.3084	0.11146\\
4.3056	0.11228\\
4.3028	0.11311\\
4.2999	0.11394\\
4.2971	0.11477\\
4.2943	0.11561\\
4.2915	0.11645\\
4.2887	0.11729\\
4.2858	0.11813\\
4.283	0.11898\\
4.2802	0.11983\\
4.2774	0.12068\\
4.2746	0.12153\\
4.2718	0.12239\\
4.2689	0.12325\\
4.2661	0.12412\\
4.2633	0.12498\\
4.2605	0.12585\\
4.2577	0.12673\\
4.2548	0.1276\\
4.252	0.12848\\
4.2492	0.12936\\
4.2464	0.13025\\
4.2436	0.13113\\
4.2407	0.13202\\
4.2379	0.13292\\
4.2351	0.13382\\
4.2323	0.13472\\
4.2295	0.13562\\
4.2267	0.13653\\
4.2238	0.13744\\
4.221	0.13835\\
4.2182	0.13926\\
4.2154	0.14018\\
4.2126	0.14111\\
4.2097	0.14203\\
4.2069	0.14296\\
4.2041	0.1439\\
4.2013	0.14484\\
4.1985	0.14578\\
4.1956	0.14672\\
4.1928	0.14767\\
4.19	0.14862\\
4.1872	0.14958\\
4.1844	0.15054\\
4.1816	0.1515\\
4.1787	0.15247\\
4.1759	0.15344\\
4.1731	0.15442\\
4.1703	0.1554\\
4.1675	0.15639\\
4.1646	0.15738\\
4.1618	0.15837\\
4.159	0.15937\\
4.1562	0.16037\\
4.1534	0.16138\\
4.1505	0.1624\\
4.1477	0.16341\\
4.1449	0.16444\\
4.1421	0.16547\\
4.1393	0.1665\\
4.1365	0.16754\\
4.1336	0.16859\\
4.1308	0.16964\\
4.128	0.1707\\
4.1252	0.17176\\
4.1224	0.17283\\
4.1195	0.1739\\
4.1167	0.17498\\
4.1139	0.17607\\
4.1111	0.17717\\
4.1083	0.17827\\
4.1054	0.17938\\
4.1026	0.18049\\
4.0998	0.18162\\
4.097	0.18275\\
4.0942	0.18388\\
4.0913	0.18503\\
4.0885	0.18618\\
4.0857	0.18734\\
4.0829	0.18851\\
4.0801	0.18969\\
4.0773	0.19088\\
4.0744	0.19207\\
4.0716	0.19328\\
4.0688	0.19449\\
4.066	0.19571\\
4.0632	0.19694\\
4.0603	0.19819\\
4.0575	0.19944\\
4.0547	0.2007\\
4.0519	0.20197\\
4.0491	0.20326\\
4.0462	0.20455\\
4.0434	0.20586\\
4.0406	0.20717\\
4.0378	0.2085\\
4.035	0.20984\\
4.0322	0.21119\\
4.0293	0.21255\\
4.0265	0.21393\\
4.0237	0.21532\\
4.0209	0.21672\\
4.0181	0.21814\\
4.0152	0.21957\\
4.0124	0.22101\\
4.0096	0.22246\\
4.0068	0.22393\\
4.004	0.22542\\
4.0011	0.22692\\
3.9983	0.22844\\
3.9955	0.22997\\
3.9927	0.23151\\
3.9899	0.23307\\
3.9871	0.23465\\
3.9842	0.23625\\
3.9814	0.23786\\
3.9786	0.23949\\
3.9758	0.24113\\
3.973	0.2428\\
3.9701	0.24448\\
3.9673	0.24618\\
3.9645	0.2479\\
3.9617	0.24964\\
3.9589	0.25139\\
3.956	0.25317\\
3.9532	0.25497\\
3.9504	0.25678\\
3.9476	0.25862\\
3.9448	0.26048\\
3.942	0.26236\\
3.9391	0.26426\\
3.9363	0.26618\\
3.9335	0.26812\\
3.9307	0.27009\\
3.9279	0.27208\\
3.925	0.27409\\
3.9222	0.27613\\
3.9194	0.27819\\
3.9166	0.28027\\
3.9138	0.28238\\
3.9109	0.28451\\
3.9081	0.28667\\
3.9053	0.28885\\
3.9025	0.29106\\
3.8997	0.2933\\
3.8969	0.29556\\
3.894	0.29785\\
3.8912	0.30016\\
3.8884	0.30251\\
3.8856	0.30488\\
3.8828	0.30727\\
3.8799	0.3097\\
3.8771	0.31216\\
3.8743	0.31464\\
3.8715	0.31716\\
3.8687	0.3197\\
3.8658	0.32228\\
3.863	0.32488\\
3.8602	0.32752\\
3.8574	0.33018\\
3.8546	0.33288\\
3.8518	0.33561\\
3.8489	0.33837\\
3.8461	0.34117\\
3.8433	0.344\\
3.8405	0.34686\\
3.8377	0.34975\\
3.8348	0.35268\\
3.832	0.35564\\
3.8292	0.35863\\
3.8264	0.36166\\
3.8236	0.36473\\
3.8207	0.36783\\
3.8179	0.37096\\
3.8151	0.37413\\
3.8123	0.37733\\
3.8095	0.38058\\
3.8067	0.38385\\
3.8038	0.38717\\
3.801	0.39052\\
3.7982	0.3939\\
3.7954	0.39733\\
3.7926	0.40079\\
3.7897	0.40429\\
3.7869	0.40782\\
3.7841	0.4114\\
3.7813	0.41501\\
3.7785	0.41866\\
3.7756	0.42235\\
3.7728	0.42607\\
3.77	0.42984\\
3.7672	0.43364\\
3.7644	0.43748\\
3.7616	0.44136\\
3.7587	0.44528\\
3.7559	0.44924\\
3.7531	0.45323\\
3.7503	0.45727\\
3.7475	0.46134\\
3.7446	0.46545\\
3.7418	0.4696\\
3.739	0.47379\\
3.7362	0.47802\\
3.7334	0.48228\\
3.7305	0.48659\\
3.7277	0.49093\\
3.7249	0.49531\\
3.7221	0.49973\\
3.7193	0.50419\\
3.7164	0.50868\\
3.7136	0.51321\\
3.7108	0.51778\\
3.708	0.52239\\
3.7052	0.52704\\
3.7024	0.53172\\
3.6995	0.53643\\
3.6967	0.54119\\
3.6939	0.54598\\
3.6911	0.5508\\
3.6883	0.55566\\
3.6854	0.56056\\
3.6826	0.56549\\
3.6798	0.57045\\
3.677	0.57545\\
3.6742	0.58049\\
3.6713	0.58555\\
3.6685	0.59065\\
3.6657	0.59578\\
3.6629	0.60094\\
3.6601	0.60613\\
3.6573	0.61135\\
3.6544	0.61661\\
3.6516	0.62189\\
3.6488	0.6272\\
3.646	0.63254\\
3.6432	0.63791\\
3.6403	0.6433\\
3.6375	0.64872\\
3.6347	0.65417\\
3.6319	0.65964\\
3.6291	0.66514\\
3.6262	0.67066\\
3.6234	0.6762\\
3.6206	0.68177\\
3.6178	0.68735\\
3.615	0.69296\\
3.6122	0.69858\\
3.6093	0.70423\\
3.6065	0.70989\\
3.6037	0.71557\\
3.6009	0.72127\\
3.5981	0.72698\\
3.5952	0.7327\\
3.5924	0.73844\\
3.5896	0.74419\\
3.5868	0.74996\\
3.584	0.75573\\
3.5811	0.76151\\
3.5783	0.7673\\
3.5755	0.7731\\
3.5727	0.77891\\
3.5699	0.78472\\
3.5671	0.79053\\
3.5642	0.79635\\
3.5614	0.80216\\
3.5586	0.80798\\
3.5558	0.8138\\
3.553	0.81962\\
3.5501	0.82543\\
3.5473	0.83124\\
3.5445	0.83705\\
3.5417	0.84285\\
3.5389	0.84864\\
3.536	0.85442\\
3.5332	0.86019\\
3.5304	0.86595\\
3.5276	0.8717\\
3.5248	0.87743\\
3.522	0.88315\\
3.5191	0.88885\\
3.5163	0.89453\\
3.5135	0.9002\\
3.5107	0.90584\\
3.5079	0.91146\\
3.505	0.91706\\
3.5022	0.92263\\
3.4994	0.92817\\
3.4966	0.93369\\
3.4938	0.93918\\
3.4909	0.94464\\
3.4881	0.95007\\
3.4853	0.95546\\
3.4825	0.96082\\
3.4797	0.96615\\
3.4769	0.97143\\
3.474	0.97668\\
3.4712	0.98189\\
3.4684	0.98705\\
3.4656	0.99218\\
3.4628	0.99726\\
3.4599	1.0023\\
3.4571	1.0073\\
3.4543	1.0122\\
3.4515	1.0171\\
3.4487	1.0219\\
3.4458	1.0267\\
3.443	1.0314\\
3.4402	1.0361\\
3.4374	1.0407\\
3.4346	1.0453\\
3.4318	1.0498\\
3.4289	1.0542\\
3.4261	1.0586\\
3.4233	1.0629\\
3.4205	1.0672\\
3.4177	1.0714\\
3.4148	1.0755\\
3.412	1.0795\\
3.4092	1.0835\\
3.4064	1.0874\\
3.4036	1.0912\\
3.4007	1.095\\
3.3979	1.0987\\
3.3951	1.1023\\
3.3923	1.1058\\
3.3895	1.1093\\
3.3866	1.1126\\
3.3838	1.1159\\
3.381	1.1192\\
3.3782	1.1223\\
3.3754	1.1253\\
3.3726	1.1283\\
3.3697	1.1312\\
3.3669	1.134\\
3.3641	1.1367\\
3.3613	1.1393\\
3.3585	1.1418\\
3.3556	1.1443\\
3.3528	1.1466\\
3.35	1.1489\\
3.3472	1.151\\
3.3444	1.1531\\
3.3415	1.1551\\
3.3387	1.157\\
3.3359	1.1588\\
3.3331	1.1605\\
3.3303	1.1621\\
3.3275	1.1636\\
3.3246	1.165\\
3.3218	1.1663\\
3.319	1.1675\\
3.3162	1.1686\\
3.3134	1.1696\\
3.3105	1.1705\\
3.3077	1.1713\\
3.3049	1.172\\
3.3021	1.1726\\
3.2993	1.1731\\
3.2964	1.1735\\
3.2936	1.1738\\
3.2908	1.174\\
3.288	1.1741\\
3.2852	1.1741\\
3.2824	1.174\\
3.2795	1.1738\\
3.2767	1.1734\\
3.2739	1.173\\
3.2711	1.1725\\
3.2683	1.1719\\
3.2654	1.1711\\
3.2626	1.1703\\
3.2598	1.1694\\
3.257	1.1683\\
3.2542	1.1672\\
3.2513	1.1659\\
3.2485	1.1646\\
3.2457	1.1631\\
3.2429	1.1616\\
3.2401	1.1599\\
3.2373	1.1582\\
3.2344	1.1563\\
3.2316	1.1544\\
3.2288	1.1523\\
3.226	1.1502\\
3.2232	1.1479\\
3.2203	1.1455\\
3.2175	1.1431\\
3.2147	1.1405\\
3.2119	1.1379\\
3.2091	1.1351\\
3.2062	1.1323\\
3.2034	1.1294\\
3.2006	1.1263\\
3.1978	1.1232\\
3.195	1.12\\
3.1922	1.1167\\
3.1893	1.1133\\
3.1865	1.1098\\
3.1837	1.1062\\
3.1809	1.1025\\
3.1781	1.0987\\
3.1752	1.0949\\
3.1724	1.091\\
3.1696	1.0869\\
3.1668	1.0828\\
3.164	1.0786\\
3.1611	1.0744\\
3.1583	1.07\\
3.1555	1.0656\\
3.1527	1.0611\\
3.1499	1.0565\\
3.1471	1.0518\\
3.1442	1.047\\
3.1414	1.0422\\
3.1386	1.0373\\
3.1358	1.0324\\
3.133	1.0273\\
3.1301	1.0222\\
3.1273	1.017\\
3.1245	1.0118\\
3.1217	1.0065\\
3.1189	1.0011\\
3.116	0.99564\\
3.1132	0.99013\\
3.1104	0.98456\\
3.1076	0.97894\\
3.1048	0.97325\\
3.102	0.96751\\
3.0991	0.96171\\
3.0963	0.95585\\
3.0935	0.94995\\
3.0907	0.94399\\
3.0879	0.93798\\
3.085	0.93192\\
3.0822	0.92581\\
3.0794	0.91966\\
3.0766	0.91346\\
3.0738	0.90722\\
3.0709	0.90094\\
3.0681	0.89462\\
3.0653	0.88825\\
3.0625	0.88185\\
3.0597	0.87542\\
3.0569	0.86895\\
3.054	0.86244\\
3.0512	0.8559\\
3.0484	0.84933\\
3.0456	0.84274\\
3.0428	0.83611\\
3.0399	0.82946\\
3.0371	0.82278\\
3.0343	0.81608\\
3.0315	0.80936\\
3.0287	0.80262\\
3.0258	0.79586\\
3.023	0.78908\\
3.0202	0.78228\\
3.0174	0.77547\\
3.0146	0.76864\\
3.0117	0.76181\\
3.0089	0.75496\\
3.0061	0.7481\\
3.0033	0.74123\\
3.0005	0.73436\\
2.9977	0.72748\\
2.9948	0.7206\\
2.992	0.71371\\
2.9892	0.70683\\
2.9864	0.69994\\
2.9836	0.69305\\
2.9807	0.68617\\
2.9779	0.67928\\
2.9751	0.67241\\
2.9723	0.66554\\
2.9695	0.65867\\
2.9666	0.65182\\
2.9638	0.64497\\
2.961	0.63814\\
2.9582	0.63132\\
2.9554	0.62451\\
2.9526	0.61771\\
2.9497	0.61093\\
2.9469	0.60417\\
2.9441	0.59743\\
2.9413	0.5907\\
2.9385	0.58399\\
2.9356	0.57731\\
2.9328	0.57064\\
2.93	0.564\\
2.9272	0.55738\\
2.9244	0.55079\\
2.9215	0.54423\\
2.9187	0.53769\\
2.9159	0.53117\\
2.9131	0.52469\\
2.9103	0.51824\\
2.9075	0.51181\\
2.9046	0.50542\\
2.9018	0.49906\\
2.899	0.49273\\
2.8962	0.48644\\
2.8934	0.48018\\
2.8905	0.47396\\
2.8877	0.46777\\
2.8849	0.46162\\
2.8821	0.4555\\
2.8793	0.44943\\
2.8764	0.44339\\
2.8736	0.43739\\
2.8708	0.43144\\
2.868	0.42552\\
2.8652	0.41964\\
2.8624	0.41381\\
2.8595	0.40801\\
2.8567	0.40227\\
2.8539	0.39656\\
2.8511	0.3909\\
2.8483	0.38528\\
2.8454	0.3797\\
2.8426	0.37418\\
2.8398	0.36869\\
2.837	0.36325\\
2.8342	0.35786\\
2.8313	0.35252\\
2.8285	0.34722\\
2.8257	0.34197\\
2.8229	0.33677\\
2.8201	0.33161\\
2.8173	0.3265\\
2.8144	0.32145\\
2.8116	0.31643\\
2.8088	0.31147\\
2.806	0.30656\\
2.8032	0.3017\\
2.8003	0.29688\\
2.7975	0.29211\\
2.7947	0.2874\\
2.7919	0.28273\\
2.7891	0.27811\\
2.7862	0.27355\\
2.7834	0.26903\\
2.7806	0.26456\\
2.7778	0.26014\\
2.775	0.25577\\
2.7722	0.25145\\
2.7693	0.24718\\
2.7665	0.24297\\
2.7637	0.2388\\
2.7609	0.23468\\
2.7581	0.2306\\
2.7552	0.22658\\
2.7524	0.22261\\
2.7496	0.21869\\
2.7468	0.21481\\
2.744	0.21099\\
2.7411	0.20721\\
2.7383	0.20349\\
2.7355	0.19981\\
2.7327	0.19618\\
2.7299	0.1926\\
2.7271	0.18906\\
2.7242	0.18557\\
2.7214	0.18213\\
2.7186	0.17874\\
2.7158	0.1754\\
2.713	0.1721\\
2.7101	0.16885\\
2.7073	0.16564\\
2.7045	0.16248\\
2.7017	0.15937\\
2.6989	0.1563\\
2.696	0.15327\\
2.6932	0.15029\\
2.6904	0.14736\\
2.6876	0.14447\\
2.6848	0.14162\\
2.682	0.13882\\
2.6791	0.13606\\
2.6763	0.13334\\
2.6735	0.13066\\
2.6707	0.12803\\
2.6679	0.12544\\
2.665	0.12289\\
2.6622	0.12038\\
2.6594	0.11791\\
2.6566	0.11548\\
2.6538	0.11309\\
2.6509	0.11074\\
2.6481	0.10843\\
2.6453	0.10616\\
2.6425	0.10393\\
2.6397	0.10173\\
2.6368	0.099572\\
2.634	0.097451\\
2.6312	0.095366\\
2.6284	0.093318\\
2.6256	0.091306\\
2.6228	0.089329\\
2.6199	0.087387\\
2.6171	0.08548\\
2.6143	0.083607\\
2.6115	0.081768\\
2.6087	0.079962\\
2.6058	0.078189\\
2.603	0.076449\\
2.6002	0.074741\\
2.5974	0.073065\\
2.5946	0.07142\\
2.5917	0.069806\\
2.5889	0.068223\\
2.5861	0.066669\\
2.5833	0.065146\\
2.5805	0.063652\\
2.5777	0.062186\\
2.5748	0.060749\\
2.572	0.05934\\
2.5692	0.057959\\
2.5664	0.056606\\
2.5636	0.055279\\
2.5607	0.053978\\
2.5579	0.052704\\
2.5551	0.051455\\
2.5523	0.050232\\
2.5495	0.049034\\
2.5466	0.04786\\
2.5438	0.04671\\
2.541	0.045585\\
2.5382	0.044482\\
2.5354	0.043403\\
2.5326	0.042346\\
2.5297	0.041312\\
2.5269	0.040299\\
2.5241	0.039308\\
2.5213	0.038338\\
2.5185	0.03739\\
2.5156	0.036461\\
2.5128	0.035553\\
2.51	0.034665\\
2.5072	0.033796\\
2.5044	0.032946\\
2.5015	0.032115\\
2.4987	0.031302\\
2.4959	0.030508\\
2.4931	0.029731\\
2.4903	0.028972\\
2.4875	0.02823\\
2.4846	0.027504\\
2.4818	0.026796\\
2.479	0.026103\\
2.4762	0.025426\\
2.4734	0.024765\\
2.4705	0.02412\\
2.4677	0.023489\\
2.4649	0.022873\\
2.4621	0.022272\\
2.4593	0.021684\\
2.4564	0.021111\\
2.4536	0.020551\\
2.4508	0.020005\\
2.448	0.019471\\
2.4452	0.018951\\
2.4424	0.018442\\
2.4395	0.017947\\
2.4367	0.017463\\
2.4339	0.016991\\
2.4311	0.016531\\
2.4283	0.016081\\
2.4254	0.015643\\
2.4226	0.015216\\
2.4198	0.0148\\
2.417	0.014393\\
2.4142	0.013997\\
2.4113	0.013611\\
2.4085	0.013235\\
2.4057	0.012868\\
2.4029	0.01251\\
2.4001	0.012162\\
2.3973	0.011822\\
2.3944	0.011492\\
2.3916	0.011169\\
2.3888	0.010855\\
}--cycle;
\addplot[ybar interval, fill=mycolor2, fill opacity=0.35, area legend, draw=none] table[row sep=crcr, x=Lower, y=Count] {%
Lower	Upper	Count\\
2.3888	2.6054	0.042103\\
2.6054	2.822	0.19943\\
2.822	3.0386	0.54352\\
3.0386	3.2553	1.0457\\
3.2553	3.4719	1.1208\\
3.4719	3.6885	0.79059\\
3.6885	3.9051	0.37708\\
3.9051	4.1217	0.23009\\
4.1217	4.3383	0.14662\\
4.3383	4.5549	0.071402\\
4.5549	4.7715	0.027576\\
4.7715	4.9882	0.0093561\\
4.9882	5.2048	0.012311\\
5.2048	5.2048	0.012311\\
};
\addlegendentry{\ApproxHistogramLegend}

\addplot [color=mycolor2, only marks, every error bar/.append style={opacity=0.45}, mark=*, mark size=0pt, draw=none, forget plot]
 plot [error bars/.cd, y dir=both, y explicit, error bar style={line width=1pt, color=mycolor2}, error mark options={mark=none,mark size=0pt}]
 table[row sep=crcr, y error plus index=2, y error minus index=3]{%
2.4971	0.042103	0.0028877	0.0028877\\
2.7137	0.19943	0.0047983	0.0047983\\
2.9303	0.54352	0.0090428	0.0090428\\
3.1469	1.0457	0.010945	0.010945\\
3.3636	1.1208	0.011946	0.011946\\
3.5802	0.79059	0.0083296	0.0083296\\
3.7968	0.37708	0.007163	0.007163\\
4.0134	0.23009	0.0062754	0.0062754\\
4.23	0.14662	0.0037889	0.0037889\\
4.4466	0.071402	0.0023802	0.0023802\\
4.6632	0.027576	0.0020575	0.0020575\\
4.8799	0.0093561	0.0012087	0.0012087\\
5.0965	0.012311	2.0756e-18	1.7347e-18\\
};
\addplot [color=mycolor2]
  table[row sep=crcr]{%
2.3888	0.010959\\
2.417	0.014497\\
2.4424	0.018547\\
2.4649	0.022979\\
2.4875	0.028338\\
2.5072	0.033909\\
2.5269	0.04042\\
2.5466	0.047993\\
2.5636	0.055426\\
2.5805	0.063819\\
2.5974	0.073258\\
2.6143	0.083833\\
2.6312	0.095633\\
2.6481	0.10875\\
2.6622	0.12075\\
2.6763	0.13377\\
2.6904	0.14786\\
2.7045	0.16306\\
2.7186	0.17941\\
2.7327	0.19695\\
2.7496	0.2196\\
2.7665	0.24404\\
2.7834	0.27028\\
2.8003	0.29833\\
2.8173	0.32819\\
2.8342	0.3598\\
2.8539	0.39882\\
2.8736	0.44001\\
2.8934	0.4832\\
2.9159	0.53469\\
2.9413	0.59484\\
2.9779	0.68441\\
3.0399	0.83634\\
3.0653	0.89578\\
3.085	0.93989\\
3.102	0.97579\\
3.1189	1.0096\\
3.133	1.036\\
3.1471	1.0606\\
3.1611	1.0832\\
3.1724	1.0997\\
3.1837	1.1149\\
3.195	1.1286\\
3.2062	1.1407\\
3.2175	1.1513\\
3.226	1.1582\\
3.2344	1.1642\\
3.2429	1.1692\\
3.2513	1.1733\\
3.2598	1.1765\\
3.2683	1.1787\\
3.2767	1.1799\\
3.2852	1.1803\\
3.2936	1.1796\\
3.3021	1.1781\\
3.3105	1.1756\\
3.319	1.1722\\
3.3275	1.1679\\
3.3359	1.1627\\
3.3444	1.1567\\
3.3556	1.1473\\
3.3669	1.1365\\
3.3782	1.1243\\
3.3895	1.1107\\
3.4007	1.096\\
3.4148	1.0759\\
3.4289	1.0541\\
3.4458	1.026\\
3.4628	0.99601\\
3.4825	0.95904\\
3.505	0.91478\\
3.536	0.85164\\
3.615	0.68973\\
3.6403	0.64016\\
3.6629	0.59794\\
3.6854	0.55776\\
3.7052	0.52445\\
3.7249	0.49297\\
3.7446	0.46338\\
3.7616	0.43953\\
3.7785	0.41709\\
3.7954	0.39602\\
3.8123	0.3763\\
3.8292	0.35787\\
3.8461	0.34067\\
3.863	0.32466\\
3.8799	0.30975\\
3.8997	0.29366\\
3.9194	0.27885\\
3.9391	0.2652\\
3.9589	0.25261\\
3.9814	0.23934\\
4.004	0.22712\\
4.0293	0.21445\\
4.0575	0.20148\\
4.0885	0.1883\\
4.1224	0.17494\\
4.1618	0.16038\\
4.2069	0.14477\\
4.2577	0.12824\\
4.3084	0.11266\\
4.3591	0.098018\\
4.4071	0.085109\\
4.455	0.073163\\
4.5001	0.062857\\
4.5452	0.053498\\
4.5903	0.045102\\
4.6354	0.037661\\
4.6805	0.031146\\
4.7284	0.025186\\
4.7763	0.020146\\
4.8271	0.015716\\
4.8834	0.011756\\
4.9426	0.0085274\\
5.0103	0.0057889\\
5.0864	0.0036481\\
5.1794	0.0019989\\
5.2048	0.0016844\\
};
\addlegendentry{\ApproxBLNFitLegend}


\addplot[area legend, draw=none, fill=mycolor2, fill opacity=0.15, forget plot]
table[row sep=crcr] {%
x	y\\
2.3888	0.011232\\
2.3916	0.011552\\
2.3944	0.01188\\
2.3973	0.012217\\
2.4001	0.012562\\
2.4029	0.012917\\
2.4057	0.013281\\
2.4085	0.013654\\
2.4113	0.014037\\
2.4142	0.01443\\
2.417	0.014833\\
2.4198	0.015246\\
2.4226	0.01567\\
2.4254	0.016104\\
2.4283	0.016549\\
2.4311	0.017006\\
2.4339	0.017474\\
2.4367	0.017953\\
2.4395	0.018445\\
2.4424	0.018949\\
2.4452	0.019465\\
2.448	0.019994\\
2.4508	0.020536\\
2.4536	0.021091\\
2.4564	0.021659\\
2.4593	0.022241\\
2.4621	0.022838\\
2.4649	0.023448\\
2.4677	0.024073\\
2.4705	0.024713\\
2.4734	0.025369\\
2.4762	0.026039\\
2.479	0.026726\\
2.4818	0.027428\\
2.4846	0.028147\\
2.4875	0.028883\\
2.4903	0.029635\\
2.4931	0.030405\\
2.4959	0.031193\\
2.4987	0.031998\\
2.5015	0.032822\\
2.5044	0.033664\\
2.5072	0.034526\\
2.51	0.035406\\
2.5128	0.036307\\
2.5156	0.037227\\
2.5185	0.038167\\
2.5213	0.039129\\
2.5241	0.040111\\
2.5269	0.041115\\
2.5297	0.04214\\
2.5326	0.043188\\
2.5354	0.044258\\
2.5382	0.045351\\
2.541	0.046467\\
2.5438	0.047607\\
2.5466	0.048771\\
2.5495	0.049959\\
2.5523	0.051172\\
2.5551	0.05241\\
2.5579	0.053673\\
2.5607	0.054963\\
2.5636	0.056278\\
2.5664	0.057621\\
2.5692	0.05899\\
2.572	0.060388\\
2.5748	0.061813\\
2.5777	0.063266\\
2.5805	0.064748\\
2.5833	0.066259\\
2.5861	0.0678\\
2.5889	0.06937\\
2.5917	0.070971\\
2.5946	0.072603\\
2.5974	0.074266\\
2.6002	0.07596\\
2.603	0.077687\\
2.6058	0.079446\\
2.6087	0.081237\\
2.6115	0.083062\\
2.6143	0.084921\\
2.6171	0.086813\\
2.6199	0.08874\\
2.6228	0.090702\\
2.6256	0.092699\\
2.6284	0.094732\\
2.6312	0.096801\\
2.634	0.098907\\
2.6368	0.10105\\
2.6397	0.10323\\
2.6425	0.10545\\
2.6453	0.1077\\
2.6481	0.11\\
2.6509	0.11233\\
2.6538	0.1147\\
2.6566	0.11711\\
2.6594	0.11956\\
2.6622	0.12206\\
2.665	0.12459\\
2.6679	0.12716\\
2.6707	0.12978\\
2.6735	0.13244\\
2.6763	0.13514\\
2.6791	0.13788\\
2.682	0.14067\\
2.6848	0.1435\\
2.6876	0.14637\\
2.6904	0.14929\\
2.6932	0.15225\\
2.696	0.15525\\
2.6989	0.1583\\
2.7017	0.1614\\
2.7045	0.16454\\
2.7073	0.16773\\
2.7101	0.17096\\
2.713	0.17424\\
2.7158	0.17757\\
2.7186	0.18095\\
2.7214	0.18437\\
2.7242	0.18784\\
2.7271	0.19135\\
2.7299	0.19492\\
2.7327	0.19853\\
2.7355	0.20219\\
2.7383	0.2059\\
2.7411	0.20966\\
2.744	0.21346\\
2.7468	0.21732\\
2.7496	0.22123\\
2.7524	0.22518\\
2.7552	0.22919\\
2.7581	0.23324\\
2.7609	0.23734\\
2.7637	0.2415\\
2.7665	0.2457\\
2.7693	0.24995\\
2.7722	0.25426\\
2.775	0.25861\\
2.7778	0.26301\\
2.7806	0.26747\\
2.7834	0.27197\\
2.7862	0.27653\\
2.7891	0.28113\\
2.7919	0.28578\\
2.7947	0.29049\\
2.7975	0.29524\\
2.8003	0.30005\\
2.8032	0.3049\\
2.806	0.3098\\
2.8088	0.31475\\
2.8116	0.31976\\
2.8144	0.32481\\
2.8173	0.32991\\
2.8201	0.33506\\
2.8229	0.34025\\
2.8257	0.3455\\
2.8285	0.35079\\
2.8313	0.35613\\
2.8342	0.36152\\
2.837	0.36695\\
2.8398	0.37244\\
2.8426	0.37796\\
2.8454	0.38354\\
2.8483	0.38916\\
2.8511	0.39482\\
2.8539	0.40053\\
2.8567	0.40629\\
2.8595	0.41209\\
2.8624	0.41793\\
2.8652	0.42381\\
2.868	0.42974\\
2.8708	0.43571\\
2.8736	0.44172\\
2.8764	0.44777\\
2.8793	0.45386\\
2.8821	0.45999\\
2.8849	0.46616\\
2.8877	0.47237\\
2.8905	0.47862\\
2.8934	0.4849\\
2.8962	0.49122\\
2.899	0.49757\\
2.9018	0.50396\\
2.9046	0.51039\\
2.9075	0.51684\\
2.9103	0.52333\\
2.9131	0.52985\\
2.9159	0.5364\\
2.9187	0.54298\\
2.9215	0.54959\\
2.9244	0.55623\\
2.9272	0.56289\\
2.93	0.56958\\
2.9328	0.57629\\
2.9356	0.58303\\
2.9385	0.5898\\
2.9413	0.59658\\
2.9441	0.60338\\
2.9469	0.61021\\
2.9497	0.61705\\
2.9526	0.62392\\
2.9554	0.63079\\
2.9582	0.63769\\
2.961	0.6446\\
2.9638	0.65152\\
2.9666	0.65845\\
2.9695	0.66539\\
2.9723	0.67235\\
2.9751	0.67931\\
2.9779	0.68628\\
2.9807	0.69325\\
2.9836	0.70023\\
2.9864	0.70721\\
2.9892	0.71419\\
2.992	0.72118\\
2.9948	0.72816\\
2.9977	0.73514\\
3.0005	0.74212\\
3.0033	0.74909\\
3.0061	0.75606\\
3.0089	0.76301\\
3.0117	0.76996\\
3.0146	0.7769\\
3.0174	0.78383\\
3.0202	0.79074\\
3.023	0.79763\\
3.0258	0.80451\\
3.0287	0.81138\\
3.0315	0.81822\\
3.0343	0.82504\\
3.0371	0.83184\\
3.0399	0.83861\\
3.0428	0.84536\\
3.0456	0.85209\\
3.0484	0.85878\\
3.0512	0.86545\\
3.054	0.87208\\
3.0569	0.87868\\
3.0597	0.88524\\
3.0625	0.89177\\
3.0653	0.89826\\
3.0681	0.90472\\
3.0709	0.91113\\
3.0738	0.9175\\
3.0766	0.92382\\
3.0794	0.93011\\
3.0822	0.93634\\
3.085	0.94253\\
3.0879	0.94866\\
3.0907	0.95475\\
3.0935	0.96078\\
3.0963	0.96676\\
3.0991	0.97269\\
3.102	0.97856\\
3.1048	0.98436\\
3.1076	0.99011\\
3.1104	0.9958\\
3.1132	1.0014\\
3.116	1.007\\
3.1189	1.0125\\
3.1217	1.0179\\
3.1245	1.0233\\
3.1273	1.0286\\
3.1301	1.0338\\
3.133	1.0389\\
3.1358	1.044\\
3.1386	1.049\\
3.1414	1.054\\
3.1442	1.0588\\
3.1471	1.0636\\
3.1499	1.0683\\
3.1527	1.0729\\
3.1555	1.0774\\
3.1583	1.0819\\
3.1611	1.0862\\
3.164	1.0905\\
3.1668	1.0947\\
3.1696	1.0988\\
3.1724	1.1028\\
3.1752	1.1067\\
3.1781	1.1106\\
3.1809	1.1143\\
3.1837	1.118\\
3.1865	1.1215\\
3.1893	1.125\\
3.1922	1.1284\\
3.195	1.1317\\
3.1978	1.1349\\
3.2006	1.1379\\
3.2034	1.1409\\
3.2062	1.1438\\
3.2091	1.1466\\
3.2119	1.1493\\
3.2147	1.1519\\
3.2175	1.1544\\
3.2203	1.1568\\
3.2232	1.1591\\
3.226	1.1613\\
3.2288	1.1633\\
3.2316	1.1653\\
3.2344	1.1672\\
3.2373	1.169\\
3.2401	1.1706\\
3.2429	1.1722\\
3.2457	1.1737\\
3.2485	1.175\\
3.2513	1.1763\\
3.2542	1.1774\\
3.257	1.1784\\
3.2598	1.1794\\
3.2626	1.1802\\
3.2654	1.1809\\
3.2683	1.1815\\
3.2711	1.182\\
3.2739	1.1824\\
3.2767	1.1827\\
3.2795	1.1829\\
3.2824	1.183\\
3.2852	1.183\\
3.288	1.1829\\
3.2908	1.1826\\
3.2936	1.1823\\
3.2964	1.1819\\
3.2993	1.1813\\
3.3021	1.1807\\
3.3049	1.1799\\
3.3077	1.1791\\
3.3105	1.1781\\
3.3134	1.1771\\
3.3162	1.1759\\
3.319	1.1747\\
3.3218	1.1733\\
3.3246	1.1719\\
3.3275	1.1703\\
3.3303	1.1686\\
3.3331	1.1669\\
3.3359	1.1651\\
3.3387	1.1631\\
3.3415	1.1611\\
3.3444	1.159\\
3.3472	1.1567\\
3.35	1.1544\\
3.3528	1.152\\
3.3556	1.1495\\
3.3585	1.1469\\
3.3613	1.1443\\
3.3641	1.1415\\
3.3669	1.1386\\
3.3697	1.1357\\
3.3726	1.1327\\
3.3754	1.1296\\
3.3782	1.1264\\
3.381	1.1231\\
3.3838	1.1198\\
3.3866	1.1164\\
3.3895	1.1129\\
3.3923	1.1093\\
3.3951	1.1056\\
3.3979	1.1019\\
3.4007	1.0981\\
3.4036	1.0942\\
3.4064	1.0903\\
3.4092	1.0863\\
3.412	1.0822\\
3.4148	1.078\\
3.4177	1.0738\\
3.4205	1.0695\\
3.4233	1.0652\\
3.4261	1.0608\\
3.4289	1.0563\\
3.4318	1.0518\\
3.4346	1.0472\\
3.4374	1.0425\\
3.4402	1.0378\\
3.443	1.0331\\
3.4458	1.0283\\
3.4487	1.0234\\
3.4515	1.0185\\
3.4543	1.0135\\
3.4571	1.0085\\
3.4599	1.0035\\
3.4628	0.99839\\
3.4656	0.99325\\
3.4684	0.98806\\
3.4712	0.98284\\
3.474	0.97757\\
3.4769	0.97227\\
3.4797	0.96693\\
3.4825	0.96156\\
3.4853	0.95615\\
3.4881	0.95071\\
3.4909	0.94523\\
3.4938	0.93973\\
3.4966	0.9342\\
3.4994	0.92864\\
3.5022	0.92305\\
3.505	0.91744\\
3.5079	0.91181\\
3.5107	0.90615\\
3.5135	0.90047\\
3.5163	0.89478\\
3.5191	0.88906\\
3.522	0.88333\\
3.5248	0.87759\\
3.5276	0.87183\\
3.5304	0.86605\\
3.5332	0.86027\\
3.536	0.85447\\
3.5389	0.84867\\
3.5417	0.84285\\
3.5445	0.83703\\
3.5473	0.83121\\
3.5501	0.82538\\
3.553	0.81954\\
3.5558	0.81371\\
3.5586	0.80787\\
3.5614	0.80204\\
3.5642	0.79621\\
3.5671	0.79037\\
3.5699	0.78455\\
3.5727	0.77872\\
3.5755	0.77291\\
3.5783	0.7671\\
3.5811	0.7613\\
3.584	0.7555\\
3.5868	0.74972\\
3.5896	0.74395\\
3.5924	0.73819\\
3.5952	0.73245\\
3.5981	0.72671\\
3.6009	0.721\\
3.6037	0.71529\\
3.6065	0.70961\\
3.6093	0.70394\\
3.6122	0.69829\\
3.615	0.69266\\
3.6178	0.68706\\
3.6206	0.68147\\
3.6234	0.6759\\
3.6262	0.67036\\
3.6291	0.66484\\
3.6319	0.65934\\
3.6347	0.65387\\
3.6375	0.64842\\
3.6403	0.643\\
3.6432	0.63761\\
3.646	0.63224\\
3.6488	0.6269\\
3.6516	0.6216\\
3.6544	0.61632\\
3.6573	0.61107\\
3.6601	0.60585\\
3.6629	0.60066\\
3.6657	0.5955\\
3.6685	0.59037\\
3.6713	0.58528\\
3.6742	0.58022\\
3.677	0.5752\\
3.6798	0.5702\\
3.6826	0.56524\\
3.6854	0.56032\\
3.6883	0.55543\\
3.6911	0.55058\\
3.6939	0.54576\\
3.6967	0.54098\\
3.6995	0.53623\\
3.7024	0.53152\\
3.7052	0.52685\\
3.708	0.52221\\
3.7108	0.51761\\
3.7136	0.51305\\
3.7164	0.50853\\
3.7193	0.50405\\
3.7221	0.4996\\
3.7249	0.49519\\
3.7277	0.49082\\
3.7305	0.48649\\
3.7334	0.48219\\
3.7362	0.47794\\
3.739	0.47372\\
3.7418	0.46955\\
3.7446	0.46541\\
3.7475	0.46131\\
3.7503	0.45725\\
3.7531	0.45323\\
3.7559	0.44925\\
3.7587	0.4453\\
3.7616	0.4414\\
3.7644	0.43753\\
3.7672	0.43371\\
3.77	0.42992\\
3.7728	0.42617\\
3.7756	0.42246\\
3.7785	0.41879\\
3.7813	0.41515\\
3.7841	0.41156\\
3.7869	0.408\\
3.7897	0.40448\\
3.7926	0.401\\
3.7954	0.39756\\
3.7982	0.39415\\
3.801	0.39078\\
3.8038	0.38745\\
3.8067	0.38416\\
3.8095	0.3809\\
3.8123	0.37768\\
3.8151	0.3745\\
3.8179	0.37135\\
3.8207	0.36824\\
3.8236	0.36516\\
3.8264	0.36212\\
3.8292	0.35912\\
3.832	0.35615\\
3.8348	0.35321\\
3.8377	0.35031\\
3.8405	0.34744\\
3.8433	0.34461\\
3.8461	0.34181\\
3.8489	0.33904\\
3.8518	0.33631\\
3.8546	0.33361\\
3.8574	0.33094\\
3.8602	0.32831\\
3.863	0.3257\\
3.8658	0.32313\\
3.8687	0.32059\\
3.8715	0.31808\\
3.8743	0.3156\\
3.8771	0.31315\\
3.8799	0.31073\\
3.8828	0.30834\\
3.8856	0.30598\\
3.8884	0.30364\\
3.8912	0.30134\\
3.894	0.29906\\
3.8969	0.29681\\
3.8997	0.29459\\
3.9025	0.2924\\
3.9053	0.29023\\
3.9081	0.28809\\
3.9109	0.28597\\
3.9138	0.28388\\
3.9166	0.28181\\
3.9194	0.27977\\
3.9222	0.27776\\
3.925	0.27576\\
3.9279	0.27379\\
3.9307	0.27185\\
3.9335	0.26992\\
3.9363	0.26802\\
3.9391	0.26614\\
3.942	0.26428\\
3.9448	0.26245\\
3.9476	0.26063\\
3.9504	0.25884\\
3.9532	0.25706\\
3.956	0.25531\\
3.9589	0.25357\\
3.9617	0.25185\\
3.9645	0.25016\\
3.9673	0.24848\\
3.9701	0.24681\\
3.973	0.24517\\
3.9758	0.24354\\
3.9786	0.24193\\
3.9814	0.24034\\
3.9842	0.23876\\
3.9871	0.2372\\
3.9899	0.23566\\
3.9927	0.23413\\
3.9955	0.23261\\
3.9983	0.23112\\
4.0011	0.22963\\
4.004	0.22816\\
4.0068	0.2267\\
4.0096	0.22526\\
4.0124	0.22383\\
4.0152	0.22241\\
4.0181	0.22101\\
4.0209	0.21962\\
4.0237	0.21824\\
4.0265	0.21687\\
4.0293	0.21552\\
4.0322	0.21417\\
4.035	0.21284\\
4.0378	0.21152\\
4.0406	0.21021\\
4.0434	0.20891\\
4.0462	0.20762\\
4.0491	0.20634\\
4.0519	0.20507\\
4.0547	0.20381\\
4.0575	0.20256\\
4.0603	0.20132\\
4.0632	0.20008\\
4.066	0.19886\\
4.0688	0.19765\\
4.0716	0.19644\\
4.0744	0.19524\\
4.0773	0.19405\\
4.0801	0.19287\\
4.0829	0.19169\\
4.0857	0.19053\\
4.0885	0.18937\\
4.0913	0.18821\\
4.0942	0.18707\\
4.097	0.18593\\
4.0998	0.1848\\
4.1026	0.18367\\
4.1054	0.18255\\
4.1083	0.18144\\
4.1111	0.18033\\
4.1139	0.17923\\
4.1167	0.17814\\
4.1195	0.17705\\
4.1224	0.17597\\
4.1252	0.17489\\
4.128	0.17382\\
4.1308	0.17275\\
4.1336	0.17169\\
4.1365	0.17063\\
4.1393	0.16958\\
4.1421	0.16853\\
4.1449	0.16749\\
4.1477	0.16646\\
4.1505	0.16542\\
4.1534	0.16439\\
4.1562	0.16337\\
4.159	0.16235\\
4.1618	0.16134\\
4.1646	0.16032\\
4.1675	0.15932\\
4.1703	0.15831\\
4.1731	0.15731\\
4.1759	0.15632\\
4.1787	0.15533\\
4.1816	0.15434\\
4.1844	0.15336\\
4.1872	0.15238\\
4.19	0.1514\\
4.1928	0.15043\\
4.1956	0.14946\\
4.1985	0.14849\\
4.2013	0.14753\\
4.2041	0.14657\\
4.2069	0.14561\\
4.2097	0.14466\\
4.2126	0.14371\\
4.2154	0.14276\\
4.2182	0.14182\\
4.221	0.14088\\
4.2238	0.13994\\
4.2267	0.13901\\
4.2295	0.13808\\
4.2323	0.13715\\
4.2351	0.13623\\
4.2379	0.1353\\
4.2407	0.13438\\
4.2436	0.13347\\
4.2464	0.13256\\
4.2492	0.13165\\
4.252	0.13074\\
4.2548	0.12984\\
4.2577	0.12893\\
4.2605	0.12804\\
4.2633	0.12714\\
4.2661	0.12625\\
4.2689	0.12536\\
4.2718	0.12447\\
4.2746	0.12359\\
4.2774	0.12271\\
4.2802	0.12183\\
4.283	0.12095\\
4.2858	0.12008\\
4.2887	0.11921\\
4.2915	0.11835\\
4.2943	0.11748\\
4.2971	0.11662\\
4.2999	0.11576\\
4.3028	0.11491\\
4.3056	0.11406\\
4.3084	0.11321\\
4.3112	0.11236\\
4.314	0.11152\\
4.3169	0.11068\\
4.3197	0.10984\\
4.3225	0.109\\
4.3253	0.10817\\
4.3281	0.10734\\
4.3309	0.10652\\
4.3338	0.1057\\
4.3366	0.10488\\
4.3394	0.10406\\
4.3422	0.10325\\
4.345	0.10244\\
4.3479	0.10163\\
4.3507	0.10082\\
4.3535	0.10002\\
4.3563	0.099224\\
4.3591	0.098429\\
4.362	0.097638\\
4.3648	0.096849\\
4.3676	0.096064\\
4.3704	0.095282\\
4.3732	0.094503\\
4.376	0.093727\\
4.3789	0.092954\\
4.3817	0.092185\\
4.3845	0.091419\\
4.3873	0.090656\\
4.3901	0.089897\\
4.393	0.089141\\
4.3958	0.088388\\
4.3986	0.087639\\
4.4014	0.086893\\
4.4042	0.086151\\
4.4071	0.085412\\
4.4099	0.084676\\
4.4127	0.083944\\
4.4155	0.083215\\
4.4183	0.08249\\
4.4211	0.081768\\
4.424	0.08105\\
4.4268	0.080335\\
4.4296	0.079624\\
4.4324	0.078916\\
4.4352	0.078212\\
4.4381	0.077511\\
4.4409	0.076814\\
4.4437	0.076121\\
4.4465	0.075431\\
4.4493	0.074745\\
4.4522	0.074063\\
4.455	0.073384\\
4.4578	0.072709\\
4.4606	0.072038\\
4.4634	0.07137\\
4.4663	0.070706\\
4.4691	0.070046\\
4.4719	0.069389\\
4.4747	0.068737\\
4.4775	0.068088\\
4.4803	0.067442\\
4.4832	0.066801\\
4.486	0.066163\\
4.4888	0.065529\\
4.4916	0.064899\\
4.4944	0.064273\\
4.4973	0.06365\\
4.5001	0.063032\\
4.5029	0.062417\\
4.5057	0.061806\\
4.5085	0.061199\\
4.5114	0.060595\\
4.5142	0.059996\\
4.517	0.0594\\
4.5198	0.058808\\
4.5226	0.05822\\
4.5254	0.057636\\
4.5283	0.057056\\
4.5311	0.056479\\
4.5339	0.055907\\
4.5367	0.055338\\
4.5395	0.054773\\
4.5424	0.054212\\
4.5452	0.053655\\
4.548	0.053101\\
4.5508	0.052552\\
4.5536	0.052006\\
4.5565	0.051464\\
4.5593	0.050926\\
4.5621	0.050392\\
4.5649	0.049861\\
4.5677	0.049335\\
4.5705	0.048812\\
4.5734	0.048293\\
4.5762	0.047778\\
4.579	0.047266\\
4.5818	0.046759\\
4.5846	0.046255\\
4.5875	0.045755\\
4.5903	0.045258\\
4.5931	0.044766\\
4.5959	0.044277\\
4.5987	0.043792\\
4.6016	0.04331\\
4.6044	0.042833\\
4.6072	0.042359\\
4.61	0.041889\\
4.6128	0.041422\\
4.6156	0.040959\\
4.6185	0.0405\\
4.6213	0.040045\\
4.6241	0.039593\\
4.6269	0.039145\\
4.6297	0.0387\\
4.6326	0.03826\\
4.6354	0.037822\\
4.6382	0.037389\\
4.641	0.036959\\
4.6438	0.036532\\
4.6467	0.03611\\
4.6495	0.03569\\
4.6523	0.035275\\
4.6551	0.034863\\
4.6579	0.034454\\
4.6607	0.034049\\
4.6636	0.033647\\
4.6664	0.033249\\
4.6692	0.032855\\
4.672	0.032464\\
4.6748	0.032076\\
4.6777	0.031692\\
4.6805	0.031311\\
4.6833	0.030934\\
4.6861	0.03056\\
4.6889	0.030189\\
4.6918	0.029822\\
4.6946	0.029458\\
4.6974	0.029098\\
4.7002	0.028741\\
4.703	0.028387\\
4.7058	0.028037\\
4.7087	0.02769\\
4.7115	0.027346\\
4.7143	0.027005\\
4.7171	0.026668\\
4.7199	0.026334\\
4.7228	0.026003\\
4.7256	0.025675\\
4.7284	0.02535\\
4.7312	0.025029\\
4.734	0.024711\\
4.7369	0.024396\\
4.7397	0.024084\\
4.7425	0.023775\\
4.7453	0.023469\\
4.7481	0.023166\\
4.7509	0.022867\\
4.7538	0.02257\\
4.7566	0.022276\\
4.7594	0.021986\\
4.7622	0.021698\\
4.765	0.021413\\
4.7679	0.021131\\
4.7707	0.020853\\
4.7735	0.020577\\
4.7763	0.020304\\
4.7791	0.020033\\
4.782	0.019766\\
4.7848	0.019502\\
4.7876	0.01924\\
4.7904	0.018981\\
4.7932	0.018725\\
4.796	0.018472\\
4.7989	0.018221\\
4.8017	0.017973\\
4.8045	0.017728\\
4.8073	0.017486\\
4.8101	0.017246\\
4.813	0.017009\\
4.8158	0.016774\\
4.8186	0.016542\\
4.8214	0.016313\\
4.8242	0.016086\\
4.8271	0.015862\\
4.8299	0.01564\\
4.8327	0.015421\\
4.8355	0.015205\\
4.8383	0.014991\\
4.8412	0.014779\\
4.844	0.01457\\
4.8468	0.014363\\
4.8496	0.014158\\
4.8524	0.013956\\
4.8552	0.013757\\
4.8581	0.013559\\
4.8609	0.013364\\
4.8637	0.013172\\
4.8665	0.012981\\
4.8693	0.012793\\
4.8722	0.012607\\
4.875	0.012424\\
4.8778	0.012242\\
4.8806	0.012063\\
4.8834	0.011886\\
4.8863	0.011711\\
4.8891	0.011538\\
4.8919	0.011368\\
4.8947	0.011199\\
4.8975	0.011033\\
4.9003	0.010869\\
4.9032	0.010706\\
4.906	0.010546\\
4.9088	0.010388\\
4.9116	0.010231\\
4.9144	0.010077\\
4.9173	0.0099247\\
4.9201	0.0097743\\
4.9229	0.0096258\\
4.9257	0.0094792\\
4.9285	0.0093344\\
4.9314	0.0091916\\
4.9342	0.0090506\\
4.937	0.0089114\\
4.9398	0.008774\\
4.9426	0.0086384\\
4.9454	0.0085046\\
4.9483	0.0083726\\
4.9511	0.0082423\\
4.9539	0.0081137\\
4.9567	0.0079868\\
4.9595	0.0078616\\
4.9624	0.0077381\\
4.9652	0.0076162\\
4.968	0.007496\\
4.9708	0.0073774\\
4.9736	0.0072604\\
4.9765	0.007145\\
4.9793	0.0070312\\
4.9821	0.0069189\\
4.9849	0.0068082\\
4.9877	0.0066989\\
4.9905	0.0065912\\
4.9934	0.006485\\
4.9962	0.0063802\\
4.999	0.0062769\\
5.0018	0.0061751\\
5.0046	0.0060746\\
5.0075	0.0059756\\
5.0103	0.005878\\
5.0131	0.0057817\\
5.0159	0.0056868\\
5.0187	0.0055933\\
5.0216	0.0055011\\
5.0244	0.0054102\\
5.0272	0.0053206\\
5.03	0.0052323\\
5.0328	0.0051453\\
5.0356	0.0050595\\
5.0385	0.004975\\
5.0413	0.0048917\\
5.0441	0.0048096\\
5.0469	0.0047287\\
5.0497	0.004649\\
5.0526	0.0045705\\
5.0554	0.0044931\\
5.0582	0.0044169\\
5.061	0.0043418\\
5.0638	0.0042678\\
5.0667	0.0041949\\
5.0695	0.0041231\\
5.0723	0.0040524\\
5.0751	0.0039828\\
5.0779	0.0039142\\
5.0807	0.0038466\\
5.0836	0.0037801\\
5.0864	0.0037146\\
5.0892	0.0036501\\
5.092	0.0035865\\
5.0948	0.003524\\
5.0977	0.0034624\\
5.1005	0.0034017\\
5.1033	0.003342\\
5.1061	0.0032832\\
5.1089	0.0032253\\
5.1118	0.0031684\\
5.1146	0.0031123\\
5.1174	0.0030571\\
5.1202	0.0030027\\
5.123	0.0029493\\
5.1258	0.0028966\\
5.1287	0.0028448\\
5.1315	0.0027938\\
5.1343	0.0027437\\
5.1371	0.0026943\\
5.1399	0.0026457\\
5.1428	0.0025979\\
5.1456	0.0025509\\
5.1484	0.0025046\\
5.1512	0.0024591\\
5.154	0.0024143\\
5.1569	0.0023702\\
5.1597	0.0023269\\
5.1625	0.0022842\\
5.1653	0.0022423\\
5.1681	0.002201\\
5.171	0.0021605\\
5.1738	0.0021205\\
5.1766	0.0020813\\
5.1794	0.0020427\\
5.1822	0.0020047\\
5.185	0.0019674\\
5.1879	0.0019307\\
5.1907	0.0018946\\
5.1935	0.0018591\\
5.1963	0.0018243\\
5.1991	0.00179\\
5.202	0.0017562\\
5.2048	0.0017231\\
5.2048	0.0016457\\
5.202	0.0016778\\
5.1991	0.0017104\\
5.1963	0.0017436\\
5.1935	0.0017774\\
5.1907	0.0018117\\
5.1879	0.0018466\\
5.185	0.0018822\\
5.1822	0.0019183\\
5.1794	0.0019551\\
5.1766	0.0019925\\
5.1738	0.0020306\\
5.171	0.0020693\\
5.1681	0.0021086\\
5.1653	0.0021486\\
5.1625	0.0021893\\
5.1597	0.0022307\\
5.1569	0.0022728\\
5.154	0.0023156\\
5.1512	0.0023591\\
5.1484	0.0024033\\
5.1456	0.0024482\\
5.1428	0.0024939\\
5.1399	0.0025404\\
5.1371	0.0025876\\
5.1343	0.0026356\\
5.1315	0.0026844\\
5.1287	0.002734\\
5.1258	0.0027844\\
5.123	0.0028357\\
5.1202	0.0028877\\
5.1174	0.0029406\\
5.1146	0.0029944\\
5.1118	0.003049\\
5.1089	0.0031046\\
5.1061	0.003161\\
5.1033	0.0032183\\
5.1005	0.0032765\\
5.0977	0.0033356\\
5.0948	0.0033957\\
5.092	0.0034567\\
5.0892	0.0035187\\
5.0864	0.0035817\\
5.0836	0.0036457\\
5.0807	0.0037106\\
5.0779	0.0037766\\
5.0751	0.0038436\\
5.0723	0.0039117\\
5.0695	0.0039808\\
5.0667	0.0040509\\
5.0638	0.0041222\\
5.061	0.0041945\\
5.0582	0.004268\\
5.0554	0.0043425\\
5.0526	0.0044182\\
5.0497	0.0044951\\
5.0469	0.0045731\\
5.0441	0.0046523\\
5.0413	0.0047327\\
5.0385	0.0048143\\
5.0356	0.0048971\\
5.0328	0.0049812\\
5.03	0.0050665\\
5.0272	0.005153\\
5.0244	0.0052409\\
5.0216	0.00533\\
5.0187	0.0054205\\
5.0159	0.0055122\\
5.0131	0.0056053\\
5.0103	0.0056998\\
5.0075	0.0057957\\
5.0046	0.0058929\\
5.0018	0.0059915\\
4.999	0.0060916\\
4.9962	0.0061931\\
4.9934	0.006296\\
4.9905	0.0064004\\
4.9877	0.0065063\\
4.9849	0.0066137\\
4.9821	0.0067227\\
4.9793	0.0068331\\
4.9765	0.0069451\\
4.9736	0.0070587\\
4.9708	0.0071738\\
4.968	0.0072906\\
4.9652	0.0074089\\
4.9624	0.007529\\
4.9595	0.0076506\\
4.9567	0.007774\\
4.9539	0.007899\\
4.9511	0.0080257\\
4.9483	0.0081541\\
4.9454	0.0082843\\
4.9426	0.0084163\\
4.9398	0.00855\\
4.937	0.0086855\\
4.9342	0.0088229\\
4.9314	0.008962\\
4.9285	0.009103\\
4.9257	0.0092459\\
4.9229	0.0093907\\
4.9201	0.0095373\\
4.9173	0.0096859\\
4.9144	0.0098365\\
4.9116	0.009989\\
4.9088	0.010143\\
4.906	0.0103\\
4.9032	0.010458\\
4.9003	0.010619\\
4.8975	0.010781\\
4.8947	0.010946\\
4.8919	0.011113\\
4.8891	0.011282\\
4.8863	0.011453\\
4.8834	0.011626\\
4.8806	0.011801\\
4.8778	0.011979\\
4.875	0.012158\\
4.8722	0.01234\\
4.8693	0.012524\\
4.8665	0.012711\\
4.8637	0.012899\\
4.8609	0.01309\\
4.8581	0.013284\\
4.8552	0.013479\\
4.8524	0.013677\\
4.8496	0.013878\\
4.8468	0.014081\\
4.844	0.014286\\
4.8412	0.014494\\
4.8383	0.014704\\
4.8355	0.014916\\
4.8327	0.015132\\
4.8299	0.015349\\
4.8271	0.015569\\
4.8242	0.015792\\
4.8214	0.016017\\
4.8186	0.016245\\
4.8158	0.016476\\
4.813	0.016709\\
4.8101	0.016945\\
4.8073	0.017183\\
4.8045	0.017424\\
4.8017	0.017668\\
4.7989	0.017915\\
4.796	0.018164\\
4.7932	0.018416\\
4.7904	0.018671\\
4.7876	0.018929\\
4.7848	0.01919\\
4.782	0.019453\\
4.7791	0.019719\\
4.7763	0.019988\\
4.7735	0.02026\\
4.7707	0.020535\\
4.7679	0.020813\\
4.765	0.021094\\
4.7622	0.021378\\
4.7594	0.021665\\
4.7566	0.021955\\
4.7538	0.022248\\
4.7509	0.022543\\
4.7481	0.022842\\
4.7453	0.023144\\
4.7425	0.02345\\
4.7397	0.023758\\
4.7369	0.024069\\
4.734	0.024384\\
4.7312	0.024702\\
4.7284	0.025023\\
4.7256	0.025347\\
4.7228	0.025674\\
4.7199	0.026005\\
4.7171	0.026339\\
4.7143	0.026676\\
4.7115	0.027016\\
4.7087	0.02736\\
4.7058	0.027707\\
4.703	0.028057\\
4.7002	0.028411\\
4.6974	0.028768\\
4.6946	0.029128\\
4.6918	0.029492\\
4.6889	0.029859\\
4.6861	0.03023\\
4.6833	0.030604\\
4.6805	0.030981\\
4.6777	0.031362\\
4.6748	0.031747\\
4.672	0.032135\\
4.6692	0.032526\\
4.6664	0.032921\\
4.6636	0.033319\\
4.6607	0.033721\\
4.6579	0.034127\\
4.6551	0.034536\\
4.6523	0.034948\\
4.6495	0.035365\\
4.6467	0.035784\\
4.6438	0.036208\\
4.641	0.036635\\
4.6382	0.037065\\
4.6354	0.0375\\
4.6326	0.037938\\
4.6297	0.038379\\
4.6269	0.038824\\
4.6241	0.039273\\
4.6213	0.039725\\
4.6185	0.040182\\
4.6156	0.040641\\
4.6128	0.041105\\
4.61	0.041572\\
4.6072	0.042043\\
4.6044	0.042518\\
4.6016	0.042996\\
4.5987	0.043478\\
4.5959	0.043964\\
4.5931	0.044453\\
4.5903	0.044946\\
4.5875	0.045443\\
4.5846	0.045943\\
4.5818	0.046448\\
4.579	0.046956\\
4.5762	0.047467\\
4.5734	0.047983\\
4.5705	0.048502\\
4.5677	0.049025\\
4.5649	0.049552\\
4.5621	0.050082\\
4.5593	0.050616\\
4.5565	0.051154\\
4.5536	0.051695\\
4.5508	0.05224\\
4.548	0.052789\\
4.5452	0.053342\\
4.5424	0.053898\\
4.5395	0.054458\\
4.5367	0.055022\\
4.5339	0.055589\\
4.5311	0.05616\\
4.5283	0.056735\\
4.5254	0.057313\\
4.5226	0.057895\\
4.5198	0.058481\\
4.517	0.05907\\
4.5142	0.059663\\
4.5114	0.060259\\
4.5085	0.06086\\
4.5057	0.061463\\
4.5029	0.062071\\
4.5001	0.062682\\
4.4973	0.063297\\
4.4944	0.063915\\
4.4916	0.064536\\
4.4888	0.065162\\
4.486	0.065791\\
4.4832	0.066423\\
4.4803	0.067059\\
4.4775	0.067699\\
4.4747	0.068342\\
4.4719	0.068988\\
4.4691	0.069638\\
4.4663	0.070292\\
4.4634	0.070949\\
4.4606	0.07161\\
4.4578	0.072274\\
4.455	0.072941\\
4.4522	0.073612\\
4.4493	0.074286\\
4.4465	0.074964\\
4.4437	0.075645\\
4.4409	0.07633\\
4.4381	0.077018\\
4.4352	0.077709\\
4.4324	0.078404\\
4.4296	0.079102\\
4.4268	0.079803\\
4.424	0.080508\\
4.4211	0.081216\\
4.4183	0.081928\\
4.4155	0.082642\\
4.4127	0.08336\\
4.4099	0.084081\\
4.4071	0.084806\\
4.4042	0.085534\\
4.4014	0.086265\\
4.3986	0.086999\\
4.3958	0.087736\\
4.393	0.088477\\
4.3901	0.08922\\
4.3873	0.089967\\
4.3845	0.090717\\
4.3817	0.09147\\
4.3789	0.092227\\
4.376	0.092986\\
4.3732	0.093748\\
4.3704	0.094514\\
4.3676	0.095283\\
4.3648	0.096054\\
4.362	0.096829\\
4.3591	0.097607\\
4.3563	0.098387\\
4.3535	0.099171\\
4.3507	0.099958\\
4.3479	0.10075\\
4.345	0.10154\\
4.3422	0.10234\\
4.3394	0.10313\\
4.3366	0.10394\\
4.3338	0.10474\\
4.3309	0.10555\\
4.3281	0.10636\\
4.3253	0.10717\\
4.3225	0.10799\\
4.3197	0.10881\\
4.3169	0.10963\\
4.314	0.11045\\
4.3112	0.11128\\
4.3084	0.11211\\
4.3056	0.11294\\
4.3028	0.11378\\
4.2999	0.11462\\
4.2971	0.11546\\
4.2943	0.11631\\
4.2915	0.11715\\
4.2887	0.118\\
4.2858	0.11886\\
4.283	0.11971\\
4.2802	0.12057\\
4.2774	0.12143\\
4.2746	0.1223\\
4.2718	0.12316\\
4.2689	0.12403\\
4.2661	0.12491\\
4.2633	0.12578\\
4.2605	0.12666\\
4.2577	0.12754\\
4.2548	0.12843\\
4.252	0.12932\\
4.2492	0.13021\\
4.2464	0.1311\\
4.2436	0.132\\
4.2407	0.1329\\
4.2379	0.1338\\
4.2351	0.1347\\
4.2323	0.13561\\
4.2295	0.13652\\
4.2267	0.13744\\
4.2238	0.13836\\
4.221	0.13928\\
4.2182	0.1402\\
4.2154	0.14113\\
4.2126	0.14206\\
4.2097	0.14299\\
4.2069	0.14393\\
4.2041	0.14487\\
4.2013	0.14582\\
4.1985	0.14676\\
4.1956	0.14772\\
4.1928	0.14867\\
4.19	0.14963\\
4.1872	0.15059\\
4.1844	0.15156\\
4.1816	0.15253\\
4.1787	0.1535\\
4.1759	0.15448\\
4.1731	0.15546\\
4.1703	0.15645\\
4.1675	0.15744\\
4.1646	0.15843\\
4.1618	0.15943\\
4.159	0.16043\\
4.1562	0.16144\\
4.1534	0.16245\\
4.1505	0.16347\\
4.1477	0.16449\\
4.1449	0.16552\\
4.1421	0.16655\\
4.1393	0.16758\\
4.1365	0.16863\\
4.1336	0.16967\\
4.1308	0.17072\\
4.128	0.17178\\
4.1252	0.17285\\
4.1224	0.17391\\
4.1195	0.17499\\
4.1167	0.17607\\
4.1139	0.17716\\
4.1111	0.17825\\
4.1083	0.17935\\
4.1054	0.18046\\
4.1026	0.18157\\
4.0998	0.18269\\
4.097	0.18381\\
4.0942	0.18495\\
4.0913	0.18609\\
4.0885	0.18724\\
4.0857	0.18839\\
4.0829	0.18956\\
4.0801	0.19073\\
4.0773	0.19191\\
4.0744	0.1931\\
4.0716	0.19429\\
4.0688	0.1955\\
4.066	0.19671\\
4.0632	0.19793\\
4.0603	0.19917\\
4.0575	0.20041\\
4.0547	0.20166\\
4.0519	0.20292\\
4.0491	0.20419\\
4.0462	0.20547\\
4.0434	0.20677\\
4.0406	0.20807\\
4.0378	0.20938\\
4.035	0.21071\\
4.0322	0.21204\\
4.0293	0.21339\\
4.0265	0.21475\\
4.0237	0.21612\\
4.0209	0.21751\\
4.0181	0.2189\\
4.0152	0.22031\\
4.0124	0.22174\\
4.0096	0.22317\\
4.0068	0.22462\\
4.004	0.22609\\
4.0011	0.22757\\
3.9983	0.22906\\
3.9955	0.23057\\
3.9927	0.23209\\
3.9899	0.23363\\
3.9871	0.23518\\
3.9842	0.23675\\
3.9814	0.23834\\
3.9786	0.23994\\
3.9758	0.24156\\
3.973	0.24319\\
3.9701	0.24485\\
3.9673	0.24652\\
3.9645	0.24821\\
3.9617	0.24991\\
3.9589	0.25164\\
3.956	0.25338\\
3.9532	0.25515\\
3.9504	0.25693\\
3.9476	0.25873\\
3.9448	0.26056\\
3.942	0.2624\\
3.9391	0.26427\\
3.9363	0.26615\\
3.9335	0.26806\\
3.9307	0.26999\\
3.9279	0.27194\\
3.925	0.27391\\
3.9222	0.2759\\
3.9194	0.27792\\
3.9166	0.27996\\
3.9138	0.28203\\
3.9109	0.28412\\
3.9081	0.28623\\
3.9053	0.28837\\
3.9025	0.29053\\
3.8997	0.29272\\
3.8969	0.29493\\
3.894	0.29717\\
3.8912	0.29944\\
3.8884	0.30173\\
3.8856	0.30405\\
3.8828	0.3064\\
3.8799	0.30877\\
3.8771	0.31117\\
3.8743	0.3136\\
3.8715	0.31606\\
3.8687	0.31855\\
3.8658	0.32107\\
3.863	0.32362\\
3.8602	0.32619\\
3.8574	0.3288\\
3.8546	0.33144\\
3.8518	0.33411\\
3.8489	0.33681\\
3.8461	0.33954\\
3.8433	0.3423\\
3.8405	0.3451\\
3.8377	0.34793\\
3.8348	0.35079\\
3.832	0.35368\\
3.8292	0.35661\\
3.8264	0.35958\\
3.8236	0.36257\\
3.8207	0.3656\\
3.8179	0.36867\\
3.8151	0.37177\\
3.8123	0.37491\\
3.8095	0.37808\\
3.8067	0.38129\\
3.8038	0.38453\\
3.801	0.38781\\
3.7982	0.39113\\
3.7954	0.39448\\
3.7926	0.39787\\
3.7897	0.4013\\
3.7869	0.40477\\
3.7841	0.40827\\
3.7813	0.41181\\
3.7785	0.41539\\
3.7756	0.41901\\
3.7728	0.42266\\
3.77	0.42636\\
3.7672	0.43009\\
3.7644	0.43386\\
3.7616	0.43767\\
3.7587	0.44152\\
3.7559	0.44541\\
3.7531	0.44933\\
3.7503	0.4533\\
3.7475	0.4573\\
3.7446	0.46135\\
3.7418	0.46543\\
3.739	0.46955\\
3.7362	0.47371\\
3.7334	0.47791\\
3.7305	0.48215\\
3.7277	0.48643\\
3.7249	0.49075\\
3.7221	0.4951\\
3.7193	0.4995\\
3.7164	0.50393\\
3.7136	0.5084\\
3.7108	0.51291\\
3.708	0.51746\\
3.7052	0.52205\\
3.7024	0.52667\\
3.6995	0.53133\\
3.6967	0.53603\\
3.6939	0.54077\\
3.6911	0.54554\\
3.6883	0.55035\\
3.6854	0.55519\\
3.6826	0.56007\\
3.6798	0.56499\\
3.677	0.56994\\
3.6742	0.57493\\
3.6713	0.57995\\
3.6685	0.585\\
3.6657	0.59009\\
3.6629	0.59521\\
3.6601	0.60037\\
3.6573	0.60555\\
3.6544	0.61077\\
3.6516	0.61602\\
3.6488	0.6213\\
3.646	0.62661\\
3.6432	0.63195\\
3.6403	0.63731\\
3.6375	0.64271\\
3.6347	0.64813\\
3.6319	0.65358\\
3.6291	0.65906\\
3.6262	0.66456\\
3.6234	0.67008\\
3.6206	0.67563\\
3.6178	0.68121\\
3.615	0.6868\\
3.6122	0.69242\\
3.6093	0.69806\\
3.6065	0.70371\\
3.6037	0.70939\\
3.6009	0.71509\\
3.5981	0.7208\\
3.5952	0.72653\\
3.5924	0.73227\\
3.5896	0.73803\\
3.5868	0.7438\\
3.584	0.74959\\
3.5811	0.75538\\
3.5783	0.76119\\
3.5755	0.767\\
3.5727	0.77283\\
3.5699	0.77866\\
3.5671	0.7845\\
3.5642	0.79034\\
3.5614	0.79619\\
3.5586	0.80204\\
3.5558	0.80789\\
3.553	0.81374\\
3.5501	0.8196\\
3.5473	0.82545\\
3.5445	0.83129\\
3.5417	0.83713\\
3.5389	0.84297\\
3.536	0.8488\\
3.5332	0.85462\\
3.5304	0.86044\\
3.5276	0.86624\\
3.5248	0.87203\\
3.522	0.87781\\
3.5191	0.88357\\
3.5163	0.88931\\
3.5135	0.89504\\
3.5107	0.90075\\
3.5079	0.90644\\
3.505	0.91211\\
3.5022	0.91776\\
3.4994	0.92338\\
3.4966	0.92898\\
3.4938	0.93455\\
3.4909	0.94009\\
3.4881	0.9456\\
3.4853	0.95108\\
3.4825	0.95653\\
3.4797	0.96194\\
3.4769	0.96732\\
3.474	0.97266\\
3.4712	0.97796\\
3.4684	0.98322\\
3.4656	0.98844\\
3.4628	0.99362\\
3.4599	0.99876\\
3.4571	1.0038\\
3.4543	1.0089\\
3.4515	1.0139\\
3.4487	1.0188\\
3.4458	1.0237\\
3.443	1.0286\\
3.4402	1.0333\\
3.4374	1.0381\\
3.4346	1.0427\\
3.4318	1.0473\\
3.4289	1.0519\\
3.4261	1.0564\\
3.4233	1.0608\\
3.4205	1.0652\\
3.4177	1.0695\\
3.4148	1.0737\\
3.412	1.0779\\
3.4092	1.082\\
3.4064	1.086\\
3.4036	1.09\\
3.4007	1.0938\\
3.3979	1.0977\\
3.3951	1.1014\\
3.3923	1.105\\
3.3895	1.1086\\
3.3866	1.1121\\
3.3838	1.1155\\
3.381	1.1189\\
3.3782	1.1221\\
3.3754	1.1253\\
3.3726	1.1284\\
3.3697	1.1314\\
3.3669	1.1343\\
3.3641	1.1371\\
3.3613	1.1398\\
3.3585	1.1425\\
3.3556	1.1451\\
3.3528	1.1475\\
3.35	1.1499\\
3.3472	1.1522\\
3.3444	1.1544\\
3.3415	1.1565\\
3.3387	1.1585\\
3.3359	1.1604\\
3.3331	1.1622\\
3.3303	1.1639\\
3.3275	1.1655\\
3.3246	1.167\\
3.3218	1.1684\\
3.319	1.1697\\
3.3162	1.1709\\
3.3134	1.172\\
3.3105	1.1731\\
3.3077	1.174\\
3.3049	1.1748\\
3.3021	1.1755\\
3.2993	1.1761\\
3.2964	1.1766\\
3.2936	1.177\\
3.2908	1.1773\\
3.288	1.1775\\
3.2852	1.1775\\
3.2824	1.1775\\
3.2795	1.1774\\
3.2767	1.1772\\
3.2739	1.1768\\
3.2711	1.1764\\
3.2683	1.1758\\
3.2654	1.1752\\
3.2626	1.1744\\
3.2598	1.1736\\
3.257	1.1726\\
3.2542	1.1715\\
3.2513	1.1704\\
3.2485	1.1691\\
3.2457	1.1677\\
3.2429	1.1662\\
3.2401	1.1646\\
3.2373	1.1629\\
3.2344	1.1612\\
3.2316	1.1593\\
3.2288	1.1573\\
3.226	1.1552\\
3.2232	1.1529\\
3.2203	1.1506\\
3.2175	1.1482\\
3.2147	1.1457\\
3.2119	1.1431\\
3.2091	1.1404\\
3.2062	1.1376\\
3.2034	1.1347\\
3.2006	1.1317\\
3.1978	1.1286\\
3.195	1.1255\\
3.1922	1.1222\\
3.1893	1.1188\\
3.1865	1.1153\\
3.1837	1.1118\\
3.1809	1.1081\\
3.1781	1.1044\\
3.1752	1.1006\\
3.1724	1.0966\\
3.1696	1.0926\\
3.1668	1.0885\\
3.164	1.0844\\
3.1611	1.0801\\
3.1583	1.0757\\
3.1555	1.0713\\
3.1527	1.0668\\
3.1499	1.0622\\
3.1471	1.0575\\
3.1442	1.0528\\
3.1414	1.048\\
3.1386	1.0431\\
3.1358	1.0381\\
3.133	1.033\\
3.1301	1.0279\\
3.1273	1.0227\\
3.1245	1.0175\\
3.1217	1.0121\\
3.1189	1.0067\\
3.116	1.0013\\
3.1132	0.99575\\
3.1104	0.99016\\
3.1076	0.98451\\
3.1048	0.9788\\
3.102	0.97303\\
3.0991	0.9672\\
3.0963	0.96132\\
3.0935	0.95538\\
3.0907	0.94939\\
3.0879	0.94335\\
3.085	0.93725\\
3.0822	0.93111\\
3.0794	0.92492\\
3.0766	0.91868\\
3.0738	0.9124\\
3.0709	0.90608\\
3.0681	0.89971\\
3.0653	0.8933\\
3.0625	0.88686\\
3.0597	0.88037\\
3.0569	0.87385\\
3.054	0.8673\\
3.0512	0.86071\\
3.0484	0.85409\\
3.0456	0.84744\\
3.0428	0.84077\\
3.0399	0.83406\\
3.0371	0.82733\\
3.0343	0.82057\\
3.0315	0.8138\\
3.0287	0.807\\
3.0258	0.80018\\
3.023	0.79334\\
3.0202	0.78648\\
3.0174	0.77961\\
3.0146	0.77273\\
3.0117	0.76583\\
3.0089	0.75892\\
3.0061	0.752\\
3.0033	0.74507\\
3.0005	0.73813\\
2.9977	0.73119\\
2.9948	0.72425\\
2.992	0.7173\\
2.9892	0.71034\\
2.9864	0.70339\\
2.9836	0.69644\\
2.9807	0.68949\\
2.9779	0.68254\\
2.9751	0.6756\\
2.9723	0.66866\\
2.9695	0.66173\\
2.9666	0.65481\\
2.9638	0.6479\\
2.961	0.641\\
2.9582	0.63411\\
2.9554	0.62724\\
2.9526	0.62037\\
2.9497	0.61353\\
2.9469	0.6067\\
2.9441	0.59989\\
2.9413	0.59309\\
2.9385	0.58632\\
2.9356	0.57957\\
2.9328	0.57284\\
2.93	0.56613\\
2.9272	0.55945\\
2.9244	0.55279\\
2.9215	0.54616\\
2.9187	0.53956\\
2.9159	0.53298\\
2.9131	0.52644\\
2.9103	0.51992\\
2.9075	0.51343\\
2.9046	0.50698\\
2.9018	0.50056\\
2.899	0.49417\\
2.8962	0.48781\\
2.8934	0.48149\\
2.8905	0.47521\\
2.8877	0.46896\\
2.8849	0.46275\\
2.8821	0.45658\\
2.8793	0.45045\\
2.8764	0.44435\\
2.8736	0.4383\\
2.8708	0.43229\\
2.868	0.42632\\
2.8652	0.42039\\
2.8624	0.4145\\
2.8595	0.40865\\
2.8567	0.40285\\
2.8539	0.3971\\
2.8511	0.39138\\
2.8483	0.38572\\
2.8454	0.3801\\
2.8426	0.37452\\
2.8398	0.36899\\
2.837	0.36351\\
2.8342	0.35807\\
2.8313	0.35268\\
2.8285	0.34734\\
2.8257	0.34205\\
2.8229	0.33681\\
2.8201	0.33161\\
2.8173	0.32646\\
2.8144	0.32137\\
2.8116	0.31632\\
2.8088	0.31132\\
2.806	0.30637\\
2.8032	0.30147\\
2.8003	0.29662\\
2.7975	0.29182\\
2.7947	0.28707\\
2.7919	0.28238\\
2.7891	0.27773\\
2.7862	0.27313\\
2.7834	0.26859\\
2.7806	0.26409\\
2.7778	0.25964\\
2.775	0.25525\\
2.7722	0.25091\\
2.7693	0.24661\\
2.7665	0.24237\\
2.7637	0.23818\\
2.7609	0.23404\\
2.7581	0.22994\\
2.7552	0.2259\\
2.7524	0.22191\\
2.7496	0.21797\\
2.7468	0.21408\\
2.744	0.21024\\
2.7411	0.20645\\
2.7383	0.2027\\
2.7355	0.19901\\
2.7327	0.19537\\
2.7299	0.19177\\
2.7271	0.18822\\
2.7242	0.18473\\
2.7214	0.18128\\
2.7186	0.17787\\
2.7158	0.17452\\
2.713	0.17121\\
2.7101	0.16795\\
2.7073	0.16474\\
2.7045	0.16157\\
2.7017	0.15845\\
2.6989	0.15538\\
2.696	0.15235\\
2.6932	0.14936\\
2.6904	0.14643\\
2.6876	0.14353\\
2.6848	0.14068\\
2.682	0.13788\\
2.6791	0.13511\\
2.6763	0.13239\\
2.6735	0.12972\\
2.6707	0.12709\\
2.6679	0.12449\\
2.665	0.12194\\
2.6622	0.11944\\
2.6594	0.11697\\
2.6566	0.11454\\
2.6538	0.11216\\
2.6509	0.10981\\
2.6481	0.1075\\
2.6453	0.10523\\
2.6425	0.103\\
2.6397	0.10081\\
2.6368	0.09866\\
2.634	0.096544\\
2.6312	0.094465\\
2.6284	0.092423\\
2.6256	0.090417\\
2.6228	0.088446\\
2.6199	0.086511\\
2.6171	0.08461\\
2.6143	0.082744\\
2.6115	0.080913\\
2.6087	0.079114\\
2.6058	0.077349\\
2.603	0.075617\\
2.6002	0.073917\\
2.5974	0.072249\\
2.5946	0.070613\\
2.5917	0.069008\\
2.5889	0.067433\\
2.5861	0.065889\\
2.5833	0.064374\\
2.5805	0.062889\\
2.5777	0.061433\\
2.5748	0.060006\\
2.572	0.058606\\
2.5692	0.057235\\
2.5664	0.055891\\
2.5636	0.054574\\
2.5607	0.053283\\
2.5579	0.052019\\
2.5551	0.05078\\
2.5523	0.049567\\
2.5495	0.048379\\
2.5466	0.047215\\
2.5438	0.046075\\
2.541	0.04496\\
2.5382	0.043867\\
2.5354	0.042798\\
2.5326	0.041752\\
2.5297	0.040727\\
2.5269	0.039725\\
2.5241	0.038744\\
2.5213	0.037784\\
2.5185	0.036846\\
2.5156	0.035927\\
2.5128	0.035029\\
2.51	0.034151\\
2.5072	0.033291\\
2.5044	0.032451\\
2.5015	0.03163\\
2.4987	0.030827\\
2.4959	0.030042\\
2.4931	0.029275\\
2.4903	0.028526\\
2.4875	0.027793\\
2.4846	0.027077\\
2.4818	0.026378\\
2.479	0.025694\\
2.4762	0.025027\\
2.4734	0.024375\\
2.4705	0.023738\\
2.4677	0.023116\\
2.4649	0.022509\\
2.4621	0.021916\\
2.4593	0.021337\\
2.4564	0.020772\\
2.4536	0.020221\\
2.4508	0.019683\\
2.448	0.019157\\
2.4452	0.018645\\
2.4424	0.018145\\
2.4395	0.017657\\
2.4367	0.01718\\
2.4339	0.016716\\
2.4311	0.016263\\
2.4283	0.015821\\
2.4254	0.015391\\
2.4226	0.014971\\
2.4198	0.014561\\
2.417	0.014162\\
2.4142	0.013772\\
2.4113	0.013393\\
2.4085	0.013023\\
2.4057	0.012663\\
2.4029	0.012311\\
2.4001	0.011969\\
2.3973	0.011636\\
2.3944	0.011311\\
2.3916	0.010994\\
2.3888	0.010686\\
}--cycle;
      % \input{../../../../.tex/20/KS/SizeDistributionFittingsfig1.tex}
  
      \nextgroupplot[%
        xmin=0.57363,xmax=5.3544,
        % xmin=1.9138,xmax=5.4719,
        % ymin=0,ymax=1.2,
      ] % This file was created by matlab2tikz.
%
\definecolor{mycolor1}{rgb}{0.00000,0.44706,0.69804}%
%
\addplot[ybar interval, fill=mycolor1, fill opacity=0.35, area legend, draw=none] table[row sep=crcr, x=Lower, y=Count] {%
Lower	Upper	Count\\
2.0755	2.3244	0\\
2.3244	2.5732	0.023471\\
2.5732	2.822	0.1882\\
2.822	3.0708	0.59696\\
3.0708	3.3196	1.0779\\
3.3196	3.5684	1.06\\
3.5684	3.8172	0.49686\\
3.8172	4.0661	0.30234\\
4.0661	4.3149	0.15915\\
4.3149	4.5637	0.072665\\
4.5637	4.8125	0.028616\\
4.8125	5.0613	0.0022507\\
5.0613	5.3101	0.010717\\
5.3101	5.3101	0.010717\\
};
\addlegendentry{\ExactHistogramLegend}

\addplot [color=mycolor1, only marks, every error bar/.append style={opacity=0.45}, mark=*, mark size=0pt, draw=none, forget plot]
 plot [error bars/.cd, y dir=both, y explicit, error bar style={line width=1pt, color=mycolor1}, error mark options={mark=none,mark size=0pt}]
 table[row sep=crcr, y error plus index=2, y error minus index=3]{%
2.2	0	0	0\\
2.4488	0.023471	0.0020432	0.0020432\\
2.6976	0.1882	0.0044353	0.0044353\\
2.9464	0.59696	0.0076556	0.0076556\\
3.1952	1.0779	0.010184	0.010184\\
3.444	1.06	0.0076436	0.0076436\\
3.6928	0.49686	0.0075777	0.0075777\\
3.9417	0.30234	0.0048672	0.0048672\\
4.1905	0.15915	0.0027928	0.0027928\\
4.4393	0.072665	0.0020629	0.0020629\\
4.6881	0.028616	0.0015429	0.0015429\\
4.9369	0.0022507	0.00087054	0.00087054\\
5.1857	0.010717	2.4216e-18	1.7347e-18\\
};
\addplot [color=mycolor1]
  table[row sep=crcr]{%
2.0755	0.00037006\\
2.1856	0.0012593\\
2.2536	0.002669\\
2.3022	0.0045137\\
2.341	0.0068051\\
2.3767	0.009823\\
2.4058	0.013166\\
2.4317	0.016977\\
2.4576	0.02176\\
2.4803	0.026898\\
2.4997	0.032129\\
2.5191	0.038231\\
2.5386	0.045317\\
2.558	0.053505\\
2.5742	0.061259\\
2.5904	0.069938\\
2.6066	0.07962\\
2.6227	0.090381\\
2.6389	0.1023\\
2.6551	0.11544\\
2.6713	0.1299\\
2.6875	0.14572\\
2.7037	0.16299\\
2.7199	0.18175\\
2.7361	0.20206\\
2.7523	0.22396\\
2.7684	0.24747\\
2.7846	0.27261\\
2.8008	0.29938\\
2.817	0.32777\\
2.8364	0.36392\\
2.8559	0.40226\\
2.8753	0.44266\\
2.898	0.49214\\
2.9206	0.54382\\
2.9498	0.61279\\
2.9919	0.71525\\
3.0372	0.82514\\
3.0631	0.8856\\
3.0825	0.92897\\
3.1019	0.97013\\
3.1181	1.0024\\
3.1343	1.0326\\
3.1473	1.055\\
3.1602	1.0759\\
3.1732	1.095\\
3.1861	1.1123\\
3.1958	1.124\\
3.2055	1.1346\\
3.2153	1.144\\
3.225	1.1523\\
3.2347	1.1595\\
3.2444	1.1654\\
3.2541	1.1701\\
3.2638	1.1735\\
3.2735	1.1758\\
3.2833	1.1768\\
3.293	1.1766\\
3.3027	1.1751\\
3.3124	1.1725\\
3.3221	1.1686\\
3.3318	1.1636\\
3.3415	1.1575\\
3.3513	1.1502\\
3.361	1.1419\\
3.3707	1.1325\\
3.3804	1.1221\\
3.3933	1.1067\\
3.4063	1.0897\\
3.4192	1.0712\\
3.4322	1.0513\\
3.4484	1.0246\\
3.4646	0.99629\\
3.484	0.96035\\
3.5067	0.91635\\
3.5358	0.85757\\
3.6265	0.67302\\
3.6524	0.62319\\
3.675	0.58154\\
3.6945	0.54751\\
3.7139	0.51516\\
3.7333	0.48457\\
3.7527	0.45579\\
3.7722	0.42884\\
3.7916	0.40371\\
3.8078	0.38413\\
3.824	0.36576\\
3.8402	0.34855\\
3.8596	0.32935\\
3.879	0.31166\\
3.8984	0.29537\\
3.9179	0.28036\\
3.9373	0.26653\\
3.9567	0.25375\\
3.9794	0.24004\\
4.0021	0.22745\\
4.0247	0.21584\\
4.0506	0.20357\\
4.0798	0.19083\\
4.1121	0.17774\\
4.1478	0.16433\\
4.1898	0.14952\\
4.2352	0.13452\\
4.287	0.11838\\
4.3388	0.10322\\
4.3874	0.08991\\
4.4359	0.077533\\
4.4813	0.066887\\
4.5266	0.057161\\
4.5719	0.048383\\
4.6172	0.040558\\
4.6626	0.033668\\
4.7079	0.027675\\
4.7565	0.022189\\
4.8083	0.017314\\
4.8633	0.013116\\
4.9216	0.0096208\\
4.9864	0.0066894\\
5.0608	0.0042965\\
5.1482	0.00247\\
5.2551	0.0011949\\
5.3101	0.00080471\\
};
\addlegendentry{\ExactBLNFitLegend}


\addplot[area legend, draw=none, fill=mycolor1, fill opacity=0.15, forget plot]
table[row sep=crcr] {%
x	y\\
2.0755	0.00039646\\
2.0788	0.0004105\\
2.082	0.00042505\\
2.0853	0.00044014\\
2.0885	0.00045578\\
2.0917	0.00047199\\
2.095	0.00048879\\
2.0982	0.00050621\\
2.1014	0.00052427\\
2.1047	0.000543\\
2.1079	0.0005624\\
2.1112	0.00058252\\
2.1144	0.00060338\\
2.1176	0.00062499\\
2.1209	0.0006474\\
2.1241	0.00067063\\
2.1274	0.00069472\\
2.1306	0.00071968\\
2.1338	0.00074555\\
2.1371	0.00077237\\
2.1403	0.00080017\\
2.1435	0.00082898\\
2.1468	0.00085884\\
2.15	0.00088979\\
2.1533	0.00092187\\
2.1565	0.00095512\\
2.1597	0.00098957\\
2.163	0.0010253\\
2.1662	0.0010623\\
2.1694	0.0011006\\
2.1727	0.0011403\\
2.1759	0.0011815\\
2.1792	0.0012242\\
2.1824	0.0012684\\
2.1856	0.0013141\\
2.1889	0.0013616\\
2.1921	0.0014107\\
2.1953	0.0014616\\
2.1986	0.0015143\\
2.2018	0.0015689\\
2.2051	0.0016255\\
2.2083	0.0016841\\
2.2115	0.0017447\\
2.2148	0.0018075\\
2.218	0.0018726\\
2.2212	0.0019399\\
2.2245	0.0020096\\
2.2277	0.0020818\\
2.231	0.0021565\\
2.2342	0.0022339\\
2.2374	0.002314\\
2.2407	0.0023968\\
2.2439	0.0024826\\
2.2472	0.0025713\\
2.2504	0.0026631\\
2.2536	0.0027581\\
2.2569	0.0028564\\
2.2601	0.0029581\\
2.2633	0.0030633\\
2.2666	0.0031721\\
2.2698	0.0032847\\
2.2731	0.003401\\
2.2763	0.0035214\\
2.2795	0.0036458\\
2.2828	0.0037744\\
2.286	0.0039074\\
2.2892	0.0040449\\
2.2925	0.004187\\
2.2957	0.0043339\\
2.299	0.0044857\\
2.3022	0.0046425\\
2.3054	0.0048046\\
2.3087	0.004972\\
2.3119	0.005145\\
2.3151	0.0053236\\
2.3184	0.0055081\\
2.3216	0.0056987\\
2.3249	0.0058954\\
2.3281	0.0060986\\
2.3313	0.0063083\\
2.3346	0.0065249\\
2.3378	0.0067484\\
2.341	0.006979\\
2.3443	0.0072171\\
2.3475	0.0074627\\
2.3508	0.0077162\\
2.354	0.0079776\\
2.3572	0.0082473\\
2.3605	0.0085255\\
2.3637	0.0088124\\
2.367	0.0091082\\
2.3702	0.0094133\\
2.3734	0.0097277\\
2.3767	0.010052\\
2.3799	0.010386\\
2.3831	0.01073\\
2.3864	0.011085\\
2.3896	0.011451\\
2.3929	0.011828\\
2.3961	0.012216\\
2.3993	0.012615\\
2.4026	0.013027\\
2.4058	0.013451\\
2.409	0.013887\\
2.4123	0.014337\\
2.4155	0.014799\\
2.4188	0.015276\\
2.422	0.015766\\
2.4252	0.01627\\
2.4285	0.016789\\
2.4317	0.017322\\
2.4349	0.017871\\
2.4382	0.018436\\
2.4414	0.019017\\
2.4447	0.019614\\
2.4479	0.020228\\
2.4511	0.020859\\
2.4544	0.021508\\
2.4576	0.022174\\
2.4608	0.02286\\
2.4641	0.023564\\
2.4673	0.024287\\
2.4706	0.02503\\
2.4738	0.025793\\
2.477	0.026577\\
2.4803	0.027381\\
2.4835	0.028207\\
2.4868	0.029055\\
2.49	0.029926\\
2.4932	0.030819\\
2.4965	0.031736\\
2.4997	0.032677\\
2.5029	0.033642\\
2.5062	0.034632\\
2.5094	0.035647\\
2.5127	0.036688\\
2.5159	0.037755\\
2.5191	0.03885\\
2.5224	0.039972\\
2.5256	0.041122\\
2.5288	0.0423\\
2.5321	0.043507\\
2.5353	0.044744\\
2.5386	0.046011\\
2.5418	0.047309\\
2.545	0.048638\\
2.5483	0.049999\\
2.5515	0.051392\\
2.5547	0.052819\\
2.558	0.054279\\
2.5612	0.055773\\
2.5645	0.057302\\
2.5677	0.058866\\
2.5709	0.060466\\
2.5742	0.062102\\
2.5774	0.063776\\
2.5806	0.065487\\
2.5839	0.067237\\
2.5871	0.069026\\
2.5904	0.070854\\
2.5936	0.072722\\
2.5968	0.074631\\
2.6001	0.076581\\
2.6033	0.078574\\
2.6066	0.080609\\
2.6098	0.082687\\
2.613	0.084809\\
2.6163	0.086975\\
2.6195	0.089186\\
2.6227	0.091443\\
2.626	0.093747\\
2.6292	0.096097\\
2.6325	0.098494\\
2.6357	0.10094\\
2.6389	0.10343\\
2.6422	0.10598\\
2.6454	0.10857\\
2.6486	0.11121\\
2.6519	0.11391\\
2.6551	0.11666\\
2.6584	0.11945\\
2.6616	0.1223\\
2.6648	0.12521\\
2.6681	0.12817\\
2.6713	0.13118\\
2.6745	0.13425\\
2.6778	0.13737\\
2.681	0.14055\\
2.6843	0.14378\\
2.6875	0.14707\\
2.6907	0.15042\\
2.694	0.15383\\
2.6972	0.1573\\
2.7004	0.16082\\
2.7037	0.16441\\
2.7069	0.16805\\
2.7102	0.17175\\
2.7134	0.17552\\
2.7166	0.17934\\
2.7199	0.18323\\
2.7231	0.18718\\
2.7263	0.19119\\
2.7296	0.19526\\
2.7328	0.19939\\
2.7361	0.20359\\
2.7393	0.20785\\
2.7425	0.21218\\
2.7458	0.21657\\
2.749	0.22102\\
2.7523	0.22554\\
2.7555	0.23012\\
2.7587	0.23476\\
2.762	0.23947\\
2.7652	0.24425\\
2.7684	0.24909\\
2.7717	0.25399\\
2.7749	0.25896\\
2.7782	0.264\\
2.7814	0.2691\\
2.7846	0.27426\\
2.7879	0.27949\\
2.7911	0.28479\\
2.7943	0.29015\\
2.7976	0.29557\\
2.8008	0.30106\\
2.8041	0.30662\\
2.8073	0.31224\\
2.8105	0.31792\\
2.8138	0.32366\\
2.817	0.32947\\
2.8202	0.33534\\
2.8235	0.34128\\
2.8267	0.34728\\
2.83	0.35334\\
2.8332	0.35946\\
2.8364	0.36564\\
2.8397	0.37188\\
2.8429	0.37819\\
2.8461	0.38455\\
2.8494	0.39097\\
2.8526	0.39745\\
2.8559	0.40399\\
2.8591	0.41058\\
2.8623	0.41724\\
2.8656	0.42394\\
2.8688	0.4307\\
2.8721	0.43752\\
2.8753	0.44439\\
2.8785	0.45131\\
2.8818	0.45828\\
2.885	0.4653\\
2.8882	0.47238\\
2.8915	0.4795\\
2.8947	0.48666\\
2.898	0.49388\\
2.9012	0.50114\\
2.9044	0.50844\\
2.9077	0.51579\\
2.9109	0.52318\\
2.9141	0.5306\\
2.9174	0.53807\\
2.9206	0.54558\\
2.9239	0.55312\\
2.9271	0.5607\\
2.9303	0.56831\\
2.9336	0.57595\\
2.9368	0.58362\\
2.94	0.59132\\
2.9433	0.59906\\
2.9465	0.60681\\
2.9498	0.61459\\
2.953	0.6224\\
2.9562	0.63022\\
2.9595	0.63807\\
2.9627	0.64593\\
2.9659	0.65381\\
2.9692	0.6617\\
2.9724	0.66961\\
2.9757	0.67752\\
2.9789	0.68545\\
2.9821	0.69338\\
2.9854	0.70132\\
2.9886	0.70926\\
2.9919	0.7172\\
2.9951	0.72514\\
2.9983	0.73308\\
3.0016	0.74101\\
3.0048	0.74894\\
3.008	0.75685\\
3.0113	0.76476\\
3.0145	0.77265\\
3.0178	0.78053\\
3.021	0.78839\\
3.0242	0.79623\\
3.0275	0.80404\\
3.0307	0.81184\\
3.0339	0.81961\\
3.0372	0.82735\\
3.0404	0.83506\\
3.0437	0.84273\\
3.0469	0.85037\\
3.0501	0.85798\\
3.0534	0.86554\\
3.0566	0.87307\\
3.0598	0.88055\\
3.0631	0.88798\\
3.0663	0.89537\\
3.0696	0.9027\\
3.0728	0.90998\\
3.076	0.91721\\
3.0793	0.92438\\
3.0825	0.93149\\
3.0857	0.93854\\
3.089	0.94552\\
3.0922	0.95244\\
3.0955	0.95929\\
3.0987	0.96607\\
3.1019	0.97278\\
3.1052	0.97942\\
3.1084	0.98597\\
3.1117	0.99245\\
3.1149	0.99885\\
3.1181	1.0052\\
3.1214	1.0114\\
3.1246	1.0175\\
3.1278	1.0236\\
3.1311	1.0296\\
3.1343	1.0354\\
3.1376	1.0412\\
3.1408	1.0469\\
3.144	1.0525\\
3.1473	1.0579\\
3.1505	1.0633\\
3.1537	1.0686\\
3.157	1.0738\\
3.1602	1.0788\\
3.1635	1.0838\\
3.1667	1.0886\\
3.1699	1.0933\\
3.1732	1.098\\
3.1764	1.1025\\
3.1796	1.1069\\
3.1829	1.1111\\
3.1861	1.1153\\
3.1894	1.1193\\
3.1926	1.1232\\
3.1958	1.127\\
3.1991	1.1307\\
3.2023	1.1342\\
3.2055	1.1376\\
3.2088	1.1409\\
3.212	1.144\\
3.2153	1.147\\
3.2185	1.1499\\
3.2217	1.1527\\
3.225	1.1553\\
3.2282	1.1578\\
3.2315	1.1602\\
3.2347	1.1624\\
3.2379	1.1645\\
3.2412	1.1665\\
3.2444	1.1683\\
3.2476	1.17\\
3.2509	1.1715\\
3.2541	1.173\\
3.2574	1.1742\\
3.2606	1.1754\\
3.2638	1.1764\\
3.2671	1.1772\\
3.2703	1.178\\
3.2735	1.1786\\
3.2768	1.179\\
3.28	1.1793\\
3.2833	1.1795\\
3.2865	1.1796\\
3.2897	1.1795\\
3.293	1.1792\\
3.2962	1.1789\\
3.2994	1.1784\\
3.3027	1.1777\\
3.3059	1.177\\
3.3092	1.1761\\
3.3124	1.175\\
3.3156	1.1739\\
3.3189	1.1726\\
3.3221	1.1711\\
3.3253	1.1696\\
3.3286	1.1679\\
3.3318	1.1661\\
3.3351	1.1641\\
3.3383	1.1621\\
3.3415	1.1599\\
3.3448	1.1575\\
3.348	1.1551\\
3.3513	1.1525\\
3.3545	1.1499\\
3.3577	1.1471\\
3.361	1.1441\\
3.3642	1.1411\\
3.3674	1.138\\
3.3707	1.1347\\
3.3739	1.1313\\
3.3772	1.1279\\
3.3804	1.1243\\
3.3836	1.1206\\
3.3869	1.1168\\
3.3901	1.1129\\
3.3933	1.1088\\
3.3966	1.1047\\
3.3998	1.1005\\
3.4031	1.0962\\
3.4063	1.0918\\
3.4095	1.0873\\
3.4128	1.0828\\
3.416	1.0781\\
3.4192	1.0733\\
3.4225	1.0685\\
3.4257	1.0636\\
3.429	1.0585\\
3.4322	1.0535\\
3.4354	1.0483\\
3.4387	1.043\\
3.4419	1.0377\\
3.4451	1.0323\\
3.4484	1.0269\\
3.4516	1.0213\\
3.4549	1.0158\\
3.4581	1.0101\\
3.4613	1.0044\\
3.4646	0.99859\\
3.4678	0.99275\\
3.471	0.98685\\
3.4743	0.9809\\
3.4775	0.9749\\
3.4808	0.96885\\
3.484	0.96275\\
3.4872	0.9566\\
3.4905	0.95041\\
3.4937	0.94418\\
3.497	0.93791\\
3.5002	0.9316\\
3.5034	0.92525\\
3.5067	0.91887\\
3.5099	0.91246\\
3.5131	0.90601\\
3.5164	0.89954\\
3.5196	0.89305\\
3.5229	0.88652\\
3.5261	0.87998\\
3.5293	0.87342\\
3.5326	0.86683\\
3.5358	0.86024\\
3.539	0.85362\\
3.5423	0.847\\
3.5455	0.84036\\
3.5488	0.83371\\
3.552	0.82706\\
3.5552	0.8204\\
3.5585	0.81374\\
3.5617	0.80707\\
3.5649	0.8004\\
3.5682	0.79374\\
3.5714	0.78708\\
3.5747	0.78042\\
3.5779	0.77377\\
3.5811	0.76713\\
3.5844	0.7605\\
3.5876	0.75388\\
3.5908	0.74727\\
3.5941	0.74067\\
3.5973	0.73409\\
3.6006	0.72753\\
3.6038	0.72099\\
3.607	0.71446\\
3.6103	0.70796\\
3.6135	0.70148\\
3.6168	0.69502\\
3.62	0.68859\\
3.6232	0.68218\\
3.6265	0.6758\\
3.6297	0.66945\\
3.6329	0.66313\\
3.6362	0.65684\\
3.6394	0.65059\\
3.6427	0.64436\\
3.6459	0.63817\\
3.6491	0.63202\\
3.6524	0.6259\\
3.6556	0.61982\\
3.6588	0.61377\\
3.6621	0.60777\\
3.6653	0.6018\\
3.6686	0.59587\\
3.6718	0.58999\\
3.675	0.58415\\
3.6783	0.57835\\
3.6815	0.57259\\
3.6847	0.56687\\
3.688	0.5612\\
3.6912	0.55558\\
3.6945	0.55\\
3.6977	0.54447\\
3.7009	0.53898\\
3.7042	0.53354\\
3.7074	0.52815\\
3.7106	0.52281\\
3.7139	0.51751\\
3.7171	0.51227\\
3.7204	0.50707\\
3.7236	0.50192\\
3.7268	0.49682\\
3.7301	0.49177\\
3.7333	0.48677\\
3.7366	0.48183\\
3.7398	0.47693\\
3.743	0.47208\\
3.7463	0.46728\\
3.7495	0.46253\\
3.7527	0.45784\\
3.756	0.45319\\
3.7592	0.4486\\
3.7625	0.44405\\
3.7657	0.43956\\
3.7689	0.43512\\
3.7722	0.43073\\
3.7754	0.42639\\
3.7786	0.4221\\
3.7819	0.41786\\
3.7851	0.41367\\
3.7884	0.40953\\
3.7916	0.40544\\
3.7948	0.4014\\
3.7981	0.39741\\
3.8013	0.39347\\
3.8045	0.38958\\
3.8078	0.38573\\
3.811	0.38194\\
3.8143	0.37819\\
3.8175	0.37449\\
3.8207	0.37084\\
3.824	0.36724\\
3.8272	0.36368\\
3.8304	0.36017\\
3.8337	0.35671\\
3.8369	0.35329\\
3.8402	0.34991\\
3.8434	0.34659\\
3.8466	0.3433\\
3.8499	0.34006\\
3.8531	0.33687\\
3.8564	0.33371\\
3.8596	0.33061\\
3.8628	0.32754\\
3.8661	0.32451\\
3.8693	0.32153\\
3.8725	0.31859\\
3.8758	0.31568\\
3.879	0.31282\\
3.8823	0.31\\
3.8855	0.30721\\
3.8887	0.30447\\
3.892	0.30176\\
3.8952	0.29909\\
3.8984	0.29646\\
3.9017	0.29386\\
3.9049	0.2913\\
3.9082	0.28878\\
3.9114	0.28628\\
3.9146	0.28383\\
3.9179	0.28141\\
3.9211	0.27902\\
3.9243	0.27666\\
3.9276	0.27433\\
3.9308	0.27204\\
3.9341	0.26978\\
3.9373	0.26754\\
3.9405	0.26534\\
3.9438	0.26317\\
3.947	0.26103\\
3.9502	0.25891\\
3.9535	0.25682\\
3.9567	0.25476\\
3.96	0.25273\\
3.9632	0.25072\\
3.9664	0.24873\\
3.9697	0.24678\\
3.9729	0.24484\\
3.9762	0.24293\\
3.9794	0.24105\\
3.9826	0.23919\\
3.9859	0.23735\\
3.9891	0.23553\\
3.9923	0.23373\\
3.9956	0.23196\\
3.9988	0.2302\\
4.0021	0.22847\\
4.0053	0.22676\\
4.0085	0.22506\\
4.0118	0.22339\\
4.015	0.22173\\
4.0182	0.22009\\
4.0215	0.21847\\
4.0247	0.21686\\
4.028	0.21527\\
4.0312	0.2137\\
4.0344	0.21215\\
4.0377	0.21061\\
4.0409	0.20908\\
4.0441	0.20757\\
4.0474	0.20608\\
4.0506	0.20459\\
4.0539	0.20313\\
4.0571	0.20167\\
4.0603	0.20023\\
4.0636	0.1988\\
4.0668	0.19739\\
4.07	0.19599\\
4.0733	0.19459\\
4.0765	0.19321\\
4.0798	0.19185\\
4.083	0.19049\\
4.0862	0.18914\\
4.0895	0.1878\\
4.0927	0.18648\\
4.0959	0.18516\\
4.0992	0.18385\\
4.1024	0.18256\\
4.1057	0.18127\\
4.1089	0.17999\\
4.1121	0.17872\\
4.1154	0.17746\\
4.1186	0.1762\\
4.1219	0.17496\\
4.1251	0.17372\\
4.1283	0.17249\\
4.1316	0.17126\\
4.1348	0.17005\\
4.138	0.16884\\
4.1413	0.16764\\
4.1445	0.16644\\
4.1478	0.16526\\
4.151	0.16408\\
4.1542	0.1629\\
4.1575	0.16173\\
4.1607	0.16057\\
4.1639	0.15941\\
4.1672	0.15826\\
4.1704	0.15712\\
4.1737	0.15598\\
4.1769	0.15484\\
4.1801	0.15371\\
4.1834	0.15259\\
4.1866	0.15147\\
4.1898	0.15036\\
4.1931	0.14925\\
4.1963	0.14814\\
4.1996	0.14705\\
4.2028	0.14595\\
4.206	0.14486\\
4.2093	0.14378\\
4.2125	0.1427\\
4.2157	0.14162\\
4.219	0.14055\\
4.2222	0.13948\\
4.2255	0.13842\\
4.2287	0.13736\\
4.2319	0.1363\\
4.2352	0.13525\\
4.2384	0.13421\\
4.2417	0.13316\\
4.2449	0.13213\\
4.2481	0.13109\\
4.2514	0.13006\\
4.2546	0.12903\\
4.2578	0.12801\\
4.2611	0.12699\\
4.2643	0.12598\\
4.2676	0.12497\\
4.2708	0.12396\\
4.274	0.12295\\
4.2773	0.12196\\
4.2805	0.12096\\
4.2837	0.11997\\
4.287	0.11898\\
4.2902	0.11799\\
4.2935	0.11701\\
4.2967	0.11604\\
4.2999	0.11506\\
4.3032	0.11409\\
4.3064	0.11313\\
4.3096	0.11217\\
4.3129	0.11121\\
4.3161	0.11026\\
4.3194	0.10931\\
4.3226	0.10836\\
4.3258	0.10742\\
4.3291	0.10648\\
4.3323	0.10554\\
4.3355	0.10461\\
4.3388	0.10369\\
4.342	0.10276\\
4.3453	0.10184\\
4.3485	0.10093\\
4.3517	0.10002\\
4.355	0.099111\\
4.3582	0.098208\\
4.3615	0.097309\\
4.3647	0.096414\\
4.3679	0.095523\\
4.3712	0.094637\\
4.3744	0.093754\\
4.3776	0.092876\\
4.3809	0.092002\\
4.3841	0.091132\\
4.3874	0.090266\\
4.3906	0.089404\\
4.3938	0.088547\\
4.3971	0.087694\\
4.4003	0.086846\\
4.4035	0.086002\\
4.4068	0.085162\\
4.41	0.084326\\
4.4133	0.083495\\
4.4165	0.082669\\
4.4197	0.081846\\
4.423	0.081029\\
4.4262	0.080215\\
4.4294	0.079407\\
4.4327	0.078603\\
4.4359	0.077803\\
4.4392	0.077008\\
4.4424	0.076218\\
4.4456	0.075432\\
4.4489	0.074651\\
4.4521	0.073875\\
4.4553	0.073103\\
4.4586	0.072336\\
4.4618	0.071574\\
4.4651	0.070816\\
4.4683	0.070063\\
4.4715	0.069315\\
4.4748	0.068572\\
4.478	0.067833\\
4.4813	0.0671\\
4.4845	0.066371\\
4.4877	0.065647\\
4.491	0.064928\\
4.4942	0.064214\\
4.4974	0.063504\\
4.5007	0.0628\\
4.5039	0.0621\\
4.5072	0.061405\\
4.5104	0.060716\\
4.5136	0.060031\\
4.5169	0.059351\\
4.5201	0.058676\\
4.5233	0.058006\\
4.5266	0.057341\\
4.5298	0.056681\\
4.5331	0.056025\\
4.5363	0.055375\\
4.5395	0.05473\\
4.5428	0.054089\\
4.546	0.053454\\
4.5492	0.052823\\
4.5525	0.052198\\
4.5557	0.051577\\
4.559	0.050962\\
4.5622	0.050351\\
4.5654	0.049745\\
4.5687	0.049144\\
4.5719	0.048548\\
4.5751	0.047957\\
4.5784	0.047371\\
4.5816	0.04679\\
4.5849	0.046214\\
4.5881	0.045643\\
4.5913	0.045076\\
4.5946	0.044515\\
4.5978	0.043958\\
4.6011	0.043406\\
4.6043	0.042859\\
4.6075	0.042317\\
4.6108	0.041779\\
4.614	0.041247\\
4.6172	0.040719\\
4.6205	0.040196\\
4.6237	0.039678\\
4.627	0.039165\\
4.6302	0.038656\\
4.6334	0.038152\\
4.6367	0.037653\\
4.6399	0.037158\\
4.6431	0.036669\\
4.6464	0.036184\\
4.6496	0.035703\\
4.6529	0.035227\\
4.6561	0.034756\\
4.6593	0.03429\\
4.6626	0.033828\\
4.6658	0.03337\\
4.669	0.032918\\
4.6723	0.032469\\
4.6755	0.032026\\
4.6788	0.031586\\
4.682	0.031152\\
4.6852	0.030721\\
4.6885	0.030296\\
4.6917	0.029874\\
4.6949	0.029457\\
4.6982	0.029044\\
4.7014	0.028636\\
4.7047	0.028232\\
4.7079	0.027833\\
4.7111	0.027437\\
4.7144	0.027046\\
4.7176	0.02666\\
4.7209	0.026277\\
4.7241	0.025899\\
4.7273	0.025524\\
4.7306	0.025154\\
4.7338	0.024788\\
4.737	0.024427\\
4.7403	0.024069\\
4.7435	0.023715\\
4.7468	0.023366\\
4.75	0.02302\\
4.7532	0.022678\\
4.7565	0.022341\\
4.7597	0.022007\\
4.7629	0.021677\\
4.7662	0.021351\\
4.7694	0.021029\\
4.7727	0.020711\\
4.7759	0.020396\\
4.7791	0.020086\\
4.7824	0.019779\\
4.7856	0.019476\\
4.7888	0.019176\\
4.7921	0.01888\\
4.7953	0.018588\\
4.7986	0.018299\\
4.8018	0.018014\\
4.805	0.017733\\
4.8083	0.017455\\
4.8115	0.01718\\
4.8147	0.016909\\
4.818	0.016642\\
4.8212	0.016378\\
4.8245	0.016117\\
4.8277	0.01586\\
4.8309	0.015606\\
4.8342	0.015355\\
4.8374	0.015107\\
4.8406	0.014863\\
4.8439	0.014622\\
4.8471	0.014384\\
4.8504	0.01415\\
4.8536	0.013918\\
4.8568	0.01369\\
4.8601	0.013464\\
4.8633	0.013242\\
4.8666	0.013023\\
4.8698	0.012806\\
4.873	0.012593\\
4.8763	0.012383\\
4.8795	0.012175\\
4.8827	0.011971\\
4.886	0.011769\\
4.8892	0.01157\\
4.8925	0.011374\\
4.8957	0.01118\\
4.8989	0.01099\\
4.9022	0.010802\\
4.9054	0.010617\\
4.9086	0.010434\\
4.9119	0.010254\\
4.9151	0.010077\\
4.9184	0.009902\\
4.9216	0.0097298\\
4.9248	0.0095601\\
4.9281	0.0093928\\
4.9313	0.0092281\\
4.9345	0.0090658\\
4.9378	0.0089059\\
4.941	0.0087483\\
4.9443	0.0085932\\
4.9475	0.0084403\\
4.9507	0.0082898\\
4.954	0.0081415\\
4.9572	0.0079955\\
4.9604	0.0078517\\
4.9637	0.0077101\\
4.9669	0.0075707\\
4.9702	0.0074335\\
4.9734	0.0072984\\
4.9766	0.0071653\\
4.9799	0.0070344\\
4.9831	0.0069055\\
4.9864	0.0067786\\
4.9896	0.0066537\\
4.9928	0.0065308\\
4.9961	0.0064098\\
4.9993	0.0062908\\
5.0025	0.0061736\\
5.0058	0.0060584\\
5.009	0.005945\\
5.0123	0.0058334\\
5.0155	0.0057237\\
5.0187	0.0056157\\
5.022	0.0055095\\
5.0252	0.005405\\
5.0284	0.0053022\\
5.0317	0.0052012\\
5.0349	0.0051018\\
5.0382	0.005004\\
5.0414	0.0049079\\
5.0446	0.0048134\\
5.0479	0.0047205\\
5.0511	0.0046291\\
5.0543	0.0045393\\
5.0576	0.004451\\
5.0608	0.0043642\\
5.0641	0.0042789\\
5.0673	0.004195\\
5.0705	0.0041126\\
5.0738	0.0040316\\
5.077	0.003952\\
5.0802	0.0038738\\
5.0835	0.0037969\\
5.0867	0.0037214\\
5.09	0.0036472\\
5.0932	0.0035743\\
5.0964	0.0035026\\
5.0997	0.0034323\\
5.1029	0.0033632\\
5.1062	0.0032953\\
5.1094	0.0032286\\
5.1126	0.0031631\\
5.1159	0.0030988\\
5.1191	0.0030357\\
5.1223	0.0029736\\
5.1256	0.0029127\\
5.1288	0.002853\\
5.1321	0.0027942\\
5.1353	0.0027366\\
5.1385	0.00268\\
5.1418	0.0026245\\
5.145	0.00257\\
5.1482	0.0025165\\
5.1515	0.002464\\
5.1547	0.0024124\\
5.158	0.0023619\\
5.1612	0.0023122\\
5.1644	0.0022635\\
5.1677	0.0022157\\
5.1709	0.0021688\\
5.1741	0.0021228\\
5.1774	0.0020777\\
5.1806	0.0020334\\
5.1839	0.00199\\
5.1871	0.0019474\\
5.1903	0.0019056\\
5.1936	0.0018647\\
5.1968	0.0018245\\
5.2	0.0017851\\
5.2033	0.0017464\\
5.2065	0.0017085\\
5.2098	0.0016714\\
5.213	0.001635\\
5.2162	0.0015992\\
5.2195	0.0015642\\
5.2227	0.0015299\\
5.226	0.0014963\\
5.2292	0.0014633\\
5.2324	0.001431\\
5.2357	0.0013993\\
5.2389	0.0013683\\
5.2421	0.0013378\\
5.2454	0.001308\\
5.2486	0.0012788\\
5.2519	0.0012502\\
5.2551	0.0012222\\
5.2583	0.0011947\\
5.2616	0.0011678\\
5.2648	0.0011414\\
5.268	0.0011156\\
5.2713	0.0010903\\
5.2745	0.0010656\\
5.2778	0.0010413\\
5.281	0.0010175\\
5.2842	0.00099427\\
5.2875	0.00097149\\
5.2907	0.00094919\\
5.2939	0.00092735\\
5.2972	0.00090596\\
5.3004	0.00088503\\
5.3037	0.00086454\\
5.3069	0.00084448\\
5.3101	0.00082484\\
5.3101	0.00078458\\
5.3069	0.00080347\\
5.3037	0.00082277\\
5.3004	0.0008425\\
5.2972	0.00086265\\
5.2939	0.00088324\\
5.2907	0.00090428\\
5.2875	0.00092577\\
5.2842	0.00094772\\
5.281	0.00097015\\
5.2778	0.00099305\\
5.2745	0.0010165\\
5.2713	0.0010403\\
5.268	0.0010647\\
5.2648	0.0010897\\
5.2616	0.0011151\\
5.2583	0.0011411\\
5.2551	0.0011676\\
5.2519	0.0011947\\
5.2486	0.0012224\\
5.2454	0.0012506\\
5.2421	0.0012794\\
5.2389	0.0013089\\
5.2357	0.0013389\\
5.2324	0.0013695\\
5.2292	0.0014008\\
5.226	0.0014327\\
5.2227	0.0014653\\
5.2195	0.0014985\\
5.2162	0.0015325\\
5.213	0.0015671\\
5.2098	0.0016024\\
5.2065	0.0016384\\
5.2033	0.0016751\\
5.2	0.0017126\\
5.1968	0.0017508\\
5.1936	0.0017898\\
5.1903	0.0018296\\
5.1871	0.0018702\\
5.1839	0.0019115\\
5.1806	0.0019537\\
5.1774	0.0019967\\
5.1741	0.0020406\\
5.1709	0.0020853\\
5.1677	0.0021309\\
5.1644	0.0021774\\
5.1612	0.0022248\\
5.158	0.002273\\
5.1547	0.0023223\\
5.1515	0.0023724\\
5.1482	0.0024236\\
5.145	0.0024757\\
5.1418	0.0025288\\
5.1385	0.0025829\\
5.1353	0.002638\\
5.1321	0.0026942\\
5.1288	0.0027514\\
5.1256	0.0028097\\
5.1223	0.0028691\\
5.1191	0.0029297\\
5.1159	0.0029913\\
5.1126	0.0030541\\
5.1094	0.003118\\
5.1062	0.0031831\\
5.1029	0.0032494\\
5.0997	0.0033169\\
5.0964	0.0033857\\
5.0932	0.0034557\\
5.09	0.003527\\
5.0867	0.0035996\\
5.0835	0.0036734\\
5.0802	0.0037486\\
5.077	0.0038252\\
5.0738	0.0039031\\
5.0705	0.0039824\\
5.0673	0.0040631\\
5.0641	0.0041452\\
5.0608	0.0042288\\
5.0576	0.0043138\\
5.0543	0.0044004\\
5.0511	0.0044884\\
5.0479	0.004578\\
5.0446	0.0046691\\
5.0414	0.0047618\\
5.0382	0.0048561\\
5.0349	0.004952\\
5.0317	0.0050496\\
5.0284	0.0051488\\
5.0252	0.0052497\\
5.022	0.0053523\\
5.0187	0.0054566\\
5.0155	0.0055627\\
5.0123	0.0056706\\
5.009	0.0057803\\
5.0058	0.0058918\\
5.0025	0.0060051\\
4.9993	0.0061203\\
4.9961	0.0062374\\
4.9928	0.0063564\\
4.9896	0.0064774\\
4.9864	0.0066003\\
4.9831	0.0067252\\
4.9799	0.0068522\\
4.9766	0.0069812\\
4.9734	0.0071122\\
4.9702	0.0072454\\
4.9669	0.0073807\\
4.9637	0.0075181\\
4.9604	0.0076577\\
4.9572	0.0077995\\
4.954	0.0079435\\
4.9507	0.0080897\\
4.9475	0.0082383\\
4.9443	0.0083891\\
4.941	0.0085423\\
4.9378	0.0086978\\
4.9345	0.0088557\\
4.9313	0.0090161\\
4.9281	0.0091788\\
4.9248	0.0093441\\
4.9216	0.0095118\\
4.9184	0.009682\\
4.9151	0.0098548\\
4.9119	0.01003\\
4.9086	0.010208\\
4.9054	0.010389\\
4.9022	0.010572\\
4.8989	0.010758\\
4.8957	0.010947\\
4.8925	0.011138\\
4.8892	0.011332\\
4.886	0.011529\\
4.8827	0.011729\\
4.8795	0.011932\\
4.8763	0.012137\\
4.873	0.012346\\
4.8698	0.012557\\
4.8666	0.012772\\
4.8633	0.012989\\
4.8601	0.01321\\
4.8568	0.013433\\
4.8536	0.01366\\
4.8504	0.013889\\
4.8471	0.014122\\
4.8439	0.014358\\
4.8406	0.014598\\
4.8374	0.01484\\
4.8342	0.015086\\
4.8309	0.015335\\
4.8277	0.015587\\
4.8245	0.015843\\
4.8212	0.016102\\
4.818	0.016365\\
4.8147	0.016631\\
4.8115	0.0169\\
4.8083	0.017173\\
4.805	0.017449\\
4.8018	0.017729\\
4.7986	0.018013\\
4.7953	0.0183\\
4.7921	0.018591\\
4.7888	0.018886\\
4.7856	0.019184\\
4.7824	0.019486\\
4.7791	0.019791\\
4.7759	0.020101\\
4.7727	0.020414\\
4.7694	0.020731\\
4.7662	0.021052\\
4.7629	0.021377\\
4.7597	0.021705\\
4.7565	0.022038\\
4.7532	0.022375\\
4.75	0.022715\\
4.7468	0.02306\\
4.7435	0.023409\\
4.7403	0.023761\\
4.737	0.024118\\
4.7338	0.024479\\
4.7306	0.024844\\
4.7273	0.025213\\
4.7241	0.025587\\
4.7209	0.025964\\
4.7176	0.026346\\
4.7144	0.026732\\
4.7111	0.027123\\
4.7079	0.027518\\
4.7047	0.027917\\
4.7014	0.02832\\
4.6982	0.028728\\
4.6949	0.02914\\
4.6917	0.029557\\
4.6885	0.029978\\
4.6852	0.030403\\
4.682	0.030833\\
4.6788	0.031267\\
4.6755	0.031706\\
4.6723	0.03215\\
4.669	0.032598\\
4.6658	0.03305\\
4.6626	0.033508\\
4.6593	0.033969\\
4.6561	0.034436\\
4.6529	0.034907\\
4.6496	0.035382\\
4.6464	0.035863\\
4.6431	0.036348\\
4.6399	0.036837\\
4.6367	0.037332\\
4.6334	0.037831\\
4.6302	0.038334\\
4.627	0.038843\\
4.6237	0.039356\\
4.6205	0.039874\\
4.6172	0.040397\\
4.614	0.040924\\
4.6108	0.041457\\
4.6075	0.041994\\
4.6043	0.042535\\
4.6011	0.043082\\
4.5978	0.043634\\
4.5946	0.04419\\
4.5913	0.044751\\
4.5881	0.045317\\
4.5849	0.045887\\
4.5816	0.046463\\
4.5784	0.047043\\
4.5751	0.047628\\
4.5719	0.048218\\
4.5687	0.048813\\
4.5654	0.049413\\
4.5622	0.050017\\
4.559	0.050626\\
4.5557	0.05124\\
4.5525	0.051859\\
4.5492	0.052483\\
4.546	0.053111\\
4.5428	0.053744\\
4.5395	0.054382\\
4.5363	0.055025\\
4.5331	0.055672\\
4.5298	0.056325\\
4.5266	0.056982\\
4.5233	0.057644\\
4.5201	0.05831\\
4.5169	0.058981\\
4.5136	0.059657\\
4.5104	0.060338\\
4.5072	0.061023\\
4.5039	0.061713\\
4.5007	0.062408\\
4.4974	0.063107\\
4.4942	0.063812\\
4.491	0.06452\\
4.4877	0.065233\\
4.4845	0.065951\\
4.4813	0.066674\\
4.478	0.067401\\
4.4748	0.068133\\
4.4715	0.068869\\
4.4683	0.06961\\
4.4651	0.070355\\
4.4618	0.071105\\
4.4586	0.071859\\
4.4553	0.072618\\
4.4521	0.073381\\
4.4489	0.074149\\
4.4456	0.074921\\
4.4424	0.075697\\
4.4392	0.076478\\
4.4359	0.077263\\
4.4327	0.078053\\
4.4294	0.078847\\
4.4262	0.079645\\
4.423	0.080447\\
4.4197	0.081254\\
4.4165	0.082065\\
4.4133	0.082881\\
4.41	0.0837\\
4.4068	0.084524\\
4.4035	0.085352\\
4.4003	0.086184\\
4.3971	0.08702\\
4.3938	0.08786\\
4.3906	0.088705\\
4.3874	0.089553\\
4.3841	0.090406\\
4.3809	0.091262\\
4.3776	0.092123\\
4.3744	0.092988\\
4.3712	0.093856\\
4.3679	0.094729\\
4.3647	0.095605\\
4.3615	0.096486\\
4.3582	0.09737\\
4.355	0.098258\\
4.3517	0.09915\\
4.3485	0.10005\\
4.3453	0.10095\\
4.342	0.10185\\
4.3388	0.10276\\
4.3355	0.10367\\
4.3323	0.10458\\
4.3291	0.1055\\
4.3258	0.10643\\
4.3226	0.10735\\
4.3194	0.10828\\
4.3161	0.10922\\
4.3129	0.11015\\
4.3096	0.1111\\
4.3064	0.11204\\
4.3032	0.11299\\
4.2999	0.11394\\
4.2967	0.1149\\
4.2935	0.11586\\
4.2902	0.11682\\
4.287	0.11779\\
4.2837	0.11876\\
4.2805	0.11974\\
4.2773	0.12072\\
4.274	0.1217\\
4.2708	0.12268\\
4.2676	0.12368\\
4.2643	0.12467\\
4.2611	0.12567\\
4.2578	0.12667\\
4.2546	0.12768\\
4.2514	0.12868\\
4.2481	0.1297\\
4.2449	0.13072\\
4.2417	0.13174\\
4.2384	0.13276\\
4.2352	0.13379\\
4.2319	0.13483\\
4.2287	0.13587\\
4.2255	0.13691\\
4.2222	0.13796\\
4.219	0.13901\\
4.2157	0.14006\\
4.2125	0.14112\\
4.2093	0.14219\\
4.206	0.14326\\
4.2028	0.14433\\
4.1996	0.14541\\
4.1963	0.14649\\
4.1931	0.14758\\
4.1898	0.14868\\
4.1866	0.14978\\
4.1834	0.15088\\
4.1801	0.15199\\
4.1769	0.1531\\
4.1737	0.15422\\
4.1704	0.15535\\
4.1672	0.15648\\
4.1639	0.15762\\
4.1607	0.15876\\
4.1575	0.15991\\
4.1542	0.16107\\
4.151	0.16223\\
4.1478	0.1634\\
4.1445	0.16458\\
4.1413	0.16576\\
4.138	0.16695\\
4.1348	0.16815\\
4.1316	0.16936\\
4.1283	0.17057\\
4.1251	0.17179\\
4.1219	0.17302\\
4.1186	0.17426\\
4.1154	0.1755\\
4.1121	0.17676\\
4.1089	0.17802\\
4.1057	0.17929\\
4.1024	0.18057\\
4.0992	0.18186\\
4.0959	0.18316\\
4.0927	0.18447\\
4.0895	0.18579\\
4.0862	0.18713\\
4.083	0.18847\\
4.0798	0.18982\\
4.0765	0.19119\\
4.0733	0.19256\\
4.07	0.19395\\
4.0668	0.19535\\
4.0636	0.19676\\
4.0603	0.19819\\
4.0571	0.19963\\
4.0539	0.20108\\
4.0506	0.20255\\
4.0474	0.20403\\
4.0441	0.20552\\
4.0409	0.20703\\
4.0377	0.20856\\
4.0344	0.2101\\
4.0312	0.21165\\
4.028	0.21323\\
4.0247	0.21482\\
4.0215	0.21642\\
4.0182	0.21805\\
4.015	0.21969\\
4.0118	0.22135\\
4.0085	0.22302\\
4.0053	0.22472\\
4.0021	0.22644\\
3.9988	0.22817\\
3.9956	0.22993\\
3.9923	0.23171\\
3.9891	0.2335\\
3.9859	0.23532\\
3.9826	0.23717\\
3.9794	0.23903\\
3.9762	0.24092\\
3.9729	0.24283\\
3.9697	0.24476\\
3.9664	0.24672\\
3.9632	0.2487\\
3.96	0.25071\\
3.9567	0.25274\\
3.9535	0.2548\\
3.9502	0.25689\\
3.947	0.259\\
3.9438	0.26114\\
3.9405	0.26331\\
3.9373	0.26551\\
3.9341	0.26773\\
3.9308	0.26999\\
3.9276	0.27228\\
3.9243	0.27459\\
3.9211	0.27694\\
3.9179	0.27932\\
3.9146	0.28173\\
3.9114	0.28417\\
3.9082	0.28665\\
3.9049	0.28916\\
3.9017	0.2917\\
3.8984	0.29428\\
3.8952	0.29689\\
3.892	0.29954\\
3.8887	0.30223\\
3.8855	0.30495\\
3.8823	0.30771\\
3.879	0.3105\\
3.8758	0.31334\\
3.8725	0.31621\\
3.8693	0.31912\\
3.8661	0.32208\\
3.8628	0.32507\\
3.8596	0.3281\\
3.8564	0.33117\\
3.8531	0.33429\\
3.8499	0.33745\\
3.8466	0.34065\\
3.8434	0.34389\\
3.8402	0.34718\\
3.8369	0.35051\\
3.8337	0.35388\\
3.8304	0.3573\\
3.8272	0.36077\\
3.824	0.36428\\
3.8207	0.36784\\
3.8175	0.37144\\
3.8143	0.37509\\
3.811	0.37879\\
3.8078	0.38253\\
3.8045	0.38632\\
3.8013	0.39017\\
3.7981	0.39405\\
3.7948	0.39799\\
3.7916	0.40198\\
3.7884	0.40602\\
3.7851	0.4101\\
3.7819	0.41424\\
3.7786	0.41843\\
3.7754	0.42266\\
3.7722	0.42695\\
3.7689	0.43129\\
3.7657	0.43567\\
3.7625	0.44011\\
3.7592	0.4446\\
3.756	0.44914\\
3.7527	0.45374\\
3.7495	0.45838\\
3.7463	0.46307\\
3.743	0.46782\\
3.7398	0.47261\\
3.7366	0.47746\\
3.7333	0.48236\\
3.7301	0.48731\\
3.7268	0.4923\\
3.7236	0.49735\\
3.7204	0.50245\\
3.7171	0.5076\\
3.7139	0.5128\\
3.7106	0.51805\\
3.7074	0.52335\\
3.7042	0.52869\\
3.7009	0.53409\\
3.6977	0.53953\\
3.6945	0.54502\\
3.6912	0.55056\\
3.688	0.55614\\
3.6847	0.56178\\
3.6815	0.56745\\
3.6783	0.57317\\
3.675	0.57894\\
3.6718	0.58475\\
3.6686	0.5906\\
3.6653	0.59649\\
3.6621	0.60243\\
3.6588	0.60841\\
3.6556	0.61442\\
3.6524	0.62048\\
3.6491	0.62658\\
3.6459	0.63271\\
3.6427	0.63888\\
3.6394	0.64508\\
3.6362	0.65132\\
3.6329	0.65759\\
3.6297	0.6639\\
3.6265	0.67023\\
3.6232	0.6766\\
3.62	0.68299\\
3.6168	0.68942\\
3.6135	0.69587\\
3.6103	0.70234\\
3.607	0.70884\\
3.6038	0.71536\\
3.6006	0.72191\\
3.5973	0.72847\\
3.5941	0.73505\\
3.5908	0.74165\\
3.5876	0.74827\\
3.5844	0.75489\\
3.5811	0.76154\\
3.5779	0.76819\\
3.5747	0.77485\\
3.5714	0.78152\\
3.5682	0.7882\\
3.5649	0.79488\\
3.5617	0.80156\\
3.5585	0.80824\\
3.5552	0.81493\\
3.552	0.82161\\
3.5488	0.82828\\
3.5455	0.83496\\
3.5423	0.84162\\
3.539	0.84827\\
3.5358	0.85491\\
3.5326	0.86154\\
3.5293	0.86815\\
3.5261	0.87474\\
3.5229	0.88132\\
3.5196	0.88787\\
3.5164	0.8944\\
3.5131	0.90091\\
3.5099	0.90738\\
3.5067	0.91383\\
3.5034	0.92024\\
3.5002	0.92663\\
3.497	0.93297\\
3.4937	0.93928\\
3.4905	0.94555\\
3.4872	0.95177\\
3.484	0.95795\\
3.4808	0.96409\\
3.4775	0.97017\\
3.4743	0.97621\\
3.471	0.98219\\
3.4678	0.98812\\
3.4646	0.99399\\
3.4613	0.9998\\
3.4581	1.0055\\
3.4549	1.0112\\
3.4516	1.0169\\
3.4484	1.0224\\
3.4451	1.0279\\
3.4419	1.0333\\
3.4387	1.0386\\
3.4354	1.0439\\
3.4322	1.0491\\
3.429	1.0542\\
3.4257	1.0592\\
3.4225	1.0642\\
3.4192	1.069\\
3.416	1.0738\\
3.4128	1.0784\\
3.4095	1.083\\
3.4063	1.0875\\
3.4031	1.0919\\
3.3998	1.0962\\
3.3966	1.1004\\
3.3933	1.1045\\
3.3901	1.1085\\
3.3869	1.1124\\
3.3836	1.1162\\
3.3804	1.1199\\
3.3772	1.1234\\
3.3739	1.1269\\
3.3707	1.1302\\
3.3674	1.1335\\
3.3642	1.1366\\
3.361	1.1396\\
3.3577	1.1425\\
3.3545	1.1452\\
3.3513	1.1479\\
3.348	1.1504\\
3.3448	1.1528\\
3.3415	1.1551\\
3.3383	1.1573\\
3.3351	1.1593\\
3.3318	1.1612\\
3.3286	1.163\\
3.3253	1.1646\\
3.3221	1.1661\\
3.3189	1.1675\\
3.3156	1.1688\\
3.3124	1.1699\\
3.3092	1.1709\\
3.3059	1.1718\\
3.3027	1.1725\\
3.2994	1.1731\\
3.2962	1.1735\\
3.293	1.1739\\
3.2897	1.1741\\
3.2865	1.1741\\
3.2833	1.174\\
3.28	1.1738\\
3.2768	1.1735\\
3.2735	1.173\\
3.2703	1.1723\\
3.2671	1.1716\\
3.2638	1.1707\\
3.2606	1.1696\\
3.2574	1.1685\\
3.2541	1.1672\\
3.2509	1.1657\\
3.2476	1.1641\\
3.2444	1.1624\\
3.2412	1.1606\\
3.2379	1.1586\\
3.2347	1.1565\\
3.2315	1.1543\\
3.2282	1.1519\\
3.225	1.1494\\
3.2217	1.1467\\
3.2185	1.1439\\
3.2153	1.141\\
3.212	1.138\\
3.2088	1.1349\\
3.2055	1.1316\\
3.2023	1.1282\\
3.1991	1.1246\\
3.1958	1.121\\
3.1926	1.1172\\
3.1894	1.1133\\
3.1861	1.1093\\
3.1829	1.1051\\
3.1796	1.1009\\
3.1764	1.0965\\
3.1732	1.092\\
3.1699	1.0874\\
3.1667	1.0827\\
3.1635	1.0779\\
3.1602	1.0729\\
3.157	1.0679\\
3.1537	1.0627\\
3.1505	1.0575\\
3.1473	1.0521\\
3.144	1.0467\\
3.1408	1.0411\\
3.1376	1.0355\\
3.1343	1.0297\\
3.1311	1.0239\\
3.1278	1.018\\
3.1246	1.012\\
3.1214	1.0059\\
3.1181	0.99966\\
3.1149	0.99339\\
3.1117	0.98703\\
3.1084	0.98059\\
3.1052	0.97407\\
3.1019	0.96748\\
3.0987	0.96081\\
3.0955	0.95408\\
3.0922	0.94727\\
3.089	0.94039\\
3.0857	0.93345\\
3.0825	0.92645\\
3.0793	0.91938\\
3.076	0.91226\\
3.0728	0.90508\\
3.0696	0.89784\\
3.0663	0.89055\\
3.0631	0.88321\\
3.0598	0.87582\\
3.0566	0.86839\\
3.0534	0.86091\\
3.0501	0.85339\\
3.0469	0.84583\\
3.0437	0.83823\\
3.0404	0.8306\\
3.0372	0.82293\\
3.0339	0.81523\\
3.0307	0.80751\\
3.0275	0.79976\\
3.0242	0.79198\\
3.021	0.78418\\
3.0178	0.77636\\
3.0145	0.76852\\
3.0113	0.76066\\
3.008	0.75279\\
3.0048	0.74491\\
3.0016	0.73702\\
2.9983	0.72912\\
2.9951	0.72122\\
2.9919	0.71331\\
2.9886	0.70539\\
2.9854	0.69748\\
2.9821	0.68957\\
2.9789	0.68167\\
2.9757	0.67377\\
2.9724	0.66587\\
2.9692	0.65799\\
2.9659	0.65012\\
2.9627	0.64226\\
2.9595	0.63442\\
2.9562	0.62659\\
2.953	0.61878\\
2.9498	0.61099\\
2.9465	0.60322\\
2.9433	0.59548\\
2.94	0.58776\\
2.9368	0.58007\\
2.9336	0.57241\\
2.9303	0.56477\\
2.9271	0.55717\\
2.9239	0.5496\\
2.9206	0.54207\\
2.9174	0.53457\\
2.9141	0.52711\\
2.9109	0.51968\\
2.9077	0.5123\\
2.9044	0.50496\\
2.9012	0.49766\\
2.898	0.4904\\
2.8947	0.48319\\
2.8915	0.47602\\
2.8882	0.46891\\
2.885	0.46183\\
2.8818	0.45481\\
2.8785	0.44784\\
2.8753	0.44092\\
2.8721	0.43406\\
2.8688	0.42724\\
2.8656	0.42048\\
2.8623	0.41378\\
2.8591	0.40713\\
2.8559	0.40053\\
2.8526	0.394\\
2.8494	0.38752\\
2.8461	0.3811\\
2.8429	0.37474\\
2.8397	0.36844\\
2.8364	0.3622\\
2.8332	0.35602\\
2.83	0.34991\\
2.8267	0.34385\\
2.8235	0.33786\\
2.8202	0.33193\\
2.817	0.32607\\
2.8138	0.32026\\
2.8105	0.31453\\
2.8073	0.30885\\
2.8041	0.30324\\
2.8008	0.2977\\
2.7976	0.29222\\
2.7943	0.2868\\
2.7911	0.28145\\
2.7879	0.27617\\
2.7846	0.27095\\
2.7814	0.2658\\
2.7782	0.26071\\
2.7749	0.25569\\
2.7717	0.25074\\
2.7684	0.24585\\
2.7652	0.24102\\
2.762	0.23626\\
2.7587	0.23157\\
2.7555	0.22694\\
2.7523	0.22238\\
2.749	0.21788\\
2.7458	0.21344\\
2.7425	0.20908\\
2.7393	0.20477\\
2.7361	0.20053\\
2.7328	0.19635\\
2.7296	0.19224\\
2.7263	0.18819\\
2.7231	0.1842\\
2.7199	0.18028\\
2.7166	0.17641\\
2.7134	0.17261\\
2.7102	0.16887\\
2.7069	0.16519\\
2.7037	0.16157\\
2.7004	0.15802\\
2.6972	0.15452\\
2.694	0.15108\\
2.6907	0.1477\\
2.6875	0.14437\\
2.6843	0.14111\\
2.681	0.1379\\
2.6778	0.13475\\
2.6745	0.13165\\
2.6713	0.12862\\
2.6681	0.12563\\
2.6648	0.1227\\
2.6616	0.11983\\
2.6584	0.117\\
2.6551	0.11423\\
2.6519	0.11152\\
2.6486	0.10885\\
2.6454	0.10624\\
2.6422	0.10367\\
2.6389	0.10116\\
2.6357	0.098696\\
2.6325	0.09628\\
2.6292	0.093912\\
2.626	0.091591\\
2.6227	0.089318\\
2.6195	0.087091\\
2.6163	0.084909\\
2.613	0.082772\\
2.6098	0.08068\\
2.6066	0.078631\\
2.6033	0.076626\\
2.6001	0.074663\\
2.5968	0.072742\\
2.5936	0.070862\\
2.5904	0.069023\\
2.5871	0.067224\\
2.5839	0.065464\\
2.5806	0.063743\\
2.5774	0.06206\\
2.5742	0.060415\\
2.5709	0.058807\\
2.5677	0.057235\\
2.5645	0.055699\\
2.5612	0.054197\\
2.558	0.052731\\
2.5547	0.051298\\
2.5515	0.049899\\
2.5483	0.048532\\
2.545	0.047198\\
2.5418	0.045895\\
2.5386	0.044623\\
2.5353	0.043382\\
2.5321	0.042171\\
2.5288	0.040989\\
2.5256	0.039835\\
2.5224	0.03871\\
2.5191	0.037613\\
2.5159	0.036543\\
2.5127	0.035499\\
2.5094	0.034482\\
2.5062	0.03349\\
2.5029	0.032523\\
2.4997	0.031581\\
2.4965	0.030663\\
2.4932	0.029768\\
2.49	0.028896\\
2.4868	0.028047\\
2.4835	0.02722\\
2.4803	0.026415\\
2.477	0.025631\\
2.4738	0.024868\\
2.4706	0.024125\\
2.4673	0.023401\\
2.4641	0.022697\\
2.4608	0.022013\\
2.4576	0.021346\\
2.4544	0.020698\\
2.4511	0.020067\\
2.4479	0.019454\\
2.4447	0.018857\\
2.4414	0.018277\\
2.4382	0.017713\\
2.4349	0.017165\\
2.4317	0.016633\\
2.4285	0.016115\\
2.4252	0.015612\\
2.422	0.015123\\
2.4188	0.014648\\
2.4155	0.014186\\
2.4123	0.013738\\
2.409	0.013303\\
2.4058	0.012881\\
2.4026	0.01247\\
2.3993	0.012072\\
2.3961	0.011685\\
2.3929	0.01131\\
2.3896	0.010946\\
2.3864	0.010592\\
2.3831	0.010249\\
2.3799	0.0099168\\
2.3767	0.0095941\\
2.3734	0.0092812\\
2.3702	0.0089776\\
2.367	0.0086833\\
2.3637	0.0083979\\
2.3605	0.0081212\\
2.3572	0.007853\\
2.354	0.007593\\
2.3508	0.0073411\\
2.3475	0.0070969\\
2.3443	0.0068603\\
2.341	0.0066311\\
2.3378	0.0064091\\
2.3346	0.006194\\
2.3313	0.0059857\\
2.3281	0.0057839\\
2.3249	0.0055885\\
2.3216	0.0053994\\
2.3184	0.0052162\\
2.3151	0.0050389\\
2.3119	0.0048673\\
2.3087	0.0047012\\
2.3054	0.0045405\\
2.3022	0.0043849\\
2.299	0.0042344\\
2.2957	0.0040888\\
2.2925	0.003948\\
2.2892	0.0038117\\
2.286	0.0036799\\
2.2828	0.0035525\\
2.2795	0.0034292\\
2.2763	0.00331\\
2.2731	0.0031948\\
2.2698	0.0030834\\
2.2666	0.0029757\\
2.2633	0.0028716\\
2.2601	0.002771\\
2.2569	0.0026738\\
2.2536	0.0025798\\
2.2504	0.0024891\\
2.2472	0.0024013\\
2.2439	0.0023166\\
2.2407	0.0022347\\
2.2374	0.0021557\\
2.2342	0.0020793\\
2.231	0.0020055\\
2.2277	0.0019343\\
2.2245	0.0018655\\
2.2212	0.0017991\\
2.218	0.001735\\
2.2148	0.0016731\\
2.2115	0.0016133\\
2.2083	0.0015556\\
2.2051	0.0015\\
2.2018	0.0014462\\
2.1986	0.0013944\\
2.1953	0.0013443\\
2.1921	0.001296\\
2.1889	0.0012494\\
2.1856	0.0012045\\
2.1824	0.0011611\\
2.1792	0.0011193\\
2.1759	0.0010789\\
2.1727	0.00104\\
2.1694	0.0010024\\
2.1662	0.00096622\\
2.163	0.0009313\\
2.1597	0.00089762\\
2.1565	0.00086514\\
2.1533	0.00083382\\
2.15	0.00080363\\
2.1468	0.00077452\\
2.1435	0.00074645\\
2.1403	0.00071939\\
2.1371	0.00069331\\
2.1338	0.00066817\\
2.1306	0.00064393\\
2.1274	0.00062057\\
2.1241	0.00059805\\
2.1209	0.00057635\\
2.1176	0.00055544\\
2.1144	0.00053528\\
2.1112	0.00051585\\
2.1079	0.00049714\\
2.1047	0.0004791\\
2.1014	0.00046172\\
2.0982	0.00044497\\
2.095	0.00042883\\
2.0917	0.00041328\\
2.0885	0.0003983\\
2.0853	0.00038387\\
2.082	0.00036996\\
2.0788	0.00035656\\
2.0755	0.00034365\\
}--cycle;
\addplot[ybar interval, fill=black, fill opacity=0.35, area legend, draw=none] table[row sep=crcr, x=Lower, y=Count] {%
Lower	Upper	Count\\
2.0755	2.3244	0.075022\\
2.3244	2.5732	0.13933\\
2.5732	2.822	0.35368\\
2.822	3.0708	0.76094\\
3.0708	3.3196	0.94314\\
3.3196	3.5684	0.75022\\
3.5684	3.8172	0.52516\\
3.8172	4.0661	0.2465\\
4.0661	4.3149	0.10717\\
4.3149	4.5637	0.053587\\
4.5637	4.8125	0.04287\\
4.8125	5.0613	0\\
5.0613	5.3101	0.021435\\
5.3101	5.3101	0.021435\\
};
\addlegendentry{\RealHistogramLegend}

\addplot [color=black]
  table[row sep=crcr]{%
2.0755	0.016864\\
2.1112	0.021504\\
2.1435	0.026645\\
2.1727	0.03214\\
2.2018	0.038572\\
2.2277	0.045168\\
2.2536	0.05268\\
2.2795	0.061193\\
2.3022	0.069533\\
2.3249	0.078766\\
2.3475	0.088951\\
2.3702	0.10014\\
2.3929	0.11239\\
2.4155	0.12576\\
2.4382	0.14028\\
2.4608	0.15599\\
2.4835	0.17292\\
2.5062	0.19111\\
2.5288	0.21056\\
2.5515	0.23127\\
2.5742	0.25324\\
2.5968	0.27644\\
2.6195	0.30084\\
2.6454	0.33012\\
2.6713	0.3608\\
2.7004	0.39681\\
2.7328	0.43844\\
2.7684	0.48576\\
2.8202	0.55625\\
2.885	0.6442\\
2.9206	0.69095\\
2.9498	0.7276\\
2.9757	0.75858\\
2.9983	0.78418\\
3.0178	0.80486\\
3.0372	0.82423\\
3.0566	0.84218\\
3.0728	0.85596\\
3.089	0.86863\\
3.1052	0.88011\\
3.1214	0.89036\\
3.1376	0.89934\\
3.1505	0.90558\\
3.1635	0.91096\\
3.1764	0.91546\\
3.1894	0.91909\\
3.2023	0.92181\\
3.2153	0.92364\\
3.2282	0.92456\\
3.2412	0.92457\\
3.2541	0.92368\\
3.2671	0.92189\\
3.28	0.91921\\
3.293	0.91564\\
3.3059	0.91119\\
3.3189	0.90589\\
3.3318	0.89975\\
3.348	0.89093\\
3.3642	0.88086\\
3.3804	0.86961\\
3.3966	0.85722\\
3.4128	0.84375\\
3.4322	0.82626\\
3.4516	0.80742\\
3.471	0.78736\\
3.4937	0.76256\\
3.5196	0.73266\\
3.5488	0.6974\\
3.5811	0.65673\\
3.6297	0.59406\\
3.7042	0.49801\\
3.7398	0.45359\\
3.7722	0.4147\\
3.8013	0.38119\\
3.8272	0.35274\\
3.8531	0.32567\\
3.879	0.30006\\
3.9017	0.27888\\
3.9243	0.25887\\
3.947	0.24004\\
3.9697	0.22239\\
3.9923	0.20591\\
4.015	0.19056\\
4.0377	0.17632\\
4.0603	0.16315\\
4.083	0.151\\
4.1057	0.13982\\
4.1283	0.12956\\
4.1542	0.1189\\
4.1801	0.10927\\
4.206	0.10061\\
4.2352	0.091914\\
4.2643	0.084204\\
4.2967	0.076663\\
4.3291	0.070062\\
4.3647	0.063729\\
4.4035	0.057748\\
4.4489	0.051762\\
4.4974	0.046278\\
4.5557	0.04065\\
4.6237	0.035031\\
4.7047	0.029269\\
4.8018	0.023311\\
4.9054	0.017882\\
5.0123	0.013203\\
5.1191	0.0094246\\
5.2324	0.006341\\
5.3101	0.0047195\\
};
\addlegendentry{\RealBLNFitLegend}

      % \input{../../../../.tex/20/KS/SizeDistributionFittingsfig2.tex}
  
      \nextgroupplot[%
        xmin=0.85621,xmax=5.4115,
        % xmin=1.9138,xmax=5.4719,
        % ymin=0,ymax=1.2,
      ] % This file was created by matlab2tikz.
%
\definecolor{mycolor1}{rgb}{0.83529,0.36863,0.00000}%
%
\addplot[ybar interval, fill=mycolor1, fill opacity=0.35, area legend, draw=none] table[row sep=crcr, x=Lower, y=Count] {%
Lower	Upper	Count\\
2.0755	2.3244	0\\
2.3244	2.5732	0.0224\\
2.5732	2.822	0.18788\\
2.822	3.0708	0.60039\\
3.0708	3.3196	1.0797\\
3.3196	3.5684	1.0557\\
3.5684	3.8172	0.49311\\
3.8172	4.0661	0.30598\\
4.0661	4.3149	0.16173\\
4.3149	4.5637	0.07095\\
4.5637	4.8125	0.029152\\
4.8125	5.0613	0.0013933\\
5.0613	5.3101	0.010717\\
5.3101	5.3101	0.010717\\
};
\addlegendentry{\ApproxHistogramLegend}

\addplot [color=mycolor1, only marks, every error bar/.append style={opacity=0.45}, mark=*, mark size=0pt, draw=none, forget plot]
 plot [error bars/.cd, y dir=both, y explicit, error bar style={line width=1pt, color=mycolor1}, error mark options={mark=none,mark size=0pt}]
 table[row sep=crcr, y error plus index=2, y error minus index=3]{%
2.2	0	0	0\\
2.4488	0.0224	0.0021606	0.0021606\\
2.6976	0.18788	0.0042688	0.0042688\\
2.9464	0.60039	0.008116	0.008116\\
3.1952	1.0797	0.0098484	0.0098484\\
3.444	1.0557	0.0088304	0.0088304\\
3.6928	0.49311	0.0073822	0.0073822\\
3.9417	0.30598	0.0046765	0.0046765\\
4.1905	0.16173	0.003128	0.003128\\
4.4393	0.07095	0.0020673	0.0020673\\
4.6881	0.029152	0.0015729	0.0015729\\
4.9369	0.0013933	0.00071878	0.00071878\\
5.1857	0.010717	2.4216e-18	1.7347e-18\\
};
\addplot [color=mycolor1]
  table[row sep=crcr]{%
2.0755	0.00039136\\
2.1889	0.0013257\\
2.2569	0.0027649\\
2.3054	0.004639\\
2.3443	0.0069624\\
2.3799	0.01002\\
2.409	0.013407\\
2.4349	0.017269\\
2.4608	0.022118\\
2.4835	0.027328\\
2.5029	0.032635\\
2.5224	0.038829\\
2.5418	0.046025\\
2.5612	0.054343\\
2.5774	0.062223\\
2.5936	0.071049\\
2.6098	0.080896\\
2.626	0.091845\\
2.6422	0.10397\\
2.6584	0.11736\\
2.6745	0.13208\\
2.6907	0.1482\\
2.7069	0.16578\\
2.7231	0.1849\\
2.7393	0.20559\\
2.7555	0.2279\\
2.7717	0.25186\\
2.7879	0.27747\\
2.8041	0.30474\\
2.8202	0.33365\\
2.8397	0.37046\\
2.8591	0.40948\\
2.8785	0.45056\\
2.9012	0.50084\\
2.9239	0.55331\\
2.953	0.62322\\
3.0016	0.74282\\
3.0404	0.83749\\
3.0663	0.8981\\
3.0857	0.94144\\
3.1052	0.98241\\
3.1214	1.0144\\
3.1376	1.0442\\
3.1505	1.0663\\
3.1635	1.0867\\
3.1764	1.1052\\
3.1894	1.1219\\
3.1991	1.1332\\
3.2088	1.1432\\
3.2185	1.1522\\
3.2282	1.1599\\
3.2379	1.1664\\
3.2476	1.1716\\
3.2574	1.1757\\
3.2671	1.1784\\
3.2768	1.18\\
3.2865	1.1802\\
3.2962	1.1793\\
3.3059	1.1771\\
3.3156	1.1737\\
3.3253	1.169\\
3.3351	1.1633\\
3.3448	1.1563\\
3.3545	1.1483\\
3.3642	1.1392\\
3.3739	1.1291\\
3.3869	1.114\\
3.3998	1.0972\\
3.4128	1.0789\\
3.4257	1.0592\\
3.4419	1.0327\\
3.4581	1.0044\\
3.4775	0.96852\\
3.5002	0.92445\\
3.5293	0.86546\\
3.6232	0.67339\\
3.6491	0.62348\\
3.6718	0.58182\\
3.6912	0.54782\\
3.7106	0.51553\\
3.7301	0.48503\\
3.7495	0.45637\\
3.7689	0.42957\\
3.7851	0.40864\\
3.8013	0.38895\\
3.8175	0.37048\\
3.8337	0.35319\\
3.8499	0.33701\\
3.8693	0.319\\
3.8887	0.30241\\
3.9082	0.28714\\
3.9276	0.27306\\
3.947	0.26006\\
3.9697	0.24611\\
3.9923	0.2333\\
4.015	0.22148\\
4.0409	0.209\\
4.07	0.19604\\
4.1024	0.1827\\
4.138	0.16904\\
4.1801	0.15393\\
4.2255	0.13861\\
4.2773	0.12211\\
4.3291	0.10658\\
4.3776	0.092923\\
4.4262	0.080215\\
4.4715	0.06927\\
4.5169	0.059261\\
4.5622	0.050217\\
4.6075	0.042145\\
4.6529	0.035029\\
4.6982	0.028832\\
4.7468	0.023152\\
4.7953	0.018383\\
4.8471	0.014197\\
4.9022	0.010639\\
4.9637	0.0075763\\
5.0349	0.0050001\\
5.1159	0.0030283\\
5.2162	0.0015575\\
5.3101	0.00080085\\
};
\addlegendentry{\ApproxBLNFitLegend}


\addplot[area legend, draw=none, fill=mycolor1, fill opacity=0.15, forget plot]
table[row sep=crcr] {%
x	y\\
2.0755	0.00042104\\
2.0788	0.00043536\\
2.082	0.0004502\\
2.0853	0.00046556\\
2.0885	0.00048147\\
2.0917	0.00049795\\
2.095	0.00051501\\
2.0982	0.00053268\\
2.1014	0.00055098\\
2.1047	0.00056993\\
2.1079	0.00058956\\
2.1112	0.00060988\\
2.1144	0.00063094\\
2.1176	0.00065274\\
2.1209	0.00067532\\
2.1241	0.00069871\\
2.1274	0.00072294\\
2.1306	0.00074803\\
2.1338	0.00077401\\
2.1371	0.00080093\\
2.1403	0.0008288\\
2.1435	0.00085768\\
2.1468	0.00088758\\
2.15	0.00091855\\
2.1533	0.00095062\\
2.1565	0.00098384\\
2.1597	0.0010182\\
2.163	0.0010539\\
2.1662	0.0010908\\
2.1694	0.001129\\
2.1727	0.0011686\\
2.1759	0.0012095\\
2.1792	0.001252\\
2.1824	0.0012959\\
2.1856	0.0013414\\
2.1889	0.0013885\\
2.1921	0.0014373\\
2.1953	0.0014878\\
2.1986	0.0015401\\
2.2018	0.0015943\\
2.2051	0.0016503\\
2.2083	0.0017083\\
2.2115	0.0017684\\
2.2148	0.0018306\\
2.218	0.0018949\\
2.2212	0.0019615\\
2.2245	0.0020305\\
2.2277	0.0021018\\
2.231	0.0021757\\
2.2342	0.0022521\\
2.2374	0.0023311\\
2.2407	0.0024129\\
2.2439	0.0024975\\
2.2472	0.0025851\\
2.2504	0.0026757\\
2.2536	0.0027694\\
2.2569	0.0028663\\
2.2601	0.0029665\\
2.2633	0.0030702\\
2.2666	0.0031774\\
2.2698	0.0032883\\
2.2731	0.0034029\\
2.2763	0.0035214\\
2.2795	0.0036439\\
2.2828	0.0037706\\
2.286	0.0039015\\
2.2892	0.0040369\\
2.2925	0.0041767\\
2.2957	0.0043213\\
2.299	0.0044706\\
2.3022	0.004625\\
2.3054	0.0047844\\
2.3087	0.0049491\\
2.3119	0.0051193\\
2.3151	0.005295\\
2.3184	0.0054765\\
2.3216	0.005664\\
2.3249	0.0058575\\
2.3281	0.0060574\\
2.3313	0.0062637\\
2.3346	0.0064767\\
2.3378	0.0066965\\
2.341	0.0069234\\
2.3443	0.0071576\\
2.3475	0.0073992\\
2.3508	0.0076485\\
2.354	0.0079057\\
2.3572	0.0081711\\
2.3605	0.0084447\\
2.3637	0.008727\\
2.367	0.0090181\\
2.3702	0.0093183\\
2.3734	0.0096278\\
2.3767	0.0099468\\
2.3799	0.010276\\
2.3831	0.010615\\
2.3864	0.010964\\
2.3896	0.011324\\
2.3929	0.011695\\
2.3961	0.012077\\
2.3993	0.012471\\
2.4026	0.012876\\
2.4058	0.013294\\
2.409	0.013724\\
2.4123	0.014167\\
2.4155	0.014622\\
2.4188	0.015092\\
2.422	0.015575\\
2.4252	0.016072\\
2.4285	0.016583\\
2.4317	0.01711\\
2.4349	0.017651\\
2.4382	0.018208\\
2.4414	0.018781\\
2.4447	0.01937\\
2.4479	0.019976\\
2.4511	0.020599\\
2.4544	0.02124\\
2.4576	0.021898\\
2.4608	0.022574\\
2.4641	0.02327\\
2.4673	0.023984\\
2.4706	0.024718\\
2.4738	0.025472\\
2.477	0.026247\\
2.4803	0.027042\\
2.4835	0.027859\\
2.4868	0.028698\\
2.49	0.029559\\
2.4932	0.030443\\
2.4965	0.03135\\
2.4997	0.032281\\
2.5029	0.033236\\
2.5062	0.034216\\
2.5094	0.035221\\
2.5127	0.036252\\
2.5159	0.03731\\
2.5191	0.038394\\
2.5224	0.039505\\
2.5256	0.040645\\
2.5288	0.041813\\
2.5321	0.04301\\
2.5353	0.044237\\
2.5386	0.045494\\
2.5418	0.046781\\
2.545	0.0481\\
2.5483	0.049451\\
2.5515	0.050834\\
2.5547	0.05225\\
2.558	0.0537\\
2.5612	0.055184\\
2.5645	0.056703\\
2.5677	0.058257\\
2.5709	0.059848\\
2.5742	0.061475\\
2.5774	0.063139\\
2.5806	0.064841\\
2.5839	0.066581\\
2.5871	0.06836\\
2.5904	0.07018\\
2.5936	0.072039\\
2.5968	0.073939\\
2.6001	0.075881\\
2.6033	0.077866\\
2.6066	0.079893\\
2.6098	0.081963\\
2.613	0.084077\\
2.6163	0.086236\\
2.6195	0.088441\\
2.6227	0.090691\\
2.626	0.092988\\
2.6292	0.095332\\
2.6325	0.097724\\
2.6357	0.10016\\
2.6389	0.10265\\
2.6422	0.10519\\
2.6454	0.10778\\
2.6486	0.11042\\
2.6519	0.11311\\
2.6551	0.11586\\
2.6584	0.11865\\
2.6616	0.1215\\
2.6648	0.1244\\
2.6681	0.12736\\
2.6713	0.13037\\
2.6745	0.13344\\
2.6778	0.13656\\
2.681	0.13974\\
2.6843	0.14298\\
2.6875	0.14628\\
2.6907	0.14963\\
2.694	0.15304\\
2.6972	0.15651\\
2.7004	0.16004\\
2.7037	0.16363\\
2.7069	0.16728\\
2.7102	0.17099\\
2.7134	0.17476\\
2.7166	0.1786\\
2.7199	0.18249\\
2.7231	0.18645\\
2.7263	0.19047\\
2.7296	0.19456\\
2.7328	0.19871\\
2.7361	0.20292\\
2.7393	0.20719\\
2.7425	0.21153\\
2.7458	0.21594\\
2.749	0.22041\\
2.7523	0.22494\\
2.7555	0.22954\\
2.7587	0.23421\\
2.762	0.23894\\
2.7652	0.24374\\
2.7684	0.2486\\
2.7717	0.25353\\
2.7749	0.25853\\
2.7782	0.26359\\
2.7814	0.26872\\
2.7846	0.27391\\
2.7879	0.27917\\
2.7911	0.2845\\
2.7943	0.28989\\
2.7976	0.29535\\
2.8008	0.30087\\
2.8041	0.30646\\
2.8073	0.31211\\
2.8105	0.31783\\
2.8138	0.32362\\
2.817	0.32946\\
2.8202	0.33538\\
2.8235	0.34135\\
2.8267	0.34739\\
2.83	0.3535\\
2.8332	0.35966\\
2.8364	0.36589\\
2.8397	0.37218\\
2.8429	0.37854\\
2.8461	0.38495\\
2.8494	0.39142\\
2.8526	0.39795\\
2.8559	0.40454\\
2.8591	0.41119\\
2.8623	0.4179\\
2.8656	0.42466\\
2.8688	0.43148\\
2.8721	0.43835\\
2.8753	0.44528\\
2.8785	0.45226\\
2.8818	0.4593\\
2.885	0.46638\\
2.8882	0.47352\\
2.8915	0.4807\\
2.8947	0.48794\\
2.898	0.49522\\
2.9012	0.50255\\
2.9044	0.50992\\
2.9077	0.51733\\
2.9109	0.52479\\
2.9141	0.53229\\
2.9174	0.53983\\
2.9206	0.54741\\
2.9239	0.55502\\
2.9271	0.56267\\
2.9303	0.57036\\
2.9336	0.57808\\
2.9368	0.58583\\
2.94	0.59361\\
2.9433	0.60142\\
2.9465	0.60925\\
2.9498	0.61711\\
2.953	0.62499\\
2.9562	0.6329\\
2.9595	0.64082\\
2.9627	0.64877\\
2.9659	0.65673\\
2.9692	0.6647\\
2.9724	0.67269\\
2.9757	0.68069\\
2.9789	0.68869\\
2.9821	0.6967\\
2.9854	0.70472\\
2.9886	0.71274\\
2.9919	0.72077\\
2.9951	0.72879\\
2.9983	0.7368\\
3.0016	0.74482\\
3.0048	0.75282\\
3.008	0.76082\\
3.0113	0.7688\\
3.0145	0.77677\\
3.0178	0.78473\\
3.021	0.79266\\
3.0242	0.80058\\
3.0275	0.80847\\
3.0307	0.81634\\
3.0339	0.82418\\
3.0372	0.83199\\
3.0404	0.83977\\
3.0437	0.84751\\
3.0469	0.85522\\
3.0501	0.86289\\
3.0534	0.87052\\
3.0566	0.87811\\
3.0598	0.88565\\
3.0631	0.89315\\
3.0663	0.90059\\
3.0696	0.90798\\
3.0728	0.91532\\
3.076	0.9226\\
3.0793	0.92982\\
3.0825	0.93698\\
3.0857	0.94408\\
3.089	0.95111\\
3.0922	0.95807\\
3.0955	0.96496\\
3.0987	0.97178\\
3.1019	0.97853\\
3.1052	0.9852\\
3.1084	0.99178\\
3.1117	0.99829\\
3.1149	1.0047\\
3.1181	1.0111\\
3.1214	1.0173\\
3.1246	1.0235\\
3.1278	1.0295\\
3.1311	1.0355\\
3.1343	1.0414\\
3.1376	1.0472\\
3.1408	1.0529\\
3.144	1.0585\\
3.1473	1.0639\\
3.1505	1.0693\\
3.1537	1.0746\\
3.157	1.0797\\
3.1602	1.0848\\
3.1635	1.0897\\
3.1667	1.0946\\
3.1699	1.0993\\
3.1732	1.1039\\
3.1764	1.1083\\
3.1796	1.1127\\
3.1829	1.1169\\
3.1861	1.1211\\
3.1894	1.125\\
3.1926	1.1289\\
3.1958	1.1327\\
3.1991	1.1363\\
3.2023	1.1398\\
3.2055	1.1431\\
3.2088	1.1463\\
3.212	1.1494\\
3.2153	1.1524\\
3.2185	1.1552\\
3.2217	1.1579\\
3.225	1.1605\\
3.2282	1.1629\\
3.2315	1.1652\\
3.2347	1.1674\\
3.2379	1.1694\\
3.2412	1.1712\\
3.2444	1.173\\
3.2476	1.1746\\
3.2509	1.1761\\
3.2541	1.1774\\
3.2574	1.1786\\
3.2606	1.1796\\
3.2638	1.1805\\
3.2671	1.1813\\
3.2703	1.1819\\
3.2735	1.1824\\
3.2768	1.1827\\
3.28	1.1829\\
3.2833	1.183\\
3.2865	1.1829\\
3.2897	1.1827\\
3.293	1.1824\\
3.2962	1.1819\\
3.2994	1.1813\\
3.3027	1.1805\\
3.3059	1.1796\\
3.3092	1.1786\\
3.3124	1.1774\\
3.3156	1.1761\\
3.3189	1.1747\\
3.3221	1.1732\\
3.3253	1.1715\\
3.3286	1.1696\\
3.3318	1.1677\\
3.3351	1.1656\\
3.3383	1.1634\\
3.3415	1.1611\\
3.3448	1.1586\\
3.348	1.1561\\
3.3513	1.1534\\
3.3545	1.1506\\
3.3577	1.1476\\
3.361	1.1446\\
3.3642	1.1414\\
3.3674	1.1381\\
3.3707	1.1347\\
3.3739	1.1312\\
3.3772	1.1276\\
3.3804	1.1239\\
3.3836	1.12\\
3.3869	1.1161\\
3.3901	1.1121\\
3.3933	1.1079\\
3.3966	1.1037\\
3.3998	1.0994\\
3.4031	1.0949\\
3.4063	1.0904\\
3.4095	1.0858\\
3.4128	1.0811\\
3.416	1.0763\\
3.4192	1.0714\\
3.4225	1.0664\\
3.4257	1.0614\\
3.429	1.0563\\
3.4322	1.051\\
3.4354	1.0458\\
3.4387	1.0404\\
3.4419	1.035\\
3.4451	1.0295\\
3.4484	1.0239\\
3.4516	1.0183\\
3.4549	1.0126\\
3.4581	1.0068\\
3.4613	1.001\\
3.4646	0.99508\\
3.4678	0.98914\\
3.471	0.98314\\
3.4743	0.9771\\
3.4775	0.971\\
3.4808	0.96485\\
3.484	0.95866\\
3.4872	0.95242\\
3.4905	0.94615\\
3.4937	0.93983\\
3.497	0.93347\\
3.5002	0.92708\\
3.5034	0.92065\\
3.5067	0.91419\\
3.5099	0.90771\\
3.5131	0.90119\\
3.5164	0.89465\\
3.5196	0.88808\\
3.5229	0.8815\\
3.5261	0.87489\\
3.5293	0.86826\\
3.5326	0.86162\\
3.5358	0.85496\\
3.539	0.8483\\
3.5423	0.84162\\
3.5455	0.83493\\
3.5488	0.82824\\
3.552	0.82154\\
3.5552	0.81484\\
3.5585	0.80813\\
3.5617	0.80143\\
3.5649	0.79473\\
3.5682	0.78804\\
3.5714	0.78134\\
3.5747	0.77466\\
3.5779	0.76799\\
3.5811	0.76132\\
3.5844	0.75467\\
3.5876	0.74803\\
3.5908	0.7414\\
3.5941	0.7348\\
3.5973	0.72821\\
3.6006	0.72164\\
3.6038	0.71509\\
3.607	0.70856\\
3.6103	0.70205\\
3.6135	0.69558\\
3.6168	0.68912\\
3.62	0.6827\\
3.6232	0.6763\\
3.6265	0.66993\\
3.6297	0.66359\\
3.6329	0.65729\\
3.6362	0.65102\\
3.6394	0.64478\\
3.6427	0.63857\\
3.6459	0.63241\\
3.6491	0.62628\\
3.6524	0.62019\\
3.6556	0.61413\\
3.6588	0.60812\\
3.6621	0.60215\\
3.6653	0.59622\\
3.6686	0.59033\\
3.6718	0.58448\\
3.675	0.57868\\
3.6783	0.57292\\
3.6815	0.5672\\
3.6847	0.56153\\
3.688	0.55591\\
3.6912	0.55033\\
3.6945	0.54481\\
3.6977	0.53932\\
3.7009	0.53389\\
3.7042	0.5285\\
3.7074	0.52317\\
3.7106	0.51788\\
3.7139	0.51264\\
3.7171	0.50745\\
3.7204	0.50232\\
3.7236	0.49723\\
3.7268	0.49219\\
3.7301	0.4872\\
3.7333	0.48227\\
3.7366	0.47738\\
3.7398	0.47255\\
3.743	0.46777\\
3.7463	0.46304\\
3.7495	0.45836\\
3.7527	0.45373\\
3.756	0.44915\\
3.7592	0.44463\\
3.7625	0.44015\\
3.7657	0.43573\\
3.7689	0.43136\\
3.7722	0.42704\\
3.7754	0.42277\\
3.7786	0.41855\\
3.7819	0.41439\\
3.7851	0.41027\\
3.7884	0.4062\\
3.7916	0.40219\\
3.7948	0.39822\\
3.7981	0.3943\\
3.8013	0.39044\\
3.8045	0.38662\\
3.8078	0.38285\\
3.811	0.37913\\
3.8143	0.37545\\
3.8175	0.37183\\
3.8207	0.36825\\
3.824	0.36472\\
3.8272	0.36124\\
3.8304	0.3578\\
3.8337	0.35441\\
3.8369	0.35106\\
3.8402	0.34776\\
3.8434	0.34451\\
3.8466	0.34129\\
3.8499	0.33813\\
3.8531	0.335\\
3.8564	0.33192\\
3.8596	0.32888\\
3.8628	0.32589\\
3.8661	0.32293\\
3.8693	0.32002\\
3.8725	0.31715\\
3.8758	0.31431\\
3.879	0.31152\\
3.8823	0.30876\\
3.8855	0.30605\\
3.8887	0.30337\\
3.892	0.30073\\
3.8952	0.29813\\
3.8984	0.29556\\
3.9017	0.29303\\
3.9049	0.29053\\
3.9082	0.28807\\
3.9114	0.28564\\
3.9146	0.28324\\
3.9179	0.28088\\
3.9211	0.27855\\
3.9243	0.27625\\
3.9276	0.27399\\
3.9308	0.27175\\
3.9341	0.26954\\
3.9373	0.26737\\
3.9405	0.26522\\
3.9438	0.2631\\
3.947	0.26101\\
3.9502	0.25894\\
3.9535	0.2569\\
3.9567	0.25489\\
3.96	0.2529\\
3.9632	0.25094\\
3.9664	0.249\\
3.9697	0.24709\\
3.9729	0.2452\\
3.9762	0.24333\\
3.9794	0.24149\\
3.9826	0.23966\\
3.9859	0.23786\\
3.9891	0.23608\\
3.9923	0.23432\\
3.9956	0.23258\\
3.9988	0.23086\\
4.0021	0.22916\\
4.0053	0.22747\\
4.0085	0.22581\\
4.0118	0.22416\\
4.015	0.22253\\
4.0182	0.22092\\
4.0215	0.21932\\
4.0247	0.21774\\
4.028	0.21618\\
4.0312	0.21463\\
4.0344	0.2131\\
4.0377	0.21158\\
4.0409	0.21007\\
4.0441	0.20858\\
4.0474	0.2071\\
4.0506	0.20564\\
4.0539	0.20419\\
4.0571	0.20275\\
4.0603	0.20132\\
4.0636	0.19991\\
4.0668	0.1985\\
4.07	0.19711\\
4.0733	0.19573\\
4.0765	0.19436\\
4.0798	0.193\\
4.083	0.19165\\
4.0862	0.19031\\
4.0895	0.18898\\
4.0927	0.18766\\
4.0959	0.18635\\
4.0992	0.18505\\
4.1024	0.18375\\
4.1057	0.18247\\
4.1089	0.18119\\
4.1121	0.17992\\
4.1154	0.17866\\
4.1186	0.17741\\
4.1219	0.17616\\
4.1251	0.17492\\
4.1283	0.17369\\
4.1316	0.17247\\
4.1348	0.17125\\
4.138	0.17004\\
4.1413	0.16883\\
4.1445	0.16764\\
4.1478	0.16644\\
4.151	0.16526\\
4.1542	0.16408\\
4.1575	0.1629\\
4.1607	0.16174\\
4.1639	0.16057\\
4.1672	0.15941\\
4.1704	0.15826\\
4.1737	0.15712\\
4.1769	0.15597\\
4.1801	0.15484\\
4.1834	0.1537\\
4.1866	0.15258\\
4.1898	0.15146\\
4.1931	0.15034\\
4.1963	0.14922\\
4.1996	0.14812\\
4.2028	0.14701\\
4.206	0.14591\\
4.2093	0.14482\\
4.2125	0.14373\\
4.2157	0.14264\\
4.219	0.14156\\
4.2222	0.14048\\
4.2255	0.1394\\
4.2287	0.13833\\
4.2319	0.13727\\
4.2352	0.1362\\
4.2384	0.13514\\
4.2417	0.13409\\
4.2449	0.13304\\
4.2481	0.13199\\
4.2514	0.13095\\
4.2546	0.12991\\
4.2578	0.12888\\
4.2611	0.12784\\
4.2643	0.12682\\
4.2676	0.12579\\
4.2708	0.12477\\
4.274	0.12376\\
4.2773	0.12275\\
4.2805	0.12174\\
4.2837	0.12073\\
4.287	0.11973\\
4.2902	0.11873\\
4.2935	0.11774\\
4.2967	0.11675\\
4.2999	0.11576\\
4.3032	0.11478\\
4.3064	0.1138\\
4.3096	0.11283\\
4.3129	0.11186\\
4.3161	0.11089\\
4.3194	0.10993\\
4.3226	0.10897\\
4.3258	0.10802\\
4.3291	0.10707\\
4.3323	0.10612\\
4.3355	0.10518\\
4.3388	0.10424\\
4.342	0.1033\\
4.3453	0.10237\\
4.3485	0.10145\\
4.3517	0.10052\\
4.355	0.099604\\
4.3582	0.098689\\
4.3615	0.097779\\
4.3647	0.096873\\
4.3679	0.095971\\
4.3712	0.095073\\
4.3744	0.094179\\
4.3776	0.09329\\
4.3809	0.092405\\
4.3841	0.091525\\
4.3874	0.090648\\
4.3906	0.089776\\
4.3938	0.088909\\
4.3971	0.088046\\
4.4003	0.087187\\
4.4035	0.086333\\
4.4068	0.085483\\
4.41	0.084638\\
4.4133	0.083798\\
4.4165	0.082962\\
4.4197	0.08213\\
4.423	0.081303\\
4.4262	0.080481\\
4.4294	0.079664\\
4.4327	0.078851\\
4.4359	0.078043\\
4.4392	0.07724\\
4.4424	0.076441\\
4.4456	0.075647\\
4.4489	0.074858\\
4.4521	0.074074\\
4.4553	0.073294\\
4.4586	0.07252\\
4.4618	0.07175\\
4.4651	0.070986\\
4.4683	0.070226\\
4.4715	0.069471\\
4.4748	0.068721\\
4.478	0.067976\\
4.4813	0.067236\\
4.4845	0.0665\\
4.4877	0.06577\\
4.491	0.065045\\
4.4942	0.064325\\
4.4974	0.06361\\
4.5007	0.0629\\
4.5039	0.062195\\
4.5072	0.061495\\
4.5104	0.0608\\
4.5136	0.06011\\
4.5169	0.059426\\
4.5201	0.058746\\
4.5233	0.058071\\
4.5266	0.057402\\
4.5298	0.056737\\
4.5331	0.056078\\
4.5363	0.055424\\
4.5395	0.054774\\
4.5428	0.05413\\
4.546	0.053491\\
4.5492	0.052857\\
4.5525	0.052228\\
4.5557	0.051604\\
4.559	0.050985\\
4.5622	0.050372\\
4.5654	0.049763\\
4.5687	0.049159\\
4.5719	0.04856\\
4.5751	0.047966\\
4.5784	0.047378\\
4.5816	0.046794\\
4.5849	0.046215\\
4.5881	0.045641\\
4.5913	0.045073\\
4.5946	0.044509\\
4.5978	0.04395\\
4.6011	0.043396\\
4.6043	0.042847\\
4.6075	0.042303\\
4.6108	0.041763\\
4.614	0.041229\\
4.6172	0.040699\\
4.6205	0.040175\\
4.6237	0.039655\\
4.627	0.03914\\
4.6302	0.03863\\
4.6334	0.038124\\
4.6367	0.037624\\
4.6399	0.037128\\
4.6431	0.036637\\
4.6464	0.03615\\
4.6496	0.035669\\
4.6529	0.035192\\
4.6561	0.03472\\
4.6593	0.034252\\
4.6626	0.033789\\
4.6658	0.033331\\
4.669	0.032877\\
4.6723	0.032428\\
4.6755	0.031983\\
4.6788	0.031543\\
4.682	0.031108\\
4.6852	0.030677\\
4.6885	0.03025\\
4.6917	0.029828\\
4.6949	0.02941\\
4.6982	0.028997\\
4.7014	0.028588\\
4.7047	0.028184\\
4.7079	0.027784\\
4.7111	0.027388\\
4.7144	0.026997\\
4.7176	0.026609\\
4.7209	0.026227\\
4.7241	0.025848\\
4.7273	0.025473\\
4.7306	0.025103\\
4.7338	0.024737\\
4.737	0.024375\\
4.7403	0.024017\\
4.7435	0.023663\\
4.7468	0.023314\\
4.75	0.022968\\
4.7532	0.022626\\
4.7565	0.022289\\
4.7597	0.021955\\
4.7629	0.021625\\
4.7662	0.021299\\
4.7694	0.020977\\
4.7727	0.020659\\
4.7759	0.020345\\
4.7791	0.020034\\
4.7824	0.019727\\
4.7856	0.019424\\
4.7888	0.019125\\
4.7921	0.018829\\
4.7953	0.018537\\
4.7986	0.018249\\
4.8018	0.017964\\
4.805	0.017683\\
4.8083	0.017405\\
4.8115	0.017131\\
4.8147	0.01686\\
4.818	0.016593\\
4.8212	0.016329\\
4.8245	0.016068\\
4.8277	0.015811\\
4.8309	0.015558\\
4.8342	0.015307\\
4.8374	0.01506\\
4.8406	0.014816\\
4.8439	0.014576\\
4.8471	0.014338\\
4.8504	0.014104\\
4.8536	0.013873\\
4.8568	0.013645\\
4.8601	0.01342\\
4.8633	0.013198\\
4.8666	0.012979\\
4.8698	0.012763\\
4.873	0.01255\\
4.8763	0.012341\\
4.8795	0.012133\\
4.8827	0.011929\\
4.886	0.011728\\
4.8892	0.01153\\
4.8925	0.011334\\
4.8957	0.011141\\
4.8989	0.010951\\
4.9022	0.010763\\
4.9054	0.010579\\
4.9086	0.010396\\
4.9119	0.010217\\
4.9151	0.01004\\
4.9184	0.0098658\\
4.9216	0.009694\\
4.9248	0.0095248\\
4.9281	0.009358\\
4.9313	0.0091937\\
4.9345	0.0090319\\
4.9378	0.0088725\\
4.941	0.0087154\\
4.9443	0.0085607\\
4.9475	0.0084083\\
4.9507	0.0082583\\
4.954	0.0081105\\
4.9572	0.0079649\\
4.9604	0.0078216\\
4.9637	0.0076805\\
4.9669	0.0075415\\
4.9702	0.0074047\\
4.9734	0.0072701\\
4.9766	0.0071375\\
4.9799	0.007007\\
4.9831	0.0068785\\
4.9864	0.0067521\\
4.9896	0.0066276\\
4.9928	0.0065051\\
4.9961	0.0063846\\
4.9993	0.006266\\
5.0025	0.0061493\\
5.0058	0.0060345\\
5.009	0.0059215\\
5.0123	0.0058103\\
5.0155	0.005701\\
5.0187	0.0055934\\
5.022	0.0054876\\
5.0252	0.0053836\\
5.0284	0.0052812\\
5.0317	0.0051805\\
5.0349	0.0050815\\
5.0382	0.0049842\\
5.0414	0.0048884\\
5.0446	0.0047943\\
5.0479	0.0047018\\
5.0511	0.0046108\\
5.0543	0.0045213\\
5.0576	0.0044334\\
5.0608	0.0043469\\
5.0641	0.004262\\
5.0673	0.0041784\\
5.0705	0.0040964\\
5.0738	0.0040157\\
5.077	0.0039364\\
5.0802	0.0038585\\
5.0835	0.003782\\
5.0867	0.0037068\\
5.09	0.0036329\\
5.0932	0.0035603\\
5.0964	0.003489\\
5.0997	0.0034189\\
5.1029	0.0033501\\
5.1062	0.0032825\\
5.1094	0.0032161\\
5.1126	0.0031509\\
5.1159	0.0030869\\
5.1191	0.003024\\
5.1223	0.0029623\\
5.1256	0.0029016\\
5.1288	0.0028421\\
5.1321	0.0027837\\
5.1353	0.0027263\\
5.1385	0.00267\\
5.1418	0.0026147\\
5.145	0.0025604\\
5.1482	0.0025071\\
5.1515	0.0024549\\
5.1547	0.0024036\\
5.158	0.0023532\\
5.1612	0.0023038\\
5.1644	0.0022553\\
5.1677	0.0022077\\
5.1709	0.0021611\\
5.1741	0.0021153\\
5.1774	0.0020703\\
5.1806	0.0020263\\
5.1839	0.001983\\
5.1871	0.0019406\\
5.1903	0.001899\\
5.1936	0.0018582\\
5.1968	0.0018182\\
5.2	0.001779\\
5.2033	0.0017405\\
5.2065	0.0017028\\
5.2098	0.0016658\\
5.213	0.0016296\\
5.2162	0.001594\\
5.2195	0.0015591\\
5.2227	0.001525\\
5.226	0.0014915\\
5.2292	0.0014587\\
5.2324	0.0014265\\
5.2357	0.0013949\\
5.2389	0.001364\\
5.2421	0.0013338\\
5.2454	0.0013041\\
5.2486	0.001275\\
5.2519	0.0012465\\
5.2551	0.0012186\\
5.2583	0.0011912\\
5.2616	0.0011644\\
5.2648	0.0011382\\
5.268	0.0011125\\
5.2713	0.0010873\\
5.2745	0.0010626\\
5.2778	0.0010385\\
5.281	0.0010148\\
5.2842	0.00099166\\
5.2875	0.00096897\\
5.2907	0.00094676\\
5.2939	0.00092501\\
5.2972	0.00090372\\
5.3004	0.00088287\\
5.3037	0.00086246\\
5.3069	0.00084248\\
5.3101	0.00082293\\
5.3101	0.00077877\\
5.3069	0.00079751\\
5.3037	0.00081667\\
5.3004	0.00083624\\
5.2972	0.00085623\\
5.2939	0.00087666\\
5.2907	0.00089754\\
5.2875	0.00091886\\
5.2842	0.00094065\\
5.281	0.0009629\\
5.2778	0.00098563\\
5.2745	0.0010088\\
5.2713	0.0010326\\
5.268	0.0010568\\
5.2648	0.0010815\\
5.2616	0.0011068\\
5.2583	0.0011325\\
5.2551	0.0011589\\
5.2519	0.0011858\\
5.2486	0.0012132\\
5.2454	0.0012412\\
5.2421	0.0012698\\
5.2389	0.001299\\
5.2357	0.0013288\\
5.2324	0.0013593\\
5.2292	0.0013903\\
5.226	0.001422\\
5.2227	0.0014543\\
5.2195	0.0014873\\
5.2162	0.001521\\
5.213	0.0015553\\
5.2098	0.0015904\\
5.2065	0.0016261\\
5.2033	0.0016626\\
5.2	0.0016998\\
5.1968	0.0017378\\
5.1936	0.0017765\\
5.1903	0.001816\\
5.1871	0.0018562\\
5.1839	0.0018973\\
5.1806	0.0019392\\
5.1774	0.0019819\\
5.1741	0.0020254\\
5.1709	0.0020698\\
5.1677	0.0021151\\
5.1644	0.0021613\\
5.1612	0.0022083\\
5.158	0.0022563\\
5.1547	0.0023052\\
5.1515	0.002355\\
5.1482	0.0024058\\
5.145	0.0024575\\
5.1418	0.0025103\\
5.1385	0.002564\\
5.1353	0.0026188\\
5.1321	0.0026746\\
5.1288	0.0027314\\
5.1256	0.0027893\\
5.1223	0.0028483\\
5.1191	0.0029085\\
5.1159	0.0029697\\
5.1126	0.0030321\\
5.1094	0.0030956\\
5.1062	0.0031603\\
5.1029	0.0032262\\
5.0997	0.0032933\\
5.0964	0.0033616\\
5.0932	0.0034311\\
5.09	0.003502\\
5.0867	0.0035741\\
5.0835	0.0036475\\
5.0802	0.0037223\\
5.077	0.0037983\\
5.0738	0.0038758\\
5.0705	0.0039546\\
5.0673	0.0040348\\
5.0641	0.0041165\\
5.0608	0.0041996\\
5.0576	0.0042841\\
5.0543	0.0043701\\
5.0511	0.0044577\\
5.0479	0.0045468\\
5.0446	0.0046374\\
5.0414	0.0047295\\
5.0382	0.0048233\\
5.0349	0.0049187\\
5.0317	0.0050157\\
5.0284	0.0051144\\
5.0252	0.0052148\\
5.022	0.0053168\\
5.0187	0.0054206\\
5.0155	0.0055261\\
5.0123	0.0056334\\
5.009	0.0057426\\
5.0058	0.0058535\\
5.0025	0.0059662\\
4.9993	0.0060809\\
4.9961	0.0061974\\
4.9928	0.0063158\\
4.9896	0.0064362\\
4.9864	0.0065586\\
4.9831	0.0066829\\
4.9799	0.0068093\\
4.9766	0.0069377\\
4.9734	0.0070681\\
4.9702	0.0072007\\
4.9669	0.0073354\\
4.9637	0.0074722\\
4.9604	0.0076112\\
4.9572	0.0077524\\
4.954	0.0078958\\
4.9507	0.0080415\\
4.9475	0.0081894\\
4.9443	0.0083396\\
4.941	0.0084922\\
4.9378	0.0086471\\
4.9345	0.0088044\\
4.9313	0.0089642\\
4.9281	0.0091263\\
4.9248	0.009291\\
4.9216	0.0094581\\
4.9184	0.0096277\\
4.9151	0.0098\\
4.9119	0.0099747\\
4.9086	0.010152\\
4.9054	0.010332\\
4.9022	0.010515\\
4.8989	0.0107\\
4.8957	0.010888\\
4.8925	0.011079\\
4.8892	0.011273\\
4.886	0.011469\\
4.8827	0.011669\\
4.8795	0.011871\\
4.8763	0.012076\\
4.873	0.012284\\
4.8698	0.012495\\
4.8666	0.012709\\
4.8633	0.012925\\
4.8601	0.013145\\
4.8568	0.013368\\
4.8536	0.013595\\
4.8504	0.013824\\
4.8471	0.014056\\
4.8439	0.014292\\
4.8406	0.014531\\
4.8374	0.014773\\
4.8342	0.015018\\
4.8309	0.015267\\
4.8277	0.015519\\
4.8245	0.015774\\
4.8212	0.016033\\
4.818	0.016295\\
4.8147	0.016561\\
4.8115	0.01683\\
4.8083	0.017103\\
4.805	0.017379\\
4.8018	0.017659\\
4.7986	0.017942\\
4.7953	0.018229\\
4.7921	0.01852\\
4.7888	0.018814\\
4.7856	0.019112\\
4.7824	0.019414\\
4.7791	0.01972\\
4.7759	0.020029\\
4.7727	0.020342\\
4.7694	0.02066\\
4.7662	0.02098\\
4.7629	0.021305\\
4.7597	0.021634\\
4.7565	0.021967\\
4.7532	0.022304\\
4.75	0.022645\\
4.7468	0.02299\\
4.7435	0.023339\\
4.7403	0.023692\\
4.737	0.024049\\
4.7338	0.02441\\
4.7306	0.024776\\
4.7273	0.025146\\
4.7241	0.02552\\
4.7209	0.025898\\
4.7176	0.02628\\
4.7144	0.026667\\
4.7111	0.027058\\
4.7079	0.027454\\
4.7047	0.027854\\
4.7014	0.028258\\
4.6982	0.028667\\
4.6949	0.02908\\
4.6917	0.029498\\
4.6885	0.02992\\
4.6852	0.030346\\
4.682	0.030778\\
4.6788	0.031213\\
4.6755	0.031654\\
4.6723	0.032099\\
4.669	0.032548\\
4.6658	0.033002\\
4.6626	0.033461\\
4.6593	0.033925\\
4.6561	0.034393\\
4.6529	0.034866\\
4.6496	0.035343\\
4.6464	0.035825\\
4.6431	0.036312\\
4.6399	0.036804\\
4.6367	0.037301\\
4.6334	0.037802\\
4.6302	0.038308\\
4.627	0.038819\\
4.6237	0.039335\\
4.6205	0.039856\\
4.6172	0.040381\\
4.614	0.040911\\
4.6108	0.041447\\
4.6075	0.041987\\
4.6043	0.042532\\
4.6011	0.043081\\
4.5978	0.043636\\
4.5946	0.044196\\
4.5913	0.04476\\
4.5881	0.04533\\
4.5849	0.045904\\
4.5816	0.046483\\
4.5784	0.047067\\
4.5751	0.047656\\
4.5719	0.04825\\
4.5687	0.048849\\
4.5654	0.049453\\
4.5622	0.050062\\
4.559	0.050675\\
4.5557	0.051294\\
4.5525	0.051917\\
4.5492	0.052545\\
4.546	0.053178\\
4.5428	0.053816\\
4.5395	0.054459\\
4.5363	0.055107\\
4.5331	0.05576\\
4.5298	0.056417\\
4.5266	0.05708\\
4.5233	0.057747\\
4.5201	0.058419\\
4.5169	0.059095\\
4.5136	0.059777\\
4.5104	0.060463\\
4.5072	0.061154\\
4.5039	0.06185\\
4.5007	0.062551\\
4.4974	0.063256\\
4.4942	0.063967\\
4.491	0.064681\\
4.4877	0.065401\\
4.4845	0.066125\\
4.4813	0.066854\\
4.478	0.067588\\
4.4748	0.068326\\
4.4715	0.069069\\
4.4683	0.069816\\
4.4651	0.070569\\
4.4618	0.071325\\
4.4586	0.072087\\
4.4553	0.072852\\
4.4521	0.073623\\
4.4489	0.074398\\
4.4456	0.075177\\
4.4424	0.075961\\
4.4392	0.07675\\
4.4359	0.077543\\
4.4327	0.07834\\
4.4294	0.079142\\
4.4262	0.079948\\
4.423	0.080758\\
4.4197	0.081573\\
4.4165	0.082393\\
4.4133	0.083216\\
4.41	0.084044\\
4.4068	0.084877\\
4.4035	0.085713\\
4.4003	0.086554\\
4.3971	0.087399\\
4.3938	0.088248\\
4.3906	0.089102\\
4.3874	0.089959\\
4.3841	0.090821\\
4.3809	0.091687\\
4.3776	0.092557\\
4.3744	0.093431\\
4.3712	0.094309\\
4.3679	0.095191\\
4.3647	0.096077\\
4.3615	0.096968\\
4.3582	0.097862\\
4.355	0.09876\\
4.3517	0.099662\\
4.3485	0.10057\\
4.3453	0.10148\\
4.342	0.10239\\
4.3388	0.10331\\
4.3355	0.10423\\
4.3323	0.10516\\
4.3291	0.10609\\
4.3258	0.10702\\
4.3226	0.10796\\
4.3194	0.1089\\
4.3161	0.10984\\
4.3129	0.11079\\
4.3096	0.11174\\
4.3064	0.1127\\
4.3032	0.11366\\
4.2999	0.11462\\
4.2967	0.11559\\
4.2935	0.11656\\
4.2902	0.11753\\
4.287	0.11851\\
4.2837	0.11949\\
4.2805	0.12048\\
4.2773	0.12147\\
4.274	0.12246\\
4.2708	0.12346\\
4.2676	0.12446\\
4.2643	0.12547\\
4.2611	0.12647\\
4.2578	0.12749\\
4.2546	0.1285\\
4.2514	0.12952\\
4.2481	0.13055\\
4.2449	0.13157\\
4.2417	0.13261\\
4.2384	0.13364\\
4.2352	0.13468\\
4.2319	0.13573\\
4.2287	0.13677\\
4.2255	0.13783\\
4.2222	0.13888\\
4.219	0.13994\\
4.2157	0.14101\\
4.2125	0.14208\\
4.2093	0.14315\\
4.206	0.14423\\
4.2028	0.14531\\
4.1996	0.1464\\
4.1963	0.14749\\
4.1931	0.14858\\
4.1898	0.14968\\
4.1866	0.15079\\
4.1834	0.1519\\
4.1801	0.15302\\
4.1769	0.15414\\
4.1737	0.15526\\
4.1704	0.1564\\
4.1672	0.15753\\
4.1639	0.15868\\
4.1607	0.15983\\
4.1575	0.16098\\
4.1542	0.16214\\
4.151	0.16331\\
4.1478	0.16448\\
4.1445	0.16566\\
4.1413	0.16684\\
4.138	0.16804\\
4.1348	0.16924\\
4.1316	0.17044\\
4.1283	0.17166\\
4.1251	0.17288\\
4.1219	0.17411\\
4.1186	0.17534\\
4.1154	0.17659\\
4.1121	0.17784\\
4.1089	0.1791\\
4.1057	0.18037\\
4.1024	0.18165\\
4.0992	0.18293\\
4.0959	0.18423\\
4.0927	0.18554\\
4.0895	0.18685\\
4.0862	0.18818\\
4.083	0.18951\\
4.0798	0.19086\\
4.0765	0.19222\\
4.0733	0.19358\\
4.07	0.19496\\
4.0668	0.19635\\
4.0636	0.19776\\
4.0603	0.19917\\
4.0571	0.2006\\
4.0539	0.20204\\
4.0506	0.20349\\
4.0474	0.20496\\
4.0441	0.20644\\
4.0409	0.20793\\
4.0377	0.20944\\
4.0344	0.21096\\
4.0312	0.2125\\
4.028	0.21406\\
4.0247	0.21563\\
4.0215	0.21721\\
4.0182	0.21881\\
4.015	0.22043\\
4.0118	0.22207\\
4.0085	0.22373\\
4.0053	0.2254\\
4.0021	0.22709\\
3.9988	0.2288\\
3.9956	0.23053\\
3.9923	0.23228\\
3.9891	0.23405\\
3.9859	0.23584\\
3.9826	0.23765\\
3.9794	0.23949\\
3.9762	0.24134\\
3.9729	0.24322\\
3.9697	0.24512\\
3.9664	0.24705\\
3.9632	0.249\\
3.96	0.25097\\
3.9567	0.25296\\
3.9535	0.25499\\
3.9502	0.25704\\
3.947	0.25911\\
3.9438	0.26121\\
3.9405	0.26334\\
3.9373	0.26549\\
3.9341	0.26768\\
3.9308	0.26989\\
3.9276	0.27213\\
3.9243	0.2744\\
3.9211	0.2767\\
3.9179	0.27903\\
3.9146	0.28139\\
3.9114	0.28379\\
3.9082	0.28621\\
3.9049	0.28867\\
3.9017	0.29116\\
3.8984	0.29368\\
3.8952	0.29624\\
3.892	0.29883\\
3.8887	0.30146\\
3.8855	0.30412\\
3.8823	0.30682\\
3.879	0.30956\\
3.8758	0.31233\\
3.8725	0.31514\\
3.8693	0.31799\\
3.8661	0.32087\\
3.8628	0.3238\\
3.8596	0.32676\\
3.8564	0.32977\\
3.8531	0.33281\\
3.8499	0.3359\\
3.8466	0.33903\\
3.8434	0.3422\\
3.8402	0.34541\\
3.8369	0.34867\\
3.8337	0.35197\\
3.8304	0.35531\\
3.8272	0.3587\\
3.824	0.36214\\
3.8207	0.36562\\
3.8175	0.36914\\
3.8143	0.37271\\
3.811	0.37633\\
3.8078	0.38\\
3.8045	0.38371\\
3.8013	0.38747\\
3.7981	0.39128\\
3.7948	0.39513\\
3.7916	0.39904\\
3.7884	0.403\\
3.7851	0.407\\
3.7819	0.41106\\
3.7786	0.41516\\
3.7754	0.41932\\
3.7722	0.42352\\
3.7689	0.42778\\
3.7657	0.43209\\
3.7625	0.43644\\
3.7592	0.44085\\
3.756	0.44531\\
3.7527	0.44983\\
3.7495	0.45439\\
3.7463	0.45901\\
3.743	0.46367\\
3.7398	0.46839\\
3.7366	0.47316\\
3.7333	0.47799\\
3.7301	0.48286\\
3.7268	0.48778\\
3.7236	0.49276\\
3.7204	0.49779\\
3.7171	0.50287\\
3.7139	0.50799\\
3.7106	0.51317\\
3.7074	0.5184\\
3.7042	0.52368\\
3.7009	0.52901\\
3.6977	0.53439\\
3.6945	0.53982\\
3.6912	0.5453\\
3.688	0.55082\\
3.6847	0.5564\\
3.6815	0.56201\\
3.6783	0.56768\\
3.675	0.57339\\
3.6718	0.57915\\
3.6686	0.58495\\
3.6653	0.5908\\
3.6621	0.59669\\
3.6588	0.60263\\
3.6556	0.6086\\
3.6524	0.61462\\
3.6491	0.62068\\
3.6459	0.62677\\
3.6427	0.63291\\
3.6394	0.63908\\
3.6362	0.64529\\
3.6329	0.65154\\
3.6297	0.65782\\
3.6265	0.66413\\
3.6232	0.67048\\
3.62	0.67686\\
3.6168	0.68327\\
3.6135	0.68971\\
3.6103	0.69617\\
3.607	0.70266\\
3.6038	0.70918\\
3.6006	0.71573\\
3.5973	0.72229\\
3.5941	0.72888\\
3.5908	0.73548\\
3.5876	0.74211\\
3.5844	0.74875\\
3.5811	0.75541\\
3.5779	0.76208\\
3.5747	0.76876\\
3.5714	0.77545\\
3.5682	0.78216\\
3.5649	0.78887\\
3.5617	0.79558\\
3.5585	0.8023\\
3.5552	0.80902\\
3.552	0.81574\\
3.5488	0.82247\\
3.5455	0.82918\\
3.5423	0.8359\\
3.539	0.8426\\
3.5358	0.8493\\
3.5326	0.85598\\
3.5293	0.86266\\
3.5261	0.86932\\
3.5229	0.87596\\
3.5196	0.88258\\
3.5164	0.88918\\
3.5131	0.89576\\
3.5099	0.90232\\
3.5067	0.90885\\
3.5034	0.91534\\
3.5002	0.92181\\
3.497	0.92825\\
3.4937	0.93465\\
3.4905	0.94101\\
3.4872	0.94733\\
3.484	0.95361\\
3.4808	0.95985\\
3.4775	0.96604\\
3.4743	0.97218\\
3.471	0.97827\\
3.4678	0.98431\\
3.4646	0.99029\\
3.4613	0.99622\\
3.4581	1.0021\\
3.4549	1.0079\\
3.4516	1.0136\\
3.4484	1.0193\\
3.4451	1.0249\\
3.4419	1.0305\\
3.4387	1.0359\\
3.4354	1.0413\\
3.4322	1.0466\\
3.429	1.0519\\
3.4257	1.057\\
3.4225	1.0621\\
3.4192	1.0671\\
3.416	1.072\\
3.4128	1.0768\\
3.4095	1.0815\\
3.4063	1.0861\\
3.4031	1.0907\\
3.3998	1.0951\\
3.3966	1.0994\\
3.3933	1.1037\\
3.3901	1.1078\\
3.3869	1.1118\\
3.3836	1.1158\\
3.3804	1.1196\\
3.3772	1.1233\\
3.3739	1.1269\\
3.3707	1.1304\\
3.3674	1.1338\\
3.3642	1.137\\
3.361	1.1402\\
3.3577	1.1432\\
3.3545	1.1461\\
3.3513	1.1489\\
3.348	1.1515\\
3.3448	1.1541\\
3.3415	1.1565\\
3.3383	1.1588\\
3.3351	1.1609\\
3.3318	1.1629\\
3.3286	1.1649\\
3.3253	1.1666\\
3.3221	1.1683\\
3.3189	1.1698\\
3.3156	1.1712\\
3.3124	1.1724\\
3.3092	1.1735\\
3.3059	1.1745\\
3.3027	1.1753\\
3.2994	1.176\\
3.2962	1.1766\\
3.293	1.1771\\
3.2897	1.1774\\
3.2865	1.1775\\
3.2833	1.1775\\
3.28	1.1774\\
3.2768	1.1772\\
3.2735	1.1768\\
3.2703	1.1762\\
3.2671	1.1756\\
3.2638	1.1748\\
3.2606	1.1738\\
3.2574	1.1727\\
3.2541	1.1715\\
3.2509	1.1702\\
3.2476	1.1687\\
3.2444	1.167\\
3.2412	1.1653\\
3.2379	1.1634\\
3.2347	1.1613\\
3.2315	1.1591\\
3.2282	1.1568\\
3.225	1.1544\\
3.2217	1.1518\\
3.2185	1.1491\\
3.2153	1.1462\\
3.212	1.1433\\
3.2088	1.1402\\
3.2055	1.1369\\
3.2023	1.1336\\
3.1991	1.1301\\
3.1958	1.1264\\
3.1926	1.1227\\
3.1894	1.1188\\
3.1861	1.1148\\
3.1829	1.1107\\
3.1796	1.1065\\
3.1764	1.1022\\
3.1732	1.0977\\
3.1699	1.0931\\
3.1667	1.0884\\
3.1635	1.0836\\
3.1602	1.0787\\
3.157	1.0736\\
3.1537	1.0685\\
3.1505	1.0633\\
3.1473	1.0579\\
3.144	1.0524\\
3.1408	1.0469\\
3.1376	1.0412\\
3.1343	1.0355\\
3.1311	1.0296\\
3.1278	1.0237\\
3.1246	1.0177\\
3.1214	1.0115\\
3.1181	1.0053\\
3.1149	0.99902\\
3.1117	0.99264\\
3.1084	0.98617\\
3.1052	0.97963\\
3.1019	0.973\\
3.0987	0.9663\\
3.0955	0.95953\\
3.0922	0.95269\\
3.089	0.94577\\
3.0857	0.93879\\
3.0825	0.93175\\
3.0793	0.92464\\
3.076	0.91747\\
3.0728	0.91024\\
3.0696	0.90295\\
3.0663	0.89561\\
3.0631	0.88822\\
3.0598	0.88078\\
3.0566	0.87329\\
3.0534	0.86576\\
3.0501	0.85818\\
3.0469	0.85056\\
3.0437	0.8429\\
3.0404	0.83521\\
3.0372	0.82748\\
3.0339	0.81972\\
3.0307	0.81193\\
3.0275	0.80411\\
3.0242	0.79626\\
3.021	0.7884\\
3.0178	0.78051\\
3.0145	0.7726\\
3.0113	0.76467\\
3.008	0.75673\\
3.0048	0.74878\\
3.0016	0.74082\\
2.9983	0.73285\\
2.9951	0.72487\\
2.9919	0.71688\\
2.9886	0.7089\\
2.9854	0.70091\\
2.9821	0.69293\\
2.9789	0.68495\\
2.9757	0.67697\\
2.9724	0.669\\
2.9692	0.66104\\
2.9659	0.6531\\
2.9627	0.64516\\
2.9595	0.63724\\
2.9562	0.62934\\
2.953	0.62145\\
2.9498	0.61359\\
2.9465	0.60574\\
2.9433	0.59792\\
2.94	0.59013\\
2.9368	0.58236\\
2.9336	0.57462\\
2.9303	0.56691\\
2.9271	0.55924\\
2.9239	0.55159\\
2.9206	0.54398\\
2.9174	0.53641\\
2.9141	0.52888\\
2.9109	0.52138\\
2.9077	0.51392\\
2.9044	0.50651\\
2.9012	0.49914\\
2.898	0.49181\\
2.8947	0.48453\\
2.8915	0.4773\\
2.8882	0.47011\\
2.885	0.46297\\
2.8818	0.45588\\
2.8785	0.44885\\
2.8753	0.44186\\
2.8721	0.43493\\
2.8688	0.42806\\
2.8656	0.42123\\
2.8623	0.41447\\
2.8591	0.40776\\
2.8559	0.40111\\
2.8526	0.39451\\
2.8494	0.38798\\
2.8461	0.3815\\
2.8429	0.37509\\
2.8397	0.36874\\
2.8364	0.36245\\
2.8332	0.35622\\
2.83	0.35005\\
2.8267	0.34395\\
2.8235	0.33791\\
2.8202	0.33193\\
2.817	0.32602\\
2.8138	0.32017\\
2.8105	0.31439\\
2.8073	0.30868\\
2.8041	0.30303\\
2.8008	0.29744\\
2.7976	0.29193\\
2.7943	0.28648\\
2.7911	0.28109\\
2.7879	0.27577\\
2.7846	0.27052\\
2.7814	0.26534\\
2.7782	0.26022\\
2.7749	0.25517\\
2.7717	0.25018\\
2.7684	0.24527\\
2.7652	0.24042\\
2.762	0.23563\\
2.7587	0.23091\\
2.7555	0.22626\\
2.7523	0.22168\\
2.749	0.21716\\
2.7458	0.2127\\
2.7425	0.20832\\
2.7393	0.20399\\
2.7361	0.19974\\
2.7328	0.19554\\
2.7296	0.19141\\
2.7263	0.18735\\
2.7231	0.18335\\
2.7199	0.17941\\
2.7166	0.17554\\
2.7134	0.17173\\
2.7102	0.16798\\
2.7069	0.16429\\
2.7037	0.16066\\
2.7004	0.1571\\
2.6972	0.15359\\
2.694	0.15015\\
2.6907	0.14676\\
2.6875	0.14344\\
2.6843	0.14017\\
2.681	0.13696\\
2.6778	0.13381\\
2.6745	0.13071\\
2.6713	0.12767\\
2.6681	0.12469\\
2.6648	0.12176\\
2.6616	0.11889\\
2.6584	0.11607\\
2.6551	0.1133\\
2.6519	0.11058\\
2.6486	0.10792\\
2.6454	0.10531\\
2.6422	0.10275\\
2.6389	0.10025\\
2.6357	0.097786\\
2.6325	0.095376\\
2.6292	0.093015\\
2.626	0.090701\\
2.6227	0.088435\\
2.6195	0.086215\\
2.6163	0.084041\\
2.613	0.081913\\
2.6098	0.079829\\
2.6066	0.077789\\
2.6033	0.075793\\
2.6001	0.07384\\
2.5968	0.071928\\
2.5936	0.070058\\
2.5904	0.068229\\
2.5871	0.06644\\
2.5839	0.064691\\
2.5806	0.06298\\
2.5774	0.061308\\
2.5742	0.059674\\
2.5709	0.058076\\
2.5677	0.056516\\
2.5645	0.05499\\
2.5612	0.053501\\
2.558	0.052045\\
2.5547	0.050624\\
2.5515	0.049236\\
2.5483	0.047881\\
2.545	0.046559\\
2.5418	0.045268\\
2.5386	0.044007\\
2.5353	0.042778\\
2.5321	0.041578\\
2.5288	0.040408\\
2.5256	0.039266\\
2.5224	0.038152\\
2.5191	0.037067\\
2.5159	0.036008\\
2.5127	0.034976\\
2.5094	0.03397\\
2.5062	0.032989\\
2.5029	0.032034\\
2.4997	0.031103\\
2.4965	0.030196\\
2.4932	0.029312\\
2.49	0.028451\\
2.4868	0.027613\\
2.4835	0.026797\\
2.4803	0.026002\\
2.477	0.025229\\
2.4738	0.024476\\
2.4706	0.023743\\
2.4673	0.02303\\
2.4641	0.022336\\
2.4608	0.021661\\
2.4576	0.021004\\
2.4544	0.020365\\
2.4511	0.019744\\
2.4479	0.01914\\
2.4447	0.018553\\
2.4414	0.017982\\
2.4382	0.017427\\
2.4349	0.016888\\
2.4317	0.016364\\
2.4285	0.015854\\
2.4252	0.015359\\
2.422	0.014879\\
2.4188	0.014412\\
2.4155	0.013958\\
2.4123	0.013518\\
2.409	0.01309\\
2.4058	0.012675\\
2.4026	0.012272\\
2.3993	0.011881\\
2.3961	0.011501\\
2.3929	0.011132\\
2.3896	0.010775\\
2.3864	0.010428\\
2.3831	0.010091\\
2.3799	0.009765\\
2.3767	0.0094484\\
2.3734	0.0091413\\
2.3702	0.0088435\\
2.367	0.0085548\\
2.3637	0.0082749\\
2.3605	0.0080035\\
2.3572	0.0077405\\
2.354	0.0074856\\
2.3508	0.0072386\\
2.3475	0.0069992\\
2.3443	0.0067673\\
2.341	0.0065426\\
2.3378	0.0063249\\
2.3346	0.0061141\\
2.3313	0.0059099\\
2.3281	0.0057122\\
2.3249	0.0055207\\
2.3216	0.0053353\\
2.3184	0.0051559\\
2.3151	0.0049821\\
2.3119	0.0048139\\
2.3087	0.0046512\\
2.3054	0.0044936\\
2.3022	0.0043412\\
2.299	0.0041937\\
2.2957	0.004051\\
2.2925	0.0039129\\
2.2892	0.0037794\\
2.286	0.0036502\\
2.2828	0.0035252\\
2.2795	0.0034044\\
2.2763	0.0032876\\
2.2731	0.0031746\\
2.2698	0.0030653\\
2.2666	0.0029597\\
2.2633	0.0028576\\
2.2601	0.0027589\\
2.2569	0.0026635\\
2.2536	0.0025713\\
2.2504	0.0024822\\
2.2472	0.0023961\\
2.2439	0.0023129\\
2.2407	0.0022325\\
2.2374	0.0021549\\
2.2342	0.0020798\\
2.231	0.0020073\\
2.2277	0.0019373\\
2.2245	0.0018697\\
2.2212	0.0018044\\
2.218	0.0017413\\
2.2148	0.0016804\\
2.2115	0.0016216\\
2.2083	0.0015648\\
2.2051	0.0015099\\
2.2018	0.001457\\
2.1986	0.0014059\\
2.1953	0.0013565\\
2.1921	0.0013089\\
2.1889	0.0012629\\
2.1856	0.0012185\\
2.1824	0.0011757\\
2.1792	0.0011344\\
2.1759	0.0010945\\
2.1727	0.001056\\
2.1694	0.0010188\\
2.1662	0.00098296\\
2.163	0.00094836\\
2.1597	0.00091499\\
2.1565	0.00088279\\
2.1533	0.00085172\\
2.15	0.00082174\\
2.1468	0.00079283\\
2.1435	0.00076493\\
2.1403	0.00073802\\
2.1371	0.00071206\\
2.1338	0.00068702\\
2.1306	0.00066287\\
2.1274	0.00063957\\
2.1241	0.0006171\\
2.1209	0.00059543\\
2.1176	0.00057452\\
2.1144	0.00055436\\
2.1112	0.00053491\\
2.1079	0.00051616\\
2.1047	0.00049807\\
2.1014	0.00048063\\
2.0982	0.00046381\\
2.095	0.00044758\\
2.0917	0.00043194\\
2.0885	0.00041685\\
2.0853	0.0004023\\
2.082	0.00038826\\
2.0788	0.00037473\\
2.0755	0.00036168\\
}--cycle;
\addplot[ybar interval, fill=black, fill opacity=0.35, area legend, draw=none] table[row sep=crcr, x=Lower, y=Count] {%
Lower	Upper	Count\\
2.0755	2.3244	0.075022\\
2.3244	2.5732	0.13933\\
2.5732	2.822	0.35368\\
2.822	3.0708	0.76094\\
3.0708	3.3196	0.94314\\
3.3196	3.5684	0.75022\\
3.5684	3.8172	0.52516\\
3.8172	4.0661	0.2465\\
4.0661	4.3149	0.10717\\
4.3149	4.5637	0.053587\\
4.5637	4.8125	0.04287\\
4.8125	5.0613	0\\
5.0613	5.3101	0.021435\\
5.3101	5.3101	0.021435\\
};
\addlegendentry{\RealHistogramLegend}

\addplot [color=black]
  table[row sep=crcr]{%
2.0755	0.016864\\
2.1112	0.021504\\
2.1435	0.026645\\
2.1727	0.03214\\
2.2018	0.038572\\
2.2277	0.045168\\
2.2536	0.05268\\
2.2795	0.061193\\
2.3022	0.069533\\
2.3249	0.078766\\
2.3475	0.088951\\
2.3702	0.10014\\
2.3929	0.11239\\
2.4155	0.12576\\
2.4382	0.14028\\
2.4608	0.15599\\
2.4835	0.17292\\
2.5062	0.19111\\
2.5288	0.21056\\
2.5515	0.23127\\
2.5742	0.25324\\
2.5968	0.27644\\
2.6195	0.30084\\
2.6454	0.33012\\
2.6713	0.3608\\
2.7004	0.39681\\
2.7328	0.43844\\
2.7684	0.48576\\
2.8202	0.55625\\
2.885	0.6442\\
2.9206	0.69095\\
2.9498	0.7276\\
2.9757	0.75858\\
2.9983	0.78418\\
3.0178	0.80486\\
3.0372	0.82423\\
3.0566	0.84218\\
3.0728	0.85596\\
3.089	0.86863\\
3.1052	0.88011\\
3.1214	0.89036\\
3.1376	0.89934\\
3.1505	0.90558\\
3.1635	0.91096\\
3.1764	0.91546\\
3.1894	0.91909\\
3.2023	0.92181\\
3.2153	0.92364\\
3.2282	0.92456\\
3.2412	0.92457\\
3.2541	0.92368\\
3.2671	0.92189\\
3.28	0.91921\\
3.293	0.91564\\
3.3059	0.91119\\
3.3189	0.90589\\
3.3318	0.89975\\
3.348	0.89093\\
3.3642	0.88086\\
3.3804	0.86961\\
3.3966	0.85722\\
3.4128	0.84375\\
3.4322	0.82626\\
3.4516	0.80742\\
3.471	0.78736\\
3.4937	0.76256\\
3.5196	0.73266\\
3.5488	0.6974\\
3.5811	0.65673\\
3.6297	0.59406\\
3.7042	0.49801\\
3.7398	0.45359\\
3.7722	0.4147\\
3.8013	0.38119\\
3.8272	0.35274\\
3.8531	0.32567\\
3.879	0.30006\\
3.9017	0.27888\\
3.9243	0.25887\\
3.947	0.24004\\
3.9697	0.22239\\
3.9923	0.20591\\
4.015	0.19056\\
4.0377	0.17632\\
4.0603	0.16315\\
4.083	0.151\\
4.1057	0.13982\\
4.1283	0.12956\\
4.1542	0.1189\\
4.1801	0.10927\\
4.206	0.10061\\
4.2352	0.091914\\
4.2643	0.084204\\
4.2967	0.076663\\
4.3291	0.070062\\
4.3647	0.063729\\
4.4035	0.057748\\
4.4489	0.051762\\
4.4974	0.046278\\
4.5557	0.04065\\
4.6237	0.035031\\
4.7047	0.029269\\
4.8018	0.023311\\
4.9054	0.017882\\
5.0123	0.013203\\
5.1191	0.0094246\\
5.2324	0.006341\\
5.3101	0.0047195\\
};
\addlegendentry{\RealBLNFitLegend}

      % \input{../../../../.tex/20/KS/SizeDistributionFittingsfig3.tex}
    \end{groupplot}
  
    \begin{groupplot}[
      group style={
        group name=Row2,
        plotGroup,
      },
      plotRow,
      row1Keys,
      yTickScaleLabelEmpty,
      %
      % xmin=1,xmax=2.3788e+05,
      ymin=0,ymax=1.0671,
    ]
      \nextgroupplot[%
        at={($(Row1 c1r1.south west)-(0,\yPlotSep em)-(0,\plotHeight cm)$)},
        xmin=0.57363,xmax=5.4115,
        % xmin=1.7165,xmax=5.4627,
        % ymin=0,ymax=1.2,
      ] % This file was created by matlab2tikz.
%
\definecolor{mycolor1}{rgb}{0.00000,0.44706,0.69804}%
\definecolor{mycolor2}{rgb}{0.83529,0.36863,0.00000}%
%
\addplot[ybar interval, fill=mycolor1, fill opacity=0.35, area legend, draw=none] table[row sep=crcr, x=Lower, y=Count] {%
Lower	Upper	Count\\
1.8868	2.1487	0.0014251\\
2.1487	2.4107	0.018221\\
2.4107	2.6727	0.12164\\
2.6727	2.9347	0.392\\
2.9347	3.1966	0.87011\\
3.1966	3.4586	1.0159\\
3.4586	3.7206	0.70572\\
3.7206	3.9826	0.35861\\
3.9826	4.2445	0.19035\\
4.2445	4.5065	0.091306\\
4.5065	4.7685	0.031861\\
4.7685	5.0305	0.0098737\\
5.0305	5.2924	0.010179\\
5.2924	5.2924	0.010179\\
};
\addlegendentry{\ExactHistogramLegend}

\addplot [color=mycolor1, only marks, every error bar/.append style={opacity=0.45}, mark=*, mark size=0pt, draw=none, forget plot]
 plot [error bars/.cd, y dir=both, y explicit, error bar style={line width=1pt, color=mycolor1}, error mark options={mark=none,mark size=0pt}]
 table[row sep=crcr, y error plus index=2, y error minus index=3]{%
2.0178	0.0014251	0.00081299	0.00081299\\
2.2797	0.018221	0.0026668	0.0026668\\
2.5417	0.12164	0.0054497	0.0054497\\
2.8037	0.392	0.010204	0.010204\\
3.0657	0.87011	0.012446	0.012446\\
3.3276	1.0159	0.012976	0.012976\\
3.5896	0.70572	0.010962	0.010962\\
3.8516	0.35861	0.0079653	0.0079653\\
4.1136	0.19035	0.0055851	0.0055851\\
4.3755	0.091306	0.0035213	0.0035213\\
4.6375	0.031861	0.0022358	0.0022358\\
4.8995	0.0098737	0.0013604	0.0013604\\
5.1615	0.010179	2.0756e-18	1.7347e-18\\
};
\addplot [color=mycolor1]
  table[row sep=crcr]{%
1.8868	0.00089722\\
2.0027	0.0020898\\
2.0811	0.0037453\\
2.1424	0.0059404\\
2.1902	0.0085109\\
2.2311	0.011572\\
2.2686	0.015307\\
2.2993	0.019201\\
2.3299	0.024026\\
2.3572	0.029248\\
2.3811	0.034661\\
2.4049	0.040976\\
2.4288	0.048314\\
2.4493	0.055516\\
2.4697	0.063649\\
2.4902	0.072803\\
2.5106	0.083067\\
2.5311	0.094533\\
2.5481	0.10507\\
2.5652	0.11657\\
2.5822	0.12906\\
2.5993	0.1426\\
2.6163	0.15724\\
2.6334	0.17302\\
2.6504	0.18997\\
2.6674	0.20811\\
2.6845	0.22749\\
2.7015	0.2481\\
2.7186	0.26995\\
2.739	0.29781\\
2.7595	0.32743\\
2.7799	0.35874\\
2.8004	0.39168\\
2.8243	0.43201\\
2.8481	0.47419\\
2.8754	0.5243\\
2.9095	0.58911\\
3.0254	0.81175\\
3.0493	0.85442\\
3.0697	0.88904\\
3.0868	0.91621\\
3.1038	0.94159\\
3.1174	0.96043\\
3.1311	0.97785\\
3.1447	0.99375\\
3.1584	1.008\\
3.1686	1.0176\\
3.1788	1.0262\\
3.189	1.0338\\
3.1993	1.0404\\
3.2095	1.0459\\
3.2197	1.0504\\
3.2299	1.0538\\
3.2402	1.0562\\
3.2504	1.0576\\
3.2606	1.0579\\
3.2709	1.0572\\
3.2811	1.0556\\
3.2913	1.053\\
3.3015	1.0494\\
3.3118	1.0449\\
3.322	1.0396\\
3.3322	1.0333\\
3.3459	1.0238\\
3.3595	1.0128\\
3.3731	1.0005\\
3.3868	0.98702\\
3.4004	0.97236\\
3.4174	0.95254\\
3.4345	0.93122\\
3.4549	0.90392\\
3.4754	0.87503\\
3.4993	0.83971\\
3.5299	0.79249\\
3.5845	0.70635\\
3.6288	0.63708\\
3.6595	0.59072\\
3.6868	0.55111\\
3.7106	0.51792\\
3.7345	0.48625\\
3.7584	0.4562\\
3.7788	0.43177\\
3.7993	0.40858\\
3.8197	0.38664\\
3.8402	0.36591\\
3.8606	0.34638\\
3.8811	0.328\\
3.9015	0.31072\\
3.922	0.29448\\
3.9459	0.27678\\
3.9697	0.26032\\
3.9936	0.245\\
4.0174	0.23073\\
4.0447	0.21557\\
4.072	0.2015\\
4.0993	0.18841\\
4.1299	0.1747\\
4.1606	0.16194\\
4.1947	0.14874\\
4.2322	0.13525\\
4.2697	0.1227\\
4.3106	0.10998\\
4.3515	0.098174\\
4.3959	0.086355\\
4.4402	0.075497\\
4.4845	0.06557\\
4.5288	0.056551\\
4.5731	0.048417\\
4.6209	0.040616\\
4.6686	0.033766\\
4.7163	0.027817\\
4.7674	0.022369\\
4.822	0.017521\\
4.8799	0.013336\\
4.9413	0.0098391\\
5.0095	0.006892\\
5.0879	0.0044704\\
5.1799	0.0026053\\
5.2924	0.0012862\\
};
\addlegendentry{\ExactBLNFitLegend}


\addplot[area legend, draw=none, fill=mycolor1, fill opacity=0.15, forget plot]
table[row sep=crcr] {%
x	y\\
1.8868	0.0011401\\
1.8902	0.001167\\
1.8936	0.0011944\\
1.897	0.0012226\\
1.9004	0.0012513\\
1.9038	0.0012808\\
1.9072	0.001311\\
1.9106	0.0013419\\
1.914	0.0013736\\
1.9174	0.001406\\
1.9209	0.0014391\\
1.9243	0.0014731\\
1.9277	0.0015079\\
1.9311	0.0015436\\
1.9345	0.0015801\\
1.9379	0.0016174\\
1.9413	0.0016557\\
1.9447	0.0016949\\
1.9481	0.0017351\\
1.9515	0.0017762\\
1.9549	0.0018183\\
1.9584	0.0018614\\
1.9618	0.0019056\\
1.9652	0.0019509\\
1.9686	0.0019973\\
1.972	0.0020447\\
1.9754	0.0020934\\
1.9788	0.0021432\\
1.9822	0.0021943\\
1.9856	0.0022466\\
1.989	0.0023002\\
1.9924	0.0023551\\
1.9959	0.0024113\\
1.9993	0.002469\\
2.0027	0.002528\\
2.0061	0.0025885\\
2.0095	0.0026505\\
2.0129	0.002714\\
2.0163	0.0027791\\
2.0197	0.0028457\\
2.0231	0.0029141\\
2.0265	0.0029841\\
2.0299	0.0030558\\
2.0334	0.0031293\\
2.0368	0.0032047\\
2.0402	0.0032819\\
2.0436	0.003361\\
2.047	0.0034421\\
2.0504	0.0035251\\
2.0538	0.0036103\\
2.0572	0.0036976\\
2.0606	0.003787\\
2.064	0.0038786\\
2.0674	0.0039726\\
2.0709	0.0040689\\
2.0743	0.0041675\\
2.0777	0.0042687\\
2.0811	0.0043723\\
2.0845	0.0044785\\
2.0879	0.0045874\\
2.0913	0.004699\\
2.0947	0.0048134\\
2.0981	0.0049306\\
2.1015	0.0050508\\
2.1049	0.005174\\
2.1084	0.0053002\\
2.1118	0.0054296\\
2.1152	0.0055622\\
2.1186	0.0056981\\
2.122	0.0058374\\
2.1254	0.0059802\\
2.1288	0.0061265\\
2.1322	0.0062765\\
2.1356	0.0064303\\
2.139	0.0065879\\
2.1424	0.0067494\\
2.1459	0.0069149\\
2.1493	0.0070845\\
2.1527	0.0072584\\
2.1561	0.0074366\\
2.1595	0.0076193\\
2.1629	0.0078065\\
2.1663	0.0079983\\
2.1697	0.0081949\\
2.1731	0.0083964\\
2.1765	0.0086029\\
2.1799	0.0088145\\
2.1834	0.0090313\\
2.1868	0.0092536\\
2.1902	0.0094813\\
2.1936	0.0097146\\
2.197	0.0099537\\
2.2004	0.010199\\
2.2038	0.01045\\
2.2072	0.010707\\
2.2106	0.010971\\
2.214	0.011241\\
2.2174	0.011517\\
2.2209	0.011801\\
2.2243	0.012091\\
2.2277	0.012389\\
2.2311	0.012693\\
2.2345	0.013006\\
2.2379	0.013325\\
2.2413	0.013653\\
2.2447	0.013988\\
2.2481	0.014332\\
2.2515	0.014684\\
2.2549	0.015045\\
2.2584	0.015414\\
2.2618	0.015792\\
2.2652	0.016179\\
2.2686	0.016575\\
2.272	0.016981\\
2.2754	0.017397\\
2.2788	0.017823\\
2.2822	0.018258\\
2.2856	0.018704\\
2.289	0.019161\\
2.2924	0.019629\\
2.2959	0.020107\\
2.2993	0.020597\\
2.3027	0.021098\\
2.3061	0.021611\\
2.3095	0.022137\\
2.3129	0.022674\\
2.3163	0.023224\\
2.3197	0.023787\\
2.3231	0.024362\\
2.3265	0.024951\\
2.3299	0.025554\\
2.3334	0.02617\\
2.3368	0.026801\\
2.3402	0.027446\\
2.3436	0.028106\\
2.347	0.028781\\
2.3504	0.029471\\
2.3538	0.030177\\
2.3572	0.030898\\
2.3606	0.031636\\
2.364	0.032391\\
2.3674	0.033162\\
2.3709	0.033951\\
2.3743	0.034757\\
2.3777	0.03558\\
2.3811	0.036422\\
2.3845	0.037283\\
2.3879	0.038162\\
2.3913	0.039061\\
2.3947	0.039979\\
2.3981	0.040917\\
2.4015	0.041875\\
2.4049	0.042854\\
2.4084	0.043854\\
2.4118	0.044876\\
2.4152	0.045919\\
2.4186	0.046984\\
2.422	0.048071\\
2.4254	0.049182\\
2.4288	0.050316\\
2.4322	0.051473\\
2.4356	0.052655\\
2.439	0.053861\\
2.4424	0.055091\\
2.4459	0.056347\\
2.4493	0.057629\\
2.4527	0.058937\\
2.4561	0.060271\\
2.4595	0.061632\\
2.4629	0.063021\\
2.4663	0.064437\\
2.4697	0.065881\\
2.4731	0.067354\\
2.4765	0.068855\\
2.4799	0.070387\\
2.4834	0.071948\\
2.4868	0.073539\\
2.4902	0.075161\\
2.4936	0.076814\\
2.497	0.078499\\
2.5004	0.080215\\
2.5038	0.081964\\
2.5072	0.083746\\
2.5106	0.085562\\
2.514	0.087411\\
2.5174	0.089294\\
2.5209	0.091211\\
2.5243	0.093164\\
2.5277	0.095152\\
2.5311	0.097176\\
2.5345	0.099237\\
2.5379	0.10133\\
2.5413	0.10347\\
2.5447	0.10564\\
2.5481	0.10785\\
2.5515	0.1101\\
2.5549	0.11239\\
2.5584	0.11471\\
2.5618	0.11708\\
2.5652	0.11948\\
2.5686	0.12193\\
2.572	0.12442\\
2.5754	0.12695\\
2.5788	0.12952\\
2.5822	0.13213\\
2.5856	0.13478\\
2.589	0.13748\\
2.5924	0.14022\\
2.5959	0.14301\\
2.5993	0.14584\\
2.6027	0.14871\\
2.6061	0.15163\\
2.6095	0.15459\\
2.6129	0.1576\\
2.6163	0.16065\\
2.6197	0.16375\\
2.6231	0.16689\\
2.6265	0.17009\\
2.6299	0.17332\\
2.6334	0.17661\\
2.6368	0.17994\\
2.6402	0.18332\\
2.6436	0.18675\\
2.647	0.19023\\
2.6504	0.19375\\
2.6538	0.19733\\
2.6572	0.20095\\
2.6606	0.20462\\
2.664	0.20834\\
2.6674	0.21211\\
2.6709	0.21593\\
2.6743	0.21979\\
2.6777	0.22371\\
2.6811	0.22768\\
2.6845	0.2317\\
2.6879	0.23576\\
2.6913	0.23988\\
2.6947	0.24405\\
2.6981	0.24827\\
2.7015	0.25254\\
2.7049	0.25685\\
2.7084	0.26122\\
2.7118	0.26564\\
2.7152	0.27011\\
2.7186	0.27463\\
2.722	0.2792\\
2.7254	0.28381\\
2.7288	0.28848\\
2.7322	0.2932\\
2.7356	0.29797\\
2.739	0.30278\\
2.7424	0.30765\\
2.7459	0.31256\\
2.7493	0.31753\\
2.7527	0.32254\\
2.7561	0.3276\\
2.7595	0.3327\\
2.7629	0.33786\\
2.7663	0.34306\\
2.7697	0.34831\\
2.7731	0.3536\\
2.7765	0.35894\\
2.7799	0.36433\\
2.7834	0.36976\\
2.7868	0.37524\\
2.7902	0.38076\\
2.7936	0.38633\\
2.797	0.39193\\
2.8004	0.39758\\
2.8038	0.40328\\
2.8072	0.40901\\
2.8106	0.41479\\
2.814	0.4206\\
2.8174	0.42645\\
2.8209	0.43235\\
2.8243	0.43828\\
2.8277	0.44425\\
2.8311	0.45025\\
2.8345	0.45629\\
2.8379	0.46237\\
2.8413	0.46848\\
2.8447	0.47462\\
2.8481	0.4808\\
2.8515	0.48701\\
2.8549	0.49324\\
2.8584	0.49951\\
2.8618	0.50581\\
2.8652	0.51213\\
2.8686	0.51848\\
2.872	0.52486\\
2.8754	0.53126\\
2.8788	0.53768\\
2.8822	0.54412\\
2.8856	0.55059\\
2.889	0.55707\\
2.8924	0.56358\\
2.8959	0.5701\\
2.8993	0.57664\\
2.9027	0.58319\\
2.9061	0.58976\\
2.9095	0.59634\\
2.9129	0.60293\\
2.9163	0.60953\\
2.9197	0.61614\\
2.9231	0.62275\\
2.9265	0.62937\\
2.9299	0.636\\
2.9334	0.64263\\
2.9368	0.64926\\
2.9402	0.65589\\
2.9436	0.66252\\
2.947	0.66914\\
2.9504	0.67577\\
2.9538	0.68238\\
2.9572	0.68899\\
2.9606	0.6956\\
2.964	0.70219\\
2.9674	0.70877\\
2.9709	0.71534\\
2.9743	0.72189\\
2.9777	0.72843\\
2.9811	0.73496\\
2.9845	0.74147\\
2.9879	0.74795\\
2.9913	0.75442\\
2.9947	0.76087\\
2.9981	0.76729\\
3.0015	0.7737\\
3.0049	0.78007\\
3.0084	0.78643\\
3.0118	0.79275\\
3.0152	0.79905\\
3.0186	0.80533\\
3.022	0.81157\\
3.0254	0.81779\\
3.0288	0.82397\\
3.0322	0.83013\\
3.0356	0.83625\\
3.039	0.84234\\
3.0424	0.8484\\
3.0459	0.85443\\
3.0493	0.86041\\
3.0527	0.86636\\
3.0561	0.87228\\
3.0595	0.87815\\
3.0629	0.88398\\
3.0663	0.88976\\
3.0697	0.8955\\
3.0731	0.90119\\
3.0765	0.90682\\
3.0799	0.91241\\
3.0834	0.91793\\
3.0868	0.92339\\
3.0902	0.92879\\
3.0936	0.93412\\
3.097	0.93937\\
3.1004	0.94456\\
3.1038	0.94967\\
3.1072	0.9547\\
3.1106	0.95964\\
3.114	0.9645\\
3.1174	0.96928\\
3.1209	0.97396\\
3.1243	0.97854\\
3.1277	0.98303\\
3.1311	0.98742\\
3.1345	0.99171\\
3.1379	0.99589\\
3.1413	0.99997\\
3.1447	1.0039\\
3.1481	1.0078\\
3.1515	1.0115\\
3.1549	1.0152\\
3.1584	1.0187\\
3.1618	1.0221\\
3.1652	1.0254\\
3.1686	1.0285\\
3.172	1.0315\\
3.1754	1.0345\\
3.1788	1.0372\\
3.1822	1.0399\\
3.1856	1.0425\\
3.189	1.0449\\
3.1924	1.0472\\
3.1959	1.0493\\
3.1993	1.0514\\
3.2027	1.0533\\
3.2061	1.055\\
3.2095	1.0567\\
3.2129	1.0582\\
3.2163	1.0596\\
3.2197	1.0609\\
3.2231	1.0621\\
3.2265	1.0631\\
3.2299	1.064\\
3.2334	1.0648\\
3.2368	1.0655\\
3.2402	1.0661\\
3.2436	1.0665\\
3.247	1.0668\\
3.2504	1.067\\
3.2538	1.0671\\
3.2572	1.0671\\
3.2606	1.067\\
3.264	1.0667\\
3.2674	1.0664\\
3.2709	1.0659\\
3.2743	1.0653\\
3.2777	1.0647\\
3.2811	1.0639\\
3.2845	1.063\\
3.2879	1.0621\\
3.2913	1.061\\
3.2947	1.0598\\
3.2981	1.0585\\
3.3015	1.0572\\
3.3049	1.0557\\
3.3084	1.0542\\
3.3118	1.0525\\
3.3152	1.0508\\
3.3186	1.0489\\
3.322	1.047\\
3.3254	1.045\\
3.3288	1.0429\\
3.3322	1.0407\\
3.3356	1.0385\\
3.339	1.0361\\
3.3424	1.0336\\
3.3459	1.0311\\
3.3493	1.0285\\
3.3527	1.0258\\
3.3561	1.023\\
3.3595	1.0201\\
3.3629	1.0172\\
3.3663	1.0142\\
3.3697	1.0111\\
3.3731	1.0079\\
3.3765	1.0046\\
3.3799	1.0013\\
3.3834	0.99786\\
3.3868	0.99437\\
3.3902	0.99082\\
3.3936	0.98719\\
3.397	0.98349\\
3.4004	0.97972\\
3.4038	0.97589\\
3.4072	0.97199\\
3.4106	0.96802\\
3.414	0.96399\\
3.4174	0.9599\\
3.4209	0.95575\\
3.4243	0.95154\\
3.4277	0.94727\\
3.4311	0.94295\\
3.4345	0.93857\\
3.4379	0.93414\\
3.4413	0.92965\\
3.4447	0.92512\\
3.4481	0.92053\\
3.4515	0.9159\\
3.4549	0.91123\\
3.4584	0.9065\\
3.4618	0.90174\\
3.4652	0.89693\\
3.4686	0.89208\\
3.472	0.8872\\
3.4754	0.88227\\
3.4788	0.87731\\
3.4822	0.87232\\
3.4856	0.86729\\
3.489	0.86223\\
3.4924	0.85715\\
3.4959	0.85203\\
3.4993	0.84688\\
3.5027	0.84171\\
3.5061	0.83651\\
3.5095	0.83129\\
3.5129	0.82605\\
3.5163	0.82079\\
3.5197	0.81551\\
3.5231	0.81021\\
3.5265	0.80489\\
3.5299	0.79956\\
3.5334	0.79421\\
3.5368	0.78885\\
3.5402	0.78348\\
3.5436	0.7781\\
3.547	0.77271\\
3.5504	0.76731\\
3.5538	0.7619\\
3.5572	0.75649\\
3.5606	0.75108\\
3.564	0.74566\\
3.5674	0.74024\\
3.5709	0.73482\\
3.5743	0.7294\\
3.5777	0.72398\\
3.5811	0.71856\\
3.5845	0.71315\\
3.5879	0.70774\\
3.5913	0.70233\\
3.5947	0.69694\\
3.5981	0.69155\\
3.6015	0.68617\\
3.6049	0.6808\\
3.6084	0.67544\\
3.6118	0.67009\\
3.6152	0.66475\\
3.6186	0.65942\\
3.622	0.65411\\
3.6254	0.64882\\
3.6288	0.64354\\
3.6322	0.63828\\
3.6356	0.63303\\
3.639	0.62781\\
3.6424	0.6226\\
3.6459	0.61741\\
3.6493	0.61224\\
3.6527	0.6071\\
3.6561	0.60197\\
3.6595	0.59687\\
3.6629	0.59179\\
3.6663	0.58674\\
3.6697	0.58171\\
3.6731	0.5767\\
3.6765	0.57172\\
3.6799	0.56677\\
3.6834	0.56184\\
3.6868	0.55694\\
3.6902	0.55207\\
3.6936	0.54723\\
3.697	0.54241\\
3.7004	0.53763\\
3.7038	0.53287\\
3.7072	0.52815\\
3.7106	0.52345\\
3.714	0.51879\\
3.7174	0.51415\\
3.7209	0.50955\\
3.7243	0.50498\\
3.7277	0.50045\\
3.7311	0.49594\\
3.7345	0.49147\\
3.7379	0.48703\\
3.7413	0.48263\\
3.7447	0.47825\\
3.7481	0.47392\\
3.7515	0.46961\\
3.7549	0.46534\\
3.7584	0.46111\\
3.7618	0.45691\\
3.7652	0.45274\\
3.7686	0.44861\\
3.772	0.44451\\
3.7754	0.44045\\
3.7788	0.43642\\
3.7822	0.43243\\
3.7856	0.42848\\
3.789	0.42455\\
3.7924	0.42067\\
3.7959	0.41682\\
3.7993	0.413\\
3.8027	0.40922\\
3.8061	0.40547\\
3.8095	0.40176\\
3.8129	0.39809\\
3.8163	0.39445\\
3.8197	0.39084\\
3.8231	0.38727\\
3.8265	0.38373\\
3.8299	0.38023\\
3.8334	0.37676\\
3.8368	0.37333\\
3.8402	0.36993\\
3.8436	0.36656\\
3.847	0.36323\\
3.8504	0.35993\\
3.8538	0.35667\\
3.8572	0.35343\\
3.8606	0.35023\\
3.864	0.34707\\
3.8674	0.34393\\
3.8709	0.34083\\
3.8743	0.33776\\
3.8777	0.33473\\
3.8811	0.33172\\
3.8845	0.32874\\
3.8879	0.3258\\
3.8913	0.32289\\
3.8947	0.32\\
3.8981	0.31715\\
3.9015	0.31433\\
3.9049	0.31153\\
3.9084	0.30877\\
3.9118	0.30603\\
3.9152	0.30333\\
3.9186	0.30065\\
3.922	0.298\\
3.9254	0.29538\\
3.9288	0.29278\\
3.9322	0.29021\\
3.9356	0.28767\\
3.939	0.28516\\
3.9424	0.28267\\
3.9459	0.28021\\
3.9493	0.27777\\
3.9527	0.27536\\
3.9561	0.27297\\
3.9595	0.27061\\
3.9629	0.26827\\
3.9663	0.26595\\
3.9697	0.26366\\
3.9731	0.2614\\
3.9765	0.25915\\
3.9799	0.25693\\
3.9834	0.25473\\
3.9868	0.25255\\
3.9902	0.2504\\
3.9936	0.24826\\
3.997	0.24615\\
4.0004	0.24406\\
4.0038	0.24199\\
4.0072	0.23993\\
4.0106	0.2379\\
4.014	0.23589\\
4.0174	0.2339\\
4.0209	0.23192\\
4.0243	0.22997\\
4.0277	0.22803\\
4.0311	0.22611\\
4.0345	0.22421\\
4.0379	0.22233\\
4.0413	0.22047\\
4.0447	0.21862\\
4.0481	0.21679\\
4.0515	0.21497\\
4.0549	0.21317\\
4.0584	0.21139\\
4.0618	0.20962\\
4.0652	0.20787\\
4.0686	0.20614\\
4.072	0.20442\\
4.0754	0.20271\\
4.0788	0.20102\\
4.0822	0.19935\\
4.0856	0.19769\\
4.089	0.19604\\
4.0924	0.1944\\
4.0959	0.19278\\
4.0993	0.19118\\
4.1027	0.18958\\
4.1061	0.188\\
4.1095	0.18643\\
4.1129	0.18488\\
4.1163	0.18333\\
4.1197	0.1818\\
4.1231	0.18028\\
4.1265	0.17878\\
4.1299	0.17728\\
4.1334	0.1758\\
4.1368	0.17433\\
4.1402	0.17287\\
4.1436	0.17141\\
4.147	0.16998\\
4.1504	0.16855\\
4.1538	0.16713\\
4.1572	0.16572\\
4.1606	0.16432\\
4.164	0.16294\\
4.1674	0.16156\\
4.1709	0.16019\\
4.1743	0.15883\\
4.1777	0.15749\\
4.1811	0.15615\\
4.1845	0.15482\\
4.1879	0.1535\\
4.1913	0.15219\\
4.1947	0.15088\\
4.1981	0.14959\\
4.2015	0.14831\\
4.2049	0.14703\\
4.2084	0.14576\\
4.2118	0.14451\\
4.2152	0.14326\\
4.2186	0.14201\\
4.222	0.14078\\
4.2254	0.13956\\
4.2288	0.13834\\
4.2322	0.13713\\
4.2356	0.13593\\
4.239	0.13473\\
4.2424	0.13355\\
4.2459	0.13237\\
4.2493	0.1312\\
4.2527	0.13003\\
4.2561	0.12888\\
4.2595	0.12773\\
4.2629	0.12659\\
4.2663	0.12546\\
4.2697	0.12433\\
4.2731	0.12321\\
4.2765	0.1221\\
4.2799	0.12099\\
4.2834	0.11989\\
4.2868	0.1188\\
4.2902	0.11772\\
4.2936	0.11664\\
4.297	0.11557\\
4.3004	0.1145\\
4.3038	0.11345\\
4.3072	0.1124\\
4.3106	0.11135\\
4.314	0.11031\\
4.3174	0.10928\\
4.3209	0.10826\\
4.3243	0.10724\\
4.3277	0.10623\\
4.3311	0.10522\\
4.3345	0.10423\\
4.3379	0.10323\\
4.3413	0.10225\\
4.3447	0.10127\\
4.3481	0.1003\\
4.3515	0.099328\\
4.3549	0.098368\\
4.3584	0.097413\\
4.3618	0.096465\\
4.3652	0.095523\\
4.3686	0.094587\\
4.372	0.093657\\
4.3754	0.092733\\
4.3788	0.091816\\
4.3822	0.090904\\
4.3856	0.089999\\
4.389	0.089099\\
4.3924	0.088206\\
4.3959	0.087319\\
4.3993	0.086437\\
4.4027	0.085562\\
4.4061	0.084692\\
4.4095	0.083829\\
4.4129	0.082971\\
4.4163	0.082119\\
4.4197	0.081274\\
4.4231	0.080434\\
4.4265	0.0796\\
4.4299	0.078771\\
4.4334	0.077949\\
4.4368	0.077132\\
4.4402	0.076322\\
4.4436	0.075517\\
4.447	0.074717\\
4.4504	0.073924\\
4.4538	0.073136\\
4.4572	0.072354\\
4.4606	0.071578\\
4.464	0.070807\\
4.4674	0.070042\\
4.4709	0.069283\\
4.4743	0.068529\\
4.4777	0.067781\\
4.4811	0.067038\\
4.4845	0.066301\\
4.4879	0.06557\\
4.4913	0.064844\\
4.4947	0.064124\\
4.4981	0.063409\\
4.5015	0.0627\\
4.5049	0.061996\\
4.5084	0.061298\\
4.5118	0.060605\\
4.5152	0.059918\\
4.5186	0.059236\\
4.522	0.058559\\
4.5254	0.057888\\
4.5288	0.057222\\
4.5322	0.056562\\
4.5356	0.055907\\
4.539	0.055257\\
4.5424	0.054613\\
4.5459	0.053973\\
4.5493	0.05334\\
4.5527	0.052711\\
4.5561	0.052087\\
4.5595	0.051469\\
4.5629	0.050856\\
4.5663	0.050248\\
4.5697	0.049646\\
4.5731	0.049048\\
4.5765	0.048455\\
4.5799	0.047868\\
4.5834	0.047286\\
4.5868	0.046709\\
4.5902	0.046137\\
4.5936	0.045569\\
4.597	0.045007\\
4.6004	0.04445\\
4.6038	0.043898\\
4.6072	0.043351\\
4.6106	0.042809\\
4.614	0.042271\\
4.6174	0.041739\\
4.6209	0.041211\\
4.6243	0.040689\\
4.6277	0.040171\\
4.6311	0.039658\\
4.6345	0.03915\\
4.6379	0.038646\\
4.6413	0.038147\\
4.6447	0.037653\\
4.6481	0.037164\\
4.6515	0.03668\\
4.6549	0.0362\\
4.6584	0.035725\\
4.6618	0.035254\\
4.6652	0.034788\\
4.6686	0.034327\\
4.672	0.03387\\
4.6754	0.033418\\
4.6788	0.03297\\
4.6822	0.032527\\
4.6856	0.032088\\
4.689	0.031654\\
4.6924	0.031224\\
4.6959	0.030799\\
4.6993	0.030378\\
4.7027	0.029961\\
4.7061	0.029549\\
4.7095	0.029141\\
4.7129	0.028737\\
4.7163	0.028338\\
4.7197	0.027943\\
4.7231	0.027552\\
4.7265	0.027165\\
4.7299	0.026783\\
4.7334	0.026405\\
4.7368	0.02603\\
4.7402	0.02566\\
4.7436	0.025294\\
4.747	0.024932\\
4.7504	0.024574\\
4.7538	0.024221\\
4.7572	0.023871\\
4.7606	0.023525\\
4.764	0.023183\\
4.7674	0.022844\\
4.7709	0.02251\\
4.7743	0.02218\\
4.7777	0.021853\\
4.7811	0.02153\\
4.7845	0.021211\\
4.7879	0.020896\\
4.7913	0.020584\\
4.7947	0.020277\\
4.7981	0.019972\\
4.8015	0.019672\\
4.8049	0.019375\\
4.8083	0.019081\\
4.8118	0.018792\\
4.8152	0.018505\\
4.8186	0.018223\\
4.822	0.017943\\
4.8254	0.017667\\
4.8288	0.017395\\
4.8322	0.017126\\
4.8356	0.01686\\
4.839	0.016598\\
4.8424	0.016339\\
4.8458	0.016083\\
4.8493	0.015831\\
4.8527	0.015581\\
4.8561	0.015335\\
4.8595	0.015092\\
4.8629	0.014853\\
4.8663	0.014616\\
4.8697	0.014383\\
4.8731	0.014152\\
4.8765	0.013925\\
4.8799	0.0137\\
4.8833	0.013479\\
4.8868	0.01326\\
4.8902	0.013045\\
4.8936	0.012832\\
4.897	0.012622\\
4.9004	0.012416\\
4.9038	0.012211\\
4.9072	0.01201\\
4.9106	0.011812\\
4.914	0.011616\\
4.9174	0.011423\\
4.9208	0.011232\\
4.9243	0.011044\\
4.9277	0.010859\\
4.9311	0.010677\\
4.9345	0.010497\\
4.9379	0.010319\\
4.9413	0.010144\\
4.9447	0.0099721\\
4.9481	0.0098022\\
4.9515	0.0096347\\
4.9549	0.0094696\\
4.9583	0.0093069\\
4.9618	0.0091466\\
4.9652	0.0089887\\
4.9686	0.008833\\
4.972	0.0086797\\
4.9754	0.0085286\\
4.9788	0.0083797\\
4.9822	0.0082331\\
4.9856	0.0080886\\
4.989	0.0079463\\
4.9924	0.0078062\\
4.9958	0.0076681\\
4.9993	0.0075322\\
5.0027	0.0073983\\
5.0061	0.0072665\\
5.0095	0.0071366\\
5.0129	0.0070088\\
5.0163	0.0068829\\
5.0197	0.006759\\
5.0231	0.006637\\
5.0265	0.0065169\\
5.0299	0.0063987\\
5.0333	0.0062824\\
5.0368	0.0061678\\
5.0402	0.0060551\\
5.0436	0.0059442\\
5.047	0.005835\\
5.0504	0.0057275\\
5.0538	0.0056218\\
5.0572	0.0055178\\
5.0606	0.0054154\\
5.064	0.0053147\\
5.0674	0.0052157\\
5.0708	0.0051182\\
5.0743	0.0050224\\
5.0777	0.0049281\\
5.0811	0.0048353\\
5.0845	0.0047441\\
5.0879	0.0046544\\
5.0913	0.0045661\\
5.0947	0.0044794\\
5.0981	0.0043941\\
5.1015	0.0043102\\
5.1049	0.0042277\\
5.1083	0.0041466\\
5.1118	0.0040669\\
5.1152	0.0039885\\
5.1186	0.0039115\\
5.122	0.0038357\\
5.1254	0.0037613\\
5.1288	0.0036881\\
5.1322	0.0036162\\
5.1356	0.0035456\\
5.139	0.0034761\\
5.1424	0.0034079\\
5.1458	0.0033409\\
5.1493	0.003275\\
5.1527	0.0032103\\
5.1561	0.0031467\\
5.1595	0.0030842\\
5.1629	0.0030228\\
5.1663	0.0029626\\
5.1697	0.0029034\\
5.1731	0.0028452\\
5.1765	0.0027881\\
5.1799	0.002732\\
5.1833	0.0026769\\
5.1868	0.0026228\\
5.1902	0.0025697\\
5.1936	0.0025175\\
5.197	0.0024663\\
5.2004	0.0024161\\
5.2038	0.0023667\\
5.2072	0.0023183\\
5.2106	0.0022707\\
5.214	0.002224\\
5.2174	0.0021782\\
5.2208	0.0021332\\
5.2243	0.0020891\\
5.2277	0.0020457\\
5.2311	0.0020032\\
5.2345	0.0019615\\
5.2379	0.0019206\\
5.2413	0.0018804\\
5.2447	0.001841\\
5.2481	0.0018024\\
5.2515	0.0017644\\
5.2549	0.0017272\\
5.2583	0.0016907\\
5.2618	0.0016549\\
5.2652	0.0016198\\
5.2686	0.0015854\\
5.272	0.0015516\\
5.2754	0.0015185\\
5.2788	0.001486\\
5.2822	0.0014542\\
5.2856	0.0014229\\
5.289	0.0013923\\
5.2924	0.0013623\\
5.2924	0.0012102\\
5.289	0.0012377\\
5.2856	0.0012658\\
5.2822	0.0012944\\
5.2788	0.0013237\\
5.2754	0.0013535\\
5.272	0.001384\\
5.2686	0.0014151\\
5.2652	0.0014468\\
5.2618	0.0014791\\
5.2583	0.0015121\\
5.2549	0.0015458\\
5.2515	0.0015801\\
5.2481	0.0016151\\
5.2447	0.0016508\\
5.2413	0.0016872\\
5.2379	0.0017244\\
5.2345	0.0017622\\
5.2311	0.0018009\\
5.2277	0.0018402\\
5.2243	0.0018804\\
5.2208	0.0019213\\
5.2174	0.0019631\\
5.214	0.0020056\\
5.2106	0.002049\\
5.2072	0.0020932\\
5.2038	0.0021382\\
5.2004	0.0021842\\
5.197	0.002231\\
5.1936	0.0022786\\
5.1902	0.0023273\\
5.1868	0.0023768\\
5.1833	0.0024272\\
5.1799	0.0024787\\
5.1765	0.002531\\
5.1731	0.0025844\\
5.1697	0.0026388\\
5.1663	0.0026941\\
5.1629	0.0027505\\
5.1595	0.002808\\
5.1561	0.0028665\\
5.1527	0.0029261\\
5.1493	0.0029868\\
5.1458	0.0030486\\
5.1424	0.0031115\\
5.139	0.0031756\\
5.1356	0.0032408\\
5.1322	0.0033073\\
5.1288	0.0033749\\
5.1254	0.0034437\\
5.122	0.0035138\\
5.1186	0.0035851\\
5.1152	0.0036577\\
5.1118	0.0037316\\
5.1083	0.0038068\\
5.1049	0.0038833\\
5.1015	0.0039612\\
5.0981	0.0040404\\
5.0947	0.004121\\
5.0913	0.004203\\
5.0879	0.0042865\\
5.0845	0.0043713\\
5.0811	0.0044577\\
5.0777	0.0045455\\
5.0743	0.0046349\\
5.0708	0.0047258\\
5.0674	0.0048182\\
5.064	0.0049122\\
5.0606	0.0050077\\
5.0572	0.0051049\\
5.0538	0.0052038\\
5.0504	0.0053042\\
5.047	0.0054064\\
5.0436	0.0055103\\
5.0402	0.0056158\\
5.0368	0.0057232\\
5.0333	0.0058323\\
5.0299	0.0059432\\
5.0265	0.0060558\\
5.0231	0.0061704\\
5.0197	0.0062868\\
5.0163	0.0064051\\
5.0129	0.0065252\\
5.0095	0.0066474\\
5.0061	0.0067714\\
5.0027	0.0068975\\
4.9993	0.0070256\\
4.9958	0.0071556\\
4.9924	0.0072878\\
4.989	0.007422\\
4.9856	0.0075583\\
4.9822	0.0076968\\
4.9788	0.0078374\\
4.9754	0.0079802\\
4.972	0.0081252\\
4.9686	0.0082724\\
4.9652	0.0084219\\
4.9618	0.0085737\\
4.9583	0.0087277\\
4.9549	0.0088842\\
4.9515	0.0090429\\
4.9481	0.0092041\\
4.9447	0.0093677\\
4.9413	0.0095337\\
4.9379	0.0097022\\
4.9345	0.0098732\\
4.9311	0.010047\\
4.9277	0.010223\\
4.9243	0.010401\\
4.9208	0.010583\\
4.9174	0.010767\\
4.914	0.010953\\
4.9106	0.011142\\
4.9072	0.011334\\
4.9038	0.011529\\
4.9004	0.011726\\
4.897	0.011927\\
4.8936	0.01213\\
4.8902	0.012336\\
4.8868	0.012545\\
4.8833	0.012756\\
4.8799	0.012971\\
4.8765	0.013189\\
4.8731	0.013409\\
4.8697	0.013633\\
4.8663	0.01386\\
4.8629	0.01409\\
4.8595	0.014323\\
4.8561	0.014559\\
4.8527	0.014798\\
4.8493	0.01504\\
4.8458	0.015286\\
4.8424	0.015535\\
4.839	0.015787\\
4.8356	0.016043\\
4.8322	0.016302\\
4.8288	0.016564\\
4.8254	0.01683\\
4.822	0.017099\\
4.8186	0.017371\\
4.8152	0.017647\\
4.8118	0.017927\\
4.8083	0.01821\\
4.8049	0.018497\\
4.8015	0.018787\\
4.7981	0.019081\\
4.7947	0.019378\\
4.7913	0.01968\\
4.7879	0.019985\\
4.7845	0.020293\\
4.7811	0.020606\\
4.7777	0.020922\\
4.7743	0.021242\\
4.7709	0.021566\\
4.7674	0.021894\\
4.764	0.022226\\
4.7606	0.022561\\
4.7572	0.022901\\
4.7538	0.023245\\
4.7504	0.023592\\
4.747	0.023944\\
4.7436	0.0243\\
4.7402	0.024659\\
4.7368	0.025023\\
4.7334	0.025391\\
4.7299	0.025764\\
4.7265	0.02614\\
4.7231	0.026521\\
4.7197	0.026906\\
4.7163	0.027295\\
4.7129	0.027689\\
4.7095	0.028086\\
4.7061	0.028489\\
4.7027	0.028895\\
4.6993	0.029306\\
4.6959	0.029721\\
4.6924	0.030141\\
4.689	0.030565\\
4.6856	0.030994\\
4.6822	0.031427\\
4.6788	0.031865\\
4.6754	0.032308\\
4.672	0.032754\\
4.6686	0.033206\\
4.6652	0.033662\\
4.6618	0.034123\\
4.6584	0.034588\\
4.6549	0.035058\\
4.6515	0.035533\\
4.6481	0.036013\\
4.6447	0.036497\\
4.6413	0.036986\\
4.6379	0.037479\\
4.6345	0.037978\\
4.6311	0.038481\\
4.6277	0.038989\\
4.6243	0.039502\\
4.6209	0.04002\\
4.6174	0.040543\\
4.614	0.04107\\
4.6106	0.041603\\
4.6072	0.04214\\
4.6038	0.042682\\
4.6004	0.04323\\
4.597	0.043782\\
4.5936	0.044339\\
4.5902	0.044901\\
4.5868	0.045468\\
4.5834	0.04604\\
4.5799	0.046617\\
4.5765	0.047199\\
4.5731	0.047786\\
4.5697	0.048379\\
4.5663	0.048976\\
4.5629	0.049578\\
4.5595	0.050185\\
4.5561	0.050798\\
4.5527	0.051415\\
4.5493	0.052038\\
4.5459	0.052665\\
4.5424	0.053298\\
4.539	0.053936\\
4.5356	0.054579\\
4.5322	0.055227\\
4.5288	0.05588\\
4.5254	0.056538\\
4.522	0.057202\\
4.5186	0.057871\\
4.5152	0.058544\\
4.5118	0.059223\\
4.5084	0.059907\\
4.5049	0.060596\\
4.5015	0.061291\\
4.4981	0.06199\\
4.4947	0.062695\\
4.4913	0.063404\\
4.4879	0.064119\\
4.4845	0.06484\\
4.4811	0.065565\\
4.4777	0.066295\\
4.4743	0.067031\\
4.4709	0.067772\\
4.4674	0.068518\\
4.464	0.069269\\
4.4606	0.070025\\
4.4572	0.070787\\
4.4538	0.071554\\
4.4504	0.072326\\
4.447	0.073103\\
4.4436	0.073885\\
4.4402	0.074673\\
4.4368	0.075466\\
4.4334	0.076264\\
4.4299	0.077067\\
4.4265	0.077876\\
4.4231	0.07869\\
4.4197	0.079509\\
4.4163	0.080333\\
4.4129	0.081163\\
4.4095	0.081998\\
4.4061	0.082838\\
4.4027	0.083684\\
4.3993	0.084535\\
4.3959	0.085391\\
4.3924	0.086253\\
4.389	0.08712\\
4.3856	0.087992\\
4.3822	0.08887\\
4.3788	0.089754\\
4.3754	0.090642\\
4.372	0.091537\\
4.3686	0.092437\\
4.3652	0.093342\\
4.3618	0.094253\\
4.3584	0.095169\\
4.3549	0.096091\\
4.3515	0.097019\\
4.3481	0.097953\\
4.3447	0.098892\\
4.3413	0.099836\\
4.3379	0.10079\\
4.3345	0.10174\\
4.3311	0.10271\\
4.3277	0.10367\\
4.3243	0.10465\\
4.3209	0.10563\\
4.3174	0.10661\\
4.314	0.10761\\
4.3106	0.1086\\
4.3072	0.10961\\
4.3038	0.11062\\
4.3004	0.11163\\
4.297	0.11266\\
4.2936	0.11369\\
4.2902	0.11472\\
4.2868	0.11576\\
4.2834	0.11681\\
4.2799	0.11787\\
4.2765	0.11893\\
4.2731	0.12\\
4.2697	0.12107\\
4.2663	0.12215\\
4.2629	0.12324\\
4.2595	0.12434\\
4.2561	0.12544\\
4.2527	0.12655\\
4.2493	0.12767\\
4.2459	0.12879\\
4.2424	0.12992\\
4.239	0.13106\\
4.2356	0.13221\\
4.2322	0.13337\\
4.2288	0.13453\\
4.2254	0.1357\\
4.222	0.13688\\
4.2186	0.13806\\
4.2152	0.13926\\
4.2118	0.14046\\
4.2084	0.14167\\
4.2049	0.14289\\
4.2015	0.14412\\
4.1981	0.14535\\
4.1947	0.1466\\
4.1913	0.14785\\
4.1879	0.14911\\
4.1845	0.15039\\
4.1811	0.15167\\
4.1777	0.15296\\
4.1743	0.15426\\
4.1709	0.15557\\
4.1674	0.15689\\
4.164	0.15822\\
4.1606	0.15956\\
4.1572	0.16092\\
4.1538	0.16228\\
4.1504	0.16365\\
4.147	0.16503\\
4.1436	0.16643\\
4.1402	0.16783\\
4.1368	0.16925\\
4.1334	0.17068\\
4.1299	0.17212\\
4.1265	0.17357\\
4.1231	0.17503\\
4.1197	0.17651\\
4.1163	0.178\\
4.1129	0.1795\\
4.1095	0.18101\\
4.1061	0.18254\\
4.1027	0.18408\\
4.0993	0.18564\\
4.0959	0.1872\\
4.0924	0.18879\\
4.089	0.19038\\
4.0856	0.19199\\
4.0822	0.19362\\
4.0788	0.19526\\
4.0754	0.19691\\
4.072	0.19858\\
4.0686	0.20026\\
4.0652	0.20197\\
4.0618	0.20368\\
4.0584	0.20541\\
4.0549	0.20716\\
4.0515	0.20893\\
4.0481	0.21071\\
4.0447	0.21251\\
4.0413	0.21433\\
4.0379	0.21617\\
4.0345	0.21802\\
4.0311	0.21989\\
4.0277	0.22178\\
4.0243	0.22369\\
4.0209	0.22561\\
4.0174	0.22756\\
4.014	0.22953\\
4.0106	0.23151\\
4.0072	0.23352\\
4.0038	0.23554\\
4.0004	0.23759\\
3.997	0.23966\\
3.9936	0.24175\\
3.9902	0.24386\\
3.9868	0.24599\\
3.9834	0.24814\\
3.9799	0.25032\\
3.9765	0.25251\\
3.9731	0.25474\\
3.9697	0.25698\\
3.9663	0.25925\\
3.9629	0.26154\\
3.9595	0.26385\\
3.9561	0.26619\\
3.9527	0.26855\\
3.9493	0.27094\\
3.9459	0.27335\\
3.9424	0.27579\\
3.939	0.27825\\
3.9356	0.28074\\
3.9322	0.28326\\
3.9288	0.2858\\
3.9254	0.28837\\
3.922	0.29096\\
3.9186	0.29358\\
3.9152	0.29623\\
3.9118	0.29891\\
3.9084	0.30161\\
3.9049	0.30434\\
3.9015	0.3071\\
3.8981	0.30989\\
3.8947	0.31271\\
3.8913	0.31556\\
3.8879	0.31843\\
3.8845	0.32134\\
3.8811	0.32427\\
3.8777	0.32724\\
3.8743	0.33024\\
3.8709	0.33326\\
3.8674	0.33632\\
3.864	0.33941\\
3.8606	0.34252\\
3.8572	0.34567\\
3.8538	0.34885\\
3.8504	0.35207\\
3.847	0.35531\\
3.8436	0.35859\\
3.8402	0.3619\\
3.8368	0.36524\\
3.8334	0.36861\\
3.8299	0.37202\\
3.8265	0.37546\\
3.8231	0.37893\\
3.8197	0.38243\\
3.8163	0.38597\\
3.8129	0.38954\\
3.8095	0.39315\\
3.8061	0.39679\\
3.8027	0.40046\\
3.7993	0.40416\\
3.7959	0.4079\\
3.7924	0.41168\\
3.789	0.41549\\
3.7856	0.41933\\
3.7822	0.4232\\
3.7788	0.42711\\
3.7754	0.43106\\
3.772	0.43503\\
3.7686	0.43905\\
3.7652	0.44309\\
3.7618	0.44717\\
3.7584	0.45129\\
3.7549	0.45544\\
3.7515	0.45962\\
3.7481	0.46383\\
3.7447	0.46808\\
3.7413	0.47236\\
3.7379	0.47668\\
3.7345	0.48103\\
3.7311	0.48541\\
3.7277	0.48983\\
3.7243	0.49427\\
3.7209	0.49875\\
3.7174	0.50327\\
3.714	0.50781\\
3.7106	0.51239\\
3.7072	0.51699\\
3.7038	0.52163\\
3.7004	0.5263\\
3.697	0.531\\
3.6936	0.53573\\
3.6902	0.54049\\
3.6868	0.54527\\
3.6834	0.55009\\
3.6799	0.55493\\
3.6765	0.55981\\
3.6731	0.56471\\
3.6697	0.56963\\
3.6663	0.57459\\
3.6629	0.57956\\
3.6595	0.58457\\
3.6561	0.5896\\
3.6527	0.59465\\
3.6493	0.59972\\
3.6459	0.60482\\
3.6424	0.60994\\
3.639	0.61508\\
3.6356	0.62024\\
3.6322	0.62542\\
3.6288	0.63062\\
3.6254	0.63584\\
3.622	0.64108\\
3.6186	0.64633\\
3.6152	0.6516\\
3.6118	0.65688\\
3.6084	0.66217\\
3.6049	0.66748\\
3.6015	0.6728\\
3.5981	0.67814\\
3.5947	0.68348\\
3.5913	0.68883\\
3.5879	0.69419\\
3.5845	0.69955\\
3.5811	0.70492\\
3.5777	0.7103\\
3.5743	0.71568\\
3.5709	0.72106\\
3.5674	0.72645\\
3.564	0.73183\\
3.5606	0.73722\\
3.5572	0.7426\\
3.5538	0.74797\\
3.5504	0.75335\\
3.547	0.75872\\
3.5436	0.76408\\
3.5402	0.76943\\
3.5368	0.77477\\
3.5334	0.7801\\
3.5299	0.78542\\
3.5265	0.79073\\
3.5231	0.79602\\
3.5197	0.8013\\
3.5163	0.80655\\
3.5129	0.81179\\
3.5095	0.81701\\
3.5061	0.82221\\
3.5027	0.82738\\
3.4993	0.83253\\
3.4959	0.83765\\
3.4924	0.84275\\
3.489	0.84782\\
3.4856	0.85286\\
3.4822	0.85786\\
3.4788	0.86284\\
3.4754	0.86778\\
3.472	0.87268\\
3.4686	0.87755\\
3.4652	0.88238\\
3.4618	0.88717\\
3.4584	0.89191\\
3.4549	0.89662\\
3.4515	0.90128\\
3.4481	0.9059\\
3.4447	0.91047\\
3.4413	0.91499\\
3.4379	0.91946\\
3.4345	0.92388\\
3.4311	0.92825\\
3.4277	0.93256\\
3.4243	0.93682\\
3.4209	0.94103\\
3.4174	0.94517\\
3.414	0.94926\\
3.4106	0.95329\\
3.4072	0.95725\\
3.4038	0.96115\\
3.4004	0.96499\\
3.397	0.96876\\
3.3936	0.97246\\
3.3902	0.9761\\
3.3868	0.97966\\
3.3834	0.98316\\
3.3799	0.98658\\
3.3765	0.98992\\
3.3731	0.99319\\
3.3697	0.99639\\
3.3663	0.9995\\
3.3629	1.0025\\
3.3595	1.0055\\
3.3561	1.0084\\
3.3527	1.0111\\
3.3493	1.0138\\
3.3459	1.0164\\
3.3424	1.019\\
3.339	1.0214\\
3.3356	1.0237\\
3.3322	1.026\\
3.3288	1.0281\\
3.3254	1.0302\\
3.322	1.0321\\
3.3186	1.034\\
3.3152	1.0357\\
3.3118	1.0373\\
3.3084	1.0389\\
3.3049	1.0403\\
3.3015	1.0416\\
3.2981	1.0428\\
3.2947	1.044\\
3.2913	1.045\\
3.2879	1.0458\\
3.2845	1.0466\\
3.2811	1.0473\\
3.2777	1.0478\\
3.2743	1.0483\\
3.2709	1.0486\\
3.2674	1.0488\\
3.264	1.0489\\
3.2606	1.0489\\
3.2572	1.0487\\
3.2538	1.0485\\
3.2504	1.0481\\
3.247	1.0477\\
3.2436	1.0471\\
3.2402	1.0464\\
3.2368	1.0456\\
3.2334	1.0447\\
3.2299	1.0436\\
3.2265	1.0425\\
3.2231	1.0412\\
3.2197	1.0399\\
3.2163	1.0384\\
3.2129	1.0368\\
3.2095	1.0351\\
3.2061	1.0333\\
3.2027	1.0314\\
3.1993	1.0294\\
3.1959	1.0273\\
3.1924	1.0251\\
3.189	1.0228\\
3.1856	1.0203\\
3.1822	1.0178\\
3.1788	1.0152\\
3.1754	1.0125\\
3.172	1.0096\\
3.1686	1.0067\\
3.1652	1.0037\\
3.1618	1.0006\\
3.1584	0.99736\\
3.1549	0.99405\\
3.1515	0.99065\\
3.1481	0.98715\\
3.1447	0.98356\\
3.1413	0.97987\\
3.1379	0.9761\\
3.1345	0.97224\\
3.1311	0.96828\\
3.1277	0.96424\\
3.1243	0.96011\\
3.1209	0.95589\\
3.1174	0.95158\\
3.114	0.94719\\
3.1106	0.94271\\
3.1072	0.93815\\
3.1038	0.9335\\
3.1004	0.92878\\
3.097	0.92396\\
3.0936	0.91907\\
3.0902	0.91409\\
3.0868	0.90904\\
3.0834	0.90391\\
3.0799	0.89869\\
3.0765	0.8934\\
3.0731	0.88803\\
3.0697	0.88259\\
3.0663	0.87707\\
3.0629	0.87148\\
3.0595	0.86582\\
3.0561	0.86008\\
3.0527	0.85428\\
3.0493	0.84842\\
3.0459	0.84249\\
3.0424	0.8365\\
3.039	0.83045\\
3.0356	0.82434\\
3.0322	0.81818\\
3.0288	0.81198\\
3.0254	0.80572\\
3.022	0.79942\\
3.0186	0.79308\\
3.0152	0.7867\\
3.0118	0.78029\\
3.0084	0.77384\\
3.0049	0.76736\\
3.0015	0.76085\\
2.9981	0.75432\\
2.9947	0.74777\\
2.9913	0.74119\\
2.9879	0.7346\\
2.9845	0.72799\\
2.9811	0.72137\\
2.9777	0.71473\\
2.9743	0.70808\\
2.9709	0.70143\\
2.9674	0.69476\\
2.964	0.6881\\
2.9606	0.68143\\
2.9572	0.67475\\
2.9538	0.66808\\
2.9504	0.66141\\
2.947	0.65474\\
2.9436	0.64807\\
2.9402	0.64141\\
2.9368	0.63475\\
2.9334	0.6281\\
2.9299	0.62146\\
2.9265	0.61483\\
2.9231	0.60822\\
2.9197	0.60161\\
2.9163	0.59502\\
2.9129	0.58844\\
2.9095	0.58188\\
2.9061	0.57533\\
2.9027	0.56881\\
2.8993	0.5623\\
2.8959	0.55581\\
2.8924	0.54934\\
2.889	0.54289\\
2.8856	0.53647\\
2.8822	0.53007\\
2.8788	0.52369\\
2.8754	0.51735\\
2.872	0.51102\\
2.8686	0.50473\\
2.8652	0.49846\\
2.8618	0.49222\\
2.8584	0.48601\\
2.8549	0.47983\\
2.8515	0.47369\\
2.8481	0.46757\\
2.8447	0.46149\\
2.8413	0.45544\\
2.8379	0.44943\\
2.8345	0.44345\\
2.8311	0.43751\\
2.8277	0.43161\\
2.8243	0.42574\\
2.8209	0.41991\\
2.8174	0.41412\\
2.814	0.40837\\
2.8106	0.40266\\
2.8072	0.39699\\
2.8038	0.39136\\
2.8004	0.38577\\
2.797	0.38023\\
2.7936	0.37473\\
2.7902	0.36927\\
2.7868	0.36385\\
2.7834	0.35848\\
2.7799	0.35315\\
2.7765	0.34787\\
2.7731	0.34264\\
2.7697	0.33744\\
2.7663	0.3323\\
2.7629	0.3272\\
2.7595	0.32215\\
2.7561	0.31715\\
2.7527	0.31219\\
2.7493	0.30728\\
2.7459	0.30242\\
2.7424	0.29761\\
2.739	0.29284\\
2.7356	0.28813\\
2.7322	0.28346\\
2.7288	0.27884\\
2.7254	0.27427\\
2.722	0.26975\\
2.7186	0.26528\\
2.7152	0.26086\\
2.7118	0.25648\\
2.7084	0.25216\\
2.7049	0.24788\\
2.7015	0.24366\\
2.6981	0.23948\\
2.6947	0.23536\\
2.6913	0.23128\\
2.6879	0.22725\\
2.6845	0.22328\\
2.6811	0.21935\\
2.6777	0.21547\\
2.6743	0.21164\\
2.6709	0.20785\\
2.6674	0.20412\\
2.664	0.20044\\
2.6606	0.1968\\
2.6572	0.19321\\
2.6538	0.18967\\
2.6504	0.18618\\
2.647	0.18273\\
2.6436	0.17934\\
2.6402	0.17599\\
2.6368	0.17268\\
2.6334	0.16943\\
2.6299	0.16622\\
2.6265	0.16305\\
2.6231	0.15994\\
2.6197	0.15686\\
2.6163	0.15384\\
2.6129	0.15086\\
2.6095	0.14792\\
2.6061	0.14503\\
2.6027	0.14218\\
2.5993	0.13937\\
2.5959	0.13661\\
2.5924	0.13389\\
2.589	0.13122\\
2.5856	0.12858\\
2.5822	0.12599\\
2.5788	0.12344\\
2.5754	0.12093\\
2.572	0.11846\\
2.5686	0.11604\\
2.5652	0.11365\\
2.5618	0.1113\\
2.5584	0.10899\\
2.5549	0.10672\\
2.5515	0.10449\\
2.5481	0.1023\\
2.5447	0.10014\\
2.5413	0.098025\\
2.5379	0.095944\\
2.5345	0.093899\\
2.5311	0.091891\\
2.5277	0.089917\\
2.5243	0.08798\\
2.5209	0.086076\\
2.5174	0.084208\\
2.514	0.082373\\
2.5106	0.080572\\
2.5072	0.078803\\
2.5038	0.077068\\
2.5004	0.075365\\
2.497	0.073693\\
2.4936	0.072053\\
2.4902	0.070444\\
2.4868	0.068866\\
2.4834	0.067317\\
2.4799	0.065799\\
2.4765	0.06431\\
2.4731	0.06285\\
2.4697	0.061418\\
2.4663	0.060014\\
2.4629	0.058638\\
2.4595	0.05729\\
2.4561	0.055968\\
2.4527	0.054673\\
2.4493	0.053403\\
2.4459	0.05216\\
2.4424	0.050941\\
2.439	0.049748\\
2.4356	0.048579\\
2.4322	0.047434\\
2.4288	0.046313\\
2.4254	0.045215\\
2.422	0.044141\\
2.4186	0.043088\\
2.4152	0.042058\\
2.4118	0.04105\\
2.4084	0.040063\\
2.4049	0.039098\\
2.4015	0.038153\\
2.3981	0.037228\\
2.3947	0.036324\\
2.3913	0.035439\\
2.3879	0.034573\\
2.3845	0.033727\\
2.3811	0.032899\\
2.3777	0.032089\\
2.3743	0.031298\\
2.3709	0.030524\\
2.3674	0.029767\\
2.364	0.029028\\
2.3606	0.028305\\
2.3572	0.027598\\
2.3538	0.026908\\
2.3504	0.026233\\
2.347	0.025573\\
2.3436	0.024929\\
2.3402	0.0243\\
2.3368	0.023685\\
2.3334	0.023085\\
2.3299	0.022498\\
2.3265	0.021926\\
2.3231	0.021366\\
2.3197	0.02082\\
2.3163	0.020287\\
2.3129	0.019767\\
2.3095	0.019258\\
2.3061	0.018762\\
2.3027	0.018278\\
2.2993	0.017806\\
2.2959	0.017345\\
2.2924	0.016895\\
2.289	0.016456\\
2.2856	0.016027\\
2.2822	0.015609\\
2.2788	0.015202\\
2.2754	0.014804\\
2.272	0.014416\\
2.2686	0.014038\\
2.2652	0.013669\\
2.2618	0.013309\\
2.2584	0.012958\\
2.2549	0.012616\\
2.2515	0.012283\\
2.2481	0.011957\\
2.2447	0.01164\\
2.2413	0.011331\\
2.2379	0.01103\\
2.2345	0.010737\\
2.2311	0.010451\\
2.2277	0.010172\\
2.2243	0.0099003\\
2.2209	0.0096356\\
2.2174	0.0093777\\
2.214	0.0091264\\
2.2106	0.0088816\\
2.2072	0.0086431\\
2.2038	0.0084109\\
2.2004	0.0081846\\
2.197	0.0079642\\
2.1936	0.0077496\\
2.1902	0.0075405\\
2.1868	0.007337\\
2.1834	0.0071387\\
2.1799	0.0069457\\
2.1765	0.0067577\\
2.1731	0.0065747\\
2.1697	0.0063966\\
2.1663	0.0062231\\
2.1629	0.0060543\\
2.1595	0.0058899\\
2.1561	0.0057299\\
2.1527	0.0055742\\
2.1493	0.0054226\\
2.1459	0.005275\\
2.1424	0.0051315\\
2.139	0.0049917\\
2.1356	0.0048557\\
2.1322	0.0047234\\
2.1288	0.0045946\\
2.1254	0.0044693\\
2.122	0.0043474\\
2.1186	0.0042288\\
2.1152	0.0041134\\
2.1118	0.0040011\\
2.1084	0.0038918\\
2.1049	0.0037856\\
2.1015	0.0036822\\
2.0981	0.0035816\\
2.0947	0.0034837\\
2.0913	0.0033886\\
2.0879	0.003296\\
2.0845	0.0032059\\
2.0811	0.0031184\\
2.0777	0.0030332\\
2.0743	0.0029503\\
2.0709	0.0028697\\
2.0674	0.0027914\\
2.064	0.0027151\\
2.0606	0.002641\\
2.0572	0.0025689\\
2.0538	0.0024988\\
2.0504	0.0024306\\
2.047	0.0023643\\
2.0436	0.0022998\\
2.0402	0.0022371\\
2.0368	0.0021761\\
2.0334	0.0021168\\
2.0299	0.0020592\\
2.0265	0.0020031\\
2.0231	0.0019486\\
2.0197	0.0018955\\
2.0163	0.001844\\
2.0129	0.0017939\\
2.0095	0.0017451\\
2.0061	0.0016977\\
2.0027	0.0016516\\
1.9993	0.0016068\\
1.9959	0.0015632\\
1.9924	0.0015208\\
1.989	0.0014796\\
1.9856	0.0014396\\
1.9822	0.0014006\\
1.9788	0.0013627\\
1.9754	0.0013259\\
1.972	0.0012901\\
1.9686	0.0012552\\
1.9652	0.0012214\\
1.9618	0.0011884\\
1.9584	0.0011564\\
1.9549	0.0011252\\
1.9515	0.001095\\
1.9481	0.0010655\\
1.9447	0.0010369\\
1.9413	0.001009\\
1.9379	0.00098192\\
1.9345	0.00095558\\
1.9311	0.00092996\\
1.9277	0.00090505\\
1.9243	0.00088081\\
1.9209	0.00085725\\
1.9174	0.00083433\\
1.914	0.00081204\\
1.9106	0.00079035\\
1.9072	0.00076927\\
1.9038	0.00074876\\
1.9004	0.0007288\\
1.897	0.0007094\\
1.8936	0.00069052\\
1.8902	0.00067216\\
1.8868	0.0006543\\
}--cycle;
\addplot[ybar interval, fill=mycolor2, fill opacity=0.35, area legend, draw=none] table[row sep=crcr, x=Lower, y=Count] {%
Lower	Upper	Count\\
1.8868	2.1487	0.0011197\\
2.1487	2.4107	0.017508\\
2.4107	2.6727	0.11991\\
2.6727	2.9347	0.40309\\
2.9347	3.1966	0.858\\
3.1966	3.4586	1.0209\\
3.4586	3.7206	0.70958\\
3.7206	3.9826	0.35189\\
3.9826	4.2445	0.19147\\
4.2445	4.5065	0.093342\\
4.5065	4.7685	0.031759\\
4.7685	5.0305	0.0084486\\
5.0305	5.2924	0.010179\\
5.2924	5.2924	0.010179\\
};
\addlegendentry{\ApproxHistogramLegend}

\addplot [color=mycolor2, only marks, every error bar/.append style={opacity=0.45}, mark=*, mark size=0pt, draw=none, forget plot]
 plot [error bars/.cd, y dir=both, y explicit, error bar style={line width=1pt, color=mycolor2}, error mark options={mark=none,mark size=0pt}]
 table[row sep=crcr, y error plus index=2, y error minus index=3]{%
2.0178	0.0011197	0.00063514	0.00063514\\
2.2797	0.017508	0.0025204	0.0025204\\
2.5417	0.11991	0.005164	0.005164\\
2.8037	0.40309	0.0098739	0.0098739\\
3.0657	0.858	0.014509	0.014509\\
3.3276	1.0209	0.015286	0.015286\\
3.5896	0.70958	0.011907	0.011907\\
3.8516	0.35189	0.0083462	0.0083462\\
4.1136	0.19147	0.005923	0.005923\\
4.3755	0.093342	0.0036882	0.0036882\\
4.6375	0.031759	0.0021344	0.0021344\\
4.8995	0.0084486	0.0014653	0.0014653\\
5.1615	0.010179	2.0756e-18	1.7347e-18\\
};
\addplot [color=mycolor2]
  table[row sep=crcr]{%
1.8868	0.00086104\\
2.0061	0.0020617\\
2.0845	0.0037066\\
2.1459	0.0058962\\
2.1936	0.008468\\
2.2345	0.011538\\
2.272	0.015292\\
2.3027	0.019215\\
2.3334	0.024081\\
2.3606	0.029357\\
2.3845	0.034831\\
2.4084	0.041227\\
2.4322	0.048667\\
2.4527	0.055976\\
2.4731	0.064238\\
2.4936	0.073544\\
2.514	0.083987\\
2.5311	0.093628\\
2.5481	0.10418\\
2.5652	0.1157\\
2.5822	0.12823\\
2.5993	0.14184\\
2.6163	0.15656\\
2.6334	0.17243\\
2.6504	0.18951\\
2.6674	0.20781\\
2.6845	0.22736\\
2.7015	0.24817\\
2.7186	0.27026\\
2.739	0.29844\\
2.7595	0.32841\\
2.7799	0.36011\\
2.8004	0.39347\\
2.8243	0.43431\\
2.8481	0.477\\
2.8754	0.52766\\
2.9095	0.59301\\
2.9981	0.76425\\
3.0254	0.81426\\
3.0459	0.85001\\
3.0663	0.88387\\
3.0834	0.91037\\
3.1004	0.93513\\
3.114	0.95355\\
3.1277	0.97066\\
3.1413	0.98635\\
3.1549	1.0006\\
3.1686	1.0133\\
3.1788	1.0217\\
3.189	1.0293\\
3.1993	1.0359\\
3.2095	1.0416\\
3.2197	1.0463\\
3.2299	1.05\\
3.2402	1.0528\\
3.2504	1.0546\\
3.2606	1.0554\\
3.2709	1.0553\\
3.2811	1.0542\\
3.2913	1.0521\\
3.3015	1.0491\\
3.3118	1.0452\\
3.322	1.0404\\
3.3322	1.0346\\
3.3424	1.0281\\
3.3561	1.0181\\
3.3697	1.0066\\
3.3834	0.99385\\
3.397	0.97983\\
3.4106	0.96465\\
3.4277	0.94417\\
3.4447	0.9222\\
3.4652	0.89412\\
3.4856	0.86447\\
3.5095	0.82834\\
3.5402	0.78022\\
3.6493	0.60748\\
3.6765	0.56665\\
3.7038	0.52755\\
3.7277	0.49496\\
3.7515	0.46399\\
3.772	0.4388\\
3.7924	0.41489\\
3.8129	0.39225\\
3.8334	0.37088\\
3.8538	0.35074\\
3.8743	0.33181\\
3.8947	0.31402\\
3.9152	0.29734\\
3.9356	0.28169\\
3.9595	0.26465\\
3.9834	0.24884\\
4.0072	0.23414\\
4.0311	0.22045\\
4.0584	0.20593\\
4.0856	0.19246\\
4.1163	0.17843\\
4.147	0.16542\\
4.1811	0.15202\\
4.2152	0.13957\\
4.2527	0.12684\\
4.2936	0.11397\\
4.3345	0.10204\\
4.3788	0.090086\\
4.4231	0.079093\\
4.4674	0.069014\\
4.5118	0.059823\\
4.5595	0.05089\\
4.6072	0.042927\\
4.6549	0.035898\\
4.7027	0.029755\\
4.7538	0.024093\\
4.8083	0.019017\\
4.8663	0.014597\\
4.9277	0.01087\\
4.9958	0.0076962\\
5.0708	0.0051518\\
5.1595	0.0031145\\
5.2652	0.0016412\\
5.2924	0.0013813\\
};
\addlegendentry{\ApproxBLNFitLegend}


\addplot[area legend, draw=none, fill=mycolor2, fill opacity=0.15, forget plot]
table[row sep=crcr] {%
x	y\\
1.8868	0.001057\\
1.8902	0.0010821\\
1.8936	0.0011077\\
1.897	0.001134\\
1.9004	0.0011609\\
1.9038	0.0011884\\
1.9072	0.0012167\\
1.9106	0.0012456\\
1.914	0.0012752\\
1.9174	0.0013055\\
1.9209	0.0013365\\
1.9243	0.0013684\\
1.9277	0.0014009\\
1.9311	0.0014343\\
1.9345	0.0014685\\
1.9379	0.0015035\\
1.9413	0.0015394\\
1.9447	0.0015762\\
1.9481	0.0016138\\
1.9515	0.0016524\\
1.9549	0.0016919\\
1.9584	0.0017325\\
1.9618	0.001774\\
1.9652	0.0018165\\
1.9686	0.00186\\
1.972	0.0019047\\
1.9754	0.0019504\\
1.9788	0.0019973\\
1.9822	0.0020454\\
1.9856	0.0020946\\
1.989	0.0021451\\
1.9924	0.0021968\\
1.9959	0.0022498\\
1.9993	0.0023041\\
2.0027	0.0023597\\
2.0061	0.0024168\\
2.0095	0.0024753\\
2.0129	0.0025352\\
2.0163	0.0025967\\
2.0197	0.0026597\\
2.0231	0.0027242\\
2.0265	0.0027904\\
2.0299	0.0028582\\
2.0334	0.0029278\\
2.0368	0.0029991\\
2.0402	0.0030722\\
2.0436	0.0031471\\
2.047	0.0032239\\
2.0504	0.0033027\\
2.0538	0.0033834\\
2.0572	0.0034662\\
2.0606	0.0035511\\
2.064	0.0036381\\
2.0674	0.0037274\\
2.0709	0.0038188\\
2.0743	0.0039127\\
2.0777	0.0040088\\
2.0811	0.0041075\\
2.0845	0.0042086\\
2.0879	0.0043123\\
2.0913	0.0044186\\
2.0947	0.0045276\\
2.0981	0.0046394\\
2.1015	0.004754\\
2.1049	0.0048716\\
2.1084	0.0049921\\
2.1118	0.0051157\\
2.1152	0.0052424\\
2.1186	0.0053724\\
2.122	0.0055057\\
2.1254	0.0056423\\
2.1288	0.0057825\\
2.1322	0.0059262\\
2.1356	0.0060735\\
2.139	0.0062246\\
2.1424	0.0063796\\
2.1459	0.0065384\\
2.1493	0.0067014\\
2.1527	0.0068684\\
2.1561	0.0070397\\
2.1595	0.0072153\\
2.1629	0.0073955\\
2.1663	0.0075801\\
2.1697	0.0077695\\
2.1731	0.0079636\\
2.1765	0.0081627\\
2.1799	0.0083668\\
2.1834	0.0085761\\
2.1868	0.0087906\\
2.1902	0.0090106\\
2.1936	0.0092361\\
2.197	0.0094674\\
2.2004	0.0097044\\
2.2038	0.0099474\\
2.2072	0.010197\\
2.2106	0.010452\\
2.214	0.010714\\
2.2174	0.010982\\
2.2209	0.011257\\
2.2243	0.011539\\
2.2277	0.011828\\
2.2311	0.012124\\
2.2345	0.012428\\
2.2379	0.012739\\
2.2413	0.013058\\
2.2447	0.013385\\
2.2481	0.01372\\
2.2515	0.014063\\
2.2549	0.014414\\
2.2584	0.014775\\
2.2618	0.015144\\
2.2652	0.015522\\
2.2686	0.01591\\
2.272	0.016307\\
2.2754	0.016713\\
2.2788	0.01713\\
2.2822	0.017557\\
2.2856	0.017994\\
2.289	0.018442\\
2.2924	0.0189\\
2.2959	0.01937\\
2.2993	0.019851\\
2.3027	0.020344\\
2.3061	0.020848\\
2.3095	0.021364\\
2.3129	0.021893\\
2.3163	0.022434\\
2.3197	0.022989\\
2.3231	0.023556\\
2.3265	0.024136\\
2.3299	0.024731\\
2.3334	0.025339\\
2.3368	0.025962\\
2.3402	0.026599\\
2.3436	0.027251\\
2.347	0.027918\\
2.3504	0.0286\\
2.3538	0.029298\\
2.3572	0.030013\\
2.3606	0.030743\\
2.364	0.031491\\
2.3674	0.032255\\
2.3709	0.033037\\
2.3743	0.033836\\
2.3777	0.034654\\
2.3811	0.03549\\
2.3845	0.036344\\
2.3879	0.037218\\
2.3913	0.038111\\
2.3947	0.039024\\
2.3981	0.039957\\
2.4015	0.04091\\
2.4049	0.041885\\
2.4084	0.04288\\
2.4118	0.043898\\
2.4152	0.044937\\
2.4186	0.045998\\
2.422	0.047083\\
2.4254	0.04819\\
2.4288	0.049321\\
2.4322	0.050476\\
2.4356	0.051656\\
2.439	0.05286\\
2.4424	0.054089\\
2.4459	0.055344\\
2.4493	0.056625\\
2.4527	0.057932\\
2.4561	0.059266\\
2.4595	0.060628\\
2.4629	0.062016\\
2.4663	0.063433\\
2.4697	0.064879\\
2.4731	0.066353\\
2.4765	0.067857\\
2.4799	0.06939\\
2.4834	0.070954\\
2.4868	0.072549\\
2.4902	0.074174\\
2.4936	0.075831\\
2.497	0.07752\\
2.5004	0.079241\\
2.5038	0.080995\\
2.5072	0.082782\\
2.5106	0.084603\\
2.514	0.086458\\
2.5174	0.088348\\
2.5209	0.090272\\
2.5243	0.092232\\
2.5277	0.094228\\
2.5311	0.09626\\
2.5345	0.098329\\
2.5379	0.10043\\
2.5413	0.10258\\
2.5447	0.10476\\
2.5481	0.10698\\
2.5515	0.10924\\
2.5549	0.11154\\
2.5584	0.11387\\
2.5618	0.11625\\
2.5652	0.11867\\
2.5686	0.12113\\
2.572	0.12363\\
2.5754	0.12617\\
2.5788	0.12875\\
2.5822	0.13138\\
2.5856	0.13405\\
2.589	0.13676\\
2.5924	0.13952\\
2.5959	0.14232\\
2.5993	0.14516\\
2.6027	0.14805\\
2.6061	0.15098\\
2.6095	0.15396\\
2.6129	0.15699\\
2.6163	0.16006\\
2.6197	0.16317\\
2.6231	0.16634\\
2.6265	0.16955\\
2.6299	0.1728\\
2.6334	0.17611\\
2.6368	0.17946\\
2.6402	0.18286\\
2.6436	0.18631\\
2.647	0.18981\\
2.6504	0.19335\\
2.6538	0.19694\\
2.6572	0.20059\\
2.6606	0.20428\\
2.664	0.20802\\
2.6674	0.21181\\
2.6709	0.21565\\
2.6743	0.21954\\
2.6777	0.22348\\
2.6811	0.22747\\
2.6845	0.23151\\
2.6879	0.2356\\
2.6913	0.23974\\
2.6947	0.24393\\
2.6981	0.24817\\
2.7015	0.25246\\
2.7049	0.25681\\
2.7084	0.2612\\
2.7118	0.26564\\
2.7152	0.27013\\
2.7186	0.27467\\
2.722	0.27926\\
2.7254	0.28391\\
2.7288	0.2886\\
2.7322	0.29334\\
2.7356	0.29813\\
2.739	0.30297\\
2.7424	0.30786\\
2.7459	0.3128\\
2.7493	0.31778\\
2.7527	0.32282\\
2.7561	0.3279\\
2.7595	0.33303\\
2.7629	0.33821\\
2.7663	0.34343\\
2.7697	0.34871\\
2.7731	0.35402\\
2.7765	0.35939\\
2.7799	0.3648\\
2.7834	0.37025\\
2.7868	0.37575\\
2.7902	0.38129\\
2.7936	0.38688\\
2.797	0.39251\\
2.8004	0.39818\\
2.8038	0.40389\\
2.8072	0.40965\\
2.8106	0.41544\\
2.814	0.42128\\
2.8174	0.42715\\
2.8209	0.43307\\
2.8243	0.43902\\
2.8277	0.44501\\
2.8311	0.45103\\
2.8345	0.45709\\
2.8379	0.46319\\
2.8413	0.46932\\
2.8447	0.47548\\
2.8481	0.48168\\
2.8515	0.48791\\
2.8549	0.49416\\
2.8584	0.50045\\
2.8618	0.50677\\
2.8652	0.51311\\
2.8686	0.51948\\
2.872	0.52588\\
2.8754	0.5323\\
2.8788	0.53874\\
2.8822	0.54521\\
2.8856	0.5517\\
2.889	0.55821\\
2.8924	0.56474\\
2.8959	0.57129\\
2.8993	0.57785\\
2.9027	0.58443\\
2.9061	0.59102\\
2.9095	0.59763\\
2.9129	0.60425\\
2.9163	0.61088\\
2.9197	0.61752\\
2.9231	0.62417\\
2.9265	0.63083\\
2.9299	0.63749\\
2.9334	0.64415\\
2.9368	0.65082\\
2.9402	0.65749\\
2.9436	0.66416\\
2.947	0.67082\\
2.9504	0.67749\\
2.9538	0.68415\\
2.9572	0.6908\\
2.9606	0.69744\\
2.964	0.70408\\
2.9674	0.7107\\
2.9709	0.71732\\
2.9743	0.72391\\
2.9777	0.73049\\
2.9811	0.73706\\
2.9845	0.7436\\
2.9879	0.75013\\
2.9913	0.75663\\
2.9947	0.76311\\
2.9981	0.76956\\
3.0015	0.77598\\
3.0049	0.78237\\
3.0084	0.78874\\
3.0118	0.79507\\
3.0152	0.80136\\
3.0186	0.80762\\
3.022	0.81384\\
3.0254	0.82002\\
3.0288	0.82615\\
3.0322	0.83225\\
3.0356	0.8383\\
3.039	0.8443\\
3.0424	0.85025\\
3.0459	0.85615\\
3.0493	0.862\\
3.0527	0.86779\\
3.0561	0.87353\\
3.0595	0.87921\\
3.0629	0.88483\\
3.0663	0.89039\\
3.0697	0.89589\\
3.0731	0.90132\\
3.0765	0.90668\\
3.0799	0.91198\\
3.0834	0.91721\\
3.0868	0.92237\\
3.0902	0.92745\\
3.0936	0.93246\\
3.097	0.9374\\
3.1004	0.94226\\
3.1038	0.94704\\
3.1072	0.95174\\
3.1106	0.95635\\
3.114	0.96089\\
3.1174	0.96534\\
3.1209	0.96971\\
3.1243	0.97398\\
3.1277	0.97817\\
3.1311	0.98227\\
3.1345	0.98628\\
3.1379	0.9902\\
3.1413	0.99402\\
3.1447	0.99775\\
3.1481	1.0014\\
3.1515	1.0049\\
3.1549	1.0084\\
3.1584	1.0117\\
3.1618	1.0149\\
3.1652	1.0181\\
3.1686	1.0211\\
3.172	1.0241\\
3.1754	1.0269\\
3.1788	1.0296\\
3.1822	1.0323\\
3.1856	1.0348\\
3.189	1.0372\\
3.1924	1.0395\\
3.1959	1.0417\\
3.1993	1.0438\\
3.2027	1.0458\\
3.2061	1.0477\\
3.2095	1.0495\\
3.2129	1.0511\\
3.2163	1.0527\\
3.2197	1.0541\\
3.2231	1.0554\\
3.2265	1.0567\\
3.2299	1.0578\\
3.2334	1.0588\\
3.2368	1.0597\\
3.2402	1.0605\\
3.2436	1.0611\\
3.247	1.0617\\
3.2504	1.0622\\
3.2538	1.0625\\
3.2572	1.0627\\
3.2606	1.0629\\
3.264	1.0629\\
3.2674	1.0628\\
3.2709	1.0626\\
3.2743	1.0623\\
3.2777	1.0618\\
3.2811	1.0613\\
3.2845	1.0607\\
3.2879	1.06\\
3.2913	1.0591\\
3.2947	1.0582\\
3.2981	1.0571\\
3.3015	1.056\\
3.3049	1.0547\\
3.3084	1.0534\\
3.3118	1.0519\\
3.3152	1.0504\\
3.3186	1.0487\\
3.322	1.0469\\
3.3254	1.0451\\
3.3288	1.0432\\
3.3322	1.0411\\
3.3356	1.039\\
3.339	1.0368\\
3.3424	1.0344\\
3.3459	1.032\\
3.3493	1.0295\\
3.3527	1.0269\\
3.3561	1.0243\\
3.3595	1.0215\\
3.3629	1.0187\\
3.3663	1.0158\\
3.3697	1.0127\\
3.3731	1.0097\\
3.3765	1.0065\\
3.3799	1.0032\\
3.3834	0.99992\\
3.3868	0.99652\\
3.3902	0.99305\\
3.3936	0.9895\\
3.397	0.98588\\
3.4004	0.9822\\
3.4038	0.97844\\
3.4072	0.97462\\
3.4106	0.97073\\
3.414	0.96677\\
3.4174	0.96275\\
3.4209	0.95867\\
3.4243	0.95453\\
3.4277	0.95033\\
3.4311	0.94607\\
3.4345	0.94175\\
3.4379	0.93738\\
3.4413	0.93295\\
3.4447	0.92847\\
3.4481	0.92394\\
3.4515	0.91936\\
3.4549	0.91473\\
3.4584	0.91006\\
3.4618	0.90533\\
3.4652	0.90057\\
3.4686	0.89576\\
3.472	0.89091\\
3.4754	0.88602\\
3.4788	0.88109\\
3.4822	0.87612\\
3.4856	0.87111\\
3.489	0.86608\\
3.4924	0.86101\\
3.4959	0.8559\\
3.4993	0.85077\\
3.5027	0.8456\\
3.5061	0.84041\\
3.5095	0.8352\\
3.5129	0.82995\\
3.5163	0.82469\\
3.5197	0.8194\\
3.5231	0.81409\\
3.5265	0.80876\\
3.5299	0.80341\\
3.5334	0.79805\\
3.5368	0.79267\\
3.5402	0.78728\\
3.5436	0.78187\\
3.547	0.77645\\
3.5504	0.77102\\
3.5538	0.76558\\
3.5572	0.76013\\
3.5606	0.75468\\
3.564	0.74922\\
3.5674	0.74376\\
3.5709	0.73829\\
3.5743	0.73282\\
3.5777	0.72735\\
3.5811	0.72188\\
3.5845	0.71641\\
3.5879	0.71095\\
3.5913	0.70548\\
3.5947	0.70003\\
3.5981	0.69458\\
3.6015	0.68913\\
3.6049	0.6837\\
3.6084	0.67827\\
3.6118	0.67285\\
3.6152	0.66744\\
3.6186	0.66205\\
3.622	0.65667\\
3.6254	0.6513\\
3.6288	0.64595\\
3.6322	0.64061\\
3.6356	0.63529\\
3.639	0.62998\\
3.6424	0.6247\\
3.6459	0.61943\\
3.6493	0.61418\\
3.6527	0.60895\\
3.6561	0.60375\\
3.6595	0.59857\\
3.6629	0.5934\\
3.6663	0.58827\\
3.6697	0.58315\\
3.6731	0.57806\\
3.6765	0.573\\
3.6799	0.56796\\
3.6834	0.56295\\
3.6868	0.55796\\
3.6902	0.55301\\
3.6936	0.54808\\
3.697	0.54318\\
3.7004	0.53831\\
3.7038	0.53347\\
3.7072	0.52866\\
3.7106	0.52388\\
3.714	0.51913\\
3.7174	0.51441\\
3.7209	0.50972\\
3.7243	0.50507\\
3.7277	0.50045\\
3.7311	0.49586\\
3.7345	0.4913\\
3.7379	0.48678\\
3.7413	0.48229\\
3.7447	0.47784\\
3.7481	0.47342\\
3.7515	0.46903\\
3.7549	0.46468\\
3.7584	0.46037\\
3.7618	0.45609\\
3.7652	0.45184\\
3.7686	0.44763\\
3.772	0.44345\\
3.7754	0.43932\\
3.7788	0.43521\\
3.7822	0.43115\\
3.7856	0.42711\\
3.789	0.42312\\
3.7924	0.41916\\
3.7959	0.41524\\
3.7993	0.41135\\
3.8027	0.4075\\
3.8061	0.40368\\
3.8095	0.3999\\
3.8129	0.39616\\
3.8163	0.39246\\
3.8197	0.38879\\
3.8231	0.38515\\
3.8265	0.38155\\
3.8299	0.37799\\
3.8334	0.37446\\
3.8368	0.37097\\
3.8402	0.36752\\
3.8436	0.3641\\
3.847	0.36071\\
3.8504	0.35736\\
3.8538	0.35405\\
3.8572	0.35077\\
3.8606	0.34752\\
3.864	0.34431\\
3.8674	0.34113\\
3.8709	0.33799\\
3.8743	0.33488\\
3.8777	0.33181\\
3.8811	0.32877\\
3.8845	0.32576\\
3.8879	0.32278\\
3.8913	0.31984\\
3.8947	0.31693\\
3.8981	0.31405\\
3.9015	0.3112\\
3.9049	0.30839\\
3.9084	0.3056\\
3.9118	0.30285\\
3.9152	0.30012\\
3.9186	0.29743\\
3.922	0.29477\\
3.9254	0.29213\\
3.9288	0.28953\\
3.9322	0.28695\\
3.9356	0.28441\\
3.939	0.28189\\
3.9424	0.2794\\
3.9459	0.27693\\
3.9493	0.2745\\
3.9527	0.27209\\
3.9561	0.2697\\
3.9595	0.26735\\
3.9629	0.26501\\
3.9663	0.26271\\
3.9697	0.26042\\
3.9731	0.25817\\
3.9765	0.25593\\
3.9799	0.25372\\
3.9834	0.25154\\
3.9868	0.24937\\
3.9902	0.24723\\
3.9936	0.24511\\
3.997	0.24302\\
4.0004	0.24094\\
4.0038	0.23889\\
4.0072	0.23686\\
4.0106	0.23484\\
4.014	0.23285\\
4.0174	0.23088\\
4.0209	0.22893\\
4.0243	0.227\\
4.0277	0.22508\\
4.0311	0.22319\\
4.0345	0.22131\\
4.0379	0.21945\\
4.0413	0.21761\\
4.0447	0.21579\\
4.0481	0.21399\\
4.0515	0.2122\\
4.0549	0.21043\\
4.0584	0.20867\\
4.0618	0.20693\\
4.0652	0.20521\\
4.0686	0.2035\\
4.072	0.20181\\
4.0754	0.20013\\
4.0788	0.19847\\
4.0822	0.19682\\
4.0856	0.19519\\
4.089	0.19357\\
4.0924	0.19196\\
4.0959	0.19037\\
4.0993	0.18879\\
4.1027	0.18723\\
4.1061	0.18567\\
4.1095	0.18414\\
4.1129	0.18261\\
4.1163	0.1811\\
4.1197	0.17959\\
4.1231	0.1781\\
4.1265	0.17663\\
4.1299	0.17516\\
4.1334	0.1737\\
4.1368	0.17226\\
4.1402	0.17083\\
4.1436	0.16941\\
4.147	0.168\\
4.1504	0.1666\\
4.1538	0.16521\\
4.1572	0.16383\\
4.1606	0.16246\\
4.164	0.1611\\
4.1674	0.15975\\
4.1709	0.15841\\
4.1743	0.15708\\
4.1777	0.15576\\
4.1811	0.15445\\
4.1845	0.15314\\
4.1879	0.15185\\
4.1913	0.15057\\
4.1947	0.14929\\
4.1981	0.14802\\
4.2015	0.14677\\
4.2049	0.14552\\
4.2084	0.14428\\
4.2118	0.14304\\
4.2152	0.14182\\
4.2186	0.1406\\
4.222	0.13939\\
4.2254	0.13819\\
4.2288	0.137\\
4.2322	0.13581\\
4.2356	0.13464\\
4.239	0.13347\\
4.2424	0.13231\\
4.2459	0.13115\\
4.2493	0.13\\
4.2527	0.12886\\
4.2561	0.12773\\
4.2595	0.1266\\
4.2629	0.12548\\
4.2663	0.12437\\
4.2697	0.12327\\
4.2731	0.12217\\
4.2765	0.12108\\
4.2799	0.11999\\
4.2834	0.11891\\
4.2868	0.11784\\
4.2902	0.11678\\
4.2936	0.11572\\
4.297	0.11467\\
4.3004	0.11362\\
4.3038	0.11259\\
4.3072	0.11155\\
4.3106	0.11053\\
4.314	0.10951\\
4.3174	0.10849\\
4.3209	0.10749\\
4.3243	0.10649\\
4.3277	0.10549\\
4.3311	0.1045\\
4.3345	0.10352\\
4.3379	0.10255\\
4.3413	0.10158\\
4.3447	0.10061\\
4.3481	0.099653\\
4.3515	0.098701\\
4.3549	0.097755\\
4.3584	0.096815\\
4.3618	0.095881\\
4.3652	0.094952\\
4.3686	0.09403\\
4.372	0.093113\\
4.3754	0.092203\\
4.3788	0.091298\\
4.3822	0.090399\\
4.3856	0.089505\\
4.389	0.088618\\
4.3924	0.087736\\
4.3959	0.08686\\
4.3993	0.08599\\
4.4027	0.085125\\
4.4061	0.084266\\
4.4095	0.083413\\
4.4129	0.082566\\
4.4163	0.081724\\
4.4197	0.080888\\
4.4231	0.080057\\
4.4265	0.079232\\
4.4299	0.078413\\
4.4334	0.0776\\
4.4368	0.076792\\
4.4402	0.075989\\
4.4436	0.075193\\
4.447	0.074402\\
4.4504	0.073616\\
4.4538	0.072836\\
4.4572	0.072062\\
4.4606	0.071293\\
4.464	0.070529\\
4.4674	0.069772\\
4.4709	0.069019\\
4.4743	0.068273\\
4.4777	0.067531\\
4.4811	0.066796\\
4.4845	0.066065\\
4.4879	0.065341\\
4.4913	0.064621\\
4.4947	0.063908\\
4.4981	0.063199\\
4.5015	0.062496\\
4.5049	0.061799\\
4.5084	0.061107\\
4.5118	0.06042\\
4.5152	0.059739\\
4.5186	0.059063\\
4.522	0.058393\\
4.5254	0.057728\\
4.5288	0.057068\\
4.5322	0.056414\\
4.5356	0.055765\\
4.539	0.055122\\
4.5424	0.054483\\
4.5459	0.05385\\
4.5493	0.053223\\
4.5527	0.0526\\
4.5561	0.051983\\
4.5595	0.051371\\
4.5629	0.050765\\
4.5663	0.050163\\
4.5697	0.049567\\
4.5731	0.048976\\
4.5765	0.04839\\
4.5799	0.047809\\
4.5834	0.047233\\
4.5868	0.046663\\
4.5902	0.046097\\
4.5936	0.045537\\
4.597	0.044981\\
4.6004	0.044431\\
4.6038	0.043885\\
4.6072	0.043345\\
4.6106	0.042809\\
4.614	0.042279\\
4.6174	0.041753\\
4.6209	0.041232\\
4.6243	0.040716\\
4.6277	0.040205\\
4.6311	0.039699\\
4.6345	0.039197\\
4.6379	0.0387\\
4.6413	0.038208\\
4.6447	0.037721\\
4.6481	0.037238\\
4.6515	0.03676\\
4.6549	0.036287\\
4.6584	0.035818\\
4.6618	0.035354\\
4.6652	0.034894\\
4.6686	0.034439\\
4.672	0.033988\\
4.6754	0.033542\\
4.6788	0.0331\\
4.6822	0.032663\\
4.6856	0.03223\\
4.689	0.031802\\
4.6924	0.031377\\
4.6959	0.030957\\
4.6993	0.030542\\
4.7027	0.030131\\
4.7061	0.029724\\
4.7095	0.029321\\
4.7129	0.028922\\
4.7163	0.028528\\
4.7197	0.028137\\
4.7231	0.027751\\
4.7265	0.027369\\
4.7299	0.026991\\
4.7334	0.026617\\
4.7368	0.026247\\
4.7402	0.025881\\
4.7436	0.025519\\
4.747	0.025161\\
4.7504	0.024807\\
4.7538	0.024457\\
4.7572	0.02411\\
4.7606	0.023768\\
4.764	0.023429\\
4.7674	0.023094\\
4.7709	0.022763\\
4.7743	0.022436\\
4.7777	0.022112\\
4.7811	0.021792\\
4.7845	0.021476\\
4.7879	0.021163\\
4.7913	0.020854\\
4.7947	0.020548\\
4.7981	0.020246\\
4.8015	0.019948\\
4.8049	0.019653\\
4.8083	0.019361\\
4.8118	0.019073\\
4.8152	0.018789\\
4.8186	0.018508\\
4.822	0.01823\\
4.8254	0.017956\\
4.8288	0.017684\\
4.8322	0.017417\\
4.8356	0.017152\\
4.839	0.016891\\
4.8424	0.016633\\
4.8458	0.016378\\
4.8493	0.016126\\
4.8527	0.015877\\
4.8561	0.015632\\
4.8595	0.01539\\
4.8629	0.01515\\
4.8663	0.014914\\
4.8697	0.014681\\
4.8731	0.014451\\
4.8765	0.014223\\
4.8799	0.013999\\
4.8833	0.013777\\
4.8868	0.013559\\
4.8902	0.013343\\
4.8936	0.01313\\
4.897	0.01292\\
4.9004	0.012713\\
4.9038	0.012508\\
4.9072	0.012306\\
4.9106	0.012107\\
4.914	0.011911\\
4.9174	0.011717\\
4.9208	0.011526\\
4.9243	0.011337\\
4.9277	0.011151\\
4.9311	0.010967\\
4.9345	0.010787\\
4.9379	0.010608\\
4.9413	0.010432\\
4.9447	0.010259\\
4.9481	0.010087\\
4.9515	0.0099187\\
4.9549	0.0097523\\
4.9583	0.0095883\\
4.9618	0.0094266\\
4.9652	0.0092673\\
4.9686	0.0091102\\
4.972	0.0089553\\
4.9754	0.0088027\\
4.9788	0.0086523\\
4.9822	0.008504\\
4.9856	0.0083579\\
4.989	0.008214\\
4.9924	0.0080721\\
4.9958	0.0079324\\
4.9993	0.0077947\\
5.0027	0.007659\\
5.0061	0.0075254\\
5.0095	0.0073937\\
5.0129	0.007264\\
5.0163	0.0071363\\
5.0197	0.0070105\\
5.0231	0.0068866\\
5.0265	0.0067646\\
5.0299	0.0066444\\
5.0333	0.0065261\\
5.0368	0.0064095\\
5.0402	0.0062948\\
5.0436	0.0061818\\
5.047	0.0060706\\
5.0504	0.0059612\\
5.0538	0.0058534\\
5.0572	0.0057473\\
5.0606	0.0056429\\
5.064	0.0055401\\
5.0674	0.0054389\\
5.0708	0.0053394\\
5.0743	0.0052414\\
5.0777	0.005145\\
5.0811	0.0050501\\
5.0845	0.0049568\\
5.0879	0.0048649\\
5.0913	0.0047746\\
5.0947	0.0046857\\
5.0981	0.0045983\\
5.1015	0.0045122\\
5.1049	0.0044276\\
5.1083	0.0043444\\
5.1118	0.0042625\\
5.1152	0.004182\\
5.1186	0.0041029\\
5.122	0.004025\\
5.1254	0.0039485\\
5.1288	0.0038732\\
5.1322	0.0037992\\
5.1356	0.0037264\\
5.139	0.0036549\\
5.1424	0.0035846\\
5.1458	0.0035154\\
5.1493	0.0034475\\
5.1527	0.0033807\\
5.1561	0.003315\\
5.1595	0.0032505\\
5.1629	0.0031871\\
5.1663	0.0031248\\
5.1697	0.0030635\\
5.1731	0.0030034\\
5.1765	0.0029443\\
5.1799	0.0028862\\
5.1833	0.0028291\\
5.1868	0.002773\\
5.1902	0.002718\\
5.1936	0.0026638\\
5.197	0.0026107\\
5.2004	0.0025585\\
5.2038	0.0025072\\
5.2072	0.0024569\\
5.2106	0.0024074\\
5.214	0.0023589\\
5.2174	0.0023112\\
5.2208	0.0022644\\
5.2243	0.0022184\\
5.2277	0.0021733\\
5.2311	0.002129\\
5.2345	0.0020855\\
5.2379	0.0020427\\
5.2413	0.0020008\\
5.2447	0.0019597\\
5.2481	0.0019193\\
5.2515	0.0018797\\
5.2549	0.0018408\\
5.2583	0.0018026\\
5.2618	0.0017651\\
5.2652	0.0017283\\
5.2686	0.0016923\\
5.272	0.0016569\\
5.2754	0.0016222\\
5.2788	0.0015881\\
5.2822	0.0015547\\
5.2856	0.0015219\\
5.289	0.0014897\\
5.2924	0.0014582\\
5.2924	0.0013044\\
5.289	0.0013335\\
5.2856	0.0013632\\
5.2822	0.0013934\\
5.2788	0.0014243\\
5.2754	0.0014558\\
5.272	0.0014879\\
5.2686	0.0015207\\
5.2652	0.0015541\\
5.2618	0.0015881\\
5.2583	0.0016229\\
5.2549	0.0016583\\
5.2515	0.0016944\\
5.2481	0.0017312\\
5.2447	0.0017688\\
5.2413	0.001807\\
5.2379	0.0018461\\
5.2345	0.0018858\\
5.2311	0.0019263\\
5.2277	0.0019677\\
5.2243	0.0020098\\
5.2208	0.0020527\\
5.2174	0.0020964\\
5.214	0.002141\\
5.2106	0.0021864\\
5.2072	0.0022326\\
5.2038	0.0022798\\
5.2004	0.0023278\\
5.197	0.0023767\\
5.1936	0.0024265\\
5.1902	0.0024773\\
5.1868	0.002529\\
5.1833	0.0025817\\
5.1799	0.0026353\\
5.1765	0.0026899\\
5.1731	0.0027456\\
5.1697	0.0028022\\
5.1663	0.0028599\\
5.1629	0.0029187\\
5.1595	0.0029785\\
5.1561	0.0030394\\
5.1527	0.0031013\\
5.1493	0.0031644\\
5.1458	0.0032287\\
5.1424	0.0032941\\
5.139	0.0033606\\
5.1356	0.0034283\\
5.1322	0.0034973\\
5.1288	0.0035674\\
5.1254	0.0036388\\
5.122	0.0037115\\
5.1186	0.0037854\\
5.1152	0.0038606\\
5.1118	0.0039371\\
5.1083	0.0040149\\
5.1049	0.0040941\\
5.1015	0.0041747\\
5.0981	0.0042566\\
5.0947	0.0043399\\
5.0913	0.0044247\\
5.0879	0.0045109\\
5.0845	0.0045986\\
5.0811	0.0046877\\
5.0777	0.0047784\\
5.0743	0.0048705\\
5.0708	0.0049642\\
5.0674	0.0050595\\
5.064	0.0051564\\
5.0606	0.0052548\\
5.0572	0.0053549\\
5.0538	0.0054567\\
5.0504	0.0055601\\
5.047	0.0056652\\
5.0436	0.005772\\
5.0402	0.0058806\\
5.0368	0.0059909\\
5.0333	0.006103\\
5.0299	0.0062169\\
5.0265	0.0063326\\
5.0231	0.0064502\\
5.0197	0.0065696\\
5.0163	0.0066909\\
5.0129	0.0068142\\
5.0095	0.0069394\\
5.0061	0.0070665\\
5.0027	0.0071957\\
4.9993	0.0073268\\
4.9958	0.00746\\
4.9924	0.0075953\\
4.989	0.0077326\\
4.9856	0.0078721\\
4.9822	0.0080136\\
4.9788	0.0081574\\
4.9754	0.0083033\\
4.972	0.0084514\\
4.9686	0.0086018\\
4.9652	0.0087544\\
4.9618	0.0089093\\
4.9583	0.0090665\\
4.9549	0.0092261\\
4.9515	0.009388\\
4.9481	0.0095523\\
4.9447	0.009719\\
4.9413	0.0098881\\
4.9379	0.01006\\
4.9345	0.010234\\
4.9311	0.01041\\
4.9277	0.01059\\
4.9243	0.010771\\
4.9208	0.010956\\
4.9174	0.011143\\
4.914	0.011332\\
4.9106	0.011525\\
4.9072	0.011719\\
4.9038	0.011917\\
4.9004	0.012118\\
4.897	0.012321\\
4.8936	0.012527\\
4.8902	0.012736\\
4.8868	0.012948\\
4.8833	0.013162\\
4.8799	0.01338\\
4.8765	0.0136\\
4.8731	0.013824\\
4.8697	0.01405\\
4.8663	0.01428\\
4.8629	0.014513\\
4.8595	0.014748\\
4.8561	0.014987\\
4.8527	0.015229\\
4.8493	0.015474\\
4.8458	0.015722\\
4.8424	0.015974\\
4.839	0.016228\\
4.8356	0.016486\\
4.8322	0.016748\\
4.8288	0.017012\\
4.8254	0.01728\\
4.822	0.017552\\
4.8186	0.017827\\
4.8152	0.018105\\
4.8118	0.018387\\
4.8083	0.018672\\
4.8049	0.01896\\
4.8015	0.019253\\
4.7981	0.019548\\
4.7947	0.019848\\
4.7913	0.020151\\
4.7879	0.020457\\
4.7845	0.020768\\
4.7811	0.021082\\
4.7777	0.0214\\
4.7743	0.021721\\
4.7709	0.022046\\
4.7674	0.022375\\
4.764	0.022708\\
4.7606	0.023045\\
4.7572	0.023386\\
4.7538	0.02373\\
4.7504	0.024078\\
4.747	0.024431\\
4.7436	0.024787\\
4.7402	0.025147\\
4.7368	0.025512\\
4.7334	0.02588\\
4.7299	0.026252\\
4.7265	0.026629\\
4.7231	0.027009\\
4.7197	0.027394\\
4.7163	0.027783\\
4.7129	0.028176\\
4.7095	0.028573\\
4.7061	0.028974\\
4.7027	0.02938\\
4.6993	0.02979\\
4.6959	0.030204\\
4.6924	0.030622\\
4.689	0.031044\\
4.6856	0.031471\\
4.6822	0.031902\\
4.6788	0.032338\\
4.6754	0.032778\\
4.672	0.033222\\
4.6686	0.033671\\
4.6652	0.034124\\
4.6618	0.034581\\
4.6584	0.035043\\
4.6549	0.035509\\
4.6515	0.03598\\
4.6481	0.036455\\
4.6447	0.036935\\
4.6413	0.037419\\
4.6379	0.037908\\
4.6345	0.038401\\
4.6311	0.038899\\
4.6277	0.039401\\
4.6243	0.039908\\
4.6209	0.040419\\
4.6174	0.040935\\
4.614	0.041455\\
4.6106	0.04198\\
4.6072	0.04251\\
4.6038	0.043044\\
4.6004	0.043583\\
4.597	0.044126\\
4.5936	0.044674\\
4.5902	0.045226\\
4.5868	0.045783\\
4.5834	0.046345\\
4.5799	0.046912\\
4.5765	0.047483\\
4.5731	0.048058\\
4.5697	0.048639\\
4.5663	0.049224\\
4.5629	0.049813\\
4.5595	0.050408\\
4.5561	0.051007\\
4.5527	0.051611\\
4.5493	0.052219\\
4.5459	0.052832\\
4.5424	0.05345\\
4.539	0.054073\\
4.5356	0.0547\\
4.5322	0.055332\\
4.5288	0.055969\\
4.5254	0.056611\\
4.522	0.057257\\
4.5186	0.057908\\
4.5152	0.058564\\
4.5118	0.059225\\
4.5084	0.05989\\
4.5049	0.060561\\
4.5015	0.061236\\
4.4981	0.061916\\
4.4947	0.062601\\
4.4913	0.063291\\
4.4879	0.063985\\
4.4845	0.064685\\
4.4811	0.065389\\
4.4777	0.066099\\
4.4743	0.066813\\
4.4709	0.067532\\
4.4674	0.068257\\
4.464	0.068986\\
4.4606	0.06972\\
4.4572	0.070459\\
4.4538	0.071203\\
4.4504	0.071952\\
4.447	0.072706\\
4.4436	0.073466\\
4.4402	0.07423\\
4.4368	0.074999\\
4.4334	0.075774\\
4.4299	0.076554\\
4.4265	0.077338\\
4.4231	0.078128\\
4.4197	0.078923\\
4.4163	0.079724\\
4.4129	0.080529\\
4.4095	0.08134\\
4.4061	0.082156\\
4.4027	0.082977\\
4.3993	0.083804\\
4.3959	0.084635\\
4.3924	0.085473\\
4.389	0.086315\\
4.3856	0.087163\\
4.3822	0.088017\\
4.3788	0.088875\\
4.3754	0.08974\\
4.372	0.090609\\
4.3686	0.091485\\
4.3652	0.092365\\
4.3618	0.093252\\
4.3584	0.094144\\
4.3549	0.095042\\
4.3515	0.095945\\
4.3481	0.096854\\
4.3447	0.097769\\
4.3413	0.09869\\
4.3379	0.099616\\
4.3345	0.10055\\
4.3311	0.10149\\
4.3277	0.10243\\
4.3243	0.10338\\
4.3209	0.10434\\
4.3174	0.1053\\
4.314	0.10627\\
4.3106	0.10724\\
4.3072	0.10822\\
4.3038	0.10921\\
4.3004	0.11021\\
4.297	0.11121\\
4.2936	0.11221\\
4.2902	0.11323\\
4.2868	0.11425\\
4.2834	0.11527\\
4.2799	0.11631\\
4.2765	0.11735\\
4.2731	0.11839\\
4.2697	0.11945\\
4.2663	0.12051\\
4.2629	0.12158\\
4.2595	0.12265\\
4.2561	0.12373\\
4.2527	0.12482\\
4.2493	0.12592\\
4.2459	0.12703\\
4.2424	0.12814\\
4.239	0.12926\\
4.2356	0.13039\\
4.2322	0.13152\\
4.2288	0.13267\\
4.2254	0.13382\\
4.222	0.13498\\
4.2186	0.13615\\
4.2152	0.13733\\
4.2118	0.13851\\
4.2084	0.13971\\
4.2049	0.14091\\
4.2015	0.14212\\
4.1981	0.14335\\
4.1947	0.14458\\
4.1913	0.14582\\
4.1879	0.14707\\
4.1845	0.14832\\
4.1811	0.14959\\
4.1777	0.15087\\
4.1743	0.15216\\
4.1709	0.15346\\
4.1674	0.15477\\
4.164	0.15609\\
4.1606	0.15742\\
4.1572	0.15876\\
4.1538	0.16011\\
4.1504	0.16148\\
4.147	0.16285\\
4.1436	0.16424\\
4.1402	0.16563\\
4.1368	0.16704\\
4.1334	0.16846\\
4.1299	0.1699\\
4.1265	0.17134\\
4.1231	0.1728\\
4.1197	0.17427\\
4.1163	0.17576\\
4.1129	0.17726\\
4.1095	0.17877\\
4.1061	0.18029\\
4.1027	0.18183\\
4.0993	0.18338\\
4.0959	0.18495\\
4.0924	0.18653\\
4.089	0.18813\\
4.0856	0.18974\\
4.0822	0.19136\\
4.0788	0.193\\
4.0754	0.19466\\
4.072	0.19633\\
4.0686	0.19802\\
4.0652	0.19973\\
4.0618	0.20145\\
4.0584	0.20318\\
4.0549	0.20494\\
4.0515	0.20671\\
4.0481	0.2085\\
4.0447	0.21031\\
4.0413	0.21213\\
4.0379	0.21398\\
4.0345	0.21584\\
4.0311	0.21772\\
4.0277	0.21962\\
4.0243	0.22153\\
4.0209	0.22347\\
4.0174	0.22543\\
4.014	0.22741\\
4.0106	0.2294\\
4.0072	0.23142\\
4.0038	0.23346\\
4.0004	0.23552\\
3.997	0.2376\\
3.9936	0.2397\\
3.9902	0.24183\\
3.9868	0.24397\\
3.9834	0.24614\\
3.9799	0.24833\\
3.9765	0.25055\\
3.9731	0.25278\\
3.9697	0.25504\\
3.9663	0.25732\\
3.9629	0.25963\\
3.9595	0.26196\\
3.9561	0.26432\\
3.9527	0.2667\\
3.9493	0.2691\\
3.9459	0.27153\\
3.9424	0.27398\\
3.939	0.27646\\
3.9356	0.27897\\
3.9322	0.2815\\
3.9288	0.28405\\
3.9254	0.28664\\
3.922	0.28925\\
3.9186	0.29189\\
3.9152	0.29455\\
3.9118	0.29724\\
3.9084	0.29996\\
3.9049	0.30271\\
3.9015	0.30548\\
3.8981	0.30829\\
3.8947	0.31112\\
3.8913	0.31398\\
3.8879	0.31687\\
3.8845	0.31979\\
3.8811	0.32274\\
3.8777	0.32572\\
3.8743	0.32873\\
3.8709	0.33177\\
3.8674	0.33484\\
3.864	0.33794\\
3.8606	0.34108\\
3.8572	0.34424\\
3.8538	0.34744\\
3.8504	0.35067\\
3.847	0.35393\\
3.8436	0.35722\\
3.8402	0.36054\\
3.8368	0.3639\\
3.8334	0.36729\\
3.8299	0.37071\\
3.8265	0.37417\\
3.8231	0.37766\\
3.8197	0.38119\\
3.8163	0.38475\\
3.8129	0.38834\\
3.8095	0.39196\\
3.8061	0.39563\\
3.8027	0.39932\\
3.7993	0.40305\\
3.7959	0.40682\\
3.7924	0.41062\\
3.789	0.41445\\
3.7856	0.41832\\
3.7822	0.42223\\
3.7788	0.42617\\
3.7754	0.43014\\
3.772	0.43415\\
3.7686	0.4382\\
3.7652	0.44228\\
3.7618	0.44639\\
3.7584	0.45054\\
3.7549	0.45473\\
3.7515	0.45895\\
3.7481	0.46321\\
3.7447	0.4675\\
3.7413	0.47182\\
3.7379	0.47618\\
3.7345	0.48058\\
3.7311	0.485\\
3.7277	0.48947\\
3.7243	0.49396\\
3.7209	0.49849\\
3.7174	0.50306\\
3.714	0.50765\\
3.7106	0.51228\\
3.7072	0.51694\\
3.7038	0.52164\\
3.7004	0.52636\\
3.697	0.53112\\
3.6936	0.53591\\
3.6902	0.54073\\
3.6868	0.54558\\
3.6834	0.55046\\
3.6799	0.55536\\
3.6765	0.5603\\
3.6731	0.56527\\
3.6697	0.57026\\
3.6663	0.57528\\
3.6629	0.58033\\
3.6595	0.5854\\
3.6561	0.5905\\
3.6527	0.59563\\
3.6493	0.60078\\
3.6459	0.60595\\
3.6424	0.61114\\
3.639	0.61636\\
3.6356	0.6216\\
3.6322	0.62686\\
3.6288	0.63213\\
3.6254	0.63743\\
3.622	0.64275\\
3.6186	0.64808\\
3.6152	0.65343\\
3.6118	0.65879\\
3.6084	0.66417\\
3.6049	0.66956\\
3.6015	0.67497\\
3.5981	0.68039\\
3.5947	0.68581\\
3.5913	0.69125\\
3.5879	0.69669\\
3.5845	0.70215\\
3.5811	0.70761\\
3.5777	0.71307\\
3.5743	0.71854\\
3.5709	0.72401\\
3.5674	0.72948\\
3.564	0.73495\\
3.5606	0.74042\\
3.5572	0.74589\\
3.5538	0.75136\\
3.5504	0.75682\\
3.547	0.76227\\
3.5436	0.76772\\
3.5402	0.77316\\
3.5368	0.77859\\
3.5334	0.78401\\
3.5299	0.78941\\
3.5265	0.7948\\
3.5231	0.80018\\
3.5197	0.80553\\
3.5163	0.81087\\
3.5129	0.81619\\
3.5095	0.82149\\
3.5061	0.82676\\
3.5027	0.83201\\
3.4993	0.83723\\
3.4959	0.84243\\
3.4924	0.84759\\
3.489	0.85273\\
3.4856	0.85783\\
3.4822	0.8629\\
3.4788	0.86793\\
3.4754	0.87292\\
3.472	0.87788\\
3.4686	0.88279\\
3.4652	0.88767\\
3.4618	0.8925\\
3.4584	0.89728\\
3.4549	0.90202\\
3.4515	0.9067\\
3.4481	0.91134\\
3.4447	0.91593\\
3.4413	0.92046\\
3.4379	0.92494\\
3.4345	0.92936\\
3.4311	0.93372\\
3.4277	0.93802\\
3.4243	0.94226\\
3.4209	0.94644\\
3.4174	0.95055\\
3.414	0.95459\\
3.4106	0.95857\\
3.4072	0.96248\\
3.4038	0.96632\\
3.4004	0.97008\\
3.397	0.97377\\
3.3936	0.97739\\
3.3902	0.98093\\
3.3868	0.98439\\
3.3834	0.98778\\
3.3799	0.99108\\
3.3765	0.9943\\
3.3731	0.99744\\
3.3697	1.0005\\
3.3663	1.0035\\
3.3629	1.0063\\
3.3595	1.0091\\
3.3561	1.0118\\
3.3527	1.0144\\
3.3493	1.017\\
3.3459	1.0194\\
3.3424	1.0217\\
3.339	1.024\\
3.3356	1.0261\\
3.3322	1.0282\\
3.3288	1.0301\\
3.3254	1.032\\
3.322	1.0338\\
3.3186	1.0354\\
3.3152	1.037\\
3.3118	1.0385\\
3.3084	1.0398\\
3.3049	1.0411\\
3.3015	1.0422\\
3.2981	1.0433\\
3.2947	1.0442\\
3.2913	1.0451\\
3.2879	1.0458\\
3.2845	1.0465\\
3.2811	1.047\\
3.2777	1.0474\\
3.2743	1.0478\\
3.2709	1.048\\
3.2674	1.0481\\
3.264	1.0481\\
3.2606	1.048\\
3.2572	1.0478\\
3.2538	1.0475\\
3.2504	1.0471\\
3.247	1.0465\\
3.2436	1.0459\\
3.2402	1.0452\\
3.2368	1.0443\\
3.2334	1.0434\\
3.2299	1.0423\\
3.2265	1.0411\\
3.2231	1.0399\\
3.2197	1.0385\\
3.2163	1.037\\
3.2129	1.0354\\
3.2095	1.0337\\
3.2061	1.0319\\
3.2027	1.03\\
3.1993	1.028\\
3.1959	1.0259\\
3.1924	1.0237\\
3.189	1.0214\\
3.1856	1.019\\
3.1822	1.0165\\
3.1788	1.0138\\
3.1754	1.0111\\
3.172	1.0083\\
3.1686	1.0054\\
3.1652	1.0024\\
3.1618	0.99929\\
3.1584	0.99609\\
3.1549	0.9928\\
3.1515	0.98941\\
3.1481	0.98593\\
3.1447	0.98235\\
3.1413	0.97869\\
3.1379	0.97493\\
3.1345	0.97109\\
3.1311	0.96716\\
3.1277	0.96314\\
3.1243	0.95903\\
3.1209	0.95485\\
3.1174	0.95057\\
3.114	0.94622\\
3.1106	0.94178\\
3.1072	0.93727\\
3.1038	0.93267\\
3.1004	0.928\\
3.097	0.92326\\
3.0936	0.91844\\
3.0902	0.91354\\
3.0868	0.90857\\
3.0834	0.90354\\
3.0799	0.89843\\
3.0765	0.89326\\
3.0731	0.88802\\
3.0697	0.88271\\
3.0663	0.87734\\
3.0629	0.87191\\
3.0595	0.86642\\
3.0561	0.86087\\
3.0527	0.85526\\
3.0493	0.8496\\
3.0459	0.84388\\
3.0424	0.83811\\
3.039	0.83228\\
3.0356	0.82641\\
3.0322	0.82049\\
3.0288	0.81452\\
3.0254	0.80851\\
3.022	0.80245\\
3.0186	0.79635\\
3.0152	0.79021\\
3.0118	0.78403\\
3.0084	0.77781\\
3.0049	0.77156\\
3.0015	0.76527\\
2.9981	0.75895\\
2.9947	0.7526\\
2.9913	0.74622\\
2.9879	0.73981\\
2.9845	0.73337\\
2.9811	0.72691\\
2.9777	0.72043\\
2.9743	0.71393\\
2.9709	0.7074\\
2.9674	0.70086\\
2.964	0.6943\\
2.9606	0.68773\\
2.9572	0.68114\\
2.9538	0.67454\\
2.9504	0.66793\\
2.947	0.66131\\
2.9436	0.65469\\
2.9402	0.64806\\
2.9368	0.64142\\
2.9334	0.63479\\
2.9299	0.62815\\
2.9265	0.62151\\
2.9231	0.61487\\
2.9197	0.60824\\
2.9163	0.60162\\
2.9129	0.595\\
2.9095	0.58838\\
2.9061	0.58178\\
2.9027	0.57519\\
2.8993	0.56861\\
2.8959	0.56205\\
2.8924	0.5555\\
2.889	0.54896\\
2.8856	0.54245\\
2.8822	0.53595\\
2.8788	0.52947\\
2.8754	0.52302\\
2.872	0.51659\\
2.8686	0.51018\\
2.8652	0.5038\\
2.8618	0.49745\\
2.8584	0.49112\\
2.8549	0.48482\\
2.8515	0.47855\\
2.8481	0.47231\\
2.8447	0.46611\\
2.8413	0.45993\\
2.8379	0.45379\\
2.8345	0.44769\\
2.8311	0.44162\\
2.8277	0.43559\\
2.8243	0.4296\\
2.8209	0.42364\\
2.8174	0.41772\\
2.814	0.41185\\
2.8106	0.40601\\
2.8072	0.40022\\
2.8038	0.39446\\
2.8004	0.38875\\
2.797	0.38309\\
2.7936	0.37746\\
2.7902	0.37189\\
2.7868	0.36635\\
2.7834	0.36087\\
2.7799	0.35543\\
2.7765	0.35003\\
2.7731	0.34469\\
2.7697	0.33939\\
2.7663	0.33414\\
2.7629	0.32894\\
2.7595	0.32378\\
2.7561	0.31868\\
2.7527	0.31362\\
2.7493	0.30862\\
2.7459	0.30366\\
2.7424	0.29876\\
2.739	0.2939\\
2.7356	0.2891\\
2.7322	0.28435\\
2.7288	0.27965\\
2.7254	0.27499\\
2.722	0.2704\\
2.7186	0.26585\\
2.7152	0.26135\\
2.7118	0.25691\\
2.7084	0.25251\\
2.7049	0.24817\\
2.7015	0.24388\\
2.6981	0.23964\\
2.6947	0.23545\\
2.6913	0.23132\\
2.6879	0.22723\\
2.6845	0.2232\\
2.6811	0.21922\\
2.6777	0.21529\\
2.6743	0.21141\\
2.6709	0.20758\\
2.6674	0.2038\\
2.664	0.20007\\
2.6606	0.19639\\
2.6572	0.19277\\
2.6538	0.18919\\
2.6504	0.18566\\
2.647	0.18218\\
2.6436	0.17875\\
2.6402	0.17537\\
2.6368	0.17204\\
2.6334	0.16876\\
2.6299	0.16552\\
2.6265	0.16234\\
2.6231	0.15919\\
2.6197	0.1561\\
2.6163	0.15306\\
2.6129	0.15006\\
2.6095	0.1471\\
2.6061	0.14419\\
2.6027	0.14133\\
2.5993	0.13851\\
2.5959	0.13574\\
2.5924	0.13301\\
2.589	0.13033\\
2.5856	0.12769\\
2.5822	0.12509\\
2.5788	0.12253\\
2.5754	0.12002\\
2.572	0.11755\\
2.5686	0.11512\\
2.5652	0.11273\\
2.5618	0.11038\\
2.5584	0.10807\\
2.5549	0.1058\\
2.5515	0.10357\\
2.5481	0.10138\\
2.5447	0.09923\\
2.5413	0.097116\\
2.5379	0.095039\\
2.5345	0.093\\
2.5311	0.090997\\
2.5277	0.08903\\
2.5243	0.087099\\
2.5209	0.085204\\
2.5174	0.083343\\
2.514	0.081517\\
2.5106	0.079724\\
2.5072	0.077965\\
2.5038	0.07624\\
2.5004	0.074546\\
2.497	0.072885\\
2.4936	0.071256\\
2.4902	0.069658\\
2.4868	0.068091\\
2.4834	0.066554\\
2.4799	0.065047\\
2.4765	0.06357\\
2.4731	0.062122\\
2.4697	0.060702\\
2.4663	0.059311\\
2.4629	0.057948\\
2.4595	0.056612\\
2.4561	0.055302\\
2.4527	0.05402\\
2.4493	0.052763\\
2.4459	0.051532\\
2.4424	0.050327\\
2.439	0.049146\\
2.4356	0.04799\\
2.4322	0.046858\\
2.4288	0.045749\\
2.4254	0.044664\\
2.422	0.043602\\
2.4186	0.042562\\
2.4152	0.041544\\
2.4118	0.040548\\
2.4084	0.039574\\
2.4049	0.03862\\
2.4015	0.037687\\
2.3981	0.036774\\
2.3947	0.035881\\
2.3913	0.035008\\
2.3879	0.034154\\
2.3845	0.033319\\
2.3811	0.032502\\
2.3777	0.031703\\
2.3743	0.030922\\
2.3709	0.030159\\
2.3674	0.029412\\
2.364	0.028683\\
2.3606	0.02797\\
2.3572	0.027273\\
2.3538	0.026592\\
2.3504	0.025927\\
2.347	0.025277\\
2.3436	0.024642\\
2.3402	0.024022\\
2.3368	0.023416\\
2.3334	0.022824\\
2.3299	0.022246\\
2.3265	0.021681\\
2.3231	0.02113\\
2.3197	0.020591\\
2.3163	0.020066\\
2.3129	0.019553\\
2.3095	0.019052\\
2.3061	0.018563\\
2.3027	0.018086\\
2.2993	0.01762\\
2.2959	0.017165\\
2.2924	0.016722\\
2.289	0.016289\\
2.2856	0.015867\\
2.2822	0.015455\\
2.2788	0.015053\\
2.2754	0.01466\\
2.272	0.014278\\
2.2686	0.013905\\
2.2652	0.013541\\
2.2618	0.013186\\
2.2584	0.01284\\
2.2549	0.012503\\
2.2515	0.012174\\
2.2481	0.011853\\
2.2447	0.011541\\
2.2413	0.011236\\
2.2379	0.010939\\
2.2345	0.010649\\
2.2311	0.010367\\
2.2277	0.010092\\
2.2243	0.0098235\\
2.2209	0.0095622\\
2.2174	0.0093077\\
2.214	0.0090596\\
2.2106	0.0088179\\
2.2072	0.0085825\\
2.2038	0.0083531\\
2.2004	0.0081296\\
2.197	0.007912\\
2.1936	0.0077\\
2.1902	0.0074935\\
2.1868	0.0072923\\
2.1834	0.0070965\\
2.1799	0.0069057\\
2.1765	0.00672\\
2.1731	0.0065391\\
2.1697	0.006363\\
2.1663	0.0061915\\
2.1629	0.0060246\\
2.1595	0.005862\\
2.1561	0.0057038\\
2.1527	0.0055498\\
2.1493	0.0053998\\
2.1459	0.0052539\\
2.1424	0.0051118\\
2.139	0.0049735\\
2.1356	0.004839\\
2.1322	0.004708\\
2.1288	0.0045805\\
2.1254	0.0044565\\
2.122	0.0043357\\
2.1186	0.0042183\\
2.1152	0.004104\\
2.1118	0.0039927\\
2.1084	0.0038845\\
2.1049	0.0037792\\
2.1015	0.0036767\\
2.0981	0.0035771\\
2.0947	0.0034801\\
2.0913	0.0033857\\
2.0879	0.0032939\\
2.0845	0.0032046\\
2.0811	0.0031178\\
2.0777	0.0030333\\
2.0743	0.0029511\\
2.0709	0.0028711\\
2.0674	0.0027933\\
2.064	0.0027177\\
2.0606	0.0026441\\
2.0572	0.0025725\\
2.0538	0.0025029\\
2.0504	0.0024352\\
2.047	0.0023693\\
2.0436	0.0023052\\
2.0402	0.0022429\\
2.0368	0.0021823\\
2.0334	0.0021234\\
2.0299	0.0020661\\
2.0265	0.0020103\\
2.0231	0.0019561\\
2.0197	0.0019034\\
2.0163	0.0018521\\
2.0129	0.0018022\\
2.0095	0.0017537\\
2.0061	0.0017065\\
2.0027	0.0016606\\
1.9993	0.001616\\
1.9959	0.0015726\\
1.9924	0.0015304\\
1.989	0.0014894\\
1.9856	0.0014494\\
1.9822	0.0014106\\
1.9788	0.0013729\\
1.9754	0.0013361\\
1.972	0.0013004\\
1.9686	0.0012657\\
1.9652	0.0012319\\
1.9618	0.001199\\
1.9584	0.0011671\\
1.9549	0.001136\\
1.9515	0.0011058\\
1.9481	0.0010764\\
1.9447	0.0010478\\
1.9413	0.0010199\\
1.9379	0.00099287\\
1.9345	0.00096655\\
1.9311	0.00094095\\
1.9277	0.00091604\\
1.9243	0.00089182\\
1.9209	0.00086825\\
1.9174	0.00084532\\
1.914	0.00082302\\
1.9106	0.00080132\\
1.9072	0.00078022\\
1.9038	0.00075968\\
1.9004	0.00073971\\
1.897	0.00072027\\
1.8936	0.00070136\\
1.8902	0.00068296\\
1.8868	0.00066506\\
}--cycle;
      % \input{../../../../.tex/20/KS/SizeDistributionFittingsfig10.tex}
      \legend{}
  
      \nextgroupplot[%
        xmin=0.57363,xmax=5.3544,
        % xmin=1.7156,xmax=5.4813,
        % ymin=0,ymax=1.2,
      ] % This file was created by matlab2tikz.
%
\definecolor{mycolor1}{rgb}{0.00000,0.44706,0.69804}%
%
\addplot[ybar interval, fill=mycolor1, fill opacity=0.35, area legend, draw=none] table[row sep=crcr, x=Lower, y=Count] {%
Lower	Upper	Count\\
1.8868	2.1501	0.0014177\\
2.1501	2.4134	0.01843\\
2.4134	2.6768	0.12354\\
2.6768	2.9401	0.39959\\
2.9401	3.2034	0.88131\\
3.2034	3.4668	1.0112\\
3.4668	3.7301	0.69032\\
3.7301	3.9935	0.35109\\
3.9935	4.2568	0.1845\\
4.2568	4.5201	0.086683\\
4.5201	4.7835	0.031392\\
4.7835	5.0468	0.0078987\\
5.0468	5.3101	0.010025\\
5.3101	5.3101	0.010025\\
};
\addlegendentry{\ExactHistogramLegend}

\addplot [color=mycolor1, only marks, every error bar/.append style={opacity=0.45}, mark=*, mark size=0pt, draw=none, forget plot]
 plot [error bars/.cd, y dir=both, y explicit, error bar style={line width=1pt, color=mycolor1}, error mark options={mark=none,mark size=0pt}]
 table[row sep=crcr, y error plus index=2, y error minus index=3]{%
2.0184	0.0014177	0.00080878	0.00080878\\
2.2818	0.01843	0.0026543	0.0026543\\
2.5451	0.12354	0.0055672	0.0055672\\
2.8084	0.39959	0.010119	0.010119\\
3.0718	0.88131	0.012292	0.012292\\
3.3351	1.0112	0.012357	0.012357\\
3.5985	0.69032	0.010758	0.010758\\
3.8618	0.35109	0.0078942	0.0078942\\
4.1251	0.1845	0.0055052	0.0055052\\
4.3885	0.086683	0.0035247	0.0035247\\
4.6518	0.031392	0.0020295	0.0020295\\
4.9151	0.0078987	0.0014159	0.0014159\\
5.1785	0.010025	0.00020093	0.00020093\\
};
\addplot [color=mycolor1]
  table[row sep=crcr]{%
1.8868	0.00089722\\
2.0033	0.0020991\\
2.0821	0.0037738\\
2.1403	0.0058472\\
2.1883	0.0083932\\
2.2294	0.011431\\
2.2671	0.015143\\
2.298	0.019021\\
2.3288	0.023831\\
2.3562	0.029043\\
2.3802	0.034451\\
2.4042	0.040767\\
2.4282	0.048113\\
2.4488	0.055328\\
2.4693	0.063482\\
2.4899	0.072664\\
2.5104	0.082968\\
2.5276	0.092479\\
2.5447	0.10289\\
2.5618	0.11424\\
2.579	0.12661\\
2.5961	0.14002\\
2.6132	0.15453\\
2.6304	0.17018\\
2.6475	0.18701\\
2.6646	0.20505\\
2.6818	0.22433\\
2.6989	0.24485\\
2.7161	0.26663\\
2.7366	0.29442\\
2.7572	0.32398\\
2.7777	0.35527\\
2.7983	0.38822\\
2.8223	0.42859\\
2.8463	0.47084\\
2.8737	0.52109\\
2.9045	0.57955\\
2.9559	0.67936\\
3.0005	0.76526\\
3.0279	0.8163\\
3.0519	0.85895\\
3.0724	0.89349\\
3.0896	0.92052\\
3.1067	0.9457\\
3.1204	0.96434\\
3.1341	0.98153\\
3.1478	0.99715\\
3.1615	1.0111\\
3.1718	1.0204\\
3.1821	1.0288\\
3.1924	1.0361\\
3.2027	1.0423\\
3.2129	1.0475\\
3.2232	1.0517\\
3.2335	1.0548\\
3.2438	1.0568\\
3.2541	1.0578\\
3.2643	1.0578\\
3.2746	1.0568\\
3.2849	1.0547\\
3.2952	1.0517\\
3.3055	1.0478\\
3.3157	1.0429\\
3.326	1.0372\\
3.3363	1.0306\\
3.35	1.0206\\
3.3637	1.0091\\
3.3774	0.9964\\
3.3911	0.98244\\
3.4048	0.96735\\
3.422	0.94702\\
3.4391	0.92522\\
3.4597	0.89738\\
3.4802	0.86801\\
3.5042	0.8322\\
3.5351	0.78449\\
3.6481	0.60767\\
3.6756	0.56718\\
3.703	0.52841\\
3.727	0.49608\\
3.7509	0.46535\\
3.7749	0.4363\\
3.7955	0.41276\\
3.8161	0.39048\\
3.8366	0.36943\\
3.8572	0.34959\\
3.8777	0.33093\\
3.8983	0.31338\\
3.9189	0.2969\\
3.9428	0.27894\\
3.9668	0.26225\\
3.9908	0.24672\\
4.0148	0.23226\\
4.0422	0.2169\\
4.0696	0.20267\\
4.097	0.18943\\
4.1279	0.17559\\
4.1587	0.16271\\
4.193	0.14938\\
4.2307	0.13577\\
4.2684	0.12313\\
4.3095	0.11031\\
4.3506	0.098424\\
4.3952	0.086526\\
4.4397	0.0756\\
4.4843	0.065615\\
4.5288	0.056547\\
4.5734	0.048374\\
4.6213	0.040539\\
4.6693	0.033667\\
4.7173	0.027702\\
4.7687	0.022246\\
4.8235	0.017397\\
4.8818	0.013217\\
4.9435	0.0097314\\
5.012	0.0067995\\
5.0908	0.0043967\\
5.1833	0.0025522\\
5.2964	0.0012533\\
5.3101	0.0011459\\
};
\addlegendentry{\ExactBLNFitLegend}


\addplot[area legend, draw=none, fill=mycolor1, fill opacity=0.15, forget plot]
table[row sep=crcr] {%
x	y\\
1.8868	0.0011401\\
1.8902	0.0011671\\
1.8936	0.0011947\\
1.897	0.001223\\
1.9005	0.0012519\\
1.9039	0.0012816\\
1.9073	0.001312\\
1.9108	0.0013431\\
1.9142	0.0013749\\
1.9176	0.0014075\\
1.921	0.0014409\\
1.9245	0.0014751\\
1.9279	0.0015101\\
1.9313	0.001546\\
1.9347	0.0015827\\
1.9382	0.0016204\\
1.9416	0.0016589\\
1.945	0.0016984\\
1.9484	0.0017389\\
1.9519	0.0017803\\
1.9553	0.0018227\\
1.9587	0.0018662\\
1.9622	0.0019107\\
1.9656	0.0019564\\
1.969	0.0020031\\
1.9724	0.002051\\
1.9759	0.0021001\\
1.9793	0.0021503\\
1.9827	0.0022018\\
1.9861	0.0022546\\
1.9896	0.0023087\\
1.993	0.0023641\\
1.9964	0.0024208\\
1.9999	0.002479\\
2.0033	0.0025386\\
2.0067	0.0025997\\
2.0101	0.0026622\\
2.0136	0.0027264\\
2.017	0.0027921\\
2.0204	0.0028594\\
2.0238	0.0029285\\
2.0273	0.0029992\\
2.0307	0.0030717\\
2.0341	0.003146\\
2.0375	0.0032221\\
2.041	0.0033002\\
2.0444	0.0033802\\
2.0478	0.0034622\\
2.0513	0.0035462\\
2.0547	0.0036323\\
2.0581	0.0037206\\
2.0615	0.0038111\\
2.065	0.0039038\\
2.0684	0.0039989\\
2.0718	0.0040963\\
2.0752	0.0041962\\
2.0787	0.0042986\\
2.0821	0.0044035\\
2.0855	0.0045111\\
2.0889	0.0046213\\
2.0924	0.0047344\\
2.0958	0.0048502\\
2.0992	0.004969\\
2.1027	0.0050908\\
2.1061	0.0052156\\
2.1095	0.0053435\\
2.1129	0.0054747\\
2.1164	0.0056091\\
2.1198	0.0057469\\
2.1232	0.0058882\\
2.1266	0.006033\\
2.1301	0.0061814\\
2.1335	0.0063336\\
2.1369	0.0064896\\
2.1403	0.0066495\\
2.1438	0.0068134\\
2.1472	0.0069814\\
2.1506	0.0071536\\
2.1541	0.0073301\\
2.1575	0.0075111\\
2.1609	0.0076965\\
2.1643	0.0078866\\
2.1678	0.0080815\\
2.1712	0.0082812\\
2.1746	0.0084859\\
2.178	0.0086957\\
2.1815	0.0089107\\
2.1849	0.0091311\\
2.1883	0.009357\\
2.1918	0.0095885\\
2.1952	0.0098257\\
2.1986	0.010069\\
2.202	0.010318\\
2.2055	0.010573\\
2.2089	0.010835\\
2.2123	0.011103\\
2.2157	0.011378\\
2.2192	0.011659\\
2.2226	0.011948\\
2.226	0.012243\\
2.2294	0.012546\\
2.2329	0.012856\\
2.2363	0.013174\\
2.2397	0.0135\\
2.2432	0.013833\\
2.2466	0.014175\\
2.25	0.014525\\
2.2534	0.014883\\
2.2569	0.015251\\
2.2603	0.015627\\
2.2637	0.016012\\
2.2671	0.016406\\
2.2706	0.01681\\
2.274	0.017224\\
2.2774	0.017648\\
2.2808	0.018082\\
2.2843	0.018526\\
2.2877	0.018981\\
2.2911	0.019446\\
2.2946	0.019923\\
2.298	0.020411\\
2.3014	0.020911\\
2.3048	0.021422\\
2.3083	0.021946\\
2.3117	0.022481\\
2.3151	0.02303\\
2.3185	0.023591\\
2.322	0.024165\\
2.3254	0.024753\\
2.3288	0.025354\\
2.3322	0.025969\\
2.3357	0.026598\\
2.3391	0.027242\\
2.3425	0.027901\\
2.346	0.028574\\
2.3494	0.029263\\
2.3528	0.029968\\
2.3562	0.030689\\
2.3597	0.031426\\
2.3631	0.03218\\
2.3665	0.03295\\
2.3699	0.033738\\
2.3734	0.034544\\
2.3768	0.035367\\
2.3802	0.036209\\
2.3837	0.037069\\
2.3871	0.037948\\
2.3905	0.038847\\
2.3939	0.039765\\
2.3974	0.040704\\
2.4008	0.041662\\
2.4042	0.042642\\
2.4076	0.043642\\
2.4111	0.044664\\
2.4145	0.045708\\
2.4179	0.046775\\
2.4213	0.047864\\
2.4248	0.048976\\
2.4282	0.050111\\
2.4316	0.05127\\
2.4351	0.052454\\
2.4385	0.053662\\
2.4419	0.054895\\
2.4453	0.056153\\
2.4488	0.057438\\
2.4522	0.058749\\
2.4556	0.060086\\
2.459	0.06145\\
2.4625	0.062842\\
2.4659	0.064262\\
2.4693	0.065711\\
2.4727	0.067188\\
2.4762	0.068694\\
2.4796	0.07023\\
2.483	0.071797\\
2.4865	0.073393\\
2.4899	0.075021\\
2.4933	0.07668\\
2.4967	0.078371\\
2.5002	0.080094\\
2.5036	0.08185\\
2.507	0.083639\\
2.5104	0.085462\\
2.5139	0.087319\\
2.5173	0.08921\\
2.5207	0.091136\\
2.5242	0.093098\\
2.5276	0.095095\\
2.531	0.097129\\
2.5344	0.099199\\
2.5379	0.10131\\
2.5413	0.10345\\
2.5447	0.10563\\
2.5481	0.10786\\
2.5516	0.11012\\
2.555	0.11242\\
2.5584	0.11475\\
2.5618	0.11713\\
2.5653	0.11955\\
2.5687	0.12201\\
2.5721	0.12452\\
2.5756	0.12706\\
2.579	0.12965\\
2.5824	0.13227\\
2.5858	0.13495\\
2.5893	0.13766\\
2.5927	0.14042\\
2.5961	0.14322\\
2.5995	0.14607\\
2.603	0.14896\\
2.6064	0.15189\\
2.6098	0.15488\\
2.6132	0.1579\\
2.6167	0.16098\\
2.6201	0.1641\\
2.6235	0.16726\\
2.627	0.17048\\
2.6304	0.17374\\
2.6338	0.17705\\
2.6372	0.18041\\
2.6407	0.18381\\
2.6441	0.18726\\
2.6475	0.19077\\
2.6509	0.19432\\
2.6544	0.19792\\
2.6578	0.20156\\
2.6612	0.20526\\
2.6646	0.20901\\
2.6681	0.21281\\
2.6715	0.21666\\
2.6749	0.22055\\
2.6784	0.2245\\
2.6818	0.2285\\
2.6852	0.23255\\
2.6886	0.23665\\
2.6921	0.2408\\
2.6955	0.245\\
2.6989	0.24925\\
2.7023	0.25355\\
2.7058	0.25791\\
2.7092	0.26231\\
2.7126	0.26676\\
2.7161	0.27127\\
2.7195	0.27582\\
2.7229	0.28043\\
2.7263	0.28508\\
2.7298	0.28979\\
2.7332	0.29454\\
2.7366	0.29935\\
2.74	0.30421\\
2.7435	0.30911\\
2.7469	0.31407\\
2.7503	0.31907\\
2.7537	0.32412\\
2.7572	0.32922\\
2.7606	0.33437\\
2.764	0.33957\\
2.7675	0.34481\\
2.7709	0.35011\\
2.7743	0.35544\\
2.7777	0.36083\\
2.7812	0.36626\\
2.7846	0.37174\\
2.788	0.37726\\
2.7914	0.38282\\
2.7949	0.38843\\
2.7983	0.39409\\
2.8017	0.39978\\
2.8051	0.40552\\
2.8086	0.4113\\
2.812	0.41712\\
2.8154	0.42298\\
2.8189	0.42889\\
2.8223	0.43483\\
2.8257	0.4408\\
2.8291	0.44682\\
2.8326	0.45287\\
2.836	0.45896\\
2.8394	0.46508\\
2.8428	0.47124\\
2.8463	0.47743\\
2.8497	0.48365\\
2.8531	0.4899\\
2.8565	0.49619\\
2.86	0.5025\\
2.8634	0.50884\\
2.8668	0.51521\\
2.8703	0.52161\\
2.8737	0.52803\\
2.8771	0.53447\\
2.8805	0.54094\\
2.884	0.54743\\
2.8874	0.55394\\
2.8908	0.56047\\
2.8942	0.56702\\
2.8977	0.57358\\
2.9011	0.58016\\
2.9045	0.58676\\
2.908	0.59337\\
2.9114	0.59998\\
2.9148	0.60661\\
2.9182	0.61325\\
2.9217	0.6199\\
2.9251	0.62655\\
2.9285	0.63321\\
2.9319	0.63987\\
2.9354	0.64654\\
2.9388	0.6532\\
2.9422	0.65987\\
2.9456	0.66653\\
2.9491	0.67319\\
2.9525	0.67984\\
2.9559	0.68649\\
2.9594	0.69313\\
2.9628	0.69976\\
2.9662	0.70638\\
2.9696	0.71299\\
2.9731	0.71958\\
2.9765	0.72616\\
2.9799	0.73273\\
2.9833	0.73927\\
2.9868	0.7458\\
2.9902	0.75231\\
2.9936	0.7588\\
2.997	0.76526\\
3.0005	0.77171\\
3.0039	0.77812\\
3.0073	0.78452\\
3.0108	0.79089\\
3.0142	0.79723\\
3.0176	0.80354\\
3.021	0.80983\\
3.0245	0.81608\\
3.0279	0.82231\\
3.0313	0.8285\\
3.0347	0.83467\\
3.0382	0.8408\\
3.0416	0.8469\\
3.045	0.85296\\
3.0484	0.85899\\
3.0519	0.86498\\
3.0553	0.87093\\
3.0587	0.87684\\
3.0622	0.88271\\
3.0656	0.88854\\
3.069	0.89431\\
3.0724	0.90004\\
3.0759	0.90572\\
3.0793	0.91134\\
3.0827	0.9169\\
3.0861	0.9224\\
3.0896	0.92784\\
3.093	0.93321\\
3.0964	0.93851\\
3.0999	0.94373\\
3.1033	0.94888\\
3.1067	0.95395\\
3.1101	0.95893\\
3.1136	0.96383\\
3.117	0.96864\\
3.1204	0.97335\\
3.1238	0.97798\\
3.1273	0.9825\\
3.1307	0.98693\\
3.1341	0.99125\\
3.1375	0.99546\\
3.141	0.99957\\
3.1444	1.0036\\
3.1478	1.0075\\
3.1513	1.0112\\
3.1547	1.0149\\
3.1581	1.0184\\
3.1615	1.0219\\
3.165	1.0252\\
3.1684	1.0283\\
3.1718	1.0314\\
3.1752	1.0343\\
3.1787	1.0371\\
3.1821	1.0398\\
3.1855	1.0424\\
3.1889	1.0448\\
3.1924	1.0471\\
3.1958	1.0493\\
3.1992	1.0513\\
3.2027	1.0533\\
3.2061	1.055\\
3.2095	1.0567\\
3.2129	1.0583\\
3.2164	1.0597\\
3.2198	1.061\\
3.2232	1.0621\\
3.2266	1.0632\\
3.2301	1.0641\\
3.2335	1.0649\\
3.2369	1.0655\\
3.2404	1.0661\\
3.2438	1.0665\\
3.2472	1.0668\\
3.2506	1.067\\
3.2541	1.0671\\
3.2575	1.0671\\
3.2609	1.0669\\
3.2643	1.0667\\
3.2678	1.0663\\
3.2712	1.0659\\
3.2746	1.0653\\
3.278	1.0646\\
3.2815	1.0638\\
3.2849	1.0629\\
3.2883	1.0619\\
3.2918	1.0608\\
3.2952	1.0596\\
3.2986	1.0584\\
3.302	1.057\\
3.3055	1.0555\\
3.3089	1.0539\\
3.3123	1.0522\\
3.3157	1.0505\\
3.3192	1.0486\\
3.3226	1.0467\\
3.326	1.0446\\
3.3294	1.0425\\
3.3329	1.0403\\
3.3363	1.038\\
3.3397	1.0356\\
3.3432	1.0331\\
3.3466	1.0306\\
3.35	1.0279\\
3.3534	1.0252\\
3.3569	1.0224\\
3.3603	1.0195\\
3.3637	1.0165\\
3.3671	1.0134\\
3.3706	1.0103\\
3.374	1.007\\
3.3774	1.0037\\
3.3808	1.0004\\
3.3843	0.99692\\
3.3877	0.9934\\
3.3911	0.9898\\
3.3946	0.98614\\
3.398	0.9824\\
3.4014	0.97859\\
3.4048	0.97472\\
3.4083	0.97078\\
3.4117	0.96677\\
3.4151	0.9627\\
3.4185	0.95857\\
3.422	0.95438\\
3.4254	0.95013\\
3.4288	0.94582\\
3.4323	0.94145\\
3.4357	0.93703\\
3.4391	0.93256\\
3.4425	0.92803\\
3.446	0.92346\\
3.4494	0.91883\\
3.4528	0.91416\\
3.4562	0.90944\\
3.4597	0.90467\\
3.4631	0.89987\\
3.4665	0.89502\\
3.4699	0.89013\\
3.4734	0.88521\\
3.4768	0.88024\\
3.4802	0.87524\\
3.4837	0.87021\\
3.4871	0.86514\\
3.4905	0.86004\\
3.4939	0.85491\\
3.4974	0.84976\\
3.5008	0.84457\\
3.5042	0.83936\\
3.5076	0.83413\\
3.5111	0.82887\\
3.5145	0.82359\\
3.5179	0.81829\\
3.5213	0.81297\\
3.5248	0.80764\\
3.5282	0.80228\\
3.5316	0.79692\\
3.5351	0.79154\\
3.5385	0.78614\\
3.5419	0.78074\\
3.5453	0.77532\\
3.5488	0.7699\\
3.5522	0.76447\\
3.5556	0.75904\\
3.559	0.75359\\
3.5625	0.74815\\
3.5659	0.7427\\
3.5693	0.73725\\
3.5727	0.7318\\
3.5762	0.72636\\
3.5796	0.72091\\
3.583	0.71547\\
3.5865	0.71003\\
3.5899	0.70459\\
3.5933	0.69916\\
3.5967	0.69374\\
3.6002	0.68833\\
3.6036	0.68293\\
3.607	0.67753\\
3.6104	0.67215\\
3.6139	0.66678\\
3.6173	0.66143\\
3.6207	0.65608\\
3.6242	0.65075\\
3.6276	0.64544\\
3.631	0.64015\\
3.6344	0.63487\\
3.6379	0.62961\\
3.6413	0.62437\\
3.6447	0.61914\\
3.6481	0.61394\\
3.6516	0.60876\\
3.655	0.6036\\
3.6584	0.59847\\
3.6618	0.59336\\
3.6653	0.58827\\
3.6687	0.5832\\
3.6721	0.57816\\
3.6756	0.57315\\
3.679	0.56816\\
3.6824	0.5632\\
3.6858	0.55827\\
3.6893	0.55337\\
3.6927	0.54849\\
3.6961	0.54364\\
3.6995	0.53883\\
3.703	0.53404\\
3.7064	0.52928\\
3.7098	0.52456\\
3.7132	0.51986\\
3.7167	0.5152\\
3.7201	0.51056\\
3.7235	0.50596\\
3.727	0.5014\\
3.7304	0.49686\\
3.7338	0.49236\\
3.7372	0.48789\\
3.7407	0.48346\\
3.7441	0.47906\\
3.7475	0.47469\\
3.7509	0.47036\\
3.7544	0.46606\\
3.7578	0.4618\\
3.7612	0.45757\\
3.7647	0.45338\\
3.7681	0.44922\\
3.7715	0.44509\\
3.7749	0.44101\\
3.7784	0.43695\\
3.7818	0.43294\\
3.7852	0.42895\\
3.7886	0.42501\\
3.7921	0.4211\\
3.7955	0.41722\\
3.7989	0.41338\\
3.8023	0.40958\\
3.8058	0.40581\\
3.8092	0.40208\\
3.8126	0.39838\\
3.8161	0.39472\\
3.8195	0.39109\\
3.8229	0.3875\\
3.8263	0.38394\\
3.8298	0.38042\\
3.8332	0.37693\\
3.8366	0.37348\\
3.84	0.37006\\
3.8435	0.36667\\
3.8469	0.36332\\
3.8503	0.36001\\
3.8537	0.35672\\
3.8572	0.35347\\
3.8606	0.35026\\
3.864	0.34707\\
3.8675	0.34392\\
3.8709	0.34081\\
3.8743	0.33772\\
3.8777	0.33467\\
3.8812	0.33165\\
3.8846	0.32866\\
3.888	0.3257\\
3.8914	0.32277\\
3.8949	0.31987\\
3.8983	0.31701\\
3.9017	0.31417\\
3.9051	0.31137\\
3.9086	0.30859\\
3.912	0.30584\\
3.9154	0.30312\\
3.9189	0.30043\\
3.9223	0.29777\\
3.9257	0.29514\\
3.9291	0.29253\\
3.9326	0.28995\\
3.936	0.2874\\
3.9394	0.28488\\
3.9428	0.28238\\
3.9463	0.27991\\
3.9497	0.27746\\
3.9531	0.27504\\
3.9566	0.27264\\
3.96	0.27027\\
3.9634	0.26792\\
3.9668	0.2656\\
3.9703	0.2633\\
3.9737	0.26103\\
3.9771	0.25877\\
3.9805	0.25654\\
3.984	0.25434\\
3.9874	0.25215\\
3.9908	0.24999\\
3.9942	0.24785\\
3.9977	0.24573\\
4.0011	0.24363\\
4.0045	0.24155\\
4.008	0.23949\\
4.0114	0.23745\\
4.0148	0.23544\\
4.0182	0.23344\\
4.0217	0.23146\\
4.0251	0.2295\\
4.0285	0.22756\\
4.0319	0.22563\\
4.0354	0.22373\\
4.0388	0.22184\\
4.0422	0.21997\\
4.0456	0.21812\\
4.0491	0.21628\\
4.0525	0.21446\\
4.0559	0.21266\\
4.0594	0.21087\\
4.0628	0.2091\\
4.0662	0.20734\\
4.0696	0.2056\\
4.0731	0.20388\\
4.0765	0.20217\\
4.0799	0.20048\\
4.0833	0.1988\\
4.0868	0.19713\\
4.0902	0.19548\\
4.0936	0.19384\\
4.097	0.19222\\
4.1005	0.19061\\
4.1039	0.18901\\
4.1073	0.18742\\
4.1108	0.18585\\
4.1142	0.18429\\
4.1176	0.18275\\
4.121	0.18121\\
4.1245	0.17969\\
4.1279	0.17818\\
4.1313	0.17668\\
4.1347	0.1752\\
4.1382	0.17372\\
4.1416	0.17226\\
4.145	0.1708\\
4.1485	0.16936\\
4.1519	0.16793\\
4.1553	0.16651\\
4.1587	0.1651\\
4.1622	0.1637\\
4.1656	0.16231\\
4.169	0.16093\\
4.1724	0.15956\\
4.1759	0.1582\\
4.1793	0.15685\\
4.1827	0.15551\\
4.1861	0.15417\\
4.1896	0.15285\\
4.193	0.15154\\
4.1964	0.15023\\
4.1999	0.14894\\
4.2033	0.14765\\
4.2067	0.14637\\
4.2101	0.14511\\
4.2136	0.14385\\
4.217	0.14259\\
4.2204	0.14135\\
4.2238	0.14011\\
4.2273	0.13889\\
4.2307	0.13767\\
4.2341	0.13646\\
4.2375	0.13525\\
4.241	0.13406\\
4.2444	0.13287\\
4.2478	0.13169\\
4.2513	0.13052\\
4.2547	0.12935\\
4.2581	0.12819\\
4.2615	0.12704\\
4.265	0.1259\\
4.2684	0.12477\\
4.2718	0.12364\\
4.2752	0.12252\\
4.2787	0.1214\\
4.2821	0.1203\\
4.2855	0.1192\\
4.2889	0.1181\\
4.2924	0.11702\\
4.2958	0.11594\\
4.2992	0.11487\\
4.3027	0.1138\\
4.3061	0.11274\\
4.3095	0.11169\\
4.3129	0.11065\\
4.3164	0.10961\\
4.3198	0.10858\\
4.3232	0.10755\\
4.3266	0.10653\\
4.3301	0.10552\\
4.3335	0.10452\\
4.3369	0.10352\\
4.3404	0.10252\\
4.3438	0.10154\\
4.3472	0.10056\\
4.3506	0.099583\\
4.3541	0.098616\\
4.3575	0.097655\\
4.3609	0.0967\\
4.3643	0.095752\\
4.3678	0.094809\\
4.3712	0.093873\\
4.3746	0.092943\\
4.378	0.09202\\
4.3815	0.091102\\
4.3849	0.090191\\
4.3883	0.089285\\
4.3918	0.088386\\
4.3952	0.087493\\
4.3986	0.086606\\
4.402	0.085724\\
4.4055	0.084849\\
4.4089	0.08398\\
4.4123	0.083117\\
4.4157	0.08226\\
4.4192	0.081409\\
4.4226	0.080564\\
4.426	0.079724\\
4.4294	0.078891\\
4.4329	0.078063\\
4.4363	0.077242\\
4.4397	0.076426\\
4.4432	0.075616\\
4.4466	0.074812\\
4.45	0.074014\\
4.4534	0.073221\\
4.4569	0.072434\\
4.4603	0.071653\\
4.4637	0.070878\\
4.4671	0.070109\\
4.4706	0.069345\\
4.474	0.068587\\
4.4774	0.067834\\
4.4809	0.067088\\
4.4843	0.066347\\
4.4877	0.065611\\
4.4911	0.064881\\
4.4946	0.064157\\
4.498	0.063438\\
4.5014	0.062725\\
4.5048	0.062018\\
4.5083	0.061316\\
4.5117	0.060619\\
4.5151	0.059928\\
4.5185	0.059243\\
4.522	0.058562\\
4.5254	0.057888\\
4.5288	0.057219\\
4.5323	0.056555\\
4.5357	0.055896\\
4.5391	0.055243\\
4.5425	0.054596\\
4.546	0.053953\\
4.5494	0.053316\\
4.5528	0.052684\\
4.5562	0.052058\\
4.5597	0.051437\\
4.5631	0.050821\\
4.5665	0.05021\\
4.5699	0.049605\\
4.5734	0.049004\\
4.5768	0.048409\\
4.5802	0.047819\\
4.5837	0.047234\\
4.5871	0.046655\\
4.5905	0.04608\\
4.5939	0.04551\\
4.5974	0.044946\\
4.6008	0.044386\\
4.6042	0.043832\\
4.6076	0.043283\\
4.6111	0.042738\\
4.6145	0.042199\\
4.6179	0.041664\\
4.6213	0.041135\\
4.6248	0.04061\\
4.6282	0.04009\\
4.6316	0.039575\\
4.6351	0.039065\\
4.6385	0.03856\\
4.6419	0.038059\\
4.6453	0.037564\\
4.6488	0.037073\\
4.6522	0.036587\\
4.6556	0.036105\\
4.659	0.035629\\
4.6625	0.035157\\
4.6659	0.034689\\
4.6693	0.034226\\
4.6728	0.033768\\
4.6762	0.033315\\
4.6796	0.032866\\
4.683	0.032421\\
4.6865	0.031982\\
4.6899	0.031546\\
4.6933	0.031115\\
4.6967	0.030689\\
4.7002	0.030267\\
4.7036	0.029849\\
4.707	0.029436\\
4.7104	0.029027\\
4.7139	0.028623\\
4.7173	0.028222\\
4.7207	0.027827\\
4.7242	0.027435\\
4.7276	0.027048\\
4.731	0.026664\\
4.7344	0.026285\\
4.7379	0.025911\\
4.7413	0.02554\\
4.7447	0.025173\\
4.7481	0.024811\\
4.7516	0.024452\\
4.755	0.024098\\
4.7584	0.023748\\
4.7618	0.023401\\
4.7653	0.023059\\
4.7687	0.02272\\
4.7721	0.022386\\
4.7756	0.022055\\
4.779	0.021728\\
4.7824	0.021405\\
4.7858	0.021086\\
4.7893	0.020771\\
4.7927	0.020459\\
4.7961	0.020151\\
4.7995	0.019847\\
4.803	0.019546\\
4.8064	0.019249\\
4.8098	0.018956\\
4.8132	0.018666\\
4.8167	0.01838\\
4.8201	0.018097\\
4.8235	0.017818\\
4.827	0.017542\\
4.8304	0.01727\\
4.8338	0.017001\\
4.8372	0.016736\\
4.8407	0.016473\\
4.8441	0.016215\\
4.8475	0.015959\\
4.8509	0.015707\\
4.8544	0.015458\\
4.8578	0.015212\\
4.8612	0.01497\\
4.8647	0.01473\\
4.8681	0.014494\\
4.8715	0.014261\\
4.8749	0.014031\\
4.8784	0.013804\\
4.8818	0.01358\\
4.8852	0.013359\\
4.8886	0.013141\\
4.8921	0.012926\\
4.8955	0.012714\\
4.8989	0.012505\\
4.9023	0.012298\\
4.9058	0.012095\\
4.9092	0.011894\\
4.9126	0.011696\\
4.9161	0.011501\\
4.9195	0.011308\\
4.9229	0.011119\\
4.9263	0.010931\\
4.9298	0.010747\\
4.9332	0.010565\\
4.9366	0.010386\\
4.94	0.010209\\
4.9435	0.010035\\
4.9469	0.0098631\\
4.9503	0.0096939\\
4.9537	0.0095271\\
4.9572	0.0093627\\
4.9606	0.0092008\\
4.964	0.0090412\\
4.9675	0.008884\\
4.9709	0.0087291\\
4.9743	0.0085765\\
4.9777	0.0084261\\
4.9812	0.008278\\
4.9846	0.0081322\\
4.988	0.0079885\\
4.9914	0.007847\\
4.9949	0.0077076\\
4.9983	0.0075704\\
5.0017	0.0074352\\
5.0051	0.0073021\\
5.0086	0.0071711\\
5.012	0.007042\\
5.0154	0.006915\\
5.0189	0.00679\\
5.0223	0.0066669\\
5.0257	0.0065457\\
5.0291	0.0064264\\
5.0326	0.006309\\
5.036	0.0061935\\
5.0394	0.0060797\\
5.0428	0.0059678\\
5.0463	0.0058577\\
5.0497	0.0057494\\
5.0531	0.0056427\\
5.0566	0.0055378\\
5.06	0.0054346\\
5.0634	0.0053331\\
5.0668	0.0052332\\
5.0703	0.005135\\
5.0737	0.0050384\\
5.0771	0.0049433\\
5.0805	0.0048498\\
5.084	0.0047579\\
5.0874	0.0046675\\
5.0908	0.0045786\\
5.0942	0.0044912\\
5.0977	0.0044052\\
5.1011	0.0043207\\
5.1045	0.0042376\\
5.108	0.0041559\\
5.1114	0.0040757\\
5.1148	0.0039967\\
5.1182	0.0039191\\
5.1217	0.0038429\\
5.1251	0.003768\\
5.1285	0.0036943\\
5.1319	0.0036219\\
5.1354	0.0035508\\
5.1388	0.0034809\\
5.1422	0.0034122\\
5.1456	0.0033448\\
5.1491	0.0032785\\
5.1525	0.0032134\\
5.1559	0.0031494\\
5.1594	0.0030866\\
5.1628	0.0030248\\
5.1662	0.0029642\\
5.1696	0.0029047\\
5.1731	0.0028462\\
5.1765	0.0027888\\
5.1799	0.0027324\\
5.1833	0.002677\\
5.1868	0.0026226\\
5.1902	0.0025692\\
5.1936	0.0025168\\
5.1971	0.0024654\\
5.2005	0.0024149\\
5.2039	0.0023653\\
5.2073	0.0023166\\
5.2108	0.0022688\\
5.2142	0.0022219\\
5.2176	0.0021759\\
5.221	0.0021307\\
5.2245	0.0020864\\
5.2279	0.0020429\\
5.2313	0.0020003\\
5.2347	0.0019584\\
5.2382	0.0019173\\
5.2416	0.001877\\
5.245	0.0018374\\
5.2485	0.0017987\\
5.2519	0.0017606\\
5.2553	0.0017233\\
5.2587	0.0016867\\
5.2622	0.0016508\\
5.2656	0.0016156\\
5.269	0.001581\\
5.2724	0.0015472\\
5.2759	0.001514\\
5.2793	0.0014814\\
5.2827	0.0014495\\
5.2861	0.0014182\\
5.2896	0.0013875\\
5.293	0.0013574\\
5.2964	0.0013279\\
5.2999	0.001299\\
5.3033	0.0012707\\
5.3067	0.0012429\\
5.3101	0.0012157\\
5.3101	0.001076\\
5.3067	0.0011009\\
5.3033	0.0011263\\
5.2999	0.0011522\\
5.2964	0.0011787\\
5.293	0.0012057\\
5.2896	0.0012333\\
5.2861	0.0012614\\
5.2827	0.0012901\\
5.2793	0.0013195\\
5.2759	0.0013494\\
5.2724	0.0013799\\
5.269	0.0014111\\
5.2656	0.0014428\\
5.2622	0.0014753\\
5.2587	0.0015084\\
5.2553	0.0015421\\
5.2519	0.0015765\\
5.2485	0.0016117\\
5.245	0.0016475\\
5.2416	0.001684\\
5.2382	0.0017213\\
5.2347	0.0017593\\
5.2313	0.0017981\\
5.2279	0.0018376\\
5.2245	0.0018779\\
5.221	0.001919\\
5.2176	0.0019609\\
5.2142	0.0020037\\
5.2108	0.0020472\\
5.2073	0.0020916\\
5.2039	0.0021369\\
5.2005	0.002183\\
5.1971	0.0022301\\
5.1936	0.002278\\
5.1902	0.0023268\\
5.1868	0.0023766\\
5.1833	0.0024273\\
5.1799	0.002479\\
5.1765	0.0025317\\
5.1731	0.0025853\\
5.1696	0.00264\\
5.1662	0.0026957\\
5.1628	0.0027524\\
5.1594	0.0028102\\
5.1559	0.0028691\\
5.1525	0.002929\\
5.1491	0.0029901\\
5.1456	0.0030523\\
5.1422	0.0031156\\
5.1388	0.0031801\\
5.1354	0.0032458\\
5.1319	0.0033126\\
5.1285	0.0033807\\
5.1251	0.00345\\
5.1217	0.0035205\\
5.1182	0.0035923\\
5.1148	0.0036654\\
5.1114	0.0037399\\
5.108	0.0038156\\
5.1045	0.0038927\\
5.1011	0.0039711\\
5.0977	0.0040509\\
5.0942	0.0041321\\
5.0908	0.0042148\\
5.0874	0.0042989\\
5.084	0.0043844\\
5.0805	0.0044714\\
5.0771	0.00456\\
5.0737	0.00465\\
5.0703	0.0047417\\
5.0668	0.0048348\\
5.0634	0.0049296\\
5.06	0.005026\\
5.0566	0.005124\\
5.0531	0.0052236\\
5.0497	0.005325\\
5.0463	0.005428\\
5.0428	0.0055328\\
5.0394	0.0056393\\
5.036	0.0057476\\
5.0326	0.0058577\\
5.0291	0.0059695\\
5.0257	0.0060833\\
5.0223	0.0061988\\
5.0189	0.0063163\\
5.0154	0.0064357\\
5.012	0.006557\\
5.0086	0.0066803\\
5.0051	0.0068055\\
5.0017	0.0069328\\
4.9983	0.0070621\\
4.9949	0.0071934\\
4.9914	0.0073269\\
4.988	0.0074624\\
4.9846	0.0076001\\
4.9812	0.0077399\\
4.9777	0.0078819\\
4.9743	0.0080262\\
4.9709	0.0081726\\
4.9675	0.0083214\\
4.964	0.0084724\\
4.9606	0.0086257\\
4.9572	0.0087814\\
4.9537	0.0089394\\
4.9503	0.0090999\\
4.9469	0.0092627\\
4.9435	0.0094281\\
4.94	0.0095959\\
4.9366	0.0097662\\
4.9332	0.009939\\
4.9298	0.010114\\
4.9263	0.010292\\
4.9229	0.010473\\
4.9195	0.010656\\
4.9161	0.010842\\
4.9126	0.011031\\
4.9092	0.011222\\
4.9058	0.011416\\
4.9023	0.011613\\
4.8989	0.011813\\
4.8955	0.012015\\
4.8921	0.012221\\
4.8886	0.012429\\
4.8852	0.01264\\
4.8818	0.012854\\
4.8784	0.013072\\
4.8749	0.013292\\
4.8715	0.013515\\
4.8681	0.013741\\
4.8647	0.013971\\
4.8612	0.014203\\
4.8578	0.014439\\
4.8544	0.014678\\
4.8509	0.01492\\
4.8475	0.015165\\
4.8441	0.015414\\
4.8407	0.015666\\
4.8372	0.015921\\
4.8338	0.01618\\
4.8304	0.016442\\
4.827	0.016708\\
4.8235	0.016976\\
4.8201	0.017249\\
4.8167	0.017525\\
4.8132	0.017804\\
4.8098	0.018087\\
4.8064	0.018374\\
4.803	0.018664\\
4.7995	0.018958\\
4.7961	0.019256\\
4.7927	0.019557\\
4.7893	0.019862\\
4.7858	0.020171\\
4.7824	0.020483\\
4.779	0.0208\\
4.7756	0.02112\\
4.7721	0.021444\\
4.7687	0.021772\\
4.7653	0.022104\\
4.7618	0.02244\\
4.7584	0.02278\\
4.755	0.023124\\
4.7516	0.023472\\
4.7481	0.023824\\
4.7447	0.024181\\
4.7413	0.024541\\
4.7379	0.024905\\
4.7344	0.025274\\
4.731	0.025647\\
4.7276	0.026024\\
4.7242	0.026405\\
4.7207	0.026791\\
4.7173	0.027181\\
4.7139	0.027575\\
4.7104	0.027974\\
4.707	0.028377\\
4.7036	0.028785\\
4.7002	0.029196\\
4.6967	0.029613\\
4.6933	0.030034\\
4.6899	0.030459\\
4.6865	0.030889\\
4.683	0.031323\\
4.6796	0.031762\\
4.6762	0.032206\\
4.6728	0.032654\\
4.6693	0.033107\\
4.6659	0.033564\\
4.6625	0.034026\\
4.659	0.034493\\
4.6556	0.034965\\
4.6522	0.035441\\
4.6488	0.035922\\
4.6453	0.036408\\
4.6419	0.036899\\
4.6385	0.037394\\
4.6351	0.037894\\
4.6316	0.038399\\
4.6282	0.038909\\
4.6248	0.039424\\
4.6213	0.039944\\
4.6179	0.040469\\
4.6145	0.040998\\
4.6111	0.041533\\
4.6076	0.042072\\
4.6042	0.042617\\
4.6008	0.043166\\
4.5974	0.043721\\
4.5939	0.04428\\
4.5905	0.044845\\
4.5871	0.045414\\
4.5837	0.045989\\
4.5802	0.046569\\
4.5768	0.047153\\
4.5734	0.047743\\
4.5699	0.048338\\
4.5665	0.048938\\
4.5631	0.049543\\
4.5597	0.050153\\
4.5562	0.050769\\
4.5528	0.051389\\
4.5494	0.052015\\
4.546	0.052645\\
4.5425	0.053281\\
4.5391	0.053922\\
4.5357	0.054568\\
4.5323	0.05522\\
4.5288	0.055876\\
4.5254	0.056538\\
4.522	0.057205\\
4.5185	0.057877\\
4.5151	0.058554\\
4.5117	0.059237\\
4.5083	0.059924\\
4.5048	0.060617\\
4.5014	0.061315\\
4.498	0.062019\\
4.4946	0.062727\\
4.4911	0.063441\\
4.4877	0.06416\\
4.4843	0.064884\\
4.4809	0.065613\\
4.4774	0.066348\\
4.474	0.067088\\
4.4706	0.067833\\
4.4671	0.068583\\
4.4637	0.069339\\
4.4603	0.0701\\
4.4569	0.070866\\
4.4534	0.071637\\
4.45	0.072414\\
4.4466	0.073195\\
4.4432	0.073983\\
4.4397	0.074775\\
4.4363	0.075573\\
4.4329	0.076376\\
4.4294	0.077184\\
4.426	0.077997\\
4.4226	0.078816\\
4.4192	0.079641\\
4.4157	0.08047\\
4.4123	0.081305\\
4.4089	0.082145\\
4.4055	0.082991\\
4.402	0.083842\\
4.3986	0.084698\\
4.3952	0.08556\\
4.3918	0.086427\\
4.3883	0.0873\\
4.3849	0.088178\\
4.3815	0.089062\\
4.378	0.089951\\
4.3746	0.090846\\
4.3712	0.091746\\
4.3678	0.092652\\
4.3643	0.093563\\
4.3609	0.09448\\
4.3575	0.095403\\
4.3541	0.096331\\
4.3506	0.097265\\
4.3472	0.098205\\
4.3438	0.099151\\
4.3404	0.1001\\
4.3369	0.10106\\
4.3335	0.10202\\
4.3301	0.10299\\
4.3266	0.10397\\
4.3232	0.10495\\
4.3198	0.10593\\
4.3164	0.10693\\
4.3129	0.10793\\
4.3095	0.10893\\
4.3061	0.10994\\
4.3027	0.11096\\
4.2992	0.11198\\
4.2958	0.11301\\
4.2924	0.11405\\
4.2889	0.11509\\
4.2855	0.11614\\
4.2821	0.1172\\
4.2787	0.11826\\
4.2752	0.11933\\
4.2718	0.12041\\
4.2684	0.12149\\
4.265	0.12258\\
4.2615	0.12368\\
4.2581	0.12478\\
4.2547	0.1259\\
4.2513	0.12701\\
4.2478	0.12814\\
4.2444	0.12927\\
4.241	0.13041\\
4.2375	0.13156\\
4.2341	0.13272\\
4.2307	0.13388\\
4.2273	0.13505\\
4.2238	0.13623\\
4.2204	0.13742\\
4.217	0.13862\\
4.2136	0.13982\\
4.2101	0.14103\\
4.2067	0.14226\\
4.2033	0.14349\\
4.1999	0.14472\\
4.1964	0.14597\\
4.193	0.14723\\
4.1896	0.14849\\
4.1861	0.14977\\
4.1827	0.15105\\
4.1793	0.15234\\
4.1759	0.15365\\
4.1724	0.15496\\
4.169	0.15628\\
4.1656	0.15762\\
4.1622	0.15896\\
4.1587	0.16031\\
4.1553	0.16168\\
4.1519	0.16305\\
4.1485	0.16444\\
4.145	0.16583\\
4.1416	0.16724\\
4.1382	0.16866\\
4.1347	0.17009\\
4.1313	0.17154\\
4.1279	0.17299\\
4.1245	0.17446\\
4.121	0.17594\\
4.1176	0.17743\\
4.1142	0.17893\\
4.1108	0.18045\\
4.1073	0.18198\\
4.1039	0.18352\\
4.1005	0.18508\\
4.097	0.18665\\
4.0936	0.18824\\
4.0902	0.18984\\
4.0868	0.19145\\
4.0833	0.19308\\
4.0799	0.19472\\
4.0765	0.19638\\
4.0731	0.19805\\
4.0696	0.19974\\
4.0662	0.20145\\
4.0628	0.20317\\
4.0594	0.2049\\
4.0559	0.20666\\
4.0525	0.20843\\
4.0491	0.21021\\
4.0456	0.21202\\
4.0422	0.21384\\
4.0388	0.21568\\
4.0354	0.21754\\
4.0319	0.21942\\
4.0285	0.22131\\
4.0251	0.22322\\
4.0217	0.22516\\
4.0182	0.22711\\
4.0148	0.22908\\
4.0114	0.23107\\
4.008	0.23308\\
4.0045	0.23512\\
4.0011	0.23717\\
3.9977	0.23924\\
3.9942	0.24134\\
3.9908	0.24345\\
3.9874	0.24559\\
3.984	0.24775\\
3.9805	0.24994\\
3.9771	0.25214\\
3.9737	0.25437\\
3.9703	0.25662\\
3.9668	0.2589\\
3.9634	0.2612\\
3.96	0.26352\\
3.9566	0.26587\\
3.9531	0.26824\\
3.9497	0.27063\\
3.9463	0.27306\\
3.9428	0.2755\\
3.9394	0.27798\\
3.936	0.28048\\
3.9326	0.283\\
3.9291	0.28555\\
3.9257	0.28813\\
3.9223	0.29074\\
3.9189	0.29337\\
3.9154	0.29603\\
3.912	0.29872\\
3.9086	0.30143\\
3.9051	0.30418\\
3.9017	0.30695\\
3.8983	0.30975\\
3.8949	0.31258\\
3.8914	0.31544\\
3.888	0.31833\\
3.8846	0.32125\\
3.8812	0.3242\\
3.8777	0.32718\\
3.8743	0.33019\\
3.8709	0.33324\\
3.8675	0.33631\\
3.864	0.33941\\
3.8606	0.34255\\
3.8572	0.34571\\
3.8537	0.34891\\
3.8503	0.35214\\
3.8469	0.3554\\
3.8435	0.3587\\
3.84	0.36202\\
3.8366	0.36538\\
3.8332	0.36878\\
3.8298	0.3722\\
3.8263	0.37566\\
3.8229	0.37915\\
3.8195	0.38268\\
3.8161	0.38624\\
3.8126	0.38983\\
3.8092	0.39346\\
3.8058	0.39712\\
3.8023	0.40081\\
3.7989	0.40454\\
3.7955	0.4083\\
3.7921	0.4121\\
3.7886	0.41593\\
3.7852	0.4198\\
3.7818	0.4237\\
3.7784	0.42763\\
3.7749	0.4316\\
3.7715	0.4356\\
3.7681	0.43964\\
3.7647	0.44371\\
3.7612	0.44782\\
3.7578	0.45196\\
3.7544	0.45614\\
3.7509	0.46035\\
3.7475	0.46459\\
3.7441	0.46887\\
3.7407	0.47318\\
3.7372	0.47752\\
3.7338	0.4819\\
3.7304	0.48631\\
3.727	0.49076\\
3.7235	0.49524\\
3.7201	0.49975\\
3.7167	0.50429\\
3.7132	0.50886\\
3.7098	0.51347\\
3.7064	0.51811\\
3.703	0.52278\\
3.6995	0.52748\\
3.6961	0.53221\\
3.6927	0.53697\\
3.6893	0.54176\\
3.6858	0.54658\\
3.6824	0.55143\\
3.679	0.55631\\
3.6756	0.56121\\
3.6721	0.56615\\
3.6687	0.57111\\
3.6653	0.57609\\
3.6618	0.58111\\
3.6584	0.58614\\
3.655	0.5912\\
3.6516	0.59629\\
3.6481	0.6014\\
3.6447	0.60653\\
3.6413	0.61168\\
3.6379	0.61686\\
3.6344	0.62205\\
3.631	0.62727\\
3.6276	0.6325\\
3.6242	0.63775\\
3.6207	0.64302\\
3.6173	0.64831\\
3.6139	0.65361\\
3.6104	0.65892\\
3.607	0.66425\\
3.6036	0.66959\\
3.6002	0.67495\\
3.5967	0.68031\\
3.5933	0.68568\\
3.5899	0.69107\\
3.5865	0.69646\\
3.583	0.70185\\
3.5796	0.70725\\
3.5762	0.71266\\
3.5727	0.71807\\
3.5693	0.72348\\
3.5659	0.72889\\
3.5625	0.7343\\
3.559	0.73972\\
3.5556	0.74512\\
3.5522	0.75053\\
3.5488	0.75593\\
3.5453	0.76132\\
3.5419	0.7667\\
3.5385	0.77208\\
3.5351	0.77744\\
3.5316	0.7828\\
3.5282	0.78814\\
3.5248	0.79346\\
3.5213	0.79877\\
3.5179	0.80407\\
3.5145	0.80934\\
3.5111	0.8146\\
3.5076	0.81983\\
3.5042	0.82504\\
3.5008	0.83023\\
3.4974	0.83539\\
3.4939	0.84053\\
3.4905	0.84564\\
3.4871	0.85071\\
3.4837	0.85576\\
3.4802	0.86077\\
3.4768	0.86575\\
3.4734	0.8707\\
3.4699	0.8756\\
3.4665	0.88047\\
3.4631	0.8853\\
3.4597	0.89009\\
3.4562	0.89484\\
3.4528	0.89954\\
3.4494	0.9042\\
3.446	0.90881\\
3.4425	0.91337\\
3.4391	0.91788\\
3.4357	0.92235\\
3.4323	0.92676\\
3.4288	0.93111\\
3.4254	0.93541\\
3.422	0.93966\\
3.4185	0.94384\\
3.4151	0.94797\\
3.4117	0.95204\\
3.4083	0.95604\\
3.4048	0.95998\\
3.4014	0.96386\\
3.398	0.96767\\
3.3946	0.97141\\
3.3911	0.97508\\
3.3877	0.97869\\
3.3843	0.98222\\
3.3808	0.98568\\
3.3774	0.98906\\
3.374	0.99237\\
3.3706	0.9956\\
3.3671	0.99874\\
3.3637	1.0018\\
3.3603	1.0048\\
3.3569	1.0077\\
3.3534	1.0105\\
3.35	1.0133\\
3.3466	1.0159\\
3.3432	1.0184\\
3.3397	1.0209\\
3.3363	1.0233\\
3.3329	1.0255\\
3.3294	1.0277\\
3.326	1.0298\\
3.3226	1.0318\\
3.3192	1.0336\\
3.3157	1.0354\\
3.3123	1.0371\\
3.3089	1.0386\\
3.3055	1.0401\\
3.302	1.0414\\
3.2986	1.0427\\
3.2952	1.0438\\
3.2918	1.0448\\
3.2883	1.0457\\
3.2849	1.0465\\
3.2815	1.0472\\
3.278	1.0478\\
3.2746	1.0482\\
3.2712	1.0486\\
3.2678	1.0488\\
3.2643	1.0489\\
3.2609	1.0489\\
3.2575	1.0488\\
3.2541	1.0485\\
3.2506	1.0482\\
3.2472	1.0477\\
3.2438	1.0471\\
3.2404	1.0464\\
3.2369	1.0456\\
3.2335	1.0447\\
3.2301	1.0437\\
3.2266	1.0425\\
3.2232	1.0413\\
3.2198	1.0399\\
3.2164	1.0384\\
3.2129	1.0368\\
3.2095	1.0351\\
3.2061	1.0333\\
3.2027	1.0314\\
3.1992	1.0294\\
3.1958	1.0273\\
3.1924	1.025\\
3.1889	1.0227\\
3.1855	1.0203\\
3.1821	1.0177\\
3.1787	1.0151\\
3.1752	1.0123\\
3.1718	1.0095\\
3.1684	1.0065\\
3.165	1.0035\\
3.1615	1.0004\\
3.1581	0.99712\\
3.1547	0.99379\\
3.1513	0.99036\\
3.1478	0.98684\\
3.1444	0.98322\\
3.141	0.97951\\
3.1375	0.9757\\
3.1341	0.97181\\
3.1307	0.96783\\
3.1273	0.96375\\
3.1238	0.95959\\
3.1204	0.95533\\
3.117	0.951\\
3.1136	0.94657\\
3.1101	0.94206\\
3.1067	0.93746\\
3.1033	0.93277\\
3.0999	0.92801\\
3.0964	0.92316\\
3.093	0.91822\\
3.0896	0.91321\\
3.0861	0.90811\\
3.0827	0.90294\\
3.0793	0.89768\\
3.0759	0.89235\\
3.0724	0.88694\\
3.069	0.88145\\
3.0656	0.87589\\
3.0622	0.87025\\
3.0587	0.86454\\
3.0553	0.85876\\
3.0519	0.85292\\
3.0484	0.84701\\
3.045	0.84103\\
3.0416	0.83499\\
3.0382	0.8289\\
3.0347	0.82275\\
3.0313	0.81655\\
3.0279	0.81029\\
3.0245	0.80399\\
3.021	0.79765\\
3.0176	0.79126\\
3.0142	0.78484\\
3.0108	0.77838\\
3.0073	0.77189\\
3.0039	0.76537\\
3.0005	0.75882\\
2.997	0.75225\\
2.9936	0.74565\\
2.9902	0.73904\\
2.9868	0.73241\\
2.9833	0.72576\\
2.9799	0.71909\\
2.9765	0.71242\\
2.9731	0.70573\\
2.9696	0.69904\\
2.9662	0.69234\\
2.9628	0.68564\\
2.9594	0.67893\\
2.9559	0.67222\\
2.9525	0.66552\\
2.9491	0.65881\\
2.9456	0.6521\\
2.9422	0.6454\\
2.9388	0.63871\\
2.9354	0.63202\\
2.9319	0.62534\\
2.9285	0.61867\\
2.9251	0.61202\\
2.9217	0.60537\\
2.9182	0.59873\\
2.9148	0.59211\\
2.9114	0.58551\\
2.908	0.57892\\
2.9045	0.57235\\
2.9011	0.5658\\
2.8977	0.55926\\
2.8942	0.55275\\
2.8908	0.54626\\
2.8874	0.53979\\
2.884	0.53334\\
2.8805	0.52692\\
2.8771	0.52053\\
2.8737	0.51416\\
2.8703	0.50781\\
2.8668	0.5015\\
2.8634	0.49522\\
2.86	0.48896\\
2.8565	0.48274\\
2.8531	0.47654\\
2.8497	0.47038\\
2.8463	0.46425\\
2.8428	0.45816\\
2.8394	0.4521\\
2.836	0.44608\\
2.8326	0.44009\\
2.8291	0.43414\\
2.8257	0.42822\\
2.8223	0.42235\\
2.8189	0.41651\\
2.8154	0.41071\\
2.812	0.40496\\
2.8086	0.39924\\
2.8051	0.39357\\
2.8017	0.38793\\
2.7983	0.38234\\
2.7949	0.3768\\
2.7914	0.37129\\
2.788	0.36583\\
2.7846	0.36042\\
2.7812	0.35505\\
2.7777	0.34972\\
2.7743	0.34444\\
2.7709	0.33921\\
2.7675	0.33402\\
2.764	0.32888\\
2.7606	0.32379\\
2.7572	0.31874\\
2.7537	0.31374\\
2.7503	0.30879\\
2.7469	0.30389\\
2.7435	0.29904\\
2.74	0.29424\\
2.7366	0.28948\\
2.7332	0.28478\\
2.7298	0.28012\\
2.7263	0.27551\\
2.7229	0.27095\\
2.7195	0.26645\\
2.7161	0.26199\\
2.7126	0.25758\\
2.7092	0.25322\\
2.7058	0.24891\\
2.7023	0.24465\\
2.6989	0.24045\\
2.6955	0.23629\\
2.6921	0.23218\\
2.6886	0.22812\\
2.6852	0.22411\\
2.6818	0.22015\\
2.6784	0.21624\\
2.6749	0.21238\\
2.6715	0.20857\\
2.6681	0.20481\\
2.6646	0.20109\\
2.6612	0.19743\\
2.6578	0.19382\\
2.6544	0.19025\\
2.6509	0.18673\\
2.6475	0.18326\\
2.6441	0.17984\\
2.6407	0.17646\\
2.6372	0.17314\\
2.6338	0.16986\\
2.6304	0.16662\\
2.627	0.16344\\
2.6235	0.1603\\
2.6201	0.15721\\
2.6167	0.15416\\
2.6132	0.15116\\
2.6098	0.1482\\
2.6064	0.14529\\
2.603	0.14242\\
2.5995	0.1396\\
2.5961	0.13682\\
2.5927	0.13408\\
2.5893	0.13139\\
2.5858	0.12874\\
2.5824	0.12613\\
2.579	0.12357\\
2.5756	0.12104\\
2.5721	0.11856\\
2.5687	0.11612\\
2.5653	0.11372\\
2.5618	0.11136\\
2.5584	0.10903\\
2.555	0.10675\\
2.5516	0.10451\\
2.5481	0.1023\\
2.5447	0.10014\\
2.5413	0.098009\\
2.5379	0.095918\\
2.5344	0.093863\\
2.531	0.091844\\
2.5276	0.089862\\
2.5242	0.087915\\
2.5207	0.086003\\
2.5173	0.084126\\
2.5139	0.082283\\
2.5104	0.080474\\
2.507	0.078699\\
2.5036	0.076956\\
2.5002	0.075247\\
2.4967	0.073569\\
2.4933	0.071923\\
2.4899	0.070308\\
2.4865	0.068724\\
2.483	0.06717\\
2.4796	0.065647\\
2.4762	0.064153\\
2.4727	0.062688\\
2.4693	0.061253\\
2.4659	0.059845\\
2.4625	0.058465\\
2.459	0.057113\\
2.4556	0.055788\\
2.4522	0.05449\\
2.4488	0.053218\\
2.4453	0.051972\\
2.4419	0.050751\\
2.4385	0.049555\\
2.4351	0.048384\\
2.4316	0.047238\\
2.4282	0.046115\\
2.4248	0.045016\\
2.4213	0.043939\\
2.4179	0.042886\\
2.4145	0.041855\\
2.4111	0.040846\\
2.4076	0.039859\\
2.4042	0.038892\\
2.4008	0.037947\\
2.3974	0.037022\\
2.3939	0.036118\\
2.3905	0.035233\\
2.3871	0.034367\\
2.3837	0.033521\\
2.3802	0.032693\\
2.3768	0.031884\\
2.3734	0.031093\\
2.3699	0.03032\\
2.3665	0.029564\\
2.3631	0.028825\\
2.3597	0.028103\\
2.3562	0.027397\\
2.3528	0.026708\\
2.3494	0.026034\\
2.346	0.025376\\
2.3425	0.024733\\
2.3391	0.024105\\
2.3357	0.023492\\
2.3322	0.022893\\
2.3288	0.022308\\
2.3254	0.021737\\
2.322	0.021179\\
2.3185	0.020635\\
2.3151	0.020103\\
2.3117	0.019584\\
2.3083	0.019078\\
2.3048	0.018584\\
2.3014	0.018101\\
2.298	0.017631\\
2.2946	0.017172\\
2.2911	0.016724\\
2.2877	0.016286\\
2.2843	0.01586\\
2.2808	0.015444\\
2.2774	0.015038\\
2.274	0.014643\\
2.2706	0.014257\\
2.2671	0.01388\\
2.2637	0.013513\\
2.2603	0.013156\\
2.2569	0.012807\\
2.2534	0.012467\\
2.25	0.012135\\
2.2466	0.011812\\
2.2432	0.011497\\
2.2397	0.01119\\
2.2363	0.010891\\
2.2329	0.0106\\
2.2294	0.010316\\
2.226	0.010039\\
2.2226	0.0097695\\
2.2192	0.0095068\\
2.2157	0.0092509\\
2.2123	0.0090016\\
2.2089	0.0087588\\
2.2055	0.0085223\\
2.202	0.008292\\
2.1986	0.0080677\\
2.1952	0.0078492\\
2.1918	0.0076365\\
2.1883	0.0074294\\
2.1849	0.0072277\\
2.1815	0.0070313\\
2.178	0.0068401\\
2.1746	0.006654\\
2.1712	0.0064728\\
2.1678	0.0062965\\
2.1643	0.0061248\\
2.1609	0.0059577\\
2.1575	0.0057951\\
2.1541	0.0056368\\
2.1506	0.0054827\\
2.1472	0.0053328\\
2.1438	0.005187\\
2.1403	0.005045\\
2.1369	0.0049069\\
2.1335	0.0047725\\
2.1301	0.0046417\\
2.1266	0.0045145\\
2.1232	0.0043907\\
2.1198	0.0042703\\
2.1164	0.0041532\\
2.1129	0.0040392\\
2.1095	0.0039284\\
2.1061	0.0038206\\
2.1027	0.0037157\\
2.0992	0.0036137\\
2.0958	0.0035145\\
2.0924	0.003418\\
2.0889	0.0033241\\
2.0855	0.0032328\\
2.0821	0.003144\\
2.0787	0.0030577\\
2.0752	0.0029738\\
2.0718	0.0028921\\
2.0684	0.0028127\\
2.065	0.0027355\\
2.0615	0.0026604\\
2.0581	0.0025874\\
2.0547	0.0025164\\
2.0513	0.0024474\\
2.0478	0.0023803\\
2.0444	0.0023151\\
2.041	0.0022516\\
2.0375	0.0021899\\
2.0341	0.0021299\\
2.0307	0.0020716\\
2.0273	0.0020149\\
2.0238	0.0019598\\
2.0204	0.0019062\\
2.017	0.0018541\\
2.0136	0.0018034\\
2.0101	0.0017541\\
2.0067	0.0017062\\
2.0033	0.0016597\\
1.9999	0.0016144\\
1.9964	0.0015704\\
1.993	0.0015276\\
1.9896	0.001486\\
1.9861	0.0014455\\
1.9827	0.0014062\\
1.9793	0.001368\\
1.9759	0.0013308\\
1.9724	0.0012947\\
1.969	0.0012595\\
1.9656	0.0012254\\
1.9622	0.0011921\\
1.9587	0.0011598\\
1.9553	0.0011284\\
1.9519	0.0010979\\
1.9484	0.0010682\\
1.945	0.0010394\\
1.9416	0.0010113\\
1.9382	0.00098401\\
1.9347	0.00095747\\
1.9313	0.00093167\\
1.9279	0.00090658\\
1.9245	0.00088218\\
1.921	0.00085846\\
1.9176	0.00083539\\
1.9142	0.00081295\\
1.9108	0.00079113\\
1.9073	0.00076992\\
1.9039	0.00074928\\
1.9005	0.00072921\\
1.897	0.0007097\\
1.8936	0.00069072\\
1.8902	0.00067225\\
1.8868	0.0006543\\
}--cycle;
\addplot[ybar interval, fill=black, fill opacity=0.35, area legend, draw=none] table[row sep=crcr, x=Lower, y=Count] {%
Lower	Upper	Count\\
1.8868	2.1501	0.020253\\
2.1501	2.4134	0.091138\\
2.4134	2.6768	0.17215\\
2.6768	2.9401	0.5367\\
2.9401	3.2034	0.90126\\
3.2034	3.4668	0.85062\\
3.4668	3.7301	0.63797\\
3.7301	3.9935	0.32405\\
3.9935	4.2568	0.14177\\
4.2568	4.5201	0.060759\\
4.5201	4.7835	0.030379\\
4.7835	5.0468	0.010126\\
5.0468	5.3101	0.020253\\
5.3101	5.3101	0.020253\\
};
\addlegendentry{\RealHistogramLegend}

\addplot [color=black]
  table[row sep=crcr]{%
1.8868	0.0040946\\
1.945	0.0064839\\
1.993	0.0093236\\
2.0341	0.01259\\
2.0718	0.016433\\
2.1061	0.020781\\
2.1403	0.026095\\
2.1712	0.031837\\
2.1986	0.037812\\
2.226	0.044707\\
2.2534	0.05262\\
2.2774	0.06046\\
2.3014	0.069229\\
2.3254	0.078996\\
2.3494	0.08983\\
2.3734	0.1018\\
2.3939	0.113\\
2.4145	0.12513\\
2.4351	0.1382\\
2.4556	0.15225\\
2.4762	0.1673\\
2.4967	0.18338\\
2.5173	0.20049\\
2.5379	0.21865\\
2.5618	0.24114\\
2.5858	0.26502\\
2.6098	0.29027\\
2.6338	0.31683\\
2.6612	0.34869\\
2.6886	0.38204\\
2.7195	0.4211\\
2.7537	0.46607\\
2.7983	0.52624\\
2.8942	0.65649\\
2.9251	0.69667\\
2.9525	0.73095\\
2.9765	0.75953\\
2.997	0.78278\\
3.0176	0.80471\\
3.0382	0.82518\\
3.0553	0.84102\\
3.0724	0.85567\\
3.0896	0.86906\\
3.1067	0.88113\\
3.1204	0.88979\\
3.1341	0.89754\\
3.1478	0.90436\\
3.1615	0.91021\\
3.1752	0.91509\\
3.1889	0.91898\\
3.2027	0.92187\\
3.2164	0.92375\\
3.2301	0.92462\\
3.2438	0.92447\\
3.2575	0.9233\\
3.2712	0.92113\\
3.2849	0.91796\\
3.2986	0.91381\\
3.3123	0.90868\\
3.326	0.90261\\
3.3397	0.8956\\
3.3534	0.88769\\
3.3671	0.87891\\
3.3843	0.86674\\
3.4014	0.85331\\
3.4185	0.8387\\
3.4357	0.82298\\
3.4562	0.80276\\
3.4768	0.7812\\
3.5008	0.75455\\
3.5248	0.72654\\
3.5522	0.69316\\
3.5865	0.64993\\
3.6413	0.579\\
3.703	0.49953\\
3.7407	0.45252\\
3.7715	0.41548\\
3.7989	0.38388\\
3.8263	0.35369\\
3.8503	0.32852\\
3.8743	0.3046\\
3.8983	0.28197\\
3.9223	0.26064\\
3.9463	0.24064\\
3.9703	0.22196\\
3.9942	0.20457\\
4.0182	0.18847\\
4.0422	0.17359\\
4.0662	0.1599\\
4.0902	0.14734\\
4.1142	0.13586\\
4.1382	0.12538\\
4.1622	0.11584\\
4.1896	0.10601\\
4.217	0.097221\\
4.2444	0.089371\\
4.2752	0.081546\\
4.3061	0.074658\\
4.3404	0.067959\\
4.378	0.061572\\
4.4192	0.055575\\
4.4637	0.049995\\
4.5151	0.044475\\
4.5768	0.038815\\
4.6488	0.033155\\
4.7344	0.027345\\
4.8338	0.021537\\
4.94	0.016265\\
5.0463	0.011904\\
5.1559	0.0083227\\
5.2724	0.0054596\\
5.3101	0.0047195\\
};
\addlegendentry{\RealBLNFitLegend}

      % \input{../../../../.tex/20/KS/SizeDistributionFittingsfig11.tex}
      \legend{}
  
      \nextgroupplot[%
        xmin=0.85621,xmax=5.4115,
        % xmin=1.8367,xmax=5.4755,
        % ymin=0,ymax=1.2,
      ] % This file was created by matlab2tikz.
%
\definecolor{mycolor1}{rgb}{0.83529,0.36863,0.00000}%
%
\addplot[ybar interval, fill=mycolor1, fill opacity=0.35, area legend, draw=none] table[row sep=crcr, x=Lower, y=Count] {%
Lower	Upper	Count\\
2.0021	2.2565	0.004087\\
2.2565	2.511	0.041918\\
2.511	2.7655	0.19859\\
2.7655	3.0199	0.55709\\
3.0199	3.2744	0.95625\\
3.2744	3.5289	0.97899\\
3.5289	3.7833	0.5924\\
3.7833	4.0378	0.30558\\
4.0378	4.2923	0.17082\\
4.2923	4.5467	0.078491\\
4.5467	4.8012	0.030286\\
4.8012	5.0557	0.0048205\\
5.0557	5.3101	0.010479\\
5.3101	5.3101	0.010479\\
};
\addlegendentry{\ApproxHistogramLegend}

\addplot [color=mycolor1, only marks, every error bar/.append style={opacity=0.45}, mark=*, mark size=0pt, draw=none, forget plot]
 plot [error bars/.cd, y dir=both, y explicit, error bar style={line width=1pt, color=mycolor1}, error mark options={mark=none,mark size=0pt}]
 table[row sep=crcr, y error plus index=2, y error minus index=3]{%
2.1293	0.004087	0.0013182	0.0013182\\
2.3838	0.041918	0.0038534	0.0038534\\
2.6382	0.19859	0.0061098	0.0061098\\
2.8927	0.55709	0.011261	0.011261\\
3.1472	0.95625	0.015779	0.015779\\
3.4016	0.97899	0.015597	0.015597\\
3.6561	0.5924	0.012018	0.012018\\
3.9106	0.30558	0.0078457	0.0078457\\
4.165	0.17082	0.0050805	0.0050805\\
4.4195	0.078491	0.0030926	0.0030926\\
4.674	0.030286	0.0020239	0.0020239\\
4.9284	0.0048205	0.0012684	0.0012684\\
5.1829	0.010479	2.4216e-18	1.7347e-18\\
};
\addplot [color=mycolor1]
  table[row sep=crcr]{%
2.0021	0.0020012\\
2.0815	0.0036252\\
2.1411	0.0056894\\
2.1908	0.0082925\\
2.2339	0.011484\\
2.2703	0.015099\\
2.3034	0.019319\\
2.3332	0.024056\\
2.3597	0.029161\\
2.3862	0.035255\\
2.4094	0.041521\\
2.4326	0.048778\\
2.4524	0.055879\\
2.4723	0.063878\\
2.4922	0.072858\\
2.512	0.082906\\
2.5319	0.094108\\
2.5518	0.10655\\
2.5683	0.11793\\
2.5849	0.13028\\
2.6014	0.14365\\
2.618	0.15807\\
2.6345	0.17359\\
2.6511	0.19023\\
2.6677	0.20804\\
2.6842	0.22703\\
2.7008	0.24721\\
2.7206	0.27301\\
2.7405	0.30053\\
2.7604	0.32974\\
2.7802	0.36059\\
2.8001	0.39299\\
2.8233	0.43261\\
2.8465	0.47398\\
2.873	0.52306\\
2.9061	0.58639\\
3.0021	0.77169\\
3.0286	0.81996\\
3.0485	0.85445\\
3.0683	0.8871\\
3.0849	0.91268\\
3.1014	0.93659\\
3.118	0.95867\\
3.1313	0.97491\\
3.1445	0.98981\\
3.1577	1.0033\\
3.171	1.0153\\
3.1842	1.0259\\
3.1942	1.0327\\
3.2041	1.0387\\
3.214	1.0438\\
3.224	1.048\\
3.2339	1.0512\\
3.2438	1.0536\\
3.2538	1.055\\
3.2637	1.0555\\
3.2736	1.0551\\
3.2836	1.0537\\
3.2935	1.0515\\
3.3034	1.0484\\
3.3134	1.0445\\
3.3233	1.0397\\
3.3332	1.034\\
3.3465	1.0253\\
3.3597	1.0151\\
3.373	1.0037\\
3.3862	0.991\\
3.3995	0.97715\\
3.416	0.95833\\
3.4326	0.93798\\
3.4491	0.91627\\
3.469	0.88866\\
3.4922	0.85468\\
3.5187	0.81408\\
3.5551	0.75639\\
3.6313	0.63521\\
3.6611	0.58961\\
3.6876	0.55062\\
3.7141	0.51337\\
3.7372	0.48235\\
3.7604	0.4529\\
3.7836	0.42508\\
3.8035	0.40254\\
3.8233	0.3812\\
3.8432	0.36103\\
3.8631	0.34203\\
3.8829	0.32413\\
3.9028	0.3073\\
3.926	0.28894\\
3.9492	0.27187\\
3.9723	0.25599\\
3.9955	0.24121\\
4.0187	0.22743\\
4.0452	0.2128\\
4.0717	0.19922\\
4.1015	0.18507\\
4.1313	0.17196\\
4.1644	0.15845\\
4.1975	0.14591\\
4.2339	0.13308\\
4.2737	0.12011\\
4.3134	0.10807\\
4.3565	0.095989\\
4.3995	0.084835\\
4.4459	0.073808\\
4.4922	0.063766\\
4.5386	0.05468\\
4.5849	0.046524\\
4.6313	0.039266\\
4.681	0.032441\\
4.7306	0.026544\\
4.7836	0.021201\\
4.8399	0.016493\\
4.8995	0.012467\\
4.9624	0.0091364\\
5.032	0.006361\\
5.1115	0.0041069\\
5.2075	0.0023409\\
5.3101	0.0012332\\
};
\addlegendentry{\ApproxBLNFitLegend}


\addplot[area legend, draw=none, fill=mycolor1, fill opacity=0.15, forget plot]
table[row sep=crcr] {%
x	y\\
2.0021	0.0023498\\
2.0054	0.002405\\
2.0087	0.0024615\\
2.012	0.0025193\\
2.0153	0.0025786\\
2.0186	0.0026393\\
2.0219	0.0027015\\
2.0253	0.0027652\\
2.0286	0.0028305\\
2.0319	0.0028973\\
2.0352	0.0029658\\
2.0385	0.003036\\
2.0418	0.0031078\\
2.0451	0.0031815\\
2.0484	0.0032569\\
2.0517	0.0033342\\
2.0551	0.0034134\\
2.0584	0.0034945\\
2.0617	0.0035776\\
2.065	0.0036628\\
2.0683	0.00375\\
2.0716	0.0038394\\
2.0749	0.003931\\
2.0782	0.0040248\\
2.0815	0.004121\\
2.0849	0.0042195\\
2.0882	0.0043204\\
2.0915	0.0044239\\
2.0948	0.0045299\\
2.0981	0.0046385\\
2.1014	0.0047498\\
2.1047	0.0048638\\
2.108	0.0049806\\
2.1113	0.0051004\\
2.1147	0.005223\\
2.118	0.0053488\\
2.1213	0.0054776\\
2.1246	0.0056096\\
2.1279	0.0057448\\
2.1312	0.0058834\\
2.1345	0.0060255\\
2.1378	0.006171\\
2.1411	0.0063201\\
2.1445	0.0064729\\
2.1478	0.0066295\\
2.1511	0.0067899\\
2.1544	0.0069543\\
2.1577	0.0071228\\
2.161	0.0072954\\
2.1643	0.0074722\\
2.1676	0.0076534\\
2.171	0.0078391\\
2.1743	0.0080293\\
2.1776	0.0082242\\
2.1809	0.0084239\\
2.1842	0.0086285\\
2.1875	0.0088381\\
2.1908	0.0090529\\
2.1941	0.0092729\\
2.1974	0.0094983\\
2.2008	0.0097293\\
2.2041	0.0099658\\
2.2074	0.010208\\
2.2107	0.010456\\
2.214	0.010711\\
2.2173	0.010971\\
2.2206	0.011238\\
2.2239	0.011511\\
2.2272	0.011791\\
2.2306	0.012078\\
2.2339	0.012372\\
2.2372	0.012672\\
2.2405	0.012981\\
2.2438	0.013296\\
2.2471	0.013619\\
2.2504	0.01395\\
2.2537	0.014289\\
2.257	0.014635\\
2.2604	0.014991\\
2.2637	0.015354\\
2.267	0.015726\\
2.2703	0.016108\\
2.2736	0.016498\\
2.2769	0.016897\\
2.2802	0.017306\\
2.2835	0.017725\\
2.2869	0.018153\\
2.2902	0.018592\\
2.2935	0.019041\\
2.2968	0.0195\\
2.3001	0.01997\\
2.3034	0.020451\\
2.3067	0.020943\\
2.31	0.021447\\
2.3133	0.021962\\
2.3167	0.022489\\
2.32	0.023029\\
2.3233	0.02358\\
2.3266	0.024145\\
2.3299	0.024722\\
2.3332	0.025313\\
2.3365	0.025917\\
2.3398	0.026534\\
2.3431	0.027166\\
2.3465	0.027811\\
2.3498	0.028472\\
2.3531	0.029147\\
2.3564	0.029837\\
2.3597	0.030543\\
2.363	0.031264\\
2.3663	0.032001\\
2.3696	0.032755\\
2.3729	0.033525\\
2.3763	0.034312\\
2.3796	0.035116\\
2.3829	0.035938\\
2.3862	0.036777\\
2.3895	0.037635\\
2.3928	0.038511\\
2.3961	0.039406\\
2.3994	0.04032\\
2.4027	0.041254\\
2.4061	0.042207\\
2.4094	0.043181\\
2.4127	0.044175\\
2.416	0.04519\\
2.4193	0.046226\\
2.4226	0.047284\\
2.4259	0.048363\\
2.4292	0.049465\\
2.4326	0.05059\\
2.4359	0.051737\\
2.4392	0.052908\\
2.4425	0.054103\\
2.4458	0.055322\\
2.4491	0.056565\\
2.4524	0.057833\\
2.4557	0.059127\\
2.459	0.060446\\
2.4624	0.061791\\
2.4657	0.063163\\
2.469	0.064561\\
2.4723	0.065987\\
2.4756	0.06744\\
2.4789	0.068921\\
2.4822	0.07043\\
2.4855	0.071969\\
2.4888	0.073536\\
2.4922	0.075133\\
2.4955	0.07676\\
2.4988	0.078417\\
2.5021	0.080105\\
2.5054	0.081825\\
2.5087	0.083575\\
2.512	0.085358\\
2.5153	0.087173\\
2.5186	0.089021\\
2.522	0.090902\\
2.5253	0.092816\\
2.5286	0.094765\\
2.5319	0.096748\\
2.5352	0.098765\\
2.5385	0.10082\\
2.5418	0.10291\\
2.5451	0.10503\\
2.5484	0.10719\\
2.5518	0.10939\\
2.5551	0.11162\\
2.5584	0.11389\\
2.5617	0.1162\\
2.565	0.11855\\
2.5683	0.12094\\
2.5716	0.12336\\
2.5749	0.12582\\
2.5783	0.12833\\
2.5816	0.13087\\
2.5849	0.13346\\
2.5882	0.13608\\
2.5915	0.13875\\
2.5948	0.14145\\
2.5981	0.1442\\
2.6014	0.14699\\
2.6047	0.14982\\
2.6081	0.1527\\
2.6114	0.15562\\
2.6147	0.15858\\
2.618	0.16159\\
2.6213	0.16463\\
2.6246	0.16773\\
2.6279	0.17086\\
2.6312	0.17405\\
2.6345	0.17727\\
2.6379	0.18054\\
2.6412	0.18386\\
2.6445	0.18722\\
2.6478	0.19063\\
2.6511	0.19408\\
2.6544	0.19758\\
2.6577	0.20113\\
2.661	0.20472\\
2.6643	0.20836\\
2.6677	0.21205\\
2.671	0.21578\\
2.6743	0.21956\\
2.6776	0.22339\\
2.6809	0.22726\\
2.6842	0.23118\\
2.6875	0.23515\\
2.6908	0.23917\\
2.6941	0.24323\\
2.6975	0.24734\\
2.7008	0.2515\\
2.7041	0.2557\\
2.7074	0.25995\\
2.7107	0.26425\\
2.714	0.2686\\
2.7173	0.273\\
2.7206	0.27744\\
2.724	0.28193\\
2.7273	0.28646\\
2.7306	0.29105\\
2.7339	0.29568\\
2.7372	0.30035\\
2.7405	0.30507\\
2.7438	0.30984\\
2.7471	0.31466\\
2.7504	0.31952\\
2.7538	0.32442\\
2.7571	0.32937\\
2.7604	0.33437\\
2.7637	0.33941\\
2.767	0.3445\\
2.7703	0.34963\\
2.7736	0.3548\\
2.7769	0.36001\\
2.7802	0.36527\\
2.7836	0.37057\\
2.7869	0.37592\\
2.7902	0.3813\\
2.7935	0.38673\\
2.7968	0.39219\\
2.8001	0.3977\\
2.8034	0.40325\\
2.8067	0.40883\\
2.81	0.41445\\
2.8134	0.42012\\
2.8167	0.42581\\
2.82	0.43155\\
2.8233	0.43732\\
2.8266	0.44313\\
2.8299	0.44897\\
2.8332	0.45484\\
2.8365	0.46075\\
2.8398	0.46669\\
2.8432	0.47266\\
2.8465	0.47867\\
2.8498	0.4847\\
2.8531	0.49076\\
2.8564	0.49685\\
2.8597	0.50297\\
2.863	0.50912\\
2.8663	0.51529\\
2.8697	0.52149\\
2.873	0.52771\\
2.8763	0.53395\\
2.8796	0.54022\\
2.8829	0.5465\\
2.8862	0.55281\\
2.8895	0.55914\\
2.8928	0.56548\\
2.8961	0.57184\\
2.8995	0.57822\\
2.9028	0.58461\\
2.9061	0.59101\\
2.9094	0.59743\\
2.9127	0.60386\\
2.916	0.6103\\
2.9193	0.61675\\
2.9226	0.62321\\
2.9259	0.62967\\
2.9293	0.63614\\
2.9326	0.64261\\
2.9359	0.64909\\
2.9392	0.65557\\
2.9425	0.66204\\
2.9458	0.66852\\
2.9491	0.67499\\
2.9524	0.68146\\
2.9557	0.68793\\
2.9591	0.69439\\
2.9624	0.70084\\
2.9657	0.70728\\
2.969	0.71371\\
2.9723	0.72012\\
2.9756	0.72653\\
2.9789	0.73291\\
2.9822	0.73928\\
2.9856	0.74563\\
2.9889	0.75196\\
2.9922	0.75827\\
2.9955	0.76456\\
2.9988	0.77082\\
3.0021	0.77705\\
3.0054	0.78326\\
3.0087	0.78944\\
3.012	0.79558\\
3.0154	0.80169\\
3.0187	0.80777\\
3.022	0.81381\\
3.0253	0.81981\\
3.0286	0.82578\\
3.0319	0.8317\\
3.0352	0.83758\\
3.0385	0.84341\\
3.0418	0.8492\\
3.0452	0.85494\\
3.0485	0.86063\\
3.0518	0.86627\\
3.0551	0.87186\\
3.0584	0.8774\\
3.0617	0.88288\\
3.065	0.8883\\
3.0683	0.89366\\
3.0716	0.89897\\
3.075	0.90421\\
3.0783	0.90938\\
3.0816	0.9145\\
3.0849	0.91954\\
3.0882	0.92452\\
3.0915	0.92943\\
3.0948	0.93427\\
3.0981	0.93904\\
3.1014	0.94374\\
3.1048	0.94836\\
3.1081	0.9529\\
3.1114	0.95737\\
3.1147	0.96175\\
3.118	0.96606\\
3.1213	0.97029\\
3.1246	0.97443\\
3.1279	0.9785\\
3.1313	0.98247\\
3.1346	0.98636\\
3.1379	0.99017\\
3.1412	0.99388\\
3.1445	0.99751\\
3.1478	1.001\\
3.1511	1.0045\\
3.1544	1.0078\\
3.1577	1.0111\\
3.1611	1.0143\\
3.1644	1.0174\\
3.1677	1.0203\\
3.171	1.0232\\
3.1743	1.026\\
3.1776	1.0287\\
3.1809	1.0313\\
3.1842	1.0338\\
3.1875	1.0362\\
3.1909	1.0384\\
3.1942	1.0406\\
3.1975	1.0427\\
3.2008	1.0447\\
3.2041	1.0466\\
3.2074	1.0484\\
3.2107	1.0501\\
3.214	1.0516\\
3.2173	1.0531\\
3.2207	1.0545\\
3.224	1.0558\\
3.2273	1.0569\\
3.2306	1.058\\
3.2339	1.0589\\
3.2372	1.0598\\
3.2405	1.0605\\
3.2438	1.0612\\
3.2471	1.0617\\
3.2505	1.0622\\
3.2538	1.0625\\
3.2571	1.0627\\
3.2604	1.0628\\
3.2637	1.0629\\
3.267	1.0628\\
3.2703	1.0626\\
3.2736	1.0623\\
3.277	1.0619\\
3.2803	1.0615\\
3.2836	1.0609\\
3.2869	1.0602\\
3.2902	1.0594\\
3.2935	1.0585\\
3.2968	1.0575\\
3.3001	1.0565\\
3.3034	1.0553\\
3.3068	1.054\\
3.3101	1.0526\\
3.3134	1.0512\\
3.3167	1.0496\\
3.32	1.048\\
3.3233	1.0462\\
3.3266	1.0444\\
3.3299	1.0425\\
3.3332	1.0405\\
3.3366	1.0384\\
3.3399	1.0362\\
3.3432	1.0339\\
3.3465	1.0316\\
3.3498	1.0291\\
3.3531	1.0266\\
3.3564	1.024\\
3.3597	1.0213\\
3.363	1.0186\\
3.3664	1.0157\\
3.3697	1.0128\\
3.373	1.0098\\
3.3763	1.0067\\
3.3796	1.0036\\
3.3829	1.0004\\
3.3862	0.99706\\
3.3895	0.9937\\
3.3928	0.99027\\
3.3962	0.98677\\
3.3995	0.98321\\
3.4028	0.97958\\
3.4061	0.97588\\
3.4094	0.97213\\
3.4127	0.96831\\
3.416	0.96443\\
3.4193	0.96049\\
3.4227	0.95649\\
3.426	0.95244\\
3.4293	0.94833\\
3.4326	0.94417\\
3.4359	0.93995\\
3.4392	0.93568\\
3.4425	0.93136\\
3.4458	0.927\\
3.4491	0.92258\\
3.4525	0.91812\\
3.4558	0.91361\\
3.4591	0.90906\\
3.4624	0.90446\\
3.4657	0.89983\\
3.469	0.89515\\
3.4723	0.89043\\
3.4756	0.88568\\
3.4789	0.88089\\
3.4823	0.87606\\
3.4856	0.8712\\
3.4889	0.86631\\
3.4922	0.86138\\
3.4955	0.85643\\
3.4988	0.85145\\
3.5021	0.84644\\
3.5054	0.8414\\
3.5087	0.83633\\
3.5121	0.83125\\
3.5154	0.82614\\
3.5187	0.82101\\
3.522	0.81586\\
3.5253	0.81069\\
3.5286	0.8055\\
3.5319	0.8003\\
3.5352	0.79508\\
3.5385	0.78984\\
3.5419	0.7846\\
3.5452	0.77934\\
3.5485	0.77407\\
3.5518	0.76879\\
3.5551	0.76351\\
3.5584	0.75821\\
3.5617	0.75291\\
3.565	0.74761\\
3.5684	0.7423\\
3.5717	0.73699\\
3.575	0.73168\\
3.5783	0.72636\\
3.5816	0.72105\\
3.5849	0.71574\\
3.5882	0.71043\\
3.5915	0.70513\\
3.5948	0.69983\\
3.5982	0.69453\\
3.6015	0.68924\\
3.6048	0.68396\\
3.6081	0.67869\\
3.6114	0.67343\\
3.6147	0.66817\\
3.618	0.66293\\
3.6213	0.6577\\
3.6246	0.65248\\
3.628	0.64728\\
3.6313	0.64209\\
3.6346	0.63692\\
3.6379	0.63176\\
3.6412	0.62662\\
3.6445	0.6215\\
3.6478	0.61639\\
3.6511	0.61131\\
3.6544	0.60624\\
3.6578	0.60119\\
3.6611	0.59617\\
3.6644	0.59117\\
3.6677	0.58618\\
3.671	0.58123\\
3.6743	0.57629\\
3.6776	0.57138\\
3.6809	0.5665\\
3.6842	0.56163\\
3.6876	0.5568\\
3.6909	0.55199\\
3.6942	0.54721\\
3.6975	0.54245\\
3.7008	0.53773\\
3.7041	0.53303\\
3.7074	0.52836\\
3.7107	0.52372\\
3.7141	0.5191\\
3.7174	0.51452\\
3.7207	0.50997\\
3.724	0.50544\\
3.7273	0.50095\\
3.7306	0.49649\\
3.7339	0.49206\\
3.7372	0.48766\\
3.7405	0.4833\\
3.7439	0.47896\\
3.7472	0.47466\\
3.7505	0.47039\\
3.7538	0.46615\\
3.7571	0.46195\\
3.7604	0.45778\\
3.7637	0.45364\\
3.767	0.44953\\
3.7703	0.44546\\
3.7737	0.44143\\
3.777	0.43742\\
3.7803	0.43345\\
3.7836	0.42952\\
3.7869	0.42561\\
3.7902	0.42174\\
3.7935	0.41791\\
3.7968	0.41411\\
3.8001	0.41034\\
3.8035	0.40661\\
3.8068	0.40291\\
3.8101	0.39925\\
3.8134	0.39562\\
3.8167	0.39203\\
3.82	0.38846\\
3.8233	0.38494\\
3.8266	0.38144\\
3.83	0.37798\\
3.8333	0.37456\\
3.8366	0.37116\\
3.8399	0.36781\\
3.8432	0.36448\\
3.8465	0.36119\\
3.8498	0.35793\\
3.8531	0.3547\\
3.8564	0.35151\\
3.8598	0.34835\\
3.8631	0.34522\\
3.8664	0.34213\\
3.8697	0.33906\\
3.873	0.33603\\
3.8763	0.33303\\
3.8796	0.33006\\
3.8829	0.32713\\
3.8862	0.32422\\
3.8896	0.32135\\
3.8929	0.3185\\
3.8962	0.31569\\
3.8995	0.31291\\
3.9028	0.31015\\
3.9061	0.30743\\
3.9094	0.30473\\
3.9127	0.30207\\
3.916	0.29943\\
3.9194	0.29682\\
3.9227	0.29424\\
3.926	0.29169\\
3.9293	0.28916\\
3.9326	0.28666\\
3.9359	0.28419\\
3.9392	0.28175\\
3.9425	0.27933\\
3.9458	0.27694\\
3.9492	0.27457\\
3.9525	0.27223\\
3.9558	0.26991\\
3.9591	0.26762\\
3.9624	0.26535\\
3.9657	0.2631\\
3.969	0.26088\\
3.9723	0.25868\\
3.9757	0.25651\\
3.979	0.25436\\
3.9823	0.25223\\
3.9856	0.25012\\
3.9889	0.24803\\
3.9922	0.24596\\
3.9955	0.24392\\
3.9988	0.24189\\
4.0021	0.23989\\
4.0055	0.2379\\
4.0088	0.23594\\
4.0121	0.23399\\
4.0154	0.23207\\
4.0187	0.23016\\
4.022	0.22827\\
4.0253	0.2264\\
4.0286	0.22455\\
4.0319	0.22271\\
4.0353	0.22089\\
4.0386	0.21909\\
4.0419	0.21731\\
4.0452	0.21554\\
4.0485	0.21379\\
4.0518	0.21205\\
4.0551	0.21033\\
4.0584	0.20863\\
4.0617	0.20694\\
4.0651	0.20526\\
4.0684	0.2036\\
4.0717	0.20196\\
4.075	0.20033\\
4.0783	0.19871\\
4.0816	0.19711\\
4.0849	0.19552\\
4.0882	0.19394\\
4.0915	0.19238\\
4.0949	0.19083\\
4.0982	0.18929\\
4.1015	0.18777\\
4.1048	0.18626\\
4.1081	0.18476\\
4.1114	0.18327\\
4.1147	0.18179\\
4.118	0.18033\\
4.1214	0.17888\\
4.1247	0.17743\\
4.128	0.176\\
4.1313	0.17458\\
4.1346	0.17318\\
4.1379	0.17178\\
4.1412	0.17039\\
4.1445	0.16901\\
4.1478	0.16764\\
4.1512	0.16629\\
4.1545	0.16494\\
4.1578	0.1636\\
4.1611	0.16227\\
4.1644	0.16095\\
4.1677	0.15964\\
4.171	0.15834\\
4.1743	0.15705\\
4.1776	0.15577\\
4.181	0.15449\\
4.1843	0.15323\\
4.1876	0.15197\\
4.1909	0.15072\\
4.1942	0.14948\\
4.1975	0.14825\\
4.2008	0.14703\\
4.2041	0.14581\\
4.2074	0.1446\\
4.2108	0.1434\\
4.2141	0.14221\\
4.2174	0.14103\\
4.2207	0.13985\\
4.224	0.13868\\
4.2273	0.13752\\
4.2306	0.13637\\
4.2339	0.13522\\
4.2372	0.13408\\
4.2406	0.13295\\
4.2439	0.13182\\
4.2472	0.1307\\
4.2505	0.12959\\
4.2538	0.12848\\
4.2571	0.12738\\
4.2604	0.12629\\
4.2637	0.12521\\
4.2671	0.12413\\
4.2704	0.12306\\
4.2737	0.12199\\
4.277	0.12093\\
4.2803	0.11988\\
4.2836	0.11883\\
4.2869	0.11779\\
4.2902	0.11676\\
4.2935	0.11573\\
4.2969	0.11471\\
4.3002	0.11369\\
4.3035	0.11269\\
4.3068	0.11168\\
4.3101	0.11068\\
4.3134	0.10969\\
4.3167	0.10871\\
4.32	0.10773\\
4.3233	0.10676\\
4.3267	0.10579\\
4.33	0.10483\\
4.3333	0.10387\\
4.3366	0.10292\\
4.3399	0.10197\\
4.3432	0.10104\\
4.3465	0.1001\\
4.3498	0.099174\\
4.3531	0.098252\\
4.3565	0.097336\\
4.3598	0.096425\\
4.3631	0.09552\\
4.3664	0.094621\\
4.3697	0.093727\\
4.373	0.092838\\
4.3763	0.091955\\
4.3796	0.091077\\
4.3829	0.090205\\
4.3863	0.089339\\
4.3896	0.088478\\
4.3929	0.087622\\
4.3962	0.086772\\
4.3995	0.085927\\
4.4028	0.085087\\
4.4061	0.084253\\
4.4094	0.083425\\
4.4128	0.082601\\
4.4161	0.081783\\
4.4194	0.080971\\
4.4227	0.080164\\
4.426	0.079362\\
4.4293	0.078565\\
4.4326	0.077774\\
4.4359	0.076988\\
4.4392	0.076207\\
4.4426	0.075431\\
4.4459	0.074661\\
4.4492	0.073896\\
4.4525	0.073137\\
4.4558	0.072382\\
4.4591	0.071633\\
4.4624	0.070889\\
4.4657	0.07015\\
4.469	0.069417\\
4.4724	0.068689\\
4.4757	0.067966\\
4.479	0.067248\\
4.4823	0.066535\\
4.4856	0.065828\\
4.4889	0.065125\\
4.4922	0.064428\\
4.4955	0.063736\\
4.4988	0.063049\\
4.5022	0.062368\\
4.5055	0.061691\\
4.5088	0.06102\\
4.5121	0.060353\\
4.5154	0.059692\\
4.5187	0.059036\\
4.522	0.058385\\
4.5253	0.057739\\
4.5287	0.057098\\
4.532	0.056462\\
4.5353	0.055832\\
4.5386	0.055206\\
4.5419	0.054585\\
4.5452	0.053969\\
4.5485	0.053359\\
4.5518	0.052753\\
4.5551	0.052152\\
4.5585	0.051556\\
4.5618	0.050965\\
4.5651	0.050379\\
4.5684	0.049798\\
4.5717	0.049222\\
4.575	0.048651\\
4.5783	0.048084\\
4.5816	0.047523\\
4.5849	0.046966\\
4.5883	0.046414\\
4.5916	0.045867\\
4.5949	0.045324\\
4.5982	0.044787\\
4.6015	0.044254\\
4.6048	0.043725\\
4.6081	0.043202\\
4.6114	0.042683\\
4.6147	0.042168\\
4.6181	0.041659\\
4.6214	0.041154\\
4.6247	0.040653\\
4.628	0.040157\\
4.6313	0.039666\\
4.6346	0.039179\\
4.6379	0.038696\\
4.6412	0.038218\\
4.6445	0.037745\\
4.6479	0.037275\\
4.6512	0.036811\\
4.6545	0.03635\\
4.6578	0.035894\\
4.6611	0.035442\\
4.6644	0.034995\\
4.6677	0.034552\\
4.671	0.034113\\
4.6744	0.033678\\
4.6777	0.033248\\
4.681	0.032822\\
4.6843	0.032399\\
4.6876	0.031981\\
4.6909	0.031568\\
4.6942	0.031158\\
4.6975	0.030752\\
4.7008	0.03035\\
4.7042	0.029953\\
4.7075	0.029559\\
4.7108	0.029169\\
4.7141	0.028784\\
4.7174	0.028402\\
4.7207	0.028024\\
4.724	0.02765\\
4.7273	0.02728\\
4.7306	0.026914\\
4.734	0.026551\\
4.7373	0.026192\\
4.7406	0.025837\\
4.7439	0.025486\\
4.7472	0.025139\\
4.7505	0.024795\\
4.7538	0.024455\\
4.7571	0.024118\\
4.7604	0.023785\\
4.7638	0.023456\\
4.7671	0.023131\\
4.7704	0.022808\\
4.7737	0.02249\\
4.777	0.022175\\
4.7803	0.021863\\
4.7836	0.021555\\
4.7869	0.02125\\
4.7902	0.020949\\
4.7936	0.020651\\
4.7969	0.020357\\
4.8002	0.020065\\
4.8035	0.019778\\
4.8068	0.019493\\
4.8101	0.019212\\
4.8134	0.018934\\
4.8167	0.018659\\
4.8201	0.018387\\
4.8234	0.018119\\
4.8267	0.017853\\
4.83	0.017591\\
4.8333	0.017332\\
4.8366	0.017076\\
4.8399	0.016823\\
4.8432	0.016573\\
4.8465	0.016326\\
4.8499	0.016082\\
4.8532	0.015842\\
4.8565	0.015603\\
4.8598	0.015368\\
4.8631	0.015136\\
4.8664	0.014907\\
4.8697	0.01468\\
4.873	0.014457\\
4.8763	0.014236\\
4.8797	0.014017\\
4.883	0.013802\\
4.8863	0.013589\\
4.8896	0.013379\\
4.8929	0.013172\\
4.8962	0.012967\\
4.8995	0.012765\\
4.9028	0.012566\\
4.9061	0.012369\\
4.9095	0.012175\\
4.9128	0.011983\\
4.9161	0.011794\\
4.9194	0.011607\\
4.9227	0.011423\\
4.926	0.011241\\
4.9293	0.011061\\
4.9326	0.010884\\
4.9359	0.01071\\
4.9393	0.010537\\
4.9426	0.010367\\
4.9459	0.0102\\
4.9492	0.010034\\
4.9525	0.0098709\\
4.9558	0.0097099\\
4.9591	0.0095511\\
4.9624	0.0093946\\
4.9658	0.0092402\\
4.9691	0.0090879\\
4.9724	0.0089378\\
4.9757	0.0087898\\
4.979	0.0086438\\
4.9823	0.0084999\\
4.9856	0.008358\\
4.9889	0.0082182\\
4.9922	0.0080803\\
4.9956	0.0079444\\
4.9989	0.0078105\\
5.0022	0.0076784\\
5.0055	0.0075483\\
5.0088	0.0074201\\
5.0121	0.0072937\\
5.0154	0.0071692\\
5.0187	0.0070464\\
5.022	0.0069255\\
5.0254	0.0068064\\
5.0287	0.006689\\
5.032	0.0065734\\
5.0353	0.0064595\\
5.0386	0.0063473\\
5.0419	0.0062367\\
5.0452	0.0061279\\
5.0485	0.0060206\\
5.0518	0.005915\\
5.0552	0.005811\\
5.0585	0.0057086\\
5.0618	0.0056077\\
5.0651	0.0055084\\
5.0684	0.0054106\\
5.0717	0.0053143\\
5.075	0.0052195\\
5.0783	0.0051262\\
5.0816	0.0050343\\
5.085	0.0049439\\
5.0883	0.0048549\\
5.0916	0.0047672\\
5.0949	0.004681\\
5.0982	0.0045961\\
5.1015	0.0045126\\
5.1048	0.0044304\\
5.1081	0.0043495\\
5.1115	0.0042698\\
5.1148	0.0041915\\
5.1181	0.0041144\\
5.1214	0.0040386\\
5.1247	0.003964\\
5.128	0.0038906\\
5.1313	0.0038184\\
5.1346	0.0037474\\
5.1379	0.0036776\\
5.1413	0.0036089\\
5.1446	0.0035413\\
5.1479	0.0034749\\
5.1512	0.0034095\\
5.1545	0.0033453\\
5.1578	0.0032821\\
5.1611	0.0032199\\
5.1644	0.0031588\\
5.1677	0.0030988\\
5.1711	0.0030397\\
5.1744	0.0029816\\
5.1777	0.0029246\\
5.181	0.0028685\\
5.1843	0.0028133\\
5.1876	0.0027591\\
5.1909	0.0027058\\
5.1942	0.0026535\\
5.1975	0.002602\\
5.2009	0.0025515\\
5.2042	0.0025018\\
5.2075	0.0024529\\
5.2108	0.002405\\
5.2141	0.0023578\\
5.2174	0.0023115\\
5.2207	0.002266\\
5.224	0.0022213\\
5.2274	0.0021774\\
5.2307	0.0021343\\
5.234	0.0020919\\
5.2373	0.0020503\\
5.2406	0.0020095\\
5.2439	0.0019693\\
5.2472	0.0019299\\
5.2505	0.0018912\\
5.2538	0.0018532\\
5.2572	0.0018159\\
5.2605	0.0017792\\
5.2638	0.0017433\\
5.2671	0.001708\\
5.2704	0.0016733\\
5.2737	0.0016392\\
5.277	0.0016058\\
5.2803	0.001573\\
5.2836	0.0015408\\
5.287	0.0015092\\
5.2903	0.0014782\\
5.2936	0.0014478\\
5.2969	0.0014179\\
5.3002	0.0013886\\
5.3035	0.0013598\\
5.3068	0.0013316\\
5.3101	0.0013039\\
5.3101	0.0011625\\
5.3068	0.0011879\\
5.3035	0.0012139\\
5.3002	0.0012403\\
5.2969	0.0012673\\
5.2936	0.0012948\\
5.2903	0.0013229\\
5.287	0.0013515\\
5.2836	0.0013807\\
5.2803	0.0014104\\
5.277	0.0014407\\
5.2737	0.0014716\\
5.2704	0.0015031\\
5.2671	0.0015352\\
5.2638	0.0015679\\
5.2605	0.0016013\\
5.2572	0.0016352\\
5.2538	0.0016699\\
5.2505	0.0017052\\
5.2472	0.0017411\\
5.2439	0.0017777\\
5.2406	0.0018151\\
5.2373	0.0018531\\
5.234	0.0018919\\
5.2307	0.0019313\\
5.2274	0.0019715\\
5.224	0.0020125\\
5.2207	0.0020542\\
5.2174	0.0020967\\
5.2141	0.00214\\
5.2108	0.0021841\\
5.2075	0.0022289\\
5.2042	0.0022746\\
5.2009	0.0023212\\
5.1975	0.0023686\\
5.1942	0.0024168\\
5.1909	0.0024659\\
5.1876	0.0025159\\
5.1843	0.0025669\\
5.181	0.0026187\\
5.1777	0.0026714\\
5.1744	0.0027251\\
5.1711	0.0027798\\
5.1677	0.0028354\\
5.1644	0.002892\\
5.1611	0.0029496\\
5.1578	0.0030082\\
5.1545	0.0030679\\
5.1512	0.0031286\\
5.1479	0.0031903\\
5.1446	0.0032532\\
5.1413	0.0033171\\
5.1379	0.0033821\\
5.1346	0.0034482\\
5.1313	0.0035155\\
5.128	0.003584\\
5.1247	0.0036536\\
5.1214	0.0037244\\
5.1181	0.0037964\\
5.1148	0.0038696\\
5.1115	0.003944\\
5.1081	0.0040197\\
5.1048	0.0040967\\
5.1015	0.004175\\
5.0982	0.0042545\\
5.0949	0.0043354\\
5.0916	0.0044177\\
5.0883	0.0045013\\
5.085	0.0045862\\
5.0816	0.0046726\\
5.0783	0.0047604\\
5.075	0.0048496\\
5.0717	0.0049403\\
5.0684	0.0050324\\
5.0651	0.005126\\
5.0618	0.0052212\\
5.0585	0.0053178\\
5.0552	0.0054161\\
5.0518	0.0055158\\
5.0485	0.0056172\\
5.0452	0.0057202\\
5.0419	0.0058248\\
5.0386	0.005931\\
5.0353	0.0060389\\
5.032	0.0061485\\
5.0287	0.0062599\\
5.0254	0.0063729\\
5.022	0.0064877\\
5.0187	0.0066043\\
5.0154	0.0067226\\
5.0121	0.0068428\\
5.0088	0.0069648\\
5.0055	0.0070887\\
5.0022	0.0072144\\
4.9989	0.0073421\\
4.9956	0.0074717\\
4.9922	0.0076032\\
4.9889	0.0077367\\
4.9856	0.0078722\\
4.9823	0.0080097\\
4.979	0.0081492\\
4.9757	0.0082908\\
4.9724	0.0084344\\
4.9691	0.0085802\\
4.9658	0.0087281\\
4.9624	0.0088782\\
4.9591	0.0090304\\
4.9558	0.0091848\\
4.9525	0.0093414\\
4.9492	0.0095003\\
4.9459	0.0096615\\
4.9426	0.0098249\\
4.9393	0.0099907\\
4.9359	0.010159\\
4.9326	0.010329\\
4.9293	0.010502\\
4.926	0.010677\\
4.9227	0.010855\\
4.9194	0.011035\\
4.9161	0.011218\\
4.9128	0.011403\\
4.9095	0.011591\\
4.9061	0.011781\\
4.9028	0.011974\\
4.8995	0.012169\\
4.8962	0.012367\\
4.8929	0.012568\\
4.8896	0.012772\\
4.8863	0.012978\\
4.883	0.013187\\
4.8797	0.013398\\
4.8763	0.013613\\
4.873	0.01383\\
4.8697	0.01405\\
4.8664	0.014273\\
4.8631	0.014499\\
4.8598	0.014727\\
4.8565	0.014959\\
4.8532	0.015193\\
4.8499	0.015431\\
4.8465	0.015672\\
4.8432	0.015915\\
4.8399	0.016162\\
4.8366	0.016412\\
4.8333	0.016664\\
4.83	0.01692\\
4.8267	0.01718\\
4.8234	0.017442\\
4.8201	0.017707\\
4.8167	0.017976\\
4.8134	0.018248\\
4.8101	0.018523\\
4.8068	0.018802\\
4.8035	0.019084\\
4.8002	0.019369\\
4.7969	0.019658\\
4.7936	0.01995\\
4.7902	0.020245\\
4.7869	0.020544\\
4.7836	0.020847\\
4.7803	0.021153\\
4.777	0.021462\\
4.7737	0.021775\\
4.7704	0.022091\\
4.7671	0.022412\\
4.7638	0.022735\\
4.7604	0.023063\\
4.7571	0.023393\\
4.7538	0.023728\\
4.7505	0.024066\\
4.7472	0.024409\\
4.7439	0.024754\\
4.7406	0.025104\\
4.7373	0.025457\\
4.734	0.025814\\
4.7306	0.026175\\
4.7273	0.02654\\
4.724	0.026909\\
4.7207	0.027281\\
4.7174	0.027658\\
4.7141	0.028038\\
4.7108	0.028422\\
4.7075	0.02881\\
4.7042	0.029203\\
4.7008	0.029599\\
4.6975	0.029999\\
4.6942	0.030403\\
4.6909	0.030811\\
4.6876	0.031224\\
4.6843	0.03164\\
4.681	0.03206\\
4.6777	0.032485\\
4.6744	0.032914\\
4.671	0.033346\\
4.6677	0.033783\\
4.6644	0.034224\\
4.6611	0.03467\\
4.6578	0.035119\\
4.6545	0.035573\\
4.6512	0.03603\\
4.6479	0.036492\\
4.6445	0.036959\\
4.6412	0.037429\\
4.6379	0.037904\\
4.6346	0.038383\\
4.6313	0.038866\\
4.628	0.039353\\
4.6247	0.039845\\
4.6214	0.040341\\
4.6181	0.040841\\
4.6147	0.041346\\
4.6114	0.041855\\
4.6081	0.042368\\
4.6048	0.042886\\
4.6015	0.043408\\
4.5982	0.043934\\
4.5949	0.044464\\
4.5916	0.044999\\
4.5883	0.045539\\
4.5849	0.046082\\
4.5816	0.04663\\
4.5783	0.047183\\
4.575	0.047739\\
4.5717	0.0483\\
4.5684	0.048866\\
4.5651	0.049436\\
4.5618	0.05001\\
4.5585	0.050589\\
4.5551	0.051172\\
4.5518	0.05176\\
4.5485	0.052352\\
4.5452	0.052948\\
4.5419	0.053549\\
4.5386	0.054155\\
4.5353	0.054765\\
4.532	0.055379\\
4.5287	0.055998\\
4.5253	0.056621\\
4.522	0.057249\\
4.5187	0.057882\\
4.5154	0.058518\\
4.5121	0.05916\\
4.5088	0.059806\\
4.5055	0.060456\\
4.5022	0.061111\\
4.4988	0.061771\\
4.4955	0.062435\\
4.4922	0.063104\\
4.4889	0.063778\\
4.4856	0.064456\\
4.4823	0.065138\\
4.479	0.065825\\
4.4757	0.066517\\
4.4724	0.067214\\
4.469	0.067915\\
4.4657	0.068621\\
4.4624	0.069332\\
4.4591	0.070047\\
4.4558	0.070767\\
4.4525	0.071492\\
4.4492	0.072221\\
4.4459	0.072956\\
4.4426	0.073695\\
4.4392	0.074439\\
4.4359	0.075187\\
4.4326	0.075941\\
4.4293	0.076699\\
4.426	0.077462\\
4.4227	0.07823\\
4.4194	0.079003\\
4.4161	0.079781\\
4.4128	0.080563\\
4.4094	0.081351\\
4.4061	0.082144\\
4.4028	0.082941\\
4.3995	0.083744\\
4.3962	0.084551\\
4.3929	0.085364\\
4.3896	0.086181\\
4.3863	0.087004\\
4.3829	0.087832\\
4.3796	0.088665\\
4.3763	0.089503\\
4.373	0.090346\\
4.3697	0.091195\\
4.3664	0.092049\\
4.3631	0.092908\\
4.3598	0.093772\\
4.3565	0.094641\\
4.3531	0.095516\\
4.3498	0.096397\\
4.3465	0.097283\\
4.3432	0.098174\\
4.3399	0.09907\\
4.3366	0.099973\\
4.3333	0.10088\\
4.33	0.10179\\
4.3267	0.10271\\
4.3233	0.10364\\
4.32	0.10457\\
4.3167	0.1055\\
4.3134	0.10645\\
4.3101	0.10739\\
4.3068	0.10835\\
4.3035	0.10931\\
4.3002	0.11027\\
4.2969	0.11125\\
4.2935	0.11222\\
4.2902	0.11321\\
4.2869	0.1142\\
4.2836	0.11519\\
4.2803	0.1162\\
4.277	0.11721\\
4.2737	0.11822\\
4.2704	0.11925\\
4.2671	0.12028\\
4.2637	0.12131\\
4.2604	0.12235\\
4.2571	0.1234\\
4.2538	0.12446\\
4.2505	0.12552\\
4.2472	0.12659\\
4.2439	0.12767\\
4.2406	0.12876\\
4.2372	0.12985\\
4.2339	0.13095\\
4.2306	0.13206\\
4.2273	0.13317\\
4.224	0.13429\\
4.2207	0.13542\\
4.2174	0.13656\\
4.2141	0.13771\\
4.2108	0.13886\\
4.2074	0.14003\\
4.2041	0.1412\\
4.2008	0.14238\\
4.1975	0.14357\\
4.1942	0.14476\\
4.1909	0.14597\\
4.1876	0.14718\\
4.1843	0.14841\\
4.181	0.14964\\
4.1776	0.15088\\
4.1743	0.15213\\
4.171	0.15339\\
4.1677	0.15467\\
4.1644	0.15595\\
4.1611	0.15724\\
4.1578	0.15854\\
4.1545	0.15985\\
4.1512	0.16117\\
4.1478	0.1625\\
4.1445	0.16385\\
4.1412	0.1652\\
4.1379	0.16657\\
4.1346	0.16794\\
4.1313	0.16933\\
4.128	0.17073\\
4.1247	0.17214\\
4.1214	0.17357\\
4.118	0.175\\
4.1147	0.17645\\
4.1114	0.17791\\
4.1081	0.17938\\
4.1048	0.18087\\
4.1015	0.18237\\
4.0982	0.18388\\
4.0949	0.18541\\
4.0915	0.18695\\
4.0882	0.1885\\
4.0849	0.19007\\
4.0816	0.19165\\
4.0783	0.19325\\
4.075	0.19486\\
4.0717	0.19648\\
4.0684	0.19813\\
4.0651	0.19978\\
4.0617	0.20145\\
4.0584	0.20314\\
4.0551	0.20485\\
4.0518	0.20657\\
4.0485	0.2083\\
4.0452	0.21005\\
4.0419	0.21182\\
4.0386	0.21361\\
4.0353	0.21542\\
4.0319	0.21724\\
4.0286	0.21908\\
4.0253	0.22094\\
4.022	0.22281\\
4.0187	0.22471\\
4.0154	0.22662\\
4.0121	0.22855\\
4.0088	0.2305\\
4.0055	0.23247\\
4.0021	0.23446\\
3.9988	0.23647\\
3.9955	0.2385\\
3.9922	0.24056\\
3.9889	0.24263\\
3.9856	0.24472\\
3.9823	0.24683\\
3.979	0.24897\\
3.9757	0.25112\\
3.9723	0.2533\\
3.969	0.2555\\
3.9657	0.25772\\
3.9624	0.25997\\
3.9591	0.26223\\
3.9558	0.26452\\
3.9525	0.26683\\
3.9492	0.26917\\
3.9458	0.27153\\
3.9425	0.27391\\
3.9392	0.27632\\
3.9359	0.27875\\
3.9326	0.28121\\
3.9293	0.28369\\
3.926	0.2862\\
3.9227	0.28873\\
3.9194	0.29128\\
3.916	0.29386\\
3.9127	0.29647\\
3.9094	0.2991\\
3.9061	0.30176\\
3.9028	0.30445\\
3.8995	0.30716\\
3.8962	0.3099\\
3.8929	0.31267\\
3.8896	0.31546\\
3.8862	0.31829\\
3.8829	0.32114\\
3.8796	0.32401\\
3.8763	0.32692\\
3.873	0.32985\\
3.8697	0.33282\\
3.8664	0.33581\\
3.8631	0.33883\\
3.8598	0.34188\\
3.8564	0.34497\\
3.8531	0.34808\\
3.8498	0.35122\\
3.8465	0.35439\\
3.8432	0.35759\\
3.8399	0.36082\\
3.8366	0.36409\\
3.8333	0.36738\\
3.83	0.37071\\
3.8266	0.37406\\
3.8233	0.37745\\
3.82	0.38088\\
3.8167	0.38433\\
3.8134	0.38781\\
3.8101	0.39133\\
3.8068	0.39488\\
3.8035	0.39846\\
3.8001	0.40208\\
3.7968	0.40573\\
3.7935	0.40941\\
3.7902	0.41312\\
3.7869	0.41687\\
3.7836	0.42065\\
3.7803	0.42446\\
3.777	0.42831\\
3.7737	0.43218\\
3.7703	0.4361\\
3.767	0.44004\\
3.7637	0.44402\\
3.7604	0.44803\\
3.7571	0.45208\\
3.7538	0.45616\\
3.7505	0.46027\\
3.7472	0.46441\\
3.7439	0.46859\\
3.7405	0.4728\\
3.7372	0.47704\\
3.7339	0.48131\\
3.7306	0.48562\\
3.7273	0.48996\\
3.724	0.49433\\
3.7207	0.49873\\
3.7174	0.50316\\
3.7141	0.50763\\
3.7107	0.51212\\
3.7074	0.51665\\
3.7041	0.52121\\
3.7008	0.52579\\
3.6975	0.53041\\
3.6942	0.53506\\
3.6909	0.53973\\
3.6876	0.54444\\
3.6842	0.54917\\
3.6809	0.55393\\
3.6776	0.55872\\
3.6743	0.56353\\
3.671	0.56837\\
3.6677	0.57324\\
3.6644	0.57813\\
3.6611	0.58305\\
3.6578	0.58799\\
3.6544	0.59295\\
3.6511	0.59794\\
3.6478	0.60295\\
3.6445	0.60799\\
3.6412	0.61304\\
3.6379	0.61811\\
3.6346	0.62321\\
3.6313	0.62832\\
3.628	0.63345\\
3.6246	0.6386\\
3.6213	0.64377\\
3.618	0.64895\\
3.6147	0.65415\\
3.6114	0.65936\\
3.6081	0.66459\\
3.6048	0.66983\\
3.6015	0.67508\\
3.5982	0.68034\\
3.5948	0.68561\\
3.5915	0.69089\\
3.5882	0.69618\\
3.5849	0.70148\\
3.5816	0.70678\\
3.5783	0.71208\\
3.575	0.71739\\
3.5717	0.72271\\
3.5684	0.72802\\
3.565	0.73334\\
3.5617	0.73865\\
3.5584	0.74397\\
3.5551	0.74928\\
3.5518	0.75458\\
3.5485	0.75988\\
3.5452	0.76518\\
3.5419	0.77047\\
3.5385	0.77575\\
3.5352	0.78101\\
3.5319	0.78627\\
3.5286	0.79151\\
3.5253	0.79675\\
3.522	0.80196\\
3.5187	0.80716\\
3.5154	0.81234\\
3.5121	0.8175\\
3.5087	0.82264\\
3.5054	0.82776\\
3.5021	0.83285\\
3.4988	0.83792\\
3.4955	0.84296\\
3.4922	0.84798\\
3.4889	0.85296\\
3.4856	0.85792\\
3.4823	0.86284\\
3.4789	0.86773\\
3.4756	0.87258\\
3.4723	0.8774\\
3.469	0.88218\\
3.4657	0.88691\\
3.4624	0.89161\\
3.4591	0.89627\\
3.4558	0.90088\\
3.4525	0.90545\\
3.4491	0.90996\\
3.4458	0.91443\\
3.4425	0.91885\\
3.4392	0.92322\\
3.4359	0.92754\\
3.4326	0.9318\\
3.4293	0.936\\
3.426	0.94015\\
3.4227	0.94424\\
3.4193	0.94827\\
3.416	0.95224\\
3.4127	0.95614\\
3.4094	0.95998\\
3.4061	0.96375\\
3.4028	0.96746\\
3.3995	0.97109\\
3.3962	0.97466\\
3.3928	0.97816\\
3.3895	0.98158\\
3.3862	0.98493\\
3.3829	0.98821\\
3.3796	0.9914\\
3.3763	0.99453\\
3.373	0.99757\\
3.3697	1.0005\\
3.3664	1.0034\\
3.363	1.0062\\
3.3597	1.0089\\
3.3564	1.0116\\
3.3531	1.0141\\
3.3498	1.0166\\
3.3465	1.019\\
3.3432	1.0212\\
3.3399	1.0234\\
3.3366	1.0256\\
3.3332	1.0276\\
3.3299	1.0295\\
3.3266	1.0313\\
3.3233	1.0331\\
3.32	1.0347\\
3.3167	1.0363\\
3.3134	1.0378\\
3.3101	1.0391\\
3.3068	1.0404\\
3.3034	1.0416\\
3.3001	1.0427\\
3.2968	1.0437\\
3.2935	1.0446\\
3.2902	1.0453\\
3.2869	1.046\\
3.2836	1.0466\\
3.2803	1.0471\\
3.277	1.0475\\
3.2736	1.0478\\
3.2703	1.048\\
3.267	1.0481\\
3.2637	1.0481\\
3.2604	1.048\\
3.2571	1.0478\\
3.2538	1.0475\\
3.2505	1.0471\\
3.2471	1.0466\\
3.2438	1.046\\
3.2405	1.0452\\
3.2372	1.0444\\
3.2339	1.0435\\
3.2306	1.0425\\
3.2273	1.0414\\
3.224	1.0402\\
3.2207	1.0389\\
3.2173	1.0375\\
3.214	1.036\\
3.2107	1.0343\\
3.2074	1.0326\\
3.2041	1.0308\\
3.2008	1.0289\\
3.1975	1.0269\\
3.1942	1.0248\\
3.1909	1.0226\\
3.1875	1.0203\\
3.1842	1.018\\
3.1809	1.0155\\
3.1776	1.0129\\
3.1743	1.0102\\
3.171	1.0075\\
3.1677	1.0046\\
3.1644	1.0017\\
3.1611	0.99863\\
3.1577	0.99551\\
3.1544	0.99229\\
3.1511	0.98899\\
3.1478	0.9856\\
3.1445	0.98212\\
3.1412	0.97855\\
3.1379	0.9749\\
3.1346	0.97117\\
3.1313	0.96735\\
3.1279	0.96346\\
3.1246	0.95948\\
3.1213	0.95542\\
3.118	0.95128\\
3.1147	0.94706\\
3.1114	0.94277\\
3.1081	0.93841\\
3.1048	0.93396\\
3.1014	0.92945\\
3.0981	0.92486\\
3.0948	0.9202\\
3.0915	0.91548\\
3.0882	0.91068\\
3.0849	0.90582\\
3.0816	0.90089\\
3.0783	0.8959\\
3.075	0.89084\\
3.0716	0.88572\\
3.0683	0.88054\\
3.065	0.8753\\
3.0617	0.87\\
3.0584	0.86465\\
3.0551	0.85924\\
3.0518	0.85378\\
3.0485	0.84826\\
3.0452	0.8427\\
3.0418	0.83708\\
3.0385	0.83142\\
3.0352	0.82571\\
3.0319	0.81995\\
3.0286	0.81415\\
3.0253	0.8083\\
3.022	0.80242\\
3.0187	0.79649\\
3.0154	0.79053\\
3.012	0.78453\\
3.0087	0.7785\\
3.0054	0.77243\\
3.0021	0.76632\\
2.9988	0.76019\\
2.9955	0.75403\\
2.9922	0.74784\\
2.9889	0.74162\\
2.9856	0.73538\\
2.9822	0.72911\\
2.9789	0.72282\\
2.9756	0.71651\\
2.9723	0.71018\\
2.969	0.70383\\
2.9657	0.69747\\
2.9624	0.69109\\
2.9591	0.6847\\
2.9557	0.6783\\
2.9524	0.67188\\
2.9491	0.66546\\
2.9458	0.65903\\
2.9425	0.65259\\
2.9392	0.64615\\
2.9359	0.6397\\
2.9326	0.63325\\
2.9293	0.62681\\
2.9259	0.62036\\
2.9226	0.61391\\
2.9193	0.60747\\
2.916	0.60104\\
2.9127	0.59461\\
2.9094	0.58818\\
2.9061	0.58177\\
2.9028	0.57537\\
2.8995	0.56898\\
2.8961	0.5626\\
2.8928	0.55624\\
2.8895	0.54989\\
2.8862	0.54356\\
2.8829	0.53724\\
2.8796	0.53095\\
2.8763	0.52467\\
2.873	0.51842\\
2.8697	0.51219\\
2.8663	0.50598\\
2.863	0.4998\\
2.8597	0.49364\\
2.8564	0.48751\\
2.8531	0.48141\\
2.8498	0.47534\\
2.8465	0.4693\\
2.8432	0.46328\\
2.8398	0.4573\\
2.8365	0.45135\\
2.8332	0.44544\\
2.8299	0.43955\\
2.8266	0.43371\\
2.8233	0.42789\\
2.82	0.42212\\
2.8167	0.41638\\
2.8134	0.41068\\
2.81	0.40502\\
2.8067	0.3994\\
2.8034	0.39382\\
2.8001	0.38827\\
2.7968	0.38277\\
2.7935	0.37731\\
2.7902	0.3719\\
2.7869	0.36652\\
2.7836	0.36119\\
2.7802	0.3559\\
2.7769	0.35066\\
2.7736	0.34546\\
2.7703	0.34031\\
2.767	0.3352\\
2.7637	0.33013\\
2.7604	0.32512\\
2.7571	0.32014\\
2.7538	0.31522\\
2.7504	0.31034\\
2.7471	0.30551\\
2.7438	0.30073\\
2.7405	0.29599\\
2.7372	0.2913\\
2.7339	0.28666\\
2.7306	0.28207\\
2.7273	0.27753\\
2.724	0.27303\\
2.7206	0.26859\\
2.7173	0.26419\\
2.714	0.25984\\
2.7107	0.25554\\
2.7074	0.25128\\
2.7041	0.24708\\
2.7008	0.24292\\
2.6975	0.23882\\
2.6941	0.23476\\
2.6908	0.23075\\
2.6875	0.22679\\
2.6842	0.22288\\
2.6809	0.21901\\
2.6776	0.21519\\
2.6743	0.21143\\
2.671	0.20771\\
2.6677	0.20403\\
2.6643	0.20041\\
2.661	0.19683\\
2.6577	0.1933\\
2.6544	0.18982\\
2.6511	0.18638\\
2.6478	0.18299\\
2.6445	0.17965\\
2.6412	0.17635\\
2.6379	0.1731\\
2.6345	0.1699\\
2.6312	0.16674\\
2.6279	0.16362\\
2.6246	0.16055\\
2.6213	0.15753\\
2.618	0.15455\\
2.6147	0.15161\\
2.6114	0.14872\\
2.6081	0.14587\\
2.6047	0.14306\\
2.6014	0.1403\\
2.5981	0.13758\\
2.5948	0.1349\\
2.5915	0.13226\\
2.5882	0.12966\\
2.5849	0.12711\\
2.5816	0.12459\\
2.5783	0.12212\\
2.5749	0.11968\\
2.5716	0.11729\\
2.5683	0.11493\\
2.565	0.11261\\
2.5617	0.11033\\
2.5584	0.10809\\
2.5551	0.10588\\
2.5518	0.10372\\
2.5484	0.10159\\
2.5451	0.099492\\
2.5418	0.097433\\
2.5385	0.09541\\
2.5352	0.093422\\
2.5319	0.091469\\
2.5286	0.08955\\
2.5253	0.087664\\
2.522	0.085813\\
2.5186	0.083994\\
2.5153	0.082208\\
2.512	0.080454\\
2.5087	0.078732\\
2.5054	0.077041\\
2.5021	0.075381\\
2.4988	0.073752\\
2.4955	0.072153\\
2.4922	0.070583\\
2.4888	0.069043\\
2.4855	0.067532\\
2.4822	0.066049\\
2.4789	0.064595\\
2.4756	0.063168\\
2.4723	0.061769\\
2.469	0.060397\\
2.4657	0.059051\\
2.4624	0.057731\\
2.459	0.056437\\
2.4557	0.055168\\
2.4524	0.053925\\
2.4491	0.052706\\
2.4458	0.051511\\
2.4425	0.05034\\
2.4392	0.049193\\
2.4359	0.048068\\
2.4326	0.046967\\
2.4292	0.045887\\
2.4259	0.04483\\
2.4226	0.043795\\
2.4193	0.04278\\
2.416	0.041787\\
2.4127	0.040814\\
2.4094	0.039862\\
2.4061	0.038929\\
2.4027	0.038016\\
2.3994	0.037122\\
2.3961	0.036247\\
2.3928	0.035391\\
2.3895	0.034553\\
2.3862	0.033733\\
2.3829	0.03293\\
2.3796	0.032145\\
2.3763	0.031376\\
2.3729	0.030625\\
2.3696	0.029889\\
2.3663	0.02917\\
2.363	0.028467\\
2.3597	0.027779\\
2.3564	0.027106\\
2.3531	0.026448\\
2.3498	0.025805\\
2.3465	0.025176\\
2.3431	0.024561\\
2.3398	0.02396\\
2.3365	0.023373\\
2.3332	0.022799\\
2.3299	0.022237\\
2.3266	0.021689\\
2.3233	0.021153\\
2.32	0.02063\\
2.3167	0.020118\\
2.3133	0.019618\\
2.31	0.01913\\
2.3067	0.018653\\
2.3034	0.018187\\
2.3001	0.017732\\
2.2968	0.017288\\
2.2935	0.016854\\
2.2902	0.01643\\
2.2869	0.016017\\
2.2835	0.015613\\
2.2802	0.015218\\
2.2769	0.014833\\
2.2736	0.014458\\
2.2703	0.014091\\
2.267	0.013733\\
2.2637	0.013383\\
2.2604	0.013043\\
2.257	0.01271\\
2.2537	0.012385\\
2.2504	0.012068\\
2.2471	0.011759\\
2.2438	0.011458\\
2.2405	0.011164\\
2.2372	0.010877\\
2.2339	0.010597\\
2.2306	0.010324\\
2.2272	0.010057\\
2.2239	0.0097978\\
2.2206	0.0095446\\
2.2173	0.0092976\\
2.214	0.0090568\\
2.2107	0.0088221\\
2.2074	0.0085932\\
2.2041	0.00837\\
2.2008	0.0081525\\
2.1974	0.0079404\\
2.1941	0.0077336\\
2.1908	0.0075321\\
2.1875	0.0073357\\
2.1842	0.0071443\\
2.1809	0.0069577\\
2.1776	0.0067759\\
2.1743	0.0065987\\
2.171	0.0064261\\
2.1676	0.0062579\\
2.1643	0.0060939\\
2.161	0.0059342\\
2.1577	0.0057786\\
2.1544	0.005627\\
2.1511	0.0054793\\
2.1478	0.0053354\\
2.1445	0.0051953\\
2.1411	0.0050587\\
2.1378	0.0049258\\
2.1345	0.0047962\\
2.1312	0.0046701\\
2.1279	0.0045472\\
2.1246	0.0044275\\
2.1213	0.004311\\
2.118	0.0041975\\
2.1147	0.0040869\\
2.1113	0.0039793\\
2.108	0.0038745\\
2.1047	0.0037724\\
2.1014	0.003673\\
2.0981	0.0035763\\
2.0948	0.003482\\
2.0915	0.0033903\\
2.0882	0.003301\\
2.0849	0.003214\\
2.0815	0.0031294\\
2.0782	0.003047\\
2.0749	0.0029667\\
2.0716	0.0028886\\
2.0683	0.0028126\\
2.065	0.0027385\\
2.0617	0.0026665\\
2.0584	0.0025963\\
2.0551	0.002528\\
2.0517	0.0024616\\
2.0484	0.0023968\\
2.0451	0.0023339\\
2.0418	0.0022726\\
2.0385	0.0022129\\
2.0352	0.0021548\\
2.0319	0.0020983\\
2.0286	0.0020432\\
2.0253	0.0019897\\
2.0219	0.0019375\\
2.0186	0.0018868\\
2.0153	0.0018374\\
2.012	0.0017893\\
2.0087	0.0017425\\
2.0054	0.001697\\
2.0021	0.0016526\\
}--cycle;
\addplot[ybar interval, fill=black, fill opacity=0.35, area legend, draw=none] table[row sep=crcr, x=Lower, y=Count] {%
Lower	Upper	Count\\
2.0021	2.2565	0.020959\\
2.2565	2.511	0.15719\\
2.511	2.7655	0.28295\\
2.7655	3.0199	0.60781\\
3.0199	3.2744	0.99555\\
3.2744	3.5289	0.80692\\
3.5289	3.7833	0.53445\\
3.7833	4.0378	0.28295\\
4.0378	4.2923	0.11527\\
4.2923	4.5467	0.062877\\
4.5467	4.8012	0.041918\\
4.8012	5.0557	0\\
5.0557	5.3101	0.020959\\
5.3101	5.3101	0.020959\\
};
\addlegendentry{\RealHistogramLegend}

\addplot [color=black]
  table[row sep=crcr]{%
2.0021	0.0099712\\
2.0451	0.01362\\
2.0815	0.017578\\
2.1147	0.022014\\
2.1478	0.027389\\
2.1776	0.033152\\
2.2074	0.039914\\
2.2339	0.046864\\
2.2604	0.054792\\
2.2835	0.062606\\
2.3067	0.071303\\
2.3299	0.080946\\
2.3531	0.091598\\
2.3763	0.10332\\
2.3994	0.11616\\
2.4226	0.13018\\
2.4458	0.14542\\
2.469	0.16192\\
2.4922	0.17971\\
2.5153	0.19882\\
2.5385	0.21924\\
2.5617	0.24099\\
2.5849	0.26404\\
2.6081	0.28837\\
2.6312	0.31392\\
2.6577	0.34454\\
2.6842	0.37657\\
2.714	0.41407\\
2.7471	0.45728\\
2.7869	0.51068\\
2.9061	0.67208\\
2.9359	0.71035\\
2.9624	0.74289\\
2.9856	0.76993\\
3.0054	0.79188\\
3.0253	0.81253\\
3.0452	0.83178\\
3.0617	0.84664\\
3.0783	0.86037\\
3.0948	0.87291\\
3.1114	0.88419\\
3.1279	0.89416\\
3.1412	0.90117\\
3.1544	0.9073\\
3.1677	0.91252\\
3.1809	0.91683\\
3.1942	0.9202\\
3.2074	0.92264\\
3.2207	0.92413\\
3.2339	0.92468\\
3.2471	0.92428\\
3.2604	0.92293\\
3.2736	0.92064\\
3.2869	0.91742\\
3.3001	0.91329\\
3.3134	0.90825\\
3.3266	0.90232\\
3.3399	0.89553\\
3.3564	0.88585\\
3.373	0.8749\\
3.3895	0.86274\\
3.4061	0.84943\\
3.4227	0.83503\\
3.4425	0.81641\\
3.4624	0.79645\\
3.4823	0.77527\\
3.5054	0.74922\\
3.5319	0.71795\\
3.5617	0.68126\\
3.5982	0.63492\\
3.7306	0.46489\\
3.7637	0.42469\\
3.7935	0.38999\\
3.82	0.36051\\
3.8465	0.33244\\
3.8697	0.30911\\
3.8929	0.28697\\
3.916	0.26606\\
3.9392	0.24638\\
3.9624	0.22793\\
3.9856	0.2107\\
4.0088	0.19467\\
4.0319	0.17981\\
4.0551	0.16608\\
4.0783	0.15343\\
4.1015	0.14181\\
4.1247	0.13116\\
4.1478	0.12143\\
4.1743	0.11134\\
4.2008	0.10228\\
4.2273	0.094159\\
4.2571	0.086023\\
4.2869	0.078834\\
4.32	0.071818\\
4.3565	0.065114\\
4.3962	0.058813\\
4.4392	0.052956\\
4.4889	0.047182\\
4.5452	0.041602\\
4.6114	0.035985\\
4.6909	0.03019\\
4.7836	0.024358\\
4.8863	0.018817\\
4.9922	0.01401\\
5.1015	0.0099863\\
5.2141	0.0067792\\
5.3101	0.0047195\\
};
\addlegendentry{\RealBLNFitLegend}

      % \input{../../../../.tex/20/KS/SizeDistributionFittingsfig12.tex}
      \legend{}
    \end{groupplot}
  
    \begin{groupplot}[
      group style={
        group name=Row3,
        plotGroup,
      },
      plotRow,
      row1Keys,
      yTickScaleLabelEmpty,
      % emptyLeged,
      %
      % xmin=1,xmax=2.3788e+05,
      ymin=0,ymax=1,
    ]
      \nextgroupplot[%
        at={($(Row2 c1r1.south west)-(0,\yPlotSep em)-(0,\plotHeight cm)$)},
        xmin=0.57363,xmax=5.4115,
        % xmin=1.2631,xmax=5.5243,
        % ymin=0,ymax=1,
      ] % This file was created by matlab2tikz.
%
\definecolor{mycolor1}{rgb}{0.00000,0.44706,0.69804}%
\definecolor{mycolor2}{rgb}{0.83529,0.36863,0.00000}%
%
\addplot[ybar interval, fill=mycolor1, fill opacity=0.35, area legend, draw=none] table[row sep=crcr, x=Lower, y=Count] {%
Lower	Upper	Count\\
1.4568	1.7548	0.00035796\\
1.7548	2.0528	0.0035796\\
2.0528	2.3508	0.033111\\
2.3508	2.6488	0.14641\\
2.6488	2.9467	0.45926\\
2.9467	3.2447	0.84317\\
3.2447	3.5427	0.88864\\
3.5427	3.8407	0.53434\\
3.8407	4.1387	0.258\\
4.1387	4.4367	0.12466\\
4.4367	4.7347	0.044656\\
4.7347	5.0326	0.011007\\
5.0326	5.3306	0.0086805\\
5.3306	5.3306	0.0086805\\
};
\addlegendentry{\ExactHistogramLegend}

\addplot [color=mycolor1, only marks, every error bar/.append style={opacity=0.45}, mark=*, mark size=0pt, draw=none, forget plot]
 plot [error bars/.cd, y dir=both, y explicit, error bar style={line width=1pt, color=mycolor1}, error mark options={mark=none,mark size=0pt}]
 table[row sep=crcr, y error plus index=2, y error minus index=3]{%
1.6058	0.00035796	0.00034971	0.00034971\\
1.9038	0.0035796	0.0011566	0.0011566\\
2.2018	0.033111	0.0031974	0.0031974\\
2.4998	0.14641	0.0066128	0.0066128\\
2.7978	0.45926	0.0096167	0.0096167\\
3.0957	0.84317	0.012799	0.012799\\
3.3937	0.88864	0.012893	0.012893\\
3.6917	0.53434	0.009257	0.009257\\
3.9897	0.258	0.0082017	0.0082017\\
4.2877	0.12466	0.004889	0.004889\\
4.5857	0.044656	0.0029701	0.0029701\\
4.8836	0.011007	0.0013079	0.0013079\\
5.1816	0.0086805	0.00030443	0.00030443\\
};
\addplot [color=mycolor1]
  table[row sep=crcr]{%
1.4568	0.00029246\\
1.6623	0.0010498\\
1.7864	0.002223\\
1.8756	0.0037956\\
1.9454	0.0057631\\
2.0036	0.0081575\\
2.054	0.011016\\
2.0966	0.014192\\
2.1354	0.017846\\
2.1703	0.021906\\
2.2013	0.02625\\
2.2324	0.031409\\
2.2595	0.036696\\
2.2866	0.042806\\
2.3099	0.04878\\
2.3332	0.055509\\
2.3564	0.063069\\
2.3797	0.07154\\
2.3991	0.079353\\
2.4185	0.087904\\
2.4379	0.097241\\
2.4573	0.10741\\
2.4767	0.11847\\
2.496	0.13045\\
2.5154	0.1434\\
2.5348	0.15737\\
2.5542	0.17239\\
2.5736	0.1885\\
2.593	0.20571\\
2.6124	0.22405\\
2.6318	0.24353\\
2.6511	0.26416\\
2.6705	0.28591\\
2.6899	0.30879\\
2.7093	0.33274\\
2.7326	0.36286\\
2.7558	0.39438\\
2.7791	0.42717\\
2.8063	0.46682\\
2.8373	0.51361\\
2.8799	0.57955\\
2.942	0.67559\\
2.9691	0.71635\\
2.9924	0.75008\\
3.0118	0.77708\\
3.0312	0.80284\\
3.0505	0.82716\\
3.0661	0.84546\\
3.0816	0.86261\\
3.0971	0.87853\\
3.1126	0.89313\\
3.1242	0.90316\\
3.1359	0.91238\\
3.1475	0.92075\\
3.1591	0.92825\\
3.1708	0.93486\\
3.1824	0.94056\\
3.194	0.94534\\
3.2057	0.94919\\
3.2173	0.9521\\
3.2289	0.95407\\
3.2406	0.95511\\
3.2522	0.9552\\
3.2638	0.95438\\
3.2755	0.95264\\
3.2871	0.95\\
3.2987	0.94649\\
3.3104	0.94212\\
3.322	0.93692\\
3.3336	0.93091\\
3.3453	0.92413\\
3.3569	0.9166\\
3.3724	0.90546\\
3.3879	0.89314\\
3.4034	0.87972\\
3.4189	0.86528\\
3.4344	0.84991\\
3.4538	0.82955\\
3.4732	0.80805\\
3.4965	0.781\\
3.5236	0.7481\\
3.5546	0.70925\\
3.6167	0.62992\\
3.6632	0.57108\\
3.6942	0.53297\\
3.7214	0.50067\\
3.7485	0.46952\\
3.7757	0.43967\\
3.7989	0.41518\\
3.8222	0.39175\\
3.8455	0.36941\\
3.8687	0.34815\\
3.892	0.32798\\
3.9153	0.30888\\
3.9385	0.29082\\
3.9618	0.27377\\
3.9851	0.25769\\
4.0083	0.24254\\
4.0316	0.22827\\
4.0549	0.21483\\
4.0781	0.20218\\
4.1053	0.18833\\
4.1324	0.17542\\
4.1596	0.16335\\
4.1867	0.15207\\
4.2177	0.14005\\
4.2488	0.12888\\
4.2798	0.11849\\
4.3147	0.10765\\
4.3496	0.097612\\
4.3845	0.08832\\
4.4232	0.078797\\
4.462	0.070053\\
4.5047	0.061271\\
4.5473	0.053311\\
4.59	0.046124\\
4.6326	0.039668\\
4.6792	0.033413\\
4.7257	0.027928\\
4.7761	0.022795\\
4.8265	0.018433\\
4.8808	0.014514\\
4.939	0.011102\\
5.001	0.0082322\\
5.0708	0.0057877\\
5.1484	0.0038387\\
5.2414	0.002286\\
5.3306	0.0013565\\
};
\addlegendentry{\ExactBLNFitLegend}


\addplot[area legend, draw=none, fill=mycolor1, fill opacity=0.15, forget plot]
table[row sep=crcr] {%
x	y\\
1.4568	0.00036823\\
1.4607	0.00037726\\
1.4646	0.00038649\\
1.4685	0.00039594\\
1.4723	0.0004056\\
1.4762	0.00041549\\
1.4801	0.00042561\\
1.484	0.00043596\\
1.4878	0.00044656\\
1.4917	0.00045739\\
1.4956	0.00046848\\
1.4995	0.00047982\\
1.5034	0.00049142\\
1.5072	0.00050329\\
1.5111	0.00051543\\
1.515	0.00052785\\
1.5189	0.00054055\\
1.5227	0.00055355\\
1.5266	0.00056684\\
1.5305	0.00058043\\
1.5344	0.00059434\\
1.5383	0.00060856\\
1.5421	0.00062311\\
1.546	0.00063799\\
1.5499	0.00065321\\
1.5538	0.00066877\\
1.5576	0.00068468\\
1.5615	0.00070096\\
1.5654	0.00071761\\
1.5693	0.00073463\\
1.5732	0.00075204\\
1.577	0.00076984\\
1.5809	0.00078805\\
1.5848	0.00080667\\
1.5887	0.0008257\\
1.5925	0.00084517\\
1.5964	0.00086507\\
1.6003	0.00088543\\
1.6042	0.00090624\\
1.6081	0.00092752\\
1.6119	0.00094928\\
1.6158	0.00097152\\
1.6197	0.00099427\\
1.6236	0.0010175\\
1.6274	0.0010413\\
1.6313	0.0010656\\
1.6352	0.0010905\\
1.6391	0.0011159\\
1.643	0.0011419\\
1.6468	0.0011684\\
1.6507	0.0011956\\
1.6546	0.0012234\\
1.6585	0.0012517\\
1.6623	0.0012808\\
1.6662	0.0013104\\
1.6701	0.0013408\\
1.674	0.0013718\\
1.6778	0.0014035\\
1.6817	0.0014359\\
1.6856	0.001469\\
1.6895	0.0015028\\
1.6934	0.0015375\\
1.6972	0.0015728\\
1.7011	0.001609\\
1.705	0.001646\\
1.7089	0.0016838\\
1.7127	0.0017224\\
1.7166	0.0017619\\
1.7205	0.0018023\\
1.7244	0.0018436\\
1.7283	0.0018858\\
1.7321	0.0019289\\
1.736	0.001973\\
1.7399	0.0020181\\
1.7438	0.0020641\\
1.7476	0.0021112\\
1.7515	0.0021594\\
1.7554	0.0022086\\
1.7593	0.0022589\\
1.7632	0.0023103\\
1.767	0.0023628\\
1.7709	0.0024165\\
1.7748	0.0024714\\
1.7787	0.0025276\\
1.7825	0.0025849\\
1.7864	0.0026436\\
1.7903	0.0027035\\
1.7942	0.0027648\\
1.7981	0.0028274\\
1.8019	0.0028914\\
1.8058	0.0029569\\
1.8097	0.0030238\\
1.8136	0.0030922\\
1.8174	0.003162\\
1.8213	0.0032335\\
1.8252	0.0033065\\
1.8291	0.0033812\\
1.833	0.0034575\\
1.8368	0.0035355\\
1.8407	0.0036152\\
1.8446	0.0036967\\
1.8485	0.00378\\
1.8523	0.0038651\\
1.8562	0.0039522\\
1.8601	0.0040412\\
1.864	0.0041321\\
1.8679	0.0042251\\
1.8717	0.0043201\\
1.8756	0.0044172\\
1.8795	0.0045165\\
1.8834	0.004618\\
1.8872	0.0047218\\
1.8911	0.0048279\\
1.895	0.0049363\\
1.8989	0.0050471\\
1.9028	0.0051604\\
1.9066	0.0052762\\
1.9105	0.0053946\\
1.9144	0.0055156\\
1.9183	0.0056393\\
1.9221	0.0057657\\
1.926	0.005895\\
1.9299	0.0060271\\
1.9338	0.0061621\\
1.9377	0.0063002\\
1.9415	0.0064413\\
1.9454	0.0065856\\
1.9493	0.006733\\
1.9532	0.0068838\\
1.957	0.0070379\\
1.9609	0.0071954\\
1.9648	0.0073564\\
1.9687	0.007521\\
1.9726	0.0076892\\
1.9764	0.0078612\\
1.9803	0.008037\\
1.9842	0.0082167\\
1.9881	0.0084004\\
1.9919	0.0085882\\
1.9958	0.0087801\\
1.9997	0.0089763\\
2.0036	0.0091769\\
2.0075	0.0093818\\
2.0113	0.0095914\\
2.0152	0.0098056\\
2.0191	0.010024\\
2.023	0.010248\\
2.0268	0.010477\\
2.0307	0.010711\\
2.0346	0.01095\\
2.0385	0.011194\\
2.0424	0.011444\\
2.0462	0.011699\\
2.0501	0.01196\\
2.054	0.012226\\
2.0579	0.012499\\
2.0617	0.012777\\
2.0656	0.013062\\
2.0695	0.013353\\
2.0734	0.01365\\
2.0773	0.013954\\
2.0811	0.014264\\
2.085	0.014581\\
2.0889	0.014906\\
2.0928	0.015237\\
2.0966	0.015576\\
2.1005	0.015922\\
2.1044	0.016275\\
2.1083	0.016637\\
2.1121	0.017006\\
2.116	0.017383\\
2.1199	0.017769\\
2.1238	0.018163\\
2.1277	0.018565\\
2.1315	0.018976\\
2.1354	0.019397\\
2.1393	0.019826\\
2.1432	0.020264\\
2.147	0.020713\\
2.1509	0.02117\\
2.1548	0.021638\\
2.1587	0.022116\\
2.1626	0.022604\\
2.1664	0.023102\\
2.1703	0.023611\\
2.1742	0.024131\\
2.1781	0.024662\\
2.1819	0.025205\\
2.1858	0.025759\\
2.1897	0.026325\\
2.1936	0.026903\\
2.1975	0.027493\\
2.2013	0.028096\\
2.2052	0.028711\\
2.2091	0.029339\\
2.213	0.029981\\
2.2168	0.030636\\
2.2207	0.031305\\
2.2246	0.031988\\
2.2285	0.032685\\
2.2324	0.033397\\
2.2362	0.034123\\
2.2401	0.034865\\
2.244	0.035621\\
2.2479	0.036394\\
2.2517	0.037182\\
2.2556	0.037987\\
2.2595	0.038808\\
2.2634	0.039646\\
2.2673	0.040501\\
2.2711	0.041374\\
2.275	0.042264\\
2.2789	0.043172\\
2.2828	0.044099\\
2.2866	0.045044\\
2.2905	0.046008\\
2.2944	0.046992\\
2.2983	0.047995\\
2.3022	0.049018\\
2.306	0.050062\\
2.3099	0.051126\\
2.3138	0.052211\\
2.3177	0.053317\\
2.3215	0.054445\\
2.3254	0.055596\\
2.3293	0.056768\\
2.3332	0.057963\\
2.3371	0.059182\\
2.3409	0.060424\\
2.3448	0.061689\\
2.3487	0.062979\\
2.3526	0.064293\\
2.3564	0.065633\\
2.3603	0.066998\\
2.3642	0.068388\\
2.3681	0.069804\\
2.372	0.071247\\
2.3758	0.072717\\
2.3797	0.074214\\
2.3836	0.075739\\
2.3875	0.077291\\
2.3913	0.078872\\
2.3952	0.080482\\
2.3991	0.082121\\
2.403	0.083789\\
2.4069	0.085488\\
2.4107	0.087216\\
2.4146	0.088976\\
2.4185	0.090767\\
2.4224	0.092589\\
2.4262	0.094443\\
2.4301	0.09633\\
2.434	0.098249\\
2.4379	0.1002\\
2.4418	0.10219\\
2.4456	0.10421\\
2.4495	0.10626\\
2.4534	0.10835\\
2.4573	0.11047\\
2.4611	0.11263\\
2.465	0.11483\\
2.4689	0.11706\\
2.4728	0.11933\\
2.4767	0.12163\\
2.4805	0.12397\\
2.4844	0.12635\\
2.4883	0.12877\\
2.4922	0.13123\\
2.496	0.13372\\
2.4999	0.13626\\
2.5038	0.13883\\
2.5077	0.14144\\
2.5116	0.1441\\
2.5154	0.14679\\
2.5193	0.14952\\
2.5232	0.1523\\
2.5271	0.15512\\
2.5309	0.15798\\
2.5348	0.16088\\
2.5387	0.16382\\
2.5426	0.16681\\
2.5464	0.16984\\
2.5503	0.17291\\
2.5542	0.17602\\
2.5581	0.17918\\
2.562	0.18238\\
2.5658	0.18563\\
2.5697	0.18892\\
2.5736	0.19226\\
2.5775	0.19564\\
2.5813	0.19906\\
2.5852	0.20253\\
2.5891	0.20605\\
2.593	0.20961\\
2.5969	0.21322\\
2.6007	0.21687\\
2.6046	0.22057\\
2.6085	0.22431\\
2.6124	0.2281\\
2.6162	0.23194\\
2.6201	0.23582\\
2.624	0.23975\\
2.6279	0.24372\\
2.6318	0.24774\\
2.6356	0.25181\\
2.6395	0.25592\\
2.6434	0.26008\\
2.6473	0.26429\\
2.6511	0.26854\\
2.655	0.27283\\
2.6589	0.27718\\
2.6628	0.28156\\
2.6667	0.286\\
2.6705	0.29048\\
2.6744	0.295\\
2.6783	0.29957\\
2.6822	0.30418\\
2.686	0.30884\\
2.6899	0.31354\\
2.6938	0.31829\\
2.6977	0.32308\\
2.7016	0.32791\\
2.7054	0.33279\\
2.7093	0.33771\\
2.7132	0.34267\\
2.7171	0.34767\\
2.7209	0.35271\\
2.7248	0.3578\\
2.7287	0.36292\\
2.7326	0.36809\\
2.7365	0.37329\\
2.7403	0.37853\\
2.7442	0.38381\\
2.7481	0.38913\\
2.752	0.39449\\
2.7558	0.39988\\
2.7597	0.40531\\
2.7636	0.41077\\
2.7675	0.41626\\
2.7714	0.42179\\
2.7752	0.42736\\
2.7791	0.43295\\
2.783	0.43857\\
2.7869	0.44423\\
2.7907	0.44991\\
2.7946	0.45563\\
2.7985	0.46137\\
2.8024	0.46713\\
2.8063	0.47293\\
2.8101	0.47874\\
2.814	0.48459\\
2.8179	0.49045\\
2.8218	0.49633\\
2.8256	0.50224\\
2.8295	0.50816\\
2.8334	0.51411\\
2.8373	0.52006\\
2.8412	0.52604\\
2.845	0.53203\\
2.8489	0.53803\\
2.8528	0.54405\\
2.8567	0.55008\\
2.8605	0.55611\\
2.8644	0.56216\\
2.8683	0.56821\\
2.8722	0.57427\\
2.8761	0.58033\\
2.8799	0.58639\\
2.8838	0.59246\\
2.8877	0.59853\\
2.8916	0.60459\\
2.8954	0.61066\\
2.8993	0.61672\\
2.9032	0.62277\\
2.9071	0.62882\\
2.911	0.63486\\
2.9148	0.64089\\
2.9187	0.64691\\
2.9226	0.65292\\
2.9265	0.65891\\
2.9303	0.66489\\
2.9342	0.67085\\
2.9381	0.6768\\
2.942	0.68272\\
2.9459	0.68862\\
2.9497	0.6945\\
2.9536	0.70036\\
2.9575	0.70619\\
2.9614	0.71199\\
2.9652	0.71776\\
2.9691	0.72351\\
2.973	0.72922\\
2.9769	0.7349\\
2.9808	0.74054\\
2.9846	0.74615\\
2.9885	0.75172\\
2.9924	0.75725\\
2.9963	0.76274\\
3.0001	0.76819\\
3.004	0.7736\\
3.0079	0.77896\\
3.0118	0.78428\\
3.0156	0.78955\\
3.0195	0.79476\\
3.0234	0.79993\\
3.0273	0.80505\\
3.0312	0.81011\\
3.035	0.81511\\
3.0389	0.82006\\
3.0428	0.82496\\
3.0467	0.82979\\
3.0505	0.83456\\
3.0544	0.83926\\
3.0583	0.84391\\
3.0622	0.84848\\
3.0661	0.85299\\
3.0699	0.85743\\
3.0738	0.8618\\
3.0777	0.8661\\
3.0816	0.87032\\
3.0854	0.87447\\
3.0893	0.87854\\
3.0932	0.88253\\
3.0971	0.88644\\
3.101	0.89027\\
3.1048	0.89401\\
3.1087	0.89767\\
3.1126	0.90125\\
3.1165	0.90473\\
3.1203	0.90813\\
3.1242	0.91143\\
3.1281	0.91465\\
3.132	0.91777\\
3.1359	0.92079\\
3.1397	0.92372\\
3.1436	0.92655\\
3.1475	0.92929\\
3.1514	0.93192\\
3.1552	0.93446\\
3.1591	0.93689\\
3.163	0.93922\\
3.1669	0.94145\\
3.1708	0.94357\\
3.1746	0.94559\\
3.1785	0.9475\\
3.1824	0.94931\\
3.1863	0.95101\\
3.1901	0.9526\\
3.194	0.95408\\
3.1979	0.95546\\
3.2018	0.95673\\
3.2057	0.95789\\
3.2095	0.95894\\
3.2134	0.95988\\
3.2173	0.96072\\
3.2212	0.96144\\
3.225	0.96206\\
3.2289	0.96257\\
3.2328	0.96297\\
3.2367	0.96326\\
3.2406	0.96344\\
3.2444	0.96352\\
3.2483	0.96349\\
3.2522	0.96336\\
3.2561	0.96312\\
3.2599	0.96277\\
3.2638	0.96233\\
3.2677	0.96178\\
3.2716	0.96112\\
3.2755	0.96037\\
3.2793	0.95951\\
3.2832	0.95856\\
3.2871	0.95751\\
3.291	0.95636\\
3.2948	0.95511\\
3.2987	0.95377\\
3.3026	0.95234\\
3.3065	0.95081\\
3.3104	0.94919\\
3.3142	0.94749\\
3.3181	0.94569\\
3.322	0.94381\\
3.3259	0.94184\\
3.3297	0.93978\\
3.3336	0.93764\\
3.3375	0.93542\\
3.3414	0.93312\\
3.3453	0.93074\\
3.3491	0.92828\\
3.353	0.92574\\
3.3569	0.92313\\
3.3608	0.92044\\
3.3646	0.91768\\
3.3685	0.91485\\
3.3724	0.91194\\
3.3763	0.90896\\
3.3802	0.90592\\
3.384	0.90281\\
3.3879	0.89963\\
3.3918	0.89638\\
3.3957	0.89307\\
3.3995	0.8897\\
3.4034	0.88627\\
3.4073	0.88277\\
3.4112	0.87922\\
3.4151	0.8756\\
3.4189	0.87193\\
3.4228	0.8682\\
3.4267	0.86441\\
3.4306	0.86057\\
3.4344	0.85668\\
3.4383	0.85274\\
3.4422	0.84874\\
3.4461	0.8447\\
3.4499	0.8406\\
3.4538	0.83646\\
3.4577	0.83228\\
3.4616	0.82805\\
3.4655	0.82377\\
3.4693	0.81946\\
3.4732	0.8151\\
3.4771	0.8107\\
3.481	0.80627\\
3.4848	0.8018\\
3.4887	0.79729\\
3.4926	0.79275\\
3.4965	0.78818\\
3.5004	0.78357\\
3.5042	0.77893\\
3.5081	0.77427\\
3.512	0.76957\\
3.5159	0.76485\\
3.5197	0.76011\\
3.5236	0.75534\\
3.5275	0.75054\\
3.5314	0.74573\\
3.5353	0.7409\\
3.5391	0.73604\\
3.543	0.73117\\
3.5469	0.72628\\
3.5508	0.72138\\
3.5546	0.71646\\
3.5585	0.71153\\
3.5624	0.70659\\
3.5663	0.70164\\
3.5702	0.69668\\
3.574	0.69171\\
3.5779	0.68673\\
3.5818	0.68175\\
3.5857	0.67677\\
3.5895	0.67177\\
3.5934	0.66678\\
3.5973	0.66179\\
3.6012	0.65679\\
3.6051	0.6518\\
3.6089	0.6468\\
3.6128	0.64181\\
3.6167	0.63683\\
3.6206	0.63184\\
3.6244	0.62687\\
3.6283	0.6219\\
3.6322	0.61693\\
3.6361	0.61198\\
3.64	0.60703\\
3.6438	0.6021\\
3.6477	0.59717\\
3.6516	0.59226\\
3.6555	0.58736\\
3.6593	0.58247\\
3.6632	0.57759\\
3.6671	0.57273\\
3.671	0.56789\\
3.6749	0.56306\\
3.6787	0.55825\\
3.6826	0.55346\\
3.6865	0.54869\\
3.6904	0.54393\\
3.6942	0.53919\\
3.6981	0.53448\\
3.702	0.52978\\
3.7059	0.52511\\
3.7098	0.52045\\
3.7136	0.51582\\
3.7175	0.51122\\
3.7214	0.50663\\
3.7253	0.50207\\
3.7291	0.49754\\
3.733	0.49303\\
3.7369	0.48854\\
3.7408	0.48408\\
3.7447	0.47965\\
3.7485	0.47524\\
3.7524	0.47086\\
3.7563	0.46651\\
3.7602	0.46218\\
3.764	0.45788\\
3.7679	0.45361\\
3.7718	0.44937\\
3.7757	0.44515\\
3.7796	0.44097\\
3.7834	0.43681\\
3.7873	0.43269\\
3.7912	0.42859\\
3.7951	0.42452\\
3.7989	0.42048\\
3.8028	0.41648\\
3.8067	0.4125\\
3.8106	0.40855\\
3.8145	0.40463\\
3.8183	0.40074\\
3.8222	0.39689\\
3.8261	0.39306\\
3.83	0.38926\\
3.8338	0.3855\\
3.8377	0.38176\\
3.8416	0.37806\\
3.8455	0.37439\\
3.8494	0.37075\\
3.8532	0.36713\\
3.8571	0.36355\\
3.861	0.36\\
3.8649	0.35648\\
3.8687	0.35299\\
3.8726	0.34953\\
3.8765	0.3461\\
3.8804	0.3427\\
3.8842	0.33934\\
3.8881	0.336\\
3.892	0.33269\\
3.8959	0.32941\\
3.8998	0.32616\\
3.9036	0.32294\\
3.9075	0.31975\\
3.9114	0.31659\\
3.9153	0.31346\\
3.9191	0.31036\\
3.923	0.30728\\
3.9269	0.30424\\
3.9308	0.30122\\
3.9347	0.29823\\
3.9385	0.29527\\
3.9424	0.29234\\
3.9463	0.28944\\
3.9502	0.28656\\
3.954	0.28371\\
3.9579	0.28089\\
3.9618	0.2781\\
3.9657	0.27533\\
3.9696	0.27259\\
3.9734	0.26987\\
3.9773	0.26718\\
3.9812	0.26452\\
3.9851	0.26188\\
3.9889	0.25927\\
3.9928	0.25668\\
3.9967	0.25412\\
4.0006	0.25158\\
4.0045	0.24907\\
4.0083	0.24658\\
4.0122	0.24412\\
4.0161	0.24168\\
4.02	0.23926\\
4.0238	0.23687\\
4.0277	0.2345\\
4.0316	0.23215\\
4.0355	0.22983\\
4.0394	0.22753\\
4.0432	0.22525\\
4.0471	0.22299\\
4.051	0.22075\\
4.0549	0.21854\\
4.0587	0.21635\\
4.0626	0.21418\\
4.0665	0.21203\\
4.0704	0.2099\\
4.0743	0.20779\\
4.0781	0.2057\\
4.082	0.20363\\
4.0859	0.20158\\
4.0898	0.19956\\
4.0936	0.19755\\
4.0975	0.19556\\
4.1014	0.19359\\
4.1053	0.19163\\
4.1092	0.1897\\
4.113	0.18779\\
4.1169	0.18589\\
4.1208	0.18401\\
4.1247	0.18215\\
4.1285	0.18031\\
4.1324	0.17848\\
4.1363	0.17667\\
4.1402	0.17488\\
4.1441	0.17311\\
4.1479	0.17135\\
4.1518	0.16961\\
4.1557	0.16788\\
4.1596	0.16617\\
4.1634	0.16448\\
4.1673	0.1628\\
4.1712	0.16114\\
4.1751	0.15949\\
4.179	0.15786\\
4.1828	0.15625\\
4.1867	0.15465\\
4.1906	0.15306\\
4.1945	0.15149\\
4.1983	0.14993\\
4.2022	0.14839\\
4.2061	0.14686\\
4.21	0.14535\\
4.2139	0.14385\\
4.2177	0.14236\\
4.2216	0.14089\\
4.2255	0.13943\\
4.2294	0.13798\\
4.2332	0.13655\\
4.2371	0.13513\\
4.241	0.13372\\
4.2449	0.13233\\
4.2488	0.13094\\
4.2526	0.12958\\
4.2565	0.12822\\
4.2604	0.12687\\
4.2643	0.12554\\
4.2681	0.12422\\
4.272	0.12291\\
4.2759	0.12162\\
4.2798	0.12033\\
4.2837	0.11906\\
4.2875	0.1178\\
4.2914	0.11655\\
4.2953	0.11531\\
4.2992	0.11408\\
4.303	0.11287\\
4.3069	0.11166\\
4.3108	0.11047\\
4.3147	0.10928\\
4.3185	0.10811\\
4.3224	0.10695\\
4.3263	0.10579\\
4.3302	0.10465\\
4.3341	0.10352\\
4.3379	0.1024\\
4.3418	0.10129\\
4.3457	0.10019\\
4.3496	0.099095\\
4.3534	0.098013\\
4.3573	0.096941\\
4.3612	0.095878\\
4.3651	0.094825\\
4.369	0.09378\\
4.3728	0.092746\\
4.3767	0.09172\\
4.3806	0.090703\\
4.3845	0.089696\\
4.3883	0.088697\\
4.3922	0.087707\\
4.3961	0.086726\\
4.4	0.085753\\
4.4039	0.084789\\
4.4077	0.083833\\
4.4116	0.082886\\
4.4155	0.081947\\
4.4194	0.081016\\
4.4232	0.080093\\
4.4271	0.079178\\
4.431	0.078271\\
4.4349	0.077372\\
4.4388	0.07648\\
4.4426	0.075597\\
4.4465	0.074721\\
4.4504	0.073852\\
4.4543	0.072991\\
4.4581	0.072138\\
4.462	0.071291\\
4.4659	0.070452\\
4.4698	0.069621\\
4.4737	0.068796\\
4.4775	0.067979\\
4.4814	0.067169\\
4.4853	0.066365\\
4.4892	0.065569\\
4.493	0.064779\\
4.4969	0.063997\\
4.5008	0.063221\\
4.5047	0.062452\\
4.5086	0.061689\\
4.5124	0.060934\\
4.5163	0.060184\\
4.5202	0.059442\\
4.5241	0.058706\\
4.5279	0.057976\\
4.5318	0.057253\\
4.5357	0.056536\\
4.5396	0.055826\\
4.5435	0.055122\\
4.5473	0.054424\\
4.5512	0.053733\\
4.5551	0.053048\\
4.559	0.052368\\
4.5628	0.051696\\
4.5667	0.051029\\
4.5706	0.050368\\
4.5745	0.049714\\
4.5784	0.049065\\
4.5822	0.048423\\
4.5861	0.047786\\
4.59	0.047156\\
4.5939	0.046531\\
4.5977	0.045912\\
4.6016	0.045299\\
4.6055	0.044692\\
4.6094	0.044091\\
4.6133	0.043496\\
4.6171	0.042906\\
4.621	0.042322\\
4.6249	0.041744\\
4.6288	0.041172\\
4.6326	0.040605\\
4.6365	0.040044\\
4.6404	0.039489\\
4.6443	0.038939\\
4.6482	0.038395\\
4.652	0.037856\\
4.6559	0.037323\\
4.6598	0.036795\\
4.6637	0.036273\\
4.6675	0.035757\\
4.6714	0.035245\\
4.6753	0.03474\\
4.6792	0.034239\\
4.6831	0.033744\\
4.6869	0.033254\\
4.6908	0.03277\\
4.6947	0.032291\\
4.6986	0.031817\\
4.7024	0.031348\\
4.7063	0.030885\\
4.7102	0.030427\\
4.7141	0.029974\\
4.718	0.029526\\
4.7218	0.029083\\
4.7257	0.028645\\
4.7296	0.028212\\
4.7335	0.027785\\
4.7373	0.027362\\
4.7412	0.026944\\
4.7451	0.026532\\
4.749	0.026124\\
4.7529	0.025721\\
4.7567	0.025323\\
4.7606	0.024929\\
4.7645	0.024541\\
4.7684	0.024157\\
4.7722	0.023778\\
4.7761	0.023404\\
4.78	0.023034\\
4.7839	0.022669\\
4.7877	0.022308\\
4.7916	0.021953\\
4.7955	0.021601\\
4.7994	0.021254\\
4.8033	0.020912\\
4.8071	0.020574\\
4.811	0.020241\\
4.8149	0.019911\\
4.8188	0.019587\\
4.8226	0.019266\\
4.8265	0.01895\\
4.8304	0.018638\\
4.8343	0.01833\\
4.8382	0.018026\\
4.842	0.017727\\
4.8459	0.017431\\
4.8498	0.01714\\
4.8537	0.016852\\
4.8575	0.016569\\
4.8614	0.016289\\
4.8653	0.016014\\
4.8692	0.015742\\
4.8731	0.015474\\
4.8769	0.01521\\
4.8808	0.01495\\
4.8847	0.014693\\
4.8886	0.01444\\
4.8924	0.014191\\
4.8963	0.013945\\
4.9002	0.013703\\
4.9041	0.013464\\
4.908	0.013229\\
4.9118	0.012998\\
4.9157	0.01277\\
4.9196	0.012545\\
4.9235	0.012323\\
4.9273	0.012105\\
4.9312	0.01189\\
4.9351	0.011679\\
4.939	0.01147\\
4.9429	0.011265\\
4.9467	0.011063\\
4.9506	0.010864\\
4.9545	0.010668\\
4.9584	0.010475\\
4.9622	0.010285\\
4.9661	0.010098\\
4.97	0.0099141\\
4.9739	0.0097329\\
4.9778	0.0095545\\
4.9816	0.009379\\
4.9855	0.0092063\\
4.9894	0.0090363\\
4.9933	0.008869\\
4.9971	0.0087044\\
5.001	0.0085424\\
5.0049	0.008383\\
5.0088	0.0082262\\
5.0127	0.008072\\
5.0165	0.0079203\\
5.0204	0.007771\\
5.0243	0.0076242\\
5.0282	0.0074798\\
5.032	0.0073378\\
5.0359	0.0071982\\
5.0398	0.0070608\\
5.0437	0.0069258\\
5.0476	0.006793\\
5.0514	0.0066624\\
5.0553	0.006534\\
5.0592	0.0064078\\
5.0631	0.0062838\\
5.0669	0.0061618\\
5.0708	0.0060419\\
5.0747	0.0059241\\
5.0786	0.0058083\\
5.0825	0.0056945\\
5.0863	0.0055827\\
5.0902	0.0054728\\
5.0941	0.0053648\\
5.098	0.0052586\\
5.1018	0.0051544\\
5.1057	0.005052\\
5.1096	0.0049513\\
5.1135	0.0048525\\
5.1174	0.0047554\\
5.1212	0.00466\\
5.1251	0.0045663\\
5.129	0.0044743\\
5.1329	0.0043839\\
5.1367	0.0042952\\
5.1406	0.0042081\\
5.1445	0.0041225\\
5.1484	0.0040385\\
5.1523	0.003956\\
5.1561	0.003875\\
5.16	0.0037955\\
5.1639	0.0037175\\
5.1678	0.0036409\\
5.1716	0.0035657\\
5.1755	0.0034919\\
5.1794	0.0034195\\
5.1833	0.0033484\\
5.1872	0.0032787\\
5.191	0.0032102\\
5.1949	0.0031431\\
5.1988	0.0030772\\
5.2027	0.0030125\\
5.2065	0.0029491\\
5.2104	0.0028869\\
5.2143	0.0028258\\
5.2182	0.0027659\\
5.222	0.0027072\\
5.2259	0.0026496\\
5.2298	0.0025931\\
5.2337	0.0025377\\
5.2376	0.0024834\\
5.2414	0.0024301\\
5.2453	0.0023779\\
5.2492	0.0023266\\
5.2531	0.0022764\\
5.2569	0.0022272\\
5.2608	0.0021789\\
5.2647	0.0021316\\
5.2686	0.0020852\\
5.2725	0.0020398\\
5.2763	0.0019952\\
5.2802	0.0019515\\
5.2841	0.0019087\\
5.288	0.0018668\\
5.2918	0.0018257\\
5.2957	0.0017854\\
5.2996	0.0017459\\
5.3035	0.0017073\\
5.3074	0.0016694\\
5.3112	0.0016323\\
5.3151	0.0015959\\
5.319	0.0015603\\
5.3229	0.0015254\\
5.3267	0.0014912\\
5.3306	0.0014578\\
5.3306	0.0012552\\
5.3267	0.0012854\\
5.3229	0.0013163\\
5.319	0.0013478\\
5.3151	0.00138\\
5.3112	0.0014129\\
5.3074	0.0014465\\
5.3035	0.0014809\\
5.2996	0.001516\\
5.2957	0.0015519\\
5.2918	0.0015885\\
5.288	0.0016259\\
5.2841	0.0016641\\
5.2802	0.0017032\\
5.2763	0.001743\\
5.2725	0.0017838\\
5.2686	0.0018253\\
5.2647	0.0018678\\
5.2608	0.0019111\\
5.2569	0.0019554\\
5.2531	0.0020006\\
5.2492	0.0020467\\
5.2453	0.0020938\\
5.2414	0.0021419\\
5.2376	0.002191\\
5.2337	0.0022411\\
5.2298	0.0022922\\
5.2259	0.0023444\\
5.222	0.0023976\\
5.2182	0.002452\\
5.2143	0.0025074\\
5.2104	0.002564\\
5.2065	0.0026217\\
5.2027	0.0026805\\
5.1988	0.0027406\\
5.1949	0.0028019\\
5.191	0.0028644\\
5.1872	0.0029281\\
5.1833	0.0029931\\
5.1794	0.0030594\\
5.1755	0.003127\\
5.1716	0.003196\\
5.1678	0.0032663\\
5.1639	0.0033379\\
5.16	0.003411\\
5.1561	0.0034855\\
5.1523	0.0035614\\
5.1484	0.0036389\\
5.1445	0.0037178\\
5.1406	0.0037982\\
5.1367	0.0038801\\
5.1329	0.0039637\\
5.129	0.0040488\\
5.1251	0.0041355\\
5.1212	0.0042238\\
5.1174	0.0043139\\
5.1135	0.0044056\\
5.1096	0.004499\\
5.1057	0.0045942\\
5.1018	0.0046911\\
5.098	0.0047898\\
5.0941	0.0048904\\
5.0902	0.0049928\\
5.0863	0.005097\\
5.0825	0.0052032\\
5.0786	0.0053113\\
5.0747	0.0054214\\
5.0708	0.0055334\\
5.0669	0.0056474\\
5.0631	0.0057635\\
5.0592	0.0058817\\
5.0553	0.0060019\\
5.0514	0.0061243\\
5.0476	0.0062488\\
5.0437	0.0063756\\
5.0398	0.0065045\\
5.0359	0.0066357\\
5.032	0.0067691\\
5.0282	0.0069048\\
5.0243	0.0070429\\
5.0204	0.0071834\\
5.0165	0.0073262\\
5.0127	0.0074715\\
5.0088	0.0076192\\
5.0049	0.0077694\\
5.001	0.0079221\\
4.9971	0.0080774\\
4.9933	0.0082352\\
4.9894	0.0083957\\
4.9855	0.0085587\\
4.9816	0.0087245\\
4.9778	0.008893\\
4.9739	0.0090642\\
4.97	0.0092382\\
4.9661	0.0094149\\
4.9622	0.0095945\\
4.9584	0.009777\\
4.9545	0.0099623\\
4.9506	0.010151\\
4.9467	0.010342\\
4.9429	0.010536\\
4.939	0.010733\\
4.9351	0.010934\\
4.9312	0.011137\\
4.9273	0.011344\\
4.9235	0.011553\\
4.9196	0.011766\\
4.9157	0.011982\\
4.9118	0.012202\\
4.908	0.012424\\
4.9041	0.01265\\
4.9002	0.012879\\
4.8963	0.013112\\
4.8924	0.013348\\
4.8886	0.013588\\
4.8847	0.013831\\
4.8808	0.014077\\
4.8769	0.014327\\
4.8731	0.014581\\
4.8692	0.014838\\
4.8653	0.015099\\
4.8614	0.015364\\
4.8575	0.015632\\
4.8537	0.015904\\
4.8498	0.01618\\
4.8459	0.01646\\
4.842	0.016744\\
4.8382	0.017031\\
4.8343	0.017322\\
4.8304	0.017618\\
4.8265	0.017917\\
4.8226	0.01822\\
4.8188	0.018528\\
4.8149	0.018839\\
4.811	0.019155\\
4.8071	0.019475\\
4.8033	0.019798\\
4.7994	0.020127\\
4.7955	0.020459\\
4.7916	0.020796\\
4.7877	0.021137\\
4.7839	0.021482\\
4.78	0.021832\\
4.7761	0.022186\\
4.7722	0.022545\\
4.7684	0.022908\\
4.7645	0.023276\\
4.7606	0.023648\\
4.7567	0.024025\\
4.7529	0.024407\\
4.749	0.024793\\
4.7451	0.025184\\
4.7412	0.02558\\
4.7373	0.02598\\
4.7335	0.026386\\
4.7296	0.026796\\
4.7257	0.027211\\
4.7218	0.027631\\
4.718	0.028056\\
4.7141	0.028486\\
4.7102	0.028921\\
4.7063	0.029361\\
4.7024	0.029806\\
4.6986	0.030256\\
4.6947	0.030712\\
4.6908	0.031173\\
4.6869	0.031639\\
4.6831	0.03211\\
4.6792	0.032586\\
4.6753	0.033068\\
4.6714	0.033555\\
4.6675	0.034048\\
4.6637	0.034546\\
4.6598	0.03505\\
4.6559	0.035559\\
4.652	0.036074\\
4.6482	0.036594\\
4.6443	0.03712\\
4.6404	0.037651\\
4.6365	0.038189\\
4.6326	0.038732\\
4.6288	0.03928\\
4.6249	0.039835\\
4.621	0.040395\\
4.6171	0.040962\\
4.6133	0.041534\\
4.6094	0.042112\\
4.6055	0.042696\\
4.6016	0.043286\\
4.5977	0.043883\\
4.5939	0.044485\\
4.59	0.045093\\
4.5861	0.045708\\
4.5822	0.046328\\
4.5784	0.046955\\
4.5745	0.047588\\
4.5706	0.048228\\
4.5667	0.048874\\
4.5628	0.049526\\
4.559	0.050184\\
4.5551	0.050849\\
4.5512	0.05152\\
4.5473	0.052198\\
4.5435	0.052882\\
4.5396	0.053573\\
4.5357	0.054271\\
4.5318	0.054975\\
4.5279	0.055685\\
4.5241	0.056403\\
4.5202	0.057127\\
4.5163	0.057858\\
4.5124	0.058595\\
4.5086	0.05934\\
4.5047	0.060091\\
4.5008	0.060849\\
4.4969	0.061614\\
4.493	0.062386\\
4.4892	0.063165\\
4.4853	0.06395\\
4.4814	0.064743\\
4.4775	0.065543\\
4.4737	0.06635\\
4.4698	0.067164\\
4.4659	0.067986\\
4.462	0.068814\\
4.4581	0.06965\\
4.4543	0.070493\\
4.4504	0.071343\\
4.4465	0.072201\\
4.4426	0.073065\\
4.4388	0.073938\\
4.4349	0.074817\\
4.431	0.075704\\
4.4271	0.076599\\
4.4232	0.077501\\
4.4194	0.078411\\
4.4155	0.079328\\
4.4116	0.080253\\
4.4077	0.081186\\
4.4039	0.082126\\
4.4	0.083074\\
4.3961	0.08403\\
4.3922	0.084994\\
4.3883	0.085965\\
4.3845	0.086945\\
4.3806	0.087933\\
4.3767	0.088928\\
4.3728	0.089932\\
4.369	0.090944\\
4.3651	0.091964\\
4.3612	0.092993\\
4.3573	0.094029\\
4.3534	0.095075\\
4.3496	0.096128\\
4.3457	0.097191\\
4.3418	0.098262\\
4.3379	0.099341\\
4.3341	0.10043\\
4.3302	0.10153\\
4.3263	0.10263\\
4.3224	0.10375\\
4.3185	0.10487\\
4.3147	0.10601\\
4.3108	0.10715\\
4.3069	0.1083\\
4.303	0.10947\\
4.2992	0.11064\\
4.2953	0.11182\\
4.2914	0.11301\\
4.2875	0.11422\\
4.2837	0.11543\\
4.2798	0.11665\\
4.2759	0.11788\\
4.272	0.11913\\
4.2681	0.12038\\
4.2643	0.12165\\
4.2604	0.12292\\
4.2565	0.12421\\
4.2526	0.12551\\
4.2488	0.12682\\
4.2449	0.12814\\
4.241	0.12948\\
4.2371	0.13082\\
4.2332	0.13218\\
4.2294	0.13355\\
4.2255	0.13493\\
4.2216	0.13633\\
4.2177	0.13774\\
4.2139	0.13916\\
4.21	0.14059\\
4.2061	0.14204\\
4.2022	0.1435\\
4.1983	0.14498\\
4.1945	0.14647\\
4.1906	0.14797\\
4.1867	0.14949\\
4.1828	0.15102\\
4.179	0.15257\\
4.1751	0.15413\\
4.1712	0.15571\\
4.1673	0.1573\\
4.1634	0.15891\\
4.1596	0.16053\\
4.1557	0.16217\\
4.1518	0.16383\\
4.1479	0.1655\\
4.1441	0.16719\\
4.1402	0.16889\\
4.1363	0.17061\\
4.1324	0.17235\\
4.1285	0.17411\\
4.1247	0.17589\\
4.1208	0.17768\\
4.1169	0.17949\\
4.113	0.18132\\
4.1092	0.18317\\
4.1053	0.18503\\
4.1014	0.18692\\
4.0975	0.18883\\
4.0936	0.19075\\
4.0898	0.1927\\
4.0859	0.19466\\
4.082	0.19664\\
4.0781	0.19865\\
4.0743	0.20068\\
4.0704	0.20272\\
4.0665	0.20479\\
4.0626	0.20688\\
4.0587	0.20899\\
4.0549	0.21112\\
4.051	0.21328\\
4.0471	0.21545\\
4.0432	0.21765\\
4.0394	0.21988\\
4.0355	0.22212\\
4.0316	0.22439\\
4.0277	0.22668\\
4.0238	0.229\\
4.02	0.23134\\
4.0161	0.2337\\
4.0122	0.23609\\
4.0083	0.2385\\
4.0045	0.24094\\
4.0006	0.2434\\
3.9967	0.24589\\
3.9928	0.2484\\
3.9889	0.25094\\
3.9851	0.25351\\
3.9812	0.2561\\
3.9773	0.25872\\
3.9734	0.26136\\
3.9696	0.26403\\
3.9657	0.26673\\
3.9618	0.26945\\
3.9579	0.2722\\
3.954	0.27498\\
3.9502	0.27778\\
3.9463	0.28062\\
3.9424	0.28348\\
3.9385	0.28637\\
3.9347	0.28929\\
3.9308	0.29223\\
3.9269	0.29521\\
3.923	0.29821\\
3.9191	0.30124\\
3.9153	0.3043\\
3.9114	0.30739\\
3.9075	0.31051\\
3.9036	0.31366\\
3.8998	0.31683\\
3.8959	0.32004\\
3.892	0.32328\\
3.8881	0.32654\\
3.8842	0.32984\\
3.8804	0.33316\\
3.8765	0.33652\\
3.8726	0.3399\\
3.8687	0.34332\\
3.8649	0.34676\\
3.861	0.35023\\
3.8571	0.35374\\
3.8532	0.35727\\
3.8494	0.36084\\
3.8455	0.36443\\
3.8416	0.36805\\
3.8377	0.37171\\
3.8338	0.37539\\
3.83	0.37911\\
3.8261	0.38285\\
3.8222	0.38662\\
3.8183	0.39043\\
3.8145	0.39426\\
3.8106	0.39812\\
3.8067	0.40201\\
3.8028	0.40593\\
3.7989	0.40988\\
3.7951	0.41386\\
3.7912	0.41787\\
3.7873	0.42191\\
3.7834	0.42597\\
3.7796	0.43006\\
3.7757	0.43418\\
3.7718	0.43833\\
3.7679	0.44251\\
3.764	0.44672\\
3.7602	0.45095\\
3.7563	0.45521\\
3.7524	0.45949\\
3.7485	0.4638\\
3.7447	0.46814\\
3.7408	0.47251\\
3.7369	0.47689\\
3.733	0.48131\\
3.7291	0.48575\\
3.7253	0.49021\\
3.7214	0.4947\\
3.7175	0.49921\\
3.7136	0.50374\\
3.7098	0.5083\\
3.7059	0.51288\\
3.702	0.51748\\
3.6981	0.5221\\
3.6942	0.52674\\
3.6904	0.53141\\
3.6865	0.53609\\
3.6826	0.54079\\
3.6787	0.54551\\
3.6749	0.55025\\
3.671	0.555\\
3.6671	0.55977\\
3.6632	0.56456\\
3.6593	0.56937\\
3.6555	0.57419\\
3.6516	0.57902\\
3.6477	0.58386\\
3.6438	0.58872\\
3.64	0.59359\\
3.6361	0.59847\\
3.6322	0.60337\\
3.6283	0.60827\\
3.6244	0.61318\\
3.6206	0.6181\\
3.6167	0.62302\\
3.6128	0.62795\\
3.6089	0.63289\\
3.6051	0.63783\\
3.6012	0.64277\\
3.5973	0.64772\\
3.5934	0.65267\\
3.5895	0.65762\\
3.5857	0.66257\\
3.5818	0.66752\\
3.5779	0.67246\\
3.574	0.67741\\
3.5702	0.68234\\
3.5663	0.68728\\
3.5624	0.6922\\
3.5585	0.69712\\
3.5546	0.70203\\
3.5508	0.70693\\
3.5469	0.71182\\
3.543	0.7167\\
3.5391	0.72156\\
3.5353	0.72641\\
3.5314	0.73124\\
3.5275	0.73606\\
3.5236	0.74086\\
3.5197	0.74564\\
3.5159	0.75039\\
3.512	0.75513\\
3.5081	0.75984\\
3.5042	0.76453\\
3.5004	0.7692\\
3.4965	0.77383\\
3.4926	0.77844\\
3.4887	0.78302\\
3.4848	0.78756\\
3.481	0.79208\\
3.4771	0.79655\\
3.4732	0.801\\
3.4693	0.8054\\
3.4655	0.80977\\
3.4616	0.8141\\
3.4577	0.81839\\
3.4538	0.82263\\
3.4499	0.82683\\
3.4461	0.83098\\
3.4422	0.83509\\
3.4383	0.83914\\
3.4344	0.84315\\
3.4306	0.8471\\
3.4267	0.851\\
3.4228	0.85484\\
3.4189	0.85863\\
3.4151	0.86235\\
3.4112	0.86602\\
3.4073	0.86962\\
3.4034	0.87316\\
3.3995	0.87664\\
3.3957	0.88004\\
3.3918	0.88338\\
3.3879	0.88665\\
3.384	0.88985\\
3.3802	0.89297\\
3.3763	0.89602\\
3.3724	0.89899\\
3.3685	0.90188\\
3.3646	0.90469\\
3.3608	0.90742\\
3.3569	0.91007\\
3.353	0.91264\\
3.3491	0.91512\\
3.3453	0.91752\\
3.3414	0.91983\\
3.3375	0.92205\\
3.3336	0.92418\\
3.3297	0.92622\\
3.3259	0.92817\\
3.322	0.93003\\
3.3181	0.9318\\
3.3142	0.93347\\
3.3104	0.93505\\
3.3065	0.93653\\
3.3026	0.93792\\
3.2987	0.93921\\
3.2948	0.9404\\
3.291	0.9415\\
3.2871	0.9425\\
3.2832	0.9434\\
3.2793	0.94421\\
3.2755	0.94491\\
3.2716	0.94552\\
3.2677	0.94602\\
3.2638	0.94643\\
3.2599	0.94674\\
3.2561	0.94694\\
3.2522	0.94705\\
3.2483	0.94706\\
3.2444	0.94696\\
3.2406	0.94677\\
3.2367	0.94647\\
3.2328	0.94608\\
3.2289	0.94558\\
3.225	0.94498\\
3.2212	0.94429\\
3.2173	0.94349\\
3.2134	0.94259\\
3.2095	0.94159\\
3.2057	0.94049\\
3.2018	0.9393\\
3.1979	0.938\\
3.194	0.9366\\
3.1901	0.9351\\
3.1863	0.93351\\
3.1824	0.93182\\
3.1785	0.93002\\
3.1746	0.92813\\
3.1708	0.92615\\
3.1669	0.92406\\
3.163	0.92188\\
3.1591	0.91961\\
3.1552	0.91723\\
3.1514	0.91477\\
3.1475	0.9122\\
3.1436	0.90955\\
3.1397	0.9068\\
3.1359	0.90396\\
3.132	0.90103\\
3.1281	0.898\\
3.1242	0.89489\\
3.1203	0.89168\\
3.1165	0.88839\\
3.1126	0.88501\\
3.1087	0.88154\\
3.1048	0.87798\\
3.101	0.87434\\
3.0971	0.87062\\
3.0932	0.86681\\
3.0893	0.86292\\
3.0854	0.85895\\
3.0816	0.8549\\
3.0777	0.85077\\
3.0738	0.84656\\
3.0699	0.84228\\
3.0661	0.83792\\
3.0622	0.83349\\
3.0583	0.82899\\
3.0544	0.82441\\
3.0505	0.81977\\
3.0467	0.81506\\
3.0428	0.81028\\
3.0389	0.80544\\
3.035	0.80053\\
3.0312	0.79557\\
3.0273	0.79054\\
3.0234	0.78546\\
3.0195	0.78032\\
3.0156	0.77512\\
3.0118	0.76988\\
3.0079	0.76458\\
3.004	0.75923\\
3.0001	0.75384\\
2.9963	0.7484\\
2.9924	0.74291\\
2.9885	0.73739\\
2.9846	0.73182\\
2.9808	0.72622\\
2.9769	0.72058\\
2.973	0.7149\\
2.9691	0.70919\\
2.9652	0.70346\\
2.9614	0.69769\\
2.9575	0.69189\\
2.9536	0.68607\\
2.9497	0.68022\\
2.9459	0.67436\\
2.942	0.66847\\
2.9381	0.66256\\
2.9342	0.65663\\
2.9303	0.65069\\
2.9265	0.64474\\
2.9226	0.63877\\
2.9187	0.63279\\
2.9148	0.6268\\
2.911	0.6208\\
2.9071	0.6148\\
2.9032	0.60879\\
2.8993	0.60278\\
2.8954	0.59676\\
2.8916	0.59075\\
2.8877	0.58473\\
2.8838	0.57872\\
2.8799	0.57271\\
2.8761	0.5667\\
2.8722	0.5607\\
2.8683	0.55471\\
2.8644	0.54872\\
2.8605	0.54275\\
2.8567	0.53678\\
2.8528	0.53083\\
2.8489	0.52489\\
2.845	0.51896\\
2.8412	0.51305\\
2.8373	0.50716\\
2.8334	0.50128\\
2.8295	0.49542\\
2.8256	0.48959\\
2.8218	0.48377\\
2.8179	0.47797\\
2.814	0.4722\\
2.8101	0.46645\\
2.8063	0.46072\\
2.8024	0.45502\\
2.7985	0.44935\\
2.7946	0.4437\\
2.7907	0.43808\\
2.7869	0.43249\\
2.783	0.42693\\
2.7791	0.42139\\
2.7752	0.41589\\
2.7714	0.41043\\
2.7675	0.40499\\
2.7636	0.39959\\
2.7597	0.39422\\
2.7558	0.38888\\
2.752	0.38358\\
2.7481	0.37832\\
2.7442	0.37309\\
2.7403	0.3679\\
2.7365	0.36275\\
2.7326	0.35764\\
2.7287	0.35256\\
2.7248	0.34753\\
2.7209	0.34253\\
2.7171	0.33757\\
2.7132	0.33266\\
2.7093	0.32778\\
2.7054	0.32295\\
2.7016	0.31815\\
2.6977	0.3134\\
2.6938	0.30869\\
2.6899	0.30403\\
2.686	0.29941\\
2.6822	0.29483\\
2.6783	0.29029\\
2.6744	0.2858\\
2.6705	0.28135\\
2.6667	0.27695\\
2.6628	0.27259\\
2.6589	0.26827\\
2.655	0.264\\
2.6511	0.25977\\
2.6473	0.25559\\
2.6434	0.25146\\
2.6395	0.24737\\
2.6356	0.24332\\
2.6318	0.23932\\
2.6279	0.23537\\
2.624	0.23146\\
2.6201	0.22759\\
2.6162	0.22377\\
2.6124	0.22\\
2.6085	0.21627\\
2.6046	0.21259\\
2.6007	0.20895\\
2.5969	0.20536\\
2.593	0.20181\\
2.5891	0.19831\\
2.5852	0.19485\\
2.5813	0.19143\\
2.5775	0.18806\\
2.5736	0.18474\\
2.5697	0.18146\\
2.5658	0.17822\\
2.562	0.17502\\
2.5581	0.17187\\
2.5542	0.16877\\
2.5503	0.1657\\
2.5464	0.16268\\
2.5426	0.1597\\
2.5387	0.15677\\
2.5348	0.15387\\
2.5309	0.15102\\
2.5271	0.14821\\
2.5232	0.14544\\
2.5193	0.14271\\
2.5154	0.14002\\
2.5116	0.13737\\
2.5077	0.13476\\
2.5038	0.1322\\
2.4999	0.12967\\
2.496	0.12718\\
2.4922	0.12473\\
2.4883	0.12231\\
2.4844	0.11994\\
2.4805	0.1176\\
2.4767	0.1153\\
2.4728	0.11304\\
2.4689	0.11081\\
2.465	0.10862\\
2.4611	0.10647\\
2.4573	0.10435\\
2.4534	0.10227\\
2.4495	0.10022\\
2.4456	0.098206\\
2.4418	0.096226\\
2.4379	0.09428\\
2.434	0.092367\\
2.4301	0.090487\\
2.4262	0.08864\\
2.4224	0.086824\\
2.4185	0.085041\\
2.4146	0.083288\\
2.4107	0.081567\\
2.4069	0.079876\\
2.403	0.078216\\
2.3991	0.076585\\
2.3952	0.074984\\
2.3913	0.073412\\
2.3875	0.071868\\
2.3836	0.070353\\
2.3797	0.068865\\
2.3758	0.067405\\
2.372	0.065972\\
2.3681	0.064566\\
2.3642	0.063187\\
2.3603	0.061833\\
2.3564	0.060505\\
2.3526	0.059202\\
2.3487	0.057924\\
2.3448	0.056671\\
2.3409	0.055442\\
2.3371	0.054236\\
2.3332	0.053054\\
2.3293	0.051895\\
2.3254	0.050759\\
2.3215	0.049645\\
2.3177	0.048553\\
2.3138	0.047483\\
2.3099	0.046434\\
2.306	0.045406\\
2.3022	0.044398\\
2.2983	0.043411\\
2.2944	0.042444\\
2.2905	0.041496\\
2.2866	0.040568\\
2.2828	0.039659\\
2.2789	0.038768\\
2.275	0.037896\\
2.2711	0.037041\\
2.2673	0.036205\\
2.2634	0.035386\\
2.2595	0.034583\\
2.2556	0.033798\\
2.2517	0.033029\\
2.2479	0.032276\\
2.244	0.03154\\
2.2401	0.030818\\
2.2362	0.030113\\
2.2324	0.029422\\
2.2285	0.028746\\
2.2246	0.028084\\
2.2207	0.027437\\
2.2168	0.026804\\
2.213	0.026184\\
2.2091	0.025578\\
2.2052	0.024985\\
2.2013	0.024405\\
2.1975	0.023837\\
2.1936	0.023283\\
2.1897	0.02274\\
2.1858	0.022209\\
2.1819	0.02169\\
2.1781	0.021182\\
2.1742	0.020686\\
2.1703	0.020201\\
2.1664	0.019727\\
2.1626	0.019263\\
2.1587	0.018809\\
2.1548	0.018366\\
2.1509	0.017933\\
2.147	0.017509\\
2.1432	0.017096\\
2.1393	0.016691\\
2.1354	0.016296\\
2.1315	0.015909\\
2.1277	0.015532\\
2.1238	0.015163\\
2.1199	0.014803\\
2.116	0.01445\\
2.1121	0.014106\\
2.1083	0.01377\\
2.1044	0.013442\\
2.1005	0.013121\\
2.0966	0.012807\\
2.0928	0.012501\\
2.0889	0.012202\\
2.085	0.01191\\
2.0811	0.011625\\
2.0773	0.011346\\
2.0734	0.011074\\
2.0695	0.010808\\
2.0656	0.010549\\
2.0617	0.010295\\
2.0579	0.010048\\
2.054	0.009806\\
2.0501	0.00957\\
2.0462	0.0093395\\
2.0424	0.0091145\\
2.0385	0.0088948\\
2.0346	0.0086804\\
2.0307	0.008471\\
2.0268	0.0082665\\
2.023	0.008067\\
2.0191	0.0078721\\
2.0152	0.0076819\\
2.0113	0.0074963\\
2.0075	0.007315\\
2.0036	0.0071381\\
1.9997	0.0069654\\
1.9958	0.0067969\\
1.9919	0.0066324\\
1.9881	0.0064718\\
1.9842	0.006315\\
1.9803	0.006162\\
1.9764	0.0060127\\
1.9726	0.005867\\
1.9687	0.0057248\\
1.9648	0.005586\\
1.9609	0.0054506\\
1.957	0.0053184\\
1.9532	0.0051894\\
1.9493	0.0050635\\
1.9454	0.0049406\\
1.9415	0.0048207\\
1.9377	0.0047037\\
1.9338	0.0045895\\
1.9299	0.0044781\\
1.926	0.0043694\\
1.9221	0.0042633\\
1.9183	0.0041598\\
1.9144	0.0040588\\
1.9105	0.0039602\\
1.9066	0.003864\\
1.9028	0.0037702\\
1.8989	0.0036786\\
1.895	0.0035892\\
1.8911	0.0035021\\
1.8872	0.003417\\
1.8834	0.003334\\
1.8795	0.003253\\
1.8756	0.0031739\\
1.8717	0.0030968\\
1.8679	0.0030216\\
1.864	0.0029482\\
1.8601	0.0028765\\
1.8562	0.0028067\\
1.8523	0.0027385\\
1.8485	0.0026719\\
1.8446	0.002607\\
1.8407	0.0025436\\
1.8368	0.0024818\\
1.833	0.0024215\\
1.8291	0.0023627\\
1.8252	0.0023053\\
1.8213	0.0022492\\
1.8174	0.0021946\\
1.8136	0.0021413\\
1.8097	0.0020892\\
1.8058	0.0020384\\
1.8019	0.0019889\\
1.7981	0.0019406\\
1.7942	0.0018934\\
1.7903	0.0018474\\
1.7864	0.0018025\\
1.7825	0.0017586\\
1.7787	0.0017159\\
1.7748	0.0016742\\
1.7709	0.0016335\\
1.767	0.0015937\\
1.7632	0.001555\\
1.7593	0.0015172\\
1.7554	0.0014803\\
1.7515	0.0014443\\
1.7476	0.0014091\\
1.7438	0.0013748\\
1.7399	0.0013414\\
1.736	0.0013087\\
1.7321	0.0012769\\
1.7283	0.0012458\\
1.7244	0.0012155\\
1.7205	0.0011859\\
1.7166	0.001157\\
1.7127	0.0011288\\
1.7089	0.0011013\\
1.705	0.0010745\\
1.7011	0.0010483\\
1.6972	0.0010227\\
1.6934	0.00099779\\
1.6895	0.00097346\\
1.6856	0.00094971\\
1.6817	0.00092653\\
1.6778	0.00090392\\
1.674	0.00088185\\
1.6701	0.00086032\\
1.6662	0.0008393\\
1.6623	0.00081879\\
1.6585	0.00079878\\
1.6546	0.00077924\\
1.6507	0.00076018\\
1.6468	0.00074158\\
1.643	0.00072343\\
1.6391	0.00070572\\
1.6352	0.00068843\\
1.6313	0.00067156\\
1.6274	0.0006551\\
1.6236	0.00063903\\
1.6197	0.00062336\\
1.6158	0.00060806\\
1.6119	0.00059312\\
1.6081	0.00057855\\
1.6042	0.00056433\\
1.6003	0.00055046\\
1.5964	0.00053691\\
1.5925	0.0005237\\
1.5887	0.0005108\\
1.5848	0.00049821\\
1.5809	0.00048593\\
1.577	0.00047394\\
1.5732	0.00046225\\
1.5693	0.00045083\\
1.5654	0.00043969\\
1.5615	0.00042882\\
1.5576	0.00041821\\
1.5538	0.00040786\\
1.5499	0.00039776\\
1.546	0.0003879\\
1.5421	0.00037828\\
1.5383	0.00036889\\
1.5344	0.00035973\\
1.5305	0.00035079\\
1.5266	0.00034207\\
1.5227	0.00033355\\
1.5189	0.00032525\\
1.515	0.00031714\\
1.5111	0.00030923\\
1.5072	0.00030152\\
1.5034	0.00029399\\
1.4995	0.00028664\\
1.4956	0.00027947\\
1.4917	0.00027247\\
1.4878	0.00026564\\
1.484	0.00025898\\
1.4801	0.00025248\\
1.4762	0.00024614\\
1.4723	0.00023996\\
1.4685	0.00023392\\
1.4646	0.00022803\\
1.4607	0.00022228\\
1.4568	0.00021668\\
}--cycle;
\addplot[ybar interval, fill=mycolor2, fill opacity=0.35, area legend, draw=none] table[row sep=crcr, x=Lower, y=Count] {%
Lower	Upper	Count\\
1.4568	1.7548	0.00017898\\
1.7548	2.0528	0.0036691\\
2.0528	2.3508	0.031769\\
2.3508	2.6488	0.14551\\
2.6488	2.9467	0.45729\\
2.9467	3.2447	0.85714\\
3.2447	3.5427	0.87557\\
3.5427	3.8407	0.53954\\
3.8407	4.1387	0.25773\\
4.1387	4.4367	0.12555\\
4.4367	4.7347	0.04206\\
4.7347	5.0326	0.010918\\
5.0326	5.3306	0.008949\\
5.3306	5.3306	0.008949\\
};
\addlegendentry{\ApproxHistogramLegend}

\addplot [color=mycolor2, only marks, every error bar/.append style={opacity=0.45}, mark=*, mark size=0pt, draw=none, forget plot]
 plot [error bars/.cd, y dir=both, y explicit, error bar style={line width=1pt, color=mycolor2}, error mark options={mark=none,mark size=0pt}]
 table[row sep=crcr, y error plus index=2, y error minus index=3]{%
1.6058	0.00017898	0.00024985	0.00017898\\
1.9038	0.0036691	0.0010435	0.0010435\\
2.2018	0.031769	0.0031561	0.0031561\\
2.4998	0.14551	0.0063329	0.0063329\\
2.7978	0.45729	0.010075	0.010075\\
3.0957	0.85714	0.016049	0.016049\\
3.3937	0.87557	0.01499	0.01499\\
3.6917	0.53954	0.0098955	0.0098955\\
3.9897	0.25773	0.0074013	0.0074013\\
4.2877	0.12555	0.0044755	0.0044755\\
4.5857	0.04206	0.0028386	0.0028386\\
4.8836	0.010918	0.0013012	0.0013012\\
5.1816	0.008949	0.00025238	0.00025238\\
};
\addplot [color=mycolor2]
  table[row sep=crcr]{%
1.4568	0.00025648\\
1.6623	0.00099426\\
1.7825	0.002132\\
1.8717	0.0037153\\
1.9415	0.005707\\
1.9997	0.0081316\\
2.0501	0.01102\\
2.0928	0.014217\\
2.1315	0.017883\\
2.1664	0.02194\\
2.1975	0.026264\\
2.2285	0.031381\\
2.2556	0.036608\\
2.2828	0.04263\\
2.306	0.048501\\
2.3293	0.0551\\
2.3526	0.062497\\
2.3758	0.07077\\
2.3952	0.078388\\
2.4146	0.086716\\
2.434	0.095799\\
2.4534	0.10569\\
2.4728	0.11642\\
2.4922	0.12806\\
2.5116	0.14064\\
2.5309	0.1542\\
2.5503	0.16879\\
2.5697	0.18444\\
2.5891	0.20118\\
2.6085	0.21903\\
2.6279	0.23802\\
2.6473	0.25816\\
2.6667	0.27943\\
2.686	0.30184\\
2.7054	0.32536\\
2.7248	0.34995\\
2.7481	0.38082\\
2.7714	0.41307\\
2.7946	0.44655\\
2.8218	0.48695\\
2.8528	0.53449\\
2.9032	0.6134\\
2.9536	0.69195\\
2.9846	0.73882\\
3.0079	0.77267\\
3.0273	0.7997\\
3.0467	0.82535\\
3.0622	0.84468\\
3.0777	0.86275\\
3.0932	0.87941\\
3.1048	0.89089\\
3.1165	0.90146\\
3.1281	0.91107\\
3.1397	0.91972\\
3.1514	0.92739\\
3.163	0.93409\\
3.1746	0.93983\\
3.1863	0.94463\\
3.1979	0.94852\\
3.2095	0.9515\\
3.2212	0.9536\\
3.2328	0.95483\\
3.2444	0.9552\\
3.2561	0.95472\\
3.2677	0.9534\\
3.2793	0.95125\\
3.291	0.94827\\
3.3026	0.94447\\
3.3142	0.93987\\
3.3259	0.93448\\
3.3375	0.92831\\
3.3491	0.92139\\
3.3608	0.91374\\
3.3763	0.90244\\
3.3918	0.88994\\
3.4073	0.87632\\
3.4228	0.86165\\
3.4383	0.84602\\
3.4577	0.82527\\
3.4771	0.80333\\
3.5004	0.77572\\
3.5275	0.74217\\
3.5624	0.69773\\
3.6438	0.59349\\
3.6749	0.5551\\
3.702	0.52256\\
3.7291	0.49116\\
3.7563	0.46101\\
3.7796	0.43623\\
3.8028	0.41245\\
3.8261	0.3897\\
3.8494	0.36797\\
3.8726	0.34726\\
3.8959	0.32756\\
3.9191	0.30885\\
3.9424	0.2911\\
3.9657	0.27427\\
3.9889	0.25834\\
4.0122	0.24326\\
4.0355	0.22899\\
4.0626	0.21332\\
4.0898	0.19864\\
4.1169	0.18488\\
4.1441	0.17198\\
4.1712	0.15989\\
4.1983	0.14855\\
4.2294	0.13643\\
4.2604	0.12516\\
4.2914	0.11467\\
4.3263	0.10374\\
4.3612	0.093645\\
4.3961	0.084338\\
4.431	0.075761\\
4.4698	0.067031\\
4.5086	0.059088\\
4.5473	0.051879\\
4.59	0.044744\\
4.6326	0.038385\\
4.6792	0.032268\\
4.7257	0.026941\\
4.7761	0.021983\\
4.8265	0.017787\\
4.8808	0.014026\\
4.939	0.010756\\
5.0049	0.0078538\\
5.0786	0.0054341\\
5.16	0.0035445\\
5.2569	0.002075\\
5.3306	0.0013557\\
};
\addlegendentry{\ApproxBLNFitLegend}


\addplot[area legend, draw=none, fill=mycolor2, fill opacity=0.15, forget plot]
table[row sep=crcr] {%
x	y\\
1.4568	0.00031411\\
1.4607	0.00032233\\
1.4646	0.00033075\\
1.4685	0.00033939\\
1.4723	0.00034824\\
1.4762	0.00035731\\
1.4801	0.0003666\\
1.484	0.00037612\\
1.4878	0.00038588\\
1.4917	0.00039587\\
1.4956	0.00040611\\
1.4995	0.00041661\\
1.5034	0.00042736\\
1.5072	0.00043837\\
1.5111	0.00044966\\
1.515	0.00046122\\
1.5189	0.00047306\\
1.5227	0.00048519\\
1.5266	0.00049762\\
1.5305	0.00051035\\
1.5344	0.00052339\\
1.5383	0.00053674\\
1.5421	0.00055042\\
1.546	0.00056443\\
1.5499	0.00057878\\
1.5538	0.00059348\\
1.5576	0.00060853\\
1.5615	0.00062395\\
1.5654	0.00063973\\
1.5693	0.0006559\\
1.5732	0.00067245\\
1.577	0.00068941\\
1.5809	0.00070677\\
1.5848	0.00072454\\
1.5887	0.00074274\\
1.5925	0.00076138\\
1.5964	0.00078046\\
1.6003	0.0008\\
1.6042	0.00082\\
1.6081	0.00084048\\
1.6119	0.00086145\\
1.6158	0.00088292\\
1.6197	0.00090489\\
1.6236	0.00092739\\
1.6274	0.00095042\\
1.6313	0.000974\\
1.6352	0.00099813\\
1.6391	0.0010228\\
1.643	0.0010481\\
1.6468	0.001074\\
1.6507	0.0011005\\
1.6546	0.0011276\\
1.6585	0.0011554\\
1.6623	0.0011838\\
1.6662	0.0012129\\
1.6701	0.0012426\\
1.674	0.0012731\\
1.6778	0.0013043\\
1.6817	0.0013362\\
1.6856	0.0013688\\
1.6895	0.0014022\\
1.6934	0.0014364\\
1.6972	0.0014714\\
1.7011	0.0015072\\
1.705	0.0015439\\
1.7089	0.0015814\\
1.7127	0.0016198\\
1.7166	0.001659\\
1.7205	0.0016992\\
1.7244	0.0017403\\
1.7283	0.0017824\\
1.7321	0.0018254\\
1.736	0.0018694\\
1.7399	0.0019145\\
1.7438	0.0019606\\
1.7476	0.0020078\\
1.7515	0.002056\\
1.7554	0.0021054\\
1.7593	0.0021559\\
1.7632	0.0022075\\
1.767	0.0022604\\
1.7709	0.0023145\\
1.7748	0.0023698\\
1.7787	0.0024264\\
1.7825	0.0024843\\
1.7864	0.0025435\\
1.7903	0.0026041\\
1.7942	0.0026661\\
1.7981	0.0027295\\
1.8019	0.0027943\\
1.8058	0.0028607\\
1.8097	0.0029285\\
1.8136	0.0029979\\
1.8174	0.0030689\\
1.8213	0.0031415\\
1.8252	0.0032158\\
1.8291	0.0032917\\
1.833	0.0033694\\
1.8368	0.0034489\\
1.8407	0.0035302\\
1.8446	0.0036133\\
1.8485	0.0036983\\
1.8523	0.0037852\\
1.8562	0.0038741\\
1.8601	0.0039651\\
1.864	0.0040581\\
1.8679	0.0041532\\
1.8717	0.0042505\\
1.8756	0.0043499\\
1.8795	0.0044516\\
1.8834	0.0045557\\
1.8872	0.004662\\
1.8911	0.0047708\\
1.895	0.004882\\
1.8989	0.0049958\\
1.9028	0.0051121\\
1.9066	0.005231\\
1.9105	0.0053526\\
1.9144	0.0054769\\
1.9183	0.0056041\\
1.9221	0.0057341\\
1.926	0.005867\\
1.9299	0.0060029\\
1.9338	0.0061419\\
1.9377	0.0062839\\
1.9415	0.0064292\\
1.9454	0.0065777\\
1.9493	0.0067295\\
1.9532	0.0068847\\
1.957	0.0070434\\
1.9609	0.0072057\\
1.9648	0.0073715\\
1.9687	0.0075411\\
1.9726	0.0077144\\
1.9764	0.0078916\\
1.9803	0.0080728\\
1.9842	0.0082579\\
1.9881	0.0084472\\
1.9919	0.0086407\\
1.9958	0.0088385\\
1.9997	0.0090407\\
2.0036	0.0092473\\
2.0075	0.0094585\\
2.0113	0.0096744\\
2.0152	0.0098951\\
2.0191	0.010121\\
2.023	0.010351\\
2.0268	0.010587\\
2.0307	0.010827\\
2.0346	0.011074\\
2.0385	0.011325\\
2.0424	0.011582\\
2.0462	0.011845\\
2.0501	0.012113\\
2.054	0.012387\\
2.0579	0.012668\\
2.0617	0.012954\\
2.0656	0.013247\\
2.0695	0.013546\\
2.0734	0.013851\\
2.0773	0.014163\\
2.0811	0.014482\\
2.085	0.014808\\
2.0889	0.015141\\
2.0928	0.015481\\
2.0966	0.015829\\
2.1005	0.016184\\
2.1044	0.016547\\
2.1083	0.016917\\
2.1121	0.017296\\
2.116	0.017682\\
2.1199	0.018077\\
2.1238	0.018481\\
2.1277	0.018893\\
2.1315	0.019314\\
2.1354	0.019744\\
2.1393	0.020183\\
2.1432	0.020631\\
2.147	0.021089\\
2.1509	0.021557\\
2.1548	0.022034\\
2.1587	0.022522\\
2.1626	0.02302\\
2.1664	0.023528\\
2.1703	0.024047\\
2.1742	0.024577\\
2.1781	0.025119\\
2.1819	0.025671\\
2.1858	0.026235\\
2.1897	0.026811\\
2.1936	0.027399\\
2.1975	0.027999\\
2.2013	0.028612\\
2.2052	0.029237\\
2.2091	0.029875\\
2.213	0.030526\\
2.2168	0.031191\\
2.2207	0.03187\\
2.2246	0.032562\\
2.2285	0.033269\\
2.2324	0.03399\\
2.2362	0.034725\\
2.2401	0.035476\\
2.244	0.036242\\
2.2479	0.037023\\
2.2517	0.03782\\
2.2556	0.038633\\
2.2595	0.039463\\
2.2634	0.040309\\
2.2673	0.041172\\
2.2711	0.042053\\
2.275	0.04295\\
2.2789	0.043866\\
2.2828	0.0448\\
2.2866	0.045752\\
2.2905	0.046723\\
2.2944	0.047713\\
2.2983	0.048723\\
2.3022	0.049752\\
2.306	0.050801\\
2.3099	0.05187\\
2.3138	0.052961\\
2.3177	0.054072\\
2.3215	0.055204\\
2.3254	0.056359\\
2.3293	0.057535\\
2.3332	0.058734\\
2.3371	0.059955\\
2.3409	0.0612\\
2.3448	0.062468\\
2.3487	0.06376\\
2.3526	0.065076\\
2.3564	0.066416\\
2.3603	0.067782\\
2.3642	0.069173\\
2.3681	0.070589\\
2.372	0.072031\\
2.3758	0.0735\\
2.3797	0.074996\\
2.3836	0.076518\\
2.3875	0.078069\\
2.3913	0.079647\\
2.3952	0.081253\\
2.3991	0.082888\\
2.403	0.084552\\
2.4069	0.086246\\
2.4107	0.087969\\
2.4146	0.089723\\
2.4185	0.091507\\
2.4224	0.093322\\
2.4262	0.095169\\
2.4301	0.097047\\
2.434	0.098958\\
2.4379	0.1009\\
2.4418	0.10288\\
2.4456	0.10489\\
2.4495	0.10693\\
2.4534	0.10901\\
2.4573	0.11112\\
2.4611	0.11326\\
2.465	0.11544\\
2.4689	0.11766\\
2.4728	0.11991\\
2.4767	0.1222\\
2.4805	0.12453\\
2.4844	0.12689\\
2.4883	0.12929\\
2.4922	0.13173\\
2.496	0.13421\\
2.4999	0.13672\\
2.5038	0.13927\\
2.5077	0.14187\\
2.5116	0.1445\\
2.5154	0.14717\\
2.5193	0.14988\\
2.5232	0.15263\\
2.5271	0.15543\\
2.5309	0.15826\\
2.5348	0.16114\\
2.5387	0.16405\\
2.5426	0.16701\\
2.5464	0.17001\\
2.5503	0.17306\\
2.5542	0.17614\\
2.5581	0.17927\\
2.562	0.18245\\
2.5658	0.18566\\
2.5697	0.18892\\
2.5736	0.19222\\
2.5775	0.19557\\
2.5813	0.19896\\
2.5852	0.2024\\
2.5891	0.20588\\
2.593	0.20941\\
2.5969	0.21298\\
2.6007	0.21659\\
2.6046	0.22026\\
2.6085	0.22396\\
2.6124	0.22771\\
2.6162	0.23151\\
2.6201	0.23535\\
2.624	0.23924\\
2.6279	0.24318\\
2.6318	0.24716\\
2.6356	0.25118\\
2.6395	0.25526\\
2.6434	0.25937\\
2.6473	0.26354\\
2.6511	0.26774\\
2.655	0.272\\
2.6589	0.2763\\
2.6628	0.28064\\
2.6667	0.28504\\
2.6705	0.28947\\
2.6744	0.29395\\
2.6783	0.29848\\
2.6822	0.30305\\
2.686	0.30766\\
2.6899	0.31232\\
2.6938	0.31703\\
2.6977	0.32177\\
2.7016	0.32657\\
2.7054	0.3314\\
2.7093	0.33628\\
2.7132	0.3412\\
2.7171	0.34616\\
2.7209	0.35116\\
2.7248	0.35621\\
2.7287	0.36129\\
2.7326	0.36642\\
2.7365	0.37159\\
2.7403	0.37679\\
2.7442	0.38204\\
2.7481	0.38732\\
2.752	0.39265\\
2.7558	0.39801\\
2.7597	0.4034\\
2.7636	0.40884\\
2.7675	0.41431\\
2.7714	0.41981\\
2.7752	0.42535\\
2.7791	0.43092\\
2.783	0.43653\\
2.7869	0.44216\\
2.7907	0.44783\\
2.7946	0.45353\\
2.7985	0.45926\\
2.8024	0.46501\\
2.8063	0.4708\\
2.8101	0.47661\\
2.814	0.48245\\
2.8179	0.48831\\
2.8218	0.4942\\
2.8256	0.50011\\
2.8295	0.50605\\
2.8334	0.512\\
2.8373	0.51798\\
2.8412	0.52397\\
2.845	0.52998\\
2.8489	0.53601\\
2.8528	0.54206\\
2.8567	0.54812\\
2.8605	0.55419\\
2.8644	0.56027\\
2.8683	0.56637\\
2.8722	0.57248\\
2.8761	0.57859\\
2.8799	0.58471\\
2.8838	0.59084\\
2.8877	0.59697\\
2.8916	0.60311\\
2.8954	0.60924\\
2.8993	0.61538\\
2.9032	0.62152\\
2.9071	0.62765\\
2.911	0.63378\\
2.9148	0.63991\\
2.9187	0.64603\\
2.9226	0.65214\\
2.9265	0.65825\\
2.9303	0.66434\\
2.9342	0.67042\\
2.9381	0.67649\\
2.942	0.68254\\
2.9459	0.68858\\
2.9497	0.6946\\
2.9536	0.7006\\
2.9575	0.70658\\
2.9614	0.71253\\
2.9652	0.71847\\
2.9691	0.72438\\
2.973	0.73027\\
2.9769	0.73613\\
2.9808	0.74196\\
2.9846	0.74777\\
2.9885	0.75354\\
2.9924	0.75929\\
2.9963	0.765\\
3.0001	0.77068\\
3.004	0.77632\\
3.0079	0.78193\\
3.0118	0.7875\\
3.0156	0.79304\\
3.0195	0.79853\\
3.0234	0.80399\\
3.0273	0.80939\\
3.0312	0.81476\\
3.035	0.82007\\
3.0389	0.82533\\
3.0428	0.83054\\
3.0467	0.83569\\
3.0505	0.84078\\
3.0544	0.8458\\
3.0583	0.85076\\
3.0622	0.85564\\
3.0661	0.86044\\
3.0699	0.86516\\
3.0738	0.86979\\
3.0777	0.87433\\
3.0816	0.87878\\
3.0854	0.88312\\
3.0893	0.88737\\
3.0932	0.89151\\
3.0971	0.89554\\
3.101	0.89945\\
3.1048	0.90325\\
3.1087	0.90694\\
3.1126	0.9105\\
3.1165	0.91394\\
3.1203	0.91726\\
3.1242	0.92046\\
3.1281	0.92354\\
3.132	0.9265\\
3.1359	0.92933\\
3.1397	0.93204\\
3.1436	0.93464\\
3.1475	0.93711\\
3.1514	0.93947\\
3.1552	0.94172\\
3.1591	0.94385\\
3.163	0.94587\\
3.1669	0.94778\\
3.1708	0.94959\\
3.1746	0.95129\\
3.1785	0.95289\\
3.1824	0.95438\\
3.1863	0.95577\\
3.1901	0.95707\\
3.194	0.95826\\
3.1979	0.95936\\
3.2018	0.96036\\
3.2057	0.96127\\
3.2095	0.96208\\
3.2134	0.96279\\
3.2173	0.96341\\
3.2212	0.96393\\
3.225	0.96436\\
3.2289	0.9647\\
3.2328	0.96494\\
3.2367	0.96508\\
3.2406	0.96514\\
3.2444	0.9651\\
3.2483	0.96497\\
3.2522	0.96474\\
3.2561	0.96442\\
3.2599	0.96402\\
3.2638	0.96352\\
3.2677	0.96293\\
3.2716	0.96225\\
3.2755	0.96149\\
3.2793	0.96064\\
3.2832	0.9597\\
3.2871	0.95868\\
3.291	0.95758\\
3.2948	0.95639\\
3.2987	0.95512\\
3.3026	0.95377\\
3.3065	0.95234\\
3.3104	0.95083\\
3.3142	0.94924\\
3.3181	0.94757\\
3.322	0.94582\\
3.3259	0.944\\
3.3297	0.9421\\
3.3336	0.94012\\
3.3375	0.93806\\
3.3414	0.93593\\
3.3453	0.93372\\
3.3491	0.93143\\
3.353	0.92906\\
3.3569	0.92662\\
3.3608	0.9241\\
3.3646	0.92151\\
3.3685	0.91883\\
3.3724	0.91608\\
3.3763	0.91326\\
3.3802	0.91035\\
3.384	0.90737\\
3.3879	0.90432\\
3.3918	0.90118\\
3.3957	0.89798\\
3.3995	0.8947\\
3.4034	0.89134\\
3.4073	0.88791\\
3.4112	0.88441\\
3.4151	0.88084\\
3.4189	0.87719\\
3.4228	0.87348\\
3.4267	0.8697\\
3.4306	0.86585\\
3.4344	0.86194\\
3.4383	0.85796\\
3.4422	0.85391\\
3.4461	0.84981\\
3.4499	0.84564\\
3.4538	0.84141\\
3.4577	0.83713\\
3.4616	0.83279\\
3.4655	0.8284\\
3.4693	0.82396\\
3.4732	0.81947\\
3.4771	0.81492\\
3.481	0.81034\\
3.4848	0.8057\\
3.4887	0.80103\\
3.4926	0.79632\\
3.4965	0.79156\\
3.5004	0.78677\\
3.5042	0.78195\\
3.5081	0.7771\\
3.512	0.77222\\
3.5159	0.7673\\
3.5197	0.76237\\
3.5236	0.75741\\
3.5275	0.75242\\
3.5314	0.74742\\
3.5353	0.74241\\
3.5391	0.73737\\
3.543	0.73232\\
3.5469	0.72726\\
3.5508	0.72219\\
3.5546	0.71712\\
3.5585	0.71203\\
3.5624	0.70694\\
3.5663	0.70185\\
3.5702	0.69676\\
3.574	0.69166\\
3.5779	0.68657\\
3.5818	0.68148\\
3.5857	0.67639\\
3.5895	0.6713\\
3.5934	0.66623\\
3.5973	0.66116\\
3.6012	0.65609\\
3.6051	0.65104\\
3.6089	0.64599\\
3.6128	0.64096\\
3.6167	0.63594\\
3.6206	0.63093\\
3.6244	0.62593\\
3.6283	0.62094\\
3.6322	0.61597\\
3.6361	0.61102\\
3.64	0.60608\\
3.6438	0.60115\\
3.6477	0.59624\\
3.6516	0.59135\\
3.6555	0.58648\\
3.6593	0.58162\\
3.6632	0.57678\\
3.6671	0.57196\\
3.671	0.56716\\
3.6749	0.56238\\
3.6787	0.55762\\
3.6826	0.55288\\
3.6865	0.54816\\
3.6904	0.54346\\
3.6942	0.53879\\
3.6981	0.53413\\
3.702	0.5295\\
3.7059	0.52489\\
3.7098	0.52031\\
3.7136	0.51575\\
3.7175	0.51121\\
3.7214	0.5067\\
3.7253	0.50221\\
3.7291	0.49774\\
3.733	0.4933\\
3.7369	0.48889\\
3.7408	0.4845\\
3.7447	0.48014\\
3.7485	0.4758\\
3.7524	0.47149\\
3.7563	0.46721\\
3.7602	0.46296\\
3.764	0.45873\\
3.7679	0.45453\\
3.7718	0.45035\\
3.7757	0.44621\\
3.7796	0.44209\\
3.7834	0.438\\
3.7873	0.43394\\
3.7912	0.42991\\
3.7951	0.4259\\
3.7989	0.42193\\
3.8028	0.41798\\
3.8067	0.41407\\
3.8106	0.41018\\
3.8145	0.40632\\
3.8183	0.40249\\
3.8222	0.39869\\
3.8261	0.39492\\
3.83	0.39117\\
3.8338	0.38746\\
3.8377	0.38378\\
3.8416	0.38012\\
3.8455	0.3765\\
3.8494	0.3729\\
3.8532	0.36934\\
3.8571	0.3658\\
3.861	0.36229\\
3.8649	0.35881\\
3.8687	0.35536\\
3.8726	0.35194\\
3.8765	0.34855\\
3.8804	0.34519\\
3.8842	0.34186\\
3.8881	0.33855\\
3.892	0.33527\\
3.8959	0.33203\\
3.8998	0.32881\\
3.9036	0.32562\\
3.9075	0.32245\\
3.9114	0.31932\\
3.9153	0.31621\\
3.9191	0.31313\\
3.923	0.31008\\
3.9269	0.30705\\
3.9308	0.30406\\
3.9347	0.30108\\
3.9385	0.29814\\
3.9424	0.29522\\
3.9463	0.29233\\
3.9502	0.28947\\
3.954	0.28663\\
3.9579	0.28382\\
3.9618	0.28103\\
3.9657	0.27827\\
3.9696	0.27553\\
3.9734	0.27282\\
3.9773	0.27013\\
3.9812	0.26747\\
3.9851	0.26483\\
3.9889	0.26222\\
3.9928	0.25963\\
3.9967	0.25706\\
4.0006	0.25452\\
4.0045	0.252\\
4.0083	0.24951\\
4.0122	0.24703\\
4.0161	0.24458\\
4.02	0.24215\\
4.0238	0.23975\\
4.0277	0.23736\\
4.0316	0.235\\
4.0355	0.23266\\
4.0394	0.23034\\
4.0432	0.22805\\
4.0471	0.22577\\
4.051	0.22351\\
4.0549	0.22128\\
4.0587	0.21906\\
4.0626	0.21687\\
4.0665	0.21469\\
4.0704	0.21254\\
4.0743	0.2104\\
4.0781	0.20829\\
4.082	0.20619\\
4.0859	0.20411\\
4.0898	0.20205\\
4.0936	0.20001\\
4.0975	0.19799\\
4.1014	0.19599\\
4.1053	0.194\\
4.1092	0.19203\\
4.113	0.19008\\
4.1169	0.18815\\
4.1208	0.18624\\
4.1247	0.18434\\
4.1285	0.18246\\
4.1324	0.18059\\
4.1363	0.17874\\
4.1402	0.17691\\
4.1441	0.1751\\
4.1479	0.1733\\
4.1518	0.17152\\
4.1557	0.16975\\
4.1596	0.168\\
4.1634	0.16626\\
4.1673	0.16454\\
4.1712	0.16284\\
4.1751	0.16115\\
4.179	0.15947\\
4.1828	0.15781\\
4.1867	0.15617\\
4.1906	0.15453\\
4.1945	0.15292\\
4.1983	0.15132\\
4.2022	0.14973\\
4.2061	0.14815\\
4.21	0.14659\\
4.2139	0.14505\\
4.2177	0.14351\\
4.2216	0.14199\\
4.2255	0.14049\\
4.2294	0.13899\\
4.2332	0.13752\\
4.2371	0.13605\\
4.241	0.13459\\
4.2449	0.13315\\
4.2488	0.13173\\
4.2526	0.13031\\
4.2565	0.12891\\
4.2604	0.12751\\
4.2643	0.12614\\
4.2681	0.12477\\
4.272	0.12341\\
4.2759	0.12207\\
4.2798	0.12074\\
4.2837	0.11942\\
4.2875	0.11811\\
4.2914	0.11682\\
4.2953	0.11553\\
4.2992	0.11426\\
4.303	0.113\\
4.3069	0.11175\\
4.3108	0.11051\\
4.3147	0.10928\\
4.3185	0.10806\\
4.3224	0.10686\\
4.3263	0.10566\\
4.3302	0.10447\\
4.3341	0.1033\\
4.3379	0.10213\\
4.3418	0.10098\\
4.3457	0.099836\\
4.3496	0.098702\\
4.3534	0.097579\\
4.3573	0.096466\\
4.3612	0.095362\\
4.3651	0.094269\\
4.369	0.093186\\
4.3728	0.092112\\
4.3767	0.091048\\
4.3806	0.089994\\
4.3845	0.088949\\
4.3883	0.087914\\
4.3922	0.086889\\
4.3961	0.085872\\
4.4	0.084865\\
4.4039	0.083868\\
4.4077	0.082879\\
4.4116	0.081899\\
4.4155	0.080929\\
4.4194	0.079968\\
4.4232	0.079015\\
4.4271	0.078071\\
4.431	0.077136\\
4.4349	0.07621\\
4.4388	0.075293\\
4.4426	0.074384\\
4.4465	0.073484\\
4.4504	0.072592\\
4.4543	0.071708\\
4.4581	0.070833\\
4.462	0.069966\\
4.4659	0.069108\\
4.4698	0.068257\\
4.4737	0.067415\\
4.4775	0.066581\\
4.4814	0.065755\\
4.4853	0.064937\\
4.4892	0.064126\\
4.493	0.063324\\
4.4969	0.062529\\
4.5008	0.061742\\
4.5047	0.060963\\
4.5086	0.060192\\
4.5124	0.059428\\
4.5163	0.058671\\
4.5202	0.057922\\
4.5241	0.05718\\
4.5279	0.056446\\
4.5318	0.055719\\
4.5357	0.054999\\
4.5396	0.054287\\
4.5435	0.053582\\
4.5473	0.052883\\
4.5512	0.052192\\
4.5551	0.051508\\
4.559	0.050831\\
4.5628	0.05016\\
4.5667	0.049497\\
4.5706	0.04884\\
4.5745	0.04819\\
4.5784	0.047547\\
4.5822	0.046911\\
4.5861	0.046281\\
4.59	0.045658\\
4.5939	0.045041\\
4.5977	0.04443\\
4.6016	0.043827\\
4.6055	0.043229\\
4.6094	0.042638\\
4.6133	0.042053\\
4.6171	0.041474\\
4.621	0.040902\\
4.6249	0.040336\\
4.6288	0.039776\\
4.6326	0.039222\\
4.6365	0.038674\\
4.6404	0.038132\\
4.6443	0.037596\\
4.6482	0.037065\\
4.652	0.036541\\
4.6559	0.036023\\
4.6598	0.03551\\
4.6637	0.035003\\
4.6675	0.034502\\
4.6714	0.034006\\
4.6753	0.033516\\
4.6792	0.033032\\
4.6831	0.032553\\
4.6869	0.032079\\
4.6908	0.031611\\
4.6947	0.031149\\
4.6986	0.030691\\
4.7024	0.03024\\
4.7063	0.029793\\
4.7102	0.029352\\
4.7141	0.028915\\
4.718	0.028484\\
4.7218	0.028058\\
4.7257	0.027638\\
4.7296	0.027222\\
4.7335	0.026811\\
4.7373	0.026405\\
4.7412	0.026004\\
4.7451	0.025608\\
4.749	0.025217\\
4.7529	0.024831\\
4.7567	0.024449\\
4.7606	0.024073\\
4.7645	0.0237\\
4.7684	0.023333\\
4.7722	0.02297\\
4.7761	0.022612\\
4.78	0.022258\\
4.7839	0.021909\\
4.7877	0.021564\\
4.7916	0.021223\\
4.7955	0.020887\\
4.7994	0.020556\\
4.8033	0.020228\\
4.8071	0.019905\\
4.811	0.019586\\
4.8149	0.019271\\
4.8188	0.018961\\
4.8226	0.018654\\
4.8265	0.018352\\
4.8304	0.018054\\
4.8343	0.017759\\
4.8382	0.017469\\
4.842	0.017182\\
4.8459	0.0169\\
4.8498	0.016621\\
4.8537	0.016346\\
4.8575	0.016075\\
4.8614	0.015808\\
4.8653	0.015544\\
4.8692	0.015284\\
4.8731	0.015028\\
4.8769	0.014775\\
4.8808	0.014526\\
4.8847	0.01428\\
4.8886	0.014038\\
4.8924	0.013799\\
4.8963	0.013564\\
4.9002	0.013332\\
4.9041	0.013103\\
4.908	0.012878\\
4.9118	0.012656\\
4.9157	0.012437\\
4.9196	0.012221\\
4.9235	0.012009\\
4.9273	0.011799\\
4.9312	0.011593\\
4.9351	0.01139\\
4.939	0.01119\\
4.9429	0.010993\\
4.9467	0.010799\\
4.9506	0.010607\\
4.9545	0.010419\\
4.9584	0.010233\\
4.9622	0.010051\\
4.9661	0.0098709\\
4.97	0.0096938\\
4.9739	0.0095194\\
4.9778	0.0093477\\
4.9816	0.0091787\\
4.9855	0.0090123\\
4.9894	0.0088485\\
4.9933	0.0086873\\
4.9971	0.0085286\\
5.001	0.0083724\\
5.0049	0.0082187\\
5.0088	0.0080674\\
5.0127	0.0079186\\
5.0165	0.0077721\\
5.0204	0.007628\\
5.0243	0.0074862\\
5.0282	0.0073467\\
5.032	0.0072095\\
5.0359	0.0070745\\
5.0398	0.0069417\\
5.0437	0.0068111\\
5.0476	0.0066826\\
5.0514	0.0065563\\
5.0553	0.006432\\
5.0592	0.0063098\\
5.0631	0.0061897\\
5.0669	0.0060715\\
5.0708	0.0059554\\
5.0747	0.0058412\\
5.0786	0.0057289\\
5.0825	0.0056185\\
5.0863	0.00551\\
5.0902	0.0054033\\
5.0941	0.0052985\\
5.098	0.0051955\\
5.1018	0.0050942\\
5.1057	0.0049947\\
5.1096	0.0048969\\
5.1135	0.0048008\\
5.1174	0.0047064\\
5.1212	0.0046136\\
5.1251	0.0045224\\
5.129	0.0044329\\
5.1329	0.0043449\\
5.1367	0.0042585\\
5.1406	0.0041736\\
5.1445	0.0040902\\
5.1484	0.0040083\\
5.1523	0.0039279\\
5.1561	0.0038489\\
5.16	0.0037713\\
5.1639	0.0036952\\
5.1678	0.0036204\\
5.1716	0.003547\\
5.1755	0.0034749\\
5.1794	0.0034041\\
5.1833	0.0033346\\
5.1872	0.0032664\\
5.191	0.0031994\\
5.1949	0.0031337\\
5.1988	0.0030692\\
5.2027	0.0030059\\
5.2065	0.0029437\\
5.2104	0.0028828\\
5.2143	0.0028229\\
5.2182	0.0027642\\
5.222	0.0027066\\
5.2259	0.0026501\\
5.2298	0.0025946\\
5.2337	0.0025402\\
5.2376	0.0024868\\
5.2414	0.0024345\\
5.2453	0.0023831\\
5.2492	0.0023328\\
5.2531	0.0022834\\
5.2569	0.0022349\\
5.2608	0.0021874\\
5.2647	0.0021408\\
5.2686	0.0020951\\
5.2725	0.0020503\\
5.2763	0.0020063\\
5.2802	0.0019633\\
5.2841	0.001921\\
5.288	0.0018796\\
5.2918	0.001839\\
5.2957	0.0017993\\
5.2996	0.0017602\\
5.3035	0.001722\\
5.3074	0.0016846\\
5.3112	0.0016478\\
5.3151	0.0016118\\
5.319	0.0015766\\
5.3229	0.001542\\
5.3267	0.0015081\\
5.3306	0.001475\\
5.3306	0.0012365\\
5.3267	0.0012658\\
5.3229	0.0012958\\
5.319	0.0013264\\
5.3151	0.0013577\\
5.3112	0.0013896\\
5.3074	0.0014222\\
5.3035	0.0014556\\
5.2996	0.0014896\\
5.2957	0.0015244\\
5.2918	0.0015599\\
5.288	0.0015962\\
5.2841	0.0016332\\
5.2802	0.001671\\
5.2763	0.0017096\\
5.2725	0.001749\\
5.2686	0.0017893\\
5.2647	0.0018304\\
5.2608	0.0018723\\
5.2569	0.0019151\\
5.2531	0.0019588\\
5.2492	0.0020034\\
5.2453	0.0020489\\
5.2414	0.0020953\\
5.2376	0.0021427\\
5.2337	0.0021911\\
5.2298	0.0022404\\
5.2259	0.0022908\\
5.222	0.0023421\\
5.2182	0.0023945\\
5.2143	0.002448\\
5.2104	0.0025025\\
5.2065	0.0025581\\
5.2027	0.0026149\\
5.1988	0.0026727\\
5.1949	0.0027317\\
5.191	0.0027919\\
5.1872	0.0028532\\
5.1833	0.0029158\\
5.1794	0.0029796\\
5.1755	0.0030446\\
5.1716	0.0031109\\
5.1678	0.0031785\\
5.1639	0.0032474\\
5.16	0.0033176\\
5.1561	0.0033892\\
5.1523	0.0034621\\
5.1484	0.0035364\\
5.1445	0.0036122\\
5.1406	0.0036894\\
5.1367	0.003768\\
5.1329	0.0038482\\
5.129	0.0039298\\
5.1251	0.004013\\
5.1212	0.0040977\\
5.1174	0.004184\\
5.1135	0.0042719\\
5.1096	0.0043614\\
5.1057	0.0044526\\
5.1018	0.0045455\\
5.098	0.00464\\
5.0941	0.0047363\\
5.0902	0.0048344\\
5.0863	0.0049342\\
5.0825	0.0050358\\
5.0786	0.0051392\\
5.0747	0.0052445\\
5.0708	0.0053517\\
5.0669	0.0054608\\
5.0631	0.0055718\\
5.0592	0.0056848\\
5.0553	0.0057998\\
5.0514	0.0059168\\
5.0476	0.0060359\\
5.0437	0.006157\\
5.0398	0.0062802\\
5.0359	0.0064056\\
5.032	0.0065331\\
5.0282	0.0066628\\
5.0243	0.0067948\\
5.0204	0.006929\\
5.0165	0.0070655\\
5.0127	0.0072042\\
5.0088	0.0073454\\
5.0049	0.0074889\\
5.001	0.0076348\\
4.9971	0.0077831\\
4.9933	0.0079339\\
4.9894	0.0080872\\
4.9855	0.008243\\
4.9816	0.0084014\\
4.9778	0.0085624\\
4.9739	0.008726\\
4.97	0.0088923\\
4.9661	0.0090612\\
4.9622	0.0092329\\
4.9584	0.0094073\\
4.9545	0.0095845\\
4.9506	0.0097645\\
4.9467	0.0099474\\
4.9429	0.010133\\
4.939	0.010322\\
4.9351	0.010514\\
4.9312	0.010708\\
4.9273	0.010906\\
4.9235	0.011107\\
4.9196	0.01131\\
4.9157	0.011517\\
4.9118	0.011727\\
4.908	0.011941\\
4.9041	0.012157\\
4.9002	0.012377\\
4.8963	0.0126\\
4.8924	0.012827\\
4.8886	0.013056\\
4.8847	0.01329\\
4.8808	0.013526\\
4.8769	0.013766\\
4.8731	0.01401\\
4.8692	0.014257\\
4.8653	0.014508\\
4.8614	0.014763\\
4.8575	0.015021\\
4.8537	0.015283\\
4.8498	0.015548\\
4.8459	0.015818\\
4.842	0.016091\\
4.8382	0.016368\\
4.8343	0.016649\\
4.8304	0.016933\\
4.8265	0.017222\\
4.8226	0.017515\\
4.8188	0.017812\\
4.8149	0.018113\\
4.811	0.018418\\
4.8071	0.018727\\
4.8033	0.01904\\
4.7994	0.019358\\
4.7955	0.019679\\
4.7916	0.020006\\
4.7877	0.020336\\
4.7839	0.020671\\
4.78	0.02101\\
4.7761	0.021354\\
4.7722	0.021702\\
4.7684	0.022055\\
4.7645	0.022412\\
4.7606	0.022774\\
4.7567	0.02314\\
4.7529	0.023512\\
4.749	0.023888\\
4.7451	0.024268\\
4.7412	0.024654\\
4.7373	0.025044\\
4.7335	0.025439\\
4.7296	0.025839\\
4.7257	0.026244\\
4.7218	0.026654\\
4.718	0.027069\\
4.7141	0.02749\\
4.7102	0.027915\\
4.7063	0.028345\\
4.7024	0.028781\\
4.6986	0.029221\\
4.6947	0.029667\\
4.6908	0.030119\\
4.6869	0.030575\\
4.6831	0.031037\\
4.6792	0.031504\\
4.6753	0.031977\\
4.6714	0.032456\\
4.6675	0.032939\\
4.6637	0.033429\\
4.6598	0.033924\\
4.6559	0.034424\\
4.652	0.03493\\
4.6482	0.035442\\
4.6443	0.03596\\
4.6404	0.036483\\
4.6365	0.037013\\
4.6326	0.037548\\
4.6288	0.038089\\
4.6249	0.038636\\
4.621	0.039188\\
4.6171	0.039747\\
4.6133	0.040312\\
4.6094	0.040883\\
4.6055	0.04146\\
4.6016	0.042043\\
4.5977	0.042633\\
4.5939	0.043229\\
4.59	0.04383\\
4.5861	0.044439\\
4.5822	0.045053\\
4.5784	0.045674\\
4.5745	0.046301\\
4.5706	0.046935\\
4.5667	0.047575\\
4.5628	0.048222\\
4.559	0.048875\\
4.5551	0.049535\\
4.5512	0.050202\\
4.5473	0.050875\\
4.5435	0.051555\\
4.5396	0.052242\\
4.5357	0.052935\\
4.5318	0.053636\\
4.5279	0.054343\\
4.5241	0.055057\\
4.5202	0.055778\\
4.5163	0.056506\\
4.5124	0.057241\\
4.5086	0.057984\\
4.5047	0.058733\\
4.5008	0.059489\\
4.4969	0.060253\\
4.493	0.061024\\
4.4892	0.061802\\
4.4853	0.062588\\
4.4814	0.063381\\
4.4775	0.064182\\
4.4737	0.06499\\
4.4698	0.065805\\
4.4659	0.066628\\
4.462	0.067459\\
4.4581	0.068297\\
4.4543	0.069144\\
4.4504	0.069998\\
4.4465	0.070859\\
4.4426	0.071729\\
4.4388	0.072607\\
4.4349	0.073493\\
4.431	0.074386\\
4.4271	0.075288\\
4.4232	0.076198\\
4.4194	0.077117\\
4.4155	0.078043\\
4.4116	0.078978\\
4.4077	0.079922\\
4.4039	0.080874\\
4.4	0.081835\\
4.3961	0.082804\\
4.3922	0.083782\\
4.3883	0.084768\\
4.3845	0.085764\\
4.3806	0.086769\\
4.3767	0.087782\\
4.3728	0.088805\\
4.369	0.089837\\
4.3651	0.090878\\
4.3612	0.091928\\
4.3573	0.092988\\
4.3534	0.094057\\
4.3496	0.095136\\
4.3457	0.096224\\
4.3418	0.097323\\
4.3379	0.098431\\
4.3341	0.099549\\
4.3302	0.10068\\
4.3263	0.10181\\
4.3224	0.10296\\
4.3185	0.10412\\
4.3147	0.10529\\
4.3108	0.10647\\
4.3069	0.10766\\
4.303	0.10886\\
4.2992	0.11007\\
4.2953	0.1113\\
4.2914	0.11253\\
4.2875	0.11378\\
4.2837	0.11503\\
4.2798	0.1163\\
4.2759	0.11758\\
4.272	0.11887\\
4.2681	0.12017\\
4.2643	0.12149\\
4.2604	0.12281\\
4.2565	0.12415\\
4.2526	0.1255\\
4.2488	0.12687\\
4.2449	0.12824\\
4.241	0.12963\\
4.2371	0.13103\\
4.2332	0.13245\\
4.2294	0.13387\\
4.2255	0.13531\\
4.2216	0.13677\\
4.2177	0.13823\\
4.2139	0.13971\\
4.21	0.14121\\
4.2061	0.14272\\
4.2022	0.14424\\
4.1983	0.14578\\
4.1945	0.14733\\
4.1906	0.14889\\
4.1867	0.15047\\
4.1828	0.15207\\
4.179	0.15368\\
4.1751	0.1553\\
4.1712	0.15694\\
4.1673	0.1586\\
4.1634	0.16027\\
4.1596	0.16196\\
4.1557	0.16366\\
4.1518	0.16538\\
4.1479	0.16712\\
4.1441	0.16887\\
4.1402	0.17064\\
4.1363	0.17242\\
4.1324	0.17423\\
4.1285	0.17605\\
4.1247	0.17788\\
4.1208	0.17974\\
4.1169	0.18161\\
4.113	0.1835\\
4.1092	0.18541\\
4.1053	0.18734\\
4.1014	0.18928\\
4.0975	0.19124\\
4.0936	0.19323\\
4.0898	0.19523\\
4.0859	0.19725\\
4.082	0.19929\\
4.0781	0.20134\\
4.0743	0.20342\\
4.0704	0.20552\\
4.0665	0.20764\\
4.0626	0.20978\\
4.0587	0.21194\\
4.0549	0.21411\\
4.051	0.21631\\
4.0471	0.21853\\
4.0432	0.22078\\
4.0394	0.22304\\
4.0355	0.22532\\
4.0316	0.22763\\
4.0277	0.22996\\
4.0238	0.23231\\
4.02	0.23468\\
4.0161	0.23707\\
4.0122	0.23949\\
4.0083	0.24192\\
4.0045	0.24439\\
4.0006	0.24687\\
3.9967	0.24938\\
3.9928	0.25191\\
3.9889	0.25446\\
3.9851	0.25704\\
3.9812	0.25964\\
3.9773	0.26226\\
3.9734	0.26491\\
3.9696	0.26758\\
3.9657	0.27028\\
3.9618	0.273\\
3.9579	0.27574\\
3.954	0.27851\\
3.9502	0.28131\\
3.9463	0.28413\\
3.9424	0.28697\\
3.9385	0.28984\\
3.9347	0.29273\\
3.9308	0.29565\\
3.9269	0.2986\\
3.923	0.30157\\
3.9191	0.30457\\
3.9153	0.30759\\
3.9114	0.31064\\
3.9075	0.31371\\
3.9036	0.31681\\
3.8998	0.31994\\
3.8959	0.3231\\
3.892	0.32628\\
3.8881	0.32948\\
3.8842	0.33272\\
3.8804	0.33598\\
3.8765	0.33926\\
3.8726	0.34258\\
3.8687	0.34592\\
3.8649	0.34929\\
3.861	0.35268\\
3.8571	0.35611\\
3.8532	0.35956\\
3.8494	0.36303\\
3.8455	0.36654\\
3.8416	0.37007\\
3.8377	0.37363\\
3.8338	0.37722\\
3.83	0.38083\\
3.8261	0.38448\\
3.8222	0.38815\\
3.8183	0.39185\\
3.8145	0.39557\\
3.8106	0.39933\\
3.8067	0.40311\\
3.8028	0.40692\\
3.7989	0.41076\\
3.7951	0.41463\\
3.7912	0.41852\\
3.7873	0.42244\\
3.7834	0.42639\\
3.7796	0.43037\\
3.7757	0.43437\\
3.7718	0.43841\\
3.7679	0.44247\\
3.764	0.44656\\
3.7602	0.45067\\
3.7563	0.45482\\
3.7524	0.45899\\
3.7485	0.46318\\
3.7447	0.46741\\
3.7408	0.47166\\
3.7369	0.47594\\
3.733	0.48024\\
3.7291	0.48458\\
3.7253	0.48893\\
3.7214	0.49332\\
3.7175	0.49772\\
3.7136	0.50216\\
3.7098	0.50662\\
3.7059	0.5111\\
3.702	0.51561\\
3.6981	0.52014\\
3.6942	0.5247\\
3.6904	0.52928\\
3.6865	0.53388\\
3.6826	0.5385\\
3.6787	0.54315\\
3.6749	0.54781\\
3.671	0.5525\\
3.6671	0.55721\\
3.6632	0.56193\\
3.6593	0.56668\\
3.6555	0.57144\\
3.6516	0.57622\\
3.6477	0.58101\\
3.6438	0.58582\\
3.64	0.59064\\
3.6361	0.59548\\
3.6322	0.60033\\
3.6283	0.60519\\
3.6244	0.61007\\
3.6206	0.61495\\
3.6167	0.61984\\
3.6128	0.62473\\
3.6089	0.62963\\
3.6051	0.63454\\
3.6012	0.63945\\
3.5973	0.64437\\
3.5934	0.64928\\
3.5895	0.6542\\
3.5857	0.65911\\
3.5818	0.66402\\
3.5779	0.66893\\
3.574	0.67384\\
3.5702	0.67873\\
3.5663	0.68363\\
3.5624	0.68851\\
3.5585	0.69338\\
3.5546	0.69825\\
3.5508	0.7031\\
3.5469	0.70794\\
3.543	0.71277\\
3.5391	0.71758\\
3.5353	0.72238\\
3.5314	0.72716\\
3.5275	0.73192\\
3.5236	0.73666\\
3.5197	0.74139\\
3.5159	0.74609\\
3.512	0.75077\\
3.5081	0.75543\\
3.5042	0.76006\\
3.5004	0.76467\\
3.4965	0.76926\\
3.4926	0.77381\\
3.4887	0.77834\\
3.4848	0.78284\\
3.481	0.7873\\
3.4771	0.79174\\
3.4732	0.79614\\
3.4693	0.80051\\
3.4655	0.80485\\
3.4616	0.80914\\
3.4577	0.8134\\
3.4538	0.81762\\
3.4499	0.8218\\
3.4461	0.82594\\
3.4422	0.83004\\
3.4383	0.83409\\
3.4344	0.8381\\
3.4306	0.84206\\
3.4267	0.84597\\
3.4228	0.84983\\
3.4189	0.85363\\
3.4151	0.85739\\
3.4112	0.86109\\
3.4073	0.86473\\
3.4034	0.86831\\
3.3995	0.87184\\
3.3957	0.8753\\
3.3918	0.8787\\
3.3879	0.88203\\
3.384	0.8853\\
3.3802	0.88849\\
3.3763	0.89162\\
3.3724	0.89467\\
3.3685	0.89765\\
3.3646	0.90055\\
3.3608	0.90337\\
3.3569	0.90612\\
3.353	0.90878\\
3.3491	0.91135\\
3.3453	0.91384\\
3.3414	0.91625\\
3.3375	0.91856\\
3.3336	0.92079\\
3.3297	0.92292\\
3.3259	0.92496\\
3.322	0.9269\\
3.3181	0.92875\\
3.3142	0.9305\\
3.3104	0.93216\\
3.3065	0.93372\\
3.3026	0.93517\\
3.2987	0.93654\\
3.2948	0.9378\\
3.291	0.93896\\
3.2871	0.94003\\
3.2832	0.94099\\
3.2793	0.94186\\
3.2755	0.94263\\
3.2716	0.9433\\
3.2677	0.94388\\
3.2638	0.94436\\
3.2599	0.94474\\
3.2561	0.94502\\
3.2522	0.94521\\
3.2483	0.94531\\
3.2444	0.9453\\
3.2406	0.94521\\
3.2367	0.94501\\
3.2328	0.94472\\
3.2289	0.94433\\
3.225	0.94385\\
3.2212	0.94327\\
3.2173	0.94259\\
3.2134	0.94181\\
3.2095	0.94092\\
3.2057	0.93994\\
3.2018	0.93886\\
3.1979	0.93767\\
3.194	0.93639\\
3.1901	0.93499\\
3.1863	0.9335\\
3.1824	0.93189\\
3.1785	0.93018\\
3.1746	0.92837\\
3.1708	0.92645\\
3.1669	0.92443\\
3.163	0.9223\\
3.1591	0.92007\\
3.1552	0.91774\\
3.1514	0.9153\\
3.1475	0.91276\\
3.1436	0.91013\\
3.1397	0.90739\\
3.1359	0.90456\\
3.132	0.90163\\
3.1281	0.8986\\
3.1242	0.89548\\
3.1203	0.89227\\
3.1165	0.88897\\
3.1126	0.88558\\
3.1087	0.8821\\
3.1048	0.87853\\
3.101	0.87488\\
3.0971	0.87114\\
3.0932	0.86731\\
3.0893	0.8634\\
3.0854	0.8594\\
3.0816	0.85533\\
3.0777	0.85117\\
3.0738	0.84692\\
3.0699	0.8426\\
3.0661	0.8382\\
3.0622	0.83372\\
3.0583	0.82916\\
3.0544	0.82452\\
3.0505	0.8198\\
3.0467	0.81501\\
3.0428	0.81015\\
3.0389	0.80521\\
3.035	0.80021\\
3.0312	0.79514\\
3.0273	0.79\\
3.0234	0.7848\\
3.0195	0.77953\\
3.0156	0.77421\\
3.0118	0.76883\\
3.0079	0.7634\\
3.004	0.75793\\
3.0001	0.7524\\
2.9963	0.74683\\
2.9924	0.74122\\
2.9885	0.73556\\
2.9846	0.72987\\
2.9808	0.72415\\
2.9769	0.7184\\
2.973	0.71261\\
2.9691	0.7068\\
2.9652	0.70096\\
2.9614	0.6951\\
2.9575	0.68921\\
2.9536	0.68331\\
2.9497	0.67738\\
2.9459	0.67144\\
2.942	0.66548\\
2.9381	0.6595\\
2.9342	0.65351\\
2.9303	0.64751\\
2.9265	0.6415\\
2.9226	0.63548\\
2.9187	0.62945\\
2.9148	0.62342\\
2.911	0.61738\\
2.9071	0.61133\\
2.9032	0.60528\\
2.8993	0.59923\\
2.8954	0.59318\\
2.8916	0.58713\\
2.8877	0.58108\\
2.8838	0.57503\\
2.8799	0.56899\\
2.8761	0.56295\\
2.8722	0.55692\\
2.8683	0.5509\\
2.8644	0.54488\\
2.8605	0.53888\\
2.8567	0.53289\\
2.8528	0.52691\\
2.8489	0.52095\\
2.845	0.515\\
2.8412	0.50907\\
2.8373	0.50316\\
2.8334	0.49726\\
2.8295	0.49139\\
2.8256	0.48553\\
2.8218	0.4797\\
2.8179	0.47389\\
2.814	0.4681\\
2.8101	0.46234\\
2.8063	0.45661\\
2.8024	0.4509\\
2.7985	0.44522\\
2.7946	0.43957\\
2.7907	0.43395\\
2.7869	0.42836\\
2.783	0.42281\\
2.7791	0.41728\\
2.7752	0.41179\\
2.7714	0.40633\\
2.7675	0.4009\\
2.7636	0.39551\\
2.7597	0.39016\\
2.7558	0.38484\\
2.752	0.37956\\
2.7481	0.37432\\
2.7442	0.36912\\
2.7403	0.36395\\
2.7365	0.35883\\
2.7326	0.35374\\
2.7287	0.3487\\
2.7248	0.3437\\
2.7209	0.33874\\
2.7171	0.33382\\
2.7132	0.32894\\
2.7093	0.3241\\
2.7054	0.31931\\
2.7016	0.31456\\
2.6977	0.30986\\
2.6938	0.3052\\
2.6899	0.30058\\
2.686	0.29601\\
2.6822	0.29148\\
2.6783	0.287\\
2.6744	0.28256\\
2.6705	0.27817\\
2.6667	0.27382\\
2.6628	0.26952\\
2.6589	0.26527\\
2.655	0.26106\\
2.6511	0.25689\\
2.6473	0.25277\\
2.6434	0.2487\\
2.6395	0.24468\\
2.6356	0.24069\\
2.6318	0.23676\\
2.6279	0.23287\\
2.624	0.22903\\
2.6201	0.22523\\
2.6162	0.22148\\
2.6124	0.21777\\
2.6085	0.21411\\
2.6046	0.21049\\
2.6007	0.20692\\
2.5969	0.20339\\
2.593	0.19991\\
2.5891	0.19648\\
2.5852	0.19308\\
2.5813	0.18974\\
2.5775	0.18643\\
2.5736	0.18317\\
2.5697	0.17996\\
2.5658	0.17678\\
2.562	0.17365\\
2.5581	0.17057\\
2.5542	0.16752\\
2.5503	0.16452\\
2.5464	0.16156\\
2.5426	0.15865\\
2.5387	0.15577\\
2.5348	0.15294\\
2.5309	0.15014\\
2.5271	0.14739\\
2.5232	0.14468\\
2.5193	0.142\\
2.5154	0.13937\\
2.5116	0.13678\\
2.5077	0.13422\\
2.5038	0.13171\\
2.4999	0.12923\\
2.496	0.12679\\
2.4922	0.12439\\
2.4883	0.12203\\
2.4844	0.1197\\
2.4805	0.11741\\
2.4767	0.11515\\
2.4728	0.11293\\
2.4689	0.11075\\
2.465	0.1086\\
2.4611	0.10649\\
2.4573	0.10441\\
2.4534	0.10237\\
2.4495	0.10036\\
2.4456	0.098378\\
2.4418	0.096433\\
2.4379	0.094521\\
2.434	0.09264\\
2.4301	0.090792\\
2.4262	0.088975\\
2.4224	0.087189\\
2.4185	0.085433\\
2.4146	0.083708\\
2.4107	0.082013\\
2.4069	0.080347\\
2.403	0.078711\\
2.3991	0.077103\\
2.3952	0.075523\\
2.3913	0.073972\\
2.3875	0.072448\\
2.3836	0.070952\\
2.3797	0.069482\\
2.3758	0.068039\\
2.372	0.066623\\
2.3681	0.065232\\
2.3642	0.063866\\
2.3603	0.062526\\
2.3564	0.06121\\
2.3526	0.059919\\
2.3487	0.058652\\
2.3448	0.057408\\
2.3409	0.056188\\
2.3371	0.054991\\
2.3332	0.053816\\
2.3293	0.052664\\
2.3254	0.051534\\
2.3215	0.050426\\
2.3177	0.049338\\
2.3138	0.048272\\
2.3099	0.047227\\
2.306	0.046201\\
2.3022	0.045196\\
2.2983	0.044211\\
2.2944	0.043245\\
2.2905	0.042297\\
2.2866	0.041369\\
2.2828	0.040459\\
2.2789	0.039568\\
2.275	0.038694\\
2.2711	0.037837\\
2.2673	0.036998\\
2.2634	0.036176\\
2.2595	0.035371\\
2.2556	0.034582\\
2.2517	0.033809\\
2.2479	0.033052\\
2.244	0.03231\\
2.2401	0.031584\\
2.2362	0.030873\\
2.2324	0.030176\\
2.2285	0.029494\\
2.2246	0.028827\\
2.2207	0.028173\\
2.2168	0.027533\\
2.213	0.026906\\
2.2091	0.026293\\
2.2052	0.025692\\
2.2013	0.025105\\
2.1975	0.02453\\
2.1936	0.023967\\
2.1897	0.023416\\
2.1858	0.022877\\
2.1819	0.02235\\
2.1781	0.021834\\
2.1742	0.021329\\
2.1703	0.020835\\
2.1664	0.020352\\
2.1626	0.019879\\
2.1587	0.019417\\
2.1548	0.018965\\
2.1509	0.018522\\
2.147	0.01809\\
2.1432	0.017667\\
2.1393	0.017253\\
2.1354	0.016848\\
2.1315	0.016453\\
2.1277	0.016066\\
2.1238	0.015688\\
2.1199	0.015318\\
2.116	0.014957\\
2.1121	0.014604\\
2.1083	0.014258\\
2.1044	0.013921\\
2.1005	0.013591\\
2.0966	0.013268\\
2.0928	0.012953\\
2.0889	0.012645\\
2.085	0.012344\\
2.0811	0.01205\\
2.0773	0.011762\\
2.0734	0.011482\\
2.0695	0.011207\\
2.0656	0.010939\\
2.0617	0.010677\\
2.0579	0.010421\\
2.054	0.010171\\
2.0501	0.0099266\\
2.0462	0.0096879\\
2.0424	0.0094548\\
2.0385	0.0092271\\
2.0346	0.0090046\\
2.0307	0.0087874\\
2.0268	0.0085752\\
2.023	0.008368\\
2.0191	0.0081656\\
2.0152	0.0079679\\
2.0113	0.0077749\\
2.0075	0.0075864\\
2.0036	0.0074023\\
1.9997	0.0072226\\
1.9958	0.0070471\\
1.9919	0.0068757\\
1.9881	0.0067084\\
1.9842	0.006545\\
1.9803	0.0063855\\
1.9764	0.0062298\\
1.9726	0.0060778\\
1.9687	0.0059293\\
1.9648	0.0057844\\
1.9609	0.005643\\
1.957	0.0055049\\
1.9532	0.0053701\\
1.9493	0.0052386\\
1.9454	0.0051102\\
1.9415	0.0049848\\
1.9377	0.0048625\\
1.9338	0.004743\\
1.9299	0.0046265\\
1.926	0.0045127\\
1.9221	0.0044017\\
1.9183	0.0042934\\
1.9144	0.0041876\\
1.9105	0.0040844\\
1.9066	0.0039837\\
1.9028	0.0038854\\
1.8989	0.0037895\\
1.895	0.0036959\\
1.8911	0.0036046\\
1.8872	0.0035155\\
1.8834	0.0034285\\
1.8795	0.0033437\\
1.8756	0.0032609\\
1.8717	0.0031801\\
1.8679	0.0031013\\
1.864	0.0030243\\
1.8601	0.0029493\\
1.8562	0.0028761\\
1.8523	0.0028046\\
1.8485	0.002735\\
1.8446	0.002667\\
1.8407	0.0026006\\
1.8368	0.0025359\\
1.833	0.0024727\\
1.8291	0.0024111\\
1.8252	0.0023511\\
1.8213	0.0022924\\
1.8174	0.0022352\\
1.8136	0.0021794\\
1.8097	0.002125\\
1.8058	0.0020719\\
1.8019	0.0020201\\
1.7981	0.0019696\\
1.7942	0.0019203\\
1.7903	0.0018722\\
1.7864	0.0018254\\
1.7825	0.0017796\\
1.7787	0.001735\\
1.7748	0.0016915\\
1.7709	0.001649\\
1.767	0.0016076\\
1.7632	0.0015672\\
1.7593	0.0015278\\
1.7554	0.0014894\\
1.7515	0.0014519\\
1.7476	0.0014154\\
1.7438	0.0013797\\
1.7399	0.0013449\\
1.736	0.001311\\
1.7321	0.001278\\
1.7283	0.0012457\\
1.7244	0.0012142\\
1.7205	0.0011835\\
1.7166	0.0011536\\
1.7127	0.0011244\\
1.7089	0.001096\\
1.705	0.0010682\\
1.7011	0.0010411\\
1.6972	0.0010147\\
1.6934	0.00098899\\
1.6895	0.00096388\\
1.6856	0.00093939\\
1.6817	0.00091552\\
1.6778	0.00089223\\
1.674	0.00086952\\
1.6701	0.00084738\\
1.6662	0.00082579\\
1.6623	0.00080473\\
1.6585	0.0007842\\
1.6546	0.00076418\\
1.6507	0.00074465\\
1.6468	0.00072562\\
1.643	0.00070705\\
1.6391	0.00068895\\
1.6352	0.00067131\\
1.6313	0.0006541\\
1.6274	0.00063732\\
1.6236	0.00062096\\
1.6197	0.00060501\\
1.6158	0.00058945\\
1.6119	0.00057429\\
1.6081	0.00055951\\
1.6042	0.00054509\\
1.6003	0.00053104\\
1.5964	0.00051733\\
1.5925	0.00050398\\
1.5887	0.00049095\\
1.5848	0.00047826\\
1.5809	0.00046588\\
1.577	0.00045381\\
1.5732	0.00044205\\
1.5693	0.00043058\\
1.5654	0.0004194\\
1.5615	0.0004085\\
1.5576	0.00039788\\
1.5538	0.00038752\\
1.5499	0.00037743\\
1.546	0.00036759\\
1.5421	0.00035799\\
1.5383	0.00034864\\
1.5344	0.00033953\\
1.5305	0.00033065\\
1.5266	0.00032199\\
1.5227	0.00031355\\
1.5189	0.00030533\\
1.515	0.00029731\\
1.5111	0.0002895\\
1.5072	0.00028188\\
1.5034	0.00027446\\
1.4995	0.00026723\\
1.4956	0.00026018\\
1.4917	0.00025331\\
1.4878	0.00024661\\
1.484	0.00024009\\
1.4801	0.00023373\\
1.4762	0.00022754\\
1.4723	0.0002215\\
1.4685	0.00021562\\
1.4646	0.00020989\\
1.4607	0.0002043\\
1.4568	0.00019886\\
}--cycle;
      % \input{../../../../.tex/20/KS/SizeDistributionFittingsfig19.tex}
      \legend{}
  
      \nextgroupplot[%
        xmin=0.57363,xmax=5.3544,
        % xmin=1.2642,xmax=5.5028,
        % ymin=0,ymax=1,
      ] % This file was created by matlab2tikz.
%
\definecolor{mycolor1}{rgb}{0.00000,0.44706,0.69804}%
%
\addplot[ybar interval, fill=mycolor1, fill opacity=0.35, area legend, draw=none] table[row sep=crcr, x=Lower, y=Count] {%
Lower	Upper	Count\\
1.4568	1.7532	0.00035986\\
1.7532	2.0496	0.0035087\\
2.0496	2.346	0.031938\\
2.346	2.6425	0.14341\\
2.6425	2.9389	0.4519\\
2.9389	3.2353	0.83146\\
3.2353	3.5317	0.8948\\
3.5317	3.8281	0.55068\\
3.8281	4.1245	0.26801\\
4.1245	4.4209	0.12964\\
4.4209	4.7173	0.045973\\
4.7173	5.0137	0.013225\\
5.0137	5.3101	0.0088167\\
5.3101	5.3101	0.0088167\\
};
\addlegendentry{\ExactHistogramLegend}

\addplot [color=mycolor1, only marks, every error bar/.append style={opacity=0.45}, mark=*, mark size=0pt, draw=none, forget plot]
 plot [error bars/.cd, y dir=both, y explicit, error bar style={line width=1pt, color=mycolor1}, error mark options={mark=none,mark size=0pt}]
 table[row sep=crcr, y error plus index=2, y error minus index=3]{%
1.605	0.00035986	0.00035157	0.00035157\\
1.9014	0.0035087	0.0011317	0.0011317\\
2.1978	0.031938	0.0030906	0.0030906\\
2.4943	0.14341	0.0064828	0.0064828\\
2.7907	0.4519	0.0094693	0.0094693\\
3.0871	0.83146	0.012868	0.012868\\
3.3835	0.8948	0.013388	0.013388\\
3.6799	0.55068	0.0092472	0.0092472\\
3.9763	0.26801	0.0083433	0.0083433\\
4.2727	0.12964	0.0050878	0.0050878\\
4.5691	0.045973	0.0028128	0.0028128\\
4.8655	0.013225	0.0012801	0.0012801\\
5.1619	0.0088167	0.00025118	0.00025118\\
};
\addplot [color=mycolor1]
  table[row sep=crcr]{%
1.4568	0.00029246\\
1.6613	0.0010428\\
1.7847	0.0021998\\
1.8734	0.0037456\\
1.9428	0.0056747\\
2.0007	0.0080179\\
2.0508	0.010811\\
2.0933	0.01391\\
2.1318	0.017473\\
2.1665	0.021427\\
2.1974	0.025656\\
2.2283	0.030676\\
2.2553	0.035817\\
2.2823	0.041758\\
2.3054	0.047564\\
2.3285	0.054104\\
2.3517	0.061451\\
2.3748	0.069683\\
2.3941	0.077276\\
2.4134	0.085586\\
2.4327	0.094662\\
2.452	0.10455\\
2.4713	0.1153\\
2.4905	0.12695\\
2.5098	0.13956\\
2.5291	0.15316\\
2.5484	0.16779\\
2.5677	0.18348\\
2.587	0.20026\\
2.6063	0.21815\\
2.6255	0.23716\\
2.6448	0.25731\\
2.6641	0.27858\\
2.6834	0.30097\\
2.7027	0.32444\\
2.722	0.34895\\
2.7451	0.37968\\
2.7683	0.41173\\
2.7953	0.4506\\
2.8223	0.4908\\
2.857	0.54391\\
2.9611	0.70447\\
2.9881	0.744\\
3.0113	0.77638\\
3.0305	0.80204\\
3.0498	0.82629\\
3.0653	0.84454\\
3.0807	0.86167\\
3.0961	0.87758\\
3.1115	0.89219\\
3.1231	0.90224\\
3.1347	0.91149\\
3.1463	0.9199\\
3.1578	0.92746\\
3.1694	0.93414\\
3.181	0.93992\\
3.1925	0.94479\\
3.2041	0.94874\\
3.2157	0.95176\\
3.2273	0.95385\\
3.2388	0.95501\\
3.2504	0.95525\\
3.262	0.95457\\
3.2735	0.95299\\
3.2851	0.95051\\
3.2967	0.94716\\
3.3083	0.94297\\
3.3198	0.93794\\
3.3314	0.93212\\
3.343	0.92552\\
3.3545	0.91817\\
3.3661	0.91012\\
3.3815	0.89833\\
3.397	0.88542\\
3.4124	0.87147\\
3.4278	0.85656\\
3.4471	0.83673\\
3.4664	0.81572\\
3.4896	0.78919\\
3.5127	0.7615\\
3.5436	0.72326\\
3.586	0.66927\\
3.6554	0.58084\\
3.6901	0.53796\\
3.721	0.50114\\
3.748	0.47014\\
3.775	0.44042\\
3.7981	0.41602\\
3.8213	0.39268\\
3.8444	0.37041\\
3.8676	0.34921\\
3.8907	0.32909\\
3.9138	0.31003\\
3.937	0.292\\
3.9601	0.27497\\
3.9833	0.2589\\
4.0064	0.24376\\
4.0296	0.22949\\
4.0527	0.21605\\
4.0758	0.20339\\
4.1028	0.18954\\
4.1298	0.17661\\
4.1568	0.16453\\
4.1838	0.15322\\
4.2147	0.14118\\
4.2456	0.13\\
4.2764	0.11958\\
4.3111	0.10871\\
4.3458	0.098647\\
4.3806	0.089327\\
4.4191	0.079771\\
4.4577	0.070991\\
4.5001	0.062168\\
4.5426	0.054163\\
4.585	0.046928\\
4.6274	0.040422\\
4.6737	0.034107\\
4.72	0.028562\\
4.7701	0.023363\\
4.8203	0.018935\\
4.8743	0.014946\\
4.9321	0.011465\\
4.9938	0.0085278\\
5.0633	0.0060171\\
5.1404	0.0040074\\
5.233	0.0023988\\
5.3101	0.0015326\\
};
\addlegendentry{\ExactBLNFitLegend}


\addplot[area legend, draw=none, fill=mycolor1, fill opacity=0.15, forget plot]
table[row sep=crcr] {%
x	y\\
1.4568	0.00036823\\
1.4607	0.00037721\\
1.4645	0.00038639\\
1.4684	0.00039578\\
1.4723	0.0004054\\
1.4761	0.00041523\\
1.48	0.00042529\\
1.4838	0.00043558\\
1.4877	0.0004461\\
1.4915	0.00045687\\
1.4954	0.00046788\\
1.4993	0.00047915\\
1.5031	0.00049068\\
1.507	0.00050246\\
1.5108	0.00051452\\
1.5147	0.00052685\\
1.5185	0.00053947\\
1.5224	0.00055237\\
1.5263	0.00056556\\
1.5301	0.00057905\\
1.534	0.00059285\\
1.5378	0.00060697\\
1.5417	0.0006214\\
1.5455	0.00063616\\
1.5494	0.00065126\\
1.5533	0.00066669\\
1.5571	0.00068247\\
1.561	0.00069862\\
1.5648	0.00071512\\
1.5687	0.000732\\
1.5725	0.00074925\\
1.5764	0.0007669\\
1.5803	0.00078494\\
1.5841	0.00080339\\
1.588	0.00082225\\
1.5918	0.00084153\\
1.5957	0.00086125\\
1.5995	0.00088141\\
1.6034	0.00090202\\
1.6073	0.00092309\\
1.6111	0.00094463\\
1.615	0.00096666\\
1.6188	0.00098917\\
1.6227	0.0010122\\
1.6265	0.0010357\\
1.6304	0.0010598\\
1.6343	0.0010844\\
1.6381	0.0011095\\
1.642	0.0011352\\
1.6458	0.0011615\\
1.6497	0.0011884\\
1.6535	0.0012158\\
1.6574	0.0012439\\
1.6613	0.0012726\\
1.6651	0.0013019\\
1.669	0.0013319\\
1.6728	0.0013625\\
1.6767	0.0013938\\
1.6805	0.0014258\\
1.6844	0.0014586\\
1.6883	0.001492\\
1.6921	0.0015262\\
1.696	0.0015612\\
1.6998	0.0015969\\
1.7037	0.0016334\\
1.7075	0.0016707\\
1.7114	0.0017089\\
1.7153	0.0017478\\
1.7191	0.0017877\\
1.723	0.0018284\\
1.7268	0.0018701\\
1.7307	0.0019126\\
1.7345	0.0019561\\
1.7384	0.0020006\\
1.7423	0.002046\\
1.7461	0.0020924\\
1.75	0.0021399\\
1.7538	0.0021884\\
1.7577	0.002238\\
1.7615	0.0022886\\
1.7654	0.0023404\\
1.7693	0.0023934\\
1.7731	0.0024475\\
1.777	0.0025028\\
1.7808	0.0025593\\
1.7847	0.0026171\\
1.7885	0.0026761\\
1.7924	0.0027364\\
1.7963	0.0027981\\
1.8001	0.0028611\\
1.804	0.0029256\\
1.8078	0.0029914\\
1.8117	0.0030587\\
1.8155	0.0031275\\
1.8194	0.0031978\\
1.8233	0.0032696\\
1.8271	0.0033431\\
1.831	0.0034181\\
1.8348	0.0034948\\
1.8387	0.0035732\\
1.8425	0.0036534\\
1.8464	0.0037353\\
1.8503	0.003819\\
1.8541	0.0039045\\
1.858	0.003992\\
1.8618	0.0040814\\
1.8657	0.0041727\\
1.8695	0.0042661\\
1.8734	0.0043615\\
1.8773	0.004459\\
1.8811	0.0045587\\
1.885	0.0046606\\
1.8888	0.0047648\\
1.8927	0.0048712\\
1.8965	0.00498\\
1.9004	0.0050912\\
1.9043	0.0052048\\
1.9081	0.005321\\
1.912	0.0054397\\
1.9158	0.0055611\\
1.9197	0.0056851\\
1.9235	0.0058119\\
1.9274	0.0059415\\
1.9313	0.0060739\\
1.9351	0.0062093\\
1.939	0.0063477\\
1.9428	0.0064891\\
1.9467	0.0066336\\
1.9505	0.0067814\\
1.9544	0.0069324\\
1.9583	0.0070867\\
1.9621	0.0072445\\
1.966	0.0074057\\
1.9698	0.0075705\\
1.9737	0.0077389\\
1.9775	0.0079111\\
1.9814	0.0080871\\
1.9853	0.0082669\\
1.9891	0.0084507\\
1.993	0.0086386\\
1.9968	0.0088306\\
2.0007	0.0090269\\
2.0045	0.0092275\\
2.0084	0.0094325\\
2.0123	0.009642\\
2.0161	0.0098562\\
2.02	0.010075\\
2.0238	0.010299\\
2.0277	0.010527\\
2.0315	0.010761\\
2.0354	0.011\\
2.0393	0.011244\\
2.0431	0.011493\\
2.047	0.011748\\
2.0508	0.012009\\
2.0547	0.012275\\
2.0585	0.012547\\
2.0624	0.012825\\
2.0663	0.013109\\
2.0701	0.013399\\
2.074	0.013696\\
2.0778	0.013999\\
2.0817	0.014309\\
2.0855	0.014626\\
2.0894	0.014949\\
2.0933	0.01528\\
2.0971	0.015617\\
2.101	0.015963\\
2.1048	0.016315\\
2.1087	0.016675\\
2.1125	0.017044\\
2.1164	0.01742\\
2.1203	0.017804\\
2.1241	0.018197\\
2.128	0.018598\\
2.1318	0.019007\\
2.1357	0.019426\\
2.1395	0.019854\\
2.1434	0.02029\\
2.1473	0.020737\\
2.1511	0.021192\\
2.155	0.021658\\
2.1588	0.022134\\
2.1627	0.022619\\
2.1665	0.023115\\
2.1704	0.023622\\
2.1743	0.02414\\
2.1781	0.024668\\
2.182	0.025208\\
2.1858	0.025759\\
2.1897	0.026322\\
2.1935	0.026897\\
2.1974	0.027484\\
2.2013	0.028083\\
2.2051	0.028695\\
2.209	0.029319\\
2.2128	0.029957\\
2.2167	0.030608\\
2.2205	0.031273\\
2.2244	0.031952\\
2.2283	0.032644\\
2.2321	0.033351\\
2.236	0.034073\\
2.2398	0.034809\\
2.2437	0.035561\\
2.2475	0.036328\\
2.2514	0.037111\\
2.2553	0.03791\\
2.2591	0.038725\\
2.263	0.039557\\
2.2668	0.040406\\
2.2707	0.041272\\
2.2745	0.042155\\
2.2784	0.043056\\
2.2823	0.043976\\
2.2861	0.044913\\
2.29	0.04587\\
2.2938	0.046845\\
2.2977	0.04784\\
2.3015	0.048855\\
2.3054	0.04989\\
2.3093	0.050945\\
2.3131	0.05202\\
2.317	0.053117\\
2.3208	0.054235\\
2.3247	0.055375\\
2.3285	0.056537\\
2.3324	0.057722\\
2.3363	0.058929\\
2.3401	0.060159\\
2.344	0.061413\\
2.3478	0.062691\\
2.3517	0.063993\\
2.3555	0.065319\\
2.3594	0.066671\\
2.3633	0.068048\\
2.3671	0.06945\\
2.371	0.070879\\
2.3748	0.072334\\
2.3787	0.073816\\
2.3825	0.075325\\
2.3864	0.076862\\
2.3903	0.078426\\
2.3941	0.08002\\
2.398	0.081641\\
2.4018	0.083292\\
2.4057	0.084973\\
2.4095	0.086683\\
2.4134	0.088424\\
2.4173	0.090196\\
2.4211	0.091998\\
2.425	0.093832\\
2.4288	0.095698\\
2.4327	0.097596\\
2.4365	0.099527\\
2.4404	0.10149\\
2.4443	0.10349\\
2.4481	0.10552\\
2.452	0.10758\\
2.4558	0.10968\\
2.4597	0.11182\\
2.4635	0.11399\\
2.4674	0.11619\\
2.4713	0.11843\\
2.4751	0.12071\\
2.479	0.12303\\
2.4828	0.12538\\
2.4867	0.12777\\
2.4905	0.1302\\
2.4944	0.13266\\
2.4983	0.13517\\
2.5021	0.13771\\
2.506	0.14029\\
2.5098	0.14291\\
2.5137	0.14557\\
2.5175	0.14828\\
2.5214	0.15102\\
2.5253	0.1538\\
2.5291	0.15663\\
2.533	0.15949\\
2.5368	0.1624\\
2.5407	0.16535\\
2.5445	0.16834\\
2.5484	0.17138\\
2.5523	0.17445\\
2.5561	0.17757\\
2.56	0.18074\\
2.5638	0.18394\\
2.5677	0.18719\\
2.5715	0.19049\\
2.5754	0.19383\\
2.5793	0.19721\\
2.5831	0.20064\\
2.587	0.20411\\
2.5908	0.20763\\
2.5947	0.21119\\
2.5985	0.2148\\
2.6024	0.21845\\
2.6063	0.22215\\
2.6101	0.22589\\
2.614	0.22968\\
2.6178	0.23352\\
2.6217	0.2374\\
2.6255	0.24132\\
2.6294	0.24529\\
2.6333	0.24931\\
2.6371	0.25337\\
2.641	0.25748\\
2.6448	0.26164\\
2.6487	0.26583\\
2.6525	0.27008\\
2.6564	0.27437\\
2.6603	0.2787\\
2.6641	0.28309\\
2.668	0.28751\\
2.6718	0.29198\\
2.6757	0.2965\\
2.6795	0.30106\\
2.6834	0.30566\\
2.6873	0.31031\\
2.6911	0.315\\
2.695	0.31973\\
2.6988	0.32451\\
2.7027	0.32933\\
2.7065	0.33419\\
2.7104	0.33909\\
2.7143	0.34404\\
2.7181	0.34903\\
2.722	0.35406\\
2.7258	0.35912\\
2.7297	0.36423\\
2.7335	0.36938\\
2.7374	0.37456\\
2.7413	0.37979\\
2.7451	0.38505\\
2.749	0.39035\\
2.7528	0.39568\\
2.7567	0.40106\\
2.7605	0.40646\\
2.7644	0.4119\\
2.7683	0.41738\\
2.7721	0.42288\\
2.776	0.42842\\
2.7798	0.43399\\
2.7837	0.43959\\
2.7875	0.44522\\
2.7914	0.45088\\
2.7953	0.45657\\
2.7991	0.46228\\
2.803	0.46803\\
2.8068	0.47379\\
2.8107	0.47958\\
2.8145	0.48539\\
2.8184	0.49123\\
2.8223	0.49709\\
2.8261	0.50296\\
2.83	0.50886\\
2.8338	0.51477\\
2.8377	0.5207\\
2.8415	0.52664\\
2.8454	0.5326\\
2.8493	0.53858\\
2.8531	0.54456\\
2.857	0.55056\\
2.8608	0.55656\\
2.8647	0.56258\\
2.8685	0.5686\\
2.8724	0.57462\\
2.8763	0.58065\\
2.8801	0.58669\\
2.884	0.59272\\
2.8878	0.59876\\
2.8917	0.60479\\
2.8955	0.61082\\
2.8994	0.61685\\
2.9033	0.62287\\
2.9071	0.62889\\
2.911	0.6349\\
2.9148	0.6409\\
2.9187	0.64689\\
2.9225	0.65286\\
2.9264	0.65882\\
2.9303	0.66477\\
2.9341	0.6707\\
2.938	0.67661\\
2.9418	0.68251\\
2.9457	0.68838\\
2.9495	0.69423\\
2.9534	0.70005\\
2.9573	0.70585\\
2.9611	0.71162\\
2.965	0.71737\\
2.9688	0.72309\\
2.9727	0.72877\\
2.9765	0.73442\\
2.9804	0.74004\\
2.9843	0.74562\\
2.9881	0.75117\\
2.992	0.75667\\
2.9958	0.76214\\
2.9997	0.76757\\
3.0035	0.77295\\
3.0074	0.77829\\
3.0113	0.78358\\
3.0151	0.78883\\
3.019	0.79402\\
3.0228	0.79917\\
3.0267	0.80427\\
3.0305	0.80931\\
3.0344	0.8143\\
3.0383	0.81923\\
3.0421	0.82411\\
3.046	0.82893\\
3.0498	0.83368\\
3.0537	0.83838\\
3.0575	0.84301\\
3.0614	0.84757\\
3.0653	0.85207\\
3.0691	0.8565\\
3.073	0.86086\\
3.0768	0.86515\\
3.0807	0.86937\\
3.0845	0.87351\\
3.0884	0.87758\\
3.0923	0.88157\\
3.0961	0.88548\\
3.1	0.8893\\
3.1038	0.89305\\
3.1077	0.89671\\
3.1115	0.90029\\
3.1154	0.90378\\
3.1193	0.90719\\
3.1231	0.9105\\
3.127	0.91372\\
3.1308	0.91685\\
3.1347	0.91989\\
3.1385	0.92284\\
3.1424	0.92568\\
3.1463	0.92843\\
3.1501	0.93109\\
3.154	0.93364\\
3.1578	0.93609\\
3.1617	0.93844\\
3.1655	0.94069\\
3.1694	0.94284\\
3.1733	0.94489\\
3.1771	0.94683\\
3.181	0.94866\\
3.1848	0.95039\\
3.1887	0.95201\\
3.1925	0.95353\\
3.1964	0.95494\\
3.2003	0.95625\\
3.2041	0.95744\\
3.208	0.95853\\
3.2118	0.95951\\
3.2157	0.96039\\
3.2195	0.96115\\
3.2234	0.96181\\
3.2273	0.96236\\
3.2311	0.96281\\
3.235	0.96314\\
3.2388	0.96338\\
3.2427	0.9635\\
3.2465	0.96352\\
3.2504	0.96343\\
3.2543	0.96324\\
3.2581	0.96295\\
3.262	0.96255\\
3.2658	0.96205\\
3.2697	0.96145\\
3.2735	0.96075\\
3.2774	0.95995\\
3.2813	0.95905\\
3.2851	0.95805\\
3.289	0.95696\\
3.2928	0.95577\\
3.2967	0.95448\\
3.3005	0.95311\\
3.3044	0.95164\\
3.3083	0.95008\\
3.3121	0.94843\\
3.316	0.94669\\
3.3198	0.94486\\
3.3237	0.94295\\
3.3275	0.94095\\
3.3314	0.93887\\
3.3353	0.93671\\
3.3391	0.93447\\
3.343	0.93215\\
3.3468	0.92974\\
3.3507	0.92727\\
3.3545	0.92471\\
3.3584	0.92208\\
3.3623	0.91938\\
3.3661	0.91661\\
3.37	0.91376\\
3.3738	0.91084\\
3.3777	0.90786\\
3.3815	0.90481\\
3.3854	0.90169\\
3.3893	0.8985\\
3.3931	0.89525\\
3.397	0.89194\\
3.4008	0.88856\\
3.4047	0.88512\\
3.4085	0.88163\\
3.4124	0.87807\\
3.4163	0.87446\\
3.4201	0.87079\\
3.424	0.86706\\
3.4278	0.86328\\
3.4317	0.85944\\
3.4355	0.85556\\
3.4394	0.85162\\
3.4433	0.84763\\
3.4471	0.84359\\
3.451	0.83951\\
3.4548	0.83538\\
3.4587	0.8312\\
3.4626	0.82699\\
3.4664	0.82272\\
3.4703	0.81842\\
3.4741	0.81408\\
3.478	0.8097\\
3.4818	0.80528\\
3.4857	0.80082\\
3.4896	0.79633\\
3.4934	0.79181\\
3.4973	0.78725\\
3.5011	0.78266\\
3.505	0.77804\\
3.5088	0.7734\\
3.5127	0.76872\\
3.5166	0.76402\\
3.5204	0.7593\\
3.5243	0.75455\\
3.5281	0.74978\\
3.532	0.74499\\
3.5358	0.74017\\
3.5397	0.73534\\
3.5436	0.7305\\
3.5474	0.72563\\
3.5513	0.72075\\
3.5551	0.71586\\
3.559	0.71096\\
3.5628	0.70604\\
3.5667	0.70111\\
3.5706	0.69618\\
3.5744	0.69123\\
3.5783	0.68628\\
3.5821	0.68133\\
3.586	0.67637\\
3.5898	0.6714\\
3.5937	0.66643\\
3.5976	0.66147\\
3.6014	0.6565\\
3.6053	0.65153\\
3.6091	0.64656\\
3.613	0.6416\\
3.6168	0.63664\\
3.6207	0.63168\\
3.6246	0.62673\\
3.6284	0.62179\\
3.6323	0.61685\\
3.6361	0.61192\\
3.64	0.607\\
3.6438	0.60209\\
3.6477	0.59719\\
3.6516	0.5923\\
3.6554	0.58743\\
3.6593	0.58257\\
3.6631	0.57772\\
3.667	0.57288\\
3.6708	0.56807\\
3.6747	0.56326\\
3.6786	0.55848\\
3.6824	0.55371\\
3.6863	0.54896\\
3.6901	0.54423\\
3.694	0.53951\\
3.6978	0.53482\\
3.7017	0.53015\\
3.7056	0.5255\\
3.7094	0.52087\\
3.7133	0.51626\\
3.7171	0.51168\\
3.721	0.50711\\
3.7248	0.50258\\
3.7287	0.49806\\
3.7326	0.49357\\
3.7364	0.48911\\
3.7403	0.48467\\
3.7441	0.48025\\
3.748	0.47586\\
3.7518	0.4715\\
3.7557	0.46717\\
3.7596	0.46286\\
3.7634	0.45858\\
3.7673	0.45433\\
3.7711	0.4501\\
3.775	0.44591\\
3.7788	0.44174\\
3.7827	0.4376\\
3.7866	0.43349\\
3.7904	0.42941\\
3.7943	0.42535\\
3.7981	0.42133\\
3.802	0.41734\\
3.8058	0.41337\\
3.8097	0.40944\\
3.8136	0.40554\\
3.8174	0.40166\\
3.8213	0.39782\\
3.8251	0.39401\\
3.829	0.39022\\
3.8328	0.38647\\
3.8367	0.38275\\
3.8406	0.37905\\
3.8444	0.37539\\
3.8483	0.37176\\
3.8521	0.36816\\
3.856	0.36459\\
3.8598	0.36105\\
3.8637	0.35754\\
3.8676	0.35406\\
3.8714	0.35061\\
3.8753	0.34718\\
3.8791	0.34379\\
3.883	0.34043\\
3.8868	0.3371\\
3.8907	0.3338\\
3.8946	0.33053\\
3.8984	0.32729\\
3.9023	0.32408\\
3.9061	0.32089\\
3.91	0.31774\\
3.9138	0.31461\\
3.9177	0.31152\\
3.9216	0.30845\\
3.9254	0.30541\\
3.9293	0.3024\\
3.9331	0.29941\\
3.937	0.29646\\
3.9408	0.29353\\
3.9447	0.29063\\
3.9486	0.28776\\
3.9524	0.28491\\
3.9563	0.28209\\
3.9601	0.2793\\
3.964	0.27654\\
3.9678	0.2738\\
3.9717	0.27108\\
3.9756	0.2684\\
3.9794	0.26574\\
3.9833	0.2631\\
3.9871	0.26049\\
3.991	0.25791\\
3.9948	0.25535\\
3.9987	0.25281\\
4.0026	0.2503\\
4.0064	0.24781\\
4.0103	0.24535\\
4.0141	0.24291\\
4.018	0.24049\\
4.0218	0.2381\\
4.0257	0.23573\\
4.0296	0.23339\\
4.0334	0.23106\\
4.0373	0.22876\\
4.0411	0.22648\\
4.045	0.22423\\
4.0488	0.22199\\
4.0527	0.21978\\
4.0566	0.21758\\
4.0604	0.21541\\
4.0643	0.21326\\
4.0681	0.21113\\
4.072	0.20902\\
4.0758	0.20693\\
4.0797	0.20487\\
4.0836	0.20282\\
4.0874	0.20079\\
4.0913	0.19878\\
4.0951	0.19678\\
4.099	0.19481\\
4.1028	0.19286\\
4.1067	0.19092\\
4.1106	0.18901\\
4.1144	0.18711\\
4.1183	0.18523\\
4.1221	0.18337\\
4.126	0.18152\\
4.1298	0.17969\\
4.1337	0.17788\\
4.1376	0.17609\\
4.1414	0.17431\\
4.1453	0.17255\\
4.1491	0.17081\\
4.153	0.16908\\
4.1568	0.16737\\
4.1607	0.16567\\
4.1646	0.164\\
4.1684	0.16233\\
4.1723	0.16068\\
4.1761	0.15905\\
4.18	0.15743\\
4.1838	0.15583\\
4.1877	0.15424\\
4.1916	0.15267\\
4.1954	0.15111\\
4.1993	0.14956\\
4.2031	0.14803\\
4.207	0.14651\\
4.2108	0.14501\\
4.2147	0.14352\\
4.2186	0.14204\\
4.2224	0.14058\\
4.2263	0.13913\\
4.2301	0.1377\\
4.234	0.13627\\
4.2378	0.13486\\
4.2417	0.13347\\
4.2456	0.13208\\
4.2494	0.13071\\
4.2533	0.12935\\
4.2571	0.128\\
4.261	0.12667\\
4.2648	0.12534\\
4.2687	0.12403\\
4.2726	0.12273\\
4.2764	0.12145\\
4.2803	0.12017\\
4.2841	0.11891\\
4.288	0.11765\\
4.2918	0.11641\\
4.2957	0.11518\\
4.2996	0.11396\\
4.3034	0.11275\\
4.3073	0.11155\\
4.3111	0.11036\\
4.315	0.10919\\
4.3188	0.10802\\
4.3227	0.10687\\
4.3266	0.10572\\
4.3304	0.10459\\
4.3343	0.10346\\
4.3381	0.10235\\
4.342	0.10124\\
4.3458	0.10014\\
4.3497	0.099059\\
4.3536	0.097983\\
4.3574	0.096917\\
4.3613	0.09586\\
4.3651	0.094812\\
4.369	0.093774\\
4.3728	0.092744\\
4.3767	0.091724\\
4.3806	0.090713\\
4.3844	0.089711\\
4.3883	0.088717\\
4.3921	0.087732\\
4.396	0.086756\\
4.3998	0.085788\\
4.4037	0.084828\\
4.4076	0.083877\\
4.4114	0.082934\\
4.4153	0.082\\
4.4191	0.081073\\
4.423	0.080155\\
4.4268	0.079244\\
4.4307	0.078341\\
4.4346	0.077446\\
4.4384	0.076559\\
4.4423	0.075679\\
4.4461	0.074807\\
4.45	0.073943\\
4.4538	0.073085\\
4.4577	0.072235\\
4.4616	0.071393\\
4.4654	0.070558\\
4.4693	0.069729\\
4.4731	0.068908\\
4.477	0.068094\\
4.4808	0.067287\\
4.4847	0.066487\\
4.4886	0.065694\\
4.4924	0.064907\\
4.4963	0.064128\\
4.5001	0.063355\\
4.504	0.062589\\
4.5078	0.061829\\
4.5117	0.061076\\
4.5156	0.06033\\
4.5194	0.05959\\
4.5233	0.058856\\
4.5271	0.058129\\
4.531	0.057408\\
4.5348	0.056694\\
4.5387	0.055986\\
4.5426	0.055284\\
4.5464	0.054589\\
4.5503	0.0539\\
4.5541	0.053216\\
4.558	0.052539\\
4.5618	0.051869\\
4.5657	0.051204\\
4.5696	0.050545\\
4.5734	0.049892\\
4.5773	0.049246\\
4.5811	0.048605\\
4.585	0.04797\\
4.5888	0.047341\\
4.5927	0.046718\\
4.5966	0.046101\\
4.6004	0.045489\\
4.6043	0.044884\\
4.6081	0.044284\\
4.612	0.04369\\
4.6158	0.043101\\
4.6197	0.042519\\
4.6236	0.041942\\
4.6274	0.04137\\
4.6313	0.040805\\
4.6351	0.040245\\
4.639	0.03969\\
4.6428	0.039141\\
4.6467	0.038598\\
4.6506	0.03806\\
4.6544	0.037528\\
4.6583	0.037001\\
4.6621	0.036479\\
4.666	0.035963\\
4.6698	0.035452\\
4.6737	0.034947\\
4.6776	0.034447\\
4.6814	0.033952\\
4.6853	0.033463\\
4.6891	0.032979\\
4.693	0.0325\\
4.6968	0.032026\\
4.7007	0.031558\\
4.7046	0.031095\\
4.7084	0.030636\\
4.7123	0.030183\\
4.7161	0.029735\\
4.72	0.029293\\
4.7238	0.028855\\
4.7277	0.028422\\
4.7316	0.027994\\
4.7354	0.027571\\
4.7393	0.027153\\
4.7431	0.02674\\
4.747	0.026332\\
4.7508	0.025929\\
4.7547	0.02553\\
4.7586	0.025136\\
4.7624	0.024747\\
4.7663	0.024363\\
4.7701	0.023983\\
4.774	0.023609\\
4.7778	0.023238\\
4.7817	0.022873\\
4.7856	0.022511\\
4.7894	0.022155\\
4.7933	0.021803\\
4.7971	0.021455\\
4.801	0.021112\\
4.8048	0.020773\\
4.8087	0.020439\\
4.8126	0.020109\\
4.8164	0.019783\\
4.8203	0.019462\\
4.8241	0.019145\\
4.828	0.018832\\
4.8318	0.018523\\
4.8357	0.018218\\
4.8396	0.017917\\
4.8434	0.017621\\
4.8473	0.017328\\
4.8511	0.01704\\
4.855	0.016755\\
4.8588	0.016475\\
4.8627	0.016198\\
4.8666	0.015925\\
4.8704	0.015656\\
4.8743	0.015391\\
4.8781	0.015129\\
4.882	0.014872\\
4.8858	0.014617\\
4.8897	0.014367\\
4.8936	0.01412\\
4.8974	0.013877\\
4.9013	0.013637\\
4.9051	0.0134\\
4.909	0.013167\\
4.9128	0.012938\\
4.9167	0.012712\\
4.9206	0.012489\\
4.9244	0.01227\\
4.9283	0.012053\\
4.9321	0.01184\\
4.936	0.011631\\
4.9398	0.011424\\
4.9437	0.011221\\
4.9476	0.01102\\
4.9514	0.010823\\
4.9553	0.010629\\
4.9591	0.010437\\
4.963	0.010249\\
4.9668	0.010064\\
4.9707	0.009881\\
4.9746	0.0097013\\
4.9784	0.0095244\\
4.9823	0.0093502\\
4.9861	0.0091789\\
4.99	0.0090102\\
4.9938	0.0088442\\
4.9977	0.0086808\\
5.0016	0.0085201\\
5.0054	0.0083619\\
5.0093	0.0082063\\
5.0131	0.0080532\\
5.017	0.0079026\\
5.0208	0.0077544\\
5.0247	0.0076086\\
5.0286	0.0074652\\
5.0324	0.0073242\\
5.0363	0.0071855\\
5.0401	0.0070491\\
5.044	0.0069149\\
5.0478	0.006783\\
5.0517	0.0066533\\
5.0556	0.0065258\\
5.0594	0.0064004\\
5.0633	0.0062771\\
5.0671	0.0061559\\
5.071	0.0060367\\
5.0748	0.0059196\\
5.0787	0.0058045\\
5.0826	0.0056914\\
5.0864	0.0055802\\
5.0903	0.0054709\\
5.0941	0.0053635\\
5.098	0.0052579\\
5.1018	0.0051542\\
5.1057	0.0050523\\
5.1096	0.0049522\\
5.1134	0.0048539\\
5.1173	0.0047573\\
5.1211	0.0046623\\
5.125	0.0045691\\
5.1288	0.0044775\\
5.1327	0.0043876\\
5.1366	0.0042993\\
5.1404	0.0042125\\
5.1443	0.0041273\\
5.1481	0.0040437\\
5.152	0.0039615\\
5.1558	0.0038809\\
5.1597	0.0038017\\
5.1636	0.0037239\\
5.1674	0.0036476\\
5.1713	0.0035727\\
5.1751	0.0034992\\
5.179	0.003427\\
5.1828	0.0033561\\
5.1867	0.0032866\\
5.1906	0.0032184\\
5.1944	0.0031514\\
5.1983	0.0030857\\
5.2021	0.0030212\\
5.206	0.002958\\
5.2098	0.0028959\\
5.2137	0.002835\\
5.2176	0.0027753\\
5.2214	0.0027167\\
5.2253	0.0026592\\
5.2291	0.0026028\\
5.233	0.0025475\\
5.2368	0.0024933\\
5.2407	0.0024401\\
5.2446	0.0023879\\
5.2484	0.0023368\\
5.2523	0.0022866\\
5.2561	0.0022374\\
5.26	0.0021892\\
5.2638	0.002142\\
5.2677	0.0020956\\
5.2716	0.0020502\\
5.2754	0.0020056\\
5.2793	0.001962\\
5.2831	0.0019192\\
5.287	0.0018773\\
5.2908	0.0018362\\
5.2947	0.0017959\\
5.2986	0.0017564\\
5.3024	0.0017178\\
5.3063	0.0016799\\
5.3101	0.0016427\\
5.3101	0.0014224\\
5.3063	0.001456\\
5.3024	0.0014904\\
5.2986	0.0015255\\
5.2947	0.0015614\\
5.2908	0.0015981\\
5.287	0.0016355\\
5.2831	0.0016737\\
5.2793	0.0017127\\
5.2754	0.0017526\\
5.2716	0.0017933\\
5.2677	0.0018349\\
5.2638	0.0018773\\
5.26	0.0019206\\
5.2561	0.0019648\\
5.2523	0.00201\\
5.2484	0.0020561\\
5.2446	0.0021031\\
5.2407	0.0021511\\
5.2368	0.0022001\\
5.233	0.0022501\\
5.2291	0.0023012\\
5.2253	0.0023532\\
5.2214	0.0024064\\
5.2176	0.0024606\\
5.2137	0.0025159\\
5.2098	0.0025723\\
5.206	0.0026299\\
5.2021	0.0026886\\
5.1983	0.0027485\\
5.1944	0.0028096\\
5.1906	0.002872\\
5.1867	0.0029355\\
5.1828	0.0030003\\
5.179	0.0030664\\
5.1751	0.0031338\\
5.1713	0.0032025\\
5.1674	0.0032725\\
5.1636	0.003344\\
5.1597	0.0034168\\
5.1558	0.003491\\
5.152	0.0035666\\
5.1481	0.0036437\\
5.1443	0.0037223\\
5.1404	0.0038023\\
5.1366	0.0038839\\
5.1327	0.0039671\\
5.1288	0.0040518\\
5.125	0.0041381\\
5.1211	0.0042261\\
5.1173	0.0043157\\
5.1134	0.0044069\\
5.1096	0.0044999\\
5.1057	0.0045945\\
5.1018	0.004691\\
5.098	0.0047892\\
5.0941	0.0048892\\
5.0903	0.004991\\
5.0864	0.0050947\\
5.0826	0.0052002\\
5.0787	0.0053077\\
5.0748	0.0054171\\
5.071	0.0055284\\
5.0671	0.0056418\\
5.0633	0.0057572\\
5.0594	0.0058746\\
5.0556	0.0059941\\
5.0517	0.0061156\\
5.0478	0.0062394\\
5.044	0.0063652\\
5.0401	0.0064933\\
5.0363	0.0066236\\
5.0324	0.0067561\\
5.0286	0.0068909\\
5.0247	0.007028\\
5.0208	0.0071674\\
5.017	0.0073092\\
5.0131	0.0074534\\
5.0093	0.0076001\\
5.0054	0.0077492\\
5.0016	0.0079007\\
4.9977	0.0080548\\
4.9938	0.0082114\\
4.99	0.0083706\\
4.9861	0.0085324\\
4.9823	0.0086969\\
4.9784	0.008864\\
4.9746	0.0090338\\
4.9707	0.0092064\\
4.9668	0.0093817\\
4.963	0.0095598\\
4.9591	0.0097408\\
4.9553	0.0099246\\
4.9514	0.010111\\
4.9476	0.010301\\
4.9437	0.010493\\
4.9398	0.010689\\
4.936	0.010888\\
4.9321	0.011089\\
4.9283	0.011294\\
4.9244	0.011502\\
4.9206	0.011713\\
4.9167	0.011927\\
4.9128	0.012144\\
4.909	0.012365\\
4.9051	0.012589\\
4.9013	0.012816\\
4.8974	0.013046\\
4.8936	0.01328\\
4.8897	0.013517\\
4.8858	0.013758\\
4.882	0.014002\\
4.8781	0.01425\\
4.8743	0.014501\\
4.8704	0.014756\\
4.8666	0.015014\\
4.8627	0.015276\\
4.8588	0.015542\\
4.855	0.015812\\
4.8511	0.016085\\
4.8473	0.016362\\
4.8434	0.016642\\
4.8396	0.016927\\
4.8357	0.017215\\
4.8318	0.017507\\
4.828	0.017804\\
4.8241	0.018104\\
4.8203	0.018408\\
4.8164	0.018716\\
4.8126	0.019029\\
4.8087	0.019345\\
4.8048	0.019666\\
4.801	0.01999\\
4.7971	0.020319\\
4.7933	0.020652\\
4.7894	0.02099\\
4.7856	0.021331\\
4.7817	0.021677\\
4.7778	0.022028\\
4.774	0.022383\\
4.7701	0.022742\\
4.7663	0.023106\\
4.7624	0.023474\\
4.7586	0.023847\\
4.7547	0.024224\\
4.7508	0.024606\\
4.747	0.024993\\
4.7431	0.025384\\
4.7393	0.02578\\
4.7354	0.026181\\
4.7316	0.026586\\
4.7277	0.026997\\
4.7238	0.027412\\
4.72	0.027832\\
4.7161	0.028257\\
4.7123	0.028687\\
4.7084	0.029122\\
4.7046	0.029562\\
4.7007	0.030007\\
4.6968	0.030458\\
4.693	0.030913\\
4.6891	0.031374\\
4.6853	0.031839\\
4.6814	0.03231\\
4.6776	0.032786\\
4.6737	0.033268\\
4.6698	0.033755\\
4.666	0.034247\\
4.6621	0.034745\\
4.6583	0.035248\\
4.6544	0.035756\\
4.6506	0.03627\\
4.6467	0.03679\\
4.6428	0.037315\\
4.639	0.037846\\
4.6351	0.038383\\
4.6313	0.038925\\
4.6274	0.039473\\
4.6236	0.040026\\
4.6197	0.040586\\
4.6158	0.041151\\
4.612	0.041722\\
4.6081	0.042299\\
4.6043	0.042882\\
4.6004	0.043471\\
4.5966	0.044066\\
4.5927	0.044667\\
4.5888	0.045274\\
4.585	0.045887\\
4.5811	0.046506\\
4.5773	0.047131\\
4.5734	0.047763\\
4.5696	0.048401\\
4.5657	0.049045\\
4.5618	0.049695\\
4.558	0.050351\\
4.5541	0.051014\\
4.5503	0.051684\\
4.5464	0.05236\\
4.5426	0.053042\\
4.5387	0.05373\\
4.5348	0.054426\\
4.531	0.055127\\
4.5271	0.055836\\
4.5233	0.056551\\
4.5194	0.057272\\
4.5156	0.058001\\
4.5117	0.058736\\
4.5078	0.059477\\
4.504	0.060226\\
4.5001	0.060981\\
4.4963	0.061743\\
4.4924	0.062512\\
4.4886	0.063288\\
4.4847	0.064071\\
4.4808	0.06486\\
4.477	0.065657\\
4.4731	0.066461\\
4.4693	0.067272\\
4.4654	0.068089\\
4.4616	0.068914\\
4.4577	0.069747\\
4.4538	0.070586\\
4.45	0.071432\\
4.4461	0.072286\\
4.4423	0.073147\\
4.4384	0.074015\\
4.4346	0.074891\\
4.4307	0.075774\\
4.4268	0.076665\\
4.423	0.077562\\
4.4191	0.078468\\
4.4153	0.079381\\
4.4114	0.080301\\
4.4076	0.081229\\
4.4037	0.082165\\
4.3998	0.083108\\
4.396	0.084059\\
4.3921	0.085018\\
4.3883	0.085985\\
4.3844	0.086959\\
4.3806	0.087942\\
4.3767	0.088932\\
4.3728	0.089931\\
4.369	0.090937\\
4.3651	0.091952\\
4.3613	0.092975\\
4.3574	0.094006\\
4.3536	0.095046\\
4.3497	0.096094\\
4.3458	0.09715\\
4.342	0.098215\\
4.3381	0.099288\\
4.3343	0.10037\\
4.3304	0.10146\\
4.3266	0.10256\\
4.3227	0.10367\\
4.3188	0.10479\\
4.315	0.10592\\
4.3111	0.10705\\
4.3073	0.1082\\
4.3034	0.10935\\
4.2996	0.11052\\
4.2957	0.11169\\
4.2918	0.11288\\
4.288	0.11407\\
4.2841	0.11528\\
4.2803	0.11649\\
4.2764	0.11772\\
4.2726	0.11895\\
4.2687	0.1202\\
4.2648	0.12146\\
4.261	0.12273\\
4.2571	0.124\\
4.2533	0.12529\\
4.2494	0.1266\\
4.2456	0.12791\\
4.2417	0.12923\\
4.2378	0.13057\\
4.234	0.13192\\
4.2301	0.13328\\
4.2263	0.13465\\
4.2224	0.13604\\
4.2186	0.13744\\
4.2147	0.13885\\
4.2108	0.14027\\
4.207	0.14171\\
4.2031	0.14316\\
4.1993	0.14462\\
4.1954	0.1461\\
4.1916	0.14759\\
4.1877	0.1491\\
4.1838	0.15062\\
4.18	0.15216\\
4.1761	0.15371\\
4.1723	0.15527\\
4.1684	0.15685\\
4.1646	0.15844\\
4.1607	0.16005\\
4.1568	0.16168\\
4.153	0.16332\\
4.1491	0.16498\\
4.1453	0.16666\\
4.1414	0.16835\\
4.1376	0.17005\\
4.1337	0.17178\\
4.1298	0.17352\\
4.126	0.17528\\
4.1221	0.17706\\
4.1183	0.17885\\
4.1144	0.18067\\
4.1106	0.1825\\
4.1067	0.18435\\
4.1028	0.18622\\
4.099	0.18811\\
4.0951	0.19001\\
4.0913	0.19194\\
4.0874	0.19389\\
4.0836	0.19585\\
4.0797	0.19784\\
4.0758	0.19985\\
4.072	0.20187\\
4.0681	0.20392\\
4.0643	0.20599\\
4.0604	0.20808\\
4.0566	0.21019\\
4.0527	0.21233\\
4.0488	0.21448\\
4.045	0.21666\\
4.0411	0.21886\\
4.0373	0.22108\\
4.0334	0.22333\\
4.0296	0.2256\\
4.0257	0.22789\\
4.0218	0.23021\\
4.018	0.23255\\
4.0141	0.23491\\
4.0103	0.2373\\
4.0064	0.23971\\
4.0026	0.24215\\
3.9987	0.24461\\
3.9948	0.2471\\
3.991	0.24961\\
3.9871	0.25214\\
3.9833	0.25471\\
3.9794	0.2573\\
3.9756	0.25991\\
3.9717	0.26255\\
3.9678	0.26522\\
3.964	0.26792\\
3.9601	0.27064\\
3.9563	0.27338\\
3.9524	0.27616\\
3.9486	0.27896\\
3.9447	0.28179\\
3.9408	0.28465\\
3.937	0.28753\\
3.9331	0.29045\\
3.9293	0.29339\\
3.9254	0.29636\\
3.9216	0.29936\\
3.9177	0.30238\\
3.9138	0.30544\\
3.91	0.30852\\
3.9061	0.31163\\
3.9023	0.31478\\
3.8984	0.31795\\
3.8946	0.32115\\
3.8907	0.32438\\
3.8868	0.32763\\
3.883	0.33092\\
3.8791	0.33424\\
3.8753	0.33758\\
3.8714	0.34096\\
3.8676	0.34437\\
3.8637	0.3478\\
3.8598	0.35127\\
3.856	0.35476\\
3.8521	0.35828\\
3.8483	0.36184\\
3.8444	0.36542\\
3.8406	0.36903\\
3.8367	0.37268\\
3.8328	0.37635\\
3.829	0.38005\\
3.8251	0.38378\\
3.8213	0.38754\\
3.8174	0.39133\\
3.8136	0.39515\\
3.8097	0.399\\
3.8058	0.40288\\
3.802	0.40678\\
3.7981	0.41072\\
3.7943	0.41468\\
3.7904	0.41867\\
3.7866	0.42269\\
3.7827	0.42674\\
3.7788	0.43082\\
3.775	0.43493\\
3.7711	0.43906\\
3.7673	0.44322\\
3.7634	0.4474\\
3.7596	0.45162\\
3.7557	0.45586\\
3.7518	0.46012\\
3.748	0.46442\\
3.7441	0.46874\\
3.7403	0.47308\\
3.7364	0.47745\\
3.7326	0.48184\\
3.7287	0.48626\\
3.7248	0.4907\\
3.721	0.49517\\
3.7171	0.49966\\
3.7133	0.50417\\
3.7094	0.50871\\
3.7056	0.51326\\
3.7017	0.51784\\
3.6978	0.52244\\
3.694	0.52706\\
3.6901	0.5317\\
3.6863	0.53636\\
3.6824	0.54103\\
3.6786	0.54573\\
3.6747	0.55044\\
3.6708	0.55517\\
3.667	0.55992\\
3.6631	0.56469\\
3.6593	0.56946\\
3.6554	0.57426\\
3.6516	0.57906\\
3.6477	0.58388\\
3.6438	0.58872\\
3.64	0.59356\\
3.6361	0.59842\\
3.6323	0.60328\\
3.6284	0.60816\\
3.6246	0.61304\\
3.6207	0.61794\\
3.6168	0.62283\\
3.613	0.62774\\
3.6091	0.63265\\
3.6053	0.63756\\
3.6014	0.64248\\
3.5976	0.6474\\
3.5937	0.65233\\
3.5898	0.65725\\
3.586	0.66217\\
3.5821	0.66709\\
3.5783	0.67201\\
3.5744	0.67693\\
3.5706	0.68184\\
3.5667	0.68675\\
3.5628	0.69165\\
3.559	0.69654\\
3.5551	0.70143\\
3.5513	0.7063\\
3.5474	0.71117\\
3.5436	0.71602\\
3.5397	0.72086\\
3.5358	0.72569\\
3.532	0.7305\\
3.5281	0.73529\\
3.5243	0.74007\\
3.5204	0.74482\\
3.5166	0.74956\\
3.5127	0.75428\\
3.5088	0.75897\\
3.505	0.76364\\
3.5011	0.76828\\
3.4973	0.7729\\
3.4934	0.77749\\
3.4896	0.78205\\
3.4857	0.78658\\
3.4818	0.79107\\
3.478	0.79553\\
3.4741	0.79996\\
3.4703	0.80436\\
3.4664	0.80871\\
3.4626	0.81302\\
3.4587	0.8173\\
3.4548	0.82153\\
3.451	0.82572\\
3.4471	0.82986\\
3.4433	0.83396\\
3.4394	0.83801\\
3.4355	0.84201\\
3.4317	0.84595\\
3.4278	0.84985\\
3.424	0.85369\\
3.4201	0.85747\\
3.4163	0.86119\\
3.4124	0.86486\\
3.4085	0.86846\\
3.4047	0.87201\\
3.4008	0.87548\\
3.397	0.8789\\
3.3931	0.88224\\
3.3893	0.88551\\
3.3854	0.88872\\
3.3815	0.89185\\
3.3777	0.89491\\
3.3738	0.89789\\
3.37	0.9008\\
3.3661	0.90363\\
3.3623	0.90638\\
3.3584	0.90904\\
3.3545	0.91163\\
3.3507	0.91413\\
3.3468	0.91655\\
3.343	0.91888\\
3.3391	0.92113\\
3.3353	0.92329\\
3.3314	0.92536\\
3.3275	0.92734\\
3.3237	0.92923\\
3.3198	0.93102\\
3.316	0.93273\\
3.3121	0.93434\\
3.3083	0.93586\\
3.3044	0.93728\\
3.3005	0.93861\\
3.2967	0.93985\\
3.2928	0.94098\\
3.289	0.94203\\
3.2851	0.94297\\
3.2813	0.94382\\
3.2774	0.94457\\
3.2735	0.94522\\
3.2697	0.94578\\
3.2658	0.94623\\
3.262	0.94659\\
3.2581	0.94685\\
3.2543	0.94701\\
3.2504	0.94707\\
3.2465	0.94703\\
3.2427	0.94689\\
3.2388	0.94665\\
3.235	0.94631\\
3.2311	0.94587\\
3.2273	0.94534\\
3.2234	0.9447\\
3.2195	0.94397\\
3.2157	0.94313\\
3.2118	0.9422\\
3.208	0.94116\\
3.2041	0.94003\\
3.2003	0.9388\\
3.1964	0.93747\\
3.1925	0.93604\\
3.1887	0.93452\\
3.1848	0.93289\\
3.181	0.93117\\
3.1771	0.92936\\
3.1733	0.92744\\
3.1694	0.92543\\
3.1655	0.92333\\
3.1617	0.92112\\
3.1578	0.91883\\
3.154	0.91644\\
3.1501	0.91395\\
3.1463	0.91137\\
3.1424	0.9087\\
3.1385	0.90594\\
3.1347	0.90309\\
3.1308	0.90014\\
3.127	0.89711\\
3.1231	0.89398\\
3.1193	0.89077\\
3.1154	0.88747\\
3.1115	0.88408\\
3.1077	0.88061\\
3.1038	0.87705\\
3.1	0.87341\\
3.0961	0.86968\\
3.0923	0.86588\\
3.0884	0.86199\\
3.0845	0.85802\\
3.0807	0.85397\\
3.0768	0.84985\\
3.073	0.84564\\
3.0691	0.84137\\
3.0653	0.83702\\
3.0614	0.8326\\
3.0575	0.8281\\
3.0537	0.82354\\
3.0498	0.81891\\
3.046	0.81421\\
3.0421	0.80944\\
3.0383	0.80462\\
3.0344	0.79973\\
3.0305	0.79478\\
3.0267	0.78977\\
3.0228	0.7847\\
3.019	0.77958\\
3.0151	0.77441\\
3.0113	0.76918\\
3.0074	0.7639\\
3.0035	0.75858\\
2.9997	0.75321\\
2.9958	0.74779\\
2.992	0.74233\\
2.9881	0.73683\\
2.9843	0.73129\\
2.9804	0.72572\\
2.9765	0.7201\\
2.9727	0.71445\\
2.9688	0.70877\\
2.965	0.70306\\
2.9611	0.69732\\
2.9573	0.69156\\
2.9534	0.68576\\
2.9495	0.67995\\
2.9457	0.67411\\
2.9418	0.66825\\
2.938	0.66237\\
2.9341	0.65648\\
2.9303	0.65057\\
2.9264	0.64465\\
2.9225	0.63871\\
2.9187	0.63276\\
2.9148	0.6268\\
2.911	0.62084\\
2.9071	0.61487\\
2.9033	0.60889\\
2.8994	0.60291\\
2.8955	0.59692\\
2.8917	0.59094\\
2.8878	0.58496\\
2.884	0.57897\\
2.8801	0.573\\
2.8763	0.56702\\
2.8724	0.56105\\
2.8685	0.55509\\
2.8647	0.54914\\
2.8608	0.54319\\
2.857	0.53726\\
2.8531	0.53134\\
2.8493	0.52543\\
2.8454	0.51953\\
2.8415	0.51365\\
2.8377	0.50779\\
2.8338	0.50194\\
2.83	0.49611\\
2.8261	0.4903\\
2.8223	0.48451\\
2.8184	0.47874\\
2.8145	0.47299\\
2.8107	0.46727\\
2.8068	0.46157\\
2.803	0.4559\\
2.7991	0.45025\\
2.7953	0.44463\\
2.7914	0.43903\\
2.7875	0.43346\\
2.7837	0.42793\\
2.7798	0.42242\\
2.776	0.41694\\
2.7721	0.4115\\
2.7683	0.40608\\
2.7644	0.4007\\
2.7605	0.39535\\
2.7567	0.39004\\
2.7528	0.38476\\
2.749	0.37952\\
2.7451	0.37431\\
2.7413	0.36914\\
2.7374	0.36401\\
2.7335	0.35891\\
2.7297	0.35385\\
2.7258	0.34883\\
2.722	0.34385\\
2.7181	0.33891\\
2.7143	0.33401\\
2.7104	0.32915\\
2.7065	0.32432\\
2.7027	0.31955\\
2.6988	0.31481\\
2.695	0.31011\\
2.6911	0.30546\\
2.6873	0.30085\\
2.6834	0.29628\\
2.6795	0.29175\\
2.6757	0.28727\\
2.6718	0.28283\\
2.668	0.27843\\
2.6641	0.27408\\
2.6603	0.26977\\
2.6564	0.26551\\
2.6525	0.26129\\
2.6487	0.25712\\
2.6448	0.25299\\
2.641	0.2489\\
2.6371	0.24486\\
2.6333	0.24086\\
2.6294	0.23691\\
2.6255	0.23301\\
2.6217	0.22914\\
2.6178	0.22533\\
2.614	0.22155\\
2.6101	0.21783\\
2.6063	0.21414\\
2.6024	0.21051\\
2.5985	0.20691\\
2.5947	0.20336\\
2.5908	0.19986\\
2.587	0.1964\\
2.5831	0.19298\\
2.5793	0.18961\\
2.5754	0.18628\\
2.5715	0.183\\
2.5677	0.17976\\
2.5638	0.17656\\
2.56	0.1734\\
2.5561	0.17029\\
2.5523	0.16722\\
2.5484	0.1642\\
2.5445	0.16121\\
2.5407	0.15827\\
2.5368	0.15537\\
2.533	0.15251\\
2.5291	0.14969\\
2.5253	0.14691\\
2.5214	0.14418\\
2.5175	0.14148\\
2.5137	0.13883\\
2.5098	0.13621\\
2.506	0.13363\\
2.5021	0.13109\\
2.4983	0.1286\\
2.4944	0.12614\\
2.4905	0.12371\\
2.4867	0.12133\\
2.4828	0.11898\\
2.479	0.11667\\
2.4751	0.1144\\
2.4713	0.11216\\
2.4674	0.10996\\
2.4635	0.1078\\
2.4597	0.10567\\
2.4558	0.10358\\
2.452	0.10152\\
2.4481	0.099493\\
2.4443	0.097501\\
2.4404	0.095544\\
2.4365	0.093619\\
2.4327	0.091728\\
2.4288	0.089869\\
2.425	0.088041\\
2.4211	0.086246\\
2.4173	0.084482\\
2.4134	0.082749\\
2.4095	0.081046\\
2.4057	0.079373\\
2.4018	0.07773\\
2.398	0.076117\\
2.3941	0.074532\\
2.3903	0.072976\\
2.3864	0.071449\\
2.3825	0.069949\\
2.3787	0.068477\\
2.3748	0.067032\\
2.371	0.065613\\
2.3671	0.064221\\
2.3633	0.062855\\
2.3594	0.061515\\
2.3555	0.0602\\
2.3517	0.05891\\
2.3478	0.057644\\
2.344	0.056403\\
2.3401	0.055185\\
2.3363	0.053991\\
2.3324	0.05282\\
2.3285	0.051671\\
2.3247	0.050545\\
2.3208	0.049441\\
2.317	0.048359\\
2.3131	0.047298\\
2.3093	0.046258\\
2.3054	0.045239\\
2.3015	0.04424\\
2.2977	0.043262\\
2.2938	0.042303\\
2.29	0.041363\\
2.2861	0.040442\\
2.2823	0.03954\\
2.2784	0.038657\\
2.2745	0.037791\\
2.2707	0.036943\\
2.2668	0.036113\\
2.263	0.0353\\
2.2591	0.034504\\
2.2553	0.033724\\
2.2514	0.032961\\
2.2475	0.032214\\
2.2437	0.031482\\
2.2398	0.030766\\
2.236	0.030065\\
2.2321	0.029379\\
2.2283	0.028707\\
2.2244	0.02805\\
2.2205	0.027407\\
2.2167	0.026777\\
2.2128	0.026162\\
2.209	0.025559\\
2.2051	0.02497\\
2.2013	0.024393\\
2.1974	0.023829\\
2.1935	0.023277\\
2.1897	0.022737\\
2.1858	0.022209\\
2.182	0.021693\\
2.1781	0.021188\\
2.1743	0.020694\\
2.1704	0.020211\\
2.1665	0.019739\\
2.1627	0.019277\\
2.1588	0.018826\\
2.155	0.018385\\
2.1511	0.017953\\
2.1473	0.017532\\
2.1434	0.01712\\
2.1395	0.016717\\
2.1357	0.016323\\
2.1318	0.015938\\
2.128	0.015562\\
2.1241	0.015194\\
2.1203	0.014835\\
2.1164	0.014484\\
2.1125	0.014141\\
2.1087	0.013805\\
2.1048	0.013478\\
2.101	0.013158\\
2.0971	0.012845\\
2.0933	0.01254\\
2.0894	0.012241\\
2.0855	0.01195\\
2.0817	0.011665\\
2.0778	0.011387\\
2.074	0.011115\\
2.0701	0.01085\\
2.0663	0.010591\\
2.0624	0.010338\\
2.0585	0.010091\\
2.0547	0.0098491\\
2.0508	0.0096133\\
2.047	0.009383\\
2.0431	0.0091582\\
2.0393	0.0089386\\
2.0354	0.0087243\\
2.0315	0.0085149\\
2.0277	0.0083105\\
2.0238	0.0081109\\
2.02	0.0079161\\
2.0161	0.0077259\\
2.0123	0.0075401\\
2.0084	0.0073588\\
2.0045	0.0071818\\
2.0007	0.0070089\\
1.9968	0.0068402\\
1.993	0.0066755\\
1.9891	0.0065147\\
1.9853	0.0063578\\
1.9814	0.0062046\\
1.9775	0.0060551\\
1.9737	0.0059091\\
1.9698	0.0057666\\
1.966	0.0056275\\
1.9621	0.0054918\\
1.9583	0.0053593\\
1.9544	0.00523\\
1.9505	0.0051038\\
1.9467	0.0049806\\
1.9428	0.0048604\\
1.939	0.004743\\
1.9351	0.0046285\\
1.9313	0.0045167\\
1.9274	0.0044077\\
1.9235	0.0043012\\
1.9197	0.0041973\\
1.9158	0.0040959\\
1.912	0.003997\\
1.9081	0.0039004\\
1.9043	0.0038062\\
1.9004	0.0037142\\
1.8965	0.0036245\\
1.8927	0.0035369\\
1.8888	0.0034514\\
1.885	0.003368\\
1.8811	0.0032866\\
1.8773	0.0032072\\
1.8734	0.0031297\\
1.8695	0.003054\\
1.8657	0.0029802\\
1.8618	0.0029082\\
1.858	0.0028379\\
1.8541	0.0027693\\
1.8503	0.0027024\\
1.8464	0.002637\\
1.8425	0.0025733\\
1.8387	0.0025111\\
1.8348	0.0024504\\
1.831	0.0023912\\
1.8271	0.0023334\\
1.8233	0.0022769\\
1.8194	0.0022219\\
1.8155	0.0021682\\
1.8117	0.0021158\\
1.8078	0.0020646\\
1.804	0.0020147\\
1.8001	0.001966\\
1.7963	0.0019185\\
1.7924	0.0018721\\
1.7885	0.0018268\\
1.7847	0.0017826\\
1.7808	0.0017395\\
1.777	0.0016974\\
1.7731	0.0016564\\
1.7693	0.0016163\\
1.7654	0.0015772\\
1.7615	0.0015391\\
1.7577	0.0015018\\
1.7538	0.0014655\\
1.75	0.00143\\
1.7461	0.0013954\\
1.7423	0.0013616\\
1.7384	0.0013287\\
1.7345	0.0012965\\
1.7307	0.0012651\\
1.7268	0.0012345\\
1.723	0.0012046\\
1.7191	0.0011754\\
1.7153	0.0011469\\
1.7114	0.0011191\\
1.7075	0.001092\\
1.7037	0.0010655\\
1.6998	0.0010397\\
1.696	0.0010145\\
1.6921	0.00098987\\
1.6883	0.00096586\\
1.6844	0.00094242\\
1.6805	0.00091954\\
1.6767	0.00089721\\
1.6728	0.00087542\\
1.669	0.00085415\\
1.6651	0.00083339\\
1.6613	0.00081313\\
1.6574	0.00079336\\
1.6535	0.00077406\\
1.6497	0.00075522\\
1.6458	0.00073684\\
1.642	0.00071889\\
1.6381	0.00070138\\
1.6343	0.00068429\\
1.6304	0.00066761\\
1.6265	0.00065133\\
1.6227	0.00063543\\
1.6188	0.00061992\\
1.615	0.00060479\\
1.6111	0.00059001\\
1.6073	0.00057559\\
1.6034	0.00056152\\
1.5995	0.00054778\\
1.5957	0.00053437\\
1.5918	0.00052128\\
1.588	0.00050851\\
1.5841	0.00049605\\
1.5803	0.00048388\\
1.5764	0.00047201\\
1.5725	0.00046042\\
1.5687	0.0004491\\
1.5648	0.00043806\\
1.561	0.00042729\\
1.5571	0.00041677\\
1.5533	0.00040651\\
1.5494	0.00039649\\
1.5455	0.00038672\\
1.5417	0.00037717\\
1.5378	0.00036786\\
1.534	0.00035877\\
1.5301	0.0003499\\
1.5263	0.00034125\\
1.5224	0.0003328\\
1.5185	0.00032455\\
1.5147	0.00031651\\
1.5108	0.00030866\\
1.507	0.00030099\\
1.5031	0.00029351\\
1.4993	0.00028622\\
1.4954	0.00027909\\
1.4915	0.00027214\\
1.4877	0.00026536\\
1.4838	0.00025874\\
1.48	0.00025228\\
1.4761	0.00024598\\
1.4723	0.00023983\\
1.4684	0.00023383\\
1.4645	0.00022797\\
1.4607	0.00022225\\
1.4568	0.00021668\\
}--cycle;
\addplot[ybar interval, fill=black, fill opacity=0.35, area legend, draw=none] table[row sep=crcr, x=Lower, y=Count] {%
Lower	Upper	Count\\
1.4568	1.7532	0\\
1.7532	2.0496	0\\
2.0496	2.346	0.071973\\
2.346	2.6425	0.14395\\
2.6425	2.9389	0.51281\\
2.9389	3.2353	0.89066\\
3.2353	3.5317	0.85468\\
3.5317	3.8281	0.51281\\
3.8281	4.1245	0.23391\\
4.1245	4.4209	0.071973\\
4.4209	4.7173	0.05398\\
4.7173	5.0137	0.0089966\\
5.0137	5.3101	0.017993\\
5.3101	5.3101	0.017993\\
};
\addlegendentry{\RealHistogramLegend}

\addplot [color=black]
  table[row sep=crcr]{%
1.4568	7.3389e-05\\
1.669	0.00061309\\
1.7731	0.001575\\
1.8464	0.0029423\\
1.9043	0.0047105\\
1.9544	0.0069684\\
1.9968	0.0095921\\
2.0354	0.012706\\
2.0701	0.016239\\
2.101	0.020075\\
2.1318	0.024674\\
2.1588	0.029417\\
2.1858	0.034919\\
2.2128	0.041268\\
2.236	0.047456\\
2.2591	0.054396\\
2.2823	0.062151\\
2.3054	0.070783\\
2.3285	0.080355\\
2.3478	0.089094\\
2.3671	0.098563\\
2.3864	0.10879\\
2.4057	0.11982\\
2.425	0.13167\\
2.4443	0.14437\\
2.4635	0.15794\\
2.4828	0.17239\\
2.5021	0.18776\\
2.5214	0.20403\\
2.5407	0.22122\\
2.56	0.23933\\
2.5793	0.25834\\
2.6024	0.28232\\
2.6255	0.30754\\
2.6487	0.33393\\
2.6718	0.36143\\
2.6988	0.39478\\
2.7258	0.42932\\
2.7567	0.47\\
2.7991	0.52736\\
2.8917	0.6531\\
2.9225	0.69342\\
2.9495	0.72734\\
2.9727	0.75511\\
2.992	0.77716\\
3.0113	0.79809\\
3.0305	0.81777\\
3.046	0.83254\\
3.0614	0.84638\\
3.0768	0.85923\\
3.0923	0.87105\\
3.1077	0.88178\\
3.1231	0.8914\\
3.1347	0.89785\\
3.1463	0.90363\\
3.1578	0.90873\\
3.1694	0.91314\\
3.181	0.91684\\
3.1925	0.91984\\
3.2041	0.92212\\
3.2157	0.92368\\
3.2273	0.92452\\
3.2388	0.92464\\
3.2504	0.92403\\
3.262	0.9227\\
3.2735	0.92066\\
3.2851	0.9179\\
3.2967	0.91445\\
3.3083	0.9103\\
3.3198	0.90547\\
3.3314	0.89996\\
3.343	0.89381\\
3.3584	0.8846\\
3.3738	0.87431\\
3.3893	0.86296\\
3.4047	0.8506\\
3.4201	0.8373\\
3.4355	0.8231\\
3.4548	0.80418\\
3.4741	0.78408\\
3.4934	0.76291\\
3.5166	0.73628\\
3.5397	0.70853\\
3.5667	0.67503\\
3.6014	0.63072\\
3.721	0.47686\\
3.7518	0.43893\\
3.7788	0.4069\\
3.8058	0.37612\\
3.829	0.35085\\
3.8521	0.32668\\
3.8753	0.30367\\
3.8984	0.28186\\
3.9216	0.26127\\
3.9408	0.24505\\
3.9601	0.22969\\
3.9794	0.21517\\
3.9987	0.20149\\
4.018	0.18863\\
4.0373	0.17656\\
4.0566	0.16527\\
4.0758	0.15473\\
4.099	0.14301\\
4.1221	0.13228\\
4.1453	0.12246\\
4.1684	0.1135\\
4.1916	0.10534\\
4.2147	0.097918\\
4.2417	0.090106\\
4.2687	0.083123\\
4.2957	0.07688\\
4.3266	0.070545\\
4.3574	0.064951\\
4.3921	0.059415\\
4.4307	0.054048\\
4.4731	0.048918\\
4.5233	0.043673\\
4.5811	0.038449\\
4.6467	0.033306\\
4.7238	0.028018\\
4.8126	0.022705\\
4.909	0.01771\\
5.0054	0.013475\\
5.1057	0.0098505\\
5.2098	0.0068842\\
5.3101	0.0047195\\
};
\addlegendentry{\RealBLNFitLegend}

      % \input{../../../../.tex/20/KS/SizeDistributionFittingsfig20.tex}
      \legend{}
  
      \nextgroupplot[%
        xmin=0.85621,xmax=5.4115,
        % xmin=1.4157,xmax=5.517,
        % ymin=0,ymax=1,
      ] % This file was created by matlab2tikz.
%
\definecolor{mycolor1}{rgb}{0.83529,0.36863,0.00000}%
%
\addplot[ybar interval, fill=mycolor1, fill opacity=0.35, area legend, draw=none] table[row sep=crcr, x=Lower, y=Count] {%
Lower	Upper	Count\\
1.6021	1.8889	0.0010228\\
1.8889	2.1757	0.0098556\\
2.1757	2.4625	0.060435\\
2.4625	2.7494	0.2291\\
2.7494	3.0362	0.59031\\
3.0362	3.323	0.91434\\
3.323	3.6098	0.81867\\
3.6098	3.8966	0.46507\\
3.8966	4.1834	0.23003\\
4.1834	4.4702	0.11157\\
4.4702	4.757	0.037191\\
4.757	5.0438	0.0099486\\
5.0438	5.3306	0.0091118\\
5.3306	5.3306	0.0091118\\
};
\addlegendentry{\ApproxHistogramLegend}

\addplot [color=mycolor1, only marks, every error bar/.append style={opacity=0.45}, mark=*, mark size=0pt, draw=none, forget plot]
 plot [error bars/.cd, y dir=both, y explicit, error bar style={line width=1pt, color=mycolor1}, error mark options={mark=none,mark size=0pt}]
 table[row sep=crcr, y error plus index=2, y error minus index=3]{%
1.7455	0.0010228	0.00058015	0.00058015\\
2.0323	0.0098556	0.0018876	0.0018876\\
2.3191	0.060435	0.0045152	0.0045152\\
2.606	0.2291	0.0085071	0.0085071\\
2.8928	0.59031	0.010329	0.010329\\
3.1796	0.91434	0.01299	0.01299\\
3.4664	0.81867	0.015655	0.015655\\
3.7532	0.46507	0.01029	0.01029\\
4.04	0.23003	0.0075845	0.0075845\\
4.3268	0.11157	0.0049267	0.0049267\\
4.6136	0.037191	0.0028242	0.0028242\\
4.9004	0.0099486	0.0015121	0.0015121\\
5.1872	0.0091118	0.00036898	0.00036898\\
};
\addplot [color=mycolor1]
  table[row sep=crcr]{%
1.6021	0.0006735\\
1.744	0.0016721\\
1.841	0.0030706\\
1.9156	0.0048694\\
1.9753	0.0070144\\
2.0276	0.0096249\\
2.0724	0.012592\\
2.1134	0.016072\\
2.1508	0.02002\\
2.1844	0.024348\\
2.2142	0.028922\\
2.2441	0.034291\\
2.2702	0.039735\\
2.2963	0.045965\\
2.3224	0.053076\\
2.3448	0.059947\\
2.3672	0.067605\\
2.3896	0.076119\\
2.412	0.085561\\
2.4307	0.09419\\
2.4493	0.10356\\
2.468	0.1137\\
2.4867	0.12467\\
2.5053	0.13649\\
2.524	0.14922\\
2.5426	0.16289\\
2.5613	0.17752\\
2.58	0.19315\\
2.5986	0.20981\\
2.6173	0.22751\\
2.636	0.24627\\
2.6546	0.26608\\
2.6733	0.28695\\
2.6919	0.30887\\
2.7106	0.3318\\
2.733	0.36063\\
2.7554	0.39079\\
2.7778	0.4222\\
2.8039	0.46022\\
2.8338	0.50519\\
2.8711	0.56299\\
2.9644	0.70842\\
2.9905	0.74751\\
3.0129	0.77977\\
3.0316	0.8055\\
3.0502	0.82989\\
3.0652	0.84825\\
3.0801	0.86542\\
3.095	0.88126\\
3.1062	0.89219\\
3.1174	0.90227\\
3.1286	0.91147\\
3.1398	0.91977\\
3.151	0.92716\\
3.1622	0.93366\\
3.1734	0.93926\\
3.1846	0.944\\
3.1958	0.94788\\
3.207	0.95092\\
3.2182	0.95314\\
3.2294	0.95456\\
3.2406	0.95517\\
3.2518	0.955\\
3.263	0.95404\\
3.2742	0.95231\\
3.2854	0.94981\\
3.2966	0.94655\\
3.3078	0.94253\\
3.319	0.93778\\
3.3301	0.9323\\
3.3413	0.92611\\
3.3525	0.91922\\
3.3675	0.909\\
3.3824	0.89764\\
3.3973	0.8852\\
3.4123	0.87174\\
3.4272	0.85733\\
3.4458	0.83811\\
3.4645	0.81769\\
3.4869	0.79185\\
3.5093	0.76481\\
3.5392	0.72745\\
3.5877	0.66517\\
3.6399	0.5984\\
3.6735	0.55673\\
3.7034	0.52094\\
3.7295	0.49076\\
3.7556	0.46173\\
3.7817	0.43394\\
3.8041	0.41113\\
3.8265	0.38927\\
3.8489	0.36835\\
3.8713	0.34839\\
3.8937	0.32935\\
3.9161	0.31124\\
3.9385	0.29402\\
3.9609	0.27766\\
3.9833	0.26213\\
4.0094	0.24503\\
4.0355	0.22896\\
4.0617	0.21386\\
4.0878	0.19968\\
4.1139	0.18635\\
4.14	0.17384\\
4.1699	0.16045\\
4.1998	0.14797\\
4.2296	0.13634\\
4.2595	0.12548\\
4.2931	0.11413\\
4.3267	0.10363\\
4.3602	0.093912\\
4.3976	0.083964\\
4.4349	0.074849\\
4.4722	0.066511\\
4.5133	0.058174\\
4.5543	0.050655\\
4.5954	0.043899\\
4.6402	0.03734\\
4.6849	0.031565\\
4.7335	0.026125\\
4.782	0.021457\\
4.8342	0.017207\\
4.8902	0.013447\\
4.9499	0.010218\\
5.0171	0.0073975\\
5.0918	0.005078\\
5.1776	0.003223\\
5.2784	0.0018364\\
5.3306	0.0013557\\
};
\addlegendentry{\ApproxBLNFitLegend}


\addplot[area legend, draw=none, fill=mycolor1, fill opacity=0.15, forget plot]
table[row sep=crcr] {%
x	y\\
1.6021	0.00080938\\
1.6059	0.00082884\\
1.6096	0.00084874\\
1.6133	0.0008691\\
1.6171	0.00088993\\
1.6208	0.00091123\\
1.6245	0.00093302\\
1.6283	0.00095531\\
1.632	0.0009781\\
1.6357	0.0010014\\
1.6394	0.0010253\\
1.6432	0.0010496\\
1.6469	0.0010746\\
1.6506	0.0011001\\
1.6544	0.0011262\\
1.6581	0.0011528\\
1.6618	0.0011801\\
1.6656	0.001208\\
1.6693	0.0012365\\
1.673	0.0012657\\
1.6768	0.0012955\\
1.6805	0.001326\\
1.6842	0.0013572\\
1.688	0.0013891\\
1.6917	0.0014217\\
1.6954	0.001455\\
1.6992	0.0014891\\
1.7029	0.001524\\
1.7066	0.0015596\\
1.7104	0.001596\\
1.7141	0.0016333\\
1.7178	0.0016713\\
1.7216	0.0017103\\
1.7253	0.0017501\\
1.729	0.0017907\\
1.7328	0.0018323\\
1.7365	0.0018749\\
1.7402	0.0019183\\
1.744	0.0019628\\
1.7477	0.0020082\\
1.7514	0.0020546\\
1.7551	0.0021021\\
1.7589	0.0021506\\
1.7626	0.0022002\\
1.7663	0.0022509\\
1.7701	0.0023027\\
1.7738	0.0023557\\
1.7775	0.0024098\\
1.7813	0.0024651\\
1.785	0.0025217\\
1.7887	0.0025795\\
1.7925	0.0026386\\
1.7962	0.002699\\
1.7999	0.0027607\\
1.8037	0.0028237\\
1.8074	0.0028882\\
1.8111	0.0029541\\
1.8149	0.0030214\\
1.8186	0.0030902\\
1.8223	0.0031606\\
1.8261	0.0032324\\
1.8298	0.0033059\\
1.8335	0.0033809\\
1.8373	0.0034576\\
1.841	0.003536\\
1.8447	0.0036161\\
1.8485	0.003698\\
1.8522	0.0037816\\
1.8559	0.0038671\\
1.8596	0.0039544\\
1.8634	0.0040436\\
1.8671	0.0041348\\
1.8708	0.004228\\
1.8746	0.0043232\\
1.8783	0.0044205\\
1.882	0.0045199\\
1.8858	0.0046214\\
1.8895	0.0047252\\
1.8932	0.0048312\\
1.897	0.0049396\\
1.9007	0.0050502\\
1.9044	0.0051633\\
1.9082	0.0052788\\
1.9119	0.0053969\\
1.9156	0.0055175\\
1.9194	0.0056407\\
1.9231	0.0057665\\
1.9268	0.0058951\\
1.9306	0.0060265\\
1.9343	0.0061606\\
1.938	0.0062977\\
1.9418	0.0064378\\
1.9455	0.0065808\\
1.9492	0.006727\\
1.953	0.0068763\\
1.9567	0.0070288\\
1.9604	0.0071845\\
1.9642	0.0073437\\
1.9679	0.0075062\\
1.9716	0.0076722\\
1.9753	0.0078418\\
1.9791	0.008015\\
1.9828	0.0081919\\
1.9865	0.0083726\\
1.9903	0.0085572\\
1.994	0.0087457\\
1.9977	0.0089382\\
2.0015	0.0091348\\
2.0052	0.0093356\\
2.0089	0.0095407\\
2.0127	0.0097502\\
2.0164	0.0099641\\
2.0201	0.010183\\
2.0239	0.010406\\
2.0276	0.010633\\
2.0313	0.010866\\
2.0351	0.011104\\
2.0388	0.011346\\
2.0425	0.011594\\
2.0463	0.011847\\
2.05	0.012105\\
2.0537	0.012369\\
2.0575	0.012638\\
2.0612	0.012913\\
2.0649	0.013194\\
2.0687	0.01348\\
2.0724	0.013773\\
2.0761	0.014072\\
2.0799	0.014377\\
2.0836	0.014688\\
2.0873	0.015006\\
2.091	0.01533\\
2.0948	0.015661\\
2.0985	0.016\\
2.1022	0.016345\\
2.106	0.016697\\
2.1097	0.017057\\
2.1134	0.017424\\
2.1172	0.017798\\
2.1209	0.018181\\
2.1246	0.018571\\
2.1284	0.018969\\
2.1321	0.019376\\
2.1358	0.019791\\
2.1396	0.020214\\
2.1433	0.020646\\
2.147	0.021087\\
2.1508	0.021537\\
2.1545	0.021996\\
2.1582	0.022464\\
2.162	0.022942\\
2.1657	0.02343\\
2.1694	0.023927\\
2.1732	0.024435\\
2.1769	0.024953\\
2.1806	0.025481\\
2.1844	0.02602\\
2.1881	0.026569\\
2.1918	0.02713\\
2.1956	0.027702\\
2.1993	0.028286\\
2.203	0.028881\\
2.2067	0.029488\\
2.2105	0.030107\\
2.2142	0.030738\\
2.2179	0.031382\\
2.2217	0.032038\\
2.2254	0.032708\\
2.2291	0.03339\\
2.2329	0.034086\\
2.2366	0.034796\\
2.2403	0.03552\\
2.2441	0.036257\\
2.2478	0.037009\\
2.2515	0.037776\\
2.2553	0.038558\\
2.259	0.039354\\
2.2627	0.040166\\
2.2665	0.040994\\
2.2702	0.041838\\
2.2739	0.042697\\
2.2777	0.043574\\
2.2814	0.044466\\
2.2851	0.045376\\
2.2889	0.046303\\
2.2926	0.047248\\
2.2963	0.048211\\
2.3001	0.049191\\
2.3038	0.05019\\
2.3075	0.051208\\
2.3112	0.052244\\
2.315	0.0533\\
2.3187	0.054376\\
2.3224	0.055471\\
2.3262	0.056586\\
2.3299	0.057722\\
2.3336	0.058879\\
2.3374	0.060057\\
2.3411	0.061256\\
2.3448	0.062477\\
2.3486	0.063721\\
2.3523	0.064986\\
2.356	0.066274\\
2.3598	0.067586\\
2.3635	0.06892\\
2.3672	0.070279\\
2.371	0.071661\\
2.3747	0.073068\\
2.3784	0.074499\\
2.3822	0.075956\\
2.3859	0.077437\\
2.3896	0.078945\\
2.3934	0.080478\\
2.3971	0.082038\\
2.4008	0.083625\\
2.4046	0.085239\\
2.4083	0.08688\\
2.412	0.088549\\
2.4158	0.090246\\
2.4195	0.091971\\
2.4232	0.093725\\
2.4269	0.095509\\
2.4307	0.097322\\
2.4344	0.099164\\
2.4381	0.10104\\
2.4419	0.10294\\
2.4456	0.10487\\
2.4493	0.10684\\
2.4531	0.10884\\
2.4568	0.11087\\
2.4605	0.11293\\
2.4643	0.11502\\
2.468	0.11715\\
2.4717	0.11931\\
2.4755	0.1215\\
2.4792	0.12373\\
2.4829	0.12599\\
2.4867	0.12828\\
2.4904	0.13061\\
2.4941	0.13298\\
2.4979	0.13538\\
2.5016	0.13782\\
2.5053	0.14029\\
2.5091	0.1428\\
2.5128	0.14535\\
2.5165	0.14793\\
2.5203	0.15055\\
2.524	0.15321\\
2.5277	0.1559\\
2.5315	0.15864\\
2.5352	0.16141\\
2.5389	0.16422\\
2.5426	0.16707\\
2.5464	0.16996\\
2.5501	0.17289\\
2.5538	0.17585\\
2.5576	0.17886\\
2.5613	0.18191\\
2.565	0.185\\
2.5688	0.18812\\
2.5725	0.19129\\
2.5762	0.1945\\
2.58	0.19775\\
2.5837	0.20104\\
2.5874	0.20438\\
2.5912	0.20775\\
2.5949	0.21117\\
2.5986	0.21462\\
2.6024	0.21812\\
2.6061	0.22167\\
2.6098	0.22525\\
2.6136	0.22887\\
2.6173	0.23254\\
2.621	0.23625\\
2.6248	0.24\\
2.6285	0.2438\\
2.6322	0.24764\\
2.636	0.25151\\
2.6397	0.25544\\
2.6434	0.2594\\
2.6472	0.26341\\
2.6509	0.26746\\
2.6546	0.27155\\
2.6583	0.27568\\
2.6621	0.27986\\
2.6658	0.28407\\
2.6695	0.28833\\
2.6733	0.29263\\
2.677	0.29698\\
2.6807	0.30136\\
2.6845	0.30579\\
2.6882	0.31025\\
2.6919	0.31476\\
2.6957	0.31931\\
2.6994	0.3239\\
2.7031	0.32853\\
2.7069	0.3332\\
2.7106	0.3379\\
2.7143	0.34265\\
2.7181	0.34744\\
2.7218	0.35226\\
2.7255	0.35713\\
2.7293	0.36203\\
2.733	0.36697\\
2.7367	0.37195\\
2.7405	0.37696\\
2.7442	0.38201\\
2.7479	0.38709\\
2.7517	0.39221\\
2.7554	0.39737\\
2.7591	0.40256\\
2.7628	0.40778\\
2.7666	0.41304\\
2.7703	0.41833\\
2.774	0.42365\\
2.7778	0.429\\
2.7815	0.43439\\
2.7852	0.4398\\
2.789	0.44524\\
2.7927	0.45071\\
2.7964	0.45621\\
2.8002	0.46174\\
2.8039	0.46729\\
2.8076	0.47287\\
2.8114	0.47847\\
2.8151	0.4841\\
2.8188	0.48975\\
2.8226	0.49542\\
2.8263	0.50111\\
2.83	0.50683\\
2.8338	0.51256\\
2.8375	0.51831\\
2.8412	0.52408\\
2.845	0.52987\\
2.8487	0.53567\\
2.8524	0.54149\\
2.8562	0.54732\\
2.8599	0.55316\\
2.8636	0.55902\\
2.8674	0.56488\\
2.8711	0.57076\\
2.8748	0.57664\\
2.8785	0.58253\\
2.8823	0.58842\\
2.886	0.59433\\
2.8897	0.60023\\
2.8935	0.60614\\
2.8972	0.61204\\
2.9009	0.61795\\
2.9047	0.62386\\
2.9084	0.62976\\
2.9121	0.63566\\
2.9159	0.64156\\
2.9196	0.64744\\
2.9233	0.65333\\
2.9271	0.6592\\
2.9308	0.66506\\
2.9345	0.67091\\
2.9383	0.67675\\
2.942	0.68258\\
2.9457	0.68839\\
2.9495	0.69418\\
2.9532	0.69996\\
2.9569	0.70572\\
2.9607	0.71145\\
2.9644	0.71717\\
2.9681	0.72287\\
2.9719	0.72854\\
2.9756	0.73419\\
2.9793	0.73981\\
2.9831	0.74541\\
2.9868	0.75098\\
2.9905	0.75652\\
2.9942	0.76204\\
2.998	0.76752\\
3.0017	0.77297\\
3.0054	0.77839\\
3.0092	0.78378\\
3.0129	0.78913\\
3.0166	0.79445\\
3.0204	0.79973\\
3.0241	0.80497\\
3.0278	0.81017\\
3.0316	0.81532\\
3.0353	0.82043\\
3.039	0.8255\\
3.0428	0.83051\\
3.0465	0.83547\\
3.0502	0.84037\\
3.054	0.84521\\
3.0577	0.84998\\
3.0614	0.85469\\
3.0652	0.85933\\
3.0689	0.86389\\
3.0726	0.86838\\
3.0764	0.87278\\
3.0801	0.87709\\
3.0838	0.88131\\
3.0876	0.88544\\
3.0913	0.88948\\
3.095	0.89341\\
3.0988	0.89724\\
3.1025	0.90096\\
3.1062	0.90458\\
3.1099	0.90808\\
3.1137	0.91148\\
3.1174	0.91476\\
3.1211	0.91793\\
3.1249	0.92099\\
3.1286	0.92393\\
3.1323	0.92677\\
3.1361	0.92949\\
3.1398	0.93209\\
3.1435	0.93459\\
3.1473	0.93698\\
3.151	0.93926\\
3.1547	0.94143\\
3.1585	0.9435\\
3.1622	0.94546\\
3.1659	0.94733\\
3.1697	0.94909\\
3.1734	0.95076\\
3.1771	0.95233\\
3.1809	0.9538\\
3.1846	0.95518\\
3.1883	0.95647\\
3.1921	0.95767\\
3.1958	0.95877\\
3.1995	0.95979\\
3.2033	0.96072\\
3.207	0.96155\\
3.2107	0.9623\\
3.2144	0.96296\\
3.2182	0.96354\\
3.2219	0.96402\\
3.2256	0.96442\\
3.2294	0.96473\\
3.2331	0.96495\\
3.2368	0.96509\\
3.2406	0.96514\\
3.2443	0.9651\\
3.248	0.96498\\
3.2518	0.96477\\
3.2555	0.96448\\
3.2592	0.9641\\
3.263	0.96363\\
3.2667	0.96309\\
3.2704	0.96246\\
3.2742	0.96175\\
3.2779	0.96096\\
3.2816	0.9601\\
3.2854	0.95915\\
3.2891	0.95812\\
3.2928	0.95702\\
3.2966	0.95584\\
3.3003	0.95458\\
3.304	0.95325\\
3.3078	0.95185\\
3.3115	0.95037\\
3.3152	0.94882\\
3.319	0.9472\\
3.3227	0.9455\\
3.3264	0.94373\\
3.3301	0.94189\\
3.3339	0.93998\\
3.3376	0.938\\
3.3413	0.93594\\
3.3451	0.93382\\
3.3488	0.93162\\
3.3525	0.92935\\
3.3563	0.92701\\
3.36	0.9246\\
3.3637	0.92212\\
3.3675	0.91956\\
3.3712	0.91694\\
3.3749	0.91424\\
3.3787	0.91147\\
3.3824	0.90863\\
3.3861	0.90572\\
3.3899	0.90274\\
3.3936	0.89969\\
3.3973	0.89658\\
3.4011	0.89339\\
3.4048	0.89013\\
3.4085	0.88681\\
3.4123	0.88342\\
3.416	0.87996\\
3.4197	0.87644\\
3.4235	0.87285\\
3.4272	0.8692\\
3.4309	0.86549\\
3.4347	0.86172\\
3.4384	0.85789\\
3.4421	0.85399\\
3.4458	0.85004\\
3.4496	0.84604\\
3.4533	0.84198\\
3.457	0.83787\\
3.4608	0.8337\\
3.4645	0.82949\\
3.4682	0.82522\\
3.472	0.82091\\
3.4757	0.81655\\
3.4794	0.81216\\
3.4832	0.80771\\
3.4869	0.80323\\
3.4906	0.79871\\
3.4944	0.79416\\
3.4981	0.78957\\
3.5018	0.78495\\
3.5056	0.7803\\
3.5093	0.77561\\
3.513	0.77091\\
3.5168	0.76617\\
3.5205	0.76142\\
3.5242	0.75664\\
3.528	0.75184\\
3.5317	0.74703\\
3.5354	0.74219\\
3.5392	0.73735\\
3.5429	0.73249\\
3.5466	0.72762\\
3.5504	0.72274\\
3.5541	0.71786\\
3.5578	0.71296\\
3.5615	0.70807\\
3.5653	0.70317\\
3.569	0.69826\\
3.5727	0.69336\\
3.5765	0.68846\\
3.5802	0.68355\\
3.5839	0.67866\\
3.5877	0.67376\\
3.5914	0.66887\\
3.5951	0.66399\\
3.5989	0.65911\\
3.6026	0.65424\\
3.6063	0.64938\\
3.6101	0.64452\\
3.6138	0.63968\\
3.6175	0.63485\\
3.6213	0.63003\\
3.625	0.62522\\
3.6287	0.62043\\
3.6325	0.61564\\
3.6362	0.61087\\
3.6399	0.60612\\
3.6437	0.60138\\
3.6474	0.59665\\
3.6511	0.59194\\
3.6549	0.58725\\
3.6586	0.58257\\
3.6623	0.57791\\
3.666	0.57327\\
3.6698	0.56864\\
3.6735	0.56403\\
3.6772	0.55944\\
3.681	0.55487\\
3.6847	0.55032\\
3.6884	0.54579\\
3.6922	0.54128\\
3.6959	0.53679\\
3.6996	0.53232\\
3.7034	0.52787\\
3.7071	0.52344\\
3.7108	0.51903\\
3.7146	0.51465\\
3.7183	0.51029\\
3.722	0.50595\\
3.7258	0.50163\\
3.7295	0.49733\\
3.7332	0.49306\\
3.737	0.48882\\
3.7407	0.48459\\
3.7444	0.48039\\
3.7482	0.47622\\
3.7519	0.47207\\
3.7556	0.46794\\
3.7594	0.46384\\
3.7631	0.45977\\
3.7668	0.45572\\
3.7706	0.45169\\
3.7743	0.44769\\
3.778	0.44372\\
3.7817	0.43977\\
3.7855	0.43585\\
3.7892	0.43196\\
3.7929	0.42809\\
3.7967	0.42425\\
3.8004	0.42043\\
3.8041	0.41664\\
3.8079	0.41288\\
3.8116	0.40915\\
3.8153	0.40544\\
3.8191	0.40176\\
3.8228	0.39811\\
3.8265	0.39448\\
3.8303	0.39088\\
3.834	0.38731\\
3.8377	0.38376\\
3.8415	0.38025\\
3.8452	0.37676\\
3.8489	0.37329\\
3.8527	0.36986\\
3.8564	0.36645\\
3.8601	0.36307\\
3.8639	0.35971\\
3.8676	0.35638\\
3.8713	0.35308\\
3.8751	0.34981\\
3.8788	0.34656\\
3.8825	0.34334\\
3.8863	0.34015\\
3.89	0.33698\\
3.8937	0.33384\\
3.8974	0.33072\\
3.9012	0.32764\\
3.9049	0.32457\\
3.9086	0.32154\\
3.9124	0.31853\\
3.9161	0.31554\\
3.9198	0.31258\\
3.9236	0.30965\\
3.9273	0.30674\\
3.931	0.30386\\
3.9348	0.301\\
3.9385	0.29817\\
3.9422	0.29536\\
3.946	0.29257\\
3.9497	0.28981\\
3.9534	0.28708\\
3.9572	0.28437\\
3.9609	0.28168\\
3.9646	0.27901\\
3.9684	0.27637\\
3.9721	0.27376\\
3.9758	0.27116\\
3.9796	0.26859\\
3.9833	0.26604\\
3.987	0.26351\\
3.9908	0.26101\\
3.9945	0.25853\\
3.9982	0.25607\\
4.002	0.25363\\
4.0057	0.25121\\
4.0094	0.24881\\
4.0131	0.24644\\
4.0169	0.24409\\
4.0206	0.24175\\
4.0243	0.23944\\
4.0281	0.23715\\
4.0318	0.23488\\
4.0355	0.23263\\
4.0393	0.23039\\
4.043	0.22818\\
4.0467	0.22599\\
4.0505	0.22382\\
4.0542	0.22166\\
4.0579	0.21953\\
4.0617	0.21741\\
4.0654	0.21531\\
4.0691	0.21323\\
4.0729	0.21117\\
4.0766	0.20913\\
4.0803	0.2071\\
4.0841	0.20509\\
4.0878	0.2031\\
4.0915	0.20113\\
4.0953	0.19917\\
4.099	0.19723\\
4.1027	0.19531\\
4.1065	0.1934\\
4.1102	0.19151\\
4.1139	0.18964\\
4.1176	0.18778\\
4.1214	0.18594\\
4.1251	0.18412\\
4.1288	0.18231\\
4.1326	0.18052\\
4.1363	0.17874\\
4.14	0.17698\\
4.1438	0.17523\\
4.1475	0.17349\\
4.1512	0.17178\\
4.155	0.17007\\
4.1587	0.16838\\
4.1624	0.16671\\
4.1662	0.16505\\
4.1699	0.16341\\
4.1736	0.16177\\
4.1774	0.16016\\
4.1811	0.15855\\
4.1848	0.15696\\
4.1886	0.15538\\
4.1923	0.15382\\
4.196	0.15227\\
4.1998	0.15073\\
4.2035	0.14921\\
4.2072	0.1477\\
4.211	0.1462\\
4.2147	0.14472\\
4.2184	0.14324\\
4.2222	0.14178\\
4.2259	0.14033\\
4.2296	0.1389\\
4.2333	0.13747\\
4.2371	0.13606\\
4.2408	0.13466\\
4.2445	0.13328\\
4.2483	0.1319\\
4.252	0.13053\\
4.2557	0.12918\\
4.2595	0.12784\\
4.2632	0.12651\\
4.2669	0.12519\\
4.2707	0.12388\\
4.2744	0.12259\\
4.2781	0.1213\\
4.2819	0.12003\\
4.2856	0.11876\\
4.2893	0.11751\\
4.2931	0.11627\\
4.2968	0.11504\\
4.3005	0.11381\\
4.3043	0.1126\\
4.308	0.1114\\
4.3117	0.11021\\
4.3155	0.10903\\
4.3192	0.10786\\
4.3229	0.1067\\
4.3267	0.10555\\
4.3304	0.10441\\
4.3341	0.10328\\
4.3379	0.10216\\
4.3416	0.10105\\
4.3453	0.099947\\
4.349	0.098855\\
4.3528	0.097772\\
4.3565	0.096699\\
4.3602	0.095635\\
4.364	0.09458\\
4.3677	0.093534\\
4.3714	0.092498\\
4.3752	0.09147\\
4.3789	0.090452\\
4.3826	0.089442\\
4.3864	0.088441\\
4.3901	0.087449\\
4.3938	0.086466\\
4.3976	0.085491\\
4.4013	0.084525\\
4.405	0.083568\\
4.4088	0.082619\\
4.4125	0.081678\\
4.4162	0.080746\\
4.42	0.079822\\
4.4237	0.078906\\
4.4274	0.077999\\
4.4312	0.0771\\
4.4349	0.076208\\
4.4386	0.075325\\
4.4424	0.07445\\
4.4461	0.073582\\
4.4498	0.072723\\
4.4536	0.071871\\
4.4573	0.071027\\
4.461	0.070191\\
4.4647	0.069363\\
4.4685	0.068542\\
4.4722	0.067728\\
4.4759	0.066922\\
4.4797	0.066124\\
4.4834	0.065333\\
4.4871	0.064549\\
4.4909	0.063773\\
4.4946	0.063004\\
4.4983	0.062242\\
4.5021	0.061487\\
4.5058	0.060739\\
4.5095	0.059998\\
4.5133	0.059265\\
4.517	0.058538\\
4.5207	0.057818\\
4.5245	0.057105\\
4.5282	0.056399\\
4.5319	0.0557\\
4.5357	0.055007\\
4.5394	0.054321\\
4.5431	0.053642\\
4.5469	0.052969\\
4.5506	0.052303\\
4.5543	0.051643\\
4.5581	0.05099\\
4.5618	0.050343\\
4.5655	0.049702\\
4.5692	0.049068\\
4.573	0.04844\\
4.5767	0.047819\\
4.5804	0.047203\\
4.5842	0.046594\\
4.5879	0.045991\\
4.5916	0.045393\\
4.5954	0.044802\\
4.5991	0.044217\\
4.6028	0.043638\\
4.6066	0.043065\\
4.6103	0.042497\\
4.614	0.041936\\
4.6178	0.04138\\
4.6215	0.04083\\
4.6252	0.040286\\
4.629	0.039747\\
4.6327	0.039214\\
4.6364	0.038686\\
4.6402	0.038164\\
4.6439	0.037648\\
4.6476	0.037137\\
4.6514	0.036632\\
4.6551	0.036131\\
4.6588	0.035637\\
4.6626	0.035147\\
4.6663	0.034663\\
4.67	0.034184\\
4.6738	0.033711\\
4.6775	0.033242\\
4.6812	0.032779\\
4.6849	0.03232\\
4.6887	0.031867\\
4.6924	0.031419\\
4.6961	0.030976\\
4.6999	0.030538\\
4.7036	0.030104\\
4.7073	0.029676\\
4.7111	0.029252\\
4.7148	0.028833\\
4.7185	0.028419\\
4.7223	0.02801\\
4.726	0.027606\\
4.7297	0.027206\\
4.7335	0.02681\\
4.7372	0.02642\\
4.7409	0.026034\\
4.7447	0.025652\\
4.7484	0.025275\\
4.7521	0.024902\\
4.7559	0.024534\\
4.7596	0.02417\\
4.7633	0.023811\\
4.7671	0.023456\\
4.7708	0.023105\\
4.7745	0.022758\\
4.7783	0.022416\\
4.782	0.022078\\
4.7857	0.021744\\
4.7895	0.021414\\
4.7932	0.021088\\
4.7969	0.020766\\
4.8006	0.020448\\
4.8044	0.020134\\
4.8081	0.019824\\
4.8118	0.019518\\
4.8156	0.019216\\
4.8193	0.018918\\
4.823	0.018624\\
4.8268	0.018333\\
4.8305	0.018046\\
4.8342	0.017763\\
4.838	0.017483\\
4.8417	0.017207\\
4.8454	0.016935\\
4.8492	0.016666\\
4.8529	0.016401\\
4.8566	0.016139\\
4.8604	0.015881\\
4.8641	0.015626\\
4.8678	0.015374\\
4.8716	0.015126\\
4.8753	0.014882\\
4.879	0.01464\\
4.8828	0.014402\\
4.8865	0.014167\\
4.8902	0.013936\\
4.894	0.013707\\
4.8977	0.013482\\
4.9014	0.01326\\
4.9052	0.01304\\
4.9089	0.012824\\
4.9126	0.012611\\
4.9163	0.012401\\
4.9201	0.012194\\
4.9238	0.01199\\
4.9275	0.011789\\
4.9313	0.01159\\
4.935	0.011395\\
4.9387	0.011202\\
4.9425	0.011012\\
4.9462	0.010825\\
4.9499	0.01064\\
4.9537	0.010459\\
4.9574	0.010279\\
4.9611	0.010103\\
4.9649	0.0099289\\
4.9686	0.0097575\\
4.9723	0.0095887\\
4.9761	0.0094223\\
4.9798	0.0092585\\
4.9835	0.0090971\\
4.9873	0.0089381\\
4.991	0.0087815\\
4.9947	0.0086273\\
4.9985	0.0084755\\
5.0022	0.0083259\\
5.0059	0.0081786\\
5.0097	0.0080336\\
5.0134	0.0078908\\
5.0171	0.0077502\\
5.0208	0.0076118\\
5.0246	0.0074756\\
5.0283	0.0073414\\
5.032	0.0072094\\
5.0358	0.0070794\\
5.0395	0.0069515\\
5.0432	0.0068256\\
5.047	0.0067017\\
5.0507	0.0065797\\
5.0544	0.0064597\\
5.0582	0.0063416\\
5.0619	0.0062254\\
5.0656	0.0061111\\
5.0694	0.0059986\\
5.0731	0.005888\\
5.0768	0.0057791\\
5.0806	0.005672\\
5.0843	0.0055667\\
5.088	0.005463\\
5.0918	0.0053611\\
5.0955	0.0052609\\
5.0992	0.0051623\\
5.103	0.0050653\\
5.1067	0.00497\\
5.1104	0.0048763\\
5.1142	0.0047841\\
5.1179	0.0046934\\
5.1216	0.0046043\\
5.1254	0.0045167\\
5.1291	0.0044306\\
5.1328	0.0043459\\
5.1365	0.0042627\\
5.1403	0.0041809\\
5.144	0.0041005\\
5.1477	0.0040215\\
5.1515	0.0039438\\
5.1552	0.0038675\\
5.1589	0.0037925\\
5.1627	0.0037188\\
5.1664	0.0036464\\
5.1701	0.0035752\\
5.1739	0.0035053\\
5.1776	0.0034366\\
5.1813	0.0033691\\
5.1851	0.0033029\\
5.1888	0.0032377\\
5.1925	0.0031738\\
5.1963	0.003111\\
5.2	0.0030493\\
5.2037	0.0029887\\
5.2075	0.0029291\\
5.2112	0.0028707\\
5.2149	0.0028133\\
5.2187	0.0027569\\
5.2224	0.0027016\\
5.2261	0.0026473\\
5.2299	0.0025939\\
5.2336	0.0025416\\
5.2373	0.0024901\\
5.2411	0.0024397\\
5.2448	0.0023901\\
5.2485	0.0023415\\
5.2522	0.0022938\\
5.256	0.0022469\\
5.2597	0.0022009\\
5.2634	0.0021558\\
5.2672	0.0021115\\
5.2709	0.0020681\\
5.2746	0.0020255\\
5.2784	0.0019836\\
5.2821	0.0019426\\
5.2858	0.0019023\\
5.2896	0.0018628\\
5.2933	0.001824\\
5.297	0.001786\\
5.3008	0.0017487\\
5.3045	0.0017121\\
5.3082	0.0016762\\
5.312	0.001641\\
5.3157	0.0016065\\
5.3194	0.0015726\\
5.3232	0.0015394\\
5.3269	0.0015069\\
5.3306	0.001475\\
5.3306	0.0012365\\
5.3269	0.0012647\\
5.3232	0.0012935\\
5.3194	0.0013229\\
5.3157	0.0013529\\
5.312	0.0013836\\
5.3082	0.0014148\\
5.3045	0.0014468\\
5.3008	0.0014793\\
5.297	0.0015126\\
5.2933	0.0015465\\
5.2896	0.0015811\\
5.2858	0.0016164\\
5.2821	0.0016525\\
5.2784	0.0016892\\
5.2746	0.0017268\\
5.2709	0.001765\\
5.2672	0.0018041\\
5.2634	0.0018439\\
5.2597	0.0018845\\
5.256	0.0019259\\
5.2522	0.0019682\\
5.2485	0.0020113\\
5.2448	0.0020552\\
5.2411	0.0021\\
5.2373	0.0021457\\
5.2336	0.0021923\\
5.2299	0.0022398\\
5.2261	0.0022882\\
5.2224	0.0023376\\
5.2187	0.0023879\\
5.2149	0.0024392\\
5.2112	0.0024915\\
5.2075	0.0025448\\
5.2037	0.0025991\\
5.2	0.0026545\\
5.1963	0.0027109\\
5.1925	0.0027684\\
5.1888	0.002827\\
5.1851	0.0028867\\
5.1813	0.0029475\\
5.1776	0.0030095\\
5.1739	0.0030726\\
5.1701	0.0031369\\
5.1664	0.0032024\\
5.1627	0.0032692\\
5.1589	0.0033371\\
5.1552	0.0034063\\
5.1515	0.0034768\\
5.1477	0.0035486\\
5.144	0.0036217\\
5.1403	0.0036962\\
5.1365	0.003772\\
5.1328	0.0038491\\
5.1291	0.0039277\\
5.1254	0.0040077\\
5.1216	0.0040891\\
5.1179	0.004172\\
5.1142	0.0042564\\
5.1104	0.0043422\\
5.1067	0.0044296\\
5.103	0.0045186\\
5.0992	0.0046091\\
5.0955	0.0047012\\
5.0918	0.0047949\\
5.088	0.0048902\\
5.0843	0.0049872\\
5.0806	0.0050859\\
5.0768	0.0051863\\
5.0731	0.0052884\\
5.0694	0.0053923\\
5.0656	0.005498\\
5.0619	0.0056054\\
5.0582	0.0057147\\
5.0544	0.0058259\\
5.0507	0.0059389\\
5.047	0.0060538\\
5.0432	0.0061707\\
5.0395	0.0062895\\
5.0358	0.0064102\\
5.032	0.006533\\
5.0283	0.0066578\\
5.0246	0.0067847\\
5.0208	0.0069137\\
5.0171	0.0070447\\
5.0134	0.0071779\\
5.0097	0.0073133\\
5.0059	0.0074509\\
5.0022	0.0075906\\
4.9985	0.0077327\\
4.9947	0.0078769\\
4.991	0.0080235\\
4.9873	0.0081725\\
4.9835	0.0083238\\
4.9798	0.0084774\\
4.9761	0.0086335\\
4.9723	0.008792\\
4.9686	0.008953\\
4.9649	0.0091165\\
4.9611	0.0092826\\
4.9574	0.0094512\\
4.9537	0.0096224\\
4.9499	0.0097962\\
4.9462	0.0099726\\
4.9425	0.010152\\
4.9387	0.010334\\
4.935	0.010518\\
4.9313	0.010706\\
4.9275	0.010896\\
4.9238	0.011089\\
4.9201	0.011285\\
4.9163	0.011483\\
4.9126	0.011685\\
4.9089	0.01189\\
4.9052	0.012097\\
4.9014	0.012308\\
4.8977	0.012521\\
4.894	0.012738\\
4.8902	0.012958\\
4.8865	0.013181\\
4.8828	0.013407\\
4.879	0.013637\\
4.8753	0.013869\\
4.8716	0.014105\\
4.8678	0.014345\\
4.8641	0.014587\\
4.8604	0.014833\\
4.8566	0.015082\\
4.8529	0.015335\\
4.8492	0.015591\\
4.8454	0.015851\\
4.8417	0.016114\\
4.838	0.016381\\
4.8342	0.016652\\
4.8305	0.016926\\
4.8268	0.017204\\
4.823	0.017485\\
4.8193	0.01777\\
4.8156	0.018059\\
4.8118	0.018352\\
4.8081	0.018649\\
4.8044	0.018949\\
4.8006	0.019254\\
4.7969	0.019562\\
4.7932	0.019874\\
4.7895	0.02019\\
4.7857	0.020511\\
4.782	0.020835\\
4.7783	0.021164\\
4.7745	0.021496\\
4.7708	0.021833\\
4.7671	0.022174\\
4.7633	0.02252\\
4.7596	0.022869\\
4.7559	0.023223\\
4.7521	0.023581\\
4.7484	0.023944\\
4.7447	0.024311\\
4.7409	0.024682\\
4.7372	0.025058\\
4.7335	0.025439\\
4.7297	0.025824\\
4.726	0.026213\\
4.7223	0.026607\\
4.7185	0.027006\\
4.7148	0.02741\\
4.7111	0.027818\\
4.7073	0.028231\\
4.7036	0.028649\\
4.6999	0.029071\\
4.6961	0.029499\\
4.6924	0.029931\\
4.6887	0.030368\\
4.6849	0.030811\\
4.6812	0.031258\\
4.6775	0.03171\\
4.6738	0.032167\\
4.67	0.032629\\
4.6663	0.033097\\
4.6626	0.03357\\
4.6588	0.034047\\
4.6551	0.03453\\
4.6514	0.035019\\
4.6476	0.035512\\
4.6439	0.036011\\
4.6402	0.036515\\
4.6364	0.037025\\
4.6327	0.03754\\
4.629	0.038061\\
4.6252	0.038587\\
4.6215	0.039118\\
4.6178	0.039655\\
4.614	0.040198\\
4.6103	0.040746\\
4.6066	0.0413\\
4.6028	0.041859\\
4.5991	0.042425\\
4.5954	0.042996\\
4.5916	0.043573\\
4.5879	0.044155\\
4.5842	0.044744\\
4.5804	0.045338\\
4.5767	0.045939\\
4.573	0.046545\\
4.5692	0.047157\\
4.5655	0.047776\\
4.5618	0.0484\\
4.5581	0.04903\\
4.5543	0.049667\\
4.5506	0.05031\\
4.5469	0.050958\\
4.5431	0.051614\\
4.5394	0.052275\\
4.5357	0.052943\\
4.5319	0.053617\\
4.5282	0.054297\\
4.5245	0.054984\\
4.5207	0.055677\\
4.517	0.056377\\
4.5133	0.057083\\
4.5095	0.057796\\
4.5058	0.058515\\
4.5021	0.059241\\
4.4983	0.059974\\
4.4946	0.060713\\
4.4909	0.061459\\
4.4871	0.062212\\
4.4834	0.062972\\
4.4797	0.063739\\
4.4759	0.064512\\
4.4722	0.065293\\
4.4685	0.06608\\
4.4647	0.066875\\
4.461	0.067677\\
4.4573	0.068485\\
4.4536	0.069301\\
4.4498	0.070124\\
4.4461	0.070955\\
4.4424	0.071793\\
4.4386	0.072638\\
4.4349	0.073491\\
4.4312	0.074351\\
4.4274	0.075218\\
4.4237	0.076093\\
4.42	0.076976\\
4.4162	0.077867\\
4.4125	0.078765\\
4.4088	0.079671\\
4.405	0.080585\\
4.4013	0.081507\\
4.3976	0.082437\\
4.3938	0.083375\\
4.3901	0.084321\\
4.3864	0.085275\\
4.3826	0.086238\\
4.3789	0.087208\\
4.3752	0.088188\\
4.3714	0.089175\\
4.3677	0.090171\\
4.364	0.091176\\
4.3602	0.092189\\
4.3565	0.093212\\
4.3528	0.094242\\
4.349	0.095282\\
4.3453	0.096331\\
4.3416	0.097389\\
4.3379	0.098456\\
4.3341	0.099532\\
4.3304	0.10062\\
4.3267	0.10171\\
4.3229	0.10282\\
4.3192	0.10393\\
4.3155	0.10505\\
4.3117	0.10619\\
4.308	0.10733\\
4.3043	0.10848\\
4.3005	0.10965\\
4.2968	0.11082\\
4.2931	0.112\\
4.2893	0.1132\\
4.2856	0.1144\\
4.2819	0.11561\\
4.2781	0.11684\\
4.2744	0.11807\\
4.2707	0.11932\\
4.2669	0.12058\\
4.2632	0.12185\\
4.2595	0.12313\\
4.2557	0.12442\\
4.252	0.12572\\
4.2483	0.12703\\
4.2445	0.12836\\
4.2408	0.1297\\
4.2371	0.13104\\
4.2333	0.13241\\
4.2296	0.13378\\
4.2259	0.13516\\
4.2222	0.13656\\
4.2184	0.13797\\
4.2147	0.13939\\
4.211	0.14083\\
4.2072	0.14228\\
4.2035	0.14374\\
4.1998	0.14521\\
4.196	0.1467\\
4.1923	0.1482\\
4.1886	0.14972\\
4.1848	0.15124\\
4.1811	0.15279\\
4.1774	0.15434\\
4.1736	0.15591\\
4.1699	0.1575\\
4.1662	0.15909\\
4.1624	0.16071\\
4.1587	0.16234\\
4.155	0.16398\\
4.1512	0.16564\\
4.1475	0.16731\\
4.1438	0.169\\
4.14	0.1707\\
4.1363	0.17242\\
4.1326	0.17415\\
4.1288	0.1759\\
4.1251	0.17767\\
4.1214	0.17945\\
4.1176	0.18125\\
4.1139	0.18307\\
4.1102	0.1849\\
4.1065	0.18675\\
4.1027	0.18862\\
4.099	0.1905\\
4.0953	0.1924\\
4.0915	0.19432\\
4.0878	0.19625\\
4.0841	0.19821\\
4.0803	0.20018\\
4.0766	0.20217\\
4.0729	0.20417\\
4.0691	0.2062\\
4.0654	0.20824\\
4.0617	0.21031\\
4.0579	0.21239\\
4.0542	0.21449\\
4.0505	0.21661\\
4.0467	0.21875\\
4.043	0.22091\\
4.0393	0.22309\\
4.0355	0.22529\\
4.0318	0.22751\\
4.0281	0.22974\\
4.0243	0.232\\
4.0206	0.23428\\
4.0169	0.23658\\
4.0131	0.2389\\
4.0094	0.24124\\
4.0057	0.24361\\
4.002	0.24599\\
3.9982	0.24839\\
3.9945	0.25082\\
3.9908	0.25327\\
3.987	0.25574\\
3.9833	0.25823\\
3.9796	0.26074\\
3.9758	0.26327\\
3.9721	0.26583\\
3.9684	0.26841\\
3.9646	0.27101\\
3.9609	0.27364\\
3.9572	0.27628\\
3.9534	0.27895\\
3.9497	0.28165\\
3.946	0.28436\\
3.9422	0.2871\\
3.9385	0.28987\\
3.9348	0.29265\\
3.931	0.29546\\
3.9273	0.29829\\
3.9236	0.30115\\
3.9198	0.30403\\
3.9161	0.30694\\
3.9124	0.30986\\
3.9086	0.31282\\
3.9049	0.31579\\
3.9012	0.31879\\
3.8974	0.32182\\
3.8937	0.32487\\
3.89	0.32794\\
3.8863	0.33104\\
3.8825	0.33417\\
3.8788	0.33732\\
3.8751	0.34049\\
3.8713	0.34369\\
3.8676	0.34691\\
3.8639	0.35016\\
3.8601	0.35344\\
3.8564	0.35674\\
3.8527	0.36006\\
3.8489	0.36341\\
3.8452	0.36679\\
3.8415	0.37019\\
3.8377	0.37362\\
3.834	0.37707\\
3.8303	0.38055\\
3.8265	0.38405\\
3.8228	0.38758\\
3.8191	0.39114\\
3.8153	0.39472\\
3.8116	0.39833\\
3.8079	0.40196\\
3.8041	0.40562\\
3.8004	0.4093\\
3.7967	0.41301\\
3.7929	0.41675\\
3.7892	0.42051\\
3.7855	0.4243\\
3.7817	0.42811\\
3.778	0.43195\\
3.7743	0.43582\\
3.7706	0.43971\\
3.7668	0.44363\\
3.7631	0.44757\\
3.7594	0.45153\\
3.7556	0.45553\\
3.7519	0.45955\\
3.7482	0.46359\\
3.7444	0.46766\\
3.7407	0.47175\\
3.737	0.47587\\
3.7332	0.48001\\
3.7295	0.48418\\
3.7258	0.48837\\
3.722	0.49258\\
3.7183	0.49682\\
3.7146	0.50109\\
3.7108	0.50537\\
3.7071	0.50968\\
3.7034	0.51401\\
3.6996	0.51836\\
3.6959	0.52274\\
3.6922	0.52714\\
3.6884	0.53156\\
3.6847	0.53599\\
3.681	0.54045\\
3.6772	0.54493\\
3.6735	0.54943\\
3.6698	0.55395\\
3.666	0.55848\\
3.6623	0.56304\\
3.6586	0.56761\\
3.6549	0.57219\\
3.6511	0.57679\\
3.6474	0.58141\\
3.6437	0.58604\\
3.6399	0.59068\\
3.6362	0.59534\\
3.6325	0.60001\\
3.6287	0.60469\\
3.625	0.60937\\
3.6213	0.61407\\
3.6175	0.61878\\
3.6138	0.62349\\
3.6101	0.6282\\
3.6063	0.63293\\
3.6026	0.63765\\
3.5989	0.64238\\
3.5951	0.64711\\
3.5914	0.65184\\
3.5877	0.65657\\
3.5839	0.6613\\
3.5802	0.66603\\
3.5765	0.67075\\
3.5727	0.67547\\
3.569	0.68018\\
3.5653	0.68489\\
3.5615	0.68959\\
3.5578	0.69428\\
3.5541	0.69896\\
3.5504	0.70363\\
3.5466	0.70828\\
3.5429	0.71293\\
3.5392	0.71756\\
3.5354	0.72218\\
3.5317	0.72678\\
3.528	0.73136\\
3.5242	0.73593\\
3.5205	0.74048\\
3.5168	0.74501\\
3.513	0.74952\\
3.5093	0.75401\\
3.5056	0.75848\\
3.5018	0.76293\\
3.4981	0.76735\\
3.4944	0.77174\\
3.4906	0.77611\\
3.4869	0.78046\\
3.4832	0.78477\\
3.4794	0.78906\\
3.4757	0.79332\\
3.472	0.79755\\
3.4682	0.80174\\
3.4645	0.8059\\
3.4608	0.81003\\
3.457	0.81412\\
3.4533	0.81818\\
3.4496	0.8222\\
3.4458	0.82618\\
3.4421	0.83012\\
3.4384	0.83402\\
3.4347	0.83788\\
3.4309	0.84169\\
3.4272	0.84546\\
3.4235	0.84918\\
3.4197	0.85286\\
3.416	0.85648\\
3.4123	0.86006\\
3.4085	0.86358\\
3.4048	0.86705\\
3.4011	0.87046\\
3.3973	0.87382\\
3.3936	0.87712\\
3.3899	0.88036\\
3.3861	0.88353\\
3.3824	0.88665\\
3.3787	0.8897\\
3.3749	0.89268\\
3.3712	0.8956\\
3.3675	0.89844\\
3.3637	0.90121\\
3.36	0.90392\\
3.3563	0.90654\\
3.3525	0.90909\\
3.3488	0.91156\\
3.3451	0.91395\\
3.3413	0.91627\\
3.3376	0.91849\\
3.3339	0.92064\\
3.3301	0.9227\\
3.3264	0.92467\\
3.3227	0.92656\\
3.319	0.92836\\
3.3152	0.93006\\
3.3115	0.93168\\
3.3078	0.93321\\
3.304	0.93465\\
3.3003	0.936\\
3.2966	0.93725\\
3.2928	0.93841\\
3.2891	0.93949\\
3.2854	0.94047\\
3.2816	0.94136\\
3.2779	0.94216\\
3.2742	0.94286\\
3.2704	0.94348\\
3.2667	0.94401\\
3.263	0.94445\\
3.2592	0.9448\\
3.2555	0.94506\\
3.2518	0.94523\\
3.248	0.94531\\
3.2443	0.9453\\
3.2406	0.94521\\
3.2368	0.94502\\
3.2331	0.94475\\
3.2294	0.94438\\
3.2256	0.94393\\
3.2219	0.94339\\
3.2182	0.94275\\
3.2144	0.94202\\
3.2107	0.9412\\
3.207	0.94029\\
3.2033	0.93928\\
3.1995	0.93818\\
3.1958	0.93699\\
3.1921	0.93569\\
3.1883	0.9343\\
3.1846	0.93282\\
3.1809	0.93123\\
3.1771	0.92955\\
3.1734	0.92777\\
3.1697	0.92589\\
3.1659	0.92392\\
3.1622	0.92185\\
3.1585	0.91968\\
3.1547	0.91742\\
3.151	0.91507\\
3.1473	0.91262\\
3.1435	0.91007\\
3.1398	0.90744\\
3.1361	0.90472\\
3.1323	0.9019\\
3.1286	0.899\\
3.1249	0.89601\\
3.1211	0.89294\\
3.1174	0.88978\\
3.1137	0.88654\\
3.1099	0.88322\\
3.1062	0.87981\\
3.1025	0.87632\\
3.0988	0.87276\\
3.095	0.86911\\
3.0913	0.86539\\
3.0876	0.86158\\
3.0838	0.8577\\
3.0801	0.85375\\
3.0764	0.84972\\
3.0726	0.84561\\
3.0689	0.84142\\
3.0652	0.83717\\
3.0614	0.83284\\
3.0577	0.82843\\
3.054	0.82396\\
3.0502	0.81941\\
3.0465	0.8148\\
3.0428	0.81012\\
3.039	0.80537\\
3.0353	0.80055\\
3.0316	0.79568\\
3.0278	0.79074\\
3.0241	0.78574\\
3.0204	0.78069\\
3.0166	0.77558\\
3.0129	0.77042\\
3.0092	0.76521\\
3.0054	0.75995\\
3.0017	0.75465\\
2.998	0.7493\\
2.9942	0.74392\\
2.9905	0.7385\\
2.9868	0.73304\\
2.9831	0.72755\\
2.9793	0.72203\\
2.9756	0.71648\\
2.9719	0.7109\\
2.9681	0.7053\\
2.9644	0.69968\\
2.9607	0.69403\\
2.9569	0.68836\\
2.9532	0.68267\\
2.9495	0.67697\\
2.9457	0.67125\\
2.942	0.66551\\
2.9383	0.65976\\
2.9345	0.654\\
2.9308	0.64823\\
2.9271	0.64244\\
2.9233	0.63665\\
2.9196	0.63085\\
2.9159	0.62504\\
2.9121	0.61923\\
2.9084	0.61341\\
2.9047	0.60759\\
2.9009	0.60176\\
2.8972	0.59594\\
2.8935	0.59011\\
2.8897	0.58429\\
2.886	0.57847\\
2.8823	0.57265\\
2.8785	0.56683\\
2.8748	0.56103\\
2.8711	0.55522\\
2.8674	0.54943\\
2.8636	0.54364\\
2.8599	0.53787\\
2.8562	0.53211\\
2.8524	0.52635\\
2.8487	0.52061\\
2.845	0.51489\\
2.8412	0.50918\\
2.8375	0.50349\\
2.8338	0.49781\\
2.83	0.49216\\
2.8263	0.48652\\
2.8226	0.4809\\
2.8188	0.4753\\
2.8151	0.46973\\
2.8114	0.46418\\
2.8076	0.45865\\
2.8039	0.45315\\
2.8002	0.44767\\
2.7964	0.44222\\
2.7927	0.4368\\
2.789	0.4314\\
2.7852	0.42603\\
2.7815	0.4207\\
2.7778	0.41539\\
2.774	0.41011\\
2.7703	0.40487\\
2.7666	0.39965\\
2.7628	0.39447\\
2.7591	0.38933\\
2.7554	0.38421\\
2.7517	0.37914\\
2.7479	0.37409\\
2.7442	0.36909\\
2.7405	0.36412\\
2.7367	0.35918\\
2.733	0.35428\\
2.7293	0.34942\\
2.7255	0.3446\\
2.7218	0.33982\\
2.7181	0.33507\\
2.7143	0.33037\\
2.7106	0.3257\\
2.7069	0.32108\\
2.7031	0.31649\\
2.6994	0.31194\\
2.6957	0.30744\\
2.6919	0.30297\\
2.6882	0.29855\\
2.6845	0.29417\\
2.6807	0.28983\\
2.677	0.28553\\
2.6733	0.28127\\
2.6695	0.27705\\
2.6658	0.27288\\
2.6621	0.26875\\
2.6583	0.26466\\
2.6546	0.26061\\
2.6509	0.25661\\
2.6472	0.25265\\
2.6434	0.24873\\
2.6397	0.24485\\
2.636	0.24102\\
2.6322	0.23723\\
2.6285	0.23348\\
2.6248	0.22977\\
2.621	0.2261\\
2.6173	0.22248\\
2.6136	0.2189\\
2.6098	0.21536\\
2.6061	0.21187\\
2.6024	0.20841\\
2.5986	0.205\\
2.5949	0.20163\\
2.5912	0.1983\\
2.5874	0.19501\\
2.5837	0.19176\\
2.58	0.18856\\
2.5762	0.18539\\
2.5725	0.18227\\
2.5688	0.17918\\
2.565	0.17614\\
2.5613	0.17313\\
2.5576	0.17017\\
2.5538	0.16724\\
2.5501	0.16436\\
2.5464	0.16151\\
2.5426	0.1587\\
2.5389	0.15593\\
2.5352	0.1532\\
2.5315	0.15051\\
2.5277	0.14785\\
2.524	0.14524\\
2.5203	0.14265\\
2.5165	0.14011\\
2.5128	0.1376\\
2.5091	0.13513\\
2.5053	0.1327\\
2.5016	0.1303\\
2.4979	0.12793\\
2.4941	0.1256\\
2.4904	0.12331\\
2.4867	0.12105\\
2.4829	0.11882\\
2.4792	0.11663\\
2.4755	0.11447\\
2.4717	0.11235\\
2.468	0.11025\\
2.4643	0.10819\\
2.4605	0.10617\\
2.4568	0.10417\\
2.4531	0.1022\\
2.4493	0.10027\\
2.4456	0.098368\\
2.4419	0.096495\\
2.4381	0.094653\\
2.4344	0.09284\\
2.4307	0.091058\\
2.4269	0.089304\\
2.4232	0.087579\\
2.4195	0.085882\\
2.4158	0.084214\\
2.412	0.082573\\
2.4083	0.08096\\
2.4046	0.079374\\
2.4008	0.077815\\
2.3971	0.076282\\
2.3934	0.074775\\
2.3896	0.073294\\
2.3859	0.071839\\
2.3822	0.070409\\
2.3784	0.069003\\
2.3747	0.067622\\
2.371	0.066265\\
2.3672	0.064932\\
2.3635	0.063623\\
2.3598	0.062337\\
2.356	0.061073\\
2.3523	0.059832\\
2.3486	0.058614\\
2.3448	0.057417\\
2.3411	0.056242\\
2.3374	0.055089\\
2.3336	0.053956\\
2.3299	0.052844\\
2.3262	0.051753\\
2.3224	0.050681\\
2.3187	0.04963\\
2.315	0.048598\\
2.3112	0.047585\\
2.3075	0.046591\\
2.3038	0.045616\\
2.3001	0.044659\\
2.2963	0.04372\\
2.2926	0.0428\\
2.2889	0.041896\\
2.2851	0.04101\\
2.2814	0.040141\\
2.2777	0.039288\\
2.2739	0.038452\\
2.2702	0.037633\\
2.2665	0.036829\\
2.2627	0.03604\\
2.259	0.035268\\
2.2553	0.03451\\
2.2515	0.033767\\
2.2478	0.033039\\
2.2441	0.032325\\
2.2403	0.031626\\
2.2366	0.03094\\
2.2329	0.030268\\
2.2291	0.02961\\
2.2254	0.028964\\
2.2217	0.028332\\
2.2179	0.027712\\
2.2142	0.027105\\
2.2105	0.026511\\
2.2067	0.025928\\
2.203	0.025358\\
2.1993	0.024799\\
2.1956	0.024251\\
2.1918	0.023715\\
2.1881	0.02319\\
2.1844	0.022675\\
2.1806	0.022172\\
2.1769	0.021679\\
2.1732	0.021196\\
2.1694	0.020723\\
2.1657	0.02026\\
2.162	0.019807\\
2.1582	0.019363\\
2.1545	0.018929\\
2.1508	0.018504\\
2.147	0.018088\\
2.1433	0.01768\\
2.1396	0.017282\\
2.1358	0.016892\\
2.1321	0.01651\\
2.1284	0.016136\\
2.1246	0.015771\\
2.1209	0.015413\\
2.1172	0.015063\\
2.1134	0.01472\\
2.1097	0.014385\\
2.106	0.014058\\
2.1022	0.013737\\
2.0985	0.013423\\
2.0948	0.013116\\
2.091	0.012816\\
2.0873	0.012523\\
2.0836	0.012235\\
2.0799	0.011954\\
2.0761	0.01168\\
2.0724	0.011411\\
2.0687	0.011148\\
2.0649	0.010891\\
2.0612	0.01064\\
2.0575	0.010395\\
2.0537	0.010154\\
2.05	0.0099195\\
2.0463	0.0096898\\
2.0425	0.0094653\\
2.0388	0.0092458\\
2.0351	0.0090312\\
2.0313	0.0088214\\
2.0276	0.0086163\\
2.0239	0.0084159\\
2.0201	0.0082199\\
2.0164	0.0080283\\
2.0127	0.0078411\\
2.0089	0.0076581\\
2.0052	0.0074792\\
2.0015	0.0073044\\
1.9977	0.0071336\\
1.994	0.0069666\\
1.9903	0.0068034\\
1.9865	0.0066439\\
1.9828	0.0064881\\
1.9791	0.0063358\\
1.9753	0.006187\\
1.9716	0.0060416\\
1.9679	0.0058995\\
1.9642	0.0057606\\
1.9604	0.005625\\
1.9567	0.0054925\\
1.953	0.005363\\
1.9492	0.0052364\\
1.9455	0.0051128\\
1.9418	0.0049921\\
1.938	0.0048741\\
1.9343	0.0047588\\
1.9306	0.0046462\\
1.9268	0.0045362\\
1.9231	0.0044288\\
1.9194	0.0043238\\
1.9156	0.0042213\\
1.9119	0.0041211\\
1.9082	0.0040233\\
1.9044	0.0039277\\
1.9007	0.0038344\\
1.897	0.0037432\\
1.8932	0.0036542\\
1.8895	0.0035672\\
1.8858	0.0034823\\
1.882	0.0033993\\
1.8783	0.0033183\\
1.8746	0.0032392\\
1.8708	0.0031619\\
1.8671	0.0030864\\
1.8634	0.0030127\\
1.8596	0.0029407\\
1.8559	0.0028704\\
1.8522	0.0028017\\
1.8485	0.0027347\\
1.8447	0.0026692\\
1.841	0.0026053\\
1.8373	0.0025429\\
1.8335	0.0024819\\
1.8298	0.0024224\\
1.8261	0.0023642\\
1.8223	0.0023075\\
1.8186	0.002252\\
1.8149	0.0021979\\
1.8111	0.0021451\\
1.8074	0.0020935\\
1.8037	0.0020431\\
1.7999	0.0019939\\
1.7962	0.0019459\\
1.7925	0.001899\\
1.7887	0.0018532\\
1.785	0.0018085\\
1.7813	0.0017648\\
1.7775	0.0017222\\
1.7738	0.0016806\\
1.7701	0.00164\\
1.7663	0.0016003\\
1.7626	0.0015616\\
1.7589	0.0015238\\
1.7551	0.0014869\\
1.7514	0.0014509\\
1.7477	0.0014157\\
1.744	0.0013814\\
1.7402	0.0013478\\
1.7365	0.0013151\\
1.7328	0.0012832\\
1.729	0.001252\\
1.7253	0.0012215\\
1.7216	0.0011918\\
1.7178	0.0011628\\
1.7141	0.0011345\\
1.7104	0.0011068\\
1.7066	0.0010798\\
1.7029	0.0010535\\
1.6992	0.0010278\\
1.6954	0.0010027\\
1.6917	0.00097816\\
1.688	0.00095424\\
1.6842	0.00093089\\
1.6805	0.0009081\\
1.6768	0.00088585\\
1.673	0.00086414\\
1.6693	0.00084295\\
1.6656	0.00082226\\
1.6618	0.00080207\\
1.6581	0.00078236\\
1.6544	0.00076312\\
1.6506	0.00074435\\
1.6469	0.00072603\\
1.6432	0.00070814\\
1.6394	0.00069069\\
1.6357	0.00067365\\
1.632	0.00065702\\
1.6283	0.00064079\\
1.6245	0.00062496\\
1.6208	0.0006095\\
1.6171	0.00059442\\
1.6133	0.00057969\\
1.6096	0.00056533\\
1.6059	0.00055131\\
1.6021	0.00053762\\
}--cycle;
\addplot[ybar interval, fill=black, fill opacity=0.35, area legend, draw=none] table[row sep=crcr, x=Lower, y=Count] {%
Lower	Upper	Count\\
1.6021	1.8889	0\\
1.8889	2.1757	0.018596\\
2.1757	2.4625	0.092978\\
2.4625	2.7494	0.28823\\
2.7494	3.0362	0.60435\\
3.0362	3.323	0.97626\\
3.323	3.6098	0.79031\\
3.6098	3.8966	0.39051\\
3.8966	4.1834	0.20455\\
4.1834	4.4702	0.037191\\
4.4702	4.757	0.055787\\
4.757	5.0438	0.0092978\\
5.0438	5.3306	0.018596\\
5.3306	5.3306	0.018596\\
};
\addlegendentry{\RealHistogramLegend}

\addplot [color=black]
  table[row sep=crcr]{%
1.6021	0.00032335\\
1.7328	0.0011012\\
1.8186	0.0023301\\
1.882	0.0039413\\
1.9343	0.0059668\\
1.9791	0.0084032\\
2.0201	0.01138\\
2.0575	0.014862\\
2.091	0.018763\\
2.1209	0.022952\\
2.1508	0.027925\\
2.1769	0.033008\\
2.203	0.038857\\
2.2291	0.045554\\
2.2515	0.052036\\
2.2739	0.059262\\
2.2963	0.067288\\
2.3187	0.076171\\
2.3411	0.085967\\
2.3635	0.096732\\
2.3822	0.10648\\
2.4008	0.11696\\
2.4195	0.12821\\
2.4381	0.14025\\
2.4568	0.1531\\
2.4755	0.16677\\
2.4941	0.18129\\
2.5128	0.19665\\
2.5315	0.21288\\
2.5501	0.22996\\
2.5725	0.25158\\
2.5949	0.27441\\
2.6173	0.29841\\
2.6397	0.32353\\
2.6621	0.34971\\
2.6882	0.3815\\
2.7143	0.41448\\
2.7442	0.45339\\
2.7815	0.50341\\
2.9121	0.67998\\
2.9383	0.71334\\
2.9607	0.74084\\
2.9831	0.76709\\
3.0017	0.78788\\
3.0204	0.80755\\
3.039	0.82601\\
3.054	0.83982\\
3.0689	0.85274\\
3.0838	0.86471\\
3.0988	0.8757\\
3.1137	0.88565\\
3.1286	0.89454\\
3.1435	0.90233\\
3.1547	0.90743\\
3.1659	0.91189\\
3.1771	0.91569\\
3.1883	0.91883\\
3.1995	0.9213\\
3.2107	0.9231\\
3.2219	0.92423\\
3.2331	0.92467\\
3.2443	0.92444\\
3.2555	0.92354\\
3.2667	0.92196\\
3.2779	0.91971\\
3.2891	0.9168\\
3.3003	0.91323\\
3.3115	0.90902\\
3.3227	0.90417\\
3.3376	0.89674\\
3.3525	0.88823\\
3.3675	0.87868\\
3.3824	0.86813\\
3.3973	0.85662\\
3.4123	0.84419\\
3.4272	0.8309\\
3.4458	0.81315\\
3.4645	0.79424\\
3.4832	0.77427\\
3.5056	0.74908\\
3.528	0.72273\\
3.5541	0.6908\\
3.5877	0.64838\\
3.6362	0.58562\\
3.6996	0.50376\\
3.7332	0.46166\\
3.7631	0.42544\\
3.7892	0.39491\\
3.8153	0.36561\\
3.8377	0.34158\\
3.8601	0.31859\\
3.8825	0.29671\\
3.9049	0.27595\\
3.9273	0.25634\\
3.9497	0.23789\\
3.9721	0.22058\\
3.9945	0.20441\\
4.0169	0.18934\\
4.0393	0.17535\\
4.0617	0.1624\\
4.0841	0.15045\\
4.1065	0.13945\\
4.1288	0.12934\\
4.1512	0.12007\\
4.1736	0.1116\\
4.196	0.10385\\
4.2222	0.095673\\
4.2483	0.088332\\
4.2744	0.081745\\
4.3043	0.075041\\
4.3341	0.069109\\
4.3677	0.063232\\
4.405	0.057535\\
4.4461	0.052103\\
4.4909	0.046972\\
4.5431	0.041793\\
4.6028	0.036667\\
4.6738	0.031371\\
4.7559	0.026015\\
4.8492	0.020717\\
4.9462	0.015987\\
5.0432	0.012016\\
5.144	0.0086684\\
5.2485	0.0059743\\
5.3306	0.0043519\\
};
\addlegendentry{\RealBLNFitLegend}

      % \input{../../../../.tex/20/KS/SizeDistributionFittingsfig21.tex}
      \legend{}
    \end{groupplot}
  
    \begin{groupplot}[
      group style={
        group name=Row4,
        plotGroup,
      },
      plotRow,
      row1Keys,
      yTickScaleLabelEmpty,
      % emptyLeged,
      %
      % xmin=1,xmax=2.3788e+05,
      ymin=0,ymax=1,
    ]
      \nextgroupplot[%
        at={($(Row3 c1r1.south west)-(0,\yPlotSep em)-(0,\plotHeight cm)$)},
        xmin=0.57363,xmax=5.4115,
        % xmin=0.89923,xmax=5.5732,
        % ymin=0,ymax=0.9,
      ] % This file was created by matlab2tikz.
%
\definecolor{mycolor1}{rgb}{0.00000,0.44706,0.69804}%
\definecolor{mycolor2}{rgb}{0.83529,0.36863,0.00000}%
%
\addplot[ybar interval, fill=mycolor1, fill opacity=0.35, area legend, draw=none] table[row sep=crcr, x=Lower, y=Count] {%
Lower	Upper	Count\\
1.1117	1.4658	0.00037656\\
1.4658	1.8198	0.0023347\\
1.8198	2.1739	0.022142\\
2.1739	2.528	0.11959\\
2.528	2.8821	0.39425\\
2.8821	3.2362	0.75951\\
3.2362	3.5903	0.79469\\
3.5903	3.9444	0.44486\\
3.9444	4.2984	0.19024\\
4.2984	4.6525	0.071094\\
4.6525	5.0066	0.017472\\
5.0066	5.3607	0.0076064\\
5.3607	5.3607	0.0076064\\
};
\addlegendentry{\ExactHistogramLegend}

\addplot [color=mycolor1, only marks, every error bar/.append style={opacity=0.45}, mark=*, mark size=0pt, draw=none, forget plot]
 plot [error bars/.cd, y dir=both, y explicit, error bar style={line width=1pt, color=mycolor1}, error mark options={mark=none,mark size=0pt}]
 table[row sep=crcr, y error plus index=2, y error minus index=3]{%
1.2887	0.00037656	0.00032732	0.00032732\\
1.6428	0.0023347	0.00075677	0.00075677\\
1.9969	0.022142	0.0026084	0.0026084\\
2.351	0.11959	0.0055398	0.0055398\\
2.7051	0.39425	0.0089796	0.0089796\\
3.0591	0.75951	0.012806	0.012806\\
3.4132	0.79469	0.012503	0.012503\\
3.7673	0.44486	0.0093146	0.0093146\\
4.1214	0.19024	0.0060323	0.0060323\\
4.4755	0.071094	0.0030582	0.0030582\\
4.8296	0.017472	0.0013072	0.0013072\\
5.1837	0.0076064	0.0002597	0.0002597\\
};
\addplot [color=mycolor1]
  table[row sep=crcr]{%
1.1117	0.00013532\\
1.3839	0.00070887\\
1.5328	0.0016796\\
1.6348	0.0029859\\
1.7156	0.004666\\
1.7837	0.0067523\\
1.8432	0.009284\\
1.8943	0.012149\\
1.9411	0.015494\\
1.9836	0.019267\\
2.0219	0.023381\\
2.0559	0.027706\\
2.0899	0.032756\\
2.1197	0.037849\\
2.1495	0.043648\\
2.1793	0.050232\\
2.2048	0.056563\\
2.2303	0.063586\\
2.2558	0.071357\\
2.2813	0.079935\\
2.3026	0.08774\\
2.3239	0.09618\\
2.3451	0.10529\\
2.3664	0.11509\\
2.3877	0.12563\\
2.4089	0.13693\\
2.4302	0.14902\\
2.4515	0.16192\\
2.4727	0.17566\\
2.494	0.19026\\
2.5153	0.20574\\
2.5365	0.2221\\
2.5578	0.23934\\
2.5791	0.25748\\
2.6003	0.2765\\
2.6216	0.2964\\
2.6429	0.31714\\
2.6684	0.34312\\
2.6939	0.37021\\
2.7194	0.39834\\
2.7492	0.43234\\
2.779	0.4674\\
2.8172	0.51366\\
2.9406	0.66395\\
2.9661	0.69337\\
2.9874	0.71688\\
3.0086	0.73927\\
3.0257	0.75625\\
3.0427	0.7723\\
3.0597	0.78736\\
3.0767	0.80138\\
3.0937	0.81436\\
3.1107	0.82627\\
3.1277	0.8371\\
3.1405	0.84449\\
3.1533	0.8512\\
3.166	0.85716\\
3.1788	0.8623\\
3.1915	0.86655\\
3.2043	0.86985\\
3.2171	0.87214\\
3.2298	0.87342\\
3.2426	0.8737\\
3.2553	0.87304\\
3.2681	0.87149\\
3.2809	0.86911\\
3.2936	0.86594\\
3.3064	0.86202\\
3.3191	0.85735\\
3.3319	0.85194\\
3.3447	0.84578\\
3.3574	0.83886\\
3.3702	0.83121\\
3.3829	0.82285\\
3.3957	0.8138\\
3.4127	0.80077\\
3.4297	0.78675\\
3.4467	0.77188\\
3.468	0.75227\\
3.4893	0.73174\\
3.5148	0.70616\\
3.5488	0.67091\\
3.6084	0.6079\\
3.6636	0.54981\\
3.7019	0.51054\\
3.736	0.47655\\
3.77	0.44359\\
3.7998	0.41571\\
3.8295	0.38885\\
3.855	0.36667\\
3.8806	0.34533\\
3.9061	0.32484\\
3.9316	0.30523\\
3.9571	0.2865\\
3.9826	0.26866\\
4.0082	0.2517\\
4.0337	0.2356\\
4.0592	0.22036\\
4.0847	0.20594\\
4.1102	0.19231\\
4.1358	0.17947\\
4.1613	0.16736\\
4.1868	0.15596\\
4.2123	0.14524\\
4.2421	0.13355\\
4.2719	0.12268\\
4.3016	0.11258\\
4.3314	0.10321\\
4.3612	0.09451\\
4.3952	0.085341\\
4.4292	0.076931\\
4.4633	0.069222\\
4.4973	0.062161\\
4.5356	0.054929\\
4.5738	0.048391\\
4.6164	0.041871\\
4.6589	0.036068\\
4.7014	0.03092\\
4.7482	0.025948\\
4.795	0.021634\\
4.8461	0.017605\\
4.9014	0.013952\\
4.9609	0.010742\\
5.0247	0.0080174\\
5.0928	0.0057875\\
5.1736	0.0038617\\
5.2671	0.0023644\\
5.3607	0.0014168\\
};
\addlegendentry{\ExactBLNFitLegend}


\addplot[area legend, draw=none, fill=mycolor1, fill opacity=0.15, forget plot]
table[row sep=crcr] {%
x	y\\
1.1117	0.000171\\
1.1159	0.00017547\\
1.1202	0.00018006\\
1.1244	0.00018477\\
1.1287	0.00018959\\
1.1329	0.00019454\\
1.1372	0.00019961\\
1.1415	0.00020481\\
1.1457	0.00021013\\
1.15	0.00021559\\
1.1542	0.00022119\\
1.1585	0.00022693\\
1.1627	0.00023281\\
1.167	0.00023884\\
1.1712	0.00024502\\
1.1755	0.00025135\\
1.1797	0.00025784\\
1.184	0.00026449\\
1.1882	0.0002713\\
1.1925	0.00027829\\
1.1967	0.00028545\\
1.201	0.00029278\\
1.2053	0.0003003\\
1.2095	0.000308\\
1.2138	0.00031589\\
1.218	0.00032398\\
1.2223	0.00033226\\
1.2265	0.00034075\\
1.2308	0.00034945\\
1.235	0.00035836\\
1.2393	0.00036749\\
1.2435	0.00037685\\
1.2478	0.00038643\\
1.252	0.00039625\\
1.2563	0.0004063\\
1.2605	0.00041661\\
1.2648	0.00042716\\
1.2691	0.00043797\\
1.2733	0.00044904\\
1.2776	0.00046039\\
1.2818	0.000472\\
1.2861	0.0004839\\
1.2903	0.00049609\\
1.2946	0.00050857\\
1.2988	0.00052136\\
1.3031	0.00053445\\
1.3073	0.00054786\\
1.3116	0.0005616\\
1.3158	0.00057566\\
1.3201	0.00059006\\
1.3243	0.00060481\\
1.3286	0.00061991\\
1.3328	0.00063537\\
1.3371	0.00065121\\
1.3414	0.00066742\\
1.3456	0.00068402\\
1.3499	0.00070102\\
1.3541	0.00071842\\
1.3584	0.00073623\\
1.3626	0.00075448\\
1.3669	0.00077315\\
1.3711	0.00079227\\
1.3754	0.00081184\\
1.3796	0.00083188\\
1.3839	0.00085239\\
1.3881	0.00087338\\
1.3924	0.00089487\\
1.3966	0.00091687\\
1.4009	0.00093939\\
1.4052	0.00096244\\
1.4094	0.00098603\\
1.4137	0.0010102\\
1.4179	0.0010349\\
1.4222	0.0010602\\
1.4264	0.0010861\\
1.4307	0.0011126\\
1.4349	0.0011397\\
1.4392	0.0011674\\
1.4434	0.0011958\\
1.4477	0.0012249\\
1.4519	0.0012546\\
1.4562	0.001285\\
1.4604	0.0013161\\
1.4647	0.001348\\
1.469	0.0013806\\
1.4732	0.0014139\\
1.4775	0.0014481\\
1.4817	0.001483\\
1.486	0.0015187\\
1.4902	0.0015552\\
1.4945	0.0015926\\
1.4987	0.0016308\\
1.503	0.00167\\
1.5072	0.00171\\
1.5115	0.0017509\\
1.5157	0.0017928\\
1.52	0.0018356\\
1.5242	0.0018794\\
1.5285	0.0019243\\
1.5328	0.0019701\\
1.537	0.002017\\
1.5413	0.002065\\
1.5455	0.002114\\
1.5498	0.0021642\\
1.554	0.0022155\\
1.5583	0.0022679\\
1.5625	0.0023216\\
1.5668	0.0023765\\
1.571	0.0024326\\
1.5753	0.0024899\\
1.5795	0.0025486\\
1.5838	0.0026086\\
1.588	0.00267\\
1.5923	0.0027327\\
1.5966	0.0027968\\
1.6008	0.0028624\\
1.6051	0.0029295\\
1.6093	0.002998\\
1.6136	0.0030681\\
1.6178	0.0031397\\
1.6221	0.003213\\
1.6263	0.0032879\\
1.6306	0.0033644\\
1.6348	0.0034427\\
1.6391	0.0035227\\
1.6433	0.0036045\\
1.6476	0.003688\\
1.6518	0.0037735\\
1.6561	0.0038608\\
1.6604	0.0039501\\
1.6646	0.0040413\\
1.6689	0.0041345\\
1.6731	0.0042298\\
1.6774	0.0043272\\
1.6816	0.0044267\\
1.6859	0.0045285\\
1.6901	0.0046324\\
1.6944	0.0047386\\
1.6986	0.0048472\\
1.7029	0.0049581\\
1.7071	0.0050714\\
1.7114	0.0051872\\
1.7156	0.0053056\\
1.7199	0.0054265\\
1.7242	0.00555\\
1.7284	0.0056762\\
1.7327	0.0058051\\
1.7369	0.0059369\\
1.7412	0.0060714\\
1.7454	0.0062089\\
1.7497	0.0063494\\
1.7539	0.0064928\\
1.7582	0.0066394\\
1.7624	0.0067891\\
1.7667	0.006942\\
1.7709	0.0070982\\
1.7752	0.0072577\\
1.7794	0.0074206\\
1.7837	0.0075871\\
1.788	0.007757\\
1.7922	0.0079306\\
1.7965	0.0081079\\
1.8007	0.0082889\\
1.805	0.0084737\\
1.8092	0.0086625\\
1.8135	0.0088553\\
1.8177	0.0090521\\
1.822	0.0092531\\
1.8262	0.0094582\\
1.8305	0.0096678\\
1.8347	0.0098817\\
1.839	0.0101\\
1.8432	0.010323\\
1.8475	0.010551\\
1.8517	0.010783\\
1.856	0.01102\\
1.8603	0.011262\\
1.8645	0.011509\\
1.8688	0.011762\\
1.873	0.012019\\
1.8773	0.012282\\
1.8815	0.01255\\
1.8858	0.012824\\
1.89	0.013103\\
1.8943	0.013389\\
1.8985	0.01368\\
1.9028	0.013976\\
1.907	0.014279\\
1.9113	0.014589\\
1.9155	0.014904\\
1.9198	0.015226\\
1.9241	0.015554\\
1.9283	0.015889\\
1.9326	0.016231\\
1.9368	0.01658\\
1.9411	0.016936\\
1.9453	0.017298\\
1.9496	0.017669\\
1.9538	0.018046\\
1.9581	0.018431\\
1.9623	0.018824\\
1.9666	0.019225\\
1.9708	0.019633\\
1.9751	0.02005\\
1.9793	0.020475\\
1.9836	0.020908\\
1.9879	0.02135\\
1.9921	0.0218\\
1.9964	0.02226\\
2.0006	0.022728\\
2.0049	0.023206\\
2.0091	0.023692\\
2.0134	0.024189\\
2.0176	0.024695\\
2.0219	0.02521\\
2.0261	0.025736\\
2.0304	0.026272\\
2.0346	0.026818\\
2.0389	0.027375\\
2.0431	0.027942\\
2.0474	0.02852\\
2.0517	0.029109\\
2.0559	0.029709\\
2.0602	0.030321\\
2.0644	0.030944\\
2.0687	0.031579\\
2.0729	0.032226\\
2.0772	0.032886\\
2.0814	0.033557\\
2.0857	0.034241\\
2.0899	0.034938\\
2.0942	0.035648\\
2.0984	0.03637\\
2.1027	0.037107\\
2.1069	0.037856\\
2.1112	0.03862\\
2.1155	0.039397\\
2.1197	0.040189\\
2.124	0.040995\\
2.1282	0.041816\\
2.1325	0.042652\\
2.1367	0.043503\\
2.141	0.044369\\
2.1452	0.04525\\
2.1495	0.046148\\
2.1537	0.047061\\
2.158	0.047991\\
2.1622	0.048937\\
2.1665	0.0499\\
2.1707	0.05088\\
2.175	0.051877\\
2.1793	0.052891\\
2.1835	0.053923\\
2.1878	0.054973\\
2.192	0.056042\\
2.1963	0.057128\\
2.2005	0.058234\\
2.2048	0.059358\\
2.209	0.060502\\
2.2133	0.061665\\
2.2175	0.062848\\
2.2218	0.064051\\
2.226	0.065274\\
2.2303	0.066517\\
2.2345	0.067781\\
2.2388	0.069067\\
2.2431	0.070373\\
2.2473	0.071702\\
2.2516	0.073052\\
2.2558	0.074424\\
2.2601	0.075819\\
2.2643	0.077236\\
2.2686	0.078676\\
2.2728	0.08014\\
2.2771	0.081627\\
2.2813	0.083137\\
2.2856	0.084672\\
2.2898	0.086231\\
2.2941	0.087815\\
2.2983	0.089423\\
2.3026	0.091057\\
2.3069	0.092716\\
2.3111	0.094401\\
2.3154	0.096112\\
2.3196	0.097849\\
2.3239	0.099613\\
2.3281	0.1014\\
2.3324	0.10322\\
2.3366	0.10507\\
2.3409	0.10694\\
2.3451	0.10884\\
2.3494	0.11077\\
2.3536	0.11272\\
2.3579	0.11471\\
2.3621	0.11672\\
2.3664	0.11877\\
2.3706	0.12084\\
2.3749	0.12294\\
2.3792	0.12508\\
2.3834	0.12724\\
2.3877	0.12944\\
2.3919	0.13166\\
2.3962	0.13392\\
2.4004	0.1362\\
2.4047	0.13852\\
2.4089	0.14087\\
2.4132	0.14325\\
2.4174	0.14567\\
2.4217	0.14812\\
2.4259	0.15059\\
2.4302	0.15311\\
2.4344	0.15565\\
2.4387	0.15823\\
2.443	0.16084\\
2.4472	0.16349\\
2.4515	0.16617\\
2.4557	0.16888\\
2.46	0.17163\\
2.4642	0.17441\\
2.4685	0.17723\\
2.4727	0.18008\\
2.477	0.18297\\
2.4812	0.18589\\
2.4855	0.18885\\
2.4897	0.19184\\
2.494	0.19487\\
2.4982	0.19793\\
2.5025	0.20103\\
2.5068	0.20417\\
2.511	0.20734\\
2.5153	0.21054\\
2.5195	0.21379\\
2.5238	0.21706\\
2.528	0.22038\\
2.5323	0.22373\\
2.5365	0.22712\\
2.5408	0.23054\\
2.545	0.234\\
2.5493	0.2375\\
2.5535	0.24103\\
2.5578	0.2446\\
2.562	0.2482\\
2.5663	0.25184\\
2.5706	0.25552\\
2.5748	0.25923\\
2.5791	0.26298\\
2.5833	0.26676\\
2.5876	0.27058\\
2.5918	0.27444\\
2.5961	0.27833\\
2.6003	0.28226\\
2.6046	0.28622\\
2.6088	0.29021\\
2.6131	0.29425\\
2.6173	0.29831\\
2.6216	0.30241\\
2.6258	0.30655\\
2.6301	0.31072\\
2.6344	0.31492\\
2.6386	0.31916\\
2.6429	0.32343\\
2.6471	0.32773\\
2.6514	0.33207\\
2.6556	0.33644\\
2.6599	0.34084\\
2.6641	0.34527\\
2.6684	0.34974\\
2.6726	0.35423\\
2.6769	0.35876\\
2.6811	0.36332\\
2.6854	0.3679\\
2.6896	0.37252\\
2.6939	0.37716\\
2.6982	0.38184\\
2.7024	0.38654\\
2.7067	0.39127\\
2.7109	0.39603\\
2.7152	0.40081\\
2.7194	0.40562\\
2.7237	0.41046\\
2.7279	0.41532\\
2.7322	0.4202\\
2.7364	0.42511\\
2.7407	0.43004\\
2.7449	0.435\\
2.7492	0.43998\\
2.7534	0.44498\\
2.7577	0.45\\
2.762	0.45504\\
2.7662	0.4601\\
2.7705	0.46518\\
2.7747	0.47028\\
2.779	0.47539\\
2.7832	0.48052\\
2.7875	0.48567\\
2.7917	0.49083\\
2.796	0.49601\\
2.8002	0.5012\\
2.8045	0.50641\\
2.8087	0.51162\\
2.813	0.51685\\
2.8172	0.52209\\
2.8215	0.52733\\
2.8258	0.53259\\
2.83	0.53785\\
2.8343	0.54312\\
2.8385	0.5484\\
2.8428	0.55368\\
2.847	0.55897\\
2.8513	0.56425\\
2.8555	0.56954\\
2.8598	0.57484\\
2.864	0.58013\\
2.8683	0.58542\\
2.8725	0.59071\\
2.8768	0.596\\
2.881	0.60128\\
2.8853	0.60656\\
2.8896	0.61183\\
2.8938	0.6171\\
2.8981	0.62235\\
2.9023	0.6276\\
2.9066	0.63284\\
2.9108	0.63806\\
2.9151	0.64327\\
2.9193	0.64847\\
2.9236	0.65365\\
2.9278	0.65881\\
2.9321	0.66395\\
2.9363	0.66908\\
2.9406	0.67418\\
2.9448	0.67926\\
2.9491	0.68431\\
2.9533	0.68934\\
2.9576	0.69433\\
2.9619	0.6993\\
2.9661	0.70423\\
2.9704	0.70913\\
2.9746	0.71399\\
2.9789	0.71881\\
2.9831	0.7236\\
2.9874	0.72834\\
2.9916	0.73303\\
2.9959	0.73768\\
3.0001	0.74228\\
3.0044	0.74683\\
3.0086	0.75132\\
3.0129	0.75576\\
3.0171	0.76014\\
3.0214	0.76446\\
3.0257	0.76872\\
3.0299	0.77292\\
3.0342	0.77705\\
3.0384	0.78112\\
3.0427	0.78511\\
3.0469	0.78904\\
3.0512	0.79289\\
3.0554	0.79667\\
3.0597	0.80038\\
3.0639	0.80401\\
3.0682	0.80756\\
3.0724	0.81104\\
3.0767	0.81444\\
3.0809	0.81776\\
3.0852	0.82101\\
3.0895	0.82417\\
3.0937	0.82726\\
3.098	0.83027\\
3.1022	0.83321\\
3.1065	0.83608\\
3.1107	0.83887\\
3.115	0.84159\\
3.1192	0.84424\\
3.1235	0.84682\\
3.1277	0.84934\\
3.132	0.8518\\
3.1362	0.85419\\
3.1405	0.85653\\
3.1447	0.85881\\
3.149	0.86103\\
3.1533	0.86319\\
3.1575	0.86529\\
3.1618	0.86733\\
3.166	0.86929\\
3.1703	0.87119\\
3.1745	0.873\\
3.1788	0.87472\\
3.183	0.87634\\
3.1873	0.87786\\
3.1915	0.87927\\
3.1958	0.88057\\
3.2	0.88174\\
3.2043	0.88278\\
3.2085	0.88368\\
3.2128	0.88445\\
3.2171	0.88509\\
3.2213	0.88559\\
3.2256	0.88596\\
3.2298	0.8862\\
3.2341	0.88631\\
3.2383	0.88631\\
3.2426	0.88619\\
3.2468	0.88597\\
3.2511	0.88565\\
3.2553	0.88525\\
3.2596	0.88476\\
3.2638	0.8842\\
3.2681	0.88357\\
3.2723	0.88286\\
3.2766	0.8821\\
3.2809	0.88127\\
3.2851	0.88038\\
3.2894	0.87941\\
3.2936	0.87838\\
3.2979	0.87728\\
3.3021	0.8761\\
3.3064	0.87485\\
3.3106	0.8735\\
3.3149	0.87207\\
3.3191	0.87055\\
3.3234	0.86893\\
3.3276	0.86721\\
3.3319	0.8654\\
3.3361	0.86348\\
3.3404	0.86146\\
3.3447	0.85933\\
3.3489	0.8571\\
3.3532	0.85477\\
3.3574	0.85233\\
3.3617	0.8498\\
3.3659	0.84717\\
3.3702	0.84444\\
3.3744	0.84161\\
3.3787	0.8387\\
3.3829	0.8357\\
3.3872	0.83262\\
3.3914	0.82947\\
3.3957	0.82623\\
3.3999	0.82293\\
3.4042	0.81956\\
3.4085	0.81613\\
3.4127	0.81263\\
3.417	0.80909\\
3.4212	0.80549\\
3.4255	0.80184\\
3.4297	0.79814\\
3.434	0.79439\\
3.4382	0.79061\\
3.4425	0.78678\\
3.4467	0.78291\\
3.451	0.779\\
3.4552	0.77505\\
3.4595	0.77106\\
3.4637	0.76703\\
3.468	0.76297\\
3.4722	0.75886\\
3.4765	0.75472\\
3.4808	0.75055\\
3.485	0.74634\\
3.4893	0.74209\\
3.4935	0.73782\\
3.4978	0.73351\\
3.502	0.72917\\
3.5063	0.7248\\
3.5105	0.7204\\
3.5148	0.71598\\
3.519	0.71153\\
3.5233	0.70706\\
3.5275	0.70257\\
3.5318	0.69806\\
3.536	0.69353\\
3.5403	0.68899\\
3.5446	0.68444\\
3.5488	0.67987\\
3.5531	0.6753\\
3.5573	0.67072\\
3.5616	0.66613\\
3.5658	0.66154\\
3.5701	0.65695\\
3.5743	0.65235\\
3.5786	0.64776\\
3.5828	0.64316\\
3.5871	0.63857\\
3.5913	0.63397\\
3.5956	0.62939\\
3.5998	0.6248\\
3.6041	0.62022\\
3.6084	0.61565\\
3.6126	0.61108\\
3.6169	0.60652\\
3.6211	0.60196\\
3.6254	0.59741\\
3.6296	0.59287\\
3.6339	0.58833\\
3.6381	0.5838\\
3.6424	0.57928\\
3.6466	0.57477\\
3.6509	0.57027\\
3.6551	0.56577\\
3.6594	0.56129\\
3.6636	0.55681\\
3.6679	0.55235\\
3.6722	0.54789\\
3.6764	0.54345\\
3.6807	0.53902\\
3.6849	0.5346\\
3.6892	0.53019\\
3.6934	0.52579\\
3.6977	0.52141\\
3.7019	0.51703\\
3.7062	0.51268\\
3.7104	0.50833\\
3.7147	0.504\\
3.7189	0.49969\\
3.7232	0.49539\\
3.7274	0.4911\\
3.7317	0.48684\\
3.736	0.48258\\
3.7402	0.47835\\
3.7445	0.47413\\
3.7487	0.46993\\
3.753	0.46574\\
3.7572	0.46158\\
3.7615	0.45743\\
3.7657	0.4533\\
3.77	0.44919\\
3.7742	0.4451\\
3.7785	0.44103\\
3.7827	0.43698\\
3.787	0.43295\\
3.7912	0.42894\\
3.7955	0.42495\\
3.7998	0.42098\\
3.804	0.41703\\
3.8083	0.4131\\
3.8125	0.4092\\
3.8168	0.40532\\
3.821	0.40146\\
3.8253	0.39762\\
3.8295	0.39381\\
3.8338	0.39001\\
3.838	0.38625\\
3.8423	0.3825\\
3.8465	0.37878\\
3.8508	0.37508\\
3.855	0.37141\\
3.8593	0.36775\\
3.8636	0.36413\\
3.8678	0.36053\\
3.8721	0.35695\\
3.8763	0.35339\\
3.8806	0.34986\\
3.8848	0.34636\\
3.8891	0.34288\\
3.8933	0.33942\\
3.8976	0.33599\\
3.9018	0.33258\\
3.9061	0.3292\\
3.9103	0.32585\\
3.9146	0.32251\\
3.9188	0.31921\\
3.9231	0.31593\\
3.9274	0.31267\\
3.9316	0.30944\\
3.9359	0.30623\\
3.9401	0.30305\\
3.9444	0.29989\\
3.9486	0.29676\\
3.9529	0.29365\\
3.9571	0.29057\\
3.9614	0.28751\\
3.9656	0.28448\\
3.9699	0.28147\\
3.9741	0.27849\\
3.9784	0.27553\\
3.9826	0.2726\\
3.9869	0.26969\\
3.9912	0.2668\\
3.9954	0.26394\\
3.9997	0.26111\\
4.0039	0.2583\\
4.0082	0.25551\\
4.0124	0.25275\\
4.0167	0.25001\\
4.0209	0.24729\\
4.0252	0.2446\\
4.0294	0.24193\\
4.0337	0.23929\\
4.0379	0.23667\\
4.0422	0.23407\\
4.0464	0.2315\\
4.0507	0.22894\\
4.0549	0.22642\\
4.0592	0.22391\\
4.0635	0.22143\\
4.0677	0.21897\\
4.072	0.21653\\
4.0762	0.21412\\
4.0805	0.21173\\
4.0847	0.20936\\
4.089	0.20701\\
4.0932	0.20469\\
4.0975	0.20238\\
4.1017	0.2001\\
4.106	0.19784\\
4.1102	0.1956\\
4.1145	0.19338\\
4.1187	0.19119\\
4.123	0.18901\\
4.1273	0.18685\\
4.1315	0.18472\\
4.1358	0.18261\\
4.14	0.18051\\
4.1443	0.17844\\
4.1485	0.17639\\
4.1528	0.17436\\
4.157	0.17234\\
4.1613	0.17035\\
4.1655	0.16838\\
4.1698	0.16642\\
4.174	0.16449\\
4.1783	0.16257\\
4.1825	0.16067\\
4.1868	0.1588\\
4.1911	0.15694\\
4.1953	0.1551\\
4.1996	0.15327\\
4.2038	0.15147\\
4.2081	0.14968\\
4.2123	0.14791\\
4.2166	0.14616\\
4.2208	0.14443\\
4.2251	0.14272\\
4.2293	0.14102\\
4.2336	0.13934\\
4.2378	0.13768\\
4.2421	0.13603\\
4.2463	0.1344\\
4.2506	0.13279\\
4.2549	0.13119\\
4.2591	0.12961\\
4.2634	0.12805\\
4.2676	0.1265\\
4.2719	0.12497\\
4.2761	0.12345\\
4.2804	0.12196\\
4.2846	0.12047\\
4.2889	0.119\\
4.2931	0.11755\\
4.2974	0.11611\\
4.3016	0.11469\\
4.3059	0.11328\\
4.3101	0.11189\\
4.3144	0.11051\\
4.3187	0.10914\\
4.3229	0.10779\\
4.3272	0.10646\\
4.3314	0.10514\\
4.3357	0.10383\\
4.3399	0.10254\\
4.3442	0.10126\\
4.3484	0.099992\\
4.3527	0.09874\\
4.3569	0.097501\\
4.3612	0.096276\\
4.3654	0.095063\\
4.3697	0.093864\\
4.3739	0.092677\\
4.3782	0.091504\\
4.3825	0.090343\\
4.3867	0.089195\\
4.391	0.088059\\
4.3952	0.086935\\
4.3995	0.085824\\
4.4037	0.084725\\
4.408	0.083638\\
4.4122	0.082563\\
4.4165	0.081499\\
4.4207	0.080447\\
4.425	0.079407\\
4.4292	0.078378\\
4.4335	0.077361\\
4.4377	0.076354\\
4.442	0.075359\\
4.4463	0.074375\\
4.4505	0.073402\\
4.4548	0.072439\\
4.459	0.071488\\
4.4633	0.070546\\
4.4675	0.069616\\
4.4718	0.068695\\
4.476	0.067785\\
4.4803	0.066885\\
4.4845	0.065995\\
4.4888	0.065115\\
4.493	0.064245\\
4.4973	0.063385\\
4.5015	0.062534\\
4.5058	0.061693\\
4.5101	0.060861\\
4.5143	0.060038\\
4.5186	0.059225\\
4.5228	0.058421\\
4.5271	0.057626\\
4.5313	0.05684\\
4.5356	0.056063\\
4.5398	0.055294\\
4.5441	0.054535\\
4.5483	0.053783\\
4.5526	0.053041\\
4.5568	0.052306\\
4.5611	0.05158\\
4.5653	0.050863\\
4.5696	0.050153\\
4.5738	0.049452\\
4.5781	0.048758\\
4.5824	0.048072\\
4.5866	0.047394\\
4.5909	0.046724\\
4.5951	0.046062\\
4.5994	0.045407\\
4.6036	0.044759\\
4.6079	0.044119\\
4.6121	0.043487\\
4.6164	0.042861\\
4.6206	0.042243\\
4.6249	0.041632\\
4.6291	0.041028\\
4.6334	0.040431\\
4.6376	0.039841\\
4.6419	0.039257\\
4.6462	0.038681\\
4.6504	0.038111\\
4.6547	0.037548\\
4.6589	0.036992\\
4.6632	0.036442\\
4.6674	0.035898\\
4.6717	0.035361\\
4.6759	0.03483\\
4.6802	0.034306\\
4.6844	0.033788\\
4.6887	0.033276\\
4.6929	0.03277\\
4.6972	0.03227\\
4.7014	0.031776\\
4.7057	0.031288\\
4.71	0.030806\\
4.7142	0.03033\\
4.7185	0.02986\\
4.7227	0.029395\\
4.727	0.028937\\
4.7312	0.028483\\
4.7355	0.028036\\
4.7397	0.027594\\
4.744	0.027158\\
4.7482	0.026727\\
4.7525	0.026301\\
4.7567	0.025881\\
4.761	0.025466\\
4.7652	0.025056\\
4.7695	0.024652\\
4.7738	0.024253\\
4.778	0.023859\\
4.7823	0.02347\\
4.7865	0.023086\\
4.7908	0.022707\\
4.795	0.022333\\
4.7993	0.021964\\
4.8035	0.0216\\
4.8078	0.021241\\
4.812	0.020887\\
4.8163	0.020537\\
4.8205	0.020192\\
4.8248	0.019852\\
4.829	0.019516\\
4.8333	0.019185\\
4.8376	0.018859\\
4.8418	0.018537\\
4.8461	0.01822\\
4.8503	0.017907\\
4.8546	0.017598\\
4.8588	0.017294\\
4.8631	0.016994\\
4.8673	0.016698\\
4.8716	0.016407\\
4.8758	0.016119\\
4.8801	0.015836\\
4.8843	0.015557\\
4.8886	0.015282\\
4.8928	0.015011\\
4.8971	0.014744\\
4.9014	0.014481\\
4.9056	0.014222\\
4.9099	0.013967\\
4.9141	0.013716\\
4.9184	0.013468\\
4.9226	0.013224\\
4.9269	0.012984\\
4.9311	0.012748\\
4.9354	0.012515\\
4.9396	0.012286\\
4.9439	0.01206\\
4.9481	0.011838\\
4.9524	0.011619\\
4.9566	0.011404\\
4.9609	0.011192\\
4.9652	0.010984\\
4.9694	0.010778\\
4.9737	0.010577\\
4.9779	0.010378\\
4.9822	0.010182\\
4.9864	0.0099902\\
4.9907	0.009801\\
4.9949	0.009615\\
4.9992	0.009432\\
5.0034	0.009252\\
5.0077	0.0090749\\
5.0119	0.0089009\\
5.0162	0.0087297\\
5.0204	0.0085614\\
5.0247	0.0083959\\
5.029	0.0082332\\
5.0332	0.0080732\\
5.0375	0.007916\\
5.0417	0.0077615\\
5.046	0.0076095\\
5.0502	0.0074602\\
5.0545	0.0073135\\
5.0587	0.0071693\\
5.063	0.0070277\\
5.0672	0.0068885\\
5.0715	0.0067517\\
5.0757	0.0066173\\
5.08	0.0064853\\
5.0842	0.0063556\\
5.0885	0.0062282\\
5.0928	0.0061031\\
5.097	0.0059802\\
5.1013	0.0058595\\
5.1055	0.005741\\
5.1098	0.0056246\\
5.114	0.0055103\\
5.1183	0.0053981\\
5.1225	0.005288\\
5.1268	0.0051798\\
5.131	0.0050736\\
5.1353	0.0049694\\
5.1395	0.0048671\\
5.1438	0.0047667\\
5.148	0.0046681\\
5.1523	0.0045714\\
5.1565	0.0044764\\
5.1608	0.0043833\\
5.1651	0.0042919\\
5.1693	0.0042022\\
5.1736	0.0041142\\
5.1778	0.0040279\\
5.1821	0.0039432\\
5.1863	0.0038601\\
5.1906	0.0037786\\
5.1948	0.0036986\\
5.1991	0.0036202\\
5.2033	0.0035433\\
5.2076	0.0034679\\
5.2118	0.003394\\
5.2161	0.0033214\\
5.2203	0.0032503\\
5.2246	0.0031806\\
5.2289	0.0031122\\
5.2331	0.0030452\\
5.2374	0.0029795\\
5.2416	0.0029151\\
5.2459	0.0028519\\
5.2501	0.00279\\
5.2544	0.0027294\\
5.2586	0.0026699\\
5.2629	0.0026116\\
5.2671	0.0025545\\
5.2714	0.0024986\\
5.2756	0.0024437\\
5.2799	0.00239\\
5.2841	0.0023373\\
5.2884	0.0022858\\
5.2927	0.0022352\\
5.2969	0.0021857\\
5.3012	0.0021372\\
5.3054	0.0020897\\
5.3097	0.0020432\\
5.3139	0.0019976\\
5.3182	0.001953\\
5.3224	0.0019093\\
5.3267	0.0018665\\
5.3309	0.0018246\\
5.3352	0.0017835\\
5.3394	0.0017433\\
5.3437	0.001704\\
5.3479	0.0016654\\
5.3522	0.0016277\\
5.3565	0.0015908\\
5.3607	0.0015546\\
5.3607	0.0012789\\
5.3565	0.0013108\\
5.3522	0.0013434\\
5.3479	0.0013767\\
5.3437	0.0014108\\
5.3394	0.0014457\\
5.3352	0.0014814\\
5.3309	0.0015179\\
5.3267	0.0015552\\
5.3224	0.0015933\\
5.3182	0.0016323\\
5.3139	0.0016722\\
5.3097	0.001713\\
5.3054	0.0017548\\
5.3012	0.0017974\\
5.2969	0.001841\\
5.2927	0.0018855\\
5.2884	0.0019311\\
5.2841	0.0019776\\
5.2799	0.0020252\\
5.2756	0.0020738\\
5.2714	0.0021235\\
5.2671	0.0021743\\
5.2629	0.0022261\\
5.2586	0.0022791\\
5.2544	0.0023333\\
5.2501	0.0023886\\
5.2459	0.0024451\\
5.2416	0.0025028\\
5.2374	0.0025617\\
5.2331	0.0026219\\
5.2289	0.0026833\\
5.2246	0.0027461\\
5.2203	0.0028102\\
5.2161	0.0028756\\
5.2118	0.0029424\\
5.2076	0.0030106\\
5.2033	0.0030802\\
5.1991	0.0031512\\
5.1948	0.0032238\\
5.1906	0.0032978\\
5.1863	0.0033733\\
5.1821	0.0034503\\
5.1778	0.0035289\\
5.1736	0.0036092\\
5.1693	0.003691\\
5.1651	0.0037745\\
5.1608	0.0038596\\
5.1565	0.0039465\\
5.1523	0.004035\\
5.148	0.0041254\\
5.1438	0.0042175\\
5.1395	0.0043114\\
5.1353	0.0044071\\
5.131	0.0045047\\
5.1268	0.0046042\\
5.1225	0.0047057\\
5.1183	0.004809\\
5.114	0.0049144\\
5.1098	0.0050218\\
5.1055	0.0051312\\
5.1013	0.0052427\\
5.097	0.0053562\\
5.0928	0.005472\\
5.0885	0.0055898\\
5.0842	0.0057099\\
5.08	0.0058322\\
5.0757	0.0059567\\
5.0715	0.0060836\\
5.0672	0.0062127\\
5.063	0.0063442\\
5.0587	0.0064781\\
5.0545	0.0066144\\
5.0502	0.0067531\\
5.046	0.0068944\\
5.0417	0.0070381\\
5.0375	0.0071844\\
5.0332	0.0073333\\
5.029	0.0074847\\
5.0247	0.0076389\\
5.0204	0.0077957\\
5.0162	0.0079552\\
5.0119	0.0081175\\
5.0077	0.0082826\\
5.0034	0.0084505\\
4.9992	0.0086212\\
4.9949	0.0087949\\
4.9907	0.0089714\\
4.9864	0.009151\\
4.9822	0.0093335\\
4.9779	0.0095191\\
4.9737	0.0097078\\
4.9694	0.0098995\\
4.9652	0.010094\\
4.9609	0.010293\\
4.9566	0.010494\\
4.9524	0.010698\\
4.9481	0.010906\\
4.9439	0.011118\\
4.9396	0.011332\\
4.9354	0.01155\\
4.9311	0.011772\\
4.9269	0.011997\\
4.9226	0.012225\\
4.9184	0.012457\\
4.9141	0.012693\\
4.9099	0.012933\\
4.9056	0.013176\\
4.9014	0.013422\\
4.8971	0.013673\\
4.8928	0.013927\\
4.8886	0.014186\\
4.8843	0.014448\\
4.8801	0.014714\\
4.8758	0.014984\\
4.8716	0.015258\\
4.8673	0.015537\\
4.8631	0.015819\\
4.8588	0.016106\\
4.8546	0.016397\\
4.8503	0.016692\\
4.8461	0.016991\\
4.8418	0.017295\\
4.8376	0.017603\\
4.8333	0.017915\\
4.829	0.018232\\
4.8248	0.018554\\
4.8205	0.01888\\
4.8163	0.019211\\
4.812	0.019546\\
4.8078	0.019886\\
4.8035	0.020231\\
4.7993	0.020581\\
4.795	0.020935\\
4.7908	0.021295\\
4.7865	0.021659\\
4.7823	0.022028\\
4.778	0.022403\\
4.7738	0.022782\\
4.7695	0.023167\\
4.7652	0.023557\\
4.761	0.023952\\
4.7567	0.024353\\
4.7525	0.024759\\
4.7482	0.02517\\
4.744	0.025586\\
4.7397	0.026009\\
4.7355	0.026436\\
4.7312	0.02687\\
4.727	0.027309\\
4.7227	0.027753\\
4.7185	0.028204\\
4.7142	0.02866\\
4.71	0.029122\\
4.7057	0.02959\\
4.7014	0.030064\\
4.6972	0.030544\\
4.6929	0.03103\\
4.6887	0.031522\\
4.6844	0.032021\\
4.6802	0.032525\\
4.6759	0.033036\\
4.6717	0.033553\\
4.6674	0.034077\\
4.6632	0.034607\\
4.6589	0.035144\\
4.6547	0.035687\\
4.6504	0.036237\\
4.6462	0.036793\\
4.6419	0.037357\\
4.6376	0.037927\\
4.6334	0.038503\\
4.6291	0.039087\\
4.6249	0.039678\\
4.6206	0.040275\\
4.6164	0.04088\\
4.6121	0.041492\\
4.6079	0.042111\\
4.6036	0.042738\\
4.5994	0.043371\\
4.5951	0.044012\\
4.5909	0.044661\\
4.5866	0.045317\\
4.5824	0.04598\\
4.5781	0.046651\\
4.5738	0.04733\\
4.5696	0.048017\\
4.5653	0.048711\\
4.5611	0.049413\\
4.5568	0.050123\\
4.5526	0.050841\\
4.5483	0.051568\\
4.5441	0.052302\\
4.5398	0.053045\\
4.5356	0.053795\\
4.5313	0.054555\\
4.5271	0.055322\\
4.5228	0.056098\\
4.5186	0.056883\\
4.5143	0.057676\\
4.5101	0.058478\\
4.5058	0.059289\\
4.5015	0.060108\\
4.4973	0.060937\\
4.493	0.061775\\
4.4888	0.062621\\
4.4845	0.063477\\
4.4803	0.064343\\
4.476	0.065217\\
4.4718	0.066101\\
4.4675	0.066995\\
4.4633	0.067898\\
4.459	0.068811\\
4.4548	0.069734\\
4.4505	0.070667\\
4.4463	0.07161\\
4.442	0.072563\\
4.4377	0.073526\\
4.4335	0.0745\\
4.4292	0.075484\\
4.425	0.076478\\
4.4207	0.077484\\
4.4165	0.0785\\
4.4122	0.079527\\
4.408	0.080565\\
4.4037	0.081614\\
4.3995	0.082674\\
4.3952	0.083746\\
4.391	0.084829\\
4.3867	0.085924\\
4.3825	0.087031\\
4.3782	0.088149\\
4.3739	0.08928\\
4.3697	0.090423\\
4.3654	0.091578\\
4.3612	0.092745\\
4.3569	0.093925\\
4.3527	0.095118\\
4.3484	0.096323\\
4.3442	0.097542\\
4.3399	0.098773\\
4.3357	0.10002\\
4.3314	0.10128\\
4.3272	0.10255\\
4.3229	0.10383\\
4.3187	0.10513\\
4.3144	0.10645\\
4.3101	0.10777\\
4.3059	0.10912\\
4.3016	0.11047\\
4.2974	0.11184\\
4.2931	0.11323\\
4.2889	0.11463\\
4.2846	0.11604\\
4.2804	0.11747\\
4.2761	0.11892\\
4.2719	0.12038\\
4.2676	0.12186\\
4.2634	0.12335\\
4.2591	0.12486\\
4.2549	0.12639\\
4.2506	0.12793\\
4.2463	0.12949\\
4.2421	0.13106\\
4.2378	0.13265\\
4.2336	0.13426\\
4.2293	0.13589\\
4.2251	0.13753\\
4.2208	0.13919\\
4.2166	0.14087\\
4.2123	0.14257\\
4.2081	0.14428\\
4.2038	0.14601\\
4.1996	0.14776\\
4.1953	0.14953\\
4.1911	0.15132\\
4.1868	0.15313\\
4.1825	0.15495\\
4.1783	0.1568\\
4.174	0.15866\\
4.1698	0.16054\\
4.1655	0.16244\\
4.1613	0.16437\\
4.157	0.16631\\
4.1528	0.16827\\
4.1485	0.17025\\
4.1443	0.17226\\
4.14	0.17428\\
4.1358	0.17632\\
4.1315	0.17839\\
4.1273	0.18047\\
4.123	0.18258\\
4.1187	0.18471\\
4.1145	0.18686\\
4.1102	0.18903\\
4.106	0.19122\\
4.1017	0.19343\\
4.0975	0.19567\\
4.0932	0.19793\\
4.089	0.20021\\
4.0847	0.20251\\
4.0805	0.20484\\
4.0762	0.20718\\
4.072	0.20955\\
4.0677	0.21195\\
4.0635	0.21436\\
4.0592	0.2168\\
4.0549	0.21926\\
4.0507	0.22175\\
4.0464	0.22426\\
4.0422	0.22679\\
4.0379	0.22934\\
4.0337	0.23192\\
4.0294	0.23452\\
4.0252	0.23715\\
4.0209	0.2398\\
4.0167	0.24247\\
4.0124	0.24517\\
4.0082	0.24789\\
4.0039	0.25064\\
3.9997	0.25341\\
3.9954	0.2562\\
3.9912	0.25902\\
3.9869	0.26186\\
3.9826	0.26473\\
3.9784	0.26762\\
3.9741	0.27053\\
3.9699	0.27347\\
3.9656	0.27644\\
3.9614	0.27943\\
3.9571	0.28244\\
3.9529	0.28548\\
3.9486	0.28854\\
3.9444	0.29162\\
3.9401	0.29473\\
3.9359	0.29787\\
3.9316	0.30102\\
3.9274	0.30421\\
3.9231	0.30741\\
3.9188	0.31064\\
3.9146	0.3139\\
3.9103	0.31718\\
3.9061	0.32048\\
3.9018	0.3238\\
3.8976	0.32715\\
3.8933	0.33053\\
3.8891	0.33393\\
3.8848	0.33735\\
3.8806	0.34079\\
3.8763	0.34426\\
3.8721	0.34775\\
3.8678	0.35126\\
3.8636	0.3548\\
3.8593	0.35835\\
3.855	0.36194\\
3.8508	0.36554\\
3.8465	0.36916\\
3.8423	0.37281\\
3.838	0.37648\\
3.8338	0.38017\\
3.8295	0.38388\\
3.8253	0.38762\\
3.821	0.39137\\
3.8168	0.39515\\
3.8125	0.39894\\
3.8083	0.40276\\
3.804	0.4066\\
3.7998	0.41045\\
3.7955	0.41433\\
3.7912	0.41822\\
3.787	0.42214\\
3.7827	0.42607\\
3.7785	0.43002\\
3.7742	0.43399\\
3.77	0.43798\\
3.7657	0.44199\\
3.7615	0.44601\\
3.7572	0.45005\\
3.753	0.45411\\
3.7487	0.45819\\
3.7445	0.46228\\
3.7402	0.46639\\
3.736	0.47051\\
3.7317	0.47465\\
3.7274	0.47881\\
3.7232	0.48298\\
3.7189	0.48716\\
3.7147	0.49136\\
3.7104	0.49558\\
3.7062	0.49981\\
3.7019	0.50405\\
3.6977	0.5083\\
3.6934	0.51257\\
3.6892	0.51685\\
3.6849	0.52115\\
3.6807	0.52546\\
3.6764	0.52978\\
3.6722	0.53411\\
3.6679	0.53845\\
3.6636	0.5428\\
3.6594	0.54717\\
3.6551	0.55154\\
3.6509	0.55592\\
3.6466	0.56032\\
3.6424	0.56472\\
3.6381	0.56913\\
3.6339	0.57354\\
3.6296	0.57796\\
3.6254	0.58239\\
3.6211	0.58683\\
3.6169	0.59126\\
3.6126	0.5957\\
3.6084	0.60014\\
3.6041	0.60459\\
3.5998	0.60903\\
3.5956	0.61347\\
3.5913	0.61791\\
3.5871	0.62234\\
3.5828	0.62678\\
3.5786	0.6312\\
3.5743	0.63562\\
3.5701	0.64003\\
3.5658	0.64444\\
3.5616	0.64883\\
3.5573	0.65322\\
3.5531	0.65759\\
3.5488	0.66195\\
3.5446	0.6663\\
3.5403	0.67064\\
3.536	0.67496\\
3.5318	0.67926\\
3.5275	0.68356\\
3.5233	0.68783\\
3.519	0.6921\\
3.5148	0.69634\\
3.5105	0.70056\\
3.5063	0.70477\\
3.502	0.70896\\
3.4978	0.71313\\
3.4935	0.71727\\
3.4893	0.72139\\
3.485	0.72549\\
3.4808	0.72955\\
3.4765	0.73359\\
3.4722	0.7376\\
3.468	0.74157\\
3.4637	0.74551\\
3.4595	0.74941\\
3.4552	0.75327\\
3.451	0.75708\\
3.4467	0.76084\\
3.4425	0.76456\\
3.4382	0.76822\\
3.434	0.77182\\
3.4297	0.77536\\
3.4255	0.77885\\
3.4212	0.78227\\
3.417	0.78562\\
3.4127	0.78891\\
3.4085	0.79213\\
3.4042	0.79528\\
3.3999	0.79836\\
3.3957	0.80137\\
3.3914	0.80431\\
3.3872	0.80719\\
3.3829	0.80999\\
3.3787	0.81272\\
3.3744	0.81539\\
3.3702	0.81799\\
3.3659	0.82052\\
3.3617	0.82299\\
3.3574	0.82539\\
3.3532	0.82773\\
3.3489	0.83001\\
3.3447	0.83222\\
3.3404	0.83437\\
3.3361	0.83646\\
3.3319	0.83848\\
3.3276	0.84044\\
3.3234	0.84233\\
3.3191	0.84415\\
3.3149	0.84591\\
3.3106	0.84759\\
3.3064	0.84919\\
3.3021	0.85071\\
3.2979	0.85215\\
3.2936	0.8535\\
3.2894	0.85475\\
3.2851	0.8559\\
3.2809	0.85695\\
3.2766	0.85788\\
3.2723	0.85871\\
3.2681	0.85942\\
3.2638	0.86001\\
3.2596	0.86048\\
3.2553	0.86083\\
3.2511	0.86107\\
3.2468	0.8612\\
3.2426	0.86121\\
3.2383	0.86112\\
3.2341	0.86093\\
3.2298	0.86063\\
3.2256	0.86024\\
3.2213	0.85976\\
3.2171	0.85918\\
3.2128	0.85852\\
3.2085	0.85776\\
3.2043	0.85692\\
3.2	0.85598\\
3.1958	0.85495\\
3.1915	0.85383\\
3.1873	0.85262\\
3.183	0.8513\\
3.1788	0.84989\\
3.1745	0.84837\\
3.1703	0.84675\\
3.166	0.84502\\
3.1618	0.84319\\
3.1575	0.84124\\
3.1533	0.8392\\
3.149	0.83705\\
3.1447	0.8348\\
3.1405	0.83245\\
3.1362	0.83001\\
3.132	0.82748\\
3.1277	0.82487\\
3.1235	0.82218\\
3.1192	0.81941\\
3.115	0.81658\\
3.1107	0.81367\\
3.1065	0.81071\\
3.1022	0.80768\\
3.098	0.80459\\
3.0937	0.80145\\
3.0895	0.79825\\
3.0852	0.795\\
3.0809	0.79169\\
3.0767	0.78832\\
3.0724	0.78491\\
3.0682	0.78144\\
3.0639	0.77792\\
3.0597	0.77434\\
3.0554	0.77071\\
3.0512	0.76702\\
3.0469	0.76328\\
3.0427	0.75949\\
3.0384	0.75564\\
3.0342	0.75174\\
3.0299	0.74778\\
3.0257	0.74378\\
3.0214	0.73972\\
3.0171	0.73561\\
3.0129	0.73144\\
3.0086	0.72723\\
3.0044	0.72296\\
3.0001	0.71865\\
2.9959	0.71429\\
2.9916	0.70988\\
2.9874	0.70543\\
2.9831	0.70093\\
2.9789	0.69638\\
2.9746	0.6918\\
2.9704	0.68717\\
2.9661	0.6825\\
2.9619	0.67779\\
2.9576	0.67305\\
2.9533	0.66827\\
2.9491	0.66345\\
2.9448	0.6586\\
2.9406	0.65372\\
2.9363	0.64881\\
2.9321	0.64387\\
2.9278	0.63891\\
2.9236	0.63392\\
2.9193	0.6289\\
2.9151	0.62386\\
2.9108	0.6188\\
2.9066	0.61372\\
2.9023	0.60863\\
2.8981	0.60352\\
2.8938	0.59839\\
2.8896	0.59325\\
2.8853	0.5881\\
2.881	0.58294\\
2.8768	0.57777\\
2.8725	0.57259\\
2.8683	0.56741\\
2.864	0.56222\\
2.8598	0.55703\\
2.8555	0.55184\\
2.8513	0.54665\\
2.847	0.54146\\
2.8428	0.53627\\
2.8385	0.53108\\
2.8343	0.5259\\
2.83	0.52072\\
2.8258	0.51555\\
2.8215	0.51039\\
2.8172	0.50524\\
2.813	0.5001\\
2.8087	0.49496\\
2.8045	0.48984\\
2.8002	0.48473\\
2.796	0.47964\\
2.7917	0.47456\\
2.7875	0.46949\\
2.7832	0.46444\\
2.779	0.4594\\
2.7747	0.45439\\
2.7705	0.44939\\
2.7662	0.44441\\
2.762	0.43945\\
2.7577	0.43451\\
2.7534	0.42959\\
2.7492	0.42469\\
2.7449	0.41982\\
2.7407	0.41497\\
2.7364	0.41014\\
2.7322	0.40533\\
2.7279	0.40055\\
2.7237	0.3958\\
2.7194	0.39107\\
2.7152	0.38637\\
2.7109	0.38169\\
2.7067	0.37704\\
2.7024	0.37242\\
2.6982	0.36783\\
2.6939	0.36326\\
2.6896	0.35873\\
2.6854	0.35422\\
2.6811	0.34974\\
2.6769	0.3453\\
2.6726	0.34088\\
2.6684	0.3365\\
2.6641	0.33214\\
2.6599	0.32782\\
2.6556	0.32353\\
2.6514	0.31927\\
2.6471	0.31504\\
2.6429	0.31085\\
2.6386	0.30669\\
2.6344	0.30256\\
2.6301	0.29847\\
2.6258	0.29441\\
2.6216	0.29038\\
2.6173	0.28638\\
2.6131	0.28242\\
2.6088	0.2785\\
2.6046	0.27461\\
2.6003	0.27075\\
2.5961	0.26693\\
2.5918	0.26314\\
2.5876	0.25939\\
2.5833	0.25567\\
2.5791	0.25198\\
2.5748	0.24834\\
2.5706	0.24472\\
2.5663	0.24114\\
2.562	0.2376\\
2.5578	0.23409\\
2.5535	0.23062\\
2.5493	0.22718\\
2.545	0.22378\\
2.5408	0.22041\\
2.5365	0.21707\\
2.5323	0.21378\\
2.528	0.21051\\
2.5238	0.20728\\
2.5195	0.20409\\
2.5153	0.20093\\
2.511	0.19781\\
2.5068	0.19472\\
2.5025	0.19167\\
2.4982	0.18865\\
2.494	0.18566\\
2.4897	0.18271\\
2.4855	0.17979\\
2.4812	0.17691\\
2.477	0.17406\\
2.4727	0.17125\\
2.4685	0.16846\\
2.4642	0.16572\\
2.46	0.163\\
2.4557	0.16032\\
2.4515	0.15768\\
2.4472	0.15506\\
2.443	0.15248\\
2.4387	0.14993\\
2.4344	0.14741\\
2.4302	0.14493\\
2.4259	0.14248\\
2.4217	0.14006\\
2.4174	0.13767\\
2.4132	0.13531\\
2.4089	0.13299\\
2.4047	0.13069\\
2.4004	0.12843\\
2.3962	0.1262\\
2.3919	0.124\\
2.3877	0.12183\\
2.3834	0.11968\\
2.3792	0.11757\\
2.3749	0.11549\\
2.3706	0.11344\\
2.3664	0.11142\\
2.3621	0.10942\\
2.3579	0.10746\\
2.3536	0.10552\\
2.3494	0.10361\\
2.3451	0.10173\\
2.3409	0.099882\\
2.3366	0.098058\\
2.3324	0.096261\\
2.3281	0.09449\\
2.3239	0.092747\\
2.3196	0.09103\\
2.3154	0.08934\\
2.3111	0.087675\\
2.3069	0.086037\\
2.3026	0.084424\\
2.2983	0.082836\\
2.2941	0.081273\\
2.2898	0.079735\\
2.2856	0.078221\\
2.2813	0.076732\\
2.2771	0.075267\\
2.2728	0.073825\\
2.2686	0.072407\\
2.2643	0.071012\\
2.2601	0.06964\\
2.2558	0.06829\\
2.2516	0.066963\\
2.2473	0.065658\\
2.2431	0.064375\\
2.2388	0.063114\\
2.2345	0.061873\\
2.2303	0.060654\\
2.226	0.059456\\
2.2218	0.058278\\
2.2175	0.057121\\
2.2133	0.055983\\
2.209	0.054865\\
2.2048	0.053767\\
2.2005	0.052688\\
2.1963	0.051628\\
2.192	0.050587\\
2.1878	0.049564\\
2.1835	0.048559\\
2.1793	0.047572\\
2.175	0.046604\\
2.1707	0.045652\\
2.1665	0.044718\\
2.1622	0.0438\\
2.158	0.0429\\
2.1537	0.042016\\
2.1495	0.041148\\
2.1452	0.040296\\
2.141	0.03946\\
2.1367	0.03864\\
2.1325	0.037834\\
2.1282	0.037044\\
2.124	0.036269\\
2.1197	0.035508\\
2.1155	0.034762\\
2.1112	0.03403\\
2.1069	0.033311\\
2.1027	0.032607\\
2.0984	0.031916\\
2.0942	0.031238\\
2.0899	0.030574\\
2.0857	0.029922\\
2.0814	0.029283\\
2.0772	0.028656\\
2.0729	0.028042\\
2.0687	0.027439\\
2.0644	0.026849\\
2.0602	0.02627\\
2.0559	0.025702\\
2.0517	0.025146\\
2.0474	0.024601\\
2.0431	0.024066\\
2.0389	0.023543\\
2.0346	0.02303\\
2.0304	0.022527\\
2.0261	0.022034\\
2.0219	0.021551\\
2.0176	0.021078\\
2.0134	0.020615\\
2.0091	0.020161\\
2.0049	0.019717\\
2.0006	0.019281\\
1.9964	0.018854\\
1.9921	0.018437\\
1.9879	0.018028\\
1.9836	0.017627\\
1.9793	0.017234\\
1.9751	0.01685\\
1.9708	0.016474\\
1.9666	0.016105\\
1.9623	0.015745\\
1.9581	0.015392\\
1.9538	0.015046\\
1.9496	0.014707\\
1.9453	0.014376\\
1.9411	0.014052\\
1.9368	0.013734\\
1.9326	0.013424\\
1.9283	0.01312\\
1.9241	0.012822\\
1.9198	0.012531\\
1.9155	0.012246\\
1.9113	0.011967\\
1.907	0.011694\\
1.9028	0.011427\\
1.8985	0.011166\\
1.8943	0.01091\\
1.89	0.01066\\
1.8858	0.010416\\
1.8815	0.010176\\
1.8773	0.0099423\\
1.873	0.0097133\\
1.8688	0.0094893\\
1.8645	0.0092703\\
1.8603	0.009056\\
1.856	0.0088464\\
1.8517	0.0086415\\
1.8475	0.008441\\
1.8432	0.008245\\
1.839	0.0080533\\
1.8347	0.0078659\\
1.8305	0.0076826\\
1.8262	0.0075034\\
1.822	0.0073282\\
1.8177	0.0071569\\
1.8135	0.0069894\\
1.8092	0.0068257\\
1.805	0.0066656\\
1.8007	0.0065091\\
1.7965	0.0063561\\
1.7922	0.0062066\\
1.788	0.0060604\\
1.7837	0.0059175\\
1.7794	0.0057779\\
1.7752	0.0056414\\
1.7709	0.005508\\
1.7667	0.0053776\\
1.7624	0.0052501\\
1.7582	0.0051256\\
1.7539	0.0050039\\
1.7497	0.004885\\
1.7454	0.0047688\\
1.7412	0.0046552\\
1.7369	0.0045443\\
1.7327	0.0044358\\
1.7284	0.0043299\\
1.7242	0.0042264\\
1.7199	0.0041252\\
1.7156	0.0040264\\
1.7114	0.0039299\\
1.7071	0.0038356\\
1.7029	0.0037435\\
1.6986	0.0036535\\
1.6944	0.0035656\\
1.6901	0.0034797\\
1.6859	0.0033958\\
1.6816	0.0033138\\
1.6774	0.0032338\\
1.6731	0.0031556\\
1.6689	0.0030792\\
1.6646	0.0030047\\
1.6604	0.0029318\\
1.6561	0.0028607\\
1.6518	0.0027912\\
1.6476	0.0027233\\
1.6433	0.0026571\\
1.6391	0.0025924\\
1.6348	0.0025292\\
1.6306	0.0024675\\
1.6263	0.0024072\\
1.6221	0.0023484\\
1.6178	0.0022909\\
1.6136	0.0022348\\
1.6093	0.00218\\
1.6051	0.0021265\\
1.6008	0.0020743\\
1.5966	0.0020233\\
1.5923	0.0019736\\
1.588	0.001925\\
1.5838	0.0018775\\
1.5795	0.0018312\\
1.5753	0.001786\\
1.571	0.0017419\\
1.5668	0.0016988\\
1.5625	0.0016567\\
1.5583	0.0016157\\
1.554	0.0015756\\
1.5498	0.0015365\\
1.5455	0.0014983\\
1.5413	0.001461\\
1.537	0.0014246\\
1.5328	0.0013891\\
1.5285	0.0013545\\
1.5242	0.0013207\\
1.52	0.0012877\\
1.5157	0.0012555\\
1.5115	0.001224\\
1.5072	0.0011933\\
1.503	0.0011634\\
1.4987	0.0011342\\
1.4945	0.0011057\\
1.4902	0.0010779\\
1.486	0.0010507\\
1.4817	0.0010243\\
1.4775	0.00099841\\
1.4732	0.0009732\\
1.469	0.00094859\\
1.4647	0.00092459\\
1.4604	0.00090117\\
1.4562	0.00087833\\
1.4519	0.00085604\\
1.4477	0.00083429\\
1.4434	0.00081307\\
1.4392	0.00079238\\
1.4349	0.00077219\\
1.4307	0.00075249\\
1.4264	0.00073328\\
1.4222	0.00071454\\
1.4179	0.00069626\\
1.4137	0.00067843\\
1.4094	0.00066104\\
1.4052	0.00064408\\
1.4009	0.00062753\\
1.3966	0.0006114\\
1.3924	0.00059566\\
1.3881	0.00058031\\
1.3839	0.00056534\\
1.3796	0.00055075\\
1.3754	0.00053651\\
1.3711	0.00052263\\
1.3669	0.00050909\\
1.3626	0.00049589\\
1.3584	0.00048302\\
1.3541	0.00047047\\
1.3499	0.00045823\\
1.3456	0.00044629\\
1.3414	0.00043466\\
1.3371	0.00042332\\
1.3328	0.00041226\\
1.3286	0.00040147\\
1.3243	0.00039096\\
1.3201	0.00038071\\
1.3158	0.00037072\\
1.3116	0.00036098\\
1.3073	0.00035149\\
1.3031	0.00034223\\
1.2988	0.00033321\\
1.2946	0.00032442\\
1.2903	0.00031585\\
1.2861	0.0003075\\
1.2818	0.00029935\\
1.2776	0.00029142\\
1.2733	0.00028369\\
1.2691	0.00027615\\
1.2648	0.0002688\\
1.2605	0.00026164\\
1.2563	0.00025467\\
1.252	0.00024787\\
1.2478	0.00024125\\
1.2435	0.00023479\\
1.2393	0.0002285\\
1.235	0.00022238\\
1.2308	0.00021641\\
1.2265	0.00021059\\
1.2223	0.00020492\\
1.218	0.0001994\\
1.2138	0.00019402\\
1.2095	0.00018878\\
1.2053	0.00018368\\
1.201	0.00017871\\
1.1967	0.00017386\\
1.1925	0.00016914\\
1.1882	0.00016455\\
1.184	0.00016007\\
1.1797	0.00015571\\
1.1755	0.00015147\\
1.1712	0.00014733\\
1.167	0.0001433\\
1.1627	0.00013938\\
1.1585	0.00013556\\
1.1542	0.00013184\\
1.15	0.00012822\\
1.1457	0.0001247\\
1.1415	0.00012126\\
1.1372	0.00011792\\
1.1329	0.00011466\\
1.1287	0.00011149\\
1.1244	0.00010841\\
1.1202	0.0001054\\
1.1159	0.00010248\\
1.1117	9.9634e-05\\
}--cycle;
\addplot[ybar interval, fill=mycolor2, fill opacity=0.35, area legend, draw=none] table[row sep=crcr, x=Lower, y=Count] {%
Lower	Upper	Count\\
1.1117	1.4658	0.00037656\\
1.4658	1.8198	0.0030125\\
1.8198	2.1739	0.026208\\
2.1739	2.528	0.11997\\
2.528	2.8821	0.38898\\
2.8821	3.2362	0.76132\\
3.2362	3.5903	0.79295\\
3.5903	3.9444	0.44705\\
3.9444	4.2984	0.19159\\
4.2984	4.6525	0.068081\\
4.6525	5.0066	0.017171\\
5.0066	5.3607	0.0074558\\
5.3607	5.3607	0.0074558\\
};
\addlegendentry{\ApproxHistogramLegend}

\addplot [color=mycolor2, only marks, every error bar/.append style={opacity=0.45}, mark=*, mark size=0pt, draw=none, forget plot]
 plot [error bars/.cd, y dir=both, y explicit, error bar style={line width=1pt, color=mycolor2}, error mark options={mark=none,mark size=0pt}]
 table[row sep=crcr, y error plus index=2, y error minus index=3]{%
1.2887	0.00037656	0.00032732	0.00032732\\
1.6428	0.0030125	0.001083	0.001083\\
1.9969	0.026208	0.0027322	0.0027322\\
2.351	0.11997	0.0050094	0.0050094\\
2.7051	0.38898	0.01049	0.01049\\
3.0591	0.76132	0.013022	0.013022\\
3.4132	0.79295	0.010231	0.010231\\
3.7673	0.44705	0.0094999	0.0094999\\
4.1214	0.19159	0.0059026	0.0059026\\
4.4755	0.068081	0.0031783	0.0031783\\
4.8296	0.017171	0.0013941	0.0013941\\
5.1837	0.0074558	0.00045031	0.00045031\\
};
\addplot [color=mycolor2]
  table[row sep=crcr]{%
1.1117	0.00018611\\
1.3626	0.000803\\
1.5072	0.0017949\\
1.6136	0.0031888\\
1.6944	0.00489\\
1.7624	0.0069661\\
1.822	0.0094493\\
1.873	0.012226\\
1.9198	0.015435\\
1.9623	0.019026\\
2.0006	0.022912\\
2.0389	0.027526\\
2.0729	0.032334\\
2.1027	0.037161\\
2.1325	0.042636\\
2.1622	0.048832\\
2.1878	0.054774\\
2.2133	0.061351\\
2.2388	0.068617\\
2.2643	0.076624\\
2.2898	0.085428\\
2.3154	0.095086\\
2.3366	0.10383\\
2.3579	0.11323\\
2.3792	0.12332\\
2.4004	0.13414\\
2.4217	0.14571\\
2.443	0.15806\\
2.4642	0.17122\\
2.4855	0.18519\\
2.5068	0.2\\
2.528	0.21567\\
2.5493	0.2322\\
2.5706	0.2496\\
2.5918	0.26786\\
2.6131	0.28699\\
2.6344	0.30696\\
2.6556	0.32776\\
2.6811	0.35377\\
2.7067	0.38087\\
2.7322	0.40898\\
2.762	0.44295\\
2.7917	0.47801\\
2.8258	0.51917\\
2.8683	0.57185\\
2.9619	0.68835\\
2.9831	0.71363\\
3.0044	0.73783\\
3.0214	0.75617\\
3.0384	0.77344\\
3.0554	0.78949\\
3.0682	0.80067\\
3.0809	0.81107\\
3.0937	0.82068\\
3.1065	0.82948\\
3.1192	0.83748\\
3.132	0.84465\\
3.1447	0.85101\\
3.1575	0.85655\\
3.1703	0.86127\\
3.183	0.86516\\
3.1958	0.86822\\
3.2085	0.87044\\
3.2213	0.87183\\
3.2341	0.87237\\
3.2468	0.87209\\
3.2596	0.87099\\
3.2723	0.8691\\
3.2851	0.86646\\
3.2979	0.86309\\
3.3106	0.85903\\
3.3234	0.85431\\
3.3361	0.84896\\
3.3489	0.84296\\
3.3617	0.83631\\
3.3744	0.82896\\
3.3872	0.82089\\
3.3999	0.81205\\
3.4127	0.80245\\
3.4255	0.79208\\
3.4425	0.77717\\
3.4595	0.76119\\
3.4808	0.74011\\
3.5105	0.70945\\
3.6041	0.61139\\
3.6679	0.54503\\
3.7104	0.50178\\
3.7445	0.46811\\
3.7742	0.43949\\
3.804	0.41177\\
3.8338	0.38504\\
3.8593	0.36297\\
3.8848	0.34171\\
3.9103	0.32128\\
3.9359	0.30171\\
3.9614	0.283\\
3.9869	0.26516\\
4.0124	0.24818\\
4.0379	0.23206\\
4.0635	0.21678\\
4.089	0.20232\\
4.1145	0.18867\\
4.14	0.17578\\
4.1655	0.16365\\
4.1911	0.15224\\
4.2166	0.14152\\
4.2421	0.13145\\
4.2719	0.1205\\
4.3016	0.11035\\
4.3314	0.10095\\
4.3612	0.092264\\
4.391	0.084237\\
4.425	0.075816\\
4.459	0.068139\\
4.493	0.061147\\
4.5313	0.054033\\
4.5696	0.047646\\
4.6079	0.041919\\
4.6504	0.036257\\
4.6929	0.031258\\
4.7397	0.026442\\
4.7908	0.021909\\
4.8461	0.01774\\
4.9056	0.013996\\
4.9694	0.010723\\
5.0375	0.0079478\\
5.1098	0.0056812\\
5.1906	0.0038257\\
5.2841	0.0023649\\
5.3607	0.001571\\
};
\addlegendentry{\ApproxBLNFitLegend}


\addplot[area legend, draw=none, fill=mycolor2, fill opacity=0.15, forget plot]
table[row sep=crcr] {%
x	y\\
1.1117	0.0002363\\
1.1159	0.00024216\\
1.1202	0.00024816\\
1.1244	0.0002543\\
1.1287	0.00026059\\
1.1329	0.00026703\\
1.1372	0.00027363\\
1.1415	0.00028038\\
1.1457	0.00028729\\
1.15	0.00029436\\
1.1542	0.0003016\\
1.1585	0.00030901\\
1.1627	0.0003166\\
1.167	0.00032437\\
1.1712	0.00033232\\
1.1755	0.00034046\\
1.1797	0.00034879\\
1.184	0.00035731\\
1.1882	0.00036604\\
1.1925	0.00037497\\
1.1967	0.00038411\\
1.201	0.00039346\\
1.2053	0.00040303\\
1.2095	0.00041283\\
1.2138	0.00042286\\
1.218	0.00043312\\
1.2223	0.00044362\\
1.2265	0.00045436\\
1.2308	0.00046536\\
1.235	0.00047661\\
1.2393	0.00048812\\
1.2435	0.00049989\\
1.2478	0.00051194\\
1.252	0.00052427\\
1.2563	0.00053689\\
1.2605	0.0005498\\
1.2648	0.000563\\
1.2691	0.00057651\\
1.2733	0.00059033\\
1.2776	0.00060447\\
1.2818	0.00061893\\
1.2861	0.00063373\\
1.2903	0.00064886\\
1.2946	0.00066434\\
1.2988	0.00068018\\
1.3031	0.00069638\\
1.3073	0.00071295\\
1.3116	0.00072989\\
1.3158	0.00074723\\
1.3201	0.00076496\\
1.3243	0.00078309\\
1.3286	0.00080163\\
1.3328	0.0008206\\
1.3371	0.00083999\\
1.3414	0.00085983\\
1.3456	0.00088011\\
1.3499	0.00090085\\
1.3541	0.00092207\\
1.3584	0.00094376\\
1.3626	0.00096594\\
1.3669	0.00098861\\
1.3711	0.0010118\\
1.3754	0.0010355\\
1.3796	0.0010598\\
1.3839	0.0010845\\
1.3881	0.0011099\\
1.3924	0.0011358\\
1.3966	0.0011623\\
1.4009	0.0011894\\
1.4052	0.0012171\\
1.4094	0.0012454\\
1.4137	0.0012743\\
1.4179	0.0013039\\
1.4222	0.0013341\\
1.4264	0.001365\\
1.4307	0.0013966\\
1.4349	0.0014289\\
1.4392	0.0014619\\
1.4434	0.0014956\\
1.4477	0.0015301\\
1.4519	0.0015653\\
1.4562	0.0016014\\
1.4604	0.0016382\\
1.4647	0.0016758\\
1.469	0.0017142\\
1.4732	0.0017535\\
1.4775	0.0017937\\
1.4817	0.0018347\\
1.486	0.0018767\\
1.4902	0.0019195\\
1.4945	0.0019633\\
1.4987	0.002008\\
1.503	0.0020537\\
1.5072	0.0021004\\
1.5115	0.0021482\\
1.5157	0.0021969\\
1.52	0.0022467\\
1.5242	0.0022976\\
1.5285	0.0023496\\
1.5328	0.0024028\\
1.537	0.002457\\
1.5413	0.0025125\\
1.5455	0.0025691\\
1.5498	0.002627\\
1.554	0.0026861\\
1.5583	0.0027465\\
1.5625	0.0028081\\
1.5668	0.0028711\\
1.571	0.0029355\\
1.5753	0.0030012\\
1.5795	0.0030683\\
1.5838	0.0031369\\
1.588	0.0032069\\
1.5923	0.0032784\\
1.5966	0.0033515\\
1.6008	0.0034261\\
1.6051	0.0035022\\
1.6093	0.00358\\
1.6136	0.0036595\\
1.6178	0.0037406\\
1.6221	0.0038234\\
1.6263	0.003908\\
1.6306	0.0039944\\
1.6348	0.0040826\\
1.6391	0.0041727\\
1.6433	0.0042646\\
1.6476	0.0043585\\
1.6518	0.0044543\\
1.6561	0.0045522\\
1.6604	0.0046521\\
1.6646	0.0047541\\
1.6689	0.0048582\\
1.6731	0.0049645\\
1.6774	0.005073\\
1.6816	0.0051838\\
1.6859	0.0052969\\
1.6901	0.0054123\\
1.6944	0.0055302\\
1.6986	0.0056504\\
1.7029	0.0057732\\
1.7071	0.0058985\\
1.7114	0.0060264\\
1.7156	0.0061569\\
1.7199	0.0062901\\
1.7242	0.0064261\\
1.7284	0.0065648\\
1.7327	0.0067064\\
1.7369	0.0068509\\
1.7412	0.0069984\\
1.7454	0.0071489\\
1.7497	0.0073024\\
1.7539	0.0074591\\
1.7582	0.007619\\
1.7624	0.0077821\\
1.7667	0.0079486\\
1.7709	0.0081184\\
1.7752	0.0082917\\
1.7794	0.0084685\\
1.7837	0.0086489\\
1.788	0.0088329\\
1.7922	0.0090206\\
1.7965	0.0092121\\
1.8007	0.0094075\\
1.805	0.0096068\\
1.8092	0.00981\\
1.8135	0.010017\\
1.8177	0.010229\\
1.822	0.010445\\
1.8262	0.010665\\
1.8305	0.010889\\
1.8347	0.011118\\
1.839	0.011352\\
1.8432	0.01159\\
1.8475	0.011832\\
1.8517	0.01208\\
1.856	0.012333\\
1.8603	0.01259\\
1.8645	0.012853\\
1.8688	0.01312\\
1.873	0.013393\\
1.8773	0.013672\\
1.8815	0.013956\\
1.8858	0.014245\\
1.89	0.01454\\
1.8943	0.014841\\
1.8985	0.015147\\
1.9028	0.01546\\
1.907	0.015779\\
1.9113	0.016103\\
1.9155	0.016435\\
1.9198	0.016772\\
1.9241	0.017116\\
1.9283	0.017467\\
1.9326	0.017824\\
1.9368	0.018189\\
1.9411	0.01856\\
1.9453	0.018938\\
1.9496	0.019324\\
1.9538	0.019717\\
1.9581	0.020117\\
1.9623	0.020525\\
1.9666	0.020941\\
1.9708	0.021364\\
1.9751	0.021796\\
1.9793	0.022236\\
1.9836	0.022684\\
1.9879	0.02314\\
1.9921	0.023605\\
1.9964	0.024079\\
2.0006	0.024562\\
2.0049	0.025053\\
2.0091	0.025554\\
2.0134	0.026064\\
2.0176	0.026584\\
2.0219	0.027113\\
2.0261	0.027652\\
2.0304	0.028201\\
2.0346	0.02876\\
2.0389	0.029329\\
2.0431	0.029908\\
2.0474	0.030499\\
2.0517	0.031099\\
2.0559	0.031711\\
2.0602	0.032334\\
2.0644	0.032969\\
2.0687	0.033614\\
2.0729	0.034272\\
2.0772	0.034941\\
2.0814	0.035622\\
2.0857	0.036316\\
2.0899	0.037022\\
2.0942	0.03774\\
2.0984	0.038472\\
2.1027	0.039216\\
2.1069	0.039973\\
2.1112	0.040744\\
2.1155	0.041528\\
2.1197	0.042326\\
2.124	0.043139\\
2.1282	0.043965\\
2.1325	0.044806\\
2.1367	0.045661\\
2.141	0.046531\\
2.1452	0.047417\\
2.1495	0.048317\\
2.1537	0.049233\\
2.158	0.050165\\
2.1622	0.051113\\
2.1665	0.052077\\
2.1707	0.053057\\
2.175	0.054054\\
2.1793	0.055068\\
2.1835	0.056099\\
2.1878	0.057147\\
2.192	0.058213\\
2.1963	0.059297\\
2.2005	0.060398\\
2.2048	0.061519\\
2.209	0.062657\\
2.2133	0.063815\\
2.2175	0.064991\\
2.2218	0.066187\\
2.226	0.067403\\
2.2303	0.068638\\
2.2345	0.069893\\
2.2388	0.071169\\
2.2431	0.072466\\
2.2473	0.073783\\
2.2516	0.075121\\
2.2558	0.076481\\
2.2601	0.077863\\
2.2643	0.079266\\
2.2686	0.080692\\
2.2728	0.08214\\
2.2771	0.083611\\
2.2813	0.085105\\
2.2856	0.086622\\
2.2898	0.088163\\
2.2941	0.089727\\
2.2983	0.091316\\
2.3026	0.092929\\
2.3069	0.094567\\
2.3111	0.09623\\
2.3154	0.097918\\
2.3196	0.099631\\
2.3239	0.10137\\
2.3281	0.10314\\
2.3324	0.10493\\
2.3366	0.10675\\
2.3409	0.10859\\
2.3451	0.11046\\
2.3494	0.11236\\
2.3536	0.11429\\
2.3579	0.11625\\
2.3621	0.11823\\
2.3664	0.12024\\
2.3706	0.12228\\
2.3749	0.12436\\
2.3792	0.12645\\
2.3834	0.12858\\
2.3877	0.13074\\
2.3919	0.13293\\
2.3962	0.13515\\
2.4004	0.1374\\
2.4047	0.13968\\
2.4089	0.142\\
2.4132	0.14434\\
2.4174	0.14672\\
2.4217	0.14912\\
2.4259	0.15156\\
2.4302	0.15403\\
2.4344	0.15654\\
2.4387	0.15907\\
2.443	0.16164\\
2.4472	0.16425\\
2.4515	0.16688\\
2.4557	0.16955\\
2.46	0.17226\\
2.4642	0.17499\\
2.4685	0.17776\\
2.4727	0.18057\\
2.477	0.18341\\
2.4812	0.18629\\
2.4855	0.18919\\
2.4897	0.19214\\
2.494	0.19512\\
2.4982	0.19813\\
2.5025	0.20118\\
2.5068	0.20427\\
2.511	0.20739\\
2.5153	0.21054\\
2.5195	0.21373\\
2.5238	0.21696\\
2.528	0.22022\\
2.5323	0.22352\\
2.5365	0.22685\\
2.5408	0.23022\\
2.545	0.23363\\
2.5493	0.23707\\
2.5535	0.24055\\
2.5578	0.24406\\
2.562	0.24761\\
2.5663	0.25119\\
2.5706	0.25481\\
2.5748	0.25847\\
2.5791	0.26216\\
2.5833	0.26588\\
2.5876	0.26964\\
2.5918	0.27344\\
2.5961	0.27727\\
2.6003	0.28114\\
2.6046	0.28504\\
2.6088	0.28898\\
2.6131	0.29295\\
2.6173	0.29696\\
2.6216	0.301\\
2.6258	0.30507\\
2.6301	0.30918\\
2.6344	0.31332\\
2.6386	0.31749\\
2.6429	0.3217\\
2.6471	0.32594\\
2.6514	0.33021\\
2.6556	0.33452\\
2.6599	0.33886\\
2.6641	0.34322\\
2.6684	0.34762\\
2.6726	0.35205\\
2.6769	0.35651\\
2.6811	0.36101\\
2.6854	0.36553\\
2.6896	0.37008\\
2.6939	0.37466\\
2.6982	0.37926\\
2.7024	0.3839\\
2.7067	0.38856\\
2.7109	0.39325\\
2.7152	0.39797\\
2.7194	0.40271\\
2.7237	0.40748\\
2.7279	0.41227\\
2.7322	0.41709\\
2.7364	0.42193\\
2.7407	0.4268\\
2.7449	0.43169\\
2.7492	0.4366\\
2.7534	0.44153\\
2.7577	0.44648\\
2.762	0.45146\\
2.7662	0.45645\\
2.7705	0.46146\\
2.7747	0.46649\\
2.779	0.47154\\
2.7832	0.47661\\
2.7875	0.48169\\
2.7917	0.48679\\
2.796	0.4919\\
2.8002	0.49703\\
2.8045	0.50217\\
2.8087	0.50733\\
2.813	0.51249\\
2.8172	0.51767\\
2.8215	0.52286\\
2.8258	0.52806\\
2.83	0.53327\\
2.8343	0.53849\\
2.8385	0.54371\\
2.8428	0.54895\\
2.847	0.55419\\
2.8513	0.55943\\
2.8555	0.56469\\
2.8598	0.56994\\
2.864	0.57521\\
2.8683	0.58047\\
2.8725	0.58574\\
2.8768	0.59102\\
2.881	0.5963\\
2.8853	0.60158\\
2.8896	0.60687\\
2.8938	0.61216\\
2.8981	0.61746\\
2.9023	0.62276\\
2.9066	0.62806\\
2.9108	0.63337\\
2.9151	0.63869\\
2.9193	0.64401\\
2.9236	0.64934\\
2.9278	0.65467\\
2.9321	0.66001\\
2.9363	0.66535\\
2.9406	0.6707\\
2.9448	0.67605\\
2.9491	0.6814\\
2.9533	0.68675\\
2.9576	0.69209\\
2.9619	0.69742\\
2.9661	0.70274\\
2.9704	0.70804\\
2.9746	0.71331\\
2.9789	0.71856\\
2.9831	0.72377\\
2.9874	0.72895\\
2.9916	0.73407\\
2.9959	0.73915\\
3.0001	0.74417\\
3.0044	0.74913\\
3.0086	0.75403\\
3.0129	0.75885\\
3.0171	0.7636\\
3.0214	0.76826\\
3.0257	0.77284\\
3.0299	0.77734\\
3.0342	0.78174\\
3.0384	0.78605\\
3.0427	0.79026\\
3.0469	0.79437\\
3.0512	0.79839\\
3.0554	0.8023\\
3.0597	0.80611\\
3.0639	0.80983\\
3.0682	0.81344\\
3.0724	0.81695\\
3.0767	0.82036\\
3.0809	0.82368\\
3.0852	0.8269\\
3.0895	0.83003\\
3.0937	0.83306\\
3.098	0.83601\\
3.1022	0.83887\\
3.1065	0.84164\\
3.1107	0.84434\\
3.115	0.84695\\
3.1192	0.84948\\
3.1235	0.85193\\
3.1277	0.85431\\
3.132	0.8566\\
3.1362	0.85883\\
3.1405	0.86097\\
3.1447	0.86304\\
3.149	0.86502\\
3.1533	0.86693\\
3.1575	0.86876\\
3.1618	0.8705\\
3.166	0.87215\\
3.1703	0.87372\\
3.1745	0.87521\\
3.1788	0.8766\\
3.183	0.87789\\
3.1873	0.8791\\
3.1915	0.8802\\
3.1958	0.88121\\
3.2	0.88212\\
3.2043	0.88293\\
3.2085	0.88364\\
3.2128	0.88425\\
3.2171	0.88475\\
3.2213	0.88515\\
3.2256	0.88545\\
3.2298	0.88564\\
3.2341	0.88573\\
3.2383	0.88571\\
3.2426	0.88559\\
3.2468	0.88536\\
3.2511	0.88503\\
3.2553	0.8846\\
3.2596	0.88407\\
3.2638	0.88343\\
3.2681	0.8827\\
3.2723	0.88187\\
3.2766	0.88094\\
3.2809	0.87991\\
3.2851	0.8788\\
3.2894	0.87759\\
3.2936	0.8763\\
3.2979	0.87492\\
3.3021	0.87346\\
3.3064	0.87193\\
3.3106	0.87032\\
3.3149	0.86865\\
3.3191	0.86691\\
3.3234	0.86512\\
3.3276	0.86328\\
3.3319	0.86138\\
3.3361	0.85945\\
3.3404	0.85747\\
3.3447	0.85546\\
3.3489	0.85342\\
3.3532	0.85133\\
3.3574	0.84921\\
3.3617	0.84705\\
3.3659	0.84484\\
3.3702	0.84258\\
3.3744	0.84025\\
3.3787	0.83785\\
3.3829	0.83537\\
3.3872	0.8328\\
3.3914	0.83014\\
3.3957	0.82737\\
3.3999	0.82448\\
3.4042	0.82148\\
3.4085	0.81836\\
3.4127	0.81511\\
3.417	0.81174\\
3.4212	0.80823\\
3.4255	0.80461\\
3.4297	0.80086\\
3.434	0.79699\\
3.4382	0.79301\\
3.4425	0.78893\\
3.4467	0.78475\\
3.451	0.78048\\
3.4552	0.77614\\
3.4595	0.77173\\
3.4637	0.76726\\
3.468	0.76275\\
3.4722	0.7582\\
3.4765	0.75363\\
3.4808	0.74905\\
3.485	0.74446\\
3.4893	0.73987\\
3.4935	0.73529\\
3.4978	0.73073\\
3.502	0.72618\\
3.5063	0.72165\\
3.5105	0.71714\\
3.5148	0.71265\\
3.519	0.70818\\
3.5233	0.70372\\
3.5275	0.69927\\
3.5318	0.69483\\
3.536	0.6904\\
3.5403	0.68598\\
3.5446	0.68155\\
3.5488	0.67713\\
3.5531	0.67271\\
3.5573	0.66829\\
3.5616	0.66386\\
3.5658	0.65944\\
3.5701	0.65501\\
3.5743	0.65057\\
3.5786	0.64613\\
3.5828	0.64169\\
3.5871	0.63725\\
3.5913	0.6328\\
3.5956	0.62835\\
3.5998	0.6239\\
3.6041	0.61944\\
3.6084	0.61499\\
3.6126	0.61053\\
3.6169	0.60607\\
3.6211	0.60161\\
3.6254	0.59715\\
3.6296	0.5927\\
3.6339	0.58824\\
3.6381	0.58379\\
3.6424	0.57934\\
3.6466	0.57489\\
3.6509	0.57044\\
3.6551	0.566\\
3.6594	0.56157\\
3.6636	0.55714\\
3.6679	0.55272\\
3.6722	0.5483\\
3.6764	0.54389\\
3.6807	0.53949\\
3.6849	0.5351\\
3.6892	0.53072\\
3.6934	0.52634\\
3.6977	0.52198\\
3.7019	0.51762\\
3.7062	0.51328\\
3.7104	0.50895\\
3.7147	0.50463\\
3.7189	0.50033\\
3.7232	0.49603\\
3.7274	0.49175\\
3.7317	0.48749\\
3.736	0.48324\\
3.7402	0.479\\
3.7445	0.47478\\
3.7487	0.47057\\
3.753	0.46639\\
3.7572	0.46221\\
3.7615	0.45806\\
3.7657	0.45392\\
3.77	0.4498\\
3.7742	0.4457\\
3.7785	0.44161\\
3.7827	0.43755\\
3.787	0.4335\\
3.7912	0.42948\\
3.7955	0.42547\\
3.7998	0.42148\\
3.804	0.41752\\
3.8083	0.41357\\
3.8125	0.40965\\
3.8168	0.40574\\
3.821	0.40186\\
3.8253	0.398\\
3.8295	0.39416\\
3.8338	0.39034\\
3.838	0.38655\\
3.8423	0.38278\\
3.8465	0.37903\\
3.8508	0.3753\\
3.855	0.3716\\
3.8593	0.36792\\
3.8636	0.36427\\
3.8678	0.36063\\
3.8721	0.35702\\
3.8763	0.35344\\
3.8806	0.34988\\
3.8848	0.34634\\
3.8891	0.34283\\
3.8933	0.33934\\
3.8976	0.33588\\
3.9018	0.33244\\
3.9061	0.32902\\
3.9103	0.32564\\
3.9146	0.32227\\
3.9188	0.31893\\
3.9231	0.31561\\
3.9274	0.31232\\
3.9316	0.30906\\
3.9359	0.30582\\
3.9401	0.3026\\
3.9444	0.29941\\
3.9486	0.29625\\
3.9529	0.29311\\
3.9571	0.28999\\
3.9614	0.2869\\
3.9656	0.28383\\
3.9699	0.28079\\
3.9741	0.27778\\
3.9784	0.27478\\
3.9826	0.27182\\
3.9869	0.26888\\
3.9912	0.26596\\
3.9954	0.26307\\
3.9997	0.2602\\
4.0039	0.25736\\
4.0082	0.25454\\
4.0124	0.25174\\
4.0167	0.24897\\
4.0209	0.24623\\
4.0252	0.24351\\
4.0294	0.24081\\
4.0337	0.23813\\
4.0379	0.23548\\
4.0422	0.23286\\
4.0464	0.23026\\
4.0507	0.22768\\
4.0549	0.22512\\
4.0592	0.22259\\
4.0635	0.22008\\
4.0677	0.2176\\
4.072	0.21513\\
4.0762	0.21269\\
4.0805	0.21028\\
4.0847	0.20788\\
4.089	0.20551\\
4.0932	0.20316\\
4.0975	0.20083\\
4.1017	0.19853\\
4.106	0.19624\\
4.1102	0.19398\\
4.1145	0.19174\\
4.1187	0.18952\\
4.123	0.18733\\
4.1273	0.18515\\
4.1315	0.183\\
4.1358	0.18086\\
4.14	0.17875\\
4.1443	0.17666\\
4.1485	0.17459\\
4.1528	0.17254\\
4.157	0.17051\\
4.1613	0.1685\\
4.1655	0.16651\\
4.1698	0.16454\\
4.174	0.16259\\
4.1783	0.16066\\
4.1825	0.15875\\
4.1868	0.15686\\
4.1911	0.15499\\
4.1953	0.15313\\
4.1996	0.1513\\
4.2038	0.14948\\
4.2081	0.14769\\
4.2123	0.14591\\
4.2166	0.14415\\
4.2208	0.14241\\
4.2251	0.14068\\
4.2293	0.13898\\
4.2336	0.13729\\
4.2378	0.13562\\
4.2421	0.13397\\
4.2463	0.13233\\
4.2506	0.13072\\
4.2549	0.12911\\
4.2591	0.12753\\
4.2634	0.12596\\
4.2676	0.12441\\
4.2719	0.12288\\
4.2761	0.12136\\
4.2804	0.11986\\
4.2846	0.11837\\
4.2889	0.11691\\
4.2931	0.11545\\
4.2974	0.11401\\
4.3016	0.11259\\
4.3059	0.11119\\
4.3101	0.10979\\
4.3144	0.10842\\
4.3187	0.10706\\
4.3229	0.10571\\
4.3272	0.10438\\
4.3314	0.10306\\
4.3357	0.10176\\
4.3399	0.10047\\
4.3442	0.099194\\
4.3484	0.097934\\
4.3527	0.096688\\
4.3569	0.095455\\
4.3612	0.094236\\
4.3654	0.093031\\
4.3697	0.091839\\
4.3739	0.09066\\
4.3782	0.089495\\
4.3825	0.088342\\
4.3867	0.087203\\
4.391	0.086076\\
4.3952	0.084962\\
4.3995	0.08386\\
4.4037	0.082771\\
4.408	0.081694\\
4.4122	0.08063\\
4.4165	0.079577\\
4.4207	0.078536\\
4.425	0.077508\\
4.4292	0.07649\\
4.4335	0.075485\\
4.4377	0.07449\\
4.442	0.073508\\
4.4463	0.072536\\
4.4505	0.071576\\
4.4548	0.070626\\
4.459	0.069687\\
4.4633	0.06876\\
4.4675	0.067842\\
4.4718	0.066936\\
4.476	0.06604\\
4.4803	0.065154\\
4.4845	0.064279\\
4.4888	0.063413\\
4.493	0.062558\\
4.4973	0.061712\\
4.5015	0.060877\\
4.5058	0.060051\\
4.5101	0.059235\\
4.5143	0.058428\\
4.5186	0.057631\\
4.5228	0.056843\\
4.5271	0.056065\\
4.5313	0.055295\\
4.5356	0.054535\\
4.5398	0.053783\\
4.5441	0.053041\\
4.5483	0.052307\\
4.5526	0.051582\\
4.5568	0.050866\\
4.5611	0.050158\\
4.5653	0.049458\\
4.5696	0.048767\\
4.5738	0.048084\\
4.5781	0.04741\\
4.5824	0.046743\\
4.5866	0.046085\\
4.5909	0.045435\\
4.5951	0.044792\\
4.5994	0.044157\\
4.6036	0.04353\\
4.6079	0.042911\\
4.6121	0.042299\\
4.6164	0.041695\\
4.6206	0.041098\\
4.6249	0.040508\\
4.6291	0.039926\\
4.6334	0.039351\\
4.6376	0.038783\\
4.6419	0.038222\\
4.6462	0.037669\\
4.6504	0.037122\\
4.6547	0.036582\\
4.6589	0.036049\\
4.6632	0.035522\\
4.6674	0.035002\\
4.6717	0.034489\\
4.6759	0.033983\\
4.6802	0.033483\\
4.6844	0.032989\\
4.6887	0.032502\\
4.6929	0.032021\\
4.6972	0.031546\\
4.7014	0.031077\\
4.7057	0.030614\\
4.71	0.030158\\
4.7142	0.029707\\
4.7185	0.029262\\
4.7227	0.028824\\
4.727	0.02839\\
4.7312	0.027963\\
4.7355	0.027541\\
4.7397	0.027125\\
4.744	0.026714\\
4.7482	0.026308\\
4.7525	0.025908\\
4.7567	0.025514\\
4.761	0.025124\\
4.7652	0.02474\\
4.7695	0.024361\\
4.7738	0.023986\\
4.778	0.023617\\
4.7823	0.023253\\
4.7865	0.022893\\
4.7908	0.022538\\
4.795	0.022188\\
4.7993	0.021843\\
4.8035	0.021502\\
4.8078	0.021166\\
4.812	0.020834\\
4.8163	0.020506\\
4.8205	0.020183\\
4.8248	0.019864\\
4.829	0.01955\\
4.8333	0.019239\\
4.8376	0.018933\\
4.8418	0.018631\\
4.8461	0.018332\\
4.8503	0.018038\\
4.8546	0.017747\\
4.8588	0.017461\\
4.8631	0.017178\\
4.8673	0.016899\\
4.8716	0.016623\\
4.8758	0.016351\\
4.8801	0.016083\\
4.8843	0.015819\\
4.8886	0.015557\\
4.8928	0.0153\\
4.8971	0.015046\\
4.9014	0.014795\\
4.9056	0.014547\\
4.9099	0.014303\\
4.9141	0.014062\\
4.9184	0.013824\\
4.9226	0.01359\\
4.9269	0.013359\\
4.9311	0.013131\\
4.9354	0.012905\\
4.9396	0.012683\\
4.9439	0.012465\\
4.9481	0.012249\\
4.9524	0.012036\\
4.9566	0.011825\\
4.9609	0.011618\\
4.9652	0.011414\\
4.9694	0.011213\\
4.9737	0.011014\\
4.9779	0.010818\\
4.9822	0.010625\\
4.9864	0.010435\\
4.9907	0.010248\\
4.9949	0.010063\\
4.9992	0.0098807\\
5.0034	0.0097013\\
5.0077	0.0095244\\
5.0119	0.0093501\\
5.0162	0.0091785\\
5.0204	0.0090093\\
5.0247	0.0088427\\
5.029	0.0086786\\
5.0332	0.008517\\
5.0375	0.0083579\\
5.0417	0.0082012\\
5.046	0.0080469\\
5.0502	0.007895\\
5.0545	0.0077455\\
5.0587	0.0075983\\
5.063	0.0074534\\
5.0672	0.0073109\\
5.0715	0.0071706\\
5.0757	0.0070326\\
5.08	0.0068968\\
5.0842	0.0067632\\
5.0885	0.0066318\\
5.0928	0.0065026\\
5.097	0.0063755\\
5.1013	0.0062506\\
5.1055	0.0061277\\
5.1098	0.0060069\\
5.114	0.0058882\\
5.1183	0.0057715\\
5.1225	0.0056568\\
5.1268	0.005544\\
5.131	0.0054333\\
5.1353	0.0053244\\
5.1395	0.0052175\\
5.1438	0.0051125\\
5.148	0.0050093\\
5.1523	0.004908\\
5.1565	0.0048084\\
5.1608	0.0047107\\
5.1651	0.0046147\\
5.1693	0.0045205\\
5.1736	0.004428\\
5.1778	0.0043372\\
5.1821	0.004248\\
5.1863	0.0041606\\
5.1906	0.0040747\\
5.1948	0.0039904\\
5.1991	0.0039077\\
5.2033	0.0038266\\
5.2076	0.003747\\
5.2118	0.0036688\\
5.2161	0.0035922\\
5.2203	0.003517\\
5.2246	0.0034433\\
5.2289	0.003371\\
5.2331	0.0033001\\
5.2374	0.0032305\\
5.2416	0.0031623\\
5.2459	0.0030954\\
5.2501	0.0030298\\
5.2544	0.0029655\\
5.2586	0.0029024\\
5.2629	0.0028406\\
5.2671	0.00278\\
5.2714	0.0027205\\
5.2756	0.0026623\\
5.2799	0.0026052\\
5.2841	0.0025492\\
5.2884	0.0024944\\
5.2927	0.0024406\\
5.2969	0.002388\\
5.3012	0.0023363\\
5.3054	0.0022857\\
5.3097	0.0022362\\
5.3139	0.0021876\\
5.3182	0.00214\\
5.3224	0.0020934\\
5.3267	0.0020477\\
5.3309	0.0020029\\
5.3352	0.001959\\
5.3394	0.0019161\\
5.3437	0.001874\\
5.3479	0.0018328\\
5.3522	0.0017924\\
5.3565	0.0017528\\
5.3607	0.0017141\\
5.3607	0.0014278\\
5.3565	0.0014623\\
5.3522	0.0014976\\
5.3479	0.0015337\\
5.3437	0.0015706\\
5.3394	0.0016083\\
5.3352	0.0016469\\
5.3309	0.0016863\\
5.3267	0.0017265\\
5.3224	0.0017677\\
5.3182	0.0018097\\
5.3139	0.0018527\\
5.3097	0.0018966\\
5.3054	0.0019414\\
5.3012	0.0019873\\
5.2969	0.002034\\
5.2927	0.0020818\\
5.2884	0.0021307\\
5.2841	0.0021805\\
5.2799	0.0022314\\
5.2756	0.0022834\\
5.2714	0.0023364\\
5.2671	0.0023906\\
5.2629	0.0024459\\
5.2586	0.0025023\\
5.2544	0.0025599\\
5.2501	0.0026187\\
5.2459	0.0026787\\
5.2416	0.0027399\\
5.2374	0.0028023\\
5.2331	0.002866\\
5.2289	0.002931\\
5.2246	0.0029973\\
5.2203	0.0030649\\
5.2161	0.0031338\\
5.2118	0.0032041\\
5.2076	0.0032758\\
5.2033	0.0033489\\
5.1991	0.0034234\\
5.1948	0.0034993\\
5.1906	0.0035767\\
5.1863	0.0036556\\
5.1821	0.003736\\
5.1778	0.0038179\\
5.1736	0.0039014\\
5.1693	0.0039865\\
5.1651	0.0040731\\
5.1608	0.0041614\\
5.1565	0.0042513\\
5.1523	0.0043428\\
5.148	0.0044361\\
5.1438	0.0045311\\
5.1395	0.0046278\\
5.1353	0.0047263\\
5.131	0.0048265\\
5.1268	0.0049286\\
5.1225	0.0050325\\
5.1183	0.0051383\\
5.114	0.0052459\\
5.1098	0.0053555\\
5.1055	0.005467\\
5.1013	0.0055805\\
5.097	0.0056959\\
5.0928	0.0058134\\
5.0885	0.005933\\
5.0842	0.0060546\\
5.08	0.0061783\\
5.0757	0.0063042\\
5.0715	0.0064322\\
5.0672	0.0065624\\
5.063	0.0066948\\
5.0587	0.0068295\\
5.0545	0.0069665\\
5.0502	0.0071058\\
5.046	0.0072474\\
5.0417	0.0073914\\
5.0375	0.0075378\\
5.0332	0.0076866\\
5.029	0.0078379\\
5.0247	0.0079917\\
5.0204	0.008148\\
5.0162	0.0083068\\
5.0119	0.0084683\\
5.0077	0.0086324\\
5.0034	0.0087992\\
4.9992	0.0089686\\
4.9949	0.0091408\\
4.9907	0.0093157\\
4.9864	0.0094934\\
4.9822	0.009674\\
4.9779	0.0098574\\
4.9737	0.010044\\
4.9694	0.010233\\
4.9652	0.010425\\
4.9609	0.01062\\
4.9566	0.010818\\
4.9524	0.01102\\
4.9481	0.011224\\
4.9439	0.011432\\
4.9396	0.011642\\
4.9354	0.011856\\
4.9311	0.012073\\
4.9269	0.012294\\
4.9226	0.012517\\
4.9184	0.012744\\
4.9141	0.012975\\
4.9099	0.013208\\
4.9056	0.013446\\
4.9014	0.013686\\
4.8971	0.013931\\
4.8928	0.014178\\
4.8886	0.01443\\
4.8843	0.014685\\
4.8801	0.014943\\
4.8758	0.015206\\
4.8716	0.015472\\
4.8673	0.015742\\
4.8631	0.016015\\
4.8588	0.016293\\
4.8546	0.016574\\
4.8503	0.016859\\
4.8461	0.017148\\
4.8418	0.017441\\
4.8376	0.017738\\
4.8333	0.018039\\
4.829	0.018344\\
4.8248	0.018654\\
4.8205	0.018967\\
4.8163	0.019285\\
4.812	0.019606\\
4.8078	0.019932\\
4.8035	0.020263\\
4.7993	0.020597\\
4.795	0.020936\\
4.7908	0.02128\\
4.7865	0.021628\\
4.7823	0.02198\\
4.778	0.022337\\
4.7738	0.022698\\
4.7695	0.023064\\
4.7652	0.023435\\
4.761	0.02381\\
4.7567	0.02419\\
4.7525	0.024575\\
4.7482	0.024965\\
4.744	0.02536\\
4.7397	0.025759\\
4.7355	0.026164\\
4.7312	0.026573\\
4.727	0.026988\\
4.7227	0.027408\\
4.7185	0.027833\\
4.7142	0.028263\\
4.71	0.028698\\
4.7057	0.029139\\
4.7014	0.029586\\
4.6972	0.030038\\
4.6929	0.030495\\
4.6887	0.030958\\
4.6844	0.031427\\
4.6802	0.031902\\
4.6759	0.032382\\
4.6717	0.032868\\
4.6674	0.033361\\
4.6632	0.033859\\
4.6589	0.034364\\
4.6547	0.034875\\
4.6504	0.035392\\
4.6462	0.035915\\
4.6419	0.036445\\
4.6376	0.036982\\
4.6334	0.037525\\
4.6291	0.038075\\
4.6249	0.038631\\
4.6206	0.039195\\
4.6164	0.039765\\
4.6121	0.040343\\
4.6079	0.040928\\
4.6036	0.041519\\
4.5994	0.042119\\
4.5951	0.042725\\
4.5909	0.043339\\
4.5866	0.043961\\
4.5824	0.04459\\
4.5781	0.045227\\
4.5738	0.045872\\
4.5696	0.046525\\
4.5653	0.047185\\
4.5611	0.047854\\
4.5568	0.048531\\
4.5526	0.049217\\
4.5483	0.04991\\
4.5441	0.050613\\
4.5398	0.051324\\
4.5356	0.052043\\
4.5313	0.052771\\
4.5271	0.053509\\
4.5228	0.054255\\
4.5186	0.05501\\
4.5143	0.055774\\
4.5101	0.056548\\
4.5058	0.057331\\
4.5015	0.058123\\
4.4973	0.058925\\
4.493	0.059737\\
4.4888	0.060558\\
4.4845	0.06139\\
4.4803	0.062231\\
4.476	0.063083\\
4.4718	0.063944\\
4.4675	0.064816\\
4.4633	0.065698\\
4.459	0.066591\\
4.4548	0.067495\\
4.4505	0.068409\\
4.4463	0.069334\\
4.442	0.07027\\
4.4377	0.071217\\
4.4335	0.072175\\
4.4292	0.073144\\
4.425	0.074125\\
4.4207	0.075118\\
4.4165	0.076122\\
4.4122	0.077138\\
4.408	0.078166\\
4.4037	0.079206\\
4.3995	0.080258\\
4.3952	0.081322\\
4.391	0.082399\\
4.3867	0.083488\\
4.3825	0.08459\\
4.3782	0.085704\\
4.3739	0.086832\\
4.3697	0.087972\\
4.3654	0.089126\\
4.3612	0.090293\\
4.3569	0.091473\\
4.3527	0.092667\\
4.3484	0.093875\\
4.3442	0.095096\\
4.3399	0.096332\\
4.3357	0.097581\\
4.3314	0.098845\\
4.3272	0.10012\\
4.3229	0.10142\\
4.3187	0.10272\\
4.3144	0.10405\\
4.3101	0.10538\\
4.3059	0.10673\\
4.3016	0.1081\\
4.2974	0.10949\\
4.2931	0.11088\\
4.2889	0.1123\\
4.2846	0.11373\\
4.2804	0.11517\\
4.2761	0.11664\\
4.2719	0.11811\\
4.2676	0.11961\\
4.2634	0.12112\\
4.2591	0.12265\\
4.2549	0.12419\\
4.2506	0.12575\\
4.2463	0.12733\\
4.2421	0.12893\\
4.2378	0.13054\\
4.2336	0.13218\\
4.2293	0.13382\\
4.2251	0.13549\\
4.2208	0.13718\\
4.2166	0.13888\\
4.2123	0.1406\\
4.2081	0.14234\\
4.2038	0.1441\\
4.1996	0.14588\\
4.1953	0.14768\\
4.1911	0.14949\\
4.1868	0.15133\\
4.1825	0.15318\\
4.1783	0.15505\\
4.174	0.15695\\
4.1698	0.15886\\
4.1655	0.16079\\
4.1613	0.16275\\
4.157	0.16472\\
4.1528	0.16671\\
4.1485	0.16873\\
4.1443	0.17076\\
4.14	0.17282\\
4.1358	0.17489\\
4.1315	0.17699\\
4.1273	0.17911\\
4.123	0.18125\\
4.1187	0.18341\\
4.1145	0.18559\\
4.1102	0.18779\\
4.106	0.19002\\
4.1017	0.19227\\
4.0975	0.19454\\
4.0932	0.19683\\
4.089	0.19914\\
4.0847	0.20147\\
4.0805	0.20383\\
4.0762	0.20621\\
4.072	0.20861\\
4.0677	0.21104\\
4.0635	0.21348\\
4.0592	0.21595\\
4.0549	0.21844\\
4.0507	0.22096\\
4.0464	0.2235\\
4.0422	0.22606\\
4.0379	0.22864\\
4.0337	0.23125\\
4.0294	0.23388\\
4.0252	0.23653\\
4.0209	0.23921\\
4.0167	0.2419\\
4.0124	0.24463\\
4.0082	0.24737\\
4.0039	0.25014\\
3.9997	0.25293\\
3.9954	0.25575\\
3.9912	0.25858\\
3.9869	0.26145\\
3.9826	0.26433\\
3.9784	0.26724\\
3.9741	0.27017\\
3.9699	0.27313\\
3.9656	0.2761\\
3.9614	0.27911\\
3.9571	0.28213\\
3.9529	0.28518\\
3.9486	0.28825\\
3.9444	0.29134\\
3.9401	0.29446\\
3.9359	0.2976\\
3.9316	0.30077\\
3.9274	0.30395\\
3.9231	0.30716\\
3.9188	0.3104\\
3.9146	0.31365\\
3.9103	0.31693\\
3.9061	0.32023\\
3.9018	0.32355\\
3.8976	0.3269\\
3.8933	0.33027\\
3.8891	0.33366\\
3.8848	0.33707\\
3.8806	0.34051\\
3.8763	0.34397\\
3.8721	0.34745\\
3.8678	0.35095\\
3.8636	0.35447\\
3.8593	0.35802\\
3.855	0.36159\\
3.8508	0.36517\\
3.8465	0.36878\\
3.8423	0.37242\\
3.838	0.37607\\
3.8338	0.37974\\
3.8295	0.38344\\
3.8253	0.38715\\
3.821	0.39089\\
3.8168	0.39464\\
3.8125	0.39842\\
3.8083	0.40222\\
3.804	0.40603\\
3.7998	0.40987\\
3.7955	0.41372\\
3.7912	0.4176\\
3.787	0.42149\\
3.7827	0.4254\\
3.7785	0.42934\\
3.7742	0.43329\\
3.77	0.43726\\
3.7657	0.44124\\
3.7615	0.44525\\
3.7572	0.44927\\
3.753	0.45331\\
3.7487	0.45737\\
3.7445	0.46144\\
3.7402	0.46553\\
3.736	0.46964\\
3.7317	0.47376\\
3.7274	0.4779\\
3.7232	0.48206\\
3.7189	0.48623\\
3.7147	0.49041\\
3.7104	0.49461\\
3.7062	0.49883\\
3.7019	0.50306\\
3.6977	0.5073\\
3.6934	0.51155\\
3.6892	0.51582\\
3.6849	0.5201\\
3.6807	0.5244\\
3.6764	0.5287\\
3.6722	0.53302\\
3.6679	0.53735\\
3.6636	0.54169\\
3.6594	0.54604\\
3.6551	0.5504\\
3.6509	0.55476\\
3.6466	0.55914\\
3.6424	0.56353\\
3.6381	0.56793\\
3.6339	0.57233\\
3.6296	0.57674\\
3.6254	0.58116\\
3.6211	0.58558\\
3.6169	0.59001\\
3.6126	0.59445\\
3.6084	0.59889\\
3.6041	0.60334\\
3.5998	0.6078\\
3.5956	0.61226\\
3.5913	0.61672\\
3.5871	0.62119\\
3.5828	0.62567\\
3.5786	0.63015\\
3.5743	0.63464\\
3.5701	0.63912\\
3.5658	0.64362\\
3.5616	0.64812\\
3.5573	0.65262\\
3.5531	0.65712\\
3.5488	0.66162\\
3.5446	0.66612\\
3.5403	0.67062\\
3.536	0.67512\\
3.5318	0.67961\\
3.5275	0.68408\\
3.5233	0.68854\\
3.519	0.69297\\
3.5148	0.69738\\
3.5105	0.70176\\
3.5063	0.70611\\
3.502	0.71041\\
3.4978	0.71467\\
3.4935	0.71887\\
3.4893	0.72303\\
3.485	0.72713\\
3.4808	0.73118\\
3.4765	0.73517\\
3.4722	0.73912\\
3.468	0.74301\\
3.4637	0.74685\\
3.4595	0.75064\\
3.4552	0.7544\\
3.451	0.7581\\
3.4467	0.76177\\
3.4425	0.7654\\
3.4382	0.76899\\
3.434	0.77255\\
3.4297	0.77607\\
3.4255	0.77956\\
3.4212	0.783\\
3.417	0.78641\\
3.4127	0.78979\\
3.4085	0.79311\\
3.4042	0.79639\\
3.3999	0.79963\\
3.3957	0.8028\\
3.3914	0.80592\\
3.3872	0.80897\\
3.3829	0.81195\\
3.3787	0.81486\\
3.3744	0.81767\\
3.3702	0.8204\\
3.3659	0.82303\\
3.3617	0.82556\\
3.3574	0.82798\\
3.3532	0.8303\\
3.3489	0.83251\\
3.3447	0.8346\\
3.3404	0.83659\\
3.3361	0.83847\\
3.3319	0.84024\\
3.3276	0.84192\\
3.3234	0.84351\\
3.3191	0.845\\
3.3149	0.84641\\
3.3106	0.84773\\
3.3064	0.84898\\
3.3021	0.85015\\
3.2979	0.85125\\
3.2936	0.85228\\
3.2894	0.85323\\
3.2851	0.85412\\
3.2809	0.85493\\
3.2766	0.85567\\
3.2723	0.85634\\
3.2681	0.85694\\
3.2638	0.85747\\
3.2596	0.85792\\
3.2553	0.85829\\
3.2511	0.85859\\
3.2468	0.85882\\
3.2426	0.85896\\
3.2383	0.85903\\
3.2341	0.85902\\
3.2298	0.85893\\
3.2256	0.85876\\
3.2213	0.8585\\
3.2171	0.85816\\
3.2128	0.85775\\
3.2085	0.85724\\
3.2043	0.85666\\
3.2	0.85598\\
3.1958	0.85522\\
3.1915	0.85438\\
3.1873	0.85345\\
3.183	0.85242\\
3.1788	0.85131\\
3.1745	0.85011\\
3.1703	0.84881\\
3.166	0.84742\\
3.1618	0.84593\\
3.1575	0.84434\\
3.1533	0.84266\\
3.149	0.84087\\
3.1447	0.83898\\
3.1405	0.83699\\
3.1362	0.8349\\
3.132	0.8327\\
3.1277	0.83039\\
3.1235	0.82799\\
3.1192	0.82547\\
3.115	0.82286\\
3.1107	0.82014\\
3.1065	0.81733\\
3.1022	0.81441\\
3.098	0.8114\\
3.0937	0.8083\\
3.0895	0.80511\\
3.0852	0.80183\\
3.0809	0.79846\\
3.0767	0.79502\\
3.0724	0.7915\\
3.0682	0.7879\\
3.0639	0.78423\\
3.0597	0.78049\\
3.0554	0.77668\\
3.0512	0.77281\\
3.0469	0.76887\\
3.0427	0.76488\\
3.0384	0.76083\\
3.0342	0.75672\\
3.0299	0.75256\\
3.0257	0.74835\\
3.0214	0.74408\\
3.0171	0.73977\\
3.0129	0.7354\\
3.0086	0.73099\\
3.0044	0.72652\\
3.0001	0.72201\\
2.9959	0.71745\\
2.9916	0.71284\\
2.9874	0.70819\\
2.9831	0.70349\\
2.9789	0.69874\\
2.9746	0.69394\\
2.9704	0.6891\\
2.9661	0.68421\\
2.9619	0.67927\\
2.9576	0.67429\\
2.9533	0.66926\\
2.9491	0.66419\\
2.9448	0.65908\\
2.9406	0.65393\\
2.9363	0.64874\\
2.9321	0.64352\\
2.9278	0.63826\\
2.9236	0.63298\\
2.9193	0.62767\\
2.9151	0.62234\\
2.9108	0.617\\
2.9066	0.61164\\
2.9023	0.60626\\
2.8981	0.60088\\
2.8938	0.5955\\
2.8896	0.59011\\
2.8853	0.58472\\
2.881	0.57934\\
2.8768	0.57396\\
2.8725	0.56858\\
2.8683	0.56322\\
2.864	0.55786\\
2.8598	0.55252\\
2.8555	0.54719\\
2.8513	0.54187\\
2.847	0.53656\\
2.8428	0.53128\\
2.8385	0.526\\
2.8343	0.52075\\
2.83	0.51551\\
2.8258	0.51029\\
2.8215	0.50508\\
2.8172	0.4999\\
2.813	0.49474\\
2.8087	0.48959\\
2.8045	0.48447\\
2.8002	0.47937\\
2.796	0.47428\\
2.7917	0.46922\\
2.7875	0.46418\\
2.7832	0.45917\\
2.779	0.45418\\
2.7747	0.44921\\
2.7705	0.44426\\
2.7662	0.43934\\
2.762	0.43444\\
2.7577	0.42957\\
2.7534	0.42472\\
2.7492	0.4199\\
2.7449	0.4151\\
2.7407	0.41033\\
2.7364	0.40559\\
2.7322	0.40088\\
2.7279	0.39619\\
2.7237	0.39153\\
2.7194	0.3869\\
2.7152	0.38229\\
2.7109	0.37772\\
2.7067	0.37317\\
2.7024	0.36866\\
2.6982	0.36417\\
2.6939	0.35971\\
2.6896	0.35529\\
2.6854	0.35089\\
2.6811	0.34653\\
2.6769	0.34219\\
2.6726	0.33789\\
2.6684	0.33362\\
2.6641	0.32938\\
2.6599	0.32517\\
2.6556	0.32099\\
2.6514	0.31685\\
2.6471	0.31274\\
2.6429	0.30866\\
2.6386	0.30461\\
2.6344	0.3006\\
2.6301	0.29662\\
2.6258	0.29267\\
2.6216	0.28876\\
2.6173	0.28487\\
2.6131	0.28103\\
2.6088	0.27721\\
2.6046	0.27343\\
2.6003	0.26968\\
2.5961	0.26597\\
2.5918	0.26229\\
2.5876	0.25864\\
2.5833	0.25503\\
2.5791	0.25145\\
2.5748	0.2479\\
2.5706	0.24439\\
2.5663	0.24091\\
2.562	0.23747\\
2.5578	0.23406\\
2.5535	0.23068\\
2.5493	0.22734\\
2.545	0.22403\\
2.5408	0.22075\\
2.5365	0.21751\\
2.5323	0.2143\\
2.528	0.21112\\
2.5238	0.20798\\
2.5195	0.20487\\
2.5153	0.20179\\
2.511	0.19875\\
2.5068	0.19574\\
2.5025	0.19276\\
2.4982	0.18982\\
2.494	0.18691\\
2.4897	0.18403\\
2.4855	0.18118\\
2.4812	0.17837\\
2.477	0.17559\\
2.4727	0.17284\\
2.4685	0.17012\\
2.4642	0.16744\\
2.46	0.16478\\
2.4557	0.16216\\
2.4515	0.15957\\
2.4472	0.15701\\
2.443	0.15448\\
2.4387	0.15199\\
2.4344	0.14952\\
2.4302	0.14709\\
2.4259	0.14468\\
2.4217	0.14231\\
2.4174	0.13996\\
2.4132	0.13765\\
2.4089	0.13536\\
2.4047	0.13311\\
2.4004	0.13088\\
2.3962	0.12869\\
2.3919	0.12652\\
2.3877	0.12438\\
2.3834	0.12227\\
2.3792	0.12019\\
2.3749	0.11814\\
2.3706	0.11612\\
2.3664	0.11412\\
2.3621	0.11215\\
2.3579	0.11021\\
2.3536	0.10829\\
2.3494	0.10641\\
2.3451	0.10455\\
2.3409	0.10271\\
2.3366	0.10091\\
2.3324	0.099124\\
2.3281	0.097368\\
2.3239	0.095638\\
2.3196	0.093934\\
2.3154	0.092254\\
2.3111	0.0906\\
2.3069	0.08897\\
2.3026	0.087365\\
2.2983	0.085784\\
2.2941	0.084227\\
2.2898	0.082694\\
2.2856	0.081185\\
2.2813	0.079699\\
2.2771	0.078235\\
2.2728	0.076795\\
2.2686	0.075377\\
2.2643	0.073982\\
2.2601	0.072608\\
2.2558	0.071257\\
2.2516	0.069927\\
2.2473	0.068618\\
2.2431	0.067331\\
2.2388	0.066064\\
2.2345	0.064818\\
2.2303	0.063592\\
2.226	0.062386\\
2.2218	0.061201\\
2.2175	0.060034\\
2.2133	0.058888\\
2.209	0.05776\\
2.2048	0.056651\\
2.2005	0.055561\\
2.1963	0.05449\\
2.192	0.053436\\
2.1878	0.052401\\
2.1835	0.051383\\
2.1793	0.050383\\
2.175	0.0494\\
2.1707	0.048434\\
2.1665	0.047484\\
2.1622	0.046552\\
2.158	0.045635\\
2.1537	0.044735\\
2.1495	0.043851\\
2.1452	0.042982\\
2.141	0.042128\\
2.1367	0.04129\\
2.1325	0.040467\\
2.1282	0.039658\\
2.124	0.038865\\
2.1197	0.038085\\
2.1155	0.037319\\
2.1112	0.036568\\
2.1069	0.03583\\
2.1027	0.035106\\
2.0984	0.034395\\
2.0942	0.033696\\
2.0899	0.033011\\
2.0857	0.032339\\
2.0814	0.031679\\
2.0772	0.031031\\
2.0729	0.030395\\
2.0687	0.029772\\
2.0644	0.02916\\
2.0602	0.028559\\
2.0559	0.02797\\
2.0517	0.027392\\
2.0474	0.026825\\
2.0431	0.026269\\
2.0389	0.025723\\
2.0346	0.025188\\
2.0304	0.024663\\
2.0261	0.024148\\
2.0219	0.023643\\
2.0176	0.023148\\
2.0134	0.022662\\
2.0091	0.022186\\
2.0049	0.021719\\
2.0006	0.021262\\
1.9964	0.020813\\
1.9921	0.020373\\
1.9879	0.019941\\
1.9836	0.019518\\
1.9793	0.019104\\
1.9751	0.018698\\
1.9708	0.018299\\
1.9666	0.017909\\
1.9623	0.017526\\
1.9581	0.017151\\
1.9538	0.016784\\
1.9496	0.016424\\
1.9453	0.016071\\
1.9411	0.015725\\
1.9368	0.015386\\
1.9326	0.015055\\
1.9283	0.014729\\
1.9241	0.014411\\
1.9198	0.014099\\
1.9155	0.013793\\
1.9113	0.013493\\
1.907	0.0132\\
1.9028	0.012912\\
1.8985	0.012631\\
1.8943	0.012355\\
1.89	0.012085\\
1.8858	0.011821\\
1.8815	0.011562\\
1.8773	0.011308\\
1.873	0.011059\\
1.8688	0.010816\\
1.8645	0.010578\\
1.8603	0.010345\\
1.856	0.010116\\
1.8517	0.0098928\\
1.8475	0.0096739\\
1.8432	0.0094596\\
1.839	0.0092498\\
1.8347	0.0090444\\
1.8305	0.0088434\\
1.8262	0.0086465\\
1.822	0.0084539\\
1.8177	0.0082653\\
1.8135	0.0080807\\
1.8092	0.0079\\
1.805	0.0077232\\
1.8007	0.0075501\\
1.7965	0.0073807\\
1.7922	0.007215\\
1.788	0.0070528\\
1.7837	0.006894\\
1.7794	0.0067387\\
1.7752	0.0065867\\
1.7709	0.0064379\\
1.7667	0.0062924\\
1.7624	0.00615\\
1.7582	0.0060107\\
1.7539	0.0058744\\
1.7497	0.005741\\
1.7454	0.0056105\\
1.7412	0.0054829\\
1.7369	0.005358\\
1.7327	0.0052359\\
1.7284	0.0051164\\
1.7242	0.0049995\\
1.7199	0.0048852\\
1.7156	0.0047733\\
1.7114	0.0046639\\
1.7071	0.0045569\\
1.7029	0.0044523\\
1.6986	0.0043499\\
1.6944	0.0042498\\
1.6901	0.0041519\\
1.6859	0.0040562\\
1.6816	0.0039626\\
1.6774	0.003871\\
1.6731	0.0037815\\
1.6689	0.0036939\\
1.6646	0.0036083\\
1.6604	0.0035246\\
1.6561	0.0034428\\
1.6518	0.0033627\\
1.6476	0.0032845\\
1.6433	0.003208\\
1.6391	0.0031331\\
1.6348	0.00306\\
1.6306	0.0029885\\
1.6263	0.0029186\\
1.6221	0.0028503\\
1.6178	0.0027835\\
1.6136	0.0027182\\
1.6093	0.0026543\\
1.6051	0.0025919\\
1.6008	0.002531\\
1.5966	0.0024713\\
1.5923	0.0024131\\
1.588	0.0023561\\
1.5838	0.0023004\\
1.5795	0.002246\\
1.5753	0.0021928\\
1.571	0.0021409\\
1.5668	0.0020901\\
1.5625	0.0020404\\
1.5583	0.0019919\\
1.554	0.0019445\\
1.5498	0.0018982\\
1.5455	0.0018529\\
1.5413	0.0018087\\
1.537	0.0017655\\
1.5328	0.0017232\\
1.5285	0.001682\\
1.5242	0.0016417\\
1.52	0.0016023\\
1.5157	0.0015638\\
1.5115	0.0015262\\
1.5072	0.0014894\\
1.503	0.0014535\\
1.4987	0.0014185\\
1.4945	0.0013842\\
1.4902	0.0013508\\
1.486	0.0013181\\
1.4817	0.0012861\\
1.4775	0.0012549\\
1.4732	0.0012245\\
1.469	0.0011947\\
1.4647	0.0011657\\
1.4604	0.0011373\\
1.4562	0.0011095\\
1.4519	0.0010824\\
1.4477	0.001056\\
1.4434	0.0010302\\
1.4392	0.0010049\\
1.4349	0.0009803\\
1.4307	0.00095624\\
1.4264	0.00093274\\
1.4222	0.0009098\\
1.4179	0.0008874\\
1.4137	0.00086552\\
1.4094	0.00084416\\
1.4052	0.0008233\\
1.4009	0.00080294\\
1.3966	0.00078305\\
1.3924	0.00076364\\
1.3881	0.00074469\\
1.3839	0.00072618\\
1.3796	0.00070812\\
1.3754	0.00069048\\
1.3711	0.00067327\\
1.3669	0.00065646\\
1.3626	0.00064006\\
1.3584	0.00062404\\
1.3541	0.00060841\\
1.3499	0.00059315\\
1.3456	0.00057826\\
1.3414	0.00056373\\
1.3371	0.00054954\\
1.3328	0.00053569\\
1.3286	0.00052218\\
1.3243	0.00050899\\
1.3201	0.00049613\\
1.3158	0.00048357\\
1.3116	0.00047131\\
1.3073	0.00045935\\
1.3031	0.00044768\\
1.2988	0.00043629\\
1.2946	0.00042518\\
1.2903	0.00041434\\
1.2861	0.00040377\\
1.2818	0.00039345\\
1.2776	0.00038338\\
1.2733	0.00037356\\
1.2691	0.00036397\\
1.2648	0.00035463\\
1.2605	0.00034551\\
1.2563	0.00033661\\
1.252	0.00032793\\
1.2478	0.00031947\\
1.2435	0.00031121\\
1.2393	0.00030316\\
1.235	0.00029531\\
1.2308	0.00028765\\
1.2265	0.00028018\\
1.2223	0.00027289\\
1.218	0.00026579\\
1.2138	0.00025886\\
1.2095	0.0002521\\
1.2053	0.00024552\\
1.201	0.00023909\\
1.1967	0.00023283\\
1.1925	0.00022672\\
1.1882	0.00022076\\
1.184	0.00021496\\
1.1797	0.0002093\\
1.1755	0.00020378\\
1.1712	0.0001984\\
1.167	0.00019315\\
1.1627	0.00018804\\
1.1585	0.00018306\\
1.1542	0.0001782\\
1.15	0.00017347\\
1.1457	0.00016885\\
1.1415	0.00016435\\
1.1372	0.00015997\\
1.1329	0.0001557\\
1.1287	0.00015153\\
1.1244	0.00014748\\
1.1202	0.00014352\\
1.1159	0.00013967\\
1.1117	0.00013591\\
}--cycle;
      % \input{../../../../.tex/20/KS/SizeDistributionFittingsfig28.tex}
      \legend{}
  
      \nextgroupplot[%
        xmin=0.57363,xmax=5.3544,
        % xmin=0.89973,xmax=5.5625,
        % ymin=0,ymax=1,
      ] % This file was created by matlab2tikz.
%
\definecolor{mycolor1}{rgb}{0.00000,0.44706,0.69804}%
%
\addplot[ybar interval, fill=mycolor1, fill opacity=0.35, area legend, draw=none] table[row sep=crcr, x=Lower, y=Count] {%
Lower	Upper	Count\\
1.1117	1.4377	0.00032713\\
1.4377	1.7638	0.001881\\
1.7638	2.0899	0.015702\\
2.0899	2.416	0.074094\\
2.416	2.742	0.27659\\
2.742	3.0681	0.61034\\
3.0681	3.3942	0.84636\\
3.3942	3.7202	0.67151\\
3.7202	4.0463	0.338\\
4.0463	4.3724	0.15211\\
4.3724	4.6985	0.057738\\
4.6985	5.0245	0.014148\\
5.0245	5.3506	0.0080146\\
5.3506	5.3506	0.0080146\\
};
\addlegendentry{\ExactHistogramLegend}

\addplot [color=mycolor1, only marks, every error bar/.append style={opacity=0.45}, mark=*, mark size=0pt, draw=none, forget plot]
 plot [error bars/.cd, y dir=both, y explicit, error bar style={line width=1pt, color=mycolor1}, error mark options={mark=none,mark size=0pt}]
 table[row sep=crcr, y error plus index=2, y error minus index=3]{%
1.2747	0.00032713	0.00031959	0.00031959\\
1.6008	0.001881	0.0007599	0.0007599\\
1.9269	0.015702	0.0022084	0.0022084\\
2.2529	0.074094	0.0042204	0.0042204\\
2.579	0.27659	0.0084725	0.0084725\\
2.9051	0.61034	0.012613	0.012613\\
3.2311	0.84636	0.012737	0.012737\\
3.5572	0.67151	0.011938	0.011938\\
3.8833	0.338	0.0089062	0.0089062\\
4.2093	0.15211	0.0057753	0.0057753\\
4.5354	0.057738	0.002979	0.002979\\
4.8615	0.014148	0.0014191	0.0014191\\
5.1876	0.0080146	0.00022833	0.00022833\\
};
\addplot [color=mycolor1]
  table[row sep=crcr]{%
1.1117	0.00013532\\
1.3917	0.00074235\\
1.5445	0.0017956\\
1.6506	0.003259\\
1.7312	0.0050794\\
1.7991	0.0073342\\
1.8585	0.010064\\
1.9094	0.013147\\
1.9561	0.016737\\
1.9985	0.020781\\
2.0367	0.02518\\
2.0706	0.029797\\
2.1046	0.035177\\
2.1343	0.040594\\
2.164	0.046752\\
2.1937	0.05373\\
2.2191	0.060429\\
2.2446	0.067848\\
2.2701	0.076044\\
2.2955	0.085074\\
2.321	0.094997\\
2.3464	0.10587\\
2.3677	0.1157\\
2.3889	0.12625\\
2.4101	0.13757\\
2.4313	0.14967\\
2.4525	0.16258\\
2.4737	0.17633\\
2.4949	0.19094\\
2.5162	0.20641\\
2.5374	0.22277\\
2.5586	0.24001\\
2.5798	0.25814\\
2.601	0.27714\\
2.6265	0.30109\\
2.6519	0.32626\\
2.6774	0.35258\\
2.7029	0.37998\\
2.7326	0.41321\\
2.7623	0.44762\\
2.7962	0.48811\\
2.8429	0.54513\\
2.9193	0.63861\\
2.949	0.67373\\
2.9744	0.70268\\
2.9956	0.72573\\
3.0169	0.74758\\
3.0338	0.76408\\
3.0508	0.77963\\
3.0678	0.79416\\
3.0847	0.80765\\
3.1017	0.8201\\
3.1187	0.83148\\
3.1357	0.84177\\
3.1484	0.84873\\
3.1611	0.85497\\
3.1739	0.86042\\
3.1866	0.86501\\
3.1993	0.86868\\
3.212	0.87136\\
3.2248	0.87303\\
3.2375	0.8737\\
3.2502	0.87342\\
3.263	0.87222\\
3.2757	0.87017\\
3.2884	0.86732\\
3.3011	0.86371\\
3.3139	0.85936\\
3.3266	0.85427\\
3.3393	0.84843\\
3.3521	0.84185\\
3.3648	0.83452\\
3.3775	0.82648\\
3.3945	0.81467\\
3.4115	0.80175\\
3.4284	0.78783\\
3.4497	0.76924\\
3.4709	0.74954\\
3.4963	0.72475\\
3.526	0.69462\\
3.5642	0.65468\\
3.6745	0.53853\\
3.7127	0.49964\\
3.7467	0.46604\\
3.7806	0.43351\\
3.8103	0.40605\\
3.84	0.37962\\
3.8697	0.35428\\
3.8952	0.33348\\
3.9207	0.31354\\
3.9461	0.29448\\
3.9716	0.2763\\
3.997	0.259\\
4.0225	0.24256\\
4.0479	0.22698\\
4.0734	0.21223\\
4.0989	0.19829\\
4.1243	0.18513\\
4.1498	0.17272\\
4.1795	0.15916\\
4.2092	0.14652\\
4.2389	0.13476\\
4.2686	0.12383\\
4.2983	0.11368\\
4.328	0.10425\\
4.3619	0.094299\\
4.3959	0.085167\\
4.4298	0.076791\\
4.468	0.068198\\
4.5062	0.060411\\
4.5444	0.053362\\
4.5868	0.046322\\
4.6293	0.040041\\
4.6759	0.033933\\
4.7226	0.028586\\
4.7735	0.023539\\
4.8244	0.01923\\
4.8796	0.015306\\
4.939	0.011842\\
5.0027	0.0088829\\
5.0748	0.0063159\\
5.1554	0.0042359\\
5.2488	0.0026081\\
5.3506	0.0014989\\
};
\addlegendentry{\ExactBLNFitLegend}


\addplot[area legend, draw=none, fill=mycolor1, fill opacity=0.15, forget plot]
table[row sep=crcr] {%
x	y\\
1.1117	0.000171\\
1.1159	0.00017546\\
1.1202	0.00018004\\
1.1244	0.00018474\\
1.1287	0.00018955\\
1.1329	0.00019448\\
1.1371	0.00019954\\
1.1414	0.00020472\\
1.1456	0.00021003\\
1.1499	0.00021548\\
1.1541	0.00022106\\
1.1584	0.00022678\\
1.1626	0.00023264\\
1.1668	0.00023865\\
1.1711	0.00024481\\
1.1753	0.00025112\\
1.1796	0.00025759\\
1.1838	0.00026422\\
1.1881	0.00027101\\
1.1923	0.00027797\\
1.1965	0.0002851\\
1.2008	0.00029241\\
1.205	0.0002999\\
1.2093	0.00030757\\
1.2135	0.00031543\\
1.2178	0.00032349\\
1.222	0.00033174\\
1.2262	0.0003402\\
1.2305	0.00034886\\
1.2347	0.00035774\\
1.239	0.00036683\\
1.2432	0.00037615\\
1.2475	0.00038569\\
1.2517	0.00039547\\
1.2559	0.00040548\\
1.2602	0.00041574\\
1.2644	0.00042625\\
1.2687	0.00043701\\
1.2729	0.00044803\\
1.2772	0.00045932\\
1.2814	0.00047089\\
1.2856	0.00048273\\
1.2899	0.00049486\\
1.2941	0.00050728\\
1.2984	0.00052001\\
1.3026	0.00053304\\
1.3069	0.00054638\\
1.3111	0.00056004\\
1.3154	0.00057403\\
1.3196	0.00058836\\
1.3238	0.00060303\\
1.3281	0.00061806\\
1.3323	0.00063344\\
1.3366	0.00064919\\
1.3408	0.00066531\\
1.3451	0.00068182\\
1.3493	0.00069873\\
1.3535	0.00071603\\
1.3578	0.00073375\\
1.362	0.00075189\\
1.3663	0.00077046\\
1.3705	0.00078947\\
1.3748	0.00080892\\
1.379	0.00082884\\
1.3832	0.00084923\\
1.3875	0.0008701\\
1.3917	0.00089147\\
1.396	0.00091333\\
1.4002	0.00093571\\
1.4045	0.00095862\\
1.4087	0.00098206\\
1.4129	0.0010061\\
1.4172	0.0010306\\
1.4214	0.0010557\\
1.4257	0.0010815\\
1.4299	0.0011078\\
1.4342	0.0011347\\
1.4384	0.0011623\\
1.4426	0.0011905\\
1.4469	0.0012193\\
1.4511	0.0012489\\
1.4554	0.0012791\\
1.4596	0.00131\\
1.4639	0.0013416\\
1.4681	0.001374\\
1.4723	0.0014071\\
1.4766	0.001441\\
1.4808	0.0014757\\
1.4851	0.0015111\\
1.4893	0.0015474\\
1.4936	0.0015845\\
1.4978	0.0016225\\
1.502	0.0016613\\
1.5063	0.001701\\
1.5105	0.0017417\\
1.5148	0.0017832\\
1.519	0.0018258\\
1.5233	0.0018692\\
1.5275	0.0019137\\
1.5318	0.0019592\\
1.536	0.0020057\\
1.5402	0.0020533\\
1.5445	0.002102\\
1.5487	0.0021518\\
1.553	0.0022027\\
1.5572	0.0022547\\
1.5615	0.0023079\\
1.5657	0.0023624\\
1.5699	0.002418\\
1.5742	0.0024749\\
1.5784	0.0025331\\
1.5827	0.0025926\\
1.5869	0.0026535\\
1.5912	0.0027157\\
1.5954	0.0027793\\
1.5996	0.0028443\\
1.6039	0.0029108\\
1.6081	0.0029788\\
1.6124	0.0030483\\
1.6166	0.0031193\\
1.6209	0.0031919\\
1.6251	0.0032661\\
1.6293	0.003342\\
1.6336	0.0034196\\
1.6378	0.0034989\\
1.6421	0.0035799\\
1.6463	0.0036628\\
1.6506	0.0037475\\
1.6548	0.003834\\
1.659	0.0039225\\
1.6633	0.0040129\\
1.6675	0.0041052\\
1.6718	0.0041997\\
1.676	0.0042962\\
1.6803	0.0043948\\
1.6845	0.0044955\\
1.6887	0.0045985\\
1.693	0.0047037\\
1.6972	0.0048113\\
1.7015	0.0049211\\
1.7057	0.0050334\\
1.71	0.0051481\\
1.7142	0.0052653\\
1.7185	0.005385\\
1.7227	0.0055074\\
1.7269	0.0056323\\
1.7312	0.00576\\
1.7354	0.0058905\\
1.7397	0.0060237\\
1.7439	0.0061598\\
1.7482	0.0062989\\
1.7524	0.0064409\\
1.7566	0.006586\\
1.7609	0.0067342\\
1.7651	0.0068856\\
1.7694	0.0070402\\
1.7736	0.0071981\\
1.7779	0.0073594\\
1.7821	0.0075241\\
1.7863	0.0076923\\
1.7906	0.0078641\\
1.7948	0.0080395\\
1.7991	0.0082186\\
1.8033	0.0084016\\
1.8076	0.0085884\\
1.8118	0.0087791\\
1.816	0.0089738\\
1.8203	0.0091727\\
1.8245	0.0093757\\
1.8288	0.009583\\
1.833	0.0097946\\
1.8373	0.010011\\
1.8415	0.010231\\
1.8457	0.010456\\
1.85	0.010686\\
1.8542	0.010921\\
1.8585	0.01116\\
1.8627	0.011405\\
1.867	0.011654\\
1.8712	0.011909\\
1.8754	0.012169\\
1.8797	0.012434\\
1.8839	0.012705\\
1.8882	0.012981\\
1.8924	0.013263\\
1.8967	0.013551\\
1.9009	0.013844\\
1.9051	0.014144\\
1.9094	0.014449\\
1.9136	0.014761\\
1.9179	0.01508\\
1.9221	0.015404\\
1.9264	0.015735\\
1.9306	0.016073\\
1.9349	0.016418\\
1.9391	0.01677\\
1.9433	0.017128\\
1.9476	0.017494\\
1.9518	0.017867\\
1.9561	0.018248\\
1.9603	0.018636\\
1.9646	0.019032\\
1.9688	0.019436\\
1.973	0.019848\\
1.9773	0.020267\\
1.9815	0.020696\\
1.9858	0.021132\\
1.99	0.021577\\
1.9943	0.022031\\
1.9985	0.022494\\
2.0027	0.022966\\
2.007	0.023447\\
2.0112	0.023937\\
2.0155	0.024437\\
2.0197	0.024946\\
2.024	0.025466\\
2.0282	0.025995\\
2.0324	0.026535\\
2.0367	0.027085\\
2.0409	0.027645\\
2.0452	0.028216\\
2.0494	0.028798\\
2.0537	0.029391\\
2.0579	0.029995\\
2.0621	0.030611\\
2.0664	0.031238\\
2.0706	0.031877\\
2.0749	0.032528\\
2.0791	0.033191\\
2.0834	0.033867\\
2.0876	0.034555\\
2.0918	0.035256\\
2.0961	0.03597\\
2.1003	0.036697\\
2.1046	0.037437\\
2.1088	0.038191\\
2.1131	0.038959\\
2.1173	0.039741\\
2.1216	0.040537\\
2.1258	0.041348\\
2.13	0.042173\\
2.1343	0.043013\\
2.1385	0.043868\\
2.1428	0.044739\\
2.147	0.045625\\
2.1513	0.046527\\
2.1555	0.047445\\
2.1597	0.048379\\
2.164	0.04933\\
2.1682	0.050297\\
2.1725	0.051281\\
2.1767	0.052283\\
2.181	0.053302\\
2.1852	0.054339\\
2.1894	0.055394\\
2.1937	0.056467\\
2.1979	0.057558\\
2.2022	0.058668\\
2.2064	0.059797\\
2.2107	0.060946\\
2.2149	0.062114\\
2.2191	0.063301\\
2.2234	0.064509\\
2.2276	0.065736\\
2.2319	0.066985\\
2.2361	0.068254\\
2.2404	0.069544\\
2.2446	0.070855\\
2.2488	0.072188\\
2.2531	0.073543\\
2.2573	0.07492\\
2.2616	0.076319\\
2.2658	0.077741\\
2.2701	0.079186\\
2.2743	0.080654\\
2.2785	0.082146\\
2.2828	0.083661\\
2.287	0.085201\\
2.2913	0.086764\\
2.2955	0.088353\\
2.2998	0.089966\\
2.304	0.091604\\
2.3082	0.093268\\
2.3125	0.094957\\
2.3167	0.096672\\
2.321	0.098414\\
2.3252	0.10018\\
2.3295	0.10198\\
2.3337	0.1038\\
2.338	0.10565\\
2.3422	0.10752\\
2.3464	0.10943\\
2.3507	0.11136\\
2.3549	0.11332\\
2.3592	0.11531\\
2.3634	0.11733\\
2.3677	0.11938\\
2.3719	0.12145\\
2.3761	0.12356\\
2.3804	0.1257\\
2.3846	0.12787\\
2.3889	0.13006\\
2.3931	0.13229\\
2.3974	0.13455\\
2.4016	0.13684\\
2.4058	0.13916\\
2.4101	0.14152\\
2.4143	0.1439\\
2.4186	0.14632\\
2.4228	0.14877\\
2.4271	0.15125\\
2.4313	0.15377\\
2.4355	0.15631\\
2.4398	0.15889\\
2.444	0.16151\\
2.4483	0.16416\\
2.4525	0.16684\\
2.4568	0.16956\\
2.461	0.17231\\
2.4652	0.17509\\
2.4695	0.17791\\
2.4737	0.18076\\
2.478	0.18365\\
2.4822	0.18657\\
2.4865	0.18953\\
2.4907	0.19252\\
2.4949	0.19555\\
2.4992	0.19862\\
2.5034	0.20172\\
2.5077	0.20485\\
2.5119	0.20802\\
2.5162	0.21123\\
2.5204	0.21447\\
2.5247	0.21775\\
2.5289	0.22106\\
2.5331	0.22441\\
2.5374	0.2278\\
2.5416	0.23122\\
2.5459	0.23468\\
2.5501	0.23817\\
2.5544	0.24171\\
2.5586	0.24527\\
2.5628	0.24887\\
2.5671	0.25251\\
2.5713	0.25619\\
2.5756	0.2599\\
2.5798	0.26364\\
2.5841	0.26742\\
2.5883	0.27124\\
2.5925	0.27509\\
2.5968	0.27898\\
2.601	0.28291\\
2.6053	0.28686\\
2.6095	0.29086\\
2.6138	0.29488\\
2.618	0.29895\\
2.6222	0.30304\\
2.6265	0.30717\\
2.6307	0.31134\\
2.635	0.31554\\
2.6392	0.31977\\
2.6435	0.32403\\
2.6477	0.32833\\
2.6519	0.33266\\
2.6562	0.33703\\
2.6604	0.34142\\
2.6647	0.34585\\
2.6689	0.3503\\
2.6732	0.35479\\
2.6774	0.35931\\
2.6816	0.36386\\
2.6859	0.36844\\
2.6901	0.37305\\
2.6944	0.37769\\
2.6986	0.38235\\
2.7029	0.38705\\
2.7071	0.39177\\
2.7113	0.39652\\
2.7156	0.40129\\
2.7198	0.40609\\
2.7241	0.41092\\
2.7283	0.41577\\
2.7326	0.42065\\
2.7368	0.42555\\
2.7411	0.43047\\
2.7453	0.43542\\
2.7495	0.44038\\
2.7538	0.44537\\
2.758	0.45038\\
2.7623	0.45541\\
2.7665	0.46046\\
2.7708	0.46553\\
2.775	0.47062\\
2.7792	0.47572\\
2.7835	0.48084\\
2.7877	0.48598\\
2.792	0.49113\\
2.7962	0.4963\\
2.8005	0.50148\\
2.8047	0.50667\\
2.8089	0.51187\\
2.8132	0.51709\\
2.8174	0.52231\\
2.8217	0.52755\\
2.8259	0.53279\\
2.8302	0.53804\\
2.8344	0.5433\\
2.8386	0.54856\\
2.8429	0.55383\\
2.8471	0.55911\\
2.8514	0.56438\\
2.8556	0.56966\\
2.8599	0.57494\\
2.8641	0.58022\\
2.8683	0.5855\\
2.8726	0.59078\\
2.8768	0.59605\\
2.8811	0.60132\\
2.8853	0.60659\\
2.8896	0.61185\\
2.8938	0.6171\\
2.898	0.62234\\
2.9023	0.62758\\
2.9065	0.6328\\
2.9108	0.63801\\
2.915	0.64321\\
2.9193	0.6484\\
2.9235	0.65356\\
2.9278	0.65872\\
2.932	0.66385\\
2.9362	0.66896\\
2.9405	0.67405\\
2.9447	0.67912\\
2.949	0.68416\\
2.9532	0.68917\\
2.9575	0.69416\\
2.9617	0.69911\\
2.9659	0.70403\\
2.9702	0.70892\\
2.9744	0.71378\\
2.9787	0.71859\\
2.9829	0.72336\\
2.9872	0.72809\\
2.9914	0.73278\\
2.9956	0.73742\\
2.9999	0.74201\\
3.0041	0.74655\\
3.0084	0.75104\\
3.0126	0.75547\\
3.0169	0.75984\\
3.0211	0.76416\\
3.0253	0.76841\\
3.0296	0.7726\\
3.0338	0.77673\\
3.0381	0.78079\\
3.0423	0.78478\\
3.0466	0.7887\\
3.0508	0.79256\\
3.055	0.79633\\
3.0593	0.80004\\
3.0635	0.80367\\
3.0678	0.80722\\
3.072	0.8107\\
3.0763	0.8141\\
3.0805	0.81742\\
3.0847	0.82066\\
3.089	0.82383\\
3.0932	0.82692\\
3.0975	0.82994\\
3.1017	0.83288\\
3.106	0.83574\\
3.1102	0.83853\\
3.1145	0.84125\\
3.1187	0.84391\\
3.1229	0.84649\\
3.1272	0.84902\\
3.1314	0.85147\\
3.1357	0.85387\\
3.1399	0.85621\\
3.1442	0.8585\\
3.1484	0.86072\\
3.1526	0.86288\\
3.1569	0.86499\\
3.1611	0.86703\\
3.1654	0.869\\
3.1696	0.8709\\
3.1739	0.87272\\
3.1781	0.87445\\
3.1823	0.87609\\
3.1866	0.87762\\
3.1908	0.87905\\
3.1951	0.88036\\
3.1993	0.88154\\
3.2036	0.8826\\
3.2078	0.88353\\
3.212	0.88433\\
3.2163	0.88499\\
3.2205	0.88551\\
3.2248	0.8859\\
3.229	0.88616\\
3.2333	0.8863\\
3.2375	0.88632\\
3.2417	0.88622\\
3.246	0.88602\\
3.2502	0.88573\\
3.2545	0.88534\\
3.2587	0.88487\\
3.263	0.88432\\
3.2672	0.8837\\
3.2714	0.88302\\
3.2757	0.88227\\
3.2799	0.88145\\
3.2842	0.88058\\
3.2884	0.87963\\
3.2927	0.87862\\
3.2969	0.87754\\
3.3011	0.87638\\
3.3054	0.87514\\
3.3096	0.87382\\
3.3139	0.87242\\
3.3181	0.87092\\
3.3224	0.86933\\
3.3266	0.86764\\
3.3309	0.86585\\
3.3351	0.86396\\
3.3393	0.86197\\
3.3436	0.85988\\
3.3478	0.85768\\
3.3521	0.85538\\
3.3563	0.85297\\
3.3606	0.85047\\
3.3648	0.84787\\
3.369	0.84517\\
3.3733	0.84238\\
3.3775	0.8395\\
3.3818	0.83653\\
3.386	0.83348\\
3.3903	0.83035\\
3.3945	0.82715\\
3.3987	0.82387\\
3.403	0.82053\\
3.4072	0.81712\\
3.4115	0.81365\\
3.4157	0.81013\\
3.42	0.80655\\
3.4242	0.80292\\
3.4284	0.79925\\
3.4327	0.79553\\
3.4369	0.79176\\
3.4412	0.78796\\
3.4454	0.78411\\
3.4497	0.78022\\
3.4539	0.77629\\
3.4581	0.77232\\
3.4624	0.76832\\
3.4666	0.76427\\
3.4709	0.76019\\
3.4751	0.75608\\
3.4794	0.75192\\
3.4836	0.74773\\
3.4878	0.74351\\
3.4921	0.73925\\
3.4963	0.73496\\
3.5006	0.73065\\
3.5048	0.7263\\
3.5091	0.72192\\
3.5133	0.71751\\
3.5176	0.71308\\
3.5218	0.70863\\
3.526	0.70416\\
3.5303	0.69967\\
3.5345	0.69516\\
3.5388	0.69063\\
3.543	0.6861\\
3.5473	0.68155\\
3.5515	0.67699\\
3.5557	0.67242\\
3.56	0.66784\\
3.5642	0.66327\\
3.5685	0.65868\\
3.5727	0.6541\\
3.577	0.64951\\
3.5812	0.64493\\
3.5854	0.64035\\
3.5897	0.63576\\
3.5939	0.63119\\
3.5982	0.62661\\
3.6024	0.62204\\
3.6067	0.61747\\
3.6109	0.61291\\
3.6151	0.60836\\
3.6194	0.60381\\
3.6236	0.59927\\
3.6279	0.59473\\
3.6321	0.5902\\
3.6364	0.58568\\
3.6406	0.58117\\
3.6448	0.57667\\
3.6491	0.57217\\
3.6533	0.56768\\
3.6576	0.56321\\
3.6618	0.55874\\
3.6661	0.55428\\
3.6703	0.54983\\
3.6745	0.54539\\
3.6788	0.54097\\
3.683	0.53655\\
3.6873	0.53215\\
3.6915	0.52775\\
3.6958	0.52337\\
3.7	0.51901\\
3.7042	0.51465\\
3.7085	0.51031\\
3.7127	0.50599\\
3.717	0.50168\\
3.7212	0.49738\\
3.7255	0.4931\\
3.7297	0.48883\\
3.734	0.48458\\
3.7382	0.48035\\
3.7424	0.47613\\
3.7467	0.47193\\
3.7509	0.46775\\
3.7552	0.46359\\
3.7594	0.45944\\
3.7637	0.45531\\
3.7679	0.4512\\
3.7721	0.44711\\
3.7764	0.44304\\
3.7806	0.43899\\
3.7849	0.43496\\
3.7891	0.43094\\
3.7934	0.42695\\
3.7976	0.42298\\
3.8018	0.41903\\
3.8061	0.41511\\
3.8103	0.4112\\
3.8146	0.40732\\
3.8188	0.40346\\
3.8231	0.39962\\
3.8273	0.3958\\
3.8315	0.392\\
3.8358	0.38823\\
3.84	0.38448\\
3.8443	0.38076\\
3.8485	0.37706\\
3.8528	0.37338\\
3.857	0.36972\\
3.8612	0.36609\\
3.8655	0.36248\\
3.8697	0.3589\\
3.874	0.35534\\
3.8782	0.35181\\
3.8825	0.3483\\
3.8867	0.34481\\
3.8909	0.34135\\
3.8952	0.33791\\
3.8994	0.3345\\
3.9037	0.33111\\
3.9079	0.32775\\
3.9122	0.32441\\
3.9164	0.3211\\
3.9207	0.31781\\
3.9249	0.31455\\
3.9291	0.31131\\
3.9334	0.30809\\
3.9376	0.30491\\
3.9419	0.30174\\
3.9461	0.2986\\
3.9504	0.29549\\
3.9546	0.2924\\
3.9588	0.28933\\
3.9631	0.28629\\
3.9673	0.28328\\
3.9716	0.28029\\
3.9758	0.27732\\
3.9801	0.27438\\
3.9843	0.27146\\
3.9885	0.26857\\
3.9928	0.2657\\
3.997	0.26286\\
4.0013	0.26004\\
4.0055	0.25724\\
4.0098	0.25447\\
4.014	0.25172\\
4.0182	0.249\\
4.0225	0.2463\\
4.0267	0.24362\\
4.031	0.24097\\
4.0352	0.23834\\
4.0395	0.23573\\
4.0437	0.23315\\
4.0479	0.23059\\
4.0522	0.22806\\
4.0564	0.22554\\
4.0607	0.22305\\
4.0649	0.22058\\
4.0692	0.21814\\
4.0734	0.21571\\
4.0776	0.21331\\
4.0819	0.21093\\
4.0861	0.20858\\
4.0904	0.20624\\
4.0946	0.20393\\
4.0989	0.20164\\
4.1031	0.19937\\
4.1073	0.19712\\
4.1116	0.19489\\
4.1158	0.19269\\
4.1201	0.1905\\
4.1243	0.18834\\
4.1286	0.1862\\
4.1328	0.18407\\
4.1371	0.18197\\
4.1413	0.17989\\
4.1455	0.17783\\
4.1498	0.17578\\
4.154	0.17376\\
4.1583	0.17176\\
4.1625	0.16978\\
4.1668	0.16781\\
4.171	0.16587\\
4.1752	0.16394\\
4.1795	0.16204\\
4.1837	0.16015\\
4.188	0.15828\\
4.1922	0.15643\\
4.1965	0.1546\\
4.2007	0.15279\\
4.2049	0.15099\\
4.2092	0.14922\\
4.2134	0.14746\\
4.2177	0.14572\\
4.2219	0.14399\\
4.2262	0.14229\\
4.2304	0.1406\\
4.2346	0.13892\\
4.2389	0.13727\\
4.2431	0.13563\\
4.2474	0.13401\\
4.2516	0.1324\\
4.2559	0.13082\\
4.2601	0.12924\\
4.2643	0.12769\\
4.2686	0.12615\\
4.2728	0.12462\\
4.2771	0.12312\\
4.2813	0.12162\\
4.2856	0.12015\\
4.2898	0.11868\\
4.294	0.11724\\
4.2983	0.11581\\
4.3025	0.11439\\
4.3068	0.11299\\
4.311	0.1116\\
4.3153	0.11023\\
4.3195	0.10887\\
4.3238	0.10753\\
4.328	0.1062\\
4.3322	0.10488\\
4.3365	0.10358\\
4.3407	0.10229\\
4.345	0.10102\\
4.3492	0.09976\\
4.3535	0.098514\\
4.3577	0.09728\\
4.3619	0.09606\\
4.3662	0.094853\\
4.3704	0.093658\\
4.3747	0.092477\\
4.3789	0.091308\\
4.3832	0.090152\\
4.3874	0.089009\\
4.3916	0.087878\\
4.3959	0.086759\\
4.4001	0.085652\\
4.4044	0.084557\\
4.4086	0.083475\\
4.4129	0.082404\\
4.4171	0.081345\\
4.4213	0.080297\\
4.4256	0.079261\\
4.4298	0.078236\\
4.4341	0.077222\\
4.4383	0.07622\\
4.4426	0.075229\\
4.4468	0.074248\\
4.451	0.073279\\
4.4553	0.07232\\
4.4595	0.071372\\
4.4638	0.070434\\
4.468	0.069507\\
4.4723	0.06859\\
4.4765	0.067683\\
4.4807	0.066786\\
4.485	0.0659\\
4.4892	0.065023\\
4.4935	0.064156\\
4.4977	0.063298\\
4.502	0.062451\\
4.5062	0.061612\\
4.5104	0.060783\\
4.5147	0.059964\\
4.5189	0.059153\\
4.5232	0.058352\\
4.5274	0.057559\\
4.5317	0.056776\\
4.5359	0.056001\\
4.5402	0.055235\\
4.5444	0.054478\\
4.5486	0.053729\\
4.5529	0.052989\\
4.5571	0.052257\\
4.5614	0.051533\\
4.5656	0.050818\\
4.5699	0.05011\\
4.5741	0.049411\\
4.5783	0.04872\\
4.5826	0.048036\\
4.5868	0.04736\\
4.5911	0.046692\\
4.5953	0.046031\\
4.5996	0.045378\\
4.6038	0.044733\\
4.608	0.044094\\
4.6123	0.043464\\
4.6165	0.04284\\
4.6208	0.042223\\
4.625	0.041614\\
4.6293	0.041012\\
4.6335	0.040416\\
4.6377	0.039828\\
4.642	0.039246\\
4.6462	0.038671\\
4.6505	0.038102\\
4.6547	0.037541\\
4.659	0.036986\\
4.6632	0.036437\\
4.6674	0.035895\\
4.6717	0.035359\\
4.6759	0.03483\\
4.6802	0.034306\\
4.6844	0.033789\\
4.6887	0.033279\\
4.6929	0.032774\\
4.6971	0.032275\\
4.7014	0.031783\\
4.7056	0.031296\\
4.7099	0.030815\\
4.7141	0.03034\\
4.7184	0.029871\\
4.7226	0.029407\\
4.7269	0.028949\\
4.7311	0.028497\\
4.7353	0.02805\\
4.7396	0.027609\\
4.7438	0.027174\\
4.7481	0.026743\\
4.7523	0.026319\\
4.7566	0.025899\\
4.7608	0.025485\\
4.765	0.025076\\
4.7693	0.024673\\
4.7735	0.024274\\
4.7778	0.023881\\
4.782	0.023492\\
4.7863	0.023109\\
4.7905	0.022731\\
4.7947	0.022358\\
4.799	0.021989\\
4.8032	0.021626\\
4.8075	0.021267\\
4.8117	0.020913\\
4.816	0.020564\\
4.8202	0.02022\\
4.8244	0.01988\\
4.8287	0.019545\\
4.8329	0.019214\\
4.8372	0.018888\\
4.8414	0.018566\\
4.8457	0.018249\\
4.8499	0.017937\\
4.8541	0.017628\\
4.8584	0.017324\\
4.8626	0.017025\\
4.8669	0.016729\\
4.8711	0.016438\\
4.8754	0.016151\\
4.8796	0.015868\\
4.8838	0.015589\\
4.8881	0.015314\\
4.8923	0.015044\\
4.8966	0.014777\\
4.9008	0.014514\\
4.9051	0.014255\\
4.9093	0.014\\
4.9136	0.013749\\
4.9178	0.013501\\
4.922	0.013257\\
4.9263	0.013017\\
4.9305	0.012781\\
4.9348	0.012548\\
4.939	0.012319\\
4.9433	0.012093\\
4.9475	0.011871\\
4.9517	0.011653\\
4.956	0.011437\\
4.9602	0.011225\\
4.9645	0.011017\\
4.9687	0.010812\\
4.973	0.01061\\
4.9772	0.010411\\
4.9814	0.010216\\
4.9857	0.010023\\
4.9899	0.0098339\\
4.9942	0.0096478\\
4.9984	0.0094646\\
5.0027	0.0092845\\
5.0069	0.0091074\\
5.0111	0.0089332\\
5.0154	0.0087619\\
5.0196	0.0085934\\
5.0239	0.0084278\\
5.0281	0.0082649\\
5.0324	0.0081048\\
5.0366	0.0079474\\
5.0408	0.0077927\\
5.0451	0.0076406\\
5.0493	0.0074911\\
5.0536	0.0073442\\
5.0578	0.0071999\\
5.0621	0.007058\\
5.0663	0.0069186\\
5.0705	0.0067816\\
5.0748	0.006647\\
5.079	0.0065148\\
5.0833	0.0063849\\
5.0875	0.0062573\\
5.0918	0.0061319\\
5.096	0.0060088\\
5.1002	0.0058879\\
5.1045	0.0057692\\
5.1087	0.0056526\\
5.113	0.005538\\
5.1172	0.0054256\\
5.1215	0.0053152\\
5.1257	0.0052068\\
5.13	0.0051004\\
5.1342	0.0049959\\
5.1384	0.0048933\\
5.1427	0.0047927\\
5.1469	0.0046939\\
5.1512	0.0045969\\
5.1554	0.0045017\\
5.1597	0.0044083\\
5.1639	0.0043166\\
5.1681	0.0042267\\
5.1724	0.0041384\\
5.1766	0.0040519\\
5.1809	0.0039669\\
5.1851	0.0038836\\
5.1894	0.0038018\\
5.1936	0.0037216\\
5.1978	0.003643\\
5.2021	0.0035658\\
5.2063	0.0034901\\
5.2106	0.0034159\\
5.2148	0.0033431\\
5.2191	0.0032718\\
5.2233	0.0032018\\
5.2275	0.0031331\\
5.2318	0.0030659\\
5.236	0.0029999\\
5.2403	0.0029352\\
5.2445	0.0028718\\
5.2488	0.0028097\\
5.253	0.0027488\\
5.2572	0.0026891\\
5.2615	0.0026306\\
5.2657	0.0025732\\
5.27	0.002517\\
5.2742	0.0024619\\
5.2785	0.002408\\
5.2827	0.0023551\\
5.2869	0.0023033\\
5.2912	0.0022525\\
5.2954	0.0022028\\
5.2997	0.002154\\
5.3039	0.0021063\\
5.3082	0.0020595\\
5.3124	0.0020137\\
5.3167	0.0019689\\
5.3209	0.0019249\\
5.3251	0.0018819\\
5.3294	0.0018398\\
5.3336	0.0017985\\
5.3379	0.0017581\\
5.3421	0.0017185\\
5.3464	0.0016798\\
5.3506	0.0016419\\
5.3506	0.0013559\\
5.3464	0.0013894\\
5.3421	0.0014237\\
5.3379	0.0014588\\
5.3336	0.0014947\\
5.3294	0.0015314\\
5.3251	0.0015689\\
5.3209	0.0016073\\
5.3167	0.0016465\\
5.3124	0.0016867\\
5.3082	0.0017277\\
5.3039	0.0017696\\
5.2997	0.0018125\\
5.2954	0.0018563\\
5.2912	0.0019011\\
5.2869	0.0019469\\
5.2827	0.0019937\\
5.2785	0.0020415\\
5.2742	0.0020903\\
5.27	0.0021402\\
5.2657	0.0021912\\
5.2615	0.0022434\\
5.2572	0.0022966\\
5.253	0.002351\\
5.2488	0.0024065\\
5.2445	0.0024633\\
5.2403	0.0025212\\
5.236	0.0025804\\
5.2318	0.0026408\\
5.2275	0.0027025\\
5.2233	0.0027656\\
5.2191	0.0028299\\
5.2148	0.0028956\\
5.2106	0.0029626\\
5.2063	0.0030311\\
5.2021	0.0031009\\
5.1978	0.0031722\\
5.1936	0.003245\\
5.1894	0.0033193\\
5.1851	0.003395\\
5.1809	0.0034723\\
5.1766	0.0035512\\
5.1724	0.0036317\\
5.1681	0.0037138\\
5.1639	0.0037975\\
5.1597	0.0038829\\
5.1554	0.00397\\
5.1512	0.0040588\\
5.1469	0.0041494\\
5.1427	0.0042418\\
5.1384	0.0043359\\
5.1342	0.0044319\\
5.13	0.0045298\\
5.1257	0.0046295\\
5.1215	0.0047312\\
5.1172	0.0048348\\
5.113	0.0049404\\
5.1087	0.005048\\
5.1045	0.0051577\\
5.1002	0.0052694\\
5.096	0.0053832\\
5.0918	0.0054991\\
5.0875	0.0056172\\
5.0833	0.0057375\\
5.079	0.00586\\
5.0748	0.0059848\\
5.0705	0.0061118\\
5.0663	0.0062412\\
5.0621	0.0063729\\
5.0578	0.0065069\\
5.0536	0.0066434\\
5.0493	0.0067824\\
5.0451	0.0069238\\
5.0408	0.0070677\\
5.0366	0.0072141\\
5.0324	0.0073632\\
5.0281	0.0075148\\
5.0239	0.0076691\\
5.0196	0.0078261\\
5.0154	0.0079857\\
5.0111	0.0081481\\
5.0069	0.0083133\\
5.0027	0.0084814\\
4.9984	0.0086522\\
4.9942	0.008826\\
4.9899	0.0090027\\
4.9857	0.0091823\\
4.9814	0.0093649\\
4.9772	0.0095506\\
4.973	0.0097393\\
4.9687	0.0099311\\
4.9645	0.010126\\
4.9602	0.010324\\
4.956	0.010526\\
4.9517	0.01073\\
4.9475	0.010938\\
4.9433	0.011149\\
4.939	0.011364\\
4.9348	0.011582\\
4.9305	0.011803\\
4.9263	0.012028\\
4.922	0.012257\\
4.9178	0.012489\\
4.9136	0.012725\\
4.9093	0.012964\\
4.9051	0.013207\\
4.9008	0.013454\\
4.8966	0.013704\\
4.8923	0.013958\\
4.8881	0.014216\\
4.8838	0.014479\\
4.8796	0.014745\\
4.8754	0.015014\\
4.8711	0.015288\\
4.8669	0.015567\\
4.8626	0.015849\\
4.8584	0.016135\\
4.8541	0.016426\\
4.8499	0.01672\\
4.8457	0.017019\\
4.8414	0.017323\\
4.8372	0.01763\\
4.8329	0.017943\\
4.8287	0.018259\\
4.8244	0.01858\\
4.8202	0.018906\\
4.816	0.019236\\
4.8117	0.019571\\
4.8075	0.019911\\
4.8032	0.020255\\
4.799	0.020605\\
4.7947	0.020959\\
4.7905	0.021317\\
4.7863	0.021681\\
4.782	0.02205\\
4.7778	0.022424\\
4.7735	0.022803\\
4.7693	0.023187\\
4.765	0.023576\\
4.7608	0.023971\\
4.7566	0.02437\\
4.7523	0.024776\\
4.7481	0.025186\\
4.7438	0.025602\\
4.7396	0.026023\\
4.7353	0.02645\\
4.7311	0.026883\\
4.7269	0.027321\\
4.7226	0.027764\\
4.7184	0.028214\\
4.7141	0.028669\\
4.7099	0.02913\\
4.7056	0.029597\\
4.7014	0.03007\\
4.6971	0.030549\\
4.6929	0.031034\\
4.6887	0.031525\\
4.6844	0.032023\\
4.6802	0.032526\\
4.6759	0.033036\\
4.6717	0.033552\\
4.6674	0.034074\\
4.6632	0.034603\\
4.659	0.035138\\
4.6547	0.03568\\
4.6505	0.036228\\
4.6462	0.036783\\
4.642	0.037345\\
4.6377	0.037914\\
4.6335	0.038489\\
4.6293	0.039071\\
4.625	0.03966\\
4.6208	0.040256\\
4.6165	0.040859\\
4.6123	0.04147\\
4.608	0.042087\\
4.6038	0.042711\\
4.5996	0.043343\\
4.5953	0.043982\\
4.5911	0.044629\\
4.5868	0.045283\\
4.5826	0.045945\\
4.5783	0.046614\\
4.5741	0.04729\\
4.5699	0.047975\\
4.5656	0.048667\\
4.5614	0.049367\\
4.5571	0.050075\\
4.5529	0.050791\\
4.5486	0.051515\\
4.5444	0.052247\\
4.5402	0.052987\\
4.5359	0.053735\\
4.5317	0.054492\\
4.5274	0.055257\\
4.5232	0.056031\\
4.5189	0.056813\\
4.5147	0.057603\\
4.5104	0.058402\\
4.5062	0.05921\\
4.502	0.060027\\
4.4977	0.060853\\
4.4935	0.061688\\
4.4892	0.062531\\
4.485	0.063384\\
4.4807	0.064247\\
4.4765	0.065118\\
4.4723	0.065999\\
4.468	0.066889\\
4.4638	0.067789\\
4.4595	0.068699\\
4.4553	0.069618\\
4.451	0.070548\\
4.4468	0.071487\\
4.4426	0.072436\\
4.4383	0.073396\\
4.4341	0.074366\\
4.4298	0.075346\\
4.4256	0.076337\\
4.4213	0.077338\\
4.4171	0.07835\\
4.4129	0.079373\\
4.4086	0.080407\\
4.4044	0.081452\\
4.4001	0.082508\\
4.3959	0.083576\\
4.3916	0.084655\\
4.3874	0.085745\\
4.3832	0.086847\\
4.3789	0.087961\\
4.3747	0.089087\\
4.3704	0.090225\\
4.3662	0.091375\\
4.3619	0.092537\\
4.3577	0.093712\\
4.3535	0.0949\\
4.3492	0.0961\\
4.345	0.097313\\
4.3407	0.098539\\
4.3365	0.099779\\
4.3322	0.10103\\
4.328	0.1023\\
4.3238	0.10358\\
4.3195	0.10487\\
4.3153	0.10618\\
4.311	0.1075\\
4.3068	0.10883\\
4.3025	0.11018\\
4.2983	0.11155\\
4.294	0.11293\\
4.2898	0.11432\\
4.2856	0.11573\\
4.2813	0.11715\\
4.2771	0.11859\\
4.2728	0.12005\\
4.2686	0.12152\\
4.2643	0.123\\
4.2601	0.12451\\
4.2559	0.12602\\
4.2516	0.12756\\
4.2474	0.12911\\
4.2431	0.13068\\
4.2389	0.13226\\
4.2346	0.13386\\
4.2304	0.13548\\
4.2262	0.13711\\
4.2219	0.13877\\
4.2177	0.14044\\
4.2134	0.14212\\
4.2092	0.14383\\
4.2049	0.14555\\
4.2007	0.14729\\
4.1965	0.14905\\
4.1922	0.15083\\
4.188	0.15263\\
4.1837	0.15444\\
4.1795	0.15628\\
4.1752	0.15813\\
4.171	0.16\\
4.1668	0.1619\\
4.1625	0.16381\\
4.1583	0.16574\\
4.154	0.16769\\
4.1498	0.16966\\
4.1455	0.17166\\
4.1413	0.17367\\
4.1371	0.1757\\
4.1328	0.17775\\
4.1286	0.17983\\
4.1243	0.18192\\
4.1201	0.18404\\
4.1158	0.18618\\
4.1116	0.18834\\
4.1073	0.19052\\
4.1031	0.19272\\
4.0989	0.19494\\
4.0946	0.19719\\
4.0904	0.19946\\
4.0861	0.20174\\
4.0819	0.20406\\
4.0776	0.20639\\
4.0734	0.20875\\
4.0692	0.21113\\
4.0649	0.21353\\
4.0607	0.21595\\
4.0564	0.2184\\
4.0522	0.22087\\
4.0479	0.22337\\
4.0437	0.22589\\
4.0395	0.22843\\
4.0352	0.23099\\
4.031	0.23358\\
4.0267	0.23619\\
4.0225	0.23882\\
4.0182	0.24148\\
4.014	0.24416\\
4.0098	0.24687\\
4.0055	0.2496\\
4.0013	0.25235\\
3.997	0.25513\\
3.9928	0.25793\\
3.9885	0.26076\\
3.9843	0.26361\\
3.9801	0.26649\\
3.9758	0.26938\\
3.9716	0.27231\\
3.9673	0.27525\\
3.9631	0.27823\\
3.9588	0.28122\\
3.9546	0.28424\\
3.9504	0.28728\\
3.9461	0.29035\\
3.9419	0.29344\\
3.9376	0.29656\\
3.9334	0.2997\\
3.9291	0.30287\\
3.9249	0.30606\\
3.9207	0.30927\\
3.9164	0.3125\\
3.9122	0.31577\\
3.9079	0.31905\\
3.9037	0.32236\\
3.8994	0.32569\\
3.8952	0.32904\\
3.8909	0.33242\\
3.8867	0.33583\\
3.8825	0.33925\\
3.8782	0.3427\\
3.874	0.34617\\
3.8697	0.34967\\
3.8655	0.35318\\
3.8612	0.35672\\
3.857	0.36028\\
3.8528	0.36387\\
3.8485	0.36748\\
3.8443	0.3711\\
3.84	0.37475\\
3.8358	0.37843\\
3.8315	0.38212\\
3.8273	0.38583\\
3.8231	0.38957\\
3.8188	0.39333\\
3.8146	0.3971\\
3.8103	0.4009\\
3.8061	0.40472\\
3.8018	0.40856\\
3.7976	0.41241\\
3.7934	0.41629\\
3.7891	0.42018\\
3.7849	0.4241\\
3.7806	0.42803\\
3.7764	0.43198\\
3.7721	0.43595\\
3.7679	0.43994\\
3.7637	0.44395\\
3.7594	0.44797\\
3.7552	0.45201\\
3.7509	0.45607\\
3.7467	0.46014\\
3.7424	0.46423\\
3.7382	0.46834\\
3.734	0.47246\\
3.7297	0.4766\\
3.7255	0.48075\\
3.7212	0.48492\\
3.717	0.4891\\
3.7127	0.49329\\
3.7085	0.4975\\
3.7042	0.50173\\
3.7	0.50597\\
3.6958	0.51022\\
3.6915	0.51448\\
3.6873	0.51876\\
3.683	0.52305\\
3.6788	0.52735\\
3.6745	0.53167\\
3.6703	0.53599\\
3.6661	0.54033\\
3.6618	0.54468\\
3.6576	0.54904\\
3.6533	0.5534\\
3.6491	0.55778\\
3.6448	0.56217\\
3.6406	0.56656\\
3.6364	0.57096\\
3.6321	0.57537\\
3.6279	0.57978\\
3.6236	0.5842\\
3.6194	0.58863\\
3.6151	0.59306\\
3.6109	0.59749\\
3.6067	0.60192\\
3.6024	0.60635\\
3.5982	0.61078\\
3.5939	0.61521\\
3.5897	0.61964\\
3.5854	0.62406\\
3.5812	0.62848\\
3.577	0.63289\\
3.5727	0.6373\\
3.5685	0.6417\\
3.5642	0.64609\\
3.56	0.65047\\
3.5557	0.65484\\
3.5515	0.6592\\
3.5473	0.66354\\
3.543	0.66788\\
3.5388	0.6722\\
3.5345	0.6765\\
3.5303	0.6808\\
3.526	0.68507\\
3.5218	0.68933\\
3.5176	0.69358\\
3.5133	0.69781\\
3.5091	0.70202\\
3.5048	0.70621\\
3.5006	0.71038\\
3.4963	0.71453\\
3.4921	0.71865\\
3.4878	0.72276\\
3.4836	0.72683\\
3.4794	0.73088\\
3.4751	0.7349\\
3.4709	0.73889\\
3.4666	0.74284\\
3.4624	0.74676\\
3.4581	0.75063\\
3.4539	0.75447\\
3.4497	0.75826\\
3.4454	0.76199\\
3.4412	0.76568\\
3.4369	0.76932\\
3.4327	0.77289\\
3.4284	0.77641\\
3.4242	0.77987\\
3.42	0.78326\\
3.4157	0.78659\\
3.4115	0.78985\\
3.4072	0.79304\\
3.403	0.79616\\
3.3987	0.79922\\
3.3945	0.8022\\
3.3903	0.80512\\
3.386	0.80796\\
3.3818	0.81074\\
3.3775	0.81345\\
3.3733	0.81609\\
3.369	0.81867\\
3.3648	0.82118\\
3.3606	0.82363\\
3.3563	0.82601\\
3.3521	0.82832\\
3.3478	0.83058\\
3.3436	0.83277\\
3.3393	0.8349\\
3.3351	0.83696\\
3.3309	0.83897\\
3.3266	0.8409\\
3.3224	0.84277\\
3.3181	0.84458\\
3.3139	0.84631\\
3.3096	0.84797\\
3.3054	0.84955\\
3.3011	0.85105\\
3.2969	0.85246\\
3.2927	0.85379\\
3.2884	0.85501\\
3.2842	0.85614\\
3.2799	0.85716\\
3.2757	0.85807\\
3.2714	0.85887\\
3.2672	0.85955\\
3.263	0.86011\\
3.2587	0.86056\\
3.2545	0.86089\\
3.2502	0.86111\\
3.246	0.86121\\
3.2417	0.8612\\
3.2375	0.86109\\
3.2333	0.86088\\
3.229	0.86057\\
3.2248	0.86016\\
3.2205	0.85966\\
3.2163	0.85907\\
3.212	0.85839\\
3.2078	0.85762\\
3.2036	0.85676\\
3.1993	0.85581\\
3.1951	0.85477\\
3.1908	0.85364\\
3.1866	0.85241\\
3.1823	0.85108\\
3.1781	0.84965\\
3.1739	0.84812\\
3.1696	0.84649\\
3.1654	0.84475\\
3.1611	0.8429\\
3.1569	0.84095\\
3.1526	0.83889\\
3.1484	0.83673\\
3.1442	0.83447\\
3.1399	0.83212\\
3.1357	0.82967\\
3.1314	0.82714\\
3.1272	0.82452\\
3.1229	0.82183\\
3.1187	0.81906\\
3.1145	0.81622\\
3.1102	0.81332\\
3.106	0.81035\\
3.1017	0.80732\\
3.0975	0.80424\\
3.0932	0.8011\\
3.089	0.7979\\
3.0847	0.79465\\
3.0805	0.79134\\
3.0763	0.78798\\
3.072	0.78457\\
3.0678	0.7811\\
3.0635	0.77758\\
3.0593	0.774\\
3.055	0.77038\\
3.0508	0.7667\\
3.0466	0.76296\\
3.0423	0.75917\\
3.0381	0.75533\\
3.0338	0.75143\\
3.0296	0.74748\\
3.0253	0.74348\\
3.0211	0.73943\\
3.0169	0.73532\\
3.0126	0.73117\\
3.0084	0.72696\\
3.0041	0.7227\\
2.9999	0.7184\\
2.9956	0.71404\\
2.9914	0.70964\\
2.9872	0.70519\\
2.9829	0.7007\\
2.9787	0.69617\\
2.9744	0.69159\\
2.9702	0.68697\\
2.9659	0.68231\\
2.9617	0.67762\\
2.9575	0.67288\\
2.9532	0.66811\\
2.949	0.66331\\
2.9447	0.65847\\
2.9405	0.6536\\
2.9362	0.6487\\
2.932	0.64377\\
2.9278	0.63882\\
2.9235	0.63384\\
2.9193	0.62883\\
2.915	0.6238\\
2.9108	0.61876\\
2.9065	0.61369\\
2.9023	0.60861\\
2.898	0.60351\\
2.8938	0.59839\\
2.8896	0.59327\\
2.8853	0.58813\\
2.8811	0.58298\\
2.8768	0.57782\\
2.8726	0.57266\\
2.8683	0.56749\\
2.8641	0.56231\\
2.8599	0.55713\\
2.8556	0.55195\\
2.8514	0.54677\\
2.8471	0.5416\\
2.8429	0.53642\\
2.8386	0.53125\\
2.8344	0.52608\\
2.8302	0.52091\\
2.8259	0.51575\\
2.8217	0.5106\\
2.8174	0.50546\\
2.8132	0.50033\\
2.8089	0.49521\\
2.8047	0.4901\\
2.8005	0.485\\
2.7962	0.47992\\
2.792	0.47485\\
2.7877	0.46979\\
2.7835	0.46475\\
2.7792	0.45973\\
2.775	0.45472\\
2.7708	0.44974\\
2.7665	0.44477\\
2.7623	0.43982\\
2.758	0.43489\\
2.7538	0.42998\\
2.7495	0.42509\\
2.7453	0.42023\\
2.7411	0.41538\\
2.7368	0.41056\\
2.7326	0.40577\\
2.7283	0.401\\
2.7241	0.39625\\
2.7198	0.39153\\
2.7156	0.38684\\
2.7113	0.38217\\
2.7071	0.37753\\
2.7029	0.37292\\
2.6986	0.36833\\
2.6944	0.36377\\
2.6901	0.35925\\
2.6859	0.35475\\
2.6816	0.35028\\
2.6774	0.34584\\
2.6732	0.34143\\
2.6689	0.33705\\
2.6647	0.3327\\
2.6604	0.32839\\
2.6562	0.3241\\
2.6519	0.31985\\
2.6477	0.31563\\
2.6435	0.31144\\
2.6392	0.30728\\
2.635	0.30316\\
2.6307	0.29907\\
2.6265	0.29501\\
2.6222	0.29099\\
2.618	0.287\\
2.6138	0.28305\\
2.6095	0.27912\\
2.6053	0.27524\\
2.601	0.27138\\
2.5968	0.26756\\
2.5925	0.26378\\
2.5883	0.26003\\
2.5841	0.25631\\
2.5798	0.25263\\
2.5756	0.24899\\
2.5713	0.24537\\
2.5671	0.2418\\
2.5628	0.23826\\
2.5586	0.23475\\
2.5544	0.23128\\
2.5501	0.22784\\
2.5459	0.22444\\
2.5416	0.22107\\
2.5374	0.21774\\
2.5331	0.21444\\
2.5289	0.21118\\
2.5247	0.20795\\
2.5204	0.20476\\
2.5162	0.2016\\
2.5119	0.19848\\
2.5077	0.19539\\
2.5034	0.19233\\
2.4992	0.18931\\
2.4949	0.18633\\
2.4907	0.18337\\
2.4865	0.18046\\
2.4822	0.17757\\
2.478	0.17472\\
2.4737	0.17191\\
2.4695	0.16913\\
2.4652	0.16638\\
2.461	0.16366\\
2.4568	0.16098\\
2.4525	0.15833\\
2.4483	0.15571\\
2.444	0.15313\\
2.4398	0.15058\\
2.4355	0.14806\\
2.4313	0.14557\\
2.4271	0.14312\\
2.4228	0.14069\\
2.4186	0.1383\\
2.4143	0.13594\\
2.4101	0.13362\\
2.4058	0.13132\\
2.4016	0.12905\\
2.3974	0.12682\\
2.3931	0.12461\\
2.3889	0.12244\\
2.3846	0.12029\\
2.3804	0.11818\\
2.3761	0.11609\\
2.3719	0.11404\\
2.3677	0.11201\\
2.3634	0.11001\\
2.3592	0.10805\\
2.3549	0.10611\\
2.3507	0.10419\\
2.3464	0.10231\\
2.3422	0.10045\\
2.338	0.098625\\
2.3337	0.096823\\
2.3295	0.095049\\
2.3252	0.093301\\
2.321	0.09158\\
2.3167	0.089885\\
2.3125	0.088216\\
2.3082	0.086573\\
2.304	0.084955\\
2.2998	0.083363\\
2.2955	0.081795\\
2.2913	0.080253\\
2.287	0.078735\\
2.2828	0.077241\\
2.2785	0.07577\\
2.2743	0.074324\\
2.2701	0.072901\\
2.2658	0.071501\\
2.2616	0.070124\\
2.2573	0.06877\\
2.2531	0.067438\\
2.2488	0.066128\\
2.2446	0.064841\\
2.2404	0.063574\\
2.2361	0.062329\\
2.2319	0.061105\\
2.2276	0.059902\\
2.2234	0.058719\\
2.2191	0.057557\\
2.2149	0.056414\\
2.2107	0.055292\\
2.2064	0.054189\\
2.2022	0.053105\\
2.1979	0.05204\\
2.1937	0.050994\\
2.1894	0.049966\\
2.1852	0.048957\\
2.181	0.047965\\
2.1767	0.046992\\
2.1725	0.046035\\
2.1682	0.045096\\
2.164	0.044174\\
2.1597	0.043269\\
2.1555	0.04238\\
2.1513	0.041508\\
2.147	0.040651\\
2.1428	0.039811\\
2.1385	0.038986\\
2.1343	0.038176\\
2.13	0.037381\\
2.1258	0.036601\\
2.1216	0.035836\\
2.1173	0.035086\\
2.1131	0.034349\\
2.1088	0.033626\\
2.1046	0.032918\\
2.1003	0.032222\\
2.0961	0.03154\\
2.0918	0.030871\\
2.0876	0.030215\\
2.0834	0.029572\\
2.0791	0.028941\\
2.0749	0.028323\\
2.0706	0.027716\\
2.0664	0.027121\\
2.0621	0.026539\\
2.0579	0.025967\\
2.0537	0.025407\\
2.0494	0.024858\\
2.0452	0.02432\\
2.0409	0.023792\\
2.0367	0.023275\\
2.0324	0.022769\\
2.0282	0.022272\\
2.024	0.021786\\
2.0197	0.021309\\
2.0155	0.020842\\
2.0112	0.020385\\
2.007	0.019937\\
2.0027	0.019498\\
1.9985	0.019068\\
1.9943	0.018647\\
1.99	0.018234\\
1.9858	0.01783\\
1.9815	0.017434\\
1.9773	0.017047\\
1.973	0.016667\\
1.9688	0.016296\\
1.9646	0.015932\\
1.9603	0.015576\\
1.9561	0.015227\\
1.9518	0.014885\\
1.9476	0.014551\\
1.9433	0.014224\\
1.9391	0.013904\\
1.9349	0.01359\\
1.9306	0.013283\\
1.9264	0.012983\\
1.9221	0.012689\\
1.9179	0.012401\\
1.9136	0.01212\\
1.9094	0.011844\\
1.9051	0.011575\\
1.9009	0.011311\\
1.8967	0.011053\\
1.8924	0.0108\\
1.8882	0.010553\\
1.8839	0.010311\\
1.8797	0.010075\\
1.8754	0.0098436\\
1.8712	0.0096173\\
1.867	0.0093959\\
1.8627	0.0091794\\
1.8585	0.0089676\\
1.8542	0.0087605\\
1.85	0.0085579\\
1.8457	0.0083598\\
1.8415	0.008166\\
1.8373	0.0079766\\
1.833	0.0077913\\
1.8288	0.0076101\\
1.8245	0.0074329\\
1.8203	0.0072597\\
1.816	0.0070903\\
1.8118	0.0069247\\
1.8076	0.0067628\\
1.8033	0.0066045\\
1.7991	0.0064497\\
1.7948	0.0062984\\
1.7906	0.0061505\\
1.7863	0.006006\\
1.7821	0.0058646\\
1.7779	0.0057265\\
1.7736	0.0055915\\
1.7694	0.0054595\\
1.7651	0.0053305\\
1.7609	0.0052045\\
1.7566	0.0050813\\
1.7524	0.0049609\\
1.7482	0.0048432\\
1.7439	0.0047282\\
1.7397	0.0046158\\
1.7354	0.004506\\
1.7312	0.0043987\\
1.7269	0.0042939\\
1.7227	0.0041915\\
1.7185	0.0040914\\
1.7142	0.0039936\\
1.71	0.003898\\
1.7057	0.0038047\\
1.7015	0.0037135\\
1.6972	0.0036244\\
1.693	0.0035373\\
1.6887	0.0034523\\
1.6845	0.0033692\\
1.6803	0.0032881\\
1.676	0.0032088\\
1.6718	0.0031314\\
1.6675	0.0030558\\
1.6633	0.0029819\\
1.659	0.0029098\\
1.6548	0.0028393\\
1.6506	0.0027705\\
1.6463	0.0027033\\
1.6421	0.0026377\\
1.6378	0.0025736\\
1.6336	0.0025109\\
1.6293	0.0024498\\
1.6251	0.0023901\\
1.6209	0.0023318\\
1.6166	0.0022749\\
1.6124	0.0022193\\
1.6081	0.002165\\
1.6039	0.002112\\
1.5996	0.0020602\\
1.5954	0.0020097\\
1.5912	0.0019604\\
1.5869	0.0019122\\
1.5827	0.0018652\\
1.5784	0.0018193\\
1.5742	0.0017744\\
1.5699	0.0017307\\
1.5657	0.001688\\
1.5615	0.0016463\\
1.5572	0.0016055\\
1.553	0.0015658\\
1.5487	0.001527\\
1.5445	0.0014891\\
1.5402	0.0014522\\
1.536	0.0014161\\
1.5318	0.0013809\\
1.5275	0.0013465\\
1.5233	0.001313\\
1.519	0.0012802\\
1.5148	0.0012483\\
1.5105	0.0012171\\
1.5063	0.0011867\\
1.502	0.001157\\
1.4978	0.001128\\
1.4936	0.0010997\\
1.4893	0.0010721\\
1.4851	0.0010451\\
1.4808	0.0010188\\
1.4766	0.0009932\\
1.4723	0.00096817\\
1.4681	0.00094375\\
1.4639	0.00091992\\
1.4596	0.00089667\\
1.4554	0.00087399\\
1.4511	0.00085185\\
1.4469	0.00083026\\
1.4426	0.00080919\\
1.4384	0.00078864\\
1.4342	0.00076859\\
1.4299	0.00074903\\
1.4257	0.00072995\\
1.4214	0.00071134\\
1.4172	0.00069318\\
1.4129	0.00067546\\
1.4087	0.00065819\\
1.4045	0.00064133\\
1.4002	0.0006249\\
1.396	0.00060886\\
1.3917	0.00059322\\
1.3875	0.00057797\\
1.3832	0.0005631\\
1.379	0.00054859\\
1.3748	0.00053444\\
1.3705	0.00052064\\
1.3663	0.00050719\\
1.362	0.00049406\\
1.3578	0.00048127\\
1.3535	0.00046879\\
1.3493	0.00045662\\
1.3451	0.00044475\\
1.3408	0.00043319\\
1.3366	0.0004219\\
1.3323	0.00041091\\
1.3281	0.00040018\\
1.3238	0.00038973\\
1.3196	0.00037953\\
1.3154	0.0003696\\
1.3111	0.00035991\\
1.3069	0.00035046\\
1.3026	0.00034126\\
1.2984	0.00033228\\
1.2941	0.00032353\\
1.2899	0.00031501\\
1.2856	0.00030669\\
1.2814	0.00029859\\
1.2772	0.00029069\\
1.2729	0.000283\\
1.2687	0.00027549\\
1.2644	0.00026818\\
1.2602	0.00026106\\
1.2559	0.00025411\\
1.2517	0.00024734\\
1.2475	0.00024075\\
1.2432	0.00023432\\
1.239	0.00022806\\
1.2347	0.00022196\\
1.2305	0.00021601\\
1.2262	0.00021022\\
1.222	0.00020458\\
1.2178	0.00019908\\
1.2135	0.00019372\\
1.2093	0.0001885\\
1.205	0.00018341\\
1.2008	0.00017846\\
1.1965	0.00017363\\
1.1923	0.00016893\\
1.1881	0.00016435\\
1.1838	0.00015989\\
1.1796	0.00015555\\
1.1753	0.00015132\\
1.1711	0.00014719\\
1.1668	0.00014318\\
1.1626	0.00013927\\
1.1584	0.00013546\\
1.1541	0.00013176\\
1.1499	0.00012815\\
1.1456	0.00012463\\
1.1414	0.00012121\\
1.1371	0.00011787\\
1.1329	0.00011462\\
1.1287	0.00011146\\
1.1244	0.00010839\\
1.1202	0.00010539\\
1.1159	0.00010247\\
1.1117	9.9634e-05\\
}--cycle;
\addplot[ybar interval, fill=black, fill opacity=0.35, area legend, draw=none] table[row sep=crcr, x=Lower, y=Count] {%
Lower	Upper	Count\\
1.1117	1.4377	0\\
1.4377	1.7638	0\\
1.7638	2.0899	0.0081782\\
2.0899	2.416	0.081782\\
2.416	2.742	0.23717\\
2.742	3.0681	0.67061\\
3.0681	3.3942	0.86689\\
3.3942	3.7202	0.7115\\
3.7202	4.0463	0.31077\\
4.0463	4.3724	0.10632\\
4.3724	4.6985	0.049069\\
4.6985	5.0245	0.0081782\\
5.0245	5.3506	0.016356\\
5.3506	5.3506	0.016356\\
};
\addlegendentry{\RealHistogramLegend}

\addplot [color=black]
  table[row sep=crcr]{%
1.1117	1.3113e-06\\
1.6209	0.00038789\\
1.7439	0.0012168\\
1.8245	0.00245\\
1.8882	0.0041414\\
1.9391	0.0061934\\
1.9858	0.0088349\\
2.024	0.0117\\
2.0621	0.01536\\
2.0961	0.01942\\
2.1258	0.02371\\
2.1555	0.028793\\
2.181	0.033866\\
2.2064	0.039678\\
2.2319	0.046307\\
2.2573	0.053834\\
2.2785	0.060851\\
2.2998	0.068597\\
2.321	0.07712\\
2.3422	0.086467\\
2.3634	0.096686\\
2.3846	0.10782\\
2.4058	0.11991\\
2.4271	0.133\\
2.4483	0.14712\\
2.4695	0.1623\\
2.4907	0.17856\\
2.5119	0.19592\\
2.5331	0.21439\\
2.5544	0.23396\\
2.5756	0.25463\\
2.5968	0.27638\\
2.618	0.29918\\
2.6392	0.32299\\
2.6647	0.35282\\
2.6901	0.3839\\
2.7156	0.4161\\
2.7453	0.45485\\
2.7835	0.50609\\
2.9065	0.67268\\
2.932	0.70545\\
2.9575	0.73697\\
2.9787	0.76206\\
2.9999	0.78589\\
3.0169	0.80393\\
3.0338	0.82099\\
3.0508	0.83697\\
3.0678	0.85181\\
3.0847	0.86542\\
3.0975	0.8748\\
3.1102	0.88343\\
3.1229	0.89129\\
3.1357	0.89836\\
3.1484	0.90462\\
3.1611	0.91005\\
3.1739	0.91464\\
3.1866	0.91838\\
3.1993	0.92126\\
3.212	0.92327\\
3.2248	0.9244\\
3.2375	0.92466\\
3.2502	0.92405\\
3.263	0.92256\\
3.2757	0.9202\\
3.2884	0.91699\\
3.3011	0.91293\\
3.3139	0.90804\\
3.3266	0.90233\\
3.3393	0.89581\\
3.3521	0.88852\\
3.3648	0.88047\\
3.3775	0.87168\\
3.3945	0.85887\\
3.4115	0.84487\\
3.4284	0.82975\\
3.4454	0.81358\\
3.4624	0.79644\\
3.4836	0.77379\\
3.5048	0.74993\\
3.5303	0.71994\\
3.5557	0.68874\\
3.5897	0.6458\\
3.6406	0.57989\\
3.7	0.50329\\
3.734	0.46076\\
3.7637	0.42477\\
3.7891	0.39503\\
3.8146	0.36646\\
3.84	0.33917\\
3.8612	0.31747\\
3.8825	0.29676\\
3.9037	0.27707\\
3.9249	0.2584\\
3.9461	0.24077\\
3.9673	0.22417\\
3.9885	0.20859\\
4.0098	0.19401\\
4.031	0.18041\\
4.0522	0.16776\\
4.0734	0.15602\\
4.0946	0.14515\\
4.1158	0.1351\\
4.1371	0.12585\\
4.1583	0.11733\\
4.1837	0.10802\\
4.2092	0.099623\\
4.2346	0.092063\\
4.2601	0.085262\\
4.2898	0.078184\\
4.3195	0.071922\\
4.3535	0.065631\\
4.3874	0.060127\\
4.4256	0.054718\\
4.468	0.049499\\
4.5147	0.044517\\
4.5699	0.039409\\
4.6335	0.034287\\
4.7099	0.028925\\
4.799	0.023471\\
4.8923	0.018518\\
4.9899	0.014105\\
5.0875	0.01045\\
5.1894	0.007407\\
5.2954	0.0049981\\
5.3506	0.004016\\
};
\addlegendentry{\RealBLNFitLegend}

      % \input{../../../../.tex/20/KS/SizeDistributionFittingsfig29.tex}
      \legend{}
  
      \nextgroupplot[%
        xmin=0.85621,xmax=5.4115,
        % xmin=0.96242,xmax=5.5701,
        % ymin=0,ymax=1,
      ] % This file was created by matlab2tikz.
%
\definecolor{mycolor1}{rgb}{0.83529,0.36863,0.00000}%
%
\addplot[ybar interval, fill=mycolor1, fill opacity=0.35, area legend, draw=none] table[row sep=crcr, x=Lower, y=Count] {%
Lower	Upper	Count\\
1.1719	1.4941	0.0004138\\
1.4941	1.8163	0.0033104\\
1.8163	2.1385	0.023173\\
2.1385	2.4607	0.095753\\
2.4607	2.783	0.31093\\
2.783	3.1052	0.6509\\
3.1052	3.4274	0.85457\\
3.4274	3.7496	0.62972\\
3.7496	4.0718	0.31904\\
4.0718	4.394	0.14177\\
4.394	4.7163	0.052304\\
4.7163	5.0385	0.013821\\
5.0385	5.3607	0.0077794\\
5.3607	5.3607	0.0077794\\
};
\addlegendentry{\ApproxHistogramLegend}

\addplot [color=mycolor1, only marks, every error bar/.append style={opacity=0.45}, mark=*, mark size=0pt, draw=none, forget plot]
 plot [error bars/.cd, y dir=both, y explicit, error bar style={line width=1pt, color=mycolor1}, error mark options={mark=none,mark size=0pt}]
 table[row sep=crcr, y error plus index=2, y error minus index=3]{%
1.333	0.0004138	0.0003597	0.0003597\\
1.6552	0.0033104	0.0011901	0.0011901\\
1.9774	0.023173	0.0027616	0.0027616\\
2.2996	0.095753	0.0051043	0.0051043\\
2.6218	0.31093	0.0094066	0.0094066\\
2.9441	0.6509	0.011955	0.011955\\
3.2663	0.85457	0.012736	0.012736\\
3.5885	0.62972	0.011456	0.011456\\
3.9107	0.31904	0.0083059	0.0083059\\
4.2329	0.14177	0.0051489	0.0051489\\
4.5552	0.052304	0.0029604	0.0029604\\
4.8774	0.013821	0.0015667	0.0015667\\
5.1996	0.0077794	0.00045618	0.00045618\\
};
\addplot [color=mycolor1]
  table[row sep=crcr]{%
1.1719	0.00026636\\
1.4025	0.001005\\
1.5408	0.0021557\\
1.6415	0.0036996\\
1.7211	0.0056241\\
1.7882	0.0079545\\
1.8469	0.010723\\
1.8973	0.013801\\
1.9434	0.017338\\
1.9853	0.021277\\
2.023	0.02552\\
2.0608	0.030537\\
2.0943	0.035742\\
2.1279	0.041746\\
2.1572	0.047735\\
2.1866	0.054485\\
2.2159	0.062072\\
2.2411	0.069303\\
2.2662	0.077259\\
2.2914	0.085995\\
2.3166	0.095563\\
2.3417	0.10602\\
2.3669	0.11741\\
2.392	0.12979\\
2.413	0.14089\\
2.434	0.15275\\
2.4549	0.16537\\
2.4759	0.17878\\
2.4969	0.193\\
2.5178	0.20805\\
2.5388	0.22393\\
2.5598	0.24066\\
2.5807	0.25822\\
2.6059	0.28041\\
2.631	0.30379\\
2.6562	0.32833\\
2.6814	0.35399\\
2.7065	0.38071\\
2.7359	0.41312\\
2.7652	0.44674\\
2.7988	0.48643\\
2.8365	0.53237\\
2.8868	0.59505\\
2.9539	0.67868\\
2.9791	0.70887\\
3	0.73297\\
3.0168	0.75131\\
3.0336	0.76864\\
3.0503	0.78482\\
3.0671	0.79976\\
3.0797	0.81008\\
3.0923	0.81964\\
3.1049	0.82841\\
3.1174	0.8364\\
3.13	0.84359\\
3.1426	0.84999\\
3.1552	0.85559\\
3.1677	0.8604\\
3.1803	0.8644\\
3.1929	0.8676\\
3.2055	0.86999\\
3.2181	0.87155\\
3.2306	0.87231\\
3.2432	0.87225\\
3.2558	0.8714\\
3.2684	0.86977\\
3.281	0.8674\\
3.2935	0.86431\\
3.3061	0.86054\\
3.3187	0.85612\\
3.3313	0.85108\\
3.3439	0.84541\\
3.3564	0.83912\\
3.369	0.83216\\
3.3816	0.82452\\
3.3942	0.81615\\
3.4067	0.80703\\
3.4193	0.79716\\
3.4361	0.78289\\
3.4529	0.76751\\
3.4738	0.74707\\
3.5032	0.71708\\
3.5703	0.64685\\
3.6625	0.55057\\
3.7086	0.50357\\
3.7464	0.46623\\
3.7799	0.43412\\
3.8093	0.40696\\
3.8386	0.38078\\
3.868	0.35564\\
3.8973	0.33158\\
3.9267	0.30864\\
3.9518	0.28989\\
3.977	0.27198\\
4.0022	0.2549\\
4.0273	0.23867\\
4.0525	0.22325\\
4.0776	0.20865\\
4.1028	0.19483\\
4.128	0.18178\\
4.1573	0.16749\\
4.1867	0.15416\\
4.216	0.14175\\
4.2454	0.13021\\
4.2747	0.11949\\
4.3041	0.10956\\
4.3334	0.10035\\
4.367	0.090658\\
4.4005	0.081797\\
4.434	0.073702\\
4.4718	0.065438\\
4.5095	0.057991\\
4.5473	0.051289\\
4.5892	0.044636\\
4.6311	0.038738\\
4.6772	0.03303\\
4.7276	0.02763\\
4.7779	0.022988\\
4.8324	0.018705\\
4.8911	0.014844\\
4.954	0.01145\\
5.0253	0.0083959\\
5.1007	0.0059301\\
5.1846	0.0039419\\
5.2852	0.0023514\\
5.3607	0.001571\\
};
\addlegendentry{\ApproxBLNFitLegend}


\addplot[area legend, draw=none, fill=mycolor1, fill opacity=0.15, forget plot]
table[row sep=crcr] {%
x	y\\
1.1719	0.00033352\\
1.1761	0.00034157\\
1.1802	0.0003498\\
1.1844	0.00035823\\
1.1886	0.00036685\\
1.1928	0.00037568\\
1.197	0.0003847\\
1.2012	0.00039394\\
1.2054	0.00040338\\
1.2096	0.00041305\\
1.2138	0.00042293\\
1.218	0.00043305\\
1.2222	0.0004434\\
1.2264	0.00045398\\
1.2306	0.00046481\\
1.2348	0.00047589\\
1.2389	0.00048722\\
1.2431	0.0004988\\
1.2473	0.00051066\\
1.2515	0.00052278\\
1.2557	0.00053518\\
1.2599	0.00054786\\
1.2641	0.00056084\\
1.2683	0.0005741\\
1.2725	0.00058767\\
1.2767	0.00060155\\
1.2809	0.00061574\\
1.2851	0.00063025\\
1.2893	0.00064509\\
1.2935	0.00066026\\
1.2977	0.00067578\\
1.3018	0.00069165\\
1.306	0.00070787\\
1.3102	0.00072446\\
1.3144	0.00074142\\
1.3186	0.00075877\\
1.3228	0.0007765\\
1.327	0.00079463\\
1.3312	0.00081317\\
1.3354	0.00083212\\
1.3396	0.0008515\\
1.3438	0.0008713\\
1.348	0.00089155\\
1.3522	0.00091225\\
1.3564	0.00093341\\
1.3605	0.00095505\\
1.3647	0.00097716\\
1.3689	0.00099976\\
1.3731	0.0010229\\
1.3773	0.0010465\\
1.3815	0.0010706\\
1.3857	0.0010953\\
1.3899	0.0011205\\
1.3941	0.0011463\\
1.3983	0.0011726\\
1.4025	0.0011995\\
1.4067	0.0012271\\
1.4109	0.0012552\\
1.4151	0.0012839\\
1.4192	0.0013133\\
1.4234	0.0013433\\
1.4276	0.0013739\\
1.4318	0.0014053\\
1.436	0.0014373\\
1.4402	0.00147\\
1.4444	0.0015034\\
1.4486	0.0015376\\
1.4528	0.0015725\\
1.457	0.0016082\\
1.4612	0.0016446\\
1.4654	0.0016818\\
1.4696	0.0017198\\
1.4738	0.0017587\\
1.478	0.0017984\\
1.4821	0.0018389\\
1.4863	0.0018803\\
1.4905	0.0019227\\
1.4947	0.0019659\\
1.4989	0.00201\\
1.5031	0.0020551\\
1.5073	0.0021012\\
1.5115	0.0021483\\
1.5157	0.0021963\\
1.5199	0.0022454\\
1.5241	0.0022956\\
1.5283	0.0023468\\
1.5325	0.0023991\\
1.5367	0.0024525\\
1.5408	0.002507\\
1.545	0.0025628\\
1.5492	0.0026197\\
1.5534	0.0026778\\
1.5576	0.0027371\\
1.5618	0.0027977\\
1.566	0.0028596\\
1.5702	0.0029227\\
1.5744	0.0029873\\
1.5786	0.0030531\\
1.5828	0.0031204\\
1.587	0.0031891\\
1.5912	0.0032592\\
1.5954	0.0033308\\
1.5995	0.0034039\\
1.6037	0.0034785\\
1.6079	0.0035547\\
1.6121	0.0036324\\
1.6163	0.0037118\\
1.6205	0.0037929\\
1.6247	0.0038756\\
1.6289	0.0039601\\
1.6331	0.0040463\\
1.6373	0.0041343\\
1.6415	0.0042242\\
1.6457	0.0043159\\
1.6499	0.0044095\\
1.6541	0.004505\\
1.6583	0.0046025\\
1.6624	0.004702\\
1.6666	0.0048036\\
1.6708	0.0049073\\
1.675	0.005013\\
1.6792	0.005121\\
1.6834	0.0052312\\
1.6876	0.0053436\\
1.6918	0.0054584\\
1.696	0.0055754\\
1.7002	0.0056949\\
1.7044	0.0058168\\
1.7086	0.0059412\\
1.7128	0.0060681\\
1.717	0.0061976\\
1.7211	0.0063298\\
1.7253	0.0064646\\
1.7295	0.0066021\\
1.7337	0.0067425\\
1.7379	0.0068856\\
1.7421	0.0070317\\
1.7463	0.0071807\\
1.7505	0.0073327\\
1.7547	0.0074877\\
1.7589	0.0076459\\
1.7631	0.0078073\\
1.7673	0.0079718\\
1.7715	0.0081397\\
1.7757	0.0083109\\
1.7799	0.0084856\\
1.784	0.0086637\\
1.7882	0.0088454\\
1.7924	0.0090307\\
1.7966	0.0092196\\
1.8008	0.0094123\\
1.805	0.0096089\\
1.8092	0.0098093\\
1.8134	0.010014\\
1.8176	0.010222\\
1.8218	0.010435\\
1.826	0.010651\\
1.8302	0.010872\\
1.8344	0.011098\\
1.8386	0.011327\\
1.8427	0.011562\\
1.8469	0.0118\\
1.8511	0.012044\\
1.8553	0.012292\\
1.8595	0.012545\\
1.8637	0.012803\\
1.8679	0.013066\\
1.8721	0.013334\\
1.8763	0.013607\\
1.8805	0.013886\\
1.8847	0.01417\\
1.8889	0.014459\\
1.8931	0.014754\\
1.8973	0.015054\\
1.9014	0.015361\\
1.9056	0.015673\\
1.9098	0.015991\\
1.914	0.016315\\
1.9182	0.016646\\
1.9224	0.016983\\
1.9266	0.017326\\
1.9308	0.017675\\
1.935	0.018032\\
1.9392	0.018395\\
1.9434	0.018764\\
1.9476	0.019141\\
1.9518	0.019525\\
1.956	0.019916\\
1.9602	0.020315\\
1.9643	0.02072\\
1.9685	0.021134\\
1.9727	0.021555\\
1.9769	0.021984\\
1.9811	0.022421\\
1.9853	0.022866\\
1.9895	0.02332\\
1.9937	0.023781\\
1.9979	0.024251\\
2.0021	0.02473\\
2.0063	0.025218\\
2.0105	0.025715\\
2.0147	0.02622\\
2.0189	0.026735\\
2.023	0.02726\\
2.0272	0.027793\\
2.0314	0.028337\\
2.0356	0.02889\\
2.0398	0.029454\\
2.044	0.030027\\
2.0482	0.030611\\
2.0524	0.031206\\
2.0566	0.031811\\
2.0608	0.032427\\
2.065	0.033054\\
2.0692	0.033692\\
2.0734	0.034341\\
2.0776	0.035002\\
2.0817	0.035675\\
2.0859	0.036359\\
2.0901	0.037056\\
2.0943	0.037764\\
2.0985	0.038486\\
2.1027	0.039219\\
2.1069	0.039966\\
2.1111	0.040726\\
2.1153	0.041499\\
2.1195	0.042285\\
2.1237	0.043085\\
2.1279	0.043898\\
2.1321	0.044726\\
2.1363	0.045568\\
2.1405	0.046424\\
2.1446	0.047295\\
2.1488	0.04818\\
2.153	0.049081\\
2.1572	0.049997\\
2.1614	0.050928\\
2.1656	0.051875\\
2.1698	0.052838\\
2.174	0.053817\\
2.1782	0.054813\\
2.1824	0.055825\\
2.1866	0.056853\\
2.1908	0.057899\\
2.195	0.058962\\
2.1992	0.060043\\
2.2033	0.061141\\
2.2075	0.062258\\
2.2117	0.063392\\
2.2159	0.064545\\
2.2201	0.065717\\
2.2243	0.066907\\
2.2285	0.068117\\
2.2327	0.069346\\
2.2369	0.070595\\
2.2411	0.071864\\
2.2453	0.073152\\
2.2495	0.074462\\
2.2537	0.075792\\
2.2579	0.077143\\
2.262	0.078515\\
2.2662	0.079908\\
2.2704	0.081324\\
2.2746	0.082761\\
2.2788	0.08422\\
2.283	0.085702\\
2.2872	0.087207\\
2.2914	0.088735\\
2.2956	0.090286\\
2.2998	0.09186\\
2.304	0.093459\\
2.3082	0.095081\\
2.3124	0.096728\\
2.3166	0.098399\\
2.3208	0.1001\\
2.3249	0.10182\\
2.3291	0.10356\\
2.3333	0.10534\\
2.3375	0.10713\\
2.3417	0.10896\\
2.3459	0.11081\\
2.3501	0.11269\\
2.3543	0.11459\\
2.3585	0.11652\\
2.3627	0.11848\\
2.3669	0.12047\\
2.3711	0.12249\\
2.3753	0.12453\\
2.3795	0.1266\\
2.3836	0.1287\\
2.3878	0.13083\\
2.392	0.13299\\
2.3962	0.13518\\
2.4004	0.1374\\
2.4046	0.13965\\
2.4088	0.14193\\
2.413	0.14424\\
2.4172	0.14658\\
2.4214	0.14895\\
2.4256	0.15135\\
2.4298	0.15378\\
2.434	0.15625\\
2.4382	0.15875\\
2.4424	0.16128\\
2.4465	0.16384\\
2.4507	0.16643\\
2.4549	0.16906\\
2.4591	0.17172\\
2.4633	0.17441\\
2.4675	0.17713\\
2.4717	0.17989\\
2.4759	0.18268\\
2.4801	0.18551\\
2.4843	0.18836\\
2.4885	0.19126\\
2.4927	0.19418\\
2.4969	0.19714\\
2.5011	0.20014\\
2.5052	0.20317\\
2.5094	0.20623\\
2.5136	0.20933\\
2.5178	0.21246\\
2.522	0.21563\\
2.5262	0.21883\\
2.5304	0.22206\\
2.5346	0.22534\\
2.5388	0.22864\\
2.543	0.23198\\
2.5472	0.23536\\
2.5514	0.23877\\
2.5556	0.24221\\
2.5598	0.24569\\
2.5639	0.2492\\
2.5681	0.25275\\
2.5723	0.25634\\
2.5765	0.25995\\
2.5807	0.26361\\
2.5849	0.26729\\
2.5891	0.27101\\
2.5933	0.27477\\
2.5975	0.27856\\
2.6017	0.28238\\
2.6059	0.28624\\
2.6101	0.29013\\
2.6143	0.29406\\
2.6185	0.29801\\
2.6227	0.30201\\
2.6268	0.30603\\
2.631	0.31009\\
2.6352	0.31418\\
2.6394	0.3183\\
2.6436	0.32245\\
2.6478	0.32664\\
2.652	0.33086\\
2.6562	0.3351\\
2.6604	0.33938\\
2.6646	0.34369\\
2.6688	0.34803\\
2.673	0.3524\\
2.6772	0.3568\\
2.6814	0.36123\\
2.6855	0.36569\\
2.6897	0.37018\\
2.6939	0.37469\\
2.6981	0.37923\\
2.7023	0.3838\\
2.7065	0.3884\\
2.7107	0.39302\\
2.7149	0.39767\\
2.7191	0.40234\\
2.7233	0.40704\\
2.7275	0.41177\\
2.7317	0.41651\\
2.7359	0.42128\\
2.7401	0.42608\\
2.7442	0.43089\\
2.7484	0.43573\\
2.7526	0.44059\\
2.7568	0.44547\\
2.761	0.45037\\
2.7652	0.45529\\
2.7694	0.46022\\
2.7736	0.46518\\
2.7778	0.47015\\
2.782	0.47514\\
2.7862	0.48015\\
2.7904	0.48517\\
2.7946	0.4902\\
2.7988	0.49525\\
2.803	0.50032\\
2.8071	0.50539\\
2.8113	0.51048\\
2.8155	0.51558\\
2.8197	0.5207\\
2.8239	0.52582\\
2.8281	0.53095\\
2.8323	0.53609\\
2.8365	0.54124\\
2.8407	0.54639\\
2.8449	0.55156\\
2.8491	0.55672\\
2.8533	0.5619\\
2.8575	0.56708\\
2.8617	0.57227\\
2.8658	0.57746\\
2.87	0.58265\\
2.8742	0.58785\\
2.8784	0.59305\\
2.8826	0.59826\\
2.8868	0.60347\\
2.891	0.60868\\
2.8952	0.6139\\
2.8994	0.61912\\
2.9036	0.62435\\
2.9078	0.62958\\
2.912	0.63481\\
2.9162	0.64006\\
2.9204	0.6453\\
2.9246	0.65056\\
2.9287	0.65582\\
2.9329	0.66108\\
2.9371	0.66635\\
2.9413	0.67162\\
2.9455	0.6769\\
2.9497	0.68217\\
2.9539	0.68744\\
2.9581	0.69271\\
2.9623	0.69796\\
2.9665	0.7032\\
2.9707	0.70842\\
2.9749	0.71362\\
2.9791	0.71879\\
2.9833	0.72393\\
2.9874	0.72903\\
2.9916	0.73409\\
2.9958	0.73909\\
3	0.74404\\
3.0042	0.74894\\
3.0084	0.75376\\
3.0126	0.75852\\
3.0168	0.76321\\
3.021	0.76781\\
3.0252	0.77234\\
3.0294	0.77678\\
3.0336	0.78113\\
3.0378	0.78539\\
3.042	0.78956\\
3.0461	0.79363\\
3.0503	0.79761\\
3.0545	0.80149\\
3.0587	0.80527\\
3.0629	0.80895\\
3.0671	0.81253\\
3.0713	0.81602\\
3.0755	0.81941\\
3.0797	0.82271\\
3.0839	0.82591\\
3.0881	0.82903\\
3.0923	0.83205\\
3.0965	0.83498\\
3.1007	0.83783\\
3.1049	0.8406\\
3.109	0.84328\\
3.1132	0.84589\\
3.1174	0.84842\\
3.1216	0.85087\\
3.1258	0.85324\\
3.13	0.85555\\
3.1342	0.85777\\
3.1384	0.85992\\
3.1426	0.862\\
3.1468	0.864\\
3.151	0.86592\\
3.1552	0.86776\\
3.1594	0.86953\\
3.1636	0.87121\\
3.1677	0.87281\\
3.1719	0.87432\\
3.1761	0.87574\\
3.1803	0.87708\\
3.1845	0.87833\\
3.1887	0.87948\\
3.1929	0.88054\\
3.1971	0.8815\\
3.2013	0.88237\\
3.2055	0.88314\\
3.2097	0.88381\\
3.2139	0.88439\\
3.2181	0.88486\\
3.2223	0.88523\\
3.2264	0.8855\\
3.2306	0.88566\\
3.2348	0.88573\\
3.239	0.8857\\
3.2432	0.88556\\
3.2474	0.88532\\
3.2516	0.88498\\
3.2558	0.88455\\
3.26	0.88401\\
3.2642	0.88338\\
3.2684	0.88264\\
3.2726	0.88182\\
3.2768	0.8809\\
3.281	0.87989\\
3.2852	0.87878\\
3.2893	0.8776\\
3.2935	0.87632\\
3.2977	0.87497\\
3.3019	0.87353\\
3.3061	0.87202\\
3.3103	0.87045\\
3.3145	0.8688\\
3.3187	0.8671\\
3.3229	0.86533\\
3.3271	0.86352\\
3.3313	0.86166\\
3.3355	0.85976\\
3.3397	0.85782\\
3.3439	0.85584\\
3.348	0.85383\\
3.3522	0.85179\\
3.3564	0.8497\\
3.3606	0.84758\\
3.3648	0.84542\\
3.369	0.8432\\
3.3732	0.84092\\
3.3774	0.83858\\
3.3816	0.83616\\
3.3858	0.83366\\
3.39	0.83106\\
3.3942	0.82837\\
3.3984	0.82557\\
3.4026	0.82265\\
3.4067	0.81962\\
3.4109	0.81647\\
3.4151	0.8132\\
3.4193	0.8098\\
3.4235	0.80628\\
3.4277	0.80264\\
3.4319	0.79888\\
3.4361	0.79501\\
3.4403	0.79104\\
3.4445	0.78696\\
3.4487	0.7828\\
3.4529	0.77856\\
3.4571	0.77425\\
3.4613	0.76987\\
3.4655	0.76545\\
3.4696	0.76099\\
3.4738	0.7565\\
3.478	0.75198\\
3.4822	0.74746\\
3.4864	0.74294\\
3.4906	0.73842\\
3.4948	0.73391\\
3.499	0.72941\\
3.5032	0.72494\\
3.5074	0.72048\\
3.5116	0.71604\\
3.5158	0.71161\\
3.52	0.70721\\
3.5242	0.70281\\
3.5283	0.69843\\
3.5325	0.69406\\
3.5367	0.68969\\
3.5409	0.68533\\
3.5451	0.68097\\
3.5493	0.67661\\
3.5535	0.67225\\
3.5577	0.66789\\
3.5619	0.66353\\
3.5661	0.65916\\
3.5703	0.65479\\
3.5745	0.65042\\
3.5787	0.64605\\
3.5829	0.64167\\
3.5871	0.63729\\
3.5912	0.6329\\
3.5954	0.62852\\
3.5996	0.62413\\
3.6038	0.61973\\
3.608	0.61534\\
3.6122	0.61095\\
3.6164	0.60655\\
3.6206	0.60216\\
3.6248	0.59776\\
3.629	0.59337\\
3.6332	0.58897\\
3.6374	0.58458\\
3.6416	0.58019\\
3.6458	0.57581\\
3.6499	0.57143\\
3.6541	0.56705\\
3.6583	0.56268\\
3.6625	0.55831\\
3.6667	0.55395\\
3.6709	0.54959\\
3.6751	0.54524\\
3.6793	0.5409\\
3.6835	0.53657\\
3.6877	0.53224\\
3.6919	0.52793\\
3.6961	0.52362\\
3.7003	0.51933\\
3.7045	0.51504\\
3.7086	0.51077\\
3.7128	0.5065\\
3.717	0.50225\\
3.7212	0.49801\\
3.7254	0.49379\\
3.7296	0.48958\\
3.7338	0.48538\\
3.738	0.4812\\
3.7422	0.47703\\
3.7464	0.47287\\
3.7506	0.46873\\
3.7548	0.46461\\
3.759	0.46051\\
3.7632	0.45642\\
3.7674	0.45234\\
3.7715	0.44829\\
3.7757	0.44425\\
3.7799	0.44023\\
3.7841	0.43623\\
3.7883	0.43225\\
3.7925	0.42828\\
3.7967	0.42434\\
3.8009	0.42042\\
3.8051	0.41651\\
3.8093	0.41263\\
3.8135	0.40876\\
3.8177	0.40492\\
3.8219	0.40109\\
3.8261	0.39729\\
3.8302	0.39351\\
3.8344	0.38975\\
3.8386	0.38602\\
3.8428	0.3823\\
3.847	0.37861\\
3.8512	0.37494\\
3.8554	0.37129\\
3.8596	0.36766\\
3.8638	0.36406\\
3.868	0.36048\\
3.8722	0.35692\\
3.8764	0.35339\\
3.8806	0.34988\\
3.8848	0.34639\\
3.8889	0.34293\\
3.8931	0.33949\\
3.8973	0.33608\\
3.9015	0.33268\\
3.9057	0.32932\\
3.9099	0.32597\\
3.9141	0.32265\\
3.9183	0.31936\\
3.9225	0.31608\\
3.9267	0.31284\\
3.9309	0.30961\\
3.9351	0.30641\\
3.9393	0.30324\\
3.9435	0.30009\\
3.9477	0.29696\\
3.9518	0.29386\\
3.956	0.29078\\
3.9602	0.28773\\
3.9644	0.2847\\
3.9686	0.28169\\
3.9728	0.27871\\
3.977	0.27576\\
3.9812	0.27282\\
3.9854	0.26992\\
3.9896	0.26703\\
3.9938	0.26417\\
3.998	0.26134\\
4.0022	0.25852\\
4.0064	0.25573\\
4.0105	0.25297\\
4.0147	0.25023\\
4.0189	0.24751\\
4.0231	0.24481\\
4.0273	0.24214\\
4.0315	0.2395\\
4.0357	0.23687\\
4.0399	0.23427\\
4.0441	0.23169\\
4.0483	0.22914\\
4.0525	0.2266\\
4.0567	0.22409\\
4.0609	0.22161\\
4.0651	0.21914\\
4.0693	0.2167\\
4.0734	0.21428\\
4.0776	0.21188\\
4.0818	0.20951\\
4.086	0.20715\\
4.0902	0.20482\\
4.0944	0.20251\\
4.0986	0.20022\\
4.1028	0.19796\\
4.107	0.19571\\
4.1112	0.19349\\
4.1154	0.19128\\
4.1196	0.1891\\
4.1238	0.18694\\
4.128	0.1848\\
4.1321	0.18268\\
4.1363	0.18058\\
4.1405	0.1785\\
4.1447	0.17644\\
4.1489	0.1744\\
4.1531	0.17238\\
4.1573	0.17038\\
4.1615	0.1684\\
4.1657	0.16644\\
4.1699	0.1645\\
4.1741	0.16257\\
4.1783	0.16067\\
4.1825	0.15879\\
4.1867	0.15692\\
4.1908	0.15508\\
4.195	0.15325\\
4.1992	0.15144\\
4.2034	0.14965\\
4.2076	0.14788\\
4.2118	0.14612\\
4.216	0.14438\\
4.2202	0.14266\\
4.2244	0.14096\\
4.2286	0.13928\\
4.2328	0.13761\\
4.237	0.13596\\
4.2412	0.13433\\
4.2454	0.13271\\
4.2496	0.13111\\
4.2537	0.12953\\
4.2579	0.12796\\
4.2621	0.12642\\
4.2663	0.12488\\
4.2705	0.12336\\
4.2747	0.12186\\
4.2789	0.12038\\
4.2831	0.11891\\
4.2873	0.11745\\
4.2915	0.11601\\
4.2957	0.11459\\
4.2999	0.11318\\
4.3041	0.11179\\
4.3083	0.11041\\
4.3124	0.10905\\
4.3166	0.1077\\
4.3208	0.10636\\
4.325	0.10504\\
4.3292	0.10374\\
4.3334	0.10245\\
4.3376	0.10117\\
4.3418	0.099904\\
4.346	0.098654\\
4.3502	0.097417\\
4.3544	0.096194\\
4.3586	0.094984\\
4.3628	0.093788\\
4.367	0.092604\\
4.3711	0.091434\\
4.3753	0.090276\\
4.3795	0.089131\\
4.3837	0.087999\\
4.3879	0.086879\\
4.3921	0.085772\\
4.3963	0.084677\\
4.4005	0.083594\\
4.4047	0.082523\\
4.4089	0.081464\\
4.4131	0.080417\\
4.4173	0.079382\\
4.4215	0.078358\\
4.4257	0.077346\\
4.4299	0.076344\\
4.434	0.075355\\
4.4382	0.074376\\
4.4424	0.073408\\
4.4466	0.072451\\
4.4508	0.071505\\
4.455	0.07057\\
4.4592	0.069645\\
4.4634	0.068731\\
4.4676	0.067827\\
4.4718	0.066934\\
4.476	0.06605\\
4.4802	0.065177\\
4.4844	0.064313\\
4.4886	0.06346\\
4.4927	0.062616\\
4.4969	0.061782\\
4.5011	0.060957\\
4.5053	0.060142\\
4.5095	0.059336\\
4.5137	0.05854\\
4.5179	0.057753\\
4.5221	0.056974\\
4.5263	0.056205\\
4.5305	0.055445\\
4.5347	0.054694\\
4.5389	0.053951\\
4.5431	0.053217\\
4.5473	0.052491\\
4.5515	0.051774\\
4.5556	0.051066\\
4.5598	0.050366\\
4.564	0.049673\\
4.5682	0.04899\\
4.5724	0.048314\\
4.5766	0.047646\\
4.5808	0.046986\\
4.585	0.046334\\
4.5892	0.04569\\
4.5934	0.045053\\
4.5976	0.044424\\
4.6018	0.043803\\
4.606	0.043189\\
4.6102	0.042582\\
4.6143	0.041983\\
4.6185	0.041391\\
4.6227	0.040806\\
4.6269	0.040229\\
4.6311	0.039658\\
4.6353	0.039094\\
4.6395	0.038538\\
4.6437	0.037988\\
4.6479	0.037445\\
4.6521	0.036908\\
4.6563	0.036379\\
4.6605	0.035856\\
4.6647	0.035339\\
4.6689	0.034829\\
4.673	0.034325\\
4.6772	0.033828\\
4.6814	0.033337\\
4.6856	0.032852\\
4.6898	0.032373\\
4.694	0.0319\\
4.6982	0.031434\\
4.7024	0.030973\\
4.7066	0.030518\\
4.7108	0.030069\\
4.715	0.029626\\
4.7192	0.029189\\
4.7234	0.028757\\
4.7276	0.028331\\
4.7318	0.02791\\
4.7359	0.027495\\
4.7401	0.027085\\
4.7443	0.02668\\
4.7485	0.026281\\
4.7527	0.025887\\
4.7569	0.025498\\
4.7611	0.025114\\
4.7653	0.024736\\
4.7695	0.024362\\
4.7737	0.023993\\
4.7779	0.023629\\
4.7821	0.023269\\
4.7863	0.022915\\
4.7905	0.022564\\
4.7946	0.022219\\
4.7988	0.021878\\
4.803	0.021542\\
4.8072	0.021209\\
4.8114	0.020882\\
4.8156	0.020558\\
4.8198	0.020239\\
4.824	0.019924\\
4.8282	0.019613\\
4.8324	0.019306\\
4.8366	0.019003\\
4.8408	0.018704\\
4.845	0.018409\\
4.8492	0.018118\\
4.8533	0.01783\\
4.8575	0.017546\\
4.8617	0.017266\\
4.8659	0.01699\\
4.8701	0.016717\\
4.8743	0.016448\\
4.8785	0.016182\\
4.8827	0.01592\\
4.8869	0.015661\\
4.8911	0.015406\\
4.8953	0.015154\\
4.8995	0.014905\\
4.9037	0.01466\\
4.9079	0.014417\\
4.9121	0.014178\\
4.9162	0.013943\\
4.9204	0.01371\\
4.9246	0.01348\\
4.9288	0.013254\\
4.933	0.01303\\
4.9372	0.01281\\
4.9414	0.012592\\
4.9456	0.012377\\
4.9498	0.012166\\
4.954	0.011957\\
4.9582	0.011751\\
4.9624	0.011547\\
4.9666	0.011347\\
4.9708	0.011149\\
4.9749	0.010955\\
4.9791	0.010762\\
4.9833	0.010573\\
4.9875	0.010386\\
4.9917	0.010202\\
4.9959	0.01002\\
5.0001	0.0098414\\
5.0043	0.009665\\
5.0085	0.0094911\\
5.0127	0.0093198\\
5.0169	0.009151\\
5.0211	0.0089846\\
5.0253	0.0088208\\
5.0295	0.0086593\\
5.0336	0.0085002\\
5.0378	0.0083436\\
5.042	0.0081893\\
5.0462	0.0080374\\
5.0504	0.0078877\\
5.0546	0.0077404\\
5.0588	0.0075954\\
5.063	0.0074526\\
5.0672	0.0073121\\
5.0714	0.0071737\\
5.0756	0.0070376\\
5.0798	0.0069037\\
5.084	0.0067718\\
5.0882	0.0066422\\
5.0924	0.0065146\\
5.0965	0.0063891\\
5.1007	0.0062657\\
5.1049	0.0061443\\
5.1091	0.0060249\\
5.1133	0.0059076\\
5.1175	0.0057922\\
5.1217	0.0056787\\
5.1259	0.0055672\\
5.1301	0.0054576\\
5.1343	0.0053499\\
5.1385	0.005244\\
5.1427	0.00514\\
5.1469	0.0050378\\
5.1511	0.0049373\\
5.1552	0.0048387\\
5.1594	0.0047418\\
5.1636	0.0046466\\
5.1678	0.0045531\\
5.172	0.0044613\\
5.1762	0.0043712\\
5.1804	0.0042827\\
5.1846	0.0041958\\
5.1888	0.0041105\\
5.193	0.0040267\\
5.1972	0.0039445\\
5.2014	0.0038639\\
5.2056	0.0037847\\
5.2098	0.0037069\\
5.214	0.0036307\\
5.2181	0.0035558\\
5.2223	0.0034824\\
5.2265	0.0034104\\
5.2307	0.0033397\\
5.2349	0.0032703\\
5.2391	0.0032023\\
5.2433	0.0031356\\
5.2475	0.0030701\\
5.2517	0.0030059\\
5.2559	0.002943\\
5.2601	0.0028812\\
5.2643	0.0028207\\
5.2685	0.0027613\\
5.2727	0.0027031\\
5.2768	0.002646\\
5.281	0.00259\\
5.2852	0.0025352\\
5.2894	0.0024814\\
5.2936	0.0024286\\
5.2978	0.0023769\\
5.302	0.0023262\\
5.3062	0.0022765\\
5.3104	0.0022278\\
5.3146	0.0021801\\
5.3188	0.0021333\\
5.323	0.0020875\\
5.3272	0.0020425\\
5.3314	0.0019985\\
5.3355	0.0019554\\
5.3397	0.0019131\\
5.3439	0.0018716\\
5.3481	0.001831\\
5.3523	0.0017913\\
5.3565	0.0017523\\
5.3607	0.0017141\\
5.3607	0.0014278\\
5.3565	0.0014618\\
5.3523	0.0014966\\
5.3481	0.0015322\\
5.3439	0.0015685\\
5.3397	0.0016056\\
5.3355	0.0016436\\
5.3314	0.0016823\\
5.3272	0.0017219\\
5.323	0.0017624\\
5.3188	0.0018037\\
5.3146	0.001846\\
5.3104	0.0018891\\
5.3062	0.0019331\\
5.302	0.0019781\\
5.2978	0.002024\\
5.2936	0.0020709\\
5.2894	0.0021188\\
5.2852	0.0021677\\
5.281	0.0022176\\
5.2768	0.0022685\\
5.2727	0.0023205\\
5.2685	0.0023736\\
5.2643	0.0024278\\
5.2601	0.002483\\
5.2559	0.0025394\\
5.2517	0.0025969\\
5.2475	0.0026556\\
5.2433	0.0027155\\
5.2391	0.0027765\\
5.2349	0.0028388\\
5.2307	0.0029023\\
5.2265	0.0029671\\
5.2223	0.0030331\\
5.2181	0.0031005\\
5.214	0.0031691\\
5.2098	0.0032391\\
5.2056	0.0033104\\
5.2014	0.0033831\\
5.1972	0.0034572\\
5.193	0.0035327\\
5.1888	0.0036096\\
5.1846	0.003688\\
5.1804	0.0037679\\
5.1762	0.0038492\\
5.172	0.0039321\\
5.1678	0.0040165\\
5.1636	0.0041024\\
5.1594	0.00419\\
5.1552	0.0042791\\
5.1511	0.0043699\\
5.1469	0.0044623\\
5.1427	0.0045564\\
5.1385	0.0046522\\
5.1343	0.0047497\\
5.1301	0.0048489\\
5.1259	0.0049499\\
5.1217	0.0050527\\
5.1175	0.0051573\\
5.1133	0.0052638\\
5.1091	0.0053721\\
5.1049	0.0054823\\
5.1007	0.0055944\\
5.0965	0.0057085\\
5.0924	0.0058245\\
5.0882	0.0059426\\
5.084	0.0060626\\
5.0798	0.0061847\\
5.0756	0.0063089\\
5.0714	0.0064351\\
5.0672	0.0065635\\
5.063	0.0066941\\
5.0588	0.0068268\\
5.0546	0.0069618\\
5.0504	0.007099\\
5.0462	0.0072385\\
5.042	0.0073803\\
5.0378	0.0075244\\
5.0336	0.0076709\\
5.0295	0.0078198\\
5.0253	0.007971\\
5.0211	0.0081248\\
5.0169	0.008281\\
5.0127	0.0084398\\
5.0085	0.0086011\\
5.0043	0.0087649\\
5.0001	0.0089314\\
4.9959	0.0091006\\
4.9917	0.0092724\\
4.9875	0.0094469\\
4.9833	0.0096241\\
4.9791	0.0098041\\
4.9749	0.0099869\\
4.9708	0.010173\\
4.9666	0.010361\\
4.9624	0.010552\\
4.9582	0.010747\\
4.954	0.010944\\
4.9498	0.011144\\
4.9456	0.011348\\
4.9414	0.011554\\
4.9372	0.011764\\
4.933	0.011976\\
4.9288	0.012192\\
4.9246	0.012411\\
4.9204	0.012633\\
4.9162	0.012859\\
4.9121	0.013087\\
4.9079	0.01332\\
4.9037	0.013555\\
4.8995	0.013794\\
4.8953	0.014036\\
4.8911	0.014282\\
4.8869	0.014531\\
4.8827	0.014784\\
4.8785	0.01504\\
4.8743	0.0153\\
4.8701	0.015564\\
4.8659	0.015831\\
4.8617	0.016102\\
4.8575	0.016377\\
4.8533	0.016655\\
4.8492	0.016937\\
4.845	0.017223\\
4.8408	0.017513\\
4.8366	0.017807\\
4.8324	0.018105\\
4.8282	0.018406\\
4.824	0.018712\\
4.8198	0.019022\\
4.8156	0.019335\\
4.8114	0.019653\\
4.8072	0.019975\\
4.803	0.020301\\
4.7988	0.020632\\
4.7946	0.020966\\
4.7905	0.021305\\
4.7863	0.021649\\
4.7821	0.021996\\
4.7779	0.022348\\
4.7737	0.022705\\
4.7695	0.023065\\
4.7653	0.023431\\
4.7611	0.023801\\
4.7569	0.024175\\
4.7527	0.024555\\
4.7485	0.024938\\
4.7443	0.025327\\
4.7401	0.02572\\
4.7359	0.026119\\
4.7318	0.026522\\
4.7276	0.02693\\
4.7234	0.027343\\
4.7192	0.027761\\
4.715	0.028184\\
4.7108	0.028613\\
4.7066	0.029046\\
4.7024	0.029485\\
4.6982	0.02993\\
4.694	0.030379\\
4.6898	0.030834\\
4.6856	0.031295\\
4.6814	0.031761\\
4.6772	0.032233\\
4.673	0.032711\\
4.6689	0.033194\\
4.6647	0.033683\\
4.6605	0.034179\\
4.6563	0.03468\\
4.6521	0.035187\\
4.6479	0.035701\\
4.6437	0.036221\\
4.6395	0.036747\\
4.6353	0.03728\\
4.6311	0.037819\\
4.6269	0.038364\\
4.6227	0.038916\\
4.6185	0.039475\\
4.6143	0.040041\\
4.6102	0.040614\\
4.606	0.041193\\
4.6018	0.04178\\
4.5976	0.042374\\
4.5934	0.042975\\
4.5892	0.043583\\
4.585	0.044199\\
4.5808	0.044822\\
4.5766	0.045453\\
4.5724	0.046091\\
4.5682	0.046737\\
4.564	0.047391\\
4.5598	0.048053\\
4.5556	0.048723\\
4.5515	0.049401\\
4.5473	0.050087\\
4.5431	0.050781\\
4.5389	0.051484\\
4.5347	0.052195\\
4.5305	0.052915\\
4.5263	0.053643\\
4.5221	0.05438\\
4.5179	0.055126\\
4.5137	0.055881\\
4.5095	0.056645\\
4.5053	0.057418\\
4.5011	0.0582\\
4.4969	0.058992\\
4.4927	0.059793\\
4.4886	0.060603\\
4.4844	0.061423\\
4.4802	0.062253\\
4.476	0.063092\\
4.4718	0.063942\\
4.4676	0.064801\\
4.4634	0.065671\\
4.4592	0.066551\\
4.455	0.067441\\
4.4508	0.068341\\
4.4466	0.069252\\
4.4424	0.070174\\
4.4382	0.071106\\
4.434	0.07205\\
4.4299	0.073004\\
4.4257	0.073969\\
4.4215	0.074946\\
4.4173	0.075934\\
4.4131	0.076933\\
4.4089	0.077944\\
4.4047	0.078966\\
4.4005	0.08\\
4.3963	0.081047\\
4.3921	0.082105\\
4.3879	0.083175\\
4.3837	0.084258\\
4.3795	0.085352\\
4.3753	0.08646\\
4.3711	0.08758\\
4.367	0.088713\\
4.3628	0.089858\\
4.3586	0.091017\\
4.3544	0.092189\\
4.3502	0.093374\\
4.346	0.094572\\
4.3418	0.095784\\
4.3376	0.09701\\
4.3334	0.098249\\
4.3292	0.099502\\
4.325	0.10077\\
4.3208	0.10205\\
4.3166	0.10335\\
4.3124	0.10466\\
4.3083	0.10598\\
4.3041	0.10732\\
4.2999	0.10868\\
4.2957	0.11005\\
4.2915	0.11143\\
4.2873	0.11283\\
4.2831	0.11425\\
4.2789	0.11568\\
4.2747	0.11712\\
4.2705	0.11859\\
4.2663	0.12007\\
4.2621	0.12156\\
4.2579	0.12307\\
4.2537	0.1246\\
4.2496	0.12614\\
4.2454	0.1277\\
4.2412	0.12928\\
4.237	0.13088\\
4.2328	0.13249\\
4.2286	0.13412\\
4.2244	0.13576\\
4.2202	0.13743\\
4.216	0.13911\\
4.2118	0.14081\\
4.2076	0.14253\\
4.2034	0.14426\\
4.1992	0.14602\\
4.195	0.14779\\
4.1908	0.14958\\
4.1867	0.15139\\
4.1825	0.15322\\
4.1783	0.15507\\
4.1741	0.15693\\
4.1699	0.15882\\
4.1657	0.16072\\
4.1615	0.16265\\
4.1573	0.16459\\
4.1531	0.16656\\
4.1489	0.16854\\
4.1447	0.17054\\
4.1405	0.17257\\
4.1363	0.17461\\
4.1321	0.17668\\
4.128	0.17876\\
4.1238	0.18087\\
4.1196	0.18299\\
4.1154	0.18514\\
4.1112	0.18731\\
4.107	0.1895\\
4.1028	0.19171\\
4.0986	0.19394\\
4.0944	0.19619\\
4.0902	0.19846\\
4.086	0.20076\\
4.0818	0.20307\\
4.0776	0.20541\\
4.0734	0.20777\\
4.0693	0.21016\\
4.0651	0.21256\\
4.0609	0.21499\\
4.0567	0.21743\\
4.0525	0.2199\\
4.0483	0.2224\\
4.0441	0.22491\\
4.0399	0.22745\\
4.0357	0.23\\
4.0315	0.23259\\
4.0273	0.23519\\
4.0231	0.23782\\
4.0189	0.24046\\
4.0147	0.24314\\
4.0105	0.24583\\
4.0064	0.24855\\
4.0022	0.25128\\
3.998	0.25405\\
3.9938	0.25683\\
3.9896	0.25964\\
3.9854	0.26247\\
3.9812	0.26532\\
3.977	0.26819\\
3.9728	0.27109\\
3.9686	0.27401\\
3.9644	0.27695\\
3.9602	0.27992\\
3.956	0.28291\\
3.9518	0.28592\\
3.9477	0.28895\\
3.9435	0.29201\\
3.9393	0.29508\\
3.9351	0.29818\\
3.9309	0.30131\\
3.9267	0.30445\\
3.9225	0.30762\\
3.9183	0.31081\\
3.9141	0.31402\\
3.9099	0.31726\\
3.9057	0.32051\\
3.9015	0.32379\\
3.8973	0.32709\\
3.8931	0.33041\\
3.8889	0.33376\\
3.8848	0.33712\\
3.8806	0.34051\\
3.8764	0.34392\\
3.8722	0.34735\\
3.868	0.3508\\
3.8638	0.35427\\
3.8596	0.35777\\
3.8554	0.36128\\
3.8512	0.36482\\
3.847	0.36838\\
3.8428	0.37195\\
3.8386	0.37555\\
3.8344	0.37917\\
3.8302	0.38281\\
3.8261	0.38647\\
3.8219	0.39015\\
3.8177	0.39384\\
3.8135	0.39756\\
3.8093	0.4013\\
3.8051	0.40506\\
3.8009	0.40883\\
3.7967	0.41263\\
3.7925	0.41644\\
3.7883	0.42028\\
3.7841	0.42413\\
3.7799	0.428\\
3.7757	0.43189\\
3.7715	0.43579\\
3.7674	0.43972\\
3.7632	0.44366\\
3.759	0.44762\\
3.7548	0.45159\\
3.7506	0.45558\\
3.7464	0.45959\\
3.7422	0.46362\\
3.738	0.46766\\
3.7338	0.47172\\
3.7296	0.47579\\
3.7254	0.47988\\
3.7212	0.48398\\
3.717	0.4881\\
3.7128	0.49223\\
3.7086	0.49638\\
3.7045	0.50054\\
3.7003	0.50471\\
3.6961	0.5089\\
3.6919	0.5131\\
3.6877	0.51731\\
3.6835	0.52154\\
3.6793	0.52578\\
3.6751	0.53002\\
3.6709	0.53428\\
3.6667	0.53855\\
3.6625	0.54283\\
3.6583	0.54712\\
3.6541	0.55142\\
3.6499	0.55573\\
3.6458	0.56005\\
3.6416	0.56438\\
3.6374	0.56871\\
3.6332	0.57305\\
3.629	0.5774\\
3.6248	0.58176\\
3.6206	0.58612\\
3.6164	0.59049\\
3.6122	0.59487\\
3.608	0.59925\\
3.6038	0.60364\\
3.5996	0.60803\\
3.5954	0.61242\\
3.5912	0.61683\\
3.5871	0.62123\\
3.5829	0.62565\\
3.5787	0.63006\\
3.5745	0.63448\\
3.5703	0.63891\\
3.5661	0.64334\\
3.5619	0.64777\\
3.5577	0.65221\\
3.5535	0.65665\\
3.5493	0.66109\\
3.5451	0.66553\\
3.5409	0.66996\\
3.5367	0.6744\\
3.5325	0.67882\\
3.5283	0.68323\\
3.5242	0.68763\\
3.52	0.69201\\
3.5158	0.69636\\
3.5116	0.70069\\
3.5074	0.70498\\
3.5032	0.70923\\
3.499	0.71344\\
3.4948	0.7176\\
3.4906	0.72172\\
3.4864	0.72578\\
3.4822	0.72979\\
3.478	0.73374\\
3.4738	0.73765\\
3.4696	0.7415\\
3.4655	0.74531\\
3.4613	0.74907\\
3.4571	0.75279\\
3.4529	0.75646\\
3.4487	0.76009\\
3.4445	0.76369\\
3.4403	0.76725\\
3.4361	0.77077\\
3.4319	0.77426\\
3.4277	0.77772\\
3.4235	0.78113\\
3.4193	0.78452\\
3.4151	0.78786\\
3.4109	0.79117\\
3.4067	0.79443\\
3.4026	0.79765\\
3.3984	0.80081\\
3.3942	0.80393\\
3.39	0.80698\\
3.3858	0.80996\\
3.3816	0.81288\\
3.3774	0.81571\\
3.3732	0.81846\\
3.369	0.82113\\
3.3648	0.82369\\
3.3606	0.82616\\
3.3564	0.82853\\
3.3522	0.83079\\
3.348	0.83294\\
3.3439	0.83498\\
3.3397	0.83692\\
3.3355	0.83876\\
3.3313	0.84049\\
3.3271	0.84213\\
3.3229	0.84368\\
3.3187	0.84515\\
3.3145	0.84653\\
3.3103	0.84783\\
3.3061	0.84905\\
3.3019	0.85021\\
3.2977	0.85129\\
3.2935	0.8523\\
3.2893	0.85324\\
3.2852	0.85411\\
3.281	0.85491\\
3.2768	0.85564\\
3.2726	0.85631\\
3.2684	0.8569\\
3.2642	0.85743\\
3.26	0.85788\\
3.2558	0.85826\\
3.2516	0.85856\\
3.2474	0.85879\\
3.2432	0.85895\\
3.239	0.85903\\
3.2348	0.85903\\
3.2306	0.85895\\
3.2264	0.8588\\
3.2223	0.85857\\
3.2181	0.85825\\
3.2139	0.85786\\
3.2097	0.85738\\
3.2055	0.85683\\
3.2013	0.85619\\
3.1971	0.85547\\
3.1929	0.85466\\
3.1887	0.85377\\
3.1845	0.85279\\
3.1803	0.85173\\
3.1761	0.85057\\
3.1719	0.84933\\
3.1677	0.84799\\
3.1636	0.84657\\
3.1594	0.84504\\
3.1552	0.84343\\
3.151	0.84171\\
3.1468	0.8399\\
3.1426	0.83798\\
3.1384	0.83597\\
3.1342	0.83385\\
3.13	0.83164\\
3.1258	0.82932\\
3.1216	0.8269\\
3.1174	0.82438\\
3.1132	0.82176\\
3.109	0.81904\\
3.1049	0.81623\\
3.1007	0.81332\\
3.0965	0.81032\\
3.0923	0.80723\\
3.0881	0.80406\\
3.0839	0.8008\\
3.0797	0.79746\\
3.0755	0.79404\\
3.0713	0.79054\\
3.0671	0.78698\\
3.0629	0.78334\\
3.0587	0.77964\\
3.0545	0.77587\\
3.0503	0.77204\\
3.0461	0.76815\\
3.042	0.76421\\
3.0378	0.7602\\
3.0336	0.75615\\
3.0294	0.75204\\
3.0252	0.74788\\
3.021	0.74367\\
3.0168	0.73941\\
3.0126	0.7351\\
3.0084	0.73074\\
3.0042	0.72634\\
3	0.72189\\
2.9958	0.7174\\
2.9916	0.71285\\
2.9874	0.70827\\
2.9833	0.70363\\
2.9791	0.69895\\
2.9749	0.69423\\
2.9707	0.68945\\
2.9665	0.68464\\
2.9623	0.67977\\
2.9581	0.67487\\
2.9539	0.66992\\
2.9497	0.66492\\
2.9455	0.65989\\
2.9413	0.65482\\
2.9371	0.64971\\
2.9329	0.64457\\
2.9287	0.63939\\
2.9246	0.63419\\
2.9204	0.62897\\
2.9162	0.62372\\
2.912	0.61845\\
2.9078	0.61317\\
2.9036	0.60788\\
2.8994	0.60257\\
2.8952	0.59727\\
2.891	0.59196\\
2.8868	0.58664\\
2.8826	0.58133\\
2.8784	0.57603\\
2.8742	0.57073\\
2.87	0.56543\\
2.8658	0.56015\\
2.8617	0.55488\\
2.8575	0.54961\\
2.8533	0.54437\\
2.8491	0.53913\\
2.8449	0.53391\\
2.8407	0.5287\\
2.8365	0.52351\\
2.8323	0.51834\\
2.8281	0.51318\\
2.8239	0.50804\\
2.8197	0.50292\\
2.8155	0.49782\\
2.8113	0.49273\\
2.8071	0.48767\\
2.803	0.48263\\
2.7988	0.4776\\
2.7946	0.4726\\
2.7904	0.46762\\
2.7862	0.46266\\
2.782	0.45772\\
2.7778	0.45281\\
2.7736	0.44791\\
2.7694	0.44304\\
2.7652	0.4382\\
2.761	0.43337\\
2.7568	0.42857\\
2.7526	0.4238\\
2.7484	0.41905\\
2.7442	0.41433\\
2.7401	0.40963\\
2.7359	0.40496\\
2.7317	0.40031\\
2.7275	0.3957\\
2.7233	0.3911\\
2.7191	0.38654\\
2.7149	0.382\\
2.7107	0.3775\\
2.7065	0.37302\\
2.7023	0.36856\\
2.6981	0.36414\\
2.6939	0.35975\\
2.6897	0.35538\\
2.6855	0.35105\\
2.6814	0.34675\\
2.6772	0.34247\\
2.673	0.33823\\
2.6688	0.33401\\
2.6646	0.32983\\
2.6604	0.32568\\
2.6562	0.32156\\
2.652	0.31747\\
2.6478	0.31341\\
2.6436	0.30938\\
2.6394	0.30539\\
2.6352	0.30142\\
2.631	0.29749\\
2.6268	0.29359\\
2.6227	0.28973\\
2.6185	0.28589\\
2.6143	0.28209\\
2.6101	0.27832\\
2.6059	0.27458\\
2.6017	0.27088\\
2.5975	0.2672\\
2.5933	0.26356\\
2.5891	0.25996\\
2.5849	0.25638\\
2.5807	0.25284\\
2.5765	0.24933\\
2.5723	0.24586\\
2.5681	0.24241\\
2.5639	0.239\\
2.5598	0.23563\\
2.5556	0.23228\\
2.5514	0.22897\\
2.5472	0.22569\\
2.543	0.22244\\
2.5388	0.21923\\
2.5346	0.21605\\
2.5304	0.2129\\
2.5262	0.20978\\
2.522	0.2067\\
2.5178	0.20364\\
2.5136	0.20062\\
2.5094	0.19764\\
2.5052	0.19468\\
2.5011	0.19176\\
2.4969	0.18887\\
2.4927	0.18601\\
2.4885	0.18318\\
2.4843	0.18038\\
2.4801	0.17762\\
2.4759	0.17488\\
2.4717	0.17218\\
2.4675	0.16951\\
2.4633	0.16687\\
2.4591	0.16426\\
2.4549	0.16168\\
2.4507	0.15913\\
2.4465	0.15661\\
2.4424	0.15413\\
2.4382	0.15167\\
2.434	0.14924\\
2.4298	0.14684\\
2.4256	0.14448\\
2.4214	0.14214\\
2.4172	0.13983\\
2.413	0.13755\\
2.4088	0.1353\\
2.4046	0.13308\\
2.4004	0.13088\\
2.3962	0.12872\\
2.392	0.12658\\
2.3878	0.12447\\
2.3836	0.12239\\
2.3795	0.12034\\
2.3753	0.11831\\
2.3711	0.11631\\
2.3669	0.11434\\
2.3627	0.1124\\
2.3585	0.11048\\
2.3543	0.10859\\
2.3501	0.10672\\
2.3459	0.10489\\
2.3417	0.10307\\
2.3375	0.10129\\
2.3333	0.099524\\
2.3291	0.097787\\
2.3249	0.096075\\
2.3208	0.094389\\
2.3166	0.092726\\
2.3124	0.091088\\
2.3082	0.089474\\
2.304	0.087884\\
2.2998	0.086318\\
2.2956	0.084775\\
2.2914	0.083255\\
2.2872	0.081758\\
2.283	0.080284\\
2.2788	0.078833\\
2.2746	0.077403\\
2.2704	0.075996\\
2.2662	0.074611\\
2.262	0.073247\\
2.2579	0.071904\\
2.2537	0.070582\\
2.2495	0.069282\\
2.2453	0.068002\\
2.2411	0.066742\\
2.2369	0.065503\\
2.2327	0.064283\\
2.2285	0.063083\\
2.2243	0.061903\\
2.2201	0.060741\\
2.2159	0.059599\\
2.2117	0.058476\\
2.2075	0.057371\\
2.2033	0.056284\\
2.1992	0.055215\\
2.195	0.054165\\
2.1908	0.053132\\
2.1866	0.052116\\
2.1824	0.051117\\
2.1782	0.050135\\
2.174	0.04917\\
2.1698	0.048222\\
2.1656	0.047289\\
2.1614	0.046373\\
2.1572	0.045473\\
2.153	0.044588\\
2.1488	0.043718\\
2.1446	0.042864\\
2.1405	0.042025\\
2.1363	0.0412\\
2.1321	0.04039\\
2.1279	0.039594\\
2.1237	0.038813\\
2.1195	0.038045\\
2.1153	0.037291\\
2.1111	0.03655\\
2.1069	0.035823\\
2.1027	0.035109\\
2.0985	0.034408\\
2.0943	0.033719\\
2.0901	0.033044\\
2.0859	0.03238\\
2.0817	0.031728\\
2.0776	0.031089\\
2.0734	0.030461\\
2.0692	0.029845\\
2.065	0.02924\\
2.0608	0.028646\\
2.0566	0.028064\\
2.0524	0.027492\\
2.0482	0.026931\\
2.044	0.026381\\
2.0398	0.025841\\
2.0356	0.025311\\
2.0314	0.024791\\
2.0272	0.024281\\
2.023	0.02378\\
2.0189	0.02329\\
2.0147	0.022808\\
2.0105	0.022336\\
2.0063	0.021873\\
2.0021	0.021418\\
1.9979	0.020973\\
1.9937	0.020536\\
1.9895	0.020107\\
1.9853	0.019687\\
1.9811	0.019275\\
1.9769	0.018871\\
1.9727	0.018475\\
1.9685	0.018087\\
1.9643	0.017706\\
1.9602	0.017333\\
1.956	0.016967\\
1.9518	0.016608\\
1.9476	0.016257\\
1.9434	0.015912\\
1.9392	0.015574\\
1.935	0.015243\\
1.9308	0.014919\\
1.9266	0.014601\\
1.9224	0.014289\\
1.9182	0.013984\\
1.914	0.013685\\
1.9098	0.013392\\
1.9056	0.013104\\
1.9014	0.012823\\
1.8973	0.012547\\
1.8931	0.012277\\
1.8889	0.012012\\
1.8847	0.011753\\
1.8805	0.011499\\
1.8763	0.01125\\
1.8721	0.011006\\
1.8679	0.010768\\
1.8637	0.010534\\
1.8595	0.010305\\
1.8553	0.01008\\
1.8511	0.0098608\\
1.8469	0.0096456\\
1.8427	0.0094349\\
1.8386	0.0092285\\
1.8344	0.0090264\\
1.8302	0.0088286\\
1.826	0.0086348\\
1.8218	0.0084451\\
1.8176	0.0082593\\
1.8134	0.0080774\\
1.8092	0.0078994\\
1.805	0.007725\\
1.8008	0.0075543\\
1.7966	0.0073873\\
1.7924	0.0072237\\
1.7882	0.0070636\\
1.784	0.0069068\\
1.7799	0.0067534\\
1.7757	0.0066032\\
1.7715	0.0064562\\
1.7673	0.0063123\\
1.7631	0.0061715\\
1.7589	0.0060337\\
1.7547	0.0058988\\
1.7505	0.0057668\\
1.7463	0.0056376\\
1.7421	0.0055111\\
1.7379	0.0053874\\
1.7337	0.0052663\\
1.7295	0.0051479\\
1.7253	0.0050319\\
1.7211	0.0049185\\
1.717	0.0048075\\
1.7128	0.0046989\\
1.7086	0.0045927\\
1.7044	0.0044887\\
1.7002	0.004387\\
1.696	0.0042875\\
1.6918	0.0041902\\
1.6876	0.0040949\\
1.6834	0.0040018\\
1.6792	0.0039106\\
1.675	0.0038215\\
1.6708	0.0037343\\
1.6666	0.003649\\
1.6624	0.0035656\\
1.6583	0.003484\\
1.6541	0.0034041\\
1.6499	0.0033261\\
1.6457	0.0032497\\
1.6415	0.003175\\
1.6373	0.003102\\
1.6331	0.0030306\\
1.6289	0.0029607\\
1.6247	0.0028924\\
1.6205	0.0028256\\
1.6163	0.0027603\\
1.6121	0.0026964\\
1.6079	0.002634\\
1.6037	0.0025729\\
1.5995	0.0025132\\
1.5954	0.0024548\\
1.5912	0.0023977\\
1.587	0.0023419\\
1.5828	0.0022873\\
1.5786	0.002234\\
1.5744	0.0021818\\
1.5702	0.0021308\\
1.566	0.0020809\\
1.5618	0.0020322\\
1.5576	0.0019846\\
1.5534	0.001938\\
1.5492	0.0018925\\
1.545	0.001848\\
1.5408	0.0018045\\
1.5367	0.0017619\\
1.5325	0.0017204\\
1.5283	0.0016797\\
1.5241	0.00164\\
1.5199	0.0016012\\
1.5157	0.0015633\\
1.5115	0.0015262\\
1.5073	0.00149\\
1.5031	0.0014546\\
1.4989	0.00142\\
1.4947	0.0013862\\
1.4905	0.0013532\\
1.4863	0.0013209\\
1.4821	0.0012893\\
1.478	0.0012585\\
1.4738	0.0012284\\
1.4696	0.001199\\
1.4654	0.0011702\\
1.4612	0.0011421\\
1.457	0.0011146\\
1.4528	0.0010878\\
1.4486	0.0010616\\
1.4444	0.001036\\
1.4402	0.001011\\
1.436	0.00098657\\
1.4318	0.0009627\\
1.4276	0.00093939\\
1.4234	0.00091661\\
1.4192	0.00089436\\
1.4151	0.00087263\\
1.4109	0.0008514\\
1.4067	0.00083067\\
1.4025	0.00081042\\
1.3983	0.00079064\\
1.3941	0.00077132\\
1.3899	0.00075245\\
1.3857	0.00073403\\
1.3815	0.00071603\\
1.3773	0.00069846\\
1.3731	0.0006813\\
1.3689	0.00066454\\
1.3647	0.00064817\\
1.3605	0.00063219\\
1.3564	0.00061659\\
1.3522	0.00060135\\
1.348	0.00058647\\
1.3438	0.00057195\\
1.3396	0.00055776\\
1.3354	0.00054392\\
1.3312	0.0005304\\
1.327	0.0005172\\
1.3228	0.00050432\\
1.3186	0.00049174\\
1.3144	0.00047946\\
1.3102	0.00046747\\
1.306	0.00045577\\
1.3018	0.00044435\\
1.2977	0.0004332\\
1.2935	0.00042232\\
1.2893	0.0004117\\
1.2851	0.00040134\\
1.2809	0.00039122\\
1.2767	0.00038135\\
1.2725	0.00037171\\
1.2683	0.00036231\\
1.2641	0.00035313\\
1.2599	0.00034417\\
1.2557	0.00033544\\
1.2515	0.00032691\\
1.2473	0.00031859\\
1.2431	0.00031047\\
1.2389	0.00030255\\
1.2348	0.00029482\\
1.2306	0.00028728\\
1.2264	0.00027992\\
1.2222	0.00027274\\
1.218	0.00026574\\
1.2138	0.00025891\\
1.2096	0.00025225\\
1.2054	0.00024575\\
1.2012	0.00023941\\
1.197	0.00023322\\
1.1928	0.00022719\\
1.1886	0.00022131\\
1.1844	0.00021557\\
1.1802	0.00020997\\
1.1761	0.00020452\\
1.1719	0.00019919\\
}--cycle;
\addplot[ybar interval, fill=black, fill opacity=0.35, area legend, draw=none] table[row sep=crcr, x=Lower, y=Count] {%
Lower	Upper	Count\\
1.1719	1.4941	0\\
1.4941	1.8163	0\\
1.8163	2.1385	0.0082759\\
2.1385	2.4607	0.091035\\
2.4607	2.783	0.28966\\
2.783	3.1052	0.77794\\
3.1052	3.4274	0.85242\\
3.4274	3.7496	0.64552\\
3.7496	4.0718	0.26483\\
4.0718	4.394	0.099311\\
4.394	4.7163	0.049656\\
4.7163	5.0385	0.0082759\\
5.0385	5.3607	0.016552\\
5.3607	5.3607	0.016552\\
};
\addlegendentry{\RealHistogramLegend}

\addplot [color=black]
  table[row sep=crcr]{%
1.1719	2.782e-06\\
1.6205	0.00038661\\
1.7463	0.0012431\\
1.826	0.00248\\
1.8889	0.0041646\\
1.9392	0.0061978\\
1.9853	0.0088046\\
2.0272	0.011981\\
2.065	0.015668\\
2.0985	0.019744\\
2.1279	0.024038\\
2.1572	0.029115\\
2.1824	0.034172\\
2.2075	0.039954\\
2.2327	0.046537\\
2.2579	0.053999\\
2.2788	0.060946\\
2.2998	0.068606\\
2.3208	0.077024\\
2.3417	0.086248\\
2.3627	0.096321\\
2.3836	0.10729\\
2.4046	0.11919\\
2.4256	0.13206\\
2.4465	0.14593\\
2.4675	0.16084\\
2.4885	0.1768\\
2.5094	0.19383\\
2.5304	0.21194\\
2.5514	0.23114\\
2.5723	0.25141\\
2.5933	0.27274\\
2.6143	0.2951\\
2.6352	0.31844\\
2.6604	0.3477\\
2.6855	0.37821\\
2.7107	0.40983\\
2.7401	0.44794\\
2.7736	0.49271\\
2.8281	0.56701\\
2.8826	0.64101\\
2.9162	0.6852\\
2.9413	0.71716\\
2.9665	0.74779\\
2.9874	0.77208\\
3.0084	0.79507\\
3.0252	0.81243\\
3.042	0.82878\\
3.0587	0.84404\\
3.0755	0.85816\\
3.0923	0.87105\\
3.1049	0.87989\\
3.1174	0.88799\\
3.13	0.89532\\
3.1426	0.90186\\
3.1552	0.90761\\
3.1677	0.91255\\
3.1803	0.91665\\
3.1929	0.91992\\
3.2055	0.92234\\
3.2181	0.92392\\
3.2306	0.92463\\
3.2432	0.92449\\
3.2558	0.9235\\
3.2684	0.92166\\
3.281	0.91898\\
3.2935	0.91546\\
3.3061	0.91112\\
3.3187	0.90597\\
3.3313	0.90003\\
3.3439	0.89331\\
3.3564	0.88584\\
3.369	0.87764\\
3.3816	0.86873\\
3.3984	0.85579\\
3.4151	0.8417\\
3.4319	0.82653\\
3.4487	0.81036\\
3.4655	0.79325\\
3.4864	0.7707\\
3.5074	0.74698\\
3.5325	0.71721\\
3.5619	0.68106\\
3.5954	0.63842\\
3.6625	0.55143\\
3.7086	0.49235\\
3.7422	0.45065\\
3.7715	0.41543\\
3.7967	0.38638\\
3.8219	0.35851\\
3.847	0.33191\\
3.868	0.31079\\
3.8889	0.29063\\
3.9099	0.27147\\
3.9309	0.25332\\
3.9518	0.23618\\
3.9728	0.22004\\
3.9938	0.2049\\
4.0147	0.19073\\
4.0357	0.17751\\
4.0567	0.16521\\
4.0776	0.15378\\
4.0986	0.1432\\
4.1196	0.13342\\
4.1405	0.1244\\
4.1615	0.11609\\
4.1867	0.10701\\
4.2118	0.098805\\
4.237	0.091411\\
4.2621	0.084751\\
4.2915	0.07781\\
4.3208	0.07166\\
4.3544	0.065472\\
4.3879	0.060048\\
4.4257	0.054708\\
4.4676	0.049548\\
4.5137	0.044615\\
4.5682	0.03955\\
4.6311	0.034467\\
4.7066	0.029141\\
4.7946	0.023719\\
4.8911	0.018579\\
4.9875	0.014204\\
5.084	0.01057\\
5.1846	0.0075325\\
5.2936	0.0050334\\
5.3607	0.003854\\
};
\addlegendentry{\RealBLNFitLegend}

      % \input{../../../../.tex/20/KS/SizeDistributionFittingsfig30.tex}
      \legend{}
    \end{groupplot}
  
    \begin{groupplot}[
      group style={
        group name=Row5,
        plotLastGroup,
      },
      plotRow,
      row1Keys,
      xlabel3,
      yTickScaleLabelEmpty,
      % emptyLeged,
      %
      % xmin=1,xmax=2.3788e+05,
      ymin=0,ymax=1,
    ]
      \nextgroupplot[%
        at={($(Row4 c1r1.south west)-(0,\yPlotSep em)-(0,\plotHeight cm)$)},
        xmin=0.57363,xmax=5.4115,
        % xmin=0.33174,xmax=5.6534,
        % ymin=0,ymax=0.9,
      ] % This file was created by matlab2tikz.
\definecolor{mycolor1}{rgb}{0.00000,0.44706,0.69804}%
%
\definecolor{mycolor2}{rgb}{0.83529,0.36863,0.00000}%
%
\addplot[ybar interval, fill=mycolor1, fill opacity=0.35, area legend, draw=none] table[row sep=crcr, x=Lower, y=Count] {%
Lower	Upper	Count\\
0.57363	0.97679	0.00026458\\
0.97679	1.3799	0.0011245\\
1.3799	1.7831	0.007739\\
1.7831	2.1863	0.043457\\
2.1863	2.5894	0.18441\\
2.5894	2.9926	0.50998\\
2.9926	3.3957	0.76894\\
3.3957	3.7989	0.58823\\
3.7989	4.202	0.25836\\
4.202	4.6052	0.091545\\
4.6052	5.0083	0.019579\\
5.0083	5.4115	0.0068129\\
5.4115	5.4115	0.0068129\\
};
\addlegendentry{\ExactHistogramLegend}

\addplot [color=mycolor1, only marks, every error bar/.append style={opacity=0.45}, mark=*, mark size=0pt, draw=none, forget plot]
 plot [error bars/.cd, y dir=both, y explicit, error bar style={line width=1pt, color=mycolor1}, error mark options={mark=none,mark size=0pt}]
 table[row sep=crcr, y error plus index=2, y error minus index=3]{%
0.77521	0.00026458	0.00025848	0.00025848\\
1.1784	0.0011245	0.00052944	0.00052944\\
1.5815	0.007739	0.0013065	0.0013065\\
1.9847	0.043457	0.0030553	0.0030553\\
2.3878	0.18441	0.0062326	0.0062326\\
2.791	0.50998	0.009613	0.009613\\
3.1941	0.76894	0.013269	0.013269\\
3.5973	0.58823	0.010197	0.010197\\
4.0004	0.25836	0.0063228	0.0063228\\
4.4036	0.091545	0.0034129	0.0034129\\
4.8068	0.019579	0.0015599	0.0015599\\
5.2099	0.0068129	0.00039374	0.00039374\\
};
\addplot [color=mycolor1]
  table[row sep=crcr]{%
0.57363	5.2469e-05\\
1.0385	0.00050662\\
1.2371	0.0013408\\
1.3678	0.0025235\\
1.4695	0.0040947\\
1.5519	0.0060202\\
1.6245	0.0084119\\
1.6875	0.011188\\
1.7407	0.014188\\
1.7892	0.017552\\
1.8327	0.021197\\
1.8763	0.025529\\
1.9151	0.030044\\
1.9538	0.035274\\
1.9877	0.04051\\
2.0216	0.046435\\
2.0555	0.053122\\
2.0846	0.059521\\
2.1136	0.066594\\
2.1427	0.074396\\
2.1717	0.082986\\
2.1959	0.090787\\
2.2202	0.099206\\
2.2444	0.10827\\
2.2686	0.11802\\
2.2928	0.12847\\
2.317	0.13963\\
2.3412	0.15153\\
2.3654	0.16419\\
2.3896	0.1776\\
2.4139	0.19179\\
2.4381	0.20676\\
2.4623	0.22253\\
2.4865	0.23911\\
2.5107	0.25652\\
2.5349	0.27475\\
2.5591	0.2938\\
2.5834	0.31365\\
2.6124	0.33848\\
2.6415	0.3643\\
2.6754	0.39548\\
2.7141	0.43217\\
2.7625	0.47906\\
2.9175	0.63\\
2.9514	0.6616\\
2.9756	0.68335\\
2.9998	0.7042\\
3.0192	0.72006\\
3.0386	0.73506\\
3.0579	0.74906\\
3.0725	0.75882\\
3.087	0.76786\\
3.1015	0.77612\\
3.1161	0.78356\\
3.1306	0.7901\\
3.1451	0.79569\\
3.1596	0.8003\\
3.1742	0.80388\\
3.1838	0.80569\\
3.1935	0.80703\\
3.2032	0.80789\\
3.2129	0.80829\\
3.2226	0.80822\\
3.2323	0.80768\\
3.2468	0.80602\\
3.2613	0.80337\\
3.2759	0.79977\\
3.2904	0.79527\\
3.3049	0.78991\\
3.3194	0.78376\\
3.334	0.77687\\
3.3485	0.7693\\
3.363	0.7611\\
3.3824	0.74929\\
3.4018	0.73658\\
3.4211	0.72307\\
3.4454	0.70523\\
3.4696	0.68649\\
3.4986	0.66304\\
3.5325	0.63464\\
3.5713	0.60124\\
3.6342	0.54585\\
3.7117	0.47782\\
3.7553	0.44047\\
3.794	0.40819\\
3.8279	0.38082\\
3.8618	0.35439\\
3.8909	0.33254\\
3.9199	0.3115\\
3.949	0.29129\\
3.9781	0.27194\\
4.0071	0.25346\\
4.0362	0.23587\\
4.0652	0.21916\\
4.0943	0.20332\\
4.1233	0.18836\\
4.1524	0.17424\\
4.1814	0.16096\\
4.2105	0.14849\\
4.2396	0.13679\\
4.2686	0.12585\\
4.2977	0.11563\\
4.3267	0.10611\\
4.3558	0.097239\\
4.3897	0.087683\\
4.4236	0.078927\\
4.4575	0.07092\\
4.4914	0.063613\\
4.5253	0.056957\\
4.564	0.050087\\
4.6028	0.04394\\
4.6415	0.038454\\
4.6851	0.032999\\
4.7287	0.028227\\
4.7771	0.023639\\
4.8255	0.019715\\
4.8788	0.016069\\
4.9369	0.01278\\
4.9999	0.0099004\\
5.0725	0.0073042\\
5.1548	0.0051094\\
5.2468	0.0033723\\
5.3582	0.0019954\\
5.4115	0.0015408\\
};
\addlegendentry{\ExactBLNFitLegend}


\addplot[area legend, draw=none, fill=mycolor1, fill opacity=0.15, forget plot]
table[row sep=crcr] {%
x	y\\
0.57363	7.5805e-05\\
0.57848	7.7445e-05\\
0.58332	7.9122e-05\\
0.58816	8.0836e-05\\
0.59301	8.2588e-05\\
0.59785	8.4378e-05\\
0.60269	8.6208e-05\\
0.60753	8.8079e-05\\
0.61238	8.9992e-05\\
0.61722	9.1947e-05\\
0.62206	9.3945e-05\\
0.6269	9.5988e-05\\
0.63175	9.8076e-05\\
0.63659	0.00010021\\
0.64143	0.00010239\\
0.64628	0.00010463\\
0.65112	0.00010691\\
0.65596	0.00010924\\
0.6608	0.00011162\\
0.66565	0.00011406\\
0.67049	0.00011655\\
0.67533	0.0001191\\
0.68017	0.0001217\\
0.68502	0.00012437\\
0.68986	0.00012709\\
0.6947	0.00012988\\
0.69954	0.00013272\\
0.70439	0.00013563\\
0.70923	0.00013861\\
0.71407	0.00014165\\
0.71892	0.00014477\\
0.72376	0.00014795\\
0.7286	0.0001512\\
0.73344	0.00015453\\
0.73829	0.00015793\\
0.74313	0.00016141\\
0.74797	0.00016497\\
0.75281	0.00016861\\
0.75766	0.00017233\\
0.7625	0.00017613\\
0.76734	0.00018002\\
0.77219	0.000184\\
0.77703	0.00018807\\
0.78187	0.00019223\\
0.78671	0.00019649\\
0.79156	0.00020084\\
0.7964	0.00020529\\
0.80124	0.00020984\\
0.80608	0.0002145\\
0.81093	0.00021926\\
0.81577	0.00022413\\
0.82061	0.00022911\\
0.82545	0.0002342\\
0.8303	0.00023941\\
0.83514	0.00024474\\
0.83998	0.00025019\\
0.84483	0.00025577\\
0.84967	0.00026147\\
0.85451	0.0002673\\
0.85935	0.00027327\\
0.8642	0.00027937\\
0.86904	0.00028561\\
0.87388	0.00029199\\
0.87872	0.00029852\\
0.88357	0.0003052\\
0.88841	0.00031203\\
0.89325	0.00031902\\
0.8981	0.00032617\\
0.90294	0.00033348\\
0.90778	0.00034096\\
0.91262	0.00034861\\
0.91747	0.00035643\\
0.92231	0.00036444\\
0.92715	0.00037263\\
0.93199	0.000381\\
0.93684	0.00038957\\
0.94168	0.00039834\\
0.94652	0.0004073\\
0.95137	0.00041647\\
0.95621	0.00042585\\
0.96105	0.00043545\\
0.96589	0.00044527\\
0.97074	0.00045531\\
0.97558	0.00046558\\
0.98042	0.00047609\\
0.98526	0.00048684\\
0.99011	0.00049784\\
0.99495	0.00050908\\
0.99979	0.00052059\\
1.0046	0.00053236\\
1.0095	0.0005444\\
1.0143	0.00055672\\
1.0192	0.00056932\\
1.024	0.00058221\\
1.0288	0.0005954\\
1.0337	0.00060889\\
1.0385	0.00062269\\
1.0434	0.0006368\\
1.0482	0.00065124\\
1.0531	0.00066601\\
1.0579	0.00068112\\
1.0627	0.00069657\\
1.0676	0.00071238\\
1.0724	0.00072856\\
1.0773	0.0007451\\
1.0821	0.00076202\\
1.087	0.00077933\\
1.0918	0.00079704\\
1.0966	0.00081515\\
1.1015	0.00083368\\
1.1063	0.00085264\\
1.1112	0.00087202\\
1.116	0.00089185\\
1.1209	0.00091214\\
1.1257	0.00093288\\
1.1305	0.00095411\\
1.1354	0.00097581\\
1.1402	0.00099802\\
1.1451	0.0010207\\
1.1499	0.001044\\
1.1548	0.0010677\\
1.1596	0.001092\\
1.1644	0.0011169\\
1.1693	0.0011423\\
1.1741	0.0011683\\
1.179	0.0011949\\
1.1838	0.0012221\\
1.1887	0.0012499\\
1.1935	0.0012784\\
1.1983	0.0013075\\
1.2032	0.0013372\\
1.208	0.0013677\\
1.2129	0.0013988\\
1.2177	0.0014306\\
1.2226	0.0014632\\
1.2274	0.0014965\\
1.2322	0.0015305\\
1.2371	0.0015653\\
1.2419	0.0016009\\
1.2468	0.0016373\\
1.2516	0.0016746\\
1.2565	0.0017126\\
1.2613	0.0017516\\
1.2661	0.0017914\\
1.271	0.0018321\\
1.2758	0.0018737\\
1.2807	0.0019162\\
1.2855	0.0019598\\
1.2904	0.0020042\\
1.2952	0.0020497\\
1.3	0.0020962\\
1.3049	0.0021438\\
1.3097	0.0021924\\
1.3146	0.0022421\\
1.3194	0.0022929\\
1.3243	0.0023449\\
1.3291	0.002398\\
1.3339	0.0024523\\
1.3388	0.0025078\\
1.3436	0.0025645\\
1.3485	0.0026225\\
1.3533	0.0026818\\
1.3582	0.0027424\\
1.363	0.0028043\\
1.3678	0.0028676\\
1.3727	0.0029323\\
1.3775	0.0029985\\
1.3824	0.0030661\\
1.3872	0.0031352\\
1.3921	0.0032058\\
1.3969	0.0032779\\
1.4017	0.0033517\\
1.4066	0.0034271\\
1.4114	0.0035041\\
1.4163	0.0035828\\
1.4211	0.0036633\\
1.4259	0.0037454\\
1.4308	0.0038294\\
1.4356	0.0039153\\
1.4405	0.004003\\
1.4453	0.0040926\\
1.4502	0.0041841\\
1.455	0.0042776\\
1.4598	0.0043732\\
1.4647	0.0044708\\
1.4695	0.0045706\\
1.4744	0.0046725\\
1.4792	0.0047766\\
1.4841	0.0048829\\
1.4889	0.0049915\\
1.4937	0.0051025\\
1.4986	0.0052158\\
1.5034	0.0053316\\
1.5083	0.0054498\\
1.5131	0.0055706\\
1.518	0.005694\\
1.5228	0.0058199\\
1.5276	0.0059486\\
1.5325	0.00608\\
1.5373	0.0062142\\
1.5422	0.0063512\\
1.547	0.0064911\\
1.5519	0.006634\\
1.5567	0.0067799\\
1.5615	0.0069288\\
1.5664	0.0070809\\
1.5712	0.0072362\\
1.5761	0.0073947\\
1.5809	0.0075566\\
1.5858	0.0077218\\
1.5906	0.0078905\\
1.5954	0.0080627\\
1.6003	0.0082385\\
1.6051	0.0084179\\
1.61	0.0086011\\
1.6148	0.008788\\
1.6197	0.0089788\\
1.6245	0.0091736\\
1.6293	0.0093723\\
1.6342	0.0095751\\
1.639	0.0097821\\
1.6439	0.0099934\\
1.6487	0.010209\\
1.6536	0.010429\\
1.6584	0.010653\\
1.6632	0.010882\\
1.6681	0.011116\\
1.6729	0.011354\\
1.6778	0.011598\\
1.6826	0.011846\\
1.6875	0.012099\\
1.6923	0.012357\\
1.6971	0.01262\\
1.702	0.012889\\
1.7068	0.013163\\
1.7117	0.013443\\
1.7165	0.013728\\
1.7214	0.014018\\
1.7262	0.014315\\
1.731	0.014617\\
1.7359	0.014926\\
1.7407	0.01524\\
1.7456	0.015561\\
1.7504	0.015888\\
1.7553	0.016221\\
1.7601	0.016561\\
1.7649	0.016908\\
1.7698	0.017261\\
1.7746	0.017621\\
1.7795	0.017988\\
1.7843	0.018362\\
1.7892	0.018744\\
1.794	0.019133\\
1.7988	0.019529\\
1.8037	0.019933\\
1.8085	0.020344\\
1.8134	0.020764\\
1.8182	0.021191\\
1.8231	0.021627\\
1.8279	0.022071\\
1.8327	0.022523\\
1.8376	0.022984\\
1.8424	0.023453\\
1.8473	0.023932\\
1.8521	0.024419\\
1.8569	0.024915\\
1.8618	0.025421\\
1.8666	0.025936\\
1.8715	0.026461\\
1.8763	0.026995\\
1.8812	0.02754\\
1.886	0.028094\\
1.8908	0.028658\\
1.8957	0.029233\\
1.9005	0.029819\\
1.9054	0.030415\\
1.9102	0.031022\\
1.9151	0.03164\\
1.9199	0.032269\\
1.9247	0.03291\\
1.9296	0.033562\\
1.9344	0.034226\\
1.9393	0.034901\\
1.9441	0.035589\\
1.949	0.036289\\
1.9538	0.037002\\
1.9586	0.037727\\
1.9635	0.038465\\
1.9683	0.039216\\
1.9732	0.039981\\
1.978	0.040759\\
1.9829	0.04155\\
1.9877	0.042355\\
1.9925	0.043175\\
1.9974	0.044008\\
2.0022	0.044856\\
2.0071	0.045719\\
2.0119	0.046597\\
2.0168	0.047489\\
2.0216	0.048398\\
2.0264	0.049321\\
2.0313	0.050261\\
2.0361	0.051216\\
2.041	0.052188\\
2.0458	0.053176\\
2.0507	0.054181\\
2.0555	0.055203\\
2.0603	0.056242\\
2.0652	0.057298\\
2.07	0.058372\\
2.0749	0.059464\\
2.0797	0.060575\\
2.0846	0.061704\\
2.0894	0.062851\\
2.0942	0.064018\\
2.0991	0.065203\\
2.1039	0.066409\\
2.1088	0.067634\\
2.1136	0.068879\\
2.1185	0.070145\\
2.1233	0.071431\\
2.1281	0.072738\\
2.133	0.074067\\
2.1378	0.075417\\
2.1427	0.076789\\
2.1475	0.078183\\
2.1524	0.079599\\
2.1572	0.081039\\
2.162	0.082501\\
2.1669	0.083987\\
2.1717	0.085496\\
2.1766	0.08703\\
2.1814	0.088588\\
2.1863	0.090171\\
2.1911	0.091778\\
2.1959	0.093411\\
2.2008	0.09507\\
2.2056	0.096755\\
2.2105	0.098466\\
2.2153	0.1002\\
2.2202	0.10197\\
2.225	0.10376\\
2.2298	0.10558\\
2.2347	0.10743\\
2.2395	0.1093\\
2.2444	0.1112\\
2.2492	0.11314\\
2.2541	0.1151\\
2.2589	0.11709\\
2.2637	0.11911\\
2.2686	0.12115\\
2.2734	0.12323\\
2.2783	0.12534\\
2.2831	0.12748\\
2.2879	0.12964\\
2.2928	0.13184\\
2.2976	0.13407\\
2.3025	0.13633\\
2.3073	0.13862\\
2.3122	0.14093\\
2.317	0.14328\\
2.3218	0.14566\\
2.3267	0.14807\\
2.3315	0.15051\\
2.3364	0.15298\\
2.3412	0.15548\\
2.3461	0.15802\\
2.3509	0.16058\\
2.3557	0.16317\\
2.3606	0.16579\\
2.3654	0.16844\\
2.3703	0.17113\\
2.3751	0.17384\\
2.38	0.17658\\
2.3848	0.17936\\
2.3896	0.18216\\
2.3945	0.18499\\
2.3993	0.18786\\
2.4042	0.19075\\
2.409	0.19368\\
2.4139	0.19663\\
2.4187	0.19962\\
2.4235	0.20264\\
2.4284	0.20569\\
2.4332	0.20877\\
2.4381	0.21188\\
2.4429	0.21502\\
2.4478	0.21819\\
2.4526	0.2214\\
2.4574	0.22463\\
2.4623	0.2279\\
2.4671	0.2312\\
2.472	0.23454\\
2.4768	0.2379\\
2.4817	0.24131\\
2.4865	0.24474\\
2.4913	0.24821\\
2.4962	0.25171\\
2.501	0.25524\\
2.5059	0.25881\\
2.5107	0.26242\\
2.5156	0.26606\\
2.5204	0.26974\\
2.5252	0.27345\\
2.5301	0.27719\\
2.5349	0.28097\\
2.5398	0.28479\\
2.5446	0.28864\\
2.5495	0.29253\\
2.5543	0.29645\\
2.5591	0.30041\\
2.564	0.3044\\
2.5688	0.30843\\
2.5737	0.3125\\
2.5785	0.31659\\
2.5834	0.32072\\
2.5882	0.32489\\
2.593	0.32908\\
2.5979	0.33331\\
2.6027	0.33757\\
2.6076	0.34186\\
2.6124	0.34619\\
2.6173	0.35053\\
2.6221	0.35491\\
2.6269	0.35932\\
2.6318	0.36375\\
2.6366	0.3682\\
2.6415	0.37268\\
2.6463	0.37718\\
2.6512	0.3817\\
2.656	0.38624\\
2.6608	0.3908\\
2.6657	0.39538\\
2.6705	0.39997\\
2.6754	0.40458\\
2.6802	0.40921\\
2.6851	0.41384\\
2.6899	0.41849\\
2.6947	0.42315\\
2.6996	0.42781\\
2.7044	0.43249\\
2.7093	0.43717\\
2.7141	0.44186\\
2.7189	0.44655\\
2.7238	0.45125\\
2.7286	0.45596\\
2.7335	0.46067\\
2.7383	0.46538\\
2.7432	0.47009\\
2.748	0.47481\\
2.7528	0.47953\\
2.7577	0.48425\\
2.7625	0.48897\\
2.7674	0.4937\\
2.7722	0.49842\\
2.7771	0.50315\\
2.7819	0.50788\\
2.7867	0.5126\\
2.7916	0.51733\\
2.7964	0.52206\\
2.8013	0.52679\\
2.8061	0.53151\\
2.811	0.53624\\
2.8158	0.54097\\
2.8206	0.5457\\
2.8255	0.55042\\
2.8303	0.55514\\
2.8352	0.55986\\
2.84	0.56458\\
2.8449	0.5693\\
2.8497	0.57401\\
2.8545	0.57872\\
2.8594	0.58342\\
2.8642	0.58811\\
2.8691	0.5928\\
2.8739	0.59749\\
2.8788	0.60216\\
2.8836	0.60683\\
2.8884	0.61149\\
2.8933	0.61613\\
2.8981	0.62077\\
2.903	0.6254\\
2.9078	0.63001\\
2.9127	0.6346\\
2.9175	0.63919\\
2.9223	0.64376\\
2.9272	0.64831\\
2.932	0.65285\\
2.9369	0.65736\\
2.9417	0.66186\\
2.9466	0.66634\\
2.9514	0.6708\\
2.9562	0.67524\\
2.9611	0.67965\\
2.9659	0.68404\\
2.9708	0.6884\\
2.9756	0.69274\\
2.9805	0.69705\\
2.9853	0.70133\\
2.9901	0.70557\\
2.995	0.70979\\
2.9998	0.71396\\
3.0047	0.7181\\
3.0095	0.7222\\
3.0144	0.72626\\
3.0192	0.73028\\
3.024	0.73424\\
3.0289	0.73816\\
3.0337	0.74202\\
3.0386	0.74583\\
3.0434	0.74958\\
3.0483	0.75327\\
3.0531	0.75689\\
3.0579	0.76045\\
3.0628	0.76393\\
3.0676	0.76735\\
3.0725	0.77068\\
3.0773	0.77394\\
3.0822	0.77711\\
3.087	0.7802\\
3.0918	0.7832\\
3.0967	0.7861\\
3.1015	0.78891\\
3.1064	0.79163\\
3.1112	0.79424\\
3.1161	0.79675\\
3.1209	0.79915\\
3.1257	0.80144\\
3.1306	0.80363\\
3.1354	0.8057\\
3.1403	0.80765\\
3.1451	0.80949\\
3.15	0.8112\\
3.1548	0.8128\\
3.1596	0.81427\\
3.1645	0.81563\\
3.1693	0.81685\\
3.1742	0.81795\\
3.179	0.81893\\
3.1838	0.81977\\
3.1887	0.82049\\
3.1935	0.82109\\
3.1984	0.82155\\
3.2032	0.82189\\
3.2081	0.82211\\
3.2129	0.8222\\
3.2177	0.82216\\
3.2226	0.822\\
3.2274	0.82172\\
3.2323	0.82131\\
3.2371	0.82079\\
3.242	0.82015\\
3.2468	0.81939\\
3.2516	0.81852\\
3.2565	0.81754\\
3.2613	0.81644\\
3.2662	0.81524\\
3.271	0.81393\\
3.2759	0.81252\\
3.2807	0.811\\
3.2855	0.80938\\
3.2904	0.80767\\
3.2952	0.80586\\
3.3001	0.80396\\
3.3049	0.80196\\
3.3098	0.79988\\
3.3146	0.7977\\
3.3194	0.79545\\
3.3243	0.79311\\
3.3291	0.79069\\
3.334	0.7882\\
3.3388	0.78562\\
3.3437	0.78298\\
3.3485	0.78026\\
3.3533	0.77747\\
3.3582	0.77462\\
3.363	0.7717\\
3.3679	0.76872\\
3.3727	0.76568\\
3.3776	0.76258\\
3.3824	0.75942\\
3.3872	0.75621\\
3.3921	0.75294\\
3.3969	0.74963\\
3.4018	0.74626\\
3.4066	0.74285\\
3.4115	0.7394\\
3.4163	0.7359\\
3.4211	0.73236\\
3.426	0.72878\\
3.4308	0.72516\\
3.4357	0.7215\\
3.4405	0.71781\\
3.4454	0.71409\\
3.4502	0.71033\\
3.455	0.70654\\
3.4599	0.70273\\
3.4647	0.69888\\
3.4696	0.695\\
3.4744	0.6911\\
3.4793	0.68718\\
3.4841	0.68322\\
3.4889	0.67925\\
3.4938	0.67525\\
3.4986	0.67123\\
3.5035	0.66719\\
3.5083	0.66313\\
3.5132	0.65905\\
3.518	0.65495\\
3.5228	0.65083\\
3.5277	0.64669\\
3.5325	0.64254\\
3.5374	0.63838\\
3.5422	0.6342\\
3.5471	0.63\\
3.5519	0.62579\\
3.5567	0.62157\\
3.5616	0.61734\\
3.5664	0.61309\\
3.5713	0.60884\\
3.5761	0.60457\\
3.581	0.6003\\
3.5858	0.59602\\
3.5906	0.59173\\
3.5955	0.58743\\
3.6003	0.58313\\
3.6052	0.57883\\
3.61	0.57451\\
3.6148	0.5702\\
3.6197	0.56588\\
3.6245	0.56156\\
3.6294	0.55724\\
3.6342	0.55291\\
3.6391	0.54859\\
3.6439	0.54426\\
3.6487	0.53994\\
3.6536	0.53562\\
3.6584	0.5313\\
3.6633	0.52699\\
3.6681	0.52267\\
3.673	0.51837\\
3.6778	0.51406\\
3.6826	0.50977\\
3.6875	0.50548\\
3.6923	0.50119\\
3.6972	0.49691\\
3.702	0.49265\\
3.7069	0.48839\\
3.7117	0.48413\\
3.7165	0.47989\\
3.7214	0.47566\\
3.7262	0.47144\\
3.7311	0.46723\\
3.7359	0.46303\\
3.7408	0.45885\\
3.7456	0.45468\\
3.7504	0.45052\\
3.7553	0.44637\\
3.7601	0.44224\\
3.765	0.43812\\
3.7698	0.43402\\
3.7747	0.42993\\
3.7795	0.42586\\
3.7843	0.42181\\
3.7892	0.41777\\
3.794	0.41375\\
3.7989	0.40974\\
3.8037	0.40575\\
3.8086	0.40179\\
3.8134	0.39784\\
3.8182	0.3939\\
3.8231	0.38999\\
3.8279	0.3861\\
3.8328	0.38222\\
3.8376	0.37837\\
3.8425	0.37453\\
3.8473	0.37072\\
3.8521	0.36692\\
3.857	0.36315\\
3.8618	0.3594\\
3.8667	0.35567\\
3.8715	0.35196\\
3.8764	0.34827\\
3.8812	0.34461\\
3.886	0.34096\\
3.8909	0.33734\\
3.8957	0.33374\\
3.9006	0.33016\\
3.9054	0.32661\\
3.9103	0.32308\\
3.9151	0.31957\\
3.9199	0.31609\\
3.9248	0.31262\\
3.9296	0.30919\\
3.9345	0.30577\\
3.9393	0.30238\\
3.9442	0.29901\\
3.949	0.29567\\
3.9538	0.29235\\
3.9587	0.28905\\
3.9635	0.28578\\
3.9684	0.28254\\
3.9732	0.27931\\
3.9781	0.27611\\
3.9829	0.27294\\
3.9877	0.26979\\
3.9926	0.26666\\
3.9974	0.26356\\
4.0023	0.26049\\
4.0071	0.25743\\
4.012	0.25441\\
4.0168	0.2514\\
4.0216	0.24842\\
4.0265	0.24547\\
4.0313	0.24254\\
4.0362	0.23964\\
4.041	0.23676\\
4.0458	0.2339\\
4.0507	0.23107\\
4.0555	0.22826\\
4.0604	0.22548\\
4.0652	0.22272\\
4.0701	0.21999\\
4.0749	0.21728\\
4.0797	0.2146\\
4.0846	0.21194\\
4.0894	0.2093\\
4.0943	0.20669\\
4.0991	0.2041\\
4.104	0.20154\\
4.1088	0.199\\
4.1136	0.19648\\
4.1185	0.19399\\
4.1233	0.19152\\
4.1282	0.18908\\
4.133	0.18666\\
4.1379	0.18426\\
4.1427	0.18189\\
4.1475	0.17954\\
4.1524	0.17721\\
4.1572	0.17491\\
4.1621	0.17263\\
4.1669	0.17037\\
4.1718	0.16814\\
4.1766	0.16593\\
4.1814	0.16374\\
4.1863	0.16157\\
4.1911	0.15943\\
4.196	0.15731\\
4.2008	0.15521\\
4.2057	0.15313\\
4.2105	0.15108\\
4.2153	0.14905\\
4.2202	0.14704\\
4.225	0.14505\\
4.2299	0.14308\\
4.2347	0.14113\\
4.2396	0.13921\\
4.2444	0.13731\\
4.2492	0.13542\\
4.2541	0.13356\\
4.2589	0.13172\\
4.2638	0.1299\\
4.2686	0.1281\\
4.2735	0.12632\\
4.2783	0.12456\\
4.2831	0.12283\\
4.288	0.12111\\
4.2928	0.11941\\
4.2977	0.11773\\
4.3025	0.11607\\
4.3074	0.11443\\
4.3122	0.11281\\
4.317	0.11121\\
4.3219	0.10962\\
4.3267	0.10806\\
4.3316	0.10651\\
4.3364	0.10499\\
4.3413	0.10348\\
4.3461	0.10199\\
4.3509	0.10051\\
4.3558	0.099058\\
4.3606	0.097621\\
4.3655	0.096202\\
4.3703	0.0948\\
4.3752	0.093415\\
4.38	0.092047\\
4.3848	0.090697\\
4.3897	0.089363\\
4.3945	0.088046\\
4.3994	0.086745\\
4.4042	0.085461\\
4.4091	0.084192\\
4.4139	0.08294\\
4.4187	0.081704\\
4.4236	0.080483\\
4.4284	0.079278\\
4.4333	0.078088\\
4.4381	0.076914\\
4.443	0.075755\\
4.4478	0.07461\\
4.4526	0.073481\\
4.4575	0.072366\\
4.4623	0.071265\\
4.4672	0.070179\\
4.472	0.069107\\
4.4768	0.06805\\
4.4817	0.067006\\
4.4865	0.065975\\
4.4914	0.064959\\
4.4962	0.063956\\
4.5011	0.062966\\
4.5059	0.061989\\
4.5107	0.061025\\
4.5156	0.060075\\
4.5204	0.059137\\
4.5253	0.058211\\
4.5301	0.057298\\
4.535	0.056398\\
4.5398	0.055509\\
4.5446	0.054633\\
4.5495	0.053768\\
4.5543	0.052915\\
4.5592	0.052074\\
4.564	0.051245\\
4.5689	0.050427\\
4.5737	0.04962\\
4.5785	0.048824\\
4.5834	0.048039\\
4.5882	0.047265\\
4.5931	0.046502\\
4.5979	0.04575\\
4.6028	0.045008\\
4.6076	0.044276\\
4.6124	0.043555\\
4.6173	0.042844\\
4.6221	0.042143\\
4.627	0.041452\\
4.6318	0.040771\\
4.6367	0.040099\\
4.6415	0.039437\\
4.6463	0.038784\\
4.6512	0.038141\\
4.656	0.037507\\
4.6609	0.036883\\
4.6657	0.036267\\
4.6706	0.03566\\
4.6754	0.035062\\
4.6802	0.034473\\
4.6851	0.033892\\
4.6899	0.03332\\
4.6948	0.032757\\
4.6996	0.032201\\
4.7045	0.031654\\
4.7093	0.031115\\
4.7141	0.030584\\
4.719	0.030061\\
4.7238	0.029546\\
4.7287	0.029039\\
4.7335	0.028539\\
4.7384	0.028047\\
4.7432	0.027562\\
4.748	0.027085\\
4.7529	0.026614\\
4.7577	0.026152\\
4.7626	0.025696\\
4.7674	0.025247\\
4.7723	0.024805\\
4.7771	0.02437\\
4.7819	0.023942\\
4.7868	0.02352\\
4.7916	0.023105\\
4.7965	0.022697\\
4.8013	0.022294\\
4.8062	0.021899\\
4.811	0.021509\\
4.8158	0.021126\\
4.8207	0.020748\\
4.8255	0.020377\\
4.8304	0.020012\\
4.8352	0.019652\\
4.8401	0.019298\\
4.8449	0.01895\\
4.8497	0.018608\\
4.8546	0.018271\\
4.8594	0.017939\\
4.8643	0.017613\\
4.8691	0.017292\\
4.874	0.016977\\
4.8788	0.016666\\
4.8836	0.016361\\
4.8885	0.01606\\
4.8933	0.015765\\
4.8982	0.015475\\
4.903	0.015189\\
4.9078	0.014908\\
4.9127	0.014632\\
4.9175	0.01436\\
4.9224	0.014093\\
4.9272	0.01383\\
4.9321	0.013572\\
4.9369	0.013318\\
4.9417	0.013068\\
4.9466	0.012822\\
4.9514	0.012581\\
4.9563	0.012344\\
4.9611	0.01211\\
4.966	0.011881\\
4.9708	0.011656\\
4.9756	0.011434\\
4.9805	0.011216\\
4.9853	0.011002\\
4.9902	0.010792\\
4.995	0.010585\\
4.9999	0.010382\\
5.0047	0.010182\\
5.0095	0.0099857\\
5.0144	0.0097928\\
5.0192	0.0096033\\
5.0241	0.009417\\
5.0289	0.009234\\
5.0338	0.0090541\\
5.0386	0.0088774\\
5.0434	0.0087038\\
5.0483	0.0085332\\
5.0531	0.0083657\\
5.058	0.008201\\
5.0628	0.0080393\\
5.0677	0.0078804\\
5.0725	0.0077244\\
5.0773	0.0075711\\
5.0822	0.0074205\\
5.087	0.0072727\\
5.0919	0.0071274\\
5.0967	0.0069848\\
5.1016	0.0068447\\
5.1064	0.0067071\\
5.1112	0.006572\\
5.1161	0.0064394\\
5.1209	0.0063091\\
5.1258	0.0061812\\
5.1306	0.0060556\\
5.1355	0.0059323\\
5.1403	0.0058113\\
5.1451	0.0056925\\
5.15	0.0055759\\
5.1548	0.0054614\\
5.1597	0.005349\\
5.1645	0.0052387\\
5.1694	0.0051304\\
5.1742	0.0050242\\
5.179	0.0049199\\
5.1839	0.0048176\\
5.1887	0.0047172\\
5.1936	0.0046187\\
5.1984	0.004522\\
5.2033	0.0044272\\
5.2081	0.0043341\\
5.2129	0.0042429\\
5.2178	0.0041533\\
5.2226	0.0040655\\
5.2275	0.0039794\\
5.2323	0.0038949\\
5.2372	0.003812\\
5.242	0.0037308\\
5.2468	0.0036511\\
5.2517	0.0035729\\
5.2565	0.0034963\\
5.2614	0.0034212\\
5.2662	0.0033476\\
5.2711	0.0032754\\
5.2759	0.0032046\\
5.2807	0.0031352\\
5.2856	0.0030673\\
5.2904	0.0030006\\
5.2953	0.0029353\\
5.3001	0.0028713\\
5.305	0.0028086\\
5.3098	0.0027472\\
5.3146	0.002687\\
5.3195	0.002628\\
5.3243	0.0025702\\
5.3292	0.0025136\\
5.334	0.0024582\\
5.3388	0.0024039\\
5.3437	0.0023507\\
5.3485	0.0022986\\
5.3534	0.0022476\\
5.3582	0.0021976\\
5.3631	0.0021487\\
5.3679	0.0021008\\
5.3727	0.002054\\
5.3776	0.0020081\\
5.3824	0.0019631\\
5.3873	0.0019192\\
5.3921	0.0018761\\
5.397	0.001834\\
5.4018	0.0017928\\
5.4066	0.0017524\\
5.4115	0.0017129\\
5.4115	0.0013688\\
5.4066	0.0014031\\
5.4018	0.0014383\\
5.397	0.0014743\\
5.3921	0.0015111\\
5.3873	0.0015487\\
5.3824	0.0015872\\
5.3776	0.0016266\\
5.3727	0.0016669\\
5.3679	0.0017081\\
5.3631	0.0017502\\
5.3582	0.0017932\\
5.3534	0.0018373\\
5.3485	0.0018823\\
5.3437	0.0019283\\
5.3388	0.0019754\\
5.334	0.0020235\\
5.3292	0.0020726\\
5.3243	0.0021229\\
5.3195	0.0021743\\
5.3146	0.0022268\\
5.3098	0.0022804\\
5.305	0.0023352\\
5.3001	0.0023913\\
5.2953	0.0024485\\
5.2904	0.002507\\
5.2856	0.0025668\\
5.2807	0.0026278\\
5.2759	0.0026902\\
5.2711	0.0027539\\
5.2662	0.002819\\
5.2614	0.0028855\\
5.2565	0.0029534\\
5.2517	0.0030227\\
5.2468	0.0030935\\
5.242	0.0031658\\
5.2372	0.0032397\\
5.2323	0.0033151\\
5.2275	0.003392\\
5.2226	0.0034706\\
5.2178	0.0035509\\
5.2129	0.0036328\\
5.2081	0.0037164\\
5.2033	0.0038017\\
5.1984	0.0038888\\
5.1936	0.0039777\\
5.1887	0.0040685\\
5.1839	0.004161\\
5.179	0.0042555\\
5.1742	0.0043519\\
5.1694	0.0044503\\
5.1645	0.0045506\\
5.1597	0.004653\\
5.1548	0.0047575\\
5.15	0.004864\\
5.1451	0.0049726\\
5.1403	0.0050835\\
5.1355	0.0051965\\
5.1306	0.0053118\\
5.1258	0.0054294\\
5.1209	0.0055492\\
5.1161	0.0056715\\
5.1112	0.0057961\\
5.1064	0.0059231\\
5.1016	0.0060527\\
5.0967	0.0061847\\
5.0919	0.0063193\\
5.087	0.0064565\\
5.0822	0.0065963\\
5.0773	0.0067388\\
5.0725	0.006884\\
5.0677	0.007032\\
5.0628	0.0071828\\
5.058	0.0073364\\
5.0531	0.007493\\
5.0483	0.0076525\\
5.0434	0.0078149\\
5.0386	0.0079804\\
5.0338	0.008149\\
5.0289	0.0083207\\
5.0241	0.0084955\\
5.0192	0.0086736\\
5.0144	0.0088549\\
5.0095	0.0090396\\
5.0047	0.0092276\\
4.9999	0.009419\\
4.995	0.0096139\\
4.9902	0.0098124\\
4.9853	0.010014\\
4.9805	0.01022\\
4.9756	0.010429\\
4.9708	0.010642\\
4.966	0.010859\\
4.9611	0.01108\\
4.9563	0.011304\\
4.9514	0.011533\\
4.9466	0.011765\\
4.9417	0.012001\\
4.9369	0.012242\\
4.9321	0.012487\\
4.9272	0.012736\\
4.9224	0.012989\\
4.9175	0.013247\\
4.9127	0.013509\\
4.9078	0.013775\\
4.903	0.014046\\
4.8982	0.014322\\
4.8933	0.014602\\
4.8885	0.014887\\
4.8836	0.015177\\
4.8788	0.015472\\
4.874	0.015772\\
4.8691	0.016076\\
4.8643	0.016386\\
4.8594	0.016701\\
4.8546	0.017021\\
4.8497	0.017346\\
4.8449	0.017676\\
4.8401	0.018012\\
4.8352	0.018354\\
4.8304	0.018701\\
4.8255	0.019053\\
4.8207	0.019412\\
4.8158	0.019776\\
4.811	0.020146\\
4.8062	0.020522\\
4.8013	0.020904\\
4.7965	0.021292\\
4.7916	0.021686\\
4.7868	0.022087\\
4.7819	0.022494\\
4.7771	0.022907\\
4.7723	0.023327\\
4.7674	0.023754\\
4.7626	0.024187\\
4.7577	0.024627\\
4.7529	0.025074\\
4.748	0.025528\\
4.7432	0.025988\\
4.7384	0.026457\\
4.7335	0.026932\\
4.7287	0.027414\\
4.7238	0.027905\\
4.719	0.028402\\
4.7141	0.028907\\
4.7093	0.02942\\
4.7045	0.029941\\
4.6996	0.03047\\
4.6948	0.031007\\
4.6899	0.031552\\
4.6851	0.032105\\
4.6802	0.032666\\
4.6754	0.033236\\
4.6706	0.033815\\
4.6657	0.034402\\
4.6609	0.034998\\
4.656	0.035603\\
4.6512	0.036216\\
4.6463	0.036839\\
4.6415	0.037471\\
4.6367	0.038113\\
4.6318	0.038764\\
4.627	0.039424\\
4.6221	0.040094\\
4.6173	0.040774\\
4.6124	0.041463\\
4.6076	0.042163\\
4.6028	0.042873\\
4.5979	0.043593\\
4.5931	0.044323\\
4.5882	0.045064\\
4.5834	0.045815\\
4.5785	0.046577\\
4.5737	0.04735\\
4.5689	0.048134\\
4.564	0.048929\\
4.5592	0.049735\\
4.5543	0.050553\\
4.5495	0.051382\\
4.5446	0.052222\\
4.5398	0.053074\\
4.535	0.053938\\
4.5301	0.054814\\
4.5253	0.055702\\
4.5204	0.056602\\
4.5156	0.057515\\
4.5107	0.05844\\
4.5059	0.059377\\
4.5011	0.060328\\
4.4962	0.061291\\
4.4914	0.062267\\
4.4865	0.063256\\
4.4817	0.064258\\
4.4768	0.065274\\
4.472	0.066304\\
4.4672	0.067347\\
4.4623	0.068403\\
4.4575	0.069474\\
4.4526	0.070559\\
4.4478	0.071658\\
4.443	0.072771\\
4.4381	0.073899\\
4.4333	0.075041\\
4.4284	0.076198\\
4.4236	0.07737\\
4.4187	0.078557\\
4.4139	0.079759\\
4.4091	0.080977\\
4.4042	0.08221\\
4.3994	0.083458\\
4.3945	0.084722\\
4.3897	0.086002\\
4.3848	0.087299\\
4.38	0.088611\\
4.3752	0.08994\\
4.3703	0.091285\\
4.3655	0.092647\\
4.3606	0.094025\\
4.3558	0.095421\\
4.3509	0.096833\\
4.3461	0.098263\\
4.3413	0.09971\\
4.3364	0.10117\\
4.3316	0.10266\\
4.3267	0.10416\\
4.3219	0.10567\\
4.317	0.10721\\
4.3122	0.10877\\
4.3074	0.11034\\
4.3025	0.11193\\
4.2977	0.11354\\
4.2928	0.11517\\
4.288	0.11682\\
4.2831	0.11848\\
4.2783	0.12017\\
4.2735	0.12188\\
4.2686	0.1236\\
4.2638	0.12535\\
4.2589	0.12711\\
4.2541	0.1289\\
4.2492	0.1307\\
4.2444	0.13253\\
4.2396	0.13438\\
4.2347	0.13624\\
4.2299	0.13813\\
4.225	0.14004\\
4.2202	0.14197\\
4.2153	0.14392\\
4.2105	0.14589\\
4.2057	0.14789\\
4.2008	0.1499\\
4.196	0.15194\\
4.1911	0.154\\
4.1863	0.15608\\
4.1814	0.15818\\
4.1766	0.16031\\
4.1718	0.16246\\
4.1669	0.16463\\
4.1621	0.16682\\
4.1572	0.16904\\
4.1524	0.17127\\
4.1475	0.17354\\
4.1427	0.17582\\
4.1379	0.17813\\
4.133	0.18046\\
4.1282	0.18281\\
4.1233	0.18519\\
4.1185	0.18759\\
4.1136	0.19002\\
4.1088	0.19247\\
4.104	0.19494\\
4.0991	0.19744\\
4.0943	0.19996\\
4.0894	0.2025\\
4.0846	0.20507\\
4.0797	0.20766\\
4.0749	0.21028\\
4.0701	0.21292\\
4.0652	0.21559\\
4.0604	0.21828\\
4.0555	0.22099\\
4.0507	0.22373\\
4.0458	0.2265\\
4.041	0.22928\\
4.0362	0.2321\\
4.0313	0.23493\\
4.0265	0.23779\\
4.0216	0.24068\\
4.0168	0.24359\\
4.012	0.24653\\
4.0071	0.24949\\
4.0023	0.25247\\
3.9974	0.25548\\
3.9926	0.25851\\
3.9877	0.26157\\
3.9829	0.26465\\
3.9781	0.26776\\
3.9732	0.27089\\
3.9684	0.27405\\
3.9635	0.27723\\
3.9587	0.28043\\
3.9538	0.28366\\
3.949	0.28691\\
3.9442	0.29018\\
3.9393	0.29348\\
3.9345	0.29681\\
3.9296	0.30015\\
3.9248	0.30352\\
3.9199	0.30691\\
3.9151	0.31033\\
3.9103	0.31377\\
3.9054	0.31723\\
3.9006	0.32071\\
3.8957	0.32422\\
3.8909	0.32775\\
3.886	0.3313\\
3.8812	0.33487\\
3.8764	0.33846\\
3.8715	0.34208\\
3.8667	0.34572\\
3.8618	0.34937\\
3.857	0.35305\\
3.8521	0.35675\\
3.8473	0.36047\\
3.8425	0.36421\\
3.8376	0.36797\\
3.8328	0.37175\\
3.8279	0.37554\\
3.8231	0.37936\\
3.8182	0.38319\\
3.8134	0.38705\\
3.8086	0.39092\\
3.8037	0.39481\\
3.7989	0.39871\\
3.794	0.40263\\
3.7892	0.40657\\
3.7843	0.41053\\
3.7795	0.4145\\
3.7747	0.41848\\
3.7698	0.42248\\
3.765	0.4265\\
3.7601	0.43053\\
3.7553	0.43457\\
3.7504	0.43863\\
3.7456	0.4427\\
3.7408	0.44678\\
3.7359	0.45087\\
3.7311	0.45498\\
3.7262	0.45909\\
3.7214	0.46322\\
3.7165	0.46736\\
3.7117	0.47151\\
3.7069	0.47566\\
3.702	0.47983\\
3.6972	0.484\\
3.6923	0.48819\\
3.6875	0.49238\\
3.6826	0.49657\\
3.6778	0.50077\\
3.673	0.50498\\
3.6681	0.50919\\
3.6633	0.51341\\
3.6584	0.51763\\
3.6536	0.52186\\
3.6487	0.52609\\
3.6439	0.53032\\
3.6391	0.53455\\
3.6342	0.53879\\
3.6294	0.54302\\
3.6245	0.54725\\
3.6197	0.55149\\
3.6148	0.55572\\
3.61	0.55995\\
3.6052	0.56418\\
3.6003	0.5684\\
3.5955	0.57262\\
3.5906	0.57684\\
3.5858	0.58105\\
3.581	0.58525\\
3.5761	0.58945\\
3.5713	0.59363\\
3.5664	0.59781\\
3.5616	0.60198\\
3.5567	0.60614\\
3.5519	0.61029\\
3.5471	0.61442\\
3.5422	0.61854\\
3.5374	0.62265\\
3.5325	0.62674\\
3.5277	0.63081\\
3.5228	0.63487\\
3.518	0.6389\\
3.5132	0.64292\\
3.5083	0.64692\\
3.5035	0.65089\\
3.4986	0.65484\\
3.4938	0.65877\\
3.4889	0.66267\\
3.4841	0.66654\\
3.4793	0.67038\\
3.4744	0.67419\\
3.4696	0.67798\\
3.4647	0.68173\\
3.4599	0.68544\\
3.455	0.68912\\
3.4502	0.69277\\
3.4454	0.69637\\
3.4405	0.69994\\
3.4357	0.70346\\
3.4308	0.70695\\
3.426	0.71039\\
3.4211	0.71378\\
3.4163	0.71713\\
3.4115	0.72043\\
3.4066	0.72369\\
3.4018	0.72689\\
3.3969	0.73004\\
3.3921	0.73313\\
3.3872	0.73618\\
3.3824	0.73916\\
3.3776	0.74209\\
3.3727	0.74496\\
3.3679	0.74776\\
3.363	0.75051\\
3.3582	0.75319\\
3.3533	0.7558\\
3.3485	0.75835\\
3.3437	0.76082\\
3.3388	0.76322\\
3.334	0.76555\\
3.3291	0.76781\\
3.3243	0.76998\\
3.3194	0.77208\\
3.3146	0.77409\\
3.3098	0.77602\\
3.3049	0.77787\\
3.3001	0.77963\\
3.2952	0.78129\\
3.2904	0.78287\\
3.2855	0.78435\\
3.2807	0.78574\\
3.2759	0.78703\\
3.271	0.78822\\
3.2662	0.78931\\
3.2613	0.7903\\
3.2565	0.79119\\
3.2516	0.79197\\
3.2468	0.79265\\
3.242	0.79322\\
3.2371	0.79369\\
3.2323	0.79404\\
3.2274	0.79429\\
3.2226	0.79443\\
3.2177	0.79446\\
3.2129	0.79438\\
3.2081	0.79419\\
3.2032	0.79389\\
3.1984	0.79348\\
3.1935	0.79297\\
3.1887	0.79234\\
3.1838	0.7916\\
3.179	0.79076\\
3.1742	0.78981\\
3.1693	0.78875\\
3.1645	0.78759\\
3.1596	0.78632\\
3.1548	0.78495\\
3.15	0.78347\\
3.1451	0.78189\\
3.1403	0.78022\\
3.1354	0.77844\\
3.1306	0.77656\\
3.1257	0.77459\\
3.1209	0.77253\\
3.1161	0.77037\\
3.1112	0.76811\\
3.1064	0.76577\\
3.1015	0.76334\\
3.0967	0.76082\\
3.0918	0.75821\\
3.087	0.75552\\
3.0822	0.75274\\
3.0773	0.74989\\
3.0725	0.74695\\
3.0676	0.74394\\
3.0628	0.74085\\
3.0579	0.73768\\
3.0531	0.73444\\
3.0483	0.73113\\
3.0434	0.72775\\
3.0386	0.7243\\
3.0337	0.72078\\
3.0289	0.7172\\
3.024	0.71356\\
3.0192	0.70985\\
3.0144	0.70608\\
3.0095	0.70225\\
3.0047	0.69837\\
2.9998	0.69443\\
2.995	0.69043\\
2.9901	0.68639\\
2.9853	0.68229\\
2.9805	0.67815\\
2.9756	0.67396\\
2.9708	0.66973\\
2.9659	0.66545\\
2.9611	0.66114\\
2.9562	0.65678\\
2.9514	0.6524\\
2.9466	0.64797\\
2.9417	0.64352\\
2.9369	0.63903\\
2.932	0.63451\\
2.9272	0.62997\\
2.9223	0.62541\\
2.9175	0.62082\\
2.9127	0.61621\\
2.9078	0.61158\\
2.903	0.60693\\
2.8981	0.60227\\
2.8933	0.59759\\
2.8884	0.5929\\
2.8836	0.58819\\
2.8788	0.58348\\
2.8739	0.57875\\
2.8691	0.57401\\
2.8642	0.56927\\
2.8594	0.56452\\
2.8545	0.55976\\
2.8497	0.55499\\
2.8449	0.55022\\
2.84	0.54545\\
2.8352	0.54067\\
2.8303	0.53589\\
2.8255	0.53111\\
2.8206	0.52633\\
2.8158	0.52155\\
2.811	0.51676\\
2.8061	0.51198\\
2.8013	0.5072\\
2.7964	0.50243\\
2.7916	0.49765\\
2.7867	0.49289\\
2.7819	0.48812\\
2.7771	0.48337\\
2.7722	0.47862\\
2.7674	0.47387\\
2.7625	0.46914\\
2.7577	0.46442\\
2.7528	0.4597\\
2.748	0.455\\
2.7432	0.45031\\
2.7383	0.44564\\
2.7335	0.44097\\
2.7286	0.43633\\
2.7238	0.43169\\
2.7189	0.42708\\
2.7141	0.42248\\
2.7093	0.4179\\
2.7044	0.41333\\
2.6996	0.40879\\
2.6947	0.40426\\
2.6899	0.39976\\
2.6851	0.39528\\
2.6802	0.39081\\
2.6754	0.38637\\
2.6705	0.38195\\
2.6657	0.37755\\
2.6608	0.37318\\
2.656	0.36883\\
2.6512	0.3645\\
2.6463	0.3602\\
2.6415	0.35592\\
2.6366	0.35167\\
2.6318	0.34744\\
2.6269	0.34324\\
2.6221	0.33906\\
2.6173	0.3349\\
2.6124	0.33078\\
2.6076	0.32668\\
2.6027	0.3226\\
2.5979	0.31856\\
2.593	0.31454\\
2.5882	0.31055\\
2.5834	0.30658\\
2.5785	0.30265\\
2.5737	0.29874\\
2.5688	0.29486\\
2.564	0.29101\\
2.5591	0.28718\\
2.5543	0.28339\\
2.5495	0.27963\\
2.5446	0.27589\\
2.5398	0.27219\\
2.5349	0.26852\\
2.5301	0.26488\\
2.5252	0.26126\\
2.5204	0.25768\\
2.5156	0.25413\\
2.5107	0.25061\\
2.5059	0.24712\\
2.501	0.24367\\
2.4962	0.24024\\
2.4913	0.23685\\
2.4865	0.23349\\
2.4817	0.23016\\
2.4768	0.22686\\
2.472	0.2236\\
2.4671	0.22036\\
2.4623	0.21716\\
2.4574	0.214\\
2.4526	0.21086\\
2.4478	0.20776\\
2.4429	0.20469\\
2.4381	0.20165\\
2.4332	0.19864\\
2.4284	0.19567\\
2.4235	0.19273\\
2.4187	0.18982\\
2.4139	0.18694\\
2.409	0.1841\\
2.4042	0.18129\\
2.3993	0.17851\\
2.3945	0.17576\\
2.3896	0.17304\\
2.3848	0.17036\\
2.38	0.1677\\
2.3751	0.16508\\
2.3703	0.16249\\
2.3654	0.15993\\
2.3606	0.1574\\
2.3557	0.1549\\
2.3509	0.15243\\
2.3461	0.14999\\
2.3412	0.14758\\
2.3364	0.1452\\
2.3315	0.14285\\
2.3267	0.14053\\
2.3218	0.13824\\
2.317	0.13598\\
2.3122	0.13374\\
2.3073	0.13154\\
2.3025	0.12936\\
2.2976	0.12721\\
2.2928	0.12509\\
2.2879	0.12299\\
2.2831	0.12093\\
2.2783	0.11889\\
2.2734	0.11687\\
2.2686	0.11489\\
2.2637	0.11293\\
2.2589	0.11099\\
2.2541	0.10908\\
2.2492	0.1072\\
2.2444	0.10535\\
2.2395	0.10351\\
2.2347	0.10171\\
2.2298	0.099929\\
2.225	0.098174\\
2.2202	0.096444\\
2.2153	0.094739\\
2.2105	0.093058\\
2.2056	0.091402\\
2.2008	0.08977\\
2.1959	0.088162\\
2.1911	0.086577\\
2.1863	0.085017\\
2.1814	0.08348\\
2.1766	0.081966\\
2.1717	0.080475\\
2.1669	0.079007\\
2.162	0.077562\\
2.1572	0.07614\\
2.1524	0.074739\\
2.1475	0.073361\\
2.1427	0.072004\\
2.1378	0.070669\\
2.133	0.069355\\
2.1281	0.068063\\
2.1233	0.066791\\
2.1185	0.06554\\
2.1136	0.064309\\
2.1088	0.063098\\
2.1039	0.061908\\
2.0991	0.060737\\
2.0942	0.059585\\
2.0894	0.058452\\
2.0846	0.057339\\
2.0797	0.056244\\
2.0749	0.055167\\
2.07	0.054109\\
2.0652	0.053069\\
2.0603	0.052046\\
2.0555	0.051041\\
2.0507	0.050052\\
2.0458	0.049081\\
2.041	0.048127\\
2.0361	0.047189\\
2.0313	0.046267\\
2.0264	0.045362\\
2.0216	0.044472\\
2.0168	0.043597\\
2.0119	0.042738\\
2.0071	0.041894\\
2.0022	0.041065\\
1.9974	0.040251\\
1.9925	0.039451\\
1.9877	0.038665\\
1.9829	0.037894\\
1.978	0.037136\\
1.9732	0.036391\\
1.9683	0.035661\\
1.9635	0.034943\\
1.9586	0.034238\\
1.9538	0.033546\\
1.949	0.032867\\
1.9441	0.0322\\
1.9393	0.031546\\
1.9344	0.030903\\
1.9296	0.030272\\
1.9247	0.029653\\
1.9199	0.029045\\
1.9151	0.028449\\
1.9102	0.027864\\
1.9054	0.027289\\
1.9005	0.026726\\
1.8957	0.026173\\
1.8908	0.02563\\
1.886	0.025098\\
1.8812	0.024575\\
1.8763	0.024063\\
1.8715	0.02356\\
1.8666	0.023067\\
1.8618	0.022584\\
1.8569	0.022109\\
1.8521	0.021644\\
1.8473	0.021188\\
1.8424	0.020741\\
1.8376	0.020302\\
1.8327	0.019872\\
1.8279	0.01945\\
1.8231	0.019036\\
1.8182	0.018631\\
1.8134	0.018233\\
1.8085	0.017843\\
1.8037	0.017461\\
1.7988	0.017087\\
1.794	0.01672\\
1.7892	0.01636\\
1.7843	0.016007\\
1.7795	0.015662\\
1.7746	0.015323\\
1.7698	0.014991\\
1.7649	0.014666\\
1.7601	0.014347\\
1.7553	0.014035\\
1.7504	0.013729\\
1.7456	0.01343\\
1.7407	0.013136\\
1.7359	0.012848\\
1.731	0.012566\\
1.7262	0.01229\\
1.7214	0.01202\\
1.7165	0.011755\\
1.7117	0.011496\\
1.7068	0.011242\\
1.702	0.010993\\
1.6971	0.010749\\
1.6923	0.010511\\
1.6875	0.010277\\
1.6826	0.010048\\
1.6778	0.0098243\\
1.6729	0.0096049\\
1.6681	0.0093902\\
1.6632	0.00918\\
1.6584	0.0089742\\
1.6536	0.0087727\\
1.6487	0.0085754\\
1.6439	0.0083824\\
1.639	0.0081934\\
1.6342	0.0080084\\
1.6293	0.0078274\\
1.6245	0.0076502\\
1.6197	0.0074768\\
1.6148	0.0073071\\
1.61	0.007141\\
1.6051	0.0069785\\
1.6003	0.0068195\\
1.5954	0.0066639\\
1.5906	0.0065116\\
1.5858	0.0063626\\
1.5809	0.0062169\\
1.5761	0.0060743\\
1.5712	0.0059348\\
1.5664	0.0057983\\
1.5615	0.0056648\\
1.5567	0.0055342\\
1.5519	0.0054065\\
1.547	0.0052815\\
1.5422	0.0051593\\
1.5373	0.0050398\\
1.5325	0.0049228\\
1.5276	0.0048085\\
1.5228	0.0046967\\
1.518	0.0045873\\
1.5131	0.0044804\\
1.5083	0.0043758\\
1.5034	0.0042736\\
1.4986	0.0041736\\
1.4937	0.0040758\\
1.4889	0.0039803\\
1.4841	0.0038868\\
1.4792	0.0037954\\
1.4744	0.0037061\\
1.4695	0.0036188\\
1.4647	0.0035334\\
1.4598	0.00345\\
1.455	0.0033684\\
1.4502	0.0032887\\
1.4453	0.0032108\\
1.4405	0.0031346\\
1.4356	0.0030601\\
1.4308	0.0029874\\
1.4259	0.0029163\\
1.4211	0.0028468\\
1.4163	0.0027788\\
1.4114	0.0027125\\
1.4066	0.0026476\\
1.4017	0.0025842\\
1.3969	0.0025223\\
1.3921	0.0024618\\
1.3872	0.0024027\\
1.3824	0.0023449\\
1.3775	0.0022884\\
1.3727	0.0022333\\
1.3678	0.0021794\\
1.363	0.0021268\\
1.3582	0.0020754\\
1.3533	0.0020252\\
1.3485	0.0019761\\
1.3436	0.0019282\\
1.3388	0.0018814\\
1.3339	0.0018357\\
1.3291	0.001791\\
1.3243	0.0017474\\
1.3194	0.0017048\\
1.3146	0.0016632\\
1.3097	0.0016226\\
1.3049	0.0015829\\
1.3	0.0015441\\
1.2952	0.0015063\\
1.2904	0.0014694\\
1.2855	0.0014333\\
1.2807	0.0013981\\
1.2758	0.0013637\\
1.271	0.0013301\\
1.2661	0.0012973\\
1.2613	0.0012653\\
1.2565	0.001234\\
1.2516	0.0012035\\
1.2468	0.0011737\\
1.2419	0.0011446\\
1.2371	0.0011163\\
1.2322	0.0010886\\
1.2274	0.0010615\\
1.2226	0.0010351\\
1.2177	0.0010093\\
1.2129	0.00098418\\
1.208	0.00095963\\
1.2032	0.00093567\\
1.1983	0.00091229\\
1.1935	0.00088947\\
1.1887	0.0008672\\
1.1838	0.00084547\\
1.179	0.00082426\\
1.1741	0.00080356\\
1.1693	0.00078337\\
1.1644	0.00076366\\
1.1596	0.00074444\\
1.1548	0.00072567\\
1.1499	0.00070737\\
1.1451	0.00068951\\
1.1402	0.00067209\\
1.1354	0.00065509\\
1.1305	0.0006385\\
1.1257	0.00062232\\
1.1209	0.00060654\\
1.116	0.00059114\\
1.1112	0.00057612\\
1.1063	0.00056147\\
1.1015	0.00054718\\
1.0966	0.00053324\\
1.0918	0.00051965\\
1.087	0.00050639\\
1.0821	0.00049345\\
1.0773	0.00048084\\
1.0724	0.00046854\\
1.0676	0.00045654\\
1.0627	0.00044484\\
1.0579	0.00043343\\
1.0531	0.0004223\\
1.0482	0.00041145\\
1.0434	0.00040087\\
1.0385	0.00039056\\
1.0337	0.0003805\\
1.0288	0.00037069\\
1.024	0.00036112\\
1.0192	0.0003518\\
1.0143	0.00034271\\
1.0095	0.00033384\\
1.0046	0.0003252\\
0.99979	0.00031678\\
0.99495	0.00030856\\
0.99011	0.00030056\\
0.98526	0.00029275\\
0.98042	0.00028515\\
0.97558	0.00027773\\
0.97074	0.0002705\\
0.96589	0.00026345\\
0.96105	0.00025659\\
0.95621	0.00024989\\
0.95137	0.00024337\\
0.94652	0.00023701\\
0.94168	0.00023081\\
0.93684	0.00022477\\
0.93199	0.00021888\\
0.92715	0.00021315\\
0.92231	0.00020756\\
0.91747	0.00020211\\
0.91262	0.0001968\\
0.90778	0.00019162\\
0.90294	0.00018658\\
0.8981	0.00018167\\
0.89325	0.00017688\\
0.88841	0.00017222\\
0.88357	0.00016768\\
0.87872	0.00016325\\
0.87388	0.00015894\\
0.86904	0.00015473\\
0.8642	0.00015064\\
0.85935	0.00014665\\
0.85451	0.00014276\\
0.84967	0.00013898\\
0.84483	0.00013529\\
0.83998	0.0001317\\
0.83514	0.0001282\\
0.8303	0.00012479\\
0.82545	0.00012147\\
0.82061	0.00011823\\
0.81577	0.00011508\\
0.81093	0.00011201\\
0.80608	0.00010903\\
0.80124	0.00010611\\
0.7964	0.00010328\\
0.79156	0.00010052\\
0.78671	9.7828e-05\\
0.78187	9.5208e-05\\
0.77703	9.2657e-05\\
0.77219	9.0173e-05\\
0.76734	8.7754e-05\\
0.7625	8.5398e-05\\
0.75766	8.3104e-05\\
0.75281	8.087e-05\\
0.74797	7.8695e-05\\
0.74313	7.6576e-05\\
0.73829	7.4514e-05\\
0.73344	7.2505e-05\\
0.7286	7.055e-05\\
0.72376	6.8646e-05\\
0.71892	6.6792e-05\\
0.71407	6.4987e-05\\
0.70923	6.323e-05\\
0.70439	6.1519e-05\\
0.69954	5.9853e-05\\
0.6947	5.8231e-05\\
0.68986	5.6653e-05\\
0.68502	5.5116e-05\\
0.68017	5.3619e-05\\
0.67533	5.2163e-05\\
0.67049	5.0745e-05\\
0.66565	4.9365e-05\\
0.6608	4.8021e-05\\
0.65596	4.6713e-05\\
0.65112	4.544e-05\\
0.64628	4.4201e-05\\
0.64143	4.2995e-05\\
0.63659	4.1821e-05\\
0.63175	4.0679e-05\\
0.6269	3.9567e-05\\
0.62206	3.8484e-05\\
0.61722	3.7431e-05\\
0.61238	3.6405e-05\\
0.60753	3.5407e-05\\
0.60269	3.4436e-05\\
0.59785	3.3491e-05\\
0.59301	3.2571e-05\\
0.58816	3.1676e-05\\
0.58332	3.0805e-05\\
0.57848	2.9958e-05\\
0.57363	2.9133e-05\\
}--cycle;
\addplot[ybar interval, fill=mycolor2, fill opacity=0.35, area legend, draw=none] table[row sep=crcr, x=Lower, y=Count] {%
Lower	Upper	Count\\
0.57363	0.97679	6.6145e-05\\
0.97679	1.3799	0.0014552\\
1.3799	1.7831	0.0080035\\
1.7831	2.1863	0.0465\\
2.1863	2.5894	0.18759\\
2.5894	2.9926	0.50052\\
2.9926	3.3957	0.76999\\
3.3957	3.7989	0.59418\\
3.7989	4.202	0.25248\\
4.202	4.6052	0.092074\\
4.6052	5.0083	0.020505\\
5.0083	5.4115	0.0070775\\
5.4115	5.4115	0.0070775\\
};
\addlegendentry{\ApproxHistogramLegend}

\addplot [color=mycolor2, only marks, every error bar/.append style={opacity=0.45}, mark=*, mark size=0pt, draw=none, forget plot]
 plot [error bars/.cd, y dir=both, y explicit, error bar style={line width=1pt, color=mycolor2}, error mark options={mark=none,mark size=0pt}]
 table[row sep=crcr, y error plus index=2, y error minus index=3]{%
0.77521	6.6145e-05	0.00013125	6.6145e-05\\
1.1784	0.0014552	0.00057739	0.00057739\\
1.5815	0.0080035	0.0013618	0.0013618\\
1.9847	0.0465	0.0036098	0.0036098\\
2.3878	0.18759	0.0063686	0.0063686\\
2.791	0.50052	0.0085242	0.0085242\\
3.1941	0.76999	0.010467	0.010467\\
3.5973	0.59418	0.0098887	0.0098887\\
4.0004	0.25248	0.0057891	0.0057891\\
4.4036	0.092074	0.0032347	0.0032347\\
4.8068	0.020505	0.0016528	0.0016528\\
5.2099	0.0070775	0.00053597	0.00053597\\
};
\addplot [color=mycolor2]
  table[row sep=crcr]{%
0.57363	4.6698e-05\\
1.0143	0.00049696\\
1.208	0.0013273\\
1.3388	0.0025151\\
1.4405	0.0040764\\
1.5228	0.0059696\\
1.5954	0.0082981\\
1.6584	0.010977\\
1.7165	0.014145\\
1.7698	0.017771\\
1.8182	0.021791\\
1.8618	0.026103\\
1.9005	0.030573\\
1.9393	0.035724\\
1.9732	0.040856\\
2.0071	0.046636\\
2.041	0.053129\\
2.0749	0.060403\\
2.1039	0.067315\\
2.133	0.0749\\
2.162	0.083205\\
2.1911	0.092278\\
2.2202	0.10217\\
2.2444	0.11107\\
2.2686	0.12059\\
2.2928	0.13077\\
2.317	0.14163\\
2.3412	0.15318\\
2.3654	0.16546\\
2.3896	0.17848\\
2.4139	0.19225\\
2.4381	0.20679\\
2.4623	0.22211\\
2.4865	0.23821\\
2.5107	0.2551\\
2.5349	0.27277\\
2.5591	0.2912\\
2.5834	0.31039\\
2.6124	0.33438\\
2.6415	0.35937\\
2.6705	0.38529\\
2.6996	0.41204\\
2.7335	0.44418\\
2.7722	0.48186\\
2.8303	0.53951\\
2.903	0.61147\\
2.9369	0.64414\\
2.9659	0.67124\\
2.9901	0.69294\\
3.0095	0.70954\\
3.0289	0.72533\\
3.0483	0.74015\\
3.0628	0.75052\\
3.0773	0.76018\\
3.0918	0.76906\\
3.1064	0.77709\\
3.1209	0.78421\\
3.1354	0.79037\\
3.15	0.79554\\
3.1645	0.7997\\
3.179	0.80283\\
3.1935	0.80496\\
3.2081	0.80608\\
3.2226	0.80624\\
3.2371	0.80546\\
3.2516	0.80379\\
3.2662	0.80127\\
3.2807	0.79792\\
3.2952	0.7938\\
3.3098	0.78892\\
3.3243	0.78333\\
3.3388	0.77706\\
3.3533	0.77012\\
3.3679	0.76255\\
3.3824	0.75437\\
3.4018	0.74258\\
3.4211	0.72984\\
3.4405	0.71622\\
3.4599	0.70181\\
3.4841	0.68281\\
3.5083	0.66285\\
3.5374	0.6379\\
3.5713	0.60774\\
3.6148	0.56795\\
3.7165	0.47471\\
3.7553	0.44026\\
3.7892	0.411\\
3.8231	0.38272\\
3.8521	0.35934\\
3.8812	0.33682\\
3.9103	0.3152\\
3.9393	0.29451\\
3.9684	0.27476\\
3.9974	0.25596\\
4.0265	0.23812\\
4.0555	0.22122\\
4.0846	0.20525\\
4.1136	0.19019\\
4.1427	0.17602\\
4.1718	0.1627\\
4.2008	0.15022\\
4.2299	0.13854\\
4.2589	0.12762\\
4.288	0.11743\\
4.317	0.10794\\
4.3461	0.099105\\
4.38	0.089588\\
4.4139	0.080867\\
4.4478	0.072887\\
4.4817	0.065594\\
4.5156	0.058938\\
4.5543	0.052049\\
4.5931	0.045861\\
4.6367	0.03966\\
4.6802	0.034183\\
4.7238	0.029359\\
4.7723	0.024683\\
4.8255	0.02028\\
4.8788	0.016561\\
4.9369	0.013185\\
4.9999	0.010217\\
5.0677	0.0076932\\
5.1451	0.0055022\\
5.2323	0.0037244\\
5.3388	0.0022715\\
5.4115	0.0016055\\
};
\addlegendentry{\ApproxBLNFitLegend}


\addplot[area legend, draw=none, fill=mycolor2, fill opacity=0.15, forget plot]
table[row sep=crcr] {%
x	y\\
0.57363	6.1595e-05\\
0.57848	6.3171e-05\\
0.58332	6.4788e-05\\
0.58816	6.6445e-05\\
0.59301	6.8145e-05\\
0.59785	6.9887e-05\\
0.60269	7.1673e-05\\
0.60753	7.3505e-05\\
0.61238	7.5383e-05\\
0.61722	7.7308e-05\\
0.62206	7.9282e-05\\
0.6269	8.1306e-05\\
0.63175	8.338e-05\\
0.63659	8.5507e-05\\
0.64143	8.7688e-05\\
0.64628	8.9923e-05\\
0.65112	9.2215e-05\\
0.65596	9.4565e-05\\
0.6608	9.6973e-05\\
0.66565	9.9442e-05\\
0.67049	0.00010197\\
0.67533	0.00010457\\
0.68017	0.00010723\\
0.68502	0.00010995\\
0.68986	0.00011275\\
0.6947	0.00011561\\
0.69954	0.00011855\\
0.70439	0.00012156\\
0.70923	0.00012464\\
0.71407	0.0001278\\
0.71892	0.00013104\\
0.72376	0.00013436\\
0.7286	0.00013777\\
0.73344	0.00014126\\
0.73829	0.00014483\\
0.74313	0.0001485\\
0.74797	0.00015225\\
0.75281	0.0001561\\
0.75766	0.00016005\\
0.7625	0.00016409\\
0.76734	0.00016823\\
0.77219	0.00017248\\
0.77703	0.00017683\\
0.78187	0.00018128\\
0.78671	0.00018585\\
0.79156	0.00019053\\
0.7964	0.00019533\\
0.80124	0.00020024\\
0.80608	0.00020528\\
0.81093	0.00021044\\
0.81577	0.00021572\\
0.82061	0.00022114\\
0.82545	0.00022669\\
0.8303	0.00023237\\
0.83514	0.0002382\\
0.83998	0.00024417\\
0.84483	0.00025028\\
0.84967	0.00025654\\
0.85451	0.00026296\\
0.85935	0.00026953\\
0.8642	0.00027626\\
0.86904	0.00028316\\
0.87388	0.00029023\\
0.87872	0.00029746\\
0.88357	0.00030488\\
0.88841	0.00031247\\
0.89325	0.00032025\\
0.8981	0.00032821\\
0.90294	0.00033637\\
0.90778	0.00034473\\
0.91262	0.00035329\\
0.91747	0.00036205\\
0.92231	0.00037103\\
0.92715	0.00038022\\
0.93199	0.00038963\\
0.93684	0.00039927\\
0.94168	0.00040915\\
0.94652	0.00041926\\
0.95137	0.00042961\\
0.95621	0.00044021\\
0.96105	0.00045106\\
0.96589	0.00046218\\
0.97074	0.00047356\\
0.97558	0.00048521\\
0.98042	0.00049714\\
0.98526	0.00050935\\
0.99011	0.00052186\\
0.99495	0.00053466\\
0.99979	0.00054776\\
1.0046	0.00056118\\
1.0095	0.00057491\\
1.0143	0.00058897\\
1.0192	0.00060337\\
1.024	0.0006181\\
1.0288	0.00063318\\
1.0337	0.00064862\\
1.0385	0.00066441\\
1.0434	0.00068059\\
1.0482	0.00069714\\
1.0531	0.00071408\\
1.0579	0.00073141\\
1.0627	0.00074916\\
1.0676	0.00076731\\
1.0724	0.0007859\\
1.0773	0.00080491\\
1.0821	0.00082437\\
1.087	0.00084428\\
1.0918	0.00086466\\
1.0966	0.00088551\\
1.1015	0.00090684\\
1.1063	0.00092867\\
1.1112	0.000951\\
1.116	0.00097385\\
1.1209	0.00099723\\
1.1257	0.0010211\\
1.1305	0.0010456\\
1.1354	0.0010706\\
1.1402	0.0010962\\
1.1451	0.0011224\\
1.1499	0.0011492\\
1.1548	0.0011766\\
1.1596	0.0012047\\
1.1644	0.0012333\\
1.1693	0.0012627\\
1.1741	0.0012927\\
1.179	0.0013234\\
1.1838	0.0013547\\
1.1887	0.0013868\\
1.1935	0.0014196\\
1.1983	0.0014532\\
1.2032	0.0014875\\
1.208	0.0015226\\
1.2129	0.0015585\\
1.2177	0.0015952\\
1.2226	0.0016327\\
1.2274	0.0016711\\
1.2322	0.0017103\\
1.2371	0.0017504\\
1.2419	0.0017914\\
1.2468	0.0018333\\
1.2516	0.0018761\\
1.2565	0.00192\\
1.2613	0.0019647\\
1.2661	0.0020105\\
1.271	0.0020573\\
1.2758	0.0021051\\
1.2807	0.002154\\
1.2855	0.002204\\
1.2904	0.002255\\
1.2952	0.0023072\\
1.3	0.0023606\\
1.3049	0.0024151\\
1.3097	0.0024708\\
1.3146	0.0025277\\
1.3194	0.0025859\\
1.3243	0.0026454\\
1.3291	0.0027061\\
1.3339	0.0027682\\
1.3388	0.0028316\\
1.3436	0.0028964\\
1.3485	0.0029626\\
1.3533	0.0030302\\
1.3582	0.0030993\\
1.363	0.0031699\\
1.3678	0.003242\\
1.3727	0.0033157\\
1.3775	0.0033909\\
1.3824	0.0034678\\
1.3872	0.0035463\\
1.3921	0.0036265\\
1.3969	0.0037084\\
1.4017	0.0037921\\
1.4066	0.0038775\\
1.4114	0.0039648\\
1.4163	0.0040539\\
1.4211	0.0041449\\
1.4259	0.0042378\\
1.4308	0.0043327\\
1.4356	0.0044296\\
1.4405	0.0045285\\
1.4453	0.0046295\\
1.4502	0.0047326\\
1.455	0.0048379\\
1.4598	0.0049454\\
1.4647	0.0050552\\
1.4695	0.0051672\\
1.4744	0.0052816\\
1.4792	0.0053984\\
1.4841	0.0055175\\
1.4889	0.0056392\\
1.4937	0.0057634\\
1.4986	0.0058901\\
1.5034	0.0060194\\
1.5083	0.0061514\\
1.5131	0.0062862\\
1.518	0.0064237\\
1.5228	0.006564\\
1.5276	0.0067072\\
1.5325	0.0068533\\
1.5373	0.0070024\\
1.5422	0.0071545\\
1.547	0.0073097\\
1.5519	0.0074681\\
1.5567	0.0076296\\
1.5615	0.0077945\\
1.5664	0.0079626\\
1.5712	0.0081342\\
1.5761	0.0083092\\
1.5809	0.0084877\\
1.5858	0.0086698\\
1.5906	0.0088555\\
1.5954	0.0090449\\
1.6003	0.0092382\\
1.6051	0.0094352\\
1.61	0.0096362\\
1.6148	0.0098411\\
1.6197	0.01005\\
1.6245	0.010263\\
1.6293	0.010481\\
1.6342	0.010702\\
1.639	0.010928\\
1.6439	0.011159\\
1.6487	0.011393\\
1.6536	0.011633\\
1.6584	0.011877\\
1.6632	0.012126\\
1.6681	0.01238\\
1.6729	0.012638\\
1.6778	0.012902\\
1.6826	0.013171\\
1.6875	0.013445\\
1.6923	0.013724\\
1.6971	0.014009\\
1.702	0.014299\\
1.7068	0.014594\\
1.7117	0.014895\\
1.7165	0.015202\\
1.7214	0.015515\\
1.7262	0.015834\\
1.731	0.016158\\
1.7359	0.016489\\
1.7407	0.016826\\
1.7456	0.01717\\
1.7504	0.01752\\
1.7553	0.017876\\
1.7601	0.018239\\
1.7649	0.018609\\
1.7698	0.018986\\
1.7746	0.019369\\
1.7795	0.01976\\
1.7843	0.020158\\
1.7892	0.020564\\
1.794	0.020976\\
1.7988	0.021397\\
1.8037	0.021825\\
1.8085	0.022261\\
1.8134	0.022705\\
1.8182	0.023157\\
1.8231	0.023617\\
1.8279	0.024086\\
1.8327	0.024563\\
1.8376	0.025048\\
1.8424	0.025543\\
1.8473	0.026046\\
1.8521	0.026558\\
1.8569	0.02708\\
1.8618	0.02761\\
1.8666	0.028151\\
1.8715	0.0287\\
1.8763	0.02926\\
1.8812	0.029829\\
1.886	0.030409\\
1.8908	0.030998\\
1.8957	0.031598\\
1.9005	0.032208\\
1.9054	0.032829\\
1.9102	0.033461\\
1.9151	0.034104\\
1.9199	0.034758\\
1.9247	0.035423\\
1.9296	0.0361\\
1.9344	0.036788\\
1.9393	0.037488\\
1.9441	0.0382\\
1.949	0.038924\\
1.9538	0.039661\\
1.9586	0.04041\\
1.9635	0.041171\\
1.9683	0.041946\\
1.9732	0.042733\\
1.978	0.043534\\
1.9829	0.044348\\
1.9877	0.045176\\
1.9925	0.046017\\
1.9974	0.046873\\
2.0022	0.047742\\
2.0071	0.048626\\
2.0119	0.049524\\
2.0168	0.050437\\
2.0216	0.051365\\
2.0264	0.052308\\
2.0313	0.053267\\
2.0361	0.054241\\
2.041	0.05523\\
2.0458	0.056236\\
2.0507	0.057258\\
2.0555	0.058296\\
2.0603	0.05935\\
2.0652	0.060422\\
2.07	0.06151\\
2.0749	0.062616\\
2.0797	0.063739\\
2.0846	0.064879\\
2.0894	0.066038\\
2.0942	0.067214\\
2.0991	0.068409\\
2.1039	0.069622\\
2.1088	0.070854\\
2.1136	0.072105\\
2.1185	0.073375\\
2.1233	0.074665\\
2.1281	0.075974\\
2.133	0.077303\\
2.1378	0.078652\\
2.1427	0.080022\\
2.1475	0.081411\\
2.1524	0.082822\\
2.1572	0.084254\\
2.162	0.085707\\
2.1669	0.087181\\
2.1717	0.088677\\
2.1766	0.090195\\
2.1814	0.091736\\
2.1863	0.093298\\
2.1911	0.094884\\
2.1959	0.096492\\
2.2008	0.098123\\
2.2056	0.099778\\
2.2105	0.10146\\
2.2153	0.10316\\
2.2202	0.10488\\
2.225	0.10663\\
2.2298	0.10841\\
2.2347	0.11021\\
2.2395	0.11203\\
2.2444	0.11388\\
2.2492	0.11576\\
2.2541	0.11766\\
2.2589	0.11959\\
2.2637	0.12154\\
2.2686	0.12352\\
2.2734	0.12552\\
2.2783	0.12756\\
2.2831	0.12961\\
2.2879	0.1317\\
2.2928	0.13381\\
2.2976	0.13596\\
2.3025	0.13813\\
2.3073	0.14032\\
2.3122	0.14255\\
2.317	0.1448\\
2.3218	0.14709\\
2.3267	0.1494\\
2.3315	0.15174\\
2.3364	0.15411\\
2.3412	0.15651\\
2.3461	0.15893\\
2.3509	0.16139\\
2.3557	0.16388\\
2.3606	0.1664\\
2.3654	0.16895\\
2.3703	0.17152\\
2.3751	0.17413\\
2.38	0.17677\\
2.3848	0.17944\\
2.3896	0.18214\\
2.3945	0.18488\\
2.3993	0.18764\\
2.4042	0.19043\\
2.409	0.19326\\
2.4139	0.19612\\
2.4187	0.199\\
2.4235	0.20192\\
2.4284	0.20488\\
2.4332	0.20786\\
2.4381	0.21088\\
2.4429	0.21392\\
2.4478	0.217\\
2.4526	0.22011\\
2.4574	0.22326\\
2.4623	0.22643\\
2.4671	0.22964\\
2.472	0.23288\\
2.4768	0.23615\\
2.4817	0.23945\\
2.4865	0.24278\\
2.4913	0.24615\\
2.4962	0.24955\\
2.501	0.25298\\
2.5059	0.25644\\
2.5107	0.25994\\
2.5156	0.26346\\
2.5204	0.26702\\
2.5252	0.27061\\
2.5301	0.27423\\
2.5349	0.27788\\
2.5398	0.28156\\
2.5446	0.28527\\
2.5495	0.28902\\
2.5543	0.29279\\
2.5591	0.2966\\
2.564	0.30043\\
2.5688	0.3043\\
2.5737	0.30819\\
2.5785	0.31212\\
2.5834	0.31607\\
2.5882	0.32006\\
2.593	0.32407\\
2.5979	0.32811\\
2.6027	0.33218\\
2.6076	0.33628\\
2.6124	0.3404\\
2.6173	0.34456\\
2.6221	0.34874\\
2.6269	0.35294\\
2.6318	0.35718\\
2.6366	0.36143\\
2.6415	0.36572\\
2.6463	0.37003\\
2.6512	0.37436\\
2.656	0.37872\\
2.6608	0.38311\\
2.6657	0.38751\\
2.6705	0.39194\\
2.6754	0.3964\\
2.6802	0.40087\\
2.6851	0.40537\\
2.6899	0.40988\\
2.6947	0.41442\\
2.6996	0.41898\\
2.7044	0.42355\\
2.7093	0.42815\\
2.7141	0.43276\\
2.7189	0.43739\\
2.7238	0.44204\\
2.7286	0.4467\\
2.7335	0.45138\\
2.7383	0.45608\\
2.7432	0.46079\\
2.748	0.46551\\
2.7528	0.47024\\
2.7577	0.47499\\
2.7625	0.47975\\
2.7674	0.48452\\
2.7722	0.48929\\
2.7771	0.49408\\
2.7819	0.49888\\
2.7867	0.50368\\
2.7916	0.50849\\
2.7964	0.5133\\
2.8013	0.51812\\
2.8061	0.52295\\
2.811	0.52778\\
2.8158	0.53261\\
2.8206	0.53744\\
2.8255	0.54227\\
2.8303	0.5471\\
2.8352	0.55193\\
2.84	0.55675\\
2.8449	0.56158\\
2.8497	0.5664\\
2.8545	0.57121\\
2.8594	0.57602\\
2.8642	0.58082\\
2.8691	0.58561\\
2.8739	0.5904\\
2.8788	0.59517\\
2.8836	0.59993\\
2.8884	0.60468\\
2.8933	0.60942\\
2.8981	0.61415\\
2.903	0.61886\\
2.9078	0.62355\\
2.9127	0.62823\\
2.9175	0.63289\\
2.9223	0.63753\\
2.9272	0.64216\\
2.932	0.64676\\
2.9369	0.65134\\
2.9417	0.65591\\
2.9466	0.66045\\
2.9514	0.66496\\
2.9562	0.66945\\
2.9611	0.67392\\
2.9659	0.67836\\
2.9708	0.68278\\
2.9756	0.68717\\
2.9805	0.69152\\
2.9853	0.69585\\
2.9901	0.70015\\
2.995	0.70441\\
2.9998	0.70864\\
3.0047	0.71283\\
3.0095	0.71699\\
3.0144	0.7211\\
3.0192	0.72517\\
3.024	0.72919\\
3.0289	0.73317\\
3.0337	0.73709\\
3.0386	0.74096\\
3.0434	0.74477\\
3.0483	0.74852\\
3.0531	0.7522\\
3.0579	0.75582\\
3.0628	0.75936\\
3.0676	0.76282\\
3.0725	0.76621\\
3.0773	0.76951\\
3.0822	0.77272\\
3.087	0.77584\\
3.0918	0.77886\\
3.0967	0.78179\\
3.1015	0.78461\\
3.1064	0.78732\\
3.1112	0.78993\\
3.1161	0.79242\\
3.1209	0.7948\\
3.1257	0.79706\\
3.1306	0.7992\\
3.1354	0.80122\\
3.1403	0.80312\\
3.1451	0.80489\\
3.15	0.80654\\
3.1548	0.80806\\
3.1596	0.80946\\
3.1645	0.81073\\
3.1693	0.81188\\
3.1742	0.8129\\
3.179	0.8138\\
3.1838	0.81458\\
3.1887	0.81524\\
3.1935	0.81578\\
3.1984	0.8162\\
3.2032	0.81651\\
3.2081	0.8167\\
3.2129	0.81678\\
3.2177	0.81676\\
3.2226	0.81663\\
3.2274	0.8164\\
3.2323	0.81607\\
3.2371	0.81563\\
3.242	0.8151\\
3.2468	0.81448\\
3.2516	0.81376\\
3.2565	0.81295\\
3.2613	0.81205\\
3.2662	0.81106\\
3.271	0.80999\\
3.2759	0.80883\\
3.2807	0.80758\\
3.2855	0.80625\\
3.2904	0.80483\\
3.2952	0.80333\\
3.3001	0.80175\\
3.3049	0.80009\\
3.3098	0.79835\\
3.3146	0.79653\\
3.3194	0.79463\\
3.3243	0.79265\\
3.3291	0.7906\\
3.334	0.78847\\
3.3388	0.78626\\
3.3437	0.78398\\
3.3485	0.78162\\
3.3533	0.77919\\
3.3582	0.77669\\
3.363	0.77412\\
3.3679	0.77148\\
3.3727	0.76877\\
3.3776	0.76599\\
3.3824	0.76314\\
3.3872	0.76023\\
3.3921	0.75726\\
3.3969	0.75422\\
3.4018	0.75112\\
3.4066	0.74796\\
3.4115	0.74474\\
3.4163	0.74146\\
3.4211	0.73812\\
3.426	0.73473\\
3.4308	0.73129\\
3.4357	0.72779\\
3.4405	0.72424\\
3.4454	0.72065\\
3.4502	0.717\\
3.455	0.71331\\
3.4599	0.70957\\
3.4647	0.70579\\
3.4696	0.70196\\
3.4744	0.69809\\
3.4793	0.69419\\
3.4841	0.69024\\
3.4889	0.68626\\
3.4938	0.68224\\
3.4986	0.67819\\
3.5035	0.67411\\
3.5083	0.66999\\
3.5132	0.66584\\
3.518	0.66167\\
3.5228	0.65747\\
3.5277	0.65324\\
3.5325	0.64899\\
3.5374	0.64471\\
3.5422	0.64041\\
3.5471	0.63609\\
3.5519	0.63175\\
3.5567	0.6274\\
3.5616	0.62302\\
3.5664	0.61863\\
3.5713	0.61423\\
3.5761	0.60981\\
3.581	0.60538\\
3.5858	0.60093\\
3.5906	0.59648\\
3.5955	0.59202\\
3.6003	0.58755\\
3.6052	0.58307\\
3.61	0.57859\\
3.6148	0.5741\\
3.6197	0.5696\\
3.6245	0.56511\\
3.6294	0.56061\\
3.6342	0.55611\\
3.6391	0.55161\\
3.6439	0.54711\\
3.6487	0.54261\\
3.6536	0.53811\\
3.6584	0.53362\\
3.6633	0.52913\\
3.6681	0.52465\\
3.673	0.52017\\
3.6778	0.51569\\
3.6826	0.51123\\
3.6875	0.50677\\
3.6923	0.50232\\
3.6972	0.49788\\
3.702	0.49344\\
3.7069	0.48902\\
3.7117	0.48461\\
3.7165	0.48021\\
3.7214	0.47583\\
3.7262	0.47145\\
3.7311	0.46709\\
3.7359	0.46275\\
3.7408	0.45841\\
3.7456	0.4541\\
3.7504	0.4498\\
3.7553	0.44551\\
3.7601	0.44124\\
3.765	0.43699\\
3.7698	0.43276\\
3.7747	0.42854\\
3.7795	0.42434\\
3.7843	0.42017\\
3.7892	0.41601\\
3.794	0.41187\\
3.7989	0.40775\\
3.8037	0.40365\\
3.8086	0.39957\\
3.8134	0.39552\\
3.8182	0.39148\\
3.8231	0.38747\\
3.8279	0.38348\\
3.8328	0.37951\\
3.8376	0.37557\\
3.8425	0.37164\\
3.8473	0.36774\\
3.8521	0.36387\\
3.857	0.36002\\
3.8618	0.35619\\
3.8667	0.35239\\
3.8715	0.34861\\
3.8764	0.34485\\
3.8812	0.34112\\
3.886	0.33742\\
3.8909	0.33374\\
3.8957	0.33009\\
3.9006	0.32646\\
3.9054	0.32286\\
3.9103	0.31928\\
3.9151	0.31573\\
3.9199	0.31221\\
3.9248	0.30871\\
3.9296	0.30523\\
3.9345	0.30179\\
3.9393	0.29837\\
3.9442	0.29498\\
3.949	0.29161\\
3.9538	0.28827\\
3.9587	0.28496\\
3.9635	0.28167\\
3.9684	0.27841\\
3.9732	0.27518\\
3.9781	0.27197\\
3.9829	0.26879\\
3.9877	0.26564\\
3.9926	0.26251\\
3.9974	0.25942\\
4.0023	0.25634\\
4.0071	0.2533\\
4.012	0.25028\\
4.0168	0.24729\\
4.0216	0.24432\\
4.0265	0.24138\\
4.0313	0.23847\\
4.0362	0.23558\\
4.041	0.23272\\
4.0458	0.22989\\
4.0507	0.22708\\
4.0555	0.2243\\
4.0604	0.22155\\
4.0652	0.21882\\
4.0701	0.21612\\
4.0749	0.21344\\
4.0797	0.21079\\
4.0846	0.20817\\
4.0894	0.20557\\
4.0943	0.203\\
4.0991	0.20045\\
4.104	0.19792\\
4.1088	0.19543\\
4.1136	0.19295\\
4.1185	0.19051\\
4.1233	0.18808\\
4.1282	0.18568\\
4.133	0.18331\\
4.1379	0.18096\\
4.1427	0.17864\\
4.1475	0.17633\\
4.1524	0.17406\\
4.1572	0.1718\\
4.1621	0.16958\\
4.1669	0.16737\\
4.1718	0.16519\\
4.1766	0.16303\\
4.1814	0.16089\\
4.1863	0.15878\\
4.1911	0.15669\\
4.196	0.15462\\
4.2008	0.15258\\
4.2057	0.15055\\
4.2105	0.14855\\
4.2153	0.14657\\
4.2202	0.14462\\
4.225	0.14268\\
4.2299	0.14077\\
4.2347	0.13888\\
4.2396	0.13701\\
4.2444	0.13516\\
4.2492	0.13333\\
4.2541	0.13152\\
4.2589	0.12973\\
4.2638	0.12797\\
4.2686	0.12622\\
4.2735	0.12449\\
4.2783	0.12279\\
4.2831	0.1211\\
4.288	0.11943\\
4.2928	0.11779\\
4.2977	0.11616\\
4.3025	0.11455\\
4.3074	0.11296\\
4.3122	0.11139\\
4.317	0.10983\\
4.3219	0.1083\\
4.3267	0.10678\\
4.3316	0.10528\\
4.3364	0.1038\\
4.3413	0.10234\\
4.3461	0.1009\\
4.3509	0.099467\\
4.3558	0.098057\\
4.3606	0.096664\\
4.3655	0.095288\\
4.3703	0.093929\\
4.3752	0.092586\\
4.38	0.09126\\
4.3848	0.08995\\
4.3897	0.088656\\
4.3945	0.087379\\
4.3994	0.086117\\
4.4042	0.084871\\
4.4091	0.083641\\
4.4139	0.082426\\
4.4187	0.081226\\
4.4236	0.080042\\
4.4284	0.078872\\
4.4333	0.077717\\
4.4381	0.076577\\
4.443	0.075452\\
4.4478	0.07434\\
4.4526	0.073243\\
4.4575	0.07216\\
4.4623	0.071091\\
4.4672	0.070036\\
4.472	0.068995\\
4.4768	0.067967\\
4.4817	0.066952\\
4.4865	0.06595\\
4.4914	0.064962\\
4.4962	0.063986\\
4.5011	0.063023\\
4.5059	0.062073\\
4.5107	0.061136\\
4.5156	0.06021\\
4.5204	0.059297\\
4.5253	0.058396\\
4.5301	0.057507\\
4.535	0.05663\\
4.5398	0.055764\\
4.5446	0.05491\\
4.5495	0.054067\\
4.5543	0.053236\\
4.5592	0.052415\\
4.564	0.051606\\
4.5689	0.050807\\
4.5737	0.050019\\
4.5785	0.049242\\
4.5834	0.048475\\
4.5882	0.047718\\
4.5931	0.046972\\
4.5979	0.046236\\
4.6028	0.045509\\
4.6076	0.044792\\
4.6124	0.044085\\
4.6173	0.043388\\
4.6221	0.0427\\
4.627	0.042021\\
4.6318	0.041351\\
4.6367	0.040691\\
4.6415	0.040039\\
4.6463	0.039397\\
4.6512	0.038763\\
4.656	0.038137\\
4.6609	0.03752\\
4.6657	0.036912\\
4.6706	0.036312\\
4.6754	0.03572\\
4.6802	0.035136\\
4.6851	0.03456\\
4.6899	0.033992\\
4.6948	0.033432\\
4.6996	0.032879\\
4.7045	0.032334\\
4.7093	0.031797\\
4.7141	0.031267\\
4.719	0.030744\\
4.7238	0.030229\\
4.7287	0.029721\\
4.7335	0.029219\\
4.7384	0.028725\\
4.7432	0.028238\\
4.748	0.027758\\
4.7529	0.027284\\
4.7577	0.026817\\
4.7626	0.026357\\
4.7674	0.025903\\
4.7723	0.025456\\
4.7771	0.025016\\
4.7819	0.024581\\
4.7868	0.024153\\
4.7916	0.023731\\
4.7965	0.023315\\
4.8013	0.022906\\
4.8062	0.022502\\
4.811	0.022104\\
4.8158	0.021713\\
4.8207	0.021327\\
4.8255	0.020947\\
4.8304	0.020572\\
4.8352	0.020203\\
4.8401	0.01984\\
4.8449	0.019483\\
4.8497	0.019131\\
4.8546	0.018784\\
4.8594	0.018443\\
4.8643	0.018107\\
4.8691	0.017777\\
4.874	0.017451\\
4.8788	0.017131\\
4.8836	0.016816\\
4.8885	0.016506\\
4.8933	0.016201\\
4.8982	0.015901\\
4.903	0.015606\\
4.9078	0.015315\\
4.9127	0.01503\\
4.9175	0.014749\\
4.9224	0.014472\\
4.9272	0.014201\\
4.9321	0.013934\\
4.9369	0.013671\\
4.9417	0.013413\\
4.9466	0.013159\\
4.9514	0.01291\\
4.9563	0.012665\\
4.9611	0.012424\\
4.966	0.012187\\
4.9708	0.011954\\
4.9756	0.011725\\
4.9805	0.0115\\
4.9853	0.011279\\
4.9902	0.011062\\
4.995	0.010849\\
4.9999	0.01064\\
5.0047	0.010434\\
5.0095	0.010232\\
5.0144	0.010033\\
5.0192	0.0098378\\
5.0241	0.0096461\\
5.0289	0.0094578\\
5.0338	0.0092729\\
5.0386	0.0090912\\
5.0434	0.0089128\\
5.0483	0.0087376\\
5.0531	0.0085655\\
5.058	0.0083965\\
5.0628	0.0082305\\
5.0677	0.0080675\\
5.0725	0.0079074\\
5.0773	0.0077502\\
5.0822	0.0075959\\
5.087	0.0074443\\
5.0919	0.0072955\\
5.0967	0.0071494\\
5.1016	0.0070059\\
5.1064	0.006865\\
5.1112	0.0067267\\
5.1161	0.006591\\
5.1209	0.0064577\\
5.1258	0.0063268\\
5.1306	0.0061984\\
5.1355	0.0060723\\
5.1403	0.0059486\\
5.1451	0.0058271\\
5.15	0.0057079\\
5.1548	0.0055909\\
5.1597	0.005476\\
5.1645	0.0053633\\
5.1694	0.0052527\\
5.1742	0.0051442\\
5.179	0.0050377\\
5.1839	0.0049332\\
5.1887	0.0048307\\
5.1936	0.0047301\\
5.1984	0.0046314\\
5.2033	0.0045346\\
5.2081	0.0044396\\
5.2129	0.0043464\\
5.2178	0.0042551\\
5.2226	0.0041654\\
5.2275	0.0040775\\
5.2323	0.0039913\\
5.2372	0.0039068\\
5.242	0.0038239\\
5.2468	0.0037426\\
5.2517	0.0036629\\
5.2565	0.0035847\\
5.2614	0.0035081\\
5.2662	0.003433\\
5.2711	0.0033594\\
5.2759	0.0032872\\
5.2807	0.0032165\\
5.2856	0.0031471\\
5.2904	0.0030792\\
5.2953	0.0030126\\
5.3001	0.0029473\\
5.305	0.0028834\\
5.3098	0.0028207\\
5.3146	0.0027593\\
5.3195	0.0026992\\
5.3243	0.0026403\\
5.3292	0.0025825\\
5.334	0.002526\\
5.3388	0.0024706\\
5.3437	0.0024164\\
5.3485	0.0023632\\
5.3534	0.0023112\\
5.3582	0.0022602\\
5.3631	0.0022103\\
5.3679	0.0021615\\
5.3727	0.0021136\\
5.3776	0.0020668\\
5.3824	0.0020209\\
5.3873	0.001976\\
5.3921	0.0019321\\
5.397	0.0018891\\
5.4018	0.001847\\
5.4066	0.0018057\\
5.4115	0.0017654\\
5.4115	0.0014456\\
5.4066	0.0014812\\
5.4018	0.0015175\\
5.397	0.0015547\\
5.3921	0.0015927\\
5.3873	0.0016316\\
5.3824	0.0016714\\
5.3776	0.0017121\\
5.3727	0.0017537\\
5.3679	0.0017963\\
5.3631	0.0018398\\
5.3582	0.0018843\\
5.3534	0.0019297\\
5.3485	0.0019762\\
5.3437	0.0020238\\
5.3388	0.0020724\\
5.334	0.0021221\\
5.3292	0.0021728\\
5.3243	0.0022247\\
5.3195	0.0022778\\
5.3146	0.002332\\
5.3098	0.0023875\\
5.305	0.0024441\\
5.3001	0.002502\\
5.2953	0.0025611\\
5.2904	0.0026216\\
5.2856	0.0026833\\
5.2807	0.0027464\\
5.2759	0.0028109\\
5.2711	0.0028768\\
5.2662	0.0029441\\
5.2614	0.0030128\\
5.2565	0.003083\\
5.2517	0.0031548\\
5.2468	0.003228\\
5.242	0.0033029\\
5.2372	0.0033793\\
5.2323	0.0034574\\
5.2275	0.0035371\\
5.2226	0.0036185\\
5.2178	0.0037016\\
5.2129	0.0037865\\
5.2081	0.0038732\\
5.2033	0.0039616\\
5.1984	0.0040519\\
5.1936	0.0041441\\
5.1887	0.0042383\\
5.1839	0.0043343\\
5.179	0.0044324\\
5.1742	0.0045325\\
5.1694	0.0046346\\
5.1645	0.0047388\\
5.1597	0.0048451\\
5.1548	0.0049537\\
5.15	0.0050644\\
5.1451	0.0051773\\
5.1403	0.0052925\\
5.1355	0.00541\\
5.1306	0.0055299\\
5.1258	0.0056521\\
5.1209	0.0057768\\
5.1161	0.0059039\\
5.1112	0.0060336\\
5.1064	0.0061657\\
5.1016	0.0063005\\
5.0967	0.0064378\\
5.0919	0.0065778\\
5.087	0.0067205\\
5.0822	0.0068659\\
5.0773	0.0070141\\
5.0725	0.0071651\\
5.0677	0.007319\\
5.0628	0.0074758\\
5.058	0.0076354\\
5.0531	0.0077981\\
5.0483	0.0079637\\
5.0434	0.0081324\\
5.0386	0.0083042\\
5.0338	0.0084791\\
5.0289	0.0086572\\
5.0241	0.0088385\\
5.0192	0.0090231\\
5.0144	0.0092109\\
5.0095	0.0094021\\
5.0047	0.0095967\\
4.9999	0.0097946\\
4.995	0.0099961\\
4.9902	0.010201\\
4.9853	0.01041\\
4.9805	0.010622\\
4.9756	0.010837\\
4.9708	0.011057\\
4.966	0.01128\\
4.9611	0.011507\\
4.9563	0.011738\\
4.9514	0.011972\\
4.9466	0.012211\\
4.9417	0.012453\\
4.9369	0.0127\\
4.9321	0.01295\\
4.9272	0.013205\\
4.9224	0.013464\\
4.9175	0.013727\\
4.9127	0.013995\\
4.9078	0.014266\\
4.903	0.014543\\
4.8982	0.014823\\
4.8933	0.015109\\
4.8885	0.015398\\
4.8836	0.015693\\
4.8788	0.015992\\
4.874	0.016296\\
4.8691	0.016605\\
4.8643	0.016918\\
4.8594	0.017237\\
4.8546	0.017561\\
4.8497	0.01789\\
4.8449	0.018224\\
4.8401	0.018563\\
4.8352	0.018908\\
4.8304	0.019258\\
4.8255	0.019614\\
4.8207	0.019975\\
4.8158	0.020342\\
4.811	0.020714\\
4.8062	0.021093\\
4.8013	0.021477\\
4.7965	0.021867\\
4.7916	0.022263\\
4.7868	0.022665\\
4.7819	0.023074\\
4.7771	0.023488\\
4.7723	0.02391\\
4.7674	0.024337\\
4.7626	0.024771\\
4.7577	0.025212\\
4.7529	0.025659\\
4.748	0.026113\\
4.7432	0.026574\\
4.7384	0.027042\\
4.7335	0.027517\\
4.7287	0.027999\\
4.7238	0.028488\\
4.719	0.028985\\
4.7141	0.029489\\
4.7093	0.03\\
4.7045	0.030519\\
4.6996	0.031046\\
4.6948	0.03158\\
4.6899	0.032122\\
4.6851	0.032672\\
4.6802	0.03323\\
4.6754	0.033797\\
4.6706	0.034371\\
4.6657	0.034954\\
4.6609	0.035544\\
4.656	0.036144\\
4.6512	0.036752\\
4.6463	0.037369\\
4.6415	0.037994\\
4.6367	0.038628\\
4.6318	0.039271\\
4.627	0.039924\\
4.6221	0.040585\\
4.6173	0.041256\\
4.6124	0.041936\\
4.6076	0.042625\\
4.6028	0.043324\\
4.5979	0.044033\\
4.5931	0.044751\\
4.5882	0.045479\\
4.5834	0.046218\\
4.5785	0.046966\\
4.5737	0.047724\\
4.5689	0.048493\\
4.564	0.049273\\
4.5592	0.050062\\
4.5543	0.050863\\
4.5495	0.051674\\
4.5446	0.052496\\
4.5398	0.05333\\
4.535	0.054174\\
4.5301	0.05503\\
4.5253	0.055897\\
4.5204	0.056775\\
4.5156	0.057666\\
4.5107	0.058568\\
4.5059	0.059482\\
4.5011	0.060408\\
4.4962	0.061346\\
4.4914	0.062297\\
4.4865	0.06326\\
4.4817	0.064236\\
4.4768	0.065224\\
4.472	0.066226\\
4.4672	0.067241\\
4.4623	0.068268\\
4.4575	0.06931\\
4.4526	0.070364\\
4.4478	0.071433\\
4.443	0.072515\\
4.4381	0.073612\\
4.4333	0.074722\\
4.4284	0.075847\\
4.4236	0.076986\\
4.4187	0.07814\\
4.4139	0.079309\\
4.4091	0.080492\\
4.4042	0.081691\\
4.3994	0.082905\\
4.3945	0.084134\\
4.3897	0.08538\\
4.3848	0.08664\\
4.38	0.087917\\
4.3752	0.08921\\
4.3703	0.090519\\
4.3655	0.091845\\
4.3606	0.093187\\
4.3558	0.094546\\
4.3509	0.095922\\
4.3461	0.097315\\
4.3413	0.098725\\
4.3364	0.10015\\
4.3316	0.1016\\
4.3267	0.10306\\
4.3219	0.10454\\
4.317	0.10604\\
4.3122	0.10756\\
4.3074	0.10909\\
4.3025	0.11065\\
4.2977	0.11222\\
4.2928	0.11381\\
4.288	0.11542\\
4.2831	0.11705\\
4.2783	0.1187\\
4.2735	0.12037\\
4.2686	0.12206\\
4.2638	0.12377\\
4.2589	0.1255\\
4.2541	0.12725\\
4.2492	0.12902\\
4.2444	0.13081\\
4.2396	0.13262\\
4.2347	0.13445\\
4.2299	0.1363\\
4.225	0.13817\\
4.2202	0.14007\\
4.2153	0.14198\\
4.2105	0.14392\\
4.2057	0.14588\\
4.2008	0.14786\\
4.196	0.14987\\
4.1911	0.15189\\
4.1863	0.15394\\
4.1814	0.15601\\
4.1766	0.1581\\
4.1718	0.16022\\
4.1669	0.16236\\
4.1621	0.16452\\
4.1572	0.1667\\
4.1524	0.16891\\
4.1475	0.17114\\
4.1427	0.1734\\
4.1379	0.17567\\
4.133	0.17798\\
4.1282	0.1803\\
4.1233	0.18265\\
4.1185	0.18503\\
4.1136	0.18743\\
4.1088	0.18985\\
4.104	0.1923\\
4.0991	0.19477\\
4.0943	0.19726\\
4.0894	0.19979\\
4.0846	0.20233\\
4.0797	0.2049\\
4.0749	0.2075\\
4.0701	0.21012\\
4.0652	0.21277\\
4.0604	0.21544\\
4.0555	0.21814\\
4.0507	0.22086\\
4.0458	0.22361\\
4.041	0.22638\\
4.0362	0.22918\\
4.0313	0.23201\\
4.0265	0.23486\\
4.0216	0.23774\\
4.0168	0.24064\\
4.012	0.24357\\
4.0071	0.24653\\
4.0023	0.24951\\
3.9974	0.25251\\
3.9926	0.25555\\
3.9877	0.25861\\
3.9829	0.26169\\
3.9781	0.26481\\
3.9732	0.26794\\
3.9684	0.27111\\
3.9635	0.2743\\
3.9587	0.27752\\
3.9538	0.28076\\
3.949	0.28403\\
3.9442	0.28732\\
3.9393	0.29065\\
3.9345	0.29399\\
3.9296	0.29737\\
3.9248	0.30077\\
3.9199	0.30419\\
3.9151	0.30764\\
3.9103	0.31112\\
3.9054	0.31462\\
3.9006	0.31815\\
3.8957	0.3217\\
3.8909	0.32528\\
3.886	0.32888\\
3.8812	0.33251\\
3.8764	0.33617\\
3.8715	0.33985\\
3.8667	0.34355\\
3.8618	0.34728\\
3.857	0.35103\\
3.8521	0.35481\\
3.8473	0.35861\\
3.8425	0.36243\\
3.8376	0.36628\\
3.8328	0.37015\\
3.8279	0.37404\\
3.8231	0.37796\\
3.8182	0.3819\\
3.8134	0.38586\\
3.8086	0.38985\\
3.8037	0.39385\\
3.7989	0.39788\\
3.794	0.40193\\
3.7892	0.40599\\
3.7843	0.41008\\
3.7795	0.41419\\
3.7747	0.41832\\
3.7698	0.42247\\
3.765	0.42663\\
3.7601	0.43082\\
3.7553	0.43502\\
3.7504	0.43924\\
3.7456	0.44347\\
3.7408	0.44772\\
3.7359	0.45199\\
3.7311	0.45627\\
3.7262	0.46057\\
3.7214	0.46488\\
3.7165	0.46921\\
3.7117	0.47354\\
3.7069	0.47789\\
3.702	0.48225\\
3.6972	0.48662\\
3.6923	0.49101\\
3.6875	0.4954\\
3.6826	0.4998\\
3.6778	0.5042\\
3.673	0.50862\\
3.6681	0.51304\\
3.6633	0.51746\\
3.6584	0.52189\\
3.6536	0.52632\\
3.6487	0.53076\\
3.6439	0.5352\\
3.6391	0.53963\\
3.6342	0.54407\\
3.6294	0.54851\\
3.6245	0.55294\\
3.6197	0.55737\\
3.6148	0.5618\\
3.61	0.56622\\
3.6052	0.57063\\
3.6003	0.57504\\
3.5955	0.57943\\
3.5906	0.58382\\
3.5858	0.5882\\
3.581	0.59256\\
3.5761	0.59691\\
3.5713	0.60125\\
3.5664	0.60557\\
3.5616	0.60987\\
3.5567	0.61416\\
3.5519	0.61842\\
3.5471	0.62267\\
3.5422	0.62689\\
3.5374	0.63109\\
3.5325	0.63526\\
3.5277	0.63941\\
3.5228	0.64353\\
3.518	0.64763\\
3.5132	0.65169\\
3.5083	0.65572\\
3.5035	0.65972\\
3.4986	0.66369\\
3.4938	0.66762\\
3.4889	0.67151\\
3.4841	0.67537\\
3.4793	0.67919\\
3.4744	0.68297\\
3.4696	0.68671\\
3.4647	0.6904\\
3.4599	0.69406\\
3.455	0.69766\\
3.4502	0.70123\\
3.4454	0.70474\\
3.4405	0.70821\\
3.4357	0.71162\\
3.4308	0.71499\\
3.426	0.7183\\
3.4211	0.72156\\
3.4163	0.72477\\
3.4115	0.72792\\
3.4066	0.73101\\
3.4018	0.73405\\
3.3969	0.73703\\
3.3921	0.73995\\
3.3872	0.74281\\
3.3824	0.74561\\
3.3776	0.74834\\
3.3727	0.75101\\
3.3679	0.75362\\
3.363	0.75616\\
3.3582	0.75864\\
3.3533	0.76104\\
3.3485	0.76338\\
3.3437	0.76565\\
3.3388	0.76785\\
3.334	0.76998\\
3.3291	0.77203\\
3.3243	0.77401\\
3.3194	0.77592\\
3.3146	0.77775\\
3.3098	0.7795\\
3.3049	0.78117\\
3.3001	0.78275\\
3.2952	0.78426\\
3.2904	0.78568\\
3.2855	0.78702\\
3.2807	0.78827\\
3.2759	0.78943\\
3.271	0.79049\\
3.2662	0.79147\\
3.2613	0.79235\\
3.2565	0.79314\\
3.2516	0.79382\\
3.2468	0.79441\\
3.242	0.7949\\
3.2371	0.79529\\
3.2323	0.79558\\
3.2274	0.79576\\
3.2226	0.79585\\
3.2177	0.79582\\
3.2129	0.79569\\
3.2081	0.79546\\
3.2032	0.79512\\
3.1984	0.79468\\
3.1935	0.79414\\
3.1887	0.79348\\
3.1838	0.79273\\
3.179	0.79187\\
3.1742	0.7909\\
3.1693	0.78983\\
3.1645	0.78866\\
3.1596	0.78739\\
3.1548	0.78602\\
3.15	0.78454\\
3.1451	0.78297\\
3.1403	0.78129\\
3.1354	0.77952\\
3.1306	0.77765\\
3.1257	0.77568\\
3.1209	0.77362\\
3.1161	0.77146\\
3.1112	0.7692\\
3.1064	0.76686\\
3.1015	0.76441\\
3.0967	0.76188\\
3.0918	0.75926\\
3.087	0.75655\\
3.0822	0.75375\\
3.0773	0.75086\\
3.0725	0.74789\\
3.0676	0.74483\\
3.0628	0.74169\\
3.0579	0.73846\\
3.0531	0.73516\\
3.0483	0.73178\\
3.0434	0.72832\\
3.0386	0.72478\\
3.0337	0.72117\\
3.0289	0.71749\\
3.024	0.71374\\
3.0192	0.70992\\
3.0144	0.70604\\
3.0095	0.7021\\
3.0047	0.69809\\
2.9998	0.69402\\
2.995	0.6899\\
2.9901	0.68572\\
2.9853	0.68149\\
2.9805	0.67722\\
2.9756	0.67289\\
2.9708	0.66853\\
2.9659	0.66412\\
2.9611	0.65967\\
2.9562	0.65519\\
2.9514	0.65067\\
2.9466	0.64612\\
2.9417	0.64155\\
2.9369	0.63694\\
2.932	0.63231\\
2.9272	0.62765\\
2.9223	0.62298\\
2.9175	0.61828\\
2.9127	0.61357\\
2.9078	0.60883\\
2.903	0.60409\\
2.8981	0.59933\\
2.8933	0.59455\\
2.8884	0.58977\\
2.8836	0.58498\\
2.8788	0.58018\\
2.8739	0.57537\\
2.8691	0.57055\\
2.8642	0.56573\\
2.8594	0.56091\\
2.8545	0.55608\\
2.8497	0.55125\\
2.8449	0.54642\\
2.84	0.54159\\
2.8352	0.53676\\
2.8303	0.53193\\
2.8255	0.5271\\
2.8206	0.52228\\
2.8158	0.51746\\
2.811	0.51265\\
2.8061	0.50784\\
2.8013	0.50304\\
2.7964	0.49824\\
2.7916	0.49346\\
2.7867	0.48868\\
2.7819	0.48392\\
2.7771	0.47916\\
2.7722	0.47442\\
2.7674	0.46969\\
2.7625	0.46497\\
2.7577	0.46026\\
2.7528	0.45557\\
2.748	0.4509\\
2.7432	0.44624\\
2.7383	0.44159\\
2.7335	0.43697\\
2.7286	0.43236\\
2.7238	0.42777\\
2.7189	0.4232\\
2.7141	0.41864\\
2.7093	0.41411\\
2.7044	0.4096\\
2.6996	0.40511\\
2.6947	0.40064\\
2.6899	0.39619\\
2.6851	0.39176\\
2.6802	0.38736\\
2.6754	0.38298\\
2.6705	0.37863\\
2.6657	0.3743\\
2.6608	0.36999\\
2.656	0.36571\\
2.6512	0.36145\\
2.6463	0.35722\\
2.6415	0.35302\\
2.6366	0.34884\\
2.6318	0.34469\\
2.6269	0.34057\\
2.6221	0.33647\\
2.6173	0.3324\\
2.6124	0.32836\\
2.6076	0.32435\\
2.6027	0.32036\\
2.5979	0.31641\\
2.593	0.31248\\
2.5882	0.30858\\
2.5834	0.30471\\
2.5785	0.30087\\
2.5737	0.29706\\
2.5688	0.29328\\
2.564	0.28953\\
2.5591	0.28581\\
2.5543	0.28212\\
2.5495	0.27846\\
2.5446	0.27482\\
2.5398	0.27122\\
2.5349	0.26765\\
2.5301	0.26411\\
2.5252	0.26061\\
2.5204	0.25713\\
2.5156	0.25368\\
2.5107	0.25026\\
2.5059	0.24688\\
2.501	0.24352\\
2.4962	0.2402\\
2.4913	0.2369\\
2.4865	0.23364\\
2.4817	0.23041\\
2.4768	0.22721\\
2.472	0.22404\\
2.4671	0.2209\\
2.4623	0.21779\\
2.4574	0.21471\\
2.4526	0.21166\\
2.4478	0.20865\\
2.4429	0.20566\\
2.4381	0.2027\\
2.4332	0.19978\\
2.4284	0.19688\\
2.4235	0.19402\\
2.4187	0.19118\\
2.4139	0.18838\\
2.409	0.18561\\
2.4042	0.18286\\
2.3993	0.18015\\
2.3945	0.17746\\
2.3896	0.17481\\
2.3848	0.17218\\
2.38	0.16959\\
2.3751	0.16702\\
2.3703	0.16448\\
2.3654	0.16197\\
2.3606	0.15949\\
2.3557	0.15704\\
2.3509	0.15462\\
2.3461	0.15223\\
2.3412	0.14986\\
2.3364	0.14752\\
2.3315	0.14521\\
2.3267	0.14293\\
2.3218	0.14068\\
2.317	0.13845\\
2.3122	0.13625\\
2.3073	0.13408\\
2.3025	0.13194\\
2.2976	0.12982\\
2.2928	0.12773\\
2.2879	0.12566\\
2.2831	0.12363\\
2.2783	0.12162\\
2.2734	0.11963\\
2.2686	0.11767\\
2.2637	0.11574\\
2.2589	0.11383\\
2.2541	0.11194\\
2.2492	0.11009\\
2.2444	0.10825\\
2.2395	0.10644\\
2.2347	0.10466\\
2.2298	0.1029\\
2.225	0.10116\\
2.2202	0.099453\\
2.2153	0.097765\\
2.2105	0.0961\\
2.2056	0.094459\\
2.2008	0.092841\\
2.1959	0.091246\\
2.1911	0.089673\\
2.1863	0.088123\\
2.1814	0.086595\\
2.1766	0.08509\\
2.1717	0.083606\\
2.1669	0.082143\\
2.162	0.080703\\
2.1572	0.079283\\
2.1524	0.077885\\
2.1475	0.076507\\
2.1427	0.07515\\
2.1378	0.073813\\
2.133	0.072496\\
2.1281	0.071199\\
2.1233	0.069922\\
2.1185	0.068665\\
2.1136	0.067427\\
2.1088	0.066208\\
2.1039	0.065007\\
2.0991	0.063826\\
2.0942	0.062663\\
2.0894	0.061518\\
2.0846	0.060391\\
2.0797	0.059282\\
2.0749	0.058191\\
2.07	0.057117\\
2.0652	0.05606\\
2.0603	0.05502\\
2.0555	0.053997\\
2.0507	0.052991\\
2.0458	0.052001\\
2.041	0.051027\\
2.0361	0.050069\\
2.0313	0.049127\\
2.0264	0.0482\\
2.0216	0.047289\\
2.0168	0.046393\\
2.0119	0.045512\\
2.0071	0.044646\\
2.0022	0.043794\\
1.9974	0.042956\\
1.9925	0.042133\\
1.9877	0.041324\\
1.9829	0.040529\\
1.978	0.039747\\
1.9732	0.038978\\
1.9683	0.038223\\
1.9635	0.037481\\
1.9586	0.036751\\
1.9538	0.036035\\
1.949	0.035331\\
1.9441	0.034639\\
1.9393	0.033959\\
1.9344	0.033291\\
1.9296	0.032636\\
1.9247	0.031991\\
1.9199	0.031358\\
1.9151	0.030737\\
1.9102	0.030126\\
1.9054	0.029527\\
1.9005	0.028938\\
1.8957	0.02836\\
1.8908	0.027792\\
1.886	0.027235\\
1.8812	0.026687\\
1.8763	0.02615\\
1.8715	0.025623\\
1.8666	0.025105\\
1.8618	0.024596\\
1.8569	0.024097\\
1.8521	0.023608\\
1.8473	0.023127\\
1.8424	0.022655\\
1.8376	0.022192\\
1.8327	0.021738\\
1.8279	0.021292\\
1.8231	0.020855\\
1.8182	0.020425\\
1.8134	0.020004\\
1.8085	0.019591\\
1.8037	0.019186\\
1.7988	0.018788\\
1.794	0.018398\\
1.7892	0.018015\\
1.7843	0.01764\\
1.7795	0.017272\\
1.7746	0.016911\\
1.7698	0.016557\\
1.7649	0.016209\\
1.7601	0.015869\\
1.7553	0.015535\\
1.7504	0.015207\\
1.7456	0.014886\\
1.7407	0.014571\\
1.7359	0.014263\\
1.731	0.01396\\
1.7262	0.013663\\
1.7214	0.013372\\
1.7165	0.013087\\
1.7117	0.012808\\
1.7068	0.012534\\
1.702	0.012265\\
1.6971	0.012002\\
1.6923	0.011744\\
1.6875	0.011491\\
1.6826	0.011244\\
1.6778	0.011001\\
1.6729	0.010763\\
1.6681	0.01053\\
1.6632	0.010301\\
1.6584	0.010078\\
1.6536	0.0098584\\
1.6487	0.0096436\\
1.6439	0.0094331\\
1.639	0.0092269\\
1.6342	0.0090249\\
1.6293	0.0088271\\
1.6245	0.0086333\\
1.6197	0.0084435\\
1.6148	0.0082575\\
1.61	0.0080754\\
1.6051	0.0078971\\
1.6003	0.0077224\\
1.5954	0.0075514\\
1.5906	0.0073838\\
1.5858	0.0072198\\
1.5809	0.0070592\\
1.5761	0.0069019\\
1.5712	0.0067479\\
1.5664	0.0065971\\
1.5615	0.0064495\\
1.5567	0.006305\\
1.5519	0.0061635\\
1.547	0.006025\\
1.5422	0.0058894\\
1.5373	0.0057567\\
1.5325	0.0056268\\
1.5276	0.0054996\\
1.5228	0.0053752\\
1.518	0.0052533\\
1.5131	0.0051341\\
1.5083	0.0050174\\
1.5034	0.0049032\\
1.4986	0.0047915\\
1.4937	0.0046821\\
1.4889	0.0045751\\
1.4841	0.0044704\\
1.4792	0.004368\\
1.4744	0.0042677\\
1.4695	0.0041696\\
1.4647	0.0040737\\
1.4598	0.0039798\\
1.455	0.003888\\
1.4502	0.0037981\\
1.4453	0.0037102\\
1.4405	0.0036243\\
1.4356	0.0035402\\
1.4308	0.0034579\\
1.4259	0.0033775\\
1.4211	0.0032988\\
1.4163	0.0032218\\
1.4114	0.0031466\\
1.4066	0.003073\\
1.4017	0.003001\\
1.3969	0.0029306\\
1.3921	0.0028618\\
1.3872	0.0027944\\
1.3824	0.0027286\\
1.3775	0.0026643\\
1.3727	0.0026014\\
1.3678	0.0025399\\
1.363	0.0024797\\
1.3582	0.0024209\\
1.3533	0.0023634\\
1.3485	0.0023073\\
1.3436	0.0022523\\
1.3388	0.0021986\\
1.3339	0.0021462\\
1.3291	0.0020949\\
1.3243	0.0020447\\
1.3194	0.0019957\\
1.3146	0.0019478\\
1.3097	0.001901\\
1.3049	0.0018553\\
1.3	0.0018106\\
1.2952	0.0017669\\
1.2904	0.0017242\\
1.2855	0.0016825\\
1.2807	0.0016418\\
1.2758	0.001602\\
1.271	0.0015631\\
1.2661	0.0015251\\
1.2613	0.0014879\\
1.2565	0.0014516\\
1.2516	0.0014162\\
1.2468	0.0013816\\
1.2419	0.0013478\\
1.2371	0.0013147\\
1.2322	0.0012825\\
1.2274	0.001251\\
1.2226	0.0012202\\
1.2177	0.0011901\\
1.2129	0.0011608\\
1.208	0.0011321\\
1.2032	0.0011041\\
1.1983	0.0010767\\
1.1935	0.00105\\
1.1887	0.001024\\
1.1838	0.00099849\\
1.179	0.00097363\\
1.1741	0.00094936\\
1.1693	0.00092566\\
1.1644	0.00090252\\
1.1596	0.00087993\\
1.1548	0.00085788\\
1.1499	0.00083635\\
1.1451	0.00081534\\
1.1402	0.00079482\\
1.1354	0.0007748\\
1.1305	0.00075526\\
1.1257	0.00073618\\
1.1209	0.00071756\\
1.116	0.00069939\\
1.1112	0.00068166\\
1.1063	0.00066435\\
1.1015	0.00064746\\
1.0966	0.00063098\\
1.0918	0.0006149\\
1.087	0.00059921\\
1.0821	0.0005839\\
1.0773	0.00056896\\
1.0724	0.00055438\\
1.0676	0.00054016\\
1.0627	0.00052629\\
1.0579	0.00051275\\
1.0531	0.00049955\\
1.0482	0.00048667\\
1.0434	0.0004741\\
1.0385	0.00046184\\
1.0337	0.00044989\\
1.0288	0.00043823\\
1.024	0.00042686\\
1.0192	0.00041577\\
1.0143	0.00040495\\
1.0095	0.0003944\\
1.0046	0.00038411\\
0.99979	0.00037408\\
0.99495	0.0003643\\
0.99011	0.00035476\\
0.98526	0.00034546\\
0.98042	0.00033639\\
0.97558	0.00032755\\
0.97074	0.00031893\\
0.96589	0.00031053\\
0.96105	0.00030233\\
0.95621	0.00029435\\
0.95137	0.00028656\\
0.94652	0.00027897\\
0.94168	0.00027157\\
0.93684	0.00026436\\
0.93199	0.00025734\\
0.92715	0.00025049\\
0.92231	0.00024381\\
0.91747	0.0002373\\
0.91262	0.00023096\\
0.90778	0.00022478\\
0.90294	0.00021876\\
0.8981	0.0002129\\
0.89325	0.00020718\\
0.88841	0.00020161\\
0.88357	0.00019618\\
0.87872	0.00019089\\
0.87388	0.00018574\\
0.86904	0.00018072\\
0.8642	0.00017583\\
0.85935	0.00017107\\
0.85451	0.00016643\\
0.84967	0.00016191\\
0.84483	0.00015751\\
0.83998	0.00015322\\
0.83514	0.00014904\\
0.8303	0.00014498\\
0.82545	0.00014101\\
0.82061	0.00013716\\
0.81577	0.0001334\\
0.81093	0.00012974\\
0.80608	0.00012618\\
0.80124	0.00012271\\
0.7964	0.00011933\\
0.79156	0.00011604\\
0.78671	0.00011284\\
0.78187	0.00010972\\
0.77703	0.00010668\\
0.77219	0.00010373\\
0.76734	0.00010085\\
0.7625	9.8052e-05\\
0.75766	9.5326e-05\\
0.75281	9.2673e-05\\
0.74797	9.009e-05\\
0.74313	8.7576e-05\\
0.73829	8.5129e-05\\
0.73344	8.2748e-05\\
0.7286	8.043e-05\\
0.72376	7.8175e-05\\
0.71892	7.598e-05\\
0.71407	7.3844e-05\\
0.70923	7.1766e-05\\
0.70439	6.9743e-05\\
0.69954	6.7776e-05\\
0.6947	6.5861e-05\\
0.68986	6.3998e-05\\
0.68502	6.2186e-05\\
0.68017	6.0423e-05\\
0.67533	5.8708e-05\\
0.67049	5.7039e-05\\
0.66565	5.5416e-05\\
0.6608	5.3837e-05\\
0.65596	5.2302e-05\\
0.65112	5.0808e-05\\
0.64628	4.9355e-05\\
0.64143	4.7942e-05\\
0.63659	4.6568e-05\\
0.63175	4.5231e-05\\
0.6269	4.3931e-05\\
0.62206	4.2667e-05\\
0.61722	4.1438e-05\\
0.61238	4.0243e-05\\
0.60753	3.9081e-05\\
0.60269	3.7951e-05\\
0.59785	3.6853e-05\\
0.59301	3.5785e-05\\
0.58816	3.4746e-05\\
0.58332	3.3737e-05\\
0.57848	3.2755e-05\\
0.57363	3.1801e-05\\
}--cycle;
      % \input{../../../../.tex/20/KS/SizeDistributionFittingsfig37.tex}
      \legend{}
  
      \nextgroupplot[%
        xmin=0.57363,xmax=5.3544,
        % xmin=0.3346,xmax=5.5934,
        % ymin=0,ymax=1,
      ] % This file was created by matlab2tikz.
%
\definecolor{mycolor1}{rgb}{0.00000,0.44706,0.69804}%
%
\addplot[ybar interval, fill=mycolor1, fill opacity=0.35, area legend, draw=none] table[row sep=crcr, x=Lower, y=Count] {%
Lower	Upper	Count\\
0.57363	0.94138	7.2513e-05\\
0.94138	1.3091	0.0013052\\
1.3091	1.6769	0.0045683\\
1.6769	2.0446	0.02567\\
2.0446	2.4124	0.11124\\
2.4124	2.7801	0.33226\\
2.7801	3.1479	0.65936\\
3.1479	3.5156	0.76472\\
3.5156	3.8834	0.4883\\
3.8834	4.2511	0.22117\\
4.2511	4.6189	0.082955\\
4.6189	4.9866	0.019796\\
4.9866	5.3544	0.0078314\\
5.3544	5.3544	0.0078314\\
};
\addlegendentry{\ExactHistogramLegend}

\addplot [color=mycolor1, only marks, every error bar/.append style={opacity=0.45}, mark=*, mark size=0pt, draw=none, forget plot]
 plot [error bars/.cd, y dir=both, y explicit, error bar style={line width=1pt, color=mycolor1}, error mark options={mark=none,mark size=0pt}]
 table[row sep=crcr, y error plus index=2, y error minus index=3]{%
0.75751	7.2513e-05	0.00014388	7.2513e-05\\
1.1253	0.0013052	0.00059201	0.00059201\\
1.493	0.0045683	0.0011323	0.0011323\\
1.8608	0.02567	0.0025413	0.0025413\\
2.2285	0.11124	0.0050518	0.0050518\\
2.5963	0.33226	0.0093187	0.0093187\\
2.964	0.65936	0.01047	0.01047\\
3.3318	0.76472	0.011784	0.011784\\
3.6995	0.4883	0.0091027	0.0091027\\
4.0673	0.22117	0.0062349	0.0062349\\
4.435	0.082955	0.0032613	0.0032613\\
4.8028	0.019796	0.0016721	0.0016721\\
5.1705	0.0078314	0.00048739	0.00048739\\
};
\addplot [color=mycolor1]
  table[row sep=crcr]{%
0.57363	5.2469e-05\\
1.057	0.00055473\\
1.258	0.0014842\\
1.392	0.0028325\\
1.4925	0.0045614\\
1.5738	0.0066645\\
1.6456	0.00926\\
1.7078	0.012256\\
1.7652	0.015807\\
1.8131	0.019475\\
1.8609	0.023916\\
1.904	0.028687\\
1.9423	0.033641\\
1.9806	0.039355\\
2.0141	0.045057\\
2.0476	0.051486\\
2.0811	0.05872\\
2.1098	0.065623\\
2.1385	0.07323\\
2.1672	0.081599\\
2.1959	0.090783\\
2.2246	0.10084\\
2.2534	0.11181\\
2.2821	0.12375\\
2.306	0.13446\\
2.3299	0.14589\\
2.3539	0.15804\\
2.3778	0.17093\\
2.4017	0.18457\\
2.4256	0.19897\\
2.4496	0.21414\\
2.4735	0.2301\\
2.4974	0.24686\\
2.5213	0.26442\\
2.5453	0.28278\\
2.574	0.30588\\
2.6027	0.33007\\
2.6314	0.35526\\
2.6649	0.38575\\
2.7032	0.42174\\
2.751	0.46787\\
2.8372	0.55224\\
2.8994	0.61274\\
2.9377	0.64895\\
2.9664	0.67517\\
2.9903	0.69614\\
3.0143	0.71609\\
3.0334	0.73115\\
3.0525	0.74527\\
3.0717	0.75831\\
3.086	0.76728\\
3.1004	0.77551\\
3.1147	0.78293\\
3.1291	0.78947\\
3.1435	0.79511\\
3.1578	0.79978\\
3.1722	0.80345\\
3.1865	0.80611\\
3.2009	0.80773\\
3.2152	0.80831\\
3.2296	0.80787\\
3.244	0.80643\\
3.2583	0.804\\
3.2727	0.80064\\
3.287	0.79639\\
3.3014	0.79129\\
3.3157	0.7854\\
3.3301	0.77878\\
3.3445	0.77148\\
3.3636	0.76077\\
3.3827	0.74907\\
3.4019	0.7365\\
3.4258	0.71971\\
3.4497	0.7019\\
3.4784	0.67942\\
3.5072	0.65598\\
3.5407	0.62769\\
3.5837	0.59034\\
3.7512	0.44391\\
3.7895	0.4119\\
3.8278	0.38093\\
3.8613	0.3548\\
3.8948	0.32967\\
3.9283	0.30561\\
3.957	0.28587\\
3.9857	0.26698\\
4.0144	0.24894\\
4.0431	0.23177\\
4.0719	0.21546\\
4.1006	0.20001\\
4.1293	0.1854\\
4.158	0.17162\\
4.1867	0.15864\\
4.2154	0.14645\\
4.2441	0.13502\\
4.2728	0.12432\\
4.3063	0.11272\\
4.3398	0.10202\\
4.3733	0.092185\\
4.4068	0.083153\\
4.4403	0.074877\\
4.4738	0.067309\\
4.5121	0.059465\\
4.5504	0.052415\\
4.5887	0.046093\\
4.6318	0.039775\\
4.6748	0.034218\\
4.7227	0.028844\\
4.7705	0.024219\\
4.8232	0.019891\\
4.8806	0.015956\\
4.9428	0.012481\\
5.0098	0.009502\\
5.0864	0.0068835\\
5.1725	0.0047231\\
5.273	0.0029871\\
5.3544	0.0020327\\
};
\addlegendentry{\ExactBLNFitLegend}


\addplot[area legend, draw=none, fill=mycolor1, fill opacity=0.15, forget plot]
table[row sep=crcr] {%
x	y\\
0.57363	7.5805e-05\\
0.57842	7.7426e-05\\
0.58321	7.9082e-05\\
0.58799	8.0774e-05\\
0.59278	8.2504e-05\\
0.59756	8.4271e-05\\
0.60235	8.6077e-05\\
0.60713	8.7923e-05\\
0.61192	8.9809e-05\\
0.6167	9.1737e-05\\
0.62149	9.3707e-05\\
0.62628	9.572e-05\\
0.63106	9.7778e-05\\
0.63585	9.9881e-05\\
0.64063	0.00010203\\
0.64542	0.00010423\\
0.6502	0.00010647\\
0.65499	0.00010877\\
0.65977	0.00011111\\
0.66456	0.00011351\\
0.66935	0.00011596\\
0.67413	0.00011846\\
0.67892	0.00012102\\
0.6837	0.00012364\\
0.68849	0.00012631\\
0.69327	0.00012905\\
0.69806	0.00013184\\
0.70284	0.0001347\\
0.70763	0.00013762\\
0.71241	0.0001406\\
0.7172	0.00014366\\
0.72199	0.00014677\\
0.72677	0.00014996\\
0.73156	0.00015322\\
0.73634	0.00015656\\
0.74113	0.00015996\\
0.74591	0.00016345\\
0.7507	0.00016701\\
0.75548	0.00017065\\
0.76027	0.00017437\\
0.76506	0.00017817\\
0.76984	0.00018206\\
0.77463	0.00018604\\
0.77941	0.00019011\\
0.7842	0.00019426\\
0.78898	0.00019851\\
0.79377	0.00020286\\
0.79855	0.0002073\\
0.80334	0.00021185\\
0.80813	0.00021649\\
0.81291	0.00022124\\
0.8177	0.0002261\\
0.82248	0.00023106\\
0.82727	0.00023614\\
0.83205	0.00024133\\
0.83684	0.00024664\\
0.84162	0.00025207\\
0.84641	0.00025762\\
0.8512	0.0002633\\
0.85598	0.0002691\\
0.86077	0.00027503\\
0.86555	0.0002811\\
0.87034	0.00028731\\
0.87512	0.00029365\\
0.87991	0.00030014\\
0.88469	0.00030678\\
0.88948	0.00031356\\
0.89426	0.0003205\\
0.89905	0.0003276\\
0.90384	0.00033485\\
0.90862	0.00034228\\
0.91341	0.00034986\\
0.91819	0.00035762\\
0.92298	0.00036556\\
0.92776	0.00037368\\
0.93255	0.00038198\\
0.93733	0.00039046\\
0.94212	0.00039914\\
0.94691	0.00040802\\
0.95169	0.0004171\\
0.95648	0.00042638\\
0.96126	0.00043587\\
0.96605	0.00044558\\
0.97083	0.00045551\\
0.97562	0.00046567\\
0.9804	0.00047605\\
0.98519	0.00048667\\
0.98998	0.00049753\\
0.99476	0.00050864\\
0.99955	0.00052\\
1.0043	0.00053162\\
1.0091	0.0005435\\
1.0139	0.00055565\\
1.0187	0.00056808\\
1.0235	0.00058078\\
1.0283	0.00059378\\
1.033	0.00060707\\
1.0378	0.00062066\\
1.0426	0.00063456\\
1.0474	0.00064878\\
1.0522	0.00066332\\
1.057	0.00067818\\
1.0618	0.00069339\\
1.0665	0.00070894\\
1.0713	0.00072484\\
1.0761	0.0007411\\
1.0809	0.00075773\\
1.0857	0.00077474\\
1.0905	0.00079213\\
1.0953	0.00080992\\
1.1	0.00082811\\
1.1048	0.00084671\\
1.1096	0.00086573\\
1.1144	0.00088518\\
1.1192	0.00090507\\
1.124	0.00092541\\
1.1288	0.00094621\\
1.1335	0.00096748\\
1.1383	0.00098923\\
1.1431	0.0010115\\
1.1479	0.0010342\\
1.1527	0.0010575\\
1.1575	0.0010813\\
1.1623	0.0011056\\
1.167	0.0011304\\
1.1718	0.0011559\\
1.1766	0.0011819\\
1.1814	0.0012084\\
1.1862	0.0012356\\
1.191	0.0012634\\
1.1958	0.0012918\\
1.2005	0.0013209\\
1.2053	0.0013506\\
1.2101	0.0013809\\
1.2149	0.001412\\
1.2197	0.0014437\\
1.2245	0.0014762\\
1.2293	0.0015094\\
1.234	0.0015433\\
1.2388	0.001578\\
1.2436	0.0016135\\
1.2484	0.0016497\\
1.2532	0.0016868\\
1.258	0.0017247\\
1.2628	0.0017634\\
1.2675	0.001803\\
1.2723	0.0018435\\
1.2771	0.0018849\\
1.2819	0.0019271\\
1.2867	0.0019704\\
1.2915	0.0020146\\
1.2962	0.0020598\\
1.301	0.0021059\\
1.3058	0.0021531\\
1.3106	0.0022014\\
1.3154	0.0022507\\
1.3202	0.0023011\\
1.325	0.0023526\\
1.3297	0.0024052\\
1.3345	0.002459\\
1.3393	0.002514\\
1.3441	0.0025702\\
1.3489	0.0026276\\
1.3537	0.0026863\\
1.3585	0.0027463\\
1.3632	0.0028076\\
1.368	0.0028702\\
1.3728	0.0029342\\
1.3776	0.0029996\\
1.3824	0.0030664\\
1.3872	0.0031347\\
1.392	0.0032044\\
1.3967	0.0032757\\
1.4015	0.0033485\\
1.4063	0.0034229\\
1.4111	0.003499\\
1.4159	0.0035766\\
1.4207	0.003656\\
1.4255	0.003737\\
1.4302	0.0038198\\
1.435	0.0039044\\
1.4398	0.0039908\\
1.4446	0.0040791\\
1.4494	0.0041693\\
1.4542	0.0042614\\
1.459	0.0043555\\
1.4637	0.0044516\\
1.4685	0.0045497\\
1.4733	0.0046499\\
1.4781	0.0047523\\
1.4829	0.0048569\\
1.4877	0.0049636\\
1.4925	0.0050727\\
1.4972	0.005184\\
1.502	0.0052977\\
1.5068	0.0054139\\
1.5116	0.0055324\\
1.5164	0.0056535\\
1.5212	0.0057771\\
1.526	0.0059033\\
1.5307	0.0060322\\
1.5355	0.0061638\\
1.5403	0.0062981\\
1.5451	0.0064353\\
1.5499	0.0065753\\
1.5547	0.0067182\\
1.5595	0.0068641\\
1.5642	0.007013\\
1.569	0.007165\\
1.5738	0.0073202\\
1.5786	0.0074786\\
1.5834	0.0076403\\
1.5882	0.0078053\\
1.593	0.0079737\\
1.5977	0.0081455\\
1.6025	0.0083209\\
1.6073	0.0084999\\
1.6121	0.0086825\\
1.6169	0.0088689\\
1.6217	0.0090591\\
1.6265	0.0092531\\
1.6312	0.0094511\\
1.636	0.0096531\\
1.6408	0.0098592\\
1.6456	0.01007\\
1.6504	0.010284\\
1.6552	0.010503\\
1.6599	0.010726\\
1.6647	0.010954\\
1.6695	0.011186\\
1.6743	0.011423\\
1.6791	0.011665\\
1.6839	0.011911\\
1.6887	0.012163\\
1.6934	0.012419\\
1.6982	0.012681\\
1.703	0.012947\\
1.7078	0.013219\\
1.7126	0.013496\\
1.7174	0.013779\\
1.7222	0.014067\\
1.7269	0.014361\\
1.7317	0.014661\\
1.7365	0.014967\\
1.7413	0.015278\\
1.7461	0.015596\\
1.7509	0.015919\\
1.7557	0.016249\\
1.7604	0.016586\\
1.7652	0.016929\\
1.77	0.017278\\
1.7748	0.017634\\
1.7796	0.017997\\
1.7844	0.018367\\
1.7892	0.018744\\
1.7939	0.019128\\
1.7987	0.01952\\
1.8035	0.019919\\
1.8083	0.020325\\
1.8131	0.020739\\
1.8179	0.021161\\
1.8227	0.021591\\
1.8274	0.022029\\
1.8322	0.022475\\
1.837	0.02293\\
1.8418	0.023393\\
1.8466	0.023864\\
1.8514	0.024344\\
1.8562	0.024833\\
1.8609	0.025332\\
1.8657	0.025839\\
1.8705	0.026356\\
1.8753	0.026882\\
1.8801	0.027417\\
1.8849	0.027963\\
1.8897	0.028518\\
1.8944	0.029084\\
1.8992	0.02966\\
1.904	0.030246\\
1.9088	0.030842\\
1.9136	0.03145\\
1.9184	0.032068\\
1.9232	0.032697\\
1.9279	0.033338\\
1.9327	0.03399\\
1.9375	0.034654\\
1.9423	0.035329\\
1.9471	0.036016\\
1.9519	0.036715\\
1.9567	0.037427\\
1.9614	0.038151\\
1.9662	0.038888\\
1.971	0.039637\\
1.9758	0.0404\\
1.9806	0.041176\\
1.9854	0.041965\\
1.9902	0.042768\\
1.9949	0.043584\\
1.9997	0.044415\\
2.0045	0.04526\\
2.0093	0.046119\\
2.0141	0.046993\\
2.0189	0.047882\\
2.0236	0.048786\\
2.0284	0.049705\\
2.0332	0.05064\\
2.038	0.051591\\
2.0428	0.052557\\
2.0476	0.05354\\
2.0524	0.054539\\
2.0571	0.055554\\
2.0619	0.056587\\
2.0667	0.057636\\
2.0715	0.058703\\
2.0763	0.059788\\
2.0811	0.06089\\
2.0859	0.062011\\
2.0906	0.06315\\
2.0954	0.064307\\
2.1002	0.065484\\
2.105	0.066679\\
2.1098	0.067894\\
2.1146	0.069129\\
2.1194	0.070383\\
2.1241	0.071658\\
2.1289	0.072953\\
2.1337	0.07427\\
2.1385	0.075607\\
2.1433	0.076965\\
2.1481	0.078346\\
2.1529	0.079748\\
2.1576	0.081173\\
2.1624	0.08262\\
2.1672	0.08409\\
2.172	0.085583\\
2.1768	0.0871\\
2.1816	0.08864\\
2.1864	0.090205\\
2.1911	0.091794\\
2.1959	0.093408\\
2.2007	0.095047\\
2.2055	0.096711\\
2.2103	0.098401\\
2.2151	0.10012\\
2.2199	0.10186\\
2.2246	0.10363\\
2.2294	0.10542\\
2.2342	0.10725\\
2.239	0.1091\\
2.2438	0.11097\\
2.2486	0.11288\\
2.2534	0.11481\\
2.2581	0.11678\\
2.2629	0.11877\\
2.2677	0.12078\\
2.2725	0.12283\\
2.2773	0.12491\\
2.2821	0.12702\\
2.2869	0.12915\\
2.2916	0.13132\\
2.2964	0.13351\\
2.3012	0.13573\\
2.306	0.13799\\
2.3108	0.14027\\
2.3156	0.14258\\
2.3204	0.14492\\
2.3251	0.1473\\
2.3299	0.1497\\
2.3347	0.15213\\
2.3395	0.15459\\
2.3443	0.15708\\
2.3491	0.1596\\
2.3539	0.16215\\
2.3586	0.16473\\
2.3634	0.16734\\
2.3682	0.16998\\
2.373	0.17264\\
2.3778	0.17534\\
2.3826	0.17807\\
2.3873	0.18083\\
2.3921	0.18361\\
2.3969	0.18643\\
2.4017	0.18927\\
2.4065	0.19215\\
2.4113	0.19505\\
2.4161	0.19799\\
2.4208	0.20095\\
2.4256	0.20395\\
2.4304	0.20697\\
2.4352	0.21003\\
2.44	0.21312\\
2.4448	0.21623\\
2.4496	0.21938\\
2.4543	0.22256\\
2.4591	0.22577\\
2.4639	0.22901\\
2.4687	0.23228\\
2.4735	0.23559\\
2.4783	0.23893\\
2.4831	0.2423\\
2.4878	0.2457\\
2.4926	0.24913\\
2.4974	0.2526\\
2.5022	0.25611\\
2.507	0.25964\\
2.5118	0.26321\\
2.5166	0.26682\\
2.5213	0.27046\\
2.5261	0.27413\\
2.5309	0.27784\\
2.5357	0.28158\\
2.5405	0.28536\\
2.5453	0.28917\\
2.5501	0.29301\\
2.5548	0.2969\\
2.5596	0.30081\\
2.5644	0.30476\\
2.5692	0.30874\\
2.574	0.31276\\
2.5788	0.31681\\
2.5836	0.3209\\
2.5883	0.32501\\
2.5931	0.32916\\
2.5979	0.33334\\
2.6027	0.33755\\
2.6075	0.34179\\
2.6123	0.34606\\
2.6171	0.35036\\
2.6218	0.35468\\
2.6266	0.35903\\
2.6314	0.36341\\
2.6362	0.36781\\
2.641	0.37223\\
2.6458	0.37667\\
2.6506	0.38114\\
2.6553	0.38563\\
2.6601	0.39013\\
2.6649	0.39465\\
2.6697	0.39919\\
2.6745	0.40374\\
2.6793	0.40831\\
2.6841	0.41288\\
2.6888	0.41747\\
2.6936	0.42207\\
2.6984	0.42668\\
2.7032	0.4313\\
2.708	0.43593\\
2.7128	0.44056\\
2.7176	0.4452\\
2.7223	0.44984\\
2.7271	0.45449\\
2.7319	0.45914\\
2.7367	0.4638\\
2.7415	0.46845\\
2.7463	0.47311\\
2.751	0.47778\\
2.7558	0.48244\\
2.7606	0.48711\\
2.7654	0.49178\\
2.7702	0.49644\\
2.775	0.50111\\
2.7798	0.50578\\
2.7845	0.51046\\
2.7893	0.51513\\
2.7941	0.5198\\
2.7989	0.52447\\
2.8037	0.52914\\
2.8085	0.53382\\
2.8133	0.53849\\
2.818	0.54316\\
2.8228	0.54783\\
2.8276	0.5525\\
2.8324	0.55716\\
2.8372	0.56183\\
2.842	0.56649\\
2.8468	0.57115\\
2.8515	0.5758\\
2.8563	0.58045\\
2.8611	0.58509\\
2.8659	0.58973\\
2.8707	0.59437\\
2.8755	0.59899\\
2.8803	0.60361\\
2.885	0.60822\\
2.8898	0.61282\\
2.8946	0.61741\\
2.8994	0.62199\\
2.9042	0.62655\\
2.909	0.63111\\
2.9138	0.63565\\
2.9185	0.64017\\
2.9233	0.64468\\
2.9281	0.64918\\
2.9329	0.65366\\
2.9377	0.65812\\
2.9425	0.66256\\
2.9473	0.66699\\
2.952	0.67139\\
2.9568	0.67577\\
2.9616	0.68013\\
2.9664	0.68447\\
2.9712	0.68878\\
2.976	0.69306\\
2.9808	0.69731\\
2.9855	0.70154\\
2.9903	0.70574\\
2.9951	0.7099\\
2.9999	0.71403\\
3.0047	0.71812\\
3.0095	0.72217\\
3.0143	0.72618\\
3.019	0.73014\\
3.0238	0.73407\\
3.0286	0.73794\\
3.0334	0.74176\\
3.0382	0.74553\\
3.043	0.74924\\
3.0478	0.75289\\
3.0525	0.75648\\
3.0573	0.76\\
3.0621	0.76345\\
3.0669	0.76684\\
3.0717	0.77014\\
3.0765	0.77338\\
3.0813	0.77653\\
3.086	0.77959\\
3.0908	0.78257\\
3.0956	0.78547\\
3.1004	0.78826\\
3.1052	0.79097\\
3.11	0.79358\\
3.1147	0.79608\\
3.1195	0.79849\\
3.1243	0.80078\\
3.1291	0.80297\\
3.1339	0.80505\\
3.1387	0.80702\\
3.1435	0.80888\\
3.1482	0.81061\\
3.153	0.81223\\
3.1578	0.81374\\
3.1626	0.81512\\
3.1674	0.81638\\
3.1722	0.81752\\
3.177	0.81853\\
3.1817	0.81942\\
3.1865	0.82019\\
3.1913	0.82083\\
3.1961	0.82135\\
3.2009	0.82175\\
3.2057	0.82202\\
3.2105	0.82217\\
3.2152	0.82219\\
3.22	0.8221\\
3.2248	0.82189\\
3.2296	0.82155\\
3.2344	0.8211\\
3.2392	0.82053\\
3.244	0.81985\\
3.2487	0.81906\\
3.2535	0.81815\\
3.2583	0.81714\\
3.2631	0.81602\\
3.2679	0.81479\\
3.2727	0.81346\\
3.2775	0.81203\\
3.2822	0.8105\\
3.287	0.80887\\
3.2918	0.80715\\
3.2966	0.80533\\
3.3014	0.80342\\
3.3062	0.80143\\
3.311	0.79935\\
3.3157	0.79718\\
3.3205	0.79493\\
3.3253	0.79261\\
3.3301	0.7902\\
3.3349	0.78772\\
3.3397	0.78516\\
3.3445	0.78254\\
3.3492	0.77984\\
3.354	0.77708\\
3.3588	0.77425\\
3.3636	0.77135\\
3.3684	0.7684\\
3.3732	0.76539\\
3.378	0.76232\\
3.3827	0.7592\\
3.3875	0.75602\\
3.3923	0.75279\\
3.3971	0.74951\\
3.4019	0.74619\\
3.4067	0.74281\\
3.4115	0.7394\\
3.4162	0.73594\\
3.421	0.73245\\
3.4258	0.72891\\
3.4306	0.72533\\
3.4354	0.72172\\
3.4402	0.71808\\
3.445	0.7144\\
3.4497	0.71069\\
3.4545	0.70695\\
3.4593	0.70318\\
3.4641	0.69938\\
3.4689	0.69556\\
3.4737	0.69171\\
3.4784	0.68783\\
3.4832	0.68393\\
3.488	0.68\\
3.4928	0.67606\\
3.4976	0.67209\\
3.5024	0.6681\\
3.5072	0.66409\\
3.5119	0.66006\\
3.5167	0.65602\\
3.5215	0.65195\\
3.5263	0.64787\\
3.5311	0.64377\\
3.5359	0.63966\\
3.5407	0.63553\\
3.5454	0.63139\\
3.5502	0.62724\\
3.555	0.62307\\
3.5598	0.61889\\
3.5646	0.6147\\
3.5694	0.6105\\
3.5742	0.60629\\
3.5789	0.60207\\
3.5837	0.59784\\
3.5885	0.59361\\
3.5933	0.58937\\
3.5981	0.58512\\
3.6029	0.58086\\
3.6077	0.57661\\
3.6124	0.57234\\
3.6172	0.56808\\
3.622	0.56381\\
3.6268	0.55954\\
3.6316	0.55526\\
3.6364	0.55099\\
3.6412	0.54672\\
3.6459	0.54245\\
3.6507	0.53817\\
3.6555	0.53391\\
3.6603	0.52964\\
3.6651	0.52538\\
3.6699	0.52112\\
3.6747	0.51686\\
3.6794	0.51261\\
3.6842	0.50837\\
3.689	0.50413\\
3.6938	0.4999\\
3.6986	0.49567\\
3.7034	0.49146\\
3.7082	0.48725\\
3.7129	0.48305\\
3.7177	0.47886\\
3.7225	0.47468\\
3.7273	0.47051\\
3.7321	0.46636\\
3.7369	0.46221\\
3.7417	0.45808\\
3.7464	0.45396\\
3.7512	0.44985\\
3.756	0.44575\\
3.7608	0.44167\\
3.7656	0.43761\\
3.7704	0.43355\\
3.7752	0.42952\\
3.7799	0.42549\\
3.7847	0.42149\\
3.7895	0.4175\\
3.7943	0.41353\\
3.7991	0.40957\\
3.8039	0.40563\\
3.8087	0.40171\\
3.8134	0.39781\\
3.8182	0.39392\\
3.823	0.39005\\
3.8278	0.3862\\
3.8326	0.38238\\
3.8374	0.37857\\
3.8421	0.37477\\
3.8469	0.371\\
3.8517	0.36725\\
3.8565	0.36352\\
3.8613	0.35981\\
3.8661	0.35612\\
3.8709	0.35246\\
3.8756	0.34881\\
3.8804	0.34518\\
3.8852	0.34158\\
3.89	0.338\\
3.8948	0.33443\\
3.8996	0.3309\\
3.9044	0.32738\\
3.9091	0.32388\\
3.9139	0.32041\\
3.9187	0.31696\\
3.9235	0.31354\\
3.9283	0.31013\\
3.9331	0.30675\\
3.9379	0.30339\\
3.9426	0.30006\\
3.9474	0.29675\\
3.9522	0.29346\\
3.957	0.29019\\
3.9618	0.28695\\
3.9666	0.28374\\
3.9714	0.28054\\
3.9761	0.27737\\
3.9809	0.27422\\
3.9857	0.2711\\
3.9905	0.268\\
3.9953	0.26493\\
4.0001	0.26188\\
4.0049	0.25885\\
4.0096	0.25585\\
4.0144	0.25287\\
4.0192	0.24991\\
4.024	0.24698\\
4.0288	0.24407\\
4.0336	0.24119\\
4.0384	0.23833\\
4.0431	0.23549\\
4.0479	0.23268\\
4.0527	0.2299\\
4.0575	0.22713\\
4.0623	0.22439\\
4.0671	0.22168\\
4.0719	0.21898\\
4.0766	0.21632\\
4.0814	0.21367\\
4.0862	0.21105\\
4.091	0.20845\\
4.0958	0.20588\\
4.1006	0.20333\\
4.1054	0.20081\\
4.1101	0.1983\\
4.1149	0.19582\\
4.1197	0.19337\\
4.1245	0.19093\\
4.1293	0.18852\\
4.1341	0.18614\\
4.1389	0.18377\\
4.1436	0.18143\\
4.1484	0.17912\\
4.1532	0.17682\\
4.158	0.17455\\
4.1628	0.1723\\
4.1676	0.17007\\
4.1724	0.16787\\
4.1771	0.16568\\
4.1819	0.16352\\
4.1867	0.16139\\
4.1915	0.15927\\
4.1963	0.15718\\
4.2011	0.1551\\
4.2058	0.15305\\
4.2106	0.15102\\
4.2154	0.14901\\
4.2202	0.14703\\
4.225	0.14506\\
4.2298	0.14312\\
4.2346	0.14119\\
4.2393	0.13929\\
4.2441	0.13741\\
4.2489	0.13555\\
4.2537	0.13371\\
4.2585	0.13189\\
4.2633	0.13009\\
4.2681	0.12831\\
4.2728	0.12655\\
4.2776	0.1248\\
4.2824	0.12308\\
4.2872	0.12138\\
4.292	0.1197\\
4.2968	0.11804\\
4.3016	0.11639\\
4.3063	0.11477\\
4.3111	0.11316\\
4.3159	0.11158\\
4.3207	0.11001\\
4.3255	0.10846\\
4.3303	0.10692\\
4.3351	0.10541\\
4.3398	0.10391\\
4.3446	0.10244\\
4.3494	0.10097\\
4.3542	0.099531\\
4.359	0.098105\\
4.3638	0.096697\\
4.3686	0.095305\\
4.3733	0.093931\\
4.3781	0.092573\\
4.3829	0.091231\\
4.3877	0.089907\\
4.3925	0.088598\\
4.3973	0.087306\\
4.4021	0.086029\\
4.4068	0.084769\\
4.4116	0.083524\\
4.4164	0.082295\\
4.4212	0.081081\\
4.426	0.079883\\
4.4308	0.0787\\
4.4356	0.077531\\
4.4403	0.076378\\
4.4451	0.075239\\
4.4499	0.074114\\
4.4547	0.073005\\
4.4595	0.071909\\
4.4643	0.070827\\
4.4691	0.069759\\
4.4738	0.068706\\
4.4786	0.067665\\
4.4834	0.066639\\
4.4882	0.065625\\
4.493	0.064625\\
4.4978	0.063638\\
4.5026	0.062664\\
4.5073	0.061703\\
4.5121	0.060754\\
4.5169	0.059818\\
4.5217	0.058895\\
4.5265	0.057983\\
4.5313	0.057084\\
4.5361	0.056197\\
4.5408	0.055321\\
4.5456	0.054458\\
4.5504	0.053606\\
4.5552	0.052765\\
4.56	0.051936\\
4.5648	0.051118\\
4.5695	0.050311\\
4.5743	0.049515\\
4.5791	0.04873\\
4.5839	0.047956\\
4.5887	0.047192\\
4.5935	0.046439\\
4.5983	0.045697\\
4.603	0.044964\\
4.6078	0.044242\\
4.6126	0.043529\\
4.6174	0.042827\\
4.6222	0.042134\\
4.627	0.041451\\
4.6318	0.040778\\
4.6365	0.040114\\
4.6413	0.03946\\
4.6461	0.038815\\
4.6509	0.038179\\
4.6557	0.037552\\
4.6605	0.036933\\
4.6653	0.036324\\
4.67	0.035724\\
4.6748	0.035132\\
4.6796	0.034549\\
4.6844	0.033974\\
4.6892	0.033407\\
4.694	0.032849\\
4.6988	0.032299\\
4.7035	0.031757\\
4.7083	0.031222\\
4.7131	0.030696\\
4.7179	0.030178\\
4.7227	0.029667\\
4.7275	0.029163\\
4.7323	0.028668\\
4.737	0.028179\\
4.7418	0.027698\\
4.7466	0.027224\\
4.7514	0.026758\\
4.7562	0.026298\\
4.761	0.025845\\
4.7658	0.0254\\
4.7705	0.024961\\
4.7753	0.024528\\
4.7801	0.024103\\
4.7849	0.023683\\
4.7897	0.023271\\
4.7945	0.022864\\
4.7993	0.022464\\
4.804	0.02207\\
4.8088	0.021683\\
4.8136	0.021301\\
4.8184	0.020925\\
4.8232	0.020556\\
4.828	0.020192\\
4.8328	0.019834\\
4.8375	0.019481\\
4.8423	0.019134\\
4.8471	0.018793\\
4.8519	0.018457\\
4.8567	0.018126\\
4.8615	0.017801\\
4.8663	0.017481\\
4.871	0.017166\\
4.8758	0.016856\\
4.8806	0.016551\\
4.8854	0.016251\\
4.8902	0.015956\\
4.895	0.015666\\
4.8998	0.01538\\
4.9045	0.0151\\
4.9093	0.014823\\
4.9141	0.014552\\
4.9189	0.014284\\
4.9237	0.014022\\
4.9285	0.013763\\
4.9332	0.013509\\
4.938	0.013259\\
4.9428	0.013013\\
4.9476	0.012771\\
4.9524	0.012534\\
4.9572	0.0123\\
4.962	0.01207\\
4.9667	0.011844\\
4.9715	0.011622\\
4.9763	0.011404\\
4.9811	0.011189\\
4.9859	0.010978\\
4.9907	0.01077\\
4.9955	0.010566\\
5.0002	0.010366\\
5.005	0.010169\\
5.0098	0.0099748\\
5.0146	0.0097843\\
5.0194	0.0095971\\
5.0242	0.0094132\\
5.029	0.0092323\\
5.0337	0.0090546\\
5.0385	0.00888\\
5.0433	0.0087084\\
5.0481	0.0085397\\
5.0529	0.008374\\
5.0577	0.0082111\\
5.0625	0.0080511\\
5.0672	0.0078939\\
5.072	0.0077395\\
5.0768	0.0075877\\
5.0816	0.0074386\\
5.0864	0.0072921\\
5.0912	0.0071482\\
5.096	0.0070069\\
5.1007	0.0068681\\
5.1055	0.0067317\\
5.1103	0.0065977\\
5.1151	0.0064662\\
5.1199	0.006337\\
5.1247	0.0062101\\
5.1295	0.0060854\\
5.1342	0.0059631\\
5.139	0.0058429\\
5.1438	0.0057249\\
5.1486	0.005609\\
5.1534	0.0054953\\
5.1582	0.0053836\\
5.163	0.0052739\\
5.1677	0.0051663\\
5.1725	0.0050606\\
5.1773	0.0049569\\
5.1821	0.0048551\\
5.1869	0.0047552\\
5.1917	0.0046571\\
5.1965	0.0045609\\
5.2012	0.0044664\\
5.206	0.0043737\\
5.2108	0.0042828\\
5.2156	0.0041935\\
5.2204	0.004106\\
5.2252	0.0040201\\
5.23	0.0039358\\
5.2347	0.0038531\\
5.2395	0.003772\\
5.2443	0.0036925\\
5.2491	0.0036145\\
5.2539	0.003538\\
5.2587	0.0034629\\
5.2635	0.0033893\\
5.2682	0.0033172\\
5.273	0.0032464\\
5.2778	0.003177\\
5.2826	0.003109\\
5.2874	0.0030424\\
5.2922	0.002977\\
5.2969	0.0029129\\
5.3017	0.0028501\\
5.3065	0.0027886\\
5.3113	0.0027282\\
5.3161	0.0026691\\
5.3209	0.0026112\\
5.3257	0.0025544\\
5.3304	0.0024988\\
5.3352	0.0024443\\
5.34	0.0023909\\
5.3448	0.0023386\\
5.3496	0.0022874\\
5.3544	0.0022372\\
5.3544	0.0018281\\
5.3496	0.0018724\\
5.3448	0.0019177\\
5.34	0.0019639\\
5.3352	0.0020112\\
5.3304	0.0020595\\
5.3257	0.0021089\\
5.3209	0.0021593\\
5.3161	0.0022109\\
5.3113	0.0022635\\
5.3065	0.0023173\\
5.3017	0.0023723\\
5.2969	0.0024285\\
5.2922	0.0024858\\
5.2874	0.0025444\\
5.2826	0.0026043\\
5.2778	0.0026654\\
5.273	0.0027278\\
5.2682	0.0027916\\
5.2635	0.0028567\\
5.2587	0.0029232\\
5.2539	0.002991\\
5.2491	0.0030603\\
5.2443	0.0031311\\
5.2395	0.0032033\\
5.2347	0.0032771\\
5.23	0.0033523\\
5.2252	0.0034292\\
5.2204	0.0035076\\
5.2156	0.0035876\\
5.2108	0.0036693\\
5.206	0.0037527\\
5.2012	0.0038378\\
5.1965	0.0039246\\
5.1917	0.0040131\\
5.1869	0.0041035\\
5.1821	0.0041957\\
5.1773	0.0042897\\
5.1725	0.0043857\\
5.1677	0.0044835\\
5.163	0.0045833\\
5.1582	0.0046851\\
5.1534	0.004789\\
5.1486	0.0048949\\
5.1438	0.0050028\\
5.139	0.005113\\
5.1342	0.0052252\\
5.1295	0.0053397\\
5.1247	0.0054564\\
5.1199	0.0055754\\
5.1151	0.0056966\\
5.1103	0.0058203\\
5.1055	0.0059463\\
5.1007	0.0060747\\
5.096	0.0062056\\
5.0912	0.006339\\
5.0864	0.0064749\\
5.0816	0.0066134\\
5.0768	0.0067545\\
5.072	0.0068983\\
5.0672	0.0070448\\
5.0625	0.007194\\
5.0577	0.0073461\\
5.0529	0.0075009\\
5.0481	0.0076586\\
5.0433	0.0078193\\
5.0385	0.0079829\\
5.0337	0.0081495\\
5.029	0.0083191\\
5.0242	0.0084919\\
5.0194	0.0086677\\
5.0146	0.0088468\\
5.0098	0.0090291\\
5.005	0.0092147\\
5.0002	0.0094037\\
4.9955	0.009596\\
4.9907	0.0097917\\
4.9859	0.0099909\\
4.9811	0.010194\\
4.9763	0.0104\\
4.9715	0.01061\\
4.9667	0.010824\\
4.962	0.011041\\
4.9572	0.011262\\
4.9524	0.011487\\
4.9476	0.011716\\
4.9428	0.011949\\
4.938	0.012186\\
4.9332	0.012426\\
4.9285	0.012672\\
4.9237	0.012921\\
4.9189	0.013174\\
4.9141	0.013432\\
4.9093	0.013694\\
4.9045	0.01396\\
4.8998	0.014231\\
4.895	0.014507\\
4.8902	0.014787\\
4.8854	0.015071\\
4.8806	0.015361\\
4.8758	0.015655\\
4.871	0.015954\\
4.8663	0.016258\\
4.8615	0.016567\\
4.8567	0.016881\\
4.8519	0.0172\\
4.8471	0.017524\\
4.8423	0.017854\\
4.8375	0.018189\\
4.8328	0.018529\\
4.828	0.018875\\
4.8232	0.019226\\
4.8184	0.019583\\
4.8136	0.019945\\
4.8088	0.020314\\
4.804	0.020688\\
4.7993	0.021068\\
4.7945	0.021454\\
4.7897	0.021846\\
4.7849	0.022244\\
4.7801	0.022649\\
4.7753	0.02306\\
4.7705	0.023477\\
4.7658	0.023901\\
4.761	0.024331\\
4.7562	0.024768\\
4.7514	0.025212\\
4.7466	0.025663\\
4.7418	0.02612\\
4.737	0.026585\\
4.7323	0.027056\\
4.7275	0.027535\\
4.7227	0.028021\\
4.7179	0.028515\\
4.7131	0.029016\\
4.7083	0.029524\\
4.7035	0.03004\\
4.6988	0.030564\\
4.694	0.031096\\
4.6892	0.031636\\
4.6844	0.032184\\
4.6796	0.03274\\
4.6748	0.033304\\
4.67	0.033876\\
4.6653	0.034458\\
4.6605	0.035047\\
4.6557	0.035645\\
4.6509	0.036253\\
4.6461	0.036869\\
4.6413	0.037494\\
4.6365	0.038128\\
4.6318	0.038771\\
4.627	0.039423\\
4.6222	0.040086\\
4.6174	0.040757\\
4.6126	0.041438\\
4.6078	0.042129\\
4.603	0.04283\\
4.5983	0.043541\\
4.5935	0.044262\\
4.5887	0.044993\\
4.5839	0.045734\\
4.5791	0.046486\\
4.5743	0.047249\\
4.5695	0.048022\\
4.5648	0.048806\\
4.56	0.049601\\
4.5552	0.050407\\
4.5504	0.051224\\
4.5456	0.052052\\
4.5408	0.052892\\
4.5361	0.053743\\
4.5313	0.054606\\
4.5265	0.05548\\
4.5217	0.056367\\
4.5169	0.057265\\
4.5121	0.058176\\
4.5073	0.059099\\
4.5026	0.060034\\
4.4978	0.060982\\
4.493	0.061942\\
4.4882	0.062915\\
4.4834	0.063901\\
4.4786	0.0649\\
4.4738	0.065912\\
4.4691	0.066938\\
4.4643	0.067977\\
4.4595	0.069029\\
4.4547	0.070095\\
4.4499	0.071175\\
4.4451	0.072269\\
4.4403	0.073377\\
4.4356	0.074499\\
4.4308	0.075635\\
4.426	0.076786\\
4.4212	0.077952\\
4.4164	0.079132\\
4.4116	0.080327\\
4.4068	0.081537\\
4.4021	0.082763\\
4.3973	0.084003\\
4.3925	0.085259\\
4.3877	0.086531\\
4.3829	0.087818\\
4.3781	0.089121\\
4.3733	0.09044\\
4.3686	0.091776\\
4.3638	0.093127\\
4.359	0.094495\\
4.3542	0.09588\\
4.3494	0.097281\\
4.3446	0.098699\\
4.3398	0.10013\\
4.3351	0.10159\\
4.3303	0.10306\\
4.3255	0.10454\\
4.3207	0.10605\\
4.3159	0.10757\\
4.3111	0.10911\\
4.3063	0.11067\\
4.3016	0.11224\\
4.2968	0.11384\\
4.292	0.11545\\
4.2872	0.11708\\
4.2824	0.11873\\
4.2776	0.1204\\
4.2728	0.12209\\
4.2681	0.1238\\
4.2633	0.12553\\
4.2585	0.12727\\
4.2537	0.12904\\
4.2489	0.13082\\
4.2441	0.13263\\
4.2393	0.13446\\
4.2346	0.1363\\
4.2298	0.13817\\
4.225	0.14005\\
4.2202	0.14196\\
4.2154	0.14389\\
4.2106	0.14584\\
4.2058	0.14781\\
4.2011	0.1498\\
4.1963	0.15181\\
4.1915	0.15384\\
4.1867	0.1559\\
4.1819	0.15797\\
4.1771	0.16007\\
4.1724	0.16219\\
4.1676	0.16433\\
4.1628	0.1665\\
4.158	0.16868\\
4.1532	0.17089\\
4.1484	0.17312\\
4.1436	0.17538\\
4.1389	0.17765\\
4.1341	0.17995\\
4.1293	0.18227\\
4.1245	0.18462\\
4.1197	0.18699\\
4.1149	0.18938\\
4.1101	0.19179\\
4.1054	0.19423\\
4.1006	0.19669\\
4.0958	0.19917\\
4.091	0.20168\\
4.0862	0.20421\\
4.0814	0.20676\\
4.0766	0.20934\\
4.0719	0.21194\\
4.0671	0.21457\\
4.0623	0.21722\\
4.0575	0.21989\\
4.0527	0.22259\\
4.0479	0.22531\\
4.0431	0.22805\\
4.0384	0.23082\\
4.0336	0.23361\\
4.0288	0.23643\\
4.024	0.23927\\
4.0192	0.24213\\
4.0144	0.24502\\
4.0096	0.24793\\
4.0049	0.25087\\
4.0001	0.25383\\
3.9953	0.25682\\
3.9905	0.25982\\
3.9857	0.26286\\
3.9809	0.26591\\
3.9761	0.26899\\
3.9714	0.27209\\
3.9666	0.27522\\
3.9618	0.27837\\
3.957	0.28155\\
3.9522	0.28474\\
3.9474	0.28796\\
3.9426	0.29121\\
3.9379	0.29448\\
3.9331	0.29777\\
3.9283	0.30108\\
3.9235	0.30441\\
3.9187	0.30777\\
3.9139	0.31115\\
3.9091	0.31456\\
3.9044	0.31798\\
3.8996	0.32143\\
3.8948	0.3249\\
3.89	0.32839\\
3.8852	0.3319\\
3.8804	0.33544\\
3.8756	0.33899\\
3.8709	0.34257\\
3.8661	0.34616\\
3.8613	0.34978\\
3.8565	0.35342\\
3.8517	0.35707\\
3.8469	0.36075\\
3.8421	0.36445\\
3.8374	0.36816\\
3.8326	0.3719\\
3.8278	0.37565\\
3.823	0.37942\\
3.8182	0.38321\\
3.8134	0.38702\\
3.8087	0.39084\\
3.8039	0.39468\\
3.7991	0.39854\\
3.7943	0.40242\\
3.7895	0.40631\\
3.7847	0.41022\\
3.7799	0.41414\\
3.7752	0.41807\\
3.7704	0.42203\\
3.7656	0.42599\\
3.7608	0.42997\\
3.756	0.43397\\
3.7512	0.43797\\
3.7464	0.44199\\
3.7417	0.44602\\
3.7369	0.45007\\
3.7321	0.45412\\
3.7273	0.45819\\
3.7225	0.46226\\
3.7177	0.46635\\
3.7129	0.47045\\
3.7082	0.47455\\
3.7034	0.47867\\
3.6986	0.48279\\
3.6938	0.48692\\
3.689	0.49106\\
3.6842	0.4952\\
3.6794	0.49935\\
3.6747	0.50351\\
3.6699	0.50767\\
3.6651	0.51184\\
3.6603	0.51601\\
3.6555	0.52018\\
3.6507	0.52436\\
3.6459	0.52854\\
3.6412	0.53272\\
3.6364	0.5369\\
3.6316	0.54109\\
3.6268	0.54527\\
3.622	0.54946\\
3.6172	0.55364\\
3.6124	0.55782\\
3.6077	0.562\\
3.6029	0.56618\\
3.5981	0.57035\\
3.5933	0.57452\\
3.5885	0.57868\\
3.5837	0.58284\\
3.5789	0.58699\\
3.5742	0.59113\\
3.5694	0.59527\\
3.5646	0.59939\\
3.5598	0.60351\\
3.555	0.60761\\
3.5502	0.61171\\
3.5454	0.61579\\
3.5407	0.61985\\
3.5359	0.62391\\
3.5311	0.62794\\
3.5263	0.63197\\
3.5215	0.63597\\
3.5167	0.63995\\
3.5119	0.64392\\
3.5072	0.64786\\
3.5024	0.65178\\
3.4976	0.65568\\
3.4928	0.65955\\
3.488	0.6634\\
3.4832	0.66722\\
3.4784	0.67102\\
3.4737	0.67478\\
3.4689	0.67851\\
3.4641	0.68221\\
3.4593	0.68588\\
3.4545	0.68951\\
3.4497	0.69311\\
3.445	0.69667\\
3.4402	0.70019\\
3.4354	0.70367\\
3.4306	0.70711\\
3.4258	0.71051\\
3.421	0.71386\\
3.4162	0.71717\\
3.4115	0.72044\\
3.4067	0.72365\\
3.4019	0.72682\\
3.3971	0.72993\\
3.3923	0.73299\\
3.3875	0.736\\
3.3827	0.73895\\
3.378	0.74185\\
3.3732	0.74469\\
3.3684	0.74747\\
3.3636	0.75019\\
3.3588	0.75285\\
3.354	0.75544\\
3.3492	0.75796\\
3.3445	0.76042\\
3.3397	0.76281\\
3.3349	0.76512\\
3.3301	0.76736\\
3.3253	0.76953\\
3.3205	0.77162\\
3.3157	0.77363\\
3.311	0.77555\\
3.3062	0.7774\\
3.3014	0.77916\\
3.2966	0.78083\\
3.2918	0.78241\\
3.287	0.78391\\
3.2822	0.78531\\
3.2775	0.78661\\
3.2727	0.78783\\
3.2679	0.78894\\
3.2631	0.78995\\
3.2583	0.79087\\
3.2535	0.79168\\
3.2487	0.79239\\
3.244	0.793\\
3.2392	0.7935\\
3.2344	0.7939\\
3.2296	0.7942\\
3.2248	0.79438\\
3.22	0.79446\\
3.2152	0.79443\\
3.2105	0.7943\\
3.2057	0.79406\\
3.2009	0.79371\\
3.1961	0.79326\\
3.1913	0.79269\\
3.1865	0.79203\\
3.1817	0.79125\\
3.177	0.79037\\
3.1722	0.78939\\
3.1674	0.7883\\
3.1626	0.78711\\
3.1578	0.78582\\
3.153	0.78442\\
3.1482	0.78293\\
3.1435	0.78134\\
3.1387	0.77965\\
3.1339	0.77786\\
3.1291	0.77598\\
3.1243	0.774\\
3.1195	0.77193\\
3.1147	0.76977\\
3.11	0.76752\\
3.1052	0.76518\\
3.1004	0.76276\\
3.0956	0.76024\\
3.0908	0.75765\\
3.086	0.75497\\
3.0813	0.75222\\
3.0765	0.74938\\
3.0717	0.74647\\
3.0669	0.74348\\
3.0621	0.74041\\
3.0573	0.73727\\
3.0525	0.73406\\
3.0478	0.73078\\
3.043	0.72744\\
3.0382	0.72402\\
3.0334	0.72054\\
3.0286	0.717\\
3.0238	0.71339\\
3.019	0.70973\\
3.0143	0.706\\
3.0095	0.70222\\
3.0047	0.69838\\
2.9999	0.69449\\
2.9951	0.69054\\
2.9903	0.68655\\
2.9855	0.6825\\
2.9808	0.67841\\
2.976	0.67427\\
2.9712	0.67009\\
2.9664	0.66587\\
2.9616	0.66161\\
2.9568	0.65731\\
2.952	0.65298\\
2.9473	0.64861\\
2.9425	0.64421\\
2.9377	0.63979\\
2.9329	0.63533\\
2.9281	0.63085\\
2.9233	0.62634\\
2.9185	0.62181\\
2.9138	0.61726\\
2.909	0.61269\\
2.9042	0.6081\\
2.8994	0.6035\\
2.8946	0.59888\\
2.8898	0.59424\\
2.885	0.5896\\
2.8803	0.58494\\
2.8755	0.58027\\
2.8707	0.57559\\
2.8659	0.57091\\
2.8611	0.56621\\
2.8563	0.56151\\
2.8515	0.55681\\
2.8468	0.55209\\
2.842	0.54738\\
2.8372	0.54266\\
2.8324	0.53794\\
2.8276	0.53321\\
2.8228	0.52849\\
2.818	0.52376\\
2.8133	0.51903\\
2.8085	0.51431\\
2.8037	0.50958\\
2.7989	0.50486\\
2.7941	0.50015\\
2.7893	0.49543\\
2.7845	0.49072\\
2.7798	0.48602\\
2.775	0.48132\\
2.7702	0.47663\\
2.7654	0.47195\\
2.7606	0.46727\\
2.7558	0.46261\\
2.751	0.45796\\
2.7463	0.45331\\
2.7415	0.44868\\
2.7367	0.44407\\
2.7319	0.43946\\
2.7271	0.43488\\
2.7223	0.4303\\
2.7176	0.42575\\
2.7128	0.42121\\
2.708	0.41668\\
2.7032	0.41218\\
2.6984	0.40769\\
2.6936	0.40323\\
2.6888	0.39878\\
2.6841	0.39435\\
2.6793	0.38995\\
2.6745	0.38556\\
2.6697	0.3812\\
2.6649	0.37686\\
2.6601	0.37254\\
2.6553	0.36824\\
2.6506	0.36397\\
2.6458	0.35972\\
2.641	0.3555\\
2.6362	0.35129\\
2.6314	0.34712\\
2.6266	0.34296\\
2.6218	0.33884\\
2.6171	0.33473\\
2.6123	0.33066\\
2.6075	0.32661\\
2.6027	0.32258\\
2.5979	0.31858\\
2.5931	0.31461\\
2.5883	0.31066\\
2.5836	0.30675\\
2.5788	0.30285\\
2.574	0.29899\\
2.5692	0.29516\\
2.5644	0.29135\\
2.5596	0.28757\\
2.5548	0.28382\\
2.5501	0.28009\\
2.5453	0.2764\\
2.5405	0.27274\\
2.5357	0.2691\\
2.5309	0.2655\\
2.5261	0.26192\\
2.5213	0.25838\\
2.5166	0.25486\\
2.5118	0.25138\\
2.507	0.24792\\
2.5022	0.2445\\
2.4974	0.24111\\
2.4926	0.23775\\
2.4878	0.23442\\
2.4831	0.23112\\
2.4783	0.22785\\
2.4735	0.22462\\
2.4687	0.22141\\
2.4639	0.21824\\
2.4591	0.2151\\
2.4543	0.21199\\
2.4496	0.20891\\
2.4448	0.20586\\
2.44	0.20285\\
2.4352	0.19986\\
2.4304	0.19691\\
2.4256	0.19399\\
2.4208	0.19111\\
2.4161	0.18825\\
2.4113	0.18542\\
2.4065	0.18263\\
2.4017	0.17987\\
2.3969	0.17714\\
2.3921	0.17443\\
2.3873	0.17176\\
2.3826	0.16913\\
2.3778	0.16652\\
2.373	0.16394\\
2.3682	0.16139\\
2.3634	0.15887\\
2.3586	0.15639\\
2.3539	0.15393\\
2.3491	0.1515\\
2.3443	0.1491\\
2.3395	0.14673\\
2.3347	0.14439\\
2.3299	0.14208\\
2.3251	0.13979\\
2.3204	0.13754\\
2.3156	0.13531\\
2.3108	0.13311\\
2.306	0.13094\\
2.3012	0.1288\\
2.2964	0.12668\\
2.2916	0.12459\\
2.2869	0.12252\\
2.2821	0.12049\\
2.2773	0.11848\\
2.2725	0.11649\\
2.2677	0.11453\\
2.2629	0.1126\\
2.2581	0.11069\\
2.2534	0.10881\\
2.2486	0.10696\\
2.2438	0.10512\\
2.239	0.10332\\
2.2342	0.10154\\
2.2294	0.09978\\
2.2246	0.098047\\
2.2199	0.096339\\
2.2151	0.094656\\
2.2103	0.092996\\
2.2055	0.09136\\
2.2007	0.089747\\
2.1959	0.088158\\
2.1911	0.086593\\
2.1864	0.08505\\
2.1816	0.083531\\
2.1768	0.082034\\
2.172	0.080559\\
2.1672	0.079108\\
2.1624	0.077678\\
2.1576	0.07627\\
2.1529	0.074884\\
2.1481	0.073519\\
2.1433	0.072176\\
2.1385	0.070854\\
2.1337	0.069553\\
2.1289	0.068272\\
2.1241	0.067012\\
2.1194	0.065772\\
2.1146	0.064552\\
2.1098	0.063351\\
2.105	0.06217\\
2.1002	0.061009\\
2.0954	0.059866\\
2.0906	0.058742\\
2.0859	0.057637\\
2.0811	0.05655\\
2.0763	0.055481\\
2.0715	0.05443\\
2.0667	0.053396\\
2.0619	0.05238\\
2.0571	0.051381\\
2.0524	0.050398\\
2.0476	0.049433\\
2.0428	0.048484\\
2.038	0.04755\\
2.0332	0.046633\\
2.0284	0.045732\\
2.0236	0.044846\\
2.0189	0.043975\\
2.0141	0.04312\\
2.0093	0.042279\\
2.0045	0.041453\\
1.9997	0.040642\\
1.9949	0.039844\\
1.9902	0.039061\\
1.9854	0.038291\\
1.9806	0.037535\\
1.9758	0.036792\\
1.971	0.036063\\
1.9662	0.035347\\
1.9614	0.034643\\
1.9567	0.033952\\
1.9519	0.033273\\
1.9471	0.032607\\
1.9423	0.031952\\
1.9375	0.03131\\
1.9327	0.030679\\
1.9279	0.03006\\
1.9232	0.029452\\
1.9184	0.028855\\
1.9136	0.028269\\
1.9088	0.027694\\
1.904	0.027129\\
1.8992	0.026575\\
1.8944	0.026031\\
1.8897	0.025498\\
1.8849	0.024974\\
1.8801	0.02446\\
1.8753	0.023956\\
1.8705	0.023461\\
1.8657	0.022976\\
1.8609	0.0225\\
1.8562	0.022032\\
1.8514	0.021574\\
1.8466	0.021125\\
1.8418	0.020684\\
1.837	0.020251\\
1.8322	0.019827\\
1.8274	0.019411\\
1.8227	0.019003\\
1.8179	0.018603\\
1.8131	0.01821\\
1.8083	0.017826\\
1.8035	0.017449\\
1.7987	0.017079\\
1.7939	0.016716\\
1.7892	0.016361\\
1.7844	0.016012\\
1.7796	0.01567\\
1.7748	0.015336\\
1.77	0.015007\\
1.7652	0.014686\\
1.7604	0.01437\\
1.7557	0.014061\\
1.7509	0.013758\\
1.7461	0.013461\\
1.7413	0.013171\\
1.7365	0.012886\\
1.7317	0.012606\\
1.7269	0.012333\\
1.7222	0.012065\\
1.7174	0.011802\\
1.7126	0.011545\\
1.7078	0.011293\\
1.703	0.011046\\
1.6982	0.010804\\
1.6934	0.010567\\
1.6887	0.010335\\
1.6839	0.010108\\
1.6791	0.0098849\\
1.6743	0.0096669\\
1.6695	0.0094534\\
1.6647	0.0092443\\
1.6599	0.0090396\\
1.6552	0.0088391\\
1.6504	0.0086428\\
1.6456	0.0084505\\
1.6408	0.0082623\\
1.636	0.0080781\\
1.6312	0.0078977\\
1.6265	0.0077211\\
1.6217	0.0075482\\
1.6169	0.007379\\
1.6121	0.0072133\\
1.6073	0.0070512\\
1.6025	0.0068925\\
1.5977	0.0067371\\
1.593	0.0065851\\
1.5882	0.0064363\\
1.5834	0.0062907\\
1.5786	0.0061482\\
1.5738	0.0060087\\
1.569	0.0058722\\
1.5642	0.0057387\\
1.5595	0.005608\\
1.5547	0.0054802\\
1.5499	0.0053551\\
1.5451	0.0052327\\
1.5403	0.005113\\
1.5355	0.0049959\\
1.5307	0.0048813\\
1.526	0.0047692\\
1.5212	0.0046595\\
1.5164	0.0045522\\
1.5116	0.0044473\\
1.5068	0.0043447\\
1.502	0.0042443\\
1.4972	0.0041462\\
1.4925	0.0040501\\
1.4877	0.0039562\\
1.4829	0.0038644\\
1.4781	0.0037746\\
1.4733	0.0036868\\
1.4685	0.0036009\\
1.4637	0.0035169\\
1.459	0.0034348\\
1.4542	0.0033546\\
1.4494	0.0032761\\
1.4446	0.0031993\\
1.4398	0.0031243\\
1.435	0.0030509\\
1.4302	0.0029792\\
1.4255	0.0029091\\
1.4207	0.0028406\\
1.4159	0.0027736\\
1.4111	0.0027081\\
1.4063	0.0026441\\
1.4015	0.0025816\\
1.3967	0.0025204\\
1.392	0.0024607\\
1.3872	0.0024023\\
1.3824	0.0023452\\
1.3776	0.0022894\\
1.3728	0.0022349\\
1.368	0.0021816\\
1.3632	0.0021295\\
1.3585	0.0020786\\
1.3537	0.0020289\\
1.3489	0.0019804\\
1.3441	0.0019329\\
1.3393	0.0018865\\
1.3345	0.0018412\\
1.3297	0.001797\\
1.325	0.0017537\\
1.3202	0.0017115\\
1.3154	0.0016702\\
1.3106	0.0016299\\
1.3058	0.0015905\\
1.301	0.001552\\
1.2962	0.0015145\\
1.2915	0.0014778\\
1.2867	0.0014419\\
1.2819	0.0014069\\
1.2771	0.0013727\\
1.2723	0.0013393\\
1.2675	0.0013067\\
1.2628	0.0012748\\
1.258	0.0012437\\
1.2532	0.0012133\\
1.2484	0.0011836\\
1.2436	0.0011547\\
1.2388	0.0011264\\
1.234	0.0010987\\
1.2293	0.0010718\\
1.2245	0.0010454\\
1.2197	0.0010197\\
1.2149	0.00099462\\
1.2101	0.00097011\\
1.2053	0.00094618\\
1.2005	0.00092282\\
1.1958	0.00090002\\
1.191	0.00087776\\
1.1862	0.00085603\\
1.1814	0.00083481\\
1.1766	0.00081411\\
1.1718	0.00079389\\
1.167	0.00077417\\
1.1623	0.00075491\\
1.1575	0.00073612\\
1.1527	0.00071778\\
1.1479	0.00069987\\
1.1431	0.0006824\\
1.1383	0.00066535\\
1.1335	0.00064872\\
1.1288	0.00063248\\
1.124	0.00061663\\
1.1192	0.00060117\\
1.1144	0.00058608\\
1.1096	0.00057136\\
1.1048	0.000557\\
1.1	0.00054298\\
1.0953	0.00052931\\
1.0905	0.00051597\\
1.0857	0.00050295\\
1.0809	0.00049025\\
1.0761	0.00047786\\
1.0713	0.00046578\\
1.0665	0.00045399\\
1.0618	0.00044249\\
1.057	0.00043127\\
1.0522	0.00042032\\
1.0474	0.00040965\\
1.0426	0.00039924\\
1.0378	0.00038908\\
1.033	0.00037917\\
1.0283	0.00036951\\
1.0235	0.00036009\\
1.0187	0.0003509\\
1.0139	0.00034193\\
1.0091	0.00033319\\
1.0043	0.00032467\\
0.99955	0.00031636\\
0.99476	0.00030825\\
0.98998	0.00030034\\
0.98519	0.00029264\\
0.9804	0.00028512\\
0.97562	0.00027779\\
0.97083	0.00027064\\
0.96605	0.00026368\\
0.96126	0.00025688\\
0.95648	0.00025026\\
0.95169	0.0002438\\
0.94691	0.0002375\\
0.94212	0.00023137\\
0.93733	0.00022538\\
0.93255	0.00021955\\
0.92776	0.00021386\\
0.92298	0.00020832\\
0.91819	0.00020292\\
0.91341	0.00019765\\
0.90862	0.00019251\\
0.90384	0.00018751\\
0.89905	0.00018263\\
0.89426	0.00017787\\
0.88948	0.00017324\\
0.88469	0.00016872\\
0.87991	0.00016432\\
0.87512	0.00016003\\
0.87034	0.00015585\\
0.86555	0.00015177\\
0.86077	0.0001478\\
0.85598	0.00014393\\
0.8512	0.00014016\\
0.84641	0.00013649\\
0.84162	0.0001329\\
0.83684	0.00012941\\
0.83205	0.00012601\\
0.82727	0.0001227\\
0.82248	0.00011947\\
0.8177	0.00011633\\
0.81291	0.00011326\\
0.80813	0.00011028\\
0.80334	0.00010737\\
0.79855	0.00010453\\
0.79377	0.00010177\\
0.78898	9.9079e-05\\
0.7842	9.6458e-05\\
0.77941	9.3905e-05\\
0.77463	9.1417e-05\\
0.76984	8.8994e-05\\
0.76506	8.6634e-05\\
0.76027	8.4334e-05\\
0.75548	8.2094e-05\\
0.7507	7.9913e-05\\
0.74591	7.7787e-05\\
0.74113	7.5717e-05\\
0.73634	7.3701e-05\\
0.73156	7.1737e-05\\
0.72677	6.9824e-05\\
0.72199	6.7962e-05\\
0.7172	6.6147e-05\\
0.71241	6.438e-05\\
0.70763	6.2659e-05\\
0.70284	6.0983e-05\\
0.69806	5.9351e-05\\
0.69327	5.7761e-05\\
0.68849	5.6213e-05\\
0.6837	5.4705e-05\\
0.67892	5.3237e-05\\
0.67413	5.1808e-05\\
0.66935	5.0416e-05\\
0.66456	4.906e-05\\
0.65977	4.774e-05\\
0.65499	4.6455e-05\\
0.6502	4.5204e-05\\
0.64542	4.3985e-05\\
0.64063	4.2799e-05\\
0.63585	4.1644e-05\\
0.63106	4.0519e-05\\
0.62628	3.9424e-05\\
0.62149	3.8358e-05\\
0.6167	3.732e-05\\
0.61192	3.631e-05\\
0.60713	3.5326e-05\\
0.60235	3.4369e-05\\
0.59756	3.3436e-05\\
0.59278	3.2529e-05\\
0.58799	3.1645e-05\\
0.58321	3.0785e-05\\
0.57842	2.9948e-05\\
0.57363	2.9133e-05\\
}--cycle;
\addplot[ybar interval, fill=black, fill opacity=0.35, area legend, draw=none] table[row sep=crcr, x=Lower, y=Count] {%
Lower	Upper	Count\\
0.57363	0.94138	0\\
0.94138	1.3091	0\\
1.3091	1.6769	0\\
1.6769	2.0446	0\\
2.0446	2.4124	0.079764\\
2.4124	2.7801	0.26105\\
2.7801	3.1479	0.74689\\
3.1479	3.5156	0.87016\\
3.5156	3.8834	0.49309\\
3.8834	4.2511	0.18128\\
4.2511	4.6189	0.065262\\
4.6189	4.9866	0.0072513\\
4.9866	5.3544	0.014503\\
5.3544	5.3544	0.014503\\
};
\addlegendentry{\RealHistogramLegend}

\addplot [color=black]
  table[row sep=crcr]{%
0.57363	6.6065e-10\\
1.6169	0.00037325\\
1.7413	0.0011889\\
1.8227	0.0024116\\
1.8849	0.0040324\\
1.9375	0.0061177\\
1.9854	0.0088083\\
2.0236	0.011674\\
2.0619	0.015337\\
2.0954	0.019333\\
2.1241	0.023452\\
2.1529	0.028307\\
2.1816	0.033998\\
2.2055	0.039456\\
2.2294	0.045633\\
2.2534	0.052595\\
2.2773	0.060412\\
2.3012	0.069152\\
2.3251	0.078885\\
2.3443	0.087432\\
2.3634	0.096692\\
2.3826	0.1067\\
2.4017	0.11748\\
2.4208	0.12907\\
2.44	0.14148\\
2.4591	0.15475\\
2.4783	0.1689\\
2.4974	0.18393\\
2.5166	0.19986\\
2.5357	0.21669\\
2.5548	0.23442\\
2.574	0.25305\\
2.5931	0.27256\\
2.6171	0.29815\\
2.641	0.32502\\
2.6649	0.3531\\
2.6888	0.38229\\
2.7176	0.41861\\
2.7463	0.45614\\
2.7845	0.50753\\
2.9042	0.6696\\
2.9329	0.7066\\
2.9568	0.73621\\
2.9808	0.76446\\
2.9999	0.7859\\
3.019	0.80618\\
3.0382	0.82519\\
3.0525	0.83854\\
3.0669	0.85106\\
3.0813	0.86272\\
3.0956	0.87347\\
3.11	0.88327\\
3.1243	0.8921\\
3.1387	0.89991\\
3.153	0.9067\\
3.1674	0.91242\\
3.1817	0.91706\\
3.1961	0.92062\\
3.2105	0.92307\\
3.2248	0.92441\\
3.2392	0.92463\\
3.2535	0.92374\\
3.2679	0.92175\\
3.2822	0.91866\\
3.2966	0.91448\\
3.311	0.90923\\
3.3253	0.90294\\
3.3397	0.89563\\
3.354	0.88733\\
3.3684	0.87807\\
3.3827	0.86788\\
3.3971	0.85681\\
3.4115	0.84489\\
3.4258	0.83217\\
3.445	0.81404\\
3.4641	0.79467\\
3.4832	0.7742\\
3.5024	0.75274\\
3.5263	0.72471\\
3.555	0.68964\\
3.5885	0.6473\\
3.6412	0.57917\\
3.6986	0.5051\\
3.7321	0.46307\\
3.7608	0.42817\\
3.7895	0.39457\\
3.8134	0.3677\\
3.8374	0.34196\\
3.8613	0.31743\\
3.8852	0.29414\\
3.9091	0.27216\\
3.9283	0.25551\\
3.9474	0.2397\\
3.9666	0.22474\\
3.9857	0.21061\\
4.0049	0.19729\\
4.024	0.18478\\
4.0431	0.17304\\
4.0623	0.16206\\
4.0814	0.15181\\
4.1054	0.13996\\
4.1293	0.12915\\
4.1532	0.1193\\
4.1771	0.11034\\
4.2011	0.10221\\
4.225	0.094836\\
4.2489	0.088161\\
4.2776	0.08098\\
4.3063	0.074603\\
4.3398	0.068051\\
4.3733	0.062319\\
4.4116	0.056608\\
4.4547	0.051058\\
4.5026	0.045747\\
4.5552	0.040697\\
4.6174	0.035518\\
4.6892	0.030307\\
4.7753	0.024846\\
4.871	0.019584\\
4.9667	0.015084\\
5.0625	0.011318\\
5.163	0.0081236\\
5.2682	0.0055474\\
5.3544	0.0039548\\
};
\addlegendentry{\RealBLNFitLegend}

      % \input{../../../../.tex/20/KS/SizeDistributionFittingsfig38.tex}
      \legend{}
  
      \nextgroupplot[%
        xmin=0.85621,xmax=5.4115,
        % xmin=0.62845,xmax=5.6393,
        % ymin=0,ymax=1,
      ] % This file was created by matlab2tikz.
%
\definecolor{mycolor1}{rgb}{0.83529,0.36863,0.00000}%
%
\addplot[ybar interval, fill=mycolor1, fill opacity=0.35, area legend, draw=none] table[row sep=crcr, x=Lower, y=Count] {%
Lower	Upper	Count\\
0.85621	1.2066	0.00053272\\
1.2066	1.557	0.0033485\\
1.557	1.9074	0.01621\\
1.9074	2.2578	0.066209\\
2.2578	2.6082	0.2105\\
2.6082	2.9586	0.49679\\
2.9586	3.3091	0.76597\\
3.3091	3.6595	0.67921\\
3.6595	4.0099	0.37153\\
4.0099	4.3603	0.16316\\
4.3603	4.7107	0.059284\\
4.7107	5.0611	0.013927\\
5.0611	5.4115	0.0071536\\
5.4115	5.4115	0.0071536\\
};
\addlegendentry{\ApproxHistogramLegend}

\addplot [color=mycolor1, only marks, every error bar/.append style={opacity=0.45}, mark=*, mark size=0pt, draw=none, forget plot]
 plot [error bars/.cd, y dir=both, y explicit, error bar style={line width=1pt, color=mycolor1}, error mark options={mark=none,mark size=0pt}]
 table[row sep=crcr, y error plus index=2, y error minus index=3]{%
1.0314	0.00053272	0.00044272	0.00044272\\
1.3818	0.0033485	0.00096749	0.00096749\\
1.7322	0.01621	0.002164	0.002164\\
2.0826	0.066209	0.0045605	0.0045605\\
2.433	0.2105	0.0078191	0.0078191\\
2.7834	0.49679	0.0093531	0.0093531\\
3.1339	0.76597	0.011309	0.011309\\
3.4843	0.67921	0.010319	0.010319\\
3.8347	0.37153	0.0081258	0.0081258\\
4.1851	0.16316	0.005594	0.005594\\
4.5355	0.059284	0.0034091	0.0034091\\
4.8859	0.013927	0.0016353	0.0016353\\
5.2363	0.0071536	0.0005605	0.0005605\\
};
\addplot [color=mycolor1]
  table[row sep=crcr]{%
0.85621	0.00021665\\
1.1344	0.00091795\\
1.294	0.0020248\\
1.4125	0.003574\\
1.5083	0.0055842\\
1.5858	0.0079457\\
1.6542	0.010775\\
1.7135	0.013959\\
1.7682	0.017651\\
1.8183	0.021803\\
1.8639	0.026334\\
1.905	0.031127\\
1.946	0.036695\\
1.9825	0.042377\\
2.019	0.04883\\
2.0554	0.056136\\
2.0874	0.063296\\
2.1193	0.071237\\
2.1512	0.080021\\
2.1786	0.088268\\
2.2059	0.09722\\
2.2333	0.10692\\
2.2606	0.1174\\
2.288	0.1287\\
2.3154	0.14087\\
2.3427	0.15392\\
2.3701	0.1679\\
2.3974	0.18282\\
2.4248	0.19872\\
2.4522	0.2156\\
2.4795	0.23349\\
2.5069	0.25237\\
2.5342	0.27225\\
2.5616	0.29311\\
2.5889	0.31493\\
2.6163	0.33768\\
2.6482	0.36532\\
2.6801	0.39406\\
2.7121	0.42377\\
2.7485	0.45872\\
2.7941	0.5035\\
2.9218	0.62975\\
2.9537	0.65999\\
2.9811	0.68494\\
3.0039	0.7048\\
3.0267	0.72359\\
3.0449	0.73768\\
3.0632	0.75079\\
3.0814	0.76277\\
3.0951	0.77094\\
3.1088	0.77833\\
3.1225	0.78492\\
3.1361	0.79065\\
3.1498	0.7955\\
3.1635	0.79945\\
3.1772	0.80249\\
3.1908	0.80464\\
3.2045	0.8059\\
3.2182	0.80629\\
3.2319	0.80585\\
3.2456	0.8046\\
3.2592	0.80258\\
3.2729	0.79981\\
3.2866	0.79634\\
3.3003	0.79218\\
3.314	0.78738\\
3.3276	0.78194\\
3.3459	0.77376\\
3.3641	0.76456\\
3.3824	0.7544\\
3.4006	0.74332\\
3.4188	0.7314\\
3.4371	0.7187\\
3.4599	0.70182\\
3.4827	0.68395\\
3.51	0.6614\\
3.5374	0.63787\\
3.5739	0.60538\\
3.624	0.55949\\
3.7152	0.47589\\
3.7563	0.4394\\
3.7927	0.40798\\
3.8247	0.38142\\
3.8566	0.35584\\
3.8885	0.3313\\
3.9204	0.30785\\
3.9478	0.28865\\
3.9751	0.27029\\
4.0025	0.25278\\
4.0299	0.23611\\
4.0572	0.22027\\
4.0846	0.20526\\
4.1119	0.19105\\
4.1393	0.17763\\
4.1667	0.16498\\
4.194	0.15307\\
4.2259	0.14007\\
4.2579	0.12801\\
4.2898	0.11683\\
4.3217	0.10648\\
4.3536	0.096925\\
4.3855	0.088113\\
4.422	0.078892\\
4.4585	0.070515\\
4.495	0.062915\\
4.536	0.055217\\
4.577	0.048342\\
4.6181	0.04221\\
4.6637	0.036184\\
4.7093	0.030901\\
4.7594	0.025855\\
4.8142	0.02116\\
4.8734	0.016907\\
4.9327	0.013408\\
5.0011	0.010165\\
5.0741	0.0074863\\
5.1561	0.0052418\\
5.2519	0.0034055\\
5.3659	0.001998\\
5.4115	0.0016055\\
};
\addlegendentry{\ApproxBLNFitLegend}


\addplot[area legend, draw=none, fill=mycolor1, fill opacity=0.15, forget plot]
table[row sep=crcr] {%
x	y\\
0.85621	0.00026525\\
0.86077	0.00027148\\
0.86533	0.00027786\\
0.86989	0.00028439\\
0.87445	0.00029107\\
0.87901	0.00029789\\
0.88357	0.00030488\\
0.88813	0.00031202\\
0.89269	0.00031933\\
0.89725	0.00032681\\
0.90181	0.00033445\\
0.90637	0.00034227\\
0.91093	0.00035027\\
0.91549	0.00035845\\
0.92005	0.00036681\\
0.92461	0.00037536\\
0.92917	0.00038411\\
0.93373	0.00039306\\
0.93829	0.00040221\\
0.94285	0.00041156\\
0.94741	0.00042113\\
0.95197	0.00043091\\
0.95653	0.00044092\\
0.96109	0.00045115\\
0.96565	0.00046161\\
0.97021	0.0004723\\
0.97477	0.00048324\\
0.97933	0.00049442\\
0.98389	0.00050585\\
0.98845	0.00051753\\
0.99301	0.00052948\\
0.99757	0.0005417\\
1.0021	0.00055419\\
1.0067	0.00056696\\
1.0112	0.00058001\\
1.0158	0.00059335\\
1.0204	0.00060699\\
1.0249	0.00062093\\
1.0295	0.00063519\\
1.034	0.00064976\\
1.0386	0.00066465\\
1.0432	0.00067987\\
1.0477	0.00069543\\
1.0523	0.00071133\\
1.0568	0.00072758\\
1.0614	0.00074419\\
1.066	0.00076117\\
1.0705	0.00077852\\
1.0751	0.00079625\\
1.0796	0.00081437\\
1.0842	0.00083289\\
1.0888	0.00085181\\
1.0933	0.00087115\\
1.0979	0.00089091\\
1.1024	0.0009111\\
1.107	0.00093173\\
1.1116	0.00095281\\
1.1161	0.00097435\\
1.1207	0.00099636\\
1.1252	0.0010188\\
1.1298	0.0010418\\
1.1344	0.0010653\\
1.1389	0.0010893\\
1.1435	0.0011138\\
1.148	0.0011388\\
1.1526	0.0011643\\
1.1572	0.0011905\\
1.1617	0.0012171\\
1.1663	0.0012444\\
1.1708	0.0012722\\
1.1754	0.0013007\\
1.18	0.0013297\\
1.1845	0.0013594\\
1.1891	0.0013896\\
1.1936	0.0014206\\
1.1982	0.0014522\\
1.2028	0.0014845\\
1.2073	0.0015174\\
1.2119	0.0015511\\
1.2164	0.0015854\\
1.221	0.0016205\\
1.2256	0.0016564\\
1.2301	0.001693\\
1.2347	0.0017303\\
1.2392	0.0017685\\
1.2438	0.0018075\\
1.2484	0.0018472\\
1.2529	0.0018878\\
1.2575	0.0019293\\
1.262	0.0019716\\
1.2666	0.0020149\\
1.2712	0.002059\\
1.2757	0.002104\\
1.2803	0.00215\\
1.2848	0.0021969\\
1.2894	0.0022448\\
1.294	0.0022937\\
1.2985	0.0023437\\
1.3031	0.0023946\\
1.3076	0.0024466\\
1.3122	0.0024997\\
1.3168	0.0025538\\
1.3213	0.0026091\\
1.3259	0.0026655\\
1.3304	0.0027231\\
1.335	0.0027819\\
1.3396	0.0028418\\
1.3441	0.002903\\
1.3487	0.0029654\\
1.3532	0.0030291\\
1.3578	0.0030941\\
1.3624	0.0031605\\
1.3669	0.0032281\\
1.3715	0.0032972\\
1.376	0.0033676\\
1.3806	0.0034395\\
1.3852	0.0035128\\
1.3897	0.0035876\\
1.3943	0.0036639\\
1.3988	0.0037417\\
1.4034	0.0038211\\
1.408	0.003902\\
1.4125	0.0039846\\
1.4171	0.0040689\\
1.4216	0.0041548\\
1.4262	0.0042425\\
1.4308	0.0043318\\
1.4353	0.004423\\
1.4399	0.004516\\
1.4444	0.0046108\\
1.449	0.0047075\\
1.4535	0.004806\\
1.4581	0.0049066\\
1.4627	0.0050091\\
1.4672	0.0051136\\
1.4718	0.0052202\\
1.4763	0.0053289\\
1.4809	0.0054397\\
1.4855	0.0055526\\
1.49	0.0056678\\
1.4946	0.0057852\\
1.4991	0.0059049\\
1.5037	0.0060269\\
1.5083	0.0061512\\
1.5128	0.006278\\
1.5174	0.0064072\\
1.5219	0.006539\\
1.5265	0.0066732\\
1.5311	0.0068101\\
1.5356	0.0069495\\
1.5402	0.0070917\\
1.5447	0.0072365\\
1.5493	0.0073842\\
1.5539	0.0075346\\
1.5584	0.0076879\\
1.563	0.0078441\\
1.5675	0.0080033\\
1.5721	0.0081655\\
1.5767	0.0083308\\
1.5812	0.0084992\\
1.5858	0.0086708\\
1.5903	0.0088456\\
1.5949	0.0090237\\
1.5995	0.0092051\\
1.604	0.0093899\\
1.6086	0.0095782\\
1.6131	0.00977\\
1.6177	0.0099653\\
1.6223	0.010164\\
1.6268	0.010367\\
1.6314	0.010573\\
1.6359	0.010784\\
1.6405	0.010998\\
1.6451	0.011216\\
1.6496	0.011438\\
1.6542	0.011664\\
1.6587	0.011895\\
1.6633	0.012129\\
1.6679	0.012368\\
1.6724	0.012611\\
1.677	0.012859\\
1.6815	0.013111\\
1.6861	0.013368\\
1.6907	0.013629\\
1.6952	0.013895\\
1.6998	0.014166\\
1.7043	0.014442\\
1.7089	0.014722\\
1.7135	0.015008\\
1.718	0.015299\\
1.7226	0.015595\\
1.7271	0.015896\\
1.7317	0.016203\\
1.7363	0.016515\\
1.7408	0.016833\\
1.7454	0.017156\\
1.7499	0.017485\\
1.7545	0.01782\\
1.7591	0.018161\\
1.7636	0.018507\\
1.7682	0.01886\\
1.7727	0.019219\\
1.7773	0.019584\\
1.7819	0.019956\\
1.7864	0.020334\\
1.791	0.020718\\
1.7955	0.02111\\
1.8001	0.021508\\
1.8047	0.021912\\
1.8092	0.022324\\
1.8138	0.022743\\
1.8183	0.023169\\
1.8229	0.023602\\
1.8275	0.024043\\
1.832	0.024491\\
1.8366	0.024947\\
1.8411	0.025411\\
1.8457	0.025882\\
1.8503	0.026361\\
1.8548	0.026849\\
1.8594	0.027344\\
1.8639	0.027848\\
1.8685	0.028361\\
1.8731	0.028881\\
1.8776	0.029411\\
1.8822	0.029949\\
1.8867	0.030497\\
1.8913	0.031053\\
1.8959	0.031618\\
1.9004	0.032193\\
1.905	0.032777\\
1.9095	0.033371\\
1.9141	0.033975\\
1.9187	0.034588\\
1.9232	0.035211\\
1.9278	0.035845\\
1.9323	0.036488\\
1.9369	0.037142\\
1.9415	0.037807\\
1.946	0.038482\\
1.9506	0.039168\\
1.9551	0.039865\\
1.9597	0.040573\\
1.9643	0.041293\\
1.9688	0.042023\\
1.9734	0.042766\\
1.9779	0.04352\\
1.9825	0.044286\\
1.9871	0.045064\\
1.9916	0.045854\\
1.9962	0.046656\\
2.0007	0.047471\\
2.0053	0.048299\\
2.0098	0.049139\\
2.0144	0.049993\\
2.019	0.050859\\
2.0235	0.051739\\
2.0281	0.052632\\
2.0326	0.053539\\
2.0372	0.05446\\
2.0418	0.055395\\
2.0463	0.056344\\
2.0509	0.057307\\
2.0554	0.058285\\
2.06	0.059277\\
2.0646	0.060284\\
2.0691	0.061307\\
2.0737	0.062344\\
2.0782	0.063397\\
2.0828	0.064465\\
2.0874	0.06555\\
2.0919	0.06665\\
2.0965	0.067766\\
2.101	0.068898\\
2.1056	0.070047\\
2.1102	0.071213\\
2.1147	0.072396\\
2.1193	0.073595\\
2.1238	0.074812\\
2.1284	0.076046\\
2.133	0.077298\\
2.1375	0.078568\\
2.1421	0.079855\\
2.1466	0.081161\\
2.1512	0.082485\\
2.1558	0.083828\\
2.1603	0.08519\\
2.1649	0.08657\\
2.1694	0.08797\\
2.174	0.089389\\
2.1786	0.090828\\
2.1831	0.092286\\
2.1877	0.093764\\
2.1922	0.095263\\
2.1968	0.096781\\
2.2014	0.098321\\
2.2059	0.099881\\
2.2105	0.10146\\
2.215	0.10306\\
2.2196	0.10469\\
2.2242	0.10633\\
2.2287	0.108\\
2.2333	0.10969\\
2.2378	0.1114\\
2.2424	0.11313\\
2.247	0.11488\\
2.2515	0.11666\\
2.2561	0.11846\\
2.2606	0.12029\\
2.2652	0.12213\\
2.2698	0.124\\
2.2743	0.1259\\
2.2789	0.12782\\
2.2834	0.12976\\
2.288	0.13172\\
2.2926	0.13371\\
2.2971	0.13573\\
2.3017	0.13777\\
2.3062	0.13983\\
2.3108	0.14192\\
2.3154	0.14403\\
2.3199	0.14617\\
2.3245	0.14834\\
2.329	0.15053\\
2.3336	0.15274\\
2.3382	0.15499\\
2.3427	0.15725\\
2.3473	0.15955\\
2.3518	0.16187\\
2.3564	0.16422\\
2.361	0.16659\\
2.3655	0.16899\\
2.3701	0.17142\\
2.3746	0.17387\\
2.3792	0.17635\\
2.3838	0.17886\\
2.3883	0.1814\\
2.3929	0.18396\\
2.3974	0.18655\\
2.402	0.18917\\
2.4066	0.19182\\
2.4111	0.19449\\
2.4157	0.1972\\
2.4202	0.19993\\
2.4248	0.20268\\
2.4294	0.20547\\
2.4339	0.20828\\
2.4385	0.21113\\
2.443	0.214\\
2.4476	0.2169\\
2.4522	0.21982\\
2.4567	0.22278\\
2.4613	0.22576\\
2.4658	0.22878\\
2.4704	0.23182\\
2.475	0.23489\\
2.4795	0.23798\\
2.4841	0.24111\\
2.4886	0.24426\\
2.4932	0.24745\\
2.4978	0.25066\\
2.5023	0.2539\\
2.5069	0.25716\\
2.5114	0.26046\\
2.516	0.26378\\
2.5206	0.26713\\
2.5251	0.27051\\
2.5297	0.27392\\
2.5342	0.27735\\
2.5388	0.28081\\
2.5434	0.2843\\
2.5479	0.28782\\
2.5525	0.29136\\
2.557	0.29494\\
2.5616	0.29853\\
2.5662	0.30216\\
2.5707	0.30581\\
2.5753	0.30949\\
2.5798	0.31319\\
2.5844	0.31692\\
2.5889	0.32068\\
2.5935	0.32446\\
2.5981	0.32827\\
2.6026	0.3321\\
2.6072	0.33596\\
2.6117	0.33984\\
2.6163	0.34374\\
2.6209	0.34767\\
2.6254	0.35163\\
2.63	0.35561\\
2.6345	0.35961\\
2.6391	0.36363\\
2.6437	0.36768\\
2.6482	0.37174\\
2.6528	0.37583\\
2.6573	0.37995\\
2.6619	0.38408\\
2.6665	0.38823\\
2.671	0.39241\\
2.6756	0.3966\\
2.6801	0.40081\\
2.6847	0.40504\\
2.6893	0.4093\\
2.6938	0.41356\\
2.6984	0.41785\\
2.7029	0.42216\\
2.7075	0.42648\\
2.7121	0.43081\\
2.7166	0.43517\\
2.7212	0.43953\\
2.7257	0.44392\\
2.7303	0.44831\\
2.7349	0.45273\\
2.7394	0.45715\\
2.744	0.46158\\
2.7485	0.46603\\
2.7531	0.47049\\
2.7577	0.47496\\
2.7622	0.47944\\
2.7668	0.48393\\
2.7713	0.48843\\
2.7759	0.49293\\
2.7805	0.49745\\
2.785	0.50197\\
2.7896	0.50649\\
2.7941	0.51103\\
2.7987	0.51556\\
2.8033	0.5201\\
2.8078	0.52465\\
2.8124	0.52919\\
2.8169	0.53374\\
2.8215	0.53829\\
2.8261	0.54284\\
2.8306	0.54739\\
2.8352	0.55193\\
2.8397	0.55648\\
2.8443	0.56102\\
2.8489	0.56556\\
2.8534	0.57009\\
2.858	0.57462\\
2.8625	0.57914\\
2.8671	0.58366\\
2.8717	0.58817\\
2.8762	0.59267\\
2.8808	0.59716\\
2.8853	0.60164\\
2.8899	0.60611\\
2.8945	0.61057\\
2.899	0.61501\\
2.9036	0.61945\\
2.9081	0.62386\\
2.9127	0.62827\\
2.9173	0.63266\\
2.9218	0.63703\\
2.9264	0.64139\\
2.9309	0.64573\\
2.9355	0.65005\\
2.9401	0.65435\\
2.9446	0.65863\\
2.9492	0.66289\\
2.9537	0.66713\\
2.9583	0.67135\\
2.9629	0.67555\\
2.9674	0.67972\\
2.972	0.68387\\
2.9765	0.688\\
2.9811	0.6921\\
2.9857	0.69617\\
2.9902	0.70021\\
2.9948	0.70423\\
2.9993	0.70821\\
3.0039	0.71217\\
3.0085	0.71608\\
3.013	0.71997\\
3.0176	0.72381\\
3.0221	0.72762\\
3.0267	0.73138\\
3.0313	0.7351\\
3.0358	0.73877\\
3.0404	0.74239\\
3.0449	0.74596\\
3.0495	0.74947\\
3.0541	0.75292\\
3.0586	0.75631\\
3.0632	0.75964\\
3.0677	0.7629\\
3.0723	0.76609\\
3.0769	0.7692\\
3.0814	0.77223\\
3.086	0.77519\\
3.0905	0.77806\\
3.0951	0.78084\\
3.0997	0.78353\\
3.1042	0.78613\\
3.1088	0.78863\\
3.1133	0.79103\\
3.1179	0.79334\\
3.1225	0.79554\\
3.127	0.79764\\
3.1316	0.79963\\
3.1361	0.80151\\
3.1407	0.80328\\
3.1453	0.80494\\
3.1498	0.8065\\
3.1544	0.80794\\
3.1589	0.80927\\
3.1635	0.81048\\
3.168	0.81159\\
3.1726	0.81259\\
3.1772	0.81348\\
3.1817	0.81425\\
3.1863	0.81493\\
3.1908	0.81549\\
3.1954	0.81595\\
3.2	0.81631\\
3.2045	0.81657\\
3.2091	0.81673\\
3.2136	0.81679\\
3.2182	0.81675\\
3.2228	0.81663\\
3.2273	0.81641\\
3.2319	0.8161\\
3.2364	0.8157\\
3.241	0.81522\\
3.2456	0.81465\\
3.2501	0.814\\
3.2547	0.81326\\
3.2592	0.81245\\
3.2638	0.81156\\
3.2684	0.81059\\
3.2729	0.80954\\
3.2775	0.80842\\
3.282	0.80722\\
3.2866	0.80594\\
3.2912	0.8046\\
3.2957	0.80318\\
3.3003	0.80168\\
3.3048	0.80012\\
3.3094	0.79848\\
3.314	0.79678\\
3.3185	0.795\\
3.3231	0.79315\\
3.3276	0.79124\\
3.3322	0.78925\\
3.3368	0.7872\\
3.3413	0.78509\\
3.3459	0.7829\\
3.3504	0.78066\\
3.355	0.77834\\
3.3596	0.77597\\
3.3641	0.77353\\
3.3687	0.77103\\
3.3732	0.76847\\
3.3778	0.76585\\
3.3824	0.76317\\
3.3869	0.76043\\
3.3915	0.75763\\
3.396	0.75478\\
3.4006	0.75187\\
3.4052	0.74891\\
3.4097	0.7459\\
3.4143	0.74283\\
3.4188	0.73971\\
3.4234	0.73655\\
3.428	0.73333\\
3.4325	0.73007\\
3.4371	0.72676\\
3.4416	0.72341\\
3.4462	0.72001\\
3.4508	0.71657\\
3.4553	0.71309\\
3.4599	0.70957\\
3.4644	0.70601\\
3.469	0.70241\\
3.4736	0.69878\\
3.4781	0.69511\\
3.4827	0.6914\\
3.4872	0.68766\\
3.4918	0.68389\\
3.4964	0.68009\\
3.5009	0.67626\\
3.5055	0.6724\\
3.51	0.66851\\
3.5146	0.6646\\
3.5192	0.66066\\
3.5237	0.6567\\
3.5283	0.65272\\
3.5328	0.64871\\
3.5374	0.64468\\
3.542	0.64064\\
3.5465	0.63657\\
3.5511	0.63249\\
3.5556	0.62839\\
3.5602	0.62427\\
3.5648	0.62014\\
3.5693	0.616\\
3.5739	0.61185\\
3.5784	0.60768\\
3.583	0.6035\\
3.5876	0.59931\\
3.5921	0.59512\\
3.5967	0.59091\\
3.6012	0.5867\\
3.6058	0.58249\\
3.6104	0.57826\\
3.6149	0.57404\\
3.6195	0.56981\\
3.624	0.56557\\
3.6286	0.56134\\
3.6332	0.5571\\
3.6377	0.55286\\
3.6423	0.54863\\
3.6468	0.54439\\
3.6514	0.54016\\
3.656	0.53592\\
3.6605	0.53169\\
3.6651	0.52747\\
3.6696	0.52325\\
3.6742	0.51903\\
3.6788	0.51482\\
3.6833	0.51062\\
3.6879	0.50642\\
3.6924	0.50223\\
3.697	0.49804\\
3.7016	0.49387\\
3.7061	0.48971\\
3.7107	0.48555\\
3.7152	0.48141\\
3.7198	0.47727\\
3.7243	0.47315\\
3.7289	0.46904\\
3.7335	0.46494\\
3.738	0.46085\\
3.7426	0.45678\\
3.7471	0.45272\\
3.7517	0.44868\\
3.7563	0.44464\\
3.7608	0.44063\\
3.7654	0.43663\\
3.7699	0.43264\\
3.7745	0.42867\\
3.7791	0.42472\\
3.7836	0.42078\\
3.7882	0.41686\\
3.7927	0.41296\\
3.7973	0.40908\\
3.8019	0.40521\\
3.8064	0.40136\\
3.811	0.39754\\
3.8155	0.39373\\
3.8201	0.38994\\
3.8247	0.38616\\
3.8292	0.38241\\
3.8338	0.37868\\
3.8383	0.37497\\
3.8429	0.37128\\
3.8475	0.36761\\
3.852	0.36396\\
3.8566	0.36033\\
3.8611	0.35673\\
3.8657	0.35314\\
3.8703	0.34958\\
3.8748	0.34604\\
3.8794	0.34252\\
3.8839	0.33902\\
3.8885	0.33555\\
3.8931	0.33209\\
3.8976	0.32866\\
3.9022	0.32526\\
3.9067	0.32187\\
3.9113	0.31851\\
3.9159	0.31517\\
3.9204	0.31185\\
3.925	0.30856\\
3.9295	0.30529\\
3.9341	0.30205\\
3.9387	0.29883\\
3.9432	0.29563\\
3.9478	0.29245\\
3.9523	0.2893\\
3.9569	0.28617\\
3.9615	0.28307\\
3.966	0.27999\\
3.9706	0.27693\\
3.9751	0.2739\\
3.9797	0.27089\\
3.9843	0.2679\\
3.9888	0.26494\\
3.9934	0.262\\
3.9979	0.25909\\
4.0025	0.25619\\
4.0071	0.25333\\
4.0116	0.25048\\
4.0162	0.24766\\
4.0207	0.24487\\
4.0253	0.2421\\
4.0299	0.23935\\
4.0344	0.23662\\
4.039	0.23392\\
4.0435	0.23124\\
4.0481	0.22858\\
4.0527	0.22595\\
4.0572	0.22334\\
4.0618	0.22076\\
4.0663	0.2182\\
4.0709	0.21566\\
4.0755	0.21314\\
4.08	0.21065\\
4.0846	0.20818\\
4.0891	0.20573\\
4.0937	0.2033\\
4.0983	0.2009\\
4.1028	0.19852\\
4.1074	0.19616\\
4.1119	0.19382\\
4.1165	0.19151\\
4.1211	0.18922\\
4.1256	0.18695\\
4.1302	0.1847\\
4.1347	0.18247\\
4.1393	0.18027\\
4.1439	0.17809\\
4.1484	0.17592\\
4.153	0.17378\\
4.1575	0.17166\\
4.1621	0.16957\\
4.1667	0.16749\\
4.1712	0.16543\\
4.1758	0.1634\\
4.1803	0.16138\\
4.1849	0.15938\\
4.1895	0.15741\\
4.194	0.15546\\
4.1986	0.15352\\
4.2031	0.15161\\
4.2077	0.14971\\
4.2123	0.14784\\
4.2168	0.14598\\
4.2214	0.14414\\
4.2259	0.14233\\
4.2305	0.14053\\
4.2351	0.13875\\
4.2396	0.13699\\
4.2442	0.13525\\
4.2487	0.13352\\
4.2533	0.13182\\
4.2579	0.13013\\
4.2624	0.12846\\
4.267	0.12681\\
4.2715	0.12518\\
4.2761	0.12356\\
4.2807	0.12197\\
4.2852	0.12039\\
4.2898	0.11882\\
4.2943	0.11728\\
4.2989	0.11575\\
4.3034	0.11424\\
4.308	0.11274\\
4.3126	0.11127\\
4.3171	0.1098\\
4.3217	0.10836\\
4.3262	0.10693\\
4.3308	0.10552\\
4.3354	0.10412\\
4.3399	0.10274\\
4.3445	0.10137\\
4.349	0.10002\\
4.3536	0.098688\\
4.3582	0.097369\\
4.3627	0.096065\\
4.3673	0.094775\\
4.3718	0.093501\\
4.3764	0.092241\\
4.381	0.090996\\
4.3855	0.089765\\
4.3901	0.088549\\
4.3946	0.087347\\
4.3992	0.086159\\
4.4038	0.084985\\
4.4083	0.083824\\
4.4129	0.082678\\
4.4174	0.081545\\
4.422	0.080425\\
4.4266	0.079319\\
4.4311	0.078226\\
4.4357	0.077146\\
4.4402	0.076079\\
4.4448	0.075024\\
4.4494	0.073983\\
4.4539	0.072954\\
4.4585	0.071937\\
4.463	0.070933\\
4.4676	0.069941\\
4.4722	0.068961\\
4.4767	0.067993\\
4.4813	0.067037\\
4.4858	0.066093\\
4.4904	0.06516\\
4.495	0.064238\\
4.4995	0.063328\\
4.5041	0.06243\\
4.5086	0.061542\\
4.5132	0.060665\\
4.5178	0.0598\\
4.5223	0.058945\\
4.5269	0.0581\\
4.5314	0.057266\\
4.536	0.056443\\
4.5406	0.05563\\
4.5451	0.054827\\
4.5497	0.054034\\
4.5542	0.053251\\
4.5588	0.052478\\
4.5634	0.051715\\
4.5679	0.050962\\
4.5725	0.050217\\
4.577	0.049483\\
4.5816	0.048757\\
4.5862	0.048041\\
4.5907	0.047334\\
4.5953	0.046635\\
4.5998	0.045946\\
4.6044	0.045265\\
4.609	0.044593\\
4.6135	0.04393\\
4.6181	0.043275\\
4.6226	0.042628\\
4.6272	0.04199\\
4.6318	0.041359\\
4.6363	0.040737\\
4.6409	0.040123\\
4.6454	0.039516\\
4.65	0.038918\\
4.6546	0.038327\\
4.6591	0.037743\\
4.6637	0.037167\\
4.6682	0.036598\\
4.6728	0.036037\\
4.6774	0.035483\\
4.6819	0.034936\\
4.6865	0.034396\\
4.691	0.033863\\
4.6956	0.033337\\
4.7002	0.032818\\
4.7047	0.032305\\
4.7093	0.031799\\
4.7138	0.0313\\
4.7184	0.030807\\
4.723	0.030321\\
4.7275	0.029841\\
4.7321	0.029368\\
4.7366	0.0289\\
4.7412	0.028439\\
4.7458	0.027984\\
4.7503	0.027535\\
4.7549	0.027092\\
4.7594	0.026654\\
4.764	0.026223\\
4.7686	0.025797\\
4.7731	0.025378\\
4.7777	0.024964\\
4.7822	0.024555\\
4.7868	0.024152\\
4.7914	0.023755\\
4.7959	0.023363\\
4.8005	0.022976\\
4.805	0.022595\\
4.8096	0.022219\\
4.8142	0.021848\\
4.8187	0.021483\\
4.8233	0.021123\\
4.8278	0.020768\\
4.8324	0.020417\\
4.837	0.020072\\
4.8415	0.019732\\
4.8461	0.019397\\
4.8506	0.019067\\
4.8552	0.018741\\
4.8597	0.01842\\
4.8643	0.018104\\
4.8689	0.017793\\
4.8734	0.017486\\
4.878	0.017184\\
4.8825	0.016886\\
4.8871	0.016593\\
4.8917	0.016304\\
4.8962	0.01602\\
4.9008	0.01574\\
4.9053	0.015465\\
4.9099	0.015193\\
4.9145	0.014926\\
4.919	0.014663\\
4.9236	0.014404\\
4.9281	0.014149\\
4.9327	0.013899\\
4.9373	0.013652\\
4.9418	0.013409\\
4.9464	0.01317\\
4.9509	0.012935\\
4.9555	0.012703\\
4.9601	0.012476\\
4.9646	0.012252\\
4.9692	0.012031\\
4.9737	0.011815\\
4.9783	0.011601\\
4.9829	0.011392\\
4.9874	0.011185\\
4.992	0.010982\\
4.9965	0.010783\\
5.0011	0.010586\\
5.0057	0.010393\\
5.0102	0.010204\\
5.0148	0.010017\\
5.0193	0.0098333\\
5.0239	0.0096528\\
5.0285	0.0094753\\
5.033	0.0093008\\
5.0376	0.0091292\\
5.0421	0.0089605\\
5.0467	0.0087946\\
5.0513	0.0086315\\
5.0558	0.0084712\\
5.0604	0.0083135\\
5.0649	0.0081585\\
5.0695	0.0080062\\
5.0741	0.0078564\\
5.0786	0.0077092\\
5.0832	0.0075645\\
5.0877	0.0074222\\
5.0923	0.0072824\\
5.0969	0.007145\\
5.1014	0.0070099\\
5.106	0.0068771\\
5.1105	0.0067466\\
5.1151	0.0066184\\
5.1197	0.0064923\\
5.1242	0.0063685\\
5.1288	0.0062468\\
5.1333	0.0061271\\
5.1379	0.0060096\\
5.1425	0.0058941\\
5.147	0.0057806\\
5.1516	0.0056691\\
5.1561	0.0055595\\
5.1607	0.0054519\\
5.1653	0.0053461\\
5.1698	0.0052422\\
5.1744	0.0051402\\
5.1789	0.0050399\\
5.1835	0.0049414\\
5.1881	0.0048447\\
5.1926	0.0047497\\
5.1972	0.0046563\\
5.2017	0.0045647\\
5.2063	0.0044747\\
5.2109	0.0043863\\
5.2154	0.0042995\\
5.22	0.0042142\\
5.2245	0.0041305\\
5.2291	0.0040483\\
5.2337	0.0039677\\
5.2382	0.0038884\\
5.2428	0.0038107\\
5.2473	0.0037343\\
5.2519	0.0036594\\
5.2565	0.0035858\\
5.261	0.0035136\\
5.2656	0.0034428\\
5.2701	0.0033732\\
5.2747	0.003305\\
5.2793	0.003238\\
5.2838	0.0031722\\
5.2884	0.0031078\\
5.2929	0.0030445\\
5.2975	0.0029824\\
5.3021	0.0029215\\
5.3066	0.0028617\\
5.3112	0.0028031\\
5.3157	0.0027456\\
5.3203	0.0026892\\
5.3249	0.0026339\\
5.3294	0.0025796\\
5.334	0.0025264\\
5.3385	0.0024742\\
5.3431	0.002423\\
5.3477	0.0023728\\
5.3522	0.0023236\\
5.3568	0.0022754\\
5.3613	0.0022281\\
5.3659	0.0021817\\
5.3705	0.0021362\\
5.375	0.0020916\\
5.3796	0.0020479\\
5.3841	0.0020051\\
5.3887	0.0019631\\
5.3933	0.001922\\
5.3978	0.0018816\\
5.4024	0.0018421\\
5.4069	0.0018034\\
5.4115	0.0017654\\
5.4115	0.0014456\\
5.4069	0.0014791\\
5.4024	0.0015132\\
5.3978	0.0015481\\
5.3933	0.0015838\\
5.3887	0.0016202\\
5.3841	0.0016574\\
5.3796	0.0016954\\
5.375	0.0017342\\
5.3705	0.0017738\\
5.3659	0.0018143\\
5.3613	0.0018556\\
5.3568	0.0018978\\
5.3522	0.0019408\\
5.3477	0.0019848\\
5.3431	0.0020297\\
5.3385	0.0020756\\
5.334	0.0021224\\
5.3294	0.0021702\\
5.3249	0.002219\\
5.3203	0.0022688\\
5.3157	0.0023196\\
5.3112	0.0023715\\
5.3066	0.0024245\\
5.3021	0.0024786\\
5.2975	0.0025337\\
5.2929	0.0025901\\
5.2884	0.0026475\\
5.2838	0.0027062\\
5.2793	0.002766\\
5.2747	0.0028271\\
5.2701	0.0028894\\
5.2656	0.002953\\
5.261	0.0030179\\
5.2565	0.003084\\
5.2519	0.0031516\\
5.2473	0.0032205\\
5.2428	0.0032907\\
5.2382	0.0033624\\
5.2337	0.0034355\\
5.2291	0.0035101\\
5.2245	0.0035862\\
5.22	0.0036637\\
5.2154	0.0037429\\
5.2109	0.0038235\\
5.2063	0.0039058\\
5.2017	0.0039897\\
5.1972	0.0040753\\
5.1926	0.0041625\\
5.1881	0.0042514\\
5.1835	0.0043421\\
5.1789	0.0044345\\
5.1744	0.0045287\\
5.1698	0.0046247\\
5.1653	0.0047226\\
5.1607	0.0048224\\
5.1561	0.004924\\
5.1516	0.0050277\\
5.147	0.0051332\\
5.1425	0.0052408\\
5.1379	0.0053504\\
5.1333	0.0054621\\
5.1288	0.0055759\\
5.1242	0.0056918\\
5.1197	0.0058098\\
5.1151	0.0059301\\
5.1105	0.0060526\\
5.106	0.0061773\\
5.1014	0.0063043\\
5.0969	0.0064336\\
5.0923	0.0065653\\
5.0877	0.0066994\\
5.0832	0.0068358\\
5.0786	0.0069748\\
5.0741	0.0071162\\
5.0695	0.0072601\\
5.0649	0.0074066\\
5.0604	0.0075556\\
5.0558	0.0077073\\
5.0513	0.0078616\\
5.0467	0.0080186\\
5.0421	0.0081783\\
5.0376	0.0083408\\
5.033	0.008506\\
5.0285	0.0086741\\
5.0239	0.008845\\
5.0193	0.0090188\\
5.0148	0.0091955\\
5.0102	0.0093751\\
5.0057	0.0095578\\
5.0011	0.0097435\\
4.9965	0.0099322\\
4.992	0.010124\\
4.9874	0.010319\\
4.9829	0.010517\\
4.9783	0.010718\\
4.9737	0.010923\\
4.9692	0.011131\\
4.9646	0.011342\\
4.9601	0.011557\\
4.9555	0.011775\\
4.9509	0.011996\\
4.9464	0.012221\\
4.9418	0.012449\\
4.9373	0.012681\\
4.9327	0.012917\\
4.9281	0.013156\\
4.9236	0.013399\\
4.919	0.013646\\
4.9145	0.013896\\
4.9099	0.01415\\
4.9053	0.014409\\
4.9008	0.014671\\
4.8962	0.014937\\
4.8917	0.015207\\
4.8871	0.015481\\
4.8825	0.01576\\
4.878	0.016042\\
4.8734	0.016329\\
4.8689	0.01662\\
4.8643	0.016916\\
4.8597	0.017216\\
4.8552	0.01752\\
4.8506	0.017829\\
4.8461	0.018142\\
4.8415	0.018461\\
4.837	0.018783\\
4.8324	0.019111\\
4.8278	0.019444\\
4.8233	0.019781\\
4.8187	0.020123\\
4.8142	0.020471\\
4.8096	0.020823\\
4.805	0.021181\\
4.8005	0.021544\\
4.7959	0.021912\\
4.7914	0.022285\\
4.7868	0.022664\\
4.7822	0.023049\\
4.7777	0.023439\\
4.7731	0.023834\\
4.7686	0.024236\\
4.764	0.024643\\
4.7594	0.025056\\
4.7549	0.025475\\
4.7503	0.025899\\
4.7458	0.02633\\
4.7412	0.026767\\
4.7366	0.02721\\
4.7321	0.02766\\
4.7275	0.028115\\
4.723	0.028577\\
4.7184	0.029046\\
4.7138	0.029521\\
4.7093	0.030003\\
4.7047	0.030491\\
4.7002	0.030987\\
4.6956	0.031489\\
4.691	0.031998\\
4.6865	0.032514\\
4.6819	0.033037\\
4.6774	0.033567\\
4.6728	0.034104\\
4.6682	0.034649\\
4.6637	0.035201\\
4.6591	0.035761\\
4.6546	0.036328\\
4.65	0.036903\\
4.6454	0.037485\\
4.6409	0.038075\\
4.6363	0.038673\\
4.6318	0.039279\\
4.6272	0.039893\\
4.6226	0.040515\\
4.6181	0.041146\\
4.6135	0.041784\\
4.609	0.042431\\
4.6044	0.043086\\
4.5998	0.04375\\
4.5953	0.044423\\
4.5907	0.045104\\
4.5862	0.045794\\
4.5816	0.046493\\
4.577	0.047201\\
4.5725	0.047918\\
4.5679	0.048644\\
4.5634	0.049379\\
4.5588	0.050124\\
4.5542	0.050878\\
4.5497	0.051642\\
4.5451	0.052416\\
4.5406	0.053199\\
4.536	0.053992\\
4.5314	0.054795\\
4.5269	0.055608\\
4.5223	0.056431\\
4.5178	0.057265\\
4.5132	0.058109\\
4.5086	0.058964\\
4.5041	0.059829\\
4.4995	0.060705\\
4.495	0.061592\\
4.4904	0.06249\\
4.4858	0.063399\\
4.4813	0.064319\\
4.4767	0.06525\\
4.4722	0.066193\\
4.4676	0.067148\\
4.463	0.068114\\
4.4585	0.069092\\
4.4539	0.070083\\
4.4494	0.071085\\
4.4448	0.072099\\
4.4402	0.073126\\
4.4357	0.074165\\
4.4311	0.075217\\
4.4266	0.076282\\
4.422	0.07736\\
4.4174	0.07845\\
4.4129	0.079554\\
4.4083	0.080671\\
4.4038	0.081801\\
4.3992	0.082945\\
4.3946	0.084103\\
4.3901	0.085275\\
4.3855	0.08646\\
4.381	0.08766\\
4.3764	0.088874\\
4.3718	0.090102\\
4.3673	0.091345\\
4.3627	0.092602\\
4.3582	0.093874\\
4.3536	0.095162\\
4.349	0.096464\\
4.3445	0.097781\\
4.3399	0.099114\\
4.3354	0.10046\\
4.3308	0.10183\\
4.3262	0.10321\\
4.3217	0.1046\\
4.3171	0.10601\\
4.3126	0.10744\\
4.308	0.10888\\
4.3034	0.11035\\
4.2989	0.11182\\
4.2943	0.11332\\
4.2898	0.11483\\
4.2852	0.11635\\
4.2807	0.1179\\
4.2761	0.11946\\
4.2715	0.12104\\
4.267	0.12264\\
4.2624	0.12425\\
4.2579	0.12589\\
4.2533	0.12754\\
4.2487	0.1292\\
4.2442	0.13089\\
4.2396	0.1326\\
4.2351	0.13432\\
4.2305	0.13606\\
4.2259	0.13782\\
4.2214	0.1396\\
4.2168	0.1414\\
4.2123	0.14322\\
4.2077	0.14506\\
4.2031	0.14691\\
4.1986	0.14879\\
4.194	0.15068\\
4.1895	0.1526\\
4.1849	0.15453\\
4.1803	0.15649\\
4.1758	0.15846\\
4.1712	0.16046\\
4.1667	0.16247\\
4.1621	0.16451\\
4.1575	0.16656\\
4.153	0.16864\\
4.1484	0.17074\\
4.1439	0.17286\\
4.1393	0.175\\
4.1347	0.17716\\
4.1302	0.17934\\
4.1256	0.18154\\
4.1211	0.18377\\
4.1165	0.18601\\
4.1119	0.18828\\
4.1074	0.19057\\
4.1028	0.19288\\
4.0983	0.19521\\
4.0937	0.19756\\
4.0891	0.19994\\
4.0846	0.20234\\
4.08	0.20476\\
4.0755	0.2072\\
4.0709	0.20967\\
4.0663	0.21215\\
4.0618	0.21466\\
4.0572	0.2172\\
4.0527	0.21975\\
4.0481	0.22233\\
4.0435	0.22493\\
4.039	0.22755\\
4.0344	0.2302\\
4.0299	0.23287\\
4.0253	0.23556\\
4.0207	0.23827\\
4.0162	0.24101\\
4.0116	0.24377\\
4.0071	0.24656\\
4.0025	0.24936\\
3.9979	0.25219\\
3.9934	0.25505\\
3.9888	0.25792\\
3.9843	0.26082\\
3.9797	0.26374\\
3.9751	0.26669\\
3.9706	0.26966\\
3.966	0.27265\\
3.9615	0.27567\\
3.9569	0.27871\\
3.9523	0.28177\\
3.9478	0.28485\\
3.9432	0.28796\\
3.9387	0.29109\\
3.9341	0.29425\\
3.9295	0.29742\\
3.925	0.30062\\
3.9204	0.30385\\
3.9159	0.30709\\
3.9113	0.31036\\
3.9067	0.31365\\
3.9022	0.31697\\
3.8976	0.32031\\
3.8931	0.32367\\
3.8885	0.32705\\
3.8839	0.33045\\
3.8794	0.33388\\
3.8748	0.33733\\
3.8703	0.3408\\
3.8657	0.34429\\
3.8611	0.34781\\
3.8566	0.35134\\
3.852	0.3549\\
3.8475	0.35848\\
3.8429	0.36208\\
3.8383	0.3657\\
3.8338	0.36934\\
3.8292	0.373\\
3.8247	0.37668\\
3.8201	0.38038\\
3.8155	0.3841\\
3.811	0.38785\\
3.8064	0.39161\\
3.8019	0.39539\\
3.7973	0.39918\\
3.7927	0.403\\
3.7882	0.40684\\
3.7836	0.41069\\
3.7791	0.41456\\
3.7745	0.41845\\
3.7699	0.42235\\
3.7654	0.42627\\
3.7608	0.43021\\
3.7563	0.43416\\
3.7517	0.43813\\
3.7471	0.44212\\
3.7426	0.44611\\
3.738	0.45013\\
3.7335	0.45415\\
3.7289	0.45819\\
3.7243	0.46224\\
3.7198	0.46631\\
3.7152	0.47038\\
3.7107	0.47447\\
3.7061	0.47857\\
3.7016	0.48267\\
3.697	0.48679\\
3.6924	0.49092\\
3.6879	0.49505\\
3.6833	0.49919\\
3.6788	0.50334\\
3.6742	0.50749\\
3.6696	0.51165\\
3.6651	0.51582\\
3.6605	0.51999\\
3.656	0.52416\\
3.6514	0.52834\\
3.6468	0.53251\\
3.6423	0.53669\\
3.6377	0.54087\\
3.6332	0.54505\\
3.6286	0.54922\\
3.624	0.5534\\
3.6195	0.55757\\
3.6149	0.56174\\
3.6104	0.5659\\
3.6058	0.57006\\
3.6012	0.57421\\
3.5967	0.57835\\
3.5921	0.58248\\
3.5876	0.58661\\
3.583	0.59072\\
3.5784	0.59482\\
3.5739	0.59891\\
3.5693	0.60299\\
3.5648	0.60705\\
3.5602	0.6111\\
3.5556	0.61513\\
3.5511	0.61914\\
3.5465	0.62313\\
3.542	0.62711\\
3.5374	0.63106\\
3.5328	0.63499\\
3.5283	0.6389\\
3.5237	0.64279\\
3.5192	0.64665\\
3.5146	0.65048\\
3.51	0.65429\\
3.5055	0.65806\\
3.5009	0.66181\\
3.4964	0.66553\\
3.4918	0.66922\\
3.4872	0.67287\\
3.4827	0.67649\\
3.4781	0.68008\\
3.4736	0.68363\\
3.469	0.68715\\
3.4644	0.69062\\
3.4599	0.69406\\
3.4553	0.69746\\
3.4508	0.70082\\
3.4462	0.70413\\
3.4416	0.7074\\
3.4371	0.71063\\
3.4325	0.71382\\
3.428	0.71696\\
3.4234	0.72005\\
3.4188	0.72309\\
3.4143	0.72609\\
3.4097	0.72903\\
3.4052	0.73193\\
3.4006	0.73478\\
3.396	0.73757\\
3.3915	0.74031\\
3.3869	0.74299\\
3.3824	0.74563\\
3.3778	0.7482\\
3.3732	0.75072\\
3.3687	0.75319\\
3.3641	0.75559\\
3.3596	0.75794\\
3.355	0.76023\\
3.3504	0.76245\\
3.3459	0.76462\\
3.3413	0.76672\\
3.3368	0.76876\\
3.3322	0.77074\\
3.3276	0.77265\\
3.3231	0.77449\\
3.3185	0.77627\\
3.314	0.77798\\
3.3094	0.77962\\
3.3048	0.78119\\
3.3003	0.78269\\
3.2957	0.78411\\
3.2912	0.78546\\
3.2866	0.78673\\
3.282	0.78793\\
3.2775	0.78905\\
3.2729	0.79008\\
3.2684	0.79104\\
3.2638	0.79191\\
3.2592	0.7927\\
3.2547	0.7934\\
3.2501	0.79402\\
3.2456	0.79455\\
3.241	0.79499\\
3.2364	0.79534\\
3.2319	0.7956\\
3.2273	0.79577\\
3.2228	0.79584\\
3.2182	0.79583\\
3.2136	0.79572\\
3.2091	0.79552\\
3.2045	0.79523\\
3.2	0.79484\\
3.1954	0.79436\\
3.1908	0.79379\\
3.1863	0.79312\\
3.1817	0.79236\\
3.1772	0.79151\\
3.1726	0.79057\\
3.168	0.78954\\
3.1635	0.78841\\
3.1589	0.7872\\
3.1544	0.78589\\
3.1498	0.7845\\
3.1453	0.78302\\
3.1407	0.78144\\
3.1361	0.77979\\
3.1316	0.77804\\
3.127	0.77621\\
3.1225	0.77429\\
3.1179	0.77229\\
3.1133	0.7702\\
3.1088	0.76803\\
3.1042	0.76578\\
3.0997	0.76345\\
3.0951	0.76103\\
3.0905	0.75854\\
3.086	0.75596\\
3.0814	0.75331\\
3.0769	0.75058\\
3.0723	0.74778\\
3.0677	0.7449\\
3.0632	0.74194\\
3.0586	0.73892\\
3.0541	0.73582\\
3.0495	0.73265\\
3.0449	0.72941\\
3.0404	0.72611\\
3.0358	0.72274\\
3.0313	0.7193\\
3.0267	0.71581\\
3.0221	0.71225\\
3.0176	0.70863\\
3.013	0.70496\\
3.0085	0.70123\\
3.0039	0.69744\\
2.9993	0.69361\\
2.9948	0.68972\\
2.9902	0.68579\\
2.9857	0.68181\\
2.9811	0.67779\\
2.9765	0.67372\\
2.972	0.66962\\
2.9674	0.66548\\
2.9629	0.6613\\
2.9583	0.6571\\
2.9537	0.65286\\
2.9492	0.64859\\
2.9446	0.64429\\
2.9401	0.63997\\
2.9355	0.63563\\
2.9309	0.63126\\
2.9264	0.62687\\
2.9218	0.62247\\
2.9173	0.61804\\
2.9127	0.6136\\
2.9081	0.60915\\
2.9036	0.60468\\
2.899	0.6002\\
2.8945	0.59571\\
2.8899	0.59121\\
2.8853	0.5867\\
2.8808	0.58218\\
2.8762	0.57765\\
2.8717	0.57312\\
2.8671	0.56859\\
2.8625	0.56405\\
2.858	0.5595\\
2.8534	0.55496\\
2.8489	0.55041\\
2.8443	0.54586\\
2.8397	0.54131\\
2.8352	0.53676\\
2.8306	0.53222\\
2.8261	0.52767\\
2.8215	0.52313\\
2.8169	0.51859\\
2.8124	0.51406\\
2.8078	0.50953\\
2.8033	0.50501\\
2.7987	0.50049\\
2.7941	0.49598\\
2.7896	0.49148\\
2.785	0.48698\\
2.7805	0.4825\\
2.7759	0.47803\\
2.7713	0.47356\\
2.7668	0.46911\\
2.7622	0.46467\\
2.7577	0.46024\\
2.7531	0.45582\\
2.7485	0.45142\\
2.744	0.44703\\
2.7394	0.44265\\
2.7349	0.43829\\
2.7303	0.43395\\
2.7257	0.42962\\
2.7212	0.4253\\
2.7166	0.42101\\
2.7121	0.41673\\
2.7075	0.41247\\
2.7029	0.40823\\
2.6984	0.404\\
2.6938	0.3998\\
2.6893	0.39561\\
2.6847	0.39145\\
2.6801	0.38731\\
2.6756	0.38318\\
2.671	0.37908\\
2.6665	0.375\\
2.6619	0.37094\\
2.6573	0.3669\\
2.6528	0.36289\\
2.6482	0.3589\\
2.6437	0.35493\\
2.6391	0.35098\\
2.6345	0.34706\\
2.63	0.34316\\
2.6254	0.33929\\
2.6209	0.33544\\
2.6163	0.33161\\
2.6117	0.32781\\
2.6072	0.32403\\
2.6026	0.32028\\
2.5981	0.31656\\
2.5935	0.31286\\
2.5889	0.30918\\
2.5844	0.30554\\
2.5798	0.30191\\
2.5753	0.29832\\
2.5707	0.29475\\
2.5662	0.2912\\
2.5616	0.28769\\
2.557	0.2842\\
2.5525	0.28073\\
2.5479	0.27729\\
2.5434	0.27388\\
2.5388	0.2705\\
2.5342	0.26714\\
2.5297	0.26381\\
2.5251	0.26051\\
2.5206	0.25724\\
2.516	0.25399\\
2.5114	0.25077\\
2.5069	0.24758\\
2.5023	0.24441\\
2.4978	0.24127\\
2.4932	0.23816\\
2.4886	0.23507\\
2.4841	0.23202\\
2.4795	0.22899\\
2.475	0.22599\\
2.4704	0.22301\\
2.4658	0.22006\\
2.4613	0.21714\\
2.4567	0.21425\\
2.4522	0.21138\\
2.4476	0.20854\\
2.443	0.20573\\
2.4385	0.20295\\
2.4339	0.20019\\
2.4294	0.19746\\
2.4248	0.19476\\
2.4202	0.19208\\
2.4157	0.18943\\
2.4111	0.1868\\
2.4066	0.18421\\
2.402	0.18164\\
2.3974	0.17909\\
2.3929	0.17658\\
2.3883	0.17408\\
2.3838	0.17162\\
2.3792	0.16918\\
2.3746	0.16677\\
2.3701	0.16438\\
2.3655	0.16202\\
2.361	0.15968\\
2.3564	0.15737\\
2.3518	0.15508\\
2.3473	0.15282\\
2.3427	0.15059\\
2.3382	0.14838\\
2.3336	0.14619\\
2.329	0.14403\\
2.3245	0.1419\\
2.3199	0.13979\\
2.3154	0.1377\\
2.3108	0.13564\\
2.3062	0.1336\\
2.3017	0.13159\\
2.2971	0.1296\\
2.2926	0.12763\\
2.288	0.12569\\
2.2834	0.12377\\
2.2789	0.12187\\
2.2743	0.12\\
2.2698	0.11815\\
2.2652	0.11632\\
2.2606	0.11451\\
2.2561	0.11273\\
2.2515	0.11097\\
2.247	0.10923\\
2.2424	0.10752\\
2.2378	0.10582\\
2.2333	0.10415\\
2.2287	0.1025\\
2.2242	0.10087\\
2.2196	0.09926\\
2.215	0.097673\\
2.2105	0.096106\\
2.2059	0.09456\\
2.2014	0.093034\\
2.1968	0.091529\\
2.1922	0.090044\\
2.1877	0.088579\\
2.1831	0.087133\\
2.1786	0.085708\\
2.174	0.084301\\
2.1694	0.082914\\
2.1649	0.081546\\
2.1603	0.080198\\
2.1558	0.078867\\
2.1512	0.077556\\
2.1466	0.076262\\
2.1421	0.074987\\
2.1375	0.07373\\
2.133	0.072491\\
2.1284	0.07127\\
2.1238	0.070066\\
2.1193	0.068879\\
2.1147	0.06771\\
2.1102	0.066557\\
2.1056	0.065421\\
2.101	0.064302\\
2.0965	0.0632\\
2.0919	0.062113\\
2.0874	0.061043\\
2.0828	0.059989\\
2.0782	0.05895\\
2.0737	0.057927\\
2.0691	0.056919\\
2.0646	0.055927\\
2.06	0.054949\\
2.0554	0.053987\\
2.0509	0.053039\\
2.0463	0.052105\\
2.0418	0.051186\\
2.0372	0.050281\\
2.0326	0.049391\\
2.0281	0.048514\\
2.0235	0.04765\\
2.019	0.0468\\
2.0144	0.045964\\
2.0098	0.045141\\
2.0053	0.04433\\
2.0007	0.043533\\
1.9962	0.042748\\
1.9916	0.041976\\
1.9871	0.041216\\
1.9825	0.040468\\
1.9779	0.039733\\
1.9734	0.039009\\
1.9688	0.038297\\
1.9643	0.037597\\
1.9597	0.036908\\
1.9551	0.03623\\
1.9506	0.035564\\
1.946	0.034908\\
1.9415	0.034263\\
1.9369	0.033629\\
1.9323	0.033006\\
1.9278	0.032392\\
1.9232	0.03179\\
1.9187	0.031197\\
1.9141	0.030614\\
1.9095	0.030041\\
1.905	0.029477\\
1.9004	0.028924\\
1.8959	0.028379\\
1.8913	0.027844\\
1.8867	0.027318\\
1.8822	0.026801\\
1.8776	0.026293\\
1.8731	0.025793\\
1.8685	0.025302\\
1.8639	0.02482\\
1.8594	0.024346\\
1.8548	0.023881\\
1.8503	0.023423\\
1.8457	0.022973\\
1.8411	0.022532\\
1.8366	0.022098\\
1.832	0.021671\\
1.8275	0.021252\\
1.8229	0.020841\\
1.8183	0.020437\\
1.8138	0.02004\\
1.8092	0.01965\\
1.8047	0.019267\\
1.8001	0.018891\\
1.7955	0.018521\\
1.791	0.018159\\
1.7864	0.017803\\
1.7819	0.017453\\
1.7773	0.017109\\
1.7727	0.016772\\
1.7682	0.016441\\
1.7636	0.016116\\
1.7591	0.015797\\
1.7545	0.015483\\
1.7499	0.015176\\
1.7454	0.014874\\
1.7408	0.014577\\
1.7363	0.014286\\
1.7317	0.014001\\
1.7271	0.013721\\
1.7226	0.013446\\
1.718	0.013176\\
1.7135	0.012911\\
1.7089	0.012651\\
1.7043	0.012395\\
1.6998	0.012145\\
1.6952	0.011899\\
1.6907	0.011658\\
1.6861	0.011422\\
1.6815	0.01119\\
1.677	0.010962\\
1.6724	0.010738\\
1.6679	0.010519\\
1.6633	0.010304\\
1.6587	0.010093\\
1.6542	0.0098865\\
1.6496	0.0096836\\
1.6451	0.0094845\\
1.6405	0.0092893\\
1.6359	0.0090978\\
1.6314	0.00891\\
1.6268	0.0087258\\
1.6223	0.0085451\\
1.6177	0.008368\\
1.6131	0.0081943\\
1.6086	0.0080239\\
1.604	0.0078569\\
1.5995	0.0076931\\
1.5949	0.0075325\\
1.5903	0.0073751\\
1.5858	0.0072207\\
1.5812	0.0070693\\
1.5767	0.006921\\
1.5721	0.0067755\\
1.5675	0.0066329\\
1.563	0.0064931\\
1.5584	0.0063561\\
1.5539	0.0062218\\
1.5493	0.0060901\\
1.5447	0.0059611\\
1.5402	0.0058346\\
1.5356	0.0057106\\
1.5311	0.0055892\\
1.5265	0.0054701\\
1.5219	0.0053534\\
1.5174	0.0052391\\
1.5128	0.005127\\
1.5083	0.0050172\\
1.5037	0.0049097\\
1.4991	0.0048042\\
1.4946	0.0047009\\
1.49	0.0045997\\
1.4855	0.0045006\\
1.4809	0.0044035\\
1.4763	0.0043083\\
1.4718	0.0042151\\
1.4672	0.0041237\\
1.4627	0.0040342\\
1.4581	0.0039466\\
1.4535	0.0038608\\
1.449	0.0037767\\
1.4444	0.0036943\\
1.4399	0.0036136\\
1.4353	0.0035346\\
1.4308	0.0034572\\
1.4262	0.0033814\\
1.4216	0.0033072\\
1.4171	0.0032345\\
1.4125	0.0031634\\
1.408	0.0030937\\
1.4034	0.0030254\\
1.3988	0.0029586\\
1.3943	0.0028931\\
1.3897	0.0028291\\
1.3852	0.0027663\\
1.3806	0.0027049\\
1.376	0.0026448\\
1.3715	0.0025859\\
1.3669	0.0025283\\
1.3624	0.0024719\\
1.3578	0.0024166\\
1.3532	0.0023626\\
1.3487	0.0023096\\
1.3441	0.0022578\\
1.3396	0.0022071\\
1.335	0.0021575\\
1.3304	0.0021089\\
1.3259	0.0020614\\
1.3213	0.0020149\\
1.3168	0.0019693\\
1.3122	0.0019248\\
1.3076	0.0018812\\
1.3031	0.0018385\\
1.2985	0.0017967\\
1.294	0.0017559\\
1.2894	0.0017159\\
1.2848	0.0016768\\
1.2803	0.0016385\\
1.2757	0.0016011\\
1.2712	0.0015644\\
1.2666	0.0015286\\
1.262	0.0014935\\
1.2575	0.0014592\\
1.2529	0.0014257\\
1.2484	0.0013928\\
1.2438	0.0013607\\
1.2392	0.0013293\\
1.2347	0.0012986\\
1.2301	0.0012686\\
1.2256	0.0012392\\
1.221	0.0012104\\
1.2164	0.0011823\\
1.2119	0.0011548\\
1.2073	0.0011279\\
1.2028	0.0011016\\
1.1982	0.0010759\\
1.1936	0.0010508\\
1.1891	0.0010262\\
1.1845	0.0010022\\
1.18	0.00097865\\
1.1754	0.00095567\\
1.1708	0.00093319\\
1.1663	0.00091123\\
1.1617	0.00088975\\
1.1572	0.00086875\\
1.1526	0.00084822\\
1.148	0.00082815\\
1.1435	0.00080854\\
1.1389	0.00078936\\
1.1344	0.00077062\\
1.1298	0.0007523\\
1.1252	0.00073439\\
1.1207	0.00071689\\
1.1161	0.00069978\\
1.1116	0.00068306\\
1.107	0.00066672\\
1.1024	0.00065075\\
1.0979	0.00063515\\
1.0933	0.0006199\\
1.0888	0.000605\\
1.0842	0.00059044\\
1.0796	0.00057621\\
1.0751	0.00056231\\
1.0705	0.00054873\\
1.066	0.00053546\\
1.0614	0.0005225\\
1.0568	0.00050983\\
1.0523	0.00049745\\
1.0477	0.00048537\\
1.0432	0.00047356\\
1.0386	0.00046202\\
1.034	0.00045075\\
1.0295	0.00043975\\
1.0249	0.00042899\\
1.0204	0.00041849\\
1.0158	0.00040824\\
1.0112	0.00039822\\
1.0067	0.00038844\\
1.0021	0.00037888\\
0.99757	0.00036955\\
0.99301	0.00036044\\
0.98845	0.00035155\\
0.98389	0.00034286\\
0.97933	0.00033437\\
0.97477	0.00032609\\
0.97021	0.000318\\
0.96565	0.0003101\\
0.96109	0.00030239\\
0.95653	0.00029487\\
0.95197	0.00028752\\
0.94741	0.00028034\\
0.94285	0.00027334\\
0.93829	0.0002665\\
0.93373	0.00025983\\
0.92917	0.00025332\\
0.92461	0.00024696\\
0.92005	0.00024075\\
0.91549	0.0002347\\
0.91093	0.00022878\\
0.90637	0.00022301\\
0.90181	0.00021738\\
0.89725	0.00021189\\
0.89269	0.00020652\\
0.88813	0.00020129\\
0.88357	0.00019618\\
0.87901	0.0001912\\
0.87445	0.00018634\\
0.86989	0.0001816\\
0.86533	0.00017697\\
0.86077	0.00017245\\
0.85621	0.00016804\\
}--cycle;
\addplot[ybar interval, fill=black, fill opacity=0.35, area legend, draw=none] table[row sep=crcr, x=Lower, y=Count] {%
Lower	Upper	Count\\
0.85621	1.2066	0\\
1.2066	1.557	0\\
1.557	1.9074	0\\
1.9074	2.2578	0.01522\\
2.2578	2.6082	0.14459\\
2.6082	2.9586	0.47183\\
2.9586	3.3091	0.93606\\
3.3091	3.6595	0.7458\\
3.6595	4.0099	0.35007\\
4.0099	4.3603	0.11415\\
4.3603	4.7107	0.053272\\
4.7107	5.0611	0.0076102\\
5.0611	5.4115	0.01522\\
5.4115	5.4115	0.01522\\
};
\addlegendentry{\RealHistogramLegend}

\addplot [color=black]
  table[row sep=crcr]{%
0.85621	4.2825e-08\\
1.6177	0.00037624\\
1.7408	0.0011837\\
1.8229	0.0024165\\
1.8867	0.0040935\\
1.9415	0.0063075\\
1.9871	0.0089199\\
2.0281	0.012055\\
2.0646	0.015623\\
2.0965	0.019473\\
2.1284	0.024123\\
2.1558	0.028843\\
2.1831	0.034332\\
2.2105	0.040681\\
2.2333	0.0467\\
2.2561	0.053442\\
2.2789	0.060967\\
2.3017	0.069334\\
2.3245	0.078604\\
2.3473	0.088835\\
2.3701	0.10008\\
2.3883	0.10985\\
2.4066	0.12034\\
2.4248	0.13156\\
2.443	0.14354\\
2.4613	0.15629\\
2.4795	0.16985\\
2.4978	0.1842\\
2.516	0.19938\\
2.5342	0.21537\\
2.557	0.23651\\
2.5798	0.25892\\
2.6026	0.28256\\
2.6254	0.30741\\
2.6482	0.3334\\
2.671	0.36046\\
2.6984	0.39422\\
2.7257	0.42921\\
2.7577	0.4713\\
2.7987	0.52679\\
2.8899	0.65072\\
2.9218	0.69249\\
2.9492	0.72688\\
2.972	0.75427\\
2.9902	0.77519\\
3.0085	0.79512\\
3.0267	0.81394\\
3.0449	0.83157\\
3.0632	0.8479\\
3.0769	0.85924\\
3.0905	0.86977\\
3.1042	0.87946\\
3.1179	0.88827\\
3.1316	0.89617\\
3.1453	0.90315\\
3.1589	0.90917\\
3.1726	0.91423\\
3.1863	0.91831\\
3.2	0.92139\\
3.2136	0.92346\\
3.2273	0.92453\\
3.241	0.92458\\
3.2547	0.92362\\
3.2684	0.92166\\
3.282	0.91871\\
3.2957	0.91476\\
3.3094	0.90985\\
3.3231	0.90399\\
3.3368	0.89719\\
3.3504	0.88949\\
3.3641	0.88091\\
3.3778	0.87149\\
3.3915	0.86124\\
3.4052	0.85021\\
3.4234	0.83435\\
3.4416	0.81726\\
3.4599	0.79904\\
3.4781	0.77977\\
3.5009	0.75441\\
3.5237	0.7278\\
3.5511	0.69453\\
3.583	0.65436\\
3.6286	0.5955\\
3.697	0.50713\\
3.7335	0.46136\\
3.7608	0.42814\\
3.7882	0.39609\\
3.811	0.3704\\
3.8338	0.34574\\
3.8566	0.32215\\
3.8794	0.2997\\
3.9022	0.27842\\
3.925	0.25832\\
3.9478	0.23942\\
3.9706	0.22171\\
3.9934	0.20518\\
4.0162	0.1898\\
4.039	0.17553\\
4.0618	0.16234\\
4.0846	0.15019\\
4.1074	0.13901\\
4.1302	0.12876\\
4.153	0.11939\\
4.1758	0.11082\\
4.1986	0.10302\\
4.2214	0.095905\\
4.2487	0.088211\\
4.2761	0.081344\\
4.3034	0.075213\\
4.3354	0.068876\\
4.3673	0.063301\\
4.4038	0.057716\\
4.4448	0.052261\\
4.4904	0.047022\\
4.5406	0.042031\\
4.5998	0.036908\\
4.6682	0.031758\\
4.7503	0.026355\\
4.8415	0.021124\\
4.9373	0.016391\\
5.033	0.012399\\
5.1333	0.0089868\\
5.2382	0.0062071\\
5.3522	0.0039897\\
5.4115	0.0031186\\
};
\addlegendentry{\RealBLNFitLegend}

      % \input{../../../../.tex/20/KS/SizeDistributionFittingsfig39.tex}
      \legend{}
    \end{groupplot}
  %\endgroup

  %\begingroup Parametri
    
%\begingroup Parametro
\providecommand\varyingParameterSymbol\lambda
\providecommand\varyingParameterList[1]{
  \ifcase#1\relax%
  \or\num{1e-2}% Caso 1
  \or\num{8.2e-2}% Caso 2
  \or\num{1.55e-1}% Caso 3
  \or\num{2.27e-1}% Caso 4
  \else\num{3e-1}% Caso 5
  \fi
}
%\endgroup

%\begingroup beta
\providecommand\betaExactList[1]{
  \ifcase#1\relax%
  \or\num{1.267+-0.007}% Caso 1
  \or\num{1.285+-0.010}% Caso 2
  \or\num{1.264+-0.015}% Caso 3
  \or\num{1.246+-0.015}% Caso 4
  \else\num{1.210+-0.015}% Caso 5
  \fi
}
\providecommand\betaAprxList[1]{
  \ifcase#1\relax%
  \or\num{1.264+-0.006}% Caso 1
  \or\num{1.270+-0.011}% Caso 2
  \or\num{1.280+-0.013}% Caso 3
  \or\num{1.250+-0.015}% Caso 4
  \else\num{1.218+-0.014}% Caso 5
  \fi
}
%\endgroup

%\begingroup xi
\providecommand\zetaExactList[1]{
  \ifcase#1\relax%
  \or\num{0.757+-0.004}% Caso 1
  \or\num{0.710+-0.026}% Caso 2
  \or\num{0.670+-0.035}% Caso 3
  \or\num{0.625+-0.052}% Caso 4
  \else\num{0.539+-0.059}% Caso 5
  \fi
}
\providecommand\zetaAprxList[1]{
  \ifcase#1\relax%
  \or\num{0.749+-0.004}% Caso 1
  \or\num{0.721+-0.022}% Caso 2
  \or\num{0.645+-0.040}% Caso 3
  \or\num{0.622+-0.051}% Caso 4
  \else\num{0.567+-0.053}% Caso 5
  \fi
}
%\endgroup

%\begingroup mI
\providecommand\mIExactList[1]{
  \ifcase#1\relax%
  \or\num{3.265+-0.001}% Caso 1
  \or\num{3.238+-0.005}% Caso 2
  \or\num{3.220+-0.008}% Caso 3
  \or\num{3.213+-0.013}% Caso 4
  \else\num{3.174+-0.023}% Caso 5
  \fi
}
\providecommand\mIAprxList[1]{
  \ifcase#1\relax%
  \or\num{3.262+-0.001}% Caso 1
  \or\num{3.240+-0.005}% Caso 2
  \or\num{3.217+-0.010}% Caso 3
  \or\num{3.217+-0.012}% Caso 4
  \else\num{3.208+-0.009}% Caso 5
  \fi
}
%\endgroup

%\begingroup sI
\providecommand\sIExactList[1]{
  \ifcase#1\relax%
  \or\num{0.284+-0.001}% Caso 1
  \or\num{0.311+-0.008}% Caso 2
  \or\num{0.343+-0.010}% Caso 3
  \or\num{0.361+-0.020}% Caso 4
  \else\num{0.398+-0.029}% Caso 5
  \fi
}
\providecommand\sIAprxList[1]{
  \ifcase#1\relax%
  \or\num{0.281+-0.001}% Caso 1
  \or\num{0.315+-0.006}% Caso 2
  \or\num{0.334+-0.013}% Caso 3
  \or\num{0.366+-0.018}% Caso 4
  \else\num{0.411+-0.025}% Caso 5
  \fi
}
%\endgroup

%\begingroup mII
\providecommand\mIIExactList[1]{
  \ifcase#1\relax%
  \or\num{3.788+-0.005}% Caso 1
  \or\num{3.730+-0.029}% Caso 2
  \or\num{3.681+-0.052}% Caso 3
  \or\num{3.650+-0.072}% Caso 4
  \else\num{3.572+-0.095}% Caso 5
  \fi
}
\providecommand\mIIAprxList[1]{
  \ifcase#1\relax%
  \or\num{3.778+-0.006}% Caso 1
  \or\num{3.735+-0.029}% Caso 2
  \or\num{3.656+-0.048}% Caso 3
  \or\num{3.658+-0.084}% Caso 4
  \else\num{3.502+-0.078}% Caso 5
  \fi
}
%\endgroup

%\begingroup sII
\providecommand\sIIExactList[1]{
  \ifcase#1\relax%
  \or\num{0.455+-0.002}% Caso 1
  \or\num{0.486+-0.007}% Caso 2
  \or\num{0.515+-0.014}% Caso 3
  \or\num{0.533+-0.021}% Caso 4
  \else\num{0.570+-0.028}% Caso 5
  \fi
}
\providecommand\sIIAprxList[1]{
  \ifcase#1\relax%
  \or\num{0.457+-0.002}% Caso 1
  \or\num{0.490+-0.007}% Caso 2
  \or\num{0.521+-0.012}% Caso 3
  \or\num{0.545+-0.021}% Caso 4
  \else\num{0.586+-0.029}% Caso 5
  \fi
}
%\endgroup

    \ifpageisodd\writeVaryingParameter{1}{west}{west}{0em}{0em}
    \else\writeVaryingParameter{3}{east}{east}{0em}{0em}\fi
  %\endgroup

  %\begingroup Didascalie
    \writeSubCaptions{north east}{north east}{0em}{0em}
    % https://tex.stackexchange.com/questions/632045/get-subcaption-for-every-plot-using-groupplot-pgf-tikz-after-2020-update
  %\endgroup

  % \node[anchor=south] at (Row1 c2r1.north) {\ifpageisodd odd\else even\fi};

  \redefineTikZbounds{Row1 c1r1}{Row5 c3r1}
\end{tikzpicture}%

% \input{Parametri.txt}
% 
%\begingroup Parametro
\providecommand\varyingParameterSymbol\lambda
\providecommand\varyingParameterList[1]{
  \ifcase#1\relax%
  \or\num{1e-2}% Caso 1
  \or\num{8.2e-2}% Caso 2
  \or\num{1.55e-1}% Caso 3
  \or\num{2.27e-1}% Caso 4
  \else\num{3e-1}% Caso 5
  \fi
}
%\endgroup

%\begingroup beta
\providecommand\betaExactList[1]{
  \ifcase#1\relax%
  \or\num{1.267+-0.007}% Caso 1
  \or\num{1.285+-0.010}% Caso 2
  \or\num{1.264+-0.015}% Caso 3
  \or\num{1.246+-0.015}% Caso 4
  \else\num{1.210+-0.015}% Caso 5
  \fi
}
\providecommand\betaAprxList[1]{
  \ifcase#1\relax%
  \or\num{1.264+-0.006}% Caso 1
  \or\num{1.270+-0.011}% Caso 2
  \or\num{1.280+-0.013}% Caso 3
  \or\num{1.250+-0.015}% Caso 4
  \else\num{1.218+-0.014}% Caso 5
  \fi
}
%\endgroup

%\begingroup xi
\providecommand\zetaExactList[1]{
  \ifcase#1\relax%
  \or\num{0.757+-0.004}% Caso 1
  \or\num{0.710+-0.026}% Caso 2
  \or\num{0.670+-0.035}% Caso 3
  \or\num{0.625+-0.052}% Caso 4
  \else\num{0.539+-0.059}% Caso 5
  \fi
}
\providecommand\zetaAprxList[1]{
  \ifcase#1\relax%
  \or\num{0.749+-0.004}% Caso 1
  \or\num{0.721+-0.022}% Caso 2
  \or\num{0.645+-0.040}% Caso 3
  \or\num{0.622+-0.051}% Caso 4
  \else\num{0.567+-0.053}% Caso 5
  \fi
}
%\endgroup

%\begingroup mI
\providecommand\mIExactList[1]{
  \ifcase#1\relax%
  \or\num{3.265+-0.001}% Caso 1
  \or\num{3.238+-0.005}% Caso 2
  \or\num{3.220+-0.008}% Caso 3
  \or\num{3.213+-0.013}% Caso 4
  \else\num{3.174+-0.023}% Caso 5
  \fi
}
\providecommand\mIAprxList[1]{
  \ifcase#1\relax%
  \or\num{3.262+-0.001}% Caso 1
  \or\num{3.240+-0.005}% Caso 2
  \or\num{3.217+-0.010}% Caso 3
  \or\num{3.217+-0.012}% Caso 4
  \else\num{3.208+-0.009}% Caso 5
  \fi
}
%\endgroup

%\begingroup sI
\providecommand\sIExactList[1]{
  \ifcase#1\relax%
  \or\num{0.284+-0.001}% Caso 1
  \or\num{0.311+-0.008}% Caso 2
  \or\num{0.343+-0.010}% Caso 3
  \or\num{0.361+-0.020}% Caso 4
  \else\num{0.398+-0.029}% Caso 5
  \fi
}
\providecommand\sIAprxList[1]{
  \ifcase#1\relax%
  \or\num{0.281+-0.001}% Caso 1
  \or\num{0.315+-0.006}% Caso 2
  \or\num{0.334+-0.013}% Caso 3
  \or\num{0.366+-0.018}% Caso 4
  \else\num{0.411+-0.025}% Caso 5
  \fi
}
%\endgroup

%\begingroup mII
\providecommand\mIIExactList[1]{
  \ifcase#1\relax%
  \or\num{3.788+-0.005}% Caso 1
  \or\num{3.730+-0.029}% Caso 2
  \or\num{3.681+-0.052}% Caso 3
  \or\num{3.650+-0.072}% Caso 4
  \else\num{3.572+-0.095}% Caso 5
  \fi
}
\providecommand\mIIAprxList[1]{
  \ifcase#1\relax%
  \or\num{3.778+-0.006}% Caso 1
  \or\num{3.735+-0.029}% Caso 2
  \or\num{3.656+-0.048}% Caso 3
  \or\num{3.658+-0.084}% Caso 4
  \else\num{3.502+-0.078}% Caso 5
  \fi
}
%\endgroup

%\begingroup sII
\providecommand\sIIExactList[1]{
  \ifcase#1\relax%
  \or\num{0.455+-0.002}% Caso 1
  \or\num{0.486+-0.007}% Caso 2
  \or\num{0.515+-0.014}% Caso 3
  \or\num{0.533+-0.021}% Caso 4
  \else\num{0.570+-0.028}% Caso 5
  \fi
}
\providecommand\sIIAprxList[1]{
  \ifcase#1\relax%
  \or\num{0.457+-0.002}% Caso 1
  \or\num{0.490+-0.007}% Caso 2
  \or\num{0.521+-0.012}% Caso 3
  \or\num{0.545+-0.021}% Caso 4
  \else\num{0.586+-0.029}% Caso 5
  \fi
}
%\endgroup


  \endgroup

  \begingroup
    \setcounter{subfigure}{0}
    
% !TEX root = ../../../../../Esperimenti/.tex/MEF.tex
% LTeX: language=it

\begin{tikzpicture}%
  %\begingroup Impostazioni
    \infopageKeysI
    \pgfplotsset{
      row1Keys/.append style={enlarge x limits=.05}
    }
  %\endgroup

  %\begingroup Righe
    \begin{groupplot}[
      group style={
        group name=Row1,
        plotGroup,
      },
      plotRow,
      row2KeysWithLabels,
      plotLegendNE,
      %
      % xmin=3637.5,xmax=2.5003e+05,
      ymin=0.0053191,ymax=1,
    ]
      \nextgroupplot[%
        xmin=3.598,xmax=5.4115,
        % xmin=3.4,xmax=5.4,
        % ymin=0.0053191,ymax=1,
      ] % This file was created by matlab2tikz.
%
\definecolor{mycolor1}{rgb}{0.00000,0.44706,0.69804}%
\definecolor{mycolor2}{rgb}{0.00000,0.44700,0.74100}%
\definecolor{mycolor3}{rgb}{0.83529,0.36863,0.00000}%
\definecolor{mycolor4}{rgb}{0.49400,0.18400,0.55600}%
%
\addplot[only marks, mark=*, mark size=1.1180pt, color=mycolor1, fill=mycolor1, opacity=0.60, draw=none, mark options={draw=none,line width=0pt}] table[row sep=crcr]{%
x	y\\
3.5893	0.99468\\
3.5929	0.98404\\
3.5962	0.9734\\
3.6001	0.96277\\
3.6042	0.95213\\
3.6081	0.94149\\
3.6124	0.93085\\
3.6162	0.92021\\
3.6204	0.90957\\
3.625	0.89894\\
3.6291	0.8883\\
3.6335	0.87766\\
3.6385	0.86702\\
3.6436	0.85638\\
3.6481	0.84574\\
3.6533	0.83511\\
3.6583	0.82447\\
3.6634	0.81383\\
3.669	0.80319\\
3.6741	0.79255\\
3.6786	0.78191\\
3.683	0.77128\\
3.6891	0.76064\\
3.6957	0.75\\
3.7019	0.73936\\
3.7088	0.72872\\
3.7138	0.71809\\
3.7192	0.70745\\
3.7256	0.69681\\
3.7321	0.68617\\
3.7405	0.67553\\
3.7483	0.66489\\
3.7548	0.65426\\
3.763	0.64362\\
3.7715	0.63298\\
3.7797	0.62234\\
3.786	0.6117\\
3.7935	0.60106\\
3.8025	0.59043\\
3.8109	0.57979\\
3.819	0.56915\\
3.8255	0.55851\\
3.8327	0.54787\\
3.8385	0.53723\\
3.8461	0.5266\\
3.852	0.51596\\
3.86	0.50532\\
3.866	0.49468\\
3.8727	0.48404\\
3.8788	0.4734\\
3.8839	0.46277\\
3.8896	0.45213\\
3.8958	0.44149\\
3.9012	0.43085\\
3.9081	0.42021\\
3.9155	0.40957\\
3.9216	0.39894\\
3.928	0.3883\\
3.9348	0.37766\\
3.9415	0.36702\\
3.9492	0.35638\\
3.9579	0.34574\\
3.9692	0.33511\\
3.9828	0.32447\\
3.9977	0.31383\\
4.0161	0.30319\\
4.0335	0.29255\\
4.0503	0.28191\\
4.0662	0.27128\\
4.0836	0.26064\\
4.0988	0.25\\
4.1112	0.23936\\
4.1243	0.22872\\
4.1342	0.21809\\
4.1454	0.20745\\
4.1576	0.19681\\
4.1705	0.18617\\
4.1844	0.17553\\
4.2007	0.16489\\
4.2194	0.15426\\
4.24	0.14362\\
4.2662	0.13298\\
4.2921	0.12234\\
4.3201	0.1117\\
4.3525	0.10106\\
4.3736	0.090426\\
4.3895	0.079787\\
4.4094	0.069149\\
4.4358	0.058511\\
4.5276	0.047872\\
4.5794	0.037234\\
4.7208	0.026596\\
4.7905	0.015957\\
5.1574	0.0053191\\
};
\addlegendentry{\ExactEmpiricCCDFLegend}

\addplot [color=mycolor1, only marks, every error bar/.append style={opacity=0.30}, mark=*, mark size=0pt, draw=none, forget plot]
 plot [error bars/.cd, x dir=both, x explicit, error bar style={line width=1pt, color=mycolor1}, error mark options={mark=none,mark size=0pt}]
 table[row sep=crcr, x error plus index=2, x error minus index=3]{%
3.5893	0.99468	0.0020617	0.0020617\\
3.5929	0.98404	0.0020269	0.0020269\\
3.5962	0.9734	0.0020561	0.0020561\\
3.6001	0.96277	0.0019457	0.0019457\\
3.6042	0.95213	0.0019054	0.0019054\\
3.6081	0.94149	0.0018168	0.0018168\\
3.6124	0.93085	0.0018416	0.0018416\\
3.6162	0.92021	0.0018589	0.0018589\\
3.6204	0.90957	0.0018898	0.0018898\\
3.625	0.89894	0.0020758	0.0020758\\
3.6291	0.8883	0.0019633	0.0019633\\
3.6335	0.87766	0.0019186	0.0019186\\
3.6385	0.86702	0.0019754	0.0019754\\
3.6436	0.85638	0.002189	0.002189\\
3.6481	0.84574	0.0020486	0.0020486\\
3.6533	0.83511	0.0020882	0.0020882\\
3.6583	0.82447	0.0020115	0.0020115\\
3.6634	0.81383	0.0020684	0.0020684\\
3.669	0.80319	0.0022894	0.0022894\\
3.6741	0.79255	0.0022671	0.0022671\\
3.6786	0.78191	0.0021921	0.0021921\\
3.683	0.77128	0.0021802	0.0021802\\
3.6891	0.76064	0.0021046	0.0021046\\
3.6957	0.75	0.0020569	0.0020569\\
3.7019	0.73936	0.0021172	0.0021172\\
3.7088	0.72872	0.0020608	0.0020608\\
3.7138	0.71809	0.0020375	0.0020375\\
3.7192	0.70745	0.0020778	0.0020778\\
3.7256	0.69681	0.0020958	0.0020958\\
3.7321	0.68617	0.0022399	0.0022399\\
3.7405	0.67553	0.0023519	0.0023519\\
3.7483	0.66489	0.0024659	0.0024659\\
3.7548	0.65426	0.0025499	0.0025499\\
3.763	0.64362	0.0027554	0.0027554\\
3.7715	0.63298	0.002923	0.002923\\
3.7797	0.62234	0.0028095	0.0028095\\
3.786	0.6117	0.0027316	0.0027316\\
3.7935	0.60106	0.002762	0.002762\\
3.8025	0.59043	0.0028749	0.0028749\\
3.8109	0.57979	0.003014	0.003014\\
3.819	0.56915	0.0027807	0.0027807\\
3.8255	0.55851	0.0025636	0.0025636\\
3.8327	0.54787	0.002599	0.002599\\
3.8385	0.53723	0.0024753	0.0024753\\
3.8461	0.5266	0.0023333	0.0023333\\
3.852	0.51596	0.0022884	0.0022884\\
3.86	0.50532	0.0022617	0.0022617\\
3.866	0.49468	0.0022885	0.0022885\\
3.8727	0.48404	0.0022411	0.0022411\\
3.8788	0.4734	0.0022106	0.0022106\\
3.8839	0.46277	0.0022586	0.0022586\\
3.8896	0.45213	0.0024622	0.0024622\\
3.8958	0.44149	0.0024482	0.0024482\\
3.9012	0.43085	0.0023713	0.0023713\\
3.9081	0.42021	0.0023255	0.0023255\\
3.9155	0.40957	0.0023465	0.0023465\\
3.9216	0.39894	0.0023832	0.0023832\\
3.928	0.3883	0.002503	0.002503\\
3.9348	0.37766	0.0025268	0.0025268\\
3.9415	0.36702	0.0027627	0.0027627\\
3.9492	0.35638	0.0027797	0.0027797\\
3.9579	0.34574	0.0028551	0.0028551\\
3.9692	0.33511	0.0029051	0.0029051\\
3.9828	0.32447	0.0031558	0.0031558\\
3.9977	0.31383	0.003395	0.003395\\
4.0161	0.30319	0.0037587	0.0037587\\
4.0335	0.29255	0.0037678	0.0037678\\
4.0503	0.28191	0.0031906	0.0031906\\
4.0662	0.27128	0.0031187	0.0031187\\
4.0836	0.26064	0.003391	0.003391\\
4.0988	0.25	0.0031882	0.0031882\\
4.1112	0.23936	0.0033285	0.0033285\\
4.1243	0.22872	0.0032658	0.0032658\\
4.1342	0.21809	0.0030368	0.0030368\\
4.1454	0.20745	0.0028444	0.0028444\\
4.1576	0.19681	0.0028028	0.0028028\\
4.1705	0.18617	0.0031669	0.0031669\\
4.1844	0.17553	0.0030161	0.0030161\\
4.2007	0.16489	0.0030855	0.0030855\\
4.2194	0.15426	0.0036315	0.0036315\\
4.24	0.14362	0.0036555	0.0036555\\
4.2662	0.13298	0.0036608	0.0036608\\
4.2921	0.12234	0.0040354	0.0040354\\
4.3201	0.1117	0.004327	0.004327\\
4.3525	0.10106	0.0037223	0.0037223\\
4.3736	0.090426	0.0034217	0.0034217\\
4.3895	0.079787	0.0036394	0.0036394\\
4.4094	0.069149	0.0039009	0.0039009\\
4.4358	0.058511	0.0043355	0.0043355\\
4.5276	0.047872	0.0054577	0.0054577\\
4.5794	0.037234	0.0042494	0.0042494\\
4.7208	0.026596	0.0050438	0.0050438\\
4.7905	0.015957	0.0050741	0.0050741\\
5.1574	0.0053191	0.0034798	0.0034798\\
};
\addplot [color=mycolor1]
  table[row sep=crcr]{%
3.5888	0.98701\\
3.6051	0.95142\\
3.6868	0.75\\
4.4703	0.076296\\
5.2048	0.0089881\\
};
\addlegendentry{\ExactParetoFitLegend}


\addplot[area legend, draw=none, fill=mycolor1, fill opacity=0.15, forget plot]
table[row sep=crcr] {%
x	y\\
3.5888	0.99033\\
3.6051	0.9567\\
3.6215	0.91248\\
3.6378	0.86978\\
3.6541	0.82908\\
3.6704	0.79029\\
3.6868	0.75332\\
3.7031	0.71808\\
3.7194	0.68449\\
3.7357	0.65247\\
3.752	0.62196\\
3.7684	0.59288\\
3.7847	0.56516\\
3.801	0.53875\\
3.8173	0.51357\\
3.8337	0.48958\\
3.85	0.46671\\
3.8663	0.44493\\
3.8826	0.42417\\
3.899	0.4044\\
3.9153	0.38556\\
3.9316	0.36762\\
3.9479	0.35053\\
3.9642	0.33424\\
3.9806	0.31873\\
3.9969	0.30395\\
4.0132	0.28986\\
4.0295	0.27643\\
4.0459	0.26363\\
4.0622	0.25142\\
4.0785	0.23979\\
4.0948	0.22869\\
4.1111	0.21811\\
4.1275	0.20802\\
4.1438	0.19839\\
4.1601	0.18922\\
4.1764	0.18047\\
4.1928	0.17212\\
4.2091	0.16416\\
4.2254	0.15657\\
4.2417	0.14933\\
4.2581	0.14243\\
4.2744	0.13584\\
4.2907	0.12956\\
4.307	0.12357\\
4.3233	0.11786\\
4.3397	0.11241\\
4.356	0.10722\\
4.3723	0.10226\\
4.3886	0.097537\\
4.405	0.09303\\
4.4213	0.088731\\
4.4376	0.084631\\
4.4539	0.08072\\
4.4703	0.07699\\
4.4866	0.073433\\
4.5029	0.07004\\
4.5192	0.066805\\
4.5355	0.063718\\
4.5519	0.060775\\
4.5682	0.057968\\
4.5845	0.05529\\
4.6008	0.052736\\
4.6172	0.0503\\
4.6335	0.047977\\
4.6498	0.045761\\
4.6661	0.043648\\
4.6824	0.041632\\
4.6988	0.039709\\
4.7151	0.037876\\
4.7314	0.036127\\
4.7477	0.034459\\
4.7641	0.032868\\
4.7804	0.03135\\
4.7967	0.029903\\
4.813	0.028522\\
4.8294	0.027205\\
4.8457	0.025949\\
4.862	0.024752\\
4.8783	0.023609\\
4.8946	0.022519\\
4.911	0.02148\\
4.9273	0.020488\\
4.9436	0.019543\\
4.9599	0.018641\\
4.9763	0.01778\\
4.9926	0.01696\\
5.0089	0.016177\\
5.0252	0.015431\\
5.0415	0.014719\\
5.0579	0.01404\\
5.0742	0.013392\\
5.0905	0.012774\\
5.1068	0.012184\\
5.1232	0.011622\\
5.1395	0.011086\\
5.1558	0.010575\\
5.1721	0.010087\\
5.1885	0.0096216\\
5.2048	0.0091777\\
5.2048	0.0087984\\
5.1885	0.0092288\\
5.1721	0.0096803\\
5.1558	0.010154\\
5.1395	0.010651\\
5.1232	0.011172\\
5.1068	0.011718\\
5.0905	0.012292\\
5.0742	0.012893\\
5.0579	0.013524\\
5.0415	0.014186\\
5.0252	0.01488\\
5.0089	0.015608\\
4.9926	0.016372\\
4.9763	0.017173\\
4.9599	0.018014\\
4.9436	0.018896\\
4.9273	0.019821\\
4.911	0.020791\\
4.8946	0.021808\\
4.8783	0.022876\\
4.862	0.023996\\
4.8457	0.025171\\
4.8294	0.026403\\
4.813	0.027696\\
4.7967	0.029052\\
4.7804	0.030474\\
4.7641	0.031967\\
4.7477	0.033532\\
4.7314	0.035174\\
4.7151	0.036897\\
4.6988	0.038704\\
4.6824	0.040599\\
4.6661	0.042588\\
4.6498	0.044673\\
4.6335	0.046862\\
4.6172	0.049157\\
4.6008	0.051565\\
4.5845	0.054091\\
4.5682	0.056741\\
4.5519	0.059521\\
4.5355	0.062437\\
4.5192	0.065496\\
4.5029	0.068705\\
4.4866	0.072071\\
4.4703	0.075602\\
4.4539	0.079307\\
4.4376	0.083193\\
4.4213	0.08727\\
4.405	0.091546\\
4.3886	0.096033\\
4.3723	0.10074\\
4.356	0.10568\\
4.3397	0.11086\\
4.3233	0.11629\\
4.307	0.12199\\
4.2907	0.12797\\
4.2744	0.13424\\
4.2581	0.14082\\
4.2417	0.14772\\
4.2254	0.15496\\
4.2091	0.16256\\
4.1928	0.17053\\
4.1764	0.17889\\
4.1601	0.18766\\
4.1438	0.19686\\
4.1275	0.20651\\
4.1111	0.21663\\
4.0948	0.22724\\
4.0785	0.23838\\
4.0622	0.25006\\
4.0459	0.26231\\
4.0295	0.27516\\
4.0132	0.28863\\
3.9969	0.30276\\
3.9806	0.31757\\
3.9642	0.33309\\
3.9479	0.34937\\
3.9316	0.36642\\
3.9153	0.38428\\
3.899	0.403\\
3.8826	0.42262\\
3.8663	0.44318\\
3.85	0.46472\\
3.8337	0.48731\\
3.8173	0.51098\\
3.801	0.5358\\
3.7847	0.56183\\
3.7684	0.58911\\
3.752	0.61771\\
3.7357	0.6477\\
3.7194	0.67914\\
3.7031	0.71211\\
3.6868	0.74668\\
3.6704	0.78293\\
3.6541	0.82093\\
3.6378	0.86079\\
3.6215	0.90257\\
3.6051	0.94613\\
3.5888	0.98368\\
}--cycle;
\addplot[only marks, mark=*, mark size=1.1180pt, color=mycolor3, fill=mycolor3, opacity=0.60, draw=none, mark options={draw=none,line width=0pt}] table[row sep=crcr]{%
x	y\\
3.5886	0.99468\\
3.5916	0.98404\\
3.5952	0.9734\\
3.5993	0.96277\\
3.6029	0.95213\\
3.6065	0.94149\\
3.6092	0.93085\\
3.6129	0.92021\\
3.6165	0.90957\\
3.6212	0.89894\\
3.6249	0.8883\\
3.6294	0.87766\\
3.6332	0.86702\\
3.6375	0.85638\\
3.642	0.84574\\
3.647	0.83511\\
3.6521	0.82447\\
3.6574	0.81383\\
3.6623	0.80319\\
3.6677	0.79255\\
3.673	0.78191\\
3.6783	0.77128\\
3.6852	0.76064\\
3.6917	0.75\\
3.6975	0.73936\\
3.7051	0.72872\\
3.7118	0.71809\\
3.7198	0.70745\\
3.7274	0.69681\\
3.7344	0.68617\\
3.7419	0.67553\\
3.7498	0.66489\\
3.7583	0.65426\\
3.7667	0.64362\\
3.7752	0.63298\\
3.7842	0.62234\\
3.7915	0.6117\\
3.7992	0.60106\\
3.8062	0.59043\\
3.8131	0.57979\\
3.8207	0.56915\\
3.8275	0.55851\\
3.835	0.54787\\
3.8424	0.53723\\
3.849	0.5266\\
3.8549	0.51596\\
3.8615	0.50532\\
3.8671	0.49468\\
3.8725	0.48404\\
3.8792	0.4734\\
3.8855	0.46277\\
3.8914	0.45213\\
3.8973	0.44149\\
3.9044	0.43085\\
3.9103	0.42021\\
3.9165	0.40957\\
3.9223	0.39894\\
3.9291	0.3883\\
3.9357	0.37766\\
3.9442	0.36702\\
3.952	0.35638\\
3.9612	0.34574\\
3.9706	0.33511\\
3.9834	0.32447\\
4.0008	0.31383\\
4.0156	0.30319\\
4.0357	0.29255\\
4.0537	0.28191\\
4.067	0.27128\\
4.0828	0.26064\\
4.0979	0.25\\
4.1112	0.23936\\
4.1231	0.22872\\
4.1341	0.21809\\
4.1468	0.20745\\
4.1597	0.19681\\
4.1722	0.18617\\
4.1856	0.17553\\
4.1984	0.16489\\
4.2174	0.15426\\
4.2391	0.14362\\
4.264	0.13298\\
4.2878	0.12234\\
4.3125	0.1117\\
4.3462	0.10106\\
4.3689	0.090426\\
4.389	0.079787\\
4.4095	0.069149\\
4.4462	0.058511\\
4.5305	0.047872\\
4.5839	0.037234\\
4.7235	0.026596\\
4.7846	0.015957\\
5.1606	0.0053191\\
};
\addlegendentry{\ApproxEmpiricCCDFLegend}

\addplot [color=mycolor3, only marks, every error bar/.append style={opacity=0.30}, mark=*, mark size=0pt, draw=none, forget plot]
 plot [error bars/.cd, x dir=both, x explicit, error bar style={line width=1pt, color=mycolor3}, error mark options={mark=none,mark size=0pt}]
 table[row sep=crcr, x error plus index=2, x error minus index=3]{%
3.5886	0.99468	0.0016453	0.0016453\\
3.5916	0.98404	0.0015936	0.0015936\\
3.5952	0.9734	0.0016864	0.0016864\\
3.5993	0.96277	0.001728	0.001728\\
3.6029	0.95213	0.0018264	0.0018264\\
3.6065	0.94149	0.0018754	0.0018754\\
3.6092	0.93085	0.0018953	0.0018953\\
3.6129	0.92021	0.0019758	0.0019758\\
3.6165	0.90957	0.0020323	0.0020323\\
3.6212	0.89894	0.0020596	0.0020596\\
3.6249	0.8883	0.0020847	0.0020847\\
3.6294	0.87766	0.0019316	0.0019316\\
3.6332	0.86702	0.0019421	0.0019421\\
3.6375	0.85638	0.002004	0.002004\\
3.642	0.84574	0.0020857	0.0020857\\
3.647	0.83511	0.0020635	0.0020635\\
3.6521	0.82447	0.0021601	0.0021601\\
3.6574	0.81383	0.0023654	0.0023654\\
3.6623	0.80319	0.0024984	0.0024984\\
3.6677	0.79255	0.0025274	0.0025274\\
3.673	0.78191	0.0022797	0.0022797\\
3.6783	0.77128	0.0023897	0.0023897\\
3.6852	0.76064	0.0024132	0.0024132\\
3.6917	0.75	0.0023622	0.0023622\\
3.6975	0.73936	0.0023958	0.0023958\\
3.7051	0.72872	0.002518	0.002518\\
3.7118	0.71809	0.0026086	0.0026086\\
3.7198	0.70745	0.0025221	0.0025221\\
3.7274	0.69681	0.0027598	0.0027598\\
3.7344	0.68617	0.0027338	0.0027338\\
3.7419	0.67553	0.0028359	0.0028359\\
3.7498	0.66489	0.0027646	0.0027646\\
3.7583	0.65426	0.0025985	0.0025985\\
3.7667	0.64362	0.0027915	0.0027915\\
3.7752	0.63298	0.0027864	0.0027864\\
3.7842	0.62234	0.0026958	0.0026958\\
3.7915	0.6117	0.0025913	0.0025913\\
3.7992	0.60106	0.0026056	0.0026056\\
3.8062	0.59043	0.0024877	0.0024877\\
3.8131	0.57979	0.0023548	0.0023548\\
3.8207	0.56915	0.0024251	0.0024251\\
3.8275	0.55851	0.0024283	0.0024283\\
3.835	0.54787	0.0023464	0.0023464\\
3.8424	0.53723	0.0024311	0.0024311\\
3.849	0.5266	0.0024573	0.0024573\\
3.8549	0.51596	0.0025555	0.0025555\\
3.8615	0.50532	0.0023886	0.0023886\\
3.8671	0.49468	0.0022532	0.0022532\\
3.8725	0.48404	0.0023285	0.0023285\\
3.8792	0.4734	0.0024835	0.0024835\\
3.8855	0.46277	0.0025563	0.0025563\\
3.8914	0.45213	0.0025453	0.0025453\\
3.8973	0.44149	0.0024626	0.0024626\\
3.9044	0.43085	0.0024371	0.0024371\\
3.9103	0.42021	0.0024808	0.0024808\\
3.9165	0.40957	0.0024173	0.0024173\\
3.9223	0.39894	0.0025136	0.0025136\\
3.9291	0.3883	0.0026601	0.0026601\\
3.9357	0.37766	0.0026816	0.0026816\\
3.9442	0.36702	0.002656	0.002656\\
3.952	0.35638	0.0027536	0.0027536\\
3.9612	0.34574	0.0027851	0.0027851\\
3.9706	0.33511	0.0028519	0.0028519\\
3.9834	0.32447	0.0028767	0.0028767\\
4.0008	0.31383	0.0031239	0.0031239\\
4.0156	0.30319	0.0032166	0.0032166\\
4.0357	0.29255	0.0032434	0.0032434\\
4.0537	0.28191	0.003496	0.003496\\
4.067	0.27128	0.00362	0.00362\\
4.0828	0.26064	0.0034719	0.0034719\\
4.0979	0.25	0.0034767	0.0034767\\
4.1112	0.23936	0.0030726	0.0030726\\
4.1231	0.22872	0.0029265	0.0029265\\
4.1341	0.21809	0.0030447	0.0030447\\
4.1468	0.20745	0.0027872	0.0027872\\
4.1597	0.19681	0.0030099	0.0030099\\
4.1722	0.18617	0.0034031	0.0034031\\
4.1856	0.17553	0.0033091	0.0033091\\
4.1984	0.16489	0.0037203	0.0037203\\
4.2174	0.15426	0.0036906	0.0036906\\
4.2391	0.14362	0.003829	0.003829\\
4.264	0.13298	0.0037271	0.0037271\\
4.2878	0.12234	0.0039995	0.0039995\\
4.3125	0.1117	0.0041657	0.0041657\\
4.3462	0.10106	0.004	0.004\\
4.3689	0.090426	0.003943	0.003943\\
4.389	0.079787	0.0037316	0.0037316\\
4.4095	0.069149	0.0040381	0.0040381\\
4.4462	0.058511	0.0052621	0.0052621\\
4.5305	0.047872	0.0061177	0.0061177\\
4.5839	0.037234	0.0060254	0.0060254\\
4.7235	0.026596	0.0050878	0.0050878\\
4.7846	0.015957	0.0047448	0.0047448\\
5.1606	0.0053191	0.0037077	0.0037077\\
};
\addplot [color=mycolor3]
  table[row sep=crcr]{%
3.5875	0.98981\\
3.6038	0.9523\\
4.6167	0.049989\\
5.2048	0.0090481\\
};
\addlegendentry{\ApproxParetoFitLegend}


\addplot[area legend, draw=none, fill=mycolor3, fill opacity=0.15, forget plot]
table[row sep=crcr] {%
x	y\\
3.5875	0.99262\\
3.6038	0.95656\\
3.6202	0.91195\\
3.6365	0.86942\\
3.6528	0.82888\\
3.6692	0.79023\\
3.6855	0.75338\\
3.7018	0.71825\\
3.7182	0.68477\\
3.7345	0.65285\\
3.7508	0.62242\\
3.7672	0.59341\\
3.7835	0.56576\\
3.7999	0.53941\\
3.8162	0.51428\\
3.8325	0.49034\\
3.8489	0.46751\\
3.8652	0.44576\\
3.8815	0.42502\\
3.8979	0.40526\\
3.9142	0.38643\\
3.9305	0.36848\\
3.9469	0.35138\\
3.9632	0.33508\\
3.9796	0.31954\\
3.9959	0.30474\\
4.0122	0.29062\\
4.0286	0.27716\\
4.0449	0.26433\\
4.0612	0.25209\\
4.0776	0.24043\\
4.0939	0.2293\\
4.1102	0.2187\\
4.1266	0.20858\\
4.1429	0.19893\\
4.1593	0.18973\\
4.1756	0.18096\\
4.1919	0.17259\\
4.2083	0.16461\\
4.2246	0.157\\
4.2409	0.14974\\
4.2573	0.14282\\
4.2736	0.13622\\
4.2899	0.12992\\
4.3063	0.12392\\
4.3226	0.11819\\
4.339	0.11273\\
4.3553	0.10752\\
4.3716	0.10255\\
4.388	0.097811\\
4.4043	0.093291\\
4.4206	0.08898\\
4.437	0.084868\\
4.4533	0.080947\\
4.4696	0.077207\\
4.486	0.07364\\
4.5023	0.070237\\
4.5187	0.066992\\
4.535	0.063897\\
4.5513	0.060945\\
4.5677	0.05813\\
4.584	0.055445\\
4.6003	0.052883\\
4.6167	0.05044\\
4.633	0.04811\\
4.6493	0.045888\\
4.6657	0.043769\\
4.682	0.041747\\
4.6984	0.039819\\
4.7147	0.03798\\
4.731	0.036226\\
4.7474	0.034553\\
4.7637	0.032957\\
4.78	0.031435\\
4.7964	0.029983\\
4.8127	0.028599\\
4.829	0.027278\\
4.8454	0.026018\\
4.8617	0.024817\\
4.8781	0.023671\\
4.8944	0.022578\\
4.9107	0.021535\\
4.9271	0.020541\\
4.9434	0.019593\\
4.9597	0.018688\\
4.9761	0.017825\\
4.9924	0.017002\\
5.0087	0.016217\\
5.0251	0.015469\\
5.0414	0.014755\\
5.0577	0.014074\\
5.0741	0.013424\\
5.0904	0.012804\\
5.1068	0.012213\\
5.1231	0.011649\\
5.1394	0.011112\\
5.1558	0.010599\\
5.1721	0.01011\\
5.1884	0.0096431\\
5.2048	0.009198\\
5.2048	0.0088982\\
5.1884	0.0093327\\
5.1721	0.0097884\\
5.1558	0.010266\\
5.1394	0.010768\\
5.1231	0.011293\\
5.1068	0.011845\\
5.0904	0.012423\\
5.0741	0.01303\\
5.0577	0.013666\\
5.0414	0.014334\\
5.0251	0.015034\\
5.0087	0.015768\\
4.9924	0.016538\\
4.9761	0.017346\\
4.9597	0.018193\\
4.9434	0.019082\\
4.9271	0.020014\\
4.9107	0.020992\\
4.8944	0.022017\\
4.8781	0.023092\\
4.8617	0.024221\\
4.8454	0.025404\\
4.829	0.026645\\
4.8127	0.027946\\
4.7964	0.029312\\
4.78	0.030744\\
4.7637	0.032246\\
4.7474	0.033821\\
4.731	0.035473\\
4.7147	0.037207\\
4.6984	0.039024\\
4.682	0.040931\\
4.6657	0.042931\\
4.6493	0.045029\\
4.633	0.047229\\
4.6167	0.049537\\
4.6003	0.051958\\
4.584	0.054496\\
4.5677	0.05716\\
4.5513	0.059953\\
4.535	0.062883\\
4.5187	0.065956\\
4.5023	0.069179\\
4.486	0.07256\\
4.4696	0.076106\\
4.4533	0.079825\\
4.437	0.083727\\
4.4206	0.087819\\
4.4043	0.092111\\
4.388	0.096613\\
4.3716	0.10133\\
4.3553	0.10629\\
4.339	0.11148\\
4.3226	0.11693\\
4.3063	0.12265\\
4.2899	0.12864\\
4.2736	0.13493\\
4.2573	0.14152\\
4.2409	0.14844\\
4.2246	0.1557\\
4.2083	0.16331\\
4.1919	0.17129\\
4.1756	0.17966\\
4.1593	0.18844\\
4.1429	0.19765\\
4.1266	0.20731\\
4.1102	0.21744\\
4.0939	0.22806\\
4.0776	0.2392\\
4.0612	0.25089\\
4.0449	0.26314\\
4.0286	0.27599\\
4.0122	0.28946\\
3.9959	0.30359\\
3.9796	0.3184\\
3.9632	0.33393\\
3.9469	0.3502\\
3.9305	0.36726\\
3.9142	0.38514\\
3.8979	0.40388\\
3.8815	0.42353\\
3.8652	0.44412\\
3.8489	0.4657\\
3.8325	0.48832\\
3.8162	0.51203\\
3.7999	0.5369\\
3.7835	0.56296\\
3.7672	0.59028\\
3.7508	0.61893\\
3.7345	0.64897\\
3.7182	0.68045\\
3.7018	0.71347\\
3.6855	0.74809\\
3.6692	0.78438\\
3.6528	0.82243\\
3.6365	0.86233\\
3.6202	0.90417\\
3.6038	0.94803\\
3.5875	0.98701\\
}--cycle;
      % \input{../../../../.tex/20/KS/SizeDistributionFittingsfig7.tex}
  
      \nextgroupplot[%
        xmin=3.5608,xmax=5.3544,
        % xmin=3.4,xmax=5.4,
        % ymin=0.0053191,ymax=1,
      ] % This file was created by matlab2tikz.
%
\definecolor{mycolor1}{rgb}{0.00000,0.44706,0.69804}%
\definecolor{mycolor2}{rgb}{0.00000,0.44700,0.74100}%
\definecolor{mycolor3}{rgb}{0.49400,0.18400,0.55600}%
%
\addplot[only marks, mark=*, mark size=1.1180pt, color=mycolor1, fill=mycolor1, opacity=0.60, draw=none, mark options={draw=none,line width=0pt}] table[row sep=crcr]{%
x	y\\
3.5893	0.99468\\
3.5929	0.98404\\
3.5962	0.9734\\
3.6001	0.96277\\
3.6042	0.95213\\
3.6081	0.94149\\
3.6124	0.93085\\
3.6162	0.92021\\
3.6204	0.90957\\
3.625	0.89894\\
3.6291	0.8883\\
3.6335	0.87766\\
3.6385	0.86702\\
3.6436	0.85638\\
3.6481	0.84574\\
3.6533	0.83511\\
3.6583	0.82447\\
3.6634	0.81383\\
3.669	0.80319\\
3.6741	0.79255\\
3.6786	0.78191\\
3.683	0.77128\\
3.6891	0.76064\\
3.6957	0.75\\
3.7019	0.73936\\
3.7088	0.72872\\
3.7138	0.71809\\
3.7192	0.70745\\
3.7256	0.69681\\
3.7321	0.68617\\
3.7405	0.67553\\
3.7483	0.66489\\
3.7548	0.65426\\
3.763	0.64362\\
3.7715	0.63298\\
3.7797	0.62234\\
3.786	0.6117\\
3.7935	0.60106\\
3.8025	0.59043\\
3.8109	0.57979\\
3.819	0.56915\\
3.8255	0.55851\\
3.8327	0.54787\\
3.8385	0.53723\\
3.8461	0.5266\\
3.852	0.51596\\
3.86	0.50532\\
3.866	0.49468\\
3.8727	0.48404\\
3.8788	0.4734\\
3.8839	0.46277\\
3.8896	0.45213\\
3.8958	0.44149\\
3.9012	0.43085\\
3.9081	0.42021\\
3.9155	0.40957\\
3.9216	0.39894\\
3.928	0.3883\\
3.9348	0.37766\\
3.9415	0.36702\\
3.9492	0.35638\\
3.9579	0.34574\\
3.9692	0.33511\\
3.9828	0.32447\\
3.9977	0.31383\\
4.0161	0.30319\\
4.0335	0.29255\\
4.0503	0.28191\\
4.0662	0.27128\\
4.0836	0.26064\\
4.0988	0.25\\
4.1112	0.23936\\
4.1243	0.22872\\
4.1342	0.21809\\
4.1454	0.20745\\
4.1576	0.19681\\
4.1705	0.18617\\
4.1844	0.17553\\
4.2007	0.16489\\
4.2194	0.15426\\
4.24	0.14362\\
4.2662	0.13298\\
4.2921	0.12234\\
4.3201	0.1117\\
4.3525	0.10106\\
4.3736	0.090426\\
4.3895	0.079787\\
4.4094	0.069149\\
4.4358	0.058511\\
4.5276	0.047872\\
4.5794	0.037234\\
4.7208	0.026596\\
4.7905	0.015957\\
5.1574	0.0053191\\
};
\addlegendentry{\ExactEmpiricCCDFLegend}

\addplot [color=mycolor1, only marks, every error bar/.append style={opacity=0.30}, mark=*, mark size=0pt, draw=none, forget plot]
 plot [error bars/.cd, x dir=both, x explicit, error bar style={line width=1pt, color=mycolor1}, error mark options={mark=none,mark size=0pt}]
 table[row sep=crcr, x error plus index=2, x error minus index=3]{%
3.5893	0.99468	0.0020617	0.0020617\\
3.5929	0.98404	0.0020269	0.0020269\\
3.5962	0.9734	0.0020561	0.0020561\\
3.6001	0.96277	0.0019457	0.0019457\\
3.6042	0.95213	0.0019054	0.0019054\\
3.6081	0.94149	0.0018168	0.0018168\\
3.6124	0.93085	0.0018416	0.0018416\\
3.6162	0.92021	0.0018589	0.0018589\\
3.6204	0.90957	0.0018898	0.0018898\\
3.625	0.89894	0.0020758	0.0020758\\
3.6291	0.8883	0.0019633	0.0019633\\
3.6335	0.87766	0.0019186	0.0019186\\
3.6385	0.86702	0.0019754	0.0019754\\
3.6436	0.85638	0.002189	0.002189\\
3.6481	0.84574	0.0020486	0.0020486\\
3.6533	0.83511	0.0020882	0.0020882\\
3.6583	0.82447	0.0020115	0.0020115\\
3.6634	0.81383	0.0020684	0.0020684\\
3.669	0.80319	0.0022894	0.0022894\\
3.6741	0.79255	0.0022671	0.0022671\\
3.6786	0.78191	0.0021921	0.0021921\\
3.683	0.77128	0.0021802	0.0021802\\
3.6891	0.76064	0.0021046	0.0021046\\
3.6957	0.75	0.0020569	0.0020569\\
3.7019	0.73936	0.0021172	0.0021172\\
3.7088	0.72872	0.0020608	0.0020608\\
3.7138	0.71809	0.0020375	0.0020375\\
3.7192	0.70745	0.0020778	0.0020778\\
3.7256	0.69681	0.0020958	0.0020958\\
3.7321	0.68617	0.0022399	0.0022399\\
3.7405	0.67553	0.0023519	0.0023519\\
3.7483	0.66489	0.0024659	0.0024659\\
3.7548	0.65426	0.0025499	0.0025499\\
3.763	0.64362	0.0027554	0.0027554\\
3.7715	0.63298	0.002923	0.002923\\
3.7797	0.62234	0.0028095	0.0028095\\
3.786	0.6117	0.0027316	0.0027316\\
3.7935	0.60106	0.002762	0.002762\\
3.8025	0.59043	0.0028749	0.0028749\\
3.8109	0.57979	0.003014	0.003014\\
3.819	0.56915	0.0027807	0.0027807\\
3.8255	0.55851	0.0025636	0.0025636\\
3.8327	0.54787	0.002599	0.002599\\
3.8385	0.53723	0.0024753	0.0024753\\
3.8461	0.5266	0.0023333	0.0023333\\
3.852	0.51596	0.0022884	0.0022884\\
3.86	0.50532	0.0022617	0.0022617\\
3.866	0.49468	0.0022885	0.0022885\\
3.8727	0.48404	0.0022411	0.0022411\\
3.8788	0.4734	0.0022106	0.0022106\\
3.8839	0.46277	0.0022586	0.0022586\\
3.8896	0.45213	0.0024622	0.0024622\\
3.8958	0.44149	0.0024482	0.0024482\\
3.9012	0.43085	0.0023713	0.0023713\\
3.9081	0.42021	0.0023255	0.0023255\\
3.9155	0.40957	0.0023465	0.0023465\\
3.9216	0.39894	0.0023832	0.0023832\\
3.928	0.3883	0.002503	0.002503\\
3.9348	0.37766	0.0025268	0.0025268\\
3.9415	0.36702	0.0027627	0.0027627\\
3.9492	0.35638	0.0027797	0.0027797\\
3.9579	0.34574	0.0028551	0.0028551\\
3.9692	0.33511	0.0029051	0.0029051\\
3.9828	0.32447	0.0031558	0.0031558\\
3.9977	0.31383	0.003395	0.003395\\
4.0161	0.30319	0.0037587	0.0037587\\
4.0335	0.29255	0.0037678	0.0037678\\
4.0503	0.28191	0.0031906	0.0031906\\
4.0662	0.27128	0.0031187	0.0031187\\
4.0836	0.26064	0.003391	0.003391\\
4.0988	0.25	0.0031882	0.0031882\\
4.1112	0.23936	0.0033285	0.0033285\\
4.1243	0.22872	0.0032658	0.0032658\\
4.1342	0.21809	0.0030368	0.0030368\\
4.1454	0.20745	0.0028444	0.0028444\\
4.1576	0.19681	0.0028028	0.0028028\\
4.1705	0.18617	0.0031669	0.0031669\\
4.1844	0.17553	0.0030161	0.0030161\\
4.2007	0.16489	0.0030855	0.0030855\\
4.2194	0.15426	0.0036315	0.0036315\\
4.24	0.14362	0.0036555	0.0036555\\
4.2662	0.13298	0.0036608	0.0036608\\
4.2921	0.12234	0.0040354	0.0040354\\
4.3201	0.1117	0.004327	0.004327\\
4.3525	0.10106	0.0037223	0.0037223\\
4.3736	0.090426	0.0034217	0.0034217\\
4.3895	0.079787	0.0036394	0.0036394\\
4.4094	0.069149	0.0039009	0.0039009\\
4.4358	0.058511	0.0043355	0.0043355\\
4.5276	0.047872	0.0054577	0.0054577\\
4.5794	0.037234	0.0042494	0.0042494\\
4.7208	0.026596	0.0050438	0.0050438\\
4.7905	0.015957	0.0050741	0.0050741\\
5.1574	0.0053191	0.0034798	0.0034798\\
};
\addplot [color=mycolor1]
  table[row sep=crcr]{%
3.5888	0.98701\\
3.6062	0.94863\\
3.7627	0.60087\\
4.5451	0.061344\\
5.3101	0.0066156\\
};
\addlegendentry{\ExactParetoFitLegend}


\addplot[area legend, draw=none, fill=mycolor1, fill opacity=0.15, forget plot]
table[row sep=crcr] {%
x	y\\
3.5888	0.99033\\
3.6062	0.95394\\
3.6236	0.9068\\
3.641	0.86167\\
3.6584	0.81879\\
3.6758	0.77804\\
3.6931	0.73933\\
3.7105	0.70255\\
3.7279	0.6676\\
3.7453	0.63439\\
3.7627	0.60284\\
3.7801	0.57287\\
3.7975	0.54439\\
3.8149	0.51733\\
3.8322	0.49162\\
3.8496	0.46721\\
3.867	0.44401\\
3.8844	0.42198\\
3.9018	0.40106\\
3.9192	0.3812\\
3.9366	0.36233\\
3.9539	0.34442\\
3.9713	0.32741\\
3.9887	0.31126\\
4.0061	0.29591\\
4.0235	0.28133\\
4.0409	0.26747\\
4.0583	0.2543\\
4.0757	0.24178\\
4.093	0.22988\\
4.1104	0.21856\\
4.1278	0.20781\\
4.1452	0.19758\\
4.1626	0.18786\\
4.18	0.17862\\
4.1974	0.16984\\
4.2148	0.16149\\
4.2321	0.15354\\
4.2495	0.14599\\
4.2669	0.13881\\
4.2843	0.13199\\
4.3017	0.1255\\
4.3191	0.11933\\
4.3365	0.11346\\
4.3538	0.10789\\
4.3712	0.10258\\
4.3886	0.097541\\
4.406	0.092747\\
4.4234	0.088188\\
4.4408	0.083854\\
4.4582	0.079733\\
4.4756	0.075815\\
4.4929	0.07209\\
4.5103	0.068547\\
4.5277	0.065179\\
4.5451	0.061977\\
4.5625	0.058932\\
4.5799	0.056037\\
4.5973	0.053284\\
4.6147	0.050666\\
4.632	0.048177\\
4.6494	0.045811\\
4.6668	0.04356\\
4.6842	0.041421\\
4.7016	0.039386\\
4.719	0.037452\\
4.7364	0.035613\\
4.7538	0.033864\\
4.7711	0.032201\\
4.7885	0.03062\\
4.8059	0.029116\\
4.8233	0.027686\\
4.8407	0.026327\\
4.8581	0.025034\\
4.8755	0.023805\\
4.8928	0.022637\\
4.9102	0.021525\\
4.9276	0.020469\\
4.945	0.019464\\
4.9624	0.018508\\
4.9798	0.0176\\
4.9972	0.016736\\
5.0146	0.015915\\
5.0319	0.015134\\
5.0493	0.014391\\
5.0667	0.013685\\
5.0841	0.013013\\
5.1015	0.012375\\
5.1189	0.011767\\
5.1363	0.01119\\
5.1537	0.010641\\
5.171	0.010119\\
5.1884	0.0096223\\
5.2058	0.0091503\\
5.2232	0.0087014\\
5.2406	0.0082745\\
5.258	0.0078686\\
5.2754	0.0074826\\
5.2927	0.0071156\\
5.3101	0.0067666\\
5.3101	0.0064646\\
5.2927	0.0068019\\
5.2754	0.0071568\\
5.258	0.0075303\\
5.2406	0.0079232\\
5.2232	0.0083367\\
5.2058	0.0087718\\
5.1884	0.0092295\\
5.171	0.0097112\\
5.1537	0.010218\\
5.1363	0.010751\\
5.1189	0.011313\\
5.1015	0.011903\\
5.0841	0.012525\\
5.0667	0.013178\\
5.0493	0.013866\\
5.0319	0.01459\\
5.0146	0.015352\\
4.9972	0.016154\\
4.9798	0.016997\\
4.9624	0.017885\\
4.945	0.018819\\
4.9276	0.019801\\
4.9102	0.020835\\
4.8928	0.021924\\
4.8755	0.023069\\
4.8581	0.024273\\
4.8407	0.025541\\
4.8233	0.026875\\
4.8059	0.028279\\
4.7885	0.029756\\
4.7711	0.031311\\
4.7538	0.032947\\
4.7364	0.034668\\
4.719	0.036479\\
4.7016	0.038385\\
4.6842	0.040391\\
4.6668	0.042501\\
4.6494	0.044722\\
4.632	0.047059\\
4.6147	0.049519\\
4.5973	0.052107\\
4.5799	0.05483\\
4.5625	0.057695\\
4.5451	0.060711\\
4.5277	0.063884\\
4.5103	0.067224\\
4.4929	0.070738\\
4.4756	0.074435\\
4.4582	0.078327\\
4.4408	0.082421\\
4.4234	0.086731\\
4.406	0.091265\\
4.3886	0.096037\\
4.3712	0.10106\\
4.3538	0.10634\\
4.3365	0.1119\\
4.3191	0.11775\\
4.3017	0.12391\\
4.2843	0.13039\\
4.2669	0.13721\\
4.2495	0.14439\\
4.2321	0.15194\\
4.2148	0.15988\\
4.1974	0.16824\\
4.18	0.17704\\
4.1626	0.1863\\
4.1452	0.19604\\
4.1278	0.2063\\
4.1104	0.21708\\
4.093	0.22843\\
4.0757	0.24038\\
4.0583	0.25294\\
4.0409	0.26616\\
4.0235	0.28007\\
4.0061	0.2947\\
3.9887	0.31008\\
3.9713	0.32626\\
3.9539	0.34327\\
3.9366	0.36115\\
3.9192	0.37994\\
3.9018	0.39969\\
3.8844	0.42045\\
3.867	0.44227\\
3.8496	0.46521\\
3.8322	0.48933\\
3.8149	0.51469\\
3.7975	0.54136\\
3.7801	0.56941\\
3.7627	0.59891\\
3.7453	0.62993\\
3.7279	0.66256\\
3.7105	0.69687\\
3.6931	0.73296\\
3.6758	0.77092\\
3.6584	0.81085\\
3.641	0.85284\\
3.6236	0.89701\\
3.6062	0.94332\\
3.5888	0.98368\\
}--cycle;
\addplot[only marks, mark=*, mark size=1.1180pt, color=black, fill=black, opacity=0.60, draw=none, mark options={draw=none,line width=0pt}] table[row sep=crcr]{%
x	y\\
3.5623	0.99468\\
3.5723	0.98404\\
3.573	0.9734\\
3.5753	0.96277\\
3.5766	0.95213\\
3.5781	0.94149\\
3.5826	0.93085\\
3.5851	0.92021\\
3.5911	0.90957\\
3.5942	0.89894\\
3.596	0.8883\\
3.6016	0.87766\\
3.603	0.86702\\
3.6047	0.85638\\
3.6086	0.84574\\
3.6086	0.83511\\
3.6097	0.82447\\
3.6142	0.81383\\
3.6261	0.80319\\
3.6383	0.79255\\
3.6464	0.78191\\
3.6535	0.77128\\
3.657	0.76064\\
3.6658	0.75\\
3.6678	0.73936\\
3.6687	0.72872\\
3.6702	0.71809\\
3.6723	0.70745\\
3.6793	0.69681\\
3.6867	0.68617\\
3.6949	0.67553\\
3.7015	0.66489\\
3.7031	0.65426\\
3.719	0.64362\\
3.7213	0.63298\\
3.7265	0.62234\\
3.7307	0.6117\\
3.737	0.60106\\
3.7374	0.59043\\
3.7423	0.57979\\
3.7487	0.56915\\
3.7677	0.55851\\
3.7685	0.54787\\
3.7768	0.53723\\
3.8005	0.5266\\
3.8015	0.51596\\
3.8022	0.50532\\
3.8032	0.49468\\
3.8039	0.48404\\
3.8108	0.4734\\
3.8214	0.46277\\
3.8405	0.45213\\
3.8605	0.44149\\
3.863	0.43085\\
3.8645	0.42021\\
3.8662	0.40957\\
3.8806	0.39894\\
3.8878	0.3883\\
3.8964	0.37766\\
3.905	0.36702\\
3.9294	0.35638\\
3.9303	0.34574\\
3.954	0.33511\\
3.9604	0.32447\\
3.9669	0.31383\\
3.9747	0.30319\\
3.9822	0.29255\\
3.9929	0.28191\\
4.0051	0.27128\\
4.0144	0.26064\\
4.0161	0.25\\
4.0433	0.23936\\
4.0578	0.22872\\
4.073	0.21809\\
4.0857	0.20745\\
4.0904	0.19681\\
4.1168	0.18617\\
4.1265	0.17553\\
4.127	0.16489\\
4.143	0.15426\\
4.1756	0.14362\\
4.2156	0.13298\\
4.3116	0.12234\\
4.3276	0.1117\\
4.3662	0.10106\\
4.4791	0.090426\\
4.4912	0.079787\\
4.517	0.069149\\
4.5743	0.058511\\
4.5914	0.047872\\
4.6138	0.037234\\
4.7898	0.026596\\
5.0876	0.015957\\
5.3101	0.0053191\\
};
\addlegendentry{\RealEmpiricCCDFLegend}

\addplot [color=black]
  table[row sep=crcr]{%
3.5608	1\\
5.3101	0.005939\\
};
\addlegendentry{\RealParetoFitLegend}

      % \input{../../../../.tex/20/KS/SizeDistributionFittingsfig8.tex}
  
      \nextgroupplot[%
        xmin=3.5608,xmax=5.4115,
        % xmin=3.4,xmax=5.4,
        % ymin=0.0053191,ymax=1,
      ] % This file was created by matlab2tikz.
%
\definecolor{mycolor1}{rgb}{0.83529,0.36863,0.00000}%
\definecolor{mycolor2}{rgb}{0.00000,0.44700,0.74100}%
\definecolor{mycolor3}{rgb}{0.49400,0.18400,0.55600}%
%
\addplot[only marks, mark=*, mark size=1.1180pt, color=mycolor1, fill=mycolor1, opacity=0.60, draw=none, mark options={draw=none,line width=0pt}] table[row sep=crcr]{%
x	y\\
3.5886	0.99468\\
3.5916	0.98404\\
3.5952	0.9734\\
3.5993	0.96277\\
3.6029	0.95213\\
3.6065	0.94149\\
3.6092	0.93085\\
3.6129	0.92021\\
3.6165	0.90957\\
3.6212	0.89894\\
3.6249	0.8883\\
3.6294	0.87766\\
3.6332	0.86702\\
3.6375	0.85638\\
3.642	0.84574\\
3.647	0.83511\\
3.6521	0.82447\\
3.6574	0.81383\\
3.6623	0.80319\\
3.6677	0.79255\\
3.673	0.78191\\
3.6783	0.77128\\
3.6852	0.76064\\
3.6917	0.75\\
3.6975	0.73936\\
3.7051	0.72872\\
3.7118	0.71809\\
3.7198	0.70745\\
3.7274	0.69681\\
3.7344	0.68617\\
3.7419	0.67553\\
3.7498	0.66489\\
3.7583	0.65426\\
3.7667	0.64362\\
3.7752	0.63298\\
3.7842	0.62234\\
3.7915	0.6117\\
3.7992	0.60106\\
3.8062	0.59043\\
3.8131	0.57979\\
3.8207	0.56915\\
3.8275	0.55851\\
3.835	0.54787\\
3.8424	0.53723\\
3.849	0.5266\\
3.8549	0.51596\\
3.8615	0.50532\\
3.8671	0.49468\\
3.8725	0.48404\\
3.8792	0.4734\\
3.8855	0.46277\\
3.8914	0.45213\\
3.8973	0.44149\\
3.9044	0.43085\\
3.9103	0.42021\\
3.9165	0.40957\\
3.9223	0.39894\\
3.9291	0.3883\\
3.9357	0.37766\\
3.9442	0.36702\\
3.952	0.35638\\
3.9612	0.34574\\
3.9706	0.33511\\
3.9834	0.32447\\
4.0008	0.31383\\
4.0156	0.30319\\
4.0357	0.29255\\
4.0537	0.28191\\
4.067	0.27128\\
4.0828	0.26064\\
4.0979	0.25\\
4.1112	0.23936\\
4.1231	0.22872\\
4.1341	0.21809\\
4.1468	0.20745\\
4.1597	0.19681\\
4.1722	0.18617\\
4.1856	0.17553\\
4.1984	0.16489\\
4.2174	0.15426\\
4.2391	0.14362\\
4.264	0.13298\\
4.2878	0.12234\\
4.3125	0.1117\\
4.3462	0.10106\\
4.3689	0.090426\\
4.389	0.079787\\
4.4095	0.069149\\
4.4462	0.058511\\
4.5305	0.047872\\
4.5839	0.037234\\
4.7235	0.026596\\
4.7846	0.015957\\
5.1606	0.0053191\\
};
\addlegendentry{\ApproxEmpiricCCDFLegend}

\addplot [color=mycolor1, only marks, every error bar/.append style={opacity=0.30}, mark=*, mark size=0pt, draw=none, forget plot]
 plot [error bars/.cd, x dir=both, x explicit, error bar style={line width=1pt, color=mycolor1}, error mark options={mark=none,mark size=0pt}]
 table[row sep=crcr, x error plus index=2, x error minus index=3]{%
3.5886	0.99468	0.0016453	0.0016453\\
3.5916	0.98404	0.0015936	0.0015936\\
3.5952	0.9734	0.0016864	0.0016864\\
3.5993	0.96277	0.001728	0.001728\\
3.6029	0.95213	0.0018264	0.0018264\\
3.6065	0.94149	0.0018754	0.0018754\\
3.6092	0.93085	0.0018953	0.0018953\\
3.6129	0.92021	0.0019758	0.0019758\\
3.6165	0.90957	0.0020323	0.0020323\\
3.6212	0.89894	0.0020596	0.0020596\\
3.6249	0.8883	0.0020847	0.0020847\\
3.6294	0.87766	0.0019316	0.0019316\\
3.6332	0.86702	0.0019421	0.0019421\\
3.6375	0.85638	0.002004	0.002004\\
3.642	0.84574	0.0020857	0.0020857\\
3.647	0.83511	0.0020635	0.0020635\\
3.6521	0.82447	0.0021601	0.0021601\\
3.6574	0.81383	0.0023654	0.0023654\\
3.6623	0.80319	0.0024984	0.0024984\\
3.6677	0.79255	0.0025274	0.0025274\\
3.673	0.78191	0.0022797	0.0022797\\
3.6783	0.77128	0.0023897	0.0023897\\
3.6852	0.76064	0.0024132	0.0024132\\
3.6917	0.75	0.0023622	0.0023622\\
3.6975	0.73936	0.0023958	0.0023958\\
3.7051	0.72872	0.002518	0.002518\\
3.7118	0.71809	0.0026086	0.0026086\\
3.7198	0.70745	0.0025221	0.0025221\\
3.7274	0.69681	0.0027598	0.0027598\\
3.7344	0.68617	0.0027338	0.0027338\\
3.7419	0.67553	0.0028359	0.0028359\\
3.7498	0.66489	0.0027646	0.0027646\\
3.7583	0.65426	0.0025985	0.0025985\\
3.7667	0.64362	0.0027915	0.0027915\\
3.7752	0.63298	0.0027864	0.0027864\\
3.7842	0.62234	0.0026958	0.0026958\\
3.7915	0.6117	0.0025913	0.0025913\\
3.7992	0.60106	0.0026056	0.0026056\\
3.8062	0.59043	0.0024877	0.0024877\\
3.8131	0.57979	0.0023548	0.0023548\\
3.8207	0.56915	0.0024251	0.0024251\\
3.8275	0.55851	0.0024283	0.0024283\\
3.835	0.54787	0.0023464	0.0023464\\
3.8424	0.53723	0.0024311	0.0024311\\
3.849	0.5266	0.0024573	0.0024573\\
3.8549	0.51596	0.0025555	0.0025555\\
3.8615	0.50532	0.0023886	0.0023886\\
3.8671	0.49468	0.0022532	0.0022532\\
3.8725	0.48404	0.0023285	0.0023285\\
3.8792	0.4734	0.0024835	0.0024835\\
3.8855	0.46277	0.0025563	0.0025563\\
3.8914	0.45213	0.0025453	0.0025453\\
3.8973	0.44149	0.0024626	0.0024626\\
3.9044	0.43085	0.0024371	0.0024371\\
3.9103	0.42021	0.0024808	0.0024808\\
3.9165	0.40957	0.0024173	0.0024173\\
3.9223	0.39894	0.0025136	0.0025136\\
3.9291	0.3883	0.0026601	0.0026601\\
3.9357	0.37766	0.0026816	0.0026816\\
3.9442	0.36702	0.002656	0.002656\\
3.952	0.35638	0.0027536	0.0027536\\
3.9612	0.34574	0.0027851	0.0027851\\
3.9706	0.33511	0.0028519	0.0028519\\
3.9834	0.32447	0.0028767	0.0028767\\
4.0008	0.31383	0.0031239	0.0031239\\
4.0156	0.30319	0.0032166	0.0032166\\
4.0357	0.29255	0.0032434	0.0032434\\
4.0537	0.28191	0.003496	0.003496\\
4.067	0.27128	0.00362	0.00362\\
4.0828	0.26064	0.0034719	0.0034719\\
4.0979	0.25	0.0034767	0.0034767\\
4.1112	0.23936	0.0030726	0.0030726\\
4.1231	0.22872	0.0029265	0.0029265\\
4.1341	0.21809	0.0030447	0.0030447\\
4.1468	0.20745	0.0027872	0.0027872\\
4.1597	0.19681	0.0030099	0.0030099\\
4.1722	0.18617	0.0034031	0.0034031\\
4.1856	0.17553	0.0033091	0.0033091\\
4.1984	0.16489	0.0037203	0.0037203\\
4.2174	0.15426	0.0036906	0.0036906\\
4.2391	0.14362	0.003829	0.003829\\
4.264	0.13298	0.0037271	0.0037271\\
4.2878	0.12234	0.0039995	0.0039995\\
4.3125	0.1117	0.0041657	0.0041657\\
4.3462	0.10106	0.004	0.004\\
4.3689	0.090426	0.003943	0.003943\\
4.389	0.079787	0.0037316	0.0037316\\
4.4095	0.069149	0.0040381	0.0040381\\
4.4462	0.058511	0.0052621	0.0052621\\
4.5305	0.047872	0.0061177	0.0061177\\
4.5839	0.037234	0.0060254	0.0060254\\
4.7235	0.026596	0.0050878	0.0050878\\
4.7846	0.015957	0.0047448	0.0047448\\
5.1606	0.0053191	0.0037077	0.0037077\\
};
\addplot [color=mycolor1]
  table[row sep=crcr]{%
3.5875	0.98981\\
3.6049	0.94935\\
4.6141	0.050361\\
5.3101	0.0066625\\
};
\addlegendentry{\ApproxParetoFitLegend}


\addplot[area legend, draw=none, fill=mycolor1, fill opacity=0.15, forget plot]
table[row sep=crcr] {%
x	y\\
3.5875	0.99262\\
3.6049	0.95359\\
3.6223	0.90629\\
3.6397	0.86134\\
3.6571	0.81863\\
3.6745	0.77803\\
3.6919	0.73945\\
3.7093	0.70279\\
3.7267	0.66794\\
3.7441	0.63483\\
3.7615	0.60337\\
3.7789	0.57347\\
3.7963	0.54505\\
3.8137	0.51805\\
3.8311	0.4924\\
3.8485	0.46802\\
3.8659	0.44486\\
3.8833	0.42285\\
3.9007	0.40195\\
3.9181	0.38209\\
3.9355	0.36322\\
3.9529	0.34529\\
3.9703	0.32826\\
3.9877	0.31208\\
4.0051	0.2967\\
4.0225	0.28208\\
4.0399	0.26819\\
4.0573	0.25499\\
4.0747	0.24244\\
4.0921	0.23051\\
4.1095	0.21917\\
4.1269	0.20839\\
4.1443	0.19814\\
4.1617	0.18839\\
4.1791	0.17913\\
4.1965	0.17032\\
4.2139	0.16194\\
4.2313	0.15398\\
4.2487	0.14641\\
4.2661	0.13921\\
4.2835	0.13237\\
4.3009	0.12586\\
4.3183	0.11968\\
4.3357	0.11379\\
4.3531	0.1082\\
4.3705	0.10288\\
4.3879	0.097826\\
4.4053	0.093018\\
4.4227	0.088447\\
4.4401	0.084101\\
4.4575	0.079968\\
4.4749	0.076038\\
4.4923	0.072302\\
4.5097	0.068749\\
4.5271	0.065371\\
4.5445	0.062159\\
4.5619	0.059105\\
4.5793	0.056202\\
4.5967	0.05344\\
4.6141	0.050815\\
4.6315	0.048319\\
4.6489	0.045945\\
4.6663	0.043688\\
4.6837	0.041542\\
4.7011	0.039501\\
4.7185	0.037561\\
4.7359	0.035716\\
4.7533	0.033962\\
4.7707	0.032294\\
4.7881	0.030708\\
4.8055	0.0292\\
4.8229	0.027765\\
4.8403	0.026402\\
4.8577	0.025105\\
4.8751	0.023872\\
4.8925	0.0227\\
4.9099	0.021585\\
4.9273	0.020525\\
4.9447	0.019518\\
4.9621	0.018559\\
4.9795	0.017648\\
4.9969	0.016781\\
5.0143	0.015958\\
5.0317	0.015174\\
5.0491	0.014429\\
5.0665	0.013721\\
5.0839	0.013047\\
5.1013	0.012407\\
5.1187	0.011798\\
5.1361	0.011218\\
5.1535	0.010668\\
5.1709	0.010144\\
5.1883	0.0096461\\
5.2057	0.0091726\\
5.2231	0.0087224\\
5.2405	0.0082943\\
5.2579	0.0078872\\
5.2753	0.0075001\\
5.2927	0.007132\\
5.3101	0.006782\\
5.3101	0.0065431\\
5.2927	0.0068839\\
5.2753	0.0072424\\
5.2579	0.0076196\\
5.2405	0.0080165\\
5.2231	0.0084341\\
5.2057	0.0088734\\
5.1883	0.0093356\\
5.1709	0.0098219\\
5.1535	0.010334\\
5.1361	0.010872\\
5.1187	0.011438\\
5.1013	0.012034\\
5.0839	0.012661\\
5.0665	0.013321\\
5.0491	0.014015\\
5.0317	0.014745\\
5.0143	0.015513\\
4.9969	0.016321\\
4.9795	0.017172\\
4.9621	0.018067\\
4.9447	0.019008\\
4.9273	0.019999\\
4.9099	0.021041\\
4.8925	0.022137\\
4.8751	0.023291\\
4.8577	0.024505\\
4.8403	0.025782\\
4.8229	0.027125\\
4.8055	0.028539\\
4.7881	0.030026\\
4.7707	0.031591\\
4.7533	0.033238\\
4.7359	0.03497\\
4.7185	0.036793\\
4.7011	0.038711\\
4.6837	0.040728\\
4.6663	0.042851\\
4.6489	0.045085\\
4.6315	0.047435\\
4.6141	0.049908\\
4.5967	0.05251\\
4.5793	0.055247\\
4.5619	0.058127\\
4.5445	0.061158\\
4.5271	0.064346\\
4.5097	0.067701\\
4.4923	0.07123\\
4.4749	0.074944\\
4.4575	0.078851\\
4.4401	0.082963\\
4.4227	0.087288\\
4.4053	0.091839\\
4.3879	0.096628\\
4.3705	0.10167\\
4.3531	0.10697\\
4.3357	0.11255\\
4.3183	0.11841\\
4.3009	0.12459\\
4.2835	0.13108\\
4.2661	0.13792\\
4.2487	0.14511\\
4.2313	0.15268\\
4.2139	0.16064\\
4.1965	0.16901\\
4.1791	0.17783\\
4.1617	0.1871\\
4.1443	0.19685\\
4.1269	0.20712\\
4.1095	0.21791\\
4.0921	0.22927\\
4.0747	0.24122\\
4.0573	0.25379\\
4.0399	0.26701\\
4.0225	0.28092\\
4.0051	0.29555\\
3.9877	0.31093\\
3.9703	0.32711\\
3.9529	0.34412\\
3.9355	0.36201\\
3.9181	0.38081\\
3.9007	0.40059\\
3.8833	0.42137\\
3.8659	0.44322\\
3.8485	0.4662\\
3.8311	0.49036\\
3.8137	0.51577\\
3.7963	0.54248\\
3.7789	0.57057\\
3.7615	0.60011\\
3.7441	0.63118\\
3.7267	0.66386\\
3.7093	0.69822\\
3.6919	0.73436\\
3.6745	0.77237\\
3.6571	0.81234\\
3.6397	0.85439\\
3.6223	0.8986\\
3.6049	0.94511\\
3.5875	0.98701\\
}--cycle;
\addplot[only marks, mark=*, mark size=1.1180pt, color=black, fill=black, opacity=0.60, draw=none, mark options={draw=none,line width=0pt}] table[row sep=crcr]{%
x	y\\
3.5623	0.99468\\
3.5723	0.98404\\
3.573	0.9734\\
3.5753	0.96277\\
3.5766	0.95213\\
3.5781	0.94149\\
3.5826	0.93085\\
3.5851	0.92021\\
3.5911	0.90957\\
3.5942	0.89894\\
3.596	0.8883\\
3.6016	0.87766\\
3.603	0.86702\\
3.6047	0.85638\\
3.6086	0.84574\\
3.6086	0.83511\\
3.6097	0.82447\\
3.6142	0.81383\\
3.6261	0.80319\\
3.6383	0.79255\\
3.6464	0.78191\\
3.6535	0.77128\\
3.657	0.76064\\
3.6658	0.75\\
3.6678	0.73936\\
3.6687	0.72872\\
3.6702	0.71809\\
3.6723	0.70745\\
3.6793	0.69681\\
3.6867	0.68617\\
3.6949	0.67553\\
3.7015	0.66489\\
3.7031	0.65426\\
3.719	0.64362\\
3.7213	0.63298\\
3.7265	0.62234\\
3.7307	0.6117\\
3.737	0.60106\\
3.7374	0.59043\\
3.7423	0.57979\\
3.7487	0.56915\\
3.7677	0.55851\\
3.7685	0.54787\\
3.7768	0.53723\\
3.8005	0.5266\\
3.8015	0.51596\\
3.8022	0.50532\\
3.8032	0.49468\\
3.8039	0.48404\\
3.8108	0.4734\\
3.8214	0.46277\\
3.8405	0.45213\\
3.8605	0.44149\\
3.863	0.43085\\
3.8645	0.42021\\
3.8662	0.40957\\
3.8806	0.39894\\
3.8878	0.3883\\
3.8964	0.37766\\
3.905	0.36702\\
3.9294	0.35638\\
3.9303	0.34574\\
3.954	0.33511\\
3.9604	0.32447\\
3.9669	0.31383\\
3.9747	0.30319\\
3.9822	0.29255\\
3.9929	0.28191\\
4.0051	0.27128\\
4.0144	0.26064\\
4.0161	0.25\\
4.0433	0.23936\\
4.0578	0.22872\\
4.073	0.21809\\
4.0857	0.20745\\
4.0904	0.19681\\
4.1168	0.18617\\
4.1265	0.17553\\
4.127	0.16489\\
4.143	0.15426\\
4.1756	0.14362\\
4.2156	0.13298\\
4.3116	0.12234\\
4.3276	0.1117\\
4.3662	0.10106\\
4.4791	0.090426\\
4.4912	0.079787\\
4.517	0.069149\\
4.5743	0.058511\\
4.5914	0.047872\\
4.6138	0.037234\\
4.7898	0.026596\\
5.0876	0.015957\\
5.3101	0.0053191\\
};
\addlegendentry{\RealEmpiricCCDFLegend}

\addplot [color=black]
  table[row sep=crcr]{%
3.5608	1\\
5.3101	0.005939\\
};
\addlegendentry{\RealParetoFitLegend}

      % \input{../../../../.tex/20/KS/SizeDistributionFittingsfig9.tex}
    \end{groupplot}
  
    \begin{groupplot}[
      group style={
        group name=Row2,
        plotGroup,
      },
      plotRow,
      row2KeysWithLabels,
      plotLegendNE,
      %
      % xmin=3637.5,xmax=2.5003e+05,
      ymin=0.0053191,ymax=1,
    ]
      \nextgroupplot[%
        at={($(Row1 c1r1.south west)-(0,\yPlotSep em)-(0,\plotHeight cm)$)},
        xmin=3.598,xmax=5.4115,
        % xmin=3.4,xmax=5.4,
        % ymin=0.0053191,ymax=1,
      ] % This file was created by matlab2tikz.
%
\definecolor{mycolor1}{rgb}{0.00000,0.44706,0.69804}%
\definecolor{mycolor2}{rgb}{0.00000,0.44700,0.74100}%
\definecolor{mycolor3}{rgb}{0.83529,0.36863,0.00000}%
\definecolor{mycolor4}{rgb}{0.49400,0.18400,0.55600}%
%
\addplot[only marks, mark=*, mark size=1.1180pt, color=mycolor1, fill=mycolor1, opacity=0.60, draw=none, mark options={draw=none,line width=0pt}] table[row sep=crcr]{%
x	y\\
3.6037	0.99468\\
3.6073	0.98404\\
3.6108	0.9734\\
3.6145	0.96277\\
3.6181	0.95213\\
3.622	0.94149\\
3.6258	0.93085\\
3.6297	0.92021\\
3.6339	0.90957\\
3.6376	0.89894\\
3.6419	0.8883\\
3.6462	0.87766\\
3.6499	0.86702\\
3.6544	0.85638\\
3.6601	0.84574\\
3.6647	0.83511\\
3.6691	0.82447\\
3.6744	0.81383\\
3.6795	0.80319\\
3.6851	0.79255\\
3.6902	0.78191\\
3.6946	0.77128\\
3.701	0.76064\\
3.7067	0.75\\
3.7119	0.73936\\
3.7176	0.72872\\
3.7232	0.71809\\
3.7287	0.70745\\
3.7342	0.69681\\
3.7398	0.68617\\
3.7458	0.67553\\
3.7524	0.66489\\
3.7597	0.65426\\
3.7658	0.64362\\
3.773	0.63298\\
3.7802	0.62234\\
3.7873	0.6117\\
3.7936	0.60106\\
3.8011	0.59043\\
3.8081	0.57979\\
3.8151	0.56915\\
3.8222	0.55851\\
3.8279	0.54787\\
3.8349	0.53723\\
3.8422	0.5266\\
3.85	0.51596\\
3.8572	0.50532\\
3.8649	0.49468\\
3.8714	0.48404\\
3.8791	0.4734\\
3.8881	0.46277\\
3.8956	0.45213\\
3.9016	0.44149\\
3.91	0.43085\\
3.9175	0.42021\\
3.9272	0.40957\\
3.9348	0.39894\\
3.9449	0.3883\\
3.9548	0.37766\\
3.9657	0.36702\\
3.9754	0.35638\\
3.9863	0.34574\\
3.9973	0.33511\\
4.0075	0.32447\\
4.0175	0.31383\\
4.0299	0.30319\\
4.0406	0.29255\\
4.0543	0.28191\\
4.0681	0.27128\\
4.0824	0.26064\\
4.0965	0.25\\
4.1083	0.23936\\
4.1212	0.22872\\
4.138	0.21809\\
4.1548	0.20745\\
4.1697	0.19681\\
4.1848	0.18617\\
4.2011	0.17553\\
4.2197	0.16489\\
4.2359	0.15426\\
4.2571	0.14362\\
4.278	0.13298\\
4.3018	0.12234\\
4.3248	0.1117\\
4.3516	0.10106\\
4.3803	0.090426\\
4.4066	0.079787\\
4.4414	0.069149\\
4.4829	0.058511\\
4.538	0.047872\\
4.6052	0.037234\\
4.7174	0.026596\\
4.8035	0.015957\\
5.159	0.0053191\\
};
\addlegendentry{\ExactEmpiricCCDFLegend}

\addplot [color=mycolor1, only marks, every error bar/.append style={opacity=0.30}, mark=*, mark size=0pt, draw=none, forget plot]
 plot [error bars/.cd, x dir=both, x explicit, error bar style={line width=1pt, color=mycolor1}, error mark options={mark=none,mark size=0pt}]
 table[row sep=crcr, x error plus index=2, x error minus index=3]{%
3.6037	0.99468	0.0031873	0.0031873\\
3.6073	0.98404	0.0032744	0.0032744\\
3.6108	0.9734	0.0033314	0.0033314\\
3.6145	0.96277	0.0034058	0.0034058\\
3.6181	0.95213	0.0034851	0.0034851\\
3.622	0.94149	0.0035	0.0035\\
3.6258	0.93085	0.0036615	0.0036615\\
3.6297	0.92021	0.0035244	0.0035244\\
3.6339	0.90957	0.0034445	0.0034445\\
3.6376	0.89894	0.0035539	0.0035539\\
3.6419	0.8883	0.0034716	0.0034716\\
3.6462	0.87766	0.0034679	0.0034679\\
3.6499	0.86702	0.0034622	0.0034622\\
3.6544	0.85638	0.0034513	0.0034513\\
3.6601	0.84574	0.0034059	0.0034059\\
3.6647	0.83511	0.003463	0.003463\\
3.6691	0.82447	0.0034186	0.0034186\\
3.6744	0.81383	0.003484	0.003484\\
3.6795	0.80319	0.0035559	0.0035559\\
3.6851	0.79255	0.0035894	0.0035894\\
3.6902	0.78191	0.0034953	0.0034953\\
3.6946	0.77128	0.0034979	0.0034979\\
3.701	0.76064	0.003521	0.003521\\
3.7067	0.75	0.0035332	0.0035332\\
3.7119	0.73936	0.0036677	0.0036677\\
3.7176	0.72872	0.0037632	0.0037632\\
3.7232	0.71809	0.0038613	0.0038613\\
3.7287	0.70745	0.0039704	0.0039704\\
3.7342	0.69681	0.0038912	0.0038912\\
3.7398	0.68617	0.0037488	0.0037488\\
3.7458	0.67553	0.0038371	0.0038371\\
3.7524	0.66489	0.0038114	0.0038114\\
3.7597	0.65426	0.0036991	0.0036991\\
3.7658	0.64362	0.0035915	0.0035915\\
3.773	0.63298	0.0034558	0.0034558\\
3.7802	0.62234	0.0035273	0.0035273\\
3.7873	0.6117	0.0037482	0.0037482\\
3.7936	0.60106	0.0040024	0.0040024\\
3.8011	0.59043	0.0041113	0.0041113\\
3.8081	0.57979	0.0041614	0.0041614\\
3.8151	0.56915	0.0042496	0.0042496\\
3.8222	0.55851	0.0043104	0.0043104\\
3.8279	0.54787	0.0043573	0.0043573\\
3.8349	0.53723	0.0041869	0.0041869\\
3.8422	0.5266	0.0040847	0.0040847\\
3.85	0.51596	0.003834	0.003834\\
3.8572	0.50532	0.0039218	0.0039218\\
3.8649	0.49468	0.0038832	0.0038832\\
3.8714	0.48404	0.003944	0.003944\\
3.8791	0.4734	0.0037933	0.0037933\\
3.8881	0.46277	0.0035871	0.0035871\\
3.8956	0.45213	0.0035147	0.0035147\\
3.9016	0.44149	0.0037251	0.0037251\\
3.91	0.43085	0.0038328	0.0038328\\
3.9175	0.42021	0.0041924	0.0041924\\
3.9272	0.40957	0.0043005	0.0043005\\
3.9348	0.39894	0.0042934	0.0042934\\
3.9449	0.3883	0.0044537	0.0044537\\
3.9548	0.37766	0.0047398	0.0047398\\
3.9657	0.36702	0.004744	0.004744\\
3.9754	0.35638	0.0046879	0.0046879\\
3.9863	0.34574	0.0045711	0.0045711\\
3.9973	0.33511	0.0047913	0.0047913\\
4.0075	0.32447	0.0047427	0.0047427\\
4.0175	0.31383	0.0050029	0.0050029\\
4.0299	0.30319	0.0051338	0.0051338\\
4.0406	0.29255	0.005174	0.005174\\
4.0543	0.28191	0.0053199	0.0053199\\
4.0681	0.27128	0.0056242	0.0056242\\
4.0824	0.26064	0.005481	0.005481\\
4.0965	0.25	0.0054067	0.0054067\\
4.1083	0.23936	0.0051031	0.0051031\\
4.1212	0.22872	0.0051613	0.0051613\\
4.138	0.21809	0.0048839	0.0048839\\
4.1548	0.20745	0.0052666	0.0052666\\
4.1697	0.19681	0.0050782	0.0050782\\
4.1848	0.18617	0.0052326	0.0052326\\
4.2011	0.17553	0.0048643	0.0048643\\
4.2197	0.16489	0.0049985	0.0049985\\
4.2359	0.15426	0.00532	0.00532\\
4.2571	0.14362	0.0053879	0.0053879\\
4.278	0.13298	0.0055106	0.0055106\\
4.3018	0.12234	0.0062277	0.0062277\\
4.3248	0.1117	0.0057529	0.0057529\\
4.3516	0.10106	0.00653	0.00653\\
4.3803	0.090426	0.0065879	0.0065879\\
4.4066	0.079787	0.0066297	0.0066297\\
4.4414	0.069149	0.0075694	0.0075694\\
4.4829	0.058511	0.0081425	0.0081425\\
4.538	0.047872	0.0091626	0.0091626\\
4.6052	0.037234	0.0096523	0.0096523\\
4.7174	0.026596	0.011213	0.011213\\
4.8035	0.015957	0.010763	0.010763\\
5.159	0.0053191	0.0090493	0.0090493\\
};
\addplot [color=mycolor1]
  table[row sep=crcr]{%
3.602	0.98157\\
3.6191	0.95011\\
3.6361	0.90544\\
4.1826	0.17977\\
4.729	0.035843\\
5.2754	0.0071764\\
5.2924	0.0068251\\
};
\addlegendentry{\ExactParetoFitLegend}


\addplot[area legend, draw=none, fill=mycolor1, fill opacity=0.15, forget plot]
table[row sep=crcr] {%
x	y\\
3.602	0.98693\\
3.6191	0.95791\\
3.6361	0.91319\\
3.6532	0.86785\\
3.6703	0.82477\\
3.6874	0.78384\\
3.7044	0.74494\\
3.7215	0.70799\\
3.7386	0.67288\\
3.7557	0.63951\\
3.7727	0.60782\\
3.7898	0.5777\\
3.8069	0.54909\\
3.824	0.52191\\
3.841	0.49608\\
3.8581	0.47155\\
3.8752	0.44824\\
3.8923	0.4261\\
3.9094	0.40506\\
3.9264	0.38508\\
3.9435	0.3661\\
3.9606	0.34807\\
3.9777	0.33094\\
3.9947	0.31467\\
4.0118	0.29921\\
4.0289	0.28451\\
4.046	0.27055\\
4.063	0.25728\\
4.0801	0.24467\\
4.0972	0.23268\\
4.1143	0.22128\\
4.1313	0.21045\\
4.1484	0.20015\\
4.1655	0.19036\\
4.1826	0.18104\\
4.1996	0.17219\\
4.2167	0.16377\\
4.2338	0.15577\\
4.2509	0.14815\\
4.2679	0.14091\\
4.285	0.13403\\
4.3021	0.12748\\
4.3192	0.12126\\
4.3362	0.11534\\
4.3533	0.1097\\
4.3704	0.10435\\
4.3875	0.099256\\
4.4045	0.094411\\
4.4216	0.089804\\
4.4387	0.085422\\
4.4558	0.081254\\
4.4728	0.07729\\
4.4899	0.07352\\
4.507	0.069934\\
4.5241	0.066524\\
4.5411	0.06328\\
4.5582	0.060194\\
4.5753	0.05726\\
4.5924	0.054468\\
4.6094	0.051813\\
4.6265	0.049288\\
4.6436	0.046886\\
4.6607	0.044601\\
4.6777	0.042428\\
4.6948	0.040361\\
4.7119	0.038394\\
4.729	0.036524\\
4.746	0.034745\\
4.7631	0.033053\\
4.7802	0.031443\\
4.7973	0.029912\\
4.8143	0.028455\\
4.8314	0.02707\\
4.8485	0.025752\\
4.8656	0.024499\\
4.8826	0.023306\\
4.8997	0.022172\\
4.9168	0.021093\\
4.9339	0.020066\\
4.9509	0.01909\\
4.968	0.018161\\
4.9851	0.017278\\
5.0022	0.016437\\
5.0192	0.015638\\
5.0363	0.014877\\
5.0534	0.014154\\
5.0705	0.013465\\
5.0875	0.012811\\
5.1046	0.012188\\
5.1217	0.011595\\
5.1388	0.011032\\
5.1558	0.010495\\
5.1729	0.0099853\\
5.19	0.0095\\
5.2071	0.0090384\\
5.2241	0.0085992\\
5.2412	0.0081814\\
5.2583	0.0077839\\
5.2754	0.0074058\\
5.2924	0.0070461\\
5.2924	0.0066041\\
5.2754	0.006947\\
5.2583	0.0073077\\
5.2412	0.0076871\\
5.2241	0.0080864\\
5.2071	0.0085063\\
5.19	0.0089481\\
5.1729	0.0094129\\
5.1558	0.0099019\\
5.1388	0.010416\\
5.1217	0.010958\\
5.1046	0.011527\\
5.0875	0.012126\\
5.0705	0.012756\\
5.0534	0.013419\\
5.0363	0.014116\\
5.0192	0.01485\\
5.0022	0.015622\\
4.9851	0.016434\\
4.968	0.017288\\
4.9509	0.018187\\
4.9339	0.019133\\
4.9168	0.020128\\
4.8997	0.021174\\
4.8826	0.022276\\
4.8656	0.023434\\
4.8485	0.024653\\
4.8314	0.025935\\
4.8143	0.027285\\
4.7973	0.028704\\
4.7802	0.030198\\
4.7631	0.031769\\
4.746	0.033422\\
4.729	0.035161\\
4.7119	0.036991\\
4.6948	0.038916\\
4.6777	0.040942\\
4.6607	0.043073\\
4.6436	0.045315\\
4.6265	0.047674\\
4.6094	0.050156\\
4.5924	0.052767\\
4.5753	0.055515\\
4.5582	0.058405\\
4.5411	0.061447\\
4.5241	0.064647\\
4.507	0.068013\\
4.4899	0.071556\\
4.4728	0.075283\\
4.4558	0.079204\\
4.4387	0.08333\\
4.4216	0.08767\\
4.4045	0.092237\\
4.3875	0.097043\\
4.3704	0.1021\\
4.3533	0.10742\\
4.3362	0.11301\\
4.3192	0.1189\\
4.3021	0.1251\\
4.285	0.13161\\
4.2679	0.13847\\
4.2509	0.14569\\
4.2338	0.15328\\
4.2167	0.16126\\
4.1996	0.16966\\
4.1826	0.1785\\
4.1655	0.1878\\
4.1484	0.19757\\
4.1313	0.20786\\
4.1143	0.21868\\
4.0972	0.23006\\
4.0801	0.24203\\
4.063	0.25462\\
4.046	0.26786\\
4.0289	0.28178\\
4.0118	0.29642\\
3.9947	0.31181\\
3.9777	0.32799\\
3.9606	0.345\\
3.9435	0.36287\\
3.9264	0.38167\\
3.9094	0.40142\\
3.8923	0.42218\\
3.8752	0.444\\
3.8581	0.46694\\
3.841	0.49105\\
3.824	0.5164\\
3.8069	0.54305\\
3.7898	0.57106\\
3.7727	0.60051\\
3.7557	0.63148\\
3.7386	0.66403\\
3.7215	0.69826\\
3.7044	0.73425\\
3.6874	0.77209\\
3.6703	0.81187\\
3.6532	0.85371\\
3.6361	0.8977\\
3.6191	0.94232\\
3.602	0.97621\\
}--cycle;
\addplot[only marks, mark=*, mark size=1.1180pt, color=mycolor3, fill=mycolor3, opacity=0.60, draw=none, mark options={draw=none,line width=0pt}] table[row sep=crcr]{%
x	y\\
3.5981	0.99468\\
3.6023	0.98404\\
3.6064	0.9734\\
3.6105	0.96277\\
3.615	0.95213\\
3.619	0.94149\\
3.6233	0.93085\\
3.6271	0.92021\\
3.6318	0.90957\\
3.6366	0.89894\\
3.6403	0.8883\\
3.6443	0.87766\\
3.6487	0.86702\\
3.6532	0.85638\\
3.6584	0.84574\\
3.6638	0.83511\\
3.6687	0.82447\\
3.6733	0.81383\\
3.6786	0.80319\\
3.6828	0.79255\\
3.6881	0.78191\\
3.6936	0.77128\\
3.6991	0.76064\\
3.7048	0.75\\
3.71	0.73936\\
3.7158	0.72872\\
3.721	0.71809\\
3.7265	0.70745\\
3.7321	0.69681\\
3.737	0.68617\\
3.7428	0.67553\\
3.7488	0.66489\\
3.7556	0.65426\\
3.763	0.64362\\
3.7694	0.63298\\
3.776	0.62234\\
3.7822	0.6117\\
3.7889	0.60106\\
3.7946	0.59043\\
3.801	0.57979\\
3.8083	0.56915\\
3.8176	0.55851\\
3.824	0.54787\\
3.8317	0.53723\\
3.8394	0.5266\\
3.8475	0.51596\\
3.8558	0.50532\\
3.8625	0.49468\\
3.8709	0.48404\\
3.8781	0.4734\\
3.8874	0.46277\\
3.8955	0.45213\\
3.9029	0.44149\\
3.9129	0.43085\\
3.9213	0.42021\\
3.9292	0.40957\\
3.9386	0.39894\\
3.9488	0.3883\\
3.9566	0.37766\\
3.9662	0.36702\\
3.9771	0.35638\\
3.9865	0.34574\\
3.9964	0.33511\\
4.0082	0.32447\\
4.0215	0.31383\\
4.0338	0.30319\\
4.0469	0.29255\\
4.0576	0.28191\\
4.0718	0.27128\\
4.0833	0.26064\\
4.0968	0.25\\
4.1097	0.23936\\
4.123	0.22872\\
4.1357	0.21809\\
4.1509	0.20745\\
4.1655	0.19681\\
4.1838	0.18617\\
4.2013	0.17553\\
4.2175	0.16489\\
4.2364	0.15426\\
4.2558	0.14362\\
4.2753	0.13298\\
4.3016	0.12234\\
4.3251	0.1117\\
4.3478	0.10106\\
4.3732	0.090426\\
4.4078	0.079787\\
4.4449	0.069149\\
4.4792	0.058511\\
4.5401	0.047872\\
4.6076	0.037234\\
4.7133	0.026596\\
4.7931	0.015957\\
5.1695	0.0053191\\
};
\addlegendentry{\ApproxEmpiricCCDFLegend}

\addplot [color=mycolor3, only marks, every error bar/.append style={opacity=0.30}, mark=*, mark size=0pt, draw=none, forget plot]
 plot [error bars/.cd, x dir=both, x explicit, error bar style={line width=1pt, color=mycolor3}, error mark options={mark=none,mark size=0pt}]
 table[row sep=crcr, x error plus index=2, x error minus index=3]{%
3.5981	0.99468	0.0033318	0.0033318\\
3.6023	0.98404	0.0032921	0.0032921\\
3.6064	0.9734	0.0032161	0.0032161\\
3.6105	0.96277	0.0033671	0.0033671\\
3.615	0.95213	0.0032769	0.0032769\\
3.619	0.94149	0.0031849	0.0031849\\
3.6233	0.93085	0.0031801	0.0031801\\
3.6271	0.92021	0.0032614	0.0032614\\
3.6318	0.90957	0.0033254	0.0033254\\
3.6366	0.89894	0.0032616	0.0032616\\
3.6403	0.8883	0.0034099	0.0034099\\
3.6443	0.87766	0.0035325	0.0035325\\
3.6487	0.86702	0.0034635	0.0034635\\
3.6532	0.85638	0.0034169	0.0034169\\
3.6584	0.84574	0.0034051	0.0034051\\
3.6638	0.83511	0.0032979	0.0032979\\
3.6687	0.82447	0.00348	0.00348\\
3.6733	0.81383	0.0035628	0.0035628\\
3.6786	0.80319	0.0034968	0.0034968\\
3.6828	0.79255	0.0034696	0.0034696\\
3.6881	0.78191	0.003529	0.003529\\
3.6936	0.77128	0.0037115	0.0037115\\
3.6991	0.76064	0.0036262	0.0036262\\
3.7048	0.75	0.0034662	0.0034662\\
3.71	0.73936	0.0034329	0.0034329\\
3.7158	0.72872	0.0033673	0.0033673\\
3.721	0.71809	0.0033237	0.0033237\\
3.7265	0.70745	0.0035375	0.0035375\\
3.7321	0.69681	0.0036761	0.0036761\\
3.737	0.68617	0.0037267	0.0037267\\
3.7428	0.67553	0.0038231	0.0038231\\
3.7488	0.66489	0.0038249	0.0038249\\
3.7556	0.65426	0.0039771	0.0039771\\
3.763	0.64362	0.0041257	0.0041257\\
3.7694	0.63298	0.0039609	0.0039609\\
3.776	0.62234	0.0039102	0.0039102\\
3.7822	0.6117	0.0039858	0.0039858\\
3.7889	0.60106	0.0042246	0.0042246\\
3.7946	0.59043	0.0041824	0.0041824\\
3.801	0.57979	0.00417	0.00417\\
3.8083	0.56915	0.0042139	0.0042139\\
3.8176	0.55851	0.0043078	0.0043078\\
3.824	0.54787	0.0042625	0.0042625\\
3.8317	0.53723	0.0043864	0.0043864\\
3.8394	0.5266	0.0043818	0.0043818\\
3.8475	0.51596	0.0043929	0.0043929\\
3.8558	0.50532	0.0045964	0.0045964\\
3.8625	0.49468	0.0045453	0.0045453\\
3.8709	0.48404	0.0045508	0.0045508\\
3.8781	0.4734	0.0044692	0.0044692\\
3.8874	0.46277	0.0044678	0.0044678\\
3.8955	0.45213	0.0044672	0.0044672\\
3.9029	0.44149	0.0044505	0.0044505\\
3.9129	0.43085	0.0043948	0.0043948\\
3.9213	0.42021	0.0042943	0.0042943\\
3.9292	0.40957	0.004448	0.004448\\
3.9386	0.39894	0.0045776	0.0045776\\
3.9488	0.3883	0.0046377	0.0046377\\
3.9566	0.37766	0.0044836	0.0044836\\
3.9662	0.36702	0.00447	0.00447\\
3.9771	0.35638	0.0045325	0.0045325\\
3.9865	0.34574	0.0047347	0.0047347\\
3.9964	0.33511	0.0049881	0.0049881\\
4.0082	0.32447	0.0053043	0.0053043\\
4.0215	0.31383	0.0054884	0.0054884\\
4.0338	0.30319	0.0055015	0.0055015\\
4.0469	0.29255	0.0056892	0.0056892\\
4.0576	0.28191	0.0055564	0.0055564\\
4.0718	0.27128	0.0057063	0.0057063\\
4.0833	0.26064	0.0058436	0.0058436\\
4.0968	0.25	0.0061932	0.0061932\\
4.1097	0.23936	0.0059823	0.0059823\\
4.123	0.22872	0.0058606	0.0058606\\
4.1357	0.21809	0.0056708	0.0056708\\
4.1509	0.20745	0.0058678	0.0058678\\
4.1655	0.19681	0.0053964	0.0053964\\
4.1838	0.18617	0.0053979	0.0053979\\
4.2013	0.17553	0.0050116	0.0050116\\
4.2175	0.16489	0.0053215	0.0053215\\
4.2364	0.15426	0.005493	0.005493\\
4.2558	0.14362	0.0055404	0.0055404\\
4.2753	0.13298	0.0054582	0.0054582\\
4.3016	0.12234	0.0058416	0.0058416\\
4.3251	0.1117	0.0059397	0.0059397\\
4.3478	0.10106	0.0058893	0.0058893\\
4.3732	0.090426	0.0064931	0.0064931\\
4.4078	0.079787	0.0070281	0.0070281\\
4.4449	0.069149	0.0074367	0.0074367\\
4.4792	0.058511	0.0080956	0.0080956\\
4.5401	0.047872	0.010773	0.010773\\
4.6076	0.037234	0.010679	0.010679\\
4.7133	0.026596	0.010343	0.010343\\
4.7931	0.015957	0.011898	0.011898\\
5.1695	0.0053191	0.0088428	0.0088428\\
};
\addplot [color=mycolor3]
  table[row sep=crcr]{%
3.5976	0.98063\\
3.6147	0.94676\\
3.6318	0.90374\\
4.1283	0.21165\\
4.6419	0.04735\\
5.1555	0.010638\\
5.2924	0.007149\\
};
\addlegendentry{\ApproxParetoFitLegend}


\addplot[area legend, draw=none, fill=mycolor3, fill opacity=0.15, forget plot]
table[row sep=crcr] {%
x	y\\
3.5976	0.98556\\
3.6147	0.95413\\
3.6318	0.9115\\
3.6489	0.86662\\
3.666	0.82396\\
3.6832	0.78341\\
3.7003	0.74487\\
3.7174	0.70823\\
3.7345	0.67341\\
3.7516	0.64032\\
3.7688	0.60887\\
3.7859	0.57897\\
3.803	0.55056\\
3.8201	0.52356\\
3.8372	0.4979\\
3.8544	0.47352\\
3.8715	0.45035\\
3.8886	0.42833\\
3.9057	0.40741\\
3.9228	0.38752\\
3.94	0.36863\\
3.9571	0.35067\\
3.9742	0.3336\\
3.9913	0.31737\\
4.0084	0.30194\\
4.0256	0.28728\\
4.0427	0.27333\\
4.0598	0.26007\\
4.0769	0.24745\\
4.094	0.23546\\
4.1112	0.22404\\
4.1283	0.21319\\
4.1454	0.20286\\
4.1625	0.19304\\
4.1796	0.1837\\
4.1968	0.17481\\
4.2139	0.16635\\
4.231	0.1583\\
4.2481	0.15064\\
4.2652	0.14336\\
4.2824	0.13643\\
4.2995	0.12983\\
4.3166	0.12355\\
4.3337	0.11758\\
4.3508	0.1119\\
4.368	0.10649\\
4.3851	0.10135\\
4.4022	0.096454\\
4.4193	0.091795\\
4.4364	0.087362\\
4.4536	0.083143\\
4.4707	0.079128\\
4.4878	0.075308\\
4.5049	0.071673\\
4.522	0.068213\\
4.5392	0.064921\\
4.5563	0.061788\\
4.5734	0.058806\\
4.5905	0.055969\\
4.6076	0.053269\\
4.6248	0.050699\\
4.6419	0.048254\\
4.659	0.045927\\
4.6761	0.043712\\
4.6932	0.041604\\
4.7104	0.039598\\
4.7275	0.037689\\
4.7446	0.035872\\
4.7617	0.034143\\
4.7788	0.032497\\
4.796	0.030931\\
4.8131	0.02944\\
4.8302	0.028022\\
4.8473	0.026672\\
4.8644	0.025387\\
4.8816	0.024164\\
4.8987	0.023\\
4.9158	0.021892\\
4.9329	0.020838\\
4.95	0.019834\\
4.9672	0.018879\\
4.9843	0.01797\\
5.0014	0.017105\\
5.0185	0.016282\\
5.0356	0.015498\\
5.0528	0.014752\\
5.0699	0.014042\\
5.087	0.013367\\
5.1041	0.012723\\
5.1212	0.012111\\
5.1384	0.011529\\
5.1555	0.010974\\
5.1726	0.010446\\
5.1897	0.0099438\\
5.2068	0.0094655\\
5.224	0.0090104\\
5.2411	0.0085771\\
5.2582	0.0081647\\
5.2753	0.0077722\\
5.2924	0.0073986\\
5.2924	0.0068993\\
5.2753	0.0072538\\
5.2582	0.0076264\\
5.2411	0.0080183\\
5.224	0.0084304\\
5.2068	0.0088636\\
5.1897	0.0093192\\
5.1726	0.0097982\\
5.1555	0.010302\\
5.1384	0.010832\\
5.1212	0.011388\\
5.1041	0.011974\\
5.087	0.01259\\
5.0699	0.013237\\
5.0528	0.013918\\
5.0356	0.014634\\
5.0185	0.015387\\
5.0014	0.016179\\
4.9843	0.017011\\
4.9672	0.017886\\
4.95	0.018807\\
4.9329	0.019775\\
4.9158	0.020793\\
4.8987	0.021863\\
4.8816	0.022989\\
4.8644	0.024172\\
4.8473	0.025417\\
4.8302	0.026726\\
4.8131	0.028102\\
4.796	0.02955\\
4.7788	0.031072\\
4.7617	0.032672\\
4.7446	0.034356\\
4.7275	0.036126\\
4.7104	0.037987\\
4.6932	0.039945\\
4.6761	0.042004\\
4.659	0.044169\\
4.6419	0.046445\\
4.6248	0.04884\\
4.6076	0.051358\\
4.5905	0.054005\\
4.5734	0.05679\\
4.5563	0.059719\\
4.5392	0.062799\\
4.522	0.066037\\
4.5049	0.069444\\
4.4878	0.073026\\
4.4707	0.076793\\
4.4536	0.080755\\
4.4364	0.084922\\
4.4193	0.089304\\
4.4022	0.093912\\
4.3851	0.098758\\
4.368	0.10386\\
4.3508	0.10922\\
4.3337	0.11485\\
4.3166	0.12078\\
4.2995	0.12702\\
4.2824	0.13357\\
4.2652	0.14047\\
4.2481	0.14772\\
4.231	0.15535\\
4.2139	0.16336\\
4.1968	0.1718\\
4.1796	0.18067\\
4.1625	0.18999\\
4.1454	0.1998\\
4.1283	0.21011\\
4.1112	0.22095\\
4.094	0.23235\\
4.0769	0.24434\\
4.0598	0.25694\\
4.0427	0.27019\\
4.0256	0.28412\\
4.0084	0.29875\\
3.9913	0.31414\\
3.9742	0.33031\\
3.9571	0.3473\\
3.94	0.36515\\
3.9228	0.38391\\
3.9057	0.40361\\
3.8886	0.42431\\
3.8715	0.44606\\
3.8544	0.46891\\
3.8372	0.49291\\
3.8201	0.51812\\
3.803	0.54461\\
3.7859	0.57244\\
3.7688	0.60168\\
3.7516	0.63241\\
3.7345	0.6647\\
3.7174	0.69862\\
3.7003	0.73428\\
3.6832	0.77174\\
3.666	0.81111\\
3.6489	0.85249\\
3.6318	0.89598\\
3.6147	0.9394\\
3.5976	0.9757\\
}--cycle;
      % \input{../../../../.tex/20/KS/SizeDistributionFittingsfig16.tex}
      \legend{}
  
      \nextgroupplot[%
        xmin=3.5608,xmax=5.3544,
        % xmin=3.4,xmax=5.4,
        % ymin=0.0053191,ymax=1,
      ] % This file was created by matlab2tikz.
%
\definecolor{mycolor1}{rgb}{0.00000,0.44706,0.69804}%
\definecolor{mycolor2}{rgb}{0.00000,0.44700,0.74100}%
\definecolor{mycolor3}{rgb}{0.49400,0.18400,0.55600}%
%
\addplot[only marks, mark=*, mark size=1.1180pt, color=mycolor1, fill=mycolor1, opacity=0.60, draw=none, mark options={draw=none,line width=0pt}] table[row sep=crcr]{%
x	y\\
3.6037	0.99468\\
3.6073	0.98404\\
3.6108	0.9734\\
3.6145	0.96277\\
3.6181	0.95213\\
3.622	0.94149\\
3.6258	0.93085\\
3.6297	0.92021\\
3.6339	0.90957\\
3.6376	0.89894\\
3.6419	0.8883\\
3.6462	0.87766\\
3.6499	0.86702\\
3.6544	0.85638\\
3.6601	0.84574\\
3.6647	0.83511\\
3.6691	0.82447\\
3.6744	0.81383\\
3.6795	0.80319\\
3.6851	0.79255\\
3.6902	0.78191\\
3.6946	0.77128\\
3.701	0.76064\\
3.7067	0.75\\
3.7119	0.73936\\
3.7176	0.72872\\
3.7232	0.71809\\
3.7287	0.70745\\
3.7342	0.69681\\
3.7398	0.68617\\
3.7458	0.67553\\
3.7524	0.66489\\
3.7597	0.65426\\
3.7658	0.64362\\
3.773	0.63298\\
3.7802	0.62234\\
3.7873	0.6117\\
3.7936	0.60106\\
3.8011	0.59043\\
3.8081	0.57979\\
3.8151	0.56915\\
3.8222	0.55851\\
3.8279	0.54787\\
3.8349	0.53723\\
3.8422	0.5266\\
3.85	0.51596\\
3.8572	0.50532\\
3.8649	0.49468\\
3.8714	0.48404\\
3.8791	0.4734\\
3.8881	0.46277\\
3.8956	0.45213\\
3.9016	0.44149\\
3.91	0.43085\\
3.9175	0.42021\\
3.9272	0.40957\\
3.9348	0.39894\\
3.9449	0.3883\\
3.9548	0.37766\\
3.9657	0.36702\\
3.9754	0.35638\\
3.9863	0.34574\\
3.9973	0.33511\\
4.0075	0.32447\\
4.0175	0.31383\\
4.0299	0.30319\\
4.0406	0.29255\\
4.0543	0.28191\\
4.0681	0.27128\\
4.0824	0.26064\\
4.0965	0.25\\
4.1083	0.23936\\
4.1212	0.22872\\
4.138	0.21809\\
4.1548	0.20745\\
4.1697	0.19681\\
4.1848	0.18617\\
4.2011	0.17553\\
4.2197	0.16489\\
4.2359	0.15426\\
4.2571	0.14362\\
4.278	0.13298\\
4.3018	0.12234\\
4.3248	0.1117\\
4.3516	0.10106\\
4.3803	0.090426\\
4.4066	0.079787\\
4.4414	0.069149\\
4.4829	0.058511\\
4.538	0.047872\\
4.6052	0.037234\\
4.7174	0.026596\\
4.8035	0.015957\\
5.159	0.0053191\\
};
\addlegendentry{\ExactEmpiricCCDFLegend}

\addplot [color=mycolor1, only marks, every error bar/.append style={opacity=0.30}, mark=*, mark size=0pt, draw=none, forget plot]
 plot [error bars/.cd, x dir=both, x explicit, error bar style={line width=1pt, color=mycolor1}, error mark options={mark=none,mark size=0pt}]
 table[row sep=crcr, x error plus index=2, x error minus index=3]{%
3.6037	0.99468	0.0031873	0.0031873\\
3.6073	0.98404	0.0032744	0.0032744\\
3.6108	0.9734	0.0033314	0.0033314\\
3.6145	0.96277	0.0034058	0.0034058\\
3.6181	0.95213	0.0034851	0.0034851\\
3.622	0.94149	0.0035	0.0035\\
3.6258	0.93085	0.0036615	0.0036615\\
3.6297	0.92021	0.0035244	0.0035244\\
3.6339	0.90957	0.0034445	0.0034445\\
3.6376	0.89894	0.0035539	0.0035539\\
3.6419	0.8883	0.0034716	0.0034716\\
3.6462	0.87766	0.0034679	0.0034679\\
3.6499	0.86702	0.0034622	0.0034622\\
3.6544	0.85638	0.0034513	0.0034513\\
3.6601	0.84574	0.0034059	0.0034059\\
3.6647	0.83511	0.003463	0.003463\\
3.6691	0.82447	0.0034186	0.0034186\\
3.6744	0.81383	0.003484	0.003484\\
3.6795	0.80319	0.0035559	0.0035559\\
3.6851	0.79255	0.0035894	0.0035894\\
3.6902	0.78191	0.0034953	0.0034953\\
3.6946	0.77128	0.0034979	0.0034979\\
3.701	0.76064	0.003521	0.003521\\
3.7067	0.75	0.0035332	0.0035332\\
3.7119	0.73936	0.0036677	0.0036677\\
3.7176	0.72872	0.0037632	0.0037632\\
3.7232	0.71809	0.0038613	0.0038613\\
3.7287	0.70745	0.0039704	0.0039704\\
3.7342	0.69681	0.0038912	0.0038912\\
3.7398	0.68617	0.0037488	0.0037488\\
3.7458	0.67553	0.0038371	0.0038371\\
3.7524	0.66489	0.0038114	0.0038114\\
3.7597	0.65426	0.0036991	0.0036991\\
3.7658	0.64362	0.0035915	0.0035915\\
3.773	0.63298	0.0034558	0.0034558\\
3.7802	0.62234	0.0035273	0.0035273\\
3.7873	0.6117	0.0037482	0.0037482\\
3.7936	0.60106	0.0040024	0.0040024\\
3.8011	0.59043	0.0041113	0.0041113\\
3.8081	0.57979	0.0041614	0.0041614\\
3.8151	0.56915	0.0042496	0.0042496\\
3.8222	0.55851	0.0043104	0.0043104\\
3.8279	0.54787	0.0043573	0.0043573\\
3.8349	0.53723	0.0041869	0.0041869\\
3.8422	0.5266	0.0040847	0.0040847\\
3.85	0.51596	0.003834	0.003834\\
3.8572	0.50532	0.0039218	0.0039218\\
3.8649	0.49468	0.0038832	0.0038832\\
3.8714	0.48404	0.003944	0.003944\\
3.8791	0.4734	0.0037933	0.0037933\\
3.8881	0.46277	0.0035871	0.0035871\\
3.8956	0.45213	0.0035147	0.0035147\\
3.9016	0.44149	0.0037251	0.0037251\\
3.91	0.43085	0.0038328	0.0038328\\
3.9175	0.42021	0.0041924	0.0041924\\
3.9272	0.40957	0.0043005	0.0043005\\
3.9348	0.39894	0.0042934	0.0042934\\
3.9449	0.3883	0.0044537	0.0044537\\
3.9548	0.37766	0.0047398	0.0047398\\
3.9657	0.36702	0.004744	0.004744\\
3.9754	0.35638	0.0046879	0.0046879\\
3.9863	0.34574	0.0045711	0.0045711\\
3.9973	0.33511	0.0047913	0.0047913\\
4.0075	0.32447	0.0047427	0.0047427\\
4.0175	0.31383	0.0050029	0.0050029\\
4.0299	0.30319	0.0051338	0.0051338\\
4.0406	0.29255	0.005174	0.005174\\
4.0543	0.28191	0.0053199	0.0053199\\
4.0681	0.27128	0.0056242	0.0056242\\
4.0824	0.26064	0.005481	0.005481\\
4.0965	0.25	0.0054067	0.0054067\\
4.1083	0.23936	0.0051031	0.0051031\\
4.1212	0.22872	0.0051613	0.0051613\\
4.138	0.21809	0.0048839	0.0048839\\
4.1548	0.20745	0.0052666	0.0052666\\
4.1697	0.19681	0.0050782	0.0050782\\
4.1848	0.18617	0.0052326	0.0052326\\
4.2011	0.17553	0.0048643	0.0048643\\
4.2197	0.16489	0.0049985	0.0049985\\
4.2359	0.15426	0.00532	0.00532\\
4.2571	0.14362	0.0053879	0.0053879\\
4.278	0.13298	0.0055106	0.0055106\\
4.3018	0.12234	0.0062277	0.0062277\\
4.3248	0.1117	0.0057529	0.0057529\\
4.3516	0.10106	0.00653	0.00653\\
4.3803	0.090426	0.0065879	0.0065879\\
4.4066	0.079787	0.0066297	0.0066297\\
4.4414	0.069149	0.0075694	0.0075694\\
4.4829	0.058511	0.0081425	0.0081425\\
4.538	0.047872	0.0091626	0.0091626\\
4.6052	0.037234	0.0096523	0.0096523\\
4.7174	0.026596	0.011213	0.011213\\
4.8035	0.015957	0.010763	0.010763\\
5.159	0.0053191	0.0090493	0.0090493\\
};
\addplot [color=mycolor1]
  table[row sep=crcr]{%
3.602	0.98157\\
3.6193	0.94969\\
3.6365	0.90448\\
4.1886	0.17657\\
4.7408	0.034618\\
5.2929	0.0068162\\
5.3101	0.0064792\\
};
\addlegendentry{\ExactParetoFitLegend}


\addplot[area legend, draw=none, fill=mycolor1, fill opacity=0.15, forget plot]
table[row sep=crcr] {%
x	y\\
3.602	0.98693\\
3.6193	0.95751\\
3.6365	0.91221\\
3.6538	0.86646\\
3.671	0.82301\\
3.6883	0.78175\\
3.7055	0.74257\\
3.7228	0.70535\\
3.74	0.67002\\
3.7573	0.63646\\
3.7745	0.60459\\
3.7918	0.57433\\
3.809	0.5456\\
3.8263	0.51831\\
3.8436	0.4924\\
3.8608	0.46781\\
3.8781	0.44445\\
3.8953	0.42227\\
3.9126	0.40122\\
3.9298	0.38123\\
3.9471	0.36225\\
3.9643	0.34423\\
3.9816	0.32712\\
3.9988	0.31087\\
4.0161	0.29544\\
4.0333	0.28079\\
4.0506	0.26687\\
4.0679	0.25365\\
4.0851	0.24109\\
4.1024	0.22916\\
4.1196	0.21782\\
4.1369	0.20705\\
4.1541	0.19681\\
4.1714	0.18709\\
4.1886	0.17784\\
4.2059	0.16906\\
4.2231	0.16071\\
4.2404	0.15277\\
4.2576	0.14523\\
4.2749	0.13806\\
4.2922	0.13125\\
4.3094	0.12477\\
4.3267	0.11862\\
4.3439	0.11277\\
4.3612	0.1072\\
4.3784	0.10192\\
4.3957	0.096893\\
4.4129	0.092116\\
4.4302	0.087575\\
4.4474	0.083258\\
4.4647	0.079155\\
4.4819	0.075254\\
4.4992	0.071546\\
4.5165	0.068021\\
4.5337	0.06467\\
4.551	0.061485\\
4.5682	0.058456\\
4.5855	0.055577\\
4.6027	0.05284\\
4.62	0.050239\\
4.6372	0.047765\\
4.6545	0.045414\\
4.6717	0.043178\\
4.689	0.041053\\
4.7062	0.039032\\
4.7235	0.037112\\
4.7408	0.035285\\
4.758	0.033549\\
4.7753	0.031899\\
4.7925	0.03033\\
4.8098	0.028838\\
4.827	0.027419\\
4.8443	0.026071\\
4.8615	0.024789\\
4.8788	0.02357\\
4.896	0.022411\\
4.9133	0.021309\\
4.9305	0.020262\\
4.9478	0.019266\\
4.9651	0.018319\\
4.9823	0.017418\\
4.9996	0.016562\\
5.0168	0.015749\\
5.0341	0.014975\\
5.0513	0.014239\\
5.0686	0.01354\\
5.0858	0.012875\\
5.1031	0.012242\\
5.1203	0.011641\\
5.1376	0.011069\\
5.1548	0.010526\\
5.1721	0.010009\\
5.1894	0.0095175\\
5.2066	0.0090503\\
5.2239	0.0086061\\
5.2411	0.0081837\\
5.2584	0.007782\\
5.2756	0.0074001\\
5.2929	0.007037\\
5.3101	0.0066917\\
5.3101	0.0062666\\
5.2929	0.0065954\\
5.2756	0.0069416\\
5.2584	0.0073059\\
5.2411	0.0076893\\
5.2239	0.0080929\\
5.2066	0.0085178\\
5.1894	0.0089649\\
5.1721	0.0094356\\
5.1548	0.009931\\
5.1376	0.010452\\
5.1203	0.011001\\
5.1031	0.011579\\
5.0858	0.012187\\
5.0686	0.012827\\
5.0513	0.013501\\
5.0341	0.01421\\
5.0168	0.014957\\
4.9996	0.015743\\
4.9823	0.01657\\
4.9651	0.017441\\
4.9478	0.018357\\
4.9305	0.019322\\
4.9133	0.020337\\
4.896	0.021406\\
4.8788	0.022532\\
4.8615	0.023716\\
4.8443	0.024963\\
4.827	0.026275\\
4.8098	0.027657\\
4.7925	0.029111\\
4.7753	0.030642\\
4.758	0.032254\\
4.7408	0.03395\\
4.7235	0.035736\\
4.7062	0.037616\\
4.689	0.039594\\
4.6717	0.041677\\
4.6545	0.04387\\
4.6372	0.046178\\
4.62	0.048608\\
4.6027	0.051166\\
4.5855	0.053859\\
4.5682	0.056693\\
4.551	0.059677\\
4.5337	0.062818\\
4.5165	0.066124\\
4.4992	0.069605\\
4.4819	0.073269\\
4.4647	0.077127\\
4.4474	0.081187\\
4.4302	0.085462\\
4.4129	0.089962\\
4.3957	0.094698\\
4.3784	0.099685\\
4.3612	0.10493\\
4.3439	0.11046\\
4.3267	0.11628\\
4.3094	0.1224\\
4.2922	0.12885\\
4.2749	0.13563\\
4.2576	0.14277\\
4.2404	0.15029\\
4.2231	0.1582\\
4.2059	0.16653\\
4.1886	0.1753\\
4.1714	0.18453\\
4.1541	0.19424\\
4.1369	0.20447\\
4.1196	0.21523\\
4.1024	0.22655\\
4.0851	0.23846\\
4.0679	0.251\\
4.0506	0.26419\\
4.0333	0.27807\\
4.0161	0.29267\\
3.9988	0.30803\\
3.9816	0.32419\\
3.9643	0.34118\\
3.9471	0.35906\\
3.9298	0.37785\\
3.9126	0.39762\\
3.8953	0.41841\\
3.8781	0.44027\\
3.8608	0.46326\\
3.8436	0.48745\\
3.8263	0.51288\\
3.809	0.53963\\
3.7918	0.56777\\
3.7745	0.59736\\
3.7573	0.62849\\
3.74	0.66124\\
3.7228	0.69569\\
3.7055	0.73193\\
3.6883	0.77006\\
3.671	0.81017\\
3.6538	0.85236\\
3.6365	0.89675\\
3.6193	0.94188\\
3.602	0.97621\\
}--cycle;
\addplot[only marks, mark=*, mark size=1.1180pt, color=black, fill=black, opacity=0.60, draw=none, mark options={draw=none,line width=0pt}] table[row sep=crcr]{%
x	y\\
3.5623	0.99468\\
3.5723	0.98404\\
3.573	0.9734\\
3.5753	0.96277\\
3.5766	0.95213\\
3.5781	0.94149\\
3.5826	0.93085\\
3.5851	0.92021\\
3.5911	0.90957\\
3.5942	0.89894\\
3.596	0.8883\\
3.6016	0.87766\\
3.603	0.86702\\
3.6047	0.85638\\
3.6086	0.84574\\
3.6086	0.83511\\
3.6097	0.82447\\
3.6142	0.81383\\
3.6261	0.80319\\
3.6383	0.79255\\
3.6464	0.78191\\
3.6535	0.77128\\
3.657	0.76064\\
3.6658	0.75\\
3.6678	0.73936\\
3.6687	0.72872\\
3.6702	0.71809\\
3.6723	0.70745\\
3.6793	0.69681\\
3.6867	0.68617\\
3.6949	0.67553\\
3.7015	0.66489\\
3.7031	0.65426\\
3.719	0.64362\\
3.7213	0.63298\\
3.7265	0.62234\\
3.7307	0.6117\\
3.737	0.60106\\
3.7374	0.59043\\
3.7423	0.57979\\
3.7487	0.56915\\
3.7677	0.55851\\
3.7685	0.54787\\
3.7768	0.53723\\
3.8005	0.5266\\
3.8015	0.51596\\
3.8022	0.50532\\
3.8032	0.49468\\
3.8039	0.48404\\
3.8108	0.4734\\
3.8214	0.46277\\
3.8405	0.45213\\
3.8605	0.44149\\
3.863	0.43085\\
3.8645	0.42021\\
3.8662	0.40957\\
3.8806	0.39894\\
3.8878	0.3883\\
3.8964	0.37766\\
3.905	0.36702\\
3.9294	0.35638\\
3.9303	0.34574\\
3.954	0.33511\\
3.9604	0.32447\\
3.9669	0.31383\\
3.9747	0.30319\\
3.9822	0.29255\\
3.9929	0.28191\\
4.0051	0.27128\\
4.0144	0.26064\\
4.0161	0.25\\
4.0433	0.23936\\
4.0578	0.22872\\
4.073	0.21809\\
4.0857	0.20745\\
4.0904	0.19681\\
4.1168	0.18617\\
4.1265	0.17553\\
4.127	0.16489\\
4.143	0.15426\\
4.1756	0.14362\\
4.2156	0.13298\\
4.3116	0.12234\\
4.3276	0.1117\\
4.3662	0.10106\\
4.4791	0.090426\\
4.4912	0.079787\\
4.517	0.069149\\
4.5743	0.058511\\
4.5914	0.047872\\
4.6138	0.037234\\
4.7898	0.026596\\
5.0876	0.015957\\
5.3101	0.0053191\\
};
\addlegendentry{\RealEmpiricCCDFLegend}

\addplot [color=black]
  table[row sep=crcr]{%
3.5608	1\\
5.3101	0.005939\\
};
\addlegendentry{\RealParetoFitLegend}

      % \input{../../../../.tex/20/KS/SizeDistributionFittingsfig17.tex}
      \legend{}
  
      \nextgroupplot[%
        xmin=3.5608,xmax=5.4115,
        % xmin=3.4,xmax=5.4,
        % ymin=0.0053191,ymax=1,
      ] % This file was created by matlab2tikz.
%
\definecolor{mycolor1}{rgb}{0.83529,0.36863,0.00000}%
\definecolor{mycolor2}{rgb}{0.00000,0.44700,0.74100}%
\definecolor{mycolor3}{rgb}{0.49400,0.18400,0.55600}%
%
\addplot[only marks, mark=*, mark size=1.1180pt, color=mycolor1, fill=mycolor1, opacity=0.60, draw=none, mark options={draw=none,line width=0pt}] table[row sep=crcr]{%
x	y\\
3.5981	0.99468\\
3.6023	0.98404\\
3.6064	0.9734\\
3.6105	0.96277\\
3.615	0.95213\\
3.619	0.94149\\
3.6233	0.93085\\
3.6271	0.92021\\
3.6318	0.90957\\
3.6366	0.89894\\
3.6403	0.8883\\
3.6443	0.87766\\
3.6487	0.86702\\
3.6532	0.85638\\
3.6584	0.84574\\
3.6638	0.83511\\
3.6687	0.82447\\
3.6733	0.81383\\
3.6786	0.80319\\
3.6828	0.79255\\
3.6881	0.78191\\
3.6936	0.77128\\
3.6991	0.76064\\
3.7048	0.75\\
3.71	0.73936\\
3.7158	0.72872\\
3.721	0.71809\\
3.7265	0.70745\\
3.7321	0.69681\\
3.737	0.68617\\
3.7428	0.67553\\
3.7488	0.66489\\
3.7556	0.65426\\
3.763	0.64362\\
3.7694	0.63298\\
3.776	0.62234\\
3.7822	0.6117\\
3.7889	0.60106\\
3.7946	0.59043\\
3.801	0.57979\\
3.8083	0.56915\\
3.8176	0.55851\\
3.824	0.54787\\
3.8317	0.53723\\
3.8394	0.5266\\
3.8475	0.51596\\
3.8558	0.50532\\
3.8625	0.49468\\
3.8709	0.48404\\
3.8781	0.4734\\
3.8874	0.46277\\
3.8955	0.45213\\
3.9029	0.44149\\
3.9129	0.43085\\
3.9213	0.42021\\
3.9292	0.40957\\
3.9386	0.39894\\
3.9488	0.3883\\
3.9566	0.37766\\
3.9662	0.36702\\
3.9771	0.35638\\
3.9865	0.34574\\
3.9964	0.33511\\
4.0082	0.32447\\
4.0215	0.31383\\
4.0338	0.30319\\
4.0469	0.29255\\
4.0576	0.28191\\
4.0718	0.27128\\
4.0833	0.26064\\
4.0968	0.25\\
4.1097	0.23936\\
4.123	0.22872\\
4.1357	0.21809\\
4.1509	0.20745\\
4.1655	0.19681\\
4.1838	0.18617\\
4.2013	0.17553\\
4.2175	0.16489\\
4.2364	0.15426\\
4.2558	0.14362\\
4.2753	0.13298\\
4.3016	0.12234\\
4.3251	0.1117\\
4.3478	0.10106\\
4.3732	0.090426\\
4.4078	0.079787\\
4.4449	0.069149\\
4.4792	0.058511\\
4.5401	0.047872\\
4.6076	0.037234\\
4.7133	0.026596\\
4.7931	0.015957\\
5.1695	0.0053191\\
};
\addlegendentry{\ApproxEmpiricCCDFLegend}

\addplot [color=mycolor1, only marks, every error bar/.append style={opacity=0.30}, mark=*, mark size=0pt, draw=none, forget plot]
 plot [error bars/.cd, x dir=both, x explicit, error bar style={line width=1pt, color=mycolor1}, error mark options={mark=none,mark size=0pt}]
 table[row sep=crcr, x error plus index=2, x error minus index=3]{%
3.5981	0.99468	0.0033318	0.0033318\\
3.6023	0.98404	0.0032921	0.0032921\\
3.6064	0.9734	0.0032161	0.0032161\\
3.6105	0.96277	0.0033671	0.0033671\\
3.615	0.95213	0.0032769	0.0032769\\
3.619	0.94149	0.0031849	0.0031849\\
3.6233	0.93085	0.0031801	0.0031801\\
3.6271	0.92021	0.0032614	0.0032614\\
3.6318	0.90957	0.0033254	0.0033254\\
3.6366	0.89894	0.0032616	0.0032616\\
3.6403	0.8883	0.0034099	0.0034099\\
3.6443	0.87766	0.0035325	0.0035325\\
3.6487	0.86702	0.0034635	0.0034635\\
3.6532	0.85638	0.0034169	0.0034169\\
3.6584	0.84574	0.0034051	0.0034051\\
3.6638	0.83511	0.0032979	0.0032979\\
3.6687	0.82447	0.00348	0.00348\\
3.6733	0.81383	0.0035628	0.0035628\\
3.6786	0.80319	0.0034968	0.0034968\\
3.6828	0.79255	0.0034696	0.0034696\\
3.6881	0.78191	0.003529	0.003529\\
3.6936	0.77128	0.0037115	0.0037115\\
3.6991	0.76064	0.0036262	0.0036262\\
3.7048	0.75	0.0034662	0.0034662\\
3.71	0.73936	0.0034329	0.0034329\\
3.7158	0.72872	0.0033673	0.0033673\\
3.721	0.71809	0.0033237	0.0033237\\
3.7265	0.70745	0.0035375	0.0035375\\
3.7321	0.69681	0.0036761	0.0036761\\
3.737	0.68617	0.0037267	0.0037267\\
3.7428	0.67553	0.0038231	0.0038231\\
3.7488	0.66489	0.0038249	0.0038249\\
3.7556	0.65426	0.0039771	0.0039771\\
3.763	0.64362	0.0041257	0.0041257\\
3.7694	0.63298	0.0039609	0.0039609\\
3.776	0.62234	0.0039102	0.0039102\\
3.7822	0.6117	0.0039858	0.0039858\\
3.7889	0.60106	0.0042246	0.0042246\\
3.7946	0.59043	0.0041824	0.0041824\\
3.801	0.57979	0.00417	0.00417\\
3.8083	0.56915	0.0042139	0.0042139\\
3.8176	0.55851	0.0043078	0.0043078\\
3.824	0.54787	0.0042625	0.0042625\\
3.8317	0.53723	0.0043864	0.0043864\\
3.8394	0.5266	0.0043818	0.0043818\\
3.8475	0.51596	0.0043929	0.0043929\\
3.8558	0.50532	0.0045964	0.0045964\\
3.8625	0.49468	0.0045453	0.0045453\\
3.8709	0.48404	0.0045508	0.0045508\\
3.8781	0.4734	0.0044692	0.0044692\\
3.8874	0.46277	0.0044678	0.0044678\\
3.8955	0.45213	0.0044672	0.0044672\\
3.9029	0.44149	0.0044505	0.0044505\\
3.9129	0.43085	0.0043948	0.0043948\\
3.9213	0.42021	0.0042943	0.0042943\\
3.9292	0.40957	0.004448	0.004448\\
3.9386	0.39894	0.0045776	0.0045776\\
3.9488	0.3883	0.0046377	0.0046377\\
3.9566	0.37766	0.0044836	0.0044836\\
3.9662	0.36702	0.00447	0.00447\\
3.9771	0.35638	0.0045325	0.0045325\\
3.9865	0.34574	0.0047347	0.0047347\\
3.9964	0.33511	0.0049881	0.0049881\\
4.0082	0.32447	0.0053043	0.0053043\\
4.0215	0.31383	0.0054884	0.0054884\\
4.0338	0.30319	0.0055015	0.0055015\\
4.0469	0.29255	0.0056892	0.0056892\\
4.0576	0.28191	0.0055564	0.0055564\\
4.0718	0.27128	0.0057063	0.0057063\\
4.0833	0.26064	0.0058436	0.0058436\\
4.0968	0.25	0.0061932	0.0061932\\
4.1097	0.23936	0.0059823	0.0059823\\
4.123	0.22872	0.0058606	0.0058606\\
4.1357	0.21809	0.0056708	0.0056708\\
4.1509	0.20745	0.0058678	0.0058678\\
4.1655	0.19681	0.0053964	0.0053964\\
4.1838	0.18617	0.0053979	0.0053979\\
4.2013	0.17553	0.0050116	0.0050116\\
4.2175	0.16489	0.0053215	0.0053215\\
4.2364	0.15426	0.005493	0.005493\\
4.2558	0.14362	0.0055404	0.0055404\\
4.2753	0.13298	0.0054582	0.0054582\\
4.3016	0.12234	0.0058416	0.0058416\\
4.3251	0.1117	0.0059397	0.0059397\\
4.3478	0.10106	0.0058893	0.0058893\\
4.3732	0.090426	0.0064931	0.0064931\\
4.4078	0.079787	0.0070281	0.0070281\\
4.4449	0.069149	0.0074367	0.0074367\\
4.4792	0.058511	0.0080956	0.0080956\\
4.5401	0.047872	0.010773	0.010773\\
4.6076	0.037234	0.010679	0.010679\\
4.7133	0.026596	0.010343	0.010343\\
4.7931	0.015957	0.011898	0.011898\\
5.1695	0.0053191	0.0088428	0.0088428\\
};
\addplot [color=mycolor1]
  table[row sep=crcr]{%
3.5976	0.98063\\
3.6149	0.94635\\
3.6322	0.90279\\
4.1338	0.20825\\
4.6355	0.04824\\
5.1544	0.01067\\
5.3101	0.0067913\\
};
\addlegendentry{\ApproxParetoFitLegend}


\addplot[area legend, draw=none, fill=mycolor1, fill opacity=0.15, forget plot]
table[row sep=crcr] {%
x	y\\
3.5976	0.98556\\
3.6149	0.95373\\
3.6322	0.91053\\
3.6495	0.86525\\
3.6668	0.82222\\
3.684	0.78135\\
3.7013	0.74252\\
3.7186	0.70563\\
3.7359	0.67059\\
3.7532	0.6373\\
3.7705	0.60567\\
3.7878	0.57564\\
3.8051	0.5471\\
3.8224	0.52\\
3.8397	0.49426\\
3.857	0.46981\\
3.8743	0.44659\\
3.8916	0.42454\\
3.9089	0.40359\\
3.9262	0.3837\\
3.9435	0.3648\\
3.9608	0.34685\\
3.9781	0.3298\\
3.9954	0.31359\\
4.0127	0.2982\\
4.03	0.28357\\
4.0473	0.26966\\
4.0646	0.25645\\
4.0819	0.24388\\
4.0992	0.23194\\
4.1165	0.22059\\
4.1338	0.20979\\
4.1511	0.19953\\
4.1684	0.18977\\
4.1857	0.18049\\
4.203	0.17167\\
4.2203	0.16328\\
4.2376	0.1553\\
4.2549	0.14771\\
4.2722	0.14049\\
4.2895	0.13363\\
4.3068	0.12711\\
4.3241	0.1209\\
4.3414	0.115\\
4.3587	0.10938\\
4.376	0.10405\\
4.3933	0.098968\\
4.4106	0.094139\\
4.4279	0.089546\\
4.4452	0.085177\\
4.4625	0.081022\\
4.4798	0.077071\\
4.4971	0.073312\\
4.5144	0.069737\\
4.5317	0.066337\\
4.549	0.063103\\
4.5663	0.060027\\
4.5836	0.057101\\
4.6009	0.054318\\
4.6182	0.051671\\
4.6355	0.049153\\
4.6528	0.046758\\
4.6701	0.04448\\
4.6874	0.042314\\
4.7047	0.040253\\
4.722	0.038292\\
4.7393	0.036427\\
4.7566	0.034653\\
4.7739	0.032966\\
4.7912	0.031361\\
4.8085	0.029835\\
4.8258	0.028382\\
4.8431	0.027001\\
4.8604	0.025687\\
4.8777	0.024437\\
4.895	0.023248\\
4.9123	0.022117\\
4.9296	0.021041\\
4.9469	0.020017\\
4.9642	0.019043\\
4.9815	0.018117\\
4.9988	0.017236\\
5.0161	0.016398\\
5.0334	0.015601\\
5.0507	0.014842\\
5.068	0.014121\\
5.0853	0.013434\\
5.1025	0.012781\\
5.1198	0.01216\\
5.1371	0.011569\\
5.1544	0.011007\\
5.1717	0.010472\\
5.189	0.0099631\\
5.2063	0.0094791\\
5.2236	0.0090186\\
5.2409	0.0085806\\
5.2582	0.0081638\\
5.2755	0.0077674\\
5.2928	0.0073902\\
5.3101	0.0070313\\
5.3101	0.0065512\\
5.2928	0.0068913\\
5.2755	0.0072491\\
5.2582	0.0076256\\
5.2409	0.0080216\\
5.2236	0.0084382\\
5.2063	0.0088765\\
5.189	0.0093377\\
5.1717	0.0098228\\
5.1544	0.010333\\
5.1371	0.01087\\
5.1198	0.011435\\
5.1025	0.012029\\
5.0853	0.012654\\
5.068	0.013312\\
5.0507	0.014004\\
5.0334	0.014732\\
5.0161	0.015498\\
4.9988	0.016304\\
4.9815	0.017152\\
4.9642	0.018044\\
4.9469	0.018983\\
4.9296	0.019971\\
4.9123	0.021009\\
4.895	0.022103\\
4.8777	0.023253\\
4.8604	0.024463\\
4.8431	0.025736\\
4.8258	0.027075\\
4.8085	0.028485\\
4.7912	0.029968\\
4.7739	0.031528\\
4.7566	0.033169\\
4.7393	0.034897\\
4.722	0.036714\\
4.7047	0.038626\\
4.6874	0.040638\\
4.6701	0.042755\\
4.6528	0.044982\\
4.6355	0.047326\\
4.6182	0.049792\\
4.6009	0.052386\\
4.5836	0.055116\\
4.5663	0.057989\\
4.549	0.061011\\
4.5317	0.064191\\
4.5144	0.067538\\
4.4971	0.071059\\
4.4798	0.074764\\
4.4625	0.078662\\
4.4452	0.082764\\
4.4279	0.08708\\
4.4106	0.091622\\
4.3933	0.0964\\
4.376	0.10143\\
4.3587	0.10672\\
4.3414	0.11229\\
4.3241	0.11814\\
4.3068	0.12431\\
4.2895	0.13079\\
4.2722	0.13762\\
4.2549	0.1448\\
4.2376	0.15235\\
4.2203	0.1603\\
4.203	0.16867\\
4.1857	0.17747\\
4.1684	0.18673\\
4.1511	0.19647\\
4.1338	0.20672\\
4.1165	0.2175\\
4.0992	0.22884\\
4.0819	0.24077\\
4.0646	0.25333\\
4.0473	0.26653\\
4.03	0.28041\\
4.0127	0.29502\\
3.9954	0.31038\\
3.9781	0.32652\\
3.9608	0.3435\\
3.9435	0.36135\\
3.9262	0.38011\\
3.9089	0.39984\\
3.8916	0.42057\\
3.8743	0.44235\\
3.857	0.46525\\
3.8397	0.48933\\
3.8224	0.51463\\
3.8051	0.54122\\
3.7878	0.56918\\
3.7705	0.59856\\
3.7532	0.62946\\
3.7359	0.66194\\
3.7186	0.69609\\
3.7013	0.73199\\
3.684	0.76974\\
3.6668	0.80943\\
3.6495	0.85116\\
3.6322	0.89505\\
3.6149	0.93897\\
3.5976	0.9757\\
}--cycle;
\addplot[only marks, mark=*, mark size=1.1180pt, color=black, fill=black, opacity=0.60, draw=none, mark options={draw=none,line width=0pt}] table[row sep=crcr]{%
x	y\\
3.5623	0.99468\\
3.5723	0.98404\\
3.573	0.9734\\
3.5753	0.96277\\
3.5766	0.95213\\
3.5781	0.94149\\
3.5826	0.93085\\
3.5851	0.92021\\
3.5911	0.90957\\
3.5942	0.89894\\
3.596	0.8883\\
3.6016	0.87766\\
3.603	0.86702\\
3.6047	0.85638\\
3.6086	0.84574\\
3.6086	0.83511\\
3.6097	0.82447\\
3.6142	0.81383\\
3.6261	0.80319\\
3.6383	0.79255\\
3.6464	0.78191\\
3.6535	0.77128\\
3.657	0.76064\\
3.6658	0.75\\
3.6678	0.73936\\
3.6687	0.72872\\
3.6702	0.71809\\
3.6723	0.70745\\
3.6793	0.69681\\
3.6867	0.68617\\
3.6949	0.67553\\
3.7015	0.66489\\
3.7031	0.65426\\
3.719	0.64362\\
3.7213	0.63298\\
3.7265	0.62234\\
3.7307	0.6117\\
3.737	0.60106\\
3.7374	0.59043\\
3.7423	0.57979\\
3.7487	0.56915\\
3.7677	0.55851\\
3.7685	0.54787\\
3.7768	0.53723\\
3.8005	0.5266\\
3.8015	0.51596\\
3.8022	0.50532\\
3.8032	0.49468\\
3.8039	0.48404\\
3.8108	0.4734\\
3.8214	0.46277\\
3.8405	0.45213\\
3.8605	0.44149\\
3.863	0.43085\\
3.8645	0.42021\\
3.8662	0.40957\\
3.8806	0.39894\\
3.8878	0.3883\\
3.8964	0.37766\\
3.905	0.36702\\
3.9294	0.35638\\
3.9303	0.34574\\
3.954	0.33511\\
3.9604	0.32447\\
3.9669	0.31383\\
3.9747	0.30319\\
3.9822	0.29255\\
3.9929	0.28191\\
4.0051	0.27128\\
4.0144	0.26064\\
4.0161	0.25\\
4.0433	0.23936\\
4.0578	0.22872\\
4.073	0.21809\\
4.0857	0.20745\\
4.0904	0.19681\\
4.1168	0.18617\\
4.1265	0.17553\\
4.127	0.16489\\
4.143	0.15426\\
4.1756	0.14362\\
4.2156	0.13298\\
4.3116	0.12234\\
4.3276	0.1117\\
4.3662	0.10106\\
4.4791	0.090426\\
4.4912	0.079787\\
4.517	0.069149\\
4.5743	0.058511\\
4.5914	0.047872\\
4.6138	0.037234\\
4.7898	0.026596\\
5.0876	0.015957\\
5.3101	0.0053191\\
};
\addlegendentry{\RealEmpiricCCDFLegend}

\addplot [color=black]
  table[row sep=crcr]{%
3.5608	1\\
5.3101	0.005939\\
};
\addlegendentry{\RealParetoFitLegend}

      % \input{../../../../.tex/20/KS/SizeDistributionFittingsfig18.tex}
      \legend{}
    \end{groupplot}
  
    \begin{groupplot}[
      group style={
        group name=Row3,
        plotGroup,
      },
      plotRow,
      row2KeysWithLabels,
      plotLegendNE,
      %
      % xmin=3637.5,xmax=2.5003e+05,
      ymin=0.0053191,ymax=1,
    ]
      \nextgroupplot[%
        at={($(Row2 c1r1.south west)-(0,\yPlotSep em)-(0,\plotHeight cm)$)},
        xmin=3.598,xmax=5.4115,
        % xmin=3.4,xmax=5.4,
        % ymin=0.0053191,ymax=1,
      ] % This file was created by matlab2tikz.
%
\definecolor{mycolor1}{rgb}{0.00000,0.44706,0.69804}%
\definecolor{mycolor2}{rgb}{0.00000,0.44700,0.74100}%
\definecolor{mycolor3}{rgb}{0.83529,0.36863,0.00000}%
\definecolor{mycolor4}{rgb}{0.49400,0.18400,0.55600}%
%
\addplot[only marks, mark=*, mark size=1.1180pt, color=mycolor1, fill=mycolor1, opacity=0.60, draw=none, mark options={draw=none,line width=0pt}] table[row sep=crcr]{%
x	y\\
3.6035	0.99468\\
3.6079	0.98404\\
3.6116	0.9734\\
3.6158	0.96277\\
3.6206	0.95213\\
3.6253	0.94149\\
3.6296	0.93085\\
3.6345	0.92021\\
3.6389	0.90957\\
3.6435	0.89894\\
3.6481	0.8883\\
3.6524	0.87766\\
3.6566	0.86702\\
3.6614	0.85638\\
3.6657	0.84574\\
3.6704	0.83511\\
3.6758	0.82447\\
3.6803	0.81383\\
3.6848	0.80319\\
3.6896	0.79255\\
3.6946	0.78191\\
3.6998	0.77128\\
3.7049	0.76064\\
3.7102	0.75\\
3.7157	0.73936\\
3.7209	0.72872\\
3.7264	0.71809\\
3.7344	0.70745\\
3.74	0.69681\\
3.7457	0.68617\\
3.7513	0.67553\\
3.7576	0.66489\\
3.7638	0.65426\\
3.7697	0.64362\\
3.7759	0.63298\\
3.7815	0.62234\\
3.7869	0.6117\\
3.7933	0.60106\\
3.8018	0.59043\\
3.8093	0.57979\\
3.8153	0.56915\\
3.8224	0.55851\\
3.8306	0.54787\\
3.8374	0.53723\\
3.844	0.5266\\
3.8513	0.51596\\
3.8589	0.50532\\
3.8663	0.49468\\
3.8746	0.48404\\
3.8818	0.4734\\
3.8898	0.46277\\
3.8988	0.45213\\
3.9079	0.44149\\
3.9163	0.43085\\
3.9249	0.42021\\
3.9333	0.40957\\
3.942	0.39894\\
3.952	0.3883\\
3.9608	0.37766\\
3.9702	0.36702\\
3.9807	0.35638\\
3.9917	0.34574\\
4.0037	0.33511\\
4.0158	0.32447\\
4.0275	0.31383\\
4.0385	0.30319\\
4.0494	0.29255\\
4.062	0.28191\\
4.0754	0.27128\\
4.0883	0.26064\\
4.1036	0.25\\
4.1184	0.23936\\
4.1331	0.22872\\
4.1471	0.21809\\
4.1608	0.20745\\
4.1772	0.19681\\
4.196	0.18617\\
4.2137	0.17553\\
4.23	0.16489\\
4.2473	0.15426\\
4.2659	0.14362\\
4.2887	0.13298\\
4.3115	0.12234\\
4.3384	0.1117\\
4.3656	0.10106\\
4.394	0.090426\\
4.4263	0.079787\\
4.4592	0.069149\\
4.5016	0.058511\\
4.5488	0.047872\\
4.6129	0.037234\\
4.7095	0.026596\\
4.8162	0.015957\\
5.1746	0.0053191\\
};
\addlegendentry{\ExactEmpiricCCDFLegend}

\addplot [color=mycolor1, only marks, every error bar/.append style={opacity=0.30}, mark=*, mark size=0pt, draw=none, forget plot]
 plot [error bars/.cd, x dir=both, x explicit, error bar style={line width=1pt, color=mycolor1}, error mark options={mark=none,mark size=0pt}]
 table[row sep=crcr, x error plus index=2, x error minus index=3]{%
3.6035	0.99468	0.0043687	0.0043687\\
3.6079	0.98404	0.0043134	0.0043134\\
3.6116	0.9734	0.0042921	0.0042921\\
3.6158	0.96277	0.0043486	0.0043486\\
3.6206	0.95213	0.0045234	0.0045234\\
3.6253	0.94149	0.0045887	0.0045887\\
3.6296	0.93085	0.0046128	0.0046128\\
3.6345	0.92021	0.0046646	0.0046646\\
3.6389	0.90957	0.0047592	0.0047592\\
3.6435	0.89894	0.004712	0.004712\\
3.6481	0.8883	0.0048417	0.0048417\\
3.6524	0.87766	0.004913	0.004913\\
3.6566	0.86702	0.0049333	0.0049333\\
3.6614	0.85638	0.0048893	0.0048893\\
3.6657	0.84574	0.0047743	0.0047743\\
3.6704	0.83511	0.0047267	0.0047267\\
3.6758	0.82447	0.0047531	0.0047531\\
3.6803	0.81383	0.0049722	0.0049722\\
3.6848	0.80319	0.0051901	0.0051901\\
3.6896	0.79255	0.0051486	0.0051486\\
3.6946	0.78191	0.0050613	0.0050613\\
3.6998	0.77128	0.0049859	0.0049859\\
3.7049	0.76064	0.0048672	0.0048672\\
3.7102	0.75	0.0049881	0.0049881\\
3.7157	0.73936	0.0048884	0.0048884\\
3.7209	0.72872	0.0049061	0.0049061\\
3.7264	0.71809	0.0048825	0.0048825\\
3.7344	0.70745	0.0048948	0.0048948\\
3.74	0.69681	0.0048439	0.0048439\\
3.7457	0.68617	0.0048029	0.0048029\\
3.7513	0.67553	0.0047493	0.0047493\\
3.7576	0.66489	0.0047029	0.0047029\\
3.7638	0.65426	0.0047906	0.0047906\\
3.7697	0.64362	0.0048818	0.0048818\\
3.7759	0.63298	0.0048859	0.0048859\\
3.7815	0.62234	0.0048079	0.0048079\\
3.7869	0.6117	0.0048096	0.0048096\\
3.7933	0.60106	0.0047568	0.0047568\\
3.8018	0.59043	0.0046697	0.0046697\\
3.8093	0.57979	0.004961	0.004961\\
3.8153	0.56915	0.0048973	0.0048973\\
3.8224	0.55851	0.0049265	0.0049265\\
3.8306	0.54787	0.0048253	0.0048253\\
3.8374	0.53723	0.0049268	0.0049268\\
3.844	0.5266	0.0047613	0.0047613\\
3.8513	0.51596	0.0048179	0.0048179\\
3.8589	0.50532	0.0048739	0.0048739\\
3.8663	0.49468	0.0046993	0.0046993\\
3.8746	0.48404	0.0049749	0.0049749\\
3.8818	0.4734	0.0052045	0.0052045\\
3.8898	0.46277	0.0049908	0.0049908\\
3.8988	0.45213	0.0050742	0.0050742\\
3.9079	0.44149	0.0051782	0.0051782\\
3.9163	0.43085	0.0052653	0.0052653\\
3.9249	0.42021	0.0052191	0.0052191\\
3.9333	0.40957	0.0050431	0.0050431\\
3.942	0.39894	0.0050961	0.0050961\\
3.952	0.3883	0.0053717	0.0053717\\
3.9608	0.37766	0.0056033	0.0056033\\
3.9702	0.36702	0.0056949	0.0056949\\
3.9807	0.35638	0.0056918	0.0056918\\
3.9917	0.34574	0.005667	0.005667\\
4.0037	0.33511	0.0055617	0.0055617\\
4.0158	0.32447	0.005337	0.005337\\
4.0275	0.31383	0.0053077	0.0053077\\
4.0385	0.30319	0.0051475	0.0051475\\
4.0494	0.29255	0.0051079	0.0051079\\
4.062	0.28191	0.0051371	0.0051371\\
4.0754	0.27128	0.0054134	0.0054134\\
4.0883	0.26064	0.0057357	0.0057357\\
4.1036	0.25	0.0058334	0.0058334\\
4.1184	0.23936	0.0059011	0.0059011\\
4.1331	0.22872	0.0059871	0.0059871\\
4.1471	0.21809	0.0064116	0.0064116\\
4.1608	0.20745	0.0066205	0.0066205\\
4.1772	0.19681	0.006858	0.006858\\
4.196	0.18617	0.0069585	0.0069585\\
4.2137	0.17553	0.0077849	0.0077849\\
4.23	0.16489	0.0082238	0.0082238\\
4.2473	0.15426	0.0081301	0.0081301\\
4.2659	0.14362	0.0082035	0.0082035\\
4.2887	0.13298	0.008017	0.008017\\
4.3115	0.12234	0.0083363	0.0083363\\
4.3384	0.1117	0.008877	0.008877\\
4.3656	0.10106	0.0079618	0.0079618\\
4.394	0.090426	0.0078787	0.0078787\\
4.4263	0.079787	0.0079578	0.0079578\\
4.4592	0.069149	0.0087748	0.0087748\\
4.5016	0.058511	0.010297	0.010297\\
4.5488	0.047872	0.010617	0.010617\\
4.6129	0.037234	0.011265	0.011265\\
4.7095	0.026596	0.013471	0.013471\\
4.8162	0.015957	0.014317	0.014317\\
5.1746	0.0053191	0.012259	0.012259\\
};
\addplot [color=mycolor1]
  table[row sep=crcr]{%
3.6023	0.97605\\
3.6198	0.94435\\
3.6372	0.90323\\
3.6547	0.86092\\
3.742	0.66762\\
4.126	0.21819\\
4.5276	0.06807\\
4.9291	0.021327\\
5.3306	0.0067093\\
};
\addlegendentry{\ExactParetoFitLegend}


\addplot[area legend, draw=none, fill=mycolor1, fill opacity=0.15, forget plot]
table[row sep=crcr] {%
x	y\\
3.6023	0.98202\\
3.6198	0.95296\\
3.6372	0.91279\\
3.6547	0.87077\\
3.6721	0.8277\\
3.6896	0.78616\\
3.707	0.74672\\
3.7245	0.70928\\
3.742	0.67373\\
3.7594	0.63997\\
3.7769	0.60792\\
3.7943	0.5775\\
3.8118	0.54861\\
3.8292	0.52119\\
3.8467	0.49516\\
3.8642	0.47045\\
3.8816	0.447\\
3.8991	0.42474\\
3.9165	0.40361\\
3.934	0.38356\\
3.9515	0.36454\\
3.9689	0.34648\\
3.9864	0.32934\\
4.0038	0.31306\\
4.0213	0.29762\\
4.0387	0.28295\\
4.0562	0.26901\\
4.0737	0.25578\\
4.0911	0.2432\\
4.1086	0.23126\\
4.126	0.2199\\
4.1435	0.20911\\
4.1609	0.19885\\
4.1784	0.18911\\
4.1959	0.17984\\
4.2133	0.17103\\
4.2308	0.16265\\
4.2482	0.15469\\
4.2657	0.14712\\
4.2832	0.13992\\
4.3006	0.13307\\
4.3181	0.12656\\
4.3355	0.12038\\
4.353	0.11449\\
4.3704	0.10889\\
4.3879	0.10357\\
4.4054	0.098512\\
4.4228	0.093699\\
4.4403	0.089122\\
4.4577	0.08477\\
4.4752	0.080631\\
4.4926	0.076695\\
4.5101	0.072951\\
4.5276	0.069391\\
4.545	0.066005\\
4.5625	0.062785\\
4.5799	0.059722\\
4.5974	0.056809\\
4.6149	0.054039\\
4.6323	0.051404\\
4.6498	0.048898\\
4.6672	0.046515\\
4.6847	0.044248\\
4.7021	0.042092\\
4.7196	0.040041\\
4.7371	0.03809\\
4.7545	0.036235\\
4.772	0.03447\\
4.7894	0.032792\\
4.8069	0.031195\\
4.8243	0.029677\\
4.8418	0.028232\\
4.8593	0.026858\\
4.8767	0.025551\\
4.8942	0.024308\\
4.9116	0.023126\\
4.9291	0.022001\\
4.9466	0.020931\\
4.964	0.019913\\
4.9815	0.018945\\
4.9989	0.018024\\
5.0164	0.017148\\
5.0338	0.016315\\
5.0513	0.015522\\
5.0688	0.014768\\
5.0862	0.01405\\
5.1037	0.013368\\
5.1211	0.012719\\
5.1386	0.012101\\
5.156	0.011513\\
5.1735	0.010954\\
5.191	0.010423\\
5.2084	0.0099168\\
5.2259	0.0094356\\
5.2433	0.0089777\\
5.2608	0.0085422\\
5.2783	0.0081278\\
5.2957	0.0077336\\
5.3132	0.0073586\\
5.3306	0.0070017\\
5.3306	0.0064169\\
5.3132	0.0067509\\
5.2957	0.0071023\\
5.2783	0.0074721\\
5.2608	0.0078612\\
5.2433	0.0082706\\
5.2259	0.0087015\\
5.2084	0.0091548\\
5.191	0.0096318\\
5.1735	0.010134\\
5.156	0.010662\\
5.1386	0.011218\\
5.1211	0.011803\\
5.1037	0.012418\\
5.0862	0.013066\\
5.0688	0.013747\\
5.0513	0.014464\\
5.0338	0.015219\\
5.0164	0.016013\\
4.9989	0.016849\\
4.9815	0.017729\\
4.964	0.018654\\
4.9466	0.019629\\
4.9291	0.020654\\
4.9116	0.021733\\
4.8942	0.022868\\
4.8767	0.024063\\
4.8593	0.02532\\
4.8418	0.026644\\
4.8243	0.028037\\
4.8069	0.029502\\
4.7894	0.031045\\
4.772	0.032669\\
4.7545	0.034378\\
4.7371	0.036176\\
4.7196	0.038069\\
4.7021	0.040062\\
4.6847	0.042159\\
4.6672	0.044366\\
4.6498	0.046689\\
4.6323	0.049134\\
4.6149	0.051707\\
4.5974	0.054415\\
4.5799	0.057266\\
4.5625	0.060267\\
4.545	0.063426\\
4.5276	0.06675\\
4.5101	0.070249\\
4.4926	0.073932\\
4.4752	0.077809\\
4.4577	0.08189\\
4.4403	0.086185\\
4.4228	0.090706\\
4.4054	0.095466\\
4.3879	0.10047\\
4.3704	0.10575\\
4.353	0.1113\\
4.3355	0.11714\\
4.3181	0.12329\\
4.3006	0.12976\\
4.2832	0.13658\\
4.2657	0.14375\\
4.2482	0.1513\\
4.2308	0.15924\\
4.2133	0.1676\\
4.1959	0.17641\\
4.1784	0.18567\\
4.1609	0.19542\\
4.1435	0.20568\\
4.126	0.21648\\
4.1086	0.22784\\
4.0911	0.23979\\
4.0737	0.25237\\
4.0562	0.2656\\
4.0387	0.27951\\
4.0213	0.29415\\
4.0038	0.30953\\
3.9864	0.32571\\
3.9689	0.34271\\
3.9515	0.36058\\
3.934	0.37937\\
3.9165	0.39911\\
3.8991	0.41985\\
3.8816	0.44166\\
3.8642	0.46457\\
3.8467	0.48866\\
3.8292	0.51399\\
3.8118	0.54061\\
3.7943	0.5686\\
3.7769	0.59803\\
3.7594	0.62898\\
3.742	0.66152\\
3.7245	0.69574\\
3.707	0.73173\\
3.6896	0.76958\\
3.6721	0.80938\\
3.6547	0.85108\\
3.6372	0.89366\\
3.6198	0.93575\\
3.6023	0.97007\\
}--cycle;
\addplot[only marks, mark=*, mark size=1.1180pt, color=mycolor3, fill=mycolor3, opacity=0.60, draw=none, mark options={draw=none,line width=0pt}] table[row sep=crcr]{%
x	y\\
3.608	0.99468\\
3.612	0.98404\\
3.6149	0.9734\\
3.6193	0.96277\\
3.6236	0.95213\\
3.6275	0.94149\\
3.6307	0.93085\\
3.6349	0.92021\\
3.6394	0.90957\\
3.6443	0.89894\\
3.6489	0.8883\\
3.6535	0.87766\\
3.6574	0.86702\\
3.6619	0.85638\\
3.6661	0.84574\\
3.6707	0.83511\\
3.6759	0.82447\\
3.6809	0.81383\\
3.6863	0.80319\\
3.6908	0.79255\\
3.6971	0.78191\\
3.7017	0.77128\\
3.7069	0.76064\\
3.7123	0.75\\
3.7172	0.73936\\
3.722	0.72872\\
3.7277	0.71809\\
3.734	0.70745\\
3.7392	0.69681\\
3.7453	0.68617\\
3.7511	0.67553\\
3.7578	0.66489\\
3.7629	0.65426\\
3.7687	0.64362\\
3.7756	0.63298\\
3.7819	0.62234\\
3.7891	0.6117\\
3.7954	0.60106\\
3.8017	0.59043\\
3.809	0.57979\\
3.8157	0.56915\\
3.8216	0.55851\\
3.8279	0.54787\\
3.835	0.53723\\
3.843	0.5266\\
3.8508	0.51596\\
3.8592	0.50532\\
3.8687	0.49468\\
3.8767	0.48404\\
3.8837	0.4734\\
3.8915	0.46277\\
3.8992	0.45213\\
3.9079	0.44149\\
3.9154	0.43085\\
3.926	0.42021\\
3.9348	0.40957\\
3.9442	0.39894\\
3.9526	0.3883\\
3.9638	0.37766\\
3.9717	0.36702\\
3.9827	0.35638\\
3.9938	0.34574\\
4.0041	0.33511\\
4.0134	0.32447\\
4.0267	0.31383\\
4.0386	0.30319\\
4.0508	0.29255\\
4.0624	0.28191\\
4.0725	0.27128\\
4.0888	0.26064\\
4.1049	0.25\\
4.117	0.23936\\
4.1302	0.22872\\
4.1455	0.21809\\
4.1599	0.20745\\
4.1742	0.19681\\
4.1902	0.18617\\
4.2074	0.17553\\
4.2233	0.16489\\
4.2398	0.15426\\
4.2602	0.14362\\
4.2789	0.13298\\
4.2994	0.12234\\
4.3276	0.1117\\
4.3546	0.10106\\
4.3817	0.090426\\
4.4171	0.079787\\
4.4522	0.069149\\
4.497	0.058511\\
4.5578	0.047872\\
4.6192	0.037234\\
4.7102	0.026596\\
4.8182	0.015957\\
5.1809	0.0053191\\
};
\addlegendentry{\ApproxEmpiricCCDFLegend}

\addplot [color=mycolor3, only marks, every error bar/.append style={opacity=0.30}, mark=*, mark size=0pt, draw=none, forget plot]
 plot [error bars/.cd, x dir=both, x explicit, error bar style={line width=1pt, color=mycolor3}, error mark options={mark=none,mark size=0pt}]
 table[row sep=crcr, x error plus index=2, x error minus index=3]{%
3.608	0.99468	0.0040359	0.0040359\\
3.612	0.98404	0.0040852	0.0040852\\
3.6149	0.9734	0.0041689	0.0041689\\
3.6193	0.96277	0.0042144	0.0042144\\
3.6236	0.95213	0.0042315	0.0042315\\
3.6275	0.94149	0.0042316	0.0042316\\
3.6307	0.93085	0.0042656	0.0042656\\
3.6349	0.92021	0.0042076	0.0042076\\
3.6394	0.90957	0.0042936	0.0042936\\
3.6443	0.89894	0.0042523	0.0042523\\
3.6489	0.8883	0.0042739	0.0042739\\
3.6535	0.87766	0.004405	0.004405\\
3.6574	0.86702	0.0044875	0.0044875\\
3.6619	0.85638	0.0044589	0.0044589\\
3.6661	0.84574	0.0043904	0.0043904\\
3.6707	0.83511	0.0043459	0.0043459\\
3.6759	0.82447	0.0044445	0.0044445\\
3.6809	0.81383	0.0044758	0.0044758\\
3.6863	0.80319	0.0044367	0.0044367\\
3.6908	0.79255	0.0043683	0.0043683\\
3.6971	0.78191	0.0043662	0.0043662\\
3.7017	0.77128	0.0041983	0.0041983\\
3.7069	0.76064	0.0042572	0.0042572\\
3.7123	0.75	0.0042455	0.0042455\\
3.7172	0.73936	0.0043061	0.0043061\\
3.722	0.72872	0.0041243	0.0041243\\
3.7277	0.71809	0.0042463	0.0042463\\
3.734	0.70745	0.0043807	0.0043807\\
3.7392	0.69681	0.0043561	0.0043561\\
3.7453	0.68617	0.0044568	0.0044568\\
3.7511	0.67553	0.0042342	0.0042342\\
3.7578	0.66489	0.0042251	0.0042251\\
3.7629	0.65426	0.0043299	0.0043299\\
3.7687	0.64362	0.0041462	0.0041462\\
3.7756	0.63298	0.0045377	0.0045377\\
3.7819	0.62234	0.0044165	0.0044165\\
3.7891	0.6117	0.0044983	0.0044983\\
3.7954	0.60106	0.0044046	0.0044046\\
3.8017	0.59043	0.0044956	0.0044956\\
3.809	0.57979	0.0046597	0.0046597\\
3.8157	0.56915	0.0048958	0.0048958\\
3.8216	0.55851	0.0048694	0.0048694\\
3.8279	0.54787	0.0047418	0.0047418\\
3.835	0.53723	0.0047159	0.0047159\\
3.843	0.5266	0.0047471	0.0047471\\
3.8508	0.51596	0.0047735	0.0047735\\
3.8592	0.50532	0.0047167	0.0047167\\
3.8687	0.49468	0.0047754	0.0047754\\
3.8767	0.48404	0.005157	0.005157\\
3.8837	0.4734	0.0052409	0.0052409\\
3.8915	0.46277	0.0054523	0.0054523\\
3.8992	0.45213	0.0055911	0.0055911\\
3.9079	0.44149	0.0055855	0.0055855\\
3.9154	0.43085	0.0056618	0.0056618\\
3.926	0.42021	0.0059361	0.0059361\\
3.9348	0.40957	0.0058921	0.0058921\\
3.9442	0.39894	0.0061974	0.0061974\\
3.9526	0.3883	0.0060931	0.0060931\\
3.9638	0.37766	0.0058931	0.0058931\\
3.9717	0.36702	0.0058103	0.0058103\\
3.9827	0.35638	0.0056547	0.0056547\\
3.9938	0.34574	0.0056375	0.0056375\\
4.0041	0.33511	0.0055906	0.0055906\\
4.0134	0.32447	0.0058442	0.0058442\\
4.0267	0.31383	0.0056071	0.0056071\\
4.0386	0.30319	0.0059954	0.0059954\\
4.0508	0.29255	0.0056373	0.0056373\\
4.0624	0.28191	0.0058582	0.0058582\\
4.0725	0.27128	0.0058151	0.0058151\\
4.0888	0.26064	0.005724	0.005724\\
4.1049	0.25	0.0057219	0.0057219\\
4.117	0.23936	0.0059065	0.0059065\\
4.1302	0.22872	0.0059432	0.0059432\\
4.1455	0.21809	0.0060705	0.0060705\\
4.1599	0.20745	0.0060167	0.0060167\\
4.1742	0.19681	0.005967	0.005967\\
4.1902	0.18617	0.0058452	0.0058452\\
4.2074	0.17553	0.0059932	0.0059932\\
4.2233	0.16489	0.0063192	0.0063192\\
4.2398	0.15426	0.0075653	0.0075653\\
4.2602	0.14362	0.0076698	0.0076698\\
4.2789	0.13298	0.0079775	0.0079775\\
4.2994	0.12234	0.0083583	0.0083583\\
4.3276	0.1117	0.008497	0.008497\\
4.3546	0.10106	0.0080207	0.0080207\\
4.3817	0.090426	0.0088382	0.0088382\\
4.4171	0.079787	0.0098774	0.0098774\\
4.4522	0.069149	0.010243	0.010243\\
4.497	0.058511	0.011325	0.011325\\
4.5578	0.047872	0.013129	0.013129\\
4.6192	0.037234	0.011595	0.011595\\
4.7102	0.026596	0.014518	0.014518\\
4.8182	0.015957	0.018109	0.018109\\
5.1809	0.0053191	0.014189	0.014189\\
};
\addplot [color=mycolor3]
  table[row sep=crcr]{%
3.606	0.97811\\
3.6234	0.94606\\
3.6408	0.90438\\
3.6931	0.77583\\
4.1112	0.22618\\
4.5293	0.066215\\
4.9474	0.019465\\
5.3306	0.0063592\\
};
\addlegendentry{\ApproxParetoFitLegend}


\addplot[area legend, draw=none, fill=mycolor3, fill opacity=0.15, forget plot]
table[row sep=crcr] {%
x	y\\
3.606	0.98437\\
3.6234	0.95474\\
3.6408	0.91396\\
3.6583	0.86904\\
3.6757	0.82509\\
3.6931	0.78338\\
3.7105	0.74379\\
3.7279	0.70622\\
3.7454	0.67056\\
3.7628	0.63672\\
3.7802	0.6046\\
3.7976	0.57412\\
3.815	0.5452\\
3.8325	0.51775\\
3.8499	0.4917\\
3.8673	0.46698\\
3.8847	0.44353\\
3.9021	0.42127\\
3.9196	0.40016\\
3.937	0.38012\\
3.9544	0.3611\\
3.9718	0.34306\\
3.9893	0.32593\\
4.0067	0.30968\\
4.0241	0.29424\\
4.0415	0.27959\\
4.0589	0.26568\\
4.0764	0.25247\\
4.0938	0.23993\\
4.1112	0.22801\\
4.1286	0.2167\\
4.146	0.20595\\
4.1635	0.19573\\
4.1809	0.18603\\
4.1983	0.17681\\
4.2157	0.16805\\
4.2331	0.15973\\
4.2506	0.15182\\
4.268	0.14431\\
4.2854	0.13717\\
4.3028	0.13038\\
4.3202	0.12394\\
4.3377	0.11781\\
4.3551	0.11198\\
4.3725	0.10645\\
4.3899	0.10119\\
4.4073	0.096188\\
4.4248	0.091436\\
4.4422	0.08692\\
4.4596	0.082628\\
4.477	0.078548\\
4.4944	0.074671\\
4.5119	0.070985\\
4.5293	0.067482\\
4.5467	0.064152\\
4.5641	0.060988\\
4.5815	0.057979\\
4.599	0.05512\\
4.6164	0.052402\\
4.6338	0.049818\\
4.6512	0.047362\\
4.6686	0.045028\\
4.6861	0.042809\\
4.7035	0.040699\\
4.7209	0.038694\\
4.7383	0.036788\\
4.7558	0.034976\\
4.7732	0.033254\\
4.7906	0.031616\\
4.808	0.03006\\
4.8254	0.02858\\
4.8429	0.027173\\
4.8603	0.025836\\
4.8777	0.024565\\
4.8951	0.023356\\
4.9125	0.022207\\
4.93	0.021115\\
4.9474	0.020077\\
4.9648	0.019089\\
4.9822	0.018151\\
4.9996	0.017259\\
5.0171	0.01641\\
5.0345	0.015604\\
5.0519	0.014837\\
5.0693	0.014108\\
5.0867	0.013415\\
5.1042	0.012756\\
5.1216	0.01213\\
5.139	0.011534\\
5.1564	0.010968\\
5.1738	0.010429\\
5.1913	0.0099173\\
5.2087	0.0094305\\
5.2261	0.0089678\\
5.2435	0.0085277\\
5.2609	0.0081094\\
5.2784	0.0077116\\
5.2958	0.0073334\\
5.3132	0.0069737\\
5.3306	0.0066318\\
5.3306	0.0060865\\
5.3132	0.006407\\
5.2958	0.0067445\\
5.2784	0.0070998\\
5.2609	0.0074738\\
5.2435	0.0078676\\
5.2261	0.0082822\\
5.2087	0.0087187\\
5.1913	0.0091782\\
5.1738	0.0096621\\
5.1564	0.010171\\
5.139	0.010708\\
5.1216	0.011273\\
5.1042	0.011867\\
5.0867	0.012493\\
5.0693	0.013152\\
5.0519	0.013846\\
5.0345	0.014577\\
5.0171	0.015346\\
4.9996	0.016156\\
4.9822	0.017009\\
4.9648	0.017907\\
4.9474	0.018853\\
4.93	0.019848\\
4.9125	0.020897\\
4.8951	0.022001\\
4.8777	0.023163\\
4.8603	0.024387\\
4.8429	0.025676\\
4.8254	0.027033\\
4.808	0.028462\\
4.7906	0.029967\\
4.7732	0.031551\\
4.7558	0.033219\\
4.7383	0.034976\\
4.7209	0.036826\\
4.7035	0.038774\\
4.6861	0.040826\\
4.6686	0.042986\\
4.6512	0.045261\\
4.6338	0.047656\\
4.6164	0.050179\\
4.599	0.052835\\
4.5815	0.055633\\
4.5641	0.058579\\
4.5467	0.061681\\
4.5293	0.064948\\
4.5119	0.068388\\
4.4944	0.072011\\
4.477	0.075826\\
4.4596	0.079843\\
4.4422	0.084074\\
4.4248	0.08853\\
4.4073	0.093222\\
4.3899	0.098163\\
4.3725	0.10337\\
4.3551	0.10885\\
4.3377	0.11462\\
4.3202	0.12069\\
4.3028	0.12709\\
4.2854	0.13383\\
4.268	0.14093\\
4.2506	0.1484\\
4.2331	0.15627\\
4.2157	0.16456\\
4.1983	0.17328\\
4.1809	0.18247\\
4.1635	0.19215\\
4.146	0.20233\\
4.1286	0.21306\\
4.1112	0.22435\\
4.0938	0.23623\\
4.0764	0.24874\\
4.0589	0.26191\\
4.0415	0.27577\\
4.0241	0.29035\\
4.0067	0.3057\\
3.9893	0.32185\\
3.9718	0.33884\\
3.9544	0.35671\\
3.937	0.37551\\
3.9196	0.39529\\
3.9021	0.41609\\
3.8847	0.43796\\
3.8673	0.46098\\
3.8499	0.48518\\
3.8325	0.51064\\
3.815	0.53742\\
3.7976	0.5656\\
3.7802	0.59524\\
3.7628	0.62642\\
3.7454	0.65923\\
3.7279	0.69375\\
3.7105	0.73007\\
3.6931	0.76828\\
3.6757	0.8085\\
3.6583	0.85081\\
3.6408	0.8948\\
3.6234	0.93738\\
3.606	0.97185\\
}--cycle;
      % \input{../../../../.tex/20/KS/SizeDistributionFittingsfig25.tex}
      \legend{}
  
      \nextgroupplot[%
        xmin=3.5608,xmax=5.3544,
        % xmin=3.4,xmax=5.4,
        % ymin=0.0053191,ymax=1,
      ] % This file was created by matlab2tikz.
%
\definecolor{mycolor1}{rgb}{0.00000,0.44706,0.69804}%
\definecolor{mycolor2}{rgb}{0.00000,0.44700,0.74100}%
\definecolor{mycolor3}{rgb}{0.49400,0.18400,0.55600}%
%
\addplot[only marks, mark=*, mark size=1.1180pt, color=mycolor1, fill=mycolor1, opacity=0.60, draw=none, mark options={draw=none,line width=0pt}] table[row sep=crcr]{%
x	y\\
3.6035	0.99468\\
3.6079	0.98404\\
3.6116	0.9734\\
3.6158	0.96277\\
3.6206	0.95213\\
3.6253	0.94149\\
3.6296	0.93085\\
3.6345	0.92021\\
3.6389	0.90957\\
3.6435	0.89894\\
3.6481	0.8883\\
3.6524	0.87766\\
3.6566	0.86702\\
3.6614	0.85638\\
3.6657	0.84574\\
3.6704	0.83511\\
3.6758	0.82447\\
3.6803	0.81383\\
3.6848	0.80319\\
3.6896	0.79255\\
3.6946	0.78191\\
3.6998	0.77128\\
3.7049	0.76064\\
3.7102	0.75\\
3.7157	0.73936\\
3.7209	0.72872\\
3.7264	0.71809\\
3.7344	0.70745\\
3.74	0.69681\\
3.7457	0.68617\\
3.7513	0.67553\\
3.7576	0.66489\\
3.7638	0.65426\\
3.7697	0.64362\\
3.7759	0.63298\\
3.7815	0.62234\\
3.7869	0.6117\\
3.7933	0.60106\\
3.8018	0.59043\\
3.8093	0.57979\\
3.8153	0.56915\\
3.8224	0.55851\\
3.8306	0.54787\\
3.8374	0.53723\\
3.844	0.5266\\
3.8513	0.51596\\
3.8589	0.50532\\
3.8663	0.49468\\
3.8746	0.48404\\
3.8818	0.4734\\
3.8898	0.46277\\
3.8988	0.45213\\
3.9079	0.44149\\
3.9163	0.43085\\
3.9249	0.42021\\
3.9333	0.40957\\
3.942	0.39894\\
3.952	0.3883\\
3.9608	0.37766\\
3.9702	0.36702\\
3.9807	0.35638\\
3.9917	0.34574\\
4.0037	0.33511\\
4.0158	0.32447\\
4.0275	0.31383\\
4.0385	0.30319\\
4.0494	0.29255\\
4.062	0.28191\\
4.0754	0.27128\\
4.0883	0.26064\\
4.1036	0.25\\
4.1184	0.23936\\
4.1331	0.22872\\
4.1471	0.21809\\
4.1608	0.20745\\
4.1772	0.19681\\
4.196	0.18617\\
4.2137	0.17553\\
4.23	0.16489\\
4.2473	0.15426\\
4.2659	0.14362\\
4.2887	0.13298\\
4.3115	0.12234\\
4.3384	0.1117\\
4.3656	0.10106\\
4.394	0.090426\\
4.4263	0.079787\\
4.4592	0.069149\\
4.5016	0.058511\\
4.5488	0.047872\\
4.6129	0.037234\\
4.7095	0.026596\\
4.8162	0.015957\\
5.1746	0.0053191\\
};
\addlegendentry{\ExactEmpiricCCDFLegend}

\addplot [color=mycolor1, only marks, every error bar/.append style={opacity=0.30}, mark=*, mark size=0pt, draw=none, forget plot]
 plot [error bars/.cd, x dir=both, x explicit, error bar style={line width=1pt, color=mycolor1}, error mark options={mark=none,mark size=0pt}]
 table[row sep=crcr, x error plus index=2, x error minus index=3]{%
3.6035	0.99468	0.0043687	0.0043687\\
3.6079	0.98404	0.0043134	0.0043134\\
3.6116	0.9734	0.0042921	0.0042921\\
3.6158	0.96277	0.0043486	0.0043486\\
3.6206	0.95213	0.0045234	0.0045234\\
3.6253	0.94149	0.0045887	0.0045887\\
3.6296	0.93085	0.0046128	0.0046128\\
3.6345	0.92021	0.0046646	0.0046646\\
3.6389	0.90957	0.0047592	0.0047592\\
3.6435	0.89894	0.004712	0.004712\\
3.6481	0.8883	0.0048417	0.0048417\\
3.6524	0.87766	0.004913	0.004913\\
3.6566	0.86702	0.0049333	0.0049333\\
3.6614	0.85638	0.0048893	0.0048893\\
3.6657	0.84574	0.0047743	0.0047743\\
3.6704	0.83511	0.0047267	0.0047267\\
3.6758	0.82447	0.0047531	0.0047531\\
3.6803	0.81383	0.0049722	0.0049722\\
3.6848	0.80319	0.0051901	0.0051901\\
3.6896	0.79255	0.0051486	0.0051486\\
3.6946	0.78191	0.0050613	0.0050613\\
3.6998	0.77128	0.0049859	0.0049859\\
3.7049	0.76064	0.0048672	0.0048672\\
3.7102	0.75	0.0049881	0.0049881\\
3.7157	0.73936	0.0048884	0.0048884\\
3.7209	0.72872	0.0049061	0.0049061\\
3.7264	0.71809	0.0048825	0.0048825\\
3.7344	0.70745	0.0048948	0.0048948\\
3.74	0.69681	0.0048439	0.0048439\\
3.7457	0.68617	0.0048029	0.0048029\\
3.7513	0.67553	0.0047493	0.0047493\\
3.7576	0.66489	0.0047029	0.0047029\\
3.7638	0.65426	0.0047906	0.0047906\\
3.7697	0.64362	0.0048818	0.0048818\\
3.7759	0.63298	0.0048859	0.0048859\\
3.7815	0.62234	0.0048079	0.0048079\\
3.7869	0.6117	0.0048096	0.0048096\\
3.7933	0.60106	0.0047568	0.0047568\\
3.8018	0.59043	0.0046697	0.0046697\\
3.8093	0.57979	0.004961	0.004961\\
3.8153	0.56915	0.0048973	0.0048973\\
3.8224	0.55851	0.0049265	0.0049265\\
3.8306	0.54787	0.0048253	0.0048253\\
3.8374	0.53723	0.0049268	0.0049268\\
3.844	0.5266	0.0047613	0.0047613\\
3.8513	0.51596	0.0048179	0.0048179\\
3.8589	0.50532	0.0048739	0.0048739\\
3.8663	0.49468	0.0046993	0.0046993\\
3.8746	0.48404	0.0049749	0.0049749\\
3.8818	0.4734	0.0052045	0.0052045\\
3.8898	0.46277	0.0049908	0.0049908\\
3.8988	0.45213	0.0050742	0.0050742\\
3.9079	0.44149	0.0051782	0.0051782\\
3.9163	0.43085	0.0052653	0.0052653\\
3.9249	0.42021	0.0052191	0.0052191\\
3.9333	0.40957	0.0050431	0.0050431\\
3.942	0.39894	0.0050961	0.0050961\\
3.952	0.3883	0.0053717	0.0053717\\
3.9608	0.37766	0.0056033	0.0056033\\
3.9702	0.36702	0.0056949	0.0056949\\
3.9807	0.35638	0.0056918	0.0056918\\
3.9917	0.34574	0.005667	0.005667\\
4.0037	0.33511	0.0055617	0.0055617\\
4.0158	0.32447	0.005337	0.005337\\
4.0275	0.31383	0.0053077	0.0053077\\
4.0385	0.30319	0.0051475	0.0051475\\
4.0494	0.29255	0.0051079	0.0051079\\
4.062	0.28191	0.0051371	0.0051371\\
4.0754	0.27128	0.0054134	0.0054134\\
4.0883	0.26064	0.0057357	0.0057357\\
4.1036	0.25	0.0058334	0.0058334\\
4.1184	0.23936	0.0059011	0.0059011\\
4.1331	0.22872	0.0059871	0.0059871\\
4.1471	0.21809	0.0064116	0.0064116\\
4.1608	0.20745	0.0066205	0.0066205\\
4.1772	0.19681	0.006858	0.006858\\
4.196	0.18617	0.0069585	0.0069585\\
4.2137	0.17553	0.0077849	0.0077849\\
4.23	0.16489	0.0082238	0.0082238\\
4.2473	0.15426	0.0081301	0.0081301\\
4.2659	0.14362	0.0082035	0.0082035\\
4.2887	0.13298	0.008017	0.008017\\
4.3115	0.12234	0.0083363	0.0083363\\
4.3384	0.1117	0.008877	0.008877\\
4.3656	0.10106	0.0079618	0.0079618\\
4.394	0.090426	0.0078787	0.0078787\\
4.4263	0.079787	0.0079578	0.0079578\\
4.4592	0.069149	0.0087748	0.0087748\\
4.5016	0.058511	0.010297	0.010297\\
4.5488	0.047872	0.010617	0.010617\\
4.6129	0.037234	0.011265	0.011265\\
4.7095	0.026596	0.013471	0.013471\\
4.8162	0.015957	0.014317	0.014317\\
5.1746	0.0053191	0.012259	0.012259\\
};
\addplot [color=mycolor1]
  table[row sep=crcr]{%
3.6023	0.97605\\
3.6195	0.9448\\
3.6368	0.90424\\
3.654	0.86244\\
3.7403	0.67086\\
4.1198	0.22216\\
4.5166	0.070267\\
4.9134	0.022318\\
5.3101	0.0071165\\
};
\addlegendentry{\ExactParetoFitLegend}


\addplot[area legend, draw=none, fill=mycolor1, fill opacity=0.15, forget plot]
table[row sep=crcr] {%
x	y\\
3.6023	0.98202\\
3.6195	0.95338\\
3.6368	0.91379\\
3.654	0.8723\\
3.6713	0.82972\\
3.6886	0.78856\\
3.7058	0.74946\\
3.7231	0.71231\\
3.7403	0.67702\\
3.7576	0.64349\\
3.7748	0.61164\\
3.7921	0.58138\\
3.8093	0.55263\\
3.8266	0.52532\\
3.8438	0.49939\\
3.8611	0.47475\\
3.8783	0.45135\\
3.8956	0.42913\\
3.9128	0.40803\\
3.9301	0.38799\\
3.9473	0.36896\\
3.9646	0.35088\\
3.9818	0.33372\\
3.9991	0.31742\\
4.0163	0.30193\\
4.0336	0.28721\\
4.0508	0.27323\\
4.0681	0.25994\\
4.0853	0.24731\\
4.1026	0.23529\\
4.1198	0.22387\\
4.1371	0.21301\\
4.1543	0.20268\\
4.1716	0.19286\\
4.1888	0.18352\\
4.2061	0.17463\\
4.2233	0.16617\\
4.2406	0.15813\\
4.2578	0.15048\\
4.2751	0.1432\\
4.2923	0.13628\\
4.3096	0.12969\\
4.3268	0.12342\\
4.3441	0.11745\\
4.3613	0.11178\\
4.3786	0.10638\\
4.3958	0.10124\\
4.4131	0.096351\\
4.4303	0.091699\\
4.4476	0.087272\\
4.4648	0.083059\\
4.4821	0.079051\\
4.4993	0.075237\\
4.5166	0.071607\\
4.5338	0.068153\\
4.5511	0.064866\\
4.5683	0.061738\\
4.5856	0.058762\\
4.6028	0.055929\\
4.6201	0.053233\\
4.6373	0.050668\\
4.6546	0.048227\\
4.6719	0.045903\\
4.6891	0.043692\\
4.7064	0.041588\\
4.7236	0.039585\\
4.7409	0.037679\\
4.7581	0.035865\\
4.7754	0.034138\\
4.7926	0.032495\\
4.8099	0.030932\\
4.8271	0.029443\\
4.8444	0.028027\\
4.8616	0.026679\\
4.8789	0.025396\\
4.8961	0.024174\\
4.9134	0.023012\\
4.9306	0.021906\\
4.9479	0.020853\\
4.9651	0.01985\\
4.9824	0.018896\\
4.9996	0.017988\\
5.0169	0.017124\\
5.0341	0.016302\\
5.0514	0.015519\\
5.0686	0.014773\\
5.0859	0.014064\\
5.1031	0.013389\\
5.1204	0.012746\\
5.1376	0.012134\\
5.1549	0.011552\\
5.1721	0.010997\\
5.1894	0.01047\\
5.2066	0.0099675\\
5.2239	0.0094894\\
5.2411	0.0090343\\
5.2584	0.0086011\\
5.2756	0.0081887\\
5.2929	0.0077961\\
5.3101	0.0074224\\
5.3101	0.0068107\\
5.2929	0.0071609\\
5.2756	0.0075292\\
5.2584	0.0079165\\
5.2411	0.0083238\\
5.2239	0.0087522\\
5.2066	0.0092026\\
5.1894	0.0096763\\
5.1721	0.010174\\
5.1549	0.010698\\
5.1376	0.011249\\
5.1204	0.011829\\
5.1031	0.012438\\
5.0859	0.013079\\
5.0686	0.013753\\
5.0514	0.014461\\
5.0341	0.015207\\
5.0169	0.015991\\
4.9996	0.016815\\
4.9824	0.017682\\
4.9651	0.018594\\
4.9479	0.019553\\
4.9306	0.020562\\
4.9134	0.021623\\
4.8961	0.022739\\
4.8789	0.023913\\
4.8616	0.025147\\
4.8444	0.026446\\
4.8271	0.027811\\
4.8099	0.029248\\
4.7926	0.030759\\
4.7754	0.032348\\
4.7581	0.034019\\
4.7409	0.035777\\
4.7236	0.037627\\
4.7064	0.039572\\
4.6891	0.041618\\
4.6719	0.04377\\
4.6546	0.046034\\
4.6373	0.048415\\
4.6201	0.05092\\
4.6028	0.053554\\
4.5856	0.056326\\
4.5683	0.059241\\
4.5511	0.062308\\
4.5338	0.065534\\
4.5166	0.068928\\
4.4993	0.072497\\
4.4821	0.076253\\
4.4648	0.080203\\
4.4476	0.084358\\
4.4303	0.08873\\
4.4131	0.093328\\
4.3958	0.098166\\
4.3786	0.10325\\
4.3613	0.10861\\
4.3441	0.11424\\
4.3268	0.12016\\
4.3096	0.12639\\
4.2923	0.13295\\
4.2751	0.13985\\
4.2578	0.1471\\
4.2406	0.15473\\
4.2233	0.16276\\
4.2061	0.1712\\
4.1888	0.18008\\
4.1716	0.18943\\
4.1543	0.19925\\
4.1371	0.20958\\
4.1198	0.22045\\
4.1026	0.23188\\
4.0853	0.2439\\
4.0681	0.25653\\
4.0508	0.26981\\
4.0336	0.28377\\
4.0163	0.29845\\
3.9991	0.31386\\
3.9818	0.33006\\
3.9646	0.34708\\
3.9473	0.36495\\
3.9301	0.38373\\
3.9128	0.40345\\
3.8956	0.42416\\
3.8783	0.44592\\
3.8611	0.46877\\
3.8438	0.49278\\
3.8266	0.51801\\
3.8093	0.54451\\
3.7921	0.57236\\
3.7748	0.60162\\
3.7576	0.63238\\
3.7403	0.6647\\
3.7231	0.69866\\
3.7058	0.73436\\
3.6886	0.77188\\
3.6713	0.81132\\
3.654	0.85259\\
3.6368	0.89468\\
3.6195	0.93622\\
3.6023	0.97007\\
}--cycle;
\addplot[only marks, mark=*, mark size=1.1180pt, color=black, fill=black, opacity=0.60, draw=none, mark options={draw=none,line width=0pt}] table[row sep=crcr]{%
x	y\\
3.5623	0.99468\\
3.5723	0.98404\\
3.573	0.9734\\
3.5753	0.96277\\
3.5766	0.95213\\
3.5781	0.94149\\
3.5826	0.93085\\
3.5851	0.92021\\
3.5911	0.90957\\
3.5942	0.89894\\
3.596	0.8883\\
3.6016	0.87766\\
3.603	0.86702\\
3.6047	0.85638\\
3.6086	0.84574\\
3.6086	0.83511\\
3.6097	0.82447\\
3.6142	0.81383\\
3.6261	0.80319\\
3.6383	0.79255\\
3.6464	0.78191\\
3.6535	0.77128\\
3.657	0.76064\\
3.6658	0.75\\
3.6678	0.73936\\
3.6687	0.72872\\
3.6702	0.71809\\
3.6723	0.70745\\
3.6793	0.69681\\
3.6867	0.68617\\
3.6949	0.67553\\
3.7015	0.66489\\
3.7031	0.65426\\
3.719	0.64362\\
3.7213	0.63298\\
3.7265	0.62234\\
3.7307	0.6117\\
3.737	0.60106\\
3.7374	0.59043\\
3.7423	0.57979\\
3.7487	0.56915\\
3.7677	0.55851\\
3.7685	0.54787\\
3.7768	0.53723\\
3.8005	0.5266\\
3.8015	0.51596\\
3.8022	0.50532\\
3.8032	0.49468\\
3.8039	0.48404\\
3.8108	0.4734\\
3.8214	0.46277\\
3.8405	0.45213\\
3.8605	0.44149\\
3.863	0.43085\\
3.8645	0.42021\\
3.8662	0.40957\\
3.8806	0.39894\\
3.8878	0.3883\\
3.8964	0.37766\\
3.905	0.36702\\
3.9294	0.35638\\
3.9303	0.34574\\
3.954	0.33511\\
3.9604	0.32447\\
3.9669	0.31383\\
3.9747	0.30319\\
3.9822	0.29255\\
3.9929	0.28191\\
4.0051	0.27128\\
4.0144	0.26064\\
4.0161	0.25\\
4.0433	0.23936\\
4.0578	0.22872\\
4.073	0.21809\\
4.0857	0.20745\\
4.0904	0.19681\\
4.1168	0.18617\\
4.1265	0.17553\\
4.127	0.16489\\
4.143	0.15426\\
4.1756	0.14362\\
4.2156	0.13298\\
4.3116	0.12234\\
4.3276	0.1117\\
4.3662	0.10106\\
4.4791	0.090426\\
4.4912	0.079787\\
4.517	0.069149\\
4.5743	0.058511\\
4.5914	0.047872\\
4.6138	0.037234\\
4.7898	0.026596\\
5.0876	0.015957\\
5.3101	0.0053191\\
};
\addlegendentry{\RealEmpiricCCDFLegend}

\addplot [color=black]
  table[row sep=crcr]{%
3.5608	1\\
5.3101	0.005939\\
};
\addlegendentry{\RealParetoFitLegend}

      % \input{../../../../.tex/20/KS/SizeDistributionFittingsfig26.tex}
      \legend{}
  
      \nextgroupplot[%
        xmin=3.5608,xmax=5.4115,
        % xmin=3.4,xmax=5.4,
        % ymin=0.0053191,ymax=1,
      ] % This file was created by matlab2tikz.
%
\definecolor{mycolor1}{rgb}{0.83529,0.36863,0.00000}%
\definecolor{mycolor2}{rgb}{0.00000,0.44700,0.74100}%
\definecolor{mycolor3}{rgb}{0.49400,0.18400,0.55600}%
%
\addplot[only marks, mark=*, mark size=1.1180pt, color=mycolor1, fill=mycolor1, opacity=0.60, draw=none, mark options={draw=none,line width=0pt}] table[row sep=crcr]{%
x	y\\
3.608	0.99468\\
3.612	0.98404\\
3.6149	0.9734\\
3.6193	0.96277\\
3.6236	0.95213\\
3.6275	0.94149\\
3.6307	0.93085\\
3.6349	0.92021\\
3.6394	0.90957\\
3.6443	0.89894\\
3.6489	0.8883\\
3.6535	0.87766\\
3.6574	0.86702\\
3.6619	0.85638\\
3.6661	0.84574\\
3.6707	0.83511\\
3.6759	0.82447\\
3.6809	0.81383\\
3.6863	0.80319\\
3.6908	0.79255\\
3.6971	0.78191\\
3.7017	0.77128\\
3.7069	0.76064\\
3.7123	0.75\\
3.7172	0.73936\\
3.722	0.72872\\
3.7277	0.71809\\
3.734	0.70745\\
3.7392	0.69681\\
3.7453	0.68617\\
3.7511	0.67553\\
3.7578	0.66489\\
3.7629	0.65426\\
3.7687	0.64362\\
3.7756	0.63298\\
3.7819	0.62234\\
3.7891	0.6117\\
3.7954	0.60106\\
3.8017	0.59043\\
3.809	0.57979\\
3.8157	0.56915\\
3.8216	0.55851\\
3.8279	0.54787\\
3.835	0.53723\\
3.843	0.5266\\
3.8508	0.51596\\
3.8592	0.50532\\
3.8687	0.49468\\
3.8767	0.48404\\
3.8837	0.4734\\
3.8915	0.46277\\
3.8992	0.45213\\
3.9079	0.44149\\
3.9154	0.43085\\
3.926	0.42021\\
3.9348	0.40957\\
3.9442	0.39894\\
3.9526	0.3883\\
3.9638	0.37766\\
3.9717	0.36702\\
3.9827	0.35638\\
3.9938	0.34574\\
4.0041	0.33511\\
4.0134	0.32447\\
4.0267	0.31383\\
4.0386	0.30319\\
4.0508	0.29255\\
4.0624	0.28191\\
4.0725	0.27128\\
4.0888	0.26064\\
4.1049	0.25\\
4.117	0.23936\\
4.1302	0.22872\\
4.1455	0.21809\\
4.1599	0.20745\\
4.1742	0.19681\\
4.1902	0.18617\\
4.2074	0.17553\\
4.2233	0.16489\\
4.2398	0.15426\\
4.2602	0.14362\\
4.2789	0.13298\\
4.2994	0.12234\\
4.3276	0.1117\\
4.3546	0.10106\\
4.3817	0.090426\\
4.4171	0.079787\\
4.4522	0.069149\\
4.497	0.058511\\
4.5578	0.047872\\
4.6192	0.037234\\
4.7102	0.026596\\
4.8182	0.015957\\
5.1809	0.0053191\\
};
\addlegendentry{\ApproxEmpiricCCDFLegend}

\addplot [color=mycolor1, only marks, every error bar/.append style={opacity=0.30}, mark=*, mark size=0pt, draw=none, forget plot]
 plot [error bars/.cd, x dir=both, x explicit, error bar style={line width=1pt, color=mycolor1}, error mark options={mark=none,mark size=0pt}]
 table[row sep=crcr, x error plus index=2, x error minus index=3]{%
3.608	0.99468	0.0040359	0.0040359\\
3.612	0.98404	0.0040852	0.0040852\\
3.6149	0.9734	0.0041689	0.0041689\\
3.6193	0.96277	0.0042144	0.0042144\\
3.6236	0.95213	0.0042315	0.0042315\\
3.6275	0.94149	0.0042316	0.0042316\\
3.6307	0.93085	0.0042656	0.0042656\\
3.6349	0.92021	0.0042076	0.0042076\\
3.6394	0.90957	0.0042936	0.0042936\\
3.6443	0.89894	0.0042523	0.0042523\\
3.6489	0.8883	0.0042739	0.0042739\\
3.6535	0.87766	0.004405	0.004405\\
3.6574	0.86702	0.0044875	0.0044875\\
3.6619	0.85638	0.0044589	0.0044589\\
3.6661	0.84574	0.0043904	0.0043904\\
3.6707	0.83511	0.0043459	0.0043459\\
3.6759	0.82447	0.0044445	0.0044445\\
3.6809	0.81383	0.0044758	0.0044758\\
3.6863	0.80319	0.0044367	0.0044367\\
3.6908	0.79255	0.0043683	0.0043683\\
3.6971	0.78191	0.0043662	0.0043662\\
3.7017	0.77128	0.0041983	0.0041983\\
3.7069	0.76064	0.0042572	0.0042572\\
3.7123	0.75	0.0042455	0.0042455\\
3.7172	0.73936	0.0043061	0.0043061\\
3.722	0.72872	0.0041243	0.0041243\\
3.7277	0.71809	0.0042463	0.0042463\\
3.734	0.70745	0.0043807	0.0043807\\
3.7392	0.69681	0.0043561	0.0043561\\
3.7453	0.68617	0.0044568	0.0044568\\
3.7511	0.67553	0.0042342	0.0042342\\
3.7578	0.66489	0.0042251	0.0042251\\
3.7629	0.65426	0.0043299	0.0043299\\
3.7687	0.64362	0.0041462	0.0041462\\
3.7756	0.63298	0.0045377	0.0045377\\
3.7819	0.62234	0.0044165	0.0044165\\
3.7891	0.6117	0.0044983	0.0044983\\
3.7954	0.60106	0.0044046	0.0044046\\
3.8017	0.59043	0.0044956	0.0044956\\
3.809	0.57979	0.0046597	0.0046597\\
3.8157	0.56915	0.0048958	0.0048958\\
3.8216	0.55851	0.0048694	0.0048694\\
3.8279	0.54787	0.0047418	0.0047418\\
3.835	0.53723	0.0047159	0.0047159\\
3.843	0.5266	0.0047471	0.0047471\\
3.8508	0.51596	0.0047735	0.0047735\\
3.8592	0.50532	0.0047167	0.0047167\\
3.8687	0.49468	0.0047754	0.0047754\\
3.8767	0.48404	0.005157	0.005157\\
3.8837	0.4734	0.0052409	0.0052409\\
3.8915	0.46277	0.0054523	0.0054523\\
3.8992	0.45213	0.0055911	0.0055911\\
3.9079	0.44149	0.0055855	0.0055855\\
3.9154	0.43085	0.0056618	0.0056618\\
3.926	0.42021	0.0059361	0.0059361\\
3.9348	0.40957	0.0058921	0.0058921\\
3.9442	0.39894	0.0061974	0.0061974\\
3.9526	0.3883	0.0060931	0.0060931\\
3.9638	0.37766	0.0058931	0.0058931\\
3.9717	0.36702	0.0058103	0.0058103\\
3.9827	0.35638	0.0056547	0.0056547\\
3.9938	0.34574	0.0056375	0.0056375\\
4.0041	0.33511	0.0055906	0.0055906\\
4.0134	0.32447	0.0058442	0.0058442\\
4.0267	0.31383	0.0056071	0.0056071\\
4.0386	0.30319	0.0059954	0.0059954\\
4.0508	0.29255	0.0056373	0.0056373\\
4.0624	0.28191	0.0058582	0.0058582\\
4.0725	0.27128	0.0058151	0.0058151\\
4.0888	0.26064	0.005724	0.005724\\
4.1049	0.25	0.0057219	0.0057219\\
4.117	0.23936	0.0059065	0.0059065\\
4.1302	0.22872	0.0059432	0.0059432\\
4.1455	0.21809	0.0060705	0.0060705\\
4.1599	0.20745	0.0060167	0.0060167\\
4.1742	0.19681	0.005967	0.005967\\
4.1902	0.18617	0.0058452	0.0058452\\
4.2074	0.17553	0.0059932	0.0059932\\
4.2233	0.16489	0.0063192	0.0063192\\
4.2398	0.15426	0.0075653	0.0075653\\
4.2602	0.14362	0.0076698	0.0076698\\
4.2789	0.13298	0.0079775	0.0079775\\
4.2994	0.12234	0.0083583	0.0083583\\
4.3276	0.1117	0.008497	0.008497\\
4.3546	0.10106	0.0080207	0.0080207\\
4.3817	0.090426	0.0088382	0.0088382\\
4.4171	0.079787	0.0098774	0.0098774\\
4.4522	0.069149	0.010243	0.010243\\
4.497	0.058511	0.011325	0.011325\\
4.5578	0.047872	0.013129	0.013129\\
4.6192	0.037234	0.011595	0.011595\\
4.7102	0.026596	0.014518	0.014518\\
4.8182	0.015957	0.018109	0.018109\\
5.1809	0.0053191	0.014189	0.014189\\
};
\addplot [color=mycolor1]
  table[row sep=crcr]{%
3.606	0.97811\\
3.6234	0.94606\\
3.6408	0.90438\\
3.6931	0.77583\\
4.1112	0.22618\\
4.5293	0.066215\\
4.9474	0.019465\\
5.3306	0.0063592\\
};
\addlegendentry{\ApproxParetoFitLegend}


\addplot[area legend, draw=none, fill=mycolor1, fill opacity=0.15, forget plot]
table[row sep=crcr] {%
x	y\\
3.606	0.98437\\
3.6234	0.95474\\
3.6408	0.91396\\
3.6583	0.86904\\
3.6757	0.82509\\
3.6931	0.78338\\
3.7105	0.74379\\
3.7279	0.70622\\
3.7454	0.67056\\
3.7628	0.63672\\
3.7802	0.6046\\
3.7976	0.57412\\
3.815	0.5452\\
3.8325	0.51775\\
3.8499	0.4917\\
3.8673	0.46698\\
3.8847	0.44353\\
3.9021	0.42127\\
3.9196	0.40016\\
3.937	0.38012\\
3.9544	0.3611\\
3.9718	0.34306\\
3.9893	0.32593\\
4.0067	0.30968\\
4.0241	0.29424\\
4.0415	0.27959\\
4.0589	0.26568\\
4.0764	0.25247\\
4.0938	0.23993\\
4.1112	0.22801\\
4.1286	0.2167\\
4.146	0.20595\\
4.1635	0.19573\\
4.1809	0.18603\\
4.1983	0.17681\\
4.2157	0.16805\\
4.2331	0.15973\\
4.2506	0.15182\\
4.268	0.14431\\
4.2854	0.13717\\
4.3028	0.13038\\
4.3202	0.12394\\
4.3377	0.11781\\
4.3551	0.11198\\
4.3725	0.10645\\
4.3899	0.10119\\
4.4073	0.096188\\
4.4248	0.091436\\
4.4422	0.08692\\
4.4596	0.082628\\
4.477	0.078548\\
4.4944	0.074671\\
4.5119	0.070985\\
4.5293	0.067482\\
4.5467	0.064152\\
4.5641	0.060988\\
4.5815	0.057979\\
4.599	0.05512\\
4.6164	0.052402\\
4.6338	0.049818\\
4.6512	0.047362\\
4.6686	0.045028\\
4.6861	0.042809\\
4.7035	0.040699\\
4.7209	0.038694\\
4.7383	0.036788\\
4.7558	0.034976\\
4.7732	0.033254\\
4.7906	0.031616\\
4.808	0.03006\\
4.8254	0.02858\\
4.8429	0.027173\\
4.8603	0.025836\\
4.8777	0.024565\\
4.8951	0.023356\\
4.9125	0.022207\\
4.93	0.021115\\
4.9474	0.020077\\
4.9648	0.019089\\
4.9822	0.018151\\
4.9996	0.017259\\
5.0171	0.01641\\
5.0345	0.015604\\
5.0519	0.014837\\
5.0693	0.014108\\
5.0867	0.013415\\
5.1042	0.012756\\
5.1216	0.01213\\
5.139	0.011534\\
5.1564	0.010968\\
5.1738	0.010429\\
5.1913	0.0099173\\
5.2087	0.0094305\\
5.2261	0.0089678\\
5.2435	0.0085277\\
5.2609	0.0081094\\
5.2784	0.0077116\\
5.2958	0.0073334\\
5.3132	0.0069737\\
5.3306	0.0066318\\
5.3306	0.0060865\\
5.3132	0.006407\\
5.2958	0.0067445\\
5.2784	0.0070998\\
5.2609	0.0074738\\
5.2435	0.0078676\\
5.2261	0.0082822\\
5.2087	0.0087187\\
5.1913	0.0091782\\
5.1738	0.0096621\\
5.1564	0.010171\\
5.139	0.010708\\
5.1216	0.011273\\
5.1042	0.011867\\
5.0867	0.012493\\
5.0693	0.013152\\
5.0519	0.013846\\
5.0345	0.014577\\
5.0171	0.015346\\
4.9996	0.016156\\
4.9822	0.017009\\
4.9648	0.017907\\
4.9474	0.018853\\
4.93	0.019848\\
4.9125	0.020897\\
4.8951	0.022001\\
4.8777	0.023163\\
4.8603	0.024387\\
4.8429	0.025676\\
4.8254	0.027033\\
4.808	0.028462\\
4.7906	0.029967\\
4.7732	0.031551\\
4.7558	0.033219\\
4.7383	0.034976\\
4.7209	0.036826\\
4.7035	0.038774\\
4.6861	0.040826\\
4.6686	0.042986\\
4.6512	0.045261\\
4.6338	0.047656\\
4.6164	0.050179\\
4.599	0.052835\\
4.5815	0.055633\\
4.5641	0.058579\\
4.5467	0.061681\\
4.5293	0.064948\\
4.5119	0.068388\\
4.4944	0.072011\\
4.477	0.075826\\
4.4596	0.079843\\
4.4422	0.084074\\
4.4248	0.08853\\
4.4073	0.093222\\
4.3899	0.098163\\
4.3725	0.10337\\
4.3551	0.10885\\
4.3377	0.11462\\
4.3202	0.12069\\
4.3028	0.12709\\
4.2854	0.13383\\
4.268	0.14093\\
4.2506	0.1484\\
4.2331	0.15627\\
4.2157	0.16456\\
4.1983	0.17328\\
4.1809	0.18247\\
4.1635	0.19215\\
4.146	0.20233\\
4.1286	0.21306\\
4.1112	0.22435\\
4.0938	0.23623\\
4.0764	0.24874\\
4.0589	0.26191\\
4.0415	0.27577\\
4.0241	0.29035\\
4.0067	0.3057\\
3.9893	0.32185\\
3.9718	0.33884\\
3.9544	0.35671\\
3.937	0.37551\\
3.9196	0.39529\\
3.9021	0.41609\\
3.8847	0.43796\\
3.8673	0.46098\\
3.8499	0.48518\\
3.8325	0.51064\\
3.815	0.53742\\
3.7976	0.5656\\
3.7802	0.59524\\
3.7628	0.62642\\
3.7454	0.65923\\
3.7279	0.69375\\
3.7105	0.73007\\
3.6931	0.76828\\
3.6757	0.8085\\
3.6583	0.85081\\
3.6408	0.8948\\
3.6234	0.93738\\
3.606	0.97185\\
}--cycle;
\addplot[only marks, mark=*, mark size=1.1180pt, color=black, fill=black, opacity=0.60, draw=none, mark options={draw=none,line width=0pt}] table[row sep=crcr]{%
x	y\\
3.5623	0.99468\\
3.5723	0.98404\\
3.573	0.9734\\
3.5753	0.96277\\
3.5766	0.95213\\
3.5781	0.94149\\
3.5826	0.93085\\
3.5851	0.92021\\
3.5911	0.90957\\
3.5942	0.89894\\
3.596	0.8883\\
3.6016	0.87766\\
3.603	0.86702\\
3.6047	0.85638\\
3.6086	0.84574\\
3.6086	0.83511\\
3.6097	0.82447\\
3.6142	0.81383\\
3.6261	0.80319\\
3.6383	0.79255\\
3.6464	0.78191\\
3.6535	0.77128\\
3.657	0.76064\\
3.6658	0.75\\
3.6678	0.73936\\
3.6687	0.72872\\
3.6702	0.71809\\
3.6723	0.70745\\
3.6793	0.69681\\
3.6867	0.68617\\
3.6949	0.67553\\
3.7015	0.66489\\
3.7031	0.65426\\
3.719	0.64362\\
3.7213	0.63298\\
3.7265	0.62234\\
3.7307	0.6117\\
3.737	0.60106\\
3.7374	0.59043\\
3.7423	0.57979\\
3.7487	0.56915\\
3.7677	0.55851\\
3.7685	0.54787\\
3.7768	0.53723\\
3.8005	0.5266\\
3.8015	0.51596\\
3.8022	0.50532\\
3.8032	0.49468\\
3.8039	0.48404\\
3.8108	0.4734\\
3.8214	0.46277\\
3.8405	0.45213\\
3.8605	0.44149\\
3.863	0.43085\\
3.8645	0.42021\\
3.8662	0.40957\\
3.8806	0.39894\\
3.8878	0.3883\\
3.8964	0.37766\\
3.905	0.36702\\
3.9294	0.35638\\
3.9303	0.34574\\
3.954	0.33511\\
3.9604	0.32447\\
3.9669	0.31383\\
3.9747	0.30319\\
3.9822	0.29255\\
3.9929	0.28191\\
4.0051	0.27128\\
4.0144	0.26064\\
4.0161	0.25\\
4.0433	0.23936\\
4.0578	0.22872\\
4.073	0.21809\\
4.0857	0.20745\\
4.0904	0.19681\\
4.1168	0.18617\\
4.1265	0.17553\\
4.127	0.16489\\
4.143	0.15426\\
4.1756	0.14362\\
4.2156	0.13298\\
4.3116	0.12234\\
4.3276	0.1117\\
4.3662	0.10106\\
4.4791	0.090426\\
4.4912	0.079787\\
4.517	0.069149\\
4.5743	0.058511\\
4.5914	0.047872\\
4.6138	0.037234\\
4.7898	0.026596\\
5.0876	0.015957\\
5.3101	0.0053191\\
};
\addlegendentry{\RealEmpiricCCDFLegend}

\addplot [color=black]
  table[row sep=crcr]{%
3.5608	1\\
5.3306	0.0055929\\
};
\addlegendentry{\RealParetoFitLegend}

      % \input{../../../../.tex/20/KS/SizeDistributionFittingsfig27.tex}
      \legend{}
    \end{groupplot}
  
    \begin{groupplot}[
      group style={
        group name=Row4,
        plotGroup,
      },
      plotRow,
      row2KeysWithLabels,
      plotLegendNE,
      %
      % xmin=3637.5,xmax=2.5003e+05,
      ymin=0.005121,ymax=1,
    ]
      \nextgroupplot[%
        at={($(Row3 c1r1.south west)-(0,\yPlotSep em)-(0,\plotHeight cm)$)},
        xmin=3.598,xmax=5.4115,
        % xmin=3.4,xmax=5.4,
        % ymin=0.0053191,ymax=1,
      ] % This file was created by matlab2tikz.
%
\definecolor{mycolor1}{rgb}{0.00000,0.44706,0.69804}%
\definecolor{mycolor2}{rgb}{0.00000,0.44700,0.74100}%
\definecolor{mycolor3}{rgb}{0.83529,0.36863,0.00000}%
\definecolor{mycolor4}{rgb}{0.49400,0.18400,0.55600}%
%
\addplot[only marks, mark=*, mark size=1.1180pt, color=mycolor1, fill=mycolor1, opacity=0.60, draw=none, mark options={draw=none,line width=0pt}] table[row sep=crcr]{%
x	y\\
3.6054	0.99468\\
3.6098	0.98404\\
3.6142	0.9734\\
3.6177	0.96277\\
3.6219	0.95213\\
3.6266	0.94149\\
3.6304	0.93085\\
3.6361	0.92021\\
3.64	0.90957\\
3.6454	0.89894\\
3.6509	0.8883\\
3.6553	0.87766\\
3.6597	0.86702\\
3.6642	0.85638\\
3.669	0.84574\\
3.6744	0.83511\\
3.6791	0.82447\\
3.6837	0.81383\\
3.6885	0.80319\\
3.6932	0.79255\\
3.699	0.78191\\
3.7045	0.77128\\
3.7096	0.76064\\
3.7159	0.75\\
3.7217	0.73936\\
3.7271	0.72872\\
3.7326	0.71809\\
3.7382	0.70745\\
3.7444	0.69681\\
3.7511	0.68617\\
3.7579	0.67553\\
3.7642	0.66489\\
3.771	0.65426\\
3.7772	0.64362\\
3.7827	0.63298\\
3.7895	0.62234\\
3.7969	0.6117\\
3.8022	0.60106\\
3.8087	0.59043\\
3.8158	0.57979\\
3.8225	0.56915\\
3.83	0.55851\\
3.8371	0.54787\\
3.8436	0.53723\\
3.8503	0.5266\\
3.858	0.51596\\
3.8673	0.50532\\
3.8759	0.49468\\
3.8835	0.48404\\
3.8911	0.4734\\
3.9004	0.46277\\
3.9082	0.45213\\
3.9162	0.44149\\
3.9249	0.43085\\
3.9328	0.42021\\
3.9406	0.40957\\
3.949	0.39894\\
3.9577	0.3883\\
3.9678	0.37766\\
3.98	0.36702\\
3.9908	0.35638\\
4.0008	0.34574\\
4.0113	0.33511\\
4.0202	0.32447\\
4.0318	0.31383\\
4.0441	0.30319\\
4.0554	0.29255\\
4.0676	0.28191\\
4.0793	0.27128\\
4.0914	0.26064\\
4.1033	0.25\\
4.1167	0.23936\\
4.1321	0.22872\\
4.1494	0.21809\\
4.1669	0.20745\\
4.1856	0.19681\\
4.2028	0.18617\\
4.2197	0.17553\\
4.2417	0.16489\\
4.2609	0.15426\\
4.2818	0.14362\\
4.3051	0.13298\\
4.33	0.12234\\
4.3521	0.1117\\
4.3782	0.10106\\
4.4054	0.090426\\
4.4416	0.079787\\
4.4781	0.069149\\
4.5216	0.058511\\
4.574	0.047872\\
4.636	0.037234\\
4.7242	0.026596\\
4.8252	0.015957\\
5.1746	0.0053191\\
};
\addlegendentry{\ExactEmpiricCCDFLegend}

\addplot [color=mycolor1, only marks, every error bar/.append style={opacity=0.30}, mark=*, mark size=0pt, draw=none, forget plot]
 plot [error bars/.cd, x dir=both, x explicit, error bar style={line width=1pt, color=mycolor1}, error mark options={mark=none,mark size=0pt}]
 table[row sep=crcr, x error plus index=2, x error minus index=3]{%
3.6054	0.99468	0.004187	0.004187\\
3.6098	0.98404	0.003977	0.003977\\
3.6142	0.9734	0.0040165	0.0040165\\
3.6177	0.96277	0.0040403	0.0040403\\
3.6219	0.95213	0.0039714	0.0039714\\
3.6266	0.94149	0.00393	0.00393\\
3.6304	0.93085	0.00398	0.00398\\
3.6361	0.92021	0.0040894	0.0040894\\
3.64	0.90957	0.0042313	0.0042313\\
3.6454	0.89894	0.004231	0.004231\\
3.6509	0.8883	0.0044287	0.0044287\\
3.6553	0.87766	0.0044395	0.0044395\\
3.6597	0.86702	0.0044097	0.0044097\\
3.6642	0.85638	0.0045406	0.0045406\\
3.669	0.84574	0.0045804	0.0045804\\
3.6744	0.83511	0.0045103	0.0045103\\
3.6791	0.82447	0.0045155	0.0045155\\
3.6837	0.81383	0.0044785	0.0044785\\
3.6885	0.80319	0.0044161	0.0044161\\
3.6932	0.79255	0.0044769	0.0044769\\
3.699	0.78191	0.004542	0.004542\\
3.7045	0.77128	0.0045765	0.0045765\\
3.7096	0.76064	0.0043954	0.0043954\\
3.7159	0.75	0.004498	0.004498\\
3.7217	0.73936	0.0045568	0.0045568\\
3.7271	0.72872	0.004502	0.004502\\
3.7326	0.71809	0.0045084	0.0045084\\
3.7382	0.70745	0.0046824	0.0046824\\
3.7444	0.69681	0.0047522	0.0047522\\
3.7511	0.68617	0.0049632	0.0049632\\
3.7579	0.67553	0.0048624	0.0048624\\
3.7642	0.66489	0.0048951	0.0048951\\
3.771	0.65426	0.0048301	0.0048301\\
3.7772	0.64362	0.0047613	0.0047613\\
3.7827	0.63298	0.0047343	0.0047343\\
3.7895	0.62234	0.004828	0.004828\\
3.7969	0.6117	0.0050184	0.0050184\\
3.8022	0.60106	0.0051359	0.0051359\\
3.8087	0.59043	0.0052829	0.0052829\\
3.8158	0.57979	0.0054389	0.0054389\\
3.8225	0.56915	0.0055371	0.0055371\\
3.83	0.55851	0.005547	0.005547\\
3.8371	0.54787	0.0055845	0.0055845\\
3.8436	0.53723	0.0056346	0.0056346\\
3.8503	0.5266	0.0055049	0.0055049\\
3.858	0.51596	0.0052181	0.0052181\\
3.8673	0.50532	0.0052516	0.0052516\\
3.8759	0.49468	0.0053721	0.0053721\\
3.8835	0.48404	0.0051237	0.0051237\\
3.8911	0.4734	0.0051448	0.0051448\\
3.9004	0.46277	0.0056128	0.0056128\\
3.9082	0.45213	0.0057038	0.0057038\\
3.9162	0.44149	0.0058207	0.0058207\\
3.9249	0.43085	0.0060439	0.0060439\\
3.9328	0.42021	0.0061768	0.0061768\\
3.9406	0.40957	0.006023	0.006023\\
3.949	0.39894	0.0062245	0.0062245\\
3.9577	0.3883	0.0064029	0.0064029\\
3.9678	0.37766	0.006833	0.006833\\
3.98	0.36702	0.0067149	0.0067149\\
3.9908	0.35638	0.0068135	0.0068135\\
4.0008	0.34574	0.0068598	0.0068598\\
4.0113	0.33511	0.0068246	0.0068246\\
4.0202	0.32447	0.0069444	0.0069444\\
4.0318	0.31383	0.0067454	0.0067454\\
4.0441	0.30319	0.0064174	0.0064174\\
4.0554	0.29255	0.0065338	0.0065338\\
4.0676	0.28191	0.0060628	0.0060628\\
4.0793	0.27128	0.0060461	0.0060461\\
4.0914	0.26064	0.0061657	0.0061657\\
4.1033	0.25	0.0062693	0.0062693\\
4.1167	0.23936	0.0063805	0.0063805\\
4.1321	0.22872	0.0066514	0.0066514\\
4.1494	0.21809	0.0069426	0.0069426\\
4.1669	0.20745	0.0072838	0.0072838\\
4.1856	0.19681	0.0076203	0.0076203\\
4.2028	0.18617	0.0080221	0.0080221\\
4.2197	0.17553	0.0079141	0.0079141\\
4.2417	0.16489	0.0083169	0.0083169\\
4.2609	0.15426	0.0082648	0.0082648\\
4.2818	0.14362	0.0083327	0.0083327\\
4.3051	0.13298	0.0078595	0.0078595\\
4.33	0.12234	0.0080214	0.0080214\\
4.3521	0.1117	0.0077859	0.0077859\\
4.3782	0.10106	0.0082103	0.0082103\\
4.4054	0.090426	0.0084553	0.0084553\\
4.4416	0.079787	0.009553	0.009553\\
4.4781	0.069149	0.0095889	0.0095889\\
4.5216	0.058511	0.0095988	0.0095988\\
4.574	0.047872	0.011021	0.011021\\
4.636	0.037234	0.013906	0.013906\\
4.7242	0.026596	0.014438	0.014438\\
4.8252	0.015957	0.016348	0.016348\\
5.1746	0.0053191	0.014408	0.014408\\
};
\addplot [color=mycolor1]
  table[row sep=crcr]{%
3.6037	0.97822\\
3.6214	0.9464\\
3.6392	0.90579\\
3.6747	0.8189\\
4.0651	0.26705\\
4.4556	0.087457\\
4.846	0.028764\\
5.2365	0.0094998\\
5.3607	0.0066835\\
};
\addlegendentry{\ExactParetoFitLegend}


\addplot[area legend, draw=none, fill=mycolor1, fill opacity=0.15, forget plot]
table[row sep=crcr] {%
x	y\\
3.6037	0.98434\\
3.6214	0.95496\\
3.6392	0.91543\\
3.6569	0.87097\\
3.6747	0.82724\\
3.6924	0.78573\\
3.7102	0.74632\\
3.7279	0.7089\\
3.7457	0.67338\\
3.7634	0.63966\\
3.7812	0.60765\\
3.7989	0.57726\\
3.8167	0.54842\\
3.8344	0.52104\\
3.8522	0.49506\\
3.8699	0.47039\\
3.8877	0.44699\\
3.9054	0.42477\\
3.9232	0.40369\\
3.9409	0.38367\\
3.9586	0.36467\\
3.9764	0.34664\\
3.9941	0.32951\\
4.0119	0.31325\\
4.0296	0.2978\\
4.0474	0.28313\\
4.0651	0.26919\\
4.0829	0.25595\\
4.1006	0.24336\\
4.1184	0.2314\\
4.1361	0.22003\\
4.1539	0.20923\\
4.1716	0.19896\\
4.1894	0.1892\\
4.2071	0.17992\\
4.2249	0.1711\\
4.2426	0.16272\\
4.2604	0.15475\\
4.2781	0.14717\\
4.2959	0.13996\\
4.3136	0.13311\\
4.3313	0.12659\\
4.3491	0.1204\\
4.3668	0.11451\\
4.3846	0.10891\\
4.4023	0.10358\\
4.4201	0.098521\\
4.4378	0.093706\\
4.4556	0.089127\\
4.4733	0.084772\\
4.4911	0.080631\\
4.5088	0.076694\\
4.5266	0.072949\\
4.5443	0.069388\\
4.5621	0.066001\\
4.5798	0.06278\\
4.5976	0.059717\\
4.6153	0.056804\\
4.6331	0.054033\\
4.6508	0.051399\\
4.6686	0.048893\\
4.6863	0.046509\\
4.704	0.044243\\
4.7218	0.042087\\
4.7395	0.040036\\
4.7573	0.038086\\
4.775	0.036231\\
4.7928	0.034467\\
4.8105	0.032789\\
4.8283	0.031193\\
4.846	0.029675\\
4.8638	0.028231\\
4.8815	0.026857\\
4.8993	0.025551\\
4.917	0.024308\\
4.9348	0.023126\\
4.9525	0.022002\\
4.9703	0.020932\\
4.988	0.019915\\
5.0058	0.018947\\
5.0235	0.018027\\
5.0412	0.017151\\
5.059	0.016318\\
5.0767	0.015526\\
5.0945	0.014772\\
5.1122	0.014055\\
5.13	0.013373\\
5.1477	0.012724\\
5.1655	0.012107\\
5.1832	0.011519\\
5.201	0.010961\\
5.2187	0.010429\\
5.2365	0.0099235\\
5.2542	0.0094424\\
5.272	0.0089848\\
5.2897	0.0085494\\
5.3075	0.0081352\\
5.3252	0.0077411\\
5.343	0.0073662\\
5.3607	0.0070095\\
5.3607	0.0063576\\
5.343	0.0066891\\
5.3252	0.007038\\
5.3075	0.0074051\\
5.2897	0.0077915\\
5.272	0.008198\\
5.2542	0.0086258\\
5.2365	0.0090761\\
5.2187	0.0095499\\
5.201	0.010048\\
5.1832	0.010573\\
5.1655	0.011125\\
5.1477	0.011707\\
5.13	0.012318\\
5.1122	0.012962\\
5.0945	0.013639\\
5.0767	0.014352\\
5.059	0.015102\\
5.0412	0.015892\\
5.0235	0.016723\\
5.0058	0.017598\\
4.988	0.018518\\
4.9703	0.019487\\
4.9525	0.020507\\
4.9348	0.02158\\
4.917	0.022709\\
4.8993	0.023898\\
4.8815	0.025149\\
4.8638	0.026466\\
4.846	0.027852\\
4.8283	0.029311\\
4.8105	0.030847\\
4.7928	0.032463\\
4.775	0.034165\\
4.7573	0.035955\\
4.7395	0.03784\\
4.7218	0.039824\\
4.704	0.041912\\
4.6863	0.044111\\
4.6686	0.046424\\
4.6508	0.04886\\
4.6331	0.051423\\
4.6153	0.054122\\
4.5976	0.056963\\
4.5798	0.059953\\
4.5621	0.0631\\
4.5443	0.066413\\
4.5266	0.069901\\
4.5088	0.073572\\
4.4911	0.077437\\
4.4733	0.081505\\
4.4556	0.085788\\
4.4378	0.090296\\
4.4201	0.095042\\
4.4023	0.10004\\
4.3846	0.1053\\
4.3668	0.11083\\
4.3491	0.11666\\
4.3313	0.12279\\
4.3136	0.12925\\
4.2959	0.13605\\
4.2781	0.14321\\
4.2604	0.15074\\
4.2426	0.15867\\
4.2249	0.16702\\
4.2071	0.1758\\
4.1894	0.18505\\
4.1716	0.19478\\
4.1539	0.20503\\
4.1361	0.21581\\
4.1184	0.22716\\
4.1006	0.23911\\
4.0829	0.25168\\
4.0651	0.2649\\
4.0474	0.27882\\
4.0296	0.29346\\
4.0119	0.30886\\
3.9941	0.32506\\
3.9764	0.3421\\
3.9586	0.36002\\
3.9409	0.37886\\
3.9232	0.39866\\
3.9054	0.41948\\
3.8877	0.44137\\
3.8699	0.46438\\
3.8522	0.48857\\
3.8344	0.51399\\
3.8167	0.54072\\
3.7989	0.56883\\
3.7812	0.59837\\
3.7634	0.62944\\
3.7457	0.66211\\
3.7279	0.69647\\
3.7102	0.7326\\
3.6924	0.77059\\
3.6747	0.81055\\
3.6569	0.85258\\
3.6392	0.89615\\
3.6214	0.93785\\
3.6037	0.97209\\
}--cycle;
\addplot[only marks, mark=*, mark size=1.1180pt, color=mycolor3, fill=mycolor3, opacity=0.60, draw=none, mark options={draw=none,line width=0pt}] table[row sep=crcr]{%
x	y\\
3.6043	0.99468\\
3.6085	0.98404\\
3.6129	0.9734\\
3.6166	0.96277\\
3.6217	0.95213\\
3.6257	0.94149\\
3.6297	0.93085\\
3.6336	0.92021\\
3.6386	0.90957\\
3.6436	0.89894\\
3.6485	0.8883\\
3.6534	0.87766\\
3.6578	0.86702\\
3.6617	0.85638\\
3.6672	0.84574\\
3.6724	0.83511\\
3.6775	0.82447\\
3.683	0.81383\\
3.6877	0.80319\\
3.6931	0.79255\\
3.6981	0.78191\\
3.7035	0.77128\\
3.7087	0.76064\\
3.7142	0.75\\
3.7195	0.73936\\
3.7257	0.72872\\
3.7321	0.71809\\
3.7384	0.70745\\
3.7448	0.69681\\
3.7505	0.68617\\
3.7557	0.67553\\
3.7608	0.66489\\
3.7667	0.65426\\
3.7732	0.64362\\
3.7794	0.63298\\
3.7861	0.62234\\
3.7926	0.6117\\
3.7986	0.60106\\
3.8051	0.59043\\
3.8111	0.57979\\
3.8186	0.56915\\
3.8258	0.55851\\
3.8344	0.54787\\
3.8405	0.53723\\
3.848	0.5266\\
3.855	0.51596\\
3.8618	0.50532\\
3.8693	0.49468\\
3.8771	0.48404\\
3.8852	0.4734\\
3.8931	0.46277\\
3.9016	0.45213\\
3.9092	0.44149\\
3.9181	0.43085\\
3.9273	0.42021\\
3.9357	0.40957\\
3.9432	0.39894\\
3.9533	0.3883\\
3.9637	0.37766\\
3.9738	0.36702\\
3.9848	0.35638\\
3.9951	0.34574\\
4.0063	0.33511\\
4.0195	0.32447\\
4.0316	0.31383\\
4.0451	0.30319\\
4.0541	0.29255\\
4.0665	0.28191\\
4.0793	0.27128\\
4.0929	0.26064\\
4.1075	0.25\\
4.1207	0.23936\\
4.1351	0.22872\\
4.1514	0.21809\\
4.1699	0.20745\\
4.1822	0.19681\\
4.1983	0.18617\\
4.2172	0.17553\\
4.2352	0.16489\\
4.2564	0.15426\\
4.2767	0.14362\\
4.2966	0.13298\\
4.3207	0.12234\\
4.3486	0.1117\\
4.3743	0.10106\\
4.4051	0.090426\\
4.439	0.079787\\
4.4763	0.069149\\
4.5194	0.058511\\
4.5656	0.047872\\
4.635	0.037234\\
4.7318	0.026596\\
4.8427	0.015957\\
5.1833	0.0053191\\
};
\addlegendentry{\ApproxEmpiricCCDFLegend}

\addplot [color=mycolor3, only marks, every error bar/.append style={opacity=0.30}, mark=*, mark size=0pt, draw=none, forget plot]
 plot [error bars/.cd, x dir=both, x explicit, error bar style={line width=1pt, color=mycolor3}, error mark options={mark=none,mark size=0pt}]
 table[row sep=crcr, x error plus index=2, x error minus index=3]{%
3.6043	0.99468	0.0045572	0.0045572\\
3.6085	0.98404	0.0043712	0.0043712\\
3.6129	0.9734	0.0045168	0.0045168\\
3.6166	0.96277	0.0044731	0.0044731\\
3.6217	0.95213	0.0044265	0.0044265\\
3.6257	0.94149	0.0044362	0.0044362\\
3.6297	0.93085	0.0044388	0.0044388\\
3.6336	0.92021	0.0043707	0.0043707\\
3.6386	0.90957	0.0046291	0.0046291\\
3.6436	0.89894	0.0046805	0.0046805\\
3.6485	0.8883	0.0046062	0.0046062\\
3.6534	0.87766	0.0047192	0.0047192\\
3.6578	0.86702	0.0047457	0.0047457\\
3.6617	0.85638	0.0047213	0.0047213\\
3.6672	0.84574	0.0048522	0.0048522\\
3.6724	0.83511	0.0047178	0.0047178\\
3.6775	0.82447	0.0045714	0.0045714\\
3.683	0.81383	0.0044815	0.0044815\\
3.6877	0.80319	0.0045796	0.0045796\\
3.6931	0.79255	0.0046427	0.0046427\\
3.6981	0.78191	0.0047601	0.0047601\\
3.7035	0.77128	0.0049339	0.0049339\\
3.7087	0.76064	0.0048633	0.0048633\\
3.7142	0.75	0.0048943	0.0048943\\
3.7195	0.73936	0.0047886	0.0047886\\
3.7257	0.72872	0.0047314	0.0047314\\
3.7321	0.71809	0.0048394	0.0048394\\
3.7384	0.70745	0.0048539	0.0048539\\
3.7448	0.69681	0.005076	0.005076\\
3.7505	0.68617	0.0050712	0.0050712\\
3.7557	0.67553	0.0050637	0.0050637\\
3.7608	0.66489	0.0051589	0.0051589\\
3.7667	0.65426	0.0051728	0.0051728\\
3.7732	0.64362	0.0050904	0.0050904\\
3.7794	0.63298	0.0050729	0.0050729\\
3.7861	0.62234	0.0049622	0.0049622\\
3.7926	0.6117	0.0050011	0.0050011\\
3.7986	0.60106	0.0051642	0.0051642\\
3.8051	0.59043	0.0050796	0.0050796\\
3.8111	0.57979	0.0051542	0.0051542\\
3.8186	0.56915	0.0054195	0.0054195\\
3.8258	0.55851	0.0056046	0.0056046\\
3.8344	0.54787	0.0058109	0.0058109\\
3.8405	0.53723	0.0059958	0.0059958\\
3.848	0.5266	0.0057197	0.0057197\\
3.855	0.51596	0.0057667	0.0057667\\
3.8618	0.50532	0.0059162	0.0059162\\
3.8693	0.49468	0.005978	0.005978\\
3.8771	0.48404	0.0059461	0.0059461\\
3.8852	0.4734	0.005987	0.005987\\
3.8931	0.46277	0.0058947	0.0058947\\
3.9016	0.45213	0.0059886	0.0059886\\
3.9092	0.44149	0.0058124	0.0058124\\
3.9181	0.43085	0.0060973	0.0060973\\
3.9273	0.42021	0.0063055	0.0063055\\
3.9357	0.40957	0.0066836	0.0066836\\
3.9432	0.39894	0.0065301	0.0065301\\
3.9533	0.3883	0.006	0.006\\
3.9637	0.37766	0.0060566	0.0060566\\
3.9738	0.36702	0.0063024	0.0063024\\
3.9848	0.35638	0.0065619	0.0065619\\
3.9951	0.34574	0.0066914	0.0066914\\
4.0063	0.33511	0.0065221	0.0065221\\
4.0195	0.32447	0.0062434	0.0062434\\
4.0316	0.31383	0.0062868	0.0062868\\
4.0451	0.30319	0.0062701	0.0062701\\
4.0541	0.29255	0.0063707	0.0063707\\
4.0665	0.28191	0.0063128	0.0063128\\
4.0793	0.27128	0.0063457	0.0063457\\
4.0929	0.26064	0.0064012	0.0064012\\
4.1075	0.25	0.0065842	0.0065842\\
4.1207	0.23936	0.0067354	0.0067354\\
4.1351	0.22872	0.0071575	0.0071575\\
4.1514	0.21809	0.0069122	0.0069122\\
4.1699	0.20745	0.0074422	0.0074422\\
4.1822	0.19681	0.0075591	0.0075591\\
4.1983	0.18617	0.0073044	0.0073044\\
4.2172	0.17553	0.0075554	0.0075554\\
4.2352	0.16489	0.0067471	0.0067471\\
4.2564	0.15426	0.0070369	0.0070369\\
4.2767	0.14362	0.0073888	0.0073888\\
4.2966	0.13298	0.0081069	0.0081069\\
4.3207	0.12234	0.0084621	0.0084621\\
4.3486	0.1117	0.0088645	0.0088645\\
4.3743	0.10106	0.0094713	0.0094713\\
4.4051	0.090426	0.0089463	0.0089463\\
4.439	0.079787	0.0095141	0.0095141\\
4.4763	0.069149	0.0092913	0.0092913\\
4.5194	0.058511	0.010771	0.010771\\
4.5656	0.047872	0.01169	0.01169\\
4.635	0.037234	0.012228	0.012228\\
4.7318	0.026596	0.015596	0.015596\\
4.8427	0.015957	0.017427	0.017427\\
5.1833	0.0053191	0.01767	0.01767\\
};
\addplot [color=mycolor3]
  table[row sep=crcr]{%
3.604	0.97462\\
3.6217	0.94187\\
3.6395	0.89995\\
3.6572	0.85753\\
3.7282	0.69936\\
4.1008	0.23927\\
4.4735	0.082203\\
4.8461	0.028358\\
5.2365	0.0093391\\
5.3607	0.0065647\\
};
\addlegendentry{\ApproxParetoFitLegend}


\addplot[area legend, draw=none, fill=mycolor3, fill opacity=0.15, forget plot]
table[row sep=crcr] {%
x	y\\
3.604	0.98135\\
3.6217	0.95082\\
3.6395	0.90965\\
3.6572	0.86747\\
3.675	0.82471\\
3.6927	0.78342\\
3.7105	0.74396\\
3.7282	0.70651\\
3.746	0.67096\\
3.7637	0.63722\\
3.7814	0.60519\\
3.7992	0.5748\\
3.8169	0.54596\\
3.8347	0.51858\\
3.8524	0.49261\\
3.8702	0.46796\\
3.8879	0.44456\\
3.9057	0.42237\\
3.9234	0.4013\\
3.9411	0.38132\\
3.9589	0.36235\\
3.9766	0.34434\\
3.9944	0.32726\\
4.0121	0.31103\\
4.0299	0.29563\\
4.0476	0.28101\\
4.0654	0.26712\\
4.0831	0.25393\\
4.1008	0.2414\\
4.1186	0.22949\\
4.1363	0.21818\\
4.1541	0.20744\\
4.1718	0.19722\\
4.1896	0.18752\\
4.2073	0.17829\\
4.2251	0.16953\\
4.2428	0.16119\\
4.2605	0.15327\\
4.2783	0.14574\\
4.296	0.13859\\
4.3138	0.13178\\
4.3315	0.12532\\
4.3493	0.11917\\
4.367	0.11332\\
4.3848	0.10776\\
4.4025	0.10248\\
4.4202	0.097458\\
4.438	0.092682\\
4.4557	0.088141\\
4.4735	0.083823\\
4.4912	0.079718\\
4.509	0.075814\\
4.5267	0.072103\\
4.5445	0.068574\\
4.5622	0.065218\\
4.5799	0.062027\\
4.5977	0.058993\\
4.6154	0.056108\\
4.6332	0.053364\\
4.6509	0.050755\\
4.6687	0.048275\\
4.6864	0.045915\\
4.7042	0.043672\\
4.7219	0.041538\\
4.7396	0.03951\\
4.7574	0.03758\\
4.7751	0.035745\\
4.7929	0.034\\
4.8106	0.032341\\
4.8284	0.030763\\
4.8461	0.029262\\
4.8639	0.027834\\
4.8816	0.026477\\
4.8993	0.025186\\
4.9171	0.023958\\
4.9348	0.02279\\
4.9526	0.021679\\
4.9703	0.020623\\
4.9881	0.019618\\
5.0058	0.018662\\
5.0236	0.017753\\
5.0413	0.016889\\
5.0591	0.016067\\
5.0768	0.015285\\
5.0945	0.014541\\
5.1123	0.013833\\
5.13	0.01316\\
5.1478	0.01252\\
5.1655	0.011911\\
5.1833	0.011332\\
5.201	0.010781\\
5.2188	0.010257\\
5.2365	0.0097583\\
5.2542	0.0092841\\
5.272	0.0088331\\
5.2897	0.008404\\
5.3075	0.0079958\\
5.3252	0.0076076\\
5.343	0.0072382\\
5.3607	0.0068869\\
5.3607	0.0062426\\
5.343	0.006569\\
5.3252	0.0069124\\
5.3075	0.0072739\\
5.2897	0.0076544\\
5.272	0.0080548\\
5.2542	0.0084763\\
5.2365	0.0089198\\
5.2188	0.0093867\\
5.201	0.009878\\
5.1833	0.010395\\
5.1655	0.01094\\
5.1478	0.011512\\
5.13	0.012116\\
5.1123	0.01275\\
5.0945	0.013418\\
5.0768	0.014122\\
5.0591	0.014862\\
5.0413	0.015641\\
5.0236	0.016461\\
5.0058	0.017324\\
4.9881	0.018233\\
4.9703	0.019189\\
4.9526	0.020196\\
4.9348	0.021256\\
4.9171	0.022372\\
4.8993	0.023546\\
4.8816	0.024782\\
4.8639	0.026084\\
4.8461	0.027454\\
4.8284	0.028896\\
4.8106	0.030414\\
4.7929	0.032012\\
4.7751	0.033694\\
4.7574	0.035465\\
4.7396	0.03733\\
4.7219	0.039293\\
4.7042	0.041359\\
4.6864	0.043534\\
4.6687	0.045825\\
4.6509	0.048236\\
4.6332	0.050774\\
4.6154	0.053446\\
4.5977	0.05626\\
4.5799	0.059222\\
4.5622	0.06234\\
4.5445	0.065623\\
4.5267	0.069079\\
4.509	0.072718\\
4.4912	0.07655\\
4.4735	0.080583\\
4.4557	0.08483\\
4.438	0.089302\\
4.4202	0.094009\\
4.4025	0.098966\\
4.3848	0.10418\\
4.367	0.10968\\
4.3493	0.11546\\
4.3315	0.12155\\
4.3138	0.12796\\
4.296	0.13472\\
4.2783	0.14182\\
4.2605	0.1493\\
4.2428	0.15718\\
4.2251	0.16548\\
4.2073	0.17421\\
4.1896	0.1834\\
4.1718	0.19307\\
4.1541	0.20326\\
4.1363	0.21398\\
4.1186	0.22526\\
4.1008	0.23714\\
4.0831	0.24963\\
4.0654	0.26278\\
4.0476	0.27662\\
4.0299	0.29118\\
4.0121	0.30649\\
3.9944	0.3226\\
3.9766	0.33954\\
3.9589	0.35735\\
3.9411	0.37608\\
3.9234	0.39577\\
3.9057	0.41648\\
3.8879	0.43825\\
3.8702	0.46114\\
3.8524	0.4852\\
3.8347	0.5105\\
3.8169	0.53711\\
3.7992	0.56508\\
3.7814	0.5945\\
3.7637	0.62545\\
3.746	0.65799\\
3.7282	0.69221\\
3.7105	0.72821\\
3.6927	0.76608\\
3.675	0.80588\\
3.6572	0.84758\\
3.6395	0.89026\\
3.6217	0.93292\\
3.604	0.9679\\
}--cycle;
      % \input{../../../../.tex/20/KS/SizeDistributionFittingsfig34.tex}
      \legend{}
  
      \nextgroupplot[%
        xmin=3.5608,xmax=5.3544,
        % xmin=3.4,xmax=5.4,
        % ymin=0.005275,ymax=1,
      ] % This file was created by matlab2tikz.
%
\definecolor{mycolor1}{rgb}{0.00000,0.44706,0.69804}%
\definecolor{mycolor2}{rgb}{0.00000,0.44700,0.74100}%
\definecolor{mycolor3}{rgb}{0.49400,0.18400,0.55600}%
%
\addplot[only marks, mark=*, mark size=1.1180pt, color=mycolor1, fill=mycolor1, opacity=0.60, draw=none, mark options={draw=none,line width=0pt}] table[row sep=crcr]{%
x	y\\
3.6054	0.99468\\
3.6098	0.98404\\
3.6142	0.9734\\
3.6177	0.96277\\
3.6219	0.95213\\
3.6266	0.94149\\
3.6304	0.93085\\
3.6361	0.92021\\
3.64	0.90957\\
3.6454	0.89894\\
3.6509	0.8883\\
3.6553	0.87766\\
3.6597	0.86702\\
3.6642	0.85638\\
3.669	0.84574\\
3.6744	0.83511\\
3.6791	0.82447\\
3.6837	0.81383\\
3.6885	0.80319\\
3.6932	0.79255\\
3.699	0.78191\\
3.7045	0.77128\\
3.7096	0.76064\\
3.7159	0.75\\
3.7217	0.73936\\
3.7271	0.72872\\
3.7326	0.71809\\
3.7382	0.70745\\
3.7444	0.69681\\
3.7511	0.68617\\
3.7579	0.67553\\
3.7642	0.66489\\
3.771	0.65426\\
3.7772	0.64362\\
3.7827	0.63298\\
3.7895	0.62234\\
3.7969	0.6117\\
3.8022	0.60106\\
3.8087	0.59043\\
3.8158	0.57979\\
3.8225	0.56915\\
3.83	0.55851\\
3.8371	0.54787\\
3.8436	0.53723\\
3.8503	0.5266\\
3.858	0.51596\\
3.8673	0.50532\\
3.8759	0.49468\\
3.8835	0.48404\\
3.8911	0.4734\\
3.9004	0.46277\\
3.9082	0.45213\\
3.9162	0.44149\\
3.9249	0.43085\\
3.9328	0.42021\\
3.9406	0.40957\\
3.949	0.39894\\
3.9577	0.3883\\
3.9678	0.37766\\
3.98	0.36702\\
3.9908	0.35638\\
4.0008	0.34574\\
4.0113	0.33511\\
4.0202	0.32447\\
4.0318	0.31383\\
4.0441	0.30319\\
4.0554	0.29255\\
4.0676	0.28191\\
4.0793	0.27128\\
4.0914	0.26064\\
4.1033	0.25\\
4.1167	0.23936\\
4.1321	0.22872\\
4.1494	0.21809\\
4.1669	0.20745\\
4.1856	0.19681\\
4.2028	0.18617\\
4.2197	0.17553\\
4.2417	0.16489\\
4.2609	0.15426\\
4.2818	0.14362\\
4.3051	0.13298\\
4.33	0.12234\\
4.3521	0.1117\\
4.3782	0.10106\\
4.4054	0.090426\\
4.4416	0.079787\\
4.4781	0.069149\\
4.5216	0.058511\\
4.574	0.047872\\
4.636	0.037234\\
4.7242	0.026596\\
4.8252	0.015957\\
5.1746	0.0053191\\
};
\addlegendentry{\ExactEmpiricCCDFLegend}

\addplot [color=mycolor1, only marks, every error bar/.append style={opacity=0.30}, mark=*, mark size=0pt, draw=none, forget plot]
 plot [error bars/.cd, x dir=both, x explicit, error bar style={line width=1pt, color=mycolor1}, error mark options={mark=none,mark size=0pt}]
 table[row sep=crcr, x error plus index=2, x error minus index=3]{%
3.6054	0.99468	0.004187	0.004187\\
3.6098	0.98404	0.003977	0.003977\\
3.6142	0.9734	0.0040165	0.0040165\\
3.6177	0.96277	0.0040403	0.0040403\\
3.6219	0.95213	0.0039714	0.0039714\\
3.6266	0.94149	0.00393	0.00393\\
3.6304	0.93085	0.00398	0.00398\\
3.6361	0.92021	0.0040894	0.0040894\\
3.64	0.90957	0.0042313	0.0042313\\
3.6454	0.89894	0.004231	0.004231\\
3.6509	0.8883	0.0044287	0.0044287\\
3.6553	0.87766	0.0044395	0.0044395\\
3.6597	0.86702	0.0044097	0.0044097\\
3.6642	0.85638	0.0045406	0.0045406\\
3.669	0.84574	0.0045804	0.0045804\\
3.6744	0.83511	0.0045103	0.0045103\\
3.6791	0.82447	0.0045155	0.0045155\\
3.6837	0.81383	0.0044785	0.0044785\\
3.6885	0.80319	0.0044161	0.0044161\\
3.6932	0.79255	0.0044769	0.0044769\\
3.699	0.78191	0.004542	0.004542\\
3.7045	0.77128	0.0045765	0.0045765\\
3.7096	0.76064	0.0043954	0.0043954\\
3.7159	0.75	0.004498	0.004498\\
3.7217	0.73936	0.0045568	0.0045568\\
3.7271	0.72872	0.004502	0.004502\\
3.7326	0.71809	0.0045084	0.0045084\\
3.7382	0.70745	0.0046824	0.0046824\\
3.7444	0.69681	0.0047522	0.0047522\\
3.7511	0.68617	0.0049632	0.0049632\\
3.7579	0.67553	0.0048624	0.0048624\\
3.7642	0.66489	0.0048951	0.0048951\\
3.771	0.65426	0.0048301	0.0048301\\
3.7772	0.64362	0.0047613	0.0047613\\
3.7827	0.63298	0.0047343	0.0047343\\
3.7895	0.62234	0.004828	0.004828\\
3.7969	0.6117	0.0050184	0.0050184\\
3.8022	0.60106	0.0051359	0.0051359\\
3.8087	0.59043	0.0052829	0.0052829\\
3.8158	0.57979	0.0054389	0.0054389\\
3.8225	0.56915	0.0055371	0.0055371\\
3.83	0.55851	0.005547	0.005547\\
3.8371	0.54787	0.0055845	0.0055845\\
3.8436	0.53723	0.0056346	0.0056346\\
3.8503	0.5266	0.0055049	0.0055049\\
3.858	0.51596	0.0052181	0.0052181\\
3.8673	0.50532	0.0052516	0.0052516\\
3.8759	0.49468	0.0053721	0.0053721\\
3.8835	0.48404	0.0051237	0.0051237\\
3.8911	0.4734	0.0051448	0.0051448\\
3.9004	0.46277	0.0056128	0.0056128\\
3.9082	0.45213	0.0057038	0.0057038\\
3.9162	0.44149	0.0058207	0.0058207\\
3.9249	0.43085	0.0060439	0.0060439\\
3.9328	0.42021	0.0061768	0.0061768\\
3.9406	0.40957	0.006023	0.006023\\
3.949	0.39894	0.0062245	0.0062245\\
3.9577	0.3883	0.0064029	0.0064029\\
3.9678	0.37766	0.006833	0.006833\\
3.98	0.36702	0.0067149	0.0067149\\
3.9908	0.35638	0.0068135	0.0068135\\
4.0008	0.34574	0.0068598	0.0068598\\
4.0113	0.33511	0.0068246	0.0068246\\
4.0202	0.32447	0.0069444	0.0069444\\
4.0318	0.31383	0.0067454	0.0067454\\
4.0441	0.30319	0.0064174	0.0064174\\
4.0554	0.29255	0.0065338	0.0065338\\
4.0676	0.28191	0.0060628	0.0060628\\
4.0793	0.27128	0.0060461	0.0060461\\
4.0914	0.26064	0.0061657	0.0061657\\
4.1033	0.25	0.0062693	0.0062693\\
4.1167	0.23936	0.0063805	0.0063805\\
4.1321	0.22872	0.0066514	0.0066514\\
4.1494	0.21809	0.0069426	0.0069426\\
4.1669	0.20745	0.0072838	0.0072838\\
4.1856	0.19681	0.0076203	0.0076203\\
4.2028	0.18617	0.0080221	0.0080221\\
4.2197	0.17553	0.0079141	0.0079141\\
4.2417	0.16489	0.0083169	0.0083169\\
4.2609	0.15426	0.0082648	0.0082648\\
4.2818	0.14362	0.0083327	0.0083327\\
4.3051	0.13298	0.0078595	0.0078595\\
4.33	0.12234	0.0080214	0.0080214\\
4.3521	0.1117	0.0077859	0.0077859\\
4.3782	0.10106	0.0082103	0.0082103\\
4.4054	0.090426	0.0084553	0.0084553\\
4.4416	0.079787	0.009553	0.009553\\
4.4781	0.069149	0.0095889	0.0095889\\
4.5216	0.058511	0.0095988	0.0095988\\
4.574	0.047872	0.011021	0.011021\\
4.636	0.037234	0.013906	0.013906\\
4.7242	0.026596	0.014438	0.014438\\
4.8252	0.015957	0.016348	0.016348\\
5.1746	0.0053191	0.014408	0.014408\\
};
\addplot [color=mycolor1]
  table[row sep=crcr]{%
3.6037	0.97822\\
3.6213	0.94662\\
3.639	0.90629\\
3.6743	0.81986\\
4.0625	0.26909\\
4.4507	0.088689\\
4.8389	0.029354\\
5.2271	0.009756\\
5.3506	0.0068775\\
};
\addlegendentry{\ExactParetoFitLegend}


\addplot[area legend, draw=none, fill=mycolor1, fill opacity=0.15, forget plot]
table[row sep=crcr] {%
x	y\\
3.6037	0.98434\\
3.6213	0.95517\\
3.639	0.91592\\
3.6566	0.87175\\
3.6743	0.82823\\
3.6919	0.7869\\
3.7096	0.74765\\
3.7272	0.71037\\
3.7449	0.67498\\
3.7625	0.64136\\
3.7801	0.60945\\
3.7978	0.57914\\
3.8154	0.55036\\
3.8331	0.52304\\
3.8507	0.4971\\
3.8684	0.47247\\
3.886	0.44909\\
3.9037	0.42689\\
3.9213	0.40582\\
3.939	0.38581\\
3.9566	0.36681\\
3.9742	0.34877\\
3.9919	0.33163\\
4.0095	0.31535\\
4.0272	0.29989\\
4.0448	0.28519\\
4.0625	0.27123\\
4.0801	0.25796\\
4.0978	0.24535\\
4.1154	0.23335\\
4.1331	0.22196\\
4.1507	0.21112\\
4.1683	0.20082\\
4.186	0.19102\\
4.2036	0.1817\\
4.2213	0.17285\\
4.2389	0.16442\\
4.2566	0.15641\\
4.2742	0.14879\\
4.2919	0.14155\\
4.3095	0.13466\\
4.3272	0.1281\\
4.3448	0.12187\\
4.3624	0.11594\\
4.3801	0.1103\\
4.3977	0.10494\\
4.4154	0.099837\\
4.433	0.094984\\
4.4507	0.090369\\
4.4683	0.085978\\
4.486	0.081802\\
4.5036	0.077829\\
4.5213	0.07405\\
4.5389	0.070455\\
4.5565	0.067035\\
4.5742	0.063782\\
4.5918	0.060687\\
4.6095	0.057743\\
4.6271	0.054943\\
4.6448	0.052278\\
4.6624	0.049744\\
4.6801	0.047332\\
4.6977	0.045038\\
4.7154	0.042856\\
4.733	0.04078\\
4.7506	0.038804\\
4.7683	0.036925\\
4.7859	0.035137\\
4.8036	0.033436\\
4.8212	0.031817\\
4.8389	0.030277\\
4.8565	0.028812\\
4.8742	0.027418\\
4.8918	0.026092\\
4.9095	0.02483\\
4.9271	0.023629\\
4.9447	0.022487\\
4.9624	0.0214\\
4.98	0.020365\\
4.9977	0.019381\\
5.0153	0.018445\\
5.033	0.017554\\
5.0506	0.016706\\
5.0683	0.015899\\
5.0859	0.015132\\
5.1036	0.014401\\
5.1212	0.013706\\
5.1388	0.013045\\
5.1565	0.012415\\
5.1741	0.011816\\
5.1918	0.011246\\
5.2094	0.010704\\
5.2271	0.010188\\
5.2447	0.0096969\\
5.2624	0.0092295\\
5.28	0.0087847\\
5.2977	0.0083615\\
5.3153	0.0079587\\
5.333	0.0075753\\
5.3506	0.0072105\\
5.3506	0.0065444\\
5.333	0.0068837\\
5.3153	0.0072406\\
5.2977	0.0076161\\
5.28	0.0080112\\
5.2624	0.0084267\\
5.2447	0.0088639\\
5.2271	0.0093239\\
5.2094	0.0098078\\
5.1918	0.010317\\
5.1741	0.010853\\
5.1565	0.011416\\
5.1388	0.012009\\
5.1212	0.012633\\
5.1036	0.013289\\
5.0859	0.013979\\
5.0683	0.014706\\
5.0506	0.01547\\
5.033	0.016274\\
5.0153	0.01712\\
4.9977	0.01801\\
4.98	0.018947\\
4.9624	0.019932\\
4.9447	0.020969\\
4.9271	0.02206\\
4.9095	0.023208\\
4.8918	0.024416\\
4.8742	0.025687\\
4.8565	0.027024\\
4.8389	0.028431\\
4.8212	0.029912\\
4.8036	0.03147\\
4.7859	0.033109\\
4.7683	0.034834\\
4.7506	0.036649\\
4.733	0.038559\\
4.7154	0.040569\\
4.6977	0.042684\\
4.6801	0.044909\\
4.6624	0.047251\\
4.6448	0.049716\\
4.6271	0.052309\\
4.6095	0.055038\\
4.5918	0.05791\\
4.5742	0.060932\\
4.5565	0.064112\\
4.5389	0.067458\\
4.5213	0.07098\\
4.5036	0.074686\\
4.486	0.078586\\
4.4683	0.082691\\
4.4507	0.08701\\
4.433	0.091556\\
4.4154	0.096339\\
4.3977	0.10137\\
4.3801	0.10667\\
4.3624	0.11225\\
4.3448	0.11811\\
4.3272	0.12429\\
4.3095	0.13079\\
4.2919	0.13762\\
4.2742	0.14482\\
4.2566	0.15239\\
4.2389	0.16036\\
4.2213	0.16875\\
4.2036	0.17757\\
4.186	0.18686\\
4.1683	0.19663\\
4.1507	0.20691\\
4.1331	0.21773\\
4.1154	0.22911\\
4.0978	0.24109\\
4.0801	0.25369\\
4.0625	0.26694\\
4.0448	0.28088\\
4.0272	0.29554\\
4.0095	0.31096\\
3.9919	0.32717\\
3.9742	0.34422\\
3.9566	0.36214\\
3.939	0.38097\\
3.9213	0.40077\\
3.9037	0.42158\\
3.886	0.44344\\
3.8684	0.46642\\
3.8507	0.49057\\
3.8331	0.51595\\
3.8154	0.54262\\
3.7978	0.57065\\
3.7801	0.60012\\
3.7625	0.63109\\
3.7449	0.66366\\
3.7272	0.69789\\
3.7096	0.73388\\
3.6919	0.77171\\
3.6743	0.8115\\
3.6566	0.85333\\
3.639	0.89665\\
3.6213	0.93807\\
3.6037	0.97209\\
}--cycle;
\addplot[only marks, mark=*, mark size=1.1180pt, color=black, fill=black, opacity=0.60, draw=none, mark options={draw=none,line width=0pt}] table[row sep=crcr]{%
x	y\\
3.5623	0.99468\\
3.5723	0.98404\\
3.573	0.9734\\
3.5753	0.96277\\
3.5766	0.95213\\
3.5781	0.94149\\
3.5826	0.93085\\
3.5851	0.92021\\
3.5911	0.90957\\
3.5942	0.89894\\
3.596	0.8883\\
3.6016	0.87766\\
3.603	0.86702\\
3.6047	0.85638\\
3.6086	0.84574\\
3.6086	0.83511\\
3.6097	0.82447\\
3.6142	0.81383\\
3.6261	0.80319\\
3.6383	0.79255\\
3.6464	0.78191\\
3.6535	0.77128\\
3.657	0.76064\\
3.6658	0.75\\
3.6678	0.73936\\
3.6687	0.72872\\
3.6702	0.71809\\
3.6723	0.70745\\
3.6793	0.69681\\
3.6867	0.68617\\
3.6949	0.67553\\
3.7015	0.66489\\
3.7031	0.65426\\
3.719	0.64362\\
3.7213	0.63298\\
3.7265	0.62234\\
3.7307	0.6117\\
3.737	0.60106\\
3.7374	0.59043\\
3.7423	0.57979\\
3.7487	0.56915\\
3.7677	0.55851\\
3.7685	0.54787\\
3.7768	0.53723\\
3.8005	0.5266\\
3.8015	0.51596\\
3.8022	0.50532\\
3.8032	0.49468\\
3.8039	0.48404\\
3.8108	0.4734\\
3.8214	0.46277\\
3.8405	0.45213\\
3.8605	0.44149\\
3.863	0.43085\\
3.8645	0.42021\\
3.8662	0.40957\\
3.8806	0.39894\\
3.8878	0.3883\\
3.8964	0.37766\\
3.905	0.36702\\
3.9294	0.35638\\
3.9303	0.34574\\
3.954	0.33511\\
3.9604	0.32447\\
3.9669	0.31383\\
3.9747	0.30319\\
3.9822	0.29255\\
3.9929	0.28191\\
4.0051	0.27128\\
4.0144	0.26064\\
4.0161	0.25\\
4.0433	0.23936\\
4.0578	0.22872\\
4.073	0.21809\\
4.0857	0.20745\\
4.0904	0.19681\\
4.1168	0.18617\\
4.1265	0.17553\\
4.127	0.16489\\
4.143	0.15426\\
4.1756	0.14362\\
4.2156	0.13298\\
4.3116	0.12234\\
4.3276	0.1117\\
4.3662	0.10106\\
4.4791	0.090426\\
4.4912	0.079787\\
4.517	0.069149\\
4.5743	0.058511\\
4.5914	0.047872\\
4.6138	0.037234\\
4.7898	0.026596\\
5.0876	0.015957\\
5.3101	0.0053191\\
};
\addlegendentry{\RealEmpiricCCDFLegend}

\addplot [color=black]
  table[row sep=crcr]{%
3.5608	1\\
5.3506	0.005275\\
};
\addlegendentry{\RealParetoFitLegend}

      % \input{../../../../.tex/20/KS/SizeDistributionFittingsfig35.tex}
      \legend{}
  
      \nextgroupplot[%
        xmin=3.5608,xmax=5.4115,
        % xmin=3.4,xmax=5.4,
        % ymin=0.005121,ymax=1,
      ] % This file was created by matlab2tikz.
%
\definecolor{mycolor1}{rgb}{0.83529,0.36863,0.00000}%
\definecolor{mycolor2}{rgb}{0.00000,0.44700,0.74100}%
\definecolor{mycolor3}{rgb}{0.49400,0.18400,0.55600}%
%
\addplot[only marks, mark=*, mark size=1.1180pt, color=mycolor1, fill=mycolor1, opacity=0.60, draw=none, mark options={draw=none,line width=0pt}] table[row sep=crcr]{%
x	y\\
3.6043	0.99468\\
3.6085	0.98404\\
3.6129	0.9734\\
3.6166	0.96277\\
3.6217	0.95213\\
3.6257	0.94149\\
3.6297	0.93085\\
3.6336	0.92021\\
3.6386	0.90957\\
3.6436	0.89894\\
3.6485	0.8883\\
3.6534	0.87766\\
3.6578	0.86702\\
3.6617	0.85638\\
3.6672	0.84574\\
3.6724	0.83511\\
3.6775	0.82447\\
3.683	0.81383\\
3.6877	0.80319\\
3.6931	0.79255\\
3.6981	0.78191\\
3.7035	0.77128\\
3.7087	0.76064\\
3.7142	0.75\\
3.7195	0.73936\\
3.7257	0.72872\\
3.7321	0.71809\\
3.7384	0.70745\\
3.7448	0.69681\\
3.7505	0.68617\\
3.7557	0.67553\\
3.7608	0.66489\\
3.7667	0.65426\\
3.7732	0.64362\\
3.7794	0.63298\\
3.7861	0.62234\\
3.7926	0.6117\\
3.7986	0.60106\\
3.8051	0.59043\\
3.8111	0.57979\\
3.8186	0.56915\\
3.8258	0.55851\\
3.8344	0.54787\\
3.8405	0.53723\\
3.848	0.5266\\
3.855	0.51596\\
3.8618	0.50532\\
3.8693	0.49468\\
3.8771	0.48404\\
3.8852	0.4734\\
3.8931	0.46277\\
3.9016	0.45213\\
3.9092	0.44149\\
3.9181	0.43085\\
3.9273	0.42021\\
3.9357	0.40957\\
3.9432	0.39894\\
3.9533	0.3883\\
3.9637	0.37766\\
3.9738	0.36702\\
3.9848	0.35638\\
3.9951	0.34574\\
4.0063	0.33511\\
4.0195	0.32447\\
4.0316	0.31383\\
4.0451	0.30319\\
4.0541	0.29255\\
4.0665	0.28191\\
4.0793	0.27128\\
4.0929	0.26064\\
4.1075	0.25\\
4.1207	0.23936\\
4.1351	0.22872\\
4.1514	0.21809\\
4.1699	0.20745\\
4.1822	0.19681\\
4.1983	0.18617\\
4.2172	0.17553\\
4.2352	0.16489\\
4.2564	0.15426\\
4.2767	0.14362\\
4.2966	0.13298\\
4.3207	0.12234\\
4.3486	0.1117\\
4.3743	0.10106\\
4.4051	0.090426\\
4.439	0.079787\\
4.4763	0.069149\\
4.5194	0.058511\\
4.5656	0.047872\\
4.635	0.037234\\
4.7318	0.026596\\
4.8427	0.015957\\
5.1833	0.0053191\\
};
\addlegendentry{\ApproxEmpiricCCDFLegend}

\addplot [color=mycolor1, only marks, every error bar/.append style={opacity=0.30}, mark=*, mark size=0pt, draw=none, forget plot]
 plot [error bars/.cd, x dir=both, x explicit, error bar style={line width=1pt, color=mycolor1}, error mark options={mark=none,mark size=0pt}]
 table[row sep=crcr, x error plus index=2, x error minus index=3]{%
3.6043	0.99468	0.0045572	0.0045572\\
3.6085	0.98404	0.0043712	0.0043712\\
3.6129	0.9734	0.0045168	0.0045168\\
3.6166	0.96277	0.0044731	0.0044731\\
3.6217	0.95213	0.0044265	0.0044265\\
3.6257	0.94149	0.0044362	0.0044362\\
3.6297	0.93085	0.0044388	0.0044388\\
3.6336	0.92021	0.0043707	0.0043707\\
3.6386	0.90957	0.0046291	0.0046291\\
3.6436	0.89894	0.0046805	0.0046805\\
3.6485	0.8883	0.0046062	0.0046062\\
3.6534	0.87766	0.0047192	0.0047192\\
3.6578	0.86702	0.0047457	0.0047457\\
3.6617	0.85638	0.0047213	0.0047213\\
3.6672	0.84574	0.0048522	0.0048522\\
3.6724	0.83511	0.0047178	0.0047178\\
3.6775	0.82447	0.0045714	0.0045714\\
3.683	0.81383	0.0044815	0.0044815\\
3.6877	0.80319	0.0045796	0.0045796\\
3.6931	0.79255	0.0046427	0.0046427\\
3.6981	0.78191	0.0047601	0.0047601\\
3.7035	0.77128	0.0049339	0.0049339\\
3.7087	0.76064	0.0048633	0.0048633\\
3.7142	0.75	0.0048943	0.0048943\\
3.7195	0.73936	0.0047886	0.0047886\\
3.7257	0.72872	0.0047314	0.0047314\\
3.7321	0.71809	0.0048394	0.0048394\\
3.7384	0.70745	0.0048539	0.0048539\\
3.7448	0.69681	0.005076	0.005076\\
3.7505	0.68617	0.0050712	0.0050712\\
3.7557	0.67553	0.0050637	0.0050637\\
3.7608	0.66489	0.0051589	0.0051589\\
3.7667	0.65426	0.0051728	0.0051728\\
3.7732	0.64362	0.0050904	0.0050904\\
3.7794	0.63298	0.0050729	0.0050729\\
3.7861	0.62234	0.0049622	0.0049622\\
3.7926	0.6117	0.0050011	0.0050011\\
3.7986	0.60106	0.0051642	0.0051642\\
3.8051	0.59043	0.0050796	0.0050796\\
3.8111	0.57979	0.0051542	0.0051542\\
3.8186	0.56915	0.0054195	0.0054195\\
3.8258	0.55851	0.0056046	0.0056046\\
3.8344	0.54787	0.0058109	0.0058109\\
3.8405	0.53723	0.0059958	0.0059958\\
3.848	0.5266	0.0057197	0.0057197\\
3.855	0.51596	0.0057667	0.0057667\\
3.8618	0.50532	0.0059162	0.0059162\\
3.8693	0.49468	0.005978	0.005978\\
3.8771	0.48404	0.0059461	0.0059461\\
3.8852	0.4734	0.005987	0.005987\\
3.8931	0.46277	0.0058947	0.0058947\\
3.9016	0.45213	0.0059886	0.0059886\\
3.9092	0.44149	0.0058124	0.0058124\\
3.9181	0.43085	0.0060973	0.0060973\\
3.9273	0.42021	0.0063055	0.0063055\\
3.9357	0.40957	0.0066836	0.0066836\\
3.9432	0.39894	0.0065301	0.0065301\\
3.9533	0.3883	0.006	0.006\\
3.9637	0.37766	0.0060566	0.0060566\\
3.9738	0.36702	0.0063024	0.0063024\\
3.9848	0.35638	0.0065619	0.0065619\\
3.9951	0.34574	0.0066914	0.0066914\\
4.0063	0.33511	0.0065221	0.0065221\\
4.0195	0.32447	0.0062434	0.0062434\\
4.0316	0.31383	0.0062868	0.0062868\\
4.0451	0.30319	0.0062701	0.0062701\\
4.0541	0.29255	0.0063707	0.0063707\\
4.0665	0.28191	0.0063128	0.0063128\\
4.0793	0.27128	0.0063457	0.0063457\\
4.0929	0.26064	0.0064012	0.0064012\\
4.1075	0.25	0.0065842	0.0065842\\
4.1207	0.23936	0.0067354	0.0067354\\
4.1351	0.22872	0.0071575	0.0071575\\
4.1514	0.21809	0.0069122	0.0069122\\
4.1699	0.20745	0.0074422	0.0074422\\
4.1822	0.19681	0.0075591	0.0075591\\
4.1983	0.18617	0.0073044	0.0073044\\
4.2172	0.17553	0.0075554	0.0075554\\
4.2352	0.16489	0.0067471	0.0067471\\
4.2564	0.15426	0.0070369	0.0070369\\
4.2767	0.14362	0.0073888	0.0073888\\
4.2966	0.13298	0.0081069	0.0081069\\
4.3207	0.12234	0.0084621	0.0084621\\
4.3486	0.1117	0.0088645	0.0088645\\
4.3743	0.10106	0.0094713	0.0094713\\
4.4051	0.090426	0.0089463	0.0089463\\
4.439	0.079787	0.0095141	0.0095141\\
4.4763	0.069149	0.0092913	0.0092913\\
4.5194	0.058511	0.010771	0.010771\\
4.5656	0.047872	0.01169	0.01169\\
4.635	0.037234	0.012228	0.012228\\
4.7318	0.026596	0.015596	0.015596\\
4.8427	0.015957	0.017427	0.017427\\
5.1833	0.0053191	0.01767	0.01767\\
};
\addplot [color=mycolor1]
  table[row sep=crcr]{%
3.604	0.97462\\
3.6217	0.94187\\
3.6395	0.89995\\
3.6572	0.85753\\
3.7282	0.69936\\
4.1008	0.23927\\
4.4735	0.082203\\
4.8461	0.028358\\
5.2365	0.0093391\\
5.3607	0.0065647\\
};
\addlegendentry{\ApproxParetoFitLegend}


\addplot[area legend, draw=none, fill=mycolor1, fill opacity=0.15, forget plot]
table[row sep=crcr] {%
x	y\\
3.604	0.98135\\
3.6217	0.95082\\
3.6395	0.90965\\
3.6572	0.86747\\
3.675	0.82471\\
3.6927	0.78342\\
3.7105	0.74396\\
3.7282	0.70651\\
3.746	0.67096\\
3.7637	0.63722\\
3.7814	0.60519\\
3.7992	0.5748\\
3.8169	0.54596\\
3.8347	0.51858\\
3.8524	0.49261\\
3.8702	0.46796\\
3.8879	0.44456\\
3.9057	0.42237\\
3.9234	0.4013\\
3.9411	0.38132\\
3.9589	0.36235\\
3.9766	0.34434\\
3.9944	0.32726\\
4.0121	0.31103\\
4.0299	0.29563\\
4.0476	0.28101\\
4.0654	0.26712\\
4.0831	0.25393\\
4.1008	0.2414\\
4.1186	0.22949\\
4.1363	0.21818\\
4.1541	0.20744\\
4.1718	0.19722\\
4.1896	0.18752\\
4.2073	0.17829\\
4.2251	0.16953\\
4.2428	0.16119\\
4.2605	0.15327\\
4.2783	0.14574\\
4.296	0.13859\\
4.3138	0.13178\\
4.3315	0.12532\\
4.3493	0.11917\\
4.367	0.11332\\
4.3848	0.10776\\
4.4025	0.10248\\
4.4202	0.097458\\
4.438	0.092682\\
4.4557	0.088141\\
4.4735	0.083823\\
4.4912	0.079718\\
4.509	0.075814\\
4.5267	0.072103\\
4.5445	0.068574\\
4.5622	0.065218\\
4.5799	0.062027\\
4.5977	0.058993\\
4.6154	0.056108\\
4.6332	0.053364\\
4.6509	0.050755\\
4.6687	0.048275\\
4.6864	0.045915\\
4.7042	0.043672\\
4.7219	0.041538\\
4.7396	0.03951\\
4.7574	0.03758\\
4.7751	0.035745\\
4.7929	0.034\\
4.8106	0.032341\\
4.8284	0.030763\\
4.8461	0.029262\\
4.8639	0.027834\\
4.8816	0.026477\\
4.8993	0.025186\\
4.9171	0.023958\\
4.9348	0.02279\\
4.9526	0.021679\\
4.9703	0.020623\\
4.9881	0.019618\\
5.0058	0.018662\\
5.0236	0.017753\\
5.0413	0.016889\\
5.0591	0.016067\\
5.0768	0.015285\\
5.0945	0.014541\\
5.1123	0.013833\\
5.13	0.01316\\
5.1478	0.01252\\
5.1655	0.011911\\
5.1833	0.011332\\
5.201	0.010781\\
5.2188	0.010257\\
5.2365	0.0097583\\
5.2542	0.0092841\\
5.272	0.0088331\\
5.2897	0.008404\\
5.3075	0.0079958\\
5.3252	0.0076076\\
5.343	0.0072382\\
5.3607	0.0068869\\
5.3607	0.0062426\\
5.343	0.006569\\
5.3252	0.0069124\\
5.3075	0.0072739\\
5.2897	0.0076544\\
5.272	0.0080548\\
5.2542	0.0084763\\
5.2365	0.0089198\\
5.2188	0.0093867\\
5.201	0.009878\\
5.1833	0.010395\\
5.1655	0.01094\\
5.1478	0.011512\\
5.13	0.012116\\
5.1123	0.01275\\
5.0945	0.013418\\
5.0768	0.014122\\
5.0591	0.014862\\
5.0413	0.015641\\
5.0236	0.016461\\
5.0058	0.017324\\
4.9881	0.018233\\
4.9703	0.019189\\
4.9526	0.020196\\
4.9348	0.021256\\
4.9171	0.022372\\
4.8993	0.023546\\
4.8816	0.024782\\
4.8639	0.026084\\
4.8461	0.027454\\
4.8284	0.028896\\
4.8106	0.030414\\
4.7929	0.032012\\
4.7751	0.033694\\
4.7574	0.035465\\
4.7396	0.03733\\
4.7219	0.039293\\
4.7042	0.041359\\
4.6864	0.043534\\
4.6687	0.045825\\
4.6509	0.048236\\
4.6332	0.050774\\
4.6154	0.053446\\
4.5977	0.05626\\
4.5799	0.059222\\
4.5622	0.06234\\
4.5445	0.065623\\
4.5267	0.069079\\
4.509	0.072718\\
4.4912	0.07655\\
4.4735	0.080583\\
4.4557	0.08483\\
4.438	0.089302\\
4.4202	0.094009\\
4.4025	0.098966\\
4.3848	0.10418\\
4.367	0.10968\\
4.3493	0.11546\\
4.3315	0.12155\\
4.3138	0.12796\\
4.296	0.13472\\
4.2783	0.14182\\
4.2605	0.1493\\
4.2428	0.15718\\
4.2251	0.16548\\
4.2073	0.17421\\
4.1896	0.1834\\
4.1718	0.19307\\
4.1541	0.20326\\
4.1363	0.21398\\
4.1186	0.22526\\
4.1008	0.23714\\
4.0831	0.24963\\
4.0654	0.26278\\
4.0476	0.27662\\
4.0299	0.29118\\
4.0121	0.30649\\
3.9944	0.3226\\
3.9766	0.33954\\
3.9589	0.35735\\
3.9411	0.37608\\
3.9234	0.39577\\
3.9057	0.41648\\
3.8879	0.43825\\
3.8702	0.46114\\
3.8524	0.4852\\
3.8347	0.5105\\
3.8169	0.53711\\
3.7992	0.56508\\
3.7814	0.5945\\
3.7637	0.62545\\
3.746	0.65799\\
3.7282	0.69221\\
3.7105	0.72821\\
3.6927	0.76608\\
3.675	0.80588\\
3.6572	0.84758\\
3.6395	0.89026\\
3.6217	0.93292\\
3.604	0.9679\\
}--cycle;
\addplot[only marks, mark=*, mark size=1.1180pt, color=black, fill=black, opacity=0.60, draw=none, mark options={draw=none,line width=0pt}] table[row sep=crcr]{%
x	y\\
3.5623	0.99468\\
3.5723	0.98404\\
3.573	0.9734\\
3.5753	0.96277\\
3.5766	0.95213\\
3.5781	0.94149\\
3.5826	0.93085\\
3.5851	0.92021\\
3.5911	0.90957\\
3.5942	0.89894\\
3.596	0.8883\\
3.6016	0.87766\\
3.603	0.86702\\
3.6047	0.85638\\
3.6086	0.84574\\
3.6086	0.83511\\
3.6097	0.82447\\
3.6142	0.81383\\
3.6261	0.80319\\
3.6383	0.79255\\
3.6464	0.78191\\
3.6535	0.77128\\
3.657	0.76064\\
3.6658	0.75\\
3.6678	0.73936\\
3.6687	0.72872\\
3.6702	0.71809\\
3.6723	0.70745\\
3.6793	0.69681\\
3.6867	0.68617\\
3.6949	0.67553\\
3.7015	0.66489\\
3.7031	0.65426\\
3.719	0.64362\\
3.7213	0.63298\\
3.7265	0.62234\\
3.7307	0.6117\\
3.737	0.60106\\
3.7374	0.59043\\
3.7423	0.57979\\
3.7487	0.56915\\
3.7677	0.55851\\
3.7685	0.54787\\
3.7768	0.53723\\
3.8005	0.5266\\
3.8015	0.51596\\
3.8022	0.50532\\
3.8032	0.49468\\
3.8039	0.48404\\
3.8108	0.4734\\
3.8214	0.46277\\
3.8405	0.45213\\
3.8605	0.44149\\
3.863	0.43085\\
3.8645	0.42021\\
3.8662	0.40957\\
3.8806	0.39894\\
3.8878	0.3883\\
3.8964	0.37766\\
3.905	0.36702\\
3.9294	0.35638\\
3.9303	0.34574\\
3.954	0.33511\\
3.9604	0.32447\\
3.9669	0.31383\\
3.9747	0.30319\\
3.9822	0.29255\\
3.9929	0.28191\\
4.0051	0.27128\\
4.0144	0.26064\\
4.0161	0.25\\
4.0433	0.23936\\
4.0578	0.22872\\
4.073	0.21809\\
4.0857	0.20745\\
4.0904	0.19681\\
4.1168	0.18617\\
4.1265	0.17553\\
4.127	0.16489\\
4.143	0.15426\\
4.1756	0.14362\\
4.2156	0.13298\\
4.3116	0.12234\\
4.3276	0.1117\\
4.3662	0.10106\\
4.4791	0.090426\\
4.4912	0.079787\\
4.517	0.069149\\
4.5743	0.058511\\
4.5914	0.047872\\
4.6138	0.037234\\
4.7898	0.026596\\
5.0876	0.015957\\
5.3101	0.0053191\\
};
\addlegendentry{\RealEmpiricCCDFLegend}

\addplot [color=black]
  table[row sep=crcr]{%
3.5608	1\\
5.3607	0.005121\\
};
\addlegendentry{\RealParetoFitLegend}

      % \input{../../../../.tex/20/KS/SizeDistributionFittingsfig36.tex}
      \legend{}
    \end{groupplot}
  
    \begin{groupplot}[
      group style={
        group name=Row5,
        plotLastGroup,
      },
      plotRow,
      row2KeysWithLabels,
      xlabel3,
      plotLegendNE,
      %
      % xmin=3637.5,xmax=2.5003e+05,
      ymin=0.0044129,ymax=1,
    ]
      \nextgroupplot[%
        at={($(Row4 c1r1.south west)-(0,\yPlotSep em)-(0,\plotHeight cm)$)},
        xmin=3.598,xmax=5.4115,
        % xmin=3.5,xmax=5.5,
        % ymin=0.0053191,ymax=1,
      ] % This file was created by matlab2tikz.
%
\definecolor{mycolor1}{rgb}{0.00000,0.44706,0.69804}%
\definecolor{mycolor2}{rgb}{0.00000,0.44700,0.74100}%
\definecolor{mycolor3}{rgb}{0.83529,0.36863,0.00000}%
\definecolor{mycolor4}{rgb}{0.49400,0.18400,0.55600}%
%
\addplot[only marks, mark=*, mark size=1.1180pt, color=mycolor1, fill=mycolor1, opacity=0.60, draw=none, mark options={draw=none,line width=0pt}] table[row sep=crcr]{%
x	y\\
3.5984	0.99468\\
3.6027	0.98404\\
3.6068	0.9734\\
3.6112	0.96277\\
3.6162	0.95213\\
3.6211	0.94149\\
3.6266	0.93085\\
3.6309	0.92021\\
3.636	0.90957\\
3.64	0.89894\\
3.6441	0.8883\\
3.6477	0.87766\\
3.6529	0.86702\\
3.6581	0.85638\\
3.6641	0.84574\\
3.6687	0.83511\\
3.674	0.82447\\
3.6794	0.81383\\
3.6854	0.80319\\
3.6908	0.79255\\
3.6952	0.78191\\
3.7011	0.77128\\
3.707	0.76064\\
3.713	0.75\\
3.7182	0.73936\\
3.7241	0.72872\\
3.73	0.71809\\
3.736	0.70745\\
3.7427	0.69681\\
3.7494	0.68617\\
3.7566	0.67553\\
3.7628	0.66489\\
3.7689	0.65426\\
3.7752	0.64362\\
3.7809	0.63298\\
3.7869	0.62234\\
3.7932	0.6117\\
3.8004	0.60106\\
3.8084	0.59043\\
3.8147	0.57979\\
3.8215	0.56915\\
3.8298	0.55851\\
3.8363	0.54787\\
3.8438	0.53723\\
3.8521	0.5266\\
3.8604	0.51596\\
3.8683	0.50532\\
3.8765	0.49468\\
3.8843	0.48404\\
3.8928	0.4734\\
3.9007	0.46277\\
3.9098	0.45213\\
3.9208	0.44149\\
3.9308	0.43085\\
3.9405	0.42021\\
3.9501	0.40957\\
3.9606	0.39894\\
3.9709	0.3883\\
3.9801	0.37766\\
3.9899	0.36702\\
3.999	0.35638\\
4.0096	0.34574\\
4.0188	0.33511\\
4.0298	0.32447\\
4.0417	0.31383\\
4.0542	0.30319\\
4.0666	0.29255\\
4.0789	0.28191\\
4.0918	0.27128\\
4.1041	0.26064\\
4.1175	0.25\\
4.131	0.23936\\
4.1454	0.22872\\
4.159	0.21809\\
4.173	0.20745\\
4.1912	0.19681\\
4.2086	0.18617\\
4.2269	0.17553\\
4.2454	0.16489\\
4.2692	0.15426\\
4.2908	0.14362\\
4.3146	0.13298\\
4.3388	0.12234\\
4.362	0.1117\\
4.3887	0.10106\\
4.4173	0.090426\\
4.4494	0.079787\\
4.4812	0.069149\\
4.5257	0.058511\\
4.5784	0.047872\\
4.6417	0.037234\\
4.741	0.026596\\
4.8583	0.015957\\
5.2036	0.0053191\\
};
\addlegendentry{\ExactEmpiricCCDFLegend}

\addplot [color=mycolor1, only marks, every error bar/.append style={opacity=0.30}, mark=*, mark size=0pt, draw=none, forget plot]
 plot [error bars/.cd, x dir=both, x explicit, error bar style={line width=1pt, color=mycolor1}, error mark options={mark=none,mark size=0pt}]
 table[row sep=crcr, x error plus index=2, x error minus index=3]{%
3.5984	0.99468	0.00473	0.00473\\
3.6027	0.98404	0.0047925	0.0047925\\
3.6068	0.9734	0.0048115	0.0048115\\
3.6112	0.96277	0.0047892	0.0047892\\
3.6162	0.95213	0.0047066	0.0047066\\
3.6211	0.94149	0.0049168	0.0049168\\
3.6266	0.93085	0.0048041	0.0048041\\
3.6309	0.92021	0.004932	0.004932\\
3.636	0.90957	0.0049959	0.0049959\\
3.64	0.89894	0.005017	0.005017\\
3.6441	0.8883	0.0051081	0.0051081\\
3.6477	0.87766	0.005056	0.005056\\
3.6529	0.86702	0.0050923	0.0050923\\
3.6581	0.85638	0.0052229	0.0052229\\
3.6641	0.84574	0.005148	0.005148\\
3.6687	0.83511	0.0052242	0.0052242\\
3.674	0.82447	0.0051186	0.0051186\\
3.6794	0.81383	0.005215	0.005215\\
3.6854	0.80319	0.0051265	0.0051265\\
3.6908	0.79255	0.0052431	0.0052431\\
3.6952	0.78191	0.0051884	0.0051884\\
3.7011	0.77128	0.0051923	0.0051923\\
3.707	0.76064	0.0050709	0.0050709\\
3.713	0.75	0.005022	0.005022\\
3.7182	0.73936	0.0050527	0.0050527\\
3.7241	0.72872	0.0052125	0.0052125\\
3.73	0.71809	0.0052284	0.0052284\\
3.736	0.70745	0.0051795	0.0051795\\
3.7427	0.69681	0.0052847	0.0052847\\
3.7494	0.68617	0.0054013	0.0054013\\
3.7566	0.67553	0.005546	0.005546\\
3.7628	0.66489	0.0056581	0.0056581\\
3.7689	0.65426	0.0055104	0.0055104\\
3.7752	0.64362	0.0056133	0.0056133\\
3.7809	0.63298	0.0055232	0.0055232\\
3.7869	0.62234	0.0054689	0.0054689\\
3.7932	0.6117	0.0054553	0.0054553\\
3.8004	0.60106	0.0057052	0.0057052\\
3.8084	0.59043	0.0058848	0.0058848\\
3.8147	0.57979	0.0057587	0.0057587\\
3.8215	0.56915	0.0059265	0.0059265\\
3.8298	0.55851	0.0061686	0.0061686\\
3.8363	0.54787	0.0061957	0.0061957\\
3.8438	0.53723	0.0060905	0.0060905\\
3.8521	0.5266	0.0059675	0.0059675\\
3.8604	0.51596	0.0058526	0.0058526\\
3.8683	0.50532	0.0057498	0.0057498\\
3.8765	0.49468	0.0057298	0.0057298\\
3.8843	0.48404	0.0056797	0.0056797\\
3.8928	0.4734	0.0056628	0.0056628\\
3.9007	0.46277	0.0055566	0.0055566\\
3.9098	0.45213	0.0053166	0.0053166\\
3.9208	0.44149	0.0054396	0.0054396\\
3.9308	0.43085	0.005428	0.005428\\
3.9405	0.42021	0.0055337	0.0055337\\
3.9501	0.40957	0.0053756	0.0053756\\
3.9606	0.39894	0.0054857	0.0054857\\
3.9709	0.3883	0.0055337	0.0055337\\
3.9801	0.37766	0.0056598	0.0056598\\
3.9899	0.36702	0.0058367	0.0058367\\
3.999	0.35638	0.0056606	0.0056606\\
4.0096	0.34574	0.0056459	0.0056459\\
4.0188	0.33511	0.0059698	0.0059698\\
4.0298	0.32447	0.0059241	0.0059241\\
4.0417	0.31383	0.0056343	0.0056343\\
4.0542	0.30319	0.0055073	0.0055073\\
4.0666	0.29255	0.0056644	0.0056644\\
4.0789	0.28191	0.005727	0.005727\\
4.0918	0.27128	0.0057123	0.0057123\\
4.1041	0.26064	0.0056946	0.0056946\\
4.1175	0.25	0.0059084	0.0059084\\
4.131	0.23936	0.0061124	0.0061124\\
4.1454	0.22872	0.0066665	0.0066665\\
4.159	0.21809	0.0068485	0.0068485\\
4.173	0.20745	0.0069575	0.0069575\\
4.1912	0.19681	0.0072717	0.0072717\\
4.2086	0.18617	0.0078688	0.0078688\\
4.2269	0.17553	0.0077805	0.0077805\\
4.2454	0.16489	0.0078623	0.0078623\\
4.2692	0.15426	0.0083156	0.0083156\\
4.2908	0.14362	0.0089906	0.0089906\\
4.3146	0.13298	0.009299	0.009299\\
4.3388	0.12234	0.0084958	0.0084958\\
4.362	0.1117	0.0078765	0.0078765\\
4.3887	0.10106	0.008669	0.008669\\
4.4173	0.090426	0.0097937	0.0097937\\
4.4494	0.079787	0.01038	0.01038\\
4.4812	0.069149	0.011015	0.011015\\
4.5257	0.058511	0.012142	0.012142\\
4.5784	0.047872	0.011936	0.011936\\
4.6417	0.037234	0.014498	0.014498\\
4.741	0.026596	0.016486	0.016486\\
4.8583	0.015957	0.020509	0.020509\\
5.2036	0.0053191	0.018099	0.018099\\
};
\addplot [color=mycolor1]
  table[row sep=crcr]{%
3.598	0.97291\\
3.6163	0.94096\\
3.6346	0.90043\\
3.6713	0.81511\\
4.0559	0.27891\\
4.4406	0.09583\\
4.8436	0.031423\\
5.2649	0.0098364\\
5.4115	0.0065737\\
};
\addlegendentry{\ExactParetoFitLegend}


\addplot[area legend, draw=none, fill=mycolor1, fill opacity=0.15, forget plot]
table[row sep=crcr] {%
x	y\\
3.598	0.97938\\
3.6163	0.9501\\
3.6346	0.91064\\
3.6529	0.86752\\
3.6713	0.82464\\
3.6896	0.78311\\
3.7079	0.74369\\
3.7262	0.70627\\
3.7445	0.67075\\
3.7628	0.63704\\
3.7812	0.60505\\
3.7995	0.57468\\
3.8178	0.54587\\
3.8361	0.51852\\
3.8544	0.49256\\
3.8728	0.46793\\
3.8911	0.44456\\
3.9094	0.42238\\
3.9277	0.40133\\
3.946	0.38135\\
3.9643	0.36239\\
3.9827	0.3444\\
4.001	0.32732\\
4.0193	0.3111\\
4.0376	0.29571\\
4.0559	0.28109\\
4.0743	0.2672\\
4.0926	0.25401\\
4.1109	0.24148\\
4.1292	0.22958\\
4.1475	0.21826\\
4.1658	0.20752\\
4.1842	0.1973\\
4.2025	0.1876\\
4.2208	0.17837\\
4.2391	0.1696\\
4.2574	0.16126\\
4.2758	0.15334\\
4.2941	0.14581\\
4.3124	0.13865\\
4.3307	0.13184\\
4.349	0.12537\\
4.3673	0.11922\\
4.3857	0.11337\\
4.404	0.10781\\
4.4223	0.10253\\
4.4406	0.0975\\
4.4589	0.092721\\
4.4773	0.088177\\
4.4956	0.083857\\
4.5139	0.079749\\
4.5322	0.075843\\
4.5505	0.072129\\
4.5688	0.068598\\
4.5872	0.06524\\
4.6055	0.062047\\
4.6238	0.05901\\
4.6421	0.056123\\
4.6604	0.053377\\
4.6788	0.050767\\
4.6971	0.048284\\
4.7154	0.045923\\
4.7337	0.043678\\
4.752	0.041543\\
4.7703	0.039512\\
4.7887	0.037582\\
4.807	0.035745\\
4.8253	0.033999\\
4.8436	0.032339\\
4.8619	0.030759\\
4.8803	0.029257\\
4.8986	0.027829\\
4.9169	0.02647\\
4.9352	0.025178\\
4.9535	0.02395\\
4.9719	0.022781\\
4.9902	0.02167\\
5.0085	0.020613\\
5.0268	0.019607\\
5.0451	0.018651\\
5.0634	0.017742\\
5.0818	0.016877\\
5.1001	0.016054\\
5.1184	0.015272\\
5.1367	0.014528\\
5.155	0.01382\\
5.1734	0.013146\\
5.1917	0.012506\\
5.21	0.011897\\
5.2283	0.011318\\
5.2466	0.010767\\
5.2649	0.010242\\
5.2833	0.0097439\\
5.3016	0.0092697\\
5.3199	0.0088186\\
5.3382	0.0083895\\
5.3565	0.0079814\\
5.3749	0.0075931\\
5.3932	0.0072238\\
5.4115	0.0068726\\
5.4115	0.0062749\\
5.3932	0.0066025\\
5.3749	0.0069473\\
5.3565	0.0073102\\
5.3382	0.007692\\
5.3199	0.0080939\\
5.3016	0.0085168\\
5.2833	0.0089619\\
5.2649	0.0094303\\
5.2466	0.0099233\\
5.2283	0.010442\\
5.21	0.010988\\
5.1917	0.011563\\
5.1734	0.012168\\
5.155	0.012804\\
5.1367	0.013474\\
5.1184	0.01418\\
5.1001	0.014922\\
5.0818	0.015703\\
5.0634	0.016525\\
5.0451	0.017391\\
5.0268	0.018302\\
5.0085	0.01926\\
4.9902	0.02027\\
4.9719	0.021332\\
4.9535	0.02245\\
4.9352	0.023627\\
4.9169	0.024865\\
4.8986	0.026169\\
4.8803	0.027542\\
4.8619	0.028987\\
4.8436	0.030507\\
4.8253	0.032108\\
4.807	0.033793\\
4.7887	0.035567\\
4.7703	0.037435\\
4.752	0.0394\\
4.7337	0.041469\\
4.7154	0.043648\\
4.6971	0.045941\\
4.6788	0.048355\\
4.6604	0.050896\\
4.6421	0.053571\\
4.6238	0.056387\\
4.6055	0.059352\\
4.5872	0.062473\\
4.5688	0.065758\\
4.5505	0.069217\\
4.5322	0.072859\\
4.5139	0.076693\\
4.4956	0.080729\\
4.4773	0.084978\\
4.4589	0.089451\\
4.4406	0.094161\\
4.4223	0.099119\\
4.404	0.10434\\
4.3857	0.10983\\
4.3673	0.11562\\
4.349	0.12171\\
4.3307	0.12812\\
4.3124	0.13487\\
4.2941	0.14198\\
4.2758	0.14946\\
4.2574	0.15734\\
4.2391	0.16563\\
4.2208	0.17436\\
4.2025	0.18355\\
4.1842	0.19322\\
4.1658	0.2034\\
4.1475	0.21412\\
4.1292	0.2254\\
4.1109	0.23727\\
4.0926	0.24976\\
4.0743	0.26291\\
4.0559	0.27674\\
4.0376	0.29129\\
4.0193	0.30659\\
4.001	0.32269\\
3.9827	0.33962\\
3.9643	0.35742\\
3.946	0.37614\\
3.9277	0.39582\\
3.9094	0.41652\\
3.8911	0.43827\\
3.8728	0.46115\\
3.8544	0.48519\\
3.8361	0.51048\\
3.8178	0.53706\\
3.7995	0.56502\\
3.7812	0.59441\\
3.7628	0.62532\\
3.7445	0.65783\\
3.7262	0.69202\\
3.7079	0.72798\\
3.6896	0.76581\\
3.6713	0.80559\\
3.6529	0.84724\\
3.6346	0.89022\\
3.6163	0.93181\\
3.598	0.96644\\
}--cycle;
\addplot[only marks, mark=*, mark size=1.1180pt, color=mycolor3, fill=mycolor3, opacity=0.60, draw=none, mark options={draw=none,line width=0pt}] table[row sep=crcr]{%
x	y\\
3.6	0.99468\\
3.6045	0.98404\\
3.6089	0.9734\\
3.6141	0.96277\\
3.6178	0.95213\\
3.6225	0.94149\\
3.6271	0.93085\\
3.6318	0.92021\\
3.6365	0.90957\\
3.6418	0.89894\\
3.6472	0.8883\\
3.6515	0.87766\\
3.657	0.86702\\
3.6615	0.85638\\
3.6662	0.84574\\
3.6709	0.83511\\
3.6761	0.82447\\
3.6804	0.81383\\
3.6852	0.80319\\
3.6904	0.79255\\
3.6957	0.78191\\
3.701	0.77128\\
3.7064	0.76064\\
3.7121	0.75\\
3.7185	0.73936\\
3.7229	0.72872\\
3.7274	0.71809\\
3.7333	0.70745\\
3.7389	0.69681\\
3.745	0.68617\\
3.7505	0.67553\\
3.7574	0.66489\\
3.7644	0.65426\\
3.7707	0.64362\\
3.7767	0.63298\\
3.7836	0.62234\\
3.7905	0.6117\\
3.7976	0.60106\\
3.8048	0.59043\\
3.812	0.57979\\
3.8198	0.56915\\
3.8269	0.55851\\
3.8352	0.54787\\
3.8439	0.53723\\
3.8506	0.5266\\
3.8568	0.51596\\
3.8635	0.50532\\
3.8723	0.49468\\
3.8796	0.48404\\
3.8878	0.4734\\
3.897	0.46277\\
3.9047	0.45213\\
3.9134	0.44149\\
3.9222	0.43085\\
3.9309	0.42021\\
3.9403	0.40957\\
3.9517	0.39894\\
3.9628	0.3883\\
3.9726	0.37766\\
3.9842	0.36702\\
3.9942	0.35638\\
4.0054	0.34574\\
4.0156	0.33511\\
4.0277	0.32447\\
4.0395	0.31383\\
4.0516	0.30319\\
4.066	0.29255\\
4.0769	0.28191\\
4.0914	0.27128\\
4.1049	0.26064\\
4.1181	0.25\\
4.1339	0.23936\\
4.1473	0.22872\\
4.1627	0.21809\\
4.1783	0.20745\\
4.1955	0.19681\\
4.2126	0.18617\\
4.231	0.17553\\
4.254	0.16489\\
4.2752	0.15426\\
4.2969	0.14362\\
4.3184	0.13298\\
4.3399	0.12234\\
4.364	0.1117\\
4.39	0.10106\\
4.4189	0.090426\\
4.4573	0.079787\\
4.4912	0.069149\\
4.5391	0.058511\\
4.5879	0.047872\\
4.6566	0.037234\\
4.7343	0.026596\\
4.8688	0.015957\\
5.1869	0.0053191\\
};
\addlegendentry{\ApproxEmpiricCCDFLegend}

\addplot [color=mycolor3, only marks, every error bar/.append style={opacity=0.30}, mark=*, mark size=0pt, draw=none, forget plot]
 plot [error bars/.cd, x dir=both, x explicit, error bar style={line width=1pt, color=mycolor3}, error mark options={mark=none,mark size=0pt}]
 table[row sep=crcr, x error plus index=2, x error minus index=3]{%
3.6	0.99468	0.0045554	0.0045554\\
3.6045	0.98404	0.0045703	0.0045703\\
3.6089	0.9734	0.0045491	0.0045491\\
3.6141	0.96277	0.0044812	0.0044812\\
3.6178	0.95213	0.0044817	0.0044817\\
3.6225	0.94149	0.0044624	0.0044624\\
3.6271	0.93085	0.0045762	0.0045762\\
3.6318	0.92021	0.0045625	0.0045625\\
3.6365	0.90957	0.004468	0.004468\\
3.6418	0.89894	0.0046379	0.0046379\\
3.6472	0.8883	0.0045394	0.0045394\\
3.6515	0.87766	0.0044728	0.0044728\\
3.657	0.86702	0.0045858	0.0045858\\
3.6615	0.85638	0.0045704	0.0045704\\
3.6662	0.84574	0.0044305	0.0044305\\
3.6709	0.83511	0.0043954	0.0043954\\
3.6761	0.82447	0.0045809	0.0045809\\
3.6804	0.81383	0.0045127	0.0045127\\
3.6852	0.80319	0.0045389	0.0045389\\
3.6904	0.79255	0.0046013	0.0046013\\
3.6957	0.78191	0.0046842	0.0046842\\
3.701	0.77128	0.0046214	0.0046214\\
3.7064	0.76064	0.0046093	0.0046093\\
3.7121	0.75	0.0045663	0.0045663\\
3.7185	0.73936	0.0045568	0.0045568\\
3.7229	0.72872	0.0046537	0.0046537\\
3.7274	0.71809	0.0046666	0.0046666\\
3.7333	0.70745	0.0046756	0.0046756\\
3.7389	0.69681	0.0047365	0.0047365\\
3.745	0.68617	0.0047282	0.0047282\\
3.7505	0.67553	0.0046776	0.0046776\\
3.7574	0.66489	0.0046232	0.0046232\\
3.7644	0.65426	0.0049389	0.0049389\\
3.7707	0.64362	0.0049657	0.0049657\\
3.7767	0.63298	0.0051864	0.0051864\\
3.7836	0.62234	0.0051426	0.0051426\\
3.7905	0.6117	0.0054154	0.0054154\\
3.7976	0.60106	0.0052834	0.0052834\\
3.8048	0.59043	0.0053245	0.0053245\\
3.812	0.57979	0.0053216	0.0053216\\
3.8198	0.56915	0.0053218	0.0053218\\
3.8269	0.55851	0.0052938	0.0052938\\
3.8352	0.54787	0.0052741	0.0052741\\
3.8439	0.53723	0.0055718	0.0055718\\
3.8506	0.5266	0.005749	0.005749\\
3.8568	0.51596	0.0058641	0.0058641\\
3.8635	0.50532	0.0058242	0.0058242\\
3.8723	0.49468	0.0060846	0.0060846\\
3.8796	0.48404	0.0058831	0.0058831\\
3.8878	0.4734	0.0060552	0.0060552\\
3.897	0.46277	0.0058784	0.0058784\\
3.9047	0.45213	0.0058689	0.0058689\\
3.9134	0.44149	0.0059804	0.0059804\\
3.9222	0.43085	0.0059113	0.0059113\\
3.9309	0.42021	0.0061887	0.0061887\\
3.9403	0.40957	0.0064169	0.0064169\\
3.9517	0.39894	0.0061866	0.0061866\\
3.9628	0.3883	0.0059948	0.0059948\\
3.9726	0.37766	0.006044	0.006044\\
3.9842	0.36702	0.0061542	0.0061542\\
3.9942	0.35638	0.0061402	0.0061402\\
4.0054	0.34574	0.006197	0.006197\\
4.0156	0.33511	0.0062917	0.0062917\\
4.0277	0.32447	0.0062543	0.0062543\\
4.0395	0.31383	0.006648	0.006648\\
4.0516	0.30319	0.0062708	0.0062708\\
4.066	0.29255	0.0064445	0.0064445\\
4.0769	0.28191	0.0065803	0.0065803\\
4.0914	0.27128	0.0068367	0.0068367\\
4.1049	0.26064	0.0067182	0.0067182\\
4.1181	0.25	0.0066727	0.0066727\\
4.1339	0.23936	0.006414	0.006414\\
4.1473	0.22872	0.0068788	0.0068788\\
4.1627	0.21809	0.0073314	0.0073314\\
4.1783	0.20745	0.0073394	0.0073394\\
4.1955	0.19681	0.0076429	0.0076429\\
4.2126	0.18617	0.007675	0.007675\\
4.231	0.17553	0.0079381	0.0079381\\
4.254	0.16489	0.0079159	0.0079159\\
4.2752	0.15426	0.0082002	0.0082002\\
4.2969	0.14362	0.0088204	0.0088204\\
4.3184	0.13298	0.00845	0.00845\\
4.3399	0.12234	0.0086796	0.0086796\\
4.364	0.1117	0.0086382	0.0086382\\
4.39	0.10106	0.0099918	0.0099918\\
4.4189	0.090426	0.010413	0.010413\\
4.4573	0.079787	0.011005	0.011005\\
4.4912	0.069149	0.011629	0.011629\\
4.5391	0.058511	0.011499	0.011499\\
4.5879	0.047872	0.012519	0.012519\\
4.6566	0.037234	0.015101	0.015101\\
4.7343	0.026596	0.016441	0.016441\\
4.8688	0.015957	0.021019	0.021019\\
5.1869	0.0053191	0.018612	0.018612\\
};
\addplot [color=mycolor3]
  table[row sep=crcr]{%
3.6008	0.97204\\
3.6191	0.94016\\
3.6374	0.89927\\
3.6922	0.77128\\
4.0763	0.26255\\
4.4604	0.089744\\
4.8445	0.0308\\
5.2286	0.010613\\
5.4115	0.0063989\\
};
\addlegendentry{\ApproxParetoFitLegend}


\addplot[area legend, draw=none, fill=mycolor3, fill opacity=0.15, forget plot]
table[row sep=crcr] {%
x	y\\
3.6008	0.97886\\
3.6191	0.94955\\
3.6374	0.90957\\
3.6557	0.86444\\
3.6739	0.82071\\
3.6922	0.77922\\
3.7105	0.73984\\
3.7288	0.70247\\
3.7471	0.667\\
3.7654	0.63335\\
3.7837	0.60141\\
3.802	0.57111\\
3.8203	0.54236\\
3.8386	0.51508\\
3.8568	0.4892\\
3.8751	0.46464\\
3.8934	0.44135\\
3.9117	0.41925\\
3.93	0.39828\\
3.9483	0.37838\\
3.9666	0.3595\\
3.9849	0.34158\\
4.0032	0.32458\\
4.0215	0.30843\\
4.0397	0.29311\\
4.058	0.27856\\
4.0763	0.26474\\
4.0946	0.25162\\
4.1129	0.23915\\
4.1312	0.22731\\
4.1495	0.21607\\
4.1678	0.20538\\
4.1861	0.19523\\
4.2044	0.18558\\
4.2226	0.17641\\
4.2409	0.1677\\
4.2592	0.15942\\
4.2775	0.15156\\
4.2958	0.14408\\
4.3141	0.13698\\
4.3324	0.13022\\
4.3507	0.1238\\
4.369	0.1177\\
4.3873	0.1119\\
4.4055	0.10639\\
4.4238	0.10115\\
4.4421	0.096173\\
4.4604	0.091439\\
4.4787	0.08694\\
4.497	0.082662\\
4.5153	0.078596\\
4.5336	0.074731\\
4.5519	0.071056\\
4.5702	0.067563\\
4.5884	0.064242\\
4.6067	0.061085\\
4.625	0.058084\\
4.6433	0.055231\\
4.6616	0.052518\\
4.6799	0.049939\\
4.6982	0.047488\\
4.7165	0.045157\\
4.7348	0.042941\\
4.7531	0.040834\\
4.7713	0.03883\\
4.7896	0.036926\\
4.8079	0.035115\\
4.8262	0.033393\\
4.8445	0.031756\\
4.8628	0.030199\\
4.8811	0.02872\\
4.8994	0.027312\\
4.9177	0.025974\\
4.936	0.024702\\
4.9542	0.023492\\
4.9725	0.022342\\
4.9908	0.021248\\
5.0091	0.020208\\
5.0274	0.019219\\
5.0457	0.018279\\
5.064	0.017385\\
5.0823	0.016534\\
5.1006	0.015726\\
5.1189	0.014957\\
5.1371	0.014226\\
5.1554	0.01353\\
5.1737	0.012869\\
5.192	0.01224\\
5.2103	0.011642\\
5.2286	0.011074\\
5.2469	0.010533\\
5.2652	0.010019\\
5.2835	0.0095295\\
5.3018	0.0090644\\
5.32	0.008622\\
5.3383	0.0082013\\
5.3566	0.0078012\\
5.3749	0.0074206\\
5.3932	0.0070587\\
5.4115	0.0067145\\
5.4115	0.0060833\\
5.3932	0.0064028\\
5.3749	0.006739\\
5.3566	0.007093\\
5.3383	0.0074657\\
5.32	0.007858\\
5.3018	0.008271\\
5.2835	0.0087057\\
5.2652	0.0091634\\
5.2469	0.0096453\\
5.2286	0.010153\\
5.2103	0.010687\\
5.192	0.011249\\
5.1737	0.011841\\
5.1554	0.012464\\
5.1371	0.01312\\
5.1189	0.01381\\
5.1006	0.014538\\
5.0823	0.015303\\
5.064	0.016109\\
5.0457	0.016958\\
5.0274	0.017851\\
5.0091	0.018792\\
4.9908	0.019782\\
4.9725	0.020825\\
4.9542	0.021923\\
4.936	0.023079\\
4.9177	0.024296\\
4.8994	0.025578\\
4.8811	0.026928\\
4.8628	0.028348\\
4.8445	0.029845\\
4.8262	0.03142\\
4.8079	0.033079\\
4.7896	0.034825\\
4.7713	0.036664\\
4.7531	0.038601\\
4.7348	0.04064\\
4.7165	0.042787\\
4.6982	0.045048\\
4.6799	0.047429\\
4.6616	0.049937\\
4.6433	0.052577\\
4.625	0.055357\\
4.6067	0.058284\\
4.5884	0.061367\\
4.5702	0.064614\\
4.5519	0.068032\\
4.5336	0.071633\\
4.5153	0.075424\\
4.497	0.079416\\
4.4787	0.08362\\
4.4604	0.088048\\
4.4421	0.09271\\
4.4238	0.09762\\
4.4055	0.10279\\
4.3873	0.10824\\
4.369	0.11397\\
4.3507	0.12001\\
4.3324	0.12637\\
4.3141	0.13306\\
4.2958	0.14012\\
4.2775	0.14754\\
4.2592	0.15536\\
4.2409	0.1636\\
4.2226	0.17227\\
4.2044	0.1814\\
4.1861	0.19102\\
4.1678	0.20114\\
4.1495	0.2118\\
4.1312	0.22302\\
4.1129	0.23483\\
4.0946	0.24727\\
4.0763	0.26036\\
4.058	0.27414\\
4.0397	0.28864\\
4.0215	0.30391\\
4.0032	0.31996\\
3.9849	0.33686\\
3.9666	0.35463\\
3.9483	0.37332\\
3.93	0.39298\\
3.9117	0.41366\\
3.8934	0.4354\\
3.8751	0.45827\\
3.8568	0.48232\\
3.8386	0.50761\\
3.8203	0.53421\\
3.802	0.56218\\
3.7837	0.59161\\
3.7654	0.62256\\
3.7471	0.65513\\
3.7288	0.68938\\
3.7105	0.72542\\
3.6922	0.76333\\
3.6739	0.80322\\
3.6557	0.8452\\
3.6374	0.88898\\
3.6191	0.93077\\
3.6008	0.96523\\
}--cycle;
      % \input{../../../../.tex/20/KS/SizeDistributionFittingsfig43.tex}
      \legend{}
  
      \nextgroupplot[%
        xmin=3.5608,xmax=5.3544,
        % xmin=3.4,xmax=5.4,
        % ymin=0.0052169,ymax=1,
      ] % This file was created by matlab2tikz.
%
\definecolor{mycolor1}{rgb}{0.00000,0.44706,0.69804}%
\definecolor{mycolor2}{rgb}{0.00000,0.44700,0.74100}%
\definecolor{mycolor3}{rgb}{0.49400,0.18400,0.55600}%
%
\addplot[only marks, mark=*, mark size=1.1180pt, color=mycolor1, fill=mycolor1, opacity=0.60, draw=none, mark options={draw=none,line width=0pt}] table[row sep=crcr]{%
x	y\\
3.5984	0.99468\\
3.6027	0.98404\\
3.6068	0.9734\\
3.6112	0.96277\\
3.6162	0.95213\\
3.6211	0.94149\\
3.6266	0.93085\\
3.6309	0.92021\\
3.636	0.90957\\
3.64	0.89894\\
3.6441	0.8883\\
3.6477	0.87766\\
3.6529	0.86702\\
3.6581	0.85638\\
3.6641	0.84574\\
3.6687	0.83511\\
3.674	0.82447\\
3.6794	0.81383\\
3.6854	0.80319\\
3.6908	0.79255\\
3.6952	0.78191\\
3.7011	0.77128\\
3.707	0.76064\\
3.713	0.75\\
3.7182	0.73936\\
3.7241	0.72872\\
3.73	0.71809\\
3.736	0.70745\\
3.7427	0.69681\\
3.7494	0.68617\\
3.7566	0.67553\\
3.7628	0.66489\\
3.7689	0.65426\\
3.7752	0.64362\\
3.7809	0.63298\\
3.7869	0.62234\\
3.7932	0.6117\\
3.8004	0.60106\\
3.8084	0.59043\\
3.8147	0.57979\\
3.8215	0.56915\\
3.8298	0.55851\\
3.8363	0.54787\\
3.8438	0.53723\\
3.8521	0.5266\\
3.8604	0.51596\\
3.8683	0.50532\\
3.8765	0.49468\\
3.8843	0.48404\\
3.8928	0.4734\\
3.9007	0.46277\\
3.9098	0.45213\\
3.9208	0.44149\\
3.9308	0.43085\\
3.9405	0.42021\\
3.9501	0.40957\\
3.9606	0.39894\\
3.9709	0.3883\\
3.9801	0.37766\\
3.9899	0.36702\\
3.999	0.35638\\
4.0096	0.34574\\
4.0188	0.33511\\
4.0298	0.32447\\
4.0417	0.31383\\
4.0542	0.30319\\
4.0666	0.29255\\
4.0789	0.28191\\
4.0918	0.27128\\
4.1041	0.26064\\
4.1175	0.25\\
4.131	0.23936\\
4.1454	0.22872\\
4.159	0.21809\\
4.173	0.20745\\
4.1912	0.19681\\
4.2086	0.18617\\
4.2269	0.17553\\
4.2454	0.16489\\
4.2692	0.15426\\
4.2908	0.14362\\
4.3146	0.13298\\
4.3388	0.12234\\
4.362	0.1117\\
4.3887	0.10106\\
4.4173	0.090426\\
4.4494	0.079787\\
4.4812	0.069149\\
4.5257	0.058511\\
4.5784	0.047872\\
4.6417	0.037234\\
4.741	0.026596\\
4.8583	0.015957\\
5.2036	0.0053191\\
};
\addlegendentry{\ExactEmpiricCCDFLegend}

\addplot [color=mycolor1, only marks, every error bar/.append style={opacity=0.30}, mark=*, mark size=0pt, draw=none, forget plot]
 plot [error bars/.cd, x dir=both, x explicit, error bar style={line width=1pt, color=mycolor1}, error mark options={mark=none,mark size=0pt}]
 table[row sep=crcr, x error plus index=2, x error minus index=3]{%
3.5984	0.99468	0.00473	0.00473\\
3.6027	0.98404	0.0047925	0.0047925\\
3.6068	0.9734	0.0048115	0.0048115\\
3.6112	0.96277	0.0047892	0.0047892\\
3.6162	0.95213	0.0047066	0.0047066\\
3.6211	0.94149	0.0049168	0.0049168\\
3.6266	0.93085	0.0048041	0.0048041\\
3.6309	0.92021	0.004932	0.004932\\
3.636	0.90957	0.0049959	0.0049959\\
3.64	0.89894	0.005017	0.005017\\
3.6441	0.8883	0.0051081	0.0051081\\
3.6477	0.87766	0.005056	0.005056\\
3.6529	0.86702	0.0050923	0.0050923\\
3.6581	0.85638	0.0052229	0.0052229\\
3.6641	0.84574	0.005148	0.005148\\
3.6687	0.83511	0.0052242	0.0052242\\
3.674	0.82447	0.0051186	0.0051186\\
3.6794	0.81383	0.005215	0.005215\\
3.6854	0.80319	0.0051265	0.0051265\\
3.6908	0.79255	0.0052431	0.0052431\\
3.6952	0.78191	0.0051884	0.0051884\\
3.7011	0.77128	0.0051923	0.0051923\\
3.707	0.76064	0.0050709	0.0050709\\
3.713	0.75	0.005022	0.005022\\
3.7182	0.73936	0.0050527	0.0050527\\
3.7241	0.72872	0.0052125	0.0052125\\
3.73	0.71809	0.0052284	0.0052284\\
3.736	0.70745	0.0051795	0.0051795\\
3.7427	0.69681	0.0052847	0.0052847\\
3.7494	0.68617	0.0054013	0.0054013\\
3.7566	0.67553	0.005546	0.005546\\
3.7628	0.66489	0.0056581	0.0056581\\
3.7689	0.65426	0.0055104	0.0055104\\
3.7752	0.64362	0.0056133	0.0056133\\
3.7809	0.63298	0.0055232	0.0055232\\
3.7869	0.62234	0.0054689	0.0054689\\
3.7932	0.6117	0.0054553	0.0054553\\
3.8004	0.60106	0.0057052	0.0057052\\
3.8084	0.59043	0.0058848	0.0058848\\
3.8147	0.57979	0.0057587	0.0057587\\
3.8215	0.56915	0.0059265	0.0059265\\
3.8298	0.55851	0.0061686	0.0061686\\
3.8363	0.54787	0.0061957	0.0061957\\
3.8438	0.53723	0.0060905	0.0060905\\
3.8521	0.5266	0.0059675	0.0059675\\
3.8604	0.51596	0.0058526	0.0058526\\
3.8683	0.50532	0.0057498	0.0057498\\
3.8765	0.49468	0.0057298	0.0057298\\
3.8843	0.48404	0.0056797	0.0056797\\
3.8928	0.4734	0.0056628	0.0056628\\
3.9007	0.46277	0.0055566	0.0055566\\
3.9098	0.45213	0.0053166	0.0053166\\
3.9208	0.44149	0.0054396	0.0054396\\
3.9308	0.43085	0.005428	0.005428\\
3.9405	0.42021	0.0055337	0.0055337\\
3.9501	0.40957	0.0053756	0.0053756\\
3.9606	0.39894	0.0054857	0.0054857\\
3.9709	0.3883	0.0055337	0.0055337\\
3.9801	0.37766	0.0056598	0.0056598\\
3.9899	0.36702	0.0058367	0.0058367\\
3.999	0.35638	0.0056606	0.0056606\\
4.0096	0.34574	0.0056459	0.0056459\\
4.0188	0.33511	0.0059698	0.0059698\\
4.0298	0.32447	0.0059241	0.0059241\\
4.0417	0.31383	0.0056343	0.0056343\\
4.0542	0.30319	0.0055073	0.0055073\\
4.0666	0.29255	0.0056644	0.0056644\\
4.0789	0.28191	0.005727	0.005727\\
4.0918	0.27128	0.0057123	0.0057123\\
4.1041	0.26064	0.0056946	0.0056946\\
4.1175	0.25	0.0059084	0.0059084\\
4.131	0.23936	0.0061124	0.0061124\\
4.1454	0.22872	0.0066665	0.0066665\\
4.159	0.21809	0.0068485	0.0068485\\
4.173	0.20745	0.0069575	0.0069575\\
4.1912	0.19681	0.0072717	0.0072717\\
4.2086	0.18617	0.0078688	0.0078688\\
4.2269	0.17553	0.0077805	0.0077805\\
4.2454	0.16489	0.0078623	0.0078623\\
4.2692	0.15426	0.0083156	0.0083156\\
4.2908	0.14362	0.0089906	0.0089906\\
4.3146	0.13298	0.009299	0.009299\\
4.3388	0.12234	0.0084958	0.0084958\\
4.362	0.1117	0.0078765	0.0078765\\
4.3887	0.10106	0.008669	0.008669\\
4.4173	0.090426	0.0097937	0.0097937\\
4.4494	0.079787	0.01038	0.01038\\
4.4812	0.069149	0.011015	0.011015\\
4.5257	0.058511	0.012142	0.012142\\
4.5784	0.047872	0.011936	0.011936\\
4.6417	0.037234	0.014498	0.014498\\
4.741	0.026596	0.016486	0.016486\\
4.8583	0.015957	0.020509	0.020509\\
5.2036	0.0053191	0.018099	0.018099\\
};
\addplot [color=mycolor1]
  table[row sep=crcr]{%
3.598	0.97291\\
3.6157	0.9421\\
3.6335	0.90312\\
3.6689	0.82039\\
4.0415	0.29033\\
4.4318	0.098196\\
4.8399	0.03175\\
5.2479	0.010308\\
5.3544	0.0076913\\
};
\addlegendentry{\ExactParetoFitLegend}


\addplot[area legend, draw=none, fill=mycolor1, fill opacity=0.15, forget plot]
table[row sep=crcr] {%
x	y\\
3.598	0.97938\\
3.6157	0.95118\\
3.6335	0.91331\\
3.6512	0.87156\\
3.6689	0.83003\\
3.6867	0.7895\\
3.7044	0.75098\\
3.7222	0.71435\\
3.7399	0.67953\\
3.7577	0.64642\\
3.7754	0.61495\\
3.7931	0.58502\\
3.8109	0.55658\\
3.8286	0.52954\\
3.8464	0.50384\\
3.8641	0.4794\\
3.8818	0.45618\\
3.8996	0.43411\\
3.9173	0.41312\\
3.9351	0.39318\\
3.9528	0.37422\\
3.9705	0.35619\\
3.9883	0.33906\\
4.006	0.32277\\
4.0238	0.30727\\
4.0415	0.29253\\
4.0593	0.27852\\
4.077	0.26518\\
4.0947	0.2525\\
4.1125	0.24042\\
4.1302	0.22894\\
4.148	0.218\\
4.1657	0.2076\\
4.1834	0.19769\\
4.2012	0.18827\\
4.2189	0.17929\\
4.2367	0.17075\\
4.2544	0.16261\\
4.2722	0.15487\\
4.2899	0.1475\\
4.3076	0.14048\\
4.3254	0.13379\\
4.3431	0.12743\\
4.3609	0.12137\\
4.3786	0.11559\\
4.3963	0.1101\\
4.4141	0.10487\\
4.4318	0.099882\\
4.4496	0.095136\\
4.4673	0.090617\\
4.485	0.086313\\
4.5028	0.082214\\
4.5205	0.078311\\
4.5383	0.074594\\
4.556	0.071053\\
4.5738	0.067682\\
4.5915	0.06447\\
4.6092	0.061412\\
4.627	0.058499\\
4.6447	0.055725\\
4.6625	0.053083\\
4.6802	0.050566\\
4.6979	0.048169\\
4.7157	0.045886\\
4.7334	0.043712\\
4.7512	0.041641\\
4.7689	0.039668\\
4.7867	0.037789\\
4.8044	0.036\\
4.8221	0.034295\\
4.8399	0.032671\\
4.8576	0.031125\\
4.8754	0.029652\\
4.8931	0.028248\\
4.9108	0.026912\\
4.9286	0.025639\\
4.9463	0.024426\\
4.9641	0.023271\\
4.9818	0.02217\\
4.9995	0.021122\\
5.0173	0.020123\\
5.035	0.019172\\
5.0528	0.018266\\
5.0705	0.017403\\
5.0883	0.01658\\
5.106	0.015797\\
5.1237	0.015051\\
5.1415	0.01434\\
5.1592	0.013663\\
5.177	0.013018\\
5.1947	0.012403\\
5.2124	0.011818\\
5.2302	0.01126\\
5.2479	0.010728\\
5.2657	0.010222\\
5.2834	0.0097399\\
5.3012	0.0092805\\
5.3189	0.0088428\\
5.3366	0.0084257\\
5.3544	0.0080284\\
5.3544	0.0073542\\
5.3366	0.0077259\\
5.3189	0.0081165\\
5.3012	0.0085269\\
5.2834	0.0089582\\
5.2657	0.0094113\\
5.2479	0.0098874\\
5.2302	0.010388\\
5.2124	0.010913\\
5.1947	0.011466\\
5.177	0.012046\\
5.1592	0.012656\\
5.1415	0.013297\\
5.1237	0.01397\\
5.106	0.014678\\
5.0883	0.015421\\
5.0705	0.016203\\
5.0528	0.017024\\
5.035	0.017887\\
5.0173	0.018794\\
4.9995	0.019746\\
4.9818	0.020748\\
4.9641	0.0218\\
4.9463	0.022906\\
4.9286	0.024068\\
4.9108	0.025289\\
4.8931	0.026572\\
4.8754	0.027921\\
4.8576	0.029339\\
4.8399	0.030828\\
4.8221	0.032394\\
4.8044	0.034039\\
4.7867	0.035768\\
4.7689	0.037585\\
4.7512	0.039495\\
4.7334	0.041503\\
4.7157	0.043612\\
4.6979	0.045829\\
4.6802	0.04816\\
4.6625	0.050609\\
4.6447	0.053183\\
4.627	0.055888\\
4.6092	0.058732\\
4.5915	0.06172\\
4.5738	0.064862\\
4.556	0.068163\\
4.5383	0.071633\\
4.5205	0.075281\\
4.5028	0.079114\\
4.485	0.083144\\
4.4673	0.087379\\
4.4496	0.091831\\
4.4318	0.09651\\
4.4141	0.10143\\
4.3963	0.1066\\
4.3786	0.11203\\
4.3609	0.11774\\
4.3431	0.12374\\
4.3254	0.13005\\
4.3076	0.13668\\
4.2899	0.14365\\
4.2722	0.15098\\
4.2544	0.15868\\
4.2367	0.16677\\
4.2189	0.17528\\
4.2012	0.18422\\
4.1834	0.19361\\
4.1657	0.20348\\
4.148	0.21386\\
4.1302	0.22476\\
4.1125	0.23622\\
4.0947	0.24825\\
4.077	0.26089\\
4.0593	0.27418\\
4.0415	0.28813\\
4.0238	0.30278\\
4.006	0.31817\\
3.9883	0.33433\\
3.9705	0.35129\\
3.9528	0.3691\\
3.9351	0.3878\\
3.9173	0.40743\\
3.8996	0.42803\\
3.8818	0.44966\\
3.8641	0.47236\\
3.8464	0.49619\\
3.8286	0.5212\\
3.8109	0.54746\\
3.7931	0.57503\\
3.7754	0.60398\\
3.7577	0.63437\\
3.7399	0.66629\\
3.7222	0.6998\\
3.7044	0.73498\\
3.6867	0.77194\\
3.6689	0.81074\\
3.6512	0.85125\\
3.6335	0.89293\\
3.6157	0.93302\\
3.598	0.96644\\
}--cycle;
\addplot[only marks, mark=*, mark size=1.1180pt, color=black, fill=black, opacity=0.60, draw=none, mark options={draw=none,line width=0pt}] table[row sep=crcr]{%
x	y\\
3.5623	0.99468\\
3.5723	0.98404\\
3.573	0.9734\\
3.5753	0.96277\\
3.5766	0.95213\\
3.5781	0.94149\\
3.5826	0.93085\\
3.5851	0.92021\\
3.5911	0.90957\\
3.5942	0.89894\\
3.596	0.8883\\
3.6016	0.87766\\
3.603	0.86702\\
3.6047	0.85638\\
3.6086	0.84574\\
3.6086	0.83511\\
3.6097	0.82447\\
3.6142	0.81383\\
3.6261	0.80319\\
3.6383	0.79255\\
3.6464	0.78191\\
3.6535	0.77128\\
3.657	0.76064\\
3.6658	0.75\\
3.6678	0.73936\\
3.6687	0.72872\\
3.6702	0.71809\\
3.6723	0.70745\\
3.6793	0.69681\\
3.6867	0.68617\\
3.6949	0.67553\\
3.7015	0.66489\\
3.7031	0.65426\\
3.719	0.64362\\
3.7213	0.63298\\
3.7265	0.62234\\
3.7307	0.6117\\
3.737	0.60106\\
3.7374	0.59043\\
3.7423	0.57979\\
3.7487	0.56915\\
3.7677	0.55851\\
3.7685	0.54787\\
3.7768	0.53723\\
3.8005	0.5266\\
3.8015	0.51596\\
3.8022	0.50532\\
3.8032	0.49468\\
3.8039	0.48404\\
3.8108	0.4734\\
3.8214	0.46277\\
3.8405	0.45213\\
3.8605	0.44149\\
3.863	0.43085\\
3.8645	0.42021\\
3.8662	0.40957\\
3.8806	0.39894\\
3.8878	0.3883\\
3.8964	0.37766\\
3.905	0.36702\\
3.9294	0.35638\\
3.9303	0.34574\\
3.954	0.33511\\
3.9604	0.32447\\
3.9669	0.31383\\
3.9747	0.30319\\
3.9822	0.29255\\
3.9929	0.28191\\
4.0051	0.27128\\
4.0144	0.26064\\
4.0161	0.25\\
4.0433	0.23936\\
4.0578	0.22872\\
4.073	0.21809\\
4.0857	0.20745\\
4.0904	0.19681\\
4.1168	0.18617\\
4.1265	0.17553\\
4.127	0.16489\\
4.143	0.15426\\
4.1756	0.14362\\
4.2156	0.13298\\
4.3116	0.12234\\
4.3276	0.1117\\
4.3662	0.10106\\
4.4791	0.090426\\
4.4912	0.079787\\
4.517	0.069149\\
4.5743	0.058511\\
4.5914	0.047872\\
4.6138	0.037234\\
4.7898	0.026596\\
5.0876	0.015957\\
5.3101	0.0053191\\
};
\addlegendentry{\RealEmpiricCCDFLegend}

\addplot [color=black]
  table[row sep=crcr]{%
3.5608	1\\
5.3544	0.0052169\\
};
\addlegendentry{\RealParetoFitLegend}

      % \input{../../../../.tex/20/KS/SizeDistributionFittingsfig44.tex}
      \legend{}
  
      \nextgroupplot[%
        xmin=3.5608,xmax=5.4115,
        % xmin=3.5,xmax=5.5,
        % ymin=0.0044129,ymax=1,
      ] % This file was created by matlab2tikz.
%
\definecolor{mycolor1}{rgb}{0.83529,0.36863,0.00000}%
\definecolor{mycolor2}{rgb}{0.00000,0.44700,0.74100}%
\definecolor{mycolor3}{rgb}{0.49400,0.18400,0.55600}%
%
\addplot[only marks, mark=*, mark size=1.1180pt, color=mycolor1, fill=mycolor1, opacity=0.60, draw=none, mark options={draw=none,line width=0pt}] table[row sep=crcr]{%
x	y\\
3.6	0.99468\\
3.6045	0.98404\\
3.6089	0.9734\\
3.6141	0.96277\\
3.6178	0.95213\\
3.6225	0.94149\\
3.6271	0.93085\\
3.6318	0.92021\\
3.6365	0.90957\\
3.6418	0.89894\\
3.6472	0.8883\\
3.6515	0.87766\\
3.657	0.86702\\
3.6615	0.85638\\
3.6662	0.84574\\
3.6709	0.83511\\
3.6761	0.82447\\
3.6804	0.81383\\
3.6852	0.80319\\
3.6904	0.79255\\
3.6957	0.78191\\
3.701	0.77128\\
3.7064	0.76064\\
3.7121	0.75\\
3.7185	0.73936\\
3.7229	0.72872\\
3.7274	0.71809\\
3.7333	0.70745\\
3.7389	0.69681\\
3.745	0.68617\\
3.7505	0.67553\\
3.7574	0.66489\\
3.7644	0.65426\\
3.7707	0.64362\\
3.7767	0.63298\\
3.7836	0.62234\\
3.7905	0.6117\\
3.7976	0.60106\\
3.8048	0.59043\\
3.812	0.57979\\
3.8198	0.56915\\
3.8269	0.55851\\
3.8352	0.54787\\
3.8439	0.53723\\
3.8506	0.5266\\
3.8568	0.51596\\
3.8635	0.50532\\
3.8723	0.49468\\
3.8796	0.48404\\
3.8878	0.4734\\
3.897	0.46277\\
3.9047	0.45213\\
3.9134	0.44149\\
3.9222	0.43085\\
3.9309	0.42021\\
3.9403	0.40957\\
3.9517	0.39894\\
3.9628	0.3883\\
3.9726	0.37766\\
3.9842	0.36702\\
3.9942	0.35638\\
4.0054	0.34574\\
4.0156	0.33511\\
4.0277	0.32447\\
4.0395	0.31383\\
4.0516	0.30319\\
4.066	0.29255\\
4.0769	0.28191\\
4.0914	0.27128\\
4.1049	0.26064\\
4.1181	0.25\\
4.1339	0.23936\\
4.1473	0.22872\\
4.1627	0.21809\\
4.1783	0.20745\\
4.1955	0.19681\\
4.2126	0.18617\\
4.231	0.17553\\
4.254	0.16489\\
4.2752	0.15426\\
4.2969	0.14362\\
4.3184	0.13298\\
4.3399	0.12234\\
4.364	0.1117\\
4.39	0.10106\\
4.4189	0.090426\\
4.4573	0.079787\\
4.4912	0.069149\\
4.5391	0.058511\\
4.5879	0.047872\\
4.6566	0.037234\\
4.7343	0.026596\\
4.8688	0.015957\\
5.1869	0.0053191\\
};
\addlegendentry{\ApproxEmpiricCCDFLegend}

\addplot [color=mycolor1, only marks, every error bar/.append style={opacity=0.30}, mark=*, mark size=0pt, draw=none, forget plot]
 plot [error bars/.cd, x dir=both, x explicit, error bar style={line width=1pt, color=mycolor1}, error mark options={mark=none,mark size=0pt}]
 table[row sep=crcr, x error plus index=2, x error minus index=3]{%
3.6	0.99468	0.0045554	0.0045554\\
3.6045	0.98404	0.0045703	0.0045703\\
3.6089	0.9734	0.0045491	0.0045491\\
3.6141	0.96277	0.0044812	0.0044812\\
3.6178	0.95213	0.0044817	0.0044817\\
3.6225	0.94149	0.0044624	0.0044624\\
3.6271	0.93085	0.0045762	0.0045762\\
3.6318	0.92021	0.0045625	0.0045625\\
3.6365	0.90957	0.004468	0.004468\\
3.6418	0.89894	0.0046379	0.0046379\\
3.6472	0.8883	0.0045394	0.0045394\\
3.6515	0.87766	0.0044728	0.0044728\\
3.657	0.86702	0.0045858	0.0045858\\
3.6615	0.85638	0.0045704	0.0045704\\
3.6662	0.84574	0.0044305	0.0044305\\
3.6709	0.83511	0.0043954	0.0043954\\
3.6761	0.82447	0.0045809	0.0045809\\
3.6804	0.81383	0.0045127	0.0045127\\
3.6852	0.80319	0.0045389	0.0045389\\
3.6904	0.79255	0.0046013	0.0046013\\
3.6957	0.78191	0.0046842	0.0046842\\
3.701	0.77128	0.0046214	0.0046214\\
3.7064	0.76064	0.0046093	0.0046093\\
3.7121	0.75	0.0045663	0.0045663\\
3.7185	0.73936	0.0045568	0.0045568\\
3.7229	0.72872	0.0046537	0.0046537\\
3.7274	0.71809	0.0046666	0.0046666\\
3.7333	0.70745	0.0046756	0.0046756\\
3.7389	0.69681	0.0047365	0.0047365\\
3.745	0.68617	0.0047282	0.0047282\\
3.7505	0.67553	0.0046776	0.0046776\\
3.7574	0.66489	0.0046232	0.0046232\\
3.7644	0.65426	0.0049389	0.0049389\\
3.7707	0.64362	0.0049657	0.0049657\\
3.7767	0.63298	0.0051864	0.0051864\\
3.7836	0.62234	0.0051426	0.0051426\\
3.7905	0.6117	0.0054154	0.0054154\\
3.7976	0.60106	0.0052834	0.0052834\\
3.8048	0.59043	0.0053245	0.0053245\\
3.812	0.57979	0.0053216	0.0053216\\
3.8198	0.56915	0.0053218	0.0053218\\
3.8269	0.55851	0.0052938	0.0052938\\
3.8352	0.54787	0.0052741	0.0052741\\
3.8439	0.53723	0.0055718	0.0055718\\
3.8506	0.5266	0.005749	0.005749\\
3.8568	0.51596	0.0058641	0.0058641\\
3.8635	0.50532	0.0058242	0.0058242\\
3.8723	0.49468	0.0060846	0.0060846\\
3.8796	0.48404	0.0058831	0.0058831\\
3.8878	0.4734	0.0060552	0.0060552\\
3.897	0.46277	0.0058784	0.0058784\\
3.9047	0.45213	0.0058689	0.0058689\\
3.9134	0.44149	0.0059804	0.0059804\\
3.9222	0.43085	0.0059113	0.0059113\\
3.9309	0.42021	0.0061887	0.0061887\\
3.9403	0.40957	0.0064169	0.0064169\\
3.9517	0.39894	0.0061866	0.0061866\\
3.9628	0.3883	0.0059948	0.0059948\\
3.9726	0.37766	0.006044	0.006044\\
3.9842	0.36702	0.0061542	0.0061542\\
3.9942	0.35638	0.0061402	0.0061402\\
4.0054	0.34574	0.006197	0.006197\\
4.0156	0.33511	0.0062917	0.0062917\\
4.0277	0.32447	0.0062543	0.0062543\\
4.0395	0.31383	0.006648	0.006648\\
4.0516	0.30319	0.0062708	0.0062708\\
4.066	0.29255	0.0064445	0.0064445\\
4.0769	0.28191	0.0065803	0.0065803\\
4.0914	0.27128	0.0068367	0.0068367\\
4.1049	0.26064	0.0067182	0.0067182\\
4.1181	0.25	0.0066727	0.0066727\\
4.1339	0.23936	0.006414	0.006414\\
4.1473	0.22872	0.0068788	0.0068788\\
4.1627	0.21809	0.0073314	0.0073314\\
4.1783	0.20745	0.0073394	0.0073394\\
4.1955	0.19681	0.0076429	0.0076429\\
4.2126	0.18617	0.007675	0.007675\\
4.231	0.17553	0.0079381	0.0079381\\
4.254	0.16489	0.0079159	0.0079159\\
4.2752	0.15426	0.0082002	0.0082002\\
4.2969	0.14362	0.0088204	0.0088204\\
4.3184	0.13298	0.00845	0.00845\\
4.3399	0.12234	0.0086796	0.0086796\\
4.364	0.1117	0.0086382	0.0086382\\
4.39	0.10106	0.0099918	0.0099918\\
4.4189	0.090426	0.010413	0.010413\\
4.4573	0.079787	0.011005	0.011005\\
4.4912	0.069149	0.011629	0.011629\\
4.5391	0.058511	0.011499	0.011499\\
4.5879	0.047872	0.012519	0.012519\\
4.6566	0.037234	0.015101	0.015101\\
4.7343	0.026596	0.016441	0.016441\\
4.8688	0.015957	0.021019	0.021019\\
5.1869	0.0053191	0.018612	0.018612\\
};
\addplot [color=mycolor1]
  table[row sep=crcr]{%
3.6008	0.97204\\
3.6191	0.94016\\
3.6374	0.89927\\
3.6922	0.77128\\
4.0763	0.26255\\
4.4604	0.089744\\
4.8445	0.0308\\
5.2286	0.010613\\
5.4115	0.0063989\\
};
\addlegendentry{\ApproxParetoFitLegend}


\addplot[area legend, draw=none, fill=mycolor1, fill opacity=0.15, forget plot]
table[row sep=crcr] {%
x	y\\
3.6008	0.97886\\
3.6191	0.94955\\
3.6374	0.90957\\
3.6557	0.86444\\
3.6739	0.82071\\
3.6922	0.77922\\
3.7105	0.73984\\
3.7288	0.70247\\
3.7471	0.667\\
3.7654	0.63335\\
3.7837	0.60141\\
3.802	0.57111\\
3.8203	0.54236\\
3.8386	0.51508\\
3.8568	0.4892\\
3.8751	0.46464\\
3.8934	0.44135\\
3.9117	0.41925\\
3.93	0.39828\\
3.9483	0.37838\\
3.9666	0.3595\\
3.9849	0.34158\\
4.0032	0.32458\\
4.0215	0.30843\\
4.0397	0.29311\\
4.058	0.27856\\
4.0763	0.26474\\
4.0946	0.25162\\
4.1129	0.23915\\
4.1312	0.22731\\
4.1495	0.21607\\
4.1678	0.20538\\
4.1861	0.19523\\
4.2044	0.18558\\
4.2226	0.17641\\
4.2409	0.1677\\
4.2592	0.15942\\
4.2775	0.15156\\
4.2958	0.14408\\
4.3141	0.13698\\
4.3324	0.13022\\
4.3507	0.1238\\
4.369	0.1177\\
4.3873	0.1119\\
4.4055	0.10639\\
4.4238	0.10115\\
4.4421	0.096173\\
4.4604	0.091439\\
4.4787	0.08694\\
4.497	0.082662\\
4.5153	0.078596\\
4.5336	0.074731\\
4.5519	0.071056\\
4.5702	0.067563\\
4.5884	0.064242\\
4.6067	0.061085\\
4.625	0.058084\\
4.6433	0.055231\\
4.6616	0.052518\\
4.6799	0.049939\\
4.6982	0.047488\\
4.7165	0.045157\\
4.7348	0.042941\\
4.7531	0.040834\\
4.7713	0.03883\\
4.7896	0.036926\\
4.8079	0.035115\\
4.8262	0.033393\\
4.8445	0.031756\\
4.8628	0.030199\\
4.8811	0.02872\\
4.8994	0.027312\\
4.9177	0.025974\\
4.936	0.024702\\
4.9542	0.023492\\
4.9725	0.022342\\
4.9908	0.021248\\
5.0091	0.020208\\
5.0274	0.019219\\
5.0457	0.018279\\
5.064	0.017385\\
5.0823	0.016534\\
5.1006	0.015726\\
5.1189	0.014957\\
5.1371	0.014226\\
5.1554	0.01353\\
5.1737	0.012869\\
5.192	0.01224\\
5.2103	0.011642\\
5.2286	0.011074\\
5.2469	0.010533\\
5.2652	0.010019\\
5.2835	0.0095295\\
5.3018	0.0090644\\
5.32	0.008622\\
5.3383	0.0082013\\
5.3566	0.0078012\\
5.3749	0.0074206\\
5.3932	0.0070587\\
5.4115	0.0067145\\
5.4115	0.0060833\\
5.3932	0.0064028\\
5.3749	0.006739\\
5.3566	0.007093\\
5.3383	0.0074657\\
5.32	0.007858\\
5.3018	0.008271\\
5.2835	0.0087057\\
5.2652	0.0091634\\
5.2469	0.0096453\\
5.2286	0.010153\\
5.2103	0.010687\\
5.192	0.011249\\
5.1737	0.011841\\
5.1554	0.012464\\
5.1371	0.01312\\
5.1189	0.01381\\
5.1006	0.014538\\
5.0823	0.015303\\
5.064	0.016109\\
5.0457	0.016958\\
5.0274	0.017851\\
5.0091	0.018792\\
4.9908	0.019782\\
4.9725	0.020825\\
4.9542	0.021923\\
4.936	0.023079\\
4.9177	0.024296\\
4.8994	0.025578\\
4.8811	0.026928\\
4.8628	0.028348\\
4.8445	0.029845\\
4.8262	0.03142\\
4.8079	0.033079\\
4.7896	0.034825\\
4.7713	0.036664\\
4.7531	0.038601\\
4.7348	0.04064\\
4.7165	0.042787\\
4.6982	0.045048\\
4.6799	0.047429\\
4.6616	0.049937\\
4.6433	0.052577\\
4.625	0.055357\\
4.6067	0.058284\\
4.5884	0.061367\\
4.5702	0.064614\\
4.5519	0.068032\\
4.5336	0.071633\\
4.5153	0.075424\\
4.497	0.079416\\
4.4787	0.08362\\
4.4604	0.088048\\
4.4421	0.09271\\
4.4238	0.09762\\
4.4055	0.10279\\
4.3873	0.10824\\
4.369	0.11397\\
4.3507	0.12001\\
4.3324	0.12637\\
4.3141	0.13306\\
4.2958	0.14012\\
4.2775	0.14754\\
4.2592	0.15536\\
4.2409	0.1636\\
4.2226	0.17227\\
4.2044	0.1814\\
4.1861	0.19102\\
4.1678	0.20114\\
4.1495	0.2118\\
4.1312	0.22302\\
4.1129	0.23483\\
4.0946	0.24727\\
4.0763	0.26036\\
4.058	0.27414\\
4.0397	0.28864\\
4.0215	0.30391\\
4.0032	0.31996\\
3.9849	0.33686\\
3.9666	0.35463\\
3.9483	0.37332\\
3.93	0.39298\\
3.9117	0.41366\\
3.8934	0.4354\\
3.8751	0.45827\\
3.8568	0.48232\\
3.8386	0.50761\\
3.8203	0.53421\\
3.802	0.56218\\
3.7837	0.59161\\
3.7654	0.62256\\
3.7471	0.65513\\
3.7288	0.68938\\
3.7105	0.72542\\
3.6922	0.76333\\
3.6739	0.80322\\
3.6557	0.8452\\
3.6374	0.88898\\
3.6191	0.93077\\
3.6008	0.96523\\
}--cycle;
\addplot[only marks, mark=*, mark size=1.1180pt, color=black, fill=black, opacity=0.60, draw=none, mark options={draw=none,line width=0pt}] table[row sep=crcr]{%
x	y\\
3.5623	0.99468\\
3.5723	0.98404\\
3.573	0.9734\\
3.5753	0.96277\\
3.5766	0.95213\\
3.5781	0.94149\\
3.5826	0.93085\\
3.5851	0.92021\\
3.5911	0.90957\\
3.5942	0.89894\\
3.596	0.8883\\
3.6016	0.87766\\
3.603	0.86702\\
3.6047	0.85638\\
3.6086	0.84574\\
3.6086	0.83511\\
3.6097	0.82447\\
3.6142	0.81383\\
3.6261	0.80319\\
3.6383	0.79255\\
3.6464	0.78191\\
3.6535	0.77128\\
3.657	0.76064\\
3.6658	0.75\\
3.6678	0.73936\\
3.6687	0.72872\\
3.6702	0.71809\\
3.6723	0.70745\\
3.6793	0.69681\\
3.6867	0.68617\\
3.6949	0.67553\\
3.7015	0.66489\\
3.7031	0.65426\\
3.719	0.64362\\
3.7213	0.63298\\
3.7265	0.62234\\
3.7307	0.6117\\
3.737	0.60106\\
3.7374	0.59043\\
3.7423	0.57979\\
3.7487	0.56915\\
3.7677	0.55851\\
3.7685	0.54787\\
3.7768	0.53723\\
3.8005	0.5266\\
3.8015	0.51596\\
3.8022	0.50532\\
3.8032	0.49468\\
3.8039	0.48404\\
3.8108	0.4734\\
3.8214	0.46277\\
3.8405	0.45213\\
3.8605	0.44149\\
3.863	0.43085\\
3.8645	0.42021\\
3.8662	0.40957\\
3.8806	0.39894\\
3.8878	0.3883\\
3.8964	0.37766\\
3.905	0.36702\\
3.9294	0.35638\\
3.9303	0.34574\\
3.954	0.33511\\
3.9604	0.32447\\
3.9669	0.31383\\
3.9747	0.30319\\
3.9822	0.29255\\
3.9929	0.28191\\
4.0051	0.27128\\
4.0144	0.26064\\
4.0161	0.25\\
4.0433	0.23936\\
4.0578	0.22872\\
4.073	0.21809\\
4.0857	0.20745\\
4.0904	0.19681\\
4.1168	0.18617\\
4.1265	0.17553\\
4.127	0.16489\\
4.143	0.15426\\
4.1756	0.14362\\
4.2156	0.13298\\
4.3116	0.12234\\
4.3276	0.1117\\
4.3662	0.10106\\
4.4791	0.090426\\
4.4912	0.079787\\
4.517	0.069149\\
4.5743	0.058511\\
4.5914	0.047872\\
4.6138	0.037234\\
4.7898	0.026596\\
5.0876	0.015957\\
5.3101	0.0053191\\
};
\addlegendentry{\RealEmpiricCCDFLegend}

\addplot [color=black]
  table[row sep=crcr]{%
3.5608	1\\
5.4115	0.0044129\\
};
\addlegendentry{\RealParetoFitLegend}

      % \input{../../../../.tex/20/KS/SizeDistributionFittingsfig45.tex}
      \legend{}
    \end{groupplot}
  %\endgroup

  %\begingroup Parametri
    
%\begingroup Parametro
\providecommand\varyingParameterSymbol\lambda
\providecommand\varyingParameterList[1]{
  \ifcase#1\relax%
  \or\num{1e-2}% Caso 1
  \or\num{8.2e-2}% Caso 2
  \or\num{1.55e-1}% Caso 3
  \or\num{2.27e-1}% Caso 4
  \else\num{3e-1}% Caso 5
  \fi
}
%\endgroup

%\begingroup beta
\providecommand\betaExactList[1]{
  \ifcase#1\relax%
  \or\num{1.267+-0.007}% Caso 1
  \or\num{1.285+-0.010}% Caso 2
  \or\num{1.264+-0.015}% Caso 3
  \or\num{1.246+-0.015}% Caso 4
  \else\num{1.210+-0.015}% Caso 5
  \fi
}
\providecommand\betaAprxList[1]{
  \ifcase#1\relax%
  \or\num{1.264+-0.006}% Caso 1
  \or\num{1.270+-0.011}% Caso 2
  \or\num{1.280+-0.013}% Caso 3
  \or\num{1.250+-0.015}% Caso 4
  \else\num{1.218+-0.014}% Caso 5
  \fi
}
%\endgroup

%\begingroup xi
\providecommand\zetaExactList[1]{
  \ifcase#1\relax%
  \or\num{0.757+-0.004}% Caso 1
  \or\num{0.710+-0.026}% Caso 2
  \or\num{0.670+-0.035}% Caso 3
  \or\num{0.625+-0.052}% Caso 4
  \else\num{0.539+-0.059}% Caso 5
  \fi
}
\providecommand\zetaAprxList[1]{
  \ifcase#1\relax%
  \or\num{0.749+-0.004}% Caso 1
  \or\num{0.721+-0.022}% Caso 2
  \or\num{0.645+-0.040}% Caso 3
  \or\num{0.622+-0.051}% Caso 4
  \else\num{0.567+-0.053}% Caso 5
  \fi
}
%\endgroup

%\begingroup mI
\providecommand\mIExactList[1]{
  \ifcase#1\relax%
  \or\num{3.265+-0.001}% Caso 1
  \or\num{3.238+-0.005}% Caso 2
  \or\num{3.220+-0.008}% Caso 3
  \or\num{3.213+-0.013}% Caso 4
  \else\num{3.174+-0.023}% Caso 5
  \fi
}
\providecommand\mIAprxList[1]{
  \ifcase#1\relax%
  \or\num{3.262+-0.001}% Caso 1
  \or\num{3.240+-0.005}% Caso 2
  \or\num{3.217+-0.010}% Caso 3
  \or\num{3.217+-0.012}% Caso 4
  \else\num{3.208+-0.009}% Caso 5
  \fi
}
%\endgroup

%\begingroup sI
\providecommand\sIExactList[1]{
  \ifcase#1\relax%
  \or\num{0.284+-0.001}% Caso 1
  \or\num{0.311+-0.008}% Caso 2
  \or\num{0.343+-0.010}% Caso 3
  \or\num{0.361+-0.020}% Caso 4
  \else\num{0.398+-0.029}% Caso 5
  \fi
}
\providecommand\sIAprxList[1]{
  \ifcase#1\relax%
  \or\num{0.281+-0.001}% Caso 1
  \or\num{0.315+-0.006}% Caso 2
  \or\num{0.334+-0.013}% Caso 3
  \or\num{0.366+-0.018}% Caso 4
  \else\num{0.411+-0.025}% Caso 5
  \fi
}
%\endgroup

%\begingroup mII
\providecommand\mIIExactList[1]{
  \ifcase#1\relax%
  \or\num{3.788+-0.005}% Caso 1
  \or\num{3.730+-0.029}% Caso 2
  \or\num{3.681+-0.052}% Caso 3
  \or\num{3.650+-0.072}% Caso 4
  \else\num{3.572+-0.095}% Caso 5
  \fi
}
\providecommand\mIIAprxList[1]{
  \ifcase#1\relax%
  \or\num{3.778+-0.006}% Caso 1
  \or\num{3.735+-0.029}% Caso 2
  \or\num{3.656+-0.048}% Caso 3
  \or\num{3.658+-0.084}% Caso 4
  \else\num{3.502+-0.078}% Caso 5
  \fi
}
%\endgroup

%\begingroup sII
\providecommand\sIIExactList[1]{
  \ifcase#1\relax%
  \or\num{0.455+-0.002}% Caso 1
  \or\num{0.486+-0.007}% Caso 2
  \or\num{0.515+-0.014}% Caso 3
  \or\num{0.533+-0.021}% Caso 4
  \else\num{0.570+-0.028}% Caso 5
  \fi
}
\providecommand\sIIAprxList[1]{
  \ifcase#1\relax%
  \or\num{0.457+-0.002}% Caso 1
  \or\num{0.490+-0.007}% Caso 2
  \or\num{0.521+-0.012}% Caso 3
  \or\num{0.545+-0.021}% Caso 4
  \else\num{0.586+-0.029}% Caso 5
  \fi
}
%\endgroup


    \foreach \r/\a/\p/\xS/\yS/\yLS in {%
      1/south west/south west/0em/0em/-.5em,
      2/south west/south west/0em/0em/-.5em,
      3/south west/south west/0em/0em/-.5em,
      4/south west/south west/0em/0em/-.5em,
      5/south west/south west/0em/0em/-.5em
    }{%
      \def\betaExact{\betaExactList\r}
      \def\betaAprx{\betaAprxList\r}

      \foreach \c/\fl/\sl in {1/0/1,2/0/2,3/1/2}{%
        \writeParetoParameters[\a]\fl\sl{Row\r\space c\c r1.\p}\xS\yS\yLS
      } % «\fl=first line», «\sl=second line», «\a=anchor»
       % «\xS=\xShift», «\yS=\yShift» e «\yLS=\yLineShift»
    }

    \ifpageisodd\writeVaryingParameter{1}{west}{west}{0em}{0em}
    \else\writeVaryingParameter{3}{east}{east}{0em}{0em}\fi
  %\endgroup

  %\begingroup Didascalie
    % \writeSubCaptions{north east}{north east}{0em}{0em}
    \ifpageisodd\writeSubCaptions{east}{east}{0em}{0em}
    \else\writeSubCaptions{west}{west}{0em}{0em}\fi
    % https://tex.stackexchange.com/questions/632045/get-subcaption-for-every-plot-using-groupplot-pgf-tikz-after-2020-update
  %\endgroup

  \redefineTikZbounds{Row1 c1r1}{Row5 c3r1}
\end{tikzpicture}%

% 
%\begingroup Parametro
\providecommand\varyingParameterSymbol\lambda
\providecommand\varyingParameterList[1]{
  \ifcase#1\relax%
  \or\num{1e-2}% Caso 1
  \or\num{8.2e-2}% Caso 2
  \or\num{1.55e-1}% Caso 3
  \or\num{2.27e-1}% Caso 4
  \else\num{3e-1}% Caso 5
  \fi
}
%\endgroup

%\begingroup beta
\providecommand\betaExactList[1]{
  \ifcase#1\relax%
  \or\num{1.267+-0.007}% Caso 1
  \or\num{1.285+-0.010}% Caso 2
  \or\num{1.264+-0.015}% Caso 3
  \or\num{1.246+-0.015}% Caso 4
  \else\num{1.210+-0.015}% Caso 5
  \fi
}
\providecommand\betaAprxList[1]{
  \ifcase#1\relax%
  \or\num{1.264+-0.006}% Caso 1
  \or\num{1.270+-0.011}% Caso 2
  \or\num{1.280+-0.013}% Caso 3
  \or\num{1.250+-0.015}% Caso 4
  \else\num{1.218+-0.014}% Caso 5
  \fi
}
%\endgroup

%\begingroup xi
\providecommand\zetaExactList[1]{
  \ifcase#1\relax%
  \or\num{0.757+-0.004}% Caso 1
  \or\num{0.710+-0.026}% Caso 2
  \or\num{0.670+-0.035}% Caso 3
  \or\num{0.625+-0.052}% Caso 4
  \else\num{0.539+-0.059}% Caso 5
  \fi
}
\providecommand\zetaAprxList[1]{
  \ifcase#1\relax%
  \or\num{0.749+-0.004}% Caso 1
  \or\num{0.721+-0.022}% Caso 2
  \or\num{0.645+-0.040}% Caso 3
  \or\num{0.622+-0.051}% Caso 4
  \else\num{0.567+-0.053}% Caso 5
  \fi
}
%\endgroup

%\begingroup mI
\providecommand\mIExactList[1]{
  \ifcase#1\relax%
  \or\num{3.265+-0.001}% Caso 1
  \or\num{3.238+-0.005}% Caso 2
  \or\num{3.220+-0.008}% Caso 3
  \or\num{3.213+-0.013}% Caso 4
  \else\num{3.174+-0.023}% Caso 5
  \fi
}
\providecommand\mIAprxList[1]{
  \ifcase#1\relax%
  \or\num{3.262+-0.001}% Caso 1
  \or\num{3.240+-0.005}% Caso 2
  \or\num{3.217+-0.010}% Caso 3
  \or\num{3.217+-0.012}% Caso 4
  \else\num{3.208+-0.009}% Caso 5
  \fi
}
%\endgroup

%\begingroup sI
\providecommand\sIExactList[1]{
  \ifcase#1\relax%
  \or\num{0.284+-0.001}% Caso 1
  \or\num{0.311+-0.008}% Caso 2
  \or\num{0.343+-0.010}% Caso 3
  \or\num{0.361+-0.020}% Caso 4
  \else\num{0.398+-0.029}% Caso 5
  \fi
}
\providecommand\sIAprxList[1]{
  \ifcase#1\relax%
  \or\num{0.281+-0.001}% Caso 1
  \or\num{0.315+-0.006}% Caso 2
  \or\num{0.334+-0.013}% Caso 3
  \or\num{0.366+-0.018}% Caso 4
  \else\num{0.411+-0.025}% Caso 5
  \fi
}
%\endgroup

%\begingroup mII
\providecommand\mIIExactList[1]{
  \ifcase#1\relax%
  \or\num{3.788+-0.005}% Caso 1
  \or\num{3.730+-0.029}% Caso 2
  \or\num{3.681+-0.052}% Caso 3
  \or\num{3.650+-0.072}% Caso 4
  \else\num{3.572+-0.095}% Caso 5
  \fi
}
\providecommand\mIIAprxList[1]{
  \ifcase#1\relax%
  \or\num{3.778+-0.006}% Caso 1
  \or\num{3.735+-0.029}% Caso 2
  \or\num{3.656+-0.048}% Caso 3
  \or\num{3.658+-0.084}% Caso 4
  \else\num{3.502+-0.078}% Caso 5
  \fi
}
%\endgroup

%\begingroup sII
\providecommand\sIIExactList[1]{
  \ifcase#1\relax%
  \or\num{0.455+-0.002}% Caso 1
  \or\num{0.486+-0.007}% Caso 2
  \or\num{0.515+-0.014}% Caso 3
  \or\num{0.533+-0.021}% Caso 4
  \else\num{0.570+-0.028}% Caso 5
  \fi
}
\providecommand\sIIAprxList[1]{
  \ifcase#1\relax%
  \or\num{0.457+-0.002}% Caso 1
  \or\num{0.490+-0.007}% Caso 2
  \or\num{0.521+-0.012}% Caso 3
  \or\num{0.545+-0.021}% Caso 4
  \else\num{0.586+-0.029}% Caso 5
  \fi
}
%\endgroup

% \input{Parametri.txt}

  \endgroup

  \begingroup
    \setcounter{subfigure}{0}
    
% !TEX root = ../../../../../Esperimenti/.tex/MEF.tex
% LTeX: language=it

\infopageKeysI

\begin{tikzpicture}%
  %\begingroup Righe
    \begin{groupplot}[
      group style={
        group name=Row1,
        plotGroup,
      },
      plotRow,
      row3Keys,
      plotLegendNE,
      %
      % legend style={%
      %   \placeLegend{0}{0}{south west}{2.5em}{0em}
      % },
      % xmin=3637.5,xmax=2.5003e+05,
      ymin=0.00088167,ymax=1.0352,
    ]
      \nextgroupplot[%
        xmin=3.5608,xmax=5.3101,
        % xmin=3.4,xmax=5.4,
        % ymin=0.0041299,ymax=1.0005,
      ] % This file was created by matlab2tikz.
%
\definecolor{mycolor1}{rgb}{0.00000,0.44700,0.74100}%
%
\addplot[only marks, mark=*, mark size=1.1180pt, color=black, fill=black, opacity=0.60, draw=none, mark options={draw=none,line width=0pt}] table[row sep=crcr]{%
x	y\\
3.5623	0.99468\\
3.5723	0.98404\\
3.573	0.9734\\
3.5753	0.96277\\
3.5766	0.95213\\
3.5781	0.94149\\
3.5826	0.93085\\
3.5851	0.92021\\
3.5911	0.90957\\
3.5942	0.89894\\
3.596	0.8883\\
3.6016	0.87766\\
3.603	0.86702\\
3.6047	0.85638\\
3.6086	0.84574\\
3.6086	0.83511\\
3.6097	0.82447\\
3.6142	0.81383\\
3.6261	0.80319\\
3.6383	0.79255\\
3.6464	0.78191\\
3.6535	0.77128\\
3.657	0.76064\\
3.6658	0.75\\
3.6678	0.73936\\
3.6687	0.72872\\
3.6702	0.71809\\
3.6723	0.70745\\
3.6793	0.69681\\
3.6867	0.68617\\
3.6949	0.67553\\
3.7015	0.66489\\
3.7031	0.65426\\
3.719	0.64362\\
3.7213	0.63298\\
3.7265	0.62234\\
3.7307	0.6117\\
3.737	0.60106\\
3.7374	0.59043\\
3.7423	0.57979\\
3.7487	0.56915\\
3.7677	0.55851\\
3.7685	0.54787\\
3.7768	0.53723\\
3.8005	0.5266\\
3.8015	0.51596\\
3.8022	0.50532\\
3.8032	0.49468\\
3.8039	0.48404\\
3.8108	0.4734\\
3.8214	0.46277\\
3.8405	0.45213\\
3.8605	0.44149\\
3.863	0.43085\\
3.8645	0.42021\\
3.8662	0.40957\\
3.8806	0.39894\\
3.8878	0.3883\\
3.8964	0.37766\\
3.905	0.36702\\
3.9294	0.35638\\
3.9303	0.34574\\
3.954	0.33511\\
3.9604	0.32447\\
3.9669	0.31383\\
3.9747	0.30319\\
3.9822	0.29255\\
3.9929	0.28191\\
4.0051	0.27128\\
4.0144	0.26064\\
4.0161	0.25\\
4.0433	0.23936\\
4.0578	0.22872\\
4.073	0.21809\\
4.0857	0.20745\\
4.0904	0.19681\\
4.1168	0.18617\\
4.1265	0.17553\\
4.127	0.16489\\
4.143	0.15426\\
4.1756	0.14362\\
4.2156	0.13298\\
4.3116	0.12234\\
4.3276	0.1117\\
4.3662	0.10106\\
4.4791	0.090426\\
4.4912	0.079787\\
4.517	0.069149\\
4.5743	0.058511\\
4.5914	0.047872\\
4.6138	0.037234\\
4.7898	0.026596\\
5.0876	0.015957\\
5.3101	0.0053191\\
};
\addlegendentry{\RealEmpiricCCDFLegend}

\addplot [color=black, mark repeat=9, mark phase=1, mark=+, mark options={solid, black}]
  table[row sep=crcr]{%
3.5608	1\\
3.5785	0.94954\\
3.5961	0.90162\\
3.6138	0.85612\\
3.6315	0.81292\\
3.6492	0.7719\\
3.6668	0.73295\\
3.6845	0.69596\\
3.7022	0.66084\\
3.7198	0.6275\\
3.7375	0.59583\\
3.7552	0.56576\\
3.7728	0.53721\\
3.7905	0.5101\\
3.8082	0.48436\\
3.8259	0.45992\\
3.8435	0.43671\\
3.8612	0.41468\\
3.8789	0.39375\\
3.8965	0.37388\\
3.9142	0.35501\\
3.9319	0.3371\\
3.9495	0.32009\\
3.9672	0.30394\\
3.9849	0.2886\\
4.0026	0.27404\\
4.0202	0.26021\\
4.0379	0.24708\\
4.0556	0.23461\\
4.0732	0.22277\\
4.0909	0.21153\\
4.1086	0.20085\\
4.1262	0.19072\\
4.1439	0.18109\\
4.1616	0.17196\\
4.1793	0.16328\\
4.1969	0.15504\\
4.2146	0.14722\\
4.2323	0.13979\\
4.2499	0.13273\\
4.2676	0.12603\\
4.2853	0.11967\\
4.3029	0.11364\\
4.3206	0.1079\\
4.3383	0.10246\\
4.356	0.097286\\
4.3736	0.092377\\
4.3913	0.087716\\
4.409	0.083289\\
4.4266	0.079086\\
4.4443	0.075095\\
4.462	0.071306\\
4.4796	0.067708\\
4.4973	0.064291\\
4.515	0.061047\\
4.5327	0.057966\\
4.5503	0.055041\\
4.568	0.052264\\
4.5857	0.049626\\
4.6033	0.047122\\
4.621	0.044744\\
4.6387	0.042486\\
4.6563	0.040342\\
4.674	0.038307\\
4.6917	0.036373\\
4.7094	0.034538\\
4.727	0.032795\\
4.7447	0.03114\\
4.7624	0.029569\\
4.78	0.028077\\
4.7977	0.02666\\
4.8154	0.025315\\
4.833	0.024037\\
4.8507	0.022824\\
4.8684	0.021672\\
4.8861	0.020579\\
4.9037	0.01954\\
4.9214	0.018554\\
4.9391	0.017618\\
4.9567	0.016729\\
4.9744	0.015885\\
4.9921	0.015083\\
5.0097	0.014322\\
5.0274	0.013599\\
5.0451	0.012913\\
5.0628	0.012261\\
5.0804	0.011643\\
5.0981	0.011055\\
5.1158	0.010497\\
5.1334	0.0099676\\
5.1511	0.0094646\\
5.1688	0.008987\\
5.1864	0.0085335\\
5.2041	0.0081029\\
5.2218	0.007694\\
5.2395	0.0073058\\
5.2571	0.0069371\\
5.2748	0.006587\\
5.2925	0.0062546\\
5.3101	0.005939\\
};
\addlegendentry{\RealParetoFitLegend}

\addplot [color=black, mark repeat=9, mark phase=1, mark=x, mark options={solid, black}]
  table[row sep=crcr]{%
3.5608	1.0005\\
3.5785	0.95316\\
3.5961	0.90742\\
3.6138	0.86328\\
3.6315	0.82076\\
3.6492	0.77986\\
3.6668	0.74057\\
3.6845	0.70289\\
3.7022	0.66682\\
3.7198	0.63232\\
3.7375	0.59938\\
3.7552	0.56796\\
3.7728	0.53805\\
3.7905	0.5096\\
3.8082	0.48257\\
3.8259	0.45692\\
3.8435	0.43262\\
3.8612	0.4096\\
3.8789	0.38783\\
3.8965	0.36726\\
3.9142	0.34784\\
3.9319	0.32951\\
3.9495	0.31222\\
3.9672	0.29594\\
3.9849	0.28059\\
4.0026	0.26614\\
4.0202	0.25254\\
4.0379	0.23974\\
4.0556	0.22769\\
4.0732	0.21634\\
4.0909	0.20566\\
4.1086	0.19561\\
4.1262	0.18613\\
4.1439	0.1772\\
4.1616	0.16878\\
4.1793	0.16082\\
4.1969	0.15331\\
4.2146	0.14621\\
4.2323	0.1395\\
4.2499	0.13313\\
4.2676	0.1271\\
4.2853	0.12137\\
4.3029	0.11592\\
4.3206	0.11074\\
4.3383	0.10581\\
4.356	0.1011\\
4.3736	0.09661\\
4.3913	0.092318\\
4.409	0.088212\\
4.4266	0.084281\\
4.4443	0.080514\\
4.462	0.076901\\
4.4796	0.073434\\
4.4973	0.070104\\
4.515	0.066905\\
4.5327	0.06383\\
4.5503	0.060873\\
4.568	0.058029\\
4.5857	0.055293\\
4.6033	0.052661\\
4.621	0.050128\\
4.6387	0.047691\\
4.6563	0.045347\\
4.674	0.043093\\
4.6917	0.040926\\
4.7094	0.038844\\
4.727	0.036843\\
4.7447	0.034921\\
4.7624	0.033077\\
4.78	0.031309\\
4.7977	0.029613\\
4.8154	0.027989\\
4.833	0.026433\\
4.8507	0.024946\\
4.8684	0.023523\\
4.8861	0.022165\\
4.9037	0.020868\\
4.9214	0.019632\\
4.9391	0.018454\\
4.9567	0.017332\\
4.9744	0.016265\\
4.9921	0.015251\\
5.0097	0.014288\\
5.0274	0.013375\\
5.0451	0.01251\\
5.0628	0.01169\\
5.0804	0.010915\\
5.0981	0.010182\\
5.1158	0.0094906\\
5.1334	0.0088382\\
5.1511	0.0082235\\
5.1688	0.0076448\\
5.1864	0.0071005\\
5.2041	0.0065892\\
5.2218	0.0061093\\
5.2395	0.0056592\\
5.2571	0.0052377\\
5.2748	0.0048432\\
5.2925	0.0044744\\
5.3101	0.0041299\\
};
\addlegendentry{\RealBLNFitLegend}

      % \input{../../../../.tex/20/KS/SizeDistributionFittingsfig4.tex}
  
      \nextgroupplot[%
        xmin=3.598,xmax=5.4115,
        % xmin=3.4,xmax=5.4,
        % ymin=0.00088795,ymax=1.035,
      ] % This file was created by matlab2tikz.
%
\definecolor{mycolor1}{rgb}{0.00000,0.44706,0.69804}%
\definecolor{mycolor2}{rgb}{0.00000,0.44700,0.74100}%
%
\addplot[only marks, mark=*, mark size=1.1180pt, color=mycolor1, fill=mycolor1, opacity=0.60, draw=none, mark options={draw=none,line width=0pt}] table[row sep=crcr]{%
x	y\\
3.5893	0.99468\\
3.5929	0.98404\\
3.5962	0.9734\\
3.6001	0.96277\\
3.6042	0.95213\\
3.6081	0.94149\\
3.6124	0.93085\\
3.6162	0.92021\\
3.6204	0.90957\\
3.625	0.89894\\
3.6291	0.8883\\
3.6335	0.87766\\
3.6385	0.86702\\
3.6436	0.85638\\
3.6481	0.84574\\
3.6533	0.83511\\
3.6583	0.82447\\
3.6634	0.81383\\
3.669	0.80319\\
3.6741	0.79255\\
3.6786	0.78191\\
3.683	0.77128\\
3.6891	0.76064\\
3.6957	0.75\\
3.7019	0.73936\\
3.7088	0.72872\\
3.7138	0.71809\\
3.7192	0.70745\\
3.7256	0.69681\\
3.7321	0.68617\\
3.7405	0.67553\\
3.7483	0.66489\\
3.7548	0.65426\\
3.763	0.64362\\
3.7715	0.63298\\
3.7797	0.62234\\
3.786	0.6117\\
3.7935	0.60106\\
3.8025	0.59043\\
3.8109	0.57979\\
3.819	0.56915\\
3.8255	0.55851\\
3.8327	0.54787\\
3.8385	0.53723\\
3.8461	0.5266\\
3.852	0.51596\\
3.86	0.50532\\
3.866	0.49468\\
3.8727	0.48404\\
3.8788	0.4734\\
3.8839	0.46277\\
3.8896	0.45213\\
3.8958	0.44149\\
3.9012	0.43085\\
3.9081	0.42021\\
3.9155	0.40957\\
3.9216	0.39894\\
3.928	0.3883\\
3.9348	0.37766\\
3.9415	0.36702\\
3.9492	0.35638\\
3.9579	0.34574\\
3.9692	0.33511\\
3.9828	0.32447\\
3.9977	0.31383\\
4.0161	0.30319\\
4.0335	0.29255\\
4.0503	0.28191\\
4.0662	0.27128\\
4.0836	0.26064\\
4.0988	0.25\\
4.1112	0.23936\\
4.1243	0.22872\\
4.1342	0.21809\\
4.1454	0.20745\\
4.1576	0.19681\\
4.1705	0.18617\\
4.1844	0.17553\\
4.2007	0.16489\\
4.2194	0.15426\\
4.24	0.14362\\
4.2662	0.13298\\
4.2921	0.12234\\
4.3201	0.1117\\
4.3525	0.10106\\
4.3736	0.090426\\
4.3895	0.079787\\
4.4094	0.069149\\
4.4358	0.058511\\
4.5276	0.047872\\
4.5794	0.037234\\
4.7208	0.026596\\
4.7905	0.015957\\
5.1574	0.0053191\\
};
\addlegendentry{\ExactEmpiricCCDFLegend}

\addplot [color=mycolor1, only marks, every error bar/.append style={opacity=0.30}, mark=*, mark size=0pt, draw=none, forget plot]
 plot [error bars/.cd, x dir=both, x explicit, error bar style={line width=1pt, color=mycolor1}, error mark options={mark=none,mark size=0pt}]
 table[row sep=crcr, x error plus index=2, x error minus index=3]{%
3.5893	0.99468	0.0020617	0.0020617\\
3.5929	0.98404	0.0020269	0.0020269\\
3.5962	0.9734	0.0020561	0.0020561\\
3.6001	0.96277	0.0019457	0.0019457\\
3.6042	0.95213	0.0019054	0.0019054\\
3.6081	0.94149	0.0018168	0.0018168\\
3.6124	0.93085	0.0018416	0.0018416\\
3.6162	0.92021	0.0018589	0.0018589\\
3.6204	0.90957	0.0018898	0.0018898\\
3.625	0.89894	0.0020758	0.0020758\\
3.6291	0.8883	0.0019633	0.0019633\\
3.6335	0.87766	0.0019186	0.0019186\\
3.6385	0.86702	0.0019754	0.0019754\\
3.6436	0.85638	0.002189	0.002189\\
3.6481	0.84574	0.0020486	0.0020486\\
3.6533	0.83511	0.0020882	0.0020882\\
3.6583	0.82447	0.0020115	0.0020115\\
3.6634	0.81383	0.0020684	0.0020684\\
3.669	0.80319	0.0022894	0.0022894\\
3.6741	0.79255	0.0022671	0.0022671\\
3.6786	0.78191	0.0021921	0.0021921\\
3.683	0.77128	0.0021802	0.0021802\\
3.6891	0.76064	0.0021046	0.0021046\\
3.6957	0.75	0.0020569	0.0020569\\
3.7019	0.73936	0.0021172	0.0021172\\
3.7088	0.72872	0.0020608	0.0020608\\
3.7138	0.71809	0.0020375	0.0020375\\
3.7192	0.70745	0.0020778	0.0020778\\
3.7256	0.69681	0.0020958	0.0020958\\
3.7321	0.68617	0.0022399	0.0022399\\
3.7405	0.67553	0.0023519	0.0023519\\
3.7483	0.66489	0.0024659	0.0024659\\
3.7548	0.65426	0.0025499	0.0025499\\
3.763	0.64362	0.0027554	0.0027554\\
3.7715	0.63298	0.002923	0.002923\\
3.7797	0.62234	0.0028095	0.0028095\\
3.786	0.6117	0.0027316	0.0027316\\
3.7935	0.60106	0.002762	0.002762\\
3.8025	0.59043	0.0028749	0.0028749\\
3.8109	0.57979	0.003014	0.003014\\
3.819	0.56915	0.0027807	0.0027807\\
3.8255	0.55851	0.0025636	0.0025636\\
3.8327	0.54787	0.002599	0.002599\\
3.8385	0.53723	0.0024753	0.0024753\\
3.8461	0.5266	0.0023333	0.0023333\\
3.852	0.51596	0.0022884	0.0022884\\
3.86	0.50532	0.0022617	0.0022617\\
3.866	0.49468	0.0022885	0.0022885\\
3.8727	0.48404	0.0022411	0.0022411\\
3.8788	0.4734	0.0022106	0.0022106\\
3.8839	0.46277	0.0022586	0.0022586\\
3.8896	0.45213	0.0024622	0.0024622\\
3.8958	0.44149	0.0024482	0.0024482\\
3.9012	0.43085	0.0023713	0.0023713\\
3.9081	0.42021	0.0023255	0.0023255\\
3.9155	0.40957	0.0023465	0.0023465\\
3.9216	0.39894	0.0023832	0.0023832\\
3.928	0.3883	0.002503	0.002503\\
3.9348	0.37766	0.0025268	0.0025268\\
3.9415	0.36702	0.0027627	0.0027627\\
3.9492	0.35638	0.0027797	0.0027797\\
3.9579	0.34574	0.0028551	0.0028551\\
3.9692	0.33511	0.0029051	0.0029051\\
3.9828	0.32447	0.0031558	0.0031558\\
3.9977	0.31383	0.003395	0.003395\\
4.0161	0.30319	0.0037587	0.0037587\\
4.0335	0.29255	0.0037678	0.0037678\\
4.0503	0.28191	0.0031906	0.0031906\\
4.0662	0.27128	0.0031187	0.0031187\\
4.0836	0.26064	0.003391	0.003391\\
4.0988	0.25	0.0031882	0.0031882\\
4.1112	0.23936	0.0033285	0.0033285\\
4.1243	0.22872	0.0032658	0.0032658\\
4.1342	0.21809	0.0030368	0.0030368\\
4.1454	0.20745	0.0028444	0.0028444\\
4.1576	0.19681	0.0028028	0.0028028\\
4.1705	0.18617	0.0031669	0.0031669\\
4.1844	0.17553	0.0030161	0.0030161\\
4.2007	0.16489	0.0030855	0.0030855\\
4.2194	0.15426	0.0036315	0.0036315\\
4.24	0.14362	0.0036555	0.0036555\\
4.2662	0.13298	0.0036608	0.0036608\\
4.2921	0.12234	0.0040354	0.0040354\\
4.3201	0.1117	0.004327	0.004327\\
4.3525	0.10106	0.0037223	0.0037223\\
4.3736	0.090426	0.0034217	0.0034217\\
4.3895	0.079787	0.0036394	0.0036394\\
4.4094	0.069149	0.0039009	0.0039009\\
4.4358	0.058511	0.0043355	0.0043355\\
4.5276	0.047872	0.0054577	0.0054577\\
4.5794	0.037234	0.0042494	0.0042494\\
4.7208	0.026596	0.0050438	0.0050438\\
4.7905	0.015957	0.0050741	0.0050741\\
5.1574	0.0053191	0.0034798	0.0034798\\
};
\addplot [color=mycolor1, mark repeat=9, mark phase=1, mark=+, mark options={solid, mycolor1}]
  table[row sep=crcr]{%
3.5888	0.98701\\
3.6051	0.95142\\
3.6215	0.90753\\
3.6378	0.86528\\
3.6541	0.82501\\
3.6704	0.78661\\
3.6868	0.75\\
3.7031	0.71509\\
3.7194	0.68182\\
3.7357	0.65009\\
3.752	0.61984\\
3.7684	0.59099\\
3.7847	0.56349\\
3.801	0.53727\\
3.8173	0.51228\\
3.8337	0.48844\\
3.85	0.46572\\
3.8663	0.44405\\
3.8826	0.4234\\
3.899	0.4037\\
3.9153	0.38492\\
3.9316	0.36702\\
3.9479	0.34995\\
3.9642	0.33367\\
3.9806	0.31815\\
3.9969	0.30335\\
4.0132	0.28925\\
4.0295	0.2758\\
4.0459	0.26297\\
4.0622	0.25074\\
4.0785	0.23908\\
4.0948	0.22797\\
4.1111	0.21737\\
4.1275	0.20726\\
4.1438	0.19763\\
4.1601	0.18844\\
4.1764	0.17968\\
4.1928	0.17132\\
4.2091	0.16336\\
4.2254	0.15577\\
4.2417	0.14853\\
4.2581	0.14162\\
4.2744	0.13504\\
4.2907	0.12877\\
4.307	0.12278\\
4.3233	0.11708\\
4.3397	0.11163\\
4.356	0.10645\\
4.3723	0.1015\\
4.3886	0.096785\\
4.405	0.092288\\
4.4213	0.088\\
4.4376	0.083912\\
4.4539	0.080013\\
4.4703	0.076296\\
4.4866	0.072752\\
4.5029	0.069373\\
4.5192	0.06615\\
4.5355	0.063078\\
4.5519	0.060148\\
4.5682	0.057354\\
4.5845	0.05469\\
4.6008	0.052151\\
4.6172	0.049729\\
4.6335	0.047419\\
4.6498	0.045217\\
4.6661	0.043118\\
4.6824	0.041116\\
4.6988	0.039207\\
4.7151	0.037386\\
4.7314	0.03565\\
4.7477	0.033995\\
4.7641	0.032417\\
4.7804	0.030912\\
4.7967	0.029477\\
4.813	0.028109\\
4.8294	0.026804\\
4.8457	0.02556\\
4.862	0.024374\\
4.8783	0.023243\\
4.8946	0.022164\\
4.911	0.021135\\
4.9273	0.020154\\
4.9436	0.019219\\
4.9599	0.018327\\
4.9763	0.017477\\
4.9926	0.016666\\
5.0089	0.015893\\
5.0252	0.015155\\
5.0415	0.014452\\
5.0579	0.013782\\
5.0742	0.013142\\
5.0905	0.012533\\
5.1068	0.011951\\
5.1232	0.011397\\
5.1395	0.010868\\
5.1558	0.010364\\
5.1721	0.0098836\\
5.1885	0.0094252\\
5.2048	0.0089881\\
};
\addlegendentry{\ExactParetoFitLegend}


\addplot[area legend, draw=none, fill=mycolor1, fill opacity=0.15, postaction={pattern={Lines[angle=90,distance={5pt},line width={0.5pt}]}, pattern color=mycolor1!60}, forget plot]
table[row sep=crcr] {%
x	y\\
3.5888	0.99033\\
3.6051	0.9567\\
3.6215	0.91248\\
3.6378	0.86978\\
3.6541	0.82908\\
3.6704	0.79029\\
3.6868	0.75332\\
3.7031	0.71808\\
3.7194	0.68449\\
3.7357	0.65247\\
3.752	0.62196\\
3.7684	0.59288\\
3.7847	0.56516\\
3.801	0.53875\\
3.8173	0.51357\\
3.8337	0.48958\\
3.85	0.46671\\
3.8663	0.44493\\
3.8826	0.42417\\
3.899	0.4044\\
3.9153	0.38556\\
3.9316	0.36762\\
3.9479	0.35053\\
3.9642	0.33424\\
3.9806	0.31873\\
3.9969	0.30395\\
4.0132	0.28986\\
4.0295	0.27643\\
4.0459	0.26363\\
4.0622	0.25142\\
4.0785	0.23979\\
4.0948	0.22869\\
4.1111	0.21811\\
4.1275	0.20802\\
4.1438	0.19839\\
4.1601	0.18922\\
4.1764	0.18047\\
4.1928	0.17212\\
4.2091	0.16416\\
4.2254	0.15657\\
4.2417	0.14933\\
4.2581	0.14243\\
4.2744	0.13584\\
4.2907	0.12956\\
4.307	0.12357\\
4.3233	0.11786\\
4.3397	0.11241\\
4.356	0.10722\\
4.3723	0.10226\\
4.3886	0.097537\\
4.405	0.09303\\
4.4213	0.088731\\
4.4376	0.084631\\
4.4539	0.08072\\
4.4703	0.07699\\
4.4866	0.073433\\
4.5029	0.07004\\
4.5192	0.066805\\
4.5355	0.063718\\
4.5519	0.060775\\
4.5682	0.057968\\
4.5845	0.05529\\
4.6008	0.052736\\
4.6172	0.0503\\
4.6335	0.047977\\
4.6498	0.045761\\
4.6661	0.043648\\
4.6824	0.041632\\
4.6988	0.039709\\
4.7151	0.037876\\
4.7314	0.036127\\
4.7477	0.034459\\
4.7641	0.032868\\
4.7804	0.03135\\
4.7967	0.029903\\
4.813	0.028522\\
4.8294	0.027205\\
4.8457	0.025949\\
4.862	0.024752\\
4.8783	0.023609\\
4.8946	0.022519\\
4.911	0.02148\\
4.9273	0.020488\\
4.9436	0.019543\\
4.9599	0.018641\\
4.9763	0.01778\\
4.9926	0.01696\\
5.0089	0.016177\\
5.0252	0.015431\\
5.0415	0.014719\\
5.0579	0.01404\\
5.0742	0.013392\\
5.0905	0.012774\\
5.1068	0.012184\\
5.1232	0.011622\\
5.1395	0.011086\\
5.1558	0.010575\\
5.1721	0.010087\\
5.1885	0.0096216\\
5.2048	0.0091777\\
5.2048	0.0087984\\
5.1885	0.0092288\\
5.1721	0.0096803\\
5.1558	0.010154\\
5.1395	0.010651\\
5.1232	0.011172\\
5.1068	0.011718\\
5.0905	0.012292\\
5.0742	0.012893\\
5.0579	0.013524\\
5.0415	0.014186\\
5.0252	0.01488\\
5.0089	0.015608\\
4.9926	0.016372\\
4.9763	0.017173\\
4.9599	0.018014\\
4.9436	0.018896\\
4.9273	0.019821\\
4.911	0.020791\\
4.8946	0.021808\\
4.8783	0.022876\\
4.862	0.023996\\
4.8457	0.025171\\
4.8294	0.026403\\
4.813	0.027696\\
4.7967	0.029052\\
4.7804	0.030474\\
4.7641	0.031967\\
4.7477	0.033532\\
4.7314	0.035174\\
4.7151	0.036897\\
4.6988	0.038704\\
4.6824	0.040599\\
4.6661	0.042588\\
4.6498	0.044673\\
4.6335	0.046862\\
4.6172	0.049157\\
4.6008	0.051565\\
4.5845	0.054091\\
4.5682	0.056741\\
4.5519	0.059521\\
4.5355	0.062437\\
4.5192	0.065496\\
4.5029	0.068705\\
4.4866	0.072071\\
4.4703	0.075602\\
4.4539	0.079307\\
4.4376	0.083193\\
4.4213	0.08727\\
4.405	0.091546\\
4.3886	0.096033\\
4.3723	0.10074\\
4.356	0.10568\\
4.3397	0.11086\\
4.3233	0.11629\\
4.307	0.12199\\
4.2907	0.12797\\
4.2744	0.13424\\
4.2581	0.14082\\
4.2417	0.14772\\
4.2254	0.15496\\
4.2091	0.16256\\
4.1928	0.17053\\
4.1764	0.17889\\
4.1601	0.18766\\
4.1438	0.19686\\
4.1275	0.20651\\
4.1111	0.21663\\
4.0948	0.22724\\
4.0785	0.23838\\
4.0622	0.25006\\
4.0459	0.26231\\
4.0295	0.27516\\
4.0132	0.28863\\
3.9969	0.30276\\
3.9806	0.31757\\
3.9642	0.33309\\
3.9479	0.34937\\
3.9316	0.36642\\
3.9153	0.38428\\
3.899	0.403\\
3.8826	0.42262\\
3.8663	0.44318\\
3.85	0.46472\\
3.8337	0.48731\\
3.8173	0.51098\\
3.801	0.5358\\
3.7847	0.56183\\
3.7684	0.58911\\
3.752	0.61771\\
3.7357	0.6477\\
3.7194	0.67914\\
3.7031	0.71211\\
3.6868	0.74668\\
3.6704	0.78293\\
3.6541	0.82093\\
3.6378	0.86079\\
3.6215	0.90257\\
3.6051	0.94613\\
3.5888	0.98368\\
}--cycle;
\addplot [color=mycolor1, mark repeat=9, mark phase=1, mark=x, mark options={solid, mycolor1}]
  table[row sep=crcr]{%
3.5888	1.0328\\
3.6051	0.98511\\
3.6215	0.93959\\
3.6378	0.89617\\
3.6541	0.85479\\
3.6704	0.81541\\
3.6868	0.77795\\
3.7031	0.74235\\
3.7194	0.70852\\
3.7357	0.67638\\
3.752	0.64586\\
3.7684	0.61687\\
3.7847	0.58932\\
3.801	0.56313\\
3.8173	0.53822\\
3.8337	0.5145\\
3.85	0.49192\\
3.8663	0.47038\\
3.8826	0.44982\\
3.899	0.43017\\
3.9153	0.41138\\
3.9316	0.39339\\
3.9479	0.37614\\
3.9642	0.35959\\
3.9806	0.34369\\
3.9969	0.32841\\
4.0132	0.3137\\
4.0295	0.29953\\
4.0459	0.28589\\
4.0622	0.27273\\
4.0785	0.26004\\
4.0948	0.2478\\
4.1111	0.23599\\
4.1275	0.2246\\
4.1438	0.21361\\
4.1601	0.203\\
4.1764	0.19278\\
4.1928	0.18293\\
4.2091	0.17344\\
4.2254	0.1643\\
4.2417	0.1555\\
4.2581	0.14705\\
4.2744	0.13892\\
4.2907	0.13113\\
4.307	0.12365\\
4.3233	0.11649\\
4.3397	0.10963\\
4.356	0.10307\\
4.3723	0.096812\\
4.3886	0.090838\\
4.405	0.085145\\
4.4213	0.079725\\
4.4376	0.074572\\
4.4539	0.069679\\
4.4703	0.065037\\
4.4866	0.06064\\
4.5029	0.056478\\
4.5192	0.052546\\
4.5355	0.048834\\
4.5519	0.045335\\
4.5682	0.04204\\
4.5845	0.038942\\
4.6008	0.036032\\
4.6172	0.033302\\
4.6335	0.030745\\
4.6498	0.028353\\
4.6661	0.026117\\
4.6824	0.024031\\
4.6988	0.022086\\
4.7151	0.020276\\
4.7314	0.018593\\
4.7477	0.017031\\
4.7641	0.015582\\
4.7804	0.01424\\
4.7967	0.012998\\
4.813	0.011852\\
4.8294	0.010794\\
4.8457	0.0098191\\
4.862	0.008922\\
4.8783	0.0080976\\
4.8946	0.0073408\\
4.911	0.0066471\\
4.9273	0.006012\\
4.9436	0.0054312\\
4.9599	0.0049008\\
4.9763	0.0044171\\
4.9926	0.0039765\\
5.0089	0.0035757\\
5.0252	0.0032115\\
5.0415	0.002881\\
5.0579	0.0025815\\
5.0742	0.0023104\\
5.0905	0.0020654\\
5.1068	0.0018442\\
5.1232	0.0016447\\
5.1395	0.0014651\\
5.1558	0.0013035\\
5.1721	0.0011584\\
5.1885	0.0010283\\
5.2048	0.00091167\\
};
\addlegendentry{\ExactBLNFitLegend}


\addplot[area legend, draw=none, fill=mycolor1, fill opacity=0.15, postaction={pattern={Lines[angle=45,distance={5pt},line width={0.5pt}]}, pattern color=mycolor1!60}, forget plot]
table[row sep=crcr] {%
x	y\\
3.5888	1.035\\
3.6051	0.98722\\
3.6215	0.94159\\
3.6378	0.89807\\
3.6541	0.85661\\
3.6704	0.81715\\
3.6868	0.77963\\
3.7031	0.74397\\
3.7194	0.71009\\
3.7357	0.67792\\
3.752	0.64737\\
3.7684	0.61835\\
3.7847	0.59078\\
3.801	0.56457\\
3.8173	0.53964\\
3.8337	0.51592\\
3.85	0.49331\\
3.8663	0.47175\\
3.8826	0.45117\\
3.899	0.4315\\
3.9153	0.41267\\
3.9316	0.39465\\
3.9479	0.37736\\
3.9642	0.36077\\
3.9806	0.34482\\
3.9969	0.32949\\
4.0132	0.31473\\
4.0295	0.30052\\
4.0459	0.28682\\
4.0622	0.27361\\
4.0785	0.26087\\
4.0948	0.24858\\
4.1111	0.23672\\
4.1275	0.22528\\
4.1438	0.21424\\
4.1601	0.2036\\
4.1764	0.19333\\
4.1928	0.18344\\
4.2091	0.17392\\
4.2254	0.16475\\
4.2417	0.15593\\
4.2581	0.14746\\
4.2744	0.13932\\
4.2907	0.13151\\
4.307	0.12402\\
4.3233	0.11685\\
4.3397	0.10998\\
4.356	0.10342\\
4.3723	0.097158\\
4.3886	0.09118\\
4.405	0.085484\\
4.4213	0.080062\\
4.4376	0.074906\\
4.4539	0.070009\\
4.4703	0.065364\\
4.4866	0.060962\\
4.5029	0.056796\\
4.5192	0.052858\\
4.5355	0.04914\\
4.5519	0.045633\\
4.5682	0.042331\\
4.5845	0.039225\\
4.6008	0.036307\\
4.6172	0.033569\\
4.6335	0.031003\\
4.6498	0.028602\\
4.6661	0.026357\\
4.6824	0.024261\\
4.6988	0.022307\\
4.7151	0.020487\\
4.7314	0.018794\\
4.7477	0.017222\\
4.7641	0.015764\\
4.7804	0.014412\\
4.7967	0.013162\\
4.813	0.012006\\
4.8294	0.010939\\
4.8457	0.009956\\
4.862	0.0090507\\
4.8783	0.0082183\\
4.8946	0.0074539\\
4.911	0.0067528\\
4.9273	0.0061106\\
4.9436	0.005523\\
4.9599	0.0049862\\
4.9763	0.0044964\\
4.9926	0.00405\\
5.0089	0.0036436\\
5.0252	0.0032743\\
5.0415	0.0029389\\
5.0579	0.0026348\\
5.0742	0.0023595\\
5.0905	0.0021104\\
5.1068	0.0018854\\
5.1232	0.0016825\\
5.1395	0.0014996\\
5.1558	0.001335\\
5.1721	0.0011871\\
5.1885	0.0010544\\
5.2048	0.00093538\\
5.2048	0.00088795\\
5.1885	0.0010022\\
5.1721	0.0011297\\
5.1558	0.001272\\
5.1395	0.0014306\\
5.1232	0.0016069\\
5.1068	0.0018029\\
5.0905	0.0020203\\
5.0742	0.0022614\\
5.0579	0.0025281\\
5.0415	0.0028231\\
5.0252	0.0031487\\
5.0089	0.0035077\\
4.9926	0.003903\\
4.9763	0.0043379\\
4.9599	0.0048155\\
4.9436	0.0053394\\
4.9273	0.0059134\\
4.911	0.0065415\\
4.8946	0.0072278\\
4.8783	0.0079769\\
4.862	0.0087933\\
4.8457	0.0096821\\
4.8294	0.010648\\
4.813	0.011697\\
4.7967	0.012835\\
4.7804	0.014067\\
4.7641	0.0154\\
4.7477	0.016839\\
4.7314	0.018392\\
4.7151	0.020066\\
4.6988	0.021866\\
4.6824	0.023801\\
4.6661	0.025878\\
4.6498	0.028104\\
4.6335	0.030487\\
4.6172	0.033035\\
4.6008	0.035756\\
4.5845	0.038658\\
4.5682	0.041748\\
4.5519	0.045036\\
4.5355	0.048529\\
4.5192	0.052234\\
4.5029	0.056161\\
4.4866	0.060317\\
4.4703	0.064711\\
4.4539	0.069349\\
4.4376	0.074239\\
4.4213	0.079389\\
4.405	0.084805\\
4.3886	0.090496\\
4.3723	0.096467\\
4.356	0.10272\\
4.3397	0.10928\\
4.3233	0.11613\\
4.307	0.12328\\
4.2907	0.13075\\
4.2744	0.13853\\
4.2581	0.14664\\
4.2417	0.15507\\
4.2254	0.16384\\
4.2091	0.17295\\
4.1928	0.18241\\
4.1764	0.19223\\
4.1601	0.20241\\
4.1438	0.21297\\
4.1275	0.22392\\
4.1111	0.23526\\
4.0948	0.24702\\
4.0785	0.25921\\
4.0622	0.27185\\
4.0459	0.28495\\
4.0295	0.29855\\
4.0132	0.31266\\
3.9969	0.32732\\
3.9806	0.34256\\
3.9642	0.35842\\
3.9479	0.37493\\
3.9316	0.39214\\
3.9153	0.41009\\
3.899	0.42885\\
3.8826	0.44847\\
3.8663	0.469\\
3.85	0.49052\\
3.8337	0.51309\\
3.8173	0.53679\\
3.801	0.56168\\
3.7847	0.58785\\
3.7684	0.61538\\
3.752	0.64435\\
3.7357	0.67485\\
3.7194	0.70694\\
3.7031	0.74073\\
3.6868	0.77628\\
3.6704	0.81367\\
3.6541	0.85297\\
3.6378	0.89426\\
3.6215	0.93759\\
3.6051	0.98301\\
3.5888	1.0306\\
}--cycle;
      % \input{../../../../.tex/20/KS/SizeDistributionFittingsfig5.tex}
  
      \nextgroupplot[%
        xmin=3.6008,xmax=5.4115,
        % xmin=3.4,xmax=5.4,
        % ymin=0.00088167,ymax=1.0352,
      ] % This file was created by matlab2tikz.
%
\definecolor{mycolor1}{rgb}{0.83529,0.36863,0.00000}%
\definecolor{mycolor2}{rgb}{0.00000,0.44700,0.74100}%
%
\addplot[only marks, mark=*, mark size=1.1180pt, color=mycolor1, fill=mycolor1, opacity=0.60, draw=none, mark options={draw=none,line width=0pt}] table[row sep=crcr]{%
x	y\\
3.5886	0.99468\\
3.5916	0.98404\\
3.5952	0.9734\\
3.5993	0.96277\\
3.6029	0.95213\\
3.6065	0.94149\\
3.6092	0.93085\\
3.6129	0.92021\\
3.6165	0.90957\\
3.6212	0.89894\\
3.6249	0.8883\\
3.6294	0.87766\\
3.6332	0.86702\\
3.6375	0.85638\\
3.642	0.84574\\
3.647	0.83511\\
3.6521	0.82447\\
3.6574	0.81383\\
3.6623	0.80319\\
3.6677	0.79255\\
3.673	0.78191\\
3.6783	0.77128\\
3.6852	0.76064\\
3.6917	0.75\\
3.6975	0.73936\\
3.7051	0.72872\\
3.7118	0.71809\\
3.7198	0.70745\\
3.7274	0.69681\\
3.7344	0.68617\\
3.7419	0.67553\\
3.7498	0.66489\\
3.7583	0.65426\\
3.7667	0.64362\\
3.7752	0.63298\\
3.7842	0.62234\\
3.7915	0.6117\\
3.7992	0.60106\\
3.8062	0.59043\\
3.8131	0.57979\\
3.8207	0.56915\\
3.8275	0.55851\\
3.835	0.54787\\
3.8424	0.53723\\
3.849	0.5266\\
3.8549	0.51596\\
3.8615	0.50532\\
3.8671	0.49468\\
3.8725	0.48404\\
3.8792	0.4734\\
3.8855	0.46277\\
3.8914	0.45213\\
3.8973	0.44149\\
3.9044	0.43085\\
3.9103	0.42021\\
3.9165	0.40957\\
3.9223	0.39894\\
3.9291	0.3883\\
3.9357	0.37766\\
3.9442	0.36702\\
3.952	0.35638\\
3.9612	0.34574\\
3.9706	0.33511\\
3.9834	0.32447\\
4.0008	0.31383\\
4.0156	0.30319\\
4.0357	0.29255\\
4.0537	0.28191\\
4.067	0.27128\\
4.0828	0.26064\\
4.0979	0.25\\
4.1112	0.23936\\
4.1231	0.22872\\
4.1341	0.21809\\
4.1468	0.20745\\
4.1597	0.19681\\
4.1722	0.18617\\
4.1856	0.17553\\
4.1984	0.16489\\
4.2174	0.15426\\
4.2391	0.14362\\
4.264	0.13298\\
4.2878	0.12234\\
4.3125	0.1117\\
4.3462	0.10106\\
4.3689	0.090426\\
4.389	0.079787\\
4.4095	0.069149\\
4.4462	0.058511\\
4.5305	0.047872\\
4.5839	0.037234\\
4.7235	0.026596\\
4.7846	0.015957\\
5.1606	0.0053191\\
};
\addlegendentry{\ApproxEmpiricCCDFLegend}

\addplot [color=mycolor1, only marks, every error bar/.append style={opacity=0.30}, mark=*, mark size=0pt, draw=none, forget plot]
 plot [error bars/.cd, x dir=both, x explicit, error bar style={line width=1pt, color=mycolor1}, error mark options={mark=none,mark size=0pt}]
 table[row sep=crcr, x error plus index=2, x error minus index=3]{%
3.5886	0.99468	0.0016453	0.0016453\\
3.5916	0.98404	0.0015936	0.0015936\\
3.5952	0.9734	0.0016864	0.0016864\\
3.5993	0.96277	0.001728	0.001728\\
3.6029	0.95213	0.0018264	0.0018264\\
3.6065	0.94149	0.0018754	0.0018754\\
3.6092	0.93085	0.0018953	0.0018953\\
3.6129	0.92021	0.0019758	0.0019758\\
3.6165	0.90957	0.0020323	0.0020323\\
3.6212	0.89894	0.0020596	0.0020596\\
3.6249	0.8883	0.0020847	0.0020847\\
3.6294	0.87766	0.0019316	0.0019316\\
3.6332	0.86702	0.0019421	0.0019421\\
3.6375	0.85638	0.002004	0.002004\\
3.642	0.84574	0.0020857	0.0020857\\
3.647	0.83511	0.0020635	0.0020635\\
3.6521	0.82447	0.0021601	0.0021601\\
3.6574	0.81383	0.0023654	0.0023654\\
3.6623	0.80319	0.0024984	0.0024984\\
3.6677	0.79255	0.0025274	0.0025274\\
3.673	0.78191	0.0022797	0.0022797\\
3.6783	0.77128	0.0023897	0.0023897\\
3.6852	0.76064	0.0024132	0.0024132\\
3.6917	0.75	0.0023622	0.0023622\\
3.6975	0.73936	0.0023958	0.0023958\\
3.7051	0.72872	0.002518	0.002518\\
3.7118	0.71809	0.0026086	0.0026086\\
3.7198	0.70745	0.0025221	0.0025221\\
3.7274	0.69681	0.0027598	0.0027598\\
3.7344	0.68617	0.0027338	0.0027338\\
3.7419	0.67553	0.0028359	0.0028359\\
3.7498	0.66489	0.0027646	0.0027646\\
3.7583	0.65426	0.0025985	0.0025985\\
3.7667	0.64362	0.0027915	0.0027915\\
3.7752	0.63298	0.0027864	0.0027864\\
3.7842	0.62234	0.0026958	0.0026958\\
3.7915	0.6117	0.0025913	0.0025913\\
3.7992	0.60106	0.0026056	0.0026056\\
3.8062	0.59043	0.0024877	0.0024877\\
3.8131	0.57979	0.0023548	0.0023548\\
3.8207	0.56915	0.0024251	0.0024251\\
3.8275	0.55851	0.0024283	0.0024283\\
3.835	0.54787	0.0023464	0.0023464\\
3.8424	0.53723	0.0024311	0.0024311\\
3.849	0.5266	0.0024573	0.0024573\\
3.8549	0.51596	0.0025555	0.0025555\\
3.8615	0.50532	0.0023886	0.0023886\\
3.8671	0.49468	0.0022532	0.0022532\\
3.8725	0.48404	0.0023285	0.0023285\\
3.8792	0.4734	0.0024835	0.0024835\\
3.8855	0.46277	0.0025563	0.0025563\\
3.8914	0.45213	0.0025453	0.0025453\\
3.8973	0.44149	0.0024626	0.0024626\\
3.9044	0.43085	0.0024371	0.0024371\\
3.9103	0.42021	0.0024808	0.0024808\\
3.9165	0.40957	0.0024173	0.0024173\\
3.9223	0.39894	0.0025136	0.0025136\\
3.9291	0.3883	0.0026601	0.0026601\\
3.9357	0.37766	0.0026816	0.0026816\\
3.9442	0.36702	0.002656	0.002656\\
3.952	0.35638	0.0027536	0.0027536\\
3.9612	0.34574	0.0027851	0.0027851\\
3.9706	0.33511	0.0028519	0.0028519\\
3.9834	0.32447	0.0028767	0.0028767\\
4.0008	0.31383	0.0031239	0.0031239\\
4.0156	0.30319	0.0032166	0.0032166\\
4.0357	0.29255	0.0032434	0.0032434\\
4.0537	0.28191	0.003496	0.003496\\
4.067	0.27128	0.00362	0.00362\\
4.0828	0.26064	0.0034719	0.0034719\\
4.0979	0.25	0.0034767	0.0034767\\
4.1112	0.23936	0.0030726	0.0030726\\
4.1231	0.22872	0.0029265	0.0029265\\
4.1341	0.21809	0.0030447	0.0030447\\
4.1468	0.20745	0.0027872	0.0027872\\
4.1597	0.19681	0.0030099	0.0030099\\
4.1722	0.18617	0.0034031	0.0034031\\
4.1856	0.17553	0.0033091	0.0033091\\
4.1984	0.16489	0.0037203	0.0037203\\
4.2174	0.15426	0.0036906	0.0036906\\
4.2391	0.14362	0.003829	0.003829\\
4.264	0.13298	0.0037271	0.0037271\\
4.2878	0.12234	0.0039995	0.0039995\\
4.3125	0.1117	0.0041657	0.0041657\\
4.3462	0.10106	0.004	0.004\\
4.3689	0.090426	0.003943	0.003943\\
4.389	0.079787	0.0037316	0.0037316\\
4.4095	0.069149	0.0040381	0.0040381\\
4.4462	0.058511	0.0052621	0.0052621\\
4.5305	0.047872	0.0061177	0.0061177\\
4.5839	0.037234	0.0060254	0.0060254\\
4.7235	0.026596	0.0050878	0.0050878\\
4.7846	0.015957	0.0047448	0.0047448\\
5.1606	0.0053191	0.0037077	0.0037077\\
};
\addplot [color=mycolor1, mark repeat=9, mark phase=1, mark=+, mark options={solid, mycolor1}]
  table[row sep=crcr]{%
3.5875	0.98981\\
3.6038	0.9523\\
3.6202	0.90806\\
3.6365	0.86588\\
3.6528	0.82566\\
3.6692	0.7873\\
3.6855	0.75073\\
3.7018	0.71586\\
3.7182	0.68261\\
3.7345	0.65091\\
3.7508	0.62068\\
3.7672	0.59185\\
3.7835	0.56436\\
3.7999	0.53815\\
3.8162	0.51316\\
3.8325	0.48933\\
3.8489	0.4666\\
3.8652	0.44494\\
3.8815	0.42427\\
3.8979	0.40457\\
3.9142	0.38579\\
3.9305	0.36787\\
3.9469	0.35079\\
3.9632	0.3345\\
3.9796	0.31897\\
3.9959	0.30416\\
4.0122	0.29004\\
4.0286	0.27658\\
4.0449	0.26373\\
4.0612	0.25149\\
4.0776	0.23982\\
4.0939	0.22868\\
4.1102	0.21807\\
4.1266	0.20794\\
4.1429	0.19829\\
4.1593	0.18909\\
4.1756	0.18031\\
4.1919	0.17194\\
4.2083	0.16396\\
4.2246	0.15635\\
4.2409	0.14909\\
4.2573	0.14217\\
4.2736	0.13557\\
4.2899	0.12928\\
4.3063	0.12328\\
4.3226	0.11756\\
4.339	0.11211\\
4.3553	0.1069\\
4.3716	0.10194\\
4.388	0.097212\\
4.4043	0.092701\\
4.4206	0.088399\\
4.437	0.084298\\
4.4533	0.080386\\
4.4696	0.076656\\
4.486	0.0731\\
4.5023	0.069708\\
4.5187	0.066474\\
4.535	0.06339\\
4.5513	0.060449\\
4.5677	0.057645\\
4.584	0.05497\\
4.6003	0.05242\\
4.6167	0.049989\\
4.633	0.04767\\
4.6493	0.045459\\
4.6657	0.04335\\
4.682	0.041339\\
4.6984	0.039422\\
4.7147	0.037593\\
4.731	0.03585\\
4.7474	0.034187\\
4.7637	0.032601\\
4.78	0.031089\\
4.7964	0.029647\\
4.8127	0.028272\\
4.829	0.026961\\
4.8454	0.025711\\
4.8617	0.024519\\
4.8781	0.023382\\
4.8944	0.022297\\
4.9107	0.021264\\
4.9271	0.020278\\
4.9434	0.019337\\
4.9597	0.018441\\
4.9761	0.017586\\
4.9924	0.01677\\
5.0087	0.015993\\
5.0251	0.015251\\
5.0414	0.014544\\
5.0577	0.01387\\
5.0741	0.013227\\
5.0904	0.012614\\
5.1068	0.012029\\
5.1231	0.011471\\
5.1394	0.01094\\
5.1558	0.010433\\
5.1721	0.009949\\
5.1884	0.0094879\\
5.2048	0.0090481\\
};
\addlegendentry{\ApproxParetoFitLegend}


\addplot[area legend, draw=none, fill=mycolor1, fill opacity=0.15, postaction={pattern={Lines[angle=90,distance={5pt},line width={0.5pt}]}, pattern color=mycolor1!60}, forget plot]
table[row sep=crcr] {%
x	y\\
3.5875	0.99262\\
3.6038	0.95656\\
3.6202	0.91195\\
3.6365	0.86942\\
3.6528	0.82888\\
3.6692	0.79023\\
3.6855	0.75338\\
3.7018	0.71825\\
3.7182	0.68477\\
3.7345	0.65285\\
3.7508	0.62242\\
3.7672	0.59341\\
3.7835	0.56576\\
3.7999	0.53941\\
3.8162	0.51428\\
3.8325	0.49034\\
3.8489	0.46751\\
3.8652	0.44576\\
3.8815	0.42502\\
3.8979	0.40526\\
3.9142	0.38643\\
3.9305	0.36848\\
3.9469	0.35138\\
3.9632	0.33508\\
3.9796	0.31954\\
3.9959	0.30474\\
4.0122	0.29062\\
4.0286	0.27716\\
4.0449	0.26433\\
4.0612	0.25209\\
4.0776	0.24043\\
4.0939	0.2293\\
4.1102	0.2187\\
4.1266	0.20858\\
4.1429	0.19893\\
4.1593	0.18973\\
4.1756	0.18096\\
4.1919	0.17259\\
4.2083	0.16461\\
4.2246	0.157\\
4.2409	0.14974\\
4.2573	0.14282\\
4.2736	0.13622\\
4.2899	0.12992\\
4.3063	0.12392\\
4.3226	0.11819\\
4.339	0.11273\\
4.3553	0.10752\\
4.3716	0.10255\\
4.388	0.097811\\
4.4043	0.093291\\
4.4206	0.08898\\
4.437	0.084868\\
4.4533	0.080947\\
4.4696	0.077207\\
4.486	0.07364\\
4.5023	0.070237\\
4.5187	0.066992\\
4.535	0.063897\\
4.5513	0.060945\\
4.5677	0.05813\\
4.584	0.055445\\
4.6003	0.052883\\
4.6167	0.05044\\
4.633	0.04811\\
4.6493	0.045888\\
4.6657	0.043769\\
4.682	0.041747\\
4.6984	0.039819\\
4.7147	0.03798\\
4.731	0.036226\\
4.7474	0.034553\\
4.7637	0.032957\\
4.78	0.031435\\
4.7964	0.029983\\
4.8127	0.028599\\
4.829	0.027278\\
4.8454	0.026018\\
4.8617	0.024817\\
4.8781	0.023671\\
4.8944	0.022578\\
4.9107	0.021535\\
4.9271	0.020541\\
4.9434	0.019593\\
4.9597	0.018688\\
4.9761	0.017825\\
4.9924	0.017002\\
5.0087	0.016217\\
5.0251	0.015469\\
5.0414	0.014755\\
5.0577	0.014074\\
5.0741	0.013424\\
5.0904	0.012804\\
5.1068	0.012213\\
5.1231	0.011649\\
5.1394	0.011112\\
5.1558	0.010599\\
5.1721	0.01011\\
5.1884	0.0096431\\
5.2048	0.009198\\
5.2048	0.0088982\\
5.1884	0.0093327\\
5.1721	0.0097884\\
5.1558	0.010266\\
5.1394	0.010768\\
5.1231	0.011293\\
5.1068	0.011845\\
5.0904	0.012423\\
5.0741	0.01303\\
5.0577	0.013666\\
5.0414	0.014334\\
5.0251	0.015034\\
5.0087	0.015768\\
4.9924	0.016538\\
4.9761	0.017346\\
4.9597	0.018193\\
4.9434	0.019082\\
4.9271	0.020014\\
4.9107	0.020992\\
4.8944	0.022017\\
4.8781	0.023092\\
4.8617	0.024221\\
4.8454	0.025404\\
4.829	0.026645\\
4.8127	0.027946\\
4.7964	0.029312\\
4.78	0.030744\\
4.7637	0.032246\\
4.7474	0.033821\\
4.731	0.035473\\
4.7147	0.037207\\
4.6984	0.039024\\
4.682	0.040931\\
4.6657	0.042931\\
4.6493	0.045029\\
4.633	0.047229\\
4.6167	0.049537\\
4.6003	0.051958\\
4.584	0.054496\\
4.5677	0.05716\\
4.5513	0.059953\\
4.535	0.062883\\
4.5187	0.065956\\
4.5023	0.069179\\
4.486	0.07256\\
4.4696	0.076106\\
4.4533	0.079825\\
4.437	0.083727\\
4.4206	0.087819\\
4.4043	0.092111\\
4.388	0.096613\\
4.3716	0.10133\\
4.3553	0.10629\\
4.339	0.11148\\
4.3226	0.11693\\
4.3063	0.12265\\
4.2899	0.12864\\
4.2736	0.13493\\
4.2573	0.14152\\
4.2409	0.14844\\
4.2246	0.1557\\
4.2083	0.16331\\
4.1919	0.17129\\
4.1756	0.17966\\
4.1593	0.18844\\
4.1429	0.19765\\
4.1266	0.20731\\
4.1102	0.21744\\
4.0939	0.22806\\
4.0776	0.2392\\
4.0612	0.25089\\
4.0449	0.26314\\
4.0286	0.27599\\
4.0122	0.28946\\
3.9959	0.30359\\
3.9796	0.3184\\
3.9632	0.33393\\
3.9469	0.3502\\
3.9305	0.36726\\
3.9142	0.38514\\
3.8979	0.40388\\
3.8815	0.42353\\
3.8652	0.44412\\
3.8489	0.4657\\
3.8325	0.48832\\
3.8162	0.51203\\
3.7999	0.5369\\
3.7835	0.56296\\
3.7672	0.59028\\
3.7508	0.61893\\
3.7345	0.64897\\
3.7182	0.68045\\
3.7018	0.71347\\
3.6855	0.74809\\
3.6692	0.78438\\
3.6528	0.82243\\
3.6365	0.86233\\
3.6202	0.90417\\
3.6038	0.94803\\
3.5875	0.98701\\
}--cycle;
\addplot [color=mycolor1, mark repeat=9, mark phase=1, mark=x, mark options={solid, mycolor1}]
  table[row sep=crcr]{%
3.5875	1.0333\\
3.6038	0.9858\\
3.6202	0.94046\\
3.6365	0.89722\\
3.6528	0.85604\\
3.6692	0.81684\\
3.6855	0.77956\\
3.7018	0.74413\\
3.7182	0.71046\\
3.7345	0.67847\\
3.7508	0.64808\\
3.7672	0.61921\\
3.7835	0.59176\\
3.7999	0.56565\\
3.8162	0.54081\\
3.8325	0.51715\\
3.8489	0.49459\\
3.8652	0.47307\\
3.8815	0.45251\\
3.8979	0.43285\\
3.9142	0.41403\\
3.9305	0.396\\
3.9469	0.37869\\
3.9632	0.36207\\
3.9796	0.3461\\
3.9959	0.33073\\
4.0122	0.31593\\
4.0286	0.30167\\
4.0449	0.28792\\
4.0612	0.27466\\
4.0776	0.26186\\
4.0939	0.24952\\
4.1102	0.2376\\
4.1266	0.2261\\
4.1429	0.215\\
4.1593	0.2043\\
4.1756	0.19398\\
4.1919	0.18403\\
4.2083	0.17444\\
4.2246	0.16521\\
4.2409	0.15633\\
4.2573	0.1478\\
4.2736	0.1396\\
4.2899	0.13173\\
4.3063	0.12419\\
4.3226	0.11696\\
4.339	0.11005\\
4.3553	0.10344\\
4.3716	0.097133\\
4.388	0.091115\\
4.4043	0.085381\\
4.4206	0.079925\\
4.437	0.07474\\
4.4533	0.069817\\
4.4696	0.065149\\
4.486	0.060728\\
4.5023	0.056547\\
4.5187	0.052596\\
4.535	0.048868\\
4.5513	0.045356\\
4.5677	0.042049\\
4.584	0.038941\\
4.6003	0.036023\\
4.6167	0.033287\\
4.633	0.030724\\
4.6493	0.028327\\
4.6657	0.026088\\
4.682	0.023999\\
4.6984	0.022053\\
4.7147	0.020241\\
4.731	0.018558\\
4.7474	0.016995\\
4.7637	0.015546\\
4.78	0.014205\\
4.7964	0.012965\\
4.8127	0.011819\\
4.829	0.010763\\
4.8454	0.0097892\\
4.8617	0.0088937\\
4.8781	0.0080709\\
4.8944	0.0073157\\
4.9107	0.0066236\\
4.9271	0.0059901\\
4.9434	0.0054109\\
4.9597	0.004882\\
4.9761	0.0043998\\
4.9924	0.0039606\\
5.0087	0.0035611\\
5.0251	0.0031982\\
5.0414	0.0028689\\
5.0577	0.0025705\\
5.0741	0.0023005\\
5.0904	0.0020565\\
5.1068	0.0018362\\
5.1231	0.0016375\\
5.1394	0.0014587\\
5.1558	0.0012979\\
5.1721	0.0011534\\
5.1884	0.0010238\\
5.2048	0.00090774\\
};
\addlegendentry{\ApproxBLNFitLegend}


\addplot[area legend, draw=none, fill=mycolor1, fill opacity=0.15, postaction={pattern={Lines[angle=45,distance={5pt},line width={0.5pt}]}, pattern color=mycolor1!60}, forget plot]
table[row sep=crcr] {%
x	y\\
3.5875	1.0352\\
3.6038	0.98762\\
3.6202	0.94218\\
3.6365	0.89886\\
3.6528	0.8576\\
3.6692	0.81834\\
3.6855	0.78102\\
3.7018	0.74555\\
3.7182	0.71186\\
3.7345	0.67987\\
3.7508	0.64948\\
3.7672	0.62061\\
3.7835	0.59317\\
3.7999	0.56707\\
3.8162	0.54223\\
3.8325	0.51856\\
3.8489	0.496\\
3.8652	0.47447\\
3.8815	0.45389\\
3.8979	0.43421\\
3.9142	0.41536\\
3.9305	0.39729\\
3.9469	0.37994\\
3.9632	0.36328\\
3.9796	0.34725\\
3.9959	0.33183\\
4.0122	0.31698\\
4.0286	0.30266\\
4.0449	0.28885\\
4.0612	0.27553\\
4.0776	0.26268\\
4.0939	0.25028\\
4.1102	0.2383\\
4.1266	0.22675\\
4.1429	0.2156\\
4.1593	0.20485\\
4.1756	0.19449\\
4.1919	0.1845\\
4.2083	0.17488\\
4.2246	0.16562\\
4.2409	0.15672\\
4.2573	0.14817\\
4.2736	0.13996\\
4.2899	0.13208\\
4.3063	0.12453\\
4.3226	0.1173\\
4.339	0.11039\\
4.3553	0.10378\\
4.3716	0.097471\\
4.388	0.091454\\
4.4043	0.085722\\
4.4206	0.080267\\
4.437	0.075082\\
4.4533	0.070158\\
4.4696	0.065489\\
4.486	0.061066\\
4.5023	0.056881\\
4.5187	0.052926\\
4.535	0.049193\\
4.5513	0.045674\\
4.5677	0.042361\\
4.584	0.039245\\
4.6003	0.036319\\
4.6167	0.033574\\
4.633	0.031002\\
4.6493	0.028596\\
4.6657	0.026347\\
4.682	0.024248\\
4.6984	0.022291\\
4.7147	0.02047\\
4.731	0.018776\\
4.7474	0.017203\\
4.7637	0.015744\\
4.78	0.014393\\
4.7964	0.013142\\
4.8127	0.011987\\
4.829	0.010921\\
4.8454	0.0099385\\
4.8617	0.0090341\\
4.8781	0.0082025\\
4.8944	0.0074391\\
4.9107	0.0067389\\
4.9271	0.0060977\\
4.9434	0.0055112\\
4.9597	0.0049754\\
4.9761	0.0044865\\
4.9924	0.0040409\\
5.0087	0.0036355\\
5.0251	0.0032669\\
5.0414	0.0029323\\
5.0577	0.002629\\
5.0741	0.0023543\\
5.0904	0.0021058\\
5.1068	0.0018814\\
5.1231	0.001679\\
5.1394	0.0014966\\
5.1558	0.0013324\\
5.1721	0.0011849\\
5.1884	0.0010525\\
5.2048	0.00093381\\
5.2048	0.00088167\\
5.1884	0.00099514\\
5.1721	0.0011219\\
5.1558	0.0012633\\
5.1394	0.0014208\\
5.1231	0.0015961\\
5.1068	0.0017909\\
5.0904	0.0020071\\
5.0741	0.0022468\\
5.0577	0.0025121\\
5.0414	0.0028055\\
5.0251	0.0031294\\
5.0087	0.0034867\\
4.9924	0.0038802\\
4.9761	0.0043131\\
4.9597	0.0047887\\
4.9434	0.0053105\\
4.9271	0.0058824\\
4.9107	0.0065083\\
4.8944	0.0071924\\
4.8781	0.0079392\\
4.8617	0.0087534\\
4.8454	0.0096399\\
4.829	0.010604\\
4.8127	0.011651\\
4.7964	0.012787\\
4.78	0.014017\\
4.7637	0.015349\\
4.7474	0.016787\\
4.731	0.01834\\
4.7147	0.020013\\
4.6984	0.021814\\
4.682	0.02375\\
4.6657	0.025829\\
4.6493	0.028058\\
4.633	0.030446\\
4.6167	0.032999\\
4.6003	0.035727\\
4.584	0.038637\\
4.5677	0.041738\\
4.5513	0.045037\\
4.535	0.048544\\
4.5187	0.052266\\
4.5023	0.056213\\
4.486	0.060391\\
4.4696	0.064809\\
4.4533	0.069476\\
4.437	0.074398\\
4.4206	0.079584\\
4.4043	0.08504\\
4.388	0.090775\\
4.3716	0.096795\\
4.3553	0.10311\\
4.339	0.10972\\
4.3226	0.11663\\
4.3063	0.12385\\
4.2899	0.13139\\
4.2736	0.13925\\
4.2573	0.14743\\
4.2409	0.15595\\
4.2246	0.1648\\
4.2083	0.174\\
4.1919	0.18356\\
4.1756	0.19347\\
4.1593	0.20375\\
4.1429	0.2144\\
4.1266	0.22545\\
4.1102	0.2369\\
4.0939	0.24876\\
4.0776	0.26105\\
4.0612	0.27379\\
4.0449	0.28699\\
4.0286	0.30068\\
4.0122	0.31488\\
3.9959	0.32963\\
3.9796	0.34495\\
3.9632	0.36087\\
3.9469	0.37745\\
3.9305	0.39471\\
3.9142	0.41271\\
3.8979	0.4315\\
3.8815	0.45114\\
3.8652	0.47168\\
3.8489	0.49319\\
3.8325	0.51574\\
3.8162	0.5394\\
3.7999	0.56424\\
3.7835	0.59035\\
3.7672	0.6178\\
3.7508	0.64668\\
3.7345	0.67707\\
3.7182	0.70905\\
3.7018	0.7427\\
3.6855	0.7781\\
3.6692	0.81534\\
3.6528	0.85447\\
3.6365	0.89559\\
3.6202	0.93873\\
3.6038	0.98397\\
3.5875	1.0313\\
}--cycle;
      % \input{../../../../.tex/20/KS/SizeDistributionFittingsfig6.tex}
    \end{groupplot}

    \begin{groupplot}[
      group style={
        group name=Row2,
        plotGroup,
      },
      plotRow,
      row3Keys,
      plotLegendSW,
      %
      % xmin=3637.5,xmax=2.5003e+05,
      ymin=0.00067767,ymax=1.0208,
    ]
      \nextgroupplot[%
        at={($(Row1 c1r1.south west)-(0,\yPlotSep em)-(0,\plotHeight cm)$)},
        xmin=3.5608,xmax=5.3101,
        % xmin=3.4,xmax=5.4,
        % ymin=0.0041299,ymax=1.0005,
      ] % This file was created by matlab2tikz.
%
\definecolor{mycolor1}{rgb}{0.00000,0.44700,0.74100}%
%
\addplot[only marks, mark=*, mark size=1.1180pt, color=black, fill=black, opacity=0.60, draw=none, mark options={draw=none,line width=0pt}] table[row sep=crcr]{%
x	y\\
3.5623	0.99468\\
3.5723	0.98404\\
3.573	0.9734\\
3.5753	0.96277\\
3.5766	0.95213\\
3.5781	0.94149\\
3.5826	0.93085\\
3.5851	0.92021\\
3.5911	0.90957\\
3.5942	0.89894\\
3.596	0.8883\\
3.6016	0.87766\\
3.603	0.86702\\
3.6047	0.85638\\
3.6086	0.84574\\
3.6086	0.83511\\
3.6097	0.82447\\
3.6142	0.81383\\
3.6261	0.80319\\
3.6383	0.79255\\
3.6464	0.78191\\
3.6535	0.77128\\
3.657	0.76064\\
3.6658	0.75\\
3.6678	0.73936\\
3.6687	0.72872\\
3.6702	0.71809\\
3.6723	0.70745\\
3.6793	0.69681\\
3.6867	0.68617\\
3.6949	0.67553\\
3.7015	0.66489\\
3.7031	0.65426\\
3.719	0.64362\\
3.7213	0.63298\\
3.7265	0.62234\\
3.7307	0.6117\\
3.737	0.60106\\
3.7374	0.59043\\
3.7423	0.57979\\
3.7487	0.56915\\
3.7677	0.55851\\
3.7685	0.54787\\
3.7768	0.53723\\
3.8005	0.5266\\
3.8015	0.51596\\
3.8022	0.50532\\
3.8032	0.49468\\
3.8039	0.48404\\
3.8108	0.4734\\
3.8214	0.46277\\
3.8405	0.45213\\
3.8605	0.44149\\
3.863	0.43085\\
3.8645	0.42021\\
3.8662	0.40957\\
3.8806	0.39894\\
3.8878	0.3883\\
3.8964	0.37766\\
3.905	0.36702\\
3.9294	0.35638\\
3.9303	0.34574\\
3.954	0.33511\\
3.9604	0.32447\\
3.9669	0.31383\\
3.9747	0.30319\\
3.9822	0.29255\\
3.9929	0.28191\\
4.0051	0.27128\\
4.0144	0.26064\\
4.0161	0.25\\
4.0433	0.23936\\
4.0578	0.22872\\
4.073	0.21809\\
4.0857	0.20745\\
4.0904	0.19681\\
4.1168	0.18617\\
4.1265	0.17553\\
4.127	0.16489\\
4.143	0.15426\\
4.1756	0.14362\\
4.2156	0.13298\\
4.3116	0.12234\\
4.3276	0.1117\\
4.3662	0.10106\\
4.4791	0.090426\\
4.4912	0.079787\\
4.517	0.069149\\
4.5743	0.058511\\
4.5914	0.047872\\
4.6138	0.037234\\
4.7898	0.026596\\
5.0876	0.015957\\
5.3101	0.0053191\\
};
\addlegendentry{\RealEmpiricCCDFLegend}

\addplot [color=black, mark repeat=9, mark phase=1, mark=+, mark options={solid, black}]
  table[row sep=crcr]{%
3.5608	1\\
3.5785	0.94954\\
3.5961	0.90162\\
3.6138	0.85612\\
3.6315	0.81292\\
3.6492	0.7719\\
3.6668	0.73295\\
3.6845	0.69596\\
3.7022	0.66084\\
3.7198	0.6275\\
3.7375	0.59583\\
3.7552	0.56576\\
3.7728	0.53721\\
3.7905	0.5101\\
3.8082	0.48436\\
3.8259	0.45992\\
3.8435	0.43671\\
3.8612	0.41468\\
3.8789	0.39375\\
3.8965	0.37388\\
3.9142	0.35501\\
3.9319	0.3371\\
3.9495	0.32009\\
3.9672	0.30394\\
3.9849	0.2886\\
4.0026	0.27404\\
4.0202	0.26021\\
4.0379	0.24708\\
4.0556	0.23461\\
4.0732	0.22277\\
4.0909	0.21153\\
4.1086	0.20085\\
4.1262	0.19072\\
4.1439	0.18109\\
4.1616	0.17196\\
4.1793	0.16328\\
4.1969	0.15504\\
4.2146	0.14722\\
4.2323	0.13979\\
4.2499	0.13273\\
4.2676	0.12603\\
4.2853	0.11967\\
4.3029	0.11364\\
4.3206	0.1079\\
4.3383	0.10246\\
4.356	0.097286\\
4.3736	0.092377\\
4.3913	0.087716\\
4.409	0.083289\\
4.4266	0.079086\\
4.4443	0.075095\\
4.462	0.071306\\
4.4796	0.067708\\
4.4973	0.064291\\
4.515	0.061047\\
4.5327	0.057966\\
4.5503	0.055041\\
4.568	0.052264\\
4.5857	0.049626\\
4.6033	0.047122\\
4.621	0.044744\\
4.6387	0.042486\\
4.6563	0.040342\\
4.674	0.038307\\
4.6917	0.036373\\
4.7094	0.034538\\
4.727	0.032795\\
4.7447	0.03114\\
4.7624	0.029569\\
4.78	0.028077\\
4.7977	0.02666\\
4.8154	0.025315\\
4.833	0.024037\\
4.8507	0.022824\\
4.8684	0.021672\\
4.8861	0.020579\\
4.9037	0.01954\\
4.9214	0.018554\\
4.9391	0.017618\\
4.9567	0.016729\\
4.9744	0.015885\\
4.9921	0.015083\\
5.0097	0.014322\\
5.0274	0.013599\\
5.0451	0.012913\\
5.0628	0.012261\\
5.0804	0.011643\\
5.0981	0.011055\\
5.1158	0.010497\\
5.1334	0.0099676\\
5.1511	0.0094646\\
5.1688	0.008987\\
5.1864	0.0085335\\
5.2041	0.0081029\\
5.2218	0.007694\\
5.2395	0.0073058\\
5.2571	0.0069371\\
5.2748	0.006587\\
5.2925	0.0062546\\
5.3101	0.005939\\
};
\addlegendentry{\RealParetoFitLegend}

\addplot [color=black, mark repeat=9, mark phase=1, mark=x, mark options={solid, black}]
  table[row sep=crcr]{%
3.5608	1.0005\\
3.5785	0.95316\\
3.5961	0.90742\\
3.6138	0.86328\\
3.6315	0.82076\\
3.6492	0.77986\\
3.6668	0.74057\\
3.6845	0.70289\\
3.7022	0.66682\\
3.7198	0.63232\\
3.7375	0.59938\\
3.7552	0.56796\\
3.7728	0.53805\\
3.7905	0.5096\\
3.8082	0.48257\\
3.8259	0.45692\\
3.8435	0.43262\\
3.8612	0.4096\\
3.8789	0.38783\\
3.8965	0.36726\\
3.9142	0.34784\\
3.9319	0.32951\\
3.9495	0.31222\\
3.9672	0.29594\\
3.9849	0.28059\\
4.0026	0.26614\\
4.0202	0.25254\\
4.0379	0.23974\\
4.0556	0.22769\\
4.0732	0.21634\\
4.0909	0.20566\\
4.1086	0.19561\\
4.1262	0.18613\\
4.1439	0.1772\\
4.1616	0.16878\\
4.1793	0.16082\\
4.1969	0.15331\\
4.2146	0.14621\\
4.2323	0.1395\\
4.2499	0.13313\\
4.2676	0.1271\\
4.2853	0.12137\\
4.3029	0.11592\\
4.3206	0.11074\\
4.3383	0.10581\\
4.356	0.1011\\
4.3736	0.09661\\
4.3913	0.092318\\
4.409	0.088212\\
4.4266	0.084281\\
4.4443	0.080514\\
4.462	0.076901\\
4.4796	0.073434\\
4.4973	0.070104\\
4.515	0.066905\\
4.5327	0.06383\\
4.5503	0.060873\\
4.568	0.058029\\
4.5857	0.055293\\
4.6033	0.052661\\
4.621	0.050128\\
4.6387	0.047691\\
4.6563	0.045347\\
4.674	0.043093\\
4.6917	0.040926\\
4.7094	0.038844\\
4.727	0.036843\\
4.7447	0.034921\\
4.7624	0.033077\\
4.78	0.031309\\
4.7977	0.029613\\
4.8154	0.027989\\
4.833	0.026433\\
4.8507	0.024946\\
4.8684	0.023523\\
4.8861	0.022165\\
4.9037	0.020868\\
4.9214	0.019632\\
4.9391	0.018454\\
4.9567	0.017332\\
4.9744	0.016265\\
4.9921	0.015251\\
5.0097	0.014288\\
5.0274	0.013375\\
5.0451	0.01251\\
5.0628	0.01169\\
5.0804	0.010915\\
5.0981	0.010182\\
5.1158	0.0094906\\
5.1334	0.0088382\\
5.1511	0.0082235\\
5.1688	0.0076448\\
5.1864	0.0071005\\
5.2041	0.0065892\\
5.2218	0.0061093\\
5.2395	0.0056592\\
5.2571	0.0052377\\
5.2748	0.0048432\\
5.2925	0.0044744\\
5.3101	0.0041299\\
};
\addlegendentry{\RealBLNFitLegend}

      % \input{../../../../.tex/20/KS/SizeDistributionFittingsfig13.tex}
      \legend{}
  
      \nextgroupplot[%
        xmin=3.598,xmax=5.4115,
        % xmin=3.6,xmax=5.4,
        % ymin=0.00067767,ymax=1.0131,
      ] % This file was created by matlab2tikz.
%
\definecolor{mycolor1}{rgb}{0.00000,0.44706,0.69804}%
\definecolor{mycolor2}{rgb}{0.00000,0.44700,0.74100}%
%
\addplot[only marks, mark=*, mark size=1.1180pt, color=mycolor1, fill=mycolor1, opacity=0.60, draw=none, mark options={draw=none,line width=0pt}] table[row sep=crcr]{%
x	y\\
3.6037	0.99468\\
3.6073	0.98404\\
3.6108	0.9734\\
3.6145	0.96277\\
3.6181	0.95213\\
3.622	0.94149\\
3.6258	0.93085\\
3.6297	0.92021\\
3.6339	0.90957\\
3.6376	0.89894\\
3.6419	0.8883\\
3.6462	0.87766\\
3.6499	0.86702\\
3.6544	0.85638\\
3.6601	0.84574\\
3.6647	0.83511\\
3.6691	0.82447\\
3.6744	0.81383\\
3.6795	0.80319\\
3.6851	0.79255\\
3.6902	0.78191\\
3.6946	0.77128\\
3.701	0.76064\\
3.7067	0.75\\
3.7119	0.73936\\
3.7176	0.72872\\
3.7232	0.71809\\
3.7287	0.70745\\
3.7342	0.69681\\
3.7398	0.68617\\
3.7458	0.67553\\
3.7524	0.66489\\
3.7597	0.65426\\
3.7658	0.64362\\
3.773	0.63298\\
3.7802	0.62234\\
3.7873	0.6117\\
3.7936	0.60106\\
3.8011	0.59043\\
3.8081	0.57979\\
3.8151	0.56915\\
3.8222	0.55851\\
3.8279	0.54787\\
3.8349	0.53723\\
3.8422	0.5266\\
3.85	0.51596\\
3.8572	0.50532\\
3.8649	0.49468\\
3.8714	0.48404\\
3.8791	0.4734\\
3.8881	0.46277\\
3.8956	0.45213\\
3.9016	0.44149\\
3.91	0.43085\\
3.9175	0.42021\\
3.9272	0.40957\\
3.9348	0.39894\\
3.9449	0.3883\\
3.9548	0.37766\\
3.9657	0.36702\\
3.9754	0.35638\\
3.9863	0.34574\\
3.9973	0.33511\\
4.0075	0.32447\\
4.0175	0.31383\\
4.0299	0.30319\\
4.0406	0.29255\\
4.0543	0.28191\\
4.0681	0.27128\\
4.0824	0.26064\\
4.0965	0.25\\
4.1083	0.23936\\
4.1212	0.22872\\
4.138	0.21809\\
4.1548	0.20745\\
4.1697	0.19681\\
4.1848	0.18617\\
4.2011	0.17553\\
4.2197	0.16489\\
4.2359	0.15426\\
4.2571	0.14362\\
4.278	0.13298\\
4.3018	0.12234\\
4.3248	0.1117\\
4.3516	0.10106\\
4.3803	0.090426\\
4.4066	0.079787\\
4.4414	0.069149\\
4.4829	0.058511\\
4.538	0.047872\\
4.6052	0.037234\\
4.7174	0.026596\\
4.8035	0.015957\\
5.159	0.0053191\\
};
\addlegendentry{\ExactEmpiricCCDFLegend}

\addplot [color=mycolor1, only marks, every error bar/.append style={opacity=0.30}, mark=*, mark size=0pt, draw=none, forget plot]
 plot [error bars/.cd, x dir=both, x explicit, error bar style={line width=1pt, color=mycolor1}, error mark options={mark=none,mark size=0pt}]
 table[row sep=crcr, x error plus index=2, x error minus index=3]{%
3.6037	0.99468	0.0031873	0.0031873\\
3.6073	0.98404	0.0032744	0.0032744\\
3.6108	0.9734	0.0033314	0.0033314\\
3.6145	0.96277	0.0034058	0.0034058\\
3.6181	0.95213	0.0034851	0.0034851\\
3.622	0.94149	0.0035	0.0035\\
3.6258	0.93085	0.0036615	0.0036615\\
3.6297	0.92021	0.0035244	0.0035244\\
3.6339	0.90957	0.0034445	0.0034445\\
3.6376	0.89894	0.0035539	0.0035539\\
3.6419	0.8883	0.0034716	0.0034716\\
3.6462	0.87766	0.0034679	0.0034679\\
3.6499	0.86702	0.0034622	0.0034622\\
3.6544	0.85638	0.0034513	0.0034513\\
3.6601	0.84574	0.0034059	0.0034059\\
3.6647	0.83511	0.003463	0.003463\\
3.6691	0.82447	0.0034186	0.0034186\\
3.6744	0.81383	0.003484	0.003484\\
3.6795	0.80319	0.0035559	0.0035559\\
3.6851	0.79255	0.0035894	0.0035894\\
3.6902	0.78191	0.0034953	0.0034953\\
3.6946	0.77128	0.0034979	0.0034979\\
3.701	0.76064	0.003521	0.003521\\
3.7067	0.75	0.0035332	0.0035332\\
3.7119	0.73936	0.0036677	0.0036677\\
3.7176	0.72872	0.0037632	0.0037632\\
3.7232	0.71809	0.0038613	0.0038613\\
3.7287	0.70745	0.0039704	0.0039704\\
3.7342	0.69681	0.0038912	0.0038912\\
3.7398	0.68617	0.0037488	0.0037488\\
3.7458	0.67553	0.0038371	0.0038371\\
3.7524	0.66489	0.0038114	0.0038114\\
3.7597	0.65426	0.0036991	0.0036991\\
3.7658	0.64362	0.0035915	0.0035915\\
3.773	0.63298	0.0034558	0.0034558\\
3.7802	0.62234	0.0035273	0.0035273\\
3.7873	0.6117	0.0037482	0.0037482\\
3.7936	0.60106	0.0040024	0.0040024\\
3.8011	0.59043	0.0041113	0.0041113\\
3.8081	0.57979	0.0041614	0.0041614\\
3.8151	0.56915	0.0042496	0.0042496\\
3.8222	0.55851	0.0043104	0.0043104\\
3.8279	0.54787	0.0043573	0.0043573\\
3.8349	0.53723	0.0041869	0.0041869\\
3.8422	0.5266	0.0040847	0.0040847\\
3.85	0.51596	0.003834	0.003834\\
3.8572	0.50532	0.0039218	0.0039218\\
3.8649	0.49468	0.0038832	0.0038832\\
3.8714	0.48404	0.003944	0.003944\\
3.8791	0.4734	0.0037933	0.0037933\\
3.8881	0.46277	0.0035871	0.0035871\\
3.8956	0.45213	0.0035147	0.0035147\\
3.9016	0.44149	0.0037251	0.0037251\\
3.91	0.43085	0.0038328	0.0038328\\
3.9175	0.42021	0.0041924	0.0041924\\
3.9272	0.40957	0.0043005	0.0043005\\
3.9348	0.39894	0.0042934	0.0042934\\
3.9449	0.3883	0.0044537	0.0044537\\
3.9548	0.37766	0.0047398	0.0047398\\
3.9657	0.36702	0.004744	0.004744\\
3.9754	0.35638	0.0046879	0.0046879\\
3.9863	0.34574	0.0045711	0.0045711\\
3.9973	0.33511	0.0047913	0.0047913\\
4.0075	0.32447	0.0047427	0.0047427\\
4.0175	0.31383	0.0050029	0.0050029\\
4.0299	0.30319	0.0051338	0.0051338\\
4.0406	0.29255	0.005174	0.005174\\
4.0543	0.28191	0.0053199	0.0053199\\
4.0681	0.27128	0.0056242	0.0056242\\
4.0824	0.26064	0.005481	0.005481\\
4.0965	0.25	0.0054067	0.0054067\\
4.1083	0.23936	0.0051031	0.0051031\\
4.1212	0.22872	0.0051613	0.0051613\\
4.138	0.21809	0.0048839	0.0048839\\
4.1548	0.20745	0.0052666	0.0052666\\
4.1697	0.19681	0.0050782	0.0050782\\
4.1848	0.18617	0.0052326	0.0052326\\
4.2011	0.17553	0.0048643	0.0048643\\
4.2197	0.16489	0.0049985	0.0049985\\
4.2359	0.15426	0.00532	0.00532\\
4.2571	0.14362	0.0053879	0.0053879\\
4.278	0.13298	0.0055106	0.0055106\\
4.3018	0.12234	0.0062277	0.0062277\\
4.3248	0.1117	0.0057529	0.0057529\\
4.3516	0.10106	0.00653	0.00653\\
4.3803	0.090426	0.0065879	0.0065879\\
4.4066	0.079787	0.0066297	0.0066297\\
4.4414	0.069149	0.0075694	0.0075694\\
4.4829	0.058511	0.0081425	0.0081425\\
4.538	0.047872	0.0091626	0.0091626\\
4.6052	0.037234	0.0096523	0.0096523\\
4.7174	0.026596	0.011213	0.011213\\
4.8035	0.015957	0.010763	0.010763\\
5.159	0.0053191	0.0090493	0.0090493\\
};
\addplot [color=mycolor1, mark repeat=9, mark phase=1, mark=+, mark options={solid, mycolor1}]
  table[row sep=crcr]{%
3.602	0.98157\\
3.6191	0.95011\\
3.6361	0.90544\\
3.6532	0.86078\\
3.6703	0.81832\\
3.6874	0.77796\\
3.7044	0.7396\\
3.7215	0.70312\\
3.7386	0.66845\\
3.7557	0.6355\\
3.7727	0.60417\\
3.7898	0.57438\\
3.8069	0.54607\\
3.824	0.51915\\
3.841	0.49357\\
3.8581	0.46924\\
3.8752	0.44612\\
3.8923	0.42414\\
3.9094	0.40324\\
3.9264	0.38337\\
3.9435	0.36449\\
3.9606	0.34653\\
3.9777	0.32947\\
3.9947	0.31324\\
4.0118	0.29781\\
4.0289	0.28315\\
4.046	0.26921\\
4.063	0.25595\\
4.0801	0.24335\\
4.0972	0.23137\\
4.1143	0.21998\\
4.1313	0.20916\\
4.1484	0.19886\\
4.1655	0.18908\\
4.1826	0.17977\\
4.1996	0.17093\\
4.2167	0.16252\\
4.2338	0.15452\\
4.2509	0.14692\\
4.2679	0.13969\\
4.285	0.13282\\
4.3021	0.12629\\
4.3192	0.12008\\
4.3362	0.11418\\
4.3533	0.10856\\
4.3704	0.10322\\
4.3875	0.098149\\
4.4045	0.093324\\
4.4216	0.088737\\
4.4387	0.084376\\
4.4558	0.080229\\
4.4728	0.076286\\
4.4899	0.072538\\
4.507	0.068974\\
4.5241	0.065585\\
4.5411	0.062363\\
4.5582	0.0593\\
4.5753	0.056387\\
4.5924	0.053618\\
4.6094	0.050984\\
4.6265	0.048481\\
4.6436	0.0461\\
4.6607	0.043837\\
4.6777	0.041685\\
4.6948	0.039638\\
4.7119	0.037693\\
4.729	0.035843\\
4.746	0.034083\\
4.7631	0.032411\\
4.7802	0.03082\\
4.7973	0.029308\\
4.8143	0.02787\\
4.8314	0.026503\\
4.8485	0.025203\\
4.8656	0.023966\\
4.8826	0.022791\\
4.8997	0.021673\\
4.9168	0.02061\\
4.9339	0.0196\\
4.9509	0.018639\\
4.968	0.017725\\
4.9851	0.016856\\
5.0022	0.01603\\
5.0192	0.015244\\
5.0363	0.014497\\
5.0534	0.013786\\
5.0705	0.013111\\
5.0875	0.012468\\
5.1046	0.011857\\
5.1217	0.011276\\
5.1388	0.010724\\
5.1558	0.010199\\
5.1729	0.0096991\\
5.19	0.0092241\\
5.2071	0.0087724\\
5.2241	0.0083428\\
5.2412	0.0079343\\
5.2583	0.0075458\\
5.2754	0.0071764\\
5.2924	0.0068251\\
};
\addlegendentry{\ExactParetoFitLegend}


\addplot[area legend, draw=none, fill=mycolor1, fill opacity=0.15, postaction={pattern={Lines[angle=90,distance={5pt},line width={0.5pt}]}, pattern color=mycolor1!60}, forget plot]
table[row sep=crcr] {%
x	y\\
3.602	0.98693\\
3.6191	0.95791\\
3.6361	0.91319\\
3.6532	0.86785\\
3.6703	0.82477\\
3.6874	0.78384\\
3.7044	0.74494\\
3.7215	0.70799\\
3.7386	0.67288\\
3.7557	0.63951\\
3.7727	0.60782\\
3.7898	0.5777\\
3.8069	0.54909\\
3.824	0.52191\\
3.841	0.49608\\
3.8581	0.47155\\
3.8752	0.44824\\
3.8923	0.4261\\
3.9094	0.40506\\
3.9264	0.38508\\
3.9435	0.3661\\
3.9606	0.34807\\
3.9777	0.33094\\
3.9947	0.31467\\
4.0118	0.29921\\
4.0289	0.28451\\
4.046	0.27055\\
4.063	0.25728\\
4.0801	0.24467\\
4.0972	0.23268\\
4.1143	0.22128\\
4.1313	0.21045\\
4.1484	0.20015\\
4.1655	0.19036\\
4.1826	0.18104\\
4.1996	0.17219\\
4.2167	0.16377\\
4.2338	0.15577\\
4.2509	0.14815\\
4.2679	0.14091\\
4.285	0.13403\\
4.3021	0.12748\\
4.3192	0.12126\\
4.3362	0.11534\\
4.3533	0.1097\\
4.3704	0.10435\\
4.3875	0.099256\\
4.4045	0.094411\\
4.4216	0.089804\\
4.4387	0.085422\\
4.4558	0.081254\\
4.4728	0.07729\\
4.4899	0.07352\\
4.507	0.069934\\
4.5241	0.066524\\
4.5411	0.06328\\
4.5582	0.060194\\
4.5753	0.05726\\
4.5924	0.054468\\
4.6094	0.051813\\
4.6265	0.049288\\
4.6436	0.046886\\
4.6607	0.044601\\
4.6777	0.042428\\
4.6948	0.040361\\
4.7119	0.038394\\
4.729	0.036524\\
4.746	0.034745\\
4.7631	0.033053\\
4.7802	0.031443\\
4.7973	0.029912\\
4.8143	0.028455\\
4.8314	0.02707\\
4.8485	0.025752\\
4.8656	0.024499\\
4.8826	0.023306\\
4.8997	0.022172\\
4.9168	0.021093\\
4.9339	0.020066\\
4.9509	0.01909\\
4.968	0.018161\\
4.9851	0.017278\\
5.0022	0.016437\\
5.0192	0.015638\\
5.0363	0.014877\\
5.0534	0.014154\\
5.0705	0.013465\\
5.0875	0.012811\\
5.1046	0.012188\\
5.1217	0.011595\\
5.1388	0.011032\\
5.1558	0.010495\\
5.1729	0.0099853\\
5.19	0.0095\\
5.2071	0.0090384\\
5.2241	0.0085992\\
5.2412	0.0081814\\
5.2583	0.0077839\\
5.2754	0.0074058\\
5.2924	0.0070461\\
5.2924	0.0066041\\
5.2754	0.006947\\
5.2583	0.0073077\\
5.2412	0.0076871\\
5.2241	0.0080864\\
5.2071	0.0085063\\
5.19	0.0089481\\
5.1729	0.0094129\\
5.1558	0.0099019\\
5.1388	0.010416\\
5.1217	0.010958\\
5.1046	0.011527\\
5.0875	0.012126\\
5.0705	0.012756\\
5.0534	0.013419\\
5.0363	0.014116\\
5.0192	0.01485\\
5.0022	0.015622\\
4.9851	0.016434\\
4.968	0.017288\\
4.9509	0.018187\\
4.9339	0.019133\\
4.9168	0.020128\\
4.8997	0.021174\\
4.8826	0.022276\\
4.8656	0.023434\\
4.8485	0.024653\\
4.8314	0.025935\\
4.8143	0.027285\\
4.7973	0.028704\\
4.7802	0.030198\\
4.7631	0.031769\\
4.746	0.033422\\
4.729	0.035161\\
4.7119	0.036991\\
4.6948	0.038916\\
4.6777	0.040942\\
4.6607	0.043073\\
4.6436	0.045315\\
4.6265	0.047674\\
4.6094	0.050156\\
4.5924	0.052767\\
4.5753	0.055515\\
4.5582	0.058405\\
4.5411	0.061447\\
4.5241	0.064647\\
4.507	0.068013\\
4.4899	0.071556\\
4.4728	0.075283\\
4.4558	0.079204\\
4.4387	0.08333\\
4.4216	0.08767\\
4.4045	0.092237\\
4.3875	0.097043\\
4.3704	0.1021\\
4.3533	0.10742\\
4.3362	0.11301\\
4.3192	0.1189\\
4.3021	0.1251\\
4.285	0.13161\\
4.2679	0.13847\\
4.2509	0.14569\\
4.2338	0.15328\\
4.2167	0.16126\\
4.1996	0.16966\\
4.1826	0.1785\\
4.1655	0.1878\\
4.1484	0.19757\\
4.1313	0.20786\\
4.1143	0.21868\\
4.0972	0.23006\\
4.0801	0.24203\\
4.063	0.25462\\
4.046	0.26786\\
4.0289	0.28178\\
4.0118	0.29642\\
3.9947	0.31181\\
3.9777	0.32799\\
3.9606	0.345\\
3.9435	0.36287\\
3.9264	0.38167\\
3.9094	0.40142\\
3.8923	0.42218\\
3.8752	0.444\\
3.8581	0.46694\\
3.841	0.49105\\
3.824	0.5164\\
3.8069	0.54305\\
3.7898	0.57106\\
3.7727	0.60051\\
3.7557	0.63148\\
3.7386	0.66403\\
3.7215	0.69826\\
3.7044	0.73425\\
3.6874	0.77209\\
3.6703	0.81187\\
3.6532	0.85371\\
3.6361	0.8977\\
3.6191	0.94232\\
3.602	0.97621\\
}--cycle;
\addplot [color=mycolor1, mark repeat=9, mark phase=1, mark=x, mark options={solid, mycolor1}]
  table[row sep=crcr]{%
3.602	1.0071\\
3.6191	0.96177\\
3.6361	0.91825\\
3.6532	0.8765\\
3.6703	0.83648\\
3.6874	0.79817\\
3.7044	0.7615\\
3.7215	0.72644\\
3.7386	0.69292\\
3.7557	0.6609\\
3.7727	0.6303\\
3.7898	0.60109\\
3.8069	0.57318\\
3.824	0.54654\\
3.841	0.52109\\
3.8581	0.49679\\
3.8752	0.47357\\
3.8923	0.45138\\
3.9094	0.43017\\
3.9264	0.40989\\
3.9435	0.3905\\
3.9606	0.37194\\
3.9777	0.35418\\
3.9947	0.33717\\
4.0118	0.32088\\
4.0289	0.30528\\
4.046	0.29032\\
4.063	0.27599\\
4.0801	0.26225\\
4.0972	0.24907\\
4.1143	0.23644\\
4.1313	0.22433\\
4.1484	0.21272\\
4.1655	0.20158\\
4.1826	0.19091\\
4.1996	0.18069\\
4.2167	0.17089\\
4.2338	0.16151\\
4.2509	0.15253\\
4.2679	0.14395\\
4.285	0.13573\\
4.3021	0.12789\\
4.3192	0.1204\\
4.3362	0.11325\\
4.3533	0.10643\\
4.3704	0.099936\\
4.3875	0.093753\\
4.4045	0.087873\\
4.4216	0.082285\\
4.4387	0.07698\\
4.4558	0.071949\\
4.4728	0.067182\\
4.4899	0.062669\\
4.507	0.058402\\
4.5241	0.054372\\
4.5411	0.050569\\
4.5582	0.046984\\
4.5753	0.04361\\
4.5924	0.040436\\
4.6094	0.037455\\
4.6265	0.034658\\
4.6436	0.032036\\
4.6607	0.029582\\
4.6777	0.027287\\
4.6948	0.025144\\
4.7119	0.023145\\
4.729	0.021282\\
4.746	0.019548\\
4.7631	0.017936\\
4.7802	0.01644\\
4.7973	0.015053\\
4.8143	0.013767\\
4.8314	0.012578\\
4.8485	0.01148\\
4.8656	0.010466\\
4.8826	0.0095315\\
4.8997	0.008671\\
4.9168	0.0078798\\
4.9339	0.0071531\\
4.9509	0.0064864\\
4.968	0.0058755\\
4.9851	0.0053165\\
5.0022	0.0048054\\
5.0192	0.0043389\\
5.0363	0.0039134\\
5.0534	0.0035259\\
5.0705	0.0031733\\
5.0875	0.0028529\\
5.1046	0.0025622\\
5.1217	0.0022986\\
5.1388	0.0020599\\
5.1558	0.0018441\\
5.1729	0.0016491\\
5.19	0.0014732\\
5.2071	0.0013146\\
5.2241	0.0011719\\
5.2412	0.0010436\\
5.2583	0.00092834\\
5.2754	0.00082495\\
5.2924	0.00073231\\
};
\addlegendentry{\ExactBLNFitLegend}


\addplot[area legend, draw=none, fill=mycolor1, fill opacity=0.15, postaction={pattern={Lines[angle=45,distance={5pt},line width={0.5pt}]}, pattern color=mycolor1!60}, forget plot]
table[row sep=crcr] {%
x	y\\
3.602	1.0131\\
3.6191	0.96758\\
3.6361	0.92391\\
3.6532	0.882\\
3.6703	0.84185\\
3.6874	0.80339\\
3.7044	0.76659\\
3.7215	0.73139\\
3.7386	0.69774\\
3.7557	0.66559\\
3.7727	0.63487\\
3.7898	0.60552\\
3.8069	0.57749\\
3.824	0.55072\\
3.841	0.52514\\
3.8581	0.5007\\
3.8752	0.47734\\
3.8923	0.45501\\
3.9094	0.43366\\
3.9264	0.41324\\
3.9435	0.3937\\
3.9606	0.375\\
3.9777	0.35709\\
3.9947	0.33994\\
4.0118	0.32351\\
4.0289	0.30777\\
4.046	0.29268\\
4.063	0.27822\\
4.0801	0.26436\\
4.0972	0.25107\\
4.1143	0.23834\\
4.1313	0.22613\\
4.1484	0.21443\\
4.1655	0.20323\\
4.1826	0.19249\\
4.1996	0.18221\\
4.2167	0.17236\\
4.2338	0.16294\\
4.2509	0.15393\\
4.2679	0.14531\\
4.285	0.13707\\
4.3021	0.1292\\
4.3192	0.12169\\
4.3362	0.11452\\
4.3533	0.10768\\
4.3704	0.10117\\
4.3875	0.094966\\
4.4045	0.089064\\
4.4216	0.083454\\
4.4387	0.078126\\
4.4558	0.07307\\
4.4728	0.068276\\
4.4899	0.063736\\
4.507	0.059441\\
4.5241	0.055381\\
4.5411	0.051547\\
4.5582	0.047931\\
4.5753	0.044524\\
4.5924	0.041317\\
4.6094	0.038303\\
4.6265	0.035472\\
4.6436	0.032817\\
4.6607	0.030329\\
4.6777	0.028001\\
4.6948	0.025825\\
4.7119	0.023793\\
4.729	0.021899\\
4.746	0.020134\\
4.7631	0.018492\\
4.7802	0.016966\\
4.7973	0.015549\\
4.8143	0.014236\\
4.8314	0.01302\\
4.8485	0.011896\\
4.8656	0.010857\\
4.8826	0.009898\\
4.8997	0.0090146\\
4.9168	0.0082013\\
4.9339	0.0074536\\
4.9509	0.006767\\
4.968	0.0061371\\
4.9851	0.00556\\
5.0022	0.0050319\\
5.0192	0.0045492\\
5.0363	0.0041085\\
5.0534	0.0037066\\
5.0705	0.0033406\\
5.0875	0.0030076\\
5.1046	0.0027049\\
5.1217	0.0024302\\
5.1388	0.0021811\\
5.1558	0.0019555\\
5.1729	0.0017515\\
5.19	0.0015671\\
5.2071	0.0014007\\
5.2241	0.0012507\\
5.2412	0.0011156\\
5.2583	0.0009941\\
5.2754	0.00088493\\
5.2924	0.00078695\\
5.2924	0.00067767\\
5.2754	0.00076497\\
5.2583	0.00086258\\
5.2412	0.00097158\\
5.2241	0.0010931\\
5.2071	0.0012286\\
5.19	0.0013793\\
5.1729	0.0015468\\
5.1558	0.0017327\\
5.1388	0.0019388\\
5.1217	0.002167\\
5.1046	0.0024195\\
5.0875	0.0026983\\
5.0705	0.003006\\
5.0534	0.0033451\\
5.0363	0.0037182\\
5.0192	0.0041285\\
5.0022	0.0045789\\
4.9851	0.0050729\\
4.968	0.0056139\\
4.9509	0.0062058\\
4.9339	0.0068525\\
4.9168	0.0075583\\
4.8997	0.0083275\\
4.8826	0.0091649\\
4.8656	0.010075\\
4.8485	0.011064\\
4.8314	0.012137\\
4.8143	0.013299\\
4.7973	0.014556\\
4.7802	0.015915\\
4.7631	0.017381\\
4.746	0.018962\\
4.729	0.020665\\
4.7119	0.022496\\
4.6948	0.024463\\
4.6777	0.026573\\
4.6607	0.028835\\
4.6436	0.031255\\
4.6265	0.033843\\
4.6094	0.036607\\
4.5924	0.039555\\
4.5753	0.042696\\
4.5582	0.046038\\
4.5411	0.049591\\
4.5241	0.053363\\
4.507	0.057363\\
4.4899	0.061602\\
4.4728	0.066087\\
4.4558	0.070828\\
4.4387	0.075835\\
4.4216	0.081116\\
4.4045	0.086682\\
4.3875	0.092541\\
4.3704	0.098704\\
4.3533	0.10518\\
4.3362	0.11198\\
4.3192	0.1191\\
4.3021	0.12658\\
4.285	0.1344\\
4.2679	0.14258\\
4.2509	0.15114\\
4.2338	0.16008\\
4.2167	0.16942\\
4.1996	0.17917\\
4.1826	0.18933\\
4.1655	0.19994\\
4.1484	0.211\\
4.1313	0.22253\\
4.1143	0.23454\\
4.0972	0.24707\\
4.0801	0.26014\\
4.063	0.27376\\
4.046	0.28797\\
4.0289	0.30279\\
4.0118	0.31826\\
3.9947	0.3344\\
3.9777	0.35127\\
3.9606	0.36888\\
3.9435	0.38729\\
3.9264	0.40654\\
3.9094	0.42668\\
3.8923	0.44774\\
3.8752	0.46979\\
3.8581	0.49287\\
3.841	0.51704\\
3.824	0.54236\\
3.8069	0.56887\\
3.7898	0.59665\\
3.7727	0.62574\\
3.7557	0.65621\\
3.7386	0.6881\\
3.7215	0.72149\\
3.7044	0.75642\\
3.6874	0.79295\\
3.6703	0.83112\\
3.6532	0.87099\\
3.6361	0.91259\\
3.6191	0.95596\\
3.602	1.0011\\
}--cycle;
      % \input{../../../../.tex/20/KS/SizeDistributionFittingsfig14.tex}
      \legend{}
  
      \nextgroupplot[%
        xmin=3.6008,xmax=5.4115,
        % xmin=3.4,xmax=5.4,
        % ymin=0.00074383,ymax=1.0208,
      ] % This file was created by matlab2tikz.
%
\definecolor{mycolor1}{rgb}{0.83529,0.36863,0.00000}%
\definecolor{mycolor2}{rgb}{0.00000,0.44700,0.74100}%
%
\addplot[only marks, mark=*, mark size=1.1180pt, color=mycolor1, fill=mycolor1, opacity=0.60, draw=none, mark options={draw=none,line width=0pt}] table[row sep=crcr]{%
x	y\\
3.5981	0.99468\\
3.6023	0.98404\\
3.6064	0.9734\\
3.6105	0.96277\\
3.615	0.95213\\
3.619	0.94149\\
3.6233	0.93085\\
3.6271	0.92021\\
3.6318	0.90957\\
3.6366	0.89894\\
3.6403	0.8883\\
3.6443	0.87766\\
3.6487	0.86702\\
3.6532	0.85638\\
3.6584	0.84574\\
3.6638	0.83511\\
3.6687	0.82447\\
3.6733	0.81383\\
3.6786	0.80319\\
3.6828	0.79255\\
3.6881	0.78191\\
3.6936	0.77128\\
3.6991	0.76064\\
3.7048	0.75\\
3.71	0.73936\\
3.7158	0.72872\\
3.721	0.71809\\
3.7265	0.70745\\
3.7321	0.69681\\
3.737	0.68617\\
3.7428	0.67553\\
3.7488	0.66489\\
3.7556	0.65426\\
3.763	0.64362\\
3.7694	0.63298\\
3.776	0.62234\\
3.7822	0.6117\\
3.7889	0.60106\\
3.7946	0.59043\\
3.801	0.57979\\
3.8083	0.56915\\
3.8176	0.55851\\
3.824	0.54787\\
3.8317	0.53723\\
3.8394	0.5266\\
3.8475	0.51596\\
3.8558	0.50532\\
3.8625	0.49468\\
3.8709	0.48404\\
3.8781	0.4734\\
3.8874	0.46277\\
3.8955	0.45213\\
3.9029	0.44149\\
3.9129	0.43085\\
3.9213	0.42021\\
3.9292	0.40957\\
3.9386	0.39894\\
3.9488	0.3883\\
3.9566	0.37766\\
3.9662	0.36702\\
3.9771	0.35638\\
3.9865	0.34574\\
3.9964	0.33511\\
4.0082	0.32447\\
4.0215	0.31383\\
4.0338	0.30319\\
4.0469	0.29255\\
4.0576	0.28191\\
4.0718	0.27128\\
4.0833	0.26064\\
4.0968	0.25\\
4.1097	0.23936\\
4.123	0.22872\\
4.1357	0.21809\\
4.1509	0.20745\\
4.1655	0.19681\\
4.1838	0.18617\\
4.2013	0.17553\\
4.2175	0.16489\\
4.2364	0.15426\\
4.2558	0.14362\\
4.2753	0.13298\\
4.3016	0.12234\\
4.3251	0.1117\\
4.3478	0.10106\\
4.3732	0.090426\\
4.4078	0.079787\\
4.4449	0.069149\\
4.4792	0.058511\\
4.5401	0.047872\\
4.6076	0.037234\\
4.7133	0.026596\\
4.7931	0.015957\\
5.1695	0.0053191\\
};
\addlegendentry{\ApproxEmpiricCCDFLegend}

\addplot [color=mycolor1, only marks, every error bar/.append style={opacity=0.30}, mark=*, mark size=0pt, draw=none, forget plot]
 plot [error bars/.cd, x dir=both, x explicit, error bar style={line width=1pt, color=mycolor1}, error mark options={mark=none,mark size=0pt}]
 table[row sep=crcr, x error plus index=2, x error minus index=3]{%
3.5981	0.99468	0.0033318	0.0033318\\
3.6023	0.98404	0.0032921	0.0032921\\
3.6064	0.9734	0.0032161	0.0032161\\
3.6105	0.96277	0.0033671	0.0033671\\
3.615	0.95213	0.0032769	0.0032769\\
3.619	0.94149	0.0031849	0.0031849\\
3.6233	0.93085	0.0031801	0.0031801\\
3.6271	0.92021	0.0032614	0.0032614\\
3.6318	0.90957	0.0033254	0.0033254\\
3.6366	0.89894	0.0032616	0.0032616\\
3.6403	0.8883	0.0034099	0.0034099\\
3.6443	0.87766	0.0035325	0.0035325\\
3.6487	0.86702	0.0034635	0.0034635\\
3.6532	0.85638	0.0034169	0.0034169\\
3.6584	0.84574	0.0034051	0.0034051\\
3.6638	0.83511	0.0032979	0.0032979\\
3.6687	0.82447	0.00348	0.00348\\
3.6733	0.81383	0.0035628	0.0035628\\
3.6786	0.80319	0.0034968	0.0034968\\
3.6828	0.79255	0.0034696	0.0034696\\
3.6881	0.78191	0.003529	0.003529\\
3.6936	0.77128	0.0037115	0.0037115\\
3.6991	0.76064	0.0036262	0.0036262\\
3.7048	0.75	0.0034662	0.0034662\\
3.71	0.73936	0.0034329	0.0034329\\
3.7158	0.72872	0.0033673	0.0033673\\
3.721	0.71809	0.0033237	0.0033237\\
3.7265	0.70745	0.0035375	0.0035375\\
3.7321	0.69681	0.0036761	0.0036761\\
3.737	0.68617	0.0037267	0.0037267\\
3.7428	0.67553	0.0038231	0.0038231\\
3.7488	0.66489	0.0038249	0.0038249\\
3.7556	0.65426	0.0039771	0.0039771\\
3.763	0.64362	0.0041257	0.0041257\\
3.7694	0.63298	0.0039609	0.0039609\\
3.776	0.62234	0.0039102	0.0039102\\
3.7822	0.6117	0.0039858	0.0039858\\
3.7889	0.60106	0.0042246	0.0042246\\
3.7946	0.59043	0.0041824	0.0041824\\
3.801	0.57979	0.00417	0.00417\\
3.8083	0.56915	0.0042139	0.0042139\\
3.8176	0.55851	0.0043078	0.0043078\\
3.824	0.54787	0.0042625	0.0042625\\
3.8317	0.53723	0.0043864	0.0043864\\
3.8394	0.5266	0.0043818	0.0043818\\
3.8475	0.51596	0.0043929	0.0043929\\
3.8558	0.50532	0.0045964	0.0045964\\
3.8625	0.49468	0.0045453	0.0045453\\
3.8709	0.48404	0.0045508	0.0045508\\
3.8781	0.4734	0.0044692	0.0044692\\
3.8874	0.46277	0.0044678	0.0044678\\
3.8955	0.45213	0.0044672	0.0044672\\
3.9029	0.44149	0.0044505	0.0044505\\
3.9129	0.43085	0.0043948	0.0043948\\
3.9213	0.42021	0.0042943	0.0042943\\
3.9292	0.40957	0.004448	0.004448\\
3.9386	0.39894	0.0045776	0.0045776\\
3.9488	0.3883	0.0046377	0.0046377\\
3.9566	0.37766	0.0044836	0.0044836\\
3.9662	0.36702	0.00447	0.00447\\
3.9771	0.35638	0.0045325	0.0045325\\
3.9865	0.34574	0.0047347	0.0047347\\
3.9964	0.33511	0.0049881	0.0049881\\
4.0082	0.32447	0.0053043	0.0053043\\
4.0215	0.31383	0.0054884	0.0054884\\
4.0338	0.30319	0.0055015	0.0055015\\
4.0469	0.29255	0.0056892	0.0056892\\
4.0576	0.28191	0.0055564	0.0055564\\
4.0718	0.27128	0.0057063	0.0057063\\
4.0833	0.26064	0.0058436	0.0058436\\
4.0968	0.25	0.0061932	0.0061932\\
4.1097	0.23936	0.0059823	0.0059823\\
4.123	0.22872	0.0058606	0.0058606\\
4.1357	0.21809	0.0056708	0.0056708\\
4.1509	0.20745	0.0058678	0.0058678\\
4.1655	0.19681	0.0053964	0.0053964\\
4.1838	0.18617	0.0053979	0.0053979\\
4.2013	0.17553	0.0050116	0.0050116\\
4.2175	0.16489	0.0053215	0.0053215\\
4.2364	0.15426	0.005493	0.005493\\
4.2558	0.14362	0.0055404	0.0055404\\
4.2753	0.13298	0.0054582	0.0054582\\
4.3016	0.12234	0.0058416	0.0058416\\
4.3251	0.1117	0.0059397	0.0059397\\
4.3478	0.10106	0.0058893	0.0058893\\
4.3732	0.090426	0.0064931	0.0064931\\
4.4078	0.079787	0.0070281	0.0070281\\
4.4449	0.069149	0.0074367	0.0074367\\
4.4792	0.058511	0.0080956	0.0080956\\
4.5401	0.047872	0.010773	0.010773\\
4.6076	0.037234	0.010679	0.010679\\
4.7133	0.026596	0.010343	0.010343\\
4.7931	0.015957	0.011898	0.011898\\
5.1695	0.0053191	0.0088428	0.0088428\\
};
\addplot [color=mycolor1, mark repeat=9, mark phase=1, mark=+, mark options={solid, mycolor1}]
  table[row sep=crcr]{%
3.5976	0.98063\\
3.6147	0.94676\\
3.6318	0.90374\\
3.6489	0.85955\\
3.666	0.81754\\
3.6832	0.77758\\
3.7003	0.73957\\
3.7174	0.70343\\
3.7345	0.66906\\
3.7516	0.63637\\
3.7688	0.60528\\
3.7859	0.57571\\
3.803	0.54759\\
3.8201	0.52084\\
3.8372	0.4954\\
3.8544	0.47121\\
3.8715	0.4482\\
3.8886	0.42632\\
3.9057	0.40551\\
3.9228	0.38572\\
3.94	0.36689\\
3.9571	0.34898\\
3.9742	0.33195\\
3.9913	0.31576\\
4.0084	0.30035\\
4.0256	0.2857\\
4.0427	0.27176\\
4.0598	0.2585\\
4.0769	0.2459\\
4.094	0.2339\\
4.1112	0.2225\\
4.1283	0.21165\\
4.1454	0.20133\\
4.1625	0.19152\\
4.1796	0.18218\\
4.1968	0.1733\\
4.2139	0.16486\\
4.231	0.15682\\
4.2481	0.14918\\
4.2652	0.14191\\
4.2824	0.135\\
4.2995	0.12842\\
4.3166	0.12217\\
4.3337	0.11622\\
4.3508	0.11056\\
4.368	0.10517\\
4.3851	0.10005\\
4.4022	0.095183\\
4.4193	0.090549\\
4.4364	0.086142\\
4.4536	0.081949\\
4.4707	0.077961\\
4.4878	0.074167\\
4.5049	0.070558\\
4.522	0.067125\\
4.5392	0.06386\\
4.5563	0.060753\\
4.5734	0.057798\\
4.5905	0.054987\\
4.6076	0.052313\\
4.6248	0.049769\\
4.6419	0.04735\\
4.659	0.045048\\
4.6761	0.042858\\
4.6932	0.040774\\
4.7104	0.038793\\
4.7275	0.036907\\
4.7446	0.035114\\
4.7617	0.033408\\
4.7788	0.031785\\
4.796	0.03024\\
4.8131	0.028771\\
4.8302	0.027374\\
4.8473	0.026044\\
4.8644	0.02478\\
4.8816	0.023576\\
4.8987	0.022432\\
4.9158	0.021343\\
4.9329	0.020306\\
4.95	0.019321\\
4.9672	0.018383\\
4.9843	0.017491\\
5.0014	0.016642\\
5.0185	0.015834\\
5.0356	0.015066\\
5.0528	0.014335\\
5.0699	0.01364\\
5.087	0.012978\\
5.1041	0.012349\\
5.1212	0.01175\\
5.1384	0.01118\\
5.1555	0.010638\\
5.1726	0.010122\\
5.1897	0.0096315\\
5.2068	0.0091646\\
5.224	0.0087204\\
5.2411	0.0082977\\
5.2582	0.0078956\\
5.2753	0.007513\\
5.2924	0.007149\\
};
\addlegendentry{\ApproxParetoFitLegend}


\addplot[area legend, draw=none, fill=mycolor1, fill opacity=0.15, postaction={pattern={Lines[angle=90,distance={5pt},line width={0.5pt}]}, pattern color=mycolor1!60}, forget plot]
table[row sep=crcr] {%
x	y\\
3.5976	0.98556\\
3.6147	0.95413\\
3.6318	0.9115\\
3.6489	0.86662\\
3.666	0.82396\\
3.6832	0.78341\\
3.7003	0.74487\\
3.7174	0.70823\\
3.7345	0.67341\\
3.7516	0.64032\\
3.7688	0.60887\\
3.7859	0.57897\\
3.803	0.55056\\
3.8201	0.52356\\
3.8372	0.4979\\
3.8544	0.47352\\
3.8715	0.45035\\
3.8886	0.42833\\
3.9057	0.40741\\
3.9228	0.38752\\
3.94	0.36863\\
3.9571	0.35067\\
3.9742	0.3336\\
3.9913	0.31737\\
4.0084	0.30194\\
4.0256	0.28728\\
4.0427	0.27333\\
4.0598	0.26007\\
4.0769	0.24745\\
4.094	0.23546\\
4.1112	0.22404\\
4.1283	0.21319\\
4.1454	0.20286\\
4.1625	0.19304\\
4.1796	0.1837\\
4.1968	0.17481\\
4.2139	0.16635\\
4.231	0.1583\\
4.2481	0.15064\\
4.2652	0.14336\\
4.2824	0.13643\\
4.2995	0.12983\\
4.3166	0.12355\\
4.3337	0.11758\\
4.3508	0.1119\\
4.368	0.10649\\
4.3851	0.10135\\
4.4022	0.096454\\
4.4193	0.091795\\
4.4364	0.087362\\
4.4536	0.083143\\
4.4707	0.079128\\
4.4878	0.075308\\
4.5049	0.071673\\
4.522	0.068213\\
4.5392	0.064921\\
4.5563	0.061788\\
4.5734	0.058806\\
4.5905	0.055969\\
4.6076	0.053269\\
4.6248	0.050699\\
4.6419	0.048254\\
4.659	0.045927\\
4.6761	0.043712\\
4.6932	0.041604\\
4.7104	0.039598\\
4.7275	0.037689\\
4.7446	0.035872\\
4.7617	0.034143\\
4.7788	0.032497\\
4.796	0.030931\\
4.8131	0.02944\\
4.8302	0.028022\\
4.8473	0.026672\\
4.8644	0.025387\\
4.8816	0.024164\\
4.8987	0.023\\
4.9158	0.021892\\
4.9329	0.020838\\
4.95	0.019834\\
4.9672	0.018879\\
4.9843	0.01797\\
5.0014	0.017105\\
5.0185	0.016282\\
5.0356	0.015498\\
5.0528	0.014752\\
5.0699	0.014042\\
5.087	0.013367\\
5.1041	0.012723\\
5.1212	0.012111\\
5.1384	0.011529\\
5.1555	0.010974\\
5.1726	0.010446\\
5.1897	0.0099438\\
5.2068	0.0094655\\
5.224	0.0090104\\
5.2411	0.0085771\\
5.2582	0.0081647\\
5.2753	0.0077722\\
5.2924	0.0073986\\
5.2924	0.0068993\\
5.2753	0.0072538\\
5.2582	0.0076264\\
5.2411	0.0080183\\
5.224	0.0084304\\
5.2068	0.0088636\\
5.1897	0.0093192\\
5.1726	0.0097982\\
5.1555	0.010302\\
5.1384	0.010832\\
5.1212	0.011388\\
5.1041	0.011974\\
5.087	0.01259\\
5.0699	0.013237\\
5.0528	0.013918\\
5.0356	0.014634\\
5.0185	0.015387\\
5.0014	0.016179\\
4.9843	0.017011\\
4.9672	0.017886\\
4.95	0.018807\\
4.9329	0.019775\\
4.9158	0.020793\\
4.8987	0.021863\\
4.8816	0.022989\\
4.8644	0.024172\\
4.8473	0.025417\\
4.8302	0.026726\\
4.8131	0.028102\\
4.796	0.02955\\
4.7788	0.031072\\
4.7617	0.032672\\
4.7446	0.034356\\
4.7275	0.036126\\
4.7104	0.037987\\
4.6932	0.039945\\
4.6761	0.042004\\
4.659	0.044169\\
4.6419	0.046445\\
4.6248	0.04884\\
4.6076	0.051358\\
4.5905	0.054005\\
4.5734	0.05679\\
4.5563	0.059719\\
4.5392	0.062799\\
4.522	0.066037\\
4.5049	0.069444\\
4.4878	0.073026\\
4.4707	0.076793\\
4.4536	0.080755\\
4.4364	0.084922\\
4.4193	0.089304\\
4.4022	0.093912\\
4.3851	0.098758\\
4.368	0.10386\\
4.3508	0.10922\\
4.3337	0.11485\\
4.3166	0.12078\\
4.2995	0.12702\\
4.2824	0.13357\\
4.2652	0.14047\\
4.2481	0.14772\\
4.231	0.15535\\
4.2139	0.16336\\
4.1968	0.1718\\
4.1796	0.18067\\
4.1625	0.18999\\
4.1454	0.1998\\
4.1283	0.21011\\
4.1112	0.22095\\
4.094	0.23235\\
4.0769	0.24434\\
4.0598	0.25694\\
4.0427	0.27019\\
4.0256	0.28412\\
4.0084	0.29875\\
3.9913	0.31414\\
3.9742	0.33031\\
3.9571	0.3473\\
3.94	0.36515\\
3.9228	0.38391\\
3.9057	0.40361\\
3.8886	0.42431\\
3.8715	0.44606\\
3.8544	0.46891\\
3.8372	0.49291\\
3.8201	0.51812\\
3.803	0.54461\\
3.7859	0.57244\\
3.7688	0.60168\\
3.7516	0.63241\\
3.7345	0.6647\\
3.7174	0.69862\\
3.7003	0.73428\\
3.6832	0.77174\\
3.666	0.81111\\
3.6489	0.85249\\
3.6318	0.89598\\
3.6147	0.9394\\
3.5976	0.9757\\
}--cycle;
\addplot [color=mycolor1, mark repeat=9, mark phase=1, mark=x, mark options={solid, mycolor1}]
  table[row sep=crcr]{%
3.5976	1.0157\\
3.6147	0.96963\\
3.6318	0.92539\\
3.6489	0.88296\\
3.666	0.84232\\
3.6832	0.80342\\
3.7003	0.76622\\
3.7174	0.73067\\
3.7345	0.6967\\
3.7516	0.66428\\
3.7688	0.63332\\
3.7859	0.60378\\
3.803	0.5756\\
3.8201	0.54871\\
3.8372	0.52305\\
3.8544	0.49856\\
3.8715	0.47519\\
3.8886	0.45288\\
3.9057	0.43157\\
3.9228	0.41121\\
3.94	0.39176\\
3.9571	0.37316\\
3.9742	0.35536\\
3.9913	0.33834\\
4.0084	0.32205\\
4.0256	0.30645\\
4.0427	0.2915\\
4.0598	0.27719\\
4.0769	0.26347\\
4.094	0.25031\\
4.1112	0.23771\\
4.1283	0.22562\\
4.1454	0.21404\\
4.1625	0.20293\\
4.1796	0.19229\\
4.1968	0.18209\\
4.2139	0.17231\\
4.231	0.16295\\
4.2481	0.15399\\
4.2652	0.14542\\
4.2824	0.13722\\
4.2995	0.12938\\
4.3166	0.12189\\
4.3337	0.11474\\
4.3508	0.10792\\
4.368	0.10142\\
4.3851	0.095223\\
4.4022	0.089329\\
4.4193	0.083723\\
4.4364	0.078398\\
4.4536	0.073343\\
4.4707	0.068549\\
4.4878	0.064009\\
4.5049	0.059711\\
4.522	0.055648\\
4.5392	0.051811\\
4.5563	0.048191\\
4.5734	0.044779\\
4.5905	0.041567\\
4.6076	0.038547\\
4.6248	0.03571\\
4.6419	0.033048\\
4.659	0.030553\\
4.6761	0.028217\\
4.6932	0.026034\\
4.7104	0.023994\\
4.7275	0.022091\\
4.7446	0.020318\\
4.7617	0.018668\\
4.7788	0.017134\\
4.796	0.015709\\
4.8131	0.014388\\
4.8302	0.013163\\
4.8473	0.012031\\
4.8644	0.010984\\
4.8816	0.010018\\
4.8987	0.0091267\\
4.9158	0.0083062\\
4.9329	0.0075514\\
4.95	0.0068579\\
4.9672	0.0062215\\
4.9843	0.0056382\\
5.0014	0.0051041\\
5.0185	0.0046157\\
5.0356	0.0041696\\
5.0528	0.0037627\\
5.0699	0.0033918\\
5.087	0.0030543\\
5.1041	0.0027474\\
5.1212	0.0024688\\
5.1384	0.002216\\
5.1555	0.0019871\\
5.1726	0.0017799\\
5.1897	0.0015926\\
5.2068	0.0014236\\
5.224	0.0012711\\
5.2411	0.0011338\\
5.2582	0.0010102\\
5.2753	0.00089921\\
5.2924	0.00079955\\
};
\addlegendentry{\ApproxBLNFitLegend}


\addplot[area legend, draw=none, fill=mycolor1, fill opacity=0.15, postaction={pattern={Lines[angle=45,distance={5pt},line width={0.5pt}]}, pattern color=mycolor1!60}, forget plot]
table[row sep=crcr] {%
x	y\\
3.5976	1.0208\\
3.6147	0.97452\\
3.6318	0.9301\\
3.6489	0.88753\\
3.666	0.84676\\
3.6832	0.80776\\
3.7003	0.77048\\
3.7174	0.73486\\
3.7345	0.70085\\
3.7516	0.66839\\
3.7688	0.63741\\
3.7859	0.60785\\
3.803	0.57964\\
3.8201	0.55272\\
3.8372	0.52703\\
3.8544	0.5025\\
3.8715	0.47908\\
3.8886	0.4567\\
3.9057	0.43531\\
3.9228	0.41487\\
3.94	0.39531\\
3.9571	0.3766\\
3.9742	0.35869\\
3.9913	0.34154\\
4.0084	0.32512\\
4.0256	0.30937\\
4.0427	0.29429\\
4.0598	0.27982\\
4.0769	0.26595\\
4.094	0.25265\\
4.1112	0.2399\\
4.1283	0.22767\\
4.1454	0.21594\\
4.1625	0.2047\\
4.1796	0.19393\\
4.1968	0.1836\\
4.2139	0.17372\\
4.231	0.16425\\
4.2481	0.1552\\
4.2652	0.14654\\
4.2824	0.13827\\
4.2995	0.13037\\
4.3166	0.12283\\
4.3337	0.11563\\
4.3508	0.10878\\
4.368	0.10225\\
4.3851	0.096039\\
4.4022	0.090129\\
4.4193	0.084513\\
4.4364	0.079179\\
4.4536	0.074117\\
4.4707	0.069316\\
4.4878	0.064769\\
4.5049	0.060463\\
4.522	0.056392\\
4.5392	0.052545\\
4.5563	0.048913\\
4.5734	0.045489\\
4.5905	0.042263\\
4.6076	0.039227\\
4.6248	0.036373\\
4.6419	0.033694\\
4.659	0.03118\\
4.6761	0.028825\\
4.6932	0.02662\\
4.7104	0.024559\\
4.7275	0.022635\\
4.7446	0.02084\\
4.7617	0.019168\\
4.7788	0.017611\\
4.796	0.016164\\
4.8131	0.014821\\
4.8302	0.013575\\
4.8473	0.012421\\
4.8644	0.011354\\
4.8816	0.010367\\
4.8987	0.0094562\\
4.9158	0.0086163\\
4.9329	0.0078429\\
4.95	0.0071314\\
4.9672	0.0064777\\
4.9843	0.0058777\\
5.0014	0.0053278\\
5.0185	0.0048243\\
5.0356	0.0043638\\
5.0528	0.0039431\\
5.0699	0.0035593\\
5.087	0.0032094\\
5.1041	0.002891\\
5.1212	0.0026014\\
5.1384	0.0023384\\
5.1555	0.0020998\\
5.1726	0.0018836\\
5.1897	0.0016879\\
5.2068	0.001511\\
5.224	0.0013512\\
5.2411	0.0012071\\
5.2582	0.0010772\\
5.2753	0.00096034\\
5.2924	0.00085526\\
5.2924	0.00074383\\
5.2753	0.00083807\\
5.2582	0.00094324\\
5.2411	0.0010605\\
5.224	0.001191\\
5.2068	0.0013361\\
5.1897	0.0014973\\
5.1726	0.0016761\\
5.1555	0.0018743\\
5.1384	0.0020936\\
5.1212	0.0023361\\
5.1041	0.0026038\\
5.087	0.0028991\\
5.0699	0.0032244\\
5.0528	0.0035822\\
5.0356	0.0039755\\
5.0185	0.0044072\\
5.0014	0.0048804\\
4.9843	0.0053986\\
4.9672	0.0059653\\
4.95	0.0065844\\
4.9329	0.0072599\\
4.9158	0.0079961\\
4.8987	0.0087973\\
4.8816	0.0096684\\
4.8644	0.010614\\
4.8473	0.01164\\
4.8302	0.012752\\
4.8131	0.013954\\
4.796	0.015253\\
4.7788	0.016656\\
4.7617	0.018168\\
4.7446	0.019796\\
4.7275	0.021547\\
4.7104	0.023428\\
4.6932	0.025447\\
4.6761	0.02761\\
4.659	0.029926\\
4.6419	0.032402\\
4.6248	0.035046\\
4.6076	0.037866\\
4.5905	0.040871\\
4.5734	0.044069\\
4.5563	0.047468\\
4.5392	0.051077\\
4.522	0.054904\\
4.5049	0.058959\\
4.4878	0.063248\\
4.4707	0.067782\\
4.4536	0.072569\\
4.4364	0.077617\\
4.4193	0.082933\\
4.4022	0.088528\\
4.3851	0.094408\\
4.368	0.10058\\
4.3508	0.10706\\
4.3337	0.11385\\
4.3166	0.12095\\
4.2995	0.12839\\
4.2824	0.13617\\
4.2652	0.1443\\
4.2481	0.15279\\
4.231	0.16165\\
4.2139	0.17091\\
4.1968	0.18057\\
4.1796	0.19065\\
4.1625	0.20116\\
4.1454	0.21213\\
4.1283	0.22358\\
4.1112	0.23552\\
4.094	0.24798\\
4.0769	0.26098\\
4.0598	0.27455\\
4.0427	0.28872\\
4.0256	0.30352\\
4.0084	0.31898\\
3.9913	0.33514\\
3.9742	0.35203\\
3.9571	0.36971\\
3.94	0.3882\\
3.9228	0.40756\\
3.9057	0.42783\\
3.8886	0.44906\\
3.8715	0.47131\\
3.8544	0.49463\\
3.8372	0.51907\\
3.8201	0.5447\\
3.803	0.57156\\
3.7859	0.59972\\
3.7688	0.62923\\
3.7516	0.66016\\
3.7345	0.69255\\
3.7174	0.72647\\
3.7003	0.76197\\
3.6832	0.79909\\
3.666	0.83788\\
3.6489	0.8784\\
3.6318	0.92067\\
3.6147	0.96474\\
3.5976	1.0106\\
}--cycle;
      % \input{../../../../.tex/20/KS/SizeDistributionFittingsfig15.tex}
      \legend{}
    \end{groupplot}

    \begin{groupplot}[
      group style={
        group name=Row3,
        plotGroup,
      },
      plotRow,
      row3Keys,
      plotLegendSW,
      %
      % xmin=3637.5,xmax=2.5003e+05,
      ymin=0.00076295,ymax=1.0188,
    ]
      \nextgroupplot[%
        at={($(Row2 c1r1.south west)-(0,\yPlotSep em)-(0,\plotHeight cm)$)},
        xmin=3.5608,xmax=5.3101,
        % xmin=3.4,xmax=5.4,
        % ymin=0.0041299,ymax=1.0005,
      ] % This file was created by matlab2tikz.
%
\definecolor{mycolor1}{rgb}{0.00000,0.44700,0.74100}%
%
\addplot[only marks, mark=*, mark size=1.1180pt, color=black, fill=black, opacity=0.60, draw=none, mark options={draw=none,line width=0pt}] table[row sep=crcr]{%
x	y\\
3.5623	0.99468\\
3.5723	0.98404\\
3.573	0.9734\\
3.5753	0.96277\\
3.5766	0.95213\\
3.5781	0.94149\\
3.5826	0.93085\\
3.5851	0.92021\\
3.5911	0.90957\\
3.5942	0.89894\\
3.596	0.8883\\
3.6016	0.87766\\
3.603	0.86702\\
3.6047	0.85638\\
3.6086	0.84574\\
3.6086	0.83511\\
3.6097	0.82447\\
3.6142	0.81383\\
3.6261	0.80319\\
3.6383	0.79255\\
3.6464	0.78191\\
3.6535	0.77128\\
3.657	0.76064\\
3.6658	0.75\\
3.6678	0.73936\\
3.6687	0.72872\\
3.6702	0.71809\\
3.6723	0.70745\\
3.6793	0.69681\\
3.6867	0.68617\\
3.6949	0.67553\\
3.7015	0.66489\\
3.7031	0.65426\\
3.719	0.64362\\
3.7213	0.63298\\
3.7265	0.62234\\
3.7307	0.6117\\
3.737	0.60106\\
3.7374	0.59043\\
3.7423	0.57979\\
3.7487	0.56915\\
3.7677	0.55851\\
3.7685	0.54787\\
3.7768	0.53723\\
3.8005	0.5266\\
3.8015	0.51596\\
3.8022	0.50532\\
3.8032	0.49468\\
3.8039	0.48404\\
3.8108	0.4734\\
3.8214	0.46277\\
3.8405	0.45213\\
3.8605	0.44149\\
3.863	0.43085\\
3.8645	0.42021\\
3.8662	0.40957\\
3.8806	0.39894\\
3.8878	0.3883\\
3.8964	0.37766\\
3.905	0.36702\\
3.9294	0.35638\\
3.9303	0.34574\\
3.954	0.33511\\
3.9604	0.32447\\
3.9669	0.31383\\
3.9747	0.30319\\
3.9822	0.29255\\
3.9929	0.28191\\
4.0051	0.27128\\
4.0144	0.26064\\
4.0161	0.25\\
4.0433	0.23936\\
4.0578	0.22872\\
4.073	0.21809\\
4.0857	0.20745\\
4.0904	0.19681\\
4.1168	0.18617\\
4.1265	0.17553\\
4.127	0.16489\\
4.143	0.15426\\
4.1756	0.14362\\
4.2156	0.13298\\
4.3116	0.12234\\
4.3276	0.1117\\
4.3662	0.10106\\
4.4791	0.090426\\
4.4912	0.079787\\
4.517	0.069149\\
4.5743	0.058511\\
4.5914	0.047872\\
4.6138	0.037234\\
4.7898	0.026596\\
5.0876	0.015957\\
5.3101	0.0053191\\
};
\addlegendentry{\RealEmpiricCCDFLegend}

\addplot [color=black, mark repeat=9, mark phase=1, mark=+, mark options={solid, black}]
  table[row sep=crcr]{%
3.5608	1\\
3.5785	0.94954\\
3.5961	0.90162\\
3.6138	0.85612\\
3.6315	0.81292\\
3.6492	0.7719\\
3.6668	0.73295\\
3.6845	0.69596\\
3.7022	0.66084\\
3.7198	0.6275\\
3.7375	0.59583\\
3.7552	0.56576\\
3.7728	0.53721\\
3.7905	0.5101\\
3.8082	0.48436\\
3.8259	0.45992\\
3.8435	0.43671\\
3.8612	0.41468\\
3.8789	0.39375\\
3.8965	0.37388\\
3.9142	0.35501\\
3.9319	0.3371\\
3.9495	0.32009\\
3.9672	0.30394\\
3.9849	0.2886\\
4.0026	0.27404\\
4.0202	0.26021\\
4.0379	0.24708\\
4.0556	0.23461\\
4.0732	0.22277\\
4.0909	0.21153\\
4.1086	0.20085\\
4.1262	0.19072\\
4.1439	0.18109\\
4.1616	0.17196\\
4.1793	0.16328\\
4.1969	0.15504\\
4.2146	0.14722\\
4.2323	0.13979\\
4.2499	0.13273\\
4.2676	0.12603\\
4.2853	0.11967\\
4.3029	0.11364\\
4.3206	0.1079\\
4.3383	0.10246\\
4.356	0.097286\\
4.3736	0.092377\\
4.3913	0.087716\\
4.409	0.083289\\
4.4266	0.079086\\
4.4443	0.075095\\
4.462	0.071306\\
4.4796	0.067708\\
4.4973	0.064291\\
4.515	0.061047\\
4.5327	0.057966\\
4.5503	0.055041\\
4.568	0.052264\\
4.5857	0.049626\\
4.6033	0.047122\\
4.621	0.044744\\
4.6387	0.042486\\
4.6563	0.040342\\
4.674	0.038307\\
4.6917	0.036373\\
4.7094	0.034538\\
4.727	0.032795\\
4.7447	0.03114\\
4.7624	0.029569\\
4.78	0.028077\\
4.7977	0.02666\\
4.8154	0.025315\\
4.833	0.024037\\
4.8507	0.022824\\
4.8684	0.021672\\
4.8861	0.020579\\
4.9037	0.01954\\
4.9214	0.018554\\
4.9391	0.017618\\
4.9567	0.016729\\
4.9744	0.015885\\
4.9921	0.015083\\
5.0097	0.014322\\
5.0274	0.013599\\
5.0451	0.012913\\
5.0628	0.012261\\
5.0804	0.011643\\
5.0981	0.011055\\
5.1158	0.010497\\
5.1334	0.0099676\\
5.1511	0.0094646\\
5.1688	0.008987\\
5.1864	0.0085335\\
5.2041	0.0081029\\
5.2218	0.007694\\
5.2395	0.0073058\\
5.2571	0.0069371\\
5.2748	0.006587\\
5.2925	0.0062546\\
5.3101	0.005939\\
};
\addlegendentry{\RealParetoFitLegend}

\addplot [color=black, mark repeat=9, mark phase=1, mark=x, mark options={solid, black}]
  table[row sep=crcr]{%
3.5608	1.0005\\
3.5785	0.95316\\
3.5961	0.90742\\
3.6138	0.86328\\
3.6315	0.82076\\
3.6492	0.77986\\
3.6668	0.74057\\
3.6845	0.70289\\
3.7022	0.66682\\
3.7198	0.63232\\
3.7375	0.59938\\
3.7552	0.56796\\
3.7728	0.53805\\
3.7905	0.5096\\
3.8082	0.48257\\
3.8259	0.45692\\
3.8435	0.43262\\
3.8612	0.4096\\
3.8789	0.38783\\
3.8965	0.36726\\
3.9142	0.34784\\
3.9319	0.32951\\
3.9495	0.31222\\
3.9672	0.29594\\
3.9849	0.28059\\
4.0026	0.26614\\
4.0202	0.25254\\
4.0379	0.23974\\
4.0556	0.22769\\
4.0732	0.21634\\
4.0909	0.20566\\
4.1086	0.19561\\
4.1262	0.18613\\
4.1439	0.1772\\
4.1616	0.16878\\
4.1793	0.16082\\
4.1969	0.15331\\
4.2146	0.14621\\
4.2323	0.1395\\
4.2499	0.13313\\
4.2676	0.1271\\
4.2853	0.12137\\
4.3029	0.11592\\
4.3206	0.11074\\
4.3383	0.10581\\
4.356	0.1011\\
4.3736	0.09661\\
4.3913	0.092318\\
4.409	0.088212\\
4.4266	0.084281\\
4.4443	0.080514\\
4.462	0.076901\\
4.4796	0.073434\\
4.4973	0.070104\\
4.515	0.066905\\
4.5327	0.06383\\
4.5503	0.060873\\
4.568	0.058029\\
4.5857	0.055293\\
4.6033	0.052661\\
4.621	0.050128\\
4.6387	0.047691\\
4.6563	0.045347\\
4.674	0.043093\\
4.6917	0.040926\\
4.7094	0.038844\\
4.727	0.036843\\
4.7447	0.034921\\
4.7624	0.033077\\
4.78	0.031309\\
4.7977	0.029613\\
4.8154	0.027989\\
4.833	0.026433\\
4.8507	0.024946\\
4.8684	0.023523\\
4.8861	0.022165\\
4.9037	0.020868\\
4.9214	0.019632\\
4.9391	0.018454\\
4.9567	0.017332\\
4.9744	0.016265\\
4.9921	0.015251\\
5.0097	0.014288\\
5.0274	0.013375\\
5.0451	0.01251\\
5.0628	0.01169\\
5.0804	0.010915\\
5.0981	0.010182\\
5.1158	0.0094906\\
5.1334	0.0088382\\
5.1511	0.0082235\\
5.1688	0.0076448\\
5.1864	0.0071005\\
5.2041	0.0065892\\
5.2218	0.0061093\\
5.2395	0.0056592\\
5.2571	0.0052377\\
5.2748	0.0048432\\
5.2925	0.0044744\\
5.3101	0.0041299\\
};
\addlegendentry{\RealBLNFitLegend}

      % \input{../../../../.tex/20/KS/SizeDistributionFittingsfig22.tex}
      \legend{}
  
      \nextgroupplot[%
        xmin=3.598,xmax=5.4115,
        % xmin=3.4,xmax=5.4,
        % ymin=0.00076361,ymax=1.0188,
      ] % This file was created by matlab2tikz.
%
\definecolor{mycolor1}{rgb}{0.00000,0.44706,0.69804}%
\definecolor{mycolor2}{rgb}{0.00000,0.44700,0.74100}%
%
\addplot[only marks, mark=*, mark size=1.1180pt, color=mycolor1, fill=mycolor1, opacity=0.60, draw=none, mark options={draw=none,line width=0pt}] table[row sep=crcr]{%
x	y\\
3.6035	0.99468\\
3.6079	0.98404\\
3.6116	0.9734\\
3.6158	0.96277\\
3.6206	0.95213\\
3.6253	0.94149\\
3.6296	0.93085\\
3.6345	0.92021\\
3.6389	0.90957\\
3.6435	0.89894\\
3.6481	0.8883\\
3.6524	0.87766\\
3.6566	0.86702\\
3.6614	0.85638\\
3.6657	0.84574\\
3.6704	0.83511\\
3.6758	0.82447\\
3.6803	0.81383\\
3.6848	0.80319\\
3.6896	0.79255\\
3.6946	0.78191\\
3.6998	0.77128\\
3.7049	0.76064\\
3.7102	0.75\\
3.7157	0.73936\\
3.7209	0.72872\\
3.7264	0.71809\\
3.7344	0.70745\\
3.74	0.69681\\
3.7457	0.68617\\
3.7513	0.67553\\
3.7576	0.66489\\
3.7638	0.65426\\
3.7697	0.64362\\
3.7759	0.63298\\
3.7815	0.62234\\
3.7869	0.6117\\
3.7933	0.60106\\
3.8018	0.59043\\
3.8093	0.57979\\
3.8153	0.56915\\
3.8224	0.55851\\
3.8306	0.54787\\
3.8374	0.53723\\
3.844	0.5266\\
3.8513	0.51596\\
3.8589	0.50532\\
3.8663	0.49468\\
3.8746	0.48404\\
3.8818	0.4734\\
3.8898	0.46277\\
3.8988	0.45213\\
3.9079	0.44149\\
3.9163	0.43085\\
3.9249	0.42021\\
3.9333	0.40957\\
3.942	0.39894\\
3.952	0.3883\\
3.9608	0.37766\\
3.9702	0.36702\\
3.9807	0.35638\\
3.9917	0.34574\\
4.0037	0.33511\\
4.0158	0.32447\\
4.0275	0.31383\\
4.0385	0.30319\\
4.0494	0.29255\\
4.062	0.28191\\
4.0754	0.27128\\
4.0883	0.26064\\
4.1036	0.25\\
4.1184	0.23936\\
4.1331	0.22872\\
4.1471	0.21809\\
4.1608	0.20745\\
4.1772	0.19681\\
4.196	0.18617\\
4.2137	0.17553\\
4.23	0.16489\\
4.2473	0.15426\\
4.2659	0.14362\\
4.2887	0.13298\\
4.3115	0.12234\\
4.3384	0.1117\\
4.3656	0.10106\\
4.394	0.090426\\
4.4263	0.079787\\
4.4592	0.069149\\
4.5016	0.058511\\
4.5488	0.047872\\
4.6129	0.037234\\
4.7095	0.026596\\
4.8162	0.015957\\
5.1746	0.0053191\\
};
\addlegendentry{\ExactEmpiricCCDFLegend}

\addplot [color=mycolor1, only marks, every error bar/.append style={opacity=0.30}, mark=*, mark size=0pt, draw=none, forget plot]
 plot [error bars/.cd, x dir=both, x explicit, error bar style={line width=1pt, color=mycolor1}, error mark options={mark=none,mark size=0pt}]
 table[row sep=crcr, x error plus index=2, x error minus index=3]{%
3.6035	0.99468	0.0043687	0.0043687\\
3.6079	0.98404	0.0043134	0.0043134\\
3.6116	0.9734	0.0042921	0.0042921\\
3.6158	0.96277	0.0043486	0.0043486\\
3.6206	0.95213	0.0045234	0.0045234\\
3.6253	0.94149	0.0045887	0.0045887\\
3.6296	0.93085	0.0046128	0.0046128\\
3.6345	0.92021	0.0046646	0.0046646\\
3.6389	0.90957	0.0047592	0.0047592\\
3.6435	0.89894	0.004712	0.004712\\
3.6481	0.8883	0.0048417	0.0048417\\
3.6524	0.87766	0.004913	0.004913\\
3.6566	0.86702	0.0049333	0.0049333\\
3.6614	0.85638	0.0048893	0.0048893\\
3.6657	0.84574	0.0047743	0.0047743\\
3.6704	0.83511	0.0047267	0.0047267\\
3.6758	0.82447	0.0047531	0.0047531\\
3.6803	0.81383	0.0049722	0.0049722\\
3.6848	0.80319	0.0051901	0.0051901\\
3.6896	0.79255	0.0051486	0.0051486\\
3.6946	0.78191	0.0050613	0.0050613\\
3.6998	0.77128	0.0049859	0.0049859\\
3.7049	0.76064	0.0048672	0.0048672\\
3.7102	0.75	0.0049881	0.0049881\\
3.7157	0.73936	0.0048884	0.0048884\\
3.7209	0.72872	0.0049061	0.0049061\\
3.7264	0.71809	0.0048825	0.0048825\\
3.7344	0.70745	0.0048948	0.0048948\\
3.74	0.69681	0.0048439	0.0048439\\
3.7457	0.68617	0.0048029	0.0048029\\
3.7513	0.67553	0.0047493	0.0047493\\
3.7576	0.66489	0.0047029	0.0047029\\
3.7638	0.65426	0.0047906	0.0047906\\
3.7697	0.64362	0.0048818	0.0048818\\
3.7759	0.63298	0.0048859	0.0048859\\
3.7815	0.62234	0.0048079	0.0048079\\
3.7869	0.6117	0.0048096	0.0048096\\
3.7933	0.60106	0.0047568	0.0047568\\
3.8018	0.59043	0.0046697	0.0046697\\
3.8093	0.57979	0.004961	0.004961\\
3.8153	0.56915	0.0048973	0.0048973\\
3.8224	0.55851	0.0049265	0.0049265\\
3.8306	0.54787	0.0048253	0.0048253\\
3.8374	0.53723	0.0049268	0.0049268\\
3.844	0.5266	0.0047613	0.0047613\\
3.8513	0.51596	0.0048179	0.0048179\\
3.8589	0.50532	0.0048739	0.0048739\\
3.8663	0.49468	0.0046993	0.0046993\\
3.8746	0.48404	0.0049749	0.0049749\\
3.8818	0.4734	0.0052045	0.0052045\\
3.8898	0.46277	0.0049908	0.0049908\\
3.8988	0.45213	0.0050742	0.0050742\\
3.9079	0.44149	0.0051782	0.0051782\\
3.9163	0.43085	0.0052653	0.0052653\\
3.9249	0.42021	0.0052191	0.0052191\\
3.9333	0.40957	0.0050431	0.0050431\\
3.942	0.39894	0.0050961	0.0050961\\
3.952	0.3883	0.0053717	0.0053717\\
3.9608	0.37766	0.0056033	0.0056033\\
3.9702	0.36702	0.0056949	0.0056949\\
3.9807	0.35638	0.0056918	0.0056918\\
3.9917	0.34574	0.005667	0.005667\\
4.0037	0.33511	0.0055617	0.0055617\\
4.0158	0.32447	0.005337	0.005337\\
4.0275	0.31383	0.0053077	0.0053077\\
4.0385	0.30319	0.0051475	0.0051475\\
4.0494	0.29255	0.0051079	0.0051079\\
4.062	0.28191	0.0051371	0.0051371\\
4.0754	0.27128	0.0054134	0.0054134\\
4.0883	0.26064	0.0057357	0.0057357\\
4.1036	0.25	0.0058334	0.0058334\\
4.1184	0.23936	0.0059011	0.0059011\\
4.1331	0.22872	0.0059871	0.0059871\\
4.1471	0.21809	0.0064116	0.0064116\\
4.1608	0.20745	0.0066205	0.0066205\\
4.1772	0.19681	0.006858	0.006858\\
4.196	0.18617	0.0069585	0.0069585\\
4.2137	0.17553	0.0077849	0.0077849\\
4.23	0.16489	0.0082238	0.0082238\\
4.2473	0.15426	0.0081301	0.0081301\\
4.2659	0.14362	0.0082035	0.0082035\\
4.2887	0.13298	0.008017	0.008017\\
4.3115	0.12234	0.0083363	0.0083363\\
4.3384	0.1117	0.008877	0.008877\\
4.3656	0.10106	0.0079618	0.0079618\\
4.394	0.090426	0.0078787	0.0078787\\
4.4263	0.079787	0.0079578	0.0079578\\
4.4592	0.069149	0.0087748	0.0087748\\
4.5016	0.058511	0.010297	0.010297\\
4.5488	0.047872	0.010617	0.010617\\
4.6129	0.037234	0.011265	0.011265\\
4.7095	0.026596	0.013471	0.013471\\
4.8162	0.015957	0.014317	0.014317\\
5.1746	0.0053191	0.012259	0.012259\\
};
\addplot [color=mycolor1, mark repeat=9, mark phase=1, mark=+, mark options={solid, mycolor1}]
  table[row sep=crcr]{%
3.6023	0.97605\\
3.6198	0.94435\\
3.6372	0.90323\\
3.6547	0.86092\\
3.6721	0.81854\\
3.6896	0.77787\\
3.707	0.73923\\
3.7245	0.70251\\
3.742	0.66762\\
3.7594	0.63448\\
3.7769	0.60298\\
3.7943	0.57305\\
3.8118	0.54461\\
3.8292	0.51759\\
3.8467	0.49191\\
3.8642	0.46751\\
3.8816	0.44433\\
3.8991	0.4223\\
3.9165	0.40136\\
3.934	0.38147\\
3.9515	0.36256\\
3.9689	0.34459\\
3.9864	0.32752\\
4.0038	0.3113\\
4.0213	0.29588\\
4.0387	0.28123\\
4.0562	0.26731\\
4.0737	0.25407\\
4.0911	0.2415\\
4.1086	0.22955\\
4.126	0.21819\\
4.1435	0.2074\\
4.1609	0.19714\\
4.1784	0.18739\\
4.1959	0.17812\\
4.2133	0.16932\\
4.2308	0.16095\\
4.2482	0.15299\\
4.2657	0.14543\\
4.2832	0.13825\\
4.3006	0.13142\\
4.3181	0.12493\\
4.3355	0.11876\\
4.353	0.11289\\
4.3704	0.10732\\
4.3879	0.10202\\
4.4054	0.096989\\
4.4228	0.092203\\
4.4403	0.087654\\
4.4577	0.08333\\
4.4752	0.07922\\
4.4926	0.075314\\
4.5101	0.0716\\
4.5276	0.06807\\
4.545	0.064715\\
4.5625	0.061526\\
4.5799	0.058494\\
4.5974	0.055612\\
4.6149	0.052873\\
4.6323	0.050269\\
4.6498	0.047793\\
4.6672	0.04544\\
4.6847	0.043203\\
4.7021	0.041077\\
4.7196	0.039055\\
4.7371	0.037133\\
4.7545	0.035306\\
4.772	0.03357\\
4.7894	0.031919\\
4.8069	0.030349\\
4.8243	0.028857\\
4.8418	0.027438\\
4.8593	0.026089\\
4.8767	0.024807\\
4.8942	0.023588\\
4.9116	0.022429\\
4.9291	0.021327\\
4.9466	0.02028\\
4.964	0.019284\\
4.9815	0.018337\\
4.9989	0.017437\\
5.0164	0.016581\\
5.0338	0.015767\\
5.0513	0.014993\\
5.0688	0.014257\\
5.0862	0.013558\\
5.1037	0.012893\\
5.1211	0.012261\\
5.1386	0.011659\\
5.156	0.011088\\
5.1735	0.010544\\
5.191	0.010027\\
5.2084	0.0095358\\
5.2259	0.0090685\\
5.2433	0.0086242\\
5.2608	0.0082017\\
5.2783	0.0078\\
5.2957	0.007418\\
5.3132	0.0070547\\
5.3306	0.0067093\\
};
\addlegendentry{\ExactParetoFitLegend}


\addplot[area legend, draw=none, fill=mycolor1, fill opacity=0.15, postaction={pattern={Lines[angle=90,distance={5pt},line width={0.5pt}]}, pattern color=mycolor1!60}, forget plot]
table[row sep=crcr] {%
x	y\\
3.6023	0.98202\\
3.6198	0.95296\\
3.6372	0.91279\\
3.6547	0.87077\\
3.6721	0.8277\\
3.6896	0.78616\\
3.707	0.74672\\
3.7245	0.70928\\
3.742	0.67373\\
3.7594	0.63997\\
3.7769	0.60792\\
3.7943	0.5775\\
3.8118	0.54861\\
3.8292	0.52119\\
3.8467	0.49516\\
3.8642	0.47045\\
3.8816	0.447\\
3.8991	0.42474\\
3.9165	0.40361\\
3.934	0.38356\\
3.9515	0.36454\\
3.9689	0.34648\\
3.9864	0.32934\\
4.0038	0.31306\\
4.0213	0.29762\\
4.0387	0.28295\\
4.0562	0.26901\\
4.0737	0.25578\\
4.0911	0.2432\\
4.1086	0.23126\\
4.126	0.2199\\
4.1435	0.20911\\
4.1609	0.19885\\
4.1784	0.18911\\
4.1959	0.17984\\
4.2133	0.17103\\
4.2308	0.16265\\
4.2482	0.15469\\
4.2657	0.14712\\
4.2832	0.13992\\
4.3006	0.13307\\
4.3181	0.12656\\
4.3355	0.12038\\
4.353	0.11449\\
4.3704	0.10889\\
4.3879	0.10357\\
4.4054	0.098512\\
4.4228	0.093699\\
4.4403	0.089122\\
4.4577	0.08477\\
4.4752	0.080631\\
4.4926	0.076695\\
4.5101	0.072951\\
4.5276	0.069391\\
4.545	0.066005\\
4.5625	0.062785\\
4.5799	0.059722\\
4.5974	0.056809\\
4.6149	0.054039\\
4.6323	0.051404\\
4.6498	0.048898\\
4.6672	0.046515\\
4.6847	0.044248\\
4.7021	0.042092\\
4.7196	0.040041\\
4.7371	0.03809\\
4.7545	0.036235\\
4.772	0.03447\\
4.7894	0.032792\\
4.8069	0.031195\\
4.8243	0.029677\\
4.8418	0.028232\\
4.8593	0.026858\\
4.8767	0.025551\\
4.8942	0.024308\\
4.9116	0.023126\\
4.9291	0.022001\\
4.9466	0.020931\\
4.964	0.019913\\
4.9815	0.018945\\
4.9989	0.018024\\
5.0164	0.017148\\
5.0338	0.016315\\
5.0513	0.015522\\
5.0688	0.014768\\
5.0862	0.01405\\
5.1037	0.013368\\
5.1211	0.012719\\
5.1386	0.012101\\
5.156	0.011513\\
5.1735	0.010954\\
5.191	0.010423\\
5.2084	0.0099168\\
5.2259	0.0094356\\
5.2433	0.0089777\\
5.2608	0.0085422\\
5.2783	0.0081278\\
5.2957	0.0077336\\
5.3132	0.0073586\\
5.3306	0.0070017\\
5.3306	0.0064169\\
5.3132	0.0067509\\
5.2957	0.0071023\\
5.2783	0.0074721\\
5.2608	0.0078612\\
5.2433	0.0082706\\
5.2259	0.0087015\\
5.2084	0.0091548\\
5.191	0.0096318\\
5.1735	0.010134\\
5.156	0.010662\\
5.1386	0.011218\\
5.1211	0.011803\\
5.1037	0.012418\\
5.0862	0.013066\\
5.0688	0.013747\\
5.0513	0.014464\\
5.0338	0.015219\\
5.0164	0.016013\\
4.9989	0.016849\\
4.9815	0.017729\\
4.964	0.018654\\
4.9466	0.019629\\
4.9291	0.020654\\
4.9116	0.021733\\
4.8942	0.022868\\
4.8767	0.024063\\
4.8593	0.02532\\
4.8418	0.026644\\
4.8243	0.028037\\
4.8069	0.029502\\
4.7894	0.031045\\
4.772	0.032669\\
4.7545	0.034378\\
4.7371	0.036176\\
4.7196	0.038069\\
4.7021	0.040062\\
4.6847	0.042159\\
4.6672	0.044366\\
4.6498	0.046689\\
4.6323	0.049134\\
4.6149	0.051707\\
4.5974	0.054415\\
4.5799	0.057266\\
4.5625	0.060267\\
4.545	0.063426\\
4.5276	0.06675\\
4.5101	0.070249\\
4.4926	0.073932\\
4.4752	0.077809\\
4.4577	0.08189\\
4.4403	0.086185\\
4.4228	0.090706\\
4.4054	0.095466\\
4.3879	0.10047\\
4.3704	0.10575\\
4.353	0.1113\\
4.3355	0.11714\\
4.3181	0.12329\\
4.3006	0.12976\\
4.2832	0.13658\\
4.2657	0.14375\\
4.2482	0.1513\\
4.2308	0.15924\\
4.2133	0.1676\\
4.1959	0.17641\\
4.1784	0.18567\\
4.1609	0.19542\\
4.1435	0.20568\\
4.126	0.21648\\
4.1086	0.22784\\
4.0911	0.23979\\
4.0737	0.25237\\
4.0562	0.2656\\
4.0387	0.27951\\
4.0213	0.29415\\
4.0038	0.30953\\
3.9864	0.32571\\
3.9689	0.34271\\
3.9515	0.36058\\
3.934	0.37937\\
3.9165	0.39911\\
3.8991	0.41985\\
3.8816	0.44166\\
3.8642	0.46457\\
3.8467	0.48866\\
3.8292	0.51399\\
3.8118	0.54061\\
3.7943	0.5686\\
3.7769	0.59803\\
3.7594	0.62898\\
3.742	0.66152\\
3.7245	0.69574\\
3.707	0.73173\\
3.6896	0.76958\\
3.6721	0.80938\\
3.6547	0.85108\\
3.6372	0.89366\\
3.6198	0.93575\\
3.6023	0.97007\\
}--cycle;
\addplot [color=mycolor1, mark repeat=9, mark phase=1, mark=x, mark options={solid, mycolor1}]
  table[row sep=crcr]{%
3.6023	1.0112\\
3.6198	0.96682\\
3.6372	0.92399\\
3.6547	0.88271\\
3.6721	0.84295\\
3.6896	0.80469\\
3.707	0.76792\\
3.7245	0.73259\\
3.742	0.69867\\
3.7594	0.66614\\
3.7769	0.63495\\
3.7943	0.60506\\
3.8118	0.57644\\
3.8292	0.54904\\
3.8467	0.52282\\
3.8642	0.49773\\
3.8816	0.47374\\
3.8991	0.4508\\
3.9165	0.42886\\
3.934	0.4079\\
3.9515	0.38786\\
3.9689	0.36871\\
3.9864	0.35041\\
4.0038	0.33293\\
4.0213	0.31622\\
4.0387	0.30025\\
4.0562	0.285\\
4.0737	0.27042\\
4.0911	0.2565\\
4.1086	0.24319\\
4.126	0.23048\\
4.1435	0.21834\\
4.1609	0.20674\\
4.1784	0.19566\\
4.1959	0.18508\\
4.2133	0.17498\\
4.2308	0.16533\\
4.2482	0.15613\\
4.2657	0.14735\\
4.2832	0.13898\\
4.3006	0.13099\\
4.3181	0.12338\\
4.3355	0.11613\\
4.353	0.10923\\
4.3704	0.10266\\
4.3879	0.096417\\
4.4054	0.09048\\
4.4228	0.08484\\
4.4403	0.079486\\
4.4577	0.074408\\
4.4752	0.069595\\
4.4926	0.065037\\
4.5101	0.060724\\
4.5276	0.056647\\
4.545	0.052795\\
4.5625	0.049161\\
4.5799	0.045734\\
4.5974	0.042507\\
4.6149	0.03947\\
4.6323	0.036615\\
4.6498	0.033935\\
4.6672	0.031421\\
4.6847	0.029065\\
4.7021	0.02686\\
4.7196	0.024799\\
4.7371	0.022873\\
4.7545	0.021077\\
4.772	0.019403\\
4.7894	0.017845\\
4.8069	0.016396\\
4.8243	0.015051\\
4.8418	0.013803\\
4.8593	0.012646\\
4.8767	0.011575\\
4.8942	0.010585\\
4.9116	0.0096706\\
4.9291	0.0088269\\
4.9466	0.0080493\\
4.964	0.0073334\\
4.9815	0.0066751\\
4.9989	0.0060703\\
5.0164	0.0055152\\
5.0338	0.0050065\\
5.0513	0.0045405\\
5.0688	0.0041143\\
5.0862	0.0037248\\
5.1037	0.0033691\\
5.1211	0.0030448\\
5.1386	0.0027493\\
5.156	0.0024803\\
5.1735	0.0022358\\
5.191	0.0020136\\
5.2084	0.0018119\\
5.2259	0.0016291\\
5.2433	0.0014635\\
5.2608	0.0013136\\
5.2783	0.0011782\\
5.2957	0.0010558\\
5.3132	0.00094534\\
5.3306	0.00084577\\
};
\addlegendentry{\ExactBLNFitLegend}


\addplot[area legend, draw=none, fill=mycolor1, fill opacity=0.15, postaction={pattern={Lines[angle=45,distance={5pt},line width={0.5pt}]}, pattern color=mycolor1!60}, forget plot]
table[row sep=crcr] {%
x	y\\
3.6023	1.0188\\
3.6198	0.97418\\
3.6372	0.93115\\
3.6547	0.88966\\
3.6721	0.84969\\
3.6896	0.81124\\
3.707	0.77426\\
3.7245	0.73872\\
3.742	0.70461\\
3.7594	0.67188\\
3.7769	0.64049\\
3.7943	0.6104\\
3.8118	0.58158\\
3.8292	0.55397\\
3.8467	0.52755\\
3.8642	0.50226\\
3.8816	0.47807\\
3.8991	0.45494\\
3.9165	0.43281\\
3.934	0.41166\\
3.9515	0.39144\\
3.9689	0.37211\\
3.9864	0.35364\\
4.0038	0.336\\
4.0213	0.31914\\
4.0387	0.30303\\
4.0562	0.28765\\
4.0737	0.27296\\
4.0911	0.25893\\
4.1086	0.24554\\
4.126	0.23275\\
4.1435	0.22054\\
4.1609	0.20888\\
4.1784	0.19775\\
4.1959	0.18713\\
4.2133	0.17699\\
4.2308	0.16732\\
4.2482	0.15808\\
4.2657	0.14927\\
4.2832	0.14087\\
4.3006	0.13285\\
4.3181	0.12521\\
4.3355	0.11793\\
4.353	0.11099\\
4.3704	0.10438\\
4.3879	0.098092\\
4.4054	0.092109\\
4.4228	0.086422\\
4.4403	0.081019\\
4.4577	0.075889\\
4.4752	0.071024\\
4.4926	0.066413\\
4.5101	0.062047\\
4.5276	0.057916\\
4.545	0.054012\\
4.5625	0.050325\\
4.5799	0.046848\\
4.5974	0.04357\\
4.6149	0.040486\\
4.6323	0.037585\\
4.6498	0.03486\\
4.6672	0.032303\\
4.6847	0.029906\\
4.7021	0.027662\\
4.7196	0.025563\\
4.7371	0.023601\\
4.7545	0.021771\\
4.772	0.020064\\
4.7894	0.018474\\
4.8069	0.016995\\
4.8243	0.015621\\
4.8418	0.014344\\
4.8593	0.01316\\
4.8767	0.012063\\
4.8942	0.011048\\
4.9116	0.010109\\
4.9291	0.0092419\\
4.9466	0.0084415\\
4.964	0.0077037\\
4.9815	0.0070242\\
4.9989	0.006399\\
5.0164	0.0058244\\
5.0338	0.0052968\\
5.0513	0.0048128\\
5.0688	0.0043693\\
5.0862	0.0039633\\
5.1037	0.003592\\
5.1211	0.0032527\\
5.1386	0.002943\\
5.156	0.0026606\\
5.1735	0.0024033\\
5.191	0.0021691\\
5.2084	0.0019561\\
5.2259	0.0017626\\
5.2433	0.001587\\
5.2608	0.0014277\\
5.2783	0.0012834\\
5.2957	0.0011528\\
5.3132	0.0010347\\
5.3306	0.00092792\\
5.3306	0.00076361\\
5.3132	0.00085601\\
5.2957	0.00095876\\
5.2783	0.0010729\\
5.2608	0.0011996\\
5.2433	0.00134\\
5.2259	0.0014956\\
5.2084	0.0016678\\
5.191	0.0018581\\
5.1735	0.0020682\\
5.156	0.0023001\\
5.1386	0.0025556\\
5.1211	0.0028369\\
5.1037	0.0031463\\
5.0862	0.0034862\\
5.0688	0.0038592\\
5.0513	0.0042682\\
5.0338	0.0047161\\
5.0164	0.0052061\\
4.9989	0.0057415\\
4.9815	0.0063259\\
4.964	0.0069631\\
4.9466	0.0076571\\
4.9291	0.0084119\\
4.9116	0.0092321\\
4.8942	0.010122\\
4.8767	0.011087\\
4.8593	0.012132\\
4.8418	0.013261\\
4.8243	0.014481\\
4.8069	0.015798\\
4.7894	0.017216\\
4.772	0.018743\\
4.7545	0.020383\\
4.7371	0.022145\\
4.7196	0.024035\\
4.7021	0.026059\\
4.6847	0.028224\\
4.6672	0.030539\\
4.6498	0.03301\\
4.6323	0.035646\\
4.6149	0.038454\\
4.5974	0.041443\\
4.5799	0.04462\\
4.5625	0.047996\\
4.545	0.051579\\
4.5276	0.055377\\
4.5101	0.059402\\
4.4926	0.063661\\
4.4752	0.068166\\
4.4577	0.072927\\
4.4403	0.077954\\
4.4228	0.083258\\
4.4054	0.08885\\
4.3879	0.094742\\
4.3704	0.10095\\
4.353	0.10747\\
4.3355	0.11434\\
4.3181	0.12155\\
4.3006	0.12913\\
4.2832	0.13708\\
4.2657	0.14543\\
4.2482	0.15418\\
4.2308	0.16335\\
4.2133	0.17296\\
4.1959	0.18302\\
4.1784	0.19356\\
4.1609	0.20459\\
4.1435	0.21613\\
4.126	0.22821\\
4.1086	0.24084\\
4.0911	0.25406\\
4.0737	0.26788\\
4.0562	0.28234\\
4.0387	0.29747\\
4.0213	0.3133\\
4.0038	0.32986\\
3.9864	0.34718\\
3.9689	0.36531\\
3.9515	0.38429\\
3.934	0.40414\\
3.9165	0.42491\\
3.8991	0.44665\\
3.8816	0.4694\\
3.8642	0.4932\\
3.8467	0.51808\\
3.8292	0.5441\\
3.8118	0.57131\\
3.7943	0.59973\\
3.7769	0.62941\\
3.7594	0.6604\\
3.742	0.69274\\
3.7245	0.72645\\
3.707	0.76158\\
3.6896	0.79815\\
3.6721	0.8362\\
3.6547	0.87576\\
3.6372	0.91683\\
3.6198	0.95945\\
3.6023	1.0036\\
}--cycle;
      % \input{../../../../.tex/20/KS/SizeDistributionFittingsfig23.tex}
      \legend{}
  
      \nextgroupplot[%
        xmin=3.6008,xmax=5.4115,
        % xmin=3.6,xmax=5.4,
        % ymin=0.00076295,ymax=1.01,
      ] % This file was created by matlab2tikz.
%
\definecolor{mycolor1}{rgb}{0.83529,0.36863,0.00000}%
\definecolor{mycolor2}{rgb}{0.00000,0.44700,0.74100}%
%
\addplot[only marks, mark=*, mark size=1.1180pt, color=mycolor1, fill=mycolor1, opacity=0.60, draw=none, mark options={draw=none,line width=0pt}] table[row sep=crcr]{%
x	y\\
3.608	0.99468\\
3.612	0.98404\\
3.6149	0.9734\\
3.6193	0.96277\\
3.6236	0.95213\\
3.6275	0.94149\\
3.6307	0.93085\\
3.6349	0.92021\\
3.6394	0.90957\\
3.6443	0.89894\\
3.6489	0.8883\\
3.6535	0.87766\\
3.6574	0.86702\\
3.6619	0.85638\\
3.6661	0.84574\\
3.6707	0.83511\\
3.6759	0.82447\\
3.6809	0.81383\\
3.6863	0.80319\\
3.6908	0.79255\\
3.6971	0.78191\\
3.7017	0.77128\\
3.7069	0.76064\\
3.7123	0.75\\
3.7172	0.73936\\
3.722	0.72872\\
3.7277	0.71809\\
3.734	0.70745\\
3.7392	0.69681\\
3.7453	0.68617\\
3.7511	0.67553\\
3.7578	0.66489\\
3.7629	0.65426\\
3.7687	0.64362\\
3.7756	0.63298\\
3.7819	0.62234\\
3.7891	0.6117\\
3.7954	0.60106\\
3.8017	0.59043\\
3.809	0.57979\\
3.8157	0.56915\\
3.8216	0.55851\\
3.8279	0.54787\\
3.835	0.53723\\
3.843	0.5266\\
3.8508	0.51596\\
3.8592	0.50532\\
3.8687	0.49468\\
3.8767	0.48404\\
3.8837	0.4734\\
3.8915	0.46277\\
3.8992	0.45213\\
3.9079	0.44149\\
3.9154	0.43085\\
3.926	0.42021\\
3.9348	0.40957\\
3.9442	0.39894\\
3.9526	0.3883\\
3.9638	0.37766\\
3.9717	0.36702\\
3.9827	0.35638\\
3.9938	0.34574\\
4.0041	0.33511\\
4.0134	0.32447\\
4.0267	0.31383\\
4.0386	0.30319\\
4.0508	0.29255\\
4.0624	0.28191\\
4.0725	0.27128\\
4.0888	0.26064\\
4.1049	0.25\\
4.117	0.23936\\
4.1302	0.22872\\
4.1455	0.21809\\
4.1599	0.20745\\
4.1742	0.19681\\
4.1902	0.18617\\
4.2074	0.17553\\
4.2233	0.16489\\
4.2398	0.15426\\
4.2602	0.14362\\
4.2789	0.13298\\
4.2994	0.12234\\
4.3276	0.1117\\
4.3546	0.10106\\
4.3817	0.090426\\
4.4171	0.079787\\
4.4522	0.069149\\
4.497	0.058511\\
4.5578	0.047872\\
4.6192	0.037234\\
4.7102	0.026596\\
4.8182	0.015957\\
5.1809	0.0053191\\
};
\addlegendentry{\ApproxEmpiricCCDFLegend}

\addplot [color=mycolor1, only marks, every error bar/.append style={opacity=0.30}, mark=*, mark size=0pt, draw=none, forget plot]
 plot [error bars/.cd, x dir=both, x explicit, error bar style={line width=1pt, color=mycolor1}, error mark options={mark=none,mark size=0pt}]
 table[row sep=crcr, x error plus index=2, x error minus index=3]{%
3.608	0.99468	0.0040359	0.0040359\\
3.612	0.98404	0.0040852	0.0040852\\
3.6149	0.9734	0.0041689	0.0041689\\
3.6193	0.96277	0.0042144	0.0042144\\
3.6236	0.95213	0.0042315	0.0042315\\
3.6275	0.94149	0.0042316	0.0042316\\
3.6307	0.93085	0.0042656	0.0042656\\
3.6349	0.92021	0.0042076	0.0042076\\
3.6394	0.90957	0.0042936	0.0042936\\
3.6443	0.89894	0.0042523	0.0042523\\
3.6489	0.8883	0.0042739	0.0042739\\
3.6535	0.87766	0.004405	0.004405\\
3.6574	0.86702	0.0044875	0.0044875\\
3.6619	0.85638	0.0044589	0.0044589\\
3.6661	0.84574	0.0043904	0.0043904\\
3.6707	0.83511	0.0043459	0.0043459\\
3.6759	0.82447	0.0044445	0.0044445\\
3.6809	0.81383	0.0044758	0.0044758\\
3.6863	0.80319	0.0044367	0.0044367\\
3.6908	0.79255	0.0043683	0.0043683\\
3.6971	0.78191	0.0043662	0.0043662\\
3.7017	0.77128	0.0041983	0.0041983\\
3.7069	0.76064	0.0042572	0.0042572\\
3.7123	0.75	0.0042455	0.0042455\\
3.7172	0.73936	0.0043061	0.0043061\\
3.722	0.72872	0.0041243	0.0041243\\
3.7277	0.71809	0.0042463	0.0042463\\
3.734	0.70745	0.0043807	0.0043807\\
3.7392	0.69681	0.0043561	0.0043561\\
3.7453	0.68617	0.0044568	0.0044568\\
3.7511	0.67553	0.0042342	0.0042342\\
3.7578	0.66489	0.0042251	0.0042251\\
3.7629	0.65426	0.0043299	0.0043299\\
3.7687	0.64362	0.0041462	0.0041462\\
3.7756	0.63298	0.0045377	0.0045377\\
3.7819	0.62234	0.0044165	0.0044165\\
3.7891	0.6117	0.0044983	0.0044983\\
3.7954	0.60106	0.0044046	0.0044046\\
3.8017	0.59043	0.0044956	0.0044956\\
3.809	0.57979	0.0046597	0.0046597\\
3.8157	0.56915	0.0048958	0.0048958\\
3.8216	0.55851	0.0048694	0.0048694\\
3.8279	0.54787	0.0047418	0.0047418\\
3.835	0.53723	0.0047159	0.0047159\\
3.843	0.5266	0.0047471	0.0047471\\
3.8508	0.51596	0.0047735	0.0047735\\
3.8592	0.50532	0.0047167	0.0047167\\
3.8687	0.49468	0.0047754	0.0047754\\
3.8767	0.48404	0.005157	0.005157\\
3.8837	0.4734	0.0052409	0.0052409\\
3.8915	0.46277	0.0054523	0.0054523\\
3.8992	0.45213	0.0055911	0.0055911\\
3.9079	0.44149	0.0055855	0.0055855\\
3.9154	0.43085	0.0056618	0.0056618\\
3.926	0.42021	0.0059361	0.0059361\\
3.9348	0.40957	0.0058921	0.0058921\\
3.9442	0.39894	0.0061974	0.0061974\\
3.9526	0.3883	0.0060931	0.0060931\\
3.9638	0.37766	0.0058931	0.0058931\\
3.9717	0.36702	0.0058103	0.0058103\\
3.9827	0.35638	0.0056547	0.0056547\\
3.9938	0.34574	0.0056375	0.0056375\\
4.0041	0.33511	0.0055906	0.0055906\\
4.0134	0.32447	0.0058442	0.0058442\\
4.0267	0.31383	0.0056071	0.0056071\\
4.0386	0.30319	0.0059954	0.0059954\\
4.0508	0.29255	0.0056373	0.0056373\\
4.0624	0.28191	0.0058582	0.0058582\\
4.0725	0.27128	0.0058151	0.0058151\\
4.0888	0.26064	0.005724	0.005724\\
4.1049	0.25	0.0057219	0.0057219\\
4.117	0.23936	0.0059065	0.0059065\\
4.1302	0.22872	0.0059432	0.0059432\\
4.1455	0.21809	0.0060705	0.0060705\\
4.1599	0.20745	0.0060167	0.0060167\\
4.1742	0.19681	0.005967	0.005967\\
4.1902	0.18617	0.0058452	0.0058452\\
4.2074	0.17553	0.0059932	0.0059932\\
4.2233	0.16489	0.0063192	0.0063192\\
4.2398	0.15426	0.0075653	0.0075653\\
4.2602	0.14362	0.0076698	0.0076698\\
4.2789	0.13298	0.0079775	0.0079775\\
4.2994	0.12234	0.0083583	0.0083583\\
4.3276	0.1117	0.008497	0.008497\\
4.3546	0.10106	0.0080207	0.0080207\\
4.3817	0.090426	0.0088382	0.0088382\\
4.4171	0.079787	0.0098774	0.0098774\\
4.4522	0.069149	0.010243	0.010243\\
4.497	0.058511	0.011325	0.011325\\
4.5578	0.047872	0.013129	0.013129\\
4.6192	0.037234	0.011595	0.011595\\
4.7102	0.026596	0.014518	0.014518\\
4.8182	0.015957	0.018109	0.018109\\
5.1809	0.0053191	0.014189	0.014189\\
};
\addplot [color=mycolor1, mark repeat=9, mark phase=1, mark=+, mark options={solid, mycolor1}]
  table[row sep=crcr]{%
3.606	0.97811\\
3.6234	0.94606\\
3.6408	0.90438\\
3.6583	0.85992\\
3.6757	0.81679\\
3.6931	0.77583\\
3.7105	0.73693\\
3.7279	0.69998\\
3.7454	0.6649\\
3.7628	0.63157\\
3.7802	0.59992\\
3.7976	0.56986\\
3.815	0.54131\\
3.8325	0.51419\\
3.8499	0.48844\\
3.8673	0.46398\\
3.8847	0.44075\\
3.9021	0.41868\\
3.9196	0.39772\\
3.937	0.37782\\
3.9544	0.35891\\
3.9718	0.34095\\
3.9893	0.32389\\
4.0067	0.30769\\
4.0241	0.2923\\
4.0415	0.27768\\
4.0589	0.2638\\
4.0764	0.25061\\
4.0938	0.23808\\
4.1112	0.22618\\
4.1286	0.21488\\
4.146	0.20414\\
4.1635	0.19394\\
4.1809	0.18425\\
4.1983	0.17505\\
4.2157	0.16631\\
4.2331	0.158\\
4.2506	0.15011\\
4.268	0.14262\\
4.2854	0.1355\\
4.3028	0.12874\\
4.3202	0.12231\\
4.3377	0.11621\\
4.3551	0.11041\\
4.3725	0.10491\\
4.3899	0.099675\\
4.4073	0.094705\\
4.4248	0.089983\\
4.4422	0.085497\\
4.4596	0.081236\\
4.477	0.077187\\
4.4944	0.073341\\
4.5119	0.069686\\
4.5293	0.066215\\
4.5467	0.062917\\
4.5641	0.059783\\
4.5815	0.056806\\
4.599	0.053978\\
4.6164	0.05129\\
4.6338	0.048737\\
4.6512	0.046312\\
4.6686	0.044007\\
4.6861	0.041817\\
4.7035	0.039737\\
4.7209	0.03776\\
4.7383	0.035882\\
4.7558	0.034098\\
4.7732	0.032402\\
4.7906	0.030791\\
4.808	0.029261\\
4.8254	0.027806\\
4.8429	0.026425\\
4.8603	0.025112\\
4.8777	0.023864\\
4.8951	0.022678\\
4.9125	0.021552\\
4.93	0.020482\\
4.9474	0.019465\\
4.9648	0.018498\\
4.9822	0.01758\\
4.9996	0.016707\\
5.0171	0.015878\\
5.0345	0.01509\\
5.0519	0.014342\\
5.0693	0.01363\\
5.0867	0.012954\\
5.1042	0.012312\\
5.1216	0.011701\\
5.139	0.011121\\
5.1564	0.01057\\
5.1738	0.010046\\
5.1913	0.0095477\\
5.2087	0.0090746\\
5.2261	0.008625\\
5.2435	0.0081977\\
5.2609	0.0077916\\
5.2784	0.0074057\\
5.2958	0.0070389\\
5.3132	0.0066904\\
5.3306	0.0063592\\
};
\addlegendentry{\ApproxParetoFitLegend}


\addplot[area legend, draw=none, fill=mycolor1, fill opacity=0.15, postaction={pattern={Lines[angle=90,distance={5pt},line width={0.5pt}]}, pattern color=mycolor1!60}, forget plot]
table[row sep=crcr] {%
x	y\\
3.606	0.98437\\
3.6234	0.95474\\
3.6408	0.91396\\
3.6583	0.86904\\
3.6757	0.82509\\
3.6931	0.78338\\
3.7105	0.74379\\
3.7279	0.70622\\
3.7454	0.67056\\
3.7628	0.63672\\
3.7802	0.6046\\
3.7976	0.57412\\
3.815	0.5452\\
3.8325	0.51775\\
3.8499	0.4917\\
3.8673	0.46698\\
3.8847	0.44353\\
3.9021	0.42127\\
3.9196	0.40016\\
3.937	0.38012\\
3.9544	0.3611\\
3.9718	0.34306\\
3.9893	0.32593\\
4.0067	0.30968\\
4.0241	0.29424\\
4.0415	0.27959\\
4.0589	0.26568\\
4.0764	0.25247\\
4.0938	0.23993\\
4.1112	0.22801\\
4.1286	0.2167\\
4.146	0.20595\\
4.1635	0.19573\\
4.1809	0.18603\\
4.1983	0.17681\\
4.2157	0.16805\\
4.2331	0.15973\\
4.2506	0.15182\\
4.268	0.14431\\
4.2854	0.13717\\
4.3028	0.13038\\
4.3202	0.12394\\
4.3377	0.11781\\
4.3551	0.11198\\
4.3725	0.10645\\
4.3899	0.10119\\
4.4073	0.096188\\
4.4248	0.091436\\
4.4422	0.08692\\
4.4596	0.082628\\
4.477	0.078548\\
4.4944	0.074671\\
4.5119	0.070985\\
4.5293	0.067482\\
4.5467	0.064152\\
4.5641	0.060988\\
4.5815	0.057979\\
4.599	0.05512\\
4.6164	0.052402\\
4.6338	0.049818\\
4.6512	0.047362\\
4.6686	0.045028\\
4.6861	0.042809\\
4.7035	0.040699\\
4.7209	0.038694\\
4.7383	0.036788\\
4.7558	0.034976\\
4.7732	0.033254\\
4.7906	0.031616\\
4.808	0.03006\\
4.8254	0.02858\\
4.8429	0.027173\\
4.8603	0.025836\\
4.8777	0.024565\\
4.8951	0.023356\\
4.9125	0.022207\\
4.93	0.021115\\
4.9474	0.020077\\
4.9648	0.019089\\
4.9822	0.018151\\
4.9996	0.017259\\
5.0171	0.01641\\
5.0345	0.015604\\
5.0519	0.014837\\
5.0693	0.014108\\
5.0867	0.013415\\
5.1042	0.012756\\
5.1216	0.01213\\
5.139	0.011534\\
5.1564	0.010968\\
5.1738	0.010429\\
5.1913	0.0099173\\
5.2087	0.0094305\\
5.2261	0.0089678\\
5.2435	0.0085277\\
5.2609	0.0081094\\
5.2784	0.0077116\\
5.2958	0.0073334\\
5.3132	0.0069737\\
5.3306	0.0066318\\
5.3306	0.0060865\\
5.3132	0.006407\\
5.2958	0.0067445\\
5.2784	0.0070998\\
5.2609	0.0074738\\
5.2435	0.0078676\\
5.2261	0.0082822\\
5.2087	0.0087187\\
5.1913	0.0091782\\
5.1738	0.0096621\\
5.1564	0.010171\\
5.139	0.010708\\
5.1216	0.011273\\
5.1042	0.011867\\
5.0867	0.012493\\
5.0693	0.013152\\
5.0519	0.013846\\
5.0345	0.014577\\
5.0171	0.015346\\
4.9996	0.016156\\
4.9822	0.017009\\
4.9648	0.017907\\
4.9474	0.018853\\
4.93	0.019848\\
4.9125	0.020897\\
4.8951	0.022001\\
4.8777	0.023163\\
4.8603	0.024387\\
4.8429	0.025676\\
4.8254	0.027033\\
4.808	0.028462\\
4.7906	0.029967\\
4.7732	0.031551\\
4.7558	0.033219\\
4.7383	0.034976\\
4.7209	0.036826\\
4.7035	0.038774\\
4.6861	0.040826\\
4.6686	0.042986\\
4.6512	0.045261\\
4.6338	0.047656\\
4.6164	0.050179\\
4.599	0.052835\\
4.5815	0.055633\\
4.5641	0.058579\\
4.5467	0.061681\\
4.5293	0.064948\\
4.5119	0.068388\\
4.4944	0.072011\\
4.477	0.075826\\
4.4596	0.079843\\
4.4422	0.084074\\
4.4248	0.08853\\
4.4073	0.093222\\
4.3899	0.098163\\
4.3725	0.10337\\
4.3551	0.10885\\
4.3377	0.11462\\
4.3202	0.12069\\
4.3028	0.12709\\
4.2854	0.13383\\
4.268	0.14093\\
4.2506	0.1484\\
4.2331	0.15627\\
4.2157	0.16456\\
4.1983	0.17328\\
4.1809	0.18247\\
4.1635	0.19215\\
4.146	0.20233\\
4.1286	0.21306\\
4.1112	0.22435\\
4.0938	0.23623\\
4.0764	0.24874\\
4.0589	0.26191\\
4.0415	0.27577\\
4.0241	0.29035\\
4.0067	0.3057\\
3.9893	0.32185\\
3.9718	0.33884\\
3.9544	0.35671\\
3.937	0.37551\\
3.9196	0.39529\\
3.9021	0.41609\\
3.8847	0.43796\\
3.8673	0.46098\\
3.8499	0.48518\\
3.8325	0.51064\\
3.815	0.53742\\
3.7976	0.5656\\
3.7802	0.59524\\
3.7628	0.62642\\
3.7454	0.65923\\
3.7279	0.69375\\
3.7105	0.73007\\
3.6931	0.76828\\
3.6757	0.8085\\
3.6583	0.85081\\
3.6408	0.8948\\
3.6234	0.93738\\
3.606	0.97185\\
}--cycle;
\addplot [color=mycolor1, mark repeat=9, mark phase=1, mark=x, mark options={solid, mycolor1}]
  table[row sep=crcr]{%
3.606	1.0025\\
3.6234	0.95866\\
3.6408	0.91639\\
3.6583	0.87564\\
3.6757	0.83639\\
3.6931	0.79862\\
3.7105	0.76228\\
3.7279	0.72736\\
3.7454	0.69381\\
3.7628	0.66161\\
3.7802	0.6307\\
3.7976	0.60106\\
3.815	0.57263\\
3.8325	0.54539\\
3.8499	0.5193\\
3.8673	0.4943\\
3.8847	0.47037\\
3.9021	0.44747\\
3.9196	0.42555\\
3.937	0.40457\\
3.9544	0.38451\\
3.9718	0.36532\\
3.9893	0.34697\\
4.0067	0.32943\\
4.0241	0.31266\\
4.0415	0.29664\\
4.0589	0.28132\\
4.0764	0.26669\\
4.0938	0.25271\\
4.1112	0.23936\\
4.1286	0.22661\\
4.146	0.21444\\
4.1635	0.20283\\
4.1809	0.19174\\
4.1983	0.18117\\
4.2157	0.17109\\
4.2331	0.16147\\
4.2506	0.15231\\
4.268	0.14359\\
4.2854	0.13528\\
4.3028	0.12737\\
4.3202	0.11984\\
4.3377	0.11269\\
4.3551	0.10589\\
4.3725	0.099431\\
4.3899	0.0933\\
4.4073	0.087483\\
4.4248	0.081968\\
4.4422	0.076743\\
4.4596	0.071795\\
4.477	0.067115\\
4.4944	0.062689\\
4.5119	0.058509\\
4.5293	0.054562\\
4.5467	0.05084\\
4.5641	0.047333\\
4.5815	0.04403\\
4.599	0.040923\\
4.6164	0.038003\\
4.6338	0.035261\\
4.6512	0.032688\\
4.6686	0.030276\\
4.6861	0.028018\\
4.7035	0.025905\\
4.7209	0.02393\\
4.7383	0.022086\\
4.7558	0.020366\\
4.7732	0.018764\\
4.7906	0.017272\\
4.808	0.015884\\
4.8254	0.014595\\
4.8429	0.013398\\
4.8603	0.012289\\
4.8777	0.011261\\
4.8951	0.01031\\
4.9125	0.0094313\\
4.93	0.0086196\\
4.9474	0.0078709\\
4.9648	0.0071808\\
4.9822	0.0065456\\
4.9996	0.0059613\\
5.0171	0.0054244\\
5.0345	0.0049317\\
5.0519	0.0044798\\
5.0693	0.0040659\\
5.0867	0.0036871\\
5.1042	0.0033408\\
5.1216	0.0030244\\
5.139	0.0027357\\
5.1564	0.0024725\\
5.1738	0.0022328\\
5.1913	0.0020147\\
5.2087	0.0018164\\
5.2261	0.0016363\\
5.2435	0.0014729\\
5.2609	0.0013247\\
5.2784	0.0011905\\
5.2958	0.0010691\\
5.3132	0.00095927\\
5.3306	0.00086007\\
};
\addlegendentry{\ApproxBLNFitLegend}


\addplot[area legend, draw=none, fill=mycolor1, fill opacity=0.15, postaction={pattern={Lines[angle=45,distance={5pt},line width={0.5pt}]}, pattern color=mycolor1!60}, forget plot]
table[row sep=crcr] {%
x	y\\
3.606	1.01\\
3.6234	0.96594\\
3.6408	0.92346\\
3.6583	0.88252\\
3.6757	0.84307\\
3.6931	0.80511\\
3.7105	0.76859\\
3.7279	0.7335\\
3.7454	0.69978\\
3.7628	0.66741\\
3.7802	0.63635\\
3.7976	0.60655\\
3.815	0.57798\\
3.8325	0.55059\\
3.8499	0.52435\\
3.8673	0.49921\\
3.8847	0.47513\\
3.9021	0.45208\\
3.9196	0.43001\\
3.937	0.40888\\
3.9544	0.38867\\
3.9718	0.36933\\
3.9893	0.35083\\
4.0067	0.33314\\
4.0241	0.31622\\
4.0415	0.30005\\
4.0589	0.28459\\
4.0764	0.26982\\
4.0938	0.25571\\
4.1112	0.24223\\
4.1286	0.22936\\
4.146	0.21707\\
4.1635	0.20534\\
4.1809	0.19416\\
4.1983	0.18349\\
4.2157	0.17332\\
4.2331	0.16363\\
4.2506	0.1544\\
4.268	0.1456\\
4.2854	0.13723\\
4.3028	0.12927\\
4.3202	0.12169\\
4.3377	0.11448\\
4.3551	0.10764\\
4.3725	0.10114\\
4.3899	0.094962\\
4.4073	0.089105\\
4.4248	0.08355\\
4.4422	0.078285\\
4.4596	0.073298\\
4.477	0.068578\\
4.4944	0.064114\\
4.5119	0.059894\\
4.5293	0.055908\\
4.5467	0.052146\\
4.5641	0.048598\\
4.5815	0.045255\\
4.599	0.042107\\
4.6164	0.039147\\
4.6338	0.036364\\
4.6512	0.033751\\
4.6686	0.031299\\
4.6861	0.029001\\
4.7035	0.026848\\
4.7209	0.024835\\
4.7383	0.022953\\
4.7558	0.021195\\
4.7732	0.019555\\
4.7906	0.018027\\
4.808	0.016604\\
4.8254	0.01528\\
4.8429	0.014049\\
4.8603	0.012907\\
4.8777	0.011848\\
4.8951	0.010866\\
4.9125	0.0099569\\
4.93	0.0091163\\
4.9474	0.0083396\\
4.9648	0.0076227\\
4.9822	0.0069616\\
4.9996	0.0063525\\
5.0171	0.0057919\\
5.0345	0.0052764\\
5.0519	0.0048028\\
5.0693	0.0043682\\
5.0867	0.0039696\\
5.1042	0.0036045\\
5.1216	0.0032703\\
5.139	0.0029648\\
5.1564	0.0026856\\
5.1738	0.0024308\\
5.1913	0.0021985\\
5.2087	0.0019868\\
5.2261	0.0017941\\
5.2435	0.0016188\\
5.2609	0.0014595\\
5.2784	0.0013149\\
5.2958	0.0011837\\
5.3132	0.0010649\\
5.3306	0.00095719\\
5.3306	0.00076295\\
5.3132	0.00085368\\
5.2958	0.00095441\\
5.2784	0.0010661\\
5.2609	0.0011899\\
5.2435	0.001327\\
5.2261	0.0014786\\
5.2087	0.0016461\\
5.1913	0.0018309\\
5.1738	0.0020348\\
5.1564	0.0022594\\
5.139	0.0025067\\
5.1216	0.0027785\\
5.1042	0.003077\\
5.0867	0.0034046\\
5.0693	0.0037637\\
5.0519	0.0041569\\
5.0345	0.004587\\
5.0171	0.005057\\
4.9996	0.0055701\\
4.9822	0.0061296\\
4.9648	0.006739\\
4.9474	0.0074022\\
4.93	0.008123\\
4.9125	0.0089058\\
4.8951	0.0097548\\
4.8777	0.010675\\
4.8603	0.01167\\
4.8429	0.012747\\
4.8254	0.01391\\
4.808	0.015164\\
4.7906	0.016516\\
4.7732	0.017972\\
4.7558	0.019538\\
4.7383	0.02122\\
4.7209	0.023026\\
4.7035	0.024961\\
4.6861	0.027035\\
4.6686	0.029253\\
4.6512	0.031625\\
4.6338	0.034157\\
4.6164	0.036859\\
4.599	0.039739\\
4.5815	0.042806\\
4.5641	0.046068\\
4.5467	0.049535\\
4.5293	0.053217\\
4.5119	0.057124\\
4.4944	0.061265\\
4.477	0.065651\\
4.4596	0.070292\\
4.4422	0.075201\\
4.4248	0.080386\\
4.4073	0.085861\\
4.3899	0.091637\\
4.3725	0.097727\\
4.3551	0.10414\\
4.3377	0.1109\\
4.3202	0.118\\
4.3028	0.12547\\
4.2854	0.13333\\
4.268	0.14157\\
4.2506	0.15023\\
4.2331	0.15932\\
4.2157	0.16885\\
4.1983	0.17885\\
4.1809	0.18933\\
4.1635	0.20031\\
4.146	0.21181\\
4.1286	0.22387\\
4.1112	0.23649\\
4.0938	0.24971\\
4.0764	0.26356\\
4.0589	0.27805\\
4.0415	0.29322\\
4.0241	0.3091\\
4.0067	0.32572\\
3.9893	0.34311\\
3.9718	0.36131\\
3.9544	0.38035\\
3.937	0.40026\\
3.9196	0.42108\\
3.9021	0.44286\\
3.8847	0.46562\\
3.8673	0.4894\\
3.8499	0.51425\\
3.8325	0.5402\\
3.815	0.56729\\
3.7976	0.59556\\
3.7802	0.62505\\
3.7628	0.6558\\
3.7454	0.68785\\
3.7279	0.72122\\
3.7105	0.75597\\
3.6931	0.79212\\
3.6757	0.82971\\
3.6583	0.86876\\
3.6408	0.90931\\
3.6234	0.95138\\
3.606	0.99499\\
}--cycle;
      % \input{../../../../.tex/20/KS/SizeDistributionFittingsfig24.tex}
      \legend{}
    \end{groupplot}

    \begin{groupplot}[
      group style={
        group name=Row4,
        plotGroup,
      },
      plotRow,
      row3Keys,
      plotLegendSW,
      %
      % xmin=3637.5,xmax=2.5003e+05,
      ymin=0.00083435,ymax=1.0175,
    ]
      \nextgroupplot[%
        at={($(Row3 c1r1.south west)-(0,\yPlotSep em)-(0,\plotHeight cm)$)},
        xmin=3.5608,xmax=5.3101,
        % xmin=3.4,xmax=5.4,
        % ymin=0.0041299,ymax=1.0005,
      ] % This file was created by matlab2tikz.
%
\definecolor{mycolor1}{rgb}{0.00000,0.44700,0.74100}%
%
\addplot[only marks, mark=*, mark size=1.1180pt, color=black, fill=black, opacity=0.60, draw=none, mark options={draw=none,line width=0pt}] table[row sep=crcr]{%
x	y\\
3.5623	0.99468\\
3.5723	0.98404\\
3.573	0.9734\\
3.5753	0.96277\\
3.5766	0.95213\\
3.5781	0.94149\\
3.5826	0.93085\\
3.5851	0.92021\\
3.5911	0.90957\\
3.5942	0.89894\\
3.596	0.8883\\
3.6016	0.87766\\
3.603	0.86702\\
3.6047	0.85638\\
3.6086	0.84574\\
3.6086	0.83511\\
3.6097	0.82447\\
3.6142	0.81383\\
3.6261	0.80319\\
3.6383	0.79255\\
3.6464	0.78191\\
3.6535	0.77128\\
3.657	0.76064\\
3.6658	0.75\\
3.6678	0.73936\\
3.6687	0.72872\\
3.6702	0.71809\\
3.6723	0.70745\\
3.6793	0.69681\\
3.6867	0.68617\\
3.6949	0.67553\\
3.7015	0.66489\\
3.7031	0.65426\\
3.719	0.64362\\
3.7213	0.63298\\
3.7265	0.62234\\
3.7307	0.6117\\
3.737	0.60106\\
3.7374	0.59043\\
3.7423	0.57979\\
3.7487	0.56915\\
3.7677	0.55851\\
3.7685	0.54787\\
3.7768	0.53723\\
3.8005	0.5266\\
3.8015	0.51596\\
3.8022	0.50532\\
3.8032	0.49468\\
3.8039	0.48404\\
3.8108	0.4734\\
3.8214	0.46277\\
3.8405	0.45213\\
3.8605	0.44149\\
3.863	0.43085\\
3.8645	0.42021\\
3.8662	0.40957\\
3.8806	0.39894\\
3.8878	0.3883\\
3.8964	0.37766\\
3.905	0.36702\\
3.9294	0.35638\\
3.9303	0.34574\\
3.954	0.33511\\
3.9604	0.32447\\
3.9669	0.31383\\
3.9747	0.30319\\
3.9822	0.29255\\
3.9929	0.28191\\
4.0051	0.27128\\
4.0144	0.26064\\
4.0161	0.25\\
4.0433	0.23936\\
4.0578	0.22872\\
4.073	0.21809\\
4.0857	0.20745\\
4.0904	0.19681\\
4.1168	0.18617\\
4.1265	0.17553\\
4.127	0.16489\\
4.143	0.15426\\
4.1756	0.14362\\
4.2156	0.13298\\
4.3116	0.12234\\
4.3276	0.1117\\
4.3662	0.10106\\
4.4791	0.090426\\
4.4912	0.079787\\
4.517	0.069149\\
4.5743	0.058511\\
4.5914	0.047872\\
4.6138	0.037234\\
4.7898	0.026596\\
5.0876	0.015957\\
5.3101	0.0053191\\
};
\addlegendentry{\RealEmpiricCCDFLegend}

\addplot [color=black, mark repeat=9, mark phase=1, mark=+, mark options={solid, black}]
  table[row sep=crcr]{%
3.5608	1\\
3.5785	0.94954\\
3.5961	0.90162\\
3.6138	0.85612\\
3.6315	0.81292\\
3.6492	0.7719\\
3.6668	0.73295\\
3.6845	0.69596\\
3.7022	0.66084\\
3.7198	0.6275\\
3.7375	0.59583\\
3.7552	0.56576\\
3.7728	0.53721\\
3.7905	0.5101\\
3.8082	0.48436\\
3.8259	0.45992\\
3.8435	0.43671\\
3.8612	0.41468\\
3.8789	0.39375\\
3.8965	0.37388\\
3.9142	0.35501\\
3.9319	0.3371\\
3.9495	0.32009\\
3.9672	0.30394\\
3.9849	0.2886\\
4.0026	0.27404\\
4.0202	0.26021\\
4.0379	0.24708\\
4.0556	0.23461\\
4.0732	0.22277\\
4.0909	0.21153\\
4.1086	0.20085\\
4.1262	0.19072\\
4.1439	0.18109\\
4.1616	0.17196\\
4.1793	0.16328\\
4.1969	0.15504\\
4.2146	0.14722\\
4.2323	0.13979\\
4.2499	0.13273\\
4.2676	0.12603\\
4.2853	0.11967\\
4.3029	0.11364\\
4.3206	0.1079\\
4.3383	0.10246\\
4.356	0.097286\\
4.3736	0.092377\\
4.3913	0.087716\\
4.409	0.083289\\
4.4266	0.079086\\
4.4443	0.075095\\
4.462	0.071306\\
4.4796	0.067708\\
4.4973	0.064291\\
4.515	0.061047\\
4.5327	0.057966\\
4.5503	0.055041\\
4.568	0.052264\\
4.5857	0.049626\\
4.6033	0.047122\\
4.621	0.044744\\
4.6387	0.042486\\
4.6563	0.040342\\
4.674	0.038307\\
4.6917	0.036373\\
4.7094	0.034538\\
4.727	0.032795\\
4.7447	0.03114\\
4.7624	0.029569\\
4.78	0.028077\\
4.7977	0.02666\\
4.8154	0.025315\\
4.833	0.024037\\
4.8507	0.022824\\
4.8684	0.021672\\
4.8861	0.020579\\
4.9037	0.01954\\
4.9214	0.018554\\
4.9391	0.017618\\
4.9567	0.016729\\
4.9744	0.015885\\
4.9921	0.015083\\
5.0097	0.014322\\
5.0274	0.013599\\
5.0451	0.012913\\
5.0628	0.012261\\
5.0804	0.011643\\
5.0981	0.011055\\
5.1158	0.010497\\
5.1334	0.0099676\\
5.1511	0.0094646\\
5.1688	0.008987\\
5.1864	0.0085335\\
5.2041	0.0081029\\
5.2218	0.007694\\
5.2395	0.0073058\\
5.2571	0.0069371\\
5.2748	0.006587\\
5.2925	0.0062546\\
5.3101	0.005939\\
};
\addlegendentry{\RealParetoFitLegend}

\addplot [color=black, mark repeat=9, mark phase=1, mark=x, mark options={solid, black}]
  table[row sep=crcr]{%
3.5608	1.0005\\
3.5785	0.95316\\
3.5961	0.90742\\
3.6138	0.86328\\
3.6315	0.82076\\
3.6492	0.77986\\
3.6668	0.74057\\
3.6845	0.70289\\
3.7022	0.66682\\
3.7198	0.63232\\
3.7375	0.59938\\
3.7552	0.56796\\
3.7728	0.53805\\
3.7905	0.5096\\
3.8082	0.48257\\
3.8259	0.45692\\
3.8435	0.43262\\
3.8612	0.4096\\
3.8789	0.38783\\
3.8965	0.36726\\
3.9142	0.34784\\
3.9319	0.32951\\
3.9495	0.31222\\
3.9672	0.29594\\
3.9849	0.28059\\
4.0026	0.26614\\
4.0202	0.25254\\
4.0379	0.23974\\
4.0556	0.22769\\
4.0732	0.21634\\
4.0909	0.20566\\
4.1086	0.19561\\
4.1262	0.18613\\
4.1439	0.1772\\
4.1616	0.16878\\
4.1793	0.16082\\
4.1969	0.15331\\
4.2146	0.14621\\
4.2323	0.1395\\
4.2499	0.13313\\
4.2676	0.1271\\
4.2853	0.12137\\
4.3029	0.11592\\
4.3206	0.11074\\
4.3383	0.10581\\
4.356	0.1011\\
4.3736	0.09661\\
4.3913	0.092318\\
4.409	0.088212\\
4.4266	0.084281\\
4.4443	0.080514\\
4.462	0.076901\\
4.4796	0.073434\\
4.4973	0.070104\\
4.515	0.066905\\
4.5327	0.06383\\
4.5503	0.060873\\
4.568	0.058029\\
4.5857	0.055293\\
4.6033	0.052661\\
4.621	0.050128\\
4.6387	0.047691\\
4.6563	0.045347\\
4.674	0.043093\\
4.6917	0.040926\\
4.7094	0.038844\\
4.727	0.036843\\
4.7447	0.034921\\
4.7624	0.033077\\
4.78	0.031309\\
4.7977	0.029613\\
4.8154	0.027989\\
4.833	0.026433\\
4.8507	0.024946\\
4.8684	0.023523\\
4.8861	0.022165\\
4.9037	0.020868\\
4.9214	0.019632\\
4.9391	0.018454\\
4.9567	0.017332\\
4.9744	0.016265\\
4.9921	0.015251\\
5.0097	0.014288\\
5.0274	0.013375\\
5.0451	0.01251\\
5.0628	0.01169\\
5.0804	0.010915\\
5.0981	0.010182\\
5.1158	0.0094906\\
5.1334	0.0088382\\
5.1511	0.0082235\\
5.1688	0.0076448\\
5.1864	0.0071005\\
5.2041	0.0065892\\
5.2218	0.0061093\\
5.2395	0.0056592\\
5.2571	0.0052377\\
5.2748	0.0048432\\
5.2925	0.0044744\\
5.3101	0.0041299\\
};
\addlegendentry{\RealBLNFitLegend}

      % \input{../../../../.tex/20/KS/SizeDistributionFittingsfig31.tex}
      \legend{}
  
      \nextgroupplot[%
        xmin=3.598,xmax=5.4115,
        % xmin=3.6,xmax=5.4,
        % ymin=0.00083435,ymax=1.0175,
      ] % This file was created by matlab2tikz.
%
\definecolor{mycolor1}{rgb}{0.00000,0.44706,0.69804}%
\definecolor{mycolor2}{rgb}{0.00000,0.44700,0.74100}%
%
\addplot[only marks, mark=*, mark size=1.1180pt, color=mycolor1, fill=mycolor1, opacity=0.60, draw=none, mark options={draw=none,line width=0pt}] table[row sep=crcr]{%
x	y\\
3.6054	0.99468\\
3.6098	0.98404\\
3.6142	0.9734\\
3.6177	0.96277\\
3.6219	0.95213\\
3.6266	0.94149\\
3.6304	0.93085\\
3.6361	0.92021\\
3.64	0.90957\\
3.6454	0.89894\\
3.6509	0.8883\\
3.6553	0.87766\\
3.6597	0.86702\\
3.6642	0.85638\\
3.669	0.84574\\
3.6744	0.83511\\
3.6791	0.82447\\
3.6837	0.81383\\
3.6885	0.80319\\
3.6932	0.79255\\
3.699	0.78191\\
3.7045	0.77128\\
3.7096	0.76064\\
3.7159	0.75\\
3.7217	0.73936\\
3.7271	0.72872\\
3.7326	0.71809\\
3.7382	0.70745\\
3.7444	0.69681\\
3.7511	0.68617\\
3.7579	0.67553\\
3.7642	0.66489\\
3.771	0.65426\\
3.7772	0.64362\\
3.7827	0.63298\\
3.7895	0.62234\\
3.7969	0.6117\\
3.8022	0.60106\\
3.8087	0.59043\\
3.8158	0.57979\\
3.8225	0.56915\\
3.83	0.55851\\
3.8371	0.54787\\
3.8436	0.53723\\
3.8503	0.5266\\
3.858	0.51596\\
3.8673	0.50532\\
3.8759	0.49468\\
3.8835	0.48404\\
3.8911	0.4734\\
3.9004	0.46277\\
3.9082	0.45213\\
3.9162	0.44149\\
3.9249	0.43085\\
3.9328	0.42021\\
3.9406	0.40957\\
3.949	0.39894\\
3.9577	0.3883\\
3.9678	0.37766\\
3.98	0.36702\\
3.9908	0.35638\\
4.0008	0.34574\\
4.0113	0.33511\\
4.0202	0.32447\\
4.0318	0.31383\\
4.0441	0.30319\\
4.0554	0.29255\\
4.0676	0.28191\\
4.0793	0.27128\\
4.0914	0.26064\\
4.1033	0.25\\
4.1167	0.23936\\
4.1321	0.22872\\
4.1494	0.21809\\
4.1669	0.20745\\
4.1856	0.19681\\
4.2028	0.18617\\
4.2197	0.17553\\
4.2417	0.16489\\
4.2609	0.15426\\
4.2818	0.14362\\
4.3051	0.13298\\
4.33	0.12234\\
4.3521	0.1117\\
4.3782	0.10106\\
4.4054	0.090426\\
4.4416	0.079787\\
4.4781	0.069149\\
4.5216	0.058511\\
4.574	0.047872\\
4.636	0.037234\\
4.7242	0.026596\\
4.8252	0.015957\\
5.1746	0.0053191\\
};
\addlegendentry{\ExactEmpiricCCDFLegend}

\addplot [color=mycolor1, only marks, every error bar/.append style={opacity=0.30}, mark=*, mark size=0pt, draw=none, forget plot]
 plot [error bars/.cd, x dir=both, x explicit, error bar style={line width=1pt, color=mycolor1}, error mark options={mark=none,mark size=0pt}]
 table[row sep=crcr, x error plus index=2, x error minus index=3]{%
3.6054	0.99468	0.004187	0.004187\\
3.6098	0.98404	0.003977	0.003977\\
3.6142	0.9734	0.0040165	0.0040165\\
3.6177	0.96277	0.0040403	0.0040403\\
3.6219	0.95213	0.0039714	0.0039714\\
3.6266	0.94149	0.00393	0.00393\\
3.6304	0.93085	0.00398	0.00398\\
3.6361	0.92021	0.0040894	0.0040894\\
3.64	0.90957	0.0042313	0.0042313\\
3.6454	0.89894	0.004231	0.004231\\
3.6509	0.8883	0.0044287	0.0044287\\
3.6553	0.87766	0.0044395	0.0044395\\
3.6597	0.86702	0.0044097	0.0044097\\
3.6642	0.85638	0.0045406	0.0045406\\
3.669	0.84574	0.0045804	0.0045804\\
3.6744	0.83511	0.0045103	0.0045103\\
3.6791	0.82447	0.0045155	0.0045155\\
3.6837	0.81383	0.0044785	0.0044785\\
3.6885	0.80319	0.0044161	0.0044161\\
3.6932	0.79255	0.0044769	0.0044769\\
3.699	0.78191	0.004542	0.004542\\
3.7045	0.77128	0.0045765	0.0045765\\
3.7096	0.76064	0.0043954	0.0043954\\
3.7159	0.75	0.004498	0.004498\\
3.7217	0.73936	0.0045568	0.0045568\\
3.7271	0.72872	0.004502	0.004502\\
3.7326	0.71809	0.0045084	0.0045084\\
3.7382	0.70745	0.0046824	0.0046824\\
3.7444	0.69681	0.0047522	0.0047522\\
3.7511	0.68617	0.0049632	0.0049632\\
3.7579	0.67553	0.0048624	0.0048624\\
3.7642	0.66489	0.0048951	0.0048951\\
3.771	0.65426	0.0048301	0.0048301\\
3.7772	0.64362	0.0047613	0.0047613\\
3.7827	0.63298	0.0047343	0.0047343\\
3.7895	0.62234	0.004828	0.004828\\
3.7969	0.6117	0.0050184	0.0050184\\
3.8022	0.60106	0.0051359	0.0051359\\
3.8087	0.59043	0.0052829	0.0052829\\
3.8158	0.57979	0.0054389	0.0054389\\
3.8225	0.56915	0.0055371	0.0055371\\
3.83	0.55851	0.005547	0.005547\\
3.8371	0.54787	0.0055845	0.0055845\\
3.8436	0.53723	0.0056346	0.0056346\\
3.8503	0.5266	0.0055049	0.0055049\\
3.858	0.51596	0.0052181	0.0052181\\
3.8673	0.50532	0.0052516	0.0052516\\
3.8759	0.49468	0.0053721	0.0053721\\
3.8835	0.48404	0.0051237	0.0051237\\
3.8911	0.4734	0.0051448	0.0051448\\
3.9004	0.46277	0.0056128	0.0056128\\
3.9082	0.45213	0.0057038	0.0057038\\
3.9162	0.44149	0.0058207	0.0058207\\
3.9249	0.43085	0.0060439	0.0060439\\
3.9328	0.42021	0.0061768	0.0061768\\
3.9406	0.40957	0.006023	0.006023\\
3.949	0.39894	0.0062245	0.0062245\\
3.9577	0.3883	0.0064029	0.0064029\\
3.9678	0.37766	0.006833	0.006833\\
3.98	0.36702	0.0067149	0.0067149\\
3.9908	0.35638	0.0068135	0.0068135\\
4.0008	0.34574	0.0068598	0.0068598\\
4.0113	0.33511	0.0068246	0.0068246\\
4.0202	0.32447	0.0069444	0.0069444\\
4.0318	0.31383	0.0067454	0.0067454\\
4.0441	0.30319	0.0064174	0.0064174\\
4.0554	0.29255	0.0065338	0.0065338\\
4.0676	0.28191	0.0060628	0.0060628\\
4.0793	0.27128	0.0060461	0.0060461\\
4.0914	0.26064	0.0061657	0.0061657\\
4.1033	0.25	0.0062693	0.0062693\\
4.1167	0.23936	0.0063805	0.0063805\\
4.1321	0.22872	0.0066514	0.0066514\\
4.1494	0.21809	0.0069426	0.0069426\\
4.1669	0.20745	0.0072838	0.0072838\\
4.1856	0.19681	0.0076203	0.0076203\\
4.2028	0.18617	0.0080221	0.0080221\\
4.2197	0.17553	0.0079141	0.0079141\\
4.2417	0.16489	0.0083169	0.0083169\\
4.2609	0.15426	0.0082648	0.0082648\\
4.2818	0.14362	0.0083327	0.0083327\\
4.3051	0.13298	0.0078595	0.0078595\\
4.33	0.12234	0.0080214	0.0080214\\
4.3521	0.1117	0.0077859	0.0077859\\
4.3782	0.10106	0.0082103	0.0082103\\
4.4054	0.090426	0.0084553	0.0084553\\
4.4416	0.079787	0.009553	0.009553\\
4.4781	0.069149	0.0095889	0.0095889\\
4.5216	0.058511	0.0095988	0.0095988\\
4.574	0.047872	0.011021	0.011021\\
4.636	0.037234	0.013906	0.013906\\
4.7242	0.026596	0.014438	0.014438\\
4.8252	0.015957	0.016348	0.016348\\
5.1746	0.0053191	0.014408	0.014408\\
};
\addplot [color=mycolor1, mark repeat=9, mark phase=1, mark=+, mark options={solid, mycolor1}]
  table[row sep=crcr]{%
3.6037	0.97822\\
3.6214	0.9464\\
3.6392	0.90579\\
3.6569	0.86178\\
3.6747	0.8189\\
3.6924	0.77816\\
3.7102	0.73946\\
3.7279	0.70268\\
3.7457	0.66775\\
3.7634	0.63455\\
3.7812	0.60301\\
3.7989	0.57304\\
3.8167	0.54457\\
3.8344	0.51752\\
3.8522	0.49181\\
3.8699	0.46739\\
3.8877	0.44418\\
3.9054	0.42213\\
3.9232	0.40117\\
3.9409	0.38126\\
3.9586	0.36235\\
3.9764	0.34437\\
3.9941	0.32729\\
4.0119	0.31106\\
4.0296	0.29563\\
4.0474	0.28097\\
4.0651	0.26705\\
4.0829	0.25381\\
4.1006	0.24123\\
4.1184	0.22928\\
4.1361	0.21792\\
4.1539	0.20713\\
4.1716	0.19687\\
4.1894	0.18713\\
4.2071	0.17786\\
4.2249	0.16906\\
4.2426	0.16069\\
4.2604	0.15274\\
4.2781	0.14519\\
4.2959	0.13801\\
4.3136	0.13118\\
4.3313	0.12469\\
4.3491	0.11853\\
4.3668	0.11267\\
4.3846	0.1071\\
4.4023	0.10181\\
4.4201	0.096781\\
4.4378	0.092001\\
4.4556	0.087457\\
4.4733	0.083139\\
4.4911	0.079034\\
4.5088	0.075133\\
4.5266	0.071425\\
4.5443	0.067901\\
4.5621	0.064551\\
4.5798	0.061366\\
4.5976	0.05834\\
4.6153	0.055463\\
4.6331	0.052728\\
4.6508	0.050129\\
4.6686	0.047658\\
4.6863	0.04531\\
4.704	0.043078\\
4.7218	0.040955\\
4.7395	0.038938\\
4.7573	0.037021\\
4.775	0.035198\\
4.7928	0.033465\\
4.8105	0.031818\\
4.8283	0.030252\\
4.846	0.028764\\
4.8638	0.027349\\
4.8815	0.026003\\
4.8993	0.024724\\
4.917	0.023509\\
4.9348	0.022353\\
4.9525	0.021254\\
4.9703	0.02021\\
4.988	0.019217\\
5.0058	0.018272\\
5.0235	0.017375\\
5.0412	0.016522\\
5.059	0.01571\\
5.0767	0.014939\\
5.0945	0.014206\\
5.1122	0.013508\\
5.13	0.012845\\
5.1477	0.012215\\
5.1655	0.011616\\
5.1832	0.011046\\
5.201	0.010505\\
5.2187	0.0099895\\
5.2365	0.0094998\\
5.2542	0.0090341\\
5.272	0.0085914\\
5.2897	0.0081704\\
5.3075	0.0077702\\
5.3252	0.0073896\\
5.343	0.0070277\\
5.3607	0.0066835\\
};
\addlegendentry{\ExactParetoFitLegend}


\addplot[area legend, draw=none, fill=mycolor1, fill opacity=0.15, postaction={pattern={Lines[angle=90,distance={5pt},line width={0.5pt}]}, pattern color=mycolor1!60}, forget plot]
table[row sep=crcr] {%
x	y\\
3.6037	0.98434\\
3.6214	0.95496\\
3.6392	0.91543\\
3.6569	0.87097\\
3.6747	0.82724\\
3.6924	0.78573\\
3.7102	0.74632\\
3.7279	0.7089\\
3.7457	0.67338\\
3.7634	0.63966\\
3.7812	0.60765\\
3.7989	0.57726\\
3.8167	0.54842\\
3.8344	0.52104\\
3.8522	0.49506\\
3.8699	0.47039\\
3.8877	0.44699\\
3.9054	0.42477\\
3.9232	0.40369\\
3.9409	0.38367\\
3.9586	0.36467\\
3.9764	0.34664\\
3.9941	0.32951\\
4.0119	0.31325\\
4.0296	0.2978\\
4.0474	0.28313\\
4.0651	0.26919\\
4.0829	0.25595\\
4.1006	0.24336\\
4.1184	0.2314\\
4.1361	0.22003\\
4.1539	0.20923\\
4.1716	0.19896\\
4.1894	0.1892\\
4.2071	0.17992\\
4.2249	0.1711\\
4.2426	0.16272\\
4.2604	0.15475\\
4.2781	0.14717\\
4.2959	0.13996\\
4.3136	0.13311\\
4.3313	0.12659\\
4.3491	0.1204\\
4.3668	0.11451\\
4.3846	0.10891\\
4.4023	0.10358\\
4.4201	0.098521\\
4.4378	0.093706\\
4.4556	0.089127\\
4.4733	0.084772\\
4.4911	0.080631\\
4.5088	0.076694\\
4.5266	0.072949\\
4.5443	0.069388\\
4.5621	0.066001\\
4.5798	0.06278\\
4.5976	0.059717\\
4.6153	0.056804\\
4.6331	0.054033\\
4.6508	0.051399\\
4.6686	0.048893\\
4.6863	0.046509\\
4.704	0.044243\\
4.7218	0.042087\\
4.7395	0.040036\\
4.7573	0.038086\\
4.775	0.036231\\
4.7928	0.034467\\
4.8105	0.032789\\
4.8283	0.031193\\
4.846	0.029675\\
4.8638	0.028231\\
4.8815	0.026857\\
4.8993	0.025551\\
4.917	0.024308\\
4.9348	0.023126\\
4.9525	0.022002\\
4.9703	0.020932\\
4.988	0.019915\\
5.0058	0.018947\\
5.0235	0.018027\\
5.0412	0.017151\\
5.059	0.016318\\
5.0767	0.015526\\
5.0945	0.014772\\
5.1122	0.014055\\
5.13	0.013373\\
5.1477	0.012724\\
5.1655	0.012107\\
5.1832	0.011519\\
5.201	0.010961\\
5.2187	0.010429\\
5.2365	0.0099235\\
5.2542	0.0094424\\
5.272	0.0089848\\
5.2897	0.0085494\\
5.3075	0.0081352\\
5.3252	0.0077411\\
5.343	0.0073662\\
5.3607	0.0070095\\
5.3607	0.0063576\\
5.343	0.0066891\\
5.3252	0.007038\\
5.3075	0.0074051\\
5.2897	0.0077915\\
5.272	0.008198\\
5.2542	0.0086258\\
5.2365	0.0090761\\
5.2187	0.0095499\\
5.201	0.010048\\
5.1832	0.010573\\
5.1655	0.011125\\
5.1477	0.011707\\
5.13	0.012318\\
5.1122	0.012962\\
5.0945	0.013639\\
5.0767	0.014352\\
5.059	0.015102\\
5.0412	0.015892\\
5.0235	0.016723\\
5.0058	0.017598\\
4.988	0.018518\\
4.9703	0.019487\\
4.9525	0.020507\\
4.9348	0.02158\\
4.917	0.022709\\
4.8993	0.023898\\
4.8815	0.025149\\
4.8638	0.026466\\
4.846	0.027852\\
4.8283	0.029311\\
4.8105	0.030847\\
4.7928	0.032463\\
4.775	0.034165\\
4.7573	0.035955\\
4.7395	0.03784\\
4.7218	0.039824\\
4.704	0.041912\\
4.6863	0.044111\\
4.6686	0.046424\\
4.6508	0.04886\\
4.6331	0.051423\\
4.6153	0.054122\\
4.5976	0.056963\\
4.5798	0.059953\\
4.5621	0.0631\\
4.5443	0.066413\\
4.5266	0.069901\\
4.5088	0.073572\\
4.4911	0.077437\\
4.4733	0.081505\\
4.4556	0.085788\\
4.4378	0.090296\\
4.4201	0.095042\\
4.4023	0.10004\\
4.3846	0.1053\\
4.3668	0.11083\\
4.3491	0.11666\\
4.3313	0.12279\\
4.3136	0.12925\\
4.2959	0.13605\\
4.2781	0.14321\\
4.2604	0.15074\\
4.2426	0.15867\\
4.2249	0.16702\\
4.2071	0.1758\\
4.1894	0.18505\\
4.1716	0.19478\\
4.1539	0.20503\\
4.1361	0.21581\\
4.1184	0.22716\\
4.1006	0.23911\\
4.0829	0.25168\\
4.0651	0.2649\\
4.0474	0.27882\\
4.0296	0.29346\\
4.0119	0.30886\\
3.9941	0.32506\\
3.9764	0.3421\\
3.9586	0.36002\\
3.9409	0.37886\\
3.9232	0.39866\\
3.9054	0.41948\\
3.8877	0.44137\\
3.8699	0.46438\\
3.8522	0.48857\\
3.8344	0.51399\\
3.8167	0.54072\\
3.7989	0.56883\\
3.7812	0.59837\\
3.7634	0.62944\\
3.7457	0.66211\\
3.7279	0.69647\\
3.7102	0.7326\\
3.6924	0.77059\\
3.6747	0.81055\\
3.6569	0.85258\\
3.6392	0.89615\\
3.6214	0.93785\\
3.6037	0.97209\\
}--cycle;
\addplot [color=mycolor1, mark repeat=9, mark phase=1, mark=x, mark options={solid, mycolor1}]
  table[row sep=crcr]{%
3.6037	1.0099\\
3.6214	0.96715\\
3.6392	0.92576\\
3.6569	0.88568\\
3.6747	0.84691\\
3.6924	0.80944\\
3.7102	0.77324\\
3.7279	0.73831\\
3.7457	0.70463\\
3.7634	0.67218\\
3.7812	0.64092\\
3.7989	0.61085\\
3.8167	0.58194\\
3.8344	0.55416\\
3.8522	0.52748\\
3.8699	0.50188\\
3.8877	0.47732\\
3.9054	0.45379\\
3.9232	0.43124\\
3.9409	0.40965\\
3.9586	0.38899\\
3.9764	0.36922\\
3.9941	0.35033\\
4.0119	0.33227\\
4.0296	0.31502\\
4.0474	0.29854\\
4.0651	0.28282\\
4.0829	0.26781\\
4.1006	0.2535\\
4.1184	0.23986\\
4.1361	0.22685\\
4.1539	0.21446\\
4.1716	0.20266\\
4.1894	0.19142\\
4.2071	0.18072\\
4.2249	0.17055\\
4.2426	0.16086\\
4.2604	0.15166\\
4.2781	0.1429\\
4.2959	0.13459\\
4.3136	0.12669\\
4.3313	0.11918\\
4.3491	0.11206\\
4.3668	0.10531\\
4.3846	0.098898\\
4.4023	0.092823\\
4.4201	0.087067\\
4.4378	0.081615\\
4.4556	0.076454\\
4.4733	0.07157\\
4.4911	0.066952\\
4.5088	0.062588\\
4.5266	0.058466\\
4.5443	0.054575\\
4.5621	0.050904\\
4.5798	0.047444\\
4.5976	0.044185\\
4.6153	0.041116\\
4.6331	0.03823\\
4.6508	0.035518\\
4.6686	0.032971\\
4.6863	0.030581\\
4.704	0.028341\\
4.7218	0.026242\\
4.7395	0.024278\\
4.7573	0.022442\\
4.775	0.020727\\
4.7928	0.019126\\
4.8105	0.017634\\
4.8283	0.016244\\
4.846	0.014951\\
4.8638	0.013749\\
4.8815	0.012633\\
4.8993	0.011598\\
4.917	0.010638\\
4.9348	0.0097497\\
4.9525	0.0089281\\
4.9703	0.008169\\
4.988	0.0074684\\
5.0058	0.0068223\\
5.0235	0.0062272\\
5.0412	0.0056796\\
5.059	0.005176\\
5.0767	0.0047135\\
5.0945	0.0042892\\
5.1122	0.0039001\\
5.13	0.0035437\\
5.1477	0.0032176\\
5.1655	0.0029194\\
5.1832	0.002647\\
5.201	0.0023984\\
5.2187	0.0021717\\
5.2365	0.0019651\\
5.2542	0.001777\\
5.272	0.0016059\\
5.2897	0.0014503\\
5.3075	0.001309\\
5.3252	0.0011807\\
5.343	0.0010644\\
5.3607	0.00095897\\
};
\addlegendentry{\ExactBLNFitLegend}


\addplot[area legend, draw=none, fill=mycolor1, fill opacity=0.15, postaction={pattern={Lines[angle=45,distance={5pt},line width={0.5pt}]}, pattern color=mycolor1!60}, forget plot]
table[row sep=crcr] {%
x	y\\
3.6037	1.0175\\
3.6214	0.97459\\
3.6392	0.93299\\
3.6569	0.89272\\
3.6747	0.85377\\
3.6924	0.81612\\
3.7102	0.77976\\
3.7279	0.74466\\
3.7457	0.71082\\
3.7634	0.6782\\
3.7812	0.64679\\
3.7989	0.61656\\
3.8167	0.58749\\
3.8344	0.55955\\
3.8522	0.53271\\
3.8699	0.50695\\
3.8877	0.48224\\
3.9054	0.45855\\
3.9232	0.43584\\
3.9409	0.41409\\
3.9586	0.39327\\
3.9764	0.37335\\
3.9941	0.35429\\
4.0119	0.33608\\
4.0296	0.31868\\
4.0474	0.30206\\
4.0651	0.2862\\
4.0829	0.27106\\
4.1006	0.25662\\
4.1184	0.24285\\
4.1361	0.22973\\
4.1539	0.21723\\
4.1716	0.20532\\
4.1894	0.19399\\
4.2071	0.1832\\
4.2249	0.17294\\
4.2426	0.16318\\
4.2604	0.1539\\
4.2781	0.14509\\
4.2959	0.13671\\
4.3136	0.12876\\
4.3313	0.1212\\
4.3491	0.11403\\
4.3668	0.10723\\
4.3846	0.10077\\
4.4023	0.094653\\
4.4201	0.088852\\
4.4378	0.083356\\
4.4556	0.07815\\
4.4733	0.073222\\
4.4911	0.068559\\
4.5088	0.064148\\
4.5266	0.05998\\
4.5443	0.056042\\
4.5621	0.052324\\
4.5798	0.048817\\
4.5976	0.04551\\
4.6153	0.042394\\
4.6331	0.039461\\
4.6508	0.036702\\
4.6686	0.034109\\
4.6863	0.031673\\
4.704	0.029388\\
4.7218	0.027246\\
4.7395	0.025239\\
4.7573	0.023362\\
4.775	0.021606\\
4.7928	0.019967\\
4.8105	0.018437\\
4.8283	0.01701\\
4.846	0.015682\\
4.8638	0.014446\\
4.8815	0.013297\\
4.8993	0.01223\\
4.917	0.011239\\
4.9348	0.010321\\
4.9525	0.0094712\\
4.9703	0.0086845\\
4.988	0.0079573\\
5.0058	0.0072856\\
5.0235	0.0066657\\
5.0412	0.0060942\\
5.059	0.0055677\\
5.0767	0.0050831\\
5.0945	0.0046375\\
5.1122	0.004228\\
5.13	0.003852\\
5.1477	0.0035072\\
5.1655	0.0031911\\
5.1832	0.0029016\\
5.201	0.0026366\\
5.2187	0.0023943\\
5.2365	0.002173\\
5.2542	0.0019709\\
5.272	0.0017865\\
5.2897	0.0016183\\
5.3075	0.0014652\\
5.3252	0.0013258\\
5.343	0.0011989\\
5.3607	0.0010836\\
5.3607	0.00083435\\
5.343	0.0009299\\
5.3252	0.0010357\\
5.3075	0.0011528\\
5.2897	0.0012822\\
5.272	0.0014252\\
5.2542	0.0015831\\
5.2365	0.0017572\\
5.2187	0.001949\\
5.201	0.0021602\\
5.1832	0.0023925\\
5.1655	0.0026477\\
5.1477	0.002928\\
5.13	0.0032354\\
5.1122	0.0035722\\
5.0945	0.0039408\\
5.0767	0.004344\\
5.059	0.0047843\\
5.0412	0.0052649\\
5.0235	0.0057887\\
5.0058	0.0063591\\
4.988	0.0069795\\
4.9703	0.0076535\\
4.9525	0.008385\\
4.9348	0.009178\\
4.917	0.010037\\
4.8993	0.010966\\
4.8815	0.011969\\
4.8638	0.013052\\
4.846	0.01422\\
4.8283	0.015478\\
4.8105	0.016831\\
4.7928	0.018286\\
4.775	0.019847\\
4.7573	0.021522\\
4.7395	0.023317\\
4.7218	0.025238\\
4.704	0.027293\\
4.6863	0.029489\\
4.6686	0.031833\\
4.6508	0.034334\\
4.6331	0.037\\
4.6153	0.039838\\
4.5976	0.042859\\
4.5798	0.046071\\
4.5621	0.049484\\
4.5443	0.053107\\
4.5266	0.056952\\
4.5088	0.061028\\
4.4911	0.065346\\
4.4733	0.069919\\
4.4556	0.074757\\
4.4378	0.079874\\
4.4201	0.085282\\
4.4023	0.090994\\
4.3846	0.097024\\
4.3668	0.10339\\
4.3491	0.1101\\
4.3313	0.11717\\
4.3136	0.12462\\
4.2959	0.13246\\
4.2781	0.14072\\
4.2604	0.14941\\
4.2426	0.15855\\
4.2249	0.16815\\
4.2071	0.17825\\
4.1894	0.18886\\
4.1716	0.2\\
4.1539	0.2117\\
4.1361	0.22398\\
4.1184	0.23687\\
4.1006	0.25039\\
4.0829	0.26457\\
4.0651	0.27944\\
4.0474	0.29502\\
4.0296	0.31135\\
4.0119	0.32845\\
3.9941	0.34636\\
3.9764	0.3651\\
3.9586	0.38471\\
3.9409	0.40521\\
3.9232	0.42664\\
3.9054	0.44903\\
3.8877	0.47241\\
3.8699	0.4968\\
3.8522	0.52225\\
3.8344	0.54877\\
3.8167	0.57639\\
3.7989	0.60515\\
3.7812	0.63506\\
3.7634	0.66615\\
3.7457	0.69845\\
3.7279	0.73197\\
3.7102	0.76673\\
3.6924	0.80275\\
3.6747	0.84005\\
3.6569	0.87863\\
3.6392	0.91852\\
3.6214	0.95972\\
3.6037	1.0022\\
}--cycle;
      % \input{../../../../.tex/20/KS/SizeDistributionFittingsfig32.tex}
      \legend{}
  
      \nextgroupplot[%
        xmin=3.6008,xmax=5.4115,
        % xmin=3.4,xmax=5.4,
        % ymin=0.00096016,ymax=1.0131,
      ] % This file was created by matlab2tikz.
%
\definecolor{mycolor1}{rgb}{0.83529,0.36863,0.00000}%
\definecolor{mycolor2}{rgb}{0.00000,0.44700,0.74100}%
%
\addplot[only marks, mark=*, mark size=1.1180pt, color=mycolor1, fill=mycolor1, opacity=0.60, draw=none, mark options={draw=none,line width=0pt}] table[row sep=crcr]{%
x	y\\
3.6043	0.99468\\
3.6085	0.98404\\
3.6129	0.9734\\
3.6166	0.96277\\
3.6217	0.95213\\
3.6257	0.94149\\
3.6297	0.93085\\
3.6336	0.92021\\
3.6386	0.90957\\
3.6436	0.89894\\
3.6485	0.8883\\
3.6534	0.87766\\
3.6578	0.86702\\
3.6617	0.85638\\
3.6672	0.84574\\
3.6724	0.83511\\
3.6775	0.82447\\
3.683	0.81383\\
3.6877	0.80319\\
3.6931	0.79255\\
3.6981	0.78191\\
3.7035	0.77128\\
3.7087	0.76064\\
3.7142	0.75\\
3.7195	0.73936\\
3.7257	0.72872\\
3.7321	0.71809\\
3.7384	0.70745\\
3.7448	0.69681\\
3.7505	0.68617\\
3.7557	0.67553\\
3.7608	0.66489\\
3.7667	0.65426\\
3.7732	0.64362\\
3.7794	0.63298\\
3.7861	0.62234\\
3.7926	0.6117\\
3.7986	0.60106\\
3.8051	0.59043\\
3.8111	0.57979\\
3.8186	0.56915\\
3.8258	0.55851\\
3.8344	0.54787\\
3.8405	0.53723\\
3.848	0.5266\\
3.855	0.51596\\
3.8618	0.50532\\
3.8693	0.49468\\
3.8771	0.48404\\
3.8852	0.4734\\
3.8931	0.46277\\
3.9016	0.45213\\
3.9092	0.44149\\
3.9181	0.43085\\
3.9273	0.42021\\
3.9357	0.40957\\
3.9432	0.39894\\
3.9533	0.3883\\
3.9637	0.37766\\
3.9738	0.36702\\
3.9848	0.35638\\
3.9951	0.34574\\
4.0063	0.33511\\
4.0195	0.32447\\
4.0316	0.31383\\
4.0451	0.30319\\
4.0541	0.29255\\
4.0665	0.28191\\
4.0793	0.27128\\
4.0929	0.26064\\
4.1075	0.25\\
4.1207	0.23936\\
4.1351	0.22872\\
4.1514	0.21809\\
4.1699	0.20745\\
4.1822	0.19681\\
4.1983	0.18617\\
4.2172	0.17553\\
4.2352	0.16489\\
4.2564	0.15426\\
4.2767	0.14362\\
4.2966	0.13298\\
4.3207	0.12234\\
4.3486	0.1117\\
4.3743	0.10106\\
4.4051	0.090426\\
4.439	0.079787\\
4.4763	0.069149\\
4.5194	0.058511\\
4.5656	0.047872\\
4.635	0.037234\\
4.7318	0.026596\\
4.8427	0.015957\\
5.1833	0.0053191\\
};
\addlegendentry{\ApproxEmpiricCCDFLegend}

\addplot [color=mycolor1, only marks, every error bar/.append style={opacity=0.30}, mark=*, mark size=0pt, draw=none, forget plot]
 plot [error bars/.cd, x dir=both, x explicit, error bar style={line width=1pt, color=mycolor1}, error mark options={mark=none,mark size=0pt}]
 table[row sep=crcr, x error plus index=2, x error minus index=3]{%
3.6043	0.99468	0.0045572	0.0045572\\
3.6085	0.98404	0.0043712	0.0043712\\
3.6129	0.9734	0.0045168	0.0045168\\
3.6166	0.96277	0.0044731	0.0044731\\
3.6217	0.95213	0.0044265	0.0044265\\
3.6257	0.94149	0.0044362	0.0044362\\
3.6297	0.93085	0.0044388	0.0044388\\
3.6336	0.92021	0.0043707	0.0043707\\
3.6386	0.90957	0.0046291	0.0046291\\
3.6436	0.89894	0.0046805	0.0046805\\
3.6485	0.8883	0.0046062	0.0046062\\
3.6534	0.87766	0.0047192	0.0047192\\
3.6578	0.86702	0.0047457	0.0047457\\
3.6617	0.85638	0.0047213	0.0047213\\
3.6672	0.84574	0.0048522	0.0048522\\
3.6724	0.83511	0.0047178	0.0047178\\
3.6775	0.82447	0.0045714	0.0045714\\
3.683	0.81383	0.0044815	0.0044815\\
3.6877	0.80319	0.0045796	0.0045796\\
3.6931	0.79255	0.0046427	0.0046427\\
3.6981	0.78191	0.0047601	0.0047601\\
3.7035	0.77128	0.0049339	0.0049339\\
3.7087	0.76064	0.0048633	0.0048633\\
3.7142	0.75	0.0048943	0.0048943\\
3.7195	0.73936	0.0047886	0.0047886\\
3.7257	0.72872	0.0047314	0.0047314\\
3.7321	0.71809	0.0048394	0.0048394\\
3.7384	0.70745	0.0048539	0.0048539\\
3.7448	0.69681	0.005076	0.005076\\
3.7505	0.68617	0.0050712	0.0050712\\
3.7557	0.67553	0.0050637	0.0050637\\
3.7608	0.66489	0.0051589	0.0051589\\
3.7667	0.65426	0.0051728	0.0051728\\
3.7732	0.64362	0.0050904	0.0050904\\
3.7794	0.63298	0.0050729	0.0050729\\
3.7861	0.62234	0.0049622	0.0049622\\
3.7926	0.6117	0.0050011	0.0050011\\
3.7986	0.60106	0.0051642	0.0051642\\
3.8051	0.59043	0.0050796	0.0050796\\
3.8111	0.57979	0.0051542	0.0051542\\
3.8186	0.56915	0.0054195	0.0054195\\
3.8258	0.55851	0.0056046	0.0056046\\
3.8344	0.54787	0.0058109	0.0058109\\
3.8405	0.53723	0.0059958	0.0059958\\
3.848	0.5266	0.0057197	0.0057197\\
3.855	0.51596	0.0057667	0.0057667\\
3.8618	0.50532	0.0059162	0.0059162\\
3.8693	0.49468	0.005978	0.005978\\
3.8771	0.48404	0.0059461	0.0059461\\
3.8852	0.4734	0.005987	0.005987\\
3.8931	0.46277	0.0058947	0.0058947\\
3.9016	0.45213	0.0059886	0.0059886\\
3.9092	0.44149	0.0058124	0.0058124\\
3.9181	0.43085	0.0060973	0.0060973\\
3.9273	0.42021	0.0063055	0.0063055\\
3.9357	0.40957	0.0066836	0.0066836\\
3.9432	0.39894	0.0065301	0.0065301\\
3.9533	0.3883	0.006	0.006\\
3.9637	0.37766	0.0060566	0.0060566\\
3.9738	0.36702	0.0063024	0.0063024\\
3.9848	0.35638	0.0065619	0.0065619\\
3.9951	0.34574	0.0066914	0.0066914\\
4.0063	0.33511	0.0065221	0.0065221\\
4.0195	0.32447	0.0062434	0.0062434\\
4.0316	0.31383	0.0062868	0.0062868\\
4.0451	0.30319	0.0062701	0.0062701\\
4.0541	0.29255	0.0063707	0.0063707\\
4.0665	0.28191	0.0063128	0.0063128\\
4.0793	0.27128	0.0063457	0.0063457\\
4.0929	0.26064	0.0064012	0.0064012\\
4.1075	0.25	0.0065842	0.0065842\\
4.1207	0.23936	0.0067354	0.0067354\\
4.1351	0.22872	0.0071575	0.0071575\\
4.1514	0.21809	0.0069122	0.0069122\\
4.1699	0.20745	0.0074422	0.0074422\\
4.1822	0.19681	0.0075591	0.0075591\\
4.1983	0.18617	0.0073044	0.0073044\\
4.2172	0.17553	0.0075554	0.0075554\\
4.2352	0.16489	0.0067471	0.0067471\\
4.2564	0.15426	0.0070369	0.0070369\\
4.2767	0.14362	0.0073888	0.0073888\\
4.2966	0.13298	0.0081069	0.0081069\\
4.3207	0.12234	0.0084621	0.0084621\\
4.3486	0.1117	0.0088645	0.0088645\\
4.3743	0.10106	0.0094713	0.0094713\\
4.4051	0.090426	0.0089463	0.0089463\\
4.439	0.079787	0.0095141	0.0095141\\
4.4763	0.069149	0.0092913	0.0092913\\
4.5194	0.058511	0.010771	0.010771\\
4.5656	0.047872	0.01169	0.01169\\
4.635	0.037234	0.012228	0.012228\\
4.7318	0.026596	0.015596	0.015596\\
4.8427	0.015957	0.017427	0.017427\\
5.1833	0.0053191	0.01767	0.01767\\
};
\addplot [color=mycolor1, mark repeat=9, mark phase=1, mark=+, mark options={solid, mycolor1}]
  table[row sep=crcr]{%
3.604	0.97462\\
3.6217	0.94187\\
3.6395	0.89995\\
3.6572	0.85753\\
3.675	0.8153\\
3.6927	0.77475\\
3.7105	0.73609\\
3.7282	0.69936\\
3.746	0.66447\\
3.7637	0.63133\\
3.7814	0.59985\\
3.7992	0.56994\\
3.8169	0.54153\\
3.8347	0.51454\\
3.8524	0.4889\\
3.8702	0.46455\\
3.8879	0.44141\\
3.9057	0.41942\\
3.9234	0.39854\\
3.9411	0.3787\\
3.9589	0.35985\\
3.9766	0.34194\\
3.9944	0.32493\\
4.0121	0.30876\\
4.0299	0.29341\\
4.0476	0.27881\\
4.0654	0.26495\\
4.0831	0.25178\\
4.1008	0.23927\\
4.1186	0.22738\\
4.1363	0.21608\\
4.1541	0.20535\\
4.1718	0.19515\\
4.1896	0.18546\\
4.2073	0.17625\\
4.2251	0.1675\\
4.2428	0.15919\\
4.2605	0.15129\\
4.2783	0.14378\\
4.296	0.13665\\
4.3138	0.12987\\
4.3315	0.12343\\
4.3493	0.11732\\
4.367	0.1115\\
4.3848	0.10597\\
4.4025	0.10072\\
4.4202	0.095734\\
4.438	0.090992\\
4.4557	0.086485\\
4.4735	0.082203\\
4.4912	0.078134\\
4.509	0.074266\\
4.5267	0.070591\\
4.5445	0.067098\\
4.5622	0.063779\\
4.5799	0.060624\\
4.5977	0.057626\\
4.6154	0.054777\\
4.6332	0.052069\\
4.6509	0.049496\\
4.6687	0.04705\\
4.6864	0.044725\\
4.7042	0.042515\\
4.7219	0.040416\\
4.7396	0.03842\\
4.7574	0.036523\\
4.7751	0.03472\\
4.7929	0.033006\\
4.8106	0.031377\\
4.8284	0.029829\\
4.8461	0.028358\\
4.8639	0.026959\\
4.8816	0.02563\\
4.8993	0.024366\\
4.9171	0.023165\\
4.9348	0.022023\\
4.9526	0.020938\\
4.9703	0.019906\\
4.9881	0.018925\\
5.0058	0.017993\\
5.0236	0.017107\\
5.0413	0.016265\\
5.0591	0.015464\\
5.0768	0.014703\\
5.0945	0.01398\\
5.1123	0.013292\\
5.13	0.012638\\
5.1478	0.012016\\
5.1655	0.011425\\
5.1833	0.010864\\
5.201	0.010329\\
5.2188	0.0098218\\
5.2365	0.0093391\\
5.2542	0.0088802\\
5.272	0.008444\\
5.2897	0.0080292\\
5.3075	0.0076349\\
5.3252	0.00726\\
5.343	0.0069036\\
5.3607	0.0065647\\
};
\addlegendentry{\ApproxParetoFitLegend}


\addplot[area legend, draw=none, fill=mycolor1, fill opacity=0.15, postaction={pattern={Lines[angle=90,distance={5pt},line width={0.5pt}]}, pattern color=mycolor1!60}, forget plot]
table[row sep=crcr] {%
x	y\\
3.604	0.98135\\
3.6217	0.95082\\
3.6395	0.90965\\
3.6572	0.86747\\
3.675	0.82471\\
3.6927	0.78342\\
3.7105	0.74396\\
3.7282	0.70651\\
3.746	0.67096\\
3.7637	0.63722\\
3.7814	0.60519\\
3.7992	0.5748\\
3.8169	0.54596\\
3.8347	0.51858\\
3.8524	0.49261\\
3.8702	0.46796\\
3.8879	0.44456\\
3.9057	0.42237\\
3.9234	0.4013\\
3.9411	0.38132\\
3.9589	0.36235\\
3.9766	0.34434\\
3.9944	0.32726\\
4.0121	0.31103\\
4.0299	0.29563\\
4.0476	0.28101\\
4.0654	0.26712\\
4.0831	0.25393\\
4.1008	0.2414\\
4.1186	0.22949\\
4.1363	0.21818\\
4.1541	0.20744\\
4.1718	0.19722\\
4.1896	0.18752\\
4.2073	0.17829\\
4.2251	0.16953\\
4.2428	0.16119\\
4.2605	0.15327\\
4.2783	0.14574\\
4.296	0.13859\\
4.3138	0.13178\\
4.3315	0.12532\\
4.3493	0.11917\\
4.367	0.11332\\
4.3848	0.10776\\
4.4025	0.10248\\
4.4202	0.097458\\
4.438	0.092682\\
4.4557	0.088141\\
4.4735	0.083823\\
4.4912	0.079718\\
4.509	0.075814\\
4.5267	0.072103\\
4.5445	0.068574\\
4.5622	0.065218\\
4.5799	0.062027\\
4.5977	0.058993\\
4.6154	0.056108\\
4.6332	0.053364\\
4.6509	0.050755\\
4.6687	0.048275\\
4.6864	0.045915\\
4.7042	0.043672\\
4.7219	0.041538\\
4.7396	0.03951\\
4.7574	0.03758\\
4.7751	0.035745\\
4.7929	0.034\\
4.8106	0.032341\\
4.8284	0.030763\\
4.8461	0.029262\\
4.8639	0.027834\\
4.8816	0.026477\\
4.8993	0.025186\\
4.9171	0.023958\\
4.9348	0.02279\\
4.9526	0.021679\\
4.9703	0.020623\\
4.9881	0.019618\\
5.0058	0.018662\\
5.0236	0.017753\\
5.0413	0.016889\\
5.0591	0.016067\\
5.0768	0.015285\\
5.0945	0.014541\\
5.1123	0.013833\\
5.13	0.01316\\
5.1478	0.01252\\
5.1655	0.011911\\
5.1833	0.011332\\
5.201	0.010781\\
5.2188	0.010257\\
5.2365	0.0097583\\
5.2542	0.0092841\\
5.272	0.0088331\\
5.2897	0.008404\\
5.3075	0.0079958\\
5.3252	0.0076076\\
5.343	0.0072382\\
5.3607	0.0068869\\
5.3607	0.0062426\\
5.343	0.006569\\
5.3252	0.0069124\\
5.3075	0.0072739\\
5.2897	0.0076544\\
5.272	0.0080548\\
5.2542	0.0084763\\
5.2365	0.0089198\\
5.2188	0.0093867\\
5.201	0.009878\\
5.1833	0.010395\\
5.1655	0.01094\\
5.1478	0.011512\\
5.13	0.012116\\
5.1123	0.01275\\
5.0945	0.013418\\
5.0768	0.014122\\
5.0591	0.014862\\
5.0413	0.015641\\
5.0236	0.016461\\
5.0058	0.017324\\
4.9881	0.018233\\
4.9703	0.019189\\
4.9526	0.020196\\
4.9348	0.021256\\
4.9171	0.022372\\
4.8993	0.023546\\
4.8816	0.024782\\
4.8639	0.026084\\
4.8461	0.027454\\
4.8284	0.028896\\
4.8106	0.030414\\
4.7929	0.032012\\
4.7751	0.033694\\
4.7574	0.035465\\
4.7396	0.03733\\
4.7219	0.039293\\
4.7042	0.041359\\
4.6864	0.043534\\
4.6687	0.045825\\
4.6509	0.048236\\
4.6332	0.050774\\
4.6154	0.053446\\
4.5977	0.05626\\
4.5799	0.059222\\
4.5622	0.06234\\
4.5445	0.065623\\
4.5267	0.069079\\
4.509	0.072718\\
4.4912	0.07655\\
4.4735	0.080583\\
4.4557	0.08483\\
4.438	0.089302\\
4.4202	0.094009\\
4.4025	0.098966\\
4.3848	0.10418\\
4.367	0.10968\\
4.3493	0.11546\\
4.3315	0.12155\\
4.3138	0.12796\\
4.296	0.13472\\
4.2783	0.14182\\
4.2605	0.1493\\
4.2428	0.15718\\
4.2251	0.16548\\
4.2073	0.17421\\
4.1896	0.1834\\
4.1718	0.19307\\
4.1541	0.20326\\
4.1363	0.21398\\
4.1186	0.22526\\
4.1008	0.23714\\
4.0831	0.24963\\
4.0654	0.26278\\
4.0476	0.27662\\
4.0299	0.29118\\
4.0121	0.30649\\
3.9944	0.3226\\
3.9766	0.33954\\
3.9589	0.35735\\
3.9411	0.37608\\
3.9234	0.39577\\
3.9057	0.41648\\
3.8879	0.43825\\
3.8702	0.46114\\
3.8524	0.4852\\
3.8347	0.5105\\
3.8169	0.53711\\
3.7992	0.56508\\
3.7814	0.5945\\
3.7637	0.62545\\
3.746	0.65799\\
3.7282	0.69221\\
3.7105	0.72821\\
3.6927	0.76608\\
3.675	0.80588\\
3.6572	0.84758\\
3.6395	0.89026\\
3.6217	0.93292\\
3.604	0.9679\\
}--cycle;
\addplot [color=mycolor1, mark repeat=9, mark phase=1, mark=x, mark options={solid, mycolor1}]
  table[row sep=crcr]{%
3.604	1.0047\\
3.6217	0.96208\\
3.6395	0.92076\\
3.6572	0.88075\\
3.675	0.84203\\
3.6927	0.80461\\
3.7105	0.76846\\
3.7282	0.73357\\
3.746	0.69992\\
3.7637	0.66749\\
3.7814	0.63627\\
3.7992	0.60622\\
3.8169	0.57734\\
3.8347	0.54958\\
3.8524	0.52293\\
3.8702	0.49735\\
3.8879	0.47282\\
3.9057	0.44932\\
3.9234	0.42681\\
3.9411	0.40526\\
3.9589	0.38464\\
3.9766	0.36493\\
3.9944	0.34609\\
4.0121	0.32809\\
4.0299	0.31091\\
4.0476	0.29452\\
4.0654	0.27888\\
4.0831	0.26398\\
4.1008	0.24977\\
4.1186	0.23624\\
4.1363	0.22335\\
4.1541	0.21108\\
4.1718	0.19941\\
4.1896	0.18831\\
4.2073	0.17776\\
4.2251	0.16773\\
4.2428	0.1582\\
4.2605	0.14915\\
4.2783	0.14056\\
4.296	0.1324\\
4.3138	0.12467\\
4.3315	0.11733\\
4.3493	0.11037\\
4.367	0.10378\\
4.3848	0.097535\\
4.4025	0.09162\\
4.4202	0.086021\\
4.438	0.080722\\
4.4557	0.07571\\
4.4735	0.07097\\
4.4912	0.066491\\
4.509	0.062259\\
4.5267	0.058262\\
4.5445	0.05449\\
4.5622	0.05093\\
4.5799	0.047574\\
4.5977	0.044409\\
4.6154	0.041429\\
4.6332	0.038622\\
4.6509	0.03598\\
4.6687	0.033495\\
4.6864	0.031159\\
4.7042	0.028964\\
4.7219	0.026903\\
4.7396	0.02497\\
4.7574	0.023156\\
4.7751	0.021457\\
4.7929	0.019867\\
4.8106	0.018378\\
4.8284	0.016987\\
4.8461	0.015687\\
4.8639	0.014474\\
4.8816	0.013343\\
4.8993	0.01229\\
4.9171	0.011309\\
4.9348	0.010398\\
4.9526	0.009551\\
4.9703	0.0087656\\
4.9881	0.0080377\\
5.0058	0.0073638\\
5.0236	0.0067407\\
5.0413	0.006165\\
5.0591	0.0056339\\
5.0768	0.0051443\\
5.0945	0.0046935\\
5.1123	0.0042789\\
5.13	0.003898\\
5.1478	0.0035483\\
5.1655	0.0032277\\
5.1833	0.002934\\
5.201	0.0026652\\
5.2188	0.0024195\\
5.2365	0.0021949\\
5.2542	0.00199\\
5.272	0.001803\\
5.2897	0.0016326\\
5.3075	0.0014774\\
5.3252	0.0013362\\
5.343	0.0012078\\
5.3607	0.0010911\\
};
\addlegendentry{\ApproxBLNFitLegend}


\addplot[area legend, draw=none, fill=mycolor1, fill opacity=0.15, postaction={pattern={Lines[angle=45,distance={5pt},line width={0.5pt}]}, pattern color=mycolor1!60}, forget plot]
table[row sep=crcr] {%
x	y\\
3.604	1.0131\\
3.6217	0.97013\\
3.6395	0.92849\\
3.6572	0.88817\\
3.675	0.84916\\
3.6927	0.81145\\
3.7105	0.77504\\
3.7282	0.7399\\
3.746	0.70603\\
3.7637	0.67339\\
3.7814	0.64197\\
3.7992	0.61174\\
3.8169	0.58268\\
3.8347	0.55476\\
3.8524	0.52795\\
3.8702	0.50223\\
3.8879	0.47756\\
3.9057	0.45392\\
3.9234	0.43127\\
3.9411	0.40958\\
3.9589	0.38883\\
3.9766	0.36899\\
3.9944	0.35002\\
4.0121	0.33189\\
4.0299	0.31458\\
4.0476	0.29806\\
4.0654	0.2823\\
4.0831	0.26726\\
4.1008	0.25293\\
4.1186	0.23927\\
4.1363	0.22627\\
4.1541	0.21388\\
4.1718	0.2021\\
4.1896	0.19089\\
4.2073	0.18023\\
4.2251	0.1701\\
4.2428	0.16047\\
4.2605	0.15133\\
4.2783	0.14265\\
4.296	0.13441\\
4.3138	0.1266\\
4.3315	0.11918\\
4.3493	0.11216\\
4.367	0.1055\\
4.3848	0.099196\\
4.4025	0.093226\\
4.4202	0.087575\\
4.438	0.082228\\
4.4557	0.077171\\
4.4735	0.07239\\
4.4912	0.067872\\
4.509	0.063603\\
4.5267	0.059573\\
4.5445	0.055768\\
4.5622	0.052178\\
4.5799	0.048792\\
4.5977	0.045599\\
4.6154	0.042589\\
4.6332	0.039754\\
4.6509	0.037084\\
4.6687	0.03457\\
4.6864	0.032205\\
4.7042	0.029981\\
4.7219	0.02789\\
4.7396	0.025925\\
4.7574	0.02408\\
4.7751	0.022349\\
4.7929	0.020725\\
4.8106	0.019203\\
4.8284	0.017778\\
4.8461	0.016444\\
4.8639	0.015197\\
4.8816	0.014033\\
4.8993	0.012946\\
4.9171	0.011933\\
4.9348	0.01099\\
4.9526	0.010113\\
4.9703	0.0092978\\
4.9881	0.0085412\\
5.0058	0.0078399\\
5.0236	0.0071904\\
5.0413	0.0065896\\
5.0591	0.0060345\\
5.0768	0.0055221\\
5.0945	0.0050496\\
5.1123	0.0046143\\
5.13	0.0042137\\
5.1478	0.0038452\\
5.1655	0.0035068\\
5.1833	0.003196\\
5.201	0.002911\\
5.2188	0.0026497\\
5.2365	0.0024104\\
5.2542	0.0021914\\
5.272	0.0019911\\
5.2897	0.001808\\
5.3075	0.0016408\\
5.3252	0.0014882\\
5.343	0.0013489\\
5.3607	0.001222\\
5.3607	0.00096016\\
5.343	0.0010667\\
5.3252	0.0011843\\
5.3075	0.0013141\\
5.2897	0.0014572\\
5.272	0.0016149\\
5.2542	0.0017885\\
5.2365	0.0019794\\
5.2188	0.0021892\\
5.201	0.0024195\\
5.1833	0.002672\\
5.1655	0.0029487\\
5.1478	0.0032514\\
5.13	0.0035823\\
5.1123	0.0039435\\
5.0945	0.0043374\\
5.0768	0.0047664\\
5.0591	0.0052332\\
5.0413	0.0057404\\
5.0236	0.0062909\\
5.0058	0.0068878\\
4.9881	0.0075341\\
4.9703	0.0082334\\
4.9526	0.0089891\\
4.9348	0.0098049\\
4.9171	0.010685\\
4.8993	0.011633\\
4.8816	0.012653\\
4.8639	0.013751\\
4.8461	0.01493\\
4.8284	0.016196\\
4.8106	0.017554\\
4.7929	0.019009\\
4.7751	0.020566\\
4.7574	0.022233\\
4.7396	0.024014\\
4.7219	0.025917\\
4.7042	0.027947\\
4.6864	0.030112\\
4.6687	0.032419\\
4.6509	0.034876\\
4.6332	0.037489\\
4.6154	0.040268\\
4.5977	0.04322\\
4.5799	0.046356\\
4.5622	0.049683\\
4.5445	0.053211\\
4.5267	0.056952\\
4.509	0.060914\\
4.4912	0.06511\\
4.4735	0.069551\\
4.4557	0.074249\\
4.438	0.079216\\
4.4202	0.084467\\
4.4025	0.090015\\
4.3848	0.095874\\
4.367	0.10206\\
4.3493	0.10859\\
4.3315	0.11548\\
4.3138	0.12274\\
4.296	0.1304\\
4.2783	0.13847\\
4.2605	0.14697\\
4.2428	0.15593\\
4.2251	0.16536\\
4.2073	0.17529\\
4.1896	0.18574\\
4.1718	0.19673\\
4.1541	0.20828\\
4.1363	0.22043\\
4.1186	0.2332\\
4.1008	0.24661\\
4.0831	0.26069\\
4.0654	0.27547\\
4.0476	0.29098\\
4.0299	0.30724\\
4.0121	0.32429\\
3.9944	0.34216\\
3.9766	0.36086\\
3.9589	0.38045\\
3.9411	0.40093\\
3.9234	0.42234\\
3.9057	0.44472\\
3.8879	0.46809\\
3.8702	0.49247\\
3.8524	0.5179\\
3.8347	0.5444\\
3.8169	0.57199\\
3.7992	0.60071\\
3.7814	0.63057\\
3.7637	0.66159\\
3.746	0.69381\\
3.7282	0.72723\\
3.7105	0.76187\\
3.6927	0.79776\\
3.675	0.83491\\
3.6572	0.87333\\
3.6395	0.91303\\
3.6217	0.95403\\
3.604	0.99633\\
}--cycle;
      % \input{../../../../.tex/20/KS/SizeDistributionFittingsfig33.tex}
      \legend{}
    \end{groupplot}

    \begin{groupplot}[
      group style={
        group name=Row5,
        plotLastGroup,
      },
      plotRow,
      row3Keys,
      plotLegendSW,
      %
      % xmin=3637.5,xmax=2.5003e+05,
      ymin=0.001,ymax=1.0148,
    ]
      \nextgroupplot[%
        at={($(Row4 c1r1.south west)-(0,\yPlotSep em)-(0,\plotHeight cm)$)},
        xmin=3.5608,xmax=5.3101,
        % xmin=3.4,xmax=5.4,
        % ymin=0.0041299,ymax=1.0005,
      ] % This file was created by matlab2tikz.
%
\definecolor{mycolor1}{rgb}{0.00000,0.44700,0.74100}%
%
\addplot[only marks, mark=*, mark size=1.1180pt, color=black, fill=black, opacity=0.60, draw=none, mark options={draw=none,line width=0pt}] table[row sep=crcr]{%
x	y\\
3.5623	0.99468\\
3.5723	0.98404\\
3.573	0.9734\\
3.5753	0.96277\\
3.5766	0.95213\\
3.5781	0.94149\\
3.5826	0.93085\\
3.5851	0.92021\\
3.5911	0.90957\\
3.5942	0.89894\\
3.596	0.8883\\
3.6016	0.87766\\
3.603	0.86702\\
3.6047	0.85638\\
3.6086	0.84574\\
3.6086	0.83511\\
3.6097	0.82447\\
3.6142	0.81383\\
3.6261	0.80319\\
3.6383	0.79255\\
3.6464	0.78191\\
3.6535	0.77128\\
3.657	0.76064\\
3.6658	0.75\\
3.6678	0.73936\\
3.6687	0.72872\\
3.6702	0.71809\\
3.6723	0.70745\\
3.6793	0.69681\\
3.6867	0.68617\\
3.6949	0.67553\\
3.7015	0.66489\\
3.7031	0.65426\\
3.719	0.64362\\
3.7213	0.63298\\
3.7265	0.62234\\
3.7307	0.6117\\
3.737	0.60106\\
3.7374	0.59043\\
3.7423	0.57979\\
3.7487	0.56915\\
3.7677	0.55851\\
3.7685	0.54787\\
3.7768	0.53723\\
3.8005	0.5266\\
3.8015	0.51596\\
3.8022	0.50532\\
3.8032	0.49468\\
3.8039	0.48404\\
3.8108	0.4734\\
3.8214	0.46277\\
3.8405	0.45213\\
3.8605	0.44149\\
3.863	0.43085\\
3.8645	0.42021\\
3.8662	0.40957\\
3.8806	0.39894\\
3.8878	0.3883\\
3.8964	0.37766\\
3.905	0.36702\\
3.9294	0.35638\\
3.9303	0.34574\\
3.954	0.33511\\
3.9604	0.32447\\
3.9669	0.31383\\
3.9747	0.30319\\
3.9822	0.29255\\
3.9929	0.28191\\
4.0051	0.27128\\
4.0144	0.26064\\
4.0161	0.25\\
4.0433	0.23936\\
4.0578	0.22872\\
4.073	0.21809\\
4.0857	0.20745\\
4.0904	0.19681\\
4.1168	0.18617\\
4.1265	0.17553\\
4.127	0.16489\\
4.143	0.15426\\
4.1756	0.14362\\
4.2156	0.13298\\
4.3116	0.12234\\
4.3276	0.1117\\
4.3662	0.10106\\
4.4791	0.090426\\
4.4912	0.079787\\
4.517	0.069149\\
4.5743	0.058511\\
4.5914	0.047872\\
4.6138	0.037234\\
4.7898	0.026596\\
5.0876	0.015957\\
5.3101	0.0053191\\
};
\addlegendentry{\RealEmpiricCCDFLegend}

\addplot [color=black, mark repeat=9, mark phase=1, mark=+, mark options={solid, black}]
  table[row sep=crcr]{%
3.5608	1\\
3.5785	0.94954\\
3.5961	0.90162\\
3.6138	0.85612\\
3.6315	0.81292\\
3.6492	0.7719\\
3.6668	0.73295\\
3.6845	0.69596\\
3.7022	0.66084\\
3.7198	0.6275\\
3.7375	0.59583\\
3.7552	0.56576\\
3.7728	0.53721\\
3.7905	0.5101\\
3.8082	0.48436\\
3.8259	0.45992\\
3.8435	0.43671\\
3.8612	0.41468\\
3.8789	0.39375\\
3.8965	0.37388\\
3.9142	0.35501\\
3.9319	0.3371\\
3.9495	0.32009\\
3.9672	0.30394\\
3.9849	0.2886\\
4.0026	0.27404\\
4.0202	0.26021\\
4.0379	0.24708\\
4.0556	0.23461\\
4.0732	0.22277\\
4.0909	0.21153\\
4.1086	0.20085\\
4.1262	0.19072\\
4.1439	0.18109\\
4.1616	0.17196\\
4.1793	0.16328\\
4.1969	0.15504\\
4.2146	0.14722\\
4.2323	0.13979\\
4.2499	0.13273\\
4.2676	0.12603\\
4.2853	0.11967\\
4.3029	0.11364\\
4.3206	0.1079\\
4.3383	0.10246\\
4.356	0.097286\\
4.3736	0.092377\\
4.3913	0.087716\\
4.409	0.083289\\
4.4266	0.079086\\
4.4443	0.075095\\
4.462	0.071306\\
4.4796	0.067708\\
4.4973	0.064291\\
4.515	0.061047\\
4.5327	0.057966\\
4.5503	0.055041\\
4.568	0.052264\\
4.5857	0.049626\\
4.6033	0.047122\\
4.621	0.044744\\
4.6387	0.042486\\
4.6563	0.040342\\
4.674	0.038307\\
4.6917	0.036373\\
4.7094	0.034538\\
4.727	0.032795\\
4.7447	0.03114\\
4.7624	0.029569\\
4.78	0.028077\\
4.7977	0.02666\\
4.8154	0.025315\\
4.833	0.024037\\
4.8507	0.022824\\
4.8684	0.021672\\
4.8861	0.020579\\
4.9037	0.01954\\
4.9214	0.018554\\
4.9391	0.017618\\
4.9567	0.016729\\
4.9744	0.015885\\
4.9921	0.015083\\
5.0097	0.014322\\
5.0274	0.013599\\
5.0451	0.012913\\
5.0628	0.012261\\
5.0804	0.011643\\
5.0981	0.011055\\
5.1158	0.010497\\
5.1334	0.0099676\\
5.1511	0.0094646\\
5.1688	0.008987\\
5.1864	0.0085335\\
5.2041	0.0081029\\
5.2218	0.007694\\
5.2395	0.0073058\\
5.2571	0.0069371\\
5.2748	0.006587\\
5.2925	0.0062546\\
5.3101	0.005939\\
};
\addlegendentry{\RealParetoFitLegend}

\addplot [color=black, mark repeat=9, mark phase=1, mark=x, mark options={solid, black}]
  table[row sep=crcr]{%
3.5608	1.0005\\
3.5785	0.95316\\
3.5961	0.90742\\
3.6138	0.86328\\
3.6315	0.82076\\
3.6492	0.77986\\
3.6668	0.74057\\
3.6845	0.70289\\
3.7022	0.66682\\
3.7198	0.63232\\
3.7375	0.59938\\
3.7552	0.56796\\
3.7728	0.53805\\
3.7905	0.5096\\
3.8082	0.48257\\
3.8259	0.45692\\
3.8435	0.43262\\
3.8612	0.4096\\
3.8789	0.38783\\
3.8965	0.36726\\
3.9142	0.34784\\
3.9319	0.32951\\
3.9495	0.31222\\
3.9672	0.29594\\
3.9849	0.28059\\
4.0026	0.26614\\
4.0202	0.25254\\
4.0379	0.23974\\
4.0556	0.22769\\
4.0732	0.21634\\
4.0909	0.20566\\
4.1086	0.19561\\
4.1262	0.18613\\
4.1439	0.1772\\
4.1616	0.16878\\
4.1793	0.16082\\
4.1969	0.15331\\
4.2146	0.14621\\
4.2323	0.1395\\
4.2499	0.13313\\
4.2676	0.1271\\
4.2853	0.12137\\
4.3029	0.11592\\
4.3206	0.11074\\
4.3383	0.10581\\
4.356	0.1011\\
4.3736	0.09661\\
4.3913	0.092318\\
4.409	0.088212\\
4.4266	0.084281\\
4.4443	0.080514\\
4.462	0.076901\\
4.4796	0.073434\\
4.4973	0.070104\\
4.515	0.066905\\
4.5327	0.06383\\
4.5503	0.060873\\
4.568	0.058029\\
4.5857	0.055293\\
4.6033	0.052661\\
4.621	0.050128\\
4.6387	0.047691\\
4.6563	0.045347\\
4.674	0.043093\\
4.6917	0.040926\\
4.7094	0.038844\\
4.727	0.036843\\
4.7447	0.034921\\
4.7624	0.033077\\
4.78	0.031309\\
4.7977	0.029613\\
4.8154	0.027989\\
4.833	0.026433\\
4.8507	0.024946\\
4.8684	0.023523\\
4.8861	0.022165\\
4.9037	0.020868\\
4.9214	0.019632\\
4.9391	0.018454\\
4.9567	0.017332\\
4.9744	0.016265\\
4.9921	0.015251\\
5.0097	0.014288\\
5.0274	0.013375\\
5.0451	0.01251\\
5.0628	0.01169\\
5.0804	0.010915\\
5.0981	0.010182\\
5.1158	0.0094906\\
5.1334	0.0088382\\
5.1511	0.0082235\\
5.1688	0.0076448\\
5.1864	0.0071005\\
5.2041	0.0065892\\
5.2218	0.0061093\\
5.2395	0.0056592\\
5.2571	0.0052377\\
5.2748	0.0048432\\
5.2925	0.0044744\\
5.3101	0.0041299\\
};
\addlegendentry{\RealBLNFitLegend}

      % \input{../../../../.tex/20/KS/SizeDistributionFittingsfig40.tex}
      \legend{}
  
      \nextgroupplot[%
        xmin=3.598,xmax=5.4115,
        % xmin=3.5,xmax=5.5,
        % ymin=0.001,ymax=1.0148,
      ] % This file was created by matlab2tikz.
%
\definecolor{mycolor1}{rgb}{0.00000,0.44706,0.69804}%
\definecolor{mycolor2}{rgb}{0.00000,0.44700,0.74100}%
%
\addplot[only marks, mark=*, mark size=1.1180pt, color=mycolor1, fill=mycolor1, opacity=0.60, draw=none, mark options={draw=none,line width=0pt}] table[row sep=crcr]{%
x	y\\
3.5984	0.99468\\
3.6027	0.98404\\
3.6068	0.9734\\
3.6112	0.96277\\
3.6162	0.95213\\
3.6211	0.94149\\
3.6266	0.93085\\
3.6309	0.92021\\
3.636	0.90957\\
3.64	0.89894\\
3.6441	0.8883\\
3.6477	0.87766\\
3.6529	0.86702\\
3.6581	0.85638\\
3.6641	0.84574\\
3.6687	0.83511\\
3.674	0.82447\\
3.6794	0.81383\\
3.6854	0.80319\\
3.6908	0.79255\\
3.6952	0.78191\\
3.7011	0.77128\\
3.707	0.76064\\
3.713	0.75\\
3.7182	0.73936\\
3.7241	0.72872\\
3.73	0.71809\\
3.736	0.70745\\
3.7427	0.69681\\
3.7494	0.68617\\
3.7566	0.67553\\
3.7628	0.66489\\
3.7689	0.65426\\
3.7752	0.64362\\
3.7809	0.63298\\
3.7869	0.62234\\
3.7932	0.6117\\
3.8004	0.60106\\
3.8084	0.59043\\
3.8147	0.57979\\
3.8215	0.56915\\
3.8298	0.55851\\
3.8363	0.54787\\
3.8438	0.53723\\
3.8521	0.5266\\
3.8604	0.51596\\
3.8683	0.50532\\
3.8765	0.49468\\
3.8843	0.48404\\
3.8928	0.4734\\
3.9007	0.46277\\
3.9098	0.45213\\
3.9208	0.44149\\
3.9308	0.43085\\
3.9405	0.42021\\
3.9501	0.40957\\
3.9606	0.39894\\
3.9709	0.3883\\
3.9801	0.37766\\
3.9899	0.36702\\
3.999	0.35638\\
4.0096	0.34574\\
4.0188	0.33511\\
4.0298	0.32447\\
4.0417	0.31383\\
4.0542	0.30319\\
4.0666	0.29255\\
4.0789	0.28191\\
4.0918	0.27128\\
4.1041	0.26064\\
4.1175	0.25\\
4.131	0.23936\\
4.1454	0.22872\\
4.159	0.21809\\
4.173	0.20745\\
4.1912	0.19681\\
4.2086	0.18617\\
4.2269	0.17553\\
4.2454	0.16489\\
4.2692	0.15426\\
4.2908	0.14362\\
4.3146	0.13298\\
4.3388	0.12234\\
4.362	0.1117\\
4.3887	0.10106\\
4.4173	0.090426\\
4.4494	0.079787\\
4.4812	0.069149\\
4.5257	0.058511\\
4.5784	0.047872\\
4.6417	0.037234\\
4.741	0.026596\\
4.8583	0.015957\\
5.2036	0.0053191\\
};
\addlegendentry{\ExactEmpiricCCDFLegend}

\addplot [color=mycolor1, only marks, every error bar/.append style={opacity=0.30}, mark=*, mark size=0pt, draw=none, forget plot]
 plot [error bars/.cd, x dir=both, x explicit, error bar style={line width=1pt, color=mycolor1}, error mark options={mark=none,mark size=0pt}]
 table[row sep=crcr, x error plus index=2, x error minus index=3]{%
3.5984	0.99468	0.00473	0.00473\\
3.6027	0.98404	0.0047925	0.0047925\\
3.6068	0.9734	0.0048115	0.0048115\\
3.6112	0.96277	0.0047892	0.0047892\\
3.6162	0.95213	0.0047066	0.0047066\\
3.6211	0.94149	0.0049168	0.0049168\\
3.6266	0.93085	0.0048041	0.0048041\\
3.6309	0.92021	0.004932	0.004932\\
3.636	0.90957	0.0049959	0.0049959\\
3.64	0.89894	0.005017	0.005017\\
3.6441	0.8883	0.0051081	0.0051081\\
3.6477	0.87766	0.005056	0.005056\\
3.6529	0.86702	0.0050923	0.0050923\\
3.6581	0.85638	0.0052229	0.0052229\\
3.6641	0.84574	0.005148	0.005148\\
3.6687	0.83511	0.0052242	0.0052242\\
3.674	0.82447	0.0051186	0.0051186\\
3.6794	0.81383	0.005215	0.005215\\
3.6854	0.80319	0.0051265	0.0051265\\
3.6908	0.79255	0.0052431	0.0052431\\
3.6952	0.78191	0.0051884	0.0051884\\
3.7011	0.77128	0.0051923	0.0051923\\
3.707	0.76064	0.0050709	0.0050709\\
3.713	0.75	0.005022	0.005022\\
3.7182	0.73936	0.0050527	0.0050527\\
3.7241	0.72872	0.0052125	0.0052125\\
3.73	0.71809	0.0052284	0.0052284\\
3.736	0.70745	0.0051795	0.0051795\\
3.7427	0.69681	0.0052847	0.0052847\\
3.7494	0.68617	0.0054013	0.0054013\\
3.7566	0.67553	0.005546	0.005546\\
3.7628	0.66489	0.0056581	0.0056581\\
3.7689	0.65426	0.0055104	0.0055104\\
3.7752	0.64362	0.0056133	0.0056133\\
3.7809	0.63298	0.0055232	0.0055232\\
3.7869	0.62234	0.0054689	0.0054689\\
3.7932	0.6117	0.0054553	0.0054553\\
3.8004	0.60106	0.0057052	0.0057052\\
3.8084	0.59043	0.0058848	0.0058848\\
3.8147	0.57979	0.0057587	0.0057587\\
3.8215	0.56915	0.0059265	0.0059265\\
3.8298	0.55851	0.0061686	0.0061686\\
3.8363	0.54787	0.0061957	0.0061957\\
3.8438	0.53723	0.0060905	0.0060905\\
3.8521	0.5266	0.0059675	0.0059675\\
3.8604	0.51596	0.0058526	0.0058526\\
3.8683	0.50532	0.0057498	0.0057498\\
3.8765	0.49468	0.0057298	0.0057298\\
3.8843	0.48404	0.0056797	0.0056797\\
3.8928	0.4734	0.0056628	0.0056628\\
3.9007	0.46277	0.0055566	0.0055566\\
3.9098	0.45213	0.0053166	0.0053166\\
3.9208	0.44149	0.0054396	0.0054396\\
3.9308	0.43085	0.005428	0.005428\\
3.9405	0.42021	0.0055337	0.0055337\\
3.9501	0.40957	0.0053756	0.0053756\\
3.9606	0.39894	0.0054857	0.0054857\\
3.9709	0.3883	0.0055337	0.0055337\\
3.9801	0.37766	0.0056598	0.0056598\\
3.9899	0.36702	0.0058367	0.0058367\\
3.999	0.35638	0.0056606	0.0056606\\
4.0096	0.34574	0.0056459	0.0056459\\
4.0188	0.33511	0.0059698	0.0059698\\
4.0298	0.32447	0.0059241	0.0059241\\
4.0417	0.31383	0.0056343	0.0056343\\
4.0542	0.30319	0.0055073	0.0055073\\
4.0666	0.29255	0.0056644	0.0056644\\
4.0789	0.28191	0.005727	0.005727\\
4.0918	0.27128	0.0057123	0.0057123\\
4.1041	0.26064	0.0056946	0.0056946\\
4.1175	0.25	0.0059084	0.0059084\\
4.131	0.23936	0.0061124	0.0061124\\
4.1454	0.22872	0.0066665	0.0066665\\
4.159	0.21809	0.0068485	0.0068485\\
4.173	0.20745	0.0069575	0.0069575\\
4.1912	0.19681	0.0072717	0.0072717\\
4.2086	0.18617	0.0078688	0.0078688\\
4.2269	0.17553	0.0077805	0.0077805\\
4.2454	0.16489	0.0078623	0.0078623\\
4.2692	0.15426	0.0083156	0.0083156\\
4.2908	0.14362	0.0089906	0.0089906\\
4.3146	0.13298	0.009299	0.009299\\
4.3388	0.12234	0.0084958	0.0084958\\
4.362	0.1117	0.0078765	0.0078765\\
4.3887	0.10106	0.008669	0.008669\\
4.4173	0.090426	0.0097937	0.0097937\\
4.4494	0.079787	0.01038	0.01038\\
4.4812	0.069149	0.011015	0.011015\\
4.5257	0.058511	0.012142	0.012142\\
4.5784	0.047872	0.011936	0.011936\\
4.6417	0.037234	0.014498	0.014498\\
4.741	0.026596	0.016486	0.016486\\
4.8583	0.015957	0.020509	0.020509\\
5.2036	0.0053191	0.018099	0.018099\\
};
\addplot [color=mycolor1, mark repeat=9, mark phase=1, mark=+, mark options={solid, mycolor1}]
  table[row sep=crcr]{%
3.598	0.97291\\
3.6163	0.94096\\
3.6346	0.90043\\
3.6529	0.85738\\
3.6713	0.81511\\
3.6896	0.77446\\
3.7079	0.73583\\
3.7262	0.69915\\
3.7445	0.66429\\
3.7628	0.63118\\
3.7812	0.59973\\
3.7995	0.56985\\
3.8178	0.54146\\
3.8361	0.5145\\
3.8544	0.48888\\
3.8728	0.46454\\
3.8911	0.44142\\
3.9094	0.41945\\
3.9277	0.39858\\
3.946	0.37875\\
3.9643	0.35991\\
3.9827	0.34201\\
4.001	0.325\\
4.0193	0.30885\\
4.0376	0.2935\\
4.0559	0.27891\\
4.0743	0.26505\\
4.0926	0.25189\\
4.1109	0.23937\\
4.1292	0.22749\\
4.1475	0.21619\\
4.1658	0.20546\\
4.1842	0.19526\\
4.2025	0.18557\\
4.2208	0.17636\\
4.2391	0.16762\\
4.2574	0.1593\\
4.2758	0.1514\\
4.2941	0.1439\\
4.3124	0.13676\\
4.3307	0.12998\\
4.349	0.12354\\
4.3673	0.11742\\
4.3857	0.1116\\
4.404	0.10608\\
4.4223	0.10082\\
4.4406	0.09583\\
4.4589	0.091086\\
4.4773	0.086578\\
4.4956	0.082293\\
4.5139	0.078221\\
4.5322	0.074351\\
4.5505	0.070673\\
4.5688	0.067178\\
4.5872	0.063856\\
4.6055	0.060699\\
4.6238	0.057699\\
4.6421	0.054847\\
4.6604	0.052137\\
4.6788	0.049561\\
4.6971	0.047112\\
4.7154	0.044785\\
4.7337	0.042574\\
4.752	0.040472\\
4.7703	0.038474\\
4.7887	0.036574\\
4.807	0.034769\\
4.8253	0.033054\\
4.8436	0.031423\\
4.8619	0.029873\\
4.8803	0.0284\\
4.8986	0.026999\\
4.9169	0.025668\\
4.9352	0.024403\\
4.9535	0.0232\\
4.9719	0.022056\\
4.9902	0.02097\\
5.0085	0.019937\\
5.0268	0.018954\\
5.0451	0.018021\\
5.0634	0.017134\\
5.0818	0.01629\\
5.1001	0.015488\\
5.1184	0.014726\\
5.1367	0.014001\\
5.155	0.013312\\
5.1734	0.012657\\
5.1917	0.012034\\
5.21	0.011443\\
5.2283	0.01088\\
5.2466	0.010345\\
5.2649	0.0098364\\
5.2833	0.0093529\\
5.3016	0.0088932\\
5.3199	0.0084562\\
5.3382	0.0080408\\
5.3565	0.0076458\\
5.3749	0.0072702\\
5.3932	0.0069132\\
5.4115	0.0065737\\
};
\addlegendentry{\ExactParetoFitLegend}


\addplot[area legend, draw=none, fill=mycolor1, fill opacity=0.15, postaction={pattern={Lines[angle=90,distance={5pt},line width={0.5pt}]}, pattern color=mycolor1!60}, forget plot]
table[row sep=crcr] {%
x	y\\
3.598	0.97938\\
3.6163	0.9501\\
3.6346	0.91064\\
3.6529	0.86752\\
3.6713	0.82464\\
3.6896	0.78311\\
3.7079	0.74369\\
3.7262	0.70627\\
3.7445	0.67075\\
3.7628	0.63704\\
3.7812	0.60505\\
3.7995	0.57468\\
3.8178	0.54587\\
3.8361	0.51852\\
3.8544	0.49256\\
3.8728	0.46793\\
3.8911	0.44456\\
3.9094	0.42238\\
3.9277	0.40133\\
3.946	0.38135\\
3.9643	0.36239\\
3.9827	0.3444\\
4.001	0.32732\\
4.0193	0.3111\\
4.0376	0.29571\\
4.0559	0.28109\\
4.0743	0.2672\\
4.0926	0.25401\\
4.1109	0.24148\\
4.1292	0.22958\\
4.1475	0.21826\\
4.1658	0.20752\\
4.1842	0.1973\\
4.2025	0.1876\\
4.2208	0.17837\\
4.2391	0.1696\\
4.2574	0.16126\\
4.2758	0.15334\\
4.2941	0.14581\\
4.3124	0.13865\\
4.3307	0.13184\\
4.349	0.12537\\
4.3673	0.11922\\
4.3857	0.11337\\
4.404	0.10781\\
4.4223	0.10253\\
4.4406	0.0975\\
4.4589	0.092721\\
4.4773	0.088177\\
4.4956	0.083857\\
4.5139	0.079749\\
4.5322	0.075843\\
4.5505	0.072129\\
4.5688	0.068598\\
4.5872	0.06524\\
4.6055	0.062047\\
4.6238	0.05901\\
4.6421	0.056123\\
4.6604	0.053377\\
4.6788	0.050767\\
4.6971	0.048284\\
4.7154	0.045923\\
4.7337	0.043678\\
4.752	0.041543\\
4.7703	0.039512\\
4.7887	0.037582\\
4.807	0.035745\\
4.8253	0.033999\\
4.8436	0.032339\\
4.8619	0.030759\\
4.8803	0.029257\\
4.8986	0.027829\\
4.9169	0.02647\\
4.9352	0.025178\\
4.9535	0.02395\\
4.9719	0.022781\\
4.9902	0.02167\\
5.0085	0.020613\\
5.0268	0.019607\\
5.0451	0.018651\\
5.0634	0.017742\\
5.0818	0.016877\\
5.1001	0.016054\\
5.1184	0.015272\\
5.1367	0.014528\\
5.155	0.01382\\
5.1734	0.013146\\
5.1917	0.012506\\
5.21	0.011897\\
5.2283	0.011318\\
5.2466	0.010767\\
5.2649	0.010242\\
5.2833	0.0097439\\
5.3016	0.0092697\\
5.3199	0.0088186\\
5.3382	0.0083895\\
5.3565	0.0079814\\
5.3749	0.0075931\\
5.3932	0.0072238\\
5.4115	0.0068726\\
5.4115	0.0062749\\
5.3932	0.0066025\\
5.3749	0.0069473\\
5.3565	0.0073102\\
5.3382	0.007692\\
5.3199	0.0080939\\
5.3016	0.0085168\\
5.2833	0.0089619\\
5.2649	0.0094303\\
5.2466	0.0099233\\
5.2283	0.010442\\
5.21	0.010988\\
5.1917	0.011563\\
5.1734	0.012168\\
5.155	0.012804\\
5.1367	0.013474\\
5.1184	0.01418\\
5.1001	0.014922\\
5.0818	0.015703\\
5.0634	0.016525\\
5.0451	0.017391\\
5.0268	0.018302\\
5.0085	0.01926\\
4.9902	0.02027\\
4.9719	0.021332\\
4.9535	0.02245\\
4.9352	0.023627\\
4.9169	0.024865\\
4.8986	0.026169\\
4.8803	0.027542\\
4.8619	0.028987\\
4.8436	0.030507\\
4.8253	0.032108\\
4.807	0.033793\\
4.7887	0.035567\\
4.7703	0.037435\\
4.752	0.0394\\
4.7337	0.041469\\
4.7154	0.043648\\
4.6971	0.045941\\
4.6788	0.048355\\
4.6604	0.050896\\
4.6421	0.053571\\
4.6238	0.056387\\
4.6055	0.059352\\
4.5872	0.062473\\
4.5688	0.065758\\
4.5505	0.069217\\
4.5322	0.072859\\
4.5139	0.076693\\
4.4956	0.080729\\
4.4773	0.084978\\
4.4589	0.089451\\
4.4406	0.094161\\
4.4223	0.099119\\
4.404	0.10434\\
4.3857	0.10983\\
4.3673	0.11562\\
4.349	0.12171\\
4.3307	0.12812\\
4.3124	0.13487\\
4.2941	0.14198\\
4.2758	0.14946\\
4.2574	0.15734\\
4.2391	0.16563\\
4.2208	0.17436\\
4.2025	0.18355\\
4.1842	0.19322\\
4.1658	0.2034\\
4.1475	0.21412\\
4.1292	0.2254\\
4.1109	0.23727\\
4.0926	0.24976\\
4.0743	0.26291\\
4.0559	0.27674\\
4.0376	0.29129\\
4.0193	0.30659\\
4.001	0.32269\\
3.9827	0.33962\\
3.9643	0.35742\\
3.946	0.37614\\
3.9277	0.39582\\
3.9094	0.41652\\
3.8911	0.43827\\
3.8728	0.46115\\
3.8544	0.48519\\
3.8361	0.51048\\
3.8178	0.53706\\
3.7995	0.56502\\
3.7812	0.59441\\
3.7628	0.62532\\
3.7445	0.65783\\
3.7262	0.69202\\
3.7079	0.72798\\
3.6896	0.76581\\
3.6713	0.80559\\
3.6529	0.84724\\
3.6346	0.89022\\
3.6163	0.93181\\
3.598	0.96644\\
}--cycle;
\addplot [color=mycolor1, mark repeat=9, mark phase=1, mark=x, mark options={solid, mycolor1}]
  table[row sep=crcr]{%
3.598	1.006\\
3.6163	0.96437\\
3.6346	0.92392\\
3.6529	0.88464\\
3.6713	0.84655\\
3.6896	0.80964\\
3.7079	0.77389\\
3.7262	0.73931\\
3.7445	0.70588\\
3.7628	0.6736\\
3.7812	0.64243\\
3.7995	0.61238\\
3.8178	0.58342\\
3.8361	0.55553\\
3.8544	0.5287\\
3.8728	0.5029\\
3.8911	0.47811\\
3.9094	0.45431\\
3.9277	0.43147\\
3.946	0.40957\\
3.9643	0.38859\\
3.9827	0.3685\\
4.001	0.34928\\
4.0193	0.33089\\
4.0376	0.31332\\
4.0559	0.29653\\
4.0743	0.28051\\
4.0926	0.26523\\
4.1109	0.25065\\
4.1292	0.23677\\
4.1475	0.22354\\
4.1658	0.21096\\
4.1842	0.19898\\
4.2025	0.1876\\
4.2208	0.17678\\
4.2391	0.16651\\
4.2574	0.15676\\
4.2758	0.1475\\
4.2941	0.13873\\
4.3124	0.13042\\
4.3307	0.12254\\
4.349	0.11509\\
4.3673	0.10803\\
4.3857	0.10136\\
4.404	0.095051\\
4.4223	0.089091\\
4.4406	0.083463\\
4.4589	0.078152\\
4.4773	0.073141\\
4.4956	0.068417\\
4.5139	0.063964\\
4.5322	0.059771\\
4.5505	0.055824\\
4.5688	0.052109\\
4.5872	0.048617\\
4.6055	0.045334\\
4.6238	0.04225\\
4.6421	0.039355\\
4.6604	0.036639\\
4.6788	0.034091\\
4.6971	0.031704\\
4.7154	0.029467\\
4.7337	0.027373\\
4.752	0.025414\\
4.7703	0.023581\\
4.7887	0.021869\\
4.807	0.02027\\
4.8253	0.018777\\
4.8436	0.017384\\
4.8619	0.016085\\
4.8803	0.014875\\
4.8986	0.013748\\
4.9169	0.012699\\
4.9352	0.011724\\
4.9535	0.010817\\
4.9719	0.0099751\\
4.9902	0.0091934\\
5.0085	0.0084681\\
5.0268	0.0077958\\
5.0451	0.0071728\\
5.0634	0.006596\\
5.0818	0.0060624\\
5.1001	0.0055689\\
5.1184	0.0051129\\
5.1367	0.0046919\\
5.155	0.0043033\\
5.1734	0.0039451\\
5.1917	0.0036149\\
5.21	0.0033108\\
5.2283	0.003031\\
5.2466	0.0027736\\
5.2649	0.0025371\\
5.2833	0.0023198\\
5.3016	0.0021204\\
5.3199	0.0019375\\
5.3382	0.0017698\\
5.3565	0.0016161\\
5.3749	0.0014754\\
5.3932	0.0013466\\
5.4115	0.0012288\\
};
\addlegendentry{\ExactBLNFitLegend}


\addplot[area legend, draw=none, fill=mycolor1, fill opacity=0.15, postaction={pattern={Lines[angle=45,distance={5pt},line width={0.5pt}]}, pattern color=mycolor1!60}, forget plot]
table[row sep=crcr] {%
x	y\\
3.598	1.0148\\
3.6163	0.97287\\
3.6346	0.93216\\
3.6529	0.89263\\
3.6713	0.85429\\
3.6896	0.81713\\
3.7079	0.78114\\
3.7262	0.74632\\
3.7445	0.71265\\
3.7628	0.68013\\
3.7812	0.64873\\
3.7995	0.61845\\
3.8178	0.58927\\
3.8361	0.56116\\
3.8544	0.53411\\
3.8728	0.5081\\
3.8911	0.48311\\
3.9094	0.4591\\
3.9277	0.43607\\
3.946	0.41399\\
3.9643	0.39282\\
3.9827	0.37256\\
4.001	0.35316\\
4.0193	0.33462\\
4.0376	0.31689\\
4.0559	0.29996\\
4.0743	0.28381\\
4.0926	0.26839\\
4.1109	0.2537\\
4.1292	0.2397\\
4.1475	0.22637\\
4.1658	0.21368\\
4.1842	0.20162\\
4.2025	0.19015\\
4.2208	0.17925\\
4.2391	0.16891\\
4.2574	0.15908\\
4.2758	0.14977\\
4.2941	0.14093\\
4.3124	0.13256\\
4.3307	0.12463\\
4.349	0.11712\\
4.3673	0.11002\\
4.3857	0.10329\\
4.404	0.096938\\
4.4223	0.090931\\
4.4406	0.085258\\
4.4589	0.079901\\
4.4773	0.074846\\
4.4956	0.070079\\
4.5139	0.065584\\
4.5322	0.061348\\
4.5505	0.057359\\
4.5688	0.053604\\
4.5872	0.050071\\
4.6055	0.046749\\
4.6238	0.043626\\
4.6421	0.040693\\
4.6604	0.037938\\
4.6788	0.035353\\
4.6971	0.032928\\
4.7154	0.030655\\
4.7337	0.028525\\
4.752	0.02653\\
4.7703	0.024662\\
4.7887	0.022914\\
4.807	0.02128\\
4.8253	0.019753\\
4.8436	0.018326\\
4.8619	0.016994\\
4.8803	0.01575\\
4.8986	0.01459\\
4.9169	0.013509\\
4.9352	0.012502\\
4.9535	0.011564\\
4.9719	0.010691\\
4.9902	0.0098785\\
5.0085	0.0091235\\
5.0268	0.0084222\\
5.0451	0.0077709\\
5.0634	0.0071667\\
5.0818	0.0066063\\
5.1001	0.0060871\\
5.1184	0.0056061\\
5.1367	0.005161\\
5.155	0.0047493\\
5.1734	0.0043688\\
5.1917	0.0040173\\
5.21	0.0036928\\
5.2283	0.0033935\\
5.2466	0.0031175\\
5.2649	0.0028632\\
5.2833	0.002629\\
5.3016	0.0024135\\
5.3199	0.0022152\\
5.3382	0.002033\\
5.3565	0.0018655\\
5.3749	0.0017116\\
5.3932	0.0015704\\
5.4115	0.0014407\\
5.4115	0.0010169\\
5.3932	0.0011228\\
5.3749	0.0012391\\
5.3565	0.0013667\\
5.3382	0.0015065\\
5.3199	0.0016597\\
5.3016	0.0018273\\
5.2833	0.0020106\\
5.2649	0.002211\\
5.2466	0.0024298\\
5.2283	0.0026685\\
5.21	0.0029288\\
5.1917	0.0032125\\
5.1734	0.0035213\\
5.155	0.0038574\\
5.1367	0.0042227\\
5.1184	0.0046197\\
5.1001	0.0050507\\
5.0818	0.0055184\\
5.0634	0.0060254\\
5.0451	0.0065747\\
5.0268	0.0071694\\
5.0085	0.0078127\\
4.9902	0.0085083\\
4.9719	0.0092597\\
4.9535	0.010071\\
4.9352	0.010946\\
4.9169	0.011889\\
4.8986	0.012906\\
4.8803	0.014\\
4.8619	0.015176\\
4.8436	0.016441\\
4.8253	0.0178\\
4.807	0.019259\\
4.7887	0.020824\\
4.7703	0.022501\\
4.752	0.024298\\
4.7337	0.026221\\
4.7154	0.028279\\
4.6971	0.030479\\
4.6788	0.03283\\
4.6604	0.03534\\
4.6421	0.038018\\
4.6238	0.040875\\
4.6055	0.043919\\
4.5872	0.047162\\
4.5688	0.050615\\
4.5505	0.054288\\
4.5322	0.058194\\
4.5139	0.062345\\
4.4956	0.066755\\
4.4773	0.071436\\
4.4589	0.076402\\
4.4406	0.081669\\
4.4223	0.087251\\
4.404	0.093163\\
4.3857	0.099423\\
4.3673	0.10605\\
4.349	0.11305\\
4.3307	0.12045\\
4.3124	0.12827\\
4.2941	0.13653\\
4.2758	0.14524\\
4.2574	0.15443\\
4.2391	0.16411\\
4.2208	0.17431\\
4.2025	0.18505\\
4.1842	0.19635\\
4.1658	0.20823\\
4.1475	0.22072\\
4.1292	0.23384\\
4.1109	0.24761\\
4.0926	0.26206\\
4.0743	0.27722\\
4.0559	0.2931\\
4.0376	0.30974\\
4.0193	0.32716\\
4.001	0.34539\\
3.9827	0.36445\\
3.9643	0.38436\\
3.946	0.40516\\
3.9277	0.42687\\
3.9094	0.44951\\
3.8911	0.47311\\
3.8728	0.4977\\
3.8544	0.52329\\
3.8361	0.5499\\
3.8178	0.57757\\
3.7995	0.60631\\
3.7812	0.63613\\
3.7628	0.66707\\
3.7445	0.69912\\
3.7262	0.73231\\
3.7079	0.76665\\
3.6896	0.80215\\
3.6713	0.83881\\
3.6529	0.87665\\
3.6346	0.91567\\
3.6163	0.95587\\
3.598	0.99725\\
}--cycle;
      % \input{../../../../.tex/20/KS/SizeDistributionFittingsfig41.tex}
      \legend{}
  
      \nextgroupplot[%
        xmin=3.6008,xmax=5.4115,
        % xmin=3.5,xmax=5.5,
        % ymin=0.001,ymax=1.0041,
      ] % This file was created by matlab2tikz.
%
\definecolor{mycolor1}{rgb}{0.83529,0.36863,0.00000}%
\definecolor{mycolor2}{rgb}{0.00000,0.44700,0.74100}%
%
\addplot[only marks, mark=*, mark size=1.1180pt, color=mycolor1, fill=mycolor1, opacity=0.60, draw=none, mark options={draw=none,line width=0pt}] table[row sep=crcr]{%
x	y\\
3.6	0.99468\\
3.6045	0.98404\\
3.6089	0.9734\\
3.6141	0.96277\\
3.6178	0.95213\\
3.6225	0.94149\\
3.6271	0.93085\\
3.6318	0.92021\\
3.6365	0.90957\\
3.6418	0.89894\\
3.6472	0.8883\\
3.6515	0.87766\\
3.657	0.86702\\
3.6615	0.85638\\
3.6662	0.84574\\
3.6709	0.83511\\
3.6761	0.82447\\
3.6804	0.81383\\
3.6852	0.80319\\
3.6904	0.79255\\
3.6957	0.78191\\
3.701	0.77128\\
3.7064	0.76064\\
3.7121	0.75\\
3.7185	0.73936\\
3.7229	0.72872\\
3.7274	0.71809\\
3.7333	0.70745\\
3.7389	0.69681\\
3.745	0.68617\\
3.7505	0.67553\\
3.7574	0.66489\\
3.7644	0.65426\\
3.7707	0.64362\\
3.7767	0.63298\\
3.7836	0.62234\\
3.7905	0.6117\\
3.7976	0.60106\\
3.8048	0.59043\\
3.812	0.57979\\
3.8198	0.56915\\
3.8269	0.55851\\
3.8352	0.54787\\
3.8439	0.53723\\
3.8506	0.5266\\
3.8568	0.51596\\
3.8635	0.50532\\
3.8723	0.49468\\
3.8796	0.48404\\
3.8878	0.4734\\
3.897	0.46277\\
3.9047	0.45213\\
3.9134	0.44149\\
3.9222	0.43085\\
3.9309	0.42021\\
3.9403	0.40957\\
3.9517	0.39894\\
3.9628	0.3883\\
3.9726	0.37766\\
3.9842	0.36702\\
3.9942	0.35638\\
4.0054	0.34574\\
4.0156	0.33511\\
4.0277	0.32447\\
4.0395	0.31383\\
4.0516	0.30319\\
4.066	0.29255\\
4.0769	0.28191\\
4.0914	0.27128\\
4.1049	0.26064\\
4.1181	0.25\\
4.1339	0.23936\\
4.1473	0.22872\\
4.1627	0.21809\\
4.1783	0.20745\\
4.1955	0.19681\\
4.2126	0.18617\\
4.231	0.17553\\
4.254	0.16489\\
4.2752	0.15426\\
4.2969	0.14362\\
4.3184	0.13298\\
4.3399	0.12234\\
4.364	0.1117\\
4.39	0.10106\\
4.4189	0.090426\\
4.4573	0.079787\\
4.4912	0.069149\\
4.5391	0.058511\\
4.5879	0.047872\\
4.6566	0.037234\\
4.7343	0.026596\\
4.8688	0.015957\\
5.1869	0.0053191\\
};
\addlegendentry{\ApproxEmpiricCCDFLegend}

\addplot [color=mycolor1, only marks, every error bar/.append style={opacity=0.30}, mark=*, mark size=0pt, draw=none, forget plot]
 plot [error bars/.cd, x dir=both, x explicit, error bar style={line width=1pt, color=mycolor1}, error mark options={mark=none,mark size=0pt}]
 table[row sep=crcr, x error plus index=2, x error minus index=3]{%
3.6	0.99468	0.0045554	0.0045554\\
3.6045	0.98404	0.0045703	0.0045703\\
3.6089	0.9734	0.0045491	0.0045491\\
3.6141	0.96277	0.0044812	0.0044812\\
3.6178	0.95213	0.0044817	0.0044817\\
3.6225	0.94149	0.0044624	0.0044624\\
3.6271	0.93085	0.0045762	0.0045762\\
3.6318	0.92021	0.0045625	0.0045625\\
3.6365	0.90957	0.004468	0.004468\\
3.6418	0.89894	0.0046379	0.0046379\\
3.6472	0.8883	0.0045394	0.0045394\\
3.6515	0.87766	0.0044728	0.0044728\\
3.657	0.86702	0.0045858	0.0045858\\
3.6615	0.85638	0.0045704	0.0045704\\
3.6662	0.84574	0.0044305	0.0044305\\
3.6709	0.83511	0.0043954	0.0043954\\
3.6761	0.82447	0.0045809	0.0045809\\
3.6804	0.81383	0.0045127	0.0045127\\
3.6852	0.80319	0.0045389	0.0045389\\
3.6904	0.79255	0.0046013	0.0046013\\
3.6957	0.78191	0.0046842	0.0046842\\
3.701	0.77128	0.0046214	0.0046214\\
3.7064	0.76064	0.0046093	0.0046093\\
3.7121	0.75	0.0045663	0.0045663\\
3.7185	0.73936	0.0045568	0.0045568\\
3.7229	0.72872	0.0046537	0.0046537\\
3.7274	0.71809	0.0046666	0.0046666\\
3.7333	0.70745	0.0046756	0.0046756\\
3.7389	0.69681	0.0047365	0.0047365\\
3.745	0.68617	0.0047282	0.0047282\\
3.7505	0.67553	0.0046776	0.0046776\\
3.7574	0.66489	0.0046232	0.0046232\\
3.7644	0.65426	0.0049389	0.0049389\\
3.7707	0.64362	0.0049657	0.0049657\\
3.7767	0.63298	0.0051864	0.0051864\\
3.7836	0.62234	0.0051426	0.0051426\\
3.7905	0.6117	0.0054154	0.0054154\\
3.7976	0.60106	0.0052834	0.0052834\\
3.8048	0.59043	0.0053245	0.0053245\\
3.812	0.57979	0.0053216	0.0053216\\
3.8198	0.56915	0.0053218	0.0053218\\
3.8269	0.55851	0.0052938	0.0052938\\
3.8352	0.54787	0.0052741	0.0052741\\
3.8439	0.53723	0.0055718	0.0055718\\
3.8506	0.5266	0.005749	0.005749\\
3.8568	0.51596	0.0058641	0.0058641\\
3.8635	0.50532	0.0058242	0.0058242\\
3.8723	0.49468	0.0060846	0.0060846\\
3.8796	0.48404	0.0058831	0.0058831\\
3.8878	0.4734	0.0060552	0.0060552\\
3.897	0.46277	0.0058784	0.0058784\\
3.9047	0.45213	0.0058689	0.0058689\\
3.9134	0.44149	0.0059804	0.0059804\\
3.9222	0.43085	0.0059113	0.0059113\\
3.9309	0.42021	0.0061887	0.0061887\\
3.9403	0.40957	0.0064169	0.0064169\\
3.9517	0.39894	0.0061866	0.0061866\\
3.9628	0.3883	0.0059948	0.0059948\\
3.9726	0.37766	0.006044	0.006044\\
3.9842	0.36702	0.0061542	0.0061542\\
3.9942	0.35638	0.0061402	0.0061402\\
4.0054	0.34574	0.006197	0.006197\\
4.0156	0.33511	0.0062917	0.0062917\\
4.0277	0.32447	0.0062543	0.0062543\\
4.0395	0.31383	0.006648	0.006648\\
4.0516	0.30319	0.0062708	0.0062708\\
4.066	0.29255	0.0064445	0.0064445\\
4.0769	0.28191	0.0065803	0.0065803\\
4.0914	0.27128	0.0068367	0.0068367\\
4.1049	0.26064	0.0067182	0.0067182\\
4.1181	0.25	0.0066727	0.0066727\\
4.1339	0.23936	0.006414	0.006414\\
4.1473	0.22872	0.0068788	0.0068788\\
4.1627	0.21809	0.0073314	0.0073314\\
4.1783	0.20745	0.0073394	0.0073394\\
4.1955	0.19681	0.0076429	0.0076429\\
4.2126	0.18617	0.007675	0.007675\\
4.231	0.17553	0.0079381	0.0079381\\
4.254	0.16489	0.0079159	0.0079159\\
4.2752	0.15426	0.0082002	0.0082002\\
4.2969	0.14362	0.0088204	0.0088204\\
4.3184	0.13298	0.00845	0.00845\\
4.3399	0.12234	0.0086796	0.0086796\\
4.364	0.1117	0.0086382	0.0086382\\
4.39	0.10106	0.0099918	0.0099918\\
4.4189	0.090426	0.010413	0.010413\\
4.4573	0.079787	0.011005	0.011005\\
4.4912	0.069149	0.011629	0.011629\\
4.5391	0.058511	0.011499	0.011499\\
4.5879	0.047872	0.012519	0.012519\\
4.6566	0.037234	0.015101	0.015101\\
4.7343	0.026596	0.016441	0.016441\\
4.8688	0.015957	0.021019	0.021019\\
5.1869	0.0053191	0.018612	0.018612\\
};
\addplot [color=mycolor1, mark repeat=9, mark phase=1, mark=+, mark options={solid, mycolor1}]
  table[row sep=crcr]{%
3.6008	0.97204\\
3.6191	0.94016\\
3.6374	0.89927\\
3.6557	0.85482\\
3.6739	0.81197\\
3.6922	0.77128\\
3.7105	0.73263\\
3.7288	0.69592\\
3.7471	0.66106\\
3.7654	0.62796\\
3.7837	0.59651\\
3.802	0.56665\\
3.8203	0.53828\\
3.8386	0.51135\\
3.8568	0.48576\\
3.8751	0.46146\\
3.8934	0.43838\\
3.9117	0.41645\\
3.93	0.39563\\
3.9483	0.37585\\
3.9666	0.35706\\
3.9849	0.33922\\
4.0032	0.32227\\
4.0215	0.30617\\
4.0397	0.29088\\
4.058	0.27635\\
4.0763	0.26255\\
4.0946	0.24944\\
4.1129	0.23699\\
4.1312	0.22517\\
4.1495	0.21393\\
4.1678	0.20326\\
4.1861	0.19312\\
4.2044	0.18349\\
4.2226	0.17434\\
4.2409	0.16565\\
4.2592	0.15739\\
4.2775	0.14955\\
4.2958	0.1421\\
4.3141	0.13502\\
4.3324	0.1283\\
4.3507	0.12191\\
4.369	0.11584\\
4.3873	0.11007\\
4.4055	0.10459\\
4.4238	0.099387\\
4.4421	0.094442\\
4.4604	0.089744\\
4.4787	0.08528\\
4.497	0.081039\\
4.5153	0.07701\\
4.5336	0.073182\\
4.5519	0.069544\\
4.5702	0.066088\\
4.5884	0.062805\\
4.6067	0.059685\\
4.625	0.05672\\
4.6433	0.053904\\
4.6616	0.051227\\
4.6799	0.048684\\
4.6982	0.046268\\
4.7165	0.043972\\
4.7348	0.04179\\
4.7531	0.039717\\
4.7713	0.037747\\
4.7896	0.035875\\
4.8079	0.034097\\
4.8262	0.032406\\
4.8445	0.0308\\
4.8628	0.029274\\
4.8811	0.027824\\
4.8994	0.026445\\
4.9177	0.025135\\
4.936	0.023891\\
4.9542	0.022708\\
4.9725	0.021584\\
4.9908	0.020515\\
5.0091	0.0195\\
5.0274	0.018535\\
5.0457	0.017618\\
5.064	0.016747\\
5.0823	0.015919\\
5.1006	0.015132\\
5.1189	0.014384\\
5.1371	0.013673\\
5.1554	0.012997\\
5.1737	0.012355\\
5.192	0.011745\\
5.2103	0.011164\\
5.2286	0.010613\\
5.2469	0.010089\\
5.2652	0.009591\\
5.2835	0.0091176\\
5.3018	0.0086677\\
5.32	0.00824\\
5.3383	0.0078335\\
5.3566	0.0074471\\
5.3749	0.0070798\\
5.3932	0.0067308\\
5.4115	0.0063989\\
};
\addlegendentry{\ApproxParetoFitLegend}


\addplot[area legend, draw=none, fill=mycolor1, fill opacity=0.15, postaction={pattern={Lines[angle=90,distance={5pt},line width={0.5pt}]}, pattern color=mycolor1!60}, forget plot]
table[row sep=crcr] {%
x	y\\
3.6008	0.97886\\
3.6191	0.94955\\
3.6374	0.90957\\
3.6557	0.86444\\
3.6739	0.82071\\
3.6922	0.77922\\
3.7105	0.73984\\
3.7288	0.70247\\
3.7471	0.667\\
3.7654	0.63335\\
3.7837	0.60141\\
3.802	0.57111\\
3.8203	0.54236\\
3.8386	0.51508\\
3.8568	0.4892\\
3.8751	0.46464\\
3.8934	0.44135\\
3.9117	0.41925\\
3.93	0.39828\\
3.9483	0.37838\\
3.9666	0.3595\\
3.9849	0.34158\\
4.0032	0.32458\\
4.0215	0.30843\\
4.0397	0.29311\\
4.058	0.27856\\
4.0763	0.26474\\
4.0946	0.25162\\
4.1129	0.23915\\
4.1312	0.22731\\
4.1495	0.21607\\
4.1678	0.20538\\
4.1861	0.19523\\
4.2044	0.18558\\
4.2226	0.17641\\
4.2409	0.1677\\
4.2592	0.15942\\
4.2775	0.15156\\
4.2958	0.14408\\
4.3141	0.13698\\
4.3324	0.13022\\
4.3507	0.1238\\
4.369	0.1177\\
4.3873	0.1119\\
4.4055	0.10639\\
4.4238	0.10115\\
4.4421	0.096173\\
4.4604	0.091439\\
4.4787	0.08694\\
4.497	0.082662\\
4.5153	0.078596\\
4.5336	0.074731\\
4.5519	0.071056\\
4.5702	0.067563\\
4.5884	0.064242\\
4.6067	0.061085\\
4.625	0.058084\\
4.6433	0.055231\\
4.6616	0.052518\\
4.6799	0.049939\\
4.6982	0.047488\\
4.7165	0.045157\\
4.7348	0.042941\\
4.7531	0.040834\\
4.7713	0.03883\\
4.7896	0.036926\\
4.8079	0.035115\\
4.8262	0.033393\\
4.8445	0.031756\\
4.8628	0.030199\\
4.8811	0.02872\\
4.8994	0.027312\\
4.9177	0.025974\\
4.936	0.024702\\
4.9542	0.023492\\
4.9725	0.022342\\
4.9908	0.021248\\
5.0091	0.020208\\
5.0274	0.019219\\
5.0457	0.018279\\
5.064	0.017385\\
5.0823	0.016534\\
5.1006	0.015726\\
5.1189	0.014957\\
5.1371	0.014226\\
5.1554	0.01353\\
5.1737	0.012869\\
5.192	0.01224\\
5.2103	0.011642\\
5.2286	0.011074\\
5.2469	0.010533\\
5.2652	0.010019\\
5.2835	0.0095295\\
5.3018	0.0090644\\
5.32	0.008622\\
5.3383	0.0082013\\
5.3566	0.0078012\\
5.3749	0.0074206\\
5.3932	0.0070587\\
5.4115	0.0067145\\
5.4115	0.0060833\\
5.3932	0.0064028\\
5.3749	0.006739\\
5.3566	0.007093\\
5.3383	0.0074657\\
5.32	0.007858\\
5.3018	0.008271\\
5.2835	0.0087057\\
5.2652	0.0091634\\
5.2469	0.0096453\\
5.2286	0.010153\\
5.2103	0.010687\\
5.192	0.011249\\
5.1737	0.011841\\
5.1554	0.012464\\
5.1371	0.01312\\
5.1189	0.01381\\
5.1006	0.014538\\
5.0823	0.015303\\
5.064	0.016109\\
5.0457	0.016958\\
5.0274	0.017851\\
5.0091	0.018792\\
4.9908	0.019782\\
4.9725	0.020825\\
4.9542	0.021923\\
4.936	0.023079\\
4.9177	0.024296\\
4.8994	0.025578\\
4.8811	0.026928\\
4.8628	0.028348\\
4.8445	0.029845\\
4.8262	0.03142\\
4.8079	0.033079\\
4.7896	0.034825\\
4.7713	0.036664\\
4.7531	0.038601\\
4.7348	0.04064\\
4.7165	0.042787\\
4.6982	0.045048\\
4.6799	0.047429\\
4.6616	0.049937\\
4.6433	0.052577\\
4.625	0.055357\\
4.6067	0.058284\\
4.5884	0.061367\\
4.5702	0.064614\\
4.5519	0.068032\\
4.5336	0.071633\\
4.5153	0.075424\\
4.497	0.079416\\
4.4787	0.08362\\
4.4604	0.088048\\
4.4421	0.09271\\
4.4238	0.09762\\
4.4055	0.10279\\
4.3873	0.10824\\
4.369	0.11397\\
4.3507	0.12001\\
4.3324	0.12637\\
4.3141	0.13306\\
4.2958	0.14012\\
4.2775	0.14754\\
4.2592	0.15536\\
4.2409	0.1636\\
4.2226	0.17227\\
4.2044	0.1814\\
4.1861	0.19102\\
4.1678	0.20114\\
4.1495	0.2118\\
4.1312	0.22302\\
4.1129	0.23483\\
4.0946	0.24727\\
4.0763	0.26036\\
4.058	0.27414\\
4.0397	0.28864\\
4.0215	0.30391\\
4.0032	0.31996\\
3.9849	0.33686\\
3.9666	0.35463\\
3.9483	0.37332\\
3.93	0.39298\\
3.9117	0.41366\\
3.8934	0.4354\\
3.8751	0.45827\\
3.8568	0.48232\\
3.8386	0.50761\\
3.8203	0.53421\\
3.802	0.56218\\
3.7837	0.59161\\
3.7654	0.62256\\
3.7471	0.65513\\
3.7288	0.68938\\
3.7105	0.72542\\
3.6922	0.76333\\
3.6739	0.80322\\
3.6557	0.8452\\
3.6374	0.88898\\
3.6191	0.93077\\
3.6008	0.96523\\
}--cycle;
\addplot [color=mycolor1, mark repeat=9, mark phase=1, mark=x, mark options={solid, mycolor1}]
  table[row sep=crcr]{%
3.6008	0.9962\\
3.6191	0.95443\\
3.6374	0.91389\\
3.6557	0.87458\\
3.6739	0.8365\\
3.6922	0.79965\\
3.7105	0.76401\\
3.7288	0.72957\\
3.7471	0.69633\\
3.7654	0.66427\\
3.7837	0.63337\\
3.802	0.6036\\
3.8203	0.57496\\
3.8386	0.54741\\
3.8568	0.52093\\
3.8751	0.4955\\
3.8934	0.4711\\
3.9117	0.44769\\
3.93	0.42525\\
3.9483	0.40375\\
3.9666	0.38316\\
3.9849	0.36347\\
4.0032	0.34463\\
4.0215	0.32663\\
4.0397	0.30944\\
4.058	0.29302\\
4.0763	0.27735\\
4.0946	0.26241\\
4.1129	0.24816\\
4.1312	0.23459\\
4.1495	0.22166\\
4.1678	0.20936\\
4.1861	0.19765\\
4.2044	0.18652\\
4.2226	0.17594\\
4.2409	0.16588\\
4.2592	0.15633\\
4.2775	0.14727\\
4.2958	0.13867\\
4.3141	0.13051\\
4.3324	0.12277\\
4.3507	0.11544\\
4.369	0.1085\\
4.3873	0.10192\\
4.4055	0.095702\\
4.4238	0.089815\\
4.4421	0.084248\\
4.4604	0.078986\\
4.4787	0.074015\\
4.497	0.06932\\
4.5153	0.064888\\
4.5336	0.060706\\
4.5519	0.056763\\
4.5702	0.053046\\
4.5884	0.049544\\
4.6067	0.046247\\
4.625	0.043144\\
4.6433	0.040226\\
4.6616	0.037482\\
4.6799	0.034905\\
4.6982	0.032485\\
4.7165	0.030214\\
4.7348	0.028085\\
4.7531	0.026091\\
4.7713	0.024222\\
4.7896	0.022474\\
4.8079	0.02084\\
4.8262	0.019312\\
4.8445	0.017886\\
4.8628	0.016555\\
4.8811	0.015314\\
4.8994	0.014158\\
4.9177	0.013081\\
4.936	0.012079\\
4.9542	0.011148\\
4.9725	0.010283\\
4.9908	0.0094796\\
5.0091	0.0087343\\
5.0274	0.0080432\\
5.0457	0.007403\\
5.064	0.0068101\\
5.0823	0.0062616\\
5.1006	0.0057543\\
5.1189	0.0052856\\
5.1371	0.0048527\\
5.1554	0.0044531\\
5.1737	0.0040846\\
5.192	0.0037449\\
5.2103	0.0034319\\
5.2286	0.0031437\\
5.2469	0.0028785\\
5.2652	0.0026346\\
5.2835	0.0024105\\
5.3018	0.0022045\\
5.32	0.0020155\\
5.3383	0.0018419\\
5.3566	0.0016828\\
5.3749	0.0015369\\
5.3932	0.0014032\\
5.4115	0.0012807\\
};
\addlegendentry{\ApproxBLNFitLegend}


\addplot[area legend, draw=none, fill=mycolor1, fill opacity=0.15, postaction={pattern={Lines[angle=45,distance={5pt},line width={0.5pt}]}, pattern color=mycolor1!60}, forget plot]
table[row sep=crcr] {%
x	y\\
3.6008	1.0041\\
3.6191	0.9621\\
3.6374	0.92133\\
3.6557	0.8818\\
3.6739	0.84349\\
3.6922	0.80642\\
3.7105	0.77056\\
3.7288	0.73591\\
3.7471	0.70246\\
3.7654	0.6702\\
3.7837	0.63909\\
3.802	0.60914\\
3.8203	0.5803\\
3.8386	0.55257\\
3.8568	0.52592\\
3.8751	0.50033\\
3.8934	0.47576\\
3.9117	0.45219\\
3.93	0.4296\\
3.9483	0.40796\\
3.9666	0.38724\\
3.9849	0.36741\\
4.0032	0.34845\\
4.0215	0.33032\\
4.0397	0.31301\\
4.058	0.29647\\
4.0763	0.2807\\
4.0946	0.26565\\
4.1129	0.2513\\
4.1312	0.23763\\
4.1495	0.2246\\
4.1678	0.21221\\
4.1861	0.20041\\
4.2044	0.18919\\
4.2226	0.17853\\
4.2409	0.16839\\
4.2592	0.15876\\
4.2775	0.14962\\
4.2958	0.14095\\
4.3141	0.13272\\
4.3324	0.12492\\
4.3507	0.11752\\
4.369	0.11052\\
4.3873	0.10388\\
4.4055	0.097599\\
4.4238	0.091655\\
4.4421	0.086032\\
4.4604	0.080715\\
4.4787	0.075689\\
4.497	0.070941\\
4.5153	0.066457\\
4.5336	0.062225\\
4.5519	0.058231\\
4.5702	0.054465\\
4.5884	0.050915\\
4.6067	0.047571\\
4.625	0.044422\\
4.6433	0.041458\\
4.6616	0.038671\\
4.6799	0.036052\\
4.6982	0.033591\\
4.7165	0.031281\\
4.7348	0.029115\\
4.7531	0.027084\\
4.7713	0.025181\\
4.7896	0.023399\\
4.8079	0.021733\\
4.8262	0.020174\\
4.8445	0.018718\\
4.8628	0.017358\\
4.8811	0.016089\\
4.8994	0.014905\\
4.9177	0.013801\\
4.936	0.012773\\
4.9542	0.011815\\
4.9725	0.010924\\
4.9908	0.010095\\
5.0091	0.0093245\\
5.0274	0.0086085\\
5.0457	0.0079436\\
5.064	0.0073265\\
5.0823	0.0067542\\
5.1006	0.0062237\\
5.1189	0.0057323\\
5.1371	0.0052774\\
5.1554	0.0048565\\
5.1737	0.0044673\\
5.192	0.0041076\\
5.2103	0.0037754\\
5.2286	0.0034688\\
5.2469	0.003186\\
5.2652	0.0029252\\
5.2835	0.0026849\\
5.3018	0.0024636\\
5.32	0.0022598\\
5.3383	0.0020723\\
5.3566	0.0018999\\
5.3749	0.0017414\\
5.3932	0.0015957\\
5.4115	0.0014619\\
5.4115	0.0010995\\
5.3932	0.0012106\\
5.3749	0.0013324\\
5.3566	0.0014657\\
5.3383	0.0016115\\
5.32	0.0017711\\
5.3018	0.0019455\\
5.2835	0.0021361\\
5.2652	0.0023441\\
5.2469	0.0025711\\
5.2286	0.0028186\\
5.2103	0.0030883\\
5.192	0.0033821\\
5.1737	0.0037019\\
5.1554	0.0040498\\
5.1371	0.004428\\
5.1189	0.0048389\\
5.1006	0.005285\\
5.0823	0.005769\\
5.064	0.0062938\\
5.0457	0.0068624\\
5.0274	0.007478\\
5.0091	0.008144\\
4.9908	0.008864\\
4.9725	0.0096416\\
4.9542	0.010481\\
4.936	0.011386\\
4.9177	0.012361\\
4.8994	0.01341\\
4.8811	0.014539\\
4.8628	0.015752\\
4.8445	0.017054\\
4.8262	0.01845\\
4.8079	0.019947\\
4.7896	0.021549\\
4.7713	0.023264\\
4.7531	0.025098\\
4.7348	0.027056\\
4.7165	0.029148\\
4.6982	0.031379\\
4.6799	0.033758\\
4.6616	0.036293\\
4.6433	0.038993\\
4.625	0.041866\\
4.6067	0.044923\\
4.5884	0.048173\\
4.5702	0.051627\\
4.5519	0.055294\\
4.5336	0.059188\\
4.5153	0.063318\\
4.497	0.067698\\
4.4787	0.072341\\
4.4604	0.077258\\
4.4421	0.082465\\
4.4238	0.087976\\
4.4055	0.093805\\
4.3873	0.099969\\
4.369	0.10648\\
4.3507	0.11336\\
4.3324	0.12063\\
4.3141	0.1283\\
4.2958	0.13638\\
4.2775	0.14491\\
4.2592	0.1539\\
4.2409	0.16338\\
4.2226	0.17335\\
4.2044	0.18385\\
4.1861	0.1949\\
4.1678	0.20651\\
4.1495	0.21872\\
4.1312	0.23155\\
4.1129	0.24503\\
4.0946	0.25917\\
4.0763	0.274\\
4.058	0.28956\\
4.0397	0.30587\\
4.0215	0.32294\\
4.0032	0.34082\\
3.9849	0.35953\\
3.9666	0.37909\\
3.9483	0.39953\\
3.93	0.42089\\
3.9117	0.44318\\
3.8934	0.46644\\
3.8751	0.49068\\
3.8568	0.51594\\
3.8386	0.54224\\
3.8203	0.56961\\
3.802	0.59807\\
3.7837	0.62764\\
3.7654	0.65835\\
3.7471	0.69021\\
3.7288	0.72324\\
3.7105	0.75745\\
3.6922	0.79288\\
3.6739	0.82951\\
3.6557	0.86737\\
3.6374	0.90645\\
3.6191	0.94676\\
3.6008	0.98831\\
}--cycle;
      % \input{../../../../.tex/20/KS/SizeDistributionFittingsfig42.tex}
      \legend{}
    \end{groupplot}
  %\endgroup

  %\begingroup Parametri
    
%\begingroup Parametro
\providecommand\varyingParameterSymbol\lambda
\providecommand\varyingParameterList[1]{
  \ifcase#1\relax%
  \or\num{1e-2}% Caso 1
  \or\num{8.2e-2}% Caso 2
  \or\num{1.55e-1}% Caso 3
  \or\num{2.27e-1}% Caso 4
  \else\num{3e-1}% Caso 5
  \fi
}
%\endgroup

%\begingroup beta
\providecommand\betaExactList[1]{
  \ifcase#1\relax%
  \or\num{1.267+-0.007}% Caso 1
  \or\num{1.285+-0.010}% Caso 2
  \or\num{1.264+-0.015}% Caso 3
  \or\num{1.246+-0.015}% Caso 4
  \else\num{1.210+-0.015}% Caso 5
  \fi
}
\providecommand\betaAprxList[1]{
  \ifcase#1\relax%
  \or\num{1.264+-0.006}% Caso 1
  \or\num{1.270+-0.011}% Caso 2
  \or\num{1.280+-0.013}% Caso 3
  \or\num{1.250+-0.015}% Caso 4
  \else\num{1.218+-0.014}% Caso 5
  \fi
}
%\endgroup

%\begingroup xi
\providecommand\zetaExactList[1]{
  \ifcase#1\relax%
  \or\num{0.757+-0.004}% Caso 1
  \or\num{0.710+-0.026}% Caso 2
  \or\num{0.670+-0.035}% Caso 3
  \or\num{0.625+-0.052}% Caso 4
  \else\num{0.539+-0.059}% Caso 5
  \fi
}
\providecommand\zetaAprxList[1]{
  \ifcase#1\relax%
  \or\num{0.749+-0.004}% Caso 1
  \or\num{0.721+-0.022}% Caso 2
  \or\num{0.645+-0.040}% Caso 3
  \or\num{0.622+-0.051}% Caso 4
  \else\num{0.567+-0.053}% Caso 5
  \fi
}
%\endgroup

%\begingroup mI
\providecommand\mIExactList[1]{
  \ifcase#1\relax%
  \or\num{3.265+-0.001}% Caso 1
  \or\num{3.238+-0.005}% Caso 2
  \or\num{3.220+-0.008}% Caso 3
  \or\num{3.213+-0.013}% Caso 4
  \else\num{3.174+-0.023}% Caso 5
  \fi
}
\providecommand\mIAprxList[1]{
  \ifcase#1\relax%
  \or\num{3.262+-0.001}% Caso 1
  \or\num{3.240+-0.005}% Caso 2
  \or\num{3.217+-0.010}% Caso 3
  \or\num{3.217+-0.012}% Caso 4
  \else\num{3.208+-0.009}% Caso 5
  \fi
}
%\endgroup

%\begingroup sI
\providecommand\sIExactList[1]{
  \ifcase#1\relax%
  \or\num{0.284+-0.001}% Caso 1
  \or\num{0.311+-0.008}% Caso 2
  \or\num{0.343+-0.010}% Caso 3
  \or\num{0.361+-0.020}% Caso 4
  \else\num{0.398+-0.029}% Caso 5
  \fi
}
\providecommand\sIAprxList[1]{
  \ifcase#1\relax%
  \or\num{0.281+-0.001}% Caso 1
  \or\num{0.315+-0.006}% Caso 2
  \or\num{0.334+-0.013}% Caso 3
  \or\num{0.366+-0.018}% Caso 4
  \else\num{0.411+-0.025}% Caso 5
  \fi
}
%\endgroup

%\begingroup mII
\providecommand\mIIExactList[1]{
  \ifcase#1\relax%
  \or\num{3.788+-0.005}% Caso 1
  \or\num{3.730+-0.029}% Caso 2
  \or\num{3.681+-0.052}% Caso 3
  \or\num{3.650+-0.072}% Caso 4
  \else\num{3.572+-0.095}% Caso 5
  \fi
}
\providecommand\mIIAprxList[1]{
  \ifcase#1\relax%
  \or\num{3.778+-0.006}% Caso 1
  \or\num{3.735+-0.029}% Caso 2
  \or\num{3.656+-0.048}% Caso 3
  \or\num{3.658+-0.084}% Caso 4
  \else\num{3.502+-0.078}% Caso 5
  \fi
}
%\endgroup

%\begingroup sII
\providecommand\sIIExactList[1]{
  \ifcase#1\relax%
  \or\num{0.455+-0.002}% Caso 1
  \or\num{0.486+-0.007}% Caso 2
  \or\num{0.515+-0.014}% Caso 3
  \or\num{0.533+-0.021}% Caso 4
  \else\num{0.570+-0.028}% Caso 5
  \fi
}
\providecommand\sIIAprxList[1]{
  \ifcase#1\relax%
  \or\num{0.457+-0.002}% Caso 1
  \or\num{0.490+-0.007}% Caso 2
  \or\num{0.521+-0.012}% Caso 3
  \or\num{0.545+-0.021}% Caso 4
  \else\num{0.586+-0.029}% Caso 5
  \fi
}
%\endgroup


    \foreach \r/\a/\p/\xS/\yS/\yLS in {%
      1/south west/south west/0em/0em/-.5em,
      2/south west/south west/0em/0em/-.5em,
      3/south west/south west/0em/0em/-.5em,
      4/south west/south west/0em/0em/-.5em,
      5/south west/south west/0em/0em/-.5em
    }{%
      \def\zetaExact{\zetaExactList\r}
      \def\zetaAprx{\zetaAprxList\r}

      \def\mIExact{\mIExactList\r}
      \def\mIAprx{\mIAprxList\r}

      \def\sIExact{\sIExactList\r}
      \def\sIAprx{\sIAprxList\r}

      \def\mIIExact{\mIIExactList\r}
      \def\mIIAprx{\mIIAprxList\r}

      \def\sIIExact{\sIIExactList\r}
      \def\sIIAprx{\sIIAprxList\r}

      \foreach \c/\t in {1/2,2/0,3/1}{%
        \writeBiLognormalParameters[\a]\t{Row\r\space c\c r1.\p}\xS\yS\yLS
      } % «\fl=first line», «\sl=second line», «\a=anchor»
       % «\xS=\xShift», «\yS=\yShift» e «\yLS=\yLineShift»
    }

    \ifpageisodd\writeVaryingParameter{1}{west}{west}{0em}{0em}
    \else\writeVaryingParameter{3}{east}{east}{0em}{0em}\fi
  %\endgroup

  %\begingroup Didascalie
    % \writeSubCaptions{south west}{south west}{0em}{0em}
    \ifpageisodd\writeSubCaptions{east}{east}{0em}{0em}
    \else\writeSubCaptions{west}{west}{0em}{0em}\fi
    % https://tex.stackexchange.com/questions/632045/get-subcaption-for-every-plot-using-groupplot-pgf-tikz-after-2020-update
  %\endgroup

  \redefineTikZbounds{Row1 c1r1}{Row4 c3r1}
\end{tikzpicture}%

% \input{Parametri.txt}
% 
%\begingroup Parametro
\providecommand\varyingParameterSymbol\lambda
\providecommand\varyingParameterList[1]{
  \ifcase#1\relax%
  \or\num{1e-2}% Caso 1
  \or\num{8.2e-2}% Caso 2
  \or\num{1.55e-1}% Caso 3
  \or\num{2.27e-1}% Caso 4
  \else\num{3e-1}% Caso 5
  \fi
}
%\endgroup

%\begingroup beta
\providecommand\betaExactList[1]{
  \ifcase#1\relax%
  \or\num{1.267+-0.007}% Caso 1
  \or\num{1.285+-0.010}% Caso 2
  \or\num{1.264+-0.015}% Caso 3
  \or\num{1.246+-0.015}% Caso 4
  \else\num{1.210+-0.015}% Caso 5
  \fi
}
\providecommand\betaAprxList[1]{
  \ifcase#1\relax%
  \or\num{1.264+-0.006}% Caso 1
  \or\num{1.270+-0.011}% Caso 2
  \or\num{1.280+-0.013}% Caso 3
  \or\num{1.250+-0.015}% Caso 4
  \else\num{1.218+-0.014}% Caso 5
  \fi
}
%\endgroup

%\begingroup xi
\providecommand\zetaExactList[1]{
  \ifcase#1\relax%
  \or\num{0.757+-0.004}% Caso 1
  \or\num{0.710+-0.026}% Caso 2
  \or\num{0.670+-0.035}% Caso 3
  \or\num{0.625+-0.052}% Caso 4
  \else\num{0.539+-0.059}% Caso 5
  \fi
}
\providecommand\zetaAprxList[1]{
  \ifcase#1\relax%
  \or\num{0.749+-0.004}% Caso 1
  \or\num{0.721+-0.022}% Caso 2
  \or\num{0.645+-0.040}% Caso 3
  \or\num{0.622+-0.051}% Caso 4
  \else\num{0.567+-0.053}% Caso 5
  \fi
}
%\endgroup

%\begingroup mI
\providecommand\mIExactList[1]{
  \ifcase#1\relax%
  \or\num{3.265+-0.001}% Caso 1
  \or\num{3.238+-0.005}% Caso 2
  \or\num{3.220+-0.008}% Caso 3
  \or\num{3.213+-0.013}% Caso 4
  \else\num{3.174+-0.023}% Caso 5
  \fi
}
\providecommand\mIAprxList[1]{
  \ifcase#1\relax%
  \or\num{3.262+-0.001}% Caso 1
  \or\num{3.240+-0.005}% Caso 2
  \or\num{3.217+-0.010}% Caso 3
  \or\num{3.217+-0.012}% Caso 4
  \else\num{3.208+-0.009}% Caso 5
  \fi
}
%\endgroup

%\begingroup sI
\providecommand\sIExactList[1]{
  \ifcase#1\relax%
  \or\num{0.284+-0.001}% Caso 1
  \or\num{0.311+-0.008}% Caso 2
  \or\num{0.343+-0.010}% Caso 3
  \or\num{0.361+-0.020}% Caso 4
  \else\num{0.398+-0.029}% Caso 5
  \fi
}
\providecommand\sIAprxList[1]{
  \ifcase#1\relax%
  \or\num{0.281+-0.001}% Caso 1
  \or\num{0.315+-0.006}% Caso 2
  \or\num{0.334+-0.013}% Caso 3
  \or\num{0.366+-0.018}% Caso 4
  \else\num{0.411+-0.025}% Caso 5
  \fi
}
%\endgroup

%\begingroup mII
\providecommand\mIIExactList[1]{
  \ifcase#1\relax%
  \or\num{3.788+-0.005}% Caso 1
  \or\num{3.730+-0.029}% Caso 2
  \or\num{3.681+-0.052}% Caso 3
  \or\num{3.650+-0.072}% Caso 4
  \else\num{3.572+-0.095}% Caso 5
  \fi
}
\providecommand\mIIAprxList[1]{
  \ifcase#1\relax%
  \or\num{3.778+-0.006}% Caso 1
  \or\num{3.735+-0.029}% Caso 2
  \or\num{3.656+-0.048}% Caso 3
  \or\num{3.658+-0.084}% Caso 4
  \else\num{3.502+-0.078}% Caso 5
  \fi
}
%\endgroup

%\begingroup sII
\providecommand\sIIExactList[1]{
  \ifcase#1\relax%
  \or\num{0.455+-0.002}% Caso 1
  \or\num{0.486+-0.007}% Caso 2
  \or\num{0.515+-0.014}% Caso 3
  \or\num{0.533+-0.021}% Caso 4
  \else\num{0.570+-0.028}% Caso 5
  \fi
}
\providecommand\sIIAprxList[1]{
  \ifcase#1\relax%
  \or\num{0.457+-0.002}% Caso 1
  \or\num{0.490+-0.007}% Caso 2
  \or\num{0.521+-0.012}% Caso 3
  \or\num{0.545+-0.021}% Caso 4
  \else\num{0.586+-0.029}% Caso 5
  \fi
}
%\endgroup


  \endgroup

  \begingroup
    \setcounter{subfigure}{0}
    
% !TEX root = ../../../../../Esperimenti/.tex/MEF.tex
% LTeX: language=it

\infopageKeysI

\begin{tikzpicture}%
  %\begingroup Righe
    \begin{groupplot}[
      group style={
        group name=Row1,
        plotGroup
      },
      plotRow,
      row4Keys,
      plotLegendNW,
      %
      xmin=10,xmax=283,
      ymin=100,ymax=2.0424e+05,
    ]
      \nextgroupplot[%
        % xmin=10,xmax=283,
        % ymin=100,ymax=1.4487e+05,
      ] % This file was created by matlab2tikz.
%
\definecolor{mycolor1}{rgb}{0.00000,0.44706,0.69804}%
\definecolor{mycolor2}{rgb}{0.00000,0.44700,0.74100}%
\definecolor{mycolor3}{rgb}{0.83529,0.36863,0.00000}%
\definecolor{mycolor4}{rgb}{0.85000,0.32500,0.09800}%
%
\addplot[only marks, mark=*, mark size=1.2748pt, color=mycolor1, fill=mycolor1, opacity=0.70, draw=none, mark options={draw=none,line width=0pt}] table[row sep=crcr]{%
x	y\\
34	2851.2\\
29	1488.4\\
96	16990\\
31	2007.1\\
27	1372.1\\
70	7626.2\\
24	1332.2\\
37	2351.4\\
42	3057.2\\
16	1215.1\\
29	1589.9\\
59	3535.1\\
69	3469.4\\
28	1682.6\\
11	360.9\\
20	1244\\
45	4011.6\\
27	1558.3\\
22	1094.6\\
27	1537.3\\
44	5285\\
16	1620\\
39	4631.6\\
19	1206\\
33	2235.4\\
37	2578.7\\
33	1433.6\\
18	623.45\\
29	2459.2\\
22	1102.8\\
33	1514.6\\
22	966.47\\
49	8453.1\\
25	1617.8\\
43	4050.1\\
26	2521.2\\
31	3217.6\\
26	1287.7\\
18	981.84\\
13	593.7\\
32	4139.3\\
38	2521\\
32	3592.8\\
29	1929.6\\
35	1658.7\\
35	3284.3\\
135	20146\\
27	4702.7\\
42	4421.1\\
39	4834.9\\
51	7571.5\\
86	10528\\
15	937.7\\
39	4873.2\\
42	3484.1\\
43	3373.7\\
58	5802.4\\
97	24886\\
50	2570.4\\
20	1602.4\\
28	1293.5\\
22	1268.1\\
47	3423.8\\
154	61793\\
30	1920\\
10	536.5\\
37	8161.1\\
31	1986.1\\
54	13976\\
76	11720\\
55	4925.9\\
26	2800\\
33	1722.9\\
34	2256.7\\
25	1771.9\\
31	3867.8\\
34	4703\\
26	2423.7\\
39	3770.4\\
27	2850.1\\
32	2139.1\\
26	2352.2\\
29	2367.4\\
23	3345.2\\
26	2487.8\\
14	914.37\\
33	2195.8\\
24	1162.2\\
20	1771.9\\
51	2172.1\\
35	3001.5\\
29	1353.1\\
46	1034.1\\
44	3716.9\\
45	3436.3\\
29	998.2\\
28	1793.6\\
77	3131.6\\
70	4049.4\\
70	3448.9\\
40	2517.7\\
90	5635.2\\
53	3698.2\\
62	1845.8\\
47	3592.4\\
34	1400.3\\
14	1125.1\\
47	2135.7\\
43	1176.2\\
20	842.46\\
28	1135.9\\
80	3228.8\\
25	595.71\\
36	1851.7\\
27	2283.8\\
82	4333.3\\
31	1132.7\\
39	1436.6\\
23	2238.8\\
27	2852.4\\
37	2272.9\\
129	8365.6\\
63	4436.6\\
64	2083.9\\
91	9555.9\\
27	1302.6\\
26	1509.8\\
17	902.04\\
33	1651.8\\
30	3363.6\\
33	1463.8\\
182	14018\\
26	1213.6\\
48	2986.2\\
46	1847.4\\
10	416.32\\
25	1175.1\\
26	835.94\\
206	24904\\
37	1338\\
43	1681\\
51	2035.7\\
57	6649.8\\
36	1806.8\\
38	1641.1\\
28	1368.6\\
19	1146.4\\
36	1679.5\\
45	3310.6\\
53	3621.7\\
48	3690.2\\
64	3215.8\\
65	3055.7\\
45	1347.5\\
54	2975.2\\
16	461.78\\
24	1108.5\\
127	8558.5\\
48	1856.5\\
64	2648.1\\
42	2749.7\\
42	1361.6\\
17	597.94\\
48	2948.6\\
42	2213.5\\
20	904.25\\
34	1259.9\\
47	1059.8\\
42	3001.8\\
50	2383.4\\
87	6514.1\\
64	3972.2\\
32	1852.1\\
39	1222.4\\
35	2584.8\\
18	1126\\
28	955.03\\
27	931.16\\
37	2131.9\\
30	2552.9\\
89	11576\\
22	1901.6\\
28	1701.4\\
25	1377.4\\
15	661.94\\
41	3504.7\\
32	1007.1\\
34	2052.9\\
16	828.35\\
88	5896.4\\
39	1033.1\\
140	26690\\
39	1350.9\\
49	1889.1\\
72	2477.6\\
37	1743.9\\
41	3578.9\\
283	1.438e+05\\
31	2923.4\\
99	18758\\
103	23535\\
39	3565.8\\
35	1571.9\\
98	13823\\
45	4384.9\\
73	8387.2\\
34	2340.2\\
56	7986.5\\
39	2947.2\\
37	3022.8\\
55	2215.8\\
25	763.38\\
29	1148.4\\
42	1240.7\\
35	3961\\
26	810.38\\
45	8358.7\\
56	6140.8\\
39	1615.2\\
62	3234.6\\
94	12172\\
112	23494\\
25	816.96\\
57	2517.7\\
50	2499.6\\
47	6505.8\\
96	7074.7\\
54	6265.4\\
28	2549.5\\
38	3812.9\\
47	3644.4\\
35	2112.1\\
45	1705.3\\
50	3764.2\\
43	1058.1\\
29	2181.7\\
42	1931.5\\
71	19926\\
68	6539.7\\
149	52722\\
67	5576.8\\
51	2915.1\\
34	1902.6\\
120	14465\\
43	7099.9\\
117	10679\\
39	1383.5\\
72	8408.7\\
40	3759.5\\
43	2443.1\\
32	2686.5\\
49	8503.8\\
41	3985\\
98	5104.3\\
96	12397\\
45	1739.9\\
129	33928\\
43	2178.4\\
93	8151.4\\
41	3386.4\\
88	9863.8\\
73	5712.5\\
108	15614\\
52	7527\\
18	497.78\\
36	1113.5\\
50	5373.2\\
37	2378.2\\
83	14111\\
40	2259.3\\
26	2719.2\\
40	1757.9\\
49	3820.3\\
27	1475.5\\
41	1434.6\\
36	933.28\\
52	5017.2\\
37	1160.8\\
68	7126.2\\
43	2855.2\\
132	13488\\
73	3056.7\\
46	4725.5\\
47	1278.6\\
44	1692.7\\
41	4987.9\\
30	1567.7\\
25	3087.5\\
24	1985.1\\
85	8156.8\\
44	3066.6\\
29	1868.2\\
34	1705.2\\
101	14451\\
18	1614.8\\
21	1428.4\\
111	16710\\
95	4784.4\\
23	1109.2\\
31	673.8\\
84	4064.8\\
23	846.48\\
55	5089.3\\
32	1184\\
27	957.19\\
48	1013.7\\
18	391.02\\
28	2036.3\\
40	1235.9\\
30	1337.6\\
13	353.76\\
33	1614.3\\
14	604.62\\
52	2477.5\\
65	7352.1\\
49	2634.2\\
35	1456.3\\
109	6953.3\\
26	985.04\\
41	1387.1\\
39	1294.3\\
61	6040.5\\
40	1520.1\\
36	2129.3\\
21	674.14\\
80	4224.6\\
36	1404\\
36	2016.1\\
38	1082.4\\
42	2146.8\\
27	783.46\\
46	2314.8\\
23	695.68\\
34	1797.3\\
200	37843\\
37	2495.7\\
16	668.81\\
50	2506.4\\
19	701.77\\
41	2959.9\\
25	866.67\\
39	2136.6\\
45	3052.8\\
83	7235.6\\
27	579.37\\
67	2635\\
46	3155.8\\
35	1676.1\\
44	2589.4\\
30	1928.3\\
29	1144.4\\
23	563.44\\
43	2045.7\\
28	811.18\\
24	1157.9\\
56	3097.1\\
33	994.96\\
13	439.83\\
41	2987\\
34	1022.6\\
14	499.21\\
83	8650.7\\
33	1702.1\\
35	1653.7\\
29	1052.6\\
54	3741.2\\
41	1343.3\\
19	753.42\\
34	2125.6\\
19	728.55\\
38	2246.1\\
43	1807.7\\
14	705.51\\
17	830.56\\
12	492.07\\
};
\addlegendentry{\ExactScatterLegend}

\addplot[only marks, mark=*, mark size=1.2748pt, color=mycolor3, fill=mycolor3, opacity=0.70, draw=none, mark options={draw=none,line width=0pt}] table[row sep=crcr]{%
x	y\\
34	2708.6\\
29	1423.3\\
96	16788\\
31	1991.2\\
27	1363.5\\
70	7601.1\\
24	1323.7\\
37	2304.2\\
42	2929.6\\
16	1214.3\\
29	1534.3\\
59	3573.2\\
69	3477.3\\
28	1623.5\\
11	345.68\\
20	1220.4\\
45	3757.7\\
27	1531.9\\
22	1092.3\\
27	1484.3\\
44	5224.2\\
16	1590.3\\
39	4459.9\\
19	1180.8\\
33	2224.1\\
37	2546.2\\
33	1438.7\\
18	610.78\\
29	2463.3\\
22	1069\\
33	1515.7\\
22	922.49\\
49	8173.5\\
25	1593.4\\
43	3932.6\\
26	2442.1\\
31	3201.4\\
26	1292\\
18	950.25\\
13	577.17\\
32	4091.3\\
38	2467.3\\
32	3493.5\\
29	1855.7\\
35	1672.7\\
35	3190\\
135	19679\\
27	4549.1\\
42	4314.1\\
39	4692.7\\
51	7190\\
86	10229\\
15	951.13\\
39	4659.8\\
42	3511.9\\
43	3263.5\\
58	5657.8\\
97	24092\\
50	2607.2\\
20	1557\\
28	1270.3\\
22	1215.5\\
47	3344.3\\
154	60852\\
30	1852.4\\
10	515.06\\
37	8050.8\\
31	1938.9\\
54	13650\\
76	11498\\
55	4712.9\\
26	2729.4\\
33	1672\\
34	2166.4\\
25	1707.7\\
31	3836.1\\
34	4507.7\\
26	2365.7\\
39	3679.1\\
27	2806.9\\
32	2117.6\\
26	2289.8\\
29	2351\\
23	3308.7\\
26	2417.2\\
14	873.66\\
33	2127.1\\
24	1133.5\\
20	1719.6\\
51	2161.3\\
35	2977.7\\
29	1374.2\\
46	1032.2\\
44	3649.5\\
45	3402\\
29	1010.9\\
28	1754.2\\
77	3114.5\\
70	4119.9\\
70	3459.3\\
40	2452.3\\
90	5508.8\\
53	3724.7\\
62	1830.6\\
47	3515.2\\
34	1398.8\\
14	1141.1\\
47	2105.7\\
43	1139.6\\
20	837.7\\
28	1115.9\\
80	3241.9\\
25	585.76\\
36	1823.2\\
27	2342.7\\
82	4291\\
31	1114.6\\
39	1413.8\\
23	2288.8\\
27	2906.9\\
37	2237.9\\
129	8493.1\\
63	4393.1\\
64	2048.5\\
91	9723.9\\
27	1289.5\\
26	1549.7\\
17	871.75\\
33	1646.7\\
30	3259.2\\
33	1433.6\\
182	13893\\
26	1236.6\\
48	2947.2\\
46	1881.8\\
10	415.77\\
25	1140.1\\
26	828.68\\
206	25071\\
37	1375.7\\
43	1761.7\\
51	2030.1\\
57	6651.6\\
36	1835.1\\
38	1643\\
28	1321.5\\
19	1171.2\\
36	1634.1\\
45	3307.9\\
53	3596.4\\
48	3747.5\\
64	3253.5\\
65	3044.1\\
45	1336.3\\
54	2967.3\\
16	455.41\\
24	1128\\
127	8559.2\\
48	1821.8\\
64	2618.1\\
42	2765.8\\
42	1378.1\\
17	586.15\\
48	2910.3\\
42	2231.6\\
20	914.71\\
34	1268.4\\
47	1045.9\\
42	3013.7\\
50	2424.2\\
87	6521.5\\
64	4006.9\\
32	1840.5\\
39	1209.7\\
35	2555.9\\
18	1133.1\\
28	973.41\\
27	940.21\\
37	2134\\
30	2497.3\\
89	11551\\
22	1878.1\\
28	1780.4\\
25	1330.8\\
15	667.92\\
41	3516.1\\
32	1002.2\\
34	2026.6\\
16	828.58\\
88	5883.1\\
39	1056\\
140	27387\\
39	1350.3\\
49	1943.7\\
72	2461.3\\
37	1763.1\\
41	3630.8\\
283	1.4487e+05\\
31	2920.1\\
99	18428\\
103	23585\\
39	3789.3\\
35	1584.9\\
98	14081\\
45	4414.9\\
73	8394.7\\
34	2362.6\\
56	8159.8\\
39	3038.7\\
37	3037.4\\
55	2240.3\\
25	753.94\\
29	1153.5\\
42	1269.3\\
35	4016.1\\
26	834.23\\
45	8300\\
56	6301.2\\
39	1615.9\\
62	3239.2\\
94	12254\\
112	23459\\
25	809.04\\
57	2507.1\\
50	2505.7\\
47	6802\\
96	7017.5\\
54	6405.3\\
28	2518.5\\
38	3907.2\\
47	3713.5\\
35	2097.6\\
45	1741.8\\
50	3782.7\\
43	1082.6\\
29	2197.2\\
42	1911\\
71	20010\\
68	6763.2\\
149	53138\\
67	5699\\
51	2876.3\\
34	1914.9\\
120	14460\\
43	7072.6\\
117	10981\\
39	1411.9\\
72	8351\\
40	3724.3\\
43	2486.6\\
32	2737.4\\
49	8641.3\\
41	4083.2\\
98	5192.2\\
96	12544\\
45	1809.2\\
129	34370\\
43	2119.8\\
93	8475.2\\
41	3465.1\\
88	9873.1\\
73	5908.1\\
108	15487\\
52	7655.5\\
18	505.46\\
36	1113.1\\
50	5454.7\\
37	2355.7\\
83	14366\\
40	2340.7\\
26	2737.3\\
40	1826.7\\
49	3905.9\\
27	1487.2\\
41	1437.4\\
36	953.53\\
52	5068.7\\
37	1157.2\\
68	7232.3\\
43	2881.6\\
132	13720\\
73	2977.2\\
46	4837.2\\
47	1324.5\\
44	1703.1\\
41	4892.5\\
30	1582.5\\
25	3061.1\\
24	2026.4\\
85	8195.2\\
44	3072.7\\
29	1926.8\\
34	1721\\
101	14583\\
18	1693.2\\
21	1458\\
111	16678\\
95	4804.9\\
23	1106.1\\
31	689.7\\
84	4100.7\\
23	870.77\\
55	5070.7\\
32	1217\\
27	968.18\\
48	1041.7\\
18	391.93\\
28	2072.4\\
40	1250.5\\
30	1350.7\\
13	356.06\\
33	1628.3\\
14	604.68\\
52	2461.4\\
65	7655.9\\
49	2550.4\\
35	1466.4\\
109	7102.6\\
26	997.27\\
41	1368.6\\
39	1301.6\\
61	6050.8\\
40	1545.7\\
36	2134.3\\
21	671.2\\
80	4229\\
36	1424.9\\
36	2036.9\\
38	1116\\
42	2180.6\\
27	784.16\\
46	2310\\
23	735.23\\
34	1825.9\\
200	38073\\
37	2527.5\\
16	655\\
50	2495.1\\
19	721.05\\
41	2929.8\\
25	870.63\\
39	2139.7\\
45	3045.6\\
83	7209.4\\
27	598.19\\
67	2609.2\\
46	3129.2\\
35	1702.9\\
44	2550.5\\
30	1894.4\\
29	1162.5\\
23	578.24\\
43	2057.2\\
28	834.06\\
24	1151.9\\
56	3158.8\\
33	1003.5\\
13	434.23\\
41	3076.6\\
34	1023\\
14	491.72\\
83	8702.6\\
33	1728.6\\
35	1659.6\\
29	1040.4\\
54	3772\\
41	1358.7\\
19	771.12\\
34	2140.3\\
19	742.96\\
38	2267.9\\
43	1822.3\\
14	711.02\\
17	831.4\\
12	482.97\\
};
\addlegendentry{\ApproxScatterLegend}

      % \input{../../../../.tex/20/KS/SizeVsDegreefig1.tex}
  
      \nextgroupplot[%
        % xmin=10,xmax=283,
        % ymin=100,ymax=1e+06,
      ] % This file was created by matlab2tikz.
%
\definecolor{mycolor1}{rgb}{0.00000,0.44706,0.69804}%
\definecolor{mycolor2}{rgb}{0.00000,0.44700,0.74100}%
\definecolor{mycolor3}{rgb}{0.85000,0.32500,0.09800}%
%
\addplot[only marks, mark=*, mark size=1.2748pt, color=mycolor1, fill=mycolor1, opacity=0.60, draw=none, mark options={draw=none,line width=0pt}] table[row sep=crcr]{%
x	y\\
34	2851.2\\
29	1488.4\\
96	16990\\
31	2007.1\\
27	1372.1\\
70	7626.2\\
24	1332.2\\
37	2351.4\\
42	3057.2\\
16	1215.1\\
29	1589.9\\
59	3535.1\\
69	3469.4\\
28	1682.6\\
11	360.9\\
20	1244\\
45	4011.6\\
27	1558.3\\
22	1094.6\\
27	1537.3\\
44	5285\\
16	1620\\
39	4631.6\\
19	1206\\
33	2235.4\\
37	2578.7\\
33	1433.6\\
18	623.45\\
29	2459.2\\
22	1102.8\\
33	1514.6\\
22	966.47\\
49	8453.1\\
25	1617.8\\
43	4050.1\\
26	2521.2\\
31	3217.6\\
26	1287.7\\
18	981.84\\
13	593.7\\
32	4139.3\\
38	2521\\
32	3592.8\\
29	1929.6\\
35	1658.7\\
35	3284.3\\
135	20146\\
27	4702.7\\
42	4421.1\\
39	4834.9\\
51	7571.5\\
86	10528\\
15	937.7\\
39	4873.2\\
42	3484.1\\
43	3373.7\\
58	5802.4\\
97	24886\\
50	2570.4\\
20	1602.4\\
28	1293.5\\
22	1268.1\\
47	3423.8\\
154	61793\\
30	1920\\
10	536.5\\
37	8161.1\\
31	1986.1\\
54	13976\\
76	11720\\
55	4925.9\\
26	2800\\
33	1722.9\\
34	2256.7\\
25	1771.9\\
31	3867.8\\
34	4703\\
26	2423.7\\
39	3770.4\\
27	2850.1\\
32	2139.1\\
26	2352.2\\
29	2367.4\\
23	3345.2\\
26	2487.8\\
14	914.37\\
33	2195.8\\
24	1162.2\\
20	1771.9\\
51	2172.1\\
35	3001.5\\
29	1353.1\\
46	1034.1\\
44	3716.9\\
45	3436.3\\
29	998.2\\
28	1793.6\\
77	3131.6\\
70	4049.4\\
70	3448.9\\
40	2517.7\\
90	5635.2\\
53	3698.2\\
62	1845.8\\
47	3592.4\\
34	1400.3\\
14	1125.1\\
47	2135.7\\
43	1176.2\\
20	842.46\\
28	1135.9\\
80	3228.8\\
25	595.71\\
36	1851.7\\
27	2283.8\\
82	4333.3\\
31	1132.7\\
39	1436.6\\
23	2238.8\\
27	2852.4\\
37	2272.9\\
129	8365.6\\
63	4436.6\\
64	2083.9\\
91	9555.9\\
27	1302.6\\
26	1509.8\\
17	902.04\\
33	1651.8\\
30	3363.6\\
33	1463.8\\
182	14018\\
26	1213.6\\
48	2986.2\\
46	1847.4\\
10	416.32\\
25	1175.1\\
26	835.94\\
206	24904\\
37	1338\\
43	1681\\
51	2035.7\\
57	6649.8\\
36	1806.8\\
38	1641.1\\
28	1368.6\\
19	1146.4\\
36	1679.5\\
45	3310.6\\
53	3621.7\\
48	3690.2\\
64	3215.8\\
65	3055.7\\
45	1347.5\\
54	2975.2\\
16	461.78\\
24	1108.5\\
127	8558.5\\
48	1856.5\\
64	2648.1\\
42	2749.7\\
42	1361.6\\
17	597.94\\
48	2948.6\\
42	2213.5\\
20	904.25\\
34	1259.9\\
47	1059.8\\
42	3001.8\\
50	2383.4\\
87	6514.1\\
64	3972.2\\
32	1852.1\\
39	1222.4\\
35	2584.8\\
18	1126\\
28	955.03\\
27	931.16\\
37	2131.9\\
30	2552.9\\
89	11576\\
22	1901.6\\
28	1701.4\\
25	1377.4\\
15	661.94\\
41	3504.7\\
32	1007.1\\
34	2052.9\\
16	828.35\\
88	5896.4\\
39	1033.1\\
140	26690\\
39	1350.9\\
49	1889.1\\
72	2477.6\\
37	1743.9\\
41	3578.9\\
283	1.438e+05\\
31	2923.4\\
99	18758\\
103	23535\\
39	3565.8\\
35	1571.9\\
98	13823\\
45	4384.9\\
73	8387.2\\
34	2340.2\\
56	7986.5\\
39	2947.2\\
37	3022.8\\
55	2215.8\\
25	763.38\\
29	1148.4\\
42	1240.7\\
35	3961\\
26	810.38\\
45	8358.7\\
56	6140.8\\
39	1615.2\\
62	3234.6\\
94	12172\\
112	23494\\
25	816.96\\
57	2517.7\\
50	2499.6\\
47	6505.8\\
96	7074.7\\
54	6265.4\\
28	2549.5\\
38	3812.9\\
47	3644.4\\
35	2112.1\\
45	1705.3\\
50	3764.2\\
43	1058.1\\
29	2181.7\\
42	1931.5\\
71	19926\\
68	6539.7\\
149	52722\\
67	5576.8\\
51	2915.1\\
34	1902.6\\
120	14465\\
43	7099.9\\
117	10679\\
39	1383.5\\
72	8408.7\\
40	3759.5\\
43	2443.1\\
32	2686.5\\
49	8503.8\\
41	3985\\
98	5104.3\\
96	12397\\
45	1739.9\\
129	33928\\
43	2178.4\\
93	8151.4\\
41	3386.4\\
88	9863.8\\
73	5712.5\\
108	15614\\
52	7527\\
18	497.78\\
36	1113.5\\
50	5373.2\\
37	2378.2\\
83	14111\\
40	2259.3\\
26	2719.2\\
40	1757.9\\
49	3820.3\\
27	1475.5\\
41	1434.6\\
36	933.28\\
52	5017.2\\
37	1160.8\\
68	7126.2\\
43	2855.2\\
132	13488\\
73	3056.7\\
46	4725.5\\
47	1278.6\\
44	1692.7\\
41	4987.9\\
30	1567.7\\
25	3087.5\\
24	1985.1\\
85	8156.8\\
44	3066.6\\
29	1868.2\\
34	1705.2\\
101	14451\\
18	1614.8\\
21	1428.4\\
111	16710\\
95	4784.4\\
23	1109.2\\
31	673.8\\
84	4064.8\\
23	846.48\\
55	5089.3\\
32	1184\\
27	957.19\\
48	1013.7\\
18	391.02\\
28	2036.3\\
40	1235.9\\
30	1337.6\\
13	353.76\\
33	1614.3\\
14	604.62\\
52	2477.5\\
65	7352.1\\
49	2634.2\\
35	1456.3\\
109	6953.3\\
26	985.04\\
41	1387.1\\
39	1294.3\\
61	6040.5\\
40	1520.1\\
36	2129.3\\
21	674.14\\
80	4224.6\\
36	1404\\
36	2016.1\\
38	1082.4\\
42	2146.8\\
27	783.46\\
46	2314.8\\
23	695.68\\
34	1797.3\\
200	37843\\
37	2495.7\\
16	668.81\\
50	2506.4\\
19	701.77\\
41	2959.9\\
25	866.67\\
39	2136.6\\
45	3052.8\\
83	7235.6\\
27	579.37\\
67	2635\\
46	3155.8\\
35	1676.1\\
44	2589.4\\
30	1928.3\\
29	1144.4\\
23	563.44\\
43	2045.7\\
28	811.18\\
24	1157.9\\
56	3097.1\\
33	994.96\\
13	439.83\\
41	2987\\
34	1022.6\\
14	499.21\\
83	8650.7\\
33	1702.1\\
35	1653.7\\
29	1052.6\\
54	3741.2\\
41	1343.3\\
19	753.42\\
34	2125.6\\
19	728.55\\
38	2246.1\\
43	1807.7\\
14	705.51\\
17	830.56\\
12	492.07\\
};
\addlegendentry{\ExactScatterLegend}

\addplot[only marks, mark=*, mark size=1.2748pt, color=black, fill=black, opacity=0.60, draw=none, mark options={draw=none,line width=0pt}] table[row sep=crcr]{%
x	y\\
34	1787\\
29	2052\\
96	39026\\
31	939\\
27	839\\
70	9435\\
24	756\\
37	2292\\
42	3353\\
16	509\\
29	1198\\
59	4061\\
69	4632\\
28	987\\
11	377\\
20	825\\
45	6367\\
27	1370\\
22	689\\
27	1107\\
44	4679\\
16	622\\
39	5236\\
19	558\\
33	1997\\
37	1291\\
33	1077\\
18	485\\
29	1640\\
22	758\\
33	1231\\
22	643\\
49	9267\\
25	1115\\
43	11048\\
26	1875\\
31	2762\\
26	927\\
18	665\\
13	140\\
32	2698\\
38	2211\\
32	775\\
29	1173\\
35	1719\\
35	3078\\
135	41095\\
27	2591\\
42	3900\\
39	3847\\
51	5607\\
86	11830\\
15	918\\
39	3169\\
42	3772\\
43	2489\\
58	4861\\
97	21264\\
50	3266\\
20	736\\
28	686\\
22	1102\\
47	4024\\
154	1.2234e+05\\
30	1543\\
10	296\\
37	7252\\
31	1142\\
54	13398\\
76	13899\\
55	3344\\
26	1499\\
33	1121\\
34	1971\\
25	1709\\
31	2961\\
34	3625\\
26	2774\\
39	3551\\
27	1922\\
32	1860\\
26	1781\\
29	1942\\
23	2014\\
26	1636\\
14	530\\
33	1466\\
24	787\\
20	1114\\
51	1692\\
35	2940\\
29	1332\\
46	1054\\
44	3996\\
45	4071\\
29	803\\
28	557\\
77	3928\\
70	3625\\
70	2531\\
40	1723\\
90	8518\\
53	3650\\
62	3213\\
47	8035\\
34	817\\
14	526\\
47	2742\\
43	759\\
20	950\\
28	533\\
80	4654\\
25	1153\\
36	1716\\
27	2299\\
82	3021\\
31	1103\\
39	1587\\
23	810\\
27	2388\\
37	2269\\
129	3241\\
63	3568\\
64	2459\\
91	6356\\
27	696\\
26	1459\\
17	558\\
33	2465\\
30	2046\\
33	1791\\
182	11424\\
26	663\\
48	2633\\
46	2157\\
10	222\\
25	802\\
26	424\\
206	37527\\
37	1069\\
43	1448\\
51	2715\\
57	7724\\
36	1800\\
38	1162\\
28	539\\
19	746\\
36	1003\\
45	3212\\
53	4779\\
48	5264\\
64	2484\\
65	3084\\
45	1597\\
54	3279\\
16	299\\
24	1131\\
127	2601\\
48	1783\\
64	2544\\
42	2057\\
42	1206\\
17	254\\
48	2507\\
42	2016\\
20	816\\
34	1805\\
47	1076\\
42	2516\\
50	2233\\
87	10377\\
64	2072\\
32	1299\\
39	1215\\
35	3735\\
18	866\\
28	626\\
27	253\\
37	2518\\
30	2667\\
89	9128\\
22	1204\\
28	1732\\
25	1513\\
15	873\\
41	3761\\
32	1140\\
34	1426\\
16	352\\
88	7596\\
39	668\\
140	20491\\
39	1086\\
49	1025\\
72	1475\\
37	1233\\
41	3002\\
283	2.0424e+05\\
31	2681\\
99	16428\\
103	32887\\
39	6629\\
35	1076\\
98	6332\\
45	4113\\
73	7877\\
34	1444\\
56	6926\\
39	2045\\
37	3241\\
55	1793\\
25	444\\
29	1020\\
42	1480\\
35	2286\\
26	591\\
45	5458\\
56	7320\\
39	1112\\
62	3063\\
94	13380\\
112	30134\\
25	317\\
57	1896\\
50	2648\\
47	5982\\
96	4539\\
54	4348\\
28	1451\\
38	3579\\
47	3124\\
35	1834\\
45	921\\
50	3106\\
43	719\\
29	1516\\
42	1249\\
71	5868\\
68	5857\\
149	61636\\
67	5463\\
51	1708\\
34	1459\\
120	10119\\
43	6342\\
117	8499\\
39	1041\\
72	6468\\
40	4061\\
43	1854\\
32	2516\\
49	12313\\
41	4009\\
98	4503\\
96	5379\\
45	1383\\
129	23237\\
43	1506\\
93	4228\\
41	2160\\
88	9837\\
73	5327\\
108	12182\\
52	5524\\
18	189\\
36	878\\
50	4430\\
37	1383\\
83	13086\\
40	2202\\
26	1338\\
40	1198\\
49	4702\\
27	1189\\
41	1263\\
36	572\\
52	3610\\
37	656\\
68	6317\\
43	2092\\
132	14984\\
73	3147\\
46	3825\\
47	739\\
44	1624\\
41	5048\\
30	1539\\
25	1574\\
24	2598\\
85	7294\\
44	1753\\
29	1546\\
34	1160\\
101	9599\\
18	1226\\
21	992\\
111	7348\\
95	2640\\
23	617\\
31	310\\
84	1691\\
23	466\\
55	3785\\
32	1352\\
27	544\\
48	517\\
18	119\\
28	1203\\
40	962\\
30	685\\
13	188\\
33	1762\\
14	209\\
52	1836\\
65	8994\\
49	3401\\
35	1176\\
109	4663\\
26	550\\
41	1006\\
39	973\\
61	4953\\
40	1238\\
36	1694\\
21	551\\
80	5029\\
36	1037\\
36	1691\\
38	832\\
42	1192\\
27	590\\
46	1533\\
23	486\\
34	1267\\
200	30990\\
37	1240\\
16	413\\
50	2688\\
19	324\\
41	2143\\
25	934\\
39	3741\\
45	2869\\
83	3945\\
27	521\\
67	2908\\
46	2419\\
35	1912\\
44	2676\\
30	2051\\
29	674\\
23	204\\
43	936\\
28	859\\
24	467\\
56	2173\\
33	679\\
13	261\\
41	2629\\
34	588\\
14	216\\
83	10336\\
33	1005\\
35	1380\\
29	687\\
54	3377\\
41	993\\
19	351\\
34	1937\\
19	459\\
38	1126\\
43	1196\\
14	412\\
17	325\\
12	184\\
};
\addlegendentry{\RealScatterLegend}

      % \input{../../../../.tex/20/KS/SizeVsDegreefig2.tex}
  
      \nextgroupplot[%
        % xmin=10,xmax=283,
        % ymin=100,ymax=1e+06,
      ] % This file was created by matlab2tikz.
%
\definecolor{mycolor1}{rgb}{0.83529,0.36863,0.00000}%
\definecolor{mycolor2}{rgb}{0.00000,0.44700,0.74100}%
\definecolor{mycolor3}{rgb}{0.85000,0.32500,0.09800}%
%
\addplot[only marks, mark=*, mark size=1.2748pt, color=mycolor1, fill=mycolor1, opacity=0.60, draw=none, mark options={draw=none,line width=0pt}] table[row sep=crcr]{%
x	y\\
34	2708.6\\
29	1423.3\\
96	16788\\
31	1991.2\\
27	1363.5\\
70	7601.1\\
24	1323.7\\
37	2304.2\\
42	2929.6\\
16	1214.3\\
29	1534.3\\
59	3573.2\\
69	3477.3\\
28	1623.5\\
11	345.68\\
20	1220.4\\
45	3757.7\\
27	1531.9\\
22	1092.3\\
27	1484.3\\
44	5224.2\\
16	1590.3\\
39	4459.9\\
19	1180.8\\
33	2224.1\\
37	2546.2\\
33	1438.7\\
18	610.78\\
29	2463.3\\
22	1069\\
33	1515.7\\
22	922.49\\
49	8173.5\\
25	1593.4\\
43	3932.6\\
26	2442.1\\
31	3201.4\\
26	1292\\
18	950.25\\
13	577.17\\
32	4091.3\\
38	2467.3\\
32	3493.5\\
29	1855.7\\
35	1672.7\\
35	3190\\
135	19679\\
27	4549.1\\
42	4314.1\\
39	4692.7\\
51	7190\\
86	10229\\
15	951.13\\
39	4659.8\\
42	3511.9\\
43	3263.5\\
58	5657.8\\
97	24092\\
50	2607.2\\
20	1557\\
28	1270.3\\
22	1215.5\\
47	3344.3\\
154	60852\\
30	1852.4\\
10	515.06\\
37	8050.8\\
31	1938.9\\
54	13650\\
76	11498\\
55	4712.9\\
26	2729.4\\
33	1672\\
34	2166.4\\
25	1707.7\\
31	3836.1\\
34	4507.7\\
26	2365.7\\
39	3679.1\\
27	2806.9\\
32	2117.6\\
26	2289.8\\
29	2351\\
23	3308.7\\
26	2417.2\\
14	873.66\\
33	2127.1\\
24	1133.5\\
20	1719.6\\
51	2161.3\\
35	2977.7\\
29	1374.2\\
46	1032.2\\
44	3649.5\\
45	3402\\
29	1010.9\\
28	1754.2\\
77	3114.5\\
70	4119.9\\
70	3459.3\\
40	2452.3\\
90	5508.8\\
53	3724.7\\
62	1830.6\\
47	3515.2\\
34	1398.8\\
14	1141.1\\
47	2105.7\\
43	1139.6\\
20	837.7\\
28	1115.9\\
80	3241.9\\
25	585.76\\
36	1823.2\\
27	2342.7\\
82	4291\\
31	1114.6\\
39	1413.8\\
23	2288.8\\
27	2906.9\\
37	2237.9\\
129	8493.1\\
63	4393.1\\
64	2048.5\\
91	9723.9\\
27	1289.5\\
26	1549.7\\
17	871.75\\
33	1646.7\\
30	3259.2\\
33	1433.6\\
182	13893\\
26	1236.6\\
48	2947.2\\
46	1881.8\\
10	415.77\\
25	1140.1\\
26	828.68\\
206	25071\\
37	1375.7\\
43	1761.7\\
51	2030.1\\
57	6651.6\\
36	1835.1\\
38	1643\\
28	1321.5\\
19	1171.2\\
36	1634.1\\
45	3307.9\\
53	3596.4\\
48	3747.5\\
64	3253.5\\
65	3044.1\\
45	1336.3\\
54	2967.3\\
16	455.41\\
24	1128\\
127	8559.2\\
48	1821.8\\
64	2618.1\\
42	2765.8\\
42	1378.1\\
17	586.15\\
48	2910.3\\
42	2231.6\\
20	914.71\\
34	1268.4\\
47	1045.9\\
42	3013.7\\
50	2424.2\\
87	6521.5\\
64	4006.9\\
32	1840.5\\
39	1209.7\\
35	2555.9\\
18	1133.1\\
28	973.41\\
27	940.21\\
37	2134\\
30	2497.3\\
89	11551\\
22	1878.1\\
28	1780.4\\
25	1330.8\\
15	667.92\\
41	3516.1\\
32	1002.2\\
34	2026.6\\
16	828.58\\
88	5883.1\\
39	1056\\
140	27387\\
39	1350.3\\
49	1943.7\\
72	2461.3\\
37	1763.1\\
41	3630.8\\
283	1.4487e+05\\
31	2920.1\\
99	18428\\
103	23585\\
39	3789.3\\
35	1584.9\\
98	14081\\
45	4414.9\\
73	8394.7\\
34	2362.6\\
56	8159.8\\
39	3038.7\\
37	3037.4\\
55	2240.3\\
25	753.94\\
29	1153.5\\
42	1269.3\\
35	4016.1\\
26	834.23\\
45	8300\\
56	6301.2\\
39	1615.9\\
62	3239.2\\
94	12254\\
112	23459\\
25	809.04\\
57	2507.1\\
50	2505.7\\
47	6802\\
96	7017.5\\
54	6405.3\\
28	2518.5\\
38	3907.2\\
47	3713.5\\
35	2097.6\\
45	1741.8\\
50	3782.7\\
43	1082.6\\
29	2197.2\\
42	1911\\
71	20010\\
68	6763.2\\
149	53138\\
67	5699\\
51	2876.3\\
34	1914.9\\
120	14460\\
43	7072.6\\
117	10981\\
39	1411.9\\
72	8351\\
40	3724.3\\
43	2486.6\\
32	2737.4\\
49	8641.3\\
41	4083.2\\
98	5192.2\\
96	12544\\
45	1809.2\\
129	34370\\
43	2119.8\\
93	8475.2\\
41	3465.1\\
88	9873.1\\
73	5908.1\\
108	15487\\
52	7655.5\\
18	505.46\\
36	1113.1\\
50	5454.7\\
37	2355.7\\
83	14366\\
40	2340.7\\
26	2737.3\\
40	1826.7\\
49	3905.9\\
27	1487.2\\
41	1437.4\\
36	953.53\\
52	5068.7\\
37	1157.2\\
68	7232.3\\
43	2881.6\\
132	13720\\
73	2977.2\\
46	4837.2\\
47	1324.5\\
44	1703.1\\
41	4892.5\\
30	1582.5\\
25	3061.1\\
24	2026.4\\
85	8195.2\\
44	3072.7\\
29	1926.8\\
34	1721\\
101	14583\\
18	1693.2\\
21	1458\\
111	16678\\
95	4804.9\\
23	1106.1\\
31	689.7\\
84	4100.7\\
23	870.77\\
55	5070.7\\
32	1217\\
27	968.18\\
48	1041.7\\
18	391.93\\
28	2072.4\\
40	1250.5\\
30	1350.7\\
13	356.06\\
33	1628.3\\
14	604.68\\
52	2461.4\\
65	7655.9\\
49	2550.4\\
35	1466.4\\
109	7102.6\\
26	997.27\\
41	1368.6\\
39	1301.6\\
61	6050.8\\
40	1545.7\\
36	2134.3\\
21	671.2\\
80	4229\\
36	1424.9\\
36	2036.9\\
38	1116\\
42	2180.6\\
27	784.16\\
46	2310\\
23	735.23\\
34	1825.9\\
200	38073\\
37	2527.5\\
16	655\\
50	2495.1\\
19	721.05\\
41	2929.8\\
25	870.63\\
39	2139.7\\
45	3045.6\\
83	7209.4\\
27	598.19\\
67	2609.2\\
46	3129.2\\
35	1702.9\\
44	2550.5\\
30	1894.4\\
29	1162.5\\
23	578.24\\
43	2057.2\\
28	834.06\\
24	1151.9\\
56	3158.8\\
33	1003.5\\
13	434.23\\
41	3076.6\\
34	1023\\
14	491.72\\
83	8702.6\\
33	1728.6\\
35	1659.6\\
29	1040.4\\
54	3772\\
41	1358.7\\
19	771.12\\
34	2140.3\\
19	742.96\\
38	2267.9\\
43	1822.3\\
14	711.02\\
17	831.4\\
12	482.97\\
};
\addlegendentry{\ApproxScatterLegend}

\addplot[only marks, mark=*, mark size=1.2748pt, color=black, fill=black, opacity=0.60, draw=none, mark options={draw=none,line width=0pt}] table[row sep=crcr]{%
x	y\\
34	1787\\
29	2052\\
96	39026\\
31	939\\
27	839\\
70	9435\\
24	756\\
37	2292\\
42	3353\\
16	509\\
29	1198\\
59	4061\\
69	4632\\
28	987\\
11	377\\
20	825\\
45	6367\\
27	1370\\
22	689\\
27	1107\\
44	4679\\
16	622\\
39	5236\\
19	558\\
33	1997\\
37	1291\\
33	1077\\
18	485\\
29	1640\\
22	758\\
33	1231\\
22	643\\
49	9267\\
25	1115\\
43	11048\\
26	1875\\
31	2762\\
26	927\\
18	665\\
13	140\\
32	2698\\
38	2211\\
32	775\\
29	1173\\
35	1719\\
35	3078\\
135	41095\\
27	2591\\
42	3900\\
39	3847\\
51	5607\\
86	11830\\
15	918\\
39	3169\\
42	3772\\
43	2489\\
58	4861\\
97	21264\\
50	3266\\
20	736\\
28	686\\
22	1102\\
47	4024\\
154	1.2234e+05\\
30	1543\\
10	296\\
37	7252\\
31	1142\\
54	13398\\
76	13899\\
55	3344\\
26	1499\\
33	1121\\
34	1971\\
25	1709\\
31	2961\\
34	3625\\
26	2774\\
39	3551\\
27	1922\\
32	1860\\
26	1781\\
29	1942\\
23	2014\\
26	1636\\
14	530\\
33	1466\\
24	787\\
20	1114\\
51	1692\\
35	2940\\
29	1332\\
46	1054\\
44	3996\\
45	4071\\
29	803\\
28	557\\
77	3928\\
70	3625\\
70	2531\\
40	1723\\
90	8518\\
53	3650\\
62	3213\\
47	8035\\
34	817\\
14	526\\
47	2742\\
43	759\\
20	950\\
28	533\\
80	4654\\
25	1153\\
36	1716\\
27	2299\\
82	3021\\
31	1103\\
39	1587\\
23	810\\
27	2388\\
37	2269\\
129	3241\\
63	3568\\
64	2459\\
91	6356\\
27	696\\
26	1459\\
17	558\\
33	2465\\
30	2046\\
33	1791\\
182	11424\\
26	663\\
48	2633\\
46	2157\\
10	222\\
25	802\\
26	424\\
206	37527\\
37	1069\\
43	1448\\
51	2715\\
57	7724\\
36	1800\\
38	1162\\
28	539\\
19	746\\
36	1003\\
45	3212\\
53	4779\\
48	5264\\
64	2484\\
65	3084\\
45	1597\\
54	3279\\
16	299\\
24	1131\\
127	2601\\
48	1783\\
64	2544\\
42	2057\\
42	1206\\
17	254\\
48	2507\\
42	2016\\
20	816\\
34	1805\\
47	1076\\
42	2516\\
50	2233\\
87	10377\\
64	2072\\
32	1299\\
39	1215\\
35	3735\\
18	866\\
28	626\\
27	253\\
37	2518\\
30	2667\\
89	9128\\
22	1204\\
28	1732\\
25	1513\\
15	873\\
41	3761\\
32	1140\\
34	1426\\
16	352\\
88	7596\\
39	668\\
140	20491\\
39	1086\\
49	1025\\
72	1475\\
37	1233\\
41	3002\\
283	2.0424e+05\\
31	2681\\
99	16428\\
103	32887\\
39	6629\\
35	1076\\
98	6332\\
45	4113\\
73	7877\\
34	1444\\
56	6926\\
39	2045\\
37	3241\\
55	1793\\
25	444\\
29	1020\\
42	1480\\
35	2286\\
26	591\\
45	5458\\
56	7320\\
39	1112\\
62	3063\\
94	13380\\
112	30134\\
25	317\\
57	1896\\
50	2648\\
47	5982\\
96	4539\\
54	4348\\
28	1451\\
38	3579\\
47	3124\\
35	1834\\
45	921\\
50	3106\\
43	719\\
29	1516\\
42	1249\\
71	5868\\
68	5857\\
149	61636\\
67	5463\\
51	1708\\
34	1459\\
120	10119\\
43	6342\\
117	8499\\
39	1041\\
72	6468\\
40	4061\\
43	1854\\
32	2516\\
49	12313\\
41	4009\\
98	4503\\
96	5379\\
45	1383\\
129	23237\\
43	1506\\
93	4228\\
41	2160\\
88	9837\\
73	5327\\
108	12182\\
52	5524\\
18	189\\
36	878\\
50	4430\\
37	1383\\
83	13086\\
40	2202\\
26	1338\\
40	1198\\
49	4702\\
27	1189\\
41	1263\\
36	572\\
52	3610\\
37	656\\
68	6317\\
43	2092\\
132	14984\\
73	3147\\
46	3825\\
47	739\\
44	1624\\
41	5048\\
30	1539\\
25	1574\\
24	2598\\
85	7294\\
44	1753\\
29	1546\\
34	1160\\
101	9599\\
18	1226\\
21	992\\
111	7348\\
95	2640\\
23	617\\
31	310\\
84	1691\\
23	466\\
55	3785\\
32	1352\\
27	544\\
48	517\\
18	119\\
28	1203\\
40	962\\
30	685\\
13	188\\
33	1762\\
14	209\\
52	1836\\
65	8994\\
49	3401\\
35	1176\\
109	4663\\
26	550\\
41	1006\\
39	973\\
61	4953\\
40	1238\\
36	1694\\
21	551\\
80	5029\\
36	1037\\
36	1691\\
38	832\\
42	1192\\
27	590\\
46	1533\\
23	486\\
34	1267\\
200	30990\\
37	1240\\
16	413\\
50	2688\\
19	324\\
41	2143\\
25	934\\
39	3741\\
45	2869\\
83	3945\\
27	521\\
67	2908\\
46	2419\\
35	1912\\
44	2676\\
30	2051\\
29	674\\
23	204\\
43	936\\
28	859\\
24	467\\
56	2173\\
33	679\\
13	261\\
41	2629\\
34	588\\
14	216\\
83	10336\\
33	1005\\
35	1380\\
29	687\\
54	3377\\
41	993\\
19	351\\
34	1937\\
19	459\\
38	1126\\
43	1196\\
14	412\\
17	325\\
12	184\\
};
\addlegendentry{\RealScatterLegend}

      % \input{../../../../.tex/20/KS/SizeVsDegreefig3.tex}
    \end{groupplot}
    
    \begin{groupplot}[
      group style={
        group name=Row2,
        plotGroup
      },
      plotRow,
      row4Keys,
      plotLegendNW,
      %
      xmin=10,xmax=283,
      ymin=100,ymax=2.0424e+05,
    ]
      \nextgroupplot[%
        at={($(Row1 c1r1.south west)-(0,\yPlotSep em)-(0,\plotHeight cm)$)},
        % xmin=10,xmax=283,
        % ymin=100,ymax=1.4849e+05,
      ] % This file was created by matlab2tikz.
%
\definecolor{mycolor1}{rgb}{0.00000,0.44706,0.69804}%
\definecolor{mycolor2}{rgb}{0.00000,0.44700,0.74100}%
\definecolor{mycolor3}{rgb}{0.83529,0.36863,0.00000}%
\definecolor{mycolor4}{rgb}{0.85000,0.32500,0.09800}%
%
\addplot[only marks, mark=*, mark size=1.2748pt, color=mycolor1, fill=mycolor1, opacity=0.70, draw=none, mark options={draw=none,line width=0pt}] table[row sep=crcr]{%
x	y\\
34	2888.8\\
29	1376\\
96	16967\\
31	1883\\
27	1354.5\\
70	7691.7\\
24	1340.6\\
37	2298.9\\
42	3033.5\\
16	1204.3\\
29	1477.2\\
59	3573.6\\
69	3634.5\\
28	1603.2\\
11	346.44\\
20	1225.1\\
45	3927.8\\
27	1479.6\\
22	1071.8\\
27	1474.6\\
44	5275.8\\
16	1583.6\\
39	4453.9\\
19	1131\\
33	2116.9\\
37	2378.8\\
33	1439\\
18	572.42\\
29	2375.2\\
22	1037.9\\
33	1405.3\\
22	948.54\\
49	7945.6\\
25	1430.5\\
43	3966.8\\
26	2258\\
31	3065.9\\
26	1246\\
18	940.15\\
13	607.27\\
32	4215.6\\
38	2417.3\\
32	3542\\
29	1834.1\\
35	1568.6\\
35	3120.1\\
135	20143\\
27	4391.8\\
42	4455.4\\
39	4621.7\\
51	7539.2\\
86	10440\\
15	971.49\\
39	4685.8\\
42	3296.1\\
43	3405.9\\
58	5638\\
97	24783\\
50	2453.8\\
20	1472.3\\
28	1205.1\\
22	1251.6\\
47	3380.9\\
154	60841\\
30	1824.2\\
10	531.37\\
37	7708.6\\
31	1947.3\\
54	13706\\
76	11783\\
55	4595.8\\
26	2685.2\\
33	1695.6\\
34	2203.9\\
25	1734\\
31	3766.6\\
34	4526.5\\
26	2339.9\\
39	3576.8\\
27	2944\\
32	1979.3\\
26	2230.2\\
29	2371.2\\
23	3394.2\\
26	2460.3\\
14	819.3\\
33	2114.2\\
24	1138.2\\
20	1670.1\\
51	2082.2\\
35	2930.9\\
29	1291.7\\
46	1048.3\\
44	3495.4\\
45	3370.4\\
29	1008.4\\
28	1803.7\\
77	3107.9\\
70	4018\\
70	3561.2\\
40	2416.8\\
90	5643.6\\
53	3404.9\\
62	1805.4\\
47	3720.8\\
34	1349.2\\
14	1080.4\\
47	2173.9\\
43	1159.5\\
20	853.89\\
28	1085.9\\
80	3216.6\\
25	565.84\\
36	1737.6\\
27	2199.1\\
82	4395.5\\
31	1042.3\\
39	1399.9\\
23	2055.7\\
27	2813.1\\
37	2136.5\\
129	8482.3\\
63	4366\\
64	2133.9\\
91	9433.2\\
27	1334.3\\
26	1490.8\\
17	874.18\\
33	1519.9\\
30	3156.6\\
33	1462.3\\
182	14035\\
26	1179.1\\
48	2761.6\\
46	1846.3\\
10	406.6\\
25	1164.1\\
26	798.26\\
206	25015\\
37	1378\\
43	1693.9\\
51	1922.8\\
57	6608.5\\
36	1762.9\\
38	1565.4\\
28	1297.7\\
19	1029.9\\
36	1717\\
45	3068.7\\
53	3669.8\\
48	3616.4\\
64	3205.7\\
65	3143.2\\
45	1277.3\\
54	3009\\
16	445.95\\
24	1070\\
127	8780.1\\
48	1786.1\\
64	2695\\
42	2634.2\\
42	1305.4\\
17	528.1\\
48	2669.7\\
42	2166.4\\
20	861.71\\
34	1148\\
47	1060.4\\
42	2893.9\\
50	2384.2\\
87	6774.2\\
64	4109.9\\
32	1801.8\\
39	1213.2\\
35	2525.2\\
18	1073.3\\
28	885.56\\
27	936.22\\
37	2159.3\\
30	2474.2\\
89	11267\\
22	1890.1\\
28	1628.2\\
25	1350.5\\
15	674.96\\
41	3477.1\\
32	977.09\\
34	1978.1\\
16	794.51\\
88	5767.3\\
39	1044.7\\
140	28786\\
39	1329.3\\
49	1848.5\\
72	2550.9\\
37	1642.1\\
41	3640.2\\
283	1.4502e+05\\
31	2918.5\\
99	18596\\
103	23894\\
39	3560.7\\
35	1471.1\\
98	13894\\
45	4542.5\\
73	8678.5\\
34	2319.7\\
56	7695.7\\
39	3045\\
37	2963.1\\
55	2168.4\\
25	769.54\\
29	1128.6\\
42	1182.2\\
35	4187.3\\
26	814.87\\
45	8433.3\\
56	6286.1\\
39	1581.3\\
62	3285.4\\
94	12909\\
112	24406\\
25	812.74\\
57	2477.9\\
50	2461\\
47	7383.3\\
96	6935.8\\
54	6360.8\\
28	2448\\
38	3985.4\\
47	3590.2\\
35	2047.1\\
45	1774.6\\
50	3648.2\\
43	1063.4\\
29	2218.6\\
42	1832.3\\
71	20091\\
68	6933.4\\
149	55336\\
67	5681.6\\
51	2996.8\\
34	1885.7\\
120	14846\\
43	7071.1\\
117	11205\\
39	1379.3\\
72	8349.7\\
40	3487\\
43	2486.6\\
32	2683.1\\
49	9046.7\\
41	4146.8\\
98	5167.5\\
96	12370\\
45	1772.5\\
129	34962\\
43	2104.4\\
93	7894.9\\
41	3312\\
88	9889.6\\
73	6025.7\\
108	15357\\
52	7620.9\\
18	496.34\\
36	1014.7\\
50	5236.1\\
37	2301.1\\
83	14638\\
40	2289.8\\
26	2752.2\\
40	1812.8\\
49	3890.7\\
27	1453.5\\
41	1484.1\\
36	956.35\\
52	5115.2\\
37	1185.1\\
68	7274.1\\
43	2802.2\\
132	13770\\
73	2948.8\\
46	4597.5\\
47	1315.5\\
44	1727.8\\
41	4974\\
30	1539.1\\
25	3085.3\\
24	1945\\
85	8178.4\\
44	3082.2\\
29	2014.3\\
34	1702\\
101	14402\\
18	1546.9\\
21	1451.7\\
111	16608\\
95	4871.3\\
23	1064.7\\
31	626.98\\
84	4122.2\\
23	782.1\\
55	5165.1\\
32	1213.3\\
27	942.64\\
48	1047.6\\
18	380.73\\
28	2074.6\\
40	1214.5\\
30	1300.5\\
13	327.71\\
33	1553.2\\
14	586.67\\
52	2391.9\\
65	7560.2\\
49	2542\\
35	1430.9\\
109	7366.4\\
26	1003\\
41	1314.3\\
39	1246.9\\
61	6105.4\\
40	1502.1\\
36	2054.8\\
21	658.2\\
80	4121.9\\
36	1344.5\\
36	1962.7\\
38	1042.5\\
42	2105.4\\
27	764.55\\
46	2326.1\\
23	687.64\\
34	1748.5\\
200	38147\\
37	2609.8\\
16	664.69\\
50	2466.9\\
19	703.39\\
41	2962.9\\
25	836.9\\
39	2082.6\\
45	3028\\
83	7391.4\\
27	559.93\\
67	2602\\
46	3083.7\\
35	1662.8\\
44	2449.7\\
30	1978\\
29	1123\\
23	563.56\\
43	1901.1\\
28	794.77\\
24	1185.6\\
56	3170.1\\
33	981.96\\
13	432.13\\
41	3034.1\\
34	950.29\\
14	464.81\\
83	8575.6\\
33	1717.2\\
35	1656.1\\
29	1141.1\\
54	3782.7\\
41	1362.9\\
19	726.37\\
34	2083.8\\
19	709.25\\
38	2121.3\\
43	1743.7\\
14	663.57\\
17	745.77\\
12	473.21\\
};
\addlegendentry{\ExactScatterLegend}

\addplot[only marks, mark=*, mark size=1.2748pt, color=mycolor3, fill=mycolor3, opacity=0.70, draw=none, mark options={draw=none,line width=0pt}] table[row sep=crcr]{%
x	y\\
34	2575.1\\
29	1412\\
96	17027\\
31	2026\\
27	1382.9\\
70	7557.9\\
24	1263.9\\
37	2264.4\\
42	2958.2\\
16	1201.8\\
29	1459.6\\
59	3567.1\\
69	3369.2\\
28	1645.4\\
11	346.55\\
20	1164.6\\
45	3948.2\\
27	1504.4\\
22	1035.2\\
27	1412.4\\
44	5048.7\\
16	1466.5\\
39	4261.8\\
19	1172.7\\
33	2277.8\\
37	2560.6\\
33	1408.3\\
18	584.11\\
29	2494.8\\
22	1042\\
33	1424.6\\
22	911.84\\
49	8130\\
25	1485.4\\
43	3912.7\\
26	2501.4\\
31	3334.2\\
26	1313.2\\
18	896.76\\
13	584.85\\
32	4198.3\\
38	2470.3\\
32	3476\\
29	1789.5\\
35	1558.5\\
35	3271.8\\
135	20059\\
27	4523.3\\
42	4266.3\\
39	4230.4\\
51	7348.6\\
86	10217\\
15	954.08\\
39	4544.5\\
42	3294.2\\
43	3336.5\\
58	5783.3\\
97	24111\\
50	2432.9\\
20	1428\\
28	1186.2\\
22	1118.3\\
47	3307.6\\
154	60390\\
30	1785.5\\
10	509.76\\
37	8330.3\\
31	1938.8\\
54	13828\\
76	11797\\
55	4758\\
26	2727.6\\
33	1700.9\\
34	2034.3\\
25	1657.9\\
31	3581.3\\
34	4486.2\\
26	2277.2\\
39	3716\\
27	2671.7\\
32	2050.2\\
26	2209\\
29	2413\\
23	3300.2\\
26	2326\\
14	910.61\\
33	2154.6\\
24	1073.6\\
20	1738.8\\
51	2077.4\\
35	3047.3\\
29	1402.2\\
46	964.76\\
44	3697.2\\
45	3441.3\\
29	964.33\\
28	1766.5\\
77	3112.2\\
70	3743.2\\
70	3331.2\\
40	2313.4\\
90	5315.9\\
53	3656.8\\
62	1751.3\\
47	3538.1\\
34	1388\\
14	1134.3\\
47	2177.9\\
43	1181.5\\
20	754.18\\
28	1106.4\\
80	3197.5\\
25	539.62\\
36	1720.5\\
27	2156.1\\
82	4326.6\\
31	1088.7\\
39	1357.8\\
23	2138.8\\
27	2858.6\\
37	2219.7\\
129	8284.4\\
63	4431.2\\
64	2037.7\\
91	10009\\
27	1260.9\\
26	1430.9\\
17	875.69\\
33	1591.5\\
30	3101.3\\
33	1458.6\\
182	14317\\
26	1185.6\\
48	2951\\
46	1732.6\\
10	387.65\\
25	1129.3\\
26	836.72\\
206	24860\\
37	1410.7\\
43	1556.8\\
51	2014.4\\
57	6608\\
36	1830.3\\
38	1583.5\\
28	1341.2\\
19	1156.6\\
36	1665.3\\
45	3222.1\\
53	3605.6\\
48	3602.5\\
64	3269.5\\
65	2931.7\\
45	1303.2\\
54	2785.5\\
16	447.12\\
24	1100.1\\
127	8673.4\\
48	1755\\
64	2506.5\\
42	2665.4\\
42	1358\\
17	607.62\\
48	2785.6\\
42	2175\\
20	918.06\\
34	1163.3\\
47	1037.8\\
42	2906.3\\
50	2447.8\\
87	6280\\
64	3963.7\\
32	1807\\
39	1184.9\\
35	2484.6\\
18	1087.3\\
28	925.94\\
27	960.91\\
37	2143.9\\
30	2581.3\\
89	11447\\
22	1773.2\\
28	1632.5\\
25	1260.6\\
15	651.77\\
41	3414.1\\
32	987.86\\
34	1942.3\\
16	850.81\\
88	6024\\
39	992.97\\
140	28024\\
39	1253.4\\
49	1839.8\\
72	2449.5\\
37	1734.5\\
41	3646.9\\
283	1.4849e+05\\
31	2972.4\\
99	19340\\
103	23817\\
39	3662.2\\
35	1559.6\\
98	14258\\
45	4398.3\\
73	8778.5\\
34	2386.2\\
56	7844.2\\
39	2975.8\\
37	3088.2\\
55	2278.7\\
25	749.27\\
29	1093.5\\
42	1224.5\\
35	4115.1\\
26	791.28\\
45	8471\\
56	6165\\
39	1606.6\\
62	3123\\
94	12539\\
112	24734\\
25	807.97\\
57	2667.2\\
50	2372.1\\
47	6685.6\\
96	6912.2\\
54	6489\\
28	2634.3\\
38	3855\\
47	3730.9\\
35	2047\\
45	1635.5\\
50	3648.1\\
43	1081.5\\
29	2191.2\\
42	1917.2\\
71	19877\\
68	6587.1\\
149	53091\\
67	5654.6\\
51	2934.7\\
34	1889.1\\
120	14412\\
43	7413.4\\
117	11148\\
39	1361\\
72	8497.1\\
40	4143.7\\
43	2477.5\\
32	2853.7\\
49	8768.1\\
41	4066.7\\
98	5220.8\\
96	12375\\
45	1693.4\\
129	34629\\
43	2028.1\\
93	8219.8\\
41	3435.5\\
88	9754.2\\
73	5996.3\\
108	15579\\
52	7853.5\\
18	464.22\\
36	1139.6\\
50	5590.6\\
37	2364\\
83	15229\\
40	2207.2\\
26	2506.4\\
40	1778.6\\
49	4042.1\\
27	1459.3\\
41	1402.6\\
36	990.8\\
52	5078.2\\
37	1150.8\\
68	6958.7\\
43	2888.3\\
132	14083\\
73	3040.9\\
46	4455.7\\
47	1220\\
44	1666.2\\
41	4740.1\\
30	1511.3\\
25	3009.3\\
24	1960.5\\
85	8343.8\\
44	3057.1\\
29	1783.9\\
34	1754.7\\
101	14322\\
18	1665.1\\
21	1396.3\\
111	16269\\
95	4741.4\\
23	1018.1\\
31	622.98\\
84	3931.6\\
23	843.66\\
55	4936.1\\
32	1210.6\\
27	906.71\\
48	1017\\
18	397.19\\
28	2000.6\\
40	1205.1\\
30	1334.7\\
13	332.05\\
33	1628.9\\
14	591.57\\
52	2431.5\\
65	7872.8\\
49	2654.1\\
35	1426.4\\
109	6961.9\\
26	980.18\\
41	1349.4\\
39	1246.5\\
61	6024.2\\
40	1503.6\\
36	2082.7\\
21	647.05\\
80	4171.2\\
36	1402.5\\
36	2026.7\\
38	1119.2\\
42	2056.1\\
27	750.24\\
46	2342.5\\
23	709.43\\
34	1804.2\\
200	39277\\
37	2532.9\\
16	619.21\\
50	2586.9\\
19	686.7\\
41	3048.8\\
25	825.51\\
39	2292.8\\
45	2841.4\\
83	7099.8\\
27	604.92\\
67	2542\\
46	3104.8\\
35	1609.9\\
44	2732.6\\
30	1790.2\\
29	1079.6\\
23	579.43\\
43	2081.9\\
28	804.3\\
24	1083.1\\
56	2959.9\\
33	974.5\\
13	418.55\\
41	3134.3\\
34	1015.9\\
14	456.95\\
83	9110\\
33	1791.7\\
35	1614.7\\
29	977.23\\
54	3730.7\\
41	1392.5\\
19	753.9\\
34	1993\\
19	726.16\\
38	2267.4\\
43	1847\\
14	695.77\\
17	863.3\\
12	457.58\\
};
\addlegendentry{\ApproxScatterLegend}

      % \input{../../../../.tex/20/KS/SizeVsDegreefig4.tex}
      \legend{}
  
      \nextgroupplot[%
        % xmin=10,xmax=283,
        % ymin=100,ymax=1e+06,
      ] % This file was created by matlab2tikz.
%
\definecolor{mycolor1}{rgb}{0.00000,0.44706,0.69804}%
\definecolor{mycolor2}{rgb}{0.00000,0.44700,0.74100}%
\definecolor{mycolor3}{rgb}{0.85000,0.32500,0.09800}%
%
\addplot[only marks, mark=*, mark size=1.2748pt, color=mycolor1, fill=mycolor1, opacity=0.60, draw=none, mark options={draw=none,line width=0pt}] table[row sep=crcr]{%
x	y\\
34	2888.8\\
29	1376\\
96	16967\\
31	1883\\
27	1354.5\\
70	7691.7\\
24	1340.6\\
37	2298.9\\
42	3033.5\\
16	1204.3\\
29	1477.2\\
59	3573.6\\
69	3634.5\\
28	1603.2\\
11	346.44\\
20	1225.1\\
45	3927.8\\
27	1479.6\\
22	1071.8\\
27	1474.6\\
44	5275.8\\
16	1583.6\\
39	4453.9\\
19	1131\\
33	2116.9\\
37	2378.8\\
33	1439\\
18	572.42\\
29	2375.2\\
22	1037.9\\
33	1405.3\\
22	948.54\\
49	7945.6\\
25	1430.5\\
43	3966.8\\
26	2258\\
31	3065.9\\
26	1246\\
18	940.15\\
13	607.27\\
32	4215.6\\
38	2417.3\\
32	3542\\
29	1834.1\\
35	1568.6\\
35	3120.1\\
135	20143\\
27	4391.8\\
42	4455.4\\
39	4621.7\\
51	7539.2\\
86	10440\\
15	971.49\\
39	4685.8\\
42	3296.1\\
43	3405.9\\
58	5638\\
97	24783\\
50	2453.8\\
20	1472.3\\
28	1205.1\\
22	1251.6\\
47	3380.9\\
154	60841\\
30	1824.2\\
10	531.37\\
37	7708.6\\
31	1947.3\\
54	13706\\
76	11783\\
55	4595.8\\
26	2685.2\\
33	1695.6\\
34	2203.9\\
25	1734\\
31	3766.6\\
34	4526.5\\
26	2339.9\\
39	3576.8\\
27	2944\\
32	1979.3\\
26	2230.2\\
29	2371.2\\
23	3394.2\\
26	2460.3\\
14	819.3\\
33	2114.2\\
24	1138.2\\
20	1670.1\\
51	2082.2\\
35	2930.9\\
29	1291.7\\
46	1048.3\\
44	3495.4\\
45	3370.4\\
29	1008.4\\
28	1803.7\\
77	3107.9\\
70	4018\\
70	3561.2\\
40	2416.8\\
90	5643.6\\
53	3404.9\\
62	1805.4\\
47	3720.8\\
34	1349.2\\
14	1080.4\\
47	2173.9\\
43	1159.5\\
20	853.89\\
28	1085.9\\
80	3216.6\\
25	565.84\\
36	1737.6\\
27	2199.1\\
82	4395.5\\
31	1042.3\\
39	1399.9\\
23	2055.7\\
27	2813.1\\
37	2136.5\\
129	8482.3\\
63	4366\\
64	2133.9\\
91	9433.2\\
27	1334.3\\
26	1490.8\\
17	874.18\\
33	1519.9\\
30	3156.6\\
33	1462.3\\
182	14035\\
26	1179.1\\
48	2761.6\\
46	1846.3\\
10	406.6\\
25	1164.1\\
26	798.26\\
206	25015\\
37	1378\\
43	1693.9\\
51	1922.8\\
57	6608.5\\
36	1762.9\\
38	1565.4\\
28	1297.7\\
19	1029.9\\
36	1717\\
45	3068.7\\
53	3669.8\\
48	3616.4\\
64	3205.7\\
65	3143.2\\
45	1277.3\\
54	3009\\
16	445.95\\
24	1070\\
127	8780.1\\
48	1786.1\\
64	2695\\
42	2634.2\\
42	1305.4\\
17	528.1\\
48	2669.7\\
42	2166.4\\
20	861.71\\
34	1148\\
47	1060.4\\
42	2893.9\\
50	2384.2\\
87	6774.2\\
64	4109.9\\
32	1801.8\\
39	1213.2\\
35	2525.2\\
18	1073.3\\
28	885.56\\
27	936.22\\
37	2159.3\\
30	2474.2\\
89	11267\\
22	1890.1\\
28	1628.2\\
25	1350.5\\
15	674.96\\
41	3477.1\\
32	977.09\\
34	1978.1\\
16	794.51\\
88	5767.3\\
39	1044.7\\
140	28786\\
39	1329.3\\
49	1848.5\\
72	2550.9\\
37	1642.1\\
41	3640.2\\
283	1.4502e+05\\
31	2918.5\\
99	18596\\
103	23894\\
39	3560.7\\
35	1471.1\\
98	13894\\
45	4542.5\\
73	8678.5\\
34	2319.7\\
56	7695.7\\
39	3045\\
37	2963.1\\
55	2168.4\\
25	769.54\\
29	1128.6\\
42	1182.2\\
35	4187.3\\
26	814.87\\
45	8433.3\\
56	6286.1\\
39	1581.3\\
62	3285.4\\
94	12909\\
112	24406\\
25	812.74\\
57	2477.9\\
50	2461\\
47	7383.3\\
96	6935.8\\
54	6360.8\\
28	2448\\
38	3985.4\\
47	3590.2\\
35	2047.1\\
45	1774.6\\
50	3648.2\\
43	1063.4\\
29	2218.6\\
42	1832.3\\
71	20091\\
68	6933.4\\
149	55336\\
67	5681.6\\
51	2996.8\\
34	1885.7\\
120	14846\\
43	7071.1\\
117	11205\\
39	1379.3\\
72	8349.7\\
40	3487\\
43	2486.6\\
32	2683.1\\
49	9046.7\\
41	4146.8\\
98	5167.5\\
96	12370\\
45	1772.5\\
129	34962\\
43	2104.4\\
93	7894.9\\
41	3312\\
88	9889.6\\
73	6025.7\\
108	15357\\
52	7620.9\\
18	496.34\\
36	1014.7\\
50	5236.1\\
37	2301.1\\
83	14638\\
40	2289.8\\
26	2752.2\\
40	1812.8\\
49	3890.7\\
27	1453.5\\
41	1484.1\\
36	956.35\\
52	5115.2\\
37	1185.1\\
68	7274.1\\
43	2802.2\\
132	13770\\
73	2948.8\\
46	4597.5\\
47	1315.5\\
44	1727.8\\
41	4974\\
30	1539.1\\
25	3085.3\\
24	1945\\
85	8178.4\\
44	3082.2\\
29	2014.3\\
34	1702\\
101	14402\\
18	1546.9\\
21	1451.7\\
111	16608\\
95	4871.3\\
23	1064.7\\
31	626.98\\
84	4122.2\\
23	782.1\\
55	5165.1\\
32	1213.3\\
27	942.64\\
48	1047.6\\
18	380.73\\
28	2074.6\\
40	1214.5\\
30	1300.5\\
13	327.71\\
33	1553.2\\
14	586.67\\
52	2391.9\\
65	7560.2\\
49	2542\\
35	1430.9\\
109	7366.4\\
26	1003\\
41	1314.3\\
39	1246.9\\
61	6105.4\\
40	1502.1\\
36	2054.8\\
21	658.2\\
80	4121.9\\
36	1344.5\\
36	1962.7\\
38	1042.5\\
42	2105.4\\
27	764.55\\
46	2326.1\\
23	687.64\\
34	1748.5\\
200	38147\\
37	2609.8\\
16	664.69\\
50	2466.9\\
19	703.39\\
41	2962.9\\
25	836.9\\
39	2082.6\\
45	3028\\
83	7391.4\\
27	559.93\\
67	2602\\
46	3083.7\\
35	1662.8\\
44	2449.7\\
30	1978\\
29	1123\\
23	563.56\\
43	1901.1\\
28	794.77\\
24	1185.6\\
56	3170.1\\
33	981.96\\
13	432.13\\
41	3034.1\\
34	950.29\\
14	464.81\\
83	8575.6\\
33	1717.2\\
35	1656.1\\
29	1141.1\\
54	3782.7\\
41	1362.9\\
19	726.37\\
34	2083.8\\
19	709.25\\
38	2121.3\\
43	1743.7\\
14	663.57\\
17	745.77\\
12	473.21\\
};
\addlegendentry{\ExactScatterLegend}

\addplot[only marks, mark=*, mark size=1.2748pt, color=black, fill=black, opacity=0.60, draw=none, mark options={draw=none,line width=0pt}] table[row sep=crcr]{%
x	y\\
34	1787\\
29	2052\\
96	39026\\
31	939\\
27	839\\
70	9435\\
24	756\\
37	2292\\
42	3353\\
16	509\\
29	1198\\
59	4061\\
69	4632\\
28	987\\
11	377\\
20	825\\
45	6367\\
27	1370\\
22	689\\
27	1107\\
44	4679\\
16	622\\
39	5236\\
19	558\\
33	1997\\
37	1291\\
33	1077\\
18	485\\
29	1640\\
22	758\\
33	1231\\
22	643\\
49	9267\\
25	1115\\
43	11048\\
26	1875\\
31	2762\\
26	927\\
18	665\\
13	140\\
32	2698\\
38	2211\\
32	775\\
29	1173\\
35	1719\\
35	3078\\
135	41095\\
27	2591\\
42	3900\\
39	3847\\
51	5607\\
86	11830\\
15	918\\
39	3169\\
42	3772\\
43	2489\\
58	4861\\
97	21264\\
50	3266\\
20	736\\
28	686\\
22	1102\\
47	4024\\
154	1.2234e+05\\
30	1543\\
10	296\\
37	7252\\
31	1142\\
54	13398\\
76	13899\\
55	3344\\
26	1499\\
33	1121\\
34	1971\\
25	1709\\
31	2961\\
34	3625\\
26	2774\\
39	3551\\
27	1922\\
32	1860\\
26	1781\\
29	1942\\
23	2014\\
26	1636\\
14	530\\
33	1466\\
24	787\\
20	1114\\
51	1692\\
35	2940\\
29	1332\\
46	1054\\
44	3996\\
45	4071\\
29	803\\
28	557\\
77	3928\\
70	3625\\
70	2531\\
40	1723\\
90	8518\\
53	3650\\
62	3213\\
47	8035\\
34	817\\
14	526\\
47	2742\\
43	759\\
20	950\\
28	533\\
80	4654\\
25	1153\\
36	1716\\
27	2299\\
82	3021\\
31	1103\\
39	1587\\
23	810\\
27	2388\\
37	2269\\
129	3241\\
63	3568\\
64	2459\\
91	6356\\
27	696\\
26	1459\\
17	558\\
33	2465\\
30	2046\\
33	1791\\
182	11424\\
26	663\\
48	2633\\
46	2157\\
10	222\\
25	802\\
26	424\\
206	37527\\
37	1069\\
43	1448\\
51	2715\\
57	7724\\
36	1800\\
38	1162\\
28	539\\
19	746\\
36	1003\\
45	3212\\
53	4779\\
48	5264\\
64	2484\\
65	3084\\
45	1597\\
54	3279\\
16	299\\
24	1131\\
127	2601\\
48	1783\\
64	2544\\
42	2057\\
42	1206\\
17	254\\
48	2507\\
42	2016\\
20	816\\
34	1805\\
47	1076\\
42	2516\\
50	2233\\
87	10377\\
64	2072\\
32	1299\\
39	1215\\
35	3735\\
18	866\\
28	626\\
27	253\\
37	2518\\
30	2667\\
89	9128\\
22	1204\\
28	1732\\
25	1513\\
15	873\\
41	3761\\
32	1140\\
34	1426\\
16	352\\
88	7596\\
39	668\\
140	20491\\
39	1086\\
49	1025\\
72	1475\\
37	1233\\
41	3002\\
283	2.0424e+05\\
31	2681\\
99	16428\\
103	32887\\
39	6629\\
35	1076\\
98	6332\\
45	4113\\
73	7877\\
34	1444\\
56	6926\\
39	2045\\
37	3241\\
55	1793\\
25	444\\
29	1020\\
42	1480\\
35	2286\\
26	591\\
45	5458\\
56	7320\\
39	1112\\
62	3063\\
94	13380\\
112	30134\\
25	317\\
57	1896\\
50	2648\\
47	5982\\
96	4539\\
54	4348\\
28	1451\\
38	3579\\
47	3124\\
35	1834\\
45	921\\
50	3106\\
43	719\\
29	1516\\
42	1249\\
71	5868\\
68	5857\\
149	61636\\
67	5463\\
51	1708\\
34	1459\\
120	10119\\
43	6342\\
117	8499\\
39	1041\\
72	6468\\
40	4061\\
43	1854\\
32	2516\\
49	12313\\
41	4009\\
98	4503\\
96	5379\\
45	1383\\
129	23237\\
43	1506\\
93	4228\\
41	2160\\
88	9837\\
73	5327\\
108	12182\\
52	5524\\
18	189\\
36	878\\
50	4430\\
37	1383\\
83	13086\\
40	2202\\
26	1338\\
40	1198\\
49	4702\\
27	1189\\
41	1263\\
36	572\\
52	3610\\
37	656\\
68	6317\\
43	2092\\
132	14984\\
73	3147\\
46	3825\\
47	739\\
44	1624\\
41	5048\\
30	1539\\
25	1574\\
24	2598\\
85	7294\\
44	1753\\
29	1546\\
34	1160\\
101	9599\\
18	1226\\
21	992\\
111	7348\\
95	2640\\
23	617\\
31	310\\
84	1691\\
23	466\\
55	3785\\
32	1352\\
27	544\\
48	517\\
18	119\\
28	1203\\
40	962\\
30	685\\
13	188\\
33	1762\\
14	209\\
52	1836\\
65	8994\\
49	3401\\
35	1176\\
109	4663\\
26	550\\
41	1006\\
39	973\\
61	4953\\
40	1238\\
36	1694\\
21	551\\
80	5029\\
36	1037\\
36	1691\\
38	832\\
42	1192\\
27	590\\
46	1533\\
23	486\\
34	1267\\
200	30990\\
37	1240\\
16	413\\
50	2688\\
19	324\\
41	2143\\
25	934\\
39	3741\\
45	2869\\
83	3945\\
27	521\\
67	2908\\
46	2419\\
35	1912\\
44	2676\\
30	2051\\
29	674\\
23	204\\
43	936\\
28	859\\
24	467\\
56	2173\\
33	679\\
13	261\\
41	2629\\
34	588\\
14	216\\
83	10336\\
33	1005\\
35	1380\\
29	687\\
54	3377\\
41	993\\
19	351\\
34	1937\\
19	459\\
38	1126\\
43	1196\\
14	412\\
17	325\\
12	184\\
};
\addlegendentry{\RealScatterLegend}

      % \input{../../../../.tex/20/KS/SizeVsDegreefig5.tex}
      \legend{}
  
      \nextgroupplot[%
        % xmin=10,xmax=283,
        % ymin=100,ymax=1e+06,
      ] % This file was created by matlab2tikz.
%
\definecolor{mycolor1}{rgb}{0.83529,0.36863,0.00000}%
\definecolor{mycolor2}{rgb}{0.00000,0.44700,0.74100}%
\definecolor{mycolor3}{rgb}{0.85000,0.32500,0.09800}%
%
\addplot[only marks, mark=*, mark size=1.2748pt, color=mycolor1, fill=mycolor1, opacity=0.60, draw=none, mark options={draw=none,line width=0pt}] table[row sep=crcr]{%
x	y\\
34	2575.1\\
29	1412\\
96	17027\\
31	2026\\
27	1382.9\\
70	7557.9\\
24	1263.9\\
37	2264.4\\
42	2958.2\\
16	1201.8\\
29	1459.6\\
59	3567.1\\
69	3369.2\\
28	1645.4\\
11	346.55\\
20	1164.6\\
45	3948.2\\
27	1504.4\\
22	1035.2\\
27	1412.4\\
44	5048.7\\
16	1466.5\\
39	4261.8\\
19	1172.7\\
33	2277.8\\
37	2560.6\\
33	1408.3\\
18	584.11\\
29	2494.8\\
22	1042\\
33	1424.6\\
22	911.84\\
49	8130\\
25	1485.4\\
43	3912.7\\
26	2501.4\\
31	3334.2\\
26	1313.2\\
18	896.76\\
13	584.85\\
32	4198.3\\
38	2470.3\\
32	3476\\
29	1789.5\\
35	1558.5\\
35	3271.8\\
135	20059\\
27	4523.3\\
42	4266.3\\
39	4230.4\\
51	7348.6\\
86	10217\\
15	954.08\\
39	4544.5\\
42	3294.2\\
43	3336.5\\
58	5783.3\\
97	24111\\
50	2432.9\\
20	1428\\
28	1186.2\\
22	1118.3\\
47	3307.6\\
154	60390\\
30	1785.5\\
10	509.76\\
37	8330.3\\
31	1938.8\\
54	13828\\
76	11797\\
55	4758\\
26	2727.6\\
33	1700.9\\
34	2034.3\\
25	1657.9\\
31	3581.3\\
34	4486.2\\
26	2277.2\\
39	3716\\
27	2671.7\\
32	2050.2\\
26	2209\\
29	2413\\
23	3300.2\\
26	2326\\
14	910.61\\
33	2154.6\\
24	1073.6\\
20	1738.8\\
51	2077.4\\
35	3047.3\\
29	1402.2\\
46	964.76\\
44	3697.2\\
45	3441.3\\
29	964.33\\
28	1766.5\\
77	3112.2\\
70	3743.2\\
70	3331.2\\
40	2313.4\\
90	5315.9\\
53	3656.8\\
62	1751.3\\
47	3538.1\\
34	1388\\
14	1134.3\\
47	2177.9\\
43	1181.5\\
20	754.18\\
28	1106.4\\
80	3197.5\\
25	539.62\\
36	1720.5\\
27	2156.1\\
82	4326.6\\
31	1088.7\\
39	1357.8\\
23	2138.8\\
27	2858.6\\
37	2219.7\\
129	8284.4\\
63	4431.2\\
64	2037.7\\
91	10009\\
27	1260.9\\
26	1430.9\\
17	875.69\\
33	1591.5\\
30	3101.3\\
33	1458.6\\
182	14317\\
26	1185.6\\
48	2951\\
46	1732.6\\
10	387.65\\
25	1129.3\\
26	836.72\\
206	24860\\
37	1410.7\\
43	1556.8\\
51	2014.4\\
57	6608\\
36	1830.3\\
38	1583.5\\
28	1341.2\\
19	1156.6\\
36	1665.3\\
45	3222.1\\
53	3605.6\\
48	3602.5\\
64	3269.5\\
65	2931.7\\
45	1303.2\\
54	2785.5\\
16	447.12\\
24	1100.1\\
127	8673.4\\
48	1755\\
64	2506.5\\
42	2665.4\\
42	1358\\
17	607.62\\
48	2785.6\\
42	2175\\
20	918.06\\
34	1163.3\\
47	1037.8\\
42	2906.3\\
50	2447.8\\
87	6280\\
64	3963.7\\
32	1807\\
39	1184.9\\
35	2484.6\\
18	1087.3\\
28	925.94\\
27	960.91\\
37	2143.9\\
30	2581.3\\
89	11447\\
22	1773.2\\
28	1632.5\\
25	1260.6\\
15	651.77\\
41	3414.1\\
32	987.86\\
34	1942.3\\
16	850.81\\
88	6024\\
39	992.97\\
140	28024\\
39	1253.4\\
49	1839.8\\
72	2449.5\\
37	1734.5\\
41	3646.9\\
283	1.4849e+05\\
31	2972.4\\
99	19340\\
103	23817\\
39	3662.2\\
35	1559.6\\
98	14258\\
45	4398.3\\
73	8778.5\\
34	2386.2\\
56	7844.2\\
39	2975.8\\
37	3088.2\\
55	2278.7\\
25	749.27\\
29	1093.5\\
42	1224.5\\
35	4115.1\\
26	791.28\\
45	8471\\
56	6165\\
39	1606.6\\
62	3123\\
94	12539\\
112	24734\\
25	807.97\\
57	2667.2\\
50	2372.1\\
47	6685.6\\
96	6912.2\\
54	6489\\
28	2634.3\\
38	3855\\
47	3730.9\\
35	2047\\
45	1635.5\\
50	3648.1\\
43	1081.5\\
29	2191.2\\
42	1917.2\\
71	19877\\
68	6587.1\\
149	53091\\
67	5654.6\\
51	2934.7\\
34	1889.1\\
120	14412\\
43	7413.4\\
117	11148\\
39	1361\\
72	8497.1\\
40	4143.7\\
43	2477.5\\
32	2853.7\\
49	8768.1\\
41	4066.7\\
98	5220.8\\
96	12375\\
45	1693.4\\
129	34629\\
43	2028.1\\
93	8219.8\\
41	3435.5\\
88	9754.2\\
73	5996.3\\
108	15579\\
52	7853.5\\
18	464.22\\
36	1139.6\\
50	5590.6\\
37	2364\\
83	15229\\
40	2207.2\\
26	2506.4\\
40	1778.6\\
49	4042.1\\
27	1459.3\\
41	1402.6\\
36	990.8\\
52	5078.2\\
37	1150.8\\
68	6958.7\\
43	2888.3\\
132	14083\\
73	3040.9\\
46	4455.7\\
47	1220\\
44	1666.2\\
41	4740.1\\
30	1511.3\\
25	3009.3\\
24	1960.5\\
85	8343.8\\
44	3057.1\\
29	1783.9\\
34	1754.7\\
101	14322\\
18	1665.1\\
21	1396.3\\
111	16269\\
95	4741.4\\
23	1018.1\\
31	622.98\\
84	3931.6\\
23	843.66\\
55	4936.1\\
32	1210.6\\
27	906.71\\
48	1017\\
18	397.19\\
28	2000.6\\
40	1205.1\\
30	1334.7\\
13	332.05\\
33	1628.9\\
14	591.57\\
52	2431.5\\
65	7872.8\\
49	2654.1\\
35	1426.4\\
109	6961.9\\
26	980.18\\
41	1349.4\\
39	1246.5\\
61	6024.2\\
40	1503.6\\
36	2082.7\\
21	647.05\\
80	4171.2\\
36	1402.5\\
36	2026.7\\
38	1119.2\\
42	2056.1\\
27	750.24\\
46	2342.5\\
23	709.43\\
34	1804.2\\
200	39277\\
37	2532.9\\
16	619.21\\
50	2586.9\\
19	686.7\\
41	3048.8\\
25	825.51\\
39	2292.8\\
45	2841.4\\
83	7099.8\\
27	604.92\\
67	2542\\
46	3104.8\\
35	1609.9\\
44	2732.6\\
30	1790.2\\
29	1079.6\\
23	579.43\\
43	2081.9\\
28	804.3\\
24	1083.1\\
56	2959.9\\
33	974.5\\
13	418.55\\
41	3134.3\\
34	1015.9\\
14	456.95\\
83	9110\\
33	1791.7\\
35	1614.7\\
29	977.23\\
54	3730.7\\
41	1392.5\\
19	753.9\\
34	1993\\
19	726.16\\
38	2267.4\\
43	1847\\
14	695.77\\
17	863.3\\
12	457.58\\
};
\addlegendentry{\ApproxScatterLegend}

\addplot[only marks, mark=*, mark size=1.2748pt, color=black, fill=black, opacity=0.60, draw=none, mark options={draw=none,line width=0pt}] table[row sep=crcr]{%
x	y\\
34	1787\\
29	2052\\
96	39026\\
31	939\\
27	839\\
70	9435\\
24	756\\
37	2292\\
42	3353\\
16	509\\
29	1198\\
59	4061\\
69	4632\\
28	987\\
11	377\\
20	825\\
45	6367\\
27	1370\\
22	689\\
27	1107\\
44	4679\\
16	622\\
39	5236\\
19	558\\
33	1997\\
37	1291\\
33	1077\\
18	485\\
29	1640\\
22	758\\
33	1231\\
22	643\\
49	9267\\
25	1115\\
43	11048\\
26	1875\\
31	2762\\
26	927\\
18	665\\
13	140\\
32	2698\\
38	2211\\
32	775\\
29	1173\\
35	1719\\
35	3078\\
135	41095\\
27	2591\\
42	3900\\
39	3847\\
51	5607\\
86	11830\\
15	918\\
39	3169\\
42	3772\\
43	2489\\
58	4861\\
97	21264\\
50	3266\\
20	736\\
28	686\\
22	1102\\
47	4024\\
154	1.2234e+05\\
30	1543\\
10	296\\
37	7252\\
31	1142\\
54	13398\\
76	13899\\
55	3344\\
26	1499\\
33	1121\\
34	1971\\
25	1709\\
31	2961\\
34	3625\\
26	2774\\
39	3551\\
27	1922\\
32	1860\\
26	1781\\
29	1942\\
23	2014\\
26	1636\\
14	530\\
33	1466\\
24	787\\
20	1114\\
51	1692\\
35	2940\\
29	1332\\
46	1054\\
44	3996\\
45	4071\\
29	803\\
28	557\\
77	3928\\
70	3625\\
70	2531\\
40	1723\\
90	8518\\
53	3650\\
62	3213\\
47	8035\\
34	817\\
14	526\\
47	2742\\
43	759\\
20	950\\
28	533\\
80	4654\\
25	1153\\
36	1716\\
27	2299\\
82	3021\\
31	1103\\
39	1587\\
23	810\\
27	2388\\
37	2269\\
129	3241\\
63	3568\\
64	2459\\
91	6356\\
27	696\\
26	1459\\
17	558\\
33	2465\\
30	2046\\
33	1791\\
182	11424\\
26	663\\
48	2633\\
46	2157\\
10	222\\
25	802\\
26	424\\
206	37527\\
37	1069\\
43	1448\\
51	2715\\
57	7724\\
36	1800\\
38	1162\\
28	539\\
19	746\\
36	1003\\
45	3212\\
53	4779\\
48	5264\\
64	2484\\
65	3084\\
45	1597\\
54	3279\\
16	299\\
24	1131\\
127	2601\\
48	1783\\
64	2544\\
42	2057\\
42	1206\\
17	254\\
48	2507\\
42	2016\\
20	816\\
34	1805\\
47	1076\\
42	2516\\
50	2233\\
87	10377\\
64	2072\\
32	1299\\
39	1215\\
35	3735\\
18	866\\
28	626\\
27	253\\
37	2518\\
30	2667\\
89	9128\\
22	1204\\
28	1732\\
25	1513\\
15	873\\
41	3761\\
32	1140\\
34	1426\\
16	352\\
88	7596\\
39	668\\
140	20491\\
39	1086\\
49	1025\\
72	1475\\
37	1233\\
41	3002\\
283	2.0424e+05\\
31	2681\\
99	16428\\
103	32887\\
39	6629\\
35	1076\\
98	6332\\
45	4113\\
73	7877\\
34	1444\\
56	6926\\
39	2045\\
37	3241\\
55	1793\\
25	444\\
29	1020\\
42	1480\\
35	2286\\
26	591\\
45	5458\\
56	7320\\
39	1112\\
62	3063\\
94	13380\\
112	30134\\
25	317\\
57	1896\\
50	2648\\
47	5982\\
96	4539\\
54	4348\\
28	1451\\
38	3579\\
47	3124\\
35	1834\\
45	921\\
50	3106\\
43	719\\
29	1516\\
42	1249\\
71	5868\\
68	5857\\
149	61636\\
67	5463\\
51	1708\\
34	1459\\
120	10119\\
43	6342\\
117	8499\\
39	1041\\
72	6468\\
40	4061\\
43	1854\\
32	2516\\
49	12313\\
41	4009\\
98	4503\\
96	5379\\
45	1383\\
129	23237\\
43	1506\\
93	4228\\
41	2160\\
88	9837\\
73	5327\\
108	12182\\
52	5524\\
18	189\\
36	878\\
50	4430\\
37	1383\\
83	13086\\
40	2202\\
26	1338\\
40	1198\\
49	4702\\
27	1189\\
41	1263\\
36	572\\
52	3610\\
37	656\\
68	6317\\
43	2092\\
132	14984\\
73	3147\\
46	3825\\
47	739\\
44	1624\\
41	5048\\
30	1539\\
25	1574\\
24	2598\\
85	7294\\
44	1753\\
29	1546\\
34	1160\\
101	9599\\
18	1226\\
21	992\\
111	7348\\
95	2640\\
23	617\\
31	310\\
84	1691\\
23	466\\
55	3785\\
32	1352\\
27	544\\
48	517\\
18	119\\
28	1203\\
40	962\\
30	685\\
13	188\\
33	1762\\
14	209\\
52	1836\\
65	8994\\
49	3401\\
35	1176\\
109	4663\\
26	550\\
41	1006\\
39	973\\
61	4953\\
40	1238\\
36	1694\\
21	551\\
80	5029\\
36	1037\\
36	1691\\
38	832\\
42	1192\\
27	590\\
46	1533\\
23	486\\
34	1267\\
200	30990\\
37	1240\\
16	413\\
50	2688\\
19	324\\
41	2143\\
25	934\\
39	3741\\
45	2869\\
83	3945\\
27	521\\
67	2908\\
46	2419\\
35	1912\\
44	2676\\
30	2051\\
29	674\\
23	204\\
43	936\\
28	859\\
24	467\\
56	2173\\
33	679\\
13	261\\
41	2629\\
34	588\\
14	216\\
83	10336\\
33	1005\\
35	1380\\
29	687\\
54	3377\\
41	993\\
19	351\\
34	1937\\
19	459\\
38	1126\\
43	1196\\
14	412\\
17	325\\
12	184\\
};
\addlegendentry{\RealScatterLegend}

      % \input{../../../../.tex/20/KS/SizeVsDegreefig6.tex}
      \legend{}
    \end{groupplot}
    
    \begin{groupplot}[
      group style={
        group name=Row3,
        plotGroup
      },
      plotRow,
      row4Keys,
      plotLegendNW,
      %
      xmin=10,xmax=283,
      ymin=100,ymax=2.0424e+05,
    ]
      \nextgroupplot[%
        at={($(Row2 c1r1.south west)-(0,\yPlotSep em)-(0,\plotHeight cm)$)},
        % xmin=10,xmax=283,
        % ymin=100,ymax=1.537e+05,
      ] % This file was created by matlab2tikz.
%
\definecolor{mycolor1}{rgb}{0.00000,0.44706,0.69804}%
\definecolor{mycolor2}{rgb}{0.00000,0.44700,0.74100}%
\definecolor{mycolor3}{rgb}{0.83529,0.36863,0.00000}%
\definecolor{mycolor4}{rgb}{0.85000,0.32500,0.09800}%
%
\addplot[only marks, mark=*, mark size=1.2748pt, color=mycolor1, fill=mycolor1, opacity=0.70, draw=none, mark options={draw=none,line width=0pt}] table[row sep=crcr]{%
x	y\\
34	2659.6\\
29	1414.7\\
96	17087\\
31	1865.1\\
27	1245.2\\
70	7969.5\\
24	1175.5\\
37	2263.9\\
42	2826.4\\
16	1123.1\\
29	1578.5\\
59	3441.5\\
69	3268.8\\
28	1565.9\\
11	337.04\\
20	1072.5\\
45	3829.7\\
27	1353.6\\
22	1064.3\\
27	1473.9\\
44	4953.6\\
16	1632.6\\
39	4333.8\\
19	1155.6\\
33	2083.2\\
37	2597.1\\
33	1420.1\\
18	615.12\\
29	2511.1\\
22	1038.8\\
33	1466.1\\
22	952.09\\
49	8841.6\\
25	1457.4\\
43	3953.8\\
26	2345.2\\
31	3343.1\\
26	1332.8\\
18	886.52\\
13	513.86\\
32	4019.9\\
38	2379.3\\
32	3443.2\\
29	1686.3\\
35	1706\\
35	3160.2\\
135	20467\\
27	4701.4\\
42	3880.7\\
39	4923.8\\
51	6896.2\\
86	10357\\
15	957.67\\
39	5028.4\\
42	3470.5\\
43	3130.2\\
58	5795.7\\
97	24183\\
50	2691.1\\
20	1456.9\\
28	1244.6\\
22	1072.2\\
47	3224.4\\
154	63393\\
30	1815.2\\
10	476.82\\
37	7531.1\\
31	1924\\
54	13514\\
76	11305\\
55	4964.6\\
26	2661.5\\
33	1634.4\\
34	2158.6\\
25	1589\\
31	3864.4\\
34	4806.9\\
26	2342.6\\
39	3443\\
27	2951.2\\
32	1946.3\\
26	2251.5\\
29	2140.4\\
23	3106\\
26	2150.7\\
14	854.42\\
33	2069.5\\
24	1114.1\\
20	1657\\
51	2229.3\\
35	2878.6\\
29	1300.8\\
46	988.34\\
44	3612\\
45	3445.2\\
29	892.64\\
28	1743.1\\
77	3217.4\\
70	3997.3\\
70	3338.8\\
40	2569\\
90	5574.9\\
53	3745.5\\
62	1604.7\\
47	3610.6\\
34	1350\\
14	959.1\\
47	2071.8\\
43	1060.6\\
20	807.84\\
28	1062.6\\
80	2845.7\\
25	552.12\\
36	1804.4\\
27	2284.9\\
82	4078\\
31	1157.3\\
39	1404.3\\
23	2079.3\\
27	2918.1\\
37	2236.7\\
129	8341\\
63	4300.2\\
64	1997.6\\
91	9282.4\\
27	1217.5\\
26	1500.1\\
17	771.8\\
33	1617.2\\
30	3277.2\\
33	1427.7\\
182	14664\\
26	1157\\
48	2931.9\\
46	1887\\
10	387.38\\
25	1142.2\\
26	784.64\\
206	26843\\
37	1387.9\\
43	1657.4\\
51	1914.5\\
57	6472.8\\
36	1809.5\\
38	1616.6\\
28	1303.5\\
19	1126.3\\
36	1644.5\\
45	3206.5\\
53	3477.4\\
48	3680.4\\
64	3226.4\\
65	2938.8\\
45	1225.4\\
54	2931.1\\
16	459.7\\
24	1120.2\\
127	8740.3\\
48	1703.2\\
64	2499.5\\
42	2597.3\\
42	1308.6\\
17	562.74\\
48	2766.7\\
42	2150.6\\
20	877.68\\
34	1199.1\\
47	923.14\\
42	2929.1\\
50	2505.7\\
87	6689.2\\
64	3904.7\\
32	1917.4\\
39	1177.3\\
35	2435.8\\
18	1026.7\\
28	945.7\\
27	903.42\\
37	2155.5\\
30	2467.3\\
89	11830\\
22	1889.6\\
28	1565.7\\
25	1328.4\\
15	625.19\\
41	3527.1\\
32	963.33\\
34	1873.5\\
16	746.94\\
88	5937.8\\
39	980.57\\
140	27483\\
39	1259\\
49	1859.2\\
72	2264.4\\
37	1552.4\\
41	3720.2\\
283	1.5097e+05\\
31	2803.7\\
99	19380\\
103	24503\\
39	3629.3\\
35	1443\\
98	12986\\
45	4464\\
73	8456.1\\
34	2322.8\\
56	7870.4\\
39	2972.8\\
37	2902.4\\
55	2161.9\\
25	773.23\\
29	1052.3\\
42	1188.1\\
35	3915.8\\
26	762.86\\
45	8590\\
56	6323\\
39	1683.5\\
62	3118.3\\
94	12804\\
112	25043\\
25	732.82\\
57	2525.8\\
50	2583.8\\
47	7345.3\\
96	6934.2\\
54	6333.5\\
28	2464.6\\
38	4195.5\\
47	3331.8\\
35	1948.8\\
45	1643.7\\
50	3800.2\\
43	1046.8\\
29	2100.3\\
42	1789.1\\
71	19782\\
68	6076.5\\
149	52274\\
67	5541.5\\
51	2942.7\\
34	1918.1\\
120	14743\\
43	6853.5\\
117	11166\\
39	1444.8\\
72	9128.1\\
40	3718.5\\
43	2233.1\\
32	2559.9\\
49	9347.7\\
41	3701.8\\
98	5270.9\\
96	12850\\
45	1718.5\\
129	34983\\
43	2086.1\\
93	7942.5\\
41	3493.8\\
88	10220\\
73	5834.7\\
108	16419\\
52	7730.4\\
18	485.3\\
36	994.07\\
50	5230.4\\
37	2378.4\\
83	14641\\
40	2223.8\\
26	2719.7\\
40	1723.2\\
49	3680.1\\
27	1455.6\\
41	1334.6\\
36	939.06\\
52	5037.1\\
37	1179.9\\
68	7085.9\\
43	2901.8\\
132	13823\\
73	2911.6\\
46	4779.2\\
47	1257.5\\
44	1735.1\\
41	4749.9\\
30	1499.2\\
25	2922.5\\
24	2038.4\\
85	8202.1\\
44	2580.5\\
29	1710.2\\
34	1545.6\\
101	15542\\
18	1635.1\\
21	1281.4\\
111	16157\\
95	4957.6\\
23	1189.4\\
31	659.22\\
84	3869.4\\
23	824.67\\
55	5344.2\\
32	1205.8\\
27	918.03\\
48	970.54\\
18	386.22\\
28	2259\\
40	1278.4\\
30	1301.7\\
13	344.21\\
33	1407.5\\
14	545.31\\
52	2329.8\\
65	7552.6\\
49	2581.5\\
35	1453.9\\
109	7506.2\\
26	902.75\\
41	1307.5\\
39	1377\\
61	5863.5\\
40	1476.7\\
36	2244.1\\
21	590.21\\
80	4183.2\\
36	1299.4\\
36	2023.4\\
38	1040\\
42	2101.6\\
27	763.39\\
46	2209.8\\
23	601.01\\
34	1740.8\\
200	38200\\
37	2321\\
16	574.46\\
50	2231.8\\
19	656.56\\
41	2758.5\\
25	812.13\\
39	1907.4\\
45	3019.7\\
83	6800.9\\
27	536.65\\
67	2545.7\\
46	3015.3\\
35	1529.3\\
44	2515\\
30	1800.7\\
29	1169.9\\
23	508.63\\
43	2142.2\\
28	705.01\\
24	1129.4\\
56	3031.9\\
33	907.22\\
13	397.3\\
41	3038.9\\
34	955.97\\
14	415.81\\
83	8462.4\\
33	1665.3\\
35	1574.2\\
29	973.7\\
54	3766.7\\
41	1428.9\\
19	678.16\\
34	2078.4\\
19	705.37\\
38	2321.5\\
43	1680\\
14	671.47\\
17	781.43\\
12	439.35\\
};
\addlegendentry{\ExactScatterLegend}

\addplot[only marks, mark=*, mark size=1.2748pt, color=mycolor3, fill=mycolor3, opacity=0.70, draw=none, mark options={draw=none,line width=0pt}] table[row sep=crcr]{%
x	y\\
34	2752.1\\
29	1447.9\\
96	17899\\
31	1897.2\\
27	1312.4\\
70	7519.4\\
24	1293.2\\
37	2217\\
42	2934.9\\
16	1196.2\\
29	1455.8\\
59	3672.2\\
69	3506.9\\
28	1430.5\\
11	328.45\\
20	1179.9\\
45	3924.6\\
27	1565.1\\
22	966.11\\
27	1377.7\\
44	4878.9\\
16	1535.8\\
39	4304.3\\
19	1016.5\\
33	2032.1\\
37	2366.7\\
33	1308.8\\
18	610.03\\
29	2329.7\\
22	1115.1\\
33	1372.5\\
22	924.94\\
49	8161.7\\
25	1350.4\\
43	3843.7\\
26	2355.5\\
31	2991.9\\
26	1322.2\\
18	910.13\\
13	524.25\\
32	4223.2\\
38	2614.1\\
32	3319.3\\
29	1690.5\\
35	1683.3\\
35	3143\\
135	19595\\
27	4534.5\\
42	4309.9\\
39	4435.6\\
51	6982.8\\
86	11144\\
15	800.43\\
39	4834.6\\
42	3437.7\\
43	3295.7\\
58	6074.7\\
97	23656\\
50	2655.5\\
20	1462.3\\
28	1227.1\\
22	1114.8\\
47	3311.4\\
154	62486\\
30	1873.3\\
10	502.74\\
37	8467.2\\
31	1883.7\\
54	13540\\
76	11832\\
55	4800.7\\
26	2723.5\\
33	1467.9\\
34	2231.5\\
25	1641.5\\
31	3774.1\\
34	4704.3\\
26	2337.3\\
39	3648.4\\
27	2584.9\\
32	2201.2\\
26	2058.9\\
29	2475.8\\
23	3275.9\\
26	2235.6\\
14	861.46\\
33	2109.5\\
24	1049.4\\
20	1658\\
51	2223.7\\
35	3137.6\\
29	1309.1\\
46	949.6\\
44	3557\\
45	3264.9\\
29	984.36\\
28	1757.2\\
77	3171\\
70	4100.3\\
70	3599.2\\
40	2546.3\\
90	5463.2\\
53	3769\\
62	1775.1\\
47	3553.9\\
34	1404.2\\
14	1081\\
47	2082.7\\
43	1081.9\\
20	788.24\\
28	1069.3\\
80	3224\\
25	568.75\\
36	1872.2\\
27	2275.1\\
82	4371.8\\
31	1083.1\\
39	1412.4\\
23	2279.4\\
27	2740.9\\
37	2116.4\\
129	8873.1\\
63	4286.9\\
64	1917.8\\
91	9929.6\\
27	1129.8\\
26	1628.7\\
17	833.66\\
33	1512.2\\
30	3272.8\\
33	1375.8\\
182	14819\\
26	1213.5\\
48	2742.3\\
46	1898.8\\
10	331.16\\
25	1145.2\\
26	744.56\\
206	25410\\
37	1287.7\\
43	1537.3\\
51	2011.3\\
57	6722.5\\
36	1672.8\\
38	1640.8\\
28	1301.3\\
19	1091.1\\
36	1630.6\\
45	3199\\
53	3491.8\\
48	3481.2\\
64	3039.5\\
65	3034.6\\
45	1160.2\\
54	2692.7\\
16	402.49\\
24	1006.1\\
127	8342.6\\
48	1814.7\\
64	2773.7\\
42	2547.1\\
42	1259.1\\
17	559.12\\
48	2914.9\\
42	2224.8\\
20	841.27\\
34	1128.3\\
47	1064.6\\
42	2961.1\\
50	2416.6\\
87	6222.2\\
64	3661.6\\
32	1823.6\\
39	1151.3\\
35	2425.1\\
18	1039.4\\
28	994.52\\
27	900.38\\
37	2226.4\\
30	2227.4\\
89	11814\\
22	1931.8\\
28	1589.8\\
25	1245\\
15	590.65\\
41	3560.7\\
32	948.29\\
34	1979.5\\
16	789.6\\
88	5367.1\\
39	1047.6\\
140	26639\\
39	1240.2\\
49	1706.1\\
72	2226.7\\
37	1609.9\\
41	3304\\
283	1.537e+05\\
31	3008.5\\
99	18581\\
103	23629\\
39	3638.2\\
35	1512.7\\
98	13472\\
45	4614.8\\
73	8967.7\\
34	2178.8\\
56	7761.4\\
39	3241.2\\
37	3173.6\\
55	2166.4\\
25	776.54\\
29	1104.3\\
42	1198.7\\
35	3878.8\\
26	762.57\\
45	8279.4\\
56	6466.5\\
39	1541.7\\
62	3102.1\\
94	11988\\
112	23458\\
25	725.17\\
57	2347.9\\
50	2470.1\\
47	7040.9\\
96	6892.2\\
54	6397.8\\
28	2418.8\\
38	3666\\
47	3740.2\\
35	1855.5\\
45	1593.2\\
50	3933.1\\
43	1021.4\\
29	2149.5\\
42	1867.7\\
71	21595\\
68	6832.4\\
149	51885\\
67	6095.8\\
51	2849.5\\
34	1741.1\\
120	14166\\
43	7024.9\\
117	10766\\
39	1385.5\\
72	8517.9\\
40	3570.3\\
43	2546.7\\
32	2440\\
49	8717\\
41	4216.1\\
98	4851.4\\
96	13082\\
45	1790.2\\
129	34070\\
43	2064.6\\
93	8244.5\\
41	3517.4\\
88	10215\\
73	6049.7\\
108	15669\\
52	7762.1\\
18	459.6\\
36	1047.3\\
50	5559\\
37	2276.1\\
83	14453\\
40	2357.3\\
26	2746.3\\
40	1747.4\\
49	3950.3\\
27	1417.2\\
41	1374.5\\
36	872.47\\
52	4928.5\\
37	1076.9\\
68	7268.5\\
43	2696.2\\
132	14335\\
73	2855.8\\
46	4938.5\\
47	1346\\
44	1660.2\\
41	5069.6\\
30	1525.9\\
25	2998\\
24	1825.6\\
85	8306.6\\
44	2949.4\\
29	1714.7\\
34	1536.1\\
101	15121\\
18	1458.8\\
21	1345.4\\
111	16738\\
95	4473.2\\
23	1047\\
31	655.26\\
84	4030.8\\
23	836.93\\
55	4995.2\\
32	1154.3\\
27	954.76\\
48	990.47\\
18	348.75\\
28	1897.4\\
40	1229\\
30	1323\\
13	334.01\\
33	1605.1\\
14	590.46\\
52	2527.6\\
65	7783.7\\
49	2447.2\\
35	1309.6\\
109	6976.6\\
26	974.29\\
41	1320.4\\
39	1273.1\\
61	5949.8\\
40	1394.4\\
36	2038.2\\
21	645.28\\
80	4376.3\\
36	1407.8\\
36	1979.2\\
38	1088.9\\
42	1970.8\\
27	765.93\\
46	2331.7\\
23	663.1\\
34	1777\\
200	40388\\
37	2389.4\\
16	607.66\\
50	2591.5\\
19	646.53\\
41	2830\\
25	942.11\\
39	2154.3\\
45	3049.6\\
83	7540.9\\
27	564.42\\
67	2554\\
46	2997.2\\
35	1591.2\\
44	2521.3\\
30	1856.8\\
29	1130.1\\
23	479.02\\
43	1924.8\\
28	732.66\\
24	1081\\
56	3305\\
33	1003.7\\
13	438.81\\
41	3042.6\\
34	955.6\\
14	441.88\\
83	8253.1\\
33	1549.1\\
35	1503.1\\
29	1021.3\\
54	3793.1\\
41	1286.9\\
19	676.49\\
34	2090.4\\
19	723.93\\
38	2349.1\\
43	1695.8\\
14	621.64\\
17	768.51\\
12	458.88\\
};
\addlegendentry{\ApproxScatterLegend}

      % \input{../../../../.tex/20/KS/SizeVsDegreefig7.tex}
      \legend{}
  
      \nextgroupplot[%
        % xmin=10,xmax=283,
        % ymin=100,ymax=1e+06,
      ] % This file was created by matlab2tikz.
%
\definecolor{mycolor1}{rgb}{0.00000,0.44706,0.69804}%
\definecolor{mycolor2}{rgb}{0.00000,0.44700,0.74100}%
\definecolor{mycolor3}{rgb}{0.85000,0.32500,0.09800}%
%
\addplot[only marks, mark=*, mark size=1.2748pt, color=mycolor1, fill=mycolor1, opacity=0.60, draw=none, mark options={draw=none,line width=0pt}] table[row sep=crcr]{%
x	y\\
34	2659.6\\
29	1414.7\\
96	17087\\
31	1865.1\\
27	1245.2\\
70	7969.5\\
24	1175.5\\
37	2263.9\\
42	2826.4\\
16	1123.1\\
29	1578.5\\
59	3441.5\\
69	3268.8\\
28	1565.9\\
11	337.04\\
20	1072.5\\
45	3829.7\\
27	1353.6\\
22	1064.3\\
27	1473.9\\
44	4953.6\\
16	1632.6\\
39	4333.8\\
19	1155.6\\
33	2083.2\\
37	2597.1\\
33	1420.1\\
18	615.12\\
29	2511.1\\
22	1038.8\\
33	1466.1\\
22	952.09\\
49	8841.6\\
25	1457.4\\
43	3953.8\\
26	2345.2\\
31	3343.1\\
26	1332.8\\
18	886.52\\
13	513.86\\
32	4019.9\\
38	2379.3\\
32	3443.2\\
29	1686.3\\
35	1706\\
35	3160.2\\
135	20467\\
27	4701.4\\
42	3880.7\\
39	4923.8\\
51	6896.2\\
86	10357\\
15	957.67\\
39	5028.4\\
42	3470.5\\
43	3130.2\\
58	5795.7\\
97	24183\\
50	2691.1\\
20	1456.9\\
28	1244.6\\
22	1072.2\\
47	3224.4\\
154	63393\\
30	1815.2\\
10	476.82\\
37	7531.1\\
31	1924\\
54	13514\\
76	11305\\
55	4964.6\\
26	2661.5\\
33	1634.4\\
34	2158.6\\
25	1589\\
31	3864.4\\
34	4806.9\\
26	2342.6\\
39	3443\\
27	2951.2\\
32	1946.3\\
26	2251.5\\
29	2140.4\\
23	3106\\
26	2150.7\\
14	854.42\\
33	2069.5\\
24	1114.1\\
20	1657\\
51	2229.3\\
35	2878.6\\
29	1300.8\\
46	988.34\\
44	3612\\
45	3445.2\\
29	892.64\\
28	1743.1\\
77	3217.4\\
70	3997.3\\
70	3338.8\\
40	2569\\
90	5574.9\\
53	3745.5\\
62	1604.7\\
47	3610.6\\
34	1350\\
14	959.1\\
47	2071.8\\
43	1060.6\\
20	807.84\\
28	1062.6\\
80	2845.7\\
25	552.12\\
36	1804.4\\
27	2284.9\\
82	4078\\
31	1157.3\\
39	1404.3\\
23	2079.3\\
27	2918.1\\
37	2236.7\\
129	8341\\
63	4300.2\\
64	1997.6\\
91	9282.4\\
27	1217.5\\
26	1500.1\\
17	771.8\\
33	1617.2\\
30	3277.2\\
33	1427.7\\
182	14664\\
26	1157\\
48	2931.9\\
46	1887\\
10	387.38\\
25	1142.2\\
26	784.64\\
206	26843\\
37	1387.9\\
43	1657.4\\
51	1914.5\\
57	6472.8\\
36	1809.5\\
38	1616.6\\
28	1303.5\\
19	1126.3\\
36	1644.5\\
45	3206.5\\
53	3477.4\\
48	3680.4\\
64	3226.4\\
65	2938.8\\
45	1225.4\\
54	2931.1\\
16	459.7\\
24	1120.2\\
127	8740.3\\
48	1703.2\\
64	2499.5\\
42	2597.3\\
42	1308.6\\
17	562.74\\
48	2766.7\\
42	2150.6\\
20	877.68\\
34	1199.1\\
47	923.14\\
42	2929.1\\
50	2505.7\\
87	6689.2\\
64	3904.7\\
32	1917.4\\
39	1177.3\\
35	2435.8\\
18	1026.7\\
28	945.7\\
27	903.42\\
37	2155.5\\
30	2467.3\\
89	11830\\
22	1889.6\\
28	1565.7\\
25	1328.4\\
15	625.19\\
41	3527.1\\
32	963.33\\
34	1873.5\\
16	746.94\\
88	5937.8\\
39	980.57\\
140	27483\\
39	1259\\
49	1859.2\\
72	2264.4\\
37	1552.4\\
41	3720.2\\
283	1.5097e+05\\
31	2803.7\\
99	19380\\
103	24503\\
39	3629.3\\
35	1443\\
98	12986\\
45	4464\\
73	8456.1\\
34	2322.8\\
56	7870.4\\
39	2972.8\\
37	2902.4\\
55	2161.9\\
25	773.23\\
29	1052.3\\
42	1188.1\\
35	3915.8\\
26	762.86\\
45	8590\\
56	6323\\
39	1683.5\\
62	3118.3\\
94	12804\\
112	25043\\
25	732.82\\
57	2525.8\\
50	2583.8\\
47	7345.3\\
96	6934.2\\
54	6333.5\\
28	2464.6\\
38	4195.5\\
47	3331.8\\
35	1948.8\\
45	1643.7\\
50	3800.2\\
43	1046.8\\
29	2100.3\\
42	1789.1\\
71	19782\\
68	6076.5\\
149	52274\\
67	5541.5\\
51	2942.7\\
34	1918.1\\
120	14743\\
43	6853.5\\
117	11166\\
39	1444.8\\
72	9128.1\\
40	3718.5\\
43	2233.1\\
32	2559.9\\
49	9347.7\\
41	3701.8\\
98	5270.9\\
96	12850\\
45	1718.5\\
129	34983\\
43	2086.1\\
93	7942.5\\
41	3493.8\\
88	10220\\
73	5834.7\\
108	16419\\
52	7730.4\\
18	485.3\\
36	994.07\\
50	5230.4\\
37	2378.4\\
83	14641\\
40	2223.8\\
26	2719.7\\
40	1723.2\\
49	3680.1\\
27	1455.6\\
41	1334.6\\
36	939.06\\
52	5037.1\\
37	1179.9\\
68	7085.9\\
43	2901.8\\
132	13823\\
73	2911.6\\
46	4779.2\\
47	1257.5\\
44	1735.1\\
41	4749.9\\
30	1499.2\\
25	2922.5\\
24	2038.4\\
85	8202.1\\
44	2580.5\\
29	1710.2\\
34	1545.6\\
101	15542\\
18	1635.1\\
21	1281.4\\
111	16157\\
95	4957.6\\
23	1189.4\\
31	659.22\\
84	3869.4\\
23	824.67\\
55	5344.2\\
32	1205.8\\
27	918.03\\
48	970.54\\
18	386.22\\
28	2259\\
40	1278.4\\
30	1301.7\\
13	344.21\\
33	1407.5\\
14	545.31\\
52	2329.8\\
65	7552.6\\
49	2581.5\\
35	1453.9\\
109	7506.2\\
26	902.75\\
41	1307.5\\
39	1377\\
61	5863.5\\
40	1476.7\\
36	2244.1\\
21	590.21\\
80	4183.2\\
36	1299.4\\
36	2023.4\\
38	1040\\
42	2101.6\\
27	763.39\\
46	2209.8\\
23	601.01\\
34	1740.8\\
200	38200\\
37	2321\\
16	574.46\\
50	2231.8\\
19	656.56\\
41	2758.5\\
25	812.13\\
39	1907.4\\
45	3019.7\\
83	6800.9\\
27	536.65\\
67	2545.7\\
46	3015.3\\
35	1529.3\\
44	2515\\
30	1800.7\\
29	1169.9\\
23	508.63\\
43	2142.2\\
28	705.01\\
24	1129.4\\
56	3031.9\\
33	907.22\\
13	397.3\\
41	3038.9\\
34	955.97\\
14	415.81\\
83	8462.4\\
33	1665.3\\
35	1574.2\\
29	973.7\\
54	3766.7\\
41	1428.9\\
19	678.16\\
34	2078.4\\
19	705.37\\
38	2321.5\\
43	1680\\
14	671.47\\
17	781.43\\
12	439.35\\
};
\addlegendentry{\ExactScatterLegend}

\addplot[only marks, mark=*, mark size=1.2748pt, color=black, fill=black, opacity=0.60, draw=none, mark options={draw=none,line width=0pt}] table[row sep=crcr]{%
x	y\\
34	1787\\
29	2052\\
96	39026\\
31	939\\
27	839\\
70	9435\\
24	756\\
37	2292\\
42	3353\\
16	509\\
29	1198\\
59	4061\\
69	4632\\
28	987\\
11	377\\
20	825\\
45	6367\\
27	1370\\
22	689\\
27	1107\\
44	4679\\
16	622\\
39	5236\\
19	558\\
33	1997\\
37	1291\\
33	1077\\
18	485\\
29	1640\\
22	758\\
33	1231\\
22	643\\
49	9267\\
25	1115\\
43	11048\\
26	1875\\
31	2762\\
26	927\\
18	665\\
13	140\\
32	2698\\
38	2211\\
32	775\\
29	1173\\
35	1719\\
35	3078\\
135	41095\\
27	2591\\
42	3900\\
39	3847\\
51	5607\\
86	11830\\
15	918\\
39	3169\\
42	3772\\
43	2489\\
58	4861\\
97	21264\\
50	3266\\
20	736\\
28	686\\
22	1102\\
47	4024\\
154	1.2234e+05\\
30	1543\\
10	296\\
37	7252\\
31	1142\\
54	13398\\
76	13899\\
55	3344\\
26	1499\\
33	1121\\
34	1971\\
25	1709\\
31	2961\\
34	3625\\
26	2774\\
39	3551\\
27	1922\\
32	1860\\
26	1781\\
29	1942\\
23	2014\\
26	1636\\
14	530\\
33	1466\\
24	787\\
20	1114\\
51	1692\\
35	2940\\
29	1332\\
46	1054\\
44	3996\\
45	4071\\
29	803\\
28	557\\
77	3928\\
70	3625\\
70	2531\\
40	1723\\
90	8518\\
53	3650\\
62	3213\\
47	8035\\
34	817\\
14	526\\
47	2742\\
43	759\\
20	950\\
28	533\\
80	4654\\
25	1153\\
36	1716\\
27	2299\\
82	3021\\
31	1103\\
39	1587\\
23	810\\
27	2388\\
37	2269\\
129	3241\\
63	3568\\
64	2459\\
91	6356\\
27	696\\
26	1459\\
17	558\\
33	2465\\
30	2046\\
33	1791\\
182	11424\\
26	663\\
48	2633\\
46	2157\\
10	222\\
25	802\\
26	424\\
206	37527\\
37	1069\\
43	1448\\
51	2715\\
57	7724\\
36	1800\\
38	1162\\
28	539\\
19	746\\
36	1003\\
45	3212\\
53	4779\\
48	5264\\
64	2484\\
65	3084\\
45	1597\\
54	3279\\
16	299\\
24	1131\\
127	2601\\
48	1783\\
64	2544\\
42	2057\\
42	1206\\
17	254\\
48	2507\\
42	2016\\
20	816\\
34	1805\\
47	1076\\
42	2516\\
50	2233\\
87	10377\\
64	2072\\
32	1299\\
39	1215\\
35	3735\\
18	866\\
28	626\\
27	253\\
37	2518\\
30	2667\\
89	9128\\
22	1204\\
28	1732\\
25	1513\\
15	873\\
41	3761\\
32	1140\\
34	1426\\
16	352\\
88	7596\\
39	668\\
140	20491\\
39	1086\\
49	1025\\
72	1475\\
37	1233\\
41	3002\\
283	2.0424e+05\\
31	2681\\
99	16428\\
103	32887\\
39	6629\\
35	1076\\
98	6332\\
45	4113\\
73	7877\\
34	1444\\
56	6926\\
39	2045\\
37	3241\\
55	1793\\
25	444\\
29	1020\\
42	1480\\
35	2286\\
26	591\\
45	5458\\
56	7320\\
39	1112\\
62	3063\\
94	13380\\
112	30134\\
25	317\\
57	1896\\
50	2648\\
47	5982\\
96	4539\\
54	4348\\
28	1451\\
38	3579\\
47	3124\\
35	1834\\
45	921\\
50	3106\\
43	719\\
29	1516\\
42	1249\\
71	5868\\
68	5857\\
149	61636\\
67	5463\\
51	1708\\
34	1459\\
120	10119\\
43	6342\\
117	8499\\
39	1041\\
72	6468\\
40	4061\\
43	1854\\
32	2516\\
49	12313\\
41	4009\\
98	4503\\
96	5379\\
45	1383\\
129	23237\\
43	1506\\
93	4228\\
41	2160\\
88	9837\\
73	5327\\
108	12182\\
52	5524\\
18	189\\
36	878\\
50	4430\\
37	1383\\
83	13086\\
40	2202\\
26	1338\\
40	1198\\
49	4702\\
27	1189\\
41	1263\\
36	572\\
52	3610\\
37	656\\
68	6317\\
43	2092\\
132	14984\\
73	3147\\
46	3825\\
47	739\\
44	1624\\
41	5048\\
30	1539\\
25	1574\\
24	2598\\
85	7294\\
44	1753\\
29	1546\\
34	1160\\
101	9599\\
18	1226\\
21	992\\
111	7348\\
95	2640\\
23	617\\
31	310\\
84	1691\\
23	466\\
55	3785\\
32	1352\\
27	544\\
48	517\\
18	119\\
28	1203\\
40	962\\
30	685\\
13	188\\
33	1762\\
14	209\\
52	1836\\
65	8994\\
49	3401\\
35	1176\\
109	4663\\
26	550\\
41	1006\\
39	973\\
61	4953\\
40	1238\\
36	1694\\
21	551\\
80	5029\\
36	1037\\
36	1691\\
38	832\\
42	1192\\
27	590\\
46	1533\\
23	486\\
34	1267\\
200	30990\\
37	1240\\
16	413\\
50	2688\\
19	324\\
41	2143\\
25	934\\
39	3741\\
45	2869\\
83	3945\\
27	521\\
67	2908\\
46	2419\\
35	1912\\
44	2676\\
30	2051\\
29	674\\
23	204\\
43	936\\
28	859\\
24	467\\
56	2173\\
33	679\\
13	261\\
41	2629\\
34	588\\
14	216\\
83	10336\\
33	1005\\
35	1380\\
29	687\\
54	3377\\
41	993\\
19	351\\
34	1937\\
19	459\\
38	1126\\
43	1196\\
14	412\\
17	325\\
12	184\\
};
\addlegendentry{\RealScatterLegend}

      % \input{../../../../.tex/20/KS/SizeVsDegreefig8.tex}
      \legend{}
  
      \nextgroupplot[%
        % xmin=10,xmax=283,
        % ymin=100,ymax=1e+06,
      ] % This file was created by matlab2tikz.
%
\definecolor{mycolor1}{rgb}{0.83529,0.36863,0.00000}%
\definecolor{mycolor2}{rgb}{0.00000,0.44700,0.74100}%
\definecolor{mycolor3}{rgb}{0.85000,0.32500,0.09800}%
%
\addplot[only marks, mark=*, mark size=1.2748pt, color=mycolor1, fill=mycolor1, opacity=0.60, draw=none, mark options={draw=none,line width=0pt}] table[row sep=crcr]{%
x	y\\
34	2752.1\\
29	1447.9\\
96	17899\\
31	1897.2\\
27	1312.4\\
70	7519.4\\
24	1293.2\\
37	2217\\
42	2934.9\\
16	1196.2\\
29	1455.8\\
59	3672.2\\
69	3506.9\\
28	1430.5\\
11	328.45\\
20	1179.9\\
45	3924.6\\
27	1565.1\\
22	966.11\\
27	1377.7\\
44	4878.9\\
16	1535.8\\
39	4304.3\\
19	1016.5\\
33	2032.1\\
37	2366.7\\
33	1308.8\\
18	610.03\\
29	2329.7\\
22	1115.1\\
33	1372.5\\
22	924.94\\
49	8161.7\\
25	1350.4\\
43	3843.7\\
26	2355.5\\
31	2991.9\\
26	1322.2\\
18	910.13\\
13	524.25\\
32	4223.2\\
38	2614.1\\
32	3319.3\\
29	1690.5\\
35	1683.3\\
35	3143\\
135	19595\\
27	4534.5\\
42	4309.9\\
39	4435.6\\
51	6982.8\\
86	11144\\
15	800.43\\
39	4834.6\\
42	3437.7\\
43	3295.7\\
58	6074.7\\
97	23656\\
50	2655.5\\
20	1462.3\\
28	1227.1\\
22	1114.8\\
47	3311.4\\
154	62486\\
30	1873.3\\
10	502.74\\
37	8467.2\\
31	1883.7\\
54	13540\\
76	11832\\
55	4800.7\\
26	2723.5\\
33	1467.9\\
34	2231.5\\
25	1641.5\\
31	3774.1\\
34	4704.3\\
26	2337.3\\
39	3648.4\\
27	2584.9\\
32	2201.2\\
26	2058.9\\
29	2475.8\\
23	3275.9\\
26	2235.6\\
14	861.46\\
33	2109.5\\
24	1049.4\\
20	1658\\
51	2223.7\\
35	3137.6\\
29	1309.1\\
46	949.6\\
44	3557\\
45	3264.9\\
29	984.36\\
28	1757.2\\
77	3171\\
70	4100.3\\
70	3599.2\\
40	2546.3\\
90	5463.2\\
53	3769\\
62	1775.1\\
47	3553.9\\
34	1404.2\\
14	1081\\
47	2082.7\\
43	1081.9\\
20	788.24\\
28	1069.3\\
80	3224\\
25	568.75\\
36	1872.2\\
27	2275.1\\
82	4371.8\\
31	1083.1\\
39	1412.4\\
23	2279.4\\
27	2740.9\\
37	2116.4\\
129	8873.1\\
63	4286.9\\
64	1917.8\\
91	9929.6\\
27	1129.8\\
26	1628.7\\
17	833.66\\
33	1512.2\\
30	3272.8\\
33	1375.8\\
182	14819\\
26	1213.5\\
48	2742.3\\
46	1898.8\\
10	331.16\\
25	1145.2\\
26	744.56\\
206	25410\\
37	1287.7\\
43	1537.3\\
51	2011.3\\
57	6722.5\\
36	1672.8\\
38	1640.8\\
28	1301.3\\
19	1091.1\\
36	1630.6\\
45	3199\\
53	3491.8\\
48	3481.2\\
64	3039.5\\
65	3034.6\\
45	1160.2\\
54	2692.7\\
16	402.49\\
24	1006.1\\
127	8342.6\\
48	1814.7\\
64	2773.7\\
42	2547.1\\
42	1259.1\\
17	559.12\\
48	2914.9\\
42	2224.8\\
20	841.27\\
34	1128.3\\
47	1064.6\\
42	2961.1\\
50	2416.6\\
87	6222.2\\
64	3661.6\\
32	1823.6\\
39	1151.3\\
35	2425.1\\
18	1039.4\\
28	994.52\\
27	900.38\\
37	2226.4\\
30	2227.4\\
89	11814\\
22	1931.8\\
28	1589.8\\
25	1245\\
15	590.65\\
41	3560.7\\
32	948.29\\
34	1979.5\\
16	789.6\\
88	5367.1\\
39	1047.6\\
140	26639\\
39	1240.2\\
49	1706.1\\
72	2226.7\\
37	1609.9\\
41	3304\\
283	1.537e+05\\
31	3008.5\\
99	18581\\
103	23629\\
39	3638.2\\
35	1512.7\\
98	13472\\
45	4614.8\\
73	8967.7\\
34	2178.8\\
56	7761.4\\
39	3241.2\\
37	3173.6\\
55	2166.4\\
25	776.54\\
29	1104.3\\
42	1198.7\\
35	3878.8\\
26	762.57\\
45	8279.4\\
56	6466.5\\
39	1541.7\\
62	3102.1\\
94	11988\\
112	23458\\
25	725.17\\
57	2347.9\\
50	2470.1\\
47	7040.9\\
96	6892.2\\
54	6397.8\\
28	2418.8\\
38	3666\\
47	3740.2\\
35	1855.5\\
45	1593.2\\
50	3933.1\\
43	1021.4\\
29	2149.5\\
42	1867.7\\
71	21595\\
68	6832.4\\
149	51885\\
67	6095.8\\
51	2849.5\\
34	1741.1\\
120	14166\\
43	7024.9\\
117	10766\\
39	1385.5\\
72	8517.9\\
40	3570.3\\
43	2546.7\\
32	2440\\
49	8717\\
41	4216.1\\
98	4851.4\\
96	13082\\
45	1790.2\\
129	34070\\
43	2064.6\\
93	8244.5\\
41	3517.4\\
88	10215\\
73	6049.7\\
108	15669\\
52	7762.1\\
18	459.6\\
36	1047.3\\
50	5559\\
37	2276.1\\
83	14453\\
40	2357.3\\
26	2746.3\\
40	1747.4\\
49	3950.3\\
27	1417.2\\
41	1374.5\\
36	872.47\\
52	4928.5\\
37	1076.9\\
68	7268.5\\
43	2696.2\\
132	14335\\
73	2855.8\\
46	4938.5\\
47	1346\\
44	1660.2\\
41	5069.6\\
30	1525.9\\
25	2998\\
24	1825.6\\
85	8306.6\\
44	2949.4\\
29	1714.7\\
34	1536.1\\
101	15121\\
18	1458.8\\
21	1345.4\\
111	16738\\
95	4473.2\\
23	1047\\
31	655.26\\
84	4030.8\\
23	836.93\\
55	4995.2\\
32	1154.3\\
27	954.76\\
48	990.47\\
18	348.75\\
28	1897.4\\
40	1229\\
30	1323\\
13	334.01\\
33	1605.1\\
14	590.46\\
52	2527.6\\
65	7783.7\\
49	2447.2\\
35	1309.6\\
109	6976.6\\
26	974.29\\
41	1320.4\\
39	1273.1\\
61	5949.8\\
40	1394.4\\
36	2038.2\\
21	645.28\\
80	4376.3\\
36	1407.8\\
36	1979.2\\
38	1088.9\\
42	1970.8\\
27	765.93\\
46	2331.7\\
23	663.1\\
34	1777\\
200	40388\\
37	2389.4\\
16	607.66\\
50	2591.5\\
19	646.53\\
41	2830\\
25	942.11\\
39	2154.3\\
45	3049.6\\
83	7540.9\\
27	564.42\\
67	2554\\
46	2997.2\\
35	1591.2\\
44	2521.3\\
30	1856.8\\
29	1130.1\\
23	479.02\\
43	1924.8\\
28	732.66\\
24	1081\\
56	3305\\
33	1003.7\\
13	438.81\\
41	3042.6\\
34	955.6\\
14	441.88\\
83	8253.1\\
33	1549.1\\
35	1503.1\\
29	1021.3\\
54	3793.1\\
41	1286.9\\
19	676.49\\
34	2090.4\\
19	723.93\\
38	2349.1\\
43	1695.8\\
14	621.64\\
17	768.51\\
12	458.88\\
};
\addlegendentry{\ApproxScatterLegend}

\addplot[only marks, mark=*, mark size=1.2748pt, color=black, fill=black, opacity=0.60, draw=none, mark options={draw=none,line width=0pt}] table[row sep=crcr]{%
x	y\\
34	1787\\
29	2052\\
96	39026\\
31	939\\
27	839\\
70	9435\\
24	756\\
37	2292\\
42	3353\\
16	509\\
29	1198\\
59	4061\\
69	4632\\
28	987\\
11	377\\
20	825\\
45	6367\\
27	1370\\
22	689\\
27	1107\\
44	4679\\
16	622\\
39	5236\\
19	558\\
33	1997\\
37	1291\\
33	1077\\
18	485\\
29	1640\\
22	758\\
33	1231\\
22	643\\
49	9267\\
25	1115\\
43	11048\\
26	1875\\
31	2762\\
26	927\\
18	665\\
13	140\\
32	2698\\
38	2211\\
32	775\\
29	1173\\
35	1719\\
35	3078\\
135	41095\\
27	2591\\
42	3900\\
39	3847\\
51	5607\\
86	11830\\
15	918\\
39	3169\\
42	3772\\
43	2489\\
58	4861\\
97	21264\\
50	3266\\
20	736\\
28	686\\
22	1102\\
47	4024\\
154	1.2234e+05\\
30	1543\\
10	296\\
37	7252\\
31	1142\\
54	13398\\
76	13899\\
55	3344\\
26	1499\\
33	1121\\
34	1971\\
25	1709\\
31	2961\\
34	3625\\
26	2774\\
39	3551\\
27	1922\\
32	1860\\
26	1781\\
29	1942\\
23	2014\\
26	1636\\
14	530\\
33	1466\\
24	787\\
20	1114\\
51	1692\\
35	2940\\
29	1332\\
46	1054\\
44	3996\\
45	4071\\
29	803\\
28	557\\
77	3928\\
70	3625\\
70	2531\\
40	1723\\
90	8518\\
53	3650\\
62	3213\\
47	8035\\
34	817\\
14	526\\
47	2742\\
43	759\\
20	950\\
28	533\\
80	4654\\
25	1153\\
36	1716\\
27	2299\\
82	3021\\
31	1103\\
39	1587\\
23	810\\
27	2388\\
37	2269\\
129	3241\\
63	3568\\
64	2459\\
91	6356\\
27	696\\
26	1459\\
17	558\\
33	2465\\
30	2046\\
33	1791\\
182	11424\\
26	663\\
48	2633\\
46	2157\\
10	222\\
25	802\\
26	424\\
206	37527\\
37	1069\\
43	1448\\
51	2715\\
57	7724\\
36	1800\\
38	1162\\
28	539\\
19	746\\
36	1003\\
45	3212\\
53	4779\\
48	5264\\
64	2484\\
65	3084\\
45	1597\\
54	3279\\
16	299\\
24	1131\\
127	2601\\
48	1783\\
64	2544\\
42	2057\\
42	1206\\
17	254\\
48	2507\\
42	2016\\
20	816\\
34	1805\\
47	1076\\
42	2516\\
50	2233\\
87	10377\\
64	2072\\
32	1299\\
39	1215\\
35	3735\\
18	866\\
28	626\\
27	253\\
37	2518\\
30	2667\\
89	9128\\
22	1204\\
28	1732\\
25	1513\\
15	873\\
41	3761\\
32	1140\\
34	1426\\
16	352\\
88	7596\\
39	668\\
140	20491\\
39	1086\\
49	1025\\
72	1475\\
37	1233\\
41	3002\\
283	2.0424e+05\\
31	2681\\
99	16428\\
103	32887\\
39	6629\\
35	1076\\
98	6332\\
45	4113\\
73	7877\\
34	1444\\
56	6926\\
39	2045\\
37	3241\\
55	1793\\
25	444\\
29	1020\\
42	1480\\
35	2286\\
26	591\\
45	5458\\
56	7320\\
39	1112\\
62	3063\\
94	13380\\
112	30134\\
25	317\\
57	1896\\
50	2648\\
47	5982\\
96	4539\\
54	4348\\
28	1451\\
38	3579\\
47	3124\\
35	1834\\
45	921\\
50	3106\\
43	719\\
29	1516\\
42	1249\\
71	5868\\
68	5857\\
149	61636\\
67	5463\\
51	1708\\
34	1459\\
120	10119\\
43	6342\\
117	8499\\
39	1041\\
72	6468\\
40	4061\\
43	1854\\
32	2516\\
49	12313\\
41	4009\\
98	4503\\
96	5379\\
45	1383\\
129	23237\\
43	1506\\
93	4228\\
41	2160\\
88	9837\\
73	5327\\
108	12182\\
52	5524\\
18	189\\
36	878\\
50	4430\\
37	1383\\
83	13086\\
40	2202\\
26	1338\\
40	1198\\
49	4702\\
27	1189\\
41	1263\\
36	572\\
52	3610\\
37	656\\
68	6317\\
43	2092\\
132	14984\\
73	3147\\
46	3825\\
47	739\\
44	1624\\
41	5048\\
30	1539\\
25	1574\\
24	2598\\
85	7294\\
44	1753\\
29	1546\\
34	1160\\
101	9599\\
18	1226\\
21	992\\
111	7348\\
95	2640\\
23	617\\
31	310\\
84	1691\\
23	466\\
55	3785\\
32	1352\\
27	544\\
48	517\\
18	119\\
28	1203\\
40	962\\
30	685\\
13	188\\
33	1762\\
14	209\\
52	1836\\
65	8994\\
49	3401\\
35	1176\\
109	4663\\
26	550\\
41	1006\\
39	973\\
61	4953\\
40	1238\\
36	1694\\
21	551\\
80	5029\\
36	1037\\
36	1691\\
38	832\\
42	1192\\
27	590\\
46	1533\\
23	486\\
34	1267\\
200	30990\\
37	1240\\
16	413\\
50	2688\\
19	324\\
41	2143\\
25	934\\
39	3741\\
45	2869\\
83	3945\\
27	521\\
67	2908\\
46	2419\\
35	1912\\
44	2676\\
30	2051\\
29	674\\
23	204\\
43	936\\
28	859\\
24	467\\
56	2173\\
33	679\\
13	261\\
41	2629\\
34	588\\
14	216\\
83	10336\\
33	1005\\
35	1380\\
29	687\\
54	3377\\
41	993\\
19	351\\
34	1937\\
19	459\\
38	1126\\
43	1196\\
14	412\\
17	325\\
12	184\\
};
\addlegendentry{\RealScatterLegend}

      % \input{../../../../.tex/20/KS/SizeVsDegreefig9.tex}
      \legend{}
    \end{groupplot}
    
    \begin{groupplot}[
      group style={
        group name=Row4,
        plotGroup
      },
      plotRow,
      row4Keys,
      plotLegendNW,
      %
      xmin=10,xmax=283,
      ymin=100,ymax=2.0424e+05,
    ]
      \nextgroupplot[%
        at={($(Row3 c1r1.south west)-(0,\yPlotSep em)-(0,\plotHeight cm)$)},
        % xmin=10,xmax=283,
        % ymin=100,ymax=1.5536e+05,
      ] % This file was created by matlab2tikz.
%
\definecolor{mycolor1}{rgb}{0.00000,0.44706,0.69804}%
\definecolor{mycolor2}{rgb}{0.00000,0.44700,0.74100}%
\definecolor{mycolor3}{rgb}{0.83529,0.36863,0.00000}%
\definecolor{mycolor4}{rgb}{0.85000,0.32500,0.09800}%
%
\addplot[only marks, mark=*, mark size=1.2748pt, color=mycolor1, fill=mycolor1, opacity=0.70, draw=none, mark options={draw=none,line width=0pt}] table[row sep=crcr]{%
x	y\\
34	2940.2\\
29	1291.1\\
96	16971\\
31	1717.7\\
27	1282.8\\
70	7732.5\\
24	1101.1\\
37	2136.9\\
42	2882.9\\
16	1252.8\\
29	1498.7\\
59	3373.1\\
69	3235.2\\
28	1441.3\\
11	325.86\\
20	1086.5\\
45	3586.9\\
27	1404.8\\
22	1002.3\\
27	1497.5\\
44	4679.6\\
16	1375.3\\
39	4395.2\\
19	1187\\
33	2235.1\\
37	2263.8\\
33	1234.3\\
18	469.22\\
29	2589.8\\
22	951.71\\
33	1325.5\\
22	917.87\\
49	8375.4\\
25	1849.5\\
43	3990\\
26	2500.2\\
31	3270.4\\
26	1074.5\\
18	867.25\\
13	545.22\\
32	3772.8\\
38	2449.8\\
32	3403.5\\
29	1777.1\\
35	1417.7\\
35	2966.1\\
135	20862\\
27	4403.8\\
42	4267.2\\
39	4914.1\\
51	7739.2\\
86	10752\\
15	905.44\\
39	4478\\
42	3506.1\\
43	3400.5\\
58	5470.6\\
97	26628\\
50	2300.4\\
20	1475.9\\
28	1132.2\\
22	1098\\
47	3094\\
154	61292\\
30	1846.2\\
10	490.72\\
37	7809.1\\
31	1847.3\\
54	14762\\
76	11584\\
55	4860.3\\
26	2444.3\\
33	1490.5\\
34	2026.2\\
25	1655.6\\
31	3337.6\\
34	4763\\
26	2168.7\\
39	3681.1\\
27	2439.5\\
32	1987.9\\
26	2043\\
29	2261.3\\
23	2913.9\\
26	2308.5\\
14	785.08\\
33	2126.2\\
24	1040.1\\
20	1536.5\\
51	2077.4\\
35	2967.7\\
29	1233.9\\
46	957.83\\
44	3317.6\\
45	3006\\
29	918.53\\
28	1788.7\\
77	2931.7\\
70	4035.6\\
70	3519.9\\
40	2535.6\\
90	5650.4\\
53	3731.3\\
62	1662.5\\
47	3632.2\\
34	1236.7\\
14	1060\\
47	2057.9\\
43	913.44\\
20	779.99\\
28	1071.4\\
80	3437.4\\
25	478.64\\
36	1833.2\\
27	2270.4\\
82	4129.7\\
31	976.91\\
39	1388.5\\
23	2194.6\\
27	2699.7\\
37	2223.4\\
129	8303.9\\
63	4338.2\\
64	1827\\
91	9821.4\\
27	1260.4\\
26	1507.9\\
17	760.59\\
33	1421.3\\
30	3139.4\\
33	1403.9\\
182	14516\\
26	1100.6\\
48	3017.9\\
46	1551.7\\
10	387.04\\
25	1148.2\\
26	817.57\\
206	26265\\
37	1282.7\\
43	1583.2\\
51	2045.1\\
57	6607.9\\
36	1844\\
38	1587.9\\
28	1151.3\\
19	1074.1\\
36	1554.3\\
45	3285.2\\
53	3520.6\\
48	3668.2\\
64	3211.3\\
65	2955.5\\
45	1237.1\\
54	2854.6\\
16	390.6\\
24	1041.2\\
127	8233\\
48	1860\\
64	2212.8\\
42	2501.3\\
42	1266.7\\
17	551.79\\
48	2673.1\\
42	2147\\
20	929.2\\
34	1237.6\\
47	896.18\\
42	3044\\
50	2265.5\\
87	6251.1\\
64	4159.1\\
32	1827.7\\
39	1112.4\\
35	2577.9\\
18	1147.8\\
28	919.89\\
27	861.42\\
37	2053.9\\
30	2412.8\\
89	11238\\
22	1696.4\\
28	1493.8\\
25	1308.4\\
15	629.19\\
41	3290.3\\
32	891.93\\
34	2063.3\\
16	745.36\\
88	5888.7\\
39	875.43\\
140	26447\\
39	1142.2\\
49	1848.4\\
72	2644.4\\
37	1767.4\\
41	3476\\
283	1.5156e+05\\
31	2849.6\\
99	19714\\
103	22455\\
39	3599.9\\
35	1640.3\\
98	13921\\
45	4180.4\\
73	8349.1\\
34	2108.4\\
56	9016\\
39	2901.2\\
37	2687.3\\
55	2139.4\\
25	718.05\\
29	1048.3\\
42	1278.7\\
35	3957.7\\
26	734.79\\
45	8275.4\\
56	6603.3\\
39	1582.9\\
62	3235.6\\
94	12780\\
112	24473\\
25	752.43\\
57	2288.6\\
50	2631.3\\
47	6745.5\\
96	6767.4\\
54	6369.8\\
28	2559.9\\
38	3834.4\\
47	3758.1\\
35	2021\\
45	1609.6\\
50	3741.8\\
43	1022.6\\
29	2162.5\\
42	1795.8\\
71	21275\\
68	6531.6\\
149	55936\\
67	5842.9\\
51	2737.1\\
34	1891.9\\
120	15638\\
43	7430.7\\
117	10629\\
39	1168.5\\
72	7968.9\\
40	3429.1\\
43	2526\\
32	2663.4\\
49	8108.9\\
41	3810.2\\
98	5261.1\\
96	12647\\
45	1658.3\\
129	36079\\
43	2109.4\\
93	8115.9\\
41	3304.3\\
88	9772.2\\
73	5831.5\\
108	17181\\
52	7850.2\\
18	425.47\\
36	1109.8\\
50	5465.4\\
37	2218.8\\
83	13731\\
40	2334.9\\
26	2418.1\\
40	1598.6\\
49	3962.6\\
27	1275\\
41	1354.5\\
36	913.21\\
52	4772\\
37	1110.5\\
68	7283.9\\
43	2413.1\\
132	15299\\
73	3129\\
46	4146.3\\
47	1117.9\\
44	1541.3\\
41	5272.3\\
30	1536\\
25	3326\\
24	1862.8\\
85	8292.9\\
44	2929.9\\
29	1746.9\\
34	1479.5\\
101	13746\\
18	1485.4\\
21	1315.5\\
111	18025\\
95	4551.7\\
23	988.81\\
31	585.61\\
84	4434.1\\
23	809.94\\
55	5486.4\\
32	1179.9\\
27	875.64\\
48	977.88\\
18	347.33\\
28	2082\\
40	1189.3\\
30	1308.5\\
13	307.16\\
33	1435.8\\
14	595.42\\
52	2770.4\\
65	7874.2\\
49	2687.7\\
35	1373.4\\
109	6696.1\\
26	872.04\\
41	1478.3\\
39	1312.5\\
61	5522.9\\
40	1370.9\\
36	2123.5\\
21	661.41\\
80	4149.8\\
36	1349.5\\
36	1949.7\\
38	1055.6\\
42	2315.6\\
27	727.45\\
46	2254.4\\
23	651.91\\
34	1779.8\\
200	40987\\
37	2145.9\\
16	620.53\\
50	2488.1\\
19	683.81\\
41	2752.5\\
25	802.5\\
39	1904.9\\
45	2679.8\\
83	7336.9\\
27	477.88\\
67	2495.2\\
46	3493.8\\
35	1444.7\\
44	2319\\
30	1943.4\\
29	1173.8\\
23	502.13\\
43	2063.1\\
28	711.1\\
24	1110.5\\
56	3015.5\\
33	939.39\\
13	387.66\\
41	2847.1\\
34	937.58\\
14	434.52\\
83	8798.4\\
33	1499.8\\
35	1486.1\\
29	908.1\\
54	3652.1\\
41	1290.2\\
19	763.71\\
34	2034.6\\
19	674.28\\
38	2256.2\\
43	1846.5\\
14	663.66\\
17	849.53\\
12	422.36\\
};
\addlegendentry{\ExactScatterLegend}

\addplot[only marks, mark=*, mark size=1.2748pt, color=mycolor3, fill=mycolor3, opacity=0.70, draw=none, mark options={draw=none,line width=0pt}] table[row sep=crcr]{%
x	y\\
34	2744.6\\
29	1352.1\\
96	18508\\
31	1972.2\\
27	1412.6\\
70	7151.7\\
24	1269.3\\
37	2171.6\\
42	3051.1\\
16	1170.4\\
29	1405.5\\
59	3376.5\\
69	3130.8\\
28	1525.7\\
11	268.8\\
20	1192.7\\
45	3650.6\\
27	1533.6\\
22	928.61\\
27	1355.4\\
44	4883.6\\
16	1490.3\\
39	4095.5\\
19	1047.5\\
33	2135.3\\
37	2527.7\\
33	1356.8\\
18	499.27\\
29	2354.9\\
22	887.35\\
33	1420.9\\
22	883.81\\
49	8402.9\\
25	1528.2\\
43	3861.1\\
26	2452.3\\
31	3038.4\\
26	1031\\
18	807.01\\
13	447.17\\
32	4157.7\\
38	2228.7\\
32	3506.7\\
29	1701\\
35	1433\\
35	2953.1\\
135	20957\\
27	4295.8\\
42	3693\\
39	4104\\
51	7517.2\\
86	10965\\
15	831.26\\
39	3901.6\\
42	3162.6\\
43	3196.8\\
58	5643.4\\
97	25859\\
50	2350.4\\
20	1365.9\\
28	1124.5\\
22	1211.7\\
47	3544.5\\
154	63804\\
30	1775.9\\
10	466.4\\
37	7501\\
31	1925.3\\
54	14002\\
76	12012\\
55	4544.8\\
26	2573.9\\
33	1501.2\\
34	2342.5\\
25	1733.6\\
31	3605.5\\
34	4708.8\\
26	2373.9\\
39	3442.2\\
27	2578.4\\
32	1848.1\\
26	2159.1\\
29	2299.2\\
23	3127.2\\
26	2579.3\\
14	856.55\\
33	1945.3\\
24	1035.2\\
20	1869.3\\
51	2174\\
35	2792\\
29	1174.5\\
46	958.99\\
44	3639.5\\
45	3103.7\\
29	881.37\\
28	1671.6\\
77	2963.3\\
70	4125.8\\
70	3446.6\\
40	2173.3\\
90	6023.9\\
53	3407.9\\
62	1779.1\\
47	3496.1\\
34	1300.6\\
14	1042.2\\
47	1952.4\\
43	1042\\
20	754.06\\
28	1028.9\\
80	3155.7\\
25	517.75\\
36	1680.5\\
27	2118.6\\
82	4432.8\\
31	1080.9\\
39	1477.1\\
23	1975.1\\
27	2684.6\\
37	2122.1\\
129	8872.1\\
63	4710.7\\
64	1961.3\\
91	10283\\
27	1192.6\\
26	1577\\
17	702.8\\
33	1628\\
30	3113.6\\
33	1392.2\\
182	14724\\
26	1010.9\\
48	2990\\
46	1538.2\\
10	361.09\\
25	1059.4\\
26	830.75\\
206	26137\\
37	1267.5\\
43	1739.2\\
51	2187.3\\
57	6934.2\\
36	1695.1\\
38	1564.5\\
28	1355.6\\
19	1091.1\\
36	1415.1\\
45	3278.8\\
53	3606.6\\
48	3690.6\\
64	3053.8\\
65	3000.2\\
45	1135.7\\
54	2964.4\\
16	400.54\\
24	971.81\\
127	9079.8\\
48	1678.6\\
64	2526.7\\
42	2497.5\\
42	1250.3\\
17	521.01\\
48	2618.5\\
42	2156.2\\
20	858.82\\
34	1203.4\\
47	965.97\\
42	2998.2\\
50	2165.9\\
87	6388\\
64	3723\\
32	1782.5\\
39	1031.4\\
35	2494.7\\
18	1010.7\\
28	949.86\\
27	1024.8\\
37	2088.9\\
30	2568.3\\
89	10616\\
22	1770\\
28	1612.6\\
25	1198.2\\
15	596.95\\
41	3503.5\\
32	975.22\\
34	2029.2\\
16	759.71\\
88	6403.9\\
39	967.84\\
140	28322\\
39	1306.9\\
49	1666.1\\
72	2412.1\\
37	1620.3\\
41	3630.4\\
283	1.5536e+05\\
31	2751.1\\
99	20066\\
103	24294\\
39	3455.7\\
35	1566.2\\
98	13600\\
45	4492.5\\
73	8693.6\\
34	2183.5\\
56	7772.4\\
39	2865.7\\
37	2844\\
55	1957.3\\
25	723.62\\
29	1064\\
42	1057.1\\
35	4251.9\\
26	810.59\\
45	8453.8\\
56	6246.8\\
39	1542\\
62	2935.7\\
94	12348\\
112	23386\\
25	786.27\\
57	2286.6\\
50	2503\\
47	6654.5\\
96	7273.5\\
54	5725\\
28	2498.9\\
38	3759.5\\
47	3726.4\\
35	1937.4\\
45	1730.2\\
50	3570.4\\
43	949.61\\
29	1982.7\\
42	1725.5\\
71	19464\\
68	6767.5\\
149	53984\\
67	5713\\
51	2940.4\\
34	1900.5\\
120	14729\\
43	6841.9\\
117	11434\\
39	1318\\
72	8418.4\\
40	3417.4\\
43	2345.3\\
32	2507.7\\
49	8243.3\\
41	4138.4\\
98	4839.7\\
96	13039\\
45	1469.2\\
129	34913\\
43	2142.6\\
93	8535.5\\
41	3591.3\\
88	9223.9\\
73	6328.5\\
108	15883\\
52	8622.1\\
18	446.75\\
36	1021.5\\
50	5862.5\\
37	2240.8\\
83	15513\\
40	2127.7\\
26	2471.3\\
40	1647.3\\
49	3903.7\\
27	1488.9\\
41	1495.2\\
36	829.09\\
52	5132\\
37	1100.3\\
68	7100.7\\
43	2679\\
132	14748\\
73	2647.1\\
46	4140.5\\
47	1156.5\\
44	1704.9\\
41	4874.5\\
30	1410.5\\
25	2614.3\\
24	2035.5\\
85	8186.2\\
44	3142.4\\
29	1643.6\\
34	1548.8\\
101	14112\\
18	1667.7\\
21	1323.5\\
111	16852\\
95	4540.8\\
23	1008.6\\
31	658.26\\
84	4168.8\\
23	703.35\\
55	5117.4\\
32	1022.6\\
27	851.13\\
48	1047.7\\
18	322.04\\
28	1872.6\\
40	1161.4\\
30	1242.6\\
13	296.76\\
33	1544.5\\
14	553.78\\
52	2464.3\\
65	7334.8\\
49	2654.8\\
35	1283.4\\
109	7332.5\\
26	839.99\\
41	1210.9\\
39	1193.9\\
61	6446.7\\
40	1261.8\\
36	2108.7\\
21	617.27\\
80	4217.9\\
36	1425.3\\
36	1970.4\\
38	1162\\
42	2142.1\\
27	722.94\\
46	2098\\
23	634.69\\
34	1656.6\\
200	40055\\
37	2526.2\\
16	552.88\\
50	2415.2\\
19	617.72\\
41	3059.7\\
25	722.37\\
39	1922.9\\
45	3190.1\\
83	7306.8\\
27	528.05\\
67	2490.3\\
46	3060\\
35	1506.2\\
44	2339.9\\
30	1796.6\\
29	1120.7\\
23	537.31\\
43	1974.8\\
28	688.91\\
24	965.84\\
56	2955.4\\
33	985.06\\
13	420.52\\
41	3164.7\\
34	931.46\\
14	442.41\\
83	8658.3\\
33	1719.3\\
35	1458.9\\
29	962.91\\
54	3339\\
41	1258\\
19	731.95\\
34	1993.9\\
19	684.41\\
38	2234.5\\
43	1515.5\\
14	625.43\\
17	775.01\\
12	448.3\\
};
\addlegendentry{\ApproxScatterLegend}

      % \input{../../../../.tex/20/KS/SizeVsDegreefig10.tex}
      \legend{}
  
      \nextgroupplot[%
        % xmin=10,xmax=283,
        % ymin=100,ymax=1e+06,
      ] % This file was created by matlab2tikz.
%
\definecolor{mycolor1}{rgb}{0.00000,0.44706,0.69804}%
\definecolor{mycolor2}{rgb}{0.00000,0.44700,0.74100}%
\definecolor{mycolor3}{rgb}{0.85000,0.32500,0.09800}%
%
\addplot[only marks, mark=*, mark size=1.2748pt, color=mycolor1, fill=mycolor1, opacity=0.60, draw=none, mark options={draw=none,line width=0pt}] table[row sep=crcr]{%
x	y\\
34	2940.2\\
29	1291.1\\
96	16971\\
31	1717.7\\
27	1282.8\\
70	7732.5\\
24	1101.1\\
37	2136.9\\
42	2882.9\\
16	1252.8\\
29	1498.7\\
59	3373.1\\
69	3235.2\\
28	1441.3\\
11	325.86\\
20	1086.5\\
45	3586.9\\
27	1404.8\\
22	1002.3\\
27	1497.5\\
44	4679.6\\
16	1375.3\\
39	4395.2\\
19	1187\\
33	2235.1\\
37	2263.8\\
33	1234.3\\
18	469.22\\
29	2589.8\\
22	951.71\\
33	1325.5\\
22	917.87\\
49	8375.4\\
25	1849.5\\
43	3990\\
26	2500.2\\
31	3270.4\\
26	1074.5\\
18	867.25\\
13	545.22\\
32	3772.8\\
38	2449.8\\
32	3403.5\\
29	1777.1\\
35	1417.7\\
35	2966.1\\
135	20862\\
27	4403.8\\
42	4267.2\\
39	4914.1\\
51	7739.2\\
86	10752\\
15	905.44\\
39	4478\\
42	3506.1\\
43	3400.5\\
58	5470.6\\
97	26628\\
50	2300.4\\
20	1475.9\\
28	1132.2\\
22	1098\\
47	3094\\
154	61292\\
30	1846.2\\
10	490.72\\
37	7809.1\\
31	1847.3\\
54	14762\\
76	11584\\
55	4860.3\\
26	2444.3\\
33	1490.5\\
34	2026.2\\
25	1655.6\\
31	3337.6\\
34	4763\\
26	2168.7\\
39	3681.1\\
27	2439.5\\
32	1987.9\\
26	2043\\
29	2261.3\\
23	2913.9\\
26	2308.5\\
14	785.08\\
33	2126.2\\
24	1040.1\\
20	1536.5\\
51	2077.4\\
35	2967.7\\
29	1233.9\\
46	957.83\\
44	3317.6\\
45	3006\\
29	918.53\\
28	1788.7\\
77	2931.7\\
70	4035.6\\
70	3519.9\\
40	2535.6\\
90	5650.4\\
53	3731.3\\
62	1662.5\\
47	3632.2\\
34	1236.7\\
14	1060\\
47	2057.9\\
43	913.44\\
20	779.99\\
28	1071.4\\
80	3437.4\\
25	478.64\\
36	1833.2\\
27	2270.4\\
82	4129.7\\
31	976.91\\
39	1388.5\\
23	2194.6\\
27	2699.7\\
37	2223.4\\
129	8303.9\\
63	4338.2\\
64	1827\\
91	9821.4\\
27	1260.4\\
26	1507.9\\
17	760.59\\
33	1421.3\\
30	3139.4\\
33	1403.9\\
182	14516\\
26	1100.6\\
48	3017.9\\
46	1551.7\\
10	387.04\\
25	1148.2\\
26	817.57\\
206	26265\\
37	1282.7\\
43	1583.2\\
51	2045.1\\
57	6607.9\\
36	1844\\
38	1587.9\\
28	1151.3\\
19	1074.1\\
36	1554.3\\
45	3285.2\\
53	3520.6\\
48	3668.2\\
64	3211.3\\
65	2955.5\\
45	1237.1\\
54	2854.6\\
16	390.6\\
24	1041.2\\
127	8233\\
48	1860\\
64	2212.8\\
42	2501.3\\
42	1266.7\\
17	551.79\\
48	2673.1\\
42	2147\\
20	929.2\\
34	1237.6\\
47	896.18\\
42	3044\\
50	2265.5\\
87	6251.1\\
64	4159.1\\
32	1827.7\\
39	1112.4\\
35	2577.9\\
18	1147.8\\
28	919.89\\
27	861.42\\
37	2053.9\\
30	2412.8\\
89	11238\\
22	1696.4\\
28	1493.8\\
25	1308.4\\
15	629.19\\
41	3290.3\\
32	891.93\\
34	2063.3\\
16	745.36\\
88	5888.7\\
39	875.43\\
140	26447\\
39	1142.2\\
49	1848.4\\
72	2644.4\\
37	1767.4\\
41	3476\\
283	1.5156e+05\\
31	2849.6\\
99	19714\\
103	22455\\
39	3599.9\\
35	1640.3\\
98	13921\\
45	4180.4\\
73	8349.1\\
34	2108.4\\
56	9016\\
39	2901.2\\
37	2687.3\\
55	2139.4\\
25	718.05\\
29	1048.3\\
42	1278.7\\
35	3957.7\\
26	734.79\\
45	8275.4\\
56	6603.3\\
39	1582.9\\
62	3235.6\\
94	12780\\
112	24473\\
25	752.43\\
57	2288.6\\
50	2631.3\\
47	6745.5\\
96	6767.4\\
54	6369.8\\
28	2559.9\\
38	3834.4\\
47	3758.1\\
35	2021\\
45	1609.6\\
50	3741.8\\
43	1022.6\\
29	2162.5\\
42	1795.8\\
71	21275\\
68	6531.6\\
149	55936\\
67	5842.9\\
51	2737.1\\
34	1891.9\\
120	15638\\
43	7430.7\\
117	10629\\
39	1168.5\\
72	7968.9\\
40	3429.1\\
43	2526\\
32	2663.4\\
49	8108.9\\
41	3810.2\\
98	5261.1\\
96	12647\\
45	1658.3\\
129	36079\\
43	2109.4\\
93	8115.9\\
41	3304.3\\
88	9772.2\\
73	5831.5\\
108	17181\\
52	7850.2\\
18	425.47\\
36	1109.8\\
50	5465.4\\
37	2218.8\\
83	13731\\
40	2334.9\\
26	2418.1\\
40	1598.6\\
49	3962.6\\
27	1275\\
41	1354.5\\
36	913.21\\
52	4772\\
37	1110.5\\
68	7283.9\\
43	2413.1\\
132	15299\\
73	3129\\
46	4146.3\\
47	1117.9\\
44	1541.3\\
41	5272.3\\
30	1536\\
25	3326\\
24	1862.8\\
85	8292.9\\
44	2929.9\\
29	1746.9\\
34	1479.5\\
101	13746\\
18	1485.4\\
21	1315.5\\
111	18025\\
95	4551.7\\
23	988.81\\
31	585.61\\
84	4434.1\\
23	809.94\\
55	5486.4\\
32	1179.9\\
27	875.64\\
48	977.88\\
18	347.33\\
28	2082\\
40	1189.3\\
30	1308.5\\
13	307.16\\
33	1435.8\\
14	595.42\\
52	2770.4\\
65	7874.2\\
49	2687.7\\
35	1373.4\\
109	6696.1\\
26	872.04\\
41	1478.3\\
39	1312.5\\
61	5522.9\\
40	1370.9\\
36	2123.5\\
21	661.41\\
80	4149.8\\
36	1349.5\\
36	1949.7\\
38	1055.6\\
42	2315.6\\
27	727.45\\
46	2254.4\\
23	651.91\\
34	1779.8\\
200	40987\\
37	2145.9\\
16	620.53\\
50	2488.1\\
19	683.81\\
41	2752.5\\
25	802.5\\
39	1904.9\\
45	2679.8\\
83	7336.9\\
27	477.88\\
67	2495.2\\
46	3493.8\\
35	1444.7\\
44	2319\\
30	1943.4\\
29	1173.8\\
23	502.13\\
43	2063.1\\
28	711.1\\
24	1110.5\\
56	3015.5\\
33	939.39\\
13	387.66\\
41	2847.1\\
34	937.58\\
14	434.52\\
83	8798.4\\
33	1499.8\\
35	1486.1\\
29	908.1\\
54	3652.1\\
41	1290.2\\
19	763.71\\
34	2034.6\\
19	674.28\\
38	2256.2\\
43	1846.5\\
14	663.66\\
17	849.53\\
12	422.36\\
};
\addlegendentry{\ExactScatterLegend}

\addplot[only marks, mark=*, mark size=1.2748pt, color=black, fill=black, opacity=0.60, draw=none, mark options={draw=none,line width=0pt}] table[row sep=crcr]{%
x	y\\
34	1787\\
29	2052\\
96	39026\\
31	939\\
27	839\\
70	9435\\
24	756\\
37	2292\\
42	3353\\
16	509\\
29	1198\\
59	4061\\
69	4632\\
28	987\\
11	377\\
20	825\\
45	6367\\
27	1370\\
22	689\\
27	1107\\
44	4679\\
16	622\\
39	5236\\
19	558\\
33	1997\\
37	1291\\
33	1077\\
18	485\\
29	1640\\
22	758\\
33	1231\\
22	643\\
49	9267\\
25	1115\\
43	11048\\
26	1875\\
31	2762\\
26	927\\
18	665\\
13	140\\
32	2698\\
38	2211\\
32	775\\
29	1173\\
35	1719\\
35	3078\\
135	41095\\
27	2591\\
42	3900\\
39	3847\\
51	5607\\
86	11830\\
15	918\\
39	3169\\
42	3772\\
43	2489\\
58	4861\\
97	21264\\
50	3266\\
20	736\\
28	686\\
22	1102\\
47	4024\\
154	1.2234e+05\\
30	1543\\
10	296\\
37	7252\\
31	1142\\
54	13398\\
76	13899\\
55	3344\\
26	1499\\
33	1121\\
34	1971\\
25	1709\\
31	2961\\
34	3625\\
26	2774\\
39	3551\\
27	1922\\
32	1860\\
26	1781\\
29	1942\\
23	2014\\
26	1636\\
14	530\\
33	1466\\
24	787\\
20	1114\\
51	1692\\
35	2940\\
29	1332\\
46	1054\\
44	3996\\
45	4071\\
29	803\\
28	557\\
77	3928\\
70	3625\\
70	2531\\
40	1723\\
90	8518\\
53	3650\\
62	3213\\
47	8035\\
34	817\\
14	526\\
47	2742\\
43	759\\
20	950\\
28	533\\
80	4654\\
25	1153\\
36	1716\\
27	2299\\
82	3021\\
31	1103\\
39	1587\\
23	810\\
27	2388\\
37	2269\\
129	3241\\
63	3568\\
64	2459\\
91	6356\\
27	696\\
26	1459\\
17	558\\
33	2465\\
30	2046\\
33	1791\\
182	11424\\
26	663\\
48	2633\\
46	2157\\
10	222\\
25	802\\
26	424\\
206	37527\\
37	1069\\
43	1448\\
51	2715\\
57	7724\\
36	1800\\
38	1162\\
28	539\\
19	746\\
36	1003\\
45	3212\\
53	4779\\
48	5264\\
64	2484\\
65	3084\\
45	1597\\
54	3279\\
16	299\\
24	1131\\
127	2601\\
48	1783\\
64	2544\\
42	2057\\
42	1206\\
17	254\\
48	2507\\
42	2016\\
20	816\\
34	1805\\
47	1076\\
42	2516\\
50	2233\\
87	10377\\
64	2072\\
32	1299\\
39	1215\\
35	3735\\
18	866\\
28	626\\
27	253\\
37	2518\\
30	2667\\
89	9128\\
22	1204\\
28	1732\\
25	1513\\
15	873\\
41	3761\\
32	1140\\
34	1426\\
16	352\\
88	7596\\
39	668\\
140	20491\\
39	1086\\
49	1025\\
72	1475\\
37	1233\\
41	3002\\
283	2.0424e+05\\
31	2681\\
99	16428\\
103	32887\\
39	6629\\
35	1076\\
98	6332\\
45	4113\\
73	7877\\
34	1444\\
56	6926\\
39	2045\\
37	3241\\
55	1793\\
25	444\\
29	1020\\
42	1480\\
35	2286\\
26	591\\
45	5458\\
56	7320\\
39	1112\\
62	3063\\
94	13380\\
112	30134\\
25	317\\
57	1896\\
50	2648\\
47	5982\\
96	4539\\
54	4348\\
28	1451\\
38	3579\\
47	3124\\
35	1834\\
45	921\\
50	3106\\
43	719\\
29	1516\\
42	1249\\
71	5868\\
68	5857\\
149	61636\\
67	5463\\
51	1708\\
34	1459\\
120	10119\\
43	6342\\
117	8499\\
39	1041\\
72	6468\\
40	4061\\
43	1854\\
32	2516\\
49	12313\\
41	4009\\
98	4503\\
96	5379\\
45	1383\\
129	23237\\
43	1506\\
93	4228\\
41	2160\\
88	9837\\
73	5327\\
108	12182\\
52	5524\\
18	189\\
36	878\\
50	4430\\
37	1383\\
83	13086\\
40	2202\\
26	1338\\
40	1198\\
49	4702\\
27	1189\\
41	1263\\
36	572\\
52	3610\\
37	656\\
68	6317\\
43	2092\\
132	14984\\
73	3147\\
46	3825\\
47	739\\
44	1624\\
41	5048\\
30	1539\\
25	1574\\
24	2598\\
85	7294\\
44	1753\\
29	1546\\
34	1160\\
101	9599\\
18	1226\\
21	992\\
111	7348\\
95	2640\\
23	617\\
31	310\\
84	1691\\
23	466\\
55	3785\\
32	1352\\
27	544\\
48	517\\
18	119\\
28	1203\\
40	962\\
30	685\\
13	188\\
33	1762\\
14	209\\
52	1836\\
65	8994\\
49	3401\\
35	1176\\
109	4663\\
26	550\\
41	1006\\
39	973\\
61	4953\\
40	1238\\
36	1694\\
21	551\\
80	5029\\
36	1037\\
36	1691\\
38	832\\
42	1192\\
27	590\\
46	1533\\
23	486\\
34	1267\\
200	30990\\
37	1240\\
16	413\\
50	2688\\
19	324\\
41	2143\\
25	934\\
39	3741\\
45	2869\\
83	3945\\
27	521\\
67	2908\\
46	2419\\
35	1912\\
44	2676\\
30	2051\\
29	674\\
23	204\\
43	936\\
28	859\\
24	467\\
56	2173\\
33	679\\
13	261\\
41	2629\\
34	588\\
14	216\\
83	10336\\
33	1005\\
35	1380\\
29	687\\
54	3377\\
41	993\\
19	351\\
34	1937\\
19	459\\
38	1126\\
43	1196\\
14	412\\
17	325\\
12	184\\
};
\addlegendentry{\RealScatterLegend}

      % \input{../../../../.tex/20/KS/SizeVsDegreefig11.tex}
      \legend{}
  
      \nextgroupplot[%
        % xmin=10,xmax=283,
        % ymin=100,ymax=1e+06,
      ] % This file was created by matlab2tikz.
%
\definecolor{mycolor1}{rgb}{0.83529,0.36863,0.00000}%
\definecolor{mycolor2}{rgb}{0.00000,0.44700,0.74100}%
\definecolor{mycolor3}{rgb}{0.85000,0.32500,0.09800}%
%
\addplot[only marks, mark=*, mark size=1.2748pt, color=mycolor1, fill=mycolor1, opacity=0.60, draw=none, mark options={draw=none,line width=0pt}] table[row sep=crcr]{%
x	y\\
34	2744.6\\
29	1352.1\\
96	18508\\
31	1972.2\\
27	1412.6\\
70	7151.7\\
24	1269.3\\
37	2171.6\\
42	3051.1\\
16	1170.4\\
29	1405.5\\
59	3376.5\\
69	3130.8\\
28	1525.7\\
11	268.8\\
20	1192.7\\
45	3650.6\\
27	1533.6\\
22	928.61\\
27	1355.4\\
44	4883.6\\
16	1490.3\\
39	4095.5\\
19	1047.5\\
33	2135.3\\
37	2527.7\\
33	1356.8\\
18	499.27\\
29	2354.9\\
22	887.35\\
33	1420.9\\
22	883.81\\
49	8402.9\\
25	1528.2\\
43	3861.1\\
26	2452.3\\
31	3038.4\\
26	1031\\
18	807.01\\
13	447.17\\
32	4157.7\\
38	2228.7\\
32	3506.7\\
29	1701\\
35	1433\\
35	2953.1\\
135	20957\\
27	4295.8\\
42	3693\\
39	4104\\
51	7517.2\\
86	10965\\
15	831.26\\
39	3901.6\\
42	3162.6\\
43	3196.8\\
58	5643.4\\
97	25859\\
50	2350.4\\
20	1365.9\\
28	1124.5\\
22	1211.7\\
47	3544.5\\
154	63804\\
30	1775.9\\
10	466.4\\
37	7501\\
31	1925.3\\
54	14002\\
76	12012\\
55	4544.8\\
26	2573.9\\
33	1501.2\\
34	2342.5\\
25	1733.6\\
31	3605.5\\
34	4708.8\\
26	2373.9\\
39	3442.2\\
27	2578.4\\
32	1848.1\\
26	2159.1\\
29	2299.2\\
23	3127.2\\
26	2579.3\\
14	856.55\\
33	1945.3\\
24	1035.2\\
20	1869.3\\
51	2174\\
35	2792\\
29	1174.5\\
46	958.99\\
44	3639.5\\
45	3103.7\\
29	881.37\\
28	1671.6\\
77	2963.3\\
70	4125.8\\
70	3446.6\\
40	2173.3\\
90	6023.9\\
53	3407.9\\
62	1779.1\\
47	3496.1\\
34	1300.6\\
14	1042.2\\
47	1952.4\\
43	1042\\
20	754.06\\
28	1028.9\\
80	3155.7\\
25	517.75\\
36	1680.5\\
27	2118.6\\
82	4432.8\\
31	1080.9\\
39	1477.1\\
23	1975.1\\
27	2684.6\\
37	2122.1\\
129	8872.1\\
63	4710.7\\
64	1961.3\\
91	10283\\
27	1192.6\\
26	1577\\
17	702.8\\
33	1628\\
30	3113.6\\
33	1392.2\\
182	14724\\
26	1010.9\\
48	2990\\
46	1538.2\\
10	361.09\\
25	1059.4\\
26	830.75\\
206	26137\\
37	1267.5\\
43	1739.2\\
51	2187.3\\
57	6934.2\\
36	1695.1\\
38	1564.5\\
28	1355.6\\
19	1091.1\\
36	1415.1\\
45	3278.8\\
53	3606.6\\
48	3690.6\\
64	3053.8\\
65	3000.2\\
45	1135.7\\
54	2964.4\\
16	400.54\\
24	971.81\\
127	9079.8\\
48	1678.6\\
64	2526.7\\
42	2497.5\\
42	1250.3\\
17	521.01\\
48	2618.5\\
42	2156.2\\
20	858.82\\
34	1203.4\\
47	965.97\\
42	2998.2\\
50	2165.9\\
87	6388\\
64	3723\\
32	1782.5\\
39	1031.4\\
35	2494.7\\
18	1010.7\\
28	949.86\\
27	1024.8\\
37	2088.9\\
30	2568.3\\
89	10616\\
22	1770\\
28	1612.6\\
25	1198.2\\
15	596.95\\
41	3503.5\\
32	975.22\\
34	2029.2\\
16	759.71\\
88	6403.9\\
39	967.84\\
140	28322\\
39	1306.9\\
49	1666.1\\
72	2412.1\\
37	1620.3\\
41	3630.4\\
283	1.5536e+05\\
31	2751.1\\
99	20066\\
103	24294\\
39	3455.7\\
35	1566.2\\
98	13600\\
45	4492.5\\
73	8693.6\\
34	2183.5\\
56	7772.4\\
39	2865.7\\
37	2844\\
55	1957.3\\
25	723.62\\
29	1064\\
42	1057.1\\
35	4251.9\\
26	810.59\\
45	8453.8\\
56	6246.8\\
39	1542\\
62	2935.7\\
94	12348\\
112	23386\\
25	786.27\\
57	2286.6\\
50	2503\\
47	6654.5\\
96	7273.5\\
54	5725\\
28	2498.9\\
38	3759.5\\
47	3726.4\\
35	1937.4\\
45	1730.2\\
50	3570.4\\
43	949.61\\
29	1982.7\\
42	1725.5\\
71	19464\\
68	6767.5\\
149	53984\\
67	5713\\
51	2940.4\\
34	1900.5\\
120	14729\\
43	6841.9\\
117	11434\\
39	1318\\
72	8418.4\\
40	3417.4\\
43	2345.3\\
32	2507.7\\
49	8243.3\\
41	4138.4\\
98	4839.7\\
96	13039\\
45	1469.2\\
129	34913\\
43	2142.6\\
93	8535.5\\
41	3591.3\\
88	9223.9\\
73	6328.5\\
108	15883\\
52	8622.1\\
18	446.75\\
36	1021.5\\
50	5862.5\\
37	2240.8\\
83	15513\\
40	2127.7\\
26	2471.3\\
40	1647.3\\
49	3903.7\\
27	1488.9\\
41	1495.2\\
36	829.09\\
52	5132\\
37	1100.3\\
68	7100.7\\
43	2679\\
132	14748\\
73	2647.1\\
46	4140.5\\
47	1156.5\\
44	1704.9\\
41	4874.5\\
30	1410.5\\
25	2614.3\\
24	2035.5\\
85	8186.2\\
44	3142.4\\
29	1643.6\\
34	1548.8\\
101	14112\\
18	1667.7\\
21	1323.5\\
111	16852\\
95	4540.8\\
23	1008.6\\
31	658.26\\
84	4168.8\\
23	703.35\\
55	5117.4\\
32	1022.6\\
27	851.13\\
48	1047.7\\
18	322.04\\
28	1872.6\\
40	1161.4\\
30	1242.6\\
13	296.76\\
33	1544.5\\
14	553.78\\
52	2464.3\\
65	7334.8\\
49	2654.8\\
35	1283.4\\
109	7332.5\\
26	839.99\\
41	1210.9\\
39	1193.9\\
61	6446.7\\
40	1261.8\\
36	2108.7\\
21	617.27\\
80	4217.9\\
36	1425.3\\
36	1970.4\\
38	1162\\
42	2142.1\\
27	722.94\\
46	2098\\
23	634.69\\
34	1656.6\\
200	40055\\
37	2526.2\\
16	552.88\\
50	2415.2\\
19	617.72\\
41	3059.7\\
25	722.37\\
39	1922.9\\
45	3190.1\\
83	7306.8\\
27	528.05\\
67	2490.3\\
46	3060\\
35	1506.2\\
44	2339.9\\
30	1796.6\\
29	1120.7\\
23	537.31\\
43	1974.8\\
28	688.91\\
24	965.84\\
56	2955.4\\
33	985.06\\
13	420.52\\
41	3164.7\\
34	931.46\\
14	442.41\\
83	8658.3\\
33	1719.3\\
35	1458.9\\
29	962.91\\
54	3339\\
41	1258\\
19	731.95\\
34	1993.9\\
19	684.41\\
38	2234.5\\
43	1515.5\\
14	625.43\\
17	775.01\\
12	448.3\\
};
\addlegendentry{\ApproxScatterLegend}

\addplot[only marks, mark=*, mark size=1.2748pt, color=black, fill=black, opacity=0.60, draw=none, mark options={draw=none,line width=0pt}] table[row sep=crcr]{%
x	y\\
34	1787\\
29	2052\\
96	39026\\
31	939\\
27	839\\
70	9435\\
24	756\\
37	2292\\
42	3353\\
16	509\\
29	1198\\
59	4061\\
69	4632\\
28	987\\
11	377\\
20	825\\
45	6367\\
27	1370\\
22	689\\
27	1107\\
44	4679\\
16	622\\
39	5236\\
19	558\\
33	1997\\
37	1291\\
33	1077\\
18	485\\
29	1640\\
22	758\\
33	1231\\
22	643\\
49	9267\\
25	1115\\
43	11048\\
26	1875\\
31	2762\\
26	927\\
18	665\\
13	140\\
32	2698\\
38	2211\\
32	775\\
29	1173\\
35	1719\\
35	3078\\
135	41095\\
27	2591\\
42	3900\\
39	3847\\
51	5607\\
86	11830\\
15	918\\
39	3169\\
42	3772\\
43	2489\\
58	4861\\
97	21264\\
50	3266\\
20	736\\
28	686\\
22	1102\\
47	4024\\
154	1.2234e+05\\
30	1543\\
10	296\\
37	7252\\
31	1142\\
54	13398\\
76	13899\\
55	3344\\
26	1499\\
33	1121\\
34	1971\\
25	1709\\
31	2961\\
34	3625\\
26	2774\\
39	3551\\
27	1922\\
32	1860\\
26	1781\\
29	1942\\
23	2014\\
26	1636\\
14	530\\
33	1466\\
24	787\\
20	1114\\
51	1692\\
35	2940\\
29	1332\\
46	1054\\
44	3996\\
45	4071\\
29	803\\
28	557\\
77	3928\\
70	3625\\
70	2531\\
40	1723\\
90	8518\\
53	3650\\
62	3213\\
47	8035\\
34	817\\
14	526\\
47	2742\\
43	759\\
20	950\\
28	533\\
80	4654\\
25	1153\\
36	1716\\
27	2299\\
82	3021\\
31	1103\\
39	1587\\
23	810\\
27	2388\\
37	2269\\
129	3241\\
63	3568\\
64	2459\\
91	6356\\
27	696\\
26	1459\\
17	558\\
33	2465\\
30	2046\\
33	1791\\
182	11424\\
26	663\\
48	2633\\
46	2157\\
10	222\\
25	802\\
26	424\\
206	37527\\
37	1069\\
43	1448\\
51	2715\\
57	7724\\
36	1800\\
38	1162\\
28	539\\
19	746\\
36	1003\\
45	3212\\
53	4779\\
48	5264\\
64	2484\\
65	3084\\
45	1597\\
54	3279\\
16	299\\
24	1131\\
127	2601\\
48	1783\\
64	2544\\
42	2057\\
42	1206\\
17	254\\
48	2507\\
42	2016\\
20	816\\
34	1805\\
47	1076\\
42	2516\\
50	2233\\
87	10377\\
64	2072\\
32	1299\\
39	1215\\
35	3735\\
18	866\\
28	626\\
27	253\\
37	2518\\
30	2667\\
89	9128\\
22	1204\\
28	1732\\
25	1513\\
15	873\\
41	3761\\
32	1140\\
34	1426\\
16	352\\
88	7596\\
39	668\\
140	20491\\
39	1086\\
49	1025\\
72	1475\\
37	1233\\
41	3002\\
283	2.0424e+05\\
31	2681\\
99	16428\\
103	32887\\
39	6629\\
35	1076\\
98	6332\\
45	4113\\
73	7877\\
34	1444\\
56	6926\\
39	2045\\
37	3241\\
55	1793\\
25	444\\
29	1020\\
42	1480\\
35	2286\\
26	591\\
45	5458\\
56	7320\\
39	1112\\
62	3063\\
94	13380\\
112	30134\\
25	317\\
57	1896\\
50	2648\\
47	5982\\
96	4539\\
54	4348\\
28	1451\\
38	3579\\
47	3124\\
35	1834\\
45	921\\
50	3106\\
43	719\\
29	1516\\
42	1249\\
71	5868\\
68	5857\\
149	61636\\
67	5463\\
51	1708\\
34	1459\\
120	10119\\
43	6342\\
117	8499\\
39	1041\\
72	6468\\
40	4061\\
43	1854\\
32	2516\\
49	12313\\
41	4009\\
98	4503\\
96	5379\\
45	1383\\
129	23237\\
43	1506\\
93	4228\\
41	2160\\
88	9837\\
73	5327\\
108	12182\\
52	5524\\
18	189\\
36	878\\
50	4430\\
37	1383\\
83	13086\\
40	2202\\
26	1338\\
40	1198\\
49	4702\\
27	1189\\
41	1263\\
36	572\\
52	3610\\
37	656\\
68	6317\\
43	2092\\
132	14984\\
73	3147\\
46	3825\\
47	739\\
44	1624\\
41	5048\\
30	1539\\
25	1574\\
24	2598\\
85	7294\\
44	1753\\
29	1546\\
34	1160\\
101	9599\\
18	1226\\
21	992\\
111	7348\\
95	2640\\
23	617\\
31	310\\
84	1691\\
23	466\\
55	3785\\
32	1352\\
27	544\\
48	517\\
18	119\\
28	1203\\
40	962\\
30	685\\
13	188\\
33	1762\\
14	209\\
52	1836\\
65	8994\\
49	3401\\
35	1176\\
109	4663\\
26	550\\
41	1006\\
39	973\\
61	4953\\
40	1238\\
36	1694\\
21	551\\
80	5029\\
36	1037\\
36	1691\\
38	832\\
42	1192\\
27	590\\
46	1533\\
23	486\\
34	1267\\
200	30990\\
37	1240\\
16	413\\
50	2688\\
19	324\\
41	2143\\
25	934\\
39	3741\\
45	2869\\
83	3945\\
27	521\\
67	2908\\
46	2419\\
35	1912\\
44	2676\\
30	2051\\
29	674\\
23	204\\
43	936\\
28	859\\
24	467\\
56	2173\\
33	679\\
13	261\\
41	2629\\
34	588\\
14	216\\
83	10336\\
33	1005\\
35	1380\\
29	687\\
54	3377\\
41	993\\
19	351\\
34	1937\\
19	459\\
38	1126\\
43	1196\\
14	412\\
17	325\\
12	184\\
};
\addlegendentry{\RealScatterLegend}

      % \input{../../../../.tex/20/KS/SizeVsDegreefig12.tex}
      \legend{}
    \end{groupplot}
    
    \begin{groupplot}[
      group style={
        group name=Row5,
        plotLastGroup
      },
      plotRow,
      row4Keys,
      plotLegendNW,
      %
      xmin=10,xmax=283,
      ymin=100,ymax=2.0424e+05,
    ]
      \nextgroupplot[%
        at={($(Row4 c1r1.south west)-(0,\yPlotSep em)-(0,\plotHeight cm)$)},
        % xmin=10,xmax=283,
        % ymin=100,ymax=1.6258e+05,
      ] % This file was created by matlab2tikz.
%
\definecolor{mycolor1}{rgb}{0.00000,0.44706,0.69804}%
\definecolor{mycolor2}{rgb}{0.00000,0.44700,0.74100}%
\definecolor{mycolor3}{rgb}{0.83529,0.36863,0.00000}%
\definecolor{mycolor4}{rgb}{0.85000,0.32500,0.09800}%
%
\addplot[only marks, mark=*, mark size=1.2748pt, color=mycolor1, fill=mycolor1, opacity=0.70, draw=none, mark options={draw=none,line width=0pt}] table[row sep=crcr]{%
x	y\\
34	2835.8\\
29	1281.2\\
96	16961\\
31	1907.1\\
27	1358.1\\
70	8228.2\\
24	1023.2\\
37	2120.4\\
42	3357\\
16	1072.2\\
29	1411.4\\
59	3467.7\\
69	3466.8\\
28	1349.4\\
11	316.13\\
20	1070.8\\
45	3977.5\\
27	1536.7\\
22	871.49\\
27	1434\\
44	5278.1\\
16	1348\\
39	4325.8\\
19	1140.6\\
33	2101.3\\
37	2463.1\\
33	1434.5\\
18	509.3\\
29	2538\\
22	966.5\\
33	1429.2\\
22	832.49\\
49	7861.9\\
25	1360.6\\
43	3705.2\\
26	2188.9\\
31	2808.2\\
26	1184.7\\
18	805.62\\
13	521.98\\
32	3942.8\\
38	2310.9\\
32	3332.7\\
29	1798\\
35	1589.4\\
35	3581.4\\
135	20513\\
27	4839.1\\
42	3923.8\\
39	4344.3\\
51	7579.7\\
86	10754\\
15	934.12\\
39	4357.7\\
42	3136.9\\
43	2987.6\\
58	5422.8\\
97	25829\\
50	2426.5\\
20	1251\\
28	1004\\
22	1077.1\\
47	3121.1\\
154	66260\\
30	1883.3\\
10	436.9\\
37	7561.7\\
31	1600.6\\
54	13758\\
76	12042\\
55	4804.4\\
26	2869.2\\
33	1433.6\\
34	2040.1\\
25	1710.4\\
31	3745.2\\
34	4730.6\\
26	2544.7\\
39	3602.9\\
27	2425.9\\
32	2074.1\\
26	2166.6\\
29	2406.6\\
23	2864.9\\
26	2225\\
14	724.95\\
33	2079.3\\
24	998.46\\
20	1488.4\\
51	1954.6\\
35	2911.2\\
29	1121.4\\
46	1013.6\\
44	3473.3\\
45	3141.6\\
29	980.02\\
28	1661.3\\
77	3030.2\\
70	3792.1\\
70	3492.7\\
40	2418.9\\
90	5416.7\\
53	3478.5\\
62	1639.1\\
47	3257.8\\
34	1229.4\\
14	1038.9\\
47	1870\\
43	1147.9\\
20	604.25\\
28	1069.4\\
80	2667.7\\
25	663.12\\
36	1750.1\\
27	2085.9\\
82	4072.5\\
31	946.95\\
39	1346.6\\
23	2212.3\\
27	2513.2\\
37	2457.5\\
129	8010.4\\
63	4342.1\\
64	1611.2\\
91	9850.3\\
27	1179\\
26	1282.1\\
17	742.34\\
33	1756.6\\
30	3368.8\\
33	1323.1\\
182	14432\\
26	1285\\
48	2838.8\\
46	1685.1\\
10	355.78\\
25	1044.5\\
26	661.88\\
206	24244\\
37	1098.6\\
43	1628.1\\
51	1818.5\\
57	6702.2\\
36	1767.4\\
38	1397.4\\
28	1335.8\\
19	1007.3\\
36	1473.6\\
45	2956.3\\
53	3572.4\\
48	3273.6\\
64	3089.7\\
65	2550.8\\
45	1043.8\\
54	2718.1\\
16	407.28\\
24	915.09\\
127	9395.8\\
48	1743.8\\
64	2407.1\\
42	2286.1\\
42	1100.4\\
17	561.03\\
48	2615.4\\
42	1986.4\\
20	762.42\\
34	1110.5\\
47	852.54\\
42	2885.8\\
50	2252.9\\
87	6718.1\\
64	3830.7\\
32	1710.1\\
39	1060.5\\
35	2293.6\\
18	951.81\\
28	807.13\\
27	888.74\\
37	1959.7\\
30	2330.3\\
89	12218\\
22	1820.9\\
28	1688.3\\
25	1235.3\\
15	593.73\\
41	3410.6\\
32	963.18\\
34	2059.8\\
16	711.15\\
88	5805.5\\
39	918.27\\
140	29149\\
39	1196.1\\
49	1680.6\\
72	2256.9\\
37	1709.1\\
41	3789.1\\
283	1.6258e+05\\
31	2570.6\\
99	19833\\
103	24083\\
39	3341.9\\
35	1415.3\\
98	14011\\
45	4659.6\\
73	7607.6\\
34	2402.8\\
56	7386.4\\
39	3202\\
37	2471.2\\
55	2327.8\\
25	753.73\\
29	1089.6\\
42	1133\\
35	3572.4\\
26	810.74\\
45	9117.4\\
56	5776.3\\
39	1481.3\\
62	3052.1\\
94	11541\\
112	24173\\
25	724.52\\
57	2267.3\\
50	2355\\
47	5943.5\\
96	6895.5\\
54	6073\\
28	2399.2\\
38	3685.8\\
47	3825.6\\
35	2212\\
45	1728.6\\
50	3652.8\\
43	938.12\\
29	2185.8\\
42	1863.9\\
71	20994\\
68	6198.2\\
149	57232\\
67	5275.3\\
51	2886.9\\
34	1880.6\\
120	14662\\
43	6350\\
117	12253\\
39	1377.6\\
72	8240.2\\
40	3189.4\\
43	2398.6\\
32	2304.8\\
49	9169.2\\
41	3863.2\\
98	4751.3\\
96	12537\\
45	1541.4\\
129	36097\\
43	2108.6\\
93	8352.9\\
41	3471.4\\
88	9257.6\\
73	6097\\
108	16282\\
52	8132.4\\
18	405.2\\
36	1056\\
50	5046.3\\
37	2260.8\\
83	15018\\
40	2486.6\\
26	2632.4\\
40	1414.9\\
49	3602\\
27	1205.2\\
41	1606.3\\
36	851.62\\
52	5228.1\\
37	1079\\
68	7623.2\\
43	2553\\
132	14101\\
73	2614.9\\
46	4257.7\\
47	1154.3\\
44	1586.3\\
41	5046.5\\
30	1314.6\\
25	2589\\
24	1950.3\\
85	8088.8\\
44	2760.4\\
29	1739.8\\
34	1492.6\\
101	14930\\
18	1598.9\\
21	1323.7\\
111	18155\\
95	4612.3\\
23	1065.9\\
31	607.24\\
84	4542.5\\
23	718.69\\
55	5654.9\\
32	1034.1\\
27	890.59\\
48	893.33\\
18	351.35\\
28	1878\\
40	1128.8\\
30	1025.5\\
13	302.16\\
33	1416.4\\
14	581.63\\
52	2385.8\\
65	7891.9\\
49	2193\\
35	1143.4\\
109	7519.4\\
26	786.2\\
41	1215\\
39	1223.7\\
61	6183.9\\
40	1292.7\\
36	1752.5\\
21	550\\
80	4196.5\\
36	1314.7\\
36	1680.6\\
38	1040.5\\
42	1858\\
27	643.41\\
46	2160.9\\
23	616.98\\
34	1636.2\\
200	38589\\
37	2263.3\\
16	617.21\\
50	2618.5\\
19	643.01\\
41	2982.5\\
25	786.71\\
39	2046.6\\
45	2793.6\\
83	6919.5\\
27	515.44\\
67	2584.9\\
46	3098.1\\
35	1406.7\\
44	2341.3\\
30	1782\\
29	1103.3\\
23	435.08\\
43	1717.7\\
28	776.53\\
24	981.93\\
56	3002.3\\
33	928.35\\
13	380.26\\
41	2886.5\\
34	1016.5\\
14	434.9\\
83	8765.7\\
33	1574\\
35	1461.2\\
29	971.3\\
54	3871.7\\
41	1304.7\\
19	628.11\\
34	1944.2\\
19	704.73\\
38	2018.7\\
43	1560.3\\
14	621.29\\
17	684.97\\
12	342.44\\
};
\addlegendentry{\ExactScatterLegend}

\addplot[only marks, mark=*, mark size=1.2748pt, color=mycolor3, fill=mycolor3, opacity=0.70, draw=none, mark options={draw=none,line width=0pt}] table[row sep=crcr]{%
x	y\\
34	2635.4\\
29	1223.2\\
96	18601\\
31	1877.2\\
27	1180.5\\
70	7587.5\\
24	1189.5\\
37	2391.8\\
42	2859.7\\
16	1059.4\\
29	1460.5\\
59	3479.2\\
69	3250.9\\
28	1479.2\\
11	297.17\\
20	1109.8\\
45	3928.6\\
27	1453.5\\
22	900.83\\
27	1493.4\\
44	5467.2\\
16	1270.6\\
39	4500.3\\
19	985.84\\
33	2402.3\\
37	2282.8\\
33	1344.8\\
18	434.16\\
29	2435.6\\
22	877.91\\
33	1419\\
22	882.09\\
49	9326.2\\
25	1470.1\\
43	3503.3\\
26	1911.2\\
31	3011.8\\
26	1080.9\\
18	788.02\\
13	499.81\\
32	4373.8\\
38	2359.2\\
32	3225.4\\
29	1651.4\\
35	1498.7\\
35	3399.5\\
135	20772\\
27	4137.3\\
42	3553.4\\
39	4324.2\\
51	7533.1\\
86	10979\\
15	815.27\\
39	4469.9\\
42	3304\\
43	3151.7\\
58	5850.3\\
97	27365\\
50	2355.8\\
20	1370.1\\
28	1231.9\\
22	1076.7\\
47	3106.5\\
154	64536\\
30	1686.2\\
10	452.29\\
37	8028.2\\
31	1972.6\\
54	13233\\
76	11599\\
55	5243.8\\
26	2483.1\\
33	1520.4\\
34	1959.8\\
25	1337.7\\
31	3740.7\\
34	4795.3\\
26	2161.9\\
39	3280.7\\
27	2613.3\\
32	1851.4\\
26	2145.5\\
29	2215.7\\
23	2839\\
26	2453.2\\
14	818.48\\
33	1855.7\\
24	866.54\\
20	1610.2\\
51	1916.7\\
35	2912.1\\
29	1361.1\\
46	853.09\\
44	3632.5\\
45	3567.5\\
29	781.34\\
28	1679.9\\
77	2729.4\\
70	3996.9\\
70	3369.9\\
40	2476\\
90	5201.7\\
53	3222.9\\
62	1458.2\\
47	3762.7\\
34	1136.1\\
14	1111.6\\
47	1787\\
43	1113\\
20	755.91\\
28	1095.8\\
80	2957.1\\
25	465.96\\
36	1622.9\\
27	2447.6\\
82	4557\\
31	981.52\\
39	1245.7\\
23	1991.5\\
27	2884.1\\
37	2217.1\\
129	9038\\
63	4524\\
64	1756.3\\
91	9810.9\\
27	1302.6\\
26	1244.2\\
17	773.4\\
33	1425\\
30	3323.6\\
33	1288\\
182	15066\\
26	895.54\\
48	2896.1\\
46	1705.9\\
10	362.84\\
25	1007\\
26	803.17\\
206	26142\\
37	1242.6\\
43	1479\\
51	1792.2\\
57	6133\\
36	1604.2\\
38	1546.2\\
28	1180\\
19	870.8\\
36	1601.1\\
45	3565.8\\
53	3855.4\\
48	3603.9\\
64	2737.7\\
65	2936.2\\
45	1185\\
54	2848.7\\
16	335.12\\
24	974.21\\
127	8274.9\\
48	1706.7\\
64	2706.4\\
42	2883.7\\
42	1279.2\\
17	445.16\\
48	3092.2\\
42	2150.2\\
20	820.71\\
34	1084.7\\
47	942.27\\
42	2971.3\\
50	2394.1\\
87	6226.2\\
64	3770.2\\
32	1642.6\\
39	1138.8\\
35	2471.6\\
18	923.64\\
28	824.67\\
27	907.09\\
37	1954.8\\
30	2420.6\\
89	10457\\
22	1561.6\\
28	1519.5\\
25	1333.2\\
15	533.96\\
41	3697.8\\
32	823.76\\
34	1873\\
16	772.42\\
88	5999.8\\
39	821.18\\
140	26954\\
39	1200.2\\
49	1758.9\\
72	2302.6\\
37	1659.7\\
41	3786.3\\
283	1.5508e+05\\
31	3021.9\\
99	20645\\
103	24986\\
39	3384\\
35	1491\\
98	14801\\
45	4404.6\\
73	9409\\
34	2150.2\\
56	8742.8\\
39	2611.4\\
37	2769.9\\
55	2215.4\\
25	593.85\\
29	1064\\
42	1177\\
35	4160.9\\
26	752.94\\
45	9726.3\\
56	5876.6\\
39	1392.3\\
62	3361.8\\
94	11691\\
112	25231\\
25	747.81\\
57	2335.1\\
50	2270.8\\
47	7085.8\\
96	6217.7\\
54	6481\\
28	2645\\
38	3696.4\\
47	3138.6\\
35	1959.4\\
45	1550.7\\
50	3602.3\\
43	1008.4\\
29	2173.2\\
42	1620.9\\
71	20431\\
68	6653.6\\
149	57012\\
67	5330.4\\
51	2393.5\\
34	1846.2\\
120	15878\\
43	6955.5\\
117	11528\\
39	1348.4\\
72	8482.6\\
40	3589.9\\
43	2304.4\\
32	2445.4\\
49	8819\\
41	3492.6\\
98	4856.4\\
96	12916\\
45	1487.7\\
129	37223\\
43	2195\\
93	7587.9\\
41	3537.6\\
88	8946.5\\
73	5132.5\\
108	16178\\
52	7523.9\\
18	405.58\\
36	986.77\\
50	5431.1\\
37	2203.2\\
83	14728\\
40	2143.6\\
26	2532.3\\
40	1566.7\\
49	3952.6\\
27	1430.5\\
41	1326.8\\
36	815.54\\
52	4681.7\\
37	998.42\\
68	7174.6\\
43	2492\\
132	15295\\
73	2899\\
46	4605\\
47	1127.4\\
44	1654.4\\
41	4434.2\\
30	1445\\
25	2921\\
24	1903.4\\
85	7414.9\\
44	3274.9\\
29	1715.2\\
34	1531.3\\
101	14465\\
18	1616\\
21	1361.4\\
111	17177\\
95	4824.5\\
23	871.66\\
31	640.1\\
84	3600.1\\
23	693.85\\
55	4952.4\\
32	1120\\
27	961.59\\
48	951.78\\
18	321.63\\
28	1868.7\\
40	1046.4\\
30	1146.9\\
13	336.26\\
33	1492.8\\
14	463.36\\
52	2182.2\\
65	7415\\
49	2526.7\\
35	1355.4\\
109	7722.8\\
26	824.04\\
41	1064.6\\
39	1122.5\\
61	6460.6\\
40	1310.4\\
36	1862.3\\
21	595.41\\
80	4382.5\\
36	1290\\
36	1789.2\\
38	1137.8\\
42	2226.3\\
27	606.61\\
46	1776.8\\
23	562.25\\
34	1842.5\\
200	40006\\
37	2588.3\\
16	569.63\\
50	2424.7\\
19	584.86\\
41	2862.2\\
25	745.66\\
39	1698.8\\
45	3012.5\\
83	7406.4\\
27	495.73\\
67	2634.1\\
46	3087.5\\
35	1679.4\\
44	2248.2\\
30	1760.2\\
29	1101.9\\
23	545.08\\
43	1930.1\\
28	795.4\\
24	1004.8\\
56	2766.9\\
33	883.61\\
13	304.86\\
41	2973.1\\
34	827.16\\
14	378.24\\
83	8853\\
33	1451\\
35	1396.9\\
29	1021.5\\
54	3541.8\\
41	1368.6\\
19	582.68\\
34	1805.8\\
19	607.18\\
38	2015.1\\
43	1495.6\\
14	524.12\\
17	659.58\\
12	392.53\\
};
\addlegendentry{\ApproxScatterLegend}

      % \input{../../../../.tex/20/KS/SizeVsDegreefig13.tex}
      \legend{}
  
      \nextgroupplot[%
        % xmin=10,xmax=283,
        % ymin=100,ymax=1e+06,
      ] % This file was created by matlab2tikz.
%
\definecolor{mycolor1}{rgb}{0.00000,0.44706,0.69804}%
\definecolor{mycolor2}{rgb}{0.00000,0.44700,0.74100}%
\definecolor{mycolor3}{rgb}{0.85000,0.32500,0.09800}%
%
\addplot[only marks, mark=*, mark size=1.2748pt, color=mycolor1, fill=mycolor1, opacity=0.60, draw=none, mark options={draw=none,line width=0pt}] table[row sep=crcr]{%
x	y\\
34	2835.8\\
29	1281.2\\
96	16961\\
31	1907.1\\
27	1358.1\\
70	8228.2\\
24	1023.2\\
37	2120.4\\
42	3357\\
16	1072.2\\
29	1411.4\\
59	3467.7\\
69	3466.8\\
28	1349.4\\
11	316.13\\
20	1070.8\\
45	3977.5\\
27	1536.7\\
22	871.49\\
27	1434\\
44	5278.1\\
16	1348\\
39	4325.8\\
19	1140.6\\
33	2101.3\\
37	2463.1\\
33	1434.5\\
18	509.3\\
29	2538\\
22	966.5\\
33	1429.2\\
22	832.49\\
49	7861.9\\
25	1360.6\\
43	3705.2\\
26	2188.9\\
31	2808.2\\
26	1184.7\\
18	805.62\\
13	521.98\\
32	3942.8\\
38	2310.9\\
32	3332.7\\
29	1798\\
35	1589.4\\
35	3581.4\\
135	20513\\
27	4839.1\\
42	3923.8\\
39	4344.3\\
51	7579.7\\
86	10754\\
15	934.12\\
39	4357.7\\
42	3136.9\\
43	2987.6\\
58	5422.8\\
97	25829\\
50	2426.5\\
20	1251\\
28	1004\\
22	1077.1\\
47	3121.1\\
154	66260\\
30	1883.3\\
10	436.9\\
37	7561.7\\
31	1600.6\\
54	13758\\
76	12042\\
55	4804.4\\
26	2869.2\\
33	1433.6\\
34	2040.1\\
25	1710.4\\
31	3745.2\\
34	4730.6\\
26	2544.7\\
39	3602.9\\
27	2425.9\\
32	2074.1\\
26	2166.6\\
29	2406.6\\
23	2864.9\\
26	2225\\
14	724.95\\
33	2079.3\\
24	998.46\\
20	1488.4\\
51	1954.6\\
35	2911.2\\
29	1121.4\\
46	1013.6\\
44	3473.3\\
45	3141.6\\
29	980.02\\
28	1661.3\\
77	3030.2\\
70	3792.1\\
70	3492.7\\
40	2418.9\\
90	5416.7\\
53	3478.5\\
62	1639.1\\
47	3257.8\\
34	1229.4\\
14	1038.9\\
47	1870\\
43	1147.9\\
20	604.25\\
28	1069.4\\
80	2667.7\\
25	663.12\\
36	1750.1\\
27	2085.9\\
82	4072.5\\
31	946.95\\
39	1346.6\\
23	2212.3\\
27	2513.2\\
37	2457.5\\
129	8010.4\\
63	4342.1\\
64	1611.2\\
91	9850.3\\
27	1179\\
26	1282.1\\
17	742.34\\
33	1756.6\\
30	3368.8\\
33	1323.1\\
182	14432\\
26	1285\\
48	2838.8\\
46	1685.1\\
10	355.78\\
25	1044.5\\
26	661.88\\
206	24244\\
37	1098.6\\
43	1628.1\\
51	1818.5\\
57	6702.2\\
36	1767.4\\
38	1397.4\\
28	1335.8\\
19	1007.3\\
36	1473.6\\
45	2956.3\\
53	3572.4\\
48	3273.6\\
64	3089.7\\
65	2550.8\\
45	1043.8\\
54	2718.1\\
16	407.28\\
24	915.09\\
127	9395.8\\
48	1743.8\\
64	2407.1\\
42	2286.1\\
42	1100.4\\
17	561.03\\
48	2615.4\\
42	1986.4\\
20	762.42\\
34	1110.5\\
47	852.54\\
42	2885.8\\
50	2252.9\\
87	6718.1\\
64	3830.7\\
32	1710.1\\
39	1060.5\\
35	2293.6\\
18	951.81\\
28	807.13\\
27	888.74\\
37	1959.7\\
30	2330.3\\
89	12218\\
22	1820.9\\
28	1688.3\\
25	1235.3\\
15	593.73\\
41	3410.6\\
32	963.18\\
34	2059.8\\
16	711.15\\
88	5805.5\\
39	918.27\\
140	29149\\
39	1196.1\\
49	1680.6\\
72	2256.9\\
37	1709.1\\
41	3789.1\\
283	1.6258e+05\\
31	2570.6\\
99	19833\\
103	24083\\
39	3341.9\\
35	1415.3\\
98	14011\\
45	4659.6\\
73	7607.6\\
34	2402.8\\
56	7386.4\\
39	3202\\
37	2471.2\\
55	2327.8\\
25	753.73\\
29	1089.6\\
42	1133\\
35	3572.4\\
26	810.74\\
45	9117.4\\
56	5776.3\\
39	1481.3\\
62	3052.1\\
94	11541\\
112	24173\\
25	724.52\\
57	2267.3\\
50	2355\\
47	5943.5\\
96	6895.5\\
54	6073\\
28	2399.2\\
38	3685.8\\
47	3825.6\\
35	2212\\
45	1728.6\\
50	3652.8\\
43	938.12\\
29	2185.8\\
42	1863.9\\
71	20994\\
68	6198.2\\
149	57232\\
67	5275.3\\
51	2886.9\\
34	1880.6\\
120	14662\\
43	6350\\
117	12253\\
39	1377.6\\
72	8240.2\\
40	3189.4\\
43	2398.6\\
32	2304.8\\
49	9169.2\\
41	3863.2\\
98	4751.3\\
96	12537\\
45	1541.4\\
129	36097\\
43	2108.6\\
93	8352.9\\
41	3471.4\\
88	9257.6\\
73	6097\\
108	16282\\
52	8132.4\\
18	405.2\\
36	1056\\
50	5046.3\\
37	2260.8\\
83	15018\\
40	2486.6\\
26	2632.4\\
40	1414.9\\
49	3602\\
27	1205.2\\
41	1606.3\\
36	851.62\\
52	5228.1\\
37	1079\\
68	7623.2\\
43	2553\\
132	14101\\
73	2614.9\\
46	4257.7\\
47	1154.3\\
44	1586.3\\
41	5046.5\\
30	1314.6\\
25	2589\\
24	1950.3\\
85	8088.8\\
44	2760.4\\
29	1739.8\\
34	1492.6\\
101	14930\\
18	1598.9\\
21	1323.7\\
111	18155\\
95	4612.3\\
23	1065.9\\
31	607.24\\
84	4542.5\\
23	718.69\\
55	5654.9\\
32	1034.1\\
27	890.59\\
48	893.33\\
18	351.35\\
28	1878\\
40	1128.8\\
30	1025.5\\
13	302.16\\
33	1416.4\\
14	581.63\\
52	2385.8\\
65	7891.9\\
49	2193\\
35	1143.4\\
109	7519.4\\
26	786.2\\
41	1215\\
39	1223.7\\
61	6183.9\\
40	1292.7\\
36	1752.5\\
21	550\\
80	4196.5\\
36	1314.7\\
36	1680.6\\
38	1040.5\\
42	1858\\
27	643.41\\
46	2160.9\\
23	616.98\\
34	1636.2\\
200	38589\\
37	2263.3\\
16	617.21\\
50	2618.5\\
19	643.01\\
41	2982.5\\
25	786.71\\
39	2046.6\\
45	2793.6\\
83	6919.5\\
27	515.44\\
67	2584.9\\
46	3098.1\\
35	1406.7\\
44	2341.3\\
30	1782\\
29	1103.3\\
23	435.08\\
43	1717.7\\
28	776.53\\
24	981.93\\
56	3002.3\\
33	928.35\\
13	380.26\\
41	2886.5\\
34	1016.5\\
14	434.9\\
83	8765.7\\
33	1574\\
35	1461.2\\
29	971.3\\
54	3871.7\\
41	1304.7\\
19	628.11\\
34	1944.2\\
19	704.73\\
38	2018.7\\
43	1560.3\\
14	621.29\\
17	684.97\\
12	342.44\\
};
\addlegendentry{\ExactScatterLegend}

\addplot[only marks, mark=*, mark size=1.2748pt, color=black, fill=black, opacity=0.60, draw=none, mark options={draw=none,line width=0pt}] table[row sep=crcr]{%
x	y\\
34	1787\\
29	2052\\
96	39026\\
31	939\\
27	839\\
70	9435\\
24	756\\
37	2292\\
42	3353\\
16	509\\
29	1198\\
59	4061\\
69	4632\\
28	987\\
11	377\\
20	825\\
45	6367\\
27	1370\\
22	689\\
27	1107\\
44	4679\\
16	622\\
39	5236\\
19	558\\
33	1997\\
37	1291\\
33	1077\\
18	485\\
29	1640\\
22	758\\
33	1231\\
22	643\\
49	9267\\
25	1115\\
43	11048\\
26	1875\\
31	2762\\
26	927\\
18	665\\
13	140\\
32	2698\\
38	2211\\
32	775\\
29	1173\\
35	1719\\
35	3078\\
135	41095\\
27	2591\\
42	3900\\
39	3847\\
51	5607\\
86	11830\\
15	918\\
39	3169\\
42	3772\\
43	2489\\
58	4861\\
97	21264\\
50	3266\\
20	736\\
28	686\\
22	1102\\
47	4024\\
154	1.2234e+05\\
30	1543\\
10	296\\
37	7252\\
31	1142\\
54	13398\\
76	13899\\
55	3344\\
26	1499\\
33	1121\\
34	1971\\
25	1709\\
31	2961\\
34	3625\\
26	2774\\
39	3551\\
27	1922\\
32	1860\\
26	1781\\
29	1942\\
23	2014\\
26	1636\\
14	530\\
33	1466\\
24	787\\
20	1114\\
51	1692\\
35	2940\\
29	1332\\
46	1054\\
44	3996\\
45	4071\\
29	803\\
28	557\\
77	3928\\
70	3625\\
70	2531\\
40	1723\\
90	8518\\
53	3650\\
62	3213\\
47	8035\\
34	817\\
14	526\\
47	2742\\
43	759\\
20	950\\
28	533\\
80	4654\\
25	1153\\
36	1716\\
27	2299\\
82	3021\\
31	1103\\
39	1587\\
23	810\\
27	2388\\
37	2269\\
129	3241\\
63	3568\\
64	2459\\
91	6356\\
27	696\\
26	1459\\
17	558\\
33	2465\\
30	2046\\
33	1791\\
182	11424\\
26	663\\
48	2633\\
46	2157\\
10	222\\
25	802\\
26	424\\
206	37527\\
37	1069\\
43	1448\\
51	2715\\
57	7724\\
36	1800\\
38	1162\\
28	539\\
19	746\\
36	1003\\
45	3212\\
53	4779\\
48	5264\\
64	2484\\
65	3084\\
45	1597\\
54	3279\\
16	299\\
24	1131\\
127	2601\\
48	1783\\
64	2544\\
42	2057\\
42	1206\\
17	254\\
48	2507\\
42	2016\\
20	816\\
34	1805\\
47	1076\\
42	2516\\
50	2233\\
87	10377\\
64	2072\\
32	1299\\
39	1215\\
35	3735\\
18	866\\
28	626\\
27	253\\
37	2518\\
30	2667\\
89	9128\\
22	1204\\
28	1732\\
25	1513\\
15	873\\
41	3761\\
32	1140\\
34	1426\\
16	352\\
88	7596\\
39	668\\
140	20491\\
39	1086\\
49	1025\\
72	1475\\
37	1233\\
41	3002\\
283	2.0424e+05\\
31	2681\\
99	16428\\
103	32887\\
39	6629\\
35	1076\\
98	6332\\
45	4113\\
73	7877\\
34	1444\\
56	6926\\
39	2045\\
37	3241\\
55	1793\\
25	444\\
29	1020\\
42	1480\\
35	2286\\
26	591\\
45	5458\\
56	7320\\
39	1112\\
62	3063\\
94	13380\\
112	30134\\
25	317\\
57	1896\\
50	2648\\
47	5982\\
96	4539\\
54	4348\\
28	1451\\
38	3579\\
47	3124\\
35	1834\\
45	921\\
50	3106\\
43	719\\
29	1516\\
42	1249\\
71	5868\\
68	5857\\
149	61636\\
67	5463\\
51	1708\\
34	1459\\
120	10119\\
43	6342\\
117	8499\\
39	1041\\
72	6468\\
40	4061\\
43	1854\\
32	2516\\
49	12313\\
41	4009\\
98	4503\\
96	5379\\
45	1383\\
129	23237\\
43	1506\\
93	4228\\
41	2160\\
88	9837\\
73	5327\\
108	12182\\
52	5524\\
18	189\\
36	878\\
50	4430\\
37	1383\\
83	13086\\
40	2202\\
26	1338\\
40	1198\\
49	4702\\
27	1189\\
41	1263\\
36	572\\
52	3610\\
37	656\\
68	6317\\
43	2092\\
132	14984\\
73	3147\\
46	3825\\
47	739\\
44	1624\\
41	5048\\
30	1539\\
25	1574\\
24	2598\\
85	7294\\
44	1753\\
29	1546\\
34	1160\\
101	9599\\
18	1226\\
21	992\\
111	7348\\
95	2640\\
23	617\\
31	310\\
84	1691\\
23	466\\
55	3785\\
32	1352\\
27	544\\
48	517\\
18	119\\
28	1203\\
40	962\\
30	685\\
13	188\\
33	1762\\
14	209\\
52	1836\\
65	8994\\
49	3401\\
35	1176\\
109	4663\\
26	550\\
41	1006\\
39	973\\
61	4953\\
40	1238\\
36	1694\\
21	551\\
80	5029\\
36	1037\\
36	1691\\
38	832\\
42	1192\\
27	590\\
46	1533\\
23	486\\
34	1267\\
200	30990\\
37	1240\\
16	413\\
50	2688\\
19	324\\
41	2143\\
25	934\\
39	3741\\
45	2869\\
83	3945\\
27	521\\
67	2908\\
46	2419\\
35	1912\\
44	2676\\
30	2051\\
29	674\\
23	204\\
43	936\\
28	859\\
24	467\\
56	2173\\
33	679\\
13	261\\
41	2629\\
34	588\\
14	216\\
83	10336\\
33	1005\\
35	1380\\
29	687\\
54	3377\\
41	993\\
19	351\\
34	1937\\
19	459\\
38	1126\\
43	1196\\
14	412\\
17	325\\
12	184\\
};
\addlegendentry{\RealScatterLegend}

      % \input{../../../../.tex/20/KS/SizeVsDegreefig14.tex}
      \legend{}
  
      \nextgroupplot[%
        % xmin=10,xmax=283,
        % ymin=100,ymax=1e+06,
      ] % This file was created by matlab2tikz.
%
\definecolor{mycolor1}{rgb}{0.83529,0.36863,0.00000}%
\definecolor{mycolor2}{rgb}{0.00000,0.44700,0.74100}%
\definecolor{mycolor3}{rgb}{0.85000,0.32500,0.09800}%
%
\addplot[only marks, mark=*, mark size=1.2748pt, color=mycolor1, fill=mycolor1, opacity=0.60, draw=none, mark options={draw=none,line width=0pt}] table[row sep=crcr]{%
x	y\\
34	2635.4\\
29	1223.2\\
96	18601\\
31	1877.2\\
27	1180.5\\
70	7587.5\\
24	1189.5\\
37	2391.8\\
42	2859.7\\
16	1059.4\\
29	1460.5\\
59	3479.2\\
69	3250.9\\
28	1479.2\\
11	297.17\\
20	1109.8\\
45	3928.6\\
27	1453.5\\
22	900.83\\
27	1493.4\\
44	5467.2\\
16	1270.6\\
39	4500.3\\
19	985.84\\
33	2402.3\\
37	2282.8\\
33	1344.8\\
18	434.16\\
29	2435.6\\
22	877.91\\
33	1419\\
22	882.09\\
49	9326.2\\
25	1470.1\\
43	3503.3\\
26	1911.2\\
31	3011.8\\
26	1080.9\\
18	788.02\\
13	499.81\\
32	4373.8\\
38	2359.2\\
32	3225.4\\
29	1651.4\\
35	1498.7\\
35	3399.5\\
135	20772\\
27	4137.3\\
42	3553.4\\
39	4324.2\\
51	7533.1\\
86	10979\\
15	815.27\\
39	4469.9\\
42	3304\\
43	3151.7\\
58	5850.3\\
97	27365\\
50	2355.8\\
20	1370.1\\
28	1231.9\\
22	1076.7\\
47	3106.5\\
154	64536\\
30	1686.2\\
10	452.29\\
37	8028.2\\
31	1972.6\\
54	13233\\
76	11599\\
55	5243.8\\
26	2483.1\\
33	1520.4\\
34	1959.8\\
25	1337.7\\
31	3740.7\\
34	4795.3\\
26	2161.9\\
39	3280.7\\
27	2613.3\\
32	1851.4\\
26	2145.5\\
29	2215.7\\
23	2839\\
26	2453.2\\
14	818.48\\
33	1855.7\\
24	866.54\\
20	1610.2\\
51	1916.7\\
35	2912.1\\
29	1361.1\\
46	853.09\\
44	3632.5\\
45	3567.5\\
29	781.34\\
28	1679.9\\
77	2729.4\\
70	3996.9\\
70	3369.9\\
40	2476\\
90	5201.7\\
53	3222.9\\
62	1458.2\\
47	3762.7\\
34	1136.1\\
14	1111.6\\
47	1787\\
43	1113\\
20	755.91\\
28	1095.8\\
80	2957.1\\
25	465.96\\
36	1622.9\\
27	2447.6\\
82	4557\\
31	981.52\\
39	1245.7\\
23	1991.5\\
27	2884.1\\
37	2217.1\\
129	9038\\
63	4524\\
64	1756.3\\
91	9810.9\\
27	1302.6\\
26	1244.2\\
17	773.4\\
33	1425\\
30	3323.6\\
33	1288\\
182	15066\\
26	895.54\\
48	2896.1\\
46	1705.9\\
10	362.84\\
25	1007\\
26	803.17\\
206	26142\\
37	1242.6\\
43	1479\\
51	1792.2\\
57	6133\\
36	1604.2\\
38	1546.2\\
28	1180\\
19	870.8\\
36	1601.1\\
45	3565.8\\
53	3855.4\\
48	3603.9\\
64	2737.7\\
65	2936.2\\
45	1185\\
54	2848.7\\
16	335.12\\
24	974.21\\
127	8274.9\\
48	1706.7\\
64	2706.4\\
42	2883.7\\
42	1279.2\\
17	445.16\\
48	3092.2\\
42	2150.2\\
20	820.71\\
34	1084.7\\
47	942.27\\
42	2971.3\\
50	2394.1\\
87	6226.2\\
64	3770.2\\
32	1642.6\\
39	1138.8\\
35	2471.6\\
18	923.64\\
28	824.67\\
27	907.09\\
37	1954.8\\
30	2420.6\\
89	10457\\
22	1561.6\\
28	1519.5\\
25	1333.2\\
15	533.96\\
41	3697.8\\
32	823.76\\
34	1873\\
16	772.42\\
88	5999.8\\
39	821.18\\
140	26954\\
39	1200.2\\
49	1758.9\\
72	2302.6\\
37	1659.7\\
41	3786.3\\
283	1.5508e+05\\
31	3021.9\\
99	20645\\
103	24986\\
39	3384\\
35	1491\\
98	14801\\
45	4404.6\\
73	9409\\
34	2150.2\\
56	8742.8\\
39	2611.4\\
37	2769.9\\
55	2215.4\\
25	593.85\\
29	1064\\
42	1177\\
35	4160.9\\
26	752.94\\
45	9726.3\\
56	5876.6\\
39	1392.3\\
62	3361.8\\
94	11691\\
112	25231\\
25	747.81\\
57	2335.1\\
50	2270.8\\
47	7085.8\\
96	6217.7\\
54	6481\\
28	2645\\
38	3696.4\\
47	3138.6\\
35	1959.4\\
45	1550.7\\
50	3602.3\\
43	1008.4\\
29	2173.2\\
42	1620.9\\
71	20431\\
68	6653.6\\
149	57012\\
67	5330.4\\
51	2393.5\\
34	1846.2\\
120	15878\\
43	6955.5\\
117	11528\\
39	1348.4\\
72	8482.6\\
40	3589.9\\
43	2304.4\\
32	2445.4\\
49	8819\\
41	3492.6\\
98	4856.4\\
96	12916\\
45	1487.7\\
129	37223\\
43	2195\\
93	7587.9\\
41	3537.6\\
88	8946.5\\
73	5132.5\\
108	16178\\
52	7523.9\\
18	405.58\\
36	986.77\\
50	5431.1\\
37	2203.2\\
83	14728\\
40	2143.6\\
26	2532.3\\
40	1566.7\\
49	3952.6\\
27	1430.5\\
41	1326.8\\
36	815.54\\
52	4681.7\\
37	998.42\\
68	7174.6\\
43	2492\\
132	15295\\
73	2899\\
46	4605\\
47	1127.4\\
44	1654.4\\
41	4434.2\\
30	1445\\
25	2921\\
24	1903.4\\
85	7414.9\\
44	3274.9\\
29	1715.2\\
34	1531.3\\
101	14465\\
18	1616\\
21	1361.4\\
111	17177\\
95	4824.5\\
23	871.66\\
31	640.1\\
84	3600.1\\
23	693.85\\
55	4952.4\\
32	1120\\
27	961.59\\
48	951.78\\
18	321.63\\
28	1868.7\\
40	1046.4\\
30	1146.9\\
13	336.26\\
33	1492.8\\
14	463.36\\
52	2182.2\\
65	7415\\
49	2526.7\\
35	1355.4\\
109	7722.8\\
26	824.04\\
41	1064.6\\
39	1122.5\\
61	6460.6\\
40	1310.4\\
36	1862.3\\
21	595.41\\
80	4382.5\\
36	1290\\
36	1789.2\\
38	1137.8\\
42	2226.3\\
27	606.61\\
46	1776.8\\
23	562.25\\
34	1842.5\\
200	40006\\
37	2588.3\\
16	569.63\\
50	2424.7\\
19	584.86\\
41	2862.2\\
25	745.66\\
39	1698.8\\
45	3012.5\\
83	7406.4\\
27	495.73\\
67	2634.1\\
46	3087.5\\
35	1679.4\\
44	2248.2\\
30	1760.2\\
29	1101.9\\
23	545.08\\
43	1930.1\\
28	795.4\\
24	1004.8\\
56	2766.9\\
33	883.61\\
13	304.86\\
41	2973.1\\
34	827.16\\
14	378.24\\
83	8853\\
33	1451\\
35	1396.9\\
29	1021.5\\
54	3541.8\\
41	1368.6\\
19	582.68\\
34	1805.8\\
19	607.18\\
38	2015.1\\
43	1495.6\\
14	524.12\\
17	659.58\\
12	392.53\\
};
\addlegendentry{\ApproxScatterLegend}

\addplot[only marks, mark=*, mark size=1.2748pt, color=black, fill=black, opacity=0.60, draw=none, mark options={draw=none,line width=0pt}] table[row sep=crcr]{%
x	y\\
34	1787\\
29	2052\\
96	39026\\
31	939\\
27	839\\
70	9435\\
24	756\\
37	2292\\
42	3353\\
16	509\\
29	1198\\
59	4061\\
69	4632\\
28	987\\
11	377\\
20	825\\
45	6367\\
27	1370\\
22	689\\
27	1107\\
44	4679\\
16	622\\
39	5236\\
19	558\\
33	1997\\
37	1291\\
33	1077\\
18	485\\
29	1640\\
22	758\\
33	1231\\
22	643\\
49	9267\\
25	1115\\
43	11048\\
26	1875\\
31	2762\\
26	927\\
18	665\\
13	140\\
32	2698\\
38	2211\\
32	775\\
29	1173\\
35	1719\\
35	3078\\
135	41095\\
27	2591\\
42	3900\\
39	3847\\
51	5607\\
86	11830\\
15	918\\
39	3169\\
42	3772\\
43	2489\\
58	4861\\
97	21264\\
50	3266\\
20	736\\
28	686\\
22	1102\\
47	4024\\
154	1.2234e+05\\
30	1543\\
10	296\\
37	7252\\
31	1142\\
54	13398\\
76	13899\\
55	3344\\
26	1499\\
33	1121\\
34	1971\\
25	1709\\
31	2961\\
34	3625\\
26	2774\\
39	3551\\
27	1922\\
32	1860\\
26	1781\\
29	1942\\
23	2014\\
26	1636\\
14	530\\
33	1466\\
24	787\\
20	1114\\
51	1692\\
35	2940\\
29	1332\\
46	1054\\
44	3996\\
45	4071\\
29	803\\
28	557\\
77	3928\\
70	3625\\
70	2531\\
40	1723\\
90	8518\\
53	3650\\
62	3213\\
47	8035\\
34	817\\
14	526\\
47	2742\\
43	759\\
20	950\\
28	533\\
80	4654\\
25	1153\\
36	1716\\
27	2299\\
82	3021\\
31	1103\\
39	1587\\
23	810\\
27	2388\\
37	2269\\
129	3241\\
63	3568\\
64	2459\\
91	6356\\
27	696\\
26	1459\\
17	558\\
33	2465\\
30	2046\\
33	1791\\
182	11424\\
26	663\\
48	2633\\
46	2157\\
10	222\\
25	802\\
26	424\\
206	37527\\
37	1069\\
43	1448\\
51	2715\\
57	7724\\
36	1800\\
38	1162\\
28	539\\
19	746\\
36	1003\\
45	3212\\
53	4779\\
48	5264\\
64	2484\\
65	3084\\
45	1597\\
54	3279\\
16	299\\
24	1131\\
127	2601\\
48	1783\\
64	2544\\
42	2057\\
42	1206\\
17	254\\
48	2507\\
42	2016\\
20	816\\
34	1805\\
47	1076\\
42	2516\\
50	2233\\
87	10377\\
64	2072\\
32	1299\\
39	1215\\
35	3735\\
18	866\\
28	626\\
27	253\\
37	2518\\
30	2667\\
89	9128\\
22	1204\\
28	1732\\
25	1513\\
15	873\\
41	3761\\
32	1140\\
34	1426\\
16	352\\
88	7596\\
39	668\\
140	20491\\
39	1086\\
49	1025\\
72	1475\\
37	1233\\
41	3002\\
283	2.0424e+05\\
31	2681\\
99	16428\\
103	32887\\
39	6629\\
35	1076\\
98	6332\\
45	4113\\
73	7877\\
34	1444\\
56	6926\\
39	2045\\
37	3241\\
55	1793\\
25	444\\
29	1020\\
42	1480\\
35	2286\\
26	591\\
45	5458\\
56	7320\\
39	1112\\
62	3063\\
94	13380\\
112	30134\\
25	317\\
57	1896\\
50	2648\\
47	5982\\
96	4539\\
54	4348\\
28	1451\\
38	3579\\
47	3124\\
35	1834\\
45	921\\
50	3106\\
43	719\\
29	1516\\
42	1249\\
71	5868\\
68	5857\\
149	61636\\
67	5463\\
51	1708\\
34	1459\\
120	10119\\
43	6342\\
117	8499\\
39	1041\\
72	6468\\
40	4061\\
43	1854\\
32	2516\\
49	12313\\
41	4009\\
98	4503\\
96	5379\\
45	1383\\
129	23237\\
43	1506\\
93	4228\\
41	2160\\
88	9837\\
73	5327\\
108	12182\\
52	5524\\
18	189\\
36	878\\
50	4430\\
37	1383\\
83	13086\\
40	2202\\
26	1338\\
40	1198\\
49	4702\\
27	1189\\
41	1263\\
36	572\\
52	3610\\
37	656\\
68	6317\\
43	2092\\
132	14984\\
73	3147\\
46	3825\\
47	739\\
44	1624\\
41	5048\\
30	1539\\
25	1574\\
24	2598\\
85	7294\\
44	1753\\
29	1546\\
34	1160\\
101	9599\\
18	1226\\
21	992\\
111	7348\\
95	2640\\
23	617\\
31	310\\
84	1691\\
23	466\\
55	3785\\
32	1352\\
27	544\\
48	517\\
18	119\\
28	1203\\
40	962\\
30	685\\
13	188\\
33	1762\\
14	209\\
52	1836\\
65	8994\\
49	3401\\
35	1176\\
109	4663\\
26	550\\
41	1006\\
39	973\\
61	4953\\
40	1238\\
36	1694\\
21	551\\
80	5029\\
36	1037\\
36	1691\\
38	832\\
42	1192\\
27	590\\
46	1533\\
23	486\\
34	1267\\
200	30990\\
37	1240\\
16	413\\
50	2688\\
19	324\\
41	2143\\
25	934\\
39	3741\\
45	2869\\
83	3945\\
27	521\\
67	2908\\
46	2419\\
35	1912\\
44	2676\\
30	2051\\
29	674\\
23	204\\
43	936\\
28	859\\
24	467\\
56	2173\\
33	679\\
13	261\\
41	2629\\
34	588\\
14	216\\
83	10336\\
33	1005\\
35	1380\\
29	687\\
54	3377\\
41	993\\
19	351\\
34	1937\\
19	459\\
38	1126\\
43	1196\\
14	412\\
17	325\\
12	184\\
};
\addlegendentry{\RealScatterLegend}

      % \input{../../../../.tex/20/KS/SizeVsDegreefig15.tex}
      \legend{}
    \end{groupplot}
  %\endgroup

  %\begingroup Parametri
    
%\begingroup Parametro
\providecommand\varyingParameterSymbol\lambda
\providecommand\varyingParameterList[1]{
  \ifcase#1\relax%
  \or\num{1e-2}% Caso 1
  \or\num{8.2e-2}% Caso 2
  \or\num{1.55e-1}% Caso 3
  \or\num{2.27e-1}% Caso 4
  \else\num{3e-1}% Caso 5
  \fi
}
%\endgroup

%\begingroup beta
\providecommand\betaExactList[1]{
  \ifcase#1\relax%
  \or\num{1.267+-0.007}% Caso 1
  \or\num{1.285+-0.010}% Caso 2
  \or\num{1.264+-0.015}% Caso 3
  \or\num{1.246+-0.015}% Caso 4
  \else\num{1.210+-0.015}% Caso 5
  \fi
}
\providecommand\betaAprxList[1]{
  \ifcase#1\relax%
  \or\num{1.264+-0.006}% Caso 1
  \or\num{1.270+-0.011}% Caso 2
  \or\num{1.280+-0.013}% Caso 3
  \or\num{1.250+-0.015}% Caso 4
  \else\num{1.218+-0.014}% Caso 5
  \fi
}
%\endgroup

%\begingroup xi
\providecommand\zetaExactList[1]{
  \ifcase#1\relax%
  \or\num{0.757+-0.004}% Caso 1
  \or\num{0.710+-0.026}% Caso 2
  \or\num{0.670+-0.035}% Caso 3
  \or\num{0.625+-0.052}% Caso 4
  \else\num{0.539+-0.059}% Caso 5
  \fi
}
\providecommand\zetaAprxList[1]{
  \ifcase#1\relax%
  \or\num{0.749+-0.004}% Caso 1
  \or\num{0.721+-0.022}% Caso 2
  \or\num{0.645+-0.040}% Caso 3
  \or\num{0.622+-0.051}% Caso 4
  \else\num{0.567+-0.053}% Caso 5
  \fi
}
%\endgroup

%\begingroup mI
\providecommand\mIExactList[1]{
  \ifcase#1\relax%
  \or\num{3.265+-0.001}% Caso 1
  \or\num{3.238+-0.005}% Caso 2
  \or\num{3.220+-0.008}% Caso 3
  \or\num{3.213+-0.013}% Caso 4
  \else\num{3.174+-0.023}% Caso 5
  \fi
}
\providecommand\mIAprxList[1]{
  \ifcase#1\relax%
  \or\num{3.262+-0.001}% Caso 1
  \or\num{3.240+-0.005}% Caso 2
  \or\num{3.217+-0.010}% Caso 3
  \or\num{3.217+-0.012}% Caso 4
  \else\num{3.208+-0.009}% Caso 5
  \fi
}
%\endgroup

%\begingroup sI
\providecommand\sIExactList[1]{
  \ifcase#1\relax%
  \or\num{0.284+-0.001}% Caso 1
  \or\num{0.311+-0.008}% Caso 2
  \or\num{0.343+-0.010}% Caso 3
  \or\num{0.361+-0.020}% Caso 4
  \else\num{0.398+-0.029}% Caso 5
  \fi
}
\providecommand\sIAprxList[1]{
  \ifcase#1\relax%
  \or\num{0.281+-0.001}% Caso 1
  \or\num{0.315+-0.006}% Caso 2
  \or\num{0.334+-0.013}% Caso 3
  \or\num{0.366+-0.018}% Caso 4
  \else\num{0.411+-0.025}% Caso 5
  \fi
}
%\endgroup

%\begingroup mII
\providecommand\mIIExactList[1]{
  \ifcase#1\relax%
  \or\num{3.788+-0.005}% Caso 1
  \or\num{3.730+-0.029}% Caso 2
  \or\num{3.681+-0.052}% Caso 3
  \or\num{3.650+-0.072}% Caso 4
  \else\num{3.572+-0.095}% Caso 5
  \fi
}
\providecommand\mIIAprxList[1]{
  \ifcase#1\relax%
  \or\num{3.778+-0.006}% Caso 1
  \or\num{3.735+-0.029}% Caso 2
  \or\num{3.656+-0.048}% Caso 3
  \or\num{3.658+-0.084}% Caso 4
  \else\num{3.502+-0.078}% Caso 5
  \fi
}
%\endgroup

%\begingroup sII
\providecommand\sIIExactList[1]{
  \ifcase#1\relax%
  \or\num{0.455+-0.002}% Caso 1
  \or\num{0.486+-0.007}% Caso 2
  \or\num{0.515+-0.014}% Caso 3
  \or\num{0.533+-0.021}% Caso 4
  \else\num{0.570+-0.028}% Caso 5
  \fi
}
\providecommand\sIIAprxList[1]{
  \ifcase#1\relax%
  \or\num{0.457+-0.002}% Caso 1
  \or\num{0.490+-0.007}% Caso 2
  \or\num{0.521+-0.012}% Caso 3
  \or\num{0.545+-0.021}% Caso 4
  \else\num{0.586+-0.029}% Caso 5
  \fi
}
%\endgroup

    \ifpageisodd\writeVaryingParameter{1}{west}{west}{0em}{1em}
    \else\writeVaryingParameter{3}{east}{east}{0em}{-1.5em}\fi
  %\endgroup

  %\begingroup Didascalie
    \writeSubCaptions{south east}{south east}{0em}{0em}
    % https://tex.stackexchange.com/questions/632045/get-subcaption-for-every-plot-using-groupplot-pgf-tikz-after-2020-update
  %\endgroup

  \redefineTikZbounds{Row1 c1r1}{Row5 c3r1}
\end{tikzpicture}%

% \input{Parametri.txt}
% 
%\begingroup Parametro
\providecommand\varyingParameterSymbol\lambda
\providecommand\varyingParameterList[1]{
  \ifcase#1\relax%
  \or\num{1e-2}% Caso 1
  \or\num{8.2e-2}% Caso 2
  \or\num{1.55e-1}% Caso 3
  \or\num{2.27e-1}% Caso 4
  \else\num{3e-1}% Caso 5
  \fi
}
%\endgroup

%\begingroup beta
\providecommand\betaExactList[1]{
  \ifcase#1\relax%
  \or\num{1.267+-0.007}% Caso 1
  \or\num{1.285+-0.010}% Caso 2
  \or\num{1.264+-0.015}% Caso 3
  \or\num{1.246+-0.015}% Caso 4
  \else\num{1.210+-0.015}% Caso 5
  \fi
}
\providecommand\betaAprxList[1]{
  \ifcase#1\relax%
  \or\num{1.264+-0.006}% Caso 1
  \or\num{1.270+-0.011}% Caso 2
  \or\num{1.280+-0.013}% Caso 3
  \or\num{1.250+-0.015}% Caso 4
  \else\num{1.218+-0.014}% Caso 5
  \fi
}
%\endgroup

%\begingroup xi
\providecommand\zetaExactList[1]{
  \ifcase#1\relax%
  \or\num{0.757+-0.004}% Caso 1
  \or\num{0.710+-0.026}% Caso 2
  \or\num{0.670+-0.035}% Caso 3
  \or\num{0.625+-0.052}% Caso 4
  \else\num{0.539+-0.059}% Caso 5
  \fi
}
\providecommand\zetaAprxList[1]{
  \ifcase#1\relax%
  \or\num{0.749+-0.004}% Caso 1
  \or\num{0.721+-0.022}% Caso 2
  \or\num{0.645+-0.040}% Caso 3
  \or\num{0.622+-0.051}% Caso 4
  \else\num{0.567+-0.053}% Caso 5
  \fi
}
%\endgroup

%\begingroup mI
\providecommand\mIExactList[1]{
  \ifcase#1\relax%
  \or\num{3.265+-0.001}% Caso 1
  \or\num{3.238+-0.005}% Caso 2
  \or\num{3.220+-0.008}% Caso 3
  \or\num{3.213+-0.013}% Caso 4
  \else\num{3.174+-0.023}% Caso 5
  \fi
}
\providecommand\mIAprxList[1]{
  \ifcase#1\relax%
  \or\num{3.262+-0.001}% Caso 1
  \or\num{3.240+-0.005}% Caso 2
  \or\num{3.217+-0.010}% Caso 3
  \or\num{3.217+-0.012}% Caso 4
  \else\num{3.208+-0.009}% Caso 5
  \fi
}
%\endgroup

%\begingroup sI
\providecommand\sIExactList[1]{
  \ifcase#1\relax%
  \or\num{0.284+-0.001}% Caso 1
  \or\num{0.311+-0.008}% Caso 2
  \or\num{0.343+-0.010}% Caso 3
  \or\num{0.361+-0.020}% Caso 4
  \else\num{0.398+-0.029}% Caso 5
  \fi
}
\providecommand\sIAprxList[1]{
  \ifcase#1\relax%
  \or\num{0.281+-0.001}% Caso 1
  \or\num{0.315+-0.006}% Caso 2
  \or\num{0.334+-0.013}% Caso 3
  \or\num{0.366+-0.018}% Caso 4
  \else\num{0.411+-0.025}% Caso 5
  \fi
}
%\endgroup

%\begingroup mII
\providecommand\mIIExactList[1]{
  \ifcase#1\relax%
  \or\num{3.788+-0.005}% Caso 1
  \or\num{3.730+-0.029}% Caso 2
  \or\num{3.681+-0.052}% Caso 3
  \or\num{3.650+-0.072}% Caso 4
  \else\num{3.572+-0.095}% Caso 5
  \fi
}
\providecommand\mIIAprxList[1]{
  \ifcase#1\relax%
  \or\num{3.778+-0.006}% Caso 1
  \or\num{3.735+-0.029}% Caso 2
  \or\num{3.656+-0.048}% Caso 3
  \or\num{3.658+-0.084}% Caso 4
  \else\num{3.502+-0.078}% Caso 5
  \fi
}
%\endgroup

%\begingroup sII
\providecommand\sIIExactList[1]{
  \ifcase#1\relax%
  \or\num{0.455+-0.002}% Caso 1
  \or\num{0.486+-0.007}% Caso 2
  \or\num{0.515+-0.014}% Caso 3
  \or\num{0.533+-0.021}% Caso 4
  \else\num{0.570+-0.028}% Caso 5
  \fi
}
\providecommand\sIIAprxList[1]{
  \ifcase#1\relax%
  \or\num{0.457+-0.002}% Caso 1
  \or\num{0.490+-0.007}% Caso 2
  \or\num{0.521+-0.012}% Caso 3
  \or\num{0.545+-0.021}% Caso 4
  \else\num{0.586+-0.029}% Caso 5
  \fi
}
%\endgroup


  \endgroup

  \begingroup
    \setcounter{subfigure}{0}
    
% !TEX root = ../../../../../Esperimenti/.tex/MEF.tex
% LTeX: language=it

\infopageKeysI

\begin{tikzpicture}%
  %\begingroup Righe
    \begin{groupplot}[
      group style={
        group name=Row1,
        plotGroup,
        group size=2 by 1,
      },
      plotRow,
      row5Keys,
      xTickScaleLabelEmpty,
      plotLegendNE,
      %
      xmin=0,xmax=1e+05,
      ymin=100,ymax=1.4472e+05,
    ]
      \nextgroupplot[%
        % xmin=0,xmax=1e+05,
        % ymin=100,ymax=1.4403e+05,
      ] % This file was created by matlab2tikz.
%
\definecolor{mycolor1}{rgb}{0.00146,0.00047,0.01387}%
\definecolor{mycolor2}{rgb}{0.01166,0.00942,0.06346}%
\definecolor{mycolor3}{rgb}{0.02943,0.02150,0.11462}%
\definecolor{mycolor4}{rgb}{0.06134,0.03659,0.17764}%
\definecolor{mycolor5}{rgb}{0.08741,0.04456,0.22481}%
\definecolor{mycolor6}{rgb}{0.12291,0.04754,0.28162}%
\definecolor{mycolor7}{rgb}{0.14907,0.04547,0.31709}%
\definecolor{mycolor8}{rgb}{0.18343,0.04033,0.35497}%
\definecolor{mycolor9}{rgb}{0.21795,0.03662,0.38352}%
\definecolor{mycolor10}{rgb}{0.24497,0.03705,0.40001}%
\definecolor{mycolor11}{rgb}{0.27135,0.04092,0.41198}%
\definecolor{mycolor12}{rgb}{0.29076,0.04564,0.41864}%
\definecolor{mycolor13}{rgb}{0.31628,0.05349,0.42512}%
\definecolor{mycolor14}{rgb}{0.33522,0.06006,0.42852}%
\definecolor{mycolor15}{rgb}{0.36028,0.06925,0.43150}%
\definecolor{mycolor16}{rgb}{0.37900,0.07625,0.43272}%
\definecolor{mycolor17}{rgb}{0.39767,0.08326,0.43318}%
\definecolor{mycolor18}{rgb}{0.41633,0.09020,0.43294}%
\definecolor{mycolor19}{rgb}{0.42877,0.09479,0.43241}%
\definecolor{mycolor20}{rgb}{0.44743,0.10160,0.43108}%
\definecolor{mycolor21}{rgb}{0.46610,0.10832,0.42912}%
\definecolor{mycolor22}{rgb}{0.47856,0.11276,0.42747}%
\definecolor{mycolor23}{rgb}{0.49726,0.11938,0.42449}%
\definecolor{mycolor24}{rgb}{0.50973,0.12377,0.42216}%
\definecolor{mycolor25}{rgb}{0.52221,0.12815,0.41955}%
\definecolor{mycolor26}{rgb}{0.53468,0.13253,0.41667}%
\definecolor{mycolor27}{rgb}{0.55339,0.13913,0.41183}%
\definecolor{mycolor28}{rgb}{0.56585,0.14357,0.40826}%
\definecolor{mycolor29}{rgb}{0.57830,0.14804,0.40441}%
\definecolor{mycolor30}{rgb}{0.59073,0.15256,0.40029}%
\definecolor{mycolor31}{rgb}{0.60314,0.15715,0.39589}%
\definecolor{mycolor32}{rgb}{0.61551,0.16182,0.39122}%
\definecolor{mycolor33}{rgb}{0.62169,0.16418,0.38878}%
\definecolor{mycolor34}{rgb}{0.63400,0.16899,0.38370}%
\definecolor{mycolor35}{rgb}{0.64626,0.17391,0.37836}%
\definecolor{mycolor36}{rgb}{0.65846,0.17896,0.37275}%
\definecolor{mycolor37}{rgb}{0.66454,0.18154,0.36985}%
\definecolor{mycolor38}{rgb}{0.67664,0.18681,0.36385}%
\definecolor{mycolor39}{rgb}{0.68865,0.19224,0.35760}%
\definecolor{mycolor40}{rgb}{0.69463,0.19502,0.35439}%
\definecolor{mycolor41}{rgb}{0.70650,0.20073,0.34778}%
\definecolor{mycolor42}{rgb}{0.71240,0.20366,0.34438}%
\definecolor{mycolor43}{rgb}{0.72410,0.20967,0.33742}%
\definecolor{mycolor44}{rgb}{0.72991,0.21276,0.33386}%
\definecolor{mycolor45}{rgb}{0.74142,0.21911,0.32658}%
\definecolor{mycolor46}{rgb}{0.74713,0.22238,0.32286}%
\definecolor{mycolor47}{rgb}{0.75279,0.22571,0.31909}%
\definecolor{mycolor48}{rgb}{0.76401,0.23255,0.31140}%
\definecolor{mycolor49}{rgb}{0.76956,0.23608,0.30749}%
\definecolor{mycolor50}{rgb}{0.77506,0.23967,0.30353}%
\definecolor{mycolor51}{rgb}{0.79129,0.25086,0.29139}%
\definecolor{mycolor52}{rgb}{0.79661,0.25473,0.28726}%
\definecolor{mycolor53}{rgb}{0.80187,0.25867,0.28310}%
\definecolor{mycolor54}{rgb}{0.81224,0.26679,0.27466}%
\definecolor{mycolor55}{rgb}{0.81734,0.27095,0.27039}%
\definecolor{mycolor56}{rgb}{0.82737,0.27952,0.26175}%
\definecolor{mycolor57}{rgb}{0.83230,0.28391,0.25738}%
\definecolor{mycolor58}{rgb}{0.83717,0.28839,0.25299}%
\definecolor{mycolor59}{rgb}{0.84671,0.29756,0.24411}%
\definecolor{mycolor60}{rgb}{0.85138,0.30226,0.23964}%
\definecolor{mycolor61}{rgb}{0.85599,0.30704,0.23513}%
\definecolor{mycolor62}{rgb}{0.86053,0.31189,0.23061}%
\definecolor{mycolor63}{rgb}{0.87374,0.32691,0.21689}%
\definecolor{mycolor64}{rgb}{0.87800,0.33206,0.21227}%
\definecolor{mycolor65}{rgb}{0.89034,0.34796,0.19829}%
\definecolor{mycolor66}{rgb}{0.89819,0.35891,0.18886}%
\definecolor{mycolor67}{rgb}{0.90200,0.36449,0.18412}%
\definecolor{mycolor68}{rgb}{0.90573,0.37014,0.17935}%
\definecolor{mycolor69}{rgb}{0.90939,0.37586,0.17456}%
\definecolor{mycolor70}{rgb}{0.91297,0.38164,0.16975}%
\definecolor{mycolor71}{rgb}{0.91646,0.38748,0.16492}%
\definecolor{mycolor72}{rgb}{0.91988,0.39339,0.16007}%
\definecolor{mycolor73}{rgb}{0.92322,0.39936,0.15519}%
\definecolor{mycolor74}{rgb}{0.92647,0.40539,0.15029}%
\definecolor{mycolor75}{rgb}{0.92964,0.41148,0.14537}%
\definecolor{mycolor76}{rgb}{0.93274,0.41763,0.14042}%
\definecolor{mycolor77}{rgb}{0.93575,0.42383,0.13544}%
\definecolor{mycolor78}{rgb}{0.93868,0.43009,0.13044}%
\definecolor{mycolor79}{rgb}{0.94152,0.43640,0.12541}%
\definecolor{mycolor80}{rgb}{0.94429,0.44277,0.12035}%
\definecolor{mycolor81}{rgb}{0.94956,0.45566,0.11016}%
\definecolor{mycolor82}{rgb}{0.95208,0.46218,0.10503}%
\definecolor{mycolor83}{rgb}{0.96129,0.48872,0.08429}%
\definecolor{mycolor84}{rgb}{0.96540,0.50225,0.07386}%
\definecolor{mycolor85}{rgb}{0.96732,0.50908,0.06866}%
\definecolor{mycolor86}{rgb}{0.97259,0.52980,0.05332}%
\definecolor{mycolor87}{rgb}{0.97568,0.54380,0.04362}%
\definecolor{mycolor88}{rgb}{0.98082,0.57221,0.02851}%
\definecolor{mycolor89}{rgb}{0.98189,0.57939,0.02625}%
\definecolor{mycolor90}{rgb}{0.98378,0.59385,0.02377}%
\definecolor{mycolor91}{rgb}{0.98532,0.60842,0.02420}%
\definecolor{mycolor92}{rgb}{0.98696,0.63048,0.03091}%
\definecolor{mycolor93}{rgb}{0.98793,0.66025,0.05175}%
\definecolor{mycolor94}{rgb}{0.98771,0.68281,0.07249}%
\definecolor{mycolor95}{rgb}{0.97750,0.78226,0.18592}%
\definecolor{mycolor96}{rgb}{0.96624,0.83619,0.26153}%
\definecolor{mycolor97}{rgb}{0.96252,0.85148,0.28555}%
\definecolor{mycolor98}{rgb}{0.98836,0.99836,0.64492}%
%
\addplot [color=mycolor1]
  table[row sep=crcr]{%
0	4395.3\\
2000	16898\\
4000	36599\\
6000	52416\\
8000	65179\\
10000	75854\\
12000	85094\\
14000	92849\\
16000	99694\\
18000	1.0511e+05\\
20000	1.0996e+05\\
22000	1.1375e+05\\
24000	1.1727e+05\\
28000	1.2297e+05\\
30000	1.256e+05\\
34000	1.2923e+05\\
40000	1.3375e+05\\
50000	1.3784e+05\\
58000	1.4027e+05\\
64000	1.414e+05\\
68000	1.415e+05\\
70000	1.4122e+05\\
74000	1.4186e+05\\
84000	1.4288e+05\\
86000	1.4258e+05\\
1e+05	1.4403e+05\\
};
\addlegendentry{\ClasskILegend}

\addplot [color=mycolor2, forget plot]
  table[row sep=crcr]{%
0	4395.3\\
2000	9768.2\\
4000	16716\\
6000	20697\\
8000	23076\\
10000	24488\\
12000	25402\\
14000	25804\\
18000	26105\\
20000	25976\\
24000	26056\\
30000	25612\\
32000	25686\\
34000	25630\\
36000	25430\\
40000	25500\\
42000	25284\\
48000	25191\\
52000	25529\\
54000	25390\\
56000	25420\\
58000	25378\\
60000	25254\\
62000	25252\\
64000	25550\\
66000	25409\\
72000	25269\\
74000	25139\\
76000	25263\\
78000	25202\\
80000	25366\\
84000	25004\\
86000	24941\\
88000	25225\\
90000	25235\\
92000	25087\\
94000	25076\\
96000	24865\\
98000	24929\\
1e+05	24908\\
};
\addplot [color=mycolor3, forget plot]
  table[row sep=crcr]{%
0	4395.3\\
2000	10818\\
4000	20017\\
6000	26123\\
8000	30120\\
10000	32957\\
12000	34678\\
14000	35900\\
16000	36743\\
18000	37301\\
22000	38084\\
28000	38366\\
30000	38236\\
34000	38179\\
38000	37758\\
42000	37849\\
44000	37969\\
48000	37541\\
52000	37661\\
56000	37530\\
60000	37846\\
64000	38105\\
68000	37794\\
70000	37913\\
74000	37816\\
76000	37791\\
78000	37979\\
82000	38147\\
84000	38348\\
88000	38226\\
94000	37672\\
96000	37757\\
98000	37974\\
1e+05	37921\\
};
\addplot [color=mycolor4, forget plot]
  table[row sep=crcr]{%
0	4395.3\\
2000	7737.2\\
4000	11788\\
6000	13776\\
8000	14728\\
10000	15093\\
12000	15209\\
14000	15180\\
16000	15213\\
18000	15141\\
20000	15200\\
22000	15045\\
24000	14794\\
26000	14611\\
28000	14344\\
30000	14269\\
36000	14433\\
38000	14383\\
40000	14210\\
42000	14140\\
44000	14197\\
46000	14179\\
48000	14214\\
50000	14147\\
52000	14248\\
54000	14246\\
56000	14127\\
58000	14183\\
60000	14102\\
62000	14138\\
64000	14028\\
66000	14054\\
68000	14023\\
70000	14138\\
72000	14175\\
76000	14163\\
78000	14053\\
82000	14272\\
84000	14220\\
86000	14019\\
88000	13912\\
90000	13926\\
92000	13895\\
94000	13981\\
98000	13876\\
1e+05	13930\\
};
\addplot [color=mycolor5, forget plot]
  table[row sep=crcr]{%
0	4395.3\\
2000	10232\\
4000	19788\\
6000	27326\\
8000	33330\\
10000	38151\\
12000	42184\\
14000	45633\\
16000	48576\\
18000	50931\\
20000	52902\\
22000	54640\\
24000	56013\\
28000	57980\\
32000	59422\\
36000	60412\\
40000	60852\\
46000	61227\\
48000	61508\\
54000	61763\\
60000	61928\\
68000	61903\\
70000	62194\\
92000	62168\\
96000	61677\\
1e+05	61667\\
};
\addplot [color=mycolor6, forget plot]
  table[row sep=crcr]{%
0	4395.3\\
2000	9443.8\\
4000	17443\\
6000	23553\\
8000	28473\\
10000	32544\\
12000	35756\\
14000	38422\\
16000	40635\\
18000	42493\\
20000	43932\\
24000	46562\\
26000	47451\\
30000	48859\\
34000	49432\\
42000	51009\\
44000	51283\\
46000	51226\\
50000	51855\\
58000	52190\\
62000	52379\\
66000	52385\\
70000	52505\\
74000	52631\\
76000	52596\\
80000	52898\\
1e+05	52491\\
};
\addplot [color=mycolor7, forget plot]
  table[row sep=crcr]{%
0	4395.3\\
2000	7979.4\\
4000	13069\\
6000	16524\\
8000	18756\\
10000	20483\\
12000	21803\\
14000	22637\\
18000	23898\\
20000	24369\\
24000	24894\\
26000	25056\\
28000	25132\\
36000	25872\\
38000	25860\\
40000	26088\\
42000	26408\\
44000	26504\\
46000	26510\\
52000	26926\\
58000	26837\\
60000	26650\\
64000	26434\\
68000	26783\\
70000	26870\\
72000	26738\\
76000	26808\\
78000	26877\\
82000	26691\\
84000	26769\\
86000	26998\\
88000	27030\\
90000	27192\\
98000	26809\\
1e+05	26753\\
};
\addplot [color=mycolor8, forget plot]
  table[row sep=crcr]{%
0	4395.3\\
2000	7799\\
4000	12612\\
6000	15801\\
8000	17816\\
10000	19120\\
12000	20020\\
14000	20707\\
16000	21070\\
18000	21292\\
22000	21475\\
24000	21456\\
26000	21584\\
28000	21567\\
30000	21487\\
34000	21101\\
36000	21200\\
40000	21139\\
44000	20850\\
46000	20714\\
48000	20739\\
50000	20594\\
52000	20572\\
54000	20751\\
56000	20642\\
58000	20628\\
60000	20501\\
62000	20508\\
64000	20374\\
66000	20381\\
68000	20176\\
74000	20215\\
76000	20348\\
78000	20378\\
80000	20323\\
82000	20149\\
86000	20236\\
88000	20339\\
96000	20234\\
98000	20269\\
1e+05	20207\\
};
\addplot [color=mycolor9, forget plot]
  table[row sep=crcr]{%
0	4395.3\\
2000	6538.2\\
4000	9223.9\\
6000	10773\\
8000	11717\\
10000	12360\\
14000	12815\\
18000	12988\\
20000	13046\\
22000	13162\\
26000	13164\\
28000	13259\\
32000	13257\\
34000	13221\\
36000	13260\\
38000	13388\\
44000	13469\\
46000	13556\\
48000	13532\\
50000	13432\\
54000	13383\\
56000	13456\\
58000	13596\\
60000	13451\\
62000	13420\\
64000	13533\\
66000	13512\\
68000	13604\\
70000	13567\\
72000	13446\\
74000	13392\\
76000	13534\\
78000	13632\\
80000	13651\\
82000	13536\\
84000	13487\\
86000	13544\\
90000	13417\\
92000	13466\\
94000	13438\\
96000	13522\\
98000	13656\\
1e+05	13603\\
};
\addplot [color=mycolor10, forget plot]
  table[row sep=crcr]{%
0	4395.3\\
2000	7037.3\\
4000	10778\\
6000	13220\\
8000	14925\\
10000	16050\\
12000	16945\\
14000	17681\\
16000	18163\\
20000	19035\\
22000	19349\\
28000	20007\\
32000	20190\\
34000	20397\\
36000	20454\\
38000	20409\\
44000	20728\\
46000	20770\\
48000	20712\\
52000	20772\\
54000	20806\\
58000	20752\\
62000	20883\\
64000	20820\\
70000	21133\\
76000	21005\\
78000	21048\\
82000	20912\\
86000	21213\\
92000	21228\\
96000	21240\\
1e+05	21079\\
};
\addplot [color=mycolor11, forget plot]
  table[row sep=crcr]{%
0	4395.3\\
2000	6301.9\\
4000	8594.4\\
6000	9645.8\\
8000	10160\\
10000	10294\\
14000	10163\\
16000	10120\\
20000	9977.4\\
22000	9781.2\\
24000	9679.1\\
28000	9316.7\\
30000	9196\\
34000	9130.9\\
36000	9151.9\\
40000	8993.6\\
42000	8961.7\\
44000	8895.4\\
46000	8924.2\\
48000	8854.4\\
50000	8738.5\\
52000	8689.3\\
54000	8720.3\\
56000	8685\\
58000	8591.7\\
62000	8522.5\\
64000	8558.8\\
66000	8622.3\\
70000	8593.3\\
72000	8704.5\\
74000	8703.1\\
76000	8599.1\\
78000	8557.4\\
82000	8599.4\\
84000	8573.3\\
86000	8607.8\\
88000	8668.9\\
90000	8650.3\\
92000	8684.1\\
94000	8622.5\\
98000	8609.1\\
1e+05	8514.8\\
};
\addplot [color=mycolor12, forget plot]
  table[row sep=crcr]{%
0	4395.3\\
2000	6449.4\\
4000	9238.7\\
6000	10892\\
8000	11793\\
10000	12399\\
12000	12919\\
14000	13231\\
16000	13425\\
18000	13546\\
22000	13665\\
24000	13625\\
28000	13624\\
34000	13971\\
36000	13855\\
40000	13764\\
42000	13912\\
44000	13969\\
46000	13972\\
48000	14022\\
52000	13867\\
54000	14035\\
60000	14161\\
62000	14319\\
74000	14148\\
82000	14225\\
84000	14289\\
86000	14285\\
88000	14375\\
90000	14367\\
92000	14318\\
94000	14223\\
98000	14210\\
1e+05	14332\\
};
\addplot [color=mycolor13, forget plot]
  table[row sep=crcr]{%
0	4395.3\\
2000	6115.6\\
4000	8168.2\\
6000	9294.1\\
8000	9881.2\\
10000	10228\\
12000	10342\\
14000	10322\\
18000	10474\\
22000	10361\\
24000	10360\\
28000	10473\\
30000	10530\\
32000	10516\\
34000	10645\\
36000	10698\\
38000	10630\\
40000	10689\\
42000	10714\\
46000	10613\\
52000	10800\\
54000	10914\\
56000	10839\\
58000	10668\\
60000	10598\\
62000	10702\\
64000	10694\\
66000	10820\\
68000	10861\\
70000	10844\\
72000	10921\\
74000	10728\\
76000	10719\\
80000	10800\\
82000	10721\\
86000	10840\\
88000	10844\\
90000	10932\\
92000	10897\\
94000	10820\\
96000	10862\\
98000	10814\\
1e+05	10726\\
};
\addplot [color=mycolor14, forget plot]
  table[row sep=crcr]{%
0	4395.3\\
2000	6916.3\\
4000	10773\\
6000	13438\\
8000	15356\\
10000	16755\\
12000	17833\\
16000	19362\\
18000	20016\\
20000	20460\\
24000	21068\\
26000	21294\\
30000	21854\\
36000	22071\\
38000	22248\\
44000	22484\\
48000	22475\\
52000	22675\\
54000	22902\\
56000	23012\\
58000	23014\\
60000	23115\\
62000	23106\\
64000	23162\\
68000	23140\\
70000	23364\\
74000	23564\\
76000	23350\\
78000	23369\\
80000	23290\\
82000	23317\\
84000	23247\\
88000	23301\\
90000	23232\\
94000	23402\\
96000	23212\\
98000	23243\\
1e+05	23481\\
};
\addplot [color=mycolor15, forget plot]
  table[row sep=crcr]{%
0	4395.3\\
2000	6346.7\\
4000	8973.1\\
6000	10851\\
8000	12225\\
10000	12936\\
14000	14046\\
18000	14613\\
20000	14770\\
22000	15030\\
24000	15232\\
26000	15327\\
28000	15328\\
30000	15442\\
32000	15457\\
36000	15770\\
40000	15849\\
42000	16028\\
46000	15808\\
48000	15959\\
50000	16167\\
54000	16283\\
56000	16251\\
58000	16366\\
60000	16356\\
62000	16406\\
66000	16401\\
68000	16315\\
72000	16341\\
74000	16291\\
76000	16358\\
78000	16374\\
80000	16518\\
82000	16527\\
86000	16333\\
88000	16414\\
92000	16389\\
1e+05	16692\\
};
\addplot [color=mycolor16, forget plot]
  table[row sep=crcr]{%
0	4395.3\\
2000	5785.4\\
4000	7500.7\\
6000	8191.2\\
8000	8412.8\\
10000	8495.1\\
12000	8444.6\\
14000	8342.9\\
16000	8183.6\\
18000	8004.3\\
20000	7858.2\\
22000	7788.2\\
26000	7548.6\\
30000	7433.2\\
32000	7322.7\\
36000	7237.8\\
38000	7089\\
40000	7013.8\\
48000	7003.1\\
50000	7063.7\\
52000	7103.4\\
56000	6989.1\\
58000	6948\\
60000	6953\\
62000	6989.6\\
64000	7003.3\\
66000	7050.3\\
68000	7018.9\\
70000	7035.3\\
72000	7002.1\\
74000	7022.3\\
76000	7007.9\\
78000	7035.1\\
80000	7082.2\\
82000	6975.9\\
88000	6920.5\\
94000	6984.3\\
96000	6965.6\\
98000	6911.8\\
1e+05	6945.9\\
};
\addplot [color=mycolor17, forget plot]
  table[row sep=crcr]{%
0	4395.3\\
2000	6312.7\\
4000	8918.5\\
6000	10647\\
8000	11869\\
10000	12752\\
12000	13325\\
14000	13666\\
16000	13817\\
18000	14149\\
22000	14329\\
24000	14398\\
26000	14596\\
30000	14679\\
32000	14794\\
34000	14778\\
36000	14880\\
40000	14922\\
42000	14987\\
44000	15000\\
46000	15098\\
48000	15130\\
50000	15098\\
52000	15011\\
54000	14973\\
56000	15113\\
58000	15152\\
60000	15100\\
62000	15245\\
70000	15390\\
80000	15335\\
82000	15428\\
84000	15367\\
88000	15542\\
90000	15558\\
92000	15692\\
94000	15660\\
96000	15676\\
98000	15632\\
1e+05	15542\\
};
\addplot [color=mycolor18, forget plot]
  table[row sep=crcr]{%
0	4395.3\\
2000	6693.7\\
4000	10292\\
6000	12987\\
8000	14935\\
10000	16465\\
12000	17822\\
14000	18735\\
16000	19515\\
18000	19934\\
20000	20607\\
22000	21073\\
28000	21748\\
30000	21816\\
36000	22318\\
38000	22620\\
44000	22919\\
46000	23108\\
48000	23172\\
56000	23010\\
58000	23145\\
62000	23141\\
64000	23225\\
66000	23168\\
68000	23351\\
70000	23406\\
72000	23337\\
76000	23465\\
84000	23707\\
86000	23442\\
88000	23482\\
90000	23428\\
92000	23553\\
94000	23527\\
96000	23635\\
98000	23512\\
1e+05	23631\\
};
\addplot [color=mycolor19, forget plot]
  table[row sep=crcr]{%
0	4395.3\\
2000	6008.7\\
4000	8265.3\\
6000	9774.9\\
8000	10889\\
10000	11793\\
12000	12335\\
14000	12713\\
16000	12859\\
18000	13231\\
20000	13476\\
22000	13445\\
24000	13560\\
28000	13950\\
34000	13798\\
36000	13949\\
38000	14143\\
40000	14238\\
44000	14271\\
46000	14388\\
48000	14288\\
50000	14126\\
52000	14100\\
56000	14214\\
58000	14203\\
62000	14373\\
64000	14189\\
66000	14153\\
68000	14059\\
70000	14029\\
72000	14081\\
74000	14230\\
82000	14333\\
86000	14314\\
88000	14327\\
90000	14417\\
96000	14383\\
98000	14308\\
1e+05	14450\\
};
\addplot [color=mycolor20, forget plot]
  table[row sep=crcr]{%
0	4395.3\\
2000	6300.5\\
4000	9090.6\\
6000	11156\\
8000	12687\\
10000	13835\\
12000	14545\\
14000	15071\\
16000	15707\\
18000	16249\\
20000	16600\\
22000	16777\\
24000	17039\\
28000	17336\\
32000	17521\\
34000	17751\\
36000	17922\\
38000	18032\\
44000	17942\\
46000	18070\\
48000	18363\\
50000	18494\\
52000	18343\\
54000	18323\\
56000	18373\\
58000	18277\\
62000	18415\\
64000	18379\\
66000	18398\\
68000	18520\\
70000	18720\\
74000	18749\\
76000	18593\\
78000	18564\\
80000	18454\\
82000	18476\\
84000	18573\\
86000	18571\\
90000	18468\\
92000	18520\\
94000	18425\\
96000	18497\\
98000	18669\\
1e+05	18710\\
};
\addplot [color=mycolor21, forget plot]
  table[row sep=crcr]{%
0	4395.3\\
2000	5212\\
4000	6368.1\\
6000	7089.1\\
8000	7577.7\\
10000	7927.5\\
14000	8379\\
16000	8504.7\\
18000	8551.9\\
20000	8738\\
26000	9022.7\\
30000	9136.6\\
32000	9201.6\\
34000	9199.6\\
38000	9133.7\\
42000	9280.1\\
44000	9340.9\\
46000	9355.7\\
48000	9230.6\\
50000	9241.7\\
52000	9290.7\\
54000	9382\\
56000	9347.5\\
58000	9397.8\\
60000	9409.8\\
62000	9300.6\\
64000	9271.3\\
66000	9308.8\\
68000	9274.5\\
72000	9325.5\\
74000	9440.9\\
76000	9530\\
78000	9464.3\\
80000	9433.2\\
84000	9438\\
86000	9457.2\\
88000	9447.5\\
90000	9400.2\\
92000	9437.5\\
1e+05	9454.4\\
};
\addplot [color=mycolor22, forget plot]
  table[row sep=crcr]{%
0	4395.3\\
2000	7191.2\\
4000	11672\\
6000	15120\\
8000	17792\\
10000	19877\\
12000	21422\\
14000	22431\\
16000	23421\\
18000	24196\\
20000	24733\\
24000	25542\\
26000	25895\\
28000	26126\\
30000	26171\\
34000	26076\\
38000	25981\\
48000	26034\\
54000	25628\\
56000	25502\\
58000	25622\\
72000	25371\\
76000	25181\\
80000	25159\\
82000	25266\\
88000	25161\\
90000	25005\\
92000	24945\\
94000	25002\\
1e+05	24905\\
};
\addplot [color=mycolor23, forget plot]
  table[row sep=crcr]{%
0	4395.3\\
2000	5759.6\\
4000	7748.6\\
6000	9004.2\\
8000	9877.4\\
10000	10483\\
12000	10923\\
14000	11275\\
16000	11494\\
24000	11980\\
30000	12051\\
32000	12076\\
34000	12161\\
38000	12159\\
40000	12232\\
46000	12214\\
48000	12307\\
52000	12212\\
54000	12177\\
56000	12188\\
58000	12129\\
60000	12134\\
62000	12102\\
64000	12166\\
68000	12179\\
70000	12122\\
72000	12106\\
76000	12157\\
82000	12137\\
84000	12162\\
92000	12106\\
1e+05	12146\\
};
\addplot [color=mycolor24, forget plot]
  table[row sep=crcr]{%
0	4395.3\\
2000	5115.1\\
4000	5909.6\\
6000	6220.2\\
8000	6305\\
10000	6207.6\\
12000	6081.5\\
14000	5915.3\\
16000	5811.5\\
20000	5575.2\\
22000	5519.5\\
24000	5436.6\\
28000	5226.3\\
30000	5175.2\\
32000	5108.1\\
34000	5012.4\\
36000	4974.1\\
38000	4880.3\\
40000	4882.7\\
42000	4914.2\\
44000	4980.5\\
46000	4969.8\\
48000	4900.3\\
50000	4855.9\\
52000	4856.5\\
54000	4839\\
56000	4838.7\\
58000	4869.2\\
62000	4826.8\\
64000	4859.5\\
66000	4816.7\\
68000	4839.4\\
70000	4880.8\\
72000	4867.8\\
74000	4829.8\\
76000	4807.6\\
82000	4872\\
84000	4850.3\\
86000	4903.5\\
88000	4936.8\\
90000	4910.2\\
92000	4807.8\\
94000	4831.4\\
96000	4832.6\\
98000	4774\\
1e+05	4795.8\\
};
\addplot [color=mycolor25, forget plot]
  table[row sep=crcr]{%
0	4395.3\\
2000	5590.3\\
4000	7443.7\\
6000	8669.9\\
8000	9443\\
10000	9915\\
14000	10611\\
16000	10870\\
18000	11048\\
20000	11264\\
22000	11401\\
24000	11411\\
26000	11498\\
28000	11434\\
30000	11566\\
32000	11738\\
34000	11836\\
36000	11829\\
38000	11890\\
42000	11728\\
44000	11693\\
46000	11698\\
48000	11661\\
54000	11965\\
56000	11922\\
60000	11932\\
64000	11819\\
66000	11962\\
70000	11834\\
76000	11997\\
78000	11919\\
80000	11924\\
82000	12058\\
84000	12143\\
86000	12060\\
92000	11968\\
94000	12050\\
1e+05	12178\\
};
\addplot [color=mycolor26, forget plot]
  table[row sep=crcr]{%
0	4395.3\\
4000	6327.3\\
6000	6844.7\\
8000	7152.7\\
10000	7326.2\\
12000	7429.2\\
14000	7565.5\\
18000	7656.5\\
20000	7640.4\\
22000	7668.8\\
24000	7675.5\\
26000	7656.6\\
30000	7763.5\\
32000	7833.3\\
34000	7944.3\\
38000	7990.2\\
40000	7972.5\\
42000	8052.9\\
44000	7897.5\\
46000	7834\\
48000	7874.8\\
54000	8047.2\\
56000	8045.4\\
58000	8068.8\\
60000	7984.2\\
62000	8078.6\\
64000	8228.2\\
68000	8169.3\\
70000	8056.6\\
72000	8105.3\\
74000	8070.2\\
76000	8085.1\\
78000	8026.1\\
80000	8029.5\\
82000	8133.4\\
84000	8118.4\\
86000	8059.9\\
90000	8067.5\\
92000	8001.1\\
96000	8083.1\\
98000	8035.7\\
1e+05	8122\\
};
\addplot [color=mycolor27, forget plot]
  table[row sep=crcr]{%
0	4395.3\\
2000	5678.5\\
4000	7474.6\\
6000	8620.7\\
8000	9275.3\\
10000	9655.4\\
12000	9899.5\\
14000	10076\\
18000	10164\\
20000	10175\\
22000	10068\\
24000	10020\\
26000	9930.2\\
28000	10007\\
30000	10028\\
32000	9959.8\\
34000	9949.7\\
36000	9904.1\\
38000	9814.6\\
40000	9647.8\\
42000	9675.5\\
44000	9779.2\\
46000	9661.3\\
48000	9694.3\\
50000	9782.8\\
54000	9742.6\\
56000	9852.7\\
58000	9843.6\\
60000	9748.4\\
66000	9659.6\\
74000	9720.5\\
76000	9771.8\\
78000	9870.3\\
80000	9807.9\\
82000	9710.7\\
84000	9702.1\\
86000	9639.3\\
88000	9609.4\\
94000	9645.6\\
96000	9612.3\\
98000	9607.2\\
1e+05	9552.8\\
};
\addplot [color=mycolor28, forget plot]
  table[row sep=crcr]{%
0	4395.3\\
2000	5057.2\\
4000	5916.3\\
6000	6387.4\\
8000	6537\\
10000	6597.8\\
16000	6469.8\\
18000	6382.5\\
20000	6233.7\\
22000	6130.3\\
24000	6117.5\\
28000	6038.4\\
30000	5944.4\\
32000	5870\\
34000	5828.1\\
36000	5813.8\\
42000	5672.9\\
44000	5640.3\\
48000	5608.3\\
50000	5559.6\\
52000	5592.6\\
54000	5585.6\\
56000	5558.7\\
60000	5570.1\\
64000	5560.8\\
66000	5563.5\\
68000	5491.1\\
70000	5515.3\\
74000	5618.8\\
76000	5617.7\\
78000	5527.7\\
80000	5484\\
86000	5545.3\\
88000	5538.7\\
92000	5551.9\\
94000	5602.1\\
96000	5620.5\\
98000	5580.5\\
1e+05	5632.7\\
};
\addplot [color=mycolor29, forget plot]
  table[row sep=crcr]{%
0	4395.3\\
2000	5804.2\\
4000	7888.6\\
6000	9317.7\\
8000	10256\\
10000	10913\\
12000	11323\\
14000	11580\\
20000	11846\\
22000	11823\\
24000	11833\\
26000	11908\\
30000	11735\\
34000	11563\\
42000	11660\\
46000	11652\\
48000	11664\\
50000	11518\\
54000	11479\\
56000	11392\\
70000	11416\\
72000	11458\\
74000	11416\\
76000	11422\\
78000	11322\\
80000	11384\\
82000	11375\\
86000	11303\\
90000	11354\\
92000	11409\\
96000	11457\\
98000	11591\\
1e+05	11573\\
};
\addplot [color=mycolor30, forget plot]
  table[row sep=crcr]{%
0	4395.3\\
2000	4900.4\\
4000	5677.5\\
6000	6169.7\\
8000	6507.3\\
10000	6767.8\\
12000	6970.7\\
14000	7082\\
18000	7241.8\\
20000	7274.2\\
22000	7280.5\\
24000	7403.4\\
26000	7482.6\\
28000	7492.4\\
32000	7617.1\\
34000	7629.1\\
36000	7578.2\\
42000	7645.2\\
44000	7618\\
46000	7640.6\\
48000	7625.3\\
50000	7558.1\\
54000	7586.5\\
56000	7678.6\\
60000	7768.2\\
62000	7734.1\\
68000	7745.6\\
70000	7671\\
74000	7717.7\\
76000	7729.5\\
78000	7819.6\\
80000	7829.9\\
82000	7771.9\\
84000	7781.4\\
90000	7728.6\\
92000	7778.6\\
94000	7905\\
96000	7924.3\\
1e+05	7891.2\\
};
\addplot [color=mycolor31, forget plot]
  table[row sep=crcr]{%
0	4395.3\\
2000	5184.7\\
4000	6235.6\\
6000	6775.2\\
8000	7068.3\\
10000	7184.6\\
12000	7347\\
14000	7227.6\\
18000	7117.3\\
20000	6994.9\\
22000	6986.1\\
30000	6763.4\\
32000	6819.7\\
34000	6842.2\\
38000	6811.2\\
44000	6666.1\\
46000	6659.3\\
50000	6750.2\\
56000	6513.9\\
58000	6501.8\\
62000	6500.9\\
64000	6479\\
66000	6529.9\\
68000	6537.8\\
70000	6512.3\\
72000	6509.7\\
76000	6545.7\\
78000	6493.6\\
80000	6518.8\\
86000	6480\\
88000	6416.5\\
90000	6408.2\\
92000	6422.1\\
94000	6467.2\\
96000	6451.2\\
1e+05	6490.5\\
};
\addplot [color=mycolor32, forget plot]
  table[row sep=crcr]{%
0	4395.3\\
2000	5764.2\\
4000	7806.6\\
6000	9058.2\\
8000	9804.6\\
10000	10269\\
12000	10636\\
16000	11020\\
18000	11114\\
20000	11095\\
24000	11153\\
26000	11075\\
32000	11170\\
34000	11121\\
38000	11266\\
40000	11102\\
46000	10943\\
48000	10813\\
50000	10770\\
54000	10937\\
56000	10937\\
58000	10830\\
66000	10774\\
70000	10671\\
72000	10628\\
74000	10634\\
76000	10540\\
80000	10549\\
82000	10458\\
88000	10611\\
90000	10497\\
94000	10528\\
96000	10607\\
1e+05	10578\\
};
\addplot [color=mycolor33, forget plot]
  table[row sep=crcr]{%
0	4395.3\\
2000	4854.2\\
4000	5540.7\\
6000	6084.3\\
8000	6483.8\\
10000	6754.4\\
12000	6983.8\\
14000	7091\\
16000	7247.4\\
18000	7275.4\\
24000	7676.9\\
26000	7647.2\\
28000	7664.2\\
30000	7734.6\\
32000	7766.6\\
36000	7919.5\\
38000	7918.2\\
44000	8052.4\\
52000	7985.1\\
54000	7973.1\\
56000	8090.1\\
60000	8065.4\\
62000	8117.8\\
66000	8116.1\\
68000	8047.4\\
70000	8014.6\\
74000	8133.4\\
78000	8026\\
80000	8033\\
82000	7986.1\\
84000	8022.5\\
88000	8028.8\\
90000	8024.8\\
92000	8172.6\\
94000	8151.7\\
96000	8212.5\\
98000	8196.5\\
1e+05	8214.9\\
};
\addplot [color=mycolor34, forget plot]
  table[row sep=crcr]{%
0	4395.3\\
2000	4912.2\\
4000	5450.3\\
6000	5532.3\\
8000	5486.6\\
10000	5364.1\\
12000	5148.7\\
14000	4956.4\\
16000	4837\\
18000	4667\\
20000	4532.8\\
22000	4427.5\\
24000	4378.3\\
26000	4342.9\\
30000	4238.5\\
36000	4154.9\\
38000	4108.3\\
40000	4100.7\\
42000	4066.9\\
46000	4041.6\\
48000	4037\\
50000	4012.2\\
52000	4037.7\\
54000	4088.1\\
58000	4021.8\\
62000	4060\\
64000	4076.9\\
68000	4063.5\\
70000	4095.3\\
72000	4108.9\\
78000	4059\\
82000	4059.9\\
84000	4041.3\\
88000	4035.7\\
92000	3991.8\\
96000	4001\\
98000	4055.4\\
1e+05	4061.4\\
};
\addplot [color=mycolor35]
  table[row sep=crcr]{%
0	4395.3\\
2000	5398.9\\
4000	6864.7\\
6000	7827.8\\
8000	8441\\
10000	8885.8\\
12000	9192.3\\
14000	9456\\
16000	9577.2\\
20000	9638.1\\
22000	9701.5\\
24000	9704.1\\
28000	9794.5\\
36000	9871.2\\
42000	9800.5\\
44000	9729.6\\
48000	9747.4\\
52000	9872.7\\
54000	9884.1\\
58000	9985\\
60000	9992.1\\
64000	9897.2\\
68000	9873.9\\
70000	9860.9\\
74000	9913.3\\
76000	9984.4\\
78000	9941.9\\
84000	9936.6\\
88000	10002\\
90000	10060\\
92000	10067\\
94000	10038\\
96000	10106\\
98000	10119\\
1e+05	10040\\
};
\addlegendentry{\ClasskIILegend}

\addplot [color=mycolor36, forget plot]
  table[row sep=crcr]{%
0	4395.3\\
2000	4799.3\\
4000	5255.9\\
6000	5485.8\\
8000	5513.8\\
10000	5465.1\\
12000	5392.9\\
22000	4879\\
24000	4869.8\\
26000	4829.7\\
28000	4726.8\\
30000	4720.1\\
32000	4666.9\\
34000	4648\\
38000	4579.2\\
40000	4555.4\\
44000	4482\\
46000	4483.4\\
48000	4442.3\\
54000	4401.6\\
56000	4418.2\\
58000	4456.6\\
60000	4450.9\\
62000	4469\\
64000	4396.8\\
66000	4409.7\\
68000	4389.9\\
70000	4400.4\\
72000	4364.9\\
74000	4384.4\\
76000	4369.7\\
78000	4391.7\\
80000	4436.5\\
82000	4380.6\\
84000	4375\\
86000	4392.3\\
88000	4444.6\\
90000	4457.8\\
92000	4404.6\\
96000	4395.5\\
98000	4368.7\\
1e+05	4360.1\\
};
\addplot [color=mycolor37, forget plot]
  table[row sep=crcr]{%
0	4395.3\\
2000	4427.5\\
4000	4427.8\\
6000	4400.8\\
8000	4355.1\\
10000	4259.6\\
12000	4203.3\\
16000	4114.3\\
18000	4054.1\\
20000	3975\\
22000	3960.1\\
24000	3960.4\\
28000	3842.9\\
30000	3830.5\\
36000	3763.6\\
38000	3786.3\\
42000	3720.3\\
44000	3698.6\\
46000	3717.7\\
48000	3719.9\\
50000	3696.4\\
52000	3708.1\\
54000	3706.5\\
56000	3672.5\\
60000	3665.3\\
62000	3687.2\\
68000	3680.9\\
70000	3693.2\\
72000	3721.6\\
74000	3716.5\\
76000	3691.6\\
82000	3691.8\\
84000	3668.2\\
86000	3693\\
90000	3689\\
94000	3665.2\\
98000	3677.4\\
1e+05	3715.9\\
};
\addplot [color=mycolor38, forget plot]
  table[row sep=crcr]{%
0	4395.3\\
2000	4200.1\\
4000	3875.4\\
6000	3720.6\\
8000	3604.6\\
10000	3565.7\\
12000	3516.7\\
14000	3445.1\\
16000	3408\\
24000	3327.2\\
26000	3314.5\\
28000	3269.2\\
30000	3236.4\\
32000	3183.7\\
34000	3198.5\\
36000	3187.1\\
38000	3210.6\\
40000	3189.9\\
42000	3208\\
44000	3193.7\\
46000	3158.6\\
52000	3166.8\\
56000	3123.6\\
58000	3159.9\\
60000	3157.3\\
62000	3177.7\\
66000	3134.6\\
68000	3143\\
70000	3162\\
72000	3119.1\\
74000	3121.6\\
76000	3094.3\\
78000	3083.6\\
80000	3102.3\\
82000	3098.9\\
84000	3114.1\\
86000	3098.4\\
88000	3100.1\\
90000	3124.7\\
92000	3090.1\\
94000	3123.6\\
96000	3120.2\\
98000	3108\\
1e+05	3129.7\\
};
\addplot [color=mycolor39, forget plot]
  table[row sep=crcr]{%
0	4395.3\\
2000	5699\\
4000	7701\\
6000	9217.8\\
8000	10285\\
10000	11105\\
12000	11561\\
14000	11938\\
16000	12149\\
18000	12295\\
20000	12301\\
22000	12264\\
24000	12313\\
26000	12485\\
30000	12658\\
32000	12712\\
34000	12699\\
36000	12584\\
38000	12538\\
40000	12425\\
46000	12299\\
48000	12304\\
52000	12437\\
58000	12219\\
62000	12154\\
66000	12075\\
70000	12117\\
74000	11962\\
76000	11819\\
78000	11847\\
82000	12026\\
84000	12145\\
86000	12093\\
88000	11937\\
90000	11916\\
92000	11998\\
94000	11966\\
98000	11759\\
1e+05	11775\\
};
\addplot [color=mycolor40, forget plot]
  table[row sep=crcr]{%
0	4395.3\\
2000	4472.7\\
4000	4646.6\\
6000	4790\\
10000	4991.3\\
12000	5094\\
16000	5216.2\\
18000	5274.1\\
26000	5367.6\\
28000	5451.2\\
34000	5448.1\\
38000	5544.7\\
40000	5538\\
50000	5603.2\\
54000	5606.3\\
56000	5631.4\\
58000	5634.7\\
60000	5620.7\\
62000	5651.3\\
66000	5655.4\\
70000	5688.7\\
72000	5733\\
74000	5748.1\\
76000	5735.2\\
82000	5768.4\\
88000	5734.5\\
90000	5693.6\\
92000	5696.1\\
94000	5737.6\\
96000	5744.1\\
1e+05	5712.5\\
};
\addplot [color=mycolor41, forget plot]
  table[row sep=crcr]{%
0	4395.3\\
2000	4343\\
4000	4351.9\\
8000	4514\\
10000	4614.9\\
12000	4688.6\\
14000	4736.8\\
16000	4829.2\\
18000	4898.8\\
22000	4998.7\\
24000	5027.1\\
26000	5083.5\\
30000	5093.5\\
32000	5126.7\\
36000	5123.9\\
38000	5147.9\\
40000	5156.9\\
42000	5201.8\\
46000	5234.5\\
48000	5211.1\\
50000	5237.3\\
52000	5187.8\\
54000	5223.3\\
56000	5226.1\\
58000	5202.9\\
62000	5263.9\\
66000	5335.1\\
68000	5317.9\\
70000	5320.6\\
72000	5290.9\\
74000	5282.6\\
78000	5335.8\\
82000	5278.2\\
86000	5437.1\\
88000	5447.9\\
96000	5381.5\\
98000	5396.4\\
1e+05	5427.3\\
};
\addplot [color=mycolor42, forget plot]
  table[row sep=crcr]{%
0	4395.3\\
2000	5814.3\\
4000	8022.8\\
6000	9754.5\\
8000	11159\\
10000	12371\\
12000	13378\\
14000	14243\\
16000	14937\\
18000	15398\\
20000	15696\\
24000	16528\\
26000	16934\\
32000	17927\\
34000	18012\\
40000	18617\\
46000	18959\\
48000	19124\\
54000	19401\\
56000	19464\\
60000	19415\\
62000	19257\\
68000	19638\\
70000	19742\\
74000	19715\\
76000	19786\\
78000	19725\\
84000	19872\\
86000	19766\\
90000	19921\\
92000	19831\\
96000	19879\\
98000	19813\\
1e+05	19922\\
};
\addplot [color=mycolor43, forget plot]
  table[row sep=crcr]{%
0	4395.3\\
2000	4680.5\\
4000	5046\\
6000	5271.1\\
8000	5449.7\\
12000	5589.4\\
14000	5594\\
16000	5540.6\\
18000	5519.7\\
22000	5514.7\\
28000	5375.9\\
30000	5385.7\\
34000	5318.4\\
36000	5278.1\\
40000	5251.2\\
42000	5255.7\\
44000	5236\\
50000	5139\\
52000	5137.3\\
56000	5115.3\\
60000	5097.6\\
62000	5083.3\\
68000	5108.7\\
74000	5056.1\\
76000	5040.8\\
78000	5067.2\\
80000	5062.3\\
82000	5075.9\\
84000	5034.9\\
86000	5039.4\\
90000	5085.1\\
1e+05	5039.1\\
};
\addplot [color=mycolor44, forget plot]
  table[row sep=crcr]{%
0	4395.3\\
4000	4585.3\\
6000	4576.9\\
10000	4425.1\\
14000	4274.7\\
16000	4164.1\\
18000	4137.9\\
22000	4003.5\\
24000	3990.5\\
26000	3952.5\\
28000	3926.8\\
30000	3846.4\\
32000	3843\\
36000	3777.2\\
40000	3698.3\\
42000	3704.4\\
44000	3728.4\\
46000	3702.1\\
52000	3587.6\\
56000	3535.8\\
60000	3563.8\\
66000	3497.2\\
70000	3509.3\\
72000	3542.9\\
74000	3503.3\\
76000	3488.9\\
78000	3490.6\\
80000	3470.9\\
90000	3478.5\\
92000	3499.9\\
94000	3498.6\\
96000	3512\\
98000	3483.6\\
1e+05	3472.6\\
};
\addplot [color=mycolor45, forget plot]
  table[row sep=crcr]{%
0	4395.3\\
2000	4551.6\\
4000	4848.1\\
8000	5348.6\\
10000	5562.1\\
12000	5715.4\\
18000	6069.3\\
22000	6236.7\\
26000	6375.4\\
30000	6440.2\\
34000	6541.2\\
40000	6622.7\\
42000	6634\\
44000	6690.3\\
46000	6706.2\\
48000	6762.1\\
52000	6743\\
54000	6745.5\\
58000	6834.9\\
60000	6852.9\\
64000	6844.9\\
66000	6878\\
70000	6794.5\\
72000	6813.6\\
74000	6773.7\\
78000	6854.3\\
82000	6872.7\\
86000	6880.5\\
90000	6809.8\\
92000	6824.8\\
94000	6866.8\\
96000	6880.8\\
98000	6859.9\\
1e+05	6806.3\\
};
\addplot [color=mycolor46, forget plot]
  table[row sep=crcr]{%
0	4395.3\\
2000	4266\\
4000	4116.8\\
6000	4060\\
10000	4001.3\\
12000	4020.2\\
14000	3999.5\\
16000	4002.1\\
22000	4044.6\\
24000	4076.7\\
26000	4037.9\\
28000	4051.7\\
30000	4081.2\\
32000	4035\\
36000	4038.9\\
40000	4054.9\\
42000	4050.2\\
44000	4075.7\\
46000	4084.1\\
48000	4058.5\\
54000	4069.3\\
56000	4066.8\\
64000	4109.6\\
66000	4077.4\\
70000	4123.5\\
72000	4095.2\\
76000	4108.6\\
80000	4098.7\\
82000	4131.1\\
84000	4122.5\\
88000	4160.3\\
90000	4161.5\\
92000	4137.8\\
94000	4098.8\\
96000	4120.5\\
1e+05	4091.3\\
};
\addplot [color=mycolor47, forget plot]
  table[row sep=crcr]{%
0	4395.3\\
2000	4559.1\\
4000	4893.2\\
6000	5203\\
8000	5400.9\\
10000	5518.4\\
12000	5549.8\\
14000	5606.8\\
18000	5603.5\\
20000	5559.5\\
22000	5557.3\\
28000	5479\\
34000	5358.1\\
36000	5365.7\\
38000	5345.9\\
42000	5351\\
44000	5332.3\\
50000	5364.7\\
52000	5346.7\\
54000	5305.6\\
56000	5319.6\\
58000	5267.9\\
60000	5250.7\\
64000	5250.5\\
66000	5230\\
68000	5296.4\\
70000	5315.1\\
72000	5308.3\\
74000	5356.9\\
76000	5355.6\\
78000	5321.2\\
80000	5272.4\\
82000	5275.9\\
84000	5250.2\\
86000	5274\\
88000	5243.5\\
94000	5253.6\\
96000	5263.5\\
1e+05	5202.8\\
};
\addplot [color=mycolor48, forget plot]
  table[row sep=crcr]{%
0	4395.3\\
2000	4218.4\\
4000	3966.2\\
6000	3792.7\\
8000	3690.1\\
12000	3523.6\\
14000	3449.8\\
20000	3296.6\\
26000	3215.9\\
28000	3196\\
30000	3156.1\\
40000	3071.1\\
42000	3036.5\\
46000	3061\\
48000	3053.1\\
50000	3055.5\\
52000	3036.3\\
54000	3034.9\\
56000	3044.8\\
60000	3014.4\\
62000	3000.8\\
74000	3018\\
76000	3008.9\\
78000	3017.7\\
80000	3002.1\\
82000	2999.8\\
86000	3025.2\\
90000	3022.7\\
92000	3010.4\\
96000	3007.1\\
1e+05	2985.7\\
};
\addplot [color=mycolor49, forget plot]
  table[row sep=crcr]{%
0	4395.3\\
2000	4562.7\\
4000	4842.7\\
6000	4961\\
10000	4981.4\\
14000	4895\\
16000	4902.6\\
18000	4855\\
20000	4776.3\\
24000	4686.8\\
26000	4619.9\\
28000	4639.1\\
30000	4590.3\\
34000	4643.5\\
36000	4635.4\\
38000	4581.2\\
40000	4505.2\\
42000	4481.2\\
48000	4478.6\\
52000	4429.4\\
58000	4495.2\\
62000	4486.6\\
64000	4456.2\\
66000	4449.7\\
68000	4413.2\\
70000	4414.2\\
72000	4437.9\\
74000	4399.5\\
76000	4414.3\\
78000	4386.3\\
80000	4416\\
84000	4393.7\\
86000	4447.5\\
88000	4408.8\\
92000	4365.3\\
96000	4366.7\\
1e+05	4424.5\\
};
\addplot [color=mycolor50, forget plot]
  table[row sep=crcr]{%
0	4395.3\\
2000	3905.5\\
4000	3305.5\\
6000	2944.3\\
8000	2763.6\\
10000	2648.6\\
12000	2561.1\\
14000	2512.2\\
16000	2477.4\\
18000	2461.7\\
20000	2457.2\\
22000	2471.7\\
26000	2463.8\\
32000	2481.6\\
34000	2480.6\\
36000	2489.3\\
38000	2511.7\\
40000	2497.7\\
42000	2493.9\\
52000	2514.9\\
56000	2524.8\\
58000	2511.9\\
60000	2483.7\\
62000	2507.9\\
64000	2499.1\\
66000	2532.5\\
70000	2501\\
72000	2523.7\\
74000	2533.7\\
80000	2513\\
82000	2496.9\\
84000	2515.3\\
86000	2522.1\\
88000	2514.6\\
90000	2528.8\\
92000	2535.5\\
94000	2518.7\\
98000	2543.5\\
1e+05	2537.1\\
};
\addplot [color=mycolor51, forget plot]
  table[row sep=crcr]{%
0	4395.3\\
2000	4638.2\\
4000	5070.1\\
6000	5391.5\\
8000	5605.6\\
10000	5764.5\\
12000	5872.7\\
14000	5954.7\\
16000	5914.4\\
18000	5914.3\\
20000	5979.4\\
22000	5983.3\\
24000	6026\\
28000	5993.3\\
30000	5944.5\\
34000	5899.8\\
36000	5848.5\\
38000	5940.8\\
40000	5979.3\\
42000	5988.8\\
44000	5980.6\\
46000	5927.7\\
50000	5885.8\\
54000	5838.7\\
56000	5807.1\\
58000	5807.8\\
62000	5951\\
64000	5938.3\\
68000	5985.9\\
70000	5952\\
72000	5875.1\\
74000	5953.7\\
76000	6012.6\\
80000	6024.1\\
82000	6007.3\\
84000	6068.4\\
86000	6090\\
88000	6024.1\\
90000	5980.6\\
92000	5955.1\\
94000	5987.7\\
98000	5989.3\\
1e+05	6029.8\\
};
\addplot [color=mycolor52, forget plot]
  table[row sep=crcr]{%
0	4395.3\\
2000	4426.5\\
4000	4437.5\\
6000	4468.6\\
8000	4423.1\\
12000	4283.3\\
14000	4254.9\\
16000	4190.1\\
20000	4135.7\\
22000	4100.8\\
24000	4030.2\\
28000	3952.4\\
30000	3969.9\\
34000	3963.3\\
40000	3863.2\\
42000	3816.1\\
44000	3808.4\\
48000	3748\\
50000	3775.3\\
56000	3742.6\\
58000	3731.8\\
60000	3684.4\\
62000	3668.7\\
68000	3701.2\\
72000	3683.5\\
74000	3711.2\\
78000	3718.1\\
80000	3671.7\\
84000	3635.5\\
86000	3657.1\\
88000	3649.6\\
90000	3678\\
92000	3665\\
94000	3637.5\\
96000	3572.9\\
98000	3542.9\\
1e+05	3540.2\\
};
\addplot [color=mycolor53, forget plot]
  table[row sep=crcr]{%
0	4395.3\\
2000	4807.6\\
4000	5442.7\\
6000	5843.1\\
8000	6096.6\\
10000	6273.9\\
12000	6340.8\\
14000	6380.2\\
16000	6354\\
18000	6375.2\\
20000	6343.9\\
24000	6394.1\\
26000	6375.5\\
30000	6384.3\\
32000	6340.2\\
34000	6258.2\\
38000	6205.4\\
40000	6196.3\\
42000	6143.7\\
46000	6119.4\\
50000	6063\\
52000	6097.3\\
54000	6081.1\\
56000	6090.3\\
58000	6117.3\\
60000	6116.9\\
62000	6047.8\\
64000	6013.8\\
66000	5941.5\\
68000	5949.1\\
70000	5939.3\\
72000	5901.8\\
74000	5889.8\\
76000	5924.6\\
80000	5899.5\\
82000	5863.7\\
84000	5849.7\\
86000	5854.6\\
88000	5836.5\\
90000	5852.8\\
92000	5794.5\\
94000	5782.4\\
96000	5815.8\\
98000	5780.4\\
1e+05	5787.7\\
};
\addplot [color=mycolor54, forget plot]
  table[row sep=crcr]{%
0	4395.3\\
2000	4436.4\\
8000	4778.1\\
10000	4815.7\\
12000	4799.1\\
14000	4817.8\\
16000	4857.1\\
18000	4831.8\\
22000	4869\\
24000	4850.8\\
26000	4819.1\\
28000	4806\\
30000	4758.1\\
34000	4750.2\\
36000	4731\\
42000	4763\\
44000	4729\\
52000	4648.2\\
54000	4642.4\\
56000	4655.7\\
58000	4617.7\\
68000	4711.9\\
70000	4704.3\\
74000	4593.2\\
76000	4583.1\\
78000	4606.3\\
82000	4582.2\\
84000	4604.3\\
92000	4568.7\\
96000	4599.4\\
1e+05	4579\\
};
\addplot [color=mycolor55, forget plot]
  table[row sep=crcr]{%
0	4395.3\\
2000	4457.5\\
8000	4900.3\\
10000	5039.9\\
12000	5140.1\\
18000	5365.6\\
20000	5382.6\\
22000	5453.3\\
26000	5473.5\\
28000	5506.7\\
30000	5513.8\\
34000	5568.5\\
38000	5565.5\\
42000	5622.7\\
44000	5669.5\\
50000	5695.5\\
54000	5652.1\\
56000	5657.9\\
58000	5618.7\\
60000	5638.8\\
62000	5675\\
64000	5682\\
66000	5706.6\\
68000	5688.4\\
70000	5633.6\\
74000	5635.1\\
76000	5678.9\\
84000	5722.8\\
88000	5700.6\\
90000	5732.4\\
98000	5759.9\\
1e+05	5738.7\\
};
\addplot [color=mycolor56, forget plot]
  table[row sep=crcr]{%
0	4395.3\\
2000	4413.1\\
4000	4462.6\\
6000	4488\\
10000	4502.8\\
12000	4489.6\\
16000	4497.7\\
18000	4429\\
20000	4409.2\\
22000	4374.6\\
24000	4353.3\\
28000	4267.6\\
34000	4230\\
36000	4207.7\\
38000	4212.3\\
40000	4204.5\\
42000	4181.4\\
44000	4145.3\\
46000	4152.4\\
60000	4091.5\\
64000	4105.1\\
66000	4090.3\\
68000	4097\\
70000	4089.4\\
72000	4106.1\\
74000	4088\\
76000	4052.8\\
80000	4087.3\\
82000	4083.8\\
86000	4036\\
88000	4032.7\\
92000	4057.5\\
96000	4052.9\\
98000	4090.6\\
1e+05	4091.9\\
};
\addplot [color=mycolor57, forget plot]
  table[row sep=crcr]{%
0	4395.3\\
2000	4601.9\\
4000	4972.7\\
6000	5305\\
8000	5609.8\\
10000	5845.5\\
12000	6041.9\\
14000	6214.1\\
16000	6365.6\\
20000	6557.1\\
24000	6696.3\\
28000	6778.5\\
46000	6824.3\\
48000	6816.6\\
52000	6852.1\\
54000	6833.7\\
60000	6856.2\\
62000	6835.6\\
68000	6859.9\\
70000	6812.6\\
74000	6784.6\\
80000	6688.2\\
92000	6706.3\\
94000	6727.1\\
96000	6716.9\\
1e+05	6736.4\\
};
\addplot [color=mycolor58, forget plot]
  table[row sep=crcr]{%
0	4395.3\\
6000	4293\\
8000	4272.9\\
10000	4235.8\\
12000	4182.2\\
14000	4150.9\\
16000	4103.8\\
18000	4083.5\\
20000	4075.1\\
22000	4026.6\\
24000	3955.5\\
28000	3849.8\\
36000	3797.7\\
38000	3776.9\\
40000	3737.8\\
42000	3739.3\\
46000	3775.3\\
48000	3763.6\\
50000	3716.6\\
52000	3706.1\\
54000	3737.9\\
56000	3715.7\\
62000	3714.9\\
70000	3732.6\\
72000	3672.3\\
74000	3692.6\\
76000	3694.8\\
78000	3686.1\\
80000	3649.2\\
82000	3637.8\\
84000	3654\\
88000	3635.8\\
90000	3635.9\\
92000	3626\\
94000	3650.1\\
96000	3663\\
98000	3664.4\\
1e+05	3680.5\\
};
\addplot [color=mycolor59, forget plot]
  table[row sep=crcr]{%
0	4395.3\\
2000	4347.2\\
4000	4370.6\\
6000	4444.1\\
14000	4615\\
16000	4640.5\\
18000	4680.7\\
20000	4741.5\\
22000	4785.7\\
28000	4833.6\\
30000	4871.2\\
32000	4889.3\\
34000	4877.6\\
36000	4910.3\\
40000	4935.5\\
42000	4947.4\\
44000	4941.2\\
46000	4912.3\\
48000	4936.4\\
52000	4869.2\\
54000	4880.1\\
56000	4935.8\\
60000	4998.1\\
64000	4993.6\\
68000	5039.3\\
70000	5042.6\\
72000	5071.4\\
84000	5028\\
86000	5067.2\\
90000	5005.8\\
94000	4992.7\\
96000	5027.7\\
1e+05	5016.5\\
};
\addplot [color=mycolor60, forget plot]
  table[row sep=crcr]{%
0	4395.3\\
2000	4221.9\\
4000	4014.4\\
6000	3902.7\\
8000	3836.1\\
10000	3818.1\\
12000	3818\\
14000	3782.5\\
16000	3781.8\\
18000	3801.1\\
24000	3796.6\\
26000	3774.3\\
28000	3767\\
30000	3744.5\\
34000	3726.3\\
36000	3756.4\\
40000	3721.5\\
44000	3732.7\\
48000	3713.1\\
52000	3714.3\\
54000	3721.6\\
56000	3709.8\\
58000	3709.4\\
62000	3667.3\\
66000	3666.6\\
70000	3624.7\\
74000	3635.1\\
76000	3646.9\\
78000	3631.2\\
82000	3631.2\\
84000	3622.5\\
1e+05	3672.8\\
};
\addplot [color=mycolor61, forget plot]
  table[row sep=crcr]{%
0	4395.3\\
2000	4106.2\\
4000	3743\\
6000	3530\\
8000	3422.6\\
10000	3334.3\\
12000	3284.6\\
18000	3235\\
26000	3194.2\\
30000	3212.7\\
32000	3211.5\\
34000	3194.9\\
38000	3214.4\\
42000	3209.9\\
48000	3189.8\\
50000	3204.8\\
52000	3174.6\\
54000	3159.4\\
58000	3173.3\\
60000	3185.9\\
62000	3178\\
64000	3187\\
66000	3180\\
70000	3203.6\\
78000	3186.3\\
80000	3184.8\\
84000	3165.4\\
88000	3170.7\\
94000	3195.1\\
1e+05	3178.7\\
};
\addplot [color=mycolor62, forget plot]
  table[row sep=crcr]{%
0	4395.3\\
2000	4348.7\\
4000	4356.7\\
10000	4553.7\\
14000	4709.7\\
20000	4891.2\\
22000	4919.2\\
24000	4911.6\\
34000	4993.5\\
36000	4990.9\\
40000	5029.3\\
42000	5005.1\\
54000	5063.3\\
56000	5089.2\\
62000	5088.4\\
68000	5050.7\\
70000	5054.8\\
80000	5006\\
84000	5022.9\\
86000	5044.7\\
96000	5035.3\\
1e+05	5064.5\\
};
\addplot [color=mycolor63, forget plot]
  table[row sep=crcr]{%
0	4395.3\\
2000	4105.2\\
4000	3688.4\\
6000	3440.4\\
8000	3258.5\\
14000	2952\\
18000	2857.6\\
20000	2820.6\\
24000	2733\\
28000	2680.6\\
32000	2654.6\\
46000	2546.4\\
48000	2538.2\\
50000	2541.5\\
56000	2506.2\\
60000	2522.6\\
62000	2523.3\\
64000	2511.5\\
70000	2504.7\\
72000	2517.7\\
78000	2509.4\\
80000	2495.4\\
82000	2490.9\\
84000	2495.8\\
92000	2485\\
94000	2488.5\\
96000	2484.2\\
1e+05	2494.1\\
};
\addplot [color=mycolor64, forget plot]
  table[row sep=crcr]{%
0	4395.3\\
2000	4059.6\\
4000	3626.7\\
6000	3389.9\\
8000	3281.2\\
10000	3209\\
12000	3151.4\\
16000	3110.3\\
20000	3108.6\\
24000	3102.2\\
28000	3102.2\\
30000	3092.1\\
34000	3097.6\\
36000	3084.6\\
38000	3088.8\\
40000	3072.1\\
44000	3101.8\\
46000	3092.1\\
48000	3102.1\\
50000	3092.8\\
52000	3098.2\\
56000	3091.3\\
60000	3104.2\\
62000	3104.7\\
64000	3115.9\\
66000	3111.3\\
70000	3119.5\\
72000	3102.4\\
74000	3097\\
80000	3108.4\\
98000	3121.4\\
1e+05	3098.2\\
};
\addplot [color=mycolor65, forget plot]
  table[row sep=crcr]{%
0	4395.3\\
2000	4043.7\\
4000	3567.4\\
6000	3276.3\\
8000	3107.6\\
10000	3012.7\\
16000	2817.2\\
18000	2771.8\\
20000	2755.3\\
24000	2705.4\\
26000	2680.7\\
32000	2637\\
34000	2635.2\\
36000	2624.6\\
46000	2628.4\\
52000	2607.8\\
54000	2600.9\\
58000	2608.9\\
60000	2613.9\\
64000	2596.7\\
74000	2616.6\\
76000	2606.6\\
80000	2599.4\\
82000	2590.3\\
86000	2596.8\\
88000	2591.6\\
90000	2606.9\\
92000	2611.6\\
96000	2601.8\\
1e+05	2613.7\\
};
\addplot [color=mycolor66, forget plot]
  table[row sep=crcr]{%
0	4395.3\\
2000	4121.9\\
4000	3788.7\\
6000	3615.7\\
8000	3531.3\\
10000	3479.2\\
12000	3465.8\\
16000	3455.5\\
18000	3436.4\\
24000	3431\\
26000	3448.2\\
30000	3453\\
40000	3442.6\\
48000	3466.1\\
50000	3476.1\\
54000	3452.2\\
56000	3462.6\\
58000	3454.3\\
64000	3470.4\\
66000	3448.4\\
68000	3449.6\\
70000	3462.8\\
80000	3460.7\\
82000	3463.6\\
86000	3452\\
98000	3484.6\\
1e+05	3486.9\\
};
\addplot [color=mycolor67, forget plot]
  table[row sep=crcr]{%
0	4395.3\\
2000	4218.7\\
4000	3990.3\\
6000	3865.1\\
8000	3762.3\\
10000	3707.1\\
14000	3647.3\\
16000	3595\\
18000	3557\\
20000	3542.4\\
22000	3508.9\\
24000	3492.5\\
26000	3464.1\\
30000	3427.4\\
32000	3397.4\\
36000	3395.4\\
40000	3375.8\\
46000	3328.3\\
50000	3308.9\\
56000	3280.5\\
62000	3290.2\\
66000	3256.2\\
70000	3285.3\\
72000	3290.2\\
74000	3284.2\\
78000	3295.2\\
80000	3289.9\\
84000	3299.4\\
92000	3251.4\\
96000	3272.5\\
1e+05	3268.4\\
};
\addplot [color=mycolor67]
  table[row sep=crcr]{%
0	4395.3\\
2000	3998.8\\
4000	3474.9\\
6000	3167.1\\
8000	2983.9\\
10000	2858.5\\
12000	2789.1\\
14000	2746.6\\
16000	2719.8\\
22000	2681.6\\
24000	2679.9\\
26000	2686.1\\
30000	2674.1\\
34000	2690.5\\
36000	2690.7\\
38000	2702.3\\
42000	2704.4\\
44000	2704.2\\
46000	2713.5\\
54000	2707.5\\
56000	2691.7\\
66000	2691.2\\
68000	2687.2\\
72000	2698.6\\
76000	2691.4\\
78000	2697.5\\
80000	2693\\
84000	2699.4\\
90000	2710.5\\
96000	2703.3\\
1e+05	2707.9\\
};
\addlegendentry{\ClasskIIILegend}

\addplot [color=mycolor68, forget plot]
  table[row sep=crcr]{%
0	4395.3\\
2000	4102\\
4000	3694.1\\
6000	3417.7\\
8000	3233.9\\
10000	3114.4\\
12000	3018.7\\
14000	2956.4\\
16000	2906\\
20000	2823.7\\
26000	2754.2\\
28000	2732.8\\
50000	2617.9\\
62000	2602.2\\
64000	2605\\
72000	2576.1\\
76000	2569.9\\
82000	2574.6\\
90000	2558.6\\
1e+05	2551.9\\
};
\addplot [color=mycolor69, forget plot]
  table[row sep=crcr]{%
0	4395.3\\
2000	4133.4\\
4000	3782.8\\
6000	3562.7\\
8000	3406.2\\
10000	3298.6\\
14000	3139.3\\
20000	3033.5\\
30000	2937.9\\
36000	2934.7\\
38000	2931\\
40000	2937.1\\
50000	2928.7\\
52000	2935.2\\
54000	2926.8\\
58000	2933.9\\
60000	2940.8\\
64000	2934\\
66000	2938.3\\
70000	2931.4\\
76000	2944.8\\
80000	2971.8\\
82000	2975.5\\
88000	2945.2\\
92000	2956.9\\
94000	2967.2\\
96000	2960.2\\
1e+05	2964.9\\
};
\addplot [color=mycolor70, forget plot]
  table[row sep=crcr]{%
0	4395.3\\
2000	3986.1\\
4000	3394.3\\
6000	3005.5\\
8000	2735.8\\
10000	2559.9\\
12000	2427.1\\
14000	2339.5\\
16000	2271.4\\
18000	2221\\
22000	2165.4\\
24000	2143.4\\
28000	2132.6\\
30000	2140.1\\
36000	2140.2\\
40000	2140.3\\
44000	2129.2\\
46000	2140.1\\
50000	2139.2\\
58000	2160\\
60000	2149.1\\
62000	2146.4\\
68000	2168\\
70000	2164.1\\
72000	2168.6\\
78000	2162.5\\
82000	2164.6\\
86000	2163.5\\
90000	2166.6\\
96000	2161.5\\
1e+05	2167.5\\
};
\addplot [color=mycolor71, forget plot]
  table[row sep=crcr]{%
0	4395.3\\
2000	4041.4\\
4000	3559.7\\
6000	3264.8\\
8000	3078.7\\
10000	2964.1\\
12000	2885\\
14000	2826.1\\
18000	2753.8\\
20000	2725.4\\
26000	2683.7\\
48000	2621.5\\
52000	2617.9\\
58000	2599.5\\
60000	2606.8\\
62000	2599.9\\
66000	2595.7\\
70000	2589.4\\
72000	2587.3\\
78000	2597.7\\
82000	2588.6\\
84000	2575.3\\
86000	2574.9\\
88000	2583.9\\
92000	2580\\
98000	2582.7\\
1e+05	2580.7\\
};
\addplot [color=mycolor72, forget plot]
  table[row sep=crcr]{%
0	4395.3\\
2000	4086.2\\
4000	3607\\
6000	3296\\
8000	3063.8\\
10000	2893.6\\
12000	2771.1\\
14000	2676.9\\
16000	2601.3\\
20000	2488\\
24000	2400\\
30000	2345.4\\
34000	2341.5\\
36000	2338.3\\
38000	2315.7\\
42000	2290.7\\
46000	2297.1\\
48000	2277.2\\
54000	2270\\
58000	2270.8\\
60000	2261.8\\
64000	2263.5\\
68000	2291\\
72000	2281.3\\
76000	2283.8\\
78000	2280.9\\
80000	2266.4\\
84000	2273.1\\
88000	2290.4\\
94000	2286.2\\
96000	2288\\
1e+05	2272.9\\
};
\addplot [color=mycolor73, forget plot]
  table[row sep=crcr]{%
0	4395.3\\
2000	4076.4\\
4000	3626.3\\
6000	3345\\
8000	3162.8\\
10000	3037.1\\
12000	2954.5\\
16000	2851.6\\
24000	2778.6\\
26000	2769.9\\
34000	2763.2\\
38000	2740.7\\
40000	2745.8\\
42000	2758.9\\
44000	2750.9\\
52000	2749.1\\
54000	2745.4\\
56000	2732.5\\
60000	2728.6\\
68000	2716.9\\
70000	2720.5\\
76000	2707.8\\
78000	2694.5\\
82000	2698.9\\
84000	2691.7\\
86000	2693.9\\
88000	2705\\
90000	2705.9\\
94000	2695\\
1e+05	2693.2\\
};
\addplot [color=mycolor74, forget plot]
  table[row sep=crcr]{%
0	4395.3\\
2000	3977.8\\
4000	3357.6\\
6000	2943.5\\
8000	2647.8\\
10000	2423.4\\
12000	2269.8\\
14000	2143.5\\
16000	2044.1\\
18000	1970.7\\
20000	1911.7\\
22000	1862.6\\
24000	1823.5\\
30000	1734\\
32000	1712\\
34000	1698.8\\
42000	1665.2\\
48000	1632.6\\
50000	1626.6\\
54000	1625.1\\
58000	1616.9\\
60000	1619.7\\
66000	1607.1\\
70000	1621.3\\
74000	1613.7\\
78000	1620.1\\
80000	1617.7\\
82000	1620.4\\
84000	1613.7\\
90000	1613.2\\
92000	1619.2\\
96000	1608.3\\
98000	1608\\
1e+05	1615.8\\
};
\addplot [color=mycolor75, forget plot]
  table[row sep=crcr]{%
0	4395.3\\
2000	4078.9\\
4000	3613.1\\
6000	3291.3\\
8000	3068.4\\
10000	2900.7\\
12000	2766.2\\
14000	2668.6\\
16000	2597.1\\
20000	2491.7\\
24000	2431.3\\
28000	2396\\
32000	2376.1\\
34000	2355.9\\
40000	2325.6\\
42000	2329.8\\
50000	2315.7\\
54000	2303.4\\
60000	2300.8\\
66000	2298.1\\
70000	2306\\
74000	2294.3\\
76000	2293.6\\
78000	2300.4\\
94000	2288.5\\
98000	2283.3\\
1e+05	2289.4\\
};
\addplot [color=mycolor76, forget plot]
  table[row sep=crcr]{%
0	4395.3\\
2000	4014.8\\
4000	3474\\
6000	3116.8\\
8000	2871.5\\
10000	2699.3\\
12000	2574\\
14000	2487.8\\
16000	2420.1\\
20000	2328.2\\
24000	2262.8\\
38000	2170\\
40000	2164.1\\
44000	2140.9\\
56000	2133\\
58000	2125.6\\
62000	2124.9\\
66000	2115.3\\
68000	2122.3\\
72000	2121.6\\
76000	2119\\
84000	2103.8\\
88000	2111.5\\
90000	2118.2\\
94000	2110.9\\
96000	2117.2\\
1e+05	2114.4\\
};
\addplot [color=mycolor77, forget plot]
  table[row sep=crcr]{%
0	4395.3\\
2000	3985.3\\
4000	3397.1\\
6000	2979.1\\
8000	2684\\
10000	2471.6\\
12000	2313\\
14000	2191.2\\
16000	2106.5\\
18000	2037.9\\
20000	1981\\
22000	1937.7\\
26000	1873.3\\
28000	1845.2\\
36000	1769.5\\
42000	1743.1\\
46000	1726.8\\
54000	1711.1\\
58000	1698.6\\
60000	1689.7\\
64000	1688.4\\
66000	1692.5\\
76000	1679.6\\
78000	1679.3\\
80000	1674.1\\
88000	1672.4\\
90000	1672.1\\
92000	1664.9\\
96000	1656.8\\
1e+05	1653.2\\
};
\addplot [color=mycolor78, forget plot]
  table[row sep=crcr]{%
0	4395.3\\
2000	4117.2\\
4000	3692.8\\
6000	3397.3\\
8000	3177.6\\
10000	3015.7\\
12000	2899\\
14000	2809.5\\
16000	2739\\
18000	2688.8\\
24000	2585.7\\
26000	2562.9\\
28000	2548.4\\
30000	2541.9\\
32000	2524.5\\
38000	2501.9\\
42000	2493.4\\
44000	2482.8\\
48000	2480.4\\
52000	2467.6\\
54000	2464\\
60000	2425.4\\
80000	2421.2\\
84000	2396.1\\
92000	2384.6\\
94000	2387.3\\
96000	2397.2\\
1e+05	2377.4\\
};
\addplot [color=mycolor79, forget plot]
  table[row sep=crcr]{%
0	4395.3\\
2000	4050.9\\
4000	3541.4\\
6000	3203.6\\
8000	2966.2\\
10000	2796.7\\
12000	2676.4\\
14000	2587.9\\
16000	2520.7\\
18000	2481.3\\
22000	2436.8\\
32000	2384.2\\
36000	2350.1\\
42000	2314.3\\
44000	2319.5\\
48000	2312.7\\
50000	2318.9\\
54000	2313.7\\
58000	2287.8\\
60000	2294.7\\
64000	2292.7\\
68000	2296.3\\
70000	2280.1\\
74000	2271.9\\
76000	2280.7\\
78000	2272.7\\
82000	2271.7\\
84000	2283.9\\
86000	2281.8\\
90000	2268.3\\
94000	2267.4\\
96000	2264.9\\
98000	2252.2\\
1e+05	2254.5\\
};
\addplot [color=mycolor79, forget plot]
  table[row sep=crcr]{%
0	4395.3\\
2000	4087.1\\
4000	3622.3\\
6000	3294.1\\
8000	3070.2\\
10000	2890\\
12000	2759.1\\
14000	2647.2\\
16000	2565.7\\
18000	2498.9\\
20000	2441.6\\
24000	2355.4\\
30000	2278.5\\
36000	2208.6\\
40000	2179.4\\
46000	2177.2\\
52000	2127\\
56000	2124.1\\
58000	2114.1\\
68000	2107.2\\
70000	2115.4\\
72000	2117.1\\
78000	2105\\
82000	2106.6\\
86000	2101.1\\
88000	2109.3\\
1e+05	2112.2\\
};
\addplot [color=mycolor80, forget plot]
  table[row sep=crcr]{%
0	4395.3\\
2000	3983.2\\
4000	3362.4\\
6000	2915\\
8000	2600.3\\
10000	2377.4\\
12000	2200.5\\
14000	2074.6\\
16000	1984.7\\
18000	1914.3\\
20000	1864\\
22000	1821.2\\
28000	1745.4\\
30000	1722.1\\
34000	1712.2\\
40000	1684.1\\
44000	1682.1\\
50000	1669.5\\
54000	1661.7\\
68000	1658.2\\
70000	1652.9\\
74000	1652\\
76000	1645.4\\
80000	1643.5\\
82000	1637.1\\
84000	1636.2\\
88000	1647.7\\
90000	1648.3\\
96000	1633.2\\
1e+05	1632.3\\
};
\addplot [color=mycolor81, forget plot]
  table[row sep=crcr]{%
0	4395.3\\
2000	3982.9\\
4000	3369.7\\
6000	2930.9\\
8000	2619.8\\
10000	2384.2\\
12000	2213.1\\
14000	2083.6\\
16000	1986.3\\
18000	1905.4\\
20000	1838.2\\
22000	1787.9\\
24000	1747.3\\
26000	1716.5\\
28000	1680.6\\
30000	1656.4\\
34000	1620\\
36000	1615.6\\
38000	1598.5\\
40000	1589.9\\
44000	1584.7\\
48000	1573.1\\
54000	1566.2\\
56000	1560.3\\
58000	1548.9\\
66000	1537.8\\
68000	1529.7\\
82000	1538.1\\
84000	1537.8\\
86000	1541.9\\
88000	1541.4\\
90000	1545.3\\
94000	1534.4\\
1e+05	1537.4\\
};
\addplot [color=mycolor82, forget plot]
  table[row sep=crcr]{%
0	4395.3\\
2000	4031.6\\
4000	3488.8\\
6000	3102.2\\
8000	2819.6\\
10000	2616.2\\
12000	2461.3\\
14000	2338.3\\
16000	2241.1\\
18000	2169.9\\
22000	2067.5\\
24000	2027.3\\
30000	1947.9\\
32000	1931.7\\
34000	1925.4\\
36000	1913\\
38000	1906.6\\
42000	1878.6\\
44000	1863.5\\
56000	1827.8\\
58000	1826.8\\
60000	1818.5\\
62000	1816.2\\
66000	1824\\
70000	1815.5\\
80000	1798.2\\
84000	1801.7\\
88000	1795.9\\
96000	1785.4\\
1e+05	1781.1\\
};
\addplot [color=mycolor83, forget plot]
  table[row sep=crcr]{%
0	4395.3\\
2000	4055.9\\
4000	3542.7\\
6000	3159.3\\
8000	2879.7\\
10000	2666.1\\
12000	2503.9\\
14000	2390.6\\
16000	2298.9\\
18000	2218.6\\
20000	2153.5\\
22000	2103.4\\
24000	2061.1\\
28000	2009.9\\
34000	1945.5\\
38000	1919.9\\
40000	1916.7\\
44000	1896.3\\
48000	1893.4\\
54000	1876.1\\
56000	1878.3\\
62000	1868.1\\
70000	1858.4\\
78000	1845.2\\
82000	1843.4\\
84000	1844.4\\
86000	1838.2\\
90000	1840.3\\
94000	1836.6\\
1e+05	1833\\
};
\addplot [color=mycolor83, forget plot]
  table[row sep=crcr]{%
0	4395.3\\
2000	3945\\
4000	3263.4\\
6000	2773.1\\
8000	2421.2\\
10000	2162.7\\
12000	1972.7\\
14000	1826.8\\
16000	1715.9\\
18000	1635.8\\
20000	1570.9\\
22000	1523.7\\
26000	1454.2\\
28000	1428.7\\
32000	1397.4\\
36000	1377.4\\
42000	1364\\
48000	1361.7\\
52000	1351.7\\
56000	1349.6\\
58000	1341.9\\
66000	1330.4\\
76000	1332.5\\
78000	1329.8\\
86000	1342.8\\
90000	1338.6\\
96000	1335.3\\
1e+05	1343.6\\
};
\addplot [color=mycolor84, forget plot]
  table[row sep=crcr]{%
0	4395.3\\
2000	3986.3\\
4000	3356.8\\
6000	2886.5\\
8000	2540.7\\
10000	2279.7\\
12000	2085.8\\
14000	1929.7\\
16000	1806.6\\
18000	1711.4\\
20000	1643.5\\
24000	1552.5\\
30000	1481.9\\
32000	1464.4\\
34000	1452.4\\
36000	1434.9\\
44000	1397.1\\
48000	1381\\
50000	1381.5\\
62000	1361\\
64000	1361.6\\
68000	1348.7\\
84000	1351.5\\
86000	1340.5\\
88000	1347.7\\
90000	1349.5\\
94000	1342.9\\
96000	1346.9\\
1e+05	1345\\
};
\addplot [color=mycolor85, forget plot]
  table[row sep=crcr]{%
0	4395.3\\
2000	4022.9\\
4000	3439.7\\
6000	3021.9\\
8000	2712.3\\
10000	2469.7\\
12000	2280.3\\
14000	2135.4\\
16000	2020.2\\
18000	1929\\
20000	1857.6\\
22000	1799.2\\
26000	1707.1\\
30000	1647.8\\
36000	1583.9\\
40000	1560.8\\
42000	1552.8\\
46000	1545.9\\
52000	1522.9\\
54000	1518.1\\
56000	1524.9\\
58000	1523.7\\
66000	1483.8\\
70000	1482.3\\
72000	1480\\
74000	1486.7\\
76000	1481.1\\
78000	1482.5\\
82000	1476.8\\
84000	1481.8\\
94000	1476.1\\
98000	1480.5\\
1e+05	1470.6\\
};
\addplot [color=mycolor86, forget plot]
  table[row sep=crcr]{%
0	4395.3\\
2000	3995.9\\
4000	3366.4\\
6000	2915.5\\
8000	2573\\
10000	2314.8\\
12000	2113.3\\
14000	1959.3\\
16000	1831.3\\
18000	1729.8\\
20000	1654.1\\
26000	1503.9\\
30000	1443.5\\
32000	1425.4\\
34000	1401.7\\
38000	1366.7\\
42000	1348.7\\
44000	1350.3\\
50000	1326.5\\
54000	1314.3\\
60000	1291.8\\
64000	1291.1\\
66000	1281\\
68000	1275.9\\
72000	1273.2\\
76000	1279.6\\
80000	1280.2\\
86000	1269.9\\
88000	1269.7\\
90000	1263\\
1e+05	1264.3\\
};
\addplot [color=mycolor87, forget plot]
  table[row sep=crcr]{%
0	4395.3\\
2000	3986\\
4000	3329.2\\
6000	2825.3\\
8000	2411.3\\
10000	2094.1\\
12000	1860.8\\
14000	1692.7\\
16000	1554.4\\
18000	1436.8\\
20000	1347.2\\
24000	1230.7\\
26000	1190.8\\
28000	1158.6\\
30000	1123.5\\
32000	1092.8\\
34000	1072.5\\
36000	1059.6\\
40000	1045.9\\
44000	1048.1\\
48000	1029.5\\
50000	1018.8\\
62000	1011.8\\
64000	1014.7\\
70000	1044.1\\
72000	1043.6\\
78000	1031.2\\
80000	1033.8\\
84000	1049.4\\
88000	1050.3\\
90000	1044.9\\
92000	1047.8\\
94000	1047.2\\
96000	1042.8\\
1e+05	1046.8\\
};
\addplot [color=mycolor88, forget plot]
  table[row sep=crcr]{%
0	4395.3\\
2000	4022.4\\
4000	3419.2\\
6000	2963.8\\
8000	2624.5\\
10000	2360.8\\
12000	2162.4\\
14000	2003.4\\
16000	1875.4\\
18000	1767.2\\
20000	1685.9\\
22000	1617.6\\
28000	1481.7\\
30000	1450\\
32000	1424.9\\
36000	1391.4\\
40000	1364\\
48000	1332.5\\
52000	1329\\
56000	1314.3\\
58000	1307.9\\
62000	1306\\
66000	1286.6\\
70000	1281.8\\
72000	1280.6\\
74000	1284\\
80000	1274.1\\
88000	1274\\
96000	1263.3\\
1e+05	1270.1\\
};
\addplot [color=mycolor89, forget plot]
  table[row sep=crcr]{%
0	4395.3\\
2000	4022.1\\
4000	3437.1\\
6000	2957.9\\
8000	2575.7\\
10000	2280.5\\
12000	2046.3\\
14000	1854.4\\
18000	1565.2\\
20000	1454.8\\
22000	1366.3\\
26000	1237.4\\
30000	1140.9\\
32000	1104.2\\
34000	1071.9\\
36000	1048.1\\
42000	999.11\\
44000	985.75\\
46000	969.74\\
48000	959.23\\
56000	942.88\\
58000	936.47\\
66000	929.25\\
68000	930.34\\
74000	922.31\\
76000	917.2\\
84000	915.47\\
86000	911.69\\
90000	914.14\\
92000	916.71\\
1e+05	906.91\\
};
\addplot [color=mycolor90, forget plot]
  table[row sep=crcr]{%
0	4395.3\\
2000	3961.5\\
4000	3275.2\\
6000	2747.1\\
8000	2347.3\\
10000	2036.4\\
12000	1796.9\\
14000	1605.1\\
16000	1461.5\\
18000	1350.7\\
20000	1262.3\\
22000	1185.8\\
24000	1126.2\\
26000	1084.4\\
30000	1018.2\\
36000	957.17\\
38000	940.75\\
40000	930.89\\
42000	924.3\\
44000	912.99\\
46000	905.44\\
50000	897.84\\
54000	897.62\\
66000	892.41\\
74000	882.19\\
76000	882.16\\
78000	879.06\\
80000	882.38\\
84000	880.93\\
86000	875.98\\
88000	875.42\\
92000	879.08\\
1e+05	873.34\\
};
\addplot [color=mycolor91, forget plot]
  table[row sep=crcr]{%
0	4395.3\\
2000	4033.2\\
4000	3430.1\\
6000	2954.2\\
8000	2569.4\\
10000	2261.8\\
12000	2007.5\\
14000	1798.9\\
16000	1628.9\\
18000	1490.9\\
20000	1372.8\\
22000	1283.5\\
26000	1141.8\\
28000	1079.9\\
30000	1031.5\\
32000	991.53\\
34000	956.46\\
36000	925.71\\
38000	905.92\\
42000	872.01\\
44000	853.04\\
46000	837.73\\
50000	821.45\\
54000	808.39\\
56000	798.88\\
58000	794.23\\
60000	795.63\\
62000	792.03\\
64000	784.54\\
70000	777.52\\
72000	775.2\\
76000	775.5\\
78000	771.11\\
84000	765.06\\
88000	766.08\\
92000	765.39\\
94000	768.06\\
96000	775.02\\
98000	776.98\\
1e+05	776.52\\
};
\addplot [color=mycolor92, forget plot]
  table[row sep=crcr]{%
0	4395.3\\
2000	4003.8\\
4000	3380.7\\
6000	2908.2\\
8000	2538.6\\
10000	2245.9\\
12000	2017.6\\
14000	1842.2\\
16000	1692.8\\
18000	1570.1\\
20000	1469.7\\
22000	1391.6\\
24000	1328.7\\
26000	1273.1\\
30000	1192.5\\
32000	1160.8\\
36000	1116\\
38000	1092\\
40000	1072.7\\
42000	1057.9\\
44000	1048.6\\
46000	1036.4\\
48000	1032.9\\
50000	1024.1\\
62000	1000.5\\
66000	999.19\\
70000	990.7\\
90000	962.44\\
94000	961.25\\
98000	962.16\\
1e+05	959.83\\
};
\addplot [color=mycolor93, forget plot]
  table[row sep=crcr]{%
0	4395.3\\
2000	4067.5\\
4000	3512.7\\
6000	3062.9\\
8000	2698.7\\
10000	2394.9\\
12000	2146.1\\
14000	1931.7\\
16000	1755.2\\
18000	1605.4\\
20000	1488.4\\
22000	1390.3\\
24000	1312\\
26000	1241.9\\
28000	1180.5\\
30000	1134.9\\
36000	1020.7\\
40000	970\\
42000	953.66\\
44000	930.59\\
46000	917.72\\
48000	900.13\\
50000	888.5\\
54000	875.77\\
56000	871.97\\
58000	864.79\\
62000	847.43\\
64000	839.78\\
68000	836.47\\
72000	830.73\\
76000	824.88\\
78000	821\\
84000	816.47\\
86000	812.44\\
90000	810.11\\
92000	807.36\\
94000	808.03\\
98000	801.64\\
1e+05	799.42\\
};
\addplot [color=mycolor94, forget plot]
  table[row sep=crcr]{%
0	4395.3\\
2000	4071.6\\
4000	3532.5\\
6000	3085.7\\
8000	2724.7\\
10000	2418.6\\
12000	2172.5\\
14000	1960.8\\
16000	1779.4\\
18000	1631.7\\
20000	1502.2\\
22000	1393\\
24000	1307.3\\
26000	1235.4\\
30000	1114\\
34000	1022.9\\
38000	958.05\\
40000	931.73\\
44000	891.49\\
50000	844.62\\
52000	833.91\\
56000	819.84\\
62000	804.15\\
64000	799.47\\
68000	796.29\\
70000	790.66\\
74000	785.45\\
76000	784.62\\
78000	779.8\\
82000	773.22\\
84000	775.92\\
94000	770.1\\
96000	771.99\\
1e+05	768.92\\
};
\addplot [color=mycolor95, forget plot]
  table[row sep=crcr]{%
0	4395.3\\
2000	4007.9\\
4000	3365.8\\
6000	2852.7\\
8000	2438.3\\
10000	2095.1\\
12000	1820.7\\
16000	1416.7\\
18000	1261.8\\
20000	1134\\
22000	1028.6\\
24000	943.6\\
28000	806.32\\
32000	710.06\\
34000	674.34\\
36000	643.68\\
38000	617.29\\
40000	595.58\\
44000	563.81\\
48000	540.35\\
52000	521.82\\
60000	492.62\\
66000	483.93\\
68000	479.76\\
72000	474.09\\
74000	475.35\\
78000	473.69\\
80000	473.43\\
90000	462.96\\
92000	462.14\\
96000	462.29\\
98000	462.65\\
1e+05	461.5\\
};
\addplot [color=mycolor96, forget plot]
  table[row sep=crcr]{%
0	4395.3\\
2000	4040.8\\
4000	3465.5\\
6000	3001.6\\
8000	2610.6\\
10000	2283\\
14000	1802.9\\
16000	1608\\
18000	1448.9\\
20000	1315.4\\
24000	1095.3\\
26000	1007.4\\
28000	931.86\\
30000	868.96\\
34000	763.44\\
36000	724.99\\
38000	690.48\\
40000	660.19\\
42000	636.19\\
44000	615.92\\
48000	591.3\\
50000	577.07\\
52000	566.62\\
54000	563.04\\
62000	531.69\\
64000	528.47\\
70000	509.76\\
72000	504.86\\
76000	505.37\\
78000	501.49\\
80000	499.26\\
82000	499.39\\
84000	495.69\\
88000	494.64\\
90000	487.63\\
92000	488.06\\
94000	489.95\\
1e+05	490.96\\
};
\addplot [color=mycolor97, forget plot]
  table[row sep=crcr]{%
0	4395.3\\
2000	3960.1\\
4000	3256\\
6000	2703.6\\
8000	2268.4\\
10000	1915.2\\
12000	1639.5\\
14000	1423.5\\
16000	1242.1\\
18000	1101.3\\
20000	983.22\\
22000	885.39\\
24000	809.48\\
26000	742.58\\
30000	634.27\\
32000	593.79\\
38000	503.52\\
42000	461.94\\
44000	446.03\\
46000	433.3\\
48000	423.34\\
50000	416.55\\
52000	408.37\\
56000	389.56\\
58000	386.04\\
60000	385.31\\
68000	376.15\\
72000	368.19\\
76000	364.62\\
78000	362.83\\
80000	363.45\\
84000	362.92\\
90000	365.48\\
92000	362.17\\
96000	361.84\\
1e+05	359.49\\
};
\addplot [color=mycolor98]
  table[row sep=crcr]{%
0	4395.3\\
2000	4073.6\\
4000	3529.1\\
6000	3070.9\\
10000	2391.1\\
12000	2122.5\\
16000	1688.4\\
18000	1530.2\\
20000	1393.1\\
22000	1272.3\\
24000	1167.4\\
26000	1076.4\\
28000	995.73\\
30000	923.77\\
32000	865.18\\
34000	818.44\\
38000	739.15\\
44000	651.14\\
46000	630.84\\
48000	614.97\\
60000	539.27\\
62000	530.21\\
64000	519.05\\
66000	513.88\\
68000	506.93\\
70000	505.53\\
72000	498.62\\
76000	490.38\\
82000	487.38\\
88000	482.68\\
90000	482.89\\
96000	476.54\\
98000	478.12\\
1e+05	476.88\\
};
\addlegendentry{\ClasskIVLegend}

      % \input{../../../../.tex/20/KS/SizeEvolutionsfig1.tex}

      \nextgroupplot[%
        % xmin=0,xmax=1e+05,
        % ymin=100,ymax=1.4472e+05,
      ] % This file was created by matlab2tikz.
%
\definecolor{mycolor1}{rgb}{0.00146,0.00047,0.01387}%
\definecolor{mycolor2}{rgb}{0.01166,0.00942,0.06346}%
\definecolor{mycolor3}{rgb}{0.02943,0.02150,0.11462}%
\definecolor{mycolor4}{rgb}{0.06134,0.03659,0.17764}%
\definecolor{mycolor5}{rgb}{0.08741,0.04456,0.22481}%
\definecolor{mycolor6}{rgb}{0.12291,0.04754,0.28162}%
\definecolor{mycolor7}{rgb}{0.14907,0.04547,0.31709}%
\definecolor{mycolor8}{rgb}{0.18343,0.04033,0.35497}%
\definecolor{mycolor9}{rgb}{0.21795,0.03662,0.38352}%
\definecolor{mycolor10}{rgb}{0.24497,0.03705,0.40001}%
\definecolor{mycolor11}{rgb}{0.27135,0.04092,0.41198}%
\definecolor{mycolor12}{rgb}{0.29076,0.04564,0.41864}%
\definecolor{mycolor13}{rgb}{0.31628,0.05349,0.42512}%
\definecolor{mycolor14}{rgb}{0.33522,0.06006,0.42852}%
\definecolor{mycolor15}{rgb}{0.36028,0.06925,0.43150}%
\definecolor{mycolor16}{rgb}{0.37900,0.07625,0.43272}%
\definecolor{mycolor17}{rgb}{0.39767,0.08326,0.43318}%
\definecolor{mycolor18}{rgb}{0.41633,0.09020,0.43294}%
\definecolor{mycolor19}{rgb}{0.42877,0.09479,0.43241}%
\definecolor{mycolor20}{rgb}{0.44743,0.10160,0.43108}%
\definecolor{mycolor21}{rgb}{0.46610,0.10832,0.42912}%
\definecolor{mycolor22}{rgb}{0.47856,0.11276,0.42747}%
\definecolor{mycolor23}{rgb}{0.49726,0.11938,0.42449}%
\definecolor{mycolor24}{rgb}{0.50973,0.12377,0.42216}%
\definecolor{mycolor25}{rgb}{0.52221,0.12815,0.41955}%
\definecolor{mycolor26}{rgb}{0.53468,0.13253,0.41667}%
\definecolor{mycolor27}{rgb}{0.55339,0.13913,0.41183}%
\definecolor{mycolor28}{rgb}{0.56585,0.14357,0.40826}%
\definecolor{mycolor29}{rgb}{0.57830,0.14804,0.40441}%
\definecolor{mycolor30}{rgb}{0.59073,0.15256,0.40029}%
\definecolor{mycolor31}{rgb}{0.60314,0.15715,0.39589}%
\definecolor{mycolor32}{rgb}{0.61551,0.16182,0.39122}%
\definecolor{mycolor33}{rgb}{0.62169,0.16418,0.38878}%
\definecolor{mycolor34}{rgb}{0.63400,0.16899,0.38370}%
\definecolor{mycolor35}{rgb}{0.64626,0.17391,0.37836}%
\definecolor{mycolor36}{rgb}{0.65846,0.17896,0.37275}%
\definecolor{mycolor37}{rgb}{0.66454,0.18154,0.36985}%
\definecolor{mycolor38}{rgb}{0.67664,0.18681,0.36385}%
\definecolor{mycolor39}{rgb}{0.68865,0.19224,0.35760}%
\definecolor{mycolor40}{rgb}{0.69463,0.19502,0.35439}%
\definecolor{mycolor41}{rgb}{0.70650,0.20073,0.34778}%
\definecolor{mycolor42}{rgb}{0.71240,0.20366,0.34438}%
\definecolor{mycolor43}{rgb}{0.72410,0.20967,0.33742}%
\definecolor{mycolor44}{rgb}{0.72991,0.21276,0.33386}%
\definecolor{mycolor45}{rgb}{0.74142,0.21911,0.32658}%
\definecolor{mycolor46}{rgb}{0.74713,0.22238,0.32286}%
\definecolor{mycolor47}{rgb}{0.75279,0.22571,0.31909}%
\definecolor{mycolor48}{rgb}{0.76401,0.23255,0.31140}%
\definecolor{mycolor49}{rgb}{0.76956,0.23608,0.30749}%
\definecolor{mycolor50}{rgb}{0.77506,0.23967,0.30353}%
\definecolor{mycolor51}{rgb}{0.79129,0.25086,0.29139}%
\definecolor{mycolor52}{rgb}{0.79661,0.25473,0.28726}%
\definecolor{mycolor53}{rgb}{0.80187,0.25867,0.28310}%
\definecolor{mycolor54}{rgb}{0.81224,0.26679,0.27466}%
\definecolor{mycolor55}{rgb}{0.81734,0.27095,0.27039}%
\definecolor{mycolor56}{rgb}{0.82737,0.27952,0.26175}%
\definecolor{mycolor57}{rgb}{0.83230,0.28391,0.25738}%
\definecolor{mycolor58}{rgb}{0.83717,0.28839,0.25299}%
\definecolor{mycolor59}{rgb}{0.84671,0.29756,0.24411}%
\definecolor{mycolor60}{rgb}{0.85138,0.30226,0.23964}%
\definecolor{mycolor61}{rgb}{0.85599,0.30704,0.23513}%
\definecolor{mycolor62}{rgb}{0.86053,0.31189,0.23061}%
\definecolor{mycolor63}{rgb}{0.87374,0.32691,0.21689}%
\definecolor{mycolor64}{rgb}{0.87800,0.33206,0.21227}%
\definecolor{mycolor65}{rgb}{0.89034,0.34796,0.19829}%
\definecolor{mycolor66}{rgb}{0.89819,0.35891,0.18886}%
\definecolor{mycolor67}{rgb}{0.90200,0.36449,0.18412}%
\definecolor{mycolor68}{rgb}{0.90573,0.37014,0.17935}%
\definecolor{mycolor69}{rgb}{0.90939,0.37586,0.17456}%
\definecolor{mycolor70}{rgb}{0.91297,0.38164,0.16975}%
\definecolor{mycolor71}{rgb}{0.91646,0.38748,0.16492}%
\definecolor{mycolor72}{rgb}{0.91988,0.39339,0.16007}%
\definecolor{mycolor73}{rgb}{0.92322,0.39936,0.15519}%
\definecolor{mycolor74}{rgb}{0.92647,0.40539,0.15029}%
\definecolor{mycolor75}{rgb}{0.92964,0.41148,0.14537}%
\definecolor{mycolor76}{rgb}{0.93274,0.41763,0.14042}%
\definecolor{mycolor77}{rgb}{0.93575,0.42383,0.13544}%
\definecolor{mycolor78}{rgb}{0.93868,0.43009,0.13044}%
\definecolor{mycolor79}{rgb}{0.94152,0.43640,0.12541}%
\definecolor{mycolor80}{rgb}{0.94429,0.44277,0.12035}%
\definecolor{mycolor81}{rgb}{0.94956,0.45566,0.11016}%
\definecolor{mycolor82}{rgb}{0.95208,0.46218,0.10503}%
\definecolor{mycolor83}{rgb}{0.96129,0.48872,0.08429}%
\definecolor{mycolor84}{rgb}{0.96540,0.50225,0.07386}%
\definecolor{mycolor85}{rgb}{0.96732,0.50908,0.06866}%
\definecolor{mycolor86}{rgb}{0.97259,0.52980,0.05332}%
\definecolor{mycolor87}{rgb}{0.97568,0.54380,0.04362}%
\definecolor{mycolor88}{rgb}{0.98082,0.57221,0.02851}%
\definecolor{mycolor89}{rgb}{0.98189,0.57939,0.02625}%
\definecolor{mycolor90}{rgb}{0.98378,0.59385,0.02377}%
\definecolor{mycolor91}{rgb}{0.98532,0.60842,0.02420}%
\definecolor{mycolor92}{rgb}{0.98696,0.63048,0.03091}%
\definecolor{mycolor93}{rgb}{0.98793,0.66025,0.05175}%
\definecolor{mycolor94}{rgb}{0.98771,0.68281,0.07249}%
\definecolor{mycolor95}{rgb}{0.97750,0.78226,0.18592}%
\definecolor{mycolor96}{rgb}{0.96624,0.83619,0.26153}%
\definecolor{mycolor97}{rgb}{0.96252,0.85148,0.28555}%
\definecolor{mycolor98}{rgb}{0.98836,0.99836,0.64492}%
%
\addplot [color=mycolor1]
  table[row sep=crcr]{%
0	4395.3\\
2000	16346\\
4000	36407\\
6000	52907\\
8000	66810\\
10000	78978\\
12000	89344\\
14000	97986\\
16000	1.0542e+05\\
18000	1.1154e+05\\
20000	1.1678e+05\\
22000	1.2118e+05\\
24000	1.2485e+05\\
26000	1.2825e+05\\
28000	1.3097e+05\\
32000	1.3542e+05\\
38000	1.3893e+05\\
40000	1.3999e+05\\
46000	1.4173e+05\\
52000	1.4271e+05\\
56000	1.4304e+05\\
62000	1.4336e+05\\
82000	1.4453e+05\\
84000	1.4472e+05\\
86000	1.4438e+05\\
1e+05	1.4466e+05\\
};
\addlegendentry{\ClasskILegend}

\addplot [color=mycolor2, forget plot]
  table[row sep=crcr]{%
0	4395.3\\
2000	8966.7\\
4000	14858\\
6000	18167\\
8000	20461\\
10000	21856\\
12000	22844\\
14000	23356\\
16000	23387\\
18000	23734\\
22000	24168\\
24000	24428\\
26000	24431\\
28000	24851\\
32000	24757\\
34000	24910\\
36000	24963\\
38000	24835\\
40000	24577\\
44000	25060\\
46000	25035\\
48000	25221\\
52000	25376\\
54000	25061\\
58000	25360\\
60000	25348\\
66000	25075\\
68000	24843\\
72000	25188\\
74000	25027\\
76000	25016\\
80000	25184\\
84000	24813\\
86000	25085\\
90000	24989\\
92000	25054\\
96000	24833\\
98000	24998\\
1e+05	24980\\
};
\addplot [color=mycolor3, forget plot]
  table[row sep=crcr]{%
0	4395.3\\
2000	10003\\
4000	18079\\
6000	23624\\
8000	27494\\
10000	30138\\
12000	32190\\
14000	33457\\
16000	34658\\
18000	35483\\
20000	36100\\
24000	36774\\
28000	37020\\
30000	37600\\
32000	37786\\
34000	37534\\
40000	38083\\
42000	38260\\
48000	37992\\
50000	37638\\
52000	37434\\
54000	37707\\
56000	37748\\
58000	37983\\
60000	38100\\
64000	38048\\
68000	38579\\
74000	38271\\
76000	38217\\
82000	38558\\
86000	38596\\
90000	38524\\
96000	38139\\
1e+05	38113\\
};
\addplot [color=mycolor4, forget plot]
  table[row sep=crcr]{%
0	4395.3\\
2000	7142\\
4000	10426\\
6000	12010\\
8000	12914\\
10000	13264\\
12000	13574\\
14000	13587\\
16000	13736\\
18000	13841\\
20000	13875\\
22000	14034\\
26000	13887\\
28000	13876\\
30000	13792\\
32000	13888\\
34000	14077\\
36000	13982\\
38000	13991\\
40000	13960\\
42000	14062\\
44000	13962\\
48000	14061\\
50000	14014\\
52000	14092\\
54000	14025\\
56000	13861\\
58000	13978\\
60000	13924\\
62000	13977\\
64000	13816\\
66000	13894\\
68000	14220\\
72000	14123\\
74000	14006\\
76000	14001\\
78000	14119\\
80000	14028\\
82000	13854\\
84000	13804\\
86000	13908\\
88000	13951\\
90000	13844\\
92000	13875\\
94000	13964\\
96000	13907\\
98000	13791\\
1e+05	13863\\
};
\addplot [color=mycolor5, forget plot]
  table[row sep=crcr]{%
0	4395.3\\
2000	9770.4\\
4000	18613\\
6000	25846\\
8000	31687\\
10000	36548\\
12000	40344\\
14000	44006\\
16000	46743\\
18000	49159\\
20000	50978\\
24000	54077\\
26000	55379\\
30000	57168\\
40000	59367\\
44000	60124\\
46000	60283\\
48000	60212\\
56000	60694\\
60000	61305\\
62000	61398\\
66000	60950\\
68000	60703\\
72000	61009\\
78000	60256\\
86000	60876\\
88000	60854\\
92000	60193\\
96000	60587\\
98000	60478\\
1e+05	60620\\
};
\addplot [color=mycolor6, forget plot]
  table[row sep=crcr]{%
0	4395.3\\
2000	9463.1\\
4000	17618\\
6000	24266\\
8000	29647\\
10000	33750\\
12000	37163\\
14000	40086\\
16000	42469\\
18000	44556\\
20000	46325\\
26000	49737\\
34000	52154\\
36000	52676\\
38000	52848\\
42000	52861\\
46000	53021\\
48000	53264\\
54000	53275\\
56000	53051\\
58000	53046\\
66000	53518\\
70000	53418\\
72000	53286\\
74000	53701\\
76000	53786\\
82000	53294\\
86000	53597\\
92000	53774\\
94000	53273\\
96000	53123\\
1e+05	53404\\
};
\addplot [color=mycolor7, forget plot]
  table[row sep=crcr]{%
0	4395.3\\
2000	7895.1\\
4000	13208\\
6000	17014\\
8000	19722\\
10000	21616\\
12000	22980\\
14000	24046\\
16000	24717\\
18000	25106\\
20000	25653\\
22000	26011\\
24000	26027\\
26000	26357\\
28000	26528\\
30000	26513\\
40000	27094\\
42000	27323\\
44000	27293\\
46000	27018\\
48000	27244\\
50000	27345\\
52000	27008\\
54000	26805\\
58000	27077\\
60000	27161\\
64000	27168\\
66000	27339\\
68000	27327\\
72000	27157\\
76000	27440\\
78000	27416\\
80000	27230\\
82000	27220\\
86000	27380\\
88000	27403\\
90000	27578\\
92000	27299\\
96000	27247\\
98000	27506\\
1e+05	27427\\
};
\addplot [color=mycolor8, forget plot]
  table[row sep=crcr]{%
0	4395.3\\
2000	7116.2\\
4000	11044\\
6000	13719\\
8000	15613\\
10000	16801\\
12000	17633\\
14000	18210\\
16000	18677\\
18000	18972\\
20000	19218\\
26000	19665\\
28000	19697\\
30000	19836\\
34000	19772\\
38000	19732\\
40000	19635\\
42000	19381\\
46000	19865\\
52000	20266\\
54000	20003\\
56000	19959\\
60000	20032\\
62000	20136\\
64000	20105\\
66000	19977\\
72000	20056\\
74000	19886\\
76000	19823\\
78000	19837\\
80000	19778\\
82000	19787\\
88000	20003\\
90000	20010\\
94000	19757\\
98000	19655\\
1e+05	19814\\
};
\addplot [color=mycolor9, forget plot]
  table[row sep=crcr]{%
0	4395.3\\
2000	6447.9\\
4000	9190.5\\
6000	10929\\
8000	11985\\
10000	12555\\
12000	12820\\
14000	13133\\
18000	13488\\
22000	13688\\
24000	13671\\
28000	13735\\
36000	13653\\
38000	13560\\
40000	13639\\
44000	13619\\
46000	13569\\
48000	13724\\
50000	13732\\
52000	13829\\
54000	13882\\
56000	13896\\
58000	13758\\
62000	13713\\
64000	13738\\
66000	13660\\
70000	13800\\
72000	13661\\
78000	13959\\
80000	13907\\
82000	13734\\
84000	13693\\
88000	13736\\
90000	13839\\
92000	13828\\
94000	13677\\
98000	13703\\
1e+05	13681\\
};
\addplot [color=mycolor10, forget plot]
  table[row sep=crcr]{%
0	4395.3\\
2000	6764.8\\
4000	10291\\
6000	12871\\
8000	14809\\
10000	16203\\
12000	17227\\
14000	18060\\
16000	18714\\
18000	19258\\
20000	19624\\
26000	20502\\
28000	20743\\
30000	20903\\
38000	21283\\
42000	21041\\
44000	21105\\
46000	21287\\
50000	21378\\
52000	21358\\
54000	21461\\
56000	21466\\
58000	21322\\
64000	21237\\
66000	21296\\
70000	21098\\
72000	21052\\
76000	21157\\
78000	21342\\
80000	21385\\
82000	21353\\
86000	21530\\
88000	21476\\
90000	21481\\
92000	21346\\
94000	21307\\
96000	21335\\
98000	21244\\
1e+05	21353\\
};
\addplot [color=mycolor11, forget plot]
  table[row sep=crcr]{%
0	4395.3\\
2000	5586.2\\
4000	7069.9\\
6000	7853.9\\
8000	8244.3\\
10000	8496.7\\
12000	8614.2\\
14000	8613.4\\
16000	8582.7\\
18000	8623.9\\
20000	8580.3\\
22000	8507.3\\
26000	8675.1\\
28000	8605.9\\
30000	8579.6\\
32000	8613.1\\
34000	8582.2\\
38000	8613\\
42000	8403.5\\
46000	8521\\
48000	8544.8\\
50000	8533.3\\
52000	8576.9\\
56000	8489.6\\
58000	8544.3\\
60000	8476.6\\
66000	8468\\
68000	8556.3\\
70000	8537\\
72000	8440.7\\
76000	8487\\
80000	8442.5\\
82000	8477.5\\
84000	8567.9\\
88000	8677.1\\
90000	8579.9\\
92000	8626\\
94000	8610.2\\
96000	8664.1\\
1e+05	8548.3\\
};
\addplot [color=mycolor12, forget plot]
  table[row sep=crcr]{%
0	4395.3\\
2000	6312.6\\
4000	9096.1\\
6000	10813\\
8000	11930\\
10000	12791\\
12000	13291\\
14000	13644\\
18000	13960\\
20000	14025\\
22000	14031\\
24000	14186\\
28000	14018\\
32000	14350\\
36000	14285\\
40000	14446\\
42000	14483\\
44000	14361\\
46000	14419\\
50000	14436\\
52000	14520\\
54000	14532\\
56000	14435\\
58000	14583\\
60000	14507\\
62000	14534\\
64000	14493\\
68000	14584\\
70000	14524\\
72000	14274\\
74000	14260\\
76000	14418\\
78000	14421\\
80000	14313\\
86000	14278\\
90000	14370\\
92000	14430\\
96000	14607\\
1e+05	14384\\
};
\addplot [color=mycolor13, forget plot]
  table[row sep=crcr]{%
0	4395.3\\
4000	7895.4\\
6000	9147.7\\
8000	9754.3\\
10000	10213\\
12000	10496\\
16000	10754\\
18000	10947\\
22000	11145\\
26000	11106\\
28000	11007\\
30000	11086\\
32000	11013\\
36000	11004\\
38000	10900\\
40000	10867\\
42000	10973\\
44000	10874\\
48000	10792\\
50000	10900\\
52000	10806\\
60000	10969\\
62000	10932\\
64000	11064\\
66000	11156\\
72000	11121\\
74000	11127\\
76000	11102\\
80000	11218\\
82000	11169\\
84000	11082\\
86000	11111\\
88000	11071\\
90000	11212\\
92000	11090\\
96000	11024\\
98000	11111\\
1e+05	11043\\
};
\addplot [color=mycolor14, forget plot]
  table[row sep=crcr]{%
0	4395.3\\
2000	7047.5\\
4000	11111\\
6000	14157\\
8000	16349\\
10000	17944\\
12000	19193\\
14000	20092\\
16000	20848\\
18000	21505\\
20000	21841\\
22000	22297\\
24000	22540\\
28000	22944\\
30000	22954\\
32000	23088\\
36000	23165\\
38000	23264\\
40000	23506\\
44000	23585\\
46000	23502\\
48000	23295\\
50000	23282\\
52000	23561\\
56000	23745\\
60000	23528\\
62000	23305\\
64000	23453\\
66000	23507\\
68000	23411\\
70000	23423\\
72000	23513\\
74000	23468\\
78000	23521\\
86000	23203\\
88000	23167\\
90000	23317\\
92000	23197\\
94000	23202\\
96000	23367\\
98000	23221\\
1e+05	23418\\
};
\addplot [color=mycolor15, forget plot]
  table[row sep=crcr]{%
0	4395.3\\
2000	6510.2\\
4000	9526.1\\
6000	11426\\
8000	12831\\
10000	13820\\
12000	14499\\
16000	15784\\
18000	16039\\
20000	16150\\
24000	16156\\
26000	16194\\
28000	16323\\
30000	16341\\
32000	16264\\
36000	16502\\
40000	16544\\
44000	16482\\
46000	16600\\
48000	16514\\
54000	16537\\
56000	16445\\
58000	16476\\
64000	16447\\
66000	16515\\
70000	16549\\
72000	16672\\
74000	16578\\
76000	16545\\
78000	16435\\
80000	16670\\
82000	16632\\
84000	16542\\
88000	16615\\
90000	16669\\
94000	16714\\
96000	16571\\
98000	16560\\
1e+05	16701\\
};
\addplot [color=mycolor16, forget plot]
  table[row sep=crcr]{%
0	4395.3\\
2000	5142.7\\
4000	6073.2\\
6000	6668.9\\
8000	6927\\
10000	7041.2\\
12000	7097.6\\
14000	7078.7\\
20000	7114.5\\
22000	7189.2\\
28000	7032.9\\
32000	7038.2\\
38000	6981.3\\
40000	7010.9\\
42000	7081.7\\
44000	7114\\
46000	7029.2\\
48000	6966.9\\
52000	7002.1\\
54000	6978.5\\
58000	7033.7\\
64000	6926.1\\
68000	6978.2\\
72000	7029.9\\
74000	7007.1\\
76000	6875.1\\
78000	7018.9\\
80000	7122.8\\
82000	7167.6\\
86000	7013.1\\
90000	6948.2\\
92000	6961.3\\
94000	7012.6\\
96000	6971.1\\
1e+05	7028.2\\
};
\addplot [color=mycolor17, forget plot]
  table[row sep=crcr]{%
0	4395.3\\
4000	9139.6\\
6000	11153\\
8000	12497\\
10000	13340\\
12000	13894\\
14000	14381\\
16000	14640\\
22000	15121\\
24000	15241\\
32000	15490\\
34000	15600\\
36000	15509\\
38000	15628\\
40000	15619\\
42000	15707\\
44000	15680\\
46000	15572\\
48000	15523\\
52000	15737\\
56000	15564\\
60000	15460\\
62000	15385\\
68000	15402\\
72000	15545\\
74000	15717\\
78000	15755\\
82000	15893\\
88000	15481\\
90000	15470\\
92000	15584\\
96000	15662\\
98000	15638\\
1e+05	15535\\
};
\addplot [color=mycolor18, forget plot]
  table[row sep=crcr]{%
0	4395.3\\
2000	6892.9\\
4000	10664\\
6000	13583\\
8000	15785\\
10000	17529\\
12000	18861\\
14000	19960\\
16000	20865\\
20000	21994\\
22000	22540\\
24000	22709\\
26000	22769\\
28000	23025\\
30000	23116\\
32000	23105\\
34000	23320\\
38000	23555\\
40000	23562\\
42000	23639\\
44000	23830\\
46000	23949\\
50000	23904\\
52000	23685\\
56000	23598\\
58000	23633\\
60000	23775\\
62000	24014\\
64000	24075\\
68000	23949\\
70000	23965\\
72000	23693\\
74000	23748\\
76000	23715\\
78000	23596\\
80000	23343\\
82000	23335\\
86000	23557\\
88000	23798\\
90000	23845\\
92000	23729\\
96000	23659\\
98000	23596\\
1e+05	23626\\
};
\addplot [color=mycolor19, forget plot]
  table[row sep=crcr]{%
0	4395.3\\
2000	6149.9\\
4000	8737.9\\
6000	10500\\
8000	11582\\
10000	12398\\
12000	12979\\
14000	13227\\
16000	13578\\
18000	13832\\
22000	14104\\
24000	14325\\
26000	14348\\
28000	14327\\
30000	14383\\
32000	14571\\
34000	14672\\
38000	14735\\
42000	14853\\
44000	14905\\
48000	14701\\
52000	14641\\
58000	14729\\
66000	14618\\
68000	14696\\
70000	14709\\
72000	14569\\
74000	14570\\
76000	14680\\
80000	14516\\
84000	14602\\
86000	14476\\
88000	14418\\
90000	14410\\
96000	14598\\
98000	14671\\
1e+05	14562\\
};
\addplot [color=mycolor20, forget plot]
  table[row sep=crcr]{%
0	4395.3\\
2000	6478.4\\
4000	9683.9\\
6000	12001\\
8000	13653\\
10000	14945\\
14000	16460\\
16000	16991\\
18000	17308\\
20000	17733\\
22000	18061\\
24000	18223\\
28000	18246\\
30000	18297\\
32000	18555\\
38000	18679\\
42000	18432\\
44000	18531\\
46000	18464\\
48000	18345\\
50000	18429\\
52000	18667\\
54000	18712\\
56000	18656\\
58000	18660\\
60000	18806\\
62000	18880\\
64000	18780\\
66000	18864\\
70000	18810\\
72000	18877\\
78000	18760\\
80000	18643\\
82000	18680\\
84000	18769\\
90000	18591\\
92000	18650\\
96000	18464\\
1e+05	18427\\
};
\addplot [color=mycolor21, forget plot]
  table[row sep=crcr]{%
0	4395.3\\
2000	5340.6\\
4000	6774.7\\
6000	7694.6\\
8000	8316.5\\
10000	8661.1\\
12000	8990.6\\
14000	9135.3\\
16000	9310.3\\
22000	9465.6\\
24000	9496.4\\
26000	9443.6\\
28000	9418\\
30000	9430.3\\
32000	9508.1\\
34000	9437.1\\
36000	9423.9\\
38000	9468.3\\
40000	9540.2\\
44000	9568.1\\
48000	9524.9\\
52000	9514.7\\
54000	9472.6\\
56000	9461.5\\
62000	9573.9\\
64000	9567.5\\
66000	9485.4\\
68000	9496.5\\
72000	9596.5\\
74000	9555.5\\
80000	9498.5\\
82000	9546.4\\
84000	9538.9\\
88000	9577.8\\
90000	9554.2\\
92000	9641.1\\
94000	9678.1\\
96000	9630.7\\
1e+05	9683.8\\
};
\addplot [color=mycolor22, forget plot]
  table[row sep=crcr]{%
0	4395.3\\
2000	6766.3\\
4000	10607\\
6000	13538\\
8000	15874\\
10000	17622\\
12000	19035\\
14000	20117\\
16000	20994\\
18000	21756\\
20000	22404\\
22000	22872\\
24000	23257\\
26000	23506\\
28000	23845\\
32000	24209\\
34000	24403\\
38000	24592\\
40000	24421\\
46000	24612\\
50000	24392\\
56000	24415\\
58000	24192\\
62000	24246\\
64000	24423\\
66000	24444\\
68000	24289\\
80000	24176\\
84000	24384\\
86000	24315\\
88000	24173\\
90000	24193\\
94000	24434\\
96000	24328\\
98000	24034\\
1e+05	24076\\
};
\addplot [color=mycolor23, forget plot]
  table[row sep=crcr]{%
0	4395.3\\
2000	5718.2\\
4000	7640.9\\
6000	8971.1\\
8000	9860.8\\
10000	10494\\
12000	10971\\
14000	11305\\
16000	11514\\
18000	11680\\
22000	11877\\
24000	11950\\
26000	11953\\
34000	12133\\
38000	12137\\
42000	11996\\
46000	12034\\
48000	12132\\
54000	12191\\
58000	12045\\
60000	11999\\
64000	12113\\
70000	12056\\
74000	12136\\
76000	12123\\
78000	12065\\
82000	12135\\
90000	11989\\
92000	12067\\
98000	12162\\
1e+05	12142\\
};
\addplot [color=mycolor24, forget plot]
  table[row sep=crcr]{%
0	4395.3\\
2000	4562.7\\
4000	4799.2\\
6000	4942.9\\
10000	4997.8\\
12000	5057.6\\
14000	5014.3\\
16000	5038.8\\
18000	5013.1\\
20000	4969.4\\
26000	4948.2\\
28000	4923.6\\
30000	4877.6\\
32000	4859.3\\
34000	4823.8\\
36000	4817\\
38000	4776.6\\
40000	4769.1\\
44000	4733.4\\
46000	4828\\
50000	4862.5\\
52000	4832.2\\
54000	4853.7\\
56000	4890.8\\
62000	4887.1\\
64000	4851.1\\
66000	4837.1\\
70000	4832.8\\
72000	4768.1\\
74000	4727.8\\
76000	4722.5\\
78000	4742.8\\
80000	4818.1\\
84000	4790.8\\
94000	4783.3\\
98000	4744\\
1e+05	4797.6\\
};
\addplot [color=mycolor25, forget plot]
  table[row sep=crcr]{%
0	4395.3\\
2000	5773.8\\
4000	7879.1\\
6000	9223.8\\
8000	10142\\
10000	10851\\
12000	11260\\
14000	11527\\
16000	11759\\
18000	11811\\
20000	11939\\
22000	11939\\
24000	12076\\
28000	12123\\
30000	12246\\
32000	12213\\
34000	12217\\
38000	11986\\
42000	12301\\
44000	12408\\
46000	12386\\
48000	12322\\
50000	12183\\
52000	12109\\
54000	12173\\
56000	12139\\
60000	12183\\
62000	12205\\
64000	12128\\
66000	12232\\
68000	12277\\
70000	12209\\
72000	12199\\
76000	12100\\
78000	12051\\
80000	12155\\
82000	12192\\
84000	12130\\
86000	11956\\
88000	12014\\
94000	12036\\
96000	12013\\
1e+05	12278\\
};
\addplot [color=mycolor26, forget plot]
  table[row sep=crcr]{%
0	4395.3\\
2000	5292\\
4000	6474.1\\
6000	7186.8\\
8000	7606.6\\
10000	7868.2\\
12000	8057.6\\
14000	8207.7\\
16000	8242.2\\
18000	8207\\
22000	8275.5\\
26000	8198\\
28000	8224.9\\
30000	8206.7\\
32000	8262.5\\
34000	8256.4\\
36000	8138.1\\
38000	8113.8\\
40000	8132\\
42000	8187.7\\
44000	8199.3\\
46000	8247.1\\
48000	8204.8\\
50000	8234.4\\
52000	8241.2\\
54000	8136.9\\
58000	8258.2\\
60000	8256.7\\
62000	8164\\
64000	8097.5\\
66000	8203.7\\
68000	8178.2\\
70000	8178.9\\
74000	8314.6\\
80000	8223.4\\
82000	8237.3\\
84000	8355.9\\
88000	8304.3\\
94000	8322.6\\
96000	8343.5\\
98000	8450.2\\
1e+05	8438.6\\
};
\addplot [color=mycolor27, forget plot]
  table[row sep=crcr]{%
0	4395.3\\
2000	5516.3\\
4000	7062.4\\
6000	7997\\
8000	8599.5\\
10000	8984\\
12000	9246.3\\
14000	9470.2\\
20000	9483.5\\
22000	9523.9\\
24000	9601.3\\
26000	9631.8\\
30000	9766.1\\
32000	9662.7\\
34000	9626.8\\
36000	9662.6\\
40000	9807.2\\
42000	9773.7\\
44000	9798.3\\
50000	9804.1\\
52000	9735.4\\
56000	9774\\
58000	9735.4\\
60000	9724.8\\
62000	9785.1\\
66000	9718.2\\
68000	9689.9\\
70000	9579\\
72000	9674\\
74000	9742.2\\
78000	9748.5\\
82000	9558.9\\
86000	9419.6\\
88000	9473.2\\
90000	9498.4\\
92000	9596.4\\
94000	9657.1\\
1e+05	9723.3\\
};
\addplot [color=mycolor28, forget plot]
  table[row sep=crcr]{%
0	4395.3\\
2000	4744.8\\
4000	5186.3\\
6000	5417.2\\
8000	5558.9\\
10000	5606.4\\
14000	5602.4\\
16000	5617\\
18000	5656.2\\
20000	5584.5\\
26000	5547.6\\
28000	5468.6\\
30000	5436.9\\
32000	5441.1\\
34000	5499.7\\
36000	5516.8\\
38000	5476.4\\
40000	5461.2\\
42000	5529.7\\
44000	5564.9\\
46000	5543.1\\
48000	5485.7\\
50000	5446.2\\
52000	5518.2\\
54000	5514.5\\
56000	5535.2\\
58000	5527.1\\
60000	5562.6\\
62000	5521.7\\
64000	5529.5\\
66000	5440\\
68000	5430.4\\
70000	5506\\
72000	5513.8\\
74000	5546.1\\
76000	5537.1\\
78000	5506\\
80000	5561.4\\
82000	5538.8\\
84000	5466.8\\
86000	5479.7\\
90000	5550.3\\
94000	5558\\
98000	5520.7\\
1e+05	5500.9\\
};
\addplot [color=mycolor29, forget plot]
  table[row sep=crcr]{%
0	4395.3\\
2000	5653.3\\
4000	7480.7\\
6000	8737\\
8000	9606.9\\
10000	10187\\
12000	10638\\
14000	10872\\
16000	11012\\
20000	11208\\
22000	11196\\
24000	11230\\
26000	11346\\
28000	11363\\
30000	11413\\
32000	11399\\
34000	11303\\
36000	11348\\
38000	11433\\
40000	11460\\
42000	11424\\
44000	11494\\
46000	11392\\
48000	11444\\
50000	11448\\
52000	11540\\
54000	11522\\
56000	11577\\
60000	11535\\
64000	11583\\
68000	11452\\
70000	11453\\
72000	11502\\
78000	11344\\
80000	11383\\
82000	11373\\
86000	11474\\
88000	11399\\
90000	11417\\
92000	11532\\
94000	11581\\
96000	11557\\
98000	11584\\
1e+05	11515\\
};
\addplot [color=mycolor30, forget plot]
  table[row sep=crcr]{%
0	4395.3\\
2000	5140.5\\
4000	6166.5\\
6000	6763.5\\
8000	7173.5\\
10000	7479.4\\
12000	7658.9\\
14000	7769.3\\
22000	7866.2\\
24000	7915\\
46000	7832.5\\
48000	7905.8\\
54000	7864.6\\
70000	7912.6\\
72000	7868.8\\
76000	7859.3\\
78000	7939.3\\
80000	7883.1\\
84000	7876.4\\
88000	7937.3\\
90000	7887.2\\
92000	7928.3\\
94000	7938.2\\
96000	7894\\
98000	7919.5\\
1e+05	7882.5\\
};
\addplot [color=mycolor31, forget plot]
  table[row sep=crcr]{%
0	4395.3\\
2000	4936.4\\
4000	5637.9\\
6000	6043.9\\
8000	6305.3\\
10000	6429.7\\
12000	6484.8\\
14000	6472.5\\
16000	6504\\
20000	6523.4\\
22000	6595\\
24000	6550.6\\
26000	6591.6\\
28000	6594.6\\
30000	6530.2\\
32000	6522.5\\
34000	6560.4\\
36000	6517.5\\
42000	6578.9\\
44000	6483.2\\
46000	6517.9\\
48000	6523.8\\
50000	6511.4\\
54000	6415.2\\
56000	6412.4\\
58000	6470.1\\
60000	6434.6\\
62000	6380.4\\
68000	6563.7\\
70000	6565.8\\
72000	6619.4\\
74000	6571\\
76000	6574.7\\
80000	6547\\
84000	6382.7\\
94000	6534.6\\
96000	6521.8\\
98000	6449.4\\
1e+05	6523.4\\
};
\addplot [color=mycolor32, forget plot]
  table[row sep=crcr]{%
0	4395.3\\
2000	5520.4\\
4000	7129.4\\
6000	8176.2\\
8000	8871.7\\
10000	9298.6\\
12000	9612.7\\
14000	9981.9\\
18000	10245\\
22000	10421\\
24000	10436\\
28000	10521\\
32000	10495\\
34000	10419\\
38000	10473\\
44000	10322\\
46000	10424\\
50000	10361\\
52000	10399\\
54000	10404\\
58000	10282\\
64000	10403\\
70000	10253\\
76000	10362\\
78000	10321\\
80000	10190\\
82000	10236\\
88000	10169\\
90000	10287\\
92000	10371\\
96000	10334\\
98000	10283\\
1e+05	10283\\
};
\addplot [color=mycolor33, forget plot]
  table[row sep=crcr]{%
0	4395.3\\
2000	5224.4\\
4000	6368.8\\
6000	7036.7\\
8000	7458.8\\
10000	7713\\
12000	7918.2\\
14000	8078.8\\
18000	8133.9\\
20000	8100.6\\
24000	8111.2\\
26000	8163.6\\
30000	8076.7\\
34000	8166.2\\
36000	8103.7\\
38000	8102\\
42000	8132.5\\
44000	8147.7\\
46000	8211.5\\
48000	8215.8\\
50000	8104\\
52000	8094.6\\
54000	8177.2\\
56000	8215.5\\
60000	8136.1\\
62000	8235.9\\
64000	8237.3\\
68000	8079\\
70000	8152.2\\
76000	8143.7\\
78000	8170\\
80000	8135.9\\
84000	8208.5\\
88000	8349.8\\
90000	8272.6\\
94000	8287.2\\
98000	8310\\
1e+05	8227.5\\
};
\addplot [color=mycolor34, forget plot]
  table[row sep=crcr]{%
0	4395.3\\
2000	4412.2\\
6000	4351.7\\
8000	4303.3\\
10000	4279.9\\
12000	4323.5\\
16000	4180\\
18000	4134.2\\
20000	4118.9\\
22000	4136.1\\
24000	4166.5\\
26000	4173.5\\
28000	4204.9\\
30000	4136.8\\
32000	4082.7\\
34000	4103.3\\
36000	4106.6\\
38000	4066.1\\
40000	4084.3\\
42000	4079.1\\
44000	4091.9\\
46000	4090.1\\
48000	4106.1\\
50000	4045.2\\
52000	4048.7\\
54000	4087.1\\
58000	4091.8\\
60000	4176.4\\
62000	4178.2\\
64000	4198\\
66000	4179.6\\
72000	4038.4\\
74000	4078.1\\
76000	4104.2\\
82000	4037.3\\
84000	4055.7\\
86000	4101.2\\
88000	4100.3\\
90000	4084\\
94000	4118.5\\
96000	4076.5\\
1e+05	4085\\
};
\addplot [color=mycolor35]
  table[row sep=crcr]{%
0	4395.3\\
2000	5383.3\\
4000	6839.9\\
6000	7812.3\\
8000	8454.8\\
10000	9030\\
12000	9350.1\\
14000	9604.3\\
16000	9790.8\\
18000	9854.4\\
26000	9981\\
30000	10044\\
32000	9968.5\\
38000	9943\\
42000	10057\\
44000	10008\\
46000	9989.5\\
48000	10064\\
52000	10067\\
54000	10159\\
58000	10113\\
60000	10070\\
64000	10119\\
68000	10050\\
70000	10055\\
72000	10120\\
80000	10187\\
84000	10144\\
86000	10151\\
90000	10060\\
98000	10175\\
1e+05	10128\\
};
\addlegendentry{\ClasskIILegend}

\addplot [color=mycolor36, forget plot]
  table[row sep=crcr]{%
0	4395.3\\
2000	4483.4\\
4000	4558.8\\
6000	4559.3\\
8000	4531.8\\
10000	4476.5\\
12000	4476.1\\
14000	4493.7\\
16000	4493.4\\
18000	4480.7\\
20000	4423.1\\
22000	4407.7\\
26000	4328.4\\
28000	4329.8\\
30000	4364.7\\
32000	4373.5\\
34000	4335.6\\
38000	4372.8\\
40000	4337.1\\
42000	4352.6\\
44000	4326.2\\
48000	4327.9\\
50000	4353.7\\
52000	4332.8\\
54000	4360.9\\
56000	4308.5\\
58000	4291.2\\
60000	4309.9\\
62000	4345.7\\
64000	4417.5\\
66000	4420.5\\
68000	4392\\
70000	4394.4\\
72000	4374.3\\
74000	4391.2\\
78000	4400.2\\
80000	4419.9\\
82000	4460.2\\
84000	4428\\
86000	4444.5\\
88000	4446.4\\
90000	4335.5\\
92000	4292.3\\
94000	4314.4\\
1e+05	4287.9\\
};
\addplot [color=mycolor37, forget plot]
  table[row sep=crcr]{%
0	4395.3\\
2000	4267.7\\
4000	4127.8\\
6000	4088.2\\
8000	4034.7\\
10000	3963.7\\
14000	3913.1\\
16000	3860.3\\
18000	3842.7\\
20000	3805.7\\
22000	3812.7\\
24000	3779.1\\
26000	3764.2\\
28000	3708.2\\
30000	3728.8\\
40000	3713.8\\
42000	3737.5\\
50000	3727.5\\
52000	3751.9\\
54000	3744\\
56000	3715.1\\
58000	3752.3\\
60000	3773.9\\
62000	3769.6\\
64000	3751.5\\
66000	3755.7\\
68000	3733.3\\
70000	3692.7\\
72000	3691.8\\
74000	3729.6\\
76000	3696.9\\
78000	3689.6\\
82000	3718.9\\
84000	3702.9\\
86000	3727.3\\
88000	3696.4\\
90000	3700.9\\
92000	3725.1\\
98000	3726.2\\
1e+05	3747.1\\
};
\addplot [color=mycolor38, forget plot]
  table[row sep=crcr]{%
0	4395.3\\
2000	4120.6\\
4000	3799.9\\
6000	3570.8\\
8000	3423.2\\
10000	3379.1\\
12000	3315\\
14000	3292\\
16000	3252.5\\
18000	3229.6\\
20000	3156.2\\
24000	3121.5\\
26000	3084.6\\
28000	3127.8\\
30000	3114.7\\
32000	3124.9\\
34000	3122.9\\
36000	3101.8\\
38000	3058.7\\
40000	3049.9\\
42000	3089.2\\
44000	3074.6\\
46000	3078.6\\
48000	3039.4\\
52000	3098.9\\
54000	3127.8\\
56000	3133.3\\
58000	3093.4\\
60000	3122.6\\
62000	3114.9\\
64000	3078.9\\
66000	3081.6\\
68000	3093.3\\
70000	3123.5\\
72000	3095.9\\
76000	3140.5\\
78000	3144.5\\
80000	3161.9\\
82000	3145.1\\
84000	3083.5\\
86000	3102.1\\
90000	3123.5\\
92000	3113.8\\
94000	3086.5\\
96000	3100.6\\
98000	3129.1\\
1e+05	3116\\
};
\addplot [color=mycolor39, forget plot]
  table[row sep=crcr]{%
0	4395.3\\
2000	5487.5\\
4000	7147.7\\
6000	8328.7\\
8000	9274.8\\
10000	9873.3\\
12000	10274\\
14000	10542\\
16000	10851\\
18000	11060\\
20000	11232\\
24000	11411\\
26000	11329\\
28000	11283\\
32000	11324\\
34000	11326\\
36000	11384\\
40000	11616\\
42000	11646\\
46000	11578\\
48000	11639\\
50000	11568\\
54000	11477\\
56000	11544\\
58000	11577\\
60000	11660\\
66000	11675\\
78000	11360\\
80000	11515\\
84000	11618\\
86000	11623\\
88000	11701\\
90000	11645\\
92000	11722\\
96000	11574\\
98000	11424\\
1e+05	11510\\
};
\addplot [color=mycolor40, forget plot]
  table[row sep=crcr]{%
0	4395.3\\
2000	4655.4\\
4000	5041.3\\
6000	5358.1\\
8000	5551.2\\
10000	5705.1\\
14000	5759.1\\
16000	5808.7\\
18000	5836.7\\
20000	5805.8\\
24000	5828.6\\
28000	5761.7\\
34000	5788.6\\
38000	5758.5\\
40000	5819.1\\
42000	5846.5\\
44000	5808.8\\
46000	5732.4\\
48000	5764.9\\
54000	5777.1\\
56000	5795.9\\
62000	5782.8\\
64000	5752.2\\
66000	5737.8\\
70000	5778.7\\
72000	5791\\
78000	5751.5\\
80000	5770.1\\
84000	5768.7\\
86000	5783.4\\
90000	5854.1\\
94000	5824.6\\
96000	5794.8\\
98000	5783.1\\
1e+05	5744.2\\
};
\addplot [color=mycolor41, forget plot]
  table[row sep=crcr]{%
0	4395.3\\
2000	4539.3\\
4000	4756.3\\
6000	4966.1\\
8000	5131\\
10000	5242.6\\
14000	5387.8\\
16000	5388.7\\
18000	5428\\
22000	5421\\
24000	5406.1\\
26000	5433.9\\
28000	5437.3\\
30000	5404.5\\
36000	5425.6\\
40000	5416.2\\
42000	5452\\
50000	5396.7\\
52000	5411.7\\
60000	5406.5\\
64000	5446.4\\
66000	5476.9\\
68000	5448\\
70000	5481.1\\
72000	5433.8\\
74000	5409.6\\
76000	5406.9\\
78000	5437.7\\
88000	5464.7\\
90000	5421.6\\
96000	5431.8\\
98000	5395.5\\
1e+05	5421\\
};
\addplot [color=mycolor42, forget plot]
  table[row sep=crcr]{%
0	4395.3\\
2000	5990.6\\
4000	8627.8\\
6000	10671\\
8000	12468\\
10000	13925\\
12000	15085\\
14000	16049\\
16000	16696\\
18000	17297\\
20000	17801\\
22000	18208\\
26000	18897\\
28000	19063\\
30000	18927\\
38000	19422\\
42000	19708\\
44000	19783\\
46000	19737\\
50000	19878\\
58000	19818\\
68000	19914\\
70000	19797\\
72000	19959\\
80000	20076\\
82000	20075\\
86000	19957\\
90000	20078\\
92000	19954\\
96000	20085\\
98000	20164\\
1e+05	20099\\
};
\addplot [color=mycolor43, forget plot]
  table[row sep=crcr]{%
0	4395.3\\
2000	4557.9\\
4000	4780.2\\
6000	4925.2\\
8000	5029.1\\
12000	5090.4\\
16000	5066.8\\
18000	5032\\
20000	5035.1\\
22000	5055.1\\
26000	5038.1\\
28000	5025.4\\
30000	5041.6\\
34000	4987.3\\
36000	5007.7\\
38000	5069.1\\
42000	4996.3\\
46000	5003.9\\
48000	4965.3\\
50000	4966.4\\
52000	4949\\
54000	4966.7\\
56000	5012.7\\
58000	5019.1\\
60000	4992.8\\
62000	5001.7\\
64000	4970.1\\
66000	4958.7\\
68000	4965.3\\
70000	4994.1\\
72000	4994.2\\
74000	4967.2\\
76000	5009.8\\
78000	4994.2\\
80000	5021.6\\
82000	5017.4\\
86000	4971\\
88000	5022.3\\
90000	5058.8\\
92000	5029\\
94000	5039\\
98000	5001.6\\
1e+05	5032.4\\
};
\addplot [color=mycolor44, forget plot]
  table[row sep=crcr]{%
0	4395.3\\
2000	4279.8\\
4000	4026.1\\
6000	3835.3\\
8000	3733\\
10000	3689.4\\
12000	3636.1\\
14000	3571.7\\
16000	3541\\
18000	3531.1\\
22000	3459.7\\
24000	3468.6\\
28000	3453.3\\
30000	3422.3\\
32000	3448.3\\
34000	3502.8\\
36000	3481.8\\
38000	3480.7\\
40000	3468.5\\
42000	3427.2\\
44000	3401.9\\
46000	3450.3\\
48000	3442.7\\
50000	3420.1\\
52000	3418.8\\
54000	3445.5\\
56000	3489.4\\
58000	3470.3\\
62000	3415.7\\
64000	3414.3\\
66000	3425.6\\
68000	3401.9\\
70000	3447.6\\
72000	3417.1\\
76000	3417.2\\
78000	3443.7\\
80000	3422\\
82000	3383.1\\
84000	3392.8\\
86000	3384.3\\
88000	3401.8\\
92000	3471.6\\
94000	3448.4\\
96000	3452.4\\
1e+05	3492\\
};
\addplot [color=mycolor45, forget plot]
  table[row sep=crcr]{%
0	4395.3\\
2000	4890.1\\
4000	5613.7\\
6000	6105\\
8000	6415.6\\
10000	6622.8\\
12000	6755.4\\
14000	6848.5\\
16000	6896\\
20000	6929.1\\
26000	6931.3\\
28000	6952.7\\
30000	6934.1\\
32000	6946.4\\
40000	6912.7\\
42000	6884.1\\
44000	6885.7\\
46000	6932.8\\
48000	6920.4\\
50000	6973\\
54000	6950.7\\
56000	6962.6\\
58000	7034.8\\
60000	6986.6\\
62000	6970.4\\
64000	7007.3\\
66000	7011.6\\
68000	6969.8\\
70000	6983.9\\
74000	6917.1\\
76000	6951.5\\
78000	6915\\
80000	6932.3\\
82000	6969.3\\
92000	6897.6\\
94000	6845.3\\
98000	6878.3\\
1e+05	6953.7\\
};
\addplot [color=mycolor46, forget plot]
  table[row sep=crcr]{%
0	4395.3\\
2000	4337.4\\
4000	4296\\
6000	4312.5\\
8000	4313.5\\
12000	4352.9\\
14000	4310.1\\
22000	4219.6\\
24000	4243.5\\
28000	4227.1\\
30000	4218.8\\
32000	4144.4\\
34000	4115.6\\
36000	4151.3\\
38000	4130.5\\
40000	4153.3\\
44000	4158.3\\
46000	4152.4\\
48000	4109.7\\
50000	4119.7\\
52000	4144.1\\
54000	4156.2\\
56000	4203.1\\
58000	4190.9\\
64000	4106.4\\
68000	4148\\
70000	4130.6\\
72000	4132.4\\
74000	4109.6\\
80000	4174.6\\
82000	4188.8\\
88000	4120.7\\
90000	4130.5\\
94000	4082\\
96000	4125\\
98000	4116.6\\
1e+05	4126.4\\
};
\addplot [color=mycolor47, forget plot]
  table[row sep=crcr]{%
0	4395.3\\
2000	4563.9\\
4000	4781\\
6000	4950.8\\
8000	5084.1\\
10000	5192.8\\
12000	5270\\
14000	5305.3\\
18000	5326.3\\
28000	5335.8\\
36000	5275.5\\
40000	5353.9\\
42000	5365.6\\
44000	5332.2\\
46000	5326.8\\
48000	5287.9\\
50000	5311.9\\
52000	5312.1\\
56000	5364.7\\
58000	5368.2\\
60000	5342.8\\
62000	5353.4\\
64000	5323.8\\
68000	5298.1\\
70000	5259.2\\
72000	5273.7\\
74000	5258.6\\
80000	5269.9\\
84000	5250.6\\
86000	5291\\
88000	5273.1\\
1e+05	5345.7\\
};
\addplot [color=mycolor48, forget plot]
  table[row sep=crcr]{%
0	4395.3\\
2000	4116\\
4000	3756.6\\
6000	3533.4\\
8000	3389.5\\
10000	3286.3\\
14000	3164.6\\
16000	3121.1\\
18000	3113.4\\
20000	3090.6\\
26000	3069.9\\
28000	3042.7\\
30000	3026.3\\
32000	3036.1\\
34000	3022\\
36000	3023.6\\
38000	3012.5\\
40000	3029.3\\
44000	3006.5\\
46000	2994.7\\
52000	3023.7\\
56000	3033\\
58000	3042.2\\
60000	3030\\
62000	2997.9\\
64000	2989\\
70000	3008.1\\
74000	3012.3\\
76000	3022.9\\
78000	3020.3\\
80000	3001.7\\
82000	2993.8\\
86000	3003.4\\
88000	3003.8\\
90000	2994\\
92000	2995.8\\
94000	2987.2\\
98000	2996.4\\
1e+05	2985.9\\
};
\addplot [color=mycolor49, forget plot]
  table[row sep=crcr]{%
0	4395.3\\
4000	4556.8\\
8000	4623.9\\
10000	4607.2\\
12000	4548.5\\
14000	4549.2\\
16000	4573.8\\
18000	4557.8\\
20000	4489.8\\
26000	4481.3\\
28000	4486.5\\
30000	4450.1\\
32000	4436.5\\
38000	4446.9\\
40000	4479.8\\
42000	4444.4\\
44000	4460.9\\
46000	4453.3\\
48000	4506.5\\
50000	4503\\
52000	4439.1\\
54000	4411.1\\
56000	4368.4\\
58000	4347.2\\
62000	4384.9\\
64000	4381.2\\
66000	4390.4\\
68000	4437.9\\
70000	4424.2\\
74000	4425.5\\
76000	4415.9\\
78000	4454.8\\
80000	4472.2\\
84000	4475.3\\
86000	4439.5\\
90000	4406.9\\
94000	4492.4\\
96000	4513.3\\
98000	4500.3\\
1e+05	4425.3\\
};
\addplot [color=mycolor50, forget plot]
  table[row sep=crcr]{%
0	4395.3\\
2000	3991.4\\
4000	3473.8\\
6000	3181.6\\
8000	3010.6\\
10000	2878.2\\
14000	2731.6\\
18000	2678.5\\
20000	2663.7\\
22000	2622.1\\
26000	2593.5\\
32000	2588.3\\
34000	2584.7\\
36000	2558.1\\
40000	2554.7\\
44000	2527.9\\
46000	2521.2\\
50000	2544.7\\
54000	2577.6\\
56000	2571.8\\
58000	2584.9\\
62000	2563.5\\
64000	2573.2\\
66000	2554.8\\
68000	2545.4\\
70000	2568.7\\
74000	2528.2\\
76000	2520.4\\
78000	2552.1\\
80000	2560.9\\
82000	2559.4\\
84000	2570.1\\
86000	2566.7\\
88000	2535.1\\
90000	2544.4\\
92000	2566.3\\
94000	2545.1\\
98000	2562.4\\
1e+05	2543.9\\
};
\addplot [color=mycolor51, forget plot]
  table[row sep=crcr]{%
0	4395.3\\
2000	4705.9\\
4000	5190\\
6000	5526.1\\
8000	5725.2\\
10000	5882.9\\
12000	6004\\
16000	5996.6\\
18000	6029.1\\
24000	5976.3\\
26000	5991.4\\
28000	6078.3\\
30000	6106.4\\
32000	6172\\
34000	6217.3\\
36000	6106.1\\
40000	6044.3\\
42000	5982.4\\
44000	5990.3\\
46000	5946.9\\
50000	5954.8\\
52000	5937.6\\
54000	5971.1\\
56000	5956.2\\
60000	5969.8\\
62000	5954.8\\
66000	6019.3\\
68000	5947.6\\
70000	6008.7\\
74000	6007.7\\
78000	5978.2\\
80000	5960.9\\
82000	5970.4\\
84000	5905.1\\
86000	5934.5\\
90000	6028.3\\
92000	6022.2\\
94000	6037.2\\
96000	5966\\
1e+05	6033.1\\
};
\addplot [color=mycolor52, forget plot]
  table[row sep=crcr]{%
0	4395.3\\
2000	4297.9\\
4000	4129.1\\
6000	4009.2\\
8000	3924.1\\
12000	3825.2\\
14000	3804.2\\
16000	3754.7\\
20000	3633.8\\
24000	3691.6\\
26000	3669.6\\
28000	3633.2\\
30000	3608.1\\
32000	3621.1\\
34000	3587.6\\
36000	3584.4\\
38000	3610\\
40000	3617.1\\
42000	3636.6\\
44000	3600.3\\
48000	3614.3\\
50000	3580\\
52000	3596.6\\
54000	3602\\
58000	3590.9\\
62000	3634.4\\
64000	3630\\
66000	3576.6\\
68000	3588.7\\
70000	3624.2\\
72000	3612.3\\
74000	3633.1\\
78000	3589.8\\
80000	3587.1\\
84000	3612.9\\
86000	3625.6\\
88000	3665.3\\
92000	3638.2\\
94000	3610.4\\
96000	3610.7\\
98000	3625.3\\
1e+05	3594.3\\
};
\addplot [color=mycolor53, forget plot]
  table[row sep=crcr]{%
0	4395.3\\
2000	4670.4\\
4000	5090\\
6000	5430\\
8000	5599.1\\
10000	5681.1\\
16000	5846.3\\
18000	5794.3\\
20000	5790.1\\
22000	5747.6\\
26000	5730.4\\
28000	5750.9\\
32000	5718.5\\
34000	5760.9\\
36000	5701.1\\
44000	5643.8\\
48000	5681.8\\
50000	5690.8\\
54000	5639.3\\
56000	5662.3\\
58000	5631.4\\
60000	5645.9\\
64000	5741.7\\
66000	5745\\
68000	5775.3\\
70000	5749.6\\
74000	5773.6\\
78000	5695.9\\
80000	5735.5\\
82000	5710.4\\
84000	5706.7\\
86000	5732\\
88000	5708.8\\
90000	5736.3\\
94000	5695.6\\
98000	5702.8\\
1e+05	5659.9\\
};
\addplot [color=mycolor54, forget plot]
  table[row sep=crcr]{%
0	4395.3\\
2000	4418.4\\
4000	4474.8\\
6000	4555.3\\
10000	4629.3\\
16000	4698.4\\
18000	4649.2\\
22000	4651.5\\
24000	4647.3\\
30000	4686.4\\
34000	4687.4\\
36000	4694.1\\
38000	4663.5\\
40000	4663.5\\
42000	4640.8\\
44000	4633.5\\
46000	4643.8\\
48000	4672\\
50000	4645\\
54000	4647.8\\
56000	4628.8\\
62000	4642.5\\
64000	4653.5\\
66000	4632.5\\
72000	4653.2\\
74000	4684.1\\
80000	4633.5\\
82000	4643.8\\
86000	4630.5\\
88000	4624.5\\
90000	4634.2\\
92000	4624.6\\
96000	4629.8\\
1e+05	4581.8\\
};
\addplot [color=mycolor55, forget plot]
  table[row sep=crcr]{%
0	4395.3\\
2000	4633.7\\
4000	5003.8\\
6000	5246.8\\
8000	5409.5\\
10000	5543.6\\
12000	5649.1\\
14000	5727.3\\
18000	5774.3\\
22000	5769.4\\
30000	5807.5\\
40000	5750.4\\
50000	5737\\
52000	5764.4\\
60000	5751.2\\
62000	5753.8\\
66000	5823.9\\
70000	5827.5\\
72000	5819\\
74000	5786\\
76000	5775.9\\
80000	5833.1\\
82000	5829.1\\
84000	5854.9\\
88000	5843.7\\
94000	5865.2\\
98000	5884.1\\
1e+05	5889.3\\
};
\addplot [color=mycolor56, forget plot]
  table[row sep=crcr]{%
0	4395.3\\
4000	4320.8\\
6000	4307.2\\
10000	4248.3\\
14000	4193.9\\
16000	4166.4\\
18000	4153.8\\
20000	4172.9\\
26000	4165\\
28000	4111.2\\
30000	4081.5\\
32000	4094.1\\
38000	4029.9\\
42000	4058.6\\
44000	4042.8\\
46000	4040.2\\
54000	4092\\
56000	4074.1\\
58000	4069.9\\
62000	4034.9\\
64000	4070.2\\
70000	4095.6\\
74000	4086\\
78000	4066.8\\
82000	4078.3\\
84000	4075.6\\
86000	4048.2\\
90000	4048.7\\
92000	4055.3\\
94000	4043.4\\
98000	4054.8\\
1e+05	4029.7\\
};
\addplot [color=mycolor57, forget plot]
  table[row sep=crcr]{%
0	4395.3\\
2000	4648.8\\
4000	5089.1\\
6000	5433.2\\
8000	5673.8\\
10000	5876.4\\
14000	6199.8\\
16000	6305.3\\
20000	6443\\
24000	6534.1\\
34000	6662.5\\
36000	6659.6\\
40000	6684.6\\
44000	6667.9\\
46000	6680.6\\
48000	6671.3\\
52000	6685.6\\
56000	6699.6\\
58000	6715.8\\
60000	6698.4\\
64000	6698.7\\
66000	6662.7\\
70000	6669.2\\
72000	6701.1\\
76000	6674.2\\
82000	6710.4\\
92000	6687.3\\
1e+05	6711.9\\
};
\addplot [color=mycolor58, forget plot]
  table[row sep=crcr]{%
0	4395.3\\
2000	4308.2\\
4000	4185.9\\
6000	4093.2\\
8000	4014.7\\
10000	3962.3\\
14000	3905.8\\
18000	3796.8\\
22000	3725.4\\
24000	3708.6\\
26000	3704\\
28000	3668\\
30000	3643\\
32000	3687\\
34000	3675.1\\
38000	3676.7\\
40000	3661.7\\
44000	3675.8\\
46000	3641.7\\
48000	3656.4\\
50000	3657\\
52000	3689.1\\
58000	3670.5\\
60000	3651.6\\
62000	3659.6\\
64000	3657.5\\
66000	3629.4\\
68000	3638.5\\
70000	3666.9\\
72000	3683\\
74000	3673.5\\
76000	3645.6\\
78000	3597.4\\
82000	3583.2\\
84000	3624.4\\
86000	3642.7\\
90000	3637.6\\
92000	3641.8\\
94000	3667\\
96000	3650.8\\
1e+05	3679.9\\
};
\addplot [color=mycolor59, forget plot]
  table[row sep=crcr]{%
0	4395.3\\
2000	4511\\
4000	4710.3\\
6000	4839.6\\
8000	4933.3\\
12000	5059\\
14000	5066.8\\
16000	5090.4\\
24000	5104.2\\
26000	5095.7\\
28000	5067.5\\
30000	5075.6\\
32000	5065.5\\
34000	5071.5\\
36000	5103.6\\
38000	5109.1\\
46000	5058.3\\
50000	5102.5\\
52000	5107.7\\
54000	5089.6\\
58000	5112.5\\
60000	5107.9\\
62000	5086.4\\
64000	5086\\
66000	5105.1\\
68000	5099.7\\
70000	5113\\
74000	5108.3\\
76000	5122.4\\
78000	5109.1\\
80000	5131.2\\
82000	5127.2\\
84000	5102.4\\
88000	5104\\
92000	5123.2\\
94000	5126.7\\
96000	5087\\
1e+05	5062.2\\
};
\addplot [color=mycolor60, forget plot]
  table[row sep=crcr]{%
0	4395.3\\
2000	4220\\
4000	3983.9\\
6000	3887.1\\
8000	3818.1\\
10000	3768.8\\
16000	3681.2\\
26000	3622.9\\
28000	3617\\
32000	3630.5\\
36000	3612.5\\
38000	3615.8\\
40000	3604.5\\
42000	3618.8\\
46000	3610.2\\
48000	3598.4\\
50000	3599.3\\
54000	3616.2\\
56000	3612.5\\
58000	3633.1\\
62000	3647.8\\
64000	3623.9\\
66000	3615.5\\
68000	3593.3\\
74000	3617.9\\
76000	3605.4\\
78000	3614.5\\
80000	3600\\
88000	3611\\
90000	3613.5\\
96000	3582.2\\
1e+05	3568.3\\
};
\addplot [color=mycolor61, forget plot]
  table[row sep=crcr]{%
0	4395.3\\
2000	4207.3\\
4000	3955\\
6000	3779.9\\
8000	3656.4\\
10000	3557.1\\
12000	3482.1\\
18000	3358.5\\
24000	3295.6\\
26000	3290.1\\
32000	3227.5\\
36000	3234.2\\
40000	3214.3\\
46000	3204.7\\
50000	3222.5\\
54000	3212.5\\
56000	3213.9\\
60000	3195.1\\
62000	3191.4\\
64000	3197.8\\
66000	3218.2\\
68000	3223.8\\
70000	3205\\
72000	3200.6\\
78000	3223.9\\
80000	3204.2\\
82000	3204.4\\
84000	3194.6\\
88000	3196.7\\
90000	3180.3\\
94000	3193.9\\
98000	3199.8\\
1e+05	3204.9\\
};
\addplot [color=mycolor62, forget plot]
  table[row sep=crcr]{%
0	4395.3\\
2000	4421.2\\
4000	4524.3\\
6000	4645.8\\
8000	4737.4\\
14000	4928.9\\
16000	4938.6\\
18000	4969.4\\
22000	4984.2\\
34000	5020.5\\
38000	5025.7\\
44000	4996.4\\
48000	4985.7\\
54000	5025.3\\
58000	5026.6\\
60000	5020.9\\
62000	5035.9\\
66000	5008.7\\
68000	5007.8\\
74000	5034.1\\
76000	5043.6\\
78000	5024.8\\
84000	5034.5\\
86000	5011.8\\
94000	5041.7\\
98000	5028\\
1e+05	5036.9\\
};
\addplot [color=mycolor63, forget plot]
  table[row sep=crcr]{%
0	4395.3\\
2000	4025.3\\
4000	3532.3\\
6000	3212.6\\
8000	2994.4\\
10000	2873.1\\
12000	2790.2\\
14000	2721.3\\
16000	2676.5\\
20000	2615.1\\
22000	2582.1\\
24000	2563.9\\
28000	2547.1\\
32000	2506.4\\
36000	2504.7\\
38000	2497.9\\
40000	2475.7\\
42000	2476\\
44000	2486.8\\
46000	2487.3\\
52000	2468\\
54000	2475.4\\
58000	2470.8\\
60000	2480.1\\
64000	2472.2\\
74000	2493.3\\
78000	2485.2\\
80000	2492\\
82000	2487.8\\
84000	2472.8\\
88000	2496.8\\
90000	2486.6\\
96000	2486.8\\
98000	2492.6\\
1e+05	2491.3\\
};
\addplot [color=mycolor64, forget plot]
  table[row sep=crcr]{%
0	4395.3\\
2000	4134.5\\
4000	3798.9\\
6000	3599.6\\
8000	3456.5\\
10000	3370.3\\
14000	3290.3\\
16000	3262.5\\
18000	3250.9\\
24000	3182.8\\
26000	3158.6\\
28000	3159.8\\
32000	3140.2\\
34000	3143.6\\
36000	3132.5\\
38000	3133.7\\
40000	3124.1\\
44000	3133.3\\
46000	3140.4\\
50000	3124.7\\
58000	3117.5\\
62000	3085.7\\
66000	3108.4\\
68000	3102.6\\
70000	3116.7\\
74000	3114.8\\
76000	3119.2\\
80000	3105.7\\
82000	3104.7\\
84000	3121.6\\
88000	3120.6\\
92000	3103.6\\
94000	3114.4\\
98000	3105.3\\
1e+05	3117.9\\
};
\addplot [color=mycolor65, forget plot]
  table[row sep=crcr]{%
0	4395.3\\
2000	4051.4\\
4000	3579.4\\
6000	3294\\
8000	3111\\
10000	3000.8\\
12000	2910.2\\
14000	2853.6\\
16000	2819.4\\
18000	2795.7\\
20000	2764.4\\
32000	2652.7\\
36000	2659.9\\
42000	2644.7\\
44000	2620.1\\
46000	2622.8\\
52000	2605.4\\
58000	2604.1\\
60000	2620.8\\
64000	2633.7\\
72000	2628\\
76000	2637.8\\
86000	2609.6\\
90000	2624.9\\
94000	2620.4\\
96000	2637.3\\
1e+05	2637.4\\
};
\addplot [color=mycolor66, forget plot]
  table[row sep=crcr]{%
0	4395.3\\
2000	4224.8\\
4000	3984.2\\
6000	3832\\
8000	3734.9\\
12000	3624.3\\
16000	3570.8\\
18000	3549.2\\
24000	3508.8\\
28000	3521.4\\
32000	3502.6\\
46000	3497.8\\
48000	3490.3\\
52000	3493.8\\
54000	3492.8\\
58000	3505\\
64000	3491.4\\
68000	3494.7\\
78000	3488.4\\
80000	3480.1\\
84000	3491.8\\
88000	3479.4\\
90000	3480.1\\
92000	3490.9\\
96000	3466.1\\
1e+05	3463.7\\
};
\addplot [color=mycolor67, forget plot]
  table[row sep=crcr]{%
0	4395.3\\
2000	4209\\
4000	3964.6\\
6000	3804.3\\
8000	3696.1\\
10000	3602.4\\
14000	3478.4\\
16000	3437.2\\
20000	3394.6\\
24000	3347.6\\
26000	3341.4\\
28000	3311\\
30000	3303.1\\
32000	3309.3\\
40000	3279\\
42000	3263.2\\
46000	3256.2\\
48000	3258.7\\
50000	3245.5\\
52000	3261.3\\
56000	3255.9\\
58000	3245.4\\
60000	3221.4\\
62000	3208\\
66000	3229.7\\
68000	3249.1\\
70000	3250\\
72000	3271\\
76000	3245.5\\
80000	3256.5\\
84000	3246.4\\
88000	3228.9\\
90000	3229.9\\
92000	3220.8\\
94000	3230.6\\
96000	3250.3\\
1e+05	3239.9\\
};
\addplot [color=mycolor67]
  table[row sep=crcr]{%
0	4395.3\\
2000	4082.6\\
4000	3650.7\\
6000	3370\\
8000	3195.1\\
10000	3076.4\\
12000	2988.4\\
14000	2915.2\\
18000	2829.4\\
20000	2808\\
24000	2781.3\\
28000	2752.6\\
40000	2709.4\\
60000	2678.1\\
66000	2673\\
68000	2688.9\\
74000	2701.5\\
76000	2698.2\\
78000	2717.1\\
84000	2693\\
88000	2716.3\\
98000	2691.8\\
1e+05	2689\\
};
\addlegendentry{\ClasskIIILegend}

\addplot [color=mycolor68, forget plot]
  table[row sep=crcr]{%
0	4395.3\\
2000	4102.2\\
4000	3659.9\\
6000	3358.9\\
8000	3155.4\\
10000	3006.8\\
12000	2894.9\\
14000	2817.6\\
16000	2760.7\\
18000	2713.3\\
24000	2621\\
26000	2601.5\\
30000	2587.6\\
34000	2584.4\\
36000	2566.8\\
42000	2552.8\\
46000	2559.8\\
48000	2556.6\\
50000	2563.5\\
54000	2554\\
62000	2550.8\\
68000	2530.6\\
70000	2528.8\\
72000	2535\\
76000	2531.4\\
82000	2540.4\\
86000	2549.3\\
88000	2545.8\\
96000	2562.3\\
98000	2561.2\\
1e+05	2551.6\\
};
\addplot [color=mycolor69, forget plot]
  table[row sep=crcr]{%
0	4395.3\\
2000	4176.5\\
4000	3853.4\\
6000	3642\\
8000	3491.2\\
10000	3383.7\\
12000	3306.3\\
14000	3243.5\\
18000	3162.1\\
24000	3054.6\\
30000	3050.7\\
32000	3028.5\\
42000	2993.5\\
46000	2993.2\\
50000	2993.3\\
58000	2973\\
62000	2977.3\\
64000	2975.4\\
68000	2997.3\\
72000	2977\\
74000	2981.2\\
80000	2959.3\\
86000	2986\\
90000	2976.7\\
92000	2987.8\\
96000	2986\\
1e+05	2974.1\\
};
\addplot [color=mycolor70, forget plot]
  table[row sep=crcr]{%
0	4395.3\\
2000	4022.7\\
4000	3494.8\\
6000	3136.4\\
8000	2893\\
10000	2714.6\\
12000	2588.5\\
14000	2494.2\\
16000	2428.1\\
18000	2375.6\\
22000	2309.8\\
24000	2274.8\\
30000	2221.3\\
34000	2217.9\\
38000	2210.7\\
40000	2213.3\\
42000	2207.2\\
46000	2185.1\\
52000	2195.4\\
54000	2184.1\\
56000	2184.2\\
58000	2190.6\\
62000	2187\\
64000	2193.4\\
70000	2188.6\\
72000	2210.2\\
74000	2194.5\\
76000	2186.2\\
80000	2187.2\\
82000	2176.8\\
84000	2173.3\\
88000	2179.5\\
94000	2188.7\\
96000	2182.5\\
1e+05	2192.8\\
};
\addplot [color=mycolor71, forget plot]
  table[row sep=crcr]{%
0	4395.3\\
2000	4080.6\\
4000	3629.9\\
6000	3335.1\\
8000	3135.6\\
10000	2992.8\\
12000	2898.6\\
14000	2821.9\\
16000	2763.9\\
18000	2719.2\\
20000	2687.9\\
26000	2627\\
34000	2572.2\\
46000	2567\\
52000	2558\\
56000	2571.5\\
64000	2548.3\\
68000	2544.1\\
74000	2561.7\\
80000	2573.7\\
84000	2577.5\\
90000	2557.5\\
94000	2561.8\\
98000	2562.7\\
1e+05	2561\\
};
\addplot [color=mycolor72, forget plot]
  table[row sep=crcr]{%
0	4395.3\\
2000	4046.3\\
4000	3548.4\\
6000	3217.6\\
8000	2982.3\\
10000	2826.6\\
12000	2702.2\\
14000	2611.5\\
16000	2534.1\\
18000	2482.8\\
20000	2440.2\\
24000	2396.5\\
28000	2367.7\\
32000	2325.6\\
36000	2314.4\\
40000	2313.8\\
48000	2288.8\\
50000	2288.7\\
52000	2273.1\\
54000	2263.9\\
56000	2266.5\\
58000	2260\\
60000	2278.5\\
62000	2275.6\\
64000	2263.3\\
66000	2271.1\\
68000	2287.4\\
72000	2277.6\\
78000	2259.7\\
90000	2263.3\\
94000	2265.2\\
96000	2266.1\\
98000	2274.1\\
1e+05	2272.1\\
};
\addplot [color=mycolor73, forget plot]
  table[row sep=crcr]{%
0	4395.3\\
2000	4113.7\\
4000	3705\\
6000	3419.6\\
8000	3228.4\\
10000	3093.7\\
12000	2999.9\\
14000	2922.2\\
16000	2866.5\\
18000	2830.3\\
24000	2745.8\\
34000	2701.7\\
36000	2704\\
42000	2688.3\\
66000	2688\\
70000	2679.1\\
74000	2677.9\\
78000	2687.8\\
84000	2678.5\\
86000	2669.6\\
92000	2672.1\\
96000	2670.8\\
1e+05	2679.7\\
};
\addplot [color=mycolor74, forget plot]
  table[row sep=crcr]{%
0	4395.3\\
2000	3928.9\\
4000	3246.8\\
6000	2790.8\\
8000	2482.3\\
10000	2247.3\\
12000	2083.1\\
14000	1969.7\\
16000	1894.2\\
18000	1831\\
20000	1776.8\\
22000	1739.7\\
26000	1691.4\\
32000	1644.6\\
34000	1637\\
38000	1641.2\\
42000	1636.7\\
44000	1630.9\\
50000	1626.8\\
52000	1615.4\\
54000	1609.6\\
66000	1619.2\\
70000	1620.6\\
74000	1621.4\\
76000	1613.7\\
80000	1614\\
84000	1611.2\\
92000	1608.8\\
1e+05	1617.7\\
};
\addplot [color=mycolor75, forget plot]
  table[row sep=crcr]{%
0	4395.3\\
2000	4092.4\\
4000	3634.4\\
6000	3317.9\\
8000	3083.1\\
10000	2905.7\\
12000	2757.1\\
14000	2658.4\\
16000	2580.8\\
20000	2464.4\\
24000	2403.6\\
28000	2366.1\\
34000	2313.6\\
46000	2290.4\\
50000	2293.6\\
52000	2287.8\\
56000	2291\\
76000	2275.1\\
84000	2285\\
92000	2292.6\\
96000	2298.1\\
1e+05	2294.8\\
};
\addplot [color=mycolor76, forget plot]
  table[row sep=crcr]{%
0	4395.3\\
2000	4056.8\\
4000	3538\\
6000	3163.7\\
8000	2905.4\\
10000	2711.1\\
12000	2565.7\\
14000	2458.6\\
16000	2378.5\\
18000	2322.6\\
22000	2245.1\\
26000	2189.6\\
30000	2155.3\\
34000	2142.4\\
38000	2117\\
42000	2102.4\\
48000	2100.4\\
50000	2096.4\\
52000	2098.6\\
54000	2094\\
56000	2098.3\\
62000	2091.2\\
66000	2087.2\\
68000	2096.9\\
76000	2091.1\\
78000	2081.6\\
82000	2085.2\\
86000	2081.7\\
88000	2087.8\\
92000	2086.9\\
94000	2081.6\\
98000	2087.2\\
1e+05	2082.2\\
};
\addplot [color=mycolor77, forget plot]
  table[row sep=crcr]{%
0	4395.3\\
2000	3974.6\\
4000	3345.7\\
6000	2900.2\\
8000	2584.1\\
10000	2354.5\\
12000	2181.4\\
14000	2052\\
16000	1965.4\\
18000	1891.8\\
20000	1837.4\\
24000	1758.5\\
26000	1730.4\\
36000	1666.3\\
40000	1655.8\\
42000	1651.5\\
46000	1653.6\\
50000	1644.2\\
52000	1645\\
54000	1640.1\\
56000	1640.5\\
58000	1647.4\\
64000	1648.8\\
66000	1640.7\\
70000	1642.4\\
72000	1640.9\\
76000	1652.5\\
80000	1642.8\\
96000	1643\\
98000	1645.9\\
1e+05	1640.7\\
};
\addplot [color=mycolor78, forget plot]
  table[row sep=crcr]{%
0	4395.3\\
2000	4113.5\\
4000	3698.9\\
6000	3380.1\\
8000	3156.1\\
10000	2981.9\\
12000	2857.3\\
16000	2668.4\\
18000	2607.3\\
20000	2559.9\\
22000	2524.3\\
30000	2422.1\\
34000	2400.3\\
40000	2407\\
42000	2396.8\\
44000	2398.2\\
46000	2382.7\\
50000	2373.7\\
54000	2358.7\\
56000	2365.1\\
62000	2361.3\\
66000	2371.8\\
68000	2365.1\\
72000	2367.6\\
76000	2353.5\\
84000	2362.3\\
88000	2360.1\\
92000	2360.1\\
1e+05	2356.7\\
};
\addplot [color=mycolor79, forget plot]
  table[row sep=crcr]{%
0	4395.3\\
2000	4081\\
4000	3612\\
6000	3276.8\\
8000	3036\\
10000	2849.3\\
12000	2719.5\\
14000	2618.1\\
16000	2538.9\\
20000	2433.5\\
22000	2393.2\\
26000	2342.9\\
28000	2317.9\\
36000	2264.9\\
38000	2255.1\\
40000	2251.7\\
42000	2254.6\\
44000	2246.5\\
46000	2250.2\\
50000	2237.2\\
52000	2243.4\\
56000	2232\\
66000	2241.9\\
80000	2223.7\\
82000	2239\\
84000	2246.8\\
90000	2242.4\\
94000	2249.7\\
1e+05	2246.6\\
};
\addplot [color=mycolor79, forget plot]
  table[row sep=crcr]{%
0	4395.3\\
2000	4105.3\\
4000	3648.8\\
6000	3291.5\\
8000	3027.6\\
10000	2826.2\\
12000	2672.5\\
14000	2561.7\\
16000	2475.9\\
18000	2401.7\\
22000	2292.5\\
24000	2247.5\\
26000	2217.6\\
32000	2156.7\\
36000	2138.7\\
40000	2128.3\\
42000	2115.7\\
44000	2112.7\\
46000	2103.5\\
48000	2104.9\\
52000	2093.7\\
54000	2091.2\\
56000	2100.7\\
58000	2099.6\\
62000	2109.5\\
66000	2100.8\\
70000	2109.4\\
78000	2101.2\\
80000	2105\\
86000	2084.8\\
88000	2071.1\\
90000	2068\\
96000	2093.7\\
1e+05	2081.6\\
};
\addplot [color=mycolor80, forget plot]
  table[row sep=crcr]{%
0	4395.3\\
2000	4013.4\\
4000	3423.1\\
6000	2998.3\\
8000	2675.2\\
10000	2438.1\\
12000	2256.5\\
14000	2114.3\\
16000	2009.3\\
18000	1928.8\\
22000	1818.9\\
24000	1786.5\\
30000	1714.9\\
34000	1688.8\\
38000	1669.9\\
54000	1635.7\\
62000	1635.6\\
68000	1623.6\\
78000	1613.9\\
84000	1623.3\\
88000	1618.8\\
90000	1627.8\\
98000	1623.9\\
1e+05	1622.8\\
};
\addplot [color=mycolor81, forget plot]
  table[row sep=crcr]{%
0	4395.3\\
2000	3986.9\\
4000	3379.7\\
6000	2941\\
8000	2608.8\\
10000	2360.1\\
12000	2173.6\\
14000	2032.6\\
16000	1928.3\\
18000	1841.7\\
20000	1779.9\\
22000	1730.4\\
26000	1656.2\\
30000	1607.1\\
32000	1588.5\\
36000	1566\\
46000	1539\\
52000	1528.6\\
58000	1525.3\\
62000	1517.8\\
66000	1529.1\\
72000	1531.7\\
78000	1529.2\\
82000	1525.9\\
84000	1531\\
88000	1527.4\\
92000	1528.2\\
96000	1521.8\\
1e+05	1527.1\\
};
\addplot [color=mycolor82, forget plot]
  table[row sep=crcr]{%
0	4395.3\\
2000	4034.6\\
4000	3479.4\\
6000	3066.5\\
8000	2763\\
10000	2542.4\\
12000	2374.9\\
14000	2241\\
16000	2141.8\\
18000	2063.1\\
20000	2003.8\\
22000	1958\\
24000	1920.6\\
28000	1865.1\\
30000	1851\\
32000	1829.9\\
38000	1803.6\\
40000	1791.8\\
44000	1783.9\\
48000	1786\\
52000	1779.4\\
90000	1771.9\\
92000	1775.9\\
94000	1774.1\\
98000	1779.8\\
1e+05	1775.1\\
};
\addplot [color=mycolor83, forget plot]
  table[row sep=crcr]{%
0	4395.3\\
2000	4063.2\\
4000	3545.7\\
6000	3162.8\\
8000	2872.7\\
10000	2647.4\\
12000	2471.7\\
14000	2334.6\\
16000	2229.1\\
18000	2148.1\\
20000	2085.6\\
22000	2032.9\\
26000	1955\\
30000	1902.4\\
34000	1865.3\\
38000	1836.8\\
42000	1819.8\\
54000	1808.6\\
56000	1813.1\\
64000	1807.4\\
70000	1799.5\\
72000	1806.4\\
80000	1801.7\\
84000	1803.1\\
86000	1805.7\\
90000	1800.9\\
94000	1795.8\\
98000	1804.1\\
1e+05	1807.6\\
};
\addplot [color=mycolor83, forget plot]
  table[row sep=crcr]{%
0	4395.3\\
2000	3970.7\\
4000	3322.8\\
6000	2842.6\\
8000	2484.1\\
10000	2223.1\\
12000	2029.5\\
14000	1872.6\\
16000	1750.6\\
18000	1663.2\\
20000	1598.4\\
22000	1544.3\\
24000	1498.7\\
28000	1436.1\\
32000	1395.7\\
34000	1380.9\\
36000	1379.3\\
46000	1347.3\\
50000	1339.2\\
52000	1341.9\\
62000	1336\\
64000	1339.5\\
72000	1331.5\\
76000	1335.9\\
78000	1336.5\\
82000	1329.3\\
94000	1318.2\\
96000	1323.6\\
1e+05	1319.7\\
};
\addplot [color=mycolor84, forget plot]
  table[row sep=crcr]{%
0	4395.3\\
2000	4027\\
4000	3427.2\\
6000	2950.5\\
8000	2611.5\\
10000	2344.5\\
12000	2125.4\\
14000	1959.3\\
16000	1830.4\\
18000	1729.3\\
20000	1662.4\\
22000	1603.4\\
26000	1511.7\\
28000	1480.2\\
30000	1460.3\\
32000	1436.1\\
34000	1408\\
38000	1381.8\\
40000	1377.1\\
44000	1372.3\\
46000	1361.7\\
50000	1351.7\\
52000	1357.3\\
54000	1358.7\\
56000	1356.1\\
62000	1360.1\\
64000	1359.6\\
66000	1355\\
70000	1355.2\\
74000	1351.2\\
76000	1350\\
80000	1361.8\\
84000	1362.1\\
86000	1358.4\\
88000	1359.1\\
92000	1349.5\\
96000	1344.7\\
1e+05	1351\\
};
\addplot [color=mycolor85, forget plot]
  table[row sep=crcr]{%
0	4395.3\\
2000	4007.1\\
4000	3408.8\\
6000	2957\\
8000	2618.5\\
10000	2369.3\\
12000	2182.2\\
14000	2032.7\\
16000	1922.3\\
18000	1832.5\\
20000	1761.9\\
24000	1653.2\\
26000	1607.5\\
30000	1553.5\\
36000	1505.3\\
44000	1484.8\\
48000	1482.1\\
50000	1486.4\\
52000	1486.3\\
58000	1465.9\\
64000	1466\\
70000	1469\\
74000	1467\\
84000	1456.3\\
90000	1464.5\\
92000	1461\\
96000	1464\\
98000	1471.9\\
1e+05	1475.7\\
};
\addplot [color=mycolor86, forget plot]
  table[row sep=crcr]{%
0	4395.3\\
2000	4025.8\\
4000	3432.9\\
6000	2981.3\\
8000	2619.7\\
10000	2333.4\\
12000	2109.5\\
14000	1935.3\\
16000	1797.5\\
18000	1691.7\\
20000	1603.5\\
22000	1534.1\\
26000	1436.4\\
28000	1403.2\\
32000	1355.7\\
34000	1338.7\\
36000	1312.8\\
40000	1287.9\\
48000	1276.1\\
50000	1277.4\\
54000	1261.5\\
56000	1261.7\\
58000	1257\\
60000	1247.5\\
62000	1255.2\\
64000	1255.9\\
66000	1248.5\\
68000	1253\\
70000	1261.4\\
78000	1249.7\\
82000	1249.4\\
84000	1240.7\\
90000	1238.1\\
94000	1237.4\\
98000	1231.8\\
1e+05	1234\\
};
\addplot [color=mycolor87, forget plot]
  table[row sep=crcr]{%
0	4395.3\\
2000	3963.4\\
4000	3319.4\\
6000	2823.4\\
8000	2455.7\\
10000	2162.5\\
12000	1936.3\\
14000	1770.1\\
16000	1629.4\\
18000	1506.4\\
20000	1414.1\\
22000	1358.6\\
24000	1296.2\\
26000	1250.7\\
28000	1217\\
30000	1191.6\\
32000	1171.1\\
36000	1122.9\\
38000	1100.8\\
44000	1075.5\\
48000	1047.6\\
50000	1040.8\\
54000	1036.7\\
58000	1037\\
62000	1055\\
64000	1057.4\\
66000	1056\\
68000	1062.5\\
74000	1062.8\\
76000	1056\\
78000	1062.7\\
80000	1064.3\\
90000	1055\\
92000	1058.2\\
96000	1058.1\\
98000	1064.7\\
1e+05	1067.3\\
};
\addplot [color=mycolor88, forget plot]
  table[row sep=crcr]{%
0	4395.3\\
2000	4056.3\\
4000	3501.7\\
6000	3067.3\\
8000	2714.5\\
10000	2427.9\\
12000	2207.2\\
14000	2031\\
16000	1887.3\\
18000	1772.9\\
20000	1672.6\\
22000	1587.6\\
24000	1528.3\\
26000	1476\\
28000	1440.2\\
32000	1385.9\\
34000	1362.2\\
40000	1308.1\\
42000	1295.7\\
44000	1288.8\\
46000	1286\\
58000	1246.9\\
62000	1247.3\\
64000	1249.2\\
68000	1242.5\\
76000	1239.5\\
78000	1242.8\\
86000	1241.6\\
88000	1242.4\\
90000	1238.9\\
92000	1240.7\\
96000	1248.7\\
1e+05	1250.9\\
};
\addplot [color=mycolor89, forget plot]
  table[row sep=crcr]{%
0	4395.3\\
2000	3987.4\\
4000	3332.9\\
6000	2824.6\\
8000	2440\\
10000	2131.2\\
12000	1880.6\\
14000	1691.8\\
16000	1540.6\\
18000	1421.8\\
20000	1323.3\\
22000	1240.5\\
24000	1172.1\\
26000	1122.7\\
30000	1055.1\\
34000	1007.3\\
38000	972.32\\
42000	951.94\\
46000	936.4\\
48000	930.84\\
50000	921.97\\
56000	912.04\\
58000	911.72\\
60000	915.2\\
66000	909.71\\
68000	911.32\\
70000	907.67\\
72000	912.22\\
74000	912.32\\
76000	908.47\\
84000	910.78\\
88000	916.49\\
96000	914.6\\
1e+05	916.51\\
};
\addplot [color=mycolor90, forget plot]
  table[row sep=crcr]{%
0	4395.3\\
2000	3991.6\\
4000	3350\\
6000	2848.4\\
8000	2449.7\\
10000	2144.7\\
12000	1902\\
14000	1707.1\\
16000	1552.7\\
18000	1426.6\\
20000	1323.3\\
22000	1242.1\\
24000	1172.8\\
26000	1114.9\\
28000	1070.3\\
30000	1039.3\\
32000	1013.9\\
34000	992.24\\
36000	974.76\\
38000	960.46\\
40000	952.58\\
44000	929.46\\
50000	905.74\\
54000	899.97\\
58000	898.38\\
62000	885.19\\
64000	882.26\\
68000	879.88\\
72000	879.32\\
76000	880.84\\
84000	873.57\\
88000	872.87\\
92000	868.01\\
94000	868.69\\
1e+05	880.89\\
};
\addplot [color=mycolor91, forget plot]
  table[row sep=crcr]{%
0	4395.3\\
2000	4007\\
4000	3372.3\\
6000	2865.8\\
8000	2458.8\\
10000	2128.1\\
12000	1862\\
14000	1656.6\\
16000	1500.7\\
18000	1363.9\\
20000	1253.8\\
22000	1168.4\\
24000	1098.6\\
26000	1038.3\\
28000	987.98\\
30000	947.1\\
32000	911.92\\
36000	859.81\\
38000	840.69\\
44000	810.29\\
50000	790.84\\
52000	788.65\\
56000	777.49\\
58000	777.27\\
60000	771.59\\
70000	770.57\\
72000	772.77\\
74000	768.76\\
78000	765.92\\
80000	771.1\\
82000	767.71\\
84000	768.12\\
86000	761.2\\
94000	765.23\\
98000	768.13\\
1e+05	765.64\\
};
\addplot [color=mycolor92, forget plot]
  table[row sep=crcr]{%
0	4395.3\\
2000	4051.3\\
4000	3481.5\\
6000	3021\\
8000	2650.8\\
10000	2350.6\\
12000	2106.8\\
14000	1904.7\\
16000	1742.9\\
18000	1609.6\\
20000	1493.2\\
22000	1400.3\\
24000	1328.4\\
26000	1265.4\\
28000	1213\\
30000	1169.8\\
36000	1079.4\\
38000	1055\\
40000	1036.3\\
44000	1011.6\\
46000	999\\
62000	956.92\\
66000	955.66\\
70000	953.1\\
82000	945.69\\
84000	950.28\\
86000	951.36\\
92000	941.91\\
96000	945.9\\
98000	942.48\\
1e+05	946.18\\
};
\addplot [color=mycolor93, forget plot]
  table[row sep=crcr]{%
0	4395.3\\
2000	4043.8\\
4000	3451.7\\
6000	2977\\
8000	2596.9\\
10000	2292.7\\
12000	2037.4\\
14000	1828.8\\
16000	1659.7\\
18000	1521.4\\
20000	1405\\
22000	1305.1\\
24000	1220.5\\
26000	1150.5\\
30000	1041.2\\
32000	1000.4\\
34000	970.24\\
38000	925.76\\
40000	900.49\\
42000	880.65\\
46000	856.08\\
48000	846.92\\
56000	831.67\\
58000	827.33\\
62000	826.51\\
64000	820.69\\
66000	818.88\\
70000	822.39\\
74000	822.74\\
76000	814.64\\
78000	809.62\\
80000	811.62\\
82000	818.41\\
84000	816.37\\
86000	808.6\\
90000	816.95\\
92000	817.58\\
96000	812.04\\
98000	808.33\\
1e+05	809.99\\
};
\addplot [color=mycolor94, forget plot]
  table[row sep=crcr]{%
0	4395.3\\
2000	4061.9\\
4000	3506.9\\
6000	3048.7\\
8000	2660\\
10000	2346.6\\
12000	2087.9\\
14000	1871.3\\
16000	1690.8\\
18000	1539.1\\
20000	1414.2\\
22000	1311\\
24000	1223.2\\
26000	1150.6\\
28000	1087.3\\
30000	1037.2\\
32000	996.14\\
36000	932.78\\
38000	906.74\\
42000	864.4\\
44000	844.44\\
46000	829.91\\
54000	796.15\\
56000	790.44\\
68000	773.4\\
72000	772.05\\
78000	765.38\\
80000	762.61\\
82000	765.34\\
86000	764.15\\
90000	762.89\\
94000	767.93\\
98000	767.95\\
1e+05	764.72\\
};
\addplot [color=mycolor95, forget plot]
  table[row sep=crcr]{%
0	4395.3\\
2000	4009.7\\
4000	3374\\
6000	2853.3\\
8000	2424.4\\
10000	2078.4\\
12000	1791.2\\
14000	1565.2\\
16000	1379.2\\
18000	1228.8\\
20000	1103\\
22000	994.9\\
24000	907.99\\
26000	832.72\\
30000	719.89\\
32000	679.26\\
34000	644.73\\
38000	588.93\\
42000	548.03\\
46000	518.41\\
52000	497.39\\
54000	493.07\\
56000	486.62\\
62000	481.51\\
66000	474.77\\
70000	476.38\\
76000	468.93\\
80000	466.05\\
82000	462.12\\
88000	463.98\\
90000	462.08\\
92000	461.48\\
94000	462.44\\
98000	455.14\\
1e+05	454.86\\
};
\addplot [color=mycolor96, forget plot]
  table[row sep=crcr]{%
0	4395.3\\
2000	4029.4\\
4000	3454.5\\
8000	2591.9\\
10000	2249.3\\
12000	1969.4\\
14000	1732.6\\
16000	1531.5\\
18000	1367.2\\
20000	1232\\
22000	1117.6\\
24000	1023.1\\
28000	868.01\\
30000	817.36\\
32000	772.11\\
36000	694.51\\
40000	641.56\\
46000	587.85\\
48000	577.21\\
50000	562.13\\
52000	550.74\\
54000	541.14\\
60000	520.77\\
62000	515.25\\
64000	508.29\\
66000	503.54\\
68000	496.8\\
70000	493.86\\
72000	497.36\\
74000	493.4\\
78000	494.29\\
80000	489.48\\
82000	486.48\\
84000	486.14\\
86000	481.76\\
92000	483.08\\
96000	489.38\\
98000	485.28\\
1e+05	484.34\\
};
\addplot [color=mycolor97, forget plot]
  table[row sep=crcr]{%
0	4395.3\\
2000	4038.5\\
4000	3423.6\\
6000	2914.7\\
8000	2496.1\\
10000	2148.9\\
12000	1857.7\\
16000	1412.2\\
18000	1241.3\\
20000	1100\\
22000	984.12\\
24000	885.16\\
26000	804.3\\
30000	679.47\\
32000	631.7\\
34000	591.28\\
36000	555.67\\
38000	526.54\\
42000	480.76\\
44000	462.35\\
46000	447.39\\
52000	410.12\\
56000	396.79\\
58000	391.5\\
62000	377.12\\
64000	373.69\\
66000	368.42\\
68000	364.89\\
70000	365.19\\
74000	362.75\\
76000	356.25\\
78000	354.62\\
82000	347.55\\
84000	344.72\\
86000	344.53\\
88000	342.64\\
90000	343.68\\
94000	342.84\\
96000	345.29\\
1e+05	344.95\\
};
\addplot [color=mycolor98]
  table[row sep=crcr]{%
0	4395.3\\
2000	4106.8\\
4000	3608.7\\
6000	3160.1\\
10000	2473.6\\
12000	2193.1\\
16000	1747.6\\
18000	1570.1\\
20000	1421.5\\
22000	1293.8\\
24000	1185.8\\
28000	1013.5\\
30000	943.19\\
32000	883.63\\
36000	781.7\\
40000	711.74\\
42000	681.32\\
44000	656.35\\
46000	634.63\\
48000	617.19\\
52000	577.38\\
54000	566.41\\
56000	552.58\\
58000	537.26\\
64000	511.56\\
68000	498.55\\
72000	488.11\\
74000	483.13\\
78000	477.18\\
82000	472.45\\
88000	470.83\\
90000	473.11\\
98000	463.82\\
1e+05	465.25\\
};
\addlegendentry{\ClasskIVLegend}

      % \input{../../../../.tex/20/KS/SizeEvolutionsfig2.tex}
    \end{groupplot}

    \begin{groupplot}[
      group style={
        group name=Row2,
        plotGroup,
        group size=2 by 1,
      },
      plotRow,
      row5Keys,
      xTickScaleLabelEmpty,
      plotLegendSE,
      %
      xmin=0,xmax=1e+05,
      ymin=100,ymax=1.5103e+05,
    ]
      \nextgroupplot[%
        at={($(Row1 c1r1.south west)-(0,\yPlotSep em)-(0,\plotHeight cm)$)},
        % xmin=0,xmax=1e+05,
        % ymin=100,ymax=1.4892e+05,
      ] % This file was created by matlab2tikz.
%
\definecolor{mycolor1}{rgb}{0.00146,0.00047,0.01387}%
\definecolor{mycolor2}{rgb}{0.01166,0.00942,0.06346}%
\definecolor{mycolor3}{rgb}{0.02943,0.02150,0.11462}%
\definecolor{mycolor4}{rgb}{0.06134,0.03659,0.17764}%
\definecolor{mycolor5}{rgb}{0.08741,0.04456,0.22481}%
\definecolor{mycolor6}{rgb}{0.12291,0.04754,0.28162}%
\definecolor{mycolor7}{rgb}{0.14907,0.04547,0.31709}%
\definecolor{mycolor8}{rgb}{0.18343,0.04033,0.35497}%
\definecolor{mycolor9}{rgb}{0.21795,0.03662,0.38352}%
\definecolor{mycolor10}{rgb}{0.24497,0.03705,0.40001}%
\definecolor{mycolor11}{rgb}{0.27135,0.04092,0.41198}%
\definecolor{mycolor12}{rgb}{0.29076,0.04564,0.41864}%
\definecolor{mycolor13}{rgb}{0.31628,0.05349,0.42512}%
\definecolor{mycolor14}{rgb}{0.33522,0.06006,0.42852}%
\definecolor{mycolor15}{rgb}{0.36028,0.06925,0.43150}%
\definecolor{mycolor16}{rgb}{0.37900,0.07625,0.43272}%
\definecolor{mycolor17}{rgb}{0.39767,0.08326,0.43318}%
\definecolor{mycolor18}{rgb}{0.41633,0.09020,0.43294}%
\definecolor{mycolor19}{rgb}{0.42877,0.09479,0.43241}%
\definecolor{mycolor20}{rgb}{0.44743,0.10160,0.43108}%
\definecolor{mycolor21}{rgb}{0.46610,0.10832,0.42912}%
\definecolor{mycolor22}{rgb}{0.47856,0.11276,0.42747}%
\definecolor{mycolor23}{rgb}{0.49726,0.11938,0.42449}%
\definecolor{mycolor24}{rgb}{0.50973,0.12377,0.42216}%
\definecolor{mycolor25}{rgb}{0.52221,0.12815,0.41955}%
\definecolor{mycolor26}{rgb}{0.53468,0.13253,0.41667}%
\definecolor{mycolor27}{rgb}{0.55339,0.13913,0.41183}%
\definecolor{mycolor28}{rgb}{0.56585,0.14357,0.40826}%
\definecolor{mycolor29}{rgb}{0.57830,0.14804,0.40441}%
\definecolor{mycolor30}{rgb}{0.59073,0.15256,0.40029}%
\definecolor{mycolor31}{rgb}{0.60314,0.15715,0.39589}%
\definecolor{mycolor32}{rgb}{0.61551,0.16182,0.39122}%
\definecolor{mycolor33}{rgb}{0.62169,0.16418,0.38878}%
\definecolor{mycolor34}{rgb}{0.63400,0.16899,0.38370}%
\definecolor{mycolor35}{rgb}{0.64626,0.17391,0.37836}%
\definecolor{mycolor36}{rgb}{0.65846,0.17896,0.37275}%
\definecolor{mycolor37}{rgb}{0.66454,0.18154,0.36985}%
\definecolor{mycolor38}{rgb}{0.67664,0.18681,0.36385}%
\definecolor{mycolor39}{rgb}{0.68865,0.19224,0.35760}%
\definecolor{mycolor40}{rgb}{0.69463,0.19502,0.35439}%
\definecolor{mycolor41}{rgb}{0.70650,0.20073,0.34778}%
\definecolor{mycolor42}{rgb}{0.71240,0.20366,0.34438}%
\definecolor{mycolor43}{rgb}{0.72410,0.20967,0.33742}%
\definecolor{mycolor44}{rgb}{0.72991,0.21276,0.33386}%
\definecolor{mycolor45}{rgb}{0.74142,0.21911,0.32658}%
\definecolor{mycolor46}{rgb}{0.74713,0.22238,0.32286}%
\definecolor{mycolor47}{rgb}{0.75279,0.22571,0.31909}%
\definecolor{mycolor48}{rgb}{0.76401,0.23255,0.31140}%
\definecolor{mycolor49}{rgb}{0.76956,0.23608,0.30749}%
\definecolor{mycolor50}{rgb}{0.77506,0.23967,0.30353}%
\definecolor{mycolor51}{rgb}{0.79129,0.25086,0.29139}%
\definecolor{mycolor52}{rgb}{0.79661,0.25473,0.28726}%
\definecolor{mycolor53}{rgb}{0.80187,0.25867,0.28310}%
\definecolor{mycolor54}{rgb}{0.81224,0.26679,0.27466}%
\definecolor{mycolor55}{rgb}{0.81734,0.27095,0.27039}%
\definecolor{mycolor56}{rgb}{0.82737,0.27952,0.26175}%
\definecolor{mycolor57}{rgb}{0.83230,0.28391,0.25738}%
\definecolor{mycolor58}{rgb}{0.83717,0.28839,0.25299}%
\definecolor{mycolor59}{rgb}{0.84671,0.29756,0.24411}%
\definecolor{mycolor60}{rgb}{0.85138,0.30226,0.23964}%
\definecolor{mycolor61}{rgb}{0.85599,0.30704,0.23513}%
\definecolor{mycolor62}{rgb}{0.86053,0.31189,0.23061}%
\definecolor{mycolor63}{rgb}{0.87374,0.32691,0.21689}%
\definecolor{mycolor64}{rgb}{0.87800,0.33206,0.21227}%
\definecolor{mycolor65}{rgb}{0.89034,0.34796,0.19829}%
\definecolor{mycolor66}{rgb}{0.89819,0.35891,0.18886}%
\definecolor{mycolor67}{rgb}{0.90200,0.36449,0.18412}%
\definecolor{mycolor68}{rgb}{0.90573,0.37014,0.17935}%
\definecolor{mycolor69}{rgb}{0.90939,0.37586,0.17456}%
\definecolor{mycolor70}{rgb}{0.91297,0.38164,0.16975}%
\definecolor{mycolor71}{rgb}{0.91646,0.38748,0.16492}%
\definecolor{mycolor72}{rgb}{0.91988,0.39339,0.16007}%
\definecolor{mycolor73}{rgb}{0.92322,0.39936,0.15519}%
\definecolor{mycolor74}{rgb}{0.92647,0.40539,0.15029}%
\definecolor{mycolor75}{rgb}{0.92964,0.41148,0.14537}%
\definecolor{mycolor76}{rgb}{0.93274,0.41763,0.14042}%
\definecolor{mycolor77}{rgb}{0.93575,0.42383,0.13544}%
\definecolor{mycolor78}{rgb}{0.93868,0.43009,0.13044}%
\definecolor{mycolor79}{rgb}{0.94152,0.43640,0.12541}%
\definecolor{mycolor80}{rgb}{0.94429,0.44277,0.12035}%
\definecolor{mycolor81}{rgb}{0.94956,0.45566,0.11016}%
\definecolor{mycolor82}{rgb}{0.95208,0.46218,0.10503}%
\definecolor{mycolor83}{rgb}{0.96129,0.48872,0.08429}%
\definecolor{mycolor84}{rgb}{0.96540,0.50225,0.07386}%
\definecolor{mycolor85}{rgb}{0.96732,0.50908,0.06866}%
\definecolor{mycolor86}{rgb}{0.97259,0.52980,0.05332}%
\definecolor{mycolor87}{rgb}{0.97568,0.54380,0.04362}%
\definecolor{mycolor88}{rgb}{0.98082,0.57221,0.02851}%
\definecolor{mycolor89}{rgb}{0.98189,0.57939,0.02625}%
\definecolor{mycolor90}{rgb}{0.98378,0.59385,0.02377}%
\definecolor{mycolor91}{rgb}{0.98532,0.60842,0.02420}%
\definecolor{mycolor92}{rgb}{0.98696,0.63048,0.03091}%
\definecolor{mycolor93}{rgb}{0.98793,0.66025,0.05175}%
\definecolor{mycolor94}{rgb}{0.98771,0.68281,0.07249}%
\definecolor{mycolor95}{rgb}{0.97750,0.78226,0.18592}%
\definecolor{mycolor96}{rgb}{0.96624,0.83619,0.26153}%
\definecolor{mycolor97}{rgb}{0.96252,0.85148,0.28555}%
\definecolor{mycolor98}{rgb}{0.98836,0.99836,0.64492}%
%
\addplot [color=mycolor1]
  table[row sep=crcr]{%
0	4395.3\\
2000	68979\\
4000	1.2115e+05\\
6000	1.3751e+05\\
8000	1.4384e+05\\
10000	1.4387e+05\\
12000	1.4331e+05\\
14000	1.4473e+05\\
16000	1.4716e+05\\
18000	1.4689e+05\\
20000	1.4766e+05\\
22000	1.4756e+05\\
24000	1.449e+05\\
26000	1.4713e+05\\
28000	1.4658e+05\\
30000	1.4762e+05\\
32000	1.46e+05\\
34000	1.4581e+05\\
38000	1.4819e+05\\
40000	1.4826e+05\\
44000	1.4687e+05\\
46000	1.4476e+05\\
48000	1.439e+05\\
50000	1.4721e+05\\
54000	1.4651e+05\\
56000	1.4789e+05\\
58000	1.4861e+05\\
62000	1.4719e+05\\
64000	1.484e+05\\
66000	1.4848e+05\\
68000	1.4662e+05\\
70000	1.4556e+05\\
72000	1.4559e+05\\
74000	1.469e+05\\
76000	1.4633e+05\\
78000	1.4663e+05\\
80000	1.4739e+05\\
82000	1.4892e+05\\
86000	1.4677e+05\\
88000	1.4716e+05\\
90000	1.4639e+05\\
92000	1.4718e+05\\
94000	1.4722e+05\\
1e+05	1.4535e+05\\
};
\addlegendentry{\ClasskILegend}

\addplot [color=mycolor2, forget plot]
  table[row sep=crcr]{%
0	4395.3\\
2000	22073\\
4000	26631\\
6000	25801\\
8000	25146\\
10000	24856\\
12000	24766\\
14000	24457\\
16000	25066\\
18000	24961\\
22000	25332\\
24000	25635\\
26000	24969\\
28000	25284\\
32000	25291\\
34000	25261\\
36000	25355\\
38000	25617\\
42000	25276\\
44000	25399\\
46000	25268\\
48000	25062\\
50000	24704\\
52000	25098\\
54000	25097\\
56000	25662\\
58000	25204\\
60000	24967\\
62000	25155\\
64000	24938\\
68000	25154\\
70000	25790\\
72000	25701\\
74000	25688\\
76000	25917\\
78000	25518\\
80000	24882\\
82000	25490\\
84000	25623\\
86000	25216\\
88000	25264\\
90000	25558\\
92000	25247\\
94000	25112\\
96000	25315\\
98000	25618\\
1e+05	25194\\
};
\addplot [color=mycolor3, forget plot]
  table[row sep=crcr]{%
0	4395.3\\
2000	29564\\
4000	38665\\
6000	38223\\
8000	37989\\
10000	38411\\
12000	37844\\
14000	38289\\
16000	38112\\
18000	38139\\
20000	38353\\
22000	39126\\
24000	39228\\
26000	38222\\
28000	38010\\
30000	37940\\
32000	38155\\
34000	38809\\
36000	37891\\
38000	38606\\
42000	37561\\
44000	38150\\
46000	38381\\
48000	38939\\
50000	39108\\
52000	38882\\
54000	38238\\
56000	39789\\
58000	38914\\
60000	38466\\
62000	38723\\
64000	39169\\
66000	38616\\
68000	39205\\
70000	38803\\
72000	38860\\
74000	37961\\
76000	38041\\
78000	38316\\
80000	38845\\
82000	39753\\
84000	39709\\
86000	38592\\
88000	37680\\
90000	37845\\
92000	37764\\
96000	38777\\
98000	38124\\
1e+05	38480\\
};
\addplot [color=mycolor4, forget plot]
  table[row sep=crcr]{%
0	4395.3\\
2000	14058\\
4000	14774\\
6000	14295\\
8000	14121\\
10000	14161\\
14000	13896\\
16000	14186\\
20000	13989\\
22000	14285\\
24000	14236\\
26000	14033\\
28000	14267\\
30000	13931\\
32000	14310\\
34000	13955\\
36000	13828\\
38000	13929\\
40000	14289\\
42000	13800\\
44000	14249\\
46000	14067\\
48000	14199\\
50000	13967\\
52000	14354\\
54000	13851\\
56000	14128\\
58000	13998\\
60000	14280\\
62000	14206\\
64000	14189\\
66000	14466\\
68000	14354\\
70000	14557\\
72000	14143\\
74000	14104\\
78000	14266\\
80000	14242\\
82000	13851\\
86000	13843\\
88000	14390\\
90000	14375\\
92000	14013\\
94000	14205\\
96000	14557\\
98000	14200\\
1e+05	14029\\
};
\addplot [color=mycolor5, forget plot]
  table[row sep=crcr]{%
0	4395.3\\
2000	36041\\
4000	58359\\
6000	63437\\
8000	63018\\
10000	63126\\
12000	61312\\
16000	62354\\
18000	63066\\
20000	62570\\
22000	62882\\
24000	62575\\
26000	62550\\
28000	63319\\
30000	62531\\
32000	62136\\
34000	61258\\
36000	61621\\
38000	59879\\
40000	60096\\
42000	62182\\
44000	62923\\
46000	62496\\
48000	61203\\
50000	60429\\
52000	60939\\
54000	62101\\
56000	62274\\
58000	61913\\
60000	62391\\
62000	62129\\
64000	60848\\
66000	60465\\
68000	60639\\
70000	61773\\
72000	62740\\
76000	61922\\
78000	61053\\
80000	61974\\
82000	61895\\
84000	60900\\
86000	62631\\
88000	62729\\
90000	61132\\
92000	60652\\
94000	61227\\
96000	63492\\
98000	62165\\
1e+05	61593\\
};
\addplot [color=mycolor6, forget plot]
  table[row sep=crcr]{%
0	4395.3\\
2000	29428\\
4000	46768\\
6000	50847\\
8000	50839\\
10000	52643\\
14000	54545\\
16000	53568\\
18000	54638\\
20000	54450\\
22000	53270\\
26000	53862\\
28000	54921\\
30000	54386\\
32000	55701\\
34000	55370\\
36000	55620\\
38000	54431\\
40000	53767\\
44000	55026\\
46000	54252\\
48000	54126\\
50000	54793\\
52000	54504\\
54000	54887\\
56000	53837\\
58000	55009\\
60000	55056\\
62000	54651\\
64000	53343\\
66000	53348\\
68000	54189\\
70000	54159\\
72000	53408\\
74000	53618\\
76000	55003\\
78000	54496\\
82000	53298\\
84000	53900\\
86000	53812\\
88000	53519\\
90000	53500\\
92000	54881\\
94000	55712\\
96000	54318\\
98000	53347\\
1e+05	54894\\
};
\addplot [color=mycolor7, forget plot]
  table[row sep=crcr]{%
0	4395.3\\
2000	18839\\
4000	24943\\
6000	26635\\
8000	27380\\
10000	27132\\
12000	27565\\
14000	27368\\
16000	27542\\
18000	27373\\
20000	26832\\
22000	27837\\
24000	28124\\
26000	27996\\
28000	28123\\
30000	27501\\
32000	27612\\
34000	27892\\
36000	27773\\
38000	28043\\
40000	27733\\
42000	28029\\
44000	28142\\
46000	27705\\
48000	27559\\
50000	27283\\
52000	27678\\
54000	27604\\
56000	26824\\
58000	27085\\
60000	26730\\
62000	27446\\
64000	26714\\
66000	26560\\
70000	27692\\
72000	27800\\
74000	27104\\
82000	27274\\
84000	27157\\
86000	28024\\
88000	27350\\
90000	27766\\
94000	27367\\
96000	27892\\
98000	27738\\
1e+05	28284\\
};
\addplot [color=mycolor8, forget plot]
  table[row sep=crcr]{%
0	4395.3\\
2000	17088\\
4000	21500\\
6000	20822\\
8000	20955\\
10000	21017\\
12000	20934\\
16000	21195\\
18000	20457\\
20000	20264\\
22000	20676\\
24000	20132\\
26000	20159\\
28000	19956\\
30000	20253\\
32000	20639\\
34000	19925\\
36000	20301\\
38000	19878\\
40000	20198\\
42000	20157\\
44000	20060\\
46000	20398\\
48000	20513\\
50000	20498\\
52000	20183\\
56000	20008\\
58000	19524\\
60000	20499\\
62000	20304\\
66000	19605\\
68000	20089\\
70000	19945\\
72000	19684\\
74000	20412\\
76000	20301\\
78000	19653\\
80000	19438\\
82000	19592\\
84000	19836\\
86000	19936\\
88000	19923\\
90000	20315\\
92000	20516\\
94000	20481\\
96000	20284\\
98000	20262\\
1e+05	19967\\
};
\addplot [color=mycolor9, forget plot]
  table[row sep=crcr]{%
0	4395.3\\
2000	11503\\
4000	13073\\
6000	13448\\
8000	13358\\
10000	13762\\
12000	13400\\
14000	13716\\
16000	13622\\
18000	13871\\
20000	13935\\
22000	13864\\
24000	13676\\
26000	13713\\
28000	13925\\
30000	13728\\
32000	13790\\
34000	13968\\
36000	14255\\
38000	13785\\
40000	13958\\
42000	13808\\
44000	13910\\
46000	14071\\
48000	14005\\
50000	13708\\
52000	14076\\
54000	13703\\
56000	14169\\
58000	14043\\
60000	14125\\
62000	13884\\
64000	14378\\
66000	14159\\
68000	14055\\
70000	14404\\
72000	14141\\
74000	13978\\
76000	14068\\
78000	14265\\
82000	13544\\
84000	13833\\
86000	13902\\
88000	14572\\
90000	13708\\
92000	13737\\
94000	13965\\
96000	13898\\
98000	13711\\
1e+05	13651\\
};
\addplot [color=mycolor10, forget plot]
  table[row sep=crcr]{%
0	4395.3\\
2000	14925\\
4000	20024\\
6000	20676\\
8000	21031\\
10000	20755\\
12000	21318\\
14000	21490\\
16000	21422\\
18000	21770\\
20000	21265\\
22000	21678\\
24000	21939\\
26000	21329\\
28000	21710\\
30000	21320\\
32000	21444\\
34000	21643\\
36000	21493\\
38000	21727\\
40000	21758\\
42000	21376\\
44000	21432\\
46000	22092\\
48000	21767\\
50000	21811\\
52000	21414\\
56000	21632\\
58000	22105\\
60000	22239\\
62000	21798\\
64000	22036\\
66000	22047\\
70000	21429\\
72000	22102\\
74000	22012\\
76000	21737\\
78000	21397\\
80000	21666\\
82000	21491\\
84000	21618\\
86000	21534\\
88000	21316\\
90000	21620\\
92000	21819\\
94000	21947\\
96000	21338\\
98000	21613\\
1e+05	21437\\
};
\addplot [color=mycolor11, forget plot]
  table[row sep=crcr]{%
0	4395.3\\
2000	9630.7\\
4000	9400.7\\
6000	8794.3\\
8000	8885\\
10000	8565.1\\
12000	8415.5\\
14000	8727.3\\
16000	8503.3\\
18000	8486.4\\
22000	8593.4\\
24000	8723.2\\
26000	8551.7\\
28000	8665.1\\
30000	8667.6\\
32000	8614.5\\
34000	8807.3\\
36000	8774.2\\
38000	8621.6\\
40000	8556.6\\
42000	8820.3\\
44000	8485.4\\
46000	8677.4\\
48000	8724.1\\
50000	8659.6\\
52000	8654\\
56000	8514.8\\
58000	8557.3\\
60000	8417.2\\
62000	8513.6\\
64000	8675.5\\
66000	8635.9\\
68000	8927.2\\
70000	8494.2\\
72000	8392.2\\
74000	8721.7\\
76000	8752.1\\
78000	8399.8\\
80000	8608.2\\
82000	8654.8\\
84000	8629.4\\
86000	8489.9\\
88000	8645.9\\
92000	8597.5\\
94000	8550\\
96000	8307.7\\
98000	8743.2\\
1e+05	8413.8\\
};
\addplot [color=mycolor12, forget plot]
  table[row sep=crcr]{%
0	4395.3\\
2000	11439\\
4000	13634\\
6000	13964\\
8000	14160\\
10000	14516\\
12000	14483\\
14000	14239\\
16000	14641\\
18000	14508\\
20000	14220\\
22000	14784\\
24000	14871\\
26000	14675\\
28000	14666\\
30000	14467\\
32000	14630\\
34000	14579\\
36000	14715\\
38000	14477\\
40000	14175\\
42000	14383\\
44000	14454\\
46000	14630\\
48000	14753\\
50000	14574\\
54000	14402\\
56000	14424\\
58000	13980\\
60000	14106\\
62000	14510\\
64000	15045\\
66000	14272\\
68000	14467\\
70000	14371\\
74000	14292\\
76000	14626\\
78000	14504\\
80000	14756\\
82000	14384\\
84000	14716\\
86000	14695\\
88000	14534\\
90000	14002\\
92000	14588\\
94000	14638\\
96000	14398\\
98000	14625\\
1e+05	14683\\
};
\addplot [color=mycolor13, forget plot]
  table[row sep=crcr]{%
0	4395.3\\
2000	9467.1\\
4000	10674\\
8000	10920\\
10000	10857\\
12000	11054\\
14000	10830\\
16000	11021\\
18000	10983\\
20000	10878\\
22000	10475\\
24000	10976\\
26000	10784\\
30000	11262\\
32000	11304\\
34000	10636\\
36000	11190\\
38000	11100\\
40000	11074\\
42000	10777\\
48000	11067\\
50000	10822\\
52000	10910\\
56000	10816\\
58000	11038\\
60000	10837\\
62000	10867\\
64000	10759\\
66000	11051\\
68000	10591\\
70000	11018\\
72000	11134\\
76000	11191\\
78000	11372\\
80000	11293\\
82000	11062\\
84000	10992\\
86000	11168\\
88000	11061\\
90000	11411\\
94000	11255\\
96000	10818\\
98000	10807\\
1e+05	11056\\
};
\addplot [color=mycolor14, forget plot]
  table[row sep=crcr]{%
0	4395.3\\
2000	15262\\
4000	21251\\
6000	22820\\
8000	23134\\
10000	23987\\
12000	23186\\
14000	23771\\
18000	23079\\
20000	23916\\
22000	23690\\
24000	23133\\
26000	24290\\
28000	24073\\
30000	23402\\
32000	23486\\
34000	24233\\
36000	23964\\
38000	23545\\
40000	23992\\
42000	24032\\
44000	23350\\
46000	23492\\
48000	24049\\
50000	24079\\
52000	23432\\
54000	23851\\
56000	23825\\
58000	23640\\
60000	24056\\
62000	23012\\
66000	23309\\
68000	23510\\
70000	24002\\
72000	23438\\
74000	24047\\
76000	23699\\
78000	24308\\
80000	24161\\
82000	24116\\
84000	24159\\
86000	24322\\
88000	23392\\
90000	24005\\
92000	23843\\
94000	23439\\
96000	23514\\
98000	23931\\
1e+05	24590\\
};
\addplot [color=mycolor15, forget plot]
  table[row sep=crcr]{%
0	4395.3\\
2000	12169\\
4000	15533\\
6000	16070\\
8000	15870\\
10000	16341\\
12000	16522\\
14000	16375\\
16000	16705\\
18000	16536\\
20000	16278\\
22000	16945\\
24000	17279\\
26000	17102\\
28000	16715\\
30000	16999\\
32000	16494\\
34000	16438\\
38000	16846\\
40000	16793\\
42000	16450\\
44000	16825\\
46000	16755\\
48000	17002\\
50000	16642\\
54000	16676\\
56000	16710\\
58000	16847\\
60000	16282\\
62000	16606\\
64000	16664\\
66000	17031\\
68000	16300\\
70000	16465\\
72000	16431\\
74000	16493\\
76000	16949\\
78000	16948\\
80000	16894\\
82000	16900\\
84000	16302\\
86000	16563\\
88000	16634\\
90000	16646\\
92000	16564\\
94000	16421\\
96000	16092\\
98000	16552\\
1e+05	16680\\
};
\addplot [color=mycolor16, forget plot]
  table[row sep=crcr]{%
0	4395.3\\
2000	7803.9\\
4000	7708.7\\
6000	7117.7\\
8000	6889.8\\
10000	6993.4\\
12000	6907.3\\
14000	7158\\
16000	6870.5\\
18000	6724.1\\
20000	7014.1\\
22000	6915.7\\
24000	6894.9\\
26000	6819\\
28000	7131.4\\
30000	6945.6\\
32000	7056.6\\
34000	6899.6\\
36000	7018.6\\
38000	6946.1\\
40000	6918.8\\
42000	6803.4\\
44000	6779.7\\
46000	6958.3\\
48000	6965.6\\
50000	6935.8\\
52000	6842.8\\
54000	6949.3\\
56000	6907.7\\
58000	6939.4\\
60000	7026.4\\
62000	6878.9\\
64000	7072\\
66000	7087.5\\
68000	7024\\
70000	7000.5\\
72000	6861.1\\
74000	7091.7\\
76000	7166.6\\
78000	7130.4\\
80000	6850.9\\
82000	6894.2\\
84000	7018.7\\
86000	7122.8\\
90000	7117.1\\
92000	6868.4\\
94000	7151.9\\
96000	6946.5\\
98000	6847\\
1e+05	6914.5\\
};
\addplot [color=mycolor17, forget plot]
  table[row sep=crcr]{%
0	4395.3\\
2000	11595\\
4000	14396\\
6000	15081\\
8000	15448\\
10000	15566\\
12000	16116\\
14000	15825\\
16000	15827\\
18000	16274\\
20000	15643\\
22000	15814\\
24000	15433\\
26000	15928\\
28000	15893\\
30000	15660\\
32000	15703\\
34000	15478\\
36000	15734\\
38000	15899\\
42000	15954\\
44000	15455\\
46000	15729\\
48000	15645\\
50000	15772\\
52000	16222\\
56000	15704\\
58000	15585\\
60000	16047\\
62000	16025\\
64000	15604\\
66000	15476\\
68000	16010\\
70000	15905\\
72000	15616\\
74000	15542\\
78000	15984\\
80000	15824\\
82000	15949\\
84000	15771\\
88000	15811\\
90000	15283\\
92000	15402\\
94000	15661\\
96000	15634\\
98000	15992\\
1e+05	15172\\
};
\addplot [color=mycolor18, forget plot]
  table[row sep=crcr]{%
0	4395.3\\
2000	14816\\
4000	21653\\
6000	22822\\
8000	23909\\
12000	24225\\
14000	24572\\
16000	24225\\
18000	23665\\
20000	24103\\
22000	23749\\
24000	24116\\
26000	23683\\
28000	23338\\
30000	23820\\
32000	23619\\
34000	23603\\
36000	24113\\
38000	24188\\
40000	24681\\
42000	24154\\
44000	23510\\
46000	24264\\
48000	24578\\
50000	24196\\
52000	24961\\
54000	24137\\
56000	24379\\
58000	24123\\
62000	24236\\
64000	23794\\
66000	24200\\
68000	24181\\
70000	23866\\
72000	24202\\
74000	24069\\
76000	23826\\
78000	24897\\
80000	24318\\
82000	24224\\
84000	24597\\
86000	24465\\
88000	24152\\
90000	23946\\
92000	24241\\
94000	24942\\
96000	24501\\
98000	23636\\
1e+05	23990\\
};
\addplot [color=mycolor19, forget plot]
  table[row sep=crcr]{%
0	4395.3\\
2000	10729\\
4000	13787\\
6000	14651\\
8000	14187\\
10000	14414\\
12000	14540\\
14000	14622\\
16000	14438\\
18000	14685\\
22000	14891\\
24000	15220\\
28000	14150\\
30000	14772\\
32000	14625\\
34000	14845\\
36000	14878\\
38000	14505\\
40000	15282\\
42000	15037\\
44000	14687\\
46000	14772\\
48000	14178\\
50000	15136\\
52000	14516\\
54000	14469\\
56000	14875\\
58000	14749\\
60000	14468\\
62000	14470\\
64000	14652\\
66000	14786\\
68000	15277\\
70000	14555\\
72000	14816\\
74000	14938\\
76000	14767\\
78000	14554\\
80000	15033\\
82000	15230\\
84000	14733\\
86000	14633\\
88000	14892\\
90000	14546\\
92000	14655\\
94000	14451\\
96000	14512\\
98000	14867\\
1e+05	14738\\
};
\addplot [color=mycolor20, forget plot]
  table[row sep=crcr]{%
0	4395.3\\
2000	12790\\
4000	17105\\
6000	17662\\
8000	18184\\
10000	18631\\
12000	18354\\
14000	18278\\
16000	18732\\
18000	19114\\
20000	19277\\
22000	18431\\
24000	18645\\
26000	19394\\
28000	18739\\
30000	18215\\
32000	18607\\
34000	18719\\
36000	18751\\
38000	19145\\
40000	18788\\
44000	18737\\
48000	18988\\
50000	19346\\
52000	18845\\
54000	19082\\
56000	18743\\
58000	18964\\
60000	18571\\
62000	18663\\
64000	19441\\
66000	19231\\
68000	18914\\
72000	19169\\
74000	18922\\
76000	19241\\
78000	18682\\
80000	18781\\
82000	18746\\
84000	18815\\
86000	18466\\
90000	18757\\
92000	19466\\
94000	18926\\
96000	19121\\
98000	19449\\
1e+05	19096\\
};
\addplot [color=mycolor21, forget plot]
  table[row sep=crcr]{%
0	4395.3\\
2000	7730\\
4000	8954\\
6000	9254.1\\
8000	9260.9\\
10000	9493.1\\
12000	9465.2\\
14000	9334.7\\
18000	9579.5\\
20000	9588.2\\
22000	9525.4\\
24000	9375.6\\
26000	9518.4\\
28000	9387.4\\
30000	9365.8\\
34000	9766\\
38000	9667.5\\
40000	9464.9\\
42000	9589.8\\
44000	9607.1\\
46000	9534.5\\
48000	9435.3\\
50000	9442.6\\
52000	9555.1\\
54000	9724.9\\
56000	9454.5\\
58000	9704.4\\
60000	9459.9\\
62000	9774.7\\
64000	9735\\
66000	9368.4\\
68000	9588.9\\
70000	9719.7\\
72000	9721.1\\
74000	9618.3\\
76000	9577.8\\
78000	9495.2\\
80000	9579\\
82000	9535.8\\
84000	9721.5\\
86000	9516.2\\
88000	9375.3\\
90000	9634.7\\
92000	9635.8\\
94000	9551.9\\
98000	9704.4\\
1e+05	9609.4\\
};
\addplot [color=mycolor22, forget plot]
  table[row sep=crcr]{%
0	4395.3\\
2000	18005\\
4000	25419\\
6000	26231\\
8000	26469\\
10000	26129\\
12000	26108\\
14000	25610\\
16000	25433\\
20000	24580\\
22000	24745\\
24000	24386\\
26000	25389\\
28000	25160\\
30000	24696\\
32000	24839\\
34000	24608\\
36000	24870\\
38000	24484\\
40000	24173\\
44000	25014\\
46000	25148\\
48000	24757\\
50000	24201\\
52000	23850\\
54000	24328\\
56000	24693\\
58000	24765\\
60000	24510\\
62000	25165\\
64000	24346\\
66000	24146\\
68000	24746\\
72000	24809\\
74000	24882\\
76000	23782\\
78000	24280\\
80000	24868\\
84000	25227\\
86000	24650\\
88000	24877\\
90000	24757\\
92000	24561\\
94000	24268\\
96000	24669\\
98000	24784\\
1e+05	24462\\
};
\addplot [color=mycolor23, forget plot]
  table[row sep=crcr]{%
0	4395.3\\
2000	9885.5\\
4000	11985\\
6000	12195\\
8000	12340\\
10000	12372\\
12000	12165\\
14000	12082\\
16000	11946\\
18000	11980\\
20000	12158\\
22000	12075\\
24000	12354\\
26000	12384\\
28000	12148\\
32000	12301\\
34000	12097\\
36000	12238\\
38000	12126\\
40000	12255\\
42000	12199\\
44000	12230\\
46000	12341\\
48000	12211\\
50000	12321\\
52000	12246\\
54000	12046\\
56000	12179\\
58000	12054\\
62000	12384\\
64000	12073\\
66000	12160\\
68000	12146\\
70000	12309\\
72000	12132\\
74000	12060\\
80000	12309\\
82000	12236\\
84000	12076\\
86000	12204\\
90000	12223\\
94000	12124\\
96000	12032\\
98000	12190\\
1e+05	12117\\
};
\addplot [color=mycolor24, forget plot]
  table[row sep=crcr]{%
0	4395.3\\
2000	5794.3\\
4000	5278.8\\
6000	4909.1\\
8000	4853.5\\
10000	4717.4\\
12000	4884.9\\
14000	4874\\
16000	4783.7\\
18000	4853.3\\
20000	4750.2\\
24000	4900\\
26000	4863.4\\
28000	4895.8\\
30000	4849.5\\
32000	4847.4\\
34000	4793.6\\
36000	4879.7\\
38000	4835.4\\
40000	4853.9\\
42000	4810\\
44000	4895.3\\
46000	4844.3\\
48000	4942\\
50000	4818.5\\
52000	4990.3\\
54000	4920.5\\
56000	4878.8\\
58000	4894.8\\
60000	4727.2\\
64000	4763.6\\
66000	4843.7\\
68000	4807.3\\
70000	4873.2\\
72000	4739.9\\
74000	4867.3\\
76000	4731\\
78000	4732.2\\
82000	4823.4\\
84000	4829.1\\
86000	4868.4\\
88000	4786.9\\
90000	4861.7\\
92000	4791.3\\
94000	4915.6\\
96000	4947.3\\
1e+05	4875.1\\
};
\addplot [color=mycolor25, forget plot]
  table[row sep=crcr]{%
0	4395.3\\
2000	9548.5\\
4000	11572\\
6000	11895\\
8000	11841\\
12000	12346\\
14000	12062\\
16000	12381\\
18000	12301\\
20000	12554\\
24000	12139\\
26000	12121\\
28000	12239\\
30000	12190\\
32000	12552\\
34000	12422\\
36000	12211\\
38000	12318\\
40000	12084\\
44000	12033\\
46000	11891\\
48000	12589\\
50000	12707\\
52000	12271\\
54000	12389\\
56000	11970\\
58000	12357\\
60000	12132\\
62000	12080\\
64000	12227\\
66000	12272\\
68000	11860\\
70000	11984\\
72000	12214\\
74000	12347\\
76000	12217\\
78000	12396\\
80000	12307\\
82000	12464\\
84000	12165\\
86000	12091\\
88000	12405\\
90000	12028\\
92000	12045\\
94000	12101\\
96000	12527\\
1e+05	12313\\
};
\addplot [color=mycolor26, forget plot]
  table[row sep=crcr]{%
0	4395.3\\
2000	6927.2\\
4000	7848.7\\
6000	7905.4\\
8000	8164.9\\
10000	8256.1\\
12000	8406.7\\
14000	8201.4\\
16000	8051.5\\
18000	8273.4\\
20000	8220.7\\
22000	8194.4\\
24000	8345.6\\
26000	8215.5\\
30000	8100.2\\
32000	8099.9\\
34000	8367.9\\
36000	8450.3\\
38000	8313.5\\
40000	8141.6\\
42000	8070\\
44000	8227.8\\
46000	8236.6\\
48000	8421.6\\
50000	8278.1\\
52000	8217.6\\
54000	8512.7\\
56000	8353.2\\
58000	8165.2\\
60000	8294.8\\
62000	8104.4\\
64000	8200.5\\
66000	8246.2\\
68000	8481\\
70000	8441.5\\
72000	8325.7\\
74000	7989\\
76000	8179.1\\
78000	8274.8\\
80000	8100.5\\
82000	8292.1\\
84000	8132.8\\
86000	8217.4\\
88000	8223.8\\
90000	8153.2\\
92000	8300.4\\
94000	8259.4\\
96000	8113.9\\
98000	8266.1\\
1e+05	8238.9\\
};
\addplot [color=mycolor27, forget plot]
  table[row sep=crcr]{%
0	4395.3\\
2000	8971.1\\
4000	10266\\
6000	9790.9\\
8000	9498.3\\
12000	9825.6\\
14000	9588.9\\
16000	9674\\
18000	9481.4\\
22000	9323.9\\
24000	9202.2\\
26000	9282.3\\
28000	9666.7\\
30000	9780.6\\
32000	9745\\
34000	9764.5\\
36000	9598.3\\
38000	9733\\
40000	9619.8\\
42000	9877.7\\
44000	9686.9\\
46000	9587.2\\
48000	9905.1\\
50000	9689\\
52000	10013\\
54000	10091\\
56000	9652\\
58000	9783\\
60000	9822.6\\
62000	9813.6\\
64000	9855\\
66000	9763.8\\
68000	9829.8\\
70000	10225\\
72000	9781.3\\
74000	9600.7\\
76000	9846.7\\
78000	9836.7\\
80000	9703.7\\
82000	9893.3\\
84000	9569.7\\
86000	9494.2\\
88000	9561.9\\
90000	9748.6\\
92000	9605.9\\
94000	9737.7\\
96000	9601.7\\
98000	9869.9\\
1e+05	9634.4\\
};
\addplot [color=mycolor28, forget plot]
  table[row sep=crcr]{%
0	4395.3\\
2000	6262.1\\
6000	5689.9\\
8000	5727.6\\
10000	5557.2\\
12000	5459.1\\
14000	5560.2\\
16000	5626.7\\
18000	5507.8\\
20000	5611.3\\
22000	5387.7\\
24000	5342.5\\
26000	5608.3\\
28000	5459.7\\
30000	5491.6\\
32000	5576.7\\
34000	5509.9\\
36000	5491.2\\
38000	5697.4\\
40000	5510.2\\
42000	5455.7\\
44000	5290.3\\
46000	5500.3\\
50000	5552.9\\
52000	5601\\
54000	5518.4\\
56000	5454.2\\
58000	5449.8\\
60000	5388.7\\
62000	5624.1\\
64000	5635\\
66000	5605.2\\
68000	5535.4\\
70000	5688.1\\
72000	5529.4\\
74000	5446.1\\
76000	5425.4\\
78000	5540.3\\
80000	5429\\
82000	5429\\
84000	5564\\
86000	5565.7\\
88000	5460.1\\
90000	5568.7\\
92000	5508.7\\
94000	5430.6\\
96000	5497.1\\
98000	5426.2\\
1e+05	5453.6\\
};
\addplot [color=mycolor29, forget plot]
  table[row sep=crcr]{%
0	4395.3\\
2000	10276\\
4000	11898\\
6000	11284\\
8000	11555\\
10000	11458\\
12000	11560\\
14000	11530\\
16000	11597\\
18000	11370\\
20000	11271\\
22000	11029\\
24000	11560\\
26000	11353\\
28000	11665\\
30000	11686\\
32000	11167\\
34000	11531\\
36000	11447\\
38000	11465\\
40000	11375\\
42000	11529\\
44000	11837\\
46000	11709\\
48000	11388\\
50000	11955\\
54000	11679\\
56000	11624\\
58000	11470\\
60000	11428\\
62000	11352\\
64000	11727\\
66000	11520\\
68000	11173\\
70000	11498\\
74000	11417\\
76000	11212\\
78000	11514\\
80000	11668\\
82000	11621\\
84000	11063\\
86000	11222\\
88000	11644\\
90000	11783\\
92000	11530\\
94000	11382\\
96000	11457\\
98000	11570\\
1e+05	11521\\
};
\addplot [color=mycolor30, forget plot]
  table[row sep=crcr]{%
0	4395.3\\
2000	6456.6\\
4000	7285.1\\
6000	7496.3\\
8000	7624.1\\
10000	7727.2\\
12000	7942.7\\
14000	7856.7\\
16000	7846.3\\
18000	7799.3\\
20000	8026\\
22000	7746.9\\
24000	7840.8\\
28000	7847.6\\
30000	7967.8\\
32000	7933.5\\
34000	8053.5\\
36000	7788.4\\
38000	7907\\
40000	7786.8\\
42000	8014.7\\
44000	7754.4\\
46000	8019.8\\
48000	7960.9\\
50000	7846.2\\
52000	7805\\
54000	7831\\
56000	7980.7\\
58000	7917.5\\
60000	7921.6\\
62000	8044.8\\
64000	7852.5\\
66000	7968\\
68000	7922.9\\
70000	7721.1\\
72000	8015\\
74000	8031\\
76000	7890\\
80000	7712.9\\
82000	7733.2\\
84000	7816.1\\
86000	8012.9\\
88000	8006.4\\
90000	7803.4\\
92000	7894.9\\
94000	7785\\
96000	7971.2\\
98000	7932\\
1e+05	8055.1\\
};
\addplot [color=mycolor31, forget plot]
  table[row sep=crcr]{%
0	4395.3\\
2000	6868\\
4000	6825.7\\
6000	6757\\
8000	6659.1\\
10000	6501.2\\
12000	6679.6\\
14000	6526.2\\
16000	6401.1\\
18000	6486.9\\
20000	6364.7\\
22000	6266\\
24000	6471.2\\
26000	6540.3\\
28000	6471.1\\
30000	6455.8\\
32000	6233.1\\
34000	6526.8\\
36000	6557\\
38000	6397.7\\
40000	6366.3\\
42000	6396.4\\
44000	6714.1\\
46000	6439.4\\
48000	6644.1\\
50000	6425.3\\
52000	6407.7\\
56000	6629.1\\
58000	6486.4\\
60000	6458.4\\
62000	6514.8\\
64000	6370\\
66000	6643\\
68000	6603.7\\
70000	6653.6\\
72000	6455.3\\
74000	6541.5\\
76000	6566\\
78000	6474.9\\
80000	6507\\
82000	6312.1\\
84000	6482.8\\
86000	6374.4\\
88000	6240.1\\
90000	6547.5\\
92000	6582\\
94000	6445.8\\
96000	6412.9\\
98000	6703.3\\
1e+05	6717.7\\
};
\addplot [color=mycolor32, forget plot]
  table[row sep=crcr]{%
0	4395.3\\
2000	9803.5\\
4000	11063\\
6000	10635\\
10000	10771\\
12000	10907\\
14000	10820\\
16000	10195\\
18000	10470\\
22000	10430\\
24000	10346\\
26000	10579\\
28000	10630\\
30000	10457\\
32000	10432\\
34000	10683\\
36000	10495\\
38000	10260\\
40000	10325\\
42000	10531\\
44000	10616\\
46000	10548\\
50000	10154\\
52000	10132\\
54000	10365\\
56000	10397\\
58000	10257\\
60000	10430\\
62000	10359\\
64000	10559\\
68000	10331\\
70000	10545\\
72000	10434\\
74000	10535\\
76000	10249\\
78000	10423\\
80000	10267\\
82000	10407\\
84000	10337\\
86000	9990.8\\
88000	10292\\
90000	10315\\
92000	10544\\
94000	10393\\
96000	10025\\
98000	10359\\
1e+05	10413\\
};
\addplot [color=mycolor33, forget plot]
  table[row sep=crcr]{%
0	4395.3\\
2000	6507.8\\
4000	7656.6\\
6000	7740.3\\
8000	7725.4\\
10000	8209.5\\
12000	8279.9\\
14000	8391\\
16000	8134.2\\
18000	7972.7\\
20000	7933.1\\
22000	8289.1\\
26000	8277.1\\
28000	8036\\
30000	8275.5\\
32000	8047.8\\
34000	8104.6\\
36000	8053.2\\
38000	8217.1\\
42000	8278.8\\
44000	8499.3\\
46000	8220.9\\
48000	8366.4\\
50000	8239.8\\
52000	7950.1\\
54000	8535.7\\
56000	8163.4\\
58000	8116.6\\
60000	7944.1\\
64000	8294.4\\
66000	8435.2\\
68000	8343.2\\
70000	8155.4\\
72000	8114.4\\
74000	7953.2\\
76000	8033.9\\
78000	8308.2\\
80000	8248.7\\
82000	8156.9\\
84000	8267.8\\
86000	8216.6\\
88000	8189.5\\
90000	8232.6\\
92000	8503.9\\
94000	8136.8\\
96000	8300\\
1e+05	8119.9\\
};
\addplot [color=mycolor34, forget plot]
  table[row sep=crcr]{%
0	4395.3\\
2000	5085.8\\
4000	4309.5\\
6000	4139.7\\
8000	4016.9\\
10000	3985.1\\
12000	4060.4\\
14000	3988.1\\
16000	3989.2\\
18000	4073.6\\
20000	3986.3\\
22000	3939.5\\
24000	3999.8\\
26000	4072.9\\
28000	4074\\
30000	4005.7\\
32000	3972.8\\
34000	4002.9\\
36000	4014.5\\
38000	4190.6\\
42000	4060.5\\
44000	4031.9\\
46000	4079.2\\
48000	4075.1\\
50000	4127.4\\
52000	4150.4\\
54000	4015.5\\
56000	4064.1\\
60000	4222.4\\
62000	4075.6\\
64000	4034.4\\
66000	4121.1\\
68000	4104.7\\
70000	4007.2\\
72000	4000.5\\
74000	4112.2\\
76000	4013.9\\
78000	3994.4\\
80000	4142.1\\
82000	4134.8\\
84000	3987.6\\
86000	4126.3\\
88000	4148.6\\
90000	4085.9\\
94000	4043.8\\
96000	4162.8\\
98000	4075.1\\
1e+05	4065.6\\
};
\addplot [color=mycolor35]
  table[row sep=crcr]{%
0	4395.3\\
2000	8360.6\\
4000	9718.6\\
6000	9679.7\\
8000	9983\\
10000	9972.1\\
12000	10147\\
14000	10156\\
16000	10097\\
18000	9979.6\\
20000	10105\\
22000	9864.6\\
24000	10123\\
26000	10119\\
28000	10017\\
30000	9966.9\\
32000	10084\\
34000	10295\\
36000	10196\\
38000	10203\\
40000	9889.7\\
42000	10009\\
44000	10030\\
46000	9998\\
48000	10083\\
50000	10121\\
52000	10208\\
58000	9913.3\\
60000	10027\\
62000	10067\\
64000	9982.6\\
68000	10435\\
70000	10294\\
72000	10231\\
74000	10127\\
76000	10106\\
78000	10265\\
80000	9971.2\\
82000	10075\\
84000	10049\\
88000	10049\\
90000	10101\\
94000	9952.4\\
96000	10103\\
98000	9958.8\\
1e+05	10048\\
};
\addlegendentry{\ClasskIILegend}

\addplot [color=mycolor36, forget plot]
  table[row sep=crcr]{%
0	4395.3\\
2000	5380.1\\
4000	4800.8\\
6000	4558\\
8000	4426.7\\
10000	4401.9\\
12000	4250.2\\
14000	4326.3\\
16000	4344\\
18000	4300.6\\
20000	4321.9\\
22000	4256.9\\
24000	4318.7\\
26000	4251.9\\
30000	4190.6\\
32000	4314.1\\
34000	4410.6\\
36000	4250.2\\
38000	4284.7\\
40000	4446\\
42000	4339.4\\
44000	4261.6\\
46000	4326.5\\
48000	4323.5\\
50000	4279.2\\
52000	4376.3\\
54000	4426.2\\
56000	4331.1\\
58000	4134.3\\
60000	4218.2\\
62000	4467.1\\
64000	4376.2\\
66000	4332\\
68000	4388.7\\
70000	4211.1\\
72000	4215\\
74000	4312.2\\
76000	4239.4\\
78000	4311.5\\
80000	4412.2\\
82000	4358.8\\
84000	4228.9\\
86000	4239.7\\
88000	4350.2\\
92000	4296.8\\
94000	4350.5\\
96000	4522.5\\
98000	4373.9\\
1e+05	4367.4\\
};
\addplot [color=mycolor37, forget plot]
  table[row sep=crcr]{%
0	4395.3\\
2000	4311.4\\
4000	3896.8\\
6000	3714.9\\
8000	3638\\
10000	3696.6\\
12000	3700.7\\
14000	3595.2\\
16000	3584.7\\
18000	3697.1\\
20000	3646.1\\
22000	3695.1\\
24000	3674.4\\
26000	3751.3\\
28000	3677.2\\
30000	3751.9\\
32000	3761.9\\
34000	3674.7\\
36000	3628.1\\
38000	3678\\
40000	3645.5\\
42000	3669\\
44000	3674.7\\
46000	3696.4\\
48000	3674.4\\
50000	3698.2\\
54000	3703.4\\
56000	3641.9\\
58000	3705.5\\
60000	3651.6\\
62000	3672.6\\
64000	3645.2\\
66000	3717\\
68000	3662.4\\
70000	3646.4\\
72000	3652.7\\
74000	3681.9\\
76000	3672.6\\
78000	3641.1\\
80000	3625.3\\
86000	3711.4\\
88000	3763.1\\
90000	3671.3\\
92000	3668.2\\
94000	3612.8\\
96000	3639\\
98000	3704.8\\
1e+05	3658.3\\
};
\addplot [color=mycolor38, forget plot]
  table[row sep=crcr]{%
0	4395.3\\
2000	3532.4\\
4000	3350.1\\
6000	3259.8\\
8000	3026.7\\
10000	3098.8\\
12000	3124.9\\
14000	3120.1\\
18000	3051.2\\
20000	3106.7\\
22000	3058.2\\
24000	2989.8\\
26000	3008.5\\
28000	3053\\
30000	3020.2\\
32000	3054.5\\
34000	3046.5\\
36000	2980.1\\
38000	3037.9\\
40000	3071\\
42000	3086.1\\
44000	3068.1\\
46000	3029.4\\
48000	3027.8\\
50000	3058.5\\
52000	3010.7\\
54000	3042.8\\
56000	3043\\
58000	3094.4\\
60000	3029.5\\
62000	3107.3\\
64000	3088.1\\
66000	3044.1\\
68000	3158.2\\
70000	2968.3\\
72000	3010.1\\
74000	2985\\
76000	3103\\
78000	3040.5\\
80000	3024.4\\
82000	3046.5\\
84000	2981\\
86000	3070.3\\
88000	3010.1\\
90000	2970.6\\
92000	3046\\
96000	2985.4\\
98000	3062.7\\
1e+05	3033.7\\
};
\addplot [color=mycolor39, forget plot]
  table[row sep=crcr]{%
0	4395.3\\
2000	10339\\
4000	12716\\
6000	12600\\
8000	12577\\
10000	11915\\
12000	11833\\
14000	11828\\
16000	12022\\
18000	11930\\
20000	12031\\
22000	11827\\
24000	11809\\
26000	12172\\
28000	11620\\
30000	11667\\
32000	11899\\
34000	11881\\
36000	11596\\
38000	11705\\
40000	12033\\
42000	12166\\
44000	11420\\
46000	11938\\
48000	11789\\
50000	11797\\
52000	11941\\
54000	11761\\
56000	11415\\
58000	11612\\
60000	11287\\
62000	11570\\
64000	11717\\
66000	11644\\
68000	11668\\
70000	11631\\
72000	11737\\
74000	12030\\
76000	12055\\
78000	11411\\
82000	11634\\
84000	11789\\
86000	11858\\
88000	11831\\
90000	11972\\
92000	11764\\
94000	11493\\
96000	11686\\
98000	11722\\
1e+05	12061\\
};
\addplot [color=mycolor40, forget plot]
  table[row sep=crcr]{%
0	4395.3\\
2000	4907.7\\
4000	5446.4\\
6000	5566.4\\
8000	5656.5\\
10000	5669.6\\
12000	5804\\
14000	5744.8\\
16000	5661.3\\
18000	5743.2\\
20000	5664.7\\
22000	5642.7\\
24000	5704\\
26000	5628.4\\
28000	5740.2\\
30000	5692.6\\
32000	5777\\
34000	5666.4\\
36000	5824.9\\
38000	5764.5\\
40000	5656.1\\
42000	5693.3\\
44000	5828.6\\
46000	5889.9\\
48000	5834.5\\
50000	5688.5\\
52000	5764.6\\
54000	5858.6\\
56000	5818.1\\
58000	5601.9\\
60000	5691.4\\
62000	5589.6\\
64000	5759.4\\
66000	5699.4\\
68000	5793.1\\
72000	5726.8\\
74000	5753.5\\
76000	5649.3\\
78000	5677.8\\
80000	5793.4\\
82000	5660.4\\
84000	5856.3\\
86000	5841.5\\
88000	5845.4\\
90000	5657.7\\
92000	5768.8\\
94000	5836.5\\
96000	5844.1\\
98000	5875.1\\
1e+05	5836.8\\
};
\addplot [color=mycolor41, forget plot]
  table[row sep=crcr]{%
0	4395.3\\
2000	4671.2\\
4000	5044.9\\
6000	5159.8\\
8000	5366.6\\
10000	5190.4\\
12000	5345.2\\
14000	5549.4\\
18000	5514\\
20000	5337.7\\
22000	5348.8\\
24000	5313.6\\
26000	5388.8\\
28000	5366.3\\
30000	5516.6\\
32000	5446.7\\
34000	5496.2\\
36000	5277.5\\
38000	5395.1\\
40000	5379.2\\
42000	5281.2\\
44000	5373.4\\
46000	5318.4\\
48000	5369.5\\
50000	5468.3\\
52000	5605.1\\
54000	5412.8\\
56000	5345.2\\
60000	5597.6\\
62000	5506.2\\
66000	5383.7\\
68000	5366.1\\
70000	5462.6\\
72000	5422.8\\
74000	5363.4\\
76000	5440.9\\
78000	5330.6\\
80000	5550.4\\
82000	5490\\
84000	5365.9\\
86000	5442.9\\
88000	5366.5\\
90000	5452.2\\
92000	5406.1\\
94000	5432\\
96000	5488.1\\
98000	5302.4\\
1e+05	5463.7\\
};
\addplot [color=mycolor42, forget plot]
  table[row sep=crcr]{%
0	4395.3\\
2000	11480\\
4000	16894\\
6000	19048\\
8000	19533\\
10000	19779\\
12000	20365\\
18000	19800\\
20000	19758\\
22000	20013\\
24000	19604\\
26000	19309\\
28000	19930\\
30000	20270\\
32000	20390\\
34000	20294\\
36000	20826\\
38000	20033\\
40000	19677\\
42000	20313\\
44000	20339\\
46000	19999\\
48000	20084\\
50000	19991\\
52000	20192\\
54000	19694\\
60000	20007\\
62000	20393\\
64000	20722\\
66000	20160\\
68000	20622\\
70000	20837\\
72000	20709\\
74000	20698\\
76000	20361\\
78000	20298\\
80000	19306\\
82000	19943\\
84000	20811\\
86000	20029\\
88000	19718\\
90000	20092\\
92000	19633\\
94000	19965\\
96000	19976\\
98000	20142\\
1e+05	20603\\
};
\addplot [color=mycolor43, forget plot]
  table[row sep=crcr]{%
0	4395.3\\
2000	5438.3\\
4000	5357.7\\
6000	5239\\
10000	5127.3\\
12000	5062.4\\
14000	4936.2\\
16000	5001.9\\
18000	5049.5\\
20000	5054.8\\
22000	5026.8\\
24000	4905.9\\
26000	4901.6\\
28000	5045.9\\
32000	4933.8\\
34000	4882.8\\
38000	5043.2\\
40000	4886.6\\
42000	4902.7\\
44000	4950.6\\
46000	4850.6\\
50000	4980.7\\
52000	4963.6\\
54000	4965.6\\
56000	5017.7\\
60000	4937.9\\
62000	4975.2\\
64000	5081.6\\
66000	4998.9\\
68000	4971.4\\
70000	4905.2\\
72000	4984.3\\
74000	4898.5\\
76000	4982.8\\
78000	4887.6\\
80000	4867.8\\
84000	4881.6\\
88000	4983.9\\
90000	4967\\
92000	4856.5\\
94000	4877.2\\
96000	5046.3\\
98000	5016.3\\
1e+05	5021.2\\
};
\addplot [color=mycolor44, forget plot]
  table[row sep=crcr]{%
0	4395.3\\
2000	4425.7\\
4000	3915.1\\
6000	3690.1\\
8000	3521.4\\
10000	3487.9\\
12000	3503.7\\
14000	3426.3\\
16000	3540.5\\
18000	3532.8\\
20000	3363.6\\
22000	3394.2\\
24000	3533.4\\
26000	3402.5\\
28000	3330.5\\
34000	3539.4\\
36000	3393.9\\
38000	3354.1\\
40000	3452.8\\
42000	3361.5\\
44000	3357.3\\
46000	3396.1\\
48000	3418.1\\
50000	3414.6\\
52000	3306.9\\
54000	3396.1\\
56000	3459.9\\
58000	3507.9\\
62000	3470\\
66000	3394.3\\
68000	3321.4\\
70000	3284.5\\
72000	3380.9\\
74000	3447.2\\
76000	3423.8\\
78000	3464.7\\
80000	3310.3\\
82000	3335.1\\
84000	3414\\
86000	3418.2\\
88000	3456.8\\
90000	3463.4\\
92000	3322.4\\
94000	3348.9\\
96000	3281.5\\
1e+05	3429.7\\
};
\addplot [color=mycolor45, forget plot]
  table[row sep=crcr]{%
0	4395.3\\
2000	5515.8\\
4000	6448\\
8000	6648.6\\
10000	6710.4\\
12000	6873.8\\
14000	6980.2\\
16000	6971.7\\
18000	7015.2\\
20000	6835.5\\
22000	6787.8\\
24000	6765.4\\
26000	6770.9\\
28000	6902.1\\
30000	7081.5\\
32000	6983\\
34000	6829.8\\
36000	6909.2\\
38000	6800.1\\
40000	6908.9\\
48000	7135.3\\
50000	6925.4\\
52000	7078.9\\
54000	6889.6\\
56000	6813.1\\
58000	6833.7\\
60000	6830.9\\
62000	6976.7\\
64000	6936.1\\
66000	7045.5\\
68000	7017.8\\
70000	6884.5\\
72000	6917.6\\
74000	6983.9\\
76000	6970.4\\
78000	7008.2\\
80000	7026.1\\
82000	6901.9\\
84000	6976.3\\
86000	6975.1\\
88000	6910.6\\
90000	6922.9\\
92000	6990.9\\
94000	6826.2\\
96000	6841.5\\
98000	6676.5\\
1e+05	6961.9\\
};
\addplot [color=mycolor46, forget plot]
  table[row sep=crcr]{%
0	4395.3\\
2000	4032.1\\
4000	4013.4\\
6000	4039\\
8000	4097.4\\
10000	4110.7\\
12000	4148.9\\
14000	4111.8\\
16000	4167.8\\
18000	4014.4\\
20000	4091.8\\
22000	4251.3\\
24000	4215.4\\
26000	4065.8\\
28000	4064.8\\
30000	4035.2\\
32000	4213.3\\
36000	4136.2\\
38000	4086.8\\
40000	4203.7\\
42000	4048.9\\
44000	4090.7\\
46000	4050.9\\
48000	4154.5\\
50000	4089.9\\
52000	4111.7\\
54000	4108.6\\
56000	4069\\
58000	4073\\
60000	4131.8\\
62000	4162.6\\
64000	4218.1\\
66000	4071.1\\
68000	4147.5\\
70000	4080.4\\
72000	4078\\
74000	4122.8\\
76000	4229.8\\
78000	4077.8\\
80000	4237.1\\
82000	4134\\
84000	4148.6\\
86000	4145.8\\
88000	4098.1\\
90000	4160.4\\
92000	4116.1\\
94000	4103\\
96000	4109.7\\
98000	4149.1\\
1e+05	4127.8\\
};
\addplot [color=mycolor47, forget plot]
  table[row sep=crcr]{%
0	4395.3\\
2000	5283.8\\
4000	5468.7\\
6000	5387.8\\
8000	5358.3\\
10000	5412.6\\
12000	5245.3\\
14000	5177.5\\
16000	5196.7\\
18000	5262.2\\
20000	5301.4\\
22000	5296.6\\
24000	5184.4\\
26000	5318.7\\
28000	5290.5\\
32000	5443.2\\
34000	5321.5\\
36000	5358.6\\
38000	5202.1\\
40000	5352.2\\
42000	5315.8\\
44000	5192.1\\
46000	5318.9\\
48000	5374.8\\
50000	5234.3\\
52000	5279.7\\
54000	5378.2\\
56000	5455.4\\
58000	5343.2\\
62000	5310.8\\
64000	5203.8\\
68000	5321.8\\
70000	5379.9\\
72000	5195.2\\
76000	5373\\
78000	5303.4\\
80000	5198.8\\
82000	5413.5\\
84000	5237.5\\
86000	5195\\
88000	5290.5\\
90000	5196\\
92000	5239.2\\
94000	5326.5\\
96000	5345.2\\
98000	5256.5\\
1e+05	5257.2\\
};
\addplot [color=mycolor48, forget plot]
  table[row sep=crcr]{%
0	4395.3\\
2000	3635.9\\
4000	3180.1\\
6000	3006.6\\
8000	2962.6\\
10000	2953.5\\
12000	2923.6\\
14000	2985.8\\
16000	2947.8\\
18000	2987\\
20000	2958.5\\
22000	2917\\
24000	2979.3\\
26000	2957.7\\
28000	2983.9\\
30000	2967.6\\
32000	2986.3\\
34000	2968.2\\
36000	2993.3\\
38000	2965.3\\
40000	2998.1\\
42000	2990.2\\
44000	2961.9\\
48000	2988.2\\
50000	2977\\
52000	2947\\
54000	2984.3\\
56000	2985.8\\
58000	2964.4\\
60000	2984.8\\
62000	3039.5\\
64000	2969.6\\
66000	2970.7\\
68000	2993.9\\
70000	2962.3\\
72000	2983.2\\
74000	2965.2\\
76000	2977.1\\
78000	2947.3\\
80000	3019.4\\
82000	2950.4\\
84000	2928\\
88000	2967.4\\
90000	2941.6\\
92000	2985.2\\
96000	2922\\
98000	2946.1\\
1e+05	2983.2\\
};
\addplot [color=mycolor49, forget plot]
  table[row sep=crcr]{%
0	4395.3\\
2000	4922.5\\
4000	4575.6\\
6000	4508.8\\
10000	4407.2\\
12000	4541.4\\
14000	4397\\
16000	4523.5\\
18000	4608.1\\
20000	4345.2\\
22000	4234.8\\
24000	4397.9\\
26000	4321.1\\
28000	4340.1\\
30000	4437.1\\
32000	4435.1\\
34000	4417\\
36000	4335.1\\
38000	4267.1\\
40000	4366.1\\
42000	4441.3\\
44000	4467\\
46000	4433.7\\
48000	4330.6\\
50000	4563.1\\
52000	4267.6\\
54000	4371.5\\
56000	4389.7\\
58000	4499.8\\
60000	4355.3\\
62000	4394.3\\
64000	4262.9\\
66000	4496.1\\
68000	4372.7\\
70000	4282.2\\
72000	4444.8\\
74000	4512.1\\
76000	4400.9\\
78000	4439.9\\
80000	4431.7\\
82000	4365.4\\
84000	4449.2\\
86000	4574.8\\
88000	4483.2\\
90000	4342.8\\
92000	4389.4\\
94000	4473\\
98000	4325.7\\
1e+05	4425.7\\
};
\addplot [color=mycolor50, forget plot]
  table[row sep=crcr]{%
0	4395.3\\
2000	2822.9\\
4000	2440.7\\
6000	2429.2\\
8000	2436.8\\
10000	2515.9\\
12000	2524.2\\
14000	2457.5\\
16000	2487.7\\
20000	2514.9\\
22000	2500.4\\
24000	2496.5\\
26000	2532.8\\
28000	2535.9\\
34000	2576.5\\
38000	2543.1\\
40000	2500.3\\
42000	2556.6\\
44000	2446.2\\
46000	2494.7\\
48000	2446.2\\
50000	2521.5\\
52000	2488.5\\
54000	2497.2\\
56000	2541\\
58000	2525.4\\
60000	2584.8\\
62000	2511.4\\
64000	2515.7\\
66000	2545.9\\
68000	2499.2\\
70000	2544.9\\
72000	2492.4\\
74000	2531.8\\
78000	2559\\
80000	2533.2\\
82000	2499.5\\
84000	2483.3\\
86000	2503.6\\
88000	2536.9\\
90000	2517\\
92000	2550.9\\
94000	2492.2\\
96000	2547.9\\
98000	2515.1\\
1e+05	2555.6\\
};
\addplot [color=mycolor51, forget plot]
  table[row sep=crcr]{%
0	4395.3\\
2000	5587.3\\
4000	5950.1\\
6000	5825\\
8000	5836.3\\
10000	5808.2\\
12000	5821.9\\
14000	5957.4\\
16000	5751.7\\
18000	5923.4\\
22000	6052.5\\
24000	6263.5\\
26000	6060.6\\
28000	6085.6\\
30000	5813\\
32000	5828\\
34000	6153.2\\
36000	6064.1\\
38000	6110.7\\
40000	6046.2\\
42000	6018.7\\
44000	6053\\
48000	5792.2\\
50000	5779\\
52000	6001.8\\
54000	5988.8\\
56000	5846.6\\
58000	5871\\
60000	5918\\
62000	5911.5\\
66000	6148.6\\
68000	6032\\
70000	5982.2\\
72000	6180.2\\
74000	5947.2\\
76000	5590.9\\
78000	5948.1\\
80000	5951.6\\
82000	5784.2\\
84000	5988.6\\
86000	5825.7\\
88000	6074.6\\
90000	5842.2\\
92000	5909.1\\
94000	6040.5\\
96000	5980\\
98000	5877.1\\
1e+05	6072.2\\
};
\addplot [color=mycolor52, forget plot]
  table[row sep=crcr]{%
0	4395.3\\
2000	4385.8\\
4000	4131.6\\
6000	3848.3\\
8000	3822.9\\
10000	3721.7\\
14000	3577\\
16000	3602.6\\
18000	3774.3\\
20000	3621.1\\
22000	3493.5\\
24000	3455.8\\
26000	3530.9\\
28000	3492\\
30000	3556.7\\
32000	3543.7\\
36000	3567.2\\
40000	3629.3\\
42000	3613.7\\
44000	3565.2\\
46000	3389\\
48000	3532.2\\
50000	3643.8\\
52000	3506.2\\
54000	3548.3\\
56000	3492.8\\
58000	3504.4\\
60000	3606.5\\
62000	3624.7\\
64000	3553.2\\
66000	3709.7\\
68000	3672.7\\
70000	3533.8\\
72000	3470.4\\
74000	3525\\
76000	3604.8\\
78000	3575.7\\
80000	3485.4\\
82000	3359.1\\
84000	3499.5\\
86000	3611.1\\
88000	3451.5\\
90000	3577.1\\
92000	3569.1\\
94000	3534.8\\
96000	3533.6\\
98000	3653\\
1e+05	3576.6\\
};
\addplot [color=mycolor53, forget plot]
  table[row sep=crcr]{%
0	4395.3\\
2000	5957.8\\
4000	6487.7\\
6000	6344.4\\
8000	6078\\
10000	5925\\
12000	5912.5\\
14000	5750\\
18000	5604.9\\
20000	5748.3\\
22000	5775.3\\
24000	5886.3\\
26000	5885.3\\
28000	5786.9\\
30000	5842.4\\
32000	5767.2\\
34000	5948.9\\
36000	5733\\
38000	5741\\
40000	5726.3\\
42000	5865.5\\
44000	5646.9\\
46000	5790.2\\
48000	5888\\
50000	5812.3\\
52000	5817\\
54000	5870.6\\
56000	5642.8\\
58000	5566.7\\
60000	5675\\
62000	5715\\
64000	5816.9\\
66000	5832.8\\
68000	5755.6\\
70000	5628.3\\
72000	5562.1\\
74000	5576.5\\
76000	5653.8\\
78000	5500.3\\
80000	5755.6\\
82000	5778.9\\
84000	5657.5\\
86000	5664.5\\
88000	5715.3\\
90000	5706.8\\
92000	5773\\
94000	5813\\
96000	5812.8\\
98000	5578\\
1e+05	5753.4\\
};
\addplot [color=mycolor54, forget plot]
  table[row sep=crcr]{%
0	4395.3\\
2000	4772.4\\
4000	4683.2\\
6000	4619.8\\
8000	4683.1\\
10000	4636.3\\
12000	4744.6\\
14000	4717.6\\
16000	4638.9\\
18000	4586.4\\
20000	4628.6\\
22000	4505\\
26000	4616.7\\
28000	4654.5\\
30000	4647.8\\
32000	4666.6\\
34000	4701.7\\
36000	4513.3\\
38000	4496.4\\
40000	4669.5\\
42000	4596.8\\
44000	4623.3\\
46000	4606.2\\
48000	4697.7\\
50000	4624.2\\
52000	4599.6\\
54000	4682.2\\
56000	4548\\
58000	4521.6\\
60000	4604.1\\
64000	4551.3\\
66000	4611.6\\
68000	4721.4\\
70000	4563.3\\
74000	4651.5\\
76000	4649.3\\
78000	4628.5\\
80000	4631.6\\
82000	4528.7\\
84000	4473.1\\
86000	4616.2\\
88000	4547.9\\
90000	4528.5\\
92000	4564.4\\
94000	4638.2\\
96000	4572.5\\
98000	4688.9\\
1e+05	4539.2\\
};
\addplot [color=mycolor55, forget plot]
  table[row sep=crcr]{%
0	4395.3\\
2000	4932.5\\
4000	5624.5\\
6000	5569.3\\
8000	5637.1\\
10000	5752.5\\
12000	5717\\
14000	5721.1\\
18000	5816\\
20000	5749.9\\
22000	5769\\
24000	5806.8\\
26000	5796.2\\
28000	5732.7\\
30000	5695.6\\
32000	5789.6\\
34000	5834.5\\
38000	5764.3\\
42000	5819.8\\
44000	5767.2\\
48000	5875\\
50000	5757.8\\
52000	5786.2\\
54000	5750.3\\
58000	5892.4\\
60000	5873.4\\
62000	5731.4\\
66000	5961.9\\
68000	5861\\
70000	5735.2\\
72000	5697.3\\
74000	5738\\
76000	5739.8\\
78000	5792.5\\
80000	5870.3\\
82000	5763.8\\
84000	5835.5\\
86000	5809.5\\
90000	5867.9\\
92000	5763.1\\
94000	5772\\
96000	5847.5\\
98000	5808.4\\
1e+05	5742.6\\
};
\addplot [color=mycolor56, forget plot]
  table[row sep=crcr]{%
0	4395.3\\
2000	4516.3\\
4000	4286.2\\
6000	4205.8\\
8000	4028.9\\
12000	4084.8\\
18000	4092.9\\
20000	4079.4\\
22000	4129.2\\
24000	4123.9\\
26000	4082.3\\
28000	4081.4\\
30000	3979.2\\
32000	4003.9\\
34000	4080.2\\
36000	4071.3\\
38000	4039.8\\
40000	4108.2\\
42000	4022.3\\
46000	3994.1\\
48000	4005.2\\
50000	4067.1\\
52000	3965.3\\
54000	3956.3\\
56000	4053.7\\
58000	4098.6\\
60000	3991\\
64000	4080.7\\
66000	4029.7\\
68000	4096.1\\
70000	4131.4\\
72000	4038\\
74000	4061.1\\
76000	4023.6\\
78000	4047.3\\
80000	4093.5\\
82000	4036.1\\
84000	4051.5\\
86000	4102.2\\
88000	4046.3\\
90000	4064.1\\
94000	4032.9\\
98000	4073.8\\
1e+05	4014.3\\
};
\addplot [color=mycolor57, forget plot]
  table[row sep=crcr]{%
0	4395.3\\
2000	5741.2\\
4000	6785.8\\
6000	6914.6\\
8000	6852.6\\
10000	6842.1\\
12000	6650.1\\
14000	6577.9\\
18000	6744.7\\
20000	6751.2\\
22000	6806.4\\
24000	6722.7\\
26000	6615.2\\
28000	6645.3\\
30000	6762.5\\
34000	6663.9\\
36000	6687.8\\
38000	6811.8\\
40000	6772.6\\
42000	6712.5\\
44000	6613.4\\
46000	6704.8\\
48000	6778.5\\
50000	6747.6\\
52000	6843.5\\
54000	6751.2\\
56000	6634\\
58000	6602.6\\
60000	6616.1\\
62000	6580.1\\
64000	6573.1\\
66000	6636.9\\
68000	6523.3\\
70000	6698.7\\
72000	6792.4\\
74000	6622.8\\
76000	6588.1\\
78000	6762.7\\
80000	6606.4\\
82000	6559.6\\
84000	6632.5\\
86000	6634.4\\
88000	6672\\
90000	6806.7\\
92000	6854.2\\
96000	6674.8\\
98000	6708.3\\
1e+05	6691.1\\
};
\addplot [color=mycolor58, forget plot]
  table[row sep=crcr]{%
0	4395.3\\
2000	4166.2\\
6000	3672.5\\
10000	3601.7\\
12000	3645.4\\
14000	3643.6\\
16000	3583.5\\
18000	3621.6\\
20000	3726.1\\
22000	3618.7\\
24000	3588.5\\
26000	3604.1\\
28000	3596.7\\
30000	3661.4\\
32000	3564.1\\
34000	3575.3\\
36000	3629\\
38000	3621.1\\
40000	3592.1\\
42000	3636.7\\
44000	3611.6\\
48000	3660.8\\
50000	3713.2\\
52000	3566.7\\
54000	3614\\
56000	3590.2\\
58000	3666.3\\
60000	3681.8\\
62000	3560\\
64000	3594.5\\
66000	3651.1\\
68000	3557.8\\
70000	3534.4\\
72000	3610.9\\
74000	3571.9\\
76000	3563.4\\
78000	3585.5\\
80000	3626.6\\
82000	3680.8\\
84000	3583.8\\
86000	3555\\
88000	3579.2\\
90000	3794.9\\
94000	3583.5\\
96000	3582.8\\
98000	3611\\
1e+05	3502.8\\
};
\addplot [color=mycolor59, forget plot]
  table[row sep=crcr]{%
0	4395.3\\
2000	4489.4\\
4000	4678.8\\
6000	4894.5\\
8000	4979.6\\
12000	4965.7\\
14000	5006.9\\
16000	4995.9\\
18000	5027.8\\
20000	5090.2\\
22000	5011.6\\
24000	5057.4\\
26000	5132.5\\
28000	5016.7\\
34000	5046.5\\
36000	4958.3\\
38000	5046.4\\
40000	5082.8\\
42000	5052.5\\
44000	5085.5\\
46000	5046.4\\
48000	5125.4\\
50000	5046.5\\
52000	5163.8\\
54000	5074.9\\
56000	5142.9\\
58000	5058.8\\
60000	5124.4\\
64000	5099.2\\
66000	5157.1\\
68000	5095.4\\
70000	5224.3\\
72000	5059.6\\
74000	5137\\
76000	5127.1\\
78000	5070.3\\
80000	5106.2\\
82000	5121.1\\
84000	5200.8\\
86000	5179.1\\
88000	5117.8\\
90000	5121.7\\
92000	5003.3\\
94000	5030.2\\
96000	5080.9\\
98000	5168.2\\
1e+05	5151\\
};
\addplot [color=mycolor60, forget plot]
  table[row sep=crcr]{%
0	4395.3\\
2000	3894.9\\
4000	3722.1\\
6000	3705.7\\
8000	3678.6\\
10000	3596.3\\
12000	3609\\
14000	3602.9\\
16000	3639.4\\
18000	3588.8\\
20000	3604.5\\
22000	3658.5\\
24000	3523.9\\
26000	3556.4\\
28000	3610.5\\
30000	3618\\
32000	3531\\
34000	3573.8\\
36000	3565.1\\
38000	3595.7\\
40000	3541.4\\
42000	3629.8\\
44000	3531.2\\
46000	3527.6\\
48000	3564.7\\
50000	3591.8\\
52000	3601.4\\
54000	3540.7\\
56000	3538.5\\
58000	3619.9\\
60000	3581.1\\
62000	3553\\
64000	3549.6\\
66000	3576\\
68000	3561.9\\
70000	3581.3\\
72000	3543.7\\
74000	3532.5\\
76000	3652.1\\
78000	3550.9\\
80000	3513.4\\
82000	3534.5\\
84000	3487.1\\
86000	3671\\
88000	3662\\
92000	3617.6\\
94000	3547.6\\
96000	3591.1\\
98000	3623.9\\
1e+05	3702.3\\
};
\addplot [color=mycolor61, forget plot]
  table[row sep=crcr]{%
0	4395.3\\
2000	3472.9\\
4000	3202.6\\
6000	3168.2\\
8000	3146.9\\
10000	3110.3\\
12000	3160.8\\
14000	3224.8\\
16000	3164.3\\
18000	3180.5\\
20000	3140.8\\
22000	3163.2\\
24000	3136.8\\
26000	3146.8\\
28000	3204.6\\
32000	3154.2\\
34000	3115.1\\
36000	3153.9\\
38000	3124.9\\
40000	3165.2\\
42000	3177.6\\
44000	3172.4\\
48000	3129.6\\
50000	3187.4\\
52000	3157.6\\
54000	3194.1\\
56000	3187.5\\
58000	3134.6\\
60000	3199.9\\
62000	3203.8\\
66000	3144.8\\
68000	3184.3\\
72000	3189.6\\
76000	3198.8\\
78000	3208.4\\
80000	3230.1\\
82000	3139.8\\
84000	3133.6\\
86000	3205.5\\
88000	3230\\
90000	3169\\
92000	3180.7\\
94000	3165\\
96000	3133.7\\
1e+05	3175.3\\
};
\addplot [color=mycolor62, forget plot]
  table[row sep=crcr]{%
0	4395.3\\
2000	4520.9\\
4000	4875.5\\
6000	5112.5\\
8000	5085.1\\
10000	4995\\
12000	5085.8\\
14000	5140.1\\
16000	5022.1\\
18000	4961.7\\
20000	5094.5\\
22000	5044.6\\
24000	5112.3\\
26000	5003.4\\
28000	5049.4\\
30000	5059.1\\
32000	5013.9\\
36000	4994.4\\
38000	4972.6\\
40000	5013.5\\
42000	5091.2\\
44000	5052.8\\
46000	4990.5\\
48000	5011.3\\
50000	5010.1\\
52000	4927.1\\
54000	5001.2\\
56000	5058\\
58000	4997.7\\
60000	5020.1\\
62000	4986.3\\
64000	4938.6\\
66000	5024.4\\
68000	4961\\
70000	5008.4\\
72000	4971.4\\
74000	4883.9\\
76000	4970.7\\
78000	5036.9\\
80000	5078.8\\
82000	5009.3\\
84000	5026.3\\
86000	4968.6\\
88000	5040.6\\
90000	5034.2\\
92000	5073.3\\
94000	4984\\
98000	5081.2\\
1e+05	5039.5\\
};
\addplot [color=mycolor63, forget plot]
  table[row sep=crcr]{%
0	4395.3\\
2000	3230.2\\
4000	2674.3\\
6000	2532.2\\
8000	2458.2\\
10000	2439.1\\
12000	2470.4\\
14000	2487\\
16000	2441.4\\
18000	2476.2\\
20000	2457.4\\
22000	2458.7\\
24000	2440\\
26000	2457.3\\
28000	2483.8\\
30000	2439.5\\
32000	2415.8\\
34000	2412.6\\
36000	2423.8\\
38000	2450.4\\
40000	2468.5\\
42000	2495.4\\
44000	2470.4\\
46000	2435.6\\
48000	2432.6\\
50000	2469.3\\
52000	2428.1\\
56000	2451.5\\
58000	2429\\
64000	2478.5\\
68000	2448.1\\
70000	2454.6\\
72000	2418.4\\
76000	2499.5\\
78000	2474.9\\
80000	2465.1\\
82000	2485.5\\
84000	2494.3\\
86000	2444.2\\
88000	2473.9\\
90000	2484.6\\
92000	2474.1\\
96000	2489.2\\
98000	2427.1\\
1e+05	2391.7\\
};
\addplot [color=mycolor64, forget plot]
  table[row sep=crcr]{%
0	4395.3\\
2000	3307.9\\
4000	3083.7\\
6000	3088.6\\
8000	3037.6\\
10000	3038.6\\
12000	3116.4\\
16000	3040\\
18000	3075.4\\
20000	3060.1\\
22000	3006.8\\
24000	3057.7\\
26000	3094.6\\
28000	3009\\
30000	3120.7\\
32000	3094.1\\
40000	3096.1\\
42000	3034.7\\
44000	3057.4\\
46000	3114.3\\
48000	3069.8\\
50000	3134.6\\
52000	3112.4\\
54000	3109.5\\
56000	3082\\
58000	3067.8\\
60000	3078.9\\
62000	3060.7\\
66000	3080.4\\
68000	3122.4\\
70000	3077.6\\
74000	3021\\
76000	3067.3\\
78000	3093.7\\
80000	3071\\
82000	3094.8\\
84000	3082.1\\
86000	3102.7\\
88000	3059.5\\
90000	3079.5\\
92000	3109.4\\
94000	3128.5\\
96000	3080.9\\
98000	3082.3\\
1e+05	3151.5\\
};
\addplot [color=mycolor65, forget plot]
  table[row sep=crcr]{%
0	4395.3\\
2000	3154.3\\
4000	2663.5\\
6000	2547.5\\
8000	2564.1\\
10000	2624.8\\
12000	2561.5\\
14000	2567.8\\
16000	2564.9\\
18000	2579.7\\
22000	2577.4\\
24000	2557.5\\
26000	2553.9\\
28000	2598.6\\
30000	2600.1\\
32000	2627\\
34000	2610.3\\
38000	2607.9\\
40000	2607.4\\
42000	2616.1\\
44000	2612.7\\
46000	2581.2\\
48000	2565.8\\
56000	2608.2\\
58000	2561.1\\
62000	2538.5\\
64000	2555.6\\
66000	2599.8\\
68000	2662.5\\
70000	2629.7\\
72000	2572.2\\
74000	2618.4\\
76000	2617.9\\
78000	2603.3\\
80000	2642\\
82000	2574.9\\
84000	2572.7\\
86000	2584.1\\
88000	2608.7\\
90000	2590.7\\
92000	2606.1\\
94000	2612.7\\
96000	2602.9\\
1e+05	2548.4\\
};
\addplot [color=mycolor66, forget plot]
  table[row sep=crcr]{%
0	4395.3\\
2000	3560.9\\
4000	3440\\
8000	3501.7\\
10000	3473\\
12000	3513\\
14000	3381.7\\
16000	3404.4\\
18000	3475.7\\
20000	3485.9\\
22000	3443.3\\
24000	3449.3\\
26000	3427\\
30000	3440.1\\
32000	3413.8\\
34000	3442.2\\
36000	3400.4\\
38000	3464.9\\
40000	3461.8\\
42000	3427.3\\
44000	3485.3\\
46000	3451.4\\
48000	3449.1\\
50000	3436.6\\
52000	3452.1\\
54000	3443.6\\
56000	3456.1\\
58000	3509.7\\
60000	3504.7\\
62000	3464.8\\
68000	3451.3\\
70000	3420.9\\
72000	3466.4\\
74000	3482.6\\
76000	3408.5\\
78000	3443\\
80000	3457.7\\
82000	3455.7\\
84000	3397.9\\
86000	3479.6\\
88000	3460.8\\
90000	3485.6\\
92000	3435.7\\
96000	3462.2\\
98000	3432.6\\
1e+05	3449.4\\
};
\addplot [color=mycolor67, forget plot]
  table[row sep=crcr]{%
0	4395.3\\
2000	3792\\
4000	3471.6\\
6000	3359.3\\
8000	3356\\
12000	3207.9\\
14000	3255.8\\
16000	3240.9\\
20000	3247.4\\
22000	3227.8\\
24000	3235.6\\
28000	3191\\
30000	3209.5\\
32000	3194\\
34000	3223.9\\
36000	3204.2\\
38000	3206.2\\
40000	3260.3\\
42000	3204.1\\
44000	3201.9\\
46000	3226.8\\
50000	3192\\
52000	3273.1\\
54000	3237.1\\
56000	3211\\
58000	3230.2\\
60000	3239.9\\
62000	3237.1\\
64000	3206.2\\
66000	3199.3\\
68000	3267.7\\
70000	3224\\
72000	3198.8\\
76000	3283.4\\
80000	3145.3\\
82000	3250.9\\
84000	3260.6\\
86000	3223.5\\
88000	3223.7\\
90000	3267.5\\
94000	3236.8\\
96000	3243.8\\
98000	3292.6\\
1e+05	3258.9\\
};
\addplot [color=mycolor67]
  table[row sep=crcr]{%
0	4395.3\\
2000	3025.4\\
4000	2688.8\\
6000	2676.6\\
8000	2698.1\\
10000	2654.1\\
12000	2651.2\\
14000	2678.2\\
16000	2654.6\\
18000	2645.9\\
20000	2678.8\\
22000	2652.7\\
24000	2651.8\\
26000	2686.8\\
30000	2648.2\\
32000	2663.1\\
36000	2637.9\\
38000	2676\\
40000	2698.8\\
42000	2712.2\\
46000	2694.5\\
48000	2676.4\\
52000	2657.6\\
54000	2680.1\\
56000	2680.6\\
58000	2695\\
60000	2664\\
62000	2652.8\\
64000	2668.6\\
66000	2706.9\\
68000	2680.7\\
70000	2666.5\\
74000	2684.2\\
76000	2637.1\\
78000	2700.3\\
80000	2640.6\\
82000	2679.9\\
84000	2691.2\\
86000	2667.7\\
88000	2682.6\\
90000	2674.2\\
92000	2658.4\\
94000	2671.9\\
96000	2668.7\\
98000	2697.3\\
1e+05	2717.3\\
};
\addlegendentry{\ClasskIIILegend}

\addplot [color=mycolor68, forget plot]
  table[row sep=crcr]{%
0	4395.3\\
2000	3230.6\\
4000	2728.3\\
6000	2614.9\\
8000	2552.6\\
10000	2513.9\\
12000	2563.5\\
14000	2540.6\\
16000	2545\\
18000	2508.9\\
20000	2524\\
24000	2523.7\\
26000	2500.8\\
28000	2507\\
30000	2523.2\\
34000	2506.7\\
36000	2519.2\\
38000	2480.4\\
40000	2508.3\\
44000	2486.5\\
46000	2511.5\\
50000	2498.7\\
54000	2496.8\\
56000	2511.1\\
60000	2509.7\\
64000	2512\\
66000	2528.7\\
68000	2475.7\\
70000	2483.8\\
74000	2543.6\\
76000	2514.2\\
78000	2515.1\\
82000	2495\\
84000	2505.6\\
86000	2507.1\\
88000	2520.3\\
90000	2504.2\\
92000	2527.5\\
94000	2506.6\\
1e+05	2506.5\\
};
\addplot [color=mycolor69, forget plot]
  table[row sep=crcr]{%
0	4395.3\\
2000	3385.2\\
4000	2988.8\\
6000	2886.5\\
8000	2928.6\\
10000	2907.2\\
12000	2861\\
14000	2909.3\\
16000	2949.5\\
18000	2918.5\\
20000	2916.9\\
22000	2939.1\\
24000	2932.4\\
28000	2883.4\\
30000	2922.9\\
32000	2945.3\\
34000	2936.2\\
36000	2914.3\\
38000	2932.9\\
40000	2915.7\\
42000	2946.1\\
44000	2957\\
46000	2935.8\\
50000	2958.7\\
52000	2903.3\\
54000	2926.3\\
56000	2934.7\\
58000	2968.1\\
60000	2947.6\\
62000	2953.8\\
64000	2977.3\\
66000	2941.5\\
68000	2926.9\\
70000	2981.2\\
74000	2950.6\\
76000	2915.3\\
78000	2933.6\\
80000	2971.3\\
82000	2957.6\\
84000	2952.7\\
86000	2977.1\\
90000	2994.8\\
92000	2955.6\\
94000	2970.6\\
96000	2941.3\\
98000	2923.5\\
1e+05	2939.8\\
};
\addplot [color=mycolor70, forget plot]
  table[row sep=crcr]{%
0	4395.3\\
2000	2755.5\\
4000	2120.4\\
6000	2089.9\\
12000	2146.6\\
18000	2121.8\\
20000	2152.7\\
22000	2161.8\\
24000	2181.2\\
26000	2149.3\\
28000	2158.3\\
30000	2178.9\\
32000	2152.2\\
36000	2164.3\\
40000	2157.8\\
42000	2175.9\\
44000	2147.4\\
46000	2179.9\\
48000	2143.7\\
52000	2185.6\\
54000	2187.8\\
56000	2167.8\\
58000	2182.6\\
60000	2161.6\\
62000	2128.1\\
64000	2120.1\\
66000	2156\\
68000	2155\\
70000	2147\\
72000	2184\\
74000	2168.4\\
76000	2145.2\\
78000	2167.7\\
80000	2162.2\\
82000	2173.6\\
84000	2151.4\\
86000	2183.2\\
88000	2132.8\\
90000	2139\\
92000	2137.6\\
94000	2156.4\\
96000	2187.8\\
98000	2186.8\\
1e+05	2168.2\\
};
\addplot [color=mycolor71, forget plot]
  table[row sep=crcr]{%
0	4395.3\\
2000	3106.5\\
4000	2674.8\\
8000	2596.2\\
10000	2601\\
12000	2563\\
14000	2586.2\\
16000	2586.8\\
18000	2550.3\\
20000	2532.5\\
22000	2554.5\\
24000	2533.5\\
26000	2552.4\\
28000	2529.4\\
32000	2537.7\\
34000	2515.7\\
36000	2526.3\\
38000	2544.4\\
40000	2530.3\\
42000	2497.1\\
44000	2494.6\\
46000	2507.3\\
48000	2485.5\\
52000	2499.2\\
54000	2530\\
56000	2543.7\\
58000	2503.8\\
60000	2539.9\\
62000	2502.4\\
64000	2556.5\\
66000	2532.8\\
68000	2567.3\\
70000	2565\\
72000	2547.5\\
76000	2549.4\\
78000	2529.6\\
82000	2530.1\\
84000	2552.1\\
86000	2525.8\\
88000	2512.4\\
90000	2550.2\\
92000	2534.2\\
94000	2529.9\\
96000	2518.5\\
98000	2526.5\\
1e+05	2515.6\\
};
\addplot [color=mycolor72, forget plot]
  table[row sep=crcr]{%
0	4395.3\\
2000	3048.7\\
4000	2375.7\\
6000	2260.6\\
8000	2256.4\\
10000	2273.5\\
12000	2262.2\\
14000	2223.4\\
16000	2229.3\\
18000	2188.8\\
20000	2244.4\\
22000	2249.8\\
24000	2218.4\\
26000	2246.8\\
28000	2250.8\\
30000	2227.8\\
32000	2238.8\\
34000	2220\\
36000	2233.3\\
38000	2216.6\\
40000	2271.3\\
42000	2232.8\\
44000	2244.2\\
46000	2224.1\\
48000	2299.2\\
54000	2247.6\\
56000	2250.3\\
58000	2243.9\\
60000	2186\\
62000	2249.9\\
64000	2241.6\\
66000	2257.1\\
68000	2230.5\\
70000	2223.4\\
72000	2237\\
74000	2199\\
76000	2248.7\\
82000	2226.6\\
84000	2209.7\\
86000	2211\\
88000	2279.4\\
90000	2292.7\\
92000	2225\\
94000	2265.2\\
96000	2218.9\\
1e+05	2240.9\\
};
\addplot [color=mycolor73, forget plot]
  table[row sep=crcr]{%
0	4395.3\\
2000	3210.2\\
4000	2745.3\\
8000	2709.6\\
10000	2693\\
22000	2645.1\\
24000	2678.7\\
26000	2682.7\\
30000	2642.9\\
32000	2635.1\\
36000	2632.2\\
38000	2687.9\\
40000	2660.9\\
42000	2654.1\\
44000	2688.7\\
46000	2662.8\\
48000	2651.4\\
50000	2613.7\\
52000	2648.5\\
54000	2651.8\\
56000	2618.5\\
60000	2660.7\\
64000	2634.4\\
66000	2604.8\\
68000	2638.4\\
70000	2653.9\\
72000	2655.9\\
74000	2615.3\\
76000	2677.6\\
78000	2653.6\\
82000	2647.6\\
84000	2658.9\\
86000	2603.9\\
88000	2612.7\\
90000	2629.6\\
92000	2630.4\\
94000	2642.4\\
96000	2642.7\\
98000	2664.6\\
1e+05	2651.6\\
};
\addplot [color=mycolor74, forget plot]
  table[row sep=crcr]{%
0	4395.3\\
2000	2606.5\\
4000	1759\\
6000	1622.9\\
8000	1575.2\\
12000	1578\\
14000	1570.9\\
16000	1588.5\\
18000	1563.5\\
20000	1576.2\\
22000	1569.1\\
24000	1577.1\\
26000	1572.9\\
28000	1593.1\\
30000	1592.4\\
32000	1569.8\\
34000	1572\\
36000	1583.1\\
38000	1577\\
42000	1589\\
44000	1604.4\\
46000	1556.8\\
48000	1573.7\\
50000	1607.4\\
52000	1584.8\\
54000	1577.9\\
56000	1557\\
58000	1584.1\\
60000	1578\\
62000	1579.1\\
64000	1594.4\\
66000	1595.6\\
68000	1570\\
70000	1561.6\\
72000	1576.4\\
74000	1582.9\\
76000	1598.5\\
78000	1580.2\\
80000	1576.6\\
82000	1561.1\\
84000	1565.9\\
86000	1554.5\\
88000	1587.6\\
90000	1587.8\\
92000	1558.1\\
94000	1559.2\\
98000	1571.4\\
1e+05	1574.6\\
};
\addplot [color=mycolor75, forget plot]
  table[row sep=crcr]{%
0	4395.3\\
2000	3064.4\\
4000	2401.8\\
6000	2302.2\\
10000	2276.5\\
12000	2279.5\\
14000	2253.3\\
18000	2283\\
22000	2240.1\\
24000	2244.6\\
26000	2234.3\\
28000	2274.3\\
30000	2290.2\\
32000	2279\\
34000	2236.5\\
36000	2220.9\\
40000	2281.8\\
42000	2286.4\\
44000	2270.1\\
46000	2278\\
48000	2264\\
50000	2237.1\\
56000	2261.5\\
58000	2246.1\\
62000	2241.6\\
64000	2264.1\\
68000	2241.6\\
70000	2257.5\\
72000	2261\\
74000	2273.9\\
76000	2262.5\\
78000	2275.3\\
82000	2248.7\\
84000	2258.8\\
86000	2281.1\\
88000	2265.5\\
90000	2270.8\\
92000	2232.5\\
94000	2265.2\\
98000	2242.4\\
1e+05	2271.8\\
};
\addplot [color=mycolor76, forget plot]
  table[row sep=crcr]{%
0	4395.3\\
2000	2869.9\\
4000	2220.8\\
6000	2137.6\\
10000	2079.7\\
12000	2062.4\\
16000	2072\\
18000	2081.7\\
20000	2069.2\\
22000	2091\\
24000	2087\\
26000	2063.7\\
28000	2064.3\\
30000	2081.5\\
32000	2053.9\\
34000	2067.1\\
36000	2070.9\\
38000	2045.7\\
40000	2062.3\\
42000	2072.3\\
44000	2061.1\\
46000	2035.8\\
50000	2056.5\\
52000	2038.7\\
54000	2066.4\\
56000	2088.5\\
58000	2074.8\\
60000	2069.8\\
62000	2078.3\\
64000	2077.2\\
66000	2050.8\\
68000	2069\\
70000	2060.8\\
72000	2074\\
74000	2069.9\\
76000	2074.9\\
78000	2040.2\\
80000	2081.5\\
82000	2061.3\\
86000	2059.2\\
88000	2074.7\\
90000	2056.1\\
92000	2064.3\\
94000	2053.5\\
96000	2054.5\\
1e+05	2068.4\\
};
\addplot [color=mycolor77, forget plot]
  table[row sep=crcr]{%
0	4395.3\\
2000	2706.4\\
4000	1850.7\\
6000	1705.7\\
8000	1663\\
10000	1649.7\\
14000	1636.6\\
16000	1639.4\\
18000	1624.1\\
20000	1616.9\\
22000	1627.5\\
24000	1608.8\\
26000	1620.3\\
28000	1603.3\\
30000	1623.9\\
32000	1625.3\\
34000	1619.2\\
40000	1618.1\\
42000	1608.9\\
44000	1619.9\\
46000	1598.4\\
48000	1601\\
50000	1596.5\\
52000	1617.8\\
54000	1618.5\\
56000	1600.5\\
58000	1607.1\\
60000	1619.1\\
62000	1617.4\\
64000	1601.8\\
66000	1609.5\\
68000	1622.9\\
72000	1614.5\\
74000	1589\\
76000	1598.8\\
78000	1613.7\\
80000	1603\\
84000	1611.4\\
86000	1595.5\\
90000	1619.8\\
92000	1601.9\\
94000	1608\\
96000	1603.3\\
98000	1609\\
1e+05	1592\\
};
\addplot [color=mycolor78, forget plot]
  table[row sep=crcr]{%
0	4395.3\\
2000	3219.7\\
4000	2558.2\\
6000	2459.2\\
8000	2391.7\\
10000	2358.1\\
14000	2367.3\\
16000	2387.3\\
18000	2341.1\\
20000	2351.8\\
22000	2375.2\\
24000	2381.3\\
26000	2339.9\\
28000	2336.2\\
30000	2346\\
32000	2297.9\\
34000	2348.7\\
36000	2320.1\\
38000	2356\\
40000	2384.5\\
42000	2353.1\\
44000	2354.2\\
46000	2332.6\\
48000	2345.5\\
50000	2314.7\\
54000	2367.1\\
56000	2353.9\\
58000	2360.7\\
60000	2310.5\\
62000	2323.5\\
64000	2368.5\\
68000	2306.9\\
70000	2316.1\\
74000	2352.2\\
76000	2330.3\\
78000	2346\\
80000	2323.3\\
82000	2370.6\\
84000	2350.2\\
86000	2351.6\\
88000	2327.7\\
90000	2348\\
92000	2350.4\\
96000	2312.4\\
98000	2346.9\\
1e+05	2334.7\\
};
\addplot [color=mycolor79, forget plot]
  table[row sep=crcr]{%
0	4395.3\\
2000	2978.6\\
4000	2407.2\\
6000	2336.7\\
8000	2250.6\\
10000	2279.4\\
12000	2231.4\\
14000	2256.9\\
16000	2205.7\\
18000	2188.4\\
20000	2240.7\\
22000	2231.5\\
24000	2206.6\\
26000	2172.6\\
28000	2232.4\\
32000	2210.5\\
34000	2240.7\\
38000	2190.2\\
40000	2203.3\\
44000	2172.5\\
46000	2219.6\\
48000	2223\\
50000	2232.6\\
52000	2221.6\\
56000	2179.9\\
58000	2190\\
60000	2239.9\\
62000	2223.3\\
64000	2198.3\\
66000	2182.8\\
68000	2191.3\\
70000	2167.2\\
72000	2197.7\\
74000	2180.7\\
76000	2233.1\\
80000	2170.3\\
82000	2164.1\\
86000	2207.7\\
88000	2230.9\\
90000	2237.8\\
92000	2202.9\\
94000	2216.3\\
96000	2193.5\\
98000	2201.6\\
1e+05	2182.2\\
};
\addplot [color=mycolor79, forget plot]
  table[row sep=crcr]{%
0	4395.3\\
2000	3028.3\\
4000	2301.6\\
6000	2143.2\\
10000	2083.2\\
12000	2084.9\\
14000	2060.6\\
16000	2087.6\\
18000	2048\\
22000	2087.2\\
24000	2059.7\\
26000	2089.3\\
28000	2064.8\\
30000	2023.3\\
32000	2088.1\\
36000	2038.3\\
38000	2075\\
40000	2073.8\\
42000	2060.6\\
44000	2031.8\\
46000	2049.7\\
48000	2081.3\\
50000	2056.9\\
52000	2057.9\\
54000	2069.1\\
56000	2052.9\\
58000	2068.7\\
60000	2054\\
62000	2032.8\\
64000	2082.1\\
66000	2083.9\\
68000	2055.4\\
72000	2068.1\\
76000	2037.1\\
78000	2061\\
80000	2061.8\\
82000	2034.9\\
84000	2052.6\\
86000	2042.2\\
88000	2055.6\\
90000	2056.6\\
92000	2043.6\\
96000	2083.1\\
98000	2082.4\\
1e+05	2061.3\\
};
\addplot [color=mycolor80, forget plot]
  table[row sep=crcr]{%
0	4395.3\\
2000	2606.7\\
4000	1742.1\\
6000	1656.6\\
8000	1636.2\\
10000	1627.2\\
12000	1634.4\\
16000	1600.2\\
18000	1607.5\\
22000	1601.8\\
24000	1618.6\\
26000	1600.2\\
28000	1577\\
30000	1592.9\\
32000	1577.4\\
34000	1595\\
38000	1600.8\\
40000	1601\\
42000	1594.3\\
44000	1606\\
46000	1623.7\\
52000	1602.6\\
54000	1609.5\\
56000	1591.8\\
58000	1606.2\\
60000	1595.1\\
62000	1597\\
64000	1591.3\\
66000	1596.3\\
70000	1586.8\\
72000	1597.7\\
74000	1577.7\\
76000	1591.4\\
80000	1602.5\\
82000	1578.4\\
86000	1611.7\\
90000	1596.7\\
94000	1628.2\\
96000	1613.9\\
98000	1585\\
1e+05	1593.3\\
};
\addplot [color=mycolor81, forget plot]
  table[row sep=crcr]{%
0	4395.3\\
2000	2592.6\\
4000	1698.3\\
6000	1563.4\\
8000	1503.8\\
10000	1519.2\\
12000	1523.3\\
14000	1504.5\\
16000	1474.7\\
18000	1493.8\\
20000	1500.1\\
22000	1521.9\\
24000	1513\\
26000	1509.7\\
30000	1525.1\\
32000	1506.6\\
34000	1508.4\\
36000	1496.8\\
38000	1502.8\\
40000	1513.6\\
42000	1493.4\\
44000	1513.1\\
46000	1498.7\\
48000	1508.8\\
50000	1493.7\\
52000	1508.3\\
54000	1493.7\\
60000	1500.9\\
62000	1513.4\\
64000	1480.7\\
66000	1495.9\\
68000	1479.1\\
70000	1478.2\\
72000	1494.6\\
76000	1517.9\\
78000	1492.4\\
80000	1482.8\\
82000	1494.5\\
84000	1515.6\\
86000	1490.5\\
88000	1504\\
90000	1483.3\\
92000	1501.5\\
94000	1488.4\\
98000	1488.3\\
1e+05	1475.3\\
};
\addplot [color=mycolor82, forget plot]
  table[row sep=crcr]{%
0	4395.3\\
2000	2841.5\\
4000	2010.5\\
6000	1863.5\\
8000	1811.7\\
10000	1789.3\\
12000	1775.6\\
14000	1783.8\\
16000	1746.5\\
18000	1758\\
20000	1713.7\\
22000	1735.3\\
24000	1749\\
26000	1727\\
28000	1712.7\\
30000	1748.6\\
32000	1735.6\\
34000	1768.1\\
36000	1737.3\\
38000	1742.8\\
40000	1736.6\\
44000	1740.4\\
46000	1767.3\\
48000	1769.6\\
50000	1763.7\\
52000	1752.7\\
54000	1727.3\\
56000	1743.1\\
58000	1747.2\\
64000	1739.9\\
66000	1756.2\\
68000	1725.5\\
70000	1734.9\\
72000	1756.7\\
74000	1769.3\\
76000	1749.9\\
78000	1752.3\\
82000	1741.4\\
84000	1731.7\\
86000	1735.3\\
88000	1745.8\\
90000	1748.3\\
92000	1732.2\\
94000	1725.9\\
96000	1734.3\\
98000	1723.8\\
1e+05	1749.2\\
};
\addplot [color=mycolor83, forget plot]
  table[row sep=crcr]{%
0	4395.3\\
2000	2836.1\\
4000	2026.3\\
6000	1874.9\\
8000	1819.7\\
10000	1829.9\\
12000	1805.4\\
14000	1790.7\\
18000	1793.3\\
20000	1781\\
22000	1790\\
24000	1789.7\\
26000	1810.4\\
30000	1781.5\\
32000	1794.4\\
34000	1758.7\\
36000	1771.4\\
38000	1776.4\\
40000	1758.4\\
42000	1760.2\\
44000	1747.1\\
46000	1753.4\\
48000	1784\\
50000	1754.3\\
52000	1771.6\\
54000	1776.2\\
56000	1758.3\\
60000	1751.7\\
62000	1765.7\\
64000	1756.6\\
66000	1774.6\\
70000	1741\\
76000	1766.4\\
78000	1763.7\\
80000	1783.8\\
82000	1778\\
84000	1765.2\\
86000	1790.8\\
88000	1774\\
90000	1778.4\\
92000	1790.3\\
94000	1750.6\\
96000	1755.4\\
98000	1777.1\\
1e+05	1761.7\\
};
\addplot [color=mycolor83, forget plot]
  table[row sep=crcr]{%
0	4395.3\\
2000	2438\\
4000	1456.3\\
6000	1327\\
8000	1308.9\\
10000	1303.8\\
12000	1306\\
14000	1322.1\\
16000	1320.6\\
18000	1309.4\\
26000	1312.5\\
28000	1324.1\\
30000	1294.7\\
34000	1298.9\\
36000	1315.2\\
38000	1309.7\\
40000	1309.3\\
42000	1286.1\\
46000	1316.9\\
48000	1319.7\\
50000	1300\\
52000	1309.6\\
62000	1290.2\\
64000	1307.8\\
66000	1304.8\\
68000	1306.9\\
70000	1314.3\\
74000	1291.4\\
76000	1286.7\\
78000	1289.8\\
80000	1302.3\\
82000	1301.6\\
84000	1293.2\\
86000	1290.5\\
88000	1304.1\\
90000	1301.4\\
92000	1321.6\\
94000	1316.1\\
96000	1300.7\\
1e+05	1286.5\\
};
\addplot [color=mycolor84, forget plot]
  table[row sep=crcr]{%
0	4395.3\\
2000	2543\\
4000	1541.3\\
6000	1376.6\\
8000	1323.5\\
10000	1330.6\\
12000	1323.2\\
14000	1309.1\\
16000	1327.1\\
18000	1300.6\\
22000	1324.4\\
26000	1313.4\\
28000	1307.4\\
30000	1341.8\\
32000	1336\\
34000	1296.6\\
38000	1308.8\\
40000	1330\\
42000	1318.5\\
44000	1332.5\\
46000	1321.7\\
48000	1301.6\\
50000	1318.8\\
52000	1311.9\\
54000	1293.7\\
58000	1325.4\\
60000	1322.4\\
62000	1310.5\\
64000	1321.7\\
66000	1325.2\\
68000	1333.6\\
70000	1317.4\\
72000	1339\\
74000	1316.3\\
76000	1337.5\\
78000	1339.3\\
80000	1348.1\\
82000	1326.5\\
84000	1317.3\\
86000	1316.7\\
88000	1336\\
90000	1305.5\\
94000	1309.3\\
96000	1338.9\\
98000	1294.3\\
1e+05	1339.9\\
};
\addplot [color=mycolor85, forget plot]
  table[row sep=crcr]{%
0	4395.3\\
2000	2682.5\\
4000	1705.6\\
6000	1562.7\\
8000	1510.8\\
10000	1471\\
14000	1442.2\\
16000	1434.7\\
20000	1465.3\\
22000	1441.8\\
24000	1448\\
26000	1431.7\\
28000	1446.7\\
30000	1431.8\\
32000	1440.8\\
34000	1428.8\\
36000	1428.4\\
38000	1432.3\\
40000	1411.8\\
42000	1415.6\\
44000	1453.8\\
46000	1420.3\\
48000	1434.6\\
52000	1427.6\\
56000	1433.4\\
58000	1429.4\\
60000	1447.6\\
62000	1447.7\\
64000	1460.4\\
66000	1431.8\\
68000	1434.9\\
70000	1430.2\\
72000	1415\\
74000	1426\\
76000	1411.4\\
78000	1442.1\\
80000	1455.1\\
82000	1453.3\\
84000	1469.1\\
86000	1428\\
88000	1411.2\\
90000	1425.7\\
92000	1405.8\\
94000	1428.8\\
96000	1446.6\\
98000	1435\\
1e+05	1416.9\\
};
\addplot [color=mycolor86, forget plot]
  table[row sep=crcr]{%
0	4395.3\\
2000	2536.5\\
4000	1485.2\\
6000	1323.1\\
8000	1264.6\\
10000	1264.3\\
12000	1247.1\\
14000	1212.4\\
16000	1225.1\\
18000	1252.5\\
20000	1216.2\\
22000	1222.8\\
26000	1219.4\\
28000	1211.8\\
30000	1227.4\\
32000	1214.6\\
34000	1227\\
36000	1213.8\\
38000	1222.1\\
40000	1210.5\\
42000	1205.3\\
44000	1219.2\\
46000	1225.5\\
48000	1227.9\\
50000	1199.1\\
52000	1189.4\\
54000	1208.1\\
56000	1188.7\\
58000	1195.9\\
60000	1192.1\\
68000	1221.3\\
70000	1219\\
74000	1228.3\\
78000	1211.4\\
80000	1219.6\\
82000	1222.6\\
84000	1240.4\\
88000	1221.4\\
90000	1233.1\\
92000	1222.8\\
94000	1223\\
96000	1239.2\\
98000	1212.6\\
1e+05	1226.3\\
};
\addplot [color=mycolor87, forget plot]
  table[row sep=crcr]{%
0	4395.3\\
2000	2347.9\\
4000	1169.3\\
6000	1030.6\\
8000	1009.4\\
10000	1045.9\\
12000	1039.2\\
14000	1021.6\\
16000	1042.1\\
18000	1022.7\\
20000	1026.2\\
22000	1047.6\\
24000	1060.5\\
26000	1018\\
30000	1020.9\\
32000	1016.4\\
34000	1017.1\\
36000	1009.6\\
38000	986.43\\
40000	1017\\
42000	1082.4\\
44000	1058.8\\
46000	1039.1\\
48000	1027.7\\
50000	1033.5\\
52000	1053.4\\
54000	1023.4\\
56000	1010.7\\
58000	1041.6\\
60000	1053.6\\
62000	1030.3\\
64000	1057.3\\
66000	1039.5\\
68000	990.03\\
70000	979.37\\
72000	1009.3\\
74000	997.63\\
76000	1030.5\\
78000	1048.4\\
80000	1043.5\\
82000	1033.7\\
84000	1052.6\\
86000	1017.6\\
88000	994.88\\
90000	1042.5\\
92000	1037\\
94000	1064.1\\
96000	1050.4\\
98000	1052.1\\
1e+05	1061.3\\
};
\addplot [color=mycolor88, forget plot]
  table[row sep=crcr]{%
0	4395.3\\
2000	2594.7\\
4000	1526.5\\
6000	1336\\
8000	1286.4\\
10000	1259.9\\
12000	1245.2\\
14000	1221\\
16000	1234.6\\
18000	1226.5\\
20000	1253.7\\
22000	1224.1\\
24000	1217.7\\
30000	1244.7\\
34000	1226.3\\
36000	1228.5\\
38000	1211.4\\
40000	1218\\
42000	1196.3\\
44000	1214.3\\
46000	1241.3\\
48000	1216.6\\
50000	1202.7\\
52000	1212.9\\
58000	1207.8\\
62000	1222.7\\
64000	1200.9\\
66000	1193.7\\
68000	1197.7\\
70000	1207.5\\
72000	1201.3\\
74000	1218.9\\
76000	1249.3\\
78000	1221.9\\
80000	1208.3\\
82000	1213\\
86000	1208.2\\
88000	1210.2\\
92000	1227.8\\
94000	1219.7\\
96000	1236.3\\
98000	1245.4\\
1e+05	1208.2\\
};
\addplot [color=mycolor89, forget plot]
  table[row sep=crcr]{%
0	4395.3\\
2000	2495.5\\
4000	1249.6\\
6000	971.09\\
8000	920.08\\
10000	893.81\\
14000	903.67\\
20000	874.62\\
22000	886.34\\
24000	891.54\\
26000	904.22\\
28000	884.87\\
30000	857\\
32000	886.78\\
34000	900.27\\
36000	909\\
38000	922.34\\
40000	897.92\\
42000	884.42\\
44000	879.57\\
46000	890.52\\
48000	894.89\\
50000	893.53\\
52000	902.05\\
54000	871.97\\
56000	864.27\\
58000	874.91\\
60000	876.85\\
62000	888.1\\
64000	880.46\\
66000	886.64\\
72000	874.05\\
74000	879.37\\
76000	890.17\\
78000	909.73\\
80000	883.5\\
84000	904.39\\
86000	915.75\\
88000	897.88\\
90000	890.12\\
92000	890\\
94000	892.7\\
96000	887.74\\
98000	890.91\\
1e+05	868.57\\
};
\addplot [color=mycolor90, forget plot]
  table[row sep=crcr]{%
0	4395.3\\
2000	2304\\
4000	1083.3\\
6000	883.55\\
8000	872.81\\
10000	873.46\\
12000	876.92\\
14000	848.83\\
16000	859.16\\
22000	865.03\\
24000	874\\
26000	871.07\\
28000	865.53\\
30000	851.81\\
32000	846.11\\
34000	861.4\\
36000	865.35\\
38000	848.12\\
40000	849.31\\
42000	867.03\\
44000	842.5\\
46000	852.85\\
48000	852.44\\
50000	840.14\\
52000	838.11\\
54000	855.19\\
56000	865.06\\
60000	856.8\\
62000	868.18\\
66000	858.28\\
68000	833.07\\
70000	837.28\\
72000	838.11\\
74000	854.84\\
76000	837.56\\
78000	831.86\\
80000	837.52\\
82000	857.97\\
84000	869.78\\
86000	844.82\\
88000	838.2\\
90000	842.34\\
92000	856.53\\
94000	867.01\\
96000	839.1\\
1e+05	852.89\\
};
\addplot [color=mycolor91, forget plot]
  table[row sep=crcr]{%
0	4395.3\\
2000	2500.7\\
4000	1154.6\\
6000	860.27\\
8000	779.17\\
10000	741.02\\
12000	739.18\\
14000	765.23\\
16000	748.67\\
18000	718\\
20000	723.54\\
22000	746.85\\
24000	735.56\\
26000	733.53\\
28000	749.65\\
32000	730.68\\
36000	752.88\\
40000	745.9\\
44000	720.49\\
46000	733.16\\
48000	731.87\\
50000	743.69\\
52000	744.39\\
54000	748\\
56000	753.83\\
58000	750.04\\
60000	718.75\\
62000	749.43\\
66000	737.61\\
68000	750.22\\
70000	730.92\\
72000	729.1\\
74000	742.56\\
76000	728.65\\
78000	740.9\\
82000	738.02\\
84000	744.15\\
86000	728.22\\
88000	723.1\\
90000	739.97\\
92000	759.48\\
94000	749.02\\
96000	732.22\\
98000	756.02\\
1e+05	740.4\\
};
\addplot [color=mycolor92, forget plot]
  table[row sep=crcr]{%
0	4395.3\\
2000	2494.9\\
4000	1302.3\\
6000	1042.6\\
8000	959.19\\
10000	962.61\\
12000	958.93\\
14000	938.61\\
18000	923.56\\
20000	919.91\\
22000	931.77\\
24000	953.27\\
26000	942.56\\
28000	927.29\\
30000	906.12\\
34000	926.59\\
36000	919.98\\
38000	927.52\\
40000	913.59\\
42000	939.24\\
44000	926.37\\
46000	918.18\\
48000	934.13\\
50000	912.87\\
52000	916.35\\
54000	938.32\\
56000	924.76\\
58000	890.99\\
60000	892.24\\
62000	923.02\\
64000	928.37\\
68000	899.77\\
70000	917.41\\
72000	908.82\\
74000	943.16\\
76000	917.86\\
78000	944.68\\
80000	937.13\\
82000	947.02\\
84000	934.4\\
88000	927.84\\
90000	931.49\\
92000	929.12\\
94000	905.39\\
96000	900.97\\
98000	903.81\\
1e+05	914.95\\
};
\addplot [color=mycolor93, forget plot]
  table[row sep=crcr]{%
0	4395.3\\
2000	2573.5\\
4000	1234.9\\
6000	901.67\\
8000	835.17\\
10000	808.34\\
12000	803.57\\
14000	795.06\\
16000	799.24\\
18000	775.72\\
20000	777.02\\
22000	786.94\\
24000	808.47\\
26000	802.9\\
30000	771.56\\
32000	780.41\\
34000	784.21\\
36000	765.93\\
38000	779.51\\
40000	784.85\\
42000	780.31\\
44000	803.33\\
46000	786.46\\
48000	767.52\\
50000	792.83\\
52000	780.57\\
56000	781.97\\
58000	763.16\\
60000	748.24\\
62000	759.05\\
64000	755.08\\
66000	770.22\\
68000	774.89\\
70000	787.07\\
72000	764.24\\
74000	797.06\\
76000	775.59\\
78000	791.87\\
80000	776.93\\
82000	770.69\\
84000	757.53\\
86000	761.7\\
88000	734.03\\
90000	784.88\\
94000	818.02\\
96000	798.8\\
98000	754.63\\
1e+05	800.75\\
};
\addplot [color=mycolor94, forget plot]
  table[row sep=crcr]{%
0	4395.3\\
2000	2651.7\\
4000	1274\\
6000	904.38\\
8000	800.45\\
10000	750.32\\
12000	736.68\\
14000	744.27\\
16000	736.58\\
20000	740.39\\
22000	735.07\\
24000	721.93\\
26000	726.86\\
28000	734.3\\
30000	739.04\\
32000	725.8\\
34000	722.78\\
36000	728.13\\
38000	738.41\\
40000	720.86\\
42000	728.27\\
44000	738.21\\
46000	724.54\\
48000	732.25\\
50000	725.3\\
52000	756.51\\
54000	738.2\\
56000	741.81\\
58000	732.33\\
60000	740\\
62000	733.85\\
68000	752.13\\
70000	730.87\\
72000	746.11\\
74000	750.16\\
76000	742.31\\
78000	744.07\\
80000	741.53\\
82000	733.85\\
86000	744.11\\
90000	741.62\\
92000	756.89\\
94000	762.76\\
96000	749.02\\
98000	741.73\\
1e+05	718.77\\
};
\addplot [color=mycolor95, forget plot]
  table[row sep=crcr]{%
0	4395.3\\
2000	2343.5\\
4000	878.28\\
6000	578.02\\
8000	488.73\\
10000	461.23\\
12000	461.54\\
14000	443.09\\
16000	447.76\\
18000	442.56\\
20000	452.13\\
22000	441.29\\
24000	451.95\\
26000	458.38\\
28000	449.88\\
30000	442.86\\
32000	467.92\\
34000	460.76\\
36000	442.12\\
40000	454.24\\
42000	448.6\\
44000	453.2\\
46000	460.27\\
48000	458.51\\
50000	441.91\\
52000	455.49\\
54000	438.88\\
56000	441.88\\
58000	457.72\\
60000	449.52\\
62000	435.43\\
64000	438.6\\
66000	435.97\\
68000	439.64\\
72000	443.13\\
74000	454.89\\
76000	453.79\\
78000	446.93\\
80000	433.12\\
84000	431.53\\
86000	447.25\\
88000	456.47\\
90000	451.05\\
92000	444.35\\
94000	433.5\\
96000	428.44\\
1e+05	449.63\\
};
\addplot [color=mycolor96, forget plot]
  table[row sep=crcr]{%
0	4395.3\\
2000	2467.8\\
4000	1007.8\\
6000	644.74\\
8000	529.93\\
10000	495.28\\
12000	488.9\\
14000	476.51\\
16000	477.73\\
18000	458.19\\
20000	462.2\\
22000	499.95\\
24000	482.05\\
26000	478.38\\
28000	460.3\\
30000	476.34\\
32000	463.72\\
34000	449.19\\
36000	482.54\\
38000	505.86\\
40000	496.92\\
44000	459\\
46000	457.82\\
48000	463.23\\
50000	459.33\\
52000	469.57\\
54000	462.64\\
56000	468.64\\
58000	471.19\\
60000	470.27\\
62000	460.67\\
64000	471.7\\
66000	485.73\\
68000	484.96\\
70000	468.21\\
72000	465.23\\
74000	473.25\\
76000	479.28\\
78000	480.6\\
80000	453.75\\
82000	488.66\\
86000	483.5\\
88000	475.04\\
90000	443.42\\
92000	469.14\\
94000	492.54\\
96000	497.2\\
98000	484.78\\
1e+05	452.66\\
};
\addplot [color=mycolor97, forget plot]
  table[row sep=crcr]{%
0	4395.3\\
2000	2176.6\\
4000	759.49\\
6000	466.66\\
8000	396.97\\
10000	361.81\\
12000	348.28\\
14000	364.45\\
16000	347.64\\
18000	335.35\\
20000	322.56\\
22000	319.52\\
24000	318.39\\
26000	344.44\\
28000	356.64\\
30000	346.47\\
32000	320.68\\
34000	327.43\\
36000	329.81\\
38000	338.96\\
40000	335.46\\
42000	348.94\\
44000	342.58\\
46000	354.29\\
48000	346.55\\
50000	355.47\\
52000	339.82\\
54000	314.88\\
56000	332.53\\
58000	336.16\\
60000	353.3\\
62000	344.26\\
64000	343.05\\
66000	321.19\\
68000	329.52\\
70000	329.15\\
72000	333.23\\
74000	328.98\\
76000	320.56\\
78000	326.23\\
80000	325.37\\
82000	326.12\\
84000	322.6\\
86000	330.42\\
88000	343.53\\
90000	333.92\\
92000	340.26\\
94000	328.27\\
96000	333.27\\
98000	351.21\\
1e+05	340.69\\
};
\addplot [color=mycolor98]
  table[row sep=crcr]{%
0	4395.3\\
2000	2558.8\\
4000	1092.2\\
6000	674.47\\
8000	530.8\\
10000	461.8\\
12000	446.18\\
14000	433.48\\
16000	435.77\\
18000	446.17\\
20000	462.8\\
22000	456.5\\
24000	458.02\\
26000	464.65\\
28000	454.07\\
30000	461.93\\
32000	453.18\\
34000	454.17\\
36000	459.39\\
38000	442.34\\
40000	447.1\\
42000	435.53\\
44000	447.08\\
46000	460.59\\
48000	451.47\\
50000	432.97\\
52000	439.86\\
56000	430.47\\
58000	433.62\\
60000	430.77\\
62000	432.93\\
64000	444.78\\
66000	455.48\\
68000	462.57\\
70000	448.72\\
74000	442.66\\
76000	452.21\\
78000	455.35\\
80000	462.45\\
82000	463.97\\
84000	453.83\\
86000	452.4\\
88000	443.14\\
90000	443.57\\
94000	435.62\\
96000	434.28\\
98000	450.54\\
1e+05	465.86\\
};
\addlegendentry{\ClasskIVLegend}

      % \input{../../../../.tex/20/KS/SizeEvolutionsfig3.tex}
      \legend{}

      \nextgroupplot[%
        % xmin=0,xmax=1e+05,
        % ymin=100,ymax=1.5103e+05,
      ] % This file was created by matlab2tikz.
%
\definecolor{mycolor1}{rgb}{0.00146,0.00047,0.01387}%
\definecolor{mycolor2}{rgb}{0.01166,0.00942,0.06346}%
\definecolor{mycolor3}{rgb}{0.02943,0.02150,0.11462}%
\definecolor{mycolor4}{rgb}{0.06134,0.03659,0.17764}%
\definecolor{mycolor5}{rgb}{0.08741,0.04456,0.22481}%
\definecolor{mycolor6}{rgb}{0.12291,0.04754,0.28162}%
\definecolor{mycolor7}{rgb}{0.14907,0.04547,0.31709}%
\definecolor{mycolor8}{rgb}{0.18343,0.04033,0.35497}%
\definecolor{mycolor9}{rgb}{0.21795,0.03662,0.38352}%
\definecolor{mycolor10}{rgb}{0.24497,0.03705,0.40001}%
\definecolor{mycolor11}{rgb}{0.27135,0.04092,0.41198}%
\definecolor{mycolor12}{rgb}{0.29076,0.04564,0.41864}%
\definecolor{mycolor13}{rgb}{0.31628,0.05349,0.42512}%
\definecolor{mycolor14}{rgb}{0.33522,0.06006,0.42852}%
\definecolor{mycolor15}{rgb}{0.36028,0.06925,0.43150}%
\definecolor{mycolor16}{rgb}{0.37900,0.07625,0.43272}%
\definecolor{mycolor17}{rgb}{0.39767,0.08326,0.43318}%
\definecolor{mycolor18}{rgb}{0.41633,0.09020,0.43294}%
\definecolor{mycolor19}{rgb}{0.42877,0.09479,0.43241}%
\definecolor{mycolor20}{rgb}{0.44743,0.10160,0.43108}%
\definecolor{mycolor21}{rgb}{0.46610,0.10832,0.42912}%
\definecolor{mycolor22}{rgb}{0.47856,0.11276,0.42747}%
\definecolor{mycolor23}{rgb}{0.49726,0.11938,0.42449}%
\definecolor{mycolor24}{rgb}{0.50973,0.12377,0.42216}%
\definecolor{mycolor25}{rgb}{0.52221,0.12815,0.41955}%
\definecolor{mycolor26}{rgb}{0.53468,0.13253,0.41667}%
\definecolor{mycolor27}{rgb}{0.55339,0.13913,0.41183}%
\definecolor{mycolor28}{rgb}{0.56585,0.14357,0.40826}%
\definecolor{mycolor29}{rgb}{0.57830,0.14804,0.40441}%
\definecolor{mycolor30}{rgb}{0.59073,0.15256,0.40029}%
\definecolor{mycolor31}{rgb}{0.60314,0.15715,0.39589}%
\definecolor{mycolor32}{rgb}{0.61551,0.16182,0.39122}%
\definecolor{mycolor33}{rgb}{0.62169,0.16418,0.38878}%
\definecolor{mycolor34}{rgb}{0.63400,0.16899,0.38370}%
\definecolor{mycolor35}{rgb}{0.64626,0.17391,0.37836}%
\definecolor{mycolor36}{rgb}{0.65846,0.17896,0.37275}%
\definecolor{mycolor37}{rgb}{0.66454,0.18154,0.36985}%
\definecolor{mycolor38}{rgb}{0.67664,0.18681,0.36385}%
\definecolor{mycolor39}{rgb}{0.68865,0.19224,0.35760}%
\definecolor{mycolor40}{rgb}{0.69463,0.19502,0.35439}%
\definecolor{mycolor41}{rgb}{0.70650,0.20073,0.34778}%
\definecolor{mycolor42}{rgb}{0.71240,0.20366,0.34438}%
\definecolor{mycolor43}{rgb}{0.72410,0.20967,0.33742}%
\definecolor{mycolor44}{rgb}{0.72991,0.21276,0.33386}%
\definecolor{mycolor45}{rgb}{0.74142,0.21911,0.32658}%
\definecolor{mycolor46}{rgb}{0.74713,0.22238,0.32286}%
\definecolor{mycolor47}{rgb}{0.75279,0.22571,0.31909}%
\definecolor{mycolor48}{rgb}{0.76401,0.23255,0.31140}%
\definecolor{mycolor49}{rgb}{0.76956,0.23608,0.30749}%
\definecolor{mycolor50}{rgb}{0.77506,0.23967,0.30353}%
\definecolor{mycolor51}{rgb}{0.79129,0.25086,0.29139}%
\definecolor{mycolor52}{rgb}{0.79661,0.25473,0.28726}%
\definecolor{mycolor53}{rgb}{0.80187,0.25867,0.28310}%
\definecolor{mycolor54}{rgb}{0.81224,0.26679,0.27466}%
\definecolor{mycolor55}{rgb}{0.81734,0.27095,0.27039}%
\definecolor{mycolor56}{rgb}{0.82737,0.27952,0.26175}%
\definecolor{mycolor57}{rgb}{0.83230,0.28391,0.25738}%
\definecolor{mycolor58}{rgb}{0.83717,0.28839,0.25299}%
\definecolor{mycolor59}{rgb}{0.84671,0.29756,0.24411}%
\definecolor{mycolor60}{rgb}{0.85138,0.30226,0.23964}%
\definecolor{mycolor61}{rgb}{0.85599,0.30704,0.23513}%
\definecolor{mycolor62}{rgb}{0.86053,0.31189,0.23061}%
\definecolor{mycolor63}{rgb}{0.87374,0.32691,0.21689}%
\definecolor{mycolor64}{rgb}{0.87800,0.33206,0.21227}%
\definecolor{mycolor65}{rgb}{0.89034,0.34796,0.19829}%
\definecolor{mycolor66}{rgb}{0.89819,0.35891,0.18886}%
\definecolor{mycolor67}{rgb}{0.90200,0.36449,0.18412}%
\definecolor{mycolor68}{rgb}{0.90573,0.37014,0.17935}%
\definecolor{mycolor69}{rgb}{0.90939,0.37586,0.17456}%
\definecolor{mycolor70}{rgb}{0.91297,0.38164,0.16975}%
\definecolor{mycolor71}{rgb}{0.91646,0.38748,0.16492}%
\definecolor{mycolor72}{rgb}{0.91988,0.39339,0.16007}%
\definecolor{mycolor73}{rgb}{0.92322,0.39936,0.15519}%
\definecolor{mycolor74}{rgb}{0.92647,0.40539,0.15029}%
\definecolor{mycolor75}{rgb}{0.92964,0.41148,0.14537}%
\definecolor{mycolor76}{rgb}{0.93274,0.41763,0.14042}%
\definecolor{mycolor77}{rgb}{0.93575,0.42383,0.13544}%
\definecolor{mycolor78}{rgb}{0.93868,0.43009,0.13044}%
\definecolor{mycolor79}{rgb}{0.94152,0.43640,0.12541}%
\definecolor{mycolor80}{rgb}{0.94429,0.44277,0.12035}%
\definecolor{mycolor81}{rgb}{0.94956,0.45566,0.11016}%
\definecolor{mycolor82}{rgb}{0.95208,0.46218,0.10503}%
\definecolor{mycolor83}{rgb}{0.96129,0.48872,0.08429}%
\definecolor{mycolor84}{rgb}{0.96540,0.50225,0.07386}%
\definecolor{mycolor85}{rgb}{0.96732,0.50908,0.06866}%
\definecolor{mycolor86}{rgb}{0.97259,0.52980,0.05332}%
\definecolor{mycolor87}{rgb}{0.97568,0.54380,0.04362}%
\definecolor{mycolor88}{rgb}{0.98082,0.57221,0.02851}%
\definecolor{mycolor89}{rgb}{0.98189,0.57939,0.02625}%
\definecolor{mycolor90}{rgb}{0.98378,0.59385,0.02377}%
\definecolor{mycolor91}{rgb}{0.98532,0.60842,0.02420}%
\definecolor{mycolor92}{rgb}{0.98696,0.63048,0.03091}%
\definecolor{mycolor93}{rgb}{0.98793,0.66025,0.05175}%
\definecolor{mycolor94}{rgb}{0.98771,0.68281,0.07249}%
\definecolor{mycolor95}{rgb}{0.97750,0.78226,0.18592}%
\definecolor{mycolor96}{rgb}{0.96624,0.83619,0.26153}%
\definecolor{mycolor97}{rgb}{0.96252,0.85148,0.28555}%
\definecolor{mycolor98}{rgb}{0.98836,0.99836,0.64492}%
%
\addplot [color=mycolor1]
  table[row sep=crcr]{%
0	4395.3\\
2000	71780\\
4000	1.2955e+05\\
6000	1.4472e+05\\
8000	1.4692e+05\\
10000	1.4473e+05\\
12000	1.4525e+05\\
14000	1.4911e+05\\
16000	1.4879e+05\\
18000	1.4999e+05\\
20000	1.4911e+05\\
22000	1.4885e+05\\
24000	1.491e+05\\
26000	1.4804e+05\\
28000	1.4743e+05\\
30000	1.4789e+05\\
32000	1.4647e+05\\
34000	1.4876e+05\\
36000	1.4931e+05\\
38000	1.4578e+05\\
40000	1.4397e+05\\
42000	1.4452e+05\\
44000	1.4572e+05\\
46000	1.4979e+05\\
48000	1.5103e+05\\
50000	1.4822e+05\\
52000	1.4641e+05\\
54000	1.4628e+05\\
56000	1.5073e+05\\
58000	1.4885e+05\\
62000	1.4597e+05\\
64000	1.475e+05\\
66000	1.5001e+05\\
70000	1.4716e+05\\
72000	1.4905e+05\\
76000	1.4789e+05\\
78000	1.489e+05\\
80000	1.4669e+05\\
82000	1.4537e+05\\
84000	1.4495e+05\\
86000	1.4791e+05\\
88000	1.4558e+05\\
90000	1.4657e+05\\
92000	1.4885e+05\\
94000	1.5002e+05\\
96000	1.4924e+05\\
98000	1.4635e+05\\
1e+05	1.4782e+05\\
};
\addlegendentry{\ClasskILegend}

\addplot [color=mycolor2, forget plot]
  table[row sep=crcr]{%
0	4395.3\\
2000	19786\\
4000	25183\\
6000	26062\\
8000	25381\\
10000	25000\\
12000	25409\\
14000	25570\\
18000	24862\\
20000	25570\\
22000	25484\\
24000	25872\\
26000	25298\\
28000	25774\\
30000	25013\\
32000	25346\\
34000	24977\\
36000	25274\\
38000	25262\\
40000	25562\\
42000	25327\\
44000	26044\\
46000	25865\\
48000	24449\\
50000	25118\\
52000	25672\\
54000	25887\\
56000	25517\\
58000	25308\\
60000	25664\\
62000	25055\\
64000	25604\\
66000	25205\\
68000	25603\\
70000	25374\\
72000	25467\\
74000	25298\\
76000	25057\\
78000	25078\\
80000	25568\\
82000	25991\\
84000	25286\\
86000	26186\\
88000	25743\\
90000	25049\\
92000	25583\\
94000	25351\\
98000	25337\\
1e+05	25306\\
};
\addplot [color=mycolor3, forget plot]
  table[row sep=crcr]{%
0	4395.3\\
2000	27626\\
4000	37322\\
6000	39122\\
8000	38762\\
10000	39655\\
12000	38182\\
14000	38178\\
16000	38399\\
18000	38481\\
20000	37464\\
22000	38723\\
24000	38423\\
26000	38536\\
28000	38233\\
32000	39390\\
34000	39346\\
36000	38521\\
38000	38522\\
40000	38749\\
42000	38838\\
44000	39413\\
46000	38350\\
48000	38123\\
50000	38076\\
52000	38181\\
54000	37629\\
56000	37976\\
58000	39078\\
60000	39173\\
62000	38774\\
64000	38789\\
66000	38473\\
70000	38490\\
72000	37677\\
74000	38079\\
76000	38160\\
78000	38867\\
80000	38919\\
82000	39590\\
84000	38848\\
86000	38493\\
88000	38785\\
90000	38485\\
92000	38034\\
98000	39547\\
1e+05	39266\\
};
\addplot [color=mycolor4, forget plot]
  table[row sep=crcr]{%
0	4395.3\\
2000	12431\\
4000	14245\\
8000	14134\\
10000	14298\\
12000	14008\\
14000	14399\\
16000	13990\\
18000	13885\\
20000	13958\\
22000	13917\\
24000	14109\\
26000	14179\\
28000	14030\\
30000	14204\\
32000	14201\\
34000	13909\\
36000	14176\\
38000	14072\\
40000	14139\\
42000	14070\\
44000	13953\\
46000	14057\\
48000	14023\\
50000	14287\\
52000	14419\\
54000	13878\\
56000	14112\\
58000	13799\\
60000	13968\\
62000	13996\\
64000	14139\\
66000	13854\\
68000	14667\\
70000	14002\\
72000	14176\\
74000	13972\\
76000	14284\\
78000	14122\\
80000	14112\\
82000	14034\\
84000	14033\\
86000	13602\\
88000	14262\\
90000	14049\\
92000	14125\\
94000	14090\\
98000	14174\\
1e+05	13994\\
};
\addplot [color=mycolor5, forget plot]
  table[row sep=crcr]{%
0	4395.3\\
2000	33577\\
4000	55951\\
6000	60552\\
8000	60329\\
10000	61012\\
12000	62738\\
14000	61076\\
16000	62688\\
18000	61298\\
20000	63226\\
22000	61914\\
24000	60945\\
28000	62223\\
30000	61409\\
32000	62206\\
34000	63483\\
36000	61848\\
38000	62532\\
40000	61887\\
42000	63625\\
46000	61207\\
48000	61283\\
50000	61778\\
52000	62090\\
54000	60433\\
56000	60497\\
58000	61646\\
60000	60624\\
62000	61310\\
64000	62197\\
66000	62179\\
70000	61657\\
72000	62160\\
74000	61694\\
76000	61543\\
78000	60954\\
80000	60807\\
82000	62631\\
84000	62997\\
86000	61884\\
88000	62907\\
90000	63144\\
92000	63040\\
94000	61699\\
96000	60568\\
98000	61261\\
1e+05	60210\\
};
\addplot [color=mycolor6, forget plot]
  table[row sep=crcr]{%
0	4395.3\\
2000	30230\\
4000	48800\\
6000	53495\\
10000	54467\\
12000	54099\\
14000	52709\\
18000	53554\\
20000	53507\\
22000	53709\\
26000	53551\\
28000	54628\\
30000	54194\\
32000	53140\\
34000	53967\\
36000	54277\\
38000	53754\\
40000	54060\\
42000	54602\\
44000	54604\\
46000	53279\\
48000	54698\\
50000	54379\\
52000	53784\\
54000	55007\\
56000	53525\\
58000	52766\\
60000	54042\\
62000	53743\\
68000	55104\\
70000	55141\\
72000	54578\\
76000	54456\\
78000	55471\\
80000	55407\\
84000	54362\\
86000	53296\\
88000	54331\\
90000	53114\\
92000	54330\\
94000	54512\\
96000	56051\\
98000	54313\\
1e+05	53630\\
};
\addplot [color=mycolor7, forget plot]
  table[row sep=crcr]{%
0	4395.3\\
2000	19381\\
4000	26791\\
6000	27554\\
8000	26869\\
10000	27654\\
12000	27824\\
14000	27729\\
16000	27515\\
20000	26696\\
22000	27504\\
24000	26870\\
28000	27751\\
30000	27558\\
32000	26976\\
34000	27508\\
36000	28217\\
38000	28066\\
40000	27715\\
42000	27123\\
44000	27135\\
46000	27381\\
48000	27189\\
50000	27165\\
52000	27660\\
54000	27560\\
56000	27795\\
58000	26772\\
60000	27398\\
62000	28157\\
64000	27625\\
66000	26526\\
68000	27204\\
70000	27408\\
72000	27770\\
74000	27932\\
76000	27601\\
80000	28371\\
82000	27545\\
84000	27277\\
86000	27614\\
90000	27639\\
92000	27246\\
94000	27577\\
96000	28231\\
98000	28038\\
1e+05	27599\\
};
\addplot [color=mycolor8, forget plot]
  table[row sep=crcr]{%
0	4395.3\\
2000	15373\\
4000	19465\\
6000	19955\\
8000	19901\\
10000	20155\\
12000	20261\\
14000	19922\\
16000	19956\\
18000	20119\\
20000	20345\\
22000	20252\\
24000	19808\\
26000	19806\\
28000	20444\\
30000	20429\\
32000	20254\\
34000	20401\\
36000	20410\\
38000	19914\\
40000	19667\\
42000	19947\\
44000	20675\\
46000	19913\\
48000	19658\\
50000	20255\\
52000	20311\\
54000	19808\\
56000	20097\\
58000	20130\\
60000	20391\\
62000	20729\\
64000	20371\\
66000	19955\\
68000	20268\\
70000	19924\\
72000	19822\\
74000	19963\\
76000	19868\\
78000	19936\\
80000	20094\\
82000	20459\\
84000	20579\\
86000	19793\\
88000	20028\\
90000	20329\\
92000	19561\\
94000	19701\\
96000	19473\\
98000	19794\\
1e+05	20270\\
};
\addplot [color=mycolor9, forget plot]
  table[row sep=crcr]{%
0	4395.3\\
2000	11641\\
4000	13824\\
6000	13792\\
8000	13650\\
10000	13741\\
12000	14288\\
14000	14374\\
16000	13970\\
18000	13824\\
20000	13379\\
22000	13862\\
24000	13925\\
28000	13851\\
30000	13925\\
32000	14039\\
34000	13600\\
36000	13564\\
38000	14108\\
40000	13870\\
42000	14035\\
44000	14047\\
46000	13720\\
48000	13767\\
50000	14153\\
52000	14146\\
54000	13967\\
56000	14156\\
58000	13970\\
60000	13710\\
62000	14183\\
64000	14099\\
66000	13769\\
68000	14114\\
70000	14284\\
72000	13863\\
74000	13663\\
76000	14302\\
78000	13477\\
80000	13770\\
82000	13905\\
84000	13978\\
86000	14234\\
90000	13834\\
92000	14128\\
94000	13830\\
96000	13439\\
98000	14064\\
1e+05	14197\\
};
\addplot [color=mycolor10, forget plot]
  table[row sep=crcr]{%
0	4395.3\\
2000	15125\\
4000	20453\\
6000	21428\\
8000	21753\\
10000	21858\\
12000	21750\\
14000	21417\\
16000	21747\\
18000	21859\\
20000	21729\\
22000	21430\\
24000	21697\\
26000	21581\\
28000	21706\\
30000	21630\\
32000	21304\\
34000	21304\\
36000	21603\\
38000	21511\\
40000	21728\\
44000	21155\\
46000	21496\\
48000	21356\\
50000	21749\\
52000	21530\\
54000	21397\\
56000	21031\\
58000	21198\\
60000	21214\\
62000	21604\\
64000	21436\\
66000	21771\\
68000	21956\\
70000	21368\\
72000	21645\\
74000	21320\\
76000	21180\\
78000	20976\\
80000	21821\\
82000	22050\\
84000	21422\\
86000	21258\\
88000	21305\\
90000	21854\\
92000	21571\\
94000	21423\\
96000	21045\\
98000	21292\\
1e+05	21145\\
};
\addplot [color=mycolor11, forget plot]
  table[row sep=crcr]{%
0	4395.3\\
2000	8023.8\\
4000	8562.4\\
6000	8327.3\\
8000	8601.1\\
10000	8648.4\\
12000	8227.6\\
14000	8591.1\\
16000	8926.7\\
18000	8620.1\\
20000	8489.7\\
22000	8407.3\\
24000	8576.8\\
26000	8589.5\\
28000	8436.4\\
30000	8492.8\\
32000	8585.2\\
34000	8651.8\\
36000	8526.7\\
38000	8631.2\\
40000	8479.3\\
42000	8844.6\\
44000	8597\\
48000	8672.8\\
50000	8547.7\\
54000	8570.5\\
56000	8677.2\\
58000	8698.2\\
60000	8535.7\\
62000	8687.4\\
64000	8634.3\\
66000	8524.4\\
68000	8575.8\\
70000	8494.8\\
72000	8692.4\\
74000	8758.9\\
76000	8413.8\\
78000	8369.6\\
82000	8738.3\\
84000	8438.1\\
86000	8613.2\\
88000	8404.4\\
90000	8712.6\\
92000	8523.3\\
94000	8696.7\\
96000	8662.4\\
98000	8579.7\\
1e+05	8543.1\\
};
\addplot [color=mycolor12, forget plot]
  table[row sep=crcr]{%
0	4395.3\\
2000	11932\\
4000	14199\\
8000	14676\\
10000	14541\\
12000	14957\\
14000	14631\\
16000	14605\\
18000	14391\\
20000	14270\\
22000	14192\\
24000	14423\\
26000	14306\\
28000	14453\\
30000	14811\\
32000	14541\\
34000	14352\\
36000	14437\\
40000	14275\\
42000	14412\\
44000	14451\\
46000	14612\\
48000	14279\\
50000	14207\\
52000	14394\\
54000	14316\\
56000	14540\\
58000	14910\\
60000	14345\\
62000	14366\\
64000	14139\\
66000	14665\\
68000	14011\\
72000	14922\\
74000	14918\\
76000	14320\\
78000	14524\\
80000	14634\\
82000	14993\\
84000	14419\\
86000	14703\\
88000	14069\\
90000	14256\\
92000	14307\\
94000	14240\\
96000	14704\\
1e+05	14360\\
};
\addplot [color=mycolor13, forget plot]
  table[row sep=crcr]{%
0	4395.3\\
2000	9522.2\\
4000	11129\\
6000	11115\\
8000	10984\\
10000	10999\\
12000	11054\\
14000	10957\\
16000	11044\\
18000	11359\\
20000	11198\\
22000	10921\\
24000	11159\\
26000	11163\\
28000	10689\\
30000	10867\\
32000	11146\\
34000	10857\\
36000	10852\\
38000	11109\\
40000	11011\\
42000	11146\\
44000	11018\\
46000	11125\\
48000	10906\\
50000	10833\\
52000	11358\\
54000	11106\\
56000	11243\\
58000	10992\\
60000	11214\\
62000	11177\\
64000	11221\\
66000	10979\\
68000	11332\\
70000	11066\\
72000	11053\\
74000	11303\\
76000	11000\\
78000	10949\\
80000	10983\\
82000	10560\\
84000	10698\\
86000	10685\\
88000	11006\\
90000	10852\\
92000	11063\\
94000	10977\\
96000	10783\\
98000	10929\\
1e+05	11035\\
};
\addplot [color=mycolor14, forget plot]
  table[row sep=crcr]{%
0	4395.3\\
2000	16277\\
4000	22442\\
6000	23212\\
8000	23465\\
12000	23619\\
14000	24052\\
16000	23860\\
18000	23249\\
20000	23786\\
22000	23359\\
24000	23934\\
26000	24308\\
28000	24240\\
30000	23427\\
32000	23361\\
34000	23586\\
36000	23557\\
38000	23604\\
40000	24614\\
42000	23343\\
44000	23840\\
46000	23966\\
48000	23757\\
50000	23627\\
52000	23828\\
54000	23406\\
56000	23334\\
58000	24162\\
60000	24344\\
62000	23868\\
64000	24025\\
68000	23631\\
70000	23549\\
72000	23690\\
74000	24212\\
76000	24172\\
78000	23233\\
80000	23566\\
82000	23985\\
84000	23947\\
86000	24265\\
88000	23885\\
90000	24483\\
92000	23707\\
94000	23823\\
96000	24110\\
98000	23892\\
1e+05	24139\\
};
\addplot [color=mycolor15, forget plot]
  table[row sep=crcr]{%
0	4395.3\\
2000	12645\\
4000	16237\\
6000	16467\\
8000	16899\\
10000	16909\\
12000	16676\\
14000	16303\\
16000	16239\\
18000	17163\\
20000	16481\\
22000	16647\\
24000	17329\\
26000	16645\\
30000	17104\\
34000	16417\\
36000	16454\\
38000	17136\\
40000	17093\\
42000	16712\\
44000	16606\\
46000	16574\\
48000	16455\\
50000	16746\\
52000	17198\\
54000	16568\\
56000	16833\\
58000	16916\\
60000	17489\\
62000	17241\\
64000	16381\\
66000	17126\\
68000	16464\\
70000	16977\\
72000	16688\\
74000	16768\\
76000	16523\\
78000	16453\\
80000	16647\\
82000	16023\\
84000	16551\\
86000	16251\\
88000	16741\\
90000	16949\\
92000	16996\\
94000	16660\\
96000	16503\\
98000	17024\\
1e+05	16599\\
};
\addplot [color=mycolor16, forget plot]
  table[row sep=crcr]{%
0	4395.3\\
2000	6559.3\\
4000	7043.2\\
6000	7130.2\\
8000	7002.4\\
10000	6951.2\\
12000	7008.9\\
14000	6777.3\\
16000	7150.8\\
18000	7078.3\\
20000	7043.2\\
22000	7038\\
24000	6871.6\\
28000	6777.9\\
30000	7059.4\\
32000	7020.1\\
34000	6842\\
36000	7061.7\\
38000	7066\\
40000	6769.3\\
42000	7104\\
44000	6957.8\\
46000	6953.6\\
48000	7004.5\\
50000	6785\\
52000	7062.7\\
54000	6959.4\\
56000	6979.5\\
60000	7223.5\\
62000	7256.3\\
64000	6987.7\\
66000	6755.8\\
68000	7117.5\\
70000	6954.6\\
72000	6831.9\\
74000	7123.4\\
76000	7005.7\\
78000	7222.5\\
80000	7112.6\\
82000	7175.4\\
84000	7007.1\\
86000	7037.3\\
88000	6811.9\\
90000	7015.1\\
92000	6938.8\\
96000	6959.4\\
98000	7127.7\\
1e+05	6958.1\\
};
\addplot [color=mycolor17, forget plot]
  table[row sep=crcr]{%
0	4395.3\\
2000	12275\\
4000	15655\\
6000	15789\\
8000	15664\\
10000	16138\\
12000	15817\\
14000	15942\\
16000	15725\\
18000	16097\\
20000	15846\\
22000	16064\\
24000	15908\\
26000	15664\\
28000	15899\\
30000	15833\\
32000	16060\\
34000	15599\\
36000	15472\\
38000	16280\\
40000	15887\\
42000	15597\\
44000	15409\\
46000	15896\\
48000	15991\\
50000	15867\\
52000	15849\\
54000	16009\\
56000	15980\\
58000	16314\\
60000	16204\\
62000	15779\\
64000	15272\\
66000	15903\\
68000	15891\\
70000	15688\\
72000	15253\\
74000	15958\\
76000	15882\\
78000	15431\\
80000	15442\\
84000	15674\\
86000	15358\\
88000	15456\\
90000	15770\\
94000	16062\\
98000	15689\\
1e+05	15564\\
};
\addplot [color=mycolor18, forget plot]
  table[row sep=crcr]{%
0	4395.3\\
2000	16360\\
4000	22677\\
6000	23294\\
8000	24102\\
10000	23860\\
12000	23399\\
14000	23918\\
16000	23474\\
18000	24253\\
20000	23829\\
22000	23775\\
24000	23966\\
26000	24608\\
28000	23798\\
30000	24179\\
32000	24189\\
34000	23893\\
36000	23777\\
40000	24283\\
42000	24386\\
44000	24130\\
46000	25107\\
48000	24160\\
50000	24012\\
54000	24339\\
56000	23828\\
58000	24560\\
60000	23858\\
62000	24285\\
64000	24234\\
66000	24489\\
68000	24890\\
70000	23876\\
72000	23589\\
74000	24707\\
76000	24501\\
80000	23702\\
82000	23707\\
84000	24323\\
86000	24691\\
88000	24119\\
90000	23212\\
92000	23829\\
94000	24306\\
96000	24418\\
98000	24772\\
1e+05	23842\\
};
\addplot [color=mycolor19, forget plot]
  table[row sep=crcr]{%
0	4395.3\\
2000	11433\\
4000	14522\\
6000	14727\\
8000	14783\\
10000	15023\\
12000	14576\\
14000	14878\\
16000	15373\\
18000	14911\\
20000	14315\\
22000	14901\\
24000	14902\\
28000	14089\\
30000	14641\\
32000	14799\\
34000	14527\\
36000	14573\\
38000	14706\\
40000	14655\\
42000	14552\\
44000	14930\\
46000	14388\\
48000	14454\\
50000	14466\\
52000	14286\\
54000	14859\\
56000	14505\\
58000	14462\\
60000	14809\\
62000	14847\\
64000	14970\\
66000	13978\\
70000	14433\\
72000	14403\\
74000	14244\\
76000	14306\\
78000	14801\\
80000	14690\\
82000	14687\\
84000	14817\\
86000	14627\\
88000	15106\\
90000	14586\\
92000	14540\\
94000	14630\\
98000	14094\\
1e+05	14531\\
};
\addplot [color=mycolor20, forget plot]
  table[row sep=crcr]{%
0	4395.3\\
2000	13511\\
4000	18193\\
6000	18913\\
8000	19151\\
10000	18469\\
12000	18724\\
14000	18743\\
16000	18640\\
18000	18961\\
20000	19420\\
22000	19408\\
24000	18571\\
26000	18805\\
28000	18915\\
30000	18473\\
34000	19127\\
36000	19215\\
38000	19020\\
40000	18758\\
42000	19008\\
44000	19479\\
48000	18767\\
50000	19034\\
52000	19031\\
54000	18750\\
56000	19394\\
58000	19259\\
60000	18555\\
62000	18765\\
64000	18639\\
66000	18089\\
68000	18704\\
70000	19666\\
72000	18736\\
74000	18394\\
76000	19083\\
78000	19185\\
80000	18813\\
82000	18673\\
84000	19036\\
86000	18908\\
88000	18569\\
90000	18780\\
92000	18259\\
96000	18294\\
98000	18747\\
1e+05	19051\\
};
\addplot [color=mycolor21, forget plot]
  table[row sep=crcr]{%
0	4395.3\\
2000	7984.1\\
4000	9215.9\\
6000	9496\\
8000	9310.4\\
10000	9621.1\\
12000	9652.1\\
14000	9742.9\\
16000	9625.8\\
18000	9776\\
20000	9309.6\\
22000	9334\\
24000	9488\\
26000	9829.6\\
28000	9655.6\\
30000	9397\\
36000	9695.3\\
38000	9598\\
40000	9652.3\\
42000	9835.4\\
46000	9572.6\\
48000	9495.1\\
52000	9793.5\\
54000	9462.5\\
56000	9558.1\\
58000	9573\\
60000	9556\\
62000	9718.4\\
64000	9727.4\\
66000	9504.1\\
68000	9467.6\\
70000	9660.7\\
72000	9613.8\\
74000	9723.5\\
76000	9578.5\\
78000	9604.9\\
80000	9602.9\\
82000	9459.5\\
84000	9585.6\\
86000	9459\\
88000	9718.8\\
92000	9586.1\\
94000	9930.8\\
96000	9677\\
98000	9617.5\\
1e+05	9729.8\\
};
\addplot [color=mycolor22, forget plot]
  table[row sep=crcr]{%
0	4395.3\\
2000	16457\\
4000	24255\\
6000	24538\\
8000	24670\\
10000	24470\\
12000	25009\\
14000	24138\\
16000	24376\\
18000	24450\\
20000	24995\\
22000	25098\\
24000	24012\\
26000	24258\\
28000	25009\\
30000	24903\\
32000	25248\\
34000	24376\\
36000	25002\\
38000	25168\\
40000	24372\\
42000	25142\\
44000	25058\\
46000	25359\\
50000	25436\\
52000	24554\\
54000	24898\\
56000	24723\\
58000	25094\\
60000	24355\\
62000	24145\\
64000	24439\\
66000	23471\\
68000	24239\\
70000	24899\\
72000	24908\\
76000	24554\\
78000	24923\\
80000	25425\\
82000	24854\\
84000	24677\\
86000	24783\\
88000	24776\\
90000	24868\\
92000	24343\\
94000	24015\\
96000	24712\\
98000	24955\\
1e+05	24799\\
};
\addplot [color=mycolor23, forget plot]
  table[row sep=crcr]{%
0	4395.3\\
2000	9917.3\\
4000	12083\\
6000	12074\\
8000	12275\\
10000	12119\\
12000	12082\\
16000	12346\\
20000	12065\\
22000	12004\\
24000	12280\\
26000	12152\\
28000	11983\\
30000	12081\\
32000	12302\\
34000	12304\\
36000	12199\\
38000	12165\\
40000	11961\\
42000	12010\\
44000	12321\\
46000	12318\\
48000	11988\\
50000	12296\\
52000	12147\\
54000	12212\\
58000	12007\\
60000	12145\\
62000	12027\\
64000	12292\\
66000	12089\\
68000	11927\\
70000	12072\\
72000	12153\\
74000	11986\\
76000	12284\\
78000	12194\\
80000	12053\\
82000	12217\\
84000	12264\\
86000	12070\\
90000	12113\\
92000	12214\\
94000	12400\\
96000	12168\\
98000	12121\\
1e+05	12142\\
};
\addplot [color=mycolor24, forget plot]
  table[row sep=crcr]{%
0	4395.3\\
2000	4928.9\\
4000	5014.3\\
6000	4892.8\\
8000	4797.2\\
10000	4850\\
12000	4809.1\\
14000	4865.5\\
16000	4870\\
18000	4807.6\\
20000	4781.9\\
22000	4854.3\\
24000	4789.5\\
26000	4749.3\\
28000	4789\\
30000	4806.7\\
32000	4855.7\\
34000	4860\\
36000	4890.1\\
38000	4778.7\\
40000	4845.2\\
42000	4844.1\\
46000	4790.4\\
50000	4850.3\\
52000	4790.1\\
54000	4912.6\\
56000	4765.6\\
58000	4893.5\\
60000	4810.3\\
62000	4656.9\\
64000	4845.9\\
66000	4946.3\\
68000	4967\\
70000	4894\\
72000	4766.6\\
74000	4616.3\\
76000	4757.9\\
78000	4725\\
80000	4774.7\\
82000	4703.1\\
84000	4841.7\\
86000	4794.8\\
88000	4768.2\\
90000	4658.6\\
92000	4973\\
94000	4785.7\\
96000	4778.8\\
98000	4798.7\\
1e+05	4857\\
};
\addplot [color=mycolor25, forget plot]
  table[row sep=crcr]{%
0	4395.3\\
2000	9971.8\\
4000	12197\\
8000	12338\\
10000	12097\\
12000	12196\\
14000	12159\\
16000	11992\\
18000	11867\\
20000	12177\\
22000	12053\\
24000	12484\\
26000	12023\\
28000	12029\\
30000	12193\\
32000	11992\\
34000	12429\\
36000	12404\\
38000	12620\\
40000	12328\\
42000	12283\\
44000	12299\\
46000	11963\\
48000	12671\\
50000	12563\\
52000	12349\\
56000	12324\\
58000	12173\\
60000	12172\\
62000	12344\\
64000	12200\\
66000	12598\\
68000	11950\\
70000	11894\\
72000	12174\\
74000	12410\\
76000	12579\\
78000	12070\\
80000	12243\\
82000	12187\\
84000	12459\\
86000	12380\\
88000	12167\\
90000	12170\\
92000	12113\\
94000	12166\\
96000	12446\\
98000	12344\\
1e+05	12549\\
};
\addplot [color=mycolor26, forget plot]
  table[row sep=crcr]{%
0	4395.3\\
2000	7490.5\\
4000	8265.7\\
6000	8207.4\\
8000	8400.6\\
10000	8407.7\\
12000	8356.3\\
14000	8204.8\\
16000	8224.1\\
18000	8084.4\\
20000	8130.2\\
22000	8245.6\\
24000	8482.1\\
26000	8254.9\\
28000	8093.3\\
30000	8393.8\\
32000	8329.2\\
34000	8154\\
36000	8089.3\\
38000	8668.4\\
40000	8407.2\\
42000	8312.7\\
44000	8268\\
46000	8305.9\\
48000	7900.2\\
50000	7732.8\\
52000	8106\\
54000	8274\\
56000	8196.7\\
58000	8298.3\\
60000	8192.8\\
62000	8353.2\\
64000	8174.6\\
66000	8150.1\\
68000	8304.6\\
70000	8269.6\\
72000	8138.3\\
76000	8282.8\\
78000	8245.7\\
80000	8157\\
82000	8488.7\\
84000	8220.2\\
86000	8069.1\\
88000	8188.6\\
90000	8265.1\\
92000	8200.1\\
94000	8402\\
96000	8076.2\\
98000	8319.4\\
1e+05	8434.1\\
};
\addplot [color=mycolor27, forget plot]
  table[row sep=crcr]{%
0	4395.3\\
2000	8342.2\\
4000	9667.3\\
6000	9545.2\\
8000	9662\\
10000	9589\\
12000	9650.3\\
14000	9555.1\\
16000	10034\\
18000	9860.1\\
20000	9608.6\\
22000	9707.4\\
24000	9494.1\\
26000	9877.1\\
28000	9467.4\\
30000	9582.9\\
32000	9538.1\\
34000	9668.5\\
36000	9714.9\\
38000	9576.1\\
40000	9733.3\\
42000	9855.3\\
44000	9727.4\\
46000	9919.4\\
48000	9715.7\\
50000	9861.1\\
52000	9867\\
54000	9663.4\\
56000	9670\\
58000	9563\\
60000	9570.5\\
62000	9614.6\\
64000	9784.1\\
66000	9767.8\\
68000	9892.3\\
70000	10142\\
72000	9349\\
74000	9384\\
78000	9763.4\\
80000	9583.9\\
82000	9736\\
84000	9941.4\\
88000	9482.9\\
90000	9851.1\\
92000	9642.3\\
94000	9711.4\\
96000	9983\\
98000	9931.7\\
1e+05	9804.7\\
};
\addplot [color=mycolor28, forget plot]
  table[row sep=crcr]{%
0	4395.3\\
2000	5353.7\\
4000	5404.5\\
6000	5433.2\\
8000	5506.2\\
10000	5468.3\\
12000	5467.3\\
14000	5574.5\\
16000	5495.6\\
18000	5618\\
20000	5603.6\\
22000	5420.5\\
24000	5471\\
26000	5539\\
28000	5532.6\\
30000	5621.5\\
32000	5358.7\\
36000	5587.6\\
38000	5507.4\\
40000	5539.5\\
42000	5607.6\\
44000	5291.5\\
46000	5312.7\\
48000	5431.7\\
50000	5751.5\\
52000	5503.4\\
54000	5483.2\\
56000	5402\\
58000	5533.1\\
60000	5551\\
62000	5447.2\\
64000	5368.9\\
66000	5406.3\\
68000	5364.6\\
70000	5470.9\\
72000	5558.5\\
74000	5345.7\\
76000	5498.7\\
78000	5571.1\\
80000	5585.1\\
82000	5496.6\\
84000	5586.5\\
88000	5560.9\\
90000	5445.4\\
94000	5431.2\\
96000	5551.6\\
98000	5494.2\\
1e+05	5263.5\\
};
\addplot [color=mycolor29, forget plot]
  table[row sep=crcr]{%
0	4395.3\\
2000	9439.2\\
4000	11204\\
6000	11235\\
8000	11672\\
10000	11724\\
14000	11200\\
16000	11398\\
20000	11979\\
22000	11811\\
24000	11588\\
26000	11322\\
28000	11822\\
30000	11535\\
32000	11350\\
34000	11258\\
36000	11528\\
38000	11496\\
46000	11536\\
48000	11123\\
50000	11436\\
52000	11567\\
54000	11639\\
56000	11769\\
58000	11441\\
60000	11620\\
64000	11576\\
66000	11336\\
68000	11265\\
70000	11532\\
72000	11285\\
74000	11104\\
76000	11807\\
78000	11811\\
80000	11640\\
82000	11287\\
84000	11903\\
86000	11641\\
88000	11667\\
90000	12027\\
92000	11560\\
94000	11553\\
96000	11440\\
98000	11386\\
1e+05	11665\\
};
\addplot [color=mycolor30, forget plot]
  table[row sep=crcr]{%
0	4395.3\\
2000	7130.9\\
4000	7902.1\\
6000	7975.2\\
8000	8107.5\\
12000	7911.7\\
14000	7855\\
16000	7841.6\\
18000	7922.7\\
20000	7789.5\\
22000	7803.1\\
24000	7882.1\\
26000	7860.6\\
28000	7805\\
30000	7897\\
32000	8078.2\\
34000	7817.8\\
36000	7914.2\\
38000	7957.3\\
40000	8101.2\\
42000	7839.4\\
46000	7677.7\\
48000	7751.4\\
52000	7772.4\\
54000	7836.2\\
56000	7808.6\\
58000	8040.6\\
62000	7973.6\\
64000	7844\\
68000	7799.1\\
70000	8050.2\\
72000	7876.5\\
74000	7822.2\\
76000	7818.9\\
78000	7881.5\\
80000	7989.4\\
82000	7765.4\\
84000	7977.5\\
86000	7792.2\\
88000	7880.4\\
90000	7681.9\\
92000	7925.8\\
94000	7880.3\\
96000	7954.1\\
98000	7815.2\\
1e+05	7881.3\\
};
\addplot [color=mycolor31, forget plot]
  table[row sep=crcr]{%
0	4395.3\\
2000	6195.7\\
4000	6602.6\\
6000	6576.2\\
8000	6581.3\\
10000	6774.5\\
12000	6466.3\\
14000	6545.2\\
16000	6573\\
18000	6529.4\\
20000	6446.8\\
22000	6417.5\\
24000	6689.3\\
28000	6289.1\\
30000	6405.3\\
32000	6563.1\\
34000	6517.7\\
36000	6711.1\\
38000	6565.4\\
40000	6573.5\\
42000	6492.8\\
44000	6624.5\\
46000	6445.6\\
48000	6466.6\\
50000	6327.7\\
52000	6303.2\\
54000	6516.7\\
56000	6405.6\\
58000	6477.5\\
60000	6483\\
62000	6641.1\\
64000	6353\\
66000	6192.6\\
68000	6378.4\\
70000	6241.1\\
72000	6500.4\\
74000	6488.1\\
76000	6594.8\\
78000	6537.8\\
80000	6372\\
82000	6592.5\\
84000	6613.3\\
86000	6696.8\\
88000	6453.1\\
90000	6337.6\\
92000	6458.7\\
94000	6367.5\\
96000	6662.8\\
98000	6366.4\\
1e+05	6058.5\\
};
\addplot [color=mycolor32, forget plot]
  table[row sep=crcr]{%
0	4395.3\\
2000	8816.1\\
4000	10434\\
6000	10437\\
8000	10348\\
10000	10455\\
12000	10400\\
14000	10850\\
16000	10471\\
18000	10618\\
20000	10265\\
22000	10402\\
24000	10134\\
28000	10519\\
30000	10094\\
32000	10273\\
34000	10262\\
36000	10291\\
38000	10156\\
40000	10270\\
42000	10484\\
44000	10575\\
46000	10266\\
48000	10364\\
50000	10161\\
52000	10431\\
56000	10571\\
58000	10575\\
60000	10449\\
62000	10561\\
64000	9889\\
66000	10242\\
68000	10673\\
70000	10545\\
72000	10728\\
74000	9989.9\\
76000	10715\\
78000	10473\\
80000	10555\\
82000	10677\\
84000	10412\\
86000	10501\\
88000	10299\\
90000	10273\\
92000	10438\\
94000	10567\\
96000	10571\\
98000	10826\\
1e+05	10591\\
};
\addplot [color=mycolor33, forget plot]
  table[row sep=crcr]{%
0	4395.3\\
2000	7327.1\\
4000	8262.3\\
6000	8530.2\\
8000	8162.2\\
10000	8171.8\\
12000	8224.3\\
14000	8027.7\\
16000	8114.6\\
18000	8092.3\\
20000	8204.7\\
22000	8232.9\\
24000	8144.6\\
26000	8152.5\\
28000	8106.3\\
30000	8263.8\\
32000	8083.8\\
34000	8231.8\\
36000	8261.3\\
38000	8238.9\\
40000	8314.2\\
44000	8106.6\\
46000	8140.3\\
48000	8342.2\\
50000	8120.8\\
52000	7940.9\\
54000	8097.3\\
56000	8150.3\\
58000	8294.7\\
60000	8175.2\\
62000	8398.7\\
64000	8226.9\\
66000	8604.7\\
68000	8221.5\\
72000	8014.6\\
74000	8226.1\\
76000	8171.6\\
78000	8046.8\\
80000	7978.6\\
82000	8045.7\\
84000	8146.9\\
86000	8368\\
88000	8276.1\\
90000	8277.2\\
92000	8093.2\\
94000	8335.9\\
96000	8267.9\\
98000	8059\\
1e+05	8222.1\\
};
\addplot [color=mycolor34, forget plot]
  table[row sep=crcr]{%
0	4395.3\\
2000	4350.6\\
6000	3977.7\\
8000	3968.3\\
10000	4006.1\\
12000	4238.5\\
14000	4253.7\\
16000	4101.6\\
18000	4142\\
20000	4119.1\\
22000	4065.2\\
24000	4122.4\\
26000	3992.5\\
28000	4038.2\\
30000	4186.4\\
32000	4117.3\\
34000	3983\\
36000	3975.6\\
38000	3918.1\\
40000	4168.2\\
42000	4079.4\\
44000	4034.8\\
46000	4149.2\\
48000	4051.4\\
50000	4081.7\\
52000	4060.4\\
54000	4102.9\\
58000	4148.5\\
60000	4070.3\\
62000	3950.1\\
64000	4068.7\\
66000	3971\\
70000	4033.7\\
72000	4031.8\\
74000	4095.3\\
76000	4083.3\\
78000	4046.8\\
80000	4080.6\\
82000	3963.9\\
84000	4049.5\\
86000	4024\\
88000	4085.3\\
90000	4061.8\\
92000	4169.4\\
94000	4077\\
96000	4069.9\\
98000	4184.6\\
1e+05	3935.5\\
};
\addplot [color=mycolor35]
  table[row sep=crcr]{%
0	4395.3\\
2000	8447.4\\
4000	9950.6\\
6000	10100\\
10000	10107\\
12000	10244\\
14000	10065\\
16000	10029\\
18000	9931.8\\
20000	10072\\
22000	10023\\
24000	10204\\
26000	10179\\
28000	10074\\
30000	10240\\
32000	10194\\
34000	9999\\
36000	10043\\
38000	9985.3\\
40000	10018\\
42000	10268\\
44000	10076\\
46000	10274\\
48000	10155\\
50000	10228\\
52000	9933.8\\
54000	10291\\
58000	10125\\
60000	10070\\
62000	10142\\
66000	10086\\
68000	10158\\
70000	10177\\
72000	10029\\
78000	10038\\
80000	9880.5\\
82000	9952.8\\
84000	10181\\
86000	10248\\
90000	10118\\
92000	10185\\
94000	10126\\
96000	10184\\
98000	10284\\
1e+05	10274\\
};
\addlegendentry{\ClasskIILegend}

\addplot [color=mycolor36, forget plot]
  table[row sep=crcr]{%
0	4395.3\\
2000	4549\\
4000	4371.8\\
6000	4381.9\\
10000	4266.9\\
12000	4258.2\\
14000	4414.9\\
16000	4236.8\\
18000	4371.2\\
20000	4306.2\\
22000	4256.2\\
24000	4330.8\\
26000	4374\\
28000	4232.3\\
30000	4334.4\\
32000	4346.6\\
34000	4183.5\\
36000	4339\\
38000	4202.6\\
40000	4390.3\\
42000	4250.1\\
44000	4304.4\\
46000	4301.6\\
48000	4329.7\\
50000	4204.7\\
52000	4289.6\\
54000	4340.3\\
58000	4285.5\\
60000	4377.7\\
64000	4324.3\\
66000	4255.9\\
68000	4267.8\\
70000	4260.9\\
72000	4310.5\\
74000	4242\\
76000	4268.3\\
78000	4227.1\\
80000	4373.8\\
82000	4234.4\\
84000	4263.1\\
88000	4275.8\\
90000	4306\\
92000	4321.8\\
96000	4296.6\\
98000	4334.4\\
1e+05	4299.5\\
};
\addplot [color=mycolor37, forget plot]
  table[row sep=crcr]{%
0	4395.3\\
2000	4031.6\\
4000	3776.7\\
6000	3671.6\\
8000	3690.4\\
10000	3655.6\\
12000	3814.1\\
14000	3727.2\\
16000	3674.4\\
22000	3743.9\\
24000	3574.5\\
26000	3614\\
28000	3731.7\\
30000	3656.2\\
32000	3717.2\\
34000	3695.4\\
36000	3724.6\\
40000	3688.3\\
42000	3736.4\\
44000	3652.1\\
46000	3701.2\\
48000	3662.5\\
50000	3676.1\\
52000	3622.5\\
54000	3703.3\\
56000	3691.5\\
58000	3701.7\\
60000	3693.9\\
64000	3755.3\\
66000	3660\\
68000	3680.8\\
70000	3669.2\\
72000	3699.6\\
76000	3718.3\\
78000	3654\\
80000	3708.6\\
82000	3678.4\\
84000	3628\\
86000	3645.1\\
88000	3703.8\\
90000	3677.5\\
92000	3692\\
94000	3686.5\\
96000	3727.7\\
98000	3684.2\\
1e+05	3746.7\\
};
\addplot [color=mycolor38, forget plot]
  table[row sep=crcr]{%
0	4395.3\\
2000	3558.1\\
4000	3104\\
6000	3155.4\\
8000	3001.5\\
10000	3026.5\\
12000	2994.8\\
14000	3047.9\\
16000	2892.9\\
18000	2981.3\\
20000	3180\\
22000	3130.5\\
24000	3033.6\\
26000	3087.6\\
28000	3068.9\\
30000	3076.8\\
32000	3133.1\\
34000	3148.2\\
36000	3123.8\\
38000	3131\\
40000	3022.3\\
42000	3122.6\\
44000	3105.6\\
46000	3001.9\\
48000	3087.1\\
50000	3135.6\\
52000	3101.4\\
54000	3034.8\\
56000	3113.3\\
58000	3039.7\\
62000	3142.4\\
64000	3116.1\\
66000	3101.3\\
68000	3127.4\\
70000	3085\\
72000	3215.9\\
74000	3161.6\\
76000	3017.1\\
78000	3040.2\\
80000	2959.8\\
82000	3035.8\\
84000	3081.3\\
86000	3155.8\\
88000	3050.4\\
90000	3057\\
92000	3023.9\\
94000	3065.7\\
96000	3050.5\\
98000	3089.1\\
1e+05	3109.3\\
};
\addplot [color=mycolor39, forget plot]
  table[row sep=crcr]{%
0	4395.3\\
2000	9407.6\\
4000	11621\\
6000	11535\\
8000	11964\\
10000	11415\\
12000	11372\\
14000	11811\\
16000	11397\\
18000	11746\\
20000	11589\\
22000	11657\\
26000	11691\\
28000	11655\\
30000	11975\\
32000	11622\\
34000	11413\\
36000	11374\\
38000	11711\\
40000	11910\\
42000	11398\\
44000	11170\\
46000	11338\\
48000	11203\\
50000	11826\\
52000	11877\\
56000	11747\\
58000	11658\\
60000	11865\\
62000	11509\\
64000	11893\\
66000	11755\\
68000	11765\\
70000	11948\\
72000	11500\\
74000	11758\\
76000	11945\\
78000	11484\\
80000	11856\\
82000	11946\\
84000	11804\\
86000	11923\\
88000	11466\\
90000	11338\\
92000	11518\\
94000	11616\\
96000	11922\\
98000	11896\\
1e+05	12177\\
};
\addplot [color=mycolor40, forget plot]
  table[row sep=crcr]{%
0	4395.3\\
2000	5423.7\\
4000	5735.3\\
6000	5840.8\\
8000	5756.7\\
10000	5867\\
12000	5841.7\\
14000	5889.1\\
16000	5740.1\\
18000	5740.9\\
20000	5808.2\\
22000	5749.7\\
24000	5809.8\\
26000	5708.7\\
28000	5803.2\\
32000	5650.9\\
34000	5704.6\\
36000	5863.9\\
38000	5767.5\\
40000	5740.2\\
42000	5782\\
44000	5800.2\\
46000	5875.8\\
48000	5864.7\\
50000	5832.6\\
52000	5769.3\\
54000	5769.2\\
56000	5877.9\\
58000	5853.1\\
60000	5753\\
62000	5857.9\\
64000	5784.6\\
66000	5792.8\\
68000	5677.6\\
70000	5794.6\\
72000	5837.8\\
74000	5722.1\\
78000	5794.4\\
80000	5777.5\\
82000	5721.6\\
84000	5816.4\\
86000	5735.1\\
88000	5789\\
90000	5755.4\\
92000	5844.1\\
96000	5784.3\\
98000	5733.8\\
1e+05	5842.3\\
};
\addplot [color=mycolor41, forget plot]
  table[row sep=crcr]{%
0	4395.3\\
2000	5324.8\\
4000	5493.8\\
6000	5467.4\\
8000	5484.7\\
10000	5402.7\\
14000	5446.3\\
16000	5255.9\\
18000	5290.4\\
20000	5255.6\\
22000	5330.8\\
24000	5277.9\\
26000	5365.7\\
28000	5336.8\\
30000	5454\\
32000	5347.1\\
34000	5308.2\\
36000	5365.3\\
38000	5391.8\\
40000	5481.6\\
44000	5282.9\\
46000	5338.1\\
48000	5483\\
52000	5388.8\\
54000	5527\\
56000	5424.1\\
58000	5447.4\\
60000	5492.1\\
62000	5475.1\\
66000	5319\\
70000	5347\\
72000	5464.7\\
74000	5395.8\\
76000	5456.2\\
78000	5380\\
80000	5375.8\\
82000	5263.4\\
84000	5485.9\\
86000	5392.6\\
88000	5363.8\\
90000	5536.2\\
92000	5538.2\\
94000	5467.9\\
96000	5478.6\\
1e+05	5341.4\\
};
\addplot [color=mycolor42, forget plot]
  table[row sep=crcr]{%
0	4395.3\\
2000	12790\\
4000	19010\\
6000	19718\\
8000	20872\\
10000	21236\\
12000	19869\\
14000	19299\\
16000	19646\\
18000	19308\\
20000	19836\\
22000	19905\\
24000	20237\\
26000	20263\\
28000	19795\\
30000	20466\\
32000	20818\\
34000	20011\\
36000	20004\\
38000	20710\\
40000	20741\\
42000	19925\\
44000	20814\\
46000	19921\\
48000	20264\\
50000	20041\\
52000	20162\\
54000	20604\\
56000	20410\\
62000	20552\\
64000	20093\\
66000	20235\\
68000	20264\\
70000	19929\\
74000	20144\\
76000	20387\\
78000	20155\\
80000	20183\\
82000	20029\\
84000	20382\\
86000	19776\\
88000	20323\\
90000	20016\\
92000	19374\\
94000	19662\\
96000	20132\\
98000	20128\\
1e+05	19846\\
};
\addplot [color=mycolor43, forget plot]
  table[row sep=crcr]{%
0	4395.3\\
2000	4880.4\\
4000	4965.6\\
6000	4963.8\\
8000	4924.8\\
10000	4809.3\\
12000	4898.7\\
14000	4900.3\\
16000	4944.2\\
18000	4907\\
20000	5040.1\\
22000	5024.7\\
24000	4920.2\\
26000	4911.5\\
28000	4944.8\\
30000	4936.7\\
32000	4959.1\\
34000	5002.5\\
36000	4958.6\\
42000	4963.1\\
44000	4896.7\\
46000	4999.4\\
50000	4987.7\\
52000	4936.6\\
54000	4979.3\\
56000	4998.2\\
58000	4990.4\\
60000	5039.1\\
62000	4921.3\\
64000	4851.2\\
66000	4998.9\\
68000	5011.8\\
70000	4964.4\\
74000	5051\\
78000	4940.4\\
80000	5054.4\\
82000	4934.4\\
84000	4939.9\\
88000	5126\\
90000	4999.5\\
92000	4964\\
94000	4853.4\\
96000	4912.8\\
98000	5018.3\\
1e+05	4972.4\\
};
\addplot [color=mycolor44, forget plot]
  table[row sep=crcr]{%
0	4395.3\\
2000	3710.9\\
4000	3525.3\\
6000	3452.2\\
10000	3470.9\\
14000	3280\\
16000	3368.4\\
18000	3273.1\\
20000	3408.6\\
22000	3411\\
24000	3390.4\\
26000	3455.8\\
30000	3404\\
32000	3373.6\\
34000	3447.1\\
36000	3456.9\\
38000	3413.2\\
40000	3349.2\\
42000	3368\\
44000	3240.7\\
46000	3385.8\\
50000	3395.6\\
52000	3375.2\\
54000	3503.4\\
56000	3381.5\\
58000	3329.6\\
60000	3443.2\\
62000	3404.8\\
64000	3337.6\\
66000	3358.5\\
68000	3307.9\\
70000	3436.4\\
72000	3318.4\\
74000	3398.9\\
76000	3491.1\\
78000	3503.9\\
80000	3382.5\\
84000	3399.1\\
86000	3378.6\\
88000	3301.2\\
92000	3426.5\\
94000	3386\\
96000	3316\\
98000	3339.6\\
1e+05	3411.8\\
};
\addplot [color=mycolor45, forget plot]
  table[row sep=crcr]{%
0	4395.3\\
2000	6280.3\\
4000	6869.9\\
6000	6903.8\\
8000	6876\\
10000	6926.8\\
12000	6900.3\\
16000	7043.5\\
18000	6924.4\\
20000	6878.5\\
22000	6897.3\\
24000	6794.9\\
26000	6847.5\\
32000	7150.5\\
34000	7104.3\\
36000	6982.6\\
40000	6858.2\\
42000	6858.2\\
44000	7005.1\\
46000	7031.5\\
50000	6878.6\\
52000	6967.7\\
56000	6815.7\\
58000	6981.6\\
60000	7057.8\\
62000	6996.5\\
64000	6870.4\\
66000	6939.8\\
68000	6850.4\\
70000	6865.5\\
72000	6913.4\\
74000	7119.6\\
76000	7030.2\\
78000	7129\\
80000	6894.9\\
82000	7076.9\\
84000	7000.9\\
86000	6828.5\\
88000	6903.8\\
90000	6853.3\\
92000	6853.6\\
94000	6897\\
96000	6810.1\\
98000	6875.5\\
1e+05	6897.3\\
};
\addplot [color=mycolor46, forget plot]
  table[row sep=crcr]{%
0	4395.3\\
2000	4261.8\\
4000	4145.2\\
6000	4083.1\\
8000	4195\\
10000	4088.3\\
14000	4031.9\\
16000	4111.2\\
18000	4133.8\\
20000	4100.2\\
22000	4128.1\\
24000	4016.9\\
26000	4116.8\\
28000	4167.6\\
30000	4152.1\\
32000	4164.2\\
34000	4152.1\\
36000	4106.7\\
38000	4226.2\\
40000	4094.5\\
42000	4193.9\\
44000	4052.9\\
46000	4146.2\\
48000	4152.3\\
50000	4136.1\\
52000	4225.7\\
54000	4116.2\\
56000	4139.3\\
58000	4090.1\\
60000	4199.6\\
62000	4065.7\\
64000	4029.1\\
66000	4065.8\\
68000	4119.6\\
70000	4154.4\\
72000	3976.7\\
74000	4061.6\\
76000	4123.8\\
78000	4224.4\\
80000	4235.9\\
82000	4110.2\\
84000	4154.3\\
86000	4108.8\\
88000	4136.2\\
90000	4151.7\\
92000	4229.4\\
94000	4074.6\\
96000	4094.5\\
98000	4148.5\\
1e+05	4136.9\\
};
\addplot [color=mycolor47, forget plot]
  table[row sep=crcr]{%
0	4395.3\\
2000	5169.5\\
4000	5367.5\\
6000	5298.1\\
8000	5304.8\\
10000	5347.6\\
12000	5284.9\\
14000	5278.7\\
16000	5152.4\\
18000	5174.5\\
20000	5332.2\\
26000	5344.7\\
28000	5243.4\\
30000	5317.9\\
32000	5209.8\\
34000	5219.4\\
36000	5142.6\\
38000	5177.1\\
40000	5469.8\\
42000	5339.2\\
44000	5433\\
46000	5233.8\\
48000	5328.9\\
50000	5312.8\\
52000	5378.1\\
54000	5383.1\\
56000	5249.1\\
58000	5288.5\\
60000	5278.8\\
62000	5230.9\\
64000	5484.2\\
66000	5481.5\\
68000	5342.4\\
70000	5339.1\\
72000	5284.8\\
74000	5301.6\\
78000	5377.2\\
80000	5274.3\\
82000	5326.2\\
84000	5324.6\\
86000	5211\\
92000	5251.5\\
94000	5348.4\\
96000	5258.6\\
98000	5343.1\\
1e+05	5334.4\\
};
\addplot [color=mycolor48, forget plot]
  table[row sep=crcr]{%
0	4395.3\\
2000	3403.3\\
4000	3042.5\\
6000	2961.5\\
8000	2932.9\\
10000	2940.6\\
12000	3020\\
14000	3026.4\\
16000	2965.8\\
18000	2936.6\\
22000	2995.3\\
24000	2939.9\\
26000	2992.3\\
30000	2986.9\\
32000	2922.9\\
34000	2987.3\\
36000	2940.1\\
38000	2986.4\\
40000	2998.2\\
42000	2937.9\\
44000	2960.6\\
46000	2993\\
48000	2948.6\\
50000	2947.9\\
52000	2958.3\\
56000	2992\\
58000	2931.7\\
60000	2924.7\\
62000	3021.4\\
64000	2909.3\\
66000	2952.6\\
70000	2999.7\\
72000	2987.9\\
74000	2947.9\\
76000	2955.4\\
78000	3009.2\\
80000	2931.5\\
84000	2948.5\\
86000	2983.8\\
88000	2993.2\\
90000	2965\\
92000	2975.4\\
96000	2953.5\\
98000	2908.6\\
1e+05	2990.9\\
};
\addplot [color=mycolor49, forget plot]
  table[row sep=crcr]{%
0	4395.3\\
2000	4643.9\\
4000	4438.3\\
6000	4367.1\\
8000	4414.3\\
10000	4276\\
12000	4428.4\\
14000	4344.4\\
16000	4352.1\\
18000	4387.6\\
20000	4515.6\\
22000	4264.9\\
24000	4371\\
26000	4458.5\\
28000	4188.5\\
30000	4398.3\\
32000	4390.7\\
34000	4359.8\\
36000	4359.2\\
38000	4312.9\\
40000	4472.3\\
44000	4365.9\\
48000	4269.1\\
50000	4337.6\\
52000	4257.6\\
54000	4451.3\\
56000	4342.5\\
58000	4310.1\\
60000	4406.1\\
62000	4380.7\\
64000	4397.4\\
66000	4344.4\\
68000	4261.1\\
70000	4410.8\\
72000	4481.7\\
74000	4322.4\\
76000	4359.8\\
78000	4353.8\\
80000	4498\\
82000	4475.1\\
84000	4401.5\\
86000	4285.5\\
88000	4550.1\\
90000	4365.3\\
94000	4333.4\\
96000	4265.1\\
98000	4258.8\\
1e+05	4330.3\\
};
\addplot [color=mycolor50, forget plot]
  table[row sep=crcr]{%
0	4395.3\\
2000	3006.4\\
4000	2578.9\\
6000	2532.9\\
8000	2551.3\\
10000	2485.7\\
12000	2505.1\\
16000	2492.3\\
18000	2527.3\\
20000	2554.2\\
22000	2485.5\\
24000	2525.9\\
26000	2485.7\\
28000	2504.5\\
36000	2511.3\\
38000	2481.6\\
40000	2502.4\\
42000	2541.5\\
44000	2523.5\\
46000	2549.4\\
48000	2532.7\\
50000	2537.6\\
52000	2490.8\\
54000	2565.8\\
56000	2495.6\\
58000	2513.7\\
60000	2486.2\\
62000	2527.4\\
64000	2506.5\\
66000	2537.5\\
68000	2581.2\\
70000	2509.8\\
72000	2493\\
74000	2567.4\\
76000	2505.7\\
78000	2475.6\\
80000	2564.8\\
82000	2532.5\\
86000	2612.4\\
88000	2563.9\\
90000	2619.1\\
92000	2513.9\\
94000	2505.8\\
96000	2580.9\\
98000	2490\\
1e+05	2510.1\\
};
\addplot [color=mycolor51, forget plot]
  table[row sep=crcr]{%
0	4395.3\\
2000	5663.1\\
4000	6125.1\\
6000	5921.1\\
8000	5903.3\\
12000	5842.8\\
14000	5992.3\\
16000	5988.2\\
18000	5959.8\\
20000	6063.1\\
22000	6032\\
24000	6031\\
26000	5916.9\\
28000	5876.8\\
30000	6007.8\\
32000	5845.3\\
34000	5873.7\\
36000	5961.8\\
38000	5886.2\\
40000	6018.8\\
42000	6036.6\\
44000	5770.4\\
46000	6010.3\\
48000	5768.6\\
50000	5947.4\\
52000	5902.9\\
54000	5924.6\\
56000	6082.5\\
58000	6226\\
60000	6143\\
62000	5993.2\\
64000	6154.6\\
66000	6132.1\\
68000	5847.1\\
70000	6030\\
72000	6000.3\\
74000	5749.8\\
76000	5852.1\\
78000	6023.5\\
80000	5833.8\\
82000	5865.2\\
84000	6045.1\\
86000	6063.1\\
88000	6215.1\\
90000	5998.6\\
92000	5980.6\\
94000	5819\\
96000	5846.8\\
98000	6257\\
1e+05	5999.5\\
};
\addplot [color=mycolor52, forget plot]
  table[row sep=crcr]{%
0	4395.3\\
2000	4007.9\\
4000	3620.3\\
6000	3476.4\\
8000	3424.6\\
10000	3487.8\\
12000	3647.7\\
14000	3584.8\\
16000	3423.9\\
18000	3464.9\\
20000	3430.5\\
22000	3508.9\\
24000	3476.3\\
26000	3573.6\\
30000	3450.8\\
32000	3575.3\\
34000	3558.9\\
36000	3483.8\\
38000	3540.1\\
40000	3554.4\\
42000	3546.2\\
44000	3624.6\\
48000	3426.1\\
50000	3485.9\\
52000	3534.8\\
54000	3428.4\\
56000	3403.5\\
58000	3627.5\\
60000	3723.1\\
64000	3549.2\\
66000	3665.8\\
70000	3489.6\\
72000	3534.5\\
74000	3522.3\\
76000	3474.4\\
78000	3480.9\\
80000	3624.2\\
82000	3611.7\\
84000	3497.2\\
86000	3579.3\\
88000	3499.9\\
90000	3613.4\\
92000	3586.9\\
94000	3504.8\\
96000	3505.5\\
98000	3583.6\\
1e+05	3467.5\\
};
\addplot [color=mycolor53, forget plot]
  table[row sep=crcr]{%
0	4395.3\\
2000	5300.8\\
4000	5706\\
6000	5608.1\\
8000	5703.5\\
12000	5741.5\\
14000	5729.7\\
16000	5527.8\\
18000	5702.2\\
20000	5704.8\\
22000	5838.1\\
24000	5769.7\\
26000	5588.7\\
28000	5780\\
30000	5758.8\\
32000	5808.7\\
34000	5639.4\\
36000	5714.7\\
38000	5921.5\\
40000	5926.1\\
42000	5725\\
44000	5887.5\\
46000	5858.4\\
48000	5812.8\\
50000	5785.4\\
52000	5849.8\\
54000	5684.2\\
56000	5599.6\\
58000	5685.6\\
60000	5632.4\\
62000	5560.8\\
64000	5766.7\\
66000	5823.6\\
68000	5967.3\\
70000	5859.2\\
72000	5936.7\\
76000	5870.8\\
78000	5804\\
80000	5771.8\\
84000	5645.5\\
86000	5748.3\\
88000	5785.8\\
90000	5706\\
92000	5786.8\\
94000	5523.1\\
96000	5590.6\\
98000	5562.7\\
1e+05	5804\\
};
\addplot [color=mycolor54, forget plot]
  table[row sep=crcr]{%
0	4395.3\\
2000	4610.9\\
4000	4674.2\\
6000	4528.2\\
8000	4598.6\\
10000	4495.3\\
12000	4681.9\\
14000	4669.7\\
16000	4687.5\\
20000	4457.8\\
22000	4611.8\\
24000	4551\\
26000	4668.3\\
28000	4677.4\\
30000	4554\\
32000	4675.8\\
34000	4631.5\\
40000	4589.5\\
44000	4613.8\\
50000	4598.8\\
52000	4634.9\\
54000	4696.5\\
56000	4628.8\\
58000	4529.7\\
60000	4587.9\\
62000	4632.7\\
64000	4717.1\\
66000	4734\\
68000	4512.9\\
70000	4643.9\\
72000	4647.7\\
74000	4682.1\\
76000	4672\\
78000	4712.2\\
80000	4677.2\\
82000	4706.1\\
86000	4571.6\\
88000	4564.7\\
90000	4626.2\\
92000	4629.2\\
94000	4604.5\\
96000	4652.9\\
98000	4579.8\\
1e+05	4488.8\\
};
\addplot [color=mycolor55, forget plot]
  table[row sep=crcr]{%
0	4395.3\\
2000	5376.5\\
4000	5744.1\\
6000	5800.2\\
10000	5650.8\\
12000	5826.3\\
16000	5713.6\\
18000	5761.6\\
20000	5852.5\\
22000	5794.8\\
24000	5777.7\\
26000	5843.3\\
28000	5797.2\\
32000	5853.4\\
34000	5814.7\\
36000	5747.8\\
40000	5731.8\\
42000	5847.5\\
44000	5792.2\\
46000	5703.4\\
48000	5885.1\\
50000	5763.2\\
52000	5856.2\\
54000	5789.9\\
56000	5842\\
58000	5727.4\\
60000	5748.9\\
62000	5822.2\\
64000	5735.3\\
66000	5694.7\\
68000	5628.9\\
70000	5679.8\\
72000	5840.6\\
74000	5842.2\\
76000	5710.3\\
78000	5771.9\\
80000	5636.3\\
82000	5771\\
84000	5784.1\\
86000	5677.3\\
88000	5903.9\\
90000	5838.1\\
92000	5828\\
94000	5920.4\\
96000	5825.5\\
98000	5802.7\\
1e+05	5667.1\\
};
\addplot [color=mycolor56, forget plot]
  table[row sep=crcr]{%
0	4395.3\\
2000	4218.8\\
6000	4032.4\\
8000	4029.8\\
10000	4066\\
12000	4049.4\\
14000	4110.4\\
16000	3950.5\\
18000	3991.7\\
20000	3952.8\\
22000	4103.7\\
24000	4112.7\\
26000	4017.9\\
28000	4008.2\\
30000	3954.1\\
32000	4006.3\\
34000	3976.8\\
38000	4073.3\\
40000	4104.8\\
42000	4023.6\\
44000	4115.7\\
46000	3997\\
48000	4001.8\\
50000	4028.5\\
52000	3992.4\\
54000	4015.5\\
56000	4001.8\\
58000	4020\\
62000	4020.9\\
64000	4071.4\\
66000	4040\\
68000	4036.3\\
70000	4015.4\\
72000	4050\\
74000	4018.3\\
76000	4083.1\\
78000	4066.7\\
80000	3963.3\\
82000	3984\\
84000	4032.2\\
88000	4026.2\\
92000	4091.3\\
94000	4065.3\\
96000	4072\\
98000	4053.6\\
1e+05	4048.9\\
};
\addplot [color=mycolor57, forget plot]
  table[row sep=crcr]{%
0	4395.3\\
2000	5675.9\\
4000	6549.7\\
6000	6651.6\\
8000	6651.3\\
10000	6672.8\\
12000	6621.4\\
14000	6667.6\\
16000	6762.4\\
18000	6735.5\\
20000	6734.4\\
22000	6640.4\\
24000	6607.2\\
26000	6793.8\\
28000	6850\\
30000	6734.2\\
32000	6728.1\\
34000	6748.2\\
36000	6592.6\\
38000	6642.3\\
42000	6798.6\\
44000	6701\\
46000	6567.3\\
48000	6611.5\\
50000	6712.5\\
52000	6733.2\\
54000	6658.8\\
56000	6618.5\\
58000	6541.5\\
60000	6659.9\\
64000	6763.5\\
66000	6641.4\\
70000	6678.3\\
76000	6695.7\\
78000	6721.3\\
80000	6847.2\\
82000	6780.8\\
84000	6685.5\\
86000	6662.2\\
88000	6719.8\\
92000	6661\\
94000	6727.7\\
96000	6628.3\\
98000	6698.3\\
1e+05	6693.9\\
};
\addplot [color=mycolor58, forget plot]
  table[row sep=crcr]{%
0	4395.3\\
2000	4013.4\\
4000	3681.3\\
6000	3606.8\\
8000	3669\\
10000	3572.6\\
12000	3631.1\\
14000	3562.6\\
16000	3620.4\\
18000	3559.2\\
20000	3639\\
22000	3570.1\\
24000	3622.6\\
26000	3543.4\\
28000	3609.3\\
30000	3741.1\\
32000	3609.3\\
36000	3614.4\\
38000	3630.3\\
44000	3534.2\\
46000	3521.7\\
48000	3634.7\\
50000	3651.6\\
52000	3550.3\\
54000	3624.8\\
56000	3625.6\\
58000	3714.4\\
60000	3644.8\\
62000	3487.4\\
64000	3661.8\\
68000	3629.3\\
70000	3688.4\\
72000	3665.9\\
74000	3672.7\\
76000	3645.5\\
78000	3602.3\\
80000	3616.9\\
82000	3657\\
84000	3647.8\\
86000	3582\\
88000	3549.8\\
90000	3706.2\\
92000	3587.2\\
96000	3555.9\\
98000	3630.7\\
1e+05	3632.1\\
};
\addplot [color=mycolor59, forget plot]
  table[row sep=crcr]{%
0	4395.3\\
2000	4855.3\\
4000	5040.8\\
6000	5039.8\\
8000	5097.9\\
10000	5117.8\\
12000	5115.6\\
14000	5074.5\\
16000	5015.8\\
18000	5046.3\\
20000	5212.6\\
22000	5055.3\\
24000	5047.8\\
26000	5070.4\\
28000	5008.6\\
30000	5015.1\\
32000	5125.8\\
34000	5090.2\\
36000	5138.2\\
38000	5106.8\\
40000	4941.8\\
42000	4944.1\\
44000	5132.6\\
46000	5087.2\\
48000	5006.8\\
52000	5055.1\\
54000	5120.4\\
56000	5041.5\\
58000	5081.8\\
60000	5021.4\\
64000	5079.1\\
66000	5111.8\\
68000	5065\\
70000	5074.8\\
72000	5158.1\\
74000	5039.3\\
76000	4992.6\\
78000	5158.5\\
80000	5220.1\\
82000	5128.5\\
84000	4981.4\\
86000	5062.7\\
88000	4982.5\\
90000	4926.1\\
92000	5018.2\\
94000	5041.1\\
98000	5170.4\\
1e+05	5099.8\\
};
\addplot [color=mycolor60, forget plot]
  table[row sep=crcr]{%
0	4395.3\\
2000	3855.1\\
4000	3611.2\\
6000	3570.1\\
10000	3665.5\\
12000	3607.8\\
14000	3649.3\\
16000	3600.7\\
18000	3522.6\\
20000	3652.7\\
22000	3615\\
24000	3706.1\\
26000	3609\\
28000	3573.9\\
30000	3584.6\\
32000	3579.7\\
34000	3596.9\\
36000	3574\\
38000	3630.1\\
42000	3523.1\\
44000	3565\\
46000	3496.5\\
48000	3583.4\\
50000	3572.4\\
52000	3530.8\\
54000	3615.5\\
56000	3556.9\\
58000	3541.2\\
60000	3569\\
62000	3628.4\\
64000	3580.5\\
66000	3653.7\\
68000	3604.4\\
70000	3652.1\\
72000	3590\\
74000	3574.3\\
76000	3519.3\\
78000	3582.2\\
80000	3604.9\\
82000	3614.1\\
84000	3611.5\\
86000	3570.4\\
88000	3616.6\\
90000	3601.4\\
92000	3523.9\\
94000	3582.5\\
96000	3504.2\\
98000	3560.7\\
1e+05	3554.7\\
};
\addplot [color=mycolor61, forget plot]
  table[row sep=crcr]{%
0	4395.3\\
2000	3578.7\\
4000	3204.1\\
6000	3175.3\\
8000	3161\\
10000	3181.8\\
12000	3189.4\\
14000	3153\\
16000	3193.3\\
18000	3139.5\\
20000	3156.4\\
22000	3188.1\\
24000	3146.7\\
26000	3144.2\\
28000	3191.9\\
30000	3192.9\\
32000	3153.6\\
34000	3167.5\\
36000	3154.3\\
38000	3187.1\\
42000	3183.5\\
44000	3156.8\\
46000	3184.3\\
48000	3170.9\\
50000	3130.1\\
52000	3187.2\\
54000	3184.1\\
56000	3205.2\\
58000	3194.5\\
60000	3155.1\\
62000	3146.5\\
64000	3181.3\\
66000	3195.1\\
68000	3169\\
70000	3185.5\\
72000	3151.5\\
74000	3144.8\\
76000	3172.1\\
78000	3155.5\\
80000	3178.7\\
82000	3124.5\\
84000	3134.6\\
86000	3191.7\\
88000	3171\\
90000	3184.4\\
92000	3152.7\\
98000	3140.7\\
1e+05	3178\\
};
\addplot [color=mycolor62, forget plot]
  table[row sep=crcr]{%
0	4395.3\\
2000	4797.7\\
4000	5010.9\\
6000	5005.3\\
8000	5094.4\\
12000	5008.5\\
14000	5099.1\\
16000	4979.8\\
20000	5036.3\\
22000	5025.3\\
24000	5030.7\\
26000	4954.7\\
28000	5105.5\\
30000	5048.8\\
32000	5079.7\\
34000	4985.8\\
36000	5013.3\\
38000	4976\\
40000	5012.4\\
42000	5084.2\\
44000	4998\\
46000	5084.9\\
48000	5013.9\\
50000	5055.3\\
52000	5035.8\\
54000	4957.4\\
56000	5025.7\\
58000	4996\\
60000	4991.9\\
62000	5099.8\\
64000	5019.7\\
66000	5000.8\\
68000	5007.9\\
70000	5029.3\\
72000	5070.7\\
74000	4961.4\\
76000	4986.6\\
78000	5106.4\\
80000	5056.6\\
82000	4971.7\\
84000	4983.7\\
86000	5047.2\\
90000	5014.9\\
92000	5015.6\\
94000	4981.4\\
96000	5000.6\\
98000	4938.2\\
1e+05	4959\\
};
\addplot [color=mycolor63, forget plot]
  table[row sep=crcr]{%
0	4395.3\\
2000	3073.2\\
4000	2548\\
6000	2482.3\\
8000	2462.9\\
10000	2469.1\\
12000	2466.7\\
14000	2489.5\\
16000	2442.5\\
18000	2463.9\\
20000	2508.5\\
22000	2408.7\\
24000	2463.3\\
26000	2465.4\\
28000	2493.7\\
30000	2440.8\\
32000	2422\\
34000	2434.3\\
36000	2470.6\\
38000	2441.1\\
40000	2425.3\\
42000	2436.8\\
44000	2471.6\\
46000	2451\\
50000	2467.9\\
52000	2449.1\\
54000	2445.4\\
56000	2413.5\\
58000	2459.8\\
60000	2476.5\\
62000	2459.7\\
64000	2472.7\\
66000	2515.9\\
68000	2460.4\\
70000	2473.7\\
72000	2449.3\\
74000	2449.6\\
76000	2418.8\\
78000	2438\\
82000	2398.5\\
84000	2458\\
86000	2483.8\\
88000	2419.5\\
90000	2424.7\\
92000	2398\\
94000	2415\\
96000	2404.2\\
1e+05	2456.8\\
};
\addplot [color=mycolor64, forget plot]
  table[row sep=crcr]{%
0	4395.3\\
2000	3491.7\\
4000	3163.4\\
6000	3078.8\\
8000	3054.2\\
10000	3089.1\\
16000	3044.9\\
18000	3070.7\\
20000	3065\\
22000	3050\\
24000	3119.3\\
26000	3113.8\\
28000	3071.1\\
30000	3050.7\\
32000	3058.3\\
34000	3043.6\\
36000	3050\\
38000	3088.1\\
40000	3074.7\\
42000	3073.4\\
44000	3057.4\\
46000	3092.7\\
48000	3094.7\\
50000	3113.7\\
52000	3048.1\\
54000	3101.1\\
56000	3076\\
60000	3105.9\\
62000	3057.4\\
64000	3084.4\\
70000	3087.8\\
72000	3126\\
74000	3116.7\\
78000	3073.9\\
80000	3118.7\\
84000	3041\\
86000	3074\\
88000	3077\\
90000	3097.5\\
92000	3047.3\\
94000	3062.4\\
96000	3102.3\\
98000	3096.4\\
1e+05	3135.9\\
};
\addplot [color=mycolor65, forget plot]
  table[row sep=crcr]{%
0	4395.3\\
2000	3134.5\\
4000	2673.4\\
6000	2614.2\\
8000	2604.8\\
10000	2582.5\\
12000	2595.5\\
14000	2623.2\\
16000	2628.4\\
18000	2572.2\\
20000	2606.2\\
24000	2583.1\\
26000	2608.1\\
28000	2564.6\\
30000	2534.6\\
32000	2586.9\\
36000	2561.3\\
40000	2611.2\\
42000	2575.2\\
44000	2613\\
46000	2609.7\\
48000	2578.6\\
50000	2627.4\\
52000	2627.4\\
54000	2564.9\\
56000	2580.8\\
58000	2607.7\\
60000	2659.8\\
62000	2597.1\\
64000	2580.2\\
68000	2575.8\\
70000	2595.5\\
72000	2603.6\\
76000	2574.8\\
78000	2500.3\\
80000	2584.8\\
82000	2574.5\\
84000	2601.6\\
86000	2611.4\\
88000	2579.8\\
90000	2605.7\\
92000	2559.7\\
94000	2598.2\\
96000	2556.1\\
98000	2565.2\\
1e+05	2561.9\\
};
\addplot [color=mycolor66, forget plot]
  table[row sep=crcr]{%
0	4395.3\\
2000	3745\\
4000	3497.8\\
6000	3481.6\\
8000	3502\\
10000	3509.6\\
12000	3502.1\\
14000	3455.3\\
16000	3491.8\\
20000	3515\\
22000	3496.3\\
24000	3461.3\\
26000	3482.9\\
28000	3462.3\\
30000	3488.7\\
32000	3451.8\\
34000	3456.2\\
36000	3445.4\\
38000	3446.5\\
40000	3501.8\\
42000	3526.1\\
44000	3452.6\\
46000	3448.1\\
48000	3463\\
50000	3437.4\\
52000	3471.9\\
54000	3477.8\\
60000	3414.8\\
62000	3466.3\\
64000	3450.3\\
66000	3376.3\\
68000	3386.2\\
70000	3422.5\\
72000	3477.4\\
74000	3473.8\\
76000	3480.8\\
78000	3462\\
80000	3466.9\\
82000	3461.9\\
84000	3390.1\\
86000	3475.6\\
88000	3416.7\\
90000	3396.9\\
92000	3433.4\\
94000	3431.3\\
96000	3470.5\\
98000	3464.5\\
1e+05	3527.3\\
};
\addplot [color=mycolor67, forget plot]
  table[row sep=crcr]{%
0	4395.3\\
2000	3705.5\\
4000	3346.4\\
6000	3187.5\\
8000	3258.3\\
12000	3249.9\\
14000	3261\\
16000	3211.2\\
18000	3226.7\\
20000	3201.6\\
22000	3251.6\\
24000	3262.2\\
28000	3138.6\\
30000	3165.1\\
32000	3230.5\\
34000	3220.8\\
36000	3181\\
38000	3252.7\\
40000	3270.5\\
42000	3235.7\\
44000	3181.9\\
46000	3235.9\\
48000	3273.1\\
50000	3298.2\\
52000	3237.9\\
54000	3216.1\\
56000	3207.7\\
58000	3242.3\\
60000	3248.2\\
62000	3217.8\\
64000	3201.1\\
66000	3227\\
68000	3180.6\\
70000	3245.7\\
74000	3193.3\\
76000	3213.9\\
78000	3250.5\\
80000	3205.4\\
86000	3266.6\\
88000	3204.3\\
92000	3227.4\\
94000	3209\\
96000	3173.4\\
98000	3235.6\\
1e+05	3247.2\\
};
\addplot [color=mycolor67]
  table[row sep=crcr]{%
0	4395.3\\
2000	3215.4\\
4000	2700.4\\
6000	2652.6\\
8000	2648.3\\
12000	2670\\
16000	2702\\
18000	2687.9\\
20000	2663.3\\
22000	2672.8\\
24000	2707.1\\
26000	2668.1\\
28000	2693.3\\
30000	2674.4\\
32000	2681.7\\
36000	2713.1\\
38000	2673.2\\
40000	2677.4\\
42000	2668.3\\
44000	2694.1\\
46000	2649.4\\
48000	2669.1\\
50000	2606.5\\
54000	2712.6\\
56000	2672.5\\
58000	2656.1\\
60000	2669.5\\
62000	2668.2\\
64000	2650.6\\
66000	2675.8\\
68000	2668.3\\
70000	2674.7\\
72000	2668.8\\
74000	2699\\
76000	2713.2\\
78000	2685.4\\
80000	2647.4\\
82000	2633.8\\
84000	2631.7\\
86000	2649.4\\
88000	2658.3\\
90000	2689.6\\
92000	2687.1\\
94000	2628.9\\
96000	2645.6\\
1e+05	2711.3\\
};
\addlegendentry{\ClasskIIILegend}

\addplot [color=mycolor68, forget plot]
  table[row sep=crcr]{%
0	4395.3\\
2000	3159.8\\
4000	2599.7\\
6000	2513.1\\
8000	2501.6\\
12000	2520.2\\
14000	2503.4\\
16000	2520.8\\
20000	2497.2\\
22000	2526.4\\
28000	2532.4\\
30000	2515.8\\
32000	2512.7\\
34000	2494.2\\
36000	2499.7\\
38000	2486.1\\
40000	2513\\
42000	2516.7\\
46000	2479.8\\
48000	2517.4\\
50000	2496.2\\
52000	2503.2\\
54000	2523.7\\
56000	2475.9\\
58000	2503.5\\
60000	2480.7\\
62000	2471.8\\
64000	2476\\
66000	2494.9\\
68000	2494.2\\
70000	2516.6\\
72000	2488.4\\
74000	2507.7\\
76000	2479.8\\
78000	2511.2\\
80000	2477.7\\
84000	2520.3\\
88000	2515.4\\
90000	2552.1\\
94000	2486.1\\
98000	2505.9\\
1e+05	2487.9\\
};
\addplot [color=mycolor69, forget plot]
  table[row sep=crcr]{%
0	4395.3\\
2000	3492.3\\
4000	3024.5\\
6000	2932.4\\
8000	2928.4\\
10000	2949.4\\
12000	2925.2\\
14000	2932\\
16000	2960.2\\
18000	2926.1\\
20000	2927.2\\
22000	2953.7\\
24000	2938.2\\
28000	2933\\
30000	2970.1\\
32000	2959.9\\
34000	2973.8\\
36000	2936.9\\
38000	2966.8\\
40000	2963.8\\
42000	2937.2\\
44000	2945.4\\
46000	2944.1\\
48000	2961\\
50000	2936.9\\
52000	2928.1\\
54000	2942\\
56000	2977.7\\
58000	2944.7\\
60000	2945.3\\
62000	2928.1\\
66000	2948\\
70000	2934\\
72000	2964\\
74000	2980.9\\
76000	2978\\
78000	2927.5\\
80000	2929.9\\
82000	2969.6\\
84000	2997.5\\
86000	3003.1\\
88000	2957.1\\
90000	2981.8\\
92000	2981\\
94000	2945.6\\
96000	2976.8\\
98000	2920.5\\
1e+05	2959.6\\
};
\addplot [color=mycolor70, forget plot]
  table[row sep=crcr]{%
0	4395.3\\
2000	2906.4\\
4000	2211.9\\
6000	2149\\
8000	2166.1\\
12000	2183.7\\
14000	2125.3\\
16000	2152.3\\
18000	2146.5\\
20000	2165.1\\
22000	2152.7\\
24000	2184.4\\
26000	2144.2\\
28000	2147.3\\
30000	2156.8\\
32000	2140.9\\
34000	2157.4\\
36000	2133.8\\
38000	2175.1\\
40000	2126.4\\
42000	2171.9\\
44000	2151.6\\
46000	2143.8\\
48000	2164.4\\
50000	2143.5\\
52000	2137.8\\
58000	2180.6\\
60000	2157.1\\
62000	2141.6\\
68000	2136.7\\
70000	2167.1\\
72000	2121.7\\
74000	2118\\
76000	2101.6\\
78000	2176.2\\
80000	2169\\
82000	2192.3\\
84000	2170.7\\
86000	2172.5\\
88000	2192.5\\
90000	2168.3\\
92000	2185.2\\
94000	2193.8\\
96000	2160.6\\
98000	2209\\
1e+05	2211\\
};
\addplot [color=mycolor71, forget plot]
  table[row sep=crcr]{%
0	4395.3\\
2000	3151.1\\
4000	2629.2\\
6000	2539\\
8000	2518.5\\
10000	2514.8\\
12000	2574.7\\
14000	2553.7\\
18000	2534.8\\
20000	2540.2\\
22000	2516.5\\
26000	2540.4\\
28000	2505.7\\
30000	2505.1\\
32000	2524.1\\
36000	2523.7\\
38000	2507\\
40000	2534\\
42000	2521.4\\
44000	2522.2\\
48000	2556.5\\
50000	2537.3\\
52000	2544.2\\
54000	2520.4\\
56000	2538.6\\
58000	2501\\
60000	2537.3\\
62000	2540.6\\
64000	2533.2\\
66000	2553\\
68000	2528.3\\
70000	2496.1\\
72000	2536.4\\
74000	2541.3\\
78000	2520.7\\
80000	2530.8\\
82000	2526.7\\
88000	2560.1\\
90000	2550.6\\
92000	2530.7\\
94000	2538.2\\
96000	2514.9\\
1e+05	2520.5\\
};
\addplot [color=mycolor72, forget plot]
  table[row sep=crcr]{%
0	4395.3\\
2000	2998.4\\
4000	2315.4\\
6000	2272.3\\
8000	2269.4\\
10000	2237.7\\
12000	2233.9\\
14000	2297.4\\
18000	2301.7\\
20000	2215.6\\
22000	2246.8\\
24000	2246.2\\
26000	2238.4\\
28000	2275.4\\
30000	2249.5\\
32000	2271.1\\
36000	2242.1\\
38000	2280.2\\
40000	2311.9\\
42000	2233.1\\
44000	2226.6\\
46000	2208.8\\
48000	2181.4\\
52000	2232.2\\
54000	2226.5\\
56000	2207.8\\
58000	2229.7\\
60000	2232.2\\
62000	2243\\
64000	2204.2\\
66000	2236.4\\
68000	2230.5\\
70000	2245.9\\
72000	2195\\
74000	2254.6\\
76000	2219.9\\
78000	2254\\
80000	2208.4\\
82000	2250.2\\
86000	2217.4\\
88000	2261.6\\
90000	2273.7\\
92000	2230.4\\
94000	2257.4\\
98000	2233.2\\
1e+05	2258.3\\
};
\addplot [color=mycolor73, forget plot]
  table[row sep=crcr]{%
0	4395.3\\
2000	3211.4\\
4000	2726.8\\
6000	2646.7\\
8000	2635.2\\
10000	2658.4\\
14000	2654.5\\
16000	2672\\
18000	2672.3\\
20000	2650.2\\
22000	2644.3\\
24000	2658.5\\
26000	2640.2\\
30000	2661\\
34000	2655.4\\
36000	2677.8\\
38000	2647.1\\
40000	2667.9\\
42000	2640.9\\
44000	2624.9\\
46000	2629.3\\
48000	2649.6\\
50000	2678.8\\
52000	2685.8\\
54000	2682\\
56000	2629.6\\
58000	2650.5\\
60000	2679.2\\
62000	2683.1\\
66000	2622\\
68000	2655.8\\
74000	2652.7\\
76000	2640.9\\
78000	2621.4\\
80000	2629.7\\
82000	2654.9\\
84000	2664.2\\
88000	2671.2\\
90000	2634.7\\
92000	2618.8\\
94000	2637\\
96000	2633.8\\
98000	2670.6\\
1e+05	2672.8\\
};
\addplot [color=mycolor74, forget plot]
  table[row sep=crcr]{%
0	4395.3\\
2000	2479.4\\
4000	1681.6\\
6000	1612.7\\
8000	1595.3\\
12000	1596.7\\
14000	1585.4\\
18000	1597.3\\
20000	1580.9\\
22000	1555.1\\
24000	1589.5\\
26000	1597.7\\
28000	1599.7\\
30000	1572.9\\
32000	1567.4\\
34000	1597.5\\
38000	1567.9\\
40000	1589.7\\
42000	1581.3\\
44000	1589.3\\
46000	1573.6\\
48000	1582.1\\
50000	1556.7\\
54000	1599.2\\
56000	1586.6\\
58000	1590.5\\
62000	1589.1\\
64000	1578\\
66000	1586.5\\
68000	1584.6\\
70000	1587.8\\
72000	1572.4\\
74000	1590.8\\
76000	1574.6\\
78000	1563.1\\
80000	1556.8\\
84000	1587.2\\
86000	1582.3\\
90000	1593\\
92000	1573.2\\
94000	1593.5\\
98000	1574.5\\
1e+05	1580.9\\
};
\addplot [color=mycolor75, forget plot]
  table[row sep=crcr]{%
0	4395.3\\
2000	3039.6\\
4000	2346.4\\
6000	2286.9\\
8000	2281.8\\
10000	2284.7\\
12000	2230.6\\
14000	2251.9\\
18000	2231.8\\
20000	2245.7\\
22000	2271\\
24000	2250.4\\
26000	2251.1\\
28000	2234.1\\
30000	2236.9\\
32000	2277.1\\
34000	2279.4\\
36000	2261.6\\
38000	2250.6\\
42000	2277.7\\
44000	2248.3\\
46000	2231.4\\
48000	2258.9\\
52000	2236.1\\
54000	2256.6\\
56000	2253.1\\
58000	2264\\
60000	2246.4\\
62000	2246.2\\
64000	2264.9\\
66000	2264.5\\
68000	2233.2\\
72000	2257.5\\
74000	2289.6\\
78000	2245.1\\
84000	2263.9\\
86000	2230.6\\
88000	2244.7\\
90000	2241.7\\
92000	2221.9\\
94000	2250.7\\
96000	2273.5\\
1e+05	2261.8\\
};
\addplot [color=mycolor76, forget plot]
  table[row sep=crcr]{%
0	4395.3\\
2000	2934.7\\
4000	2187.2\\
6000	2087.8\\
8000	2058.3\\
10000	2072.5\\
12000	2074\\
14000	2061\\
16000	2087.5\\
18000	2084.9\\
20000	2053.1\\
22000	2081.3\\
24000	2042.4\\
26000	2072.1\\
28000	2055.1\\
30000	2069.1\\
32000	2068.2\\
34000	2048.9\\
36000	2050.5\\
38000	2061.1\\
40000	2055.2\\
42000	2070.3\\
46000	2038.7\\
48000	2041.5\\
50000	2063.2\\
52000	2066.3\\
54000	2087\\
56000	2073.8\\
58000	2042.4\\
62000	2055.5\\
64000	2088.5\\
66000	2091.7\\
68000	2075.8\\
70000	2047\\
74000	2064.9\\
78000	2054.3\\
82000	2084.7\\
84000	2057.4\\
90000	2049.3\\
92000	2055.5\\
94000	2048\\
96000	2062\\
1e+05	2042.1\\
};
\addplot [color=mycolor77, forget plot]
  table[row sep=crcr]{%
0	4395.3\\
2000	2572.7\\
4000	1731.4\\
6000	1636.4\\
8000	1609\\
10000	1600.4\\
12000	1616.7\\
14000	1625.5\\
16000	1596.3\\
18000	1599.6\\
20000	1625.6\\
22000	1623.8\\
24000	1606.6\\
26000	1615.1\\
28000	1600\\
30000	1626.7\\
32000	1625.7\\
34000	1603.9\\
36000	1606.7\\
38000	1633.3\\
42000	1636.3\\
44000	1618.5\\
46000	1612.5\\
48000	1584.7\\
50000	1623.7\\
52000	1637.7\\
54000	1596.4\\
56000	1599.3\\
58000	1616.6\\
60000	1591.3\\
62000	1579.2\\
64000	1609\\
66000	1624\\
68000	1618.1\\
70000	1598.9\\
72000	1602.2\\
74000	1621.4\\
76000	1609.9\\
78000	1618.8\\
80000	1621\\
82000	1606.8\\
84000	1630.4\\
86000	1632\\
88000	1603.3\\
90000	1603.4\\
92000	1598.1\\
94000	1586.2\\
96000	1608.8\\
98000	1614.2\\
1e+05	1624.2\\
};
\addplot [color=mycolor78, forget plot]
  table[row sep=crcr]{%
0	4395.3\\
2000	3133.7\\
4000	2464\\
6000	2376.2\\
8000	2319.7\\
12000	2309.2\\
14000	2316.1\\
18000	2368.3\\
20000	2338.2\\
24000	2322\\
28000	2361.7\\
30000	2330.1\\
32000	2361.5\\
34000	2295.7\\
36000	2324.6\\
38000	2306.1\\
40000	2374\\
42000	2356.9\\
44000	2350.6\\
46000	2364.7\\
48000	2327.9\\
50000	2312\\
56000	2361.8\\
58000	2320.6\\
60000	2331\\
62000	2278.1\\
64000	2311.9\\
66000	2331.1\\
68000	2374.3\\
70000	2357.5\\
74000	2312.5\\
76000	2320.8\\
78000	2337.7\\
80000	2287.1\\
82000	2317.7\\
84000	2312.7\\
86000	2321.7\\
88000	2370.9\\
90000	2368.8\\
92000	2385.9\\
94000	2349.8\\
96000	2324.5\\
98000	2359.6\\
1e+05	2368.2\\
};
\addplot [color=mycolor79, forget plot]
  table[row sep=crcr]{%
0	4395.3\\
2000	3030.3\\
4000	2327.5\\
6000	2253.9\\
8000	2192.9\\
16000	2214.9\\
18000	2173.2\\
20000	2205\\
22000	2174\\
24000	2197.9\\
26000	2204.7\\
28000	2170.5\\
30000	2174.8\\
32000	2123.8\\
34000	2189.4\\
40000	2203.6\\
42000	2172.7\\
44000	2172.3\\
46000	2221.3\\
50000	2188.7\\
52000	2219.7\\
54000	2187.3\\
56000	2208.8\\
58000	2197.2\\
60000	2174\\
62000	2201.2\\
64000	2195.7\\
66000	2175.7\\
68000	2188.2\\
70000	2177.6\\
78000	2194.5\\
80000	2205.9\\
82000	2193.1\\
86000	2204.1\\
88000	2177.7\\
90000	2171.8\\
92000	2245.8\\
94000	2191.4\\
96000	2181.3\\
98000	2188.9\\
1e+05	2232.8\\
};
\addplot [color=mycolor79, forget plot]
  table[row sep=crcr]{%
0	4395.3\\
2000	2967.2\\
4000	2191.8\\
6000	2084.1\\
8000	2092.5\\
10000	2054.1\\
12000	2068\\
14000	2095.2\\
16000	2113.6\\
20000	2043.5\\
22000	2090.1\\
24000	2061.8\\
28000	2060.6\\
30000	2031.8\\
32000	2021.6\\
34000	2056.4\\
38000	2063.9\\
40000	2081.5\\
42000	2087.8\\
44000	2072.7\\
46000	2048.5\\
48000	2077.6\\
50000	2044.5\\
52000	2032.4\\
54000	2069.1\\
56000	2063.9\\
60000	2039.2\\
62000	2072.1\\
64000	2076.5\\
66000	2063\\
68000	2029.2\\
70000	2076.4\\
72000	2054.1\\
74000	2081.3\\
76000	2072.6\\
78000	2082.9\\
80000	2061.8\\
84000	2042.8\\
86000	2065.8\\
88000	2064.8\\
90000	2123.1\\
92000	2073.3\\
94000	2091.9\\
96000	2089\\
98000	2057.7\\
1e+05	2039.6\\
};
\addplot [color=mycolor80, forget plot]
  table[row sep=crcr]{%
0	4395.3\\
2000	2656.9\\
4000	1741.8\\
6000	1606.4\\
8000	1584.2\\
10000	1603.4\\
12000	1605.4\\
16000	1585.5\\
18000	1589.7\\
20000	1610.4\\
22000	1598.3\\
24000	1591.4\\
26000	1609.4\\
28000	1588.9\\
30000	1579.5\\
32000	1593.4\\
36000	1582.6\\
38000	1594.2\\
42000	1588.7\\
44000	1607.7\\
46000	1605.9\\
48000	1614.1\\
50000	1592.7\\
52000	1594.4\\
54000	1607.7\\
56000	1578.8\\
58000	1574.6\\
60000	1627.4\\
62000	1603.7\\
64000	1598.4\\
66000	1607.5\\
68000	1595.7\\
70000	1577.9\\
72000	1590.5\\
74000	1578\\
76000	1596.7\\
78000	1587\\
80000	1604.2\\
82000	1604.7\\
84000	1593.5\\
86000	1597.5\\
88000	1616.6\\
90000	1596.8\\
92000	1603.8\\
94000	1590.7\\
96000	1586.2\\
98000	1595.7\\
1e+05	1595.2\\
};
\addplot [color=mycolor81, forget plot]
  table[row sep=crcr]{%
0	4395.3\\
2000	2600.6\\
4000	1649.4\\
6000	1525.4\\
8000	1511.1\\
14000	1500.9\\
16000	1476.5\\
18000	1522.2\\
20000	1512.7\\
22000	1488.9\\
24000	1485.1\\
26000	1487.8\\
28000	1497.1\\
30000	1483.9\\
32000	1497.5\\
34000	1493.8\\
36000	1484.4\\
40000	1504.3\\
42000	1499.4\\
44000	1481.3\\
46000	1471.2\\
48000	1478.5\\
50000	1470.6\\
54000	1503.9\\
56000	1501.3\\
58000	1482\\
60000	1499.4\\
62000	1493.4\\
64000	1512.8\\
66000	1511.8\\
68000	1499.9\\
70000	1493.3\\
72000	1481.4\\
74000	1481.9\\
76000	1498.5\\
78000	1491.6\\
80000	1489.3\\
82000	1481.6\\
86000	1489.4\\
88000	1484.2\\
90000	1493.9\\
92000	1492.2\\
94000	1498.5\\
1e+05	1500.5\\
};
\addplot [color=mycolor82, forget plot]
  table[row sep=crcr]{%
0	4395.3\\
2000	2755.9\\
4000	1896.3\\
6000	1732.7\\
8000	1737.4\\
10000	1763.7\\
12000	1741.5\\
14000	1740.5\\
18000	1724.4\\
20000	1731\\
22000	1748.8\\
24000	1745.2\\
26000	1751.4\\
28000	1765.7\\
30000	1730.1\\
32000	1740.6\\
34000	1760.8\\
36000	1769.7\\
38000	1753.7\\
40000	1705.3\\
42000	1708.1\\
44000	1731.9\\
48000	1750\\
50000	1752.9\\
54000	1701.6\\
56000	1728.6\\
62000	1723.6\\
64000	1734.7\\
66000	1739.8\\
68000	1732.8\\
70000	1715.5\\
72000	1742.4\\
74000	1732.4\\
76000	1743\\
78000	1727.2\\
80000	1723.7\\
86000	1736.2\\
88000	1758.3\\
90000	1757.4\\
92000	1743.7\\
94000	1745.5\\
96000	1726.3\\
1e+05	1721.7\\
};
\addplot [color=mycolor83, forget plot]
  table[row sep=crcr]{%
0	4395.3\\
2000	2866.3\\
4000	1954.5\\
6000	1816.8\\
8000	1770.6\\
10000	1780.9\\
12000	1755.7\\
18000	1790.6\\
20000	1750.9\\
22000	1768.6\\
24000	1780.8\\
26000	1765.7\\
28000	1786\\
30000	1776.9\\
32000	1776.8\\
34000	1791.2\\
36000	1742.1\\
38000	1745.3\\
40000	1766\\
42000	1761.3\\
46000	1762.1\\
48000	1788.1\\
50000	1758.8\\
52000	1782.5\\
54000	1772\\
56000	1752.7\\
58000	1781.3\\
60000	1769.8\\
64000	1778.5\\
66000	1759.7\\
68000	1758.8\\
70000	1781.9\\
72000	1766.9\\
74000	1765\\
76000	1750.3\\
80000	1798.2\\
82000	1784.2\\
88000	1761.9\\
90000	1764.3\\
92000	1777.4\\
94000	1769.7\\
96000	1782.3\\
98000	1762.2\\
1e+05	1752.9\\
};
\addplot [color=mycolor83, forget plot]
  table[row sep=crcr]{%
0	4395.3\\
2000	2500.7\\
4000	1461.5\\
6000	1344.4\\
8000	1294.9\\
10000	1301.4\\
12000	1296.8\\
14000	1281.1\\
16000	1291.5\\
18000	1289.1\\
20000	1299.5\\
22000	1305.7\\
24000	1298.3\\
26000	1311.1\\
28000	1296.2\\
30000	1299.9\\
32000	1298.6\\
34000	1305.4\\
36000	1303.6\\
38000	1296.9\\
40000	1300.8\\
42000	1286.1\\
44000	1299.8\\
46000	1290\\
48000	1294.2\\
50000	1347.5\\
52000	1319.3\\
54000	1317.5\\
56000	1300\\
58000	1301\\
62000	1291.8\\
64000	1278.4\\
68000	1318.7\\
72000	1321\\
74000	1305.1\\
76000	1307.8\\
78000	1297.2\\
80000	1302.8\\
82000	1292.6\\
86000	1293.5\\
88000	1282.4\\
90000	1297.2\\
92000	1316.8\\
94000	1310.4\\
96000	1298.1\\
1e+05	1286.9\\
};
\addplot [color=mycolor84, forget plot]
  table[row sep=crcr]{%
0	4395.3\\
2000	2563\\
4000	1516.3\\
6000	1384.9\\
8000	1322.5\\
10000	1329\\
12000	1313.6\\
14000	1303.7\\
16000	1316.3\\
20000	1334.6\\
22000	1326.6\\
24000	1303.3\\
26000	1324.7\\
28000	1323.1\\
30000	1311.7\\
32000	1332.1\\
34000	1320.4\\
36000	1293.3\\
38000	1295.3\\
42000	1324.7\\
44000	1316.9\\
46000	1314.7\\
48000	1341.7\\
50000	1311.2\\
52000	1333.3\\
54000	1335.3\\
56000	1330.2\\
58000	1347.7\\
60000	1347.2\\
62000	1315.5\\
64000	1319.3\\
66000	1315.9\\
68000	1325.8\\
70000	1294.8\\
72000	1299.2\\
76000	1301.5\\
78000	1313.1\\
80000	1317.5\\
82000	1317.3\\
86000	1305.9\\
90000	1328.2\\
92000	1317.3\\
94000	1292.5\\
96000	1304.1\\
98000	1311.4\\
1e+05	1308.4\\
};
\addplot [color=mycolor85, forget plot]
  table[row sep=crcr]{%
0	4395.3\\
2000	2616\\
4000	1645.6\\
6000	1466.6\\
8000	1485.5\\
10000	1442.5\\
12000	1432.1\\
14000	1443.2\\
16000	1419.4\\
18000	1434.2\\
20000	1417.7\\
24000	1417.6\\
26000	1463.9\\
28000	1433.4\\
30000	1440.4\\
32000	1430.7\\
34000	1429.1\\
36000	1422.2\\
38000	1445.1\\
40000	1437.6\\
44000	1431.7\\
46000	1451.5\\
48000	1465.6\\
50000	1395.6\\
52000	1425.4\\
54000	1469.3\\
58000	1456.1\\
60000	1420.3\\
62000	1463.7\\
64000	1436.7\\
66000	1429.1\\
70000	1430.7\\
72000	1433.2\\
74000	1419.8\\
76000	1427.9\\
78000	1415.3\\
80000	1424.1\\
82000	1420.2\\
84000	1455.3\\
86000	1460.5\\
88000	1457\\
92000	1414.7\\
94000	1401.5\\
96000	1424.2\\
98000	1426\\
1e+05	1417.8\\
};
\addplot [color=mycolor86, forget plot]
  table[row sep=crcr]{%
0	4395.3\\
2000	2601\\
4000	1458.2\\
6000	1263.2\\
10000	1188\\
12000	1210.5\\
14000	1192.7\\
16000	1197.8\\
18000	1215.3\\
20000	1248.6\\
22000	1240.9\\
24000	1213.5\\
26000	1213\\
28000	1219.2\\
32000	1205.7\\
34000	1221.5\\
38000	1227.4\\
40000	1223.7\\
42000	1214.6\\
44000	1217.6\\
46000	1216.5\\
48000	1204.2\\
52000	1231\\
54000	1218.9\\
56000	1216.6\\
58000	1224.5\\
60000	1244.9\\
62000	1216.5\\
64000	1211.4\\
66000	1210.6\\
68000	1219.8\\
70000	1223.8\\
72000	1201.8\\
74000	1205.7\\
76000	1204.3\\
80000	1232.7\\
82000	1203.4\\
86000	1203.1\\
88000	1223.1\\
90000	1211.5\\
92000	1234.8\\
96000	1219\\
98000	1225.7\\
1e+05	1206.8\\
};
\addplot [color=mycolor87, forget plot]
  table[row sep=crcr]{%
0	4395.3\\
2000	2415.7\\
4000	1234.2\\
6000	1053.9\\
8000	1051.3\\
10000	1037.9\\
14000	1029.7\\
16000	1036.1\\
18000	1070.8\\
20000	1042\\
22000	1024.7\\
24000	1016.4\\
26000	1043.7\\
28000	1065.5\\
30000	1042.7\\
32000	1047.6\\
34000	1022.8\\
36000	1046\\
38000	1053.3\\
40000	1040.8\\
42000	1044.4\\
44000	1020.9\\
46000	1029.5\\
48000	1058.4\\
52000	1039.3\\
54000	991.93\\
56000	1012.3\\
58000	1007.1\\
60000	1024.7\\
62000	1021.5\\
64000	1034.6\\
66000	1040.4\\
68000	1057.8\\
70000	1041.7\\
72000	1021.6\\
74000	1042.6\\
76000	1020.9\\
78000	1022.2\\
80000	1049.4\\
82000	1036.7\\
86000	1041.5\\
88000	1053.7\\
90000	1038.3\\
92000	1036\\
94000	1052.1\\
96000	1040.6\\
98000	1055.7\\
1e+05	1033.2\\
};
\addplot [color=mycolor88, forget plot]
  table[row sep=crcr]{%
0	4395.3\\
2000	2642.7\\
4000	1474.1\\
6000	1246.5\\
8000	1202.4\\
10000	1222.4\\
12000	1224.6\\
14000	1211.2\\
16000	1216.1\\
18000	1210.1\\
20000	1188.3\\
22000	1203.1\\
24000	1201.6\\
26000	1229.8\\
28000	1227.3\\
30000	1213.5\\
32000	1205.5\\
34000	1203.6\\
36000	1209.4\\
38000	1207.2\\
40000	1221.6\\
44000	1228.6\\
46000	1227.9\\
48000	1210.1\\
50000	1212.3\\
52000	1195.9\\
54000	1202.6\\
56000	1199\\
62000	1248.6\\
64000	1240\\
66000	1209.6\\
68000	1203.5\\
70000	1222.7\\
72000	1223.9\\
74000	1209\\
76000	1215.9\\
80000	1202.4\\
82000	1228.9\\
84000	1199.5\\
90000	1177.9\\
94000	1241.2\\
96000	1219.4\\
1e+05	1193.2\\
};
\addplot [color=mycolor89, forget plot]
  table[row sep=crcr]{%
0	4395.3\\
2000	2399.9\\
4000	1123.8\\
6000	919.33\\
8000	895.92\\
10000	890.86\\
12000	890.33\\
14000	902.53\\
16000	888.89\\
20000	890.31\\
22000	902.16\\
24000	882.92\\
26000	871.29\\
28000	865.87\\
30000	891.02\\
36000	856.95\\
40000	884\\
42000	881.81\\
44000	890.7\\
46000	866.48\\
48000	867.42\\
50000	864.61\\
52000	887.92\\
54000	886.87\\
56000	903.71\\
58000	896.57\\
60000	867.89\\
62000	864.07\\
66000	884.99\\
68000	884.44\\
70000	879.12\\
72000	866.68\\
74000	888.37\\
76000	873.62\\
82000	876.7\\
88000	881.44\\
90000	878.29\\
92000	895.42\\
94000	888.11\\
98000	897.14\\
1e+05	892.86\\
};
\addplot [color=mycolor90, forget plot]
  table[row sep=crcr]{%
0	4395.3\\
2000	2377.1\\
4000	1113.3\\
6000	897.67\\
8000	858.07\\
10000	842.42\\
14000	845.65\\
16000	838.07\\
18000	847.21\\
20000	853.52\\
22000	839.89\\
24000	853.42\\
26000	834.31\\
28000	843.83\\
30000	833.69\\
32000	845.05\\
34000	850.35\\
36000	846.8\\
38000	848.63\\
40000	853.9\\
44000	844.92\\
46000	839.48\\
48000	850.2\\
50000	856.35\\
52000	858.41\\
54000	844.98\\
56000	848.23\\
60000	860.47\\
62000	847.01\\
64000	845.82\\
66000	832.28\\
68000	855.79\\
70000	859.01\\
72000	842.79\\
74000	862.88\\
76000	855.95\\
78000	860.39\\
80000	845.26\\
84000	861.96\\
86000	865.54\\
88000	883.05\\
90000	875.47\\
92000	860.21\\
94000	851.76\\
96000	870.01\\
98000	839.86\\
1e+05	837.39\\
};
\addplot [color=mycolor91, forget plot]
  table[row sep=crcr]{%
0	4395.3\\
2000	2365.3\\
4000	1070.3\\
6000	807.88\\
8000	780.76\\
10000	744.11\\
12000	732.29\\
14000	755.96\\
16000	756.53\\
18000	751.13\\
20000	741.77\\
22000	740.6\\
24000	719.6\\
26000	725.21\\
28000	727.15\\
30000	743.54\\
32000	729.95\\
34000	742.58\\
36000	737.51\\
38000	728.47\\
40000	752.86\\
42000	728.61\\
44000	750.35\\
46000	753.27\\
48000	738.73\\
50000	770.77\\
52000	756.37\\
54000	763.02\\
58000	745.64\\
60000	720.39\\
62000	732.18\\
64000	725.87\\
66000	742.13\\
68000	741.82\\
70000	754.72\\
72000	748\\
74000	729.8\\
76000	747.29\\
78000	752.15\\
80000	750.74\\
82000	764.93\\
84000	739.89\\
88000	759.09\\
90000	734.61\\
92000	719.28\\
94000	744.92\\
96000	738.61\\
98000	742.6\\
1e+05	754.48\\
};
\addplot [color=mycolor92, forget plot]
  table[row sep=crcr]{%
0	4395.3\\
2000	2550.5\\
4000	1277.9\\
6000	992.29\\
8000	924.61\\
10000	897.96\\
12000	922.37\\
14000	913.57\\
16000	938.51\\
18000	924.95\\
20000	941.16\\
22000	925.1\\
30000	925.77\\
32000	935.49\\
36000	913.6\\
40000	933.69\\
44000	903.04\\
46000	906.3\\
48000	939.47\\
50000	934.59\\
52000	935.27\\
54000	922.1\\
56000	916.81\\
58000	901.28\\
60000	888.77\\
62000	925.24\\
66000	905.78\\
68000	918.32\\
70000	928.36\\
72000	931.79\\
74000	909.27\\
76000	914.24\\
78000	915.68\\
80000	924.11\\
82000	921.28\\
84000	924.68\\
86000	916.25\\
88000	901.1\\
92000	915.2\\
94000	907.49\\
98000	924.43\\
1e+05	928.18\\
};
\addplot [color=mycolor93, forget plot]
  table[row sep=crcr]{%
0	4395.3\\
2000	2551.8\\
4000	1143.9\\
6000	854.17\\
8000	792.62\\
10000	784.27\\
12000	772.29\\
14000	764.84\\
16000	761.17\\
18000	764.16\\
20000	747.37\\
22000	755.37\\
24000	777.74\\
26000	814.21\\
28000	814.95\\
30000	777.79\\
32000	770.66\\
34000	792.66\\
36000	803.83\\
38000	772.69\\
40000	770.72\\
42000	792.7\\
44000	778.88\\
46000	779.14\\
48000	745.34\\
54000	769.54\\
56000	769.26\\
60000	746.78\\
62000	745.96\\
66000	799.62\\
68000	780.01\\
72000	765.39\\
74000	774.23\\
76000	791.96\\
78000	776.27\\
80000	801.35\\
82000	824.65\\
84000	786.08\\
86000	782.37\\
88000	788.16\\
90000	784.76\\
92000	750.45\\
94000	746.71\\
96000	771.54\\
98000	780.19\\
1e+05	807.5\\
};
\addplot [color=mycolor94, forget plot]
  table[row sep=crcr]{%
0	4395.3\\
2000	2579.7\\
4000	1188.1\\
6000	871.85\\
8000	781.95\\
12000	726.32\\
20000	741.87\\
22000	754.58\\
26000	750.58\\
28000	738.77\\
30000	743.85\\
32000	739.48\\
34000	725.15\\
36000	724.65\\
38000	719.2\\
40000	720.28\\
42000	717.74\\
44000	718.96\\
46000	738.2\\
48000	741.47\\
50000	736.76\\
52000	742.15\\
54000	740.81\\
56000	715.66\\
58000	729.64\\
60000	728.35\\
62000	755.08\\
64000	745.37\\
66000	733.28\\
68000	732.84\\
70000	727.09\\
72000	736.7\\
74000	734.74\\
76000	742.42\\
78000	734.04\\
80000	730.84\\
82000	738.42\\
84000	729.89\\
86000	750.13\\
88000	738.53\\
90000	730.66\\
92000	756.19\\
94000	754.97\\
96000	736.08\\
98000	740.53\\
1e+05	748.58\\
};
\addplot [color=mycolor95, forget plot]
  table[row sep=crcr]{%
0	4395.3\\
2000	2343.2\\
4000	835.61\\
6000	549.55\\
8000	442.83\\
10000	439\\
12000	437.36\\
14000	441.95\\
16000	440.6\\
18000	440.53\\
20000	438.06\\
24000	449.24\\
28000	445.2\\
30000	448.45\\
32000	443.4\\
34000	440.07\\
36000	446.97\\
38000	440.21\\
40000	443.6\\
42000	451.89\\
44000	446.76\\
46000	451.18\\
48000	448.46\\
50000	438.13\\
52000	444.28\\
54000	454.78\\
56000	441.49\\
58000	445.54\\
60000	440.8\\
62000	432.12\\
64000	435.67\\
66000	447.06\\
68000	445.99\\
72000	439.56\\
74000	432.79\\
76000	442.63\\
78000	437.23\\
80000	433.87\\
82000	440.63\\
84000	431.02\\
88000	442.92\\
90000	447.19\\
92000	449.45\\
94000	429.27\\
96000	422.86\\
1e+05	448.72\\
};
\addplot [color=mycolor96, forget plot]
  table[row sep=crcr]{%
0	4395.3\\
2000	2486.9\\
4000	955.93\\
6000	622.55\\
8000	504.73\\
10000	468.37\\
12000	467.99\\
14000	478.39\\
16000	490.41\\
18000	489.45\\
20000	479.27\\
22000	475.21\\
24000	465.67\\
26000	471.3\\
28000	479.98\\
30000	461.43\\
32000	485.82\\
34000	497.88\\
36000	502.48\\
38000	482.89\\
40000	472.39\\
42000	478.81\\
44000	457.8\\
48000	491.8\\
50000	491.23\\
52000	484.31\\
54000	446.61\\
56000	462.36\\
58000	489.21\\
60000	480.93\\
62000	456.26\\
64000	456.9\\
66000	476.96\\
68000	488.95\\
70000	455.53\\
72000	469.45\\
74000	476.28\\
76000	473.37\\
78000	482.45\\
80000	472.52\\
82000	446.71\\
84000	444.41\\
86000	481.15\\
88000	476.91\\
90000	463.83\\
92000	461.47\\
94000	450.11\\
96000	445.84\\
98000	454.5\\
1e+05	446.42\\
};
\addplot [color=mycolor97, forget plot]
  table[row sep=crcr]{%
0	4395.3\\
2000	2393.8\\
4000	853.12\\
6000	478.62\\
8000	361.14\\
10000	352.44\\
12000	357.53\\
14000	332.28\\
16000	324.32\\
18000	304.12\\
20000	326.82\\
22000	332.69\\
24000	333.92\\
26000	324.31\\
28000	330.81\\
30000	322.16\\
32000	326.38\\
34000	326.42\\
36000	322.79\\
38000	324.55\\
42000	321.81\\
44000	335.07\\
46000	340.41\\
48000	330.23\\
50000	326.28\\
52000	311.96\\
54000	319.42\\
56000	332.32\\
58000	318.92\\
60000	314.57\\
62000	323.3\\
64000	339.78\\
66000	335.14\\
68000	328.87\\
70000	323.71\\
72000	314.11\\
74000	326.5\\
76000	321.22\\
78000	327.18\\
80000	330.56\\
82000	327.9\\
84000	334.86\\
86000	350.88\\
88000	327.43\\
90000	336.55\\
92000	319\\
94000	321.67\\
96000	343.23\\
98000	335.88\\
1e+05	331.4\\
};
\addplot [color=mycolor98]
  table[row sep=crcr]{%
0	4395.3\\
2000	2672.6\\
4000	1116.9\\
6000	661.14\\
8000	525.46\\
10000	473.85\\
12000	452.69\\
14000	459.12\\
16000	444.79\\
18000	463.29\\
20000	451.08\\
22000	464.81\\
24000	449.22\\
26000	453.05\\
28000	466.23\\
30000	460.47\\
32000	448.05\\
34000	445.59\\
36000	456.54\\
38000	453.11\\
40000	456.32\\
42000	443.71\\
44000	438.37\\
46000	448.83\\
48000	451.65\\
50000	440.47\\
52000	446.38\\
54000	439.39\\
56000	440.76\\
58000	437.83\\
60000	443.22\\
62000	460.31\\
66000	456.41\\
68000	458.12\\
70000	448.12\\
72000	451.37\\
74000	460.74\\
76000	452.19\\
78000	457.68\\
80000	442.02\\
82000	446.32\\
84000	445.01\\
86000	458.34\\
88000	455.19\\
90000	458.73\\
92000	459.42\\
94000	466.39\\
96000	455.2\\
98000	432.46\\
1e+05	448.61\\
};
\addlegendentry{\ClasskIVLegend}

      % \input{../../../../.tex/20/KS/SizeEvolutionsfig4.tex}
      \legend{}
    \end{groupplot}

    \begin{groupplot}[
      group style={
        group name=Row3,
        plotGroup,
        group size=2 by 1,
      },
      plotRow,
      row5Keys,
      xTickScaleLabelEmpty,
      plotLegendSE,
      %
      xmin=0,xmax=1e+05,
      ymin=100,ymax=1.539e+05,
    ]
      \nextgroupplot[%
        at={($(Row2 c1r1.south west)-(0,\yPlotSep em)-(0,\plotHeight cm)$)},
        % xmin=0,xmax=1e+05,
        % ymin=100,ymax=1.5248e+05,
      ] % This file was created by matlab2tikz.
%
\definecolor{mycolor1}{rgb}{0.00146,0.00047,0.01387}%
\definecolor{mycolor2}{rgb}{0.01166,0.00942,0.06346}%
\definecolor{mycolor3}{rgb}{0.02943,0.02150,0.11462}%
\definecolor{mycolor4}{rgb}{0.06134,0.03659,0.17764}%
\definecolor{mycolor5}{rgb}{0.08741,0.04456,0.22481}%
\definecolor{mycolor6}{rgb}{0.12291,0.04754,0.28162}%
\definecolor{mycolor7}{rgb}{0.14907,0.04547,0.31709}%
\definecolor{mycolor8}{rgb}{0.18343,0.04033,0.35497}%
\definecolor{mycolor9}{rgb}{0.21795,0.03662,0.38352}%
\definecolor{mycolor10}{rgb}{0.24497,0.03705,0.40001}%
\definecolor{mycolor11}{rgb}{0.27135,0.04092,0.41198}%
\definecolor{mycolor12}{rgb}{0.29076,0.04564,0.41864}%
\definecolor{mycolor13}{rgb}{0.31628,0.05349,0.42512}%
\definecolor{mycolor14}{rgb}{0.33522,0.06006,0.42852}%
\definecolor{mycolor15}{rgb}{0.36028,0.06925,0.43150}%
\definecolor{mycolor16}{rgb}{0.37900,0.07625,0.43272}%
\definecolor{mycolor17}{rgb}{0.39767,0.08326,0.43318}%
\definecolor{mycolor18}{rgb}{0.41633,0.09020,0.43294}%
\definecolor{mycolor19}{rgb}{0.42877,0.09479,0.43241}%
\definecolor{mycolor20}{rgb}{0.44743,0.10160,0.43108}%
\definecolor{mycolor21}{rgb}{0.46610,0.10832,0.42912}%
\definecolor{mycolor22}{rgb}{0.47856,0.11276,0.42747}%
\definecolor{mycolor23}{rgb}{0.49726,0.11938,0.42449}%
\definecolor{mycolor24}{rgb}{0.50973,0.12377,0.42216}%
\definecolor{mycolor25}{rgb}{0.52221,0.12815,0.41955}%
\definecolor{mycolor26}{rgb}{0.53468,0.13253,0.41667}%
\definecolor{mycolor27}{rgb}{0.55339,0.13913,0.41183}%
\definecolor{mycolor28}{rgb}{0.56585,0.14357,0.40826}%
\definecolor{mycolor29}{rgb}{0.57830,0.14804,0.40441}%
\definecolor{mycolor30}{rgb}{0.59073,0.15256,0.40029}%
\definecolor{mycolor31}{rgb}{0.60314,0.15715,0.39589}%
\definecolor{mycolor32}{rgb}{0.61551,0.16182,0.39122}%
\definecolor{mycolor33}{rgb}{0.62169,0.16418,0.38878}%
\definecolor{mycolor34}{rgb}{0.63400,0.16899,0.38370}%
\definecolor{mycolor35}{rgb}{0.64626,0.17391,0.37836}%
\definecolor{mycolor36}{rgb}{0.65846,0.17896,0.37275}%
\definecolor{mycolor37}{rgb}{0.66454,0.18154,0.36985}%
\definecolor{mycolor38}{rgb}{0.67664,0.18681,0.36385}%
\definecolor{mycolor39}{rgb}{0.68865,0.19224,0.35760}%
\definecolor{mycolor40}{rgb}{0.69463,0.19502,0.35439}%
\definecolor{mycolor41}{rgb}{0.70650,0.20073,0.34778}%
\definecolor{mycolor42}{rgb}{0.71240,0.20366,0.34438}%
\definecolor{mycolor43}{rgb}{0.72410,0.20967,0.33742}%
\definecolor{mycolor44}{rgb}{0.72991,0.21276,0.33386}%
\definecolor{mycolor45}{rgb}{0.74142,0.21911,0.32658}%
\definecolor{mycolor46}{rgb}{0.74713,0.22238,0.32286}%
\definecolor{mycolor47}{rgb}{0.75279,0.22571,0.31909}%
\definecolor{mycolor48}{rgb}{0.76401,0.23255,0.31140}%
\definecolor{mycolor49}{rgb}{0.76956,0.23608,0.30749}%
\definecolor{mycolor50}{rgb}{0.77506,0.23967,0.30353}%
\definecolor{mycolor51}{rgb}{0.79129,0.25086,0.29139}%
\definecolor{mycolor52}{rgb}{0.79661,0.25473,0.28726}%
\definecolor{mycolor53}{rgb}{0.80187,0.25867,0.28310}%
\definecolor{mycolor54}{rgb}{0.81224,0.26679,0.27466}%
\definecolor{mycolor55}{rgb}{0.81734,0.27095,0.27039}%
\definecolor{mycolor56}{rgb}{0.82737,0.27952,0.26175}%
\definecolor{mycolor57}{rgb}{0.83230,0.28391,0.25738}%
\definecolor{mycolor58}{rgb}{0.83717,0.28839,0.25299}%
\definecolor{mycolor59}{rgb}{0.84671,0.29756,0.24411}%
\definecolor{mycolor60}{rgb}{0.85138,0.30226,0.23964}%
\definecolor{mycolor61}{rgb}{0.85599,0.30704,0.23513}%
\definecolor{mycolor62}{rgb}{0.86053,0.31189,0.23061}%
\definecolor{mycolor63}{rgb}{0.87374,0.32691,0.21689}%
\definecolor{mycolor64}{rgb}{0.87800,0.33206,0.21227}%
\definecolor{mycolor65}{rgb}{0.89034,0.34796,0.19829}%
\definecolor{mycolor66}{rgb}{0.89819,0.35891,0.18886}%
\definecolor{mycolor67}{rgb}{0.90200,0.36449,0.18412}%
\definecolor{mycolor68}{rgb}{0.90573,0.37014,0.17935}%
\definecolor{mycolor69}{rgb}{0.90939,0.37586,0.17456}%
\definecolor{mycolor70}{rgb}{0.91297,0.38164,0.16975}%
\definecolor{mycolor71}{rgb}{0.91646,0.38748,0.16492}%
\definecolor{mycolor72}{rgb}{0.91988,0.39339,0.16007}%
\definecolor{mycolor73}{rgb}{0.92322,0.39936,0.15519}%
\definecolor{mycolor74}{rgb}{0.92647,0.40539,0.15029}%
\definecolor{mycolor75}{rgb}{0.92964,0.41148,0.14537}%
\definecolor{mycolor76}{rgb}{0.93274,0.41763,0.14042}%
\definecolor{mycolor77}{rgb}{0.93575,0.42383,0.13544}%
\definecolor{mycolor78}{rgb}{0.93868,0.43009,0.13044}%
\definecolor{mycolor79}{rgb}{0.94152,0.43640,0.12541}%
\definecolor{mycolor80}{rgb}{0.94429,0.44277,0.12035}%
\definecolor{mycolor81}{rgb}{0.94956,0.45566,0.11016}%
\definecolor{mycolor82}{rgb}{0.95208,0.46218,0.10503}%
\definecolor{mycolor83}{rgb}{0.96129,0.48872,0.08429}%
\definecolor{mycolor84}{rgb}{0.96540,0.50225,0.07386}%
\definecolor{mycolor85}{rgb}{0.96732,0.50908,0.06866}%
\definecolor{mycolor86}{rgb}{0.97259,0.52980,0.05332}%
\definecolor{mycolor87}{rgb}{0.97568,0.54380,0.04362}%
\definecolor{mycolor88}{rgb}{0.98082,0.57221,0.02851}%
\definecolor{mycolor89}{rgb}{0.98189,0.57939,0.02625}%
\definecolor{mycolor90}{rgb}{0.98378,0.59385,0.02377}%
\definecolor{mycolor91}{rgb}{0.98532,0.60842,0.02420}%
\definecolor{mycolor92}{rgb}{0.98696,0.63048,0.03091}%
\definecolor{mycolor93}{rgb}{0.98793,0.66025,0.05175}%
\definecolor{mycolor94}{rgb}{0.98771,0.68281,0.07249}%
\definecolor{mycolor95}{rgb}{0.97750,0.78226,0.18592}%
\definecolor{mycolor96}{rgb}{0.96624,0.83619,0.26153}%
\definecolor{mycolor97}{rgb}{0.96252,0.85148,0.28555}%
\definecolor{mycolor98}{rgb}{0.98836,0.99836,0.64492}%
%
\addplot [color=mycolor1]
  table[row sep=crcr]{%
0	4395.3\\
2000	94531\\
4000	1.3812e+05\\
6000	1.4618e+05\\
8000	1.4644e+05\\
10000	1.518e+05\\
12000	1.5096e+05\\
14000	1.4711e+05\\
16000	1.486e+05\\
22000	1.4849e+05\\
24000	1.4743e+05\\
26000	1.5247e+05\\
28000	1.4977e+05\\
30000	1.5046e+05\\
32000	1.5041e+05\\
34000	1.5086e+05\\
36000	1.4748e+05\\
38000	1.4779e+05\\
42000	1.5019e+05\\
44000	1.5024e+05\\
48000	1.4928e+05\\
52000	1.5211e+05\\
54000	1.5087e+05\\
56000	1.4886e+05\\
58000	1.4852e+05\\
60000	1.5199e+05\\
62000	1.5248e+05\\
64000	1.5038e+05\\
66000	1.5023e+05\\
68000	1.4946e+05\\
70000	1.5143e+05\\
74000	1.4911e+05\\
76000	1.4704e+05\\
78000	1.4751e+05\\
80000	1.4988e+05\\
86000	1.4752e+05\\
88000	1.4856e+05\\
90000	1.5016e+05\\
92000	1.4887e+05\\
96000	1.4997e+05\\
98000	1.5196e+05\\
1e+05	1.5023e+05\\
};
\addlegendentry{\ClasskILegend}

\addplot [color=mycolor2, forget plot]
  table[row sep=crcr]{%
0	4395.3\\
2000	24749\\
4000	25274\\
6000	25731\\
8000	25914\\
10000	25398\\
12000	25490\\
14000	25957\\
16000	25833\\
18000	25441\\
20000	26182\\
26000	25182\\
30000	25732\\
32000	25641\\
34000	25018\\
36000	25548\\
38000	25684\\
40000	25621\\
42000	25749\\
44000	25556\\
46000	25646\\
50000	25503\\
52000	26027\\
54000	24895\\
56000	25976\\
58000	25879\\
60000	25349\\
62000	25505\\
64000	25920\\
66000	25613\\
68000	25893\\
72000	25818\\
74000	25624\\
76000	25708\\
80000	26152\\
82000	25776\\
86000	25598\\
88000	25387\\
90000	25987\\
92000	26370\\
94000	25515\\
96000	25741\\
98000	25505\\
1e+05	25974\\
};
\addplot [color=mycolor3, forget plot]
  table[row sep=crcr]{%
0	4395.3\\
2000	34244\\
4000	38146\\
6000	38613\\
8000	38358\\
10000	38762\\
12000	39421\\
14000	39933\\
16000	39944\\
18000	39549\\
20000	39778\\
24000	39446\\
26000	39017\\
28000	38744\\
30000	37950\\
32000	38942\\
34000	39014\\
36000	39448\\
38000	39232\\
42000	38574\\
44000	39458\\
46000	39557\\
48000	39397\\
50000	38527\\
52000	39192\\
54000	39004\\
56000	40154\\
58000	39139\\
60000	38609\\
62000	39503\\
64000	38568\\
66000	38964\\
68000	38821\\
70000	38281\\
72000	39510\\
74000	39566\\
76000	38725\\
78000	38850\\
80000	39375\\
82000	39557\\
84000	37957\\
86000	38615\\
88000	39857\\
90000	40296\\
92000	39002\\
96000	38537\\
1e+05	38462\\
};
\addplot [color=mycolor4, forget plot]
  table[row sep=crcr]{%
0	4395.3\\
2000	14833\\
4000	14334\\
6000	14174\\
8000	14434\\
10000	14447\\
12000	14184\\
14000	14242\\
16000	14402\\
18000	14159\\
20000	14309\\
22000	14404\\
24000	14125\\
26000	14179\\
28000	14126\\
30000	14209\\
32000	14374\\
34000	13965\\
36000	14177\\
38000	13880\\
40000	14237\\
42000	13943\\
44000	13807\\
46000	14442\\
48000	14138\\
52000	14188\\
54000	14320\\
56000	14030\\
58000	14149\\
60000	14039\\
62000	14020\\
64000	14257\\
66000	13957\\
68000	14225\\
70000	14233\\
72000	14091\\
74000	14046\\
76000	14088\\
78000	14592\\
80000	14570\\
82000	13916\\
84000	14200\\
86000	14137\\
88000	14025\\
90000	14212\\
92000	14449\\
94000	14529\\
96000	14063\\
98000	14065\\
1e+05	14114\\
};
\addplot [color=mycolor5, forget plot]
  table[row sep=crcr]{%
0	4395.3\\
2000	46056\\
4000	62952\\
6000	64504\\
8000	63313\\
10000	61933\\
12000	62389\\
14000	62191\\
16000	63089\\
18000	63796\\
20000	63335\\
22000	61418\\
24000	62529\\
26000	60443\\
28000	61917\\
30000	62368\\
32000	62477\\
34000	62150\\
36000	63688\\
38000	63944\\
40000	62237\\
42000	64235\\
44000	62510\\
46000	62207\\
48000	60784\\
50000	61103\\
52000	59948\\
54000	61339\\
56000	63325\\
60000	62524\\
62000	62565\\
64000	63379\\
66000	62065\\
68000	63041\\
70000	62395\\
72000	63483\\
74000	64011\\
76000	64111\\
78000	63881\\
80000	62437\\
82000	63916\\
84000	64036\\
86000	61751\\
88000	62922\\
92000	61379\\
94000	62029\\
96000	62926\\
1e+05	62250\\
};
\addplot [color=mycolor6, forget plot]
  table[row sep=crcr]{%
0	4395.3\\
2000	38102\\
4000	53152\\
6000	52998\\
8000	53900\\
10000	53481\\
12000	55058\\
14000	55439\\
16000	55079\\
18000	56693\\
20000	53741\\
22000	55019\\
26000	55196\\
28000	54363\\
30000	54780\\
32000	55945\\
34000	55767\\
36000	55979\\
38000	56491\\
40000	55815\\
42000	54410\\
44000	53273\\
46000	54675\\
48000	55195\\
50000	55300\\
52000	56354\\
54000	54835\\
56000	54826\\
58000	55415\\
62000	56142\\
64000	56112\\
66000	54009\\
68000	55184\\
70000	54963\\
72000	54970\\
74000	56949\\
76000	57163\\
78000	58004\\
80000	57167\\
82000	56541\\
84000	56442\\
86000	55383\\
88000	54007\\
90000	56854\\
92000	55447\\
94000	55491\\
96000	55939\\
98000	54465\\
1e+05	54757\\
};
\addplot [color=mycolor7, forget plot]
  table[row sep=crcr]{%
0	4395.3\\
2000	22170\\
4000	26878\\
6000	28427\\
8000	27404\\
10000	27052\\
12000	27603\\
14000	26733\\
16000	27795\\
18000	27727\\
20000	28318\\
22000	27572\\
24000	28339\\
26000	27254\\
28000	27671\\
30000	27667\\
32000	27037\\
34000	27656\\
36000	28154\\
38000	28028\\
40000	27554\\
42000	27898\\
44000	27577\\
46000	28416\\
48000	28234\\
50000	28932\\
52000	28026\\
54000	27749\\
56000	28528\\
58000	27959\\
60000	28444\\
62000	27496\\
64000	28132\\
66000	27929\\
68000	27576\\
70000	28360\\
72000	27890\\
74000	27564\\
76000	27849\\
78000	27437\\
80000	27739\\
82000	27354\\
84000	28473\\
86000	28106\\
88000	27887\\
92000	27236\\
94000	28275\\
96000	28887\\
98000	28871\\
1e+05	27891\\
};
\addplot [color=mycolor8, forget plot]
  table[row sep=crcr]{%
0	4395.3\\
2000	20042\\
4000	21604\\
6000	20441\\
8000	20388\\
10000	20716\\
12000	20451\\
14000	20291\\
16000	20474\\
18000	19994\\
20000	20395\\
22000	20537\\
24000	20309\\
26000	20645\\
28000	19691\\
30000	20784\\
32000	20712\\
36000	20236\\
38000	20209\\
40000	20588\\
42000	19722\\
44000	20389\\
46000	20682\\
48000	20266\\
50000	20238\\
52000	19863\\
54000	19430\\
56000	20399\\
58000	19993\\
60000	20292\\
62000	20381\\
64000	20828\\
66000	20476\\
68000	19900\\
70000	19827\\
74000	20040\\
76000	19974\\
78000	19423\\
82000	20115\\
84000	20411\\
86000	20510\\
88000	20422\\
90000	20107\\
92000	20226\\
94000	20452\\
96000	20442\\
98000	20193\\
1e+05	20397\\
};
\addplot [color=mycolor9, forget plot]
  table[row sep=crcr]{%
0	4395.3\\
2000	13068\\
4000	13855\\
6000	13923\\
8000	14098\\
10000	13796\\
12000	13609\\
14000	13936\\
16000	14312\\
18000	14090\\
20000	13950\\
22000	13737\\
24000	13884\\
26000	13871\\
28000	14070\\
30000	13967\\
32000	14121\\
34000	13975\\
38000	13611\\
40000	14432\\
42000	14116\\
44000	14509\\
46000	14117\\
48000	14164\\
52000	13656\\
54000	14097\\
56000	14250\\
58000	13593\\
60000	14350\\
62000	13969\\
64000	14112\\
66000	14087\\
68000	14133\\
70000	13761\\
72000	14353\\
74000	14394\\
76000	13725\\
78000	14140\\
80000	13651\\
82000	13927\\
86000	14090\\
88000	13778\\
90000	13524\\
92000	13936\\
94000	14421\\
96000	13878\\
98000	13771\\
1e+05	13785\\
};
\addplot [color=mycolor10, forget plot]
  table[row sep=crcr]{%
0	4395.3\\
2000	17670\\
4000	21156\\
6000	21204\\
8000	21666\\
10000	21459\\
12000	21921\\
14000	22535\\
16000	22223\\
18000	21430\\
20000	21392\\
22000	21948\\
24000	21465\\
26000	21892\\
28000	22087\\
30000	21824\\
32000	21673\\
34000	21843\\
36000	21564\\
38000	21698\\
40000	21649\\
42000	21710\\
44000	22102\\
46000	21487\\
48000	21670\\
50000	21674\\
52000	21973\\
54000	21920\\
56000	21624\\
58000	22101\\
60000	22079\\
62000	21973\\
64000	21640\\
66000	22081\\
68000	22129\\
70000	21808\\
72000	21412\\
74000	21779\\
76000	21883\\
78000	21656\\
80000	22002\\
82000	21893\\
84000	21648\\
86000	21856\\
88000	21557\\
90000	21541\\
92000	22258\\
94000	22128\\
96000	21569\\
98000	22169\\
1e+05	21420\\
};
\addplot [color=mycolor11, forget plot]
  table[row sep=crcr]{%
0	4395.3\\
2000	9728\\
4000	9050\\
6000	8800.4\\
8000	8716.4\\
10000	8537.9\\
14000	8682.7\\
16000	8909.5\\
18000	8779.6\\
22000	8603.3\\
24000	8672.6\\
26000	8829.4\\
28000	8527.8\\
30000	8409.7\\
32000	8469.5\\
34000	8571.1\\
36000	8738\\
38000	8882.9\\
40000	8838\\
42000	8573.7\\
44000	8561.6\\
46000	8681.5\\
50000	8652.4\\
52000	8418.4\\
54000	8469.5\\
56000	8806.4\\
58000	8825.5\\
60000	8674.6\\
62000	8604.5\\
64000	8403.9\\
66000	8686\\
68000	8767.2\\
70000	8335.8\\
72000	8605.4\\
76000	8323.8\\
80000	8397.3\\
82000	8571.9\\
84000	8650.3\\
88000	8361.2\\
90000	8591.6\\
92000	8741.9\\
96000	8661.1\\
98000	8431.7\\
1e+05	8602.1\\
};
\addplot [color=mycolor12, forget plot]
  table[row sep=crcr]{%
0	4395.3\\
2000	12907\\
4000	14668\\
8000	14521\\
10000	14758\\
12000	14551\\
16000	14581\\
18000	14841\\
20000	14645\\
22000	14730\\
24000	14650\\
26000	14334\\
28000	14741\\
30000	14816\\
32000	14190\\
34000	14340\\
36000	14748\\
38000	14933\\
40000	14670\\
42000	14943\\
44000	14534\\
46000	14864\\
48000	14301\\
50000	14950\\
52000	14163\\
54000	14855\\
56000	14858\\
58000	14252\\
60000	14277\\
62000	14505\\
64000	14601\\
66000	14440\\
68000	14552\\
70000	15055\\
72000	14842\\
74000	14852\\
78000	14308\\
80000	14523\\
82000	14507\\
84000	14861\\
86000	14929\\
88000	14792\\
90000	14467\\
92000	14815\\
94000	14610\\
96000	14505\\
98000	13846\\
1e+05	14588\\
};
\addplot [color=mycolor13, forget plot]
  table[row sep=crcr]{%
0	4395.3\\
2000	10105\\
4000	11064\\
6000	10937\\
8000	11401\\
10000	11049\\
12000	11163\\
14000	11143\\
16000	10893\\
18000	10689\\
20000	11102\\
22000	11227\\
24000	11070\\
26000	11265\\
28000	10898\\
30000	11212\\
32000	11253\\
34000	11357\\
36000	11329\\
38000	11396\\
40000	11316\\
44000	11326\\
46000	11065\\
48000	10896\\
50000	11145\\
52000	10972\\
54000	11234\\
56000	11101\\
58000	11036\\
60000	10806\\
62000	11099\\
64000	11107\\
68000	11019\\
70000	11174\\
72000	11013\\
74000	11207\\
76000	11074\\
78000	10974\\
80000	10712\\
82000	10922\\
84000	11052\\
86000	11133\\
88000	11001\\
90000	10966\\
92000	11129\\
94000	10988\\
96000	10790\\
98000	11053\\
1e+05	10742\\
};
\addplot [color=mycolor14, forget plot]
  table[row sep=crcr]{%
0	4395.3\\
2000	18238\\
4000	23005\\
6000	23837\\
8000	24215\\
10000	24825\\
14000	23907\\
16000	23689\\
18000	24239\\
20000	23634\\
22000	24199\\
24000	24395\\
26000	24520\\
28000	24052\\
30000	24129\\
32000	24649\\
34000	23657\\
36000	24233\\
38000	24252\\
40000	24454\\
42000	23671\\
44000	24541\\
46000	23315\\
48000	24801\\
50000	24272\\
52000	24106\\
54000	24384\\
56000	23627\\
58000	24199\\
60000	24081\\
62000	24037\\
64000	24461\\
66000	24512\\
68000	23765\\
70000	23838\\
72000	24302\\
74000	24466\\
78000	23803\\
80000	25237\\
82000	24188\\
84000	24204\\
86000	25136\\
88000	24383\\
90000	23245\\
92000	24068\\
94000	24685\\
96000	23853\\
1e+05	24592\\
};
\addplot [color=mycolor15, forget plot]
  table[row sep=crcr]{%
0	4395.3\\
2000	13963\\
4000	16386\\
6000	16230\\
8000	16500\\
10000	16897\\
12000	16732\\
14000	16421\\
16000	17019\\
18000	16673\\
20000	16712\\
22000	16672\\
24000	17082\\
26000	16936\\
28000	17159\\
30000	16611\\
32000	17211\\
34000	16436\\
36000	16392\\
38000	17017\\
40000	16884\\
42000	16256\\
44000	16510\\
46000	17240\\
48000	17158\\
50000	16794\\
52000	16686\\
54000	16959\\
56000	16957\\
58000	16889\\
60000	17012\\
62000	16798\\
64000	16760\\
68000	17204\\
70000	16737\\
74000	16991\\
76000	16588\\
78000	17406\\
80000	16568\\
84000	17391\\
86000	16558\\
88000	17148\\
90000	17278\\
92000	17036\\
94000	16964\\
96000	17161\\
98000	17059\\
1e+05	16848\\
};
\addplot [color=mycolor16, forget plot]
  table[row sep=crcr]{%
0	4395.3\\
2000	8107.8\\
4000	7151.7\\
6000	7067.3\\
8000	6899.7\\
10000	6811.1\\
12000	6912.1\\
14000	6803\\
16000	6856.2\\
18000	7029.2\\
22000	7039.6\\
24000	6993.3\\
26000	7116.8\\
28000	7090.3\\
32000	6971.1\\
34000	7125.8\\
36000	7015.6\\
38000	7068.8\\
40000	6860.2\\
42000	6921.7\\
44000	7050.7\\
46000	6900.1\\
48000	7061.5\\
50000	7082.8\\
52000	6983.7\\
54000	7107.5\\
56000	7019.7\\
58000	6852.8\\
60000	7003.8\\
62000	6851\\
64000	7009.3\\
66000	7039.7\\
68000	7154.8\\
70000	7099\\
72000	6984.3\\
74000	7171.8\\
76000	6982.3\\
78000	7081.7\\
80000	6894.5\\
82000	7001.4\\
84000	6821\\
86000	7187.6\\
88000	7113.9\\
90000	7007.1\\
92000	6863.4\\
94000	7106.6\\
96000	6852.5\\
98000	6928.7\\
1e+05	6858\\
};
\addplot [color=mycolor17, forget plot]
  table[row sep=crcr]{%
0	4395.3\\
2000	13239\\
4000	14941\\
6000	15982\\
8000	15755\\
10000	16092\\
12000	15821\\
14000	16128\\
16000	16006\\
18000	15717\\
20000	16440\\
22000	15796\\
24000	15717\\
26000	15841\\
28000	16090\\
30000	15671\\
32000	15426\\
34000	16396\\
38000	15899\\
40000	15970\\
42000	15868\\
44000	15989\\
46000	16228\\
48000	16353\\
50000	15804\\
52000	15866\\
54000	16218\\
56000	15392\\
58000	15863\\
60000	16140\\
62000	16303\\
64000	15808\\
68000	16143\\
70000	15727\\
72000	16271\\
74000	16052\\
76000	16242\\
80000	16256\\
82000	16105\\
84000	16495\\
86000	15813\\
88000	15654\\
90000	15282\\
92000	16097\\
94000	15708\\
96000	15939\\
98000	15957\\
1e+05	16100\\
};
\addplot [color=mycolor18, forget plot]
  table[row sep=crcr]{%
0	4395.3\\
2000	18488\\
4000	24692\\
8000	24435\\
10000	24538\\
12000	24080\\
14000	24433\\
16000	23665\\
18000	24330\\
20000	24550\\
22000	24982\\
24000	24581\\
26000	24884\\
28000	23898\\
30000	23985\\
32000	24339\\
34000	24114\\
36000	23216\\
38000	23750\\
40000	24442\\
42000	24125\\
44000	24757\\
46000	24121\\
48000	24196\\
50000	24025\\
52000	23573\\
56000	24195\\
58000	24129\\
60000	24303\\
62000	24158\\
64000	24564\\
68000	23743\\
70000	24197\\
72000	24281\\
74000	23437\\
76000	23739\\
78000	24619\\
80000	23970\\
82000	25131\\
84000	24872\\
86000	23963\\
88000	23993\\
90000	23191\\
92000	23359\\
94000	24108\\
96000	24467\\
98000	23203\\
1e+05	24037\\
};
\addplot [color=mycolor19, forget plot]
  table[row sep=crcr]{%
0	4395.3\\
2000	12378\\
4000	14795\\
6000	14856\\
8000	14499\\
10000	15106\\
12000	14845\\
14000	14915\\
18000	14431\\
20000	15033\\
22000	14952\\
24000	14735\\
26000	14848\\
28000	14649\\
30000	14831\\
32000	14603\\
34000	14858\\
36000	14436\\
38000	14679\\
40000	14818\\
42000	14444\\
44000	14851\\
46000	14606\\
48000	14781\\
54000	14837\\
56000	14684\\
58000	15152\\
60000	15061\\
62000	14881\\
64000	14634\\
68000	14443\\
70000	14675\\
72000	14637\\
74000	14685\\
76000	14645\\
78000	14665\\
80000	14811\\
82000	14784\\
84000	14609\\
86000	15230\\
88000	15178\\
90000	14764\\
92000	15062\\
94000	14654\\
96000	14173\\
98000	14312\\
1e+05	15362\\
};
\addplot [color=mycolor20, forget plot]
  table[row sep=crcr]{%
0	4395.3\\
2000	15163\\
4000	18098\\
6000	18463\\
8000	18900\\
10000	18899\\
12000	18551\\
14000	18617\\
16000	18870\\
18000	19044\\
20000	18846\\
22000	19176\\
24000	19257\\
26000	19265\\
28000	18935\\
30000	19311\\
32000	19126\\
34000	19020\\
36000	19804\\
38000	18946\\
40000	18487\\
42000	19712\\
44000	18969\\
46000	18616\\
48000	19091\\
50000	19096\\
52000	19415\\
56000	18933\\
58000	18753\\
60000	19545\\
62000	18800\\
64000	19546\\
68000	19022\\
70000	19153\\
72000	19069\\
74000	18743\\
76000	19309\\
78000	19531\\
80000	18347\\
82000	19380\\
84000	18986\\
86000	19299\\
88000	19114\\
90000	19237\\
92000	19120\\
94000	19309\\
96000	19305\\
98000	19715\\
1e+05	19002\\
};
\addplot [color=mycolor21, forget plot]
  table[row sep=crcr]{%
0	4395.3\\
2000	8330.8\\
4000	9755.8\\
6000	9683.8\\
8000	9909.4\\
12000	9569.5\\
14000	9778.8\\
16000	9751.3\\
18000	9602\\
20000	9904.1\\
22000	9763.5\\
24000	9472.6\\
26000	9484.2\\
28000	9753.1\\
30000	9623.8\\
32000	9777\\
36000	9616.3\\
38000	9645.8\\
40000	9844.6\\
42000	9549.1\\
44000	9565.6\\
46000	9692.6\\
48000	9689.5\\
50000	9616.1\\
52000	9673\\
54000	9625.9\\
56000	9744.2\\
58000	9394.5\\
60000	9513.7\\
62000	9490.3\\
64000	9672.4\\
66000	9637.7\\
68000	9695.5\\
70000	9574.3\\
72000	9739.6\\
74000	9468.3\\
76000	9626.5\\
78000	9510.2\\
80000	9569\\
84000	9591.4\\
86000	9794.8\\
88000	9746.9\\
90000	9770.7\\
92000	9734.6\\
94000	9611.7\\
98000	9716.5\\
1e+05	9583\\
};
\addplot [color=mycolor22, forget plot]
  table[row sep=crcr]{%
0	4395.3\\
2000	21781\\
4000	26239\\
6000	26291\\
10000	25670\\
12000	26220\\
14000	25319\\
16000	25135\\
18000	25504\\
20000	24166\\
22000	24585\\
24000	25095\\
26000	25225\\
28000	25172\\
30000	25719\\
32000	25536\\
34000	25600\\
36000	25999\\
38000	25029\\
40000	24555\\
42000	25067\\
44000	24977\\
46000	24068\\
48000	24803\\
50000	25434\\
52000	25319\\
54000	25014\\
56000	25097\\
58000	24928\\
60000	24392\\
62000	23549\\
64000	23978\\
66000	23892\\
68000	25280\\
70000	25297\\
72000	24587\\
74000	24432\\
76000	25612\\
78000	25332\\
80000	25239\\
82000	24427\\
84000	24793\\
86000	25093\\
88000	24514\\
90000	25336\\
92000	24904\\
94000	24919\\
98000	24616\\
1e+05	25033\\
};
\addplot [color=mycolor23, forget plot]
  table[row sep=crcr]{%
0	4395.3\\
2000	10887\\
4000	12471\\
6000	12649\\
10000	12204\\
12000	12217\\
14000	12287\\
16000	12225\\
18000	12355\\
22000	12012\\
24000	12313\\
26000	12245\\
28000	12301\\
30000	12136\\
32000	12226\\
34000	12366\\
36000	11973\\
38000	12041\\
40000	12279\\
42000	12265\\
44000	12418\\
46000	12381\\
48000	12508\\
50000	12468\\
52000	12066\\
54000	12078\\
60000	12322\\
62000	12046\\
64000	12287\\
66000	12467\\
68000	12217\\
70000	12157\\
72000	12054\\
74000	12195\\
76000	12415\\
78000	12151\\
80000	12047\\
82000	12149\\
84000	11955\\
86000	12036\\
90000	12576\\
92000	12348\\
94000	12217\\
96000	12265\\
98000	12375\\
1e+05	12248\\
};
\addplot [color=mycolor24, forget plot]
  table[row sep=crcr]{%
0	4395.3\\
2000	5526.5\\
4000	4904.5\\
6000	4690.3\\
8000	4738.4\\
10000	4860.1\\
12000	4828.4\\
14000	4860.9\\
16000	4933.9\\
18000	5056\\
20000	4801.1\\
22000	4778.1\\
24000	4707.9\\
26000	4814.1\\
28000	4890.7\\
30000	4807.7\\
32000	4812.4\\
36000	4851.9\\
38000	4923.4\\
40000	4825.2\\
42000	4929.6\\
44000	4812.1\\
46000	4931.6\\
48000	4746.1\\
50000	4996.5\\
52000	5015.9\\
54000	4886.4\\
56000	4854.8\\
58000	4809.2\\
60000	4898.3\\
62000	4738.4\\
64000	4682.9\\
66000	4708.3\\
68000	4802.4\\
70000	4973.2\\
72000	4767.9\\
74000	4877.5\\
76000	4661.1\\
78000	4813.2\\
80000	4860.5\\
82000	4646.5\\
84000	4745.4\\
86000	4754.4\\
88000	4608.6\\
90000	4868.6\\
92000	4863.6\\
94000	4807.4\\
96000	4884.4\\
98000	4717.4\\
1e+05	4766.7\\
};
\addplot [color=mycolor25, forget plot]
  table[row sep=crcr]{%
0	4395.3\\
2000	10710\\
4000	12175\\
6000	12209\\
8000	12097\\
10000	12395\\
12000	12180\\
14000	12052\\
16000	12193\\
22000	12065\\
24000	12404\\
26000	12140\\
28000	12324\\
30000	12324\\
32000	12650\\
34000	12038\\
36000	11940\\
38000	11982\\
40000	12382\\
42000	12103\\
44000	12470\\
46000	12494\\
48000	12339\\
50000	12426\\
52000	12806\\
54000	12273\\
56000	12641\\
58000	12786\\
60000	12718\\
62000	12465\\
64000	12177\\
66000	12635\\
68000	12186\\
70000	12523\\
72000	12520\\
74000	12827\\
76000	12445\\
78000	12746\\
80000	12434\\
84000	11982\\
86000	12877\\
88000	12532\\
90000	12157\\
92000	12360\\
94000	12471\\
96000	12408\\
98000	12539\\
1e+05	12468\\
};
\addplot [color=mycolor26, forget plot]
  table[row sep=crcr]{%
0	4395.3\\
2000	7429.1\\
4000	8022.3\\
6000	8113.5\\
8000	8307.5\\
12000	8134\\
14000	8161.8\\
16000	8516.5\\
18000	8167.4\\
20000	8389.2\\
22000	8316.4\\
24000	8182.8\\
26000	8354.2\\
28000	8493\\
30000	8231.1\\
32000	8136.6\\
34000	8173.7\\
36000	8339.3\\
38000	8348.6\\
40000	8395.6\\
42000	8124.6\\
44000	8275.1\\
46000	8258.1\\
48000	8602.6\\
50000	8429.7\\
52000	8568.8\\
54000	8325.9\\
56000	8124.4\\
58000	8607.3\\
60000	8205.3\\
62000	8192\\
64000	8029\\
66000	8388.1\\
68000	8510.6\\
72000	8699.1\\
74000	8313.9\\
76000	8167.3\\
78000	8205.1\\
80000	8080.6\\
82000	8206.6\\
84000	8282\\
86000	8275.8\\
88000	8358.2\\
90000	8025.6\\
94000	8254.4\\
96000	8353.6\\
98000	8228.3\\
1e+05	8352.6\\
};
\addplot [color=mycolor27, forget plot]
  table[row sep=crcr]{%
0	4395.3\\
2000	9509.9\\
4000	9515.8\\
6000	9924.2\\
8000	9984.1\\
10000	9603.8\\
12000	9934.6\\
14000	9553.2\\
16000	9843.1\\
18000	9857.7\\
20000	9598.9\\
22000	9723.1\\
24000	9516.1\\
26000	9649.6\\
28000	9409.7\\
30000	9888\\
32000	9666.9\\
34000	9960.9\\
36000	10071\\
38000	9574.5\\
40000	9743.3\\
42000	9751.6\\
44000	9511.7\\
46000	9427.9\\
48000	9601\\
50000	9549.9\\
52000	9959.8\\
54000	9978.4\\
56000	9760.2\\
58000	9452.3\\
60000	9488.4\\
62000	9667.5\\
64000	9779\\
66000	9728.1\\
68000	9752.7\\
70000	9697.4\\
72000	10103\\
74000	9659\\
76000	9986.1\\
78000	9457.7\\
80000	9998.3\\
82000	9967.2\\
84000	9615\\
86000	9784.6\\
88000	10009\\
90000	9850.2\\
92000	9983\\
94000	9750.9\\
96000	9861.2\\
98000	9748.9\\
1e+05	9767.4\\
};
\addplot [color=mycolor28, forget plot]
  table[row sep=crcr]{%
0	4395.3\\
2000	6256.2\\
4000	5678\\
6000	5617.4\\
8000	5635.8\\
10000	5556.4\\
12000	5359.8\\
14000	5503.4\\
16000	5355\\
18000	5689.1\\
20000	5321.9\\
22000	5440.2\\
24000	5492\\
26000	5633\\
30000	5523.9\\
32000	5661\\
34000	5436.1\\
36000	5535.9\\
38000	5422.5\\
40000	5590.9\\
42000	5500.2\\
44000	5538.9\\
46000	5456.9\\
48000	5327.8\\
50000	5221.5\\
52000	5514.2\\
54000	5420.4\\
56000	5432.2\\
58000	5641.5\\
60000	5469.5\\
62000	5513\\
64000	5452.6\\
66000	5533\\
68000	5571.3\\
70000	5352.8\\
72000	5475.4\\
74000	5351.8\\
76000	5481.9\\
78000	5647.6\\
80000	5433.3\\
84000	5309.8\\
86000	5520.9\\
88000	5594.2\\
90000	5304.1\\
92000	5476.3\\
94000	5484.8\\
96000	5388.4\\
98000	5342.6\\
1e+05	5419.2\\
};
\addplot [color=mycolor29, forget plot]
  table[row sep=crcr]{%
0	4395.3\\
2000	11020\\
4000	11509\\
6000	11662\\
8000	11417\\
10000	11559\\
12000	11315\\
14000	11437\\
16000	11364\\
18000	11417\\
20000	11006\\
24000	11379\\
26000	11498\\
28000	11895\\
30000	11511\\
32000	11752\\
34000	11356\\
36000	11453\\
38000	11760\\
42000	11624\\
44000	11633\\
46000	11831\\
48000	11927\\
50000	11584\\
52000	11571\\
54000	11772\\
56000	11495\\
58000	11146\\
60000	11704\\
62000	12065\\
64000	11290\\
66000	11702\\
68000	11263\\
70000	11488\\
72000	11509\\
74000	11671\\
76000	11505\\
78000	11176\\
80000	10995\\
82000	11531\\
84000	12004\\
86000	11479\\
88000	11632\\
90000	11724\\
92000	11390\\
94000	11259\\
96000	11727\\
98000	11460\\
1e+05	11509\\
};
\addplot [color=mycolor30, forget plot]
  table[row sep=crcr]{%
0	4395.3\\
2000	6968.9\\
4000	7736.2\\
6000	7923\\
8000	7900.5\\
10000	7773.7\\
12000	8052.6\\
14000	7806.4\\
16000	7668.6\\
18000	7697.5\\
20000	7949.2\\
22000	8042.3\\
24000	7803.3\\
26000	7856.6\\
28000	7828.8\\
30000	7755.3\\
34000	7970.9\\
36000	7875.2\\
38000	7941.4\\
40000	7834.3\\
46000	8003.7\\
48000	7835.7\\
50000	7870.6\\
52000	7788.1\\
54000	8034\\
56000	7804\\
58000	8033.4\\
62000	7826.3\\
64000	7962.7\\
66000	7780\\
70000	7960.6\\
74000	7935.2\\
76000	8015\\
78000	7986.6\\
80000	7804.2\\
82000	7741.7\\
84000	8016.6\\
86000	7678\\
88000	7744.3\\
90000	7756.3\\
92000	7793.4\\
94000	7752.8\\
96000	8002.9\\
98000	7906.2\\
1e+05	7889.8\\
};
\addplot [color=mycolor31, forget plot]
  table[row sep=crcr]{%
0	4395.3\\
2000	7008.5\\
4000	6774.8\\
6000	6262.1\\
8000	6346\\
10000	6375.4\\
12000	6271.4\\
14000	6470.6\\
16000	6321\\
18000	6460.9\\
20000	6285.8\\
22000	6243.7\\
24000	6526.4\\
26000	6435.1\\
28000	6476.2\\
30000	6402.9\\
32000	6353.8\\
34000	6495.7\\
36000	6341.2\\
38000	6263.3\\
40000	6534.4\\
42000	6446.4\\
44000	6527.2\\
46000	6436.8\\
50000	6458.5\\
52000	6332.4\\
54000	6360.2\\
56000	6542.9\\
58000	6682.2\\
60000	6546.3\\
62000	6377.9\\
64000	6582.3\\
66000	6570.9\\
68000	6378.4\\
70000	6818.4\\
72000	6632.1\\
76000	6369.7\\
78000	6379.3\\
80000	6504.1\\
82000	6400.6\\
84000	6507.9\\
86000	6559.4\\
88000	6312.1\\
90000	6713.4\\
92000	6663.2\\
94000	6524.4\\
96000	6586.2\\
98000	6268.7\\
1e+05	6578.1\\
};
\addplot [color=mycolor32, forget plot]
  table[row sep=crcr]{%
0	4395.3\\
2000	10345\\
4000	11184\\
6000	10630\\
8000	10759\\
10000	10711\\
12000	10980\\
14000	10376\\
16000	10466\\
18000	10620\\
20000	10161\\
22000	10852\\
24000	10743\\
26000	10864\\
28000	10480\\
30000	10281\\
32000	10485\\
34000	10468\\
36000	10693\\
38000	10608\\
42000	10240\\
44000	10724\\
46000	10636\\
48000	10266\\
50000	10086\\
52000	10567\\
54000	10672\\
56000	10368\\
58000	10325\\
60000	10564\\
62000	10273\\
64000	10592\\
66000	10206\\
70000	10403\\
72000	10338\\
74000	10402\\
76000	10574\\
78000	10528\\
80000	10362\\
82000	10457\\
84000	10355\\
86000	10593\\
88000	10541\\
90000	10398\\
92000	10446\\
94000	10375\\
96000	10598\\
98000	10670\\
1e+05	10149\\
};
\addplot [color=mycolor33, forget plot]
  table[row sep=crcr]{%
0	4395.3\\
2000	6887.6\\
4000	7934.1\\
6000	7942.6\\
8000	7881.6\\
10000	8294.5\\
12000	8255\\
14000	8626.3\\
16000	7920.1\\
18000	8151.1\\
20000	8085\\
22000	8390.9\\
24000	8317.3\\
26000	8012.8\\
28000	8075.8\\
30000	8263.1\\
32000	8376.2\\
34000	7823.4\\
36000	8015.7\\
38000	8256.7\\
40000	7910\\
42000	8247.7\\
44000	8302.8\\
46000	8121\\
50000	8367\\
52000	8242\\
54000	8276.6\\
56000	8102.9\\
58000	8483.5\\
60000	8123.5\\
62000	8226.4\\
64000	8114.2\\
66000	8032.6\\
68000	8175.3\\
70000	8101.2\\
72000	8295.8\\
74000	8215.9\\
76000	8057.2\\
78000	8197.8\\
80000	8393.1\\
82000	7865.5\\
84000	8166.3\\
86000	8428.8\\
88000	8189.9\\
90000	8104.2\\
92000	8052.7\\
94000	8254.6\\
96000	8241.5\\
98000	8129.4\\
1e+05	8159.8\\
};
\addplot [color=mycolor34, forget plot]
  table[row sep=crcr]{%
0	4395.3\\
2000	4866.4\\
4000	4038.3\\
6000	3930.5\\
8000	4057.1\\
10000	4120.5\\
12000	4046.4\\
16000	4044.3\\
18000	3951.6\\
20000	4131.6\\
22000	4208.2\\
26000	4028.3\\
30000	4153.1\\
32000	4077.4\\
34000	4104.7\\
36000	3993.3\\
38000	4002.9\\
40000	4144.2\\
42000	4142.2\\
44000	4099.4\\
46000	3994.1\\
48000	4145.1\\
50000	4002.9\\
52000	4200.7\\
54000	4003.1\\
56000	4174.2\\
58000	3941.7\\
60000	4078.7\\
62000	4056.6\\
64000	4093.9\\
66000	3991.3\\
68000	4032.9\\
70000	4265.2\\
72000	4092.9\\
74000	4090.3\\
76000	3948.2\\
78000	3911.2\\
80000	4167.3\\
82000	3980.7\\
84000	4004.7\\
86000	3972.9\\
88000	4032.5\\
90000	4073.9\\
92000	4101.7\\
94000	4071.7\\
96000	4081.9\\
98000	4057.6\\
1e+05	3922.5\\
};
\addplot [color=mycolor35]
  table[row sep=crcr]{%
0	4395.3\\
2000	9096.1\\
4000	9866.2\\
8000	10261\\
12000	10087\\
14000	10206\\
16000	10116\\
20000	10137\\
22000	10222\\
24000	10400\\
26000	10073\\
28000	10193\\
32000	10087\\
34000	10240\\
36000	10296\\
38000	10065\\
40000	10258\\
44000	9898\\
46000	10105\\
48000	10112\\
50000	10289\\
52000	10185\\
54000	10255\\
58000	10327\\
60000	9995.6\\
62000	10221\\
64000	9962.3\\
66000	10034\\
68000	10151\\
70000	10104\\
72000	10208\\
74000	9800.3\\
76000	9996.1\\
78000	9962.8\\
82000	10319\\
84000	10347\\
88000	10248\\
90000	10163\\
92000	10373\\
94000	10222\\
98000	10224\\
1e+05	10336\\
};
\addlegendentry{\ClasskIILegend}

\addplot [color=mycolor36, forget plot]
  table[row sep=crcr]{%
0	4395.3\\
2000	5047.5\\
4000	4455.1\\
6000	4365\\
8000	4181.5\\
10000	4341.6\\
12000	4098.7\\
14000	4127.7\\
16000	4203\\
18000	4333.1\\
20000	4220.3\\
22000	4291.1\\
24000	4254.2\\
26000	4435.8\\
28000	4312.7\\
30000	4422.8\\
32000	4473.9\\
34000	4143.4\\
36000	4298.1\\
38000	4315.8\\
40000	4594.8\\
42000	4451\\
44000	4453.6\\
46000	4227.7\\
48000	4388\\
50000	4341.7\\
52000	4229.4\\
54000	4330\\
56000	4456.5\\
58000	4310.6\\
60000	4214.2\\
64000	4276.3\\
66000	4329.5\\
68000	4370.9\\
70000	4325.5\\
72000	4187.2\\
74000	4346.9\\
76000	4262.5\\
78000	4194.9\\
80000	4424.2\\
82000	4362.4\\
84000	4239.6\\
86000	4199.3\\
88000	4346\\
90000	4285.6\\
92000	4334.7\\
94000	4216.9\\
96000	4315.8\\
98000	4326.4\\
1e+05	4288.8\\
};
\addplot [color=mycolor37, forget plot]
  table[row sep=crcr]{%
0	4395.3\\
2000	4089.2\\
4000	3689.2\\
6000	3604.9\\
8000	3667.5\\
10000	3557.1\\
12000	3574.7\\
14000	3648.5\\
16000	3817.3\\
18000	3685.6\\
20000	3628.1\\
22000	3714.3\\
24000	3666.4\\
26000	3690.5\\
28000	3579.3\\
30000	3674.5\\
32000	3677.2\\
34000	3620.8\\
38000	3618.1\\
40000	3528.4\\
42000	3677.9\\
44000	3610\\
46000	3578.5\\
48000	3598.3\\
50000	3735.3\\
52000	3615.4\\
54000	3607.5\\
56000	3643.5\\
58000	3693.4\\
60000	3584.2\\
62000	3697.1\\
64000	3595.3\\
66000	3648.1\\
68000	3614.3\\
70000	3658.6\\
72000	3765.2\\
74000	3674.6\\
76000	3759.3\\
78000	3632.2\\
80000	3681.4\\
82000	3657.4\\
86000	3639.5\\
88000	3680.4\\
90000	3670.9\\
92000	3631.5\\
94000	3671.2\\
96000	3614.7\\
1e+05	3676.9\\
};
\addplot [color=mycolor38, forget plot]
  table[row sep=crcr]{%
0	4395.3\\
2000	3427.6\\
4000	3172.3\\
6000	3008.2\\
8000	3125.1\\
10000	3131.6\\
12000	2968.4\\
14000	3037.5\\
16000	3000.5\\
18000	3106.6\\
20000	3048.6\\
22000	3009.8\\
24000	2995.4\\
26000	3109\\
28000	3018.7\\
30000	3138\\
32000	3054.6\\
36000	3013.1\\
38000	3064.3\\
40000	2976.4\\
42000	3019.7\\
44000	3108.1\\
46000	2992.1\\
50000	2908\\
52000	3005.8\\
54000	3001.8\\
56000	2984.2\\
60000	3079.2\\
62000	2976.5\\
64000	2894.6\\
66000	3048.4\\
68000	3047.2\\
70000	3082.8\\
72000	3037.7\\
74000	2941.5\\
76000	3023.9\\
78000	3002.1\\
80000	3041.4\\
82000	3046\\
84000	3030.8\\
86000	3073.4\\
88000	3129.1\\
90000	3129.4\\
92000	2951.3\\
94000	2941.6\\
96000	3025.3\\
98000	3016.4\\
1e+05	3097.7\\
};
\addplot [color=mycolor39, forget plot]
  table[row sep=crcr]{%
0	4395.3\\
2000	11633\\
4000	12676\\
6000	12278\\
10000	11656\\
16000	11892\\
18000	11830\\
20000	11831\\
22000	11899\\
24000	12233\\
26000	11938\\
28000	11709\\
30000	12010\\
34000	11787\\
36000	11878\\
40000	11688\\
42000	11678\\
44000	12035\\
46000	11752\\
48000	11593\\
50000	11298\\
52000	11551\\
54000	11534\\
56000	12013\\
58000	11648\\
60000	11480\\
62000	11548\\
64000	11769\\
66000	11713\\
68000	11400\\
70000	11498\\
72000	11768\\
74000	12083\\
76000	11876\\
78000	11807\\
80000	12157\\
82000	11727\\
84000	11610\\
86000	12047\\
88000	11948\\
90000	12090\\
92000	12089\\
94000	11875\\
96000	11774\\
98000	12013\\
1e+05	11612\\
};
\addplot [color=mycolor40, forget plot]
  table[row sep=crcr]{%
0	4395.3\\
2000	5184.7\\
4000	5655.9\\
6000	5589.1\\
8000	5677.6\\
10000	5609\\
14000	5918.2\\
16000	5842.2\\
18000	5648.8\\
20000	5851.4\\
22000	5763.1\\
24000	5844.8\\
26000	5734.8\\
28000	5739.3\\
30000	5776.7\\
32000	5861.3\\
34000	5770.9\\
36000	5771\\
38000	5752.9\\
40000	5681.2\\
42000	5850.8\\
44000	5708.3\\
46000	5737.3\\
48000	5664.5\\
50000	5933.1\\
52000	5812.6\\
54000	5619.3\\
56000	5704.7\\
58000	5663.4\\
62000	5702.9\\
64000	5828.2\\
66000	5654.6\\
70000	5761.1\\
74000	5686.3\\
76000	5680.3\\
80000	5747.7\\
82000	5729.9\\
84000	5796.1\\
86000	5767.7\\
88000	5768.4\\
90000	5734\\
92000	5816.2\\
94000	5794.1\\
96000	5741.4\\
98000	5810.7\\
1e+05	5810.9\\
};
\addplot [color=mycolor41, forget plot]
  table[row sep=crcr]{%
0	4395.3\\
2000	4898.5\\
4000	5188.9\\
6000	5224.4\\
8000	5193.7\\
10000	5230.7\\
12000	5344.5\\
14000	5437.2\\
16000	5435.4\\
20000	5299\\
22000	5438.4\\
24000	5350.3\\
26000	5569.1\\
28000	5415.3\\
30000	5526.6\\
32000	5248\\
34000	5367.7\\
36000	5318\\
38000	5340.5\\
40000	5384.2\\
42000	5288.1\\
44000	5494.1\\
46000	5316.9\\
48000	5391.5\\
50000	5310.5\\
52000	5267.2\\
54000	5477.2\\
56000	5308.4\\
58000	5606.6\\
60000	5387.8\\
62000	5334.1\\
64000	5425.3\\
68000	5239.2\\
70000	5579.9\\
72000	5481.9\\
74000	5303.1\\
76000	5317\\
78000	5550.8\\
80000	5374.6\\
82000	5412.9\\
84000	5416.8\\
86000	5354.4\\
88000	5356.2\\
90000	5438.9\\
92000	5366.4\\
94000	5240\\
96000	5377\\
98000	5294.6\\
1e+05	5455.3\\
};
\addplot [color=mycolor42, forget plot]
  table[row sep=crcr]{%
0	4395.3\\
2000	14373\\
4000	18951\\
6000	19658\\
8000	20767\\
10000	19913\\
12000	21007\\
14000	20871\\
18000	20385\\
20000	20963\\
22000	20882\\
24000	20363\\
26000	20041\\
28000	20494\\
30000	20318\\
32000	20624\\
34000	20523\\
38000	20635\\
40000	20944\\
44000	20917\\
46000	19985\\
48000	20359\\
50000	20654\\
52000	20039\\
54000	21024\\
56000	20141\\
58000	20610\\
60000	20427\\
62000	20340\\
64000	20424\\
66000	19866\\
68000	20694\\
70000	20397\\
72000	20649\\
74000	20699\\
76000	20514\\
78000	20203\\
80000	20643\\
82000	20328\\
84000	20236\\
86000	21154\\
88000	20302\\
90000	20219\\
92000	20597\\
94000	20464\\
96000	20208\\
98000	21031\\
1e+05	20312\\
};
\addplot [color=mycolor43, forget plot]
  table[row sep=crcr]{%
0	4395.3\\
2000	5446.8\\
4000	5192\\
6000	5029.6\\
8000	5075.7\\
10000	5063.6\\
12000	5036.8\\
14000	4910.6\\
16000	4867\\
18000	5079.2\\
20000	5084.3\\
22000	4889.1\\
24000	4907.1\\
26000	4963\\
28000	4802.1\\
30000	4882.3\\
32000	4847.1\\
34000	5066.9\\
36000	4914.3\\
38000	4949.8\\
40000	4929.2\\
42000	4947.7\\
44000	4857.7\\
46000	4886.3\\
48000	4890.3\\
50000	4994.2\\
52000	4916.6\\
54000	5005.1\\
56000	5030.4\\
58000	5017.7\\
60000	4814.4\\
64000	4899.5\\
66000	4907.6\\
68000	5032.8\\
70000	4894.8\\
72000	4908.7\\
74000	5030.2\\
76000	4949.3\\
78000	4825.3\\
80000	4833\\
82000	4942.8\\
84000	4959.8\\
86000	5031.9\\
88000	4950.9\\
90000	4913.6\\
92000	5067.9\\
94000	4972\\
96000	5052\\
98000	4987.4\\
1e+05	4954.2\\
};
\addplot [color=mycolor44, forget plot]
  table[row sep=crcr]{%
0	4395.3\\
2000	4133.6\\
4000	3589.6\\
6000	3509.6\\
8000	3457.4\\
10000	3427.3\\
12000	3411.4\\
14000	3322\\
16000	3436.4\\
18000	3407.3\\
20000	3524.1\\
22000	3387.8\\
24000	3329.5\\
26000	3429.4\\
28000	3320.6\\
30000	3379.9\\
32000	3310.3\\
34000	3292.4\\
36000	3388.5\\
38000	3409.7\\
40000	3407.2\\
42000	3225.7\\
44000	3357.4\\
46000	3478.7\\
48000	3382\\
50000	3191.1\\
52000	3179.4\\
54000	3382.2\\
56000	3199.5\\
58000	3369\\
60000	3200.4\\
62000	3401.2\\
64000	3341.4\\
66000	3353\\
68000	3396.9\\
70000	3321.6\\
72000	3362.9\\
74000	3289.9\\
76000	3293.1\\
78000	3426\\
80000	3278.9\\
82000	3327\\
84000	3395.6\\
88000	3512.4\\
90000	3417\\
92000	3349\\
94000	3434.6\\
96000	3353.4\\
1e+05	3391.1\\
};
\addplot [color=mycolor45, forget plot]
  table[row sep=crcr]{%
0	4395.3\\
2000	5672.7\\
4000	6696.3\\
6000	6967.2\\
8000	7011.5\\
10000	6941.8\\
12000	7011.5\\
14000	6919.3\\
16000	7025.2\\
18000	6768.7\\
20000	6929.9\\
22000	6720.4\\
24000	6916.5\\
26000	6818.4\\
28000	7136.9\\
30000	6788.2\\
32000	6853.4\\
34000	6860.4\\
36000	6977.7\\
38000	7062.2\\
40000	6974.6\\
42000	6807.4\\
44000	6848\\
46000	6953.3\\
48000	6998.5\\
50000	6934.2\\
52000	7110.9\\
54000	6873.9\\
56000	6867.6\\
58000	6890.1\\
60000	6814.7\\
64000	7051.7\\
66000	6932.5\\
68000	6967\\
70000	6771.1\\
72000	7112.4\\
74000	7059\\
76000	6893.4\\
78000	6876.7\\
80000	6791.1\\
82000	7098.5\\
84000	6977.7\\
86000	7091.1\\
88000	7002.8\\
90000	6874.4\\
92000	6826.2\\
94000	6955\\
96000	6994\\
98000	6872\\
1e+05	6725.6\\
};
\addplot [color=mycolor46, forget plot]
  table[row sep=crcr]{%
0	4395.3\\
2000	4056.6\\
4000	3985.4\\
6000	3998.5\\
8000	4153.8\\
10000	4006.6\\
12000	4146\\
14000	4169.5\\
16000	4148\\
18000	4244.4\\
20000	4021\\
22000	4100.3\\
24000	4041.1\\
26000	4108\\
28000	4262.7\\
30000	4245.2\\
32000	4242.1\\
34000	4141.3\\
38000	4035.6\\
42000	4123.2\\
44000	3904.3\\
46000	4178.5\\
48000	4101\\
50000	4084.3\\
52000	4026.9\\
54000	4216.4\\
56000	4182.7\\
58000	4239.1\\
60000	4089.7\\
62000	4089.6\\
64000	4150.5\\
66000	4039.9\\
68000	4114.6\\
70000	4132\\
72000	4008.1\\
74000	4069.7\\
78000	4034.2\\
80000	4156.1\\
82000	4186.9\\
84000	4048.7\\
86000	4050.6\\
90000	4027\\
92000	4054.5\\
94000	4160.8\\
96000	4080.7\\
98000	4136.6\\
1e+05	4087.4\\
};
\addplot [color=mycolor47, forget plot]
  table[row sep=crcr]{%
0	4395.3\\
2000	5509.7\\
4000	5366.4\\
6000	5330.8\\
8000	5368.8\\
10000	5241.4\\
12000	5247.3\\
14000	5373.4\\
16000	5286\\
18000	5394.1\\
20000	5228.5\\
22000	5273.4\\
24000	5376.6\\
26000	5347.7\\
28000	5211.6\\
30000	5310.2\\
32000	5351.5\\
34000	5287.8\\
38000	5296.5\\
40000	5082.2\\
42000	5182.2\\
44000	5395.7\\
46000	5427.3\\
48000	5432.7\\
50000	5399.5\\
52000	5141.4\\
54000	5281.4\\
56000	5356.3\\
58000	5392.9\\
60000	5203.3\\
62000	5289.1\\
64000	5230.1\\
66000	5360.7\\
68000	5377.2\\
70000	5320.9\\
72000	5227.9\\
74000	5224.1\\
76000	5281.6\\
78000	5370.3\\
80000	5329.6\\
82000	5371.7\\
84000	5336.5\\
88000	5387.8\\
90000	5211.3\\
92000	5370.1\\
94000	5271.6\\
96000	5208.6\\
98000	5248.9\\
1e+05	5251\\
};
\addplot [color=mycolor48, forget plot]
  table[row sep=crcr]{%
0	4395.3\\
2000	3417.5\\
4000	3000.6\\
8000	2944.1\\
12000	2987.1\\
14000	2890.3\\
16000	2916.4\\
18000	2955.3\\
20000	2945\\
22000	2917.1\\
24000	2955.9\\
26000	2937.7\\
28000	2896.4\\
30000	2950.1\\
32000	2947.6\\
34000	2904.2\\
36000	2942.4\\
38000	2954.5\\
40000	2945.7\\
42000	2982.9\\
44000	2972.5\\
46000	2917.6\\
48000	2953.1\\
50000	2954.4\\
52000	2931.1\\
60000	2930.5\\
64000	2925.4\\
66000	2964.4\\
68000	2978.3\\
70000	2922\\
72000	2937.7\\
76000	2914.3\\
78000	2946.6\\
80000	2953.4\\
82000	2972.5\\
84000	2922.5\\
86000	2959.5\\
88000	2939.4\\
90000	2984.9\\
92000	2913.4\\
94000	2900.4\\
1e+05	2968.4\\
};
\addplot [color=mycolor49, forget plot]
  table[row sep=crcr]{%
0	4395.3\\
2000	4900.7\\
4000	4475.4\\
6000	4306.4\\
8000	4319.8\\
10000	4498.3\\
12000	4357.5\\
14000	4369.1\\
16000	4363.1\\
18000	4504.8\\
20000	4369.6\\
22000	4401.1\\
24000	4344.8\\
26000	4318.4\\
28000	4312\\
30000	4187.7\\
32000	4405.8\\
34000	4451.5\\
36000	4205.1\\
40000	4414.8\\
42000	4225\\
44000	4229.6\\
46000	4409.1\\
48000	4436.2\\
50000	4241.7\\
52000	4441.8\\
54000	4530.3\\
56000	4298.8\\
58000	4446.7\\
60000	4380.4\\
62000	4521.2\\
64000	4370.8\\
66000	4454.7\\
68000	4279.6\\
70000	4340.1\\
72000	4228.1\\
74000	4407.9\\
76000	4326.5\\
78000	4373.4\\
80000	4494.2\\
82000	4419.6\\
84000	4470.4\\
86000	4302.7\\
88000	4345.3\\
90000	4406.8\\
92000	4287.9\\
94000	4319.7\\
96000	4392.5\\
98000	4419.9\\
1e+05	4326.4\\
};
\addplot [color=mycolor50, forget plot]
  table[row sep=crcr]{%
0	4395.3\\
2000	2603\\
4000	2481.6\\
6000	2411.1\\
12000	2494.1\\
14000	2481.8\\
16000	2483.2\\
18000	2399.9\\
20000	2482.5\\
22000	2549.9\\
24000	2561\\
26000	2440\\
28000	2431.1\\
30000	2472.7\\
32000	2468.4\\
36000	2443.7\\
38000	2522\\
40000	2403.8\\
42000	2467.6\\
44000	2506.3\\
46000	2444.8\\
50000	2528.4\\
52000	2521.5\\
54000	2482.2\\
56000	2542.7\\
58000	2478.1\\
60000	2463.7\\
62000	2486.3\\
64000	2450.6\\
66000	2525.2\\
68000	2544.5\\
70000	2476.5\\
72000	2439.2\\
74000	2475.6\\
76000	2441.5\\
78000	2475.2\\
80000	2431.1\\
82000	2498.2\\
86000	2520.7\\
88000	2479.2\\
90000	2459.8\\
92000	2490.1\\
94000	2534.7\\
96000	2501.6\\
98000	2482.9\\
1e+05	2404.2\\
};
\addplot [color=mycolor51, forget plot]
  table[row sep=crcr]{%
0	4395.3\\
2000	5451.9\\
4000	5811.3\\
6000	5775.6\\
8000	5771.7\\
10000	5853.1\\
12000	5877.8\\
14000	6024.9\\
16000	5735.5\\
18000	5975\\
20000	6002.3\\
22000	5979.9\\
24000	6041.1\\
28000	5991.2\\
30000	5945.9\\
32000	5871\\
34000	5868.6\\
36000	5992.3\\
38000	5787.3\\
40000	6058.2\\
42000	5875\\
44000	6124.2\\
46000	6054.3\\
48000	5823.8\\
50000	6047.4\\
52000	6038.8\\
54000	6056.1\\
56000	5670.2\\
58000	6053.5\\
60000	5937.1\\
62000	6351.3\\
64000	5939.9\\
66000	6002.8\\
68000	5953.8\\
70000	5827.8\\
72000	5889.3\\
74000	5777.1\\
76000	5947.1\\
78000	6032.5\\
80000	5854.1\\
82000	5966.9\\
84000	5854\\
86000	6019.2\\
88000	5891.2\\
90000	6013.6\\
92000	6107.7\\
94000	6121.9\\
96000	6107.6\\
98000	6131.7\\
1e+05	5970.6\\
};
\addplot [color=mycolor52, forget plot]
  table[row sep=crcr]{%
0	4395.3\\
2000	4061.7\\
4000	3718.7\\
6000	3636.7\\
8000	3588\\
10000	3502.7\\
12000	3540.4\\
14000	3477.5\\
16000	3468.7\\
18000	3431.9\\
20000	3502.6\\
22000	3475.4\\
24000	3507\\
26000	3589.5\\
28000	3572.8\\
30000	3527.7\\
32000	3573.9\\
34000	3699.9\\
36000	3501.3\\
38000	3561.6\\
42000	3671\\
44000	3435.3\\
46000	3584.3\\
48000	3404.8\\
50000	3578.6\\
52000	3525.6\\
54000	3409.2\\
56000	3437.9\\
58000	3481.4\\
60000	3556.9\\
62000	3540.1\\
64000	3502.9\\
66000	3594.3\\
68000	3504.9\\
70000	3436.3\\
72000	3556.4\\
74000	3547.3\\
76000	3667.8\\
78000	3435.9\\
80000	3476.7\\
82000	3553.2\\
84000	3542.4\\
86000	3609.1\\
88000	3584.5\\
90000	3574.1\\
92000	3354.2\\
94000	3468.4\\
96000	3572.9\\
98000	3655.2\\
1e+05	3565.5\\
};
\addplot [color=mycolor53, forget plot]
  table[row sep=crcr]{%
0	4395.3\\
2000	6101.8\\
4000	6333.5\\
6000	6063.2\\
8000	5730.2\\
10000	5669.4\\
14000	5776.1\\
16000	5777.5\\
18000	5547\\
20000	5673.5\\
22000	5763.4\\
24000	5963.6\\
26000	5765\\
28000	5719.6\\
32000	5737.5\\
34000	5619.3\\
36000	5805.3\\
38000	5709.2\\
40000	5672.6\\
42000	5874\\
44000	5494.9\\
46000	5829.9\\
48000	5773.5\\
50000	5795\\
52000	5661.2\\
54000	5586.6\\
56000	5553\\
58000	5805.8\\
60000	5790.2\\
62000	5554.4\\
64000	5501.8\\
66000	5556.4\\
68000	5795.1\\
70000	5494.3\\
72000	5324\\
74000	5527.7\\
76000	5799.3\\
78000	5776.6\\
80000	5728.3\\
82000	5504.2\\
84000	5775.2\\
86000	5738.8\\
88000	5868.8\\
90000	5628.5\\
92000	5722\\
94000	5544.6\\
96000	5759.1\\
98000	5869.7\\
1e+05	5896.1\\
};
\addplot [color=mycolor54, forget plot]
  table[row sep=crcr]{%
0	4395.3\\
2000	4703.4\\
4000	4629.4\\
6000	4588.4\\
8000	4488.3\\
10000	4549.9\\
12000	4515.9\\
14000	4504.2\\
18000	4522.7\\
20000	4638.4\\
22000	4613.4\\
24000	4539.6\\
26000	4580.5\\
28000	4521.6\\
30000	4648.2\\
32000	4625\\
34000	4555.8\\
36000	4660.5\\
38000	4492.8\\
40000	4620.4\\
42000	4685.7\\
44000	4543.7\\
46000	4575.2\\
48000	4474.4\\
50000	4569.2\\
52000	4542.2\\
54000	4546.3\\
56000	4591.7\\
58000	4516.8\\
60000	4570.9\\
62000	4557.3\\
64000	4646.8\\
66000	4704.6\\
68000	4778.7\\
70000	4684.5\\
72000	4571.9\\
74000	4616.4\\
76000	4736.2\\
78000	4570\\
80000	4539.2\\
82000	4693\\
84000	4602.6\\
88000	4522.4\\
90000	4658\\
92000	4428.3\\
94000	4434.5\\
96000	4505.3\\
98000	4550.4\\
1e+05	4516.5\\
};
\addplot [color=mycolor55, forget plot]
  table[row sep=crcr]{%
0	4395.3\\
2000	5193.5\\
4000	5634.1\\
6000	5748.6\\
8000	5637.3\\
10000	5921\\
12000	5892.2\\
16000	5746.4\\
18000	5649\\
20000	5866.9\\
22000	5732.1\\
24000	5824.6\\
26000	5668.8\\
28000	5787.5\\
32000	5696.1\\
34000	5690.8\\
36000	5720.5\\
38000	5665.5\\
40000	5665.5\\
42000	5875\\
44000	5787.7\\
48000	5729.8\\
52000	5799.9\\
56000	5693\\
58000	5849.5\\
60000	5828.8\\
62000	5884.9\\
64000	5830.3\\
68000	5693.9\\
70000	5776.2\\
72000	5796.3\\
74000	5781.9\\
76000	5817.1\\
78000	5895.6\\
80000	5802.5\\
82000	5738.1\\
84000	5753.5\\
86000	5837.4\\
88000	5719.2\\
90000	5801.2\\
92000	5790.1\\
94000	5642.6\\
96000	5655.7\\
98000	5788.3\\
1e+05	5750.4\\
};
\addplot [color=mycolor56, forget plot]
  table[row sep=crcr]{%
0	4395.3\\
2000	4248.5\\
4000	4118.9\\
6000	4077.1\\
8000	4065\\
10000	3998.3\\
12000	3911.7\\
14000	4108\\
24000	4015.1\\
26000	3998.6\\
30000	4052.6\\
32000	3979.3\\
34000	4046.9\\
38000	4023.3\\
40000	3920.8\\
42000	3970.7\\
44000	4064.6\\
46000	4048.8\\
48000	4129.3\\
50000	4032.8\\
52000	3950.6\\
54000	4069.3\\
56000	4022.8\\
58000	4068.4\\
60000	3971.9\\
62000	4007.4\\
64000	3948.9\\
66000	4029.2\\
68000	3968.8\\
70000	4018.8\\
72000	4113.1\\
74000	3986.5\\
76000	4028.8\\
78000	4014.5\\
80000	4092.6\\
82000	3933.4\\
84000	3934.6\\
86000	4024.3\\
88000	4016.3\\
90000	4032.7\\
92000	3993.4\\
94000	4101.3\\
96000	4061.1\\
98000	4040.8\\
1e+05	4079\\
};
\addplot [color=mycolor57, forget plot]
  table[row sep=crcr]{%
0	4395.3\\
2000	6092\\
4000	6766.6\\
6000	6594.5\\
8000	6528.9\\
10000	6700.4\\
12000	6453.8\\
14000	6639.1\\
16000	6716.8\\
18000	6644.6\\
20000	6625.8\\
22000	6672.6\\
24000	6819.8\\
26000	6768.2\\
28000	6872.6\\
30000	6704.4\\
32000	6677.9\\
34000	6701.8\\
36000	6680.9\\
38000	6583.7\\
40000	6653.2\\
42000	6801.9\\
44000	6663.5\\
48000	6771.9\\
50000	6612.8\\
52000	6743.3\\
54000	6569.7\\
56000	6772.7\\
58000	6653.8\\
60000	6766.5\\
62000	6719.3\\
64000	6809.8\\
68000	6575.5\\
70000	6826.9\\
72000	6659.2\\
74000	6639.3\\
76000	6655.4\\
78000	6572.5\\
80000	6656.1\\
82000	6567.1\\
84000	6709.5\\
86000	6780\\
88000	6759.6\\
90000	6619.6\\
92000	6723\\
94000	6638.5\\
96000	6581\\
98000	6738.4\\
1e+05	6639.2\\
};
\addplot [color=mycolor58, forget plot]
  table[row sep=crcr]{%
0	4395.3\\
2000	3950.3\\
4000	3741.6\\
6000	3624\\
8000	3580.3\\
10000	3648.7\\
12000	3572.1\\
16000	3595.3\\
18000	3586.4\\
20000	3648.6\\
22000	3591.1\\
24000	3516.8\\
26000	3557.9\\
28000	3516.4\\
30000	3648.8\\
34000	3529.9\\
36000	3658.6\\
38000	3553\\
40000	3651.3\\
42000	3643.7\\
44000	3578.1\\
46000	3618.5\\
48000	3506.6\\
50000	3537\\
52000	3588.1\\
54000	3535.3\\
56000	3512.3\\
58000	3562.5\\
60000	3440.5\\
62000	3598.4\\
64000	3568.4\\
66000	3727.6\\
68000	3709.8\\
70000	3574.4\\
72000	3574.5\\
76000	3600.8\\
78000	3563.9\\
80000	3556.2\\
82000	3645.8\\
84000	3528.4\\
88000	3588.8\\
90000	3581.1\\
92000	3547.4\\
94000	3622.2\\
98000	3563.7\\
1e+05	3549.7\\
};
\addplot [color=mycolor59, forget plot]
  table[row sep=crcr]{%
0	4395.3\\
2000	4669.1\\
4000	4898.9\\
6000	5012.6\\
8000	4974.4\\
10000	5149.3\\
12000	5012.2\\
14000	4985.1\\
18000	5128.3\\
20000	5110.2\\
22000	5028.9\\
24000	5159.8\\
26000	4958\\
28000	5125.3\\
30000	4904.2\\
32000	5049\\
34000	5095.5\\
36000	4973.2\\
38000	5170.9\\
40000	5065.8\\
44000	5064.4\\
46000	4928.8\\
48000	4955.9\\
50000	5128.7\\
52000	4961.2\\
54000	5062.5\\
56000	5061.3\\
58000	5025.3\\
60000	5123.2\\
62000	5139\\
64000	5087.3\\
66000	5052.8\\
68000	5103.9\\
70000	5115.3\\
72000	4996.1\\
74000	5067.8\\
76000	5192\\
78000	5108.3\\
80000	4987.2\\
82000	5112.9\\
84000	5014.4\\
86000	4948.9\\
88000	4999.2\\
90000	5073.2\\
92000	5033.4\\
94000	5098.2\\
96000	5111.4\\
98000	4962.5\\
1e+05	5104.2\\
};
\addplot [color=mycolor60, forget plot]
  table[row sep=crcr]{%
0	4395.3\\
2000	3868.7\\
6000	3627.8\\
10000	3563.9\\
12000	3738.3\\
14000	3642.5\\
16000	3588.1\\
18000	3487.4\\
20000	3528.2\\
22000	3615.6\\
24000	3479.5\\
26000	3574.3\\
28000	3569.3\\
30000	3668.5\\
32000	3575\\
34000	3473.9\\
36000	3543.8\\
38000	3538.4\\
40000	3652.4\\
48000	3568.9\\
50000	3513.1\\
52000	3607.8\\
54000	3513.2\\
56000	3494.5\\
58000	3607.7\\
60000	3428\\
62000	3524.5\\
64000	3495.2\\
66000	3498.1\\
68000	3564.9\\
70000	3535.5\\
72000	3524.3\\
74000	3489.2\\
76000	3516.4\\
78000	3592.3\\
80000	3546.8\\
86000	3562.3\\
88000	3446.4\\
90000	3632.2\\
92000	3568.9\\
94000	3482.5\\
96000	3571.2\\
98000	3601.5\\
1e+05	3513\\
};
\addplot [color=mycolor61, forget plot]
  table[row sep=crcr]{%
0	4395.3\\
2000	3327.2\\
4000	3152\\
6000	3111.6\\
8000	3135.1\\
10000	3127.5\\
12000	3006.4\\
14000	3126\\
16000	3137.5\\
22000	3128.8\\
26000	3113.5\\
28000	3155.6\\
30000	3131.9\\
32000	3122.4\\
36000	3181.1\\
40000	3129.3\\
42000	3203.2\\
44000	3155\\
46000	3130.3\\
48000	3141.3\\
50000	3114.5\\
54000	3149.7\\
56000	3232.5\\
58000	3180.1\\
60000	3161.8\\
62000	3106.8\\
64000	3106\\
66000	3093.7\\
68000	3145.8\\
70000	3167.2\\
72000	3132.1\\
74000	3174.7\\
76000	3135.3\\
78000	3198.4\\
80000	3093.4\\
82000	3168.7\\
84000	3123.7\\
86000	3121.8\\
88000	3130.5\\
90000	3104.2\\
92000	3140\\
94000	3104.4\\
96000	3135.5\\
98000	3138.4\\
1e+05	3156.2\\
};
\addplot [color=mycolor62, forget plot]
  table[row sep=crcr]{%
0	4395.3\\
2000	4743.2\\
4000	4965.1\\
6000	5070.6\\
8000	5081.9\\
12000	4956.1\\
14000	5095.4\\
16000	5035.6\\
18000	4886.9\\
20000	4972.7\\
22000	5028.5\\
24000	4986.5\\
26000	5096.3\\
30000	4977.1\\
32000	4900.5\\
34000	5026\\
36000	4922.1\\
38000	4977.7\\
40000	5001.5\\
42000	5048.6\\
44000	4918.9\\
46000	5022.7\\
48000	5044\\
50000	4962.4\\
52000	4996.6\\
54000	5085.5\\
56000	4912.8\\
58000	4982.4\\
60000	4928.1\\
62000	5054.4\\
64000	5026.1\\
66000	5084.7\\
68000	4942.1\\
70000	4980.6\\
72000	5069.7\\
74000	4983.3\\
76000	5032.4\\
82000	5016.1\\
84000	4932.5\\
86000	4964.6\\
88000	5019.7\\
90000	5006.5\\
92000	5018.5\\
94000	4933.4\\
96000	4973.6\\
98000	4885.4\\
1e+05	5049.7\\
};
\addplot [color=mycolor63, forget plot]
  table[row sep=crcr]{%
0	4395.3\\
2000	2983.7\\
4000	2495.2\\
6000	2419.1\\
8000	2431.3\\
10000	2404.3\\
12000	2339.1\\
14000	2436.2\\
16000	2422.5\\
18000	2440.6\\
20000	2426.1\\
22000	2447.2\\
24000	2404.5\\
26000	2369.3\\
28000	2424.8\\
30000	2465.1\\
32000	2426.4\\
34000	2395.4\\
38000	2407.3\\
40000	2445.4\\
42000	2412.6\\
44000	2425.3\\
46000	2418.3\\
50000	2442.4\\
52000	2439.4\\
54000	2471.4\\
56000	2492.1\\
58000	2417\\
62000	2393.1\\
64000	2364.7\\
66000	2415.7\\
68000	2390.2\\
70000	2430.7\\
72000	2449.5\\
74000	2451.7\\
76000	2437.4\\
78000	2439.4\\
80000	2462.6\\
82000	2373.4\\
84000	2402.4\\
86000	2474.1\\
88000	2388.8\\
90000	2451.6\\
92000	2459.3\\
96000	2434.8\\
98000	2379.2\\
1e+05	2399.4\\
};
\addplot [color=mycolor64, forget plot]
  table[row sep=crcr]{%
0	4395.3\\
2000	3204.9\\
4000	3061.8\\
6000	3056.6\\
8000	3019.1\\
10000	3072.4\\
12000	3036.2\\
14000	3083.3\\
16000	3041.7\\
18000	3040.6\\
20000	3061.5\\
22000	3117.6\\
24000	3038.9\\
26000	3009.8\\
28000	3034.8\\
30000	3007.6\\
32000	3113.8\\
34000	3044.8\\
36000	3123.7\\
38000	3067.5\\
40000	3138.1\\
42000	3066.3\\
44000	3052.8\\
46000	3068.1\\
48000	3127\\
50000	3082.9\\
52000	3112\\
54000	3068.3\\
56000	3049.6\\
58000	3053.7\\
60000	3039.2\\
62000	3046.2\\
64000	3070.1\\
66000	3016\\
68000	3033.7\\
70000	3039\\
72000	3075.3\\
74000	3049.7\\
76000	3096.8\\
80000	3016\\
82000	3058.7\\
84000	3114.6\\
88000	3057.6\\
90000	3057.5\\
92000	3089.3\\
94000	3046.1\\
96000	3073.2\\
98000	2995\\
1e+05	3062.5\\
};
\addplot [color=mycolor65, forget plot]
  table[row sep=crcr]{%
0	4395.3\\
2000	2869.6\\
4000	2560.2\\
6000	2492.9\\
8000	2508.8\\
10000	2517.8\\
12000	2556.6\\
14000	2522.1\\
16000	2530.4\\
18000	2581.3\\
20000	2608.7\\
22000	2560.6\\
24000	2578.1\\
26000	2580.6\\
28000	2555\\
30000	2538.4\\
32000	2586.2\\
34000	2555\\
36000	2543\\
40000	2556.5\\
44000	2502.5\\
46000	2601.4\\
48000	2547.7\\
50000	2594.1\\
52000	2547.4\\
54000	2555.7\\
56000	2576.9\\
58000	2542.9\\
60000	2497.5\\
62000	2565.8\\
66000	2563.6\\
68000	2552.4\\
70000	2529.4\\
72000	2569.8\\
74000	2579.3\\
76000	2518.3\\
80000	2552.7\\
82000	2589.7\\
84000	2560.8\\
86000	2509.7\\
88000	2579.6\\
90000	2544.9\\
92000	2574\\
94000	2594.8\\
96000	2534\\
98000	2586.7\\
1e+05	2565.4\\
};
\addplot [color=mycolor66, forget plot]
  table[row sep=crcr]{%
0	4395.3\\
2000	3537.5\\
4000	3407.1\\
6000	3461.5\\
8000	3478.8\\
10000	3444.4\\
12000	3445\\
14000	3398.3\\
16000	3453.6\\
18000	3435\\
20000	3401\\
22000	3456.3\\
24000	3441.7\\
26000	3404.6\\
28000	3459\\
30000	3465.5\\
32000	3413.7\\
34000	3461.7\\
36000	3438.7\\
38000	3510.6\\
42000	3464.2\\
44000	3352.6\\
46000	3392.6\\
48000	3395.9\\
50000	3378.3\\
52000	3410.1\\
54000	3377.2\\
56000	3443.6\\
58000	3467.4\\
60000	3433.5\\
62000	3378.2\\
64000	3395.1\\
66000	3457.2\\
68000	3403.3\\
70000	3412.3\\
76000	3492.5\\
78000	3457.6\\
80000	3451.6\\
82000	3384\\
86000	3395.5\\
88000	3450.4\\
90000	3411.4\\
92000	3418.8\\
94000	3439.2\\
96000	3382.5\\
98000	3456.6\\
1e+05	3472.1\\
};
\addplot [color=mycolor67, forget plot]
  table[row sep=crcr]{%
0	4395.3\\
2000	3625.2\\
4000	3227.7\\
6000	3206.4\\
8000	3203.1\\
10000	3126.9\\
12000	3205.5\\
14000	3139\\
16000	3147.1\\
18000	3232.8\\
20000	3218.4\\
22000	3189.9\\
24000	3184.7\\
26000	3213.3\\
28000	3184\\
30000	3211.6\\
32000	3187.9\\
34000	3193.3\\
36000	3208\\
38000	3235.6\\
40000	3197.8\\
42000	3174.6\\
44000	3235.4\\
46000	3232.4\\
48000	3181.5\\
50000	3196.2\\
52000	3161.8\\
54000	3290.9\\
56000	3249.9\\
58000	3237.9\\
60000	3166\\
62000	3197.5\\
64000	3190.1\\
66000	3153.3\\
68000	3225.5\\
70000	3253\\
72000	3183.8\\
74000	3160.6\\
76000	3123.3\\
78000	3182.7\\
80000	3186.2\\
86000	3155\\
88000	3232.8\\
90000	3247.5\\
92000	3209.1\\
94000	3235.1\\
96000	3164\\
98000	3187.8\\
1e+05	3190.5\\
};
\addplot [color=mycolor67]
  table[row sep=crcr]{%
0	4395.3\\
2000	2828.3\\
4000	2619.7\\
6000	2660.2\\
12000	2632.7\\
16000	2651.6\\
18000	2634.3\\
20000	2655\\
22000	2695.7\\
24000	2679.7\\
26000	2674.4\\
30000	2612.2\\
32000	2629.5\\
34000	2613.7\\
36000	2659.5\\
38000	2672.8\\
40000	2646.3\\
42000	2658.4\\
44000	2657.8\\
46000	2686.9\\
50000	2609.8\\
52000	2610\\
54000	2680.3\\
56000	2666.3\\
58000	2640\\
60000	2599.8\\
62000	2651.3\\
64000	2629\\
66000	2658.6\\
68000	2650.6\\
72000	2619.5\\
74000	2631.2\\
76000	2624.1\\
78000	2663.3\\
80000	2680.8\\
82000	2635\\
84000	2625.4\\
86000	2679.3\\
88000	2644.2\\
90000	2576.2\\
92000	2627.6\\
96000	2634.5\\
98000	2636.5\\
1e+05	2626.6\\
};
\addlegendentry{\ClasskIIILegend}

\addplot [color=mycolor68, forget plot]
  table[row sep=crcr]{%
0	4395.3\\
2000	2924\\
4000	2605.9\\
6000	2498.9\\
8000	2506.2\\
10000	2476.6\\
12000	2482.2\\
14000	2508\\
16000	2492.1\\
20000	2481.9\\
22000	2451.6\\
24000	2476\\
26000	2509.2\\
28000	2466.9\\
30000	2530\\
32000	2480.6\\
36000	2486.7\\
38000	2468.4\\
40000	2494.2\\
44000	2487.4\\
46000	2473.3\\
48000	2477\\
50000	2514.7\\
54000	2460.7\\
56000	2464.7\\
58000	2483.2\\
60000	2489.5\\
64000	2521.4\\
66000	2499\\
68000	2499.5\\
70000	2464.7\\
72000	2498.9\\
74000	2467.7\\
76000	2481.6\\
78000	2447.4\\
80000	2496.5\\
82000	2493.4\\
84000	2481.6\\
86000	2493.7\\
88000	2491.9\\
90000	2467.2\\
92000	2457\\
94000	2483.1\\
96000	2496.4\\
98000	2476.5\\
1e+05	2496.9\\
};
\addplot [color=mycolor69, forget plot]
  table[row sep=crcr]{%
0	4395.3\\
2000	3134.6\\
4000	2859.6\\
6000	2865.1\\
8000	2935.4\\
12000	2898.4\\
14000	2936.2\\
16000	2920.6\\
18000	2883.6\\
20000	2927.9\\
22000	2928.6\\
24000	2919.4\\
26000	2852.6\\
28000	2908.2\\
30000	2951.8\\
32000	2914.8\\
36000	2886.4\\
38000	2939.9\\
40000	2901.6\\
42000	2929.5\\
44000	2930.6\\
46000	2916.7\\
48000	2932.2\\
50000	2906.2\\
52000	2919.9\\
54000	2886.3\\
56000	2877.7\\
58000	2909.1\\
60000	2895.7\\
62000	2922.7\\
64000	2936.2\\
66000	2968.4\\
68000	2934.5\\
70000	2981\\
76000	2918.3\\
78000	2931.8\\
80000	2898.4\\
82000	2930.3\\
84000	2925.1\\
86000	2931.3\\
88000	2916.8\\
92000	2925.9\\
94000	2946.7\\
96000	2938.3\\
98000	2897.7\\
1e+05	2926\\
};
\addplot [color=mycolor70, forget plot]
  table[row sep=crcr]{%
0	4395.3\\
2000	2394.4\\
4000	2082.5\\
6000	2113.4\\
8000	2118.7\\
10000	2133.7\\
12000	2114.1\\
16000	2122.6\\
18000	2140.8\\
20000	2106.2\\
22000	2125.2\\
24000	2136.2\\
26000	2072.4\\
28000	2141.6\\
30000	2128.9\\
32000	2100.7\\
34000	2097.7\\
36000	2116.8\\
38000	2096.8\\
40000	2124.5\\
42000	2118\\
44000	2151.9\\
46000	2110.1\\
48000	2117.1\\
50000	2136.8\\
52000	2109.1\\
54000	2130\\
56000	2095.6\\
58000	2129.3\\
60000	2126.8\\
62000	2149.4\\
64000	2157.3\\
66000	2130.4\\
68000	2136.6\\
70000	2166.3\\
72000	2119\\
74000	2116\\
76000	2101.1\\
82000	2155\\
84000	2093.5\\
90000	2107.7\\
92000	2071.2\\
94000	2115.9\\
96000	2140.9\\
98000	2152.1\\
1e+05	2184.1\\
};
\addplot [color=mycolor71, forget plot]
  table[row sep=crcr]{%
0	4395.3\\
2000	2896.1\\
4000	2609\\
6000	2544.9\\
8000	2510.6\\
10000	2520\\
12000	2486.7\\
14000	2504.7\\
16000	2498.2\\
18000	2509.8\\
20000	2478.7\\
24000	2515.7\\
26000	2487.2\\
28000	2483.1\\
30000	2512.3\\
34000	2549.4\\
36000	2498.2\\
38000	2496.3\\
40000	2477.9\\
42000	2498.6\\
46000	2518.1\\
48000	2534.9\\
50000	2490.1\\
54000	2524.3\\
56000	2504.1\\
58000	2519.9\\
60000	2504.8\\
62000	2500.2\\
64000	2465.7\\
66000	2473.1\\
68000	2495.9\\
70000	2503.4\\
72000	2469.5\\
74000	2487.4\\
78000	2495.9\\
80000	2493.4\\
82000	2468.7\\
84000	2516\\
86000	2469.9\\
88000	2514.5\\
94000	2472\\
96000	2482.2\\
1e+05	2513.5\\
};
\addplot [color=mycolor72, forget plot]
  table[row sep=crcr]{%
0	4395.3\\
2000	2704.6\\
4000	2182.9\\
6000	2239.6\\
8000	2197.4\\
10000	2220.1\\
12000	2223.6\\
14000	2217.8\\
18000	2249.1\\
20000	2201.8\\
22000	2202.3\\
26000	2234.7\\
28000	2219.7\\
30000	2196\\
32000	2209.6\\
34000	2177.1\\
36000	2226.7\\
38000	2192.8\\
40000	2232.6\\
42000	2184.6\\
44000	2177.4\\
46000	2230.1\\
48000	2186.4\\
50000	2211.4\\
52000	2209.7\\
54000	2158.2\\
56000	2244.9\\
58000	2212.2\\
60000	2230.4\\
62000	2242.1\\
64000	2196.5\\
66000	2192.7\\
68000	2229.8\\
72000	2183.8\\
76000	2240.8\\
78000	2201.4\\
80000	2209\\
82000	2157.1\\
84000	2212.7\\
86000	2211.1\\
88000	2258\\
90000	2240.3\\
92000	2215.6\\
94000	2163.3\\
96000	2276.2\\
98000	2231.4\\
1e+05	2203.4\\
};
\addplot [color=mycolor73, forget plot]
  table[row sep=crcr]{%
0	4395.3\\
2000	2962\\
4000	2679.2\\
6000	2642.2\\
8000	2617.8\\
10000	2631.4\\
12000	2657.9\\
14000	2601.5\\
16000	2604.2\\
18000	2617.9\\
20000	2662.2\\
24000	2626.5\\
26000	2629.7\\
28000	2692.6\\
30000	2657.6\\
32000	2590.3\\
34000	2626.2\\
38000	2638.6\\
40000	2600.2\\
42000	2619.3\\
44000	2606.2\\
46000	2620.2\\
50000	2662.8\\
52000	2682.2\\
54000	2660.3\\
56000	2583.8\\
58000	2573.7\\
60000	2633.8\\
62000	2603.9\\
64000	2595.8\\
66000	2661.9\\
68000	2601.3\\
70000	2593.8\\
74000	2609.8\\
76000	2611.3\\
78000	2632\\
80000	2630.5\\
82000	2616.5\\
84000	2667.2\\
86000	2628.3\\
88000	2656.3\\
90000	2628.8\\
92000	2624\\
94000	2607.3\\
96000	2605.4\\
1e+05	2615.9\\
};
\addplot [color=mycolor74, forget plot]
  table[row sep=crcr]{%
0	4395.3\\
2000	2198.3\\
4000	1590.4\\
6000	1573.1\\
8000	1586.4\\
10000	1562.2\\
12000	1550.8\\
14000	1526.3\\
16000	1544.1\\
18000	1550.3\\
20000	1567.6\\
22000	1571.9\\
24000	1530.2\\
26000	1559.9\\
28000	1550.6\\
30000	1564.9\\
32000	1543.7\\
36000	1548.3\\
38000	1574\\
40000	1572.9\\
42000	1566\\
44000	1567.4\\
46000	1549.2\\
48000	1554.7\\
50000	1567.8\\
52000	1548.1\\
54000	1549.5\\
56000	1557.6\\
58000	1525.6\\
60000	1536.6\\
62000	1541.5\\
64000	1562.4\\
66000	1561.2\\
68000	1534.4\\
70000	1537.7\\
72000	1564.5\\
74000	1556.7\\
78000	1556.1\\
80000	1578.6\\
82000	1548.5\\
84000	1545.2\\
86000	1529.1\\
88000	1565.5\\
90000	1549.6\\
92000	1544.3\\
94000	1555.7\\
96000	1557.7\\
98000	1526.3\\
1e+05	1552.8\\
};
\addplot [color=mycolor75, forget plot]
  table[row sep=crcr]{%
0	4395.3\\
2000	2726.7\\
4000	2276.9\\
6000	2244.6\\
8000	2280.4\\
10000	2210.4\\
12000	2256.1\\
14000	2226.6\\
16000	2228.9\\
18000	2197.2\\
20000	2216.4\\
22000	2223.8\\
24000	2238.9\\
26000	2219.1\\
28000	2260.3\\
30000	2209.8\\
32000	2227.3\\
34000	2210.4\\
38000	2208.3\\
40000	2241.4\\
42000	2220.6\\
44000	2212.3\\
46000	2228.4\\
48000	2205.2\\
50000	2200\\
54000	2242.3\\
56000	2212\\
58000	2238.7\\
62000	2210.8\\
64000	2166.7\\
66000	2223.3\\
68000	2246.4\\
70000	2210.2\\
72000	2198.7\\
74000	2227.2\\
76000	2196.3\\
78000	2238.8\\
80000	2226.3\\
82000	2227.5\\
84000	2220.3\\
86000	2204\\
88000	2209.2\\
90000	2198.8\\
92000	2223.7\\
94000	2274\\
96000	2199.6\\
98000	2235.3\\
1e+05	2226.9\\
};
\addplot [color=mycolor76, forget plot]
  table[row sep=crcr]{%
0	4395.3\\
2000	2567.6\\
4000	2096.5\\
8000	2032.5\\
10000	2039.8\\
12000	2018.6\\
20000	2029.4\\
22000	2042.2\\
24000	2009.7\\
28000	2031.7\\
32000	2005.7\\
34000	2036.6\\
36000	2006.9\\
40000	2022.9\\
42000	1999.2\\
46000	2043.7\\
50000	2031.2\\
54000	2058.6\\
56000	2024.8\\
58000	2009.5\\
60000	2031.7\\
62000	2012\\
64000	1998.9\\
66000	2027.2\\
70000	2025.7\\
72000	1999.7\\
74000	2020.5\\
76000	2032.1\\
78000	2026.4\\
80000	2002.9\\
82000	2038.9\\
84000	2020.7\\
86000	2051.6\\
88000	2049.4\\
90000	2014\\
92000	2009.4\\
94000	2040.8\\
96000	2015.5\\
98000	2022.8\\
1e+05	2064.2\\
};
\addplot [color=mycolor77, forget plot]
  table[row sep=crcr]{%
0	4395.3\\
2000	2250.8\\
4000	1675.1\\
6000	1624.4\\
10000	1574.5\\
12000	1579\\
14000	1572.1\\
16000	1575.3\\
18000	1591.8\\
20000	1560.7\\
24000	1580.8\\
26000	1566.6\\
28000	1596.4\\
30000	1563.3\\
32000	1557.1\\
34000	1580.7\\
36000	1571.4\\
38000	1570.9\\
40000	1579.8\\
42000	1569.8\\
44000	1593.2\\
46000	1553.4\\
48000	1592.5\\
50000	1565.3\\
54000	1573.9\\
56000	1587.2\\
58000	1556.4\\
60000	1570.5\\
62000	1561.6\\
64000	1607.4\\
66000	1577.9\\
68000	1590.1\\
70000	1583.3\\
72000	1567.3\\
74000	1580.3\\
76000	1568.4\\
78000	1568.3\\
80000	1573.3\\
82000	1566.6\\
84000	1570.2\\
86000	1579.8\\
88000	1567.3\\
90000	1563.5\\
92000	1569.2\\
94000	1591.2\\
96000	1596.5\\
98000	1571.6\\
1e+05	1569.7\\
};
\addplot [color=mycolor78, forget plot]
  table[row sep=crcr]{%
0	4395.3\\
2000	2840.2\\
4000	2385.9\\
6000	2343.7\\
8000	2319.4\\
10000	2339.9\\
12000	2318\\
14000	2311.5\\
18000	2274.2\\
20000	2319.8\\
22000	2271.9\\
24000	2316.3\\
26000	2337.6\\
28000	2348.1\\
30000	2293.4\\
32000	2326.1\\
34000	2335.6\\
36000	2308.4\\
38000	2340.9\\
40000	2325.2\\
42000	2270.3\\
44000	2277.5\\
46000	2341.9\\
48000	2340.1\\
50000	2287.8\\
52000	2300.5\\
56000	2261.1\\
58000	2319.2\\
60000	2310.5\\
62000	2320.5\\
64000	2257.7\\
66000	2278.1\\
68000	2318.7\\
70000	2307\\
72000	2319.9\\
74000	2317.2\\
76000	2343.8\\
78000	2296.1\\
80000	2259.6\\
82000	2273\\
84000	2279.3\\
86000	2277.5\\
88000	2322.4\\
90000	2304.8\\
92000	2277.4\\
94000	2289.3\\
96000	2319.6\\
98000	2315.2\\
1e+05	2317.9\\
};
\addplot [color=mycolor79, forget plot]
  table[row sep=crcr]{%
0	4395.3\\
2000	2689.7\\
4000	2316.9\\
6000	2225.1\\
10000	2218.4\\
12000	2195.6\\
14000	2233.2\\
16000	2204.4\\
18000	2212.2\\
22000	2149.7\\
24000	2190.8\\
26000	2194.1\\
28000	2172.6\\
30000	2158\\
32000	2116.9\\
34000	2172\\
36000	2197.9\\
42000	2150.4\\
44000	2192.3\\
46000	2156.7\\
48000	2162.1\\
50000	2149.5\\
52000	2168.6\\
54000	2159.7\\
56000	2166.2\\
58000	2124.8\\
60000	2170.4\\
62000	2146\\
66000	2217.6\\
68000	2154.6\\
70000	2181.3\\
72000	2161.4\\
74000	2160.8\\
76000	2138.2\\
78000	2166.2\\
80000	2168.9\\
82000	2129.1\\
84000	2173.2\\
86000	2156.9\\
88000	2171.9\\
90000	2208.8\\
92000	2221.3\\
94000	2178.3\\
96000	2198\\
98000	2207.1\\
1e+05	2207.2\\
};
\addplot [color=mycolor79, forget plot]
  table[row sep=crcr]{%
0	4395.3\\
2000	2672.3\\
4000	2067.4\\
6000	1994.7\\
8000	2014.2\\
10000	1990.4\\
12000	2023\\
14000	2040.8\\
16000	1997.8\\
18000	2028.8\\
20000	2070\\
22000	2026.1\\
24000	2034.5\\
26000	2013.1\\
28000	2060.6\\
30000	2007.6\\
32000	2010.7\\
34000	2040.3\\
36000	2008.6\\
38000	2051.8\\
40000	2062.5\\
42000	2048.6\\
44000	2078\\
46000	2010.6\\
48000	2013.1\\
50000	2005.7\\
52000	2026.2\\
54000	2022.6\\
56000	2031\\
58000	2007.6\\
60000	2029.2\\
62000	2008.4\\
64000	2050.1\\
66000	2050.3\\
68000	2027.5\\
70000	1994.3\\
72000	1993.3\\
74000	2007.5\\
76000	1983.8\\
78000	1992.2\\
82000	2045.7\\
84000	2012.4\\
86000	2057.9\\
88000	2044.9\\
90000	2064.5\\
92000	2017.9\\
94000	2025.9\\
96000	2052.2\\
98000	1982\\
1e+05	1989.9\\
};
\addplot [color=mycolor80, forget plot]
  table[row sep=crcr]{%
0	4395.3\\
2000	2172.5\\
4000	1637.8\\
6000	1585.9\\
8000	1577.5\\
12000	1587.6\\
14000	1602.6\\
16000	1568.8\\
18000	1553.3\\
20000	1567.7\\
22000	1576.3\\
24000	1553.6\\
26000	1585.6\\
28000	1561.7\\
30000	1565\\
32000	1576.3\\
34000	1578\\
36000	1591\\
38000	1555.7\\
40000	1560.9\\
44000	1549.1\\
46000	1585.7\\
50000	1554.7\\
52000	1564.6\\
56000	1597.5\\
58000	1561.2\\
62000	1578.8\\
64000	1571.3\\
66000	1571.9\\
68000	1588.9\\
72000	1557.1\\
74000	1557.4\\
76000	1570.3\\
78000	1551.7\\
80000	1579.1\\
82000	1566.7\\
84000	1571.4\\
86000	1563.3\\
88000	1572.5\\
90000	1574.5\\
92000	1557.8\\
96000	1568.8\\
98000	1547.4\\
1e+05	1557\\
};
\addplot [color=mycolor81, forget plot]
  table[row sep=crcr]{%
0	4395.3\\
2000	2130.1\\
4000	1490.3\\
6000	1456.5\\
8000	1461.4\\
10000	1475.9\\
12000	1453.9\\
14000	1450.9\\
16000	1452.9\\
18000	1444.6\\
20000	1475.8\\
22000	1468.2\\
24000	1484.7\\
28000	1453.5\\
30000	1453.4\\
32000	1441.5\\
34000	1452.3\\
36000	1484.2\\
38000	1469.1\\
40000	1467.7\\
44000	1448.4\\
48000	1474.9\\
50000	1492.6\\
52000	1490.7\\
54000	1484.5\\
58000	1480.5\\
60000	1488.7\\
62000	1455.7\\
66000	1483.9\\
68000	1511.5\\
70000	1457.8\\
72000	1487.4\\
76000	1458.4\\
80000	1462.4\\
84000	1469.4\\
86000	1454.3\\
88000	1467.9\\
90000	1453.3\\
92000	1481\\
94000	1459.6\\
96000	1429.5\\
98000	1464.7\\
1e+05	1477.2\\
};
\addplot [color=mycolor82, forget plot]
  table[row sep=crcr]{%
0	4395.3\\
2000	2370.7\\
4000	1818.8\\
6000	1734.9\\
10000	1711.7\\
12000	1717.8\\
14000	1712.4\\
16000	1697\\
18000	1716.1\\
20000	1716.5\\
22000	1688\\
24000	1715.8\\
26000	1701.4\\
28000	1707.6\\
30000	1725.3\\
32000	1703.9\\
34000	1706.6\\
38000	1723.4\\
40000	1670.2\\
42000	1714.5\\
46000	1747.4\\
48000	1705.8\\
50000	1699.8\\
52000	1713.6\\
54000	1687.9\\
56000	1695.3\\
58000	1717.8\\
64000	1717.3\\
66000	1727.7\\
68000	1705.6\\
72000	1677.9\\
74000	1695.9\\
76000	1724\\
78000	1724.2\\
80000	1706.7\\
84000	1681.4\\
86000	1689.6\\
88000	1709.1\\
90000	1701.4\\
92000	1706.6\\
94000	1717.2\\
96000	1710\\
98000	1710.3\\
1e+05	1743.1\\
};
\addplot [color=mycolor83, forget plot]
  table[row sep=crcr]{%
0	4395.3\\
2000	2453.2\\
4000	1858\\
6000	1809.4\\
8000	1773.5\\
10000	1746.8\\
16000	1739.3\\
18000	1746.5\\
20000	1736.1\\
26000	1757.2\\
28000	1753.2\\
30000	1769.4\\
32000	1731.2\\
36000	1757.3\\
38000	1709.9\\
40000	1717.9\\
42000	1718.4\\
44000	1750.4\\
48000	1731.7\\
52000	1732\\
54000	1734\\
56000	1718.3\\
58000	1737.5\\
60000	1712.2\\
62000	1734.2\\
64000	1727.1\\
66000	1748.2\\
68000	1749.8\\
72000	1727.1\\
78000	1730.8\\
80000	1730\\
82000	1741.3\\
86000	1736\\
88000	1745.7\\
90000	1770\\
92000	1764\\
96000	1727.4\\
98000	1742.4\\
1e+05	1724.1\\
};
\addplot [color=mycolor83, forget plot]
  table[row sep=crcr]{%
0	4395.3\\
2000	1920.4\\
4000	1289.2\\
6000	1292.1\\
8000	1275.5\\
10000	1272.2\\
12000	1262.7\\
14000	1271.8\\
16000	1273\\
18000	1283\\
20000	1264.3\\
22000	1280.3\\
24000	1253.7\\
26000	1270.9\\
28000	1272.2\\
30000	1264.4\\
32000	1260.7\\
34000	1266.7\\
36000	1303.4\\
42000	1268.7\\
44000	1284.7\\
46000	1276.3\\
48000	1281.7\\
50000	1259.9\\
52000	1252\\
54000	1272.8\\
56000	1284.6\\
58000	1279.9\\
60000	1287.4\\
62000	1256.8\\
64000	1281.2\\
66000	1331.8\\
68000	1281.7\\
70000	1270.2\\
72000	1264.7\\
76000	1286.2\\
78000	1275.8\\
80000	1284.2\\
82000	1285.2\\
84000	1250.8\\
86000	1273.9\\
88000	1306.7\\
90000	1294.4\\
92000	1304\\
94000	1290.7\\
96000	1269.9\\
98000	1269.9\\
1e+05	1276.4\\
};
\addplot [color=mycolor84, forget plot]
  table[row sep=crcr]{%
0	4395.3\\
2000	1990\\
4000	1371.1\\
6000	1270.2\\
10000	1292\\
12000	1284.3\\
14000	1291.5\\
16000	1291.6\\
18000	1300.1\\
20000	1267.5\\
22000	1242.9\\
24000	1298.4\\
26000	1264.5\\
28000	1266.9\\
30000	1282\\
34000	1295.6\\
36000	1322\\
38000	1283.7\\
40000	1302.8\\
42000	1293\\
44000	1258.9\\
46000	1282.6\\
48000	1261.4\\
50000	1297.9\\
52000	1314.9\\
54000	1267.7\\
56000	1271\\
58000	1252.6\\
60000	1285.9\\
62000	1284.4\\
64000	1278.4\\
66000	1279.3\\
68000	1269.7\\
70000	1283.6\\
72000	1263.1\\
74000	1302.1\\
76000	1322.6\\
78000	1283.5\\
80000	1273.3\\
82000	1279.3\\
84000	1308.4\\
86000	1290.6\\
88000	1246.4\\
90000	1295.1\\
92000	1275.7\\
94000	1312.9\\
96000	1302.4\\
98000	1306.5\\
1e+05	1286.3\\
};
\addplot [color=mycolor85, forget plot]
  table[row sep=crcr]{%
0	4395.3\\
2000	2212.2\\
4000	1490.5\\
6000	1444.3\\
8000	1455.3\\
10000	1405.5\\
12000	1435.2\\
14000	1422.4\\
16000	1394.6\\
18000	1402.2\\
20000	1457.4\\
22000	1407.4\\
26000	1362.3\\
28000	1414.7\\
30000	1384.8\\
32000	1430.2\\
34000	1400.6\\
38000	1422.6\\
40000	1388.8\\
42000	1403.6\\
44000	1438\\
48000	1385.1\\
50000	1439.2\\
52000	1404.9\\
54000	1412.3\\
56000	1428.6\\
60000	1405.9\\
62000	1408\\
64000	1439.5\\
66000	1414.8\\
68000	1417.1\\
70000	1398.6\\
72000	1445.8\\
74000	1436.2\\
76000	1433.8\\
78000	1424\\
80000	1407.4\\
82000	1438.5\\
84000	1410.7\\
88000	1388.6\\
92000	1413\\
94000	1404.9\\
96000	1430.8\\
98000	1379.5\\
1e+05	1414\\
};
\addplot [color=mycolor86, forget plot]
  table[row sep=crcr]{%
0	4395.3\\
2000	1980.1\\
4000	1262.3\\
6000	1200.4\\
8000	1228.3\\
10000	1180.1\\
12000	1177.2\\
16000	1188.2\\
18000	1202.6\\
20000	1198.1\\
22000	1202.3\\
24000	1215.2\\
26000	1191.7\\
28000	1175.3\\
30000	1175\\
32000	1204.7\\
34000	1179.9\\
36000	1185\\
38000	1174.9\\
40000	1203.7\\
42000	1243.5\\
44000	1202.2\\
46000	1198.4\\
48000	1187\\
50000	1186.1\\
52000	1189.8\\
54000	1207.5\\
56000	1205.3\\
58000	1182\\
60000	1178.9\\
62000	1191.3\\
64000	1196.1\\
66000	1181.3\\
68000	1186.7\\
70000	1181.6\\
72000	1184.2\\
74000	1223.5\\
76000	1225.2\\
78000	1199.1\\
80000	1154\\
82000	1184\\
84000	1190.1\\
86000	1186.6\\
88000	1206.5\\
90000	1185.8\\
92000	1171.4\\
94000	1177.2\\
96000	1176.3\\
98000	1214\\
1e+05	1181.2\\
};
\addplot [color=mycolor87, forget plot]
  table[row sep=crcr]{%
0	4395.3\\
2000	1808.9\\
4000	996.57\\
6000	1026.8\\
8000	1027\\
10000	960.79\\
12000	986.53\\
14000	1029.1\\
16000	1043.8\\
18000	1005.2\\
20000	991.6\\
22000	1032.6\\
26000	983.22\\
30000	992.74\\
32000	972.98\\
34000	978.71\\
36000	994.95\\
38000	986.05\\
40000	990.8\\
42000	1001\\
44000	1005.2\\
46000	1023.2\\
48000	1028\\
50000	1019.8\\
54000	1011.8\\
56000	1002.8\\
58000	1049\\
60000	1022.6\\
64000	1003.7\\
66000	1058.3\\
68000	1036.2\\
70000	1030.9\\
72000	1042.1\\
74000	1006.2\\
78000	980.49\\
80000	1035.1\\
82000	1005.9\\
84000	1012.3\\
86000	1026.4\\
88000	997.64\\
90000	1020.4\\
92000	1002.9\\
94000	1006.4\\
96000	981.8\\
98000	1011.1\\
1e+05	960.27\\
};
\addplot [color=mycolor88, forget plot]
  table[row sep=crcr]{%
0	4395.3\\
2000	2033.5\\
4000	1269\\
6000	1200.5\\
8000	1212\\
10000	1188.2\\
12000	1186.9\\
14000	1201.8\\
16000	1190.3\\
18000	1189.4\\
20000	1160.5\\
24000	1190.3\\
26000	1176.7\\
28000	1176.5\\
30000	1204.5\\
32000	1203.9\\
34000	1182.8\\
36000	1205\\
38000	1208.1\\
40000	1188.2\\
42000	1174\\
44000	1198.9\\
46000	1172.2\\
48000	1153.7\\
50000	1169.4\\
52000	1196\\
54000	1187.7\\
56000	1183.9\\
58000	1172.1\\
60000	1201.5\\
62000	1193.7\\
64000	1180.2\\
66000	1195.1\\
68000	1174.7\\
70000	1173.7\\
72000	1156.6\\
74000	1149.7\\
78000	1211.5\\
80000	1167.1\\
82000	1171.9\\
84000	1194.2\\
86000	1183.3\\
88000	1176.2\\
90000	1201.9\\
92000	1238.2\\
94000	1168.2\\
96000	1197.7\\
98000	1219.4\\
1e+05	1193.3\\
};
\addplot [color=mycolor89, forget plot]
  table[row sep=crcr]{%
0	4395.3\\
2000	1940.9\\
4000	948.01\\
6000	855.39\\
8000	871.42\\
10000	846.68\\
12000	855.41\\
14000	847.31\\
16000	863.03\\
18000	894.33\\
20000	877.19\\
22000	856.51\\
24000	868.24\\
28000	862.99\\
30000	848.05\\
32000	855.28\\
34000	856.2\\
36000	851.43\\
38000	834.69\\
40000	835.11\\
42000	856.88\\
44000	844.13\\
48000	877.64\\
50000	862.84\\
52000	839.15\\
54000	856.15\\
56000	858.39\\
58000	835.46\\
60000	875.68\\
62000	849.53\\
64000	867.38\\
66000	851.13\\
68000	848.83\\
72000	855.62\\
74000	872.27\\
76000	867.41\\
78000	850.02\\
80000	857.22\\
82000	856.36\\
84000	840.75\\
86000	857.07\\
88000	862.3\\
90000	871.11\\
92000	862.31\\
94000	862.18\\
96000	857\\
1e+05	856.63\\
};
\addplot [color=mycolor90, forget plot]
  table[row sep=crcr]{%
0	4395.3\\
2000	1686.3\\
4000	858.74\\
8000	830.3\\
10000	836.2\\
12000	823.74\\
14000	831.61\\
16000	844.62\\
18000	836.81\\
20000	818.04\\
22000	829.01\\
24000	822.85\\
26000	828.09\\
28000	820.06\\
30000	835.9\\
32000	842.94\\
34000	835.12\\
36000	840.15\\
38000	836.32\\
40000	822.66\\
42000	823.76\\
44000	835.9\\
46000	830.45\\
48000	832.14\\
50000	839.5\\
52000	839.82\\
56000	823.64\\
58000	852.2\\
60000	827.95\\
62000	825.48\\
64000	839.7\\
66000	831.32\\
68000	832.75\\
70000	808.45\\
72000	842.31\\
74000	825.2\\
76000	826.49\\
78000	806.33\\
80000	819.04\\
82000	837.64\\
84000	837.37\\
86000	847.08\\
90000	824.02\\
92000	821.84\\
94000	839.58\\
96000	808.89\\
98000	813.87\\
1e+05	829.33\\
};
\addplot [color=mycolor91, forget plot]
  table[row sep=crcr]{%
0	4395.3\\
2000	1863.3\\
4000	811.12\\
6000	720.23\\
8000	712.4\\
10000	717.99\\
12000	705.68\\
14000	735.49\\
16000	726.02\\
18000	751.62\\
20000	753.25\\
22000	745.99\\
24000	724.42\\
26000	732.48\\
28000	721.21\\
30000	716.77\\
32000	717.5\\
34000	737.5\\
36000	702.65\\
38000	717.53\\
40000	723.53\\
42000	688.36\\
44000	699.97\\
46000	753.63\\
48000	735.15\\
50000	712.58\\
52000	725.21\\
56000	696.95\\
60000	724.66\\
62000	702.93\\
64000	732.72\\
66000	744.26\\
68000	736.6\\
70000	715.89\\
72000	729.59\\
74000	710.88\\
76000	709.22\\
78000	727.73\\
80000	706.52\\
82000	734.53\\
84000	715.09\\
86000	740.19\\
88000	740.99\\
90000	723.14\\
92000	721.3\\
94000	691.27\\
96000	699.27\\
98000	711.13\\
1e+05	720.57\\
};
\addplot [color=mycolor92, forget plot]
  table[row sep=crcr]{%
0	4395.3\\
2000	1880.9\\
4000	1004.3\\
6000	919.7\\
8000	910.02\\
10000	911.14\\
12000	876.57\\
14000	897\\
16000	888.46\\
18000	931.41\\
20000	926.47\\
22000	905.48\\
24000	882.09\\
26000	889.99\\
28000	892.21\\
30000	916.01\\
32000	888.99\\
34000	926.91\\
36000	898.07\\
38000	898.62\\
40000	880.72\\
42000	907.77\\
44000	900.75\\
46000	888.87\\
48000	895.21\\
50000	905.63\\
52000	902.29\\
54000	888.54\\
56000	919.63\\
58000	885.24\\
60000	904.17\\
64000	912.88\\
66000	913.68\\
68000	883.97\\
70000	908.53\\
72000	905.59\\
74000	895.21\\
76000	923.77\\
78000	905.08\\
80000	902.31\\
82000	892.61\\
84000	900.37\\
86000	905.31\\
88000	907.26\\
90000	902.62\\
92000	911.56\\
94000	917.25\\
98000	871.29\\
1e+05	904.89\\
};
\addplot [color=mycolor93, forget plot]
  table[row sep=crcr]{%
0	4395.3\\
2000	1995.2\\
4000	887.41\\
6000	783.05\\
8000	765.48\\
10000	756.07\\
12000	770.36\\
14000	798.91\\
18000	757.01\\
20000	747.05\\
22000	748.53\\
24000	742.21\\
26000	751.73\\
28000	751.33\\
30000	753.77\\
32000	749.62\\
34000	750.13\\
36000	771.32\\
38000	762.69\\
40000	726.4\\
42000	761.8\\
44000	774.36\\
46000	756.42\\
48000	767.45\\
50000	744.65\\
52000	746.45\\
54000	744.84\\
56000	770.42\\
58000	728.58\\
60000	774.08\\
62000	764.83\\
64000	768\\
66000	780.78\\
68000	735.89\\
70000	783.81\\
72000	773.71\\
74000	770.47\\
76000	733.06\\
78000	737.75\\
80000	764.64\\
82000	711.1\\
84000	735.85\\
86000	776.17\\
88000	748.7\\
90000	743.54\\
92000	761.04\\
94000	745.85\\
96000	749.31\\
98000	748.22\\
1e+05	767.2\\
};
\addplot [color=mycolor94, forget plot]
  table[row sep=crcr]{%
0	4395.3\\
2000	1955.8\\
4000	842.05\\
6000	751.87\\
8000	745.71\\
10000	733.81\\
14000	718.91\\
16000	708.29\\
18000	729.12\\
20000	733.88\\
22000	720.74\\
24000	701.82\\
26000	728.79\\
28000	727.51\\
30000	706.63\\
32000	697.34\\
34000	697.58\\
36000	725.32\\
38000	726.24\\
40000	717\\
44000	723.05\\
46000	719.4\\
48000	731.77\\
50000	721.95\\
52000	724.23\\
54000	716.79\\
56000	725.58\\
58000	713.6\\
62000	708.19\\
64000	710\\
66000	695.14\\
68000	721.39\\
70000	715.65\\
72000	738.39\\
74000	734.6\\
76000	708.64\\
78000	710.34\\
80000	702.73\\
82000	703.4\\
84000	722.26\\
86000	705.01\\
88000	729.23\\
90000	679.9\\
92000	688.26\\
94000	692.06\\
96000	726.28\\
98000	713.4\\
1e+05	709.72\\
};
\addplot [color=mycolor95, forget plot]
  table[row sep=crcr]{%
0	4395.3\\
2000	1552.9\\
4000	516.54\\
6000	439.22\\
8000	431.49\\
10000	426.97\\
12000	431.52\\
14000	429.19\\
16000	424.92\\
18000	423.06\\
20000	428.05\\
22000	436.24\\
24000	425.04\\
26000	430.12\\
28000	452.52\\
32000	422.81\\
34000	414.58\\
36000	409.76\\
38000	422.94\\
40000	442.73\\
42000	410.3\\
44000	437\\
46000	414.51\\
48000	407.56\\
50000	441.58\\
52000	436.73\\
54000	422.2\\
56000	423.29\\
58000	425.8\\
60000	432.57\\
62000	423.22\\
64000	412.21\\
66000	436.18\\
68000	437.66\\
70000	426.87\\
72000	423.32\\
74000	422.84\\
76000	434.06\\
78000	419.74\\
80000	446.39\\
82000	422.24\\
84000	426.05\\
86000	437.46\\
88000	438.06\\
90000	431.87\\
92000	441.88\\
94000	424.55\\
96000	398.06\\
98000	418.25\\
1e+05	424.99\\
};
\addplot [color=mycolor96, forget plot]
  table[row sep=crcr]{%
0	4395.3\\
2000	1757.6\\
4000	588.79\\
6000	478.22\\
8000	457.06\\
10000	463.44\\
12000	477.42\\
14000	446.87\\
16000	453.12\\
18000	465.24\\
20000	444.77\\
22000	444.77\\
24000	455.92\\
26000	442.6\\
28000	431.88\\
30000	426.76\\
32000	452.59\\
34000	460.03\\
36000	441.05\\
38000	444.51\\
40000	469.85\\
42000	429.53\\
44000	464.8\\
46000	447.11\\
48000	442.68\\
50000	435.41\\
52000	467.61\\
54000	462.07\\
56000	454.04\\
58000	420.67\\
60000	445.82\\
62000	450.53\\
64000	485.22\\
66000	456.07\\
68000	449.39\\
70000	458.78\\
72000	434.52\\
74000	470.72\\
76000	466.84\\
78000	437.88\\
80000	451.33\\
82000	470.49\\
84000	455.45\\
86000	456.75\\
88000	441.92\\
90000	460.36\\
92000	467.64\\
94000	448.1\\
96000	421.58\\
98000	456.28\\
1e+05	463.68\\
};
\addplot [color=mycolor97, forget plot]
  table[row sep=crcr]{%
0	4395.3\\
2000	1521.5\\
4000	434.76\\
6000	341.12\\
8000	320.9\\
10000	335.5\\
12000	334.1\\
14000	313.01\\
16000	318.78\\
18000	327.23\\
20000	317.4\\
22000	313.46\\
24000	328.24\\
26000	322.21\\
28000	323.91\\
30000	321.42\\
32000	329.1\\
34000	315.6\\
36000	298.29\\
38000	319.29\\
40000	308.14\\
42000	317.96\\
44000	324.33\\
46000	322.87\\
48000	348.7\\
50000	323.06\\
52000	320.29\\
54000	313.02\\
56000	316.47\\
58000	322.89\\
60000	306.47\\
62000	304.57\\
64000	325.17\\
66000	327.03\\
68000	317.74\\
70000	323.57\\
72000	319.94\\
74000	328.5\\
76000	327.71\\
78000	331.91\\
80000	323.94\\
82000	346.28\\
84000	352.96\\
86000	329.96\\
88000	339.47\\
90000	322.57\\
92000	318.45\\
94000	326.83\\
96000	323.94\\
98000	317.2\\
1e+05	317.22\\
};
\addplot [color=mycolor98]
  table[row sep=crcr]{%
0	4395.3\\
2000	1874.7\\
4000	596.96\\
6000	467.86\\
8000	416.13\\
10000	431.56\\
12000	444.21\\
14000	428.81\\
16000	424.92\\
18000	431.13\\
20000	452.75\\
22000	434.23\\
24000	422.36\\
26000	426.79\\
28000	415.67\\
30000	444.75\\
32000	444.08\\
34000	436.01\\
36000	445.58\\
38000	429.29\\
40000	439.02\\
42000	435.31\\
44000	439.29\\
46000	434.63\\
50000	429.9\\
52000	437.56\\
54000	413.5\\
56000	420.11\\
58000	416.25\\
60000	428.21\\
62000	412.51\\
64000	401.85\\
66000	429.89\\
68000	457.54\\
70000	434.12\\
72000	400.6\\
74000	409.32\\
76000	429.35\\
78000	445.65\\
82000	412.58\\
84000	432.44\\
88000	436.31\\
90000	418\\
92000	414.12\\
94000	435.71\\
96000	432.64\\
98000	450.87\\
1e+05	454.61\\
};
\addlegendentry{\ClasskIVLegend}

      % \input{../../../../.tex/20/KS/SizeEvolutionsfig5.tex}
      \legend{}

      \nextgroupplot[%
        % xmin=0,xmax=1e+05,
        % ymin=100,ymax=1.539e+05,
      ] % This file was created by matlab2tikz.
%
\definecolor{mycolor1}{rgb}{0.00146,0.00047,0.01387}%
\definecolor{mycolor2}{rgb}{0.01166,0.00942,0.06346}%
\definecolor{mycolor3}{rgb}{0.02943,0.02150,0.11462}%
\definecolor{mycolor4}{rgb}{0.06134,0.03659,0.17764}%
\definecolor{mycolor5}{rgb}{0.08741,0.04456,0.22481}%
\definecolor{mycolor6}{rgb}{0.12291,0.04754,0.28162}%
\definecolor{mycolor7}{rgb}{0.14907,0.04547,0.31709}%
\definecolor{mycolor8}{rgb}{0.18343,0.04033,0.35497}%
\definecolor{mycolor9}{rgb}{0.21795,0.03662,0.38352}%
\definecolor{mycolor10}{rgb}{0.24497,0.03705,0.40001}%
\definecolor{mycolor11}{rgb}{0.27135,0.04092,0.41198}%
\definecolor{mycolor12}{rgb}{0.29076,0.04564,0.41864}%
\definecolor{mycolor13}{rgb}{0.31628,0.05349,0.42512}%
\definecolor{mycolor14}{rgb}{0.33522,0.06006,0.42852}%
\definecolor{mycolor15}{rgb}{0.36028,0.06925,0.43150}%
\definecolor{mycolor16}{rgb}{0.37900,0.07625,0.43272}%
\definecolor{mycolor17}{rgb}{0.39767,0.08326,0.43318}%
\definecolor{mycolor18}{rgb}{0.41633,0.09020,0.43294}%
\definecolor{mycolor19}{rgb}{0.42877,0.09479,0.43241}%
\definecolor{mycolor20}{rgb}{0.44743,0.10160,0.43108}%
\definecolor{mycolor21}{rgb}{0.46610,0.10832,0.42912}%
\definecolor{mycolor22}{rgb}{0.47856,0.11276,0.42747}%
\definecolor{mycolor23}{rgb}{0.49726,0.11938,0.42449}%
\definecolor{mycolor24}{rgb}{0.50973,0.12377,0.42216}%
\definecolor{mycolor25}{rgb}{0.52221,0.12815,0.41955}%
\definecolor{mycolor26}{rgb}{0.53468,0.13253,0.41667}%
\definecolor{mycolor27}{rgb}{0.55339,0.13913,0.41183}%
\definecolor{mycolor28}{rgb}{0.56585,0.14357,0.40826}%
\definecolor{mycolor29}{rgb}{0.57830,0.14804,0.40441}%
\definecolor{mycolor30}{rgb}{0.59073,0.15256,0.40029}%
\definecolor{mycolor31}{rgb}{0.60314,0.15715,0.39589}%
\definecolor{mycolor32}{rgb}{0.61551,0.16182,0.39122}%
\definecolor{mycolor33}{rgb}{0.62169,0.16418,0.38878}%
\definecolor{mycolor34}{rgb}{0.63400,0.16899,0.38370}%
\definecolor{mycolor35}{rgb}{0.64626,0.17391,0.37836}%
\definecolor{mycolor36}{rgb}{0.65846,0.17896,0.37275}%
\definecolor{mycolor37}{rgb}{0.66454,0.18154,0.36985}%
\definecolor{mycolor38}{rgb}{0.67664,0.18681,0.36385}%
\definecolor{mycolor39}{rgb}{0.68865,0.19224,0.35760}%
\definecolor{mycolor40}{rgb}{0.69463,0.19502,0.35439}%
\definecolor{mycolor41}{rgb}{0.70650,0.20073,0.34778}%
\definecolor{mycolor42}{rgb}{0.71240,0.20366,0.34438}%
\definecolor{mycolor43}{rgb}{0.72410,0.20967,0.33742}%
\definecolor{mycolor44}{rgb}{0.72991,0.21276,0.33386}%
\definecolor{mycolor45}{rgb}{0.74142,0.21911,0.32658}%
\definecolor{mycolor46}{rgb}{0.74713,0.22238,0.32286}%
\definecolor{mycolor47}{rgb}{0.75279,0.22571,0.31909}%
\definecolor{mycolor48}{rgb}{0.76401,0.23255,0.31140}%
\definecolor{mycolor49}{rgb}{0.76956,0.23608,0.30749}%
\definecolor{mycolor50}{rgb}{0.77506,0.23967,0.30353}%
\definecolor{mycolor51}{rgb}{0.79129,0.25086,0.29139}%
\definecolor{mycolor52}{rgb}{0.79661,0.25473,0.28726}%
\definecolor{mycolor53}{rgb}{0.80187,0.25867,0.28310}%
\definecolor{mycolor54}{rgb}{0.81224,0.26679,0.27466}%
\definecolor{mycolor55}{rgb}{0.81734,0.27095,0.27039}%
\definecolor{mycolor56}{rgb}{0.82737,0.27952,0.26175}%
\definecolor{mycolor57}{rgb}{0.83230,0.28391,0.25738}%
\definecolor{mycolor58}{rgb}{0.83717,0.28839,0.25299}%
\definecolor{mycolor59}{rgb}{0.84671,0.29756,0.24411}%
\definecolor{mycolor60}{rgb}{0.85138,0.30226,0.23964}%
\definecolor{mycolor61}{rgb}{0.85599,0.30704,0.23513}%
\definecolor{mycolor62}{rgb}{0.86053,0.31189,0.23061}%
\definecolor{mycolor63}{rgb}{0.87374,0.32691,0.21689}%
\definecolor{mycolor64}{rgb}{0.87800,0.33206,0.21227}%
\definecolor{mycolor65}{rgb}{0.89034,0.34796,0.19829}%
\definecolor{mycolor66}{rgb}{0.89819,0.35891,0.18886}%
\definecolor{mycolor67}{rgb}{0.90200,0.36449,0.18412}%
\definecolor{mycolor68}{rgb}{0.90573,0.37014,0.17935}%
\definecolor{mycolor69}{rgb}{0.90939,0.37586,0.17456}%
\definecolor{mycolor70}{rgb}{0.91297,0.38164,0.16975}%
\definecolor{mycolor71}{rgb}{0.91646,0.38748,0.16492}%
\definecolor{mycolor72}{rgb}{0.91988,0.39339,0.16007}%
\definecolor{mycolor73}{rgb}{0.92322,0.39936,0.15519}%
\definecolor{mycolor74}{rgb}{0.92647,0.40539,0.15029}%
\definecolor{mycolor75}{rgb}{0.92964,0.41148,0.14537}%
\definecolor{mycolor76}{rgb}{0.93274,0.41763,0.14042}%
\definecolor{mycolor77}{rgb}{0.93575,0.42383,0.13544}%
\definecolor{mycolor78}{rgb}{0.93868,0.43009,0.13044}%
\definecolor{mycolor79}{rgb}{0.94152,0.43640,0.12541}%
\definecolor{mycolor80}{rgb}{0.94429,0.44277,0.12035}%
\definecolor{mycolor81}{rgb}{0.94956,0.45566,0.11016}%
\definecolor{mycolor82}{rgb}{0.95208,0.46218,0.10503}%
\definecolor{mycolor83}{rgb}{0.96129,0.48872,0.08429}%
\definecolor{mycolor84}{rgb}{0.96540,0.50225,0.07386}%
\definecolor{mycolor85}{rgb}{0.96732,0.50908,0.06866}%
\definecolor{mycolor86}{rgb}{0.97259,0.52980,0.05332}%
\definecolor{mycolor87}{rgb}{0.97568,0.54380,0.04362}%
\definecolor{mycolor88}{rgb}{0.98082,0.57221,0.02851}%
\definecolor{mycolor89}{rgb}{0.98189,0.57939,0.02625}%
\definecolor{mycolor90}{rgb}{0.98378,0.59385,0.02377}%
\definecolor{mycolor91}{rgb}{0.98532,0.60842,0.02420}%
\definecolor{mycolor92}{rgb}{0.98696,0.63048,0.03091}%
\definecolor{mycolor93}{rgb}{0.98793,0.66025,0.05175}%
\definecolor{mycolor94}{rgb}{0.98771,0.68281,0.07249}%
\definecolor{mycolor95}{rgb}{0.97750,0.78226,0.18592}%
\definecolor{mycolor96}{rgb}{0.96624,0.83619,0.26153}%
\definecolor{mycolor97}{rgb}{0.96252,0.85148,0.28555}%
\definecolor{mycolor98}{rgb}{0.98836,0.99836,0.64492}%
%
\addplot [color=mycolor1]
  table[row sep=crcr]{%
0	4395.3\\
2000	1.0042e+05\\
4000	1.4503e+05\\
6000	1.5328e+05\\
8000	1.539e+05\\
10000	1.515e+05\\
12000	1.5364e+05\\
14000	1.5117e+05\\
16000	1.4768e+05\\
18000	1.502e+05\\
20000	1.4923e+05\\
22000	1.5048e+05\\
24000	1.5031e+05\\
26000	1.5137e+05\\
28000	1.4932e+05\\
30000	1.5265e+05\\
34000	1.5062e+05\\
36000	1.4862e+05\\
38000	1.4798e+05\\
42000	1.4871e+05\\
44000	1.4946e+05\\
46000	1.5223e+05\\
50000	1.4947e+05\\
52000	1.4971e+05\\
56000	1.5184e+05\\
58000	1.4985e+05\\
60000	1.5021e+05\\
62000	1.5191e+05\\
64000	1.5319e+05\\
66000	1.4962e+05\\
68000	1.4905e+05\\
70000	1.4895e+05\\
72000	1.5263e+05\\
74000	1.4832e+05\\
76000	1.4869e+05\\
78000	1.4981e+05\\
80000	1.5217e+05\\
82000	1.4732e+05\\
84000	1.4901e+05\\
86000	1.5318e+05\\
88000	1.5059e+05\\
90000	1.4976e+05\\
92000	1.5164e+05\\
94000	1.5196e+05\\
96000	1.5065e+05\\
1e+05	1.5342e+05\\
};
\addlegendentry{\ClasskILegend}

\addplot [color=mycolor2, forget plot]
  table[row sep=crcr]{%
0	4395.3\\
2000	22455\\
4000	25445\\
6000	24818\\
8000	25261\\
10000	25787\\
12000	25461\\
14000	25757\\
16000	26145\\
18000	25558\\
20000	25290\\
24000	25908\\
26000	25745\\
28000	25781\\
30000	25404\\
32000	25924\\
34000	25170\\
36000	25922\\
38000	26434\\
40000	25351\\
42000	25177\\
44000	25571\\
46000	24852\\
48000	25989\\
50000	25430\\
52000	25454\\
54000	25400\\
56000	26224\\
58000	25792\\
60000	25960\\
62000	25436\\
64000	26122\\
66000	25647\\
70000	25917\\
72000	25486\\
74000	26193\\
76000	25623\\
78000	25474\\
80000	25190\\
82000	25976\\
84000	26190\\
86000	25369\\
88000	25968\\
92000	25377\\
94000	25280\\
96000	25615\\
98000	26094\\
1e+05	26124\\
};
\addplot [color=mycolor3, forget plot]
  table[row sep=crcr]{%
0	4395.3\\
2000	33116\\
4000	38356\\
6000	38542\\
8000	37842\\
10000	39241\\
12000	39351\\
14000	39641\\
16000	38620\\
18000	38782\\
20000	40676\\
22000	39201\\
24000	38808\\
26000	39899\\
28000	38962\\
30000	39558\\
34000	38250\\
36000	38515\\
38000	39233\\
40000	39458\\
42000	38716\\
44000	38894\\
46000	39205\\
48000	39028\\
50000	39800\\
52000	38882\\
54000	39088\\
56000	40448\\
58000	38230\\
60000	39196\\
62000	37808\\
64000	38710\\
66000	38355\\
68000	39057\\
72000	39596\\
74000	39433\\
76000	38392\\
78000	38510\\
82000	39955\\
84000	40128\\
86000	39295\\
88000	40213\\
90000	39402\\
92000	39289\\
94000	39016\\
96000	40009\\
98000	39309\\
1e+05	38923\\
};
\addplot [color=mycolor4, forget plot]
  table[row sep=crcr]{%
0	4395.3\\
2000	13544\\
4000	14456\\
6000	14496\\
8000	13996\\
10000	13917\\
12000	14432\\
14000	14289\\
16000	14618\\
18000	14265\\
22000	14563\\
24000	14315\\
26000	14210\\
28000	14328\\
30000	14009\\
32000	14012\\
34000	14103\\
36000	14293\\
38000	14394\\
40000	14303\\
42000	14171\\
44000	14515\\
46000	14114\\
48000	14526\\
50000	14066\\
54000	14261\\
56000	13873\\
58000	14405\\
60000	14550\\
62000	14237\\
64000	14036\\
66000	14353\\
68000	14374\\
70000	13909\\
72000	14026\\
76000	14553\\
78000	14522\\
80000	14274\\
82000	14235\\
84000	14388\\
86000	14275\\
88000	13686\\
90000	14337\\
92000	14486\\
94000	14364\\
96000	14113\\
98000	14243\\
1e+05	14733\\
};
\addplot [color=mycolor5, forget plot]
  table[row sep=crcr]{%
0	4395.3\\
2000	43921\\
4000	60906\\
6000	61867\\
8000	63661\\
10000	62974\\
12000	61481\\
14000	62424\\
16000	63023\\
18000	61782\\
20000	62786\\
22000	64323\\
24000	62133\\
26000	63838\\
28000	63103\\
30000	62045\\
32000	63000\\
34000	63470\\
36000	63672\\
38000	63500\\
40000	65021\\
42000	62638\\
46000	61851\\
48000	61775\\
50000	62406\\
52000	62167\\
54000	63005\\
56000	62972\\
58000	65128\\
60000	62278\\
62000	62658\\
64000	62728\\
68000	61548\\
70000	62610\\
72000	62508\\
74000	64182\\
76000	62737\\
78000	62153\\
80000	62861\\
82000	64287\\
84000	62135\\
86000	61671\\
88000	64097\\
90000	62076\\
92000	60918\\
94000	63090\\
96000	63142\\
1e+05	62787\\
};
\addplot [color=mycolor6, forget plot]
  table[row sep=crcr]{%
0	4395.3\\
2000	39344\\
4000	54514\\
6000	55046\\
8000	53470\\
10000	52798\\
12000	53775\\
14000	54148\\
18000	56754\\
20000	55260\\
22000	55452\\
24000	55196\\
26000	55162\\
28000	55820\\
30000	54615\\
32000	54581\\
34000	56537\\
36000	56692\\
40000	54697\\
42000	55383\\
44000	55372\\
46000	56089\\
48000	53908\\
50000	55616\\
52000	55624\\
54000	55186\\
56000	54383\\
58000	54902\\
60000	55121\\
62000	53914\\
64000	53774\\
66000	53798\\
68000	55597\\
70000	53879\\
72000	55227\\
74000	54498\\
78000	55775\\
80000	55146\\
82000	55441\\
86000	55246\\
88000	54435\\
90000	54655\\
92000	55223\\
94000	54083\\
96000	55485\\
98000	55455\\
1e+05	53217\\
};
\addplot [color=mycolor7, forget plot]
  table[row sep=crcr]{%
0	4395.3\\
2000	23881\\
4000	28372\\
6000	28110\\
8000	27300\\
10000	27380\\
12000	27839\\
14000	28049\\
16000	27526\\
18000	28431\\
20000	27555\\
22000	28741\\
24000	27938\\
26000	27655\\
28000	28208\\
30000	27621\\
34000	28049\\
36000	28400\\
38000	28234\\
42000	28533\\
44000	27757\\
46000	27995\\
48000	28111\\
50000	27528\\
52000	27713\\
54000	27586\\
56000	27938\\
58000	27476\\
60000	28466\\
62000	28910\\
66000	28049\\
68000	27836\\
70000	28055\\
72000	27933\\
74000	28657\\
76000	27373\\
78000	28826\\
80000	28533\\
82000	27605\\
84000	27798\\
86000	27198\\
88000	27795\\
90000	28514\\
94000	28443\\
96000	28225\\
98000	27241\\
1e+05	26809\\
};
\addplot [color=mycolor8, forget plot]
  table[row sep=crcr]{%
0	4395.3\\
2000	17743\\
4000	20215\\
6000	19874\\
8000	20917\\
10000	20503\\
12000	19964\\
14000	20523\\
16000	20145\\
18000	20483\\
20000	20334\\
22000	19773\\
24000	20750\\
26000	20703\\
28000	20859\\
30000	20420\\
32000	19877\\
34000	20395\\
36000	20562\\
38000	20401\\
40000	19939\\
42000	20240\\
44000	20611\\
46000	20467\\
48000	19761\\
50000	20232\\
52000	20443\\
54000	20885\\
56000	20191\\
58000	20782\\
60000	20234\\
62000	20648\\
64000	20181\\
66000	20584\\
68000	20535\\
72000	20925\\
74000	20586\\
76000	20151\\
78000	20560\\
80000	20571\\
82000	20283\\
84000	20413\\
86000	20239\\
88000	20783\\
90000	19962\\
92000	20982\\
94000	20829\\
96000	20265\\
98000	20977\\
1e+05	20159\\
};
\addplot [color=mycolor9, forget plot]
  table[row sep=crcr]{%
0	4395.3\\
2000	12738\\
4000	14449\\
6000	13986\\
8000	14142\\
10000	13970\\
12000	13560\\
14000	13588\\
16000	14060\\
18000	14296\\
20000	13920\\
22000	13958\\
24000	13847\\
26000	14030\\
28000	14285\\
30000	13730\\
32000	14235\\
34000	14180\\
36000	14348\\
38000	13708\\
40000	14180\\
42000	13908\\
44000	13978\\
46000	13819\\
48000	14175\\
50000	14189\\
52000	14332\\
54000	13792\\
56000	14052\\
58000	14022\\
60000	13670\\
62000	14064\\
64000	13882\\
66000	13976\\
68000	14329\\
70000	13914\\
72000	14032\\
74000	14517\\
76000	14225\\
78000	14335\\
80000	13875\\
82000	14228\\
84000	14189\\
86000	14007\\
88000	14136\\
90000	14079\\
92000	14434\\
94000	13417\\
96000	14049\\
98000	13876\\
1e+05	13862\\
};
\addplot [color=mycolor10, forget plot]
  table[row sep=crcr]{%
0	4395.3\\
2000	17740\\
4000	22124\\
6000	22078\\
8000	22108\\
10000	22002\\
12000	21821\\
14000	21808\\
16000	22307\\
18000	21505\\
20000	22678\\
22000	21823\\
24000	22338\\
28000	22066\\
32000	21569\\
34000	21496\\
38000	22057\\
40000	21597\\
42000	22214\\
44000	22107\\
46000	21691\\
48000	21475\\
50000	21877\\
52000	21087\\
54000	21767\\
56000	21523\\
58000	21879\\
60000	21097\\
62000	21870\\
64000	21558\\
66000	22170\\
68000	21469\\
70000	22328\\
74000	21709\\
76000	21812\\
78000	22285\\
80000	21292\\
84000	21577\\
86000	21973\\
88000	21897\\
90000	21979\\
92000	21375\\
94000	22071\\
96000	21496\\
98000	21866\\
1e+05	21684\\
};
\addplot [color=mycolor11, forget plot]
  table[row sep=crcr]{%
0	4395.3\\
2000	8383.2\\
4000	8409.4\\
6000	8582\\
10000	8596.9\\
12000	8462.2\\
14000	8672\\
16000	8488.6\\
18000	8576.2\\
20000	8923.6\\
22000	8695.7\\
24000	8280.7\\
26000	8807.7\\
30000	8704.6\\
32000	8434.1\\
34000	8499.3\\
36000	8589.2\\
38000	8449.6\\
40000	8877.5\\
42000	8647\\
44000	8918.9\\
46000	8643\\
48000	8781.9\\
50000	8510.1\\
52000	8756.1\\
56000	8657.1\\
58000	8510.2\\
60000	8499.3\\
62000	8754\\
64000	8819.2\\
66000	8971.7\\
68000	8656.9\\
74000	8530.7\\
76000	8759.4\\
80000	8505.7\\
82000	8477.2\\
84000	8609\\
86000	8467.5\\
88000	8702.6\\
90000	8685.9\\
92000	8549.6\\
94000	8514.1\\
96000	8686.7\\
98000	8732\\
1e+05	8538.8\\
};
\addplot [color=mycolor12, forget plot]
  table[row sep=crcr]{%
0	4395.3\\
2000	13075\\
4000	14618\\
6000	14044\\
8000	14789\\
10000	14393\\
12000	14548\\
14000	14585\\
16000	14523\\
18000	13974\\
20000	15133\\
22000	15146\\
24000	14923\\
26000	14392\\
28000	14777\\
30000	14841\\
32000	14313\\
34000	14593\\
36000	14409\\
38000	14278\\
40000	14678\\
44000	14922\\
46000	14564\\
48000	14523\\
50000	14618\\
52000	14927\\
54000	14540\\
56000	14555\\
58000	14150\\
60000	14350\\
62000	15017\\
64000	14622\\
66000	14850\\
68000	14582\\
70000	14673\\
72000	14853\\
74000	14572\\
76000	14396\\
78000	14493\\
80000	14504\\
82000	14859\\
84000	15123\\
86000	14622\\
88000	14966\\
90000	14475\\
92000	14249\\
94000	14755\\
96000	14769\\
98000	14563\\
1e+05	14595\\
};
\addplot [color=mycolor13, forget plot]
  table[row sep=crcr]{%
0	4395.3\\
2000	10313\\
4000	11172\\
6000	11063\\
8000	11210\\
10000	11167\\
12000	11295\\
14000	11125\\
16000	11295\\
18000	10965\\
20000	10902\\
22000	10888\\
24000	10796\\
26000	10793\\
28000	11105\\
30000	11012\\
32000	11178\\
34000	11233\\
36000	11202\\
38000	10890\\
40000	11127\\
42000	10794\\
44000	10913\\
46000	11317\\
48000	11094\\
50000	10990\\
52000	11046\\
54000	10959\\
56000	10962\\
58000	10909\\
60000	11147\\
62000	10963\\
64000	11004\\
66000	11373\\
70000	11041\\
72000	11338\\
74000	11216\\
76000	10857\\
78000	10835\\
80000	11090\\
82000	10846\\
84000	10960\\
86000	11033\\
88000	11170\\
90000	11218\\
92000	10958\\
94000	11614\\
96000	11285\\
98000	10924\\
1e+05	11221\\
};
\addplot [color=mycolor14, forget plot]
  table[row sep=crcr]{%
0	4395.3\\
2000	19926\\
4000	23504\\
6000	24541\\
8000	23731\\
10000	24623\\
12000	24163\\
14000	24430\\
16000	24918\\
18000	23705\\
20000	22879\\
22000	23940\\
24000	23744\\
26000	23848\\
28000	23429\\
30000	23669\\
32000	23991\\
34000	23458\\
36000	24166\\
38000	24542\\
40000	23520\\
42000	24032\\
44000	23667\\
46000	24478\\
48000	24759\\
50000	24835\\
52000	23705\\
54000	24290\\
56000	23343\\
58000	24056\\
60000	24347\\
62000	23842\\
64000	24435\\
66000	23931\\
68000	24460\\
70000	23876\\
72000	24305\\
74000	23480\\
76000	23862\\
78000	23626\\
80000	24046\\
82000	24656\\
84000	24495\\
86000	23542\\
88000	24318\\
92000	24554\\
94000	23453\\
96000	24095\\
98000	25246\\
1e+05	23886\\
};
\addplot [color=mycolor15, forget plot]
  table[row sep=crcr]{%
0	4395.3\\
2000	14590\\
4000	16762\\
6000	16895\\
8000	17174\\
10000	16489\\
12000	16271\\
14000	16430\\
16000	17453\\
18000	17192\\
20000	17006\\
22000	17424\\
24000	16548\\
26000	16784\\
28000	16925\\
30000	16687\\
32000	16969\\
34000	16773\\
36000	17296\\
38000	17146\\
40000	17247\\
42000	16827\\
44000	16482\\
46000	17331\\
48000	16767\\
50000	16815\\
52000	16992\\
54000	17228\\
56000	16177\\
58000	16778\\
60000	17072\\
62000	16619\\
64000	16790\\
66000	17196\\
68000	16945\\
70000	17006\\
72000	16786\\
74000	15933\\
76000	16553\\
78000	17429\\
82000	16924\\
84000	17180\\
86000	16780\\
88000	17063\\
90000	16766\\
92000	17129\\
94000	16956\\
96000	16874\\
98000	16886\\
1e+05	17152\\
};
\addplot [color=mycolor16, forget plot]
  table[row sep=crcr]{%
0	4395.3\\
2000	6917.7\\
4000	7016.8\\
6000	7269.4\\
8000	6920.7\\
10000	7073.1\\
12000	6991.1\\
14000	6817.2\\
16000	7090.4\\
20000	7068.8\\
22000	7085.5\\
24000	6892.8\\
26000	6974.9\\
28000	7022.2\\
30000	6960.2\\
32000	6961.5\\
36000	6915.7\\
38000	7220\\
40000	6920.2\\
42000	6967.3\\
44000	6891.1\\
46000	6944.9\\
48000	6931.5\\
50000	6868\\
52000	6993\\
54000	6944\\
56000	7106.5\\
58000	7012.5\\
60000	6977.5\\
62000	6963.1\\
64000	6984.5\\
66000	6931.6\\
70000	7073.2\\
72000	6896.3\\
74000	6883.3\\
76000	6974.7\\
78000	6999.7\\
82000	6991.1\\
84000	6690.1\\
86000	6979.8\\
88000	6894.8\\
90000	6972.5\\
92000	6893.1\\
94000	7300.9\\
96000	6999.2\\
98000	6924.8\\
1e+05	7079.7\\
};
\addplot [color=mycolor17, forget plot]
  table[row sep=crcr]{%
0	4395.3\\
2000	14516\\
4000	16127\\
6000	15708\\
8000	15588\\
10000	15827\\
12000	16222\\
14000	15828\\
16000	16317\\
18000	15560\\
20000	15712\\
22000	16267\\
24000	15935\\
26000	15722\\
28000	16244\\
30000	16081\\
32000	15765\\
34000	15886\\
36000	15874\\
38000	16070\\
40000	15922\\
42000	16359\\
46000	15726\\
48000	16509\\
50000	15850\\
52000	15618\\
54000	15835\\
56000	15906\\
58000	15603\\
60000	15629\\
62000	16028\\
64000	16195\\
66000	16032\\
68000	15613\\
70000	15728\\
72000	16037\\
76000	15481\\
78000	15325\\
80000	16458\\
82000	16004\\
84000	15959\\
86000	15823\\
88000	15202\\
90000	16082\\
92000	15386\\
94000	15784\\
96000	15802\\
98000	16118\\
1e+05	16046\\
};
\addplot [color=mycolor18, forget plot]
  table[row sep=crcr]{%
0	4395.3\\
2000	19774\\
4000	24442\\
6000	24181\\
8000	24337\\
10000	24341\\
12000	24255\\
14000	24676\\
16000	24031\\
18000	24496\\
20000	24368\\
22000	24439\\
24000	23709\\
26000	24116\\
28000	24609\\
30000	24044\\
32000	24279\\
34000	24391\\
36000	24326\\
38000	23764\\
40000	24592\\
42000	24219\\
44000	24540\\
48000	23853\\
50000	24472\\
52000	24535\\
54000	24080\\
56000	23983\\
58000	24417\\
60000	25088\\
62000	24745\\
64000	24242\\
66000	24869\\
68000	23692\\
70000	24236\\
72000	22702\\
74000	24332\\
76000	24446\\
78000	25339\\
80000	23622\\
82000	23537\\
84000	23930\\
86000	24727\\
88000	24364\\
90000	24933\\
92000	24696\\
94000	24579\\
96000	24543\\
98000	23535\\
1e+05	24145\\
};
\addplot [color=mycolor19, forget plot]
  table[row sep=crcr]{%
0	4395.3\\
2000	12808\\
4000	14554\\
6000	14597\\
8000	15029\\
10000	14679\\
12000	14849\\
14000	14352\\
16000	14624\\
18000	14661\\
20000	15065\\
22000	14647\\
24000	14911\\
26000	14462\\
28000	15031\\
30000	14651\\
32000	14353\\
34000	14721\\
36000	14838\\
38000	15085\\
40000	14581\\
42000	14678\\
46000	15100\\
48000	14942\\
50000	14962\\
52000	15074\\
54000	14786\\
56000	14634\\
58000	14920\\
60000	14420\\
62000	15030\\
64000	15131\\
66000	14762\\
68000	15259\\
70000	14716\\
72000	14759\\
74000	14695\\
76000	14418\\
80000	14948\\
82000	14621\\
84000	14598\\
86000	14701\\
88000	14871\\
90000	14716\\
92000	14359\\
94000	14428\\
96000	14624\\
98000	14247\\
1e+05	14514\\
};
\addplot [color=mycolor20, forget plot]
  table[row sep=crcr]{%
0	4395.3\\
2000	16132\\
4000	19029\\
6000	18752\\
8000	18936\\
10000	19041\\
12000	18818\\
14000	18803\\
16000	19277\\
18000	18867\\
20000	19358\\
22000	18628\\
24000	18953\\
26000	18451\\
28000	19058\\
30000	19536\\
32000	19133\\
34000	18911\\
36000	18767\\
38000	19114\\
40000	19229\\
42000	19017\\
44000	19446\\
46000	18714\\
48000	18895\\
50000	19391\\
52000	19234\\
54000	18690\\
56000	19042\\
58000	19176\\
60000	18727\\
64000	19600\\
68000	19659\\
70000	19210\\
72000	19021\\
74000	18641\\
76000	18776\\
78000	19014\\
80000	18982\\
82000	19620\\
84000	18786\\
86000	18885\\
88000	19177\\
90000	19416\\
92000	19528\\
94000	19124\\
96000	19031\\
98000	19072\\
1e+05	18990\\
};
\addplot [color=mycolor21, forget plot]
  table[row sep=crcr]{%
0	4395.3\\
2000	8821.7\\
4000	9727.5\\
6000	9774.4\\
8000	9454.8\\
10000	9550.8\\
12000	9390.1\\
14000	9793.8\\
16000	9371.7\\
20000	9605.9\\
24000	9595.2\\
28000	9758.5\\
30000	9531.9\\
32000	9549.9\\
34000	9743.8\\
36000	9728.5\\
38000	9761.8\\
40000	9593.9\\
42000	9489.6\\
44000	9640.3\\
46000	9689.2\\
48000	9860.5\\
50000	9677\\
54000	9810.2\\
56000	9646.3\\
58000	9571.8\\
60000	9968.8\\
62000	9664.5\\
64000	9528.2\\
66000	9693.5\\
68000	9689.2\\
70000	9413.4\\
72000	9592.2\\
74000	9695.9\\
76000	9726.4\\
78000	9390.5\\
80000	9558.1\\
84000	9757.5\\
86000	9559.5\\
88000	9570.5\\
90000	9903.7\\
94000	9602\\
96000	9534.4\\
98000	9436.2\\
1e+05	9417\\
};
\addplot [color=mycolor22, forget plot]
  table[row sep=crcr]{%
0	4395.3\\
2000	19457\\
4000	24628\\
6000	25551\\
8000	24795\\
10000	25618\\
12000	24882\\
14000	25177\\
16000	24964\\
18000	24659\\
20000	24801\\
22000	24161\\
24000	25243\\
26000	24530\\
28000	24516\\
30000	24815\\
32000	25513\\
34000	25785\\
36000	24491\\
38000	24665\\
40000	25393\\
42000	25254\\
44000	24796\\
46000	25403\\
48000	24950\\
50000	24687\\
52000	25064\\
54000	24893\\
56000	25110\\
58000	24695\\
60000	25450\\
62000	24868\\
64000	24176\\
68000	25108\\
70000	25912\\
72000	23816\\
74000	25605\\
76000	24555\\
78000	25325\\
80000	25006\\
82000	25202\\
84000	26074\\
86000	25406\\
88000	25225\\
90000	24869\\
92000	25352\\
94000	25323\\
98000	25067\\
1e+05	24213\\
};
\addplot [color=mycolor23, forget plot]
  table[row sep=crcr]{%
0	4395.3\\
2000	10958\\
4000	12118\\
6000	12150\\
8000	12072\\
10000	12115\\
12000	12457\\
14000	12402\\
16000	12430\\
18000	12272\\
20000	12263\\
22000	12196\\
26000	12286\\
28000	12087\\
30000	12074\\
32000	12506\\
34000	12243\\
36000	12083\\
38000	12339\\
40000	12243\\
42000	12446\\
44000	12338\\
46000	12129\\
48000	12211\\
50000	12464\\
54000	12462\\
56000	12314\\
58000	12310\\
60000	12236\\
62000	12120\\
64000	12291\\
66000	12175\\
68000	12172\\
72000	12301\\
74000	12486\\
76000	12548\\
78000	11889\\
80000	12305\\
82000	12437\\
84000	12048\\
86000	12026\\
88000	12368\\
90000	12032\\
94000	12328\\
96000	12233\\
98000	12338\\
1e+05	12399\\
};
\addplot [color=mycolor24, forget plot]
  table[row sep=crcr]{%
0	4395.3\\
2000	5044.2\\
4000	4813.9\\
6000	4768.5\\
8000	4816.3\\
10000	4803.5\\
12000	4679.9\\
14000	4717.2\\
16000	4729.8\\
18000	4979.8\\
20000	4616.5\\
22000	4711.3\\
24000	4772.1\\
26000	4764.1\\
28000	4774.6\\
30000	4760.5\\
32000	4908.9\\
34000	4811.2\\
36000	4775.4\\
38000	4869.4\\
40000	4930.6\\
42000	4893.8\\
44000	4750.2\\
46000	4794.8\\
48000	4825.6\\
50000	4575.8\\
52000	4948.7\\
54000	4782.5\\
56000	4752.6\\
58000	4809.9\\
60000	4694.4\\
62000	4793.1\\
64000	4708\\
66000	4813.2\\
68000	4902.7\\
70000	4799.1\\
72000	4750.2\\
74000	4903.2\\
76000	4769.5\\
78000	4673\\
80000	4711.7\\
82000	4804\\
84000	4685.1\\
86000	4783\\
88000	4792.6\\
90000	4886\\
92000	4753.8\\
94000	4842.4\\
96000	4707.8\\
98000	4623.2\\
1e+05	4718.9\\
};
\addplot [color=mycolor25, forget plot]
  table[row sep=crcr]{%
0	4395.3\\
2000	11221\\
4000	12318\\
6000	12466\\
8000	12415\\
10000	12540\\
12000	12364\\
14000	12576\\
16000	12344\\
18000	12463\\
20000	12500\\
24000	12396\\
26000	12522\\
28000	12325\\
30000	12384\\
32000	12279\\
34000	12520\\
36000	11988\\
38000	12184\\
40000	12562\\
42000	12391\\
44000	12574\\
46000	12905\\
48000	12079\\
50000	12557\\
52000	12246\\
54000	12366\\
56000	12296\\
58000	12336\\
60000	12155\\
62000	12263\\
64000	12787\\
66000	12650\\
68000	12849\\
70000	12365\\
72000	12171\\
76000	12856\\
78000	12705\\
80000	12215\\
82000	12338\\
84000	12201\\
86000	12320\\
88000	12660\\
90000	12356\\
92000	12508\\
94000	12111\\
96000	12315\\
98000	12458\\
1e+05	12328\\
};
\addplot [color=mycolor26, forget plot]
  table[row sep=crcr]{%
0	4395.3\\
2000	7928.5\\
4000	8275.7\\
6000	8225.3\\
8000	8407.3\\
10000	8550.4\\
12000	8477.2\\
14000	8250.5\\
16000	8469.6\\
18000	8246.1\\
20000	8238.9\\
22000	8146.5\\
24000	8126.3\\
26000	8243\\
28000	8335.3\\
30000	8388\\
32000	8239.6\\
34000	8355.2\\
38000	8073.7\\
40000	8190.6\\
42000	7979.9\\
44000	8208.8\\
46000	8212.5\\
48000	8466.1\\
50000	8318.1\\
52000	8295.4\\
54000	8186.3\\
56000	8101.4\\
58000	8403.7\\
60000	8159.5\\
62000	8081.7\\
64000	8180.9\\
66000	7834.8\\
68000	8428.8\\
70000	8207.2\\
72000	8075.2\\
74000	8160.4\\
76000	8186.2\\
78000	7960\\
80000	8175.7\\
82000	8339.1\\
84000	8260.5\\
86000	8019.4\\
88000	8463.5\\
90000	8507.8\\
92000	8355.3\\
94000	7998.2\\
96000	8374\\
1e+05	8221.4\\
};
\addplot [color=mycolor27, forget plot]
  table[row sep=crcr]{%
0	4395.3\\
2000	9092.8\\
4000	9767.8\\
6000	10082\\
8000	9636\\
10000	10065\\
12000	9835.3\\
14000	9378.9\\
16000	9683.2\\
18000	9948.7\\
20000	9520.9\\
22000	9449.5\\
24000	9597.5\\
26000	9537.9\\
28000	9697.3\\
30000	9453.4\\
32000	9878.4\\
34000	9732\\
36000	9949.8\\
38000	9641.6\\
40000	9414.4\\
42000	9680.4\\
44000	9623.4\\
46000	9715.2\\
50000	9568.6\\
52000	9911.7\\
54000	9777\\
56000	9739.9\\
58000	9784.8\\
60000	9903.2\\
62000	9643.9\\
64000	9878.8\\
66000	9648.4\\
68000	9776.1\\
70000	9873.9\\
72000	9703.6\\
74000	9683.9\\
76000	9931.2\\
78000	10078\\
80000	9658.1\\
82000	9582.3\\
84000	9904.4\\
86000	9682.9\\
88000	9962.9\\
92000	9826.1\\
94000	9591.9\\
96000	9902.5\\
98000	9703.4\\
1e+05	9546.7\\
};
\addplot [color=mycolor28, forget plot]
  table[row sep=crcr]{%
0	4395.3\\
2000	5599.4\\
4000	5426.6\\
6000	5492.2\\
8000	5434.7\\
10000	5621.4\\
12000	5511.1\\
14000	5521.7\\
16000	5293\\
18000	5395\\
20000	5369.1\\
22000	5462.3\\
24000	5487.9\\
26000	5477.1\\
28000	5435.3\\
30000	5656.8\\
32000	5471.4\\
34000	5540.9\\
36000	5453.9\\
38000	5433.2\\
40000	5658.8\\
42000	5464.3\\
44000	5402.4\\
46000	5507.3\\
48000	5467.9\\
50000	5408.3\\
52000	5372.6\\
54000	5307.6\\
56000	5403.4\\
58000	5592.4\\
60000	5355.2\\
62000	5494.8\\
64000	5410.1\\
66000	5452.5\\
68000	5377.6\\
70000	5496\\
72000	5535.4\\
74000	5289.5\\
76000	5413.3\\
78000	5647.9\\
80000	5495.7\\
82000	5583.9\\
84000	5596.1\\
86000	5396.1\\
88000	5466.2\\
90000	5342.5\\
92000	5442.4\\
94000	5306.7\\
98000	5389.1\\
1e+05	5352.3\\
};
\addplot [color=mycolor29, forget plot]
  table[row sep=crcr]{%
0	4395.3\\
2000	10430\\
4000	11307\\
6000	11355\\
8000	11621\\
10000	11634\\
12000	11405\\
14000	11574\\
16000	11292\\
18000	11444\\
20000	11292\\
22000	11355\\
24000	11381\\
26000	11138\\
28000	11786\\
30000	11789\\
32000	11587\\
34000	11606\\
38000	11124\\
40000	11526\\
42000	11589\\
46000	11473\\
48000	11520\\
50000	11844\\
52000	11333\\
54000	11460\\
56000	11749\\
58000	11676\\
60000	11712\\
62000	11496\\
64000	11336\\
66000	11401\\
68000	11359\\
70000	11877\\
72000	11385\\
74000	11453\\
76000	11404\\
78000	11428\\
80000	11717\\
82000	11569\\
84000	11743\\
86000	11683\\
88000	11766\\
90000	11980\\
92000	11486\\
94000	11638\\
96000	11364\\
98000	11801\\
1e+05	11712\\
};
\addplot [color=mycolor30, forget plot]
  table[row sep=crcr]{%
0	4395.3\\
2000	7713.9\\
4000	7919.3\\
8000	7691.9\\
10000	7894.3\\
12000	7727.7\\
16000	7973.6\\
18000	7712\\
20000	7906\\
22000	7852.5\\
24000	7929.6\\
26000	7919.5\\
28000	7529.2\\
30000	7860.9\\
32000	7895.7\\
34000	8047.6\\
36000	7964\\
38000	7798.3\\
40000	7949.9\\
42000	7907.7\\
44000	7834.8\\
46000	7975.5\\
48000	7953.2\\
50000	7864.8\\
52000	8013.4\\
54000	7977\\
56000	7975\\
58000	7939.7\\
60000	7735.3\\
62000	7849.4\\
64000	7650\\
66000	7868.9\\
68000	7809\\
70000	7872.7\\
72000	7613.5\\
74000	7978.6\\
76000	8071.1\\
78000	7893.8\\
80000	8029.8\\
82000	7870.1\\
86000	7860.6\\
88000	7964.6\\
90000	7766.7\\
92000	7893.7\\
94000	7833.7\\
96000	7814.3\\
98000	7891.7\\
1e+05	7701.7\\
};
\addplot [color=mycolor31, forget plot]
  table[row sep=crcr]{%
0	4395.3\\
2000	6462\\
4000	6666.6\\
6000	6369.5\\
8000	6433.7\\
10000	6369.1\\
12000	6474.2\\
14000	6710.6\\
16000	6395.4\\
18000	6524.8\\
20000	6505.9\\
22000	6255\\
24000	6517.6\\
28000	6360.2\\
30000	6126.1\\
32000	6399.1\\
34000	6340.2\\
36000	6668.6\\
38000	6509.3\\
44000	6570.1\\
46000	6462.2\\
48000	6527.5\\
50000	6646.9\\
52000	6452.8\\
54000	6404.7\\
56000	6466.4\\
58000	6433.2\\
60000	6530.1\\
62000	6492.6\\
64000	6331.5\\
66000	6575\\
68000	6507.1\\
70000	6512\\
72000	6560.5\\
74000	6659.7\\
76000	6399.4\\
78000	6546.2\\
80000	6286.3\\
82000	6437.9\\
84000	6547.5\\
86000	6367.1\\
88000	6298.5\\
90000	6479.4\\
92000	6441.3\\
96000	6448.5\\
98000	6318.5\\
1e+05	6513.9\\
};
\addplot [color=mycolor32, forget plot]
  table[row sep=crcr]{%
0	4395.3\\
2000	9571.7\\
4000	10502\\
6000	10336\\
8000	10687\\
10000	10458\\
12000	10551\\
14000	10448\\
16000	10468\\
18000	10828\\
20000	10281\\
22000	10321\\
24000	10709\\
26000	10177\\
28000	10698\\
30000	10557\\
32000	10201\\
34000	10665\\
36000	10390\\
38000	10242\\
40000	10167\\
42000	10565\\
44000	10519\\
46000	10317\\
48000	10539\\
50000	10290\\
54000	10594\\
56000	10859\\
58000	10846\\
60000	10438\\
62000	10787\\
66000	10551\\
68000	10641\\
70000	10438\\
74000	10801\\
76000	10707\\
78000	10838\\
80000	10649\\
82000	10362\\
84000	10603\\
86000	10491\\
88000	10218\\
90000	10254\\
92000	10182\\
94000	10689\\
96000	10527\\
98000	10291\\
1e+05	10857\\
};
\addplot [color=mycolor33, forget plot]
  table[row sep=crcr]{%
0	4395.3\\
2000	7976.6\\
4000	8466.6\\
6000	7956.3\\
8000	8346.1\\
10000	8258.3\\
12000	8276.8\\
14000	8062\\
16000	8297.5\\
18000	8155.6\\
20000	8565\\
22000	8315.4\\
24000	8578.2\\
26000	8437.4\\
28000	8248\\
30000	8369.4\\
32000	8088\\
34000	8159\\
36000	8077.4\\
38000	8405.4\\
40000	8283.4\\
42000	8285.9\\
44000	8143.3\\
46000	8269.4\\
48000	8299.3\\
50000	8104.9\\
52000	8214.2\\
54000	8300.1\\
56000	7817.7\\
58000	8114.1\\
60000	8220.9\\
62000	8419.4\\
64000	8079.4\\
66000	8452.9\\
68000	8082.6\\
70000	8442.2\\
72000	8279\\
74000	8152.4\\
76000	8270.5\\
78000	8344\\
80000	8171.9\\
82000	8089.2\\
84000	8060.2\\
86000	8166.8\\
88000	8458.7\\
90000	8431.1\\
92000	8263.7\\
94000	8262.6\\
96000	8133.3\\
98000	8401.9\\
1e+05	8266.2\\
};
\addplot [color=mycolor34, forget plot]
  table[row sep=crcr]{%
0	4395.3\\
2000	4228.9\\
4000	4190.7\\
6000	4042.3\\
8000	4028.7\\
10000	4049.3\\
12000	3994.8\\
14000	4045.4\\
16000	3919.7\\
18000	4062.9\\
20000	4126.7\\
22000	4144.4\\
24000	4139.6\\
28000	4008.5\\
30000	3975.7\\
32000	4103\\
34000	3972.5\\
36000	4080.9\\
38000	4051.8\\
40000	4042.1\\
42000	4096.4\\
44000	4064.1\\
46000	3981\\
48000	4172.4\\
50000	4042.8\\
52000	4068.5\\
54000	3910\\
56000	4024.8\\
60000	4205.1\\
62000	4010.6\\
64000	4015\\
66000	4052.5\\
68000	3983.6\\
70000	4065.3\\
72000	4116.1\\
74000	4208.7\\
76000	4059.5\\
78000	3998.1\\
80000	4060.1\\
82000	4009.9\\
84000	4034.2\\
86000	3912.7\\
88000	4020.5\\
90000	3939\\
92000	4176.7\\
94000	4069.3\\
96000	4123.2\\
98000	4138.3\\
1e+05	4064.3\\
};
\addplot [color=mycolor35]
  table[row sep=crcr]{%
0	4395.3\\
2000	9126.3\\
4000	10183\\
6000	10289\\
10000	10203\\
12000	10004\\
14000	10171\\
16000	10148\\
18000	9992.5\\
20000	10022\\
22000	10095\\
24000	10315\\
26000	10168\\
28000	9860.1\\
30000	10161\\
32000	10308\\
34000	10225\\
36000	10009\\
38000	10141\\
42000	10229\\
44000	10150\\
46000	10041\\
48000	10350\\
50000	9963.6\\
52000	10207\\
54000	10223\\
56000	10112\\
58000	10178\\
60000	10127\\
62000	10162\\
64000	10232\\
66000	10021\\
68000	10089\\
70000	10238\\
72000	9895.9\\
74000	10095\\
76000	10221\\
78000	10096\\
80000	10357\\
82000	10051\\
84000	10291\\
86000	10264\\
88000	9954.6\\
90000	10243\\
92000	10175\\
94000	10036\\
96000	10014\\
98000	10161\\
1e+05	10154\\
};
\addlegendentry{\ClasskIILegend}

\addplot [color=mycolor36, forget plot]
  table[row sep=crcr]{%
0	4395.3\\
2000	4507.9\\
4000	4308.1\\
6000	4280.1\\
8000	4209.8\\
10000	4182.9\\
12000	4271.2\\
14000	4396.1\\
16000	4266.8\\
20000	4306.5\\
22000	4269.7\\
28000	4431.2\\
30000	4290.7\\
32000	4296.8\\
34000	4321.2\\
36000	4248\\
38000	4273.2\\
40000	4250.4\\
42000	4351.5\\
44000	4372.6\\
46000	4306.9\\
48000	4347.5\\
50000	4328.2\\
52000	4286.6\\
54000	4335.3\\
56000	4210.9\\
58000	4442.9\\
60000	4259.1\\
62000	4384.6\\
64000	4216.3\\
66000	4359.8\\
68000	4192.3\\
70000	4224.1\\
72000	4227.1\\
74000	4280.3\\
76000	4238.4\\
78000	4353.2\\
80000	4263.2\\
82000	4330.7\\
84000	4312.2\\
86000	4226.9\\
90000	4243\\
92000	4283.9\\
94000	4219.5\\
96000	4325.1\\
98000	4385.1\\
1e+05	4275.7\\
};
\addplot [color=mycolor37, forget plot]
  table[row sep=crcr]{%
0	4395.3\\
2000	3956.3\\
4000	3683.3\\
6000	3675.5\\
8000	3651.5\\
10000	3685.1\\
12000	3612.1\\
14000	3747.8\\
16000	3680.9\\
18000	3644\\
20000	3662.5\\
22000	3692.3\\
24000	3664.3\\
26000	3619.6\\
28000	3668.7\\
30000	3619.9\\
32000	3649.7\\
34000	3695\\
38000	3620.8\\
40000	3673.8\\
42000	3583\\
44000	3624.5\\
46000	3703.9\\
48000	3641.9\\
50000	3680\\
52000	3661.8\\
54000	3620.7\\
58000	3635.2\\
60000	3670.9\\
62000	3660.9\\
64000	3714.2\\
66000	3627.5\\
68000	3584.7\\
70000	3618.9\\
72000	3715.2\\
74000	3653\\
76000	3654.7\\
78000	3608.8\\
80000	3608.6\\
82000	3641.5\\
84000	3614.4\\
86000	3680.1\\
88000	3635.2\\
90000	3633.9\\
92000	3741.4\\
94000	3633\\
96000	3675.1\\
98000	3729\\
1e+05	3602.7\\
};
\addplot [color=mycolor38, forget plot]
  table[row sep=crcr]{%
0	4395.3\\
2000	3227.1\\
4000	3075.6\\
6000	2997.3\\
8000	3097.9\\
10000	3011.7\\
12000	2997.2\\
14000	3031.4\\
16000	2973.8\\
18000	2984.4\\
20000	2979.2\\
22000	2937.3\\
24000	3116.5\\
26000	2954.8\\
28000	3043\\
30000	2978.2\\
32000	3131.4\\
34000	3203\\
36000	2949.8\\
38000	3087.9\\
40000	3017.3\\
42000	2920.7\\
44000	3120.2\\
46000	2979.1\\
48000	3059.6\\
50000	3017.5\\
52000	2946.3\\
54000	3084.4\\
56000	3005.8\\
58000	2990.9\\
60000	3043.3\\
62000	3073.8\\
64000	2902.7\\
66000	3063.9\\
68000	3031.1\\
70000	3062.5\\
72000	3128.5\\
74000	3000.5\\
76000	3023.1\\
78000	3001\\
80000	3011.8\\
84000	3073.6\\
86000	2972.9\\
88000	2985.4\\
90000	2985.4\\
92000	3145.1\\
94000	2878.6\\
96000	2983.7\\
98000	3042.2\\
1e+05	3018\\
};
\addplot [color=mycolor39, forget plot]
  table[row sep=crcr]{%
0	4395.3\\
2000	10790\\
4000	11536\\
6000	11837\\
8000	11592\\
10000	11726\\
14000	11611\\
18000	12053\\
20000	11590\\
22000	11950\\
24000	11826\\
26000	12062\\
28000	11479\\
30000	11594\\
32000	11578\\
34000	11403\\
36000	12160\\
38000	11795\\
40000	11674\\
42000	11685\\
44000	11595\\
46000	11879\\
48000	11479\\
50000	11594\\
52000	11966\\
54000	11499\\
56000	11483\\
58000	11870\\
60000	12068\\
62000	11828\\
64000	11386\\
66000	11956\\
70000	11944\\
74000	11663\\
76000	12063\\
78000	12064\\
80000	11547\\
82000	11440\\
84000	11636\\
86000	11759\\
88000	11742\\
90000	12123\\
92000	11939\\
94000	11100\\
96000	11564\\
98000	11465\\
1e+05	11804\\
};
\addplot [color=mycolor40, forget plot]
  table[row sep=crcr]{%
0	4395.3\\
2000	5620.1\\
4000	5851.6\\
6000	5906.4\\
8000	5777.4\\
10000	5717.4\\
12000	5742\\
14000	5793.5\\
16000	5742.3\\
18000	5939.5\\
20000	5672.3\\
22000	5663.4\\
24000	5699.9\\
26000	5607.5\\
28000	5779.2\\
30000	5835.5\\
32000	5694\\
34000	5737.5\\
36000	5685.5\\
38000	5696.4\\
40000	5785.8\\
44000	5746.8\\
46000	5828.2\\
48000	5719.3\\
50000	5821.8\\
52000	5739.1\\
54000	5849.2\\
56000	5770.1\\
58000	5656.4\\
60000	5740.8\\
62000	5862.8\\
64000	5716.9\\
66000	5740.1\\
70000	5838.2\\
72000	5773.8\\
74000	5685.3\\
78000	5651.2\\
80000	5733.1\\
82000	5755.7\\
86000	5699.5\\
88000	5676.3\\
90000	5760.2\\
92000	5591.2\\
94000	5827.3\\
96000	5786.6\\
98000	5694.1\\
1e+05	5817\\
};
\addplot [color=mycolor41, forget plot]
  table[row sep=crcr]{%
0	4395.3\\
2000	5214.7\\
4000	5667.3\\
6000	5379.6\\
8000	5328.6\\
10000	5388.1\\
12000	5515.8\\
14000	5432.8\\
16000	5279.3\\
18000	5373.6\\
20000	5423.3\\
22000	5452.8\\
24000	5327.3\\
26000	5382.8\\
28000	5393.4\\
30000	5377.9\\
32000	5322.7\\
34000	5229.5\\
36000	5310\\
38000	5459.8\\
40000	5466.1\\
42000	5343.5\\
44000	5465.5\\
48000	5359\\
50000	5487.5\\
52000	5720.2\\
54000	5398.5\\
56000	5333.8\\
58000	5290\\
60000	5431.1\\
62000	5334.8\\
64000	5292.6\\
66000	5274.2\\
68000	5493.5\\
70000	5454.8\\
72000	5574.9\\
74000	5430.5\\
76000	5419.5\\
78000	5445.7\\
80000	5549.4\\
82000	5370\\
84000	5512.4\\
86000	5502.5\\
88000	5317.9\\
92000	5472.4\\
94000	5504.3\\
96000	5362.6\\
98000	5112.3\\
1e+05	5353.5\\
};
\addplot [color=mycolor42, forget plot]
  table[row sep=crcr]{%
0	4395.3\\
2000	15830\\
4000	19513\\
6000	19844\\
8000	20575\\
10000	20367\\
12000	20089\\
14000	20311\\
16000	20066\\
18000	20697\\
20000	20746\\
22000	20162\\
24000	20774\\
26000	20698\\
28000	20444\\
30000	21414\\
32000	19788\\
34000	20251\\
36000	20629\\
38000	20921\\
40000	20339\\
42000	20095\\
44000	20779\\
46000	20178\\
48000	20576\\
50000	19858\\
52000	20698\\
54000	20882\\
56000	20441\\
58000	19320\\
60000	20783\\
62000	20378\\
64000	20221\\
66000	20186\\
68000	19964\\
70000	20382\\
72000	21017\\
74000	20841\\
76000	20903\\
78000	21061\\
80000	19953\\
82000	20430\\
84000	20104\\
86000	20704\\
88000	20352\\
90000	20161\\
92000	20747\\
94000	19935\\
98000	20185\\
1e+05	20437\\
};
\addplot [color=mycolor43, forget plot]
  table[row sep=crcr]{%
0	4395.3\\
2000	4947.1\\
6000	4898.4\\
8000	4944.6\\
10000	4876.9\\
12000	4921.5\\
14000	4987.7\\
16000	4988.5\\
18000	4960.3\\
20000	4995\\
22000	4966.1\\
24000	4909.8\\
26000	4941.9\\
28000	4848.6\\
30000	5029.7\\
32000	4942.7\\
34000	4890.5\\
36000	5059.7\\
38000	4986.8\\
40000	4966.3\\
42000	4849.6\\
44000	4952.5\\
46000	4932.7\\
48000	4855.4\\
50000	4972.2\\
52000	5043.2\\
54000	5025.4\\
56000	4934.4\\
58000	5002.6\\
60000	5016.3\\
62000	4985.6\\
64000	5063.3\\
66000	5039.4\\
68000	4890.3\\
70000	5015.2\\
72000	4895.3\\
74000	5027.3\\
76000	4971.4\\
80000	4898.6\\
82000	4966.1\\
84000	4899.8\\
86000	5089\\
88000	5028.9\\
90000	4924.3\\
92000	5029\\
96000	4950.5\\
98000	4867.4\\
1e+05	4915.7\\
};
\addplot [color=mycolor44, forget plot]
  table[row sep=crcr]{%
0	4395.3\\
2000	3563.9\\
4000	3401.4\\
6000	3278.8\\
8000	3377.4\\
10000	3356.8\\
14000	3353.2\\
16000	3360.8\\
18000	3273.8\\
20000	3257\\
22000	3439.7\\
24000	3318.7\\
26000	3331.7\\
28000	3358.4\\
30000	3349.1\\
32000	3415.2\\
34000	3371.4\\
36000	3428.5\\
38000	3360\\
40000	3271.4\\
42000	3304.9\\
46000	3413.5\\
48000	3361\\
52000	3396.3\\
54000	3403\\
56000	3374\\
58000	3359.8\\
60000	3369.8\\
62000	3320.6\\
66000	3381.4\\
68000	3376.9\\
70000	3440.4\\
72000	3459.6\\
74000	3379.4\\
76000	3313.8\\
78000	3449.9\\
80000	3398.1\\
82000	3460.7\\
84000	3266.2\\
86000	3383.4\\
88000	3256\\
90000	3314.5\\
92000	3276.7\\
94000	3285.4\\
96000	3329\\
98000	3324.6\\
1e+05	3410.2\\
};
\addplot [color=mycolor45, forget plot]
  table[row sep=crcr]{%
0	4395.3\\
2000	6790\\
4000	6898\\
6000	6958\\
8000	6864.9\\
10000	6973.9\\
12000	6931.5\\
14000	7015.2\\
16000	6790.4\\
18000	6999\\
20000	6972.2\\
22000	7049.9\\
24000	6942.3\\
26000	6989.1\\
28000	7097\\
30000	6868.4\\
32000	7048.8\\
34000	6868.2\\
36000	6872.9\\
38000	7004.5\\
40000	6909.8\\
44000	6890.1\\
46000	6973.4\\
48000	7151.5\\
50000	6948.2\\
52000	6797.7\\
54000	6885.8\\
56000	7000.4\\
58000	6942.1\\
60000	6955.9\\
62000	6812\\
64000	6959.4\\
66000	6725\\
68000	6876.1\\
70000	6890.6\\
72000	6943.8\\
74000	6949.1\\
76000	6890.2\\
80000	6935.5\\
82000	6815.7\\
84000	6907.4\\
86000	6963.6\\
88000	6794.3\\
90000	7013.5\\
92000	7022.9\\
94000	7061.9\\
96000	6946.4\\
98000	7051.2\\
1e+05	6916.6\\
};
\addplot [color=mycolor46, forget plot]
  table[row sep=crcr]{%
0	4395.3\\
4000	4069.6\\
6000	4166.8\\
8000	4168\\
10000	4015.4\\
12000	4126.9\\
14000	4100.6\\
16000	4058.4\\
18000	4228.8\\
20000	4196.3\\
22000	4120.9\\
24000	4005\\
26000	3967.7\\
30000	4164.2\\
32000	4185.4\\
34000	4140.1\\
36000	4147.8\\
38000	4003.4\\
40000	4012.9\\
42000	4192.3\\
44000	3981.3\\
46000	4016.3\\
48000	4174\\
52000	4089.6\\
54000	4141\\
56000	4124\\
58000	3996.8\\
60000	4126.9\\
62000	4091.4\\
64000	4189\\
66000	4065.1\\
68000	4232.2\\
70000	4147\\
72000	4126.3\\
74000	4008.3\\
76000	4092.9\\
78000	4048.8\\
80000	4051.7\\
82000	4169.9\\
84000	4110.2\\
86000	4031.8\\
88000	4188.9\\
90000	4222.5\\
92000	4069.9\\
94000	4133.3\\
96000	3998.7\\
98000	4104.1\\
1e+05	4153.8\\
};
\addplot [color=mycolor47, forget plot]
  table[row sep=crcr]{%
0	4395.3\\
2000	5226.7\\
4000	5225.2\\
6000	5370.1\\
8000	5166.5\\
10000	5125.2\\
12000	5367.4\\
14000	5437.8\\
16000	5352\\
18000	5354.8\\
20000	5330.5\\
22000	5287.7\\
24000	5135.6\\
26000	5325.1\\
28000	5236.6\\
30000	5185.3\\
32000	5475.7\\
34000	5306.9\\
36000	5265\\
38000	5271\\
40000	5346.8\\
42000	5142.5\\
44000	5263.4\\
46000	5261.6\\
48000	5339.2\\
50000	5398.2\\
52000	5387.4\\
54000	5229.1\\
56000	5270\\
58000	5396.1\\
62000	5220.1\\
64000	5442.6\\
66000	5299.5\\
68000	5248.4\\
70000	5381.9\\
72000	5430.4\\
76000	5486.3\\
78000	5305.2\\
82000	5393.1\\
84000	5299.7\\
86000	5333\\
88000	5246.8\\
90000	5209.7\\
94000	5464.1\\
96000	5478.5\\
98000	5345.9\\
1e+05	5242.6\\
};
\addplot [color=mycolor48, forget plot]
  table[row sep=crcr]{%
0	4395.3\\
2000	3215.6\\
4000	2966.1\\
6000	2937.1\\
8000	2930.6\\
10000	2910.3\\
12000	2903.3\\
14000	2948.6\\
16000	2879.8\\
18000	2968.8\\
20000	2975.3\\
24000	2927.3\\
26000	2916\\
28000	2943.9\\
30000	2864.1\\
32000	2917.2\\
34000	2953.2\\
36000	2959.3\\
40000	2945.1\\
44000	2905\\
46000	2940.8\\
48000	2922.6\\
50000	2945.4\\
52000	2899.6\\
54000	2950.6\\
56000	2927.3\\
58000	2956.3\\
60000	2901.5\\
62000	2910.4\\
64000	2934.1\\
66000	2948.3\\
68000	2978.8\\
70000	2849.3\\
72000	2948.3\\
74000	2970.3\\
76000	2907.9\\
78000	2892.8\\
80000	2967.5\\
82000	2980.2\\
84000	2899.3\\
86000	2934.5\\
88000	2867.4\\
92000	2932.5\\
94000	2912.9\\
96000	2919.4\\
98000	2875.3\\
1e+05	2922.6\\
};
\addplot [color=mycolor49, forget plot]
  table[row sep=crcr]{%
0	4395.3\\
2000	4489.2\\
4000	4209\\
6000	4216.7\\
8000	4386.7\\
10000	4286.9\\
12000	4526.3\\
14000	4279.1\\
16000	4461.2\\
18000	4251.1\\
20000	4422.8\\
22000	4347.7\\
24000	4546.1\\
26000	4486\\
28000	4546.9\\
30000	4431.3\\
32000	4380.9\\
34000	4345.2\\
36000	4279.4\\
38000	4335.4\\
40000	4232.8\\
42000	4445.9\\
44000	4296.5\\
48000	4373.1\\
50000	4221.4\\
52000	4384\\
54000	4483.1\\
56000	4240.6\\
58000	4332.1\\
60000	4261.2\\
62000	4306.8\\
64000	4429.8\\
68000	4304.3\\
70000	4328\\
72000	4415.8\\
74000	4421.8\\
76000	4390.1\\
78000	4441.5\\
80000	4565.3\\
82000	4345.9\\
84000	4439.8\\
86000	4409.6\\
88000	4397.8\\
90000	4298.3\\
92000	4300.1\\
94000	4384\\
96000	4339.1\\
98000	4327.3\\
1e+05	4418.9\\
};
\addplot [color=mycolor50, forget plot]
  table[row sep=crcr]{%
0	4395.3\\
2000	2729.3\\
4000	2443.2\\
6000	2441\\
8000	2543\\
10000	2528.8\\
12000	2501.2\\
14000	2484.9\\
16000	2453.5\\
18000	2516\\
20000	2481.1\\
22000	2507.2\\
24000	2495.1\\
26000	2457.5\\
28000	2536.5\\
30000	2519.7\\
32000	2421.7\\
34000	2527.7\\
36000	2543.8\\
38000	2510.6\\
40000	2529.1\\
42000	2533.2\\
44000	2452.1\\
46000	2493\\
48000	2492.9\\
50000	2507.4\\
52000	2487.6\\
54000	2458.8\\
56000	2498.9\\
58000	2506.9\\
60000	2524.3\\
62000	2485.7\\
64000	2455.7\\
66000	2540.6\\
68000	2479.6\\
70000	2586.7\\
72000	2459.1\\
74000	2495\\
76000	2519.8\\
78000	2514.2\\
80000	2528.2\\
82000	2483.1\\
84000	2518.7\\
86000	2485.6\\
88000	2412.2\\
90000	2478\\
92000	2483\\
94000	2407.8\\
96000	2466.5\\
98000	2488\\
1e+05	2487.7\\
};
\addplot [color=mycolor51, forget plot]
  table[row sep=crcr]{%
0	4395.3\\
2000	5918.8\\
4000	6142\\
6000	5990.1\\
8000	6028.4\\
10000	5930.1\\
12000	5903.5\\
14000	5972.7\\
16000	6136.5\\
18000	5914.6\\
20000	5911.2\\
22000	6173\\
24000	6009.9\\
26000	5944\\
28000	5691.2\\
30000	5897.5\\
32000	6054\\
34000	6024.1\\
36000	6269\\
40000	5888.4\\
42000	6088.7\\
44000	6178.1\\
46000	5933.3\\
48000	6212.5\\
50000	6059.9\\
52000	6347.3\\
54000	6114.5\\
56000	6011.1\\
58000	5782.8\\
60000	5937\\
62000	5915.1\\
64000	6016.8\\
66000	6027.4\\
68000	6339.3\\
70000	5958.1\\
72000	5734.8\\
74000	5950.5\\
76000	5952\\
78000	6165.2\\
80000	5905.8\\
82000	5957.3\\
84000	6055.5\\
86000	6317.6\\
88000	6038.8\\
90000	5946.9\\
92000	6085.5\\
94000	6025.1\\
96000	6021.5\\
98000	5860.6\\
1e+05	5916.6\\
};
\addplot [color=mycolor52, forget plot]
  table[row sep=crcr]{%
0	4395.3\\
2000	3771.7\\
4000	3541.4\\
6000	3734.3\\
8000	3518.6\\
10000	3522.2\\
12000	3348.4\\
14000	3617.2\\
16000	3467.4\\
18000	3526.2\\
20000	3431.7\\
22000	3505.7\\
24000	3598.9\\
26000	3485.3\\
28000	3462.9\\
30000	3414\\
32000	3542.6\\
34000	3594.8\\
36000	3437.3\\
38000	3543.3\\
40000	3446.6\\
44000	3538.8\\
46000	3649.9\\
48000	3428.8\\
52000	3460.8\\
54000	3400.4\\
56000	3431.7\\
58000	3514.3\\
60000	3333.6\\
62000	3518.1\\
64000	3557.7\\
66000	3539.9\\
68000	3579.9\\
70000	3454.5\\
72000	3521\\
74000	3791.3\\
76000	3520.2\\
78000	3473.8\\
80000	3491.2\\
82000	3553.6\\
84000	3560\\
86000	3476.5\\
88000	3469.6\\
90000	3497.1\\
92000	3450.1\\
94000	3475.6\\
96000	3665.5\\
98000	3402.2\\
1e+05	3653.6\\
};
\addplot [color=mycolor53, forget plot]
  table[row sep=crcr]{%
0	4395.3\\
2000	5690.8\\
4000	5730.5\\
6000	5692.1\\
8000	5602.2\\
10000	5709.7\\
12000	5785.9\\
14000	5691.5\\
16000	5950.8\\
18000	5755.1\\
20000	5838.7\\
22000	5652.2\\
24000	5615.2\\
26000	5668.6\\
28000	5635.8\\
30000	5732.2\\
32000	5859.7\\
34000	5622.6\\
36000	5758\\
38000	5606.8\\
40000	5829.1\\
42000	5771.3\\
44000	5423.9\\
46000	5580.7\\
48000	5557.3\\
50000	6034.3\\
52000	5495.3\\
54000	5695\\
56000	5563.6\\
58000	5657.3\\
60000	5642.7\\
62000	5849.1\\
64000	5632\\
66000	5662.5\\
68000	5611.5\\
70000	5539.7\\
72000	5612.8\\
74000	5582.8\\
76000	5704.1\\
78000	5801.8\\
80000	5825.2\\
82000	5599.2\\
84000	5612.8\\
86000	5848.4\\
88000	5612.5\\
90000	5599.8\\
92000	5751.7\\
94000	5572.3\\
96000	5612.5\\
98000	5702.7\\
1e+05	5750\\
};
\addplot [color=mycolor54, forget plot]
  table[row sep=crcr]{%
0	4395.3\\
2000	4630.7\\
6000	4638.4\\
8000	4515.9\\
10000	4692.1\\
12000	4641.5\\
14000	4703.3\\
16000	4553.7\\
18000	4527.6\\
20000	4597.3\\
22000	4566.7\\
24000	4655.6\\
26000	4612.7\\
28000	4652.1\\
30000	4638.7\\
34000	4496\\
36000	4589.9\\
38000	4605.8\\
40000	4544.4\\
42000	4700.6\\
44000	4672.2\\
46000	4532.9\\
48000	4607.9\\
50000	4657.3\\
52000	4666.4\\
54000	4652.3\\
58000	4708.3\\
60000	4646.5\\
62000	4570.5\\
64000	4552.5\\
66000	4741.7\\
68000	4784.4\\
70000	4656.4\\
72000	4572.8\\
76000	4671.8\\
78000	4485.4\\
80000	4501\\
82000	4754.9\\
84000	4619.9\\
86000	4392.1\\
88000	4734.9\\
90000	4456.1\\
92000	4507.5\\
96000	4645.1\\
98000	4507.5\\
1e+05	4723.4\\
};
\addplot [color=mycolor55, forget plot]
  table[row sep=crcr]{%
0	4395.3\\
2000	5534.3\\
4000	5925.1\\
6000	5808.2\\
8000	5777.1\\
12000	5858.4\\
14000	5833.3\\
16000	5772.1\\
18000	5829.8\\
20000	5834.1\\
22000	5887.3\\
24000	5818.4\\
26000	5806.8\\
28000	5724.3\\
30000	5868.4\\
32000	5857.9\\
34000	5806.6\\
36000	5717.9\\
38000	5891.3\\
40000	5775.2\\
42000	5820.6\\
46000	5805.9\\
48000	5868.4\\
50000	5868.6\\
52000	5749.6\\
54000	5825.3\\
56000	5808.7\\
58000	5708.9\\
62000	5837.6\\
64000	5828.3\\
66000	5677.4\\
68000	5816.8\\
70000	5677.7\\
74000	5834.2\\
78000	5768.6\\
80000	5719.1\\
82000	5703.6\\
86000	5912.7\\
88000	5744.2\\
90000	5932.4\\
92000	5838\\
94000	5723.8\\
96000	5713\\
98000	5877.5\\
1e+05	5727.5\\
};
\addplot [color=mycolor56, forget plot]
  table[row sep=crcr]{%
0	4395.3\\
2000	4119.9\\
4000	4004.7\\
6000	3979.9\\
8000	3936.5\\
10000	3986.6\\
12000	3953.6\\
14000	3948.5\\
16000	4041.5\\
20000	4105.2\\
22000	4019.1\\
24000	3962\\
26000	4083\\
28000	4020.8\\
30000	4002\\
32000	4071.4\\
34000	3942.9\\
36000	4002.3\\
38000	3943.1\\
40000	3925.5\\
42000	3961.2\\
44000	4037.2\\
46000	3938.8\\
48000	3989.7\\
50000	4058.2\\
52000	3987.8\\
54000	3957.1\\
56000	3862.6\\
58000	3979.4\\
60000	3988\\
62000	4083\\
64000	4114.3\\
66000	4074.1\\
68000	4016\\
70000	4026.4\\
72000	3962.5\\
74000	4042.8\\
76000	3973.2\\
80000	3980\\
82000	3975.2\\
84000	4021\\
86000	3995.5\\
88000	4059.2\\
90000	4078.8\\
92000	4061.7\\
94000	3924\\
96000	4052.1\\
98000	4034.7\\
1e+05	3972.9\\
};
\addplot [color=mycolor57, forget plot]
  table[row sep=crcr]{%
0	4395.3\\
2000	6156.2\\
4000	6868.9\\
6000	6715.8\\
8000	6747.7\\
10000	6670.3\\
12000	6450\\
14000	6442.6\\
16000	6526.7\\
18000	6592.7\\
20000	6617.1\\
22000	6779.2\\
24000	6814.8\\
26000	6708.6\\
28000	6582\\
30000	6572.2\\
32000	6647\\
34000	6550.6\\
36000	6566.1\\
38000	6619.5\\
40000	6802.8\\
42000	6635.9\\
44000	6698\\
46000	6603.6\\
48000	6745.2\\
50000	6619.9\\
52000	6585.9\\
54000	6361.8\\
56000	6613.3\\
58000	6773.5\\
60000	6671\\
62000	6626.2\\
64000	6628.6\\
66000	6712.5\\
68000	6690.3\\
70000	6763.7\\
72000	6701.1\\
74000	6785.4\\
76000	6813.8\\
78000	6673.5\\
80000	6716.8\\
82000	6812.7\\
84000	6854.3\\
86000	6619.9\\
88000	6602.2\\
90000	6757.9\\
94000	6651.8\\
96000	6659.9\\
98000	6582.8\\
1e+05	6749.4\\
};
\addplot [color=mycolor58, forget plot]
  table[row sep=crcr]{%
0	4395.3\\
2000	3897.2\\
4000	3604.6\\
6000	3561.8\\
8000	3537.5\\
10000	3594.8\\
12000	3599\\
14000	3538.8\\
16000	3601.9\\
18000	3631.7\\
22000	3533.3\\
24000	3643.8\\
28000	3501.5\\
30000	3498.6\\
32000	3632.5\\
34000	3573\\
36000	3486.3\\
38000	3617.5\\
40000	3663.5\\
42000	3559.5\\
44000	3528.2\\
46000	3527.5\\
50000	3598.7\\
52000	3492\\
54000	3518.8\\
56000	3582.1\\
58000	3533.1\\
60000	3576.6\\
62000	3596.8\\
68000	3478\\
70000	3555.4\\
72000	3583.4\\
74000	3514.7\\
78000	3576.2\\
80000	3577.2\\
82000	3639\\
84000	3474.1\\
86000	3629.1\\
88000	3582.5\\
90000	3590.2\\
92000	3541.2\\
96000	3625\\
98000	3512.9\\
1e+05	3477.7\\
};
\addplot [color=mycolor59, forget plot]
  table[row sep=crcr]{%
0	4395.3\\
2000	5009.3\\
4000	5113.8\\
6000	5052.7\\
8000	5232.1\\
12000	5104.2\\
14000	5250.5\\
16000	5031.1\\
18000	4893.8\\
20000	5014.7\\
22000	4944.4\\
24000	5035.9\\
26000	4960.9\\
28000	5168.7\\
30000	5152.3\\
32000	5028.2\\
34000	4922.6\\
36000	5049.9\\
38000	5046.8\\
40000	5145.3\\
42000	5007.1\\
44000	5042.4\\
46000	4923.4\\
48000	5068.2\\
50000	5111.8\\
52000	4957.7\\
54000	5072.8\\
56000	5003.8\\
58000	4960.1\\
60000	5101\\
62000	5047.4\\
64000	5057.5\\
66000	5127.5\\
68000	5128.8\\
70000	4965.6\\
72000	5032\\
74000	4981.1\\
76000	5035.1\\
78000	4973.1\\
80000	5005.5\\
82000	5129.2\\
84000	4983.1\\
86000	5081.3\\
90000	5070.5\\
92000	4982.8\\
94000	5189.1\\
96000	4945.2\\
98000	5160.7\\
1e+05	5143.7\\
};
\addplot [color=mycolor60, forget plot]
  table[row sep=crcr]{%
0	4395.3\\
2000	3705.8\\
4000	3575\\
6000	3482.9\\
8000	3566.2\\
10000	3562\\
12000	3533.9\\
14000	3484.9\\
16000	3574\\
18000	3616.7\\
20000	3475.3\\
22000	3566.3\\
24000	3547.2\\
26000	3584.7\\
28000	3525.1\\
30000	3509.8\\
32000	3540.2\\
34000	3590.9\\
36000	3597.8\\
38000	3626.7\\
40000	3595.3\\
42000	3599.3\\
44000	3537.9\\
46000	3580.7\\
48000	3457.6\\
50000	3576.3\\
52000	3546.2\\
54000	3566.5\\
56000	3546.5\\
58000	3476.4\\
60000	3564.8\\
64000	3511.1\\
66000	3634.7\\
68000	3513.1\\
72000	3517.6\\
74000	3542.8\\
76000	3552.5\\
78000	3532.1\\
82000	3575.8\\
84000	3536.5\\
86000	3538.7\\
88000	3524.1\\
90000	3529\\
92000	3504.1\\
96000	3601.8\\
98000	3607.3\\
1e+05	3553.9\\
};
\addplot [color=mycolor61, forget plot]
  table[row sep=crcr]{%
0	4395.3\\
2000	3413.3\\
4000	3160.8\\
6000	3086\\
8000	3112.1\\
10000	3191.5\\
12000	3099.4\\
14000	3153.4\\
16000	3116.2\\
18000	3160.8\\
22000	3112.1\\
24000	3123.9\\
26000	3154.3\\
28000	3136.8\\
32000	3083.9\\
34000	3141.5\\
36000	3186.7\\
38000	3115\\
40000	3118.9\\
42000	3195.1\\
44000	3119.3\\
46000	3121.5\\
48000	3092.8\\
50000	3170.2\\
52000	3101.2\\
56000	3149\\
60000	3140.3\\
62000	3106\\
64000	3143.1\\
66000	3161.8\\
68000	3142.6\\
70000	3147.5\\
72000	3138.2\\
74000	3143.1\\
76000	3127.3\\
78000	3094.5\\
82000	3145.1\\
84000	3113.5\\
86000	3158.3\\
88000	3135.8\\
90000	3154.6\\
92000	3118\\
98000	3090.1\\
1e+05	3116.9\\
};
\addplot [color=mycolor62, forget plot]
  table[row sep=crcr]{%
0	4395.3\\
2000	4837.4\\
4000	4950.8\\
6000	4882.6\\
8000	4977.8\\
10000	5046.4\\
12000	5031.7\\
14000	5032.6\\
16000	5105.7\\
18000	5003.7\\
20000	4966.6\\
22000	5062.1\\
24000	4963.3\\
26000	4922.8\\
28000	5032.9\\
30000	5069.4\\
32000	4981.6\\
34000	4970.7\\
36000	5020.8\\
38000	4946.5\\
40000	5114.8\\
44000	5030.7\\
46000	4887\\
48000	5067.5\\
50000	5069.9\\
54000	4947.5\\
56000	5056\\
58000	5049.6\\
60000	5058.6\\
62000	5048.9\\
64000	5006\\
66000	4995\\
68000	4903.3\\
70000	4936.4\\
72000	4869\\
74000	4896.3\\
76000	5016.6\\
78000	5008\\
80000	5057.3\\
82000	4968.2\\
86000	4953\\
88000	4961.4\\
90000	4990.4\\
92000	4979.4\\
94000	5012.7\\
96000	4948.5\\
98000	4959.5\\
1e+05	4943.3\\
};
\addplot [color=mycolor63, forget plot]
  table[row sep=crcr]{%
0	4395.3\\
2000	2784.1\\
4000	2474.8\\
6000	2443.1\\
8000	2434.3\\
10000	2406.8\\
12000	2413\\
14000	2392.4\\
16000	2381.7\\
18000	2401.1\\
20000	2442.8\\
22000	2435\\
24000	2434.8\\
26000	2415.5\\
30000	2408\\
32000	2382.1\\
34000	2401.4\\
36000	2434.7\\
38000	2504.5\\
40000	2451.1\\
44000	2404.1\\
46000	2373.4\\
48000	2427.2\\
50000	2430.4\\
52000	2463.2\\
54000	2420.5\\
56000	2413.7\\
58000	2450.8\\
60000	2455.3\\
62000	2446.9\\
64000	2425.6\\
66000	2411.9\\
68000	2433.7\\
70000	2404.3\\
72000	2434.7\\
74000	2435.2\\
76000	2385.2\\
78000	2415.3\\
80000	2385.2\\
82000	2455.3\\
86000	2386.3\\
88000	2400.6\\
92000	2443\\
94000	2457.4\\
96000	2407.6\\
98000	2366.1\\
1e+05	2398.3\\
};
\addplot [color=mycolor64, forget plot]
  table[row sep=crcr]{%
0	4395.3\\
2000	3329.1\\
4000	3080.3\\
6000	3075.2\\
8000	3104.2\\
10000	3041.3\\
14000	3094.6\\
16000	3108.3\\
18000	3079.4\\
20000	3102.6\\
22000	3069.5\\
24000	3089\\
26000	3068\\
28000	3059.4\\
30000	3095.9\\
32000	3029.2\\
34000	3095.6\\
36000	3082.9\\
38000	3035.4\\
40000	3047.4\\
42000	3083.1\\
44000	3047.8\\
48000	3065.3\\
50000	3084.4\\
52000	3032.4\\
54000	3073.7\\
56000	3090.9\\
58000	3071.7\\
60000	3016\\
64000	3067.4\\
66000	3060.8\\
68000	3069.9\\
70000	3045.7\\
72000	3088.4\\
74000	3047.1\\
76000	3066\\
78000	3103.2\\
82000	3054.8\\
84000	3080.5\\
90000	3040.6\\
94000	3089.6\\
96000	3102.7\\
98000	3066.9\\
1e+05	3053.2\\
};
\addplot [color=mycolor65, forget plot]
  table[row sep=crcr]{%
0	4395.3\\
2000	2882.4\\
4000	2617.1\\
6000	2622.4\\
8000	2591\\
10000	2546.8\\
12000	2568.4\\
14000	2529\\
16000	2558\\
18000	2547.5\\
22000	2554.2\\
24000	2608.2\\
26000	2569.5\\
28000	2591.5\\
30000	2530\\
32000	2549.8\\
34000	2539.5\\
38000	2603.7\\
40000	2554\\
42000	2575.5\\
44000	2535\\
46000	2543.2\\
48000	2526.3\\
50000	2531\\
52000	2618.6\\
54000	2550.8\\
56000	2564.1\\
58000	2563.3\\
60000	2519.7\\
62000	2569.5\\
64000	2544.7\\
68000	2515.6\\
70000	2569.4\\
72000	2611.3\\
74000	2546.3\\
76000	2584.2\\
78000	2581.4\\
80000	2534.2\\
82000	2530\\
86000	2575.7\\
88000	2546.3\\
90000	2533.7\\
92000	2618.8\\
94000	2481.7\\
96000	2596.6\\
98000	2559.7\\
1e+05	2568.7\\
};
\addplot [color=mycolor66, forget plot]
  table[row sep=crcr]{%
0	4395.3\\
2000	3628.4\\
4000	3428\\
6000	3452\\
8000	3448.7\\
10000	3482\\
12000	3458\\
14000	3404.5\\
16000	3447.8\\
18000	3448.5\\
20000	3420.8\\
22000	3458.6\\
24000	3450.6\\
26000	3454.7\\
28000	3439.2\\
30000	3402.6\\
32000	3399.1\\
34000	3422\\
36000	3422.5\\
38000	3455\\
40000	3455.1\\
42000	3404.5\\
44000	3443\\
46000	3439.2\\
50000	3447.1\\
52000	3418.5\\
54000	3400.3\\
56000	3462.1\\
58000	3468.7\\
60000	3455.6\\
64000	3493.4\\
66000	3420.6\\
68000	3453.5\\
70000	3440\\
72000	3416.3\\
74000	3429.9\\
76000	3407.2\\
78000	3419\\
80000	3407.2\\
82000	3439.4\\
84000	3438\\
86000	3463.3\\
88000	3436.5\\
90000	3356.1\\
92000	3364.9\\
94000	3480.2\\
96000	3484.6\\
98000	3451.6\\
1e+05	3451.9\\
};
\addplot [color=mycolor67, forget plot]
  table[row sep=crcr]{%
0	4395.3\\
2000	3509.4\\
4000	3256.7\\
6000	3226.3\\
8000	3145.1\\
10000	3176.9\\
12000	3281.3\\
14000	3183.7\\
16000	3296\\
18000	3189.1\\
20000	3146.6\\
22000	3146.3\\
24000	3186.3\\
26000	3145.1\\
30000	3237.2\\
34000	3187\\
36000	3144.7\\
40000	3184.6\\
42000	3176.7\\
44000	3205.5\\
46000	3193.2\\
48000	3112.6\\
50000	3200.3\\
52000	3206.6\\
54000	3134.3\\
56000	3149.9\\
58000	3127.6\\
60000	3217.4\\
62000	3211.9\\
64000	3165.6\\
66000	3231.8\\
68000	3198.1\\
70000	3178.6\\
72000	3218.8\\
74000	3218\\
76000	3236.1\\
78000	3156.4\\
80000	3184.1\\
82000	3221.3\\
84000	3157.5\\
86000	3136.3\\
88000	3162.2\\
90000	3177.9\\
92000	3250.8\\
94000	3157.7\\
96000	3155.6\\
98000	3205.7\\
1e+05	3148.4\\
};
\addplot [color=mycolor67]
  table[row sep=crcr]{%
0	4395.3\\
2000	2961.5\\
4000	2720.5\\
6000	2686.8\\
10000	2650\\
12000	2656.7\\
14000	2598.1\\
18000	2648.7\\
20000	2641.8\\
22000	2609.8\\
24000	2645.3\\
26000	2638.3\\
28000	2680.4\\
32000	2610.1\\
34000	2609.9\\
36000	2656.6\\
38000	2679.9\\
40000	2636.6\\
44000	2662.3\\
50000	2647.6\\
52000	2604.7\\
54000	2658.9\\
56000	2653.8\\
58000	2617.3\\
60000	2649.2\\
62000	2704.7\\
64000	2646.6\\
66000	2636.6\\
68000	2674.3\\
70000	2661.1\\
74000	2622.6\\
78000	2676.8\\
80000	2650.8\\
82000	2649\\
84000	2655.5\\
86000	2613.9\\
88000	2642.4\\
90000	2637.3\\
92000	2607.7\\
96000	2661.6\\
98000	2638.5\\
1e+05	2661.8\\
};
\addlegendentry{\ClasskIIILegend}

\addplot [color=mycolor68, forget plot]
  table[row sep=crcr]{%
0	4395.3\\
2000	2817.8\\
4000	2447.2\\
6000	2468.9\\
8000	2520.1\\
10000	2474.2\\
12000	2451.4\\
14000	2449.1\\
16000	2499\\
18000	2486.6\\
20000	2432.3\\
22000	2491.8\\
24000	2440.6\\
26000	2474\\
28000	2474.3\\
30000	2454.2\\
32000	2487.9\\
34000	2479.9\\
38000	2440.6\\
40000	2485.3\\
42000	2454\\
44000	2456.9\\
46000	2490.6\\
48000	2464.9\\
50000	2459\\
52000	2510.9\\
54000	2472.9\\
56000	2491.9\\
58000	2454\\
60000	2493.9\\
62000	2488.7\\
64000	2446.5\\
66000	2500.9\\
68000	2515.4\\
70000	2486.6\\
72000	2482.9\\
74000	2455.4\\
76000	2495.1\\
78000	2490.5\\
80000	2477.3\\
82000	2475.7\\
84000	2496.3\\
86000	2453\\
88000	2483.4\\
90000	2441\\
92000	2412.7\\
94000	2448.6\\
96000	2454.1\\
98000	2495\\
1e+05	2493.3\\
};
\addplot [color=mycolor69, forget plot]
  table[row sep=crcr]{%
0	4395.3\\
2000	3234.7\\
4000	2919.4\\
6000	2895.9\\
8000	2900.6\\
10000	2940.8\\
12000	2953.6\\
14000	2943.8\\
16000	2921.8\\
18000	2951.6\\
22000	2888.9\\
24000	2969.4\\
26000	2949.4\\
28000	2894\\
30000	2931.7\\
32000	2945.9\\
34000	2895.5\\
36000	2913.2\\
38000	2961.4\\
40000	2908.1\\
42000	2943.4\\
46000	2880.6\\
48000	2934.1\\
50000	2903.2\\
52000	2932\\
54000	2941\\
56000	2915.8\\
58000	2922.9\\
60000	2919.4\\
62000	2905.8\\
64000	2926.6\\
66000	2926\\
68000	2910.2\\
70000	2917.9\\
72000	2916.2\\
74000	2902.4\\
76000	2915.6\\
78000	2831.9\\
80000	2938.7\\
84000	2889.4\\
86000	2891.5\\
88000	2928.6\\
90000	2945.4\\
92000	2924.1\\
94000	2927\\
96000	2880.5\\
98000	2868.7\\
1e+05	2915.6\\
};
\addplot [color=mycolor70, forget plot]
  table[row sep=crcr]{%
0	4395.3\\
2000	2603\\
4000	2126.6\\
8000	2101\\
10000	2099.5\\
12000	2079.6\\
14000	2120.4\\
16000	2128.7\\
18000	2105.8\\
20000	2067.7\\
22000	2115.8\\
24000	2105.8\\
26000	2142.3\\
28000	2123.1\\
30000	2144.3\\
32000	2116.6\\
34000	2152.2\\
36000	2152.5\\
38000	2098.4\\
40000	2135.4\\
42000	2111\\
44000	2150.3\\
46000	2119.3\\
50000	2073.6\\
52000	2134.7\\
54000	2151.5\\
56000	2132.2\\
58000	2144.2\\
60000	2117.4\\
62000	2145.2\\
66000	2093.4\\
68000	2089\\
70000	2140.7\\
74000	2147.2\\
76000	2139.5\\
78000	2149\\
80000	2098\\
84000	2128.7\\
86000	2111.3\\
88000	2123.6\\
90000	2054.6\\
92000	2086.7\\
94000	2104.3\\
96000	2107.4\\
98000	2152.5\\
1e+05	2129.4\\
};
\addplot [color=mycolor71, forget plot]
  table[row sep=crcr]{%
0	4395.3\\
2000	2873.7\\
4000	2503.9\\
6000	2468.3\\
8000	2468.1\\
10000	2492\\
12000	2523.6\\
14000	2521.8\\
16000	2464.9\\
18000	2447.9\\
20000	2482.8\\
22000	2470.6\\
26000	2464.5\\
28000	2502.6\\
30000	2478.4\\
34000	2508.8\\
36000	2468.6\\
38000	2467.2\\
42000	2529.9\\
46000	2502.9\\
50000	2503.8\\
52000	2514.3\\
54000	2500.4\\
56000	2502\\
60000	2460.1\\
62000	2509\\
66000	2476.7\\
68000	2491.7\\
70000	2472.2\\
72000	2492.8\\
74000	2492.9\\
76000	2481.3\\
80000	2482.1\\
82000	2522.7\\
84000	2500.5\\
92000	2444.7\\
94000	2494.9\\
98000	2511.2\\
1e+05	2514.6\\
};
\addplot [color=mycolor72, forget plot]
  table[row sep=crcr]{%
0	4395.3\\
2000	2607.4\\
4000	2268.7\\
6000	2152\\
8000	2207.8\\
10000	2193.8\\
12000	2189.9\\
14000	2204.1\\
16000	2197.5\\
18000	2213.3\\
20000	2258.5\\
22000	2239.4\\
28000	2213.4\\
30000	2210.2\\
32000	2245.4\\
34000	2228\\
36000	2161.7\\
40000	2201.9\\
42000	2262.4\\
44000	2230.5\\
46000	2217.6\\
50000	2215.8\\
52000	2163.1\\
54000	2180.8\\
56000	2143\\
58000	2172.1\\
62000	2144.7\\
64000	2214.4\\
66000	2251.9\\
68000	2276.3\\
70000	2228.3\\
72000	2161.6\\
74000	2209.8\\
76000	2168.2\\
82000	2215.1\\
84000	2202.5\\
86000	2262.3\\
88000	2162.9\\
90000	2225.1\\
92000	2249.4\\
94000	2189.3\\
96000	2193.9\\
98000	2164.3\\
1e+05	2242.2\\
};
\addplot [color=mycolor73, forget plot]
  table[row sep=crcr]{%
0	4395.3\\
2000	2941.8\\
4000	2617.1\\
6000	2655.7\\
10000	2633.5\\
12000	2636.1\\
16000	2575.9\\
18000	2584.5\\
20000	2642.9\\
22000	2581.3\\
26000	2572.6\\
28000	2610\\
30000	2565.6\\
32000	2599.5\\
34000	2624\\
36000	2664\\
38000	2614.8\\
40000	2604.4\\
42000	2647.4\\
46000	2580.8\\
48000	2617.8\\
50000	2592.7\\
54000	2642.7\\
56000	2607.1\\
58000	2635.8\\
60000	2614.5\\
62000	2578.1\\
64000	2609.7\\
66000	2650.7\\
70000	2587.4\\
72000	2574.9\\
74000	2585.3\\
76000	2604.5\\
78000	2581.3\\
80000	2606.6\\
82000	2582.3\\
84000	2636.5\\
86000	2647.3\\
88000	2571.3\\
90000	2605.3\\
92000	2593.3\\
94000	2652.9\\
96000	2634.3\\
1e+05	2665.2\\
};
\addplot [color=mycolor74, forget plot]
  table[row sep=crcr]{%
0	4395.3\\
2000	2067.8\\
4000	1579.2\\
6000	1558.4\\
8000	1567.9\\
10000	1548.7\\
12000	1524.4\\
14000	1541.6\\
16000	1535\\
20000	1555\\
22000	1556.8\\
26000	1535\\
28000	1543.2\\
30000	1544\\
32000	1563\\
36000	1542.5\\
40000	1551.9\\
42000	1557.4\\
44000	1538.9\\
46000	1535.9\\
48000	1524.6\\
50000	1526.9\\
52000	1557\\
54000	1544.2\\
56000	1554.4\\
58000	1544\\
60000	1555.4\\
62000	1535.4\\
64000	1560.4\\
66000	1546.3\\
68000	1558.9\\
70000	1559.3\\
72000	1539.2\\
74000	1543\\
76000	1531.3\\
78000	1548.2\\
80000	1520.9\\
82000	1548.1\\
84000	1545\\
88000	1550.9\\
90000	1569.4\\
92000	1540.5\\
94000	1551.3\\
96000	1539.1\\
98000	1521.8\\
1e+05	1541.2\\
};
\addplot [color=mycolor75, forget plot]
  table[row sep=crcr]{%
0	4395.3\\
2000	2668.8\\
4000	2236.7\\
8000	2246.7\\
10000	2242.7\\
14000	2202.6\\
16000	2226.9\\
20000	2211.1\\
22000	2228\\
24000	2217.1\\
26000	2231.2\\
28000	2209.4\\
30000	2226.3\\
32000	2222.1\\
34000	2262\\
36000	2252.2\\
38000	2220.1\\
40000	2224.2\\
42000	2209.8\\
44000	2246.8\\
46000	2268.8\\
50000	2235.2\\
52000	2199.1\\
54000	2229.8\\
56000	2215\\
58000	2236.5\\
62000	2238.2\\
64000	2223.3\\
66000	2230\\
68000	2209.6\\
70000	2221.7\\
74000	2233\\
76000	2250.4\\
78000	2223.8\\
80000	2224.4\\
82000	2204.1\\
84000	2218.2\\
86000	2221.3\\
88000	2217.9\\
90000	2245.8\\
92000	2220.9\\
94000	2234.2\\
96000	2214.8\\
98000	2239.7\\
1e+05	2216.7\\
};
\addplot [color=mycolor76, forget plot]
  table[row sep=crcr]{%
0	4395.3\\
2000	2499.3\\
4000	2013\\
6000	2009.7\\
8000	2050\\
10000	2061\\
14000	2051.9\\
16000	2012.9\\
18000	2013\\
20000	2046\\
22000	2019.5\\
24000	2034.5\\
26000	2026.5\\
28000	2041.5\\
30000	2036.1\\
34000	2051.1\\
36000	2020.3\\
42000	2017.2\\
44000	2002.9\\
46000	2049.2\\
48000	2039.6\\
50000	2003.2\\
52000	2035.3\\
54000	2021.8\\
56000	2047.8\\
58000	2041.7\\
60000	2004.1\\
64000	2017.1\\
66000	2016\\
68000	2032.1\\
70000	2024\\
72000	2041.6\\
74000	2050\\
76000	2064.8\\
78000	2044.8\\
80000	2011.3\\
82000	2049.4\\
84000	2015.9\\
86000	2037.3\\
88000	2018.1\\
90000	2011.9\\
92000	2015.5\\
94000	1990.5\\
96000	2032.9\\
98000	2024.6\\
1e+05	2041.8\\
};
\addplot [color=mycolor77, forget plot]
  table[row sep=crcr]{%
0	4395.3\\
2000	2152.1\\
4000	1596.4\\
6000	1577\\
10000	1568.9\\
12000	1544.3\\
14000	1542.3\\
16000	1567.5\\
20000	1572.2\\
22000	1565.3\\
24000	1548.3\\
28000	1598\\
30000	1579\\
32000	1590.4\\
34000	1581.7\\
36000	1578\\
38000	1589.7\\
40000	1569.4\\
42000	1559.4\\
44000	1594.4\\
46000	1562.4\\
52000	1574.5\\
56000	1571.6\\
58000	1597.9\\
60000	1586.8\\
62000	1551.5\\
66000	1574.1\\
68000	1592.3\\
70000	1558.5\\
72000	1581.5\\
74000	1565.8\\
76000	1600.8\\
78000	1577.9\\
80000	1565.5\\
82000	1568.2\\
86000	1535.9\\
88000	1573.3\\
90000	1557.7\\
92000	1590.1\\
94000	1562.3\\
96000	1571.6\\
98000	1546.9\\
1e+05	1563\\
};
\addplot [color=mycolor78, forget plot]
  table[row sep=crcr]{%
0	4395.3\\
2000	2766.5\\
4000	2287.2\\
8000	2281.7\\
12000	2360\\
14000	2288.1\\
16000	2329.8\\
18000	2336.1\\
20000	2312.8\\
22000	2342.1\\
26000	2287.8\\
28000	2335.3\\
30000	2287.4\\
32000	2307.5\\
34000	2337.7\\
36000	2285.8\\
38000	2300.5\\
40000	2283.4\\
42000	2319.3\\
44000	2318.4\\
46000	2377.6\\
48000	2300.9\\
50000	2254.2\\
52000	2280\\
54000	2277.8\\
56000	2358.5\\
58000	2316\\
60000	2289.6\\
64000	2299.3\\
66000	2325.9\\
68000	2279.8\\
70000	2323.3\\
72000	2282.4\\
74000	2309.6\\
76000	2290.4\\
78000	2291.6\\
80000	2268.3\\
82000	2270.1\\
84000	2286.5\\
86000	2310.8\\
88000	2290.5\\
90000	2333.6\\
92000	2309.2\\
94000	2313.7\\
96000	2296.5\\
1e+05	2284\\
};
\addplot [color=mycolor79, forget plot]
  table[row sep=crcr]{%
0	4395.3\\
2000	2684.7\\
4000	2160.5\\
6000	2151.6\\
10000	2182.7\\
12000	2178.5\\
14000	2196.1\\
16000	2161.8\\
20000	2164.4\\
22000	2189.1\\
24000	2199.4\\
28000	2127.3\\
30000	2187.7\\
32000	2144.8\\
34000	2163.8\\
36000	2175.2\\
42000	2148.1\\
44000	2154.2\\
46000	2153.8\\
48000	2216.5\\
50000	2173.7\\
54000	2175.2\\
56000	2145\\
58000	2139.7\\
60000	2151.8\\
64000	2140.2\\
66000	2159.6\\
68000	2195.7\\
70000	2169\\
72000	2157.2\\
74000	2167.1\\
76000	2194.9\\
82000	2131.5\\
86000	2174.4\\
88000	2147\\
90000	2162.3\\
92000	2144\\
96000	2189.4\\
98000	2154.7\\
1e+05	2146.9\\
};
\addplot [color=mycolor79, forget plot]
  table[row sep=crcr]{%
0	4395.3\\
2000	2600\\
4000	2069.3\\
6000	1990.6\\
8000	2004.1\\
10000	1987.9\\
12000	2038.1\\
14000	2048.6\\
16000	2007.1\\
20000	2027.5\\
22000	2024.6\\
24000	2006.3\\
26000	2041.5\\
28000	2031\\
30000	1996.6\\
32000	1998.1\\
34000	1964.9\\
36000	1990.8\\
38000	2029\\
40000	2010.3\\
44000	2075.4\\
46000	2010.9\\
48000	2024.8\\
50000	2011.3\\
52000	2035.8\\
54000	2003.1\\
56000	2015.1\\
58000	2043.8\\
60000	2017.1\\
62000	2049.4\\
64000	1974.4\\
68000	2059.2\\
70000	2085.1\\
72000	2018.8\\
74000	1999.9\\
76000	2057.4\\
78000	2062\\
80000	2032.8\\
82000	2015.9\\
84000	2014\\
86000	2050.6\\
88000	1989.9\\
90000	2006.8\\
92000	1994.1\\
98000	2034.5\\
1e+05	2046.2\\
};
\addplot [color=mycolor80, forget plot]
  table[row sep=crcr]{%
0	4395.3\\
2000	2203.5\\
4000	1596.5\\
6000	1567.5\\
10000	1538.6\\
12000	1562.7\\
16000	1560.3\\
18000	1566.7\\
20000	1558.3\\
22000	1572.7\\
24000	1575\\
26000	1566.1\\
28000	1564.1\\
30000	1585.6\\
34000	1557.3\\
38000	1558.8\\
40000	1553.6\\
42000	1565.9\\
44000	1558.6\\
46000	1564.1\\
48000	1556.4\\
50000	1562.1\\
52000	1590.7\\
54000	1556.1\\
56000	1573.2\\
58000	1572.8\\
60000	1586.4\\
62000	1562.4\\
64000	1558.2\\
66000	1571.1\\
68000	1572.2\\
70000	1556.3\\
72000	1582.5\\
74000	1561.2\\
76000	1559.3\\
78000	1573.3\\
80000	1574.3\\
82000	1559.9\\
84000	1572.3\\
86000	1555.3\\
88000	1570.5\\
90000	1573\\
92000	1551\\
94000	1569.3\\
96000	1557.5\\
98000	1559.7\\
1e+05	1579.1\\
};
\addplot [color=mycolor81, forget plot]
  table[row sep=crcr]{%
0	4395.3\\
2000	2118.8\\
4000	1496.9\\
6000	1462.8\\
8000	1445.8\\
10000	1454\\
12000	1475.6\\
14000	1438.2\\
16000	1472.3\\
18000	1446\\
20000	1474.6\\
22000	1446\\
24000	1462.6\\
26000	1462.5\\
28000	1473.9\\
30000	1474\\
32000	1443.7\\
34000	1478.3\\
36000	1466.4\\
38000	1432.5\\
40000	1469.6\\
44000	1465.6\\
46000	1450.8\\
50000	1462.3\\
52000	1465.9\\
54000	1431.8\\
56000	1464.9\\
58000	1449.9\\
60000	1479.2\\
62000	1445.6\\
64000	1477.9\\
66000	1461.3\\
68000	1488.8\\
70000	1479\\
74000	1449.7\\
76000	1455.8\\
78000	1476.6\\
80000	1488.2\\
82000	1470.2\\
84000	1489\\
86000	1472.8\\
90000	1466.7\\
92000	1477.1\\
94000	1475.9\\
98000	1462.9\\
1e+05	1449.9\\
};
\addplot [color=mycolor82, forget plot]
  table[row sep=crcr]{%
0	4395.3\\
2000	2329.5\\
4000	1724.7\\
6000	1676.6\\
8000	1671.8\\
12000	1726.5\\
18000	1754.8\\
20000	1710.6\\
22000	1687.9\\
24000	1682.6\\
26000	1700.6\\
28000	1694.1\\
30000	1720.5\\
32000	1710.7\\
34000	1739.2\\
36000	1716.8\\
38000	1730.8\\
40000	1698.8\\
42000	1676\\
44000	1688.1\\
46000	1730.9\\
48000	1707.4\\
50000	1695.9\\
52000	1734\\
56000	1697\\
58000	1713.3\\
60000	1714.4\\
62000	1697.6\\
64000	1704.5\\
72000	1708.5\\
74000	1706.3\\
76000	1715.8\\
78000	1703.7\\
80000	1722.1\\
82000	1701\\
86000	1717.3\\
88000	1685.2\\
90000	1708.7\\
92000	1722.5\\
94000	1708.9\\
96000	1727.5\\
98000	1707.9\\
1e+05	1695.2\\
};
\addplot [color=mycolor83, forget plot]
  table[row sep=crcr]{%
0	4395.3\\
2000	2358.7\\
4000	1745.9\\
6000	1728.2\\
8000	1719.1\\
10000	1736.5\\
12000	1725.3\\
14000	1729.7\\
16000	1756.3\\
18000	1773.5\\
20000	1732.9\\
22000	1722.4\\
24000	1733.4\\
26000	1712.5\\
28000	1740.2\\
30000	1754.2\\
32000	1756.5\\
34000	1723.8\\
36000	1709.5\\
38000	1716.7\\
40000	1711\\
42000	1735.7\\
44000	1736.1\\
48000	1751.4\\
50000	1731\\
52000	1728.4\\
54000	1753.6\\
56000	1761.4\\
60000	1720.3\\
62000	1734.8\\
64000	1718\\
66000	1751.3\\
68000	1714.6\\
70000	1749.7\\
74000	1738.5\\
76000	1751.9\\
78000	1725.8\\
80000	1712.1\\
82000	1750.7\\
84000	1766.7\\
86000	1736.3\\
90000	1733.1\\
92000	1759.9\\
94000	1712.7\\
96000	1705.3\\
1e+05	1745.6\\
};
\addplot [color=mycolor83, forget plot]
  table[row sep=crcr]{%
0	4395.3\\
2000	1959.4\\
4000	1285.4\\
6000	1287.2\\
8000	1251.8\\
10000	1256.1\\
12000	1284.8\\
14000	1266.9\\
16000	1279.4\\
20000	1257.7\\
22000	1241\\
24000	1285.6\\
26000	1276.6\\
28000	1239.2\\
30000	1262.3\\
32000	1236.7\\
34000	1251.1\\
36000	1251.1\\
38000	1270.3\\
40000	1267.1\\
42000	1292.6\\
44000	1278.5\\
46000	1277.2\\
48000	1261.7\\
50000	1291.2\\
52000	1290.8\\
54000	1268.8\\
56000	1291.7\\
58000	1266.3\\
64000	1255\\
66000	1257.1\\
68000	1271.1\\
70000	1264.9\\
72000	1288.2\\
74000	1285.8\\
76000	1277.9\\
78000	1266.3\\
80000	1284.6\\
82000	1280.8\\
84000	1266.6\\
86000	1281.8\\
88000	1270.1\\
90000	1297.7\\
94000	1239.6\\
96000	1262.9\\
98000	1256\\
1e+05	1278.9\\
};
\addplot [color=mycolor84, forget plot]
  table[row sep=crcr]{%
0	4395.3\\
2000	2038.2\\
4000	1313.6\\
6000	1305.2\\
8000	1258.2\\
10000	1269.9\\
14000	1299.1\\
16000	1287.4\\
18000	1285.2\\
22000	1267.6\\
24000	1297.8\\
26000	1308.8\\
30000	1277.2\\
32000	1305.7\\
34000	1297\\
36000	1305.6\\
38000	1260.6\\
40000	1279.3\\
42000	1286\\
44000	1275.7\\
46000	1272.3\\
48000	1264.3\\
50000	1309.3\\
52000	1293.3\\
54000	1289.8\\
56000	1297.7\\
58000	1284.5\\
60000	1277.7\\
62000	1291\\
64000	1266.4\\
66000	1231.3\\
68000	1271.1\\
70000	1288.8\\
74000	1310.1\\
76000	1299\\
78000	1298.5\\
80000	1287.7\\
82000	1291.1\\
84000	1276.3\\
86000	1289\\
88000	1246.5\\
90000	1272.8\\
92000	1295.6\\
94000	1301.4\\
96000	1259.5\\
98000	1281.4\\
1e+05	1295.5\\
};
\addplot [color=mycolor85, forget plot]
  table[row sep=crcr]{%
0	4395.3\\
2000	2098.9\\
4000	1455\\
6000	1379.5\\
8000	1375.8\\
10000	1389.1\\
14000	1430.4\\
16000	1432.7\\
18000	1393.8\\
20000	1368.8\\
22000	1383.4\\
24000	1433.9\\
26000	1399.5\\
28000	1440.8\\
30000	1407.4\\
32000	1433.8\\
34000	1399.7\\
36000	1414\\
38000	1372\\
40000	1368.7\\
42000	1407.2\\
44000	1430.4\\
46000	1375.4\\
48000	1402.9\\
50000	1409.2\\
52000	1433.3\\
54000	1415.3\\
56000	1414\\
58000	1449.2\\
60000	1416.4\\
62000	1400.5\\
64000	1371.3\\
66000	1430.5\\
68000	1404.5\\
70000	1395.5\\
72000	1403.9\\
74000	1370.6\\
76000	1434.9\\
78000	1419.9\\
80000	1368.8\\
82000	1387.5\\
84000	1355.8\\
88000	1409.2\\
90000	1408.1\\
92000	1447.5\\
94000	1411.9\\
96000	1369.5\\
98000	1400.4\\
1e+05	1392.4\\
};
\addplot [color=mycolor86, forget plot]
  table[row sep=crcr]{%
0	4395.3\\
2000	1994.4\\
4000	1210.2\\
6000	1171.9\\
8000	1173.9\\
10000	1203.6\\
12000	1209.6\\
14000	1198.2\\
16000	1177.6\\
18000	1210.6\\
20000	1196.4\\
26000	1176.7\\
28000	1158.5\\
30000	1176.3\\
32000	1223.4\\
34000	1168.3\\
36000	1167.4\\
38000	1185.5\\
40000	1174.7\\
42000	1177\\
44000	1200.1\\
46000	1189.1\\
48000	1198.5\\
50000	1183.8\\
56000	1158.5\\
58000	1180.8\\
60000	1185.1\\
62000	1147.5\\
64000	1177.2\\
66000	1211.6\\
68000	1212.4\\
70000	1174\\
72000	1148.4\\
74000	1142.9\\
76000	1166.6\\
78000	1175.1\\
80000	1201.3\\
82000	1209.9\\
84000	1180.6\\
86000	1199.8\\
88000	1186.3\\
90000	1195.1\\
92000	1184.8\\
94000	1169.3\\
96000	1170.3\\
98000	1190.2\\
1e+05	1194.6\\
};
\addplot [color=mycolor87, forget plot]
  table[row sep=crcr]{%
0	4395.3\\
2000	1849.1\\
4000	1058.3\\
6000	1037\\
8000	988.57\\
10000	987.08\\
12000	989.46\\
14000	999.69\\
16000	1018.5\\
18000	982.01\\
20000	977.99\\
22000	1021\\
26000	1001.4\\
28000	1007.8\\
30000	997.27\\
32000	1025.3\\
34000	995.07\\
36000	1032\\
38000	993.35\\
40000	986.98\\
42000	1059.5\\
44000	1017.3\\
46000	950.07\\
48000	1017.3\\
50000	1030.8\\
52000	994.09\\
54000	1000.6\\
56000	1032.5\\
58000	1001.6\\
60000	1006.5\\
62000	1019.3\\
64000	1005.7\\
66000	999.44\\
70000	978.27\\
74000	1010.3\\
78000	983.21\\
80000	1010.8\\
84000	965.61\\
86000	995.18\\
88000	1016.8\\
90000	1020\\
92000	1020.3\\
94000	1014.7\\
96000	991.14\\
98000	984.67\\
1e+05	1031\\
};
\addplot [color=mycolor88, forget plot]
  table[row sep=crcr]{%
0	4395.3\\
2000	2077.9\\
4000	1245.4\\
6000	1207\\
8000	1147.8\\
10000	1188.6\\
12000	1222.5\\
14000	1182.9\\
16000	1212.4\\
18000	1204.8\\
20000	1208.4\\
22000	1142.7\\
24000	1175.2\\
26000	1201.7\\
28000	1175.2\\
30000	1201.5\\
32000	1193.9\\
34000	1180.9\\
36000	1186.9\\
38000	1176.8\\
40000	1195\\
42000	1175.4\\
44000	1161.1\\
46000	1183.7\\
48000	1196.9\\
50000	1172.9\\
54000	1183.3\\
56000	1200.3\\
58000	1240.5\\
60000	1202.4\\
62000	1208.9\\
64000	1193.3\\
66000	1182.5\\
68000	1202.9\\
70000	1191.1\\
72000	1174.9\\
74000	1162.5\\
76000	1175.9\\
78000	1151.6\\
80000	1171\\
82000	1197.9\\
84000	1187.9\\
86000	1218.9\\
88000	1197.2\\
90000	1207.4\\
92000	1189.5\\
94000	1203\\
96000	1193.5\\
1e+05	1198.6\\
};
\addplot [color=mycolor89, forget plot]
  table[row sep=crcr]{%
0	4395.3\\
2000	1755.5\\
4000	889.93\\
6000	870.62\\
8000	857.85\\
10000	858.93\\
12000	868.15\\
14000	870.27\\
16000	845.69\\
18000	837.35\\
20000	843.62\\
22000	861.24\\
24000	849.34\\
26000	864.44\\
28000	848.26\\
30000	875.13\\
32000	863.01\\
34000	847.73\\
36000	851.15\\
38000	869.04\\
40000	863.58\\
42000	841.13\\
44000	857.3\\
46000	831.1\\
48000	843.59\\
50000	853.09\\
52000	854.03\\
54000	860.4\\
56000	862.38\\
58000	854.42\\
60000	876.03\\
62000	832.11\\
64000	850.08\\
66000	856.59\\
68000	828.42\\
70000	837.45\\
72000	868.4\\
74000	844.73\\
76000	865.33\\
78000	837.63\\
80000	871.68\\
82000	863.46\\
84000	839.42\\
86000	844.68\\
88000	853.94\\
90000	860.09\\
92000	874.07\\
94000	884.82\\
96000	862.23\\
98000	860.41\\
1e+05	867.72\\
};
\addplot [color=mycolor90, forget plot]
  table[row sep=crcr]{%
0	4395.3\\
2000	1773.8\\
4000	890.68\\
6000	844.47\\
8000	820.14\\
10000	811.6\\
12000	825.12\\
14000	826.37\\
16000	856.92\\
18000	828.92\\
22000	818.52\\
24000	828.36\\
26000	817.32\\
28000	816.79\\
30000	819.13\\
32000	847.16\\
34000	839.11\\
36000	836.49\\
38000	824.3\\
40000	819.4\\
42000	805.37\\
44000	802.68\\
46000	836.47\\
48000	846.65\\
50000	818.31\\
52000	818.1\\
56000	831.26\\
58000	831.04\\
60000	838.6\\
62000	815.48\\
66000	847.64\\
68000	817.61\\
70000	845.16\\
72000	846.28\\
76000	835.43\\
78000	839.02\\
80000	822.63\\
82000	825.51\\
84000	837.57\\
86000	804.51\\
88000	821.54\\
90000	824.11\\
92000	835.59\\
94000	838.06\\
96000	859.36\\
98000	841.26\\
1e+05	819.61\\
};
\addplot [color=mycolor91, forget plot]
  table[row sep=crcr]{%
0	4395.3\\
2000	1725.3\\
4000	757.04\\
6000	693.74\\
8000	699.45\\
10000	716.62\\
12000	731.56\\
14000	718.65\\
16000	734.9\\
18000	736.68\\
20000	721.6\\
22000	677.28\\
24000	700.07\\
26000	718.58\\
28000	730.42\\
30000	687.83\\
32000	733.43\\
34000	704.17\\
36000	728.35\\
38000	730.66\\
40000	717.09\\
42000	709.24\\
44000	713.82\\
46000	727.27\\
48000	724.51\\
50000	714.01\\
52000	718.15\\
54000	710.72\\
56000	719.39\\
58000	704.23\\
60000	721.39\\
62000	711.9\\
64000	704.7\\
66000	692.52\\
68000	731.39\\
70000	714.65\\
72000	719.88\\
74000	728.85\\
76000	712.23\\
78000	760.65\\
80000	706.32\\
82000	728.23\\
84000	720.08\\
86000	741.37\\
88000	704.03\\
92000	731.53\\
94000	729.77\\
96000	714.68\\
98000	716.55\\
1e+05	693.94\\
};
\addplot [color=mycolor92, forget plot]
  table[row sep=crcr]{%
0	4395.3\\
2000	1966.9\\
4000	968.61\\
6000	898.91\\
8000	925.77\\
10000	915.34\\
12000	868.7\\
14000	867.31\\
16000	889.86\\
18000	888.14\\
20000	870.85\\
22000	884.61\\
24000	910.14\\
26000	895.75\\
28000	903.55\\
30000	899.43\\
32000	890.3\\
34000	893.29\\
36000	901.65\\
38000	921.85\\
40000	899.07\\
42000	921.52\\
44000	898.9\\
46000	869.91\\
48000	892.81\\
50000	881.29\\
52000	848.58\\
54000	898.8\\
58000	887.87\\
60000	909.32\\
62000	886.14\\
64000	897.8\\
66000	903.88\\
68000	869.75\\
70000	908.56\\
72000	923.91\\
74000	896.45\\
76000	913.26\\
78000	884.1\\
80000	909.28\\
82000	896.34\\
84000	895.55\\
86000	888.39\\
88000	919.23\\
90000	903.43\\
92000	921.19\\
94000	922.63\\
96000	914.7\\
98000	889.57\\
1e+05	887.56\\
};
\addplot [color=mycolor93, forget plot]
  table[row sep=crcr]{%
0	4395.3\\
2000	1878\\
4000	860.9\\
6000	767.84\\
8000	711\\
10000	733.66\\
12000	711.09\\
14000	718.31\\
16000	763.56\\
18000	758.53\\
20000	757.24\\
22000	723.32\\
24000	724.58\\
26000	753.14\\
28000	750.83\\
30000	756.83\\
32000	753.36\\
34000	722.98\\
36000	753.28\\
38000	731.78\\
40000	731.31\\
42000	783.79\\
44000	763.49\\
46000	751.79\\
48000	774.85\\
50000	749.97\\
52000	694.78\\
54000	755.21\\
56000	748.52\\
58000	733.32\\
60000	739.98\\
62000	753.61\\
64000	755.93\\
66000	734.02\\
68000	774.5\\
70000	753.23\\
72000	788.76\\
74000	759.69\\
76000	788.86\\
78000	747.3\\
82000	713.95\\
84000	726.8\\
86000	736.74\\
88000	742.89\\
90000	720.92\\
92000	716.61\\
94000	739.76\\
96000	771.98\\
1e+05	715.14\\
};
\addplot [color=mycolor94, forget plot]
  table[row sep=crcr]{%
0	4395.3\\
2000	1876.1\\
4000	773.9\\
6000	725.47\\
8000	745.68\\
10000	717.23\\
12000	724.73\\
14000	705.9\\
16000	710.39\\
18000	717.73\\
20000	709.77\\
22000	732.59\\
24000	710.97\\
26000	731.34\\
28000	724.75\\
30000	739.92\\
32000	724.53\\
34000	716.71\\
36000	688.58\\
38000	733.1\\
40000	722.49\\
42000	722.58\\
44000	716.08\\
46000	715.14\\
48000	703.78\\
50000	695.55\\
52000	711.14\\
54000	721.69\\
56000	715.01\\
58000	711.71\\
60000	704.26\\
62000	701.61\\
64000	702.54\\
66000	718.31\\
68000	716.97\\
70000	700.13\\
72000	695.12\\
74000	712.96\\
76000	701.04\\
78000	698.08\\
80000	704.4\\
82000	705.02\\
84000	714.31\\
86000	700.73\\
88000	724.58\\
92000	728.97\\
94000	711.7\\
96000	710.13\\
98000	697.94\\
1e+05	708.94\\
};
\addplot [color=mycolor95, forget plot]
  table[row sep=crcr]{%
0	4395.3\\
2000	1648.3\\
4000	522.61\\
6000	441.65\\
8000	431.41\\
12000	443.32\\
14000	426.82\\
16000	415.61\\
18000	414.59\\
20000	415.55\\
22000	426.6\\
24000	418.97\\
26000	420.29\\
28000	431.09\\
30000	434.72\\
32000	425.29\\
34000	422.08\\
36000	439.76\\
38000	412.82\\
40000	424.28\\
42000	431.56\\
44000	422.23\\
46000	418.07\\
48000	433.14\\
50000	419.93\\
52000	429.93\\
54000	428.87\\
56000	414.41\\
58000	405.94\\
60000	415.9\\
62000	433.27\\
64000	424.99\\
66000	428.37\\
68000	423.04\\
70000	425.91\\
74000	417.69\\
76000	425.34\\
78000	424.98\\
80000	434.42\\
82000	422.68\\
86000	408.71\\
88000	428.04\\
90000	422.38\\
92000	435.81\\
94000	432.75\\
96000	431.52\\
98000	418.22\\
1e+05	437.61\\
};
\addplot [color=mycolor96, forget plot]
  table[row sep=crcr]{%
0	4395.3\\
2000	1740.2\\
4000	541.95\\
6000	445.05\\
8000	457.36\\
10000	479.32\\
12000	436.11\\
14000	431.44\\
16000	436.69\\
18000	470.71\\
20000	462.33\\
22000	490.25\\
24000	471.14\\
26000	459.64\\
28000	476.46\\
30000	430.06\\
32000	469.69\\
34000	450.61\\
36000	457.64\\
38000	436.48\\
40000	446.88\\
42000	443.77\\
44000	483.5\\
46000	466.18\\
48000	456.83\\
50000	437.36\\
52000	464.34\\
54000	463.49\\
56000	429.04\\
58000	435.33\\
60000	453.66\\
62000	463.6\\
64000	489.73\\
66000	484.13\\
68000	464.42\\
70000	442.98\\
72000	431.81\\
74000	453.19\\
76000	509.4\\
78000	462.4\\
80000	466.65\\
82000	453.65\\
84000	457.35\\
86000	462.39\\
88000	457.52\\
90000	494.75\\
92000	441.3\\
94000	434.19\\
96000	430.66\\
98000	458.53\\
1e+05	470.21\\
};
\addplot [color=mycolor97, forget plot]
  table[row sep=crcr]{%
0	4395.3\\
2000	1621.3\\
4000	436.77\\
6000	324.54\\
8000	321.34\\
12000	298.01\\
14000	302.14\\
16000	312.94\\
18000	326.23\\
20000	342.68\\
22000	303.91\\
24000	304.99\\
26000	298.96\\
28000	305.9\\
30000	320.36\\
32000	297.71\\
34000	305.9\\
36000	316.07\\
38000	316.08\\
40000	317.02\\
42000	310.77\\
44000	312.86\\
46000	319.53\\
48000	290.18\\
50000	312.78\\
54000	331.82\\
56000	325.91\\
58000	333.68\\
60000	351.75\\
64000	344.58\\
66000	324.18\\
68000	299.03\\
70000	318.75\\
72000	307.01\\
74000	313.31\\
80000	314.59\\
82000	333.51\\
84000	326.31\\
86000	316.35\\
88000	331.79\\
90000	335.1\\
92000	314.14\\
94000	308.51\\
96000	320.21\\
98000	321.23\\
1e+05	339.76\\
};
\addplot [color=mycolor98]
  table[row sep=crcr]{%
0	4395.3\\
2000	1892.4\\
4000	580.88\\
6000	431.84\\
8000	433.59\\
10000	428.64\\
12000	437.76\\
14000	435.52\\
16000	411.08\\
18000	418.51\\
20000	413.81\\
22000	420.24\\
24000	435.65\\
26000	421.24\\
28000	420.61\\
30000	429.41\\
32000	440.14\\
34000	437.51\\
36000	441.32\\
38000	437.22\\
40000	450.99\\
42000	436.26\\
44000	439.49\\
46000	420.04\\
48000	417.6\\
50000	427.05\\
52000	406.13\\
54000	438.13\\
56000	429.56\\
58000	427.87\\
60000	435.79\\
62000	449.63\\
64000	446.36\\
66000	421.43\\
68000	430.01\\
70000	432.5\\
72000	425.03\\
74000	406.99\\
76000	419.2\\
78000	405.09\\
80000	405.99\\
82000	416.41\\
84000	416.04\\
86000	410.7\\
88000	421.31\\
90000	415.1\\
92000	411.98\\
94000	426.76\\
96000	426.74\\
98000	430.37\\
1e+05	421.55\\
};
\addlegendentry{\ClasskIVLegend}

      % \input{../../../../.tex/20/KS/SizeEvolutionsfig6.tex}
      \legend{}
    \end{groupplot}

    \begin{groupplot}[
      group style={
        group name=Row4,
        plotGroup,
        group size=2 by 1,
      },
      plotRow,
      row5Keys,
      xTickScaleLabelEmpty,
      plotLegendSE,
      %
      xmin=0,xmax=1e+05,
      ymin=100,ymax=1.5886e+05,
    ]
      \nextgroupplot[%
        at={($(Row3 c1r1.south west)-(0,\yPlotSep em)-(0,\plotHeight cm)$)},
        % xmin=0,xmax=1e+05,
        % ymin=100,ymax=1.5592e+05,
      ] % This file was created by matlab2tikz.
%
\definecolor{mycolor1}{rgb}{0.00146,0.00047,0.01387}%
\definecolor{mycolor2}{rgb}{0.01166,0.00942,0.06346}%
\definecolor{mycolor3}{rgb}{0.02943,0.02150,0.11462}%
\definecolor{mycolor4}{rgb}{0.06134,0.03659,0.17764}%
\definecolor{mycolor5}{rgb}{0.08741,0.04456,0.22481}%
\definecolor{mycolor6}{rgb}{0.12291,0.04754,0.28162}%
\definecolor{mycolor7}{rgb}{0.14907,0.04547,0.31709}%
\definecolor{mycolor8}{rgb}{0.18343,0.04033,0.35497}%
\definecolor{mycolor9}{rgb}{0.21795,0.03662,0.38352}%
\definecolor{mycolor10}{rgb}{0.24497,0.03705,0.40001}%
\definecolor{mycolor11}{rgb}{0.27135,0.04092,0.41198}%
\definecolor{mycolor12}{rgb}{0.29076,0.04564,0.41864}%
\definecolor{mycolor13}{rgb}{0.31628,0.05349,0.42512}%
\definecolor{mycolor14}{rgb}{0.33522,0.06006,0.42852}%
\definecolor{mycolor15}{rgb}{0.36028,0.06925,0.43150}%
\definecolor{mycolor16}{rgb}{0.37900,0.07625,0.43272}%
\definecolor{mycolor17}{rgb}{0.39767,0.08326,0.43318}%
\definecolor{mycolor18}{rgb}{0.41633,0.09020,0.43294}%
\definecolor{mycolor19}{rgb}{0.42877,0.09479,0.43241}%
\definecolor{mycolor20}{rgb}{0.44743,0.10160,0.43108}%
\definecolor{mycolor21}{rgb}{0.46610,0.10832,0.42912}%
\definecolor{mycolor22}{rgb}{0.47856,0.11276,0.42747}%
\definecolor{mycolor23}{rgb}{0.49726,0.11938,0.42449}%
\definecolor{mycolor24}{rgb}{0.50973,0.12377,0.42216}%
\definecolor{mycolor25}{rgb}{0.52221,0.12815,0.41955}%
\definecolor{mycolor26}{rgb}{0.53468,0.13253,0.41667}%
\definecolor{mycolor27}{rgb}{0.55339,0.13913,0.41183}%
\definecolor{mycolor28}{rgb}{0.56585,0.14357,0.40826}%
\definecolor{mycolor29}{rgb}{0.57830,0.14804,0.40441}%
\definecolor{mycolor30}{rgb}{0.59073,0.15256,0.40029}%
\definecolor{mycolor31}{rgb}{0.60314,0.15715,0.39589}%
\definecolor{mycolor32}{rgb}{0.61551,0.16182,0.39122}%
\definecolor{mycolor33}{rgb}{0.62169,0.16418,0.38878}%
\definecolor{mycolor34}{rgb}{0.63400,0.16899,0.38370}%
\definecolor{mycolor35}{rgb}{0.64626,0.17391,0.37836}%
\definecolor{mycolor36}{rgb}{0.65846,0.17896,0.37275}%
\definecolor{mycolor37}{rgb}{0.66454,0.18154,0.36985}%
\definecolor{mycolor38}{rgb}{0.67664,0.18681,0.36385}%
\definecolor{mycolor39}{rgb}{0.68865,0.19224,0.35760}%
\definecolor{mycolor40}{rgb}{0.69463,0.19502,0.35439}%
\definecolor{mycolor41}{rgb}{0.70650,0.20073,0.34778}%
\definecolor{mycolor42}{rgb}{0.71240,0.20366,0.34438}%
\definecolor{mycolor43}{rgb}{0.72410,0.20967,0.33742}%
\definecolor{mycolor44}{rgb}{0.72991,0.21276,0.33386}%
\definecolor{mycolor45}{rgb}{0.74142,0.21911,0.32658}%
\definecolor{mycolor46}{rgb}{0.74713,0.22238,0.32286}%
\definecolor{mycolor47}{rgb}{0.75279,0.22571,0.31909}%
\definecolor{mycolor48}{rgb}{0.76401,0.23255,0.31140}%
\definecolor{mycolor49}{rgb}{0.76956,0.23608,0.30749}%
\definecolor{mycolor50}{rgb}{0.77506,0.23967,0.30353}%
\definecolor{mycolor51}{rgb}{0.79129,0.25086,0.29139}%
\definecolor{mycolor52}{rgb}{0.79661,0.25473,0.28726}%
\definecolor{mycolor53}{rgb}{0.80187,0.25867,0.28310}%
\definecolor{mycolor54}{rgb}{0.81224,0.26679,0.27466}%
\definecolor{mycolor55}{rgb}{0.81734,0.27095,0.27039}%
\definecolor{mycolor56}{rgb}{0.82737,0.27952,0.26175}%
\definecolor{mycolor57}{rgb}{0.83230,0.28391,0.25738}%
\definecolor{mycolor58}{rgb}{0.83717,0.28839,0.25299}%
\definecolor{mycolor59}{rgb}{0.84671,0.29756,0.24411}%
\definecolor{mycolor60}{rgb}{0.85138,0.30226,0.23964}%
\definecolor{mycolor61}{rgb}{0.85599,0.30704,0.23513}%
\definecolor{mycolor62}{rgb}{0.86053,0.31189,0.23061}%
\definecolor{mycolor63}{rgb}{0.87374,0.32691,0.21689}%
\definecolor{mycolor64}{rgb}{0.87800,0.33206,0.21227}%
\definecolor{mycolor65}{rgb}{0.89034,0.34796,0.19829}%
\definecolor{mycolor66}{rgb}{0.89819,0.35891,0.18886}%
\definecolor{mycolor67}{rgb}{0.90200,0.36449,0.18412}%
\definecolor{mycolor68}{rgb}{0.90573,0.37014,0.17935}%
\definecolor{mycolor69}{rgb}{0.90939,0.37586,0.17456}%
\definecolor{mycolor70}{rgb}{0.91297,0.38164,0.16975}%
\definecolor{mycolor71}{rgb}{0.91646,0.38748,0.16492}%
\definecolor{mycolor72}{rgb}{0.91988,0.39339,0.16007}%
\definecolor{mycolor73}{rgb}{0.92322,0.39936,0.15519}%
\definecolor{mycolor74}{rgb}{0.92647,0.40539,0.15029}%
\definecolor{mycolor75}{rgb}{0.92964,0.41148,0.14537}%
\definecolor{mycolor76}{rgb}{0.93274,0.41763,0.14042}%
\definecolor{mycolor77}{rgb}{0.93575,0.42383,0.13544}%
\definecolor{mycolor78}{rgb}{0.93868,0.43009,0.13044}%
\definecolor{mycolor79}{rgb}{0.94152,0.43640,0.12541}%
\definecolor{mycolor80}{rgb}{0.94429,0.44277,0.12035}%
\definecolor{mycolor81}{rgb}{0.94956,0.45566,0.11016}%
\definecolor{mycolor82}{rgb}{0.95208,0.46218,0.10503}%
\definecolor{mycolor83}{rgb}{0.96129,0.48872,0.08429}%
\definecolor{mycolor84}{rgb}{0.96540,0.50225,0.07386}%
\definecolor{mycolor85}{rgb}{0.96732,0.50908,0.06866}%
\definecolor{mycolor86}{rgb}{0.97259,0.52980,0.05332}%
\definecolor{mycolor87}{rgb}{0.97568,0.54380,0.04362}%
\definecolor{mycolor88}{rgb}{0.98082,0.57221,0.02851}%
\definecolor{mycolor89}{rgb}{0.98189,0.57939,0.02625}%
\definecolor{mycolor90}{rgb}{0.98378,0.59385,0.02377}%
\definecolor{mycolor91}{rgb}{0.98532,0.60842,0.02420}%
\definecolor{mycolor92}{rgb}{0.98696,0.63048,0.03091}%
\definecolor{mycolor93}{rgb}{0.98793,0.66025,0.05175}%
\definecolor{mycolor94}{rgb}{0.98771,0.68281,0.07249}%
\definecolor{mycolor95}{rgb}{0.97750,0.78226,0.18592}%
\definecolor{mycolor96}{rgb}{0.96624,0.83619,0.26153}%
\definecolor{mycolor97}{rgb}{0.96252,0.85148,0.28555}%
\definecolor{mycolor98}{rgb}{0.98836,0.99836,0.64492}%
%
\addplot [color=mycolor1]
  table[row sep=crcr]{%
0	4395.3\\
2000	1.0922e+05\\
4000	1.4937e+05\\
6000	1.4897e+05\\
10000	1.5192e+05\\
16000	1.5095e+05\\
18000	1.4784e+05\\
22000	1.5219e+05\\
24000	1.5168e+05\\
26000	1.4951e+05\\
28000	1.5306e+05\\
30000	1.517e+05\\
32000	1.4985e+05\\
34000	1.5123e+05\\
36000	1.4854e+05\\
40000	1.5167e+05\\
42000	1.5442e+05\\
44000	1.5563e+05\\
48000	1.524e+05\\
50000	1.4893e+05\\
52000	1.5078e+05\\
54000	1.548e+05\\
56000	1.5525e+05\\
58000	1.5318e+05\\
60000	1.5036e+05\\
62000	1.499e+05\\
64000	1.5018e+05\\
66000	1.5184e+05\\
68000	1.5263e+05\\
70000	1.5079e+05\\
74000	1.5182e+05\\
76000	1.5437e+05\\
78000	1.5592e+05\\
80000	1.5145e+05\\
82000	1.5071e+05\\
84000	1.5138e+05\\
86000	1.5099e+05\\
88000	1.489e+05\\
90000	1.5143e+05\\
1e+05	1.5125e+05\\
};
\addlegendentry{\ClasskILegend}

\addplot [color=mycolor2, forget plot]
  table[row sep=crcr]{%
0	4395.3\\
2000	26004\\
4000	26062\\
6000	26198\\
8000	26238\\
10000	26018\\
12000	25367\\
14000	25707\\
16000	25879\\
18000	26202\\
22000	25215\\
24000	26450\\
26000	25984\\
28000	26379\\
30000	25791\\
32000	25861\\
34000	25663\\
36000	25818\\
38000	26159\\
40000	26039\\
42000	26224\\
44000	26529\\
46000	25081\\
48000	26003\\
50000	26843\\
52000	25807\\
56000	26240\\
58000	26145\\
60000	26534\\
62000	26752\\
64000	25947\\
66000	26477\\
68000	25707\\
70000	26660\\
72000	25724\\
74000	25842\\
76000	26938\\
78000	26323\\
80000	25810\\
82000	25537\\
84000	26490\\
86000	26058\\
88000	25995\\
90000	25699\\
92000	26689\\
94000	25652\\
96000	26479\\
98000	26627\\
1e+05	25629\\
};
\addplot [color=mycolor3, forget plot]
  table[row sep=crcr]{%
0	4395.3\\
2000	36378\\
4000	38746\\
6000	39325\\
8000	39317\\
14000	39973\\
16000	39186\\
20000	39715\\
22000	38938\\
24000	39494\\
26000	39331\\
28000	39547\\
30000	39166\\
32000	39510\\
34000	39242\\
36000	40227\\
38000	39507\\
40000	40344\\
42000	40256\\
44000	38651\\
46000	40396\\
48000	40168\\
50000	39361\\
52000	39419\\
54000	40037\\
56000	39102\\
58000	38777\\
60000	39661\\
62000	40221\\
64000	38656\\
66000	39313\\
68000	40457\\
70000	40217\\
72000	40242\\
74000	39736\\
76000	39484\\
78000	40155\\
80000	40281\\
82000	39915\\
84000	39401\\
86000	39642\\
90000	39369\\
92000	38178\\
94000	39950\\
96000	38788\\
98000	40212\\
1e+05	40092\\
};
\addplot [color=mycolor4, forget plot]
  table[row sep=crcr]{%
0	4395.3\\
2000	14990\\
4000	14178\\
6000	14494\\
8000	14236\\
10000	14185\\
12000	14612\\
14000	14381\\
16000	14597\\
18000	14583\\
20000	14307\\
22000	14580\\
24000	14249\\
26000	14462\\
28000	14388\\
30000	14259\\
32000	14612\\
34000	14481\\
36000	13718\\
38000	14415\\
42000	14605\\
44000	14490\\
46000	14536\\
48000	14875\\
50000	14419\\
54000	14128\\
56000	14602\\
58000	14437\\
60000	14187\\
62000	14687\\
64000	14320\\
66000	14706\\
68000	14782\\
70000	14204\\
72000	14702\\
74000	14740\\
76000	14396\\
78000	14417\\
80000	14332\\
82000	14518\\
84000	14310\\
86000	14425\\
88000	14475\\
90000	14283\\
92000	14485\\
94000	14185\\
96000	14582\\
98000	14229\\
1e+05	13954\\
};
\addplot [color=mycolor5, forget plot]
  table[row sep=crcr]{%
0	4395.3\\
2000	52470\\
4000	65974\\
6000	64872\\
8000	64454\\
10000	65927\\
12000	61548\\
14000	63524\\
16000	63585\\
18000	65300\\
20000	64230\\
24000	63991\\
26000	64793\\
28000	64726\\
30000	64066\\
32000	62297\\
34000	62786\\
36000	61771\\
38000	64882\\
40000	63569\\
42000	62740\\
44000	62437\\
46000	63412\\
48000	64010\\
50000	65372\\
52000	63250\\
54000	62764\\
56000	65057\\
58000	63588\\
60000	64221\\
62000	63107\\
66000	62829\\
68000	63292\\
70000	62946\\
72000	65760\\
74000	65390\\
76000	64023\\
78000	62184\\
80000	60721\\
82000	63591\\
84000	63060\\
86000	62103\\
88000	63489\\
90000	62487\\
92000	63567\\
94000	64974\\
96000	63175\\
98000	63833\\
1e+05	62730\\
};
\addplot [color=mycolor6, forget plot]
  table[row sep=crcr]{%
0	4395.3\\
2000	41984\\
4000	53801\\
6000	55176\\
8000	57020\\
10000	56768\\
12000	57189\\
14000	57014\\
16000	57293\\
18000	54324\\
22000	56785\\
24000	54749\\
28000	54767\\
30000	56756\\
32000	55245\\
34000	55037\\
36000	55970\\
38000	55783\\
40000	56256\\
42000	54767\\
44000	55272\\
46000	55483\\
48000	55001\\
50000	56925\\
52000	55287\\
54000	54803\\
56000	56202\\
58000	55120\\
60000	56524\\
62000	55165\\
64000	56245\\
66000	56414\\
68000	56036\\
70000	56621\\
72000	54052\\
74000	53475\\
76000	54630\\
78000	55042\\
80000	57973\\
82000	56285\\
84000	56782\\
86000	57838\\
88000	55124\\
90000	54701\\
92000	55978\\
94000	56232\\
96000	56719\\
1e+05	55387\\
};
\addplot [color=mycolor7, forget plot]
  table[row sep=crcr]{%
0	4395.3\\
2000	24178\\
4000	27776\\
6000	27988\\
8000	27869\\
10000	26898\\
12000	28821\\
14000	28999\\
16000	28546\\
18000	27042\\
20000	28622\\
22000	28074\\
24000	27058\\
26000	28109\\
28000	28258\\
30000	27721\\
32000	27860\\
34000	28484\\
36000	28688\\
38000	29218\\
40000	27953\\
42000	28387\\
44000	27901\\
46000	27653\\
48000	28745\\
50000	28189\\
52000	28670\\
54000	28883\\
56000	28896\\
58000	27478\\
60000	28085\\
62000	28536\\
64000	28633\\
66000	27832\\
68000	28355\\
70000	28222\\
72000	28200\\
74000	28728\\
76000	26910\\
78000	27995\\
80000	28309\\
82000	29010\\
84000	27587\\
86000	28506\\
88000	28573\\
90000	27868\\
92000	28452\\
94000	28660\\
96000	29045\\
98000	28762\\
1e+05	27798\\
};
\addplot [color=mycolor8, forget plot]
  table[row sep=crcr]{%
0	4395.3\\
2000	20527\\
4000	21197\\
6000	22136\\
8000	20139\\
12000	20716\\
14000	19703\\
16000	20632\\
18000	20151\\
20000	20774\\
22000	21608\\
24000	20726\\
26000	21079\\
28000	20500\\
30000	20509\\
32000	20318\\
36000	20757\\
38000	20405\\
40000	20478\\
42000	20639\\
44000	20679\\
46000	20943\\
48000	20192\\
50000	20807\\
52000	20518\\
54000	20138\\
56000	20173\\
58000	20845\\
60000	20571\\
62000	20785\\
64000	20336\\
66000	20598\\
68000	20579\\
70000	20382\\
72000	20512\\
74000	20749\\
76000	21654\\
78000	21057\\
80000	20604\\
82000	19796\\
84000	19945\\
86000	20186\\
88000	20538\\
90000	21484\\
92000	20802\\
94000	20507\\
96000	20766\\
98000	20492\\
1e+05	21313\\
};
\addplot [color=mycolor9, forget plot]
  table[row sep=crcr]{%
0	4395.3\\
2000	13069\\
4000	14159\\
6000	13959\\
8000	14182\\
10000	14078\\
12000	14149\\
14000	14096\\
16000	14582\\
18000	14404\\
22000	13727\\
24000	14030\\
26000	13914\\
28000	13936\\
30000	13707\\
32000	14073\\
34000	13700\\
36000	14311\\
38000	13770\\
40000	14123\\
42000	14368\\
44000	14348\\
46000	14621\\
48000	13921\\
50000	14258\\
52000	13886\\
54000	14158\\
56000	14095\\
58000	14375\\
60000	14057\\
62000	13957\\
64000	14226\\
66000	13877\\
68000	14210\\
70000	13882\\
72000	13760\\
74000	14185\\
78000	14102\\
80000	13843\\
82000	14048\\
84000	13984\\
86000	14034\\
88000	14181\\
90000	14063\\
92000	13889\\
94000	14256\\
98000	13971\\
1e+05	14236\\
};
\addplot [color=mycolor10, forget plot]
  table[row sep=crcr]{%
0	4395.3\\
2000	18934\\
4000	21114\\
6000	22337\\
8000	22355\\
10000	22178\\
12000	22120\\
14000	22230\\
16000	21954\\
18000	22668\\
20000	22051\\
24000	22231\\
28000	22107\\
30000	22248\\
32000	21925\\
34000	22249\\
36000	22355\\
38000	22223\\
40000	21725\\
42000	22101\\
44000	21748\\
46000	21553\\
48000	21263\\
50000	22579\\
52000	22037\\
54000	21868\\
56000	21457\\
58000	22156\\
60000	21917\\
62000	21996\\
64000	22429\\
66000	22101\\
68000	22583\\
72000	22201\\
74000	21623\\
76000	21506\\
80000	22470\\
82000	22432\\
84000	22314\\
86000	22423\\
88000	22170\\
90000	22364\\
92000	21559\\
96000	22105\\
98000	22019\\
1e+05	22130\\
};
\addplot [color=mycolor11, forget plot]
  table[row sep=crcr]{%
0	4395.3\\
2000	9660.1\\
4000	8903.8\\
6000	8767.3\\
8000	8710\\
10000	8552.7\\
14000	8586.8\\
16000	8334.4\\
18000	8993.9\\
20000	8823.3\\
24000	8381.4\\
26000	8872.9\\
28000	8629.9\\
30000	8480.2\\
32000	8655.6\\
34000	8725.6\\
36000	8975.7\\
38000	8810.3\\
40000	8953.5\\
42000	8706.1\\
44000	8873.5\\
46000	8879.2\\
48000	8457.9\\
50000	8642.6\\
52000	8693\\
54000	8798.7\\
56000	8452.2\\
58000	8687.9\\
60000	8662.9\\
62000	8731.1\\
64000	8568.6\\
66000	8707.8\\
68000	8762.4\\
70000	8727.4\\
72000	8773.6\\
74000	8787.5\\
76000	8589.8\\
78000	8525\\
80000	8705.6\\
82000	8713.8\\
84000	8484.3\\
86000	8684.5\\
88000	8785.7\\
90000	8681.4\\
92000	8753.5\\
94000	8494.7\\
96000	8654.4\\
98000	8955.1\\
1e+05	8827.6\\
};
\addplot [color=mycolor12, forget plot]
  table[row sep=crcr]{%
0	4395.3\\
2000	13202\\
4000	14046\\
6000	14331\\
8000	14919\\
10000	14943\\
12000	15508\\
16000	14780\\
18000	14906\\
20000	14834\\
22000	15046\\
24000	14301\\
26000	14806\\
28000	14901\\
30000	15173\\
32000	14346\\
34000	14876\\
36000	14960\\
38000	14553\\
40000	14784\\
42000	14903\\
44000	14689\\
46000	14955\\
48000	14827\\
50000	14843\\
52000	14807\\
54000	14709\\
56000	13724\\
58000	14461\\
60000	15311\\
62000	14757\\
64000	14929\\
66000	14481\\
68000	14551\\
70000	15227\\
72000	14121\\
74000	14961\\
76000	14315\\
78000	15109\\
80000	14514\\
82000	15081\\
84000	14876\\
86000	14165\\
88000	14896\\
90000	15062\\
92000	14676\\
94000	14990\\
96000	14589\\
98000	14773\\
1e+05	14909\\
};
\addplot [color=mycolor13, forget plot]
  table[row sep=crcr]{%
0	4395.3\\
2000	10761\\
4000	11030\\
6000	11060\\
10000	11711\\
12000	11151\\
14000	11165\\
16000	10978\\
18000	11332\\
20000	11444\\
22000	11192\\
24000	11216\\
26000	11171\\
28000	11459\\
30000	11177\\
32000	10942\\
34000	10938\\
36000	11064\\
38000	11422\\
40000	10785\\
42000	11351\\
44000	11165\\
46000	10642\\
48000	11263\\
50000	11126\\
54000	11263\\
56000	11176\\
58000	11040\\
60000	11294\\
62000	11004\\
64000	11291\\
66000	10889\\
68000	11163\\
70000	11342\\
72000	11174\\
74000	11100\\
76000	10741\\
78000	10789\\
80000	11395\\
82000	11351\\
84000	11108\\
86000	11181\\
88000	11053\\
90000	11451\\
92000	11233\\
94000	11085\\
96000	11168\\
98000	11000\\
1e+05	10919\\
};
\addplot [color=mycolor14, forget plot]
  table[row sep=crcr]{%
0	4395.3\\
2000	20144\\
4000	23893\\
6000	24466\\
8000	24205\\
10000	23879\\
12000	24547\\
14000	24782\\
16000	24513\\
18000	23608\\
20000	23504\\
22000	24381\\
24000	24698\\
26000	24427\\
28000	24779\\
32000	26047\\
34000	24009\\
36000	25021\\
38000	24966\\
40000	24530\\
42000	23834\\
44000	23774\\
46000	24589\\
48000	24167\\
50000	23652\\
52000	24663\\
54000	23563\\
56000	24073\\
58000	24488\\
60000	24145\\
62000	24640\\
64000	24707\\
66000	23463\\
68000	23894\\
70000	24029\\
72000	24635\\
74000	23986\\
76000	24382\\
78000	23531\\
80000	23854\\
82000	23931\\
84000	24821\\
86000	24165\\
88000	25072\\
92000	23800\\
94000	24164\\
96000	23335\\
98000	24114\\
1e+05	24279\\
};
\addplot [color=mycolor15, forget plot]
  table[row sep=crcr]{%
0	4395.3\\
2000	14867\\
4000	16337\\
6000	16640\\
8000	16452\\
10000	16911\\
12000	17250\\
14000	17191\\
16000	17265\\
18000	17461\\
20000	17344\\
22000	17064\\
24000	16848\\
26000	17131\\
28000	17001\\
30000	17612\\
32000	17173\\
34000	17444\\
36000	17390\\
38000	17061\\
40000	17175\\
42000	16823\\
44000	16884\\
48000	16861\\
50000	16550\\
52000	17074\\
54000	17187\\
56000	16622\\
58000	16852\\
60000	17318\\
62000	17565\\
64000	17250\\
66000	16678\\
68000	16914\\
70000	17068\\
72000	16251\\
74000	17083\\
76000	17194\\
78000	17243\\
80000	16425\\
82000	16404\\
84000	17302\\
86000	17085\\
88000	17114\\
90000	17461\\
92000	17173\\
94000	17228\\
96000	16562\\
98000	17183\\
1e+05	16990\\
};
\addplot [color=mycolor16, forget plot]
  table[row sep=crcr]{%
0	4395.3\\
2000	7696.4\\
4000	7064.4\\
6000	6817.4\\
8000	7124.8\\
10000	6976.6\\
12000	7200.9\\
14000	6811.9\\
16000	6766.6\\
18000	7155.1\\
20000	7062.2\\
22000	6866.5\\
24000	7010.3\\
26000	6892\\
28000	7018.7\\
30000	7010.9\\
32000	6962.3\\
34000	6956.1\\
36000	6871.8\\
38000	6936.7\\
40000	6901.5\\
42000	7193.1\\
44000	6970.1\\
46000	6847.4\\
48000	6964.2\\
50000	7015.5\\
52000	7198.8\\
54000	6815.8\\
56000	7097.7\\
58000	7264.1\\
60000	7154\\
62000	6871\\
64000	7193.6\\
66000	7045.1\\
68000	6718.1\\
72000	7125.1\\
74000	6933.6\\
76000	7174.1\\
78000	6797.9\\
80000	6948.3\\
82000	6896.2\\
84000	7011.7\\
86000	6880.3\\
88000	7093.9\\
90000	7132.3\\
94000	7059.2\\
96000	7057\\
98000	6931.8\\
1e+05	7091.2\\
};
\addplot [color=mycolor17, forget plot]
  table[row sep=crcr]{%
0	4395.3\\
2000	13848\\
4000	15391\\
6000	15818\\
8000	16193\\
10000	15744\\
12000	16315\\
14000	16088\\
16000	16374\\
20000	15895\\
22000	15720\\
24000	16801\\
26000	16168\\
28000	16603\\
30000	16121\\
32000	15941\\
34000	15810\\
36000	15749\\
38000	15751\\
40000	15928\\
42000	15747\\
44000	16095\\
46000	15572\\
48000	16378\\
50000	16029\\
52000	15780\\
54000	16080\\
56000	15736\\
58000	16213\\
60000	16028\\
62000	16324\\
64000	15981\\
66000	15520\\
68000	15691\\
70000	15988\\
72000	15941\\
74000	16181\\
76000	16363\\
78000	15760\\
80000	16498\\
84000	16228\\
86000	16475\\
88000	16623\\
90000	16411\\
92000	16666\\
94000	16392\\
96000	16294\\
98000	16072\\
1e+05	16054\\
};
\addplot [color=mycolor18, forget plot]
  table[row sep=crcr]{%
0	4395.3\\
2000	20309\\
4000	24038\\
6000	23924\\
8000	24278\\
10000	24277\\
12000	24472\\
14000	24249\\
16000	24209\\
18000	25295\\
22000	24094\\
24000	24844\\
26000	24666\\
28000	24684\\
30000	24055\\
32000	23850\\
34000	24283\\
36000	24503\\
38000	25127\\
40000	24931\\
42000	25790\\
44000	23864\\
46000	25258\\
48000	25091\\
50000	24819\\
52000	24893\\
54000	24779\\
56000	24257\\
58000	24684\\
60000	24552\\
62000	25072\\
64000	24475\\
66000	24970\\
68000	24353\\
70000	23987\\
72000	24531\\
74000	25494\\
76000	24515\\
78000	25035\\
82000	25089\\
84000	24165\\
86000	24912\\
88000	24730\\
90000	24418\\
92000	24980\\
94000	24415\\
96000	24586\\
98000	23883\\
1e+05	24856\\
};
\addplot [color=mycolor19, forget plot]
  table[row sep=crcr]{%
0	4395.3\\
2000	13376\\
4000	14291\\
6000	15370\\
8000	14929\\
10000	14814\\
14000	15380\\
16000	14996\\
18000	15198\\
20000	15354\\
22000	15156\\
24000	15026\\
26000	15318\\
28000	14896\\
30000	14716\\
32000	15350\\
34000	15177\\
36000	15133\\
38000	14768\\
40000	15189\\
42000	14420\\
44000	15250\\
46000	14777\\
48000	14554\\
50000	14993\\
52000	14735\\
54000	14534\\
56000	14568\\
58000	15095\\
60000	14842\\
62000	15146\\
64000	14991\\
66000	14960\\
68000	15436\\
70000	14504\\
72000	14726\\
74000	14848\\
76000	15089\\
78000	14546\\
80000	14522\\
82000	14947\\
86000	15071\\
88000	15049\\
90000	15465\\
92000	14723\\
94000	14933\\
96000	14822\\
98000	14467\\
1e+05	14792\\
};
\addplot [color=mycolor20, forget plot]
  table[row sep=crcr]{%
0	4395.3\\
2000	16583\\
4000	19214\\
6000	18510\\
8000	18914\\
10000	18849\\
12000	19507\\
14000	18809\\
16000	19188\\
18000	19652\\
20000	19138\\
22000	19278\\
24000	19342\\
26000	19767\\
28000	19247\\
30000	19607\\
32000	19728\\
34000	19331\\
36000	19428\\
38000	19695\\
40000	19188\\
42000	19146\\
44000	19505\\
46000	18913\\
48000	19316\\
50000	18933\\
52000	19298\\
54000	19745\\
56000	19660\\
58000	19708\\
60000	19960\\
64000	19429\\
66000	19403\\
68000	19310\\
70000	19125\\
72000	18997\\
74000	19807\\
76000	19308\\
78000	19405\\
80000	19089\\
82000	18587\\
84000	19145\\
88000	19981\\
90000	19162\\
92000	19708\\
94000	19143\\
96000	19917\\
98000	18851\\
1e+05	19213\\
};
\addplot [color=mycolor21, forget plot]
  table[row sep=crcr]{%
0	4395.3\\
2000	8698.8\\
4000	9470.2\\
6000	9559\\
8000	9546.4\\
10000	9799\\
12000	9861.2\\
14000	9633.2\\
16000	10096\\
18000	9657.8\\
20000	9759.1\\
22000	9899.1\\
24000	9673.4\\
26000	9606\\
28000	9862.7\\
32000	9635.1\\
34000	9608.2\\
36000	9806.9\\
38000	9716.6\\
40000	9662.4\\
42000	9563.6\\
44000	9702.1\\
46000	9679.6\\
48000	9512.2\\
50000	9780.1\\
52000	9958.3\\
54000	9480.4\\
56000	9910.5\\
58000	9476.4\\
60000	9650.9\\
62000	9573.8\\
64000	9610\\
66000	9700.3\\
68000	9749.1\\
72000	9550.6\\
74000	9563.8\\
76000	9722.3\\
78000	9614.9\\
80000	9604\\
82000	9716.8\\
84000	9774.7\\
86000	9733.1\\
88000	9550\\
90000	9770.2\\
92000	9603.5\\
94000	9802.4\\
96000	9909.2\\
98000	9568.5\\
1e+05	9588.6\\
};
\addplot [color=mycolor22, forget plot]
  table[row sep=crcr]{%
0	4395.3\\
2000	23968\\
4000	26223\\
6000	25431\\
8000	25447\\
10000	25870\\
12000	25265\\
14000	24433\\
16000	24606\\
18000	25284\\
20000	25649\\
22000	24934\\
24000	25240\\
26000	25063\\
28000	24796\\
30000	25598\\
32000	25944\\
34000	25908\\
36000	24957\\
38000	25321\\
40000	25334\\
42000	25078\\
44000	25401\\
46000	25182\\
48000	26574\\
50000	23958\\
52000	25481\\
54000	25597\\
56000	25333\\
58000	24809\\
60000	26064\\
62000	25136\\
64000	25407\\
66000	24802\\
68000	25101\\
70000	26063\\
72000	25779\\
74000	25420\\
76000	24503\\
78000	25989\\
80000	25914\\
82000	25603\\
84000	25417\\
86000	25005\\
88000	25602\\
90000	24806\\
92000	25501\\
94000	25359\\
96000	24517\\
98000	25677\\
1e+05	25444\\
};
\addplot [color=mycolor23, forget plot]
  table[row sep=crcr]{%
0	4395.3\\
2000	11830\\
4000	12431\\
6000	12439\\
8000	12147\\
10000	12123\\
12000	12180\\
16000	12168\\
18000	12378\\
20000	12323\\
22000	12547\\
24000	12593\\
26000	12385\\
28000	12480\\
30000	12341\\
32000	12309\\
34000	12356\\
36000	12468\\
40000	12261\\
42000	12341\\
44000	12306\\
46000	12385\\
48000	11986\\
50000	12503\\
52000	12401\\
54000	12095\\
56000	12418\\
58000	12238\\
60000	12107\\
62000	12331\\
64000	12204\\
68000	12518\\
70000	12064\\
72000	12441\\
74000	12163\\
76000	12117\\
78000	12254\\
80000	12290\\
82000	12588\\
84000	12540\\
86000	12370\\
88000	12433\\
92000	12212\\
98000	12402\\
1e+05	12372\\
};
\addplot [color=mycolor24, forget plot]
  table[row sep=crcr]{%
0	4395.3\\
2000	5455.6\\
4000	4848.6\\
6000	4722.4\\
10000	4807\\
12000	4649.9\\
14000	4745.1\\
18000	4797.6\\
20000	4879.7\\
22000	4772.1\\
24000	4785.2\\
26000	4940.2\\
28000	4752\\
30000	4694.4\\
32000	4865.6\\
34000	4800.1\\
36000	4824.2\\
38000	4662.8\\
40000	4749.1\\
42000	5000.6\\
44000	4859.2\\
46000	4657.4\\
48000	4753.8\\
50000	4788.6\\
52000	4936.9\\
54000	4800.1\\
56000	4711\\
58000	4599\\
60000	4854.6\\
64000	4935.9\\
66000	4717.9\\
68000	4894\\
70000	4752.3\\
72000	4896\\
74000	4955.2\\
76000	4897.9\\
78000	4894.1\\
80000	4678\\
82000	4693\\
84000	4683.8\\
86000	4969\\
88000	4842.6\\
90000	4910.8\\
92000	4607.1\\
94000	4785.4\\
96000	4713.1\\
98000	4711\\
1e+05	4790.8\\
};
\addplot [color=mycolor25, forget plot]
  table[row sep=crcr]{%
0	4395.3\\
2000	11037\\
4000	12261\\
6000	12600\\
8000	12750\\
10000	12810\\
12000	12921\\
14000	12176\\
18000	12452\\
20000	12309\\
22000	12251\\
24000	12542\\
26000	13015\\
28000	12277\\
30000	12582\\
32000	12825\\
34000	12535\\
36000	12844\\
38000	12573\\
40000	12490\\
42000	12046\\
44000	12500\\
46000	12626\\
48000	12791\\
50000	12574\\
52000	12488\\
54000	12036\\
56000	12211\\
58000	12515\\
60000	12310\\
62000	12420\\
64000	12489\\
66000	12003\\
68000	12831\\
70000	12503\\
72000	12345\\
74000	11971\\
76000	12054\\
78000	12195\\
80000	12664\\
82000	12276\\
84000	12649\\
86000	12396\\
90000	12439\\
92000	12542\\
94000	12375\\
96000	12365\\
98000	12402\\
1e+05	12303\\
};
\addplot [color=mycolor26, forget plot]
  table[row sep=crcr]{%
0	4395.3\\
2000	7323.2\\
4000	8219.1\\
6000	8077.4\\
8000	8271.1\\
10000	8367.4\\
12000	8110.4\\
14000	8393.8\\
16000	8436.4\\
18000	8201.9\\
20000	8362.9\\
22000	8243.9\\
24000	8199.5\\
26000	8414.4\\
28000	7813.5\\
30000	8199\\
32000	8212.1\\
34000	8164.5\\
36000	8147.7\\
38000	8406.3\\
40000	8263.4\\
42000	8319.3\\
44000	8338.7\\
46000	8123.6\\
48000	8657.7\\
50000	8300.9\\
52000	8152.3\\
56000	8106.1\\
58000	8167.7\\
60000	8418.4\\
62000	8231\\
64000	8073\\
66000	7958.7\\
68000	8179.6\\
70000	8150.6\\
72000	7992.5\\
74000	8084.7\\
76000	8297.5\\
78000	8087.3\\
80000	8195.6\\
82000	8172.3\\
84000	8500.5\\
86000	8520.9\\
88000	8296.8\\
90000	7981.8\\
94000	8450.3\\
96000	8169.7\\
98000	8716.7\\
1e+05	8280.1\\
};
\addplot [color=mycolor27, forget plot]
  table[row sep=crcr]{%
0	4395.3\\
2000	9980.2\\
4000	9699.8\\
6000	9357.8\\
8000	9780.4\\
10000	10090\\
12000	9686.4\\
14000	9785.7\\
16000	9449.1\\
18000	9518.5\\
20000	9718\\
22000	9764.8\\
24000	9921.1\\
26000	9891.3\\
28000	9388.8\\
30000	9511.2\\
32000	9589.4\\
34000	10065\\
36000	9740.8\\
38000	9854.7\\
40000	9627.9\\
42000	9732.7\\
44000	9929\\
46000	9835.3\\
48000	9645.2\\
50000	9616.8\\
52000	9657.1\\
54000	10131\\
56000	9669.2\\
58000	9614.8\\
60000	9814\\
62000	9775.1\\
64000	9631.7\\
66000	10206\\
70000	9867.7\\
72000	9782.1\\
74000	9664.2\\
76000	9918.5\\
80000	9834.3\\
82000	9862.4\\
84000	9608.9\\
86000	10138\\
88000	9894.8\\
90000	9900.1\\
92000	10113\\
94000	9847.8\\
96000	9679\\
98000	9732.3\\
1e+05	9726.3\\
};
\addplot [color=mycolor28, forget plot]
  table[row sep=crcr]{%
0	4395.3\\
2000	5997\\
4000	5599.2\\
6000	5672.7\\
8000	5512.1\\
10000	5490.7\\
12000	5340.3\\
14000	5389.1\\
16000	5415.5\\
18000	5578.7\\
20000	5435.6\\
22000	5549\\
24000	5361.1\\
26000	5682.6\\
28000	5425.6\\
30000	5523.6\\
32000	5430.1\\
34000	5509.2\\
36000	5450.4\\
38000	5431.9\\
40000	5436.4\\
42000	5662.7\\
44000	5550.4\\
46000	5341.1\\
48000	5490.6\\
50000	5390.6\\
52000	5458.9\\
56000	5376.4\\
58000	5492.4\\
60000	5551.7\\
64000	5403\\
66000	5475.7\\
68000	5608.3\\
70000	5593.2\\
72000	5298.8\\
76000	5763.9\\
78000	5288.4\\
80000	5706.9\\
82000	5478\\
84000	5232\\
86000	5473\\
88000	5556.1\\
90000	5554.2\\
92000	5511.8\\
94000	5750\\
96000	5430.8\\
98000	5335.1\\
1e+05	5598.5\\
};
\addplot [color=mycolor29, forget plot]
  table[row sep=crcr]{%
0	4395.3\\
2000	11052\\
4000	11811\\
6000	11698\\
8000	11844\\
10000	11260\\
12000	11556\\
14000	11681\\
16000	11410\\
18000	12097\\
20000	11799\\
22000	11636\\
24000	11975\\
26000	11505\\
28000	11614\\
30000	12095\\
32000	11688\\
34000	11509\\
36000	11609\\
38000	11464\\
40000	11872\\
42000	11630\\
44000	11484\\
46000	11759\\
48000	11794\\
50000	11724\\
52000	11990\\
54000	11795\\
56000	11333\\
58000	12220\\
60000	11108\\
64000	11896\\
66000	11495\\
68000	11403\\
70000	11675\\
72000	11599\\
74000	11741\\
76000	11678\\
78000	11768\\
80000	11434\\
82000	11956\\
84000	11588\\
86000	11888\\
88000	11971\\
90000	11896\\
92000	11901\\
94000	11481\\
98000	11535\\
1e+05	11824\\
};
\addplot [color=mycolor30, forget plot]
  table[row sep=crcr]{%
0	4395.3\\
2000	7037.5\\
4000	7740.5\\
6000	7663.5\\
8000	7902.7\\
10000	7841.3\\
12000	7973.5\\
14000	7898.6\\
16000	8029.9\\
18000	7765.3\\
20000	7877.5\\
22000	7616.9\\
24000	7963.8\\
26000	7791.6\\
28000	7952.8\\
30000	7841\\
32000	7842.9\\
34000	7894.9\\
36000	7825.6\\
38000	7903.9\\
40000	7877.2\\
42000	7684.6\\
44000	7828.4\\
46000	7875.3\\
48000	7789.7\\
50000	7778\\
52000	7862.8\\
54000	7979.6\\
56000	7606.1\\
58000	7864.5\\
60000	7776.1\\
62000	7993.6\\
64000	7979.6\\
66000	8023.9\\
68000	7916.8\\
70000	7863.7\\
72000	7737.2\\
74000	7931.3\\
76000	7784.2\\
78000	7873.7\\
80000	7872.1\\
82000	8128.2\\
84000	7675.1\\
86000	7744.7\\
88000	8004.4\\
90000	7874.9\\
92000	7794.7\\
94000	8039.9\\
96000	7750.8\\
98000	7966.1\\
1e+05	7942.7\\
};
\addplot [color=mycolor31, forget plot]
  table[row sep=crcr]{%
0	4395.3\\
2000	6806.2\\
4000	6467.1\\
6000	6690.2\\
8000	6576.7\\
10000	6646.1\\
12000	6399.8\\
14000	6681.3\\
16000	6745.6\\
18000	6503.7\\
20000	6532\\
22000	6333.5\\
26000	6796.3\\
28000	6515.7\\
30000	6665.7\\
32000	6489.2\\
34000	6540.3\\
36000	6488.4\\
38000	6653.7\\
40000	6669.5\\
42000	6447.8\\
44000	6301.6\\
46000	6714.1\\
48000	6279.5\\
50000	6613.1\\
52000	6711.3\\
54000	6625.8\\
56000	6560.3\\
58000	6661.7\\
60000	6568.6\\
62000	6363\\
64000	6545.3\\
66000	6544.3\\
68000	6280\\
70000	6483.7\\
72000	6510.5\\
74000	6630.6\\
76000	6605.3\\
78000	6473.9\\
80000	6550.1\\
82000	6303.5\\
84000	6580.3\\
86000	6468.5\\
88000	6637.3\\
90000	6532.2\\
94000	6787.1\\
96000	6456.2\\
98000	6564.3\\
1e+05	6346.3\\
};
\addplot [color=mycolor32, forget plot]
  table[row sep=crcr]{%
0	4395.3\\
2000	10978\\
4000	11213\\
6000	11073\\
8000	10759\\
10000	10545\\
12000	10529\\
14000	10180\\
16000	10306\\
18000	10699\\
20000	10753\\
22000	10412\\
24000	10572\\
26000	10533\\
28000	10295\\
30000	10644\\
32000	10693\\
34000	10421\\
36000	10451\\
38000	10595\\
40000	10799\\
42000	10589\\
44000	10564\\
46000	10783\\
48000	10576\\
50000	10411\\
52000	10730\\
54000	10639\\
56000	10320\\
58000	10704\\
60000	10533\\
62000	10198\\
64000	10831\\
66000	10722\\
68000	10356\\
70000	10641\\
72000	10733\\
74000	10433\\
76000	10570\\
78000	10800\\
80000	10859\\
82000	10667\\
84000	10084\\
86000	10316\\
88000	10901\\
90000	11044\\
92000	10552\\
94000	10218\\
96000	10551\\
98000	10711\\
1e+05	10914\\
};
\addplot [color=mycolor33, forget plot]
  table[row sep=crcr]{%
0	4395.3\\
2000	7403.1\\
4000	8148.2\\
6000	7903.5\\
8000	8203.3\\
10000	7887.1\\
12000	8327.8\\
14000	8071.6\\
16000	8191.9\\
18000	8396\\
20000	8394.6\\
22000	8062.3\\
24000	8290.2\\
26000	8195.2\\
28000	8009.2\\
30000	7868.2\\
32000	8546.9\\
34000	8466.6\\
36000	8189.3\\
38000	8453.5\\
40000	8167.5\\
42000	7997.9\\
44000	8458.1\\
46000	8112.6\\
48000	8254.8\\
50000	8218.5\\
52000	8257.6\\
54000	8116.2\\
56000	8226.4\\
58000	8296.9\\
60000	8145.3\\
62000	8059.3\\
64000	8503.7\\
66000	8133.5\\
68000	8130.6\\
70000	8195.7\\
72000	8297.7\\
74000	8160.6\\
76000	7679.7\\
78000	8180.4\\
80000	8151.6\\
82000	8379.5\\
84000	8280.3\\
86000	7966.2\\
90000	8093.1\\
92000	8000.6\\
94000	8311.3\\
96000	8586.8\\
98000	8292.1\\
1e+05	8474\\
};
\addplot [color=mycolor34, forget plot]
  table[row sep=crcr]{%
0	4395.3\\
2000	4684.5\\
4000	3847.9\\
6000	3910.2\\
10000	4096.8\\
12000	4057.2\\
14000	4099.5\\
16000	4000.4\\
18000	3979.4\\
20000	3924.6\\
22000	3998.9\\
24000	3927.7\\
26000	4085.8\\
28000	3902.3\\
30000	3882.7\\
32000	4152.1\\
34000	3998.5\\
36000	4011.4\\
38000	4171.5\\
40000	3986.6\\
42000	4025.4\\
44000	4104.4\\
48000	4002.4\\
50000	3991.2\\
52000	4102.1\\
54000	3985.5\\
56000	4020.3\\
58000	3942.8\\
60000	4040\\
62000	3950.5\\
64000	4037\\
66000	3904.6\\
68000	3827.5\\
70000	4022.5\\
72000	3963.7\\
74000	4260.2\\
76000	4070.2\\
78000	3943.2\\
80000	4048.7\\
82000	4098.1\\
84000	4177.2\\
86000	3971.7\\
88000	4001.5\\
90000	4199.9\\
92000	4028.7\\
94000	4150.8\\
96000	3962.3\\
98000	3974.3\\
1e+05	4160\\
};
\addplot [color=mycolor35]
  table[row sep=crcr]{%
0	4395.3\\
2000	9482.2\\
4000	9999.5\\
6000	10302\\
8000	10110\\
10000	10269\\
12000	10164\\
14000	10324\\
16000	10072\\
18000	10279\\
20000	10445\\
22000	10261\\
24000	10429\\
26000	10079\\
28000	9973\\
32000	10136\\
34000	10240\\
36000	10120\\
38000	10231\\
42000	10235\\
44000	10440\\
46000	10328\\
48000	10108\\
50000	10445\\
52000	10310\\
54000	10337\\
56000	10024\\
58000	9889.6\\
60000	10458\\
64000	10206\\
66000	10251\\
68000	10054\\
72000	10011\\
74000	10162\\
76000	10070\\
78000	10192\\
80000	10160\\
84000	10217\\
86000	10178\\
88000	9988.6\\
90000	10065\\
92000	10378\\
94000	9982.4\\
96000	10281\\
98000	10250\\
1e+05	10148\\
};
\addlegendentry{\ClasskIILegend}

\addplot [color=mycolor36, forget plot]
  table[row sep=crcr]{%
0	4395.3\\
2000	4794.2\\
4000	4321.3\\
6000	4388.4\\
8000	4288.4\\
10000	4473\\
12000	4354.7\\
14000	4376.6\\
16000	4295.7\\
20000	4444.5\\
22000	4418.3\\
24000	4281.1\\
28000	4174.7\\
30000	4342.4\\
32000	4222.8\\
34000	4445.2\\
36000	4179.8\\
38000	4149.5\\
40000	4172.5\\
42000	4388.7\\
44000	4243.5\\
46000	4332.7\\
48000	4346.7\\
50000	4432.2\\
52000	4431.8\\
54000	4203.3\\
56000	4324.9\\
58000	4368.7\\
60000	4439.9\\
62000	4124.3\\
64000	4172.9\\
66000	4388.6\\
68000	4366.8\\
70000	4372.7\\
72000	4130.6\\
74000	4362\\
76000	4414.2\\
78000	4280.9\\
80000	4313.3\\
82000	4289.8\\
84000	4356.2\\
86000	4288.5\\
88000	4177.9\\
94000	4305.8\\
96000	4193.6\\
1e+05	4207.5\\
};
\addplot [color=mycolor37, forget plot]
  table[row sep=crcr]{%
0	4395.3\\
2000	3874.2\\
4000	3553.8\\
6000	3494.7\\
8000	3599.5\\
10000	3638.9\\
12000	3601.1\\
14000	3678.2\\
16000	3601.1\\
20000	3610.9\\
22000	3592.3\\
24000	3608.8\\
26000	3606.2\\
28000	3637.3\\
32000	3632.4\\
36000	3682.4\\
38000	3543.4\\
40000	3531.4\\
42000	3580.6\\
44000	3676.9\\
46000	3631.1\\
48000	3560.8\\
50000	3565.1\\
52000	3697.9\\
54000	3640.1\\
56000	3663.5\\
58000	3611.4\\
60000	3615.1\\
62000	3659.2\\
64000	3627.7\\
66000	3655.4\\
68000	3695\\
70000	3710\\
72000	3569.3\\
74000	3666.3\\
76000	3623.1\\
78000	3613.9\\
80000	3577.5\\
82000	3602.6\\
84000	3590.7\\
86000	3642.4\\
88000	3552\\
92000	3662.8\\
94000	3688.8\\
98000	3617.9\\
1e+05	3610.9\\
};
\addplot [color=mycolor38, forget plot]
  table[row sep=crcr]{%
0	4395.3\\
2000	3246.8\\
4000	3062.6\\
6000	2867.6\\
8000	3022\\
10000	3002.9\\
12000	2968\\
14000	2871.8\\
16000	3003.5\\
20000	2933.5\\
22000	3047.9\\
24000	2959.8\\
26000	2967.6\\
28000	3153.5\\
30000	2883.5\\
32000	3020.7\\
34000	2935.9\\
36000	2958\\
38000	3051.2\\
40000	3025.9\\
42000	2942.8\\
44000	3095.5\\
46000	3005.3\\
50000	3010.5\\
52000	2926.3\\
54000	2933\\
56000	2857.7\\
58000	3071.2\\
60000	2948.9\\
62000	2962.5\\
64000	2943.4\\
66000	2936.4\\
68000	2965.8\\
70000	3070.8\\
72000	3062.1\\
74000	2903.8\\
76000	3044.6\\
78000	3022.3\\
80000	2978\\
82000	3024.9\\
84000	2977.3\\
86000	3044.6\\
88000	2966.3\\
90000	2995.8\\
92000	2941.6\\
94000	2876.1\\
96000	3115.7\\
98000	3076.7\\
1e+05	2903\\
};
\addplot [color=mycolor39, forget plot]
  table[row sep=crcr]{%
0	4395.3\\
2000	12034\\
4000	12593\\
6000	11733\\
8000	11563\\
10000	11844\\
12000	11853\\
14000	11976\\
16000	11970\\
18000	11882\\
20000	11211\\
22000	11732\\
24000	11443\\
26000	11571\\
28000	11904\\
30000	12078\\
32000	12137\\
34000	12329\\
38000	11804\\
40000	11787\\
42000	11456\\
44000	11983\\
46000	11694\\
48000	11683\\
50000	12108\\
52000	11919\\
54000	11321\\
56000	12125\\
58000	12423\\
60000	11838\\
62000	11846\\
64000	11915\\
66000	12042\\
68000	11539\\
70000	11482\\
72000	12005\\
74000	11797\\
76000	12041\\
78000	11686\\
80000	11739\\
82000	11695\\
84000	11822\\
86000	12157\\
88000	11831\\
92000	11858\\
94000	11781\\
96000	11426\\
98000	11202\\
1e+05	11508\\
};
\addplot [color=mycolor40, forget plot]
  table[row sep=crcr]{%
0	4395.3\\
2000	5149\\
4000	5539.8\\
6000	5717.7\\
8000	5773.1\\
10000	5701.1\\
12000	5768.4\\
14000	5781.5\\
16000	5728.8\\
18000	5749.6\\
20000	5734.8\\
22000	5787.2\\
24000	5811.7\\
26000	5671.6\\
28000	5822.9\\
30000	5622.9\\
34000	5837.8\\
36000	5648.9\\
38000	5562.2\\
40000	5730.9\\
42000	5657.8\\
44000	5653.3\\
46000	5717.5\\
48000	5660.8\\
50000	5734.9\\
52000	5755.5\\
54000	5680.4\\
56000	5717.9\\
58000	5738.8\\
60000	5467\\
62000	5567.1\\
64000	5733.2\\
66000	5741.1\\
68000	5890.1\\
70000	5793.5\\
72000	5817.9\\
76000	5718.1\\
78000	5609.2\\
80000	5690.6\\
82000	5872.9\\
84000	5705.2\\
86000	5811.9\\
88000	5792.7\\
90000	5738.7\\
92000	5829.2\\
94000	5690\\
96000	5744.7\\
98000	5759.9\\
1e+05	5739.9\\
};
\addplot [color=mycolor41, forget plot]
  table[row sep=crcr]{%
0	4395.3\\
2000	4728.5\\
4000	5312.9\\
6000	5265.3\\
8000	5380\\
10000	5453.9\\
12000	5358.5\\
14000	5287.5\\
16000	5445.7\\
18000	5383.6\\
20000	5247.4\\
22000	5399\\
24000	5262.6\\
26000	5348.9\\
28000	5415.3\\
30000	5346.5\\
32000	5428.4\\
34000	5426.2\\
36000	5382.9\\
38000	5146.5\\
40000	5482\\
42000	5238.9\\
44000	5461.9\\
46000	5435.4\\
48000	5537.8\\
50000	5251.4\\
52000	5530.3\\
54000	5535.8\\
56000	5408.1\\
58000	5463.3\\
60000	5256.9\\
62000	5505.9\\
64000	5389\\
68000	5421.5\\
70000	5354\\
72000	5258.4\\
74000	5303.8\\
76000	5304.9\\
78000	5365.7\\
80000	5324.3\\
82000	5318\\
84000	5584.8\\
86000	5274.5\\
88000	5299.1\\
90000	5438.2\\
92000	5540.6\\
94000	5167.1\\
96000	5397.5\\
98000	5603\\
1e+05	5416.7\\
};
\addplot [color=mycolor42, forget plot]
  table[row sep=crcr]{%
0	4395.3\\
2000	15774\\
4000	19877\\
6000	21919\\
8000	20773\\
10000	20623\\
12000	20278\\
14000	21019\\
16000	21150\\
18000	20983\\
20000	20711\\
22000	20755\\
24000	20708\\
26000	21351\\
28000	20451\\
30000	20453\\
32000	20215\\
36000	20704\\
38000	20245\\
40000	20950\\
42000	20850\\
44000	20851\\
46000	21239\\
48000	20185\\
50000	20500\\
52000	21194\\
56000	20611\\
58000	20436\\
60000	20606\\
62000	21082\\
64000	21073\\
68000	20499\\
70000	20636\\
72000	20894\\
74000	21091\\
76000	21593\\
78000	20882\\
80000	20752\\
82000	21330\\
84000	20370\\
86000	21019\\
88000	20494\\
90000	20233\\
92000	20848\\
94000	20554\\
96000	21917\\
98000	20337\\
1e+05	20712\\
};
\addplot [color=mycolor43, forget plot]
  table[row sep=crcr]{%
0	4395.3\\
2000	5433.4\\
4000	5084.7\\
6000	4998.3\\
8000	5059.8\\
10000	4855.1\\
12000	5013.3\\
14000	5040.3\\
16000	5025.6\\
18000	5066.6\\
20000	5043.9\\
22000	4798.7\\
24000	4995.2\\
26000	4930.3\\
28000	4888.4\\
30000	4971.3\\
32000	4974.5\\
34000	5021.2\\
36000	5024.5\\
38000	4896.8\\
40000	4907.4\\
42000	4968.1\\
44000	4974.4\\
46000	4947.4\\
48000	4875.2\\
50000	4963.1\\
52000	4892.4\\
54000	5018.5\\
56000	4981.3\\
58000	5012.2\\
60000	4814.8\\
62000	4952.6\\
64000	4939.1\\
66000	4991.1\\
68000	4896.3\\
70000	4876.4\\
72000	5050.8\\
74000	4870.6\\
76000	4956\\
78000	4940.9\\
80000	5035.8\\
82000	4882.5\\
84000	4869.1\\
86000	5003.1\\
88000	4892.9\\
90000	4893.7\\
92000	5075.9\\
96000	4981.2\\
98000	4917\\
1e+05	5016.3\\
};
\addplot [color=mycolor44, forget plot]
  table[row sep=crcr]{%
0	4395.3\\
2000	4032.9\\
4000	3512.3\\
6000	3447.8\\
8000	3411\\
10000	3412.5\\
14000	3373.5\\
16000	3272.2\\
18000	3338.7\\
20000	3345.4\\
22000	3294.5\\
24000	3320.5\\
26000	3388.1\\
28000	3342.1\\
30000	3456.6\\
32000	3341\\
34000	3400.1\\
36000	3315.9\\
38000	3349.7\\
40000	3371.2\\
42000	3437.6\\
44000	3349.1\\
46000	3375.3\\
48000	3554.5\\
50000	3466.4\\
52000	3317.4\\
54000	3416.1\\
56000	3404.5\\
58000	3316.3\\
60000	3390.4\\
62000	3327.2\\
64000	3386.6\\
66000	3296\\
68000	3417\\
70000	3339.4\\
72000	3406.1\\
74000	3458.2\\
76000	3357.8\\
78000	3339.4\\
80000	3369.6\\
82000	3303.3\\
84000	3318.2\\
88000	3315.1\\
90000	3431.4\\
92000	3382.1\\
94000	3468.8\\
96000	3450.3\\
98000	3441.6\\
1e+05	3490.4\\
};
\addplot [color=mycolor45, forget plot]
  table[row sep=crcr]{%
0	4395.3\\
2000	6011.2\\
4000	6604.4\\
6000	6823.2\\
8000	6923.5\\
10000	6848.2\\
12000	7039.1\\
14000	6876.7\\
16000	6932.9\\
18000	6824.4\\
20000	7064.8\\
22000	6782.2\\
24000	6962.9\\
26000	6843.2\\
28000	6904.9\\
30000	6924.5\\
32000	7060.1\\
34000	7049.4\\
38000	6970.3\\
40000	7020.7\\
42000	6823.3\\
44000	7059.7\\
46000	6840.6\\
48000	7023.3\\
50000	6806.1\\
52000	6854.8\\
54000	6734.8\\
56000	6946.6\\
58000	6892.4\\
60000	6737.3\\
62000	6969.2\\
64000	6759.4\\
66000	7044\\
68000	6811.8\\
70000	6882\\
72000	6904.9\\
74000	6725.7\\
76000	6851.2\\
78000	6796.3\\
80000	6864.8\\
82000	6990.2\\
84000	6851.5\\
88000	6780.5\\
92000	7074.6\\
94000	6961.7\\
96000	6940.8\\
98000	6947.1\\
1e+05	6904.3\\
};
\addplot [color=mycolor46, forget plot]
  table[row sep=crcr]{%
0	4395.3\\
2000	3859.1\\
4000	3933.6\\
6000	4124.9\\
8000	4067.2\\
10000	3992.1\\
12000	4013\\
14000	3993.5\\
16000	4079.1\\
20000	4067\\
22000	4076.9\\
24000	4230.4\\
26000	4040.5\\
28000	4032.1\\
30000	4112.8\\
32000	4162.7\\
34000	3972\\
36000	4070.4\\
38000	4135.7\\
40000	4015.4\\
42000	4012.4\\
44000	4056.4\\
46000	4047.3\\
48000	4096.4\\
50000	4071.5\\
52000	4219.3\\
54000	4205.6\\
56000	4017\\
58000	3937.9\\
60000	4017.9\\
62000	3964.9\\
64000	4067.7\\
66000	3980.8\\
68000	4093.2\\
70000	4041\\
72000	4021\\
74000	4070.6\\
76000	3997.8\\
78000	4180.1\\
80000	4117.9\\
82000	4131.1\\
84000	4052.5\\
86000	4115.5\\
88000	4097.4\\
90000	4115.3\\
92000	4188.3\\
94000	4043.4\\
96000	3995.3\\
98000	3982.4\\
1e+05	4105.7\\
};
\addplot [color=mycolor47, forget plot]
  table[row sep=crcr]{%
0	4395.3\\
2000	5456.5\\
4000	5342.6\\
6000	5339.2\\
8000	5404.3\\
10000	5263.7\\
12000	5297.5\\
14000	5372.5\\
16000	5178.4\\
18000	5211\\
20000	5325.8\\
22000	5489.2\\
24000	5230.2\\
26000	5311.2\\
28000	5135.1\\
30000	5207.9\\
32000	5322.7\\
34000	5204.4\\
36000	5356.2\\
38000	5341.6\\
40000	5214.8\\
42000	5429.4\\
44000	5273.6\\
46000	5302.7\\
48000	5372.6\\
50000	5340.1\\
52000	5400.7\\
54000	5343\\
56000	5434.9\\
58000	5302.3\\
60000	5406.7\\
62000	5327.3\\
64000	5289.3\\
66000	5376.6\\
68000	5320.6\\
70000	5336.2\\
72000	5287.1\\
74000	5475.7\\
76000	5253.8\\
78000	5466.1\\
80000	5420.1\\
82000	5490.2\\
84000	5061.4\\
86000	5289\\
88000	5343.6\\
90000	5472\\
92000	5226.2\\
94000	5249.9\\
96000	5175.6\\
98000	5313.7\\
1e+05	5234.5\\
};
\addplot [color=mycolor48, forget plot]
  table[row sep=crcr]{%
0	4395.3\\
2000	3281.2\\
4000	2922\\
8000	2956.3\\
10000	2897.7\\
12000	2849.8\\
14000	2901.8\\
16000	2900\\
18000	2927.8\\
20000	2967.3\\
22000	2900.7\\
24000	2871.6\\
26000	2959.8\\
28000	2867.7\\
30000	2858.5\\
32000	2903\\
34000	2886.5\\
36000	2958.5\\
38000	2954.2\\
40000	2902.4\\
42000	2918.8\\
46000	2909.5\\
48000	2915.2\\
50000	2956.4\\
52000	2874.8\\
54000	2887.4\\
58000	2959.5\\
60000	2861.8\\
62000	2880.5\\
64000	2910.6\\
66000	2949.9\\
70000	2930.8\\
72000	2871\\
74000	2958.2\\
76000	2915.1\\
78000	2853.5\\
80000	2922.9\\
84000	2886.8\\
86000	2935.6\\
88000	2953.5\\
90000	2856\\
94000	2946.4\\
96000	2994.9\\
98000	2905.4\\
1e+05	2874.5\\
};
\addplot [color=mycolor49, forget plot]
  table[row sep=crcr]{%
0	4395.3\\
2000	4708.4\\
4000	4266.9\\
6000	4340.8\\
8000	4661.8\\
10000	4324\\
12000	4368.5\\
14000	4455.6\\
16000	4277.2\\
18000	4268\\
20000	4372\\
22000	4242\\
24000	4483\\
26000	4333.8\\
30000	4350.7\\
32000	4455.2\\
34000	4339.1\\
36000	4485\\
38000	4281.2\\
40000	4186.6\\
42000	4272.1\\
44000	4390.8\\
46000	4455.2\\
48000	4493.6\\
50000	4372.4\\
52000	4209.9\\
54000	4359.6\\
56000	4280\\
58000	4427.3\\
60000	4283\\
62000	4506\\
66000	4405.1\\
70000	4435\\
72000	4406.4\\
74000	4441\\
76000	4279.3\\
78000	4377.9\\
80000	4357.2\\
82000	4160.6\\
84000	4304.2\\
86000	4519.3\\
88000	4363.8\\
90000	4367.1\\
92000	4496.8\\
96000	4206.3\\
98000	4386.2\\
1e+05	4349.1\\
};
\addplot [color=mycolor50, forget plot]
  table[row sep=crcr]{%
0	4395.3\\
2000	2533.8\\
4000	2393.4\\
6000	2386\\
8000	2488.9\\
10000	2434.5\\
12000	2525.4\\
14000	2413.7\\
16000	2420.5\\
18000	2407.3\\
20000	2481.6\\
22000	2433.7\\
24000	2468\\
26000	2491.1\\
28000	2389.1\\
30000	2439.2\\
32000	2478.8\\
34000	2501.5\\
36000	2471.1\\
38000	2373.7\\
40000	2438.9\\
42000	2546.6\\
44000	2457.9\\
46000	2439.2\\
48000	2365.4\\
50000	2468.6\\
52000	2501.2\\
54000	2478.5\\
56000	2480.6\\
58000	2491.7\\
60000	2390.7\\
62000	2483.2\\
64000	2483.9\\
66000	2441.3\\
68000	2471.6\\
70000	2470.9\\
72000	2501.8\\
74000	2430.2\\
76000	2461.3\\
78000	2423\\
80000	2441.9\\
82000	2545.2\\
84000	2522.9\\
86000	2469.5\\
88000	2407.6\\
90000	2397.8\\
92000	2484.1\\
94000	2451.6\\
96000	2427.5\\
98000	2484.2\\
1e+05	2474.7\\
};
\addplot [color=mycolor51, forget plot]
  table[row sep=crcr]{%
0	4395.3\\
2000	5865.6\\
4000	5890.1\\
6000	5814.7\\
8000	5958.7\\
10000	5973.9\\
12000	6138.9\\
14000	5798.1\\
16000	5709.1\\
18000	5873.1\\
20000	5871\\
22000	5817.6\\
24000	6033.6\\
28000	6168.9\\
30000	6102.5\\
34000	5912.1\\
36000	5959.8\\
38000	6264.3\\
40000	5862.5\\
42000	5856.5\\
46000	6041.4\\
48000	6018.1\\
52000	5890.4\\
54000	6055.4\\
58000	5876.6\\
60000	5954.2\\
62000	5966.4\\
64000	6149.1\\
66000	6017.6\\
68000	5830.9\\
70000	6139.7\\
72000	5691.6\\
74000	6090.1\\
76000	6088.7\\
78000	5991.8\\
80000	5745.9\\
84000	5975.8\\
86000	5780.1\\
88000	5895.7\\
90000	5967.6\\
92000	6112.4\\
94000	5935.6\\
96000	6077.5\\
98000	6162.7\\
1e+05	5574.4\\
};
\addplot [color=mycolor52, forget plot]
  table[row sep=crcr]{%
0	4395.3\\
4000	3663.3\\
6000	3580.3\\
8000	3401.7\\
10000	3386.6\\
12000	3475\\
14000	3404.4\\
16000	3447.1\\
18000	3424.7\\
20000	3565.5\\
22000	3429.8\\
24000	3403.9\\
26000	3564\\
30000	3432.7\\
34000	3653.6\\
36000	3537\\
38000	3587.1\\
40000	3409.1\\
42000	3540.6\\
44000	3531\\
46000	3536.2\\
48000	3621.4\\
50000	3456.3\\
52000	3518.5\\
54000	3537.2\\
56000	3427.2\\
58000	3496.3\\
60000	3586.9\\
62000	3416.7\\
64000	3549.9\\
66000	3559.2\\
70000	3482.7\\
72000	3462.8\\
74000	3589.7\\
76000	3593.7\\
78000	3656.2\\
80000	3514.8\\
82000	3399.2\\
84000	3568.5\\
86000	3403\\
88000	3604.6\\
90000	3536\\
92000	3517.8\\
94000	3445.4\\
96000	3626.9\\
98000	3701.7\\
1e+05	3492.7\\
};
\addplot [color=mycolor53, forget plot]
  table[row sep=crcr]{%
0	4395.3\\
2000	6394.4\\
4000	6354.5\\
6000	5950.3\\
8000	5651\\
10000	5583.8\\
12000	5466\\
14000	5571.7\\
16000	5737.4\\
18000	5523.4\\
20000	6115.8\\
22000	5813.4\\
24000	5780.4\\
26000	5790.1\\
28000	5724.9\\
30000	5590.1\\
32000	5921.6\\
34000	5741.4\\
36000	5951.7\\
38000	5655.7\\
40000	5645.6\\
42000	5785.6\\
44000	5685.2\\
46000	5739.3\\
48000	5880.9\\
50000	5602.2\\
52000	5843.3\\
54000	5640.5\\
56000	5919.7\\
58000	5922.3\\
60000	5728.8\\
62000	5706.7\\
64000	5712.7\\
66000	5555\\
68000	5712.3\\
70000	5512.2\\
72000	5824.9\\
74000	5712.6\\
76000	5799.1\\
78000	5638.1\\
80000	5823.4\\
82000	5781.4\\
84000	5804.2\\
86000	5693.9\\
88000	5663.5\\
90000	5764.8\\
92000	5733.2\\
94000	5666.2\\
96000	5576.6\\
98000	5699.5\\
1e+05	5801.5\\
};
\addplot [color=mycolor54, forget plot]
  table[row sep=crcr]{%
0	4395.3\\
2000	4520.9\\
4000	4455.7\\
6000	4732.1\\
8000	4585.7\\
10000	4473.7\\
12000	4529.2\\
14000	4444.6\\
16000	4642.4\\
18000	4677\\
20000	4686.8\\
22000	4617\\
24000	4512.8\\
26000	4538.8\\
28000	4620.9\\
30000	4627.2\\
32000	4800.1\\
34000	4510.2\\
36000	4458.7\\
38000	4685.7\\
40000	4374\\
42000	4621.8\\
44000	4627.2\\
46000	4727\\
48000	4651.2\\
52000	4769.3\\
54000	4404.4\\
56000	4648.9\\
58000	4600\\
60000	4655.3\\
62000	4657.1\\
64000	4694.6\\
66000	4529.3\\
70000	4628.5\\
72000	4634.9\\
74000	4605.2\\
76000	4718.3\\
78000	4522.4\\
80000	4634.5\\
82000	4641.2\\
84000	4693.6\\
86000	4478.6\\
88000	4554\\
90000	4708.8\\
92000	4622\\
94000	4589.8\\
96000	4727\\
98000	4743\\
1e+05	4464\\
};
\addplot [color=mycolor55, forget plot]
  table[row sep=crcr]{%
0	4395.3\\
2000	5383.3\\
4000	5736.2\\
8000	5846.3\\
10000	5653.1\\
12000	5910.3\\
14000	5823.2\\
16000	5711\\
18000	5761.8\\
20000	5794.9\\
22000	5637.6\\
24000	5975.9\\
26000	5593.1\\
28000	5789.7\\
32000	5847.5\\
36000	5638\\
40000	5729.6\\
42000	5858.3\\
46000	5525.1\\
48000	5751.6\\
50000	5715.6\\
52000	5781.8\\
54000	5732.5\\
56000	5771.4\\
58000	5742\\
60000	5875\\
62000	5750.2\\
64000	5683.5\\
66000	5924.7\\
68000	5839.9\\
70000	5724.1\\
72000	5770\\
74000	5693.4\\
76000	5737.6\\
78000	5732.3\\
80000	5938.1\\
82000	5699.9\\
86000	5661.6\\
88000	5760.3\\
90000	5917\\
92000	5570.5\\
94000	5687.4\\
96000	5658.7\\
98000	5740.5\\
1e+05	5868.7\\
};
\addplot [color=mycolor56, forget plot]
  table[row sep=crcr]{%
0	4395.3\\
2000	4273.7\\
4000	4041.1\\
6000	3995.9\\
8000	3985.3\\
10000	3945\\
12000	4006\\
14000	3946.9\\
16000	3980.9\\
18000	3988.1\\
20000	3907.6\\
22000	3962.5\\
24000	4031.3\\
26000	3966.8\\
28000	3867.5\\
30000	3918.2\\
32000	4041.7\\
34000	4152.6\\
36000	3966.6\\
38000	3992.5\\
40000	3988.3\\
42000	4046.7\\
44000	3967.6\\
46000	4006.7\\
48000	3969.5\\
50000	4066.5\\
52000	3999.6\\
54000	3962.1\\
56000	3943.6\\
58000	3824.5\\
60000	3970.8\\
62000	4051.4\\
64000	4055.5\\
66000	4003.4\\
68000	3918.4\\
70000	3957.7\\
72000	4047.9\\
74000	3929.7\\
76000	3955\\
78000	3994.4\\
80000	3922.2\\
82000	4033.2\\
84000	3941.1\\
86000	3926.7\\
88000	4001.1\\
90000	3869.8\\
92000	4020.7\\
94000	3929.7\\
96000	4005.5\\
98000	4037.5\\
1e+05	4026.5\\
};
\addplot [color=mycolor57, forget plot]
  table[row sep=crcr]{%
0	4395.3\\
2000	6459.4\\
4000	6957.6\\
6000	6753.7\\
8000	6802.9\\
10000	6502.3\\
12000	6598.8\\
14000	6919.1\\
16000	6892.8\\
18000	6818.4\\
20000	6694.3\\
22000	6767.8\\
24000	6710.3\\
26000	6751.7\\
28000	7002.6\\
30000	6830.2\\
32000	6877.9\\
34000	6884\\
38000	6676.5\\
40000	6752.7\\
42000	6724.9\\
44000	6536.3\\
46000	6663\\
48000	6613.8\\
50000	6626.2\\
52000	6499.5\\
54000	6656.6\\
56000	6595.5\\
58000	6753\\
60000	6745.3\\
62000	6637.6\\
64000	6683.1\\
66000	6681.5\\
68000	6606.7\\
70000	6790.6\\
72000	6720.2\\
74000	6744\\
76000	6531.6\\
78000	6529.6\\
80000	6715.3\\
82000	6711.5\\
84000	6865.3\\
86000	6636.3\\
88000	6745.6\\
90000	6835.5\\
92000	6699.5\\
94000	6723.8\\
96000	6845.6\\
98000	6830.6\\
1e+05	6775.2\\
};
\addplot [color=mycolor58, forget plot]
  table[row sep=crcr]{%
0	4395.3\\
2000	3856.5\\
4000	3541.6\\
6000	3541.7\\
8000	3569.1\\
10000	3800.7\\
12000	3586.1\\
14000	3541.7\\
16000	3612.2\\
20000	3554.8\\
22000	3648.9\\
24000	3647\\
26000	3580.4\\
28000	3542\\
30000	3692.3\\
32000	3703.6\\
34000	3622.6\\
36000	3615.5\\
38000	3647.9\\
40000	3646.7\\
42000	3497.3\\
44000	3565.8\\
46000	3496.3\\
48000	3592.2\\
50000	3537.7\\
52000	3553.7\\
54000	3582.3\\
58000	3677.7\\
60000	3636.2\\
62000	3496.6\\
64000	3554.3\\
66000	3572.4\\
68000	3458.4\\
70000	3574.7\\
72000	3535.7\\
74000	3572.9\\
76000	3598.7\\
78000	3581.3\\
80000	3513.8\\
82000	3704.4\\
84000	3451.8\\
86000	3500.5\\
88000	3618.1\\
90000	3513.5\\
92000	3571.1\\
94000	3577.2\\
96000	3676.1\\
98000	3526.7\\
1e+05	3527.5\\
};
\addplot [color=mycolor59, forget plot]
  table[row sep=crcr]{%
0	4395.3\\
2000	4638\\
4000	4958\\
6000	4836.5\\
8000	5106.7\\
10000	5072.9\\
12000	5023.3\\
14000	5057.3\\
16000	5116.1\\
18000	5158.6\\
20000	4987.4\\
22000	5076.9\\
26000	5018.3\\
28000	5046\\
30000	5000.3\\
32000	5057.9\\
34000	5053.9\\
36000	5111.7\\
38000	5083.8\\
40000	4996.9\\
42000	5076.9\\
44000	5120.8\\
46000	5061.5\\
48000	5173.9\\
50000	4945.1\\
52000	5041.4\\
54000	5197\\
56000	5055.2\\
58000	5284.4\\
60000	4956.8\\
62000	5189\\
64000	4923.4\\
66000	5054.7\\
68000	5075.3\\
70000	5036.6\\
72000	5050.1\\
74000	5024.1\\
76000	4969.1\\
78000	4953.3\\
80000	5170.8\\
82000	5007.9\\
84000	5182.8\\
86000	5012.3\\
88000	5064.3\\
90000	4867.3\\
92000	5023.1\\
94000	5071.5\\
96000	5187.2\\
98000	5048.1\\
1e+05	5046.3\\
};
\addplot [color=mycolor60, forget plot]
  table[row sep=crcr]{%
0	4395.3\\
2000	3790.8\\
4000	3637.3\\
6000	3627\\
8000	3469.8\\
10000	3573.7\\
14000	3535.9\\
16000	3550.5\\
18000	3508.4\\
20000	3519.7\\
22000	3548.7\\
24000	3515.3\\
26000	3523.5\\
28000	3463.6\\
30000	3522.1\\
32000	3533.3\\
34000	3623.8\\
36000	3520.9\\
38000	3607.3\\
40000	3596.6\\
42000	3556.8\\
44000	3489.5\\
46000	3512.8\\
48000	3509.8\\
50000	3569.2\\
54000	3458.5\\
56000	3526.5\\
58000	3523\\
60000	3611.3\\
64000	3559.4\\
66000	3466.7\\
68000	3441\\
70000	3534.3\\
72000	3505.6\\
74000	3600.9\\
76000	3539.2\\
78000	3573.7\\
80000	3589.6\\
82000	3427.2\\
84000	3537.7\\
88000	3564.8\\
90000	3591\\
92000	3602.9\\
94000	3632\\
96000	3574.7\\
98000	3561.5\\
1e+05	3563.1\\
};
\addplot [color=mycolor61, forget plot]
  table[row sep=crcr]{%
0	4395.3\\
2000	3261.9\\
4000	3076.1\\
6000	3119.5\\
8000	3085.4\\
10000	3131.6\\
12000	3058\\
14000	3108.9\\
16000	3150.7\\
18000	3110.6\\
20000	3088.2\\
22000	3025.2\\
24000	3116.8\\
26000	3123.7\\
28000	3115.1\\
30000	3153.6\\
32000	3122\\
34000	3128.5\\
36000	3190\\
38000	3070.1\\
40000	3085.1\\
42000	3114.3\\
44000	3094.1\\
46000	3128.8\\
48000	3096.2\\
50000	3095.9\\
52000	3114.5\\
54000	3143.9\\
56000	3037.6\\
58000	3034.1\\
60000	3129.6\\
62000	3101.2\\
64000	3163.5\\
66000	3079.1\\
70000	3082.2\\
72000	3120.1\\
74000	3063.5\\
76000	3089.2\\
80000	3111.7\\
84000	3115.4\\
86000	3075.7\\
88000	3165.3\\
90000	3070.2\\
92000	3149.4\\
94000	3097.6\\
96000	3139.6\\
98000	3118.7\\
1e+05	3156.1\\
};
\addplot [color=mycolor62, forget plot]
  table[row sep=crcr]{%
0	4395.3\\
2000	4743.2\\
4000	5054.4\\
6000	5003.4\\
8000	4915.7\\
10000	4878.2\\
12000	4925\\
14000	5111.8\\
16000	4996.3\\
18000	5028.5\\
20000	5045.6\\
22000	4963.4\\
24000	4908.4\\
26000	5042.3\\
28000	4936.1\\
30000	5024.8\\
32000	4975.9\\
34000	5014.1\\
38000	4973.1\\
40000	4953.6\\
42000	4829\\
44000	4907.9\\
46000	4921.9\\
48000	4974.6\\
50000	4978.9\\
52000	5013.7\\
54000	4910.5\\
56000	5029.7\\
58000	4928.1\\
60000	4945.2\\
62000	4990.7\\
64000	4990.3\\
66000	4962.2\\
68000	4890.7\\
70000	4998.8\\
72000	4924.6\\
76000	4990.6\\
78000	4963.5\\
80000	4841.5\\
82000	4900\\
86000	5107.6\\
88000	5010.6\\
90000	4945\\
92000	4978.2\\
94000	4969\\
96000	4813.3\\
98000	5006.8\\
1e+05	4995.1\\
};
\addplot [color=mycolor63, forget plot]
  table[row sep=crcr]{%
0	4395.3\\
2000	2729.5\\
4000	2405\\
6000	2394.4\\
10000	2433.8\\
12000	2447.4\\
18000	2370.4\\
20000	2416.1\\
22000	2416.3\\
24000	2439.4\\
26000	2395.4\\
28000	2360.3\\
30000	2393.3\\
32000	2374\\
34000	2373\\
36000	2440.6\\
38000	2434.2\\
40000	2362.4\\
42000	2361\\
44000	2393.9\\
46000	2414.5\\
48000	2388\\
50000	2424.6\\
52000	2389.8\\
54000	2391.6\\
56000	2352.7\\
58000	2369.9\\
60000	2432.9\\
62000	2421.5\\
64000	2361.8\\
66000	2358.7\\
68000	2419.8\\
70000	2391.3\\
72000	2417.4\\
74000	2393.4\\
76000	2437.2\\
78000	2394\\
80000	2388.4\\
82000	2395.9\\
84000	2433.7\\
86000	2385.7\\
88000	2388.6\\
90000	2377.7\\
92000	2410.3\\
94000	2418.2\\
98000	2368.6\\
1e+05	2407\\
};
\addplot [color=mycolor64, forget plot]
  table[row sep=crcr]{%
0	4395.3\\
2000	3066\\
4000	2997.4\\
8000	3044.5\\
12000	3044.3\\
14000	3011.7\\
16000	3064.2\\
18000	3037.4\\
22000	3023.5\\
24000	3061.3\\
28000	3037.2\\
32000	3046\\
34000	3003.9\\
36000	2993.7\\
38000	3051.7\\
40000	3036.3\\
42000	3003.4\\
44000	2934.9\\
48000	3002.5\\
50000	3074\\
52000	3074.2\\
56000	3006.3\\
58000	3032.4\\
64000	3014.8\\
66000	3059\\
68000	3025.1\\
70000	3046.1\\
72000	2999.5\\
74000	3015.4\\
76000	2983.5\\
78000	2995.7\\
80000	3032.3\\
82000	3107.9\\
84000	3108.4\\
88000	3025.1\\
90000	3027.6\\
92000	2978.9\\
94000	3020\\
96000	3071.4\\
98000	3105.4\\
1e+05	2928\\
};
\addplot [color=mycolor65, forget plot]
  table[row sep=crcr]{%
0	4395.3\\
2000	2756.5\\
4000	2502.8\\
6000	2488.9\\
8000	2540\\
10000	2567.8\\
12000	2575.2\\
14000	2496.2\\
16000	2471.6\\
18000	2490\\
20000	2569.8\\
22000	2536.8\\
30000	2537.1\\
32000	2521.6\\
34000	2533.1\\
36000	2498.5\\
38000	2487.4\\
40000	2524.2\\
44000	2516.9\\
46000	2558.2\\
48000	2511.6\\
50000	2546.1\\
52000	2528.1\\
54000	2563.3\\
56000	2458.9\\
60000	2557.2\\
64000	2560.2\\
66000	2563.5\\
68000	2543.7\\
72000	2550.5\\
74000	2483.9\\
76000	2498.3\\
78000	2544.6\\
80000	2488.3\\
82000	2539.9\\
84000	2561.1\\
86000	2530.7\\
88000	2546\\
90000	2499.4\\
92000	2514.2\\
94000	2519.9\\
98000	2497.1\\
1e+05	2507\\
};
\addplot [color=mycolor66, forget plot]
  table[row sep=crcr]{%
0	4395.3\\
2000	3430.6\\
4000	3417.8\\
6000	3427.3\\
8000	3449.7\\
10000	3451\\
12000	3401.9\\
14000	3366.1\\
16000	3411.9\\
20000	3428.2\\
22000	3407\\
24000	3413.7\\
26000	3469.7\\
28000	3372\\
30000	3426.7\\
32000	3403.6\\
34000	3396.7\\
36000	3453.3\\
38000	3433.9\\
40000	3443.4\\
42000	3387.7\\
44000	3373.2\\
46000	3414\\
48000	3439.5\\
50000	3425.8\\
52000	3397\\
56000	3437.1\\
58000	3414.2\\
60000	3360.5\\
62000	3424.2\\
64000	3351.7\\
66000	3477.5\\
68000	3387.7\\
72000	3411.2\\
74000	3378.5\\
76000	3381.7\\
78000	3359.9\\
80000	3496.1\\
82000	3413.7\\
84000	3383\\
86000	3402.8\\
88000	3358.9\\
90000	3423\\
92000	3438.9\\
94000	3475.6\\
96000	3371.3\\
98000	3413.4\\
1e+05	3391.1\\
};
\addplot [color=mycolor67, forget plot]
  table[row sep=crcr]{%
0	4395.3\\
2000	3475.3\\
4000	3227.8\\
6000	3131.1\\
8000	3197.3\\
10000	3150.9\\
12000	3139.5\\
14000	3139.8\\
16000	3174.4\\
18000	3125.7\\
22000	3171.9\\
24000	3168.6\\
26000	3216.3\\
30000	3147.4\\
32000	3182.1\\
34000	3140.8\\
36000	3219.8\\
38000	3106.8\\
40000	3248.5\\
42000	3218.5\\
44000	3203.7\\
46000	3204.4\\
48000	3168.6\\
52000	3175.7\\
54000	3216.8\\
56000	3156.7\\
58000	3172.9\\
60000	3107.6\\
62000	3159.6\\
64000	3163.3\\
66000	3234\\
68000	3148.2\\
70000	3217.8\\
72000	3253.7\\
74000	3179.4\\
78000	3182.7\\
80000	3150.6\\
82000	3165.4\\
84000	3208.6\\
86000	3151\\
88000	3216.4\\
90000	3180.9\\
94000	3138.6\\
96000	3137.1\\
98000	3144.9\\
1e+05	3207.9\\
};
\addplot [color=mycolor67]
  table[row sep=crcr]{%
0	4395.3\\
2000	2747.2\\
4000	2625.6\\
6000	2573.2\\
8000	2582.1\\
10000	2557.7\\
14000	2620.8\\
16000	2607.1\\
18000	2640.8\\
20000	2626.1\\
22000	2655.3\\
24000	2617\\
26000	2587.2\\
28000	2635.8\\
30000	2553.4\\
32000	2630.4\\
34000	2598\\
38000	2609.6\\
40000	2642.8\\
42000	2606.3\\
44000	2587.1\\
48000	2586.3\\
50000	2599.3\\
52000	2653.4\\
54000	2611.2\\
56000	2608.6\\
58000	2586.9\\
60000	2593.1\\
62000	2577.7\\
64000	2609.1\\
66000	2627.9\\
72000	2561.9\\
74000	2623.3\\
76000	2618.9\\
78000	2638.8\\
80000	2599.9\\
82000	2572.9\\
84000	2632.5\\
86000	2649.2\\
88000	2608.7\\
90000	2612.4\\
92000	2575.6\\
94000	2563\\
96000	2564.6\\
98000	2606.3\\
1e+05	2614.5\\
};
\addlegendentry{\ClasskIIILegend}

\addplot [color=mycolor68, forget plot]
  table[row sep=crcr]{%
0	4395.3\\
2000	2845\\
4000	2550.1\\
6000	2483.8\\
8000	2454.2\\
10000	2409.8\\
12000	2440.1\\
14000	2416\\
16000	2446.2\\
18000	2426.4\\
22000	2448.3\\
24000	2429.4\\
26000	2463.2\\
28000	2472.2\\
30000	2465.8\\
34000	2438.3\\
36000	2442.5\\
38000	2478.2\\
40000	2455.7\\
42000	2464.6\\
44000	2457.1\\
48000	2471.2\\
50000	2442.8\\
52000	2431.8\\
54000	2467.5\\
56000	2424.3\\
58000	2444.1\\
60000	2438.3\\
62000	2415.6\\
64000	2479.7\\
66000	2438.3\\
68000	2445.6\\
70000	2480\\
74000	2466\\
76000	2427.4\\
80000	2503.7\\
82000	2413.6\\
84000	2457.6\\
86000	2419.3\\
90000	2454.2\\
92000	2448.1\\
94000	2465.6\\
96000	2445.6\\
98000	2445.9\\
1e+05	2454.5\\
};
\addplot [color=mycolor69, forget plot]
  table[row sep=crcr]{%
0	4395.3\\
2000	3065.8\\
4000	2789.9\\
6000	2842.7\\
10000	2923.8\\
12000	2895\\
14000	2905.7\\
16000	2811.1\\
18000	2872.3\\
20000	2872.1\\
24000	2894.9\\
26000	2857.3\\
28000	2893.4\\
30000	2879.2\\
32000	2941.5\\
34000	2901.7\\
36000	2906.6\\
38000	2889.8\\
40000	2890.8\\
42000	2922.1\\
44000	2890.2\\
50000	2855.1\\
52000	2881.8\\
54000	2863.1\\
56000	2910.2\\
58000	2908.6\\
60000	2883.9\\
62000	2879.3\\
64000	2910\\
66000	2864.1\\
68000	2933.1\\
70000	2909.9\\
72000	2920.3\\
74000	2848.1\\
76000	2919.8\\
78000	2887.7\\
80000	2901.7\\
82000	2945.9\\
84000	2852.1\\
86000	2901.4\\
88000	2867.6\\
90000	2847.2\\
92000	2850.1\\
94000	2879.5\\
96000	2858.2\\
98000	2880.4\\
1e+05	2894.2\\
};
\addplot [color=mycolor70, forget plot]
  table[row sep=crcr]{%
0	4395.3\\
2000	2252.8\\
4000	2041.9\\
6000	2078\\
8000	2075.4\\
10000	2053.9\\
16000	2101.3\\
18000	2087.8\\
20000	2039.9\\
22000	2053.1\\
24000	2046.7\\
26000	2123.2\\
28000	2034.2\\
30000	2065.6\\
32000	2065.6\\
34000	2056.3\\
36000	2104\\
38000	2086.8\\
40000	2042.5\\
42000	2085\\
44000	2055.7\\
46000	2121.8\\
48000	2136.5\\
50000	2096.8\\
52000	2035.2\\
54000	2054.9\\
56000	2066.2\\
58000	2035\\
60000	2106.7\\
62000	2117.8\\
64000	2090.8\\
66000	2078.4\\
68000	2133.4\\
70000	2087.8\\
72000	2095\\
74000	2043.8\\
76000	2042.4\\
78000	2095.5\\
80000	2086.6\\
82000	2147.6\\
84000	2053.2\\
86000	2097.4\\
88000	2105.6\\
90000	2071.9\\
92000	2110.2\\
94000	2065.2\\
96000	2129.6\\
98000	2135.4\\
1e+05	2120.4\\
};
\addplot [color=mycolor71, forget plot]
  table[row sep=crcr]{%
0	4395.3\\
2000	2763.6\\
4000	2512.6\\
8000	2480.4\\
10000	2491.5\\
12000	2482.7\\
14000	2454.5\\
16000	2462.1\\
18000	2456.3\\
20000	2459.6\\
26000	2439.4\\
28000	2458.8\\
30000	2492.5\\
34000	2462.3\\
36000	2455.7\\
40000	2469.9\\
42000	2497.2\\
46000	2469\\
48000	2483.2\\
50000	2490.2\\
52000	2433.4\\
54000	2443.5\\
56000	2480.3\\
60000	2482.4\\
62000	2497.8\\
64000	2445.9\\
66000	2467.9\\
68000	2480.7\\
70000	2457.4\\
72000	2485.9\\
74000	2485\\
76000	2505.4\\
80000	2458.3\\
82000	2441.3\\
84000	2484.3\\
86000	2468.2\\
88000	2463.7\\
90000	2500.8\\
92000	2468.6\\
1e+05	2483.8\\
};
\addplot [color=mycolor72, forget plot]
  table[row sep=crcr]{%
0	4395.3\\
2000	2574.5\\
4000	2160.4\\
6000	2217\\
8000	2184.6\\
10000	2215.2\\
12000	2171.6\\
16000	2165.4\\
18000	2180.1\\
20000	2167.6\\
22000	2226.4\\
24000	2092.7\\
26000	2140.8\\
28000	2196.8\\
30000	2138.5\\
32000	2137.5\\
34000	2190.1\\
36000	2209.9\\
38000	2183\\
40000	2185.5\\
42000	2154.9\\
44000	2164.3\\
46000	2194.3\\
48000	2176.3\\
50000	2148.7\\
52000	2178.7\\
54000	2175.3\\
58000	2222.9\\
60000	2187.8\\
62000	2161\\
66000	2183.4\\
68000	2163.7\\
70000	2181.3\\
72000	2207.4\\
74000	2159.7\\
76000	2215.3\\
78000	2183.6\\
80000	2251.8\\
82000	2132\\
84000	2184.9\\
86000	2193.8\\
88000	2187.2\\
90000	2207\\
92000	2148.8\\
94000	2169.6\\
96000	2144.7\\
98000	2134.3\\
1e+05	2192.4\\
};
\addplot [color=mycolor73, forget plot]
  table[row sep=crcr]{%
0	4395.3\\
2000	2865\\
4000	2685.7\\
6000	2661.7\\
8000	2589.6\\
10000	2617.9\\
12000	2587.4\\
14000	2610.7\\
16000	2605.2\\
18000	2580.9\\
20000	2573.5\\
22000	2611.5\\
24000	2605.4\\
26000	2580.3\\
28000	2592.2\\
32000	2560.1\\
34000	2608.4\\
36000	2614.1\\
38000	2590.5\\
40000	2618.9\\
42000	2592.9\\
44000	2593.3\\
46000	2625.5\\
48000	2635.6\\
50000	2623.2\\
52000	2564.3\\
54000	2594.8\\
56000	2592.2\\
58000	2618.9\\
60000	2585\\
62000	2578.2\\
64000	2593.1\\
66000	2627.9\\
68000	2585.5\\
70000	2637\\
72000	2596.7\\
74000	2605.2\\
76000	2595.7\\
78000	2551.4\\
80000	2591.6\\
82000	2576.2\\
84000	2610\\
86000	2586.9\\
88000	2572.1\\
92000	2587.4\\
94000	2586\\
96000	2635.9\\
98000	2601.3\\
1e+05	2599.8\\
};
\addplot [color=mycolor74, forget plot]
  table[row sep=crcr]{%
0	4395.3\\
2000	1951.7\\
4000	1505.7\\
6000	1516.9\\
8000	1514.6\\
10000	1537.6\\
12000	1502.7\\
14000	1537.3\\
16000	1541.7\\
18000	1537.7\\
20000	1504.8\\
22000	1496.7\\
24000	1517.8\\
26000	1558.9\\
28000	1524\\
32000	1514.5\\
34000	1524.6\\
36000	1508.5\\
38000	1531.7\\
40000	1502.2\\
42000	1533.4\\
44000	1506.8\\
46000	1533.5\\
48000	1515.7\\
50000	1518.8\\
52000	1534.6\\
54000	1538.9\\
56000	1531.9\\
58000	1490.8\\
60000	1493\\
62000	1508.7\\
64000	1504.4\\
66000	1537.9\\
68000	1536.4\\
70000	1518\\
72000	1542.4\\
74000	1500.2\\
76000	1531.8\\
78000	1539.5\\
82000	1498.7\\
84000	1521.9\\
88000	1536.5\\
90000	1500.5\\
92000	1497.2\\
94000	1512.3\\
96000	1512.9\\
98000	1504.3\\
1e+05	1558.6\\
};
\addplot [color=mycolor75, forget plot]
  table[row sep=crcr]{%
0	4395.3\\
2000	2553.4\\
4000	2229.4\\
6000	2177.2\\
8000	2193.2\\
10000	2201.9\\
12000	2217.2\\
14000	2140.6\\
16000	2178.5\\
18000	2244.6\\
20000	2181.9\\
22000	2218.1\\
24000	2208.9\\
26000	2167\\
28000	2203.6\\
32000	2184.8\\
34000	2183.7\\
36000	2213.8\\
38000	2155.8\\
42000	2159.1\\
44000	2191.3\\
48000	2185.6\\
50000	2181.3\\
52000	2220.9\\
56000	2190.4\\
58000	2195.9\\
60000	2218.3\\
62000	2167.6\\
64000	2189\\
66000	2193.3\\
68000	2185\\
70000	2191.3\\
74000	2228.5\\
76000	2228.6\\
78000	2213.5\\
80000	2177.8\\
82000	2191.7\\
84000	2184\\
86000	2155.2\\
88000	2207.4\\
90000	2176.2\\
92000	2166.7\\
94000	2196.5\\
96000	2181\\
98000	2186.9\\
1e+05	2250.7\\
};
\addplot [color=mycolor76, forget plot]
  table[row sep=crcr]{%
0	4395.3\\
2000	2371.3\\
4000	2043.2\\
6000	1992.7\\
8000	1965.4\\
10000	2003\\
12000	1997.3\\
14000	1998.1\\
16000	1951.6\\
18000	2019.2\\
20000	1996.3\\
24000	1995.2\\
26000	1970.6\\
28000	1991.8\\
30000	1993.6\\
32000	2005.4\\
34000	2009.1\\
36000	1992.7\\
38000	1969.1\\
40000	1998.8\\
44000	2000\\
46000	1963\\
50000	2002.9\\
52000	1980.9\\
54000	2003.1\\
56000	2017.4\\
58000	1978.7\\
62000	2031.6\\
64000	2025.6\\
66000	2054.4\\
68000	1971.8\\
70000	2018.2\\
74000	1981.9\\
76000	2007.5\\
78000	2002.4\\
80000	1968.4\\
82000	2007.1\\
84000	1997.3\\
86000	1935\\
88000	1986.7\\
90000	1986.8\\
92000	1995.8\\
94000	1990.2\\
96000	2008.7\\
1e+05	1986.5\\
};
\addplot [color=mycolor77, forget plot]
  table[row sep=crcr]{%
0	4395.3\\
2000	2076.3\\
4000	1592.1\\
8000	1554.2\\
12000	1525.5\\
14000	1550\\
16000	1548.6\\
20000	1569.3\\
22000	1535\\
24000	1561.8\\
26000	1570.4\\
30000	1540.4\\
32000	1553\\
34000	1571.2\\
36000	1555.4\\
40000	1559.9\\
48000	1531\\
50000	1547.3\\
52000	1555.4\\
54000	1568.1\\
56000	1547.4\\
62000	1556.7\\
64000	1543.9\\
66000	1559.9\\
68000	1538.8\\
70000	1553.3\\
72000	1526.4\\
74000	1533.1\\
76000	1534\\
80000	1566\\
82000	1534.2\\
84000	1541.7\\
86000	1541.2\\
88000	1524.7\\
90000	1550.5\\
92000	1554\\
94000	1543.4\\
96000	1570.1\\
98000	1554.6\\
1e+05	1548.3\\
};
\addplot [color=mycolor78, forget plot]
  table[row sep=crcr]{%
0	4395.3\\
2000	2681.1\\
4000	2340.9\\
6000	2339.1\\
8000	2295.7\\
10000	2261.4\\
12000	2288.5\\
14000	2293.9\\
16000	2277.1\\
18000	2302.7\\
22000	2311\\
24000	2294.6\\
26000	2258.9\\
28000	2271.5\\
30000	2229.3\\
32000	2286.1\\
34000	2234.9\\
36000	2275.4\\
42000	2257\\
44000	2276.5\\
48000	2288.5\\
50000	2249.7\\
52000	2261.3\\
54000	2295.6\\
56000	2268.4\\
58000	2270.2\\
60000	2246.1\\
62000	2251.8\\
64000	2285.5\\
68000	2228.2\\
72000	2306.4\\
74000	2308.9\\
76000	2330.2\\
78000	2324.1\\
80000	2247.3\\
82000	2275.1\\
86000	2287.7\\
88000	2267.1\\
90000	2323.3\\
92000	2259.8\\
94000	2240.6\\
96000	2249.9\\
98000	2242.4\\
1e+05	2284.5\\
};
\addplot [color=mycolor79, forget plot]
  table[row sep=crcr]{%
0	4395.3\\
2000	2482.2\\
4000	2206.4\\
6000	2123.8\\
8000	2111.5\\
10000	2150.3\\
12000	2133.6\\
14000	2172.3\\
16000	2185.9\\
18000	2167.4\\
20000	2105.4\\
22000	2125.2\\
24000	2099.6\\
26000	2135.6\\
28000	2161.5\\
30000	2134.8\\
32000	2133.2\\
34000	2155.5\\
36000	2144.5\\
38000	2160.9\\
42000	2105.3\\
44000	2152.6\\
46000	2124.4\\
48000	2180\\
52000	2110.4\\
54000	2091.8\\
56000	2154.8\\
58000	2188.1\\
60000	2155.4\\
62000	2138.1\\
64000	2135.2\\
66000	2121.2\\
68000	2097.5\\
70000	2103.9\\
72000	2195.7\\
76000	2158.5\\
78000	2123.1\\
80000	2140.8\\
82000	2109.6\\
84000	2124.7\\
86000	2157.2\\
88000	2147\\
90000	2185.8\\
92000	2149.3\\
94000	2159.3\\
96000	2113.7\\
98000	2124.7\\
1e+05	2150.9\\
};
\addplot [color=mycolor79, forget plot]
  table[row sep=crcr]{%
0	4395.3\\
2000	2456.4\\
4000	2056.6\\
6000	1971.8\\
8000	1980.5\\
12000	2026.1\\
14000	1967\\
16000	2032.9\\
18000	1986.1\\
20000	1967\\
22000	1992.3\\
24000	1989.4\\
26000	2010.5\\
28000	2003.8\\
30000	1956.7\\
32000	1962.3\\
34000	2014.8\\
36000	2001.1\\
38000	1993.7\\
40000	1967.2\\
42000	1983.7\\
44000	2034.1\\
46000	1980.8\\
48000	2006.9\\
50000	2005.9\\
52000	1989\\
54000	1993\\
56000	2011.3\\
58000	1965.3\\
60000	2004.3\\
62000	1942.3\\
64000	2001.9\\
66000	1993.6\\
68000	2029.3\\
70000	2079.3\\
72000	2003\\
74000	1987.9\\
78000	1985\\
80000	1997.3\\
82000	1960.8\\
84000	1982.9\\
86000	2031.3\\
88000	2053.6\\
90000	1953.5\\
92000	2007.6\\
94000	2016.8\\
96000	1999.6\\
98000	1998\\
1e+05	2026.8\\
};
\addplot [color=mycolor80, forget plot]
  table[row sep=crcr]{%
0	4395.3\\
2000	1978.3\\
4000	1546.2\\
6000	1538.3\\
8000	1515.2\\
10000	1534.9\\
12000	1526.5\\
14000	1525.4\\
16000	1550.2\\
18000	1546.5\\
20000	1555\\
22000	1527\\
24000	1521.8\\
26000	1527.6\\
28000	1525.5\\
30000	1510\\
32000	1520.6\\
34000	1525.1\\
36000	1544\\
38000	1517.5\\
40000	1537.9\\
42000	1518.2\\
46000	1539.1\\
48000	1511.4\\
50000	1502.5\\
54000	1547.3\\
56000	1536.5\\
58000	1536.5\\
60000	1529.6\\
62000	1552.2\\
64000	1542.8\\
66000	1520.6\\
68000	1532.9\\
72000	1511.3\\
74000	1546.1\\
78000	1524.9\\
80000	1550.2\\
82000	1520.2\\
84000	1529.2\\
86000	1520.8\\
88000	1543.1\\
90000	1536.1\\
92000	1518.7\\
94000	1523.2\\
96000	1539.9\\
98000	1539.2\\
1e+05	1543.6\\
};
\addplot [color=mycolor81, forget plot]
  table[row sep=crcr]{%
0	4395.3\\
2000	1930.9\\
4000	1430.2\\
6000	1450.1\\
8000	1447.9\\
12000	1420.2\\
14000	1449.4\\
18000	1420.4\\
20000	1398.3\\
22000	1451.8\\
24000	1416.3\\
26000	1462.5\\
28000	1441.6\\
30000	1438.4\\
32000	1416.5\\
34000	1413.7\\
36000	1435.2\\
38000	1418.5\\
40000	1427\\
42000	1418.8\\
44000	1466.3\\
46000	1408.8\\
48000	1444\\
50000	1460.8\\
52000	1432.2\\
54000	1422.5\\
56000	1446.1\\
58000	1433.5\\
60000	1456.8\\
62000	1409.9\\
64000	1434.7\\
68000	1421.7\\
70000	1427.7\\
72000	1422.9\\
74000	1432.3\\
76000	1415.3\\
78000	1424.5\\
80000	1408.2\\
82000	1436.9\\
84000	1426.9\\
86000	1407.6\\
88000	1467.6\\
90000	1467.2\\
92000	1455.7\\
94000	1425.8\\
98000	1424.6\\
1e+05	1441.8\\
};
\addplot [color=mycolor82, forget plot]
  table[row sep=crcr]{%
0	4395.3\\
2000	2207.9\\
4000	1726.9\\
8000	1667.4\\
10000	1660.1\\
12000	1693.5\\
14000	1666.8\\
16000	1684.5\\
18000	1649.2\\
20000	1671.2\\
22000	1715.3\\
24000	1714.9\\
26000	1695.3\\
30000	1695.2\\
32000	1711.1\\
34000	1668.9\\
36000	1679.6\\
38000	1660.3\\
40000	1666.7\\
42000	1694.7\\
44000	1689.5\\
46000	1674.3\\
48000	1686.2\\
50000	1688.1\\
52000	1666\\
54000	1694.8\\
56000	1682.9\\
58000	1702.1\\
60000	1711.3\\
62000	1685.5\\
66000	1695.3\\
68000	1722.5\\
70000	1686\\
72000	1728.7\\
74000	1691\\
76000	1693.8\\
78000	1671.5\\
80000	1683.7\\
82000	1700.8\\
84000	1704.7\\
86000	1688.8\\
88000	1684.8\\
90000	1690.5\\
92000	1683.2\\
94000	1653.2\\
96000	1651.1\\
98000	1655.2\\
1e+05	1648.3\\
};
\addplot [color=mycolor83, forget plot]
  table[row sep=crcr]{%
0	4395.3\\
2000	2200.4\\
4000	1773.8\\
6000	1739.9\\
8000	1722.3\\
10000	1716\\
12000	1690.1\\
14000	1734.2\\
16000	1705.1\\
18000	1719.8\\
20000	1725\\
22000	1691.8\\
24000	1695.7\\
26000	1713.2\\
32000	1722.3\\
34000	1705.7\\
36000	1719.2\\
38000	1690.6\\
40000	1705.6\\
42000	1696\\
44000	1714\\
46000	1725.3\\
48000	1713.5\\
50000	1695.8\\
52000	1753.4\\
54000	1718.6\\
56000	1706.3\\
58000	1712\\
60000	1690.3\\
62000	1718.5\\
64000	1728.3\\
66000	1729.7\\
68000	1685\\
70000	1718.8\\
72000	1718.1\\
74000	1705.4\\
76000	1712.3\\
78000	1687.9\\
80000	1686.1\\
82000	1662.7\\
84000	1677\\
86000	1729.2\\
92000	1699.7\\
94000	1714.7\\
96000	1715.2\\
98000	1729.7\\
1e+05	1707.8\\
};
\addplot [color=mycolor83, forget plot]
  table[row sep=crcr]{%
0	4395.3\\
2000	1700.8\\
4000	1281\\
6000	1239.3\\
8000	1248.3\\
10000	1269.2\\
12000	1211\\
14000	1238.7\\
16000	1254.3\\
18000	1252.2\\
20000	1241.8\\
22000	1244\\
24000	1229.8\\
26000	1251\\
28000	1229.5\\
30000	1224.4\\
32000	1250.3\\
34000	1254.2\\
36000	1208.2\\
38000	1232.8\\
40000	1226\\
42000	1247.6\\
44000	1242.4\\
46000	1225.6\\
48000	1280.4\\
50000	1239.5\\
52000	1227.5\\
54000	1252.6\\
56000	1233.6\\
58000	1220.5\\
60000	1247.4\\
62000	1236.6\\
64000	1245.6\\
66000	1235.8\\
68000	1238.6\\
70000	1231.7\\
72000	1246.4\\
74000	1230.4\\
76000	1250.8\\
78000	1238.6\\
80000	1248.9\\
82000	1245.1\\
84000	1255.3\\
86000	1195.9\\
88000	1240.2\\
92000	1258.3\\
94000	1234.5\\
96000	1239.7\\
98000	1230\\
1e+05	1249\\
};
\addplot [color=mycolor84, forget plot]
  table[row sep=crcr]{%
0	4395.3\\
2000	1763.8\\
4000	1254.2\\
6000	1253.6\\
8000	1242.1\\
12000	1255.3\\
14000	1244.6\\
16000	1269.4\\
18000	1279.2\\
20000	1263.8\\
24000	1245.4\\
26000	1238.9\\
28000	1236.8\\
30000	1250.5\\
32000	1240.6\\
34000	1258.9\\
36000	1234.4\\
42000	1278\\
44000	1248.1\\
46000	1273\\
48000	1246\\
50000	1240.9\\
52000	1256.9\\
54000	1248.3\\
56000	1247.3\\
58000	1255\\
60000	1226.7\\
62000	1247.6\\
64000	1253.4\\
66000	1213\\
68000	1233.5\\
70000	1258.5\\
72000	1258.3\\
74000	1253.4\\
76000	1269.5\\
78000	1257.2\\
80000	1258.6\\
82000	1273.4\\
84000	1243.6\\
86000	1271\\
88000	1244\\
90000	1241.8\\
92000	1236.1\\
94000	1264.7\\
96000	1263.8\\
98000	1243.5\\
1e+05	1262.3\\
};
\addplot [color=mycolor85, forget plot]
  table[row sep=crcr]{%
0	4395.3\\
2000	2004.9\\
4000	1462.3\\
6000	1346.3\\
8000	1391.5\\
10000	1400.4\\
12000	1389.6\\
14000	1344.3\\
16000	1388.1\\
18000	1379.1\\
20000	1383.2\\
22000	1357.3\\
24000	1382.6\\
26000	1392.8\\
28000	1368.4\\
30000	1382.9\\
32000	1359.1\\
34000	1357.1\\
36000	1411.9\\
38000	1392.1\\
42000	1378.1\\
44000	1407.8\\
46000	1383.3\\
48000	1327.5\\
50000	1389.1\\
52000	1334.4\\
54000	1375\\
56000	1370.7\\
58000	1429.3\\
60000	1393.7\\
62000	1382.8\\
64000	1398.4\\
66000	1377.8\\
68000	1391.8\\
70000	1389.8\\
72000	1370.2\\
76000	1393.4\\
78000	1385\\
80000	1400.5\\
82000	1352.3\\
84000	1384.4\\
86000	1409.5\\
88000	1378.7\\
90000	1390.9\\
92000	1383.3\\
94000	1379.7\\
96000	1383.3\\
98000	1378.5\\
1e+05	1363.5\\
};
\addplot [color=mycolor86, forget plot]
  table[row sep=crcr]{%
0	4395.3\\
2000	1769.9\\
4000	1231.5\\
6000	1156\\
8000	1144.6\\
10000	1115.6\\
12000	1135.2\\
14000	1196.4\\
18000	1123.7\\
20000	1155.9\\
22000	1141.1\\
24000	1148.2\\
26000	1144.7\\
28000	1130.6\\
30000	1142.6\\
32000	1165.2\\
34000	1170.3\\
36000	1190.5\\
38000	1158.2\\
40000	1165.5\\
42000	1159.8\\
44000	1147.9\\
48000	1160.1\\
50000	1134.4\\
52000	1166.7\\
54000	1167.3\\
56000	1143.6\\
58000	1159.9\\
60000	1147.4\\
62000	1163.2\\
64000	1189.8\\
68000	1142.8\\
70000	1129.3\\
72000	1149.8\\
74000	1164.8\\
76000	1110.7\\
78000	1156.9\\
80000	1144.7\\
82000	1163.3\\
84000	1124.2\\
86000	1146.9\\
88000	1139.3\\
90000	1166.3\\
92000	1142.4\\
94000	1155.9\\
96000	1130.4\\
98000	1141.3\\
1e+05	1160.3\\
};
\addplot [color=mycolor87, forget plot]
  table[row sep=crcr]{%
0	4395.3\\
2000	1481.1\\
4000	957.75\\
6000	1007.3\\
8000	953.5\\
10000	938.25\\
12000	932.45\\
16000	1012.4\\
18000	995.45\\
20000	986.97\\
22000	1005.9\\
24000	932.48\\
26000	928.68\\
28000	950.54\\
30000	968.15\\
32000	994.5\\
34000	959.87\\
36000	985.04\\
38000	967.52\\
40000	1024.8\\
42000	1012.1\\
44000	949.11\\
46000	961.93\\
48000	986.83\\
50000	963.99\\
52000	994.03\\
56000	968.23\\
58000	966.97\\
60000	999.37\\
62000	968.47\\
64000	961.15\\
66000	1008.1\\
68000	981.57\\
70000	951.97\\
72000	951.37\\
74000	978.1\\
76000	990.02\\
78000	950.98\\
80000	979.68\\
82000	967.34\\
84000	899.47\\
86000	1001.9\\
88000	953.74\\
90000	972.14\\
92000	971.72\\
96000	984.6\\
98000	974.72\\
1e+05	973.29\\
};
\addplot [color=mycolor88, forget plot]
  table[row sep=crcr]{%
0	4395.3\\
2000	1809.2\\
4000	1218.2\\
6000	1203.8\\
8000	1174.8\\
10000	1169.3\\
12000	1171\\
14000	1133.9\\
16000	1119.5\\
18000	1167.9\\
20000	1184.6\\
22000	1171.2\\
24000	1147.2\\
26000	1167.7\\
28000	1142.7\\
30000	1158.2\\
32000	1154.5\\
34000	1184.2\\
36000	1174.4\\
38000	1154.6\\
40000	1172.9\\
42000	1140.4\\
44000	1177.8\\
46000	1165.7\\
48000	1180.2\\
50000	1158.1\\
52000	1164.6\\
54000	1159.5\\
56000	1143.8\\
58000	1186.4\\
60000	1206.2\\
62000	1179.8\\
64000	1160.6\\
68000	1201.3\\
70000	1214.4\\
72000	1165.3\\
74000	1171.9\\
76000	1172.9\\
78000	1214.9\\
80000	1176.5\\
82000	1158.4\\
84000	1181.8\\
88000	1154\\
90000	1158.7\\
92000	1149.2\\
94000	1165.1\\
96000	1172.7\\
98000	1141.3\\
1e+05	1157.2\\
};
\addplot [color=mycolor89, forget plot]
  table[row sep=crcr]{%
0	4395.3\\
2000	1548.4\\
4000	851.28\\
6000	823.94\\
8000	811.73\\
10000	842.52\\
12000	841.85\\
14000	834.8\\
16000	817.15\\
18000	830.7\\
22000	809.12\\
24000	806.17\\
26000	820.63\\
28000	811.16\\
30000	831.22\\
32000	842.72\\
34000	850.61\\
36000	828.17\\
40000	840.51\\
42000	831.83\\
44000	816.8\\
46000	812.3\\
48000	844.54\\
50000	823.1\\
52000	828.8\\
54000	848\\
56000	835.6\\
58000	828.76\\
60000	832.75\\
62000	852.85\\
64000	857.69\\
66000	839.83\\
68000	811.34\\
70000	835.88\\
72000	845.6\\
74000	837.37\\
76000	835.89\\
78000	825.16\\
80000	822.89\\
82000	842.26\\
84000	821.16\\
86000	851.77\\
88000	855.47\\
90000	827.69\\
92000	840.97\\
94000	837.37\\
96000	824.76\\
98000	805.27\\
1e+05	847.1\\
};
\addplot [color=mycolor90, forget plot]
  table[row sep=crcr]{%
0	4395.3\\
2000	1356.8\\
4000	810.3\\
6000	797.45\\
8000	796.55\\
10000	807.54\\
12000	815.7\\
14000	817.8\\
16000	833.19\\
18000	803.29\\
20000	814.89\\
22000	811.36\\
24000	787.25\\
26000	808.02\\
28000	782.83\\
30000	818.94\\
32000	809.94\\
34000	811.07\\
36000	796.41\\
38000	807.2\\
40000	791.28\\
42000	791.06\\
44000	812.2\\
46000	813.2\\
48000	830.96\\
50000	804.48\\
52000	811.43\\
54000	771.91\\
56000	797.56\\
58000	819.78\\
60000	823.12\\
62000	804.33\\
64000	808.87\\
66000	810.8\\
68000	839.22\\
70000	791.02\\
72000	807.98\\
74000	794.65\\
76000	806.46\\
78000	803.06\\
80000	790.49\\
82000	794.04\\
84000	807.02\\
86000	816.12\\
88000	817.38\\
90000	840.71\\
92000	828.62\\
94000	824.24\\
96000	794.46\\
98000	792.38\\
1e+05	786.62\\
};
\addplot [color=mycolor91, forget plot]
  table[row sep=crcr]{%
0	4395.3\\
2000	1543.7\\
4000	730.11\\
6000	689.45\\
8000	709.25\\
10000	688.66\\
12000	706.88\\
14000	682.55\\
16000	693.44\\
18000	689.8\\
20000	678.17\\
22000	681.19\\
24000	680.8\\
26000	691.46\\
28000	715.49\\
30000	706.81\\
32000	706.11\\
34000	696.5\\
36000	743.17\\
38000	706.91\\
40000	704.39\\
42000	682.41\\
44000	666.22\\
46000	729.05\\
48000	703.96\\
50000	704.82\\
52000	710.25\\
54000	684.9\\
56000	683.2\\
58000	708.05\\
60000	683.61\\
64000	695.27\\
66000	695.38\\
68000	684.62\\
70000	679.41\\
74000	714.74\\
76000	717.97\\
78000	688.37\\
80000	684.87\\
82000	688.03\\
84000	729.83\\
86000	712.44\\
88000	689.4\\
90000	692.9\\
92000	701.62\\
94000	687.08\\
96000	701.61\\
98000	694.67\\
1e+05	707.49\\
};
\addplot [color=mycolor92, forget plot]
  table[row sep=crcr]{%
0	4395.3\\
2000	1601.8\\
4000	893.99\\
6000	889.1\\
8000	893.13\\
10000	860.65\\
12000	847.45\\
14000	853.44\\
16000	862.27\\
20000	873.44\\
22000	873.32\\
24000	842.41\\
26000	851.78\\
28000	854.1\\
30000	876.91\\
32000	892.33\\
34000	895.85\\
36000	878.1\\
38000	844.06\\
40000	852.88\\
42000	888.94\\
46000	868.3\\
48000	850.36\\
50000	874.55\\
52000	873.45\\
54000	865.33\\
56000	853.49\\
58000	899.48\\
62000	849.64\\
64000	857.04\\
68000	848.79\\
70000	875.96\\
72000	885.99\\
74000	873.84\\
76000	853.14\\
78000	841.51\\
80000	859.94\\
82000	858.76\\
84000	881.43\\
86000	862.95\\
88000	849.32\\
90000	874.19\\
92000	870.4\\
94000	907.33\\
96000	931.53\\
98000	878.74\\
1e+05	874.24\\
};
\addplot [color=mycolor93, forget plot]
  table[row sep=crcr]{%
0	4395.3\\
2000	1572.7\\
4000	762.87\\
6000	740.17\\
8000	794.9\\
10000	765.42\\
12000	752.14\\
14000	746.92\\
16000	745.01\\
18000	735.46\\
20000	747.86\\
22000	725.06\\
24000	705.47\\
26000	749.01\\
28000	715.53\\
30000	743.84\\
32000	699.09\\
34000	730.73\\
36000	751.27\\
38000	746.49\\
40000	771.02\\
42000	741.85\\
44000	728.39\\
46000	701.09\\
48000	751.96\\
50000	715.05\\
52000	756.46\\
54000	737.44\\
56000	772.86\\
58000	710.64\\
60000	717.04\\
62000	738.38\\
64000	706.64\\
66000	749.37\\
70000	715.22\\
72000	722.35\\
74000	717.15\\
76000	719.82\\
78000	740.92\\
80000	704.37\\
82000	702.8\\
84000	723.85\\
86000	734.41\\
88000	719.2\\
90000	768.87\\
92000	713.29\\
94000	740.22\\
96000	696.05\\
98000	732.26\\
1e+05	737.61\\
};
\addplot [color=mycolor94, forget plot]
  table[row sep=crcr]{%
0	4395.3\\
2000	1590.8\\
4000	724.76\\
6000	690.3\\
8000	714.17\\
10000	703.38\\
12000	706.33\\
14000	690.1\\
16000	676.68\\
18000	697.91\\
20000	701.14\\
22000	698.24\\
24000	680.43\\
26000	683.12\\
28000	701.86\\
30000	703.44\\
32000	687.32\\
34000	695.89\\
36000	695.52\\
38000	698.52\\
40000	711.24\\
42000	695.64\\
44000	676.18\\
46000	701.48\\
48000	686.32\\
50000	684.62\\
54000	702.01\\
56000	706.79\\
60000	683.61\\
62000	694.75\\
64000	681.92\\
66000	676.21\\
68000	693.33\\
70000	681.97\\
72000	703.07\\
74000	713.91\\
76000	702.14\\
78000	707.34\\
80000	682.43\\
82000	675.6\\
84000	690.82\\
86000	689.64\\
88000	713.08\\
92000	698.35\\
94000	675.05\\
96000	659.15\\
98000	695.19\\
1e+05	692.18\\
};
\addplot [color=mycolor95, forget plot]
  table[row sep=crcr]{%
0	4395.3\\
2000	1256\\
4000	450.52\\
8000	415.8\\
10000	429.45\\
12000	407.2\\
14000	405.66\\
16000	396.78\\
18000	408.97\\
22000	401.98\\
24000	406.24\\
26000	405.97\\
28000	392.42\\
30000	400.72\\
32000	403.06\\
34000	408.52\\
36000	407.69\\
38000	430.91\\
40000	427.25\\
42000	411.65\\
44000	412.44\\
46000	418.6\\
48000	396.53\\
50000	396.12\\
52000	382.9\\
54000	391.06\\
56000	423.3\\
58000	409.54\\
60000	427.04\\
62000	402.72\\
64000	397.35\\
66000	400.83\\
68000	410.14\\
70000	405.77\\
72000	410.5\\
74000	412.45\\
76000	405.69\\
78000	418.33\\
80000	393.49\\
82000	409.58\\
84000	400.24\\
86000	425.31\\
88000	413.52\\
90000	416.12\\
94000	409.34\\
96000	394.06\\
98000	420.71\\
1e+05	418.83\\
};
\addplot [color=mycolor96, forget plot]
  table[row sep=crcr]{%
0	4395.3\\
2000	1442.1\\
4000	457.41\\
6000	443.51\\
8000	411.21\\
10000	397.04\\
12000	407.35\\
14000	414.21\\
16000	400.72\\
18000	421.67\\
20000	419.04\\
22000	406.94\\
24000	447.47\\
26000	441.07\\
28000	398.14\\
30000	416.3\\
32000	439.17\\
34000	449.57\\
36000	453.52\\
38000	432\\
40000	452.93\\
42000	427.6\\
44000	414.57\\
46000	421.62\\
48000	416.35\\
50000	469.45\\
52000	452.1\\
54000	455.34\\
56000	506.39\\
58000	447.54\\
60000	448.21\\
62000	460.82\\
64000	412.81\\
66000	431.65\\
68000	428.43\\
70000	456.67\\
72000	450.86\\
76000	410.46\\
78000	477.65\\
80000	430.25\\
82000	434.88\\
84000	440.79\\
86000	422.26\\
88000	418.54\\
90000	460.29\\
92000	454.61\\
94000	433.51\\
96000	446.56\\
98000	423.73\\
1e+05	411.91\\
};
\addplot [color=mycolor97, forget plot]
  table[row sep=crcr]{%
0	4395.3\\
2000	1080.2\\
4000	334.46\\
6000	339.5\\
8000	323.95\\
10000	328.51\\
12000	335.07\\
14000	307.21\\
16000	308.08\\
18000	333.21\\
20000	309.79\\
22000	316.19\\
24000	301.49\\
26000	284.06\\
28000	293.5\\
30000	310.94\\
32000	350.01\\
34000	339.8\\
36000	293.45\\
38000	307.92\\
40000	291.77\\
42000	286.22\\
44000	300.56\\
46000	325.7\\
48000	321.43\\
50000	323.18\\
52000	285.24\\
54000	279.95\\
58000	302.3\\
60000	299.12\\
62000	292.29\\
64000	278.89\\
66000	317.33\\
68000	322.6\\
70000	314.04\\
72000	323.11\\
74000	318.89\\
76000	308.42\\
78000	323.88\\
80000	321.24\\
82000	299.9\\
84000	300.1\\
86000	303.75\\
88000	325.19\\
90000	337.43\\
92000	337.94\\
94000	327.91\\
96000	313.43\\
98000	318.23\\
1e+05	316.45\\
};
\addplot [color=mycolor98]
  table[row sep=crcr]{%
0	4395.3\\
2000	1517\\
4000	494.38\\
6000	417.61\\
8000	408.83\\
10000	398.74\\
12000	403.76\\
14000	415.34\\
16000	421.94\\
18000	409.26\\
20000	420.12\\
22000	414.74\\
24000	432.97\\
26000	448.65\\
28000	426.05\\
30000	414.77\\
32000	397.59\\
34000	389.49\\
36000	420.27\\
38000	409.68\\
40000	428.89\\
42000	415.45\\
44000	432.72\\
46000	404.4\\
48000	403.08\\
50000	423.53\\
52000	424.72\\
54000	421.95\\
56000	412.52\\
58000	411.02\\
60000	407.78\\
62000	430.14\\
64000	414.8\\
66000	405.83\\
68000	406.59\\
70000	402.67\\
72000	423.37\\
74000	421.35\\
76000	411.46\\
78000	413\\
80000	406.17\\
82000	428.11\\
84000	419.59\\
86000	424.23\\
88000	422.48\\
90000	411.89\\
92000	430.87\\
96000	441.5\\
98000	422.33\\
1e+05	436.98\\
};
\addlegendentry{\ClasskIVLegend}

      % \input{../../../../.tex/20/KS/SizeEvolutionsfig7.tex}
      \legend{}

      \nextgroupplot[%
        % xmin=0,xmax=1e+05,
        % ymin=100,ymax=1.5886e+05,
      ] % This file was created by matlab2tikz.
%
\definecolor{mycolor1}{rgb}{0.00146,0.00047,0.01387}%
\definecolor{mycolor2}{rgb}{0.01166,0.00942,0.06346}%
\definecolor{mycolor3}{rgb}{0.02943,0.02150,0.11462}%
\definecolor{mycolor4}{rgb}{0.06134,0.03659,0.17764}%
\definecolor{mycolor5}{rgb}{0.08741,0.04456,0.22481}%
\definecolor{mycolor6}{rgb}{0.12291,0.04754,0.28162}%
\definecolor{mycolor7}{rgb}{0.14907,0.04547,0.31709}%
\definecolor{mycolor8}{rgb}{0.18343,0.04033,0.35497}%
\definecolor{mycolor9}{rgb}{0.21795,0.03662,0.38352}%
\definecolor{mycolor10}{rgb}{0.24497,0.03705,0.40001}%
\definecolor{mycolor11}{rgb}{0.27135,0.04092,0.41198}%
\definecolor{mycolor12}{rgb}{0.29076,0.04564,0.41864}%
\definecolor{mycolor13}{rgb}{0.31628,0.05349,0.42512}%
\definecolor{mycolor14}{rgb}{0.33522,0.06006,0.42852}%
\definecolor{mycolor15}{rgb}{0.36028,0.06925,0.43150}%
\definecolor{mycolor16}{rgb}{0.37900,0.07625,0.43272}%
\definecolor{mycolor17}{rgb}{0.39767,0.08326,0.43318}%
\definecolor{mycolor18}{rgb}{0.41633,0.09020,0.43294}%
\definecolor{mycolor19}{rgb}{0.42877,0.09479,0.43241}%
\definecolor{mycolor20}{rgb}{0.44743,0.10160,0.43108}%
\definecolor{mycolor21}{rgb}{0.46610,0.10832,0.42912}%
\definecolor{mycolor22}{rgb}{0.47856,0.11276,0.42747}%
\definecolor{mycolor23}{rgb}{0.49726,0.11938,0.42449}%
\definecolor{mycolor24}{rgb}{0.50973,0.12377,0.42216}%
\definecolor{mycolor25}{rgb}{0.52221,0.12815,0.41955}%
\definecolor{mycolor26}{rgb}{0.53468,0.13253,0.41667}%
\definecolor{mycolor27}{rgb}{0.55339,0.13913,0.41183}%
\definecolor{mycolor28}{rgb}{0.56585,0.14357,0.40826}%
\definecolor{mycolor29}{rgb}{0.57830,0.14804,0.40441}%
\definecolor{mycolor30}{rgb}{0.59073,0.15256,0.40029}%
\definecolor{mycolor31}{rgb}{0.60314,0.15715,0.39589}%
\definecolor{mycolor32}{rgb}{0.61551,0.16182,0.39122}%
\definecolor{mycolor33}{rgb}{0.62169,0.16418,0.38878}%
\definecolor{mycolor34}{rgb}{0.63400,0.16899,0.38370}%
\definecolor{mycolor35}{rgb}{0.64626,0.17391,0.37836}%
\definecolor{mycolor36}{rgb}{0.65846,0.17896,0.37275}%
\definecolor{mycolor37}{rgb}{0.66454,0.18154,0.36985}%
\definecolor{mycolor38}{rgb}{0.67664,0.18681,0.36385}%
\definecolor{mycolor39}{rgb}{0.68865,0.19224,0.35760}%
\definecolor{mycolor40}{rgb}{0.69463,0.19502,0.35439}%
\definecolor{mycolor41}{rgb}{0.70650,0.20073,0.34778}%
\definecolor{mycolor42}{rgb}{0.71240,0.20366,0.34438}%
\definecolor{mycolor43}{rgb}{0.72410,0.20967,0.33742}%
\definecolor{mycolor44}{rgb}{0.72991,0.21276,0.33386}%
\definecolor{mycolor45}{rgb}{0.74142,0.21911,0.32658}%
\definecolor{mycolor46}{rgb}{0.74713,0.22238,0.32286}%
\definecolor{mycolor47}{rgb}{0.75279,0.22571,0.31909}%
\definecolor{mycolor48}{rgb}{0.76401,0.23255,0.31140}%
\definecolor{mycolor49}{rgb}{0.76956,0.23608,0.30749}%
\definecolor{mycolor50}{rgb}{0.77506,0.23967,0.30353}%
\definecolor{mycolor51}{rgb}{0.79129,0.25086,0.29139}%
\definecolor{mycolor52}{rgb}{0.79661,0.25473,0.28726}%
\definecolor{mycolor53}{rgb}{0.80187,0.25867,0.28310}%
\definecolor{mycolor54}{rgb}{0.81224,0.26679,0.27466}%
\definecolor{mycolor55}{rgb}{0.81734,0.27095,0.27039}%
\definecolor{mycolor56}{rgb}{0.82737,0.27952,0.26175}%
\definecolor{mycolor57}{rgb}{0.83230,0.28391,0.25738}%
\definecolor{mycolor58}{rgb}{0.83717,0.28839,0.25299}%
\definecolor{mycolor59}{rgb}{0.84671,0.29756,0.24411}%
\definecolor{mycolor60}{rgb}{0.85138,0.30226,0.23964}%
\definecolor{mycolor61}{rgb}{0.85599,0.30704,0.23513}%
\definecolor{mycolor62}{rgb}{0.86053,0.31189,0.23061}%
\definecolor{mycolor63}{rgb}{0.87374,0.32691,0.21689}%
\definecolor{mycolor64}{rgb}{0.87800,0.33206,0.21227}%
\definecolor{mycolor65}{rgb}{0.89034,0.34796,0.19829}%
\definecolor{mycolor66}{rgb}{0.89819,0.35891,0.18886}%
\definecolor{mycolor67}{rgb}{0.90200,0.36449,0.18412}%
\definecolor{mycolor68}{rgb}{0.90573,0.37014,0.17935}%
\definecolor{mycolor69}{rgb}{0.90939,0.37586,0.17456}%
\definecolor{mycolor70}{rgb}{0.91297,0.38164,0.16975}%
\definecolor{mycolor71}{rgb}{0.91646,0.38748,0.16492}%
\definecolor{mycolor72}{rgb}{0.91988,0.39339,0.16007}%
\definecolor{mycolor73}{rgb}{0.92322,0.39936,0.15519}%
\definecolor{mycolor74}{rgb}{0.92647,0.40539,0.15029}%
\definecolor{mycolor75}{rgb}{0.92964,0.41148,0.14537}%
\definecolor{mycolor76}{rgb}{0.93274,0.41763,0.14042}%
\definecolor{mycolor77}{rgb}{0.93575,0.42383,0.13544}%
\definecolor{mycolor78}{rgb}{0.93868,0.43009,0.13044}%
\definecolor{mycolor79}{rgb}{0.94152,0.43640,0.12541}%
\definecolor{mycolor80}{rgb}{0.94429,0.44277,0.12035}%
\definecolor{mycolor81}{rgb}{0.94956,0.45566,0.11016}%
\definecolor{mycolor82}{rgb}{0.95208,0.46218,0.10503}%
\definecolor{mycolor83}{rgb}{0.96129,0.48872,0.08429}%
\definecolor{mycolor84}{rgb}{0.96540,0.50225,0.07386}%
\definecolor{mycolor85}{rgb}{0.96732,0.50908,0.06866}%
\definecolor{mycolor86}{rgb}{0.97259,0.52980,0.05332}%
\definecolor{mycolor87}{rgb}{0.97568,0.54380,0.04362}%
\definecolor{mycolor88}{rgb}{0.98082,0.57221,0.02851}%
\definecolor{mycolor89}{rgb}{0.98189,0.57939,0.02625}%
\definecolor{mycolor90}{rgb}{0.98378,0.59385,0.02377}%
\definecolor{mycolor91}{rgb}{0.98532,0.60842,0.02420}%
\definecolor{mycolor92}{rgb}{0.98696,0.63048,0.03091}%
\definecolor{mycolor93}{rgb}{0.98793,0.66025,0.05175}%
\definecolor{mycolor94}{rgb}{0.98771,0.68281,0.07249}%
\definecolor{mycolor95}{rgb}{0.97750,0.78226,0.18592}%
\definecolor{mycolor96}{rgb}{0.96624,0.83619,0.26153}%
\definecolor{mycolor97}{rgb}{0.96252,0.85148,0.28555}%
\definecolor{mycolor98}{rgb}{0.98836,0.99836,0.64492}%
%
\addplot [color=mycolor1]
  table[row sep=crcr]{%
0	4395.3\\
2000	1.1384e+05\\
4000	1.5122e+05\\
6000	1.5666e+05\\
8000	1.5499e+05\\
12000	1.5264e+05\\
14000	1.5495e+05\\
16000	1.525e+05\\
18000	1.5302e+05\\
20000	1.5419e+05\\
22000	1.4994e+05\\
26000	1.5656e+05\\
28000	1.5385e+05\\
30000	1.5687e+05\\
32000	1.5405e+05\\
34000	1.5071e+05\\
36000	1.5163e+05\\
38000	1.5189e+05\\
40000	1.5354e+05\\
42000	1.5632e+05\\
44000	1.5239e+05\\
46000	1.5481e+05\\
48000	1.4999e+05\\
50000	1.5439e+05\\
52000	1.5099e+05\\
54000	1.5423e+05\\
56000	1.5301e+05\\
58000	1.5238e+05\\
60000	1.5463e+05\\
62000	1.5313e+05\\
64000	1.5566e+05\\
66000	1.5748e+05\\
68000	1.5378e+05\\
70000	1.5566e+05\\
74000	1.5181e+05\\
76000	1.517e+05\\
78000	1.5437e+05\\
82000	1.5122e+05\\
84000	1.5364e+05\\
86000	1.5015e+05\\
88000	1.5886e+05\\
90000	1.5237e+05\\
92000	1.5206e+05\\
94000	1.4989e+05\\
96000	1.5415e+05\\
98000	1.5443e+05\\
1e+05	1.5258e+05\\
};
\addlegendentry{\ClasskILegend}

\addplot [color=mycolor2, forget plot]
  table[row sep=crcr]{%
0	4395.3\\
2000	23745\\
4000	25582\\
6000	26418\\
8000	25371\\
10000	25181\\
12000	26533\\
14000	26132\\
20000	26370\\
22000	25867\\
26000	26631\\
28000	24803\\
30000	26520\\
32000	26450\\
34000	25604\\
36000	25553\\
38000	26777\\
40000	26345\\
42000	25841\\
44000	25663\\
46000	25715\\
48000	25284\\
50000	25679\\
52000	25927\\
54000	25720\\
56000	26062\\
58000	25476\\
60000	25831\\
62000	26681\\
64000	25703\\
68000	25974\\
70000	26363\\
72000	25579\\
74000	25179\\
76000	26424\\
78000	26105\\
80000	25312\\
82000	26050\\
84000	26165\\
86000	26057\\
88000	26316\\
90000	25833\\
92000	26568\\
94000	25823\\
96000	25677\\
98000	26484\\
1e+05	26492\\
};
\addplot [color=mycolor3, forget plot]
  table[row sep=crcr]{%
0	4395.3\\
2000	35376\\
4000	40530\\
6000	40323\\
8000	39829\\
10000	39737\\
12000	40201\\
14000	39702\\
16000	41258\\
18000	40121\\
20000	40639\\
22000	40792\\
24000	39606\\
26000	41033\\
28000	40105\\
30000	38825\\
32000	39569\\
34000	41205\\
36000	41114\\
38000	39698\\
42000	39728\\
44000	39599\\
46000	41243\\
48000	40597\\
50000	39266\\
52000	38580\\
54000	39428\\
56000	41058\\
58000	40261\\
60000	39017\\
62000	40139\\
66000	39288\\
68000	40623\\
70000	39617\\
72000	41061\\
76000	38942\\
78000	40698\\
82000	39881\\
84000	38938\\
86000	39985\\
88000	39241\\
92000	40148\\
94000	40405\\
96000	40262\\
98000	39565\\
1e+05	40343\\
};
\addplot [color=mycolor4, forget plot]
  table[row sep=crcr]{%
0	4395.3\\
2000	13548\\
4000	14493\\
6000	13903\\
8000	14721\\
10000	14156\\
12000	14246\\
14000	14146\\
16000	14456\\
18000	14514\\
20000	14670\\
24000	14405\\
26000	14287\\
28000	14329\\
30000	13846\\
32000	14500\\
34000	14566\\
36000	14356\\
38000	14414\\
40000	14524\\
42000	14701\\
44000	14096\\
46000	14484\\
48000	14465\\
50000	14139\\
52000	14421\\
54000	14371\\
56000	14393\\
58000	14507\\
60000	14330\\
62000	14090\\
68000	14157\\
70000	14318\\
72000	14618\\
74000	14461\\
76000	14419\\
78000	14595\\
80000	14528\\
82000	14952\\
84000	14512\\
86000	13971\\
88000	14303\\
90000	14212\\
92000	14290\\
94000	14446\\
96000	14851\\
98000	14042\\
1e+05	14305\\
};
\addplot [color=mycolor5, forget plot]
  table[row sep=crcr]{%
0	4395.3\\
2000	52381\\
4000	64410\\
6000	63712\\
8000	62146\\
10000	65373\\
12000	63446\\
14000	63025\\
16000	63746\\
18000	65140\\
20000	65094\\
22000	66836\\
24000	63026\\
26000	62952\\
28000	64008\\
30000	62493\\
32000	65742\\
36000	62845\\
40000	64551\\
42000	64040\\
44000	64725\\
46000	62819\\
48000	64306\\
52000	66004\\
54000	63826\\
56000	64844\\
58000	67305\\
60000	64733\\
62000	63967\\
64000	65576\\
66000	62084\\
68000	63872\\
70000	62968\\
72000	63328\\
74000	61607\\
76000	62478\\
78000	64959\\
80000	65152\\
82000	63768\\
84000	65764\\
86000	63960\\
88000	62444\\
90000	63870\\
92000	64007\\
94000	65083\\
96000	63273\\
98000	63511\\
1e+05	65000\\
};
\addplot [color=mycolor6, forget plot]
  table[row sep=crcr]{%
0	4395.3\\
2000	45780\\
4000	54472\\
6000	55119\\
8000	55336\\
12000	56353\\
14000	57458\\
16000	56223\\
18000	56292\\
20000	54723\\
22000	57654\\
24000	57495\\
26000	56451\\
30000	57174\\
32000	56054\\
34000	56314\\
36000	57080\\
38000	57648\\
40000	54799\\
42000	55926\\
44000	56296\\
46000	56117\\
48000	57531\\
50000	55576\\
52000	56720\\
54000	56887\\
56000	57539\\
58000	57914\\
60000	56154\\
62000	56820\\
64000	55463\\
66000	55403\\
68000	53050\\
70000	55408\\
72000	55249\\
74000	56458\\
76000	56539\\
78000	55177\\
80000	57097\\
82000	57195\\
84000	55062\\
86000	55038\\
88000	56286\\
90000	58009\\
92000	57413\\
94000	55591\\
96000	57451\\
98000	55529\\
1e+05	56008\\
};
\addplot [color=mycolor7, forget plot]
  table[row sep=crcr]{%
0	4395.3\\
2000	25581\\
4000	28186\\
6000	28063\\
8000	28814\\
10000	29120\\
12000	29193\\
14000	28589\\
16000	28854\\
18000	27885\\
22000	28752\\
24000	28444\\
26000	27598\\
28000	28645\\
30000	27424\\
32000	27410\\
34000	27124\\
36000	26759\\
38000	28976\\
40000	28572\\
42000	27701\\
48000	28485\\
50000	29118\\
52000	28508\\
54000	28269\\
56000	27934\\
58000	28253\\
60000	28314\\
62000	27638\\
64000	27527\\
66000	29117\\
68000	28410\\
70000	27640\\
72000	28134\\
74000	28369\\
76000	28911\\
78000	28480\\
80000	27596\\
82000	28652\\
84000	27994\\
86000	27940\\
88000	28645\\
90000	27863\\
92000	28172\\
94000	29631\\
96000	28045\\
98000	28112\\
1e+05	29104\\
};
\addplot [color=mycolor8, forget plot]
  table[row sep=crcr]{%
0	4395.3\\
2000	19222\\
4000	21165\\
6000	20074\\
8000	20351\\
10000	20494\\
12000	21245\\
14000	19767\\
16000	19680\\
18000	20403\\
20000	20452\\
22000	20399\\
24000	20773\\
26000	20150\\
28000	20952\\
30000	21006\\
32000	20488\\
34000	21293\\
36000	20935\\
38000	20460\\
40000	20736\\
42000	20785\\
44000	20399\\
46000	20535\\
48000	20172\\
50000	20599\\
52000	20045\\
54000	21216\\
56000	19964\\
58000	20270\\
60000	19777\\
62000	20988\\
64000	20217\\
66000	20079\\
68000	20828\\
70000	20201\\
72000	20820\\
74000	21126\\
76000	20670\\
78000	20381\\
80000	19749\\
82000	20922\\
84000	21333\\
86000	20619\\
88000	20142\\
90000	21080\\
92000	20823\\
94000	20917\\
96000	19975\\
98000	20608\\
1e+05	20496\\
};
\addplot [color=mycolor9, forget plot]
  table[row sep=crcr]{%
0	4395.3\\
2000	13595\\
4000	14123\\
6000	14047\\
8000	14308\\
10000	14394\\
12000	13637\\
14000	13947\\
16000	14480\\
18000	14143\\
20000	13740\\
22000	14323\\
24000	14379\\
26000	13939\\
28000	14065\\
30000	14584\\
32000	13863\\
34000	14023\\
36000	13822\\
38000	14346\\
40000	13935\\
42000	14565\\
44000	14168\\
46000	14104\\
48000	14151\\
50000	13877\\
52000	14469\\
54000	13916\\
58000	13717\\
62000	14881\\
64000	14302\\
66000	14582\\
68000	13933\\
70000	13775\\
72000	14245\\
74000	14288\\
76000	14477\\
78000	14016\\
80000	13975\\
84000	14317\\
86000	14504\\
88000	14385\\
90000	14078\\
92000	14035\\
96000	14455\\
98000	13927\\
1e+05	14021\\
};
\addplot [color=mycolor10, forget plot]
  table[row sep=crcr]{%
0	4395.3\\
2000	19406\\
4000	22181\\
6000	21064\\
8000	22367\\
12000	21987\\
14000	22370\\
16000	21396\\
18000	22180\\
20000	22570\\
22000	21963\\
24000	21827\\
26000	21981\\
28000	21786\\
30000	22190\\
32000	22424\\
34000	21950\\
36000	21721\\
38000	21873\\
40000	22203\\
44000	22065\\
46000	21810\\
48000	22580\\
52000	21824\\
54000	21924\\
56000	22280\\
58000	22019\\
60000	22623\\
62000	21822\\
64000	22334\\
66000	22309\\
68000	21761\\
70000	21633\\
72000	22162\\
74000	22614\\
76000	22720\\
78000	21638\\
80000	21254\\
82000	21985\\
84000	21964\\
88000	22807\\
90000	22131\\
92000	22227\\
94000	21912\\
96000	22296\\
98000	22002\\
1e+05	21922\\
};
\addplot [color=mycolor11, forget plot]
  table[row sep=crcr]{%
0	4395.3\\
2000	8304.5\\
4000	8699.5\\
8000	8404.8\\
10000	8664.1\\
12000	8740.6\\
14000	8312.1\\
16000	8630.7\\
18000	8843.5\\
20000	8561.3\\
22000	8483.7\\
24000	8508.4\\
26000	8364.6\\
28000	8699.9\\
30000	8884.1\\
32000	8392.9\\
34000	8580.7\\
36000	8447.2\\
38000	8367.1\\
40000	8832.1\\
42000	8553.9\\
44000	8642.3\\
46000	8381.7\\
48000	8704.7\\
50000	8448.3\\
52000	8717.7\\
54000	8853.6\\
56000	8458.7\\
58000	8652.7\\
60000	8722.4\\
62000	8618.1\\
64000	8559.5\\
66000	8346.3\\
68000	8651.4\\
70000	8732\\
72000	8363.2\\
74000	8635.4\\
76000	8468.9\\
78000	8760.7\\
80000	8880.1\\
82000	8594.8\\
84000	8673.5\\
88000	8541.6\\
90000	8550.1\\
92000	8679.2\\
94000	8747.8\\
98000	8631.6\\
1e+05	8716.9\\
};
\addplot [color=mycolor12, forget plot]
  table[row sep=crcr]{%
0	4395.3\\
2000	13521\\
4000	15055\\
6000	14903\\
8000	15020\\
10000	14594\\
12000	14388\\
14000	14882\\
16000	14164\\
18000	14139\\
20000	14783\\
22000	14382\\
24000	14861\\
26000	15094\\
28000	14592\\
30000	15007\\
32000	15038\\
34000	15290\\
36000	14709\\
38000	14867\\
40000	14645\\
42000	14548\\
46000	15225\\
48000	15158\\
50000	14632\\
52000	15109\\
54000	15221\\
56000	14384\\
58000	14343\\
60000	14465\\
62000	14692\\
64000	14614\\
66000	14686\\
68000	14470\\
70000	14767\\
72000	14251\\
74000	14823\\
76000	14632\\
78000	14521\\
80000	14651\\
84000	14357\\
86000	14454\\
88000	15244\\
90000	14768\\
92000	14627\\
94000	14961\\
96000	14810\\
98000	14936\\
1e+05	14683\\
};
\addplot [color=mycolor13, forget plot]
  table[row sep=crcr]{%
0	4395.3\\
2000	10472\\
4000	11377\\
6000	10879\\
8000	11218\\
10000	10957\\
12000	10800\\
14000	11064\\
16000	10839\\
18000	10937\\
20000	11392\\
22000	11523\\
24000	11389\\
26000	11517\\
28000	11367\\
30000	11154\\
32000	10871\\
36000	11319\\
38000	10888\\
40000	10999\\
42000	10927\\
44000	11259\\
46000	11231\\
48000	11159\\
50000	10900\\
52000	11199\\
54000	11070\\
56000	11183\\
58000	10916\\
60000	11188\\
62000	11115\\
64000	11213\\
66000	11255\\
70000	11191\\
72000	11128\\
74000	10779\\
76000	11125\\
78000	10959\\
80000	11148\\
82000	10876\\
84000	10950\\
86000	11136\\
88000	11196\\
90000	10767\\
94000	11285\\
96000	11166\\
98000	10970\\
1e+05	11005\\
};
\addplot [color=mycolor14, forget plot]
  table[row sep=crcr]{%
0	4395.3\\
2000	20769\\
4000	24544\\
6000	24438\\
8000	24229\\
10000	24394\\
12000	24987\\
14000	24070\\
16000	24998\\
18000	23927\\
20000	24862\\
22000	23859\\
24000	24855\\
26000	24810\\
28000	25073\\
30000	24081\\
32000	24029\\
34000	25374\\
36000	24796\\
38000	24069\\
40000	24044\\
42000	24173\\
44000	24846\\
46000	24104\\
48000	23740\\
50000	23690\\
52000	25005\\
54000	24451\\
56000	24136\\
58000	24143\\
60000	24260\\
62000	25083\\
64000	24939\\
66000	23939\\
68000	24379\\
70000	24638\\
72000	24427\\
74000	24712\\
76000	24234\\
78000	24703\\
80000	24386\\
82000	23997\\
84000	24630\\
86000	24503\\
88000	24987\\
90000	24111\\
92000	23719\\
94000	23885\\
96000	24241\\
98000	25614\\
1e+05	23508\\
};
\addplot [color=mycolor15, forget plot]
  table[row sep=crcr]{%
0	4395.3\\
2000	15648\\
4000	17665\\
6000	16882\\
8000	16855\\
10000	16741\\
12000	16718\\
14000	17503\\
16000	16710\\
18000	16340\\
20000	17742\\
22000	16992\\
24000	16637\\
26000	17176\\
28000	17012\\
30000	16696\\
32000	16503\\
34000	16696\\
36000	17034\\
38000	17248\\
40000	17085\\
42000	17215\\
44000	16588\\
46000	17298\\
48000	16900\\
50000	16933\\
52000	17561\\
54000	17242\\
56000	17523\\
58000	17163\\
60000	17094\\
62000	16735\\
64000	17221\\
66000	17032\\
68000	17184\\
70000	16415\\
72000	17118\\
74000	16728\\
76000	16564\\
78000	16854\\
80000	17343\\
82000	17381\\
84000	17521\\
86000	16606\\
88000	16502\\
90000	17047\\
92000	16773\\
94000	17440\\
96000	16702\\
98000	17274\\
1e+05	16534\\
};
\addplot [color=mycolor16, forget plot]
  table[row sep=crcr]{%
0	4395.3\\
2000	6838.6\\
4000	7250\\
8000	6999.1\\
10000	7148\\
14000	6802.5\\
16000	7008.6\\
18000	6972.3\\
20000	6880.5\\
22000	7060.2\\
24000	6892.8\\
26000	6955.5\\
28000	6760.2\\
30000	7177.7\\
32000	6898.1\\
34000	6916.2\\
36000	7055.8\\
38000	7156.2\\
40000	7061.2\\
42000	7356.8\\
44000	7011.8\\
46000	7070.8\\
48000	6908.1\\
50000	7076.9\\
54000	6814\\
56000	7177.4\\
58000	7383.5\\
60000	7040.9\\
62000	6950\\
64000	7045.4\\
66000	6795\\
68000	6987.4\\
70000	7049.4\\
72000	7044.5\\
74000	6994.1\\
76000	7016.5\\
80000	6891.3\\
82000	7109.4\\
84000	6962.3\\
86000	7107.3\\
88000	7042.7\\
90000	7082.6\\
92000	6914\\
94000	6903.5\\
96000	7170.1\\
98000	7240.7\\
1e+05	7138.6\\
};
\addplot [color=mycolor17, forget plot]
  table[row sep=crcr]{%
0	4395.3\\
2000	15368\\
4000	16295\\
6000	15846\\
8000	15762\\
10000	16214\\
12000	15692\\
14000	16316\\
16000	15760\\
18000	15783\\
20000	15438\\
24000	16175\\
26000	15875\\
28000	15828\\
30000	16339\\
32000	15283\\
34000	15836\\
36000	16225\\
38000	16172\\
40000	15729\\
42000	15606\\
44000	16000\\
46000	15788\\
48000	15814\\
50000	15908\\
52000	16491\\
54000	16251\\
56000	16824\\
58000	16703\\
60000	15819\\
62000	15727\\
64000	16298\\
66000	16375\\
68000	15751\\
70000	16068\\
72000	16590\\
74000	16045\\
76000	16696\\
78000	16541\\
80000	16268\\
82000	15790\\
84000	15957\\
86000	16508\\
88000	15662\\
90000	15391\\
92000	16132\\
94000	15801\\
96000	16534\\
98000	15842\\
1e+05	15874\\
};
\addplot [color=mycolor18, forget plot]
  table[row sep=crcr]{%
0	4395.3\\
2000	21368\\
4000	24197\\
6000	24330\\
8000	24753\\
10000	25072\\
14000	24327\\
16000	25527\\
18000	23779\\
20000	25131\\
22000	23652\\
24000	24599\\
26000	25064\\
28000	23852\\
30000	24530\\
32000	24724\\
34000	23854\\
36000	25045\\
38000	24444\\
40000	24902\\
42000	24972\\
44000	24553\\
46000	24229\\
48000	25262\\
50000	25004\\
54000	23942\\
56000	23604\\
58000	23602\\
60000	23497\\
62000	23733\\
64000	24534\\
66000	24575\\
68000	24712\\
70000	24186\\
72000	23974\\
74000	25848\\
76000	24618\\
78000	24749\\
80000	24493\\
82000	24873\\
84000	24249\\
86000	24435\\
88000	25208\\
90000	24195\\
92000	24973\\
94000	25232\\
96000	24532\\
98000	24854\\
1e+05	23763\\
};
\addplot [color=mycolor19, forget plot]
  table[row sep=crcr]{%
0	4395.3\\
2000	14148\\
4000	14741\\
6000	14897\\
8000	14943\\
10000	15103\\
12000	15147\\
16000	14662\\
18000	15386\\
20000	14692\\
22000	15081\\
24000	15155\\
26000	14033\\
28000	14627\\
30000	14790\\
32000	14644\\
34000	14803\\
38000	15246\\
40000	15413\\
42000	15050\\
44000	14997\\
46000	14744\\
48000	14862\\
50000	14623\\
52000	14670\\
54000	14496\\
58000	14505\\
62000	14846\\
64000	14692\\
66000	14855\\
68000	14534\\
70000	15108\\
72000	15173\\
76000	15015\\
78000	15112\\
80000	15004\\
82000	14765\\
84000	15217\\
86000	14902\\
88000	14703\\
90000	14665\\
92000	15272\\
94000	14568\\
96000	15207\\
1e+05	14603\\
};
\addplot [color=mycolor20, forget plot]
  table[row sep=crcr]{%
0	4395.3\\
2000	17593\\
4000	19310\\
6000	19526\\
8000	18578\\
12000	19020\\
14000	18603\\
16000	19718\\
18000	19662\\
20000	19227\\
22000	19122\\
24000	19723\\
26000	19823\\
28000	19117\\
30000	19931\\
32000	19070\\
34000	19526\\
36000	18557\\
38000	19464\\
40000	19685\\
42000	19703\\
44000	20083\\
46000	20010\\
48000	20346\\
50000	19299\\
52000	19666\\
54000	19084\\
56000	19654\\
58000	19190\\
60000	19131\\
62000	19899\\
64000	19130\\
66000	19177\\
68000	19289\\
70000	19587\\
72000	20299\\
74000	19196\\
76000	19392\\
78000	19429\\
80000	19863\\
82000	19682\\
84000	19691\\
86000	19766\\
88000	18898\\
90000	20426\\
92000	19828\\
94000	18970\\
96000	19866\\
98000	19075\\
1e+05	19924\\
};
\addplot [color=mycolor21, forget plot]
  table[row sep=crcr]{%
0	4395.3\\
2000	9040.8\\
4000	9239.2\\
6000	9780.6\\
8000	9808.3\\
10000	9691.8\\
12000	9901.6\\
14000	9844.7\\
16000	9678\\
18000	9763.4\\
20000	9429.7\\
22000	9578.6\\
24000	9582.6\\
26000	9763.1\\
28000	9714.3\\
30000	9755.8\\
32000	9564.1\\
34000	9417.7\\
36000	9756.4\\
38000	9793.2\\
40000	9528.7\\
42000	9828.1\\
44000	9501.5\\
46000	9700.9\\
48000	9738.1\\
50000	9713\\
52000	9539.8\\
54000	9664.1\\
56000	9502.3\\
60000	9709.9\\
62000	9889\\
64000	9848\\
66000	9760.4\\
68000	10069\\
70000	9607.6\\
72000	9937.2\\
74000	9869.6\\
76000	10019\\
78000	9290.3\\
80000	9754.1\\
82000	9619.7\\
84000	9575.4\\
86000	9688.8\\
88000	9664.4\\
90000	9424.8\\
92000	9633.3\\
94000	9755.6\\
96000	9817\\
98000	9736.1\\
1e+05	9725.3\\
};
\addplot [color=mycolor22, forget plot]
  table[row sep=crcr]{%
0	4395.3\\
2000	21893\\
4000	24904\\
6000	25903\\
8000	25616\\
12000	26173\\
14000	25735\\
16000	25754\\
18000	25272\\
20000	24434\\
22000	24570\\
24000	25554\\
26000	25155\\
30000	25311\\
32000	26318\\
34000	26311\\
36000	25715\\
38000	24985\\
40000	25166\\
42000	26087\\
44000	24933\\
46000	25110\\
48000	26138\\
50000	25328\\
52000	25682\\
54000	25223\\
56000	25087\\
58000	25580\\
60000	25584\\
62000	25725\\
64000	25524\\
66000	25239\\
68000	25870\\
70000	25005\\
72000	24812\\
74000	26237\\
76000	25704\\
78000	24930\\
80000	25759\\
82000	26735\\
84000	25236\\
86000	26151\\
88000	25141\\
90000	25398\\
92000	25457\\
94000	25321\\
96000	24990\\
98000	25182\\
1e+05	25910\\
};
\addplot [color=mycolor23, forget plot]
  table[row sep=crcr]{%
0	4395.3\\
2000	11469\\
4000	12505\\
6000	12330\\
8000	12504\\
10000	12285\\
12000	12142\\
14000	12088\\
16000	12109\\
18000	12354\\
20000	12479\\
22000	12053\\
24000	12409\\
26000	12341\\
28000	12415\\
30000	12167\\
32000	12047\\
34000	12446\\
38000	12289\\
40000	12384\\
42000	12254\\
44000	12323\\
46000	12232\\
52000	12458\\
54000	12432\\
56000	12249\\
58000	12198\\
60000	12414\\
62000	12236\\
64000	12348\\
66000	12258\\
70000	12304\\
72000	12356\\
74000	12276\\
76000	12329\\
78000	12527\\
80000	12092\\
82000	12270\\
84000	12254\\
86000	12394\\
88000	12311\\
90000	12115\\
94000	12352\\
96000	12330\\
98000	12241\\
1e+05	12429\\
};
\addplot [color=mycolor24, forget plot]
  table[row sep=crcr]{%
0	4395.3\\
2000	4915.4\\
4000	4851.1\\
6000	4644.5\\
8000	4879.5\\
10000	4799\\
14000	4779\\
16000	4966.4\\
18000	4781.6\\
20000	4821.9\\
22000	4773.8\\
24000	4624.7\\
26000	4865\\
30000	4588.8\\
32000	4767.9\\
34000	4766\\
36000	4739.5\\
38000	4735.3\\
40000	4790.2\\
42000	4610.7\\
44000	4774.8\\
46000	4771.5\\
48000	4704\\
50000	4611.8\\
52000	4805.8\\
54000	4661.7\\
56000	4666.1\\
60000	4777\\
62000	4617.1\\
64000	4875.1\\
66000	4701\\
68000	4786.2\\
70000	4746.2\\
72000	4844.2\\
74000	4740.4\\
76000	4968.2\\
78000	4726.9\\
80000	4768.3\\
82000	5021.9\\
84000	4833.8\\
86000	4693.8\\
88000	4860.7\\
92000	4713.1\\
94000	4817.7\\
96000	4730.7\\
98000	4772.1\\
1e+05	4795\\
};
\addplot [color=mycolor25, forget plot]
  table[row sep=crcr]{%
0	4395.3\\
2000	12120\\
4000	12321\\
6000	11708\\
8000	12420\\
10000	12694\\
12000	12789\\
14000	12932\\
16000	12430\\
18000	12580\\
20000	12538\\
24000	12235\\
26000	12387\\
28000	12343\\
30000	12011\\
32000	12710\\
34000	12499\\
36000	12748\\
38000	12076\\
40000	12633\\
42000	12408\\
44000	12275\\
46000	12522\\
48000	12845\\
50000	12168\\
52000	12079\\
54000	12400\\
56000	12214\\
58000	12245\\
60000	12426\\
62000	12032\\
64000	12397\\
66000	12030\\
68000	12486\\
70000	12789\\
72000	12049\\
74000	12263\\
76000	12795\\
78000	12639\\
80000	12737\\
82000	12089\\
84000	12272\\
86000	12169\\
88000	12479\\
90000	12576\\
92000	12016\\
94000	12120\\
96000	12469\\
98000	12182\\
1e+05	12376\\
};
\addplot [color=mycolor26, forget plot]
  table[row sep=crcr]{%
0	4395.3\\
2000	8268.6\\
4000	8142.6\\
6000	8224\\
8000	8457\\
10000	8158.3\\
12000	8421.5\\
14000	8278\\
16000	8584.7\\
18000	8632.1\\
22000	8463.2\\
24000	8157.6\\
26000	8266.1\\
28000	8205.3\\
30000	8063.8\\
32000	7973.1\\
34000	8558.8\\
36000	8075.8\\
38000	8436.1\\
40000	7950.6\\
42000	8034.5\\
44000	8014.4\\
46000	8257.2\\
48000	7789.1\\
50000	8099.7\\
52000	8367\\
54000	8361.6\\
56000	8192\\
58000	8206.9\\
60000	8327.6\\
62000	8144.4\\
64000	8342.4\\
66000	8228.8\\
68000	8414.1\\
70000	8446.3\\
72000	8047.6\\
76000	8204.3\\
78000	8476.8\\
80000	8070.5\\
82000	8199.1\\
84000	8147.1\\
86000	8310.6\\
88000	8333\\
90000	8218.5\\
92000	8429.1\\
94000	8488.7\\
96000	7960.8\\
98000	7971.6\\
1e+05	8192.3\\
};
\addplot [color=mycolor27, forget plot]
  table[row sep=crcr]{%
0	4395.3\\
2000	9661.6\\
4000	9646.5\\
6000	9716.8\\
8000	9686.5\\
10000	9706.5\\
12000	9840.5\\
14000	9590.7\\
16000	9887.4\\
18000	9594.3\\
20000	9634\\
22000	10063\\
24000	10120\\
26000	9793.7\\
28000	9809.7\\
30000	9622.4\\
32000	9641.9\\
34000	9872.3\\
36000	9475.4\\
38000	9584.2\\
40000	9908.9\\
42000	9653\\
44000	9747.6\\
46000	10009\\
48000	9603\\
50000	9916.4\\
52000	9891.2\\
54000	9661.3\\
56000	9729.1\\
58000	9955.7\\
60000	9909.2\\
62000	9658.1\\
64000	9492.1\\
66000	9850.3\\
68000	9743.6\\
70000	9702.7\\
72000	9915.9\\
74000	9721.3\\
76000	10096\\
78000	9588.5\\
80000	9704.2\\
82000	9948.9\\
84000	9786.6\\
86000	9791.1\\
88000	9523.3\\
90000	9618.3\\
94000	9937.2\\
96000	9643.6\\
1e+05	9582\\
};
\addplot [color=mycolor28, forget plot]
  table[row sep=crcr]{%
0	4395.3\\
2000	5403.8\\
4000	5392.6\\
6000	5563.5\\
8000	5627.4\\
10000	5410.4\\
12000	5504.6\\
14000	5445.3\\
16000	5405.9\\
18000	5495.1\\
20000	5510.8\\
22000	5442.3\\
24000	5448.7\\
26000	5365\\
28000	5598\\
30000	5536.2\\
32000	5555.3\\
34000	5467.7\\
36000	5610.3\\
38000	5408.4\\
40000	5368.5\\
42000	5731.9\\
44000	5290.1\\
48000	5540.3\\
50000	5368.2\\
52000	5250\\
54000	5362.7\\
56000	5351.8\\
58000	5439.7\\
60000	5419.2\\
62000	5286.3\\
64000	5503.1\\
66000	5375\\
68000	5388.9\\
70000	5520.3\\
72000	5448.4\\
74000	5323.6\\
76000	5337.6\\
78000	5534.9\\
82000	5486\\
84000	5349.8\\
86000	5483.2\\
88000	5522.8\\
90000	5370.2\\
92000	5205.4\\
94000	5427.1\\
96000	5361.9\\
98000	5455.3\\
1e+05	5445.2\\
};
\addplot [color=mycolor29, forget plot]
  table[row sep=crcr]{%
0	4395.3\\
2000	11197\\
4000	11924\\
6000	11691\\
8000	11680\\
10000	11741\\
12000	11922\\
14000	10917\\
16000	11920\\
18000	11439\\
20000	11507\\
22000	11665\\
24000	11493\\
26000	11363\\
28000	11755\\
30000	11555\\
32000	11786\\
34000	11761\\
36000	12144\\
38000	11601\\
40000	11250\\
42000	11595\\
44000	11482\\
46000	11998\\
48000	11613\\
50000	11165\\
52000	11705\\
54000	11373\\
56000	11470\\
58000	11725\\
60000	10866\\
62000	11353\\
64000	11473\\
66000	11531\\
68000	11232\\
70000	11688\\
72000	11601\\
74000	11657\\
76000	11537\\
78000	11599\\
80000	11588\\
82000	11540\\
84000	11662\\
86000	12090\\
88000	11326\\
90000	11592\\
98000	11570\\
1e+05	11578\\
};
\addplot [color=mycolor30, forget plot]
  table[row sep=crcr]{%
0	4395.3\\
2000	7652.7\\
4000	7982.4\\
6000	7974\\
8000	7889\\
10000	7953.4\\
12000	7798.7\\
14000	7887.4\\
16000	7843.9\\
18000	7966.4\\
20000	7760\\
22000	7798.6\\
24000	8004.4\\
26000	8027.3\\
28000	7917.2\\
30000	7710.3\\
32000	8032.1\\
34000	7779.5\\
36000	7853.2\\
38000	7779.9\\
40000	7758.6\\
42000	7909.1\\
44000	8003.1\\
46000	7885\\
48000	7988.3\\
50000	7842.6\\
52000	7733\\
54000	7824.8\\
56000	8030.8\\
58000	8000\\
62000	8015.6\\
64000	8096.9\\
66000	7820.7\\
68000	7883.1\\
70000	7984.4\\
72000	7911.3\\
74000	7886.7\\
76000	7902.7\\
78000	7965.1\\
80000	7883\\
84000	7950.7\\
86000	8137.3\\
88000	7852.1\\
90000	7972.2\\
92000	7870.2\\
94000	8017.2\\
96000	7789.1\\
98000	7888.1\\
1e+05	7872.7\\
};
\addplot [color=mycolor31, forget plot]
  table[row sep=crcr]{%
0	4395.3\\
2000	6429.8\\
4000	6483.1\\
6000	6238.8\\
8000	6342\\
12000	6615.5\\
14000	6370.8\\
16000	6408.1\\
18000	6571.3\\
20000	6184.2\\
22000	6327.5\\
24000	6514.3\\
26000	6851.4\\
28000	6446.2\\
30000	6480.3\\
32000	6192.3\\
34000	6154.7\\
36000	6668.9\\
38000	6349.8\\
40000	6388.6\\
42000	6450.7\\
44000	6364.1\\
46000	6451.8\\
48000	6489\\
52000	6485.6\\
54000	6307.1\\
56000	6569.3\\
58000	6442.3\\
60000	6376.9\\
62000	6505.9\\
64000	6421.2\\
66000	6428.7\\
68000	6669.8\\
70000	6334.1\\
72000	6497.7\\
74000	6750.6\\
76000	6408.9\\
78000	6593.1\\
80000	6698.1\\
82000	6442.4\\
84000	6514.1\\
86000	6369.3\\
88000	6517\\
90000	6310.5\\
92000	6618.7\\
94000	6739.9\\
96000	6344.1\\
1e+05	6308\\
};
\addplot [color=mycolor32, forget plot]
  table[row sep=crcr]{%
0	4395.3\\
2000	9862.5\\
4000	10862\\
6000	10537\\
8000	10457\\
10000	10284\\
14000	10650\\
16000	10355\\
18000	10706\\
20000	10436\\
22000	10749\\
24000	10675\\
26000	10463\\
28000	10655\\
32000	10619\\
34000	11043\\
36000	10261\\
38000	10530\\
40000	10301\\
42000	10618\\
44000	10139\\
46000	10651\\
48000	10384\\
50000	10614\\
52000	10509\\
54000	10119\\
56000	10363\\
58000	10772\\
60000	10394\\
62000	10591\\
64000	10127\\
66000	10271\\
68000	10732\\
70000	10607\\
74000	10786\\
76000	10339\\
78000	10474\\
80000	10407\\
82000	10156\\
84000	10439\\
86000	10146\\
88000	9895.9\\
90000	10487\\
92000	10632\\
94000	10448\\
96000	10525\\
98000	10190\\
1e+05	10834\\
};
\addplot [color=mycolor33, forget plot]
  table[row sep=crcr]{%
0	4395.3\\
2000	7975.8\\
4000	8193.7\\
6000	8298.6\\
8000	8174.8\\
10000	8155.7\\
12000	8348.8\\
14000	8319\\
16000	8510.4\\
18000	8016.5\\
20000	8205.3\\
22000	8273.5\\
24000	8299.2\\
26000	8155.8\\
28000	8128.3\\
32000	8353.9\\
34000	8388.2\\
36000	8324.6\\
38000	8426.6\\
40000	8234.9\\
42000	8252\\
44000	8204.7\\
46000	8292.5\\
48000	8497.2\\
50000	8343.5\\
52000	8348.7\\
54000	8062.8\\
56000	8094\\
58000	8196.7\\
60000	8339.1\\
62000	8443.9\\
64000	8310.2\\
66000	8034.8\\
68000	7938.3\\
70000	8547.4\\
72000	8408\\
74000	8459.8\\
76000	8338.9\\
78000	8358\\
80000	8093.1\\
82000	8285.4\\
84000	7915.2\\
86000	8021.6\\
88000	8167.1\\
90000	8503.2\\
92000	8527.4\\
94000	8226.7\\
96000	8095.9\\
98000	8078\\
1e+05	8138.1\\
};
\addplot [color=mycolor34, forget plot]
  table[row sep=crcr]{%
0	4395.3\\
2000	4147.7\\
4000	3973.2\\
6000	4070.2\\
8000	4065.3\\
10000	4150.7\\
12000	3979.7\\
14000	4033\\
16000	4028\\
18000	3852.7\\
20000	3988.4\\
24000	4005.4\\
26000	4033.8\\
28000	3993.2\\
30000	3912.1\\
32000	4194.6\\
34000	3986.3\\
36000	3931.5\\
38000	4033.2\\
40000	4044\\
42000	3937.8\\
46000	4007.1\\
48000	4130.1\\
50000	3923\\
52000	3925.7\\
54000	3893.7\\
56000	4050.9\\
58000	3999.6\\
60000	3994.7\\
62000	3910.9\\
64000	3905.6\\
66000	3958.5\\
68000	4027.3\\
70000	3978.3\\
74000	4108.1\\
76000	4061.4\\
78000	4055.8\\
80000	3967\\
82000	3918\\
84000	4117.2\\
86000	3947.8\\
88000	4100.3\\
90000	4054\\
92000	4077.5\\
96000	3944\\
98000	4064.1\\
1e+05	3908.9\\
};
\addplot [color=mycolor35]
  table[row sep=crcr]{%
0	4395.3\\
2000	9545\\
4000	10356\\
6000	10260\\
8000	10010\\
10000	10339\\
12000	10356\\
14000	10249\\
16000	9987.3\\
18000	10364\\
20000	10301\\
22000	10363\\
24000	10015\\
26000	9998.4\\
28000	10307\\
30000	10274\\
32000	10277\\
34000	10206\\
36000	10221\\
38000	10398\\
40000	10141\\
42000	10192\\
48000	10083\\
52000	10293\\
54000	10121\\
56000	10115\\
58000	10190\\
60000	10345\\
62000	10116\\
64000	9988.2\\
66000	10159\\
68000	10161\\
70000	10095\\
72000	10525\\
74000	10193\\
76000	10247\\
80000	10296\\
82000	10075\\
84000	10213\\
86000	10228\\
88000	9781.4\\
90000	10015\\
92000	10027\\
94000	10336\\
96000	9938.5\\
98000	10117\\
1e+05	10171\\
};
\addlegendentry{\ClasskIILegend}

\addplot [color=mycolor36, forget plot]
  table[row sep=crcr]{%
0	4395.3\\
2000	4423\\
4000	4241.9\\
6000	4190.3\\
8000	4316.2\\
10000	4343\\
12000	4116.4\\
14000	4398.4\\
16000	4180.1\\
18000	4183.3\\
20000	4165\\
22000	4268.7\\
24000	4156.4\\
26000	4399.6\\
28000	4327.7\\
30000	4118\\
32000	4350.4\\
34000	4261.8\\
36000	4262.3\\
38000	4298.9\\
40000	4384\\
42000	4148.6\\
44000	4276.3\\
46000	4146.9\\
48000	4285.4\\
50000	4164.7\\
52000	4188.9\\
54000	4199.6\\
56000	4090.5\\
58000	4406.1\\
60000	4274.8\\
62000	4267.5\\
64000	4193.8\\
66000	4213.5\\
68000	4306\\
70000	4253.7\\
72000	4378.8\\
74000	4372.8\\
76000	4320\\
78000	4115.5\\
80000	4206.2\\
82000	4219.2\\
84000	4364.9\\
86000	4348.9\\
88000	4269\\
90000	4430.3\\
92000	4235.9\\
94000	4316\\
96000	4214.8\\
98000	4182.8\\
1e+05	4321.7\\
};
\addplot [color=mycolor37, forget plot]
  table[row sep=crcr]{%
0	4395.3\\
2000	3831\\
4000	3620.6\\
6000	3693.2\\
8000	3521.5\\
10000	3560.6\\
12000	3648.6\\
14000	3544.9\\
18000	3551.2\\
22000	3619.6\\
24000	3684.8\\
26000	3616.5\\
28000	3576.4\\
30000	3609.4\\
32000	3675.9\\
34000	3648.4\\
36000	3581.9\\
38000	3654.8\\
40000	3571.4\\
42000	3563.1\\
44000	3606.9\\
46000	3620.4\\
48000	3601.7\\
50000	3560.4\\
52000	3596.1\\
54000	3648.7\\
56000	3502.6\\
58000	3626.5\\
60000	3713.7\\
62000	3773.9\\
64000	3572.3\\
66000	3673.9\\
68000	3700.3\\
70000	3709.6\\
72000	3613.9\\
74000	3570.8\\
76000	3648.2\\
78000	3701.9\\
80000	3629.8\\
82000	3638.8\\
84000	3604.6\\
86000	3629.4\\
90000	3607.3\\
92000	3599.7\\
96000	3650.3\\
98000	3588.1\\
1e+05	3625.9\\
};
\addplot [color=mycolor38, forget plot]
  table[row sep=crcr]{%
0	4395.3\\
2000	3046\\
6000	2924.8\\
8000	2992.6\\
10000	3071.9\\
12000	2974.9\\
14000	2895\\
16000	2993.8\\
18000	2863.8\\
20000	3031.9\\
22000	2871.6\\
24000	3033.3\\
26000	2990.3\\
30000	2943.8\\
34000	2986.3\\
38000	3116.7\\
40000	3073.1\\
42000	2951.2\\
44000	3012.6\\
46000	2897.1\\
48000	2935.7\\
50000	2912.3\\
52000	3028.6\\
54000	2960\\
56000	2940.2\\
60000	3024.2\\
62000	2925\\
64000	2924.6\\
66000	2939.4\\
70000	2930.8\\
72000	3062.8\\
74000	3014.9\\
76000	2982.5\\
78000	2961.8\\
80000	3008.3\\
82000	3075.6\\
84000	3043.7\\
86000	2906\\
88000	3023.3\\
90000	3060.4\\
92000	3067.9\\
96000	2921.9\\
98000	3006.4\\
1e+05	2968.3\\
};
\addplot [color=mycolor39, forget plot]
  table[row sep=crcr]{%
0	4395.3\\
2000	10709\\
4000	11864\\
6000	12038\\
8000	11846\\
10000	11471\\
12000	11515\\
14000	11457\\
16000	11908\\
18000	11554\\
20000	11509\\
22000	12044\\
24000	11863\\
26000	11776\\
28000	11797\\
30000	11780\\
32000	11135\\
34000	11779\\
36000	11761\\
38000	11604\\
40000	11594\\
42000	11773\\
44000	12017\\
46000	11413\\
48000	12082\\
50000	11829\\
52000	12066\\
54000	12261\\
56000	11795\\
58000	11314\\
62000	11591\\
64000	11806\\
66000	11858\\
70000	11446\\
72000	11713\\
74000	11582\\
76000	11529\\
78000	11739\\
80000	12033\\
82000	11610\\
84000	11638\\
86000	11503\\
90000	12175\\
92000	11481\\
96000	12026\\
1e+05	11573\\
};
\addplot [color=mycolor40, forget plot]
  table[row sep=crcr]{%
0	4395.3\\
2000	5582.8\\
4000	5785\\
6000	5782.6\\
8000	5704.2\\
10000	5672.1\\
12000	5789.2\\
14000	5728.1\\
16000	5809.9\\
18000	5789\\
20000	5627.2\\
22000	5696.4\\
24000	5708.9\\
26000	5753.7\\
28000	5648.5\\
30000	5581.9\\
32000	5692.1\\
34000	5661.4\\
36000	5729.2\\
38000	5669.2\\
40000	5648.2\\
44000	5849.5\\
46000	5744.4\\
48000	5827.7\\
50000	5644.7\\
52000	5624.3\\
54000	5827.3\\
56000	5818.2\\
58000	5750.4\\
60000	5831.2\\
62000	5697.9\\
64000	5673.9\\
66000	5700.9\\
68000	5806.2\\
70000	5511.3\\
72000	5610\\
74000	5657.2\\
76000	5818.9\\
78000	5681.4\\
80000	5725.8\\
82000	5735.7\\
84000	5601.1\\
86000	5690.1\\
92000	5793.8\\
94000	5715\\
96000	5701.1\\
98000	5811.9\\
1e+05	5858\\
};
\addplot [color=mycolor41, forget plot]
  table[row sep=crcr]{%
0	4395.3\\
2000	5370.2\\
4000	5399.8\\
6000	5494.5\\
8000	5371.2\\
10000	5501.1\\
12000	5262.8\\
14000	5252.7\\
16000	5502.8\\
18000	5437.8\\
20000	5098.4\\
22000	5210.2\\
24000	5360.5\\
26000	5385.5\\
28000	5593\\
30000	5458.4\\
32000	5535.5\\
34000	5462\\
36000	5527.2\\
38000	5333.7\\
40000	5381.4\\
42000	5308.6\\
44000	5354.5\\
46000	5355.4\\
48000	5445.9\\
50000	5339.4\\
52000	5418.3\\
54000	5306.8\\
56000	5319.2\\
58000	5421.7\\
60000	5269.7\\
62000	5351.4\\
64000	5228.6\\
66000	5318\\
70000	5449.4\\
72000	5414.3\\
74000	5421.4\\
76000	5350.4\\
78000	5330.1\\
80000	5528.4\\
84000	5469.9\\
86000	5484.6\\
94000	5421.3\\
96000	5391.9\\
98000	5396.6\\
1e+05	5446.2\\
};
\addplot [color=mycolor42, forget plot]
  table[row sep=crcr]{%
0	4395.3\\
2000	17077\\
4000	20391\\
6000	21374\\
8000	20441\\
10000	19720\\
12000	20146\\
14000	20671\\
16000	20892\\
18000	21192\\
20000	20854\\
22000	20698\\
26000	20915\\
28000	20620\\
30000	21132\\
32000	20357\\
34000	20108\\
36000	20230\\
38000	20964\\
40000	21874\\
42000	19861\\
44000	20965\\
46000	21237\\
48000	21024\\
50000	21188\\
52000	21115\\
56000	20781\\
58000	19937\\
60000	20714\\
62000	21689\\
64000	20701\\
66000	20060\\
68000	20538\\
70000	20964\\
72000	20284\\
74000	20670\\
76000	20243\\
78000	19743\\
80000	19959\\
82000	20480\\
86000	21137\\
88000	20824\\
90000	21089\\
92000	21041\\
94000	20864\\
98000	20328\\
1e+05	20547\\
};
\addplot [color=mycolor43, forget plot]
  table[row sep=crcr]{%
0	4395.3\\
2000	4918.5\\
4000	4930.2\\
6000	4989.4\\
8000	5014.1\\
10000	4897.2\\
14000	5008.4\\
16000	4953.8\\
18000	4848.4\\
20000	4816.2\\
22000	4906.7\\
24000	4904.4\\
26000	4984.1\\
28000	5037.7\\
30000	4973.2\\
32000	5023.4\\
34000	4934.3\\
40000	4985.5\\
42000	4965.1\\
46000	4861.9\\
48000	4862.7\\
50000	4915.8\\
52000	4829.2\\
54000	4846.6\\
58000	4928\\
60000	4988.9\\
62000	4885.2\\
64000	4906.8\\
66000	4999.2\\
68000	4855.5\\
70000	4850\\
72000	4942.5\\
74000	4983.6\\
76000	4827.4\\
82000	4995.9\\
84000	5000.2\\
88000	4920.9\\
90000	4817.1\\
92000	5028.8\\
94000	4918.7\\
98000	4996.8\\
1e+05	4851.6\\
};
\addplot [color=mycolor44, forget plot]
  table[row sep=crcr]{%
0	4395.3\\
2000	3545.8\\
4000	3353.8\\
6000	3252.9\\
8000	3339.8\\
10000	3203.3\\
12000	3297.3\\
14000	3372.1\\
16000	3248\\
18000	3320.5\\
20000	3409.8\\
22000	3314.9\\
24000	3324.8\\
28000	3306.9\\
30000	3359.8\\
32000	3237.4\\
34000	3239.1\\
36000	3263.1\\
38000	3357.3\\
40000	3427.5\\
42000	3385.9\\
44000	3291.1\\
46000	3307.9\\
48000	3248.7\\
50000	3263.5\\
52000	3337.3\\
54000	3381.9\\
56000	3275.6\\
58000	3363.5\\
60000	3391.4\\
62000	3318.7\\
64000	3389.7\\
66000	3241.1\\
68000	3263.4\\
70000	3243\\
72000	3317.9\\
74000	3434.5\\
76000	3275.5\\
78000	3227.5\\
80000	3219.5\\
82000	3373.4\\
84000	3356.2\\
86000	3218.2\\
88000	3236.5\\
90000	3271.1\\
92000	3383.5\\
94000	3280.9\\
96000	3299.1\\
98000	3410.9\\
1e+05	3298.4\\
};
\addplot [color=mycolor45, forget plot]
  table[row sep=crcr]{%
0	4395.3\\
2000	6728.7\\
4000	6926.1\\
6000	6789.8\\
8000	6963.1\\
10000	6797.9\\
12000	6940.3\\
14000	6952.9\\
16000	7126.8\\
18000	7015.9\\
20000	7008.4\\
22000	7109\\
24000	6892.3\\
26000	6990.5\\
28000	6840.2\\
30000	6659.4\\
32000	6749.3\\
34000	6753.5\\
36000	6993.7\\
38000	7127.1\\
40000	6858.1\\
42000	6655.7\\
44000	6926.2\\
46000	7174.4\\
48000	7075.1\\
50000	6779.3\\
54000	6962.3\\
56000	7072.9\\
58000	7063\\
60000	6870.4\\
62000	7028.4\\
64000	6878.2\\
66000	6970.4\\
68000	6803.3\\
70000	7153.2\\
72000	7026.6\\
74000	7157.3\\
76000	6889\\
78000	6927.4\\
80000	6844.1\\
82000	6947\\
84000	7022.7\\
86000	6908.9\\
88000	7119.7\\
90000	7007\\
92000	7224.3\\
96000	6886.3\\
98000	6928.8\\
1e+05	6924.7\\
};
\addplot [color=mycolor46, forget plot]
  table[row sep=crcr]{%
0	4395.3\\
2000	4089.9\\
4000	4153.5\\
6000	3967.3\\
8000	4087.2\\
10000	4101.5\\
12000	4082.6\\
14000	3968.8\\
16000	4050.9\\
18000	3999\\
20000	4078.8\\
22000	4109.3\\
24000	3992.8\\
26000	4010.4\\
28000	4064.8\\
30000	3972.5\\
32000	4113.6\\
34000	4175.5\\
36000	4022.7\\
38000	4170.9\\
40000	4014.9\\
42000	4160.7\\
44000	4064.5\\
46000	4042.1\\
48000	3992.4\\
50000	4041.8\\
52000	4025.9\\
54000	3962.4\\
56000	4169.7\\
58000	4028.1\\
60000	4229.8\\
62000	3885.4\\
64000	4034.5\\
66000	4009.9\\
68000	4071\\
70000	4091.3\\
72000	3977.5\\
74000	4019.5\\
76000	4145.7\\
80000	4104.8\\
82000	4133.7\\
84000	4117.8\\
86000	4174.1\\
88000	4132.8\\
90000	4246.2\\
98000	4124.7\\
1e+05	3985.1\\
};
\addplot [color=mycolor47, forget plot]
  table[row sep=crcr]{%
0	4395.3\\
2000	5257.8\\
4000	5270.1\\
6000	5369.5\\
8000	5188.3\\
10000	5142.7\\
12000	5317.1\\
14000	5235.8\\
16000	5060.4\\
18000	5120.4\\
20000	5320.2\\
22000	5286\\
24000	5025.4\\
26000	5228.8\\
28000	5360.1\\
30000	5330.6\\
32000	5434.7\\
34000	5218.8\\
36000	5464.2\\
38000	5324.2\\
40000	5359.2\\
42000	5295.5\\
44000	5309.8\\
46000	5239.8\\
48000	5092.3\\
50000	5282.2\\
52000	5202.4\\
54000	5222.2\\
56000	5343.2\\
58000	5337.4\\
60000	5315.2\\
62000	5155\\
64000	5258.6\\
66000	5245.7\\
68000	5376.9\\
70000	5230.7\\
72000	5166.3\\
76000	5283.9\\
78000	5362.3\\
80000	5192.4\\
82000	5137.3\\
84000	5104.4\\
86000	5238.9\\
88000	5103.2\\
90000	5279.6\\
92000	5446.5\\
94000	5128.4\\
96000	5216.3\\
98000	5492.9\\
1e+05	5344.8\\
};
\addplot [color=mycolor48, forget plot]
  table[row sep=crcr]{%
0	4395.3\\
2000	3053.7\\
4000	2896.1\\
6000	2830\\
8000	2849.7\\
10000	2911.6\\
12000	2922.8\\
14000	2917.6\\
16000	2859.4\\
18000	2888.2\\
20000	2978.1\\
22000	2934.3\\
24000	2880.5\\
30000	2861.1\\
32000	2928.3\\
34000	2901.1\\
36000	2923.4\\
38000	2904.7\\
42000	2823\\
44000	2903.4\\
46000	2887.8\\
48000	2949.5\\
50000	2897.4\\
52000	2856.2\\
54000	2916.6\\
56000	2916.6\\
58000	2882.8\\
60000	2913.9\\
62000	2879.2\\
64000	2856.6\\
66000	2851.2\\
68000	2933.9\\
70000	2880.2\\
72000	2975.4\\
74000	2922.7\\
76000	2926.5\\
78000	2941.9\\
80000	2912.8\\
82000	2840\\
84000	2810.3\\
88000	2937\\
90000	2889.7\\
98000	2929.9\\
1e+05	2863.7\\
};
\addplot [color=mycolor49, forget plot]
  table[row sep=crcr]{%
0	4395.3\\
2000	4531.6\\
4000	4409.2\\
6000	4461.6\\
8000	4378.6\\
10000	4380.7\\
12000	4205.6\\
14000	4439.8\\
16000	4566.6\\
18000	4276.3\\
22000	4088.9\\
24000	4283.8\\
26000	4259.4\\
28000	4279.2\\
30000	4274.9\\
32000	4307.9\\
34000	4192.5\\
36000	4348.1\\
38000	4396\\
40000	4308.9\\
44000	4342.7\\
46000	4086.6\\
48000	4316.7\\
50000	4495.5\\
52000	4322.9\\
54000	4323.3\\
56000	4445.8\\
58000	4169.6\\
60000	4129.1\\
62000	4395\\
64000	4348.3\\
66000	4319.4\\
68000	4370.3\\
70000	4455.2\\
72000	4415.5\\
74000	4417.5\\
76000	4299.6\\
78000	4451.4\\
80000	4395.5\\
82000	4256.5\\
84000	4339.3\\
86000	4325.4\\
88000	4349.9\\
90000	4442.3\\
94000	4263.6\\
98000	4350.2\\
1e+05	4492.5\\
};
\addplot [color=mycolor50, forget plot]
  table[row sep=crcr]{%
0	4395.3\\
2000	2646.9\\
4000	2383.7\\
6000	2421.6\\
8000	2452.1\\
10000	2442.6\\
12000	2408.4\\
14000	2459.8\\
16000	2529.2\\
18000	2511.8\\
20000	2420.4\\
22000	2494.2\\
24000	2404.3\\
26000	2383.1\\
28000	2379.7\\
32000	2434.4\\
34000	2417.4\\
36000	2447.5\\
38000	2506.4\\
40000	2425.5\\
42000	2549.6\\
44000	2466.4\\
46000	2440.4\\
48000	2443.9\\
50000	2455.7\\
52000	2474.6\\
54000	2483.2\\
56000	2474.7\\
58000	2390.1\\
60000	2471.6\\
62000	2377.8\\
64000	2441.3\\
66000	2517.3\\
68000	2535.9\\
70000	2461.6\\
72000	2475.3\\
74000	2392.8\\
76000	2449.8\\
78000	2367.5\\
80000	2537.2\\
82000	2431.8\\
84000	2497.6\\
86000	2422.7\\
88000	2409.8\\
90000	2451.8\\
92000	2388.6\\
94000	2415.6\\
96000	2458.2\\
98000	2539.2\\
1e+05	2482\\
};
\addplot [color=mycolor51, forget plot]
  table[row sep=crcr]{%
0	4395.3\\
2000	5932.5\\
4000	5909.3\\
6000	5807\\
8000	6211.1\\
10000	6238.5\\
12000	6062.4\\
14000	6025.5\\
16000	6330.6\\
18000	5660.4\\
20000	5814.2\\
22000	5934.8\\
24000	5813.5\\
26000	5828.5\\
28000	5914\\
30000	6076.9\\
32000	5742.1\\
34000	6230.6\\
36000	6243.8\\
38000	5870.5\\
40000	5949.9\\
42000	5979.9\\
44000	6133.3\\
46000	5741\\
48000	5959\\
50000	5826\\
52000	5767\\
54000	5913.3\\
56000	6106.1\\
58000	5742.3\\
60000	5904.5\\
62000	6004.7\\
64000	6019.4\\
66000	6011.7\\
70000	5926.2\\
72000	5968\\
74000	5873.6\\
76000	5854.7\\
78000	5817.1\\
80000	6033.8\\
82000	5835.8\\
84000	5933.5\\
86000	5955.5\\
88000	6010\\
90000	5884.1\\
92000	6394.5\\
94000	6002.7\\
96000	5845.6\\
98000	5822.2\\
1e+05	5715.7\\
};
\addplot [color=mycolor52, forget plot]
  table[row sep=crcr]{%
0	4395.3\\
2000	3711.8\\
4000	3516.7\\
6000	3445.1\\
8000	3591.8\\
10000	3481\\
12000	3433.2\\
14000	3370.4\\
16000	3498.5\\
18000	3414.2\\
20000	3447.6\\
22000	3456.7\\
24000	3347.9\\
26000	3486.7\\
28000	3452.7\\
30000	3508.1\\
32000	3451\\
34000	3458.7\\
36000	3256.3\\
38000	3394.1\\
40000	3498.7\\
42000	3492.2\\
44000	3501.3\\
46000	3364.8\\
48000	3545.9\\
50000	3514.4\\
52000	3496.5\\
54000	3434.2\\
56000	3499.7\\
58000	3443.9\\
60000	3547.5\\
62000	3472.8\\
64000	3495.4\\
66000	3380.8\\
68000	3619.2\\
70000	3552.7\\
72000	3504.8\\
74000	3539.7\\
76000	3334.7\\
78000	3491.4\\
80000	3445.1\\
82000	3488.2\\
84000	3404.7\\
86000	3438.6\\
88000	3428.6\\
90000	3336.7\\
92000	3268.1\\
94000	3274.4\\
96000	3456.2\\
98000	3401.3\\
1e+05	3277.4\\
};
\addplot [color=mycolor53, forget plot]
  table[row sep=crcr]{%
0	4395.3\\
2000	5583\\
4000	5862.3\\
6000	5488.3\\
8000	6012.1\\
10000	5846\\
12000	5725.8\\
14000	5547.3\\
16000	5733.4\\
18000	5850.4\\
20000	5507.2\\
22000	5629\\
26000	5561.9\\
28000	5766.9\\
30000	5825.1\\
34000	5828.2\\
36000	5674.1\\
38000	5420\\
40000	5552.8\\
42000	5477\\
44000	5447.3\\
46000	5785.4\\
48000	5726.1\\
50000	5550.1\\
52000	5835.7\\
54000	5724.1\\
56000	5548.3\\
58000	5612.1\\
60000	5458.6\\
62000	5621\\
64000	5572.2\\
66000	5908.2\\
70000	5640.3\\
72000	5587.3\\
74000	5551.8\\
76000	5684.1\\
78000	5468.2\\
80000	5647.9\\
84000	5489.9\\
86000	5567.9\\
88000	5485.4\\
90000	5760.3\\
92000	5840.7\\
94000	5541.2\\
96000	5509.3\\
98000	5646.6\\
1e+05	5646.1\\
};
\addplot [color=mycolor54, forget plot]
  table[row sep=crcr]{%
0	4395.3\\
2000	4662.1\\
4000	4489.2\\
6000	4520.9\\
8000	4632.6\\
10000	4559.9\\
12000	4685.3\\
14000	4686.1\\
16000	4705.5\\
18000	4583.6\\
20000	4677.4\\
22000	4706.1\\
24000	4512.3\\
26000	4562.6\\
28000	4547.3\\
30000	4462\\
32000	4342.2\\
34000	4441.5\\
36000	4530\\
38000	4484.6\\
40000	4702.1\\
42000	4613.9\\
44000	4493.7\\
46000	4563.9\\
48000	4432.1\\
50000	4576.6\\
52000	4623.2\\
54000	4588.2\\
56000	4672.2\\
58000	4475.9\\
60000	4656.1\\
62000	4419.3\\
64000	4656.9\\
66000	4666.6\\
68000	4600.6\\
70000	4673.1\\
72000	4474.8\\
74000	4563.9\\
76000	4598.8\\
78000	4613.9\\
80000	4591.4\\
82000	4500.8\\
84000	4488.3\\
86000	4547.9\\
88000	4571.5\\
90000	4635.9\\
92000	4574.5\\
94000	4614.6\\
96000	4549\\
98000	4669.9\\
1e+05	4490.3\\
};
\addplot [color=mycolor55, forget plot]
  table[row sep=crcr]{%
0	4395.3\\
2000	5639.6\\
4000	5725.7\\
6000	5759.3\\
10000	5793.9\\
12000	5767.2\\
14000	5641.8\\
16000	5686.3\\
18000	5878.4\\
20000	5819.2\\
22000	5711.5\\
24000	5738.2\\
26000	5732.8\\
28000	5650.1\\
30000	5751\\
32000	5808.1\\
34000	5830.1\\
36000	5753.2\\
38000	5714.4\\
40000	5927.5\\
42000	5630.5\\
44000	5748.9\\
46000	5732.2\\
48000	5758.1\\
50000	5728.6\\
52000	5793.2\\
54000	5741.7\\
56000	5795.2\\
58000	5814.2\\
62000	5657.1\\
64000	5723.4\\
66000	5962.7\\
68000	5837.6\\
70000	5674.6\\
72000	5718.4\\
74000	5609.7\\
76000	5731.1\\
80000	5861.3\\
82000	5876\\
84000	5696\\
86000	5813\\
88000	5527.1\\
90000	5941.3\\
92000	5859.6\\
94000	5669.2\\
98000	5936.5\\
1e+05	5749.6\\
};
\addplot [color=mycolor56, forget plot]
  table[row sep=crcr]{%
0	4395.3\\
2000	4093.7\\
6000	3934.1\\
8000	3961.4\\
10000	4034.8\\
12000	3898.9\\
14000	4014.3\\
16000	4051.6\\
18000	4018\\
20000	4062.5\\
22000	3994.9\\
24000	4041.5\\
26000	4113.4\\
28000	4080.8\\
30000	3997.1\\
32000	4034.3\\
34000	4085\\
36000	4001.7\\
38000	3961\\
42000	3973.1\\
44000	4041.8\\
46000	3885.7\\
48000	3901\\
50000	4004.9\\
52000	4004.9\\
54000	4096.3\\
56000	4087.1\\
58000	4054.6\\
60000	4040.8\\
62000	3969.8\\
64000	3878.9\\
68000	3997.3\\
70000	4015.5\\
72000	3948.6\\
74000	3950.8\\
76000	4028\\
78000	3895.9\\
80000	4069.6\\
82000	3990.4\\
84000	3981.7\\
86000	4083.9\\
88000	3983.9\\
90000	3927.4\\
92000	3997.1\\
96000	4065.8\\
98000	3874.1\\
1e+05	3920.2\\
};
\addplot [color=mycolor57, forget plot]
  table[row sep=crcr]{%
0	4395.3\\
2000	6327.5\\
4000	6683.3\\
6000	6872.9\\
8000	6700.8\\
10000	6608.1\\
12000	6676.6\\
16000	6688.9\\
18000	6642.5\\
20000	6756.9\\
22000	6741.5\\
24000	6580.7\\
26000	6586.3\\
28000	6817.3\\
30000	6883.3\\
32000	6802.8\\
34000	6702.5\\
36000	6639.9\\
40000	6677.9\\
42000	6826.8\\
44000	6800\\
46000	6706.8\\
48000	6795.2\\
52000	6639.7\\
54000	6877.9\\
58000	6761.4\\
60000	6961.1\\
62000	6887.4\\
64000	6862\\
66000	6733.4\\
68000	6817.5\\
70000	6521.7\\
72000	6647.1\\
74000	6570.1\\
76000	6595.7\\
78000	6815.4\\
80000	6678.7\\
82000	6667.9\\
84000	6817\\
86000	6871.8\\
88000	6651.4\\
90000	6884.7\\
92000	6653.7\\
94000	6760.8\\
98000	6663.3\\
1e+05	6810.9\\
};
\addplot [color=mycolor58, forget plot]
  table[row sep=crcr]{%
0	4395.3\\
2000	3757.6\\
4000	3561\\
6000	3536.8\\
8000	3485.7\\
10000	3530.2\\
12000	3486.8\\
14000	3566.6\\
18000	3604.7\\
20000	3594.7\\
22000	3550.2\\
24000	3610.2\\
26000	3529\\
28000	3584.9\\
30000	3482.4\\
32000	3437.2\\
34000	3601.6\\
36000	3509.6\\
40000	3622.6\\
42000	3537.7\\
44000	3530.8\\
46000	3487.2\\
48000	3426.2\\
50000	3556.3\\
52000	3603.1\\
54000	3531.8\\
56000	3518.1\\
58000	3486.2\\
62000	3532.9\\
64000	3600.2\\
66000	3465.7\\
68000	3597\\
70000	3507\\
72000	3535\\
74000	3433.9\\
78000	3484.5\\
80000	3536.1\\
82000	3540.1\\
84000	3559.7\\
86000	3563.8\\
88000	3583.1\\
90000	3512.1\\
92000	3563.5\\
94000	3519.5\\
96000	3506.3\\
98000	3535.7\\
1e+05	3423.2\\
};
\addplot [color=mycolor59, forget plot]
  table[row sep=crcr]{%
0	4395.3\\
2000	4965.3\\
4000	4926.6\\
6000	4967.3\\
8000	5036\\
10000	5017.6\\
12000	5058.1\\
14000	4887.6\\
16000	5055.9\\
18000	5081.2\\
20000	5182.3\\
22000	5111\\
24000	5132.9\\
26000	5102.5\\
30000	4988.7\\
32000	5049.3\\
34000	4980.8\\
36000	5131.6\\
38000	5083.6\\
40000	4999.5\\
42000	4997\\
44000	4902.7\\
46000	5064.3\\
48000	4975.6\\
50000	5245.2\\
52000	5201.3\\
54000	4946.2\\
56000	5006.2\\
58000	4965.2\\
62000	5082.5\\
64000	5021\\
66000	5080.6\\
68000	5063.2\\
70000	5156.2\\
74000	4998.6\\
76000	5128.1\\
78000	4998.9\\
80000	5107.6\\
82000	4956.5\\
84000	5106.1\\
86000	5132.5\\
88000	5032.4\\
90000	4983.4\\
92000	4994.7\\
94000	4932.4\\
96000	5036\\
98000	5059.4\\
1e+05	5143.5\\
};
\addplot [color=mycolor60, forget plot]
  table[row sep=crcr]{%
0	4395.3\\
2000	3734.1\\
4000	3566\\
6000	3421.2\\
8000	3581.3\\
10000	3500.6\\
12000	3466.4\\
14000	3487.7\\
16000	3475.5\\
18000	3570.1\\
20000	3450.3\\
22000	3573.8\\
24000	3487.4\\
28000	3578.7\\
30000	3575\\
34000	3499.3\\
36000	3496.3\\
38000	3559.8\\
40000	3470.7\\
42000	3530.6\\
44000	3414\\
46000	3534\\
48000	3457\\
50000	3489\\
52000	3493\\
54000	3523.6\\
56000	3491.7\\
58000	3586.1\\
60000	3540.9\\
62000	3460.4\\
64000	3458\\
66000	3578.3\\
68000	3482.3\\
70000	3538.6\\
72000	3571.3\\
74000	3552.7\\
76000	3473.7\\
78000	3606.2\\
80000	3444.8\\
82000	3564.7\\
84000	3546.9\\
86000	3642.2\\
88000	3554.6\\
90000	3481.5\\
92000	3472.5\\
94000	3554.5\\
96000	3607\\
98000	3580\\
1e+05	3570.2\\
};
\addplot [color=mycolor61, forget plot]
  table[row sep=crcr]{%
0	4395.3\\
2000	3351.3\\
4000	3083.6\\
6000	3103.4\\
12000	3077.8\\
14000	3123.6\\
18000	3078.6\\
20000	3103\\
22000	3062\\
24000	3123.6\\
28000	3094.6\\
30000	3142.2\\
34000	3116.8\\
36000	3116.8\\
38000	3055.3\\
40000	3122.6\\
42000	3138\\
44000	3071.7\\
46000	3041.2\\
50000	3107.1\\
52000	3113.3\\
54000	3106.9\\
56000	3047.4\\
58000	3079.9\\
60000	3086.3\\
62000	3118\\
64000	3112\\
68000	3066.8\\
70000	3139\\
72000	3122.4\\
76000	3124.9\\
78000	3123.5\\
80000	3155.4\\
84000	3051.6\\
86000	3110.9\\
88000	3062.2\\
90000	3026.7\\
92000	3128.4\\
94000	3126.8\\
96000	3084.4\\
98000	3031.1\\
1e+05	3194.3\\
};
\addplot [color=mycolor62, forget plot]
  table[row sep=crcr]{%
0	4395.3\\
2000	4875.9\\
4000	4970.7\\
6000	4929\\
8000	4982.3\\
10000	5054\\
12000	5001.5\\
14000	5014.2\\
16000	4955.7\\
18000	5084.6\\
20000	4860.2\\
22000	4948.5\\
24000	4986.4\\
26000	4979.1\\
28000	5024.3\\
30000	4938.8\\
32000	4956.8\\
34000	5076.4\\
36000	5016.7\\
38000	4887.6\\
42000	4935\\
44000	5065.4\\
46000	4886.1\\
48000	5003.3\\
50000	4938.4\\
52000	5054.7\\
54000	4974.8\\
56000	5023.1\\
58000	4992.7\\
60000	4933.4\\
62000	4998.1\\
64000	4879.2\\
66000	4950.9\\
68000	5047.8\\
70000	5084.4\\
72000	5008.8\\
74000	4964.5\\
76000	5010.3\\
78000	4962.7\\
80000	5023.9\\
82000	4946.6\\
84000	4917.3\\
86000	4945.1\\
88000	4899.9\\
90000	5090.6\\
92000	4940.8\\
94000	5004.3\\
96000	4890\\
98000	4943.8\\
1e+05	5051.2\\
};
\addplot [color=mycolor63, forget plot]
  table[row sep=crcr]{%
0	4395.3\\
2000	2635.4\\
4000	2443.4\\
6000	2388.1\\
8000	2460.3\\
10000	2396.9\\
12000	2344.9\\
14000	2314.7\\
16000	2406.1\\
18000	2347\\
20000	2380.4\\
22000	2372.9\\
24000	2449\\
26000	2353.7\\
32000	2440.8\\
34000	2404.5\\
36000	2404.4\\
38000	2396.7\\
40000	2418.8\\
42000	2350\\
44000	2371.3\\
46000	2321.8\\
48000	2303.8\\
50000	2334.8\\
52000	2385.9\\
54000	2377.5\\
56000	2387.2\\
58000	2365\\
60000	2390\\
62000	2356.7\\
64000	2359.4\\
66000	2381.2\\
68000	2422.4\\
72000	2402.3\\
74000	2430.5\\
76000	2484.3\\
78000	2448\\
80000	2432.8\\
82000	2376.2\\
84000	2356.1\\
86000	2360.1\\
88000	2325.8\\
92000	2361.1\\
94000	2395\\
96000	2349.9\\
98000	2406.1\\
1e+05	2376.5\\
};
\addplot [color=mycolor64, forget plot]
  table[row sep=crcr]{%
0	4395.3\\
2000	3140.5\\
4000	3045.9\\
6000	3004.9\\
8000	2997.6\\
10000	3007.7\\
12000	2977.1\\
14000	3014.9\\
16000	3036.6\\
18000	2962.2\\
20000	2966.5\\
22000	3004.2\\
24000	3021.6\\
26000	2995.8\\
28000	3028\\
30000	3026.1\\
32000	3038.3\\
34000	3040.9\\
36000	3111.6\\
38000	3000.1\\
40000	3074.1\\
42000	2988.5\\
44000	3059.3\\
46000	3031.6\\
48000	3101.5\\
50000	3080.8\\
52000	3003.9\\
54000	3021.6\\
56000	2964.9\\
58000	2968.1\\
60000	3031.2\\
62000	3049.8\\
64000	3033.2\\
66000	2989.6\\
68000	3072.1\\
72000	3004.8\\
74000	3069.1\\
76000	3031.1\\
78000	3041\\
80000	2972.5\\
86000	3061.9\\
88000	3008.4\\
90000	3003.4\\
92000	2972.1\\
94000	3027.8\\
96000	2943.9\\
98000	3108.3\\
1e+05	3016.1\\
};
\addplot [color=mycolor65, forget plot]
  table[row sep=crcr]{%
0	4395.3\\
2000	2708.3\\
4000	2484.5\\
6000	2556.3\\
8000	2520.9\\
10000	2606.6\\
12000	2563\\
14000	2557.7\\
16000	2502.9\\
18000	2567.4\\
20000	2427.7\\
22000	2474.9\\
24000	2512.3\\
28000	2541.1\\
30000	2526.7\\
32000	2502.6\\
34000	2560.5\\
36000	2529.6\\
38000	2523.5\\
40000	2533\\
42000	2507.2\\
46000	2515.2\\
48000	2557.9\\
50000	2540.3\\
52000	2553.9\\
54000	2550.6\\
56000	2448\\
58000	2479.6\\
60000	2470.2\\
62000	2561.1\\
64000	2483.8\\
66000	2501.3\\
68000	2538.7\\
70000	2539.2\\
72000	2517.9\\
76000	2504\\
78000	2510.9\\
80000	2584.5\\
82000	2570.1\\
84000	2501.1\\
86000	2526\\
88000	2543.3\\
90000	2478.7\\
92000	2463\\
94000	2524\\
96000	2518\\
98000	2493.9\\
1e+05	2495.5\\
};
\addplot [color=mycolor66, forget plot]
  table[row sep=crcr]{%
0	4395.3\\
2000	3517.4\\
4000	3380.4\\
6000	3412.1\\
8000	3431.1\\
10000	3383.5\\
12000	3457.4\\
14000	3392.5\\
18000	3410.8\\
20000	3432.4\\
22000	3438.5\\
26000	3387.8\\
28000	3415\\
30000	3335.6\\
34000	3491.1\\
36000	3425.8\\
38000	3436.1\\
40000	3376.7\\
42000	3457.2\\
44000	3496.9\\
46000	3437.8\\
48000	3396.4\\
50000	3414.2\\
52000	3395\\
54000	3399.5\\
56000	3447.8\\
58000	3462.1\\
60000	3375.9\\
62000	3408.1\\
64000	3367.5\\
66000	3399.7\\
68000	3420\\
70000	3427.6\\
82000	3387.8\\
84000	3493.5\\
86000	3477.1\\
88000	3390.7\\
90000	3448.6\\
96000	3367.5\\
98000	3391\\
1e+05	3484.2\\
};
\addplot [color=mycolor67, forget plot]
  table[row sep=crcr]{%
0	4395.3\\
2000	3387.2\\
4000	3199.2\\
6000	3176\\
8000	3163.2\\
10000	3123.2\\
12000	3211.5\\
14000	3221.3\\
16000	3197.2\\
18000	3160.1\\
20000	3203.7\\
22000	3072.7\\
24000	3126.5\\
26000	3142\\
28000	3147\\
30000	3128.7\\
32000	3156.7\\
34000	3160.7\\
36000	3233.6\\
38000	3196.9\\
40000	3086.1\\
42000	3132.1\\
44000	3155.8\\
46000	3101\\
48000	3097.8\\
50000	3163.6\\
52000	3125.7\\
54000	3114.9\\
56000	3151.7\\
58000	3121.4\\
60000	3198.7\\
62000	3172.9\\
64000	3179.5\\
68000	3213.9\\
70000	3175.5\\
72000	3176.6\\
74000	3155.9\\
76000	3170.3\\
78000	3118.7\\
80000	3189\\
82000	3143.3\\
84000	3085.9\\
86000	3169.2\\
88000	3106.8\\
90000	3186.6\\
92000	3178\\
94000	3193.6\\
96000	3154.9\\
98000	3157.5\\
1e+05	3213\\
};
\addplot [color=mycolor67]
  table[row sep=crcr]{%
0	4395.3\\
2000	2806.6\\
4000	2611\\
10000	2592.5\\
12000	2614.1\\
16000	2600.3\\
20000	2626.5\\
22000	2615\\
24000	2574.9\\
26000	2557.5\\
28000	2634.3\\
30000	2614.5\\
32000	2582.1\\
34000	2622.1\\
38000	2609.4\\
40000	2624.8\\
42000	2573.9\\
46000	2624.2\\
48000	2566.9\\
50000	2602.7\\
52000	2594.1\\
54000	2576.9\\
56000	2587.1\\
58000	2581.7\\
60000	2591.2\\
62000	2580.2\\
68000	2625\\
70000	2617.1\\
72000	2595.3\\
74000	2621.3\\
76000	2631.6\\
78000	2595\\
82000	2610.7\\
84000	2610.4\\
86000	2635.2\\
88000	2652.4\\
90000	2637.7\\
92000	2612.6\\
94000	2595.6\\
98000	2642.1\\
1e+05	2635.8\\
};
\addlegendentry{\ClasskIIILegend}

\addplot [color=mycolor68, forget plot]
  table[row sep=crcr]{%
0	4395.3\\
2000	2692.4\\
4000	2425.7\\
6000	2435.1\\
8000	2412.2\\
10000	2446.1\\
12000	2416.5\\
14000	2411.2\\
18000	2441.7\\
20000	2481.6\\
22000	2435.3\\
26000	2444.4\\
28000	2361.3\\
32000	2457.3\\
34000	2428.6\\
36000	2475.8\\
38000	2411.5\\
42000	2456.9\\
44000	2426.4\\
46000	2421.6\\
48000	2405.1\\
50000	2429\\
54000	2442.5\\
56000	2429.4\\
58000	2395.3\\
60000	2455\\
62000	2430\\
64000	2454.9\\
66000	2428.9\\
68000	2443.1\\
72000	2432.7\\
74000	2437.5\\
76000	2421.2\\
78000	2375\\
80000	2382.4\\
82000	2417.5\\
84000	2467.6\\
86000	2433.7\\
88000	2417.8\\
90000	2470.6\\
92000	2430.5\\
94000	2475.5\\
96000	2435.7\\
98000	2453.2\\
1e+05	2418.4\\
};
\addplot [color=mycolor69, forget plot]
  table[row sep=crcr]{%
0	4395.3\\
2000	3109.6\\
4000	2844.6\\
6000	2903.7\\
8000	2886.3\\
10000	2909.3\\
12000	2882.9\\
16000	2885.8\\
18000	2853.3\\
22000	2875.5\\
24000	2908.3\\
26000	2850.9\\
30000	2874.5\\
32000	2914.7\\
34000	2896.2\\
40000	2912.3\\
42000	2875.3\\
46000	2934.1\\
48000	2886.6\\
52000	2905.4\\
54000	2888.5\\
56000	2886.3\\
58000	2852.7\\
60000	2900.3\\
62000	2872.2\\
64000	2885.2\\
66000	2868.2\\
68000	2882.8\\
72000	2875.5\\
74000	2900.2\\
76000	2857.1\\
80000	2908.3\\
84000	2881.7\\
86000	2884.5\\
88000	2856.9\\
92000	2841.4\\
96000	2910.5\\
98000	2927.3\\
1e+05	2912.1\\
};
\addplot [color=mycolor70, forget plot]
  table[row sep=crcr]{%
0	4395.3\\
2000	2406.2\\
4000	2030.2\\
6000	2081.5\\
8000	2058.7\\
10000	2089.7\\
12000	2111.9\\
14000	2095.9\\
16000	2059.9\\
18000	2089.7\\
20000	2110.6\\
22000	2062.1\\
24000	2097\\
26000	2095.2\\
28000	2047.3\\
30000	2069.2\\
32000	2048.3\\
34000	2102.2\\
36000	2078.9\\
38000	2031.3\\
40000	2109.8\\
42000	2068\\
44000	2088.6\\
46000	2040.2\\
48000	2077.6\\
50000	2091.3\\
52000	2064.4\\
54000	2082.8\\
56000	2112.2\\
60000	2049.1\\
62000	2082.5\\
64000	2123.9\\
66000	2096.4\\
68000	2091.3\\
70000	2099.7\\
72000	2098.6\\
74000	2065.4\\
76000	2055\\
80000	2079.7\\
82000	2032.1\\
84000	2086.4\\
86000	2127.2\\
88000	2101.1\\
90000	2097.2\\
92000	2054.3\\
94000	2056.6\\
98000	2087.6\\
1e+05	2081\\
};
\addplot [color=mycolor71, forget plot]
  table[row sep=crcr]{%
0	4395.3\\
2000	2730.9\\
4000	2428.4\\
6000	2451.1\\
8000	2454.9\\
10000	2444.8\\
12000	2463.9\\
14000	2520.9\\
16000	2488\\
18000	2441.9\\
22000	2463.5\\
28000	2466.2\\
32000	2451.5\\
34000	2488.6\\
36000	2474.7\\
38000	2433.4\\
40000	2474.5\\
42000	2427.9\\
44000	2510.4\\
46000	2470.4\\
48000	2497.5\\
50000	2487.1\\
52000	2433.7\\
56000	2467.9\\
58000	2463.4\\
60000	2433.2\\
64000	2462.3\\
66000	2471.2\\
68000	2435.4\\
70000	2458.9\\
72000	2435.6\\
74000	2469.2\\
78000	2494.1\\
80000	2459.3\\
82000	2483\\
86000	2443.9\\
88000	2469.4\\
90000	2479.1\\
92000	2505.4\\
96000	2444.8\\
98000	2438.3\\
1e+05	2420.4\\
};
\addplot [color=mycolor72, forget plot]
  table[row sep=crcr]{%
0	4395.3\\
2000	2482.2\\
4000	2212\\
6000	2167.8\\
8000	2139.2\\
10000	2194.2\\
12000	2177.7\\
14000	2181.3\\
16000	2193.3\\
18000	2170.7\\
20000	2136.3\\
22000	2168.5\\
24000	2193.6\\
26000	2202.4\\
30000	2171.3\\
32000	2153.4\\
34000	2179.2\\
36000	2140.4\\
38000	2196.7\\
40000	2185.1\\
42000	2141.2\\
44000	2181.5\\
46000	2110.2\\
48000	2162.8\\
50000	2189.6\\
52000	2164\\
54000	2216\\
56000	2164.8\\
58000	2156.2\\
60000	2183.8\\
62000	2235.1\\
64000	2191.9\\
66000	2134.5\\
68000	2239\\
70000	2198.7\\
72000	2146.6\\
74000	2184.9\\
76000	2192.5\\
78000	2162.8\\
80000	2165.2\\
82000	2155\\
84000	2137.5\\
86000	2178.7\\
88000	2201.5\\
90000	2165.8\\
92000	2183.5\\
94000	2229.4\\
96000	2144.2\\
98000	2147.3\\
1e+05	2207\\
};
\addplot [color=mycolor73, forget plot]
  table[row sep=crcr]{%
0	4395.3\\
2000	2859.5\\
4000	2610.2\\
6000	2615.4\\
8000	2565.8\\
10000	2577.2\\
12000	2581\\
14000	2608.2\\
16000	2553.3\\
18000	2550.4\\
20000	2559.2\\
22000	2550.3\\
24000	2587.5\\
26000	2559.2\\
28000	2635.3\\
30000	2608.8\\
32000	2618.2\\
34000	2537.7\\
36000	2546.5\\
38000	2591.3\\
40000	2531.7\\
42000	2589\\
44000	2586.1\\
46000	2634.7\\
48000	2619.4\\
50000	2656.9\\
52000	2579.2\\
54000	2614.6\\
56000	2600\\
58000	2597.3\\
60000	2564.2\\
62000	2590\\
64000	2566.4\\
66000	2592.5\\
68000	2584.3\\
70000	2603.5\\
72000	2557.9\\
74000	2597\\
76000	2566.8\\
78000	2615.5\\
80000	2570.8\\
82000	2614.8\\
84000	2573.9\\
86000	2571.1\\
88000	2491.9\\
90000	2565.8\\
92000	2543.9\\
96000	2553.5\\
98000	2598.3\\
1e+05	2546.7\\
};
\addplot [color=mycolor74, forget plot]
  table[row sep=crcr]{%
0	4395.3\\
2000	1884.1\\
4000	1527.1\\
6000	1532.7\\
8000	1521.1\\
12000	1517\\
14000	1492\\
16000	1511.9\\
18000	1527.7\\
20000	1497.1\\
22000	1537.4\\
24000	1527.8\\
26000	1469.8\\
28000	1504.2\\
30000	1488.8\\
32000	1500.4\\
36000	1514.2\\
38000	1504.9\\
40000	1530.5\\
44000	1517.1\\
46000	1522.2\\
48000	1501.5\\
50000	1488\\
52000	1512.1\\
54000	1496.3\\
56000	1510.3\\
62000	1536\\
64000	1505.6\\
68000	1502.9\\
70000	1547.3\\
72000	1541.7\\
74000	1515.4\\
76000	1506.3\\
78000	1502.2\\
80000	1506.8\\
82000	1535.3\\
84000	1531.7\\
86000	1489.1\\
88000	1510.7\\
90000	1524.1\\
92000	1468.4\\
94000	1490.6\\
96000	1494.8\\
98000	1488.4\\
1e+05	1504.1\\
};
\addplot [color=mycolor75, forget plot]
  table[row sep=crcr]{%
0	4395.3\\
2000	2518.4\\
4000	2206.1\\
6000	2184.2\\
8000	2199.2\\
10000	2177\\
12000	2171.7\\
14000	2226.2\\
16000	2184.2\\
18000	2176.9\\
20000	2196.6\\
22000	2224.6\\
26000	2151.8\\
28000	2177.5\\
34000	2168.5\\
36000	2194.4\\
38000	2199.9\\
40000	2191.3\\
42000	2222\\
44000	2218.8\\
46000	2169.5\\
48000	2174.7\\
50000	2223.7\\
54000	2188.8\\
56000	2166.1\\
58000	2228.3\\
62000	2169\\
64000	2190.7\\
66000	2184.1\\
68000	2192.9\\
70000	2154.3\\
72000	2180.4\\
74000	2230.8\\
76000	2216.4\\
78000	2181.4\\
82000	2168.2\\
84000	2201.1\\
86000	2175.9\\
88000	2167.3\\
90000	2196.2\\
92000	2161.5\\
94000	2145.3\\
96000	2189.6\\
98000	2155.6\\
1e+05	2182.4\\
};
\addplot [color=mycolor76, forget plot]
  table[row sep=crcr]{%
0	4395.3\\
2000	2352.1\\
4000	2014.5\\
6000	1968.9\\
8000	1978.4\\
10000	1939\\
14000	1969.7\\
16000	1980.8\\
18000	2018\\
20000	2021.1\\
22000	1995.3\\
24000	1980.1\\
26000	1970.9\\
28000	1982.6\\
30000	1984\\
32000	1969.7\\
34000	1995.8\\
36000	1983.4\\
38000	1977.4\\
40000	1991.6\\
42000	1935.1\\
44000	1989.9\\
46000	1989.4\\
48000	1962.3\\
50000	1981.5\\
52000	1989.9\\
58000	1992.6\\
60000	1972.5\\
62000	1994.6\\
64000	2008.7\\
66000	1971.4\\
68000	1979.6\\
70000	1953.6\\
72000	1960.2\\
74000	1980.2\\
76000	1970.5\\
78000	1982.9\\
80000	2002.1\\
82000	2002.6\\
84000	1957.4\\
86000	1980.4\\
88000	1972\\
90000	1994.6\\
94000	1953.3\\
96000	1978.1\\
98000	2011.3\\
1e+05	1999.6\\
};
\addplot [color=mycolor77, forget plot]
  table[row sep=crcr]{%
0	4395.3\\
2000	1932.9\\
4000	1534.3\\
6000	1524.7\\
8000	1540\\
10000	1530.6\\
12000	1501.7\\
14000	1542.8\\
16000	1529.7\\
18000	1527\\
20000	1533.1\\
22000	1581\\
24000	1520.1\\
26000	1541.5\\
28000	1513.4\\
30000	1524.8\\
32000	1522.1\\
36000	1544.8\\
38000	1519.5\\
40000	1539.7\\
42000	1527.6\\
44000	1526.6\\
46000	1555.5\\
50000	1534\\
52000	1535.8\\
54000	1576.4\\
56000	1522.2\\
58000	1532.3\\
60000	1564.6\\
64000	1534.3\\
66000	1536.4\\
68000	1517\\
70000	1555\\
72000	1549.6\\
74000	1559.1\\
78000	1520\\
82000	1553.9\\
84000	1534.9\\
86000	1553\\
90000	1524\\
92000	1520.8\\
94000	1540.4\\
96000	1554.9\\
98000	1538.9\\
1e+05	1552.4\\
};
\addplot [color=mycolor78, forget plot]
  table[row sep=crcr]{%
0	4395.3\\
2000	2568.5\\
4000	2297.3\\
6000	2286.6\\
8000	2253.9\\
10000	2268.1\\
12000	2213.8\\
14000	2242\\
16000	2283.7\\
18000	2270.4\\
20000	2240.1\\
22000	2307.2\\
26000	2231.1\\
28000	2304.9\\
30000	2243.2\\
34000	2249.5\\
36000	2283.5\\
40000	2292.6\\
42000	2276.9\\
44000	2237.8\\
50000	2266.8\\
52000	2306.4\\
56000	2236.6\\
58000	2271.2\\
60000	2237.6\\
62000	2262\\
64000	2215.6\\
66000	2223.5\\
68000	2237.9\\
70000	2307.3\\
72000	2222.5\\
74000	2267.5\\
76000	2287.7\\
78000	2275.3\\
80000	2289.5\\
82000	2269.3\\
84000	2257.3\\
88000	2286.9\\
90000	2257.1\\
92000	2238.9\\
96000	2281.1\\
98000	2223.7\\
1e+05	2247.1\\
};
\addplot [color=mycolor79, forget plot]
  table[row sep=crcr]{%
0	4395.3\\
2000	2519.5\\
4000	2103.5\\
6000	2157.7\\
8000	2115.8\\
10000	2117.6\\
12000	2160.8\\
14000	2197.4\\
16000	2076.6\\
22000	2113.1\\
24000	2138.2\\
26000	2086.2\\
28000	2186.4\\
30000	2174.9\\
32000	2133\\
34000	2130.1\\
36000	2064.1\\
38000	2152.2\\
40000	2131\\
42000	2144\\
44000	2116.2\\
48000	2118.8\\
50000	2097.8\\
52000	2127.4\\
54000	2137.4\\
56000	2128.2\\
58000	2159\\
60000	2152.4\\
62000	2156.4\\
64000	2110.2\\
66000	2138.8\\
68000	2136.2\\
70000	2146.6\\
72000	2135.5\\
74000	2104.5\\
76000	2120.4\\
78000	2099.2\\
80000	2137.8\\
82000	2113.5\\
84000	2128.4\\
86000	2173.7\\
88000	2173.3\\
90000	2153.9\\
92000	2155.7\\
94000	2147.9\\
98000	2114.3\\
1e+05	2089.5\\
};
\addplot [color=mycolor79, forget plot]
  table[row sep=crcr]{%
0	4395.3\\
2000	2398.3\\
4000	2016.5\\
6000	1992.1\\
8000	1997.2\\
10000	1996.4\\
12000	2023.1\\
14000	2007.1\\
16000	1982.5\\
18000	2038.5\\
20000	1992.6\\
22000	2015.9\\
24000	1973.9\\
26000	2008.7\\
28000	1985.4\\
30000	2029.8\\
34000	1950.6\\
36000	1998.8\\
38000	1991.2\\
40000	1945.4\\
42000	1943.2\\
44000	1975\\
46000	2039.9\\
48000	1962.6\\
50000	1966.9\\
52000	1963.2\\
54000	2022\\
56000	1964\\
58000	1991.4\\
60000	1979\\
62000	2002.3\\
64000	1987.3\\
66000	2024\\
68000	1985.8\\
70000	1961.4\\
74000	1962.8\\
76000	1967.5\\
78000	1986.2\\
80000	1972.7\\
82000	1991.7\\
86000	1989.6\\
88000	1954.6\\
90000	1977.2\\
92000	1953.7\\
94000	2010.5\\
96000	2014.5\\
98000	1971.2\\
1e+05	1969.3\\
};
\addplot [color=mycolor80, forget plot]
  table[row sep=crcr]{%
0	4395.3\\
2000	1989.6\\
4000	1552.9\\
6000	1529.9\\
8000	1544.5\\
10000	1508.2\\
12000	1510.4\\
14000	1525\\
16000	1520.6\\
18000	1522.4\\
20000	1512.1\\
22000	1522.5\\
24000	1511.7\\
26000	1514.7\\
28000	1524.3\\
30000	1518.9\\
32000	1528.3\\
34000	1516\\
36000	1535.6\\
38000	1533.3\\
40000	1554\\
42000	1547.9\\
44000	1557.4\\
46000	1545.3\\
48000	1539.2\\
50000	1500\\
54000	1549.5\\
56000	1516.8\\
58000	1509\\
60000	1523.2\\
62000	1526.7\\
66000	1550.1\\
68000	1533.5\\
70000	1550.7\\
72000	1534.9\\
74000	1493.9\\
76000	1520.6\\
78000	1509.4\\
80000	1521.3\\
82000	1568.1\\
84000	1518.8\\
86000	1522.1\\
88000	1543.2\\
92000	1542.5\\
94000	1533.5\\
96000	1513\\
98000	1571.9\\
1e+05	1501.7\\
};
\addplot [color=mycolor81, forget plot]
  table[row sep=crcr]{%
0	4395.3\\
2000	1857.4\\
4000	1420.1\\
6000	1430.3\\
10000	1462.2\\
12000	1432.2\\
14000	1419.7\\
16000	1429.3\\
18000	1428.2\\
20000	1412.5\\
24000	1440.3\\
26000	1433.6\\
30000	1440.3\\
32000	1424.1\\
34000	1423.3\\
36000	1411.9\\
38000	1426.3\\
40000	1412.2\\
42000	1407.6\\
44000	1413.3\\
46000	1438.8\\
48000	1425.6\\
50000	1431.4\\
52000	1459.1\\
54000	1435.8\\
56000	1420.3\\
58000	1430.5\\
60000	1446.3\\
62000	1442.2\\
68000	1443.3\\
70000	1440.8\\
72000	1448.6\\
76000	1432.7\\
78000	1420.7\\
82000	1449.6\\
84000	1472.1\\
86000	1458.8\\
88000	1431\\
90000	1420.8\\
92000	1439.9\\
94000	1431.4\\
1e+05	1433.3\\
};
\addplot [color=mycolor82, forget plot]
  table[row sep=crcr]{%
0	4395.3\\
2000	2102.4\\
4000	1691.3\\
6000	1682.6\\
8000	1680.3\\
10000	1686.2\\
14000	1671.8\\
16000	1689.8\\
18000	1696.8\\
20000	1672\\
24000	1703.4\\
26000	1676.8\\
28000	1639.6\\
30000	1644.5\\
32000	1674.1\\
34000	1668.4\\
36000	1701.6\\
38000	1681.2\\
40000	1675.3\\
42000	1648.2\\
44000	1692.2\\
46000	1659.9\\
50000	1681.9\\
52000	1659.9\\
54000	1681\\
58000	1692\\
62000	1638.1\\
64000	1681.4\\
66000	1701.8\\
68000	1710.4\\
70000	1681.8\\
72000	1677.3\\
74000	1667.2\\
76000	1673.8\\
78000	1667.4\\
80000	1694.2\\
82000	1687.3\\
84000	1709.7\\
88000	1625\\
90000	1663.3\\
92000	1661.8\\
94000	1693.3\\
96000	1705.6\\
98000	1669.6\\
1e+05	1683.2\\
};
\addplot [color=mycolor83, forget plot]
  table[row sep=crcr]{%
0	4395.3\\
2000	2156.1\\
4000	1711\\
6000	1681.9\\
8000	1726.3\\
10000	1680.7\\
12000	1711.4\\
14000	1703.5\\
16000	1680.6\\
18000	1701.6\\
20000	1706.7\\
22000	1682.2\\
24000	1697.4\\
26000	1703.4\\
28000	1723\\
30000	1699.8\\
32000	1682.5\\
34000	1730.4\\
38000	1680.7\\
40000	1677.7\\
42000	1706.8\\
44000	1709\\
46000	1680.7\\
48000	1706.1\\
50000	1723.1\\
52000	1683\\
54000	1669.5\\
56000	1683.8\\
58000	1707.9\\
60000	1699.7\\
62000	1682\\
66000	1711.2\\
68000	1719.9\\
70000	1691.4\\
72000	1703.9\\
74000	1694.9\\
76000	1696.9\\
78000	1670.1\\
80000	1667.6\\
82000	1711.9\\
84000	1706.3\\
86000	1686.1\\
90000	1722.2\\
92000	1700.6\\
96000	1725.8\\
98000	1704.2\\
1e+05	1699.8\\
};
\addplot [color=mycolor83, forget plot]
  table[row sep=crcr]{%
0	4395.3\\
2000	1706\\
4000	1252.6\\
6000	1239.7\\
8000	1238.6\\
10000	1260\\
12000	1220.4\\
14000	1228.4\\
16000	1246.2\\
18000	1253.5\\
20000	1230.7\\
22000	1257.8\\
24000	1264.4\\
26000	1248.7\\
28000	1221.9\\
30000	1223.7\\
32000	1252.5\\
34000	1213.7\\
36000	1253.1\\
40000	1228.1\\
42000	1270.6\\
44000	1242\\
46000	1233.6\\
48000	1232.4\\
50000	1246.8\\
52000	1249.1\\
54000	1230.8\\
56000	1271.8\\
58000	1226.4\\
60000	1246.9\\
62000	1214.4\\
64000	1225.5\\
66000	1257.6\\
68000	1244\\
70000	1288.5\\
72000	1236.4\\
74000	1254.7\\
76000	1248.3\\
78000	1253.2\\
82000	1246.7\\
84000	1211.2\\
86000	1249.1\\
88000	1260.8\\
90000	1218.3\\
92000	1226.4\\
94000	1252.3\\
96000	1235.3\\
98000	1246.1\\
1e+05	1272.4\\
};
\addplot [color=mycolor84, forget plot]
  table[row sep=crcr]{%
0	4395.3\\
2000	1757.2\\
4000	1242.4\\
6000	1228.6\\
8000	1230\\
10000	1238\\
14000	1240.1\\
16000	1259.3\\
18000	1273.5\\
20000	1211.1\\
22000	1267.5\\
26000	1233\\
28000	1227.1\\
30000	1293.5\\
32000	1254.9\\
34000	1233.2\\
36000	1275.6\\
40000	1229.1\\
42000	1233.4\\
44000	1224.8\\
46000	1251.4\\
48000	1250\\
50000	1266.6\\
52000	1276.5\\
54000	1261\\
56000	1236.9\\
58000	1269\\
60000	1224.9\\
62000	1261.4\\
64000	1245.5\\
66000	1259.7\\
68000	1259.6\\
70000	1240.1\\
72000	1264.7\\
76000	1233.3\\
78000	1266.9\\
80000	1290.7\\
82000	1229.5\\
84000	1244.3\\
86000	1279.9\\
88000	1247.2\\
90000	1254.6\\
94000	1216.7\\
96000	1245.2\\
98000	1245.8\\
1e+05	1208.4\\
};
\addplot [color=mycolor85, forget plot]
  table[row sep=crcr]{%
0	4395.3\\
2000	1863.8\\
4000	1378.2\\
8000	1361.3\\
10000	1340.5\\
12000	1403.6\\
14000	1356.2\\
16000	1382.1\\
18000	1371.6\\
20000	1346.1\\
22000	1388.3\\
24000	1356.8\\
26000	1363\\
28000	1382.1\\
30000	1373.1\\
32000	1374.4\\
34000	1338.9\\
36000	1334.1\\
38000	1415.5\\
40000	1393\\
42000	1363.6\\
44000	1362.9\\
46000	1383\\
48000	1379.9\\
50000	1384.1\\
52000	1367.6\\
56000	1366.1\\
58000	1369.6\\
60000	1344.4\\
62000	1337.4\\
64000	1359.6\\
66000	1412.9\\
68000	1377.7\\
70000	1337.9\\
72000	1370.3\\
74000	1371.3\\
76000	1380.7\\
78000	1362.1\\
80000	1380.6\\
82000	1357\\
84000	1346.9\\
86000	1382.2\\
88000	1362.9\\
90000	1370.7\\
92000	1374.4\\
94000	1389.2\\
96000	1424.9\\
98000	1336.6\\
1e+05	1353.5\\
};
\addplot [color=mycolor86, forget plot]
  table[row sep=crcr]{%
0	4395.3\\
2000	1718.9\\
4000	1145.3\\
6000	1177.2\\
8000	1181.3\\
10000	1177.4\\
12000	1166.4\\
14000	1130.6\\
16000	1163\\
18000	1137.7\\
22000	1146.8\\
24000	1144.7\\
26000	1159.9\\
28000	1147\\
30000	1148.8\\
34000	1173.4\\
36000	1149.1\\
38000	1143.1\\
40000	1117.5\\
42000	1157.5\\
44000	1155.7\\
46000	1175\\
48000	1150.6\\
50000	1142.9\\
52000	1118.5\\
54000	1130.5\\
56000	1159.7\\
58000	1166.6\\
60000	1142.4\\
62000	1154.4\\
64000	1143.7\\
66000	1158.8\\
68000	1128.9\\
70000	1175.8\\
72000	1162.5\\
74000	1187.2\\
76000	1157.6\\
78000	1168.2\\
80000	1142.6\\
82000	1180.3\\
84000	1164.4\\
86000	1138.2\\
88000	1142\\
90000	1107.5\\
92000	1142.5\\
94000	1132\\
96000	1127.1\\
98000	1171.8\\
1e+05	1137.4\\
};
\addplot [color=mycolor87, forget plot]
  table[row sep=crcr]{%
0	4395.3\\
2000	1608.2\\
4000	975.53\\
6000	966.68\\
8000	1084.5\\
10000	961.62\\
12000	1004.1\\
14000	996.95\\
18000	991.82\\
20000	983.9\\
22000	960.51\\
24000	975.07\\
26000	949.04\\
28000	960.36\\
30000	1022.2\\
34000	979.16\\
36000	1003.4\\
38000	960.56\\
42000	1009.9\\
44000	1017.7\\
46000	992.25\\
48000	987.35\\
50000	967.6\\
54000	970.91\\
56000	986.07\\
58000	983.33\\
60000	1007.7\\
62000	987.6\\
64000	956.85\\
66000	981.89\\
68000	965.89\\
70000	1004.9\\
72000	943.74\\
74000	972.58\\
76000	966.27\\
78000	969.2\\
80000	984.07\\
82000	978.31\\
84000	990.81\\
86000	985.74\\
88000	957.92\\
90000	984.38\\
92000	1016.3\\
94000	998.47\\
96000	976.89\\
98000	967.63\\
1e+05	993.08\\
};
\addplot [color=mycolor88, forget plot]
  table[row sep=crcr]{%
0	4395.3\\
2000	1786.7\\
4000	1143.2\\
6000	1182.4\\
8000	1137.3\\
10000	1111.6\\
12000	1133.6\\
14000	1152.4\\
16000	1108.9\\
18000	1144.5\\
20000	1174.5\\
22000	1162.1\\
24000	1164.5\\
26000	1132.6\\
30000	1122.4\\
32000	1163.5\\
34000	1165\\
36000	1149.1\\
38000	1177.5\\
40000	1141.5\\
42000	1133\\
44000	1119.8\\
46000	1162\\
50000	1133.1\\
52000	1138\\
54000	1132.5\\
56000	1169\\
58000	1127.1\\
60000	1154.5\\
62000	1141.8\\
64000	1147\\
66000	1128.8\\
68000	1140.7\\
70000	1171.8\\
72000	1117.5\\
74000	1157\\
76000	1142.7\\
78000	1167.9\\
80000	1147.9\\
84000	1137.7\\
86000	1141\\
88000	1133.8\\
92000	1153.1\\
94000	1142.6\\
96000	1161.9\\
98000	1120.5\\
1e+05	1166.3\\
};
\addplot [color=mycolor89, forget plot]
  table[row sep=crcr]{%
0	4395.3\\
2000	1461\\
4000	809.63\\
6000	808.56\\
8000	803.92\\
10000	836.04\\
12000	809.21\\
14000	813.17\\
18000	813.09\\
20000	817.95\\
22000	845.15\\
24000	824.05\\
26000	821.51\\
28000	806.86\\
30000	823.53\\
32000	834.35\\
34000	812.55\\
36000	836.52\\
38000	832.02\\
40000	814.57\\
42000	819.74\\
44000	822.26\\
46000	804.05\\
48000	829.19\\
50000	817.49\\
54000	815.61\\
56000	816.52\\
58000	809.55\\
60000	825.07\\
62000	848.82\\
64000	833.5\\
66000	821.57\\
68000	834.25\\
70000	829.51\\
72000	816.92\\
74000	836.76\\
76000	814.12\\
78000	816.98\\
80000	825.52\\
82000	820.39\\
84000	835.47\\
86000	835.58\\
90000	821.77\\
92000	835.72\\
94000	828.62\\
98000	838.05\\
1e+05	857.17\\
};
\addplot [color=mycolor90, forget plot]
  table[row sep=crcr]{%
0	4395.3\\
2000	1417.1\\
4000	822.68\\
8000	805.95\\
12000	792.53\\
14000	814.38\\
16000	814.94\\
18000	820.25\\
20000	791.92\\
22000	800.58\\
24000	791.73\\
26000	811.81\\
28000	796.02\\
30000	800.44\\
32000	817.87\\
34000	786.8\\
36000	785.5\\
38000	807.72\\
40000	815.05\\
42000	785.6\\
44000	808.87\\
46000	796.49\\
48000	815.69\\
50000	795.68\\
52000	782\\
54000	782.65\\
56000	790.44\\
58000	784.1\\
60000	821.54\\
62000	819.72\\
64000	805.06\\
66000	795.73\\
68000	807.61\\
70000	791.56\\
74000	816.58\\
78000	807.86\\
80000	811.87\\
82000	802.1\\
84000	805.88\\
86000	763.65\\
88000	835.32\\
92000	786.48\\
94000	814.9\\
98000	780.8\\
1e+05	805.65\\
};
\addplot [color=mycolor91, forget plot]
  table[row sep=crcr]{%
0	4395.3\\
2000	1395.7\\
4000	707.14\\
6000	693.64\\
8000	688.15\\
10000	706.37\\
12000	680.73\\
14000	663.32\\
16000	704.98\\
18000	681.54\\
20000	715.69\\
22000	724.96\\
24000	689.75\\
26000	674.34\\
28000	706.81\\
30000	679.75\\
32000	704.89\\
34000	696.06\\
36000	724.52\\
38000	691.73\\
40000	663.61\\
42000	666.35\\
44000	704.22\\
46000	737.61\\
48000	690.14\\
50000	664.08\\
52000	700.42\\
54000	676.07\\
56000	663.47\\
58000	657.2\\
60000	698.2\\
64000	675.56\\
66000	684.8\\
68000	677.27\\
70000	680.92\\
72000	676.32\\
74000	682.02\\
76000	680.17\\
78000	704.1\\
80000	722.88\\
82000	703.59\\
84000	696.61\\
86000	675.26\\
88000	700.81\\
90000	672.45\\
94000	718.43\\
96000	679.75\\
98000	693.29\\
1e+05	701.81\\
};
\addplot [color=mycolor92, forget plot]
  table[row sep=crcr]{%
0	4395.3\\
2000	1583.6\\
4000	897.28\\
6000	887.73\\
8000	908.1\\
10000	856.77\\
12000	887.21\\
16000	859.92\\
18000	857.43\\
20000	867.68\\
22000	856.93\\
24000	850.92\\
26000	863.96\\
28000	859.27\\
30000	880.1\\
32000	885.69\\
34000	861.15\\
36000	848.62\\
40000	861.32\\
42000	859.86\\
44000	883.31\\
46000	885.45\\
48000	844.53\\
50000	861.53\\
52000	867.19\\
54000	887.24\\
56000	853.82\\
58000	883.3\\
60000	857.31\\
62000	869.78\\
64000	872.81\\
66000	872.38\\
68000	828.25\\
70000	864.93\\
72000	872.22\\
74000	866.21\\
76000	872.17\\
78000	842.49\\
80000	873.18\\
82000	854.26\\
84000	886.8\\
86000	894.48\\
88000	858.47\\
90000	835.96\\
94000	879.67\\
96000	862.47\\
98000	876.76\\
1e+05	844.81\\
};
\addplot [color=mycolor93, forget plot]
  table[row sep=crcr]{%
0	4395.3\\
2000	1486.3\\
4000	754.95\\
6000	722.73\\
8000	732.68\\
10000	757.22\\
12000	747.97\\
14000	761.42\\
16000	755.58\\
18000	743.99\\
20000	709.72\\
22000	713.51\\
24000	743.74\\
28000	721.78\\
30000	696.9\\
32000	710.76\\
34000	746.88\\
36000	737.73\\
38000	696.53\\
40000	722.52\\
42000	734.5\\
44000	707.36\\
46000	734.56\\
48000	769.24\\
50000	711.65\\
52000	716.94\\
54000	705.64\\
56000	734.35\\
60000	714.39\\
62000	752.44\\
64000	720.62\\
66000	714.43\\
68000	718.15\\
70000	703.53\\
72000	703.08\\
74000	726.54\\
76000	740.22\\
78000	721\\
80000	746.01\\
82000	695.7\\
84000	726.53\\
86000	735.25\\
88000	769.96\\
90000	701.74\\
92000	750.07\\
94000	681.98\\
96000	724.02\\
98000	759.89\\
1e+05	738.82\\
};
\addplot [color=mycolor94, forget plot]
  table[row sep=crcr]{%
0	4395.3\\
2000	1534.2\\
4000	711.37\\
8000	673.28\\
10000	661.13\\
12000	662.38\\
14000	673.33\\
16000	689.5\\
18000	684.3\\
20000	687.69\\
22000	687.49\\
24000	694.62\\
26000	681.81\\
28000	657.65\\
30000	694.28\\
32000	690.61\\
34000	699.82\\
36000	703\\
38000	677.4\\
40000	697.02\\
42000	667.92\\
44000	702.38\\
46000	696.08\\
48000	685.38\\
50000	693.63\\
52000	687.22\\
54000	693.89\\
58000	700.81\\
60000	677.23\\
62000	683.13\\
64000	695.92\\
66000	699.45\\
68000	695.42\\
70000	703.41\\
72000	681.87\\
74000	670.72\\
76000	683.88\\
78000	703.52\\
80000	696.35\\
82000	667.17\\
84000	681.39\\
86000	662.29\\
90000	688.95\\
92000	680.16\\
94000	685\\
96000	680.79\\
98000	680.28\\
1e+05	676.23\\
};
\addplot [color=mycolor95, forget plot]
  table[row sep=crcr]{%
0	4395.3\\
2000	1283.2\\
4000	438.22\\
6000	390.97\\
8000	405.09\\
10000	421.4\\
12000	401.84\\
14000	400.28\\
16000	417.79\\
18000	416.41\\
20000	406.91\\
22000	395.51\\
24000	414.64\\
26000	385.55\\
28000	401.97\\
30000	400\\
32000	410.79\\
34000	398.41\\
36000	405.15\\
38000	385.87\\
40000	403.42\\
42000	398.71\\
46000	409\\
48000	391.29\\
50000	403.05\\
52000	396.11\\
54000	399.25\\
56000	384.91\\
58000	383.67\\
60000	397.71\\
62000	410.16\\
64000	404.12\\
66000	412.64\\
68000	401.22\\
70000	398.71\\
72000	409.27\\
74000	398.79\\
76000	422.2\\
78000	423.25\\
80000	416.92\\
82000	404.64\\
84000	406.92\\
88000	398.82\\
90000	414.82\\
92000	395.23\\
94000	405.31\\
96000	402.06\\
98000	411.47\\
1e+05	398.56\\
};
\addplot [color=mycolor96, forget plot]
  table[row sep=crcr]{%
0	4395.3\\
2000	1383.7\\
4000	446.81\\
6000	447.96\\
8000	431.4\\
10000	458.48\\
12000	428.13\\
14000	406.77\\
16000	418.86\\
20000	440.24\\
22000	411.72\\
24000	398.37\\
26000	397.15\\
28000	412.1\\
30000	421.74\\
32000	422.23\\
34000	444.39\\
36000	436\\
38000	447.83\\
40000	444.7\\
42000	418\\
44000	444.77\\
46000	456.32\\
48000	415.18\\
50000	412.65\\
52000	444.47\\
54000	442.99\\
56000	463.57\\
58000	438.22\\
60000	405.36\\
62000	431.28\\
64000	430.34\\
66000	457.89\\
68000	445.66\\
70000	439.95\\
72000	425.4\\
74000	452.97\\
76000	444.43\\
78000	442.01\\
80000	428.13\\
82000	447.64\\
84000	426.82\\
86000	400.98\\
88000	400\\
90000	411.62\\
92000	414.12\\
94000	450.09\\
96000	443.6\\
98000	423.33\\
1e+05	449.56\\
};
\addplot [color=mycolor97, forget plot]
  table[row sep=crcr]{%
0	4395.3\\
2000	1269.8\\
4000	338.83\\
6000	314.02\\
8000	304.22\\
10000	305.48\\
12000	299.15\\
14000	291.06\\
16000	331.07\\
18000	293.48\\
20000	286.28\\
22000	313.72\\
24000	307.08\\
26000	296.64\\
28000	278.8\\
30000	297.75\\
32000	305.83\\
34000	312.28\\
36000	291.68\\
38000	309.68\\
40000	291.26\\
42000	288.76\\
44000	323.01\\
46000	296.71\\
48000	304.89\\
50000	319.07\\
52000	300\\
54000	290.08\\
56000	289.21\\
58000	303.88\\
60000	288.7\\
62000	322.28\\
64000	324.2\\
66000	315.84\\
68000	315.22\\
70000	295.88\\
72000	300.43\\
74000	279.69\\
76000	310.52\\
78000	331.79\\
80000	317.79\\
82000	292.97\\
84000	274.81\\
86000	306.59\\
88000	326.85\\
90000	307.55\\
92000	325.89\\
94000	303.4\\
96000	315.75\\
98000	299.49\\
1e+05	300.1\\
};
\addplot [color=mycolor98]
  table[row sep=crcr]{%
0	4395.3\\
2000	1459\\
4000	458.7\\
6000	417.41\\
8000	405.78\\
10000	391.17\\
12000	405.56\\
14000	406.99\\
16000	409.62\\
18000	389.94\\
20000	419.22\\
22000	378.92\\
24000	402.7\\
26000	424.74\\
28000	407.89\\
30000	414.77\\
32000	429.8\\
34000	416.48\\
36000	396.95\\
38000	414.05\\
40000	418.58\\
42000	417.27\\
44000	397.9\\
46000	399.1\\
48000	398.82\\
50000	408.99\\
52000	400.11\\
56000	416.74\\
58000	405.38\\
60000	417.42\\
62000	416.22\\
64000	407.78\\
66000	395.92\\
68000	412.47\\
70000	415.41\\
72000	398.43\\
74000	440.98\\
76000	436.86\\
78000	416.38\\
80000	398.65\\
82000	430.69\\
84000	390.69\\
86000	412.24\\
88000	409.79\\
90000	391.91\\
92000	401.87\\
94000	416.99\\
96000	421.58\\
1e+05	409.58\\
};
\addlegendentry{\ClasskIVLegend}

      % \input{../../../../.tex/20/KS/SizeEvolutionsfig8.tex}
      \legend{}
    \end{groupplot}

    \begin{groupplot}[
      group style={
        group name=Row5,
        plotLastGroup,
        group size=2 by 1,
      },
      plotRow,
      row5Keys,
      plotLegendSE,
      %
      xmin=0,xmax=1e+05,
      ymin=100,ymax=1.6107e+05,
    ]
      \nextgroupplot[%
        at={($(Row4 c1r1.south west)-(0,\yPlotSep em)-(0,\plotHeight cm)$)},
        % xmin=0,xmax=1e+05,
        % ymin=100,ymax=1.6107e+05,
      ] % This file was created by matlab2tikz.
%
\definecolor{mycolor1}{rgb}{0.00146,0.00047,0.01387}%
\definecolor{mycolor2}{rgb}{0.01166,0.00942,0.06346}%
\definecolor{mycolor3}{rgb}{0.02943,0.02150,0.11462}%
\definecolor{mycolor4}{rgb}{0.06134,0.03659,0.17764}%
\definecolor{mycolor5}{rgb}{0.08741,0.04456,0.22481}%
\definecolor{mycolor6}{rgb}{0.12291,0.04754,0.28162}%
\definecolor{mycolor7}{rgb}{0.14907,0.04547,0.31709}%
\definecolor{mycolor8}{rgb}{0.18343,0.04033,0.35497}%
\definecolor{mycolor9}{rgb}{0.21795,0.03662,0.38352}%
\definecolor{mycolor10}{rgb}{0.24497,0.03705,0.40001}%
\definecolor{mycolor11}{rgb}{0.27135,0.04092,0.41198}%
\definecolor{mycolor12}{rgb}{0.29076,0.04564,0.41864}%
\definecolor{mycolor13}{rgb}{0.31628,0.05349,0.42512}%
\definecolor{mycolor14}{rgb}{0.33522,0.06006,0.42852}%
\definecolor{mycolor15}{rgb}{0.36028,0.06925,0.43150}%
\definecolor{mycolor16}{rgb}{0.37900,0.07625,0.43272}%
\definecolor{mycolor17}{rgb}{0.39767,0.08326,0.43318}%
\definecolor{mycolor18}{rgb}{0.41633,0.09020,0.43294}%
\definecolor{mycolor19}{rgb}{0.42877,0.09479,0.43241}%
\definecolor{mycolor20}{rgb}{0.44743,0.10160,0.43108}%
\definecolor{mycolor21}{rgb}{0.46610,0.10832,0.42912}%
\definecolor{mycolor22}{rgb}{0.47856,0.11276,0.42747}%
\definecolor{mycolor23}{rgb}{0.49726,0.11938,0.42449}%
\definecolor{mycolor24}{rgb}{0.50973,0.12377,0.42216}%
\definecolor{mycolor25}{rgb}{0.52221,0.12815,0.41955}%
\definecolor{mycolor26}{rgb}{0.53468,0.13253,0.41667}%
\definecolor{mycolor27}{rgb}{0.55339,0.13913,0.41183}%
\definecolor{mycolor28}{rgb}{0.56585,0.14357,0.40826}%
\definecolor{mycolor29}{rgb}{0.57830,0.14804,0.40441}%
\definecolor{mycolor30}{rgb}{0.59073,0.15256,0.40029}%
\definecolor{mycolor31}{rgb}{0.60314,0.15715,0.39589}%
\definecolor{mycolor32}{rgb}{0.61551,0.16182,0.39122}%
\definecolor{mycolor33}{rgb}{0.62169,0.16418,0.38878}%
\definecolor{mycolor34}{rgb}{0.63400,0.16899,0.38370}%
\definecolor{mycolor35}{rgb}{0.64626,0.17391,0.37836}%
\definecolor{mycolor36}{rgb}{0.65846,0.17896,0.37275}%
\definecolor{mycolor37}{rgb}{0.66454,0.18154,0.36985}%
\definecolor{mycolor38}{rgb}{0.67664,0.18681,0.36385}%
\definecolor{mycolor39}{rgb}{0.68865,0.19224,0.35760}%
\definecolor{mycolor40}{rgb}{0.69463,0.19502,0.35439}%
\definecolor{mycolor41}{rgb}{0.70650,0.20073,0.34778}%
\definecolor{mycolor42}{rgb}{0.71240,0.20366,0.34438}%
\definecolor{mycolor43}{rgb}{0.72410,0.20967,0.33742}%
\definecolor{mycolor44}{rgb}{0.72991,0.21276,0.33386}%
\definecolor{mycolor45}{rgb}{0.74142,0.21911,0.32658}%
\definecolor{mycolor46}{rgb}{0.74713,0.22238,0.32286}%
\definecolor{mycolor47}{rgb}{0.75279,0.22571,0.31909}%
\definecolor{mycolor48}{rgb}{0.76401,0.23255,0.31140}%
\definecolor{mycolor49}{rgb}{0.76956,0.23608,0.30749}%
\definecolor{mycolor50}{rgb}{0.77506,0.23967,0.30353}%
\definecolor{mycolor51}{rgb}{0.79129,0.25086,0.29139}%
\definecolor{mycolor52}{rgb}{0.79661,0.25473,0.28726}%
\definecolor{mycolor53}{rgb}{0.80187,0.25867,0.28310}%
\definecolor{mycolor54}{rgb}{0.81224,0.26679,0.27466}%
\definecolor{mycolor55}{rgb}{0.81734,0.27095,0.27039}%
\definecolor{mycolor56}{rgb}{0.82737,0.27952,0.26175}%
\definecolor{mycolor57}{rgb}{0.83230,0.28391,0.25738}%
\definecolor{mycolor58}{rgb}{0.83717,0.28839,0.25299}%
\definecolor{mycolor59}{rgb}{0.84671,0.29756,0.24411}%
\definecolor{mycolor60}{rgb}{0.85138,0.30226,0.23964}%
\definecolor{mycolor61}{rgb}{0.85599,0.30704,0.23513}%
\definecolor{mycolor62}{rgb}{0.86053,0.31189,0.23061}%
\definecolor{mycolor63}{rgb}{0.87374,0.32691,0.21689}%
\definecolor{mycolor64}{rgb}{0.87800,0.33206,0.21227}%
\definecolor{mycolor65}{rgb}{0.89034,0.34796,0.19829}%
\definecolor{mycolor66}{rgb}{0.89819,0.35891,0.18886}%
\definecolor{mycolor67}{rgb}{0.90200,0.36449,0.18412}%
\definecolor{mycolor68}{rgb}{0.90573,0.37014,0.17935}%
\definecolor{mycolor69}{rgb}{0.90939,0.37586,0.17456}%
\definecolor{mycolor70}{rgb}{0.91297,0.38164,0.16975}%
\definecolor{mycolor71}{rgb}{0.91646,0.38748,0.16492}%
\definecolor{mycolor72}{rgb}{0.91988,0.39339,0.16007}%
\definecolor{mycolor73}{rgb}{0.92322,0.39936,0.15519}%
\definecolor{mycolor74}{rgb}{0.92647,0.40539,0.15029}%
\definecolor{mycolor75}{rgb}{0.92964,0.41148,0.14537}%
\definecolor{mycolor76}{rgb}{0.93274,0.41763,0.14042}%
\definecolor{mycolor77}{rgb}{0.93575,0.42383,0.13544}%
\definecolor{mycolor78}{rgb}{0.93868,0.43009,0.13044}%
\definecolor{mycolor79}{rgb}{0.94152,0.43640,0.12541}%
\definecolor{mycolor80}{rgb}{0.94429,0.44277,0.12035}%
\definecolor{mycolor81}{rgb}{0.94956,0.45566,0.11016}%
\definecolor{mycolor82}{rgb}{0.95208,0.46218,0.10503}%
\definecolor{mycolor83}{rgb}{0.96129,0.48872,0.08429}%
\definecolor{mycolor84}{rgb}{0.96540,0.50225,0.07386}%
\definecolor{mycolor85}{rgb}{0.96732,0.50908,0.06866}%
\definecolor{mycolor86}{rgb}{0.97259,0.52980,0.05332}%
\definecolor{mycolor87}{rgb}{0.97568,0.54380,0.04362}%
\definecolor{mycolor88}{rgb}{0.98082,0.57221,0.02851}%
\definecolor{mycolor89}{rgb}{0.98189,0.57939,0.02625}%
\definecolor{mycolor90}{rgb}{0.98378,0.59385,0.02377}%
\definecolor{mycolor91}{rgb}{0.98532,0.60842,0.02420}%
\definecolor{mycolor92}{rgb}{0.98696,0.63048,0.03091}%
\definecolor{mycolor93}{rgb}{0.98793,0.66025,0.05175}%
\definecolor{mycolor94}{rgb}{0.98771,0.68281,0.07249}%
\definecolor{mycolor95}{rgb}{0.97750,0.78226,0.18592}%
\definecolor{mycolor96}{rgb}{0.96624,0.83619,0.26153}%
\definecolor{mycolor97}{rgb}{0.96252,0.85148,0.28555}%
\definecolor{mycolor98}{rgb}{0.98836,0.99836,0.64492}%
%
\addplot [color=mycolor1]
  table[row sep=crcr]{%
0	4395.3\\
2000	1.2078e+05\\
4000	1.5161e+05\\
8000	1.5984e+05\\
10000	1.543e+05\\
12000	1.5569e+05\\
14000	1.5392e+05\\
16000	1.5558e+05\\
20000	1.5363e+05\\
22000	1.5877e+05\\
26000	1.5247e+05\\
28000	1.5379e+05\\
30000	1.5411e+05\\
32000	1.5501e+05\\
34000	1.5384e+05\\
36000	1.5615e+05\\
40000	1.5417e+05\\
42000	1.5435e+05\\
44000	1.5527e+05\\
48000	1.5312e+05\\
50000	1.567e+05\\
52000	1.5292e+05\\
54000	1.5447e+05\\
56000	1.5788e+05\\
58000	1.5628e+05\\
62000	1.5444e+05\\
64000	1.5021e+05\\
68000	1.5701e+05\\
70000	1.5877e+05\\
72000	1.5369e+05\\
74000	1.5625e+05\\
76000	1.538e+05\\
78000	1.5345e+05\\
80000	1.5489e+05\\
84000	1.5446e+05\\
88000	1.5657e+05\\
90000	1.5513e+05\\
92000	1.5515e+05\\
94000	1.5732e+05\\
96000	1.5684e+05\\
98000	1.5248e+05\\
1e+05	1.6107e+05\\
};
\addlegendentry{\ClasskILegend}

\addplot [color=mycolor2, forget plot]
  table[row sep=crcr]{%
0	4395.3\\
2000	25895\\
4000	26367\\
6000	25508\\
8000	26497\\
10000	26869\\
12000	26125\\
14000	26367\\
16000	27041\\
18000	25825\\
20000	27057\\
22000	26078\\
24000	26675\\
26000	26832\\
28000	25802\\
30000	25863\\
32000	25709\\
34000	26274\\
36000	26308\\
38000	25882\\
40000	26627\\
44000	26370\\
46000	26550\\
48000	25967\\
50000	25871\\
52000	27278\\
54000	26824\\
56000	26237\\
58000	26634\\
60000	26730\\
62000	26497\\
64000	26540\\
66000	26110\\
68000	26452\\
70000	26227\\
74000	27000\\
76000	27184\\
78000	26039\\
80000	26818\\
82000	25873\\
84000	26400\\
88000	26041\\
90000	26510\\
92000	26537\\
94000	26097\\
96000	26205\\
98000	25762\\
1e+05	26370\\
};
\addplot [color=mycolor3, forget plot]
  table[row sep=crcr]{%
0	4395.3\\
2000	37117\\
4000	39886\\
6000	39722\\
8000	40134\\
10000	39410\\
12000	40770\\
14000	39716\\
16000	40909\\
18000	41689\\
20000	40816\\
22000	41626\\
24000	41431\\
26000	40531\\
28000	39404\\
30000	40485\\
32000	40557\\
34000	40367\\
36000	41184\\
38000	39689\\
40000	39438\\
42000	40052\\
44000	40171\\
46000	40855\\
48000	41084\\
50000	39921\\
52000	40605\\
54000	40263\\
56000	39249\\
58000	40297\\
60000	39701\\
62000	40517\\
64000	41108\\
66000	40297\\
68000	41385\\
70000	40727\\
72000	40661\\
74000	39421\\
76000	39948\\
80000	39823\\
82000	39076\\
84000	39988\\
88000	40584\\
90000	41017\\
92000	39515\\
94000	40600\\
98000	40699\\
1e+05	39804\\
};
\addplot [color=mycolor4, forget plot]
  table[row sep=crcr]{%
0	4395.3\\
2000	15391\\
4000	14478\\
6000	14432\\
8000	14662\\
12000	14187\\
14000	14732\\
16000	14575\\
18000	14341\\
20000	14874\\
22000	14323\\
24000	14233\\
26000	14871\\
28000	14313\\
30000	14533\\
32000	14696\\
34000	14323\\
36000	14789\\
40000	14344\\
42000	14423\\
44000	14628\\
46000	14379\\
48000	14490\\
50000	14546\\
52000	14462\\
54000	14494\\
58000	14855\\
60000	14565\\
64000	14730\\
66000	14388\\
68000	14369\\
70000	14639\\
72000	14191\\
74000	14551\\
76000	14228\\
78000	14337\\
80000	14182\\
82000	14634\\
84000	14151\\
86000	14253\\
88000	14466\\
90000	14778\\
92000	14428\\
94000	14371\\
96000	13965\\
98000	14673\\
1e+05	13994\\
};
\addplot [color=mycolor5, forget plot]
  table[row sep=crcr]{%
0	4395.3\\
2000	55786\\
4000	67056\\
6000	65878\\
8000	63845\\
10000	70136\\
12000	65075\\
16000	66694\\
18000	64640\\
20000	63709\\
22000	63826\\
24000	66009\\
26000	64765\\
28000	64564\\
30000	66008\\
32000	63033\\
34000	63172\\
36000	63540\\
38000	65408\\
40000	65794\\
42000	67233\\
44000	66216\\
46000	64465\\
48000	64186\\
50000	66814\\
52000	65773\\
54000	65216\\
56000	64329\\
58000	64484\\
60000	65341\\
62000	65897\\
64000	62756\\
66000	63278\\
68000	64559\\
70000	64408\\
74000	63036\\
76000	64287\\
80000	65875\\
82000	63164\\
84000	65429\\
86000	65462\\
88000	64779\\
90000	65684\\
92000	64803\\
94000	61562\\
96000	63232\\
98000	64344\\
1e+05	64606\\
};
\addplot [color=mycolor6, forget plot]
  table[row sep=crcr]{%
0	4395.3\\
2000	48092\\
4000	56897\\
8000	55040\\
10000	58671\\
12000	57272\\
14000	57382\\
16000	57041\\
18000	57243\\
20000	57856\\
22000	57516\\
24000	54690\\
26000	56816\\
28000	58141\\
30000	55410\\
32000	57068\\
34000	56828\\
36000	55885\\
38000	57968\\
40000	56292\\
44000	57264\\
48000	55425\\
50000	56528\\
52000	58225\\
54000	60178\\
56000	56223\\
58000	59638\\
60000	55330\\
62000	56301\\
64000	58493\\
66000	57134\\
68000	56540\\
70000	56640\\
72000	58363\\
74000	57045\\
76000	57211\\
78000	57582\\
80000	57501\\
82000	57790\\
84000	56238\\
86000	55695\\
88000	56428\\
90000	56491\\
92000	55677\\
94000	56820\\
96000	55881\\
98000	57637\\
1e+05	55528\\
};
\addplot [color=mycolor7, forget plot]
  table[row sep=crcr]{%
0	4395.3\\
2000	24599\\
4000	28766\\
6000	27997\\
8000	28127\\
10000	29072\\
14000	29931\\
16000	28515\\
18000	28460\\
20000	28493\\
24000	28275\\
26000	28641\\
28000	28494\\
30000	29258\\
32000	28279\\
34000	29254\\
36000	28360\\
38000	28567\\
40000	28330\\
42000	28536\\
44000	29714\\
46000	27977\\
48000	28916\\
52000	27904\\
56000	28389\\
58000	27559\\
62000	29591\\
64000	28970\\
66000	29072\\
68000	28791\\
70000	28746\\
72000	28834\\
74000	28234\\
76000	27737\\
78000	28594\\
80000	28590\\
82000	29482\\
84000	29138\\
86000	28304\\
88000	27724\\
90000	27737\\
92000	29506\\
94000	28617\\
98000	28215\\
1e+05	28663\\
};
\addplot [color=mycolor8, forget plot]
  table[row sep=crcr]{%
0	4395.3\\
2000	20749\\
4000	21326\\
6000	21110\\
8000	20731\\
10000	20724\\
12000	20636\\
14000	20487\\
20000	20831\\
22000	20166\\
24000	20091\\
26000	21002\\
28000	20889\\
30000	21401\\
32000	20834\\
36000	20271\\
38000	20633\\
40000	20794\\
42000	20821\\
44000	20404\\
46000	20344\\
48000	20942\\
50000	21750\\
52000	20735\\
54000	20323\\
56000	20888\\
58000	21095\\
60000	21216\\
62000	20841\\
64000	20778\\
68000	20451\\
70000	21065\\
72000	20335\\
74000	20970\\
76000	20907\\
78000	20716\\
80000	19870\\
82000	20852\\
84000	20768\\
86000	20831\\
88000	20680\\
90000	20592\\
92000	20896\\
94000	21042\\
96000	20776\\
98000	21045\\
1e+05	20198\\
};
\addplot [color=mycolor9, forget plot]
  table[row sep=crcr]{%
0	4395.3\\
2000	13061\\
4000	13539\\
6000	14381\\
8000	14210\\
10000	14484\\
12000	13852\\
14000	13998\\
16000	14314\\
18000	14025\\
20000	14286\\
22000	14379\\
24000	14065\\
26000	14161\\
28000	14307\\
30000	14499\\
32000	14289\\
34000	14145\\
36000	14208\\
38000	13779\\
40000	14370\\
42000	14179\\
44000	14077\\
46000	14207\\
48000	13759\\
50000	13723\\
52000	14270\\
54000	14745\\
56000	14378\\
58000	14325\\
60000	14434\\
62000	14165\\
64000	14490\\
66000	14380\\
68000	14081\\
72000	14523\\
74000	14293\\
76000	14766\\
78000	14587\\
80000	14257\\
82000	14235\\
84000	14057\\
86000	14530\\
88000	14400\\
90000	14202\\
92000	14202\\
94000	14084\\
96000	14310\\
98000	14392\\
1e+05	14153\\
};
\addplot [color=mycolor10, forget plot]
  table[row sep=crcr]{%
0	4395.3\\
2000	19833\\
4000	22494\\
6000	22298\\
8000	22191\\
10000	21792\\
12000	22499\\
14000	22222\\
16000	22212\\
18000	21801\\
20000	22169\\
22000	22410\\
26000	22399\\
28000	22129\\
30000	22232\\
32000	21975\\
34000	22925\\
36000	22807\\
38000	21246\\
40000	22358\\
42000	22988\\
44000	22335\\
46000	22343\\
48000	22714\\
50000	22299\\
52000	22683\\
54000	22014\\
56000	22255\\
58000	21972\\
60000	21587\\
62000	21851\\
64000	22666\\
66000	22777\\
68000	22788\\
72000	22245\\
76000	22001\\
78000	22413\\
80000	22152\\
82000	22390\\
84000	22786\\
86000	22327\\
88000	22539\\
92000	22051\\
94000	22925\\
96000	22723\\
98000	22245\\
1e+05	22498\\
};
\addplot [color=mycolor11, forget plot]
  table[row sep=crcr]{%
0	4395.3\\
2000	9493\\
4000	8998.1\\
6000	8449.2\\
8000	8713.2\\
10000	8620.7\\
12000	8782.6\\
14000	8462.7\\
16000	8856\\
18000	8818.2\\
20000	8601.3\\
24000	8857.8\\
26000	8848.5\\
30000	8607.2\\
32000	8725.7\\
34000	8816.1\\
36000	9027.4\\
38000	8366.3\\
40000	8617.7\\
44000	8727.9\\
46000	8650\\
48000	8710\\
50000	8538.6\\
52000	8479.5\\
56000	8691.8\\
60000	8555.1\\
62000	8742.8\\
64000	8692.5\\
66000	8983.8\\
68000	8734.5\\
70000	8894.4\\
72000	8587\\
74000	8737.5\\
76000	8933.8\\
78000	8661.1\\
80000	8372.2\\
84000	8878.6\\
86000	8637\\
88000	8683.2\\
92000	8552.5\\
94000	8579.1\\
96000	8486.5\\
1e+05	8960.2\\
};
\addplot [color=mycolor12, forget plot]
  table[row sep=crcr]{%
0	4395.3\\
2000	13464\\
4000	14944\\
6000	15406\\
8000	15175\\
10000	15264\\
12000	14848\\
14000	15263\\
16000	14614\\
18000	14771\\
20000	14421\\
22000	15230\\
24000	14675\\
26000	15214\\
28000	15490\\
30000	14820\\
32000	14746\\
34000	14632\\
36000	15187\\
38000	14256\\
40000	15280\\
42000	14631\\
44000	14715\\
46000	15220\\
48000	15079\\
50000	14771\\
52000	15068\\
54000	14794\\
56000	14652\\
58000	14817\\
60000	14784\\
62000	14974\\
64000	14605\\
66000	15040\\
68000	14910\\
70000	15120\\
72000	14936\\
76000	14672\\
78000	15177\\
80000	14810\\
82000	14890\\
84000	14803\\
86000	15152\\
88000	14656\\
90000	14999\\
92000	14839\\
94000	14823\\
96000	14998\\
98000	15317\\
1e+05	14709\\
};
\addplot [color=mycolor13, forget plot]
  table[row sep=crcr]{%
0	4395.3\\
2000	10628\\
6000	11270\\
8000	10769\\
12000	11074\\
14000	11048\\
16000	10665\\
18000	11414\\
20000	11405\\
22000	11488\\
24000	11254\\
26000	11941\\
28000	11409\\
30000	11028\\
32000	11144\\
34000	11087\\
36000	11602\\
38000	11066\\
40000	10793\\
42000	11442\\
44000	11441\\
46000	11087\\
48000	10826\\
50000	11352\\
52000	10973\\
54000	11247\\
56000	10855\\
58000	10718\\
60000	11290\\
62000	11401\\
64000	11404\\
66000	11824\\
68000	11132\\
70000	11259\\
74000	11206\\
76000	11663\\
78000	10993\\
82000	11299\\
84000	10771\\
86000	11560\\
88000	11236\\
90000	10986\\
92000	11159\\
94000	11519\\
96000	11149\\
98000	11690\\
1e+05	11260\\
};
\addplot [color=mycolor14, forget plot]
  table[row sep=crcr]{%
0	4395.3\\
2000	22563\\
4000	24733\\
6000	25319\\
8000	24898\\
10000	23707\\
14000	24847\\
16000	24302\\
20000	25960\\
22000	24993\\
24000	24259\\
26000	25550\\
28000	24308\\
30000	24975\\
32000	24341\\
34000	25147\\
36000	24482\\
38000	24477\\
40000	24543\\
42000	24336\\
44000	23868\\
48000	25263\\
50000	24536\\
52000	24203\\
54000	25336\\
56000	24587\\
58000	24049\\
60000	23851\\
62000	24253\\
64000	25391\\
66000	25524\\
68000	25014\\
70000	23666\\
72000	24548\\
74000	25058\\
76000	24588\\
78000	25041\\
80000	25317\\
82000	24868\\
84000	24714\\
86000	25409\\
88000	24560\\
92000	25400\\
94000	24587\\
96000	25494\\
98000	25677\\
1e+05	24078\\
};
\addplot [color=mycolor15, forget plot]
  table[row sep=crcr]{%
0	4395.3\\
2000	15140\\
4000	16937\\
8000	17179\\
10000	16852\\
12000	17159\\
16000	16828\\
18000	17510\\
20000	17033\\
22000	17438\\
24000	17125\\
26000	17075\\
28000	17282\\
30000	17051\\
32000	16530\\
34000	17292\\
36000	17620\\
38000	16939\\
40000	17174\\
42000	17043\\
44000	17700\\
46000	17469\\
48000	17142\\
50000	17518\\
52000	17229\\
54000	17360\\
56000	16375\\
58000	17125\\
60000	17431\\
62000	17396\\
64000	17815\\
66000	17311\\
68000	17575\\
70000	17036\\
72000	17369\\
74000	16945\\
76000	17518\\
78000	17416\\
80000	17101\\
82000	17480\\
84000	16792\\
86000	17363\\
88000	17662\\
90000	17783\\
92000	16972\\
94000	17195\\
96000	17597\\
98000	17231\\
1e+05	17214\\
};
\addplot [color=mycolor16, forget plot]
  table[row sep=crcr]{%
0	4395.3\\
2000	7521.4\\
4000	6961.6\\
6000	7131.9\\
8000	7009.7\\
10000	7029.8\\
12000	7027.1\\
14000	7073.2\\
16000	6806.5\\
18000	6965.7\\
20000	7196.3\\
22000	6950.1\\
24000	7093.6\\
26000	6984.2\\
28000	7079.4\\
30000	7295.1\\
34000	6967.8\\
36000	7170.3\\
38000	7133.3\\
40000	7020.7\\
44000	7054.9\\
46000	7075.5\\
48000	7128.1\\
50000	7101.4\\
52000	6942.8\\
54000	7236.4\\
56000	6711.2\\
58000	7115.9\\
60000	7014.3\\
62000	6968.1\\
64000	7194.9\\
66000	7086\\
68000	7087.9\\
70000	6837.6\\
72000	7455.1\\
74000	7123.8\\
76000	7175.9\\
78000	7152.9\\
80000	6813.5\\
82000	7005\\
84000	7012.5\\
86000	7248.3\\
88000	7460.3\\
90000	7229.2\\
92000	7236\\
94000	6947.7\\
96000	7051.4\\
98000	6917.2\\
1e+05	7172\\
};
\addplot [color=mycolor17, forget plot]
  table[row sep=crcr]{%
0	4395.3\\
2000	14146\\
4000	15699\\
6000	16397\\
8000	15723\\
10000	16308\\
12000	16418\\
14000	16589\\
16000	16062\\
18000	16982\\
20000	16722\\
22000	16203\\
24000	15748\\
26000	16427\\
28000	16158\\
30000	16464\\
32000	16360\\
34000	15713\\
36000	16024\\
38000	16605\\
40000	16149\\
42000	15922\\
44000	15294\\
46000	16792\\
48000	16592\\
50000	15978\\
52000	16036\\
54000	16436\\
56000	16060\\
58000	16549\\
60000	16040\\
62000	15898\\
64000	16165\\
66000	16201\\
68000	16482\\
70000	16467\\
72000	16376\\
74000	16537\\
76000	16891\\
78000	16814\\
80000	16523\\
82000	16329\\
86000	16057\\
88000	16050\\
90000	15983\\
92000	16557\\
94000	16274\\
96000	16316\\
98000	16688\\
1e+05	15969\\
};
\addplot [color=mycolor18, forget plot]
  table[row sep=crcr]{%
0	4395.3\\
2000	21824\\
4000	23777\\
6000	24590\\
8000	24773\\
10000	25291\\
12000	24953\\
14000	24880\\
16000	24588\\
18000	25670\\
20000	25028\\
22000	23918\\
24000	24685\\
26000	24210\\
28000	24785\\
30000	24570\\
32000	26505\\
34000	25937\\
36000	23811\\
38000	24746\\
40000	24642\\
42000	23401\\
44000	24013\\
46000	25056\\
48000	24637\\
50000	24811\\
52000	24428\\
54000	24622\\
56000	25502\\
58000	24515\\
60000	24721\\
62000	24806\\
64000	25047\\
66000	24513\\
68000	25366\\
70000	25534\\
72000	25166\\
74000	24701\\
76000	24994\\
78000	25166\\
80000	25484\\
82000	25264\\
84000	24808\\
86000	24750\\
88000	23940\\
90000	24628\\
94000	25372\\
96000	25287\\
98000	25513\\
1e+05	24469\\
};
\addplot [color=mycolor19, forget plot]
  table[row sep=crcr]{%
0	4395.3\\
2000	14046\\
4000	14658\\
6000	15190\\
8000	15075\\
10000	15141\\
12000	14925\\
14000	14936\\
16000	15102\\
18000	14995\\
20000	15213\\
22000	14837\\
24000	15067\\
26000	14659\\
28000	14932\\
30000	14946\\
32000	14797\\
34000	15241\\
38000	15365\\
40000	14811\\
42000	15484\\
44000	14967\\
48000	15351\\
50000	14985\\
52000	14512\\
54000	14522\\
56000	14225\\
58000	15042\\
60000	14949\\
62000	15386\\
64000	14913\\
66000	14603\\
68000	15181\\
70000	15442\\
72000	14692\\
74000	15314\\
76000	15396\\
78000	14995\\
80000	14850\\
82000	14935\\
84000	14930\\
86000	14770\\
88000	15405\\
90000	14793\\
92000	15202\\
94000	15545\\
96000	15176\\
1e+05	15020\\
};
\addplot [color=mycolor20, forget plot]
  table[row sep=crcr]{%
0	4395.3\\
2000	17048\\
4000	19447\\
6000	19625\\
10000	19564\\
12000	19458\\
14000	20069\\
16000	19413\\
18000	19714\\
22000	19258\\
24000	19233\\
26000	20093\\
28000	19685\\
30000	18885\\
34000	19658\\
36000	19279\\
38000	19782\\
40000	19172\\
42000	19199\\
44000	19412\\
46000	19068\\
48000	20049\\
50000	18582\\
52000	19823\\
54000	19355\\
56000	19775\\
58000	19234\\
60000	19957\\
62000	19215\\
64000	19442\\
66000	18999\\
68000	18991\\
70000	19553\\
74000	19548\\
76000	19903\\
78000	19313\\
80000	19468\\
82000	20019\\
84000	19244\\
86000	20141\\
88000	19524\\
90000	19740\\
92000	19196\\
94000	19193\\
96000	19259\\
98000	19583\\
1e+05	19065\\
};
\addplot [color=mycolor21, forget plot]
  table[row sep=crcr]{%
0	4395.3\\
2000	8940.2\\
4000	9580.8\\
6000	9758.6\\
10000	9738\\
12000	9871.7\\
14000	9729.8\\
18000	9528.2\\
20000	9679.5\\
22000	10028\\
28000	9682.3\\
30000	9992.8\\
32000	9868.1\\
34000	9957.2\\
36000	9735\\
38000	9773.8\\
40000	9658.8\\
42000	9812.4\\
44000	9810.4\\
46000	9770\\
48000	9924.1\\
50000	9634.8\\
52000	9899.6\\
54000	9583.1\\
56000	9692.1\\
58000	9633.5\\
60000	9862.7\\
62000	9951.9\\
64000	9724.3\\
66000	9733.2\\
68000	9542.4\\
70000	9505.4\\
72000	9770.4\\
74000	9778.9\\
76000	9867.8\\
78000	9780.5\\
80000	9777.1\\
82000	9945.2\\
84000	9530.8\\
86000	9682.6\\
88000	9754.3\\
90000	9578.9\\
94000	9858.1\\
96000	9557.1\\
98000	9888.6\\
1e+05	9831.4\\
};
\addplot [color=mycolor22, forget plot]
  table[row sep=crcr]{%
0	4395.3\\
2000	24856\\
4000	26899\\
6000	25561\\
8000	25998\\
10000	24801\\
12000	24907\\
14000	26281\\
16000	25483\\
18000	25745\\
20000	24894\\
22000	24730\\
24000	26065\\
26000	25447\\
28000	26543\\
32000	25138\\
34000	24777\\
36000	25867\\
38000	27379\\
40000	25947\\
42000	25516\\
44000	26257\\
46000	26618\\
48000	25930\\
50000	25973\\
52000	26621\\
54000	26037\\
56000	25364\\
62000	26142\\
64000	25933\\
70000	25547\\
72000	26270\\
74000	24748\\
76000	25175\\
78000	25283\\
82000	25200\\
84000	27182\\
86000	25964\\
88000	25157\\
90000	25472\\
92000	25912\\
94000	25649\\
96000	24998\\
98000	25837\\
1e+05	26576\\
};
\addplot [color=mycolor23, forget plot]
  table[row sep=crcr]{%
0	4395.3\\
2000	12005\\
4000	12466\\
6000	12600\\
8000	12205\\
10000	12576\\
12000	12484\\
14000	12658\\
16000	12562\\
18000	12525\\
20000	12297\\
22000	12189\\
24000	12256\\
26000	12268\\
28000	12860\\
30000	12592\\
32000	12437\\
34000	12478\\
36000	12169\\
38000	12482\\
40000	12617\\
42000	12536\\
46000	12620\\
48000	12412\\
50000	12258\\
52000	12580\\
56000	12374\\
58000	12566\\
60000	12698\\
62000	12325\\
64000	12559\\
68000	12236\\
70000	12335\\
72000	12533\\
74000	12512\\
76000	12435\\
78000	12631\\
80000	12478\\
82000	12431\\
84000	12605\\
86000	12414\\
88000	12623\\
90000	12283\\
94000	12688\\
96000	12559\\
98000	12323\\
1e+05	12334\\
};
\addplot [color=mycolor24, forget plot]
  table[row sep=crcr]{%
0	4395.3\\
2000	5243.1\\
4000	4833.8\\
6000	4776.6\\
8000	4774.5\\
10000	4878.1\\
12000	4835.2\\
14000	4809.3\\
16000	4864.2\\
18000	4940.1\\
20000	4667.6\\
22000	4794.4\\
24000	4792.3\\
26000	4904.3\\
28000	4780\\
30000	4776.7\\
32000	4804.6\\
34000	4636.6\\
36000	4860.1\\
38000	4880.7\\
40000	4834.8\\
42000	4634.3\\
44000	4741.4\\
46000	5052.2\\
48000	4619.1\\
50000	4751.8\\
52000	4579\\
54000	4864.7\\
56000	4826\\
58000	4651.5\\
60000	4664.5\\
62000	4869.6\\
64000	4765.6\\
66000	4904.1\\
68000	4754.2\\
70000	4692.2\\
72000	4920.7\\
74000	4525.8\\
76000	4803\\
78000	4910.3\\
80000	4549.3\\
82000	4811.5\\
84000	4810\\
86000	4724.8\\
88000	4790.4\\
92000	4767\\
94000	4832\\
96000	4834.1\\
98000	4739.9\\
1e+05	4760.1\\
};
\addplot [color=mycolor25, forget plot]
  table[row sep=crcr]{%
0	4395.3\\
2000	11974\\
4000	12173\\
6000	12965\\
8000	12628\\
10000	12968\\
12000	12559\\
14000	12759\\
16000	12707\\
18000	12484\\
20000	12823\\
22000	12572\\
24000	12102\\
26000	12620\\
28000	12746\\
30000	12719\\
32000	12951\\
34000	12273\\
36000	12562\\
38000	12675\\
40000	12549\\
42000	12342\\
44000	12082\\
46000	12817\\
48000	12276\\
50000	12429\\
52000	12098\\
54000	12377\\
56000	12754\\
58000	12963\\
60000	12541\\
62000	11970\\
64000	12257\\
66000	12830\\
68000	12277\\
70000	12671\\
72000	12300\\
74000	12286\\
76000	12744\\
78000	12238\\
80000	12431\\
82000	12191\\
84000	12742\\
86000	12507\\
88000	11995\\
90000	12183\\
92000	12882\\
94000	12165\\
96000	12888\\
98000	12769\\
1e+05	12374\\
};
\addplot [color=mycolor26, forget plot]
  table[row sep=crcr]{%
0	4395.3\\
2000	7883.5\\
4000	8173.1\\
6000	8334.9\\
8000	8442.4\\
10000	8178.3\\
12000	8807.9\\
16000	7979.9\\
18000	8173.1\\
20000	8319\\
22000	8432.3\\
26000	8019\\
28000	8150.7\\
30000	8253.3\\
32000	8259.4\\
34000	8235.9\\
36000	8189.1\\
38000	8102.8\\
40000	8261.4\\
42000	8053.2\\
44000	8185.5\\
46000	8110.8\\
48000	8288.4\\
50000	8280.2\\
52000	8196.5\\
54000	8020.2\\
56000	8200.4\\
58000	8220.1\\
60000	8093\\
62000	8185.4\\
64000	8433.6\\
66000	8260.9\\
68000	8359.6\\
70000	8071.1\\
72000	8445.5\\
74000	8394.1\\
76000	8259.6\\
78000	8297.6\\
80000	8229.1\\
82000	8342.2\\
84000	8037\\
86000	8313.4\\
88000	8316.7\\
90000	8293\\
92000	8734.8\\
94000	8164.1\\
96000	8402.2\\
98000	8357.5\\
1e+05	8366.4\\
};
\addplot [color=mycolor27, forget plot]
  table[row sep=crcr]{%
0	4395.3\\
2000	10048\\
6000	10054\\
8000	10004\\
10000	9767.4\\
12000	10056\\
14000	9676.4\\
16000	9934.1\\
18000	9850.7\\
20000	10035\\
22000	9266.1\\
24000	9993.7\\
26000	10018\\
28000	9314.8\\
30000	9755.6\\
32000	9894.1\\
34000	9763.5\\
36000	9694.8\\
38000	9778.4\\
40000	10034\\
42000	9538.3\\
44000	9972.3\\
46000	9883.2\\
48000	9857.9\\
50000	9407.2\\
52000	9824.7\\
54000	9690.8\\
56000	9925.5\\
58000	9722.7\\
60000	10268\\
62000	9658.7\\
64000	9905.4\\
66000	9874.5\\
68000	9906.4\\
70000	9557.7\\
72000	9512.7\\
74000	10008\\
76000	9925.9\\
78000	9736.6\\
80000	9941\\
82000	10284\\
84000	9816.1\\
86000	9953.4\\
88000	9826.5\\
90000	9812.2\\
96000	9998\\
98000	9656.4\\
1e+05	9705.4\\
};
\addplot [color=mycolor28, forget plot]
  table[row sep=crcr]{%
0	4395.3\\
2000	5842\\
4000	5401.8\\
6000	5271.8\\
8000	5591.7\\
10000	5470.7\\
12000	5650.2\\
14000	5458.5\\
16000	5577\\
18000	5434.2\\
20000	5689.1\\
22000	5460.1\\
24000	5543.8\\
26000	5358.9\\
28000	5512\\
30000	5579.1\\
32000	5627.1\\
34000	5379.6\\
36000	5559.8\\
38000	5580.5\\
40000	5507.9\\
42000	5643.6\\
44000	5446.8\\
46000	5465.9\\
48000	5520.3\\
50000	5401.2\\
52000	5397.7\\
54000	5417.7\\
56000	5473.1\\
58000	5579.8\\
60000	5573\\
62000	5549.7\\
64000	5632.9\\
66000	5621\\
68000	5454.9\\
72000	5409\\
74000	5546.8\\
76000	5467.8\\
78000	5547.4\\
80000	5324.9\\
82000	5424.1\\
84000	5379.3\\
86000	5503.4\\
88000	5262.3\\
90000	5520.2\\
92000	5428\\
94000	5415\\
96000	5512.1\\
98000	5591.9\\
1e+05	5390\\
};
\addplot [color=mycolor29, forget plot]
  table[row sep=crcr]{%
0	4395.3\\
2000	11753\\
6000	11639\\
8000	11758\\
10000	11781\\
12000	11755\\
14000	11545\\
16000	11566\\
18000	11722\\
20000	11549\\
22000	11918\\
24000	12236\\
26000	12285\\
28000	11923\\
30000	12315\\
32000	12532\\
34000	11997\\
36000	11569\\
38000	12124\\
40000	12005\\
42000	11599\\
44000	11324\\
46000	11628\\
48000	11622\\
50000	12247\\
54000	11321\\
56000	11854\\
58000	11473\\
60000	11589\\
62000	11232\\
64000	11782\\
66000	11851\\
68000	11974\\
70000	11287\\
72000	11459\\
74000	12156\\
76000	12102\\
78000	11780\\
80000	11761\\
82000	11515\\
84000	11898\\
86000	11689\\
88000	12152\\
90000	11802\\
92000	12123\\
94000	12176\\
96000	11932\\
98000	12092\\
1e+05	11684\\
};
\addplot [color=mycolor30, forget plot]
  table[row sep=crcr]{%
0	4395.3\\
2000	7318.5\\
4000	7754.7\\
6000	7925.4\\
10000	7803\\
12000	7826.6\\
14000	7946.7\\
16000	7672.7\\
18000	7803\\
20000	7855.4\\
22000	8040.7\\
24000	7979.5\\
26000	8016.7\\
28000	7958.3\\
30000	7995.1\\
32000	7857.7\\
34000	8019.7\\
36000	7998.1\\
38000	7925.4\\
40000	7691.8\\
42000	8250.9\\
44000	7931.2\\
46000	7848\\
48000	8132.3\\
50000	7943.6\\
52000	8078.1\\
54000	7785.5\\
56000	7963.6\\
58000	7808.1\\
60000	7821\\
62000	7793.7\\
64000	7727.3\\
66000	7878.6\\
68000	7903.5\\
70000	7793.8\\
72000	8160.7\\
76000	8239.4\\
78000	7945.2\\
80000	8022.7\\
82000	7716.7\\
84000	7844.8\\
86000	7831.3\\
88000	7957.4\\
90000	8157.3\\
92000	8098.2\\
94000	7992.7\\
96000	7944.1\\
98000	7986.5\\
1e+05	7793.9\\
};
\addplot [color=mycolor31, forget plot]
  table[row sep=crcr]{%
0	4395.3\\
2000	6968.7\\
4000	6747.9\\
6000	6364.7\\
8000	6580.5\\
10000	6720.7\\
12000	6436.9\\
14000	6435.6\\
16000	6481.1\\
18000	6448.7\\
20000	6504.9\\
22000	6333.4\\
24000	6780.9\\
26000	6294.9\\
28000	6481\\
32000	6370.5\\
34000	6584.6\\
36000	6378.6\\
38000	6559\\
40000	6576.6\\
42000	6895\\
44000	6348.8\\
46000	6262.3\\
48000	6363.9\\
50000	6432\\
52000	6457.4\\
54000	6737.2\\
56000	6353.7\\
58000	6639\\
60000	6388.4\\
64000	6436.4\\
66000	6308.4\\
68000	6638.8\\
70000	6587.3\\
72000	6602.7\\
74000	6576.4\\
76000	6623\\
78000	6467.8\\
80000	6599.7\\
82000	6462.1\\
84000	6495.8\\
86000	6833.9\\
88000	6392.7\\
90000	6484.6\\
92000	6317.3\\
94000	6533.3\\
96000	6353.6\\
98000	6380.9\\
1e+05	6445.8\\
};
\addplot [color=mycolor32, forget plot]
  table[row sep=crcr]{%
0	4395.3\\
2000	10738\\
4000	10625\\
6000	10584\\
8000	10872\\
10000	10556\\
12000	10760\\
14000	10705\\
16000	10829\\
18000	10351\\
20000	10722\\
22000	11067\\
24000	10545\\
26000	10378\\
28000	10509\\
30000	10872\\
32000	10694\\
34000	10717\\
36000	10408\\
38000	10232\\
40000	10594\\
42000	10532\\
44000	11007\\
46000	10764\\
48000	10356\\
50000	10742\\
52000	10861\\
54000	10472\\
56000	10705\\
58000	10348\\
60000	10612\\
62000	10787\\
64000	10615\\
66000	10573\\
68000	10746\\
70000	9967.5\\
72000	10747\\
74000	10324\\
76000	10791\\
78000	10944\\
80000	10374\\
82000	10485\\
84000	10692\\
86000	10739\\
88000	10717\\
90000	10621\\
94000	10625\\
96000	10234\\
98000	10636\\
1e+05	10536\\
};
\addplot [color=mycolor33, forget plot]
  table[row sep=crcr]{%
0	4395.3\\
2000	7330.1\\
4000	8150.7\\
6000	8004.4\\
8000	8336\\
10000	7819.1\\
12000	8481.2\\
14000	8477\\
16000	7828.9\\
18000	8315.6\\
20000	8129.6\\
22000	8161.5\\
24000	8135.4\\
26000	8449.8\\
28000	8456.2\\
30000	8101\\
32000	8230.1\\
34000	8115.5\\
36000	8183\\
38000	8189.1\\
42000	8146.1\\
44000	8269.3\\
46000	8048.3\\
48000	8369.1\\
50000	7628.6\\
52000	8220.6\\
54000	8272.1\\
56000	8082.4\\
58000	8216.8\\
60000	8614.3\\
62000	8352.9\\
64000	8448.9\\
66000	8317.1\\
68000	8272.9\\
70000	7987.9\\
72000	8467.5\\
74000	8380.7\\
76000	8378.5\\
78000	8060.9\\
80000	8098.1\\
82000	8375.7\\
84000	8268.8\\
86000	8238.2\\
88000	8267.1\\
90000	8214.1\\
92000	8473.8\\
94000	8284.2\\
96000	8070.5\\
98000	8142.9\\
1e+05	8124.5\\
};
\addplot [color=mycolor34, forget plot]
  table[row sep=crcr]{%
0	4395.3\\
2000	4370\\
4000	3915.5\\
6000	4013.4\\
8000	3999\\
10000	4095.1\\
12000	4035.2\\
14000	3889.1\\
16000	3856.8\\
18000	4181.8\\
20000	4183.5\\
22000	4084.3\\
24000	4082.9\\
26000	4165.6\\
28000	4038.4\\
30000	3959.7\\
34000	3985.8\\
36000	4041.4\\
38000	3915.1\\
40000	3924.5\\
42000	4014.5\\
44000	3968.6\\
46000	3976.3\\
48000	4020.9\\
50000	3957\\
52000	4019.8\\
54000	3914.3\\
56000	3921.9\\
58000	3832.2\\
60000	4007.4\\
62000	4099.4\\
64000	4033.6\\
66000	4131.4\\
68000	4089.9\\
70000	3964.1\\
72000	3909\\
74000	3792.7\\
76000	3869.6\\
80000	3987.9\\
82000	3942.9\\
84000	4039.4\\
86000	4082.6\\
88000	3914.2\\
90000	4016.7\\
92000	3916.7\\
94000	4005.8\\
96000	4011.9\\
98000	3948.9\\
1e+05	3976.5\\
};
\addplot [color=mycolor35]
  table[row sep=crcr]{%
0	4395.3\\
2000	9399.8\\
4000	10121\\
6000	10071\\
8000	10251\\
12000	10454\\
14000	10225\\
16000	10095\\
18000	10263\\
20000	10335\\
22000	10264\\
24000	10268\\
26000	10366\\
30000	10122\\
32000	10447\\
34000	10530\\
36000	10506\\
38000	10364\\
40000	10142\\
42000	10210\\
44000	10151\\
46000	10321\\
48000	10282\\
50000	10512\\
52000	10195\\
54000	10220\\
56000	10169\\
58000	10449\\
60000	10150\\
62000	10411\\
64000	10258\\
66000	10293\\
68000	10119\\
70000	10492\\
72000	10105\\
74000	10551\\
76000	10398\\
78000	9991.4\\
80000	10079\\
82000	10268\\
84000	10014\\
88000	10329\\
90000	10132\\
92000	10228\\
94000	10170\\
96000	10240\\
98000	10061\\
1e+05	10015\\
};
\addlegendentry{\ClasskIILegend}

\addplot [color=mycolor36, forget plot]
  table[row sep=crcr]{%
0	4395.3\\
2000	4803.7\\
4000	4269\\
6000	4302.5\\
8000	4197.3\\
10000	4247.1\\
12000	4396.7\\
14000	4276.4\\
16000	4280.9\\
18000	4180.7\\
20000	4178.7\\
22000	4164.5\\
24000	4385.3\\
26000	4186.2\\
28000	4287.7\\
30000	4319.3\\
32000	4034.8\\
34000	4301.4\\
36000	4252.8\\
38000	4371.5\\
42000	4330\\
44000	4338.7\\
46000	4327.9\\
48000	4160.3\\
50000	4398.3\\
52000	4183.5\\
56000	4207.9\\
58000	4216.1\\
60000	4253.8\\
62000	4399\\
64000	4200.3\\
66000	4301.5\\
68000	4444.1\\
70000	4373.4\\
72000	4409\\
74000	4458.5\\
76000	4335.4\\
78000	4243.9\\
80000	4211.9\\
82000	4202.2\\
84000	4407\\
86000	4280\\
88000	4232.4\\
90000	4307.9\\
92000	4270.9\\
96000	4277.5\\
98000	4389.3\\
1e+05	4150.1\\
};
\addplot [color=mycolor37, forget plot]
  table[row sep=crcr]{%
0	4395.3\\
2000	3805.7\\
4000	3507.6\\
6000	3476.2\\
8000	3481.5\\
10000	3553.9\\
12000	3559.7\\
14000	3630.5\\
16000	3511.6\\
18000	3576.5\\
20000	3698.3\\
22000	3701.3\\
24000	3637.8\\
26000	3598.6\\
30000	3583.4\\
32000	3612.2\\
34000	3569.4\\
36000	3481.7\\
38000	3592.7\\
40000	3475.7\\
42000	3618.2\\
44000	3612\\
46000	3663.7\\
48000	3529.2\\
50000	3673.2\\
52000	3476.2\\
54000	3584.3\\
56000	3535.1\\
58000	3600.2\\
62000	3663.2\\
64000	3525.1\\
68000	3646.5\\
70000	3585.9\\
72000	3592.1\\
74000	3521\\
76000	3692\\
78000	3576.2\\
82000	3557.8\\
84000	3633\\
86000	3564.5\\
88000	3591\\
90000	3658.8\\
92000	3556.2\\
94000	3590.6\\
96000	3756.9\\
98000	3730.5\\
1e+05	3550\\
};
\addplot [color=mycolor38, forget plot]
  table[row sep=crcr]{%
0	4395.3\\
2000	3239.3\\
4000	2982.8\\
6000	3026.9\\
8000	2982\\
10000	2961.4\\
12000	2930.5\\
14000	2924.4\\
16000	3021.9\\
18000	2901.1\\
20000	2949.7\\
22000	2928.3\\
24000	2830.3\\
26000	3022.5\\
30000	2853.5\\
32000	2797.7\\
34000	2986.2\\
36000	2844.2\\
38000	2793.2\\
40000	3031\\
42000	3080.2\\
44000	2941.3\\
46000	2883.7\\
48000	2970.8\\
50000	2986.9\\
52000	2976.8\\
54000	2898.8\\
56000	2848.5\\
58000	3155.5\\
60000	2987.3\\
62000	2915.5\\
64000	2879.1\\
66000	2989\\
68000	2943.9\\
70000	2884.6\\
72000	2844.5\\
74000	2933.7\\
76000	2902.7\\
78000	3002.6\\
80000	2992.9\\
82000	3011.3\\
84000	3072.4\\
86000	2876\\
88000	2816.9\\
90000	3037.8\\
92000	2961.6\\
94000	2868.8\\
96000	3052.9\\
98000	3032.9\\
1e+05	2938.5\\
};
\addplot [color=mycolor39, forget plot]
  table[row sep=crcr]{%
0	4395.3\\
2000	12376\\
4000	12792\\
6000	12136\\
8000	12222\\
10000	11577\\
12000	11850\\
14000	12062\\
18000	12342\\
20000	12041\\
22000	11887\\
24000	12005\\
26000	12052\\
28000	11874\\
30000	12236\\
34000	11935\\
36000	11934\\
38000	12197\\
40000	12228\\
42000	12173\\
44000	12192\\
46000	11840\\
48000	12097\\
50000	11741\\
52000	11901\\
54000	12541\\
56000	12447\\
58000	11768\\
60000	12023\\
62000	12046\\
64000	11846\\
66000	11218\\
68000	12345\\
70000	12055\\
72000	12395\\
74000	11895\\
76000	11755\\
78000	11943\\
80000	11853\\
82000	11847\\
84000	12040\\
86000	11612\\
88000	12155\\
90000	12095\\
92000	11788\\
94000	12287\\
96000	11816\\
98000	12172\\
1e+05	12434\\
};
\addplot [color=mycolor40, forget plot]
  table[row sep=crcr]{%
0	4395.3\\
2000	5233.5\\
4000	5593.1\\
6000	5762\\
8000	5744.1\\
10000	5665\\
12000	5814\\
16000	5672\\
18000	5743.9\\
20000	5619.1\\
22000	5673.7\\
24000	5753.2\\
26000	5627.2\\
28000	5743.5\\
30000	5655.7\\
32000	5726.3\\
34000	5629.5\\
36000	5578.4\\
38000	5571.8\\
40000	5710.6\\
42000	5693.7\\
44000	5622\\
46000	5761.7\\
48000	5610.9\\
50000	5713.9\\
54000	5629.9\\
58000	5739\\
60000	5662.9\\
62000	5798.7\\
68000	5768.1\\
70000	5748.9\\
72000	5673.4\\
74000	5663.4\\
78000	5819.4\\
80000	5656.4\\
82000	5884.1\\
84000	5635.2\\
86000	5645.3\\
88000	5549.5\\
90000	5692.3\\
94000	5711.4\\
96000	5820.1\\
1e+05	5675.6\\
};
\addplot [color=mycolor41, forget plot]
  table[row sep=crcr]{%
0	4395.3\\
2000	4859.1\\
4000	5206.9\\
6000	5243.2\\
8000	5301.1\\
10000	5394.3\\
12000	5341\\
14000	5319\\
16000	5356.8\\
18000	5126.2\\
20000	5435.3\\
22000	5288\\
24000	5265.7\\
26000	5279.3\\
28000	5159.6\\
30000	5226.4\\
36000	5485.2\\
38000	5380.2\\
40000	5294.7\\
42000	5400.3\\
44000	5389.4\\
46000	5411.7\\
48000	5510.7\\
50000	5370.4\\
52000	5418.7\\
54000	5170.7\\
56000	5170.8\\
58000	5387.4\\
60000	5316.8\\
62000	5268.5\\
64000	5401.6\\
66000	5346\\
68000	5257.3\\
70000	5185.6\\
72000	5365.4\\
74000	5480\\
76000	5431.6\\
78000	5422.3\\
80000	5395.3\\
82000	5337.3\\
84000	5469.4\\
86000	5325.1\\
88000	5464.9\\
90000	5330.3\\
92000	5154.3\\
94000	5304.6\\
96000	5375.1\\
98000	5394.1\\
1e+05	5350.3\\
};
\addplot [color=mycolor42, forget plot]
  table[row sep=crcr]{%
0	4395.3\\
2000	17392\\
4000	20799\\
6000	20694\\
8000	21471\\
10000	20723\\
12000	20252\\
14000	19398\\
16000	20267\\
18000	19809\\
20000	20897\\
22000	20442\\
24000	21126\\
26000	20570\\
28000	20701\\
30000	20429\\
32000	21055\\
34000	21541\\
36000	22136\\
38000	21486\\
40000	20370\\
42000	20711\\
44000	21149\\
46000	20368\\
48000	20142\\
50000	21833\\
52000	21575\\
54000	21155\\
56000	21551\\
58000	20639\\
60000	20611\\
62000	20332\\
64000	21846\\
68000	21532\\
70000	21274\\
72000	21099\\
74000	21002\\
76000	21177\\
78000	21030\\
80000	21617\\
82000	21972\\
84000	21004\\
86000	20876\\
88000	20494\\
90000	21386\\
92000	20572\\
94000	20412\\
96000	21311\\
98000	20663\\
1e+05	21564\\
};
\addplot [color=mycolor43, forget plot]
  table[row sep=crcr]{%
0	4395.3\\
2000	5459.9\\
4000	5143.5\\
6000	4901.9\\
8000	4930\\
10000	4905.5\\
12000	5020.7\\
14000	4908.5\\
16000	4908.9\\
18000	5007.5\\
20000	5010.5\\
22000	4992.8\\
24000	4878.1\\
26000	4873\\
28000	5074.3\\
30000	4904.5\\
32000	4855.7\\
34000	4916.9\\
36000	5061.7\\
38000	4991.1\\
40000	4978.7\\
42000	5084.9\\
44000	5001.4\\
46000	5035.2\\
48000	4988.3\\
50000	4991.8\\
52000	4915.2\\
54000	4932.4\\
56000	4897.3\\
58000	4934.8\\
60000	4889.7\\
62000	4888.2\\
64000	4969\\
66000	4887.8\\
68000	4925.7\\
70000	5026.5\\
72000	4871.9\\
74000	4879.8\\
76000	4901.7\\
78000	5092.7\\
80000	5114.2\\
82000	4929.8\\
84000	4947.6\\
86000	5029.8\\
90000	4942.5\\
92000	4835.6\\
94000	4853.1\\
96000	4983.6\\
98000	5052.1\\
1e+05	5076\\
};
\addplot [color=mycolor44, forget plot]
  table[row sep=crcr]{%
0	4395.3\\
2000	3956.2\\
4000	3422.6\\
6000	3244.4\\
8000	3422.8\\
10000	3330\\
12000	3338.3\\
14000	3296.7\\
16000	3460.7\\
18000	3465.1\\
20000	3126.6\\
22000	3201.4\\
26000	3422.3\\
28000	3237.1\\
30000	3282.5\\
32000	3397\\
34000	3383.1\\
36000	3308.7\\
38000	3184.2\\
40000	3350.7\\
42000	3243.5\\
44000	3348.1\\
46000	3366.7\\
48000	3300.1\\
50000	3289.2\\
52000	3390.3\\
54000	3360.6\\
56000	3292.5\\
58000	3396.2\\
60000	3387\\
62000	3406.3\\
64000	3254.7\\
66000	3243.2\\
68000	3368.1\\
70000	3284.7\\
72000	3220.3\\
74000	3128.8\\
76000	3356.8\\
78000	3304.8\\
80000	3455.9\\
82000	3177.7\\
84000	3327.7\\
86000	3299.4\\
88000	3322\\
90000	3281\\
94000	3280.4\\
96000	3234.9\\
98000	3309\\
1e+05	3255.7\\
};
\addplot [color=mycolor45, forget plot]
  table[row sep=crcr]{%
0	4395.3\\
2000	6140.6\\
4000	6808\\
6000	6728.1\\
8000	6757.8\\
10000	6920.1\\
12000	6907.6\\
14000	6862.8\\
16000	6770.9\\
18000	6840.5\\
20000	7005.1\\
24000	6930.5\\
26000	7010.9\\
28000	6721.9\\
30000	6934.3\\
32000	7036.7\\
34000	6831.1\\
36000	6786.2\\
38000	6903.2\\
40000	6974.4\\
42000	6888.4\\
44000	6920\\
46000	6854.1\\
48000	7086.7\\
50000	6811.8\\
52000	6846.6\\
56000	6949.4\\
58000	6780.3\\
60000	6978.5\\
62000	6822.3\\
64000	6820.7\\
66000	7083.5\\
68000	6841.7\\
70000	6728.7\\
72000	6869\\
74000	7041.9\\
76000	6916.9\\
78000	7105.1\\
80000	6855.6\\
82000	7000\\
84000	6620.2\\
86000	6843.1\\
88000	6918\\
90000	7047\\
92000	6965.4\\
94000	6960.6\\
96000	6856.3\\
98000	6941.2\\
1e+05	6786.3\\
};
\addplot [color=mycolor46, forget plot]
  table[row sep=crcr]{%
0	4395.3\\
2000	3851.8\\
4000	3962.8\\
6000	3922\\
8000	4099.4\\
10000	4157\\
12000	3855.8\\
14000	3947.1\\
16000	4010.8\\
18000	3931.6\\
20000	4054.9\\
22000	4062.4\\
24000	3993.7\\
28000	4089.2\\
30000	4075.7\\
32000	3902\\
34000	3970.9\\
36000	4017\\
38000	3983.8\\
40000	3982.1\\
42000	4137.1\\
44000	4003.9\\
46000	4015.1\\
48000	4134.1\\
52000	3957.1\\
54000	3988.1\\
56000	3870.9\\
58000	3956.1\\
64000	4024.5\\
66000	4168.1\\
68000	3976.4\\
70000	4026.1\\
72000	4110.8\\
74000	3967.1\\
76000	4121.1\\
78000	4043\\
80000	3954\\
82000	3991.8\\
84000	4070.2\\
86000	3997.3\\
88000	4045.7\\
90000	3985.2\\
92000	4067.2\\
94000	3902.7\\
96000	4019.3\\
98000	3979.9\\
1e+05	4016.4\\
};
\addplot [color=mycolor47, forget plot]
  table[row sep=crcr]{%
0	4395.3\\
2000	5277.6\\
4000	5294.6\\
6000	5362\\
8000	5193.8\\
10000	5224.8\\
12000	5213.9\\
14000	5100.1\\
16000	5294.2\\
18000	5280.8\\
20000	5345.5\\
22000	5472.4\\
24000	5133.6\\
26000	5212.7\\
28000	5370.8\\
30000	5265.4\\
32000	5283.8\\
34000	5260.8\\
36000	5366.4\\
38000	5239.5\\
40000	5255.5\\
42000	5350.7\\
46000	5265.3\\
48000	5336\\
50000	5319\\
52000	5264.9\\
54000	5153.8\\
56000	5202.6\\
58000	5163.7\\
60000	5199.7\\
62000	5359.5\\
64000	5319.5\\
66000	5414.9\\
68000	5308.7\\
70000	5348\\
72000	5269.4\\
74000	5298.3\\
76000	5346.1\\
78000	5219.8\\
80000	5320.5\\
82000	5302.8\\
84000	5019.5\\
86000	5313.3\\
88000	5411.8\\
90000	5416.1\\
92000	5274.2\\
94000	5226.1\\
98000	5286.6\\
1e+05	5132.4\\
};
\addplot [color=mycolor48, forget plot]
  table[row sep=crcr]{%
0	4395.3\\
2000	3144.4\\
4000	2832.9\\
6000	2887\\
8000	2796.4\\
10000	2861.6\\
12000	2851\\
14000	2815\\
16000	2867.8\\
18000	2891.1\\
20000	2946\\
22000	2861\\
24000	2944.8\\
26000	2888.5\\
30000	2813.2\\
32000	2883.3\\
34000	2884.1\\
36000	2922.7\\
38000	2822.8\\
42000	2844.4\\
44000	2901.7\\
46000	2888.5\\
48000	2897.5\\
50000	2888.6\\
54000	2845.6\\
56000	2853\\
58000	2933.1\\
60000	2890.9\\
62000	2923.3\\
64000	2908.9\\
66000	2865.9\\
68000	2884.6\\
70000	2823.3\\
72000	2823.5\\
74000	2840.4\\
76000	2903.2\\
78000	2806.4\\
82000	2872.3\\
84000	2881.1\\
86000	2901.2\\
90000	2869.3\\
92000	2901.8\\
94000	2847.1\\
96000	2873.3\\
98000	2875.2\\
1e+05	2835.9\\
};
\addplot [color=mycolor49, forget plot]
  table[row sep=crcr]{%
0	4395.3\\
2000	4599.8\\
4000	4402.9\\
6000	4356\\
8000	4245.8\\
10000	4489.1\\
12000	4369.6\\
14000	4374.2\\
16000	4104\\
18000	4181.1\\
20000	4351.1\\
22000	4412.9\\
24000	4395.5\\
26000	4509.4\\
28000	4326.3\\
30000	4402.2\\
32000	4320\\
34000	4392.3\\
36000	4315.7\\
38000	4441.4\\
40000	4266.7\\
42000	4244.9\\
44000	4351.9\\
46000	4338.3\\
48000	4512.1\\
50000	4334.1\\
52000	4389.2\\
54000	4425.6\\
56000	4289.6\\
58000	4200.7\\
60000	4298.9\\
62000	4424.3\\
64000	4338.2\\
66000	4275.1\\
68000	4631\\
70000	4443.6\\
72000	4399\\
74000	4499.1\\
76000	4341.7\\
78000	4309.5\\
80000	4351.1\\
82000	4440.6\\
84000	4358\\
86000	4253.2\\
88000	4375.4\\
90000	4426\\
92000	4370.7\\
94000	4336\\
96000	4339.9\\
98000	4356.9\\
1e+05	4263.5\\
};
\addplot [color=mycolor50, forget plot]
  table[row sep=crcr]{%
0	4395.3\\
2000	2409.6\\
4000	2323.6\\
6000	2376.2\\
8000	2358.4\\
10000	2395.7\\
12000	2366.9\\
14000	2403.4\\
16000	2336.2\\
18000	2376.8\\
20000	2385\\
22000	2426.8\\
24000	2447.6\\
26000	2387.7\\
28000	2397\\
30000	2358.7\\
32000	2446.6\\
34000	2391.3\\
36000	2449.2\\
38000	2313.5\\
40000	2397.8\\
42000	2410.8\\
44000	2455\\
46000	2390.5\\
48000	2374.5\\
50000	2399\\
52000	2354.5\\
54000	2377\\
56000	2424.4\\
58000	2374.8\\
60000	2363.9\\
62000	2432\\
64000	2422\\
66000	2400.9\\
68000	2436.4\\
70000	2459.4\\
72000	2417.2\\
74000	2398.3\\
76000	2464.2\\
78000	2424.7\\
80000	2468.7\\
82000	2390\\
84000	2436.4\\
86000	2433\\
88000	2480.7\\
90000	2362.8\\
92000	2454\\
94000	2444.7\\
98000	2371.3\\
1e+05	2346.3\\
};
\addplot [color=mycolor51, forget plot]
  table[row sep=crcr]{%
0	4395.3\\
2000	5727.8\\
4000	5717.1\\
6000	5916.1\\
8000	5686.8\\
12000	5928.8\\
14000	5817\\
16000	6137.3\\
18000	6009.3\\
20000	6115.3\\
22000	5814.4\\
24000	6184.4\\
26000	5789.2\\
28000	6063.5\\
30000	5801.8\\
32000	6026.3\\
34000	5911.1\\
36000	5725\\
38000	6098.4\\
40000	6082\\
42000	5794.6\\
44000	6114.7\\
46000	5861.1\\
48000	5917.7\\
52000	5771.1\\
54000	5754.6\\
56000	5754.4\\
58000	6106.5\\
60000	6150.2\\
62000	6012.9\\
64000	5818.3\\
66000	6075.8\\
68000	5769.4\\
70000	5862.1\\
72000	5768.1\\
74000	5715.3\\
76000	5943.5\\
78000	5966.8\\
80000	6134.7\\
84000	5830.8\\
86000	6087.5\\
88000	6037.5\\
90000	5843.7\\
92000	6269.3\\
94000	5822.8\\
96000	5982.5\\
98000	5836.3\\
1e+05	5935.4\\
};
\addplot [color=mycolor52, forget plot]
  table[row sep=crcr]{%
0	4395.3\\
2000	4046.7\\
4000	3368.6\\
6000	3424.3\\
8000	3407.8\\
10000	3340.1\\
12000	3524.3\\
14000	3470.6\\
16000	3497.1\\
18000	3361.3\\
20000	3326.6\\
22000	3427.2\\
24000	3491.8\\
26000	3422.5\\
28000	3422.9\\
30000	3411.7\\
32000	3358.3\\
34000	3460.4\\
36000	3516.5\\
38000	3452.3\\
40000	3588.5\\
42000	3567.6\\
44000	3561.4\\
46000	3637.9\\
48000	3541.9\\
50000	3261.5\\
52000	3395.8\\
54000	3609.5\\
56000	3555.2\\
58000	3531.7\\
60000	3538.6\\
62000	3532.5\\
64000	3569.6\\
66000	3545.9\\
68000	3389.4\\
70000	3382.7\\
72000	3572.2\\
74000	3493.1\\
76000	3319.2\\
78000	3482.3\\
80000	3500\\
82000	3495.1\\
84000	3521.7\\
86000	3470.3\\
88000	3697.6\\
90000	3401.7\\
92000	3487.8\\
94000	3509.3\\
96000	3414.2\\
98000	3423.8\\
1e+05	3450.3\\
};
\addplot [color=mycolor53, forget plot]
  table[row sep=crcr]{%
0	4395.3\\
2000	6132.8\\
6000	5820.9\\
8000	5775.9\\
10000	5559.1\\
12000	5458.9\\
14000	5975.4\\
16000	5702.8\\
18000	5713\\
20000	5545.8\\
24000	5852.1\\
26000	5616.4\\
28000	5867.9\\
30000	5746.4\\
32000	5658.1\\
34000	5437.8\\
36000	5636.8\\
38000	5692.1\\
40000	5861.4\\
42000	5877.8\\
44000	5861.1\\
46000	5827.3\\
48000	5696.4\\
50000	5817.3\\
52000	5748.4\\
54000	5799.1\\
56000	5683.6\\
58000	5658.6\\
60000	5518\\
62000	5847.3\\
64000	5993\\
66000	5953.6\\
68000	5476.1\\
70000	5803.2\\
72000	5435.9\\
74000	5903.7\\
76000	5733.5\\
78000	5924.2\\
80000	5460.1\\
82000	5507\\
84000	6045.4\\
86000	5632.1\\
88000	5615.9\\
90000	5527.9\\
94000	5627.2\\
96000	5719.6\\
98000	5739\\
1e+05	5711.1\\
};
\addplot [color=mycolor54, forget plot]
  table[row sep=crcr]{%
0	4395.3\\
2000	4739\\
4000	4526.5\\
6000	4484.7\\
8000	4779.3\\
10000	4570.1\\
12000	4532.1\\
14000	4569.3\\
16000	4592.6\\
18000	4726.4\\
20000	4590.3\\
22000	4538.8\\
24000	4577.9\\
26000	4548.2\\
28000	4446.8\\
30000	4752.2\\
32000	4513.3\\
34000	4356.8\\
36000	4621\\
38000	4713.2\\
40000	4697.9\\
42000	4500.7\\
44000	4668.5\\
46000	4708.8\\
48000	4560.6\\
50000	4675.1\\
52000	4576.6\\
54000	4531.2\\
56000	4602.1\\
60000	4573.8\\
62000	4605.8\\
64000	4456.7\\
66000	4546.8\\
68000	4605.3\\
70000	4599\\
72000	4400.8\\
74000	4507.1\\
76000	4431.7\\
78000	4544.8\\
80000	4580\\
82000	4554.1\\
84000	4484.2\\
86000	4587.9\\
88000	4621\\
90000	4498.5\\
92000	4522.7\\
94000	4781.1\\
96000	4541.2\\
1e+05	4540.2\\
};
\addplot [color=mycolor55, forget plot]
  table[row sep=crcr]{%
0	4395.3\\
2000	5354\\
4000	5605.8\\
6000	5841.4\\
8000	5834.2\\
10000	5619.3\\
12000	5746.2\\
14000	5690.9\\
16000	5670.4\\
18000	5891.2\\
20000	5806.4\\
22000	5767\\
24000	5881.8\\
26000	5868\\
30000	5709.4\\
32000	5698.2\\
34000	5655.1\\
38000	5698.9\\
40000	5721.8\\
42000	5661.9\\
44000	5718.3\\
46000	5757.8\\
50000	5696.4\\
52000	5702.3\\
54000	5586.2\\
56000	5770.6\\
58000	5692.5\\
60000	5851.2\\
62000	5712.3\\
64000	5654\\
66000	5796.2\\
68000	5666.8\\
70000	5805.3\\
72000	5780.1\\
74000	5604\\
76000	5658.3\\
78000	5609\\
80000	5639.2\\
82000	5872.1\\
86000	5707.9\\
88000	5518.2\\
90000	5641.4\\
92000	5535.3\\
94000	5886.3\\
96000	5884.3\\
98000	5739.3\\
1e+05	5716.1\\
};
\addplot [color=mycolor56, forget plot]
  table[row sep=crcr]{%
0	4395.3\\
2000	4293.3\\
4000	4017.7\\
10000	4068.8\\
12000	3906.6\\
14000	4013.4\\
16000	4061.9\\
18000	4057.2\\
20000	3917.2\\
22000	3811.3\\
24000	3997.5\\
26000	3927\\
28000	3964.5\\
30000	3926.5\\
32000	4056.8\\
34000	3963.6\\
36000	4011.8\\
38000	4048.5\\
40000	4069\\
42000	3955.5\\
44000	3923.6\\
46000	3926.8\\
48000	3962.6\\
50000	3945.3\\
52000	3854.2\\
54000	3891.8\\
56000	4025.9\\
58000	3921.8\\
60000	4006.7\\
62000	3910.9\\
64000	4040\\
66000	3893.9\\
68000	3895.8\\
70000	3883.9\\
72000	3930.9\\
74000	3804.3\\
76000	4062.8\\
78000	3938.7\\
84000	3989.3\\
86000	4046.1\\
90000	3840.4\\
92000	4041\\
94000	3946.6\\
96000	3948.3\\
1e+05	4100.7\\
};
\addplot [color=mycolor57, forget plot]
  table[row sep=crcr]{%
0	4395.3\\
2000	6585.6\\
4000	6936.1\\
6000	6905.5\\
8000	6648.2\\
10000	6758.1\\
12000	6763.8\\
14000	6705.3\\
16000	6812.5\\
20000	6651.7\\
22000	6749.3\\
24000	6703.7\\
26000	6513.8\\
28000	6610.9\\
30000	6817.5\\
32000	6826\\
34000	6854\\
36000	6934.2\\
40000	6904\\
42000	6759.2\\
44000	6779.6\\
46000	6588.1\\
48000	6716.9\\
50000	6986.6\\
52000	6851.5\\
54000	6821.5\\
56000	6756.8\\
58000	6648.6\\
60000	6965.5\\
62000	6683.8\\
64000	6508.7\\
66000	6788.8\\
68000	6672.1\\
70000	6730.9\\
72000	6815.7\\
74000	6764.9\\
76000	6463.1\\
78000	6662.5\\
80000	6936.4\\
82000	6799\\
84000	6804.7\\
86000	6750.4\\
88000	6761.5\\
90000	6739\\
92000	6740\\
94000	6629.5\\
96000	7045.4\\
98000	6658.1\\
1e+05	6537.5\\
};
\addplot [color=mycolor58, forget plot]
  table[row sep=crcr]{%
0	4395.3\\
2000	3919\\
4000	3616.5\\
6000	3468.7\\
8000	3459.3\\
10000	3593\\
12000	3540.4\\
16000	3488.5\\
18000	3555\\
20000	3482.5\\
22000	3536.9\\
24000	3602.6\\
26000	3511.6\\
28000	3498\\
30000	3514.3\\
32000	3548.3\\
34000	3510.7\\
36000	3443.1\\
38000	3677.4\\
40000	3515.2\\
42000	3537.8\\
44000	3650.2\\
46000	3561.5\\
48000	3444.6\\
50000	3600.3\\
52000	3464.7\\
54000	3553.6\\
56000	3453.1\\
58000	3525.5\\
60000	3440.4\\
62000	3545.8\\
64000	3551.5\\
66000	3484\\
68000	3564.8\\
70000	3475.2\\
72000	3588\\
74000	3443.7\\
76000	3490.3\\
78000	3499.6\\
80000	3580.7\\
82000	3528.3\\
84000	3358.6\\
86000	3477.8\\
88000	3469.5\\
90000	3739.2\\
92000	3540.3\\
94000	3580.5\\
98000	3441.3\\
1e+05	3540.7\\
};
\addplot [color=mycolor59, forget plot]
  table[row sep=crcr]{%
0	4395.3\\
2000	4942.9\\
4000	4948.6\\
6000	5076.5\\
10000	4944.5\\
12000	4999.9\\
14000	4911.6\\
18000	5019\\
20000	4903.8\\
22000	5060.4\\
24000	4948.8\\
26000	4993.7\\
28000	4885.5\\
30000	4952.1\\
34000	4922.3\\
36000	5044\\
38000	5024.6\\
40000	4937.2\\
42000	5070.7\\
44000	4948.9\\
46000	5012.2\\
48000	5053\\
50000	4949.5\\
52000	5138.5\\
54000	5017\\
56000	4951.7\\
58000	4811.5\\
60000	5022\\
62000	5039.7\\
64000	4998.7\\
66000	5025.6\\
68000	5072.3\\
70000	4890.8\\
72000	5048\\
74000	4997.4\\
76000	5005.8\\
78000	5072.5\\
80000	4990.4\\
82000	5107.4\\
84000	4990.4\\
86000	5020\\
88000	4898.9\\
90000	5022\\
92000	5011.2\\
94000	5095.9\\
96000	5007.6\\
98000	5048.3\\
1e+05	4975.6\\
};
\addplot [color=mycolor60, forget plot]
  table[row sep=crcr]{%
0	4395.3\\
2000	3763.8\\
4000	3597.7\\
6000	3503.7\\
8000	3538.7\\
10000	3523.5\\
14000	3554.4\\
16000	3489.4\\
18000	3564.3\\
20000	3557.5\\
22000	3434.5\\
24000	3433\\
26000	3533.1\\
28000	3430.8\\
30000	3629.8\\
32000	3599.1\\
34000	3476.1\\
36000	3457.9\\
38000	3499.6\\
40000	3600.3\\
42000	3561.4\\
44000	3506.2\\
46000	3560.2\\
48000	3580.2\\
50000	3517.9\\
52000	3537.6\\
54000	3503\\
56000	3489.8\\
58000	3572.4\\
60000	3541.5\\
62000	3467.7\\
64000	3642.2\\
66000	3463.4\\
68000	3511\\
70000	3584\\
72000	3522.4\\
74000	3540.6\\
76000	3586.8\\
78000	3500.9\\
80000	3532.1\\
82000	3547.1\\
84000	3530.9\\
86000	3488.2\\
88000	3611.4\\
90000	3505\\
92000	3519.8\\
94000	3463.5\\
96000	3530.2\\
98000	3531.6\\
1e+05	3557.7\\
};
\addplot [color=mycolor61, forget plot]
  table[row sep=crcr]{%
0	4395.3\\
2000	3079.3\\
4000	3027.3\\
6000	3065\\
8000	3024.5\\
10000	3044.4\\
12000	3102.7\\
14000	3101.2\\
18000	3054.5\\
20000	3017.3\\
22000	3124.4\\
24000	3002.1\\
26000	3096.2\\
28000	3112.3\\
30000	3025.6\\
32000	3099.7\\
34000	3025.3\\
36000	3012.1\\
38000	3036\\
40000	3157.1\\
42000	3043.8\\
44000	3038.9\\
46000	3073.2\\
48000	3083.9\\
50000	3021.7\\
52000	3040.2\\
54000	3081.8\\
56000	3035.1\\
58000	3023\\
60000	3059.3\\
64000	3103.5\\
66000	3090\\
70000	3022.3\\
72000	3061.8\\
74000	3033.6\\
76000	3084\\
78000	3079.7\\
80000	3045.9\\
82000	3023\\
84000	3057.2\\
86000	3029.3\\
88000	3054.8\\
90000	3039.7\\
92000	3004.3\\
94000	3126.1\\
96000	3031.6\\
98000	3064.8\\
1e+05	3036\\
};
\addplot [color=mycolor62, forget plot]
  table[row sep=crcr]{%
0	4395.3\\
2000	4811\\
4000	5012.9\\
6000	4892.4\\
8000	4967.5\\
10000	4829.9\\
12000	4948.5\\
14000	4936.6\\
16000	5083.1\\
18000	4979.1\\
20000	4959\\
22000	4967.1\\
24000	5044.4\\
26000	5016.8\\
28000	5022.8\\
30000	5049.2\\
32000	4912.7\\
34000	5049.6\\
36000	4967.6\\
38000	4946.1\\
40000	4993.4\\
42000	4849.1\\
44000	4891.3\\
46000	5024.1\\
50000	4914.7\\
52000	5049.5\\
54000	5103.1\\
58000	4985.6\\
60000	5022.5\\
62000	4865.1\\
64000	4885.6\\
66000	5060.3\\
68000	5074.7\\
70000	4959.5\\
72000	4944.2\\
74000	5019.6\\
76000	4927.8\\
78000	4991.7\\
80000	4932.6\\
84000	4914.2\\
86000	4955.1\\
88000	4938.8\\
90000	4969.1\\
92000	5113.3\\
94000	5023\\
98000	4940.9\\
1e+05	4946.5\\
};
\addplot [color=mycolor63, forget plot]
  table[row sep=crcr]{%
0	4395.3\\
2000	2706.8\\
4000	2346\\
6000	2370.3\\
10000	2369.4\\
12000	2350.3\\
14000	2355.3\\
16000	2338\\
18000	2379.9\\
22000	2305.1\\
24000	2410.1\\
26000	2418.2\\
28000	2400.4\\
30000	2357.2\\
32000	2374.9\\
34000	2362.5\\
36000	2385\\
38000	2331.4\\
40000	2367\\
42000	2372.4\\
44000	2400.9\\
46000	2314.7\\
48000	2404.2\\
50000	2293.9\\
52000	2339.4\\
54000	2367.1\\
56000	2327\\
60000	2402\\
62000	2363.2\\
64000	2349.1\\
66000	2361\\
68000	2308.1\\
70000	2364.5\\
72000	2326.5\\
76000	2284.9\\
78000	2337.5\\
80000	2346.9\\
82000	2338\\
84000	2356.8\\
86000	2324.6\\
88000	2410.3\\
90000	2344.4\\
92000	2334.3\\
94000	2284\\
96000	2326.5\\
98000	2351.9\\
1e+05	2362.1\\
};
\addplot [color=mycolor64, forget plot]
  table[row sep=crcr]{%
0	4395.3\\
2000	3032.4\\
4000	2989.9\\
6000	3018.3\\
8000	2913.8\\
10000	2966.7\\
12000	2997.1\\
14000	2983.2\\
16000	3006.9\\
20000	3022.5\\
24000	2980\\
28000	3001.2\\
30000	2940.9\\
32000	3048.8\\
36000	2991.9\\
38000	2976.5\\
40000	3004.2\\
42000	2936.6\\
44000	2980.6\\
46000	3034.6\\
52000	2933.9\\
54000	2968\\
56000	2990.7\\
58000	3023.2\\
60000	2941.1\\
62000	3006.5\\
64000	3043.8\\
66000	3045.6\\
68000	2949.4\\
70000	2982.2\\
72000	2904.1\\
74000	3030.6\\
76000	2990.4\\
78000	2959.9\\
80000	3028.3\\
82000	3033.3\\
84000	2974.9\\
86000	2966.6\\
88000	3056.1\\
90000	2951.3\\
92000	3000.2\\
94000	2891.7\\
96000	2953.4\\
98000	2977.6\\
1e+05	2974.5\\
};
\addplot [color=mycolor65, forget plot]
  table[row sep=crcr]{%
0	4395.3\\
2000	2629.3\\
4000	2440.1\\
6000	2470.2\\
8000	2509.3\\
10000	2492.3\\
12000	2485.2\\
14000	2449.5\\
18000	2463.7\\
20000	2539.9\\
22000	2525.6\\
24000	2522.7\\
28000	2486.3\\
30000	2553.8\\
32000	2457.5\\
34000	2459.6\\
36000	2521.8\\
38000	2479.1\\
40000	2473.3\\
42000	2447.9\\
44000	2535.4\\
46000	2522.3\\
48000	2467.7\\
50000	2436\\
52000	2445.2\\
54000	2525.2\\
58000	2451.6\\
60000	2536.6\\
64000	2467.1\\
66000	2567.3\\
68000	2457.4\\
70000	2475.5\\
72000	2527.3\\
74000	2472.4\\
76000	2477.1\\
78000	2435.6\\
80000	2468.6\\
82000	2485.5\\
84000	2532\\
86000	2434.3\\
88000	2465.1\\
92000	2493.3\\
94000	2483\\
96000	2446.6\\
98000	2438.5\\
1e+05	2477.3\\
};
\addplot [color=mycolor66, forget plot]
  table[row sep=crcr]{%
0	4395.3\\
2000	3399.5\\
4000	3331.8\\
6000	3381.9\\
8000	3388.8\\
10000	3348.7\\
12000	3406.6\\
14000	3355.6\\
16000	3379\\
18000	3353.8\\
20000	3406.7\\
22000	3348.9\\
24000	3431.8\\
26000	3352.1\\
28000	3455.4\\
30000	3295.2\\
32000	3397.1\\
34000	3359.8\\
36000	3312.2\\
40000	3353.6\\
42000	3326.1\\
44000	3366.3\\
46000	3365\\
48000	3381.7\\
50000	3319.6\\
52000	3314.7\\
54000	3320.8\\
56000	3371.7\\
58000	3277.4\\
60000	3386.1\\
64000	3404.3\\
66000	3454.4\\
68000	3361\\
70000	3354.4\\
72000	3373.8\\
76000	3346.2\\
78000	3396.2\\
80000	3359.9\\
82000	3371.3\\
84000	3400.3\\
86000	3380.6\\
88000	3311.4\\
90000	3368.9\\
92000	3379.7\\
96000	3366.6\\
98000	3332.6\\
1e+05	3377\\
};
\addplot [color=mycolor67, forget plot]
  table[row sep=crcr]{%
0	4395.3\\
2000	3375.1\\
4000	3149.9\\
6000	3167.3\\
10000	3103.2\\
12000	3132.2\\
14000	3108.9\\
16000	3174.5\\
18000	3122.8\\
20000	3103.5\\
22000	3102.8\\
24000	3188.2\\
26000	3127.2\\
28000	3136.4\\
30000	3094.9\\
32000	3119.7\\
34000	3116.8\\
36000	3139.3\\
38000	3220\\
40000	3126.4\\
42000	3116.5\\
44000	3085.4\\
46000	3150.7\\
50000	3101.7\\
52000	3084.3\\
54000	3157.2\\
56000	3140.6\\
58000	3169.2\\
60000	3157\\
62000	3160.4\\
64000	3145\\
66000	3167.6\\
68000	3079.1\\
70000	3141.3\\
72000	3173.8\\
74000	3187.2\\
76000	3151\\
78000	3174.6\\
80000	3135.7\\
82000	3178\\
86000	3197.4\\
88000	3158.4\\
90000	3064.4\\
92000	3151.8\\
94000	3171.3\\
96000	3136.8\\
98000	3170.2\\
1e+05	3099.7\\
};
\addplot [color=mycolor67]
  table[row sep=crcr]{%
0	4395.3\\
2000	2668.7\\
4000	2609.8\\
6000	2590.3\\
8000	2538\\
16000	2570.9\\
18000	2559.2\\
20000	2633.6\\
22000	2565.7\\
24000	2562.3\\
26000	2532\\
28000	2552.5\\
30000	2543.1\\
32000	2598.8\\
34000	2595.9\\
38000	2534.4\\
42000	2588.6\\
44000	2546.3\\
48000	2597.1\\
50000	2562.2\\
52000	2590.1\\
54000	2568.7\\
56000	2583.2\\
58000	2619.8\\
60000	2559.4\\
64000	2541.1\\
66000	2589.4\\
68000	2580.4\\
70000	2530.6\\
72000	2574.9\\
74000	2608.5\\
76000	2533.4\\
78000	2598.6\\
80000	2565.6\\
82000	2565.7\\
84000	2557.5\\
86000	2606.9\\
88000	2536\\
90000	2566.3\\
92000	2512\\
94000	2594\\
96000	2591.2\\
98000	2561.5\\
1e+05	2546.3\\
};
\addlegendentry{\ClasskIIILegend}

\addplot [color=mycolor68, forget plot]
  table[row sep=crcr]{%
0	4395.3\\
2000	2733.6\\
4000	2429.5\\
6000	2454.1\\
8000	2411.9\\
10000	2401.4\\
12000	2401\\
14000	2408.9\\
16000	2395\\
18000	2409.9\\
20000	2394.2\\
22000	2410.7\\
28000	2388.1\\
30000	2419.5\\
32000	2372.5\\
34000	2406.2\\
36000	2403.4\\
40000	2383.1\\
42000	2431.4\\
44000	2387.4\\
46000	2390.9\\
48000	2447.9\\
50000	2398.4\\
52000	2389.6\\
54000	2454.4\\
56000	2409.6\\
58000	2427.5\\
60000	2422.7\\
62000	2389.8\\
64000	2441.6\\
66000	2392.4\\
68000	2386.1\\
70000	2430.5\\
72000	2385.3\\
74000	2369.7\\
76000	2381\\
78000	2401.8\\
80000	2465.5\\
82000	2456.2\\
84000	2433.2\\
86000	2369.8\\
88000	2423.7\\
90000	2408.1\\
92000	2402\\
94000	2373.4\\
96000	2358.8\\
98000	2404.6\\
1e+05	2387.5\\
};
\addplot [color=mycolor69, forget plot]
  table[row sep=crcr]{%
0	4395.3\\
2000	2936.5\\
4000	2772.6\\
6000	2826.2\\
8000	2840.3\\
14000	2837.8\\
16000	2824.2\\
18000	2863.6\\
20000	2835\\
22000	2862.2\\
24000	2872.8\\
26000	2867.6\\
28000	2871.4\\
30000	2825.2\\
32000	2882.5\\
34000	2879.3\\
36000	2781.7\\
40000	2882.7\\
42000	2836\\
44000	2819.7\\
48000	2810.3\\
50000	2787.2\\
52000	2813.6\\
54000	2813.1\\
56000	2847.2\\
58000	2846.4\\
60000	2808.5\\
62000	2859.4\\
64000	2834.6\\
66000	2880.7\\
68000	2828.9\\
70000	2866.3\\
72000	2882.9\\
74000	2862.7\\
76000	2852.4\\
78000	2850.9\\
80000	2859.2\\
86000	2784.6\\
88000	2878.4\\
90000	2842.3\\
92000	2855\\
94000	2805.7\\
96000	2832.4\\
1e+05	2875.3\\
};
\addplot [color=mycolor70, forget plot]
  table[row sep=crcr]{%
0	4395.3\\
2000	2168\\
4000	1999.3\\
6000	2033.1\\
8000	2026.4\\
10000	2100.2\\
12000	2070.8\\
14000	1974.2\\
16000	2002.4\\
18000	2042.1\\
20000	2021.6\\
22000	2076\\
24000	2020\\
26000	2067.8\\
28000	2040\\
30000	2050.1\\
32000	2043.2\\
34000	2029.6\\
38000	2032.5\\
40000	2060.9\\
42000	2070.5\\
44000	2042.2\\
46000	2024.1\\
52000	2036.8\\
54000	2030\\
56000	2082.2\\
58000	2113.1\\
60000	1988.7\\
62000	1979.9\\
64000	2039.3\\
66000	2016.5\\
68000	2085.1\\
70000	2034.8\\
72000	2016.5\\
74000	2055.3\\
76000	2017.8\\
78000	2044.1\\
80000	2020.4\\
82000	2047.6\\
84000	2058\\
86000	2019.3\\
88000	2035.6\\
90000	2031.4\\
92000	2049.4\\
94000	2054.3\\
96000	2028.2\\
98000	2049.6\\
1e+05	2059.3\\
};
\addplot [color=mycolor71, forget plot]
  table[row sep=crcr]{%
0	4395.3\\
2000	2678.8\\
4000	2529.2\\
6000	2456.9\\
8000	2429.1\\
10000	2443.9\\
12000	2419.8\\
14000	2460.8\\
16000	2452.5\\
18000	2414\\
20000	2402.9\\
22000	2441.2\\
24000	2428\\
26000	2452.9\\
28000	2409.4\\
30000	2427.7\\
32000	2454\\
34000	2446.2\\
36000	2427.4\\
38000	2439.2\\
40000	2440.8\\
42000	2479.9\\
44000	2424.7\\
46000	2416.1\\
48000	2473.6\\
52000	2450.6\\
54000	2418.2\\
56000	2428.3\\
58000	2421.5\\
62000	2462.6\\
64000	2417.1\\
66000	2385.7\\
70000	2392.2\\
74000	2490.7\\
76000	2425.5\\
78000	2395.2\\
80000	2422.9\\
82000	2462.5\\
84000	2404.8\\
86000	2467.9\\
88000	2469.7\\
90000	2412\\
92000	2425.6\\
94000	2432.2\\
96000	2398.9\\
98000	2410.2\\
1e+05	2443.3\\
};
\addplot [color=mycolor72, forget plot]
  table[row sep=crcr]{%
0	4395.3\\
2000	2432.7\\
4000	2088\\
6000	2125.9\\
8000	2109.1\\
10000	2100.9\\
12000	2131.3\\
14000	2132.9\\
16000	2102\\
18000	2098.3\\
20000	2160.3\\
22000	2118.1\\
24000	2168.8\\
26000	2122\\
28000	2183.6\\
30000	2106.7\\
32000	2173.7\\
34000	2139.3\\
36000	2114.4\\
38000	2132.9\\
40000	2170.2\\
42000	2154.4\\
44000	2117.9\\
46000	2089.8\\
48000	2115.8\\
50000	2102.2\\
54000	2183.6\\
56000	2125\\
58000	2154.3\\
60000	2113.6\\
62000	2188.3\\
64000	2180.8\\
66000	2112.6\\
68000	2140.6\\
70000	2114\\
72000	2096.3\\
74000	2153.5\\
76000	2112.6\\
80000	2153.1\\
82000	2104.6\\
84000	2110\\
86000	2157.4\\
88000	2124.1\\
90000	2166.3\\
92000	2145.6\\
94000	2095.6\\
96000	2156.3\\
98000	2138.5\\
1e+05	2162\\
};
\addplot [color=mycolor73, forget plot]
  table[row sep=crcr]{%
0	4395.3\\
2000	2767.1\\
4000	2586.7\\
6000	2564.1\\
10000	2570.8\\
12000	2542.8\\
14000	2572.3\\
16000	2524\\
18000	2528.9\\
20000	2512.9\\
22000	2540.1\\
24000	2537.6\\
26000	2605.5\\
28000	2583.5\\
30000	2532.9\\
32000	2517.6\\
34000	2570.7\\
36000	2528.9\\
40000	2560.7\\
42000	2512.8\\
46000	2583.7\\
48000	2582.6\\
50000	2636.5\\
52000	2576.7\\
54000	2552\\
56000	2616.3\\
58000	2566.7\\
60000	2585.3\\
64000	2557.6\\
66000	2521.6\\
68000	2516.3\\
70000	2548.7\\
72000	2540.8\\
74000	2560.9\\
76000	2554.1\\
78000	2574.2\\
80000	2546.8\\
82000	2527.8\\
84000	2569.6\\
86000	2584\\
88000	2575.5\\
90000	2578.3\\
92000	2544.4\\
94000	2545\\
96000	2565.9\\
98000	2575.5\\
1e+05	2572.8\\
};
\addplot [color=mycolor74, forget plot]
  table[row sep=crcr]{%
0	4395.3\\
2000	1826\\
4000	1490.3\\
6000	1470.6\\
8000	1471.8\\
10000	1485.8\\
12000	1470.7\\
14000	1510.1\\
16000	1476.1\\
18000	1500.1\\
20000	1471\\
22000	1505.6\\
24000	1510.5\\
26000	1484\\
28000	1492.2\\
34000	1467.1\\
36000	1504.6\\
38000	1490.7\\
40000	1496.5\\
42000	1479.1\\
44000	1475.9\\
46000	1501.2\\
48000	1504.9\\
50000	1462.5\\
52000	1443.6\\
56000	1490.5\\
60000	1482.6\\
62000	1490\\
64000	1468.3\\
66000	1484.2\\
68000	1466.9\\
72000	1495\\
74000	1494.2\\
76000	1481\\
78000	1461.4\\
82000	1481.5\\
84000	1469.4\\
88000	1471.9\\
90000	1480\\
94000	1453.1\\
96000	1487.7\\
98000	1499.7\\
1e+05	1499.2\\
};
\addplot [color=mycolor75, forget plot]
  table[row sep=crcr]{%
0	4395.3\\
2000	2475.5\\
4000	2129.1\\
6000	2117.3\\
8000	2175.3\\
10000	2143.4\\
12000	2118.7\\
14000	2155.3\\
16000	2132.2\\
18000	2181\\
20000	2119.1\\
22000	2131.1\\
24000	2167.2\\
26000	2140.7\\
28000	2128.3\\
30000	2171\\
32000	2143.2\\
34000	2138.2\\
36000	2158.9\\
38000	2170.3\\
40000	2152.4\\
42000	2120.7\\
44000	2154.9\\
46000	2151.4\\
48000	2140.5\\
50000	2175.5\\
52000	2189.2\\
54000	2103.7\\
56000	2176.6\\
58000	2163.3\\
60000	2169.1\\
62000	2148.2\\
66000	2142.9\\
68000	2101.7\\
70000	2127.4\\
72000	2164\\
74000	2154.6\\
76000	2138.9\\
78000	2115.9\\
82000	2150.4\\
84000	2184.6\\
86000	2164\\
88000	2121.9\\
90000	2166\\
92000	2154.7\\
94000	2173.9\\
98000	2133.7\\
1e+05	2153.8\\
};
\addplot [color=mycolor76, forget plot]
  table[row sep=crcr]{%
0	4395.3\\
2000	2265.4\\
4000	1969.2\\
6000	1964.6\\
8000	1954.1\\
12000	1948.3\\
14000	1955.8\\
16000	1983.3\\
18000	1948.8\\
22000	1966.1\\
24000	1969.4\\
26000	1956.7\\
28000	1953.4\\
30000	1983\\
32000	1951.6\\
34000	1954.6\\
36000	1945\\
40000	1959.9\\
44000	1949.6\\
46000	1983.2\\
48000	1977.7\\
50000	1954.3\\
52000	1938.2\\
56000	1993.9\\
58000	1952.2\\
60000	1973.3\\
62000	1987.1\\
64000	1971.8\\
66000	1936.3\\
68000	1937.9\\
70000	1908\\
72000	1929.7\\
74000	1968.6\\
76000	1952.8\\
78000	1989.2\\
82000	1935.4\\
84000	1959.4\\
86000	1971.2\\
88000	1949.6\\
90000	1956.9\\
92000	1981.6\\
94000	1970.3\\
96000	1952.7\\
98000	1983.2\\
1e+05	1973.2\\
};
\addplot [color=mycolor77, forget plot]
  table[row sep=crcr]{%
0	4395.3\\
2000	1897.9\\
4000	1561.9\\
6000	1519\\
8000	1536.7\\
10000	1518.6\\
12000	1508.5\\
14000	1491.2\\
16000	1516.3\\
18000	1493.9\\
20000	1517.6\\
22000	1470.2\\
24000	1517.9\\
26000	1507.9\\
28000	1506.3\\
30000	1548.6\\
32000	1504.2\\
36000	1506.9\\
38000	1508.7\\
40000	1526.7\\
42000	1499.4\\
44000	1551\\
46000	1543.1\\
48000	1521.9\\
50000	1491.3\\
52000	1520.2\\
54000	1494.2\\
56000	1532.9\\
58000	1518.9\\
60000	1533.4\\
62000	1505.9\\
64000	1490.5\\
66000	1509\\
68000	1500.3\\
70000	1496.7\\
72000	1518.5\\
74000	1504.9\\
76000	1520.9\\
78000	1496\\
80000	1506.3\\
82000	1493.3\\
84000	1529.4\\
86000	1528.2\\
88000	1533.2\\
90000	1500.5\\
92000	1518.2\\
94000	1476.9\\
96000	1508.3\\
98000	1513.6\\
1e+05	1505.1\\
};
\addplot [color=mycolor78, forget plot]
  table[row sep=crcr]{%
0	4395.3\\
2000	2567.3\\
4000	2259.5\\
6000	2246.8\\
8000	2271.1\\
10000	2198.2\\
12000	2244.1\\
14000	2279.3\\
16000	2252.5\\
18000	2208.9\\
22000	2257.2\\
24000	2145.4\\
26000	2198.3\\
28000	2236.7\\
30000	2245.1\\
32000	2219\\
34000	2296.9\\
36000	2254.9\\
38000	2237.4\\
40000	2277.8\\
42000	2280.4\\
44000	2210.3\\
46000	2216.4\\
48000	2246.9\\
50000	2241.9\\
52000	2298.1\\
54000	2253.9\\
56000	2257.9\\
58000	2230.3\\
62000	2232.2\\
64000	2247.1\\
66000	2187.3\\
68000	2218.8\\
70000	2226.6\\
72000	2261.7\\
76000	2277.4\\
78000	2282.6\\
80000	2227\\
82000	2215.7\\
84000	2262.7\\
86000	2206.8\\
88000	2240.7\\
90000	2288.4\\
92000	2228.6\\
94000	2257.7\\
96000	2194.7\\
98000	2234.9\\
1e+05	2235.2\\
};
\addplot [color=mycolor79, forget plot]
  table[row sep=crcr]{%
0	4395.3\\
2000	2379.1\\
4000	2122.3\\
6000	2115.7\\
8000	2128.2\\
10000	2133.8\\
12000	2115.1\\
14000	2126.8\\
16000	2097.7\\
18000	2086.6\\
20000	2068.1\\
22000	2066.9\\
24000	2104.3\\
26000	2033.2\\
28000	2116.2\\
30000	2133.8\\
32000	2121.2\\
34000	2156.8\\
36000	2114.5\\
38000	2147\\
40000	2152.6\\
42000	2081\\
44000	2083.6\\
46000	2067.4\\
48000	2122.5\\
50000	2038.4\\
52000	2107.1\\
54000	2166.1\\
56000	2121.5\\
58000	2086.9\\
60000	2121.3\\
62000	2058.2\\
64000	2134.2\\
66000	2094.5\\
68000	2129.6\\
70000	2117.6\\
72000	2079.1\\
74000	2053.2\\
76000	2132.3\\
80000	2082.6\\
82000	2112.8\\
84000	2135.7\\
86000	2110.8\\
88000	2114.3\\
90000	2107.3\\
94000	2070.4\\
96000	2105.9\\
98000	2150.8\\
1e+05	2085.2\\
};
\addplot [color=mycolor79, forget plot]
  table[row sep=crcr]{%
0	4395.3\\
2000	2331.8\\
4000	1985.3\\
6000	1940\\
8000	2008.8\\
10000	1922.8\\
12000	1959.5\\
14000	1952.8\\
18000	1996.2\\
20000	1964.3\\
22000	1898.3\\
24000	1980.6\\
28000	1953\\
30000	2020.7\\
32000	2028.7\\
34000	1957.4\\
36000	1991\\
38000	1954.4\\
40000	1962.8\\
42000	1945.2\\
44000	1934.9\\
46000	1955.7\\
48000	1943.3\\
50000	1937.7\\
52000	1992.7\\
54000	1950\\
56000	1945.5\\
58000	1921.7\\
60000	1941.4\\
62000	1911.8\\
64000	1977.2\\
66000	1971.8\\
70000	1894.8\\
72000	1998.9\\
74000	1935.9\\
76000	1970.1\\
78000	1909.5\\
80000	1943.2\\
82000	1949.5\\
84000	2011.3\\
86000	1947.9\\
88000	1944.9\\
90000	1973.7\\
92000	1921.7\\
96000	1963.8\\
98000	1919.5\\
1e+05	1977.6\\
};
\addplot [color=mycolor80, forget plot]
  table[row sep=crcr]{%
0	4395.3\\
2000	1846\\
4000	1505\\
6000	1485.8\\
8000	1450.9\\
10000	1483.6\\
12000	1499.6\\
14000	1508.3\\
16000	1502\\
18000	1509.2\\
20000	1485.9\\
22000	1484.5\\
24000	1503.3\\
26000	1515.9\\
28000	1499.1\\
30000	1510\\
32000	1496.2\\
34000	1502.5\\
36000	1503.7\\
38000	1481.8\\
40000	1470.3\\
42000	1493.9\\
44000	1496.8\\
46000	1489.5\\
48000	1492.6\\
50000	1489.8\\
52000	1464.4\\
54000	1469.3\\
56000	1500.1\\
58000	1503.6\\
60000	1479.9\\
62000	1522.9\\
66000	1475.7\\
68000	1475\\
70000	1479\\
72000	1511.7\\
74000	1501.9\\
76000	1481.2\\
78000	1496\\
80000	1488.8\\
82000	1496.9\\
84000	1499.5\\
86000	1497.3\\
88000	1483.3\\
90000	1487.1\\
92000	1525.2\\
94000	1508.5\\
96000	1464\\
98000	1477.6\\
1e+05	1470.9\\
};
\addplot [color=mycolor81, forget plot]
  table[row sep=crcr]{%
0	4395.3\\
2000	1761.6\\
4000	1425.3\\
6000	1382\\
8000	1402.7\\
10000	1393.5\\
12000	1401.8\\
14000	1400.8\\
16000	1408.3\\
18000	1391.2\\
20000	1403.1\\
22000	1406.9\\
24000	1395.6\\
26000	1402.2\\
28000	1391.4\\
32000	1404.8\\
34000	1430.9\\
36000	1379.5\\
38000	1412.5\\
40000	1426.7\\
42000	1404.6\\
44000	1394.7\\
46000	1379\\
48000	1389.9\\
50000	1372.6\\
52000	1376.5\\
54000	1370\\
56000	1393.2\\
58000	1406.9\\
62000	1404.9\\
64000	1388.1\\
66000	1348.3\\
70000	1427.4\\
72000	1428.4\\
74000	1407.9\\
76000	1414.9\\
78000	1413.8\\
82000	1381.7\\
84000	1374\\
86000	1344.5\\
88000	1390.7\\
90000	1393.9\\
92000	1419.2\\
94000	1407.4\\
98000	1428.6\\
1e+05	1380.2\\
};
\addplot [color=mycolor82, forget plot]
  table[row sep=crcr]{%
0	4395.3\\
2000	2018.9\\
4000	1649.7\\
6000	1659.8\\
8000	1627.5\\
10000	1621.4\\
12000	1652.8\\
14000	1653.5\\
16000	1682.4\\
20000	1639.3\\
24000	1617.8\\
28000	1683.3\\
30000	1639.8\\
32000	1642.7\\
34000	1620.1\\
36000	1660.1\\
38000	1659\\
40000	1627.8\\
42000	1640.6\\
44000	1621.2\\
46000	1624.8\\
48000	1644.7\\
52000	1646.1\\
54000	1655.6\\
56000	1645.2\\
58000	1649\\
60000	1658.5\\
62000	1648.1\\
64000	1654.8\\
66000	1668.6\\
68000	1662.6\\
70000	1665.7\\
72000	1620\\
74000	1637.2\\
76000	1640.4\\
78000	1663.7\\
80000	1653.1\\
82000	1655\\
86000	1669.1\\
88000	1648.2\\
90000	1678.5\\
92000	1671.9\\
94000	1637.4\\
96000	1613\\
98000	1648\\
1e+05	1644.8\\
};
\addplot [color=mycolor83, forget plot]
  table[row sep=crcr]{%
0	4395.3\\
2000	2064.4\\
4000	1711.1\\
6000	1655\\
8000	1639.2\\
14000	1666.1\\
16000	1697.2\\
18000	1654\\
20000	1676.6\\
22000	1660.6\\
24000	1680.1\\
26000	1686.2\\
28000	1678.6\\
30000	1658.3\\
32000	1679.5\\
34000	1654.7\\
36000	1676\\
38000	1676.3\\
40000	1659.6\\
42000	1681.3\\
44000	1691.5\\
48000	1657.2\\
50000	1702.1\\
52000	1710\\
54000	1692.9\\
56000	1693.4\\
58000	1664.5\\
60000	1649.5\\
62000	1661\\
64000	1666.1\\
66000	1685.6\\
68000	1647.3\\
70000	1646.3\\
72000	1653.1\\
74000	1688.1\\
76000	1671.2\\
78000	1663.3\\
80000	1674.1\\
84000	1677\\
86000	1652.9\\
90000	1687.1\\
92000	1662.3\\
94000	1680.7\\
96000	1689.3\\
98000	1681.6\\
1e+05	1683.9\\
};
\addplot [color=mycolor83, forget plot]
  table[row sep=crcr]{%
0	4395.3\\
2000	1560.9\\
4000	1230.8\\
6000	1182\\
8000	1208.3\\
10000	1216.9\\
12000	1164.1\\
14000	1209.3\\
16000	1188.3\\
18000	1195.8\\
20000	1219\\
22000	1200.4\\
24000	1221.4\\
26000	1196.6\\
28000	1210.6\\
32000	1214.1\\
34000	1183.5\\
36000	1214.3\\
38000	1212\\
40000	1203.2\\
42000	1221.5\\
46000	1216.3\\
48000	1195\\
50000	1213.6\\
54000	1191.2\\
56000	1238.6\\
58000	1225.8\\
60000	1237.9\\
64000	1203.8\\
66000	1199.3\\
68000	1189.4\\
70000	1226.1\\
72000	1204.7\\
74000	1199.8\\
78000	1212.4\\
80000	1192.7\\
82000	1211.9\\
84000	1211.1\\
86000	1238.2\\
88000	1202\\
90000	1206.3\\
92000	1195.1\\
94000	1224\\
96000	1187.1\\
98000	1179.8\\
1e+05	1231.3\\
};
\addplot [color=mycolor84, forget plot]
  table[row sep=crcr]{%
0	4395.3\\
2000	1618\\
4000	1205.7\\
6000	1226.7\\
8000	1221.4\\
10000	1225.2\\
12000	1176.9\\
14000	1214.2\\
16000	1213.9\\
18000	1239.7\\
20000	1235.1\\
22000	1213.8\\
24000	1200.7\\
26000	1222.3\\
28000	1211.2\\
30000	1231.9\\
32000	1222.4\\
34000	1229\\
36000	1214.3\\
38000	1180.5\\
40000	1242.2\\
42000	1203\\
44000	1216.2\\
46000	1207.5\\
48000	1207.1\\
50000	1211.6\\
52000	1207.7\\
54000	1190.2\\
56000	1234.7\\
58000	1193.9\\
60000	1200.9\\
62000	1227.1\\
64000	1208.7\\
66000	1179.4\\
68000	1212.3\\
70000	1228.1\\
72000	1221.2\\
74000	1203\\
76000	1237.9\\
78000	1194.7\\
80000	1235.5\\
82000	1255\\
84000	1210.3\\
86000	1210.1\\
88000	1225.2\\
90000	1236.8\\
92000	1222.4\\
94000	1251.8\\
96000	1178.6\\
98000	1196.7\\
1e+05	1228.3\\
};
\addplot [color=mycolor85, forget plot]
  table[row sep=crcr]{%
0	4395.3\\
2000	1772\\
4000	1331.7\\
6000	1381.3\\
8000	1362.1\\
10000	1296.9\\
12000	1382.2\\
14000	1351.7\\
16000	1335.7\\
18000	1357.6\\
20000	1366\\
22000	1329.5\\
24000	1355.4\\
26000	1377.2\\
28000	1350.5\\
30000	1377.5\\
32000	1357\\
34000	1385.1\\
36000	1387.3\\
38000	1332.4\\
42000	1350.9\\
44000	1376.2\\
46000	1379.8\\
48000	1313.1\\
50000	1341\\
52000	1345.9\\
54000	1319.5\\
56000	1331\\
60000	1349.6\\
62000	1315.8\\
64000	1330.5\\
66000	1313.1\\
68000	1354.8\\
70000	1331.2\\
74000	1318.7\\
76000	1341.1\\
78000	1339.7\\
80000	1357.8\\
82000	1311.3\\
84000	1373.5\\
86000	1339.8\\
88000	1376.3\\
90000	1344.1\\
92000	1376.3\\
94000	1339\\
96000	1333.4\\
98000	1346.8\\
1e+05	1320.7\\
};
\addplot [color=mycolor86, forget plot]
  table[row sep=crcr]{%
0	4395.3\\
2000	1562.1\\
4000	1141.5\\
6000	1139.6\\
8000	1152.6\\
10000	1138\\
12000	1102.2\\
14000	1106.7\\
16000	1156.3\\
18000	1118.9\\
20000	1106.5\\
22000	1111.7\\
24000	1089.6\\
26000	1128.2\\
28000	1113.6\\
30000	1110.1\\
32000	1113.6\\
34000	1138.5\\
36000	1097.7\\
38000	1143.7\\
40000	1114.4\\
42000	1127.6\\
44000	1112.3\\
46000	1155.3\\
48000	1100.3\\
52000	1147.4\\
54000	1097.5\\
56000	1122.5\\
58000	1131\\
60000	1143.1\\
62000	1119.1\\
64000	1143.2\\
66000	1081.8\\
70000	1097.7\\
72000	1137\\
74000	1113.6\\
76000	1145.6\\
82000	1105.1\\
84000	1119.5\\
86000	1141.3\\
88000	1152.1\\
90000	1107.9\\
92000	1109.5\\
94000	1145.5\\
96000	1129.6\\
98000	1143.9\\
1e+05	1095.3\\
};
\addplot [color=mycolor87, forget plot]
  table[row sep=crcr]{%
0	4395.3\\
2000	1366.6\\
4000	948.48\\
6000	931.3\\
8000	920.36\\
10000	901.13\\
12000	984.78\\
14000	935.33\\
16000	902.85\\
18000	929.34\\
20000	991.98\\
22000	931.52\\
24000	892.23\\
26000	952.95\\
28000	896.45\\
30000	903.22\\
32000	929.29\\
34000	923.47\\
36000	966.03\\
38000	942.73\\
40000	931.69\\
42000	912.9\\
44000	974.51\\
46000	955.65\\
48000	947.52\\
50000	930.69\\
52000	927.97\\
54000	904.61\\
58000	946.44\\
60000	902.59\\
62000	923.99\\
64000	934.81\\
66000	939.72\\
68000	914.61\\
70000	993.72\\
72000	951.53\\
74000	964.67\\
76000	891.35\\
78000	927.46\\
80000	972.47\\
82000	896.83\\
84000	948.52\\
86000	938.54\\
88000	948.48\\
90000	983.97\\
92000	912.17\\
94000	967.27\\
96000	926.74\\
98000	938.62\\
1e+05	926.58\\
};
\addplot [color=mycolor88, forget plot]
  table[row sep=crcr]{%
0	4395.3\\
2000	1596.1\\
4000	1168.8\\
6000	1134.1\\
8000	1115\\
10000	1065.1\\
12000	1105.2\\
14000	1152.4\\
16000	1162\\
18000	1100.9\\
20000	1104.7\\
22000	1134.6\\
24000	1123.1\\
26000	1131.2\\
28000	1108.9\\
30000	1110.6\\
32000	1092.2\\
34000	1119.3\\
36000	1125.4\\
38000	1126.8\\
40000	1119.6\\
42000	1171.5\\
44000	1098.9\\
46000	1143\\
48000	1142.8\\
50000	1126.8\\
52000	1159.1\\
54000	1141.4\\
56000	1118.3\\
58000	1146.6\\
60000	1143\\
62000	1157.6\\
66000	1153.3\\
68000	1115.1\\
70000	1161.4\\
72000	1129\\
74000	1121.8\\
76000	1148.5\\
78000	1135.6\\
80000	1163.9\\
82000	1122.9\\
84000	1097.2\\
86000	1166.4\\
88000	1146.3\\
90000	1132\\
92000	1185.1\\
94000	1146\\
96000	1156\\
98000	1105.6\\
1e+05	1126.1\\
};
\addplot [color=mycolor89, forget plot]
  table[row sep=crcr]{%
0	4395.3\\
2000	1382.8\\
4000	809.62\\
6000	775.52\\
8000	829.62\\
10000	786.88\\
12000	814.95\\
14000	818.53\\
16000	767.26\\
18000	787.71\\
22000	810.89\\
24000	814.5\\
26000	788.94\\
28000	817.89\\
30000	825.52\\
32000	797.82\\
34000	807.98\\
36000	830.12\\
38000	813.81\\
40000	824.95\\
42000	831.08\\
44000	790.16\\
46000	776.91\\
48000	786.32\\
50000	774.58\\
52000	801.6\\
54000	809.29\\
56000	800.52\\
58000	801.81\\
60000	838.91\\
62000	802.64\\
64000	818.95\\
66000	803.26\\
68000	793.45\\
70000	779.69\\
72000	820.81\\
74000	821.85\\
76000	832.22\\
78000	793.11\\
80000	801.21\\
82000	778.95\\
84000	792.19\\
86000	797.07\\
88000	774.59\\
90000	799.9\\
92000	806.3\\
94000	805.74\\
96000	792.85\\
98000	792.17\\
1e+05	816.85\\
};
\addplot [color=mycolor90, forget plot]
  table[row sep=crcr]{%
0	4395.3\\
2000	1223.5\\
4000	775.37\\
6000	761.29\\
8000	742.48\\
10000	756.12\\
12000	766.57\\
14000	781.66\\
16000	786.49\\
18000	786.26\\
20000	790.6\\
22000	792.18\\
24000	782.13\\
26000	765.45\\
28000	783.52\\
30000	773\\
32000	819.93\\
34000	803.7\\
36000	760.22\\
38000	772.3\\
40000	769.52\\
42000	776.07\\
44000	796.48\\
46000	779.06\\
48000	780.98\\
50000	774.45\\
52000	763.96\\
54000	786.92\\
56000	801.07\\
58000	741.14\\
60000	762.87\\
62000	779.18\\
64000	787.46\\
66000	784.09\\
68000	777.96\\
70000	782.11\\
72000	796.89\\
74000	780.56\\
76000	767.38\\
78000	787.28\\
82000	781.31\\
84000	763.78\\
86000	752.23\\
88000	764.31\\
90000	791.47\\
92000	756.43\\
94000	750.32\\
96000	781.22\\
98000	775.34\\
1e+05	784.19\\
};
\addplot [color=mycolor91, forget plot]
  table[row sep=crcr]{%
0	4395.3\\
2000	1287.4\\
4000	691.86\\
6000	662.25\\
10000	645.11\\
14000	673.65\\
16000	683.57\\
18000	679.99\\
20000	668.66\\
22000	667.43\\
24000	672.37\\
26000	685.13\\
28000	677.69\\
30000	663.21\\
32000	692.7\\
34000	655.98\\
36000	674.39\\
38000	655\\
40000	660.21\\
42000	623.1\\
44000	658.7\\
46000	683.6\\
48000	688.32\\
50000	684.08\\
52000	688.6\\
56000	675.56\\
58000	690.07\\
60000	694.59\\
62000	727.45\\
64000	680.17\\
66000	669.7\\
68000	683.23\\
70000	694.63\\
72000	691.82\\
74000	655.56\\
76000	635.17\\
78000	676.49\\
80000	692.07\\
82000	644.62\\
84000	650.75\\
86000	666.91\\
88000	688.59\\
90000	684.91\\
92000	653.9\\
94000	679.55\\
96000	677.02\\
98000	681.23\\
1e+05	660.55\\
};
\addplot [color=mycolor92, forget plot]
  table[row sep=crcr]{%
0	4395.3\\
2000	1373.7\\
4000	883.62\\
6000	844.24\\
8000	851.9\\
10000	838.45\\
12000	810.45\\
14000	809.16\\
16000	861.93\\
18000	859.26\\
20000	850.42\\
22000	788.85\\
24000	842.73\\
26000	870.99\\
28000	820.39\\
30000	862.27\\
32000	828.08\\
34000	827.34\\
36000	843.55\\
38000	843.14\\
42000	829.21\\
44000	844.52\\
46000	837.57\\
48000	847.98\\
50000	829.85\\
52000	824.59\\
54000	844.24\\
56000	874.04\\
58000	843.56\\
60000	841.22\\
62000	853.06\\
64000	860.27\\
66000	850.22\\
70000	870.34\\
72000	851.73\\
74000	826.49\\
76000	840.36\\
78000	865.57\\
80000	865.85\\
82000	825.13\\
84000	850.3\\
86000	814.37\\
88000	835.67\\
90000	832.81\\
92000	845.8\\
94000	836.82\\
96000	848.15\\
98000	818.02\\
1e+05	835.3\\
};
\addplot [color=mycolor93, forget plot]
  table[row sep=crcr]{%
0	4395.3\\
2000	1368\\
4000	692.65\\
6000	717.96\\
8000	686.99\\
10000	695.77\\
12000	685.01\\
14000	718.05\\
16000	737.75\\
18000	646.04\\
20000	723.28\\
22000	719.85\\
24000	727.53\\
26000	732.68\\
28000	714.85\\
30000	747.93\\
32000	713.14\\
34000	691.8\\
36000	681.28\\
38000	711.71\\
40000	704.42\\
42000	678.19\\
44000	705.98\\
46000	719.05\\
48000	726.02\\
50000	687.26\\
52000	697.49\\
54000	688.45\\
58000	712.79\\
60000	730.45\\
62000	702.04\\
64000	689.45\\
66000	679.24\\
68000	734.01\\
70000	702.89\\
72000	711.72\\
74000	733.39\\
76000	659.82\\
78000	722.33\\
80000	708.08\\
82000	705.51\\
84000	732.98\\
86000	673.06\\
88000	700.77\\
90000	748.99\\
92000	686.68\\
94000	713.39\\
96000	727.31\\
1e+05	714.76\\
};
\addplot [color=mycolor94, forget plot]
  table[row sep=crcr]{%
0	4395.3\\
2000	1363\\
4000	682.09\\
6000	662.62\\
8000	666.96\\
10000	676.58\\
12000	669.2\\
14000	643.82\\
16000	671.93\\
18000	692.68\\
20000	672.52\\
22000	667.62\\
24000	667.33\\
26000	654.97\\
28000	664.74\\
30000	648.67\\
32000	674.58\\
34000	660.31\\
36000	654.38\\
38000	673.58\\
40000	661.2\\
42000	671.43\\
44000	657.53\\
46000	666.58\\
48000	667.46\\
50000	661.2\\
54000	664.45\\
56000	638.1\\
58000	655.04\\
60000	679.41\\
62000	672.95\\
64000	658.93\\
66000	673.93\\
68000	701.86\\
70000	681.95\\
72000	672.75\\
74000	657.89\\
76000	647.41\\
78000	681.63\\
80000	640.28\\
82000	657.5\\
84000	643.5\\
86000	688.78\\
88000	672.11\\
90000	661.65\\
92000	676.46\\
94000	664.35\\
96000	675.21\\
98000	665.37\\
1e+05	668.37\\
};
\addplot [color=mycolor95, forget plot]
  table[row sep=crcr]{%
0	4395.3\\
2000	1031.7\\
4000	408.12\\
6000	410.12\\
8000	383.93\\
10000	376.99\\
12000	391.57\\
14000	380.51\\
16000	402.7\\
18000	393.63\\
20000	407.49\\
22000	383.02\\
24000	383.85\\
26000	392.72\\
28000	391.09\\
30000	402.4\\
32000	399.81\\
34000	383.86\\
36000	412.38\\
38000	418.35\\
40000	413.22\\
42000	386.23\\
44000	382.28\\
46000	398.08\\
48000	401.97\\
50000	387.49\\
52000	384.03\\
54000	386.09\\
56000	403.59\\
58000	417.41\\
60000	387.71\\
62000	398.35\\
64000	413.33\\
66000	386.96\\
68000	387.06\\
70000	399.5\\
72000	401.41\\
74000	401.03\\
76000	381.77\\
78000	407.2\\
80000	391.4\\
82000	391.24\\
84000	387.42\\
86000	394.72\\
88000	379.92\\
90000	374.54\\
92000	404.34\\
94000	394.62\\
96000	382.25\\
98000	383.88\\
1e+05	408.29\\
};
\addplot [color=mycolor96, forget plot]
  table[row sep=crcr]{%
0	4395.3\\
2000	1227.7\\
4000	421.61\\
6000	440.56\\
8000	419.83\\
10000	411.76\\
12000	420.21\\
14000	403.86\\
16000	427.66\\
18000	432.37\\
20000	425.04\\
22000	403.93\\
24000	357.85\\
26000	413.56\\
28000	406.31\\
30000	404.07\\
32000	369.9\\
34000	407.01\\
36000	400.7\\
38000	412.63\\
40000	398.9\\
42000	396.71\\
44000	427.11\\
46000	403.91\\
48000	398.47\\
50000	421.71\\
52000	462.07\\
54000	417.59\\
56000	388.58\\
58000	409.28\\
60000	395.62\\
62000	390.85\\
64000	423.11\\
66000	426.45\\
68000	422.77\\
70000	388.06\\
72000	428.22\\
74000	421.89\\
76000	420.47\\
78000	422.58\\
80000	446.09\\
82000	374.74\\
84000	402.36\\
86000	403.58\\
88000	430.2\\
90000	430.45\\
92000	406.14\\
94000	412.7\\
96000	422.15\\
98000	420.54\\
1e+05	392.18\\
};
\addplot [color=mycolor97, forget plot]
  table[row sep=crcr]{%
0	4395.3\\
2000	862.18\\
4000	312.85\\
6000	306.71\\
8000	312.81\\
10000	306.75\\
12000	278.77\\
16000	311.35\\
22000	298.13\\
24000	288.06\\
26000	300.71\\
28000	311.89\\
30000	289.6\\
32000	296.01\\
34000	318.75\\
36000	321.06\\
38000	320.31\\
40000	280.27\\
42000	307.3\\
44000	300.05\\
46000	285.06\\
48000	293.92\\
50000	288.81\\
52000	296.98\\
54000	310.06\\
56000	307.7\\
58000	299.95\\
60000	307.83\\
62000	312.08\\
64000	285.77\\
66000	279.61\\
68000	283.54\\
70000	285\\
72000	303.1\\
74000	268.52\\
76000	288.63\\
78000	281.4\\
80000	263.43\\
82000	317.21\\
84000	311.63\\
86000	300.5\\
88000	304.03\\
90000	313.11\\
92000	296.75\\
94000	335.08\\
96000	291.79\\
98000	297.64\\
1e+05	311.36\\
};
\addplot [color=mycolor98]
  table[row sep=crcr]{%
0	4395.3\\
2000	1189.1\\
4000	426.86\\
6000	384.51\\
8000	391.22\\
10000	394.08\\
12000	401.71\\
14000	397.1\\
16000	403.95\\
18000	407.52\\
20000	384.33\\
22000	376.74\\
24000	379.63\\
26000	391.31\\
28000	389.81\\
32000	410.34\\
34000	367.87\\
36000	389.46\\
38000	406.1\\
40000	409.45\\
42000	399.08\\
44000	415.67\\
46000	414.37\\
48000	392.46\\
50000	381.86\\
52000	393.92\\
54000	395.49\\
56000	423.27\\
58000	421.55\\
60000	392.23\\
62000	395.81\\
64000	417.7\\
66000	383.24\\
68000	397.28\\
70000	425.53\\
72000	398.5\\
74000	391.13\\
76000	425.02\\
78000	399.43\\
80000	372.35\\
82000	359.12\\
84000	413.7\\
86000	397.59\\
88000	379.1\\
90000	398.84\\
94000	383.69\\
98000	414.89\\
1e+05	399.51\\
};
\addlegendentry{\ClasskIVLegend}

      % \input{../../../../.tex/20/KS/SizeEvolutionsfig9.tex}
      \legend{}

      \nextgroupplot[%
        % xmin=0,xmax=1e+05,
        % ymin=100,ymax=1.6117e+05,
      ] % This file was created by matlab2tikz.
%
\definecolor{mycolor1}{rgb}{0.00146,0.00047,0.01387}%
\definecolor{mycolor2}{rgb}{0.01166,0.00942,0.06346}%
\definecolor{mycolor3}{rgb}{0.02943,0.02150,0.11462}%
\definecolor{mycolor4}{rgb}{0.06134,0.03659,0.17764}%
\definecolor{mycolor5}{rgb}{0.08741,0.04456,0.22481}%
\definecolor{mycolor6}{rgb}{0.12291,0.04754,0.28162}%
\definecolor{mycolor7}{rgb}{0.14907,0.04547,0.31709}%
\definecolor{mycolor8}{rgb}{0.18343,0.04033,0.35497}%
\definecolor{mycolor9}{rgb}{0.21795,0.03662,0.38352}%
\definecolor{mycolor10}{rgb}{0.24497,0.03705,0.40001}%
\definecolor{mycolor11}{rgb}{0.27135,0.04092,0.41198}%
\definecolor{mycolor12}{rgb}{0.29076,0.04564,0.41864}%
\definecolor{mycolor13}{rgb}{0.31628,0.05349,0.42512}%
\definecolor{mycolor14}{rgb}{0.33522,0.06006,0.42852}%
\definecolor{mycolor15}{rgb}{0.36028,0.06925,0.43150}%
\definecolor{mycolor16}{rgb}{0.37900,0.07625,0.43272}%
\definecolor{mycolor17}{rgb}{0.39767,0.08326,0.43318}%
\definecolor{mycolor18}{rgb}{0.41633,0.09020,0.43294}%
\definecolor{mycolor19}{rgb}{0.42877,0.09479,0.43241}%
\definecolor{mycolor20}{rgb}{0.44743,0.10160,0.43108}%
\definecolor{mycolor21}{rgb}{0.46610,0.10832,0.42912}%
\definecolor{mycolor22}{rgb}{0.47856,0.11276,0.42747}%
\definecolor{mycolor23}{rgb}{0.49726,0.11938,0.42449}%
\definecolor{mycolor24}{rgb}{0.50973,0.12377,0.42216}%
\definecolor{mycolor25}{rgb}{0.52221,0.12815,0.41955}%
\definecolor{mycolor26}{rgb}{0.53468,0.13253,0.41667}%
\definecolor{mycolor27}{rgb}{0.55339,0.13913,0.41183}%
\definecolor{mycolor28}{rgb}{0.56585,0.14357,0.40826}%
\definecolor{mycolor29}{rgb}{0.57830,0.14804,0.40441}%
\definecolor{mycolor30}{rgb}{0.59073,0.15256,0.40029}%
\definecolor{mycolor31}{rgb}{0.60314,0.15715,0.39589}%
\definecolor{mycolor32}{rgb}{0.61551,0.16182,0.39122}%
\definecolor{mycolor33}{rgb}{0.62169,0.16418,0.38878}%
\definecolor{mycolor34}{rgb}{0.63400,0.16899,0.38370}%
\definecolor{mycolor35}{rgb}{0.64626,0.17391,0.37836}%
\definecolor{mycolor36}{rgb}{0.65846,0.17896,0.37275}%
\definecolor{mycolor37}{rgb}{0.66454,0.18154,0.36985}%
\definecolor{mycolor38}{rgb}{0.67664,0.18681,0.36385}%
\definecolor{mycolor39}{rgb}{0.68865,0.19224,0.35760}%
\definecolor{mycolor40}{rgb}{0.69463,0.19502,0.35439}%
\definecolor{mycolor41}{rgb}{0.70650,0.20073,0.34778}%
\definecolor{mycolor42}{rgb}{0.71240,0.20366,0.34438}%
\definecolor{mycolor43}{rgb}{0.72410,0.20967,0.33742}%
\definecolor{mycolor44}{rgb}{0.72991,0.21276,0.33386}%
\definecolor{mycolor45}{rgb}{0.74142,0.21911,0.32658}%
\definecolor{mycolor46}{rgb}{0.74713,0.22238,0.32286}%
\definecolor{mycolor47}{rgb}{0.75279,0.22571,0.31909}%
\definecolor{mycolor48}{rgb}{0.76401,0.23255,0.31140}%
\definecolor{mycolor49}{rgb}{0.76956,0.23608,0.30749}%
\definecolor{mycolor50}{rgb}{0.77506,0.23967,0.30353}%
\definecolor{mycolor51}{rgb}{0.79129,0.25086,0.29139}%
\definecolor{mycolor52}{rgb}{0.79661,0.25473,0.28726}%
\definecolor{mycolor53}{rgb}{0.80187,0.25867,0.28310}%
\definecolor{mycolor54}{rgb}{0.81224,0.26679,0.27466}%
\definecolor{mycolor55}{rgb}{0.81734,0.27095,0.27039}%
\definecolor{mycolor56}{rgb}{0.82737,0.27952,0.26175}%
\definecolor{mycolor57}{rgb}{0.83230,0.28391,0.25738}%
\definecolor{mycolor58}{rgb}{0.83717,0.28839,0.25299}%
\definecolor{mycolor59}{rgb}{0.84671,0.29756,0.24411}%
\definecolor{mycolor60}{rgb}{0.85138,0.30226,0.23964}%
\definecolor{mycolor61}{rgb}{0.85599,0.30704,0.23513}%
\definecolor{mycolor62}{rgb}{0.86053,0.31189,0.23061}%
\definecolor{mycolor63}{rgb}{0.87374,0.32691,0.21689}%
\definecolor{mycolor64}{rgb}{0.87800,0.33206,0.21227}%
\definecolor{mycolor65}{rgb}{0.89034,0.34796,0.19829}%
\definecolor{mycolor66}{rgb}{0.89819,0.35891,0.18886}%
\definecolor{mycolor67}{rgb}{0.90200,0.36449,0.18412}%
\definecolor{mycolor68}{rgb}{0.90573,0.37014,0.17935}%
\definecolor{mycolor69}{rgb}{0.90939,0.37586,0.17456}%
\definecolor{mycolor70}{rgb}{0.91297,0.38164,0.16975}%
\definecolor{mycolor71}{rgb}{0.91646,0.38748,0.16492}%
\definecolor{mycolor72}{rgb}{0.91988,0.39339,0.16007}%
\definecolor{mycolor73}{rgb}{0.92322,0.39936,0.15519}%
\definecolor{mycolor74}{rgb}{0.92647,0.40539,0.15029}%
\definecolor{mycolor75}{rgb}{0.92964,0.41148,0.14537}%
\definecolor{mycolor76}{rgb}{0.93274,0.41763,0.14042}%
\definecolor{mycolor77}{rgb}{0.93575,0.42383,0.13544}%
\definecolor{mycolor78}{rgb}{0.93868,0.43009,0.13044}%
\definecolor{mycolor79}{rgb}{0.94152,0.43640,0.12541}%
\definecolor{mycolor80}{rgb}{0.94429,0.44277,0.12035}%
\definecolor{mycolor81}{rgb}{0.94956,0.45566,0.11016}%
\definecolor{mycolor82}{rgb}{0.95208,0.46218,0.10503}%
\definecolor{mycolor83}{rgb}{0.96129,0.48872,0.08429}%
\definecolor{mycolor84}{rgb}{0.96540,0.50225,0.07386}%
\definecolor{mycolor85}{rgb}{0.96732,0.50908,0.06866}%
\definecolor{mycolor86}{rgb}{0.97259,0.52980,0.05332}%
\definecolor{mycolor87}{rgb}{0.97568,0.54380,0.04362}%
\definecolor{mycolor88}{rgb}{0.98082,0.57221,0.02851}%
\definecolor{mycolor89}{rgb}{0.98189,0.57939,0.02625}%
\definecolor{mycolor90}{rgb}{0.98378,0.59385,0.02377}%
\definecolor{mycolor91}{rgb}{0.98532,0.60842,0.02420}%
\definecolor{mycolor92}{rgb}{0.98696,0.63048,0.03091}%
\definecolor{mycolor93}{rgb}{0.98793,0.66025,0.05175}%
\definecolor{mycolor94}{rgb}{0.98771,0.68281,0.07249}%
\definecolor{mycolor95}{rgb}{0.97750,0.78226,0.18592}%
\definecolor{mycolor96}{rgb}{0.96624,0.83619,0.26153}%
\definecolor{mycolor97}{rgb}{0.96252,0.85148,0.28555}%
\definecolor{mycolor98}{rgb}{0.98836,0.99836,0.64492}%
%
\addplot [color=mycolor1]
  table[row sep=crcr]{%
0	4395.3\\
2000	1.3026e+05\\
4000	1.584e+05\\
6000	1.5438e+05\\
8000	1.5703e+05\\
10000	1.5205e+05\\
12000	1.5447e+05\\
14000	1.5472e+05\\
16000	1.5719e+05\\
20000	1.5215e+05\\
24000	1.5403e+05\\
26000	1.605e+05\\
28000	1.5573e+05\\
30000	1.5639e+05\\
32000	1.5766e+05\\
34000	1.5412e+05\\
36000	1.5633e+05\\
40000	1.5678e+05\\
42000	1.5652e+05\\
44000	1.573e+05\\
46000	1.5528e+05\\
48000	1.5937e+05\\
52000	1.5349e+05\\
54000	1.531e+05\\
56000	1.587e+05\\
58000	1.5975e+05\\
60000	1.5549e+05\\
62000	1.5666e+05\\
64000	1.6117e+05\\
66000	1.5579e+05\\
68000	1.5845e+05\\
70000	1.5573e+05\\
72000	1.5846e+05\\
74000	1.5338e+05\\
76000	1.5831e+05\\
78000	1.6075e+05\\
82000	1.5537e+05\\
84000	1.5777e+05\\
86000	1.5569e+05\\
92000	1.5969e+05\\
94000	1.5483e+05\\
96000	1.5409e+05\\
98000	1.5584e+05\\
1e+05	1.5627e+05\\
};
\addlegendentry{\ClasskILegend}

\addplot [color=mycolor2, forget plot]
  table[row sep=crcr]{%
0	4395.3\\
2000	23864\\
4000	26265\\
6000	25755\\
8000	26826\\
10000	26215\\
12000	26327\\
14000	26859\\
16000	26384\\
18000	26430\\
20000	26235\\
22000	26416\\
24000	26687\\
26000	26273\\
28000	26828\\
30000	27180\\
32000	26450\\
34000	26112\\
38000	26024\\
40000	25910\\
42000	26700\\
44000	26080\\
46000	25965\\
48000	26197\\
50000	26187\\
52000	26617\\
54000	26563\\
56000	26611\\
58000	26019\\
60000	25297\\
62000	25970\\
64000	25858\\
66000	26400\\
68000	26571\\
70000	26006\\
72000	26359\\
74000	26604\\
76000	26407\\
78000	26460\\
80000	26746\\
82000	26911\\
84000	26353\\
86000	26993\\
88000	25351\\
94000	26651\\
96000	26337\\
98000	26644\\
1e+05	26554\\
};
\addplot [color=mycolor3, forget plot]
  table[row sep=crcr]{%
0	4395.3\\
2000	37023\\
4000	41341\\
6000	40309\\
8000	41296\\
10000	41462\\
14000	41328\\
16000	40407\\
20000	40038\\
22000	40599\\
24000	40082\\
26000	41144\\
30000	39988\\
32000	39279\\
34000	39275\\
36000	40399\\
38000	41272\\
42000	40579\\
44000	40694\\
46000	40964\\
48000	39788\\
50000	40196\\
52000	40398\\
54000	38959\\
56000	38918\\
58000	40667\\
60000	41078\\
62000	40682\\
64000	39848\\
66000	39986\\
68000	40842\\
70000	41366\\
72000	40644\\
74000	40858\\
76000	39951\\
78000	40456\\
80000	40353\\
82000	40872\\
86000	41104\\
88000	39773\\
90000	40563\\
92000	39576\\
94000	41109\\
96000	40378\\
98000	40928\\
1e+05	40025\\
};
\addplot [color=mycolor4, forget plot]
  table[row sep=crcr]{%
0	4395.3\\
2000	13777\\
4000	14467\\
8000	14335\\
10000	14591\\
12000	14477\\
14000	14296\\
16000	14480\\
18000	14621\\
20000	14618\\
22000	14745\\
24000	14527\\
26000	14030\\
28000	14466\\
30000	14596\\
32000	14540\\
34000	14976\\
36000	14393\\
38000	14794\\
40000	14534\\
42000	14815\\
44000	14155\\
46000	14150\\
48000	14798\\
50000	14377\\
52000	14544\\
54000	14438\\
58000	14520\\
60000	14710\\
62000	14565\\
64000	14234\\
66000	14570\\
68000	14569\\
70000	15018\\
72000	14797\\
74000	14415\\
76000	14687\\
78000	14326\\
80000	14091\\
82000	14659\\
84000	14015\\
86000	14725\\
88000	14477\\
90000	14436\\
92000	14141\\
94000	14719\\
96000	14770\\
98000	14459\\
1e+05	14868\\
};
\addplot [color=mycolor5, forget plot]
  table[row sep=crcr]{%
0	4395.3\\
2000	52592\\
4000	63494\\
6000	65062\\
8000	65745\\
10000	67464\\
12000	64945\\
14000	62217\\
16000	64151\\
18000	63054\\
20000	65715\\
22000	66305\\
24000	66245\\
26000	64685\\
30000	65701\\
32000	67349\\
34000	65310\\
36000	65060\\
38000	63970\\
40000	64055\\
42000	64979\\
44000	67326\\
46000	66353\\
48000	66167\\
50000	66240\\
52000	65289\\
54000	66043\\
56000	65023\\
58000	63731\\
60000	64782\\
62000	65150\\
64000	65003\\
66000	64186\\
68000	64033\\
70000	63584\\
72000	62042\\
74000	64875\\
76000	65056\\
78000	65436\\
80000	64079\\
82000	68033\\
84000	66190\\
86000	63488\\
88000	66128\\
90000	63348\\
92000	63630\\
94000	65650\\
96000	64921\\
98000	65232\\
1e+05	66109\\
};
\addplot [color=mycolor6, forget plot]
  table[row sep=crcr]{%
0	4395.3\\
2000	50301\\
4000	58900\\
6000	58251\\
8000	58119\\
10000	57702\\
12000	58277\\
14000	57346\\
16000	57455\\
18000	56557\\
20000	59402\\
22000	57954\\
24000	57996\\
26000	57595\\
28000	58628\\
30000	59005\\
32000	55026\\
34000	57559\\
36000	55976\\
38000	55001\\
40000	58073\\
42000	59815\\
44000	56712\\
46000	58463\\
48000	57132\\
50000	58293\\
52000	58693\\
54000	60249\\
56000	58231\\
58000	57347\\
60000	59351\\
62000	55741\\
64000	56556\\
66000	57703\\
68000	55944\\
70000	57109\\
72000	56484\\
74000	57149\\
76000	58448\\
78000	58434\\
80000	59194\\
82000	57757\\
84000	58425\\
86000	57337\\
88000	59004\\
90000	58919\\
92000	58516\\
94000	58324\\
96000	57145\\
98000	59070\\
1e+05	55302\\
};
\addplot [color=mycolor7, forget plot]
  table[row sep=crcr]{%
0	4395.3\\
2000	26463\\
4000	28437\\
6000	29571\\
8000	28561\\
10000	28446\\
12000	28582\\
14000	29250\\
16000	28827\\
18000	28241\\
20000	27446\\
22000	27855\\
24000	29286\\
26000	28805\\
30000	29331\\
32000	28912\\
34000	28904\\
36000	28532\\
38000	29704\\
40000	29234\\
42000	28402\\
44000	28803\\
46000	28757\\
48000	28360\\
50000	28910\\
52000	29179\\
54000	28885\\
56000	28922\\
58000	28207\\
60000	28571\\
64000	28421\\
66000	28226\\
68000	28559\\
70000	29641\\
72000	28806\\
74000	28307\\
76000	28254\\
78000	28932\\
80000	28576\\
82000	27699\\
84000	27893\\
86000	28809\\
88000	28527\\
90000	28766\\
92000	27706\\
94000	28760\\
96000	29647\\
98000	28722\\
1e+05	29464\\
};
\addplot [color=mycolor8, forget plot]
  table[row sep=crcr]{%
0	4395.3\\
2000	19557\\
4000	20854\\
6000	21086\\
8000	20781\\
10000	21242\\
12000	20685\\
14000	21038\\
16000	19957\\
18000	20186\\
20000	21236\\
22000	20969\\
24000	20874\\
26000	20997\\
28000	21248\\
30000	20984\\
32000	21445\\
34000	20710\\
36000	21474\\
38000	21100\\
40000	20807\\
42000	19860\\
44000	19947\\
46000	20410\\
48000	20326\\
50000	21113\\
52000	20949\\
54000	21192\\
56000	20748\\
58000	20989\\
60000	20724\\
62000	20626\\
64000	21134\\
66000	20997\\
68000	20441\\
70000	20958\\
72000	21114\\
76000	20763\\
80000	20160\\
82000	20601\\
84000	20495\\
86000	20546\\
88000	20736\\
90000	20514\\
92000	20915\\
94000	21078\\
96000	21379\\
98000	22108\\
1e+05	20280\\
};
\addplot [color=mycolor9, forget plot]
  table[row sep=crcr]{%
0	4395.3\\
2000	13784\\
4000	13891\\
6000	14301\\
8000	14123\\
10000	14909\\
12000	13984\\
14000	14183\\
16000	14524\\
18000	14672\\
20000	14095\\
22000	14275\\
24000	14082\\
26000	13840\\
28000	14141\\
30000	14712\\
32000	14429\\
34000	14067\\
36000	13838\\
38000	14348\\
40000	14264\\
42000	14034\\
44000	14133\\
46000	14053\\
48000	14226\\
50000	14086\\
52000	14939\\
56000	13955\\
58000	13963\\
60000	14209\\
62000	14278\\
64000	13984\\
66000	13927\\
68000	14361\\
70000	14163\\
72000	14195\\
74000	14300\\
76000	14452\\
78000	13624\\
82000	14158\\
84000	14251\\
86000	14441\\
88000	14437\\
90000	13883\\
92000	14079\\
94000	14332\\
96000	14418\\
98000	13721\\
1e+05	14040\\
};
\addplot [color=mycolor10, forget plot]
  table[row sep=crcr]{%
0	4395.3\\
2000	20560\\
4000	22761\\
6000	22970\\
8000	22518\\
10000	22905\\
12000	22697\\
14000	22738\\
16000	22987\\
18000	22183\\
20000	22401\\
24000	22376\\
26000	22699\\
28000	22333\\
30000	22329\\
32000	22540\\
34000	22865\\
36000	21634\\
38000	22724\\
40000	22785\\
42000	22341\\
44000	22185\\
46000	22659\\
48000	22467\\
52000	22438\\
54000	23009\\
56000	22789\\
58000	22673\\
62000	22269\\
64000	22222\\
66000	23193\\
68000	23149\\
70000	22286\\
72000	22579\\
74000	21783\\
76000	22513\\
78000	22212\\
80000	22566\\
82000	22061\\
84000	22321\\
86000	23127\\
88000	22829\\
90000	22169\\
92000	22271\\
94000	22112\\
96000	22194\\
98000	22477\\
1e+05	22131\\
};
\addplot [color=mycolor11, forget plot]
  table[row sep=crcr]{%
0	4395.3\\
2000	8402.9\\
4000	8774\\
6000	8635.1\\
8000	8626.5\\
10000	8878.3\\
12000	9095.2\\
14000	8768.4\\
16000	8660.4\\
18000	8663.2\\
20000	8526.2\\
22000	8577.9\\
24000	8538.1\\
26000	8675.7\\
28000	8359.4\\
30000	8778.8\\
32000	8463.1\\
34000	8801.7\\
36000	8714.2\\
38000	8386.3\\
40000	8640.2\\
42000	8696.4\\
44000	8426.5\\
46000	8612.8\\
50000	8519\\
52000	8681.9\\
56000	8532.8\\
58000	8768.7\\
60000	8844.7\\
62000	8562.4\\
64000	8539\\
66000	8767.1\\
68000	8452\\
70000	8729.5\\
72000	8720.8\\
74000	8585\\
76000	8390.6\\
78000	8778.5\\
80000	8884.9\\
82000	8788.6\\
84000	8301\\
86000	8481.8\\
88000	8733.5\\
90000	8635.2\\
92000	8713.4\\
94000	8420.8\\
96000	8652.4\\
98000	8826.5\\
1e+05	8639.4\\
};
\addplot [color=mycolor12, forget plot]
  table[row sep=crcr]{%
0	4395.3\\
2000	14301\\
4000	15223\\
6000	14662\\
8000	14817\\
10000	14929\\
12000	14979\\
14000	14271\\
16000	15555\\
20000	14792\\
22000	14675\\
24000	14818\\
26000	14486\\
28000	14710\\
30000	15032\\
32000	14829\\
34000	15052\\
36000	15407\\
38000	14879\\
40000	15493\\
42000	14946\\
44000	14608\\
46000	15017\\
48000	14521\\
50000	15133\\
52000	15239\\
54000	15124\\
56000	14318\\
58000	15216\\
60000	14715\\
62000	14968\\
64000	14934\\
66000	15107\\
68000	14970\\
70000	15108\\
72000	14932\\
74000	14937\\
76000	14800\\
78000	14604\\
80000	15230\\
82000	14970\\
84000	14654\\
86000	15042\\
88000	14537\\
90000	14919\\
92000	14909\\
94000	14627\\
96000	14681\\
98000	15190\\
1e+05	15090\\
};
\addplot [color=mycolor13, forget plot]
  table[row sep=crcr]{%
0	4395.3\\
2000	10797\\
4000	10997\\
6000	11052\\
8000	11365\\
10000	11218\\
12000	11148\\
14000	11374\\
16000	11113\\
18000	11091\\
20000	11326\\
22000	11340\\
24000	11238\\
26000	11274\\
28000	11432\\
30000	11060\\
32000	11231\\
36000	11095\\
38000	11331\\
40000	10901\\
42000	11537\\
44000	10988\\
46000	10973\\
48000	11027\\
50000	11271\\
52000	11108\\
54000	11125\\
56000	11383\\
60000	11171\\
64000	11659\\
66000	11025\\
68000	10948\\
70000	11353\\
72000	10979\\
74000	11319\\
76000	10986\\
78000	10882\\
80000	11282\\
82000	11233\\
84000	11311\\
86000	11076\\
88000	11313\\
90000	11046\\
92000	11437\\
94000	11085\\
96000	11629\\
98000	10946\\
1e+05	11337\\
};
\addplot [color=mycolor14, forget plot]
  table[row sep=crcr]{%
0	4395.3\\
2000	22508\\
4000	24782\\
6000	24182\\
8000	24130\\
10000	24576\\
12000	25393\\
14000	24424\\
16000	23573\\
18000	24949\\
20000	25124\\
22000	24842\\
24000	25257\\
26000	24324\\
28000	24671\\
30000	24877\\
32000	24636\\
34000	25573\\
36000	24762\\
38000	24844\\
40000	24471\\
42000	24214\\
44000	25433\\
46000	24409\\
48000	23899\\
50000	24859\\
52000	25280\\
54000	25012\\
56000	24360\\
58000	24682\\
60000	24680\\
62000	24372\\
64000	24721\\
66000	25002\\
68000	24876\\
70000	25280\\
72000	25020\\
74000	25118\\
76000	24281\\
78000	24488\\
80000	24487\\
82000	23894\\
84000	25159\\
86000	24642\\
88000	24380\\
90000	25544\\
92000	24896\\
94000	24150\\
96000	25471\\
98000	24114\\
1e+05	24423\\
};
\addplot [color=mycolor15, forget plot]
  table[row sep=crcr]{%
0	4395.3\\
2000	15973\\
4000	16806\\
6000	17369\\
8000	18133\\
10000	17728\\
12000	17518\\
14000	17240\\
16000	16883\\
18000	17146\\
20000	17548\\
22000	17342\\
24000	17996\\
26000	16978\\
28000	17461\\
30000	16976\\
32000	17159\\
34000	16487\\
36000	16935\\
38000	16823\\
40000	17169\\
42000	16412\\
44000	17570\\
46000	17283\\
48000	17236\\
50000	16335\\
52000	17306\\
54000	16746\\
56000	16811\\
58000	17062\\
60000	17076\\
62000	16940\\
64000	17512\\
66000	17801\\
68000	16650\\
70000	17398\\
72000	16805\\
74000	17259\\
76000	17016\\
78000	17287\\
80000	17508\\
82000	17365\\
84000	16967\\
86000	16881\\
88000	17656\\
90000	16962\\
92000	17268\\
94000	16930\\
96000	16910\\
98000	16656\\
1e+05	17303\\
};
\addplot [color=mycolor16, forget plot]
  table[row sep=crcr]{%
0	4395.3\\
2000	7113.6\\
4000	6987.2\\
6000	6987.8\\
8000	6789.1\\
10000	6949\\
12000	7072.3\\
14000	6938.9\\
18000	6916.1\\
20000	7196.3\\
22000	6914.9\\
24000	6810\\
26000	7038.6\\
28000	6848.4\\
30000	6920.8\\
32000	7273.6\\
34000	6927.6\\
36000	6721.3\\
38000	7003.8\\
40000	7015.3\\
42000	7323.9\\
44000	7091.9\\
46000	7017.5\\
48000	7131.8\\
50000	7139.8\\
52000	6742.3\\
54000	6794.4\\
56000	6950.6\\
58000	6848.3\\
60000	7039.2\\
62000	6868.5\\
64000	6811.4\\
66000	6789.5\\
68000	7224\\
70000	6952.4\\
72000	7030.7\\
74000	7088.8\\
76000	6959.1\\
78000	6906.6\\
80000	7274.9\\
82000	7038.8\\
84000	7008.9\\
86000	7206.4\\
88000	7004.5\\
90000	6908.6\\
92000	7077.9\\
94000	7030.3\\
96000	7007\\
98000	7005.9\\
1e+05	7032.1\\
};
\addplot [color=mycolor17, forget plot]
  table[row sep=crcr]{%
0	4395.3\\
2000	15368\\
4000	16320\\
6000	16446\\
8000	16048\\
10000	16222\\
12000	16228\\
14000	15825\\
16000	16249\\
18000	16273\\
20000	15832\\
22000	16583\\
24000	15613\\
28000	16057\\
30000	16427\\
32000	16087\\
34000	16446\\
36000	16272\\
38000	16477\\
40000	16600\\
42000	16044\\
44000	15830\\
46000	16404\\
48000	15767\\
50000	15438\\
52000	16033\\
54000	16142\\
56000	16097\\
58000	16184\\
60000	16372\\
62000	16079\\
64000	16117\\
66000	16851\\
68000	16516\\
70000	16527\\
72000	16198\\
74000	16469\\
76000	16218\\
78000	15350\\
80000	16665\\
82000	16256\\
86000	16176\\
88000	16317\\
90000	16563\\
92000	16571\\
94000	16459\\
96000	16472\\
98000	15901\\
1e+05	16502\\
};
\addplot [color=mycolor18, forget plot]
  table[row sep=crcr]{%
0	4395.3\\
2000	22736\\
4000	24789\\
6000	24698\\
8000	24794\\
10000	24024\\
12000	25973\\
14000	25264\\
16000	25511\\
18000	25123\\
20000	25942\\
22000	24791\\
24000	25363\\
26000	25353\\
30000	24459\\
32000	25410\\
34000	24107\\
36000	25227\\
38000	24518\\
40000	25424\\
42000	25696\\
44000	25211\\
46000	24892\\
48000	25000\\
50000	24646\\
52000	25304\\
54000	25023\\
56000	25598\\
58000	24967\\
60000	24647\\
62000	24794\\
64000	23989\\
66000	25278\\
68000	25245\\
70000	24764\\
72000	24473\\
74000	24991\\
76000	25138\\
78000	24831\\
80000	25094\\
82000	25593\\
84000	24343\\
86000	24792\\
88000	25880\\
90000	24775\\
92000	25374\\
94000	26138\\
96000	24439\\
98000	24109\\
1e+05	25152\\
};
\addplot [color=mycolor19, forget plot]
  table[row sep=crcr]{%
0	4395.3\\
2000	14170\\
4000	14752\\
6000	14480\\
8000	15191\\
10000	15122\\
12000	14848\\
14000	15284\\
16000	15340\\
18000	15320\\
20000	15345\\
22000	15450\\
24000	15414\\
26000	14976\\
28000	14763\\
30000	14418\\
32000	15377\\
34000	15060\\
36000	14526\\
38000	15416\\
40000	15088\\
42000	14553\\
44000	14863\\
46000	14775\\
48000	15347\\
50000	14716\\
52000	14935\\
54000	14714\\
56000	15491\\
58000	15446\\
60000	15099\\
62000	14998\\
64000	15395\\
66000	14932\\
68000	15114\\
70000	15390\\
72000	14738\\
74000	15144\\
76000	15082\\
78000	14820\\
80000	14821\\
82000	15113\\
84000	15250\\
86000	15467\\
88000	15245\\
90000	15792\\
92000	15480\\
94000	15314\\
96000	14909\\
98000	14593\\
1e+05	15236\\
};
\addplot [color=mycolor20, forget plot]
  table[row sep=crcr]{%
0	4395.3\\
2000	17742\\
4000	19716\\
6000	19872\\
8000	18864\\
10000	19255\\
12000	19360\\
14000	20113\\
16000	19851\\
18000	19953\\
22000	19517\\
24000	19386\\
26000	18931\\
28000	19321\\
30000	19436\\
32000	19755\\
34000	19472\\
36000	19747\\
38000	19717\\
40000	19358\\
42000	19616\\
44000	20398\\
46000	18549\\
48000	19705\\
50000	19395\\
52000	19819\\
54000	19150\\
56000	19363\\
58000	19503\\
60000	19397\\
62000	19629\\
64000	19238\\
66000	18948\\
68000	19052\\
70000	19850\\
72000	19625\\
74000	18640\\
76000	19658\\
78000	18776\\
80000	20058\\
82000	19804\\
84000	19465\\
86000	19291\\
88000	19638\\
90000	19607\\
92000	19176\\
94000	19091\\
96000	20417\\
98000	18982\\
1e+05	19226\\
};
\addplot [color=mycolor21, forget plot]
  table[row sep=crcr]{%
0	4395.3\\
2000	9289.6\\
4000	9585.1\\
6000	9761.6\\
8000	9742.6\\
10000	9815\\
12000	9617.7\\
14000	9709.1\\
16000	9829.8\\
18000	9821.6\\
20000	9727.1\\
22000	9874.1\\
24000	9803.9\\
26000	9945.5\\
28000	9640.9\\
30000	9782.9\\
32000	9709.9\\
36000	9482.9\\
38000	9634.8\\
40000	9532.3\\
42000	9560.2\\
44000	9731.9\\
46000	9713.5\\
50000	9986.1\\
52000	9844.2\\
54000	9841.3\\
56000	9608.1\\
58000	9793.5\\
60000	9461.3\\
62000	9832.8\\
68000	9688.6\\
70000	9690.5\\
72000	9521.6\\
74000	9737\\
76000	9837\\
82000	9513.6\\
86000	9720.8\\
88000	9867.4\\
90000	9575.7\\
92000	9691.6\\
94000	9626.2\\
96000	9972.1\\
1e+05	9721.8\\
};
\addplot [color=mycolor22, forget plot]
  table[row sep=crcr]{%
0	4395.3\\
2000	22993\\
4000	25362\\
6000	26194\\
8000	26497\\
10000	26721\\
12000	24536\\
14000	25889\\
16000	26134\\
18000	25310\\
20000	25543\\
22000	26032\\
24000	25592\\
26000	25695\\
28000	25310\\
30000	26048\\
32000	25114\\
34000	25582\\
36000	26230\\
38000	26081\\
40000	26431\\
42000	26107\\
44000	25630\\
46000	25724\\
48000	27544\\
50000	26381\\
52000	26912\\
54000	26277\\
56000	25135\\
58000	25702\\
60000	26656\\
62000	26444\\
64000	25564\\
66000	25022\\
68000	26016\\
70000	25929\\
74000	26164\\
76000	25298\\
78000	26672\\
80000	26426\\
82000	26427\\
84000	25988\\
86000	25786\\
90000	26198\\
92000	26127\\
94000	27673\\
96000	26110\\
98000	26206\\
1e+05	25941\\
};
\addplot [color=mycolor23, forget plot]
  table[row sep=crcr]{%
0	4395.3\\
2000	11636\\
4000	12429\\
6000	12384\\
10000	12583\\
12000	12115\\
14000	12373\\
16000	12550\\
18000	12169\\
20000	12636\\
22000	12555\\
24000	12404\\
26000	12595\\
28000	12350\\
30000	12565\\
32000	12234\\
34000	12450\\
36000	12352\\
38000	12401\\
40000	12604\\
42000	12400\\
44000	12258\\
46000	12552\\
48000	12272\\
50000	12041\\
52000	12467\\
54000	12391\\
56000	12550\\
58000	12375\\
60000	12397\\
62000	12821\\
64000	12426\\
66000	12386\\
68000	12618\\
70000	12545\\
72000	12287\\
74000	12316\\
76000	12564\\
78000	12156\\
80000	12434\\
82000	12516\\
84000	12524\\
86000	12691\\
88000	12471\\
90000	12307\\
92000	12541\\
94000	12589\\
96000	12108\\
98000	12550\\
1e+05	12614\\
};
\addplot [color=mycolor24, forget plot]
  table[row sep=crcr]{%
0	4395.3\\
2000	4866.2\\
4000	4697.7\\
6000	4602.9\\
8000	4813\\
10000	4924.3\\
12000	4956\\
14000	4720\\
16000	4570.5\\
18000	4636.4\\
20000	4806.7\\
22000	4785.2\\
24000	4843.9\\
26000	4660.3\\
28000	4836.1\\
30000	4741.4\\
32000	4716.6\\
34000	4739.8\\
36000	4745.6\\
38000	4631.5\\
40000	4846.5\\
42000	4667.1\\
44000	4654\\
46000	4535.3\\
48000	4656.7\\
50000	4839.1\\
54000	4827.3\\
56000	4759.2\\
58000	4758.8\\
60000	4723\\
62000	4853.4\\
64000	4575.1\\
66000	4766.6\\
68000	4551.9\\
70000	4729.5\\
72000	5025.4\\
74000	4742\\
76000	4786.8\\
78000	4672.3\\
80000	4619.9\\
82000	4744.6\\
84000	4638.2\\
86000	4862.9\\
88000	4771.5\\
90000	4700.4\\
92000	4825.3\\
94000	4797.2\\
96000	4722.2\\
1e+05	4656.8\\
};
\addplot [color=mycolor25, forget plot]
  table[row sep=crcr]{%
0	4395.3\\
2000	11887\\
4000	12761\\
6000	12677\\
8000	12003\\
10000	12349\\
12000	12293\\
14000	13258\\
16000	12998\\
18000	13054\\
20000	12278\\
22000	12490\\
24000	12455\\
26000	12519\\
28000	12924\\
32000	12583\\
34000	12827\\
36000	12271\\
38000	11888\\
40000	11742\\
42000	12518\\
44000	12236\\
46000	12420\\
48000	12172\\
50000	12599\\
52000	12456\\
54000	12353\\
56000	12705\\
58000	12766\\
60000	12511\\
62000	12312\\
64000	12449\\
66000	12404\\
68000	12509\\
72000	12817\\
74000	12904\\
76000	12654\\
78000	11768\\
80000	12573\\
82000	12584\\
84000	12848\\
86000	12113\\
88000	12593\\
90000	12390\\
92000	12712\\
94000	12463\\
96000	12678\\
98000	12141\\
1e+05	12246\\
};
\addplot [color=mycolor26, forget plot]
  table[row sep=crcr]{%
0	4395.3\\
2000	7978.5\\
4000	8535.1\\
6000	8135.4\\
8000	8321.4\\
10000	8213\\
12000	8462.9\\
14000	8392.5\\
16000	8486.3\\
18000	8265.6\\
20000	8353.7\\
22000	8171.7\\
24000	8438.2\\
26000	8013.4\\
28000	8152.4\\
30000	8119.3\\
32000	8566.1\\
34000	8417.1\\
36000	8173.6\\
38000	8345.6\\
40000	8070.9\\
42000	8160.5\\
44000	8196.1\\
46000	8119.1\\
48000	8236.1\\
50000	8384.6\\
52000	8049.2\\
54000	8371.9\\
56000	8254\\
58000	8188.5\\
60000	8267\\
62000	8299.1\\
64000	8303.2\\
66000	8130.1\\
68000	8326.6\\
70000	8383.3\\
72000	8291.1\\
74000	8358.8\\
76000	8022.6\\
78000	8439.7\\
80000	7958.6\\
82000	8293.4\\
84000	8009.3\\
86000	8533\\
88000	8172.9\\
90000	8578.7\\
92000	7859.3\\
96000	8432.9\\
98000	8367.3\\
1e+05	8212.7\\
};
\addplot [color=mycolor27, forget plot]
  table[row sep=crcr]{%
0	4395.3\\
2000	9376.3\\
4000	9643\\
6000	9881.3\\
8000	9926\\
10000	9619.8\\
12000	9899\\
14000	9534.4\\
16000	9869.3\\
18000	10014\\
20000	9625.3\\
22000	10033\\
24000	9776.9\\
26000	9861.3\\
28000	9618.7\\
30000	10038\\
32000	9829.2\\
34000	10143\\
36000	10167\\
38000	9704.2\\
40000	9821.1\\
44000	9846.1\\
46000	9945\\
48000	9502.9\\
50000	10014\\
52000	9873.9\\
54000	9799.8\\
56000	9407\\
58000	9921.4\\
60000	9595.4\\
62000	9701.8\\
64000	9585.9\\
66000	9779.1\\
68000	9654.2\\
70000	9822.6\\
72000	10266\\
74000	9789.9\\
76000	9749.8\\
78000	10224\\
80000	9938.9\\
82000	10085\\
84000	9583.2\\
86000	9911.7\\
88000	10056\\
90000	9757.5\\
92000	9406\\
94000	9776\\
96000	9704.6\\
98000	9921.8\\
1e+05	9697.2\\
};
\addplot [color=mycolor28, forget plot]
  table[row sep=crcr]{%
0	4395.3\\
2000	5413.3\\
4000	5565.3\\
6000	5701.6\\
8000	5280.2\\
10000	5308.6\\
12000	5464.4\\
14000	5264.6\\
16000	5348.8\\
18000	5342.4\\
20000	5463\\
22000	5204.1\\
24000	5357.4\\
26000	5446.5\\
28000	5510.4\\
30000	5411.1\\
32000	5491.7\\
34000	5298.8\\
36000	5470.3\\
38000	5297.8\\
40000	5479.6\\
42000	5464\\
44000	5320\\
46000	5378.7\\
48000	5403.2\\
50000	5474.2\\
52000	5502.7\\
54000	5359.9\\
56000	5413.3\\
58000	5406.9\\
60000	5550.2\\
62000	5504.9\\
64000	5640.4\\
66000	5548.5\\
68000	5359.3\\
70000	5420.9\\
72000	5587.7\\
74000	5404\\
76000	5638.9\\
78000	5188.9\\
80000	5132.4\\
82000	5431.6\\
84000	5435.6\\
86000	5263.8\\
88000	5572.1\\
90000	5490.9\\
92000	5234.7\\
94000	5542.6\\
96000	5354.7\\
98000	5463.1\\
1e+05	5626.7\\
};
\addplot [color=mycolor29, forget plot]
  table[row sep=crcr]{%
0	4395.3\\
2000	11046\\
4000	11909\\
6000	11976\\
8000	11490\\
10000	12128\\
14000	11434\\
16000	11828\\
18000	11727\\
20000	12083\\
22000	11933\\
24000	11968\\
26000	11483\\
28000	11697\\
30000	11688\\
32000	11924\\
34000	11855\\
36000	12038\\
38000	11611\\
42000	11686\\
44000	11843\\
46000	11363\\
48000	11866\\
50000	12094\\
52000	11644\\
54000	12126\\
56000	11947\\
58000	11704\\
60000	12159\\
62000	11984\\
64000	11958\\
68000	11624\\
70000	11596\\
72000	11836\\
74000	11419\\
76000	11598\\
78000	11910\\
80000	11535\\
82000	11673\\
84000	11992\\
86000	12278\\
88000	12033\\
90000	11873\\
92000	11571\\
94000	11461\\
96000	12044\\
98000	11486\\
1e+05	11253\\
};
\addplot [color=mycolor30, forget plot]
  table[row sep=crcr]{%
0	4395.3\\
2000	7847.5\\
4000	7944.7\\
6000	7905.4\\
8000	7812.2\\
10000	7791.9\\
12000	8051.2\\
16000	7913.1\\
18000	8166\\
20000	7931.5\\
22000	8008.8\\
24000	8113.3\\
26000	7872.7\\
28000	7706.3\\
30000	8102.5\\
34000	7682.5\\
36000	7827.8\\
38000	7838.3\\
40000	8007\\
42000	7707.4\\
44000	7964.5\\
46000	7973.9\\
48000	7946.9\\
50000	7853.7\\
52000	7785.7\\
54000	8055.5\\
56000	7873.7\\
58000	7786.3\\
62000	7936.8\\
64000	7788.5\\
66000	8014.9\\
68000	8057\\
70000	7865\\
72000	8050.1\\
74000	7800.9\\
76000	7920.8\\
82000	7784.7\\
84000	7708.8\\
86000	8185.2\\
88000	7988\\
90000	7913.9\\
92000	7999.9\\
94000	7654.4\\
96000	8132.7\\
98000	7886.1\\
1e+05	8014.5\\
};
\addplot [color=mycolor31, forget plot]
  table[row sep=crcr]{%
0	4395.3\\
2000	6432.2\\
4000	6392.3\\
6000	6564.8\\
8000	6388.9\\
10000	6272.8\\
12000	6696\\
14000	6540.2\\
16000	6635.5\\
18000	6326.9\\
20000	6340.8\\
22000	6607.1\\
24000	6555.8\\
26000	6455.2\\
28000	6323.3\\
30000	6247.7\\
32000	6527.2\\
34000	6666.9\\
36000	6215.6\\
38000	6421.2\\
40000	6389.6\\
42000	6551.9\\
44000	6322\\
46000	6380.3\\
48000	6287.4\\
50000	6270\\
52000	6560.3\\
54000	6423.7\\
56000	6463.7\\
58000	6320.8\\
60000	6267\\
62000	6667.5\\
64000	6433.4\\
66000	6306.5\\
68000	6752.9\\
70000	6402.2\\
72000	6702.6\\
74000	6357.8\\
76000	6340.2\\
78000	6584.5\\
80000	6676.7\\
82000	6342.6\\
84000	6293.4\\
86000	6638.4\\
88000	6339.7\\
90000	6674.1\\
92000	6714.3\\
94000	6449\\
96000	6294.5\\
98000	6518.3\\
1e+05	6344\\
};
\addplot [color=mycolor32, forget plot]
  table[row sep=crcr]{%
0	4395.3\\
2000	10254\\
4000	10790\\
6000	10503\\
8000	10784\\
10000	10732\\
12000	10363\\
14000	10743\\
16000	10795\\
18000	10624\\
20000	10353\\
22000	10928\\
24000	10439\\
26000	10617\\
28000	10394\\
30000	10647\\
32000	10554\\
34000	10813\\
38000	10607\\
40000	10319\\
42000	10725\\
44000	10737\\
46000	10388\\
48000	10458\\
50000	10689\\
52000	10549\\
54000	10355\\
56000	10701\\
58000	10819\\
60000	10720\\
62000	10321\\
68000	10912\\
70000	11096\\
72000	10429\\
74000	10255\\
76000	10358\\
78000	10764\\
80000	10729\\
82000	10078\\
84000	10278\\
86000	10929\\
88000	10450\\
90000	10695\\
92000	10314\\
94000	10306\\
96000	10392\\
1e+05	10782\\
};
\addplot [color=mycolor33, forget plot]
  table[row sep=crcr]{%
0	4395.3\\
2000	7895.4\\
4000	8012.8\\
10000	8233.8\\
12000	7908.6\\
14000	7948.8\\
16000	8430.4\\
18000	8564.7\\
20000	8000.1\\
22000	8240.3\\
24000	8112\\
26000	8487.2\\
28000	8385.3\\
30000	8202.4\\
32000	8297.4\\
34000	8293.4\\
36000	8701.7\\
38000	8392.4\\
40000	8303.6\\
42000	8278\\
44000	8163.5\\
48000	8444\\
50000	8586.4\\
52000	8770.3\\
54000	8132.5\\
56000	8557.7\\
58000	8089.5\\
60000	8292.7\\
62000	7985\\
64000	8395.5\\
66000	8365.6\\
68000	8063.2\\
70000	8508.2\\
72000	8105.5\\
74000	8121.3\\
76000	8238.6\\
78000	8333.3\\
80000	8132.3\\
84000	8271.5\\
86000	8280.3\\
88000	8551\\
90000	8361.3\\
92000	8213.4\\
94000	8114.4\\
96000	8201.7\\
98000	8157.4\\
1e+05	7973\\
};
\addplot [color=mycolor34, forget plot]
  table[row sep=crcr]{%
0	4395.3\\
2000	4057.5\\
4000	3945.9\\
6000	4121.5\\
8000	3941.4\\
10000	4053.1\\
12000	4050.6\\
14000	4091.5\\
16000	4111.1\\
18000	4082.7\\
20000	3836.5\\
22000	3947.2\\
24000	4013.7\\
26000	3876.3\\
28000	3926.2\\
30000	3932\\
32000	4044.5\\
34000	4058.6\\
36000	3894.4\\
38000	4038.7\\
40000	3833.1\\
42000	3936\\
44000	3970.5\\
46000	3931.2\\
48000	3938.8\\
50000	4019.2\\
52000	3825.6\\
54000	3893\\
56000	3985\\
58000	4018.7\\
60000	4000.6\\
62000	3903.5\\
64000	3981.7\\
66000	3978.9\\
68000	4126\\
70000	3891.7\\
72000	3939.9\\
74000	3876.3\\
76000	3877.2\\
78000	3891\\
80000	3878\\
82000	3926\\
84000	3907\\
86000	4002.6\\
88000	3925.1\\
90000	3961.5\\
92000	3764.3\\
94000	3937.7\\
96000	4017.3\\
98000	3984.4\\
1e+05	3964\\
};
\addplot [color=mycolor35]
  table[row sep=crcr]{%
0	4395.3\\
2000	9796.1\\
4000	10464\\
6000	10090\\
8000	10031\\
10000	10092\\
12000	10088\\
14000	10253\\
16000	10458\\
18000	10373\\
20000	10341\\
22000	10211\\
24000	10124\\
26000	10219\\
28000	10154\\
30000	10353\\
32000	10477\\
34000	10329\\
36000	10609\\
38000	10210\\
40000	10231\\
42000	10195\\
44000	10204\\
48000	10173\\
52000	10252\\
54000	10332\\
56000	10051\\
58000	10305\\
60000	10351\\
62000	10080\\
64000	10025\\
66000	10283\\
68000	10145\\
70000	10126\\
72000	10424\\
78000	10243\\
80000	10418\\
82000	10367\\
84000	10201\\
86000	10198\\
88000	9919.5\\
90000	10170\\
92000	10305\\
94000	10202\\
96000	10329\\
98000	10216\\
1e+05	10259\\
};
\addlegendentry{\ClasskIILegend}

\addplot [color=mycolor36, forget plot]
  table[row sep=crcr]{%
0	4395.3\\
2000	4430.5\\
4000	4249.3\\
6000	4261.9\\
8000	4313.1\\
10000	4095.7\\
12000	4286.8\\
14000	4333.1\\
16000	4129.4\\
18000	4438.4\\
20000	4052.8\\
22000	4315.9\\
24000	4110\\
26000	4149.9\\
30000	4165.5\\
32000	4267.5\\
34000	4227.5\\
36000	4089.8\\
38000	4237.2\\
40000	4240.1\\
42000	4343.8\\
44000	4190.9\\
46000	4124.9\\
48000	4144.1\\
50000	4296.6\\
54000	4300.7\\
56000	4235.3\\
58000	4190.4\\
60000	4344.5\\
62000	4316.2\\
64000	4174.7\\
66000	4299.7\\
68000	4047.4\\
70000	4048.5\\
72000	4226.7\\
74000	4277.5\\
76000	4015.8\\
78000	4265.1\\
80000	4125\\
82000	4296.8\\
84000	4303.8\\
86000	4120.1\\
88000	4146.8\\
90000	4242.9\\
92000	4288.5\\
94000	4362.2\\
96000	4288.2\\
98000	4333.7\\
1e+05	4229.3\\
};
\addplot [color=mycolor37, forget plot]
  table[row sep=crcr]{%
0	4395.3\\
2000	3700.4\\
4000	3567.4\\
6000	3598.3\\
8000	3517.7\\
10000	3551.5\\
12000	3603.7\\
14000	3570.1\\
16000	3624\\
18000	3591.3\\
20000	3635.9\\
22000	3645.2\\
24000	3538\\
26000	3642.8\\
28000	3496.6\\
32000	3494.7\\
34000	3611.5\\
36000	3564.6\\
38000	3581.8\\
40000	3556.9\\
44000	3611.2\\
46000	3565.9\\
48000	3614.2\\
50000	3530.2\\
52000	3578.5\\
54000	3526.1\\
56000	3594.8\\
58000	3500.5\\
60000	3496.1\\
62000	3502.6\\
64000	3553.8\\
68000	3587\\
72000	3511.7\\
74000	3589.8\\
76000	3652.1\\
78000	3517.1\\
80000	3566.4\\
82000	3601.3\\
84000	3604.4\\
86000	3623.3\\
90000	3523.5\\
92000	3560\\
94000	3541.3\\
96000	3621.4\\
98000	3527.7\\
1e+05	3526.3\\
};
\addplot [color=mycolor38, forget plot]
  table[row sep=crcr]{%
0	4395.3\\
2000	2880.6\\
4000	2905.7\\
6000	2964.5\\
8000	3007.6\\
10000	2949.5\\
12000	2938.4\\
14000	2824.5\\
16000	2970.7\\
18000	2919.3\\
20000	2918.3\\
22000	2973.2\\
24000	3000.9\\
26000	2932.4\\
28000	3026.7\\
30000	3013.8\\
32000	2864.1\\
34000	2996.3\\
36000	2962.3\\
38000	2908.2\\
40000	2863.9\\
42000	2951.9\\
44000	2925.9\\
46000	2932.3\\
48000	2994.6\\
50000	2995.7\\
52000	2946.3\\
54000	2967.3\\
56000	2913.3\\
58000	2920.4\\
60000	2999.6\\
62000	3011.7\\
64000	2975.8\\
66000	2851.7\\
68000	3018.7\\
70000	2868.8\\
72000	2931\\
74000	2903.9\\
76000	2919.1\\
78000	2775.5\\
80000	2846.1\\
82000	3035.8\\
84000	2938.8\\
86000	2862.5\\
88000	2828.3\\
90000	2946.7\\
92000	2997.8\\
94000	2882.4\\
96000	2964.2\\
98000	2941.2\\
1e+05	2990.5\\
};
\addplot [color=mycolor39, forget plot]
  table[row sep=crcr]{%
0	4395.3\\
2000	11111\\
4000	11743\\
6000	11726\\
8000	11881\\
10000	11745\\
12000	11557\\
14000	12239\\
16000	11140\\
18000	11704\\
20000	11783\\
22000	12237\\
24000	11991\\
26000	12126\\
28000	11988\\
30000	11268\\
32000	11493\\
34000	11911\\
36000	11695\\
38000	11663\\
40000	11715\\
42000	11424\\
44000	11771\\
46000	11836\\
48000	11702\\
50000	11686\\
52000	11591\\
54000	12298\\
56000	11965\\
58000	11786\\
60000	11868\\
62000	11750\\
64000	12262\\
66000	11493\\
68000	12157\\
72000	11750\\
74000	12229\\
76000	11507\\
78000	11738\\
80000	12056\\
82000	11634\\
84000	11702\\
86000	11733\\
88000	11925\\
90000	11722\\
92000	11467\\
96000	11677\\
98000	11928\\
1e+05	12457\\
};
\addplot [color=mycolor40, forget plot]
  table[row sep=crcr]{%
0	4395.3\\
2000	5811.7\\
4000	5547.3\\
6000	5644.9\\
8000	5661.7\\
10000	5645.2\\
12000	5697.6\\
14000	5915.3\\
16000	5609.6\\
18000	5738\\
20000	5707.2\\
22000	5718.3\\
24000	5770\\
26000	5653.9\\
28000	5681.6\\
30000	5803\\
34000	5654.5\\
36000	5922.1\\
38000	5659.2\\
42000	5786.1\\
44000	5791.6\\
48000	5711.1\\
50000	5815.8\\
52000	5665.4\\
54000	5643.5\\
56000	5726.1\\
58000	5828.2\\
60000	5638\\
62000	5823.8\\
64000	5725.3\\
66000	5797.3\\
68000	5761.1\\
72000	5832.3\\
74000	5704.5\\
76000	5684.4\\
78000	5784.2\\
80000	5620.3\\
82000	5698.1\\
84000	5678.3\\
86000	5769\\
88000	5767.3\\
90000	5707.5\\
92000	5623.3\\
94000	5797.4\\
96000	5881.6\\
98000	5877.8\\
1e+05	5672.5\\
};
\addplot [color=mycolor41, forget plot]
  table[row sep=crcr]{%
0	4395.3\\
2000	5529.8\\
4000	5463.8\\
8000	5409\\
10000	5503.9\\
12000	5426.6\\
14000	5407.5\\
16000	5557.8\\
18000	5476.9\\
20000	5322.7\\
22000	5266.3\\
24000	5437.1\\
26000	5409.5\\
28000	5446.6\\
30000	5200.1\\
32000	5354.2\\
34000	5326\\
36000	5501.6\\
38000	5301.6\\
40000	5419.6\\
42000	5501.2\\
44000	5473.1\\
46000	5491.1\\
48000	5244\\
50000	5353.8\\
52000	5345.9\\
54000	5521\\
56000	5347.5\\
58000	5263.4\\
60000	5368.7\\
62000	5433.4\\
64000	5475.7\\
66000	5229.1\\
68000	5360.7\\
70000	5446.5\\
72000	5407.6\\
74000	5474.4\\
76000	5211.5\\
78000	5420.7\\
82000	5328.1\\
84000	5295.8\\
86000	5446.7\\
88000	5191.4\\
90000	5442.8\\
92000	5470.5\\
94000	5301\\
96000	5353.8\\
98000	5622\\
1e+05	5361.6\\
};
\addplot [color=mycolor42, forget plot]
  table[row sep=crcr]{%
0	4395.3\\
2000	18116\\
4000	20924\\
6000	21844\\
8000	22079\\
10000	20335\\
12000	20979\\
14000	21155\\
16000	20878\\
18000	20980\\
20000	20683\\
22000	20192\\
24000	21018\\
26000	20907\\
28000	21133\\
30000	20603\\
32000	20605\\
34000	21126\\
36000	21086\\
38000	20574\\
40000	21238\\
42000	20213\\
44000	20631\\
46000	21257\\
48000	20333\\
50000	21990\\
52000	21619\\
54000	20824\\
56000	21250\\
58000	20990\\
60000	21206\\
62000	21226\\
64000	20939\\
66000	20202\\
68000	20420\\
70000	22029\\
72000	21201\\
74000	21500\\
76000	21094\\
78000	21644\\
80000	20350\\
82000	20425\\
84000	21606\\
86000	20839\\
90000	20037\\
92000	20709\\
94000	20690\\
96000	21182\\
1e+05	20127\\
};
\addplot [color=mycolor43, forget plot]
  table[row sep=crcr]{%
0	4395.3\\
2000	4925.8\\
4000	5119.3\\
6000	4847.4\\
8000	4876.6\\
10000	4866.4\\
12000	4956.3\\
14000	4886.2\\
16000	4847.8\\
18000	4973.1\\
20000	4971\\
22000	5041.5\\
24000	4973.3\\
26000	4984.2\\
28000	4898.2\\
30000	4764.7\\
32000	5026.4\\
34000	5020.7\\
36000	4874.9\\
38000	5069.5\\
40000	4760.6\\
42000	4858.2\\
44000	4977.2\\
46000	4864.6\\
48000	4916.6\\
50000	4902.8\\
52000	4913.9\\
54000	4875.9\\
56000	4879.5\\
58000	4940\\
60000	4893.6\\
62000	4992.1\\
66000	5069.7\\
68000	4841.5\\
72000	4889.5\\
74000	4790.9\\
76000	4958.9\\
78000	4911.9\\
80000	4882.5\\
82000	4769.8\\
84000	4737.5\\
86000	4940.1\\
90000	4951.7\\
92000	4761.3\\
94000	4921.3\\
96000	5002.7\\
98000	4861\\
1e+05	4872.4\\
};
\addplot [color=mycolor44, forget plot]
  table[row sep=crcr]{%
0	4395.3\\
2000	3460.8\\
4000	3238.1\\
6000	3285.3\\
8000	3245.1\\
10000	3418.1\\
12000	3283.3\\
14000	3347.4\\
16000	3301.2\\
18000	3234.8\\
20000	3157.4\\
22000	3319.3\\
24000	3143.5\\
26000	3169.3\\
28000	3310.4\\
32000	3335.5\\
34000	3268.1\\
36000	3354.6\\
38000	3262.2\\
40000	3203.7\\
42000	3125\\
44000	3302.2\\
46000	3252.8\\
48000	3275.8\\
50000	3356.6\\
52000	3191.4\\
54000	3240\\
56000	3191.4\\
58000	3344\\
60000	3289.9\\
62000	3302.3\\
66000	3228.3\\
68000	3219.5\\
70000	3110.1\\
72000	3162.1\\
74000	3285\\
76000	3306.7\\
78000	3223.5\\
80000	3255.6\\
82000	3246.6\\
84000	3175.5\\
86000	3229.4\\
88000	3294.2\\
90000	3221.9\\
92000	3072.7\\
94000	3389.6\\
96000	3383.6\\
98000	3170.3\\
1e+05	3275.1\\
};
\addplot [color=mycolor45, forget plot]
  table[row sep=crcr]{%
0	4395.3\\
2000	6759.1\\
4000	7102.1\\
6000	6831.8\\
8000	6957.9\\
10000	6971.6\\
16000	6892.4\\
18000	7243.4\\
20000	7232.8\\
22000	6899.3\\
24000	6959.3\\
26000	6735.5\\
28000	7053.8\\
30000	6841.4\\
32000	7060.6\\
34000	7112.2\\
36000	7117\\
38000	6929\\
40000	6805\\
42000	6925.9\\
44000	7123.8\\
46000	6997.9\\
48000	6929.6\\
50000	6722.2\\
52000	6989.2\\
54000	6974.6\\
56000	6824.2\\
58000	6861.8\\
60000	6775.6\\
62000	6923\\
64000	6912.1\\
68000	6788.3\\
70000	6862\\
72000	6956.7\\
74000	7177.5\\
76000	6791.9\\
78000	6763.8\\
82000	6982.2\\
84000	7151.6\\
88000	6874.2\\
90000	7012.6\\
92000	7036\\
94000	6857\\
96000	6914.5\\
98000	7074.9\\
1e+05	6980.7\\
};
\addplot [color=mycolor46, forget plot]
  table[row sep=crcr]{%
0	4395.3\\
2000	4133\\
4000	4022.7\\
6000	4130.9\\
8000	4100.1\\
10000	4048.5\\
12000	4012.4\\
14000	3998.5\\
16000	3942.9\\
18000	4029.5\\
20000	4018.2\\
22000	4124.2\\
24000	3941\\
26000	4009.6\\
28000	4006.4\\
30000	3963.9\\
34000	4080.3\\
36000	4118.9\\
38000	4016.7\\
42000	3988.1\\
44000	4010.2\\
46000	3992.5\\
48000	4034.1\\
50000	4168\\
54000	3950.7\\
56000	4036.1\\
60000	4055.8\\
62000	4096.7\\
64000	3998.1\\
66000	4056.8\\
68000	3877.6\\
70000	4105.9\\
72000	4018.7\\
74000	4162.1\\
76000	4012.6\\
80000	3895.5\\
82000	4105.6\\
84000	4148\\
86000	4159.9\\
88000	3944.3\\
90000	4004.8\\
92000	4027.7\\
94000	4175.2\\
96000	4070.9\\
98000	4015.8\\
1e+05	4054.6\\
};
\addplot [color=mycolor47, forget plot]
  table[row sep=crcr]{%
0	4395.3\\
2000	5277.9\\
4000	5261.4\\
8000	5291.3\\
10000	5220.5\\
12000	5121.4\\
14000	5188.7\\
16000	5115.6\\
18000	5275.3\\
20000	5327.2\\
22000	5179.4\\
24000	5367.4\\
26000	5202.2\\
28000	5259.3\\
30000	5270.4\\
32000	5160.5\\
34000	5242.9\\
36000	5300.6\\
40000	5299.6\\
42000	5282.6\\
44000	5132.4\\
46000	5387.6\\
48000	5251.7\\
50000	5394.2\\
52000	5094.9\\
54000	5245.2\\
56000	5338.1\\
58000	5133.3\\
60000	5205.5\\
62000	5171.2\\
64000	5316.7\\
66000	5356.4\\
68000	5121.4\\
70000	5140.6\\
74000	5264.6\\
76000	5237.7\\
78000	5324\\
80000	5220.6\\
82000	5301.5\\
84000	5213\\
86000	5289.7\\
88000	5411.7\\
90000	5391.3\\
92000	5318.3\\
94000	5151.7\\
96000	5186.9\\
98000	5146.6\\
1e+05	5381.6\\
};
\addplot [color=mycolor48, forget plot]
  table[row sep=crcr]{%
0	4395.3\\
2000	2985.1\\
4000	2865.3\\
6000	2842\\
8000	2800.3\\
10000	2889.7\\
12000	2835.1\\
14000	2851.8\\
16000	2910.4\\
18000	2846.9\\
20000	2810.1\\
24000	2827.1\\
26000	2834.7\\
28000	2813.9\\
30000	2838.6\\
32000	2847.7\\
34000	2871.7\\
36000	2810.1\\
38000	2873.5\\
40000	2878.5\\
42000	2779.1\\
44000	2814.9\\
46000	2913.5\\
48000	2851.2\\
50000	2904.2\\
52000	2839.9\\
54000	2871.1\\
56000	2823.1\\
58000	2864.6\\
60000	2866.5\\
62000	2855\\
66000	2817.9\\
68000	2866\\
70000	2773.3\\
72000	2876.7\\
74000	2823.8\\
76000	2853.3\\
78000	2809.3\\
80000	2864.8\\
82000	2827.2\\
84000	2807.3\\
86000	2817.9\\
88000	2849.8\\
90000	2838.3\\
92000	2856\\
94000	2933.2\\
96000	2864.3\\
98000	2845.1\\
1e+05	2895.1\\
};
\addplot [color=mycolor49, forget plot]
  table[row sep=crcr]{%
0	4395.3\\
4000	4297.6\\
6000	4044.2\\
8000	4233.6\\
10000	4249.5\\
12000	4158.6\\
14000	4365.5\\
16000	4275.5\\
18000	4421\\
20000	4155.1\\
22000	4274.3\\
24000	4334.7\\
26000	4442\\
28000	4371.7\\
30000	4360.5\\
32000	4444.2\\
34000	4441.7\\
36000	4326.7\\
38000	4236.5\\
40000	4123\\
42000	4337.8\\
44000	4192.9\\
46000	4323.6\\
48000	4233.2\\
50000	4118.4\\
52000	4303.8\\
54000	4340.9\\
56000	4206.9\\
58000	4429.4\\
60000	4325.9\\
62000	4527.5\\
64000	4394.7\\
66000	4316\\
68000	4222.5\\
70000	4195.5\\
72000	4349.1\\
74000	4380.8\\
76000	4226.4\\
78000	4378.2\\
80000	4320.5\\
82000	4371.9\\
84000	4286.3\\
86000	4278.5\\
88000	4320.3\\
90000	4376.6\\
92000	4122.5\\
94000	4161.8\\
96000	4368.3\\
98000	4290.4\\
1e+05	4327.4\\
};
\addplot [color=mycolor50, forget plot]
  table[row sep=crcr]{%
0	4395.3\\
2000	2567.5\\
4000	2370.3\\
6000	2382.9\\
8000	2418.5\\
10000	2386.5\\
12000	2375.4\\
14000	2410.1\\
16000	2335.8\\
18000	2471.7\\
20000	2394.3\\
22000	2357.4\\
26000	2417.1\\
28000	2369.5\\
30000	2428.4\\
32000	2510.6\\
34000	2437.9\\
36000	2458.5\\
38000	2366\\
40000	2394.8\\
42000	2300.5\\
44000	2420.7\\
46000	2442.2\\
48000	2439.1\\
50000	2356.8\\
52000	2437.1\\
54000	2363.4\\
56000	2436.1\\
58000	2410.3\\
60000	2418.9\\
62000	2326.9\\
64000	2411.1\\
66000	2348.6\\
72000	2418.1\\
74000	2379.9\\
76000	2397.8\\
78000	2372.9\\
80000	2380.1\\
82000	2372.5\\
84000	2474.8\\
86000	2428.7\\
88000	2405.7\\
90000	2418.4\\
92000	2405.7\\
94000	2373.7\\
96000	2430.1\\
98000	2399\\
1e+05	2386.6\\
};
\addplot [color=mycolor51, forget plot]
  table[row sep=crcr]{%
0	4395.3\\
2000	5828.5\\
4000	5564\\
6000	6066.9\\
8000	6091.9\\
10000	5818.6\\
12000	6130.8\\
14000	5786.3\\
16000	6299.3\\
18000	5998.4\\
20000	5953.9\\
22000	5939.2\\
24000	5874.5\\
26000	6126.6\\
28000	5809.1\\
30000	5702.2\\
32000	5937.8\\
34000	5949.8\\
36000	5979\\
38000	5884.3\\
40000	5875.8\\
42000	5909.5\\
44000	5742.1\\
46000	5820.6\\
48000	5833.3\\
50000	6107.2\\
52000	5882.5\\
54000	5994.2\\
56000	5841.7\\
58000	5888.6\\
60000	5720.5\\
62000	5903\\
64000	5877.3\\
66000	6004.3\\
68000	6031.3\\
70000	6143\\
72000	5981.1\\
74000	6300.1\\
76000	5878.1\\
78000	5936.1\\
80000	6110.9\\
82000	5928.5\\
84000	5792.9\\
86000	5836.8\\
88000	5712.2\\
90000	5731.1\\
92000	6280.9\\
94000	5910.6\\
96000	6021.2\\
98000	5979.3\\
1e+05	6235.3\\
};
\addplot [color=mycolor52, forget plot]
  table[row sep=crcr]{%
0	4395.3\\
2000	3651\\
4000	3446.3\\
6000	3570.4\\
8000	3594.9\\
10000	3389.5\\
12000	3374.6\\
14000	3432.3\\
16000	3501.3\\
18000	3520.9\\
20000	3421.4\\
22000	3486.4\\
26000	3379.3\\
28000	3402.8\\
30000	3468.5\\
32000	3301.4\\
34000	3359.4\\
36000	3393.7\\
38000	3518\\
40000	3565.3\\
42000	3486.2\\
44000	3591\\
48000	3452.7\\
50000	3270.7\\
52000	3417.6\\
54000	3450.1\\
56000	3415.2\\
58000	3318.7\\
60000	3540\\
62000	3337.1\\
64000	3507.6\\
66000	3513.3\\
68000	3403.9\\
70000	3345.5\\
72000	3298.5\\
74000	3337.8\\
76000	3605.6\\
78000	3452.4\\
80000	3428.6\\
82000	3341.3\\
84000	3411\\
86000	3671.5\\
88000	3669.7\\
90000	3437.5\\
92000	3463.3\\
94000	3385.4\\
96000	3378\\
98000	3548.5\\
1e+05	3395.3\\
};
\addplot [color=mycolor53, forget plot]
  table[row sep=crcr]{%
0	4395.3\\
2000	5572.6\\
4000	5593.4\\
6000	5567.2\\
8000	5942\\
10000	5680\\
12000	5635.3\\
14000	5493.8\\
16000	5608.6\\
18000	5453.2\\
20000	5506.3\\
22000	5700.3\\
24000	5654.8\\
26000	5573.4\\
28000	5652.1\\
30000	5811.7\\
32000	5572.2\\
34000	5658.5\\
36000	5977.8\\
38000	5798.2\\
40000	5765.4\\
42000	5545.5\\
44000	5755.9\\
46000	5716.1\\
48000	5727\\
50000	5479.7\\
52000	5769.2\\
54000	5527.5\\
56000	5855.1\\
58000	5652.3\\
60000	5616\\
64000	5625.3\\
66000	5712.7\\
68000	5835.4\\
70000	5719.2\\
72000	5662.4\\
74000	5755.1\\
76000	5661.1\\
78000	5641.2\\
80000	5244.8\\
82000	5878.8\\
84000	5659.3\\
86000	5688.9\\
92000	5651.9\\
94000	5679.4\\
96000	5604\\
98000	5402.3\\
1e+05	5733\\
};
\addplot [color=mycolor54, forget plot]
  table[row sep=crcr]{%
0	4395.3\\
2000	4709.9\\
4000	4439.5\\
6000	4404.6\\
8000	4533.1\\
10000	4555.5\\
12000	4660\\
16000	4530.5\\
18000	4708.9\\
20000	4512.3\\
22000	4565.2\\
24000	4462.8\\
26000	4681.1\\
28000	4490.3\\
30000	4685.6\\
32000	4509.7\\
34000	4485.9\\
36000	4485.2\\
38000	4671.8\\
40000	4462.3\\
42000	4532.3\\
44000	4407.6\\
46000	4502.2\\
48000	4511.4\\
50000	4547.6\\
52000	4413\\
54000	4625.6\\
56000	4480.1\\
58000	4560.1\\
60000	4522.4\\
62000	4757.5\\
66000	4562.7\\
68000	4572.2\\
70000	4429.4\\
72000	4574.3\\
74000	4597\\
76000	4545.4\\
78000	4684.6\\
80000	4671.7\\
82000	4576\\
86000	4694.1\\
88000	4553.9\\
90000	4671.7\\
92000	4658.4\\
94000	4389.8\\
96000	4557.7\\
98000	4692.5\\
1e+05	4636.4\\
};
\addplot [color=mycolor55, forget plot]
  table[row sep=crcr]{%
0	4395.3\\
2000	5665.7\\
4000	5722.6\\
6000	5671.5\\
8000	5704.9\\
10000	5850.3\\
14000	5807.8\\
16000	5720.7\\
18000	5844.1\\
20000	5615.1\\
22000	5739\\
24000	5610.2\\
26000	5704.8\\
28000	5916.9\\
30000	5748.5\\
32000	5844.5\\
34000	5856.2\\
36000	5767.5\\
38000	5920.5\\
40000	5907.1\\
42000	5808.5\\
44000	5820.4\\
46000	5808.6\\
48000	5731.1\\
50000	5850.8\\
52000	5706.1\\
54000	5832.3\\
56000	5912.4\\
58000	5741.3\\
62000	5858.6\\
64000	5869.3\\
66000	5567.5\\
68000	5758.7\\
70000	5679.7\\
72000	5756.6\\
74000	5806.3\\
76000	5776.4\\
78000	5828.1\\
80000	5800.7\\
82000	6026.2\\
84000	5866.6\\
86000	5841.5\\
88000	5650.8\\
90000	5726.2\\
92000	5759.8\\
94000	5777.1\\
98000	5716.3\\
1e+05	5770.8\\
};
\addplot [color=mycolor56, forget plot]
  table[row sep=crcr]{%
0	4395.3\\
2000	3993.3\\
4000	3919.5\\
10000	3991.9\\
14000	3880.5\\
16000	4060.8\\
18000	3974.3\\
20000	3996.8\\
22000	3896.9\\
24000	4000.5\\
26000	3924.2\\
28000	3860.3\\
30000	3883.2\\
32000	3917.8\\
34000	3928.8\\
36000	4040.8\\
38000	4038.7\\
40000	3845.7\\
42000	3967.4\\
44000	4001\\
46000	3990\\
48000	3860.8\\
50000	3874.8\\
52000	3855.1\\
54000	4042.5\\
56000	3983.2\\
58000	4005.5\\
60000	4092.4\\
62000	3899.1\\
64000	3998.4\\
66000	3965.7\\
68000	3857.6\\
70000	3969.4\\
72000	4107.8\\
74000	4043.6\\
76000	3941.1\\
78000	3856.8\\
80000	4058.5\\
82000	3942.6\\
84000	3843.6\\
86000	3842.9\\
88000	3945.7\\
90000	3968.1\\
92000	3867.1\\
94000	3974.9\\
96000	3882.8\\
98000	3975.9\\
1e+05	3996.5\\
};
\addplot [color=mycolor57, forget plot]
  table[row sep=crcr]{%
0	4395.3\\
2000	6500.2\\
4000	6760.2\\
6000	6858.6\\
8000	6511.4\\
10000	6850.5\\
12000	6825.3\\
14000	6999.4\\
16000	6780.7\\
18000	6689.8\\
22000	6728.9\\
24000	6631.1\\
26000	6568.3\\
28000	6591.7\\
30000	6713.9\\
32000	6871.3\\
34000	6677.5\\
36000	6728.7\\
38000	6607.6\\
40000	6725.7\\
42000	6754.7\\
44000	6754.9\\
46000	6839.7\\
48000	6804.7\\
50000	6593.2\\
52000	6711.8\\
54000	6538.6\\
56000	6826.4\\
58000	6673.7\\
62000	6830.6\\
64000	6490.6\\
66000	6715.7\\
70000	6723.2\\
72000	6906.3\\
74000	6617.5\\
76000	6885.7\\
78000	6847.8\\
80000	6830.9\\
82000	6643.2\\
86000	6731.7\\
88000	6628.3\\
90000	6792.3\\
92000	6639.6\\
94000	6648.4\\
96000	6754.5\\
98000	6642.2\\
1e+05	6753.2\\
};
\addplot [color=mycolor58, forget plot]
  table[row sep=crcr]{%
0	4395.3\\
2000	3617.9\\
4000	3453.3\\
6000	3517.6\\
8000	3450.1\\
10000	3498.4\\
12000	3531\\
14000	3522.5\\
16000	3412.1\\
18000	3604.9\\
20000	3571.9\\
22000	3495.1\\
24000	3616.4\\
26000	3418.6\\
28000	3439.7\\
30000	3546.5\\
32000	3442.5\\
34000	3627\\
36000	3572.4\\
38000	3477.5\\
40000	3500.5\\
42000	3562.5\\
44000	3457.2\\
46000	3574.4\\
48000	3478.5\\
50000	3442.6\\
52000	3527.7\\
54000	3522.6\\
56000	3439.8\\
58000	3542.5\\
60000	3504.1\\
62000	3538\\
64000	3587.4\\
66000	3516.9\\
68000	3622.1\\
70000	3450.4\\
72000	3629.2\\
74000	3605.2\\
76000	3468.1\\
78000	3616.4\\
80000	3518.7\\
82000	3399.2\\
84000	3527.5\\
86000	3502.7\\
88000	3521.9\\
90000	3474.1\\
92000	3554.8\\
94000	3398.7\\
96000	3549.6\\
98000	3502.8\\
1e+05	3513.6\\
};
\addplot [color=mycolor59, forget plot]
  table[row sep=crcr]{%
0	4395.3\\
2000	5011.4\\
4000	4960.3\\
6000	5178.4\\
8000	4978.3\\
10000	4985.4\\
12000	4818.6\\
14000	5114.3\\
16000	5008.7\\
18000	5066.4\\
20000	4994.7\\
22000	4983.3\\
24000	5091.8\\
26000	4866.8\\
28000	5173.7\\
30000	5113.7\\
32000	4888.1\\
34000	4859.6\\
36000	5126.2\\
38000	5061.3\\
40000	5032.9\\
42000	5081.7\\
44000	5098.4\\
46000	4896.6\\
48000	5040.6\\
50000	4916.5\\
52000	5005\\
54000	5013.4\\
56000	5124\\
58000	4941.3\\
60000	5034\\
62000	4960.6\\
64000	5009.8\\
68000	5034.8\\
70000	4953.2\\
72000	5106.8\\
74000	5149.5\\
76000	5037\\
78000	5085.9\\
80000	4974.1\\
82000	4901.9\\
84000	5031.1\\
86000	4862.3\\
88000	4985.4\\
90000	4999.6\\
92000	5104.9\\
94000	5058\\
96000	5142\\
98000	5080.9\\
1e+05	4980.5\\
};
\addplot [color=mycolor60, forget plot]
  table[row sep=crcr]{%
0	4395.3\\
2000	3549.7\\
4000	3525.5\\
6000	3457.8\\
8000	3596.1\\
10000	3510.7\\
12000	3519.6\\
14000	3466\\
16000	3475.7\\
18000	3523.6\\
20000	3355.1\\
22000	3477.6\\
24000	3453\\
26000	3522.8\\
28000	3509.8\\
30000	3450.2\\
32000	3484.5\\
34000	3552.8\\
36000	3554.1\\
38000	3614.9\\
40000	3498.3\\
44000	3551.1\\
46000	3475.6\\
48000	3467.6\\
50000	3523.5\\
52000	3430.4\\
54000	3469.6\\
56000	3474.5\\
58000	3506.6\\
62000	3473.2\\
64000	3435.7\\
66000	3497.3\\
68000	3448.6\\
70000	3524.6\\
72000	3554.3\\
74000	3560.5\\
76000	3579.7\\
78000	3536.7\\
80000	3540.7\\
82000	3386.2\\
84000	3411.5\\
86000	3408\\
88000	3533.6\\
90000	3412.7\\
92000	3466.5\\
94000	3572.7\\
96000	3469.2\\
98000	3495.2\\
1e+05	3452.4\\
};
\addplot [color=mycolor61, forget plot]
  table[row sep=crcr]{%
0	4395.3\\
2000	3203.7\\
4000	3103.8\\
6000	3115.8\\
8000	3089.2\\
10000	3007.1\\
12000	3054.5\\
14000	3063.7\\
16000	3030.1\\
18000	3065\\
20000	3071.2\\
22000	3064.1\\
24000	3073.3\\
26000	3098.7\\
28000	3065.3\\
30000	3001.6\\
32000	3032.2\\
34000	3091\\
36000	3105\\
38000	3102.5\\
40000	3066.3\\
42000	3044.4\\
44000	3130.9\\
46000	3068\\
48000	3153.6\\
50000	3070.5\\
52000	3082.6\\
54000	3070.3\\
56000	3045.8\\
58000	3131.4\\
60000	3054.4\\
62000	3022.7\\
64000	3104\\
66000	3104.4\\
70000	3054.5\\
72000	3085.6\\
74000	3077.8\\
76000	3048.2\\
78000	3049.1\\
80000	2997.8\\
82000	3040.5\\
84000	3113.4\\
86000	3057.5\\
88000	3030.1\\
90000	3023.1\\
92000	3069.9\\
94000	3012.1\\
98000	3074.5\\
1e+05	3141\\
};
\addplot [color=mycolor62, forget plot]
  table[row sep=crcr]{%
0	4395.3\\
2000	4942.3\\
4000	4886.9\\
6000	5036.8\\
8000	4915.4\\
10000	5040\\
12000	4974.9\\
14000	4929.5\\
18000	4965.9\\
20000	4909.6\\
22000	5002.8\\
24000	4896.2\\
26000	4932.7\\
28000	4952.1\\
30000	5020.1\\
32000	4812.4\\
34000	4925.3\\
36000	5150.6\\
38000	4943.8\\
40000	4984.8\\
42000	4968.3\\
44000	5010.1\\
46000	5018.6\\
48000	4928.3\\
50000	4971.9\\
52000	4866.1\\
54000	4861.6\\
56000	4962.7\\
58000	4950.6\\
60000	5078.6\\
62000	5024.2\\
64000	4954\\
66000	4918.7\\
68000	4899.5\\
70000	4850.1\\
74000	5034.1\\
76000	4976.5\\
78000	5142.7\\
80000	4809.9\\
82000	5095.9\\
84000	5035.1\\
86000	5028.6\\
88000	5039.4\\
90000	4963.1\\
92000	5119.1\\
94000	4880.9\\
96000	4952.6\\
98000	4959.7\\
1e+05	4943.1\\
};
\addplot [color=mycolor63, forget plot]
  table[row sep=crcr]{%
0	4395.3\\
2000	2510.1\\
4000	2280.5\\
6000	2360.4\\
8000	2351.3\\
10000	2334.2\\
12000	2351.6\\
14000	2337.3\\
16000	2334.5\\
18000	2318.4\\
20000	2338.7\\
22000	2384.8\\
26000	2376.4\\
28000	2317\\
30000	2382.3\\
32000	2337.3\\
34000	2410.8\\
36000	2325\\
38000	2347.1\\
40000	2358.1\\
44000	2342.3\\
46000	2272.8\\
48000	2330.1\\
50000	2354.6\\
52000	2323.5\\
56000	2312.5\\
58000	2272.3\\
60000	2368.8\\
62000	2334.3\\
64000	2351.1\\
66000	2341\\
68000	2399.9\\
70000	2353.2\\
72000	2393.4\\
74000	2278.4\\
76000	2329.5\\
80000	2340.3\\
82000	2357.5\\
84000	2314.5\\
88000	2304.3\\
90000	2326.1\\
94000	2315.6\\
96000	2272.8\\
98000	2338.2\\
1e+05	2305.3\\
};
\addplot [color=mycolor64, forget plot]
  table[row sep=crcr]{%
0	4395.3\\
2000	3132.3\\
4000	2991.2\\
6000	2919.9\\
8000	2960.8\\
10000	2943.7\\
12000	3044.1\\
14000	2983.7\\
16000	3004.4\\
18000	3053.3\\
20000	3062.8\\
22000	2993.8\\
24000	3020.2\\
28000	3010.4\\
30000	2970.7\\
32000	2972.9\\
34000	2987.7\\
36000	2992.7\\
38000	3022\\
40000	2994.7\\
42000	2929.1\\
44000	2966.9\\
46000	3055.1\\
48000	2951.4\\
50000	3049.9\\
52000	2973.2\\
54000	2985.9\\
56000	2972.8\\
58000	2980.8\\
60000	2961.1\\
62000	3018.4\\
64000	3042.8\\
66000	2973.3\\
68000	2943.9\\
70000	2964.8\\
72000	2994.3\\
74000	3079.5\\
76000	2981.6\\
78000	2925.4\\
80000	2966\\
82000	2970.2\\
86000	3013.7\\
88000	2876.1\\
90000	3015.2\\
94000	2989.9\\
96000	2969.7\\
98000	3064.9\\
1e+05	3002.6\\
};
\addplot [color=mycolor65, forget plot]
  table[row sep=crcr]{%
0	4395.3\\
2000	2637.5\\
4000	2434.6\\
6000	2437.8\\
8000	2487.9\\
10000	2506.4\\
12000	2553\\
14000	2493.3\\
16000	2510.6\\
18000	2452.5\\
20000	2453.7\\
22000	2516.5\\
24000	2430.8\\
28000	2477.5\\
30000	2464.3\\
32000	2481.2\\
34000	2552.7\\
36000	2375.5\\
38000	2455.9\\
40000	2468.4\\
42000	2525.6\\
44000	2430.2\\
48000	2478.8\\
50000	2440.9\\
52000	2485.9\\
54000	2467.9\\
56000	2524.8\\
58000	2467\\
60000	2435.6\\
64000	2485.9\\
66000	2455.1\\
68000	2449.9\\
70000	2467\\
72000	2421.1\\
74000	2544.4\\
76000	2446.6\\
78000	2477.8\\
80000	2444.4\\
82000	2472.9\\
84000	2527.3\\
86000	2458.1\\
88000	2445.7\\
90000	2412.2\\
92000	2495\\
94000	2521.6\\
96000	2527.6\\
98000	2492.8\\
1e+05	2513\\
};
\addplot [color=mycolor66, forget plot]
  table[row sep=crcr]{%
0	4395.3\\
2000	3501.5\\
4000	3317\\
8000	3364.3\\
12000	3360.2\\
14000	3368.8\\
16000	3346.2\\
18000	3443.2\\
20000	3396.2\\
22000	3398\\
24000	3372.1\\
26000	3408.6\\
28000	3341.1\\
30000	3339.9\\
32000	3324.5\\
36000	3364.7\\
38000	3448.8\\
40000	3287.7\\
42000	3377.9\\
44000	3406.9\\
46000	3410.5\\
48000	3377.4\\
50000	3404.9\\
52000	3344.7\\
54000	3391\\
56000	3381.4\\
58000	3313.2\\
60000	3382\\
62000	3255\\
64000	3325.1\\
66000	3428.8\\
68000	3380.1\\
70000	3371.6\\
72000	3352.9\\
74000	3348.7\\
76000	3367.5\\
78000	3367.5\\
82000	3393.4\\
84000	3343.7\\
86000	3399.3\\
88000	3354.5\\
90000	3374.9\\
92000	3383.7\\
94000	3325.9\\
96000	3295.1\\
98000	3395\\
1e+05	3419.8\\
};
\addplot [color=mycolor67, forget plot]
  table[row sep=crcr]{%
0	4395.3\\
2000	3279.3\\
4000	3205.5\\
6000	3179.3\\
8000	3080.4\\
10000	3058.2\\
12000	3094.7\\
14000	3060.4\\
16000	3123.2\\
18000	3223.9\\
20000	3161.6\\
22000	3133.8\\
24000	3128.9\\
26000	3054.8\\
28000	3116\\
30000	3111.1\\
32000	3077.7\\
34000	3137.1\\
36000	3131.9\\
38000	3149.8\\
40000	3129.7\\
42000	3135.3\\
44000	3128.3\\
46000	3103.1\\
48000	3184.5\\
50000	3104.2\\
52000	3073.9\\
54000	3073.8\\
56000	3010.9\\
58000	3115.4\\
60000	3197.9\\
62000	3145.6\\
64000	3103.8\\
66000	3146.5\\
68000	3179.8\\
70000	3148.5\\
72000	3096.7\\
74000	3071.9\\
76000	3068.1\\
78000	3054.2\\
80000	3120.1\\
82000	3158\\
84000	3155\\
86000	3083.5\\
88000	3162.8\\
90000	3195.2\\
92000	3079.1\\
94000	3137.4\\
96000	3121\\
98000	3117.6\\
1e+05	3156.7\\
};
\addplot [color=mycolor67]
  table[row sep=crcr]{%
0	4395.3\\
2000	2736.8\\
4000	2558.8\\
6000	2592.9\\
8000	2553.8\\
10000	2559.1\\
12000	2508.6\\
14000	2614.6\\
16000	2559.2\\
18000	2527.8\\
20000	2573.3\\
22000	2551.2\\
24000	2586.5\\
26000	2521.4\\
28000	2580.3\\
30000	2545.6\\
32000	2598.3\\
34000	2616.7\\
36000	2584.6\\
40000	2553.9\\
42000	2614.5\\
44000	2563\\
46000	2588.4\\
48000	2559.1\\
52000	2595.6\\
54000	2554.2\\
56000	2598.5\\
60000	2544.9\\
62000	2568.2\\
64000	2579.9\\
66000	2556.6\\
68000	2566.2\\
70000	2565.6\\
72000	2547.1\\
74000	2582.5\\
76000	2600.5\\
78000	2565.9\\
80000	2554.3\\
82000	2558.4\\
84000	2548.8\\
88000	2475.6\\
90000	2556\\
92000	2627.7\\
94000	2624.6\\
96000	2548.9\\
98000	2616.3\\
1e+05	2549.4\\
};
\addlegendentry{\ClasskIIILegend}

\addplot [color=mycolor68, forget plot]
  table[row sep=crcr]{%
0	4395.3\\
2000	2611.7\\
4000	2395.9\\
6000	2402.2\\
8000	2400.2\\
10000	2367.2\\
12000	2376\\
14000	2364.6\\
18000	2398.6\\
20000	2398.8\\
22000	2408.8\\
24000	2377\\
26000	2414.6\\
28000	2468.1\\
30000	2404.2\\
32000	2392.6\\
34000	2394.3\\
36000	2373.6\\
38000	2344.8\\
40000	2416.3\\
42000	2400.7\\
44000	2399.8\\
46000	2384\\
50000	2399.6\\
52000	2436.4\\
54000	2444.2\\
56000	2400.4\\
58000	2402.7\\
60000	2380\\
62000	2413\\
64000	2391.6\\
66000	2397.5\\
68000	2382.9\\
70000	2387\\
72000	2410\\
74000	2417.7\\
76000	2400.8\\
78000	2353.3\\
80000	2375.1\\
82000	2412.8\\
84000	2366.1\\
86000	2386.4\\
88000	2452\\
90000	2398.8\\
92000	2412\\
94000	2379.6\\
96000	2375.6\\
98000	2358.9\\
1e+05	2379.3\\
};
\addplot [color=mycolor69, forget plot]
  table[row sep=crcr]{%
0	4395.3\\
2000	3044\\
4000	2886\\
6000	2880\\
8000	2860.1\\
10000	2874.5\\
12000	2826.4\\
14000	2870.1\\
16000	2854.2\\
20000	2869.9\\
22000	2851.2\\
24000	2854.4\\
26000	2817\\
28000	2893.4\\
30000	2854.8\\
32000	2857.5\\
34000	2835.1\\
36000	2835.6\\
38000	2815.9\\
40000	2826.3\\
42000	2844.8\\
44000	2836.2\\
48000	2867.2\\
50000	2790.7\\
52000	2902.8\\
54000	2849\\
56000	2833.2\\
60000	2842.2\\
62000	2816.3\\
64000	2849.1\\
68000	2885.9\\
70000	2868.7\\
72000	2814.7\\
74000	2851.2\\
76000	2823.5\\
78000	2835.2\\
80000	2861.6\\
82000	2910.9\\
84000	2829.8\\
86000	2792.5\\
88000	2845.6\\
90000	2812\\
92000	2826.5\\
94000	2860.5\\
96000	2866.3\\
98000	2822.9\\
1e+05	2820.6\\
};
\addplot [color=mycolor70, forget plot]
  table[row sep=crcr]{%
0	4395.3\\
2000	2289.6\\
4000	2044.6\\
6000	2075.1\\
8000	2049.4\\
10000	2004.4\\
12000	2011.6\\
14000	2066.8\\
16000	1976.6\\
18000	2032.7\\
20000	2058.4\\
22000	2017\\
24000	2031.3\\
26000	2023.6\\
28000	2051.9\\
30000	2009.1\\
32000	2040.8\\
34000	2014.1\\
36000	2026.3\\
38000	2075.6\\
40000	2013.2\\
42000	2024.1\\
44000	2010.8\\
46000	2021.7\\
48000	2042.9\\
50000	2031.9\\
52000	2049.9\\
54000	2027.5\\
56000	2049.9\\
58000	2026.6\\
62000	2009.2\\
64000	2018.6\\
66000	2048.1\\
68000	2040.3\\
70000	2041.7\\
72000	2061.9\\
74000	2054.5\\
76000	1992.9\\
78000	2054.1\\
80000	2041.7\\
82000	1976.2\\
84000	2034\\
86000	2041.9\\
88000	2024.7\\
90000	2059.7\\
92000	2043.7\\
94000	2085.5\\
96000	2010.8\\
98000	2015.6\\
1e+05	2040.3\\
};
\addplot [color=mycolor71, forget plot]
  table[row sep=crcr]{%
0	4395.3\\
2000	2674.9\\
4000	2406.6\\
6000	2419.1\\
8000	2451.4\\
10000	2451.8\\
12000	2444.6\\
14000	2396.9\\
16000	2416.8\\
18000	2415.6\\
20000	2438.4\\
22000	2428.6\\
24000	2409\\
26000	2422.5\\
28000	2453\\
30000	2456.1\\
32000	2439.8\\
34000	2436.2\\
36000	2401.1\\
38000	2397.8\\
40000	2449.7\\
42000	2428.7\\
44000	2399.9\\
46000	2403.3\\
48000	2427.8\\
50000	2399.1\\
52000	2408.5\\
54000	2438.9\\
56000	2422.6\\
60000	2442.8\\
62000	2478.9\\
64000	2440\\
66000	2430\\
68000	2469.6\\
70000	2403.2\\
72000	2396.5\\
74000	2451.1\\
76000	2423.7\\
78000	2406.1\\
80000	2404.3\\
82000	2390.6\\
84000	2415.9\\
86000	2379.6\\
88000	2405.3\\
90000	2439\\
92000	2399.4\\
94000	2444.5\\
96000	2393.7\\
98000	2418.9\\
1e+05	2435.6\\
};
\addplot [color=mycolor72, forget plot]
  table[row sep=crcr]{%
0	4395.3\\
2000	2393.4\\
4000	2108.1\\
6000	2148.6\\
8000	2097.3\\
10000	2143.4\\
12000	2202.6\\
14000	2106\\
16000	2119.3\\
18000	2117.3\\
20000	2180\\
22000	2172.5\\
24000	2129.5\\
26000	2111.9\\
30000	2100.7\\
34000	2141.2\\
36000	2055.7\\
38000	2128.4\\
40000	2083.7\\
42000	2149.8\\
44000	2069.2\\
46000	2132.1\\
48000	2114.5\\
50000	2118.9\\
52000	2155.1\\
54000	2118.2\\
56000	2093\\
58000	2170.9\\
60000	2122.5\\
62000	2150\\
64000	2081.1\\
66000	2123.6\\
68000	2111.2\\
70000	2114.8\\
72000	2204.7\\
74000	2132.6\\
76000	2093.6\\
78000	2147\\
82000	2132.9\\
84000	2091\\
86000	2154\\
88000	2082.9\\
90000	2145.7\\
92000	2048\\
96000	2130.9\\
98000	2151.2\\
1e+05	2142.5\\
};
\addplot [color=mycolor73, forget plot]
  table[row sep=crcr]{%
0	4395.3\\
2000	2755\\
4000	2484.3\\
6000	2525.5\\
8000	2521.2\\
10000	2541.6\\
12000	2508.2\\
14000	2500.5\\
16000	2506.5\\
18000	2621.5\\
20000	2600.8\\
22000	2595\\
24000	2579.9\\
26000	2524.2\\
28000	2564.4\\
30000	2547.9\\
32000	2563.5\\
36000	2561.2\\
38000	2562\\
40000	2501.2\\
42000	2552.5\\
44000	2538.4\\
48000	2589.3\\
50000	2557.6\\
54000	2576.7\\
56000	2529.4\\
58000	2516.3\\
60000	2494.5\\
62000	2529.9\\
64000	2583.5\\
66000	2532.6\\
68000	2549\\
70000	2540.8\\
72000	2555.9\\
74000	2594.4\\
76000	2553.9\\
78000	2542.1\\
80000	2564.6\\
82000	2563.7\\
84000	2571\\
86000	2525.3\\
90000	2527.3\\
92000	2551.1\\
94000	2485.2\\
98000	2559.6\\
1e+05	2569.4\\
};
\addplot [color=mycolor74, forget plot]
  table[row sep=crcr]{%
0	4395.3\\
2000	1749.8\\
4000	1479.6\\
6000	1464.2\\
8000	1471.1\\
10000	1472.2\\
12000	1482.1\\
14000	1456.3\\
16000	1485.7\\
18000	1475.9\\
20000	1457.8\\
22000	1460.3\\
24000	1474.9\\
26000	1479.3\\
28000	1471.1\\
30000	1458.5\\
32000	1478.1\\
34000	1469.1\\
36000	1472.7\\
38000	1507.9\\
40000	1481.3\\
42000	1459.3\\
44000	1446.7\\
46000	1481.9\\
48000	1468.4\\
50000	1474.2\\
52000	1458.4\\
54000	1489.2\\
56000	1486.4\\
58000	1467.1\\
62000	1471\\
64000	1442.8\\
66000	1492.3\\
70000	1475.6\\
72000	1476.6\\
74000	1498.1\\
76000	1461.6\\
78000	1445.2\\
80000	1454.4\\
82000	1434.4\\
84000	1451.5\\
86000	1474.8\\
88000	1456.6\\
90000	1496.9\\
92000	1481.5\\
94000	1447.8\\
96000	1477.4\\
98000	1453.4\\
1e+05	1471.2\\
};
\addplot [color=mycolor75, forget plot]
  table[row sep=crcr]{%
0	4395.3\\
2000	2445.5\\
4000	2143.4\\
6000	2146\\
8000	2133.5\\
10000	2113.3\\
14000	2139.5\\
16000	2113.2\\
18000	2176.9\\
20000	2196.2\\
22000	2158.5\\
24000	2206.3\\
26000	2166.3\\
28000	2117.2\\
30000	2105.5\\
32000	2186\\
34000	2122.9\\
36000	2141.9\\
38000	2071.9\\
40000	2173.4\\
42000	2169.7\\
44000	2194.8\\
46000	2130.9\\
48000	2122.5\\
50000	2159\\
52000	2121.7\\
54000	2153.4\\
56000	2099.7\\
58000	2146.2\\
60000	2157.1\\
62000	2129.2\\
66000	2173.7\\
68000	2163.5\\
70000	2160.3\\
72000	2118.4\\
76000	2153.5\\
80000	2129.1\\
82000	2151.4\\
84000	2136.9\\
86000	2151.6\\
88000	2152.3\\
90000	2169.1\\
92000	2118.4\\
94000	2143.7\\
96000	2135.1\\
98000	2112.3\\
1e+05	2168.9\\
};
\addplot [color=mycolor76, forget plot]
  table[row sep=crcr]{%
0	4395.3\\
2000	2200.2\\
4000	1919\\
6000	1950.4\\
8000	1925.9\\
10000	1952.3\\
12000	1937\\
14000	1971.6\\
16000	1927.5\\
20000	1969\\
22000	1938.1\\
24000	1939.3\\
26000	1919.9\\
28000	1946.8\\
30000	1931.5\\
34000	1968.4\\
36000	1930.9\\
38000	1959\\
44000	1931.6\\
46000	1953.6\\
48000	1907.3\\
50000	1940.8\\
52000	1943.4\\
54000	1938.6\\
56000	1955.5\\
58000	1911.4\\
60000	1983.1\\
62000	1938.9\\
64000	1924.4\\
66000	1951.2\\
68000	1923.8\\
70000	1949.6\\
72000	1939\\
76000	1931.5\\
78000	1900.7\\
80000	1955.1\\
82000	1943.4\\
84000	1961.4\\
86000	1919.7\\
88000	1957.6\\
92000	1936.9\\
94000	1955.3\\
98000	1947.6\\
1e+05	1977.4\\
};
\addplot [color=mycolor77, forget plot]
  table[row sep=crcr]{%
0	4395.3\\
2000	1813.7\\
4000	1494.7\\
6000	1477.3\\
8000	1481\\
12000	1541.6\\
14000	1524.9\\
16000	1513.3\\
20000	1473.2\\
24000	1501.8\\
26000	1507.7\\
28000	1485.9\\
30000	1471.3\\
32000	1485\\
34000	1503.9\\
36000	1474.2\\
38000	1522\\
40000	1497.7\\
42000	1490.3\\
44000	1496.9\\
46000	1495.1\\
48000	1501.4\\
52000	1482.2\\
54000	1482.8\\
56000	1520.6\\
58000	1496.9\\
60000	1490.3\\
62000	1520.6\\
64000	1493.2\\
66000	1543.5\\
68000	1484.5\\
72000	1511.2\\
74000	1512.6\\
76000	1497.9\\
78000	1510.6\\
80000	1502.1\\
82000	1474.3\\
84000	1504.9\\
86000	1508.1\\
88000	1464.2\\
90000	1524.3\\
94000	1498.3\\
96000	1502.1\\
98000	1490.9\\
1e+05	1502.7\\
};
\addplot [color=mycolor78, forget plot]
  table[row sep=crcr]{%
0	4395.3\\
2000	2454.2\\
4000	2217.1\\
6000	2269.4\\
8000	2167.4\\
10000	2220.6\\
12000	2241.1\\
14000	2208.1\\
16000	2196.6\\
20000	2235.3\\
22000	2226.1\\
24000	2233.8\\
26000	2226.8\\
28000	2281.7\\
30000	2250\\
32000	2263.6\\
34000	2209.8\\
36000	2252.9\\
42000	2250.6\\
44000	2239.7\\
48000	2296.1\\
50000	2202.6\\
58000	2236.4\\
60000	2262.1\\
62000	2233.9\\
64000	2238.2\\
66000	2173\\
68000	2241.1\\
70000	2190.7\\
72000	2184.9\\
74000	2245.2\\
76000	2281.4\\
78000	2209.6\\
80000	2216.7\\
84000	2201.4\\
86000	2172.9\\
88000	2285.6\\
90000	2250.8\\
92000	2205.5\\
96000	2269.8\\
98000	2241.9\\
1e+05	2259.4\\
};
\addplot [color=mycolor79, forget plot]
  table[row sep=crcr]{%
0	4395.3\\
2000	2330\\
4000	2081.9\\
6000	2047.4\\
8000	2069.3\\
10000	2079.7\\
12000	2133.2\\
16000	2067.7\\
18000	2074.4\\
20000	2104.3\\
22000	2077.7\\
24000	2065.4\\
26000	2080.1\\
28000	2054.1\\
30000	2125.3\\
32000	2057.4\\
34000	2102.8\\
36000	2119.3\\
38000	2192.2\\
40000	2145.4\\
42000	2110.9\\
44000	2091.6\\
46000	2024.1\\
48000	2051.9\\
50000	2105.3\\
52000	2029.7\\
56000	2095.9\\
58000	2118.5\\
60000	2073.3\\
62000	2143.8\\
64000	2083.6\\
66000	2115.5\\
68000	2102.2\\
70000	2108.2\\
72000	2060.7\\
74000	2122.1\\
76000	2095\\
78000	2078.7\\
82000	2063.9\\
84000	2110\\
86000	2079.1\\
88000	2122.2\\
92000	2108\\
94000	2117.7\\
96000	2155.6\\
98000	2090.5\\
1e+05	2122.2\\
};
\addplot [color=mycolor79, forget plot]
  table[row sep=crcr]{%
0	4395.3\\
2000	2215.4\\
4000	1936.1\\
6000	1919.2\\
8000	1986.6\\
10000	1968.2\\
12000	1988\\
14000	1937.5\\
16000	1926.9\\
20000	1988.5\\
24000	1978.9\\
26000	1912.7\\
28000	1939.5\\
30000	1907.1\\
32000	1889.9\\
34000	1950.2\\
40000	1930.3\\
42000	1968.6\\
44000	1900.3\\
46000	1899.9\\
48000	1935.1\\
50000	1920.7\\
52000	1935.6\\
58000	1956.3\\
60000	1946\\
62000	1984.8\\
64000	1971.9\\
66000	1940\\
68000	1890.2\\
72000	1956.2\\
74000	1974.3\\
76000	1920.3\\
78000	1899.2\\
80000	1974.4\\
82000	1949.2\\
84000	1985\\
86000	1938\\
88000	1971.3\\
90000	1909\\
92000	1962.6\\
94000	1964.8\\
96000	1945.6\\
98000	1947.9\\
1e+05	2003.8\\
};
\addplot [color=mycolor80, forget plot]
  table[row sep=crcr]{%
0	4395.3\\
2000	1799.6\\
4000	1502.2\\
6000	1521.8\\
8000	1499\\
10000	1497.4\\
12000	1462.8\\
14000	1495.7\\
16000	1470.6\\
18000	1526.9\\
20000	1472\\
22000	1489.7\\
24000	1487.7\\
26000	1507.4\\
28000	1473\\
30000	1484.2\\
32000	1470.5\\
34000	1469.1\\
36000	1492.1\\
38000	1471.9\\
40000	1466.9\\
42000	1468.9\\
46000	1489.5\\
48000	1467.4\\
50000	1493.1\\
52000	1474.4\\
56000	1477.5\\
58000	1531.5\\
60000	1472.3\\
62000	1495.8\\
64000	1488.1\\
66000	1526.4\\
68000	1486.3\\
70000	1507\\
72000	1493\\
74000	1489.9\\
76000	1478.6\\
78000	1491.3\\
82000	1499.4\\
84000	1476.6\\
86000	1496.9\\
88000	1494.9\\
92000	1479.8\\
94000	1490.8\\
96000	1521.3\\
98000	1492.9\\
1e+05	1502.8\\
};
\addplot [color=mycolor81, forget plot]
  table[row sep=crcr]{%
0	4395.3\\
2000	1707.1\\
4000	1393.7\\
6000	1369.2\\
8000	1385.9\\
10000	1393.8\\
12000	1422.4\\
14000	1410.3\\
16000	1404.6\\
18000	1384\\
20000	1386.1\\
22000	1399.7\\
24000	1393.3\\
26000	1381\\
28000	1395.2\\
30000	1367.8\\
32000	1384.3\\
34000	1408\\
36000	1384.7\\
38000	1399.9\\
40000	1396.5\\
42000	1369.1\\
44000	1384.2\\
46000	1442.2\\
48000	1397.9\\
50000	1379.4\\
52000	1384.3\\
54000	1396.8\\
56000	1420.3\\
58000	1401.7\\
60000	1379\\
62000	1377.8\\
64000	1413.1\\
66000	1387.2\\
68000	1395.8\\
72000	1387.2\\
74000	1400.5\\
76000	1402.2\\
78000	1388.1\\
80000	1408.3\\
82000	1400.3\\
84000	1401.5\\
86000	1355.6\\
88000	1375.3\\
90000	1409\\
92000	1379\\
94000	1370.1\\
96000	1408.2\\
98000	1431.3\\
1e+05	1405.3\\
};
\addplot [color=mycolor82, forget plot]
  table[row sep=crcr]{%
0	4395.3\\
2000	1952.6\\
4000	1642.4\\
6000	1653.6\\
8000	1638.4\\
12000	1647.6\\
14000	1640.6\\
16000	1628\\
18000	1628.3\\
20000	1666.2\\
22000	1605.5\\
24000	1624.7\\
28000	1638.3\\
32000	1629\\
34000	1632.2\\
36000	1653.9\\
38000	1646.4\\
40000	1620.6\\
42000	1610.1\\
44000	1634.7\\
46000	1619.9\\
48000	1635.6\\
52000	1644.5\\
54000	1620.4\\
56000	1638.9\\
58000	1613.3\\
60000	1627.5\\
62000	1624.9\\
64000	1613\\
66000	1635.7\\
68000	1623.6\\
70000	1637.5\\
72000	1637.2\\
74000	1629.6\\
76000	1642.2\\
78000	1623.3\\
80000	1655.1\\
82000	1635.3\\
84000	1663.7\\
86000	1645.2\\
88000	1615.6\\
90000	1611.4\\
92000	1628.7\\
94000	1637.8\\
96000	1639.7\\
98000	1616.4\\
1e+05	1615\\
};
\addplot [color=mycolor83, forget plot]
  table[row sep=crcr]{%
0	4395.3\\
2000	1996.1\\
4000	1633.1\\
6000	1660.7\\
8000	1616.8\\
10000	1663.2\\
12000	1697.7\\
14000	1641.6\\
16000	1654.2\\
18000	1676.3\\
20000	1646.9\\
22000	1649.8\\
24000	1629.6\\
26000	1662.8\\
30000	1644.6\\
34000	1648.4\\
36000	1652.1\\
38000	1664.1\\
40000	1665.3\\
42000	1683.5\\
44000	1669.1\\
46000	1674.9\\
48000	1636.5\\
50000	1637.9\\
52000	1647.1\\
54000	1674.5\\
56000	1631.4\\
58000	1621.1\\
60000	1669.2\\
62000	1668.6\\
66000	1657.3\\
68000	1677.5\\
70000	1664.2\\
72000	1645.4\\
74000	1654.2\\
76000	1668.5\\
78000	1652.5\\
80000	1627.9\\
82000	1658.8\\
84000	1637\\
88000	1689.7\\
90000	1667.5\\
92000	1665.1\\
94000	1674.5\\
96000	1626.6\\
98000	1648.6\\
1e+05	1642.2\\
};
\addplot [color=mycolor83, forget plot]
  table[row sep=crcr]{%
0	4395.3\\
2000	1575.5\\
4000	1191.9\\
6000	1210.8\\
8000	1210.1\\
10000	1178\\
12000	1226.4\\
14000	1240.7\\
16000	1195.8\\
18000	1201\\
20000	1214.4\\
22000	1223.8\\
24000	1209.3\\
26000	1181.2\\
28000	1227.1\\
30000	1166\\
32000	1213.7\\
34000	1228.3\\
36000	1233.9\\
38000	1205.8\\
40000	1228.6\\
42000	1212.6\\
46000	1232.3\\
48000	1212.4\\
50000	1212.7\\
52000	1190\\
54000	1210.7\\
56000	1180.8\\
58000	1181.4\\
60000	1201.6\\
62000	1226.3\\
66000	1210.1\\
68000	1187.1\\
72000	1225.7\\
74000	1236.4\\
76000	1203.8\\
78000	1218.4\\
80000	1193.2\\
82000	1190.7\\
84000	1208.6\\
86000	1235.3\\
88000	1212.7\\
90000	1201.7\\
96000	1199.7\\
98000	1208.7\\
1e+05	1189\\
};
\addplot [color=mycolor84, forget plot]
  table[row sep=crcr]{%
0	4395.3\\
2000	1618.6\\
4000	1229.3\\
6000	1205.9\\
8000	1189.3\\
10000	1191.1\\
12000	1236.4\\
14000	1231.7\\
16000	1213.4\\
18000	1201\\
20000	1202.1\\
22000	1191.7\\
24000	1207\\
26000	1181.6\\
28000	1186.3\\
30000	1202.9\\
32000	1202.2\\
34000	1215.7\\
36000	1221.7\\
38000	1211\\
40000	1178\\
42000	1194.9\\
44000	1246.2\\
46000	1195.5\\
48000	1234.7\\
50000	1210.2\\
52000	1177.7\\
54000	1229.1\\
56000	1251.2\\
58000	1216.9\\
60000	1173.2\\
62000	1201.3\\
64000	1180.7\\
66000	1242.5\\
68000	1202.3\\
70000	1193\\
72000	1175.1\\
74000	1195\\
76000	1196.3\\
78000	1206.5\\
80000	1202.5\\
82000	1215.6\\
84000	1266.3\\
86000	1238.4\\
88000	1216.4\\
90000	1188.3\\
92000	1197\\
94000	1212.5\\
96000	1203.6\\
98000	1182\\
1e+05	1205.4\\
};
\addplot [color=mycolor85, forget plot]
  table[row sep=crcr]{%
0	4395.3\\
2000	1693.3\\
4000	1356.4\\
6000	1354.6\\
8000	1369.7\\
10000	1312\\
14000	1417.9\\
16000	1392.6\\
18000	1355.7\\
20000	1344.1\\
22000	1314\\
24000	1322.7\\
26000	1341.6\\
28000	1350.4\\
30000	1354.8\\
32000	1318.7\\
34000	1347.5\\
36000	1361.2\\
38000	1328.3\\
40000	1335.9\\
42000	1348.4\\
44000	1330.1\\
46000	1346.5\\
48000	1342.6\\
52000	1357.3\\
54000	1358.2\\
56000	1322.2\\
58000	1367.5\\
60000	1301.7\\
62000	1341\\
66000	1354.2\\
68000	1335.5\\
70000	1325.1\\
72000	1335\\
74000	1308.7\\
76000	1310.8\\
78000	1336\\
80000	1352.4\\
82000	1319.3\\
86000	1334.2\\
88000	1312.9\\
90000	1359.4\\
92000	1329.7\\
94000	1387.1\\
96000	1368.4\\
98000	1359.8\\
1e+05	1303\\
};
\addplot [color=mycolor86, forget plot]
  table[row sep=crcr]{%
0	4395.3\\
2000	1562.2\\
4000	1108.1\\
6000	1092.4\\
8000	1118.6\\
10000	1161.5\\
12000	1094.1\\
14000	1100\\
16000	1127.3\\
18000	1120\\
20000	1090.8\\
22000	1116.5\\
24000	1121.9\\
28000	1138.1\\
30000	1104.1\\
32000	1137.7\\
34000	1151.5\\
36000	1112.5\\
38000	1132.5\\
40000	1128.6\\
42000	1102.9\\
44000	1117.2\\
46000	1105.5\\
50000	1110.1\\
52000	1109.7\\
54000	1105.8\\
56000	1113.4\\
58000	1109.2\\
60000	1114\\
62000	1129.7\\
64000	1102.1\\
66000	1140.7\\
68000	1103.3\\
72000	1094.5\\
74000	1121.4\\
76000	1133.4\\
78000	1083.1\\
80000	1083.5\\
82000	1101.4\\
84000	1131.3\\
86000	1141.5\\
88000	1109.8\\
90000	1104.9\\
94000	1141.9\\
96000	1105.2\\
1e+05	1101.9\\
};
\addplot [color=mycolor87, forget plot]
  table[row sep=crcr]{%
0	4395.3\\
2000	1378\\
4000	930.3\\
8000	925.59\\
10000	875.46\\
12000	989.72\\
14000	992.72\\
16000	981.1\\
18000	957.04\\
20000	979.35\\
22000	912.08\\
24000	928.62\\
26000	903.56\\
28000	944.67\\
30000	958.61\\
32000	905.8\\
34000	940.67\\
36000	962.45\\
38000	970.28\\
40000	930.91\\
42000	940.37\\
44000	914.53\\
46000	976.35\\
48000	989\\
50000	918.4\\
52000	960.21\\
54000	911.92\\
60000	1004.4\\
62000	1011.3\\
64000	918.45\\
66000	873.2\\
68000	959.92\\
70000	975\\
74000	965.56\\
76000	974.19\\
78000	971.51\\
80000	947.46\\
82000	917.98\\
84000	913.41\\
86000	903.93\\
88000	941.37\\
90000	924.78\\
92000	930.59\\
94000	993.4\\
96000	930.68\\
1e+05	949.69\\
};
\addplot [color=mycolor88, forget plot]
  table[row sep=crcr]{%
0	4395.3\\
2000	1538.8\\
4000	1086.2\\
6000	1114.4\\
8000	1103.7\\
10000	1128.1\\
12000	1118.6\\
14000	1129.2\\
16000	1151.8\\
18000	1154.7\\
22000	1104.9\\
26000	1125.7\\
28000	1153.4\\
30000	1126.1\\
34000	1126.8\\
36000	1143.7\\
38000	1116.5\\
40000	1135.6\\
42000	1167.1\\
44000	1102.5\\
46000	1123.8\\
50000	1103.3\\
52000	1112.7\\
56000	1087.3\\
58000	1117.8\\
60000	1109.7\\
62000	1134.9\\
64000	1098.8\\
66000	1125.8\\
68000	1097.4\\
70000	1117.4\\
72000	1112.2\\
74000	1117.8\\
76000	1129.7\\
78000	1128.4\\
80000	1131.8\\
82000	1100.7\\
84000	1087.4\\
86000	1150\\
88000	1087.5\\
90000	1130.3\\
92000	1094.8\\
94000	1089.9\\
96000	1108.1\\
98000	1076.6\\
1e+05	1105\\
};
\addplot [color=mycolor89, forget plot]
  table[row sep=crcr]{%
0	4395.3\\
2000	1277.2\\
4000	801\\
6000	793.95\\
8000	781.79\\
12000	801.41\\
14000	767.46\\
16000	759.3\\
18000	800.93\\
20000	801.27\\
22000	778\\
24000	777.85\\
26000	799.11\\
28000	790.35\\
30000	811.74\\
32000	796.13\\
34000	801.75\\
36000	787.2\\
38000	807.15\\
40000	821.98\\
42000	796.97\\
44000	812.18\\
46000	776.65\\
48000	782.22\\
50000	797.19\\
52000	783.85\\
54000	788.87\\
58000	767.93\\
60000	793.14\\
62000	810.29\\
64000	803.27\\
66000	799.7\\
68000	801.56\\
70000	798.57\\
72000	785.53\\
74000	838.3\\
76000	770.55\\
78000	798.93\\
80000	762.01\\
82000	780.35\\
84000	804.39\\
86000	800.92\\
88000	805.75\\
90000	788.88\\
92000	809.18\\
94000	784.94\\
96000	772.1\\
98000	791.28\\
1e+05	771.06\\
};
\addplot [color=mycolor90, forget plot]
  table[row sep=crcr]{%
0	4395.3\\
2000	1249.2\\
4000	802.63\\
6000	762.24\\
8000	769.76\\
10000	794.37\\
12000	775.38\\
14000	804.67\\
16000	787.9\\
18000	768.9\\
20000	760.15\\
22000	776.75\\
28000	777.1\\
30000	769.74\\
32000	792.83\\
34000	768.45\\
38000	780.32\\
40000	767.35\\
42000	763.87\\
44000	754.42\\
46000	782.73\\
48000	795.03\\
50000	786.25\\
52000	781.22\\
54000	764.34\\
56000	758.31\\
58000	787.77\\
60000	782.7\\
62000	786.48\\
64000	760.27\\
66000	778.42\\
68000	784.81\\
70000	755.05\\
72000	782.66\\
74000	759.77\\
76000	769.62\\
78000	753.19\\
80000	790.91\\
82000	764.38\\
84000	768.65\\
86000	783.35\\
88000	766.74\\
92000	769.52\\
94000	764.86\\
96000	804.96\\
98000	773.4\\
1e+05	761.18\\
};
\addplot [color=mycolor91, forget plot]
  table[row sep=crcr]{%
0	4395.3\\
2000	1212.2\\
4000	667.76\\
6000	658.14\\
8000	644.49\\
10000	651.51\\
12000	671.11\\
14000	642.67\\
16000	648.25\\
18000	658.77\\
20000	658.08\\
22000	673.92\\
24000	658.62\\
26000	676.05\\
28000	661.05\\
30000	659.08\\
32000	665.51\\
34000	660.12\\
36000	646.17\\
38000	670.44\\
42000	670.64\\
44000	613.87\\
46000	665.27\\
48000	669.72\\
50000	668.33\\
52000	669.83\\
54000	683.73\\
56000	673.89\\
58000	678.78\\
62000	650.46\\
64000	650.31\\
66000	690.81\\
68000	669.21\\
70000	680.8\\
72000	687.78\\
74000	654.72\\
76000	651.09\\
78000	690.01\\
80000	621.83\\
82000	685.26\\
84000	672.57\\
86000	675.02\\
88000	646.34\\
90000	656.07\\
92000	644.94\\
94000	648.64\\
96000	679.29\\
98000	675.41\\
1e+05	662.37\\
};
\addplot [color=mycolor92, forget plot]
  table[row sep=crcr]{%
0	4395.3\\
2000	1420.2\\
4000	838.28\\
6000	798.45\\
8000	827.74\\
10000	825.79\\
12000	816.58\\
14000	815.66\\
16000	801.82\\
18000	815.93\\
20000	858.43\\
22000	831.08\\
24000	843.45\\
26000	818.78\\
28000	831.34\\
30000	835.09\\
32000	820.96\\
34000	822.02\\
36000	818.66\\
38000	834.57\\
40000	831.72\\
42000	855.47\\
44000	837.45\\
46000	807.71\\
48000	839.82\\
50000	824.63\\
52000	835.5\\
54000	838.07\\
56000	807.7\\
58000	820.73\\
60000	846.4\\
62000	846.61\\
64000	818.96\\
66000	830.3\\
68000	823.63\\
70000	838.89\\
72000	833.17\\
74000	818.44\\
76000	832.26\\
78000	828.83\\
80000	851.02\\
82000	839.89\\
84000	834.74\\
86000	846.56\\
88000	815.4\\
90000	832.36\\
94000	828.8\\
96000	835.27\\
98000	824.75\\
1e+05	827.98\\
};
\addplot [color=mycolor93, forget plot]
  table[row sep=crcr]{%
0	4395.3\\
2000	1288\\
4000	714.28\\
6000	714.93\\
8000	708.21\\
10000	679.19\\
12000	678.83\\
14000	692.77\\
16000	742.1\\
18000	713.43\\
20000	701.35\\
22000	661.45\\
24000	683.65\\
26000	689.75\\
28000	685.56\\
30000	710.84\\
32000	710.17\\
34000	703.75\\
36000	736.01\\
38000	683.83\\
40000	672.37\\
42000	688.34\\
44000	715.27\\
46000	729.25\\
48000	645.61\\
50000	684.41\\
52000	718.47\\
54000	708.13\\
56000	708.9\\
58000	680.01\\
60000	706.08\\
62000	694.02\\
64000	710.74\\
66000	704.67\\
68000	754.5\\
70000	675.93\\
72000	712.9\\
74000	748.65\\
76000	692.64\\
78000	649.72\\
80000	660.74\\
82000	701.5\\
84000	688.24\\
86000	712.45\\
88000	721.7\\
90000	703.71\\
92000	662.42\\
94000	709.19\\
96000	692.92\\
98000	714.21\\
1e+05	701.06\\
};
\addplot [color=mycolor94, forget plot]
  table[row sep=crcr]{%
0	4395.3\\
2000	1266.6\\
4000	668.84\\
6000	660.77\\
8000	676.74\\
10000	660.18\\
12000	636.24\\
14000	684.23\\
16000	674.97\\
18000	657.15\\
20000	667.87\\
22000	673.93\\
24000	664.51\\
26000	661.68\\
28000	643.08\\
30000	666.96\\
32000	666.91\\
34000	648.32\\
36000	684.36\\
38000	690.31\\
40000	666.46\\
42000	670.7\\
44000	662.34\\
46000	656.04\\
48000	668.25\\
50000	656.25\\
52000	669.58\\
54000	670.47\\
56000	656.74\\
58000	654.41\\
60000	660.1\\
62000	653.18\\
64000	661.66\\
66000	655.87\\
68000	665.52\\
70000	642.69\\
72000	653.09\\
74000	633.77\\
76000	651.9\\
78000	660.5\\
80000	654.46\\
82000	654.88\\
84000	652.31\\
88000	668.32\\
90000	672.17\\
92000	671.23\\
94000	656.64\\
96000	665.56\\
98000	647.43\\
1e+05	657.72\\
};
\addplot [color=mycolor95, forget plot]
  table[row sep=crcr]{%
0	4395.3\\
2000	993.86\\
4000	397.92\\
6000	384.14\\
8000	368.9\\
10000	383.07\\
12000	406.43\\
14000	384\\
16000	382.78\\
18000	390.1\\
20000	376.31\\
24000	386.44\\
26000	389.81\\
28000	404.2\\
30000	384.52\\
32000	400.94\\
34000	398.49\\
36000	376.17\\
38000	396.66\\
40000	396.17\\
42000	386.2\\
44000	400.5\\
46000	402.64\\
48000	376.68\\
50000	377.23\\
52000	389.62\\
54000	377.18\\
56000	387.28\\
58000	399.51\\
60000	380.44\\
62000	384.93\\
64000	375.54\\
66000	389.18\\
68000	385.33\\
70000	397.14\\
72000	387.81\\
74000	385.07\\
76000	385.43\\
78000	391.21\\
80000	375.71\\
82000	379.33\\
84000	391.3\\
86000	377.66\\
88000	388.84\\
90000	380.62\\
92000	405.82\\
94000	389.37\\
96000	387.6\\
98000	396.05\\
1e+05	387.32\\
};
\addplot [color=mycolor96, forget plot]
  table[row sep=crcr]{%
0	4395.3\\
2000	1174.9\\
4000	430.04\\
6000	401.4\\
8000	432.54\\
10000	404.56\\
12000	402.66\\
14000	367.87\\
16000	421.09\\
18000	382.48\\
20000	402.15\\
22000	381.91\\
24000	360.96\\
26000	401.29\\
28000	414.46\\
30000	440.52\\
32000	427.31\\
34000	391.92\\
36000	411.58\\
38000	390.36\\
40000	407.99\\
42000	441.16\\
44000	398.49\\
46000	411.54\\
48000	386.63\\
50000	413.5\\
52000	428.33\\
54000	404.61\\
56000	391.22\\
58000	417.46\\
60000	388.28\\
62000	408.77\\
64000	431.83\\
66000	393.73\\
68000	425.24\\
72000	421.12\\
74000	406.6\\
76000	449.1\\
78000	403.13\\
80000	396.16\\
82000	435.96\\
84000	398.01\\
86000	451.83\\
88000	421.95\\
90000	405.84\\
92000	373.45\\
94000	443.58\\
96000	412.07\\
98000	385.2\\
1e+05	392.2\\
};
\addplot [color=mycolor97, forget plot]
  table[row sep=crcr]{%
0	4395.3\\
2000	1014.3\\
4000	270.68\\
6000	257.61\\
8000	274.83\\
10000	292.09\\
12000	286.67\\
14000	292.1\\
16000	267.02\\
18000	324.12\\
20000	274.13\\
22000	265.26\\
24000	269.93\\
26000	275.63\\
28000	267.9\\
30000	268.08\\
32000	261.17\\
34000	289.77\\
36000	263.64\\
38000	308.45\\
40000	284.55\\
44000	277.84\\
46000	292.9\\
48000	274.84\\
52000	276.32\\
54000	266.08\\
56000	276.19\\
58000	291\\
60000	276.19\\
62000	287.96\\
64000	269.39\\
66000	283.83\\
68000	281.43\\
70000	275.1\\
72000	265.63\\
74000	267.44\\
76000	295.29\\
78000	284.62\\
80000	272.96\\
82000	291.93\\
84000	286.56\\
86000	292.66\\
88000	289.95\\
90000	288.78\\
92000	279.91\\
94000	299.66\\
96000	287.99\\
98000	263.36\\
1e+05	278.92\\
};
\addplot [color=mycolor98]
  table[row sep=crcr]{%
0	4395.3\\
2000	1238.3\\
4000	404\\
6000	390.41\\
8000	383.62\\
10000	382.41\\
12000	391.72\\
14000	402.46\\
16000	373\\
18000	370.05\\
20000	382.58\\
22000	399.47\\
24000	400.2\\
26000	379.72\\
28000	391.43\\
30000	406.08\\
32000	377.42\\
34000	394.16\\
36000	395.3\\
38000	405.03\\
40000	392.86\\
42000	392.29\\
44000	400.05\\
46000	414.28\\
48000	390.98\\
50000	416.05\\
52000	405.15\\
54000	402.45\\
56000	379.4\\
58000	404.55\\
60000	393.34\\
62000	401.03\\
64000	385.29\\
66000	399.45\\
68000	388.07\\
72000	382.11\\
74000	352.61\\
76000	372.64\\
78000	368.41\\
80000	389.99\\
82000	401.24\\
84000	365.87\\
86000	369.46\\
88000	389.95\\
90000	401.86\\
92000	382.88\\
94000	403.96\\
96000	390.38\\
98000	384.16\\
1e+05	382.08\\
};
\addlegendentry{\ClasskIVLegend}

      % \input{../../../../.tex/20/KS/SizeEvolutionsfig10.tex}
      \legend{}
    \end{groupplot}
  %\endgroup

  %\begingroup Parametri
    
%\begingroup Parametro
\providecommand\varyingParameterSymbol\lambda
\providecommand\varyingParameterList[1]{
  \ifcase#1\relax%
  \or\num{1e-2}% Caso 1
  \or\num{8.2e-2}% Caso 2
  \or\num{1.55e-1}% Caso 3
  \or\num{2.27e-1}% Caso 4
  \else\num{3e-1}% Caso 5
  \fi
}
%\endgroup

%\begingroup beta
\providecommand\betaExactList[1]{
  \ifcase#1\relax%
  \or\num{1.267+-0.007}% Caso 1
  \or\num{1.285+-0.010}% Caso 2
  \or\num{1.264+-0.015}% Caso 3
  \or\num{1.246+-0.015}% Caso 4
  \else\num{1.210+-0.015}% Caso 5
  \fi
}
\providecommand\betaAprxList[1]{
  \ifcase#1\relax%
  \or\num{1.264+-0.006}% Caso 1
  \or\num{1.270+-0.011}% Caso 2
  \or\num{1.280+-0.013}% Caso 3
  \or\num{1.250+-0.015}% Caso 4
  \else\num{1.218+-0.014}% Caso 5
  \fi
}
%\endgroup

%\begingroup xi
\providecommand\zetaExactList[1]{
  \ifcase#1\relax%
  \or\num{0.757+-0.004}% Caso 1
  \or\num{0.710+-0.026}% Caso 2
  \or\num{0.670+-0.035}% Caso 3
  \or\num{0.625+-0.052}% Caso 4
  \else\num{0.539+-0.059}% Caso 5
  \fi
}
\providecommand\zetaAprxList[1]{
  \ifcase#1\relax%
  \or\num{0.749+-0.004}% Caso 1
  \or\num{0.721+-0.022}% Caso 2
  \or\num{0.645+-0.040}% Caso 3
  \or\num{0.622+-0.051}% Caso 4
  \else\num{0.567+-0.053}% Caso 5
  \fi
}
%\endgroup

%\begingroup mI
\providecommand\mIExactList[1]{
  \ifcase#1\relax%
  \or\num{3.265+-0.001}% Caso 1
  \or\num{3.238+-0.005}% Caso 2
  \or\num{3.220+-0.008}% Caso 3
  \or\num{3.213+-0.013}% Caso 4
  \else\num{3.174+-0.023}% Caso 5
  \fi
}
\providecommand\mIAprxList[1]{
  \ifcase#1\relax%
  \or\num{3.262+-0.001}% Caso 1
  \or\num{3.240+-0.005}% Caso 2
  \or\num{3.217+-0.010}% Caso 3
  \or\num{3.217+-0.012}% Caso 4
  \else\num{3.208+-0.009}% Caso 5
  \fi
}
%\endgroup

%\begingroup sI
\providecommand\sIExactList[1]{
  \ifcase#1\relax%
  \or\num{0.284+-0.001}% Caso 1
  \or\num{0.311+-0.008}% Caso 2
  \or\num{0.343+-0.010}% Caso 3
  \or\num{0.361+-0.020}% Caso 4
  \else\num{0.398+-0.029}% Caso 5
  \fi
}
\providecommand\sIAprxList[1]{
  \ifcase#1\relax%
  \or\num{0.281+-0.001}% Caso 1
  \or\num{0.315+-0.006}% Caso 2
  \or\num{0.334+-0.013}% Caso 3
  \or\num{0.366+-0.018}% Caso 4
  \else\num{0.411+-0.025}% Caso 5
  \fi
}
%\endgroup

%\begingroup mII
\providecommand\mIIExactList[1]{
  \ifcase#1\relax%
  \or\num{3.788+-0.005}% Caso 1
  \or\num{3.730+-0.029}% Caso 2
  \or\num{3.681+-0.052}% Caso 3
  \or\num{3.650+-0.072}% Caso 4
  \else\num{3.572+-0.095}% Caso 5
  \fi
}
\providecommand\mIIAprxList[1]{
  \ifcase#1\relax%
  \or\num{3.778+-0.006}% Caso 1
  \or\num{3.735+-0.029}% Caso 2
  \or\num{3.656+-0.048}% Caso 3
  \or\num{3.658+-0.084}% Caso 4
  \else\num{3.502+-0.078}% Caso 5
  \fi
}
%\endgroup

%\begingroup sII
\providecommand\sIIExactList[1]{
  \ifcase#1\relax%
  \or\num{0.455+-0.002}% Caso 1
  \or\num{0.486+-0.007}% Caso 2
  \or\num{0.515+-0.014}% Caso 3
  \or\num{0.533+-0.021}% Caso 4
  \else\num{0.570+-0.028}% Caso 5
  \fi
}
\providecommand\sIIAprxList[1]{
  \ifcase#1\relax%
  \or\num{0.457+-0.002}% Caso 1
  \or\num{0.490+-0.007}% Caso 2
  \or\num{0.521+-0.012}% Caso 3
  \or\num{0.545+-0.021}% Caso 4
  \else\num{0.586+-0.029}% Caso 5
  \fi
}
%\endgroup

    \ifpageisodd\writeVaryingParameter{1}{west}{east}{-.25em}{0em}
    \else\writeVaryingParameter{2}{east}{west}{+.25em}{0em}\fi
  %\endgroup

  %\begingroup Didascalie
    \writeSubCaptions[2]{south west}{south west}{0em}{0em}
    % https://tex.stackexchange.com/questions/632045/get-subcaption-for-every-plot-using-groupplot-pgf-tikz-after-2020-update
    \writeSizeEvolutionDescriptions[1]
  %\endgroup

  %\begingroup Titoli
    % \node[title,overlay] at (Row1 c1r1.north) {Esatto};
    % \node[title,overlay] at (Row1 c2r1.north) {Approssimato};
  %\endgroup

  \redefineTikZbounds{Row1 c1r1}{Row5 c2r1}
\end{tikzpicture}%

% \input{Parametri.txt}
% 
%\begingroup Parametro
\providecommand\varyingParameterSymbol\lambda
\providecommand\varyingParameterList[1]{
  \ifcase#1\relax%
  \or\num{1e-2}% Caso 1
  \or\num{8.2e-2}% Caso 2
  \or\num{1.55e-1}% Caso 3
  \or\num{2.27e-1}% Caso 4
  \else\num{3e-1}% Caso 5
  \fi
}
%\endgroup

%\begingroup beta
\providecommand\betaExactList[1]{
  \ifcase#1\relax%
  \or\num{1.267+-0.007}% Caso 1
  \or\num{1.285+-0.010}% Caso 2
  \or\num{1.264+-0.015}% Caso 3
  \or\num{1.246+-0.015}% Caso 4
  \else\num{1.210+-0.015}% Caso 5
  \fi
}
\providecommand\betaAprxList[1]{
  \ifcase#1\relax%
  \or\num{1.264+-0.006}% Caso 1
  \or\num{1.270+-0.011}% Caso 2
  \or\num{1.280+-0.013}% Caso 3
  \or\num{1.250+-0.015}% Caso 4
  \else\num{1.218+-0.014}% Caso 5
  \fi
}
%\endgroup

%\begingroup xi
\providecommand\zetaExactList[1]{
  \ifcase#1\relax%
  \or\num{0.757+-0.004}% Caso 1
  \or\num{0.710+-0.026}% Caso 2
  \or\num{0.670+-0.035}% Caso 3
  \or\num{0.625+-0.052}% Caso 4
  \else\num{0.539+-0.059}% Caso 5
  \fi
}
\providecommand\zetaAprxList[1]{
  \ifcase#1\relax%
  \or\num{0.749+-0.004}% Caso 1
  \or\num{0.721+-0.022}% Caso 2
  \or\num{0.645+-0.040}% Caso 3
  \or\num{0.622+-0.051}% Caso 4
  \else\num{0.567+-0.053}% Caso 5
  \fi
}
%\endgroup

%\begingroup mI
\providecommand\mIExactList[1]{
  \ifcase#1\relax%
  \or\num{3.265+-0.001}% Caso 1
  \or\num{3.238+-0.005}% Caso 2
  \or\num{3.220+-0.008}% Caso 3
  \or\num{3.213+-0.013}% Caso 4
  \else\num{3.174+-0.023}% Caso 5
  \fi
}
\providecommand\mIAprxList[1]{
  \ifcase#1\relax%
  \or\num{3.262+-0.001}% Caso 1
  \or\num{3.240+-0.005}% Caso 2
  \or\num{3.217+-0.010}% Caso 3
  \or\num{3.217+-0.012}% Caso 4
  \else\num{3.208+-0.009}% Caso 5
  \fi
}
%\endgroup

%\begingroup sI
\providecommand\sIExactList[1]{
  \ifcase#1\relax%
  \or\num{0.284+-0.001}% Caso 1
  \or\num{0.311+-0.008}% Caso 2
  \or\num{0.343+-0.010}% Caso 3
  \or\num{0.361+-0.020}% Caso 4
  \else\num{0.398+-0.029}% Caso 5
  \fi
}
\providecommand\sIAprxList[1]{
  \ifcase#1\relax%
  \or\num{0.281+-0.001}% Caso 1
  \or\num{0.315+-0.006}% Caso 2
  \or\num{0.334+-0.013}% Caso 3
  \or\num{0.366+-0.018}% Caso 4
  \else\num{0.411+-0.025}% Caso 5
  \fi
}
%\endgroup

%\begingroup mII
\providecommand\mIIExactList[1]{
  \ifcase#1\relax%
  \or\num{3.788+-0.005}% Caso 1
  \or\num{3.730+-0.029}% Caso 2
  \or\num{3.681+-0.052}% Caso 3
  \or\num{3.650+-0.072}% Caso 4
  \else\num{3.572+-0.095}% Caso 5
  \fi
}
\providecommand\mIIAprxList[1]{
  \ifcase#1\relax%
  \or\num{3.778+-0.006}% Caso 1
  \or\num{3.735+-0.029}% Caso 2
  \or\num{3.656+-0.048}% Caso 3
  \or\num{3.658+-0.084}% Caso 4
  \else\num{3.502+-0.078}% Caso 5
  \fi
}
%\endgroup

%\begingroup sII
\providecommand\sIIExactList[1]{
  \ifcase#1\relax%
  \or\num{0.455+-0.002}% Caso 1
  \or\num{0.486+-0.007}% Caso 2
  \or\num{0.515+-0.014}% Caso 3
  \or\num{0.533+-0.021}% Caso 4
  \else\num{0.570+-0.028}% Caso 5
  \fi
}
\providecommand\sIIAprxList[1]{
  \ifcase#1\relax%
  \or\num{0.457+-0.002}% Caso 1
  \or\num{0.490+-0.007}% Caso 2
  \or\num{0.521+-0.012}% Caso 3
  \or\num{0.545+-0.021}% Caso 4
  \else\num{0.586+-0.029}% Caso 5
  \fi
}
%\endgroup


  \endgroup
\end{document}
